% !TEX program = xelatex
% UTF-8/LF. Self-contained source for the public dictionary PDF.
\documentclass[11pt]{article}
\usepackage{fontspec}
\IfFontExistsTF{Libertinus Serif}{\setmainfont{Libertinus Serif}}{\IfFontExistsTF{TeX Gyre Pagella}{\setmainfont{TeX Gyre Pagella}}{\setmainfont{Times New Roman}}}
\IfFontExistsTF{Libertinus Sans}{\setsansfont{Libertinus Sans}}{\setsansfont{Arial}}
\IfFontExistsTF{DejaVu Sans Mono}{\setmonofont{DejaVu Sans Mono}}{\setmonofont{Consolas}}
\usepackage[a4paper,top=2.05cm,bottom=2.1cm,left=2.05cm,right=2.05cm]{geometry}
\usepackage{microtype}
\usepackage{amsmath,amssymb}
\usepackage{array,booktabs,longtable,tabularx}
\usepackage{enumitem}
\usepackage{xcolor}
\usepackage{fancyhdr}
\usepackage{needspace}
\usepackage{xurl}
\usepackage[hidelinks,unicode]{hyperref}
\usepackage{bookmark}
\definecolor{NoetherBlue}{HTML}{173E63}
\definecolor{NoetherGold}{HTML}{9B6B18}
\definecolor{SoftBlue}{HTML}{EAF2F8}
\definecolor{SoftGold}{HTML}{F8F2E6}
\definecolor{MidGray}{HTML}{5B6570}
\hypersetup{
  pdftitle={Matematičny medžuslovjansky: slovnik i metoda / Mathematical Interslavic: dictionary and method},
  pdfauthor={AI-stvorjena studija medžuslovjanskoj terminologije / AI-created Interslavic terminology study},
  pdfsubject={An auditable method, terminology dictionary, and decision atlas for AI-translated mathematical Interslavic},
  pdfkeywords={Interslavic, Medžuslovjansky, Emmy Noether, mathematical terminology, language planning, artificial intelligence, translation, dictionary, reproducibility}
}
\setlength{\parindent}{1.05em}
\setlength{\parskip}{0.28em}
\setlist{nosep,leftmargin=1.55em}
\setlength{\LTpre}{0.6em}
\setlength{\LTpost}{0.8em}
\renewcommand{\arraystretch}{1.18}
\emergencystretch=4em
\widowpenalty=10000
\clubpenalty=10000
\raggedbottom
\pagestyle{fancy}
\fancyhf{}
\fancyhead[L]{\small\sffamily Matematičny medžuslovjansky / Mathematical Interslavic}
\fancyhead[R]{\small\sffamily Terminologičny slovnik / Terminology dictionary}
\fancyfoot[C]{\thepage}
\newcommand{\code}[1]{\texttt{#1}}
\newcommand{\doiurl}[1]{\href{https://doi.org/#1}{#1}}
\newcommand{\repo}{\href{https://kokunoyumeto.github.io/emmy-noether-isv/}{onlajnovy čitateljnik / online reader}}
\newcommand{\evstrong}{\textcolor{NoetherBlue}{\ensuremath{\bullet}}}
\newcommand{\partialev}{\textcolor{NoetherGold}{\ensuremath{\circ}}}
\newcommand{\noev}{\textcolor{MidGray}{--}}
\newcommand{\profile}[5]{\mbox{S\,#1\;O\,#2\;V\,#3\;K\,#4\;R\,#5}}
\newcolumntype{Y}{>{\raggedright\arraybackslash}X}
\newcolumntype{P}[1]{>{\raggedright\arraybackslash}p{#1}}
\renewcommand{\abstractname}{Sažetak / Abstract}
\renewcommand{\contentsname}{Sodržanje / Contents}

\title{\sffamily\bfseries\color{NoetherBlue}
Matematičny medžuslovjansky\\[0.25em]
\Large Slovnik i metoda izbora v prěvodu děl Emmy Noether\\[0.85em]
\large Mathematical Interslavic\\
A dictionary and account of linguistic choices in the translation of Emmy Noether}
\author{AI-stvorjena studija medžuslovjanskoj terminologije / AI-created Interslavic terminology study}
\date{24 oktobra / October 2026\\
Iztočny korpus / Source corpus: \doiurl{10.5281/zenodo.21926382}}

\begin{document}
\maketitle

\begin{abstract}
\noindent\textbf{Medžuslovjansky.}
Kako nazvati matematičny pojem v jezyku, ktory hoče byti razumlivy govoriteljam različnyh slovjanskyh jezykov? Samo podobnost slov ne dostaje: prěvod mora zachovati točny smysl definicije, tvrdženja ili dokaza. Tut slovnik objasnjaje izbory napravjene pri AI-prěvodu děl Emmy Noether. Pri terminu dava iztočny pojem, prědloženu medžuslovjansku formu, sravniteljne formy, dostupny razlog izbora i granice dokaza. Uvod objasnjaje, kako upotrěbiti slovnik, gramatiku i odbornu literaturu pri takoj rabotě, a takože kako prověriti prědlog modela. Kde ranějši razlog ne jest zachovany, to jest izrično označeno; nova analiza ne jest predstavljena kako starše svědočstvo. Slovnik ne zaměnjaje matematičny tekst i ne tvrdi, že vsaka prědložena forma jest nova, pravilna ili potvrđena medžuslovjanskoju občinoju.

\medskip
\noindent\textbf{English.}
How should a mathematical concept be named in a language intended to be understood across the Slavic-speaking world? Resemblance between words is not enough: a translation must preserve the precise meaning of a definition, assertion, or proof. This dictionary explains choices made in an AI translation of Emmy Noether's works. Entries identify the source concept, proposed Interslavic form, comparative forms, available reasons for selection, and limits of the evidence. The introduction explains how dictionaries, grammars, and disciplinary literature can guide this work, and how to examine a model's proposal. Where an earlier reason has not survived, the entry says so; a new analysis is not presented as earlier testimony. The dictionary complements the mathematics without claiming that every proposal is a new coinage, a correct solution, or an Interslavic community standard.
\end{abstract}

\tableofcontents
\clearpage

\part{Metoda na medžuslovjanskom}

\section{Naučno pytanje i glavno tvrdženje}

Medžuslovjansky jest planovany zonalny jezyk, čija praktična cělj jest razumljivost medžu govoriteljami slovjanskyh jezykov, a ne samo interesna demonstracija jezyčnogo konstruovanja. Ranějša literatura go razmatraje kako orudje receptivnogo mnogojezyčja i nadnacionalnogo dostupa v slovjanskom prostoru \cite{merunka2019}. Iz toga slijedi konkretny problem: obyčny slovnik može dostati za běžnu komunikaciju, ale nauka trěbuje nazvanja, smyslove granice, gramatične konstrukcije i redakcijne konvencije, ktore nikogda ne sut ravnoměrno zasvědočene v vseh slovjanskyh tradicijah.

Tuto izslědovanje pyta:

\begin{enumerate}
  \item može li veliki historijno-matematičny korpus služiti kako kontrolovano mjesto za razvitje tehničnogo medžuslovjanskogo registra;
  \item može li jezyčny model sravniti više vrst dokazov i prědlagati rěšenja bez toga, da njegov prědlog sam stane jezyčnoju autoritetom;
  \item može li vsako terminologično, gramatično i zapisno rěšenje ostati dostatočno jasno dokumentovano, da ekspert, ktory ne zna orudja AI, može go prověriti i popraviti;
  \item može li taka metoda poslě jednorazovogo pripravjenja kanona byti prěnesena na drugy odborny korpus bez treniranja novogo modela.
\end{enumerate}

Praktična hipoteza jest, že jezyčny model može ulehčiti tvorjenje kandidatov, sravnjenje korpusov i doslědnu upotrěbu terminov, ako iztočny smysl i jezyčne izvori ostavajut osnovoju ocěny. Tuto izslědovanje dokumentuje tak metod, ale ne měri uštedu vrěmena nasuprot ljudskomu prěvodu. Korpus Emmy Noether ne dokazuje optimalnost vsakoj formy, razumljivost za vse govoritelje ili uspěh metody v drugoj disciplini.

\subsection{Pytanja, s ktorymi čitatelj može početi}

Za lingvista glavno pytanje jest: ktore formy mogut byti koristne medžuslovjanskomu jezyku, na kakom dokazě one stojet i kde dokaz ne dostaje? Za matematika: zachovuje li prěvod iste objekty, uslovja i argumenty? Za izslědovatelja, ktory hoče povtoriti metodu: kakve izvori potrěbuje, što dava modelu i kako razpoznaje nepravilny odgovor? Slovnik i metoda sut organizovane okolo tųh pytanj. Znanje programovanja ne jest potrěbno za čitanje obrazloženij.

Pri vsakom terminu trěba razděliti četyre stvari: matematičny pojem, izbranu formu, dokaz za tu formu i granicu tvrdženja. Forma iz ranogo glosara može byti kandidat, ne rekomendacija. Vysoka čestota slova ne dokazuje pravilnost njegova smysla. Ako razlog ne jest zachovany, zapis to govori; novo objasnjenje ne směje byti predstavljeno kako historijno svědočstvo.

\section{Što čitatelj dostavaje i kako časti se vežut}

Publikacija naměrno razděljuje predmet čitanja od dokaza o njegovom nastanku.

\begin{longtable}{@{}P{0.24\textwidth}P{0.34\textwidth}P{0.34\textwidth}@{}}
\toprule
Predmet & Funkcija & Kako go upotrěbiti\\
\midrule
Glavny čitatelj děl Emmy Noether & Latiničny i deterministično izvedeny kiriličny medžuslovjansky prěvod avtorskogo matematičnogo teksta & Čitati děla; prověrjati formuly, strukturu i kontinuity korpusa. Tut dokument go ne zaměnjaje i ne stavaje prěd njim.\\
\addlinespace
Aktualny terminologičny slovnik & Preferovane formy i smyslove granice, ktore upravljajut aktualnym čitateljem & Iskati aktualnu formu, definovany smysl, alternativy i razlog izbora.\\
\addlinespace
Polny arhivny atlas & Vsih 2574 izvornyh terminologičnyh zapisov, agregovanyh v 2016 različnyh parov & Rekonstruovati, kde forma byla prědložena, kako byla obrazložena i či ona jest točno zasvědočena v zamrznjenom korpusu. Historijna forma ne jest avtomatično aktualna norma.\\
\addlinespace
Atlas jezyčnyh i redakcijnyh odluk & Leksika, semantika, morfologija, sintaksa, ortografija, imena, citaty, matematika i granice transliteracije & Prověriti pravilo jednokratno, potom sledovati vsako pojavjenje prěz mašinno čitljivu mapu.\\
\addlinespace
JSON i CSV tablice & Polny zapisy, lokatory, pojavjenja i kvalitativne ili stvarno čislene vage & Reanalizovati podatky, graditi drugy slovnik ili sravniti buduću reviziju bez čitanja internyh produkcijnyh zapisov.\\
\addlinespace
LaTeX i iztočny ZIP & Polno redagovateljne izdanje s instrukcijami sestavjenja & Povtoriti dokument, citovati stabilne identifikatory i prědložiti popravy nad točno opreděljenym zapisom.\\
\bottomrule
\end{longtable}

Izraz \emph{kanon} v toj publikaciji znači aktualny, doslědny redakcijny osnov toj edicije. On ne znači oficialnu odluku medžuslovjanske občiny. Ekspertna ili obćinska revizija može změniti formu; atlas tada čini jasno, koje smysly i pojavjenja ta změna zasěgaje.

\section{Rolja jezyčnogo modela}

Generativny jezyčny model jest program, ktory tvori tekst na osnově statističnyh obrazcov naučenyh iz velikogo množstva tekstov. On může dobro sravnjivati formulacije, predlagati gramatične parafrazy i sledovati složene instrukcije. On ne jest slovnik, jezyčna akademija, matematičny iztočnik niti svijestny svědok togo, odkud vsaka njegova forma prihodi. Može uvěreno prědložiti nezasvědočenu riječ, směšati dvě nacionalne normy ili proměniti matematičny odnos, osoblivo v malo zastupljenom jezyku. Izslědovanja nizko-resursnoj translacije potvrđajut, že dostavnost slovnika, gramatiky i podobnyh priměrov sama po sobě ne odstranjaje problemy izvlačenja i razuměvanja \cite{court2024}; terminologične ograničenja pomagajut, ale mogu škoditi tečnosti, ako se priměnjujut mehanično \cite{bane2023}.

Zato model v toj metodi ima tri dozvoljene role:

\begin{enumerate}
  \item \textbf{prědlagatelj}: tvori kandidatny prěvod ili termin;
  \item \textbf{sravnitelj}: stavja kandidaty ob iztočnik, gramatiku, slovnik, matematičnu literaturu i slovjanske tradicije;
  \item \textbf{upravitelj doslědnosti}: nahodit sva pojavjenja jednogo pravila i prověrjaje, že poprava ne zasěgaje imena, formule, citaty ili drugi smysl.
\end{enumerate}

Model ne dostavaje čtvrtu rolu jezyčnogo zakonodatelja. Rěšenje mora pokazati dokaz, granicu i odklonjene alternativy. Tuto približava tok konceptno orientovanoj terminologičnoj rabotě: predmet, koncept, definicija i nazvanje razděljujut se prěd izborom preferovanogo termina \cite{iso704}. Jednako važno, jedan napisany oblik s različnymi smyslami dostavaje različne zapisy, a ne jedno globalno zaměnjenje.

\subsection{Upotrěba modela bez novogo treniranja}

Metoda upotrěbjaje gotovy jezyčny model i ne trěbuje treniranje novoj sistemy. Izslědovatelj dava modelu tekst, relevantne pasaže gramatiky i slovnika, odbornu literaturu i obrazložene terminologične odluke. Tuto jest važna praktična razlika: potrěbno jest pripraviti dobre izvory i jasnu zadaču, ne konstruovati novy model. Izraz \emph{zero-shot} zato ne jest točny opis: prěvod ne nastaje bez priměrov, izvorov i ranějših rěšenij. Povtorljivost znači, že drugy izslědovatelj može povtoriti sravnjenje i prověrku; ne znači, že vsaky model mora dati doslovno isti prědlog.

\section{Kako byl dokaz koncepta stvarno sformovany}

Dostupny arhiv dozvoljaje razpoznati nekoliko vrst raboty. Tekst byl razděljeny na opreděljene iztočne jedinice s 215 glosarami. Oni sodržet 2574 zapisy: německy termin ili pojem, rusku i ukrajinsku sravniteljnu formu, medžuslovjansky kandidat, razlog i, kde bylo potrěbno, označenje nejasnosti. To jest izvorna terminologična baza. Prisutnost sravniteljnogo slova v glosaru ne dokazuje sama po sobě, že konkretny odborny izvor byl konsultovany.

Druga faza izprobovala različne sposoby prěvoda: prěz anglijsky posrědny tekst, s odvojenjem prozy od TeX-a i matematiky i neposrědnje iz německogo. Zachovanje vseh matematičnyh oznak jest potrěbna, ale ne dostatočna prověrka: pravilno čislo oznak ne dokazuje pravilnost prozy ili logičnyh odnosov. Trideset jedna jedinica knižnogo teksta byla prěvedena v rusky, ukrajinsky i medžuslovjansky s pomočju modela. Dostupny dokaz ne ustanavjaje pouzdano točnu identičnost modela i nivo usilja pri tom prěvodu; ljudska redakcija ne jest potvrđena.

Jezyčna analiza razděljuje leksiku, semantiku, gramatiku, ortografiju i věrnost iztočniku. Osobno pytanje jest granica medžu prozoju, ktora može byti transliterovana, i imenami, citatami, rimskimi čislami i matematičnymi oznakami, ktore trěbujut zachovati svoju identičnost. Tuto ne sut samo različne načiny zapisa jednogo slova: rolja znaka v argumentu odlučaje, či pravilo prěpisivanja jest priměnimo. Pojavjenjske tablice omogožnjajut prověrjati tu razliku v konkretnyh pasažah.

Kde zachovany zapis izrično dava razlog, možno go objasniti i prověriti. Kde razlog ne jest zachovany, možno izslědovati formu znova, ale novo opravdanje ne jest svědočstvo o tom, kako izvorna odluka byla napravjena. Ta razlika jest nužna za naučnu rekonstrukciju: tuta kniga opisujuče dokumentovane izbory, ne skryte misli modela.

\section{Povtorljiva metoda za novy odborny korpus}

Slědujuči tok jest napisan za lingvista ili terminologa, ktory ne mora znati programovanje ili arhitekturu AI.

\begin{longtable}{@{}P{0.08\textwidth}P{0.24\textwidth}P{0.60\textwidth}@{}}
\toprule
Krok & Produkt & Što se děla i što može propasti\\
\midrule
0 & Izslědovatelsky mandat & Opreděliti publiku, registr, pismo, oblast, prava, iztočnu ediciju i granice tvrdženja. Bez toga model sam izmysli nepostojanu cělj.\\
\addlinespace
1 & Iztočna avtoriteta & Zamrznuti polny iztočnik, numeraciju, formuly, odkazy i bibliografiju. Razděliti diplomatsky prěvod od oddělnyh opravljenyh izdanj.\\
\addlinespace
2 & Jezyčny i disciplinarny kanon & Pročitati oficialnu ciljnu gramatiku i slovnik, potom relevantnu matematiku ili drugu oblast v reprezentativnyh slovjanskyh jezykah. Zapisati točne pasaže, ne samo spisok knig.\\
\addlinespace
3 & Jedinice i zaščitene ostrovy & Razděliti korpus na stabilne sekcije. Zaščititi formuly, komandy, identifikatory, citaty, polna imena i druge objekty, ktore model ne směje parafrazovati.\\
\addlinespace
4 & Kandidatny prěvod & Dati modelu jednu omeđenu jedinicu, relevantny kanon i jasny izhodny format. Po mogućnosti prěvoditi prozu odvojeno i programski vratiti matematiku.\\
\addlinespace
5 & Konceptny zapis & Za vsaky netrivialny izbor zabilježiti iztočny koncept, ciljnu formu, gramatičnu funkciju, ruske i ukrajinske ili druge sravniteljne formy, oficialny dokaz, disciplinarnu upotrěbu i kontekst. Polisemiju razděliti po smyslu.\\
\addlinespace
6 & Izbor kandidatov & Najprvo upotrěbiti iztočnu i predmetnu věrnost kako tvrd prag. Potom sravniti oficialny ISV dokaz, vzhodnu--zapadnu--južnu razpoznavnost, tehničnu upotrěbu, doslědnost i nejasnost. Broj upotrěbiti samo ako historijno pravilo stvarno jest čisleno.\\
\addlinespace
7 & Odluka i odklonjene alternativy & Zapisati ne samo pobjědnika, ale realne alternativy, razlog njihova odklonjenja, ostavajuču neuvěrjenost i granicu pravila. Jedno dokazano pravilo može upravljati mnogimi pojavjenjami.\\
\addlinespace
8 & Pojavjenjska mapa & Najti vsako relevantno pojavjenje i prověriti njegov lokalny smysl. Ne priměnjivati globalnu zaměnu na homografy, citaty, imena ili matematične simbole.\\
\addlinespace
9 & Jezyčna i strukturna prověrka & Sravniti cěly cilj s iztočnikom: tvrdženja, kvantifikaciju, znaky, odkazy, padež, soglasovanje, aspekt i terminologiju. Kompilacija dokazuje samo tehničnu sestavljivost.\\
\addlinespace
10 & Pisma i izdanje & Zamrznuti jednu redagovateljnu jezyčnu površinu. Drugo pismo izvoditi deterministično s jasnymi zaštitami. Sestaviti, vizualno pregledati i publikovati čitatelj, neposrědnji LaTeX, iztočny paket, slovnik i podatky.\\
\addlinespace
11 & Ekspertna revizija & Dati ekspertu kratky aktualny zapis plus polny dokazny prilog. Poprava imaje stabilny identifikator, smyslovu oblast i spisok pojavjenj, tako da ne trěba povtoriti cělu produkciju.\\
\bottomrule
\end{longtable}

Poluavtomatično tvorjenje terminološkyh baz obično razděljuje obrabotku korpusa, izvlačenje terminov, napolnjenje bazy i upravljanje zapisami \cite{nazar2023}. Tut tok dodava dvě stvari nužne za planovany medžujezyk: ravnověsno sravnjenje jezyčnyh vetvij i dokaznu granicu medžu prědlogom modela i prijetym kanonom.

\section{Kako ekspert može čitati i osporiti rěšenje}

\subsection{Priměnjenje metody na jednu pasažu}

Izslědovatelj može početi s jednim polnym matematičnym argumentom, na priměr s opreděljenjem i jego prvim priměrom. Iztočna pasaža mora iměti stalny lokator: avtor, dělo, odděl i strana ili numerovany paragraf. S njom trěba čitati toliko okružajučego teksta, koliko jest potrěbno za značenje oznak i predpostavok. Izbrana pasaža jest rabotna jedinica, ne zaměna za polny naručeny korpus.

Priprava sodržuje tri kratke česti: sam iztočnik, relevantne jezyčne i predmetne izvory i spisok raněje obrazloženyh terminov. Bibliografija bez konkretnyh pasaž ne dostaje. Pri novom matematičnom pojmu trěba razděliti oficialno zasvědočenje obyčnogo slova od zasvědočenja tehničnogo smysla. Pri polisemii trěba najprvo napisati, ktory odnos věta vyražaje. Model može prědložiti formu i obrazloženje, ale citat iz knjigy mora byti najdeny v samoj knjigě, ne prijety na osnově modelovogo tvrdženja.

Rezultat za jednu pasažu ima tri čitateljske česti: prěvod, spisok terminologičnyh ili gramatičnyh rěšenij i označene nejasnosti. Za vsako rěšenje dovoljen jest kratak zapis, ako on sodržuje pojem, vybranu formu, relevantny izvor, realne alternativy i razlog. Ako ne jest nalezena opravdana alternativa, to trěba napisati; ne trěba izmysljati drugogo kandidata radi simetričnoj tablicy. Ako izvor ne podpiraje točny smysl, ta granica ostavaje vidima.

Prověrka prěvoda ne jest isto što ponovno pytanje, či on jest dobar. Trěba sravniti konkretne tvrdženja: kto ili što jest podmet, na ktore objekty odnosi se kvantifikacija, što jest predpostavka i što zaključok, kde počinaje uslov i kde negacija. Pri dvoh kandidatnyh formach trěba prověriti takože rod, padež i upravljanje v cěloj věte. Slučaj \emph{prsten} pokazuje, čemu proměna imennika bez soglasovanja jego okoline ne dostaje.

Poslě izbora pravila izslědovatelj nahodi druge pojavjenja togo pojma. Iskanje po napisu može pomoći, ale isti napis može iměti drugy smysl. Zato pojavjenjska tabela jest konkordancija: spisok pasaž, v ktoryh trěba prověriti priměnimost pravila. Ona ne jest dokaz pravilnosti sama po sobě. CSV-tablica može byti otvorena v običnom tabličnom programu; JSON dava iste svęzi v strukturě za avtomatičnu obrabotku. Čitatelj ne trěbuje čitati JSON, da razuměl razlog izbora.

\subsection{Što rekonstrukcija može dokazati}

\begin{description}
  \item[Izvorno dokumentovany razlog.] Zachovany zapis govori, zašto kandidat byl izbrany. Tuto ustanavjaje dokumentovany razlog, ne avtomatično njegovu pravilnost.
  \item[Pozdnějša prověrka.] Forma jest sravnjena s izvorom, ktory ne jest dokumentovany pri prvom izborě. Tuto može podpirati ili osporiti formu, ale ne dokazuje ranějšu konsultaciju.
  \item[Praktično pravilo za povtorjenje.] Metoda govori, kako izbor trěba dokumentovati i prověrjati v novom korpusu. Takvo pravilo ne jest tvrdženje, že vsaka historijna odluka byla napravjena idealno.
\end{description}

Tute razlike omogućnjajut ekspertu prověrjati i dobru ideju, i slabost dokaza. Otsustvje čislene vagi ne jest otsustvje razloga: mnoge izbory zaviset od razlike v smyslu ili od gramatičnogo pravila, ktore ne potrěbujut glasovanje. Ale čisla, ako sut upotrěbjene, morajut iměti opreděljene jedinice, polny spisok prinosov i dostupnu aritmetiku.

\subsection{Pytanja za odbornu ocěnu}

Ekspert ne trěba prověrjati ``AI'' kako cělinu. Za vsaky zapis može pytati šest obyčnyh naučnyh pytanj:

\begin{enumerate}
  \item Jest li iztočny koncept opreděljen pravilno i jest li německa pasaža dobro pročitana?
  \item Jest li ciljna forma stvarno medžuslovjanska po deklarovanoj gramatikě i ortografiji?
  \item Sut li sravniteljne jezyčne tradicije reprezentativne, ili jedna vetv jest prěsilno zastupljena?
  \item Jest li disciplinarny smysl zasvědočeny, osoblivo gdje obyčny slovnik dava samo běžny smysl?
  \item Jesu li realne alternativy i opasnost mnogoznačnosti prikazane čestno?
  \item Sut li vsaka pojavjenja istogo pravila najdena, a različne smysly ostavljene odvojeno?
\end{enumerate}

Taka revizija slijedi izslědovanja, ktore pokazuju, že profesionalni ocěnitelji s polnym dokumentnym kontekstom i točno kategorisanymi pogrěškami dajut drugačiju i koristnějšu ocěnu od prostyh agregatnyh bodov \cite{freitag2021}. Avtomatične metrike takože ne razpoznavajut pouzdano vse male, ale smyslovo važne razlike medžu silnymi prěvodi \cite{agrawal2024}. Zato tuto izdanje ne zaměnjaje eksperta jednim rezultatom; ono mu dava predmet, dokaz i lokalizovanu popravljivu tvrdnju.

\section{Čemu služi tut slovnik}

Korpus jest prěvod matematičnyh tekstov, ne zbirka izolovanyh slov. Zato jedinica izbora često ne jest jedno slovo ili lemma, ale odnos medžu terminom, větoju i matematičnym predmetom. Německo \emph{gilt}, na priměr, može govoriti, že tvrdženje jest pravdivo, že uslov jest izpolnjeny, že teorem može byti priměnjeny ili že pravilo ostavaje v silě. Jedna navodno elegantna zaměna za vse te funkcije by stvorila sistematičnu pogrěšku. Slovnik zato daje samostojne terminy i važne smyslove razděljenja.

Latinična medžuslovjanska površina jest redakcijny jezyčny kanon. Kirilična površina jest deterministično izvedeno pismo togo samogo jezyčnogo teksta, a ne drugi prěvodny svědok. Slovnik citira latinične zaglavne formy, ibo one upravljajut leksičnymi i morfologičnymi rěšenjami; kiriličny čitateljnik ostavaje polno dostupny v glavnom izdanje.

\section{Kako byl izbor napravjen}

Rekonstruovana metoda razděljuje pěť vrst dokazov. Ne vsaky ranějši zapis sodržuje vse pěť; njihovo otsustvje ne směje byti zaměnjeno domyslom o ranějšej konsultaciji:

\begin{enumerate}
  \item \textbf{Iztočny i matematičny smysl.} Německa věta, formula, argumentna struktura, kvantifikacija i terminološka funkcija sut tvrd prag. Kandidat, ktory měnja predmet ili logičny odnos, odpada bez daljnogo bodovanja.
  \item \textbf{Oficialny medžuslovjansky dokaz.} Slovnik i gramatika dajut lemmy, paradigmy, aspekt, upravljanje i standardne vezy. Oni ne mogut sami dokazati novy strogo matematičny smysl, ale sut glavno jezyčno sidro.
  \item \textbf{Sravnjenje slovjanskyh tradicij.} Vzhodne, zapadne i južne jezyčne grupy sut gledane ravnověsno. Tuto sravnjenje služi za razpoznavanje i za odkryvanje lažnyh prijateljev; ono ne jest prosty plebiscit.
  \item \textbf{Matematičny i unutrašnji korpus.} Upotrěba v ciljnoj disciplini i doslědnost medžu statjami sut važne. Termin ne dostavaje prednost samo zato, že byl slučajno čest v naslědženom tekstu.
  \item \textbf{Čitateljska jasnost.} Kratkost, razbirljivost bez anglijskogo teksta i mala opasnost od mnogoznačnosti služe kako izbor medžu kandidatami, ktore sut već smyslovo dopustime.
\end{enumerate}

Tut porędok jest važen: četvrta ili pěta vrstva nikogda ne popravjaje kandidata, ktory ne projde prvu. Rěšenje takože ne stavaje kanonom zato, že AI go prědlaže. Ono mora byti lokalizovano, sravnjeno s dokazom, ograničeno po smyslu i prověrjeno v vsakom relevantnom pojavjenju.

\section{Vaga i profil dokaza}

Slědujuča tabela jest čitateljska sinteza metode, ne avtomatičny algoritm, ktory byl naknadno puščeny nad slovami.

\begin{table}[ht]
\centering
\caption{Funkcija dokaznyh slojev.}
\begin{tabularx}{\textwidth}{@{}P{0.25\textwidth}P{0.16\textwidth}Y@{}}
\toprule
Sloj & Funkcija & Pytanje\\
\midrule
Iztočna/matematična věrnost & obvezny prag & Ohranjaje li kandidat točny predmet, odnos, opseg i logiku?\\
Oficialny ISV slovnik i gramatika & vysoka dokazna sila & Jest li lemma ili konstrukcija zasvědočena, i jest li morfologija pravilna?\\
Vzhod--Zapad--Jug & ravnověsna diagnostika & Kde forma pomaga ili škodi razpoznavanju medžu grupami?\\
Disciplinarny i unutrašnji korpus & silny veznik & Funkcionuje li termin u matematičnom registru i ostavaje li doslědny?\\
Razbirljivost i mnogoznačnost & poslědnji razlikovnik & Ktory dopustimy kandidat najjasněje vodi čitatelja?\\
\bottomrule
\end{tabularx}
\end{table}

V slovničnyh tabelah profil \profile{$\checkmark$}{\evstrong}{\evstrong}{\evstrong}{\evstrong} znači: \textbf{S} -- iztočny smysl projde tvrd prag; \textbf{O} -- oficialny medžuslovjansky dokaz; \textbf{V} -- medžugrupno slovjansko sravnjenje; \textbf{K} -- matematičny ili unutrašnji korpus; \textbf{R} -- čitateljska razbirljivost. \evstrong\ označaje direktno i silno podporjenje, \partialev\ ograničeno, posrědno ili smyslovo užše podporjenje, a \noev\ govori, že izbrano rěšenje ne počivaje na toj vrstvě. Simboly ne sut kalibrovane věrojatnosti.

\section{Četyri priměry, gdje samo lemma ne dostaje}

\subsection{Algebraičny \emph{prsten}: jediny stvarno čisleny izbor}

Za německy \emph{Ring} byli matematične tradicije razdělene. Dany dokazy sut byli grupovane v šest ravnyh glasovnyh grup, tako da jezyk s mnogimi dostupnymi knjigami ne može sam dominovati. Rezultat byl:

\begin{table}[ht]
\centering
\caption{Historijno bodovanje kandidatov za algebraičny prsten.}
\begin{tabular}{@{}lcl@{}}
\toprule
Kandidatny korenj & Jedinice & Glavne zasvědočene tradicije\\
\midrule
\textbf{prsten / pierścień} & \textbf{2.5} & poljska; hrvatska, srbska, makedonska; blizky bulgarsky oblik\\
kolco & 2.0 & ruska, ukrajinska i blizke vzhodne tradicije\\
okruh & 1.0 & češska i slovačska\\
kolobar & 0.5 & slovenska\\
\bottomrule
\end{tabular}
\end{table}

Zato izdanje upotrěbjaje \textbf{prsten}, \textbf{podprsten}, \textbf{nadprsten} i \textbf{medžuprsten}. Oficialny slovnik potvrđaje široku razpoznavnost fizičnogo prstena, ale ne sam algebraičny smysl; matematične knjigy dajut tehničny dokaz. Pri prvom pojavjenju čitateljski most imenuje rusko \emph{кольцо} i ukrajinsko \emph{кільце}. Tuto jest dobro ograničeny slučaj, kde čisla stvarno pripadajut rěšenju.

\noindent Polny razklad tųh šest jedinic jest slědujuči; grupovanje slijedi kriterije jezyčnogo projekta \cite{designcriteria}. Tabela rekonstruuje zapisane summy iz grup i zasvědočenyh korenjev, ne iz novogo opyta s čitateljami.

\begin{center}\small
\begin{tabularx}{\textwidth}{@{}Yrrrr@{}}
\toprule
Jezyčna grupa & prsten & kolco & okruh & kolobar\\
\midrule
Ruska & 0 & 1 & 0 & 0\\
Ukrajinska i běloruska & 0 & 1 & 0 & 0\\
Poljska & 1 & 0 & 0 & 0\\
Češska i slovačska & 0 & 0 & 1 & 0\\
Slovenska i srbska/hrvatska/bosanska & 0.5 & 0 & 0 & 0.5\\
Bulgarska i makedonska & 1 & 0 & 0 & 0\\
\midrule
Suma & 2.5 & 2 & 1 & 0.5\\
\bottomrule
\end{tabularx}
\end{center}

Prsten ima največši děl, ale ne absolutnu bolšinu šest jedinic. Broj knjig ili pojavjenj ne dava dodatny glas; grupovanje s poljskim pierścień jest jezyčny sud o srodnosti. Spisok iztočnyh korpusnyh brojov ne dava osobne broje za bělorusky i bosansky. Tabela zato ne směje byti čitana kako ravnoměrny empiricky opyt vseh nazvanyh jezykov. Polny čisleny zapis jest v dodatku podatkov \texttt{RING\_COMPARISON\_WEIGHTS.json}.

\subsection{Německo \emph{gilt}: smyslovo razděljenje vměsto pobjědnika}

Češsky i slovačsky često imajut \emph{platí}; hrvatsky \emph{vrijedi}; poljsky \emph{zachodzi} ili \emph{ma miejsce}; slovensky \emph{velja}; bulgarske matematične věty často izbirajut konstrukciju izpolnjenja. Jedno nacionalno slovo zato ne daje ravnověsny medžuslovjansky predikat. Izdanje razděljaje funkcije: \textbf{imaje město} za prosty matematičny predikat, \textbf{jest pravdivo} za pravdivost, \textbf{izpolnjaje se} za uslov, \textbf{može se priměniti} za priměnjivost i \textbf{ostavaje v silě} za trajuču silu pravila. Oficialne \emph{važiti} i \emph{platiti} ne sut upotrěbjene v izmišljeno proširjenom smyslu.

\Needspace{0.38\textheight}
\subsection{Vezy dokaza: jedna površina za jednu funkciju}

\begin{longtable}{@{}P{0.20\textwidth}P{0.25\textwidth}P{0.18\textwidth}P{0.27\textwidth}@{}}
\caption{Glavne logične i diskursne vezy.}
\\
\toprule
Funkcija & Izbrana površina & Tipičny německy odnos & Važna granica\\
\midrule
uzrok & \textbf{zato že} & weil, da & ne \emph{zato} bez že\\
zaključak / rezultat & \textbf{zato} & also, daher, somit & ne vremensko \emph{togda}\\
ekvivalencija / objasnjenje & \textbf{to jest} & das heißt & ne zaključny \emph{zato}\\
točna imenovka & \textbf{imenno} & nämlich & užše od opčego \emph{to jest}\\
uslovna apodoza & \textbf{ako \ldots{} to} & wenn \ldots{} so & ne vremensko \emph{togda}\\
dodavanje & \textbf{takože} & auch & jedna riječ\\
poslědica ili cělj & \textbf{tako že} & so, so dass & dvě riječi\\
\bottomrule
\end{longtable}

\subsection{Lažni prijatelji: podobnost ne jest dokaz}

Nekoliko velikyh familij byli razdělene po smyslu: \textbf{sodrživati} znači matematično ``sodržati/contain'', \textbf{sodržanje} znači ``content'', a \textbf{zadržati} ostavaje ``hold back''. \textbf{Obstajati} označaje egzistenciju ili trajuče obstajanje, \textbf{sučstveny} esencialnost, a \textbf{v sučnosti} ``in essence''. \textbf{Prědpokladati} jest nesvršeno ``assume'', dok \textbf{prědpoložiti/prědpoloženje} nosi svršenu, pasivnu ili imensku funkciju. Takva razděljenja sut važnějša od površnoj podobnosti s jednym iz slovjanskyh jezykov.

\clearpage
\part{Terminologičny slovnik / Terminology dictionary}

\section{Kako čitati zapisy / How to read the entries}

Zapisy razděljujut historijne prědlogy, formy zasvědočene v prěvodu i obrazložene rekomendacije. To ne sut sinonimy: prědlog može byti zavrgnjeny, a upotrěbljena forma može potrěbovati popravu. ``Alternativy'' sut stvarno razgledane ili zasvědočene možnosti, ne izmyslene protivniki izbranogo slova. Polno medžuslovjansko objasnjenje imaje anglijsku paralelu. Alfabetičny indeks vodi k konkretnym zapisam; identifikatory dozvoljajut svęzati zapis s polnymi podatkami.
\par\emph{EN:} Entries distinguish historical proposals, forms attested in the translation, and justified recommendations. These are not synonyms: a proposal may be rejected, and an attested form may need correction. Alternatives are documented or attested possibilities, not invented competitors. Complete Interslavic explanations have English parallels. The alphabetical index leads to individual entries; stable identifiers connect them to the full supporting data.

Prva tabela dava orientaciju v glavnyh smyslovyh razděljenjah; ona ne jest polny slovnik. Polny spisok potom sodržuje 2016 različnyh iztočno-medžuslovjanskyh parov iz 2574 zachovanyh zapisov, 168 podrobno objasnjenyh jezyčnyh i redakcijnyh izborov i dvanadset gramatičnyh ili pravopisnyh konvencij. Pisemne role i pojavjenja imajut osobne tablice. Tute čisla měręt obseg dokumentacije, ne čislo novyh slov ili dokazanyh jezyčnyh uspěhov.
\par\emph{EN:} The opening table is a guide to major distinctions, not the whole dictionary. The complete entries follow: 2,016 distinct source--Interslavic pairs from 2,574 surviving records, 168 detailed linguistic and editorial choices, and twelve grammatical or orthographic conventions. Separate tables identify script functions and occurrences. These figures measure documentary coverage, not new coinages or demonstrated linguistic successes.

\small
\begin{longtable}{@{}P{0.18\textwidth}P{0.23\textwidth}P{0.51\textwidth}@{}}
\caption{Glavne leksične i terminologične izbory / Principal lexical and terminological choices.}\label{tab:dictionary}\\
\toprule
\textbf{Forma} & \textbf{Smysl i iztočna funkcija} & \textbf{Dokaz, izbor i alternativy / Evidence, selection, and alternatives}\\
\midrule
\endfirsthead
\multicolumn{3}{l}{\small\itshape Prodolženje slovnika / Dictionary continued}\\
\toprule
\textbf{Forma} & \textbf{Smysl i iztočna funkcija} & \textbf{Dokaz, izbor i alternativy / Evidence, selection, and alternatives}\\
\midrule
\endhead
\midrule
\multicolumn{3}{r}{\small\itshape Prodolženje na slědujučej stranici / Continued on next page}\\
\endfoot
\bottomrule
\endlastfoot

\textbf{analogično}\\\emph{prislov / adverb}
& analogično, na podobny način; discourse ``similarly''
& \profile{$\checkmark$}{\evstrong}{\partialev}{\evstrong}{\evstrong}. Izbrano za větnu funkciju, gdje iztočnik pravi analogiju; ne za odnos odgovaranja medžu objektami. Zato \emph{odgovarjajuče} i \emph{korespondencija} ostavajut odvojene.\\[-0.1em]
&& \emph{EN:} Used when the source advances an analogous argument, not for correspondence between objects; the latter remains \emph{odgovarjajuče/korespondencija}.\\
\addlinespace

\textbf{byti}\\\emph{glagol / verb}
& kopula i infinitiv ``to be''; German \emph{sein}
& \profile{$\checkmark$}{\evstrong}{\evstrong}{\evstrong}{\evstrong}. Oficialny infinitiv daje jednu površinu za kopularne, pasivne, modalne i predikativne konstrukcije. Nacionalno oblikovano \emph{biti} ne nosi drugy smysl i zato ne ostavaje kako slobodny dublet.\\[-0.1em]
&& \emph{EN:} The official infinitive is used in copular, passive, modal, and predicative constructions. The nationally shaped doublet \emph{biti} had no distinct function.\\
\addlinespace

\textbf{cělkom}\\\emph{prislov / adverb}
& cělo, sovsim; German \emph{ganz, vollständig} in ordinary adverbial use
& \profile{$\checkmark$}{\evstrong}{\partialev}{\evstrong}{\evstrong}. Semantično razděljenje odklanjaje slepu zaměnu: obyčno ``completely'' stavaje \emph{cělkom}, tehnična complete reducibility upotrěbjaje \emph{polna reducibilnost}, a perfect field \emph{doskonalo polje}.\\[-0.1em]
&& \emph{EN:} Ordinary ``completely'' is \emph{cělkom}; complete reducibility and perfect fields have separate technical terms.\\
\addlinespace

\textbf{cěločislovy}\\\emph{pridavnik / adjective}
& cěločislny; German \emph{ganzzahlig}, integral
& \profile{$\checkmark$}{\evstrong}{\evstrong}{\evstrong}{\evstrong}. Koren \emph{cěločisl-} jest stabilny v korpusu; finalna pridavnička familija jest \emph{cěločislov-}. Starše \emph{cěločislovn-} směšalo dvě derivacijske površiny.\\[-0.1em]
&& \emph{EN:} The corpus keeps the transparent \emph{cělo + čislo} root and uses the final adjectival family \emph{cěločislov-}; the older \emph{-ovn-} layer was not retained.\\
\addlinespace

\textbf{cělj} / \textbf{cěly}
& \emph{cělj}: cělj, naměr; \emph{cěly}: cěly, polny
& \profile{$\checkmark$}{\evstrong}{\evstrong}{\evstrong}{\evstrong}. Formy sut podobne, ale pripadajut različnym lemmam. \emph{Za toj cělj} označaje purpose; \emph{cěly red} označaje whole series. Konstrukcija \emph{celěsoobrazno} jest razložena smyslovo kako \emph{koristno jest}.\\[-0.1em]
&& \emph{EN:} \emph{Cělj} is ``aim/purpose'' and \emph{cěly} is ``whole/entire.'' Similar spelling is not allowed to collapse the distinction.\\
\addlinespace

\textbf{cělost}\\\emph{ženski rod / feminine noun}
& cělost, integriteta; wholeness or integrity
& \profile{$\checkmark$}{\evstrong}{\partialev}{\evstrong}{\evstrong}. Oficialny korenj jest \emph{cělosť}; osnovna ortografija izdanja piše produktivne imenice na \emph{-ost} bez finalnogo znaka měkkosti. Neutrudženo \emph{celost} ne pokazuje etimologičny ě.\\[-0.1em]
&& \emph{EN:} The official root is \emph{cělosť}; the edition's basic orthography writes productive \emph{-ost} nouns without final softness marking.\\
\addlinespace

\textbf{cěnnost}\\\emph{ženski rod / feminine noun}
& vrědnost, value; German \emph{Wert} in mathematical value contexts
& \profile{$\checkmark$}{\evstrong}{\evstrong}{\evstrong}{\evstrong}. Oficialny korenj \emph{cěnnosť} razlučaje imensku matematičnu ``value'' od pridavnika \emph{vrědnostny}, ktory může imati drugi smysl. Izdanjska površina piše \emph{cěnnost}.\\[-0.1em]
&& \emph{EN:} Mathematical noun uses of German \emph{Wert} take \emph{cěnnost}. The distinct adjective \emph{vrědnostny} is not swept into the same replacement.\\
\addlinespace

\textbf{curišsko}\\\emph{pridavničko prislovje / adjectival adverb}
& curiško, ``in the Zurich lecture''; German \emph{Züricher}
& \profile{$\checkmark$}{\noev}{\partialev}{\evstrong}{\partialev}. Tuto jest jasno označeno domovo rěšenje, ne oficialny slovničny zaglavnik. Ono adaptuje početno /ts/ i \emph{h+sk} po zasvědočenom medžuslovjanskom obrazcu, a odklanjaje německo \emph{ü}. Odklonjeno \emph{švejcarsko} měnja grad v državu. Uvěrjenost ostavaje ograničena.\\[-0.1em]
&& \emph{EN:} A declared house adaptation for the Zurich city adjective, not an official headword. It adapts initial /ts/ and the \emph{h+sk} sequence; \emph{švejcarsko} was rejected because it changes city to country.\\
\addlinespace

\textbf{čest}\\\emph{ženski rod / feminine noun}
& čest, dio; German \emph{Teil}
& \profile{$\checkmark$}{\evstrong}{\evstrong}{\evstrong}{\evstrong}. Aktualna osnovna ortografija piše \emph{čest, česti, čestično}. Nosna opcionalna površina \emph{čęst} ne jest drugi termin i zato ne služi kako dublet. Homografi kako \emph{često} i drugi korenji ostavajut odvojeni.\\[-0.1em]
&& \emph{EN:} The reader's basic spelling uses \emph{čest/česti/čestično}. Etymological \emph{čęst} is not a second term, and unrelated look-alikes are excluded.\\
\addlinespace

\textbf{čto se tyče}\\\emph{ustaljena veza / fixed phrase}
& čto se tyče, glede; ``as regards, concerning''
& \profile{$\checkmark$}{\evstrong}{\evstrong}{\evstrong}{\evstrong}. Oficialna leksikalizovana prepozicijska veza jest ohranjena kako cělina. Samostojno južnoslovjansko \emph{što} ne jest globalno zaměnjeno; ne-vprašalne konstrukcije sut rěšene věta po větě.\\[-0.1em]
&& \emph{EN:} The official lexicalized phrase is kept as a unit. Standalone \emph{što} was not globally replaced; non-interrogative constructions require sentence-level analysis.\\
\addlinespace

\textbf{dělitelj}\\\emph{mužski rod / masculine noun}
& dělitelj, divisor; German \emph{Teiler}
& \profile{$\checkmark$}{\evstrong}{\evstrong}{\evstrong}{\evstrong}. Transparentny derivat od \emph{děliti}; forma jest zasvědočena i masovno podprta v matematičnom korpusu. Neakcentovano \emph{delitelj} jest ortografičny ostatok, ne konkurentny termin.\\[-0.1em]
&& \emph{EN:} The transparent derivative of \emph{děliti} is used for divisor. Unaccented \emph{delitelj} was an orthographic residue, not a competing concept.\\
\addlinespace

\textbf{dobyti} / \textbf{dobyvati}
& dobiti, izvesti, obtain/derive; German \emph{erhalten, gewinnen}
& \profile{$\checkmark$}{\evstrong}{\evstrong}{\evstrong}{\evstrong}. Aspekt i větna struktura upravljajut izborom: \emph{dobyti} jest svršeno, \emph{dobyvati} nesvršeno. \emph{Dostaviti} ostavaje samo gdje iztočnik stvarno znači supply, submit ili deliver; ne služi za matematično dobyvanje rezultata.\\[-0.1em]
&& \emph{EN:} Perfective \emph{dobyti} and imperfective \emph{dobyvati} cover mathematical obtaining or derivation. \emph{Dostaviti} is reserved for genuine supply or submission.\\
\addlinespace

\textbf{dopuščati}
& dopustiti, dozvoliti; German \emph{zulassen}
& \profile{$\checkmark$}{\evstrong}{\evstrong}{\evstrong}{\evstrong}. Oficialny etimologičny zapis \emph{dopušćati} jest prikazany v osnovnoj čitateljskoj ortografiji kako \emph{dopuščati}. Tuto jest pisemna projekcija istej lemmy, ne smyslova zaměna.\\[-0.1em]
&& \emph{EN:} The official etymological form \emph{dopušćati} appears as \emph{dopuščati} under the edition's declared basic spelling; the lemma and sense do not change.\\
\addlinespace

\textbf{doskonalo polje} / \textbf{nedoskonalo polje}
& perfect field / imperfect field; German \emph{vollkommener Körper / unvollkommener Körper}
& \profile{$\checkmark$}{\evstrong}{\partialev}{\evstrong}{\evstrong}. Tehničny termin upotrěbjaje oficialnu tvrdopridavničku familiju \emph{doskonaly/nedoskonaly}. Obyčna polnost jest odvojena; \emph{polno polje} by směšalo algebraičnu perfektnost s opčim ``complete/full''.\\[-0.1em]
&& \emph{EN:} Field-theoretic perfection is \emph{doskonalo polje}; ordinary completeness uses other terms. \emph{Polno polje} was rejected as semantically misleading.\\
\addlinespace

\textbf{dva}: \textbf{dvoh, dvom, dvoma, dvě}
& kosy padeži čislovnika dva / oblique forms of two
& \profile{$\checkmark$}{\evstrong}{\evstrong}{\evstrong}{\evstrong}. Oficialna paradigma dava genitiv/lokativ \emph{dvoh}, dativ \emph{dvom}, instrumental \emph{dvoma} i ženski nominativ/akuzativ \emph{dvě}. Varijanty \emph{dvuh, dveh, dvaju, dvěma, dvije} ne ostavajut kako nesemantične nacionalne sloje.\\[-0.1em]
&& \emph{EN:} The official paradigm supplies \emph{dvoh, dvom, dvoma}, and feminine \emph{dvě}. Nationally shaped variants are not retained as functionless layers.\\
\addlinespace

\textbf{egzistovati}\\\emph{glagol / verb}
& egzistovati; international-register ``exist''
& \profile{$\checkmark$}{\evstrong}{\evstrong}{\evstrong}{\evstrong}. Internacionalna familija jest \emph{egzist-}, soglasno oficialnym lemmam. Ona jest stilističny alternativ k \emph{obstajati}, ale ne směšaje se s \emph{sučstveny} ``essential''. Starše \emph{eksist-} jest samo ortografičny sloj.\\[-0.1em]
&& \emph{EN:} The official international family is \emph{egzist-}. It can alternate by register with \emph{obstajati}, but not with the distinct concept ``essential.''\\
\addlinespace

\textbf{hoti}
& hoti, premda; concessive ``although/even if''
& \profile{$\checkmark$}{\evstrong}{\partialev}{\evstrong}{\evstrong}. Upotrěbljeno samo v stvarno dopustnyh klauzah. Jedinica, kde površina označava nižnju granicu, jest \emph{prinajmenje}; zato homograf ne upravljaje globalno.\\[-0.1em]
&& \emph{EN:} Reserved for genuine concessive clauses. A look-alike meaning ``at least'' is rendered \emph{prinajmenje}, not forced into the conjunction.\\
\addlinespace

\textbf{identičnost}
& točna identiteta; German \emph{Identität, identisch}
& \profile{$\checkmark$}{\evstrong}{\evstrong}{\evstrong}{\evstrong}. Upotrěbljeno za točnu jednakost ili identitet objektov. \emph{Soglasnost} jest užše za podudaranje definicij, a \emph{sběgati se} za glagolno coinciding/converging. Tuto razděljenje odklanjaje jeden nejasny \emph{sovpad-} korenj.\\[-0.1em]
&& \emph{EN:} Used for exact identity. Agreement of definitions is \emph{soglasnost}; verbal coincidence or convergence is \emph{sběgati se}.\\
\addlinespace

\textbf{imaje město} / \textbf{imajut město}
& analitičny matematičny predikat; ``holds/is the case''
& \profile{$\checkmark$}{\partialev}{\evstrong}{\evstrong}{\evstrong}. Transparentna analitična konstrukcija pokryvaje gole predikaty, gdje nacionalne tradicije razhodet se. Ona ne zaměnja preciznějše \emph{jest pravdivo}, \emph{izpolnjaje se}, \emph{može se priměniti} ili \emph{ostavaje v silě}, kogda iztočnik imaje tu užšu funkciju.\\[-0.1em]
&& \emph{EN:} A transparent analytic predicate where national mathematical verbs diverge. More specific truth, satisfaction, applicability, and force constructions remain distinct.\\
\addlinespace

\textbf{imenno}\\\emph{prislov / adverb}
& imenno, točno nazvano; German \emph{nämlich}
& \profile{$\checkmark$}{\evstrong}{\evstrong}{\evstrong}{\evstrong}. Užša appozicijska veza: ona imenuje ili točno razkriva prědhodny predmet. \emph{To jest} pokryvaje širšu ekvivalenciju i parafrazu. Imenica \emph{imeno} jest homografična granica i ne podlěži toj normalizaciji.\\[-0.1em]
&& \emph{EN:} The narrower marker ``namely,'' used to name or unpack the preceding object. Broader restatement takes \emph{to jest}; noun \emph{imeno} is a protected homograph.\\
\addlinespace

\textbf{izpolnjati} / \textbf{izpolniti}
& izpolniti uslov, relaciju ili ravnanje; satisfy/fulfil
& \profile{$\checkmark$}{\evstrong}{\evstrong}{\evstrong}{\evstrong}. Oficialna familija daje nesvršeny i svršeny aspekt za formalno izpolnjenje. Dativno upravljanje naslědženogo \emph{udovletvor-} korenja ne jest kopirano: matematičny objekt stavaje direktny akuzativny dopolnitelj, gdje věta to trěbuje.\\[-0.1em]
&& \emph{EN:} The official aspect pair expresses satisfaction or fulfilment of equations, relations, and conditions. Case government is recast rather than copied from the rejected Russian-shaped family.\\
\addlinespace

\textbf{jasno} / \textbf{jednoznačny}
& jasno/evidentno nasuprot explicitno ili jednoznačno
& \profile{$\checkmark$}{\evstrong}{\partialev}{\evstrong}{\evstrong}. Njemački ``öffentlich/deutlich/eindeutig''-podobny konteksty ne sut jedna kategorija. Gdje smysl jest unambiguous/explicit, izbrano jest \emph{jednoznačny}; gdje zaključak jest clear/evident, \emph{jasno}. Površina \emph{public-} ne jest dopuščena kako lažni prijatelj.\\[-0.1em]
&& \emph{EN:} Explicit or unambiguous readings take \emph{jednoznačny}; clear or evident readings take \emph{jasno}. A public/clear false friend is not allowed to collapse them.\\
\addlinespace

\textbf{ješče}
& ješče, još, still/yet/also
& \profile{$\checkmark$}{\evstrong}{\evstrong}{\evstrong}{\evstrong}. Oficialny zapis imaje podporu iz vseh glavnyh slovjanskyh tradicij i veže južnyh \emph{še, još/jošte, ušte, ošte}. Zato južna površina \emph{još} ne ostavaje kako nerazlučeny dublet.\\[-0.1em]
&& \emph{EN:} The official bridge connects East and West forms with South Slavic \emph{še, još/jošte, ušte, ošte}; \emph{još} is therefore not kept as an unmotivated duplicate.\\
\addlinespace

\textbf{jediničny}
& jediničny, unit-; German compounds with \emph{Einheits-}
& \profile{$\checkmark$}{\evstrong}{\evstrong}{\evstrong}{\evstrong}. Produktivny pridavnik od \emph{jedinica}; aktualna forma piše \emph{č}. \emph{Jedinicny} bez palatalizacije jest ortografičny ostatok. Termin ostavaje odvojen od \emph{pojediny} ``individual''.\\[-0.1em]
&& \emph{EN:} Productive adjective from \emph{jedinica} for unit constructions. It is distinct from \emph{pojediny} ``individual.''\\
\addlinespace

\textbf{kde}
& gdje/kde; German \emph{wo}
& \profile{$\checkmark$}{\evstrong}{\evstrong}{\evstrong}{\evstrong}. Oficialna lemma imaje pozitivny dokaz kroz vse navedene slovjanske standardy. Formy \emph{gde, gdje, gdě} ne noset smyslovu razliky i zato ne ostavajut kako paralelny sloj.\\[-0.1em]
&& \emph{EN:} The official headword has unusually broad cross-Slavic support. \emph{Gde/gdje/gdě} had no semantic distinction and were not retained as a dialect mosaic.\\
\addlinespace

\textbf{kontrakcija}\\\emph{ženski rod / feminine noun}
& tehnična kontrakcija; Noetherino \emph{Faltung/Faltungsprozess}
& \profile{$\checkmark$}{\evstrong}{\partialev}{\evstrong}{\evstrong}. Internacionalna oficialna imenica jest upotrěbljena doslědno v cěloj paradigmi. Glagolne konteksty imajut \emph{izvršiti kontrakciju}; samokontrakcija se opisuje eksplicitno. Kandidaty iz običnogo ``folding'' registra ne dokazujut Noetherin tehničny smysl.\\[-0.1em]
&& \emph{EN:} The official international noun consistently renders Noether's technical \emph{Faltung}. Verbal contexts use \emph{izvršiti kontrakciju}; ordinary folding words were not assumed to carry the technical sense.\\
\addlinespace

\textbf{korespondencija}
& korespondencija, prirědženje; German \emph{Zuordnung, Korrespondenz}
& \profile{$\checkmark$}{\evstrong}{\evstrong}{\evstrong}{\evstrong}. Imenska funkcija jest odvojena od glagola \emph{odgovarjati} i pridavnika \emph{odgovarjajuči}. Diskursno ``analogously'' ostavaje \emph{analogično}; ordered/respective struktury sut raděje eksplicitno razpisane.\\[-0.1em]
&& \emph{EN:} The noun covers correspondence or assignment. Verbal and adjectival functions use \emph{odgovarjati/odgovarjajuči}; discourse analogy and ordered relations are treated separately.\\
\addlinespace

\textbf{kromě}
& izključajuči, osim; German \emph{außer}
& \profile{$\checkmark$}{\evstrong}{\evstrong}{\evstrong}{\evstrong}. Oficialna forma pokryvaje izključnu prepozicijsku funkciju. Ona ne směšaje se s \emph{pored} ``beside/in addition to'', čije površine mogut izgledati blizko v nekterem nacionalnom jezyku.\\[-0.1em]
&& \emph{EN:} The official exclusion preposition. It is kept distinct from \emph{pored} ``beside/in addition to.''\\
\addlinespace

\textbf{libovoljny}
& libovoljny, arbitrary/any in mathematical quantification
& \profile{$\checkmark$}{\evstrong}{\evstrong}{\evstrong}{\evstrong}. Izbrana familija razprostranjaje razpoznavnost izvan naslědženogo vzhodnogo \emph{proizvolj-} i južnogo \emph{ljub-} sloja. Padež, rod i čislo sut ohranjene pri vsakoj upotrěbě; termin ne zaměnjaje običny ``beloved/dear'' smysl.\\[-0.1em]
&& \emph{EN:} The official family gives a shared mathematical quantifier without defaulting to the inherited East-Slavic or South-Slavic layers. Inflection is preserved occurrence by occurrence.\\
\addlinespace

\textbf{matričny}
& matričny, matrix-; German \emph{Matrix-}
& \profile{$\checkmark$}{\evstrong}{\evstrong}{\evstrong}{\evstrong}. Produktivny pridavnik od oficialnoj \emph{matrica}. Završno \emph{-ny} mora soglasovati se s mužskym \emph{prsten, komponent} i drugymi zaglavnikami; neutro \emph{matrično} ne može byti nesklanjany tehničny etiketa.\\[-0.1em]
&& \emph{EN:} Productive adjective from \emph{matrica}. Its ending agrees with the noun; neuter \emph{matrično} is not treated as an indeclinable label.\\
\addlinespace

\textbf{najmenše obče mnogokratno}
& najmalo obče mnogokratno; least common multiple
& \profile{$\checkmark$}{\evstrong}{\evstrong}{\evstrong}{\evstrong}. Tehnična fraza jest odvojena od prislovnogo \emph{prinajmenje} ``at least''. Aktualna osnovna ortografija piše \emph{obče}; to jest ta sama lemma kako etimologično \emph{obće}.\\[-0.1em]
&& \emph{EN:} The technical noun phrase is distinct from adverbial \emph{prinajmenje} ``at least.'' Basic \emph{obče} and etymological \emph{obće} are spelling layers of the same lemma.\\
\addlinespace

\textbf{najprvo}
& najprvo, first of all; German \emph{zunächst}
& \profile{$\checkmark$}{\evstrong}{\evstrong}{\evstrong}{\evstrong}. Oficialny diskursny prislov upravljaje početkom argumenta. Arhaično ili regionalno \emph{najprěd} ne daje drugu logičnu funkciju i zato ne ostavaje kako dublet.\\[-0.1em]
&& \emph{EN:} Official discourse adverb for the first step of an argument. Regional \emph{najprěd} carried no distinct logical function.\\
\addlinespace

\textbf{Neter-}\\\emph{integrirany eponim / integrated eponym}
& sklanjany eponimi v běžnom medžuslovjanskom tekstu
& \profile{$\checkmark$}{\noev}{\partialev}{\evstrong}{\evstrong}. Kompaktny korenj dovoljava oblike kako \emph{Neterov, Netera} i determinističnu kirilicu bez privilegovanja ruskoj ili ukrajinskoj adaptaciji. Točno \emph{Noether} ostavaje netknuto v avtorskyh kreditah, citatah, naslovah i bibliografičnyh identitetah.\\[-0.1em]
&& \emph{EN:} The compact stem is used only for morphologically integrated running-text eponyms. Exact \emph{Noether} is preserved in credits, citations, titles, and bibliographic identities.\\
\addlinespace

\textbf{někoj} / \textbf{několiko}
& někoj, izvěstny; nekoliko, několiko
& \profile{$\checkmark$}{\evstrong}{\evstrong}{\evstrong}{\evstrong}. Oficialne nezaimenne i količinske formy sut pisane s \emph{ě}. One popravljajut ne samo ortografiju, ale i razlučajut ``někoj/certain'' od jussivnogo \emph{nehaj}. Zato jedna površina ne može avtomatično zaměniti drugu.\\[-0.1em]
&& \emph{EN:} Official indefinite forms. They are semantically distinct from jussive \emph{nehaj}; apparent surface similarity is not a license for global replacement.\\
\addlinespace

\textbf{nyně}
& sada, nyní, now; German \emph{nun, jetzt}
& \profile{$\checkmark$}{\evstrong}{\evstrong}{\evstrong}{\evstrong}. Oficialna lemma veže vzhodne formy s češskym \emph{nyní} i zapadnym književnym cognatom. Naslědženo \emph{teper} bylo silno vzhodno skoncentrovano, a južno \emph{sada} by samo preměstilo tu neravnověsnost.\\[-0.1em]
&& \emph{EN:} The official bridge connects East forms with Czech \emph{nyní} and related West usage. \emph{Teper} was East-heavy; replacing it by South-heavy \emph{sada} would only move the imbalance.\\
\addlinespace

\textbf{obči} / \textbf{občnost}
& obči, opči, common/general; generality
& \profile{$\checkmark$}{\evstrong}{\evstrong}{\evstrong}{\evstrong}. Oficialny etimologičny korenj jest \emph{obć-}; izdanje pokazuje osnovno \emph{obč-}. Komunikacijske lažne srodnosti ne sut zadržane: po smyslu sut \emph{publikacija, informacija, komunikacija, prědstaviti} ili \emph{izvěstiti}.\\[-0.1em]
&& \emph{EN:} Basic \emph{obč-} represents official etymological \emph{obć-}. Communication false friends are recast as publication, information, communication, presentation, or notification according to the source.\\
\addlinespace

\textbf{obstajati} / \textbf{obstajanje}
& obstajati, byti prisutny; exist, subsist, persist
& \profile{$\checkmark$}{\evstrong}{\evstrong}{\evstrong}{\evstrong}. Oficialna familija pokryvaje egzistenciju i trajuče obstajanje. Ona ne znači \emph{sostojati se iz} ``consist of'' i ne směšaje se s \emph{sučstveny} ``essential''. Internacionalno \emph{egzistovati} može služiti drugomu registru.\\[-0.1em]
&& \emph{EN:} The official family expresses existence or persistence. It is distinct from composition \emph{sostojati se iz} and essence \emph{sučstveny}; international \emph{egzistovati} is a register alternative.\\
\addlinespace

\textbf{odgovarjati} / \textbf{odgovarjajuči}
& odgovarjati čemu; corresponding, associated
& \profile{$\checkmark$}{\evstrong}{\evstrong}{\evstrong}{\evstrong}. Glagolna i pridavnička funkcija jest odvojena od imenice \emph{korespondencija}. Naslědženy \emph{sootv-} sloj byl razumliv, ale vzhodno skoncentrovany i morfologično nejednaky; globalna zaměna by takože pogrěšno zahvatila ``respectively''.\\[-0.1em]
&& \emph{EN:} Verb and adjective for mathematical correspondence. The East-heavy \emph{sootv-} layer was not retained, and ordered ``respectively'' constructions are written explicitly rather than forced into this family.\\
\addlinespace

\textbf{oprěděliti} / \textbf{oprědělenny}
& oprěděliti, determinovati; define/determine, definite
& \profile{$\checkmark$}{\evstrong}{\evstrong}{\evstrong}{\evstrong}. Oficialna familija može nositi nekoliko blizkyh smyslov, ktore věta razlučaje. Zato ortografična popravka ohranjaje derivaciju, a malformed participi, imenice i glagolne formy sut rěšene po svojoj sintaksi; univerzalno \emph{definovati} by izbrisalo ``determine/fixed''.\\[-0.1em]
&& \emph{EN:} The official family spans define, determine, definite, and indeterminate functions. Sentence grammar selects the form; global \emph{definovati} would erase real distinctions.\\
\addlinespace

\textbf{po odnošenju k}
& po odnošenju k, vzhledem k; with respect to, relative to
& \profile{$\checkmark$}{\evstrong}{\evstrong}{\evstrong}{\evstrong}. Oficialna analitična prepozicija ima sravnjivatelje v vzhodnyh, zapadnyh i južnyh jezykah. Ona jasno uvodi bazu, polje, operator, grupu, rang ili invariant. Hibridno \emph{vzhodno k} ne jest zasvědočena standardna lemma.\\[-0.1em]
&& \emph{EN:} The official analytic preposition has comparators across all three branches and transparently marks mathematical relativity. Hybrid \emph{vzhodno k} lacked standard support.\\
\addlinespace

\textbf{početkovy} / \textbf{odpravny punkt}
& početny, izhodny punkt; initial, starting point
& \profile{$\checkmark$}{\evstrong}{\evstrong}{\evstrong}{\evstrong}. Jedna naslědžena korenna familija byla razdělena: početny/startni smysl daje \emph{početkovy} ili \emph{odpravny punkt/točka}; stvarno porijeklo \emph{proizhoditi iz}; logična poslědica \emph{slěduje}; nastavanje \emph{nastavaje}.\\[-0.1em]
&& \emph{EN:} Initial point, origin, consequence, and arising are not collapsed: \emph{početkovy/odpravny punkt}, \emph{proizhoditi iz}, \emph{slěduje}, and \emph{nastavaje} divide the functions.\\
\addlinespace

\textbf{pojediny}
& pojediny, individual/single; German \emph{einzeln}
& \profile{$\checkmark$}{\evstrong}{\evstrong}{\evstrong}{\evstrong}. Oficialna lemma i već stabilna korpusna familija dajut jednu površinu. Češsko-slovačko \emph{jednotliv-} jest razumno, ale bi ostalo regionalny sloj bez smyslove razlike.\\[-0.1em]
&& \emph{EN:} Official and already established corpus term for individual/single. Czech/Slovak-shaped \emph{jednotliv-} carried no distinct meaning.\\
\addlinespace

\textbf{polna reducibilnost}
& complete reducibility; German \emph{vollständige Reduzibilität}
& \profile{$\checkmark$}{\partialev}{\partialev}{\evstrong}{\evstrong}. Tehničny kolokat jest stabilny v korpusu. On mora ostati odvojen od obyčnogo prislova \emph{cělkom} i od field-theoretic \emph{doskonalo}. Jedna globalna zaměna ``complete/perfect'' by razrušila terminologiju.\\[-0.1em]
&& \emph{EN:} Stable technical collocation, kept distinct from ordinary \emph{cělkom} and field-theoretic \emph{doskonalo}.\\
\addlinespace

\textbf{pomočny} / \textbf{s pomočju}
& pomočny, pomocny; auxiliary, with the help of
& \profile{$\checkmark$}{\evstrong}{\evstrong}{\evstrong}{\evstrong}. Oficialny etimologičny korenj \emph{pomoć-} jest prikazany v osnovnoj ortografiji kako \emph{pomoč-}. Termin pokryvaje matematične pomočny konstrukcije i instrumentalnu vezu; ne uvodi novy smysl.\\[-0.1em]
&& \emph{EN:} Basic \emph{pomoč-} represents official etymological \emph{pomoć-}. It covers auxiliary constructions and ``with the help of'' without changing the lemma.\\
\addlinespace

\textbf{ponovno}
& ponovno, opet; again/anew
& \profile{$\checkmark$}{\evstrong}{\evstrong}{\evstrong}{\evstrong}. Oficialna forma jest ravnověsny čitateljski izbor za povtorjenje operacije ili argumenta. Varijanty \emph{znovu/znovo} ne sut ohranjene kako nerazlučeny regionalny dublety.\\[-0.1em]
&& \emph{EN:} Official recurrence adverb for repeating an operation or argument. Regional \emph{znovu/znovo} variants had no separate function.\\
\addlinespace

\textbf{pored} / \textbf{poredok}
& pored: obok/dodatno; poredok: order/sequence
& \profile{$\checkmark$}{\evstrong}{\evstrong}{\evstrong}{\evstrong}. Aktualna osnovna ortografija prikazuje etimologične \emph{poręd/porędok} bez nosnogo znaka. Dvě lemmy ostavajut razlučene po značenju i paradigmi. \emph{Kromě} znači ``except'' i ne jest zaměna za \emph{pored}.\\[-0.1em]
&& \emph{EN:} Basic spelling of etymological \emph{poręd/porędok}. The preposition and order noun remain distinct, and exclusion \emph{kromě} is not substituted.\\
\addlinespace

\textbf{poslě}
& poslě, nakon; after/afterwards
& \profile{$\checkmark$}{\evstrong}{\evstrong}{\evstrong}{\evstrong}. Oficialna forma ohranjaje etimologično \emph{ě}. Neakcentovano \emph{posle} jest ortografičny ostatok. Logično ``therefore'' ne pripada toj lemmi, ale \emph{zato}.\\[-0.1em]
&& \emph{EN:} Official form for temporal ``after.'' Logical consequence remains \emph{zato}; unaccented \emph{posle} is only an orthographic residue.\\
\addlinespace

\textbf{potrěbny} / \textbf{potrěbno} / \textbf{trěba}
& potrěbny, nužno; necessary/needed, it is necessary
& \profile{$\checkmark$}{\evstrong}{\evstrong}{\evstrong}{\evstrong}. Familija nosi imensku, pridavničku i predikativnu potrěbnost. Negacija jest rěšena oprezno: ``ne nužno $P$'' ne znači ``$P$ ne jest potrěbno''. \emph{Neizběžno} bylo odklonjeno, ibo znači silněju neizběžnost.\\[-0.1em]
&& \emph{EN:} The family expresses mathematical necessity while preserving the scope of negation. \emph{Neizběžno} was rejected because inevitability is stronger than necessity.\\
\addlinespace

\textbf{prědpokladati} / \textbf{prědpoložiti, prědpoloženje}
& prědpokladati; assume/presume, supposition
& \profile{$\checkmark$}{\evstrong}{\evstrong}{\evstrong}{\evstrong}. Nesvršena radnja upotrěbjaje \emph{prědpokladati}; svršena, pasivna, pridavnička i imenska funkcija \emph{prědpolož-}. Rod i soglasovanje se měnjajut, kogda ženska imenica ustupa neutru \emph{prědpoloženje}.\\[-0.1em]
&& \emph{EN:} Imperfective assumption takes \emph{prědpokladati}; perfective, passive, adjectival, and supposition-noun functions use \emph{prědpolož-}. Agreement follows the selected noun.\\
\addlinespace

\textbf{prědstaviti} / \textbf{prědstavjenje}
& prědstaviti, reprezentovati; representation/presentation
& \profile{$\checkmark$}{\evstrong}{\evstrong}{\evstrong}{\evstrong}. Oficialna familija daje tip-1 morfologiju: \emph{prědstavja, prědstavjajut, prědstavjenje}. Nacionalno oblikovano \emph{predstavlj-} i zero-ending genitiv množiny ne sut zadržane. Mentalna \emph{prědstava} ne zaměnjaje matematično \emph{prědstavjenje}.\\[-0.1em]
&& \emph{EN:} Official type-1 family for mathematical presentation or representation. South-shaped \emph{-lj-}, zero genitive plurals, and the mental-image noun are not conflated with it.\\
\addlinespace

\textbf{prěz}
& prěz, kroz, via; through/across/by
& \profile{$\checkmark$}{\evstrong}{\evstrong}{\evstrong}{\evstrong}. Oficialna lemma imaje široku bělorusku, poljsku, češsku, bulgarsku i drugu zapadno/južnu cognatnu podporu. Ona ravnověsi naslědženo vzhodno koncentrirano \emph{črez/črěz}; poljsko \emph{przez} jest normalizovano k istoj površini.\\[-0.1em]
&& \emph{EN:} The official form has broad Belarusian, Polish, Czech, Bulgarian, and related West/South support. It balances the inherited East-heavy \emph{črez/črěz} layer.\\
\addlinespace

\textbf{prinaležati}
& prinaležati, belong to; German \emph{angehören}
& \profile{$\checkmark$}{\evstrong}{\evstrong}{\evstrong}{\evstrong}. Oficialna familija nosi matematično členstvo i pripadnost. Množinske glagolne formy prěmaju pravilno \emph{-et}; naslědženo \emph{prinadlež-} ne ostavaje paralelny vzhodny sloj.\\[-0.1em]
&& \emph{EN:} Official family for membership or belonging, with number-correct verb morphology. East-shaped \emph{prinadlež-} is not retained as a duplicate.\\
\addlinespace

\textbf{prsten}; \textbf{podprsten, nadprsten, medžuprsten}
& algebraičny ring i njegovi odnosi; German \emph{Ring, Unterring, Oberring, Zwischenring}
& \profile{$\checkmark$}{\partialev}{\evstrong}{\evstrong}{\evstrong}. Šest-grupno matematično sravnjenje dalo korenju \emph{prsten/pierścień} 2.5 jedinice, prěd \emph{kolco} 2.0, \emph{okruh} 1.0 i \emph{kolobar} 0.5. Oficialny fizičny smysl podpiraje razpoznavanje, a matematične knjigy tehničny smysl. Prvy pojav ima vzhodny most.\\[-0.1em]
&& \emph{EN:} Six-group comparison scored the \emph{prsten/pierścień} root 2.5, ahead of \emph{kolco} 2.0, \emph{okruh} 1.0, and \emph{kolobar} 0.5. The first occurrence gives an East-Slavic bridge.\\
\addlinespace

\textbf{pytanje}
& pytanje, problem/question; German \emph{Frage}
& \profile{$\checkmark$}{\evstrong}{\evstrong}{\evstrong}{\evstrong}. Oficialna type-1 familija jest čitateljski širša od vzhodno skoncentrovanogo \emph{vopros}. Padežna paradigma ostavaje regularna: \emph{pytanje, pytanja, pytanjah, pytanjam}.\\[-0.1em]
&& \emph{EN:} The official family is used instead of East-heavy \emph{vopros}, with the regular case paradigm preserved.\\
\addlinespace

\textbf{razgledati} / \textbf{razgled}
& razgledati, razmotriti, examine/consider
& \profile{$\checkmark$}{\evstrong}{\evstrong}{\evstrong}{\evstrong}. Oficialna familija pokryvaje pregled, ispitivanje i větnu komandu ``razgledajmo''. Naslědžene \emph{smatr-} i \emph{razsmotr-} familije sut regionalne i mogle bi prizvati perceptivny smysl. Eksplicitno rezonovanje ostavaje \emph{razvažanje}.\\[-0.1em]
&& \emph{EN:} Official family for examination or consideration, including ``let us consider.'' Explicit reasoning/deliberation remains \emph{razvažanje}.\\
\addlinespace

\textbf{razvažanje}
& razvažanje, promišljanje; German \emph{Überlegung}
& \profile{$\checkmark$}{\evstrong}{\evstrong}{\evstrong}{\evstrong}. Upotrěbljeno samo gdje iztočnik imenuje čin ili rezultat rezonovanja. German \emph{Betrachtung} i ordinary ``consider'' větinom ostavajut v \emph{razgled-} familiji.\\[-0.1em]
&& \emph{EN:} Reserved for explicit reasoning or deliberation. Ordinary examination/consideration remains in the \emph{razgled-} family.\\
\addlinespace

\textbf{razširjenje}
& razširjenje, extension/expansion
& \profile{$\checkmark$}{\evstrong}{\evstrong}{\evstrong}{\evstrong}. Oficialna derivacija od \emph{razširiti} daje participialnu i imensku familiju \emph{razširjen-}. Hrvatsko oblikovano \emph{proširenje} jest razumljivo, ale ne ostavaje regionalny dublet bez smyslove razlike.\\[-0.1em]
&& \emph{EN:} Official derivational family from \emph{razširiti}. Croatian-shaped \emph{proširenje} carried no separate meaning.\\
\addlinespace

\textbf{red}
& red, porędok, series/rank/order according to context
& \profile{$\checkmark$}{\evstrong}{\evstrong}{\evstrong}{\partialev}. Osnovna ortografija prikazuje etimologično \emph{ręd} kako \emph{red}. Lažni pridavnički derivaty ne sut mehanično ohranjeni: kontekst može zahtěvati \emph{serija, rang} ili \emph{poredok}. Točne matematične identifikatory \code{red} ostavajut identiteti, ne riječi.\\[-0.1em]
&& \emph{EN:} Basic spelling of etymological \emph{ręd}. Context separates series, rank, and order; literal mathematical identifiers remain protected identities.\\
\addlinespace

\textbf{sběgati se}
& sběgati se, sovpadati ili konvergirati po kontekstu
& \profile{$\checkmark$}{\evstrong}{\evstrong}{\evstrong}{\evstrong}. Oficialna familija pokryvaje glagolnu funkciju coinciding/converging. Imenična točna identiteta jest \emph{identičnost}; soglasje definicij \emph{soglasnost}. Jedna vzhodno oblikovana \emph{sovpad-} familija by skrila tu trojnu razliku.\\[-0.1em]
&& \emph{EN:} Official verbal family for coincidence or convergence. Exact identity and agreement of definitions use separate nouns.\\
\addlinespace

\textbf{shranjati} / \textbf{shranjenje}
& ohranjati, conservation/preservation; German \emph{erhalten, Erhaltung}
& \profile{$\checkmark$}{\evstrong}{\evstrong}{\evstrong}{\evstrong}. Oficialna familija pokryvaje ohranjanje relacij, rangov, regularnosti, svojstv, form i fizičnyh zakonov. \emph{Ohranjati} jest užše guard/protect; \emph{ostavati} ne može imenovati zakon ohranjanja.\\[-0.1em]
&& \emph{EN:} Official family for preservation and conservation in mathematics and physics. Guard/protect and remain constructions are too narrow for the full range.\\
\addlinespace

\textbf{slěduje} / \textbf{slědujuči}
& slěduje: follows; slědujuči: following
& \profile{$\checkmark$}{\evstrong}{\evstrong}{\evstrong}{\evstrong}. Logična poslědica, větny glagol i pridavnička orientacija sut morfologično razlučene. Aktualna osnovna ortografija piše etimologično \emph{slědujući} kako \emph{slědujuči}. Vremensko ``then'' jest \emph{togda}, zaključny marker \emph{zato}.\\[-0.1em]
&& \emph{EN:} Consequence, finite verb, and following adjective are morphologically distinguished. Temporal ``then'' is \emph{togda}; inferential ``therefore'' is \emph{zato}.\\
\addlinespace

\textbf{sodrživati} / \textbf{sodržanje}
& sodrživati element/podobjekt; content
& \profile{$\checkmark$}{\evstrong}{\evstrong}{\evstrong}{\evstrong}. Oficialny nesvršeny glagol nosi matematično containment, imenica \emph{sodržanje} content. \emph{Sodržati} oficialno znači hold back/inhibit i zato ne jest containment. \emph{Zadržati} ostavaje za svoju samostojnu funkciju.\\[-0.1em]
&& \emph{EN:} Official \emph{sodrživati} expresses containment and \emph{sodržanje} content. \emph{Sodržati} means hold back/inhibit and is not repurposed as a false friend.\\
\addlinespace

\textbf{soglasnost}
& soglasnost ili podudaranje definicij; German \emph{Übereinstimmung}
& \profile{$\checkmark$}{\evstrong}{\evstrong}{\evstrong}{\evstrong}. Upotrěbljeno v dvuh zaglavjah, kde iztočnik eksplicitno govori o soglasnosti definicij. Točna identiteta objektov ostavaje \emph{identičnost}; glagolno sběganje \emph{sběgati se}.\\[-0.1em]
&& \emph{EN:} Used specifically for agreement of definitions. Exact identity and verbal coincidence remain separate.\\
\addlinespace

\textbf{sostojati se iz}
& sostojati se iz, consist of
& \profile{$\checkmark$}{\evstrong}{\evstrong}{\evstrong}{\evstrong}. Oficialna pan-slovjanska veza imaje pozitivny dokaz kroz vse grupy. Refleksivno \emph{se} i upravljanje s \emph{iz} sut obvezne pri kompoziciji. Konstrukcije ``problem jest v tom'' ili ``sučnost leži v'' ne sut silom ugurane pod tu lemmu.\\[-0.1em]
&& \emph{EN:} Official broadly supported composition phrase. Reflexive \emph{se} and \emph{iz} government are required; non-composition senses are recast instead of forced into the lemma.\\
\addlinespace

\textbf{sumand}
& sumand, pribrojenik; German \emph{Summand}
& \profile{$\checkmark$}{\partialev}{\evstrong}{\evstrong}{\evstrong}. Internacionalny matematičny termin jest pisan s jednoj \emph{m}. Naslědženo \emph{sčestnik} bylo leksično nejasno, a dvojno \emph{summand} nepotrěbno ohranjalo německu grafiju.\\[-0.1em]
&& \emph{EN:} International mathematical term with one \emph{m}. The opaque inherited calque and German double consonant were not retained.\\
\addlinespace

\textbf{sučstveny} / \textbf{v sučnosti}
& esencialny / v sučnosti; German \emph{wesentlich / im Wesen}
& \profile{$\checkmark$}{\evstrong}{\evstrong}{\evstrong}{\evstrong}. Pridavnik, imenica i idiom sut razlučeni od prostoj egzistencije. Oficialny etimologičny zapis može imati opcionalne znaky; osnovna površina izdanja piše \emph{sučstveny, sučnost}.\\[-0.1em]
&& \emph{EN:} Essentiality and ``in essence'' are kept distinct from existence. The reader uses the declared basic spelling \emph{sučstveny/sučnost}.\\
\addlinespace

\textbf{tako že} / \textbf{takože}
& tako že: so/so that; takože: also
& \profile{$\checkmark$}{\evstrong}{\evstrong}{\evstrong}{\evstrong}. Razmak nosi smyslovu funkciju. Dvě riječi uvodet poslědicu ili cělj; jedna riječ dodava paralelny element. Površiny \emph{tak že, takže, takođe} sut klasifikovane po větnom odnosu, ne slepo po slovu.\\[-0.1em]
&& \emph{EN:} Spacing marks a real distinction: two words introduce result or purpose, while one word means ``also.'' Older variants were classified by clause relation.\\
\addlinespace

\textbf{to jest}
& to jest, in other words; German \emph{das heißt}
& \profile{$\checkmark$}{\partialev}{\evstrong}{\evstrong}{\evstrong}. Široky marker ekvivalencije, reformulacije ili objasnjenja. \emph{Imenno} jest užše ``namely''; \emph{zato} označaje zaključak. Na linijah s oběma funkcijama sam lokálny větny odnos rěšaje izbor.\\[-0.1em]
&& \emph{EN:} Broad marker of equivalence or restatement. Narrow \emph{imenno} and inferential \emph{zato} remain distinct, including where functions occur on the same line.\\
\addlinespace

\textbf{togda}
& togda, then in time or condition
& \profile{$\checkmark$}{\evstrong}{\evstrong}{\evstrong}{\evstrong}. Oficialna vremenska/uslovna funkcija jest odvojena od zaključnogo \emph{zato}. Naslědženo \emph{tada} ne dava drugy smysl i zato ne ostavaje paralelny regionalny sloj.\\[-0.1em]
&& \emph{EN:} Temporal or conditional ``then,'' distinct from inferential \emph{zato}. Regional \emph{tada} had no separate function.\\
\addlinespace

\textbf{toj, ta, to} / \textbf{tu}
& demonstrativy / lokativny prislov
& \profile{$\checkmark$}{\evstrong}{\evstrong}{\evstrong}{\evstrong}. Atributivny njemački \emph{dieser/diese/dieses} prěmaje soglasujuči demonstrativ; prislovno \emph{hier} prěmaje \emph{tu}. Arhaično \emph{sej} i neanalizovano \emph{tuta} ne sut upotrěbjene kako univerzalne zaměny.\\[-0.1em]
&& \emph{EN:} Agreeing demonstratives render adjectival ``this/that''; adverbial ``here'' is \emph{tu}. The two functions are not collapsed.\\
\addlinespace

\textbf{tym samym}
& tym samym, thereby/thus
& \profile{$\checkmark$}{\partialev}{\evstrong}{\evstrong}{\evstrong}. Upotrěbljeno kako collision-safe zaključna veza, osobito kogda \emph{zato} bi neposrědno stalo prěd uzročno \emph{zato že}. Tuto ne jest slobodny ukrasny sinonim, ale sposob ohranjanja logične struktury.\\[-0.1em]
&& \emph{EN:} A collision-safe inferential phrase, especially where plain \emph{zato} would stand beside causal \emph{zato že}. It preserves clause logic rather than adding stylistic variation.\\
\addlinespace

\textbf{uslovje}
& uslov, condition; neuter noun
& \profile{$\checkmark$}{\evstrong}{\evstrong}{\evstrong}{\evstrong}. Oficialna lemma jest neutrum; pridavnik, zaimennik i glagol morajut soglasovati se s njim. Zato leksičny izbor uključaje lokalne popravky roda, ne samo zaměnu korenja.\\[-0.1em]
&& \emph{EN:} The official noun is neuter. Adjectives, pronouns, and predicates must agree with it, so the lexical decision includes local gender repair.\\
\addlinespace

\textbf{vsegda}
& vsegda, always
& \profile{$\checkmark$}{\evstrong}{\evstrong}{\evstrong}{\evstrong}. Oficialna i već dominantna korpusna forma. \emph{Vsagda} i etimologično preopterečeno \emph{vsěgda} ne sut samostojni smyslovi; aktualny čitateljski zapis jest \emph{vsegda}.\\[-0.1em]
&& \emph{EN:} Official and corpus-dominant form for ``always.'' Other spellings do not mark separate meanings.\\
\addlinespace

\textbf{vměsto}
& vměsto, instead of
& \profile{$\checkmark$}{\evstrong}{\evstrong}{\evstrong}{\evstrong}. Oficialna forma jest jedna riječ. \emph{Na město} znači doslovno ``to a place'' i može byti druga konstrukcija; zato samo stvarno prepozicijske upotrěby sut spojene.\\[-0.1em]
&& \emph{EN:} The official preposition is one word. Literal ``to a place'' remains a different construction, so only genuine prepositional uses are joined.\\
\addlinespace

\textbf{vsaky}
& vsaky, each/every
& \profile{$\checkmark$}{\evstrong}{\evstrong}{\evstrong}{\evstrong}. Osnovna površina oficialnogo \emph{vśaky} jest osobito razpoznatljiva v južnom \emph{svak/svaki} i ostavaje jasna vzhodu i zapadu. \emph{Každy} jest oficialny, ale bi zadržal besmyslny iztočny sloj i slabšu južnu razpoznavnost.\\[-0.1em]
&& \emph{EN:} Basic form of official \emph{vśaky}, with particularly direct South-Slavic recognition and good East/West access. Official synonym \emph{každy} was not kept as a source-stratum duplicate.\\
\addlinespace

\textbf{vzajemny}
& vzajemny, mutual/reciprocal
& \profile{$\checkmark$}{\evstrong}{\evstrong}{\evstrong}{\evstrong}. Oficialna familija daje stabilny korenj za uzajemne odnose. Malformirano \emph{vzaim-} jest regionalno-ortografičny ostatok i ne nosi posebny matematičny smysl.\\[-0.1em]
&& \emph{EN:} Official family for mutual or reciprocal relations. Malformed \emph{vzaim-} had no separate mathematical meaning.\\
\addlinespace

\textbf{zato}
& zato, slědovateljno; therefore/hence
& \profile{$\checkmark$}{\evstrong}{\evstrong}{\evstrong}{\evstrong}. Glavny marker matematičnogo zaključka ili rezultata. Zapadne \emph{tedy/teda} i drugi iztočni slojevi sut klasifikovani po klauznoj funkciji; ekvivalencija stavaje \emph{to jest}, uzrok \emph{zato že}, a vremenska poslědovateljnost ne jest avtomatično zaključak.\\[-0.1em]
&& \emph{EN:} Principal marker of mathematical inference or result. Restatement, cause, and temporal sequence are assigned their own forms.\\
\addlinespace

\textbf{zato že}
& zato že, because
& \profile{$\checkmark$}{\evstrong}{\evstrong}{\evstrong}{\evstrong}. Oficialna analitična veza daje najravnověsniji vzhodno-zapadno-južny profil medžu razgledanymi uzročnymi kandidatami. Kratke nacionalne čestice \emph{bo/ibo, jer, poneže} ne sut nametnute cělomu korpusu; postpozitivne i objasnjajuče funkcije sut recast po věti.\\[-0.1em]
&& \emph{EN:} The official analytic conjunction gives the broadest balanced profile among considered causal candidates. National short particles were not imposed corpus-wide.\\
\addlinespace

\textbf{zaradi}
& zaradi, zbog; because of/for the sake of
& \profile{$\checkmark$}{\evstrong}{\evstrong}{\evstrong}{\evstrong}. Oficialna prepozicija jest upotrěbljena v uzročnyh i motivacijskyh konstrukcijah s padežno ohranjenoju dopolniteljnoju frazoju. \emph{Zbog} jest razumljivo, ale ostavaje južny regionalny sloj bez funkcionalnoj potrěby.\\[-0.1em]
&& \emph{EN:} Official preposition for causal or purposive relations, with case preserved. South-shaped \emph{zbog} carried no necessary distinct function.\\
\addlinespace

\textbf{nerazložimy} / \textbf{razložimy}
& irreducible / reducible; German \emph{irreduzibel/reduzibel}
& \profile{$\checkmark$}{\evstrong}{\evstrong}{\evstrong}{\evstrong}. Oficialny korenj \emph{razložiti} i stabilna korpusna familija dajut \emph{razložimost, nerazložimost}. ASCII-sloj \emph{razloz-} ne jest drugi termin. Rod, padež i čislo ostavajut vidljivi v završetku.\\[-0.1em]
&& \emph{EN:} Official root and stable corpus family for reducibility. ASCII \emph{razloz-} is not a competing term; inflection remains explicit.\\
\addlinespace

\textbf{odpravny punkt} / \textbf{proizhoditi iz} / \textbf{nastavati}
& starting point / originate from / arise
& \profile{$\checkmark$}{\evstrong}{\evstrong}{\evstrong}{\evstrong}. Tri blizke iztočne metafory sut razdělene po argumentnoj strukturi. \emph{Odpravny punkt} imenuje bazu; \emph{proizhoditi iz} stvarno porijeklo ili pripisivanje; \emph{nastavati} pojavjenje ili generovanje.\\[-0.1em]
&& \emph{EN:} Starting basis, genuine origin, and arising/generation use distinct constructions, selected by argument structure.\\
\addlinespace

\textbf{pravdivy} / \textbf{jest pravdivo}
& pravdivy, true
& \profile{$\checkmark$}{\evstrong}{\evstrong}{\evstrong}{\evstrong}. Oficialna pridavnička familija imaje široku slovjansku podporu. Ona se upotrěbjaje, kogda iztočnik govori o pravdivosti tvrdženja, ne kako univerzalna zaměna za priměnjivost, izpolnjenje ili ostavanje v silě.\\[-0.1em]
&& \emph{EN:} Broadly supported official truth adjective, used only for actual truth conditions rather than as a universal validity verb.\\
\addlinespace

\textbf{prinajmenje}
& prinajmenje, at least
& \profile{$\checkmark$}{\evstrong}{\evstrong}{\evstrong}{\evstrong}. Upotrěbljeno za nižnju granicu. Ono ne směšaje se s tehničnym \emph{najmenše obče mnogokratno} i ne interpretiraje dopustnu česticu \emph{hoti} mimo větnogo konteksta.\\[-0.1em]
&& \emph{EN:} Adverbial lower-bound marker, distinct from the technical least-common-multiple phrase and from concessive \emph{hoti}.\\
\addlinespace

\textbf{prirěditi} / \textbf{priredženje}
& assign, assignment
& \profile{$\checkmark$}{\evstrong}{\partialev}{\evstrong}{\evstrong}. Upotrěbljeno za eksplicitno pridavanje jednogo objekta drugomu. Aktualna osnovna ortografija piše \emph{e} vměsto opcionalnogo nosnogo znaka. Korespondencija kako odnos ostavaje \emph{korespondencija/odgovarjati}.\\[-0.1em]
&& \emph{EN:} Used for explicit assignment of one object to another. General correspondence remains a separate family.\\
\addlinespace

\textbf{v silě}
& v silě, in force/remains valid
& \profile{$\checkmark$}{\evstrong}{\evstrong}{\evstrong}{\evstrong}. Izbrano za pravilo, zakon ili rezultat, ktory prodolžaje dějstvovati pod novymi uslovjami. Tuto jest jedna od kontekstnyh razgranjenij německogo \emph{gilt}, ne generičny sinonim pravdivosti.\\[-0.1em]
&& \emph{EN:} Used when a rule or result remains in force under changed conditions. It is one context-sensitive branch of German \emph{gilt}, not a synonym for truth.\\
\addlinespace

\textbf{očevidny} / \textbf{očevidno}
& očevidny, evident/obvious
& \profile{$\checkmark$}{\evstrong}{\evstrong}{\evstrong}{\evstrong}. Oficialna familija ima \emph{očevid-}; \emph{očivid-} jest ortografičny nacionalny sloj. Smysl ostavaje dokazna očevidnost i ne zaměnjaje eksplicitnu jednoznačnost \emph{jednoznačny}.\\[-0.1em]
&& \emph{EN:} Official family for evident/obvious. Nationally shaped \emph{očivid-} is not retained, and explicit unambiguity remains \emph{jednoznačny}.\\
\addlinespace

\textbf{napriměr} / \textbf{priměr}
& na priměr / example
& \profile{$\checkmark$}{\evstrong}{\evstrong}{\evstrong}{\evstrong}. Prislov jest jedna riječ \emph{napriměr}; imenica ostavaje \emph{priměr} s regularnym padežom. Nacionalne površiny \emph{na primjer, naprimer} ne označavajut razny register v tom izdanje.\\[-0.1em]
&& \emph{EN:} Adverb \emph{napriměr} and noun \emph{priměr} are distinguished. National spellings do not encode a separate register here.\\
\addlinespace

\textbf{nehaj}
& nehaj, let; mathematical jussive
& \profile{$\checkmark$}{\evstrong}{\evstrong}{\evstrong}{\evstrong}. Upotrěbljeno samo za jussiv ``nehaj $x$ bude''. Konteksty, gdje iztočnik znači ``někoj/some'', sut recast s \emph{někoj}, ne globalno normalizovane k jussivu.\\[-0.1em]
&& \emph{EN:} Reserved for mathematical ``let.'' Indefinite ``some'' contexts use \emph{někoj} and are not swallowed by token-level normalization.\\
\addlinespace

\textbf{už}
& už, already
& \profile{$\checkmark$}{\evstrong}{\evstrong}{\evstrong}{\evstrong}. Oficialna kratka forma jest korpusny standard. \emph{Uže} ne nosi drugy smysl i zato ne ostavaje kako regionalny dublet.\\[-0.1em]
&& \emph{EN:} Official short form for ``already''; \emph{uže} carried no separate function.\\
\addlinespace

\textbf{ljubezno}
& ljubezno, kindly/courteously
& \profile{$\checkmark$}{\evstrong}{\evstrong}{\partialev}{\evstrong}. Oficialny prislov služi v redkom netehničnom diskursnom kontekstu. On ne jest kvantifikator \emph{libovoljny}; podobnost početnogo korenja ne směšaje funkcije.\\[-0.1em]
&& \emph{EN:} Official courtesy adverb in a rare nontechnical context, kept distinct from arbitrary-quantifier \emph{libovoljny}.\\
\addlinespace

\textbf{najkratši} / \textbf{največši} / \textbf{najčestěje}
& superlativy: shortest, greatest, most often
& \profile{$\checkmark$}{\evstrong}{\evstrong}{\evstrong}{\evstrong}. Vsaka forma slijeduje svoju oficialnu pridavničku ili prislovnu paradigmu. One sut vključene kako sistemny čitateljski ključi, ibo matematične granice zaviset od točnogo komparativa i soglasovanja.\\[-0.1em]
&& \emph{EN:} Official superlative families. Exact comparison and agreement matter because mathematical bounds depend on them.\\
\addlinespace

\textbf{poprědny} / \textbf{raněje}
& previous/preceding versus earlier
& \profile{$\checkmark$}{\evstrong}{\evstrong}{\evstrong}{\evstrong}. Pridavničko ``previous/preceding'' jest \emph{poprědny}; prislovno ``earlier'' jest \emph{raněje}. Glagol \emph{prědhoditi} ostavaje samostojna oficijalna familija.\\[-0.1em]
&& \emph{EN:} Adjectival previous/preceding is \emph{poprědny}; adverbial earlier is \emph{raněje}; verb \emph{prědhoditi} remains distinct.\\
\addlinespace

\textbf{treťji}
& treti, third
& \profile{$\checkmark$}{\evstrong}{\evstrong}{\evstrong}{\evstrong}. Oficialna ordinalna familija imaje měkko pridavničko sklanjanje. Južny genitiv \emph{tretjega} ne jest nesklanjany standard za cěly korpus.\\[-0.1em]
&& \emph{EN:} Official ordinal with soft-adjective inflection; a regional genitive is not generalized corpus-wide.\\
\addlinespace

\textbf{vyše} / \textbf{najvyše}; \textbf{bolje}
& above/at most; better
& \profile{$\checkmark$}{\evstrong}{\evstrong}{\evstrong}{\evstrong}. Prislovne familije sut ohranjene odvojeno od pridavničkyh korenjev. \emph{Vyše/najvyše} ne popravjajut avtomatično vse \emph{visok-} ili glagolne formy; \emph{bolje} ne zaměnjaje pravilny pridavnik \emph{bolši}.\\[-0.1em]
&& \emph{EN:} Adverb families are kept distinct from adjectives and verbs; normalization is grammatical, not a blind root substitution.\\
\addlinespace

\textbf{grupovati} / \textbf{pojmati}
& grupovati nasuprot razuměti ili konceptualizovati
& \profile{$\checkmark$}{\evstrong}{\evstrong}{\evstrong}{\evstrong}. Naslědžena \emph{shvat-} familija byla lažni prijatelj. Německo \emph{zusammenfassen/Zusammenfassung} prěmaje \emph{grupovati/grupovanje}; abstraktno \emph{fassen/auffassen} prěmaje \emph{pojmati}. Oficialno \emph{shvatiti} znači fizično grab/seize i ne jest tehničny most.\\[-0.1em]
&& \emph{EN:} The inherited \emph{shvat-} family was split by source sense: grouping is \emph{grupovati/grupovanje}; abstract conception is \emph{pojmati}. Official \emph{shvatiti} means physical grasping or seizing.\\
\addlinespace

\textbf{jednako} / \textbf{no}
& jednako: however/nevertheless; no: but
& \profile{$\checkmark$}{\evstrong}{\evstrong}{\evstrong}{\evstrong}. Upotrěbljene samo pri stvarnom protivnom odnosu. Naslědžene \emph{však, vendar} i homografično \emph{jednako} sut rěšene po klauznoj funkciji; ne-protivne upotrěby dostavajut eksplicitnu ekvivalenciju, način ili količinu.\\[-0.1em]
&& \emph{EN:} Reserved for genuine contrast. Czech/Slovak-shaped \emph{však}, Slovene \emph{vendar}, and ambiguous \emph{jednako} are classified by clause function rather than globally substituted.\\
\addlinespace

\textbf{dozvaljati}
& dozvoljavati, allow/permit
& \profile{$\checkmark$}{\evstrong}{\evstrong}{\partialev}{\evstrong}. Oficialna leksična površina odklanjaje poljsko oblikovano \emph{pozwalja}. Ona ostavaje odvojena od \emph{dopuščati}, ktoro može imati užšu formalnu funkciju ``admit/allow'' v matematičnoj konstrukciji.\\[-0.1em]
&& \emph{EN:} Official lexical surface replacing Polish-shaped \emph{pozwalja}. It remains distinct from formal \emph{dopuščati} where the source construction requires it.\\
\addlinespace

\textbf{jednočasno}
& jednočasno, simultaneously
& \profile{$\checkmark$}{\evstrong}{\evstrong}{\evstrong}{\evstrong}. Izbrano za stvarnu vremensku simultanost. Češsko-poljsko \emph{zároveň} může takože označavati diskursno ``at the same time/in addition''; zato kontekst, ne površina, upravljaje izborom.\\[-0.1em]
&& \emph{EN:} Used for genuine simultaneity. Czech/Polish-shaped \emph{zároveň} can have broader discourse functions, so context controls the choice.\\
\addlinespace

\textbf{prědtym} / \textbf{točno}
& prědtym: previously; točno: precisely/exactly
& \profile{$\checkmark$}{\evstrong}{\evstrong}{\evstrong}{\evstrong}. Površina \emph{právě} skryvala dvě funkcije. Německo \emph{eben behandelt/eben bewiesen} jest temporalno \emph{prědtym}; fokusno \emph{gerade} ili identifikacija jest \emph{točno}. Slepa zaměna by proměšala čas i fokus.\\[-0.1em]
&& \emph{EN:} The surface \emph{právě} hid two functions: temporal/anaphoric uses become \emph{prědtym}; focus-identifying or precise uses become \emph{točno}.\\
\addlinespace

\textbf{upotrěba} / \textbf{upotrěboju}
& upotrěba, use; instrumental ``by use of''
& \profile{$\checkmark$}{\evstrong}{\evstrong}{\evstrong}{\evstrong}. Oficialna imenica i korpusna instrumentalna paralela upravljajut oblikom \emph{upotrěboju}. Poljsko \emph{użyvanjem} ne jest ohranjeno kako samostojny morfologičny sloj.\\[-0.1em]
&& \emph{EN:} Official noun and corpus-attested instrumental yield \emph{upotrěboju}; Polish-shaped \emph{użyvanjem} is not retained as a separate layer.\\
\addlinespace

\textbf{sravnjenje} / \textbf{identifikovati}
& compare nasuprot identify
& \profile{$\checkmark$}{\evstrong}{\evstrong}{\evstrong}{\evstrong}. Stvarno matematično sravnjenje upotrěbjaje oficialnu \emph{sravn-} familiju i imperativ \emph{sravni}. Konteksty, gdje objekty sut proglašene identičnymi ili prirědžene, dostavajut \emph{identifikovati} ili \emph{prirěditi}; površna podobnost ne opravdava ``compare'' svuda.\\[-0.1em]
&& \emph{EN:} Genuine comparison uses the official \emph{sravn-} family. Identification and assignment contexts use \emph{identifikovati} or \emph{prirěditi}; one source-shaped family is not forced across all senses.\\
\addlinespace

\textbf{daljši}
& daljši, further/additional/subsequent
& \profile{$\checkmark$}{\evstrong}{\evstrong}{\evstrong}{\evstrong}. Oficialna komparativna familija jest sklanjana po rodu, čislu i padežu. Prislovno \emph{dalje} ostavaje samostojno; pridavnik ne može byti zaměnjen prislovom samo zato, že oba znače ``further''.\\[-0.1em]
&& \emph{EN:} Official inflected adjective for further/additional/subsequent. Adverb \emph{dalje} remains grammatically distinct.\\
\addlinespace

\textbf{poslědnji}
& poslědnji, last/final
& \profile{$\checkmark$}{\evstrong}{\evstrong}{\evstrong}{\evstrong}. Oficialna familija ohranjaje \emph{ě} i soglasovanje. Neakcentovano \emph{posledn-} jest pisemny sloj, ne drugi matematičny odnos.\\[-0.1em]
&& \emph{EN:} Official inflected family for last/final; unaccented \emph{posledn-} is a spelling layer, not a separate relation.\\
\addlinespace

\textbf{prěhoditi} / \textbf{prěhod}
& transition, pass to
& \profile{$\checkmark$}{\evstrong}{\evstrong}{\evstrong}{\evstrong}. Oficialny korenj jest \emph{prěhod-}; glagolna osoba ostavaje vidljiva, napriměr tretja osoba jedininy \emph{prěhodi}. Termin jest odvojen od \emph{poprědny} ``previous'' i \emph{prědtym} ``previously''.\\[-0.1em]
&& \emph{EN:} Official transition family with explicit person marking, distinct from the adjective ``previous'' and adverb ``previously.''\\
\addlinespace

\textbf{različny}
& različny, different/distinct
& \profile{$\checkmark$}{\evstrong}{\evstrong}{\evstrong}{\evstrong}. Zaglavna lemma jest stabilna; materijalny izbor jest soglasovanje: \emph{različnyh, različne, različnoj} po sintaktičnoj funkciji. Blizko \emph{razny} može imati opči ``various'' smysl i ne jest avtomatična zaměna.\\[-0.1em]
&& \emph{EN:} Stable lemma whose material issue is agreement. Related \emph{razny} can mean broader ``various'' and is not automatically substituted.\\
\addlinespace

\textbf{jih}
& njihovy, their
& \profile{$\checkmark$}{\evstrong}{\evstrong}{\evstrong}{\evstrong}. Oficialny preferovany tretje-osobny množinski posesiv jest nesklanjano \emph{jih}. Pridavničko \emph{njihov-} ne jest ohranjeno samo po nacionalnoj navykě.\\[-0.1em]
&& \emph{EN:} Official preferred indeclinable third-person plural possessive. Adjectival \emph{njihov-} is not retained merely from national habit.\\
\addlinespace

\textbf{vsi}: \textbf{vsih, vsim, vsimi}
& vsi/all with oblique cases
& \profile{$\checkmark$}{\evstrong}{\evstrong}{\evstrong}{\evstrong}. Čitateljnik upotrěbjaje jednu deklarovanu paradigmu. Zapadno oblikovane \emph{vseh, vsem, vsemi} ne označavajut drugi kvantifikator i zato ne ostavajut paralelny morfologičny sloj.\\[-0.1em]
&& \emph{EN:} One declared oblique paradigm for ``all.'' West-shaped variants do not encode a different quantifier and are not kept as a parallel layer.\\
\addlinespace

\textbf{kako}
& kako, as/how
& \profile{$\checkmark$}{\evstrong}{\evstrong}{\evstrong}{\evstrong}. Oficialna vezna i prislovna forma jest \emph{kako}. Samostojno \emph{kak} jest regionalno skračenje bez deklarovane smyslove funkcije v tom korpusu.\\[-0.1em]
&& \emph{EN:} Official conjunction/adverb. Bare \emph{kak} had no declared distinct function in the corpus.\\
\addlinespace

\end{longtable}
\normalsize

\section{Sistemne jezyčne konvencije / System-level language conventions}

Ne vsako rěšenje jest dobro predstavjeno kako zaglavno slovo. Slědujuče konvencije djelajut nad mnogimi terminami i zato sut opisane odvojeno.

\begin{longtable}{@{}P{0.22\textwidth}P{0.35\textwidth}P{0.35\textwidth}@{}}
\caption{Konvencije, ktore upravljajut slovnikom i korpusom.}\\
\toprule
\textbf{Oblast} & \textbf{Medžuslovjansko objasnjenje} & \textbf{English explanation}\\
\midrule
\endfirsthead
\toprule
\textbf{Oblast} & \textbf{Medžuslovjansko objasnjenje} & \textbf{English explanation}\\
\midrule
\endhead
\bottomrule
\endlastfoot
Osnovna ortografija
& Čitateljski tekst upotrěbjaje deklarovanu osnovnu površinu za opcionalne bukvy. Tuto ne jest generično odstranjenje diakritiky: \emph{ě} i drugi razlikujuči znaky ostavajut, a tuđe imena i citaty sut zaščitene.
& The reader uses a declared basic surface for optional letters. This is not generic diacritic stripping: contrastive signs such as \emph{ě} remain, and foreign names and quotations are protected.\\
\addlinespace
Morfologija
& Padež, rod, čislo, aspekt i glagolna osoba sut dio terminološkogo izbora. Genitiv množiny glagolnyh imen ima \emph{-anij/-enij}; tretja osoba množiny ima završno \emph{-t}; pridavničko soglasovanje slijedi stvarny zaglavnik.
& Case, gender, number, aspect, and person are part of the term choice. Verbal-noun genitive plurals take \emph{-anij/-enij}; third-person plurals take final \emph{-t}; adjective agreement follows the actual head noun.\\
\addlinespace
Polisemija
& Jedna iztočna riječ može zahtěvati nekoliko ciljnyh konstrukcij, a jedna podobna ciljna familija može skryvati nekoliko iztočnyh smyslov. Zato slovnik preferuje smyslove razděljenja prěd globalnym zaměnjenjem.
& One source word may require several target constructions, while one similar-looking target family may hide several source senses. Sense partition therefore precedes global replacement.\\
\addlinespace
Imena i identitety
& Polny avtorsky, citatny, naslovny, vydavateljski i bibliografičny identitet ostavaje v iztočnoj površini. Samo ime, ktoro stvarno prěma medžuslovjansku morfologiju v běžnom tekstu, může imati integrirany eponimny korenj.
& Complete author, citation, title, publisher, and bibliographic identities remain in their source surface. Only names functioning with Interslavic morphology in running prose may use an integrated eponymic stem.\\
\addlinespace
Latinica i kirilica
& Latinica jest jediny redagovateljny jezyčny kanon. Kirilica jest deterministična projekcija; ona ne jest nezavisny prěvod i ne može sama potvrditi termin.
& Latin is the sole editable language canon. Cyrillic is a deterministic projection, not an independent translation and not independent evidence for a term.\\
\addlinespace
Matematika i TeX
& Formuly, komandy, oznaky, odkazy, ključe citat i drugi ne-jezyčni objekty ne sut slovničny material. Jezyčna redakcija mora ohranjati njihovu strukturu i čitateljsku funkciju.
& Formulae, commands, labels, references, citation keys, and other nonlanguage objects are not dictionary material. Language editing must preserve their structure and reader function.\\
\end{longtable}

\section{Gramatične i pravopisne konvencije / Grammatical and orthographic conventions}
Leksičny izbor, gramatična funkcija i pravopis trěbujut različne razloge. Slědujuče zapisy objasnjajut, kako tute razlike upravljajut čitanjem terminov. Etimologične i historijne formy sut sravnenja, ne avtomatično preporučene formy glavnogo prěvoda.
\par\smallskip\noindent\emph{EN:} Lexical selection, grammatical function, and spelling require different reasons. These entries explain how those distinctions affect the reading of terms. Etymological and historical forms are comparanda, not automatically preferred forms for the translated reader.
\Needspace{8\baselineskip}
\subsection{Sklanjanje čislovnika dva / Declension of the numeral two}
\hypertarget{LING-01}{}\noindent\texttt{LING-01}\par\smallskip
Oficialna gramatika čislovnikov dava dva i žensku formu dvě, genitiv i lokativ dvoh, dativ dvom i instrumental dvoma. Izbor zavisi od funkcije v věte, ne samo od podobnosti napisa. Posebno, ranějša forma dvěma ne može iměti jednu zaměnu v vseh kontekstah: dativ trěbuje dvom, instrumental dvoma. Formy dvuh, dveh, dvěh i dvaju v razmotrenyh genitivnyh ili lokativnyh konstrukcijah ne imajut prednost nad deklarovanoju paradigmatoju. To jest gramatičny razlog, ne glasovanje jezykov. Pojavjenjska tabela zachovuje razliku medžu padežami; ona ne dava razrěšenje zaměnjati čisla, indeksi ili citaty.
\par\smallskip\noindent\emph{EN:} The official numeral grammar gives dva and feminine dvě, dvoh in the genitive and locative, dvom in the dative, and dvoma in the instrumental. The selection depends on the sentence, not spelling resemblance alone. In particular, the earlier form dvěma cannot receive one replacement everywhere: a dative context requires dvom, an instrumental context dvoma. The forms dvuh, dveh, dvěh, and dvaju in the examined genitive or locative constructions have no priority over the declared paradigm. This is a grammatical judgment, not a vote among languages. The occurrence table preserves the case distinction; it does not license altering numbers, indices, or quotations.
\par\nopagebreak[4]\smallskip
\noindent{\footnotesize\emph{Izvory / Sources:} \href{https://interslavic.fun/learn/grammar/numerals/}{Čislovniki / Numerals}, sklanjanje / declension.}\par\medskip
\Needspace{8\baselineskip}
\subsection{Leksična osnova i njezina pisemna forma / Lexical identity and its written form}
\hypertarget{LING-02}{}\noindent\texttt{LING-02}\par\smallskip
Oficialny slovnik podpira osnovy pomoč-, cělost-, cěnnost- i prislovnik imenno. Etimologične zaglavja pomoć, cělosť i cěnnosť trěba razlikovati od osnovnogo pravopisa togo izdanja: pomoč, cělost i cěnnost. Imenno znači točno imenovanje, ne dějanje imenovanja, ktore sugeruje imenujuče. Pri matematičnoj cěnnosti izbrana jest imennikova familija cěnnost-, ne automatno vsako slovo s podobnoju osnovoju: pridavnik vrědnostny ne byl rěšeny tym pravilom. Pět složenyh pridavnikov s cěločislovn- imajut zapis s cěločislov-; samostojno čislovno ne spada pod tu ograničenu odluku. Razlog za leksičnu osnovu, razlog za sklanjanje i razlog za pravopis sut različne pytanja; slovnično zaglavje samo ne dokazuje vsaku matematičnu upotrěbu.
\par\smallskip\noindent\emph{EN:} The dictionary supports the lexical families pomoč-, cělost-, cěnnost-, and the adverb imenno. Its etymological headwords pomoć, cělosť, and cěnnosť must be distinguished from this edition's basic spellings pomoč, cělost, and cěnnost. Imenno introduces an exact designation rather than an act of naming suggested by imenujuče. In mathematical value contexts, the noun family cěnnost- was selected; this did not settle every related word, including the adjective vrědnostny. Five compound-adjective forms with cěločislovn- were assigned cěločislov-, without extending that decision to unrelated bare čislovno. Lexical identity, inflection, and spelling therefore require separate reasons. A dictionary headword alone does not establish every technical use.
\par\nopagebreak[4]\smallskip
\noindent{\footnotesize\emph{Izvory / Sources:} \href{https://interslavic-dictionary.com/}{Medžuslovjansky slovnik / Interslavic dictionary}, zapisy / entries 3945, 100, 17486, 2281, 19667; \href{https://interslavic.fun/learn/orthography/}{pravopis / orthography}.}\par\medskip
\Needspace{8\baselineskip}
\subsection{Smysl, upravljanje i soglasovanje v věte / Meaning, government, and agreement in the sentence}
\hypertarget{LING-03}{}\noindent\texttt{LING-03}\par\smallskip
Německy argument upravljaje izborom cěloj konstrukcije. Uzimanje teorema v obzir vyražaje se s obzirom na ili vzame v obzir, ne slovesom, ktore označaje uređenje. K prostějšej zadače zachovuje dativ po k; pridavnik postavjenoj ne trěbuje proměnu samo zato, že susědny oblik jest nepravilny. Jedinym uslovjem jest instrumental, a taky stepen trěbuje pridavnik, ne prislovnik tak ili tako. Komponent jest mužskogo roda: primarny komponent soglasuje se s tym imennikom i s genitivom množiny komponentov. Nahodet i govoret slědujut paradigmě nastoječego vrěmena; etimologična bukva ę ne jest sama dokazom za slovesny konec. Tute odluke objasnjajut gramatičnu i smyslovu funkciju; one ne opravjajut formuly ili nepravilnosti v avtorskih matematičnyh znakah.
\par\smallskip\noindent\emph{EN:} The German argument determines the whole construction. Taking a theorem into account is expressed by s obzirom na or vzame v obzir, not by a verb meaning to arrange. K prostějšej zadače preserves the dative governed by k; the adjective postavjenoj does not require alteration merely because a neighbouring form is defective. Jedinym uslovjem is instrumental, while taky stepen requires an adjective, not the adverbs tak or tako. Komponent is masculine: primarny komponent agrees with that noun and with genitive plural komponentov. Nahodet and govoret follow the present-tense paradigm; the etymological letter ę is not by itself evidence for a verbal ending. These decisions concern grammatical function and meaning. They do not emend formulas or irregularities in the author's mathematical notation.
\par\nopagebreak[4]\smallskip
\noindent{\footnotesize\emph{Izvory / Sources:} Korpus Emmy Noether / Emmy Noether corpus, dělo / paper 45; \href{https://interslavic.fun/learn/grammar/verbs/}{slovesa / verbs}.}\par\medskip
\Needspace{8\baselineskip}
\subsection{Gramatičny rod termina upravljaje jego okolinom / A term's grammatical gender governs its surroundings}
\hypertarget{LING-04}{}\noindent\texttt{LING-04}\par\smallskip
Prsten i ideal sut mužske imenniki. Izbor termina zato imaje poslědstva za pridavniki, zaimenniki i prědikativne formy, ne samo za samo zaglavno slovo: grupovy prsten, izomorfny ideal i vsaky ideal trěbujut soglasovanje. Slična kontrola dava žensku formu linearna, ne linearnaja. Někoj- s običnoju paradigmoju jest osnova za někoje. Zaglavje prěkriženy trěbuje isto imenovanje v sodržanju i v samom oddělu, ako obě označajut isti pojem. Izomorfizm jest jedno slovo; pečatna granica reda v německom Iso-/morphismus ne stvara dve slova. Tute principy ne dokazyvajut, že vsako soglasovanje v korpusu jest pravilno. Konsultacija konkretnogo pojavjenja ostavaje nužna, osoblivo pri terminah, čij rod ne odgovarja rodu ranějšej kandidatnoj formy.
\par\smallskip\noindent\emph{EN:} Prsten and ideal are masculine nouns. Selecting the term consequently affects adjectives, pronouns, and predicative forms, not just the headword: grupovy prsten, izomorfny ideal, and vsaky ideal require agreement. The same kind of grammatical comparison supports feminine linearna rather than linearnaja and někoje from the ordinary někoj- paradigm. A heading using prěkriženy should name the same concept in the contents and the corresponding section. Izomorfizm is one word; the German printed line division Iso-/morphismus does not create two words. These principles do not establish that every agreement in the corpus is correct. Individual occurrences still require examination, especially when the selected noun has a different gender from an earlier candidate.
\par\nopagebreak[4]\smallskip
\noindent{\footnotesize\emph{Izvory / Sources:} \href{https://interslavic.fun/learn/grammar/}{Medžuslovjanska gramatika / Interslavic grammar}; \texttt{TEXT-COMP-01}, \texttt{CORPUS\_FORM\_COMPARISONS.json}.}\par\medskip
\Needspace{8\baselineskip}
\subsection{Razloženje i genitivne konce / Decomposition vocabulary and genitive endings}
\hypertarget{LING-05}{}\noindent\texttt{LING-05}\par\smallskip
Osnova razlož- veže slova o razloženju s oficialnym razložiti, zapisanym etimologično råzložiti. V familijah razlož- i nerazlož- bukva ž jest čestju leksičnoj formy, ne voljnym okrašeniem slova. Prosty z v razloz- ne jest zasnovany drugim smyslom v razmotrenyh pasažah. Genitiv mužskogo ili srědnjego roda ima konec -ogo v deklarovanoj paradigme: ktory, dodatny i svezany ne trěbujut paralelnu paradigmu na -oho. Rozhodujuči razlog jest slovnična osnova i gramatika; velika čestota istej formy v korpusu služi kako dodatno sravnjenje, ne kako zaměna za normativny dokaz.
\par\smallskip\noindent\emph{EN:} The stem razlož- connects the decomposition vocabulary with dictionary razložiti, written etymologically råzložiti. In the razlož- and nerazlož- families, ž belongs to the lexical form rather than being optional decoration. Plain z in razloz- has no distinct sense in the examined passages. Masculine or neuter genitives take -ogo under the declared paradigm; ktory, dodatny, and svezany do not require a parallel -oho paradigm. The grounds are lexical evidence and grammar. Frequent use of the same form in the corpus supplies a consistency comparison, not a substitute for that evidence.
\par\nopagebreak[4]\smallskip
\noindent{\footnotesize\emph{Izvory / Sources:} \href{https://interslavic-dictionary.com/}{Medžuslovjansky slovnik / Interslavic dictionary}, zapisy / entries 18608, 30231, 2643, 36196, 3512.}\par\medskip
\Needspace{8\baselineskip}
\subsection{Slovesna forma ne može byti opreděljena samo po koncu slova / An ending alone does not identify a verb form}
\hypertarget{LING-06}{}\noindent\texttt{LING-06}\par\smallskip
Oficialna gramatika dava -et za tretju osobu množiny drugoj konjugacije. Zato nahodet, govoret, tvoret i naležet trěba sravnjati s toj paradigmoju, ne izvoditi konec -ęt iz etimologičnogo pravopisa drugyh slov. Ale vzęt v izrazu jest vzęt označaje trpny pričestnik od vzeti, ne osobnu formu nastoječego vrěmena. V osnovnom pravopisě on piše se vzet. Razlikovanje jest morfologično: pričestnik ne može byti obražen kako slučaj paradigmy tretjej osoby. Pričestna i osobna forma mogut iměti podobny konec bez istej gramatičnoj funkcije. Prěpis bukvy po osnovnom pravopisě ne měnja tu funkciju.
\par\smallskip\noindent\emph{EN:} The official grammar gives -et for the third-person plural of the second conjugation. Nahodet, govoret, tvoret, and naležet must therefore be compared with that paradigm, rather than deriving an ending -ęt from etymological spellings elsewhere. But vzęt in jest vzęt is a past passive participle of vzeti, not a finite present-tense form. Its basic spelling is vzet. The distinction is morphological: a participle cannot be treated as an instance of the third-person paradigm. Participles and finite forms may have similar endings without sharing a grammatical function. Converting an optional letter to basic spelling does not change that function.
\par\nopagebreak[4]\smallskip
\noindent{\footnotesize\emph{Izvory / Sources:} \href{https://interslavic.fun/learn/grammar/verbs/}{Slovesa / Verbs}, nastoječe vrěme i pričestniki / present tense and participles.}\par\medskip
\Needspace{8\baselineskip}
\subsection{Čest i srodne formy v osnovnom pravopisě / Čest and related forms in basic spelling}
\hypertarget{LING-07}{}\noindent\texttt{LING-07}\par\smallskip
Pravilo osnovnogo pravopisa dava e vměsto opcionalnogo ę. Tym objasnjajut se čest, čestično i srodne formy; to ne jest izbor novogo korenja. Oficialny pravopis razděljuje etimologično objasnjenje od běžnogo pisanja, a slovnične priměry takože upotrěbjajut osnovne formy. Iz toga ne slijedi, že vsako podobno slovo jest pravilno: sčestnik trěbuje osobnu leksičnu analizu, a padež i tvorjenje učestjem sut odvojeno pytanje. Odstranjenje opcionalnogo znaka ne može služiti kako dokaz za lemma, smysl ili gramatičnu konstrukciju.
\par\smallskip\noindent\emph{EN:} Basic spelling uses e in place of optional ę. This accounts for čest, čestično, and related forms without selecting a new lexical root. The orthography distinguishes etymological explanation from ordinary writing, and dictionary examples also use basic forms. It does not follow that every similar word is correct: sčestnik requires a separate lexical analysis, while the inflection and formation of učestjem raise another question. Removing an optional mark cannot establish a lemma, a sense, or a grammatical construction.
\par\nopagebreak[4]\smallskip
\noindent{\footnotesize\emph{Izvory / Sources:} \href{https://interslavic.fun/learn/orthography/}{Pravopis / Orthography}, etimologičny alfabet / etymological alphabet.}\par\medskip
\Needspace{8\baselineskip}
\subsection{Sumand: tehničny termin bez samostojnogo slovničnogo zaglavja / Sumand: a technical term without its own dictionary headword}
\hypertarget{LING-08}{}\noindent\texttt{LING-08}\par\smallskip
Německy iztočnik govori o konečnom čislu Summanden. Sumand označaje čest sumy; po mnogo trěbuje se genitiv množiny sumandov. Izbor jednogo m opiraje se na sumandy i sumandov v sravnimyh pasažah istogo děla, ne na samostojno oficialno zaglavje. V konsultovanom slovniku jest suma, ale ne jest točno sumand ni summand. Srodnost s suma i doslědnost v tom samom matematičnom smyslu podpirajut projektno rěšenje; ne ustanavjajut občinskej normy. Dvojno m ne označaje drugy pojem, citat ili gramatičnu funkciju v razmotrenyh pojavjenjah. Slovnik zato prezentuje sumand kako obrazloženy tehničny prědlog s izrično omejenym dokazom.
\par\smallskip\noindent\emph{EN:} The German source refers to finitely many Summanden. Sumand denotes a term of a sum, and mnogo governs genitive plural sumandov. The single-m spelling rests on sumandy and sumandov in comparable passages of the same work, not on an independent official headword. The consulted dictionary contains suma but neither exact sumand nor summand. Its relation to suma and consistent use in the same mathematical sense support an editorial proposal; they do not establish community adoption. Double m marks no different concept, quotation, or grammatical function in the examined occurrences. Sumand is therefore presented as a reasoned technical proposal with explicitly limited evidence.
\par\nopagebreak[4]\smallskip
\noindent{\footnotesize\emph{Izvory / Sources:} Korpus Emmy Noether / Emmy Noether corpus; \href{https://interslavic-dictionary.com/}{Medžuslovjansky slovnik / Interslavic dictionary}, \emph{suma}.}\par\medskip
\Needspace{8\baselineskip}
\subsection{Slovnične osnovy, tvorjenje slov i kontekstualny izbor / Dictionary stems, word formation, and contextual selection}
\hypertarget{LING-09}{}\noindent\texttt{LING-09}\par\smallskip
Slovničny dokaz podpiraje někoj, několiko, imenno, izslědovati, poslě, těsno, trěbovati, upotrěbiti i vsegda. Pri někojih drugyh familijah trěba objasniti tvorjenje: matričny jest pridavnik od matrica, dělitelj veže se s děliti, a jediničny jest pridavnik, ne zaměna za imennik jedinica. Izčezati dava osnovu za izčez-, zatymčas dopušćati v etimologičnom zapisu odgovarja osnovnomu dopuščati. Pomoč- i cěl- isto trěbujut razlikovanje osnove i njezinoj formy v věte. V izrazu cěly red form izbran jest cěly, ne zato, že cělkovity jest nepravilen, ale radi točnogo lokalnogo smysla i doslědnosti s drugymi pojavjenjami. Korisno jest izjasnjaje koristnost razmatranja; uměstno ima užši smysl prikladnosti i ne dava identičny razlog. Predloženo cěljesȯobrazno ne imělo točnogo slovničnogo dokaza. Soglasovanje pri cělj trěbuje gramatičnu prověrku odvojenu od izbora korenja. Tute različne razloge ne trěba svesti na jednu nezasnovanu čislenu vagu.
\par\smallskip\noindent\emph{EN:} Dictionary evidence supports někoj, několiko, imenno, izslědovati, poslě, těsno, trěbovati, upotrěbiti, and vsegda. Other families also require a word-formation argument: matričny is an adjective from matrica, dělitelj relates to děliti, and jediničny is an adjective rather than a replacement for the noun jedinica. Izčezati supplies izčez-, while etymological dopušćati corresponds to basic dopuščati. The pomoč- and cěl- families likewise require lexical identity to be distinguished from the form used in a sentence. In cěly red form, cěly was preferred for its exact local sense and established usage, not because dictionary cělkovity is inherently wrong. Korisno jest states the usefulness of a line of thought; uměstno expresses appropriateness and is not identical in force. Proposed cěljesȯobrazno lacked an exact dictionary attestation. Agreement around cělj requires grammatical examination separately from root selection. These different reasons cannot responsibly be collapsed into an invented numerical weight.
\par\nopagebreak[4]\smallskip
\noindent{\footnotesize\emph{Izvory / Sources:} \href{https://interslavic-dictionary.com/}{Medžuslovjansky slovnik / Interslavic dictionary}, zapisy / entries 17486, 21308, 5696, 2396, 10975, 1942, 3894, 3092, 780, 35975, 1754, 5174, 4143, 30950.}\par\medskip
\Needspace{8\baselineskip}
\subsection{Etimologično zaglavje i osnovny čitatelsky pravopis / Etymological headwords and ordinary reading forms}
\hypertarget{LING-10}{}\noindent\texttt{LING-10}\par\smallskip
Oficialny slovnik može pokazovati etimologično zaglavje i osnovnu formu istogo slova. Pri dopuščati, skračati, sučstveny i vratnyh slovesah s vrač- bukva č v čitatelskom tekstu odgovarja etimologičnomu ć. Tuto ne jest rěšenje po prostoj čestotě: oficialny pravopis izrično dava ć → č, a slovničny prikaz potvrđuje formy. Vrač v smyslu lěčitelja ili čarovnika jest drugy imennik, ne dokaz za značenje vračati. Sravnjenje trěbuje istu pravopisnu vrstvu i isti smysl; inače možut dvě saglasne česti slovnika izgledati kako protivorěčje. Prednost osnovnoj formě ne izključaje etimologičnu formu iz naučnogo sravnjenja, kde njezina funkcija jest jasno objasnjena.
\par\smallskip\noindent\emph{EN:} The dictionary may supply an etymological headword and a basic reading form of the same word. In dopuščati, skračati, sučstveny, and the return verbs with vrač-, reader-facing č corresponds to etymological ć. This is not a decision by raw frequency: the official orthography explicitly prescribes ć → č, and the dictionary's displayed forms corroborate it. Vrač meaning a healer or sorcerer is a separate noun, not evidence for the sense of vračati. Comparisons must use the same orthographic level and the same sense; otherwise two compatible parts of a dictionary can appear to contradict each other. Preferring basic spelling in the reader does not exclude etymological forms from scholarly comparisons where their purpose is stated.
\par\nopagebreak[4]\smallskip
\noindent{\footnotesize\emph{Izvory / Sources:} \href{https://interslavic.fun/learn/orthography/}{Pravopis / Orthography}; \href{https://interslavic-dictionary.com/}{slovnik / dictionary}, zapisy / entries 7208, 20859, 3139, 953, 17039, 1839, 22658, 22646, 18360, 20192, 20193.}\par\medskip
\Needspace{8\baselineskip}
\subsection{Dvě osobne pravila: ć → č i đ → dž / Two special correspondences: ć → č and đ → dž}
\hypertarget{LING-11}{}\noindent\texttt{LING-11}\par\smallskip
Za osnovny latiničny pravopis oficialna norma dava dve osobne odgovaranja: ć piše se č, a đ piše se dž. To ne jest opče odstranjenje vseh diakritičnyh znakov. Napriměr, č ostavaje č. Pravilo opisujuče zapis ne izbira novo slovo i ne měnja jego matematičny smysl. Priměnjaje se na medžuslovjansku prozu; ne daje pravo proměniti slovo v točnom inojezyčnom citatu, ime v bibliografiji ili identifikator v formule. Takve razlike morajut byti naznačene prěd avtomatičnym prěpisom. V sravniteljnoj tablici etimologična forma može ostati neproměnjena, ibo čitatelj tam izučaje samo odgovaranje medžu pravopisami.
\par\smallskip\noindent\emph{EN:} For basic Latin spelling, the official orthography specifies two special correspondences: ć is written č, and đ is written dž. This is not the indiscriminate removal of diacritics; č, for example, remains č. A spelling rule does not select a new lexeme or alter its mathematical sense. It applies to Interslavic prose, not to a word in an exact foreign quotation, a bibliographical name, or an identifier in a formula. Those functions must be marked before automatic conversion. An etymological form can remain unchanged in a comparison table because the correspondence between the orthographies is itself the object being studied.
\par\nopagebreak[4]\smallskip
\noindent{\footnotesize\emph{Izvory / Sources:} \href{https://interslavic.fun/learn/orthography/}{Pravopis / Orthography}, etimologičny alfabet / etymological alphabet.}\par\medskip
\Needspace{8\baselineskip}
\subsection{Opcionalne bukvy i granicy pravopisnogo uproščenja / Optional letters and the limits of spelling simplification}
\hypertarget{LING-12}{}\noindent\texttt{LING-12}\par\smallskip
Osnovny latiničny pravopis izdanja ne upotrěbjaje opcionalne etimologične znaky v běžnoj medžuslovjanskoj proze. Odgovaranja sut å → a, ę → e, ų → u, ȯ → o, ė → e, ŕ → r, ĺ → l, ń → n, ť → t, ď → d, ś → s i ź → z. Dopustime alternativne zapisy è i ľ odgovarjajut e i l; samo inojazyčno è ne jest medžuslovjansky kandidat za proměnu. Č, š, ž i ě sut čestju osnovnogo alfabeta i ostavajut. Drugoj poredok pravopisnyh informaciji ne može služiti kako nova leksična ili gramatična odluka. Pravilo zato razděljuje čitatelsku prozu, originalne citaty i etimologične sravnjenja; vsaka od tih funkcij potrěbuje svoj zapis. Naučno objasnjenje može pokazati opcionalne bukvy bez vnucanja ih čitatelju glavnogo prěvoda.
\par\smallskip\noindent\emph{EN:} The edition's basic Latin spelling omits optional etymological marks in ordinary Interslavic prose. The correspondences are å → a, ę → e, ų → u, ȯ → o, ė → e, ŕ → r, ĺ → l, ń → n, ť → t, ď → d, ś → s, and ź → z. Permitted alternative spellings è and ľ correspond to e and l; a foreign è is not thereby an Interslavic candidate for conversion. Č, š, ž, and ě belong to the basic alphabet and remain. Presenting orthographic information differently cannot itself settle a lexical or grammatical question. The rule therefore distinguishes reader prose, original quotations, and etymological comparisons, each with its own purpose. A scholarly explanation can display optional letters without requiring them throughout the translated work.
\par\nopagebreak[4]\smallskip
\noindent{\footnotesize\emph{Izvory / Sources:} \href{https://interslavic.fun/learn/orthography/}{Pravopis / Orthography}, etimologičny alfabet i jego upotrěba / etymological alphabet and its usage.}\par\medskip

\section{Granice tvrdženja}

Slovnik može silno dokumentovati aktualnu izbranu površinu, iztočny smysl, zasvědočenu jezyčnu osnovu, realno odklonjene alternativy i unutrašnju doslědnost. On ne dokazuje, že vsaky govoritelj slovjanskogo jezyka bez trudnosti razuměje vsaku větu, niti že medžuslovjanska občina ili rodny govoritelji potvrdili vse izbory. Profili sut redakcijne sinteze dokazov, ne psiholingvistične mjery. Kde oficialny slovnik potvrđaje samo obyčny smysl, a matematičny smysl dolazi iz disciplinarnoj literatury, ta granica jest imenovana, kako pri \emph{prsten}.

\section*{Povtorljivost, AI-prinos, prava i citovanje}

Javny slovnik jest čitateljska sinteza redakcijnyh dokazov, ktore stvarno upravljajut korpusom. Njegov neposrědnje dostavny LaTeX jest samosodrživy, a iztočny ZIP sodrživaje toj samy iztočnik i instrukcije sestavjenja. Stabilny koncept korpusa jest \doiurl{10.5281/zenodo.21926382}; historijska arhivna versija korpusa jest \doiurl{10.5281/zenodo.23496017}; latiničny i kiriličny webovy čitatelj jest dostavny prěz \repo. Polny čitatelj avtorskogo matematičnogo korpusa ostavaje prvy dokument i PDF-prědgled; slovnik jest oddělny naučny pomoćnik.

Metodološko izloženje i sinteza slovnika sut izrabotane s \textbf{OpenAI Codex --- GPT-5.6 Sol, nivo usilja Ultra}; rekonstrukcija terminologičnyh zapisov i daljnje dvojezyčno obrazloženja s \textbf{OpenAI Codex --- GPT-6 Astra, nivo usilja Ultra}. Točna identičnost modela i nivo usilja pri naslědženom prěvodu korpusa ne sut pouzdano dokumentovane i ne sut retrospektivno pripisane. Tuto ne podrazuměvaje ljudsku lingvističnu recenziju.

Tekst slovnika a jeho tipografija sut izrabotane AI; matematični originaly, citovane formy, slovniki, gramatiky i sravnjiva literatura ohranjajut svoje avtorstvo i prava. Izdanje jest ponudženo bezplatno za nekomercialne, čisto obrazovateljne cěle na prědpoloženju izdatelja, že nositelji prav ne budu protiv. To jest prědpoloženje, ne ustanovjeno dozvoljenje: bezplatny dostop k originalu sam ne dozvaljaje publikaciju prěklada. Nositelj prav može \href{https://github.com/KokunoYumeto/emmy-noether-isv/issues}{podati prigovor}; izdatelj se obvezuje odstraniti dotičny material na jego zahtěv. Prava na originaly i zaščitjene izvedene prispevky ostavajut pri njihovih nositeljah. Historijske licencijne izjavy sut ohranjene; toj dokument ne davaje novogo dozvoljenja za děla tretjih stran ni za cělo izvedeno izdanje.

Prědloženo citovanje: \emph{AI-stvorjena studija medžuslovjanskoj terminologije. Matematičny medžuslovjansky v korpusu Emmy Noether: terminologičny slovnik i metoda izbora / Mathematical Interslavic in the Emmy Noether corpus: terminology dictionary and selection method.} 2026. Historijski DOI korpusa: \doiurl{10.5281/zenodo.21926382}.

\clearpage
\part{Method: English version}

\section{Purpose and object}

This dictionary describes the mathematical Interslavic used in the complete Noether corpus. Its unit of analysis is not automatically a word token. A lexical decision may concern a lemma, an inflectional family, a clause relation, a source-language polysemy, or the boundary between translatable prose and an identity that must remain exact. The Latin text is the editable language canon. The Cyrillic text is deterministically projected from that canon and therefore cannot serve as an independent translation witness.

The dictionary is designed for two readers at once. A Slavic-language reader can use the primary Interslavic account without English. A linguist or mathematician can use the English counterpart to inspect the distinction between source fidelity, cross-Slavic evidence, and editorial judgment. The alphabetical table is intentionally compact; expanded discussion is reserved for decisions whose structure cannot be understood from a gloss alone.

\section{Research questions and the limited claim}

Interslavic is a planned zonal language intended for practical comprehension among speakers of Slavic languages, not merely as an exercise in language construction. Previous scholarship has discussed it as an instrument of receptive multilingualism and cross-border access in the Slavic region \cite{merunka2019}. The present study asks whether a broad historical mathematical corpus can serve as a controlled site for developing its technical register; whether a language model can compare several kinds of evidence without itself becoming the authority; whether every consequential choice can be documented in a form that a non-technical language specialist can review; and whether the method can transfer to another domain without training a new model.

The practical hypothesis is that a language model can assist candidate generation, corpus comparison, and consistent terminology use when source meaning and linguistic evidence remain the basis of assessment. This study documents such a method; it does not measure time saved against a human-translation baseline. The Noether corpus does not establish the optimality of every form, comprehension by every Slavic reader, or successful transfer to another discipline.

\subsection{The reader's starting questions}

A linguist needs to know which proposals might be useful to Interslavic, what supports them, and where the evidence is insufficient. A mathematician needs to know whether the same objects, conditions, and arguments survive translation. A researcher wishing to repeat the method needs to know which sources to collect, what to give the model, and how to recognize a wrong answer. The dictionary and methodological account are organized around these questions. Reading the explanations requires no programming knowledge.

Each entry must distinguish four things: the mathematical concept, the selected expression, the evidence for it, and the limits of the claim. An early glossary form may be a proposal rather than a recommendation. Frequent occurrence does not establish the correctness of a sense. Where a reason has not survived, the entry says so; a new explanation cannot be passed off as historical testimony.

\section{How to use the materials}

The translated author-work reader and the account of its making are separate objects. The reader remains the first document and PDF preview. The present publication contains:

\begin{itemize}
  \item the current canonical dictionary governing the frozen reader;
  \item all 2,574 recovered original terminology records, aggregated without loss into 2,016 distinct source--Interslavic pairs;
  \item 168 public lexical, semantic, grammatical, orthographic, source-fidelity, notation, and identity decisions from the main normalization lineage;
  \item the complete role-policy account for 4,560 reviewed projection records, including 364 corrected classifications and eight added identity-protection spans;
  \item all twelve final normalization decisions and all 4,053 occurrence-level changes;
  \item machine-readable JSON and CSV crosswalks, direct LaTeX, and a reproducible source package.
\end{itemize}

``Canon'' here means the internally consistent editorial basis of this edition. It does not mean an official decision by the Interslavic community. The stable identifiers and occurrence maps are designed precisely so that specialist review can revise a rule without reconstructing the entire production history.

\section{What the AI did---and did not do}

A generative language model is software that produces text from statistical patterns learned from large collections of text. It can compare formulations, propose grammatical paraphrases, and follow elaborate instructions. It is not a dictionary, a language academy, a mathematical source, or a conscious witness to the provenance of every form it proposes. It can confidently invent an unattested word, blend national standards, or alter a mathematical relation, especially in a low-resource language. Research on low-resource translation likewise shows that supplying dictionaries, grammar descriptions, and examples does not by itself solve retrieval and understanding failures \cite{court2024}; terminology constraints help, but mechanical enforcement can reduce fluency \cite{bane2023}.

The model therefore has three legitimate roles here: it proposes candidates, compares them against supplied evidence, and manages consistency across occurrences. It has no fourth role as a language legislator. A decision must expose its evidence, scope, rejected alternatives, and uncertainty. This accords with concept-oriented terminology work, where object, concept, definition, and designation are distinguished before a preferred term is chosen \cite{iso704}. Polysemous source forms receive separate sense records rather than one global replacement.

The method uses an existing language model without training a new system. The researcher supplies the text, relevant passages from grammars and dictionaries, subject literature, and reasoned terminology decisions. The practical requirement is to prepare suitable sources and a well-defined assignment, not to construct a new model. Calling this ``zero-shot'' would obscure the role of those sources, examples, and prior decisions. Repeatability means that another researcher can repeat the comparisons and tests; it does not mean that every model must propose identical wording.

\section{Reconstruction of the production method}

The surviving evidence distinguishes several kinds of work. The source was divided into identified units accompanied by 215 glossaries containing 2,574 records: a German term or concept, Russian and Ukrainian comparanda, an Interslavic candidate, a rationale, and an uncertainty annotation where present. These files constitute the original terminology base. The presence of a comparison word in a glossary does not by itself establish consultation of a particular scholarly source.

The translation trials compared an English intermediate text, separation of prose from TeX and mathematics, and direct translation from German. Preserving every mathematical identifier is necessary but insufficient: the right number of identifiers does not establish correct prose or logical relations. Thirty-one units of the lecture-book text were translated into Russian, Ukrainian, and Interslavic with model assistance. The surviving evidence does not reliably establish the exact model identity or effort level for that translation, nor does it establish human editing.

The linguistic analysis distinguishes vocabulary, sense, grammar, orthography, and fidelity to the source. A separate question concerns the boundary between prose that can be transliterated and names, quotations, Roman numerals, and mathematical identifiers whose identities must be preserved. These are not merely different spellings of words: the function of a sign in the argument determines whether a transliteration rule applies. The occurrence tables make that distinction inspectable in particular passages. Where only later evidence survives, the account is explicitly reconstructive; a later explanation is not presented as contemporaneous testimony.

\section{A repeatable workflow for another domain}

\begin{enumerate}
  \item \textbf{Define the mandate.} State audience, register, script, domain, rights, source edition, and claim boundary.
  \item \textbf{Identify and preserve the source edition.} Retain the whole text, numbering, formulae, references, and bibliography. Keep a translation of the author's text separate from any editorially corrected edition.
  \item \textbf{Read a language and disciplinary canon.} Consult the target grammar and dictionary and relevant subject literature in representative Slavic languages; record passages actually used.
  \item \textbf{Number passages and identify material that must remain exact.} Divide the text into meaningful sections. Mark formulas, typesetting instructions, identifiers, quotations, and names whose identity must be preserved.
  \item \textbf{Request a translation of one identified passage.} Supply the passage, relevant evidence, and an explicit description of the required answer. If formulas are separated from prose, restore them exactly and check their order and relation to the translated sentences.
  \item \textbf{Create concept records.} For every non-trivial choice record source concept, target form, grammatical role, branch comparanda, official evidence, disciplinary use, and context. Split polysemy by sense.
  \item \textbf{Compare the candidates.} First exclude expressions that change the source's mathematical meaning. Compare the remaining candidates for support in Interslavic resources, recognition across Slavic branches, technical usage, consistency, and ambiguity. Report historical numerical weights only where they were actually used.
  \item \textbf{Record the decision.} Preserve the selected form, genuine alternatives, reasons for rejection, uncertainty, and scope.
  \item \textbf{Map every occurrence.} Check each local sense and exclude homographs, citations, names, and mathematical symbols from a global replacement.
  \item \textbf{Verify language and structure.} Compare source and target claims, quantifier scope, signs, references, morphology, aspect, and terminology. Compilation proves only technical buildability.
  \item \textbf{Maintain one editable text.} Derive the other script by stated transliteration rules, preserving names, quotations, and mathematical identifiers where required. Inspect the resulting pages and provide the reader, editable source, dictionary, and supporting data together.
  \item \textbf{Make particular choices open to criticism.} Number the decisions and identify their senses, evidence, and occurrences. A specialist should be able to challenge one choice without having to reconstruct the whole project.
\end{enumerate}

Semi-automatic terminology systems commonly separate corpus processing, term extraction, database population, and database management \cite{nazar2023}. This workflow adds two requirements specific to a planned interlanguage: balanced comparison among language branches and an auditable boundary between a model proposal and the accepted editorial canon.

\section{How a specialist can review the result}

\subsection{Applying the method to one passage}

A researcher can begin with one complete mathematical argument, such as a definition and its first example. Identify the source passage by author, work, section, and page or numbered paragraph. Read enough surrounding text to establish the notation and assumptions. The passage is a working unit, not a substitute for the full corpus to be translated.

Prepare three short components: the passage itself, relevant language and subject sources, and previously justified terminology. A bibliography without the pertinent passages is insufficient. For a new mathematical concept, distinguish attestation of an ordinary word from attestation of its technical sense. For a word with several meanings, first state which relation the sentence expresses. The model may propose a form and a reason, but a quotation from a book must be found in the book, not accepted on the strength of the model's assertion.

The answer for one passage has three readable components: the translation, its terminology and grammatical decisions, and unresolved questions. A compact decision record is sufficient if it gives the concept, selected form, relevant source, genuine alternatives, and reason. When no defensible alternative has been found, say so; do not invent a second candidate to fill a comparison table. When a source supports only a broader or different sense, state that limit.

Checking the translation is not the same as asking again whether it is good. Compare particular claims: the subject and referent, the objects quantified over, the assumptions and conclusions, and the scope of conditions and negation. Compare candidate forms within whole sentences, including gender, case, and government. The \emph{prsten} example shows why replacing a noun without checking agreement around it is insufficient.

After selecting a rule, locate other occurrences of the concept. Searching for a spelling helps, but the same spelling may have another sense. The occurrence table is therefore a concordance: a list of passages in which the rule's applicability can be examined. It is not proof of correctness by itself. A CSV table opens in an ordinary spreadsheet program; JSON represents relationships for automatic processing. Reading JSON is not necessary to understand a linguistic decision.

\subsection{What the reconstruction can establish}

\begin{description}
  \item[Contemporaneously recorded reason.] A surviving record states why a candidate was selected. This establishes the documented reason, not automatically its correctness.
  \item[Later examination.] A form is compared with a source not documented at the original selection. That comparison may support or challenge the form, but does not establish earlier consultation.
  \item[Procedure for repetition.] The method specifies how decisions should be documented and checked in a new corpus. This is not a claim that every historical decision followed an ideal procedure.
\end{description}

These distinctions let a specialist examine both a useful proposal and weaknesses in its evidence. The absence of a numerical weight is not the absence of a reason: many choices turn on a difference of meaning or a grammatical rule, neither of which requires a vote. Where numbers are used, however, their units, complete contributions, and arithmetic must be available. The reconstruction concerns documented decisions, not inaccessible internal thoughts of the model.

\subsection{Questions for scholarly assessment}

A specialist need not review ``the AI'' as a single object. Each entry supports six ordinary scholarly questions: Was the source concept identified correctly? Is the target form supported by the declared grammar and orthography? Are the compared Slavic traditions representative? Is the technical sense attested rather than inferred from an ordinary-language cognate? Are alternatives and ambiguity stated honestly? Were all and only the relevant occurrences found?

This design follows evidence that professional evaluators using full document context and explicit error categories produce more informative judgments than aggregate crowd scores \cite{freitag2021}. Automatic metrics also remain unreliable for distinguishing all small but meaningful differences among strong translations \cite{agrawal2024}. The publication therefore does not replace expert judgment with a score. It supplies a bounded claim, its evidence, and a revisable occurrence set.

\section{Decision hierarchy}

The method has five layers, but they are not five votes of equal status.

\begin{enumerate}
  \item \textbf{Preserving mathematical meaning is indispensable.} A candidate must preserve the object referred to, the structure of the argument, the scope of quantifiers, and the relations expressed in the proof. Widespread usage or an attractive etymology cannot justify a term that changes the claim.
  \item \textbf{Official Interslavic evidence anchors the language.} The declared dictionary and grammar provide lemmas, paradigms, aspect, government, and connective structure. An ordinary-language headword does not by itself establish a specialized mathematical sense.
  \item \textbf{East, West, and South Slavic comparison is diagnostic.} The branches are considered in balance to identify recognizability, concentration, and false friends. This is not a raw population vote, and the number of available books in one language is not allowed to dominate the result.
  \item \textbf{Disciplinary and internal corpus evidence resolves technical use.} Mathematical usage can establish a technical extension that an ordinary dictionary does not record. Existing corpus frequency matters only after inherited source bias and error have been considered.
  \item \textbf{Readability and ambiguity break ties among admissible candidates.} A form that can be parsed without English and does not silently merge two mathematical relations is preferred.
\end{enumerate}

Different choices require different kinds of examination. A spelling convention may apply to a family of forms, whereas a word with several meanings requires attention to each sentence. A citation or identifying label must remain recognizable; a mathematical sign is not a word to translate. The surviving records do not establish that every earlier choice received equally thorough examination. Alternatives are reported where they were genuinely considered, not invented to make the selected form look inevitable.

\section{Weighting without false precision}

The evidence profile in the dictionary uses five symbols: source fidelity (S), official Interslavic support (O), branch comparison (V), corpus or disciplinary fit (K), and readability (R). A filled mark means direct strong support, an open mark limited or indirect support, and a dash that the decision did not rely on that layer. The marks are an explanatory synopsis, not calibrated probabilities and not a hidden numerical model.

Numerical weights are shown only where the historical decision actually used them. The algebraic-ring decision divided evidence into six equally weighted Slavic groups. The \emph{prsten/pierścień} root received 2.5 units, \emph{kolco} 2.0, \emph{okruh} 1.0, and \emph{kolobar} 0.5. Official evidence established broad recognition of the physical-ring root, while mathematical books established its technical use. Because East Slavic uses a different root, the first corpus occurrence supplies Russian \emph{кольцо} and Ukrainian \emph{кільце} as a reader bridge. No comparable numerical history is fabricated for terms whose decisions were qualitative.

The complete arithmetic is shown below. The six-group structure follows the language project's design criteria \cite{designcriteria}; the contributions reconstruct the recorded totals from the documented root assignments.

\begin{center}\small
\begin{tabularx}{\textwidth}{@{}Yrrrr@{}}
\toprule
Language group & prsten & kolco & okruh & kolobar\\
\midrule
Russian & 0 & 1 & 0 & 0\\
Ukrainian and Belarusian & 0 & 1 & 0 & 0\\
Polish & 1 & 0 & 0 & 0\\
Czech and Slovak & 0 & 0 & 1 & 0\\
Slovene and Serbian/Croatian/Bosnian & 0.5 & 0 & 0 & 0.5\\
Bulgarian and Macedonian & 1 & 0 & 0 & 0\\
\midrule
Total & 2.5 & 2 & 1 & 0.5\\
\bottomrule
\end{tabularx}
\end{center}

Prsten has the largest share, not an absolute majority of the six units. More books or occurrences confer no extra vote; grouping it with Polish pierścień is a linguistic judgment about cognacy. The reported corpus observations contain no separate Belarusian or Bosnian counts. This is consequently not an equally sampled empirical study of every named language, still less a comprehension experiment. The complete numerical record is supplied in \texttt{RING\_COMPARISON\_WEIGHTS.json}.

\section{Why semantic partitions matter}

Three families illustrate the method particularly well.

First, German \emph{gilt} cannot be represented responsibly by one borrowed national predicate. Czech and Slovak use \emph{platí}; Croatian uses \emph{vrijedi}; Polish uses \emph{zachodzi} or \emph{ma miejsce}; Slovene uses \emph{velja}; and Bulgarian frequently uses fulfilment constructions. The edition therefore selects among \emph{imaje město}, \emph{jest pravdivo}, \emph{izpolnjaje se}, \emph{može se priměniti}, and \emph{ostavaje v silě} according to the source relation.

Second, discourse connectives are selected by clause function. Causal \emph{zato že}, inferential \emph{zato}, equivalential \emph{to jest}, naming \emph{imenno}, conditional \emph{ako\ldots{}to}, additive \emph{takože}, and result/purpose \emph{tako že} are not stylistic variants. Substituting one for another can reverse or blur the proof structure.

Third, apparent cognates are tested against sense. Mathematical containment is \emph{sodrživati}; content is \emph{sodržanje}; hold back is \emph{zadržati}. Existence is \emph{obstajati}; essentiality is \emph{sučstveny}; composition is \emph{sostojati se iz}. Imperfective assumption is \emph{prědpokladati}; perfective, passive, adjectival, and nominal supposition uses belong to the \emph{prědpolož-} family. Similar spelling is evidence to investigate, not a conclusion.

\section{Reproducibility and limitations}

The public dictionary is a human-readable synthesis of the exact editorial evidence used for the corpus. Its directly downloadable LaTeX is self-contained, and the source archive contains the same source plus build instructions. The stable corpus concept DOI is \doiurl{10.5281/zenodo.21926382}; a historical archived version of the corpus is \doiurl{10.5281/zenodo.23496017}; the live Latin/Cyrillic reader is available through \repo. The complete author-work reader remains the first file and PDF preview; this dictionary is a separate scholarly aid.

The dictionary does not claim native-speaker validation, community endorsement, or independent peer review. A historical proposal, a form attested in the translation, and a linguistically justified recommendation establish different things. For example, the presence of masculine prsten does not certify the agreement of neighbouring words; the coexistence of zamknjen- and zatvorjen- in passages on algebraic closedness does not establish a single preferred form. Such differences require contextual examination, not an inference from word counts. A scholarly objection should identify the affected sense and supply contrary or stronger evidence, while distinguishing a revised language choice from an unchanged mathematical source.

\section*{AI-work disclosure, rights, and citation}

The mathematical authors are Emmy Noether and the named co-authors: R. Brauer in Paper 32; R. Brauer and H. Hasse in Paper 38. M. Deuring prepared the account of Noether's winter 1929/30 lectures in work 44. Paper 45 is by H. Kapferer, with an appendix jointly written with E. Noether. Their mathematical authorship is distinct from the AI translation and the present linguistic study.
\par\emph{Medžuslovjansky:} Matematični avtori sut Emmy Noether i navedeni soavtori: R. Brauer v statji 32; R. Brauer i H. Hasse v statji 38. M. Deuring obrabotal lekcije Noether iz zimnogo semestra 1929/30 v dělu 44. Statja 45 jest od H. Kapferera, s dodatkom napisanym sovmestno s E. Noether. Njihovo matematično avtorstvo jest različno od AI-prěvoda i tutoj jezyčnoj studije.

The methodological exposition and dictionary synthesis were produced with \textbf{OpenAI Codex --- GPT-5.6 Sol, Ultra effort}; reconstruction of terminology records and further bilingual explanations were produced with \textbf{OpenAI Codex --- GPT-6 Astra, Ultra effort}. The exact model identity and effort level for the inherited corpus translation are not reliably documented and are not retrospectively attributed. The corpus provides its editable source and scholarly source notices. No human linguistic review is implied.

The dictionary text and typesetting are AI-created. The original mathematics, quoted forms, dictionaries, grammars, and comparative literature retain their human authorship and respective rights. The edition is offered free of charge for non-commercial, purely educational purposes, on the publisher's presumption that rights holders would not object. This is a presumption, not established permission: free access to an original does not itself authorize derivative publication. Rights holders may \href{https://github.com/KokunoYumeto/emmy-noether-isv/issues}{raise an objection}; the publisher undertakes to take down the affected material upon a rights holder's request. Rights in originals and protected derivative contributions remain with their respective holders. Historical licence statements are preserved; this document makes no new licence grant for third-party works or the derivative edition as a whole.

Suggested citation: \emph{AI-created Interslavic terminology study. Matematičny medžuslovjansky v korpusu Emmy Noether: terminologičny slovnik i metoda izbora / Mathematical Interslavic in the Emmy Noether corpus: terminology dictionary and selection method.} 2026. Historical corpus DOI: \doiurl{10.5281/zenodo.21926382}.

\section*{Primary public sources}

{\small\noindent
\textbf{Corpus:} \doiurl{10.5281/zenodo.21926382}.\quad
\textbf{Reader:} \repo.\quad
\textbf{Grammar:} \href{https://interslavic.fun/learn/grammar/}{Interslavic grammar and learning materials}.\quad
\textbf{Dictionary:} \href{https://interslavic-dictionary.com/}{Interslavic dictionary project}.\par}

\clearpage\part{Polny slovnik i obrazloženja / Complete dictionary and explanations}
\clearpage\section{Alfabetičny indeks / Alphabetical index}
Indeks uključaje i historijne prědlogy. Odkaz vodi k obrazloženju, ne sam po sobě k rekomendaciji. / The index includes historical proposals. A link leads to the explanation, not by itself to a recommendation.
\small\begin{longtable}{@{}P{0.66\textwidth}P{0.28\textwidth}@{}}
\toprule Forma ili pytanje / Form or question & Zapis / Entry\\\midrule\endhead
\bottomrule\endfoot
$\lambda$ jest proizvoljno & \hyperlink{TERM-0002}{TERM-0002}\\
$\Omega$-proces & \hyperlink{TERM-1278}{TERM-1278}\\
$n$-e korenje iz jedinice & \hyperlink{TERM-0003}{TERM-0003}\\
(p)\textasciicircum{}alpha abo jediničny ideal & \hyperlink{TERM-0005}{TERM-0005}\\
-1 kako suma dvoh kvadratov & \hyperlink{TERM-0007}{TERM-0007}\\
-ajut i -aju: kto jest subjekt? & \hyperlink{DEC-0059}{DEC-0059}\\
-yh v pridavnikah, zaimenah i participah & \hyperlink{DEC-0096}{DEC-0096}\\
11. poręd / 12. poręd / 13. poręd / 15. poręd & \hyperlink{TERM-0008}{TERM-0008}\\
117 form & \hyperlink{TERM-0009}{TERM-0009}\\
30. marca 1926. & \hyperlink{TERM-0010}{TERM-0010}\\
5./6./... poręd & \hyperlink{TERM-0011}{TERM-0011}\\
\textbackslash{}D & \hyperlink{TERM-0012}{TERM-0012}\\
\textbackslash{}mathfrak h / \textbackslash{}mathfrak t & \hyperlink{TERM-0600}{TERM-0600}\\
\textbackslash{}mathrm\{A)\}, \textbackslash{}mathrm\{B)\}, \textbackslash{}mathrm\{C)\} & \hyperlink{TERM-0021}{TERM-0021}\\
\textbackslash{}mathrm\{a)\}, \textbackslash{}mathrm\{b)\}, \textbackslash{}mathrm\{c)\} & \hyperlink{TERM-0599}{TERM-0599}\\
\textbackslash{}pmod\{(\textbackslash{}sigma,s)\} & \hyperlink{TERM-0014}{TERM-0014}\\
\textbackslash{}pmod\{(\textbackslash{}varrho,t)\} & \hyperlink{TERM-0015}{TERM-0015}\\
a jest děljiva črez t / ideal a, kratny od t & \hyperlink{TERM-0016}{TERM-0016}\\
a priori & \hyperlink{TERM-0017}{TERM-0017}\\
A) / B) / C) / D) & \hyperlink{TERM-0018}{TERM-0018}\\
A) / B) / C) / D), a) / b) & \hyperlink{TERM-0019}{TERM-0019}\\
A) / B) / C), a) / b) / c) & \hyperlink{TERM-0020}{TERM-0020}\\
Abelove grupy/polja & \hyperlink{TERM-0028}{TERM-0028}\\
absolutna nerazložimost / absolutno nerazložimy & \hyperlink{TERM-0042}{TERM-0042}\\
absolutna oblast množiteljev & \hyperlink{TERM-0044}{TERM-0044}\\
absolutna prosta funkcija / absolutno nerazložimy polinom & \hyperlink{TERM-0043}{TERM-0043}\\
absolutno algebraično polje & \hyperlink{TERM-0037}{TERM-0037}\\
absolutno ireducibilno predstavjenje & \hyperlink{TERM-0041}{TERM-0041}\\
absolutno ireducibilny & \hyperlink{TERM-0040}{TERM-0040}\\
absolutno nerazložime sostavne česti & \hyperlink{TERM-0038}{TERM-0038}\\
absolutno polny sistem & \hyperlink{TERM-0039}{TERM-0039}\\
absolutny prost ideal & \hyperlink{TERM-0045}{TERM-0045}\\
absolutny prosty ideal & \hyperlink{TERM-0046}{TERM-0046}\\
abstraktna izgradnja teorije idealov & \hyperlink{TERM-0048}{TERM-0048}\\
abstraktna konstrukcija teorije idealov & \hyperlink{TERM-0049}{TERM-0049}\\
abstraktna teorija idealov & \hyperlink{TERM-0047}{TERM-0047}\\
adjungovane formy/funkcije/krive & \hyperlink{TERM-0058}{TERM-0058}\\
adjungovano kako modul & \hyperlink{TERM-0128}{TERM-0128}\\
adjungovany modul & \hyperlink{TERM-0054}{TERM-0054}\\
adjungovati funkciju log x k oblasti racionalnosti & \hyperlink{TERM-0057}{TERM-0057}\\
adjungujemo & \hyperlink{TERM-0055}{TERM-0055}\\
adjungujut se kako moduly & \hyperlink{TERM-0129}{TERM-0129}\\
adjunkcija & \hyperlink{TERM-0059}{TERM-0059}\\
adjunkcija / adjungovati & \hyperlink{TERM-0060}{TERM-0060}\\
adjunkcija sqrt(x) / modulna baza & \hyperlink{TERM-0064}{TERM-0064}\\
afekt / bez afekta & \hyperlink{TERM-0065}{TERM-0065}\\
afinna grupa & \hyperlink{TERM-0067}{TERM-0067}\\
Ako … to: uslovje i posledica & \hyperlink{DEC-0072}{DEC-0072}\\
algebrai?na baza & \hyperlink{TERM-0104}{TERM-0104}\\
algebrai?no-drobno ?islo & \hyperlink{TERM-0102}{TERM-0102}\\
algebraična baza & \hyperlink{TERM-0105}{TERM-0105}\\
algebraična baza / racionalna baza & \hyperlink{TERM-0106}{TERM-0106}\\
algebraična baza H iz Theta & \hyperlink{TERM-0107}{TERM-0107}\\
algebraične čislove polja & \hyperlink{TERM-0111}{TERM-0111}\\
algebraične elementy / cěle elementy & \hyperlink{TERM-0109}{TERM-0109}\\
algebraične i diferencialne invarianty & \hyperlink{TERM-0110}{TERM-0110}\\
algebraično cěla i celočislena vzhodno k & \hyperlink{TERM-0093}{TERM-0093}\\
algebraično cěla vzhodno k G & \hyperlink{TERM-0092}{TERM-0092}\\
algebraično cěle čisla iz Omega / cěle čisla & \hyperlink{TERM-0099}{TERM-0099}\\
algebraično cělo čislo & \hyperlink{TERM-0098}{TERM-0098}\\
algebraično cěly nad R & \hyperlink{TERM-0078}{TERM-0078}\\
algebraično čislovo polje & \hyperlink{TERM-0117}{TERM-0117}\\
algebraično nezavisne & \hyperlink{TERM-0080}{TERM-0080}\\
algebraično nezavisne / transcendentne & \hyperlink{TERM-0082}{TERM-0082}\\
algebraično nezavisne / zavisne & \hyperlink{TERM-0081}{TERM-0081}\\
algebraično nezavisne funkcije & \hyperlink{TERM-0083}{TERM-0083}\\
algebraično obrazovanje & \hyperlink{TERM-0068}{TERM-0068}\\
algebraično razširenje & \hyperlink{TERM-0108}{TERM-0108}\\
algebraično razširjenje polja / podpolje & \hyperlink{TERM-0112}{TERM-0112}\\
algebraično zamknjeno podkolco P & \hyperlink{TERM-0086}{TERM-0086}\\
algebraično zamknjeno polje & \hyperlink{TERM-0069}{TERM-0069}, \hyperlink{TERM-0070}{TERM-0070}\\
algebraično zatvorjeno polje & \hyperlink{TERM-0085}{TERM-0085}\\
algebraično zatvorjeno polje / algebraično razširiteljno polje & \hyperlink{TERM-0084}{TERM-0084}\\
algebraično zatvorjeno polje izvedeno iz P & \hyperlink{TERM-0254}{TERM-0254}\\
algebraično zavisimy / nezavisimy & \hyperlink{TERM-0071}{TERM-0071}\\
algebraično zavisne funkcije & \hyperlink{TERM-0073}{TERM-0073}\\
algebraično zavisny / algebraično nezavisny & \hyperlink{TERM-0087}{TERM-0087}\\
algebraično zavisny / nezavisny & \hyperlink{TERM-0072}{TERM-0072}\\
algebraično-cěla adjunkcija & \hyperlink{TERM-0097}{TERM-0097}\\
algebraično-cěla zatvorenost & \hyperlink{TERM-0096}{TERM-0096}\\
algebraično-cěle / algebraično-drobne & \hyperlink{TERM-0088}{TERM-0088}\\
algebraično-cěle substitucije / cěle algebraične čisla & \hyperlink{TERM-0101}{TERM-0101}\\
algebraično-cělo zatvorjena & \hyperlink{TERM-0090}{TERM-0090}\\
algebraično-cělo zavisi od [H,xi\_i] & \hyperlink{TERM-0091}{TERM-0091}\\
algebraično-celočislene polinomy & \hyperlink{TERM-0100}{TERM-0100}\\
algebraično-cěly / racionalno-cěly & \hyperlink{TERM-0089}{TERM-0089}\\
algebraično-drobno čislo & \hyperlink{TERM-0079}{TERM-0079}, \hyperlink{TERM-0103}{TERM-0103}\\
algebraičny rang & \hyperlink{TERM-0113}{TERM-0113}\\
algebraičny rang (stepen transcendencije) & \hyperlink{TERM-0114}{TERM-0114}\\
algebraičny rang / polje ostatkovyh klasov & \hyperlink{TERM-0115}{TERM-0115}\\
algebraičny rang ostavaje sohranjen & \hyperlink{TERM-0116}{TERM-0116}\\
antistrofna matrica & \hyperlink{TERM-0139}{TERM-0139}\\
Aritmetična formulacija pojma dimenzije & \hyperlink{TERM-0142}{TERM-0142}\\
aritmetična teorija algebraičnyh funkcij jednoj proměnnoj & \hyperlink{TERM-0267}{TERM-0267}\\
asociovany prost ideal & \hyperlink{TERM-1939}{TERM-1939}\\
asociovany prosty ideal / naležati k p & \hyperlink{TERM-1941}{TERM-1941}\\
baza ... prirěđena posrědstvom izomorfizma & \hyperlink{TERM-1812}{TERM-1812}\\
baza direktnoj sumy / blokovo-diagonalna diskriminantna matrica & \hyperlink{TERM-0161}{TERM-0161}\\
baza diskriminantnogo ideala & \hyperlink{TERM-0860}{TERM-0860}\\
baza ideala & \hyperlink{TERM-0858}{TERM-0858}\\
baza razširjenogo ideala diskriminantnogo ideala & \hyperlink{TERM-0163}{TERM-0163}\\
baza za tvorjenje slěda & \hyperlink{TERM-0164}{TERM-0164}\\
bazis modula / konečny bazis modula & \hyperlink{TERM-1173}{TERM-1173}\\
bazisna funkcija & \hyperlink{TERM-0169}{TERM-0169}\\
bazova veličina & \hyperlink{TERM-0170}{TERM-0170}\\
bazovo svojstvo & \hyperlink{TERM-0166}{TERM-0166}\\
beskonečno blizke točke / točke razvětvenja & \hyperlink{TERM-0144}{TERM-0144}\\
beskonečno mnogo razkladov & \hyperlink{TERM-1769}{TERM-1769}\\
bez povtorjenja / jedin raz i samo jedin raz & \hyperlink{TERM-1277}{TERM-1277}\\
bezbazisna forma & \hyperlink{TERM-0168}{TERM-0168}\\
bibliografija & \hyperlink{TERM-0186}{TERM-0186}\\
bilinearne kombinacije & \hyperlink{TERM-0190}{TERM-0190}\\
bimodul & \hyperlink{TERM-0341}{TERM-0341}\\
binarna forma / binarna linearna forma & \hyperlink{TERM-0191}{TERM-0191}\\
binarne rědy / rědy proměnnyh & \hyperlink{TERM-0192}{TERM-0192}\\
biracionalna transformacija & \hyperlink{TERM-0193}{TERM-0193}\\
blokova matrica / blokovo-diagonalna matrica & \hyperlink{TERM-0196}{TERM-0196}\\
blokovo-trikutna forma / matrično predstavjenje & \hyperlink{TERM-0195}{TERM-0195}\\
Bolje ne jest bolši & \hyperlink{DEC-0082}{DEC-0082}\\
Brauerova grupa & \hyperlink{TERM-1430}{TERM-1430}\\
Brauerova grupa klasov algeber & \hyperlink{TERM-0197}{TERM-0197}\\
byly poznane kako reducibilne & \hyperlink{TERM-0131}{TERM-0131}\\
byti algebraično cěly nad \textbackslash{}mathfrak F & \hyperlink{TERM-0094}{TERM-0094}\\
Byti i biti & \hyperlink{DEC-0105}{DEC-0105}\\
c?la algebrai?na funkcija & \hyperlink{TERM-0639}{TERM-0639}\\
častne derivaty 2-go reda & \hyperlink{TERM-0282}{TERM-0282}\\
cěla racionalna funkcija & \hyperlink{TERM-0651}{TERM-0651}\\
cěla racionalna kombinacija & \hyperlink{TERM-0656}{TERM-0656}\\
cěla simetrična funkcija & \hyperlink{TERM-0660}{TERM-0660}\\
cěle algebraične funkcije & \hyperlink{TERM-0640}{TERM-0640}\\
cěle celočislene funkcije & \hyperlink{TERM-0641}{TERM-0641}\\
cěle celočislene simetrične funkcije & \hyperlink{TERM-0642}{TERM-0642}\\
cěle celočislene simetrične funkcije od n redov veličin & \hyperlink{TERM-0643}{TERM-0643}\\
cěle invariantne tvorjenja & \hyperlink{TERM-0647}{TERM-0647}\\
cěle racionalne celočislene kombinacije & \hyperlink{TERM-0659}{TERM-0659}\\
cěle racionalne čisla & \hyperlink{TERM-0658}{TERM-0658}\\
cěle racionalne funkcije & \hyperlink{TERM-0652}{TERM-0652}\\
cěle racionalne funkcije (polinomy) & \hyperlink{TERM-0653}{TERM-0653}\\
cěle racionalne funkcije, to jest polinomy & \hyperlink{TERM-0654}{TERM-0654}\\
cěle racionalne projektivne invarianty kvadratičnoj formy & \hyperlink{TERM-0655}{TERM-0655}\\
cěle transcendentne čisla & \hyperlink{TERM-0661}{TERM-0661}\\
cěle veličiny & \hyperlink{TERM-0644}{TERM-0644}\\
cěle veličiny / cěle racionalne čisla & \hyperlink{TERM-0645}{TERM-0645}\\
cěle veličiny libovoljnogo polja & \hyperlink{TERM-0646}{TERM-0646}\\
Cělo čislo; cěločislny & \hyperlink{DEC-0037}{DEC-0037}\\
cělo racionalno čislo & \hyperlink{TERM-0657}{TERM-0657}\\
cělo racionalno predstavjanje & \hyperlink{TERM-0649}{TERM-0649}\\
cělo racionalno predstavjenje & \hyperlink{TERM-0650}{TERM-0650}\\
cělo-racionalne čislove koeficienty & \hyperlink{TERM-0663}{TERM-0663}\\
cělo-racionalny & \hyperlink{TERM-0634}{TERM-0634}\\
cěločiselne matrice & \hyperlink{TERM-0672}{TERM-0672}\\
celočislena integralna baza & \hyperlink{TERM-0668}{TERM-0668}\\
celočislena integralna oblast & \hyperlink{TERM-0674}{TERM-0674}\\
celočislena linearna familija & \hyperlink{TERM-0671}{TERM-0671}\\
celočislena relativno cěla oblast & \hyperlink{TERM-0675}{TERM-0675}\\
celočislena relativno cěla oblast prvogo roda & \hyperlink{TERM-0676}{TERM-0676}\\
celočislene koeficienty & \hyperlink{TERM-0670}{TERM-0670}\\
celočisleno-regularny & \hyperlink{TERM-0666}{TERM-0666}\\
celočisleny & \hyperlink{TERM-0665}{TERM-0665}\\
celočisleny polinom & \hyperlink{TERM-0677}{TERM-0677}\\
celočisleny sistem S\_\{n rho\} & \hyperlink{TERM-0678}{TERM-0678}\\
cěločislne funkcije & \hyperlink{TERM-0667}{TERM-0667}\\
celočislno grupovo kolco & \hyperlink{TERM-0673}{TERM-0673}\\
celočislovost & \hyperlink{TERM-0679}{TERM-0679}\\
cěločislovost & \hyperlink{TERM-0680}{TERM-0680}\\
cěločislovy invariant & \hyperlink{TERM-0669}{TERM-0669}\\
cělokupnost & \hyperlink{TERM-0703}{TERM-0703}\\
cělokupnost polinomov & \hyperlink{TERM-0705}{TERM-0705}\\
cělost vseh funkcij & \hyperlink{TERM-0704}{TERM-0704}\\
cěly i primitivny po sootvětnyh neopreděljenyh & \hyperlink{TERM-0635}{TERM-0635}\\
cěly po neopreděljenyh & \hyperlink{TERM-0633}{TERM-0633}\\
cěly racionalny (absolutny) invariant & \hyperlink{TERM-0648}{TERM-0648}\\
cěly racionalny funkcional & \hyperlink{TERM-0662}{TERM-0662}\\
cěly ręd form & \hyperlink{TERM-0700}{TERM-0700}\\
cěly vzhodno k R & \hyperlink{TERM-0631}{TERM-0631}\\
cěly vzhodno k R / elementy cěle vzhodno k R & \hyperlink{TERM-0632}{TERM-0632}\\
centar & \hyperlink{TERM-1896}{TERM-1896}\\
centar / predstavjenja centra & \hyperlink{TERM-1897}{TERM-1897}\\
centralny element & \hyperlink{TERM-1898}{TERM-1898}\\
centralny element / centr & \hyperlink{TERM-1899}{TERM-1899}\\
cěp děliteljev & \hyperlink{TERM-1726}{TERM-1726}\\
cěp děliteljev ili kratnikov & \hyperlink{TERM-1711}{TERM-1711}\\
cěp idealov po dělimosti / cěp mnogokratnyh modulo ideal & \hyperlink{TERM-1727}{TERM-1727}\\
cěp kratnyh & \hyperlink{TERM-1828}{TERM-1828}\\
cepi prostyh idealov / cepova definicija & \hyperlink{TERM-0975}{TERM-0975}\\
čest I / čest II & \hyperlink{TERM-1701}{TERM-1701}\\
četverna grupa / grupa osmogo reda & \hyperlink{TERM-1831}{TERM-1831}\\
Četyrěh i šesti & \hyperlink{DEC-0112}{DEC-0112}\\
ciklična / Dicksonov tip & \hyperlink{TERM-1983}{TERM-1983}\\
ciklična polja Omega\_n stepena 2\textasciicircum{}n & \hyperlink{TERM-1985}{TERM-1985}\\
ciklična polja stepena 2\textasciicircum{}n & \hyperlink{TERM-1986}{TERM-1986}\\
ciklična zaměna prěměnnyh & \hyperlink{TERM-1987}{TERM-1987}\\
ciklične grupy & \hyperlink{TERM-1984}{TERM-1984}\\
ciklotomične polja & \hyperlink{TERM-1058}{TERM-1058}\\
Čirnhausova transformacija & \hyperlink{TERM-1755}{TERM-1755}\\
čisla mnogokratnosti & \hyperlink{TERM-1187}{TERM-1187}\\
čislo rędov & \hyperlink{TERM-1511}{TERM-1511}\\
čislove i funkcionalne polja & \hyperlink{TERM-1889}{TERM-1889}\\
čislove koeficienty & \hyperlink{TERM-1271}{TERM-1271}\\
čislove polja i polja funkcij & \hyperlink{TERM-1890}{TERM-1890}\\
čislove polja najnižšego stepena / stepen nad P & \hyperlink{TERM-1894}{TERM-1894}\\
čislovo fiksovati & \hyperlink{TERM-1891}{TERM-1891}\\
čislovo polje & \hyperlink{TERM-1893}{TERM-1893}\\
čislovo-teoretičny teorem sušstvovanja & \hyperlink{TERM-1892}{TERM-1892}\\
čisto imaginarny & \hyperlink{TERM-1512}{TERM-1512}\\
Citacija v statji i v zbirnoj knigě & \hyperlink{DEC-0171}{DEC-0171}\\
citation preserved with page marker str. & \hyperlink{TERM-1638}{TERM-1638}\\
Citovane W. Schmeidler i sklonjeno V. Šmajdlera & \hyperlink{DEC-0170}{DEC-0170}\\
Citovany jezyk ostaje citovanym jezykom & \hyperlink{DEC-0148}{DEC-0148}\\
členy najvyššej dimenzije & \hyperlink{TERM-0719}{TERM-0719}\\
črez dvojnu redukciju & \hyperlink{TERM-0359}{TERM-0359}\\
črez homomorfizm diskriminant prěhodi & \hyperlink{TERM-0356}{TERM-0356}\\
črez tvorjenje kvocientov poraždaje medžupolje & \hyperlink{TERM-0365}{TERM-0365}\\
črez vraćanje ... k nižšim formam & \hyperlink{TERM-0366}{TERM-0366}\\
črez vraćanje od $j$ i \$\textbackslash{}A\$ k nižšim formam & \hyperlink{TERM-0367}{TERM-0367}\\
črez zaměnu ... & \hyperlink{TERM-0363}{TERM-0363}\\
črez zaměnu formy \textbackslash{}A nižšimi formami ręda form & \hyperlink{TERM-0362}{TERM-0362}\\
Čto se tyče: o čem govorit fraza & \hyperlink{DEC-0127}{DEC-0127}\\
Čto v pytanju; sut v prozě & \hyperlink{DEC-0020}{DEC-0020}\\
Curišsko: pridavnik od Zürich & \hyperlink{DEC-0153}{DEC-0153}\\
D.–W. i neměnjane bibliografične frazy & \hyperlink{DEC-0169}{DEC-0169}\\
d?litvena algebra; prir??ena d?litvena algebra & \hyperlink{TERM-0334}{TERM-0334}\\
dalej ne budut spominjane & \hyperlink{TERM-1640}{TERM-1640}\\
Daljši, dalje, daljny & \hyperlink{DEC-0040}{DEC-0040}\\
Daljši, ranši, pozdnějši, najobčejši & \hyperlink{DEC-0109}{DEC-0109}\\
Danyh i oddělnyh: -nyh pri tvrdem n & \hyperlink{DEC-0091}{DEC-0091}\\
Dedekind--Mertensov teorem & \hyperlink{TERM-0246}{TERM-0246}\\
Dedekindov poslědok I/II & \hyperlink{TERM-0249}{TERM-0249}\\
Dedekindov poslědok II & \hyperlink{TERM-0250}{TERM-0250}\\
Dedekindovo definovanje / Dedekind-Henselovogo tipa & \hyperlink{TERM-0248}{TERM-0248}\\
defekt svijanja & \hyperlink{TERM-0251}{TERM-0251}\\
definicija Clebscha / ekvivalentne & \hyperlink{TERM-0218}{TERM-0218}\\
definicija se pokryva s definicijeju v §7, 1 & \hyperlink{TERM-0268}{TERM-0268}\\
definirano, kako v §1, 2, kako kolco togo samogo ranga & \hyperlink{TERM-1877}{TERM-1877}\\
definujuči ideal & \hyperlink{TERM-0252}{TERM-0252}, \hyperlink{TERM-1131}{TERM-1131}\\
definujųcy ideal & \hyperlink{TERM-0253}{TERM-0253}\\
dělimo črez p / izčezaje v P & \hyperlink{TERM-0364}{TERM-0364}\\
dělimost / dělimy čerez modul & \hyperlink{TERM-1703}{TERM-1703}\\
dělimost D\_T črez p & \hyperlink{TERM-1704}{TERM-1704}\\
dělimy črez (x\_i-a) & \hyperlink{TERM-0353}{TERM-0353}\\
dělimy črez ideal / f ≡ 0 (M) & \hyperlink{TERM-0361}{TERM-0361}\\
dělimy nerazložimym uravnenjem & \hyperlink{TERM-1702}{TERM-1702}\\
dělitelj & \hyperlink{TERM-1706}{TERM-1706}\\
dělitelj jednakogo ranga & \hyperlink{TERM-1708}{TERM-1708}\\
dělitelj polja K & \hyperlink{TERM-1710}{TERM-1710}\\
dělitelj v idealnom smyslu & \hyperlink{TERM-1709}{TERM-1709}\\
determinant matrice P jest jedinica & \hyperlink{TERM-0255}{TERM-0255}\\
determinant porjadka rho & \hyperlink{TERM-1587}{TERM-1587}\\
determinant transformacije & \hyperlink{TERM-1739}{TERM-1739}\\
determinant transformacije / jedinica & \hyperlink{TERM-1740}{TERM-1740}\\
determinant |a\_\{ik\}| & \hyperlink{TERM-0257}{TERM-0257}\\
determinant |S(tau\_i tau\_j)| & \hyperlink{TERM-1990}{TERM-1990}\\
determinanta hiperkompleksnogo sistema & \hyperlink{TERM-0256}{TERM-0256}\\
determinanta reda najmanje dva & \hyperlink{TERM-1150}{TERM-1150}\\
determinantny dělitelj & \hyperlink{TERM-0261}{TERM-0261}\\
determinanty / stupenj & \hyperlink{TERM-0258}{TERM-0258}\\
determinanty iz vsakih n baznyh elementov i jih konjugovanyh elementov & \hyperlink{TERM-0259}{TERM-0259}\\
diagonalna matrica / transformovane čerez diagonalne matrice & \hyperlink{TERM-0263}{TERM-0263}\\
diagonalna shema & \hyperlink{TERM-0265}{TERM-0265}\\
diagonalne matrice / nulova matrica & \hyperlink{TERM-0264}{TERM-0264}\\
diagonalny člen & \hyperlink{TERM-0262}{TERM-0262}\\
diedralne grupy & \hyperlink{TERM-0273}{TERM-0273}\\
diferencialne invarianty & \hyperlink{TERM-0279}{TERM-0279}\\
diferencialne izrazy & \hyperlink{TERM-0277}{TERM-0277}, \hyperlink{TERM-0278}{TERM-0278}\\
diferencialne simboly & \hyperlink{TERM-0285}{TERM-0285}\\
diferencialny izraz & \hyperlink{TERM-0276}{TERM-0276}\\
diferencialny kvocient & \hyperlink{TERM-0280}{TERM-0280}\\
diferencialny kvocient / razlikovy kvocient & \hyperlink{TERM-0281}{TERM-0281}\\
diferenciranje idealov & \hyperlink{TERM-0861}{TERM-0861}\\
diferenta & \hyperlink{TERM-0275}{TERM-0275}\\
dimenzija v sobstvenom smyslu & \hyperlink{TERM-0289}{TERM-0289}\\
dimenzijne čisla & \hyperlink{TERM-0290}{TERM-0290}\\
direktna dvustranna suma & \hyperlink{TERM-0294}{TERM-0294}\\
direktna suma & \hyperlink{TERM-0300}{TERM-0300}\\
direktna suma / direktny sumand & \hyperlink{TERM-0301}{TERM-0301}\\
direktna suma direktno nerazložimyh jednostavnyh idealov & \hyperlink{TERM-0302}{TERM-0302}\\
direktna suma kolc ranga jedan & \hyperlink{TERM-0304}{TERM-0304}\\
direktna suma konečno mnogyh polj & \hyperlink{TERM-0303}{TERM-0303}\\
direktne sumy klasov & \hyperlink{TERM-0305}{TERM-0305}\\
direktno nerazložima & \hyperlink{TERM-0291}{TERM-0291}\\
direktno nerazložime faktory & \hyperlink{TERM-0293}{TERM-0293}\\
direktno nerazložimy & \hyperlink{TERM-0292}{TERM-0292}, \hyperlink{TERM-0295}{TERM-0295}\\
direktno nerazložimy / razložimy & \hyperlink{TERM-0296}{TERM-0296}\\
direktno razloženje na prave idealy & \hyperlink{TERM-0308}{TERM-0308}\\
direktno sumno razloženje & \hyperlink{TERM-0307}{TERM-0307}\\
direktny faktor & \hyperlink{TERM-0310}{TERM-0310}\\
direktny prěsěk & \hyperlink{TERM-0309}{TERM-0309}\\
direktny produkt & \hyperlink{TERM-0313}{TERM-0313}\\
direktny produkt / tvorjenje direktnogo produkta & \hyperlink{TERM-0314}{TERM-0314}\\
direktny sumand & \hyperlink{TERM-0311}{TERM-0311}\\
direktny summand & \hyperlink{TERM-0312}{TERM-0312}\\
diskontinuirny / ekstremno diskontinuirny / totalno diskontinuirny & \hyperlink{TERM-1777}{TERM-1777}\\
diskontinuirny / kontinuirny / ekstremno diskontinuirny & \hyperlink{TERM-1778}{TERM-1778}\\
diskriminant & \hyperlink{TERM-0315}{TERM-0315}\\
diskriminant D\_T / kvadrat determinanta konjugovanyh baznyh elementov & \hyperlink{TERM-0316}{TERM-0316}\\
diskriminant D\_\{T\_P\} & \hyperlink{TERM-0317}{TERM-0317}\\
diskriminant ideala vzhodno k klasu & \hyperlink{TERM-0318}{TERM-0318}\\
diskriminant kako kvadrat determinanta konjugovanyh baznyh elementov & \hyperlink{TERM-1988}{TERM-1988}\\
diskriminant R izčezaje & \hyperlink{TERM-0319}{TERM-0319}\\
diskriminant, opreděljeny do kvadrata jedinicy, jest različny od nuly & \hyperlink{TERM-0194}{TERM-0194}\\
diskriminantna matrica & \hyperlink{TERM-0330}{TERM-0330}\\
diskriminantne idealy primarnyh kolc & \hyperlink{TERM-0328}{TERM-0328}\\
diskriminantny ideal & \hyperlink{TERM-0320}{TERM-0320}\\
diskriminantny ideal ... nulovy ideal & \hyperlink{TERM-0321}{TERM-0321}\\
diskriminantny ideal jest jediničny ideal v P & \hyperlink{TERM-0324}{TERM-0324}\\
diskriminantny ideal kolca & \hyperlink{TERM-0323}{TERM-0323}\\
diskriminantny ideal komponentnogo ideala & \hyperlink{TERM-0322}{TERM-0322}\\
diskriminantny ideal od R & \hyperlink{TERM-0325}{TERM-0325}\\
diskriminantny ideal porjadka T & \hyperlink{TERM-0326}{TERM-0326}\\
diskriminantny ideal pri prěhodu k kolcu kvocientov & \hyperlink{TERM-0241}{TERM-0241}\\
diskriminantny ideal T vzhodno k o & \hyperlink{TERM-0327}{TERM-0327}\\
diskriminantny kriterij polnoj razložimosti prvogo roda & \hyperlink{TERM-0329}{TERM-0329}\\
divergencija / divergencijne relacije & \hyperlink{TERM-0333}{TERM-0333}\\
dlžina i & \hyperlink{TERM-1128}{TERM-1128}\\
dobro uporjadkovana množina / dobro uporjadkovanje & \hyperlink{TERM-1880}{TERM-1880}\\
dobro uporjadkovanje & \hyperlink{TERM-1881}{TERM-1881}\\
dobro uporjadkovanje / porjadok elementov & \hyperlink{TERM-1882}{TERM-1882}\\
dobro uporjadkovanje / teorema o dobrom uporjadkovanju & \hyperlink{TERM-1883}{TERM-1883}\\
dobro uporjadkovanje Omega prime & \hyperlink{TERM-1884}{TERM-1884}\\
Dobyti i dobyvati; dobyty rezultat & \hyperlink{DEC-0022}{DEC-0022}\\
Dobyvati i dobyty: dobivanje rezultata & \hyperlink{DEC-0073}{DEC-0073}\\
Dobyvati ne jest dostavjati & \hyperlink{DEC-0075}{DEC-0075}\\
Dodatna primětka & \hyperlink{TERM-1202}{TERM-1202}\\
dodatne členy n\_i xi\_i / dodatne cěločiselne kratne & \hyperlink{TERM-1967}{TERM-1967}\\
dodatno uslovje & \hyperlink{TERM-1205}{TERM-1205}\\
dodatok & \hyperlink{TERM-1966}{TERM-1966}\\
dokaz konečnosti sistema form & \hyperlink{TERM-0495}{TERM-0495}\\
dolžina ideala / kompozicijna serija & \hyperlink{TERM-1079}{TERM-1079}\\
dopolnitelno predstavjenje & \hyperlink{TERM-0505}{TERM-0505}\\
dopolnjenje algebrai?noj bazy do racionalnoj & \hyperlink{TERM-0504}{TERM-0504}\\
dopustima podgrupa & \hyperlink{TERM-1949}{TERM-1949}\\
dopustimy ideal & \hyperlink{TERM-1950}{TERM-1950}\\
dopustimy sistem & \hyperlink{TERM-1948}{TERM-1948}, \hyperlink{TERM-1951}{TERM-1951}\\
Doskonalo polje; nedoskonalo polje & \hyperlink{DEC-0048}{DEC-0048}\\
Dozvaljati, jednočasno i doslědny pravopis & \hyperlink{DEC-0150}{DEC-0150}\\
drobne funkcionaly & \hyperlink{TERM-0684}{TERM-0684}\\
drobno algebraično čislo & \hyperlink{TERM-0683}{TERM-0683}\\
drobno-linearne preobraženje & \hyperlink{TERM-1119}{TERM-1119}\\
drobny ideal & \hyperlink{TERM-0685}{TERM-0685}\\
drobny ideal / pravo drobny element & \hyperlink{TERM-0686}{TERM-0686}\\
drobny racionalny funkcional & \hyperlink{TERM-0687}{TERM-0687}\\
druga baza togo samogo ideala & \hyperlink{TERM-0133}{TERM-0133}\\
drugi teorem o izomorfizmu & \hyperlink{TERM-1975}{TERM-1975}\\
drugy teorem izomorfizma & \hyperlink{TERM-1976}{TERM-1976}\\
drugy teorem o izomorfizmu & \hyperlink{TERM-1977}{TERM-1977}\\
dualistična razkladna identičnost & \hyperlink{TERM-0352}{TERM-0352}\\
dualistične formy & \hyperlink{TERM-0351}{TERM-0351}\\
dualistične identičnosti & \hyperlink{TERM-0350}{TERM-0350}\\
dualistično protivopoložne & \hyperlink{TERM-0348}{TERM-0348}\\
dualistično protivopoložny & \hyperlink{TERM-0349}{TERM-0349}\\
dvojna oblast (Σ,T) & \hyperlink{TERM-0335}{TERM-0335}\\
dvojna redukcija & \hyperlink{TERM-0346}{TERM-0346}\\
dvojno izračunanje ... daje povod k redukciji & \hyperlink{TERM-0269}{TERM-0269}\\
dvojno uslovje cěpov & \hyperlink{TERM-0337}{TERM-0337}\\
dvojno uslovje cěpov / uslovje cěpov děliteljev / uslovje cěpov kratnyh & \hyperlink{TERM-0340}{TERM-0340}\\
dvojno uslovje konečnoj cěpi & \hyperlink{TERM-0338}{TERM-0338}\\
dvojnorazova polarizacija & \hyperlink{TERM-1968}{TERM-1968}\\
dvojny indeks & \hyperlink{TERM-0336}{TERM-0336}\\
dvojny modul & \hyperlink{TERM-0342}{TERM-0342}\\
dvojny modul odnosno bimodul & \hyperlink{TERM-0343}{TERM-0343}\\
dvomodul & \hyperlink{TERM-0344}{TERM-0344}\\
dvustranne idealy & \hyperlink{TERM-1973}{TERM-1973}\\
dvustranne klasy & \hyperlink{TERM-1974}{TERM-1974}\\
dvustranno nerazložime kolca & \hyperlink{TERM-1972}{TERM-1972}\\
dvustranno proste kolca & \hyperlink{TERM-1970}{TERM-1970}\\
dvustranno prosto kolco & \hyperlink{TERM-1971}{TERM-1971}\\
dvustranno prosty & \hyperlink{TERM-1969}{TERM-1969}\\
E. Noether, O. Ore & \hyperlink{TERM-0370}{TERM-0370}\\
e\_1,\textbackslash{}ldots,e\_n kako modulna baza / diagonalna matrica & \hyperlink{TERM-0371}{TERM-0371}\\
Egzistovati i obstajati & \hyperlink{DEC-0135}{DEC-0135}\\
eksplicitno predstavjenje & \hyperlink{TERM-0536}{TERM-0536}\\
eksplicitno predstavjenje invariantnyh tvorjenj & \hyperlink{TERM-0537}{TERM-0537}\\
eksplicitno razloženje & \hyperlink{TERM-0538}{TERM-0538}\\
eksponent & \hyperlink{TERM-0539}{TERM-0539}\\
eksponent = indeks & \hyperlink{TERM-0540}{TERM-0540}\\
eksponent primarnogo ideala & \hyperlink{TERM-0541}{TERM-0541}, \hyperlink{TERM-0542}{TERM-0542}\\
ekstremaly & \hyperlink{TERM-0543}{TERM-0543}\\
ekvivalentne kvadratne diferencialne formy & \hyperlink{TERM-1993}{TERM-1993}\\
ekvivalentne predstavjenja / klas predstavjenj & \hyperlink{TERM-1992}{TERM-1992}\\
ekvivalentno četverokratno simbolično predstavjenje & \hyperlink{TERM-1994}{TERM-1994}\\
ekvivalentno razširjeno polje / razširjeno kolco & \hyperlink{TERM-1995}{TERM-1995}\\
ekvivalentnost črez formy vyššego svijanja & \hyperlink{TERM-1997}{TERM-1997}\\
ekvivalentny sistem & \hyperlink{TERM-1996}{TERM-1996}\\
element drugogo roda / eksponenta & \hyperlink{TERM-0462}{TERM-0462}\\
element p, čij glavny ideal jest dělimy točno prvoju stepenju & \hyperlink{TERM-0460}{TERM-0460}\\
element prosty k a / vzajemno prosty v idealnom smyslu & \hyperlink{TERM-1934}{TERM-1934}\\
element s nulevoju normoju & \hyperlink{TERM-0461}{TERM-0461}\\
element, ktory poraždaje izvedeny prosty ideal & \hyperlink{TERM-0162}{TERM-0162}\\
elementarne dělitelje & \hyperlink{TERM-0464}{TERM-0464}\\
elementarne kovarianty & \hyperlink{TERM-0463}{TERM-0463}\\
elementarne simetrične funkcije & \hyperlink{TERM-1683}{TERM-1683}\\
elementarny dělitelj & \hyperlink{TERM-0465}{TERM-0465}\\
elementno komutuju?e podkolca; komutantny teorem & \hyperlink{TERM-0471}{TERM-0471}\\
elementy cěle vzhodno k mo & \hyperlink{TERM-0890}{TERM-0890}\\
elementy iz Ω & \hyperlink{TERM-0753}{TERM-0753}\\
eliminacija & \hyperlink{TERM-0472}{TERM-0472}\\
energijske relacije / energijske zakony / energijske komponenty & \hyperlink{TERM-0500}{TERM-0500}\\
faktor proporcionalnosti & \hyperlink{TERM-1405}{TERM-1405}\\
faktor-grupy & \hyperlink{TERM-1547}{TERM-1547}\\
faktor-kolco R/p ranga jedan & \hyperlink{TERM-1566}{TERM-1566}\\
faktorgrupa & \hyperlink{TERM-0548}{TERM-0548}\\
faktorna sistema & \hyperlink{TERM-0545}{TERM-0545}\\
faktorna sistema jedinica & \hyperlink{TERM-0546}{TERM-0546}\\
faktorne sistemy & \hyperlink{TERM-0547}{TERM-0547}\\
fiksovany klas & \hyperlink{TERM-0567}{TERM-0567}\\
foreign journal/title fragments protected with \textbackslash{}emph; running Interslavic prose transliterated & \hyperlink{TERM-1404}{TERM-1404}\\
forma $\mu$-go stepena po $x$ i $\lambda$-go stepena po $u$ & \hyperlink{TERM-0001}{TERM-0001}\\
forma elementarnyh děliteljev & \hyperlink{TERM-0466}{TERM-0466}\\
forma elementarnyh děliteljev / eksponenta & \hyperlink{TERM-0467}{TERM-0467}\\
forma elementarnyh děliteljev / norma & \hyperlink{TERM-0468}{TERM-0468}\\
forma kriviny & \hyperlink{TERM-0571}{TERM-0571}, \hyperlink{TERM-1064}{TERM-1064}\\
forma reducibilnosti & \hyperlink{TERM-1490}{TERM-1490}\\
formalne linearne formy & \hyperlink{TERM-0572}{TERM-0572}\\
formalny variacionny račun & \hyperlink{TERM-0573}{TERM-0573}\\
Formula v nazvanju statje & \hyperlink{DEC-0167}{DEC-0167}\\
formy 10. poręda & \hyperlink{TERM-0589}{TERM-0589}\\
formy 11. poręda & \hyperlink{TERM-0574}{TERM-0574}\\
formy 11., 12., 13. i 15. poręda & \hyperlink{TERM-0575}{TERM-0575}\\
formy 12., 13., 14., 15. poręda & \hyperlink{TERM-0576}{TERM-0576}\\
formy 7. poręda (sistem $s_{III}$) & \hyperlink{TERM-0577}{TERM-0577}\\
formy 8. poręda & \hyperlink{TERM-0578}{TERM-0578}\\
formy 9. poręda & \hyperlink{TERM-0583}{TERM-0583}\\
formy četvrtogo poręda & \hyperlink{TERM-0587}{TERM-0587}\\
formy desętogo poręda & \hyperlink{TERM-0590}{TERM-0590}\\
formy devętogo poręda & \hyperlink{TERM-0584}{TERM-0584}\\
formy inercije & \hyperlink{TERM-1753}{TERM-1753}\\
formy osmogo poręda & \hyperlink{TERM-0579}{TERM-0579}\\
formy petogo poręda & \hyperlink{TERM-0581}{TERM-0581}\\
formy s vyšimi simbolami & \hyperlink{TERM-0582}{TERM-0582}\\
formy sedmogo poręda & \hyperlink{TERM-0586}{TERM-0586}\\
formy šestogo poręda & \hyperlink{TERM-0585}{TERM-0585}\\
formy tretjego poręda & \hyperlink{TERM-0580}{TERM-0580}\\
formy v n prěměnnyh & \hyperlink{TERM-0588}{TERM-0588}\\
formy, koje se razkladųt na faktory & \hyperlink{TERM-1909}{TERM-1909}\\
Frobeniusova teorija grupnyh karakterov & \hyperlink{TERM-0602}{TERM-0602}\\
fundamentalna točka & \hyperlink{TERM-0607}{TERM-0607}\\
fundamentalne točke involucije & \hyperlink{TERM-0608}{TERM-0608}\\
fundamentalno jednačenje / fundamentalna forma & \hyperlink{TERM-0605}{TERM-0605}\\
fundamentalno uravnenje & \hyperlink{TERM-0604}{TERM-0604}\\
fundamentalno uravnenje / fundamentalna forma & \hyperlink{TERM-0606}{TERM-0606}\\
fundamentalny teorem teorije eliminacije & \hyperlink{TERM-0609}{TERM-0609}\\
funkcija prvogo roda / polje prvogo roda & \hyperlink{TERM-0610}{TERM-0610}\\
funkcije nultogo stepena & \hyperlink{TERM-0617}{TERM-0617}\\
funkcije odobraženja & \hyperlink{TERM-0024}{TERM-0024}\\
funkcionalna matrica & \hyperlink{TERM-0616}{TERM-0616}\\
funkcionalna oblast & \hyperlink{TERM-0611}{TERM-0611}\\
funkcionalne determinanty & \hyperlink{TERM-0614}{TERM-0614}\\
funkcionalne uravnenja & \hyperlink{TERM-0615}{TERM-0615}\\
funkcionalny determinant & \hyperlink{TERM-0612}{TERM-0612}\\
funkcionalny determinant / funkcionalna matrica & \hyperlink{TERM-0613}{TERM-0613}\\
Galoisov modul & \hyperlink{TERM-0622}{TERM-0622}\\
Galoisova forma grupy & \hyperlink{TERM-0624}{TERM-0624}\\
Galoisova polja / Galoisovo razširjeno polje & \hyperlink{TERM-0625}{TERM-0625}\\
Galoisova rezolventa & \hyperlink{TERM-0621}{TERM-0621}, \hyperlink{TERM-0626}{TERM-0626}\\
Galoisova teorija nekomutativnyh těl & \hyperlink{TERM-0627}{TERM-0627}\\
Galoisovo kolco razširenja & \hyperlink{TERM-0628}{TERM-0628}\\
Galoisovo polje & \hyperlink{TERM-0629}{TERM-0629}\\
Galoisovy modul & \hyperlink{TERM-0623}{TERM-0623}\\
Genitiv množiny glagolnyh imenovnikov & \hyperlink{DEC-0058}{DEC-0058}\\
Glava III & \hyperlink{TERM-0973}{TERM-0973}\\
glavna diagonala & \hyperlink{TERM-0783}{TERM-0783}\\
glavna jedinica & \hyperlink{TERM-0784}{TERM-0784}\\
glavna teorema & \hyperlink{TERM-0803}{TERM-0803}\\
glavny ideal & \hyperlink{TERM-0786}{TERM-0786}\\
glavny ideal / bazny polinom & \hyperlink{TERM-0787}{TERM-0787}\\
glavny ideal Tp / po parah vzajemno proste primarne komponenty & \hyperlink{TERM-0788}{TERM-0788}\\
glavny ideal, izvedeny iz p & \hyperlink{TERM-0031}{TERM-0031}\\
glavny klas & \hyperlink{TERM-0793}{TERM-0793}\\
glavny modul & \hyperlink{TERM-0794}{TERM-0794}\\
glavny porjadok & \hyperlink{TERM-0795}{TERM-0795}\\
glavny porjadok / glavna oblast & \hyperlink{TERM-0797}{TERM-0797}\\
glavny porjadok / kvadrat prostogo ideala & \hyperlink{TERM-0798}{TERM-0798}\\
glavny porjadok čislovogo abo funkcionalnogo polja & \hyperlink{TERM-0800}{TERM-0800}\\
glavny porjadok čislovogo polja & \hyperlink{TERM-0801}{TERM-0801}\\
glavny porjadok relativnogo polja & \hyperlink{TERM-0799}{TERM-0799}\\
glavny rjad & \hyperlink{TERM-0802}{TERM-0802}\\
glavny slěd & \hyperlink{TERM-0805}{TERM-0805}\\
glavny teorem / pomočny teorem / dodatok / zamětka & \hyperlink{TERM-0804}{TERM-0804}\\
gorny / dolny indeks vagy & \hyperlink{TERM-1276}{TERM-1276}\\
Göttingen & \hyperlink{TERM-0778}{TERM-0778}\\
granične členy / graničny integral & \hyperlink{TERM-1434}{TERM-1434}\\
granične oblasti & \hyperlink{TERM-0726}{TERM-0726}\\
grupa analitičnyh transformacij & \hyperlink{TERM-0741}{TERM-0741}\\
grupa posunov & \hyperlink{TERM-1814}{TERM-1814}\\
grupa razlo?enja; razv?tvjene m?sta; nerazv?tvjeny & \hyperlink{TERM-1925}{TERM-1925}\\
grupa s oblastju operatorov & \hyperlink{TERM-0742}{TERM-0742}\\
grupa transformacij & \hyperlink{TERM-1742}{TERM-1742}\\
grupne karaktere & \hyperlink{TERM-0743}{TERM-0743}\\
grupno kolco & \hyperlink{TERM-0749}{TERM-0749}\\
grupno kolco / grupne čisla & \hyperlink{TERM-0751}{TERM-0751}\\
grupno kolco / kvaternionno tělo & \hyperlink{TERM-0752}{TERM-0752}\\
grupno razděljenje po nevzaimnoj prostosti & \hyperlink{TERM-0747}{TERM-0747}\\
grupno razděljenje pridruženih prostyh idealov & \hyperlink{TERM-0746}{TERM-0746}\\
grupny determinant & \hyperlink{TERM-0744}{TERM-0744}\\
grupova determinanta & \hyperlink{TERM-0745}{TERM-0745}\\
Grupovati i pojmati: dva smysly shvat- & \hyperlink{DEC-0139}{DEC-0139}\\
grupovo kolco & \hyperlink{TERM-0750}{TERM-0750}\\
grupy ostatkov sut izomorfne & \hyperlink{TERM-1548}{TERM-1548}\\
harakteristična mnogokratnost & \hyperlink{TERM-0214}{TERM-0214}\\
harakteristika p & \hyperlink{TERM-0211}{TERM-0211}\\
harakterizujuča oblast odobraženja & \hyperlink{TERM-0206}{TERM-0206}\\
Harakterne svojstva i stepenne redy & \hyperlink{DEC-0107}{DEC-0107}\\
harakternye svojstva & \hyperlink{TERM-0212}{TERM-0212}\\
Hilbertov teorem II o celočislenej modulnoj bazi & \hyperlink{TERM-0814}{TERM-0814}\\
Hilbertov teorem o bazisu modula & \hyperlink{TERM-0812}{TERM-0812}\\
Hilbertov teorem o modulnom bazisu & \hyperlink{TERM-0815}{TERM-0815}\\
Hilbertov teorem o nerazložimosti & \hyperlink{TERM-0810}{TERM-0810}\\
Hilbertova teorema o nerazložimosti & \hyperlink{TERM-0811}{TERM-0811}\\
Hilbertove čisla & \hyperlink{TERM-0808}{TERM-0808}\\
Hilbertove funkcije & \hyperlink{TERM-0807}{TERM-0807}\\
Hilbertovo tvrđenje & \hyperlink{TERM-0809}{TERM-0809}\\
Hilbertovy teorem I o modulnoj bazi & \hyperlink{TERM-0813}{TERM-0813}\\
hiperkompleksna sistema & \hyperlink{TERM-0843}{TERM-0843}\\
hiperkompleksne sistemy, algebry & \hyperlink{TERM-0842}{TERM-0842}\\
hiperkompleksne veličiny & \hyperlink{TERM-0841}{TERM-0841}\\
hiperkompleksny sistem & \hyperlink{TERM-0844}{TERM-0844}\\
homogena forma & \hyperlink{TERM-0823}{TERM-0823}\\
homogena oblast cělosti & \hyperlink{TERM-0833}{TERM-0833}\\
homogena rezultanta členov najvyššego stepena & \hyperlink{TERM-0831}{TERM-0831}\\
homogene formy / nehomogene polinomy & \hyperlink{TERM-0827}{TERM-0827}\\
homogene formy / nehomogenny polinomy & \hyperlink{TERM-0828}{TERM-0828}\\
homogene formy i-toj dimenzije & \hyperlink{TERM-0829}{TERM-0829}\\
homogene polinomy & \hyperlink{TERM-0830}{TERM-0830}\\
homogenna forma & \hyperlink{TERM-0824}{TERM-0824}\\
homogenna forma f(dx) & \hyperlink{TERM-0825}{TERM-0825}\\
homogenny formy & \hyperlink{TERM-0826}{TERM-0826}\\
homogenny izobarny invarianty & \hyperlink{TERM-0832}{TERM-0832}\\
homomorfizm centra & \hyperlink{TERM-0840}{TERM-0840}\\
homomorfizmy iz centra v Omega & \hyperlink{TERM-0839}{TERM-0839}\\
homomorfizmy prehodet v izomorfizmy & \hyperlink{TERM-0838}{TERM-0838}\\
homomorfno matrično kolco ... stvorjeno idealom & \hyperlink{TERM-0836}{TERM-0836}\\
homomorfny / modulno-homomorfny / kolcevo-homomorfny & \hyperlink{TERM-0835}{TERM-0835}\\
Hurwitz v citaciji, Hurvica v prozě & \hyperlink{DEC-0160}{DEC-0160}\\
i podobne & \hyperlink{TERM-1802}{TERM-1802}\\
ideal / glavny ideal / prosty ideal / primarny ideal & \hyperlink{TERM-0855}{TERM-0855}\\
ideal / modul / modul form & \hyperlink{TERM-0856}{TERM-0856}\\
ideal dvojnyh toček & \hyperlink{TERM-0345}{TERM-0345}\\
ideal iz polinomov & \hyperlink{TERM-0857}{TERM-0857}\\
ideal razny od nulovogo ideala & \hyperlink{TERM-1858}{TERM-1858}\\
ideal razširenja & \hyperlink{TERM-0520}{TERM-0520}\\
ideal razvětvenja (osnovny ideal, diferenta) & \hyperlink{TERM-1820}{TERM-1820}\\
ideal suženja & \hyperlink{TERM-1808}{TERM-1808}\\
ideal-dělitelj nuly & \hyperlink{TERM-1270}{TERM-1270}\\
idealna baza & \hyperlink{TERM-0859}{TERM-0859}\\
idealna klasa; glavna klasa; glavne idealy & \hyperlink{TERM-0863}{TERM-0863}\\
idealne kvocienty / glavny ideal / prost vzhodno k t & \hyperlink{TERM-0868}{TERM-0868}\\
idealny kvocient & \hyperlink{TERM-0865}{TERM-0865}\\
idealny kvocient nulovogo ideala črez a & \hyperlink{TERM-0866}{TERM-0866}\\
idealny kvocient nulovogo ideala črez c & \hyperlink{TERM-1419}{TERM-1419}\\
Idealtheorie, §7, 3 & \hyperlink{TERM-0870}{TERM-0870}\\
Idealtheorie, §9 i §3, 1 & \hyperlink{TERM-0872}{TERM-0872}\\
Idealtheorie, §9, 7 & \hyperlink{TERM-0871}{TERM-0871}\\
idealy različne od nulovogo i jediničnogo idealov & \hyperlink{TERM-1856}{TERM-1856}\\
idealy razvětvenja & \hyperlink{TERM-1821}{TERM-1821}\\
idempotentne elementy & \hyperlink{TERM-0875}{TERM-0875}\\
idempotentny element & \hyperlink{TERM-0877}{TERM-0877}\\
idempotenty; komplementarne bazy & \hyperlink{TERM-0876}{TERM-0876}\\
identične razklady & \hyperlink{TERM-0879}{TERM-0879}\\
identično v Y & \hyperlink{TERM-0878}{TERM-0878}\\
identičnost medžu polinomami & \hyperlink{TERM-0881}{TERM-0881}\\
identičnost polj & \hyperlink{TERM-0880}{TERM-0880}\\
identičnosti ne nastavajut & \hyperlink{TERM-0882}{TERM-0882}\\
Iduči i poprědny & \hyperlink{DEC-0111}{DEC-0111}\\
Idučih: otnošenje k Gaussu & \hyperlink{DEC-0113}{DEC-0113}\\
imaje rang jedan & \hyperlink{TERM-1859}{TERM-1859}\\
Ime v kurzivě ostaje cělo ime & \hyperlink{DEC-0159}{DEC-0159}\\
Imena v prěkladanoj prozě & \hyperlink{DEC-0158}{DEC-0158}\\
Imenno i ime & \hyperlink{DEC-0146}{DEC-0146}\\
Imenno: pojasnjenje termina & \hyperlink{DEC-0068}{DEC-0068}\\
imenovatelj & \hyperlink{TERM-1206}{TERM-1206}\\
Indeks pri koeficientu & \hyperlink{DEC-0138}{DEC-0138}\\
indeks vagi & \hyperlink{TERM-0713}{TERM-0713}\\
indeksy vagy & \hyperlink{TERM-0714}{TERM-0714}\\
Individualna norma i i zato & \hyperlink{DEC-0130}{DEC-0130}\\
indukcija & \hyperlink{TERM-0896}{TERM-0896}\\
indukovane transformacije / totalna diferencialna forma & \hyperlink{TERM-0898}{TERM-0898}\\
indukovano prěobraženje & \hyperlink{TERM-0897}{TERM-0897}\\
inercijno polje / razvětvenno polje & \hyperlink{TERM-1754}{TERM-1754}\\
Inicijaly i partikuly v imenah & \hyperlink{DEC-0161}{DEC-0161}\\
integralna baza & \hyperlink{TERM-0905}{TERM-0905}\\
integralna oblast & \hyperlink{TERM-0906}{TERM-0906}\\
integralna oblast iz polinomov & \hyperlink{TERM-0912}{TERM-0912}\\
integralna zamknjenost v polju kvocientov & \hyperlink{TERM-0637}{TERM-0637}\\
integralna zamknutost o v Omega & \hyperlink{TERM-0638}{TERM-0638}\\
integralne oblasti & \hyperlink{TERM-0913}{TERM-0913}\\
integralne oblasti iz polinomov & \hyperlink{TERM-0914}{TERM-0914}\\
integralne oblasti iz relativno cělyh funkcij & \hyperlink{TERM-0915}{TERM-0915}\\
integralno zamknjene vzhodno k svojemu polju kvocientov & \hyperlink{TERM-0074}{TERM-0074}\\
integralno zamknjeno / integralna zamknjenost & \hyperlink{TERM-0636}{TERM-0636}\\
integralno zamknjeny / integralna zamknjenost & \hyperlink{TERM-0630}{TERM-0630}\\
integralnost predstavjenja & \hyperlink{TERM-0681}{TERM-0681}\\
integraly sistema diferencialnyh uravnenj & \hyperlink{TERM-0901}{TERM-0901}\\
integritetna oblast / kolco bez nulovyh děliteljev & \hyperlink{TERM-0908}{TERM-0908}\\
invarianta grupy / invariantny integral & \hyperlink{TERM-0917}{TERM-0917}\\
invariantna variacijska zadača & \hyperlink{TERM-0923}{TERM-0923}\\
invariantne diferencialne operacije & \hyperlink{TERM-0918}{TERM-0918}\\
invariantne procesy diferenciranja & \hyperlink{TERM-0919}{TERM-0919}\\
invariantne tvorjenja & \hyperlink{TERM-0916}{TERM-0916}\\
invariantne variacijske problemy & \hyperlink{TERM-0920}{TERM-0920}\\
invariantno polje & \hyperlink{TERM-0922}{TERM-0922}\\
invarianty libovoljnyh diferencialnyh izrazov & \hyperlink{TERM-0921}{TERM-0921}\\
involucija v linearnom prostoru & \hyperlink{TERM-0924}{TERM-0924}\\
involucijska baza & \hyperlink{TERM-0925}{TERM-0925}\\
involucijska forma & \hyperlink{TERM-0926}{TERM-0926}\\
ireducibilna funkcija / ireducibilne faktory & \hyperlink{TERM-1345}{TERM-1345}\\
ireducibilne klasy predstavjenj & \hyperlink{TERM-0934}{TERM-0934}\\
ireducibilne sostavne česti & \hyperlink{TERM-0932}{TERM-0932}\\
iteracija & \hyperlink{TERM-0963}{TERM-0963}\\
iz predstavjenja v obliku produkta stepenjev & \hyperlink{TERM-0148}{TERM-0148}\\
iz punktov 3 i 4 slěduje & \hyperlink{TERM-0147}{TERM-0147}\\
iz?etna mno?ina / oblasti G\_nu & \hyperlink{TERM-0050}{TERM-0050}\\
izčetno beskonečno mnogo klasov & \hyperlink{TERM-0052}{TERM-0052}\\
izčisljenje ostatkovyh klasov idealov & \hyperlink{TERM-0053}{TERM-0053}\\
izčrpane & \hyperlink{TERM-0506}{TERM-0506}\\
izhodna forma & \hyperlink{TERM-0152}{TERM-0152}\\
izključene prvočisla / izključene proste idealy & \hyperlink{TERM-0155}{TERM-0155}\\
izolovana komponenta & \hyperlink{TERM-0950}{TERM-0950}\\
izolovana komponenta razloženja & \hyperlink{TERM-0951}{TERM-0951}\\
izolovane idealy & \hyperlink{TERM-0949}{TERM-0949}\\
izolovane najveće primarne idealy & \hyperlink{TERM-0948}{TERM-0948}\\
izolovany ideal / neizolovany ideal & \hyperlink{TERM-0952}{TERM-0952}\\
izomorfija / izomorfno prirěditi & \hyperlink{TERM-0960}{TERM-0960}\\
izomorfizm kolc klasov ostatkov & \hyperlink{TERM-0961}{TERM-0961}\\
izomorfna prirědba / nulovy klas / dělitelji nuly & \hyperlink{TERM-0956}{TERM-0956}\\
izomorfne / izomorfna sootvětnost / izomorfno odobraženje & \hyperlink{TERM-0953}{TERM-0953}\\
izomorfne kako R-moduly & \hyperlink{TERM-0145}{TERM-0145}\\
izomorfno odobraženje & \hyperlink{TERM-0955}{TERM-0955}\\
izomorfny / izomorfija & \hyperlink{TERM-0954}{TERM-0954}\\
Izpolniti, izpolnjati, izpolnjenje & \hyperlink{DEC-0031}{DEC-0031}\\
Izpolnjati uslovje & \hyperlink{DEC-0038}{DEC-0038}\\
Izpolnjati: objekt i padež & \hyperlink{DEC-0054}{DEC-0054}\\
izpolnjenje redukcij & \hyperlink{TERM-0151}{TERM-0151}\\
izvedena funkcionalna oblast / polje kvocientov & \hyperlink{TERM-0029}{TERM-0029}\\
izvedeny glavny ideal & \hyperlink{TERM-0032}{TERM-0032}\\
izvodženje formul & \hyperlink{TERM-0035}{TERM-0035}\\
Izvor mora byti dohvatny & \hyperlink{DEC-0145}{DEC-0145}\\
Javno, jasno i jednoznačno & \hyperlink{DEC-0143}{DEC-0143}\\
jedine koeficienty pri t v |tE-C| & \hyperlink{TERM-0457}{TERM-0457}\\
jedini?ny klas & \hyperlink{TERM-0436}{TERM-0436}\\
jedinica iz [Omega] & \hyperlink{TERM-0426}{TERM-0426}\\
jedinice & \hyperlink{TERM-0427}{TERM-0427}\\
jedinice grupy ostatkov / ra?unanje s jedinicami & \hyperlink{TERM-0428}{TERM-0428}\\
jedinična matrica & \hyperlink{TERM-0438}{TERM-0438}\\
jediničny element & \hyperlink{TERM-0429}{TERM-0429}\\
jediničny element e / modulna baza & \hyperlink{TERM-0430}{TERM-0430}\\
jediničny ideal & \hyperlink{TERM-0431}{TERM-0431}\\
jediničny ideal / jediničny element & \hyperlink{TERM-0432}{TERM-0432}\\
jediničny ideal / nulovy ideal / \textbackslash{}frako kako jediničny ideal & \hyperlink{TERM-0433}{TERM-0433}\\
jediničny ideal \textbackslash{}mo / jediničny element / jediničny klas & \hyperlink{TERM-0434}{TERM-0434}\\
jediničny ideal D & \hyperlink{TERM-0435}{TERM-0435}\\
jediničny klas / jedinica & \hyperlink{TERM-0437}{TERM-0437}\\
jediničny modul & \hyperlink{TERM-0439}{TERM-0439}\\
jediničny operator & \hyperlink{TERM-0440}{TERM-0440}\\
jednačenja "bez afekta" & \hyperlink{TERM-0066}{TERM-0066}\\
Jednako: protivopostavjenje & \hyperlink{DEC-0137}{DEC-0137}\\
jednakogo vida & \hyperlink{TERM-1860}{TERM-1860}, \hyperlink{TERM-1861}{TERM-1861}\\
Jednaky tekst v samostojnoj i zbirnoj formě & \hyperlink{DEC-0126}{DEC-0126}\\
Jedno izdanje, ne směs dvuh redakcij & \hyperlink{DEC-0132}{DEC-0132}\\
jedno- i dvojnokratna svijanja & \hyperlink{TERM-0390}{TERM-0390}\\
jedno-jednozna?ny / jedno-jednozna?no prir??enje & \hyperlink{TERM-0411}{TERM-0411}\\
jedno-jednoznačno & \hyperlink{TERM-0406}{TERM-0406}\\
jedno-jednoznačno / jedno-jednoznačno pridružiti & \hyperlink{TERM-0410}{TERM-0410}\\
jedno-jednoznačno i izomorfno odobraziti & \hyperlink{TERM-0414}{TERM-0414}\\
jedno-jednoznačno odgovarjanje & \hyperlink{TERM-0416}{TERM-0416}\\
jedno-jednoznačno pridruženje / izomorfny & \hyperlink{TERM-0415}{TERM-0415}\\
jednoelementny / mnogočlenny bazis modula & \hyperlink{TERM-0425}{TERM-0425}\\
jednoformny slučaj & \hyperlink{TERM-0424}{TERM-0424}\\
Jednoj: ženska forma čislovnika jedin & \hyperlink{DEC-0147}{DEC-0147}\\
jednokratna transvekcija & \hyperlink{TERM-0444}{TERM-0444}\\
jednokratno svijanje & \hyperlink{TERM-0442}{TERM-0442}\\
Jednolike formy i kontekstne izključki & \hyperlink{DEC-0008}{DEC-0008}\\
jednorazovo svijanje & \hyperlink{TERM-0443}{TERM-0443}\\
jednorazovy prěnos & \hyperlink{TERM-0445}{TERM-0445}\\
jednoredna matrica & \hyperlink{TERM-0453}{TERM-0453}\\
jednostavne idealy & \hyperlink{TERM-0420}{TERM-0420}\\
jednostranne / dvustranne idealy & \hyperlink{TERM-0449}{TERM-0449}\\
jednostranne moduly & \hyperlink{TERM-0450}{TERM-0450}\\
jednostranno prosty ideal & \hyperlink{TERM-0448}{TERM-0448}\\
jednozna?no opred?ljeno preslikanje / obratno jednozna?no & \hyperlink{TERM-0396}{TERM-0396}\\
jednoznačna sootvětnost / jednoznačnost & \hyperlink{TERM-0401}{TERM-0401}\\
jednoznačno byti predstavjen & \hyperlink{TERM-0395}{TERM-0395}\\
jednoznačno oprěděljene / mnogoznačnost & \hyperlink{TERM-0394}{TERM-0394}\\
jednoznačno produktno predstavjenje črez nerazložime polinomy oblasti & \hyperlink{TERM-0399}{TERM-0399}\\
jednoznačno razloženje & \hyperlink{TERM-0400}{TERM-0400}\\
jednoznačno vzajemna sootvětnost & \hyperlink{TERM-0417}{TERM-0417}\\
jednoznačno vzajemno ... razlika i množenje tym samym elementom & \hyperlink{TERM-0409}{TERM-0409}\\
jednoznačnost & \hyperlink{TERM-0398}{TERM-0398}\\
jednoznačny additivny razklad grupy ostatkov & \hyperlink{TERM-0397}{TERM-0397}\\
Ješče & \hyperlink{DEC-0007}{DEC-0007}\\
Ješče i ide & \hyperlink{DEC-0117}{DEC-0117}\\
jest algebraično cěla nad & \hyperlink{TERM-0075}{TERM-0075}\\
jest algebraično cěla nad T & \hyperlink{TERM-0095}{TERM-0095}\\
jest algebraično cěly nad & \hyperlink{TERM-0076}{TERM-0076}, \hyperlink{TERM-0077}{TERM-0077}\\
jest dělimy kvadratom prostogo ideala & \hyperlink{TERM-0355}{TERM-0355}\\
jest ispolnjen diskriminantny teorem (§7, 3) & \hyperlink{TERM-0332}{TERM-0332}\\
Jezyk bibliografije ne zavisi od jezyka knigy & \hyperlink{DEC-0164}{DEC-0164}\\
Jih: čije objekty? & \hyperlink{DEC-0102}{DEC-0102}\\
Jordan–Hölder: citovano ime teoremy & \hyperlink{DEC-0165}{DEC-0165}\\
K aksiomatice ... Idealtheorie & \hyperlink{TERM-1954}{TERM-1954}\\
K čemu odnosi se citacija teorema & \hyperlink{DEC-0129}{DEC-0129}\\
K jest razsmatrjany kako Omega-modul & \hyperlink{TERM-0968}{TERM-0968}\\
K\_i konjugovane polja v Galoisovom razširjenom polju Q & \hyperlink{TERM-0971}{TERM-0971}\\
k\_i-rędne determinanty & \hyperlink{TERM-0972}{TERM-0972}\\
kako nerazložima veličina & \hyperlink{TERM-0132}{TERM-0132}\\
kako produkty sistemov I i II ostavaju & \hyperlink{TERM-0130}{TERM-0130}\\
karakter & \hyperlink{TERM-0202}{TERM-0202}\\
karakter / sootvětny nerazložimy karakter & \hyperlink{TERM-0203}{TERM-0203}\\
karaktere & \hyperlink{TERM-0204}{TERM-0204}\\
karaktere / rang & \hyperlink{TERM-0205}{TERM-0205}\\
karakteristična funkcija / tvorjača funkcija & \hyperlink{TERM-0213}{TERM-0213}\\
karakterističny rjad & \hyperlink{TERM-0215}{TERM-0215}\\
karakteristika & \hyperlink{TERM-0207}{TERM-0207}\\
karakteristika 0 / p & \hyperlink{TERM-0208}{TERM-0208}\\
karakteristika nul & \hyperlink{TERM-0201}{TERM-0201}\\
karakteristika nul / prosto polje & \hyperlink{TERM-0210}{TERM-0210}\\
karakteristika nula & \hyperlink{TERM-0209}{TERM-0209}\\
Kde & \hyperlink{DEC-0006}{DEC-0006}\\
Kde, nehaj, kako: funkcija ostaje važna & \hyperlink{DEC-0018}{DEC-0018}\\
klas idealov / klas predstavjenj & \hyperlink{TERM-0862}{TERM-0862}\\
klasa grupovogo elementa & \hyperlink{TERM-0978}{TERM-0978}\\
klasa ostatkov / ostatkova klasa & \hyperlink{TERM-1549}{TERM-1549}\\
klasa R\_i / podmnožstvo R\_i\textasciicircum{}* & \hyperlink{TERM-0979}{TERM-0979}\\
klasno polje / prim-formy & \hyperlink{TERM-0982}{TERM-0982}\\
klasovy invariant & \hyperlink{TERM-0981}{TERM-0981}\\
klasy idealov & \hyperlink{TERM-0864}{TERM-0864}\\
klasy nerazložimyh idealov & \hyperlink{TERM-1797}{TERM-1797}\\
klasy ostatkov & \hyperlink{TERM-1552}{TERM-1552}\\
klasy ostatkov / grupy ostatkov & \hyperlink{TERM-1553}{TERM-1553}\\
klasy ostatkov; normovanje klasov ostatkov; normovany polinom & \hyperlink{TERM-1555}{TERM-1555}\\
klasy podobnyh algeber; grupa klasov & \hyperlink{TERM-0980}{TERM-0980}\\
klasy predstavjenj & \hyperlink{TERM-0236}{TERM-0236}\\
Klebš--Gordanovo razloženje po redah & \hyperlink{TERM-0216}{TERM-0216}\\
knjigy, sostavjene s učestjem Emmy Noether & \hyperlink{TERM-0200}{TERM-0200}\\
koeficient pri (-t\textasciicircum{}\{s-1\}) & \hyperlink{TERM-0994}{TERM-0994}\\
koeficientna oblast Omega & \hyperlink{TERM-0999}{TERM-0999}\\
koeficientno podružje & \hyperlink{TERM-0996}{TERM-0996}\\
koeficientno polje Ω & \hyperlink{TERM-1000}{TERM-1000}\\
koeficienty kako simetrične funkcije integralov & \hyperlink{TERM-0995}{TERM-0995}\\
kogredientna transformacija & \hyperlink{TERM-1008}{TERM-1008}\\
kogredientne prěměnne & \hyperlink{TERM-1010}{TERM-1010}\\
kogredientne varijable & \hyperlink{TERM-1009}{TERM-1009}\\
kogredientno svijanje & \hyperlink{TERM-1006}{TERM-1006}\\
kogredientno svijanje s defektom rho & \hyperlink{TERM-1007}{TERM-1007}\\
kogredientnost & \hyperlink{TERM-1011}{TERM-1011}\\
kogredientny & \hyperlink{TERM-1003}{TERM-1003}\\
kogredientny / kontragredientny & \hyperlink{TERM-1004}{TERM-1004}, \hyperlink{TERM-1005}{TERM-1005}\\
kolca konečnogo ranga & \hyperlink{TERM-1602}{TERM-1602}\\
kolca ranga jedan & \hyperlink{TERM-1605}{TERM-1605}\\
kolca s jediničnym elementom / kolco s jedinico & \hyperlink{TERM-1603}{TERM-1603}\\
kolca s jediničnym elementom i bez nuljevyh děliteljev & \hyperlink{TERM-1604}{TERM-1604}\\
kolceva oblast & \hyperlink{TERM-1601}{TERM-1601}\\
kolco / ideal / P-modul izvedeny iz S v T & \hyperlink{TERM-0149}{TERM-0149}\\
kolco avtomorfizmov & \hyperlink{TERM-0159}{TERM-0159}\\
kolco avtomorfizmov / tělo avtomorfizmov & \hyperlink{TERM-0160}{TERM-0160}\\
kolco bez děliteljev nula i s jediničnym elementom & \hyperlink{TERM-1597}{TERM-1597}\\
kolco bez jediničnogo elementa & \hyperlink{TERM-1594}{TERM-1594}\\
kolco bez nulovyh děliteljev & \hyperlink{TERM-1595}{TERM-1595}\\
kolco bez radikala & \hyperlink{TERM-1598}{TERM-1598}\\
kolco glavnyh idealov & \hyperlink{TERM-0789}{TERM-0789}\\
kolco glavnyh idealov / izhodno kolco & \hyperlink{TERM-0790}{TERM-0790}\\
kolco glavnyh idealov bez děliteljev nuly i s jediniceju & \hyperlink{TERM-0792}{TERM-0792}\\
kolco glavnyh idealov o\_p & \hyperlink{TERM-0791}{TERM-0791}\\
kolco izvedeno iz & \hyperlink{TERM-0030}{TERM-0030}\\
kolco klasov ostatkov & \hyperlink{TERM-1561}{TERM-1561}\\
kolco klasov ostatkov / polje klasov ostatkov & \hyperlink{TERM-1562}{TERM-1562}\\
kolco klasov ostatkov bez děliteljev nuly / dělitelj nuly & \hyperlink{TERM-1565}{TERM-1565}\\
kolco klasov ostatkov modulo p\textasciicircum{}t & \hyperlink{TERM-1563}{TERM-1563}\\
kolco klasov ostatkov po radikalu & \hyperlink{TERM-1564}{TERM-1564}\\
kolco klasov ostatkov T\_P/T\_P p & \hyperlink{TERM-1567}{TERM-1567}\\
kolco konečnogo ranga vzhodno k podpolju P & \hyperlink{TERM-1599}{TERM-1599}\\
kolco kvocientov & \hyperlink{TERM-1425}{TERM-1425}\\
kolco kvocientov T\_a & \hyperlink{TERM-1426}{TERM-1426}\\
kolco matric / matrične jedinice & \hyperlink{TERM-1592}{TERM-1592}\\
kolco produkta & \hyperlink{TERM-1393}{TERM-1393}\\
kolco raz?irenja; modul raz?irenja & \hyperlink{TERM-0532}{TERM-0532}\\
kolco v abstraktnoj definiciji & \hyperlink{TERM-1593}{TERM-1593}\\
kolco vsih cělyh algebraičnyh čisel & \hyperlink{TERM-1591}{TERM-1591}\\
kolcova adjunkcija & \hyperlink{TERM-1600}{TERM-1600}\\
kolcove operacije iz elementov kolca & \hyperlink{TERM-1607}{TERM-1607}\\
kolcovy homomorfizm od R k R\_a & \hyperlink{TERM-1606}{TERM-1606}\\
kombinatoričny/progresivny/regresivny/směšany produkt & \hyperlink{TERM-1012}{TERM-1012}\\
kompleksne koordinaty / linijne koordinaty & \hyperlink{TERM-1025}{TERM-1025}\\
komplement & \hyperlink{TERM-1018}{TERM-1018}\\
komplement / ortogonalne relacije / direktny prěsěk & \hyperlink{TERM-1019}{TERM-1019}\\
komplement Q & \hyperlink{TERM-1020}{TERM-1020}\\
komplementarna baza / komplementarny modul & \hyperlink{TERM-1023}{TERM-1023}\\
komplementarne bazy / komplementarne moduly & \hyperlink{TERM-1022}{TERM-1022}\\
komplementarny modul & \hyperlink{TERM-1024}{TERM-1024}\\
komponenta & \hyperlink{TERM-1026}{TERM-1026}\\
komponentne klasy & \hyperlink{TERM-1030}{TERM-1030}\\
komponentny klas & \hyperlink{TERM-1029}{TERM-1029}\\
komponenty kako kolca same po sebě, a ne kako idealy od \textbackslash{}overline\{\textbackslash{}mathfrak R\} & \hyperlink{TERM-1027}{TERM-1027}\\
kompozicijna serija modula predstavjenja & \hyperlink{TERM-1036}{TERM-1036}\\
kompozicijne faktory & \hyperlink{TERM-1031}{TERM-1031}\\
kompozicijny red & \hyperlink{TERM-1032}{TERM-1032}\\
kompozicijny red / kompozicijny faktor & \hyperlink{TERM-1034}{TERM-1034}\\
kompozicijny rjad & \hyperlink{TERM-1033}{TERM-1033}\\
kompozicijny rjad / kompozicijny faktor & \hyperlink{TERM-1035}{TERM-1035}\\
kompozitum polj & \hyperlink{TERM-1807}{TERM-1807}\\
komutativno polje K & \hyperlink{TERM-1015}{TERM-1015}\\
komutativno povezano s ... & \hyperlink{TERM-1013}{TERM-1013}\\
komutativno povezany s P & \hyperlink{TERM-1014}{TERM-1014}\\
komutativno tělo P & \hyperlink{TERM-1016}{TERM-1016}\\
komutativny zakon / zam?nljivost & \hyperlink{TERM-1017}{TERM-1017}\\
konečna / beskonečna grupa ostatkov & \hyperlink{TERM-0476}{TERM-0476}\\
konečna / beskonečna kontinuirna grupa & \hyperlink{TERM-0475}{TERM-0475}\\
konečna Abelova grupa & \hyperlink{TERM-0477}{TERM-0477}\\
konečna formula & \hyperlink{TERM-1626}{TERM-1626}\\
konečna grupa & \hyperlink{TERM-0480}{TERM-0480}\\
konečna integralna oblast & \hyperlink{TERM-0487}{TERM-0487}\\
konečna linearna grupa karakteristiki p & \hyperlink{TERM-0481}{TERM-0481}\\
konečne izrazy / početne izrazy & \hyperlink{TERM-1624}{TERM-1624}\\
konečne porjadky algebraičnogo čislovogo polja & \hyperlink{TERM-0483}{TERM-0483}\\
konečne vrědnosti najdene rekursivno & \hyperlink{TERM-1514}{TERM-1514}\\
konečno algebraično čislovo polje & \hyperlink{TERM-0486}{TERM-0486}\\
konečno algebraično razširenje & \hyperlink{TERM-0478}{TERM-0478}\\
konečno algebraično razširenje nad & \hyperlink{TERM-0004}{TERM-0004}\\
konečno algebraično razširjenje polja & \hyperlink{TERM-0485}{TERM-0485}\\
konečno mnogo ostatkovyh klasov linearno nezavisnyh & \hyperlink{TERM-0474}{TERM-0474}\\
konečnost sistema form & \hyperlink{TERM-0490}{TERM-0490}\\
konečnost v užšem smyslu & \hyperlink{TERM-0491}{TERM-0491}\\
konečny bazis diferencialnyh izrazov & \hyperlink{TERM-0479}{TERM-0479}\\
konečny ideal / idealny bazis & \hyperlink{TERM-0489}{TERM-0489}\\
konečny p-porjadok & \hyperlink{TERM-0484}{TERM-0484}\\
konečny porjadok & \hyperlink{TERM-0482}{TERM-0482}\\
konečny, to jest konečno porodženy modul & \hyperlink{TERM-0488}{TERM-0488}\\
Kongruencija, ne nekongruencija & \hyperlink{DEC-0140}{DEC-0140}\\
kongruentny po modulu & \hyperlink{TERM-1037}{TERM-1037}\\
Kongruenzkomposition & \hyperlink{TERM-1038}{TERM-1038}\\
konjugacija črez P & \hyperlink{TERM-1304}{TERM-1304}\\
konjugacijna klasa & \hyperlink{TERM-1045}{TERM-1045}\\
konjugat / konjugaty vzhodno k (H) & \hyperlink{TERM-1040}{TERM-1040}\\
konjugaty & \hyperlink{TERM-1039}{TERM-1039}\\
konjugaty imenovatelja & \hyperlink{TERM-1041}{TERM-1041}\\
konjugovane elementy & \hyperlink{TERM-1042}{TERM-1042}\\
konjugovane faktory & \hyperlink{TERM-0220}{TERM-0220}\\
konjugovane redy veličin & \hyperlink{TERM-1044}{TERM-1044}\\
konjugovane veličiny & \hyperlink{TERM-1043}{TERM-1043}\\
konstanty & \hyperlink{TERM-1046}{TERM-1046}\\
konstrukcija f(z) & \hyperlink{TERM-1048}{TERM-1048}\\
konstrukcija po zadannom dobrom uporjadkovanju & \hyperlink{TERM-1047}{TERM-1047}\\
konstrukcije & \hyperlink{TERM-0188}{TERM-0188}\\
kontinuirna grupa v smislu Lija & \hyperlink{TERM-1049}{TERM-1049}\\
kontragradientny / kogradientny & \hyperlink{TERM-1050}{TERM-1050}\\
kontragredientne formy & \hyperlink{TERM-1051}{TERM-1051}\\
kontragredientne simboly & \hyperlink{TERM-1052}{TERM-1052}\\
Kontrakcija; ne konvolucija & \hyperlink{DEC-0042}{DEC-0042}\\
kontravariant/kovariant & \hyperlink{TERM-1053}{TERM-1053}\\
kontravarianty & \hyperlink{TERM-1054}{TERM-1054}\\
korenevo čislo & \hyperlink{TERM-1887}{TERM-1887}\\
korenj uravnenja & \hyperlink{TERM-1264}{TERM-1264}\\
korenjevo polje & \hyperlink{TERM-1886}{TERM-1886}\\
kovariantna derivacija & \hyperlink{TERM-1056}{TERM-1056}\\
kovariantne derivacije & \hyperlink{TERM-1057}{TERM-1057}\\
krajeva zadacha / vsude konečne integraly & \hyperlink{TERM-1435}{TERM-1435}\\
kratke komunikacije & \hyperlink{TERM-1065}{TERM-1065}\\
kriterij konečnosti & \hyperlink{TERM-0494}{TERM-0494}\\
Kromě: iskľučenje i dopolnjenje & \hyperlink{DEC-0095}{DEC-0095}\\
Kroneckerov fundamentalny sistem & \hyperlink{TERM-1063}{TERM-1063}\\
Kroneckerova substitucija & \hyperlink{TERM-1061}{TERM-1061}\\
kružno polje & \hyperlink{TERM-1059}{TERM-1059}\\
kružno-dělitelno polje & \hyperlink{TERM-1060}{TERM-1060}\\
kvadrat determinanta soprjaženyh bazovyh elementov (§6, 4) & \hyperlink{TERM-0260}{TERM-0260}\\
kvadrat jedinice & \hyperlink{TERM-1407}{TERM-1407}\\
kvadratične relacije & \hyperlink{TERM-1408}{TERM-1408}\\
kvadratične relacije / kvadratične formy & \hyperlink{TERM-1409}{TERM-1409}\\
kvaternionno tělo & \hyperlink{TERM-1411}{TERM-1411}\\
kvaternionno tělo / kvaternionna grupa & \hyperlink{TERM-1413}{TERM-1413}\\
kvaternionsko tělo & \hyperlink{TERM-1412}{TERM-1412}\\
kvazi-leksikografičny porjadok & \hyperlink{TERM-1410}{TERM-1410}\\
kvocient ... črez osnovny ideal & \hyperlink{TERM-1416}{TERM-1416}\\
kvocient ... neopreděljenym = 0/0 & \hyperlink{TERM-1414}{TERM-1414}\\
kvocient a/D\textasciicircum{}\{(i)\} & \hyperlink{TERM-1415}{TERM-1415}\\
kvocient dvoh modulov & \hyperlink{TERM-1417}{TERM-1417}\\
kvocientno polje & \hyperlink{TERM-1420}{TERM-1420}\\
kvocienty idealov & \hyperlink{TERM-0867}{TERM-0867}\\
Lagrangeove izrazy & \hyperlink{TERM-1081}{TERM-1081}\\
Lagrangeva interpolacijska formula & \hyperlink{TERM-1084}{TERM-1084}\\
Lagrangeva rodova oblast & \hyperlink{TERM-1085}{TERM-1085}\\
Lagrangeve derivacije & \hyperlink{TERM-1080}{TERM-1080}\\
Lagrangeve rodove oblasti & \hyperlink{TERM-1083}{TERM-1083}\\
Lagranžev rodovy domen & \hyperlink{TERM-1086}{TERM-1086}\\
Lagranževe izrazy & \hyperlink{TERM-1082}{TERM-1082}\\
Landauove produktne teoremy & \hyperlink{TERM-1394}{TERM-1394}\\
Laplaceovo razloženje & \hyperlink{TERM-1087}{TERM-1087}\\
leksikografičny red & \hyperlink{TERM-1089}{TERM-1089}\\
leksikografičny red / najvyšše členy & \hyperlink{TERM-1088}{TERM-1088}\\
letopisy Njemeckogo matematičnogo družstva & \hyperlink{TERM-0964}{TERM-0964}\\
lěva jedinica & \hyperlink{TERM-1124}{TERM-1124}\\
lěva jedinica / prava jedinica & \hyperlink{TERM-1125}{TERM-1125}\\
lěve idealy / pravy ideal & \hyperlink{TERM-1127}{TERM-1127}\\
lěvy dělitelj nulja & \hyperlink{TERM-1122}{TERM-1122}\\
lěvy gorny kut & \hyperlink{TERM-1123}{TERM-1123}\\
lěvy ideal / pravy ideal / recipročny & \hyperlink{TERM-1126}{TERM-1126}\\
Leži ili ležet & \hyperlink{DEC-0065}{DEC-0065}\\
libovoljna oblast racionalnosti & \hyperlink{TERM-0177}{TERM-0177}\\
libovoljno dobro aproksimovati & \hyperlink{TERM-0175}{TERM-0175}\\
Libovoljny i ljubezno: razne familije & \hyperlink{DEC-0103}{DEC-0103}\\
Libovoljny: libovoljno izbrany objekt & \hyperlink{DEC-0071}{DEC-0071}\\
libovolne veličiny & \hyperlink{TERM-1878}{TERM-1878}\\
libovolny sistem & \hyperlink{TERM-0178}{TERM-0178}\\
linearna / racionalna / algebraična baza & \hyperlink{TERM-1096}{TERM-1096}\\
linearna baza & \hyperlink{TERM-1099}{TERM-1099}\\
linearna baza grupy ostatkov & \hyperlink{TERM-1100}{TERM-1100}\\
linearna familija & \hyperlink{TERM-1105}{TERM-1105}\\
linearna familija form & \hyperlink{TERM-1101}{TERM-1101}\\
linearna mnogovidnost & \hyperlink{TERM-1104}{TERM-1104}\\
linearna nezavisimost / definicija ravnosti & \hyperlink{TERM-1107}{TERM-1107}\\
linearna nezavisnost (alpha\_i) & \hyperlink{TERM-1121}{TERM-1121}\\
linearna transformacija & \hyperlink{TERM-1106}{TERM-1106}\\
linearna zavisnost / relacija & \hyperlink{TERM-1097}{TERM-1097}\\
linearne faktory & \hyperlink{TERM-1112}{TERM-1112}\\
linearne faktory / produktno predstavjenje & \hyperlink{TERM-1113}{TERM-1113}\\
linearne faktory; odnorodnjene korenji & \hyperlink{TERM-1114}{TERM-1114}\\
linearne formy & \hyperlink{TERM-1115}{TERM-1115}\\
linearne kombinacije & \hyperlink{TERM-1098}{TERM-1098}, \hyperlink{TERM-1109}{TERM-1109}\\
linearne sistemy jednačenj & \hyperlink{TERM-1103}{TERM-1103}\\
linearne snope form & \hyperlink{TERM-1102}{TERM-1102}\\
linearno nezavisna modulna baza / klasy ostatkov / kongruencija & \hyperlink{TERM-1093}{TERM-1093}\\
linearno nezavisna Z-modulna baza & \hyperlink{TERM-1095}{TERM-1095}\\
linearno nezavisne medžu soboju i od ... & \hyperlink{TERM-1092}{TERM-1092}\\
linearno nezavisne relacije & \hyperlink{TERM-1094}{TERM-1094}\\
linearno nezavisny vzhodno k P & \hyperlink{TERM-1091}{TERM-1091}\\
linearno nezavisnyh nad Omega elementov & \hyperlink{TERM-0891}{TERM-0891}\\
linearno sočinenje / linearna forma & \hyperlink{TERM-1108}{TERM-1108}\\
linearno-drobna transformacija & \hyperlink{TERM-1120}{TERM-1120}\\
linearno-drobna zavisnost & \hyperlink{TERM-1090}{TERM-1090}\\
linearno-drobno & \hyperlink{TERM-1118}{TERM-1118}\\
linearny faktor & \hyperlink{TERM-1111}{TERM-1111}\\
Ljubezno: priznanije čužoj pomoci & \hyperlink{DEC-0104}{DEC-0104}\\
lu? & \hyperlink{TERM-1662}{TERM-1662}\\
lu?ne klasy & \hyperlink{TERM-1663}{TERM-1663}\\
Lürothova funkcija & \hyperlink{TERM-1130}{TERM-1130}\\
m?sta razv?tvjenja & \hyperlink{TERM-1822}{TERM-1822}\\
m?sto & \hyperlink{TERM-1659}{TERM-1659}\\
maksimalna oblast & \hyperlink{TERM-1146}{TERM-1146}\\
maksimalne porjadky & \hyperlink{TERM-1148}{TERM-1148}\\
maksimalny nilpotentny ideal & \hyperlink{TERM-1147}{TERM-1147}\\
maksimalny porjadok & \hyperlink{TERM-0796}{TERM-0796}\\
matematične proměnne i rimske označenja ostavajut zaštićene & \hyperlink{TERM-0221}{TERM-0221}\\
matematične tablice i moduly ostavajut v matematičnom režimu & \hyperlink{TERM-1207}{TERM-1207}\\
matematične tablice ostavajut v matematičnom režimu & \hyperlink{TERM-1208}{TERM-1208}\\
matrice & \hyperlink{TERM-1142}{TERM-1142}\\
matrice nad K & \hyperlink{TERM-1143}{TERM-1143}\\
matrične jedinice & \hyperlink{TERM-1135}{TERM-1135}\\
matrične kolca homomorfne k R & \hyperlink{TERM-1937}{TERM-1937}\\
matrične produkty & \hyperlink{TERM-1139}{TERM-1139}\\
matrično kolco & \hyperlink{TERM-1140}{TERM-1140}\\
matrično kolco, izomorfno ... pri tom ekvivalentno & \hyperlink{TERM-0957}{TERM-0957}\\
matrično predstavjenje / matrične predstavjenja & \hyperlink{TERM-1134}{TERM-1134}\\
matričny produkt & \hyperlink{TERM-1136}{TERM-1136}\\
matričny produkt / produktna matrica & \hyperlink{TERM-1137}{TERM-1137}\\
matričny produkt tretjega vida & \hyperlink{TERM-1138}{TERM-1138}\\
medžu soboju različne izomorfizmy & \hyperlink{TERM-1783}{TERM-1783}\\
medžukolco & \hyperlink{TERM-1982}{TERM-1982}\\
medžupolja / medžupolj & \hyperlink{TERM-1979}{TERM-1979}\\
medžupolje & \hyperlink{TERM-1980}{TERM-1980}\\
medžupolje / polje prěsěka & \hyperlink{TERM-1981}{TERM-1981}\\
medžupredstavjenja & \hyperlink{TERM-1978}{TERM-1978}\\
medžusobno prosty & \hyperlink{TERM-1712}{TERM-1712}\\
město & \hyperlink{TERM-1660}{TERM-1660}\\
mimobežna prjama & \hyperlink{TERM-1879}{TERM-1879}\\
minimalna baza & \hyperlink{TERM-1151}{TERM-1151}\\
minimalna rezolventa & \hyperlink{TERM-1154}{TERM-1154}\\
minimalne razpadne polja kvaternionnogo těla & \hyperlink{TERM-1153}{TERM-1153}\\
minory & \hyperlink{TERM-1784}{TERM-1784}\\
mnogokratnost & \hyperlink{TERM-1193}{TERM-1193}\\
mnogokratnost prěseka / rezultanta-mnogokratnost & \hyperlink{TERM-1628}{TERM-1628}\\
mnogokratny teorem o cěpah / minimalno uslovje & \hyperlink{TERM-1830}{TERM-1830}\\
mnogovidnost & \hyperlink{TERM-1132}{TERM-1132}\\
Mnogyh: množina, ne novy kvantifikator & \hyperlink{DEC-0092}{DEC-0092}\\
množinno-teoretičny prěsěk & \hyperlink{TERM-1149}{TERM-1149}\\
množitelje & \hyperlink{TERM-1190}{TERM-1190}\\
močnosť & \hyperlink{TERM-1194}{TERM-1194}\\
močnosť / kontinuum / dobra uporjadkovanja & \hyperlink{TERM-1195}{TERM-1195}\\
močnosť kontinuuma & \hyperlink{TERM-1196}{TERM-1196}\\
mod $(i,J)$ & \hyperlink{TERM-1161}{TERM-1161}\\
modul & \hyperlink{TERM-1162}{TERM-1162}, \hyperlink{TERM-1163}{TERM-1163}\\
modul / moduly & \hyperlink{TERM-1166}{TERM-1166}\\
modul / osnovna oblast cělosti & \hyperlink{TERM-1164}{TERM-1164}\\
modul / vyši modul & \hyperlink{TERM-1165}{TERM-1165}\\
modul iz linearnyh form & \hyperlink{TERM-1168}{TERM-1168}\\
modul klasov ostatkov / kolco klasov ostatkov / ostatkova klasa & \hyperlink{TERM-1560}{TERM-1560}\\
modul linearnyh form & \hyperlink{TERM-1116}{TERM-1116}\\
modul linearnyh form / modul iz linearnyh form & \hyperlink{TERM-1117}{TERM-1117}\\
modul M\_sigma / modulno svojstvo & \hyperlink{TERM-1170}{TERM-1170}\\
modul ostatkovyh klasov & \hyperlink{TERM-1559}{TERM-1559}\\
modul predstavjenja & \hyperlink{TERM-0237}{TERM-0237}\\
modul predstavjenja / klas modulov & \hyperlink{TERM-0238}{TERM-0238}\\
modul v (Σ,T) & \hyperlink{TERM-1169}{TERM-1169}\\
modul vseh identičnyh relacij & \hyperlink{TERM-1167}{TERM-1167}\\
modulna baza & \hyperlink{TERM-1172}{TERM-1172}\\
modulna baza / fundamentalna forma / fundamentalno uravnenje & \hyperlink{TERM-1174}{TERM-1174}\\
modulna baza / linearno nezavisna modulna baza & \hyperlink{TERM-1175}{TERM-1175}\\
modulna baza iz linearno nezavisnyh elementov & \hyperlink{TERM-1176}{TERM-1176}\\
modulna homomorfnost ... k sistemě slědov (§4, 3) & \hyperlink{TERM-1180}{TERM-1180}\\
modulna oblast & \hyperlink{TERM-1177}{TERM-1177}\\
modulne bazy & \hyperlink{TERM-1171}{TERM-1171}\\
modulne svojstvo & \hyperlink{TERM-1179}{TERM-1179}\\
modulno predstavjenje & \hyperlink{TERM-1178}{TERM-1178}\\
modulny kvocient & \hyperlink{TERM-1183}{TERM-1183}\\
modulny kvocient / idealny kvocient & \hyperlink{TERM-1184}{TERM-1184}\\
modulny produkt & \hyperlink{TERM-1182}{TERM-1182}\\
modulny teorem & \hyperlink{TERM-1185}{TERM-1185}\\
modulo & \hyperlink{TERM-1160}{TERM-1160}\\
moduly & \hyperlink{TERM-1181}{TERM-1181}\\
moduly predstavjenja & \hyperlink{TERM-0239}{TERM-0239}\\
Mogu i mogut: čislo predikata & \hyperlink{DEC-0061}{DEC-0061}\\
možno prisojedinjenje neopreděljenyh & \hyperlink{TERM-0535}{TERM-0535}\\
multiplikacijsko kolco / regularne multiplikacijske kolca & \hyperlink{TERM-1188}{TERM-1188}\\
načelo prěnosa & \hyperlink{TERM-2014}{TERM-2014}\\
najbližši vyši modul & \hyperlink{TERM-1273}{TERM-1273}\\
Najčestěje: čestota, ne veličina & \hyperlink{DEC-0077}{DEC-0077}\\
najkrajše predstavjenje & \hyperlink{TERM-1077}{TERM-1077}\\
najkratše predstavjenja & \hyperlink{TERM-1078}{TERM-1078}\\
Najkratši: najmenša dolžina & \hyperlink{DEC-0108}{DEC-0108}\\
najmene obče mnogokratno & \hyperlink{TERM-0989}{TERM-0989}\\
najmenjše polje, ktore soderži sistem & \hyperlink{TERM-0984}{TERM-0984}\\
najmenša soderžeča integralna oblast & \hyperlink{TERM-0983}{TERM-0983}\\
najmenša soderžeča relativno cěla oblast & \hyperlink{TERM-0986}{TERM-0986}\\
najmenše obče mnogokratno & \hyperlink{TERM-0990}{TERM-0990}\\
najmenše obče mnogokratno / največši obči dělitelj & \hyperlink{TERM-0992}{TERM-0992}\\
najmenše obče mnogokratno stavaje jednako produktu & \hyperlink{TERM-0993}{TERM-0993}\\
najmenše razširjeno polje & \hyperlink{TERM-0987}{TERM-0987}\\
najmenše soderžeče polje & \hyperlink{TERM-0985}{TERM-0985}\\
najmenši obči imenovatelj & \hyperlink{TERM-0988}{TERM-0988}\\
najmenši obči kratnik & \hyperlink{TERM-0991}{TERM-0991}\\
najnižše po koeficientah formy & \hyperlink{TERM-0892}{TERM-0892}\\
najobčejše jednoznačne rěšenja & \hyperlink{TERM-0125}{TERM-0125}\\
najobčejše razloženja po redah & \hyperlink{TERM-0126}{TERM-0126}\\
Najprvo: porędok izlaganja i prvenstvo & \hyperlink{DEC-0087}{DEC-0087}\\
najve??i ob?i d?litelj / pr?s?k & \hyperlink{TERM-0720}{TERM-0720}\\
najve??i ob?i d?litelj ili pr?s?k & \hyperlink{TERM-0721}{TERM-0721}\\
največa komutativna podpolja & \hyperlink{TERM-0758}{TERM-0758}\\
najveće možno invariantno ograničenje & \hyperlink{TERM-0777}{TERM-0777}\\
najveće primarne idealy & \hyperlink{TERM-0759}{TERM-0759}\\
največe primarne komponenty & \hyperlink{TERM-0760}{TERM-0760}\\
najveće primarne komponenty / idealy & \hyperlink{TERM-0763}{TERM-0763}\\
največi obči dělitelj & \hyperlink{TERM-0767}{TERM-0767}\\
najveći obšty dělitelj idealov & \hyperlink{TERM-0773}{TERM-0773}\\
najveći primarny dělitelj / primarny dělitelj & \hyperlink{TERM-0775}{TERM-0775}\\
najveći primarny faktor normy & \hyperlink{TERM-0774}{TERM-0774}\\
največša primarna komponenta & \hyperlink{TERM-0764}{TERM-0764}\\
največše primarne komponenty & \hyperlink{TERM-0761}{TERM-0761}\\
največše primarne komponenty / izolovane komponenty & \hyperlink{TERM-0765}{TERM-0765}\\
največše primarne komponenty / največše primarne faktory & \hyperlink{TERM-0762}{TERM-0762}\\
največše primarne komponenty / primarna komponenta naležeča k \textbackslash{}mo & \hyperlink{TERM-0766}{TERM-0766}\\
največši obči dělitelj & \hyperlink{TERM-0768}{TERM-0768}\\
največši obči dělitelj / najmene obče mnogokratno & \hyperlink{TERM-0727}{TERM-0727}\\
največši obči dělitelj / najmenše obče mnogokratno & \hyperlink{TERM-0770}{TERM-0770}\\
največši obči dělitelj / prěsěk & \hyperlink{TERM-0769}{TERM-0769}\\
največši obči dělitelj jest jednaky jediničnomu idealu D=(ε) & \hyperlink{TERM-0771}{TERM-0771}\\
največši obči dělitelj v polinomnom smyslu & \hyperlink{TERM-0772}{TERM-0772}\\
največši popolno reducibilny modul & \hyperlink{TERM-0776}{TERM-0776}\\
Največši: nepravilno stepenovanje & \hyperlink{DEC-0076}{DEC-0076}\\
najvyšša dimenzija & \hyperlink{TERM-0845}{TERM-0845}\\
najvyšša dimenzija libovoljnogo ideala & \hyperlink{TERM-0846}{TERM-0846}\\
najvyšši elementarny dělitelj & \hyperlink{TERM-0806}{TERM-0806}, \hyperlink{TERM-0847}{TERM-0847}\\
najvyšši koeficient & \hyperlink{TERM-0848}{TERM-0848}\\
najvyšši koeficient jest jedinica & \hyperlink{TERM-0849}{TERM-0849}\\
nauka o razširjenju & \hyperlink{TERM-0150}{TERM-0150}\\
Nazvanje knigy i nomer glavy & \hyperlink{DEC-0166}{DEC-0166}\\
ne imajut nikakogo vlastnogo dělitelja raznogo od jediničnogo ideala & \hyperlink{TERM-0974}{TERM-0974}\\
ne izčezaje identično / identično izčezaje & \hyperlink{TERM-1210}{TERM-1210}\\
ne jest dostatočno obosnovany & \hyperlink{TERM-1212}{TERM-1212}\\
ne jest popolno reducibilno prvogo roda & \hyperlink{TERM-1213}{TERM-1213}\\
ne jest zatvorjena vzhodno k tvorenju pr?s?kov & \hyperlink{TERM-1209}{TERM-1209}\\
ne nužno celočislene funkcije & \hyperlink{TERM-1211}{TERM-1211}\\
ne vhodi v svijanje & \hyperlink{TERM-1750}{TERM-1750}\\
ne zavisi od klasa položenogo v osnovu & \hyperlink{TERM-1764}{TERM-1764}\\
necěly element & \hyperlink{TERM-1219}{TERM-1219}\\
neizolovane (imbedded) primarne idealy & \hyperlink{TERM-1214}{TERM-1214}\\
neizomorfne ili suštno različne & \hyperlink{TERM-1215}{TERM-1215}\\
nekomutativna kolceva oblast / nekomutativno kolco & \hyperlink{TERM-1216}{TERM-1216}\\
nekomutativne kolca ... Sigma komutuje so vsimi elementami & \hyperlink{TERM-1223}{TERM-1223}\\
nekomutativne oblasti & \hyperlink{TERM-1220}{TERM-1220}\\
nekomutativne těla & \hyperlink{TERM-1222}{TERM-1222}\\
nekomutativno tělo & \hyperlink{TERM-1224}{TERM-1224}\\
nekomutativno vstavjenje & \hyperlink{TERM-1221}{TERM-1221}\\
nekomutativny dvustranny slu?aj & \hyperlink{TERM-1225}{TERM-1225}\\
něktore neopreděljene zaměniti čislami & \hyperlink{TERM-0441}{TERM-0441}\\
Něktory i nehaj: něšto jest, ili nehaj bude & \hyperlink{DEC-0010}{DEC-0010}\\
nelinearne identične relacije & \hyperlink{TERM-1226}{TERM-1226}\\
neměrime funkcije & \hyperlink{TERM-1227}{TERM-1227}\\
nemožnost identičnosti & \hyperlink{TERM-1776}{TERM-1776}\\
nenegativne eksponenty & \hyperlink{TERM-1228}{TERM-1228}\\
nenulovost diskriminanta & \hyperlink{TERM-1232}{TERM-1232}\\
nenulovy element / někoja stepenj izčezaje & \hyperlink{TERM-1865}{TERM-1865}\\
nenulovy, skoro tau sut linearno nezavisne & \hyperlink{TERM-1864}{TERM-1864}\\
neobstajanje eksplicitnogo predstavjenja & \hyperlink{TERM-1218}{TERM-1218}\\
neopreděljena & \hyperlink{TERM-1766}{TERM-1766}\\
neopreděljene & \hyperlink{TERM-1767}{TERM-1767}\\
neopreděljene / stati se neopreděljenoju & \hyperlink{TERM-1768}{TERM-1768}\\
neparne čislo defekta & \hyperlink{TERM-1770}{TERM-1770}\\
neparno prosto čislo & \hyperlink{TERM-1771}{TERM-1771}\\
neparny relativny stepen & \hyperlink{TERM-1772}{TERM-1772}\\
neposrědnje izračunanje / neposrědnje razvitje & \hyperlink{TERM-0297}{TERM-0297}\\
neposrědnje svijanje & \hyperlink{TERM-0298}{TERM-0298}\\
neposrědno svijanje s reducentami & \hyperlink{TERM-0299}{TERM-0299}\\
neprimarny & \hyperlink{TERM-1229}{TERM-1229}\\
neprimarny / primarny & \hyperlink{TERM-1230}{TERM-1230}\\
neracionalna involucija & \hyperlink{TERM-1217}{TERM-1217}\\
nerazložima algebraična mnogovidnost & \hyperlink{TERM-0946}{TERM-0946}\\
nerazložima cěla funkcija & \hyperlink{TERM-0939}{TERM-0939}\\
nerazložime dvustranne klasy & \hyperlink{TERM-0945}{TERM-0945}\\
nerazložime funkcije & \hyperlink{TERM-0938}{TERM-0938}\\
nerazložime identičnosti & \hyperlink{TERM-1798}{TERM-1798}\\
nerazložime komponenty & \hyperlink{TERM-1799}{TERM-1799}\\
nerazložime komponenty / podsistemy & \hyperlink{TERM-0942}{TERM-0942}\\
nerazložime komponenty / polja izomorfne s \textbackslash{}overline P & \hyperlink{TERM-1800}{TERM-1800}\\
nerazložime predstavjenja & \hyperlink{TERM-0933}{TERM-0933}\\
nerazložimo jednačenje & \hyperlink{TERM-0940}{TERM-0940}\\
nerazložimo uravnenje vzhodno k Omega\_i & \hyperlink{TERM-0941}{TERM-0941}\\
nerazložimo vzhodno k [H] & \hyperlink{TERM-0885}{TERM-0885}\\
nerazložimy / nerazložimo & \hyperlink{TERM-0929}{TERM-0929}\\
nerazložimy / razložimy & \hyperlink{TERM-0930}{TERM-0930}, \hyperlink{TERM-0931}{TERM-0931}\\
nerazložimy ideal & \hyperlink{TERM-0947}{TERM-0947}\\
nerazložimy polinom & \hyperlink{TERM-1343}{TERM-1343}\\
nerazložimy polinom drugogo roda & \hyperlink{TERM-1349}{TERM-1349}\\
nerazložimy polinom prvogo roda & \hyperlink{TERM-1347}{TERM-1347}\\
nerazložimy polinomy & \hyperlink{TERM-0928}{TERM-0928}\\
nereducibilne formy & \hyperlink{TERM-0935}{TERM-0935}\\
nereducibilne formy / nereducibilna forma & \hyperlink{TERM-0936}{TERM-0936}\\
nereducibilne formy / poręd & \hyperlink{TERM-0937}{TERM-0937}\\
nereduktibilne nerazložime identičnosti & \hyperlink{TERM-0944}{TERM-0944}\\
neregularne sistemy & \hyperlink{TERM-1231}{TERM-1231}\\
Neter- v běgu prozy; Noether v citacijah & \hyperlink{DEC-0001}{DEC-0001}\\
nezavisna modulna baza & \hyperlink{TERM-1765}{TERM-1765}\\
ni??a mnogovidnost & \hyperlink{TERM-1234}{TERM-1234}\\
nilpotentny ideal & \hyperlink{TERM-1236}{TERM-1236}\\
nižše formy rěda form & \hyperlink{TERM-1233}{TERM-1233}\\
nižšego stepena & \hyperlink{TERM-1235}{TERM-1235}\\
Noetherov fundamentalny teorem algebraičnyh funkcij & \hyperlink{TERM-1237}{TERM-1237}\\
Nomery formul: čitatelj mora viděti cěly blok & \hyperlink{DEC-0134}{DEC-0134}\\
norma & \hyperlink{TERM-1238}{TERM-1238}\\
norma / elementarny dělitelj & \hyperlink{TERM-1239}{TERM-1239}\\
norma / obrazovanje normy & \hyperlink{TERM-1240}{TERM-1240}\\
norma G vzhodno k A & \hyperlink{TERM-1243}{TERM-1243}\\
norma osnovnogo ideala vzhodno k m & \hyperlink{TERM-1241}{TERM-1241}\\
norma po tom modulu / po Q & \hyperlink{TERM-1242}{TERM-1242}\\
normalna baza & \hyperlink{TERM-1244}{TERM-1244}\\
normalna dělitvena algebra & \hyperlink{TERM-1245}{TERM-1245}\\
normalna forma & \hyperlink{TERM-1247}{TERM-1247}\\
normalna forma ρ-tej variacije & \hyperlink{TERM-1248}{TERM-1248}\\
normalna prosta algebra & \hyperlink{TERM-1246}{TERM-1246}\\
normalne formy & \hyperlink{TERM-1249}{TERM-1249}\\
normalny dělitelj & \hyperlink{TERM-1250}{TERM-1250}\\
normalny podgrup & \hyperlink{TERM-1251}{TERM-1251}\\
normna teorema Hasse & \hyperlink{TERM-0782}{TERM-0782}\\
normno-ostatkovy symbol & \hyperlink{TERM-1252}{TERM-1252}\\
normny teorem & \hyperlink{TERM-1253}{TERM-1253}\\
normovati ręd form & \hyperlink{TERM-0593}{TERM-0593}\\
nula / jedinica / vlastne oblasti cělosti & \hyperlink{TERM-1255}{TERM-1255}\\
nula / polje nuljev & \hyperlink{TERM-1263}{TERM-1263}\\
nula prostogo ideala & \hyperlink{TERM-1265}{TERM-1265}\\
nulje & \hyperlink{TERM-1266}{TERM-1266}\\
nulje / polje nuljev & \hyperlink{TERM-1267}{TERM-1267}\\
nuljeva točka & \hyperlink{TERM-1262}{TERM-1262}\\
nuljevy operator & \hyperlink{TERM-1261}{TERM-1261}\\
nulovy dělitelj & \hyperlink{TERM-1269}{TERM-1269}\\
nulovy determinant & \hyperlink{TERM-1818}{TERM-1818}\\
nulovy i jediničny ideal & \hyperlink{TERM-1257}{TERM-1257}\\
nulovy ideal & \hyperlink{TERM-1258}{TERM-1258}\\
nulovy ideal / jediničny ideal & \hyperlink{TERM-1259}{TERM-1259}\\
nulovy ideal / prosty ideal & \hyperlink{TERM-1260}{TERM-1260}\\
nulovy ili jediničny ideal & \hyperlink{TERM-1256}{TERM-1256}\\
Nyně & \hyperlink{DEC-0005}{DEC-0005}\\
o\_p imaje samo jedin prosty ideal, različny od nulovogo i jediničnogo & \hyperlink{TERM-1275}{TERM-1275}\\
obča linearna grupa & \hyperlink{TERM-0122}{TERM-0122}\\
obča teorija idealov & \hyperlink{TERM-0121}{TERM-0121}\\
obče nule & \hyperlink{TERM-0693}{TERM-0693}\\
obče nuly & \hyperlink{TERM-0694}{TERM-0694}\\
občeje |tE-C| & \hyperlink{TERM-0124}{TERM-0124}\\
obči dělitelj & \hyperlink{TERM-0697}{TERM-0697}\\
obči dělitelj v G & \hyperlink{TERM-0698}{TERM-0698}\\
obći imenovatelj & \hyperlink{TERM-0695}{TERM-0695}\\
obči imenovatelj & \hyperlink{TERM-0696}{TERM-0696}\\
obći ręd form & \hyperlink{TERM-0701}{TERM-0701}\\
obči teorem o diskriminantu & \hyperlink{TERM-0123}{TERM-0123}\\
Obči, obobčiti i soobčenje: razny smysl & \hyperlink{DEC-0015}{DEC-0015}\\
občna nuleva točka & \hyperlink{TERM-0692}{TERM-0692}\\
obglavjena relacija & \hyperlink{TERM-0266}{TERM-0266}\\
obična v literaturě definicija polnoj razložimosti & \hyperlink{TERM-0270}{TERM-0270}\\
oblast / čislova oblast & \hyperlink{TERM-0179}{TERM-0179}\\
oblast c?losti [H] & \hyperlink{TERM-0902}{TERM-0902}\\
oblast c?losti [Theta] & \hyperlink{TERM-0903}{TERM-0903}\\
oblast cělosti & \hyperlink{TERM-0904}{TERM-0904}, \hyperlink{TERM-0907}{TERM-0907}, \hyperlink{TERM-0909}{TERM-0909}, \hyperlink{TERM-1596}{TERM-1596}\\
oblast cělosti [H] & \hyperlink{TERM-0910}{TERM-0910}\\
oblast cělosti [Theta] & \hyperlink{TERM-0911}{TERM-0911}\\
oblast H iz algebraično-cělyh transcendentnyh čisel & \hyperlink{TERM-0182}{TERM-0182}\\
oblast koeficientov & \hyperlink{TERM-0997}{TERM-0997}\\
oblast koeficientov / polje P(u) & \hyperlink{TERM-0998}{TERM-0998}\\
oblast koeficientov Ω & \hyperlink{TERM-1001}{TERM-1001}\\
oblast množiteljev & \hyperlink{TERM-1191}{TERM-1191}\\
oblast odobraženja & \hyperlink{TERM-0023}{TERM-0023}\\
oblast operatorov & \hyperlink{TERM-1280}{TERM-1280}\\
oblast prěsěka / zatvorjena vzhodno k tvorenju prěsěkov & \hyperlink{TERM-1627}{TERM-1627}\\
oblast racionalnosti & \hyperlink{TERM-1461}{TERM-1461}\\
oblast vseh homogenyh form & \hyperlink{TERM-0181}{TERM-0181}\\
oblast vseh parnyh polinomov & \hyperlink{TERM-0180}{TERM-0180}\\
oblast/kolco klasov ostatkov & \hyperlink{TERM-1556}{TERM-1556}\\
oblasti & \hyperlink{TERM-0682}{TERM-0682}\\
oblasti iz c?lyh transcendentnyh ?isel & \hyperlink{TERM-0183}{TERM-0183}\\
obrat & \hyperlink{TERM-1760}{TERM-1760}\\
obrat (1) po (1') & \hyperlink{TERM-1762}{TERM-1762}\\
obrat lemy I & \hyperlink{TERM-1763}{TERM-1763}\\
obrat sistema ravnanj & \hyperlink{TERM-1761}{TERM-1761}\\
obratna matrica & \hyperlink{TERM-1585}{TERM-1585}\\
obratno jednoznačny / vzaimno jednoznačny & \hyperlink{TERM-1759}{TERM-1759}\\
obratno-jednoznačna funkcija & \hyperlink{TERM-1758}{TERM-1758}\\
obrubljenje & \hyperlink{TERM-1608}{TERM-1608}\\
Obstajanje: dokaz, že objekt jest & \hyperlink{DEC-0047}{DEC-0047}\\
Obstajati, suči, sučstveny, sučnost & \hyperlink{DEC-0025}{DEC-0025}\\
Obstajati: bytije matematičnogo objekta & \hyperlink{DEC-0123}{DEC-0123}\\
obyčno linearno diferencialno uravnenje 2-go reda & \hyperlink{TERM-0708}{TERM-0708}\\
obzor redukcij & \hyperlink{TERM-2011}{TERM-2011}\\
Obzor tvorjenja sistema & \hyperlink{TERM-1999}{TERM-1999}\\
Očevidny i očevidno & \hyperlink{DEC-0144}{DEC-0144}\\
od nuly različny i konečny & \hyperlink{TERM-1863}{TERM-1863}\\
od nyně vse formuly trěba razuměti mod $(\varrho,t)$ & \hyperlink{TERM-1862}{TERM-1862}\\
odděliti & \hyperlink{TERM-0051}{TERM-0051}\\
oddělne formy & \hyperlink{TERM-0455}{TERM-0455}\\
odgovarjajuče formy & \hyperlink{TERM-0502}{TERM-0502}\\
odgovarjajuče matrice & \hyperlink{TERM-1055}{TERM-1055}\\
odgovarjajut ... jednoznačno vzajemno & \hyperlink{TERM-0412}{TERM-0412}\\
Odgovarjati, analogično i po redu & \hyperlink{DEC-0051}{DEC-0051}\\
odnosi ravnosti i odnosi operacij & \hyperlink{TERM-0716}{TERM-0716}\\
odnosno polny sistem & \hyperlink{TERM-1524}{TERM-1524}\\
odobražajuči sistem & \hyperlink{TERM-0026}{TERM-0026}\\
odobraženje s neopreděljenymi koeficientami & \hyperlink{TERM-0022}{TERM-0022}\\
odpovědanje vlastnyh energijskih zakonov / obča relativnost & \hyperlink{TERM-1813}{TERM-1813}\\
Odpovědati, odgovarjajuči, korespondencija, neuspěh & \hyperlink{DEC-0021}{DEC-0021}\\
ograničenje čisla & \hyperlink{TERM-0140}{TERM-0140}\\
ograničeny po x\_i / ne dosega stepena k & \hyperlink{TERM-0184}{TERM-0184}\\
okolina & \hyperlink{TERM-1757}{TERM-1757}\\
operacija prěryvaje se & \hyperlink{TERM-1279}{TERM-1279}\\
operacije / vrsty spojivanja: složenje i množenje & \hyperlink{TERM-1811}{TERM-1811}\\
operatorno izomorfny & \hyperlink{TERM-1283}{TERM-1283}\\
operatorno izomorfny / operatorny homomorfizm & \hyperlink{TERM-1285}{TERM-1285}\\
operatorno-izomorfny & \hyperlink{TERM-1284}{TERM-1284}\\
operatorny homomorfizm & \hyperlink{TERM-1282}{TERM-1282}\\
operatorny homomorfizm / operatorno-homomorfny & \hyperlink{TERM-1281}{TERM-1281}\\
operatorny izomorfizm & \hyperlink{TERM-1286}{TERM-1286}\\
Oprěděliti, oprěděljati: definicija i opreděljenje & \hyperlink{DEC-0066}{DEC-0066}\\
Oprěděljeny: oblik zavisi od konstrukcije & \hyperlink{DEC-0067}{DEC-0067}\\
Originalno nazvanje i město & \hyperlink{DEC-0163}{DEC-0163}\\
ortogonalnostne uslovja & \hyperlink{TERM-1296}{TERM-1296}\\
osnovna forma & \hyperlink{TERM-0731}{TERM-0731}\\
osnovna forma / sistema osnovnyh form & \hyperlink{TERM-0732}{TERM-0732}\\
osnovne formy & \hyperlink{TERM-0733}{TERM-0733}\\
osnovne funkcije & \hyperlink{TERM-0734}{TERM-0734}\\
osnovne oblasti & \hyperlink{TERM-0728}{TERM-0728}\\
osnovne svijanja & \hyperlink{TERM-0730}{TERM-0730}\\
osnovno polje & \hyperlink{TERM-0737}{TERM-0737}\\
osnovno polje / nadpolje & \hyperlink{TERM-0738}{TERM-0738}\\
osnovno svijanje tipa rho & \hyperlink{TERM-0729}{TERM-0729}\\
osnovny ideal & \hyperlink{TERM-0735}{TERM-0735}\\
osnovny ideal i-go nivoja & \hyperlink{TERM-0736}{TERM-0736}\\
osnovny modul & \hyperlink{TERM-0739}{TERM-0739}\\
osnovny modul / osnovny ideal & \hyperlink{TERM-0740}{TERM-0740}\\
osoblive točke / točka vozvrata & \hyperlink{TERM-1636}{TERM-1636}\\
ostale formy & \hyperlink{TERM-0272}{TERM-0272}\\
ostatkova klasa & \hyperlink{TERM-1550}{TERM-1550}\\
ostatkova klasa / sistem ostatkovyh klasov & \hyperlink{TERM-1551}{TERM-1551}\\
ostatkove klase modulo p & \hyperlink{TERM-1554}{TERM-1554}\\
ostatny sistem G-[Theta] & \hyperlink{TERM-1575}{TERM-1575}\\
ostrej?e relacije (27), (28) & \hyperlink{TERM-1621}{TERM-1621}\\
ovladanje tym svęzom medžu formami & \hyperlink{TERM-0174}{TERM-0174}\\
označenja i opreděljenja & \hyperlink{TERM-0185}{TERM-0185}\\
P = T p kako razširjeny ideal & \hyperlink{TERM-1298}{TERM-1298}\\
p sam poraždaje tot prosty ideal & \hyperlink{TERM-1300}{TERM-1300}\\
p-adične komponenty / razvětvjene města & \hyperlink{TERM-1302}{TERM-1302}\\
p-adično razširjenje & \hyperlink{TERM-1301}{TERM-1301}\\
p-togo razmjera & \hyperlink{TERM-1303}{TERM-1303}\\
paralelno prěnešenje vektora & \hyperlink{TERM-1311}{TERM-1311}\\
parametrično predstavjenje & \hyperlink{TERM-1312}{TERM-1312}\\
parametrno predstavjenje & \hyperlink{TERM-1313}{TERM-1313}\\
Peirceovo razloženje & \hyperlink{TERM-1314}{TERM-1314}\\
Pet kolon i dva znaka primětky & \hyperlink{DEC-0133}{DEC-0133}\\
pitanja o bazah & \hyperlink{TERM-0167}{TERM-0167}\\
po dualistično protivopoložnom razmatrjenju & \hyperlink{TERM-0360}{TERM-0360}\\
Po odnošenju k: odnos i padež & \hyperlink{DEC-0009}{DEC-0009}\\
po parah vzajemno proste & \hyperlink{TERM-1305}{TERM-1305}\\
po parah vzajemno proste primarne ideale & \hyperlink{TERM-1307}{TERM-1307}\\
po parah vzajemno proste primarne komponenty & \hyperlink{TERM-1309}{TERM-1309}\\
po punktu 1 vsaky determinant ... jest dělimy črez & \hyperlink{TERM-1197}{TERM-1197}\\
po slědujućej shemě & \hyperlink{TERM-1198}{TERM-1198}\\
po slědujućej začetnoj shemě & \hyperlink{TERM-1199}{TERM-1199}\\
po sootvětnom računu & \hyperlink{TERM-1201}{TERM-1201}\\
po začętnoj shemě & \hyperlink{TERM-1200}{TERM-1200}\\
početna forma & \hyperlink{TERM-0135}{TERM-0135}\\
početne členy & \hyperlink{TERM-0136}{TERM-0136}\\
početno uslovje & \hyperlink{TERM-0134}{TERM-0134}\\
podgrupa grupy ostatkov & \hyperlink{TERM-1785}{TERM-1785}\\
podgrupy & \hyperlink{TERM-1786}{TERM-1786}\\
podgrupy projektivnoj grupy & \hyperlink{TERM-1787}{TERM-1787}\\
podkolca, ktore jednovrěmenno sut idealy & \hyperlink{TERM-1794}{TERM-1794}\\
podkolco & \hyperlink{TERM-1792}{TERM-1792}\\
podkolco / nadkolco / nadpolje & \hyperlink{TERM-1793}{TERM-1793}\\
podmodul & \hyperlink{TERM-1791}{TERM-1791}\\
podpolje & \hyperlink{TERM-1707}{TERM-1707}, \hyperlink{TERM-1788}{TERM-1788}\\
podpolje ostatkovyh klasov & \hyperlink{TERM-1572}{TERM-1572}\\
podpolje P, izomorfno k polju o/op & \hyperlink{TERM-1790}{TERM-1790}\\
podrazděljenje & \hyperlink{TERM-1796}{TERM-1796}\\
Podrobno predstavjenje imaje pojaviti se v Math. Ann. & \hyperlink{TERM-0403}{TERM-0403}\\
podsistemy & \hyperlink{TERM-1733}{TERM-1733}\\
podsistemy $s_I$, $s_{II}$, $s_{III}$, $s_{IV}$ & \hyperlink{TERM-1734}{TERM-1734}\\
podtělo & \hyperlink{TERM-1789}{TERM-1789}\\
pojedina norma & \hyperlink{TERM-0458}{TERM-0458}\\
pojedina suma & \hyperlink{TERM-0459}{TERM-0459}\\
Pojedinny: jedinica razsmatrana oddělno & \hyperlink{DEC-0136}{DEC-0136}\\
pojem produkta ne prěnosi se na moduly & \hyperlink{TERM-1388}{TERM-1388}\\
pojem točke / grupa toček & \hyperlink{TERM-1406}{TERM-1406}\\
pojmy bazy & \hyperlink{TERM-0165}{TERM-0165}\\
pojmy děliteljnosti & \hyperlink{TERM-1705}{TERM-1705}\\
polarizacija $x$ po $\nu_1$ & \hyperlink{TERM-1319}{TERM-1319}\\
polarizacija / polarny členy & \hyperlink{TERM-1320}{TERM-1320}\\
polarne procesy / polarne operacije & \hyperlink{TERM-1322}{TERM-1322}\\
polarne procesy / polary & \hyperlink{TERM-1321}{TERM-1321}\\
polary & \hyperlink{TERM-1316}{TERM-1316}\\
polary / polarne operacije / polarny proces & \hyperlink{TERM-1318}{TERM-1318}\\
polary / tvorjenje polarov / polarny procesy & \hyperlink{TERM-1317}{TERM-1317}\\
poligon (odnosno divizor) / klas poligonov & \hyperlink{TERM-1323}{TERM-1323}\\
polinomialno kolco & \hyperlink{TERM-1324}{TERM-1324}\\
polinomna oblast & \hyperlink{TERM-1325}{TERM-1325}\\
polinomna oblast / polinomy & \hyperlink{TERM-1326}{TERM-1326}\\
polinomna oblast iz izčetno mnogih neopreděljenyh & \hyperlink{TERM-1327}{TERM-1327}\\
polinomne idealy / idealy iz polinomov & \hyperlink{TERM-1329}{TERM-1329}\\
polja funkcij & \hyperlink{TERM-0618}{TERM-0618}\\
Polja i sistemy racionalnyh funkcij & \hyperlink{TERM-1076}{TERM-1076}\\
polja izomorfne k Q & \hyperlink{TERM-1936}{TERM-1936}\\
polja K\_\{n ho\} & \hyperlink{TERM-1074}{TERM-1074}\\
polja racionalnyh funkcij / polje racionalnyh funkcij & \hyperlink{TERM-1454}{TERM-1454}\\
polje & \hyperlink{TERM-1066}{TERM-1066}\\
polje / kolco klasov ostatkov bez děliteljev nuly & \hyperlink{TERM-1071}{TERM-1071}\\
polje / polje funkcij / čislovo polje & \hyperlink{TERM-1070}{TERM-1070}\\
polje klasov ostatkov / sovršeno polje & \hyperlink{TERM-1558}{TERM-1558}\\
polje kvocientov & \hyperlink{TERM-1421}{TERM-1421}\\
polje kvocientov / kolco kvocientov & \hyperlink{TERM-1422}{TERM-1422}\\
polje kvocientov integralnoj oblasti & \hyperlink{TERM-1423}{TERM-1423}\\
polje kvocientov od T & \hyperlink{TERM-1424}{TERM-1424}\\
polje nuljev & \hyperlink{TERM-1268}{TERM-1268}\\
polje obrazov / oblast obrazov & \hyperlink{TERM-0025}{TERM-0025}\\
polje ostatkovyh klasov & \hyperlink{TERM-1557}{TERM-1557}\\
polje početnogo odrězka K(vartheta) & \hyperlink{TERM-0036}{TERM-0036}\\
polje racionalnyh funkcij & \hyperlink{TERM-1075}{TERM-1075}\\
polje razloženja & \hyperlink{TERM-1918}{TERM-1918}\\
polje vsih kompleksnyh čisel & \hyperlink{TERM-1072}{TERM-1072}\\
polje vsih kompleksnyh čisel / čislovo polje / tělo & \hyperlink{TERM-1073}{TERM-1073}\\
polna linearna sistema & \hyperlink{TERM-1110}{TERM-1110}\\
polna matrična algebra & \hyperlink{TERM-1852}{TERM-1852}\\
polna razložimost prvogo roda & \hyperlink{TERM-1853}{TERM-1853}\\
polne sistemy form & \hyperlink{TERM-1851}{TERM-1851}\\
polne sistemy invariantov & \hyperlink{TERM-1832}{TERM-1832}\\
polne transvekcije nad kovariantami sistema III & \hyperlink{TERM-1833}{TERM-1833}\\
Polno ime i město izdavanja & \hyperlink{DEC-0152}{DEC-0152}\\
polno razložimy & \hyperlink{TERM-1842}{TERM-1842}\\
polno razložimy drugogo roda & \hyperlink{TERM-1846}{TERM-1846}\\
polno razložimy prvogo roda & \hyperlink{TERM-1845}{TERM-1845}\\
polno redukovane idealy & \hyperlink{TERM-1847}{TERM-1847}, \hyperlink{TERM-1850}{TERM-1850}\\
polnost identičnostij & \hyperlink{TERM-1854}{TERM-1854}\\
polny sistem invariantov & \hyperlink{TERM-1834}{TERM-1834}\\
Polny, cělkom, doskonaly & \hyperlink{DEC-0062}{DEC-0062}\\
poluprosto kolco & \hyperlink{TERM-0779}{TERM-0779}\\
poluprosty sistem & \hyperlink{TERM-0780}{TERM-0780}\\
pomnoženyh konstantami & \hyperlink{TERM-1159}{TERM-1159}\\
pomocna lema & \hyperlink{TERM-0816}{TERM-0816}\\
pomočna lema & \hyperlink{TERM-0818}{TERM-0818}\\
pomocne prěměnne & \hyperlink{TERM-0821}{TERM-0821}\\
pomoćne prěměnne & \hyperlink{TERM-0822}{TERM-0822}\\
pomoćny teorem & \hyperlink{TERM-0817}{TERM-0817}\\
pomočny teorem & \hyperlink{TERM-0819}{TERM-0819}\\
pomoćny teorem o algebraično cěloj zavisnosti & \hyperlink{TERM-0820}{TERM-0820}\\
ponašanje slěda i diskriminanta pri direktnoj sumě klasov & \hyperlink{TERM-1810}{TERM-1810}\\
poneže po punktu 3 to kolco klasov ostatkov jest izomorfno & \hyperlink{TERM-0222}{TERM-0222}\\
Ponovno: povtorjenje & \hyperlink{DEC-0085}{DEC-0085}\\
poparno vzaimno proste & \hyperlink{TERM-1306}{TERM-1306}\\
poparno vzaimno proste primarne idealy & \hyperlink{TERM-1308}{TERM-1308}\\
popolno reducibilna grupa & \hyperlink{TERM-1841}{TERM-1841}\\
popolno reducibilno kolco & \hyperlink{TERM-1848}{TERM-1848}\\
popolno reducibilno kolco / grupno kolco & \hyperlink{TERM-1849}{TERM-1849}\\
popolno reducibilny & \hyperlink{TERM-1840}{TERM-1840}, \hyperlink{TERM-1843}{TERM-1843}\\
popolno reducibilny / popolna reducibilnost & \hyperlink{TERM-1844}{TERM-1844}\\
Poprědny i raněje & \hyperlink{DEC-0114}{DEC-0114}\\
Pored i poredok; poręd i porędok & \hyperlink{DEC-0056}{DEC-0056}\\
poręd v koeficientah & \hyperlink{TERM-1292}{TERM-1292}\\
porjadky ranga n & \hyperlink{TERM-1294}{TERM-1294}\\
porjadky relativnyh polj & \hyperlink{TERM-1295}{TERM-1295}\\
porjadky v čislovyh i funkcionalnyh poljah & \hyperlink{TERM-1293}{TERM-1293}\\
porjadok & \hyperlink{TERM-1287}{TERM-1287}\\
porjadok / glavny porjadok & \hyperlink{TERM-1289}{TERM-1289}\\
porjadok / konečny porjadok & \hyperlink{TERM-1288}{TERM-1288}\\
porjadok grupy / porjadkov h\_i & \hyperlink{TERM-0748}{TERM-0748}\\
porjadok vzhodno k o & \hyperlink{TERM-1291}{TERM-1291}\\
porodžuje se nad P črez adjunkciju jednogo/dvuh elementov & \hyperlink{TERM-0354}{TERM-0354}\\
porodžujųće formy / konstrukcije & \hyperlink{TERM-0189}{TERM-0189}\\
poslědki & \hyperlink{TERM-0570}{TERM-0570}\\
Poslědnji, poslědny & \hyperlink{DEC-0036}{DEC-0036}\\
poslědok & \hyperlink{TERM-0569}{TERM-0569}\\
poslědovateljno jednorazovo svijanje & \hyperlink{TERM-1667}{TERM-1667}\\
postulat izbora & \hyperlink{TERM-0156}{TERM-0156}\\
potencijalne produkty & \hyperlink{TERM-1335}{TERM-1335}\\
potrěbne i dostatočne uslovja & \hyperlink{TERM-1254}{TERM-1254}\\
Potrěbny i neobhodimy & \hyperlink{DEC-0052}{DEC-0052}\\
Potrěbny, potrěbno; ne potrěbno & \hyperlink{DEC-0039}{DEC-0039}\\
povsudu razložena algebra & \hyperlink{TERM-1998}{TERM-1998}\\
povtor $H_j^2$, $H_j^2$ sohranjen & \hyperlink{TERM-1641}{TERM-1641}\\
Povtorjeno ime ne znači isto upotrěbjenje & \hyperlink{DEC-0168}{DEC-0168}\\
prava / lěva obratna matrica & \hyperlink{TERM-1467}{TERM-1467}\\
Pravdivost, izpolnjenje i priměnimost & \hyperlink{DEC-0011}{DEC-0011}\\
prave odnosno lěve klasy & \hyperlink{TERM-1465}{TERM-1465}\\
pravo predstavjenje / lěvo predstavjenje & \hyperlink{TERM-1466}{TERM-1466}\\
pravostranne / lěvostranne / dvostranne moduly i idealy & \hyperlink{TERM-1469}{TERM-1469}\\
pravostranny / lěvostranny / dvustranny modul & \hyperlink{TERM-1468}{TERM-1468}\\
pravy K-modul & \hyperlink{TERM-0969}{TERM-0969}\\
Prědhoditi: objekt prědhodi drugomu & \hyperlink{DEC-0115}{DEC-0115}\\
Predloženo g. Šurom & \hyperlink{TERM-1867}{TERM-1867}\\
Prědpokladati i prědpoloženje & \hyperlink{DEC-0116}{DEC-0116}\\
Prědpokladati i prědpoložiti & \hyperlink{DEC-0043}{DEC-0043}\\
predpostavjeno razloženje & \hyperlink{TERM-1866}{TERM-1866}\\
prědšedši paragrafy & \hyperlink{TERM-0603}{TERM-0603}\\
Prědstav- v složenyh slovah i koncovkah & \hyperlink{DEC-0017}{DEC-0017}\\
predstavime s koeficientami iz o & \hyperlink{TERM-1158}{TERM-1158}\\
Prědstaviti, prědstavjati, prědstavjenje & \hyperlink{DEC-0016}{DEC-0016}\\
Prědstavjenij; tym jest dokazano & \hyperlink{DEC-0019}{DEC-0019}\\
predstavjenja konečnyh grup & \hyperlink{TERM-0233}{TERM-0233}\\
predstavjenje / klas / transformovane & \hyperlink{TERM-0229}{TERM-0229}\\
predstavjenje / klas predstavjenj & \hyperlink{TERM-0226}{TERM-0226}\\
predstavjenje / klasy predstavjenj & \hyperlink{TERM-0227}{TERM-0227}\\
predstavjenje / teorija predstavjenj & \hyperlink{TERM-0228}{TERM-0228}\\
predstavjenje jediničnogo elementa kako sumy & \hyperlink{TERM-1673}{TERM-1673}\\
predstavjenje kako direktna suma primarnyh kolc & \hyperlink{TERM-0230}{TERM-0230}\\
predstavjenje kako najmenše obče mnogokratno & \hyperlink{TERM-0231}{TERM-0231}\\
predstavjenje kvaternionnoj grupy & \hyperlink{TERM-0232}{TERM-0232}\\
predstavjenje racionalno nad poljem K & \hyperlink{TERM-0894}{TERM-0894}\\
predstavjenje razširjenyh kolc kako direktne sumy & \hyperlink{TERM-0306}{TERM-0306}\\
predstavjenje s stalym imenovateljem & \hyperlink{TERM-0234}{TERM-0234}\\
predstavjenje slěda, normy i diskriminanta & \hyperlink{TERM-0235}{TERM-0235}\\
Prědstavjenje tipa Weierstrass & \hyperlink{DEC-0002}{DEC-0002}\\
Prědtym ili točno; izpolnjeny ili dějajuči & \hyperlink{DEC-0151}{DEC-0151}\\
prěhod k kolcu kvocientov & \hyperlink{TERM-2003}{TERM-2003}\\
prěhod k konečnym razširjenjam polj & \hyperlink{TERM-2001}{TERM-2001}\\
prěhod k razširjenym kolcam & \hyperlink{TERM-2002}{TERM-2002}\\
prěhod od porjadka k kolcam konečnogo ranga vzhodno k polju & \hyperlink{TERM-2000}{TERM-2000}\\
Prěhodi: pravopis i finitna forma & \hyperlink{DEC-0118}{DEC-0118}\\
prehoditi v sebe / sootvětovati samoj sobě & \hyperlink{TERM-0895}{TERM-0895}\\
prěhodno prěobraženje & \hyperlink{TERM-2004}{TERM-2004}\\
prěkrivanja & \hyperlink{TERM-2009}{TERM-2009}\\
pręmougolna shema & \hyperlink{TERM-1464}{TERM-1464}\\
prěnos & \hyperlink{TERM-2012}{TERM-2012}\\
prěnos iz binarnoj oblasti & \hyperlink{TERM-2013}{TERM-2013}\\
prěobraženje matric & \hyperlink{TERM-1141}{TERM-1141}\\
prěsěčno polje & \hyperlink{TERM-0369}{TERM-0369}\\
prevedeno kako 'A i B' & \hyperlink{TERM-1642}{TERM-1642}\\
Prěz i črěz & \hyperlink{DEC-0003}{DEC-0003}\\
pri jednovrěmennom razsmotrenju & \hyperlink{TERM-1781}{TERM-1781}\\
pri tom drži & \hyperlink{TERM-0223}{TERM-0223}\\
pridavati & \hyperlink{TERM-0056}{TERM-0056}\\
pridruženy prosty ideal & \hyperlink{TERM-1940}{TERM-1940}, \hyperlink{TERM-1944}{TERM-1944}\\
pridruženy prosty ideal / pridruženo prosto čislo & \hyperlink{TERM-1942}{TERM-1942}\\
pridruženy prosty ideal v nekomutativnom slučaju & \hyperlink{TERM-1943}{TERM-1943}\\
primarna funkcija & \hyperlink{TERM-1379}{TERM-1379}\\
primarna funkcija / prosta funkcija & \hyperlink{TERM-1380}{TERM-1380}\\
primarna komponenta naležeča k \textbackslash{}mo & \hyperlink{TERM-1933}{TERM-1933}\\
primarna komponenta prinadležeča k p & \hyperlink{TERM-1935}{TERM-1935}\\
primarne faktory & \hyperlink{TERM-1378}{TERM-1378}\\
primarne idealy / primarny faktor & \hyperlink{TERM-1375}{TERM-1375}\\
primarne kolca & \hyperlink{TERM-1376}{TERM-1376}\\
primarne komponenty stavajut prostymi idealami & \hyperlink{TERM-1383}{TERM-1383}\\
primarno kolco & \hyperlink{TERM-1377}{TERM-1377}\\
primarny & \hyperlink{TERM-1374}{TERM-1374}\\
primarny ideal & \hyperlink{TERM-1381}{TERM-1381}\\
primarny ideal / asociovany prosty ideal & \hyperlink{TERM-1382}{TERM-1382}\\
Priměr i napriměr & \hyperlink{DEC-0097}{DEC-0097}\\
Priměry i priměry: kako čitati čisla pojavjenij & \hyperlink{DEC-0098}{DEC-0098}\\
primětka & \hyperlink{TERM-1203}{TERM-1203}\\
primětka ne jest duplicirovana & \hyperlink{TERM-0219}{TERM-0219}\\
primitivne idempotentne elementy & \hyperlink{TERM-1367}{TERM-1367}\\
primitivne polinomy & \hyperlink{TERM-1369}{TERM-1369}\\
primitivny element / primitivna funkcija & \hyperlink{TERM-1370}{TERM-1370}\\
primitivny element kone?nogo polja raz?irjenja & \hyperlink{TERM-1371}{TERM-1371}\\
primitivny p-ty korenj iz jedinice & \hyperlink{TERM-1368}{TERM-1368}\\
Prinajmenje i najmenše obče mnogokratno & \hyperlink{DEC-0078}{DEC-0078}\\
Prinaležati; prinaležet & \hyperlink{DEC-0045}{DEC-0045}\\
prinaležne sistemi veličin / sistem nul & \hyperlink{TERM-1959}{TERM-1959}\\
prinaležny sistem nuljev & \hyperlink{TERM-1960}{TERM-1960}\\
prinaležny sistem veličin & \hyperlink{TERM-0143}{TERM-0143}\\
princip prěnosa & \hyperlink{TERM-2015}{TERM-2015}\\
princip prěnosa za racionalnu bazu & \hyperlink{TERM-2016}{TERM-2016}\\
prirědžene idealy & \hyperlink{TERM-1945}{TERM-1945}\\
prirędžene rędy & \hyperlink{TERM-1947}{TERM-1947}\\
prirędženje & \hyperlink{TERM-1953}{TERM-1953}\\
prirěđene matrice & \hyperlink{TERM-1946}{TERM-1946}\\
prirodne iracionalnosti prinadležeče grupě & \hyperlink{TERM-1204}{TERM-1204}\\
prisojedinjenje neopreděljenyh u\_\{mu nu\} & \hyperlink{TERM-0062}{TERM-0062}\\
prisojedinjenje primitivnogo elementa prvogo roda & \hyperlink{TERM-0063}{TERM-0063}\\
prisojedinjenje u\_\{\textbackslash{}mu\textbackslash{}nu\} & \hyperlink{TERM-0061}{TERM-0061}\\
proces diferenciacije, ktory točno izčerpava rędy & \hyperlink{TERM-0271}{TERM-0271}\\
proces Poissonovyh skobok & \hyperlink{TERM-1315}{TERM-1315}\\
proces svijanja & \hyperlink{TERM-0563}{TERM-0563}\\
proces svijanja / osnovne svijanja & \hyperlink{TERM-0566}{TERM-0566}\\
proces svijanja / svijanja & \hyperlink{TERM-0565}{TERM-0565}\\
proces svijanja / svijanje & \hyperlink{TERM-0564}{TERM-0564}\\
procesy diferenciacije & \hyperlink{TERM-0284}{TERM-0284}\\
procesy diferenciranja i variiranja & \hyperlink{TERM-0283}{TERM-0283}\\
produkt diskriminantnyh idealov komponentov & \hyperlink{TERM-1384}{TERM-1384}\\
produkt idealov & \hyperlink{TERM-1387}{TERM-1387}\\
produkt konečno mnogih po parah vzajemno prostyh primarnyh idealov & \hyperlink{TERM-1386}{TERM-1386}\\
produkt konečno mnogih poparno vzaimno prostyh idealov & \hyperlink{TERM-1385}{TERM-1385}\\
produkt stepenjev & \hyperlink{TERM-1332}{TERM-1332}\\
produkt stepenjev prostyh idealov / vlastny dělitelj & \hyperlink{TERM-1334}{TERM-1334}\\
produktna matrica & \hyperlink{TERM-1392}{TERM-1392}\\
produktno predstavjenje & \hyperlink{TERM-1389}{TERM-1389}\\
produktno predstavjenje jego nulovogo ideala & \hyperlink{TERM-1390}{TERM-1390}\\
produkty & \hyperlink{TERM-1391}{TERM-1391}\\
produkty stepenev & \hyperlink{TERM-1336}{TERM-1336}\\
produkty stepenev v imenovatelju / častne derivaty & \hyperlink{TERM-1338}{TERM-1338}\\
Proizhoditi iz, slědovati, nastavati, početny & \hyperlink{DEC-0032}{DEC-0032}\\
proizvoljne formy & \hyperlink{TERM-0176}{TERM-0176}\\
projektivna invarianta & \hyperlink{TERM-1400}{TERM-1400}\\
projektivna invariantna polja & \hyperlink{TERM-1401}{TERM-1401}\\
projektivna invariantna teorija & \hyperlink{TERM-1402}{TERM-1402}\\
prost ideal & \hyperlink{TERM-1353}{TERM-1353}\\
prost ideal / primarny ideal / asociovany prost ideal & \hyperlink{TERM-1359}{TERM-1359}\\
prost ideal / vlastny primarny ideal & \hyperlink{TERM-1355}{TERM-1355}\\
prost vzhodno k / vzaimno prosty & \hyperlink{TERM-1341}{TERM-1341}\\
prost vzhodno k a / vzaimno proste v tom smyslu & \hyperlink{TERM-1342}{TERM-1342}\\
prosta funkcija & \hyperlink{TERM-1344}{TERM-1344}\\
prosta funkcija / primarna funkcija & \hyperlink{TERM-1346}{TERM-1346}\\
prosta funkcija prvogo roda & \hyperlink{TERM-1348}{TERM-1348}\\
prosta grupa & \hyperlink{TERM-0418}{TERM-0418}, \hyperlink{TERM-0419}{TERM-0419}, \hyperlink{TERM-1350}{TERM-1350}\\
prosta normalna algebra / dělitevna algebra & \hyperlink{TERM-0421}{TERM-0421}\\
proste idealy & \hyperlink{TERM-1365}{TERM-1365}, \hyperlink{TERM-1366}{TERM-1366}\\
prosto izomorfna & \hyperlink{TERM-0451}{TERM-0451}\\
prosto slěd / nezavisno od klasa & \hyperlink{TERM-1657}{TERM-1657}\\
prosty element & \hyperlink{TERM-1351}{TERM-1351}\\
prosty element osnovnoj oblasti & \hyperlink{TERM-1352}{TERM-1352}\\
prosty element p iz o\_p vhodi v diskriminant D\_\{T\_P\} & \hyperlink{TERM-1299}{TERM-1299}\\
prosty element vhodi v diskriminant & \hyperlink{TERM-0404}{TERM-0404}\\
prosty ideal & \hyperlink{TERM-1354}{TERM-1354}\\
prosty ideal / primarny ideal / asociovany prosty ideal & \hyperlink{TERM-1360}{TERM-1360}\\
prosty ideal / primarny ideal / primarna komponenta & \hyperlink{TERM-1356}{TERM-1356}, \hyperlink{TERM-1358}{TERM-1358}\\
prosty ideal / proste idealy & \hyperlink{TERM-1357}{TERM-1357}\\
prosty ideal drugogo roda & \hyperlink{TERM-1364}{TERM-1364}\\
prosty ideal p iz multiplikacijskogo kolca o & \hyperlink{TERM-1363}{TERM-1363}\\
prosty ideal prvogo / drugogo roda & \hyperlink{TERM-1361}{TERM-1361}\\
prosty ideal prvogo i drugogo roda & \hyperlink{TERM-1362}{TERM-1362}\\
prosty ideal, asociovany s nulovym idealom & \hyperlink{TERM-0245}{TERM-0245}\\
prosty lěvy ideal & \hyperlink{TERM-0423}{TERM-0423}\\
Prosty lěvy ideal i jego indeks & \hyperlink{DEC-0131}{DEC-0131}\\
prosty modul & \hyperlink{TERM-0422}{TERM-0422}, \hyperlink{TERM-1373}{TERM-1373}\\
provodimo tabelarno & \hyperlink{TERM-1699}{TERM-1699}\\
provodnik & \hyperlink{TERM-0544}{TERM-0544}, \hyperlink{TERM-0619}{TERM-0619}\\
provodnik / porjadok / oblast cělosti & \hyperlink{TERM-0620}{TERM-0620}\\
Prsten: algebraična struktura & \hyperlink{DEC-0013}{DEC-0013}\\
prve integraly & \hyperlink{TERM-0507}{TERM-0507}\\
prvogo roda & \hyperlink{TERM-0509}{TERM-0509}\\
prvotno polje & \hyperlink{TERM-1372}{TERM-1372}\\
prvotno razkladanje & \hyperlink{TERM-1801}{TERM-1801}\\
prvy / drugi teorem o izomorfiji & \hyperlink{TERM-0508}{TERM-0508}\\
prvy teorem o izomorfizmu & \hyperlink{TERM-0510}{TERM-0510}\\
Pytanje; vopros i zapytanje & \hyperlink{DEC-0023}{DEC-0023}\\
R jest homomorfno k K\_i / izraziti črez elementy iz K\_i & \hyperlink{TERM-1427}{TERM-1427}\\
R samo kako polje ekvivalentno k K\_i, a ne kako polje samo ležeče v Q & \hyperlink{TERM-1428}{TERM-1428}\\
R-porjadok & \hyperlink{TERM-1429}{TERM-1429}\\
racionalna baza & \hyperlink{TERM-1447}{TERM-1447}, \hyperlink{TERM-1448}{TERM-1448}, \hyperlink{TERM-1449}{TERM-1449}\\
racionalna baza / algebraična baza & \hyperlink{TERM-1450}{TERM-1450}\\
racionalna baza s dodatnym uslovjem & \hyperlink{TERM-1452}{TERM-1452}\\
racionalna baza vsih čisel & \hyperlink{TERM-1451}{TERM-1451}\\
racionalna baza, jednoznačno opreděljena poljem & \hyperlink{TERM-0357}{TERM-0357}\\
racionalna kombinacija & \hyperlink{TERM-1457}{TERM-1457}\\
racionalne relacije & \hyperlink{TERM-1456}{TERM-1456}\\
racionalno / algebraično / transcendentno čislo & \hyperlink{TERM-1458}{TERM-1458}\\
racionalno nezavisimy & \hyperlink{TERM-1441}{TERM-1441}\\
racionalno predstavjenje & \hyperlink{TERM-1453}{TERM-1453}\\
racionalno prosto čislo & \hyperlink{TERM-1455}{TERM-1455}\\
racionalno zavisny / nezavisny & \hyperlink{TERM-1440}{TERM-1440}\\
racionalno- i algebraično-drobne funkcije & \hyperlink{TERM-1442}{TERM-1442}\\
racionalno-c?la adjunkcija & \hyperlink{TERM-1444}{TERM-1444}\\
racionalno-c?lo adjungovati & \hyperlink{TERM-1443}{TERM-1443}\\
racionalno-cěla adjunkcija & \hyperlink{TERM-1445}{TERM-1445}\\
racionalno-cěle koeficienty / cělo-racionalno-čislove invarianty & \hyperlink{TERM-0664}{TERM-0664}\\
racionalno-čislove koeficienty / funkcije & \hyperlink{TERM-1462}{TERM-1462}\\
racionalno-drobny / algebraično-drobny & \hyperlink{TERM-1446}{TERM-1446}\\
racionalny / cěly Galoisov modul & \hyperlink{TERM-1459}{TERM-1459}\\
racionalny funkcional & \hyperlink{TERM-1460}{TERM-1460}\\
Racionalny koeficient i zavisimy indeks & \hyperlink{DEC-0128}{DEC-0128}\\
računno sredstvo & \hyperlink{TERM-1463}{TERM-1463}\\
radikal & \hyperlink{TERM-1432}{TERM-1432}\\
radikal / kolco bez radikala & \hyperlink{TERM-1433}{TERM-1433}\\
rang & \hyperlink{TERM-1436}{TERM-1436}\\
rang četyri/tri/dva/jedin & \hyperlink{TERM-1438}{TERM-1438}\\
rang vzhodno k K & \hyperlink{TERM-1437}{TERM-1437}\\
rangi & \hyperlink{TERM-1439}{TERM-1439}\\
ravnost čisla komponentov & \hyperlink{TERM-0141}{TERM-0141}\\
ravnost koeficientov & \hyperlink{TERM-1002}{TERM-1002}\\
razděljenje oblastij ... na klasy & \hyperlink{TERM-0452}{TERM-0452}\\
Razgledajmo; razgleděti; razgledanje & \hyperlink{DEC-0049}{DEC-0049}\\
Razgledati i razvažati & \hyperlink{DEC-0030}{DEC-0030}\\
Razgledati objekt kako něšto & \hyperlink{DEC-0041}{DEC-0041}\\
razklad Klebš--Gordana & \hyperlink{TERM-0217}{TERM-0217}\\
razklad v k sostavnyh ?esti / k podgrup & \hyperlink{TERM-1923}{TERM-1923}\\
razklad v rědy v teoriji form & \hyperlink{TERM-1508}{TERM-1508}\\
razklada se / razkladna forma & \hyperlink{TERM-1917}{TERM-1917}\\
razkladaje se / razkladanje & \hyperlink{TERM-1916}{TERM-1916}\\
razkladajut se / razkladne formy & \hyperlink{TERM-1903}{TERM-1903}\\
razkladanje ili reducibilnost ostalyh form & \hyperlink{TERM-1904}{TERM-1904}\\
razkladanje transvekcij s funkcionalnymi determinantami & \hyperlink{TERM-1905}{TERM-1905}\\
razkladati se / razkladna forma & \hyperlink{TERM-1902}{TERM-1902}\\
razkladna identičnost & \hyperlink{TERM-1926}{TERM-1926}\\
razkladne formy & \hyperlink{TERM-1910}{TERM-1910}\\
razkladne identičnosti & \hyperlink{TERM-1927}{TERM-1927}\\
razkladne identičnosti v eksplicitnoj formě & \hyperlink{TERM-1928}{TERM-1928}\\
razkladny ręd form & \hyperlink{TERM-1912}{TERM-1912}\\
različny od nulovogo i jediničnogo ideala & \hyperlink{TERM-1855}{TERM-1855}\\
različny od nulovogo ideala & \hyperlink{TERM-1857}{TERM-1857}\\
Različnyh, različne, različnoj & \hyperlink{DEC-0088}{DEC-0088}\\
razlika med cělym i drobnym razmyvaje se & \hyperlink{TERM-1795}{TERM-1795}\\
razlikove izrazy & \hyperlink{TERM-0286}{TERM-0286}\\
razlo?ena algebra & \hyperlink{TERM-1906}{TERM-1906}\\
razlo?ena algebra; povsudu razlo?ena & \hyperlink{TERM-1907}{TERM-1907}\\
razložene skrěžene produkty & \hyperlink{TERM-1914}{TERM-1914}\\
razloženje ideala / razloženje nulovogo ideala & \hyperlink{TERM-0873}{TERM-0873}\\
razloženje ideala T p & \hyperlink{TERM-0874}{TERM-0874}\\
razloženje na dvustranne idealy & \hyperlink{TERM-1924}{TERM-1924}\\
razloženje obćih form po normalnyh formah & \hyperlink{TERM-0503}{TERM-0503}\\
razložima identičnost & \hyperlink{TERM-1922}{TERM-1922}\\
razložima/razložime formy & \hyperlink{TERM-1908}{TERM-1908}\\
razložime formy & \hyperlink{TERM-1911}{TERM-1911}\\
razložime rědy form ili formy & \hyperlink{TERM-1913}{TERM-1913}\\
razložimy / nerazložimy & \hyperlink{TERM-1486}{TERM-1486}\\
razložne teoremy & \hyperlink{TERM-1930}{TERM-1930}\\
razmatrivana kako kovariant & \hyperlink{TERM-0127}{TERM-0127}\\
Raznyh razkladov & \hyperlink{DEC-0090}{DEC-0090}\\
razpadne polja & \hyperlink{TERM-1915}{TERM-1915}, \hyperlink{TERM-1919}{TERM-1919}\\
razpadno polje & \hyperlink{TERM-1900}{TERM-1900}, \hyperlink{TERM-1920}{TERM-1920}\\
razpadno polje; odc?pno polje & \hyperlink{TERM-1901}{TERM-1901}\\
razpadno polje; plural title form: razpadne polja / locative razpadnyh poljah & \hyperlink{TERM-1921}{TERM-1921}\\
razširena Gaussova teorema & \hyperlink{TERM-0511}{TERM-0511}\\
razširenje Gaussova teorema na polinomy s algebraično-celočislenymi koeficientami & \hyperlink{TERM-0514}{TERM-0514}\\
razširenje Gaussovej teoremy & \hyperlink{TERM-0513}{TERM-0513}\\
razširenje osnovnogo těla & \hyperlink{TERM-0515}{TERM-0515}\\
razširenje prvogo roda & \hyperlink{TERM-0516}{TERM-0516}\\
razširitelno polje & \hyperlink{TERM-0524}{TERM-0524}\\
Razširiti, razširjeny, razširjenje & \hyperlink{DEC-0120}{DEC-0120}\\
razširjati oddělne komponenty & \hyperlink{TERM-0456}{TERM-0456}\\
razširjene idealy & \hyperlink{TERM-0523}{TERM-0523}\\
razširjene kolca polj & \hyperlink{TERM-0533}{TERM-0533}\\
razširjene polja s konjugatami y i Y & \hyperlink{TERM-1310}{TERM-1310}\\
razširjenja ... sovpadajut v vsih kolcah kvocientov & \hyperlink{TERM-0519}{TERM-0519}\\
razširjenje / suženje oblasti koeficientov & \hyperlink{TERM-0512}{TERM-0512}\\
razširjenje po pečatnom skanu & \hyperlink{TERM-1620}{TERM-1620}\\
razširjenje prvogo roda & \hyperlink{TERM-0517}{TERM-0517}\\
razširjenje prvogo roda / vtorogo roda & \hyperlink{TERM-0518}{TERM-0518}\\
Razširjenje: razširenje objekta & \hyperlink{DEC-0121}{DEC-0121}\\
razširjeno kolco & \hyperlink{TERM-0528}{TERM-0528}\\
razširjeno kolco \textbackslash{}overline\{\textbackslash{}mathfrak R\}=\textbackslash{}mathfrak R[\textbackslash{}overline P] & \hyperlink{TERM-0529}{TERM-0529}\\
razširjeno kolco od Z bez nulovyh děliteljev & \hyperlink{TERM-0531}{TERM-0531}\\
razširjeno kolco polja / razširjeno polje & \hyperlink{TERM-0530}{TERM-0530}\\
razširjeno polje / razširjeno kolco & \hyperlink{TERM-0525}{TERM-0525}\\
razširjeno polje K polja kvocientov kolca o\_p & \hyperlink{TERM-0527}{TERM-0527}\\
razširjeno polje prvogo roda podpolja P & \hyperlink{TERM-0526}{TERM-0526}\\
razširjeny ideal & \hyperlink{TERM-0521}{TERM-0521}\\
razširjeny ideal libovoljnogo ideala c & \hyperlink{TERM-0522}{TERM-0522}\\
Razvažanje i v razgled & \hyperlink{DEC-0050}{DEC-0050}\\
razvitja v ręd po polarah & \hyperlink{TERM-1510}{TERM-1510}\\
razvitje v rěd & \hyperlink{TERM-1506}{TERM-1506}\\
razvitje v ręd po polarah & \hyperlink{TERM-1509}{TERM-1509}\\
recenzije & \hyperlink{TERM-0198}{TERM-0198}\\
recipro?no predstavjenje; direktno predstavjenje & \hyperlink{TERM-1584}{TERM-1584}\\
ręd form & \hyperlink{TERM-0591}{TERM-0591}\\
ręd form / rędy form & \hyperlink{TERM-0592}{TERM-0592}, \hyperlink{TERM-0594}{TERM-0594}\\
red grupy & \hyperlink{TERM-1290}{TERM-1290}\\
ręd jest razložen & \hyperlink{TERM-0405}{TERM-0405}\\
ręd prěměnnyh & \hyperlink{TERM-1803}{TERM-1803}\\
red veličin / redy veličin & \hyperlink{TERM-0755}{TERM-0755}\\
Red, porędok i rang & \hyperlink{DEC-0119}{DEC-0119}\\
rędovo razloženje & \hyperlink{TERM-1507}{TERM-1507}\\
reducent & \hyperlink{TERM-1482}{TERM-1482}\\
reducenty & \hyperlink{TERM-1483}{TERM-1483}\\
reducenty / redukcija & \hyperlink{TERM-1484}{TERM-1484}\\
reducibilne predstavjenja & \hyperlink{TERM-1491}{TERM-1491}\\
reducibilny & \hyperlink{TERM-1485}{TERM-1485}\\
reducibilny / ireducibilny & \hyperlink{TERM-1487}{TERM-1487}\\
reducirana norma / reducirana diferenta & \hyperlink{TERM-0006}{TERM-0006}\\
redukcija & \hyperlink{TERM-1470}{TERM-1470}\\
redukcija abo razklad & \hyperlink{TERM-1472}{TERM-1472}\\
redukcija modula $(ss'u)^4$ i sistema od $s$ & \hyperlink{TERM-1471}{TERM-1471}\\
redukcije & \hyperlink{TERM-1473}{TERM-1473}\\
redukcije / reducent & \hyperlink{TERM-1474}{TERM-1474}\\
redukcije slědujut neposrědno ... črez jednokratno svijanje & \hyperlink{TERM-1475}{TERM-1475}\\
redukcijne formuly & \hyperlink{TERM-1476}{TERM-1476}\\
redukcijne teoremy & \hyperlink{TERM-1479}{TERM-1479}\\
redukcijny teorem & \hyperlink{TERM-1477}{TERM-1477}\\
redukcijske teoremy & \hyperlink{TERM-1480}{TERM-1480}\\
redukovalny / reducent & \hyperlink{TERM-1488}{TERM-1488}\\
redukovano predstavjenje & \hyperlink{TERM-1493}{TERM-1493}\\
redukovano v odnosu k komplementu / k R & \hyperlink{TERM-1492}{TERM-1492}\\
redukovany diskriminant & \hyperlink{TERM-1494}{TERM-1494}\\
redukovany slěd & \hyperlink{TERM-1495}{TERM-1495}\\
redukovany stepen & \hyperlink{TERM-1496}{TERM-1496}\\
redukovany stepenj / eksponenta & \hyperlink{TERM-1497}{TERM-1497}\\
redy po n varijablov & \hyperlink{TERM-1505}{TERM-1505}\\
rędy prěměnnyh / rędy simbolov & \hyperlink{TERM-1805}{TERM-1805}\\
rędy prěměnnyh / rędy simbolov vyšego nivoja & \hyperlink{TERM-1804}{TERM-1804}\\
rędy simbolov & \hyperlink{TERM-1681}{TERM-1681}\\
rędy simbolov / rędy prěměnnyh & \hyperlink{TERM-1682}{TERM-1682}\\
redy veličin & \hyperlink{TERM-0756}{TERM-0756}\\
rędy veličin / rędy & \hyperlink{TERM-0757}{TERM-0757}\\
regularna & \hyperlink{TERM-1499}{TERM-1499}\\
regularne sistemy iz polinomov & \hyperlink{TERM-1501}{TERM-1501}\\
regularne sistemy; bez fundamentalnyh toček & \hyperlink{TERM-1502}{TERM-1502}\\
regularno predstavjenje & \hyperlink{TERM-1500}{TERM-1500}\\
regularny element / nedělitelj nuly & \hyperlink{TERM-1503}{TERM-1503}\\
regularny polinom vzhodno k x\_i & \hyperlink{TERM-1504}{TERM-1504}\\
regularny sistem vrědnostij / singularnosti & \hyperlink{TERM-1498}{TERM-1498}\\
rekursivna konstrukcija specialnyh razvitij v rędy & \hyperlink{TERM-1515}{TERM-1515}\\
rekursivny sistem ravnanj & \hyperlink{TERM-1516}{TERM-1516}\\
rekurzivno označene bazove veličiny & \hyperlink{TERM-1513}{TERM-1513}\\
relacija & \hyperlink{TERM-1517}{TERM-1517}\\
relacija medžu razkladnymi formami & \hyperlink{TERM-1518}{TERM-1518}\\
relacije & \hyperlink{TERM-1519}{TERM-1519}\\
relacije ortogonalnosti & \hyperlink{TERM-1297}{TERM-1297}\\
relacije za razložime formy & \hyperlink{TERM-1520}{TERM-1520}\\
relativna invariantnost & \hyperlink{TERM-1538}{TERM-1538}\\
relativna norma & \hyperlink{TERM-1540}{TERM-1540}\\
relativno algebraično cěla & \hyperlink{TERM-1529}{TERM-1529}\\
relativno cěla oblast & \hyperlink{TERM-1535}{TERM-1535}\\
relativno cěle funkcije & \hyperlink{TERM-1521}{TERM-1521}, \hyperlink{TERM-1533}{TERM-1533}\\
relativno cěle funkcije / oblasti prvogo roda & \hyperlink{TERM-1534}{TERM-1534}\\
relativno cěle oblasti prvogo roda & \hyperlink{TERM-1531}{TERM-1531}\\
relativno polny sistem & \hyperlink{TERM-1523}{TERM-1523}\\
relativno polny sistem mod $(\varrho,t)$ & \hyperlink{TERM-1525}{TERM-1525}\\
relativno polny sistem mod (A,ν) & \hyperlink{TERM-1526}{TERM-1526}\\
relativno polny sistem mod (s,ρ,t) & \hyperlink{TERM-1527}{TERM-1527}\\
relativno polny sistem po modulu ... / mod ... & \hyperlink{TERM-1528}{TERM-1528}\\
relativno proste & \hyperlink{TERM-1541}{TERM-1541}\\
relativno prosty k & \hyperlink{TERM-1542}{TERM-1542}\\
relativno-cěle oblasti prvogo roda & \hyperlink{TERM-1532}{TERM-1532}\\
relativno-Galoisove ?islove polja & \hyperlink{TERM-1530}{TERM-1530}\\
relativno-prosto-nerazlozimy & \hyperlink{TERM-1543}{TERM-1543}\\
relativno-prosto-nerazložimy & \hyperlink{TERM-1544}{TERM-1544}\\
relativny / modularny / formalny invariant & \hyperlink{TERM-1539}{TERM-1539}\\
relativny diskriminant & \hyperlink{TERM-1536}{TERM-1536}\\
relativny diskriminant razširjenogo polja & \hyperlink{TERM-1537}{TERM-1537}\\
residualny modul / kvocient dvoh modulov & \hyperlink{TERM-1546}{TERM-1546}\\
rěšimo polje razloženja & \hyperlink{TERM-0146}{TERM-0146}\\
rezultant & \hyperlink{TERM-1576}{TERM-1576}\\
rezultant členov najvyššej dimenzije & \hyperlink{TERM-1577}{TERM-1577}\\
rezultant i-go nivoja / elementarny dělitelj i-go nivoja & \hyperlink{TERM-1578}{TERM-1578}\\
rezultanta forma & \hyperlink{TERM-1580}{TERM-1580}\\
rezultantna forma / forma elementarnyh děliteljev & \hyperlink{TERM-1581}{TERM-1581}\\
rezultantna forma / rezultant & \hyperlink{TERM-1582}{TERM-1582}\\
rezultantna forma i teorija eliminacije & \hyperlink{TERM-1583}{TERM-1583}\\
rezultantno značenje & \hyperlink{TERM-1579}{TERM-1579}\\
Riemannova površina / absolutna Riemannova površina & \hyperlink{TERM-1589}{TERM-1589}\\
Rimanove normalne koordinaty & \hyperlink{TERM-1590}{TERM-1590}\\
Rimsky nomer ne jest slovo & \hyperlink{DEC-0154}{DEC-0154}\\
running names Slavicized; foreign citations protected in Latin italic & \hyperlink{TERM-1403}{TERM-1403}\\
s dodatkom, společně s E. Noether & \hyperlink{TERM-1156}{TERM-1156}\\
s izključenjem & \hyperlink{TERM-1155}{TERM-1155}\\
s koeficientami iz K & \hyperlink{TERM-1157}{TERM-1157}\\
s učestjem & \hyperlink{TERM-1782}{TERM-1782}\\
s učtenjem redukcij togo paragrafa & \hyperlink{TERM-1780}{TERM-1780}\\
S(\textbackslash{}gamma) slěd / N(\textbackslash{}gamma) norma vzhodno k glavnomu klasu & \hyperlink{TERM-1609}{TERM-1609}\\
samo jedin prosty ideal & \hyperlink{TERM-1272}{TERM-1272}\\
Sběgati se, identičnost i soglasnost & \hyperlink{DEC-0027}{DEC-0027}\\
shema & \hyperlink{TERM-0137}{TERM-0137}\\
shema svijanja & \hyperlink{TERM-1622}{TERM-1622}\\
Shranjati, shraniti, shranjenje & \hyperlink{DEC-0053}{DEC-0053}\\
shvatiti komponenty v grupy & \hyperlink{TERM-1028}{TERM-1028}\\
simbolična predstavimost & \hyperlink{TERM-1677}{TERM-1677}\\
simbolične identičnosti & \hyperlink{TERM-1678}{TERM-1678}\\
simboličny metod & \hyperlink{TERM-1679}{TERM-1679}\\
simboličny ręd & \hyperlink{TERM-1680}{TERM-1680}\\
simboly theta-H / f-H & \hyperlink{TERM-1676}{TERM-1676}\\
simboly v \textbackslash{}pmod\{...\} ostavajut latinske & \hyperlink{TERM-0013}{TERM-0013}\\
simetrične funkcije redov veličin & \hyperlink{TERM-1684}{TERM-1684}\\
simultanna sistema & \hyperlink{TERM-1632}{TERM-1632}\\
simultanne invarianty & \hyperlink{TERM-1634}{TERM-1634}\\
simultanny sistem dvuh kvadratičnyh form & \hyperlink{TERM-1633}{TERM-1633}\\
singularne primarne čisla / ostatkove karaktere / stepenne ostatky & \hyperlink{TERM-1635}{TERM-1635}\\
singularne sistemy vrědnostij & \hyperlink{TERM-1637}{TERM-1637}\\
sistem $s_I$ & \hyperlink{TERM-1685}{TERM-1685}\\
sistem $s_{II}$ & \hyperlink{TERM-1686}{TERM-1686}\\
sistem $s_{IV}$ & \hyperlink{TERM-1687}{TERM-1687}\\
sistem elementarnyh děliteljev & \hyperlink{TERM-0469}{TERM-0469}\\
sistem form & \hyperlink{TERM-0595}{TERM-0595}\\
sistem form; where explanatory, sistem porodžujųćih form & \hyperlink{TERM-0596}{TERM-0596}\\
Sistem III & \hyperlink{TERM-1689}{TERM-1689}\\
sistem klasov ostatkov / komplement & \hyperlink{TERM-1688}{TERM-1688}\\
sistem koeficientov Galoisovoj rezolventy & \hyperlink{TERM-1690}{TERM-1690}\\
sistem obrazov & \hyperlink{TERM-0027}{TERM-0027}\\
sistem ostatkov / kolco ostatkovyh klasov / ostatkove klasy & \hyperlink{TERM-1574}{TERM-1574}\\
sistem ostatkovyh klasov & \hyperlink{TERM-1568}{TERM-1568}\\
sistem ostatkovyh klasov / kolco bez nulovyh děliteljev & \hyperlink{TERM-1570}{TERM-1570}\\
sistem ostatkovyh klasov / polje ostatkovyh klasov & \hyperlink{TERM-1569}{TERM-1569}\\
sistem ostatkovyh klasov prostogo čisla p & \hyperlink{TERM-1571}{TERM-1571}\\
sistem predstaviteljev & \hyperlink{TERM-1545}{TERM-1545}\\
sistem vrědnostij & \hyperlink{TERM-1872}{TERM-1872}, \hyperlink{TERM-1873}{TERM-1873}\\
sistema form & \hyperlink{TERM-0597}{TERM-0597}\\
sistemna determinanta & \hyperlink{TERM-1691}{TERM-1691}\\
sistemna matrica / grupova matrica & \hyperlink{TERM-1697}{TERM-1697}\\
sistemy \textbackslash{}Ssys\_\{n\textbackslash{}rho\} / \textbackslash{}Ssys\textasciicircum{}*\_\{\textbackslash{}rho\textbackslash{}rho\} & \hyperlink{TERM-1692}{TERM-1692}\\
sistemy blokovyh matric R\textasciicircum{}b i R\textasciicircum{}c & \hyperlink{TERM-1431}{TERM-1431}\\
sistemy linearnyh diferencialnyh uravnenj v častnyh derivatah & \hyperlink{TERM-1693}{TERM-1693}\\
sistemy parcialnyh ili obyčnyh linearnyh diferencialnyh izrazov & \hyperlink{TERM-1694}{TERM-1694}\\
sistemy rěšenj & \hyperlink{TERM-1129}{TERM-1129}\\
sistemy S\_\{n ho\} & \hyperlink{TERM-1695}{TERM-1695}\\
sistemy S\_\{n1\} & \hyperlink{TERM-1696}{TERM-1696}\\
sistemy vrědnostij & \hyperlink{TERM-1874}{TERM-1874}\\
skladanje & \hyperlink{TERM-1965}{TERM-1965}\\
sklop rěšetkovyh toček & \hyperlink{TERM-1133}{TERM-1133}\\
skobkovy faktor & \hyperlink{TERM-0977}{TERM-0977}\\
skr??eno predstavjenje & \hyperlink{TERM-1815}{TERM-1815}\\
skr??eny produkt & \hyperlink{TERM-1816}{TERM-1816}\\
skrežene produkty & \hyperlink{TERM-1817}{TERM-1817}\\
sl?d & \hyperlink{TERM-1651}{TERM-1651}\\
slaby / silny primarny ideal & \hyperlink{TERM-1631}{TERM-1631}\\
slaby / silny prosty ideal & \hyperlink{TERM-1630}{TERM-1630}\\
slěd & \hyperlink{TERM-1652}{TERM-1652}\\
slěd / matrična predstavjenost / konjugovane elementy & \hyperlink{TERM-1653}{TERM-1653}\\
slěd / slědy & \hyperlink{TERM-1654}{TERM-1654}\\
slěd elementa vzhodno k klasu & \hyperlink{TERM-1655}{TERM-1655}\\
slěd elementa, vzety vzhodno k klasu & \hyperlink{TERM-1656}{TERM-1656}\\
slěd, vzety vzhodno k glavnomu klasu & \hyperlink{TERM-0888}{TERM-0888}\\
slědujuće produkty form trěba razmatrivati & \hyperlink{TERM-0568}{TERM-0568}\\
Slědujuči i slědny & \hyperlink{DEC-0034}{DEC-0034}\\
slědujući vyši relativno polny sistem & \hyperlink{TERM-1274}{TERM-1274}\\
složene polary & \hyperlink{TERM-1962}{TERM-1962}\\
slu?aj jednozna?nogo razklada & \hyperlink{TERM-0389}{TERM-0389}\\
směšana grupa & \hyperlink{TERM-0699}{TERM-0699}\\
Sodrživati vměsto obsah- i vsebu- & \hyperlink{DEC-0028}{DEC-0028}\\
Sodrživati; sodržanje; držati & \hyperlink{DEC-0024}{DEC-0024}\\
sohranjena jest prazdnina v označenjah izvora & \hyperlink{TERM-0943}{TERM-0943}\\
sootvětni prosti idealy & \hyperlink{TERM-1938}{TERM-1938}\\
sootvětstvo medžu razloženjem Tp i razloženjem nulovogo ideala v R & \hyperlink{TERM-0501}{TERM-0501}\\
sostavimost iz osnovnyh svijanj & \hyperlink{TERM-1964}{TERM-1964}\\
sostavne identičnosti & \hyperlink{TERM-1961}{TERM-1961}\\
Sostojati se iz: cělo i jego česti & \hyperlink{DEC-0142}{DEC-0142}\\
Sostojati se iz: sostav & \hyperlink{DEC-0014}{DEC-0014}\\
Sovmestno s R. Brauerom & \hyperlink{TERM-0691}{TERM-0691}\\
sovpadajuča razširjenja v vsih kolcah kvocientov daju identične idealy & \hyperlink{TERM-1610}{TERM-1610}\\
sovršeno / nesovršeno polje & \hyperlink{TERM-1835}{TERM-1835}\\
sovršeno / nesovršeno polje koeficientov & \hyperlink{TERM-1839}{TERM-1839}\\
sovršeno osnovno polje / sovršena osnovna oblast & \hyperlink{TERM-1836}{TERM-1836}\\
sovršeno polje & \hyperlink{TERM-1837}{TERM-1837}\\
sovršeno/nesovršeno polje & \hyperlink{TERM-1838}{TERM-1838}\\
specializacija & \hyperlink{TERM-1643}{TERM-1643}\\
specializacija / specializovati & \hyperlink{TERM-1644}{TERM-1644}\\
specializacija na relativne diskriminanty čislovyh polj & \hyperlink{TERM-1645}{TERM-1645}\\
specializacija někojih neopreděljenyh & \hyperlink{TERM-1646}{TERM-1646}\\
specializacija u\_\{mu nu\}=ubar\_\{mu nu\} & \hyperlink{TERM-1647}{TERM-1647}\\
specialna modulna baza / dopolniti libovoljnu bazu od p & \hyperlink{TERM-1648}{TERM-1648}\\
specialne P-modulne bazy & \hyperlink{TERM-1649}{TERM-1649}\\
specialno dvojnorazovo svijanje & \hyperlink{TERM-1650}{TERM-1650}\\
srav. konec § 5 & \hyperlink{TERM-1825}{TERM-1825}\\
Sravniti; identifikovati; prirěditi & \hyperlink{DEC-0026}{DEC-0026}\\
Srodne formy ne imajut jednu gramatiku & \hyperlink{DEC-0046}{DEC-0046}\\
stavaje reducibilny / reducent & \hyperlink{TERM-1489}{TERM-1489}\\
stavajut reducibilne & \hyperlink{TERM-1871}{TERM-1871}\\
Štejnicova terminologija / razširjenje prvogo ili drugogo roda & \hyperlink{TERM-1658}{TERM-1658}\\
stepen & \hyperlink{TERM-0287}{TERM-0287}, \hyperlink{TERM-0288}{TERM-0288}\\
stepen komplementa / dělimy črez primarny ideal & \hyperlink{TERM-1021}{TERM-1021}\\
stepeni / ? kako najve?a med stepenami & \hyperlink{TERM-0725}{TERM-0725}\\
stepeni v rędah prěmennyh & \hyperlink{TERM-0724}{TERM-0724}\\
stepenj po prěměnnyh & \hyperlink{TERM-0723}{TERM-0723}\\
stepenj prostogo ideala & \hyperlink{TERM-0722}{TERM-0722}\\
stepenj svojego asociovanogo prostogo ideala & \hyperlink{TERM-1330}{TERM-1330}\\
stepenj svojego asociovanogo prostogo ideala / produkt stepenjev & \hyperlink{TERM-1331}{TERM-1331}\\
stepenj transcendencije & \hyperlink{TERM-1746}{TERM-1746}\\
stepenna suma & \hyperlink{TERM-1340}{TERM-1340}\\
stepenove produkty & \hyperlink{TERM-1337}{TERM-1337}\\
stepenove produkty q-oj stepeni & \hyperlink{TERM-1339}{TERM-1339}\\
stepenovy produkt & \hyperlink{TERM-1333}{TERM-1333}\\
što jest vsegda možno & \hyperlink{TERM-1868}{TERM-1868}\\
Što, ktory, to i tym: referencija v prědloženju & \hyperlink{DEC-0124}{DEC-0124}\\
strukturne teoremy & \hyperlink{TERM-1664}{TERM-1664}\\
substitucija (35) & \hyperlink{TERM-1665}{TERM-1665}\\
substitucija novih rędov & \hyperlink{TERM-1666}{TERM-1666}\\
sukcesivna čislova specializacija & \hyperlink{TERM-1670}{TERM-1670}\\
sukcesivna eliminacija & \hyperlink{TERM-1668}{TERM-1668}\\
sukcesivna substitucija & \hyperlink{TERM-1669}{TERM-1669}\\
sukcesivno razloženje & \hyperlink{TERM-1671}{TERM-1671}\\
sukcesivny razklad v dv? sostavne ?esti & \hyperlink{TERM-1672}{TERM-1672}\\
suplementne čisla & \hyperlink{TERM-1674}{TERM-1674}\\
Šurov indeks & \hyperlink{TERM-1629}{TERM-1629}\\
suštno različne & \hyperlink{TERM-1875}{TERM-1875}\\
suštstvene parametry / libovoljne funkcije & \hyperlink{TERM-1876}{TERM-1876}\\
sut dělimy tymi samymi prostymi idealami p i v tyh samyh stepenjah & \hyperlink{TERM-0358}{TERM-0358}\\
sut izpušćene & \hyperlink{TERM-0153}{TERM-0153}\\
Sut pravdive: istinnost tvrdženij & \hyperlink{DEC-0004}{DEC-0004}\\
suženy ideal & \hyperlink{TERM-1809}{TERM-1809}\\
svedenje modula (A,ν) k modulu (ν) & \hyperlink{TERM-1957}{TERM-1957}\\
svedenje modula ... k ... & \hyperlink{TERM-1958}{TERM-1958}\\
svęz medžu formami & \hyperlink{TERM-0598}{TERM-0598}\\
svijanja s modulom & \hyperlink{TERM-0562}{TERM-0562}\\
svijanje & \hyperlink{TERM-0549}{TERM-0549}, \hyperlink{TERM-0554}{TERM-0554}\\
svijanje $s$ s $\nu; H; L; H^2$ & \hyperlink{TERM-0552}{TERM-0552}\\
svijanje $s$ s ... & \hyperlink{TERM-0553}{TERM-0553}\\
svijanje / skruženje & \hyperlink{TERM-0550}{TERM-0550}\\
svijanje f s ν & \hyperlink{TERM-0555}{TERM-0555}\\
svijanje K,j,\textbackslash{}A s \textbackslash{}nu; f s H & \hyperlink{TERM-0556}{TERM-0556}\\
svijanje s $s$ samo kako polarizacija & \hyperlink{TERM-0551}{TERM-0551}\\
svijanje s s formami 4. poręda & \hyperlink{TERM-0557}{TERM-0557}\\
svijanje s s formami 5. poręda i s produktom & \hyperlink{TERM-0559}{TERM-0559}\\
svijanje s s formami 6. poręda & \hyperlink{TERM-0560}{TERM-0560}\\
svijanje s s formami 7. poręda & \hyperlink{TERM-0558}{TERM-0558}\\
svijanje θ s ν & \hyperlink{TERM-0561}{TERM-0561}\\
Svjazane imena dvuh avtorov & \hyperlink{DEC-0162}{DEC-0162}\\
svobodne od x\_1,...,x\_\{i-1\} & \hyperlink{TERM-0601}{TERM-0601}\\
Sylowovo polje & \hyperlink{TERM-1675}{TERM-1675}\\
syzygije & \hyperlink{TERM-1698}{TERM-1698}\\
t?lo; komutativno polje & \hyperlink{TERM-0970}{TERM-0970}\\
tabela & \hyperlink{TERM-1700}{TERM-1700}\\
tablica množenja & \hyperlink{TERM-1189}{TERM-1189}\\
Tako že i takože: slědstvo i dodavanje & \hyperlink{DEC-0122}{DEC-0122}\\
te same razmysljenja važe za multiplikacijsko kolco & \hyperlink{TERM-0274}{TERM-0274}\\
tělo & \hyperlink{TERM-1067}{TERM-1067}, \hyperlink{TERM-1068}{TERM-1068}, \hyperlink{TERM-1069}{TERM-1069}\\
tělo avtomorfizmov & \hyperlink{TERM-0157}{TERM-0157}\\
tělo avtomorfizmov K\_ν & \hyperlink{TERM-0158}{TERM-0158}\\
ten sam ostatok / polny prěsěk & \hyperlink{TERM-0715}{TERM-0715}\\
teorem & \hyperlink{TERM-1611}{TERM-1611}\\
teorem Dedekinda--Mertensa & \hyperlink{TERM-0247}{TERM-0247}\\
teorem Jordan--Hölder & \hyperlink{TERM-0967}{TERM-0967}\\
teorem konečnosti invariantov & \hyperlink{TERM-0499}{TERM-0499}\\
teorem Kronekera / relativno-Abelove polja & \hyperlink{TERM-1062}{TERM-1062}\\
teorem Lurota / Kastelnuova / Enrikesa & \hyperlink{TERM-1616}{TERM-1616}\\
teorem o cěpah děliteljev & \hyperlink{TERM-1728}{TERM-1728}\\
teorem o diskriminantu & \hyperlink{TERM-0331}{TERM-0331}\\
teorem o dvojnoj cěpi & \hyperlink{TERM-0339}{TERM-0339}\\
teorem o glavnom rodu & \hyperlink{TERM-0785}{TERM-0785}\\
teorem o homomorfizmu za kolca & \hyperlink{TERM-0837}{TERM-0837}\\
teorem o izomorfiji & \hyperlink{TERM-1615}{TERM-1615}\\
teorem o konečnoj cěpi & \hyperlink{TERM-1614}{TERM-1614}\\
teorem o konečnosti & \hyperlink{TERM-0498}{TERM-0498}\\
teorem o ostatku / ostatkova grupa / koresidualny & \hyperlink{TERM-1573}{TERM-1573}\\
teorem obstajanja Hassea & \hyperlink{TERM-0781}{TERM-0781}\\
teorem Ostrowski & \hyperlink{TERM-1617}{TERM-1617}\\
teorem redukcije & \hyperlink{TERM-1478}{TERM-1478}, \hyperlink{TERM-1481}{TERM-1481}\\
teorem Riemanna--Rocha / redukcijny teorem & \hyperlink{TERM-1588}{TERM-1588}\\
Teorem VI / Teorem VII & \hyperlink{TERM-1613}{TERM-1613}\\
Teorem XI / Teorem XII & \hyperlink{TERM-1619}{TERM-1619}\\
teorema Burnsida / obobčena teorema Burnsida & \hyperlink{TERM-0199}{TERM-0199}\\
Teorema I & \hyperlink{TERM-1612}{TERM-1612}\\
teorema identičnosti & \hyperlink{TERM-0883}{TERM-0883}\\
teorema jednoznačnosti i uporjadženja & \hyperlink{TERM-0402}{TERM-0402}\\
teorema o dobrom uporjadkovanju & \hyperlink{TERM-1885}{TERM-1885}\\
teorema produkta & \hyperlink{TERM-1395}{TERM-1395}\\
teorema produkta v najobćejšej formě & \hyperlink{TERM-1399}{TERM-1399}\\
teorema produkta za determinanty & \hyperlink{TERM-1396}{TERM-1396}\\
teorema produkta za determinanty / matrice & \hyperlink{TERM-1397}{TERM-1397}\\
teorema produkta za matrice & \hyperlink{TERM-1398}{TERM-1398}\\
teorema razloženja & \hyperlink{TERM-1929}{TERM-1929}\\
Teorema X & \hyperlink{TERM-1618}{TERM-1618}\\
teoremy konečnosti & \hyperlink{TERM-0496}{TERM-0496}\\
teoremy konečnosti / Hilbertov teorem o bazisu modula & \hyperlink{TERM-0497}{TERM-0497}\\
teoremy o bazisu / fundamentalna sistema & \hyperlink{TERM-0171}{TERM-0171}\\
teoremy o izomorfiji pri aditivnom razkladu & \hyperlink{TERM-0958}{TERM-0958}\\
teoremy o izomorfiji; direktne sumy & \hyperlink{TERM-0959}{TERM-0959}\\
teoremy o razloženju norm i modulov & \hyperlink{TERM-1931}{TERM-1931}\\
teorija cělyh elementov & \hyperlink{TERM-1737}{TERM-1737}\\
teorija čislovyh polj / algebraično čislovo polje & \hyperlink{TERM-1895}{TERM-1895}\\
teorija elementarnyh děliteljev & \hyperlink{TERM-0470}{TERM-0470}\\
teorija eliminacije & \hyperlink{TERM-0473}{TERM-0473}\\
teorija idealov v T & \hyperlink{TERM-0869}{TERM-0869}\\
teorija modulov / bazis modula / prosty modul / nulovy modul & \hyperlink{TERM-1186}{TERM-1186}\\
teorija predstavjenj & \hyperlink{TERM-0240}{TERM-0240}\\
teorija razvětvenja / ideal razvětvenja & \hyperlink{TERM-1823}{TERM-1823}\\
teorija razvětvljenja / ideal razvětvljenja & \hyperlink{TERM-1824}{TERM-1824}\\
ternarna bikvadratična forma; explanatory first title phrase: ternarna forma četvrtogo stepena & \hyperlink{TERM-1735}{TERM-1735}\\
ternarna forma četvrtogo stepena; kratko: ternarna kvartika & \hyperlink{TERM-1736}{TERM-1736}\\
to jest & \hyperlink{TERM-0242}{TERM-0242}\\
To jest pri formulah i definicijah & \hyperlink{DEC-0079}{DEC-0079}\\
To jest v algebričnyh definicijah & \hyperlink{DEC-0070}{DEC-0070}\\
To jest: ekvivalencija i preformulovanje & \hyperlink{DEC-0069}{DEC-0069}\\
togda i samo togda & \hyperlink{TERM-0225}{TERM-0225}\\
togda o jest kolco glavnyh idealov & \hyperlink{TERM-1639}{TERM-1639}\\
Toj, togo; sej i tutoj & \hyperlink{DEC-0035}{DEC-0035}\\
totalno vzaimno prosta sistema & \hyperlink{TERM-1738}{TERM-1738}\\
transcendencijny stepen & \hyperlink{TERM-1747}{TERM-1747}\\
transcendentny stepen / medžupolje & \hyperlink{TERM-1748}{TERM-1748}\\
transformacijne veli?iny & \hyperlink{TERM-1741}{TERM-1741}\\
transformovane koeficienty & \hyperlink{TERM-1744}{TERM-1744}\\
transformovany ideal & \hyperlink{TERM-1743}{TERM-1743}, \hyperlink{TERM-1745}{TERM-1745}\\
transvekcija & \hyperlink{TERM-2006}{TERM-2006}\\
transvekcija črez funkcionalne determinanty & \hyperlink{TERM-2008}{TERM-2008}\\
transvekcija črez funkcionalny determinant & \hyperlink{TERM-2007}{TERM-2007}\\
transvekcije nad funkcionalnymi determinantami & \hyperlink{TERM-2010}{TERM-2010}\\
Tretji, treťji, trětji & \hyperlink{DEC-0110}{DEC-0110}\\
tridimensionalny prostor & \hyperlink{TERM-0347}{TERM-0347}\\
trivialno predstavjenje & \hyperlink{TERM-1752}{TERM-1752}\\
Tu, ta i takože & \hyperlink{DEC-0106}{DEC-0106}\\
Tvoret, zaviset, vhodet, vodet & \hyperlink{DEC-0063}{DEC-0063}\\
tvorjenje kvocientov & \hyperlink{TERM-1418}{TERM-1418}\\
tvorjenje prěsěkov & \hyperlink{TERM-0368}{TERM-0368}\\
tvorjut Z-modul / modulno svojstvo & \hyperlink{TERM-0187}{TERM-0187}\\
tvrdženje & \hyperlink{TERM-1756}{TERM-1756}\\
Tyh: pokazateljno, ne tihy & \hyperlink{DEC-0094}{DEC-0094}\\
tym je dokazana & \hyperlink{TERM-0962}{TERM-0962}\\
Tym samym pri pričinnom objašnjenju & \hyperlink{DEC-0089}{DEC-0089}\\
Tym teorem jest dokazany & \hyperlink{TERM-0224}{TERM-0224}\\
umenšenje vagi & \hyperlink{TERM-0710}{TERM-0710}\\
umenšenje vagi / vkupna vaga & \hyperlink{TERM-0711}{TERM-0711}\\
unimodularna matrica & \hyperlink{TERM-1774}{TERM-1774}\\
unimodularna transformacija & \hyperlink{TERM-1775}{TERM-1775}\\
unimodularny & \hyperlink{TERM-1773}{TERM-1773}\\
uravnenja s zadanoj grupoj & \hyperlink{TERM-0718}{TERM-0718}\\
uravnenje najnižšego stepena & \hyperlink{TERM-0717}{TERM-0717}\\
uravnenje nerazlo?imo vzhodno k [H] & \hyperlink{TERM-0886}{TERM-0886}\\
uslovja cěpov v konečnyh modulnyh oblastjah & \hyperlink{TERM-0976}{TERM-0976}\\
uslovja integrabilnosti & \hyperlink{TERM-0900}{TERM-0900}\\
uslovja mnogokratnosti / mnogokratnost & \hyperlink{TERM-1192}{TERM-1192}\\
uslovja věsy / uslovja simetrije & \hyperlink{TERM-0712}{TERM-0712}\\
uslovja vyšljego defekta & \hyperlink{TERM-0172}{TERM-0172}\\
uslovje cěpov děliteljev & \hyperlink{TERM-1729}{TERM-1729}\\
uslovje cěpov děliteljev / dobry porjadok / postulat izbora & \hyperlink{TERM-1732}{TERM-1732}\\
uslovje cěpov děliteljev / fiksovany dobry porjadok / konečna baza ideala & \hyperlink{TERM-1730}{TERM-1730}\\
uslovje cěpov děliteljev / uslovje cěpov kratnyh & \hyperlink{TERM-1731}{TERM-1731}\\
uslovje cěpov kratnyh & \hyperlink{TERM-1829}{TERM-1829}\\
uslovje cěpov kratnyh modulo vsaky nenulovy ideal & \hyperlink{TERM-1827}{TERM-1827}\\
uslovje homogennosti / homogenizovanje & \hyperlink{TERM-0834}{TERM-0834}\\
uslovje konečnosti / idealny bazis & \hyperlink{TERM-0492}{TERM-0492}\\
uslovje konečnosti; konečny modul; modulny bazis & \hyperlink{TERM-0493}{TERM-0493}\\
uslovje maksimalnosti & \hyperlink{TERM-1144}{TERM-1144}\\
uslovje maksimalnosti / minimalnosti za prave idealy & \hyperlink{TERM-1145}{TERM-1145}\\
uslovje minimalnosti & \hyperlink{TERM-1152}{TERM-1152}\\
Uslovje: neutrum i soglasje & \hyperlink{DEC-0149}{DEC-0149}\\
uslovne uravnenja & \hyperlink{TERM-0173}{TERM-0173}\\
Uvrstjenje hiperkompleksnyh sistemov v teoriju & \hyperlink{TERM-0447}{TERM-0447}\\
v imenovatelju & \hyperlink{TERM-0884}{TERM-0884}\\
v polinomnom smyslu & \hyperlink{TERM-1328}{TERM-1328}\\
v rabotě citovanoj v primětce 6, str. 83 & \hyperlink{TERM-0893}{TERM-0893}\\
v stepenjah ne vhodi v svijanje & \hyperlink{TERM-1751}{TERM-1751}\\
V tom slučaju: uslovna uloga Dann & \hyperlink{DEC-0086}{DEC-0086}\\
vaga rěda / indeks vagy & \hyperlink{TERM-0709}{TERM-0709}\\
Veber / E. Štejnic & \hyperlink{TERM-1869}{TERM-1869}\\
veličiny y, z, vartheta & \hyperlink{TERM-0754}{TERM-0754}\\
vez medžu idealami iz polinomov i modulami iz linearnyh form & \hyperlink{TERM-1963}{TERM-1963}\\
vhodet retained beside vhodi & \hyperlink{TERM-1826}{TERM-1826}\\
vhodi v diskriminantny ideal & \hyperlink{TERM-0690}{TERM-0690}\\
vhodi v svijanje & \hyperlink{TERM-1749}{TERM-1749}\\
vkupna vaga formy & \hyperlink{TERM-0702}{TERM-0702}\\
vlastna podgrupa & \hyperlink{TERM-0372}{TERM-0372}\\
vlastna teorija eliminacije & \hyperlink{TERM-0382}{TERM-0382}\\
vlastne / nevlastne divergencijne relacije & \hyperlink{TERM-0381}{TERM-0381}\\
vlastne vrědnosti i vlastne funkcije & \hyperlink{TERM-0388}{TERM-0388}\\
vlastno podkolco & \hyperlink{TERM-0378}{TERM-0378}\\
vlastno primarna komponenta & \hyperlink{TERM-0379}{TERM-0379}, \hyperlink{TERM-0384}{TERM-0384}\\
vlastno primarna komponenta / prosty ideal drugogo roda & \hyperlink{TERM-0385}{TERM-0385}\\
vlastno primarne idealy abo proste idealy & \hyperlink{TERM-0383}{TERM-0383}\\
vlastno primarno kolco & \hyperlink{TERM-0380}{TERM-0380}\\
vlastny dělitelj & \hyperlink{TERM-0373}{TERM-0373}\\
vlastny dělitelj / čislo členov & \hyperlink{TERM-0376}{TERM-0376}\\
vlastny dělitelj / prěsěčna klasa & \hyperlink{TERM-0374}{TERM-0374}\\
vlastny dělitelj / vlastno kratno & \hyperlink{TERM-0375}{TERM-0375}\\
vlastny dělitelj jednakogo ranga & \hyperlink{TERM-0377}{TERM-0377}\\
vlastny primarny ideal & \hyperlink{TERM-0386}{TERM-0386}\\
vlastny primarny ideal / vlastno primarna komponenta & \hyperlink{TERM-0387}{TERM-0387}\\
Vměsto: zaměna, ne město & \hyperlink{DEC-0100}{DEC-0100}\\
vnutrišnje avtomorfizmy & \hyperlink{TERM-0899}{TERM-0899}\\
vrědnosti diskontinuirnosti & \hyperlink{TERM-1779}{TERM-1779}\\
vsaka potenca indeksa / vsaky parny stepen & \hyperlink{TERM-0966}{TERM-0966}\\
Vsaky i každy; each ne jest prosto all & \hyperlink{DEC-0060}{DEC-0060}\\
vsaky prosty element p, ktory vhodi v D\_T, i samo takovy element & \hyperlink{TERM-0965}{TERM-0965}\\
vse parne cěle racionalne čisla & \hyperlink{TERM-0118}{TERM-0118}\\
vse sistemy iz n elementov iz T & \hyperlink{TERM-0119}{TERM-0119}\\
Vsegda & \hyperlink{DEC-0033}{DEC-0033}\\
vsegda obstojna o-modulna baza & \hyperlink{TERM-1661}{TERM-1661}\\
vsi ostali idealy sut glavne idealy & \hyperlink{TERM-0120}{TERM-0120}\\
Vsih, vsim, vsimi: vseobče otnošenje & \hyperlink{DEC-0057}{DEC-0057}\\
Vstavjenje osnovnyh idealov i jih komponentov v obču teoriju idealov & \hyperlink{TERM-0446}{TERM-0446}\\
vyčerpana jest cělost nerazložimyh identičnostij & \hyperlink{TERM-0706}{TERM-0706}\\
vyša forma & \hyperlink{TERM-0850}{TERM-0850}\\
vyše formy & \hyperlink{TERM-0851}{TERM-0851}\\
Vyše i najvyše: sravnenje i gorna granica & \hyperlink{DEC-0080}{DEC-0080}\\
vyše podana shema ... dokazana & \hyperlink{TERM-0243}{TERM-0243}\\
vyše razvětvjenje & \hyperlink{TERM-0853}{TERM-0853}\\
vyšše formy / vyšše samosvijanje & \hyperlink{TERM-0852}{TERM-0852}\\
vyšši modul & \hyperlink{TERM-0854}{TERM-0854}\\
Vyšših: korenj komparativa & \hyperlink{DEC-0101}{DEC-0101}\\
Vyvod i fraza na že & \hyperlink{DEC-0083}{DEC-0083}\\
Vyvod, ekvivalencija, uslovje: tri ulogi & \hyperlink{DEC-0093}{DEC-0093}\\
vyznačena baza & \hyperlink{TERM-0154}{TERM-0154}\\
vzaimna prostota & \hyperlink{TERM-1725}{TERM-1725}\\
vzaimna zavisimost form & \hyperlink{TERM-0689}{TERM-0689}\\
vzaimno jednoznačno & \hyperlink{TERM-0391}{TERM-0391}\\
vzaimno jednoznačno odobraženje & \hyperlink{TERM-0393}{TERM-0393}\\
vzaimno jednoznačno prirędženy & \hyperlink{TERM-0392}{TERM-0392}\\
vzaimno jednoznačno sootvětujuča baza & \hyperlink{TERM-0413}{TERM-0413}\\
vzaimno jednoznačny & \hyperlink{TERM-0407}{TERM-0407}\\
vzaimno proste i vzaimno-prosto-nerazložime klasy & \hyperlink{TERM-1724}{TERM-1724}\\
vzaimno proste idealy / jediničny ideal & \hyperlink{TERM-1721}{TERM-1721}\\
vzaimno proste polinomy & \hyperlink{TERM-1722}{TERM-1722}\\
vzaimno proste v polinomnom smyslu & \hyperlink{TERM-1522}{TERM-1522}\\
vzaimno prosty & \hyperlink{TERM-1713}{TERM-1713}\\
vzaimno prosty / pravi dělitelji & \hyperlink{TERM-1714}{TERM-1714}\\
vzaimno prosty s modulom & \hyperlink{TERM-1716}{TERM-1716}\\
vzaimno relativno proste & \hyperlink{TERM-0688}{TERM-0688}\\
vzaimno-jednoznačno & \hyperlink{TERM-0408}{TERM-0408}\\
vzaimno-prosto nerazložimy / relativno-prosto nerazložimy & \hyperlink{TERM-1719}{TERM-1719}\\
vzaimno-prosto-nerazlozime idealy v nekomutativnom kolcu s jedinicoju & \hyperlink{TERM-1720}{TERM-1720}\\
vzaimno-prosto-nerazlozimy & \hyperlink{TERM-1717}{TERM-1717}\\
vzaimno-prosto-nerazložimy & \hyperlink{TERM-1718}{TERM-1718}\\
vzaimnosti & \hyperlink{TERM-1586}{TERM-1586}\\
vzajemno proste primitivne polinomy & \hyperlink{TERM-1723}{TERM-1723}\\
vzajemno prosty v J & \hyperlink{TERM-1715}{TERM-1715}\\
Vzajemny, vzajemno, vzajemnost & \hyperlink{DEC-0029}{DEC-0029}\\
vzete formy & \hyperlink{TERM-2005}{TERM-2005}\\
vzete pojedinačno & \hyperlink{TERM-0454}{TERM-0454}\\
vzhodno k glavnomu klasu & \hyperlink{TERM-0887}{TERM-0887}\\
vzhodno k H & \hyperlink{TERM-0889}{TERM-0889}\\
Y gubi svojstvo 1 & \hyperlink{TERM-1888}{TERM-1888}\\
Za dokaz sravn. H. Grell. & \hyperlink{TERM-1952}{TERM-1952}\\
za kontrolu računa & \hyperlink{TERM-1955}{TERM-1955}\\
začetna shema & \hyperlink{TERM-0138}{TERM-0138}\\
zahvaljujem Hentzeltu za ... zaměčenja & \hyperlink{TERM-1806}{TERM-1806}\\
zaključenje & \hyperlink{TERM-1625}{TERM-1625}\\
zaključny izraz & \hyperlink{TERM-1623}{TERM-1623}\\
zakon neograničenogo i jednoznačnogo odčitanja & \hyperlink{TERM-0707}{TERM-0707}\\
zaměna simbolov & \hyperlink{TERM-1819}{TERM-1819}\\
zaradi izomorfizma kolc klasov ostatkov & \hyperlink{TERM-1870}{TERM-1870}\\
Zaradi: příčina, ne cilj & \hyperlink{DEC-0141}{DEC-0141}\\
Zaštita citatnogo prezimena pri proměně pisma & \hyperlink{DEC-0055}{DEC-0055}\\
Zato kako vyvod & \hyperlink{DEC-0064}{DEC-0064}\\
Zato najprvo treba razsmotriti diskriminantne idealy primarnyh kolc. & \hyperlink{TERM-0534}{TERM-0534}\\
Zato požadanje znači, že & \hyperlink{DEC-0084}{DEC-0084}\\
Zato v dokazě & \hyperlink{DEC-0081}{DEC-0081}\\
Zato že; pričina ne jest dopuščenje & \hyperlink{DEC-0012}{DEC-0012}\\
Zato, to jest i tym samym & \hyperlink{DEC-0074}{DEC-0074}\\
Zato; poslědovateljno; zato že & \hyperlink{DEC-0044}{DEC-0044}\\
zatvorjena vzhodno k tvorenju pr?s?kov & \hyperlink{TERM-0033}{TERM-0033}\\
zavisnosti medžu Lagranževymi izrazami & \hyperlink{TERM-0034}{TERM-0034}\\
Zermelova oblast & \hyperlink{TERM-1932}{TERM-1932}\\
Znak \textbackslash{}equiv znači ... & \hyperlink{TERM-0244}{TERM-0244}\\
Znak prěd jednoj tretinoju & \hyperlink{DEC-0125}{DEC-0125}\\
Znakah, krokah, izrazah: lokativ množiny & \hyperlink{DEC-0099}{DEC-0099}\\
zvedenje procesa svijanja na osnovne svijanja & \hyperlink{TERM-1956}{TERM-1956}\\
|S(alpha\_i alpha\_j)| & \hyperlink{TERM-1989}{TERM-1989}\\
§ 12 C) & \hyperlink{TERM-1991}{TERM-1991}\\
«iracionalny» / «racionalny» směr & \hyperlink{TERM-0927}{TERM-0927}\\
\end{longtable}\normalsize

\section{Termini i dokumentovane alternativy / Terms and documented alternatives}
Zapisy razděljujut iztočny pojem, zabilěženu medžuslovjansku formu, ruske i ukrajinske sravnjenja i razlog izbora. Formy iz ranogo prěvoda sut historijne prědlogy, ne avtomatično aktualne rekomendacije. Sravniteljna forma v glosaru ne dokazuje sama po sobě konsultaciju konkretnogo odbornogo iztočnika. Kde razlog ili grafija ne sut zachovane, objasnjenje to izrično govori.
\par\smallskip\noindent\emph{EN:} Each entry distinguishes the source concept, recorded Interslavic proposal, Russian and Ukrainian comparanda, and reason for the choice. Early proposals are not automatically current recommendations. A comparison form in a glossary does not by itself establish consultation of a particular scholarly source. Missing reasoning and damaged spelling are explicitly identified.
\small
\par\medskip\Needspace{7\baselineskip}
\noindent\begin{minipage}{\linewidth}
\hypertarget{TERM-0001}{}
\noindent\textbf{forma $\mu$-go stepena po $x$ i $\lambda$-go stepena po $u$}\quad{\footnotesize\texttt{TERM-0001}}\par
\noindent\emph{Iztočnik / Source:} $(\mu,\lambda)$ eine Form $\mu$ten Grades in $x$, $\lambda$ten Grades in $u$\par
\noindent\emph{Rusky / Russian:} форма $\mu$-й степени по $x$ и $\lambda$-й степени по $u$\par
\noindent\emph{Ukrajinsky / Ukrainian:} форма $\mu$-го степеня за $x$ і $\lambda$-го степеня за $u$\par
\smallskip\noindent\emph{Razlog:} Forma ima dva nezavisna stepena: μ po proměnnoj x i λ po proměnnoj u. Prěvod izgovarjaje oba odnosa, aby čitatelj ne směšal bidegree s jednym ukupnym stepenem. Zapis -go pri simbolu označaje porědkovo čislo; arhivny razlog go predstavlja kako praktičnu pisemnu konvenciju, ne kako dokazanu jedinu mogućnost.\par
\noindent\emph{Reason:} The form has two independent degrees: μ in x and λ in u. The wording states both relations rather than suggesting one total degree. The suffix -go after a symbol was recorded as a practical way of writing the ordinal, not as an independently established unique solution.\par
\noindent{\footnotesize Izvorne zapisy / Original records: HIST-g019-012.}\par
\end{minipage}\par
\par\medskip\Needspace{7\baselineskip}
\noindent\begin{minipage}{\linewidth}
\hypertarget{TERM-0002}{}
\noindent\textbf{$\lambda$ jest proizvoljno}\quad{\footnotesize\texttt{TERM-0002}}\par
\noindent\emph{Iztočnik / Source:} $\lambda$ beliebig\par
\noindent\emph{Rusky / Russian:} $\lambda$ произвольно\par
\noindent\emph{Ukrajinsky / Ukrainian:} $\lambda$ довільне\par
\smallskip\noindent\emph{Razlog:} Proizvoljno označaje, že λ ne jest ograničeno osobnym izborom. Zapis se opira na obyčno matematično vyraženje v ruskom i ukrajinskom i na doslědnost s upotrěbljenjem v oddělu 18. To jest razlog registra i doslědnosti, ne čisleno sravnjenje vseh slovjanskyh jezykov.\par
\noindent\emph{Reason:} Proizvoljno expresses that λ is arbitrary. The recorded reasons are ordinary Russian and Ukrainian mathematical usage and consistency with Section 18. This is a register-and-consistency argument, not a numerical comparison across all Slavic languages.\par
\noindent{\footnotesize Izvorne zapisy / Original records: HIST-g020-009.}\par
\end{minipage}\par
\par\medskip\Needspace{7\baselineskip}
\noindent\begin{minipage}{\linewidth}
\hypertarget{TERM-0003}{}
\noindent\textbf{$n$-e korenje iz jedinice}\quad{\footnotesize\texttt{TERM-0003}}\par
\noindent\emph{Iztočnik / Source:} $n$te Einheitswurzeln\par
\noindent\emph{Rusky / Russian:} $n$-е корни из единицы\par
\noindent\emph{Ukrajinsky / Ukrainian:} $n$-ті корені з одиниці\par
\smallskip\noindent\emph{Razlog:} Korenje iz jedinice jest prěmy algebraičny izraz. Zapis upotrěbjaje poznate sostavne slova i ne vidi potrěbu tvoriti osobny novy termin. Čislo n ostavaje čestju matematičnogo opreděljenja; zapis sam ne dokazuje, že vsaka padežna forma jest normativno prověrena.\par
\noindent\emph{Reason:} The proposal uses the familiar algebraic expression for roots of unity rather than coining a separate term. The index n remains part of the mathematical specification. The short original reason does not establish that every inflected form has been independently checked.\par
\noindent{\footnotesize Izvorne zapisy / Original records: HIST-g039-016.}\par
\end{minipage}\par
\par\medskip\Needspace{7\baselineskip}
\noindent\begin{minipage}{\linewidth}
\hypertarget{TERM-0004}{}
\noindent\textbf{konečno algebraično razširenje nad}\quad{\footnotesize\texttt{TERM-0004}}\par
\noindent\emph{Iztočnik / Source:} (endlicher) algebraischer Körper über\par
\noindent\emph{Rusky / Russian:} конечное алгебраическое расширение над\par
\noindent\emph{Ukrajinsky / Ukrainian:} скінченне алгебраїчне розширення над\par
\smallskip\noindent\emph{Razlog:} Razširenje nad izjasnjaje odnos algebraičnogo polja k osnovnomu polju v sučasnom matematičnom registru. Konečno zachovuje ograničenje iz iztočnika. To jest izrično terminologično osuvrěmenjenje: v diplomatskom prěvodu trěba prověriti, že ono ne pribavjaje ili ne gubi nijednu hipotezu.\par
\noindent\emph{Reason:} The recorded choice expresses the field-over-base relation as a finite algebraic extension, using modern mathematical terminology. It is an interpretive modernization, so a diplomatic translation must check that it neither adds nor removes a hypothesis.\par
\noindent{\footnotesize Izvorne zapisy / Original records: HIST-g044-010.}\par
\end{minipage}\par
\par\medskip\Needspace{7\baselineskip}
\noindent\begin{minipage}{\linewidth}
\hypertarget{TERM-0005}{}
\noindent\textbf{(p)\textasciicircum{}alpha abo jediničny ideal}\quad{\footnotesize\texttt{TERM-0005}}\par
\noindent\emph{Iztočnik / Source:} (p)\textasciicircum{}alpha oder Einheitsideal\par
\noindent\emph{Rusky / Russian:} (p)\textasciicircum{}alpha или единичный идеал\par
\noindent\emph{Ukrajinsky / Ukrainian:} (p)\textasciicircum{}alpha або одиничний ідеал\par
\smallskip\noindent\emph{Razlog:} Izraz zachovuje dvě alternativy: stepen lokalnogo prostogo generatora i jediničny ideal. Razlog se odnosi na tų matematičnu razliku. Sam zapis ne dava nezavisnogo opravdanja za veznik abo; on zato ne dostaje za rekomendaciju vsej historijnoj formy.\par
\noindent\emph{Reason:} The reason preserves the dichotomy between a power of the local prime generator and the unit ideal. It concerns that mathematical distinction. It does not independently justify the conjunction abo, so the record alone cannot recommend the entire historical wording.\par
\noindent{\footnotesize Izvorne zapisy / Original records: HIST-g160-009.}\par
\end{minipage}\par
\par\medskip\Needspace{7\baselineskip}
\noindent\begin{minipage}{\linewidth}
\hypertarget{TERM-0006}{}
\noindent\textbf{reducirana norma / reducirana diferenta}\quad{\footnotesize\texttt{TERM-0006}}\par
\noindent\emph{Iztočnik / Source:} (red.) Norm / (red.) Differente\par
\noindent\emph{Rusky / Russian:} приведенная норма / приведенная дифферента\par
\noindent\emph{Ukrajinsky / Ukrainian:} зведена норма / зведена диферента\par
\smallskip\noindent\emph{Razlog:} Norma i diferenta imajut osobny algebraičny smysl; pridavnik reducirana go odliča od drugyh vrst normy ili diferenty. Arhivny izbor zachovuje internacionalny korenj reduciran- radi prepoznajemości. Vagu ili polno medžujezyčno sravnjenje zapis ne dava.\par
\noindent\emph{Reason:} Norm and different are technical algebraic terms; reduced distinguishes the intended constructions. The recorded proposal retains the international reduciran- root for recognizability. No numerical weighting or complete cross-language comparison is recorded.\par
\noindent{\footnotesize Izvorne zapisy / Original records: HIST-g208-014.}\par
\end{minipage}\par
\par\medskip\Needspace{7\baselineskip}
\noindent\begin{minipage}{\linewidth}
\hypertarget{TERM-0007}{}
\noindent\textbf{-1 kako suma dvoh kvadratov}\quad{\footnotesize\texttt{TERM-0007}}\par
\noindent\emph{Iztočnik / Source:} -1 als Summe von zwei Quadraten\par
\noindent\emph{Rusky / Russian:} -1 как сумма двух квадратов\par
\noindent\emph{Ukrajinsky / Ukrainian:} -1 як сума двох квадратів\par
\smallskip\noindent\emph{Razlog:} Vyraženje jest kriterij, ktory se pojavjaje v oběh napravjenjah dokaza. Zato -1, čislo dva i kvadraty morajut ostati jasno razpoznajeme. Razlog podpira semantičnu doslědnost; padežna forma dvoh trěbuje samostojnu gramatičnu prověrku.\par
\noindent\emph{Reason:} The phrase states the criterion used in both directions of the proof. The value -1, the number two, and the squares must remain identifiable. The reason supports semantic consistency; the inflected form dvoh requires a separate grammatical assessment.\par
\noindent{\footnotesize Izvorne zapisy / Original records: HIST-g176-004.}\par
\end{minipage}\par
\par\medskip\Needspace{7\baselineskip}
\noindent\begin{minipage}{\linewidth}
\hypertarget{TERM-0008}{}
\noindent\textbf{11. poręd / 12. poręd / 13. poręd / 15. poręd}\quad{\footnotesize\texttt{TERM-0008}}\par
\noindent\emph{Iztočnik / Source:} 11. Ordnung / 12. Ordnung / 13. Ordnung / 15. Ordnung\par
\noindent\emph{Rusky / Russian:} 11-й порядок / 12-й порядок / 13-й порядок / 15-й порядок\par
\noindent\emph{Ukrajinsky / Ukrainian:} 11-й порядок / 12-й порядок / 13-й порядок / 15-й порядок\par
\smallskip\noindent\emph{Razlog:} Kratke podzaglavja zachovujut numeričny poręd izloženja. Poręd byl prědložen radi kratkosti, ale sam izvorny zapis ostavjaje otvorjeno sravnjenje s poręda i s vyraženjem čisla slovom. Historijna kandidatna forma zato ne znači zaključen normativny izbor.\par
\noindent\emph{Reason:} The short subsection labels preserve the enumeration. Poręd was proposed for compactness, but the record expressly leaves comparison with poręda and a written-out ordinal unresolved. The historical proposal is not a settled normative recommendation.\par
\noindent{\footnotesize Izvorne zapisy / Original records: HIST-g017-002.}\par
\end{minipage}\par
\par\medskip\Needspace{7\baselineskip}
\noindent\begin{minipage}{\linewidth}
\hypertarget{TERM-0009}{}
\noindent\textbf{117 form}\quad{\footnotesize\texttt{TERM-0009}}\par
\noindent\emph{Iztočnik / Source:} 117 Formen\par
\noindent\emph{Rusky / Russian:} 117 форм\par
\noindent\emph{Ukrajinsky / Ukrainian:} 117 форм\par
\smallskip\noindent\emph{Razlog:} Čislo 117 jest čestju matematičnogo zaključka, ne redakcijna ocěna. Prěvod zachovuje točno čislo i predmet brojanja, formy. Zapis objasnjaje věrnost zaključku, ne novo tvorjenje terminologije.\par
\noindent\emph{Reason:} The number 117 belongs to the mathematical conclusion, not to an editorial estimate. The translation preserves the exact count and the objects counted, forms. The reason concerns fidelity rather than a newly coined term.\par
\noindent{\footnotesize Izvorne zapisy / Original records: HIST-g017-013.}\par
\end{minipage}\par
\par\medskip\Needspace{7\baselineskip}
\noindent\begin{minipage}{\linewidth}
\hypertarget{TERM-0010}{}
\noindent\textbf{30. marca 1926.}\quad{\footnotesize\texttt{TERM-0010}}\par
\noindent\emph{Iztočnik / Source:} 30. März 1926\par
\noindent\emph{Rusky / Russian:} 30 марта 1926 г.\par
\noindent\emph{Ukrajinsky / Ukrainian:} 30 березня 1926 р.\par
\smallskip\noindent\emph{Razlog:} Nazvanje měsęca jest prěvedeno, a denj i godina ostavajut neizměnjene. Medžuslovjanska forma upotrěbjaje poręd denj–měsęc–godina; rusky i ukrajinsky zapisy imajut svoje obyčne kratke oznaky godiny. Zapis neopravdava směnu samogo datuma.\par
\noindent\emph{Reason:} The month name is localized while the day and year remain unchanged. The Interslavic proposal uses day–month–year order; the Russian and Ukrainian comparanda use their ordinary year abbreviations. Nothing in the decision licenses a change to the date itself.\par
\noindent{\footnotesize Izvorne zapisy / Original records: HIST-g165-002.}\par
\end{minipage}\par
\par\medskip\Needspace{7\baselineskip}
\noindent\begin{minipage}{\linewidth}
\hypertarget{TERM-0011}{}
\noindent\textbf{5./6./... poręd}\quad{\footnotesize\texttt{TERM-0011}}\par
\noindent\emph{Iztočnik / Source:} 5./6./7./8./10./11. Ordnung\par
\noindent\emph{Rusky / Russian:} 5-й/6-й/... порядок\par
\noindent\emph{Ukrajinsky / Ukrainian:} 5-й/6-й/... порядок\par
\smallskip\noindent\emph{Razlog:} Terminom poręd se zachovuje nazvanje reda upotrěbljeno pri invariantah v drugom dělu. Izvorny zapis podpira doslědnost, ale izrično ostavjaje ortografiju za jezyčnu prověrku. Skračenje s tri točkami jest zapis skupiny zaglavij, ne zaměna za vsako polno zaglavje.\par
\noindent\emph{Reason:} Poręd continues the invariant-order vocabulary used in Paper 2. The reason is consistency, with orthography explicitly left for linguistic assessment. The ellipsis summarizes a group of headings; it is not a substitute for their complete text.\par
\noindent{\footnotesize Izvorne zapisy / Original records: HIST-g019-004.}\par
\end{minipage}\par
\par\medskip\Needspace{7\baselineskip}
\noindent\begin{minipage}{\linewidth}
\hypertarget{TERM-0012}{}
\noindent\textbf{\textbackslash{}D}\quad{\footnotesize\texttt{TERM-0012}}\par
\noindent\emph{Iztočnik / Source:} \textbackslash{}D\par
\noindent\emph{Rusky / Russian:} \textbackslash{}D\par
\noindent\emph{Ukrajinsky / Ukrainian:} \textbackslash{}D\par
\smallskip\noindent\emph{Razlog:} Komanda \textbackslash{}D označaje isti matematičny znak kako \textbackslash{}mathfrak D. Opreděljenje komandy omogožnjaje nabrati znak bez změny njegova značenja. Tuto jest pravilo matematičnogo zapisa, ne slovnično zaglavno slovo.\par
\noindent\emph{Reason:} The command \textbackslash{}D denotes the same mathematical sign as \textbackslash{}mathfrak D. Defining the command supplies the typesetting instruction without changing the sign's meaning. This is a notation convention, not a dictionary headword.\par
\noindent{\footnotesize Izvorne zapisy / Original records: HIST-g025-012.}\par
\end{minipage}\par
\par\medskip\Needspace{7\baselineskip}
\noindent\begin{minipage}{\linewidth}
\hypertarget{TERM-0013}{}
\noindent\textbf{simboly v \textbackslash{}pmod\{...\} ostavajut latinske}\quad{\footnotesize\texttt{TERM-0013}}\par
\noindent\emph{Iztočnik / Source:} \textbackslash{}modu\{(...)\} / \textbackslash{}pmod\{(...)\} symbolic modulus list\par
\noindent\emph{Rusky / Russian:} математические символы остаются латинскими\par
\noindent\emph{Ukrajinsky / Ukrainian:} математичні символи лишаються латинськими\par
\smallskip\noindent\emph{Razlog:} Simboly v uslovju kongruencije ostavajut matematične oznaky i ne sut prěvedene ili prěpisane kak slova. To uključaje latinske s, t, i i J vnųtri \textbackslash{}pmod\{...\}. Pravilo zachovuje identičnost proměnnyh v latiničnom i kiriličnom čitanju.\par
\noindent\emph{Reason:} Symbols in a congruence condition remain mathematical identifiers rather than words to translate or transliterate. This includes Latin s, t, i, and J inside \textbackslash{}pmod\{...\}. The rule preserves variable identity in both the Latin and Cyrillic readers.\par
\noindent{\footnotesize Izvorne zapisy / Original records: HIST-g013-015.}\par
\end{minipage}\par
\par\medskip\Needspace{7\baselineskip}
\noindent\begin{minipage}{\linewidth}
\hypertarget{TERM-0014}{}
\noindent\textbf{\textbackslash{}pmod\{(\textbackslash{}sigma,s)\}}\quad{\footnotesize\texttt{TERM-0014}}\par
\noindent\emph{Iztočnik / Source:} \textbackslash{}modu\{(\textbackslash{}sigma,s)\}\par
\noindent\emph{Rusky / Russian:} \textbackslash{}pmod\{(\textbackslash{}sigma,s)\}\par
\noindent\emph{Ukrajinsky / Ukrainian:} \textbackslash{}pmod\{(\textbackslash{}sigma,s)\}\par
\smallskip\noindent\emph{Razlog:} \textbackslash{}pmod\{(\textbackslash{}sigma,s)\} jest tipografsko vyraženje istogo uslovja kongruencije kako iztočna komanda \textbackslash{}modu\{(\textbackslash{}sigma,s)\}. σ i s ostavajut točne matematične oznaky. Sadržina skobok ne podleži prěpisivanju v drugy alfabet.\par
\noindent\emph{Reason:} The standard command \textbackslash{}pmod\{(\textbackslash{}sigma,s)\} expresses the same congruence condition as the source-specific \textbackslash{}modu\{(\textbackslash{}sigma,s)\}. The identifiers σ and s are preserved exactly. The parenthesized mathematical content is not transliterated.\par
\noindent{\footnotesize Izvorne zapisy / Original records: HIST-g017-015.}\par
\end{minipage}\par
\par\medskip\Needspace{7\baselineskip}
\noindent\begin{minipage}{\linewidth}
\hypertarget{TERM-0015}{}
\noindent\textbf{\textbackslash{}pmod\{(\textbackslash{}varrho,t)\}}\quad{\footnotesize\texttt{TERM-0015}}\par
\noindent\emph{Iztočnik / Source:} \textbackslash{}modu\{(\textbackslash{}varrho,t)\}\par
\noindent\emph{Rusky / Russian:} \textbackslash{}pmod\{(\textbackslash{}varrho,t)\}\par
\noindent\emph{Ukrajinsky / Ukrainian:} \textbackslash{}pmod\{(\textbackslash{}varrho,t)\}\par
\smallskip\noindent\emph{Razlog:} \textbackslash{}pmod\{(\textbackslash{}varrho,t)\} zachovuje matematičny uslov iztočnoj komandy \textbackslash{}modu\{(\textbackslash{}varrho,t)\}. Izbor se odnosi na nabirku, ne na proměnu kongruencije; ϱ i t ostavajut iste oznaky v oběh alfabeth.\par
\noindent\emph{Reason:} The standard command \textbackslash{}pmod\{(\textbackslash{}varrho,t)\} preserves the condition expressed by \textbackslash{}modu\{(\textbackslash{}varrho,t)\}. The choice concerns typesetting, not a change to the congruence. The identifiers ϱ and t remain the same in both scripts.\par
\noindent{\footnotesize Izvorne zapisy / Original records: HIST-g018-013.}\par
\end{minipage}\par
\par\medskip\Needspace{7\baselineskip}
\noindent\begin{minipage}{\linewidth}
\hypertarget{TERM-0016}{}
\noindent\textbf{a jest děljiva črez t / ideal a, kratny od t}\quad{\footnotesize\texttt{TERM-0016}}\par
\noindent\emph{Iztočnik / Source:} a ist teilbar durch t / Vielfaches a von t\par
\noindent\emph{Rusky / Russian:} a делится на t / идеал a, кратный t\par
\noindent\emph{Ukrajinsky / Ukrainian:} a ділиться на t / ідеал a, кратний t\par
\smallskip\noindent\emph{Razlog:} V glavnom tekstu prěvod izrično govori o děljivosti; v podčrtnoj primětce upotrěbjaje smysl kratnogo objekta. Taka razlika slijedi različna vyraženja iztočnika i ne zaměnjaje jih novym abstraktnym pojmom teorije idealov. Historijna slova črez i kratny od trěbajut odvojen jezyčny sud.\par
\noindent\emph{Reason:} The main text states divisibility, while the footnote uses the wording of a multiple. This follows the source's distinct expressions without introducing a new ideal-theoretic abstraction. The historical expressions črez and kratny od need separate linguistic assessment.\par
\noindent{\footnotesize Izvorne zapisy / Original records: HIST-g141-002.}\par
\end{minipage}\par
\par\medskip\Needspace{7\baselineskip}
\noindent\begin{minipage}{\linewidth}
\hypertarget{TERM-0017}{}
\noindent\textbf{a priori}\quad{\footnotesize\texttt{TERM-0017}}\par
\noindent\emph{Iztočnik / Source:} a priori\par
\noindent\emph{Rusky / Russian:} a priori\par
\noindent\emph{Ukrajinsky / Ukrainian:} a priori\par
\smallskip\noindent\emph{Razlog:} A priori jest zachovano kako internacionalno naučno latinsko vyraženje. To samo po sobě ne rěšaje, kako vsaky latinizm trěba pisati v kiriličnom tekstu. V originalnom razlogu nijedno opče pravilo za vse latinizmy ne bylo ustanovjeno.\par
\noindent\emph{Reason:} A priori is retained as an international scholarly Latin expression. That does not, by itself, determine how every Latinism should be written in Cyrillic text. The original reason establishes no general rule covering all Latinisms.\par
\noindent{\footnotesize Izvorne zapisy / Original records: HIST-g012-010.}\par
\end{minipage}\par
\par\medskip\Needspace{7\baselineskip}
\noindent\begin{minipage}{\linewidth}
\hypertarget{TERM-0018}{}
\noindent\textbf{A) / B) / C) / D)}\quad{\footnotesize\texttt{TERM-0018}}\par
\noindent\emph{Iztočnik / Source:} A) / B) / C) / D)\par
\noindent\emph{Rusky / Russian:} A) / B) / C) / D)\par
\noindent\emph{Ukrajinsky / Ukrainian:} A) / B) / C) / D)\par
\smallskip\noindent\emph{Razlog:} A), B), C) i D) sut oznaky česti argumenta, ne jezyčne jedinice. Njihova latinska identičnost ostavaje neizměnjena pri kiriličnom prěpisu; \textbackslash{}mbox\{...\} jest jedan način tipografskoj zaščity. Osobno C) ne směje stati Ц).\par
\noindent\emph{Reason:} A), B), C), and D) identify parts of the argument rather than words. Their Latin identities remain unchanged in Cyrillic text; \textbackslash{}mbox\{...\} is one means of protecting them in the source. In particular, C) must not become Ц).\par
\noindent{\footnotesize Izvorne zapisy / Original records: HIST-g023-012.}\par
\end{minipage}\par
\par\medskip\Needspace{7\baselineskip}
\noindent\begin{minipage}{\linewidth}
\hypertarget{TERM-0019}{}
\noindent\textbf{A) / B) / C) / D), a) / b)}\quad{\footnotesize\texttt{TERM-0019}}\par
\noindent\emph{Iztočnik / Source:} A) / B) / C) / D), a) / b)\par
\noindent\emph{Rusky / Russian:} A) / B) / C) / D), a) / b)\par
\noindent\emph{Ukrajinsky / Ukrainian:} A) / B) / C) / D), a) / b)\par
\smallskip\noindent\emph{Razlog:} Velike A), B), C), D) i male a), b) organizujut iztočny argument. Prěpis v kirilicu ne menja te oznaky. Zaščita oznak jest nezavisna od prěvoda obkružajučej prozy.\par
\noindent\emph{Reason:} The uppercase labels A), B), C), D) and lowercase a), b) organize the source argument. Conversion to Cyrillic does not change these identifiers. Their protection is independent of the translation of surrounding prose.\par
\noindent{\footnotesize Izvorne zapisy / Original records: HIST-g025-013.}\par
\end{minipage}\par
\par\medskip\Needspace{7\baselineskip}
\noindent\begin{minipage}{\linewidth}
\hypertarget{TERM-0020}{}
\noindent\textbf{A) / B) / C), a) / b) / c)}\quad{\footnotesize\texttt{TERM-0020}}\par
\noindent\emph{Iztočnik / Source:} A) / B) / C), a) / b) / c)\par
\noindent\emph{Rusky / Russian:} A) / B) / C), a) / b) / c)\par
\noindent\emph{Ukrajinsky / Ukrainian:} A) / B) / C), a) / b) / c)\par
\smallskip\noindent\emph{Razlog:} Oznaky A), B), C) i a), b), c) ostavajut latinske, ibo upučujut na iste česti dokaza v vseh jezyčnyh verzijah. Tuto pravilo se priměnjaje na strukturne oznaky; ne govori, že vsaka latinska bukva v običnoj prozě ima tų rolu.\par
\noindent\emph{Reason:} The labels A), B), C) and a), b), c) remain Latin because they refer to the same parts of the proof across language versions. The rule applies to structural labels, not indiscriminately to every Latin letter in ordinary prose.\par
\noindent{\footnotesize Izvorne zapisy / Original records: HIST-g024-014.}\par
\end{minipage}\par
\par\medskip\Needspace{7\baselineskip}
\noindent\begin{minipage}{\linewidth}
\hypertarget{TERM-0021}{}
\noindent\textbf{\textbackslash{}mathrm\{A)\}, \textbackslash{}mathrm\{B)\}, \textbackslash{}mathrm\{C)\}}\quad{\footnotesize\texttt{TERM-0021}}\par
\noindent\emph{Iztočnik / Source:} A), B), C) structural labels\par
\noindent\emph{Rusky / Russian:} A), B), C)\par
\noindent\emph{Ukrajinsky / Ukrainian:} A), B), C)\par
\smallskip\noindent\emph{Razlog:} \textbackslash{}mathrm\{A)\}, \textbackslash{}mathrm\{B)\} i \textbackslash{}mathrm\{C)\} nabirajut strukturne oznaky kako nepreměnjenu matematičnu latinicu. Cělj jest odlišiti oznaku od jezyčnogo teksta i zachovati upučivanja. To jest redakcijna konvencija, ne novy leksičny termin.\par
\noindent\emph{Reason:} The commands \textbackslash{}mathrm\{A)\}, \textbackslash{}mathrm\{B)\}, and \textbackslash{}mathrm\{C)\} typeset structural labels in upright mathematical Latin. The purpose is to distinguish identifiers from prose and preserve references. This is an editorial convention rather than a lexical innovation.\par
\noindent{\footnotesize Izvorne zapisy / Original records: HIST-g019-017.}\par
\end{minipage}\par
\par\medskip\Needspace{7\baselineskip}
\noindent\begin{minipage}{\linewidth}
\hypertarget{TERM-0022}{}
\noindent\textbf{odobraženje s neopreděljenymi koeficientami}\quad{\footnotesize\texttt{TERM-0022}}\par
\noindent\emph{Iztočnik / Source:} Abbildung mit unbestimmten Koeffizienten\par
\noindent\emph{Rusky / Russian:} отображение с неопределенными коэффициентами\par
\noindent\emph{Ukrajinsky / Ukrainian:} відображення з невизначеними коефіцієнтами\par
\smallskip\noindent\emph{Razlog:} Neopreděljene koeficienty sut važna čest uslovja: odobraženje može byti izbrano s takymi koeficientami. Prěvod ne redukuje izraz do opčego odobraženja i zachovuje ograničenje pri tvrdženju o jedinstvenosti.\par
\noindent\emph{Reason:} The indeterminate coefficients are an essential qualification: the map may be chosen with such coefficients. The translation does not reduce the expression to an unqualified map and preserves the caveat in the uniqueness discussion.\par
\noindent{\footnotesize Izvorne zapisy / Original records: HIST-g045-014.}\par
\end{minipage}\par
\par\medskip\Needspace{7\baselineskip}
\noindent\begin{minipage}{\linewidth}
\hypertarget{TERM-0023}{}
\noindent\textbf{oblast odobraženja}\quad{\footnotesize\texttt{TERM-0023}}\par
\noindent\emph{Iztočnik / Source:} Abbildungsbereich\par
\noindent\emph{Rusky / Russian:} область отображения\par
\noindent\emph{Ukrajinsky / Ukrainian:} область відображення\par
\smallskip\noindent\emph{Razlog:} Oblast odobraženja povezuje isti pojem v oddělah 3, 7 i 8, v opreděljenju VII i v finalnom kriteriju. Izvorne razloge podpira doslědna upotrěba v argumentah prěnosa. To ne razrěšaje avtomatično razliku medžu oblastju opreděljenja i obrazom; tę razliku trěba čitati v každom kontekstu.\par
\noindent\emph{Reason:} The phrase maintains continuity across Sections 3, 7, and 8, Definition VII, and the final criterion. The recorded reasons concern consistent use in transfer arguments. They do not automatically settle the distinction between domain and image; that must be read in context.\par
\noindent{\footnotesize Izvorne zapisy / Original records: HIST-g048-011, HIST-g049-016, HIST-g050-008.}\par
\end{minipage}\par
\par\medskip\Needspace{7\baselineskip}
\noindent\begin{minipage}{\linewidth}
\hypertarget{TERM-0024}{}
\noindent\textbf{funkcije odobraženja}\quad{\footnotesize\texttt{TERM-0024}}\par
\noindent\emph{Iztočnik / Source:} Abbildungsfunktionen\par
\noindent\emph{Rusky / Russian:} функции отображения\par
\noindent\emph{Ukrajinsky / Ukrainian:} функції відображення\par
\smallskip\noindent\emph{Razlog:} Podčrtna primětka opreděljaje tę funkcije kako funkcije proměnnyh x₁, …, xρ. Prěmy termin funkcije odobraženja jest zato dopolnjen iztočnym objasnjenjem, vměsto toga da bi sam termin nesl neizrečene nove hipotezy.\par
\noindent\emph{Reason:} The footnote defines these as functions of x₁, …, xρ. The direct term is therefore accompanied by the source's explanation rather than being made to carry additional unstated assumptions.\par
\noindent{\footnotesize Izvorne zapisy / Original records: HIST-g043-005.}\par
\end{minipage}\par
\par\medskip\Needspace{7\baselineskip}
\noindent\begin{minipage}{\linewidth}
\hypertarget{TERM-0025}{}
\noindent\textbf{polje obrazov / oblast obrazov}\quad{\footnotesize\texttt{TERM-0025}}\par
\noindent\emph{Iztočnik / Source:} Abbildungskörper / Abbildungsbereich\par
\noindent\emph{Rusky / Russian:} поле образов / область образов\par
\noindent\emph{Ukrajinsky / Ukrainian:} поле образів / область образів\par
\smallskip\noindent\emph{Razlog:} Polje obrazov i oblast obrazov upučujut na rezultaty odobraženja, ne na sam proces. Izbor prodolžuje smysl izraza sistem obrazov i izběgaje dvojsmysl, ktory by nastal pri doslovnom prevodu vsakego Abbildung jednym slovom.\par
\noindent\emph{Reason:} The expressions name the image field or range, not the act of mapping. They continue the interpretation of sistem obrazov and avoid the ambiguity that would result from translating every occurrence of Abbildung with one undifferentiated term.\par
\noindent{\footnotesize Izvorne zapisy / Original records: HIST-g070-008.}\par
\end{minipage}\par
\par\medskip\Needspace{7\baselineskip}
\noindent\begin{minipage}{\linewidth}
\hypertarget{TERM-0026}{}
\noindent\textbf{odobražajuči sistem}\quad{\footnotesize\texttt{TERM-0026}}\par
\noindent\emph{Iztočnik / Source:} Abbildungssystem\par
\noindent\emph{Rusky / Russian:} система отображения\par
\noindent\emph{Ukrajinsky / Ukrainian:} система відображення\par
\smallskip\noindent\emph{Razlog:} Odobražajuči sistem byl kandidat za sistem S*, ktory služi v argumentu prěnosa i v dokazu o racionalnoj osnově. Izvorny zapis izrično ostavjaje aktivnu pridavničnu tvorbu za jezyčnu prověrku. Ne trěba go směšati s konačno potvrđenym nazvanjem sistemy obrazov.\par
\noindent\emph{Reason:} The proposal names the system S* used in transfer arguments and the rational-basis proof. The original record expressly leaves its active adjectival formation for linguistic review. It should not be confused with an established recommendation for the image system.\par
\noindent{\footnotesize Izvorne zapisy / Original records: HIST-g043-010, HIST-g044-005.}\par
\end{minipage}\par
\par\medskip\Needspace{7\baselineskip}
\noindent\begin{minipage}{\linewidth}
\hypertarget{TERM-0027}{}
\noindent\textbf{sistem obrazov}\quad{\footnotesize\texttt{TERM-0027}}\par
\noindent\emph{Iztočnik / Source:} Abbildungssystem\par
\noindent\emph{Rusky / Russian:} система образов\par
\noindent\emph{Ukrajinsky / Ukrainian:} система образів\par
\smallskip\noindent\emph{Razlog:} Sistem obrazov jest izbrano vměsto doslovnogo nazvanja procesa odobraženja, ibo věta govori o sistemu vseh cěnnostij f(z). Presudny jest matematičny referent: kolekcija obrazov, ne dejanje odobražanja.\par
\noindent\emph{Reason:} Sistem obrazov was preferred to a literal mapping-system expression because the sentence concerns the system of all values f(z). The decisive point is the referent: the collection of images, not the process of mapping.\par
\noindent{\footnotesize Izvorne zapisy / Original records: HIST-g069-006.}\par
\end{minipage}\par
\par\medskip\Needspace{7\baselineskip}
\noindent\begin{minipage}{\linewidth}
\hypertarget{TERM-0028}{}
\noindent\textbf{Abelove grupy/polja}\quad{\footnotesize\texttt{TERM-0028}}\par
\noindent\emph{Iztočnik / Source:} Abelsche Gruppen/Körper\par
\noindent\emph{Rusky / Russian:} абелевы группы/поля\par
\noindent\emph{Ukrajinsky / Ukrainian:} абелеві групи/поля\par
\smallskip\noindent\emph{Razlog:} Abelove zachovuje běžno algebraično nazvanje, sravnjeno s ruskoju i ukrajinskoju pridavničnomu formoju. Izbor se takože opira na doslědnost s ranějšimi pojavjenjami tutogo pridavnika. Sam zapis ne daje opče pravilo za vse imenske pridavniki.\par
\noindent\emph{Reason:} Abelove follows conventional algebraic adjective formation, with Russian and Ukrainian comparanda and consistency with earlier occurrences. The record does not supply a general rule for every adjective derived from a personal name.\par
\noindent{\footnotesize Izvorne zapisy / Original records: HIST-g073-021.}\par
\end{minipage}\par
\par\medskip\Needspace{7\baselineskip}
\noindent\begin{minipage}{\linewidth}
\hypertarget{TERM-0029}{}
\noindent\textbf{izvedena funkcionalna oblast / polje kvocientov}\quad{\footnotesize\texttt{TERM-0029}}\par
\noindent\emph{Iztočnik / Source:} abgeleiteter Funktionalbereich / Quotientenkörper\par
\noindent\emph{Rusky / Russian:} выведенная функциональная область / поле частных\par
\noindent\emph{Ukrajinsky / Ukrainian:} виведена функціональна область / поле часток\par
\smallskip\noindent\emph{Razlog:} Izvedena funkcionalna oblast i polje kvocientov prodolžujut terminologiju odděla 7. Izvedena izjasnjaje odnos konstrukcije, a polje kvocientov imenuje algebraičny objekt. Razlog jest transparentnost i doslědnost tųh izrazov v tom samom argumentu.\par
\noindent\emph{Reason:} The two expressions continue the terminology of Section 7. Izvedena conveys the derivational relation; polje kvocientov names the algebraic object. The recorded reasons are transparency and consistency within the argument.\par
\noindent{\footnotesize Izvorne zapisy / Original records: HIST-g153-005.}\par
\end{minipage}\par
\par\medskip\Needspace{7\baselineskip}
\noindent\begin{minipage}{\linewidth}
\hypertarget{TERM-0030}{}
\noindent\textbf{kolco izvedeno iz}\quad{\footnotesize\texttt{TERM-0030}}\par
\noindent\emph{Iztočnik / Source:} abgeleiteter Ring\par
\noindent\emph{Rusky / Russian:} кольцо, выведенное из\par
\noindent\emph{Ukrajinsky / Ukrainian:} кільце, виведене з\par
\smallskip\noindent\emph{Razlog:} Izvedeno iz razkriva odnos izraženy německym participom vměsto težkoj imennoj konstrukcije. To jest sintaktičny razlog. Historijna forma kolco ostavaje kandidatnym zapisom; ten razlog ne dokazuje, že ona jest aktualny preferovany termin za ring.\par
\noindent\emph{Reason:} Izvedeno iz makes the relation expressed by the German participle explicit instead of using a dense nominal construction. This is a syntactic reason. The historical kolco remains a recorded candidate; the reason does not establish it as the current preferred term for ring.\par
\noindent{\footnotesize Izvorne zapisy / Original records: HIST-g123-004.}\par
\smallskip\noindent\emph{Sravnjenje s tekstom:} V dělu 31, oddělu 2, izraz prsten izvedeno iz ima izbrany imennik prsten, ale pridavny oblik izvedeno ostavaje srědnjego roda. Mužsky rod prsten trěbuje izvedeny. Zato točna prisutnost termina ne dokazuje pravilnost cěloj konstrukcije. Sravniteljna pasaža s instrumentalom prstenom izvedenym iz pokazuje soglasovanje v drugom padežu. Tuto jest gramatična analiza, ne tvrdženje, že tekst korpusa ima ispravjenu formu.\par
\noindent\emph{Textual comparison:} In Paper 31, section 2, prsten izvedeno iz contains the selected noun prsten, but izvedeno retains a neuter ending. Masculine prsten calls for izvedeny. Exact occurrence of the term therefore does not establish the correctness of the whole construction. A comparison passage with instrumental prstenom izvedenym iz exhibits agreement in another case. This is a grammatical analysis, not a claim that the corpus text contains the corrected form.\par
\noindent{\footnotesize TEXT-COMP-01; pasaže / passages: 17561, 17314; \texttt{CORPUS\_FORM\_COMPARISONS.json}.}\par
\end{minipage}\par
\par\medskip\Needspace{7\baselineskip}
\noindent\begin{minipage}{\linewidth}
\hypertarget{TERM-0031}{}
\noindent\textbf{glavny ideal, izvedeny iz p}\quad{\footnotesize\texttt{TERM-0031}}\par
\noindent\emph{Iztočnik / Source:} abgeleitetes Hauptideal\par
\noindent\emph{Rusky / Russian:} главный идеал, выведенный из p\par
\noindent\emph{Ukrajinsky / Ukrainian:} головний ідеал, виведений із p\par
\smallskip\noindent\emph{Razlog:} Izvedeny iz p zachovuje iztočno vyraženje o izvedenom glavnom idealu i upotrěbjaje produktivnu familiju izvedeny. Vyraženje ne zaměnjaje odnosa k p neizrečenym sučasnym opreděljenjem.\par
\noindent\emph{Reason:} The expression preserves the source's derived-principal-ideal wording and uses the productive izvedeny family. The relation to p is retained rather than replaced by an unstated modern definition.\par
\noindent{\footnotesize Izvorne zapisy / Original records: HIST-g151-003.}\par
\end{minipage}\par
\par\medskip\Needspace{7\baselineskip}
\noindent\begin{minipage}{\linewidth}
\hypertarget{TERM-0032}{}
\noindent\textbf{izvedeny glavny ideal}\quad{\footnotesize\texttt{TERM-0032}}\par
\noindent\emph{Iztočnik / Source:} abgeleitetes Hauptideal\par
\noindent\emph{Rusky / Russian:} выведенный главный идеал\par
\noindent\emph{Ukrajinsky / Ukrainian:} виведений головний ідеал\par
\smallskip\noindent\emph{Razlog:} Izvedeny glavny ideal zachovuje imenovanje iztočnika. Zapis izrično odklanjaje tihu zaměnu opisom ideala generovanogo diskriminantom: tak opis by byl interpretacija, ne sam prěvod iztočnoj konstrukcije.\par
\noindent\emph{Reason:} The expression retains the source's designation. The record explicitly rejects silently replacing it with an ideal-generated-by-the-discriminant paraphrase: that would be an interpretation rather than a translation of the construction itself.\par
\noindent{\footnotesize Izvorne zapisy / Original records: HIST-g137-002.}\par
\end{minipage}\par
\par\medskip\Needspace{7\baselineskip}
\noindent\begin{minipage}{\linewidth}
\hypertarget{TERM-0033}{}
\noindent\textbf{zatvorjena vzhodno k tvorenju pr?s?kov}\quad{\footnotesize\texttt{TERM-0033}}\par
\noindent\emph{Iztočnik / Source:} abgeschlossen in bezug auf Durchschnittbildung\par
\noindent\emph{Rusky / Russian:} [Bukvy ne sut zachovane / Lettering not recoverable]\par
\noindent\emph{Ukrajinsky / Ukrainian:} [Bukvy ne sut zachovane / Lettering not recoverable]\par
\smallskip\noindent\emph{Razlog:} Razlog govori o zatvorjenosti pri tvorjenju prěsekov i o kontrastu s nezatvorjenostju v oddělu 5. Ale zachovana forma ima znaky ? vnųtri slova i ne može služiti kako pouzdan leksičny obrazec. Matematičny razlog i poškođeny zapis morajut byti razdvojene.\par
\noindent\emph{Reason:} The rationale concerns closure under intersections and its contrast with the non-closure statement in Section 5. The preserved candidate contains question marks inside a word and cannot serve as a reliable lexical form. Its mathematical rationale and its damaged spelling must be distinguished.\par
\noindent{\footnotesize Izvorne zapisy / Original records: HIST-g064-009.}\par
\end{minipage}\par
\par\medskip\Needspace{7\baselineskip}
\noindent\begin{minipage}{\linewidth}
\hypertarget{TERM-0034}{}
\noindent\textbf{zavisnosti medžu Lagranževymi izrazami}\quad{\footnotesize\texttt{TERM-0034}}\par
\noindent\emph{Iztočnik / Source:} Abhängigkeiten zwischen den Lagrangeschen Ausdrücken\par
\noindent\emph{Rusky / Russian:} зависимости между лагранжевыми выражениями\par
\noindent\emph{Ukrajinsky / Ukrainian:} залежності між лагранжевими виразами\par
\smallskip\noindent\emph{Razlog:} Zavisnosti sut identičnosti medžu izrazami Eulera–Lagrangea i jih derivacijami. Izbor zachovuje razliku od relacij divergencije, s ktorymi ih Noether sravnjaje. Opče slovo relacije by tu smyslovu razliku oslabilo.\par
\noindent\emph{Reason:} The dependencies are identities among the Euler–Lagrange expressions and their derivatives. The wording preserves Noether's contrast with divergence relations. Replacing both by the generic relations would obscure that distinction.\par
\noindent{\footnotesize Izvorne zapisy / Original records: HIST-g075-010.}\par
\end{minipage}\par
\par\medskip\Needspace{7\baselineskip}
\noindent\begin{minipage}{\linewidth}
\hypertarget{TERM-0035}{}
\noindent\textbf{izvodženje formul}\quad{\footnotesize\texttt{TERM-0035}}\par
\noindent\emph{Iztočnik / Source:} Ableitung der Formeln\par
\noindent\emph{Rusky / Russian:} вывод формул\par
\noindent\emph{Ukrajinsky / Ukrainian:} виведення формул\par
\smallskip\noindent\emph{Razlog:} Izvodženje jest kandidat za proces vyvodženja formul, sravnjany s obyčnym ruskim i ukrajinskim matematičnym vyraženjem. Zapis navęzuje na izvod v drugom dělu, ale izrično označaje morfologiju i prirodnost konstrukcije kako otvorjene pytanja.\par
\noindent\emph{Reason:} The candidate denotes the derivation of formulas, compared with ordinary Russian and Ukrainian mathematical phrasing and the earlier izvod family in Paper 2. The record explicitly leaves its morphology and idiomatic naturalness open to assessment.\par
\noindent{\footnotesize Izvorne zapisy / Original records: HIST-g011-011.}\par
\end{minipage}\par
\par\medskip\Needspace{7\baselineskip}
\noindent\begin{minipage}{\linewidth}
\hypertarget{TERM-0036}{}
\noindent\textbf{polje početnogo odrězka K(vartheta)}\quad{\footnotesize\texttt{TERM-0036}}\par
\noindent\emph{Iztočnik / Source:} Abschnittskörper K(vartheta)\par
\noindent\emph{Rusky / Russian:} поле начального отрезка K(vartheta)\par
\noindent\emph{Ukrajinsky / Ukrainian:} поле початкового відрізка K(vartheta)\par
\smallskip\noindent\emph{Razlog:} Opis polje početnogo odrězka jest izbrany vměsto kratkoj doslovnoj skovanice. On izjasnjaje konstrukciju: polje generovano početnym odrězkom prěd ϑ v dobrom uporędčenju. Dodatna dolžina služi točnosti referenta.\par
\noindent\emph{Reason:} The descriptive term is preferred to a compact literal segment-field calque. It identifies the construction: the field generated by the initial segment preceding ϑ in the well-order. The additional length serves referential precision.\par
\noindent{\footnotesize Izvorne zapisy / Original records: HIST-g070-005.}\par
\end{minipage}\par
\par\medskip\Needspace{7\baselineskip}
\noindent\begin{minipage}{\linewidth}
\hypertarget{TERM-0037}{}
\noindent\textbf{absolutno algebraično polje}\quad{\footnotesize\texttt{TERM-0037}}\par
\noindent\emph{Iztočnik / Source:} absolut algebraischer Körper\par
\noindent\emph{Rusky / Russian:} абсолютно алгебраическое поле\par
\noindent\emph{Ukrajinsky / Ukrainian:} абсолютно алгебраїчне поле\par
\smallskip\noindent\emph{Razlog:} Absolutno zachovuje historijno imenovanje algebraičnogo polja nad prostym poljem. Izvorny razlog upučuje na pozdnějše opreděljenja algebraično cělyh i drobnyh veličin. Objem pojma trěba čitati v tųh opreděljenjah, ne izvedati samo iz sučasnogo značenja pridavnika.\par
\noindent\emph{Reason:} Absolutno preserves the historical designation of the algebraic field over the prime field, to which later definitions of algebraically integral and fractional quantities refer. Its scope must be read in those definitions, not inferred solely from a modern adjective.\par
\noindent{\footnotesize Izvorne zapisy / Original records: HIST-g068-006.}\par
\end{minipage}\par
\par\medskip\Needspace{7\baselineskip}
\noindent\begin{minipage}{\linewidth}
\hypertarget{TERM-0038}{}
\noindent\textbf{absolutno nerazložime sostavne česti}\quad{\footnotesize\texttt{TERM-0038}}\par
\noindent\emph{Iztočnik / Source:} absolut irreduzible Bestandteile\par
\noindent\emph{Rusky / Russian:} абсолютно неприводимые составляющие\par
\noindent\emph{Ukrajinsky / Ukrainian:} абсолютно незвідні складові\par
\smallskip\noindent\emph{Razlog:} Sostavne česti razlikuje sastojki od predstavjenij i klasov. Pridavnik vyražaje absolutnu nerazložimost. Zapis se opira na rusko i ukrajinsko odborno vyraženje i na doslědnost s dělom 17, ale označaje slabu podkrepu medžuslovjanskogo registra teorije predstavjenij.\par
\noindent\emph{Reason:} Sostavne česti distinguishes constituents from representations and classes; the adjective expresses absolute irreducibility. The record invokes Russian and Ukrainian mathematical phrasing and consistency with Paper 17, while acknowledging limited evidence for an Interslavic representation-theory register.\par
\noindent{\footnotesize Izvorne zapisy / Original records: HIST-g167-002.}\par
\end{minipage}\par
\par\medskip\Needspace{7\baselineskip}
\noindent\begin{minipage}{\linewidth}
\hypertarget{TERM-0039}{}
\noindent\textbf{absolutno polny sistem}\quad{\footnotesize\texttt{TERM-0039}}\par
\noindent\emph{Iztočnik / Source:} absolut vollständiges System\par
\noindent\emph{Rusky / Russian:} абсолютно полная система\par
\noindent\emph{Ukrajinsky / Ukrainian:} абсолютно повна система\par
\smallskip\noindent\emph{Razlog:} V zachovanom terminologičnom zapisu ne jest osobno objasnjeno, zašto byla izbrana tuta forma. Par absolut vollständiges System — absolutno polny sistem jest dokumentovany, ale njegove istorijne vage, alternativy i razlog ne možut byti izmysljene. Vsako novo opravdanje trěba jasno odlišiti od izvornogo dokaza.\par
\noindent\emph{Reason:} No separate rationale survives in this terminology record. The source–candidate pair is documented, but its historical weights, alternatives, and justification cannot be invented. Any new justification must be explicitly distinguished from the original evidence.\par
\noindent{\footnotesize Izvorne zapisy / Original records: HIST-g001-005.}\par
\end{minipage}\par
\par\medskip\Needspace{7\baselineskip}
\noindent\begin{minipage}{\linewidth}
\hypertarget{TERM-0040}{}
\noindent\textbf{absolutno ireducibilny}\quad{\footnotesize\texttt{TERM-0040}}\par
\noindent\emph{Iztočnik / Source:} absolut-irreduzibel\par
\noindent\emph{Rusky / Russian:} абсолютно неприводимый\par
\noindent\emph{Ukrajinsky / Ukrainian:} абсолютно незвідний\par
\smallskip\noindent\emph{Razlog:} Internacionalna forma ireducibilny jest tu objasnjena doslědnostju s terminologijeju reducibilnosti v oddělah 16–20 děla 34. Zapis ne tvrdi, že internacionalny korenj mora iměti prědnost v vseh kontekstah; druge zapisy upotrěbjajut familiju nerazložim-.\par
\noindent\emph{Reason:} The international ireducibilny is justified here by consistency with the reducibility vocabulary of Paper 34, Sections 16–20. The record does not establish a universal preference for that root; other records use the nerazložim- family.\par
\noindent{\footnotesize Izvorne zapisy / Original records: HIST-g199-002.}\par
\end{minipage}\par
\par\medskip\Needspace{7\baselineskip}
\noindent\begin{minipage}{\linewidth}
\hypertarget{TERM-0041}{}
\noindent\textbf{absolutno ireducibilno predstavjenje}\quad{\footnotesize\texttt{TERM-0041}}\par
\noindent\emph{Iztočnik / Source:} absolut-irreduzible Darstellung\par
\noindent\emph{Rusky / Russian:} абсолютно неприводимое представление\par
\noindent\emph{Ukrajinsky / Ukrainian:} абсолютно незвідне представлення\par
\smallskip\noindent\emph{Razlog:} Predstavjenje i ireducibilno prodolžujut tehničnu terminologiju děla 34 i oddělov 21–23. Rusko i ukrajinsko sravnjenje podpira matematičny smysl nerazložimosti. Zapis dokumentuje doslědnost jednoj kandidatnoj familije, ne pobědu nad vsemi alternativami.\par
\noindent\emph{Reason:} Predstavjenje and ireducibilno continue the terminology of Paper 34 and Sections 21–23. Russian and Ukrainian comparanda support the mathematical sense of irreducibility. The record documents consistency within a candidate family, not a demonstrated victory over every alternative.\par
\noindent{\footnotesize Izvorne zapisy / Original records: HIST-g202-005, HIST-g204-009.}\par
\end{minipage}\par
\par\medskip\Needspace{7\baselineskip}
\noindent\begin{minipage}{\linewidth}
\hypertarget{TERM-0042}{}
\noindent\textbf{absolutna nerazložimost / absolutno nerazložimy}\quad{\footnotesize\texttt{TERM-0042}}\par
\noindent\emph{Iztočnik / Source:} absolute Irreduzibilität / absolut irreduzibel\par
\noindent\emph{Rusky / Russian:} абсолютная неприводимость / абсолютно неприводимый\par
\noindent\emph{Ukrajinsky / Ukrainian:} абсолютна незвідність / абсолютно незвідний\par
\smallskip\noindent\emph{Razlog:} Familija razložiti/nerazložimy jest izbrana vměsto čisto internacionalnogo pridavnika, ibo argument povtorno sravnjaje nerazloženost s faktičnoju faktorizacijeju. Presudny razlog jest prozirnost toj smyslovoj vezy. Rusky i ukrajinsky izrazy služet sravnjenju tehničnogo smysla.\par
\noindent\emph{Reason:} The razložiti/nerazložimy family is preferred to a purely international adjective because the argument repeatedly contrasts remaining undecomposed with actual factorization. The decisive reason is that semantic connection. Russian and Ukrainian terms provide comparanda for the technical sense.\par
\noindent{\footnotesize Izvorne zapisy / Original records: HIST-g083-001.}\par
\end{minipage}\par
\par\medskip\Needspace{7\baselineskip}
\noindent\begin{minipage}{\linewidth}
\hypertarget{TERM-0043}{}
\noindent\textbf{absolutna prosta funkcija / absolutno nerazložimy polinom}\quad{\footnotesize\texttt{TERM-0043}}\par
\noindent\emph{Iztočnik / Source:} absolute Primfunktion / absolut irreduzibler Polynom\par
\noindent\emph{Rusky / Russian:} абсолютная простая функция / абсолютно неприводимый полином\par
\noindent\emph{Ukrajinsky / Ukrainian:} абсолютна проста функція / абсолютно незвідний поліном\par
\smallskip\noindent\emph{Razlog:} Prosta funkcija zachovuje historijno slovo funkcija; polinom v objasnjenju ukazuje sučasny algebraičny smysl. Taka para ne dozvoljaje avtomatično zaměniti vsako iztočno funkcija slovom polinom. Ona razděljuje avtorsko imenovanje od objasnjenja.\par
\noindent\emph{Reason:} Prosta funkcija retains the historical function-language, while the polynomial expression clarifies the modern algebraic reading. The pairing does not license replacing every source occurrence of function by polynomial. It distinguishes the author's designation from an explanatory gloss.\par
\noindent{\footnotesize Izvorne zapisy / Original records: HIST-g101-002.}\par
\end{minipage}\par
\par\medskip\Needspace{7\baselineskip}
\noindent\begin{minipage}{\linewidth}
\hypertarget{TERM-0044}{}
\noindent\textbf{absolutna oblast množiteljev}\quad{\footnotesize\texttt{TERM-0044}}\par
\noindent\emph{Iztočnik / Source:} absoluter Multiplikatorenbereich\par
\noindent\emph{Rusky / Russian:} абсолютная область множителей\par
\noindent\emph{Ukrajinsky / Ukrainian:} абсолютна область множників\par
\smallskip\noindent\emph{Razlog:} Oblast množiteljev zachovuje nazvanje upotrěbljeno v oddělah 1 i 8. Oblast ne jest skryto zaměnjena slovom algebra, ibo iztočny pojem jest vezany s avtorskim opisom operatornoj oblasti. Absolutna ostavaje čestju tehničnogo izraza.\par
\noindent\emph{Reason:} The term continues the designation used in Sections 1 and 8. Oblast is not silently replaced by algebra, because the source concept is tied to Noether's operator-domain formulation. Absolutna remains part of the technical expression.\par
\noindent{\footnotesize Izvorne zapisy / Original records: HIST-g196-007, HIST-g198-006.}\par
\end{minipage}\par
\par\medskip\Needspace{7\baselineskip}
\noindent\begin{minipage}{\linewidth}
\hypertarget{TERM-0045}{}
\noindent\textbf{absolutny prost ideal}\quad{\footnotesize\texttt{TERM-0045}}\par
\noindent\emph{Iztočnik / Source:} absolutes Primideal\par
\noindent\emph{Rusky / Russian:} абсолютный простой идеал\par
\noindent\emph{Ukrajinsky / Ukrainian:} абсолютний простий ідеал\par
\smallskip\noindent\emph{Razlog:} Prost ideal zachovuje razliku medžu prostym i primarnym idealom; absolutny jest internacionalny pridavnik. Izvorny razlog govori o toj algebraičnoj razlicě i o doslědnosti. On ne odlučaje sam po sobě medžu kratkoju formoju prost i formoju prosty.\par
\noindent\emph{Reason:} The proposal preserves the algebraic distinction between prime and primary ideals and combines it with the international adjective absolutny. The reason concerns that distinction and consistency; it does not independently decide between prost and prosty.\par
\noindent{\footnotesize Izvorne zapisy / Original records: HIST-g101-001.}\par
\end{minipage}\par
\par\medskip\Needspace{7\baselineskip}
\noindent\begin{minipage}{\linewidth}
\hypertarget{TERM-0046}{}
\noindent\textbf{absolutny prosty ideal}\quad{\footnotesize\texttt{TERM-0046}}\par
\noindent\emph{Iztočnik / Source:} absolutes Primideal\par
\noindent\emph{Rusky / Russian:} абсолютный простой идеал\par
\noindent\emph{Ukrajinsky / Ukrainian:} абсолютний простий ідеал\par
\smallskip\noindent\emph{Razlog:} Izraz neposrědno veže absolutny prosty ideal s diskusijeju o absolutnoj nerazložimosti. Tuto jest zabilěženy smyslovy razlog. Zapis ne dava osobnu morfologičnu argumentaciju za formu prosty nasuprot prost.\par
\noindent\emph{Reason:} The expression maintains the link between an absolute prime ideal and the discussion of absolute irreducibility. This is the recorded semantic reason. No separate morphological argument for prosty over prost is supplied.\par
\noindent{\footnotesize Izvorne zapisy / Original records: HIST-g205-008.}\par
\end{minipage}\par
\par\medskip\Needspace{7\baselineskip}
\noindent\begin{minipage}{\linewidth}
\hypertarget{TERM-0047}{}
\noindent\textbf{abstraktna teorija idealov}\quad{\footnotesize\texttt{TERM-0047}}\par
\noindent\emph{Iztočnik / Source:} abstrakte Idealtheorie\par
\noindent\emph{Rusky / Russian:} абстрактная теория идеалов\par
\noindent\emph{Ukrajinsky / Ukrainian:} абстрактна теорія ідеалів\par
\smallskip\noindent\emph{Razlog:} Prěma internacionalna konstrukcija označaje abstraktnu teoriju idealov, v ktorej Noether nanovo formuluje rezultat Hentzelta o eliminaciji. Razlog izbora jest jasno nazvanje teorijnoj perspektivy, ne zaměna samogo rezultata.\par
\noindent\emph{Reason:} The direct international expression names the abstract ideal-theoretic perspective in which Noether formulates Hentzelt's elimination result. The choice identifies that perspective clearly without replacing the result itself.\par
\noindent{\footnotesize Izvorne zapisy / Original records: HIST-g080-010.}\par
\end{minipage}\par
\par\medskip\Needspace{7\baselineskip}
\noindent\begin{minipage}{\linewidth}
\hypertarget{TERM-0048}{}
\noindent\textbf{abstraktna izgradnja teorije idealov}\quad{\footnotesize\texttt{TERM-0048}}\par
\noindent\emph{Iztočnik / Source:} Abstrakter Aufbau der Idealtheorie\par
\noindent\emph{Rusky / Russian:} абстрактное построение теории идеалов\par
\noindent\emph{Ukrajinsky / Ukrainian:} абстрактна побудова теорії ідеалів\par
\smallskip\noindent\emph{Razlog:} Izgradnja zachovuje iztočnu metaforu tvorenja teorije. Zapis ne zaměnjaje Aufbau pozdnějšim nazvanjem aksiomatizacije, ktoro by pribavilo interpretaciju k zaglavju.\par
\noindent\emph{Reason:} Izgradnja preserves the source's construction metaphor. The record avoids substituting the later label axiomatization, which would add an interpretation to the title.\par
\noindent{\footnotesize Izvorne zapisy / Original records: HIST-g103-001.}\par
\end{minipage}\par
\par\medskip\Needspace{7\baselineskip}
\noindent\begin{minipage}{\linewidth}
\hypertarget{TERM-0049}{}
\noindent\textbf{abstraktna konstrukcija teorije idealov}\quad{\footnotesize\texttt{TERM-0049}}\par
\noindent\emph{Iztočnik / Source:} Abstrakter Aufbau der Idealtheorie\par
\noindent\emph{Rusky / Russian:} абстрактное построение теории идеалов\par
\noindent\emph{Ukrajinsky / Ukrainian:} абстрактна побудова теорії ідеалів\par
\smallskip\noindent\emph{Razlog:} Konstrukcija takože tumači Aufbau kako organizaciju teorije, ale dava prědnost poznatomu internacionalnomu imeniku prěd doslovnoju arhitektoničnoju skovanicoju. To jest druga zabilěžena możliwość uz izgradnja, ne dokaz, že oba zapisy sut jedna jedina odluka.\par
\noindent\emph{Reason:} Konstrukcija likewise interprets Aufbau as the organization of a theory, but prefers a familiar international noun to a literal architectural calque. It is a separately recorded alternative to izgradnja, not evidence that the two records constitute one settled decision.\par
\noindent{\footnotesize Izvorne zapisy / Original records: HIST-g107-001.}\par
\end{minipage}\par
\par\medskip\Needspace{7\baselineskip}
\noindent\begin{minipage}{\linewidth}
\hypertarget{TERM-0050}{}
\noindent\textbf{iz?etna mno?ina / oblasti G\_nu}\quad{\footnotesize\texttt{TERM-0050}}\par
\noindent\emph{Iztočnik / Source:} abz?hlbare Menge / abz?hlbare Bereiche G\_nu\par
\noindent\emph{Rusky / Russian:} ??????? ????????? / ??????? G\_nu\par
\noindent\emph{Ukrajinsky / Ukrainian:} ???????? ??????? / ??????? G\_nu\par
\smallskip\noindent\emph{Razlog:} Zapis govori o izčetnoj familiji oblastij Gν i navęzuje na ranějše vyraženja o velikosti množin. Ale i iztočno i kandidatno slovo imajut izgubljene znaky, zapisane kako ?. Semantično objasnjenje se može čitati; oštećena grafija ne jest preporučena forma.\par
\noindent\emph{Reason:} The reason concerns the countable family of domains Gν and consistency with earlier set-size terminology. Both the source and candidate contain lost characters represented by question marks. The semantic explanation is recoverable; the damaged spelling is not a recommended form.\par
\noindent{\footnotesize Izvorne zapisy / Original records: HIST-g063-005.}\par
\end{minipage}\par
\par\medskip\Needspace{7\baselineskip}
\noindent\begin{minipage}{\linewidth}
\hypertarget{TERM-0051}{}
\noindent\textbf{odděliti}\quad{\footnotesize\texttt{TERM-0051}}\par
\noindent\emph{Iztočnik / Source:} abzuspalten\par
\noindent\emph{Rusky / Russian:} выделить\par
\noindent\emph{Ukrajinsky / Ukrainian:} відокремити\par
\smallskip\noindent\emph{Razlog:} Odděliti vyražaje izdvojenje najviše jednogo zajedničnogo dělitelja iz vseh polinomov osnovy. Izbor ne zaměnjaje to ograničeno dejanje tvrdženjem o jednoznačnoj faktorizaciji v cělom prstenu.\par
\noindent\emph{Reason:} Odděliti denotes splitting off at most one common divisor from all basis polynomials. The choice does not turn this limited operation into a claim about unique factorization throughout a ring.\par
\noindent{\footnotesize Izvorne zapisy / Original records: HIST-g049-015.}\par
\end{minipage}\par
\par\medskip\Needspace{7\baselineskip}
\noindent\begin{minipage}{\linewidth}
\hypertarget{TERM-0052}{}
\noindent\textbf{izčetno beskonečno mnogo klasov}\quad{\footnotesize\texttt{TERM-0052}}\par
\noindent\emph{Iztočnik / Source:} abzählbar unendlich viele Klassen\par
\noindent\emph{Rusky / Russian:} счётно бесконечно много классов\par
\noindent\emph{Ukrajinsky / Ukrainian:} зліченно нескінченно багато класів\par
\smallskip\noindent\emph{Razlog:} Izčetno i beskonečno sut oba nužne česti zaključka. Priměr daje beskonečno mnogo neizomorfnyh klasov, ale familija jest izčetna. Izostavjenje jednoj česti bi oslabilo ili změnilo tvrdženje.\par
\noindent\emph{Reason:} Both countable and infinite are essential. The example supplies infinitely many non-isomorphic classes, with a countable family. Omitting either qualification would weaken or alter the claim.\par
\noindent{\footnotesize Izvorne zapisy / Original records: HIST-g067-001.}\par
\end{minipage}\par
\par\medskip\Needspace{7\baselineskip}
\noindent\begin{minipage}{\linewidth}
\hypertarget{TERM-0053}{}
\noindent\textbf{izčisljenje ostatkovyh klasov idealov}\quad{\footnotesize\texttt{TERM-0053}}\par
\noindent\emph{Iztočnik / Source:} Abzählung von Restklassen der Ideale\par
\noindent\emph{Rusky / Russian:} подсчёт классов вычетов идеалов\par
\noindent\emph{Ukrajinsky / Ukrainian:} підрахунок класів лишків ідеалів\par
\smallskip\noindent\emph{Razlog:} Ostatkove klasy prodolžujut terminologiju děl 24 i 25. Izčisljenje označaje brojanje tųh klasov; zapis odklanjaje automaticku zaměnu izrazom kvocientne klasy, ibo iztočnik govori o ostatkovyh klasah i jih broju.\par
\noindent\emph{Reason:} Ostatkove klasy continues the terminology of Papers 24 and 25. Izčisljenje denotes their enumeration. The record avoids automatically substituting quotient-class terminology because the source speaks of residue classes and their count.\par
\noindent{\footnotesize Izvorne zapisy / Original records: HIST-g104-002.}\par
\end{minipage}\par
\par\medskip\Needspace{7\baselineskip}
\noindent\begin{minipage}{\linewidth}
\hypertarget{TERM-0054}{}
\noindent\textbf{adjungovany modul}\quad{\footnotesize\texttt{TERM-0054}}\par
\noindent\emph{Iztočnik / Source:} adj. Modul\par
\noindent\emph{Rusky / Russian:} присоед. модуль\par
\noindent\emph{Ukrajinsky / Ukrainian:} приєдн. модуль\par
\smallskip\noindent\emph{Razlog:} Kratka oznaka pri formule jest razvinuta v adjungovany modul i vezana s terminologijeju adjunkcije v oddělu 7. Sam izvor označaje familiju adjungovati kako predmet daljnjej jezyčnoj prověrky; kratkost ne zaměnjaje tę prověrku.\par
\noindent\emph{Reason:} The short formula annotation is expanded as adjungovany modul and linked to the adjoining terminology of Section 7. The record treats the adjungovati family as needing further linguistic assessment; brevity is not a substitute for that assessment.\par
\noindent{\footnotesize Izvorne zapisy / Original records: HIST-g010-006.}\par
\end{minipage}\par
\par\medskip\Needspace{7\baselineskip}
\noindent\begin{minipage}{\linewidth}
\hypertarget{TERM-0055}{}
\noindent\textbf{adjungujemo}\quad{\footnotesize\texttt{TERM-0055}}\par
\noindent\emph{Iztočnik / Source:} adjungieren\par
\noindent\emph{Rusky / Russian:} присоединяем\par
\noindent\emph{Ukrajinsky / Ukrainian:} ад'юнктуємо\par
\smallskip\noindent\emph{Razlog:} Adjungujemo zachovuje odborny korenj za dodavanje neopreděljenyh elementov k polju. Ukrajinsky i medžuslovjansky kandidaty upotrěbjajut adjunkcijnu terminologiju, rusky sravniteljny izraz jest присоединять. Zapis zachovuje tu razliku vměsto pretvarjanja vseh sravnjenij v jedan korenj.\par
\noindent\emph{Reason:} Adjungujemo preserves the technical root for adjoining indeterminates to a field. The Ukrainian and Interslavic proposals use adjunction vocabulary, while the Russian comparandum is присоединять. The comparison retains this difference rather than pretending that all witnesses share one root.\par
\noindent{\footnotesize Izvorne zapisy / Original records: HIST-g079-010.}\par
\end{minipage}\par
\par\medskip\Needspace{7\baselineskip}
\noindent\begin{minipage}{\linewidth}
\hypertarget{TERM-0056}{}
\noindent\textbf{pridavati}\quad{\footnotesize\texttt{TERM-0056}}\par
\noindent\emph{Iztočnik / Source:} adjungieren\par
\noindent\emph{Rusky / Russian:} присоединять\par
\noindent\emph{Ukrajinsky / Ukrainian:} приєднувати\par
\smallskip\noindent\emph{Razlog:} Pridavati ima izjasniti dodavanje objektov k sistemu. Zapis izrično izběgaje sugerovati adjungovany funktor ili operaciju adjugaty. Izbor zato zavisi od konkretnogo algebraičnogo dejanja, ne samo od podobnosti slov.\par
\noindent\emph{Reason:} Pridavati is intended to express adding objects to a system. The reason explicitly avoids suggesting an adjoint functor or an adjugate operation. The choice depends on the particular algebraic operation, not merely on word resemblance.\par
\noindent{\footnotesize Izvorne zapisy / Original records: HIST-g002-006.}\par
\end{minipage}\par
\par\medskip\Needspace{7\baselineskip}
\noindent\begin{minipage}{\linewidth}
\hypertarget{TERM-0057}{}
\noindent\textbf{adjungovati funkciju log x k oblasti racionalnosti}\quad{\footnotesize\texttt{TERM-0057}}\par
\noindent\emph{Iztočnik / Source:} adjungieren von log x zum Rationalitaetsbereich\par
\noindent\emph{Rusky / Russian:} присоединить \textbackslash{}log x к области рациональности\par
\noindent\emph{Ukrajinsky / Ukrainian:} ад'юнктувати \textbackslash{}log x до області раціональності\par
\smallskip\noindent\emph{Razlog:} V četvrtom priměru dodavanje log x menja oblast koeficientov oziroma racionalnosti. Medžuslovjansky i ukrajinsky izrazy zachovujut korenj adjunkcije. V zachovanom razlogu rusko slovo ima samo znaky ?; njegov točny zapis ne može byti citovany iz tutogo dokaza.\par
\noindent\emph{Reason:} In Example 4, adjoining log x changes the coefficient or rationality domain. The Interslavic and Ukrainian wording retains the adjunction root. The preserved rationale renders the Russian word only as question marks; its exact spelling cannot be quoted from this evidence.\par
\noindent{\footnotesize Izvorne zapisy / Original records: HIST-g079-077.}\par
\end{minipage}\par
\par\medskip\Needspace{7\baselineskip}
\noindent\begin{minipage}{\linewidth}
\hypertarget{TERM-0058}{}
\noindent\textbf{adjungovane formy/funkcije/krive}\quad{\footnotesize\texttt{TERM-0058}}\par
\noindent\emph{Iztočnik / Source:} adjungierte Formen/Funktionen/Kurven\par
\noindent\emph{Rusky / Russian:} адъюнктные формы/функции/кривые\par
\noindent\emph{Ukrajinsky / Ukrainian:} ад'юнктні форми/функції/криві\par
\smallskip\noindent\emph{Razlog:} Adjungovane zachovuje historijno algebraično-geometrično nazvanje form, funkcij i krivyh. Presudna jest veza s geometričnoju teorijeju, s ktoroju tekst sravnjaje konstrukciju. Izvor ne rěšaje vsih drugyh smyslov německogo adjungiert tym samym slovom.\par
\noindent\emph{Reason:} Adjungovane preserves the older algebraic-geometric designation of forms, functions, and curves. The reason is its role in the comparison with geometric theory. The record does not resolve every other sense of German adjungiert with the same word.\par
\noindent{\footnotesize Izvorne zapisy / Original records: HIST-g076-018.}\par
\end{minipage}\par
\par\medskip\Needspace{7\baselineskip}
\noindent\begin{minipage}{\linewidth}
\hypertarget{TERM-0059}{}
\noindent\textbf{adjunkcija}\quad{\footnotesize\texttt{TERM-0059}}\par
\noindent\emph{Iztočnik / Source:} Adjunktion\par
\noindent\emph{Rusky / Russian:} присоединение\par
\noindent\emph{Ukrajinsky / Ukrainian:} приєднання\par
\smallskip\noindent\emph{Razlog:} Adjunkcija jest internacionalny termin za algebraično dodavanje, doslědno upotrěbljany v dělu 9. Zapisy go sravnjajut s ruskim i ukrajinskim odbornym vyraženjem i smatrajut polno slovjansku skovanicu mene stabilnoju. Odděl 4 takože dava objasnjenje pojma. Jezyčna prověrka ostavaje izrično otvorjena.\par
\noindent\emph{Reason:} Adjunkcija is the international candidate for algebraic adjoining, used consistently in Paper 9. The records compare Russian and Ukrainian mathematical phrasing and regard a fully Slavic calque as less stable. Section 4 also explains the concept. Linguistic assessment remains expressly open.\par
\noindent{\footnotesize Izvorne zapisy / Original records: HIST-g041-008, HIST-g044-011, HIST-g058-011, HIST-g071-010.}\par
\end{minipage}\par
\par\medskip\Needspace{7\baselineskip}
\noindent\begin{minipage}{\linewidth}
\hypertarget{TERM-0060}{}
\noindent\textbf{adjunkcija / adjungovati}\quad{\footnotesize\texttt{TERM-0060}}\par
\noindent\emph{Iztočnik / Source:} Adjunktion / adjungieren\par
\noindent\emph{Rusky / Russian:} присоединение / присоединять\par
\noindent\emph{Ukrajinsky / Ukrainian:} приєднання / приєднувати\par
\smallskip\noindent\emph{Razlog:} Para adjunkcija/adjungovati zachovuje jedan korenj za imenik i glagol v teoriji razširenij polja. Izvorny razlog preferuje tu kratku internacionalnu familiju prěd dolžšim opisnym kalkom; ne dava čislene ocěny razumljivosti.\par
\noindent\emph{Reason:} The noun–verb pair retains one root in field-extension discourse. The recorded reason prefers the compact international family to a longer descriptive calque. It provides no numerical estimate of intelligibility.\par
\noindent{\footnotesize Izvorne zapisy / Original records: HIST-g068-011.}\par
\end{minipage}\par
\par\medskip\Needspace{7\baselineskip}
\noindent\begin{minipage}{\linewidth}
\hypertarget{TERM-0061}{}
\noindent\textbf{prisojedinjenje u\_\{\textbackslash{}mu\textbackslash{}nu\}}\quad{\footnotesize\texttt{TERM-0061}}\par
\noindent\emph{Iztočnik / Source:} Adjunktion der u\_\{\textbackslash{}mu\textbackslash{}nu\}\par
\noindent\emph{Rusky / Russian:} присоединение u\_\{\textbackslash{}mu\textbackslash{}nu\}\par
\noindent\emph{Ukrajinsky / Ukrainian:} приєднання u\_\{\textbackslash{}mu\textbackslash{}nu\}\par
\smallskip\noindent\emph{Razlog:} Prisojedinjenje označaje algebraično dodavanje elementov k polju, ne obyčno pribavjanje. Izvorny razlog se opira na podobnost s ruskim i ukrajinskim izrazom, predpolaganu razumljivost za vzhodne i južne Slovjane i doslědnost s drugymi pasažami o razširenjah polja.\par
\noindent\emph{Reason:} Prisojedinjenje names algebraic adjoining, not everyday addition. The recorded reasons are the Russian and Ukrainian comparanda, proposed East/South Slavic recognizability, and consistency with other field-extension passages. Recognizability was an editorial judgment, not a measured result.\par
\noindent{\footnotesize Izvorne zapisy / Original records: HIST-g096-002.}\par
\end{minipage}\par
\par\medskip\Needspace{7\baselineskip}
\noindent\begin{minipage}{\linewidth}
\hypertarget{TERM-0062}{}
\noindent\textbf{prisojedinjenje neopreděljenyh u\_\{mu nu\}}\quad{\footnotesize\texttt{TERM-0062}}\par
\noindent\emph{Iztočnik / Source:} Adjunktion der Unbestimmten u\_\{mu nu\}\par
\noindent\emph{Rusky / Russian:} присоединение неопределенных u\_\{mu nu\}\par
\noindent\emph{Ukrajinsky / Ukrainian:} приєднання невизначених u\_\{mu nu\}\par
\smallskip\noindent\emph{Razlog:} Neopreděljene u\_\{mu nu\} sut prisojedinjene k polju koeficientov. Izraz mora izjasniti razširenje polja i ne dopuščati čitanje kako obyčnu substituciju. To jest matematičny razlog za izbor specializovanogo glagolnogo imenika.\par
\noindent\emph{Reason:} The indeterminates u\_\{mu nu\} are adjoined to the coefficient field. The expression must indicate a field extension, not a mere substitution. That mathematical distinction motivates the specialized adjoining terminology.\par
\noindent{\footnotesize Izvorne zapisy / Original records: HIST-g087-003, HIST-g092-023.}\par
\end{minipage}\par
\par\medskip\Needspace{7\baselineskip}
\noindent\begin{minipage}{\linewidth}
\hypertarget{TERM-0063}{}
\noindent\textbf{prisojedinjenje primitivnogo elementa prvogo roda}\quad{\footnotesize\texttt{TERM-0063}}\par
\noindent\emph{Iztočnik / Source:} Adjunktion eines primitiven Elementes erster Art\par
\noindent\emph{Rusky / Russian:} присоединение примитивного элемента первого рода\par
\noindent\emph{Ukrajinsky / Ukrainian:} приєднання примітивного елемента першого роду\par
\smallskip\noindent\emph{Razlog:} Prisojedinjenje primitivnogo elementa prvogo roda zachovuje konkretny krok avtorskoj redukcije. Zapis izrično odklanjaje zaměnu opčim vyraženjem o generovanom razširenju, ktoro by skrylo taj krok i klasifikaciju elementa.\par
\noindent\emph{Reason:} The wording retains the particular step in Noether's reduction: adjoining a primitive element of the first kind. The record rejects replacing it with an abstract generated-extension expression that would obscure the operation and the element's classification.\par
\noindent{\footnotesize Izvorne zapisy / Original records: HIST-g127-004.}\par
\end{minipage}\par
\par\medskip\Needspace{7\baselineskip}
\noindent\begin{minipage}{\linewidth}
\hypertarget{TERM-0064}{}
\noindent\textbf{adjunkcija sqrt(x) / modulna baza}\quad{\footnotesize\texttt{TERM-0064}}\par
\noindent\emph{Iztočnik / Source:} Adjunktion von sqrt(x) / Modulbasis\par
\noindent\emph{Rusky / Russian:} присоединение sqrt(x) / модульный базис\par
\noindent\emph{Ukrajinsky / Ukrainian:} приєднання sqrt(x) / модульна база\par
\smallskip\noindent\emph{Razlog:} Adjunkcija prodolžuje terminologiju razširenij; internacionalny korenj byl izbrany radi kratkosti i radi izběganja nejasnosti nove skovanice. Modulna baza ostavaje svazana s konkretnym priměrom dodavanja sqrt(x).\par
\noindent\emph{Reason:} Adjunkcija continues the extension terminology. The international root was preferred for compactness and to avoid the ambiguity of a new calque. Modulna baza remains tied to the specific example adjoining sqrt(x).\par
\noindent{\footnotesize Izvorne zapisy / Original records: HIST-g153-006.}\par
\end{minipage}\par
\par\medskip\Needspace{7\baselineskip}
\noindent\begin{minipage}{\linewidth}
\hypertarget{TERM-0065}{}
\noindent\textbf{afekt / bez afekta}\quad{\footnotesize\texttt{TERM-0065}}\par
\noindent\emph{Iztočnik / Source:} Affekt / affektlos\par
\noindent\emph{Rusky / Russian:} аффект / без аффекта\par
\noindent\emph{Ukrajinsky / Ukrainian:} афект / без афекту\par
\smallskip\noindent\emph{Razlog:} Afekt jest zachovany kak historijno algebraično nazvanje iz literatury o kubičnyh i kvartičnyh jednačenjah. Zapis ostavjaje njegovu točnu interpretaciju za odbornu prověrku i ne zaměnjaje go nepodkrepljenym sučasnym ekvivalentom.\par
\noindent\emph{Reason:} Afekt preserves a historical algebraic term from work on cubic and quartic equations. The record leaves its precise interpretation for subject review rather than substituting an unsupported modern equivalent.\par
\noindent{\footnotesize Izvorne zapisy / Original records: HIST-g073-017.}\par
\end{minipage}\par
\par\medskip\Needspace{7\baselineskip}
\noindent\begin{minipage}{\linewidth}
\hypertarget{TERM-0066}{}
\noindent\textbf{jednačenja "bez afekta"}\quad{\footnotesize\texttt{TERM-0066}}\par
\noindent\emph{Iztočnik / Source:} affektlose Gleichungen\par
\noindent\emph{Rusky / Russian:} уравнения "без аффекта"\par
\noindent\emph{Ukrajinsky / Ukrainian:} рівняння "без афекту"\par
\smallskip\noindent\emph{Razlog:} Izraz bez afekta označaje historijnu klasu jednačenij. On ne jest tiho zaměnjen slovom nerazložime: v argumentu teorem Hilberta o nerazložimosti služi za prěhod od prostogo parametričnogo jednačenja k toj klasě. Prěvod zachovuje razliku medžu nazvanjem i mehanizmom dokaza.\par
\noindent\emph{Reason:} Without affect is retained as a historical class designation. It is not silently replaced by irreducible: Hilbert's irreducibility theorem supplies the mechanism connecting the simple parametric equation to that class. The wording distinguishes the designation from the proof mechanism.\par
\noindent{\footnotesize Izvorne zapisy / Original records: HIST-g039-013.}\par
\end{minipage}\par
\par\medskip\Needspace{7\baselineskip}
\noindent\begin{minipage}{\linewidth}
\hypertarget{TERM-0067}{}
\noindent\textbf{afinna grupa}\quad{\footnotesize\texttt{TERM-0067}}\par
\noindent\emph{Iztočnik / Source:} affine Gruppe\par
\noindent\emph{Rusky / Russian:} аффинная группа\par
\noindent\emph{Ukrajinsky / Ukrainian:} афінна група\par
\smallskip\noindent\emph{Razlog:} Afinna grupa jest internacionalno nazvanje grupe v završnoj redukciji nehomogennyh funkcij pomoču homogenizacije. Rusky i ukrajinsky izrazy dajut neposrědne odborné sravnjenja; izvorny zapis ne prědlaga novy slovjansky korenj.\par
\noindent\emph{Reason:} Afinna grupa supplies the international group-theoretic term in the concluding reduction of nonhomogeneous functions by homogenization. Russian and Ukrainian provide direct technical comparanda. The record proposes no new Slavic root.\par
\noindent{\footnotesize Izvorne zapisy / Original records: HIST-g074-017.}\par
\end{minipage}\par
\par\medskip\Needspace{7\baselineskip}
\noindent\begin{minipage}{\linewidth}
\hypertarget{TERM-0068}{}
\noindent\textbf{algebraično obrazovanje}\quad{\footnotesize\texttt{TERM-0068}}\par
\noindent\emph{Iztočnik / Source:} Algebraic Gebilde\par
\noindent\emph{Rusky / Russian:} алгебраическое образование\par
\noindent\emph{Ukrajinsky / Ukrainian:} алгебраїчне утворення\par
\smallskip\noindent\emph{Razlog:} Obrazovanje jest doslovny kandidat za opče strukturno Gebilde. Izvorny zapis izrično dopuščaje, že sučasny matematičny kontekst može trěbovati točnjeje nazvanje raznovidnosti. Kandidat zato ne ustanavjaje univerzalny prěvod slova Gebilde.\par
\noindent\emph{Reason:} Obrazovanje is a literal candidate for the broad structural word Gebilde. The record expressly leaves open whether the mathematical context requires terminology specific to a variety. It does not establish a universal translation of Gebilde.\par
\noindent{\footnotesize Izvorne zapisy / Original records: HIST-g091-008, HIST-g092-077.}\par
\end{minipage}\par
\par\medskip\Needspace{7\baselineskip}
\noindent\begin{minipage}{\linewidth}
\hypertarget{TERM-0069}{}
\noindent\textbf{algebraično zamknjeno polje}\quad{\footnotesize\texttt{TERM-0069}}\par
\noindent\emph{Iztočnik / Source:} algebraisch abgeschlossener Körper\par
\noindent\emph{Rusky / Russian:} алгебраически замкнутое поле\par
\noindent\emph{Ukrajinsky / Ukrainian:} алгебрично замкнене поле\par
\smallskip\noindent\emph{Razlog:} Zamknjeno polje slijedi rusku i ukrajinsku konstrukciju za algebraičnu zatvorjenost. Termin jest važny v podčrtnoj primětce o razširenjah polja Steinitza; razlog izbora jest i jasno sravnjenje istogo pojma medžu prěvodami.\par
\noindent\emph{Reason:} Zamknjeno polje follows the Russian and Ukrainian construction for algebraic closedness. The term anchors the Steinitz field-extension footnote, so cross-language identification of the same concept is part of the reason.\par
\noindent{\footnotesize Izvorne zapisy / Original records: HIST-g083-002.}\par
\smallskip\noindent\emph{Sravnjenje s tekstom:} Prěvod upotrěbjaje i algebraično zamknjeno polje, i algebraično zatvorjeno polje v pasažah o algebraičnoj zamknjenosti. Kontekst s razloženjem na linearne faktory podpiraje isti matematičny smysl, ale samo prisutnost oběh form ne izbira jednu iz njih kako normu. Historijne zapisy trěba zato čitati kako dokumentovane alternativy. Te tri pasaže ne dokazyvajut, že vsako pojavjenje zamknjen- ili zatvorjen- v korpusu ima isti smysl.\par
\noindent\emph{Textual comparison:} The translation uses both algebraično zamknjeno polje and algebraično zatvorjeno polje in passages concerning algebraic closedness. Context involving factorization into linear factors supports the same mathematical sense, but the coexistence of the forms does not select either as the norm. The historical records must therefore be read as documented alternatives. These three passages do not establish that every occurrence of zamknjen- or zatvorjen- in the corpus has the same sense.\par
\noindent{\footnotesize TEXT-COMP-02; pasaže / passages: 13803, 14930, 15719; \texttt{CORPUS\_FORM\_COMPARISONS.json}.}\par
\end{minipage}\par
\par\medskip\Needspace{7\baselineskip}
\noindent\begin{minipage}{\linewidth}
\hypertarget{TERM-0070}{}
\noindent\textbf{algebraično zamknjeno polje}\quad{\footnotesize\texttt{TERM-0070}}\par
\noindent\emph{Iztočnik / Source:} algebraisch abgeschlossenes Feld\par
\noindent\emph{Rusky / Russian:} алгебраически замкнутое поле\par
\noindent\emph{Ukrajinsky / Ukrainian:} алгебраїчно замкнене поле\par
\smallskip\noindent\emph{Razlog:} Zamknjeno jest izbrano radi predpolaganoj medžuslovjanskoj razumljivosti i podobnosti s ruskim i ukrajinskim odbornym pridavnikom. Zatvorjeno jest zabilěženo kako stvarna alternativa. Razlog ne dava empiricku ocěnu razumljivosti ni konačno opče razrěšenje medžu oběma formami.\par
\noindent\emph{Reason:} Zamknjeno is preferred for proposed pan-Slavic transparency and similarity to the Russian and Ukrainian technical adjective. Zatvorjeno is a documented alternative. The reason supplies neither an empirical comprehension estimate nor a universal resolution between them.\par
\noindent{\footnotesize Izvorne zapisy / Original records: HIST-g101-003.}\par
\smallskip\noindent\emph{Sravnjenje s tekstom:} Prěvod upotrěbjaje i algebraično zamknjeno polje, i algebraično zatvorjeno polje v pasažah o algebraičnoj zamknjenosti. Kontekst s razloženjem na linearne faktory podpiraje isti matematičny smysl, ale samo prisutnost oběh form ne izbira jednu iz njih kako normu. Historijne zapisy trěba zato čitati kako dokumentovane alternativy. Te tri pasaže ne dokazyvajut, že vsako pojavjenje zamknjen- ili zatvorjen- v korpusu ima isti smysl.\par
\noindent\emph{Textual comparison:} The translation uses both algebraično zamknjeno polje and algebraično zatvorjeno polje in passages concerning algebraic closedness. Context involving factorization into linear factors supports the same mathematical sense, but the coexistence of the forms does not select either as the norm. The historical records must therefore be read as documented alternatives. These three passages do not establish that every occurrence of zamknjen- or zatvorjen- in the corpus has the same sense.\par
\noindent{\footnotesize TEXT-COMP-02; pasaže / passages: 13803, 14930, 15719; \texttt{CORPUS\_FORM\_COMPARISONS.json}.}\par
\end{minipage}\par
\par\medskip\Needspace{7\baselineskip}
\noindent\begin{minipage}{\linewidth}
\hypertarget{TERM-0071}{}
\noindent\textbf{algebraično zavisimy / nezavisimy}\quad{\footnotesize\texttt{TERM-0071}}\par
\noindent\emph{Iztočnik / Source:} algebraisch abhängig / unabhängig\par
\noindent\emph{Rusky / Russian:} алгебраически зависимый / независимый\par
\noindent\emph{Ukrajinsky / Ukrainian:} алгебраїчно залежний / незалежний\par
\smallskip\noindent\emph{Razlog:} Algebraična zavisimost jest izrično odlišena od racionalnoj zavisimosti iz odděla 6. Zapis objasnjaje smyslovu razliku, ale ne dava osobno sravnjenje morfologičnyh kandidatov zavisimy i zavisny.\par
\noindent\emph{Reason:} The expression explicitly distinguishes algebraic from the rational dependence discussed in Section 6. The record explains that semantic contrast, but supplies no separate morphological comparison of zavisimy and zavisny.\par
\noindent{\footnotesize Izvorne zapisy / Original records: HIST-g065-009.}\par
\end{minipage}\par
\par\medskip\Needspace{7\baselineskip}
\noindent\begin{minipage}{\linewidth}
\hypertarget{TERM-0072}{}
\noindent\textbf{algebraično zavisny / nezavisny}\quad{\footnotesize\texttt{TERM-0072}}\par
\noindent\emph{Iztočnik / Source:} algebraisch abhängig / unabhängig\par
\noindent\emph{Rusky / Russian:} алгебраически зависимый / независимый\par
\noindent\emph{Ukrajinsky / Ukrainian:} алгебраїчно залежний / незалежний\par
\smallskip\noindent\emph{Razlog:} Zavisny/nezavisny prodolžuje terminologiju drugej česti i razlikuje dva slučaje argumenta. Zachovany razlog jest doslědnost konkretnogo izloženja; sam ne dokazuje prědnost toj morfologije nad kandidatami s -im-.\par
\noindent\emph{Reason:} Zavisny/nezavisny continues the vocabulary of Part 2 and distinguishes the argument's two cases. The surviving reason is continuity within the exposition, not a demonstrated preference over the competing -im- forms.\par
\noindent{\footnotesize Izvorne zapisy / Original records: HIST-g071-008.}\par
\end{minipage}\par
\par\medskip\Needspace{7\baselineskip}
\noindent\begin{minipage}{\linewidth}
\hypertarget{TERM-0073}{}
\noindent\textbf{algebraično zavisne funkcije}\quad{\footnotesize\texttt{TERM-0073}}\par
\noindent\emph{Iztočnik / Source:} algebraisch abhängige Funktionen\par
\noindent\emph{Rusky / Russian:} алгебраически зависимые функции\par
\noindent\emph{Ukrajinsky / Ukrainian:} алгебраїчно залежні функції\par
\smallskip\noindent\emph{Razlog:} Algebraično zavisne funkcije sut vezane s opreděljenjem algebraičnogo ranga v prvom oddělu. Izbor upotrěbjaje sravniteljne ruske i ukrajinske nazvanja zavisimosti i zachovuje funkcionovanje termina v tom opreděljenju.\par
\noindent\emph{Reason:} The term is tied to the definition of algebraic rank in Section 1. Russian and Ukrainian dependence terminology provides comparanda; the wording preserves the term's role in that definition.\par
\noindent{\footnotesize Izvorne zapisy / Original records: HIST-g042-002.}\par
\end{minipage}\par
\par\medskip\Needspace{7\baselineskip}
\noindent\begin{minipage}{\linewidth}
\hypertarget{TERM-0074}{}
\noindent\textbf{integralno zamknjene vzhodno k svojemu polju kvocientov}\quad{\footnotesize\texttt{TERM-0074}}\par
\noindent\emph{Iztočnik / Source:} algebraisch ganz abgeschlossen in bezug auf ihren Quotientenkörper\par
\noindent\emph{Rusky / Russian:} интегрально замкнуты относительно своего поля частных\par
\noindent\emph{Ukrajinsky / Ukrainian:} інтегрально замкнені відносно свого поля часток\par
\smallskip\noindent\emph{Razlog:} Integralno zamknjene izjasnjaje zatvorjenost vzhledem na algebraičnu cělost v polju kvocientov. To jest sučasno čitanje německogo algebraisch ganz. Izvorny zapis ostavjaje konstrukciju vzhodno k i izraz polje kvocientov za jezyčnu prověrku; matematičny smysl ne rěšaje sam jih idiomatičnost.\par
\noindent\emph{Reason:} Integralno zamknjene gives the modern integrally-closed reading relative to the field of fractions. The record leaves vzhodno k and polje kvocientov for linguistic assessment. Identifying the mathematical sense does not by itself establish their idiomatic suitability.\par
\noindent{\footnotesize Izvorne zapisy / Original records: HIST-g103-006.}\par
\end{minipage}\par
\par\medskip\Needspace{7\baselineskip}
\noindent\begin{minipage}{\linewidth}
\hypertarget{TERM-0075}{}
\noindent\textbf{jest algebraično cěla nad}\quad{\footnotesize\texttt{TERM-0075}}\par
\noindent\emph{Iztočnik / Source:} algebraisch ganz abhängt\par
\noindent\emph{Rusky / Russian:} является алгебраически целой над\par
\noindent\emph{Ukrajinsky / Ukrainian:} є алгебраїчно цілою над\par
\smallskip\noindent\emph{Razlog:} Algebraično cěla nad vyražaje odnos integralnosti v argumentu o konečnosti. Prědlog nad čini osnovu odnosa izričnoju i odliča tuto od obyčnoj algebraičnoj zavisimosti. Forma cěla jest kontekstna ženska forma, ne druga definicija integralnosti.\par
\noindent\emph{Reason:} The phrase expresses integrality over a base in the finiteness argument. Nad makes the base relation explicit rather than suggesting ordinary algebraic dependence. Cěla is the contextual feminine form, not a different definition of integrality.\par
\noindent{\footnotesize Izvorne zapisy / Original records: HIST-g094-010.}\par
\end{minipage}\par
\par\medskip\Needspace{7\baselineskip}
\noindent\begin{minipage}{\linewidth}
\hypertarget{TERM-0076}{}
\noindent\textbf{jest algebraično cěly nad}\quad{\footnotesize\texttt{TERM-0076}}\par
\noindent\emph{Iztočnik / Source:} algebraisch ganz abhängt\par
\noindent\emph{Rusky / Russian:} алгебраически цел над\par
\noindent\emph{Ukrajinsky / Ukrainian:} є алгебраїчно цілим над\par
\smallskip\noindent\emph{Razlog:} Německo vyraženje zavisimosti jest tumačeno kako byti integralnym nad osnovoju. Izbor izběgaje směšenje s obyčnoju algebraičnoju zavisimostju ili nezavisimostju. Tuto jest izrično terminologično osuvrěmenjenje, ktoro mora ostati svazano s iztočnym uslovjem.\par
\noindent\emph{Reason:} The German dependence wording is interpreted as being integral over a base. The reason is to avoid confusion with ordinary algebraic dependence or independence. This is an explicit terminological modernization that must remain tied to the source condition.\par
\noindent{\footnotesize Izvorne zapisy / Original records: HIST-g050-004.}\par
\end{minipage}\par
\par\medskip\Needspace{7\baselineskip}
\noindent\begin{minipage}{\linewidth}
\hypertarget{TERM-0077}{}
\noindent\textbf{jest algebraično cěly nad}\quad{\footnotesize\texttt{TERM-0077}}\par
\noindent\emph{Iztočnik / Source:} algebraisch ganz von ... abhängt / algebraisch-ganz abhängig\par
\noindent\emph{Rusky / Russian:} алгебраически цел над\par
\noindent\emph{Ukrajinsky / Ukrainian:} є алгебраїчно цілим над\par
\smallskip\noindent\emph{Razlog:} Izraz prodolžuje čitanje integralno nad iz odděla 10. Obě německe formulacije dostavajut isti odnos, ale ne slivajut se s opčim pojmom algebraičnoj zavisimosti. Razlog jest semantična razlika i doslědnost susědnyh oddělov.\par
\noindent\emph{Reason:} The expression continues the integral-over interpretation of Section 10. Both German formulations denote that relation without being merged into ordinary algebraic dependence. Semantic distinction and continuity between sections motivate the choice.\par
\noindent{\footnotesize Izvorne zapisy / Original records: HIST-g051-002.}\par
\end{minipage}\par
\par\medskip\Needspace{7\baselineskip}
\noindent\begin{minipage}{\linewidth}
\hypertarget{TERM-0078}{}
\noindent\textbf{algebraično cěly nad R}\quad{\footnotesize\texttt{TERM-0078}}\par
\noindent\emph{Iztočnik / Source:} algebraisch ganz von R abhängt\par
\noindent\emph{Rusky / Russian:} алгебраически целый над R\par
\noindent\emph{Ukrajinsky / Ukrainian:} алгебраїчно цілий над R\par
\smallskip\noindent\emph{Razlog:} R jest izrično nazvana osnova integralnosti. Vyraženje nad R izběgaje čitanje německogo abhängt kako obyčnu algebraičnu zavisimost. Izbor zachovuje odnos k konkretnomu R, ne samo pridavnik cěly.\par
\noindent\emph{Reason:} R is explicitly named as the base of integrality. Nad R avoids reading abhängt as ordinary algebraic dependence. The wording preserves the relation to this particular R, not merely the adjective integral.\par
\noindent{\footnotesize Izvorne zapisy / Original records: HIST-g106-006.}\par
\end{minipage}\par
\par\medskip\Needspace{7\baselineskip}
\noindent\begin{minipage}{\linewidth}
\hypertarget{TERM-0079}{}
\noindent\textbf{algebraično-drobno čislo}\quad{\footnotesize\texttt{TERM-0079}}\par
\noindent\emph{Iztočnik / Source:} algebraisch gebrochene Zahl\par
\noindent\emph{Rusky / Russian:} алгебраически-дробное число\par
\noindent\emph{Ukrajinsky / Ukrainian:} алгебраїчно-дробове число\par
\smallskip\noindent\emph{Razlog:} Algebraično-drobno ostavjaje vidimu avtorsku razliku medžu algebraično cělymi i drobnymi čislami v uslovju dopustimosti. Spojnik veže obě česti v tehnično označenje; slovo drobno ne jest izostavjeno kako nepotrěbna stilistična čest.\par
\noindent\emph{Reason:} The compound preserves Noether's contrast between algebraic integers and fractional algebraic numbers in the admissibility condition. Its fractional qualification is substantive, not a stylistic element that can be dropped.\par
\noindent{\footnotesize Izvorne zapisy / Original records: HIST-g062-009.}\par
\end{minipage}\par
\par\medskip\Needspace{7\baselineskip}
\noindent\begin{minipage}{\linewidth}
\hypertarget{TERM-0080}{}
\noindent\textbf{algebraično nezavisne}\quad{\footnotesize\texttt{TERM-0080}}\par
\noindent\emph{Iztočnik / Source:} algebraisch unabhängig\par
\noindent\emph{Rusky / Russian:} алгебраически независимые\par
\noindent\emph{Ukrajinsky / Ukrainian:} алгебраїчно незалежні\par
\smallskip\noindent\emph{Razlog:} Nezavisne imenuje algebraičnu nezavisimost, ktora v opreděljenju IV slijedi iz minimalnosti. Izbor zachovuje tu logičnu rolu pojma; ne govori o obyčnoj nezavisimosti ili o nezavisimosti slučajnyh veličin.\par
\noindent\emph{Reason:} Nezavisne denotes the algebraic independence that Definition IV derives from minimality. The wording preserves that logical role; it is not ordinary independence or probabilistic independence.\par
\noindent{\footnotesize Izvorne zapisy / Original records: HIST-g046-003.}\par
\end{minipage}\par
\par\medskip\Needspace{7\baselineskip}
\noindent\begin{minipage}{\linewidth}
\hypertarget{TERM-0081}{}
\noindent\textbf{algebraično nezavisne / zavisne}\quad{\footnotesize\texttt{TERM-0081}}\par
\noindent\emph{Iztočnik / Source:} algebraisch unabhängig / abhängig\par
\noindent\emph{Rusky / Russian:} алгебраически независимые / зависимые\par
\noindent\emph{Ukrajinsky / Ukrainian:} алгебраїчно незалежні / залежні\par
\smallskip\noindent\emph{Razlog:} Izvorny razlog upučuje na obyčnu algebraičnu terminologiju zavisimosti i nezavisimosti. Ruske i ukrajinske formy sut zabilěžene za sravnjenje. Posebna odklonjena alternativa ili čislena vaga ne jest navedena.\par
\noindent\emph{Reason:} The surviving reason invokes ordinary algebraic dependence/independence terminology, with Russian and Ukrainian forms for comparison. No particular rejected alternative or numerical weighting is recorded.\par
\noindent{\footnotesize Izvorne zapisy / Original records: HIST-g041-011.}\par
\end{minipage}\par
\par\medskip\Needspace{7\baselineskip}
\noindent\begin{minipage}{\linewidth}
\hypertarget{TERM-0082}{}
\noindent\textbf{algebraično nezavisne / transcendentne}\quad{\footnotesize\texttt{TERM-0082}}\par
\noindent\emph{Iztočnik / Source:} algebraisch unabhängig / transzendent\par
\noindent\emph{Rusky / Russian:} алгебраически независимые / трансцендентные\par
\noindent\emph{Ukrajinsky / Ukrainian:} алгебраїчно незалежні / трансцендентні\par
\smallskip\noindent\emph{Razlog:} Nezavisne i transcendentne zachovujut sučasny matematičny registar. Internacionalny korenj transcend- byl smatrany dostatočno prepoznajemy, zato zapis ne tvori novu slovjansku zaměnu. Taka ocěna jest terminologičny sud, ne rezultat opyta s čitateljami.\par
\noindent\emph{Reason:} The expressions retain modern mathematical usage. The international transcend- root was considered recognizable enough that no new Slavic replacement was proposed. This is a terminological judgment, not a reader-study result.\par
\noindent{\footnotesize Izvorne zapisy / Original records: HIST-g097-006.}\par
\end{minipage}\par
\par\medskip\Needspace{7\baselineskip}
\noindent\begin{minipage}{\linewidth}
\hypertarget{TERM-0083}{}
\noindent\textbf{algebraično nezavisne funkcije}\quad{\footnotesize\texttt{TERM-0083}}\par
\noindent\emph{Iztočnik / Source:} algebraisch unabhängige Funktionen\par
\noindent\emph{Rusky / Russian:} алгебраически независимые функции\par
\noindent\emph{Ukrajinsky / Ukrainian:} алгебраїчно незалежні функції\par
\smallskip\noindent\emph{Razlog:} Algebraično nezavisne funkcije prodolžujut terminologiju ranga iz prvogo odděla děla 6. Razlog jest zachovanje istogo pojma v svazanyh čestjah dokaza, ne zavedenje novogo termina.\par
\noindent\emph{Reason:} The phrase continues the rank terminology of Paper 6, Section 1. The reason is to retain the same concept across connected parts of the proof, not to introduce a new term.\par
\noindent{\footnotesize Izvorne zapisy / Original records: HIST-g044-006.}\par
\end{minipage}\par
\par\medskip\Needspace{7\baselineskip}
\noindent\begin{minipage}{\linewidth}
\hypertarget{TERM-0084}{}
\noindent\textbf{algebraično zatvorjeno polje / algebraično razširiteljno polje}\quad{\footnotesize\texttt{TERM-0084}}\par
\noindent\emph{Iztočnik / Source:} algebraisch-abgeschlossener / algebraischer Erweiterungskörper\par
\noindent\emph{Rusky / Russian:} алгебраически замкнутое поле / алгебраическое расширение поля\par
\noindent\emph{Ukrajinsky / Ukrainian:} алгебраїчно замкнене поле / алгебраїчне розширення поля\par
\smallskip\noindent\emph{Razlog:} Obě nazvanja izrično govoret o polju: jedno o algebraičnoj zatvorjenosti, drugo o razširenju. Ta razlika jest potrěbna, ibo dokaz dodava proměnne i izdvaja linearne faktory. Zapis ne opravdava zamenjanje tųh dvuh svojstv jednym opčim pridavnikom.\par
\noindent\emph{Reason:} Both designations explicitly concern fields, distinguishing algebraic closedness from an extension. This matters because the proof adjoins variables and extracts linear factors. The record does not justify collapsing the two properties into one general adjective.\par
\noindent{\footnotesize Izvorne zapisy / Original records: HIST-g090-011, HIST-g092-065.}\par
\end{minipage}\par
\par\medskip\Needspace{7\baselineskip}
\noindent\begin{minipage}{\linewidth}
\hypertarget{TERM-0085}{}
\noindent\textbf{algebraično zatvorjeno polje}\quad{\footnotesize\texttt{TERM-0085}}\par
\noindent\emph{Iztočnik / Source:} algebraisch-abgeschlossener Körper\par
\noindent\emph{Rusky / Russian:} алгебраически замкнутое поле\par
\noindent\emph{Ukrajinsky / Ukrainian:} алгебраїчно замкнене поле\par
\smallskip\noindent\emph{Razlog:} Zatvorjeno jest kandidat v konstrukciji algebraično zatvorjeno polje. Izvorny razlog govori o poznatoj slovjanskoj algebraičnoj terminologiji, ale ne sravnjaje podrobno alternativu zamknjeno. To ograničenje trěba zachovati pri čitanju tutogo zapisa.\par
\noindent\emph{Reason:} Zatvorjeno is the candidate in the algebraically-closed-field construction. The original reason appeals to familiar Slavic algebraic terminology but does not compare zamknjeno in detail. That evidentiary limit remains relevant.\par
\noindent{\footnotesize Izvorne zapisy / Original records: HIST-g125-006.}\par
\smallskip\noindent\emph{Sravnjenje s tekstom:} Prěvod upotrěbjaje i algebraično zamknjeno polje, i algebraično zatvorjeno polje v pasažah o algebraičnoj zamknjenosti. Kontekst s razloženjem na linearne faktory podpiraje isti matematičny smysl, ale samo prisutnost oběh form ne izbira jednu iz njih kako normu. Historijne zapisy trěba zato čitati kako dokumentovane alternativy. Te tri pasaže ne dokazyvajut, že vsako pojavjenje zamknjen- ili zatvorjen- v korpusu ima isti smysl.\par
\noindent\emph{Textual comparison:} The translation uses both algebraično zamknjeno polje and algebraično zatvorjeno polje in passages concerning algebraic closedness. Context involving factorization into linear factors supports the same mathematical sense, but the coexistence of the forms does not select either as the norm. The historical records must therefore be read as documented alternatives. These three passages do not establish that every occurrence of zamknjen- or zatvorjen- in the corpus has the same sense.\par
\noindent{\footnotesize TEXT-COMP-02; pasaže / passages: 13803, 14930, 15719; \texttt{CORPUS\_FORM\_COMPARISONS.json}.}\par
\end{minipage}\par
\par\medskip\Needspace{7\baselineskip}
\noindent\begin{minipage}{\linewidth}
\hypertarget{TERM-0086}{}
\noindent\textbf{algebraično zamknjeno podkolco P}\quad{\footnotesize\texttt{TERM-0086}}\par
\noindent\emph{Iztočnik / Source:} algebraisch-abgeschlossener Unterring P\par
\noindent\emph{Rusky / Russian:} алгебраически замкнутое подкольцо P\par
\noindent\emph{Ukrajinsky / Ukrainian:} алгебраїчно замкнене підкільце P\par
\smallskip\noindent\emph{Razlog:} Zamknjeno bylo preferovano prěd zatvorjeno radi predpolaganoj jasnoće matematičnoj zatvorjenosti v vzhodnyh i zapadnyh jezykah. Tuto jest razlog konkretnogo historijnogo izbora. Podkolco jest takože historijna forma i ne směje bez prověrky stati aktualnym rekomendovanym slovom.\par
\noindent\emph{Reason:} Zamknjeno was preferred to zatvorjeno for its proposed fit with East and West Slavic mathematical closure terminology. This explains a historical choice. Podkolco is also historical and must not be promoted to a current recommendation without checking its subsequent status.\par
\noindent{\footnotesize Izvorne zapisy / Original records: HIST-g142-002.}\par
\end{minipage}\par
\par\medskip\Needspace{7\baselineskip}
\noindent\begin{minipage}{\linewidth}
\hypertarget{TERM-0087}{}
\noindent\textbf{algebraično zavisny / algebraično nezavisny}\quad{\footnotesize\texttt{TERM-0087}}\par
\noindent\emph{Iztočnik / Source:} algebraisch-abhängig / algebraisch-unabhängig\par
\noindent\emph{Rusky / Russian:} алгебраически зависимый / алгебраически независимый\par
\noindent\emph{Ukrajinsky / Ukrainian:} алгебраїчно залежний / алгебраїчно незалежний\par
\smallskip\noindent\emph{Razlog:} Para zavisny/nezavisny razlikuje elementy racionalnoj bazy od algebraičnoj bazy H. Razlog jest rola tųh dvuh baz v argumentu, a ne samo suprotna značenja dvuh pridavnikov.\par
\noindent\emph{Reason:} The pair distinguishes rational-basis elements from the algebraic basis H. The reason concerns the roles of those two bases in the argument, not merely the opposition between the adjectives.\par
\noindent{\footnotesize Izvorne zapisy / Original records: HIST-g070-007.}\par
\end{minipage}\par
\par\medskip\Needspace{7\baselineskip}
\noindent\begin{minipage}{\linewidth}
\hypertarget{TERM-0088}{}
\noindent\textbf{algebraično-cěle / algebraično-drobne}\quad{\footnotesize\texttt{TERM-0088}}\par
\noindent\emph{Iztočnik / Source:} algebraisch-ganz / algebraisch-gebrochen\par
\noindent\emph{Rusky / Russian:} алгебраически целые / алгебраически дробные\par
\noindent\emph{Ukrajinsky / Ukrainian:} алгебраїчно цілі / алгебраїчно дробові\par
\smallskip\noindent\emph{Razlog:} Cěly/drobny jest osnovna smyslova opozicija pasaže. Spojnik v algebraično-cěle i algebraično-drobne čini tehničnu pridavničnu konstrukciju vidimoju i zachovuje doslědnost s ranějšimi oddělami děla 9.\par
\noindent\emph{Reason:} Integral/fractional is the passage's central contrast. Hyphenation makes the technical compound adjectives visible and maintains consistency with earlier sections of Paper 9.\par
\noindent{\footnotesize Izvorne zapisy / Original records: HIST-g068-007.}\par
\end{minipage}\par
\par\medskip\Needspace{7\baselineskip}
\noindent\begin{minipage}{\linewidth}
\hypertarget{TERM-0089}{}
\noindent\textbf{algebraično-cěly / racionalno-cěly}\quad{\footnotesize\texttt{TERM-0089}}\par
\noindent\emph{Iztočnik / Source:} algebraisch-ganz / rational-ganz\par
\noindent\emph{Rusky / Russian:} алгебраически-целый / рационально-целый\par
\noindent\emph{Ukrajinsky / Ukrainian:} алгебраїчно-цілий / раціонально-цілий\par
\smallskip\noindent\emph{Razlog:} Algebraično-cěly i racionalno-cěly ostavajut odvojene, ibo iztočnik razlikuje dva pojma. Spojnik pokazuje, že cěly jest čestju tehničnogo termina, ne obyčnym stilističnym usiljenjem.\par
\noindent\emph{Reason:} The compounds remain distinct because the source distinguishes two concepts. Hyphenation indicates that cěly belongs to a technical designation rather than serving as an ordinary intensifier.\par
\noindent{\footnotesize Izvorne zapisy / Original records: HIST-g058-008.}\par
\end{minipage}\par
\par\medskip\Needspace{7\baselineskip}
\noindent\begin{minipage}{\linewidth}
\hypertarget{TERM-0090}{}
\noindent\textbf{algebraično-cělo zatvorjena}\quad{\footnotesize\texttt{TERM-0090}}\par
\noindent\emph{Iztočnik / Source:} algebraisch-ganz abgeschlossen\par
\noindent\emph{Rusky / Russian:} алгебраически цело замкнутая; интегрально замкнута\par
\noindent\emph{Ukrajinsky / Ukrainian:} алгебраїчно-ціло замкнена; інтегрально замкнена\par
\smallskip\noindent\emph{Razlog:} Algebraično-cělo ograničaje smysl zatvorjenosti i odliča go od topologičnoj zatvorjenosti. V dokazu protivrěčjem jest potrěbna imenno integralna zatvorjenost T. Ruske i ukrajinske zapisy dajut i doslovne i sučasne integralne vyraženja.\par
\noindent\emph{Reason:} Algebraično-cělo specifies integral closedness rather than topological closure. The contradiction argument requires precisely this property of T. The Russian and Ukrainian records include both more literal and modern integral terminology.\par
\noindent{\footnotesize Izvorne zapisy / Original records: HIST-g062-002, HIST-g067-008.}\par
\end{minipage}\par
\par\medskip\Needspace{7\baselineskip}
\noindent\begin{minipage}{\linewidth}
\hypertarget{TERM-0091}{}
\noindent\textbf{algebraično-cělo zavisi od [H,xi\_i]}\quad{\footnotesize\texttt{TERM-0091}}\par
\noindent\emph{Iztočnik / Source:} algebraisch-ganz abhängig von [H,xi\_i]\par
\noindent\emph{Rusky / Russian:} алгебраически-цело зависит от [H,xi\_i]\par
\noindent\emph{Ukrajinsky / Ukrainian:} алгебраїчно-ціло залежить від [H,xi\_i]\par
\smallskip\noindent\emph{Razlog:} Vyraženje zachovuje avtorsku leksiku algebraičnoj cělosti, ale označaje integralnu zavisimost nad pokazanoju oblastju [H,xi\_i]. Oznaka oblasti jest čestju odnosa i ne može byti izostavjena.\par
\noindent\emph{Reason:} The phrase retains the paper's algebraically-whole vocabulary while denoting integral dependence over the displayed domain [H,xi\_i]. The domain identifier is part of the relation and cannot be omitted.\par
\noindent{\footnotesize Izvorne zapisy / Original records: HIST-g065-014.}\par
\end{minipage}\par
\par\medskip\Needspace{7\baselineskip}
\noindent\begin{minipage}{\linewidth}
\hypertarget{TERM-0092}{}
\noindent\textbf{algebraično cěla vzhodno k G}\quad{\footnotesize\texttt{TERM-0092}}\par
\noindent\emph{Iztočnik / Source:} algebraisch-ganz in bezug auf G\par
\noindent\emph{Rusky / Russian:} алгебраически целая относительно G\par
\noindent\emph{Ukrajinsky / Ukrainian:} алгебраїчно ціла відносно G\par
\smallskip\noindent\emph{Razlog:} Forma cěla soglasuje se s ženskym imenikom veličina v opreděljenju V. Odnos k G označaje osnovu algebraičnoj cělosti. Izvorny razlog podpira tu gramatičnu i semantičnu strukturu, ne avtomatično vse upotrěby vzhodno k.\par
\noindent\emph{Reason:} Cěla agrees with the feminine noun veličina in Definition V. The relation to G specifies the base of algebraic integrality. The reason supports that grammatical and semantic structure, not automatically every use of vzhodno k.\par
\noindent{\footnotesize Izvorne zapisy / Original records: HIST-g049-009.}\par
\end{minipage}\par
\par\medskip\Needspace{7\baselineskip}
\noindent\begin{minipage}{\linewidth}
\hypertarget{TERM-0093}{}
\noindent\textbf{algebraično cěla i celočislena vzhodno k}\quad{\footnotesize\texttt{TERM-0093}}\par
\noindent\emph{Iztočnik / Source:} algebraisch-ganz und ganzzahlig in Bezug auf\par
\noindent\emph{Rusky / Russian:} алгебраически целая и целочисленная относительно\par
\noindent\emph{Ukrajinsky / Ukrainian:} алгебраїчно ціла й цілочисельна відносно\par
\smallskip\noindent\emph{Razlog:} Dva uslovja ostavajut izrično odvojene: integralna zavisimost i uslovje cělosti v prstenu koeficientov. Jedna opča zaměna slovom cěly by skryla razliku, na ktorej počivaje formulacija.\par
\noindent\emph{Reason:} Two conditions are kept explicit: integral dependence and the integrality condition in the coefficient ring. Replacing both by an undifferentiated integral would conceal the distinction on which the formulation depends.\par
\noindent{\footnotesize Izvorne zapisy / Original records: HIST-g055-002.}\par
\end{minipage}\par
\par\medskip\Needspace{7\baselineskip}
\noindent\begin{minipage}{\linewidth}
\hypertarget{TERM-0094}{}
\noindent\textbf{byti algebraično cěly nad \textbackslash{}mathfrak F}\quad{\footnotesize\texttt{TERM-0094}}\par
\noindent\emph{Iztočnik / Source:} algebraisch-ganz von \textbackslash{}mathfrak F abhängen\par
\noindent\emph{Rusky / Russian:} быть алгебраически целым над \textbackslash{}mathfrak F\par
\noindent\emph{Ukrajinsky / Ukrainian:} бути алгебраїчно цілим над \textbackslash{}mathfrak F\par
\smallskip\noindent\emph{Razlog:} Nad \textbackslash{}mathfrak F označaje osnovu integralnosti i prodolžuje čitanje iz oddělov 10–11 děla 6. Objasnjajuča věta zachovuje avtorsky kriterij prěz jednačenje; sučasny termin go ne zaměnjaje.\par
\noindent\emph{Reason:} Nad \textbackslash{}mathfrak F names the base of integrality and continues the interpretation in Paper 6, Sections 10–11. The explanatory clause retains Noether's equation criterion; the modern term does not replace it.\par
\noindent{\footnotesize Izvorne zapisy / Original records: HIST-g060-006.}\par
\end{minipage}\par
\par\medskip\Needspace{7\baselineskip}
\noindent\begin{minipage}{\linewidth}
\hypertarget{TERM-0095}{}
\noindent\textbf{jest algebraično cěla nad T}\quad{\footnotesize\texttt{TERM-0095}}\par
\noindent\emph{Iztočnik / Source:} algebraisch-ganz von T abhängen\par
\noindent\emph{Rusky / Russian:} алгебраически цела над T\par
\noindent\emph{Ukrajinsky / Ukrainian:} є алгебраїчно цілою над T\par
\smallskip\noindent\emph{Razlog:} Nad T dava sučasno algebraično vyraženje odnosa integralnosti. Avtorsko vyraženje zavisimosti ostavaje objasnjeno v věte, tako da terminologično osuvrěmenjenje ne izbriše obsah iztočnogo uslovja.\par
\noindent\emph{Reason:} Nad T gives the modern algebraic expression of integrality. The explanatory sentence preserves the author's dependence formulation, so terminological modernization does not erase the source condition.\par
\noindent{\footnotesize Izvorne zapisy / Original records: HIST-g062-001.}\par
\end{minipage}\par
\par\medskip\Needspace{7\baselineskip}
\noindent\begin{minipage}{\linewidth}
\hypertarget{TERM-0096}{}
\noindent\textbf{algebraično-cěla zatvorenost}\quad{\footnotesize\texttt{TERM-0096}}\par
\noindent\emph{Iztočnik / Source:} algebraisch-ganze Abgeschlossenheit\par
\noindent\emph{Rusky / Russian:} алгебраически-целая замкнутость\par
\noindent\emph{Ukrajinsky / Ukrainian:} алгебраїчно-ціла замкненість\par
\smallskip\noindent\emph{Razlog:} Zatvorenost jest imenik za svojstvo zatvorjenosti; algebraično-cěla izrično ograničaje go na integralnost. Razlog jest zachovati tehničny modifikator, ne predstaviti opču zatvorjenost bez opreděljenogo smysla.\par
\noindent\emph{Reason:} Zatvorenost names the closure property, while algebraično-cěla restricts it to integrality. The reason is to preserve that technical modifier rather than suggest an unspecified form of closure.\par
\noindent{\footnotesize Izvorne zapisy / Original records: HIST-g058-010.}\par
\end{minipage}\par
\par\medskip\Needspace{7\baselineskip}
\noindent\begin{minipage}{\linewidth}
\hypertarget{TERM-0097}{}
\noindent\textbf{algebraično-cěla adjunkcija}\quad{\footnotesize\texttt{TERM-0097}}\par
\noindent\emph{Iztočnik / Source:} algebraisch-ganze Adjunktion\par
\noindent\emph{Rusky / Russian:} алгебраически-целое присоединение\par
\noindent\emph{Ukrajinsky / Ukrainian:} алгебраїчно-ціле приєднання\par
\smallskip\noindent\emph{Razlog:} Algebraično-cěla adjunkcija jest obrazovana paralelno k racionalno-cěloj adjunkciji. Spojnik čini pridavnik čitljivym kako jednu tehničnu kvalifikaciju, zachovajuči razliku medžu dvuma konstrukcijami.\par
\noindent\emph{Reason:} The expression parallels rationally integral adjunction. The hyphenated adjective is intended to read as one technical qualification while preserving the distinction between the two constructions.\par
\noindent{\footnotesize Izvorne zapisy / Original records: HIST-g062-008.}\par
\end{minipage}\par
\par\medskip\Needspace{7\baselineskip}
\noindent\begin{minipage}{\linewidth}
\hypertarget{TERM-0098}{}
\noindent\textbf{algebraično cělo čislo}\quad{\footnotesize\texttt{TERM-0098}}\par
\noindent\emph{Iztočnik / Source:} algebraisch-ganze Zahl\par
\noindent\emph{Rusky / Russian:} алгебраически целое число\par
\noindent\emph{Ukrajinsky / Ukrainian:} алгебраїчно ціле число\par
\smallskip\noindent\emph{Razlog:} Algebraično cělo čislo imenuje algebraičny integer i ostavaje odvojeno od uslovja, že funkcija ima cěločislene koeficienty. Zapis podpira to razděljenje pomoču ruskogo i ukrajinskogo odbornogo izraza.\par
\noindent\emph{Reason:} The expression names an algebraic integer and remains distinct from the condition that a function have integer coefficients. Russian and Ukrainian technical expressions support this distinction in the record.\par
\noindent{\footnotesize Izvorne zapisy / Original records: HIST-g061-006.}\par
\end{minipage}\par
\par\medskip\Needspace{7\baselineskip}
\noindent\begin{minipage}{\linewidth}
\hypertarget{TERM-0099}{}
\noindent\textbf{algebraično cěle čisla iz Omega / cěle čisla}\quad{\footnotesize\texttt{TERM-0099}}\par
\noindent\emph{Iztočnik / Source:} algebraisch-ganze Zahlen aus Omega / ganze Zahlen\par
\noindent\emph{Rusky / Russian:} алгебраически целые числа из Omega / целые числа\par
\noindent\emph{Ukrajinsky / Ukrainian:} алгебраїчно цілі числа з Omega / цілі числа\par
\smallskip\noindent\emph{Razlog:} Prěvod zachovuje avtorsku lokalnu konvenciju: algebraično cěle čisla iz Omega mogut v tom oddělu byti nazvane kratko cěle čisla. Kratko nazvanje ne ustanavjaje totožnost vseh različnyh pojmov cělosti izvan tutogo konteksta.\par
\noindent\emph{Reason:} The translation preserves the author's local convention: algebraic integers in Omega may be called simply integers within the section. This abbreviation does not identify every different notion of integrality outside that context.\par
\noindent{\footnotesize Izvorne zapisy / Original records: HIST-g054-002.}\par
\end{minipage}\par
\par\medskip\Needspace{7\baselineskip}
\noindent\begin{minipage}{\linewidth}
\hypertarget{TERM-0100}{}
\noindent\textbf{algebraično-celočislene polinomy}\quad{\footnotesize\texttt{TERM-0100}}\par
\noindent\emph{Iztočnik / Source:} algebraisch-ganzzahlige Polynome\par
\noindent\emph{Rusky / Russian:} алгебраически-целочисленные полиномы\par
\noindent\emph{Ukrajinsky / Ukrainian:} алгебраїчно-цілочисельні поліноми\par
\smallskip\noindent\emph{Razlog:} Tehničny pridavnik se odnosi na integralnost polinomialnyh koeficientov vzhledem na algebraično cěle čisla. Zapis izrično odliča to od obyčnoj racionalnoj cělosti; prěvod ne směje ztratiti osnovu tųh koeficientov.\par
\noindent\emph{Reason:} The technical adjective concerns polynomial coefficient integrality over algebraic integers. The record explicitly distinguishes this from ordinary rational integrality; the coefficient setting must remain identifiable.\par
\noindent{\footnotesize Izvorne zapisy / Original records: HIST-g055-006.}\par
\end{minipage}\par
\par\medskip\Needspace{7\baselineskip}
\noindent\begin{minipage}{\linewidth}
\hypertarget{TERM-0101}{}
\noindent\textbf{algebraično-cěle substitucije / cěle algebraične čisla}\quad{\footnotesize\texttt{TERM-0101}}\par
\noindent\emph{Iztočnik / Source:} algebraisch-ganzzahlige Substitutionen / algebraische ganze Zahlen\par
\noindent\emph{Rusky / Russian:} алгебраически-целые подстановки / алгебраические целые числа\par
\noindent\emph{Ukrajinsky / Ukrainian:} алгебраїчно-цілі підстановки / алгебраїчно цілі числа\par
\smallskip\noindent\emph{Razlog:} Substitucija jest kratko algebraično nazvanje transformacije. V oddělu o konečnyh grupah čisla v koeficientnyh poljah sut algebraično cěle; pridavnik cěle zato ima konkretny odborny smysl, ne opče značenje polnost.\par
\noindent\emph{Reason:} Substitucija is retained as a compact algebraic transformation term. The finite-group section concerns algebraic integers in coefficient fields, so integral has a specific technical meaning rather than general completeness.\par
\noindent{\footnotesize Izvorne zapisy / Original records: HIST-g077-012.}\par
\end{minipage}\par
\par\medskip\Needspace{7\baselineskip}
\noindent\begin{minipage}{\linewidth}
\hypertarget{TERM-0102}{}
\noindent\textbf{algebrai?no-drobno ?islo}\quad{\footnotesize\texttt{TERM-0102}}\par
\noindent\emph{Iztočnik / Source:} algebraisch-gebrochene Zahl\par
\noindent\emph{Rusky / Russian:} ?????????????-??????? ?????\par
\noindent\emph{Ukrajinsky / Ukrainian:} ???????????-??????? ?????\par
\smallskip\noindent\emph{Razlog:} Zachovany razlog govori, že drobny status jest nužno dopolnitelno uslovje i ne směje byti oslaben do neodređenoj necělosti. Ale medžuslovjanska forma i obe sravniteljne formy sodržet znaky ?. Poškođeny zapis ne jest pouzdan obrazec pravopisa.\par
\noindent\emph{Reason:} The preserved reason treats fractional status as an essential side condition, not an unspecified lack of integrality. The Interslavic candidate and both comparanda contain question marks replacing characters. They are not reliable spelling witnesses.\par
\noindent{\footnotesize Izvorne zapisy / Original records: HIST-g064-007.}\par
\end{minipage}\par
\par\medskip\Needspace{7\baselineskip}
\noindent\begin{minipage}{\linewidth}
\hypertarget{TERM-0103}{}
\noindent\textbf{algebraično-drobno čislo}\quad{\footnotesize\texttt{TERM-0103}}\par
\noindent\emph{Iztočnik / Source:} algebraisch-gebrochene Zahl\par
\noindent\emph{Rusky / Russian:} алгебраически-дробное число\par
\noindent\emph{Ukrajinsky / Ukrainian:} алгебраїчно-дробове число\par
\smallskip\noindent\emph{Razlog:} Iztočno uslovje izključaje imenno algebraično-drobne čisla. Prěvod zato zachovuje tu kvalifikaciju, vměsto opčej negacije integralnosti bez pojmovogo konteksta.\par
\noindent\emph{Reason:} The condition excludes specifically algebraically fractional numbers. The wording retains this qualification rather than giving a context-free negation of integrality.\par
\noindent{\footnotesize Izvorne zapisy / Original records: HIST-g065-006.}\par
\end{minipage}\par
\par\medskip\Needspace{7\baselineskip}
\noindent\begin{minipage}{\linewidth}
\hypertarget{TERM-0104}{}
\noindent\textbf{algebrai?na baza}\quad{\footnotesize\texttt{TERM-0104}}\par
\noindent\emph{Iztočnik / Source:} algebraische Basis\par
\noindent\emph{Rusky / Russian:} [Bukvy ne sut zachovane / Lettering not recoverable]\par
\noindent\emph{Ukrajinsky / Ukrainian:} [Bukvy ne sut zachovane / Lettering not recoverable]\par
\smallskip\noindent\emph{Razlog:} Baza H jest fiksovana osnova, vzhledem na ktoru se izslěduje familija M(G). Razlog takože trěbuje doslědnu vezu s ranějšimi svojstvami bazy. Grafija algebrai?na i sravniteljne formy sut poškođene; točny znak prěd 5 v razlogu ne jest zachovany.\par
\noindent\emph{Reason:} H is the fixed basis relative to which the family M(G) is studied, with continuity to earlier basis-property terminology. The candidate and comparanda are damaged, and the sign before 5 in the rationale has not survived. These losses must not be silently reconstructed.\par
\noindent{\footnotesize Izvorne zapisy / Original records: HIST-g063-001.}\par
\end{minipage}\par
\par\medskip\Needspace{7\baselineskip}
\noindent\begin{minipage}{\linewidth}
\hypertarget{TERM-0105}{}
\noindent\textbf{algebraična baza}\quad{\footnotesize\texttt{TERM-0105}}\par
\noindent\emph{Iztočnik / Source:} algebraische Basis\par
\noindent\emph{Rusky / Russian:} алгебраическая база\par
\noindent\emph{Ukrajinsky / Ukrainian:} алгебраїчна база\par
\smallskip\noindent\emph{Razlog:} Algebraična baza zachovuje imenovanje iz děla 9. Odděl potrěbuje zaměnu jednoj bazy drugoju i zachovanje svojstva bazy pri izomorfizmu; upotrěba istogo termina čini tu vezu jasnoju.\par
\noindent\emph{Reason:} The phrase continues Paper 9's terminology. The section replaces one algebraic basis by another and uses preservation of basis status under isomorphism; consistent naming makes that connection explicit.\par
\noindent{\footnotesize Izvorne zapisy / Original records: HIST-g066-007.}\par
\end{minipage}\par
\par\medskip\Needspace{7\baselineskip}
\noindent\begin{minipage}{\linewidth}
\hypertarget{TERM-0106}{}
\noindent\textbf{algebraična baza / racionalna baza}\quad{\footnotesize\texttt{TERM-0106}}\par
\noindent\emph{Iztočnik / Source:} algebraische Basis / rationale Basis\par
\noindent\emph{Rusky / Russian:} алгебраическая база / рациональная база\par
\noindent\emph{Ukrajinsky / Ukrainian:} алгебраїчна база / раціональна база\par
\smallskip\noindent\emph{Razlog:} Baza jest preferovana prěd osnova, ibo tekst opreděljaje osobne generujuče sistemy čisel. Pridavniki algebraična i racionalna razlikujut dva nazvane objekty. Razlog jest točnost definovanogo pojma, ne opča zabrana slova osnova.\par
\noindent\emph{Reason:} Baza is preferred to osnova because the text defines particular generating systems of numbers. Algebraična and racionalna distinguish the named objects. This is a reason concerning a defined concept, not a general prohibition on osnova.\par
\noindent{\footnotesize Izvorne zapisy / Original records: HIST-g058-004.}\par
\end{minipage}\par
\par\medskip\Needspace{7\baselineskip}
\noindent\begin{minipage}{\linewidth}
\hypertarget{TERM-0107}{}
\noindent\textbf{algebraična baza H iz Theta}\quad{\footnotesize\texttt{TERM-0107}}\par
\noindent\emph{Iztočnik / Source:} algebraische Basis H aus Theta\par
\noindent\emph{Rusky / Russian:} алгебраическая база H из Theta\par
\noindent\emph{Ukrajinsky / Ukrainian:} алгебраїчна база H з Theta\par
\smallskip\noindent\emph{Razlog:} H jest izbrana iz veličin racionalnoj bazy Theta. Izraz zachovuje taj odnos, ibo dokaz prěvodi racionalnu zavisimost na konečno mnogo primitivnyh elementov. H i Theta ostavajut točne oznaky različnyh objektov.\par
\noindent\emph{Reason:} H is chosen from the quantities of the rational basis Theta. The wording preserves that relation because the proof reduces rational dependence to finitely many primitive elements. H and Theta remain exact identifiers of different objects.\par
\noindent{\footnotesize Izvorne zapisy / Original records: HIST-g065-010.}\par
\end{minipage}\par
\par\medskip\Needspace{7\baselineskip}
\noindent\begin{minipage}{\linewidth}
\hypertarget{TERM-0108}{}
\noindent\textbf{algebraično razširenje}\quad{\footnotesize\texttt{TERM-0108}}\par
\noindent\emph{Iztočnik / Source:} algebraische Erweiterung\par
\noindent\emph{Rusky / Russian:} алгебраическое расширение\par
\noindent\emph{Ukrajinsky / Ukrainian:} алгебраїчне розширення\par
\smallskip\noindent\emph{Razlog:} Razširenje jest odborny izraz za razširenje osnovnogo polja v oddělu 6. Konečnost i algebraičnost morajut byti čitane kako osobne kvalifikacije iztočnogo argumenta, ne izmysljene iz samogo imenika.\par
\noindent\emph{Reason:} Razširenje denotes a base-field extension in Section 6. Finiteness and algebraicity must be read as the argument's particular qualifications, not inferred from the noun extension alone.\par
\noindent{\footnotesize Izvorne zapisy / Original records: HIST-g046-008.}\par
\end{minipage}\par
\par\medskip\Needspace{7\baselineskip}
\noindent\begin{minipage}{\linewidth}
\hypertarget{TERM-0109}{}
\noindent\textbf{algebraične elementy / cěle elementy}\quad{\footnotesize\texttt{TERM-0109}}\par
\noindent\emph{Iztočnik / Source:} algebraische Größen / ganze Größen\par
\noindent\emph{Rusky / Russian:} алгебраические элементы / целые элементы\par
\noindent\emph{Ukrajinsky / Ukrainian:} алгебраїчні елементи / цілі елементи\par
\smallskip\noindent\emph{Razlog:} Elementy jest izbrano vměsto veličiny, ibo dělo 14 razmatraje polja i idealy, ne skalarny račun. Cěle zachovuje integralnost. Tuto jest kontekstno osuvrěmenjenje historijnogo Größe, ne pravilo za vsako pojavjenje německogo slova.\par
\noindent\emph{Reason:} Elementy is preferred to veličiny because Paper 14 concerns fields and ideals rather than scalar calculus. Cěle preserves integrality. This is a contextual modernization of historical Größe, not a rule for every occurrence of that German word.\par
\noindent{\footnotesize Izvorne zapisy / Original records: HIST-g076-003.}\par
\end{minipage}\par
\par\medskip\Needspace{7\baselineskip}
\noindent\begin{minipage}{\linewidth}
\hypertarget{TERM-0110}{}
\noindent\textbf{algebraične i diferencialne invarianty}\quad{\footnotesize\texttt{TERM-0110}}\par
\noindent\emph{Iztočnik / Source:} Algebraische und Differentialinvarianten\par
\noindent\emph{Rusky / Russian:} алгебраические и дифференциальные инварианты\par
\noindent\emph{Ukrajinsky / Ukrainian:} алгебраїчні та диференціальні інваріанти\par
\smallskip\noindent\emph{Razlog:} Internacionalny korenj invariant- i pridavniki algebraične/diferencialne zachovujut poznatu odbornu terminologiju. Zapis ne tvori slovjansku zaměnu samo radi jezyčnoj čistoty; para pridavnikov ostavjaje obe vrsty invariantov vidime.\par
\noindent\emph{Reason:} The international invariant- root and the algebraic/differential adjective pair retain familiar technical terminology. The record proposes no Slavic replacement merely for linguistic purity and keeps both classes of invariants explicit.\par
\noindent{\footnotesize Izvorne zapisy / Original records: HIST-g094-001.}\par
\end{minipage}\par
\par\medskip\Needspace{7\baselineskip}
\noindent\begin{minipage}{\linewidth}
\hypertarget{TERM-0111}{}
\noindent\textbf{algebraične čislove polja}\quad{\footnotesize\texttt{TERM-0111}}\par
\noindent\emph{Iztočnik / Source:} algebraische Zahlkörper\par
\noindent\emph{Rusky / Russian:} алгебраические числовые поля\par
\noindent\emph{Ukrajinsky / Ukrainian:} алгебраїчні числові поля\par
\smallskip\noindent\emph{Razlog:} Polja jest izbrano vměsto těla, ibo tute čislove polja sut komutativne. Pridavniki algebraične i čislove zachovujut obě kvalifikacije. Razlog se odnosi na konkretny algebraičny objekt.\par
\noindent\emph{Reason:} Polja is preferred to těla because these number fields are commutative. Both algebraic and number qualifications remain explicit. The reason concerns the particular algebraic object.\par
\noindent{\footnotesize Izvorne zapisy / Original records: HIST-g174-003.}\par
\end{minipage}\par
\par\medskip\Needspace{7\baselineskip}
\noindent\begin{minipage}{\linewidth}
\hypertarget{TERM-0112}{}
\noindent\textbf{algebraično razširjenje polja / podpolje}\quad{\footnotesize\texttt{TERM-0112}}\par
\noindent\emph{Iztočnik / Source:} algebraischer Erweiterungskörper / Teilkörper\par
\noindent\emph{Rusky / Russian:} алгебраическое расширение поля / подполе\par
\noindent\emph{Ukrajinsky / Ukrainian:} алгебраїчне розширення поля / підполе\par
\smallskip\noindent\emph{Razlog:} Razširjenje polja izjasnjaje konstrukciju i izběgaje směšenje konečnogo polja s konečnym razširenjem polja. Podpolje odvojeno imenuje podobjekt. Medžuslovjansky izraz slijedi tu istu strukturu kak rusky i ukrajinsky prěvod.\par
\noindent\emph{Reason:} Extension of a field makes the construction explicit and avoids confusing a finite field with a finite extension. Podpolje separately names the subfield. The Interslavic expression follows the same structure as the Russian and Ukrainian comparanda.\par
\noindent{\footnotesize Izvorne zapisy / Original records: HIST-g097-007.}\par
\end{minipage}\par
\par\medskip\Needspace{7\baselineskip}
\noindent\begin{minipage}{\linewidth}
\hypertarget{TERM-0113}{}
\noindent\textbf{algebraičny rang}\quad{\footnotesize\texttt{TERM-0113}}\par
\noindent\emph{Iztočnik / Source:} algebraischer Rang\par
\noindent\emph{Rusky / Russian:} алгебраический ранг\par
\noindent\emph{Ukrajinsky / Ukrainian:} алгебраїчний ранг\par
\smallskip\noindent\emph{Razlog:} Rang jest karakteristično čislo iz opreděljenja II. Tri zapisy podpirajut doslědnost termina v dělu 6, osoblivo v oddělah 1 i 4. Zapis ne prědlaga zaměnu togo imenovanja novym opisnym slovom.\par
\noindent\emph{Reason:} Rank is the characteristic number introduced by Definition II. Three records support consistent use across Paper 6, particularly Sections 1 and 4. They propose no new descriptive replacement for the designation.\par
\noindent{\footnotesize Izvorne zapisy / Original records: HIST-g041-010, HIST-g044-007, HIST-g047-007.}\par
\end{minipage}\par
\par\medskip\Needspace{7\baselineskip}
\noindent\begin{minipage}{\linewidth}
\hypertarget{TERM-0114}{}
\noindent\textbf{algebraičny rang (stepen transcendencije)}\quad{\footnotesize\texttt{TERM-0114}}\par
\noindent\emph{Iztočnik / Source:} algebraischer Rang (Transzendenzgrad)\par
\noindent\emph{Rusky / Russian:} алгебраический ранг (степень трансцендентности)\par
\noindent\emph{Ukrajinsky / Ukrainian:} алгебраїчний ранг (степінь трансцендентності)\par
\smallskip\noindent\emph{Razlog:} Zaglavje zachovuje avtorsko objasnjenje v skobkah: algebraičny rang odpovědaje stepenju transcendencije. Internacionalna familija transcendencija jest izbrana vměsto novogo domorodnogo abstraktnogo imenika.\par
\noindent\emph{Reason:} The wording retains the author's parenthetical equivalence between algebraic rank and transcendence degree. It uses the international transcendencija family rather than coining a new native abstract noun.\par
\noindent{\footnotesize Izvorne zapisy / Original records: HIST-g097-004.}\par
\end{minipage}\par
\par\medskip\Needspace{7\baselineskip}
\noindent\begin{minipage}{\linewidth}
\hypertarget{TERM-0115}{}
\noindent\textbf{algebraičny rang / polje ostatkovyh klasov}\quad{\footnotesize\texttt{TERM-0115}}\par
\noindent\emph{Iztočnik / Source:} algebraischer Rang / Restklassenkörper\par
\noindent\emph{Rusky / Russian:} алгебраический ранг / поле классов вычетов\par
\noindent\emph{Ukrajinsky / Ukrainian:} алгебраїчний ранг / поле класів лишків\par
\smallskip\noindent\emph{Razlog:} Polje ostatkovyh klasov prodolžuje terminologiju odděla 3. Svaz s algebraičnym rangom jest potrěbny za aritmetično opreděljenje dimenzije v teoremu V. Oba termina ostavajut razpoznajeme v toj konstrukciji.\par
\noindent\emph{Reason:} The residue-class-field expression continues Section 3's terminology. Its connection with algebraic rank is required for Theorem V's arithmetic definition of dimension. Both terms remain identifiable within that construction.\par
\noindent{\footnotesize Izvorne zapisy / Original records: HIST-g098-003.}\par
\end{minipage}\par
\par\medskip\Needspace{7\baselineskip}
\noindent\begin{minipage}{\linewidth}
\hypertarget{TERM-0116}{}
\noindent\textbf{algebraičny rang ostavaje sohranjen}\quad{\footnotesize\texttt{TERM-0116}}\par
\noindent\emph{Iztočnik / Source:} algebraischer Rang bleibt erhalten\par
\noindent\emph{Rusky / Russian:} алгебраический ранг сохраняется\par
\noindent\emph{Ukrajinsky / Ukrainian:} алгебраїчний ранг зберігається\par
\smallskip\noindent\emph{Razlog:} Ostavaje sohranjen izrično govori, že dozvoljene substitucije ne menjajut algebraičny rang. Izvorny razlog podpira doslědnost s drugymi vyraženjami zachovanja, ale ne dava osobnu jezyčnu argumentaciju za sohranjen nasuprot drugym sinonimam.\par
\noindent\emph{Reason:} The expression states explicitly that the permitted substitutions preserve algebraic rank. The reason invokes consistency with other preservation wording but provides no separate linguistic argument for sohranjen over its alternatives.\par
\noindent{\footnotesize Izvorne zapisy / Original records: HIST-g042-008.}\par
\end{minipage}\par
\par\medskip\Needspace{7\baselineskip}
\noindent\begin{minipage}{\linewidth}
\hypertarget{TERM-0117}{}
\noindent\textbf{algebraično čislovo polje}\quad{\footnotesize\texttt{TERM-0117}}\par
\noindent\emph{Iztočnik / Source:} algebraischer Zahlkörper\par
\noindent\emph{Rusky / Russian:} алгебраическое числовое поле\par
\noindent\emph{Ukrajinsky / Ukrainian:} алгебраїчне числове поле\par
\smallskip\noindent\emph{Razlog:} Čislovo polje zachovuje razliku medžu opčim poljem i algebraičnym čislovym poljem v zaglavji i tekstu. Četyre zapisy se opirajut na sučasny vzhodnoslovjansky registar i doslědnost s dělom 14 i uvodom děla 24; Körper ne jest doslovno prěvedeno kak tělo.\par
\noindent\emph{Reason:} The phrase preserves the distinction between a field generally and an algebraic number field in both title and body. Four records invoke modern East Slavic usage and continuity with Paper 14 and Paper 24's introduction. Körper is not translated literally as body.\par
\noindent{\footnotesize Izvorne zapisy / Original records: HIST-g077-013, HIST-g098-012, HIST-g101-005, HIST-g103-002.}\par
\end{minipage}\par
\par\medskip\Needspace{7\baselineskip}
\noindent\begin{minipage}{\linewidth}
\hypertarget{TERM-0118}{}
\noindent\textbf{vse parne cěle racionalne čisla}\quad{\footnotesize\texttt{TERM-0118}}\par
\noindent\emph{Iztočnik / Source:} alle geraden, ganzen rationalen Zahlen\par
\noindent\emph{Rusky / Russian:} все четные целые рациональные числа\par
\noindent\emph{Ukrajinsky / Ukrainian:} усі парні цілі раціональні числа\par
\smallskip\noindent\emph{Razlog:} Parne, cěle i racionalne sut vse potrebne kvalifikacije. Priměr odděla 11 razmatraje oblast bez jedinice; n tam označaje kratky zapis konečnyh sum, ne nužno element prstena. Prěvod ne směje dodati edinicu ili izostaviti parnost.\par
\noindent\emph{Reason:} Even, integral, and rational are all substantive qualifications. Section 11 uses a domain without identity; n abbreviates finite sums rather than necessarily denoting a ring element. The translation must neither supply an identity nor omit parity.\par
\noindent{\footnotesize Izvorne zapisy / Original records: HIST-g081-088.}\par
\end{minipage}\par
\par\medskip\Needspace{7\baselineskip}
\noindent\begin{minipage}{\linewidth}
\hypertarget{TERM-0119}{}
\noindent\textbf{vse sistemy iz n elementov iz T}\quad{\footnotesize\texttt{TERM-0119}}\par
\noindent\emph{Iztočnik / Source:} alle Systeme von je n Elementen aus T\par
\noindent\emph{Rusky / Russian:} все системы из n элементов из T\par
\noindent\emph{Ukrajinsky / Ukrainian:} усі системи з n елементів із T\par
\smallskip\noindent\emph{Razlog:} Vse zachovuje kvantifikaciju nad vsemi sistemami, a iz n elementov razkriva německo je n. Izraz govori o n elementah v každom sistemu, ne o jednom zbirnom čislu za vse sistemy.\par
\noindent\emph{Reason:} Vse preserves quantification over all systems; iz n elementov makes German je n explicit. Each system has n elements; n is not a combined count across all systems.\par
\noindent{\footnotesize Izvorne zapisy / Original records: HIST-g158-002.}\par
\end{minipage}\par
\par\medskip\Needspace{7\baselineskip}
\noindent\begin{minipage}{\linewidth}
\hypertarget{TERM-0120}{}
\noindent\textbf{vsi ostali idealy sut glavne idealy}\quad{\footnotesize\texttt{TERM-0120}}\par
\noindent\emph{Iztočnik / Source:} alle übrigen Ideale Hauptideale\par
\noindent\emph{Rusky / Russian:} все остальные идеалы являются главными\par
\noindent\emph{Ukrajinsky / Ukrainian:} усі інші ідеали є головними\par
\smallskip\noindent\emph{Razlog:} Zaključok ostavaje větoju o vseh ostalyh idealah, ktore sut glavne. Zapis izběgaje skračenje do ponovnogo nazvanja prstena glavnyh idealov, ibo konkretna formulacija izjasnjaje krok dokaza.\par
\noindent\emph{Reason:} The conclusion remains a statement that all the remaining ideals are principal. It avoids compressing the claim into another mention of a principal ideal ring, because the explicit formulation makes the proof step clearer.\par
\noindent{\footnotesize Izvorne zapisy / Original records: HIST-g156-007.}\par
\end{minipage}\par
\par\medskip\Needspace{7\baselineskip}
\noindent\begin{minipage}{\linewidth}
\hypertarget{TERM-0121}{}
\noindent\textbf{obča teorija idealov}\quad{\footnotesize\texttt{TERM-0121}}\par
\noindent\emph{Iztočnik / Source:} allgemeine Idealtheorie\par
\noindent\emph{Rusky / Russian:} общая теория идеалов\par
\noindent\emph{Ukrajinsky / Ukrainian:} загальна теорія ідеалів\par
\smallskip\noindent\emph{Razlog:} Obča teorija idealov zachovuje ime matematičnoj oblasti i veže formulaciju s terminologijeju idealov v prědhodnyh dvuh rabotah. Obča označaje obči opseg teorije; ne dodavaje tvrđenje, že vsaky možny slučaj jest rěšeny. Zapis podpira doslědnost imenovanja v korpusu, ale ne prědstavjaje samostojno izslědovanje razumělivosti termina.\par
\noindent\emph{Reason:} Obča teorija idealov names the mathematical subject and continues the ideal-theoretic terminology of the two preceding papers. Obča marks the general scope of the theory, not a claim that every possible case has been solved. The recorded reason supports consistency across the corpus; it is not an independent study of the term's intelligibility.\par
\noindent{\footnotesize Izvorne zapisy / Original records: HIST-g095-002.}\par
\end{minipage}\par
\par\medskip\Needspace{7\baselineskip}
\noindent\begin{minipage}{\linewidth}
\hypertarget{TERM-0122}{}
\noindent\textbf{obča linearna grupa}\quad{\footnotesize\texttt{TERM-0122}}\par
\noindent\emph{Iztočnik / Source:} allgemeine lineare Gruppe\par
\noindent\emph{Rusky / Russian:} общая линейная группа\par
\noindent\emph{Ukrajinsky / Ukrainian:} загальна лінійна група\par
\smallskip\noindent\emph{Razlog:} Obča linearna grupa zachovuje ustalenu strukturu matematičnogo nazvanja: obča odliča tu grupu od osobnyh linearnih grup. Izbor obča slěduje prěvodnomu odgovaranju za allgemein v tom korpusu. Podobnost s ruskoju i ukrajinskoju konstrukcijeju jest sravniteljnym razlogom, ne dokazom, že medžuslovjansky termin byl nezavisno potvrđeny čitateljami.\par
\noindent\emph{Reason:} Obča linearna grupa preserves the structure of the mathematical name: obča distinguishes the general group from particular linear groups. The choice continues the corpus's treatment of allgemein. The parallel Russian and Ukrainian constructions provide comparative support, not evidence that readers independently validated the Interslavic term.\par
\noindent{\footnotesize Izvorne zapisy / Original records: HIST-g094-002.}\par
\end{minipage}\par
\par\medskip\Needspace{7\baselineskip}
\noindent\begin{minipage}{\linewidth}
\hypertarget{TERM-0123}{}
\noindent\textbf{obči teorem o diskriminantu}\quad{\footnotesize\texttt{TERM-0123}}\par
\noindent\emph{Iztočnik / Source:} Allgemeiner Diskriminantensatz\par
\noindent\emph{Rusky / Russian:} общая теорема о дискриминанте\par
\noindent\emph{Ukrajinsky / Ukrainian:} загальна теорема про дискримінант\par
\smallskip\noindent\emph{Razlog:} Obči teorem o diskriminantu prodolžaje imenovanje teorema iz odděla 7. Pridavnik obči jest sučstveny: nazvanje ne omejaje rezultat na lokalny slučaj prstena glavnyh idealov. Izbor zato služi razlikovanju opsega tvrđenja, ne samo doslědnomu prěvodu nazvanja.\par
\noindent\emph{Reason:} Obči teorem o diskriminantu continues the theorem terminology of section 7. Obči matters because the title is not restricted to the local principal-ideal-ring case. The choice distinguishes the scope of the assertion as well as maintaining a consistent theorem name.\par
\noindent{\footnotesize Izvorne zapisy / Original records: HIST-g163-001.}\par
\end{minipage}\par
\par\medskip\Needspace{7\baselineskip}
\noindent\begin{minipage}{\linewidth}
\hypertarget{TERM-0124}{}
\noindent\textbf{občeje |tE-C|}\quad{\footnotesize\texttt{TERM-0124}}\par
\noindent\emph{Iztočnik / Source:} allgemeiner |tE-C|\par
\noindent\emph{Rusky / Russian:} в более общем виде |tE-C|\par
\noindent\emph{Ukrajinsky / Ukrainian:} загальніше |tE-C|\par
\smallskip\noindent\emph{Razlog:} Občeje jest prislovnik, ktory govori, že slědujuči zapis jest občejši. On ne stvara novo ime objekta: $|tE-C|$ ostavaje ista karakteristična determinanta. Pridavnik v izrazu obča determinanta mogl by sugerovati iny tip determinanty. Kratkost byla dodatnym razlogom za občeje; ona ne jest izměrenoju razumělivostju.\par
\noindent\emph{Reason:} The adverb občeje says that the following formulation is more general. It does not name a new object: $|tE-C|$ remains the characteristic determinant. An adjective in a phrase such as general determinant could suggest a different kind of determinant. Brevity was an additional reason for občeje, not a measured intelligibility result.\par
\noindent{\footnotesize Izvorne zapisy / Original records: HIST-g135-002.}\par
\end{minipage}\par
\par\medskip\Needspace{7\baselineskip}
\noindent\begin{minipage}{\linewidth}
\hypertarget{TERM-0125}{}
\noindent\textbf{najobčejše jednoznačne rěšenja}\quad{\footnotesize\texttt{TERM-0125}}\par
\noindent\emph{Iztočnik / Source:} allgemeinste eindeutige Lösungen\par
\noindent\emph{Rusky / Russian:} самые общие однозначные решения\par
\noindent\emph{Ukrajinsky / Ukrainian:} найзагальніші однозначні розв'язки\par
\smallskip\noindent\emph{Razlog:} Najobčejše i jednoznačne zachovujut dve različne vlastnosti rěšenij: njih občnost i jednoznačnost. Izraz ne trěba skratiti do najobčejše rěšenja, niti razuměti kako tvrđenje, že postoji samo jedno rěšenje. Zapis izbira běžny matematičny register bez dodavanja starinnogo stilja.\par
\noindent\emph{Reason:} Najobčejše and jednoznačne preserve two distinct properties of the solutions: generality and single-valuedness. The phrase should neither lose the second property nor be read as asserting that only one solution exists. The recorded preference is for ordinary mathematical prose rather than an artificially archaic style.\par
\noindent{\footnotesize Izvorne zapisy / Original records: HIST-g069-007.}\par
\end{minipage}\par
\par\medskip\Needspace{7\baselineskip}
\noindent\begin{minipage}{\linewidth}
\hypertarget{TERM-0126}{}
\noindent\textbf{najobčejše razloženja po redah}\quad{\footnotesize\texttt{TERM-0126}}\par
\noindent\emph{Iztočnik / Source:} allgemeinste Reihenentwicklungen\par
\noindent\emph{Rusky / Russian:} наиболее общие разложения по рядам\par
\noindent\emph{Ukrajinsky / Ukrainian:} найзагальніші розклади за рядами\par
\smallskip\noindent\emph{Razlog:} Razloženja po redah zachovuje formulaciju v kontekstu redov, svezanom s metodami Mertensa, Capellija i Deruytsa. Najobčejše ostavaje čestju tvrđenja o opsegu teh razloženij. Zapis ne daje razlog avtomatično razuměti vsako Reihenentwicklung kako stepenni red; konkretny matematičny objekt trěba čitati iz pasaža.\par
\noindent\emph{Reason:} Razloženja po redah retains the row-expansion formulation in the context associated with Mertens, Capelli, and Deruyts. Najobčejše preserves the asserted generality of those expansions. The record does not justify automatically interpreting every occurrence of Reihenentwicklung as a power-series expansion; the mathematical object must be identified from the passage.\par
\noindent{\footnotesize Izvorne zapisy / Original records: HIST-g057-013.}\par
\end{minipage}\par
\par\medskip\Needspace{7\baselineskip}
\noindent\begin{minipage}{\linewidth}
\hypertarget{TERM-0127}{}
\noindent\textbf{razmatrivana kako kovariant}\quad{\footnotesize\texttt{TERM-0127}}\par
\noindent\emph{Iztočnik / Source:} als Kovariante betrachtet\par
\noindent\emph{Rusky / Russian:} рассматриваемая как ковариант\par
\noindent\emph{Ukrajinsky / Ukrainian:} розглянута як коваріант\par
\smallskip\noindent\emph{Razlog:} Razmatrivana kako kovariant govori o načinu razmatranja formy, ne o proměně samej formy. Kovariant prodolžaje medžunarodnu terminologiju kovariantov i kontravariantov iz uvoda i odděla 11. Ta veza jest razlogom za izbor, ale tvrdnja o lehkoj razumělivosti za vse slovjanske čitatelje ne slěduje iz samogo medžunarodnogo korenja.\par
\noindent\emph{Reason:} Razmatrivana kako kovariant states how the form is regarded, not that the form itself changes. Kovariant continues the international covariant/contravariant terminology of the introduction and section 11. That continuity explains the choice; an international root alone does not establish easy comprehension by all Slavic readers.\par
\noindent{\footnotesize Izvorne zapisy / Original records: HIST-g020-008.}\par
\end{minipage}\par
\par\medskip\Needspace{7\baselineskip}
\noindent\begin{minipage}{\linewidth}
\hypertarget{TERM-0128}{}
\noindent\textbf{adjungovano kako modul}\quad{\footnotesize\texttt{TERM-0128}}\par
\noindent\emph{Iztočnik / Source:} als Modul adjungiert\par
\noindent\emph{Rusky / Russian:} присоединено как модуль\par
\noindent\emph{Ukrajinsky / Ukrainian:} приєднано як модуль\par
\smallskip\noindent\emph{Razlog:} Adjungovano kako modul izrično zachovuje kako modul pri formuli 11 v oddělu 10. Izraz pripada familiji adjungovati, upotrěbjanoj v oddělah 7 i 9. Rusky i ukrajinsky prěvody imajut obične tehnične slovesne odgovaranja; medžuslovjansky zapis izbira medžunarodnu osnovu. Ta leksična odluka byla označena za daljše jezyčno razmotrenje, ne za dokazanu normu.\par
\noindent\emph{Reason:} Adjungovano kako modul retains the explicit qualification as a module beside formula 11 in section 10. It belongs to the adjungovati family used in sections 7 and 9. The Russian and Ukrainian versions use ordinary technical verbs, whereas the Interslavic candidate uses an international stem. The lexical choice was left for further linguistic examination, not established as a norm.\par
\noindent{\footnotesize Izvorne zapisy / Original records: HIST-g011-013.}\par
\end{minipage}\par
\par\medskip\Needspace{7\baselineskip}
\noindent\begin{minipage}{\linewidth}
\hypertarget{TERM-0129}{}
\noindent\textbf{adjungujut se kako moduly}\quad{\footnotesize\texttt{TERM-0129}}\par
\noindent\emph{Iztočnik / Source:} als Moduln adjungiert\par
\noindent\emph{Rusky / Russian:} присоединяются как модули\par
\noindent\emph{Ukrajinsky / Ukrainian:} приєднуються як модулі\par
\smallskip\noindent\emph{Razlog:} Adjungujut se kako moduly zachovuje slovesnu konstrukciju i množinu modulov. Zapis dava prednost adjungovati radi tehničnoj točnosti, zatymčas rusky prisojedinjat i ukrajinsky pryjednuvaty pokazuju možnost običnogo slovjanskogo slovesa. To jest stvarno sravnjenje alternativnyh načinov imenovanja; zapis ne zaključaje, že medžunarodna osnova jest jedina pravilna možnost.\par
\noindent\emph{Reason:} Adjungujut se kako moduly preserves a verbal construction and plural modules. The record prefers adjungovati for technical precision, while the Russian and Ukrainian equivalents illustrate the alternative of an ordinary Slavic verb meaning to adjoin. This is a comparison between ways of naming the operation, not a finding that the international stem is the only correct choice.\par
\noindent{\footnotesize Izvorne zapisy / Original records: HIST-g008-005.}\par
\end{minipage}\par
\par\medskip\Needspace{7\baselineskip}
\noindent\begin{minipage}{\linewidth}
\hypertarget{TERM-0130}{}
\noindent\textbf{kako produkty sistemov I i II ostavaju}\quad{\footnotesize\texttt{TERM-0130}}\par
\noindent\emph{Iztočnik / Source:} Als Produkte von System I und II bleiben\par
\noindent\emph{Rusky / Russian:} как произведения систем I и II остаются\par
\noindent\emph{Ukrajinsky / Ukrainian:} як добутки систем I і II залишаються\par
\smallskip\noindent\emph{Razlog:} Kako produkty sistemov I i II ostavaju zachovuje omejenje zaključka na produkty dvuh imenovanyh sistemov i veže větu s dvoma izpisanymi produktami. I i II sut oznaky sistemov; ostavajut rimske čisla i v kiriličnom tekstu. Njih funkcija jest matematična identifikacija, ne običny pravopis slova.\par
\noindent\emph{Reason:} Kako produkty sistemov I i II ostavaju preserves the restriction to products of the two named systems and connects the sentence with the two displayed products. I and II identify the systems and remain Roman numerals in Cyrillic prose. They function as mathematical labels, not words to be transliterated.\par
\noindent{\footnotesize Izvorne zapisy / Original records: HIST-g020-013.}\par
\end{minipage}\par
\par\medskip\Needspace{7\baselineskip}
\noindent\begin{minipage}{\linewidth}
\hypertarget{TERM-0131}{}
\noindent\textbf{byly poznane kako reducibilne}\quad{\footnotesize\texttt{TERM-0131}}\par
\noindent\emph{Iztočnik / Source:} als reduzibel erkannt\par
\noindent\emph{Rusky / Russian:} были признаны приводимыми\par
\noindent\emph{Ukrajinsky / Ukrainian:} були визнані звідними\par
\smallskip\noindent\emph{Razlog:} Byly poznane kako reducibilne jest konstrukcija soglasovana s imennikom formy v ženskoj množině. Sučstveny razlog zapisa jest gramatično soglasovanje; on ne jest novym sravnjenjem korenj za reducibilnost. Zachovuje se smysl, že reducibilnost byla ustanovjena, ne že formy stali reducibilnymi v toj chvíli.\par
\noindent\emph{Reason:} Byly poznane kako reducibilne agrees with feminine plural formy. The recorded point is grammatical agreement, not a new comparison of lexical roots for reducibility. The sentence says that reducibility was recognized, not that the forms became reducible at that moment.\par
\noindent{\footnotesize Izvorne zapisy / Original records: HIST-g021-011.}\par
\end{minipage}\par
\par\medskip\Needspace{7\baselineskip}
\noindent\begin{minipage}{\linewidth}
\hypertarget{TERM-0132}{}
\noindent\textbf{kako nerazložima veličina}\quad{\footnotesize\texttt{TERM-0132}}\par
\noindent\emph{Iztočnik / Source:} als unzerlegbare Größe\par
\noindent\emph{Rusky / Russian:} как неразложимая величина\par
\noindent\emph{Ukrajinsky / Ukrainian:} як нерозкладна величина\par
\smallskip\noindent\emph{Razlog:} Veličina zachovuje obči matematičny smysl německogo Größe v izrazu kako nerazložima veličina. Zapis ne sužaje objekt do skalarnoj měry ili fizičnoj velikosti. Nerazložima označaje matematičnu vlastnost objekta; konkretna vrsta objekta ostavaje opreděljena kontekstom.\par
\noindent\emph{Reason:} Veličina preserves the broad mathematical use of Größe in kako nerazložima veličina. It does not narrow the object to a scalar measure or physical size. Nerazložima states a mathematical property; the kind of object remains determined by the context.\par
\noindent{\footnotesize Izvorne zapisy / Original records: HIST-g032-005.}\par
\end{minipage}\par
\par\medskip\Needspace{7\baselineskip}
\noindent\begin{minipage}{\linewidth}
\hypertarget{TERM-0133}{}
\noindent\textbf{druga baza togo samogo ideala}\quad{\footnotesize\texttt{TERM-0133}}\par
\noindent\emph{Iztočnik / Source:} andere Basis desselben Ideals\par
\noindent\emph{Rusky / Russian:} другая база того же идеала\par
\noindent\emph{Ukrajinsky / Ukrainian:} інша база того самого ідеалу\par
\smallskip\noindent\emph{Razlog:} Togo samogo ideala jest ključno omejenje: měnjaje se baza, ne ideal. Druga baza zato ne može byti izjasnjena kako baza drugogo ideala ili drugoj klasy. Zapis zachovuje razliku medžu proměnoju predstavjenja odnogo objekta i sravnjenjem dvuh objektov.\par
\noindent\emph{Reason:} Togo samogo ideala is the crucial restriction: the basis changes, not the ideal. Druga baza must not become the basis of another ideal or another class. The wording distinguishes a change of presentation of one object from a comparison between two objects.\par
\noindent{\footnotesize Izvorne zapisy / Original records: HIST-g132-001.}\par
\end{minipage}\par
\par\medskip\Needspace{7\baselineskip}
\noindent\begin{minipage}{\linewidth}
\hypertarget{TERM-0134}{}
\noindent\textbf{početno uslovje}\quad{\footnotesize\texttt{TERM-0134}}\par
\noindent\emph{Iztočnik / Source:} Anfangsbedingung\par
\noindent\emph{Rusky / Russian:} начальное условие\par
\noindent\emph{Ukrajinsky / Ukrainian:} початкова умова\par
\smallskip\noindent\emph{Razlog:} Početno uslovje jest neposrědno odgovaranje za načelno ili početno uslovje v matematičnom izlaganju. Ruska i ukrajinska paralela pokazuju istu imennikovu konstrukciju. Kratky izvorni razlog ne sodržuje osobny citat iz medžuslovjanskogo matematičnogo iztočnika; izbor zato ne trěba predstavljati kako nezavisno potvrđenu tehničnu normu.\par
\noindent\emph{Reason:} Početno uslovje directly expresses an initial condition in mathematical exposition. The Russian and Ukrainian parallels have the same noun-phrase construction. The brief original rationale supplies no separate citation from Interslavic mathematical usage, so the choice should not be presented as an independently established technical norm.\par
\noindent{\footnotesize Izvorne zapisy / Original records: HIST-g069-012.}\par
\end{minipage}\par
\par\medskip\Needspace{7\baselineskip}
\noindent\begin{minipage}{\linewidth}
\hypertarget{TERM-0135}{}
\noindent\textbf{početna forma}\quad{\footnotesize\texttt{TERM-0135}}\par
\noindent\emph{Iztočnik / Source:} Anfangsform\par
\noindent\emph{Rusky / Russian:} начальная форма\par
\noindent\emph{Ukrajinsky / Ukrainian:} початкова форма\par
\smallskip\noindent\emph{Razlog:} Početna forma označaje formu, od ktorej začne razmatranje. Izbor prodolžaje istu familiju početn- v pozdnějših tvrđenjah o redukciji. Doslědnost jest ovde lokalny razlog; početna ne vnašaje dodatnu definiciju, ktora ne jest dana v avtorskym kontekstu.\par
\noindent\emph{Reason:} Početna forma names the form from which the discussion starts. The choice maintains the početn- family in later reduction statements. This is a local consistency argument; početna does not introduce an additional definition absent from the author's context.\par
\noindent{\footnotesize Izvorne zapisy / Original records: HIST-g003-009.}\par
\end{minipage}\par
\par\medskip\Needspace{7\baselineskip}
\noindent\begin{minipage}{\linewidth}
\hypertarget{TERM-0136}{}
\noindent\textbf{početne členy}\quad{\footnotesize\texttt{TERM-0136}}\par
\noindent\emph{Iztočnik / Source:} Anfangsglieder\par
\noindent\emph{Rusky / Russian:} начальные члены\par
\noindent\emph{Ukrajinsky / Ukrainian:} початкові члени\par
\smallskip\noindent\emph{Razlog:} Početne členy veže se s početna forma črěz isti pridavnik, ale imenniki označajut različne objekty: členy ne sut sinonimom cěloj formy. Izbor zachovuje tu razliku, zatymčas daje čitatelju poznavajemu familiju izrazov o početku razmatranja.\par
\noindent\emph{Reason:} Početne členy shares its adjective with početna forma, but the nouns denote different objects: terms are not synonymous with the whole form. The choice preserves that distinction while giving readers a recognizable family of expressions for the starting part of a construction.\par
\noindent{\footnotesize Izvorne zapisy / Original records: HIST-g004-009.}\par
\end{minipage}\par
\par\medskip\Needspace{7\baselineskip}
\noindent\begin{minipage}{\linewidth}
\hypertarget{TERM-0137}{}
\noindent\textbf{shema}\quad{\footnotesize\texttt{TERM-0137}}\par
\noindent\emph{Iztočnik / Source:} Ansatz\par
\noindent\emph{Rusky / Russian:} схема\par
\noindent\emph{Ukrajinsky / Ukrainian:} схема\par
\smallskip\noindent\emph{Razlog:} Shema prědaje funkciju německogo Ansatz v danom vyčisljenju vměsto prějemanja německogo slova. Razlog jest, že čitatelj trěbuje poznati razmatrany obrazec postupka, ne novo inojezyčno nazvanje. To ne jest obče pravilo prěvoditi Ansatz kako shema v vsakoj matematičnoj oblasti.\par
\noindent\emph{Reason:} Shema conveys the function of Ansatz in the calculation instead of borrowing the German word. The reader needs to recognize the proposed pattern of calculation, not learn an additional foreign label. This is not a universal rule that Ansatz means shema in every mathematical subject.\par
\noindent{\footnotesize Izvorne zapisy / Original records: HIST-g023-004.}\par
\end{minipage}\par
\par\medskip\Needspace{7\baselineskip}
\noindent\begin{minipage}{\linewidth}
\hypertarget{TERM-0138}{}
\noindent\textbf{začetna shema}\quad{\footnotesize\texttt{TERM-0138}}\par
\noindent\emph{Iztočnik / Source:} Ansatz\par
\noindent\emph{Rusky / Russian:} схема\par
\noindent\emph{Ukrajinsky / Ukrainian:} схема\par
\smallskip\noindent\emph{Razlog:} Začetna shema izjasnjaje, že izpisany obrazec jest početnym postavjenjem vyčisljenja. Zapis izběgaje německy zaimani termin Ansatz, ale sam izrično ostavjaje pytanje doslědnogo napisa i jezyčnoj formy za daljše razmotrenje. Proto začetna shema ostavaje dokumentovanym kandidatom, ne avtomatično predpisanymi slovami za vsaky pojav Ansatz.\par
\noindent\emph{Reason:} Začetna shema explains that the displayed pattern is the starting setup of the calculation. The record avoids the German loan Ansatz but explicitly leaves spelling consistency and linguistic form for further examination. The phrase is therefore a documented candidate, not an automatic prescription for every occurrence of Ansatz.\par
\noindent{\footnotesize Izvorne zapisy / Original records: HIST-g014-013.}\par
\end{minipage}\par
\par\medskip\Needspace{7\baselineskip}
\noindent\begin{minipage}{\linewidth}
\hypertarget{TERM-0139}{}
\noindent\textbf{antistrofna matrica}\quad{\footnotesize\texttt{TERM-0139}}\par
\noindent\emph{Iztočnik / Source:} antistrophe Matrix\par
\noindent\emph{Rusky / Russian:} антистрофная матрица\par
\noindent\emph{Ukrajinsky / Ukrainian:} антистрофна матриця\par
\smallskip\noindent\emph{Razlog:} Antistrofna matrica zachovuje istorično nazvanje svezano s Frobeniusom. Medžunarodna osnova dopuščaje čitatelju povezati prěvod s terminom v originalu. Zapis ne dava osnovu tiho zaměniti termin nekym sovrěmennym imenom matrice: takvo utočnjenje trěbovalo by zasebno sravnjenje definicij.\par
\noindent\emph{Reason:} Antistrofna matrica retains the historical name associated with Frobenius. The international stem helps readers connect the translation with the original term. The record does not justify silently replacing it with a modern matrix name; that identification would require a separate comparison of definitions.\par
\noindent{\footnotesize Izvorne zapisy / Original records: HIST-g201-008.}\par
\end{minipage}\par
\par\medskip\Needspace{7\baselineskip}
\noindent\begin{minipage}{\linewidth}
\hypertarget{TERM-0140}{}
\noindent\textbf{ograničenje čisla}\quad{\footnotesize\texttt{TERM-0140}}\par
\noindent\emph{Iztočnik / Source:} Anzahlbeschränkung\par
\noindent\emph{Rusky / Russian:} ограничение числа\par
\noindent\emph{Ukrajinsky / Ukrainian:} обмеження кількості\par
\smallskip\noindent\emph{Razlog:} Ograničenje čisla govori o broju racionalnyh bazovyh funkcij. V pasažu jest važna razlika medžu granicami $\rho+1$ i $\rho$; predmet ograničenja zato ne može byti izpuščeny. Izraz ne označaje ograničenost cěnnosti samyh funkcij.\par
\noindent\emph{Reason:} Ograničenje čisla concerns the number of rational basis functions. The distinction between bounds $\rho+1$ and $\rho$ matters in the passage, so the object of the bound must remain explicit. The phrase does not mean boundedness of the values of the functions.\par
\noindent{\footnotesize Izvorne zapisy / Original records: HIST-g047-012.}\par
\end{minipage}\par
\par\medskip\Needspace{7\baselineskip}
\noindent\begin{minipage}{\linewidth}
\hypertarget{TERM-0141}{}
\noindent\textbf{ravnost čisla komponentov}\quad{\footnotesize\texttt{TERM-0141}}\par
\noindent\emph{Iztočnik / Source:} Anzahlgleichheit der Komponenten\par
\noindent\emph{Rusky / Russian:} равенство числа компонентов\par
\noindent\emph{Ukrajinsky / Ukrainian:} рівність кількості компонент\par
\smallskip\noindent\emph{Razlog:} Ravnost čisla komponentov vyražaje točno tvrđenje odděla 3: najkratše razloženja imajut isto čislo komponentov. Izbor ne zaměnjaje tu větu občejšim nazvanjem teorema o jednoznačnosti. Ravnost broja sama ne tvrdi ravnost komponentov po jedinomu.\par
\noindent\emph{Reason:} Ravnost čisla komponentov states the particular assertion in section 3: shortest decompositions have the same number of components. It does not replace that assertion with a broader uniqueness-theorem label. Equality of the number alone does not assert component-by-component equality.\par
\noindent{\footnotesize Izvorne zapisy / Original records: HIST-g081-035, HIST-g082-035.}\par
\end{minipage}\par
\par\medskip\Needspace{7\baselineskip}
\noindent\begin{minipage}{\linewidth}
\hypertarget{TERM-0142}{}
\noindent\textbf{Aritmetična formulacija pojma dimenzije}\quad{\footnotesize\texttt{TERM-0142}}\par
\noindent\emph{Iztočnik / Source:} Arithmetische Fassung des Dimensionsbegriffes\par
\noindent\emph{Rusky / Russian:} Арифметическая формулировка понятия размерности\par
\noindent\emph{Ukrajinsky / Ukrainian:} Арифметичне формулювання поняття розмірності\par
\smallskip\noindent\emph{Razlog:} Aritmetična formulacija pojma dimenzije prědaje Fassung kako način formulovanja pojma, ne kako doslovnu metaforu. Kontekst veže dimenziju s lancami prostyh idealov. Naslov zato zachovuje i aritmetičny pristup, i to, že predmetom jest pojem dimenzije.\par
\noindent\emph{Reason:} Aritmetična formulacija pojma dimenzije renders Fassung as a formulation of a concept rather than a literal metaphor. The context relates dimension to chains of prime ideals. The title preserves both the arithmetic approach and the fact that the subject is the concept of dimension.\par
\noindent{\footnotesize Izvorne zapisy / Original records: HIST-g098-001.}\par
\end{minipage}\par
\par\medskip\Needspace{7\baselineskip}
\noindent\begin{minipage}{\linewidth}
\hypertarget{TERM-0143}{}
\noindent\textbf{prinaležny sistem veličin}\quad{\footnotesize\texttt{TERM-0143}}\par
\noindent\emph{Iztočnik / Source:} Associated system of values\par
\noindent\emph{Rusky / Russian:} связанная система значений\par
\noindent\emph{Ukrajinsky / Ukrainian:} пов'язана система значень\par
\smallskip\noindent\emph{Razlog:} Prinaležny imaje označiti odnos medžu danim sistemom i veličinami, ne vlasničstvo v običnom smyslu. Matematičny kontekst omejaje značenje, ale izvorni zapis označaje pridavnik kako važno pytanje za jezyčno razmotrenje. Ruska svjazannaja i ukrajinska povjazana pokazuju alternativno izražanje črěz svezanost. Sama matematična jasnoća odnosa ne rěšaje, ktory medžuslovjansky pridavnik jest najprikladnějši.\par
\noindent\emph{Reason:} Prinaležny is intended to express the relation between the system and its values, not ordinary ownership. The mathematics constrains the sense, but the original record marks the adjective as a significant linguistic review question. Russian and Ukrainian express the relation through connectedness. A mathematically clear relation does not itself settle the best Interslavic adjective.\par
\noindent{\footnotesize Izvorne zapisy / Original records: HIST-g091-007, HIST-g092-076.}\par
\end{minipage}\par
\par\medskip\Needspace{7\baselineskip}
\noindent\begin{minipage}{\linewidth}
\hypertarget{TERM-0144}{}
\noindent\textbf{beskonečno blizke točke / točke razvětvenja}\quad{\footnotesize\texttt{TERM-0144}}\par
\noindent\emph{Iztočnik / Source:} aufeinanderfolgende Punkte / Verzweigungspunkte\par
\noindent\emph{Rusky / Russian:} бесконечно близкие точки / точки ветвления\par
\noindent\emph{Ukrajinsky / Ukrainian:} нескінченно близькі точки / точки розгалуження\par
\smallskip\noindent\emph{Razlog:} Zapis razuměje aufeinanderfolgende Punkte v tom geometriskom kontekstu kako beskonečno blizke točke, ne prosto točke, ktore slědujut jedna po drugoj. To jest kontekstualno tolkovanje svezano s razmatranjem osoblivostij v tradiciji Brilla i Noether. Točke razvětvenja zachovujut drugu čest para. Takvo tolkovanje ne može služiti kako obči slovničny ekvivalent za vsako aufeinanderfolgend.\par
\noindent\emph{Reason:} The record interprets aufeinanderfolgende Punkte in this geometric context as infinitely near points, not merely points in a sequence. This is a contextual reading tied to the discussion of singularities in the Brill-Noether tradition. Točke razvětvenja retains the second term, branching points. The interpretation is not a general dictionary equivalent of aufeinanderfolgend.\par
\noindent{\footnotesize Izvorne zapisy / Original records: HIST-g076-016.}\par
\end{minipage}\par
\par\medskip\Needspace{7\baselineskip}
\noindent\begin{minipage}{\linewidth}
\hypertarget{TERM-0145}{}
\noindent\textbf{izomorfne kako R-moduly}\quad{\footnotesize\texttt{TERM-0145}}\par
\noindent\emph{Iztočnik / Source:} aufgefaßt als R-Moduln -- isomorph\par
\noindent\emph{Rusky / Russian:} изоморфны как R-модули\par
\noindent\emph{Ukrajinsky / Ukrainian:} ізоморфні як R-модулі\par
\smallskip\noindent\emph{Razlog:} Izomorfne kako R-moduly izrično imenuje strukturu, v ktorej se tvrdi izomorfizm. To ne jest samo ravnost klas idealov, ni proizvoljno odgovaranje množstv, ni nužno izomorfizm prstenov. Kako R-moduly zato nosi matematičnu informaciju i ne može byti izpuščeno radi kratkosti.\par
\noindent\emph{Reason:} Izomorfne kako R-moduly explicitly names the structure in which the isomorphism holds. It is not merely equality of ideal classes, an arbitrary correspondence of sets, or necessarily a ring isomorphism. As R-modules carries mathematical information and cannot be dropped for brevity.\par
\noindent{\footnotesize Izvorne zapisy / Original records: HIST-g133-003.}\par
\end{minipage}\par
\par\medskip\Needspace{7\baselineskip}
\noindent\begin{minipage}{\linewidth}
\hypertarget{TERM-0146}{}
\noindent\textbf{rěšimo polje razloženja}\quad{\footnotesize\texttt{TERM-0146}}\par
\noindent\emph{Iztočnik / Source:} auflösbarer Zerfällungskörper\par
\noindent\emph{Rusky / Russian:} разрешимое поле разложения\par
\noindent\emph{Ukrajinsky / Ukrainian:} розв'язне поле розщеплення\par
\smallskip\noindent\emph{Razlog:} Rěšimo v izrazu rěšimo polje razloženja odnosi se k tehničnomu pojmu rěšimosti iz teorije grup. Ono ne govori samo, že neka zadača može byti rěšena. Polje razloženja ostavaje imenovanym objektom; zapis zachovuje oba sostavnika termina bez dodavanja drugoj definicije rěšimosti.\par
\noindent\emph{Reason:} Rěšimo in rěšimo polje razloženja concerns the technical, group-theoretic sense of solvability. It does not merely say that some problem can be solved. Polje razloženja remains the named object, with both parts of the term retained and no additional definition of solvability supplied by the translation.\par
\noindent{\footnotesize Izvorne zapisy / Original records: HIST-g208-008.}\par
\end{minipage}\par
\par\medskip\Needspace{7\baselineskip}
\noindent\begin{minipage}{\linewidth}
\hypertarget{TERM-0147}{}
\noindent\textbf{iz punktov 3 i 4 slěduje}\quad{\footnotesize\texttt{TERM-0147}}\par
\noindent\emph{Iztočnik / Source:} Aus 3 und 4 folgt\par
\noindent\emph{Rusky / Russian:} из пунктов 3 и 4 следует\par
\noindent\emph{Ukrajinsky / Ukrainian:} з пунктів 3 і 4 випливає\par
\smallskip\noindent\emph{Razlog:} Iz punktov 3 i 4 slěduje imenuje zavisimost zaključka od dvuh prědhodnyh tvrđenij: o lokalnom faktor-prstenu i o diskriminantnom idealu. Dodavanje punktov čini jasnym, že 3 i 4 sut oznaky čestij dokaza, ne čisla v aritmetičnoj operaciji. Izvorne oznaky ostavajut neproměnjene.\par
\noindent\emph{Reason:} Iz punktov 3 i 4 slěduje makes the conclusion's dependence on two preceding assertions explicit: the local quotient-ring statement and the discriminant-ideal argument. Adding punktov identifies 3 and 4 as proof-item labels, not operands in an arithmetic calculation. The original labels remain unchanged.\par
\noindent{\footnotesize Izvorne zapisy / Original records: HIST-g164-001.}\par
\end{minipage}\par
\par\medskip\Needspace{7\baselineskip}
\noindent\begin{minipage}{\linewidth}
\hypertarget{TERM-0148}{}
\noindent\textbf{iz predstavjenja v obliku produkta stepenjev}\quad{\footnotesize\texttt{TERM-0148}}\par
\noindent\emph{Iztočnik / Source:} aus der Potenzdarstellung\par
\noindent\emph{Rusky / Russian:} из представления в виде произведения степеней\par
\noindent\emph{Ukrajinsky / Ukrainian:} з подання у вигляді добутку степенів\par
\smallskip\noindent\emph{Razlog:} Predstavjenje v obliku produkta stepenjev razkriva skračenost německogo Potenzdarstellung. Razlogom jest veza s prědhodnym Potenzprodukt: govori se o produktu stepenjev, ne neizběžno o jednoj stepeni. Dolga konstrukcija zato ne dodavaje novu operaciju, ale izjasnjaje onu, ktora jest už dana kontekstom.\par
\noindent\emph{Reason:} Predstavjenje v obliku produkta stepenjev unpacks the compressed Potenzdarstellung. Its reason is the connection with the preceding Potenzprodukt: the representation is a product of powers, not necessarily a single power. The longer phrase clarifies the operation already given by the context rather than adding a new one.\par
\noindent{\footnotesize Izvorne zapisy / Original records: HIST-g156-006.}\par
\end{minipage}\par
\par\medskip\Needspace{7\baselineskip}
\noindent\begin{minipage}{\linewidth}
\hypertarget{TERM-0149}{}
\noindent\textbf{kolco / ideal / P-modul izvedeny iz S v T}\quad{\footnotesize\texttt{TERM-0149}}\par
\noindent\emph{Iztočnik / Source:} aus S in T abgeleiteter Ring / Ideal / P-Modul\par
\noindent\emph{Rusky / Russian:} кольцо / идеал / P-модуль, выведенный из S в T\par
\noindent\emph{Ukrajinsky / Ukrainian:} кільце / ідеал / P-модуль, виведений із S у T\par
\smallskip\noindent\emph{Razlog:} Izvedeny iz S v T zachovuje smer konstrukcije i njezinu definiciju črěz prěsek prěd osobnym idealno-teoretičnym upotrěbjenjem. Skupny zapis s kosymi črtami objedinjuje tri objekty, ne jednu gotovu větu: kolco jest srědnjego roda, ideal i modul mužskogo, zato pridavnik trěbuje odgovarjajuče soglasovanje. Ranějše kolco takože ne ustanavjaje samo po sobě današnji termin vměsto prsten.\par
\noindent\emph{Reason:} Izvedeny iz S v T preserves the direction of construction and its intersection definition before the ideal-theoretic specialization. The slash-separated entry groups three objects; it is not one ready-made grammatical phrase. Kolco is neuter, whereas ideal and modul are masculine, so the adjective requires the appropriate agreement. Historical kolco also does not by itself establish the current preference over prsten.\par
\noindent{\footnotesize Izvorne zapisy / Original records: HIST-g119-003.}\par
\end{minipage}\par
\par\medskip\Needspace{7\baselineskip}
\noindent\begin{minipage}{\linewidth}
\hypertarget{TERM-0150}{}
\noindent\textbf{nauka o razširjenju}\quad{\footnotesize\texttt{TERM-0150}}\par
\noindent\emph{Iztočnik / Source:} Ausdehnungslehre\par
\noindent\emph{Rusky / Russian:} учение о протяженности\par
\noindent\emph{Ukrajinsky / Ukrainian:} теорія протяжності\par
\smallskip\noindent\emph{Razlog:} Nauka o razširjenju služi objasnjenju Grassmannovogo pojmovogo aparata v proze. Bibliografično nazvanje Ausdehnungslehre ostavaje v německoj formě, kako ime citovanogo děla. Odluka zato razděljuje razjasnjujuči prěvod pojma i točno imenovanje knigy; jedno ne zaměnjaje drugo.\par
\noindent\emph{Reason:} Nauka o razširjenju explains Grassmann's conceptual framework in the prose. The bibliographical title Ausdehnungslehre remains in German as the name of the cited work. The decision distinguishes an explanatory rendering of the concept from the exact identification of the book; neither substitutes for the other.\par
\noindent{\footnotesize Izvorne zapisy / Original records: HIST-g029-015.}\par
\end{minipage}\par
\par\medskip\Needspace{7\baselineskip}
\noindent\begin{minipage}{\linewidth}
\hypertarget{TERM-0151}{}
\noindent\textbf{izpolnjenje redukcij}\quad{\footnotesize\texttt{TERM-0151}}\par
\noindent\emph{Iztočnik / Source:} Ausführung der Reduktionen\par
\noindent\emph{Rusky / Russian:} выполнение редукций\par
\noindent\emph{Ukrajinsky / Ukrainian:} виконання редукцій\par
\smallskip\noindent\emph{Razlog:} Izpolnjenje redukcij označaje izvedenje postupkov redukcije, ne novu vrstu redukcije. Ruska i ukrajinska paralela podpirajut običnu proceduralnu konstrukciju. Izbor medžuslovjanskogo izpolnjenje byl obrazloženy razumělivostju, ale sam zapis ostavjaje jego idiomatičnost i oporu v jezyčnom upotrěbjenju za daljše razmotrenje.\par
\noindent\emph{Reason:} Izpolnjenje redukcij denotes carrying out the reduction procedures, not a new kind of reduction. The Russian and Ukrainian parallels support an ordinary procedural construction. The Interslavic noun was preferred on grounds of intelligibility, but the record itself leaves its idiomaticity and usage support for further examination.\par
\noindent{\footnotesize Izvorne zapisy / Original records: HIST-g012-005.}\par
\end{minipage}\par
\par\medskip\Needspace{7\baselineskip}
\noindent\begin{minipage}{\linewidth}
\hypertarget{TERM-0152}{}
\noindent\textbf{izhodna forma}\quad{\footnotesize\texttt{TERM-0152}}\par
\noindent\emph{Iztočnik / Source:} Ausgangsform\par
\noindent\emph{Rusky / Russian:} исходная форма\par
\noindent\emph{Ukrajinsky / Ukrainian:} вихідна форма\par
\smallskip\noindent\emph{Razlog:} Izhodna forma jest forma, od ktorej začne razmatrany red form. Izbor zachovuje tu početnu funkciju, ne sugerujuči teoriju zadač s početnymi uslovjami. Sličnost običnyh slov o početku ne daje razlog proměniti matematičnu oblast ili vrstu objekta.\par
\noindent\emph{Reason:} Izhodna forma is the starting form for the sequence of forms under discussion. The wording preserves that starting role without suggesting the theory of initial-value problems. Similar everyday vocabulary about a starting point is not a reason to change the mathematical setting or type of object.\par
\noindent{\footnotesize Izvorne zapisy / Original records: HIST-g038-003.}\par
\end{minipage}\par
\par\medskip\Needspace{7\baselineskip}
\noindent\begin{minipage}{\linewidth}
\hypertarget{TERM-0153}{}
\noindent\textbf{sut izpušćene}\quad{\footnotesize\texttt{TERM-0153}}\par
\noindent\emph{Iztočnik / Source:} ausgelassen sind\par
\noindent\emph{Rusky / Russian:} опущены\par
\noindent\emph{Ukrajinsky / Ukrainian:} опущено\par
\smallskip\noindent\emph{Razlog:} Sut izpušćene jest čestju podnožnoj primětky, ktora objasnjaje znak ekvivalentnosti i govori, ktore česti sut izpuščene. Izbor dava prednost izričnomu objasnjenju. Izvorni zapis ostavjaje pravopis i idiomatičnost kandidata izpušćene za razmotrenje: odgovaranje ć → č v osnovnom pravopisě samo ne dokazuje pravilnost cěloj konstrukcije.\par
\noindent\emph{Reason:} Sut izpušćene belongs to the footnote explaining an equivalence sign and identifying what is omitted. The choice favours an explicit explanation. The original record leaves the candidate's spelling and idiomaticity open to examination: the basic-spelling correspondence ć → č does not by itself establish the correctness of the whole construction.\par
\noindent{\footnotesize Izvorne zapisy / Original records: HIST-g024-010.}\par
\end{minipage}\par
\par\medskip\Needspace{7\baselineskip}
\noindent\begin{minipage}{\linewidth}
\hypertarget{TERM-0154}{}
\noindent\textbf{vyznačena baza}\quad{\footnotesize\texttt{TERM-0154}}\par
\noindent\emph{Iztočnik / Source:} ausgezeichnete Basis\par
\noindent\emph{Rusky / Russian:} выделенный базис\par
\noindent\emph{Ukrajinsky / Ukrainian:} виділений базис\par
\smallskip\noindent\emph{Razlog:} Vyznačena baza jest osobno izbrana baza: v danom dobře uporjadkovanom množstvu ona jest prva. Vyznačena zato ne znači pohvalu ili vyšu kvalitetu bazy. Razlog izbora jest njezina opreděljena funkcija v konstrukciji, ne stilistična zaměna za običny pridavnik.\par
\noindent\emph{Reason:} Vyznačena baza is a distinguished, selected basis: it is the first one in the given well-order. Vyznačena is not praise or a claim of superior quality. The reason for the adjective is the basis's specified role in the construction, not stylistic variation for an ordinary adjective.\par
\noindent{\footnotesize Izvorne zapisy / Original records: HIST-g113-007.}\par
\end{minipage}\par
\par\medskip\Needspace{7\baselineskip}
\noindent\begin{minipage}{\linewidth}
\hypertarget{TERM-0155}{}
\noindent\textbf{izključene prvočisla / izključene proste idealy}\quad{\footnotesize\texttt{TERM-0155}}\par
\noindent\emph{Iztočnik / Source:} Ausnahmeprimzahlen / Ausnahme-Primideale\par
\noindent\emph{Rusky / Russian:} исключительные простые / исключительные простые идеалы\par
\noindent\emph{Ukrajinsky / Ukrainian:} виняткові прості / виняткові прості ідеали\par
\smallskip\noindent\emph{Razlog:} Izključene prvočisla i izključene proste idealy byli izbrani kako opisne izrazy za avtorske iznimke. Ruska i ukrajinska paralela govore o iznimnosti, zatymčas medžuslovjansky kandidat naglašaje izključenje iz razmatranogo slučaja. Ta razlika trěbuje kontekst: zapis ne dokazuje, že izključeny i iznimny sut zaměnime vo vseh upotrěbjenjah.\par
\noindent\emph{Reason:} Izključene prvočisla and izključene proste idealy were selected as descriptive expressions for the author's exceptional cases. The Russian and Ukrainian parallels express exceptionality, whereas the Interslavic candidate stresses exclusion from the case under consideration. That difference needs contextual scrutiny; the record does not establish that excluded and exceptional are interchangeable in every use.\par
\noindent{\footnotesize Izvorne zapisy / Original records: HIST-g205-011.}\par
\end{minipage}\par
\par\medskip\Needspace{7\baselineskip}
\noindent\begin{minipage}{\linewidth}
\hypertarget{TERM-0156}{}
\noindent\textbf{postulat izbora}\quad{\footnotesize\texttt{TERM-0156}}\par
\noindent\emph{Iztočnik / Source:} Auswahlpostulat\par
\noindent\emph{Rusky / Russian:} постулат выбора\par
\noindent\emph{Ukrajinsky / Ukrainian:} постулат вибору\par
\smallskip\noindent\emph{Razlog:} Postulat izbora zachovuje izričnu zavisimost argumenta od postavky izbora. Ta zavisimost jest čestju avtorskogo izlaganja, ne slobodnym izborom prěvoditelja. Termin ne trěba izpuščati, da by dokaz izgledal bezuslovny, niti dodavati novu ili silnějšu postavku, ktora ne jest v iztočniku.\par
\noindent\emph{Reason:} Postulat izbora preserves the argument's explicit dependence on a choice postulate. That dependence belongs to the author's exposition, not the translator's discretion. The term must neither disappear so that the proof appears unconditional nor be expanded into a new or stronger assumption absent from the source.\par
\noindent{\footnotesize Izvorne zapisy / Original records: HIST-g081-028, HIST-g082-028.}\par
\end{minipage}\par
\par\medskip\Needspace{7\baselineskip}
\noindent\begin{minipage}{\linewidth}
\hypertarget{TERM-0157}{}
\noindent\textbf{tělo avtomorfizmov}\quad{\footnotesize\texttt{TERM-0157}}\par
\noindent\emph{Iztočnik / Source:} Automorphismenkörper\par
\noindent\emph{Rusky / Russian:} тело автоморфизмов\par
\noindent\emph{Ukrajinsky / Ukrainian:} тіло автоморфізмів\par
\smallskip\noindent\emph{Razlog:} Tělo avtomorfizmov razlikuje možlivo nekomutativno tělo od komutativnogo polja. Německo Körper v tom kontekstu ne opravdava avtomatičnu zaměnu s polje. Odlučujuči razlog jest algebraična struktura; podobnost slov v ruskoj i ukrajinskoj paralelě podpiraje imenovanje, ale ne zaměnjaje sravnjenje matematičnyh uslovij.\par
\noindent\emph{Reason:} Tělo avtomorfizmov distinguishes a potentially noncommutative division ring from a commutative field. Körper in this context does not license automatic substitution of polje. The decisive issue is the algebraic structure. Russian and Ukrainian parallels support the naming choice but do not replace examination of the mathematical assumptions.\par
\noindent{\footnotesize Izvorne zapisy / Original records: HIST-g177-007, HIST-g197-005.}\par
\end{minipage}\par
\par\medskip\Needspace{7\baselineskip}
\noindent\begin{minipage}{\linewidth}
\hypertarget{TERM-0158}{}
\noindent\textbf{tělo avtomorfizmov $K_\nu$}\quad{\footnotesize\texttt{TERM-0158}}\par
\noindent\emph{Iztočnik / Source:} Automorphismenkörper $K_\nu$\par
\noindent\emph{Rusky / Russian:} тело автоморфизмов $K_\nu$\par
\noindent\emph{Ukrajinsky / Ukrainian:} тіло автоморфізмів $K_\nu$\par
\smallskip\noindent\emph{Razlog:} Tělo avtomorfizmov $K_\nu$ zachovuje razliku medžu možlivo nekomutativnym tělom i komutativnym poljem koeficientov. Oznaka $K_\nu$ ostavaje točno svezana s avtorskym objektom. Izbor tělo ne jest prosto alternativnym nazvanjem za polje: on izběgaje dodavanje uslovja komutativnosti, ktore danomu objektu ne bylo pripisano.\par
\noindent\emph{Reason:} Tělo avtomorfizmov $K_\nu$ preserves the distinction between a potentially noncommutative division ring and the commutative coefficient field. $K_\nu$ remains attached to the author's object. Tělo is not merely a stylistic alternative to polje: it avoids imposing commutativity where that assumption has not been made.\par
\noindent{\footnotesize Izvorne zapisy / Original records: HIST-g198-003.}\par
\end{minipage}\par
\par\medskip\Needspace{7\baselineskip}
\noindent\begin{minipage}{\linewidth}
\hypertarget{TERM-0159}{}
\noindent\textbf{kolco avtomorfizmov}\quad{\footnotesize\texttt{TERM-0159}}\par
\noindent\emph{Iztočnik / Source:} Automorphismenring\par
\noindent\emph{Rusky / Russian:} кольцо автоморфизмов\par
\noindent\emph{Ukrajinsky / Ukrainian:} кільце автоморфізмів\par
\smallskip\noindent\emph{Razlog:} Avtomorfizmov zachovuje avtorske ime Automorphismenring daže tam, kde konstrukcija jest opisana kako sistem endomorfizmov. Slědujuči teorem utočnjaje, že nenulove endomorfizmy sut avtomorfizmy; prěvod ne smije vnositi taj rezultat v prědhodnu definiciju. Kolco jest forma ranějšego kandidata; aktualno terminologično pravilo prsten trěbuje odvojeno sravnjenje konkretnyh pojavjenij.\par
\noindent\emph{Reason:} Avtomorfizmov preserves the author's label Automorphismenring even where the construction is described as a system of endomorphisms. The following theorem establishes that the nonzero endomorphisms are automorphisms; the translation must not insert that result into the preceding definition. Kolco is the historical candidate's form, while the current prsten policy requires separate comparison with actual occurrences.\par
\noindent{\footnotesize Izvorne zapisy / Original records: HIST-g179-006, HIST-g180-005.}\par
\end{minipage}\par
\par\medskip\Needspace{7\baselineskip}
\noindent\begin{minipage}{\linewidth}
\hypertarget{TERM-0160}{}
\noindent\textbf{kolco avtomorfizmov / tělo avtomorfizmov}\quad{\footnotesize\texttt{TERM-0160}}\par
\noindent\emph{Iztočnik / Source:} Automorphismenring / Automorphismenkörper\par
\noindent\emph{Rusky / Russian:} кольцо автоморфизмов / тело автоморфизмов\par
\noindent\emph{Ukrajinsky / Ukrainian:} кільце автоморфізмів / тіло автоморфізмів\par
\smallskip\noindent\emph{Razlog:} Para kolco avtomorfizmov / tělo avtomorfizmov zachovuje razliku medžu prstenom i osobnym slučajem těla v dokazě matričnogo predstavjenja po Wedderburnu. Izravnanje obojih nazvanij by skrylo matematičny krok. Istorično kolco jest dokumentovano kako zapis kandidata, ne kako nova odluka protiv termina prsten.\par
\noindent\emph{Reason:} The pair kolco avtomorfizmov / tělo avtomorfizmov preserves the distinction between a ring and the division-ring specialization in the Wedderburn matrix-representation argument. Flattening the two names would hide a mathematical step. Historical kolco is documented as the candidate's wording, not as a new decision against prsten.\par
\noindent{\footnotesize Izvorne zapisy / Original records: HIST-g178-005, HIST-g192-003.}\par
\end{minipage}\par
\par\medskip\Needspace{7\baselineskip}
\noindent\begin{minipage}{\linewidth}
\hypertarget{TERM-0161}{}
\noindent\textbf{baza direktnoj sumy / blokovo-diagonalna diskriminantna matrica}\quad{\footnotesize\texttt{TERM-0161}}\par
\noindent\emph{Iztočnik / Source:} Basis der direkten Summe / blockdiagonale Diskriminantenmatrix\par
\noindent\emph{Rusky / Russian:} база прямой суммы / блочно-диагональная дискриминантная матрица\par
\noindent\emph{Ukrajinsky / Ukrainian:} база прямої суми / блоково-діагональна дискримінантна матриця\par
\smallskip\noindent\emph{Razlog:} Baza direktnoj sumy i blokovo-diagonalna diskriminantna matrica imenujut svezane, ale različne objekty dokaza. Matrica opiraje se na izbranu bazu; početno oznaka $S_{(a)}$ nastupaje v obojih blokah, poslě čego slědy komponentov sut identifikovane. Prěvod ne trěbuje prědčasno zaměniti avtorske oznaky gotovym zaključkom o komponentah.\par
\noindent\emph{Reason:} Baza direktnoj sumy and blokovo-diagonalna diskriminantna matrica name related but distinct objects in the proof. The matrix depends on the chosen basis. The notation initially uses $S_{(a)}$ in both blocks before identifying the component traces. The translation should not replace the author's notation prematurely with the conclusion about the components.\par
\noindent{\footnotesize Izvorne zapisy / Original records: HIST-g139-006.}\par
\end{minipage}\par
\par\medskip\Needspace{7\baselineskip}
\noindent\begin{minipage}{\linewidth}
\hypertarget{TERM-0162}{}
\noindent\textbf{element, ktory poraždaje izvedeny prosty ideal}\quad{\footnotesize\texttt{TERM-0162}}\par
\noindent\emph{Iztočnik / Source:} Basis des abgeleiteten Primideals\par
\noindent\emph{Rusky / Russian:} элемент, порождающий выведенный простой идеал\par
\noindent\emph{Ukrajinsky / Ukrainian:} елемент, що породжує виведений простий ідеал\par
\smallskip\noindent\emph{Razlog:} V danom prstenu glavnyh idealov Basis označaje poraždajuči element ideala. Zato element, ktory poraždaje izvedeny prosty ideal izjasnjaje funkciju objekta i svez s prědhodnym zaključkom. To jest kontekstualno razkrytje smysla, ne pravilo, že vsaka matematična baza jest jedinym elementom.\par
\noindent\emph{Reason:} In this principal ideal ring, Basis denotes an ideal generator. Element, ktory poraždaje izvedeny prosty ideal makes the object's function and its connection with the preceding conclusion explicit. This is contextual disambiguation, not a rule that every mathematical basis is a single element.\par
\noindent{\footnotesize Izvorne zapisy / Original records: HIST-g160-006.}\par
\end{minipage}\par
\par\medskip\Needspace{7\baselineskip}
\noindent\begin{minipage}{\linewidth}
\hypertarget{TERM-0163}{}
\noindent\textbf{baza razširjenogo ideala diskriminantnogo ideala}\quad{\footnotesize\texttt{TERM-0163}}\par
\noindent\emph{Iztočnik / Source:} Basis des Erweiterungsideals des Diskriminantenideals\par
\noindent\emph{Rusky / Russian:} базис идеала расширения дискриминантного идеала\par
\noindent\emph{Ukrajinsky / Ukrainian:} базис ідеалу розширення дискримінантного ідеалу\par
\smallskip\noindent\emph{Razlog:} Baza razširjenogo ideala diskriminantnogo ideala zachovuje vse stupeni imennikovej konstrukcije: govori se o bazě ideala dobivenogo razširjenjem diskriminantnogo ideala. V lokalnom prstenu toj ideal jest glavny, zato poraždajuči element mogl by razjasniti funkciju bazy. Izvorna odluka odnak zachovuje baza; takvo objasnjenje ne trěba predstavljati kako už priměnjenu zaměnu.\par
\noindent\emph{Reason:} Baza razširjenogo ideala diskriminantnogo ideala retains the nested construction: the basis belongs to the extension of the discriminant ideal. In the local ring that ideal is principal, so generator could clarify the function of Basis. The recorded choice nevertheless retains baza; the explanatory possibility must not be represented as an already adopted replacement.\par
\noindent{\footnotesize Izvorne zapisy / Original records: HIST-g162-002.}\par
\end{minipage}\par
\par\medskip\Needspace{7\baselineskip}
\noindent\begin{minipage}{\linewidth}
\hypertarget{TERM-0164}{}
\noindent\textbf{baza za tvorjenje slěda}\quad{\footnotesize\texttt{TERM-0164}}\par
\noindent\emph{Iztočnik / Source:} Basis zur Erzeugung der Spur\par
\noindent\emph{Rusky / Russian:} базис для построения следа\par
\noindent\emph{Ukrajinsky / Ukrainian:} базис для побудови сліду\par
\smallskip\noindent\emph{Razlog:} Tvorjenje slěda govori o konstrukciji slěda s pomočju bazy. Izbor tvorjenje izběgaje avtomatično upotrěbjenje poraždati, ktore v susědnem izlaganju ima osobnu funkciju pri idealah. Dve operacije ne dobivajut isto nazvanje samo zato, že německo Erzeugung može nastupati pri obojih.\par
\noindent\emph{Reason:} Tvorjenje slěda refers to constructing the trace using a basis. Tvorjenje avoids automatically reusing poraždati, which has a particular role in the nearby ideal-theoretic discussion. The two operations do not receive the same name merely because German Erzeugung can be used for both.\par
\noindent{\footnotesize Izvorne zapisy / Original records: HIST-g157-005.}\par
\end{minipage}\par
\par\medskip\Needspace{7\baselineskip}
\noindent\begin{minipage}{\linewidth}
\hypertarget{TERM-0165}{}
\noindent\textbf{pojmy bazy}\quad{\footnotesize\texttt{TERM-0165}}\par
\noindent\emph{Iztočnik / Source:} Basisbegriffe\par
\noindent\emph{Rusky / Russian:} понятия базы\par
\noindent\emph{Ukrajinsky / Ukrainian:} поняття бази\par
\smallskip\noindent\emph{Razlog:} Pojmy bazy zachovuje množinu različnyh načinov razuměnja bazy, ktore tekst organizuje i sravnjaje. Izraz ne svodi razmatranje do jednoj sovrěmennoj definicije bazy. Množina pojmy jest zato sadržajna, ne slučajna gramatična forma.\par
\noindent\emph{Reason:} Pojmy bazy preserves the plurality of basis notions organized and compared in the text. It does not collapse the discussion into one modern definition of a basis. The plural notions is substantive, not an incidental grammatical form.\par
\noindent{\footnotesize Izvorne zapisy / Original records: HIST-g044-002.}\par
\end{minipage}\par
\par\medskip\Needspace{7\baselineskip}
\noindent\begin{minipage}{\linewidth}
\hypertarget{TERM-0166}{}
\noindent\textbf{bazovo svojstvo}\quad{\footnotesize\texttt{TERM-0166}}\par
\noindent\emph{Iztočnik / Source:} Basiseigenschaft\par
\noindent\emph{Rusky / Russian:} базовое свойство\par
\noindent\emph{Ukrajinsky / Ukrainian:} базова властивість\par
\smallskip\noindent\emph{Razlog:} Bazovo svojstvo upučaje na svojstva 1–5, ktore sut dane v uvodu pri opisu bazy. Bazovo zato trěba čitati v svezě s bazoju, ne kako svobodno tvrđenje, že svojstvo jest fundamentalno ili posebno važno. Točna lokalna referencija razděljuje tehničny i običny smysl pridavnika.\par
\noindent\emph{Reason:} Bazovo svojstvo refers to properties 1–5 used in the introduction to characterize the basis. Bazovo is therefore connected with a basis, not a free assertion that the property is fundamental or especially important. The local reference distinguishes the technical from the everyday sense of the adjective.\par
\noindent{\footnotesize Izvorne zapisy / Original records: HIST-g059-001.}\par
\end{minipage}\par
\par\medskip\Needspace{7\baselineskip}
\noindent\begin{minipage}{\linewidth}
\hypertarget{TERM-0167}{}
\noindent\textbf{pitanja o bazah}\quad{\footnotesize\texttt{TERM-0167}}\par
\noindent\emph{Iztočnik / Source:} Basisfragen\par
\noindent\emph{Rusky / Russian:} вопросы о базисах\par
\noindent\emph{Ukrajinsky / Ukrainian:} питання про бази\par
\smallskip\noindent\emph{Razlog:} Pitanja o bazah zachovuje otvorenu formu matematičnyh pytanj i problemov. Izraz ne naziva jeden konkretni teorem i ne tvrdi, že vse tute pytanja sut rěšene. Predložna konstrukcija jasno dava predmet razmatranja bez stvarjanja novogo složenogo termina.\par
\noindent\emph{Reason:} Pitanja o bazah retains the open-ended reference to mathematical questions and problems about bases. It neither names a particular theorem nor claims that all those questions have been solved. The prepositional construction makes the subject explicit without creating a new compound technical term.\par
\noindent{\footnotesize Izvorne zapisy / Original records: HIST-g040-002.}\par
\end{minipage}\par
\par\medskip\Needspace{7\baselineskip}
\noindent\begin{minipage}{\linewidth}
\hypertarget{TERM-0168}{}
\noindent\textbf{bezbazisna forma}\quad{\footnotesize\texttt{TERM-0168}}\par
\noindent\emph{Iztočnik / Source:} basisfreie Form\par
\noindent\emph{Rusky / Russian:} безбазисная форма\par
\noindent\emph{Ukrajinsky / Ukrainian:} безбазисна форма\par
\smallskip\noindent\emph{Razlog:} Bezbazisna forma označaje formulaciju, ktora ne trěbuje izričnogo opiranja na konečnu bazu ideala. Zapis veže ju s možnostju formulacije črěz teorem o konečnosti lancov. Termin ne govori, že razmatrany objekt nema bazy; govori o načinu izlaganja tvrđenja.\par
\noindent\emph{Reason:} Bezbazisna forma names a formulation that does not explicitly rely on a finite ideal basis. The record relates it to a formulation through the finite-chain theorem. The term does not say that the object has no basis; it describes the way the assertion is expressed.\par
\noindent{\footnotesize Izvorne zapisy / Original records: HIST-g081-029, HIST-g082-029.}\par
\end{minipage}\par
\par\medskip\Needspace{7\baselineskip}
\noindent\begin{minipage}{\linewidth}
\hypertarget{TERM-0169}{}
\noindent\textbf{bazisna funkcija}\quad{\footnotesize\texttt{TERM-0169}}\par
\noindent\emph{Iztočnik / Source:} Basisfunktion\par
\noindent\emph{Rusky / Russian:} базисная функция\par
\noindent\emph{Ukrajinsky / Ukrainian:} базисна функція\par
\smallskip\noindent\emph{Razlog:} Bazisna funkcija imenuje funkciju, ktora nastupaje v konstrukciji determinanty $\Delta$ i kvadratnyh form v modulu $M$. Izbor zachovuje vezu s bazoju, ne zaměnjaje funkciju neoprěděljenym elementom. Konkretna upotrěba jest važnějša od samej podobnosti internacionalnyh korenj.\par
\noindent\emph{Reason:} Bazisna funkcija names a function used in constructing the determinant $\Delta$ and the quadratic forms in the module $M$. The choice preserves the connection with a basis without replacing function by an unspecified element. The particular use matters more than resemblance between international roots.\par
\noindent{\footnotesize Izvorne zapisy / Original records: HIST-g078-008.}\par
\end{minipage}\par
\par\medskip\Needspace{7\baselineskip}
\noindent\begin{minipage}{\linewidth}
\hypertarget{TERM-0170}{}
\noindent\textbf{bazova veličina}\quad{\footnotesize\texttt{TERM-0170}}\par
\noindent\emph{Iztočnik / Source:} Basisgröße\par
\noindent\emph{Rusky / Russian:} базовая величина\par
\noindent\emph{Ukrajinsky / Ukrainian:} базова величина\par
\smallskip\noindent\emph{Razlog:} Bazova veličina zachovuje avtorske različenje medžu veličinoju i čislom v priměru zaměny. Zapis prodolžaje upotrěbjenje iz devetoj raboty i ne modernizuje avtomatično veličina do element. Takva modernizacija mogla by skryti to, kako sam avtor organizuje svoje objekty.\par
\noindent\emph{Reason:} Bazova veličina preserves the author's distinction between a quantity and a number in the substitution example. It continues the usage of the ninth paper rather than automatically modernizing quantity to element. That modernization could hide the author's own organization of the objects.\par
\noindent{\footnotesize Izvorne zapisy / Original records: HIST-g059-002, HIST-g061-005.}\par
\end{minipage}\par
\par\medskip\Needspace{7\baselineskip}
\noindent\begin{minipage}{\linewidth}
\hypertarget{TERM-0171}{}
\noindent\textbf{teoremy o bazisu / fundamentalna sistema}\quad{\footnotesize\texttt{TERM-0171}}\par
\noindent\emph{Iztočnik / Source:} Basissätze / Fundamentalsystem\par
\noindent\emph{Rusky / Russian:} теоремы о базисе / фундаментальная система\par
\noindent\emph{Ukrajinsky / Ukrainian:} теореми про базис / фундаментальна система\par
\smallskip\noindent\emph{Razlog:} Teoremy o bazisu i fundamentalna sistema zachovujut dva različne, hoti svezane izrazy. Prvi imenuje tvrđenja o bazah; vtory označaje sistem, ktory se pojavjaje pri izvodženju cěločislenyh baz iz modulnyh argumentov. Skupny zapis ne znači, že oba nazvanja sut sinonimy.\par
\noindent\emph{Reason:} Teoremy o bazisu and fundamentalna sistema retain two related but distinct expressions. The first names basis theorems; the second names the system occurring in the derivation of integral bases from module arguments. Their inclusion in one entry does not make them synonyms.\par
\noindent{\footnotesize Izvorne zapisy / Original records: HIST-g076-007.}\par
\end{minipage}\par
\par\medskip\Needspace{7\baselineskip}
\noindent\begin{minipage}{\linewidth}
\hypertarget{TERM-0172}{}
\noindent\textbf{uslovja vyšljego defekta}\quad{\footnotesize\texttt{TERM-0172}}\par
\noindent\emph{Iztočnik / Source:} Bedingungen von höherem Defekt\par
\noindent\emph{Rusky / Russian:} условия более высокого дефекта\par
\noindent\emph{Ukrajinsky / Ukrainian:} умови вищого дефекту\par
\smallskip\noindent\emph{Razlog:} Uslovja vyšljego defekta zachovuje porjadok po defektu v razmatranju redukcije pri neparnom defektu. Vyšljego ne govori o silě ili kvalitetě uslovij v običnom smyslu. Pridavnik upučaje na konkretny parametr, od ktorego zavisi zaključok.\par
\noindent\emph{Reason:} Uslovja vyšljego defekta retains the ordering by defect in the reduction argument for odd defect. Higher does not refer to the strength or quality of the conditions in an everyday sense. It identifies the particular parameter on which the conclusion depends.\par
\noindent{\footnotesize Izvorne zapisy / Original records: HIST-g034-012.}\par
\end{minipage}\par
\par\medskip\Needspace{7\baselineskip}
\noindent\begin{minipage}{\linewidth}
\hypertarget{TERM-0173}{}
\noindent\textbf{uslovne uravnenja}\quad{\footnotesize\texttt{TERM-0173}}\par
\noindent\emph{Iztočnik / Source:} Bedingungsgleichungen\par
\noindent\emph{Rusky / Russian:} уравнения условий\par
\noindent\emph{Ukrajinsky / Ukrainian:} рівняння умов\par
\smallskip\noindent\emph{Razlog:} Uslovne uravnenja označaje jednačenja, ktore služe kako uslovja. Pridavnik ne znači, že njih istinnost jest samo hipotetična ili proizvoljno dopustima. U zapisu jest sučstven odnos medžu jednačenjem i uslovjem; ranějša forma uravnenja sama ne rěšaje izbor aktualnogo medžuslovjanskogo termina za jednačenje.\par
\noindent\emph{Reason:} Uslovne uravnenja denotes equations that serve as conditions. The adjective does not mean that their truth is merely hypothetical or discretionary. The recorded point is the relation between an equation and a condition; historical uravnenja does not by itself settle the current Interslavic term for equation.\par
\noindent{\footnotesize Izvorne zapisy / Original records: HIST-g030-006.}\par
\end{minipage}\par
\par\medskip\Needspace{7\baselineskip}
\noindent\begin{minipage}{\linewidth}
\hypertarget{TERM-0174}{}
\noindent\textbf{ovladanje tym svęzom medžu formami}\quad{\footnotesize\texttt{TERM-0174}}\par
\noindent\emph{Iztočnik / Source:} Beherrschung dieses Formenzusammenhangs\par
\noindent\emph{Rusky / Russian:} овладение этой связью между формами\par
\noindent\emph{Ukrajinsky / Ukrainian:} опанування цього зв'язку між формами\par
\smallskip\noindent\emph{Razlog:} Ovladanje tym svęzom medžu formami prědaje Beherrschung kako pojmově ovladanje odnosom medžu formami. To jest metodologična proza, ne tvrđenje o čislenom ograničenju. Etimologičny napis svęzom pripada citovanoj kandidatnoj formě; on ne ustanavjaje obvezno upotrěbjenje ę v osnovnom pravopisě čitatelskogo teksta.\par
\noindent\emph{Reason:} Ovladanje tym svęzom medžu formami renders Beherrschung as conceptual command of the relation between forms. This is methodological prose, not an assertion of a numerical bound. Etymological svęzom belongs to the quoted candidate; it does not require ę throughout the reader's basic-spelling text.\par
\noindent{\footnotesize Izvorne zapisy / Original records: HIST-g029-005.}\par
\end{minipage}\par
\par\medskip\Needspace{7\baselineskip}
\noindent\begin{minipage}{\linewidth}
\hypertarget{TERM-0175}{}
\noindent\textbf{libovoljno dobro aproksimovati}\quad{\footnotesize\texttt{TERM-0175}}\par
\noindent\emph{Iztočnik / Source:} beliebig gut approximieren\par
\noindent\emph{Rusky / Russian:} сколь угодно хорошо аппроксимировать\par
\noindent\emph{Ukrajinsky / Ukrainian:} як завгодно добре апроксимувати\par
\smallskip\noindent\emph{Razlog:} Libovoljno dobro aproksimovati zachovuje neograničenu točnost približenja, važnu v argumentu o gustotě svezanom s Hamelom. Libovoljno dobro ne znači samo dovoljno dobro za neko jednu cělj. Zapis ks v aproksimovati dava pravilno pisemno odgovaranje s kiriličnym кс; pravopisna odluka ne zaměnjaje matematičny argument o približenju.\par
\noindent\emph{Reason:} Libovoljno dobro aproksimovati retains arbitrarily good approximation in the density argument associated with Hamel. It does not mean merely good enough for one purpose. Writing ks in aproksimovati gives the corresponding Cyrillic кс; that spelling decision is separate from the mathematical approximation argument.\par
\noindent{\footnotesize Izvorne zapisy / Original records: HIST-g072-013.}\par
\end{minipage}\par
\par\medskip\Needspace{7\baselineskip}
\noindent\begin{minipage}{\linewidth}
\hypertarget{TERM-0176}{}
\noindent\textbf{proizvoljne formy}\quad{\footnotesize\texttt{TERM-0176}}\par
\noindent\emph{Iztočnik / Source:} beliebige Formen\par
\noindent\emph{Rusky / Russian:} произвольные формы\par
\noindent\emph{Ukrajinsky / Ukrainian:} довільні форми\par
\smallskip\noindent\emph{Razlog:} Proizvoljne formy zachovuje proizvoljnost izbranyh objektov v dokazě. Zapis sravnjaje proizvoljny s libokaky i dava prednost prvomu radi formalnogo matematičnogo registra. To jest stilistično obrazloženje, ne izměrena prednost razumělivosti niti tvrđenje, že libokaky ne može byti upotrěbjen v drugom kontekstu.\par
\noindent\emph{Reason:} Proizvoljne formy preserves the arbitrariness of the objects in the proof. The record compares proizvoljny with libokaky and prefers the former for formal mathematical register. This is a register argument, not a measured intelligibility advantage or a claim that libokaky is unusable elsewhere.\par
\noindent{\footnotesize Izvorne zapisy / Original records: HIST-g019-013.}\par
\end{minipage}\par
\par\medskip\Needspace{7\baselineskip}
\noindent\begin{minipage}{\linewidth}
\hypertarget{TERM-0177}{}
\noindent\textbf{libovoljna oblast racionalnosti}\quad{\footnotesize\texttt{TERM-0177}}\par
\noindent\emph{Iztočnik / Source:} beliebiger Rationalitätsbereich\par
\noindent\emph{Rusky / Russian:} произвольная область рациональности\par
\noindent\emph{Ukrajinsky / Ukrainian:} довільна область раціональності\par
\smallskip\noindent\emph{Razlog:} Libovoljna oblast racionalnosti prodolžaje istoričnu terminologiju iz raboty 11. Opseg oblasti opisan jest avtorskym utočnjenjem črěz konečno algebraično čislovo polje, moguće konečno čislo pridanyh parametrov i medžupolja. Konečno v tom opisu ne trěba čitati kako tvrđenje o konečnom čislu elementov polja. Izraz oblast ne smije samostojno zaměniti vse te uslovja.\par
\noindent\emph{Reason:} Libovoljna oblast racionalnosti continues the historical terminology of the eleventh paper. The domain's scope is governed by the author's qualification involving a finite algebraic number field, possibly finitely many adjoined parameters, and intermediate fields. Finite here must not be misread as saying that the field has finitely many elements. The word domain cannot replace those conditions.\par
\noindent{\footnotesize Izvorne zapisy / Original records: HIST-g078-016.}\par
\end{minipage}\par
\par\medskip\Needspace{7\baselineskip}
\noindent\begin{minipage}{\linewidth}
\hypertarget{TERM-0178}{}
\noindent\textbf{libovolny sistem}\quad{\footnotesize\texttt{TERM-0178}}\par
\noindent\emph{Iztočnik / Source:} beliebiges System\par
\noindent\emph{Rusky / Russian:} произвольная система\par
\noindent\emph{Ukrajinsky / Ukrainian:} довільна система\par
\smallskip\noindent\emph{Razlog:} Libovolny sistem označaje proizvoljnu zbirku prěd uvedenjem operacij, pod ktorymi ona može byti zamknjena. Prěvod ne trěbuje unapred pripisati sistemu algebraičnu strukturu, ktora se teprvo razmatraje. Izbor zachovuje tu logičnu poslědovatelnost; sam napis libovolny ostavaje dokumentovanoju formoju kandidata.\par
\noindent\emph{Reason:} Libovolny sistem denotes an arbitrary collection before the operations under which it may be closed are introduced. The translation must not assume an algebraic structure that the argument is only about to consider. The choice preserves that logical order, while libovolny remains the candidate's documented spelling.\par
\noindent{\footnotesize Izvorne zapisy / Original records: HIST-g047-001.}\par
\end{minipage}\par
\par\medskip\Needspace{7\baselineskip}
\noindent\begin{minipage}{\linewidth}
\hypertarget{TERM-0179}{}
\noindent\textbf{oblast / čislova oblast}\quad{\footnotesize\texttt{TERM-0179}}\par
\noindent\emph{Iztočnik / Source:} Bereich / Zahlbereich\par
\noindent\emph{Rusky / Russian:} область / числовая область\par
\noindent\emph{Ukrajinsky / Ukrainian:} область / числова область\par
\smallskip\noindent\emph{Razlog:} Oblast i čislova oblast zachovujut široko istorično upotrěbjenje Bereich i Zahlbereich. Čislova dodavaje uslovje, ktore jest prědstavjeno v drugom německom slově; oba izrazy ne sut avtomatično imenami prstenov. Prednost doslovnoj širini ovde izběgaje tiho dodavanje sovrěmennoj struktury, ktora ne jest v každom pojavjenju dana.\par
\noindent\emph{Reason:} Oblast and čislova oblast preserve the broad historical usage of Bereich and Zahlbereich. Čislova adds the qualification present in the second German word; neither expression is automatically the name of a ring. Preserving that breadth avoids silently imposing a modern structure not specified in every occurrence.\par
\noindent{\footnotesize Izvorne zapisy / Original records: HIST-g058-001.}\par
\end{minipage}\par
\par\medskip\Needspace{7\baselineskip}
\noindent\begin{minipage}{\linewidth}
\hypertarget{TERM-0180}{}
\noindent\textbf{oblast vseh parnyh polinomov}\quad{\footnotesize\texttt{TERM-0180}}\par
\noindent\emph{Iztočnik / Source:} Bereich aller geraden Polynome\par
\noindent\emph{Rusky / Russian:} область всех четных полиномов\par
\noindent\emph{Ukrajinsky / Ukrainian:} область усіх парних поліномів\par
\smallskip\noindent\emph{Razlog:} Parnyh jest sučstvena čest opisa oblasti, v ktorej ravnost $x^2y^2=(xy)^2$ dava priměr protiv jednoznačnosti produkta nerazložimyh polinomov. Prěvod ne zaměnjaje tu osobnu oblast neoprěděljenym podprstenom. Iz samogo zaglavja ne trěba dodavati novu definiciju parnosti po každoj proměnnoj; takvo čitanje trěbovalo by sravnjenje s cělym priměrom.\par
\noindent\emph{Reason:} Parnyh is essential to the description of the domain in which $x^2y^2=(xy)^2$ supplies the counterexample to unique factorization into irreducible polynomials. The translation does not replace that particular domain with an unspecified subring. The headword alone should not be expanded into a definition requiring evenness in each variable; that reading would need comparison with the complete example.\par
\noindent{\footnotesize Izvorne zapisy / Original records: HIST-g081-086.}\par
\end{minipage}\par
\par\medskip\Needspace{7\baselineskip}
\noindent\begin{minipage}{\linewidth}
\hypertarget{TERM-0181}{}
\noindent\textbf{oblast vseh homogenyh form}\quad{\footnotesize\texttt{TERM-0181}}\par
\noindent\emph{Iztočnik / Source:} Bereich aller homogenen Formen\par
\noindent\emph{Rusky / Russian:} область всех однородных форм\par
\noindent\emph{Ukrajinsky / Ukrainian:} область усіх однорідних форм\par
\smallskip\noindent\emph{Razlog:} Homogenyh form zachovuje vrstu objektov, za ktore Noether priměnjaje obče teoremy. V tom priměru dodavanje opreděljeno jest samo za formy istogo stepena. Oblast zato ne trěba bez objasnjenja zaměniti sovrěmennym nazvanjem struktury s neograničenym dodavanjem. Medžunarodne osnovy homogen- i form- podpirajut doslědno imenovanje tih objektov.\par
\noindent\emph{Reason:} Homogenyh form retains the objects to which Noether applies the general theorems. In this example addition is defined only for forms of equal degree. Oblast must therefore not be replaced without explanation by a modern structure name implying unrestricted addition. The international stems homogen- and form- provide consistent names for the objects.\par
\noindent{\footnotesize Izvorne zapisy / Original records: HIST-g081-085.}\par
\end{minipage}\par
\par\medskip\Needspace{7\baselineskip}
\noindent\begin{minipage}{\linewidth}
\hypertarget{TERM-0182}{}
\noindent\textbf{oblast H iz algebraično-cělyh transcendentnyh čisel}\quad{\footnotesize\texttt{TERM-0182}}\par
\noindent\emph{Iztočnik / Source:} Bereich H aus algebraisch-ganzen transzendenten Zahlen\par
\noindent\emph{Rusky / Russian:} область H из алгебраически целых трансцендентных чисел\par
\noindent\emph{Ukrajinsky / Ukrainian:} область H з алгебраїчно-цілих трансцендентних чисел\par
\smallskip\noindent\emph{Razlog:} Izraz zachovuje vse kvalifikacije nazvanja oblasti $H$: algebraično-cělyh i transcendentnyh. Na prvi pogled neobičnu kombinaciju ne trěba tiho uproščati do običnogo čislovogo polja. Jej smysl zavisi od avtorske konstrukcije oblasti; $H$ ostavaje matematičnoju oznakoju, ne slovom za prěvod.\par
\noindent\emph{Reason:} The phrase retains all the qualifications in the name of the domain $H$: algebraically integral and transcendental. Their apparently unusual combination should not be silently simplified to an ordinary number field. Its meaning depends on the author's construction of the domain, and $H$ remains a mathematical label rather than a word to translate.\par
\noindent{\footnotesize Izvorne zapisy / Original records: HIST-g062-005.}\par
\end{minipage}\par
\par\medskip\Needspace{7\baselineskip}
\noindent\begin{minipage}{\linewidth}
\hypertarget{TERM-0183}{}
\noindent\textbf{oblasti iz c?lyh transcendentnyh ?isel}\quad{\footnotesize\texttt{TERM-0183}}\par
\noindent\emph{Iztočnik / Source:} Bereiche aus ganzen transzendenten Zahlen\par
\noindent\emph{Rusky / Russian:} [Bukvy ne sut zachovane / Lettering not recoverable]\par
\noindent\emph{Ukrajinsky / Ukrainian:} [Bukvy ne sut zachovane / Lettering not recoverable]\par
\smallskip\noindent\emph{Razlog:} Izvorni razlog dava prednost istoričnomu oblast vměsto avtomatičnogo sovrěmennogo prsten i veže terminologiju cělosti s ranějšim razmatranjem Zermelovy oblasti v rabotě 9. Ale medžuslovjansky zapis ima izgubljene bukvy, a v ruskom i ukrajinskom polju bukvy sut zaměnjene znakom ?. Zato poznajemo razlog izbora, ale ne možemo iz toga zapisa ustanoviti točnu grafiju vseh trěh prěvodov.\par
\noindent\emph{Reason:} The original rationale prefers historical oblast to an automatic modernization as ring and connects the integral/whole terminology with the earlier Zermelo-domain discussion in Paper 9. However, letters are lost in the Interslavic form, and the Russian and Ukrainian fields are damaged by question marks. The reason survives, but the exact spelling of all three translations cannot be established from this record.\par
\noindent{\footnotesize Izvorne zapisy / Original records: HIST-g063-002.}\par
\end{minipage}\par
\par\medskip\Needspace{7\baselineskip}
\noindent\begin{minipage}{\linewidth}
\hypertarget{TERM-0184}{}
\noindent\textbf{ograničeny po $x_i$ / ne dosega stepena k}\quad{\footnotesize\texttt{TERM-0184}}\par
\noindent\emph{Iztočnik / Source:} beschränkt in $x_i$ / den Grad k nicht erreichen\par
\noindent\emph{Rusky / Russian:} ограниченный по $x_i$ / не достигает степени k\par
\noindent\emph{Ukrajinsky / Ukrainian:} обмежений за $x_i$ / не досягає степеня k\par
\smallskip\noindent\emph{Razlog:} Ograničeny po $x_i$ govori o ograničenju stepena predstavitelja v toj proměnnoj. Ne dosega stepena $k$ zachovuje strogu granicu: stepen jest menši od $k$, ne samo menši ili ravny. Tuto jest vlastnost, ktora se zachovuje v dokazě, ne neoprěděljeno kvalitativno ograničenje.\par
\noindent\emph{Reason:} Ograničeny po $x_i$ concerns the degree of the representative in that variable. Ne dosega stepena $k$ preserves a strict bound: the degree is below $k$, not merely at most $k$. This is a property maintained in the proof, not an unspecified qualitative restriction.\par
\noindent{\footnotesize Izvorne zapisy / Original records: HIST-g088-005, HIST-g092-036.}\par
\end{minipage}\par
\par\medskip\Needspace{7\baselineskip}
\noindent\begin{minipage}{\linewidth}
\hypertarget{TERM-0185}{}
\noindent\textbf{označenja i opreděljenja}\quad{\footnotesize\texttt{TERM-0185}}\par
\noindent\emph{Iztočnik / Source:} Bezeichnungen und Definitionen\par
\noindent\emph{Rusky / Russian:} обозначения и определения\par
\noindent\emph{Ukrajinsky / Ukrainian:} позначення й означення\par
\smallskip\noindent\emph{Razlog:} Označenja i opreděljenja razděljuje imena ili znaky od definicij pojmov. Za naslov odděla izbrany sut obične matematične imenniki, bez naměrnoj arhaizacije. Takvy naslov pomaga čitatelju razuměti funkciju odděla prěd čitanjem njegovih podrobnostij.\par
\noindent\emph{Reason:} Označenja i opreděljenja distinguishes names or notation from definitions of concepts. The section title uses ordinary mathematical nouns without deliberate archaism, making the section's purpose clear before its details are read.\par
\noindent{\footnotesize Izvorne zapisy / Original records: HIST-g030-001.}\par
\end{minipage}\par
\par\medskip\Needspace{7\baselineskip}
\noindent\begin{minipage}{\linewidth}
\hypertarget{TERM-0186}{}
\noindent\textbf{bibliografija}\quad{\footnotesize\texttt{TERM-0186}}\par
\noindent\emph{Iztočnik / Source:} Bibliographie\par
\noindent\emph{Rusky / Russian:} библиография\par
\noindent\emph{Ukrajinsky / Ukrainian:} бібліографія\par
\smallskip\noindent\emph{Razlog:} Bibliografija imenuje funkciju završnogo spiska citovanyh děl. Izbor medžunarodnogo imennika odgovarja ruskoj i ukrajinskoj paralelě i ne trěbuje dolgu opisnu zaměnu. Izvorni razlog govori o izběganju dvojsmyslenosti, ne o posebnom slovotvornom novotvoru.\par
\noindent\emph{Reason:} Bibliografija identifies the function of the concluding list of cited works. The international noun matches the Russian and Ukrainian parallels without a lengthy descriptive substitute. The recorded reason concerns avoiding ambiguity, not creating a novel technical word.\par
\noindent{\footnotesize Izvorne zapisy / Original records: HIST-g215-001.}\par
\end{minipage}\par
\par\medskip\Needspace{7\baselineskip}
\noindent\begin{minipage}{\linewidth}
\hypertarget{TERM-0187}{}
\noindent\textbf{tvorjut Z-modul / modulno svojstvo}\quad{\footnotesize\texttt{TERM-0187}}\par
\noindent\emph{Iztočnik / Source:} bilden einen Z-Modul / Moduleigenschaft\par
\noindent\emph{Rusky / Russian:} образуют Z-модуль / модульное свойство\par
\noindent\emph{Ukrajinsky / Ukrainian:} утворюють Z-модуль / модульна властивість\par
\smallskip\noindent\emph{Razlog:} Tvorjut Z-modul zachovuje tvrđenje, že slědy za fiksovanu klasu imajut modulnu strukturu. Modulno svojstvo odnosi se k identičnostjam, ktore tu strukturu dokazyvajut. Prěvod zato ne ostavaje pri običnom zbirnom nazvanju: algebraične operacije i njih svojstva sut sučstvena čest argumenta.\par
\noindent\emph{Reason:} Tvorjut Z-modul preserves the assertion that the traces for a fixed class have a module structure. Modulno svojstvo refers to the displayed identities establishing that structure. The wording therefore goes beyond a collective name: the algebraic operations and their properties carry the argument.\par
\noindent{\footnotesize Izvorne zapisy / Original records: HIST-g136-002.}\par
\end{minipage}\par
\par\medskip\Needspace{7\baselineskip}
\noindent\begin{minipage}{\linewidth}
\hypertarget{TERM-0188}{}
\noindent\textbf{konstrukcije}\quad{\footnotesize\texttt{TERM-0188}}\par
\noindent\emph{Iztočnik / Source:} Bildungen\par
\noindent\emph{Rusky / Russian:} построения\par
\noindent\emph{Ukrajinsky / Ukrainian:} побудови\par
\smallskip\noindent\emph{Razlog:} Konstrukcije prodolžaje izbor iz raboty 1. Izvorni zapis odnak razlikuje konstrukciju od formy i ostavjaje možnost, že konkretni zapisy v tablicah trěbujut forma. Izbor ne jest tvrdo pravilo za německo Bildungen: odlučuje, čto imenovany objekt děla v konkretnoj pasaži.\par
\noindent\emph{Reason:} Konstrukcije continues the choice in Paper 1. The original record nevertheless distinguishes a construction from a form and leaves open whether particular table entries require form. Bildungen has no fixed replacement here; the object's role in the passage determines the wording.\par
\noindent{\footnotesize Izvorne zapisy / Original records: HIST-g002-011.}\par
\end{minipage}\par
\par\medskip\Needspace{7\baselineskip}
\noindent\begin{minipage}{\linewidth}
\hypertarget{TERM-0189}{}
\noindent\textbf{porodžujųće formy / konstrukcije}\quad{\footnotesize\texttt{TERM-0189}}\par
\noindent\emph{Iztočnik / Source:} Bildungen\par
\noindent\emph{Rusky / Russian:} порождающие формы / построения\par
\noindent\emph{Ukrajinsky / Ukrainian:} породжувальні форми / конструкції\par
\smallskip\noindent\emph{Razlog:} Porodžujųće formy i konstrukcije sut dvě zapisane možnosti, ne dva obvezno zaměnime nazvanja. Kratka primětka opreděljaje kontekst kako kriterij izbora, ale ne dava odvojenu čislenu ocěnu kandidatom. Etimologične bukvy v citovanoj formě sut zachovane radi točnosti zapisa; on sam ne rěšaje osnovny pravopis pozdnějšego teksta.\par
\noindent\emph{Reason:} Porodžujųće formy and konstrukcije are two recorded possibilities, not necessarily interchangeable names. The brief note makes context the selection criterion but assigns no numerical scores. The etymological letters are retained to quote the candidate accurately; the record does not itself settle the basic spelling of a later text.\par
\noindent{\footnotesize Izvorne zapisy / Original records: HIST-g001-003.}\par
\end{minipage}\par
\par\medskip\Needspace{7\baselineskip}
\noindent\begin{minipage}{\linewidth}
\hypertarget{TERM-0190}{}
\noindent\textbf{bilinearne kombinacije}\quad{\footnotesize\texttt{TERM-0190}}\par
\noindent\emph{Iztočnik / Source:} bilineare Verbindungen\par
\noindent\emph{Rusky / Russian:} билинейные комбинации\par
\noindent\emph{Ukrajinsky / Ukrainian:} білінійні комбінації\par
\smallskip\noindent\emph{Razlog:} Kombinacije jest točnějše od neoprěděljenyh svezij, ibo formuly zaměnjajut dane izrazy bilinearnymi izrazami v $u$ i $v$. Bilinearne zachovuje važnu algebraičnu kvalifikaciju. Razlog izbora zato leži v funkciji formuly, ne samo v podobnosti medžunarodnogo termina.\par
\noindent\emph{Reason:} Kombinacije is more specific than vague connections because the formulas replace the given expressions by bilinear expressions in $u$ and $v$. Bilinearne retains the essential algebraic qualification. The choice is grounded in what the formulas do, not simply in the international term's familiar appearance.\par
\noindent{\footnotesize Izvorne zapisy / Original records: HIST-g097-015.}\par
\end{minipage}\par
\par\medskip\Needspace{7\baselineskip}
\noindent\begin{minipage}{\linewidth}
\hypertarget{TERM-0191}{}
\noindent\textbf{binarna forma / binarna linearna forma}\quad{\footnotesize\texttt{TERM-0191}}\par
\noindent\emph{Iztočnik / Source:} binäre Form / binäre Linearform\par
\noindent\emph{Rusky / Russian:} бинарная форма / бинарная линейная форма\par
\noindent\emph{Ukrajinsky / Ukrainian:} бінарна форма / бінарна лінійна форма\par
\smallskip\noindent\emph{Razlog:} Binarna forma i binarna linearna forma imajut istu strukturu nazvanja; linearna izrično odliča užši slučaj. Izvorni razlog daje prednost kratkoj medžunarodnoj terminologiji i sravnjaje ju s ruskoju i ukrajinskoju matematičnoju formoju. On ne dava novy lokalny sinonim, ktory by trěboval dodatno objasnjenje.\par
\noindent\emph{Reason:} Binarna forma and binarna linearna forma share the same naming structure, with linearna explicitly identifying the narrower case. The recorded reason favours compact international terminology, compared with the Russian and Ukrainian mathematical forms, rather than introducing a new local synonym requiring explanation.\par
\noindent{\footnotesize Izvorne zapisy / Original records: HIST-g077-004.}\par
\end{minipage}\par
\par\medskip\Needspace{7\baselineskip}
\noindent\begin{minipage}{\linewidth}
\hypertarget{TERM-0192}{}
\noindent\textbf{binarne rědy / rědy proměnnyh}\quad{\footnotesize\texttt{TERM-0192}}\par
\noindent\emph{Iztočnik / Source:} binäre Reihen / Reihen von Variabeln\par
\noindent\emph{Rusky / Russian:} бинарные ряды / ряды переменных\par
\noindent\emph{Ukrajinsky / Ukrainian:} бінарні ряди / ряди змінних\par
\smallskip\noindent\emph{Razlog:} Rědy proměnnyh trěba čitati v simboličnom jazyku teorije invariantov. Izbor rědy zachovuje avtorske grupovanje proměnnyh i ne vnosi avtomatično analitičny smysl beskonečnogo sumovanja. Binarne ostavaje konkretnoju kvalifikacijeju teh redov, ne občim imenovanjem vsakoj poslědovatelnosti.\par
\noindent\emph{Reason:} Rědy proměnnyh must be read in the symbolic language of invariant theory. Rědy preserves the author's grouping of variables without automatically importing the analytic sense of infinite summation. Binarne remains the particular qualification of those rows, not a generic name for every sequence.\par
\noindent{\footnotesize Izvorne zapisy / Original records: HIST-g078-002.}\par
\end{minipage}\par
\par\medskip\Needspace{7\baselineskip}
\noindent\begin{minipage}{\linewidth}
\hypertarget{TERM-0193}{}
\noindent\textbf{biracionalna transformacija}\quad{\footnotesize\texttt{TERM-0193}}\par
\noindent\emph{Iztočnik / Source:} birationale Transformation\par
\noindent\emph{Rusky / Russian:} бирациональное преобразование\par
\noindent\emph{Ukrajinsky / Ukrainian:} біраціональне перетворення\par
\smallskip\noindent\emph{Razlog:} Biracionalna transformacija zachovuje obostrannu racionalnu obnovimost baz: jedna racionalna baza može byti vyražena črěz drugu i obratno. Izbor medžunarodnogo termina drži tu obostrannost v kratkom nazvanju. Prosto racionalna transformacija ne by izrično zachovala istu informaciju.\par
\noindent\emph{Reason:} Biracionalna transformacija preserves mutual rational recoverability of the bases: one rational basis can be expressed through the other and conversely. The international term retains that two-way condition in a short name. Merely rational transformation would not express the same information.\par
\noindent{\footnotesize Izvorne zapisy / Original records: HIST-g044-012.}\par
\end{minipage}\par
\par\medskip\Needspace{7\baselineskip}
\noindent\begin{minipage}{\linewidth}
\hypertarget{TERM-0194}{}
\noindent\textbf{diskriminant, opreděljeny do kvadrata jedinicy, jest različny od nuly}\quad{\footnotesize\texttt{TERM-0194}}\par
\noindent\emph{Iztočnik / Source:} bis auf das Quadrat einer Einheit bestimmte Diskriminante von Null verschieden\par
\noindent\emph{Rusky / Russian:} дискриминант, определенный с точностью до квадрата единицы, отличен от нуля\par
\noindent\emph{Ukrajinsky / Ukrainian:} дискримінант, визначений з точністю до квадрата одиниці, відмінний від нуля\par
\smallskip\noindent\emph{Razlog:} Opreděljeny do kvadrata jedinicy vyražaje nejednoznačnost izbora skalarnogo diskriminanta, zatymčas različny od nuly zachovuje tvrđenje o njegovoj nenulovosti. Tuto ne jest istovětna formulacija kriteriju črěz diskriminantny ideal. Kratka konstrukcija opreděljeny do jest zabilěženy jezyčny izbor, čiju priměnu trěba sravnjati v istom smyslu pri drugyh pojavjenjah.\par
\noindent\emph{Reason:} Opreděljeny do kvadrata jedinicy expresses the ambiguity in choosing a scalar discriminant, while različny od nuly preserves its nonvanishing. This is not identical in formulation to a criterion stated through the discriminant ideal. The compact construction opreděljeny do is a recorded linguistic choice whose other occurrences require comparison in the same sense.\par
\noindent{\footnotesize Izvorne zapisy / Original records: HIST-g144-006.}\par
\end{minipage}\par
\par\medskip\Needspace{7\baselineskip}
\noindent\begin{minipage}{\linewidth}
\hypertarget{TERM-0195}{}
\noindent\textbf{blokovo-trikutna forma / matrično predstavjenje}\quad{\footnotesize\texttt{TERM-0195}}\par
\noindent\emph{Iztočnik / Source:} Block-Dreiecksform / Matrizendarstellung\par
\noindent\emph{Rusky / Russian:} блочно-треугольная форма / матричное представление\par
\noindent\emph{Ukrajinsky / Ukrainian:} блоково-трикутна форма / матричне подання\par
\smallskip\noindent\emph{Razlog:} Blokovo-trikutna prodolžaje familiju blokovo-diagonalna, ale zachovuje drugy tip matričnoj formy. Medžu trikutna i treugolna zapis daje prednost prozračnomu složenju tri i kut, da by ne zavisel samo od ruskoj formy. To jest slovotvorny i sravniteljny razlog; posebna izměrena vseslovjanska razumělivost ne jest dana.\par
\noindent\emph{Reason:} Blokovo-trikutna continues the blokovo-diagonalna family while naming a different matrix shape. Between trikutna and treugolna, the record prefers the transparent tri-plus-kut formation to avoid reliance on the Russian form alone. This is a word-formation and comparative argument, not a reported measurement of pan-Slavic comprehension.\par
\noindent{\footnotesize Izvorne zapisy / Original records: HIST-g141-005.}\par
\end{minipage}\par
\par\medskip\Needspace{7\baselineskip}
\noindent\begin{minipage}{\linewidth}
\hypertarget{TERM-0196}{}
\noindent\textbf{blokova matrica / blokovo-diagonalna matrica}\quad{\footnotesize\texttt{TERM-0196}}\par
\noindent\emph{Iztočnik / Source:} Blockmatrix / Blockdiagonalmatrix\par
\noindent\emph{Rusky / Russian:} блочная матрица / блочно-диагональная матрица\par
\noindent\emph{Ukrajinsky / Ukrainian:} блокова матриця / блоково-діагональна матриця\par
\smallskip\noindent\emph{Razlog:} Blokova matrica i blokovo-diagonalna matrica razděljujut obče grupovanje elementov v bloky od osobnogo diagonalnogo raspoloženja. Kratke nazvanja služe dokazě o konečnom podpolju bez ponovjenogo dolgogo opisa. Slične ruske i ukrajinske formy podpirajut izbor medžunarodnoj familije blok- i matric-.\par
\noindent\emph{Reason:} Blokova matrica and blokovo-diagonalna matrica distinguish a general partition into blocks from the particular diagonal arrangement. The compact names serve the finite-subfield proof without repeated lengthy descriptions. The corresponding Russian and Ukrainian forms support the international block/matrix family.\par
\noindent{\footnotesize Izvorne zapisy / Original records: HIST-g196-003.}\par
\end{minipage}\par
\par\medskip\Needspace{7\baselineskip}
\noindent\begin{minipage}{\linewidth}
\hypertarget{TERM-0197}{}
\noindent\textbf{Brauerova grupa klasov algeber}\quad{\footnotesize\texttt{TERM-0197}}\par
\noindent\emph{Iztočnik / Source:} Brauersche Gruppe der Algebrenklassen\par
\noindent\emph{Rusky / Russian:} [Bukvy ne sut zachovane / Lettering not recoverable]\par
\noindent\emph{Ukrajinsky / Ukrainian:} [Bukvy ne sut zachovane / Lettering not recoverable]\par
\smallskip\noindent\emph{Razlog:} Brauerova grupa klasov algeber izrično imenuje grupu i objekty njezinyh klas, da by istorična opisna formulacija ne ostala nejasna. Medžuslovjanska forma i razlog sut čitateljne. Ruska i ukrajinska polja odnak imajut samo znaky ? vměsto bukv, zato ta zapis ne može byti dokazom točnoj grafije njih terminov.\par
\noindent\emph{Reason:} Brauerova grupa klasov algeber explicitly names the group and the objects being classified so that the historical descriptive wording is not left opaque. The Interslavic form and its rationale are readable. The Russian and Ukrainian fields, however, contain only damaged question marks and cannot establish those languages' exact spellings.\par
\noindent{\footnotesize Izvorne zapisy / Original records: HIST-g209-015.}\par
\end{minipage}\par
\par\medskip\Needspace{7\baselineskip}
\noindent\begin{minipage}{\linewidth}
\hypertarget{TERM-0198}{}
\noindent\textbf{recenzije}\quad{\footnotesize\texttt{TERM-0198}}\par
\noindent\emph{Iztočnik / Source:} Buchbesprechungen\par
\noindent\emph{Rusky / Russian:} рецензии\par
\noindent\emph{Ukrajinsky / Ukrainian:} рецензії\par
\smallskip\noindent\emph{Razlog:} Recenzije imenuje bibliografičny žanr vměsto doslovnoj konstrukcije besědy o knigah. Izbor odgovarjaje funkciji rubriky i ruskoj i ukrajinskoj paralelě. Kratkost ne vypuščaje informaciju o žanru; ona izběgaje neprikladnu doslovnu metaforu.\par
\noindent\emph{Reason:} Recenzije names the bibliographical genre rather than calquing book discussions. The choice fits the heading's function and the Russian and Ukrainian parallels. Its brevity does not omit the genre information; it avoids a misleading literal formulation.\par
\noindent{\footnotesize Izvorne zapisy / Original records: HIST-g215-003.}\par
\end{minipage}\par
\par\medskip\Needspace{7\baselineskip}
\noindent\begin{minipage}{\linewidth}
\hypertarget{TERM-0199}{}
\noindent\textbf{teorema Burnsida / obobčena teorema Burnsida}\quad{\footnotesize\texttt{TERM-0199}}\par
\noindent\emph{Iztočnik / Source:} Burnsidescher Satz / verallgemeinerter Burnsidescher Satz\par
\noindent\emph{Rusky / Russian:} теорема Бернсайда / обобщенная теорема Бернсайда\par
\noindent\emph{Ukrajinsky / Ukrainian:} теорема Бернсайда / узагальнена теорема Бернсайда\par
\smallskip\noindent\emph{Razlog:} Teorema Burnsida zachovuje imenovany matematičny rezultat, a obobčena dava točnu razliku njegovoj občejšej varianty. Paralelna konstrukcija pomaga čitatelju viděti svez medžu dvoma rezultatami. Zapis opravdava to imenovanje, ale ne ustanavjaje osobno pravilo za vsake vlastno ime ili za rod slova teorem v cělom korpusu.\par
\noindent\emph{Reason:} Teorema Burnsida preserves the named mathematical result, while obobčena identifies its generalized variant. The parallel construction makes the relationship between the two results visible. The record explains this naming decision but does not independently settle every proper-name adaptation or the gender of theorem throughout the corpus.\par
\noindent{\footnotesize Izvorne zapisy / Original records: HIST-g198-005.}\par
\end{minipage}\par
\par\medskip\Needspace{7\baselineskip}
\noindent\begin{minipage}{\linewidth}
\hypertarget{TERM-0200}{}
\noindent\textbf{knjigy, sostavjene s učestjem Emmy Noether}\quad{\footnotesize\texttt{TERM-0200}}\par
\noindent\emph{Iztočnik / Source:} Bücher, entstanden unter Mitwirkung von Emmy Noether\par
\noindent\emph{Rusky / Russian:} книги, созданные при участии Эмми Нётер\par
\noindent\emph{Ukrajinsky / Ukrainian:} книги, створені за участю Еммі Нетер\par
\smallskip\noindent\emph{Razlog:} S učestjem Emmy Noether izrično označaje učestje v sostavjenju knig. Iztočnik ne pripisuje jej jedino avtorstvo, zato prěvod ne smije proměniti učestje v takvu atribuciju. To jest odluka o točnosti bibliografičnogo opisa, ne samo o leksičnoj podobnosti slova Mitwirkung.\par
\noindent\emph{Reason:} S učestjem Emmy Noether explicitly states participation in producing the books. The source does not assign sole authorship, so the translation must not turn participation into that attribution. This is a decision about bibliographical accuracy, not merely lexical resemblance to Mitwirkung.\par
\noindent{\footnotesize Izvorne zapisy / Original records: HIST-g215-004.}\par
\end{minipage}\par
\par\medskip\Needspace{7\baselineskip}
\noindent\begin{minipage}{\linewidth}
\hypertarget{TERM-0201}{}
\noindent\textbf{karakteristika nul}\quad{\footnotesize\texttt{TERM-0201}}\par
\noindent\emph{Iztočnik / Source:} Characteristic zero\par
\noindent\emph{Rusky / Russian:} характеристика нуль\par
\noindent\emph{Ukrajinsky / Ukrainian:} характеристика нуль\par
\smallskip\noindent\emph{Razlog:} Karakteristika zachovuje medžunarodny termin teorije poljev, a nul dava čislenu kvalifikaciju. Izvorni razlog jest kratky i odnosi se k tehničnomu registru. On ne sravnjaje samostojno formy nul i nula, ktore sut dokumentovane v drugyh zapisah.\par
\noindent\emph{Reason:} Karakteristika retains the international field-theoretic term, and nul supplies its numerical qualification. The brief original reason concerns technical register. It does not independently compare nul with nula, both of which occur in other records.\par
\noindent{\footnotesize Izvorne zapisy / Original records: HIST-g091-011, HIST-g092-080.}\par
\end{minipage}\par
\par\medskip\Needspace{7\baselineskip}
\noindent\begin{minipage}{\linewidth}
\hypertarget{TERM-0202}{}
\noindent\textbf{karakter}\quad{\footnotesize\texttt{TERM-0202}}\par
\noindent\emph{Iztočnik / Source:} Charakter\par
\noindent\emph{Rusky / Russian:} характер\par
\noindent\emph{Ukrajinsky / Ukrainian:} характер\par
\smallskip\noindent\emph{Razlog:} Karakter označaje tehničny pojem teorije predstavjenij, ne znak ili oznaku v običnom smyslu. Dva zapisy opirajut se na medžunarodno imenovanje i doslědnost s oddělom 24. Ta identičnost pojma jest razlogom, da se ne izbira običny sinonim bez njegovej matematičnoj funkcije.\par
\noindent\emph{Reason:} Karakter denotes the representation-theoretic concept, not a sign or mark in the ordinary sense. Two records appeal to the international name and consistency with section 24. Preserving that concept explains why an everyday synonym without the technical function was not chosen.\par
\noindent{\footnotesize Izvorne zapisy / Original records: HIST-g202-004, HIST-g204-010.}\par
\end{minipage}\par
\par\medskip\Needspace{7\baselineskip}
\noindent\begin{minipage}{\linewidth}
\hypertarget{TERM-0203}{}
\noindent\textbf{karakter / sootvětny nerazložimy karakter}\quad{\footnotesize\texttt{TERM-0203}}\par
\noindent\emph{Iztočnik / Source:} Charakter / zugehöriger einfacher Charakter\par
\noindent\emph{Rusky / Russian:} характер / соответствующий неприводимый характер\par
\noindent\emph{Ukrajinsky / Ukrainian:} характер / відповідний незвідний характер\par
\smallskip\noindent\emph{Razlog:} Nerazložimy karakter jest kontekstualno tolkovanje starějšego einfacher Charakter, obrazloženo razmatranjem absolutno nerazložimyh sostavnikov. Izbor približaje nazvanje sovrěmennomu tehničnomu jazyku; ne jest prosto doslovny prěvod einfach. Slovnik zato zachovuje originalnu formulaciju prěd očami čitatelja, da by toj mogl ocěniti to tolkovanje.\par
\noindent\emph{Reason:} Nerazložimy karakter is a contextual interpretation of the older einfacher Charakter, justified in the record by the discussion of absolutely irreducible constituents. It brings the name into modern technical vocabulary rather than translating einfach literally. The dictionary therefore keeps the original wording visible so that readers can assess the interpretation.\par
\noindent{\footnotesize Izvorne zapisy / Original records: HIST-g167-004.}\par
\end{minipage}\par
\par\medskip\Needspace{7\baselineskip}
\noindent\begin{minipage}{\linewidth}
\hypertarget{TERM-0204}{}
\noindent\textbf{karaktere}\quad{\footnotesize\texttt{TERM-0204}}\par
\noindent\emph{Iztočnik / Source:} Charaktere\par
\noindent\emph{Rusky / Russian:} характеры\par
\noindent\emph{Ukrajinsky / Ukrainian:} характери\par
\smallskip\noindent\emph{Razlog:} Karaktere zachovuje medžunarodnu osnovu za grupove karaktery i prodolžaje njezino upotrěbjenje v korpusu. Ruska i ukrajinska paralela podpirajut tehničny smysl, ne avtomatično medžuslovjansky konec množiny. Točny kandidat jest dokumentovany; izbor korenja i njegovo sklanjanje ostavajut različnymi pytanjami.\par
\noindent\emph{Reason:} Karaktere retains the international stem for group characters and continues its use in the corpus. The Russian and Ukrainian parallels support the technical sense, not automatically the Interslavic plural ending. The exact candidate is documented; selecting the stem and establishing its inflection remain distinct questions.\par
\noindent{\footnotesize Izvorne zapisy / Original records: HIST-g200-010.}\par
\end{minipage}\par
\par\medskip\Needspace{7\baselineskip}
\noindent\begin{minipage}{\linewidth}
\hypertarget{TERM-0205}{}
\noindent\textbf{karaktere / rang}\quad{\footnotesize\texttt{TERM-0205}}\par
\noindent\emph{Iztočnik / Source:} Charaktere / Rang\par
\noindent\emph{Rusky / Russian:} характеры / ранг\par
\noindent\emph{Ukrajinsky / Ukrainian:} характери / ранг\par
\smallskip\noindent\emph{Razlog:} Karaktere i rang sut dva kratke algebraične termina, ne sinonimy. Izvorni razlog posebno zachovuje rang kako medžunarodno nazvanje vměsto dolgogo opisnogo izraza. On ne sodržuje samostojno sravnjenje drugih korenj za oba pojmy.\par
\noindent\emph{Reason:} Karaktere and rang are two compact algebraic terms, not synonyms. The original reason specifically retains rang as an international name rather than a lengthy descriptive phrase. It contains no separate comparison of alternative roots for both concepts.\par
\noindent{\footnotesize Izvorne zapisy / Original records: HIST-g105-011.}\par
\end{minipage}\par
\par\medskip\Needspace{7\baselineskip}
\noindent\begin{minipage}{\linewidth}
\hypertarget{TERM-0206}{}
\noindent\textbf{harakterizujuča oblast odobraženja}\quad{\footnotesize\texttt{TERM-0206}}\par
\noindent\emph{Iztočnik / Source:} charakterisierender Abbildungsbereich\par
\noindent\emph{Rusky / Russian:} характеризующая область отображения\par
\noindent\emph{Ukrajinsky / Ukrainian:} характеризувальна область відображення\par
\smallskip\noindent\emph{Razlog:} Oblast odobraženja prodolžaje imenovanje Abbildungsbereich, a harakterizujuča označaje oblast, črěz ktoru se karakterizuje objekt prvogo roda. Pridavnik zato ne jest neobvezno ukrašenje: on razlikuje konkretnu funkciju oblasti. Ranějši napis harakterizujuča jest citovana forma, ne novo pravilo protiv korenja karakter- v drugyh tehničnyh slovah.\par
\noindent\emph{Reason:} Oblast odobraženja continues the rendering of Abbildungsbereich, while harakterizujuča identifies the domain that characterizes the first-kind object. The adjective is not optional decoration; it distinguishes the domain's particular function. Historical harakterizujuča is the quoted candidate, not a new rule against karakter- in other technical words.\par
\noindent{\footnotesize Izvorne zapisy / Original records: HIST-g050-006.}\par
\end{minipage}\par
\par\medskip\Needspace{7\baselineskip}
\noindent\begin{minipage}{\linewidth}
\hypertarget{TERM-0207}{}
\noindent\textbf{karakteristika}\quad{\footnotesize\texttt{TERM-0207}}\par
\noindent\emph{Iztočnik / Source:} Charakteristik\par
\noindent\emph{Rusky / Russian:} характеристика\par
\noindent\emph{Ukrajinsky / Ukrainian:} характеристика\par
\smallskip\noindent\emph{Razlog:} Karakteristika jest vlastnost polja, od ktorej zavisi razmatrana konstrukcija korenja stepena $p$. Medžunarodny termin odgovarja ruskoj i ukrajinskoj paralelě. Matematičny razlog za njegovo točno zachovanje jest zavisimost konstrukcije od karakteristiky $p$, ne samo terminologična doslědnost.\par
\noindent\emph{Reason:} Karakteristika is the field property on which the construction of a $p$th root depends. The international term matches the Russian and Ukrainian parallels. Its mathematical importance is the construction's dependence on characteristic $p$, not merely terminological consistency.\par
\noindent{\footnotesize Izvorne zapisy / Original records: HIST-g128-006.}\par
\end{minipage}\par
\par\medskip\Needspace{7\baselineskip}
\noindent\begin{minipage}{\linewidth}
\hypertarget{TERM-0208}{}
\noindent\textbf{karakteristika 0 / p}\quad{\footnotesize\texttt{TERM-0208}}\par
\noindent\emph{Iztočnik / Source:} Charakteristik 0 / p\par
\noindent\emph{Rusky / Russian:} характеристика 0 / p\par
\noindent\emph{Ukrajinsky / Ukrainian:} характеристика 0 / p\par
\smallskip\noindent\emph{Razlog:} Karakteristika 0 i karakteristika $p$ zachovujut avtorske razděljenje na nulevu i prostu pozitivnu karakteristiku. Oznaky 0 i $p$ ne sut dva proizvoljne priměry istogo slučaja. Skupna terminologija pomaga viděti rozděl, ne odstraniti uslovja osobnyh argumentov.\par
\noindent\emph{Reason:} Karakteristika 0 and karakteristika $p$ preserve the author's distinction between zero and positive prime characteristic. The symbols 0 and $p$ are not merely two arbitrary examples of the same case. Shared terminology makes the distinction visible rather than removing the conditions of the separate arguments.\par
\noindent{\footnotesize Izvorne zapisy / Original records: HIST-g068-004.}\par
\end{minipage}\par
\par\medskip\Needspace{7\baselineskip}
\noindent\begin{minipage}{\linewidth}
\hypertarget{TERM-0209}{}
\noindent\textbf{karakteristika nula}\quad{\footnotesize\texttt{TERM-0209}}\par
\noindent\emph{Iztočnik / Source:} Charakteristik Null\par
\noindent\emph{Rusky / Russian:} характеристика нуль\par
\noindent\emph{Ukrajinsky / Ukrainian:} характеристика нуль\par
\smallskip\noindent\emph{Razlog:} Karakteristika nula jest kratka algebraična konstrukcija pri formulě. Zapis dava prednost kratkosti, ne dolgomu objasnjajučemu obratovi. On dokumentuje formu nula, ale ne daje razloga za jej izključivu prednost nad istorijno zapisanym nul v drugih pasažah.\par
\noindent\emph{Reason:} Karakteristika nula is the concise algebraic expression used beside the formula. The record favours brevity over an explanatory circumlocution. It documents nula but supplies no argument for its exclusive preference over the historical nul found elsewhere.\par
\noindent{\footnotesize Izvorne zapisy / Original records: HIST-g203-010.}\par
\end{minipage}\par
\par\medskip\Needspace{7\baselineskip}
\noindent\begin{minipage}{\linewidth}
\hypertarget{TERM-0210}{}
\noindent\textbf{karakteristika nul / prosto polje}\quad{\footnotesize\texttt{TERM-0210}}\par
\noindent\emph{Iztočnik / Source:} Charakteristik Null / Primkörper\par
\noindent\emph{Rusky / Russian:} характеристика нуль / простое поле\par
\noindent\emph{Ukrajinsky / Ukrainian:} характеристика нуль / просте поле\par
\smallskip\noindent\emph{Razlog:} Karakteristika nul i prosto polje imenujut različne česti argumenta. Slučaj prostogo polja razděljuje se od argumenta specializacije pri karakteristikě nula na koncu odděla 4. Kratke nazvanja ne smiju spojiti tute dva koraka ili prěnesti uslovje odnogo na drugy bez avtorskogo osnovanja.\par
\noindent\emph{Reason:} Karakteristika nul and prosto polje name distinct parts of the argument. The prime-field case is kept separate from the characteristic-zero specialization at the end of section 4. Compact names must not merge those steps or transfer one step's assumptions to the other without support from the source.\par
\noindent{\footnotesize Izvorne zapisy / Original records: HIST-g088-009, HIST-g092-040.}\par
\end{minipage}\par
\par\medskip\Needspace{7\baselineskip}
\noindent\begin{minipage}{\linewidth}
\hypertarget{TERM-0211}{}
\noindent\textbf{harakteristika p}\quad{\footnotesize\texttt{TERM-0211}}\par
\noindent\emph{Iztočnik / Source:} Charakteristik p\par
\noindent\emph{Rusky / Russian:} характеристика p\par
\noindent\emph{Ukrajinsky / Ukrainian:} характеристика p\par
\smallskip\noindent\emph{Razlog:} $p$ jest latinična matematična oznaka i ostavaje takva i pri prěpisu okolnoj prozy v kirilicu. Izvorni razlog opreděljaje tu granicu medžu tekstom i oznakoju. On ne sravnjaje h i k na početku kandidata harakteristika, zato ne daje odvojeno opravdanje za razliku od karakteristika v drugyh zapisah.\par
\noindent\emph{Reason:} $p$ is a Latin mathematical symbol and remains so when the surrounding prose is written in Cyrillic. The original reason establishes that boundary between prose and notation. It does not compare initial h and k in harakteristika and therefore gives no separate justification for the difference from karakteristika elsewhere.\par
\noindent{\footnotesize Izvorne zapisy / Original records: HIST-g040-019.}\par
\end{minipage}\par
\par\medskip\Needspace{7\baselineskip}
\noindent\begin{minipage}{\linewidth}
\hypertarget{TERM-0212}{}
\noindent\textbf{harakternye svojstva}\quad{\footnotesize\texttt{TERM-0212}}\par
\noindent\emph{Iztočnik / Source:} charakteristische Eigenschaften\par
\noindent\emph{Rusky / Russian:} характеристические свойства\par
\noindent\emph{Ukrajinsky / Ukrainian:} характеристичні властивості\par
\smallskip\noindent\emph{Razlog:} Harakternye svojstva veže zaključok odděla 7 s jezykom prěnosa i odobraženja iz odděla 3. Zapis objasnjaje kontekstualnu doslědnost, ale ne dava osobnu gramatičnu argumentaciju za citovany pridavnik. Njegovu točnu formu trěba razlikovati od razloga za zachovanje istogo pojma v obojih pasažah.\par
\noindent\emph{Reason:} Harakternye svojstva connects the conclusion of section 7 with the transfer and mapping language of section 3. The record explains contextual consistency but supplies no separate grammatical argument for the quoted adjective. Its exact form must be distinguished from the reason for preserving the same concept in both passages.\par
\noindent{\footnotesize Izvorne zapisy / Original records: HIST-g047-013.}\par
\end{minipage}\par
\par\medskip\Needspace{7\baselineskip}
\noindent\begin{minipage}{\linewidth}
\hypertarget{TERM-0213}{}
\noindent\textbf{karakteristična funkcija / tvorjača funkcija}\quad{\footnotesize\texttt{TERM-0213}}\par
\noindent\emph{Iztočnik / Source:} charakteristische Funktion / erzeugende Funktion\par
\noindent\emph{Rusky / Russian:} характеристическая функция / производящая функция\par
\noindent\emph{Ukrajinsky / Ukrainian:} характеристична функція / твірна функція\par
\smallskip\noindent\emph{Razlog:} Tvorjača funkcija jest izbrana vměsto generujuča kako kratka konstrukcija so slovjanskoju osnovoju. Ruska i ukrajinska paralela podpirajut tehničny smysl tvorjačej funkcije. Izvorni zapis odnak izrično ostavjaje medžuslovjansku formu za sravnjenje s jezyčnymi iztočnikami. Para s karakteristična funkcija ne znači, že dva německe nazvanja sut obče sinonimy.\par
\noindent\emph{Reason:} Tvorjača funkcija is preferred to generujuča as a compact expression with a Slavic stem. The Russian and Ukrainian parallels support the technical generating-function sense. The original record nevertheless explicitly leaves the Interslavic form for examination against linguistic sources. Pairing it with karakteristična funkcija does not make the two German names general synonyms.\par
\noindent{\footnotesize Izvorne zapisy / Original records: HIST-g104-003.}\par
\end{minipage}\par
\par\medskip\Needspace{7\baselineskip}
\noindent\begin{minipage}{\linewidth}
\hypertarget{TERM-0214}{}
\noindent\textbf{harakteristična mnogokratnost}\quad{\footnotesize\texttt{TERM-0214}}\par
\noindent\emph{Iztočnik / Source:} charakteristische Multiplizität\par
\noindent\emph{Rusky / Russian:} характеристическая кратность\par
\noindent\emph{Ukrajinsky / Ukrainian:} характеристична кратність\par
\smallskip\noindent\emph{Razlog:} Mnogokratnost zachovuje algebraičny smysl kratnosti i prodolžaje ranějše imenovanje takvyh čislenyh odnosov v korpusu. Izvorni razlog daje prednost razumělivosti slovjanskoj osnove v sravnjenju s ruskoju i ukrajinskoju kratnostju. To jest obrazloženy leksičny izbor, ne rezultat izměrjenja razumělivosti medžu govornikami vseh slovjanskyh jezykov.\par
\noindent\emph{Reason:} Mnogokratnost retains algebraic multiplicity and continues the corpus's earlier naming of such counts. The recorded reason favours the intelligibility of the Slavic stem, compared with Russian and Ukrainian multiplicity terms. This is a reasoned lexical preference, not a measured comprehension result across Slavic-speaking populations.\par
\noindent{\footnotesize Izvorne zapisy / Original records: HIST-g080-004.}\par
\end{minipage}\par
\par\medskip\Needspace{7\baselineskip}
\noindent\begin{minipage}{\linewidth}
\hypertarget{TERM-0215}{}
\noindent\textbf{karakterističny rjad}\quad{\footnotesize\texttt{TERM-0215}}\par
\noindent\emph{Iztočnik / Source:} charakteristische Reihe\par
\noindent\emph{Rusky / Russian:} характеристический ряд\par
\noindent\emph{Ukrajinsky / Ukrainian:} характеристичний ряд\par
\smallskip\noindent\emph{Razlog:} Karakterističny označaje invariantnost pod vsemi avtomorfizmami. Pridavnik zato ne govori samo, že red jest tipičny ili osoblivo poznavajemy. Izvorni razlog opravdava tehničnu funkciju pridavnika, ale ne daje nezavisnu argumentaciju za napis rjad vměsto drugih dokumentovanyh form slova za red.\par
\noindent\emph{Reason:} Karakterističny refers to invariance under all automorphisms. It does not merely mean that the series is typical or especially recognizable. The original reason supports the adjective's technical function but supplies no independent argument for rjad over other recorded forms of the word for a series.\par
\noindent{\footnotesize Izvorne zapisy / Original records: HIST-g181-004.}\par
\end{minipage}\par
\par\medskip\Needspace{7\baselineskip}
\noindent\begin{minipage}{\linewidth}
\hypertarget{TERM-0216}{}
\noindent\textbf{Klebš--Gordanovo razloženje po redah}\quad{\footnotesize\texttt{TERM-0216}}\par
\noindent\emph{Iztočnik / Source:} Clebsch-Gordansche Reihenentwicklung\par
\noindent\emph{Rusky / Russian:} разложение в ряд Клебша--Гордана\par
\noindent\emph{Ukrajinsky / Ukrainian:} розклад у ряд Клебша--Гордана\par
\smallskip\noindent\emph{Razlog:} Razloženje po redah zachovuje vezu s istoričnym teoremom o razloženju, zatymčas Klebš i Gordan ostavajut poznavajemyma imenami v medžuslovjanskoj proze. Prěpis imena v proze ne měnja originalny napis inojezyčnogo bibliografičnogo nazvanja. Ta razlika dopuščaje i čitatelsku dostupnost, i točno nahodženje iztočnika.\par
\noindent\emph{Reason:} Razloženje po redah maintains the connection with the historical expansion theorem, while Klebš and Gordan remain recognizable names in Interslavic prose. Adapting a name in running text does not change the original spelling of a foreign bibliographical title. The distinction supports both readability and exact source identification.\par
\noindent{\footnotesize Izvorne zapisy / Original records: HIST-g057-009.}\par
\end{minipage}\par
\par\medskip\Needspace{7\baselineskip}
\noindent\begin{minipage}{\linewidth}
\hypertarget{TERM-0217}{}
\noindent\textbf{razklad Klebš--Gordana}\quad{\footnotesize\texttt{TERM-0217}}\par
\noindent\emph{Iztočnik / Source:} Clebsch-Gordansche Reihenentwicklung\par
\noindent\emph{Rusky / Russian:} разложение Клебша--Гордана\par
\noindent\emph{Ukrajinsky / Ukrainian:} розклад Клебша--Гордана\par
\smallskip\noindent\emph{Razlog:} Klebš zaměnjaje německy napis Clebsch v běžnoj proze, da by ne trěbal kombinaciju bukv, ktora ne odpovada običnomu medžuslovjanskomu čitanju. Razklad Klebš--Gordana zachovuje eponimičnu konstrukciju. Izvorni razlog jest osoblivo razlogom za napis imena; ne ustanavjaje samostojno prednost razklad nad razloženje.\par
\noindent\emph{Reason:} Klebš adapts Clebsch in running prose so that the name does not depend on a German letter cluster outside ordinary Interslavic reading conventions. Razklad Klebš--Gordana retains the eponymic construction. The recorded reason concerns the name's spelling in particular; it does not independently establish a preference for razklad over razloženje.\par
\noindent{\footnotesize Izvorne zapisy / Original records: HIST-g078-003.}\par
\end{minipage}\par
\par\medskip\Needspace{7\baselineskip}
\noindent\begin{minipage}{\linewidth}
\hypertarget{TERM-0218}{}
\noindent\textbf{definicija Clebscha / ekvivalentne}\quad{\footnotesize\texttt{TERM-0218}}\par
\noindent\emph{Iztočnik / Source:} Clebschsche Definition / äquivalent\par
\noindent\emph{Rusky / Russian:} определение Клебша / эквивалентны\par
\noindent\emph{Ukrajinsky / Ukrainian:} означення Клебша / еквівалентні\par
\smallskip\noindent\emph{Razlog:} Definicija Clebscha zachovuje istorično ime, a ekvivalentne označaje matematičny odnos medžu formami (41) i (42). Taj zapis ostavjaje německy napis imena, različny od Klebš v drugih kandidatnyh izrazah. Slovnik zachovuje obě odluke; sama njih prisutnost ne dokazuje jednu obču politiku napisa imen.\par
\noindent\emph{Reason:} Definicija Clebscha preserves the historical name, while ekvivalentne states the mathematical relation between forms (41) and (42). This record retains the German spelling, unlike Klebš in other candidates. The dictionary preserves both decisions; their coexistence does not establish one consistent policy for names.\par
\noindent{\footnotesize Izvorne zapisy / Original records: HIST-g035-009.}\par
\end{minipage}\par
\par\medskip\Needspace{7\baselineskip}
\noindent\begin{minipage}{\linewidth}
\hypertarget{TERM-0220}{}
\noindent\textbf{konjugovane faktory}\quad{\footnotesize\texttt{TERM-0220}}\par
\noindent\emph{Iztočnik / Source:} Conjugate factors\par
\noindent\emph{Rusky / Russian:} сопряженные множители\par
\noindent\emph{Ukrajinsky / Ukrainian:} спряжені множники\par
\smallskip\noindent\emph{Razlog:} Konjugovane faktory upotrěbjaje medžunarodnu osnovu konjug- v kontekstu algebraičnogo razširjenja. Ruska i ukrajinska paralela upotrěbjajut slovjanske osnovy za isti odnos. Izvorni razlog daje prednost jasnoći medžunarodnogo termina, bez čislenogo rangovanja tih dvuh načinov imenovanja.\par
\noindent\emph{Reason:} Konjugovane faktory uses the international conjug- stem in the context of an algebraic extension. The Russian and Ukrainian parallels use Slavic stems for the same relation. The original reason favours the clarity of the international term without numerically ranking the two approaches to naming it.\par
\noindent{\footnotesize Izvorne zapisy / Original records: HIST-g091-009, HIST-g092-078.}\par
\end{minipage}\par
\par\medskip\Needspace{7\baselineskip}
\noindent\begin{minipage}{\linewidth}
\hypertarget{TERM-0222}{}
\noindent\textbf{poneže po punktu 3 to kolco klasov ostatkov jest izomorfno}\quad{\footnotesize\texttt{TERM-0222}}\par
\noindent\emph{Iztočnik / Source:} da nach 3 dieser Restklassenring isomorph ist\par
\noindent\emph{Rusky / Russian:} поскольку по пункту 3 это кольцо классов вычетов изоморфно\par
\noindent\emph{Ukrajinsky / Ukrainian:} оскільки за пунктом 3 це кільце класів лишків ізоморфне\par
\smallskip\noindent\emph{Razlog:} Po punktu 3 jest neobhodima referencija v dokazě: ona upučaje na prědhodno odgovaranje prstenov klasov ostatkov. Samo poneže to kolco ... by skrylo, otkud slěduje izomorfizm. Izvorni zapis izrično veže tu referenciju s pečatnym iztočnikom; ovde jest objasnjena njezina argumentativna funkcija, ne ustanovjena nova prověrka skana.\par
\noindent\emph{Reason:} Po punktu 3 is a necessary proof reference to the preceding correspondence of residue-class rings. Merely saying because this ring ... would hide the stated basis for the isomorphism. The original record explicitly relates the reference to the printed source; this explanation identifies its argumentative function rather than claiming a new examination of the scan.\par
\noindent{\footnotesize Izvorne zapisy / Original records: HIST-g164-006.}\par
\smallskip\noindent\emph{Sravnjenje s tekstom:} V dělu 31, oddělu 8, punktu 3, prsten kvocientov jest definovany črěz elementy polja kvocientov s dopuščenyh imenovateljev. V tom samom punktu sravnjavajut se prsteny klasov ostatkov pri korespondenciji idealov. To sut različne konstrukcije. Nazvanje v punktu 4, Diskriminantny ideal pri prěhodu k prstenu kvocientov, odkazuje na prvu konstrukciju, ne na faktor-prsten po idealu. Anglijsky izraz quotient ring v ranom objasnjenju sam ne razlučaje tute smysly; definicija jest odlučajuča. Tekst takojže ima prsten, kde historijni kandidat ima kolco. Tuto sravnjenje ustanavjaje lokalny smysl i aktualnu formu, ne dokazuje věrnost vsih pasaž izdanja.\par
\noindent\emph{Textual comparison:} In Paper 31, section 8, point 3, the ring is defined through elements of a fraction field with specified admissible denominators. The same point discusses residue-class rings in a correspondence between ideals. These are different constructions. The heading in point 4, Diskriminantny ideal pri prěhodu k prstenu kvocientov, refers to the first construction, not to a quotient by an ideal. The early explanation's English label quotient ring cannot by itself disambiguate the senses; the definition is decisive. The text also has prsten where the historical candidate has kolco. This comparison establishes a local sense and current wording, not the fidelity of every passage in the edition.\par
\noindent{\footnotesize TEXT-COMP-03; pasaže / passages: 18932, 19003; \texttt{CORPUS\_FORM\_COMPARISONS.json}.}\par
\end{minipage}\par
\par\medskip\Needspace{7\baselineskip}
\noindent\begin{minipage}{\linewidth}
\hypertarget{TERM-0223}{}
\noindent\textbf{pri tom drži}\quad{\footnotesize\texttt{TERM-0223}}\par
\noindent\emph{Iztočnik / Source:} Dabei gilt\par
\noindent\emph{Rusky / Russian:} при этом выполняется\par
\noindent\emph{Ukrajinsky / Ukrainian:} при цьому виконується\par
\smallskip\noindent\emph{Razlog:} Pri tom drži jest kratka konstrukcija za tvrđenje, ktore važi pri razmatranyh uslovjah. Izvorni razlog daje prednost běžnoj proze dokaza bez proměny sadržanja. On ne sodržuje sravnjenje drugih medžuslovjanskyh konstrukcij za platnost tvrđenja, zato ne opravdava izključivu prednost glagola držati.\par
\noindent\emph{Reason:} Pri tom drži is a short construction for an assertion holding under the conditions being discussed. The original reason favours ordinary proof prose without altering the assertion. It provides no comparison with other Interslavic constructions for validity and therefore does not establish an exclusive preference for držati.\par
\noindent{\footnotesize Izvorne zapisy / Original records: HIST-g021-010.}\par
\end{minipage}\par
\par\medskip\Needspace{7\baselineskip}
\noindent\begin{minipage}{\linewidth}
\hypertarget{TERM-0224}{}
\noindent\textbf{Tym teorem jest dokazany}\quad{\footnotesize\texttt{TERM-0224}}\par
\noindent\emph{Iztočnik / Source:} Damit ist der Satz bewiesen\par
\noindent\emph{Rusky / Russian:} Тем самым теорема доказана\par
\noindent\emph{Ukrajinsky / Ukrainian:} Цим теорему доведено\par
\smallskip\noindent\emph{Razlog:} Tym svazuje završnu větu s prědhodnym argumentom. Teorem i dokazany zachovujut imennik i rezultativnu trpnu konstrukciju už upotrěbjanu pri drugih završenkah dokazov. Kratkost jest uměstna, ibo věta ne dodavaje novu postavku, ale konstatiruje dovršenje dokaza.\par
\noindent\emph{Reason:} Tym links the closing sentence to the preceding argument. Teorem and dokazany retain the noun and resultative passive construction used in other proof endings. Brevity fits the function: the sentence introduces no new assumption but states that the proof is complete.\par
\noindent{\footnotesize Izvorne zapisy / Original records: HIST-g145-007.}\par
\end{minipage}\par
\par\medskip\Needspace{7\baselineskip}
\noindent\begin{minipage}{\linewidth}
\hypertarget{TERM-0225}{}
\noindent\textbf{togda i samo togda}\quad{\footnotesize\texttt{TERM-0225}}\par
\noindent\emph{Iztočnik / Source:} dann und nur dann\par
\noindent\emph{Rusky / Russian:} тогда и только тогда\par
\noindent\emph{Ukrajinsky / Ukrainian:} тоді й лише тоді\par
\smallskip\noindent\emph{Razlog:} Togda i samo togda zachovuje oba směra kriterija. Samo togda ne jest izpuščena čest, ktoru možno nadoměstiti običnym ako bez proměny logičnogo sadržanja. Izbor slěduje ranějšim větam o ekvivalentnosti i čini neobhodimost i dostatočnost izričnymi.\par
\noindent\emph{Reason:} Togda i samo togda preserves both directions of the criterion. The only-then component cannot be replaced by an ordinary if without changing the logical content. The choice continues earlier equivalence statements and keeps necessity and sufficiency explicit.\par
\noindent{\footnotesize Izvorne zapisy / Original records: HIST-g144-004.}\par
\end{minipage}\par
\par\medskip\Needspace{7\baselineskip}
\noindent\begin{minipage}{\linewidth}
\hypertarget{TERM-0226}{}
\noindent\textbf{predstavjenje / klas predstavjenj}\quad{\footnotesize\texttt{TERM-0226}}\par
\noindent\emph{Iztočnik / Source:} Darstellung / Darstellungsklasse\par
\noindent\emph{Rusky / Russian:} представление / класс представлений\par
\noindent\emph{Ukrajinsky / Ukrainian:} зображення / клас зображень\par
\smallskip\noindent\emph{Razlog:} Predstavjenje i klas predstavjenj imenujut objekt i njegovu klasu, ne istu stvar dvoma imenami. Izbor prodolžaje terminologiju raboty 32. Rusko predstavlenie i ukrajinsko zobražennja imajut različne leksične osnovy; medžuslovjansky prěvod zachovuje svoju familiju predstavjen-, ne podrazuměvajuči identičny korenj vo vseh jezykah.\par
\noindent\emph{Reason:} Predstavjenje and klas predstavjenj name an object and its class, not the same thing twice. The choice continues Paper 32. Russian predstavlenie and Ukrainian zobražennja have different lexical stems; Interslavic retains its predstavjen- family without implying that all the comparison languages share one root.\par
\noindent{\footnotesize Izvorne zapisy / Original records: HIST-g170-001.}\par
\end{minipage}\par
\par\medskip\Needspace{7\baselineskip}
\noindent\begin{minipage}{\linewidth}
\hypertarget{TERM-0227}{}
\noindent\textbf{predstavjenje / klasy predstavjenj}\quad{\footnotesize\texttt{TERM-0227}}\par
\noindent\emph{Iztočnik / Source:} Darstellung / Darstellungsklassen\par
\noindent\emph{Rusky / Russian:} представление / классы представлений\par
\noindent\emph{Ukrajinsky / Ukrainian:} представлення / класи представлень\par
\smallskip\noindent\emph{Razlog:} Predstavjenje jest izbrano kako široko razumějemy slovjansky izraz vměsto novogo užšego novotvora. Klasy predstavjenj zachovuje množinu klas. Ukrajinska paralela v tom zapisu jest predstavlennja, različna od zobražennja v drugih zapisah; sravnjenje zato ne može byti opisano kako jednolično vo vseh pasažah.\par
\noindent\emph{Reason:} Predstavjenje is chosen as a broadly recognizable Slavic expression rather than a new, narrower coinage. Klasy predstavjenj preserves plural classes. This record's Ukrainian parallel uses predstavlennja, unlike zobražennja elsewhere, so the comparative evidence cannot be described as uniform across passages.\par
\noindent{\footnotesize Izvorne zapisy / Original records: HIST-g173-003.}\par
\end{minipage}\par
\par\medskip\Needspace{7\baselineskip}
\noindent\begin{minipage}{\linewidth}
\hypertarget{TERM-0228}{}
\noindent\textbf{predstavjenje / teorija predstavjenj}\quad{\footnotesize\texttt{TERM-0228}}\par
\noindent\emph{Iztočnik / Source:} Darstellung / Darstellungstheorie\par
\noindent\emph{Rusky / Russian:} представление / теория представлений\par
\noindent\emph{Ukrajinsky / Ukrainian:} зображення / теорія зображень\par
\smallskip\noindent\emph{Razlog:} Teorija predstavjenj tvori nazvanje oblasti iz istogo imennika kako predstavjenje. Izvorni zapis izrično daje prednost toj familiji pred odvojenym reprezentacija. Ruska i ukrajinska paralela so svojimi vlastnymi osnovami podpirajut istu matematičnu funkciju bez potrěby za identičnym korenjem.\par
\noindent\emph{Reason:} Teorija predstavjenj forms the discipline's name from the same noun as predstavjenje. The record explicitly prefers that family to introducing a separate reprezentacija. Russian and Ukrainian support the same mathematical function through their respective stems, without requiring identical roots.\par
\noindent{\footnotesize Izvorne zapisy / Original records: HIST-g193-001.}\par
\end{minipage}\par
\par\medskip\Needspace{7\baselineskip}
\noindent\begin{minipage}{\linewidth}
\hypertarget{TERM-0229}{}
\noindent\textbf{predstavjenje / klas / transformovane}\quad{\footnotesize\texttt{TERM-0229}}\par
\noindent\emph{Iztočnik / Source:} Darstellung / Klasse / Transformierten\par
\noindent\emph{Rusky / Russian:} представление / класс / преобразованные\par
\noindent\emph{Ukrajinsky / Ukrainian:} зображення / клас / перетворені\par
\smallskip\noindent\emph{Razlog:} Predstavjenje, klas i transformovane označajut tri različne urovnje opisa: dano predstavjenje, njegovu klasu ekvivalentnosti ili podobnosti, i transformovane objekty, ktore naležet k toj klasě. Izbor ne svodi član klase i samu klasu na isto ime. Tuto različenje jest neobhodimo za čitanje argumenta o predstavjenjah.\par
\noindent\emph{Reason:} Predstavjenje, klas, and transformovane distinguish three levels of description: a representation, its equivalence or similarity class, and the transformed objects belonging to that class. A member and the class itself are not collapsed into one name. The distinction is needed to follow the representation-theoretic argument.\par
\noindent{\footnotesize Izvorne zapisy / Original records: HIST-g167-003.}\par
\end{minipage}\par
\par\medskip\Needspace{7\baselineskip}
\noindent\begin{minipage}{\linewidth}
\hypertarget{TERM-0230}{}
\noindent\textbf{predstavjenje kako direktna suma primarnyh kolc}\quad{\footnotesize\texttt{TERM-0230}}\par
\noindent\emph{Iztočnik / Source:} Darstellung als direkte Summe primärer Ringe\par
\noindent\emph{Rusky / Russian:} представление в виде прямой суммы первичных колец\par
\noindent\emph{Ukrajinsky / Ukrainian:} подання як пряма сума первинних кілець\par
\smallskip\noindent\emph{Razlog:} Predstavjenje kako direktna suma imenuje način razloženja, ne nužno predstavjenje grupy v smislu teorije predstavjenij. Primarnyh zachovuje paralelu s primarnymi idealami, a direktna suma prodolžaje termin iz odděla 2. Istorična forma kolc ostavaje zapisom kandidata; njezina prisutnost ne ustanavjaje aktualnu prednost pred prstenov.\par
\noindent\emph{Reason:} Predstavjenje kako direktna suma names a mode of decomposition, not necessarily a group representation in the representation-theoretic sense. Primarnyh parallels primary ideals, and direktna suma continues section 2. Historical kolc remains the candidate's wording rather than evidence of a current preference over prstenov.\par
\noindent{\footnotesize Izvorne zapisy / Original records: HIST-g122-001.}\par
\end{minipage}\par
\par\medskip\Needspace{7\baselineskip}
\noindent\begin{minipage}{\linewidth}
\hypertarget{TERM-0231}{}
\noindent\textbf{predstavjenje kako najmenše obče mnogokratno}\quad{\footnotesize\texttt{TERM-0231}}\par
\noindent\emph{Iztočnik / Source:} Darstellung als kleinstes gemeinsames Vielfaches\par
\noindent\emph{Rusky / Russian:} представление как наименьшее общее кратное\par
\noindent\emph{Ukrajinsky / Ukrainian:} подання як найменше спільне кратне\par
\smallskip\noindent\emph{Razlog:} Najmenše obče mnogokratno zachovuje jezyk aritmetiky idealov, v ktorej dokaz organizovan jest po děliteljnosti. Tiha zaměna slovom prěsek mogla by skryti, kako Noether sama organizuje odnose medžu idealami. Prěvod zato zachovuje avtorske nazvanje operacije, iako pozdnějše objasnjenje može sravniti jeho s drugym jezykom.\par
\noindent\emph{Reason:} Najmenše obče mnogokratno preserves the ideal-arithmetic vocabulary in which the proof is organized by divisibility. Silently replacing it with intersection could obscure Noether's own ordering of the relations between ideals. The translation therefore retains the author's operation name, even where a separate explanation may compare it with later terminology.\par
\noindent{\footnotesize Izvorne zapisy / Original records: HIST-g113-008.}\par
\end{minipage}\par
\par\medskip\Needspace{7\baselineskip}
\noindent\begin{minipage}{\linewidth}
\hypertarget{TERM-0232}{}
\noindent\textbf{predstavjenje kvaternionnoj grupy}\quad{\footnotesize\texttt{TERM-0232}}\par
\noindent\emph{Iztočnik / Source:} Darstellung der Quaternionengruppe\par
\noindent\emph{Rusky / Russian:} представление группы кватернионов\par
\noindent\emph{Ukrajinsky / Ukrainian:} представлення кватерніонної групи\par
\smallskip\noindent\emph{Razlog:} Predstavjenje kvaternionnoj grupy izrično govori, že matrice davajut predstavjenje grupy. Izraz ešče ne tvrdi ničto novo o polju razloženja. Zachovanje predstavjenje razděljuje dannu konstrukciju od drugyh rezultatov, ktore možut byti s njeju svezane.\par
\noindent\emph{Reason:} Predstavjenje kvaternionnoj grupy explicitly says that the matrices give a group representation. It does not yet state a result about a splitting field. Retaining representation distinguishes the construction at hand from other results that may be associated with it.\par
\noindent{\footnotesize Izvorne zapisy / Original records: HIST-g176-001.}\par
\end{minipage}\par
\par\medskip\Needspace{7\baselineskip}
\noindent\begin{minipage}{\linewidth}
\hypertarget{TERM-0233}{}
\noindent\textbf{predstavjenja konečnyh grup}\quad{\footnotesize\texttt{TERM-0233}}\par
\noindent\emph{Iztočnik / Source:} Darstellung endlicher Gruppen\par
\noindent\emph{Rusky / Russian:} представления конечных групп\par
\noindent\emph{Ukrajinsky / Ukrainian:} представлення скінченних груп\par
\smallskip\noindent\emph{Razlog:} Predstavjenja konečnyh grup upotrěbjaje predstavjenje v tehničnom smyslu teorije predstavjenij. Konečnyh kvalifikuje grupy i ostavaje izrično v nazvanju. Izbor jest prodolženjem istej familije terminov, ne novym sinonimom za običny opis ili izobraženje.\par
\noindent\emph{Reason:} Predstavjenja konečnyh grup uses representation in its technical sense. Konečnyh qualifies the groups and remains explicit in the name. The choice continues the same terminology family rather than introducing a synonym for an ordinary description or picture.\par
\noindent{\footnotesize Izvorne zapisy / Original records: HIST-g105-003.}\par
\end{minipage}\par
\par\medskip\Needspace{7\baselineskip}
\noindent\begin{minipage}{\linewidth}
\hypertarget{TERM-0234}{}
\noindent\textbf{predstavjenje s stalym imenovateljem}\quad{\footnotesize\texttt{TERM-0234}}\par
\noindent\emph{Iztočnik / Source:} Darstellung mit festem Nenner\par
\noindent\emph{Rusky / Russian:} представление с фиксированным знаменателем\par
\noindent\emph{Ukrajinsky / Ukrainian:} подання зі сталим знаменником\par
\smallskip\noindent\emph{Razlog:} S stalym imenovateljem zachovuje ograničenje, ktore jest silnějše od samej racionalnoj predstavimosti. Važno jest ne samo moči napisati racionalny izraz, ale imati fiksovany imenovatelj v razmatranom predstavjenju. Izbor ne dopuščaje izpuščenje togo uslovja kako stilistično nenadobnogo.\par
\noindent\emph{Reason:} S stalym imenovateljem retains a constraint stronger than mere rational representability. The point is not simply that a rational expression can be written, but that the representation has a fixed denominator. That qualification cannot be dropped as stylistically dispensable.\par
\noindent{\footnotesize Izvorne zapisy / Original records: HIST-g040-010.}\par
\end{minipage}\par
\par\medskip\Needspace{7\baselineskip}
\noindent\begin{minipage}{\linewidth}
\hypertarget{TERM-0235}{}
\noindent\textbf{predstavjenje slěda, normy i diskriminanta}\quad{\footnotesize\texttt{TERM-0235}}\par
\noindent\emph{Iztočnik / Source:} Darstellung von Spur, Norm und Diskriminante\par
\noindent\emph{Rusky / Russian:} представление следа, нормы и дискриминанта\par
\noindent\emph{Ukrajinsky / Ukrainian:} подання сліду, норми й дискримінанта\par
\smallskip\noindent\emph{Razlog:} Slěd, norma i diskriminant ostavajut vidimoju trojicoju objektov. Prva dva nazvanja prodolžajut terminologiju odděla 4; diskriminant zachovuje medžunarodnu osnovu. Sadržaj naslova ne trěba skratiti do neoprěděljenyh invariantov, zato že to by odstranilo točny seznam razmatranyh veličin.\par
\noindent\emph{Reason:} Slěd, norma, and diskriminant remain a visible triad of objects. The first two continue section 4, while diskriminant retains its international stem. The heading should not be shortened to unspecified invariants, which would remove the exact list of quantities under discussion.\par
\noindent{\footnotesize Izvorne zapisy / Original records: HIST-g146-001.}\par
\end{minipage}\par
\par\medskip\Needspace{7\baselineskip}
\noindent\begin{minipage}{\linewidth}
\hypertarget{TERM-0236}{}
\noindent\textbf{klasy predstavjenj}\quad{\footnotesize\texttt{TERM-0236}}\par
\noindent\emph{Iztočnik / Source:} Darstellungsklassen\par
\noindent\emph{Rusky / Russian:} классы представлений\par
\noindent\emph{Ukrajinsky / Ukrainian:} класи представлень\par
\smallskip\noindent\emph{Razlog:} Klasy predstavjenj prodolžaje terminologiju klas iz raboty 28. Izraz imenuje pojem teorije predstavjenij, ne obične vizualne obrazy. Izvorni zapis odnak dopuščaje ukrajinsko matryčne zobražennja pri matričnom predstavjenju; ta dopustimost ne měnja medžuslovjansky izbor korenja.\par
\noindent\emph{Reason:} Klasy predstavjenj continues the class terminology of Paper 28. It names a representation-theoretic concept, not ordinary visual images. The original record nevertheless permits Ukrainian matryčne zobražennja for matrix representation; that does not change the Interslavic stem choice.\par
\noindent{\footnotesize Izvorne zapisy / Original records: HIST-g133-002.}\par
\end{minipage}\par
\par\medskip\Needspace{7\baselineskip}
\noindent\begin{minipage}{\linewidth}
\hypertarget{TERM-0237}{}
\noindent\textbf{modul predstavjenja}\quad{\footnotesize\texttt{TERM-0237}}\par
\noindent\emph{Iztočnik / Source:} Darstellungsmodul\par
\noindent\emph{Rusky / Russian:} модуль представления\par
\noindent\emph{Ukrajinsky / Ukrainian:} модуль зображення\par
\smallskip\noindent\emph{Razlog:} Modul predstavjenja razkriva německe složeno slovo dvoma imennikami, ktorih odnos jest čitatelju vidimy. Odděl 15 opreděljaje taj objekt i potom sravnjaje njegove klasy s klasami predstavjenj. Drugi zapis podpira istu konstrukciju terminologijeju tretjej glavy raboty 34; odvojeno novo korenje ne jest nužno.\par
\noindent\emph{Reason:} Modul predstavjenja unpacks the German compound into two nouns whose relation is visible. Section 15 defines the object and then compares its classes with representation classes. A second record supports the same construction through chapter III of Paper 34; no separate new lexical root is required.\par
\noindent{\footnotesize Izvorne zapisy / Original records: HIST-g193-002, HIST-g196-002.}\par
\end{minipage}\par
\par\medskip\Needspace{7\baselineskip}
\noindent\begin{minipage}{\linewidth}
\hypertarget{TERM-0238}{}
\noindent\textbf{modul predstavjenja / klas modulov}\quad{\footnotesize\texttt{TERM-0238}}\par
\noindent\emph{Iztočnik / Source:} Darstellungsmodul / Modulklasse\par
\noindent\emph{Rusky / Russian:} модуль представления / класс модулей\par
\noindent\emph{Ukrajinsky / Ukrainian:} модуль зображення / клас модулів\par
\smallskip\noindent\emph{Razlog:} Modul predstavjenja označaje objekt, a klas modulov klasifikaciju takvyh objektov. Jezyk klas modulov jest važno pojmove posrědstvo k razmatranju klas idealov. On zato ne može byti izpuščeny ili spojen s nazvanjem samogo modula.\par
\noindent\emph{Reason:} Modul predstavjenja names an object; klas modulov names a class of such objects. The language of module classes provides the conceptual connection to ideal classes and cannot be omitted or merged with the name of the module itself.\par
\noindent{\footnotesize Izvorne zapisy / Original records: HIST-g177-004.}\par
\end{minipage}\par
\par\medskip\Needspace{7\baselineskip}
\noindent\begin{minipage}{\linewidth}
\hypertarget{TERM-0239}{}
\noindent\textbf{moduly predstavjenja}\quad{\footnotesize\texttt{TERM-0239}}\par
\noindent\emph{Iztočnik / Source:} Darstellungsmoduln\par
\noindent\emph{Rusky / Russian:} модули представлений\par
\noindent\emph{Ukrajinsky / Ukrainian:} модулі зображення\par
\smallskip\noindent\emph{Razlog:} Moduly predstavjenja daje množinu osnovnoj konstrukciji modul predstavjenja. Izbor drži istu medžuslovjansku familiju predstavjen-, pri čem ruska i ukrajinska paralela zachovujut svoje imena predstavjenja. Zapis podpira kontinuitet terminologije, ne potrěbu za doslovno identičnymi morfemami vo vseh jezykah.\par
\noindent\emph{Reason:} Moduly predstavjenja supplies the plural of modul predstavjenja. It retains the Interslavic predstavjen- family while the Russian and Ukrainian parallels keep their respective names for representation. The record supports terminological continuity, not a requirement for identical morphemes in every language.\par
\noindent{\footnotesize Izvorne zapisy / Original records: HIST-g195-004.}\par
\end{minipage}\par
\par\medskip\Needspace{7\baselineskip}
\noindent\begin{minipage}{\linewidth}
\hypertarget{TERM-0240}{}
\noindent\textbf{teorija predstavjenj}\quad{\footnotesize\texttt{TERM-0240}}\par
\noindent\emph{Iztočnik / Source:} Darstellungstheorie\par
\noindent\emph{Rusky / Russian:} теория представлений\par
\noindent\emph{Ukrajinsky / Ukrainian:} теорія зображень\par
\smallskip\noindent\emph{Razlog:} Teorija predstavjenj prodolžaje familiju iz rabot 32 i 33. Dva izvorne zapisy zachovujut razliku medžu ruskoju i ukrajinskoju leksikoju i ostavjajut lokalne matrične upotrěby za kontekstualno sravnjenje. Zato jedna obča nazva oblasti ne jest dovoljnym razlogom avtomatično izravnati vsaku konstrukciju s Darstellung.\par
\noindent\emph{Reason:} Teorija predstavjenj continues the family used in Papers 32 and 33. The two original records preserve the lexical difference between Russian and Ukrainian and leave local matrix usages for contextual comparison. A common discipline name is therefore not enough to flatten every construction involving Darstellung into one wording.\par
\noindent{\footnotesize Izvorne zapisy / Original records: HIST-g177-001, HIST-g178-001.}\par
\end{minipage}\par
\par\medskip\Needspace{7\baselineskip}
\noindent\begin{minipage}{\linewidth}
\hypertarget{TERM-0241}{}
\noindent\textbf{diskriminantny ideal pri prěhodu k kolcu kvocientov}\quad{\footnotesize\texttt{TERM-0241}}\par
\noindent\emph{Iztočnik / Source:} Das Diskriminantenideal beim Übergang zum Quotientenring\par
\noindent\emph{Rusky / Russian:} дискриминантный идеал при переходе к кольцу частных\par
\noindent\emph{Ukrajinsky / Ukrainian:} дискримінантний ідеал при переході до кільця часток\par
\smallskip\noindent\emph{Razlog:} Naslov zachovuje diskriminantny ideal i imenuje operaciju prěhoda k prstenu kvocientov. Izvorni razlog prodolžaje prědhodno imenovanje togo prstena i daje prednost kratkoj naslovnoj konstrukciji. Smysl kvocientov trěba ustanoviti po definiciji v dělu, ne samo po podobnosti s anglijskym quotient ring; ruska i ukrajinska paralela ovde izrično govore o častnyh ili častkah.\par
\noindent\emph{Reason:} The heading retains discriminant ideal and names passage to the ring of quotients. The recorded reason continues the earlier name for that ring and favours a compact heading. The sense of quotients must be determined from the work's definition, not merely from resemblance to English quotient ring; the Russian and Ukrainian parallels here explicitly use their vocabulary for fractions or quotients.\par
\noindent{\footnotesize Izvorne zapisy / Original records: HIST-g161-001.}\par
\smallskip\noindent\emph{Sravnjenje s tekstom:} V dělu 31, oddělu 8, punktu 3, prsten kvocientov jest definovany črěz elementy polja kvocientov s dopuščenyh imenovateljev. V tom samom punktu sravnjavajut se prsteny klasov ostatkov pri korespondenciji idealov. To sut različne konstrukcije. Nazvanje v punktu 4, Diskriminantny ideal pri prěhodu k prstenu kvocientov, odkazuje na prvu konstrukciju, ne na faktor-prsten po idealu. Anglijsky izraz quotient ring v ranom objasnjenju sam ne razlučaje tute smysly; definicija jest odlučajuča. Tekst takojže ima prsten, kde historijni kandidat ima kolco. Tuto sravnjenje ustanavjaje lokalny smysl i aktualnu formu, ne dokazuje věrnost vsih pasaž izdanja.\par
\noindent\emph{Textual comparison:} In Paper 31, section 8, point 3, the ring is defined through elements of a fraction field with specified admissible denominators. The same point discusses residue-class rings in a correspondence between ideals. These are different constructions. The heading in point 4, Diskriminantny ideal pri prěhodu k prstenu kvocientov, refers to the first construction, not to a quotient by an ideal. The early explanation's English label quotient ring cannot by itself disambiguate the senses; the definition is decisive. The text also has prsten where the historical candidate has kolco. This comparison establishes a local sense and current wording, not the fidelity of every passage in the edition.\par
\noindent{\footnotesize TEXT-COMP-03; pasaže / passages: 18932, 19003; \texttt{CORPUS\_FORM\_COMPARISONS.json}.}\par
\end{minipage}\par
\par\medskip\Needspace{7\baselineskip}
\noindent\begin{minipage}{\linewidth}
\hypertarget{TERM-0243}{}
\noindent\textbf{vyše podana shema ... dokazana}\quad{\footnotesize\texttt{TERM-0243}}\par
\noindent\emph{Iztočnik / Source:} Das oben stehende Schema ... ist somit bewiesen\par
\noindent\emph{Rusky / Russian:} доказана приведенная выше схема ...\par
\noindent\emph{Ukrajinsky / Ukrainian:} наведену вище схему ... доведено\par
\smallskip\noindent\emph{Razlog:} Vyše podana shema ... dokazana veže zaključok s konkretnoju shemoju, ktora jest dana raněje. Dokazana soglasuje se s ženskim imennikom shema. Izvorni razlog daje prednost običnoj završnoj věte dokaza bez odstranjenja referencije k tomu, čto bylo dokazano.\par
\noindent\emph{Reason:} Vyše podana shema ... dokazana links the conclusion to the particular scheme given earlier. Dokazana agrees with feminine shema. The recorded reason favours an ordinary proof-closing sentence while retaining the reference to what has been proved.\par
\noindent{\footnotesize Izvorne zapisy / Original records: HIST-g019-016.}\par
\end{minipage}\par
\par\medskip\Needspace{7\baselineskip}
\noindent\begin{minipage}{\linewidth}
\hypertarget{TERM-0244}{}
\noindent\textbf{Znak $\equiv$ znači ...}\quad{\footnotesize\texttt{TERM-0244}}\par
\noindent\emph{Iztočnik / Source:} Das Zeichen $\equiv$ bedeutet ...\par
\noindent\emph{Rusky / Russian:} Знак $\equiv$ означает ...\par
\noindent\emph{Ukrajinsky / Ukrainian:} Знак $\equiv$ означає ...\par
\smallskip\noindent\emph{Razlog:} Znak $\equiv$ znači ... zachovuje avtorske izrično objasnjenje znaka. Objasnjenje pripada podnožnoj primětkě pri odgovarjajučej formulě; samo obče znanje znaka kongruencije ili ekvivalentnosti ne može njegovo lokalno značenje nadoměstiti. Prěvod zachovuje i simbol, i jeho vezu s konkretnym opreděljenjem.\par
\noindent\emph{Reason:} Znak $\equiv$ znači ... preserves the author's explicit explanation of the symbol. The explanation belongs to the footnote attached to the relevant display; general familiarity with congruence or equivalence notation cannot replace its local meaning. The translation retains both the symbol and its connection to the particular definition.\par
\noindent{\footnotesize Izvorne zapisy / Original records: HIST-g024-011.}\par
\end{minipage}\par
\par\medskip\Needspace{7\baselineskip}
\noindent\begin{minipage}{\linewidth}
\hypertarget{TERM-0245}{}
\noindent\textbf{prosty ideal, asociovany s nulovym idealom}\quad{\footnotesize\texttt{TERM-0245}}\par
\noindent\emph{Iztočnik / Source:} das zum Nullideal gehörige Primideal\par
\noindent\emph{Rusky / Russian:} простой идеал, ассоциированный с нулевым идеалом\par
\noindent\emph{Ukrajinsky / Ukrainian:} простий ідеал, асоційований із нульовим ідеалом\par
\smallskip\noindent\emph{Razlog:} Asociovany s nulovym idealom izjasnjaje tehničny odnos asocijacije. Doslovno prinaležeči nulovomu idealu moglo by byti razuměno kako obična pripadnost ili vključennost. Izvorni izbor medžunarodnogo asociovany izběgaje takvu nejasnost bez stvarjanja novogo domorodnogo složenogo slova.\par
\noindent\emph{Reason:} Asociovany s nulovym idealom makes the technical association relation explicit. A literal belonging to the zero ideal could suggest ordinary membership or inclusion. The recorded preference for international asociovany avoids that ambiguity without inventing a new native compound.\par
\noindent{\footnotesize Izvorne zapisy / Original records: HIST-g141-003.}\par
\end{minipage}\par
\par\medskip\Needspace{7\baselineskip}
\noindent\begin{minipage}{\linewidth}
\hypertarget{TERM-0246}{}
\noindent\textbf{Dedekind--Mertensov teorem}\quad{\footnotesize\texttt{TERM-0246}}\par
\noindent\emph{Iztočnik / Source:} Dedekind-Mertensscher Satz\par
\noindent\emph{Rusky / Russian:} теорема Дедекинда--Мертенса\par
\noindent\emph{Ukrajinsky / Ukrainian:} теорема Дедекінда--Мертенса\par
\smallskip\noindent\emph{Razlog:} Dedekind--Mertensov teorem zachovuje ime teorema vměsto opisnoj parafrazy. Zapis razlikuje imena v běžnoj proze, ktore kirilična verzija prěpisuje, od originalnyh bibliografičnyh nazvanij i citatnyh oznak, ktore ostavajut v latiničnom zapisu. Izbor skripta zavisi od funkcije imena.\par
\noindent\emph{Reason:} Dedekind--Mertensov teorem retains the theorem's name rather than paraphrasing it. The record distinguishes names in running prose, which the Cyrillic version transliterates, from original bibliographical titles and citation labels, which remain in Latin script. Script selection depends on the name's function.\par
\noindent{\footnotesize Izvorne zapisy / Original records: HIST-g087-009, HIST-g092-029.}\par
\end{minipage}\par
\par\medskip\Needspace{7\baselineskip}
\noindent\begin{minipage}{\linewidth}
\hypertarget{TERM-0247}{}
\noindent\textbf{teorem Dedekinda--Mertensa}\quad{\footnotesize\texttt{TERM-0247}}\par
\noindent\emph{Iztočnik / Source:} Dedekind-Mertensscher Satz\par
\noindent\emph{Rusky / Russian:} теорема Дедекинда--Мертенса\par
\noindent\emph{Ukrajinsky / Ukrainian:} теорема Дедекінда--Мертенса\par
\smallskip\noindent\emph{Razlog:} Teorem Dedekinda--Mertensa upotrěbjaje genitivy imen za imenovanje istogo teorema. Latiničny prěvod zachovuje latiničny napis; kiriličny tekst prěpisuje genitivne formy v běžnom tekstu, ne originalne bibliografične citaty. Ta konstrukcija jest dokumentovana alternativa pridavnikovomu Dedekind--Mertensov teorem.\par
\noindent\emph{Reason:} Teorem Dedekinda--Mertensa uses genitives of the names for the same theorem. The Latin version retains Latin spelling, while Cyrillic prose transliterates those running-text forms rather than original bibliographical quotations. This is a documented alternative to the adjectival Dedekind--Mertensov teorem.\par
\noindent{\footnotesize Izvorne zapisy / Original records: HIST-g096-006.}\par
\end{minipage}\par
\par\medskip\Needspace{7\baselineskip}
\noindent\begin{minipage}{\linewidth}
\hypertarget{TERM-0248}{}
\noindent\textbf{Dedekindovo definovanje / Dedekind-Henselovogo tipa}\quad{\footnotesize\texttt{TERM-0248}}\par
\noindent\emph{Iztočnik / Source:} Dedekindsche Definition / Dedekind-Henselscher Art\par
\noindent\emph{Rusky / Russian:} дедекиндово определение / дедекинд-гензелевского типа\par
\noindent\emph{Ukrajinsky / Ukrainian:} дедекіндове означення / дедекінд-гензелівського типу\par
\smallskip\noindent\emph{Razlog:} Zapis daje prednost poznavajemosti imen Dedekind i Hensel. Za medžuslovjansky tekst on naměrno zachovuje latiničny Dedekind daže v kiriličnom okruženju. To jest druga odluka než prěpis imen v běžnom tekstu v zapisah o Dedekind--Mertensovom teoremu; razlika trěbuje izrično objasnjenje, a ne tvrdnju, že istorijni zapisy imali jednu politiku.\par
\noindent\emph{Reason:} The record prioritizes recognizability of Dedekind and Hensel. For Interslavic it deliberately retains Latin Dedekind even in Cyrillic surroundings. That differs from transliterating running-text names in the Dedekind-Mertens entries. The difference requires explicit explanation rather than a claim that the historical records followed one policy.\par
\noindent{\footnotesize Izvorne zapisy / Original records: HIST-g213-016.}\par
\end{minipage}\par
\par\medskip\Needspace{7\baselineskip}
\noindent\begin{minipage}{\linewidth}
\hypertarget{TERM-0249}{}
\noindent\textbf{Dedekindov poslědok I/II}\quad{\footnotesize\texttt{TERM-0249}}\par
\noindent\emph{Iztočnik / Source:} Dedekindsche Folgerung I/II\par
\noindent\emph{Rusky / Russian:} следствие Dedekind I/II\par
\noindent\emph{Ukrajinsky / Ukrainian:} наслідок Dedekind I/II\par
\smallskip\noindent\emph{Razlog:} Poslědok zachovuje funkciju lokalno izvedenogo tvrđenja; ne pretvarjaje jeho v novo ime samostojnogo teorema. I i II identifikujut dva poslědka. Dedekind ostavaje poznavajemy v paralelnyh izrazah, ale točno gramatično oblikovanje imena ne jest identično vo vseh trěh jezykah.\par
\noindent\emph{Reason:} Poslědok preserves the role of a locally derived consequence rather than turning it into the title of a separate named theorem. I and II identify the two corollaries. Dedekind remains recognizable in the parallels, although the grammatical treatment of the name is not identical in all three languages.\par
\noindent{\footnotesize Izvorne zapisy / Original records: HIST-g108-009.}\par
\end{minipage}\par
\par\medskip\Needspace{7\baselineskip}
\noindent\begin{minipage}{\linewidth}
\hypertarget{TERM-0250}{}
\noindent\textbf{Dedekindov poslědok II}\quad{\footnotesize\texttt{TERM-0250}}\par
\noindent\emph{Iztočnik / Source:} Dedekindsche Folgerung II\par
\noindent\emph{Rusky / Russian:} следствие Dedekind II\par
\noindent\emph{Ukrajinsky / Ukrainian:} наслідок Dedekind II\par
\smallskip\noindent\emph{Razlog:} Dedekindov poslědok II odpovada imenovanju v prvom oddělu raboty 30, kde jest tvrđenje formulovano. II ostavaje oznakoju togo samogo poslědka. Zapis takože dokumentuje latiničny Dedekind v ruskoj i ukrajinskoj paralelě; ne trěba tu vlastnost pripisati vsakoj upotrěbě imena v korpusu.\par
\noindent\emph{Reason:} Dedekindov poslědok II matches the naming in Paper 30, section 1, where the assertion is stated. II continues to identify that same corollary. The record also documents Latin Dedekind in the Russian and Ukrainian parallels; that feature should not be generalized to every use of the name in the corpus.\par
\noindent{\footnotesize Izvorne zapisy / Original records: HIST-g115-010.}\par
\end{minipage}\par
\par\medskip\Needspace{7\baselineskip}
\noindent\begin{minipage}{\linewidth}
\hypertarget{TERM-0251}{}
\noindent\textbf{defekt svijanja}\quad{\footnotesize\texttt{TERM-0251}}\par
\noindent\emph{Iztočnik / Source:} Defekt der Faltung\par
\noindent\emph{Rusky / Russian:} дефект свертки\par
\noindent\emph{Ukrajinsky / Ukrainian:} дефект згортання\par
\smallskip\noindent\emph{Razlog:} Defekt jest cěločisleny parametr $\lambda$, ne obična ocěna nedostatka. Tri zapisy vežut jego s formulami (41), (44) i (62), i s indukcijnym krokom od $\lambda$ k $\lambda-1$. Defekt svijanja zato drži istu čislenu veličinu vidimoju v proze i formulah.\par
\noindent\emph{Reason:} Defekt is the integer parameter $\lambda$, not an ordinary judgment of deficiency. Three records connect it with formulas (41), (44), and (62), and with induction from $\lambda$ to $\lambda-1$. Defekt svijanja keeps the same numerical quantity visible in both prose and formulas.\par
\noindent{\footnotesize Izvorne zapisy / Original records: HIST-g031-007, HIST-g035-005, HIST-g037-005.}\par
\end{minipage}\par
\par\medskip\Needspace{7\baselineskip}
\noindent\begin{minipage}{\linewidth}
\hypertarget{TERM-0252}{}
\noindent\textbf{definujuči ideal}\quad{\footnotesize\texttt{TERM-0252}}\par
\noindent\emph{Iztočnik / Source:} definierendes Ideal\par
\noindent\emph{Rusky / Russian:} определяющий идеал\par
\noindent\emph{Ukrajinsky / Ukrainian:} визначальний ідеал\par
\smallskip\noindent\emph{Razlog:} Definujuči ideal jest ideal relacij, ktorymi jest opreděljen prsten ili faktor. Pričestny pridavnik zachovuje jego definujuču funkciju, na ktoru iztočnik ponovjeno upučaje. Izraz relacijny ideal ne jest izbrany, ibo by odstranil tu izričnu funkciju iz nazvanja.\par
\noindent\emph{Reason:} Definujuči ideal is the ideal of relations defining the ring or quotient. The participial adjective retains the defining function repeatedly emphasized in the source. Relational ideal is not chosen because it would remove that explicit function from the name.\par
\noindent{\footnotesize Izvorne zapisy / Original records: HIST-g213-004.}\par
\end{minipage}\par
\par\medskip\Needspace{7\baselineskip}
\noindent\begin{minipage}{\linewidth}
\hypertarget{TERM-0253}{}
\noindent\textbf{definujųcy ideal}\quad{\footnotesize\texttt{TERM-0253}}\par
\noindent\emph{Iztočnik / Source:} definierendes Ideal\par
\noindent\emph{Rusky / Russian:} определяющий идеал\par
\noindent\emph{Ukrajinsky / Ukrainian:} визначальний ідеал\par
\smallskip\noindent\emph{Razlog:} Definujųcy ideal prodolžaje familiju definujučego ideala iz raboty 34. Razlog dava prednost medžunarodnoj osnově s pričestnoju funkcijeju vměsto novogo doslovnogo slovotvora. To ne jest osobny dokaz za etimologičnu bukvu ų ili za konec -cy v zachovanom kandidatu; leksičny izbor i gramatična forma trěbujut različne osnovanja.\par
\noindent\emph{Reason:} Definujųcy ideal continues Paper 34's defining-ideal family. The reason favours an international stem with a participial function rather than a new native calque. It does not separately establish the etymological ų or ending -cy in the preserved candidate; lexical choice and grammatical form require different grounds.\par
\noindent{\footnotesize Izvorne zapisy / Original records: HIST-g206-004.}\par
\end{minipage}\par
\par\medskip\Needspace{7\baselineskip}
\noindent\begin{minipage}{\linewidth}
\hypertarget{TERM-0254}{}
\noindent\textbf{algebraično zatvorjeno polje izvedeno iz P}\quad{\footnotesize\texttt{TERM-0254}}\par
\noindent\emph{Iztočnik / Source:} der aus P abgeleitete algebraisch-abgeschlossene Körper\par
\noindent\emph{Rusky / Russian:} алгебраически замкнутое поле, выведенное из P\par
\noindent\emph{Ukrajinsky / Ukrainian:} алгебраїчно замкнене поле, виведене з P\par
\smallskip\noindent\emph{Razlog:} Izvedeno iz $P$ zachovuje avtorske izrično odnosenje polja k $P$. Zaměna cěloj konstrukcije samo simbolom algebraičnogo zamknjenja odstranila by to slovno objasnjenje. Algebraično zatvorjeno ostavaje kvalifikacijeju polja, a ne zaměnoju za odnos, ktory vyražaje izvedeno.\par
\noindent\emph{Reason:} Izvedeno iz $P$ retains the author's explicit relation between the field and $P$. Replacing the whole construction by algebraic-closure notation alone would remove that verbal explanation. Algebraično zatvorjeno qualifies the field but does not replace the relation expressed by derived from.\par
\noindent{\footnotesize Izvorne zapisy / Original records: HIST-g125-007.}\par
\end{minipage}\par
\par\medskip\Needspace{7\baselineskip}
\noindent\begin{minipage}{\linewidth}
\hypertarget{TERM-0255}{}
\noindent\textbf{determinant matrice P jest jedinica}\quad{\footnotesize\texttt{TERM-0255}}\par
\noindent\emph{Iztočnik / Source:} Determinante der Matrix P eine Einheit\par
\noindent\emph{Rusky / Russian:} определитель матрицы P является единицей\par
\noindent\emph{Ukrajinsky / Ukrainian:} визначник матриці P є одиницею\par
\smallskip\noindent\emph{Razlog:} Jest jedinica zachovuje uslovje dopustimosti proměny bazy: determinant matrice $P$ trěbuje byti invertibilny, ne prosto nenulovy. Jedinica ovde imaje tehničny smysl; izraz ne tvrdi, že determinant mora byti doslovno čislo 1. Izvorni zapis svazuje determinant i jedinica s už danymi uslovjami odděla.\par
\noindent\emph{Reason:} Jest jedinica retains the admissibility condition for a basis change: the determinant of $P$ must be a unit, not merely nonzero. Unit has its technical sense here; the phrase does not say that the determinant must literally equal the number 1. The record connects determinant and unit with the section's existing assumptions.\par
\noindent{\footnotesize Izvorne zapisy / Original records: HIST-g132-003.}\par
\end{minipage}\par
\par\medskip\Needspace{7\baselineskip}
\noindent\begin{minipage}{\linewidth}
\hypertarget{TERM-0256}{}
\noindent\textbf{determinanta hiperkompleksnogo sistema}\quad{\footnotesize\texttt{TERM-0256}}\par
\noindent\emph{Iztočnik / Source:} Determinante eines hyperkomplexen Systems\par
\noindent\emph{Rusky / Russian:} определитель гиперкомплексной системы\par
\noindent\emph{Ukrajinsky / Ukrainian:} визначник гіперкомплексної системи\par
\smallskip\noindent\emph{Razlog:} Determinanta jest izbrana kako medžunarodno nazvanje, dokole ruska i ukrajinska paralela upotrěbjajut svoje slovjanske nazvanja determinanty. Razlog izběgaje prednost jednoj nacionalnoj formě. Hiperkompleksnogo sistema zachovuje avtorsky tip objekta; sama medžunarodna osnova ne rěšaje vsaku gramatičnu razliku medžu determinant i determinanta v drugih zapisah.\par
\noindent\emph{Reason:} Determinanta is selected as an international name while Russian and Ukrainian use their respective Slavic determinant terms. The reason avoids privileging one national form. Hiperkompleksnogo sistema preserves the author's kind of object; the international stem alone does not settle every grammatical difference between determinant and determinanta elsewhere.\par
\noindent{\footnotesize Izvorne zapisy / Original records: HIST-g201-001.}\par
\end{minipage}\par
\par\medskip\Needspace{7\baselineskip}
\noindent\begin{minipage}{\linewidth}
\hypertarget{TERM-0257}{}
\noindent\textbf{determinant $|a_{ik}|$}\quad{\footnotesize\texttt{TERM-0257}}\par
\noindent\emph{Iztočnik / Source:} Determinante $|a_{ik}|$\par
\noindent\emph{Rusky / Russian:} детерминант $|a_{ik}|$\par
\noindent\emph{Ukrajinsky / Ukrainian:} детермінант $|a_{ik}|$\par
\smallskip\noindent\emph{Razlog:} Determinant imenuje objekt, ktory jest zapisan kako $|a_{ik}|$. Točno zachovanje indeksov i determinantnyh črt jest čestju točnosti termina v toj pasaži. Izvorni razlog dava neposrědno tehnično odgovaranje, bez odvojenogo sravnjenja sinonimov.\par
\noindent\emph{Reason:} Determinant names the object written $|a_{ik}|$. Preserving the indices and determinant bars is part of naming it accurately in this passage. The original reason gives a direct technical equivalent without a separate comparison of synonyms.\par
\noindent{\footnotesize Izvorne zapisy / Original records: HIST-g053-009.}\par
\end{minipage}\par
\par\medskip\Needspace{7\baselineskip}
\noindent\begin{minipage}{\linewidth}
\hypertarget{TERM-0258}{}
\noindent\textbf{determinanty / stupenj}\quad{\footnotesize\texttt{TERM-0258}}\par
\noindent\emph{Iztočnik / Source:} Determinanten / Grad\par
\noindent\emph{Rusky / Russian:} детерминанты / степень\par
\noindent\emph{Ukrajinsky / Ukrainian:} детермінанти / степінь\par
\smallskip\noindent\emph{Razlog:} Stupenj jest v tom zapisu upotrěbjen za porjadok determinanty. Jego smysl zato ne slěduje avtomatično iz drugih pasažov, kde stupenj označaje stepen polinoma. Determinanty zachovuje medžunarodno nazvanje objektov; različenje značenij slova Grad mora ostati kontekstualno.\par
\noindent\emph{Reason:} Stupenj is used here for the order of a determinant. Its sense must not be imported automatically from passages where degree concerns a polynomial. Determinanty retains the international object name, while the senses of Grad remain context-dependent.\par
\noindent{\footnotesize Izvorne zapisy / Original records: HIST-g030-005.}\par
\end{minipage}\par
\par\medskip\Needspace{7\baselineskip}
\noindent\begin{minipage}{\linewidth}
\hypertarget{TERM-0259}{}
\noindent\textbf{determinanty iz vsakih n baznyh elementov i jih konjugovanyh elementov}\quad{\footnotesize\texttt{TERM-0259}}\par
\noindent\emph{Iztočnik / Source:} Determinanten aus je n Basiselementen und ihren Konjugierten\par
\noindent\emph{Rusky / Russian:} определители из каждых n базисных элементов и их сопряженных\par
\noindent\emph{Ukrajinsky / Ukrainian:} визначники з кожних n базисних елементів та їхніх спряжених\par
\smallskip\noindent\emph{Razlog:} Iz vsakih $n$ baznyh elementov i jih konjugovanyh elementov zachovuje točno sostavjenje Hilbertovy determinanty. Vsakih $n$ ne može byti skračeno do neoprěděljenogo množenstva elementov. Izbor zachovuje opis konstrukcije vměsto jeho zaměny pozdnějšoju kratkoju formuloju za diskriminant.\par
\noindent\emph{Reason:} Iz vsakih $n$ baznyh elementov i jih konjugovanyh elementov preserves the exact construction of Hilbert's determinant. Every selection of $n$ must not become an unspecified collection of elements. The description of the construction is retained rather than replaced by a later compact discriminant formula.\par
\noindent{\footnotesize Izvorne zapisy / Original records: HIST-g162-007.}\par
\end{minipage}\par
\par\medskip\Needspace{7\baselineskip}
\noindent\begin{minipage}{\linewidth}
\hypertarget{TERM-0260}{}
\noindent\textbf{kvadrat determinanta soprjaženyh bazovyh elementov (§6, 4)}\quad{\footnotesize\texttt{TERM-0260}}\par
\noindent\emph{Iztočnik / Source:} Determinantenquadrat der konjugierten Basiselemente (§6, 4)\par
\noindent\emph{Rusky / Russian:} квадрат определителя сопряженных базисных элементов (§6, 4)\par
\noindent\emph{Ukrajinsky / Ukrainian:} квадрат визначника спряжених базисних елементів (§6, 4)\par
\smallskip\noindent\emph{Razlog:} Kvadrat determinanta zachovuje operaciju kvadriranja cěloj determinanty, a referencija k oddělu 6, punktu 4 veže ju s ranějšim opreděljenjem. Za konjugovane elementy toj zapis izbira soprjaženyh, s obrazloženjem prepoznavajemosti vo vostočnoj i južnoj slovjanskoj terminologiji. To jest drugy dokumentovany kandidat než konjugovane; izvorni razlog ne dava čislenu ocěnu nih razumělivosti.\par
\noindent\emph{Reason:} Kvadrat determinanta retains squaring the whole determinant, and the reference to section 6, item 4 connects it with the earlier definition. This record chooses soprjaženyh for conjugate, invoking recognizability in East and South Slavic terminology. It is a different documented candidate from konjugovane; the rationale supplies no numerical intelligibility comparison.\par
\noindent{\footnotesize Izvorne zapisy / Original records: HIST-g157-008.}\par
\end{minipage}\par
\par\medskip\Needspace{7\baselineskip}
\noindent\begin{minipage}{\linewidth}
\hypertarget{TERM-0261}{}
\noindent\textbf{determinantny dělitelj}\quad{\footnotesize\texttt{TERM-0261}}\par
\noindent\emph{Iztočnik / Source:} Determinantenteiler\par
\noindent\emph{Rusky / Russian:} детерминантный делитель\par
\noindent\emph{Ukrajinsky / Ukrainian:} детермінантний дільник\par
\smallskip\noindent\emph{Razlog:} Determinantny dělitelj imenuje objekt teorije elementarnyh děliteljev, opreděljany kako največši obči dělitelj odgovarjajučih minorov. Oba zachovane razloga davajut prednost imenu objekta vměsto ponovjenogo prěpisovanja definicije pri každom pojavjenju. Definicija ostavaje potrěbna za razuměnje termina, ale ne služi kako njegova obvezna dolga zaměna.\par
\noindent\emph{Reason:} Determinantny dělitelj names the elementary-divisor-theory object defined as the greatest common divisor of the relevant minors. Both preserved rationales favour the object's name instead of repeating the definition at every occurrence. The definition remains necessary for understanding the term, but is not its mandatory long-form replacement.\par
\noindent{\footnotesize Izvorne zapisy / Original records: HIST-g085-004, HIST-g092-004, HIST-g093-008.}\par
\end{minipage}\par
\par\medskip\Needspace{7\baselineskip}
\noindent\begin{minipage}{\linewidth}
\hypertarget{TERM-0262}{}
\noindent\textbf{diagonalny člen}\quad{\footnotesize\texttt{TERM-0262}}\par
\noindent\emph{Iztočnik / Source:} Diagonalglied\par
\noindent\emph{Rusky / Russian:} диагональный член\par
\noindent\emph{Ukrajinsky / Ukrainian:} діагональний член\par
\smallskip\noindent\emph{Razlog:} Diagonalny člen svezaje medžunarodny pridavnik so slovjanskim imenovanjem člena. Dva zapisy opirajut se na diagonalno raspoloženje v izpisanoj tablici. Zato izraz ostavaje kratky: sama formula pokazuje odnos, ktory pridavnik imenuje.\par
\noindent\emph{Reason:} Diagonalny člen combines an international adjective with a Slavic noun for a member or term. Both records refer to the diagonal arrangement in the displayed array. The phrase can remain short because the formula shows the relation named by the adjective.\par
\noindent{\footnotesize Izvorne zapisy / Original records: HIST-g003-016, HIST-g012-012.}\par
\end{minipage}\par
\par\medskip\Needspace{7\baselineskip}
\noindent\begin{minipage}{\linewidth}
\hypertarget{TERM-0263}{}
\noindent\textbf{diagonalna matrica / transformovane čerez diagonalne matrice}\quad{\footnotesize\texttt{TERM-0263}}\par
\noindent\emph{Iztočnik / Source:} Diagonalmatrix / durch Diagonalmatrizen transformiert\par
\noindent\emph{Rusky / Russian:} диагональная матрица / преобразованы диагональными матрицами\par
\noindent\emph{Ukrajinsky / Ukrainian:} діагональна матриця / перетворені діагональними матрицями\par
\smallskip\noindent\emph{Razlog:} Diagonalna matrica i transformovanje s pomočju diagonalnyh matric sut dva različne uslovja. Odděl 10 razděljuje izbory parametrov, ktore ne tvorjut diagonalnu matricu i ne mogut byti prěvedene jeden v drugy diagonalnymi matricami. Zachovanje obojih konstrukcij čini vidimym kriterij za istinno različne razloženja.\par
\noindent\emph{Reason:} Being a diagonal matrix and being transformable by diagonal matrices are distinct conditions. Section 10 distinguishes parameter choices that do not form a diagonal matrix and cannot be transformed into one another by diagonal matrices. Keeping both constructions explicit preserves the criterion for genuinely different decompositions.\par
\noindent{\footnotesize Izvorne zapisy / Original records: HIST-g079-058.}\par
\end{minipage}\par
\par\medskip\Needspace{7\baselineskip}
\noindent\begin{minipage}{\linewidth}
\hypertarget{TERM-0264}{}
\noindent\textbf{diagonalne matrice / nulova matrica}\quad{\footnotesize\texttt{TERM-0264}}\par
\noindent\emph{Iztočnik / Source:} Diagonalmatrizen / Nullmatrix\par
\noindent\emph{Rusky / Russian:} диагональные матрицы / нулевая матрица\par
\noindent\emph{Ukrajinsky / Ukrainian:} діагональні матриці / нульова матриця\par
\smallskip\noindent\emph{Razlog:} Diagonalne matrice i nulova matrica zachovujut dva točne matrične nazvanja. Medžunarodna familija matric- i diagonal- odgovarja paralelnym ruskim i ukrajinskim izrazam. Nulova takože ostavaje izrična kvalifikacija; skupny zapis ne izravnjava dva pojmy.\par
\noindent\emph{Reason:} Diagonalne matrice and nulova matrica retain two exact matrix names. The international matrix/diagonal family corresponds to the Russian and Ukrainian parallels. Nulova remains an explicit qualification; grouping the expressions in one entry does not equate the concepts.\par
\noindent{\footnotesize Izvorne zapisy / Original records: HIST-g131-007.}\par
\end{minipage}\par
\par\medskip\Needspace{7\baselineskip}
\noindent\begin{minipage}{\linewidth}
\hypertarget{TERM-0265}{}
\noindent\textbf{diagonalna shema}\quad{\footnotesize\texttt{TERM-0265}}\par
\noindent\emph{Iztočnik / Source:} Diagonalschema\par
\noindent\emph{Rusky / Russian:} [Bukvy ne sut zachovane / Lettering not recoverable]\par
\noindent\emph{Ukrajinsky / Ukrainian:} [Bukvy ne sut zachovane / Lettering not recoverable]\par
\smallskip\noindent\emph{Razlog:} Diagonalna shema opisuje raspoloženje nenulovyh produktov po diagonali. Medžunarodna osnova diagonal- daje kratko nazvanje matrično podobnomu obrazcu. Ruska i ukrajinska grafija v zachovanom zapisu ne sut čitateljne: ostali sut samo znaky ?. Razlog izbora jest zachovany, ale točne sravniteljne formy ne mogut byti citovane iz tih poljev.\par
\noindent\emph{Reason:} Diagonalna shema describes a pattern in which the nonzero products lie on the diagonal. The international stem gives the matrix-like pattern a compact name. The preserved Russian and Ukrainian spellings are unreadable: only question marks remain. The rationale survives, but those fields cannot supply exact comparison forms.\par
\noindent{\footnotesize Izvorne zapisy / Original records: HIST-g079-038.}\par
\end{minipage}\par
\par\medskip\Needspace{7\baselineskip}
\noindent\begin{minipage}{\linewidth}
\hypertarget{TERM-0266}{}
\noindent\textbf{obglavjena relacija}\quad{\footnotesize\texttt{TERM-0266}}\par
\noindent\emph{Iztočnik / Source:} die angekündigte Relation\par
\noindent\emph{Rusky / Russian:} объявленное соотношение\par
\noindent\emph{Ukrajinsky / Ukrainian:} оголошене співвідношення\par
\smallskip\noindent\emph{Razlog:} Izvorni zapis razuměje konstrukciju kako nazvanje raněje objavjenoj relacije i označaje obglavjena kako možno previše doslovno. Njegova ocěna razumělivosti ne jest podkrěpljena citatom jezyčnogo iztočnika. Proto historijni razlog trěba zachovati odvojeno od pozdnějšego slovničnogo sravnjenja: on dokumentuje naměr, ne dokazuje, že kandidat vyražaje taj naměr pravilno.\par
\noindent\emph{Reason:} The original record intends a previously announced relation and marks obglavjena as possibly too literal. Its intelligibility judgment is not supported there by a cited linguistic source. The historical reason must therefore remain distinct from the later dictionary comparison: it records the intention, not proof that the candidate expresses that intention correctly.\par
\noindent{\footnotesize Izvorne zapisy / Original records: HIST-g018-007.}\par
\smallskip\noindent\emph{Sravnjenje s tekstom:} V dělu 2, oddělu 17, prěd formuloju (22), tekst ima Obglavjena relacija jest. Konsultovana kopija oficialnogo slovnika ne dava medžuslovjansko zaglavje obglaviti. Podobny napis obglaviti nahodi se v slovenskom, ne medžuslovjanskom, stolbcu pri obezglåviti ‘odrězati glavu’ (773); to ne jest dokazom za medžuslovjansky smysl ‘objaviti’. Za potrěbny smysl slovnik dava oglåsiti (173) i obvěstiti (22225). Prvi kandidat dava osnovno oglasiti, a pravilo za prošle trpne pričestje sloves na -siti dava oglašeny; ženska forma oglašena soglasuje se s relacija. Proto oglašena relacija jest leksično i morfologično podprěty prědlog za tuto město. On ne jest potvrđeny slovnikom kako cěla fraza, ne jest priměnjen v citovanom tekstu i ne izključaje drugu pravilnu formulaciju. Razlika medžu izvornoju ocěnoju razumělivosti i tutym pozdnějšim dokazom mora ostati vidima.\par
\noindent\emph{Textual comparison:} In Paper 2, section 17, immediately before equation (22), the text reads Obglavjena relacija jest. The consulted official dictionary snapshot supplies no Interslavic headword obglaviti. A similar spelling occurs in the Slovenian comparison column, not the Interslavic column, under obezglåviti ‘decapitate’ (773); it does not support an Interslavic sense ‘announce’. For the required sense the dictionary supplies oglåsiti (173) and obvěstiti (22225). The former gives basic-spelling oglasiti; the rule for the past passive participle of verbs in -siti gives oglašeny, with feminine oglašena agreeing with relacija. Oglašena relacija is therefore a lexically and morphologically supported proposal for this passage. It is not a dictionary-attested phrase, is not applied in the quoted text, and does not exclude other well-formed formulations. The original intelligibility judgment must remain distinguishable from this later evidence.\par
\noindent{\footnotesize TEXT-COMP-04; pasaže / passages: 2453; \texttt{CORPUS\_FORM\_COMPARISONS.json}.}\par
\noindent{\footnotesize\href{https://interslavic-dictionary.com/}{Medžuslovjansky slovnik / Interslavic Dictionary}: 173, 773, 1481, 5078, 22224, 22225; kopija / snapshot 2026-08-16.}\par
\noindent{\footnotesize\href{https://interslavic.fun/learn/grammar/verbs/}{Slovesa / Verbs}: prošle trpne pričestje / past passive participle.}\par
\noindent{\footnotesize\href{https://interslavic.fun/learn/orthography/}{Pravopis / Orthography}: opcionalne bukvy / optional letters.}\par
\end{minipage}\par
\par\medskip\Needspace{7\baselineskip}
\noindent\begin{minipage}{\linewidth}
\hypertarget{TERM-0267}{}
\noindent\textbf{aritmetična teorija algebraičnyh funkcij jednoj proměnnoj}\quad{\footnotesize\texttt{TERM-0267}}\par
\noindent\emph{Iztočnik / Source:} Die arithmetische Theorie der algebraischen Funktionen einer Veränderlichen\par
\noindent\emph{Rusky / Russian:} арифметическая теория алгебраических функций одной переменной\par
\noindent\emph{Ukrajinsky / Ukrainian:} арифметична теорія алгебраїчних функцій однієї змінної\par
\smallskip\noindent\emph{Razlog:} Naslov zachovuje aritmetičnu teoriju, algebraične funkcije i omejenje do jednoj proměnnoj. Poredok slov jest prilagodženy običnoj matematičnoj konstrukciji, ne německoj sintaksi. Medžunarodne aritmetična i algebraična so slovjanskim proměnna davajut čitatelju tri glavne informacije o opsegu děla.\par
\noindent\emph{Reason:} The title retains arithmetic theory, algebraic functions, and the restriction to one variable. Its word order follows an ordinary mathematical construction rather than German syntax. International aritmetična and algebraična together with Slavic proměnna carry the three principal scope qualifications.\par
\noindent{\footnotesize Izvorne zapisy / Original records: HIST-g076-001.}\par
\end{minipage}\par
\par\medskip\Needspace{7\baselineskip}
\noindent\begin{minipage}{\linewidth}
\hypertarget{TERM-0268}{}
\noindent\textbf{definicija se pokryva s definicijeju v §7, 1}\quad{\footnotesize\texttt{TERM-0268}}\par
\noindent\emph{Iztočnik / Source:} die Definition deckt sich mit der in §7, 1\par
\noindent\emph{Rusky / Russian:} определение совпадает с определением в §7, 1\par
\noindent\emph{Ukrajinsky / Ukrainian:} означення збігається з означенням у §7, 1\par
\smallskip\noindent\emph{Razlog:} Definicija se pokryva s definicijeju v oddělu 7, punktu 1 jest tvrđenje o sovpadanju definicij, ne samo bibliografična ukazka. Referencija mora ostati svezana s tym tvrđenjem. Samo sravni odděl 7, 1 by oslabilo odnos, ktory iztočnik izrično ustanavjaje.\par
\noindent\emph{Reason:} Definicija se pokryva s definicijeju v §7, 1 asserts agreement between definitions, not merely a bibliographical pointer. The reference must remain attached to that assertion. Simply writing see section 7, item 1 would weaken the relation explicitly stated by the source.\par
\noindent{\footnotesize Izvorne zapisy / Original records: HIST-g158-010.}\par
\end{minipage}\par
\par\medskip\Needspace{7\baselineskip}
\noindent\begin{minipage}{\linewidth}
\hypertarget{TERM-0269}{}
\noindent\textbf{dvojno izračunanje ... daje povod k redukciji}\quad{\footnotesize\texttt{TERM-0269}}\par
\noindent\emph{Iztočnik / Source:} Die doppelte Berechnung ... gibt Anlaß zur Reduktion\par
\noindent\emph{Rusky / Russian:} двойное вычисление ... дает основание для редукции\par
\noindent\emph{Ukrajinsky / Ukrainian:} подвійне обчислення ... дає підставу для редукції\par
\smallskip\noindent\emph{Razlog:} Dvojno izračunanje ... daje povod k redukciji svazuje dva vyčisljenja s redukcijnym argumentom. Ruska i ukrajinska paralela govore o osnovanji, ne kopirujut doslovno německy obrat s Anlaß. Povod byl izbrany radi razumělivosti, ale sam zapis ostavjaje njegovu vhodnost za sravnjenje s drugimi povtorjenymi konstrukcijami dokazov.\par
\noindent\emph{Reason:} Dvojno izračunanje ... daje povod k redukciji connects the two calculations with the reduction argument. Russian and Ukrainian speak of grounds rather than copying the German Anlaß construction literally. Povod was selected for intelligibility, but the record itself leaves its suitability for comparison with recurring proof-language constructions.\par
\noindent{\footnotesize Izvorne zapisy / Original records: HIST-g011-009.}\par
\end{minipage}\par
\par\medskip\Needspace{7\baselineskip}
\noindent\begin{minipage}{\linewidth}
\hypertarget{TERM-0270}{}
\noindent\textbf{obična v literaturě definicija polnoj razložimosti}\quad{\footnotesize\texttt{TERM-0270}}\par
\noindent\emph{Iztočnik / Source:} die in der Literatur übliche Definition der vollständigen Reduzibilität\par
\noindent\emph{Rusky / Russian:} обычное в литературе определение полной приводимости\par
\noindent\emph{Ukrajinsky / Ukrainian:} звичайне в літературі означення цілковитої звідності\par
\smallskip\noindent\emph{Razlog:} Polna razložimost prodolžaje terminologiju iz odděla 3 i početka odděla 6. Podnožna primětka svazuje Noetherinu formulaciju črěz matrice i klasy s definicijeju, običnoju v literaturě. To ne jest nova tvrdnja prěvoditelja o vsej literaturě: prěvod zachovuje odnos, ktory sama primětka ustanavjaje.\par
\noindent\emph{Reason:} Polna razložimost continues the terminology of section 3 and the opening of section 6. The footnote connects Noether's matrix-and-class formulation with the definition usual in the literature. This is not a new translator's claim about the whole literature; the translation preserves the relationship asserted by the footnote.\par
\noindent{\footnotesize Izvorne zapisy / Original records: HIST-g141-006.}\par
\end{minipage}\par
\par\medskip\Needspace{7\baselineskip}
\noindent\begin{minipage}{\linewidth}
\hypertarget{TERM-0271}{}
\noindent\textbf{proces diferenciacije, ktory točno izčerpava rędy}\quad{\footnotesize\texttt{TERM-0271}}\par
\noindent\emph{Iztočnik / Source:} die Reihen genau erschöpfender Differentiationsprozeß\par
\noindent\emph{Rusky / Russian:} процесс дифференцирования, точно исчерпывающий строки\par
\noindent\emph{Ukrajinsky / Ukrainian:} процес диференціювання, що точно вичерпує рядки\par
\smallskip\noindent\emph{Razlog:} Točno izčerpava zachovuje uslovje, že operacija diferenciacije ispolzuje imenovane redy točno, ne samo njih čest ili neoprěděljeno mnoho redov. Zapis veže proces s redami označenymi črěz xi v konstrukciji. Etimologični napis rędy jest zachovan kako forma kandidata, ne kako dokaz za nužnost ę v čitatelskom tekstu.\par
\noindent\emph{Reason:} Točno izčerpava preserves the condition that differentiation consumes exactly the specified rows, not merely some of them or an unspecified number. The record relates the process to the xi rows in the construction. Etymological rędy is retained as the candidate's form, not as a requirement to use ę throughout the reader.\par
\noindent{\footnotesize Izvorne zapisy / Original records: HIST-g036-010.}\par
\end{minipage}\par
\par\medskip\Needspace{7\baselineskip}
\noindent\begin{minipage}{\linewidth}
\hypertarget{TERM-0272}{}
\noindent\textbf{ostale formy}\quad{\footnotesize\texttt{TERM-0272}}\par
\noindent\emph{Iztočnik / Source:} die übrigen Formen\par
\noindent\emph{Rusky / Russian:} остальные формы\par
\noindent\emph{Ukrajinsky / Ukrainian:} решта форм\par
\smallskip\noindent\emph{Razlog:} Ostale formy jest kratko imenovanje form, ktore ostavajut poslě už razmotrenyh slučajev. Izvorni razlog daje prednost običnoj proze dokaza i prirodnomu poredku slov. Ne vnašaje se nova matematična kategorija tamo, kde avtor samo razděljuje razmatrane slučaje.\par
\noindent\emph{Reason:} Ostale formy briefly names the forms remaining after the cases already considered. The recorded reason favours ordinary proof prose and natural word order. It introduces no new mathematical category where the author is simply partitioning the cases under discussion.\par
\noindent{\footnotesize Izvorne zapisy / Original records: HIST-g016-006.}\par
\end{minipage}\par
\par\medskip\Needspace{7\baselineskip}
\noindent\begin{minipage}{\linewidth}
\hypertarget{TERM-0273}{}
\noindent\textbf{diedralne grupy}\quad{\footnotesize\texttt{TERM-0273}}\par
\noindent\emph{Iztočnik / Source:} Diedergruppen\par
\noindent\emph{Rusky / Russian:} диэдральные группы\par
\noindent\emph{Ukrajinsky / Ukrainian:} діедральні групи\par
\smallskip\noindent\emph{Razlog:} Diedralne grupy zachovuje medžunarodno ime togo tipa grup s pridavnikovym zakončenjem, soglasovanym s grupy. Ruska i ukrajinska paralela imajut istu terminologičnu osnovu. Kratky izvorni zapis ne sodržuje zasebno sravnjenje slovotvornyh alternativ.\par
\noindent\emph{Reason:} Diedralne grupy retains the international name for this type of group with an adjectival ending agreeing with grupy. The Russian and Ukrainian parallels share the terminological stem. The brief original record contains no separate comparison of alternative formations.\par
\noindent{\footnotesize Izvorne zapisy / Original records: HIST-g039-015.}\par
\end{minipage}\par
\par\medskip\Needspace{7\baselineskip}
\noindent\begin{minipage}{\linewidth}
\hypertarget{TERM-0274}{}
\noindent\textbf{te same razmysljenja važe za multiplikacijsko kolco}\quad{\footnotesize\texttt{TERM-0274}}\par
\noindent\emph{Iztočnik / Source:} Dieselben Überlegungen gelten für den Multiplikationsring\par
\noindent\emph{Rusky / Russian:} те же рассуждения применимы к мультипликационному кольцу\par
\noindent\emph{Ukrajinsky / Ukrainian:} ті самі міркування застосовні до мультиплікаційного кільця\par
\smallskip\noindent\emph{Razlog:} Te same razmysljenja važe ... izrično prěnose ranějši argument na multiplikacijsko kolco. Izbor zachovuje Krullov imenovany pojem iz početka odděla 8 i ne skračaje cěly prěhod do podobno. Čitatelj može viděti, ktory argument i na ktory objekt se priměnjaje.\par
\noindent\emph{Reason:} Te same razmysljenja važe ... explicitly transfers the earlier reasoning to the multiplication ring. It preserves Krull's named concept from the start of section 8 rather than reducing the whole transition to similarly. Readers can see which argument is being applied to which object.\par
\noindent{\footnotesize Izvorne zapisy / Original records: HIST-g160-001.}\par
\end{minipage}\par
\par\medskip\Needspace{7\baselineskip}
\noindent\begin{minipage}{\linewidth}
\hypertarget{TERM-0275}{}
\noindent\textbf{diferenta}\quad{\footnotesize\texttt{TERM-0275}}\par
\noindent\emph{Iztočnik / Source:} Differente\par
\noindent\emph{Rusky / Russian:} дифферента\par
\noindent\emph{Ukrajinsky / Ukrainian:} диферент; диферента\par
\smallskip\noindent\emph{Razlog:} Diferenta jest osobny objekt algebraičnoj teorije čisel, ne obična razlika i ne sinonim diskriminanta. Tri zapisy zachovujut medžunarodno nazvanje i njegovu vezu s rabotami 36 i 37. Ukrajinske formy v sravnjenju imajut i diferent, i diferenta; medžuslovjansky žensky imennik jest obrazloženy kandidat, čiju morfologiju ne dokazuje samo prisutnost medžunarodnogo korenja.\par
\noindent\emph{Reason:} Diferenta names a particular algebraic-number-theory object, not an ordinary difference or a synonym for discriminant. Three records retain the international name and its continuity with Papers 36 and 37. The Ukrainian comparanda include both diferent and diferenta; the Interslavic feminine noun is a reasoned candidate whose morphology is not established merely by the shared international root.\par
\noindent{\footnotesize Izvorne zapisy / Original records: HIST-g206-002, HIST-g212-012, HIST-g213-002.}\par
\end{minipage}\par
\par\medskip\Needspace{7\baselineskip}
\noindent\begin{minipage}{\linewidth}
\hypertarget{TERM-0276}{}
\noindent\textbf{diferencialny izraz}\quad{\footnotesize\texttt{TERM-0276}}\par
\noindent\emph{Iztočnik / Source:} Differentialausdruck\par
\noindent\emph{Rusky / Russian:} дифференциальное выражение\par
\noindent\emph{Ukrajinsky / Ukrainian:} диференціальний вираз\par
\smallskip\noindent\emph{Razlog:} Diferencialny izraz zachovuje široky objekt: funkciju koordinatov i diferencialov ili derivacij. Izvorni razlog izrično izběgaje suženje do diferencialnoj formy v užšem sovrěmennom smyslu. Izraz zato jest sadržajno važny, ne samo stilistična zaměna slova forma.\par
\noindent\emph{Reason:} Diferencialny izraz preserves a broad object: a function of coordinates and differentials or derivatives. The recorded reason explicitly avoids narrowing it to a differential form in the restricted modern sense. Expression is therefore substantive, not merely a stylistic substitute for form.\par
\noindent{\footnotesize Izvorne zapisy / Original records: HIST-g074-002.}\par
\end{minipage}\par
\par\medskip\Needspace{7\baselineskip}
\noindent\begin{minipage}{\linewidth}
\hypertarget{TERM-0277}{}
\noindent\textbf{diferencialne izrazy}\quad{\footnotesize\texttt{TERM-0277}}\par
\noindent\emph{Iztočnik / Source:} Differentialausdruecke\par
\noindent\emph{Rusky / Russian:} дифференциальные выражения\par
\noindent\emph{Ukrajinsky / Ukrainian:} диференціальні вирази\par
\smallskip\noindent\emph{Razlog:} Diferencialne izrazy upotrěbjaje medžunarodnu osnovu diferencial- i običny imennik izraz, paralelno s ruskim i ukrajinskim imenovanjami. To jest množinna formulacija o vyraženjah. Izvorni napis Differentialausdruecke s ue jest zachovan kako osobna dokumentovana forma německogo zaglavja, ne kako drugy matematičny pojem.\par
\noindent\emph{Reason:} Diferencialne izrazy combines international diferencial- with the ordinary noun izraz, parallel to Russian and Ukrainian. It is the plural expression wording. Historical Differentialausdruecke with ue remains a separately recorded German headword form, not a different mathematical concept.\par
\noindent{\footnotesize Izvorne zapisy / Original records: HIST-g079-004.}\par
\end{minipage}\par
\par\medskip\Needspace{7\baselineskip}
\noindent\begin{minipage}{\linewidth}
\hypertarget{TERM-0278}{}
\noindent\textbf{diferencialne izrazy}\quad{\footnotesize\texttt{TERM-0278}}\par
\noindent\emph{Iztočnik / Source:} Differentialausdrücke\par
\noindent\emph{Rusky / Russian:} дифференциальные выражения\par
\noindent\emph{Ukrajinsky / Ukrainian:} диференціальні вирази\par
\smallskip\noindent\emph{Razlog:} Diferencialne izrazy zachovuje širši opseg pregleda, ktory ne omejaje se diferencialnymi jednačenjami. Zaměna izrazy slovom jednačenja by nepravilno suzila imenovane objekty. Ista medžuslovjanska forma jest dokumentovana i pod německim napisom s ue; pravopisna razlika iztočnika ne stvara novu terminologičnu odluku.\par
\noindent\emph{Reason:} Diferencialne izrazy preserves the survey's scope beyond differential equations. Replacing expressions with equations would wrongly narrow the named objects. The same Interslavic form is recorded under the German spelling with ue; the source-spelling difference does not create a new terminological decision.\par
\noindent{\footnotesize Izvorne zapisy / Original records: HIST-g094-023.}\par
\end{minipage}\par
\par\medskip\Needspace{7\baselineskip}
\noindent\begin{minipage}{\linewidth}
\hypertarget{TERM-0279}{}
\noindent\textbf{diferencialne invarianty}\quad{\footnotesize\texttt{TERM-0279}}\par
\noindent\emph{Iztočnik / Source:} Differentialinvarianten\par
\noindent\emph{Rusky / Russian:} дифференциальные инварианты\par
\noindent\emph{Ukrajinsky / Ukrainian:} диференціальні інваріанти\par
\smallskip\noindent\emph{Razlog:} Diferencialne invarianty zachovuje medžunarodnu familiju invariant-, ktora čitatelju povezuje termin s teorijeju invariantov. Izvorni razlog ne daje prednost novoj domorodnoj parafraze, ibo ona mogla by učiniti tu vezu menje poznavajemoju. To jest razlog terminologičnoj prepoznavajemosti, ne izměreno testovanje publike.\par
\noindent\emph{Reason:} Diferencialne invarianty retains the international invariant family, making the connection with invariant theory recognizable. The recorded reason rejects introducing a native paraphrase that could obscure that connection. This is an argument about terminological recognizability, not a reported audience test.\par
\noindent{\footnotesize Izvorne zapisy / Original records: HIST-g084-002.}\par
\end{minipage}\par
\par\medskip\Needspace{7\baselineskip}
\noindent\begin{minipage}{\linewidth}
\hypertarget{TERM-0280}{}
\noindent\textbf{diferencialny kvocient}\quad{\footnotesize\texttt{TERM-0280}}\par
\noindent\emph{Iztočnik / Source:} Differentialquotient\par
\noindent\emph{Rusky / Russian:} дифференциальное частное\par
\noindent\emph{Ukrajinsky / Ukrainian:} диференціальний частник\par
\smallskip\noindent\emph{Razlog:} Diferencialny kvocient zachovuje německo imenovanje v argumentu, kde diferenta jest svezana s definujučim idealom. Izvorni zapis naměrno ne zaměnjaje termin avtomatično sovrěmennym derivacija. Opravdanje leži v lokalnom objasnjenju kvocienta; ono ne jest obče pravilo prěvoditi vsako Differentialquotient istim sposobom.\par
\noindent\emph{Reason:} Diferencialny kvocient retains the German naming within the argument relating the different to a defining ideal. The original record deliberately avoids automatically modernizing the term as derivative. The justification lies in the local explanation of the quotient; it is not a universal rule for every occurrence of Differentialquotient.\par
\noindent{\footnotesize Izvorne zapisy / Original records: HIST-g206-003.}\par
\end{minipage}\par
\par\medskip\Needspace{7\baselineskip}
\noindent\begin{minipage}{\linewidth}
\hypertarget{TERM-0281}{}
\noindent\textbf{diferencialny kvocient / razlikovy kvocient}\quad{\footnotesize\texttt{TERM-0281}}\par
\noindent\emph{Iztočnik / Source:} Differentialquotient / Differenzenquotient\par
\noindent\emph{Rusky / Russian:} дифференциальное частное / разностное частное\par
\noindent\emph{Ukrajinsky / Ukrainian:} диференціальний частник / різницевий частник\par
\smallskip\noindent\emph{Razlog:} Diferencialny kvocient i razlikovy kvocient morajut ostati različne nazvanja, ibo argument sravnjava derivaciji podobny objekt s kvocientom razlik. Kvocient jest izbran vměsto častka, da by poslednje slovo ne nosilo dodatny smysl. Ruska i ukrajinska paralela imaju svoja imena kvocienta; jejich funkcija podpiraje različenje objektov, ne tvrdnju, že vsi jezyky imajut isti korenj.\par
\noindent\emph{Reason:} Diferencialny kvocient and razlikovy kvocient must remain distinct because the argument compares a derivative-like object with a difference quotient. Kvocient was chosen over častka to avoid giving the latter another sense. The Russian and Ukrainian comparanda use their own quotient nouns; they support the distinction between the objects, not a claim of a shared lexical root.\par
\noindent{\footnotesize Izvorne zapisy / Original records: HIST-g213-003.}\par
\end{minipage}\par
\par\medskip\Needspace{7\baselineskip}
\noindent\begin{minipage}{\linewidth}
\hypertarget{TERM-0282}{}
\noindent\textbf{častne derivaty 2-go reda}\quad{\footnotesize\texttt{TERM-0282}}\par
\noindent\emph{Iztočnik / Source:} Differentialquotienten 2. Ordnung\par
\noindent\emph{Rusky / Russian:} частные производные 2-го порядка\par
\noindent\emph{Ukrajinsky / Ukrainian:} частинні похідні 2-го порядку\par
\smallskip\noindent\emph{Razlog:} Častne derivaty 2-go reda jest naměrna modernizacija staršego Differentialquotienten v kontekstu častnyh derivatov. Razlog sostavja slovjansky pridavnik s medžunarodnym imennikom i zachovuje vtory red. Sam historijni zapis ostavja termin za prověrku po avtoritetnom iztočniku. Zato tuto ne jest dokaz, že modernizacija podhodit vsakoj upotrěbě istogo německogo slova.\par
\noindent\emph{Reason:} Častne derivaty 2-go reda deliberately modernizes the older Differentialquotienten in a context of partial derivatives. The reason combines a Slavic adjective with an international noun and retains the second order. The historical record itself leaves the term for comparison with an authoritative source. It therefore does not establish that the modernization suits every occurrence of the German word.\par
\noindent{\footnotesize Izvorne zapisy / Original records: HIST-g035-008.}\par
\end{minipage}\par
\par\medskip\Needspace{7\baselineskip}
\noindent\begin{minipage}{\linewidth}
\hypertarget{TERM-0283}{}
\noindent\textbf{procesy diferenciranja i variiranja}\quad{\footnotesize\texttt{TERM-0283}}\par
\noindent\emph{Iztočnik / Source:} Differentiations- und Variationsprozesse\par
\noindent\emph{Rusky / Russian:} процессы дифференцирования и варьирования\par
\noindent\emph{Ukrajinsky / Ukrainian:} процеси диференціювання та варіювання\par
\smallskip\noindent\emph{Razlog:} Procesy diferenciranja i variiranja zachovujut dvě operacije, ktorymi dokaz tvori invarianty. Slovo procesy ne jest izlišnja dobavka: ono čini vidimym konstrukcijny postupok, ne samo tvrdženje o suščestvovanju rezultata. Izbor sovrěmennogo računskogo jezyka ne daje razloga izpustiti tu proceduralnu čest.\par
\noindent\emph{Reason:} Procesy diferenciranja i variiranja preserves the two operations by which the proof constructs invariants. Processes is not a dispensable addition: it identifies a construction, rather than merely asserting existence. Choosing modern calculus vocabulary does not justify dropping that procedural content.\par
\noindent{\footnotesize Izvorne zapisy / Original records: HIST-g074-008.}\par
\end{minipage}\par
\par\medskip\Needspace{7\baselineskip}
\noindent\begin{minipage}{\linewidth}
\hypertarget{TERM-0284}{}
\noindent\textbf{procesy diferenciacije}\quad{\footnotesize\texttt{TERM-0284}}\par
\noindent\emph{Iztočnik / Source:} Differentiationsprozesse\par
\noindent\emph{Rusky / Russian:} процессы дифференцирования\par
\noindent\emph{Ukrajinsky / Ukrainian:} процеси диференціювання\par
\smallskip\noindent\emph{Razlog:} Procesy diferenciacije zachovujut paralelu s nazvanjem procesa svijanja. Diferenciacija jest sovrěmenno ime operacije, a procesy ukazyvajut na njezino priměnjenje v dokazě. Drugy zapis ima diferenciranje; izvorni razlog ovde ne sravnjava te dvě tvorjenja i ne ustanavjaje jednu opču normu.\par
\noindent\emph{Reason:} Procesy diferenciacije keeps the parallel with the naming of the folding process. Diferenciacija names the operation in modern vocabulary, while procesy identifies its use in the argument. Another record uses diferenciranje; the present rationale neither compares the two formations nor establishes a general preference.\par
\noindent{\footnotesize Izvorne zapisy / Original records: HIST-g035-006.}\par
\end{minipage}\par
\par\medskip\Needspace{7\baselineskip}
\noindent\begin{minipage}{\linewidth}
\hypertarget{TERM-0285}{}
\noindent\textbf{diferencialne simboly}\quad{\footnotesize\texttt{TERM-0285}}\par
\noindent\emph{Iztočnik / Source:} Differentiationssymbole\par
\noindent\emph{Rusky / Russian:} дифференциальные символы\par
\noindent\emph{Ukrajinsky / Ukrainian:} диференціальні символи\par
\smallskip\noindent\emph{Razlog:} Diferencialne simboly imenujut simbolične redy, ktore v formulě (45) zaměnjajut $S^\sigma$ i $T_{n-\tau}$. Ne imenujut derivat sam po sobě. Izraz zato zachovuje razliku medžu formalnym zapisom operacije i objektom, na ktory ona dějstvuje.\par
\noindent\emph{Reason:} Diferencialne simboly names the symbolic rows replacing $S^\sigma$ and $T_{n-\tau}$ in formula (45), not a derivative object by itself. The wording preserves the distinction between the formal expression of an operation and the object on which it acts.\par
\noindent{\footnotesize Izvorne zapisy / Original records: HIST-g035-007.}\par
\end{minipage}\par
\par\medskip\Needspace{7\baselineskip}
\noindent\begin{minipage}{\linewidth}
\hypertarget{TERM-0286}{}
\noindent\textbf{razlikove izrazy}\quad{\footnotesize\texttt{TERM-0286}}\par
\noindent\emph{Iztočnik / Source:} Differenzenausdruecke\par
\noindent\emph{Rusky / Russian:} разностные выражения\par
\noindent\emph{Ukrajinsky / Ukrainian:} різницеві вирази\par
\smallskip\noindent\emph{Razlog:} Razlikove izrazy sut izrazy s operacijami konečnyh razlik, postavjene paralelno diferencialnym izrazam. Izraz ne jest avtomatično jednačenje. Zaměna slovom jednačenja suzila by obseg i skryla by operatornu funkciju, na ktoru ukazyvaje izvorni razlog.\par
\noindent\emph{Reason:} Razlikove izrazy denotes finite-difference operator expressions, paired with differential expressions. An expression is not automatically an equation. Replacing the noun with equations would narrow the scope and obscure the operator role identified by the recorded reason.\par
\noindent{\footnotesize Izvorne zapisy / Original records: HIST-g079-003.}\par
\end{minipage}\par
\par\medskip\Needspace{7\baselineskip}
\noindent\begin{minipage}{\linewidth}
\hypertarget{TERM-0287}{}
\noindent\textbf{stepen}\quad{\footnotesize\texttt{TERM-0287}}\par
\noindent\emph{Iztočnik / Source:} Dimension (of a row in the homogeneous-form context)\par
\noindent\emph{Rusky / Russian:} степень\par
\noindent\emph{Ukrajinsky / Ukrainian:} степінь\par
\smallskip\noindent\emph{Razlog:} Stepen prědava lokalny smysl německogo Dimension za odnorodnost po jednom redu proměnnyh. Razlog izrično razlučaje taj stepen od dimenzije vektornogo prostora. Leksično podoben prěvod dimenzija mogl by zato navesti čitatelja na drugy matematičny pojem; kontekst ima prednost nad podobnostju slov.\par
\noindent\emph{Reason:} Stepen renders the local sense of Dimension as the degree of homogeneity in one row of variables. The reason explicitly distinguishes this from vector-space dimension. The lexically similar dimenzija could therefore suggest a different mathematical concept; contextual meaning takes priority over resemblance.\par
\noindent{\footnotesize Izvorne zapisy / Original records: HIST-g057-008.}\par
\end{minipage}\par
\par\medskip\Needspace{7\baselineskip}
\noindent\begin{minipage}{\linewidth}
\hypertarget{TERM-0288}{}
\noindent\textbf{stepen}\quad{\footnotesize\texttt{TERM-0288}}\par
\noindent\emph{Iztočnik / Source:} Dimension / Grad in the homogeneous-form argument\par
\noindent\emph{Rusky / Russian:} степень\par
\noindent\emph{Ukrajinsky / Ukrainian:} степінь\par
\smallskip\noindent\emph{Razlog:} V argumentu o odnorodnoj formě Dimension i Grad imajut smysl stepena. Zapis nastavlja isto kontekstualno različenje iz raboty 8, ne pravilo zaměnjati vsako Dimension slovom stepen. Takojže ruska i ukrajinska paralela ovde govoret o stepeni, ne o razměrnosti prostora.\par
\noindent\emph{Reason:} In this homogeneous-form argument, Dimension and Grad have the sense of degree. The record continues the contextual distinction used in Paper 8, not a rule replacing every Dimension with stepen. Here the Russian and Ukrainian comparanda likewise use degree vocabulary rather than space dimension.\par
\noindent{\footnotesize Izvorne zapisy / Original records: HIST-g061-009.}\par
\end{minipage}\par
\par\medskip\Needspace{7\baselineskip}
\noindent\begin{minipage}{\linewidth}
\hypertarget{TERM-0289}{}
\noindent\textbf{dimenzija v sobstvenom smyslu}\quad{\footnotesize\texttt{TERM-0289}}\par
\noindent\emph{Iztočnik / Source:} Dimension schlechthin\par
\noindent\emph{Rusky / Russian:} размерность в собственном смысле\par
\noindent\emph{Ukrajinsky / Ukrainian:} розмірність у власному сенсі\par
\smallskip\noindent\emph{Razlog:} Dimenzija v sobstvenom smyslu razkriva funkciju schlechthin kako ograničenja načina upotrěby termina. Konstrukcija trěbuje kontrast s najvyšša dimenzija; prosto ponovjenje imennika dimenzija skrylo by tu razliku. Izvorni razlog zato vybira objasnjujuču frazu vměsto doslovnogo pridevka k imenu.\par
\noindent\emph{Reason:} Dimenzija v sobstvenom smyslu makes schlechthin's qualification of the term explicit. The expression needs its contrast with najvyšša dimenzija; simply repeating dimenzija would conceal that distinction. The recorded reason therefore favours an explanatory phrase over a literal adjective attached to the noun.\par
\noindent{\footnotesize Izvorne zapisy / Original records: HIST-g096-012.}\par
\end{minipage}\par
\par\medskip\Needspace{7\baselineskip}
\noindent\begin{minipage}{\linewidth}
\hypertarget{TERM-0290}{}
\noindent\textbf{dimenzijne čisla}\quad{\footnotesize\texttt{TERM-0290}}\par
\noindent\emph{Iztočnik / Source:} Dimensionszahlen\par
\noindent\emph{Rusky / Russian:} числа размерности\par
\noindent\emph{Ukrajinsky / Ukrainian:} числа розмірності\par
\smallskip\noindent\emph{Razlog:} Dimenzijne čisla zachovujut množinu čisel, ktore teorem V sravnjava za asociovane proste idealy. Samo dimenzija mogla by izgledati kako jedna veličina i skryti seznam invariantov. Imennik čisla jest zato čestju matematičnogo sodržanja, ne izpustima razširenost nazvanja.\par
\noindent\emph{Reason:} Dimenzijne čisla preserves the plural numerical data that Theorem V compares for associated primes. Dimension alone could suggest one value and obscure the list of invariants. Čisla is therefore part of the mathematical content, not an expendable elaboration of the name.\par
\noindent{\footnotesize Izvorne zapisy / Original records: HIST-g098-002.}\par
\end{minipage}\par
\par\medskip\Needspace{7\baselineskip}
\noindent\begin{minipage}{\linewidth}
\hypertarget{TERM-0291}{}
\noindent\textbf{direktno nerazložima}\quad{\footnotesize\texttt{TERM-0291}}\par
\noindent\emph{Iztočnik / Source:} direkt unzerlegbar\par
\noindent\emph{Rusky / Russian:} прямо неразложимая\par
\noindent\emph{Ukrajinsky / Ukrainian:} прямо нерозкладна\par
\smallskip\noindent\emph{Razlog:} Direktno ograničaje tip razloženja, ktory jest izključeny; nerazložimost ne jest tym samym nerazložimostju v smyslu neprivedimosti. Ženska forma nerazložima naleži k konkretnoj konstrukciji iz odděla 4. Jej gramatičny rod ne tvori novy pojem v sravnjenju s mužskym nerazložimy.\par
\noindent\emph{Reason:} Direktno specifies the kind of decomposition being excluded; this is not thereby irreducibility. Feminine nerazložima belongs to the particular construction in section 4. Its grammatical gender does not create a different concept from masculine nerazložimy.\par
\noindent{\footnotesize Izvorne zapisy / Original records: HIST-g182-003.}\par
\end{minipage}\par
\par\medskip\Needspace{7\baselineskip}
\noindent\begin{minipage}{\linewidth}
\hypertarget{TERM-0292}{}
\noindent\textbf{direktno nerazložimy}\quad{\footnotesize\texttt{TERM-0292}}\par
\noindent\emph{Iztočnik / Source:} direkt unzerlegbar\par
\noindent\emph{Rusky / Russian:} прямо неразложимый\par
\noindent\emph{Ukrajinsky / Ukrainian:} прямо нерозкладний\par
\smallskip\noindent\emph{Razlog:} Direktno nerazložimy označaje nemogućnost razloženja v direktnu sumu, ne neprivedimost. Zapis svezava termin s ranějšim jezykom teorije idealov. Zachovanje priloga direktno štiti upravo to matematično ograničenje; kratak prěvod nerazložimy bez konteksta mogl by jeho izgubiti.\par
\noindent\emph{Reason:} Direktno nerazložimy excludes a direct-sum decomposition, not reducibility in general. The record connects the term to earlier ideal-theoretic usage. Keeping direktno protects that mathematical restriction, which a context-free nerazložimy could lose.\par
\noindent{\footnotesize Izvorne zapisy / Original records: HIST-g194-005.}\par
\end{minipage}\par
\par\medskip\Needspace{7\baselineskip}
\noindent\begin{minipage}{\linewidth}
\hypertarget{TERM-0293}{}
\noindent\textbf{direktno nerazložime faktory}\quad{\footnotesize\texttt{TERM-0293}}\par
\noindent\emph{Iztočnik / Source:} direkt unzerlegbare Faktoren\par
\noindent\emph{Rusky / Russian:} прямо неразложимые факторы\par
\noindent\emph{Ukrajinsky / Ukrainian:} прямо нерозкладні фактори\par
\smallskip\noindent\emph{Razlog:} Direktno nerazložime faktory zachovujut uslovje direktnogo produkta iz odděla 4. Faktory ne sut izrično nazvane neprivedimymi v drugom smyslu. Prilog zato mora ostati pri pridavniku i v množině; skračenje do nerazložime faktory moglo by izpustiti nužno ograničenje.\par
\noindent\emph{Reason:} Direktno nerazložime faktory retains the direct-product qualification from section 4. It does not assert that the factors are irreducible in another sense. The adverb therefore remains attached to the adjective in the plural; shortening the phrase could remove the necessary restriction.\par
\noindent{\footnotesize Izvorne zapisy / Original records: HIST-g183-004.}\par
\end{minipage}\par
\par\medskip\Needspace{7\baselineskip}
\noindent\begin{minipage}{\linewidth}
\hypertarget{TERM-0294}{}
\noindent\textbf{direktna dvustranna suma}\quad{\footnotesize\texttt{TERM-0294}}\par
\noindent\emph{Iztočnik / Source:} direkt zweiseitig\par
\noindent\emph{Rusky / Russian:} прямая двусторонняя сумма\par
\noindent\emph{Ukrajinsky / Ukrainian:} пряма двобічна сума\par
\smallskip\noindent\emph{Razlog:} Direktna dvustranna suma zachovuje obe kvalifikacije, od ktoryh zavisi blokovo-diagonalna forma diskriminantnoj matrice. Dvustranna nastavlja terminologiju stran v teoriji idealov i prstenov. Razlog ne daje pravo skratiti izraz do direktna suma, ibo strannost jest relevantna konstrukciji.\par
\noindent\emph{Reason:} Direktna dvustranna suma preserves both qualifications governing the block-diagonal form of the discriminant matrix. Dvustranna continues the vocabulary for sides in ideal and ring theory. The rationale does not warrant shortening the expression to direct sum, because sidedness matters to the construction.\par
\noindent{\footnotesize Izvorne zapisy / Original records: HIST-g203-008.}\par
\end{minipage}\par
\par\medskip\Needspace{7\baselineskip}
\noindent\begin{minipage}{\linewidth}
\hypertarget{TERM-0295}{}
\noindent\textbf{direktno nerazložimy}\quad{\footnotesize\texttt{TERM-0295}}\par
\noindent\emph{Iztočnik / Source:} direkt-unzerlegbar\par
\noindent\emph{Rusky / Russian:} прямо неразложимый\par
\noindent\emph{Ukrajinsky / Ukrainian:} прямо нерозкладний\par
\smallskip\noindent\emph{Razlog:} Direktno nerazložimy izključaje direktnu sumu, ne vsako možno vnutrišnje razděljenje. Německa spojka v direkt-unzerlegbar ne měnja togo smysla. Zapis zachovuje prilog kako nositelja uslovja razloženja, a ne kako stilističnu čest, ktoru možno odstraniti.\par
\noindent\emph{Reason:} Direktno nerazložimy excludes a direct-sum decomposition, not every possible internal division. The German hyphen in direkt-unzerlegbar does not change that sense. The adverb carries the decomposition condition rather than serving as dispensable emphasis.\par
\noindent{\footnotesize Izvorne zapisy / Original records: HIST-g191-003.}\par
\end{minipage}\par
\par\medskip\Needspace{7\baselineskip}
\noindent\begin{minipage}{\linewidth}
\hypertarget{TERM-0296}{}
\noindent\textbf{direktno nerazložimy / razložimy}\quad{\footnotesize\texttt{TERM-0296}}\par
\noindent\emph{Iztočnik / Source:} direkt-unzerlegbar / zerlegbar\par
\noindent\emph{Rusky / Russian:} прямо неразложимый / разложимый\par
\noindent\emph{Ukrajinsky / Ukrainian:} прямо нерозкладний / розкладний\par
\smallskip\noindent\emph{Razlog:} Para direktno nerazložimy i razložimy čini vidimym različenje dekompozicije od neprivedimosti v argumentah o predstavjenjah. Izvorni zapis veže to različenje s rabotami 28 i 34. Samo podoben opis ‘ne možno razložiti’ ne ustanavjaje, ktora matematična operacija jest izključena.\par
\noindent\emph{Reason:} The pair direktno nerazložimy and razložimy keeps decomposition distinct from irreducibility in the representation-theoretic arguments. The recorded reason links the distinction to Papers 28 and 34. A general description such as cannot be broken up does not determine which mathematical operation is excluded.\par
\noindent{\footnotesize Izvorne zapisy / Original records: HIST-g187-004.}\par
\end{minipage}\par
\par\medskip\Needspace{7\baselineskip}
\noindent\begin{minipage}{\linewidth}
\hypertarget{TERM-0297}{}
\noindent\textbf{neposrědnje izračunanje / neposrědnje razvitje}\quad{\footnotesize\texttt{TERM-0297}}\par
\noindent\emph{Iztočnik / Source:} direkte Berechnung / direkte Entwicklung\par
\noindent\emph{Rusky / Russian:} прямое вычисление / прямое разложение\par
\noindent\emph{Ukrajinsky / Ukrainian:} пряме обчислення / прямий розклад\par
\smallskip\noindent\emph{Razlog:} Neposrědnje izračunanje i neposrědnje razvitje sut frazy běžnogo dokazovogo jezyka. Izvorni zapis ocěnjuje jih kako razumělive, ale ostavja prirodnost neposrědnje za stilistično sravnjenje. To jest zachovana sumnja o izraženju, ne novy tehničny smysl slova direktny.\par
\noindent\emph{Reason:} Neposrědnje izračunanje and neposrědnje razvitje belong to ordinary proof prose. The original record regards them as intelligible but leaves the naturalness of neposrědnje for stylistic comparison. That is a recorded uncertainty about expression, not a new technical sense of direct.\par
\noindent{\footnotesize Izvorne zapisy / Original records: HIST-g013-011.}\par
\end{minipage}\par
\par\medskip\Needspace{7\baselineskip}
\noindent\begin{minipage}{\linewidth}
\hypertarget{TERM-0298}{}
\noindent\textbf{neposrědnje svijanje}\quad{\footnotesize\texttt{TERM-0298}}\par
\noindent\emph{Iztočnik / Source:} direkte Faltung\par
\noindent\emph{Rusky / Russian:} пряма свертка\par
\noindent\emph{Ukrajinsky / Ukrainian:} безпосереднє згортання\par
\smallskip\noindent\emph{Razlog:} Neposrědnje svijanje ukazyvaje, že redukcije slědujut odmah iz operacije svijanja, ne iz pozdnějše izvedenoj redukcije. Neposrědnje zato nosi informaciju o argumentu. Historijna ruska paralela zapisana jest kako пряма свертка; ta grafija ostavaje dokazom zapisa, ne potvrđenoju ruskoju normoju.\par
\noindent\emph{Reason:} Neposrědnje svijanje says that the reductions arise immediately from the folding operation rather than from a later derived reduction. Neposrědnje thus carries information about the argument. The historical Russian comparison is written пряма свертка; that spelling is retained as evidence of the record, not certified Russian usage.\par
\noindent{\footnotesize Izvorne zapisy / Original records: HIST-g026-002.}\par
\end{minipage}\par
\par\medskip\Needspace{7\baselineskip}
\noindent\begin{minipage}{\linewidth}
\hypertarget{TERM-0299}{}
\noindent\textbf{neposrědno svijanje s reducentami}\quad{\footnotesize\texttt{TERM-0299}}\par
\noindent\emph{Iztočnik / Source:} direkte Faltung mit den Reduzenten\par
\noindent\emph{Rusky / Russian:} прямая свертка с редуцентами\par
\noindent\emph{Ukrajinsky / Ukrainian:} пряме згортання з редуцентами\par
\smallskip\noindent\emph{Razlog:} Neposrědno svijanje s reducentami zachovuje i bezposrědnost operacije, i rolu reducentov. Zapis smatra izraz jasnym, ale ostavja formu neposrědno za sravnjenje. Odstranjenje s reducentami izpustilo by učestvujuči objekty, a samo svijanje ne bi izrazilo načina dobyvanja rezultata.\par
\noindent\emph{Reason:} Neposrědno svijanje s reducentami retains both the operation's immediacy and the reducers' role. The record regards the expression as clear but leaves neposrědno for comparison. Omitting s reducentami would remove the objects involved; svijanje alone would not convey how the result is obtained.\par
\noindent{\footnotesize Izvorne zapisy / Original records: HIST-g021-004.}\par
\end{minipage}\par
\par\medskip\Needspace{7\baselineskip}
\noindent\begin{minipage}{\linewidth}
\hypertarget{TERM-0300}{}
\noindent\textbf{direktna suma}\quad{\footnotesize\texttt{TERM-0300}}\par
\noindent\emph{Iztočnik / Source:} direkte Summe\par
\noindent\emph{Rusky / Russian:} прямая сумма\par
\noindent\emph{Ukrajinsky / Ukrainian:} пряма сума\par
\smallskip\noindent\emph{Razlog:} Direktna suma jest izbrana kako prepoznavajemo algebraično nazvanje. Dva razloga povezuju njegovo priměnjenje v rabotah 31 i 34 i razlučajut sumu kako aditivnu konstrukciju od direktnogo produkta. Tvrđenje o širokoj razumělivosti jest razlogom izbora, ne rezultatom izměrenogo čitatelskogo testa.\par
\noindent\emph{Reason:} Direktna suma is selected as a recognizable algebraic name. Two rationales connect its use in Papers 31 and 34 and distinguish the additive construction from direct product. Broad intelligibility is a reason advanced for the choice, not the result of a measured reader test.\par
\noindent{\footnotesize Izvorne zapisy / Original records: HIST-g121-010, HIST-g182-004.}\par
\end{minipage}\par
\par\medskip\Needspace{7\baselineskip}
\noindent\begin{minipage}{\linewidth}
\hypertarget{TERM-0301}{}
\noindent\textbf{direktna suma / direktny sumand}\quad{\footnotesize\texttt{TERM-0301}}\par
\noindent\emph{Iztočnik / Source:} direkte Summe / direkter Summand\par
\noindent\emph{Rusky / Russian:} прямая сумма / прямое слагаемое\par
\noindent\emph{Ukrajinsky / Ukrainian:} пряма сума / прямий доданок\par
\smallskip\noindent\emph{Razlog:} Direktna suma imenuje cělost, direktny sumand jednu njezinu komponentu. Izvorni razlog daje prednost medžunarodnym osnovam direkt- i sumand i govori o sravnjenju s češskym, poljskym, ruskym i ukrajinskym jezykom. V zachovanyh poljah sut točne formy samo za poslednje dva; bez citatov češskogo i poljskogo iztočnika ne možno prědstaviti tu tvrdnju kako polno potvrđeno sravnjenje.\par
\noindent\emph{Reason:} Direktna suma names the whole and direktny sumand a component. The recorded reason favours international direkt- and sumand and invokes Czech, Polish, Russian, and Ukrainian comparison. The preserved fields supply exact forms only for the last two; without the cited Czech and Polish sources, that assertion cannot be presented as a fully evidenced comparison.\par
\noindent{\footnotesize Izvorne zapisy / Original records: HIST-g184-004.}\par
\end{minipage}\par
\par\medskip\Needspace{7\baselineskip}
\noindent\begin{minipage}{\linewidth}
\hypertarget{TERM-0302}{}
\noindent\textbf{direktna suma direktno nerazložimyh jednostavnyh idealov}\quad{\footnotesize\texttt{TERM-0302}}\par
\noindent\emph{Iztočnik / Source:} direkte Summe direkt unzerlegbarer, einfacher Ideale\par
\noindent\emph{Rusky / Russian:} прямая сумма прямо неразложимых простых идеалов\par
\noindent\emph{Ukrajinsky / Ukrainian:} пряма сума прямо нерозкладних простих ідеалів\par
\smallskip\noindent\emph{Razlog:} Jednostavny ideal ovde znači simple ideal, a prosty ideal ostavaje nazvanjem prime ideal v sosednyh rabotah. To jest naměrno razděljenje dvuh smyslov, ktore ruska i ukrajinska paralela ne morajut razděliti različnymi osnovami. Cěly izraz takojže zachovuje direktnu sumu i direktnu nerazložimost, ne zaměnjajuči jih neprivedimostju.\par
\noindent\emph{Reason:} Jednostavny ideal here means a simple ideal, while prosty ideal remains the name for a prime ideal in the neighbouring works. This deliberately separates senses that the Russian and Ukrainian comparanda need not distinguish by different stems. The full expression also preserves direct sum and direct indecomposability rather than substituting irreducibility.\par
\noindent{\footnotesize Izvorne zapisy / Original records: HIST-g105-006.}\par
\end{minipage}\par
\par\medskip\Needspace{7\baselineskip}
\noindent\begin{minipage}{\linewidth}
\hypertarget{TERM-0303}{}
\noindent\textbf{direktna suma konečno mnogyh polj}\quad{\footnotesize\texttt{TERM-0303}}\par
\noindent\emph{Iztočnik / Source:} direkte Summe von endlich vielen Körpern\par
\noindent\emph{Rusky / Russian:} прямая сумма конечного числа полей\par
\noindent\emph{Ukrajinsky / Ukrainian:} пряма сума скінченно багатьох полів\par
\smallskip\noindent\emph{Razlog:} Direktna suma konečno mnogyh polj zachovuje vtoroju harakterizaciju iz definicije. Konečnost odnašaje se k čislu polj, ne avtomatično k čislu elementov v každom polju. Obyčny slovoporedok čini tuto ograničenje čitatelnym bez skračenja matematičnogo sodržanja.\par
\noindent\emph{Reason:} Direktna suma konečno mnogyh polj retains the definition's second characterization. Finiteness qualifies the number of fields, not automatically the number of elements in each field. Ordinary word order makes that restriction readable without abbreviating the mathematics.\par
\noindent{\footnotesize Izvorne zapisy / Original records: HIST-g125-004.}\par
\end{minipage}\par
\par\medskip\Needspace{7\baselineskip}
\noindent\begin{minipage}{\linewidth}
\hypertarget{TERM-0304}{}
\noindent\textbf{direktna suma kolc ranga jedan}\quad{\footnotesize\texttt{TERM-0304}}\par
\noindent\emph{Iztočnik / Source:} direkte Summe von Ringen vom Rang eins\par
\noindent\emph{Rusky / Russian:} прямая сумма колец ранга один\par
\noindent\emph{Ukrajinsky / Ukrainian:} пряма сума кілець рангу один\par
\smallskip\noindent\emph{Razlog:} Direktna suma kolc ranga jedan nastavlja terminologiju razloženja iz odděla 3. Rang jedan kvalifikuje prsteny, a ne čislo sumandov; tu razliku ne možno izgubiti pri skračenju. Kolc jest zachovana historijna forma, ne dokaz aktualnoj prednosti kolco nad prsten.\par
\noindent\emph{Reason:} Direktna suma kolc ranga jedan continues the decomposition vocabulary of section 3. Rank one qualifies the rings, not the number of summands; shortening must not lose that distinction. Kolc is the preserved historical form, not evidence of a current preference for kolco over prsten.\par
\noindent{\footnotesize Izvorne zapisy / Original records: HIST-g146-004.}\par
\end{minipage}\par
\par\medskip\Needspace{7\baselineskip}
\noindent\begin{minipage}{\linewidth}
\hypertarget{TERM-0305}{}
\noindent\textbf{direktne sumy klasov}\quad{\footnotesize\texttt{TERM-0305}}\par
\noindent\emph{Iztočnik / Source:} Direkte Summen der Klassen\par
\noindent\emph{Rusky / Russian:} прямые суммы классов\par
\noindent\emph{Ukrajinsky / Ukrainian:} прямі суми класів\par
\smallskip\noindent\emph{Razlog:} Direktne sumy klasov zachovujut klasy kako objekty razmatranja v oddělu 5. Izraz nastavlja direktnu sumu iz oddělov 2 i 3, ale ne podměnjaje klas konkretnym njegovym prědstavitelem. Ta razlika jest čestju nazvanja, ne samo okolnogo objasnjenja.\par
\noindent\emph{Reason:} Direktne sumy klasov keeps classes as the objects under discussion in section 5. The expression continues the direct-sum terminology of sections 2 and 3 without silently substituting a particular representative for a class. That distinction belongs to the naming itself.\par
\noindent{\footnotesize Izvorne zapisy / Original records: HIST-g138-001.}\par
\end{minipage}\par
\par\medskip\Needspace{7\baselineskip}
\noindent\begin{minipage}{\linewidth}
\hypertarget{TERM-0306}{}
\noindent\textbf{predstavjenje razširjenyh kolc kako direktne sumy}\quad{\footnotesize\texttt{TERM-0306}}\par
\noindent\emph{Iztočnik / Source:} Direkte Summendarstellung von Erweiterungsringen\par
\noindent\emph{Rusky / Russian:} представление колец расширения в виде прямых сумм\par
\noindent\emph{Ukrajinsky / Ukrainian:} подання кілець розширення як прямих сум\par
\smallskip\noindent\emph{Razlog:} Predstavjenje ... kako direktne sumy razkriva německo složeno slovo čitatelnoj konstrukcijeju s kako. Razširjene prsteny i direktne sumy ostavajut izrično imenovane. Razlog daje prednost takvoj slovosvezi vměsto zbitogo složenogo nazvanja; ne zaměnjaje algebraično razloženje svobodnym opisom.\par
\noindent\emph{Reason:} Predstavjenje ... kako direktne sumy unpacks the German compound into a readable as-construction. The extension rings and direct sums remain explicitly named. The reason favours that syntax over a compressed compound without replacing the algebraic decomposition by a loose description.\par
\noindent{\footnotesize Izvorne zapisy / Original records: HIST-g123-001.}\par
\end{minipage}\par
\par\medskip\Needspace{7\baselineskip}
\noindent\begin{minipage}{\linewidth}
\hypertarget{TERM-0307}{}
\noindent\textbf{direktno sumno razloženje}\quad{\footnotesize\texttt{TERM-0307}}\par
\noindent\emph{Iztočnik / Source:} direkte Summenzerlegung\par
\noindent\emph{Rusky / Russian:} прямое суммарное разложение\par
\noindent\emph{Ukrajinsky / Ukrainian:} прямий сумарний розклад\par
\smallskip\noindent\emph{Razlog:} Direktno sumno razloženje zachovuje i pojem razloženja, i njegovu vezu s direktnoju sumoju. Izvorni razlog daje prednost medžunarodnomu direktno prěd neimenovanoju domorodnoju alternativoju. On ne soderži točny kandidat ili čislenu ocěnu toj alternativy; ih ne trěba dodavati kako by byli historijno zapisane.\par
\noindent\emph{Reason:} Direktno sumno razloženje retains decomposition and its connection with direct sum. The recorded reason favours international direktno over an unnamed native alternative. It supplies neither an exact alternative form nor a numerical rating; neither should be invented as historical evidence.\par
\noindent{\footnotesize Izvorne zapisy / Original records: HIST-g200-009.}\par
\end{minipage}\par
\par\medskip\Needspace{7\baselineskip}
\noindent\begin{minipage}{\linewidth}
\hypertarget{TERM-0308}{}
\noindent\textbf{direktno razloženje na prave idealy}\quad{\footnotesize\texttt{TERM-0308}}\par
\noindent\emph{Iztočnik / Source:} direkte Summenzerlegung in Rechtsideale\par
\noindent\emph{Rusky / Russian:} прямое разложение на правые идеалы\par
\noindent\emph{Ukrajinsky / Ukrainian:} прямий розклад на праві ідеали\par
\smallskip\noindent\emph{Razlog:} Direktno razloženje na prave idealy zachovuje konkretny tip idealov, na ktore objekt razkladaje se. Izraz nastavlja direktnu sumu iz oddělov 4–6, ne novo opče značenje direktnogo razloženja. Prave ovde odnašaje se k pravoj straně množenja, ne k ocěně pravilnosti ideala.\par
\noindent\emph{Reason:} Direktno razloženje na prave idealy retains the particular kind of ideals into which the object is decomposed. It continues the direct-sum terminology of sections 4–6 rather than introducing a general sense of direct decomposition. Prave concerns the right side of multiplication, not whether an ideal is correct.\par
\noindent{\footnotesize Izvorne zapisy / Original records: HIST-g187-002.}\par
\end{minipage}\par
\par\medskip\Needspace{7\baselineskip}
\noindent\begin{minipage}{\linewidth}
\hypertarget{TERM-0309}{}
\noindent\textbf{direktny prěsěk}\quad{\footnotesize\texttt{TERM-0309}}\par
\noindent\emph{Iztočnik / Source:} direkter Durchschnitt\par
\noindent\emph{Rusky / Russian:} прямое пересечение\par
\noindent\emph{Ukrajinsky / Ukrainian:} прямий перетин\par
\smallskip\noindent\emph{Razlog:} Direktny prěsěk zachovuje odnos k direktnoj sumě, o ktorom govori Prüferova podnožna primětka, i nastavlja terminologiju raboty 30 v rabotě 34. Vtory razlog izrično upominaje češsku, poljsku i južnoslovjansku paralelu, ale ne dava jejich točne formy ili citaty. Zapisany kiriličny oblik пресек jest dokazom izbora pisanja v tom zapisu, ne dodatnym dokazom leksičnogo značenja.\par
\noindent\emph{Reason:} Direktny prěsěk preserves the relation to direct sum discussed in Prüfer's footnote and continues Paper 30 terminology in Paper 34. The second reason invokes Czech, Polish, and South Slavic parallels but supplies neither their exact forms nor citations. The recorded Cyrillic пресек is evidence of that record's spelling choice, not additional evidence for lexical meaning.\par
\noindent{\footnotesize Izvorne zapisy / Original records: HIST-g111-009, HIST-g182-002.}\par
\end{minipage}\par
\par\medskip\Needspace{7\baselineskip}
\noindent\begin{minipage}{\linewidth}
\hypertarget{TERM-0310}{}
\noindent\textbf{direktny faktor}\quad{\footnotesize\texttt{TERM-0310}}\par
\noindent\emph{Iztočnik / Source:} direkter Faktor\par
\noindent\emph{Rusky / Russian:} прямой фактор\par
\noindent\emph{Ukrajinsky / Ukrainian:} прямий фактор\par
\smallskip\noindent\emph{Razlog:} Direktny faktor nastavlja terminologiju direktnogo produkta iz odděla 4. Izvorni razlog izrično ne dopuščaje zaměnu slovom dopolnjenje: Noether tvrdži strukturu direktnogo faktora, ne samo suščestvovanje nekogo dopolnjujučego objekta. Prěvod zato zachovuje imenovanu algebraičnu rolu.\par
\noindent\emph{Reason:} Direktny faktor continues the direct-product vocabulary of section 4. The recorded reason explicitly rejects complement as a paraphrase: Noether states direct-factor structure, not merely the presence of some complementary object. The translation therefore retains the named algebraic role.\par
\noindent{\footnotesize Izvorne zapisy / Original records: HIST-g183-003.}\par
\end{minipage}\par
\par\medskip\Needspace{7\baselineskip}
\noindent\begin{minipage}{\linewidth}
\hypertarget{TERM-0311}{}
\noindent\textbf{direktny sumand}\quad{\footnotesize\texttt{TERM-0311}}\par
\noindent\emph{Iztočnik / Source:} direkter Summand\par
\noindent\emph{Rusky / Russian:} прямое слагаемое\par
\noindent\emph{Ukrajinsky / Ukrainian:} прямий доданок\par
\smallskip\noindent\emph{Razlog:} Direktny sumand zachovuje vezu s direktna suma. Dva zapisy iz oddělov 13 i 24 davajut prednost kratkomu medžunarodnomu nazvanju komponenta, dokolě ruska i ukrajinska paralela imajut svoje obične algebraične imena. To jest kontinuitet terminologičnoj familije, ne novy pojem pri každom pojavjenju.\par
\noindent\emph{Reason:} Direktny sumand keeps its connection with direktna suma. Records from sections 13 and 24 favour a compact international name for the component, while the Russian and Ukrainian comparanda use their own algebraic nouns. This is continuity within a terminology family, not a new concept at each occurrence.\par
\noindent{\footnotesize Izvorne zapisy / Original records: HIST-g191-002, HIST-g202-007.}\par
\end{minipage}\par
\par\medskip\Needspace{7\baselineskip}
\noindent\begin{minipage}{\linewidth}
\hypertarget{TERM-0312}{}
\noindent\textbf{direktny summand}\quad{\footnotesize\texttt{TERM-0312}}\par
\noindent\emph{Iztočnik / Source:} direkter Summand\par
\noindent\emph{Rusky / Russian:} прямое слагаемое\par
\noindent\emph{Ukrajinsky / Ukrainian:} прямий доданок\par
\smallskip\noindent\emph{Razlog:} Direktny summand vyražaje isti izbor medžunarodnogo imennika vměsto novogo domorodnogo tvorjenja. Izvorni razlog govori o prepoznavajemosti v algebraičnom kontekstu, ale ne objašnjaje dvojno m. Proto summand s dvojnym m i sumand s jednym m trěba razmatrjati kako pravopisny problem, ne kako dva matematične termina.\par
\noindent\emph{Reason:} Direktny summand expresses the same preference for an international noun over a new native formation. The recorded reason invokes recognizability in algebraic contexts but does not justify the doubled m. Summand and sumand therefore present an orthographic question, not two mathematical concepts.\par
\noindent{\footnotesize Izvorne zapisy / Original records: HIST-g194-004.}\par
\end{minipage}\par
\par\medskip\Needspace{7\baselineskip}
\noindent\begin{minipage}{\linewidth}
\hypertarget{TERM-0313}{}
\noindent\textbf{direktny produkt}\quad{\footnotesize\texttt{TERM-0313}}\par
\noindent\emph{Iztočnik / Source:} direktes Produkt\par
\noindent\emph{Rusky / Russian:} прямое произведение\par
\noindent\emph{Ukrajinsky / Ukrainian:} прямий добуток\par
\smallskip\noindent\emph{Razlog:} Direktny produkt zachovuje medžunarodno nazvanje algebraičnoj konstrukcije. Rusko proizvedenie i ukrajinsko dobutok upotrěbjajut druge osnovy; razlog dava prednost produkt za medžuslovjansku tehničnu prozu. To ne jest tvrdženje, že priměr ne trěbuje definicije ili že vsaky čitatelj uže znaje pojem.\par
\noindent\emph{Reason:} Direktny produkt retains an international name for the algebraic construction. Russian proizvedenie and Ukrainian dobutok use different stems; the reason favours produkt for Interslavic technical prose. This does not establish that the construction needs no definition or that every reader already knows it.\par
\noindent{\footnotesize Izvorne zapisy / Original records: HIST-g182-001.}\par
\end{minipage}\par
\par\medskip\Needspace{7\baselineskip}
\noindent\begin{minipage}{\linewidth}
\hypertarget{TERM-0314}{}
\noindent\textbf{direktny produkt / tvorjenje direktnogo produkta}\quad{\footnotesize\texttt{TERM-0314}}\par
\noindent\emph{Iztočnik / Source:} direktes Produkt / direkte Produktbildung\par
\noindent\emph{Rusky / Russian:} прямое произведение / образование прямого произведения\par
\noindent\emph{Ukrajinsky / Ukrainian:} прямий добуток / утворення прямого добутку\par
\smallskip\noindent\emph{Razlog:} Direktny produkt jest objekt, tvorjenje direktnogo produkta njegova konstrukcija. Oba morajut ostati vidime, ibo pasaža ne govori samo o gotovom rezultatu. Produkt nastavlja imennik iz terminologije skrěščenogo produkta; direktny označaje drugy tip konstrukcije bez novogo nazvanja samogo produkta.\par
\noindent\emph{Reason:} Direktny produkt names the object; tvorjenje direktnogo produkta names its construction. Both remain explicit because the passage concerns more than the finished result. Produkt continues the noun used for crossed product, while direktny specifies a different construction without renaming product itself.\par
\noindent{\footnotesize Izvorne zapisy / Original records: HIST-g213-006.}\par
\end{minipage}\par
\par\medskip\Needspace{7\baselineskip}
\noindent\begin{minipage}{\linewidth}
\hypertarget{TERM-0315}{}
\noindent\textbf{diskriminant}\quad{\footnotesize\texttt{TERM-0315}}\par
\noindent\emph{Iztočnik / Source:} Diskriminante\par
\noindent\emph{Rusky / Russian:} дискриминант\par
\noindent\emph{Ukrajinsky / Ukrainian:} дискримінант\par
\smallskip\noindent\emph{Razlog:} Diskriminant zachovuje medžunarodnu osnovu i mužsky rod, už upotrěbjeny v rabotě 31. Vtory razlog izrično odbacaje vvedenje konkurentnoj ženske formy; prvi odbacaje novy opisny prěvod invariantnogo nazvanja. Oba razloge zato podpirajut i leksičny kontinuitet, i konkretnu imennikovu formu.\par
\noindent\emph{Reason:} Diskriminant retains the international stem and the masculine form already used in Paper 31. The second rationale explicitly rejects introducing a competing feminine form; the first rejects a new descriptive paraphrase for the invariant. Together they support lexical continuity and a particular noun form.\par
\noindent{\footnotesize Izvorne zapisy / Original records: HIST-g130-003, HIST-g203-001.}\par
\end{minipage}\par
\par\medskip\Needspace{7\baselineskip}
\noindent\begin{minipage}{\linewidth}
\hypertarget{TERM-0316}{}
\noindent\textbf{diskriminant $D_T$ / kvadrat determinanta konjugovanyh baznyh elementov}\quad{\footnotesize\texttt{TERM-0316}}\par
\noindent\emph{Iztočnik / Source:} Diskriminante $D_T$ / Determinantenquadrat konjugierter Basiselemente\par
\noindent\emph{Rusky / Russian:} дискриминант $D_T$ / квадрат определителя сопряженных базисных элементов\par
\noindent\emph{Ukrajinsky / Ukrainian:} дискримінант $D_T$ / квадрат визначника спряжених базисних елементів\par
\smallskip\noindent\emph{Razlog:} Diskriminant $D_T$ jest svezany s kvadratom determinanta konjugovanyh baznyh elementov. Zapis prěnosi terminologiju iz odděla 6, punkta 4 do odděla 7 pri prěhodu od prstenov konečnogo ranga k poredkam. Kvadrat zato ostavaje izričny: determinanta bez kvadrata ne jest taj sam objekt.\par
\noindent\emph{Reason:} Diskriminant $D_T$ is connected with the square of the determinant of conjugate basis elements. The record carries the vocabulary from section 6, point 4 into section 7 as the discussion passes from finite-rank rings to orders. Square remains explicit because the unsquared determinant is not the same object.\par
\noindent{\footnotesize Izvorne zapisy / Original records: HIST-g149-008.}\par
\end{minipage}\par
\par\medskip\Needspace{7\baselineskip}
\noindent\begin{minipage}{\linewidth}
\hypertarget{TERM-0317}{}
\noindent\textbf{diskriminant $D_{T_P}$}\quad{\footnotesize\texttt{TERM-0317}}\par
\noindent\emph{Iztočnik / Source:} Diskriminante $D_{T_P}$\par
\noindent\emph{Rusky / Russian:} дискриминант $D_{T_P}$\par
\noindent\emph{Ukrajinsky / Ukrainian:} дискримінант $D_{T_P}$\par
\smallskip\noindent\emph{Razlog:} Diskriminant $D_{T_P}$ imenuje skalarny diskriminant, ne diskriminantny ideal. Indeks identifikuje odgovarjajuči objekt i ne može byti izgubjen pri skračenju nazvanja. Isti korenj diskriminant- v oběh terminah ne svodi skalar i ideal na jednu věšč.\par
\noindent\emph{Reason:} Diskriminant $D_{T_P}$ names the scalar discriminant, not the discriminant ideal. Its index identifies the relevant object and cannot be lost in shortening the name. The shared diskriminant- stem does not collapse scalar and ideal into one entity.\par
\noindent{\footnotesize Izvorne zapisy / Original records: HIST-g162-001.}\par
\end{minipage}\par
\par\medskip\Needspace{7\baselineskip}
\noindent\begin{minipage}{\linewidth}
\hypertarget{TERM-0318}{}
\noindent\textbf{diskriminant ideala vzhodno k klasu}\quad{\footnotesize\texttt{TERM-0318}}\par
\noindent\emph{Iztočnik / Source:} Diskriminante eines Ideals bezüglich einer Klasse\par
\noindent\emph{Rusky / Russian:} дискриминант идеала относительно класса\par
\noindent\emph{Ukrajinsky / Ukrainian:} дискримінант ідеалу відносно класу\par
\smallskip\noindent\emph{Razlog:} Diskriminant ideala ... k klasu zachovuje klas kako izričny parametr, paralelno konstrukcijam slěda i normy. Historijni kandidat ima vzhodno k; zapis objašnjaje potrěbu vyraziti odnos, ale ne dava jezyčnogo dokaza za tuto predložnu konstrukciju. Razlog za matematičnu funkciju proto trěba razlučiti od opravdanja konkretnyh slov.\par
\noindent\emph{Reason:} Diskriminant ideala ... k klasu retains the class as an explicit parameter, parallel to the trace and norm constructions. The historical candidate has vzhodno k. Its rationale explains why the relation must be expressed but supplies no linguistic evidence for that prepositional phrase. The mathematical function and the justification of the wording are distinct questions.\par
\noindent{\footnotesize Izvorne zapisy / Original records: HIST-g137-001.}\par
\smallskip\noindent\emph{Sravnjenje s tekstom:} Aktualny tekst v dělu 31 ima po odnošenju k klasu i po odnošenju k osnovnomu prstenu. Kopija slovnika citovana v rěšenju o toj konstrukciji davaje oba zaglavja po odnošeńju k (18227) i v odnošeńju k (19113), s dativom i smyslom ‘po odnošenju k’. Zato argument, že samo prva forma jest oficialna, ne podpiraje izbora; konsistentna prednost jednoj formě može byti redakcijnoju konvencijeju. Osobny prilog vȯzhodno (4185) znači ‘na vȯzhod’, ne ‘po odnošenju k’. To ne opravdaje ranějše vzhodno k. Kratke odnosno (4107) ima rodni padež i zato ne jest mehanična zaměna frazy s dativom. Čitatelsko po odnošenju k slěduje osnovnomu pravopisu, kde ń prěhodi v n. Sravnjenje podpira značenje i padežno upravjanje, ne dokaz izměrenoj prednosti v razumělivosti.\par
\noindent\emph{Textual comparison:} The current Paper 31 text has po odnošenju k klasu and po odnošenju k with the base ring. The dictionary copy cited by the decision supplies both po odnošeńju k (18227) and v odnošeńju k (19113), each with the dative and the meaning with regard to. The claim that only the former is official therefore cannot justify the choice; consistently preferring one remains a possible editorial convention. The separate adverb vȯzhodno (4185) means eastwards, not with respect to, and does not justify historical vzhodno k. Shorter odnosno (4107) governs the genitive, so it is not a mechanical replacement for a dative construction. Reader spelling po odnošenju k follows the basic-orthography treatment of ń as n. The comparison supports meaning and government, not measured superiority in intelligibility.\par
\noindent{\footnotesize TEXT-COMP-06; pasaže / passages: 18094, 18787; \texttt{CORPUS\_FORM\_COMPARISONS.json}.}\par
\noindent{\footnotesize\href{https://interslavic-dictionary.com/}{Medžuslovjansky slovnik / Interslavic Dictionary}: 18227, 19113, 4107, 4185; točne kopije / exact snapshots in supporting data.}\par
\noindent{\footnotesize\href{https://interslavic.fun/learn/orthography/}{Pravopis / Orthography}: opcionalne bukvy / optional letters.}\par
\end{minipage}\par
\par\medskip\Needspace{7\baselineskip}
\noindent\begin{minipage}{\linewidth}
\hypertarget{TERM-0319}{}
\noindent\textbf{diskriminant R izčezaje}\quad{\footnotesize\texttt{TERM-0319}}\par
\noindent\emph{Iztočnik / Source:} Diskriminante von R verschwindet\par
\noindent\emph{Rusky / Russian:} дискриминант R обращается в нуль\par
\noindent\emph{Ukrajinsky / Ukrainian:} дискримінант R зникає\par
\smallskip\noindent\emph{Razlog:} Diskriminant $R$ izčezaje zachovuje uslovje nulovosti diskriminanta. Izvorni razlog naměrno ne měnja to na vyraženje o nenulovoj determinantě. Prěvod mora zachovati i imenovany invariant, i směr kriterija; podobnost algebraičnyh objektov ne opravdaje zaměnu.\par
\noindent\emph{Reason:} Diskriminant $R$ izčezaje preserves vanishing of the discriminant. The recorded reason deliberately avoids changing it to a statement about a nonzero determinant. The translation must retain both the named invariant and the direction of the criterion; related algebraic objects are not interchangeable.\par
\noindent{\footnotesize Izvorne zapisy / Original records: HIST-g154-006.}\par
\end{minipage}\par
\par\medskip\Needspace{7\baselineskip}
\noindent\begin{minipage}{\linewidth}
\hypertarget{TERM-0320}{}
\noindent\textbf{diskriminantny ideal}\quad{\footnotesize\texttt{TERM-0320}}\par
\noindent\emph{Iztočnik / Source:} Diskriminantenideal\par
\noindent\emph{Rusky / Russian:} дискриминантный идеал\par
\noindent\emph{Ukrajinsky / Ukrainian:} дискримінантний ідеал\par
\smallskip\noindent\emph{Razlog:} Diskriminantny ideal tvori se pridavnikom i imennikom, ne zaměnnoju frazoju o generovanji. Definicija opreděljaje ideal generovany diskriminantnymi determinantami, ale vsako daljše imenovanje ne mora ponoviti cělu definiciju. Dva zapisy povezuju isto kratko nazvanje s formalnym tvorjenjem ideala.\par
\noindent\emph{Reason:} Diskriminantny ideal uses an adjective and a noun rather than replacing the name with a generating description. The definition specifies the ideal generated by discriminant determinants, but later references need not repeat the entire definition. Two records connect the same short name with that formal construction.\par
\noindent{\footnotesize Izvorne zapisy / Original records: HIST-g157-002, HIST-g158-001.}\par
\end{minipage}\par
\par\medskip\Needspace{7\baselineskip}
\noindent\begin{minipage}{\linewidth}
\hypertarget{TERM-0321}{}
\noindent\textbf{diskriminantny ideal ... nulovy ideal}\quad{\footnotesize\texttt{TERM-0321}}\par
\noindent\emph{Iztočnik / Source:} Diskriminantenideal ... Nullideal\par
\noindent\emph{Rusky / Russian:} дискриминантный идеал ... нулевой идеал\par
\noindent\emph{Ukrajinsky / Ukrainian:} дискримінантний ідеал ... нульовий ідеал\par
\smallskip\noindent\emph{Razlog:} Diskriminantny ideal i nulovy ideal zachovujut dve nazvanja nužne za konkretno tvrdženje. Jednina jest naměrna: paragraf govori o jednom diskriminantnom idealu, iako naslov razmatra familiju idealov. Čislo imennika mora slědovati objektu pasaže, ne mehanično formě naslova.\par
\noindent\emph{Reason:} Diskriminantny ideal and nulovy ideal retain the two names needed by the statement. The singular is deliberate: the paragraph concerns one discriminant ideal even though the heading treats a family. Grammatical number follows the passage's object, not mechanically the heading's form.\par
\noindent{\footnotesize Izvorne zapisy / Original records: HIST-g143-002.}\par
\end{minipage}\par
\par\medskip\Needspace{7\baselineskip}
\noindent\begin{minipage}{\linewidth}
\hypertarget{TERM-0322}{}
\noindent\textbf{diskriminantny ideal komponentnogo ideala}\quad{\footnotesize\texttt{TERM-0322}}\par
\noindent\emph{Iztočnik / Source:} Diskriminantenideal des Komponentenideals\par
\noindent\emph{Rusky / Russian:} дискриминантный идеал компонентного идеала\par
\noindent\emph{Ukrajinsky / Ukrainian:} дискримінантний ідеал компонентного ідеалу\par
\smallskip\noindent\emph{Razlog:} Diskriminantny ideal komponentnogo ideala zachovuje dva urovnja opisa: komponent jest ideal, a od njega tvori se diskriminantny ideal. Odgovarjajuči klas ostavaje parametrom konstrukcije. Skračenje do diskriminant komponenta moglo by skryti i tip objekta, i njegov odnos k klasu.\par
\noindent\emph{Reason:} Diskriminantny ideal komponentnogo ideala preserves two levels of description: the component is an ideal, and its discriminant ideal is constructed. The relevant class remains a parameter. Shortening to discriminant of the component could obscure both the type of object and its relation to the class.\par
\noindent{\footnotesize Izvorne zapisy / Original records: HIST-g139-003.}\par
\end{minipage}\par
\par\medskip\Needspace{7\baselineskip}
\noindent\begin{minipage}{\linewidth}
\hypertarget{TERM-0323}{}
\noindent\textbf{diskriminantny ideal kolca}\quad{\footnotesize\texttt{TERM-0323}}\par
\noindent\emph{Iztočnik / Source:} Diskriminantenideal des Ringes\par
\noindent\emph{Rusky / Russian:} дискриминантный идеал кольца\par
\noindent\emph{Ukrajinsky / Ukrainian:} дискримінантний ідеал кільця\par
\smallskip\noindent\emph{Razlog:} Diskriminantny ideal kolca imenuje konstrukciju na urovnju prstena, polagajuču glavny klas v osnovu. Izvorni zapis svezava nazvanje s konečnoju větoju iz odděla 5 i njegovoju upotrěboju v oddělu 6. Njegovo pozvanje na pečatny iztočnik jest čestju istorijnogo dokaza; ne predstavlja novu prověrku skana pri pisanju ovogo objasnjenja.\par
\noindent\emph{Reason:} Diskriminantny ideal kolca names the ring-level construction based on the principal class. The record connects the term with the final sentence of section 5 and its use in section 6. Its appeal to the printed source belongs to the historical evidence; it is not a fresh examination of the scan for this explanation.\par
\noindent{\footnotesize Izvorne zapisy / Original records: HIST-g140-005.}\par
\end{minipage}\par
\par\medskip\Needspace{7\baselineskip}
\noindent\begin{minipage}{\linewidth}
\hypertarget{TERM-0324}{}
\noindent\textbf{diskriminantny ideal jest jediničny ideal v P}\quad{\footnotesize\texttt{TERM-0324}}\par
\noindent\emph{Iztočnik / Source:} Diskriminantenideal ist das Einheitsideal in P\par
\noindent\emph{Rusky / Russian:} дискриминантный идеал является единичным идеалом в P\par
\noindent\emph{Ukrajinsky / Ukrainian:} дискримінантний ідеал є одиничним ідеалом у P\par
\smallskip\noindent\emph{Razlog:} Diskriminantny ideal jest jediničny ideal v $P$ stavja jednakost idealov na prvo město. Tek potom slěduje ekvivalentna formulacija s skalarnym diskriminantom. Prěvod zachovuje toj poredok i ne podměnjaje ideal čislom; jediničny ideal jest tehnično nazvanje, ne prosto jedinica v proze.\par
\noindent\emph{Reason:} Diskriminantny ideal jest jediničny ideal v $P$ puts equality of ideals first, before the equivalent scalar-discriminant formulation. The translation preserves that order and does not substitute a number for an ideal. Jediničny ideal is a technical name, not simply the number one in prose.\par
\noindent{\footnotesize Izvorne zapisy / Original records: HIST-g144-005.}\par
\end{minipage}\par
\par\medskip\Needspace{7\baselineskip}
\noindent\begin{minipage}{\linewidth}
\hypertarget{TERM-0325}{}
\noindent\textbf{diskriminantny ideal od R}\quad{\footnotesize\texttt{TERM-0325}}\par
\noindent\emph{Iztočnik / Source:} Diskriminantenideal von R\par
\noindent\emph{Rusky / Russian:} дискриминантный идеал R\par
\noindent\emph{Ukrajinsky / Ukrainian:} дискримінантний ідеал R\par
\smallskip\noindent\emph{Razlog:} Diskriminantny ideal od $R$ prěnosi nazvanje prstenovogo objekta iz odděla 5, punkta 3 do formulacije kriterija. Identifikator $R$ ostavaje izričny. Historijni razlog govori o oporě v pečatnom iztočniku; tuto objasnjenje zachovuje jeho argument, ne tvrdži novu nezavisnu kolaciju.\par
\noindent\emph{Reason:} Diskriminantny ideal od $R$ carries the ring-level naming from section 5, point 3 into the criterion. The identifier $R$ remains explicit. The historical rationale invokes the printed source; this explanation preserves that reasoning without claiming a new independent collation.\par
\noindent{\footnotesize Izvorne zapisy / Original records: HIST-g148-004.}\par
\end{minipage}\par
\par\medskip\Needspace{7\baselineskip}
\noindent\begin{minipage}{\linewidth}
\hypertarget{TERM-0326}{}
\noindent\textbf{diskriminantny ideal porjadka T}\quad{\footnotesize\texttt{TERM-0326}}\par
\noindent\emph{Iztočnik / Source:} Diskriminantenideal von T\par
\noindent\emph{Rusky / Russian:} дискриминантный идеал порядка T\par
\noindent\emph{Ukrajinsky / Ukrainian:} дискримінантний ідеал порядку T\par
\smallskip\noindent\emph{Razlog:} Diskriminantny ideal porjadka $T$ dodaje nazvanje tipa objekta, ktory jest už definovany prěd tym. Porjadka zato ima objasnjujuču funkciju; ne smije proměniti referent $T$ ili vnesti novo prědpoloženje. Izbor jest lokalno razkrytje informacije, ne dozvoljenje proizvoljno razširjati avtorske tvrdženja.\par
\noindent\emph{Reason:} Diskriminantny ideal porjadka $T$ supplies the kind of object already defined in the preceding text. Porjadka is explanatory: it must not change the referent of $T$ or introduce a new hypothesis. The choice is local explicitation, not permission to expand the author's claims arbitrarily.\par
\noindent{\footnotesize Izvorne zapisy / Original records: HIST-g164-002.}\par
\end{minipage}\par
\par\medskip\Needspace{7\baselineskip}
\noindent\begin{minipage}{\linewidth}
\hypertarget{TERM-0327}{}
\noindent\textbf{diskriminantny ideal T vzhodno k o}\quad{\footnotesize\texttt{TERM-0327}}\par
\noindent\emph{Iztočnik / Source:} Diskriminantenideal von T in bezug auf o\par
\noindent\emph{Rusky / Russian:} дискриминантный идеал T относительно o\par
\noindent\emph{Ukrajinsky / Ukrainian:} дискримінантний ідеал T відносно o\par
\smallskip\noindent\emph{Razlog:} Diskriminantny ideal $T$ ... k $o$ zachovuje oba parametry konstrukcije i pridavnik diskriminantny. Zapis nazyvaje vzhodno k ustaljenym sposobom vyraziti ‘po odnošenju k’, ale opakovana upotrěba sama ne jest jezyčnym dokazom. Tuto tvrdženje o ranějšej praksě mora ostati odvojeno od prověrky, čto konstrukcija pravilno znači.\par
\noindent\emph{Reason:} Diskriminantny ideal $T$ ... k $o$ retains both parameters and the adjective diskriminantny. The record calls vzhodno k an established way to express with respect to, but repeated use is not linguistic evidence by itself. That account of earlier practice must remain separate from checking what the expression actually means.\par
\noindent{\footnotesize Izvorne zapisy / Original records: HIST-g155-005.}\par
\smallskip\noindent\emph{Sravnjenje s tekstom:} Aktualny tekst v dělu 31 ima po odnošenju k klasu i po odnošenju k osnovnomu prstenu. Kopija slovnika citovana v rěšenju o toj konstrukciji davaje oba zaglavja po odnošeńju k (18227) i v odnošeńju k (19113), s dativom i smyslom ‘po odnošenju k’. Zato argument, že samo prva forma jest oficialna, ne podpiraje izbora; konsistentna prednost jednoj formě može byti redakcijnoju konvencijeju. Osobny prilog vȯzhodno (4185) znači ‘na vȯzhod’, ne ‘po odnošenju k’. To ne opravdaje ranějše vzhodno k. Kratke odnosno (4107) ima rodni padež i zato ne jest mehanična zaměna frazy s dativom. Čitatelsko po odnošenju k slěduje osnovnomu pravopisu, kde ń prěhodi v n. Sravnjenje podpira značenje i padežno upravjanje, ne dokaz izměrenoj prednosti v razumělivosti.\par
\noindent\emph{Textual comparison:} The current Paper 31 text has po odnošenju k klasu and po odnošenju k with the base ring. The dictionary copy cited by the decision supplies both po odnošeńju k (18227) and v odnošeńju k (19113), each with the dative and the meaning with regard to. The claim that only the former is official therefore cannot justify the choice; consistently preferring one remains a possible editorial convention. The separate adverb vȯzhodno (4185) means eastwards, not with respect to, and does not justify historical vzhodno k. Shorter odnosno (4107) governs the genitive, so it is not a mechanical replacement for a dative construction. Reader spelling po odnošenju k follows the basic-orthography treatment of ń as n. The comparison supports meaning and government, not measured superiority in intelligibility.\par
\noindent{\footnotesize TEXT-COMP-06; pasaže / passages: 18094, 18787; \texttt{CORPUS\_FORM\_COMPARISONS.json}.}\par
\noindent{\footnotesize\href{https://interslavic-dictionary.com/}{Medžuslovjansky slovnik / Interslavic Dictionary}: 18227, 19113, 4107, 4185; točne kopije / exact snapshots in supporting data.}\par
\noindent{\footnotesize\href{https://interslavic.fun/learn/orthography/}{Pravopis / Orthography}: opcionalne bukvy / optional letters.}\par
\end{minipage}\par
\par\medskip\Needspace{7\baselineskip}
\noindent\begin{minipage}{\linewidth}
\hypertarget{TERM-0328}{}
\noindent\textbf{diskriminantne idealy primarnyh kolc}\quad{\footnotesize\texttt{TERM-0328}}\par
\noindent\emph{Iztočnik / Source:} Diskriminantenideale primärer Ringe\par
\noindent\emph{Rusky / Russian:} дискриминантные идеалы первичных колец\par
\noindent\emph{Ukrajinsky / Ukrainian:} дискримінантні ідеали первинних кілець\par
\smallskip\noindent\emph{Razlog:} Diskriminantne idealy primarnyh kolc zachovujut familiju diskriminantnyh idealov i kvalifikaciju primarnyh prstenov. Razlog daje prednost medžunarodnomu diskriminantny prěd opisnym prěvodom radi tehničnoj prepoznavajemosti. On ne objašnjaje novu etimologiju i ne dostavjaje izměrenogo sravnjenja razumělivosti.\par
\noindent\emph{Reason:} Diskriminantne idealy primarnyh kolc retains the discriminant-ideal family and the qualification of the rings as primary. The rationale favours international diskriminantny over a paraphrase for technical recognizability. It neither establishes a new etymology nor supplies a measured intelligibility comparison.\par
\noindent{\footnotesize Izvorne zapisy / Original records: HIST-g142-001.}\par
\end{minipage}\par
\par\medskip\Needspace{7\baselineskip}
\noindent\begin{minipage}{\linewidth}
\hypertarget{TERM-0329}{}
\noindent\textbf{diskriminantny kriterij polnoj razložimosti prvogo roda}\quad{\footnotesize\texttt{TERM-0329}}\par
\noindent\emph{Iztočnik / Source:} Diskriminantenkriterium der vollständigen Reduzibilität erster Art\par
\noindent\emph{Rusky / Russian:} дискриминантный критерий полной приводимости первого рода\par
\noindent\emph{Ukrajinsky / Ukrainian:} дискримінантний критерій цілковитої звідності першого роду\par
\smallskip\noindent\emph{Razlog:} Diskriminantny kriterij polnoj razložimosti prvogo roda zachovuje cělost kvalifikacij. Razložimost nastavlja termin iz odděla 3; ukrajinska paralela naměrno izbira zvìdnìstʹ vměsto zajma reducibelʹnistʹ, dokolě ruska ima privodimostʹ. To sravnjenje podpiraje značenje redukcije, ne dozvoljaje izpustiti polna ili prvogo roda.\par
\noindent\emph{Reason:} Diskriminantny kriterij polnoj razložimosti prvogo roda preserves the full set of qualifications. Razložimost continues section 3 terminology; the Ukrainian choice deliberately prefers звідність over borrowed редуцибельність, while Russian has приводимость. The comparison concerns the reduction sense, not permission to omit complete or of the first kind.\par
\noindent{\footnotesize Izvorne zapisy / Original records: HIST-g148-001.}\par
\end{minipage}\par
\par\medskip\Needspace{7\baselineskip}
\noindent\begin{minipage}{\linewidth}
\hypertarget{TERM-0330}{}
\noindent\textbf{diskriminantna matrica}\quad{\footnotesize\texttt{TERM-0330}}\par
\noindent\emph{Iztočnik / Source:} Diskriminantenmatrix\par
\noindent\emph{Rusky / Russian:} дискриминантная матрица\par
\noindent\emph{Ukrajinsky / Ukrainian:} дискримінантна матриця\par
\smallskip\noindent\emph{Razlog:} Diskriminantna matrica jest pridavnikova konstrukcija, paralelna ruskoj i ukrajinskoj. Matrica ostavaje imenovanjem objekta, ne zaměnjaje se slovom tablica; diskriminantna ukazyvaje na njegovu matematičnu funkciju. Izvorni razlog podpiraje kratko tehnično nazvanje bez novogo slovotvorstva.\par
\noindent\emph{Reason:} Diskriminantna matrica is an adjective–noun construction parallel to the Russian and Ukrainian forms. Matrica remains the object name rather than being replaced by table; diskriminantna identifies its mathematical role. The recorded reason supports a short technical name without a new coinage.\par
\noindent{\footnotesize Izvorne zapisy / Original records: HIST-g203-002.}\par
\end{minipage}\par
\par\medskip\Needspace{7\baselineskip}
\noindent\begin{minipage}{\linewidth}
\hypertarget{TERM-0331}{}
\noindent\textbf{teorem o diskriminantu}\quad{\footnotesize\texttt{TERM-0331}}\par
\noindent\emph{Iztočnik / Source:} Diskriminantensatz\par
\noindent\emph{Rusky / Russian:} теорема о дискриминанте\par
\noindent\emph{Ukrajinsky / Ukrainian:} дискримінантна теорема; теорема про дискримінант\par
\smallskip\noindent\emph{Razlog:} Teorem o diskriminantu razuměje Satz kako matematičny teorem, ne običnu větu. Konstrukcija teorem o nastavlja imenovanje iz uvoda raboty 31 i daje jasno ime rezultata. Ukrajinske zapisy imajo i pridavnikovu, i predložnu formu; točna paralela ne jest ve vseh zapisah jedna i ta sama.\par
\noindent\emph{Reason:} Teorem o diskriminantu reads Satz as a mathematical theorem, not an ordinary sentence. Teorem o continues the naming in Paper 31's introduction and makes the result's title clear. The Ukrainian records contain both adjectival and prepositional forms, so the exact comparison wording is not uniform across records.\par
\noindent{\footnotesize Izvorne zapisy / Original records: HIST-g118-001, HIST-g153-001.}\par
\end{minipage}\par
\par\medskip\Needspace{7\baselineskip}
\noindent\begin{minipage}{\linewidth}
\hypertarget{TERM-0332}{}
\noindent\textbf{jest ispolnjen diskriminantny teorem (§7, 3)}\quad{\footnotesize\texttt{TERM-0332}}\par
\noindent\emph{Iztočnik / Source:} Diskriminantensatz (§7, 3) erfüllt\par
\noindent\emph{Rusky / Russian:} выполнена дискриминантная теорема (§7, 3)\par
\noindent\emph{Ukrajinsky / Ukrainian:} виконується дискримінантна теорема (§7, 3)\par
\smallskip\noindent\emph{Razlog:} Jest ispolnjen diskriminantny teorem (§7, 3) zachovuje izrično pozvanje na opreděljany rezultat. Smysl jest upotrěba ranějšego teorema, ne njegovo novo prěformulovanje. Pridavnikova forma diskriminantny teorem jest historijna alternativa konstrukciji teorem o diskriminantu; sam zapis ne razrěšaje jejich stilistično sravnjenje.\par
\noindent\emph{Reason:} Jest ispolnjen diskriminantny teorem (§7, 3) preserves the explicit reference to a particular result. The point is to invoke an earlier theorem, not restate it. Adjectival diskriminantny teorem is a historical alternative to teorem o diskriminantu; this record alone does not settle their stylistic comparison.\par
\noindent{\footnotesize Izvorne zapisy / Original records: HIST-g161-007.}\par
\end{minipage}\par
\par\medskip\Needspace{7\baselineskip}
\noindent\begin{minipage}{\linewidth}
\hypertarget{TERM-0333}{}
\noindent\textbf{divergencija / divergencijne relacije}\quad{\footnotesize\texttt{TERM-0333}}\par
\noindent\emph{Iztočnik / Source:} Divergenz / Divergenzrelationen\par
\noindent\emph{Rusky / Russian:} дивергенция / дивергентные соотношения\par
\noindent\emph{Ukrajinsky / Ukrainian:} дивергенція / дивергентні співвідношення\par
\smallskip\noindent\emph{Razlog:} Divergencija i divergencijne relacije zachovujut organizaciju dokaza okolo divergencijnyh identičnostij. Relacija imenuje odnos na urovnju formuly; divergencija ostavaje medžunarodno ime operacije. Izbor ne pretvarjaje každu relaciju v odvojenu novu definiciju divergencije.\par
\noindent\emph{Reason:} Divergencija and divergencijne relacije preserve the proof's organization around divergence identities. Relacija names a formula-level relation, while divergencija remains the international operation name. The choice does not turn each relation into a separate new definition of divergence.\par
\noindent{\footnotesize Izvorne zapisy / Original records: HIST-g075-009.}\par
\end{minipage}\par
\par\medskip\Needspace{7\baselineskip}
\noindent\begin{minipage}{\linewidth}
\hypertarget{TERM-0334}{}
\noindent\textbf{d?litvena algebra; prir??ena d?litvena algebra}\quad{\footnotesize\texttt{TERM-0334}}\par
\noindent\emph{Iztočnik / Source:} Divisionsalgebra; zugeordnete Divisionsalgebra\par
\noindent\emph{Rusky / Russian:} [Bukvy ne sut zachovane / Lettering not recoverable]\par
\noindent\emph{Ukrajinsky / Ukrainian:} [Bukvy ne sut zachovane / Lettering not recoverable]\par
\smallskip\noindent\emph{Razlog:} Zapis naměravaje pridavnik od děljenja i konstrukciju označajuču pridruženu algebru, v kontinuitetu terminologije centralno prostyh algebr. Česti medžuslovjanskyh slov i cěle ruske i ukrajinske formy sut povrědžene. Razlog možno izložiti, ale iz samogo zapisa ne možno pošteno vosstanoviti točne nedostajuče bukvy ili proglasiti kandidat normalnoju slovničnoju formoju.\par
\noindent\emph{Reason:} The record intends an adjective derived from division and an expression for the associated algebra, continuing central-simple-algebra terminology. Parts of the Interslavic spellings and the entire Russian and Ukrainian forms are damaged. The rationale can be explained, but the record alone cannot justify reconstructing the missing letters or presenting the candidate as a regular dictionary form.\par
\noindent{\footnotesize Izvorne zapisy / Original records: HIST-g210-008.}\par
\end{minipage}\par
\par\medskip\Needspace{7\baselineskip}
\noindent\begin{minipage}{\linewidth}
\hypertarget{TERM-0335}{}
\noindent\textbf{dvojna oblast $(\Sigma,T)$}\quad{\footnotesize\texttt{TERM-0335}}\par
\noindent\emph{Iztočnik / Source:} Doppelbereich $(\Sigma,T)$\par
\noindent\emph{Rusky / Russian:} двойная область $(\Sigma,T)$\par
\noindent\emph{Ukrajinsky / Ukrainian:} подвійна область $(\Sigma,T)$\par
\smallskip\noindent\emph{Razlog:} Dvojna oblast $(\Sigma,T)$ nastavlja oblast za Bereich v rabotě 19. Dvojna razkriva složeno Doppelbereich bez vvedenja konkurentnogo domena. Par oznak ostavaje čestju identifikacije; sam nazvanny smysl oblasti ne ustanavjaje sovrěmennu definiciju integralnogo prstena.\par
\noindent\emph{Reason:} Dvojna oblast $(\Sigma,T)$ continues oblast for Bereich in Paper 19. Dvojna unpacks Doppelbereich without introducing competing domena. The pair of symbols remains part of the identification; calling the object a domain does not itself supply the modern definition of an integral domain.\par
\noindent{\footnotesize Izvorne zapisy / Original records: HIST-g081-069.}\par
\end{minipage}\par
\par\medskip\Needspace{7\baselineskip}
\noindent\begin{minipage}{\linewidth}
\hypertarget{TERM-0336}{}
\noindent\textbf{dvojny indeks}\quad{\footnotesize\texttt{TERM-0336}}\par
\noindent\emph{Iztočnik / Source:} Doppelindex\par
\noindent\emph{Rusky / Russian:} двойной индекс\par
\noindent\emph{Ukrajinsky / Ukrainian:} подвійний індекс\par
\smallskip\noindent\emph{Razlog:} Dvojny indeks jest kratko nazvanje notacijnogo objekta. Zapis takojže dava instrumental dvojnym indeksom za upotrěbu v proze. Ta forma pokazuje, kako lemma vhodži v konkretnu větu; ona ne predstavlja odvojene tvorjenje termina.\par
\noindent\emph{Reason:} Dvojny indeks is a compact notation term. The record also supplies instrumental dvojnym indeksom for running prose. That inflection shows how the headword enters a sentence; it is not a separate terminological coinage.\par
\noindent{\footnotesize Izvorne zapisy / Original records: HIST-g041-012.}\par
\end{minipage}\par
\par\medskip\Needspace{7\baselineskip}
\noindent\begin{minipage}{\linewidth}
\hypertarget{TERM-0337}{}
\noindent\textbf{dvojno uslovje cěpov}\quad{\footnotesize\texttt{TERM-0337}}\par
\noindent\emph{Iztočnik / Source:} Doppelkettensatz\par
\noindent\emph{Rusky / Russian:} двойное условие цепей\par
\noindent\emph{Ukrajinsky / Ukrainian:} подвійна умова ланцюгів\par
\smallskip\noindent\emph{Razlog:} Dvojno uslovje cěpov obično v těch zapisah znači sovměstno uslovje cěpov děliteljev i uslovje cěpov kratnyh. Tri razloga davajut prednost jezyku uslovij vměsto novogo nazvanja teorema. Četvrty govori o ekvivalentnosti dvuh napravljenij cěpov; ta formulacija ne slěduje iz ostalyh razlogov i trěbuje sravnjenje s samym argumentom. Ne možno tiho svesti vse četyri zapisy k jednoj bezprotivorěčnoj motivaciji.\par
\noindent\emph{Reason:} Dvojno uslovje cěpov ordinarily denotes the divisor-chain and multiple-chain conditions together in these records. Three rationales favour condition language over a new theorem name. A fourth speaks of equivalence between the two chain directions; that wording does not follow from the other reasons and requires comparison with the argument. The four records must not be silently harmonized into one unproblematic rationale.\par
\noindent{\footnotesize Izvorne zapisy / Original records: HIST-g107-004, HIST-g114-006, HIST-g117-001, HIST-g177-005.}\par
\smallskip\noindent\emph{Sravnjenje s tekstom:} V dělu 30, oddělu 10, oba směry tvrdženja sut izrične: dvojno uslovje cěpov davaje kompozicijny red, a kompozicijny red davaje dvojno uslovje. Definicija dvojnego uslovja požaduje, aby vsaka cěp děliteljev i vsaka cěp kratnyh končila se. Zato zabilěženo objasnjenje o ekvivalentnosti narastajučego i opadajučego uslovja jest netočno: ekvivalentnost jest medžu njihovym sovměstnym izpolnjenjem i kompozicijnym redom, pri prědpoloženjah pasaže, vključajuči dobro uporjadkovanje. Tuto sravnjenje opravdaje razlikovanje uslovja od nazvanja teorema, ale ne proglašaje razne upotrěby slova Satz vzajemno zaměnnymi.\par
\noindent\emph{Textual comparison:} Paper 30, section 10 states both directions explicitly: the double-chain condition gives a composition series, and a composition series gives the double-chain condition. The definition requires termination of every divisor chain and every multiple chain. The recorded explanation about equivalence of the ascending and descending conditions is therefore inaccurate: the equivalence is between their joint satisfaction and a composition series, under the passage's stated assumptions, including well-ordering. This comparison supports distinguishing a condition from a theorem name; it does not make every use of Satz interchangeable.\par
\noindent{\footnotesize TEXT-COMP-05; pasaže / passages: 17189, 17230; \texttt{CORPUS\_FORM\_COMPARISONS.json}.}\par
\end{minipage}\par
\par\medskip\Needspace{7\baselineskip}
\noindent\begin{minipage}{\linewidth}
\hypertarget{TERM-0338}{}
\noindent\textbf{dvojno uslovje konečnoj cěpi}\quad{\footnotesize\texttt{TERM-0338}}\par
\noindent\emph{Iztočnik / Source:} Doppelkettensatz\par
\noindent\emph{Rusky / Russian:} условие двойной цепи\par
\noindent\emph{Ukrajinsky / Ukrainian:} умова подвійного ланцюга\par
\smallskip\noindent\emph{Razlog:} Dvojno uslovje konečnoj cěpi prědava Satz kako uslovje, ibo věta harakterizuje prsteny črěz prědpoloženja. Cěp nastavlja izbrano ime cěpi idealov. Konečnost odnašaje se k prěkidanju cěpi pri danom uslovju; samo nazvanje ne zaměnjaje točne kvantifikatory definicije.\par
\noindent\emph{Reason:} Dvojno uslovje konečnoj cěpi renders Satz as a condition because the sentence characterizes rings by hypotheses. Cěp continues the chosen name for a chain of ideals. Finiteness concerns termination under the stated condition; the name cannot replace the definition's exact quantifiers.\par
\noindent{\footnotesize Izvorne zapisy / Original records: HIST-g103-004.}\par
\end{minipage}\par
\par\medskip\Needspace{7\baselineskip}
\noindent\begin{minipage}{\linewidth}
\hypertarget{TERM-0339}{}
\noindent\textbf{teorem o dvojnoj cěpi}\quad{\footnotesize\texttt{TERM-0339}}\par
\noindent\emph{Iztočnik / Source:} Doppelkettensatz\par
\noindent\emph{Rusky / Russian:} теорема о двойной цепи\par
\noindent\emph{Ukrajinsky / Ukrainian:} теорема про подвійний ланцюг\par
\smallskip\noindent\emph{Razlog:} Teorem o dvojnoj cěpi jest naměrno bližši doslovnomu německomu nazvanju než dvojno uslovje cěpov. Razlog zachovuje vezu so sovměstnymi uslovjami děliteljnyh i mnogokratnyh cěpov. To jest dokumentovana druga interpretacija formy Satz, ne dokaz, že vsi ranějše zapisy imajut jednu jezyčnu politiku.\par
\noindent\emph{Reason:} Teorem o dvojnoj cěpi is deliberately closer to the German title form than dvojno uslovje cěpov. Its rationale retains the connection with simultaneous divisor-chain and multiple-chain conditions. This is a documented alternative treatment of Satz, not evidence that all earlier entries followed one wording policy.\par
\noindent{\footnotesize Izvorne zapisy / Original records: HIST-g121-004.}\par
\end{minipage}\par
\par\medskip\Needspace{7\baselineskip}
\noindent\begin{minipage}{\linewidth}
\hypertarget{TERM-0340}{}
\noindent\textbf{dvojno uslovje cěpov / uslovje cěpov děliteljev / uslovje cěpov kratnyh}\quad{\footnotesize\texttt{TERM-0340}}\par
\noindent\emph{Iztočnik / Source:} Doppelkettensatz / Teilerkettensatz / Vielfachenkettensatz\par
\noindent\emph{Rusky / Russian:} двойное условие цепей / условие цепей делителей / условие цепей кратных\par
\noindent\emph{Ukrajinsky / Ukrainian:} подвійна умова ланцюгів / умова ланцюгів дільників / умова ланцюгів кратних\par
\smallskip\noindent\emph{Razlog:} Trojično imenovanje zachovuje cělost i obe sostavne česti: dvojno uslovje cěpov znači obě uslovja zajedno, a ostale dve frazy razlučajut dělitelje i kratne. Razlog nastavlja terminologiju uvoda raboty 30. Ne utvrždaje, že možno jednu polovinu izpustiti kakoby druga ju avtomatično vlekla.\par
\noindent\emph{Reason:} The three names preserve the whole and its two components: the double-chain condition means both conditions together, while the other phrases distinguish divisors from multiples. The rationale continues Paper 30's introduction. It does not establish that one component may be dropped because the other automatically entails it.\par
\noindent{\footnotesize Izvorne zapisy / Original records: HIST-g114-001.}\par
\smallskip\noindent\emph{Sravnjenje s tekstom:} V dělu 30, oddělu 10, oba směry tvrdženja sut izrične: dvojno uslovje cěpov davaje kompozicijny red, a kompozicijny red davaje dvojno uslovje. Definicija dvojnego uslovja požaduje, aby vsaka cěp děliteljev i vsaka cěp kratnyh končila se. Zato zabilěženo objasnjenje o ekvivalentnosti narastajučego i opadajučego uslovja jest netočno: ekvivalentnost jest medžu njihovym sovměstnym izpolnjenjem i kompozicijnym redom, pri prědpoloženjah pasaže, vključajuči dobro uporjadkovanje. Tuto sravnjenje opravdaje razlikovanje uslovja od nazvanja teorema, ale ne proglašaje razne upotrěby slova Satz vzajemno zaměnnymi.\par
\noindent\emph{Textual comparison:} Paper 30, section 10 states both directions explicitly: the double-chain condition gives a composition series, and a composition series gives the double-chain condition. The definition requires termination of every divisor chain and every multiple chain. The recorded explanation about equivalence of the ascending and descending conditions is therefore inaccurate: the equivalence is between their joint satisfaction and a composition series, under the passage's stated assumptions, including well-ordering. This comparison supports distinguishing a condition from a theorem name; it does not make every use of Satz interchangeable.\par
\noindent{\footnotesize TEXT-COMP-05; pasaže / passages: 17189, 17230; \texttt{CORPUS\_FORM\_COMPARISONS.json}.}\par
\end{minipage}\par
\par\medskip\Needspace{7\baselineskip}
\noindent\begin{minipage}{\linewidth}
\hypertarget{TERM-0341}{}
\noindent\textbf{bimodul}\quad{\footnotesize\texttt{TERM-0341}}\par
\noindent\emph{Iztočnik / Source:} Doppelmodul\par
\noindent\emph{Rusky / Russian:} бимодуль\par
\noindent\emph{Ukrajinsky / Ukrainian:} бімодуль\par
\smallskip\noindent\emph{Razlog:} Bimodul jest medžunarodno nazvanje, izbrano vměsto doslovnogo dvojny modul radi kratkosti i tehničnoj prepoznavajemosti. Dva razloga izrično sravnjavajut ty dve formy. Ih izbor razlikuje se od ranějšego dvomodul i od definicijnogo dvojny modul; ne možno pripisati vsim zapisam jednu prednost bez dalšego sravnjenja.\par
\noindent\emph{Reason:} Bimodul is the international name chosen over literal dvojny modul for brevity and technical recognizability. Two rationales explicitly compare those alternatives. Their choice differs from earlier dvomodul and from dvojny modul in the defining passage; the records cannot all be assigned one preference without further comparison.\par
\noindent{\footnotesize Izvorne zapisy / Original records: HIST-g196-008, HIST-g197-006.}\par
\end{minipage}\par
\par\medskip\Needspace{7\baselineskip}
\noindent\begin{minipage}{\linewidth}
\hypertarget{TERM-0342}{}
\noindent\textbf{dvojny modul}\quad{\footnotesize\texttt{TERM-0342}}\par
\noindent\emph{Iztočnik / Source:} Doppelmodul\par
\noindent\emph{Rusky / Russian:} бимодуль\par
\noindent\emph{Ukrajinsky / Ukrainian:} бімодуль\par
\smallskip\noindent\emph{Razlog:} Dvojny modul jest naměrny doslovny prěvod v definicijnoj pasažě. Razlog dava prednost vidimoj dvustrannoj operatornoj strukturě, dokolě ruska i ukrajinska paralela imaju bimodul. Tuto objasnjuje lokalny izbor, ale ne dokazuje, že doslovna konstrukcija dolžna zaměniti medžunarodno nazvanje v každom kontekstu.\par
\noindent\emph{Reason:} Dvojny modul is a deliberate calque in the defining passage. The rationale favours making the two-sided operator structure visible, while the Russian and Ukrainian comparanda have bimodule. That explains a local choice but does not establish that the calque should replace the international name in every context.\par
\noindent{\footnotesize Izvorne zapisy / Original records: HIST-g193-003.}\par
\end{minipage}\par
\par\medskip\Needspace{7\baselineskip}
\noindent\begin{minipage}{\linewidth}
\hypertarget{TERM-0343}{}
\noindent\textbf{dvojny modul odnosno bimodul}\quad{\footnotesize\texttt{TERM-0343}}\par
\noindent\emph{Iztočnik / Source:} Doppelmodul\par
\noindent\emph{Rusky / Russian:} [Bukvy ne sut zachovane / Lettering not recoverable]\par
\noindent\emph{Ukrajinsky / Ukrainian:} [Bukvy ne sut zachovane / Lettering not recoverable]\par
\smallskip\noindent\emph{Razlog:} Dvojny modul odnosno bimodul naměravaje pri prvom pojavjenju spojiti historijno doslovno nazvanje so sovrěmennym terminom, a potom dopustiti kratku formu. To jest funkcija dvuh imen v jednoj frazě, ne dva različne objekty. Ruske i ukrajinske bukvy v zachovanyh sravniteljnyh poljah sut nečitateljne i ne mogut byti vydavane za točne citaty.\par
\noindent\emph{Reason:} Dvojny modul odnosno bimodul is intended to connect the historical literal name with the modern term at first occurrence, allowing a shorter form thereafter. The two names serve that explanatory purpose, not two different objects. The preserved Russian and Ukrainian lettering is unreadable and cannot be presented as exact quotations.\par
\noindent{\footnotesize Izvorne zapisy / Original records: HIST-g210-002.}\par
\end{minipage}\par
\par\medskip\Needspace{7\baselineskip}
\noindent\begin{minipage}{\linewidth}
\hypertarget{TERM-0344}{}
\noindent\textbf{dvomodul}\quad{\footnotesize\texttt{TERM-0344}}\par
\noindent\emph{Iztočnik / Source:} Doppelmodul\par
\noindent\emph{Rusky / Russian:} бимодуль\par
\noindent\emph{Ukrajinsky / Ukrainian:} бімодуль\par
\smallskip\noindent\emph{Razlog:} Dvomodul vyražaje ‘dva’ i ‘modul’ v jednom slově. Tri zapisy davajut prednost toj morfologičnoj prozračnosti i nastavljajut termin iz odděla 2. Jedin iz njih nazyvaje medžunarodnu alternativu neprozračnoju, kde drugy zapisy hvalet kratkost bimodul. To sut različne ocěny avtorov zapisov, ne izměreny rezultaty sravnjenja čitateljev.\par
\noindent\emph{Reason:} Dvomodul expresses two and module in one word. Three records favour that morphological transparency and continue section 2 terminology. One calls the international alternative opaque, whereas other records favour bimodul for brevity. These are differing judgments in the records, not measured results from a reader comparison.\par
\noindent{\footnotesize Izvorne zapisy / Original records: HIST-g179-007, HIST-g180-007, HIST-g186-004.}\par
\smallskip\noindent\emph{Sravnjenje s tekstom:} Definicija v dělu 34, oddělu 1, imenuje dvomodul objekt s lěvym i pravym modulnym dějstvovanjem, svazanymi ravnostju $a\cdot\alpha a'=a\alpha\cdot a'$. Samo čislo dva v nazvanju ne zaměnjaje to uslovje soglasnosti. Pozdnějša pasaža v oddělu 15 ima dvojny modul s izričnym odkazom na tu definiciju; druga pasaža o tělahi avtomorfizmov ima bimoduly. Ty priměry pokazyvajut različne zapisy za svazany strukturny pojem. Ne dokazujut, že različie jest naměrna universalna politika ili že vsako pojavjenje ima isti smysl.\par
\noindent\emph{Textual comparison:} The definition in Paper 34, section 1 calls an object a dvomodul when its left and right module actions satisfy $a\cdot\alpha a'=a\alpha\cdot a'$. Two in the name cannot replace that compatibility condition. The later section 15 passage uses dvojny modul with an explicit reference to this definition; another passage concerning automorphism division rings uses bimoduly. These examples demonstrate different names for the connected structural concept. They do not establish a deliberate universal policy or prove that every occurrence has the same sense.\par
\noindent{\footnotesize TEXT-COMP-07; pasaže / passages: 19840, 19842, 19844, 21339, 21685; \texttt{CORPUS\_FORM\_COMPARISONS.json}.}\par
\end{minipage}\par
\par\medskip\Needspace{7\baselineskip}
\noindent\begin{minipage}{\linewidth}
\hypertarget{TERM-0345}{}
\noindent\textbf{ideal dvojnyh toček}\quad{\footnotesize\texttt{TERM-0345}}\par
\noindent\emph{Iztočnik / Source:} Doppelpunktsideal\par
\noindent\emph{Rusky / Russian:} идеал двойных точек\par
\noindent\emph{Ukrajinsky / Ukrainian:} ідеал подвійних точок\par
\smallskip\noindent\emph{Razlog:} Ideal dvojnyh toček zachovuje doslovno nazvanje geometrickogo objekta. V pregledu on jest identifikovany s konduktorom, ale to ne daje razloga tiho zaměniti avtorske ime sovrěmennym idealom singularnostij. Izvorni razlog zato zachovuje istorijno imenovanje i ego objašnjenu vezu s drugym objektom.\par
\noindent\emph{Reason:} Ideal dvojnyh toček retains the literal geometric name. The survey identifies it with the conductor, but that does not warrant silently replacing the author's name with a modern singularity ideal. The rationale preserves the historical naming and its stated connection with another object.\par
\noindent{\footnotesize Izvorne zapisy / Original records: HIST-g076-019.}\par
\end{minipage}\par
\par\medskip\Needspace{7\baselineskip}
\noindent\begin{minipage}{\linewidth}
\hypertarget{TERM-0346}{}
\noindent\textbf{dvojna redukcija}\quad{\footnotesize\texttt{TERM-0346}}\par
\noindent\emph{Iztočnik / Source:} Doppelte Reduktion\par
\noindent\emph{Rusky / Russian:} двойная редукция\par
\noindent\emph{Ukrajinsky / Ukrainian:} подвійна редукція\par
\smallskip\noindent\emph{Razlog:} Dvojna redukcija jest imenovana tehnična operacija, ne prosto dvojno uproščenje. Vosem zapisov pokazyvajut kontinuitet od ranih oddělov raboty 2 do oddělov 20–25; kavyčky zachovujut funkciju nazvanja metody. Jedin zapis izrično govori o računskoj prověrkě, a drugy svazuje formulu (31) s kombinacijeju redukcij 1 i 5. Pridavnik dvojna proto zachovuje relevantno čislo operacij, a imennik redukcija konkretny vid operacije.\par
\noindent\emph{Reason:} Dvojna redukcija is a named technical operation, not merely a double simplification. Eight records show continuity from the early sections of Paper 2 through sections 20–25; quotation marks preserve its role as a method name. One record explicitly identifies a computational check, and another connects formula (31) with combining reductions 1 and 5. Dvojna therefore retains the relevant operation count, while redukcija identifies the particular operation.\par
\noindent{\footnotesize Izvorne zapisy / Original records: HIST-g005-005, HIST-g008-009, HIST-g013-006, HIST-g021-008, HIST-g022-005, HIST-g023-008, HIST-g024-007, HIST-g026-007.}\par
\end{minipage}\par
\par\medskip\Needspace{7\baselineskip}
\noindent\begin{minipage}{\linewidth}
\hypertarget{TERM-0347}{}
\noindent\textbf{tridimensionalny prostor}\quad{\footnotesize\texttt{TERM-0347}}\par
\noindent\emph{Iztočnik / Source:} dreidimensionaler Raum\par
\noindent\emph{Rusky / Russian:} трёхмерное пространство\par
\noindent\emph{Ukrajinsky / Ukrainian:} тривимірний простір\par
\smallskip\noindent\emph{Razlog:} Tridimensionalny prostor zachovuje točno tri dimenzije v sovrěmennom matematičnom registru. Izvorni razlog ne vnašaje drugu staršu terminologiju linearnogo prostora, ibo iztočnik ovde neposrědno imenuje tridimenzionalny prostor. Pridavnik jest čestju danogo objekta, ne stilistično ukrašenje.\par
\noindent\emph{Reason:} Tridimensionalny prostor retains exactly three dimensions in a modern mathematical register. The rationale does not introduce a different, older linear-space terminology because the source directly names three-dimensional space here. The adjective belongs to the specified object, not to stylistic decoration.\par
\noindent{\footnotesize Izvorne zapisy / Original records: HIST-g046-013.}\par
\end{minipage}\par
\par\medskip\Needspace{7\baselineskip}
\noindent\begin{minipage}{\linewidth}
\hypertarget{TERM-0348}{}
\noindent\textbf{dualistično protivopoložne}\quad{\footnotesize\texttt{TERM-0348}}\par
\noindent\emph{Iztočnik / Source:} dualistisch entgegengesetzt\par
\noindent\emph{Rusky / Russian:} дуалистически противоположные\par
\noindent\emph{Ukrajinsky / Ukrainian:} дуалістично протилежні\par
\smallskip\noindent\emph{Razlog:} Dualistično protivopoložne nastavlja jezyk dualnosti iz odděla 12, ale zapis izrično priznava stilističnu těžkost. On upominaje krače možnosti dualne i protivne za sravnjenje. Te alternativy sut izrično zapisane; nijedna ne jest tamo podprěta točnym citatom, ktory by sam razrěšil izbor.\par
\noindent\emph{Reason:} Dualistično protivopoložne continues section 12's duality language, but the record explicitly acknowledges stylistic heaviness. It proposes shorter dualne and protivne for comparison. Those alternatives are genuinely recorded, although neither is supported there by a precise citation that would settle the choice.\par
\noindent{\footnotesize Izvorne zapisy / Original records: HIST-g014-015.}\par
\end{minipage}\par
\par\medskip\Needspace{7\baselineskip}
\noindent\begin{minipage}{\linewidth}
\hypertarget{TERM-0349}{}
\noindent\textbf{dualistično protivopoložny}\quad{\footnotesize\texttt{TERM-0349}}\par
\noindent\emph{Iztočnik / Source:} dualistisch entgegengesetzt\par
\noindent\emph{Rusky / Russian:} дуалистически противоположным\par
\noindent\emph{Ukrajinsky / Ukrainian:} дуалістично протилежним\par
\smallskip\noindent\emph{Razlog:} Dualistično protivopoložny zachovuje Noetherino izrično imenovanje dualnosti, vměsto jego odstranjenja prěformulovanjem. Medžunarodno dualistično i slovjansko protivopoložny imajut različne funkcije v opisu. Izvorni razlog opravdaje zachovanje oběh, ale ne dokazuje, že dolga konstrukcija jest najprirodnějša možnost.\par
\noindent\emph{Reason:} Dualistično protivopoložny keeps Noether's explicit duality language instead of paraphrasing it away. International dualistično and Slavic protivopoložny contribute different parts of the description. The rationale supports retaining both functions but does not establish the long construction as the most idiomatic option.\par
\noindent{\footnotesize Izvorne zapisy / Original records: HIST-g011-012.}\par
\end{minipage}\par
\par\medskip\Needspace{7\baselineskip}
\noindent\begin{minipage}{\linewidth}
\hypertarget{TERM-0350}{}
\noindent\textbf{dualistične identičnosti}\quad{\footnotesize\texttt{TERM-0350}}\par
\noindent\emph{Iztočnik / Source:} dualistische Identitäten\par
\noindent\emph{Rusky / Russian:} дуалистические тождества\par
\noindent\emph{Ukrajinsky / Ukrainian:} дуалістичні тотожності\par
\smallskip\noindent\emph{Razlog:} Dualistične identičnosti zachovujut specifičny dualističny jezyk iztočnika, ne zaměnjajuči jego občeju frazoju dualne identičnosti. Razlog daje prednost istorijnoj terminologičnoj točnosti. On ne ustanavjaje novy obči smysl slova identičnost za vsaku upotrěbu v sovrěmennoj matematikě.\par
\noindent\emph{Reason:} Dualistične identičnosti retains the source's specifically dualistic terminology rather than replacing it with generic dual identities. The reason favours historical terminological fidelity. It does not establish a new general sense of identity for every use in modern mathematics.\par
\noindent{\footnotesize Izvorne zapisy / Original records: HIST-g032-011.}\par
\end{minipage}\par
\par\medskip\Needspace{7\baselineskip}
\noindent\begin{minipage}{\linewidth}
\hypertarget{TERM-0351}{}
\noindent\textbf{dualistične formy}\quad{\footnotesize\texttt{TERM-0351}}\par
\noindent\emph{Iztočnik / Source:} dualistischen Formen\par
\noindent\emph{Rusky / Russian:} дуалистические формы\par
\noindent\emph{Ukrajinsky / Ukrainian:} дуалістичні форми\par
\smallskip\noindent\emph{Razlog:} Dualistične formy zachovujut klasicky jezyk teorije invariantov. Medžunarodna osnova dualističn- prodolžaje točno imenovanje iztočnika, ne prosto obče prědloženje o dvoh formah. Kratky izvorni razlog ne dava zasebno sravnjenje krače alternativy dualne.\par
\noindent\emph{Reason:} Dualistične formy preserves the classical invariant-theoretic register. International dualističn- continues the source's particular naming rather than merely saying that two forms are involved. The brief recorded reason supplies no separate comparison with shorter dualne.\par
\noindent{\footnotesize Izvorne zapisy / Original records: HIST-g005-009.}\par
\end{minipage}\par
\par\medskip\Needspace{7\baselineskip}
\noindent\begin{minipage}{\linewidth}
\hypertarget{TERM-0352}{}
\noindent\textbf{dualistična razkladna identičnost}\quad{\footnotesize\texttt{TERM-0352}}\par
\noindent\emph{Iztočnik / Source:} dualististische Zerlegungsidentität\par
\noindent\emph{Rusky / Russian:} дуалистическое тождество разложения\par
\noindent\emph{Ukrajinsky / Ukrainian:} дуалістична тотожність розкладання\par
\smallskip\noindent\emph{Razlog:} Dualistična razkladna identičnost nastavlja dualističnu identičnost iz odděla 3 i dodaje specifičnu funkciju razklada. Pridavnik razkladna ne jest svobodno objasnjenje: on odlišaje tutu identičnost od ostalyh dualističnyh identičnostij. Neobyčna grafija německogo zaglavja ostavaje v izvornih podatkah, ne tihi popravok citata.\par
\noindent\emph{Reason:} Dualistična razkladna identičnost continues section 3's dualistic-identity terminology and adds the specific decomposition function. Razkladna distinguishes this identity from other dualistic identities; it is not a free explanatory addition. The unusual spelling of the recorded German headword remains in the source data rather than being silently corrected as a quotation.\par
\noindent{\footnotesize Izvorne zapisy / Original records: HIST-g033-005.}\par
\end{minipage}\par
\par\medskip\Needspace{7\baselineskip}
\noindent\begin{minipage}{\linewidth}
\hypertarget{TERM-0353}{}
\noindent\textbf{dělimy črez $(x_i-a)$}\quad{\footnotesize\texttt{TERM-0353}}\par
\noindent\emph{Iztočnik / Source:} durch $(x_i-a)$ teilbar\par
\noindent\emph{Rusky / Russian:} делится на $(x_i-a)$\par
\noindent\emph{Ukrajinsky / Ukrainian:} ділиться на $(x_i-a)$\par
\smallskip\noindent\emph{Razlog:} Dělimy črez $(x_i-a)$ imenuje točny algebraičny odnos, na ktorom zasnovano jest protivorěčje. Izvorni razlog odbacaje neoprěděljeno ‘ima faktor’, ibo takva parafraza može skryti, ktory izraz dělji ktory. Indeks i oduzimanje v faktoru ostavajut neproměnjene.\par
\noindent\emph{Reason:} Dělimy črez $(x_i-a)$ names the precise algebraic relation used in the contradiction. The recorded reason rejects a vague has a factor because it can obscure which expression divides which. The subscript and subtraction in the factor remain unchanged.\par
\noindent{\footnotesize Izvorne zapisy / Original records: HIST-g042-007.}\par
\end{minipage}\par
\par\medskip\Needspace{7\baselineskip}
\noindent\begin{minipage}{\linewidth}
\hypertarget{TERM-0354}{}
\noindent\textbf{porodžuje se nad P črez adjunkciju jednogo/dvuh elementov}\quad{\footnotesize\texttt{TERM-0354}}\par
\noindent\emph{Iztočnik / Source:} durch Adjunktion eines/zwei Elementes/Elementen zu P erzeugt\par
\noindent\emph{Rusky / Russian:} порождается над P присоединением одного/двух элементов\par
\noindent\emph{Ukrajinsky / Ukrainian:} породжується над P приєднанням одного/двох елементів\par
\smallskip\noindent\emph{Razlog:} Porodžuje se nad $P$ črez adjunkciju jednogo ili dvuh elementov zachovuje osnovno polje, operaciju adjunkcije i čislo dodanyh elementov. Glagol nastavlja ranějše porodženy. Razlika medžu jednym i dvoma ne može byti skryta občeju frazoju ‘pridavanjem elementov’, ibo ona jest čestju matematičnogo sravnjenja.\par
\noindent\emph{Reason:} Porodžuje se nad $P$ črez adjunkciju jednogo or dvuh elementov retains the base field, adjunction, and the number of adjoined elements. The verb continues earlier porodženy. One versus two cannot be replaced by the indefinite adjoining elements, because the count belongs to the mathematical comparison.\par
\noindent{\footnotesize Izvorne zapisy / Original records: HIST-g153-008.}\par
\end{minipage}\par
\par\medskip\Needspace{7\baselineskip}
\noindent\begin{minipage}{\linewidth}
\hypertarget{TERM-0355}{}
\noindent\textbf{jest dělimy kvadratom prostogo ideala}\quad{\footnotesize\texttt{TERM-0355}}\par
\noindent\emph{Iztočnik / Source:} durch das Quadrat eines Primideals teilbar\par
\noindent\emph{Rusky / Russian:} делится на квадрат простого идеала\par
\noindent\emph{Ukrajinsky / Ukrainian:} ділиться на квадрат простого ідеалу\par
\smallskip\noindent\emph{Razlog:} Jest dělimy kvadratom prostogo ideala nastavlja formulaciju o děliteljnosti na kvadrat iz odděla 7. Kvadrat i prosty ideal ostavajut izrične, da by specialny slučaj ostal svazany s kriterijem vlastno primarnoj komponenty. Prosta děliteljnost idealom ne vyražaje tuto silnějše uslovje.\par
\noindent\emph{Reason:} Jest dělimy kvadratom prostogo ideala continues section 7's square-divisibility formulation. Square and prime ideal remain explicit so that the specialization stays connected to the properly primary component criterion. Divisibility by the ideal alone would express a weaker condition.\par
\noindent{\footnotesize Izvorne zapisy / Original records: HIST-g164-010.}\par
\end{minipage}\par
\par\medskip\Needspace{7\baselineskip}
\noindent\begin{minipage}{\linewidth}
\hypertarget{TERM-0356}{}
\noindent\textbf{črez homomorfizm diskriminant prěhodi}\quad{\footnotesize\texttt{TERM-0356}}\par
\noindent\emph{Iztočnik / Source:} Durch den Homomorphismus geht die Diskriminante ... über\par
\noindent\emph{Rusky / Russian:} при гомоморфизме дискриминант переходит\par
\noindent\emph{Ukrajinsky / Ukrainian:} під дією гомоморфізму дискримінант переходить\par
\smallskip\noindent\emph{Razlog:} Črez homomorfizm diskriminant prěhodi zachovuje napravjenje prěnosa: dany homomorfizm posylaje diskriminant v drugy objekt. Prěhodi ne smije byti zaměnjeno bezsměrnoju podobnostju. Izvorni razlog opiraje se na uže upotrěbljene homomorfizm, diskriminant i glagol prěhoda v dokazě o konečnom rangu.\par
\noindent\emph{Reason:} Črez homomorfizm diskriminant prěhodi retains the direction of the transfer: the given homomorphism sends the discriminant to another object. Passes to must not become an undirected statement of resemblance. The rationale continues the homomorphism, discriminant, and transition vocabulary of the finite-rank proof.\par
\noindent{\footnotesize Izvorne zapisy / Original records: HIST-g152-001.}\par
\end{minipage}\par
\par\medskip\Needspace{7\baselineskip}
\noindent\begin{minipage}{\linewidth}
\hypertarget{TERM-0357}{}
\noindent\textbf{racionalna baza, jednoznačno opreděljena poljem}\quad{\footnotesize\texttt{TERM-0357}}\par
\noindent\emph{Iztočnik / Source:} durch den Körper eindeutig bestimmte Rationalbasis\par
\noindent\emph{Rusky / Russian:} рациональная база, однозначно определенная полем\par
\noindent\emph{Ukrajinsky / Ukrainian:} раціональна база, однозначно визначена полем\par
\smallskip\noindent\emph{Razlog:} Racionalna baza, jednoznačno opreděljena poljem, zachovuje kontrast s proizvoljnymi bazami iz odděla 4. Polje jest to, čto oprěděljaje bazu; glavnaja informacija ne jest samo suščestvovanje bazy. Izpustiti jednoznačno ili poljem značilo by oslabiti upravo tu kvalifikaciju, ktoru zapis podčrkiva.\par
\noindent\emph{Reason:} Racionalna baza, jednoznačno opreděljena poljem retains the contrast with arbitrary bases in section 4. The field determines the basis; the point is not merely that a basis exists. Dropping jednoznačno or poljem would weaken the very qualification emphasized by the record.\par
\noindent{\footnotesize Izvorne zapisy / Original records: HIST-g045-003.}\par
\end{minipage}\par
\par\medskip\Needspace{7\baselineskip}
\noindent\begin{minipage}{\linewidth}
\hypertarget{TERM-0358}{}
\noindent\textbf{sut dělimy tymi samymi prostymi idealami p i v tyh samyh stepenjah}\quad{\footnotesize\texttt{TERM-0358}}\par
\noindent\emph{Iztočnik / Source:} durch dieselben Primideale p und zur gleichen Potenz teilbar\par
\noindent\emph{Rusky / Russian:} делятся одними и теми же простыми идеалами p и в одинаковых степенях\par
\noindent\emph{Ukrajinsky / Ukrainian:} діляться тими самими простими ідеалами p і в тих самих степенях\par
\smallskip\noindent\emph{Razlog:} Děliteljnost tymi samymi prostymi idealami v tyh samyh stepenjah objašnjaje, začto lokalne jednakosti dajut globalnu jednakost. Morajut ostati i identičnost prostyh idealov, i jednakost njihovih eksponentov. Izvorni zapis pripisuje tuto obrazloženje pečatnomu iztočniku; njegovo prědavanje ovde ne jest novo sravnjenje skana.\par
\noindent\emph{Reason:} Divisibility by the same prime ideals to the same powers explains why the local equalities imply global equality. Both the identities of the prime ideals and the equality of their exponents must remain. The record attributes this reasoning to the printed source; recounting it here is not a fresh scan comparison.\par
\noindent{\footnotesize Izvorne zapisy / Original records: HIST-g160-011.}\par
\end{minipage}\par
\par\medskip\Needspace{7\baselineskip}
\noindent\begin{minipage}{\linewidth}
\hypertarget{TERM-0359}{}
\noindent\textbf{črez dvojnu redukciju}\quad{\footnotesize\texttt{TERM-0359}}\par
\noindent\emph{Iztočnik / Source:} durch doppelte Reduktion\par
\noindent\emph{Rusky / Russian:} посредством двойной редукции\par
\noindent\emph{Ukrajinsky / Ukrainian:} шляхом подвійної редукції\par
\smallskip\noindent\emph{Razlog:} Črez dvojnu redukciju označaje sredstvo dokazovogo kroka. Izraz nastavlja imenovanje dvojnoj redukcije iz oddělov 3 i 7 i kratku predložnu konstrukciju, upotrěbjenu za postupky. Črez ne měnja čislo ili vid operacije: ona ostavaje ta sama imenovana dvojna redukcija.\par
\noindent\emph{Reason:} Črez dvojnu redukciju identifies the means of the proof step. It continues the double-reduction terminology of sections 3 and 7 and the compact prepositional construction used for procedures. Črez changes neither the operation's count nor its kind: it remains the same named double reduction.\par
\noindent{\footnotesize Izvorne zapisy / Original records: HIST-g010-004.}\par
\end{minipage}\par
\par\medskip\Needspace{7\baselineskip}
\noindent\begin{minipage}{\linewidth}
\hypertarget{TERM-0360}{}
\noindent\textbf{po dualistično protivopoložnom razmatrjenju}\quad{\footnotesize\texttt{TERM-0360}}\par
\noindent\emph{Iztočnik / Source:} durch dualistisch entgegengesetzte Betrachtung\par
\noindent\emph{Rusky / Russian:} посредством дуалистически противоположного рассмотрения\par
\noindent\emph{Ukrajinsky / Ukrainian:} через дуалістично протилежний розгляд\par
\smallskip\noindent\emph{Razlog:} Po dualistično protivopoložnom razmatrjenju označaje dokazovy prěhod po dualnosti. Izvorni razlog zachovuje vyraženu dualnost vměsto neoprěděljenogo ‘podobno’. On smatra konstrukciju razumělivoju, ale ne daje citat jezyčnogo iztočnika za njezinu idiomatičnost. Opravdanje dokazovoj funkcije i prověrka konkretnogo vyraženja zato sut različne zadače.\par
\noindent\emph{Reason:} Po dualistično protivopoložnom razmatrjenju names a proof move by duality. The rationale preserves explicit duality rather than replacing it with vaguely similarly. It judges the construction intelligible but cites no linguistic source for its idiomaticity. Preserving the proof function and checking the wording are therefore separate tasks.\par
\noindent{\footnotesize Izvorne zapisy / Original records: HIST-g016-009.}\par
\end{minipage}\par
\par\medskip\Needspace{7\baselineskip}
\noindent\begin{minipage}{\linewidth}
\hypertarget{TERM-0361}{}
\noindent\textbf{dělimy črez ideal / $f\equiv0\;(M)$}\quad{\footnotesize\texttt{TERM-0361}}\par
\noindent\emph{Iztočnik / Source:} durch ein Ideal teilbar / $f\equiv0\;(M)$\par
\noindent\emph{Rusky / Russian:} делится на идеал / $f\equiv0\;(M)$\par
\noindent\emph{Ukrajinsky / Ukrainian:} ділиться на ідеал / $f\equiv0\;(M)$\par
\smallskip\noindent\emph{Razlog:} Dělimy črez ideal zachovuje Noetherin jezyk děliteljnosti idealov i njegovu vezu s kongruencijeju $f\equiv0\pmod{M}$. Směr děliteljnosti idealov ne može byti pročitan kako směr vključjenja množin. Izvorni formula ima zapis $f\equiv0\;(M)$; drugy tipografski zapis kongruencije služi objasnjenju, ne zaměně iztočnoj notacije.\par
\noindent\emph{Reason:} Dělimy črez ideal retains Noether's ideal-divisibility language and its relation to the congruence $f\equiv0\pmod{M}$. The direction of ideal divisibility must not be read as the direction of set inclusion. The recorded formula uses $f\equiv0\;(M)$; the alternative congruence typography explains the sense rather than replacing the source notation.\par
\noindent{\footnotesize Izvorne zapisy / Original records: HIST-g081-023, HIST-g082-023.}\par
\end{minipage}\par
\par\medskip\Needspace{7\baselineskip}
\noindent\begin{minipage}{\linewidth}
\hypertarget{TERM-0362}{}
\noindent\textbf{črez zaměnu formy $\mathfrak{A}$ nižšimi formami ręda form}\quad{\footnotesize\texttt{TERM-0362}}\par
\noindent\emph{Iztočnik / Source:} durch Ersetzen der Form $\mathfrak{A}$ durch niedere Formen der Formenreihe\par
\noindent\emph{Rusky / Russian:} благодаря замене формы $\mathfrak{A}$ нижними формами ряда форм\par
\noindent\emph{Ukrajinsky / Ukrainian:} завдяки заміні форми $\mathfrak{A}$ нижчими формами ряду форм\par
\smallskip\noindent\emph{Razlog:} Črez zaměnu formy nižšimi formami reda form čini postupok redukcije izričnym. Zaměna, nižši i red jesut izbrane kako deljene slovjanske osnovy. Identifikator zaměnjanoj formy ostavaje nužny za sledovanje dokaza; nižši imenuje njezino město v danom redu, ne neoprěděljenu matematičnu prostotu.\par
\noindent\emph{Reason:} Črez zaměnu formy nižšimi formami reda form makes the reduction procedure explicit. Zaměna, nižši, and red were selected as shared Slavic stems. The identifier of the replaced form remains necessary for following the proof; lower concerns its place in the stated series, not an unspecified mathematical simplicity.\par
\noindent{\footnotesize Izvorne zapisy / Original records: HIST-g011-008.}\par
\end{minipage}\par
\par\medskip\Needspace{7\baselineskip}
\noindent\begin{minipage}{\linewidth}
\hypertarget{TERM-0363}{}
\noindent\textbf{črez zaměnu ...}\quad{\footnotesize\texttt{TERM-0363}}\par
\noindent\emph{Iztočnik / Source:} durch Ersetzen von ...\par
\noindent\emph{Rusky / Russian:} через замену ...\par
\noindent\emph{Ukrajinsky / Ukrainian:} через заміну ...\par
\smallskip\noindent\emph{Razlog:} Črez zaměnu jest kratka proceduralna konstrukcija za zaměnu izraza v redukciji. Ruska i ukrajinska paralela imaju odgovarjajuče imennikove frazy. Razlog podpiraje kontinuitet dokazovogo jezyka, ale ne dopuščaje izpustiti objekty zaměny, kogda oni sut nužni za točnost kroka.\par
\noindent\emph{Reason:} Črez zaměnu is a compact procedural construction for replacement in a reduction. Russian and Ukrainian supply corresponding nominal phrases. The reason supports continuity in proof prose, not omission of the objects being replaced when they are needed to specify the step.\par
\noindent{\footnotesize Izvorne zapisy / Original records: HIST-g015-012.}\par
\end{minipage}\par
\par\medskip\Needspace{7\baselineskip}
\noindent\begin{minipage}{\linewidth}
\hypertarget{TERM-0364}{}
\noindent\textbf{dělimo črez p / izčezaje v P}\quad{\footnotesize\texttt{TERM-0364}}\par
\noindent\emph{Iztočnik / Source:} durch p teilbar / verschwindet in P\par
\noindent\emph{Rusky / Russian:} делится на p / обращается в нуль в P\par
\noindent\emph{Ukrajinsky / Ukrainian:} ділиться на p / зникає в P\par
\smallskip\noindent\emph{Razlog:} Dělimo črez $p$ i izčezaje v $P$ svazuju děliteljnost s nulovostju klasy koeficienta v ostatkovom polju. Ruska paralela upotrěbjaje obrat ‘obraščaje se v nulu’, dokolě ukrajinska i medžuslovjanska imaju glagol izčezanja. On ovde znači nulu v opreděljenom polju, ne nestajanje samogo formalnogo izraza.\par
\noindent\emph{Reason:} Dělimo črez $p$ and izčezaje v $P$ connect divisibility with the coefficient class being zero in the residue field. Russian uses the turns-to-zero idiom, while Ukrainian and Interslavic use vanishing. Here vanishing means zero in the specified field, not disappearance of the formal expression itself.\par
\noindent{\footnotesize Izvorne zapisy / Original records: HIST-g151-008.}\par
\end{minipage}\par
\par\medskip\Needspace{7\baselineskip}
\noindent\begin{minipage}{\linewidth}
\hypertarget{TERM-0365}{}
\noindent\textbf{črez tvorjenje kvocientov poraždaje medžupolje}\quad{\footnotesize\texttt{TERM-0365}}\par
\noindent\emph{Iztočnik / Source:} durch Quotientenbildung einen Zwischenkörper erzeugt\par
\noindent\emph{Rusky / Russian:} посредством образования частных порождает промежуточное поле\par
\noindent\emph{Ukrajinsky / Ukrainian:} через утворення часток породжує проміжне поле\par
\smallskip\noindent\emph{Razlog:} Črez tvorjenje kvocientov poraždaje medžupolje zachovuje operaciju i rezultat v argumentu o poredkah nižšego ranga. Medžupolje jest kompaktny kandidat za pole medžu danymi poljami, ne proizvoljno novo razširjenje. Razlog svazuje pasažu s dostatočnostju ranga $n$; jego pozvanje na pečatny iztočnik ostavaje historijnoju tvrdnjeju, ne novoju kolacijeju.\par
\noindent\emph{Reason:} Črez tvorjenje kvocientov poraždaje medžupolje retains both operation and result in the argument about lower-rank orders. Medžupolje is a compact candidate for an intermediate field, not an arbitrary new extension. The rationale connects the passage with the sufficiency of rank $n$; its appeal to the printed source remains a historical claim rather than a fresh collation.\par
\noindent{\footnotesize Izvorne zapisy / Original records: HIST-g149-007.}\par
\end{minipage}\par
\par\medskip\Needspace{7\baselineskip}
\noindent\begin{minipage}{\linewidth}
\hypertarget{TERM-0366}{}
\noindent\textbf{črez vraćanje ... k nižšim formam}\quad{\footnotesize\texttt{TERM-0366}}\par
\noindent\emph{Iztočnik / Source:} durch Zurückgehen ... auf niedrigere Formen\par
\noindent\emph{Rusky / Russian:} сведением ... к нижним формам\par
\noindent\emph{Ukrajinsky / Ukrainian:} зводячи ... до нижчих форм\par
\smallskip\noindent\emph{Razlog:} Črez vraćanje ... k nižšim formam nastavlja izraz iz odděla 18 za svedenje k nižšim formam. Etimologično ć jest čestju zachovanogo kandidata, ne dokaz nužnosti toj bukvy v osnovnom pravopisu. Izvorni razlog ne dava sravnjenje, ktoro by dokončatelno vybralo glagol vračanja vměsto drugoj formulacije redukcije.\par
\noindent\emph{Reason:} Črez vraćanje ... k nižšim formam continues section 18's expression for reduction to lower forms. Etymological ć belongs to the recorded candidate, not a requirement to use it in basic spelling. The reason contains no comparison that conclusively selects returning over another reduction construction.\par
\noindent{\footnotesize Izvorne zapisy / Original records: HIST-g021-007.}\par
\end{minipage}\par
\par\medskip\Needspace{7\baselineskip}
\noindent\begin{minipage}{\linewidth}
\hypertarget{TERM-0367}{}
\noindent\textbf{črez vraćanje od $j$ i $\mathfrak{A}$ k nižšim formam}\quad{\footnotesize\texttt{TERM-0367}}\par
\noindent\emph{Iztočnik / Source:} durch Zurückgehen von $j$ und $\mathfrak{A}$ auf niedrigere Formen\par
\noindent\emph{Rusky / Russian:} при возвращении от $j$ и $\mathfrak{A}$ к более низким формам\par
\noindent\emph{Ukrajinsky / Ukrainian:} через повернення від $j$ і $\mathfrak{A}$ до нижчих форм\par
\smallskip\noindent\emph{Razlog:} Fraza zachovuje oba imenovane početkove objekty i směr k nižšim formam. Izvorni razlog govori o razkryvanju zbitogo německogo vyraženja i prědlagaje sravniti črez vraćanje s konstrukcijeju nazad ili k. Takva alternativa jest izrično zapisana, ale ne jest dano gotovo drugo vyraženje, ktoro by možno citovati kako prijety izbor.\par
\noindent\emph{Reason:} The phrase retains both named starting objects and the direction towards lower forms. The rationale describes unpacking the compressed German expression and proposes comparing črez vraćanje with a construction using nazad or k. That alternative is genuinely recorded, but no complete second expression is supplied that could be quoted as an adopted choice.\par
\noindent{\footnotesize Izvorne zapisy / Original records: HIST-g019-007.}\par
\end{minipage}\par
\par\medskip\Needspace{7\baselineskip}
\noindent\begin{minipage}{\linewidth}
\hypertarget{TERM-0368}{}
\noindent\textbf{tvorjenje prěsěkov}\quad{\footnotesize\texttt{TERM-0368}}\par
\noindent\emph{Iztočnik / Source:} Durchschnittsbildung\par
\noindent\emph{Rusky / Russian:} взятие пересечений\par
\noindent\emph{Ukrajinsky / Ukrainian:} утворення перетинів\par
\smallskip\noindent\emph{Razlog:} Tvorjenje prěsěkov imenuje operaciju prěhoda iz razširjenogo prstena nazad k $R$ črěz prěsěk. Imennik v proze davaje značenje, a formula zachovuje avtorske skobky. Tute dve funkcije ne trěba směšati: objasniti operaciju ne znači modernizovati njezinu iztočnu notaciju.\par
\noindent\emph{Reason:} Tvorjenje prěsěkov names the intersection or contraction step from the extended ring back to $R$. The prose noun supplies the meaning while the formula retains the author's brackets. These functions must remain distinct: explaining the operation is not the same as modernizing its source notation.\par
\noindent{\footnotesize Izvorne zapisy / Original records: HIST-g126-005.}\par
\end{minipage}\par
\par\medskip\Needspace{7\baselineskip}
\noindent\begin{minipage}{\linewidth}
\hypertarget{TERM-0369}{}
\noindent\textbf{prěsěčno polje}\quad{\footnotesize\texttt{TERM-0369}}\par
\noindent\emph{Iztočnik / Source:} Durchschnittskörper\par
\noindent\emph{Rusky / Russian:} поле-пересечение\par
\noindent\emph{Ukrajinsky / Ukrainian:} поле-перетин\par
\smallskip\noindent\emph{Razlog:} Prěsěčno polje imenuje polje polučeno kak prěsěk, ne samo proizvoljny prěsěk množin. Imennik polje zachovuje algebraičny tip objekta. Kratko složeno nazvanje podpirajut i ruska, i ukrajinska paralela; zapis ne dava konkurentnu medžuslovjansku konstrukciju s odvojenym genitivom.\par
\noindent\emph{Reason:} Prěsěčno polje names a field obtained as an intersection, not merely an arbitrary set intersection. Polje preserves the algebraic type of object. The compact naming has Russian and Ukrainian parallels; the record provides no competing Interslavic genitive construction.\par
\noindent{\footnotesize Izvorne zapisy / Original records: HIST-g065-015.}\par
\end{minipage}\par
\par\medskip\Needspace{7\baselineskip}
\noindent\begin{minipage}{\linewidth}
\hypertarget{TERM-0371}{}
\noindent\textbf{$e_1,\ldots,e_n$ kako modulna baza / diagonalna matrica}\quad{\footnotesize\texttt{TERM-0371}}\par
\noindent\emph{Iztočnik / Source:} $e_1,\ldots,e_n$ als Modulbasis / Diagonalmatrix\par
\noindent\emph{Rusky / Russian:} $e_1,\ldots,e_n$ как модульная база / диагональная матрица\par
\noindent\emph{Ukrajinsky / Ukrainian:} $e_1,\ldots,e_n$ як модульна база / діагональна матриця\par
\smallskip\noindent\emph{Razlog:} Modulna baza i diagonalna matrica imenujut dve česti linearno-algebraičnoj konstrukcije: izbrany bazny seznam i matricu pridruženu každomu elementu. Izložene formuly ustanavjajut indeksny zapis. Proza zato ne zaměnjaje formuly, ale objašnjaje rolu točno imenovanyh objektov.\par
\noindent\emph{Reason:} Modulna baza and diagonalna matrica name two parts of the linear-algebra construction: the chosen basis list and the matrix assigned to each element. The displays specify the indexed notation. The prose does not replace those formulas; it explains the roles of the named objects.\par
\noindent{\footnotesize Izvorne zapisy / Original records: HIST-g147-001.}\par
\end{minipage}\par
\par\medskip\Needspace{7\baselineskip}
\noindent\begin{minipage}{\linewidth}
\hypertarget{TERM-0372}{}
\noindent\textbf{vlastna podgrupa}\quad{\footnotesize\texttt{TERM-0372}}\par
\noindent\emph{Iztočnik / Source:} echte Untergruppe\par
\noindent\emph{Rusky / Russian:} собственная подгруппа\par
\noindent\emph{Ukrajinsky / Ukrainian:} справжня підгрупа\par
\smallskip\noindent\emph{Razlog:} Vlastna podgrupa označaje strogo vključenu podgrupu, različnu od cěloj grupy. Izvorni razlog dava prednost slovjanskomu vlastna i ostavja njegove stilistične vlastnosti za sravnjenje. Ta kvalifikacija sama ne isključaje trivialnu podgrupu; uslovja konkretnogo dokaza trěba čitati odvojeno.\par
\noindent\emph{Reason:} Vlastna podgrupa marks a subgroup strictly contained in, and different from, the whole group. The rationale favours Slavic vlastna while leaving its stylistic qualities for comparison. Properness alone does not exclude the trivial subgroup; the proof's specific conditions must be read separately.\par
\noindent{\footnotesize Izvorne zapisy / Original records: HIST-g079-048.}\par
\end{minipage}\par
\par\medskip\Needspace{7\baselineskip}
\noindent\begin{minipage}{\linewidth}
\hypertarget{TERM-0373}{}
\noindent\textbf{vlastny dělitelj}\quad{\footnotesize\texttt{TERM-0373}}\par
\noindent\emph{Iztočnik / Source:} echter Teiler\par
\noindent\emph{Rusky / Russian:} собственный делитель\par
\noindent\emph{Ukrajinsky / Ukrainian:} власний дільник\par
\smallskip\noindent\emph{Razlog:} Vlastny dělitelj zachovuje strogi odnos děliteljnosti idealov v cěpnom argumentu. Tri razloga podpirajut kontinuitet togo jezyka vměsto tihoj zaměny vključjenjem podmnožin. Jedin govori široko o vlastnyh ili netrivialnyh podobektah; ty kvalifikacije ne sut avtomatično sinonimy. V vsakom pojavjenju trěba ustanoviti, ktora jednakost ili trivialny slučaj jest izključen.\par
\noindent\emph{Reason:} Vlastny dělitelj retains strict ideal divisibility in the chain argument. Three rationales support continuity of that language instead of silently replacing it with subset terminology. One loosely groups proper and nontrivial subobjects; those qualifications are not automatically synonymous. Each occurrence requires identifying the equality or trivial case that is actually excluded.\par
\noindent{\footnotesize Izvorne zapisy / Original records: HIST-g081-024, HIST-g082-024, HIST-g096-013, HIST-g113-003.}\par
\end{minipage}\par
\par\medskip\Needspace{7\baselineskip}
\noindent\begin{minipage}{\linewidth}
\hypertarget{TERM-0374}{}
\noindent\textbf{vlastny dělitelj / prěsěčna klasa}\quad{\footnotesize\texttt{TERM-0374}}\par
\noindent\emph{Iztočnik / Source:} echter Teiler / Durchschnittsklasse\par
\noindent\emph{Rusky / Russian:} собственный делитель / класс пересечения\par
\noindent\emph{Ukrajinsky / Ukrainian:} власний дільник / клас перетину\par
\smallskip\noindent\emph{Razlog:} Vlastny dělitelj i prěsěčna klasa ostavajut dva odvojene pojmy v završnoj primětkě. Pervy služi opisu vloženja v sebe s vlastnym děliteljem, drugy imenovanju klasy, čija dobra opreděljenost jest otricana. Slična blizost dvuh terminov v tekstu ne dava osnovy svesti jih k jednoj konstrukciji.\par
\noindent\emph{Reason:} Vlastny dělitelj and prěsěčna klasa remain separate concepts in the closing remark. The first belongs to the proper-divisor self-embedding discussion; the second names an intersection class whose well-definedness is denied. Their proximity in the passage does not justify collapsing them into one construction.\par
\noindent{\footnotesize Izvorne zapisy / Original records: HIST-g067-012.}\par
\end{minipage}\par
\par\medskip\Needspace{7\baselineskip}
\noindent\begin{minipage}{\linewidth}
\hypertarget{TERM-0375}{}
\noindent\textbf{vlastny dělitelj / vlastno kratno}\quad{\footnotesize\texttt{TERM-0375}}\par
\noindent\emph{Iztočnik / Source:} echter Teiler / echtes Vielfaches\par
\noindent\emph{Rusky / Russian:} собственный делитель / собственное кратное\par
\noindent\emph{Ukrajinsky / Ukrainian:} власний дільник / власне кратне\par
\smallskip\noindent\emph{Razlog:} Vlastny dělitelj i vlastno kratno označajut strogost kroka v cěpi. Izbor zachovuje Noetherino simetrično imenovanje dělitelja i kratnogo vměsto vvedenja novyh podobektov samo radi sovrěmennogo jezyka. Strogost odnašaje se k danomu odnosu; ona ne měnja jeho směr.\par
\noindent\emph{Reason:} Vlastny dělitelj and vlastno kratno mark the strictness of a chain step. The choice preserves Noether's paired divisor and multiple vocabulary rather than introducing new subobject terminology merely to sound modern. Strictness qualifies the given relation; it does not reverse its direction.\par
\noindent{\footnotesize Izvorne zapisy / Original records: HIST-g109-006.}\par
\end{minipage}\par
\par\medskip\Needspace{7\baselineskip}
\noindent\begin{minipage}{\linewidth}
\hypertarget{TERM-0376}{}
\noindent\textbf{vlastny dělitelj / čislo členov}\quad{\footnotesize\texttt{TERM-0376}}\par
\noindent\emph{Iztočnik / Source:} echter Teiler / Gliederzahl\par
\noindent\emph{Rusky / Russian:} собственный делитель / число членов\par
\noindent\emph{Ukrajinsky / Ukrainian:} власний дільник / кількість членів\par
\smallskip\noindent\emph{Razlog:} Vlastny dělitelj i čislo členov svazuju idealnu děliteljnost s veličinoju cěpi. Pri Noetherinoj konvenciji dělitelj ideala može byti boljšim po uključjenju, a ne menšim. Izvorni razlog zato nastavlja ranějši cěpny jezyk i ne prepisuje njegovo napravjenje po intuitivnoj veličině množiny.\par
\noindent\emph{Reason:} Vlastny dělitelj and čislo členov connect ideal divisibility with the size of a chain. Under Noether's convention, an ideal divisor can be larger by inclusion rather than smaller. The reason therefore continues the earlier chain vocabulary instead of reversing its direction according to an intuitive notion of set size.\par
\noindent{\footnotesize Izvorne zapisy / Original records: HIST-g098-007.}\par
\end{minipage}\par
\par\medskip\Needspace{7\baselineskip}
\noindent\begin{minipage}{\linewidth}
\hypertarget{TERM-0377}{}
\noindent\textbf{vlastny dělitelj jednakogo ranga}\quad{\footnotesize\texttt{TERM-0377}}\par
\noindent\emph{Iztočnik / Source:} echter Teiler gleichen Ranges\par
\noindent\emph{Rusky / Russian:} собственный делитель того же ранга\par
\noindent\emph{Ukrajinsky / Ukrainian:} справжній дільник того самого рангу\par
\smallskip\noindent\emph{Razlog:} Vlastny dělitelj jednakogo ranga zachovuje i strogu děliteljnost modulov, i jednakost njihovih rangov. Podnožna primětka objašnjaje odnos děliteljnosti k obratnomu uključjenju. Jednak rang ne daje prava zaměniti oba moduly jednim objektom; upravo strogost i jednakost ranga sut ovde sovměstne informacije.\par
\noindent\emph{Reason:} Vlastny dělitelj jednakogo ranga retains both strict module divisibility and equality of ranks. The footnote explains the reverse relation between divisibility and inclusion. Equal rank does not justify identifying the two modules; strictness and equal rank are jointly relevant information here.\par
\noindent{\footnotesize Izvorne zapisy / Original records: HIST-g085-007, HIST-g092-007.}\par
\end{minipage}\par
\par\medskip\Needspace{7\baselineskip}
\noindent\begin{minipage}{\linewidth}
\hypertarget{TERM-0378}{}
\noindent\textbf{vlastno podkolco}\quad{\footnotesize\texttt{TERM-0378}}\par
\noindent\emph{Iztočnik / Source:} echter Unterring\par
\noindent\emph{Rusky / Russian:} собственное подкольцо\par
\noindent\emph{Ukrajinsky / Ukrainian:} власне підкільце\par
\smallskip\noindent\emph{Razlog:} Vlastno podkolco označaje strogo vključjenje pri sravnjenju centra faktor-prstena s faktorom centra. Pridavnik ne znači ‘pravilno’ ili ‘stvarno’, ale izključaje jednakost v tvrđenom uključjenju. Podkolco jest historijno imenovanje objekta; jego prisutnost sama ne razrěšaje aktualny izbor prsten ili kolco.\par
\noindent\emph{Reason:} Vlastno podkolco marks strict inclusion in the comparison between the centre of a quotient ring and the quotient of the centre. The adjective means neither correct nor real; it excludes equality in the stated inclusion. Podkolco is the historical object name, not by itself a resolution of the current prsten/kolco preference.\par
\noindent{\footnotesize Izvorne zapisy / Original records: HIST-g190-006.}\par
\end{minipage}\par
\par\medskip\Needspace{7\baselineskip}
\noindent\begin{minipage}{\linewidth}
\hypertarget{TERM-0379}{}
\noindent\textbf{vlastno primarna komponenta}\quad{\footnotesize\texttt{TERM-0379}}\par
\noindent\emph{Iztočnik / Source:} eigentlich primäre Komponente\par
\noindent\emph{Rusky / Russian:} собственно первичная компонента\par
\noindent\emph{Ukrajinsky / Ukrainian:} власне первинна компонента\par
\smallskip\noindent\emph{Razlog:} Vlastno primarna komponenta ovde odlišaje primarny slučaj od prostogo. Zapis svazuje razliku s komponentoju, ktora ne jest polje, i s terminologijeju uvoda. Vlastno ne jest prosto pojačanje poput ‘napravdu’: ono nosi točno matematično izključenje.\par
\noindent\emph{Reason:} Vlastno primarna komponenta distinguishes the primary case from the prime case. The record connects that distinction with a component that is not a field and with the introduction's terminology. Vlastno is not emphasis such as genuinely; it carries a specific mathematical exclusion.\par
\noindent{\footnotesize Izvorne zapisy / Original records: HIST-g128-002.}\par
\end{minipage}\par
\par\medskip\Needspace{7\baselineskip}
\noindent\begin{minipage}{\linewidth}
\hypertarget{TERM-0380}{}
\noindent\textbf{vlastno primarno kolco}\quad{\footnotesize\texttt{TERM-0380}}\par
\noindent\emph{Iztočnik / Source:} eigentlich primärer Ring\par
\noindent\emph{Rusky / Russian:} собственно первичное кольцо\par
\noindent\emph{Ukrajinsky / Ukrainian:} власне первинне кільце\par
\smallskip\noindent\emph{Razlog:} Vlastno primarno kolco znači ovde primarny, ale ne prosty slučaj. Razlog odbacaje čitanje vlastno kako ‘realno’ i daje mu tehničnu kvalifikacijsku funkciju. Ta kvalifikacija ne govori samo, že prsten jest različny od nekogo okružajučego prstena; ne trěba směšati ju s vlastnym podprstenom.\par
\noindent\emph{Reason:} Vlastno primarno kolco means the primary but non-prime case here. The reason rejects reading vlastno as real and assigns it a technical qualifying role. This is not merely a statement that the ring differs from an ambient ring; it must not be confused with proper subring.\par
\noindent{\footnotesize Izvorne zapisy / Original records: HIST-g143-001.}\par
\end{minipage}\par
\par\medskip\Needspace{7\baselineskip}
\noindent\begin{minipage}{\linewidth}
\hypertarget{TERM-0381}{}
\noindent\textbf{vlastne / nevlastne divergencijne relacije}\quad{\footnotesize\texttt{TERM-0381}}\par
\noindent\emph{Iztočnik / Source:} eigentliche / uneigentliche Divergenzrelationen\par
\noindent\emph{Rusky / Russian:} собственные / несобственные дивергентные соотношения\par
\noindent\emph{Ukrajinsky / Ukrainian:} власні / невласні дивергентні співвідношення\par
\smallskip\noindent\emph{Razlog:} Vlastne i nevlastne divergencijne relacije imenujut klasifikaciju po tomu, dali relacije polučajut se linearno iz zavisimostij. Izvorni razlog izrično odbacaje pravdive i ložne: nevlastna relacija ne jest zato matematično netočna. Vlastny i nevlastny ostavajut slovotvorno sravnime kandidaty, ne rezultat nezavisnogo ispitovanja čitateljev.\par
\noindent\emph{Reason:} Vlastne and nevlastne divergencijne relacije name a classification involving whether the relations arise linearly from dependencies. The rationale explicitly rejects true and false: an improper relation is not thereby mathematically incorrect. Vlastny and nevlastny remain comparable lexical candidates, not the outcome of an independent reader study.\par
\noindent{\footnotesize Izvorne zapisy / Original records: HIST-g075-017.}\par
\end{minipage}\par
\par\medskip\Needspace{7\baselineskip}
\noindent\begin{minipage}{\linewidth}
\hypertarget{TERM-0382}{}
\noindent\textbf{vlastna teorija eliminacije}\quad{\footnotesize\texttt{TERM-0382}}\par
\noindent\emph{Iztočnik / Source:} Eigentliche Eliminationstheorie\par
\noindent\emph{Rusky / Russian:} собственно теория исключения\par
\noindent\emph{Ukrajinsky / Ukrainian:} власне теорія елімінування\par
\smallskip\noindent\emph{Razlog:} Vlastna teorija eliminacije razlučaje vlastny eliminacijny krok odděla 6 od teorije eliminacije črěz rezultantne formy v oddělu 7. Vlastna trěba čitati po toj strukturnoj opoziciji, ne kako obecnu ocěnu pravdivosti ili važnosti teorije. Dva zapisy lokuju termin v toj samoj rabotě.\par
\noindent\emph{Reason:} Vlastna teorija eliminacije distinguishes section 6's elimination step proper from section 7's resultant-form elimination theory. Vlastna must be read through that structural contrast, not as a general judgment of the theory's truth or importance. Two records locate the term within the same work.\par
\noindent{\footnotesize Izvorne zapisy / Original records: HIST-g090-001, HIST-g092-055.}\par
\end{minipage}\par
\par\medskip\Needspace{7\baselineskip}
\noindent\begin{minipage}{\linewidth}
\hypertarget{TERM-0383}{}
\noindent\textbf{vlastno primarne idealy abo proste idealy}\quad{\footnotesize\texttt{TERM-0383}}\par
\noindent\emph{Iztočnik / Source:} eigentliche Primärideale oder Primideale\par
\noindent\emph{Rusky / Russian:} собственно первичные идеалы или простые идеалы\par
\noindent\emph{Ukrajinsky / Ukrainian:} власне первинні ідеали або прості ідеали\par
\smallskip\noindent\emph{Razlog:} Vlastno primarne idealy abo proste idealy zachovuje dva člena razděljenja. Prvi ne uključaje vtory v smyslu konkretnoj pasaže: vlastno primarny izrično odlišaje primarno od prostogo. Izraz zato ne može byti skračen do ‘primarne idealy’ bez prověrky, čto izključenje jest zachovano v drugom městu věty.\par
\noindent\emph{Reason:} Vlastno primarne idealy abo proste idealy retains both parts of the distinction. In this passage the first is explicitly primary but non-prime. Shortening to primary ideals would lose the contrast unless another part of the sentence preserved the exclusion.\par
\noindent{\footnotesize Izvorne zapisy / Original records: HIST-g152-007.}\par
\end{minipage}\par
\par\medskip\Needspace{7\baselineskip}
\noindent\begin{minipage}{\linewidth}
\hypertarget{TERM-0384}{}
\noindent\textbf{vlastno primarna komponenta}\quad{\footnotesize\texttt{TERM-0384}}\par
\noindent\emph{Iztočnik / Source:} eigentliche Primärkomponente\par
\noindent\emph{Rusky / Russian:} собственно первичная компонента\par
\noindent\emph{Ukrajinsky / Ukrainian:} власне первинна компонента\par
\smallskip\noindent\emph{Razlog:} Vlastno primarna komponenta zachovuje tehničny kontrast s prostymi komponentami v ranějšej klasifikaciji prstenov ostatkovyh klasov. Prilog ne jest ukrasny. Izvorni razlog opravdaje imenno kontinuitet togo različenja, ne opču zaměnu vsih značenij vlastny jednim modernym anglijskym proper.\par
\noindent\emph{Reason:} Vlastno primarna komponenta preserves the technical contrast with prime components in the earlier residue-class-ring taxonomy. The adverb is not decorative. The reason supports continuity of that distinction, not a universal identification of every vlastny sense with modern English proper.\par
\noindent{\footnotesize Izvorne zapisy / Original records: HIST-g163-007.}\par
\end{minipage}\par
\par\medskip\Needspace{7\baselineskip}
\noindent\begin{minipage}{\linewidth}
\hypertarget{TERM-0385}{}
\noindent\textbf{vlastno primarna komponenta / prosty ideal drugogo roda}\quad{\footnotesize\texttt{TERM-0385}}\par
\noindent\emph{Iztočnik / Source:} eigentliche Primärkomponente / Primideal zweiter Art\par
\noindent\emph{Rusky / Russian:} собственно первичная компонента / простой идеал второго рода\par
\noindent\emph{Ukrajinsky / Ukrainian:} власне первинна компонента / простий ідеал другого роду\par
\smallskip\noindent\emph{Razlog:} Vlastno primarna komponenta i prosty ideal drugogo roda sut alternative v kriteriju diskriminanta. Drugogo roda ostavaje doslovno, ibo rabota 31 tehnično razlučaje dva typy povedenja ostatkovogo polja. Tuto ne jest redni broj ideala v seznamu; izpustiti kvalifikaciju značilo by proměniti kriterij.\par
\noindent\emph{Reason:} Vlastno primarna komponenta and prosty ideal drugogo roda are alternatives in the discriminant criterion. Of the second kind is retained because Paper 31 technically distinguishes two kinds of residue-field behaviour. It is not the ideal's position in a list; dropping the qualification would change the criterion.\par
\noindent{\footnotesize Izvorne zapisy / Original records: HIST-g153-004.}\par
\end{minipage}\par
\par\medskip\Needspace{7\baselineskip}
\noindent\begin{minipage}{\linewidth}
\hypertarget{TERM-0386}{}
\noindent\textbf{vlastny primarny ideal}\quad{\footnotesize\texttt{TERM-0386}}\par
\noindent\emph{Iztočnik / Source:} eigentliches Primärideal\par
\noindent\emph{Rusky / Russian:} собственный первичный идеал\par
\noindent\emph{Ukrajinsky / Ukrainian:} власний первинний ідеал\par
\smallskip\noindent\emph{Razlog:} Vlastny primarny ideal v tom zapisu znači primarny ideal, ktory ne jest prosty. To jest silnějša informacija než samo različie od cělogo prstena. Tu točnu funkciju pridavnika trěba zachovati, da by čitatelj ne prěnesl sem drugo sovrěmenno značenje ‘proper ideal’.\par
\noindent\emph{Reason:} Vlastny primarny ideal in this record means a primary ideal that is not prime. That says more than merely being different from the whole ring. The adjective's precise function must be retained so that a reader does not import the different modern sense of proper ideal.\par
\noindent{\footnotesize Izvorne zapisy / Original records: HIST-g095-008.}\par
\end{minipage}\par
\par\medskip\Needspace{7\baselineskip}
\noindent\begin{minipage}{\linewidth}
\hypertarget{TERM-0387}{}
\noindent\textbf{vlastny primarny ideal / vlastno primarna komponenta}\quad{\footnotesize\texttt{TERM-0387}}\par
\noindent\emph{Iztočnik / Source:} eigentliches Primärideal / eigentlich primäre Komponente\par
\noindent\emph{Rusky / Russian:} собственный первичный идеал / собственно первичная компонента\par
\noindent\emph{Ukrajinsky / Ukrainian:} власний первинний ідеал / власне первинна компонента\par
\smallskip\noindent\emph{Razlog:} Para vlastny primarny ideal i vlastno primarna komponenta vyražaje isto izključenje prostogo slučaja različnymi gramatičnymi konstrukcijami. Pridavnik soglasuje se s ideal, prilog kvalifikuje primarna. Izvorni razlog izrično odbacaje neoprěděljeno ‘stvarny’: izbor slova jest upravljany matematičnym kontrastom.\par
\noindent\emph{Reason:} Vlastny primarny ideal and vlastno primarna komponenta express the same exclusion of the prime case through different grammatical constructions. The adjective agrees with ideal; the adverb qualifies primarna. The recorded reason explicitly rejects vague actual: the mathematical contrast governs the wording.\par
\noindent{\footnotesize Izvorne zapisy / Original records: HIST-g118-005.}\par
\end{minipage}\par
\par\medskip\Needspace{7\baselineskip}
\noindent\begin{minipage}{\linewidth}
\hypertarget{TERM-0388}{}
\noindent\textbf{vlastne vrědnosti i vlastne funkcije}\quad{\footnotesize\texttt{TERM-0388}}\par
\noindent\emph{Iztočnik / Source:} Eigenwerte und Eigenfunktionen\par
\noindent\emph{Rusky / Russian:} собственные значения и собственные функции\par
\noindent\emph{Ukrajinsky / Ukrainian:} власні значення і власні функції\par
\smallskip\noindent\emph{Razlog:} Vlastne vrědnosti i vlastne funkcije sledujut slovjanskoj konstrukciji za objekty teorije linearnyh operatorov. Vlastne ne znači tu vlastne podobekty ili primarny-ne-prosty slučaj. Zapis ostavja formy otvorene za sravnjenje s dalšimi operatornymi pasažami; kontinuitet korenja sam ne jest dovoljen za pravilnost vseh jeho smyslov.\par
\noindent\emph{Reason:} Vlastne vrědnosti and vlastne funkcije follow the Slavic construction for linear-operator concepts. Here vlastne does not mean proper subobjects or the primary-but-non-prime case. The record leaves the forms open to comparison with later operator passages; continuity of the stem alone cannot establish every sense as correct.\par
\noindent{\footnotesize Izvorne zapisy / Original records: HIST-g078-011.}\par
\end{minipage}\par
\par\medskip\Needspace{7\baselineskip}
\noindent\begin{minipage}{\linewidth}
\hypertarget{TERM-0389}{}
\noindent\textbf{slu?aj jednozna?nogo razklada}\quad{\footnotesize\texttt{TERM-0389}}\par
\noindent\emph{Iztočnik / Source:} ein Fall eindeutiger Zerlegung\par
\noindent\emph{Rusky / Russian:} [Bukvy ne sut zachovane / Lettering not recoverable]\par
\noindent\emph{Ukrajinsky / Ukrainian:} [Bukvy ne sut zachovane / Lettering not recoverable]\par
\smallskip\noindent\emph{Razlog:} Izvorni razlog ograničaje tvrdženje do jednogo slučaja jednoznačnogo razloženja, upravljanogo diagonalnoju shemoju produktov. Ne predstavlja jeho kak obči teorem Krull–Schmidt. Čest medžuslovjanske grafije i cěle ruske i ukrajinske formy sut povrědžene; možno zachovati ograničenje smysla, ale ne izmysliti nedostajuče bukvy.\par
\noindent\emph{Reason:} The rationale restricts the claim to one uniqueness-of-decomposition case governed by a diagonal product scheme. It does not present it as a general Krull–Schmidt theorem. Part of the Interslavic spelling and the entire Russian and Ukrainian forms are damaged; the sense restriction can be retained without inventing the missing letters.\par
\noindent{\footnotesize Izvorne zapisy / Original records: HIST-g079-037.}\par
\end{minipage}\par
\par\medskip\Needspace{7\baselineskip}
\noindent\begin{minipage}{\linewidth}
\hypertarget{TERM-0390}{}
\noindent\textbf{jedno- i dvojnokratna svijanja}\quad{\footnotesize\texttt{TERM-0390}}\par
\noindent\emph{Iztočnik / Source:} ein- und zweimalige Faltungen\par
\noindent\emph{Rusky / Russian:} одно- и двукратные свертки\par
\noindent\emph{Ukrajinsky / Ukrainian:} одно- і дворазові згортання\par
\smallskip\noindent\emph{Razlog:} Jedno- i dvojnokratna svijanja zachovuje dva točne čisla operacij v kratkoj paralelnoj konstrukciji. Prva čest s crticoju deli svoje poslednje nazvanje so vtoroju. Izvorni razlog daje prednost takvoj kompaktnosti i slovjanskoj prozračnosti, ale ne daje rezultat sravniteljnogo ispitovanja publiky.\par
\noindent\emph{Reason:} Jedno- i dvojnokratna svijanja retains two exact operation counts in a compact parallel construction. The hyphenated first component shares the following name with the second. The reason favours compactness and Slavic transparency but supplies no comparative audience-test result.\par
\noindent{\footnotesize Izvorne zapisy / Original records: HIST-g023-010.}\par
\end{minipage}\par
\par\medskip\Needspace{7\baselineskip}
\noindent\begin{minipage}{\linewidth}
\hypertarget{TERM-0391}{}
\noindent\textbf{vzaimno jednoznačno}\quad{\footnotesize\texttt{TERM-0391}}\par
\noindent\emph{Iztočnik / Source:} ein-eindeutig\par
\noindent\emph{Rusky / Russian:} взаимно однозначно\par
\noindent\emph{Ukrajinsky / Ukrainian:} взаємно однозначно\par
\smallskip\noindent\emph{Razlog:} Vzaimno jednoznačno zachovuje odgovaranje v obeh směrah: vsakomu elementu jednoj strany odgovarjaje točno jedin element drugoj strany, i obratno. Dva razloga povezuju izraz s bijekcijeju, od togo jedin s definicijeju izomorfizma. Samo jednoznačno ne vyražaje vse to. Tvrdženje zapisa, že opisna fraza jest točnėjša od bijektivno, ne ustanovjaje matematičnu razliku medžu pravilno definovanymi sinonimami; jde o čitatelsku prozračnost. Zapis priznaje potrěbu prověrky idiomatičnosti, ne dokazuje vsegslovjansku razumělivost.\par
\noindent\emph{Reason:} Vzaimno jednoznačno retains a correspondence in both directions: each element on either side corresponds to exactly one on the other. Both rationales connect the expression with bijectivity, one specifically in the definition of isomorphism. Jednoznačno alone does not express all of this. The record's claim that the descriptive phrase is more precise than bijective establishes no mathematical difference between properly defined synonyms; the issue is transparency to readers. The record acknowledges a need to examine idiom, rather than demonstrating pan-Slavic intelligibility.\par
\noindent{\footnotesize Izvorne zapisy / Original records: HIST-g028-012, HIST-g066-004.}\par
\end{minipage}\par
\par\medskip\Needspace{7\baselineskip}
\noindent\begin{minipage}{\linewidth}
\hypertarget{TERM-0392}{}
\noindent\textbf{vzaimno jednoznačno prirędženy}\quad{\footnotesize\texttt{TERM-0392}}\par
\noindent\emph{Iztočnik / Source:} ein-eindeutig zugeordnet\par
\noindent\emph{Rusky / Russian:} взаимно однозначно сопоставлена\par
\noindent\emph{Ukrajinsky / Ukrainian:} взаємно однозначно зіставлено\par
\smallskip\noindent\emph{Razlog:} Vzaimno jednoznačno prirędženy imenuje ne samo jedinstvenost, ale to, že objekty sut prirędžene v obeh směrah. Razlog izbora jest jasny opis bijektivnogo odgovaranja vměsto bukvalnogo prěnosa německogo zdvojenja ein-ein. Ruska i ukrajinska fraza davajut sravnjenje konstrukcije. Zapis ne dava samostalny slovnikovy iztočnik za prirędženy ani doklad, že imenovane alternativy byly čitatelsky testovane.\par
\noindent\emph{Reason:} Vzaimno jednoznačno prirędženy names not merely uniqueness but an assignment between objects in both directions. The stated reason is a clear description of bijective correspondence rather than a literal reproduction of German ein-ein. The Russian and Ukrainian phrases provide structural comparisons. The record supplies neither a separate lexical source for prirędženy nor evidence that the alternatives were tested with readers.\par
\noindent{\footnotesize Izvorne zapisy / Original records: HIST-g030-004.}\par
\end{minipage}\par
\par\medskip\Needspace{7\baselineskip}
\noindent\begin{minipage}{\linewidth}
\hypertarget{TERM-0393}{}
\noindent\textbf{vzaimno jednoznačno odobraženje}\quad{\footnotesize\texttt{TERM-0393}}\par
\noindent\emph{Iztočnik / Source:} Ein-eindeutige Abbildung\par
\noindent\emph{Rusky / Russian:} взаимно однозначное отображение\par
\noindent\emph{Ukrajinsky / Ukrainian:} взаємно однозначне відображення\par
\smallskip\noindent\emph{Razlog:} Odobraženje dodavaje ime objekta k frazě vzaimno jednoznačno. Razlog odkazuje na ranějše upotrěby bijektivnogo odgovaranja v dělah 3–5 i na potrěbu govoriti o odobraženju sistemov. Takva doslědnost podderžuje edin izraz v prěvodu; sama po sobě ne jest novym jezyčnym dokazom. Forma prědavava bijektivnost, ne samo to, že odobraženje ima točno oprěděljene značenja.\par
\noindent\emph{Reason:} Odobraženje adds the name of the object to vzaimno jednoznačno. The rationale refers to earlier correspondence terminology in Papers 3–5 and to the need to name a mapping of systems. This continuity supports consistent usage within the translation; it is not, by itself, fresh linguistic evidence. The expression conveys bijectivity, not merely that a mapping has well-defined values.\par
\noindent{\footnotesize Izvorne zapisy / Original records: HIST-g043-001.}\par
\end{minipage}\par
\par\medskip\Needspace{7\baselineskip}
\noindent\begin{minipage}{\linewidth}
\hypertarget{TERM-0394}{}
\noindent\textbf{jednoznačno oprěděljene / mnogoznačnost}\quad{\footnotesize\texttt{TERM-0394}}\par
\noindent\emph{Iztočnik / Source:} eindeutig bestimmt / Vieldeutigkeit\par
\noindent\emph{Rusky / Russian:} однозначно определены / неоднозначность\par
\noindent\emph{Ukrajinsky / Ukrainian:} однозначно визначені / неоднозначність\par
\smallskip\noindent\emph{Razlog:} Jednoznačno oprěděljene i mnogoznačnost razděljujut dva slučaja v teoremu XIII: izolovane komponenty sut jedinstveno oprěděljene, a nejednoznačnost ostavaje pri neizolovanyh komponentah. Slučaje ne možno zameniti slabšim tvrdženjem, že komponenty možno izbrati. Dva zapisa zachovuju jedin razlog. Mnogoznačnost tuta govori o nejedinstvenosti razklada, ne o mnogoznačnom slově ili o nejasnosti dokaza.\par
\noindent\emph{Reason:} Jednoznačno oprěděljene and mnogoznačnost separate two cases in Theorem XIII: isolated components are uniquely determined, while nonuniqueness is confined to non-isolated components. This cannot be weakened to saying that components can be chosen. Two records preserve one rationale. Here mnogoznačnost concerns nonuniqueness of a decomposition, not a word with several senses or an unclear proof.\par
\noindent{\footnotesize Izvorne zapisy / Original records: HIST-g081-062, HIST-g082-062.}\par
\end{minipage}\par
\par\medskip\Needspace{7\baselineskip}
\noindent\begin{minipage}{\linewidth}
\hypertarget{TERM-0395}{}
\noindent\textbf{jednoznačno byti predstavjen}\quad{\footnotesize\texttt{TERM-0395}}\par
\noindent\emph{Iztočnik / Source:} eindeutig darstellen\par
\noindent\emph{Rusky / Russian:} единственным образом представляется\par
\noindent\emph{Ukrajinsky / Ukrainian:} єдиним способом подати\par
\smallskip\noindent\emph{Razlog:} Jednoznačno byti predstavjen zachovuje dva obvezateljne elementy: predstavjenje postoji i jest jedinstveno v smyslu danogo tvrdženja. Razlog izbira obšču osnovu predstav- takože za slědujuče produktne i razkladne vyraženja. Takva svęz medžu izrazami ne dava prava vynechati množenje, additivnost ili druge uslovja tam, kde iztočnik ih ustanovjaje. Ruska i ukrajinska sravnjenja vyražaju jedinstveny sposob, ne prosto jasnost opisa.\par
\noindent\emph{Reason:} Jednoznačno byti predstavjen retains two necessary elements: a representation exists and is unique in the sense of the particular statement. The rationale chooses the predstav- family for use also in subsequent product and decomposition expressions. This lexical continuity does not license omitting multiplication, additivity, or other qualifications where the source specifies them. The Russian and Ukrainian comparanda express a unique manner of representation, not simply clarity of description.\par
\noindent{\footnotesize Izvorne zapisy / Original records: HIST-g150-003.}\par
\end{minipage}\par
\par\medskip\Needspace{7\baselineskip}
\noindent\begin{minipage}{\linewidth}
\hypertarget{TERM-0396}{}
\noindent\textbf{jednozna?no opred?ljeno preslikanje / obratno jednozna?no}\quad{\footnotesize\texttt{TERM-0396}}\par
\noindent\emph{Iztočnik / Source:} eindeutige Abbildung / umkehrbar eindeutig\par
\noindent\emph{Rusky / Russian:} однозначно определенное отображение / взаимно однозначное\par
\noindent\emph{Ukrajinsky / Ukrainian:} цілком визначене відображення / взаємно однозначне\par
\smallskip\noindent\emph{Razlog:} Zapis razlikuje dobro oprěděljeno odobraženje od obratno jednoznačnogo odobraženja. Po razlogu, dokaz najprvo konstruuje odobraženja medžu grupami ostatkov, a potom osobno dokazuje obrativnost s pomoču konečnoj modulnoj bazy. Izraz vzaimno jednoznačno zato ne treba prědnesti v ranějši krok. Několike bukvy medžuslovjanskoj formy sut utratjene; tablica ih ne rekonstruuje po domyslu. Opis odobraženij kako podobnyh surjektivnym ne jest dokaz surjektivnosti.\par
\noindent\emph{Reason:} The record distinguishes a well-defined mapping from an invertibly one-to-one mapping. Its rationale says that the proof first constructs maps between residue groups, then separately proves invertibility using a finite module basis. Bijective terminology must therefore not be introduced into the earlier step. Several letters of the Interslavic form are lost; the table does not reconstruct them by conjecture. The rationale's description of the maps as surjective-looking is not evidence of surjectivity.\par
\noindent{\footnotesize Izvorne zapisy / Original records: HIST-g079-052.}\par
\end{minipage}\par
\par\medskip\Needspace{7\baselineskip}
\noindent\begin{minipage}{\linewidth}
\hypertarget{TERM-0397}{}
\noindent\textbf{jednoznačny additivny razklad grupy ostatkov}\quad{\footnotesize\texttt{TERM-0397}}\par
\noindent\emph{Iztočnik / Source:} eindeutige additive Zerlegung der Restgruppe\par
\noindent\emph{Rusky / Russian:} единственное аддитивное разложение группы остатков\par
\noindent\emph{Ukrajinsky / Ukrainian:} єдиний адитивний розклад групи лишків\par
\smallskip\noindent\emph{Razlog:} Jednoznačny additivny razklad imenuje jedinstveny zapis vsakoj klasy ostatkov kak summy iz dvoh komponentnyh sistemov. Razlog privezuje izraz k formuli (11) i naglašaje, že množenje klasov bylo izključeno v oddělu 3. Additivny jest zato matematično suščestvenny, ne izbytočny pridevnik. Fraza ne tvrdži jedinstvenost vsakoj vozmožnoj razkladnoj struktury grupy; ona ostavaje svęzana s danoj formuloju.\par
\noindent\emph{Reason:} Jednoznačny additivny razklad names the unique expression of each residue class as a sum from the two component systems. The rationale connects the wording with formula (11) and stresses that multiplication of residue classes was excluded in Section 3. Additive is therefore mathematically substantive, not a redundant adjective. The phrase does not assert uniqueness of every possible decomposition structure of the group; it remains tied to the stated formula.\par
\noindent{\footnotesize Izvorne zapisy / Original records: HIST-g079-022.}\par
\end{minipage}\par
\par\medskip\Needspace{7\baselineskip}
\noindent\begin{minipage}{\linewidth}
\hypertarget{TERM-0398}{}
\noindent\textbf{jednoznačnost}\quad{\footnotesize\texttt{TERM-0398}}\par
\noindent\emph{Iztočnik / Source:} eindeutige Bestimmtheit\par
\noindent\emph{Rusky / Russian:} единственность\par
\noindent\emph{Ukrajinsky / Ukrainian:} єдиність\par
\smallskip\noindent\emph{Razlog:} Jednoznačnost prědavava jedinstvenost objektov $R$ i $S$, ktora jest cěljom dokaza. Rusky i ukrajinsky izraz davajut imena jedinstvenosti, a medžuslovjansky izbor pokračuje ranějšu skupinu jednoznačno. Dva zapisa zachovuju jedin razlog. Česte pojavjenje slova v korpusu ne dokazuje, že vsako pojavjenje ima točno toj smysl; jedinstvenost objektov treba razlikovati od jednoznačnosti značenja funkcije.\par
\noindent\emph{Reason:} Jednoznačnost conveys the uniqueness of the objects $R$ and $S$ that the proof seeks to establish. The Russian and Ukrainian forms name uniqueness, while the Interslavic choice continues the existing jednoznačno family. Two records preserve one rationale. Frequent occurrence in the corpus does not show that every occurrence has this precise sense; uniqueness of objects must be distinguished from single-valuedness of a function.\par
\noindent{\footnotesize Izvorne zapisy / Original records: HIST-g086-010, HIST-g092-020.}\par
\end{minipage}\par
\par\medskip\Needspace{7\baselineskip}
\noindent\begin{minipage}{\linewidth}
\hypertarget{TERM-0399}{}
\noindent\textbf{jednoznačno produktno predstavjenje črez nerazložime polinomy oblasti}\quad{\footnotesize\texttt{TERM-0399}}\par
\noindent\emph{Iztočnik / Source:} eindeutige Produktdarstellung durch irreduzible Polynome des Bereichs\par
\noindent\emph{Rusky / Russian:} однозначное произведенное представление через неприводимые полиномы области\par
\noindent\emph{Ukrajinsky / Ukrainian:} єдине добуткове подання через незвідні поліноми області\par
\smallskip\noindent\emph{Razlog:} Produktno predstavjenje zachovuje množenje nerazložimyh polinomov oblasti kako izričny predmet tvrdženja. Razlog povezuje frazu s vtorym priměrom i s teorijeju eliminacije; skračenje do razklada moglo by poměšati razklad polinoma na faktory s razkladom ideala. Jednoznačnost ne daje bez dodatnyh uslovij doslovnu ravnost vsakoj zapisanoj listy faktorov. Pravila o poredku ili jedinicah moraju izhoditi iz matematičnogo teksta, ne iz samogo slovničnogo nazvanja.\par
\noindent\emph{Reason:} Produktno predstavjenje keeps multiplication of irreducible polynomials in the domain as the explicit subject of the statement. The rationale links the phrase to the second example and elimination theory; shortening it to decomposition could conflate polynomial factorization with ideal decomposition. Uniqueness does not, without further qualifications, mean literal equality of every written list of factors. Conventions about order or units must come from the mathematical text, not from the dictionary label alone.\par
\noindent{\footnotesize Izvorne zapisy / Original records: HIST-g081-087.}\par
\end{minipage}\par
\par\medskip\Needspace{7\baselineskip}
\noindent\begin{minipage}{\linewidth}
\hypertarget{TERM-0400}{}
\noindent\textbf{jednoznačno razloženje}\quad{\footnotesize\texttt{TERM-0400}}\par
\noindent\emph{Iztočnik / Source:} eindeutige Zerlegbarkeit\par
\noindent\emph{Rusky / Russian:} единственное разложение\par
\noindent\emph{Ukrajinsky / Ukrainian:} єдиний розклад\par
\smallskip\noindent\emph{Razlog:} Jednoznačno razloženje jest izbrano v kontekstu zakonov razklada na faktory. Razlog izrično govori, že ti zakony mogut ne važiti v $J$. Slovnična oznaka zato ne može proměniti negativno ili uslovno tvrdženje v obšče tvrdženje o jedinstvenoj faktorizaciji. Německo Zerlegbarkeit imenuje svojstvo možnosti razklada, kdežto razloženje može imenovati i rezultat; točna rečenica mora zachovati razlikovanje.\par
\noindent\emph{Reason:} Jednoznačno razloženje is chosen in the context of factorization laws. The rationale explicitly says that these laws may fail in $J$. The dictionary label must therefore not turn a negative or conditional statement into a general assertion of unique factorization. German Zerlegbarkeit names a property of decomposability, whereas razloženje can also name the resulting decomposition; the surrounding sentence must preserve the distinction.\par
\noindent{\footnotesize Izvorne zapisy / Original records: HIST-g048-007.}\par
\end{minipage}\par
\par\medskip\Needspace{7\baselineskip}
\noindent\begin{minipage}{\linewidth}
\hypertarget{TERM-0401}{}
\noindent\textbf{jednoznačna sootvětnost / jednoznačnost}\quad{\footnotesize\texttt{TERM-0401}}\par
\noindent\emph{Iztočnik / Source:} eindeutige Zuordnung / Eindeutigkeit\par
\noindent\emph{Rusky / Russian:} однозначное соответствие / однозначность\par
\noindent\emph{Ukrajinsky / Ukrainian:} однозначна відповідність / однозначність\par
\smallskip\noindent\emph{Razlog:} Jednoznačna sootvětnost imenuje jednoznačno zadano odgovaranje prěd dokazom injektivnosti. Razlog zato izrično odbacaje prědrěmenny prěvod slovom izomorfizm. Jednoznačnost rezultata za danu vhodnu veličinu sama ne znači, že različnym vhodam odgovarjaju različne izhody. Ruska i ukrajinska paralela podpiraju to razlikovanje; zapis ne razsudžaje leksičnu prednost sootvětnost prěd odgovaranje.\par
\noindent\emph{Reason:} Jednoznačna sootvětnost names a single-valued assignment before injectivity has been proved. The rationale therefore explicitly rejects prematurely calling it an isomorphism. A unique output for each input does not itself imply different outputs for different inputs. The Russian and Ukrainian comparanda support that distinction; the record does not settle a lexical preference for sootvětnost over odgovaranje.\par
\noindent{\footnotesize Izvorne zapisy / Original records: HIST-g069-003.}\par
\end{minipage}\par
\par\medskip\Needspace{7\baselineskip}
\noindent\begin{minipage}{\linewidth}
\hypertarget{TERM-0402}{}
\noindent\textbf{teorema jednoznačnosti i uporjadženja}\quad{\footnotesize\texttt{TERM-0402}}\par
\noindent\emph{Iztočnik / Source:} Eindeutigkeits- und Anordnungssatz\par
\noindent\emph{Rusky / Russian:} теорема единственности и упорядочения\par
\noindent\emph{Ukrajinsky / Ukrainian:} теорема однозначності й упорядкування\par
\smallskip\noindent\emph{Razlog:} Nazvanje teorema jednoznačnosti i uporjadženja zachovuje oba člena německogo složenogo nazvanja. Razlog povezuje prvo s bijektivnym odgovaranjem, a vtoro s poredkom po vključjenju. Kratko nazvanje teorema jednoznačnosti izgubilo by vtory rezultat. To ne znači, že sam imenik jednoznačnost v vsakom drugom kontekstu uključaje bijektivnost.\par
\noindent\emph{Reason:} Teorema jednoznačnosti i uporjadženja retains both parts of the German compound title. The rationale relates the first to bijective correspondence and the second to the inclusion order. Shortening the title to uniqueness theorem would lose the second component. This does not mean that the noun jednoznačnost itself includes bijectivity in every other context.\par
\noindent{\footnotesize Izvorne zapisy / Original records: HIST-g208-017.}\par
\end{minipage}\par
\par\medskip\Needspace{7\baselineskip}
\noindent\begin{minipage}{\linewidth}
\hypertarget{TERM-0404}{}
\noindent\textbf{prosty element vhodi v diskriminant}\quad{\footnotesize\texttt{TERM-0404}}\par
\noindent\emph{Iztočnik / Source:} Eine Primgröße p ... geht ... in der Diskriminante auf\par
\noindent\emph{Rusky / Russian:} простой элемент входит в дискриминант\par
\noindent\emph{Ukrajinsky / Ukrainian:} простий елемент входить до дискримінанта\par
\smallskip\noindent\emph{Razlog:} Prosty element vhodi v diskriminant zachovuje obraz vhoda iz uvoda raboty. Izvorni razlog dava mu prednost prěd samym děli, aby ne izčezl avtorsky algebraičny idiom. Tuta ne jde o vključjenje elementa v množinu zvanu diskriminant: matematičny smysl ostavaje svęzany s děliteljnostju. Zapis dava razlog stilističnogo izbora, ale ne měrenje razumělivosti toj frazy.\par
\noindent\emph{Reason:} Prosty element vhodi v diskriminant retains the introduction's enters idiom. The stated reason prefers it to simply divides in order to preserve the author's algebraic manner of expression. This is not set membership in a set called the discriminant: the mathematical sense remains connected with divisibility. The record provides a stylistic rationale, not a measurement of the phrase's intelligibility.\par
\noindent{\footnotesize Izvorne zapisy / Original records: HIST-g153-002.}\par
\end{minipage}\par
\par\medskip\Needspace{7\baselineskip}
\noindent\begin{minipage}{\linewidth}
\hypertarget{TERM-0405}{}
\noindent\textbf{ręd jest razložen}\quad{\footnotesize\texttt{TERM-0405}}\par
\noindent\emph{Iztočnik / Source:} eine Reihe zerlegt\par
\noindent\emph{Rusky / Russian:} строка разложена\par
\noindent\emph{Ukrajinsky / Ukrainian:} рядок розкладено\par
\smallskip\noindent\emph{Razlog:} Ręd jest razložen govori o razkladu reda, ne o faktorizaciji skalara ili polinoma. Razlog nastavlja imena redov iz raboty 4. Rusko строка i ukrajinsko рядок podpiraju točny tip objekta i upozornjaju, že německo Reihe tuta ne trěbuje avtomatičny prěvod kako beskonečny red. Sam glagol razložiti bez svojego subjekta ne bi zachoval toj razlike.\par
\noindent\emph{Reason:} Ręd jest razložen concerns decomposition of a row, not factorization of a scalar or polynomial. The rationale continues the row terminology of Paper 4. Russian строка and Ukrainian рядок support the object type and caution against automatically reading German Reihe as an infinite series. The verb decompose alone, without its subject, would not preserve that distinction.\par
\noindent{\footnotesize Izvorne zapisy / Original records: HIST-g033-003.}\par
\end{minipage}\par
\par\medskip\Needspace{7\baselineskip}
\noindent\begin{minipage}{\linewidth}
\hypertarget{TERM-0406}{}
\noindent\textbf{jedno-jednoznačno}\quad{\footnotesize\texttt{TERM-0406}}\par
\noindent\emph{Iztočnik / Source:} eineindeutig\par
\noindent\emph{Rusky / Russian:} взаимно однозначно\par
\noindent\emph{Ukrajinsky / Ukrainian:} взаємно однозначно\par
\smallskip\noindent\emph{Razlog:} Jedno-jednoznačno vyražaje bijektivno odgovaranje medžu prinaležečimi prostymi objektami i največšimi primarnymi normovymi faktorami. Dvojna forma naměrava učiniti odnos v obeh směrah vidimym bez dolge dodatne frazy. Zapis ne dava čislenu ocěnku ani sravnjenje s vzaimno jednoznačno. Povtorjenje jedne osnovy jest zabilěžena tvorba, ne dokaz, že ona jest obvezateljna jezyčna norma.\par
\noindent\emph{Reason:} Jedno-jednoznačno expresses the bijective correspondence between the associated primes and largest primary norm factors. The doubled form is intended to make the two-way relation explicit without a long additional phrase. The record supplies neither a numerical rating nor a comparison with vzaimno jednoznačno. The repetition of a stem is a recorded formation, not evidence that it is a binding linguistic norm.\par
\noindent{\footnotesize Izvorne zapisy / Original records: HIST-g099-008.}\par
\end{minipage}\par
\par\medskip\Needspace{7\baselineskip}
\noindent\begin{minipage}{\linewidth}
\hypertarget{TERM-0407}{}
\noindent\textbf{vzaimno jednoznačny}\quad{\footnotesize\texttt{TERM-0407}}\par
\noindent\emph{Iztočnik / Source:} eineindeutig\par
\noindent\emph{Rusky / Russian:} взаимно однозначный\par
\noindent\emph{Ukrajinsky / Ukrainian:} взаємно однозначний\par
\smallskip\noindent\emph{Razlog:} Vzaimno jednoznačny jest pridevni oblik za odnos v obeh směrah. On jest naměrno odlišeny od samoj jednoznačnosti v toj že rabotě; ne treba izbrisati vzaimno pri skračenju. Ruska i ukrajinska pridevna oblika dajut neposrědne sravniteljne formy. Označenje termina kako standardnogo v zapisu ne zaměnjaje točny citat jezyčnogo iztočnika.\par
\noindent\emph{Reason:} Vzaimno jednoznačny is the adjectival form for a two-way correspondence. It is deliberately distinguished from mere jednoznačnost within the same work; vzaimno must not disappear through shortening. The Russian and Ukrainian adjectives supply direct comparanda. Calling the expression standard in the record does not substitute for a precise linguistic citation.\par
\noindent{\footnotesize Izvorne zapisy / Original records: HIST-g069-005.}\par
\end{minipage}\par
\par\medskip\Needspace{7\baselineskip}
\noindent\begin{minipage}{\linewidth}
\hypertarget{TERM-0408}{}
\noindent\textbf{vzaimno-jednoznačno}\quad{\footnotesize\texttt{TERM-0408}}\par
\noindent\emph{Iztočnik / Source:} eineindeutig\par
\noindent\emph{Rusky / Russian:} взаимно однозначно\par
\noindent\emph{Ukrajinsky / Ukrainian:} взаємно однозначно\par
\smallskip\noindent\emph{Razlog:} Vzaimno-jednoznačno svęzuje dvustranne razklady s razkladami centra. Razlog izrično uključaje oba směry odgovaranja, ne samo jedinstvenost objekta na jednoj straně. Spojnica odlišaje zapis od vzaimno jednoznačno, ale ne ustanovjaje novy matematičny pojem. Njezina pravopisna prednost ne jest odvojeno argumentovana.\par
\noindent\emph{Reason:} Vzaimno-jednoznačno connects two-sided decompositions with decompositions of the centre. The reason explicitly includes both directions of correspondence, not merely uniqueness on one side. The hyphen distinguishes the spelling from vzaimno jednoznačno without creating a new mathematical concept. No separate argument for that orthographic preference is recorded.\par
\noindent{\footnotesize Izvorne zapisy / Original records: HIST-g189-005.}\par
\end{minipage}\par
\par\medskip\Needspace{7\baselineskip}
\noindent\begin{minipage}{\linewidth}
\hypertarget{TERM-0409}{}
\noindent\textbf{jednoznačno vzajemno ... razlika i množenje tym samym elementom}\quad{\footnotesize\texttt{TERM-0409}}\par
\noindent\emph{Iztočnik / Source:} eineindeutig ... Differenz und Multiplikation mit demselben Element\par
\noindent\emph{Rusky / Russian:} взаимно однозначно ... разность и умножение на один и тот же элемент\par
\noindent\emph{Ukrajinsky / Ukrainian:} взаємно однозначно ... різниця і множення на той самий елемент\par
\smallskip\noindent\emph{Razlog:} Fraza zachovuje i bijektivnost, i soglasje s razlikoju i množenjem tym samym elementom. Razlog izrično dava uslovje: iz odgovaranja $\alpha\sim\beta$ slěduje $\gamma\alpha\sim\gamma\beta$. Poredok faktorov ne možno zaměniti. Bijekcija sama po sobě ne jest operatorny izomorfizm. Pozvanje na pečatny iztočnik jest historijna dokazova tvrdnja zapisa, ne nova kolacija izvršena pri pisanju togo objasnjenja.\par
\noindent\emph{Reason:} The phrase preserves both bijectivity and compatibility with difference and multiplication by the same element. The rationale explicitly gives the condition: $\alpha\sim\beta$ implies $\gamma\alpha\sim\gamma\beta$. The order of the factors must not be reversed. A bijection alone is not an operator isomorphism. The appeal to the printed witness is the record's historical evidential claim, not a fresh collation performed for this explanation.\par
\noindent{\footnotesize Izvorne zapisy / Original records: HIST-g133-004.}\par
\end{minipage}\par
\par\medskip\Needspace{7\baselineskip}
\noindent\begin{minipage}{\linewidth}
\hypertarget{TERM-0410}{}
\noindent\textbf{jedno-jednoznačno / jedno-jednoznačno pridružiti}\quad{\footnotesize\texttt{TERM-0410}}\par
\noindent\emph{Iztočnik / Source:} eineindeutig / eineindeutig zuordnen\par
\noindent\emph{Rusky / Russian:} взаимно однозначно / взаимно однозначно сопоставить\par
\noindent\emph{Ukrajinsky / Ukrainian:} взаємно однозначно / взаємно однозначно зіставити\par
\smallskip\noindent\emph{Razlog:} Jedno-jednoznačno pridružiti govori o bijektivnom prirědženju komponentov v dvoh redukovanyh prědstavjenjah. Samo sravniti bylo by slabše tvrdženje. Dva zapisa zachovuju jedin razlog, a ruska i ukrajinska fraza vyražaju vzajemnu jednoznačnost. Razlog opravdaje nužny smysl prirědženja, ale ne dokazuje, že dvojna forma jest boljša od vseh drugyh medžuslovjanskyh konstrukcij.\par
\noindent\emph{Reason:} Jedno-jednoznačno pridružiti concerns a bijective assignment of components in two reduced presentations. Merely compare would be a weaker claim. Two records preserve one rationale, and the Russian and Ukrainian phrases express mutual one-to-one correspondence. The rationale establishes the required sense of the assignment, not the superiority of the doubled form over every other Interslavic construction.\par
\noindent{\footnotesize Izvorne zapisy / Original records: HIST-g081-048, HIST-g082-048.}\par
\end{minipage}\par
\par\medskip\Needspace{7\baselineskip}
\noindent\begin{minipage}{\linewidth}
\hypertarget{TERM-0411}{}
\noindent\textbf{jedno-jednozna?ny / jedno-jednozna?no prir??enje}\quad{\footnotesize\texttt{TERM-0411}}\par
\noindent\emph{Iztočnik / Source:} eineindeutig / eineindeutige Zuordnung\par
\noindent\emph{Rusky / Russian:} [Bukvy ne sut zachovane / Lettering not recoverable]\par
\noindent\emph{Ukrajinsky / Ukrainian:} [Bukvy ne sut zachovane / Lettering not recoverable]\par
\smallskip\noindent\emph{Razlog:} Izvorni razlog trěbuje injektivnost i obrativnost prirědženja medžu sistemami klasov ostatkov. Dvojno jedno- imaje odlišiti to od prostej jedinstvenosti. Medžuslovjansky zapis ima utratjene bukvy, a ruske i ukrajinske formy ne sut zachovane čitatelno. Razlog smyslovogo razlikovanja ostavaje dostupny, ale iz njega ne možno izmysliti točne historijne grafije.\par
\noindent\emph{Reason:} The rationale requires injectivity and reversibility of the assignment between residue-class systems. Doubled jedno- is intended to distinguish this from simple uniqueness. Letters are missing in the Interslavic wording, and the Russian and Ukrainian forms have not survived legibly. The semantic reason remains recoverable, but it cannot establish the exact historical spellings.\par
\noindent{\footnotesize Izvorne zapisy / Original records: HIST-g079-040.}\par
\end{minipage}\par
\par\medskip\Needspace{7\baselineskip}
\noindent\begin{minipage}{\linewidth}
\hypertarget{TERM-0412}{}
\noindent\textbf{odgovarjajut ... jednoznačno vzajemno}\quad{\footnotesize\texttt{TERM-0412}}\par
\noindent\emph{Iztočnik / Source:} eineindeutig entsprechen\par
\noindent\emph{Rusky / Russian:} взаимно однозначно соответствуют\par
\noindent\emph{Ukrajinsky / Ukrainian:} відповідають ... взаємно однозначно\par
\smallskip\noindent\emph{Razlog:} Odgovarjajut ... jednoznačno vzajemno zachovuje glagol odgovaranja i obosměrnu jedinstvenost. Razlog izrično pokračuje zapis bez spojnice iz ranějšej pasaže raboty 31. To jest zabilěžena pravopisna i frazeologična doslědnost, ne drugi tip bijekcije. Sama doslědnost ne razrěšaje izbor medžu vzajemno i vzaimno v cělom prěvodu.\par
\noindent\emph{Reason:} Odgovarjajut ... jednoznačno vzajemno retains the verb of correspondence and uniqueness in both directions. The rationale explicitly continues the unhyphenated wording of an earlier passage in Paper 31. This is a recorded orthographic and phraseological preference, not a different kind of bijection. Consistency alone does not settle the corpus-wide choice between vzajemno and vzaimno.\par
\noindent{\footnotesize Izvorne zapisy / Original records: HIST-g134-001.}\par
\end{minipage}\par
\par\medskip\Needspace{7\baselineskip}
\noindent\begin{minipage}{\linewidth}
\hypertarget{TERM-0413}{}
\noindent\textbf{vzaimno jednoznačno sootvětujuča baza}\quad{\footnotesize\texttt{TERM-0413}}\par
\noindent\emph{Iztočnik / Source:} eineindeutig entsprechende Basis\par
\noindent\emph{Rusky / Russian:} взаимно однозначно соответствующий базис\par
\noindent\emph{Ukrajinsky / Ukrainian:} взаємно однозначно відповідна база\par
\smallskip\noindent\emph{Razlog:} Vzaimno jednoznačno sootvětujuča baza svezuje bijektivno odgovaranje s konkretnym objektom, bazoju. Kratky razlog odkazuje na ranějše dva odděly raboty 10. On podpiraje kontinuitet nazvanja, ale ne dava novy dokaz za pričestny oblik sootvětujuča. Izpustiti baza značilo by izgubiti tip objekta; zaměniti vzaimno jednoznačno samym podobna značilo by oslabiti odnos.\par
\noindent\emph{Reason:} Vzaimno jednoznačno sootvětujuča baza applies bijective correspondence to a specific object, the basis. The short rationale refers to the first two parts of Paper 10. It supports terminological continuity but gives no new evidence for the participial form sootvětujuča. Omitting baza would lose the object type; replacing vzaimno jednoznačno by similar would weaken the relation.\par
\noindent{\footnotesize Izvorne zapisy / Original records: HIST-g071-002.}\par
\end{minipage}\par
\par\medskip\Needspace{7\baselineskip}
\noindent\begin{minipage}{\linewidth}
\hypertarget{TERM-0414}{}
\noindent\textbf{jedno-jednoznačno i izomorfno odobraziti}\quad{\footnotesize\texttt{TERM-0414}}\par
\noindent\emph{Iztočnik / Source:} eineindeutig und isomorph abbilden\par
\noindent\emph{Rusky / Russian:} взаимно однозначно и изоморфно отобразить\par
\noindent\emph{Ukrajinsky / Ukrainian:} взаємно однозначно й ізоморфно відобразити\par
\smallskip\noindent\emph{Razlog:} Jedno-jednoznačno i izomorfno odobraziti zachovuje oba německe opisy odobraženja. Razlog izrično odbacaje skračenje do sovrěmennogo izomorfne kako prsteny. Zadača jest věrno podati avtorsku formulaciju, ne ukloniti ono, što v pozdnějšej definiciji može izgledati izbytočno. O tom, kakve operacije sohranjaju se, odlučuje iztočna definicija, ne samo slovo izomorfno.\par
\noindent\emph{Reason:} Jedno-jednoznačno i izomorfno odobraziti retains both German descriptions of the mapping. The rationale explicitly rejects replacing them with the modern shorthand isomorphic as rings. The purpose is to reproduce the author's formulation, not to delete what a later definition might make redundant. Which operations are preserved is determined by the source definition, not by the word izomorfno alone.\par
\noindent{\footnotesize Izvorne zapisy / Original records: HIST-g058-013.}\par
\end{minipage}\par
\par\medskip\Needspace{7\baselineskip}
\noindent\begin{minipage}{\linewidth}
\hypertarget{TERM-0415}{}
\noindent\textbf{jedno-jednoznačno pridruženje / izomorfny}\quad{\footnotesize\texttt{TERM-0415}}\par
\noindent\emph{Iztočnik / Source:} eineindeutige Zuordnung / isomorph\par
\noindent\emph{Rusky / Russian:} взаимно однозначное сопоставление / изоморфный\par
\noindent\emph{Ukrajinsky / Ukrainian:} взаємно однозначна відповідність / ізоморфний\par
\smallskip\noindent\emph{Razlog:} Jedno-jednoznačno pridruženje jest čestju definicije izomorfizma sistemov klasov ostatkov, v ktorej treba zachovati addiciju i dějstvo polinomov. Ne dostačaje prosto spočtati objekty na obeh stranah. Zapis naměrava prozračnu frazu, ale ostavjaje njezin idiom za sravnjenje. Takva otvorjena ocěnka ne može byti predstavljena kako potvrđeno čitatelsko prěimuščestvo.\par
\noindent\emph{Reason:} Jedno-jednoznačno pridruženje belongs to a definition of isomorphism for systems of residue classes that requires preservation of addition and polynomial action. Merely counting the objects on either side is insufficient. The record aims at a transparent expression but leaves its idiom open to comparison. That provisional judgment cannot be presented as a demonstrated advantage for readers.\par
\noindent{\footnotesize Izvorne zapisy / Original records: HIST-g079-023.}\par
\end{minipage}\par
\par\medskip\Needspace{7\baselineskip}
\noindent\begin{minipage}{\linewidth}
\hypertarget{TERM-0416}{}
\noindent\textbf{jedno-jednoznačno odgovarjanje}\quad{\footnotesize\texttt{TERM-0416}}\par
\noindent\emph{Iztočnik / Source:} eineindeutiges Entsprechen\par
\noindent\emph{Rusky / Russian:} взаимно однозначное соответствие\par
\noindent\emph{Ukrajinsky / Ukrainian:} взаємно однозначна відповідність\par
\smallskip\noindent\emph{Razlog:} Jedno-jednoznačno odgovarjanje imenuje odnos medžu děliteljami i odgovarjajučimi idealami pri prěhodu modulo prosty ideal. Razlog dava prednost izričnoj konstrukciji prěd skračenim nazvanjem. On ne nabraja alternativne kratke termine i ne dava težiny za njih. Zato slovnik objasnjaje argument za eksplicitnost, ne izmysljaje nezachovano natěkanje kandidatov.\par
\noindent\emph{Reason:} Jedno-jednoznačno odgovarjanje names the correspondence between divisors and the associated ideals in passing modulo a prime ideal. The rationale prefers an explicit expression to a compressed name. It lists neither alternative short terms nor weights for them. The dictionary therefore explains the argument for explicitness without inventing a lost candidate comparison.\par
\noindent{\footnotesize Izvorne zapisy / Original records: HIST-g114-012.}\par
\end{minipage}\par
\par\medskip\Needspace{7\baselineskip}
\noindent\begin{minipage}{\linewidth}
\hypertarget{TERM-0417}{}
\noindent\textbf{jednoznačno vzajemna sootvětnost}\quad{\footnotesize\texttt{TERM-0417}}\par
\noindent\emph{Iztočnik / Source:} eineindeutiges Entsprechen\par
\noindent\emph{Rusky / Russian:} взаимно однозначное соответствие\par
\noindent\emph{Ukrajinsky / Ukrainian:} взаємно однозначна відповідність\par
\smallskip\noindent\emph{Razlog:} Jednoznačno vzajemna sootvětnost jest zabilěžena fraza za bijektivno odgovaranje, postavjena pokraj izomorfizma prstenov klasov ostatkov. Ta dva izraza ne sut prosti sinonimi: vtory imenuje strukturno soglasje, ne samo odgovaranje množin. Razlog govori o sovrěmennom idiomu, ale ne podpiraje osobnym citatom slovo sootvětnost ni toj poredok slov.\par
\noindent\emph{Reason:} Jednoznačno vzajemna sootvětnost is the recorded phrase for bijective correspondence, paired with isomorphism of residue-class rings. The expressions are not simply synonyms: the latter names structural compatibility, not just a correspondence of sets. The rationale invokes modern idiom but supplies no separate citation for sootvětnost or this word order.\par
\noindent{\footnotesize Izvorne zapisy / Original records: HIST-g159-008.}\par
\end{minipage}\par
\par\medskip\Needspace{7\baselineskip}
\noindent\begin{minipage}{\linewidth}
\hypertarget{TERM-0418}{}
\noindent\textbf{prosta grupa}\quad{\footnotesize\texttt{TERM-0418}}\par
\noindent\emph{Iztočnik / Source:} einfach\par
\noindent\emph{Rusky / Russian:} простая группа\par
\noindent\emph{Ukrajinsky / Ukrainian:} проста група\par
\smallskip\noindent\emph{Razlog:} Prosta grupa prědavava tehničny smysl einfach v grupovom kontekstu, ne běžno značenje legka. Ruska i ukrajinska paralela imajut istu pridavničnu osnovu. Razlog zato ne vidi potrěbu tvoriti novy kandidat. To ne jest universalno pravilo prěvoditi vsako einfach ili vsaky prosty tojže osnovoj nezavisno od matematičnogo objekta.\par
\noindent\emph{Reason:} Prosta grupa renders technical einfach in the group context, not ordinary easy. The Russian and Ukrainian comparanda share the adjectival stem. The reason therefore sees no need to construct a new candidate. This is not a universal rule for every occurrence of einfach or prosty regardless of the mathematical object.\par
\noindent{\footnotesize Izvorne zapisy / Original records: HIST-g180-004.}\par
\end{minipage}\par
\par\medskip\Needspace{7\baselineskip}
\noindent\begin{minipage}{\linewidth}
\hypertarget{TERM-0419}{}
\noindent\textbf{prosta grupa}\quad{\footnotesize\texttt{TERM-0419}}\par
\noindent\emph{Iztočnik / Source:} einfache Gruppe\par
\noindent\emph{Rusky / Russian:} простая группа\par
\noindent\emph{Ukrajinsky / Ukrainian:} проста група\par
\smallskip\noindent\emph{Razlog:} Prosta grupa prodolžuje nazvanje prostoj grupy iz ranějšego odděla tojže raboty. Zapis govori o soglasju s ruskoju i ukrajinskoju terminologijeju; on ne opisuva novy slovotvorny postupok. Forma ostavaje odvojena od jednostavny ideal, ktory v drugom kontekstu služi razlikovanju prostosti ideala v dvoh smyslah.\par
\noindent\emph{Reason:} Prosta grupa continues the simple-group term from an earlier section of the same work. The record cites agreement with the Russian and Ukrainian terminology rather than a new word-formation procedure. It remains distinct from jednostavny ideal, which elsewhere serves to distinguish two meanings of simplicity or primality for ideals.\par
\noindent{\footnotesize Izvorne zapisy / Original records: HIST-g183-002.}\par
\end{minipage}\par
\par\medskip\Needspace{7\baselineskip}
\noindent\begin{minipage}{\linewidth}
\hypertarget{TERM-0420}{}
\noindent\textbf{jednostavne idealy}\quad{\footnotesize\texttt{TERM-0420}}\par
\noindent\emph{Iztočnik / Source:} einfache Ideale\par
\noindent\emph{Rusky / Russian:} элементарные идеалы\par
\noindent\emph{Ukrajinsky / Ukrainian:} елементарні ідеали\par
\smallskip\noindent\emph{Razlog:} Jednostavne idealy jest naměrny izbor za Noetherino uslovje, že ideal ne ima vlastnogo dělitelja v danom smyslu. Proste idealy už imenuju prime idealy, i dva razloga izrično trěbuju izběgnuti to směšenje. Ruska i ukrajinska paralela izbiraju elementarny, ne prosty. To pokazuje konkretnu strategiju razlikovanja smyslov, ne obšče izključjenje slova prosty iz vseh nazvanj prostyh algebraičnyh struktur.\par
\noindent\emph{Reason:} Jednostavne idealy deliberately names Noether's condition that an ideal have no proper divisor in the specified sense. Proste idealy is already used for prime ideals, and both rationales expressly seek to avoid that conflation. The Russian and Ukrainian comparanda choose elementary rather than simple. This is a context-specific strategy for separating senses, not a general exclusion of prosty from all names of simple algebraic structures.\par
\noindent{\footnotesize Izvorne zapisy / Original records: HIST-g116-001, HIST-g116-008.}\par
\end{minipage}\par
\par\medskip\Needspace{7\baselineskip}
\noindent\begin{minipage}{\linewidth}
\hypertarget{TERM-0421}{}
\noindent\textbf{prosta normalna algebra / dělitevna algebra}\quad{\footnotesize\texttt{TERM-0421}}\par
\noindent\emph{Iztočnik / Source:} einfache normale Algebra / Divisionsalgebra\par
\noindent\emph{Rusky / Russian:} простая нормальная алгебра / алгебра с делением\par
\noindent\emph{Ukrajinsky / Ukrainian:} проста нормальна алгебра / алгебра з діленням\par
\smallskip\noindent\emph{Razlog:} Prosta normalna algebra i dělitevna algebra zachovuju dva struktur nazyvajuče izraza iz raboty 42. Razlog prodolžuje izbor dělitevna algebra iz raboty 40. On ne govori, že vsaka prosta normalna algebra jest sama dělitevna, ni ne dava prava izpustiti normalna. Označenje izbora kako kanoničnogo opisuva projektnu doslědnost, ne nezavisnu jezyčnu ili matematičnu certifikaciju.\par
\noindent\emph{Reason:} Prosta normalna algebra and dělitevna algebra retain two structural expressions from Paper 42. The rationale continues the division-algebra terminology of Paper 40. It does not say that every simple normal algebra is itself a division algebra, or permit normalna to be omitted. Describing the choice as canonical records a project convention, not independent linguistic or mathematical certification.\par
\noindent{\footnotesize Izvorne zapisy / Original records: HIST-g212-013.}\par
\end{minipage}\par
\par\medskip\Needspace{7\baselineskip}
\noindent\begin{minipage}{\linewidth}
\hypertarget{TERM-0422}{}
\noindent\textbf{prosty modul}\quad{\footnotesize\texttt{TERM-0422}}\par
\noindent\emph{Iztočnik / Source:} einfacher Modul\par
\noindent\emph{Rusky / Russian:} простой модуль\par
\noindent\emph{Ukrajinsky / Ukrainian:} простий модуль\par
\smallskip\noindent\emph{Razlog:} Prosty modul jest izbran v jezyku kompozicijnyh redov. Izvorni razlog izrično razlikuje jeho od jednostavne idealy v prědhodnom oddělu 9. Pridevnik zatym ne može byti izbran samostatno od imenika: ranějše izběganje prosty pri jednom idealnom pojmu ne zabrana jego upotrěby pri modulah. Ruska i ukrajinska paralela zdes takože imaju prosty.\par
\noindent\emph{Reason:} Prosty modul is selected in composition-series terminology. The rationale explicitly distinguishes it from jednostavne idealy in the preceding Section 9. The adjective therefore cannot be chosen independently of its noun: avoiding prosty for one ideal-theoretic concept does not prohibit it for modules. The Russian and Ukrainian comparanda likewise use their simple-module terms.\par
\noindent{\footnotesize Izvorne zapisy / Original records: HIST-g117-004.}\par
\end{minipage}\par
\par\medskip\Needspace{7\baselineskip}
\noindent\begin{minipage}{\linewidth}
\hypertarget{TERM-0423}{}
\noindent\textbf{prosty lěvy ideal}\quad{\footnotesize\texttt{TERM-0423}}\par
\noindent\emph{Iztočnik / Source:} einfaches Linksideal\par
\noindent\emph{Rusky / Russian:} простой левый идеал\par
\noindent\emph{Ukrajinsky / Ukrainian:} простий лівий ідеал\par
\smallskip\noindent\emph{Razlog:} Prosty lěvy ideal zachovuje kratko nazvanje objekta, na ktorom zasnovana jest konstrukcija prědstavjenja. Lěvy jest obvezateljna strukturna kvalifikacija, ne oznaka prostorovogo položjenja. Razlog podpiraje svęz s ranějšimi lěvymi i pravymi idealami; on ne razrěšaje drugih smyslov prosty pri idealah izvan toga konteksta.\par
\noindent\emph{Reason:} Prosty lěvy ideal keeps a compact name for the object underlying the representation construction. Lěvy is a necessary structural qualification, not a spatial description. The rationale supports continuity with earlier left- and right-ideal terminology; it does not resolve other senses of prosty applied to ideals outside this context.\par
\noindent{\footnotesize Izvorne zapisy / Original records: HIST-g197-004.}\par
\end{minipage}\par
\par\medskip\Needspace{7\baselineskip}
\noindent\begin{minipage}{\linewidth}
\hypertarget{TERM-0424}{}
\noindent\textbf{jednoformny slučaj}\quad{\footnotesize\texttt{TERM-0424}}\par
\noindent\emph{Iztočnik / Source:} einförmiger Fall\par
\noindent\emph{Rusky / Russian:} одноформный случай\par
\noindent\emph{Ukrajinsky / Ukrainian:} одноформний випадок\par
\smallskip\noindent\emph{Razlog:} Jednoformny slučaj jest bukvalny kandidat za lokalno uslovje, svęzano s jednoju formoju ili prostym redom. Razlog dava prednost vidimoj strukturi německogo nazvanja prěd sovrěmennoju parafrazoju, ktora mogla by skryti toj lokalny pojem. Slučaj ne jest opisan kako prosto jednoobrazny stil; konkretny matematičny smysl treba čitati iz definicije raboty.\par
\noindent\emph{Reason:} Jednoformny slučaj is a literal candidate for a local condition involving one form or a simple row. The rationale prefers the visible structure of the German name to a modern paraphrase that might obscure the local concept. This is not merely a case of uniform style; its precise mathematical sense must be read from the work's definition.\par
\noindent{\footnotesize Izvorne zapisy / Original records: HIST-g056-005.}\par
\end{minipage}\par
\par\medskip\Needspace{7\baselineskip}
\noindent\begin{minipage}{\linewidth}
\hypertarget{TERM-0425}{}
\noindent\textbf{jednoelementny / mnogočlenny bazis modula}\quad{\footnotesize\texttt{TERM-0425}}\par
\noindent\emph{Iztočnik / Source:} eingliedrige / mehrgliedrige Modulbasis\par
\noindent\emph{Rusky / Russian:} одноэлементный / многоэлементный базис модуля\par
\noindent\emph{Ukrajinsky / Ukrainian:} одноелементний / багатоелементний базис модуля\par
\smallskip\noindent\emph{Razlog:} Jednoelementny i mnogočlenny imenuju čislo generatorov izbranoj modulnoj bazy. Razlog izrično odbacaje zaměnu gotovym sovrěmennym nazvanjem glavnogo modula, ibo predmetom sravnjenja jest sama baza. Imenovanje elementov ne daje dopolniteljne tvrdnje o ih linearnoj nezavisnosti. Različne osnovy element- i člen- jesut zabilěžene, ale njihov stilističny kontrast ne jest odvojeno opravdany.\par
\noindent\emph{Reason:} Jednoelementny and mnogočlenny name the number of generators in the chosen module basis. The rationale explicitly rejects substituting a modern principal-module shorthand because the basis itself is under discussion. Naming its members does not add an assertion of linear independence. The different stems element- and člen- are recorded, but their stylistic contrast is not separately justified.\par
\noindent{\footnotesize Izvorne zapisy / Original records: HIST-g079-073.}\par
\end{minipage}\par
\par\medskip\Needspace{7\baselineskip}
\noindent\begin{minipage}{\linewidth}
\hypertarget{TERM-0426}{}
\noindent\textbf{jedinica iz [Omega]}\quad{\footnotesize\texttt{TERM-0426}}\par
\noindent\emph{Iztočnik / Source:} Einheit aus [Omega]\par
\noindent\emph{Rusky / Russian:} единица из [Omega]\par
\noindent\emph{Ukrajinsky / Ukrainian:} одиниця з [Omega]\par
\smallskip\noindent\emph{Razlog:} Jedinica iz [Omega] povezuje naziv jedinicy s označenoju oblastju. Ediny zachovany razlog govori, gde izraz jest potrěbny: v definicijah V′ i VIII′ i v zaoštrjenom pomoćnom teoremu. On ne dava sravnjenje leksičnyh alternativ ani samostatnu definiciju jedinicy. Oznaka [Omega] ostavaje tako, kako jest zapisana; iz samej zaglavnoj formy ne možno ustanoviti njezin izvorni tipografski oblik.\par
\noindent\emph{Reason:} Jedinica iz [Omega] relates the unit term to a designated domain. The sole preserved rationale says where the expression is needed: Definitions V′ and VIII′ and the sharpened auxiliary theorem. It supplies neither a comparison of lexical alternatives nor a separate definition of unit. [Omega] remains as recorded; the headword alone does not determine its original typography.\par
\noindent{\footnotesize Izvorne zapisy / Original records: HIST-g055-003.}\par
\end{minipage}\par
\par\medskip\Needspace{7\baselineskip}
\noindent\begin{minipage}{\linewidth}
\hypertarget{TERM-0427}{}
\noindent\textbf{jedinice}\quad{\footnotesize\texttt{TERM-0427}}\par
\noindent\emph{Iztočnik / Source:} Einheiten\par
\noindent\emph{Rusky / Russian:} единицы\par
\noindent\emph{Ukrajinsky / Ukrainian:} одиниці\par
\smallskip\noindent\emph{Razlog:} Jedinice jest kratko nazvanje obrativyh elementov, postavjeno v kontrast s nulovymi děliteljami. Razlog izrično dava toj kratkoj formě prednost prěd dolgom frazoju obrativne elementy. To jest argument za ekonomiju nazvanja pri sohranjenu smyslu, ne izključjenje možnosti objasniti termin pri prvom pojavjenju. Jedinica ne znači v tom množinskom kontekstu samo jediničny element $1$.\par
\noindent\emph{Reason:} Jedinice is the short name for invertible elements, contrasted with zero divisors. The rationale expressly prefers it to the longer invertible-elements phrase. This is an argument for economy without loss of sense, not against explaining the term at first occurrence. In this plural context, unit does not mean only the identity element $1$.\par
\noindent{\footnotesize Izvorne zapisy / Original records: HIST-g121-007.}\par
\end{minipage}\par
\par\medskip\Needspace{7\baselineskip}
\noindent\begin{minipage}{\linewidth}
\hypertarget{TERM-0428}{}
\noindent\textbf{jedinice grupy ostatkov / ra?unanje s jedinicami}\quad{\footnotesize\texttt{TERM-0428}}\par
\noindent\emph{Iztočnik / Source:} Einheiten der Restgruppe / Rechnen mit Einheiten\par
\noindent\emph{Rusky / Russian:} [Bukvy ne sut zachovane / Lettering not recoverable]\par
\noindent\emph{Ukrajinsky / Ukrainian:} [Bukvy ne sut zachovane / Lettering not recoverable]\par
\smallskip\noindent\emph{Razlog:} Jedinice grupy ostatkov ovde imenovaju izbrane klasy, ktore v računju razkladov imaju idempotentnu, projekciji podobnu funkciju. Izvorni razlog ne poistovjetnjaje jih s libovoljnymi obrativymi elementami, i zato sovrěmenno značenje unit ne treba vnesti avtomatično. Bukvy ruskoj i ukrajinskoj frazy sut utratjene, a čest medžuslovjanskoj grafije jest povrědžena. Funkcija opisana v razlogu ne rekonstruuje te bukvy.\par
\noindent\emph{Reason:} Jedinice grupy ostatkov here names selected classes with an idempotent, projector-like role in decomposition calculations. The rationale does not identify them with arbitrary invertible elements, so the modern sense of unit must not be imported automatically. The Russian and Ukrainian lettering is lost, and part of the Interslavic spelling is damaged. The function described in the rationale does not reconstruct those letters.\par
\noindent{\footnotesize Izvorne zapisy / Original records: HIST-g079-026.}\par
\end{minipage}\par
\par\medskip\Needspace{7\baselineskip}
\noindent\begin{minipage}{\linewidth}
\hypertarget{TERM-0429}{}
\noindent\textbf{jediničny element}\quad{\footnotesize\texttt{TERM-0429}}\par
\noindent\emph{Iztočnik / Source:} Einheitselement\par
\noindent\emph{Rusky / Russian:} единичный элемент\par
\noindent\emph{Ukrajinsky / Ukrainian:} одиничний елемент\par
\smallskip\noindent\emph{Razlog:} Jediničny element imenuje konkretno multiplicativno jedinstvo, od ktorogo addicijeju i odčitanjem tvorjat se cěle veličiny prostogo polja v danoj konstrukciji. Dva razloga povezuju frazu v rabotah 9 i 10. Oni ne tvrdžu, že samo addicija i odčitanje davajut vse elementy vsakogo prostogo polja; ne treba razširiti iztočno tvrdženje zamenjanjem cělyh veličin cělym poljem.\par
\noindent\emph{Reason:} Jediničny element names the particular multiplicative identity from which the whole quantities of the prime field are generated by addition and subtraction in the construction. Two rationales connect its use in Papers 9 and 10. They do not assert that addition and subtraction alone generate every element of every prime field; the source claim must not be enlarged by replacing its whole quantities with the entire field.\par
\noindent{\footnotesize Izvorne zapisy / Original records: HIST-g068-003, HIST-g071-012.}\par
\end{minipage}\par
\par\medskip\Needspace{7\baselineskip}
\noindent\begin{minipage}{\linewidth}
\hypertarget{TERM-0430}{}
\noindent\textbf{jediničny element e / modulna baza}\quad{\footnotesize\texttt{TERM-0430}}\par
\noindent\emph{Iztočnik / Source:} Einheitselement e / Modulbasis\par
\noindent\emph{Rusky / Russian:} единичный элемент e / модульный базис\par
\noindent\emph{Ukrajinsky / Ukrainian:} одиничний елемент e / модульна база\par
\smallskip\noindent\emph{Razlog:} Jediničny element $e$ i modulna baza ostavaju dva odvojene algebraične nazvanja. Razlog izrično odlišaje element od jediničnogo ideala, iako oba imajut osnovu jediničn-. Ruska paralela imaje базис, a ukrajinska база; takva razlika ne čini dva različne matematične objekty. Identifikator $e$ služi točnomu odkazivanju, ne novomu prěvodnomu nazvanju.\par
\noindent\emph{Reason:} Jediničny element $e$ and modulna baza remain separate algebraic names. The rationale expressly distinguishes the element from the unit ideal although both use the jediničn- stem. Russian has базис while Ukrainian has база; this lexical difference does not produce different mathematical objects. The identifier $e$ provides precise reference, not a new translated name.\par
\noindent{\footnotesize Izvorne zapisy / Original records: HIST-g142-005.}\par
\end{minipage}\par
\par\medskip\Needspace{7\baselineskip}
\noindent\begin{minipage}{\linewidth}
\hypertarget{TERM-0431}{}
\noindent\textbf{jediničny ideal}\quad{\footnotesize\texttt{TERM-0431}}\par
\noindent\emph{Iztočnik / Source:} Einheitsideal\par
\noindent\emph{Rusky / Russian:} единичный идеал\par
\noindent\emph{Ukrajinsky / Ukrainian:} одиничний ідеал\par
\smallskip\noindent\emph{Razlog:} Jediničny ideal prodolžuje ranějšu familiju jedinica i jediničny. Tri razloga naglašaju doslědnost idealnoj terminologije i razlikovanje ideala od elementa. Točny imenik jest suščestvenny: jedin ideal moglo by značiti samo čislo idealov, a jediničny element jest drugi objekt. Kakvu jediničnu funkciju ideal ima v danej teoriji, treba ustanoviti iz definicije i operacije, ne samo iz pridevnika.\par
\noindent\emph{Reason:} Jediničny ideal continues the earlier jedinica/jediničny family. Three rationales emphasize consistent ideal terminology and the distinction between ideal and element. The noun matters: one ideal might merely count ideals, while identity element denotes a different object. Which identity role the ideal has in the theory must be established from its definition and operation, not from the adjective alone.\par
\noindent{\footnotesize Izvorne zapisy / Original records: HIST-g124-003, HIST-g134-005, HIST-g142-006.}\par
\end{minipage}\par
\par\medskip\Needspace{7\baselineskip}
\noindent\begin{minipage}{\linewidth}
\hypertarget{TERM-0432}{}
\noindent\textbf{jediničny ideal / jediničny element}\quad{\footnotesize\texttt{TERM-0432}}\par
\noindent\emph{Iztočnik / Source:} Einheitsideal / Einheitselement\par
\noindent\emph{Rusky / Russian:} единичный идеал / единичный элемент\par
\noindent\emph{Ukrajinsky / Ukrainian:} одиничний ідеал / одиничний елемент\par
\smallskip\noindent\emph{Razlog:} Para jediničny ideal i jediničny element zachovuje obšču osnovu za jedinstvo, ale imenik razděljuje tipy objektov. Izvorni razlog odlišaje ih takože od obyčnyh jedinic, to jest obrativyh elementov v odgovarjajučem kontekstu. Slovnična povezanost ne daje prava utožniti ideal, njegov element i množenstvenu grupu jedinic.\par
\noindent\emph{Reason:} The pair jediničny ideal and jediničny element retains a common unity stem while the nouns distinguish object types. The rationale also separates them from ordinary units, meaning invertible elements in the relevant context. Lexical relatedness does not license identifying an ideal, one of its elements, and a group of units.\par
\noindent{\footnotesize Izvorne zapisy / Original records: HIST-g121-005.}\par
\end{minipage}\par
\par\medskip\Needspace{7\baselineskip}
\noindent\begin{minipage}{\linewidth}
\hypertarget{TERM-0433}{}
\noindent\textbf{jediničny ideal / nulovy ideal / $\mathfrak{o}$ kako jediničny ideal}\quad{\footnotesize\texttt{TERM-0433}}\par
\noindent\emph{Iztočnik / Source:} Einheitsideal / Nullideal / $\mathfrak{o}$\par
\noindent\emph{Rusky / Russian:} единичный идеал / нулевой идеал / $\mathfrak{o}$ как единичный идеал\par
\noindent\emph{Ukrajinsky / Ukrainian:} одиничний ідеал / нульовий ідеал / $\mathfrak{o}$ як одиничний ідеал\par
\smallskip\noindent\emph{Razlog:} Jediničny ideal i nulovy ideal ne směšajut se, iako simboličny zapis možet čitatelju vspomněti nulu. Razlog odkazuje na definiciju iz odděla 1 raboty 24 i odvojeno spominanje nulovogo ideala v primětce odděla 4. Makro imenuje tipografsku oznaku; smysl ideala mora slědovati definiciji, ne vněšnej podobnosti znakov. Tuto jest konkretna notacijna opasnost, ne nova definicija nuly.\par
\noindent\emph{Reason:} Jediničny ideal and nulovy ideal must not be conflated, even if a symbol visually suggests zero. The rationale refers to Paper 24's Section 1 definition and a separate mention of the zero ideal in the Section 4 footnote. The macro names a typographical symbol; the ideal's meaning must follow the definition, not a resemblance between characters. This is a specific notational hazard, not a new definition of zero.\par
\noindent{\footnotesize Izvorne zapisy / Original records: HIST-g098-009.}\par
\end{minipage}\par
\par\medskip\Needspace{7\baselineskip}
\noindent\begin{minipage}{\linewidth}
\hypertarget{TERM-0434}{}
\noindent\textbf{jediničny ideal $\mathfrak{o}$ / jediničny element / jediničny klas}\quad{\footnotesize\texttt{TERM-0434}}\par
\noindent\emph{Iztočnik / Source:} Einheitsideal $\mathfrak{o}$ / Einheitselement / Einheitsklasse\par
\noindent\emph{Rusky / Russian:} единичный идеал $\mathfrak{o}$ / единичный элемент / единичный класс\par
\noindent\emph{Ukrajinsky / Ukrainian:} одиничний ідеал $\mathfrak{o}$ / одиничний елемент / одиничний клас\par
\smallskip\noindent\emph{Razlog:} Jediničny ideal, jediničny element i jediničny klas moraju ostati tri odvojene izraza. Oba razloga iz odděla 7 govorju, že ideal vseh elementov može obstajati i pri otsutnosti multiplicativnogo jediničnogo elementa v prstenu, dokolě odgovarjajuče polje klasov ima jediničny klas. Zato ne možno iz slova jediničny ideal samo po sobě zaključiti, že prsten imaje $1$. Tuta lokalna avtorska konvencija trěbuje osobno objasnjenje.\par
\noindent\emph{Reason:} Jediničny ideal, jediničny element and jediničny klas must remain three distinct expressions. Both Section 7 rationales say that the ideal of all elements can exist even when the ring has no multiplicative identity element, while the associated field of classes has an identity class. Unit ideal alone therefore does not justify assuming that the ring has $1$. The author's local convention needs explicit explanation.\par
\noindent{\footnotesize Izvorne zapisy / Original records: HIST-g114-002, HIST-g114-008.}\par
\smallskip\noindent\emph{Sravnjenje s tekstom:} Podnožna primětka raboty 30, odděla 7, izrično nazva $\mathfrak{o}$ idealom vseh elementov i dopušča, že v prstenu ne ima jediničnogo elementa, iako polje klasov ostatkov ima jediničny klas. Rabota 24, odděl 1, takože definira $\mathfrak{o}$ kako ideal vseh elementov. Primětka raboty 19, odděla 6, dodava drugo ograničenje: $\mathfrak{D}$ igra jediničnu rolu pri děliteljnosti i najmenšem občem kratnom, ne pri tvorjenju produkta. To sut različne konkretne upotrěby tojže familije slov. Prěvod ne smije vnesti sovrěmenny predpoklad, že vsak prsten ili modul iz teksta už jest zadany s jedinicoju i jediničnym dějstvom. Pasaže podpiraju razloge v TERM-0433–0435; ne dokazujut avtomatično smysl vsakego drugogo pojavjenja jedinica.\par
\noindent\emph{Textual comparison:} The footnote in Paper 30, section 7 expressly calls $\mathfrak{o}$ the ideal of all elements and allows the ring to lack an identity element even though the residue-class field has an identity class. Paper 24, section 1 also defines $\mathfrak{o}$ as the ideal of all elements. Paper 19, section 6 adds a different restriction: $\mathfrak{D}$ has an identity role for divisibility and least common multiple, not for ideal multiplication. These are distinct contextual uses of the same word family. Translation must not import a modern assumption that every ring or module in the text already comes with an identity and a unital action. The passages substantiate the reasons in TERM-0433–0435, not automatically every other occurrence of jedinica.\par
\noindent{\footnotesize TEXT-COMP-08; pasaže / passages: 13100, 15283, 16778; \texttt{CORPUS\_FORM\_COMPARISONS.json}.}\par
\end{minipage}\par
\par\medskip\Needspace{7\baselineskip}
\noindent\begin{minipage}{\linewidth}
\hypertarget{TERM-0435}{}
\noindent\textbf{jediničny ideal D}\quad{\footnotesize\texttt{TERM-0435}}\par
\noindent\emph{Iztočnik / Source:} Einheitsideal D\par
\noindent\emph{Rusky / Russian:} единичный идеал D\par
\noindent\emph{Ukrajinsky / Ukrainian:} одиничний ідеал D\par
\smallskip\noindent\emph{Razlog:} Jediničny ideal $D$ ima v toj primětce omeženu jediničnu funkciju: pri děliteljnosti i pri najmenšem občem kratnom, ne avtomatično pri produktu idealov. Dva zapisa zachovuju jedin razlog s tym izričnym ograničenjem. Skračenje do prosto jedinica bez operacije moglo by proměniti matematično tvrdženje. Latinična oznaka $D$ ne jest nova leksična čest termina.\par
\noindent\emph{Reason:} Jediničny ideal $D$ has a restricted identity role in the cited footnote: for divisibility and least common multiple, not automatically for ideal multiplication. Two records preserve one rationale with that explicit limitation. Shortening this to unit without naming the operation could change the mathematical claim. The identifier $D$ is not a new lexical part of the term.\par
\noindent{\footnotesize Izvorne zapisy / Original records: HIST-g081-054, HIST-g082-054.}\par
\end{minipage}\par
\par\medskip\Needspace{7\baselineskip}
\noindent\begin{minipage}{\linewidth}
\hypertarget{TERM-0436}{}
\noindent\textbf{jedini?ny klas}\quad{\footnotesize\texttt{TERM-0436}}\par
\noindent\emph{Iztočnik / Source:} Einheitsklasse\par
\noindent\emph{Rusky / Russian:} [Bukvy ne sut zachovane / Lettering not recoverable]\par
\noindent\emph{Ukrajinsky / Ukrainian:} [Bukvy ne sut zachovane / Lettering not recoverable]\par
\smallskip\noindent\emph{Razlog:} Jediničny klas prodolžuje imena klasov iz oddělov 3 i 5 v formulah (22)–(24). Razlog povezuje jego s jediniceju jednoj komponentnoj podgrupy. On ne opravdaje zaměnu imenika klas slovom element bez očuvanja kvocientnoj struktury. Cěla ruska i ukrajinska grafija i jedna čest medžuslovjanskoj formy ne sut zachovane; slovnična rekonstrukcija to naznačaje.\par
\noindent\emph{Reason:} Jediničny klas continues the class terminology of Sections 3 and 5 in formulas (22)–(24). The rationale relates it to the unit of a constituent subgroup. It does not justify replacing class by element without retaining the quotient structure. The Russian and Ukrainian lettering and part of the Interslavic spelling have not survived; the dictionary makes this explicit.\par
\noindent{\footnotesize Izvorne zapisy / Original records: HIST-g079-039.}\par
\end{minipage}\par
\par\medskip\Needspace{7\baselineskip}
\noindent\begin{minipage}{\linewidth}
\hypertarget{TERM-0437}{}
\noindent\textbf{jediničny klas / jedinica}\quad{\footnotesize\texttt{TERM-0437}}\par
\noindent\emph{Iztočnik / Source:} Einheitsklasse / Einheit\par
\noindent\emph{Rusky / Russian:} единичный класс / единица\par
\noindent\emph{Ukrajinsky / Ukrainian:} одиничний клас / одиниця\par
\smallskip\noindent\emph{Razlog:} Jediničny klas jest klas, ktory sodrži $1$, i v toj konstrukciji jest označeny $e$. Jedinica imenuje jego funkciju, klas jego status. Razlog naglašaje, že oba aspekty sut potrěbne v oddělu 3 za popis cěloj grupy ostatkov. To ne jest dozvoljenje utožniti predstavitelja $1$ s njegovoju klasoju bez ukazanja prěhoda k klasam.\par
\noindent\emph{Reason:} Jediničny klas is the class containing $1$, denoted $e$ in this construction. Jedinica names its role and klas its status. The rationale stresses that both aspects matter to Section 3's account of the entire residue group. This does not license identifying the representative $1$ with its class without acknowledging the passage to classes.\par
\noindent{\footnotesize Izvorne zapisy / Original records: HIST-g079-017.}\par
\end{minipage}\par
\par\medskip\Needspace{7\baselineskip}
\noindent\begin{minipage}{\linewidth}
\hypertarget{TERM-0438}{}
\noindent\textbf{jedinična matrica}\quad{\footnotesize\texttt{TERM-0438}}\par
\noindent\emph{Iztočnik / Source:} Einheitsmatrix\par
\noindent\emph{Rusky / Russian:} единичная матрица\par
\noindent\emph{Ukrajinsky / Ukrainian:} одинична матриця\par
\smallskip\noindent\emph{Razlog:} Jedinična matrica svęzuje matričny zapis s jediničnym operatorom, ale ostavaje imenikovoj frazoj za matricu. Razlog daje prednost poznavajemoj osnově matric- prěd novym užšim tvorjenjem. On ne govori, že operator i njegova matrica jesut jeden i toj samy objekt nezavisno od izbora prědstavjenja. Paralelne ruske i ukrajinske formy podpiraju oblik nazvanja.\par
\noindent\emph{Reason:} Jedinična matrica connects the matrix expression with the identity operator but remains a noun phrase naming a matrix. The rationale favours the recognizable matric- family over a narrower new coinage. It does not identify an operator and its matrix as the same object independently of a representation. The Russian and Ukrainian comparanda support the form of the name.\par
\noindent{\footnotesize Izvorne zapisy / Original records: HIST-g195-006.}\par
\end{minipage}\par
\par\medskip\Needspace{7\baselineskip}
\noindent\begin{minipage}{\linewidth}
\hypertarget{TERM-0439}{}
\noindent\textbf{jediničny modul}\quad{\footnotesize\texttt{TERM-0439}}\par
\noindent\emph{Iztočnik / Source:} Einheitsmodul\par
\noindent\emph{Rusky / Russian:} единичный модуль\par
\noindent\emph{Ukrajinsky / Ukrainian:} одиничний модуль\par
\smallskip\noindent\emph{Razlog:} Jediničny modul imenuje vtory dopustimy dělitelj v lokalnoj definiciji Primmodul. Razlog prodolžuje terminologiju jediničnogo klasa, no ostavjaje modul izričnym tipom objekta. Iz samej frazy ne slěduje sovrěmenna definicija prostogo ili prime modula bez čitanja iztočnogo uslovja o děliteljah.\par
\noindent\emph{Reason:} Jediničny modul names the second permitted divisor in the local definition of Primmodul. The rationale continues the unit-class vocabulary while keeping module explicit as the object type. The phrase alone does not supply a modern definition of simple or prime module without reading the source's divisor condition.\par
\noindent{\footnotesize Izvorne zapisy / Original records: HIST-g079-045.}\par
\end{minipage}\par
\par\medskip\Needspace{7\baselineskip}
\noindent\begin{minipage}{\linewidth}
\hypertarget{TERM-0440}{}
\noindent\textbf{jediničny operator}\quad{\footnotesize\texttt{TERM-0440}}\par
\noindent\emph{Iztočnik / Source:} Einheitsoperator\par
\noindent\emph{Rusky / Russian:} единичный оператор\par
\noindent\emph{Ukrajinsky / Ukrainian:} одиничний оператор\par
\smallskip\noindent\emph{Razlog:} Jediničny operator imenuje dějstvo, ne samo element prstena. Razlog naglašaje razklad, pri ktorom toj samy jediničny element dějstvuje kako identičnost na jednom sumandu i kako nula na drugom. Avtomatično poistovjetnjenje jedinicy prstena s identičnym operatorom na cělom modulu by izbrisalo toj slučaj. Prěvod mora zachovati razliku medžu elementom i jego dějstvom.\par
\noindent\emph{Reason:} Jediničny operator names an action, not merely a ring element. The rationale emphasizes a decomposition in which the same ring identity acts as identity on one summand and as zero on the other. Automatically identifying the ring's identity element with the identity operator on the entire module would erase that case. The translation must retain the distinction between an element and its action.\par
\noindent{\footnotesize Izvorne zapisy / Original records: HIST-g195-002.}\par
\smallskip\noindent\emph{Sravnjenje s tekstom:} Rabota 34, odděl 17, počinaje s modulom nad prstenom s jedinicoju, ale ne prědpoklada, že ta jedinica dějstvuje identično na cělom modulu. Ona razklada modul na $\mathfrak{M}_e+\mathfrak{M}_0$ s pomoču $m=em+(m-em)$. Na prvoj česti $e$ jest jediničny operator, na vtoroj nulovy. Tekst potom ograničuje razmatranje prědstavjenij na prvu čest, kde jedinici odgovarjaje jedinična matrica. Ta izrična poslědovatelnost pokazuje, začto element, operator i matrica ne mogut byti zaměnjene jednim neoprěděljenym nazvanjem jedinica.\par
\noindent\emph{Textual comparison:} Paper 34, section 17 begins with a module over a ring with identity, without assuming that the identity acts identically on the whole module. It decomposes the module as $\mathfrak{M}_e+\mathfrak{M}_0$ using $m=em+(m-em)$. On the first part $e$ is the identity operator; on the second it is zero. The text then restricts its representation discussion to the first part, where the ring identity corresponds to the identity matrix. This explicit sequence explains why element, operator and matrix cannot be replaced by one undifferentiated unity term.\par
\noindent{\footnotesize TEXT-COMP-09; pasaže / passages: 21587, 21589, 21591, 21595, 21605, 21607; \texttt{CORPUS\_FORM\_COMPARISONS.json}.}\par
\end{minipage}\par
\par\medskip\Needspace{7\baselineskip}
\noindent\begin{minipage}{\linewidth}
\hypertarget{TERM-0441}{}
\noindent\textbf{něktore neopreděljene zaměniti čislami}\quad{\footnotesize\texttt{TERM-0441}}\par
\noindent\emph{Iztočnik / Source:} einige der Unbestimmten durch Zahlen ersetzen\par
\noindent\emph{Rusky / Russian:} некоторые неопределенные заменить числами\par
\noindent\emph{Ukrajinsky / Ukrainian:} деякі невизначені замінити числами\par
\smallskip\noindent\emph{Razlog:} Fraza označaje zaměnu něktoryh neoprěděljenyh čislami, ne zaměnu vseh proměnnyh bez razlike. Razlog povezuje odděl 3 s postupnym čislovym specializovanjem iz lemy odděla 2. Něktore nosi kvantifikacijno ograničenje. Ta formulacija govori o matematičnoj operaciji specializacije, ne samo o približnom čislovom izračunanju.\par
\noindent\emph{Reason:} The phrase specifies replacing some indeterminates by numbers, not indiscriminately replacing all variables. The rationale connects Section 3 with the successive numerical specialization lemma in Section 2. Něktore carries a quantifier restriction. This describes a mathematical operation of specialization, not merely an approximate numerical calculation.\par
\noindent{\footnotesize Izvorne zapisy / Original records: HIST-g043-003.}\par
\end{minipage}\par
\par\medskip\Needspace{7\baselineskip}
\noindent\begin{minipage}{\linewidth}
\hypertarget{TERM-0442}{}
\noindent\textbf{jednokratno svijanje}\quad{\footnotesize\texttt{TERM-0442}}\par
\noindent\emph{Iztočnik / Source:} einmalige Faltung\par
\noindent\emph{Rusky / Russian:} однократная свертка\par
\noindent\emph{Ukrajinsky / Ukrainian:} одноразове згортання\par
\smallskip\noindent\emph{Razlog:} Jednokratno svijanje zachovuje čislo operacij: svijanje izvršaje se jedin raz. Razlog prodolžuje tu kratku frazu iz ranějših oddělov. Jednokratno ne znači prosto neponovimo događanje; v dokazě jest važno, kolikokrat operacija priměnjaje se. Zapis ne sravnjava jeho s drugym historijnym oblikom jednorazovo.\par
\noindent\emph{Reason:} Jednokratno svijanje retains the operation count: folding is performed once. The rationale continues the short phrase from earlier sections. Jednokratno does not simply describe a unique event; the number of applications matters to the proof. The record does not compare it with the other historical form jednorazovo.\par
\noindent{\footnotesize Izvorne zapisy / Original records: HIST-g025-005.}\par
\end{minipage}\par
\par\medskip\Needspace{7\baselineskip}
\noindent\begin{minipage}{\linewidth}
\hypertarget{TERM-0443}{}
\noindent\textbf{jednorazovo svijanje}\quad{\footnotesize\texttt{TERM-0443}}\par
\noindent\emph{Iztočnik / Source:} einmalige Faltung\par
\noindent\emph{Rusky / Russian:} однократная свертка\par
\noindent\emph{Ukrajinsky / Ukrainian:} одноразове згортання\par
\smallskip\noindent\emph{Razlog:} Jednorazovo svijanje vyražaje isto čislo operacij, ktoro v drugom zapisu imenuje jednokratno svijanje. Ruska i ukrajinska paralela takože izrično govorju o jednom priměnjenju. Razlog naglašaje kontinuitet vyraženja, ale ne dava razlog za prednost osnovy -raz- prěd -krat-. Tute zapisy zato dokumentuju formnu variaciju, ne dva razne matematične postupky.\par
\noindent\emph{Reason:} Jednorazovo svijanje expresses the same operation count named elsewhere by jednokratno svijanje. The Russian and Ukrainian comparanda also explicitly denote one application. The rationale stresses continuity of wording but gives no reason to prefer the -raz- stem over -krat-. These records document formal variation, not two different mathematical procedures.\par
\noindent{\footnotesize Izvorne zapisy / Original records: HIST-g016-007.}\par
\end{minipage}\par
\par\medskip\Needspace{7\baselineskip}
\noindent\begin{minipage}{\linewidth}
\hypertarget{TERM-0444}{}
\noindent\textbf{jednokratna transvekcija}\quad{\footnotesize\texttt{TERM-0444}}\par
\noindent\emph{Iztočnik / Source:} einmalige Überschiebung\par
\noindent\emph{Rusky / Russian:} однократная трансвекция\par
\noindent\emph{Ukrajinsky / Ukrainian:} одноразова трансвекція\par
\smallskip\noindent\emph{Razlog:} Jednokratna transvekcija stavjaje jedin korak v kontrast s blizkimi povtorjenymi ili dvojnymi redukcijami. Čislo operacij jest glavni zabilěženy razlog. On ne ustanovjaje, že vsaka redukcija jest transvekcija, ni ne opravdaje avtomatično prěvoditi vsako Überschiebung tojže imeniceju bez konteksta. Sravniteljne formy zachovuju meždunarodnu osnovu transvekc-.\par
\noindent\emph{Reason:} Jednokratna transvekcija contrasts one step with nearby repeated or double reductions. Operation count is the principal recorded reason. It does not establish that every reduction is a transvection, or justify translating every Überschiebung this way regardless of context. The comparanda retain the international transvection stem.\par
\noindent{\footnotesize Izvorne zapisy / Original records: HIST-g022-006.}\par
\end{minipage}\par
\par\medskip\Needspace{7\baselineskip}
\noindent\begin{minipage}{\linewidth}
\hypertarget{TERM-0445}{}
\noindent\textbf{jednorazovy prěnos}\quad{\footnotesize\texttt{TERM-0445}}\par
\noindent\emph{Iztočnik / Source:} einmalige Überschiebung\par
\noindent\emph{Rusky / Russian:} однократное перекладывание\par
\noindent\emph{Ukrajinsky / Ukrainian:} одноразове перекладання\par
\smallskip\noindent\emph{Razlog:} Jednorazovy prěnos jest drugy historijny kandidat za Überschiebung. Razlog imenuje jeho prozračnym, ale izrično ostavjaje izbor otvorjenym do sravnjenja s terminologijeju teorije invariantov. Takva napomena ne jest dokaz, že prěnos už byl potvrđeny autoritetnym iztočnikom. Razlika od transvekcija v TERM-0444 potrěbuje kontekstualno i terminologično sravnjenje, ne tihu zaměnu istorije izbora.\par
\noindent\emph{Reason:} Jednorazovy prěnos is another historical candidate for Überschiebung. The rationale calls it transparent but expressly leaves the choice open pending comparison with invariant-theory terminology. That reservation is not evidence that an authoritative source had already confirmed prěnos. Its relationship to transvekcija in TERM-0444 requires contextual and terminological comparison, not silent rewriting of the choice's history.\par
\noindent{\footnotesize Izvorne zapisy / Original records: HIST-g013-003.}\par
\end{minipage}\par
\par\medskip\Needspace{7\baselineskip}
\noindent\begin{minipage}{\linewidth}
\hypertarget{TERM-0446}{}
\noindent\textbf{Vstavjenje osnovnyh idealov i jih komponentov v obču teoriju idealov}\quad{\footnotesize\texttt{TERM-0446}}\par
\noindent\emph{Iztočnik / Source:} Einordnung der Grundideale und ihrer Komponenten in die allgemeine Idealtheorie\par
\noindent\emph{Rusky / Russian:} Включение основных идеалов и их компонент в общую теорию идеалов\par
\noindent\emph{Ukrajinsky / Ukrainian:} Вбудування основних ідеалів та їхніх компонент у загальну теорію ідеалів\par
\smallskip\noindent\emph{Razlog:} Vstavjenje ... v obču teoriju idealov opisuje, kako osnovne idealy i njihove komponenty nahodet město v širšej teoriji. Razlog odbacaje slab opis prostogo rasporedženja. V tom nazvanju vstavjenje ne mora samo po sobě imenovati injektivny homomorfizm: predmetom jest teoretično uključjenje Hentzelt–Noetherovyh objektov. Konkretnu matematičnu konstrukciju mora pokazati odděl, ne metafora naslova.\par
\noindent\emph{Reason:} Vstavjenje ... v obču teoriju idealov describes how the basic ideals and their components fit within the broader theory. The rationale rejects a loose description of mere arrangement. In this title, vstavjenje does not by itself name an injective homomorphism: it concerns integrating the Hentzelt–Noether objects into the theory. The section must supply the particular mathematical construction; the title's metaphor cannot do so.\par
\noindent{\footnotesize Izvorne zapisy / Original records: HIST-g099-001.}\par
\end{minipage}\par
\par\medskip\Needspace{7\baselineskip}
\noindent\begin{minipage}{\linewidth}
\hypertarget{TERM-0447}{}
\noindent\textbf{Uvrstjenje hiperkompleksnyh sistemov v teoriju}\quad{\footnotesize\texttt{TERM-0447}}\par
\noindent\emph{Iztočnik / Source:} Einordnung der hyperkomplexen Systeme\par
\noindent\emph{Rusky / Russian:} Встраивание гиперкомплексных систем в теорию\par
\noindent\emph{Ukrajinsky / Ukrainian:} Впорядкування гіперкомплексних систем у теорії\par
\smallskip\noindent\emph{Razlog:} Uvrstjenje hiperkompleksnyh sistemov v teoriju opisuva njihovo město v teoriji prědstavjenij. Zapis razděljuje německo Einordnung iz navedenej pečatnej strani od Einbeziehung v ranějšem TeX-tekstu i dava prednost prvomu smyslu. Tuta jezyčna odluka jest svęzana s izborom iztočnogo čitanja. Nova objasnjujuča proza zachovuje tvrdnju zapisa, ale ne predstavlja ju kako novu vizualnu prověrku originala.\par
\noindent\emph{Reason:} Uvrstjenje hiperkompleksnyh sistemov v teoriju describes their place within representation theory. The record distinguishes Einordnung in the cited printed witness from Einbeziehung in an earlier TeX text and favours the former sense. The linguistic choice is thus linked to a source-reading decision. This explanation preserves the record's claim without presenting it as a fresh visual examination of the original.\par
\noindent{\footnotesize Izvorne zapisy / Original records: HIST-g198-001.}\par
\end{minipage}\par
\par\medskip\Needspace{7\baselineskip}
\noindent\begin{minipage}{\linewidth}
\hypertarget{TERM-0448}{}
\noindent\textbf{jednostranno prosty ideal}\quad{\footnotesize\texttt{TERM-0448}}\par
\noindent\emph{Iztočnik / Source:} einseitig einfaches Ideal\par
\noindent\emph{Rusky / Russian:} односторонне простой идеал\par
\noindent\emph{Ukrajinsky / Ukrainian:} однобічно простий ідеал\par
\smallskip\noindent\emph{Razlog:} Jednostranno prosty ideal zachovuje i prostost, i stranu idealnogo uslovja. Razlog povezuje izraz s prěhodom od idempotentov k razščepljenju v teoriji prědstavjenij. Izpustiti jednostranno moglo by pretvoriti uslovje o lěvyh ili pravyh idealah v uslovje o dvustrannyh idealah. Dve kvalifikacije imajut različne funkcije i ne zaměnjaju jedna drugu.\par
\noindent\emph{Reason:} Jednostranno prosty ideal retains both simplicity and the side on which the ideal condition applies. The rationale connects it with the passage from idempotents to splitting in representation theory. Omitting jednostranno could turn a left- or right-ideal condition into a two-sided-ideal condition. The two qualifications have different functions and cannot replace one another.\par
\noindent{\footnotesize Izvorne zapisy / Original records: HIST-g173-001.}\par
\end{minipage}\par
\par\medskip\Needspace{7\baselineskip}
\noindent\begin{minipage}{\linewidth}
\hypertarget{TERM-0449}{}
\noindent\textbf{jednostranne / dvustranne idealy}\quad{\footnotesize\texttt{TERM-0449}}\par
\noindent\emph{Iztočnik / Source:} einseitige / zweiseitige Ideale\par
\noindent\emph{Rusky / Russian:} односторонние / двусторонние идеалы\par
\noindent\emph{Ukrajinsky / Ukrainian:} односторонні / двосторонні ідеали\par
\smallskip\noindent\emph{Razlog:} Jednostranne i dvustranne idealy prodolžuju terminologiju nekomutativnyh uslovij iz ranějših rabot. Strana govori o množenju, ne o prostorovoj formě ideala. Ruska i ukrajinska paralela tak samo zachovuju kontrast jednoj i dvoh stran. Kratky razlog ne dava nove definicije, ale trěbuje ne izgubiti tu informaciju pri prěvodu.\par
\noindent\emph{Reason:} Jednostranne and dvustranne idealy continue the terminology of noncommutative side conditions from earlier works. Side concerns multiplication, not an ideal's spatial shape. The Russian and Ukrainian comparanda likewise preserve the one-sided/two-sided contrast. The short rationale adds no new definition but requires that the information survive translation.\par
\noindent{\footnotesize Izvorne zapisy / Original records: HIST-g105-007.}\par
\end{minipage}\par
\par\medskip\Needspace{7\baselineskip}
\noindent\begin{minipage}{\linewidth}
\hypertarget{TERM-0450}{}
\noindent\textbf{jednostranne moduly}\quad{\footnotesize\texttt{TERM-0450}}\par
\noindent\emph{Iztočnik / Source:} einseitige Moduln\par
\noindent\emph{Rusky / Russian:} односторонние модули\par
\noindent\emph{Ukrajinsky / Ukrainian:} однобічні модулі\par
\smallskip\noindent\emph{Razlog:} Jednostranne moduly jesut objekty zatvorjene pri množenju s jednoj strany v smyslu odděla 3. Izvorni razlog izrično zabrana poistovjetniti jih s obyčnymi dvustrannymi idealami. Imenik modul i pridevnik jednostranny zato oba nosat matematično sodržanje. Obšče nazvanje ideal bez kvalifikacije izgubilo by upotrěbjenu strukturnu razliku.\par
\noindent\emph{Reason:} Jednostranne moduly are objects closed under multiplication on one side in the sense of Section 3. The rationale explicitly rejects collapsing them into ordinary two-sided ideals. Both modul and jednostranny therefore carry mathematical content. An unqualified ideal label would lose the structural distinction used in the passage.\par
\noindent{\footnotesize Izvorne zapisy / Original records: HIST-g079-013.}\par
\end{minipage}\par
\par\medskip\Needspace{7\baselineskip}
\noindent\begin{minipage}{\linewidth}
\hypertarget{TERM-0451}{}
\noindent\textbf{prosto izomorfna}\quad{\footnotesize\texttt{TERM-0451}}\par
\noindent\emph{Iztočnik / Source:} einstufig isomorph\par
\noindent\emph{Rusky / Russian:} просто изоморфна\par
\noindent\emph{Ukrajinsky / Ukrainian:} просто ізоморфна\par
\smallskip\noindent\emph{Razlog:} Prosto izomorfna jest zabilěženo čitanje staršego einstufig isomorph v podnožnoj primětce. Razlog odbacaje bukvalno jednostupenno izomorfna kako zavodljivo v sovrěmennoj proze i svęzuje frazu s redukcijeju k abstraktnoj konečnoj grupě. To jest kontekstualna interpretacija staršego izraza, ne obšče leksično pravilo o slově einstufig. Zapis ne dava odvojenu historijno-terminologičnu citaciju, ktora by zatvorila pitanje.\par
\noindent\emph{Reason:} Prosto izomorfna is the recorded reading of older einstufig isomorph in a footnote. The rationale rejects literal one-stage isomorphic as misleading in modern prose and connects the phrase with reduction to an abstract finite group. This is a contextual interpretation of an older expression, not a general lexical rule for einstufig. The record supplies no separate historical-terminological citation that would settle the question.\par
\noindent{\footnotesize Izvorne zapisy / Original records: HIST-g077-014.}\par
\end{minipage}\par
\par\medskip\Needspace{7\baselineskip}
\noindent\begin{minipage}{\linewidth}
\hypertarget{TERM-0452}{}
\noindent\textbf{razděljenje oblastij ... na klasy}\quad{\footnotesize\texttt{TERM-0452}}\par
\noindent\emph{Iztočnik / Source:} Einteilung der Bereiche ... in Klassen\par
\noindent\emph{Rusky / Russian:} разбиение областей ... на классы\par
\noindent\emph{Ukrajinsky / Ukrainian:} поділ областей ... на класи\par
\smallskip\noindent\emph{Razlog:} Razděljenje oblastij ... na klasy opisuje klasifikaciju po izomorfizmu, ne prosto tvorjenje novih oblastij. Rusko разбиение i ukrajinsko поділ podpiraju razděljenje; ukrajinsky izbor izrično protivopostavjen jest формування. Medžuslovjansko klasy zachovuje rezultat razděljenja. Točny kriterij členstva v klasu treba ostaviti vidimym v tekstu.\par
\noindent\emph{Reason:} Razděljenje oblastij ... na klasy describes classification by isomorphism, not merely forming new domains. Russian разбиение and Ukrainian поділ support partition vocabulary; the Ukrainian choice is explicitly contrasted with формування. Interslavic klasy retains the outcome of the partition. The exact criterion for class membership must remain visible in the text.\par
\noindent{\footnotesize Izvorne zapisy / Original records: HIST-g066-001.}\par
\end{minipage}\par
\par\medskip\Needspace{7\baselineskip}
\noindent\begin{minipage}{\linewidth}
\hypertarget{TERM-0453}{}
\noindent\textbf{jednoredna matrica}\quad{\footnotesize\texttt{TERM-0453}}\par
\noindent\emph{Iztočnik / Source:} einzeilige Matrix\par
\noindent\emph{Rusky / Russian:} однострочная матрица\par
\noindent\emph{Ukrajinsky / Ukrainian:} однорядкова матриця\par
\smallskip\noindent\emph{Razlog:} Jednoredna matrica imenuje matricu s jednim redom. Razlog daje toj kratkoj tvorbě prednost prěd težšim meždunarodnym vyraženjem, ale ne zachovuje jeho točnu formu. Jednoredna ne znači matrica ranga jedan: čislo redov i rang sut različne informaciji. Kratkost izbora zato ne smije izbrisati konkretnu dimenzionalnu kvalifikaciju.\par
\noindent\emph{Reason:} Jednoredna matrica names a matrix with one row. The reason prefers this compact formation to a heavier international expression, but does not preserve the alternative's exact wording. One-row does not mean rank-one matrix: row count and rank are different information. The choice's brevity must not erase that specific size qualification.\par
\noindent{\footnotesize Izvorne zapisy / Original records: HIST-g131-003.}\par
\end{minipage}\par
\par\medskip\Needspace{7\baselineskip}
\noindent\begin{minipage}{\linewidth}
\hypertarget{TERM-0454}{}
\noindent\textbf{vzete pojedinačno}\quad{\footnotesize\texttt{TERM-0454}}\par
\noindent\emph{Iztočnik / Source:} einzeln genommen\par
\noindent\emph{Rusky / Russian:} взятые по отдельности\par
\noindent\emph{Ukrajinsky / Ukrainian:} взяті окремо\par
\smallskip\noindent\emph{Razlog:} Vzete pojedinačno ustanovjaje, že postupky razmatrajut se po jednomu, pri zaključku o njihovoj nezavisnosti. Razlog govori, že one ne reproduciraju jedna drugu. To ne jest tvrdženje o vsakoj vozmožnoj kombinaciji postupkov; kvantifikaciju ne treba razširjati. Ruska i ukrajinska paralela takože imenuju odvojeno razmatranje.\par
\noindent\emph{Reason:} Vzete pojedinačno specifies considering the procedures one at a time in the concluding independence statement. The rationale says that they do not reproduce each other. This is not a claim about every possible combination of procedures; the quantification must not be enlarged. The Russian and Ukrainian comparanda likewise name separate consideration.\par
\noindent{\footnotesize Izvorne zapisy / Original records: HIST-g033-010.}\par
\end{minipage}\par
\par\medskip\Needspace{7\baselineskip}
\noindent\begin{minipage}{\linewidth}
\hypertarget{TERM-0455}{}
\noindent\textbf{oddělne formy}\quad{\footnotesize\texttt{TERM-0455}}\par
\noindent\emph{Iztočnik / Source:} einzelne Formen\par
\noindent\emph{Rusky / Russian:} отдельные формы\par
\noindent\emph{Ukrajinsky / Ukrainian:} окремі форми\par
\smallskip\noindent\emph{Razlog:} Oddělne formy razlikuje pojedine formy od ih produktov ili većih redov form. Razlog jest različenje urovnja objekta, ne tvrdnja, že formy sut matematično nezavisne. Oddělne zato ne treba zameniti samym nezavisne bez iztočnogo osnova.\par
\noindent\emph{Reason:} Oddělne formy distinguishes individual forms from their products or larger sequences. The reason concerns the level of the object, not an assertion that the forms are mathematically independent. Oddělne must therefore not become independent without support from the source.\par
\noindent{\footnotesize Izvorne zapisy / Original records: HIST-g026-010.}\par
\end{minipage}\par
\par\medskip\Needspace{7\baselineskip}
\noindent\begin{minipage}{\linewidth}
\hypertarget{TERM-0456}{}
\noindent\textbf{razširjati oddělne komponenty}\quad{\footnotesize\texttt{TERM-0456}}\par
\noindent\emph{Iztočnik / Source:} einzelne Komponenten erweitern\par
\noindent\emph{Rusky / Russian:} расширять отдельные компоненты\par
\noindent\emph{Ukrajinsky / Ukrainian:} розширювати окремі компоненти\par
\smallskip\noindent\emph{Razlog:} Razširjati oddělne komponenty vyražaje postupok po komponentah. Razlog bira obyčnu glagolnu frazu, ibo ona podaje zaključok paragrafa: razširjenja možno razmatrjati oddělno dlja komponentov. Fraza ne prinosi novy abstraktny pojem i ne izpustjaje razširjenje kako operaciju.\par
\noindent\emph{Reason:} Razširjati oddělne komponenty expresses a componentwise procedure. The reason favours an ordinary verb phrase because it states the paragraph's conclusion: extensions can be considered separately for the components. The expression neither introduces a new abstract concept nor omits extension as the operation.\par
\noindent{\footnotesize Izvorne zapisy / Original records: HIST-g123-007.}\par
\end{minipage}\par
\par\medskip\Needspace{7\baselineskip}
\noindent\begin{minipage}{\linewidth}
\hypertarget{TERM-0457}{}
\noindent\textbf{jedine koeficienty pri t v $|tE-C|$}\quad{\footnotesize\texttt{TERM-0457}}\par
\noindent\emph{Iztočnik / Source:} einzelnen Koeffizienten von t in $|tE-C|$\par
\noindent\emph{Rusky / Russian:} отдельные коэффициенты при t в $|tE-C|$\par
\noindent\emph{Ukrajinsky / Ukrainian:} окремі коефіцієнти при t у $|tE-C|$\par
\smallskip\noindent\emph{Razlog:} Koeficienty pri $t$ jest zabilěžena fraza, izbrana blizko ruskoj i ukrajinskoj upotrěbě. Vse tri jezyčne razloga insistuju na množinskom tvrdženju o koeficientah determinanta $|tE-C|$. Německo einzelnen znači oddělne koeficienty, ne jedinstveny koeficient; slovo jedine v kandidatu ne treba čitati kak tvrdženje o unikalnosti. Izvorni polja s jezyčnymi imenami sodržuju anglijske razloge, ne citaty na tyh jezykah.\par
\noindent\emph{Reason:} Koeficienty pri $t$ is the recorded phrase, chosen close to Russian and Ukrainian usage. All three language-specific rationales insist on a plural statement about the coefficients of $|tE-C|$. German einzelnen refers to the individual coefficients, not a unique coefficient; jedine in the proposal must not be read as a uniqueness claim. The source fields labelled by language contain English rationales, not quotations in those languages.\par
\noindent{\footnotesize Izvorne zapisy / Original records: HIST-g135-004.}\par
\end{minipage}\par
\par\medskip\Needspace{7\baselineskip}
\noindent\begin{minipage}{\linewidth}
\hypertarget{TERM-0458}{}
\noindent\textbf{pojedina norma}\quad{\footnotesize\texttt{TERM-0458}}\par
\noindent\emph{Iztočnik / Source:} Einzelnorm\par
\noindent\emph{Rusky / Russian:} отдельная норма\par
\noindent\emph{Ukrajinsky / Ukrainian:} окрема норма\par
\smallskip\noindent\emph{Razlog:} Pojedina norma zachovuje konkretno nazvanje pojedinogo člena $R^{(i)}$ v Noetherinoj notaciji. Razlog odbacaje previše abstraktno novo nazvanje, ktoro skrylo by tu indeksovanu familiju. Pojedina ne znači vsegda jedina vozmožna norma; ona izdvaja člen familije v danom kontekstu.\par
\noindent\emph{Reason:} Pojedina norma retains the specific name of an individual $R^{(i)}$ in Noether's notation. The rationale rejects over-abstracting the name in a way that would obscure the indexed family. Pojedina does not generally mean the only possible norm; it singles out a member of the family in context.\par
\noindent{\footnotesize Izvorne zapisy / Original records: HIST-g095-015.}\par
\end{minipage}\par
\par\medskip\Needspace{7\baselineskip}
\noindent\begin{minipage}{\linewidth}
\hypertarget{TERM-0459}{}
\noindent\textbf{pojedina suma}\quad{\footnotesize\texttt{TERM-0459}}\par
\noindent\emph{Iztočnik / Source:} Einzelsumme\par
\noindent\emph{Rusky / Russian:} отдельная сумма\par
\noindent\emph{Ukrajinsky / Ukrainian:} окрема сума\par
\smallskip\noindent\emph{Razlog:} Pojedina suma imenuje jednu iz sum izdvajanih v dokazě iz formul (25) i (30), indeksovanih $\alpha$ ili $\beta$. Razlog govorit o razkladu argumenta na take česti. Imenik suma ne treba avtomatično zameniti slovo summand: jedna iz čestej sama može byti sumoju, i o tom rozhoduje konkretny formulovy zapis.\par
\noindent\emph{Reason:} Pojedina suma names one of the sums separated in the proof from formulas (25) and (30), indexed by $\alpha$ or $\beta$. The rationale concerns dividing the argument into these parts. Suma must not automatically become summand: a part can itself be a sum, as determined by the particular formula.\par
\noindent{\footnotesize Izvorne zapisy / Original records: HIST-g034-006.}\par
\end{minipage}\par
\par\medskip\Needspace{7\baselineskip}
\noindent\begin{minipage}{\linewidth}
\hypertarget{TERM-0460}{}
\noindent\textbf{element p, čij glavny ideal jest dělimy točno prvoju stepenju}\quad{\footnotesize\texttt{TERM-0460}}\par
\noindent\emph{Iztočnik / Source:} Element p, dessen Hauptideal genau durch die erste Potenz ... teilbar ist\par
\noindent\emph{Rusky / Russian:} элемент p, главный идеал которого делится ровно на первую степень\par
\noindent\emph{Ukrajinsky / Ukrainian:} елемент p, головний ідеал якого ділиться точно на перший степінь\par
\smallskip\noindent\emph{Razlog:} Element $p$, čij glavny ideal jest dělimy točno prvoju stepenju, razděljuje element od ideala, ktory on generuje, i zachovuje točny eksponent. Izvorni zapis pripisuje tu formulaciju redakcijnomu popravjenju, ne prostomu prěvodu neproměnjenoj pečatnej frazy. Zato ona mora byti razpoznana kako čitanje popravjenogo teksta. Bez sravnjenja samogo pečatnogo města i argumenta ne možno proglasiti ju potvrđenym ispravjenjem; diplomatičny prěvod ne može priměniti ju tiho.\par
\noindent\emph{Reason:} Element $p$, čij glavny ideal jest dělimy točno prvoju stepenju distinguishes an element from the principal ideal it generates and retains the exact exponent. The record attributes this wording to an editorial correction, not simply to translation of an unchanged printed phrase. It must therefore be identified as a reading of corrected text. Without examining the printed passage and argument, it cannot be certified as a valid emendation; a diplomatic translation must not adopt it silently.\par
\noindent{\footnotesize Izvorne zapisy / Original records: HIST-g156-004.}\par
\smallskip\noindent\emph{Sravnjenje s tekstom:} Aktualny prěvod raboty 31, odděla 8, izrično izbira element $p$ i govori o glavnom idealu, ktory on poraždaje. To potvrđaje, že zachovany leksičny razlog ne govori samo o nepriměnjenom prědlogu. Ale toj prěvod sam ne jest nezavisnym dokazom pečatnogo čitanja ili opravdanosti redakcijnogo popravjenja. Treba odvojeno zachovati autorovu formu i argument za popravjenu formu; iz dostupnogo slovničnogo zapisa ne slěduje završenost takvoj kolacije. Tuto sravnjenje utvrđuje čitanje prěvoda, ne potvrđaje věrnost pečatnomu iztočniku.\par
\noindent\emph{Textual comparison:} The current translation of Paper 31, section 8 explicitly chooses an element $p$ and refers to the principal ideal it generates. This confirms that the preserved lexical rationale concerns an adopted reading, not merely an unused proposal. But the translation is not independent evidence for the printed reading or the validity of the emendation. The author's wording and the argument for an emended wording must be preserved separately; the available lexical record does not establish that this collation is complete. This comparison identifies the translation's reading without certifying fidelity to the printed source.\par
\noindent{\footnotesize TEXT-COMP-10; pasaže / passages: 18823; \texttt{CORPUS\_FORM\_COMPARISONS.json}.}\par
\end{minipage}\par
\par\medskip\Needspace{7\baselineskip}
\noindent\begin{minipage}{\linewidth}
\hypertarget{TERM-0461}{}
\noindent\textbf{element s nulevoju normoju}\quad{\footnotesize\texttt{TERM-0461}}\par
\noindent\emph{Iztočnik / Source:} Element verschwindender Norm\par
\noindent\emph{Rusky / Russian:} элемент нулевой нормы\par
\noindent\emph{Ukrajinsky / Ukrainian:} елемент нульової норми\par
\smallskip\noindent\emph{Razlog:} Element s nulevoju normoju zaměnjaje slovnu frazu o izčezajučej normě kratkym imennym vyraženjem. Uslovje ostavaje, že norma jest nula; to ne označaje avtomatično, že sam element jest nula. Ruska i ukrajinska forma takože kvalifikujut normu, ne element. Razlog jest jasnost kratkoj formy v podnožnoj primětce, ne proměna matematičnogo kriterija.\par
\noindent\emph{Reason:} Element s nulevoju normoju replaces the verbal expression about a vanishing norm with a compact noun phrase. The condition is that the norm is zero; it does not automatically say that the element itself is zero. The Russian and Ukrainian forms also qualify the norm rather than the element. The recorded reason concerns concise wording in a footnote, not a different mathematical criterion.\par
\noindent{\footnotesize Izvorne zapisy / Original records: HIST-g174-004.}\par
\end{minipage}\par
\par\medskip\Needspace{7\baselineskip}
\noindent\begin{minipage}{\linewidth}
\hypertarget{TERM-0462}{}
\noindent\textbf{element drugogo roda / eksponenta}\quad{\footnotesize\texttt{TERM-0462}}\par
\noindent\emph{Iztočnik / Source:} Element zweiter Art / Exponent\par
\noindent\emph{Rusky / Russian:} элемент второго рода / экспонента\par
\noindent\emph{Ukrajinsky / Ukrainian:} елемент другого роду / експонента\par
\smallskip\noindent\emph{Razlog:} Element drugogo roda i eksponenta pokračujut terminologiju Steinitza iz odděla 2. Izvorni razlog razděljuje eksponentu neseparabilnosti v karakteristikě $p$ od obyčnogo pokazatelja stepeni. Zaglavje spojuje dva termina, ne predlaga ih kako sinonimy. Samo slovo eksponenta bez definicije ne pokazuje, ktory iz tih smyslov jest potrěbny.\par
\noindent\emph{Reason:} Element drugogo roda and eksponenta continue the Steinitz terminology of Section 2. The record distinguishes an inseparability exponent in characteristic $p$ from an ordinary power index. The combined heading groups two terms; it does not propose them as synonyms. Eksponenta alone, without its definition, cannot establish which sense is intended.\par
\noindent{\footnotesize Izvorne zapisy / Original records: HIST-g128-004.}\par
\end{minipage}\par
\par\medskip\Needspace{7\baselineskip}
\noindent\begin{minipage}{\linewidth}
\hypertarget{TERM-0463}{}
\noindent\textbf{elementarne kovarianty}\quad{\footnotesize\texttt{TERM-0463}}\par
\noindent\emph{Iztočnik / Source:} Elementarkovarianten\par
\noindent\emph{Rusky / Russian:} элементарные коварианты\par
\noindent\emph{Ukrajinsky / Ukrainian:} елементарні коваріанти\par
\smallskip\noindent\emph{Razlog:} Elementarne kovarianty zachovuju medžunarodnu terminologiju teorije invariantov. Tri zapisy povezuju izbor s konkretnym slučajem Clebschovoj formuly i s doslědnostju medžu oddělami. Razlog odbacaje dulže slovjansko opisanje kako nepomagajuče jasnosti, ale ne dava točnu alternativnu frazu ni čitatelsky test. Ruska i ukrajinska forma upotrěbjajut iste medžunarodne osnovy.\par
\noindent\emph{Reason:} Elementarne kovarianty retains international invariant-theory terminology. Three records connect the choice with the particular Clebsch-formula case and continuity across sections. The rationale rejects a longer Slavic paraphrase as unhelpful to clarity, but supplies neither its exact wording nor a reader test. The Russian and Ukrainian forms use the same international stems.\par
\noindent{\footnotesize Izvorne zapisy / Original records: HIST-g004-003, HIST-g034-010, HIST-g038-009.}\par
\end{minipage}\par
\par\medskip\Needspace{7\baselineskip}
\noindent\begin{minipage}{\linewidth}
\hypertarget{TERM-0464}{}
\noindent\textbf{elementarne dělitelje}\quad{\footnotesize\texttt{TERM-0464}}\par
\noindent\emph{Iztočnik / Source:} Elementarteiler\par
\noindent\emph{Rusky / Russian:} элементарные делители\par
\noindent\emph{Ukrajinsky / Ukrainian:} елементарні дільники\par
\smallskip\noindent\emph{Razlog:} Elementarne dělitelje jest množinska forma v teoremu o netrivialnyh elementarnyh děliteljah po pridavanju proměnnyh. Razlog pokračuje terminologiju ranějših rabot; ne zaměnjaje elementarny dělitelj proizvoljnym děliteljem. Dva izvorne zapisy zachovuju kandidat, ale samo jedin zachovany razlog objasnjava jego izbor. Množina sama ne dava novu matematičnu kategoriju.\par
\noindent\emph{Reason:} Elementarne dělitelje is the plural wording in a theorem about nontrivial elementary divisors after adjoining variables. The rationale continues the terminology of earlier papers; it does not replace an elementary divisor with an arbitrary divisor. Two original records preserve the candidate, but only one surviving rationale explains it. The plural is not itself a new mathematical category.\par
\noindent{\footnotesize Izvorne zapisy / Original records: HIST-g088-008, HIST-g092-039.}\par
\end{minipage}\par
\par\medskip\Needspace{7\baselineskip}
\noindent\begin{minipage}{\linewidth}
\hypertarget{TERM-0465}{}
\noindent\textbf{elementarny dělitelj}\quad{\footnotesize\texttt{TERM-0465}}\par
\noindent\emph{Iztočnik / Source:} Elementarteiler\par
\noindent\emph{Rusky / Russian:} элементарный делитель\par
\noindent\emph{Ukrajinsky / Ukrainian:} елементарний дільник\par
\smallskip\noindent\emph{Razlog:} Elementarny dělitelj jest jedinstvena forma iste terminologične semje kako TERM-0464. Dva zachovane razloga izrično povezuju ju s teoremom o normalnoj formě matrice i s ranějšej upotrěboju. Tri zapisy ne sut tri nezavisne jezyčne potvrđanja: jeden kandidat može byti prěneseny medžu slovničnymi popisami. Matematičny smysl trěbuje lokalnu definiciju, ne samo slovotvorno sravnjenje.\par
\noindent\emph{Reason:} Elementarny dělitelj is the singular member of the same terminological family as TERM-0464. Two surviving rationales connect it explicitly with the matrix normal-form theorem and earlier usage. Three records are not three independent linguistic attestations: a candidate can be carried between glossary lists. Its mathematical sense requires the local definition, not merely comparison of word formation.\par
\noindent{\footnotesize Izvorne zapisy / Original records: HIST-g085-005, HIST-g092-005, HIST-g093-009.}\par
\end{minipage}\par
\par\medskip\Needspace{7\baselineskip}
\noindent\begin{minipage}{\linewidth}
\hypertarget{TERM-0466}{}
\noindent\textbf{forma elementarnyh děliteljev}\quad{\footnotesize\texttt{TERM-0466}}\par
\noindent\emph{Iztočnik / Source:} Elementarteilerform\par
\noindent\emph{Rusky / Russian:} форма элементарных делителей\par
\noindent\emph{Ukrajinsky / Ukrainian:} форма елементарних дільників\par
\smallskip\noindent\emph{Razlog:} Forma elementarnyh děliteljev imenuje cěly invariant v rabotě 24, ne jedin elementarny dělitelj. Genitivna konstrukcija zachovuje odnos formy k děliteljam i odpověda ruskomu i ukrajinskomu sravnjenju. Kratky razlog potvrđuje kontinuitet nazvanja, ale ne dava pravo zaměniti formu proizvoljnoju drugoju normalnoju formoju.\par
\noindent\emph{Reason:} Forma elementarnyh děliteljev names the whole invariant in Paper 24, not one elementary divisor. The genitive construction retains the relation of the form to its divisors and parallels the Russian and Ukrainian comparanda. The brief rationale establishes continuity of naming, not permission to substitute an arbitrary different normal form.\par
\noindent{\footnotesize Izvorne zapisy / Original records: HIST-g095-014.}\par
\end{minipage}\par
\par\medskip\Needspace{7\baselineskip}
\noindent\begin{minipage}{\linewidth}
\hypertarget{TERM-0467}{}
\noindent\textbf{forma elementarnyh děliteljev / eksponenta}\quad{\footnotesize\texttt{TERM-0467}}\par
\noindent\emph{Iztočnik / Source:} Elementarteilerform / Exponent\par
\noindent\emph{Rusky / Russian:} форма элементарных делителей / экспонента\par
\noindent\emph{Ukrajinsky / Ukrainian:} форма елементарних дільників / експонента\par
\smallskip\noindent\emph{Razlog:} Forma elementarnyh děliteljev i eksponenta sut dva povezane termina iz raboty 24. Razlog ostavjaje eksponenta za osobny Steinitzov smysl v karakteristikě $p$, a za obyčne stepeni upominaje potęga. To jest zabilěžena smyslova razlika, ne obče pravilo zaměny vsake pojavy Exponent. Etimologični zapis potęga jest zachovany v sravnjenju; sam po sobě ne ustanavjaje aktualny osnovny pravopis v cělom prěvodu.\par
\noindent\emph{Reason:} Forma elementarnyh děliteljev and eksponenta are two associated terms from Paper 24. The rationale reserves eksponenta for the particular Steinitz sense in characteristic $p$ and mentions potęga for ordinary powers. This is a recorded semantic distinction, not a universal replacement rule for every occurrence of Exponent. The etymological spelling potęga is retained in the comparison; it does not by itself establish the current basic spelling throughout the translation.\par
\noindent{\footnotesize Izvorne zapisy / Original records: HIST-g097-009.}\par
\end{minipage}\par
\par\medskip\Needspace{7\baselineskip}
\noindent\begin{minipage}{\linewidth}
\hypertarget{TERM-0468}{}
\noindent\textbf{forma elementarnyh děliteljev / norma}\quad{\footnotesize\texttt{TERM-0468}}\par
\noindent\emph{Iztočnik / Source:} Elementarteilerform / Norm\par
\noindent\emph{Rusky / Russian:} форма элементарных делителей / норма\par
\noindent\emph{Ukrajinsky / Ukrainian:} форма елементарних дільників / норма\par
\smallskip\noindent\emph{Razlog:} Forma elementarnyh děliteljev i norma ostavajut dva različne invarianty. Dva razloga povezuju osnovu dělitelj s prozračnostju slovotvorjenja, a norma s medžunarodnoju terminologijeju. Odděl sravnjava, kogda tute invarianty sut proste, primarne, ravne ili povezane stepenju. Zapis ih v jednom zaglavju ne smije zatriti te razlike ni predstaviti ih kako sinonimy.\par
\noindent\emph{Reason:} Forma elementarnyh děliteljev and norma remain two different invariants. The two rationales associate dělitelj with transparent word formation and norma with international terminology. The section compares when these invariants are simple, primary, equal, or related by a power. Grouping them under one heading must not erase those distinctions or present them as synonyms.\par
\noindent{\footnotesize Izvorne zapisy / Original records: HIST-g096-011, HIST-g100-002.}\par
\end{minipage}\par
\par\medskip\Needspace{7\baselineskip}
\noindent\begin{minipage}{\linewidth}
\hypertarget{TERM-0469}{}
\noindent\textbf{sistem elementarnyh děliteljev}\quad{\footnotesize\texttt{TERM-0469}}\par
\noindent\emph{Iztočnik / Source:} Elementarteilersystem\par
\noindent\emph{Rusky / Russian:} система элементарных делителей\par
\noindent\emph{Ukrajinsky / Ukrainian:} система елементарних дільників\par
\smallskip\noindent\emph{Razlog:} Sistem elementarnyh děliteljev zachovuje informaciju, že klasu opisuju uporedžene elementarne dělitelje. Izostaviti sistem ili govoriti samo o děliteljah moglo by skryti strukturu popisa. Razlog izrično upominaje poredok; zato sravnjenje ne može zaměniti sistem prostym neuporedženym množestvom bez prověrky definicije.\par
\noindent\emph{Reason:} Sistem elementarnyh děliteljev retains the information that a class is described by an ordered system of elementary divisors. Omitting sistem or mentioning only divisors could conceal the structure of that list. The rationale explicitly mentions ordering; the system must not be replaced with an unordered set without checking the definition.\par
\noindent{\footnotesize Izvorne zapisy / Original records: HIST-g081-098.}\par
\end{minipage}\par
\par\medskip\Needspace{7\baselineskip}
\noindent\begin{minipage}{\linewidth}
\hypertarget{TERM-0470}{}
\noindent\textbf{teorija elementarnyh děliteljev}\quad{\footnotesize\texttt{TERM-0470}}\par
\noindent\emph{Iztočnik / Source:} Elementarteilertheorie\par
\noindent\emph{Rusky / Russian:} теория элементарных делителей\par
\noindent\emph{Ukrajinsky / Ukrainian:} теорія елементарних дільників\par
\smallskip\noindent\emph{Razlog:} Teorija elementarnyh děliteljev jest ime matematičnej oblasti, iz ktorej odděl 12 dava razrabotany priměr. Dva razloga davajut prednost ustaljenomu nazvanju prěd opisnym prěkazom. Oni ne zachovuju točnu odbacenu alternativu. Izbor prodolžuje semju elementarny dělitelj i ne stvara novu teoriju samym prěvodom nazvanja.\par
\noindent\emph{Reason:} Teorija elementarnyh děliteljev names the mathematical subject from which Section 12 takes a worked example. Two rationales prefer the established name to a descriptive paraphrase. They do not preserve the exact rejected alternative. The choice continues the elementary-divisor family; translating the title does not introduce a new theory.\par
\noindent{\footnotesize Izvorne zapisy / Original records: HIST-g081-017, HIST-g082-017.}\par
\end{minipage}\par
\par\medskip\Needspace{7\baselineskip}
\noindent\begin{minipage}{\linewidth}
\hypertarget{TERM-0471}{}
\noindent\textbf{elementno komutuju?e podkolca; komutantny teorem}\quad{\footnotesize\texttt{TERM-0471}}\par
\noindent\emph{Iztočnik / Source:} elementenweise vertauschbare Unterringe; Kommutantensatz\par
\noindent\emph{Rusky / Russian:} [Bukvy ne sut zachovane / Lettering not recoverable]\par
\noindent\emph{Ukrajinsky / Ukrainian:} [Bukvy ne sut zachovane / Lettering not recoverable]\par
\smallskip\noindent\emph{Razlog:} Zachovany kandidat elementno komutuju?e podkolca ima poškodenju bukvu, a ruske i ukrajinske formy sut izgubljene. Razlog podpiraje medžunarodnu osnovu komut- v tehničnom smyslu komutanta ili centralizatora, ne obyčno putovanje. Elementno označaje uslovje na elementah podprstenov; ne znači, že dva podprstena sut ravne. Tuto objasnjenje ne rekonstruuje izgubljene bukvy ni ne proglašaje kandidat konečnym pravopisnym izborom.\par
\noindent\emph{Reason:} The preserved candidate elementno komutuju?e podkolca contains damaged lettering, and the Russian and Ukrainian forms are lost. The rationale supports the international komut- stem in the technical commutant or centralizer sense, not ordinary commuting or travel. Elementno concerns a condition on elements of the subrings; it does not say the two subrings are equal. This explanation neither reconstructs the lost letters nor certifies the candidate as the final spelling.\par
\noindent{\footnotesize Izvorne zapisy / Original records: HIST-g210-007.}\par
\end{minipage}\par
\par\medskip\Needspace{7\baselineskip}
\noindent\begin{minipage}{\linewidth}
\hypertarget{TERM-0472}{}
\noindent\textbf{eliminacija}\quad{\footnotesize\texttt{TERM-0472}}\par
\noindent\emph{Iztočnik / Source:} Elimination\par
\noindent\emph{Rusky / Russian:} исключение; элиминация\par
\noindent\emph{Ukrajinsky / Ukrainian:} елімінація\par
\smallskip\noindent\emph{Razlog:} Eliminacija zachovuje medžunarodno ime algebraičnej operacije izlučenja proměnnyh. Dva zapisy davajut ruske formy исключение i элиминация, pri čem ranějši razlog upominaje prvu kako obyčno matematično ime dějanja. To ne jest dokaz, že vse sravniteljne jezyky izbrali isti koren. Ukrajinska i medžuslovjanska forma zbližajut se medžunarodnoju osnovoju.\par
\noindent\emph{Reason:} Eliminacija retains the international name for the algebraic operation of eliminating variables. The two records supply Russian исключение and элиминация; the earlier reason describes the former as the ordinary mathematical action noun. This does not establish that all comparison languages selected the same root. The Ukrainian and Interslavic forms share the international stem.\par
\noindent{\footnotesize Izvorne zapisy / Original records: HIST-g007-004, HIST-g044-008.}\par
\end{minipage}\par
\par\medskip\Needspace{7\baselineskip}
\noindent\begin{minipage}{\linewidth}
\hypertarget{TERM-0473}{}
\noindent\textbf{teorija eliminacije}\quad{\footnotesize\texttt{TERM-0473}}\par
\noindent\emph{Iztočnik / Source:} Eliminationstheorie\par
\noindent\emph{Rusky / Russian:} теория исключения; теория элиминации\par
\noindent\emph{Ukrajinsky / Ukrainian:} теорія елімінації; теорія елімінування\par
\smallskip\noindent\emph{Razlog:} Teorija eliminacije identifikuje imenovanu algebraičnu oblast. Četyri razloga podpirajut medžunarodnu osnovu v medžuslovjanskom, ale različno opisuju rusku i ukrajinsku upotrěbu: zapisy imajut исключение i элиминация, takože елімінація i елімінування. Zato historiju ne možno predstaviti kako jedinstveno glasovanje za isti koren. Tvrdženje, že slovjansko opisanje bylo by menje jasno, jest zabilěženo razsudžanje, ne izměren rezultat čitanja.\par
\noindent\emph{Reason:} Teorija eliminacije identifies a named algebraic subject. Four rationales support the international stem in Interslavic but describe Russian and Ukrainian usage inconsistently: the records contain исключение and элиминация, as well as елімінація and елімінування. Their history cannot therefore be presented as a unanimous vote for one root. The claim that a Slavic paraphrase would be less clear is recorded judgment, not a measured reading result.\par
\noindent{\footnotesize Izvorne zapisy / Original records: HIST-g080-001, HIST-g093-003, HIST-g095-001, HIST-g102-001.}\par
\end{minipage}\par
\par\medskip\Needspace{7\baselineskip}
\noindent\begin{minipage}{\linewidth}
\hypertarget{TERM-0474}{}
\noindent\textbf{konečno mnogo ostatkovyh klasov linearno nezavisnyh}\quad{\footnotesize\texttt{TERM-0474}}\par
\noindent\emph{Iztočnik / Source:} endlich viele linear unabhängige Restklassen\par
\noindent\emph{Rusky / Russian:} конечное число линейно независимых классов вычетов\par
\noindent\emph{Ukrajinsky / Ukrainian:} скінченно багато лінійно незалежних класів лишків\par
\smallskip\noindent\emph{Razlog:} Konečno mnogo ostatkovyh klasov linearno nezavisnyh zachovuje i čisleno ograničenje i linearnu nezavisnost. Postavjenje kvalifikacije po imenniku bylo izbrano radi čitanja dolgih dokaznyh rečenij. To ne dava pravo prenesti nezavisnost na drugy objekt ili izpustiti ju. Razlog povezuje frazu s konečnorazměrnym linearnym argumentom v sistemu ostatkovyh klasov, ne s konečnostju vsih elementov.\par
\noindent\emph{Reason:} Konečno mnogo ostatkovyh klasov linearno nezavisnyh retains both the finite-number restriction and linear independence. Placing the qualification after the noun was chosen for readability in long proof sentences. This does not permit independence to be attached to a different object or omitted. The rationale connects the phrase with finite-dimensional linear algebra in a residue-class system, not finiteness of all its elements.\par
\noindent{\footnotesize Izvorne zapisy / Original records: HIST-g096-015.}\par
\end{minipage}\par
\par\medskip\Needspace{7\baselineskip}
\noindent\begin{minipage}{\linewidth}
\hypertarget{TERM-0475}{}
\noindent\textbf{konečna / beskonečna kontinuirna grupa}\quad{\footnotesize\texttt{TERM-0475}}\par
\noindent\emph{Iztočnik / Source:} endliche / unendliche kontinuierliche Gruppe\par
\noindent\emph{Rusky / Russian:} конечная / бесконечная непрерывная группа\par
\noindent\emph{Ukrajinsky / Ukrainian:} скінченна / нескінченна неперервна група\par
\smallskip\noindent\emph{Razlog:} Konečna i beskonečna kontinuirna grupa trěbujut historijni smysl definicije. Sam razlog upozornjava, že razlika tyče se parametrov ili funkcij, ne prosto čisla elementov grupy. Njegovo nazvanje vyraženja sovrěmennoju terminologijeju ne dostavaje za odčitanje konečna kako konečna množina. V slovniku mora ostati razlika medžu toj upotrěboju i običnym konečnym grupovym porjadkom.\par
\noindent\emph{Reason:} Konečna and beskonečna kontinuirna grupa require the historical sense supplied by the definition. The rationale itself warns that the distinction concerns parameters or functions, not simply the number of group elements. Its description of the wording as modern terminology does not justify reading konečna as a finite set. The dictionary must distinguish this usage from ordinary finite group order.\par
\noindent{\footnotesize Izvorne zapisy / Original records: HIST-g075-004.}\par
\smallskip\noindent\emph{Sravnjenje s tekstom:} Rabota 13, odděl 1, definije konečnu kontinuirnu grupu analitičnoju zavisimostju od $\rho$ sučstvenyh parametrov. Beskonečna kontinuirna grupa zavisi od $\rho$ libovoljnyh funkcij i ih izvodnyh. Tekst osobno razlikuje i graničny slučaj beskonečno mnogyh parametrov bez libovoljnyh funkcij, i směšany slučaj. Zato kratka para konečna/beskonečna ne znači prosto dva slučaja čislenoj moči množiny transformacij. Sravnjenje definicije utvrđuje osobny smysl frazy; ne zaměnjaje prověrku vseh jej pojavjenj.\par
\noindent\emph{Textual comparison:} Paper 13, section 1 defines a finite continuous group by analytic dependence on $\rho$ essential parameters. An infinite continuous group depends on $\rho$ arbitrary functions and their derivatives. The text separately distinguishes the boundary case of infinitely many parameters without arbitrary functions, and a mixed case. Thus the short finite/infinite pair does not simply distinguish two cardinalities of the set of transformations. Comparing the definition establishes the phrase's particular sense; it does not yet validate every occurrence.\par
\noindent{\footnotesize TEXT-COMP-11; pasaže / passages: 9201; \texttt{CORPUS\_FORM\_COMPARISONS.json}.}\par
\end{minipage}\par
\par\medskip\Needspace{7\baselineskip}
\noindent\begin{minipage}{\linewidth}
\hypertarget{TERM-0476}{}
\noindent\textbf{konečna / beskonečna grupa ostatkov}\quad{\footnotesize\texttt{TERM-0476}}\par
\noindent\emph{Iztočnik / Source:} endliche / unendliche Restgruppe\par
\noindent\emph{Rusky / Russian:} конечная / бесконечная группа остатков\par
\noindent\emph{Ukrajinsky / Ukrainian:} скінченна / нескінченна група лишків\par
\smallskip\noindent\emph{Razlog:} Konečna grupa ostatkov v tutom zapisu znači grupu s konečnym linearnym bazisom ostatkovyh klasov. Razlog izrično upira konečnost na bazis, ne na čislo vsih klasov. Nazvanje grupa ostatkov ostavaje vidimo vměsto tihoj zaměny sovrěmennym opisom faktor-prostranstva. Takvo zachovanje avtorovogo objekta trěbuje i objasnjenje osobnogo smysla konečnosti.\par
\noindent\emph{Reason:} Konečna grupa ostatkov here means a group with a finite linear basis of residue classes. The reason explicitly locates finiteness in the basis, not in the number of all classes. Grupa ostatkov keeps the author's named object visible instead of silently replacing it with modern quotient-vector-space terminology. Preserving that object also requires explaining the particular sense of finiteness.\par
\noindent{\footnotesize Izvorne zapisy / Original records: HIST-g079-062.}\par
\end{minipage}\par
\par\medskip\Needspace{7\baselineskip}
\noindent\begin{minipage}{\linewidth}
\hypertarget{TERM-0477}{}
\noindent\textbf{konečna Abelova grupa}\quad{\footnotesize\texttt{TERM-0477}}\par
\noindent\emph{Iztočnik / Source:} endliche Abelsche Gruppe\par
\noindent\emph{Rusky / Russian:} конечная абелева группа\par
\noindent\emph{Ukrajinsky / Ukrainian:} скінченна абелева група\par
\smallskip\noindent\emph{Razlog:} Konečna Abelova grupa zachovuje oba kvalifikatora i raspoznavajemo ime Abela. Razlog bira imenovany priděvnik prěd menje prozračnym kalkom, ale ne zachovuje točnu odbacenu formu. Rusky i ukrajinsky sravnitelny izraz imajut takovy poredok priděvnikov. To jest razlog imenovanja i registra, ne novo tvrđenje o struktuře vseh konečnyh grup.\par
\noindent\emph{Reason:} Konečna Abelova grupa retains both qualifications and the recognizable name Abel. The reason prefers the named adjective to a less transparent calque but does not preserve the rejected wording. The Russian and Ukrainian comparanda use the corresponding adjective order. This is an argument about naming and register, not a new assertion about the structure of all finite groups.\par
\noindent{\footnotesize Izvorne zapisy / Original records: HIST-g200-001.}\par
\end{minipage}\par
\par\medskip\Needspace{7\baselineskip}
\noindent\begin{minipage}{\linewidth}
\hypertarget{TERM-0478}{}
\noindent\textbf{konečno algebraično razširenje}\quad{\footnotesize\texttt{TERM-0478}}\par
\noindent\emph{Iztočnik / Source:} endliche algebraische Erweiterung\par
\noindent\emph{Rusky / Russian:} конечное алгебраическое расширение\par
\noindent\emph{Ukrajinsky / Ukrainian:} скінченне алгебраїчне розширення\par
\smallskip\noindent\emph{Razlog:} Konečno algebraično razširenje izrično zachovuje uslovje, pod ktorym zabilěženy princip prěnosa ostavaje važny. Izpustiti konečno ili algebraično by razširilo tvrđenje izvan njegovyh hipotez. V jezyku razširenij konečnost ne možno bez konteksta razuměti kako konečno čislo elementov cělogo polja. Zapisu razširenje odgovarja i druga zabilěžena forma razširjenje v TERM-0485; ovde nema sravniteljnogo razloga za izbor medžu njima.\par
\noindent\emph{Reason:} Konečno algebraično razširenje explicitly retains the condition under which the recorded transfer principle remains valid. Omitting either finite or algebraic would extend the assertion beyond its hypotheses. In extension terminology, finiteness cannot without context be read as a finite number of elements in the entire field. This record's razširenje also has the recorded variant razširjenje in TERM-0485; no comparative reason for choosing between them is supplied here.\par
\noindent{\footnotesize Izvorne zapisy / Original records: HIST-g046-009.}\par
\end{minipage}\par
\par\medskip\Needspace{7\baselineskip}
\noindent\begin{minipage}{\linewidth}
\hypertarget{TERM-0479}{}
\noindent\textbf{konečny bazis diferencialnyh izrazov}\quad{\footnotesize\texttt{TERM-0479}}\par
\noindent\emph{Iztočnik / Source:} endliche Basis der Differentialausdruecke\par
\noindent\emph{Rusky / Russian:} конечный базис дифференциальных выражений\par
\noindent\emph{Ukrajinsky / Ukrainian:} скінченний базис диференціальних виразів\par
\smallskip\noindent\emph{Razlog:} Konečny bazis diferencialnyh izrazov povezuje teorem o konečnom bazisu s operatornymi izrazami. Razlog izrično bira izraz, ne ravnanje: objekty generujut sistem i ne morajut sami byti ravnanja. Konečny opisuje bazis. Imennik bazis ne ustanavjaje sam po sobě linearnu nezavisnost v vsakom historijnom kontekstu; potrěbno jest uslovje dano v tekstu.\par
\noindent\emph{Reason:} Konečny bazis diferencialnyh izrazov connects a finite-basis theorem with operator expressions. The rationale explicitly selects expression rather than equation: the objects generate a system and need not themselves be equations. Konečny qualifies the basis. Bazis does not by itself establish linear independence in every historical context; the condition stated in the text must govern.\par
\noindent{\footnotesize Izvorne zapisy / Original records: HIST-g079-067.}\par
\end{minipage}\par
\par\medskip\Needspace{7\baselineskip}
\noindent\begin{minipage}{\linewidth}
\hypertarget{TERM-0480}{}
\noindent\textbf{konečna grupa}\quad{\footnotesize\texttt{TERM-0480}}\par
\noindent\emph{Iztočnik / Source:} endliche Gruppe\par
\noindent\emph{Rusky / Russian:} конечная группа\par
\noindent\emph{Ukrajinsky / Ukrainian:} скінченна група\par
\smallskip\noindent\emph{Razlog:} Konečna grupa zachovuje obyčno nazvanje grupy s konečnym porjadkom; prvi razlog povezuje ju s oznakoju $h$ za porjadok. Drugy razlog daje prednost obščej slovjanskoj osnově koneč-. Tuta upotrěba ne zaměnjaje konečnu kontinuirnu grupu v historijnom parametričnom smyslu iz TERM-0475. Jednako priděvnično nazvanje ne znači jednaku definiciju vo vseh rabotah.\par
\noindent\emph{Reason:} Konečna grupa retains the ordinary name for a group of finite order; the first rationale connects it with the order symbol $h$. The second favours the shared Slavic koneč- stem. This usage does not replace finite continuous group in the historical parameter sense of TERM-0475. An identical adjective does not imply an identical definition across the papers.\par
\noindent{\footnotesize Izvorne zapisy / Original records: HIST-g056-002, HIST-g204-002.}\par
\end{minipage}\par
\par\medskip\Needspace{7\baselineskip}
\noindent\begin{minipage}{\linewidth}
\hypertarget{TERM-0481}{}
\noindent\textbf{konečna linearna grupa karakteristiki p}\quad{\footnotesize\texttt{TERM-0481}}\par
\noindent\emph{Iztočnik / Source:} endliche lineare Gruppe der Charakteristik p\par
\noindent\emph{Rusky / Russian:} конечная линейная группа характеристики p\par
\noindent\emph{Ukrajinsky / Ukrainian:} скінченна лінійна група характеристики p\par
\smallskip\noindent\emph{Razlog:} Konečna linearna grupa karakteristiki $p$ zachovuje osobnu informaciju o osnovnom polju ili prstenu. Karakteristika ne jest karakter grupy, hoti oba termina stojat blizko v ranějšej rabotě. Razlog zato podpiraje smyslovu granicu medžu dvoma sličnymi medžunarodnymi slovami. Svojstvo osnovnyh skalarov ne trěba prěpisyvati grupovomu karakteru.\par
\noindent\emph{Reason:} Konečna linearna grupa karakteristiki $p$ retains information about the underlying field or ring. Characteristic is not a group character, although the two terms occur near each other in the preceding work. The rationale therefore supports a semantic boundary between similar international words. A property of the underlying scalars must not be assigned to the group character.\par
\noindent{\footnotesize Izvorne zapisy / Original records: HIST-g106-002.}\par
\end{minipage}\par
\par\medskip\Needspace{7\baselineskip}
\noindent\begin{minipage}{\linewidth}
\hypertarget{TERM-0482}{}
\noindent\textbf{konečny porjadok}\quad{\footnotesize\texttt{TERM-0482}}\par
\noindent\emph{Iztočnik / Source:} endliche Ordnung\par
\noindent\emph{Rusky / Russian:} конечный порядок\par
\noindent\emph{Ukrajinsky / Ukrainian:} скінченний порядок\par
\smallskip\noindent\emph{Razlog:} Konečny porjadok zachovuje historijno nazvanje algebraičnogo porjadka v Dedekindovom kontekstu. Razlog izrično odbacaje tihu zaměnu maksimalnym porjadkom ili prstenom cělyh čisel, ktoryh nazvanja iztočnik ne dava. Porjadok ovde ne trěba avtonomno razuměti kako čislo elementov grupy. Kratky zapis ne ustanavjaje sam polnu definiciju priděvnika konečny; ju mora dati odgovarjajuča pasaža.\par
\noindent\emph{Reason:} Konečny porjadok retains the historical name for an algebraic order in the Dedekind context. The rationale explicitly rejects silently substituting maximal order or ring of integers, names not supplied by the source. Porjadok here must not independently be interpreted as the number of elements of a group. The brief record does not itself give the full definition of konečny; the relevant passage must supply it.\par
\noindent{\footnotesize Izvorne zapisy / Original records: HIST-g103-007.}\par
\end{minipage}\par
\par\medskip\Needspace{7\baselineskip}
\noindent\begin{minipage}{\linewidth}
\hypertarget{TERM-0483}{}
\noindent\textbf{konečne porjadky algebraičnogo čislovogo polja}\quad{\footnotesize\texttt{TERM-0483}}\par
\noindent\emph{Iztočnik / Source:} endliche Ordnungen eines algebraischen Zahlkörpers\par
\noindent\emph{Rusky / Russian:} конечные порядки алгебраического числового поля\par
\noindent\emph{Ukrajinsky / Ukrainian:} скінченні порядки алгебраїчного числового поля\par
\smallskip\noindent\emph{Razlog:} Konečne porjadky algebraičnogo čislovogo polja prodolžuju terminologiju raboty 26 v množinskom priměru iz raboty 27. Razlog povezuje priměr s neuspehom predstavjenja primarnogo ideala kako stepeni prostogo ideala. Zato primarny i prosty ne mogut byti sliti v jedno nazvanje. Sam plural ne dava osnovu preimenovati porjadky v maksimalne porjadky ili cěla čislova polja.\par
\noindent\emph{Reason:} Konečne porjadky algebraičnogo čislovogo polja continues Paper 26's terminology in a plural example in Paper 27. The rationale associates the example with failure of a primary ideal's representation as a power of a prime ideal. Primary and prime must therefore remain distinct. The plural does not justify renaming the orders as maximal orders or entire number fields.\par
\noindent{\footnotesize Izvorne zapisy / Original records: HIST-g104-009.}\par
\end{minipage}\par
\par\medskip\Needspace{7\baselineskip}
\noindent\begin{minipage}{\linewidth}
\hypertarget{TERM-0484}{}
\noindent\textbf{konečny p-porjadok}\quad{\footnotesize\texttt{TERM-0484}}\par
\noindent\emph{Iztočnik / Source:} endliche p-Ordnung\par
\noindent\emph{Rusky / Russian:} конечный p-порядок\par
\noindent\emph{Ukrajinsky / Ukrainian:} скінченний p-порядок\par
\smallskip\noindent\emph{Razlog:} Konečny $p$-porjadok prodolžuje osnovu porjadok za algebraičny objekt. Razlog opravdava defis, ktory vidimo veže $p$ s imennikom i ulegčaje čitanje. Točna uloga $p$ i osobny smysl konečnosti zavisat od lokalnej definicije, ne od samogo defisa.\par
\noindent\emph{Reason:} Konečny $p$-porjadok continues porjadok as the name of an algebraic object. The rationale justifies a hyphen that visibly connects $p$ with the noun and assists reading. The precise role of $p$ and the particular sense of finiteness depend on the local definition, not on punctuation alone.\par
\noindent{\footnotesize Izvorne zapisy / Original records: HIST-g161-005.}\par
\end{minipage}\par
\par\medskip\Needspace{7\baselineskip}
\noindent\begin{minipage}{\linewidth}
\hypertarget{TERM-0485}{}
\noindent\textbf{konečno algebraično razširjenje polja}\quad{\footnotesize\texttt{TERM-0485}}\par
\noindent\emph{Iztočnik / Source:} endlicher algebraischer Erweiterungskörper\par
\noindent\emph{Rusky / Russian:} конечное алгебраическое расширение поля\par
\noindent\emph{Ukrajinsky / Ukrainian:} скінченне алгебраїчне розширення поля\par
\smallskip\noindent\emph{Razlog:} Konečno algebraično razširjenje polja prědava německy složeny imennik frazoju s razširjenjem. Razlog identifikuje vloženu kopiju $(P)$ kako osnovno polje; tuta relacija mora ostati vidima. Predmetom konečnosti jest razširjenje, ne proizvoljno čislo elementov razširenogo polja. Izbor sovrěmennoj konstrukcije ne opravdava izgubiti odnos k osnovnomu polju.\par
\noindent\emph{Reason:} Konečno algebraično razširjenje polja renders the German compound noun through an extension phrase. The rationale identifies the embedded copy $(P)$ as the base field; that relation must remain visible. Finiteness concerns the extension, not an arbitrary count of the elements of the larger field. Modernizing the construction does not justify losing its relation to the base.\par
\noindent{\footnotesize Izvorne zapisy / Original records: HIST-g124-006.}\par
\end{minipage}\par
\par\medskip\Needspace{7\baselineskip}
\noindent\begin{minipage}{\linewidth}
\hypertarget{TERM-0486}{}
\noindent\textbf{konečno algebraično čislovo polje}\quad{\footnotesize\texttt{TERM-0486}}\par
\noindent\emph{Iztočnik / Source:} endlicher algebraischer Zahlkörper\par
\noindent\emph{Rusky / Russian:} конечное алгебраическое числовое поле\par
\noindent\emph{Ukrajinsky / Ukrainian:} скінченне алгебраїчне числове поле\par
\smallskip\noindent\emph{Razlog:} Konečno algebraično čislovo polje zachovuje cělu kvalifikaciju, ne samo obšče racionalnostno nazvanje. Zapis pripisuje tu užšu formulaciju pečatnomu popravjenju raboty 11 v rabotě 16. To razlikuje avtorovo popravjenje, kako go opisuje zapis, od novoj redakcijnoj domysly. Tuto objasnjenje ne tvrdi, že pečatno popravjenje bylo ponovno vizualno sravnjeno; ono zachovuje njegov zabilěženy smysl i hipotezy.\par
\noindent\emph{Reason:} Konečno algebraično čislovo polje retains the full qualification rather than a general rationality-domain label. The record attributes the narrower formulation to a printed correction of Paper 11 in Paper 16. This distinguishes an author's correction, as described by the record, from a new editorial conjecture. This explanation does not claim a fresh visual collation of the printed correction; it preserves its reported meaning and hypotheses.\par
\noindent{\footnotesize Izvorne zapisy / Original records: HIST-g078-015.}\par
\end{minipage}\par
\par\medskip\Needspace{7\baselineskip}
\noindent\begin{minipage}{\linewidth}
\hypertarget{TERM-0487}{}
\noindent\textbf{konečna integralna oblast}\quad{\footnotesize\texttt{TERM-0487}}\par
\noindent\emph{Iztočnik / Source:} endlicher Integritätsbereich\par
\noindent\emph{Rusky / Russian:} конечная область целостности\par
\noindent\emph{Ukrajinsky / Ukrainian:} скінченна інтегральна область\par
\smallskip\noindent\emph{Razlog:} Konečna integralna oblast v oběh zachovanyh razlogah znači oblast s konečnym bazisom cělosti, ne oblast s konečnym čislom elementov. Bez toj razlike sovrěmenny čitatelj može dobiti drugy matematičny objekt. Sličnost ruskoj i ukrajinskoj frazy ne odstranja tu dvoznačnost. Historijno nazvanje mora byti objasnjeno njegovym uslovjem, ne mehanično prěvedeno kako konečna množina.\par
\noindent\emph{Reason:} Konečna integralna oblast means, in both surviving rationales, a domain with a finite integrality basis, not a domain with a finite number of elements. Without that distinction a modern reader could infer a different mathematical object. Similar Russian and Ukrainian phrasing does not remove the ambiguity. The historical name must be explained through its condition, not mechanically read as a finite set.\par
\noindent{\footnotesize Izvorne zapisy / Original records: HIST-g050-003, HIST-g052-003.}\par
\smallskip\noindent\emph{Sravnjenje s tekstom:} V rabotě 6, oddělu 10, definicija VI trěbuje konečno mnogo polinomov, iz ktoryh vse polinomy oblasti dobyvajut se cělymi racionalnymi kombinacijami nad koeficientnoju oblastju. Taka oblast jest nazvana konečnoju integralnoju oblastju. V rabotě 30, oddělu 1, konečny modul jest izrično objasnjen kako konečno porodženy; formula $M=(\mu_1,\ldots,\mu_n)$ i slědujuča rečenija davajut linearnu kombinaciju $r_1\mu_1+\cdots+r_n\mu_n$. V oběh slučajah konečnost odnosi se k generujučemu sistemu, ale operacije generovanja ne sut iste. Nijedna definicija ne ustanavjaje konečno čislo vseh elementov ni linearnu nezavisnost proizvoljnogo bazisa.\par
\noindent\emph{Textual comparison:} In Paper 6, section 10, Definition VI requires finitely many polynomials from which every polynomial in the domain is obtained by integral rational combinations over the coefficient domain. Such a domain is called a finite integral domain. In Paper 30, section 1, finite module is explicitly explained as finitely generated; $M=(\mu_1,\ldots,\mu_n)$ and the following sentence give the linear combination $r_1\mu_1+\cdots+r_n\mu_n$. In both cases finiteness concerns a generating system, but the generating operations differ. Neither definition asserts a finite number of all elements or the linear independence of an arbitrary basis.\par
\noindent{\footnotesize TEXT-COMP-12; pasaže / passages: 6061, 16190, 16192, 16194; \texttt{CORPUS\_FORM\_COMPARISONS.json}.}\par
\end{minipage}\par
\par\medskip\Needspace{7\baselineskip}
\noindent\begin{minipage}{\linewidth}
\hypertarget{TERM-0488}{}
\noindent\textbf{konečny, to jest konečno porodženy modul}\quad{\footnotesize\texttt{TERM-0488}}\par
\noindent\emph{Iztočnik / Source:} endlicher Modul\par
\noindent\emph{Rusky / Russian:} конечный, то есть конечно порожденный модуль\par
\noindent\emph{Ukrajinsky / Ukrainian:} скінченний, тобто скінченно породжений модуль\par
\smallskip\noindent\emph{Razlog:} Konečny, to jest konečno porodženy modul jest izrično razjasnjenje historijnogo konečny. Razlog upira definiciju na konečno mnogo generujučih elementov $u_i$, ne na čislo vsih elementov modula. Dodana fraza vysvětljaje smysl pri prvoj pojavě, vměsto tihoj proměny ponjatja. Ona ne tvrdi linearnu nezavisnost generatorov; konečno generovanje i slobodny bazis sut različne uslovja.\par
\noindent\emph{Reason:} Konečny, to jest konečno porodženy modul explicitly explains historical finite. The rationale grounds the definition in finitely many generating elements $u_i$, not in the number of all module elements. The added phrase explains the sense at first occurrence instead of silently changing the concept. It does not assert the generators' linear independence: finite generation and a free basis are different conditions.\par
\noindent{\footnotesize Izvorne zapisy / Original records: HIST-g108-003.}\par
\smallskip\noindent\emph{Sravnjenje s tekstom:} V rabotě 6, oddělu 10, definicija VI trěbuje konečno mnogo polinomov, iz ktoryh vse polinomy oblasti dobyvajut se cělymi racionalnymi kombinacijami nad koeficientnoju oblastju. Taka oblast jest nazvana konečnoju integralnoju oblastju. V rabotě 30, oddělu 1, konečny modul jest izrično objasnjen kako konečno porodženy; formula $M=(\mu_1,\ldots,\mu_n)$ i slědujuča rečenija davajut linearnu kombinaciju $r_1\mu_1+\cdots+r_n\mu_n$. V oběh slučajah konečnost odnosi se k generujučemu sistemu, ale operacije generovanja ne sut iste. Nijedna definicija ne ustanavjaje konečno čislo vseh elementov ni linearnu nezavisnost proizvoljnogo bazisa.\par
\noindent\emph{Textual comparison:} In Paper 6, section 10, Definition VI requires finitely many polynomials from which every polynomial in the domain is obtained by integral rational combinations over the coefficient domain. Such a domain is called a finite integral domain. In Paper 30, section 1, finite module is explicitly explained as finitely generated; $M=(\mu_1,\ldots,\mu_n)$ and the following sentence give the linear combination $r_1\mu_1+\cdots+r_n\mu_n$. In both cases finiteness concerns a generating system, but the generating operations differ. Neither definition asserts a finite number of all elements or the linear independence of an arbitrary basis.\par
\noindent{\footnotesize TEXT-COMP-12; pasaže / passages: 6061, 16190, 16192, 16194; \texttt{CORPUS\_FORM\_COMPARISONS.json}.}\par
\end{minipage}\par
\par\medskip\Needspace{7\baselineskip}
\noindent\begin{minipage}{\linewidth}
\hypertarget{TERM-0489}{}
\noindent\textbf{konečny ideal / idealny bazis}\quad{\footnotesize\texttt{TERM-0489}}\par
\noindent\emph{Iztočnik / Source:} endliches Ideal / Idealbasis\par
\noindent\emph{Rusky / Russian:} конечный идеал / идеальный базис\par
\noindent\emph{Ukrajinsky / Ukrainian:} скінченний ідеал / ідеальний базис\par
\smallskip\noindent\emph{Razlog:} Konečny ideal i idealny bazis zachovuju izričnu hipotezu odděla 1. Razlog odbacaje zaměnu poznějšim nazvanjem Noetherov prsten, ibo ono by potisnulo avtorovo predstavljanje uslovja. To ne zabranjaje sovrěmenno vysvětlenje v osobnom komentaru, ale ono ne može zaměniti iztočnu frazu. Dva zapisy zachovuju kandidat, pri jednom zachovanom razlogu.\par
\noindent\emph{Reason:} Konečny ideal and idealny bazis retain Section 1's explicit hypothesis. The rationale rejects replacement by the later label Noetherian ring because that would displace the author's presentation of the condition. This does not preclude a modern explanation in separate commentary, but it cannot replace the source phrase. Two records preserve the candidate, with one surviving rationale.\par
\noindent{\footnotesize Izvorne zapisy / Original records: HIST-g081-026, HIST-g082-026.}\par
\end{minipage}\par
\par\medskip\Needspace{7\baselineskip}
\noindent\begin{minipage}{\linewidth}
\hypertarget{TERM-0490}{}
\noindent\textbf{konečnost sistema form}\quad{\footnotesize\texttt{TERM-0490}}\par
\noindent\emph{Iztočnik / Source:} Endlichkeit des Formensystems\par
\noindent\emph{Rusky / Russian:} конечность системы форм\par
\noindent\emph{Ukrajinsky / Ukrainian:} скінченність системи форм\par
\smallskip\noindent\emph{Razlog:} Konečnost sistema form zachovuje i svojstvo i objekt, na ktory ono odnosi se. Razlog prodolžuje terminologiju sistema form i poredok slov blizky ruskomu i ukrajinskomu. On ne razviva sam definiciju konečnosti. Zato naslovna fraza ne može sama ustanoviti, je li rěč o čisle form ili o konečnom generovanju; to mora byti povezano s konkretnym tvrđenjem.\par
\noindent\emph{Reason:} Konečnost sistema form retains both the property and the object to which it applies. The rationale continues the form-system terminology and a word order close to Russian and Ukrainian. It does not itself develop the definition of finiteness. The heading alone therefore cannot decide between a finite number of forms and finite generation; the specific assertion must supply that connection.\par
\noindent{\footnotesize Izvorne zapisy / Original records: HIST-g028-002.}\par
\end{minipage}\par
\par\medskip\Needspace{7\baselineskip}
\noindent\begin{minipage}{\linewidth}
\hypertarget{TERM-0491}{}
\noindent\textbf{konečnost v užšem smyslu}\quad{\footnotesize\texttt{TERM-0491}}\par
\noindent\emph{Iztočnik / Source:} Endlichkeit im engeren Sinn\par
\noindent\emph{Rusky / Russian:} конечность в более узком смысле\par
\noindent\emph{Ukrajinsky / Ukrainian:} скінченність у вужчому сенсі\par
\smallskip\noindent\emph{Razlog:} Konečnost v užšem smyslu zachovuje avtorovo izrično ograničenje smysla. Razlog protistavjaje tu užšu konečnost postojanju konečnyh bazisov. Izostaviti v užšem smyslu by skrylo imenovanu razliku, ne prosto skratilo rečenju. Sam zapis ne dava cělu definiciju užšego uslovja, zato ju ne dopolnjamo domysloju.\par
\noindent\emph{Reason:} Konečnost v užšem smyslu retains the author's explicit restriction of meaning. The reason contrasts this narrower finiteness with the existence of finite bases. Omitting v užšem smyslu would conceal a named distinction, not merely shorten the sentence. The record does not supply the narrower condition's complete definition, so none is invented here.\par
\noindent{\footnotesize Izvorne zapisy / Original records: HIST-g040-012.}\par
\end{minipage}\par
\par\medskip\Needspace{7\baselineskip}
\noindent\begin{minipage}{\linewidth}
\hypertarget{TERM-0492}{}
\noindent\textbf{uslovje konečnosti / idealny bazis}\quad{\footnotesize\texttt{TERM-0492}}\par
\noindent\emph{Iztočnik / Source:} Endlichkeitsbedingung / Idealbasis\par
\noindent\emph{Rusky / Russian:} условие конечности / идеальный базис\par
\noindent\emph{Ukrajinsky / Ukrainian:} умова скінченності / ідеальний базис\par
\smallskip\noindent\emph{Razlog:} Uslovje konečnosti i idealny bazis vyražajut hipotezu konečnogo bazisa ideala. Razlog izrično čuva bazisnu formulaciju vměsto retrospektivnogo nazvanja noetherovosti. Važno jest i čije idealne bazisy ispolnjajut uslovje: skračenje ne smije izgubiti kvantifikaciju iz iztočnogo tvrđenja. Zapis podpiraje vernost formulaciji, ne zabranju vsi sovrěmenne objasnjujuče paralely.\par
\noindent\emph{Reason:} Uslovje konečnosti and idealny bazis express the finite ideal-basis hypothesis. The rationale expressly preserves the basis formulation instead of retrospectively naming Noetherianity. Which ideals must satisfy the condition also matters: abbreviation must not lose the source assertion's quantification. The record supports fidelity to the formulation, not a prohibition on all explanatory modern parallels.\par
\noindent{\footnotesize Izvorne zapisy / Original records: HIST-g081-012, HIST-g082-012.}\par
\end{minipage}\par
\par\medskip\Needspace{7\baselineskip}
\noindent\begin{minipage}{\linewidth}
\hypertarget{TERM-0493}{}
\noindent\textbf{uslovje konečnosti; konečny modul; modulny bazis}\quad{\footnotesize\texttt{TERM-0493}}\par
\noindent\emph{Iztočnik / Source:} Endlichkeitsbedingung; endlicher Modul; Modulbasis\par
\noindent\emph{Rusky / Russian:} условие конечности; конечный модуль; базис модуля\par
\noindent\emph{Ukrajinsky / Ukrainian:} умова скінченності; скінченний модуль; база модуля\par
\smallskip\noindent\emph{Razlog:} Uslovje konečnosti; konečny modul; modulny bazis prodolžuju idealno uslovje v modulnom kontekstu odděla 9. Razlog bira medžunarodno bazis, pri ruskom базис модуля i ukrajinskom база модуля. Tute sravniteljne formy podpirajut imennik, ale ne dokazujut, že vsak takovy bazis mora byti linearno nezavisny. Smysl generovanja i točne hipotezy mora nositi definicija modula v tekstu.\par
\noindent\emph{Reason:} Uslovje konečnosti; konečny modul; modulny bazis carries the ideal condition into the module setting of Section 9. The rationale chooses international bazis alongside Russian базис модуля and Ukrainian база модуля. These comparanda support the noun but do not establish that every such basis must be linearly independent. The text's definition must supply the generating sense and exact hypotheses.\par
\noindent{\footnotesize Izvorne zapisy / Original records: HIST-g081-072.}\par
\end{minipage}\par
\par\medskip\Needspace{7\baselineskip}
\noindent\begin{minipage}{\linewidth}
\hypertarget{TERM-0494}{}
\noindent\textbf{kriterij konečnosti}\quad{\footnotesize\texttt{TERM-0494}}\par
\noindent\emph{Iztočnik / Source:} Endlichkeitskriterium\par
\noindent\emph{Rusky / Russian:} критерий конечности\par
\noindent\emph{Ukrajinsky / Ukrainian:} критерій скінченності\par
\smallskip\noindent\emph{Razlog:} Kriterij konečnosti odgovarja imenovanju teorema i označaje uslovja za konečno generovanje. Razlog upominaje podprsten i cělost nad njim; ne daje v toj kratkoj primětce cělu teoremu so vsimi kvantifikatorami. Zato konečny podprsten iz englijskogo opisa ne možno samovolno čitati kako konečnu množinu. Slovnična formula mora voditi k polnomu tvrđenju, ne zaměnjati go.\par
\noindent\emph{Reason:} Kriterij konečnosti parallels the theorem's name and denotes conditions for finite generation. The rationale mentions a subring and integrality over it; its short note is not the whole theorem with all quantifiers. Its English phrase finite subring therefore cannot independently be read as a finite set. The dictionary wording must lead to the full assertion, not replace it.\par
\noindent{\footnotesize Izvorne zapisy / Original records: HIST-g106-003.}\par
\end{minipage}\par
\par\medskip\Needspace{7\baselineskip}
\noindent\begin{minipage}{\linewidth}
\hypertarget{TERM-0495}{}
\noindent\textbf{dokaz konečnosti sistema form}\quad{\footnotesize\texttt{TERM-0495}}\par
\noindent\emph{Iztočnik / Source:} Endlichkeitsnachweis des Formensystems\par
\noindent\emph{Rusky / Russian:} доказательство конечности системы форм\par
\noindent\emph{Ukrajinsky / Ukrainian:} доведення скінченності системи форм\par
\smallskip\noindent\emph{Razlog:} Dokaz konečnosti sistema form zachovuje dokaznu funkciju imennika Nachweis. Ne imenuje samo sistem form ni samo svojstvo konečnosti, ale argument za to svojstvo. Razlog prodolžuje ranějše nazvanje sistema, bez predlaganja nove smyslove varijante. Vrsta konečnosti ostavaje oprědělena samym dokazom.\par
\noindent\emph{Reason:} Dokaz konečnosti sistema form retains the argumentative function of Nachweis. It names neither just the system of forms nor just finiteness, but the proof of that property. The rationale continues the earlier form-system name without proposing a new semantic variant. The proof itself determines the relevant kind of finiteness.\par
\noindent{\footnotesize Izvorne zapisy / Original records: HIST-g029-003.}\par
\end{minipage}\par
\par\medskip\Needspace{7\baselineskip}
\noindent\begin{minipage}{\linewidth}
\hypertarget{TERM-0496}{}
\noindent\textbf{teoremy konečnosti}\quad{\footnotesize\texttt{TERM-0496}}\par
\noindent\emph{Iztočnik / Source:} Endlichkeitssaetze\par
\noindent\emph{Rusky / Russian:} теоремы конечности\par
\noindent\emph{Ukrajinsky / Ukrainian:} теореми скінченності\par
\smallskip\noindent\emph{Razlog:} Teoremy konečnosti obsežajut v oddělu 6 tvrđenja o konečnom generovanju i konečnom razloženju. Razlog bira obšču slovjansku osnovu konec-/koneč- prěd latinizovanym kalkom, ale ne zachovuje točno ime odbacene alternativy. Dvě vrste tvrđenij ne sut ista teorema; množina v zaglavju zachovuje ih obem bez stiranija različnyh hipotez.\par
\noindent\emph{Reason:} Teoremy konečnosti covers statements about finite generation and finite decomposition in Section 6. The rationale prefers the shared Slavic konec-/koneč- stem to a Latinized calque but does not preserve the exact rejected alternative. The two kinds of assertion are not the same theorem; the plural heading accommodates them without erasing their different hypotheses.\par
\noindent{\footnotesize Izvorne zapisy / Original records: HIST-g079-031.}\par
\end{minipage}\par
\par\medskip\Needspace{7\baselineskip}
\noindent\begin{minipage}{\linewidth}
\hypertarget{TERM-0497}{}
\noindent\textbf{teoremy konečnosti / Hilbertov teorem o bazisu modula}\quad{\footnotesize\texttt{TERM-0497}}\par
\noindent\emph{Iztočnik / Source:} Endlichkeitssaetze / Hilbertscher Modulsatz\par
\noindent\emph{Rusky / Russian:} теоремы конечности / теорема Гильберта о базисе модуля\par
\noindent\emph{Ukrajinsky / Ukrainian:} теореми скінченності / теорема Гільберта про базис модуля\par
\smallskip\noindent\emph{Razlog:} Teoremy konečnosti i Hilbertov teorem o bazisu modula sut v izvornoj primětce predložne nazvanja iz uvoda, čije dokazy dolžny byti pročitane pozněje. Razlog izrično ostavjaje ih dočasnymi do sravnjenja s argumentami o konečnom generovanju. Samo povtorjenje v slovniku ne proměnja dočasny izbor v prověreny. Potrěbna jest veza s točnym tvrđenjem i njegovym smyslom bazisa.\par
\noindent\emph{Reason:} Teoremy konečnosti and Hilbertov teorem o bazisu modula are provisional introductory names in the original note, whose proofs were to be examined later. The rationale expressly leaves them provisional until comparison with the finite-generation arguments. Repetition in a dictionary does not turn a provisional choice into a verified one. The names need a link to the exact statement and its sense of basis.\par
\noindent{\footnotesize Izvorne zapisy / Original records: HIST-g079-012.}\par
\end{minipage}\par
\par\medskip\Needspace{7\baselineskip}
\noindent\begin{minipage}{\linewidth}
\hypertarget{TERM-0498}{}
\noindent\textbf{teorem o konečnosti}\quad{\footnotesize\texttt{TERM-0498}}\par
\noindent\emph{Iztočnik / Source:} Endlichkeitssatz\par
\noindent\emph{Rusky / Russian:} теорема конечности\par
\noindent\emph{Ukrajinsky / Ukrainian:} теорема про скінченність\par
\smallskip\noindent\emph{Razlog:} Teorem o konečnosti zachovuje ime teorema o konečnom generovanju invariantov. Razlog daje prednost toj prěmoj frazě prěd slobodnym opisom s konečnym bazisom. To ne znači, že vse invarianty obrazuju konečnu množinu. Nazvanje i točno svojstvo mora čitatelj viděti zajedno, osobno kogda sam priděvnik konečny dopuščaje više čitanij.\par
\noindent\emph{Reason:} Teorem o konečnosti retains the name of a theorem about finite generation of invariants. The rationale favours this direct phrase over a loose finite-basis paraphrase. It does not say that all invariants form a finite set. The name and precise property must be presented together, particularly where finite admits several readings.\par
\noindent{\footnotesize Izvorne zapisy / Original records: HIST-g056-001.}\par
\end{minipage}\par
\par\medskip\Needspace{7\baselineskip}
\noindent\begin{minipage}{\linewidth}
\hypertarget{TERM-0499}{}
\noindent\textbf{teorem konečnosti invariantov}\quad{\footnotesize\texttt{TERM-0499}}\par
\noindent\emph{Iztočnik / Source:} Endlichkeitssatz der Invarianten\par
\noindent\emph{Rusky / Russian:} теорема конечности для инвариантов\par
\noindent\emph{Ukrajinsky / Ukrainian:} теорема скінченності для інваріантів\par
\smallskip\noindent\emph{Razlog:} Teorem konečnosti invariantov jest izrično razjasnjen kako konečno generovanje sistema invariantov, ne konečno čislo vsih invariantov. Okolna terminologija cělogo bazisa nosi tuto razliku. Razlog zato podpiraje smyslovu preciznost, ne samo zajedničny koren koneč-. Bez uslovja generovanja kratko zaglavje može byti dvoznačno za sovrěmennogo čitatelja.\par
\noindent\emph{Reason:} Teorem konečnosti invariantov is explicitly explained as finite generation of an invariant system, not a finite number of all invariants. The surrounding integral-basis terminology carries that distinction. The rationale therefore supports semantic precision, not merely a shared koneč- root. Without the generating condition, the short heading can be ambiguous to a modern reader.\par
\noindent{\footnotesize Izvorne zapisy / Original records: HIST-g106-001.}\par
\end{minipage}\par
\par\medskip\Needspace{7\baselineskip}
\noindent\begin{minipage}{\linewidth}
\hypertarget{TERM-0500}{}
\noindent\textbf{energijske relacije / energijske zakony / energijske komponenty}\quad{\footnotesize\texttt{TERM-0500}}\par
\noindent\emph{Iztočnik / Source:} Energierelationen / Energiesätze / Energiekomponenten\par
\noindent\emph{Rusky / Russian:} энергетические соотношения / законы энергии / энергетические компоненты\par
\noindent\emph{Ukrajinsky / Ukrainian:} енергетичні співвідношення / закони енергії / енергетичні компоненти\par
\smallskip\noindent\emph{Razlog:} Energijske relacije, energijske zakony i energijske komponenty zachovuju energiju izrično vo vseh trěh konstrukcijah. Razlog povezuje ih s Hilbertovym tvrđenjem i s Noetherinoju razlicoju medžu sobstvennymi i nesobstvennymi zakonami ohranjenja. Relacija, zakon i komponenta ne sut vzaimno zaměnime imenniki. Obšča priděvnična osnova svjazuje tematiku, ale ne može zatriti matematičnu ulogu vsakoj frazy.\par
\noindent\emph{Reason:} Energijske relacije, energijske zakony, and energijske komponenty keep energy explicit in all three constructions. The rationale links them to Hilbert's claim and Noether's distinction between proper and improper conservation laws. Relation, law, and component are not interchangeable nouns. A shared adjectival stem connects the subject matter without erasing each phrase's mathematical role.\par
\noindent{\footnotesize Izvorne zapisy / Original records: HIST-g075-016.}\par
\end{minipage}\par
\par\medskip\Needspace{7\baselineskip}
\noindent\begin{minipage}{\linewidth}
\hypertarget{TERM-0501}{}
\noindent\textbf{sootvětstvo medžu razloženjem $\mathfrak{T}p$ i razloženjem nulovogo ideala v $\mathfrak{R}$}\quad{\footnotesize\texttt{TERM-0501}}\par
\noindent\emph{Iztočnik / Source:} Entsprechen der Idealzerlegung von $\mathfrak{T}p$ und des Nullideals in $\mathfrak{R}$\par
\noindent\emph{Rusky / Russian:} соответствие между разложением $\mathfrak{T}p$ и разложением нулевого идеала в $\mathfrak{R}$\par
\noindent\emph{Ukrajinsky / Ukrainian:} відповідність між розкладом $\mathfrak{T}p$ і розкладом нульового ідеалу в $\mathfrak{R}$\par
\smallskip\noindent\emph{Razlog:} Sootvětstvo medžu razloženjem $\mathfrak{T}p$ i razloženjem nulovogo ideala v $\mathfrak{R}$ označaje odnos medžu dvoma razloženjami. Razlog povezuje frazu s dokazom iz odděla 7, punkta 2, i prodolžuje terminologiju homomorfizmov. To ne jest tvrdženje, že $\mathfrak{T}p$ i nulovy ideal sut isti objekt v istom prstenu. Identitet prstenov i napravljanje sootvětstva morajut ostati vidime v argumentu.\par
\noindent\emph{Reason:} Sootvětstvo medžu razloženjem $\mathfrak{T}p$ i razloženjem nulovogo ideala v $\mathfrak{R}$ names a relation between two decompositions. The rationale connects it with the proof in Section 7, item 2 and continues the homomorphism terminology. It does not assert that $\mathfrak{T}p$ and the zero ideal are the same object in the same ring. The rings' identities and the direction of the correspondence must remain visible in the argument.\par
\noindent{\footnotesize Izvorne zapisy / Original records: HIST-g154-004.}\par
\smallskip\noindent\emph{Sravnjenje s tekstom:} Rabota 31, odděl 7, izrično prěnosi razloženje ideala $\mathfrak{T}p$ v razloženje nulovogo ideala v $\mathfrak{R}$. Poznějša pasaža piše $\mathfrak{R}=\mathfrak{T}/\mathfrak{T}p$ i upotrěbjaje korespondencija, ne sootvětstvo iz historijnogo glosara. Kratko arhivno Tp zato označaje $\mathfrak{T}p$, ne indeks $T_p$. Lokalni TeX-oprěděljenja pišut $\mathfrak{T}$ i $\mathfrak{R}$ frakturoju. Sravnjenje zachovuje i matematičny odnos i razliku medžu kandidatnym nazvanjem i točno pročitanom aktualnom frazoju.\par
\noindent\emph{Textual comparison:} Paper 31, section 7 explicitly transfers the decomposition of $\mathfrak{T}p$ to the decomposition of the zero ideal in $\mathfrak{R}$. A later passage writes $\mathfrak{R}=\mathfrak{T}/\mathfrak{T}p$ and uses korespondencija, not the historical glossary's sootvětstvo. The archival shorthand Tp therefore denotes $\mathfrak{T}p$, not a subscript $T_p$. The local TeX definitions set $\mathfrak{T}$ and $\mathfrak{R}$ in fraktur. This comparison preserves both the mathematical relation and the difference between the candidate name and the current wording actually examined.\par
\noindent{\footnotesize TEXT-COMP-13; pasaže / passages: 18671, 18751; \texttt{CORPUS\_FORM\_COMPARISONS.json}.}\par
\end{minipage}\par
\par\medskip\Needspace{7\baselineskip}
\noindent\begin{minipage}{\linewidth}
\hypertarget{TERM-0502}{}
\noindent\textbf{odgovarjajuče formy}\quad{\footnotesize\texttt{TERM-0502}}\par
\noindent\emph{Iztočnik / Source:} entsprechenden Formen\par
\noindent\emph{Rusky / Russian:} соответствующие формы\par
\noindent\emph{Ukrajinsky / Ukrainian:} відповідні форми\par
\smallskip\noindent\emph{Razlog:} Odgovarjajuče formy prědava matematično odgovaranje, ne odgovornost. Izvorni zapis izrično protistavjaje tomu ranějše odgovornymi, ktoro moglo dati smysl odgovornosti; za konkretnu rečeniju on dava instrumental odgovarjajučimi formami. Zato izbor vključaje i leksičny smysl i padežnu formu. Kratko slovnično zaglavje ne zaměnjaje formu potrěbnu v rečeniji.\par
\noindent\emph{Reason:} Odgovarjajuče formy expresses mathematical correspondence, not responsibility. The original record explicitly contrasts it with earlier odgovornymi, which could suggest responsibility; for the particular sentence it gives instrumental odgovarjajučimi formami. The choice therefore concerns both lexical meaning and case form. The short dictionary heading does not replace the inflection required in the sentence.\par
\noindent{\footnotesize Izvorne zapisy / Original records: HIST-g025-004.}\par
\smallskip\noindent\emph{Sravnjenje s tekstom:} Dve redukcije v rabotě 2, oddělu 24, imajut točnu konstrukciju s odgovarjajučimi formami: jedna pri simbolah $\theta\cdot H$, druga pri $f\cdot H$. Tute pasaže potvrđajut formu opisanu v zachovanom razlogu, uključajuči padež po s. One ne sut samostatny dokaz ranějšej historije proměny ni cělokorpusna prověrka vsakeho odgovarjajuči. Upotrěby mimo matematičnogo odgovaranja trěbujut vlastny kontekst.\par
\noindent\emph{Textual comparison:} Two reductions in Paper 2, section 24 contain the exact construction s odgovarjajučimi formami: one in the symbols $\theta\cdot H$, the other in $f\cdot H$. These passages confirm the form described by the preserved rationale, including the case after s. They are not independent evidence for the earlier history of the change or a corpus-wide check of every corresponding-form expression. Uses outside mathematical correspondence still require their own context.\par
\noindent{\footnotesize TEXT-COMP-14; pasaže / passages: 3176, 3192; \texttt{CORPUS\_FORM\_COMPARISONS.json}.}\par
\end{minipage}\par
\par\medskip\Needspace{7\baselineskip}
\noindent\begin{minipage}{\linewidth}
\hypertarget{TERM-0503}{}
\noindent\textbf{razloženje obćih form po normalnyh formah}\quad{\footnotesize\texttt{TERM-0503}}\par
\noindent\emph{Iztočnik / Source:} Entwicklung der allgemeinen Formen nach Normalformen\par
\noindent\emph{Rusky / Russian:} разложение общих форм по нормальным формам\par
\noindent\emph{Ukrajinsky / Ukrainian:} розклад загальних форм за нормальними формами\par
\smallskip\noindent\emph{Razlog:} Razloženje obćih form po normalnyh formah čita Entwicklung nach kako matematično razloženje po danih formah, ne kako historijno razvijanje. Razlog bira razloženje i normalne formy radi prozračnosti, ale ostavjaje izbor dočasnym do sravnjenja s avtoritetnym iztočnikom. Tuto objasnjenje zachovuje tu granicu. Ono ne dava novy dokaz za historijnu grafiju obćih v kandidatu.\par
\noindent\emph{Reason:} Razloženje obćih form po normalnyh formah reads Entwicklung nach as mathematical expansion in terms of given forms, not historical development. The rationale selects razloženje and normalne formy for transparency but leaves the choice provisional pending an authoritative comparison. This explanation preserves that qualification. It supplies no new evidence for the historical candidate's spelling obćih.\par
\noindent{\footnotesize Izvorne zapisy / Original records: HIST-g036-001.}\par
\end{minipage}\par
\par\medskip\Needspace{7\baselineskip}
\noindent\begin{minipage}{\linewidth}
\hypertarget{TERM-0504}{}
\noindent\textbf{dopolnjenje algebrai?noj bazy do racionalnoj}\quad{\footnotesize\texttt{TERM-0504}}\par
\noindent\emph{Iztočnik / Source:} Erg?nzung einer algebraischen Basis zu einer rationalen\par
\noindent\emph{Rusky / Russian:} [Bukvy ne sut zachovane / Lettering not recoverable]\par
\noindent\emph{Ukrajinsky / Ukrainian:} [Bukvy ne sut zachovane / Lettering not recoverable]\par
\smallskip\noindent\emph{Razlog:} Dopolnjenje algebrai?noj bazy do racionalnoj ima poškodenju bukvu; německo zaglavje takože jest poškodeno, a ruske i ukrajinske bukvy ne sut zachovane. Davajemy razlog ipak razděljuje dva kroky: dopolniti $H$ do racionalnoj bazy i pridati dopustimy ostatkovy sistem. To ne jest prosta zaměna slova algebraična slovom racionalna. Matematično objasnjenje ne rekonstruuje izgubljenu grafiju.\par
\noindent\emph{Reason:} Dopolnjenje algebrai?noj bazy do racionalnoj contains damaged lettering; the German heading is also damaged, and the Russian and Ukrainian lettering is lost. The surviving reason nevertheless distinguishes two steps: completing $H$ to a rational basis and adjoining an admissible residual system. This is not simply replacing the word algebraic with rational. Explaining the construction does not recover the lost spelling.\par
\noindent{\footnotesize Izvorne zapisy / Original records: HIST-g064-013.}\par
\end{minipage}\par
\par\medskip\Needspace{7\baselineskip}
\noindent\begin{minipage}{\linewidth}
\hypertarget{TERM-0505}{}
\noindent\textbf{dopolnitelno predstavjenje}\quad{\footnotesize\texttt{TERM-0505}}\par
\noindent\emph{Iztočnik / Source:} Ergänzungsdarstellung\par
\noindent\emph{Rusky / Russian:} дополнительное представление\par
\noindent\emph{Ukrajinsky / Ukrainian:} доповнювальне подання\par
\smallskip\noindent\emph{Razlog:} Dopolnitelno predstavjenje označaje dodatne predstavjenje v završnoj kvalifikaciji teksta. Razlog izrično odbacaje zaměnu drugoju algebraičnoju konstrukcijeju bez tekstnogo dokaza. Samo dopolnitelno ne ustanavjaje, že jest dana direktna suma, komplement ili posebna univerzalna konstrukcija. Točna uloga predstavjenja mora ostati vezana s toj pasažoju.\par
\noindent\emph{Reason:} Dopolnitelno predstavjenje denotes an additional representation in the text's closing qualification. The rationale explicitly rejects substituting a different algebraic construction without textual evidence. Dopolnitelno alone does not establish a direct sum, a complement, or a particular universal construction. The representation's precise role must remain tied to the passage.\par
\noindent{\footnotesize Izvorne zapisy / Original records: HIST-g073-015.}\par
\end{minipage}\par
\par\medskip\Needspace{7\baselineskip}
\noindent\begin{minipage}{\linewidth}
\hypertarget{TERM-0506}{}
\noindent\textbf{izčrpane}\quad{\footnotesize\texttt{TERM-0506}}\par
\noindent\emph{Iztočnik / Source:} erschöpft\par
\noindent\emph{Rusky / Russian:} исчерпаны\par
\noindent\emph{Ukrajinsky / Ukrainian:} вичерпано\par
\smallskip\noindent\emph{Razlog:} Izčrpane označaje, že popis nerazložimyh form Sistema III jest dovršeny. Tuto jest tvrdženje o polnotě popisa, ne prosto o ustalosti ili prestanku računa. Razlog smatra osnovu prozračnoju, ale izrično ostavjaje konečnu morfologiju k prověrkě. Zato izbor značenja ne možno predstaviti kako završen dokaz soglasovanja v rečeniji.\par
\noindent\emph{Reason:} Izčrpane marks the completeness of the enumeration of irreducible forms in System III. It asserts exhaustion of the list, not merely fatigue or the stopping of a calculation. The rationale considers the stem transparent but explicitly leaves the final morphology for checking. Selecting the meaning must not be represented as a completed justification of agreement in the sentence.\par
\noindent{\footnotesize Izvorne zapisy / Original records: HIST-g017-012.}\par
\end{minipage}\par
\par\medskip\Needspace{7\baselineskip}
\noindent\begin{minipage}{\linewidth}
\hypertarget{TERM-0507}{}
\noindent\textbf{prve integraly}\quad{\footnotesize\texttt{TERM-0507}}\par
\noindent\emph{Iztočnik / Source:} erste Integrale\par
\noindent\emph{Rusky / Russian:} первые интегралы\par
\noindent\emph{Ukrajinsky / Ukrainian:} перші інтеграли\par
\smallskip\noindent\emph{Razlog:} Prve integraly jest termin variacijskogo računa i diferencialnyh uravnenij, v množině. Razlog čuva to imenovanje v poslědku Noetherinogo teorema, vměsto zaměny obščim nazvanjem zakonov ohranjenja. Svez s ohranjenjem ne čini dva vyraženja vsegda vzaimno zaměnimymi. Ruska i ukrajinska para zachovuje i porědkovy priděvnik i množinu.\par
\noindent\emph{Reason:} Prve integraly is the plural term from variational calculus and differential equations. The rationale retains it in a consequence of Noether's theorem instead of replacing it with the general label conservation laws. A connection with conservation does not make the expressions universally interchangeable. The Russian and Ukrainian forms preserve both the ordinal adjective and the plural.\par
\noindent{\footnotesize Izvorne zapisy / Original records: HIST-g075-014.}\par
\end{minipage}\par
\par\medskip\Needspace{7\baselineskip}
\noindent\begin{minipage}{\linewidth}
\hypertarget{TERM-0508}{}
\noindent\textbf{prvy / drugi teorem o izomorfiji}\quad{\footnotesize\texttt{TERM-0508}}\par
\noindent\emph{Iztočnik / Source:} Erster / Zweiter Isomorphiesatz\par
\noindent\emph{Rusky / Russian:} первая / вторая теорема об изоморфизме\par
\noindent\emph{Ukrajinsky / Ukrainian:} перша / друга теорема про ізоморфізм\par
\smallskip\noindent\emph{Razlog:} Prvy / drugi teorem o izomorfiji zachovuje čislovano imenovanje i jedinstvenu formu teorema. Razlog ne bira obšči naslov oblasti teorije izomorfizmov. Numeracija mora sledovati konkretnomu iztočniku, ne avtomatično sovrěmennomu učebniku. Historijna forma o izomorfiji ostavaje razlikovana od o izomorfizmu v TERM-0510; zapis ne dava sravnitelny razlog za odklon medžu těmi formami.\par
\noindent\emph{Reason:} Prvy / drugi teorem o izomorfiji retains numbered naming and the singular theorem form. The rationale does not substitute a general heading for the theory of isomorphisms. Numbering must follow the particular source, not automatically a modern textbook. Historical o izomorfiji remains distinct from o izomorfizmu in TERM-0510; the record does not explain the difference between those forms.\par
\noindent{\footnotesize Izvorne zapisy / Original records: HIST-g111-007.}\par
\end{minipage}\par
\par\medskip\Needspace{7\baselineskip}
\noindent\begin{minipage}{\linewidth}
\hypertarget{TERM-0509}{}
\noindent\textbf{prvogo roda}\quad{\footnotesize\texttt{TERM-0509}}\par
\noindent\emph{Iztočnik / Source:} erster Art\par
\noindent\emph{Rusky / Russian:} первого рода\par
\noindent\emph{Ukrajinsky / Ukrainian:} першого роду\par
\smallskip\noindent\emph{Razlog:} Prvogo roda zachovuje klasifikaciju algebraičnyh razširenij iz uvoda raboty 30. Razlog podpiraje širu slovjansku osnovu rod- prěd užšim novym nazvanjem, ale ne zachovuje jeho točnu formu. Tuto ne jest samostatna definicija prvogo roda. Isti slovny vzor v drugom predmetu mora byti razuměny po vlastnom kontekstu, ne po jednom avtomatičnom prěvodu.\par
\noindent\emph{Reason:} Prvogo roda retains the classification of algebraic extensions introduced in Paper 30. The reason supports the broad Slavic rod- family over a narrower new name, whose exact form is not preserved. This is not an independent definition of first kind. The same verbal pattern in a different subject must be interpreted in its own context, not through one automatic substitution.\par
\noindent{\footnotesize Izvorne zapisy / Original records: HIST-g110-006.}\par
\end{minipage}\par
\par\medskip\Needspace{7\baselineskip}
\noindent\begin{minipage}{\linewidth}
\hypertarget{TERM-0510}{}
\noindent\textbf{prvy teorem o izomorfizmu}\quad{\footnotesize\texttt{TERM-0510}}\par
\noindent\emph{Iztočnik / Source:} erster Isomorphiesatz\par
\noindent\emph{Rusky / Russian:} первая теорема об изоморфизме\par
\noindent\emph{Ukrajinsky / Ukrainian:} перша теорема про ізоморфізм\par
\smallskip\noindent\emph{Razlog:} Prvy teorem o izomorfizmu prodolžuje čislovany styl, pri čem zapis izrično dava prvy, ne drugi, po konkretnomu iztočnomu segmentu. Čislo jest čest identifikacije teorema, ne slobodna stilna proměna. Fraza uporablja o izomorfizmu, dokolě TERM-0508 zachovuje o izomorfiji. Sam zapis prodolženja čislovanja ne rěšaje tu terminologičnu varijaciju.\par
\noindent\emph{Reason:} Prvy teorem o izomorfizmu continues the numbered style, with the record specifying first rather than second from the particular source segment. The ordinal identifies the theorem; it is not a free stylistic change. This phrase uses o izomorfizmu, whereas TERM-0508 preserves o izomorfiji. The recorded continuation of numbering does not resolve that terminological variation.\par
\noindent{\footnotesize Izvorne zapisy / Original records: HIST-g152-006.}\par
\end{minipage}\par
\par\medskip\Needspace{7\baselineskip}
\noindent\begin{minipage}{\linewidth}
\hypertarget{TERM-0511}{}
\noindent\textbf{razširena Gaussova teorema}\quad{\footnotesize\texttt{TERM-0511}}\par
\noindent\emph{Iztočnik / Source:} erweiterter Gaußscher Satz\par
\noindent\emph{Rusky / Russian:} расширенная теорема Гаусса\par
\noindent\emph{Ukrajinsky / Ukrainian:} розширена теорема Ґаусса\par
\smallskip\noindent\emph{Razlog:} Razširena Gaussova teorema zachovuje status imenovanogo tvrđenja kako teoremy. Razlog upozornjava, že Noether iz nje oddělno izvodi dvě pomoćne tvrđenja; zaměniti teoremu lemmoju by skrylo tu strukturu. Razširena označaje razširenje obsega tvrđenja, ne dolžinu teksta. Zapis čuva avtorovo razděljenje matematičnyh ulog.\par
\noindent\emph{Reason:} Razširena Gaussova teorema retains the named assertion's status as a theorem. The reason notes that Noether separately derives two auxiliary propositions from it; renaming the theorem a lemma would obscure that structure. Razširena concerns extension of the assertion's scope, not the length of the text. The record preserves the author's distinction of mathematical roles.\par
\noindent{\footnotesize Izvorne zapisy / Original records: HIST-g060-008.}\par
\end{minipage}\par
\par\medskip\Needspace{7\baselineskip}
\noindent\begin{minipage}{\linewidth}
\hypertarget{TERM-0512}{}
\noindent\textbf{razširjenje / suženje oblasti koeficientov}\quad{\footnotesize\texttt{TERM-0512}}\par
\noindent\emph{Iztočnik / Source:} Erweiterung / Verengung des Koeffizientenbereichs\par
\noindent\emph{Rusky / Russian:} расширение / сужение области коэффициентов\par
\noindent\emph{Ukrajinsky / Ukrainian:} розширення / звуження області коефіцієнтів\par
\smallskip\noindent\emph{Razlog:} Razširjenje i suženje oblasti koeficientov označajut protivne napravljanja proměny koeficientnoj oblasti. Izvorni razlog nazyvaje ih maksimalno prozračnymi medžu vostočnymi, zapadnymi i južnymi slovjanskymi jezykami, ale sam zapis dava samo rusku i ukrajinsku paralelu, bez izměrenja razumělivosti. Zato takvu ocěnu trěba čitati kako razsudžanje o kandidatu, ne dokaz vseh trěh skupin. Točna operacija nad objektami zavisi od pasaže.\par
\noindent\emph{Reason:} Razširjenje and suženje oblasti koeficientov name opposite directions of change in the coefficient domain. The original reason calls them maximally transparent across East, West, and South Slavic languages, but the record itself supplies only Russian and Ukrainian parallels and no intelligibility measurement. That assessment must therefore be read as judgment about the candidate, not evidence from all three groups. The precise operation on the objects depends on the passage.\par
\noindent{\footnotesize Izvorne zapisy / Original records: HIST-g213-007.}\par
\end{minipage}\par
\par\medskip\Needspace{7\baselineskip}
\noindent\begin{minipage}{\linewidth}
\hypertarget{TERM-0513}{}
\noindent\textbf{razširenje Gaussovej teoremy}\quad{\footnotesize\texttt{TERM-0513}}\par
\noindent\emph{Iztočnik / Source:} Erweiterung des Gaußschen Satzes\par
\noindent\emph{Rusky / Russian:} обобщение теоремы Гаусса\par
\noindent\emph{Ukrajinsky / Ukrainian:} узагальнення теореми Гауса\par
\smallskip\noindent\emph{Razlog:} Razširenje Gaussovej teoremy imenuje obobčenje, upotrěbljeno za zaključok o polinomialnyh koeficientah iz dělimosti. Ruska i ukrajinska paralela izrično davajut obobčenje, dokolě medžuslovjanska forma prodolžuje osnovu razšir-. Tute kandidatne frazy opisujut odnos medžu tvrđenjami, ne razširenje osnovnogo polja. Samo německo Erweiterung ne odlučuje medžu tymi matematičnymi operacijami.\par
\noindent\emph{Reason:} Razširenje Gaussovej teoremy names the generalization used to infer polynomial coefficients from divisibility. Russian and Ukrainian explicitly use generalization wording, while the Interslavic form continues the razšir- stem. These candidate phrases describe a relation between assertions, not an extension of the ground field. German Erweiterung alone cannot decide between those mathematical operations.\par
\noindent{\footnotesize Izvorne zapisy / Original records: HIST-g049-013.}\par
\end{minipage}\par
\par\medskip\Needspace{7\baselineskip}
\noindent\begin{minipage}{\linewidth}
\hypertarget{TERM-0514}{}
\noindent\textbf{razširenje Gaussova teorema na polinomy s algebraično-celočislenymi koeficientami}\quad{\footnotesize\texttt{TERM-0514}}\par
\noindent\emph{Iztočnik / Source:} Erweiterung des Gaußschen Satzes auf Polynome mit algebraisch-ganzzahligen Koeffizienten\par
\noindent\emph{Rusky / Russian:} расширение теоремы Гаусса на полиномы с алгебраически-целочисленными коэффициентами\par
\noindent\emph{Ukrajinsky / Ukrainian:} розширення теореми Гаусса на поліноми з алгебраїчно-цілочисельними коефіцієнтами\par
\smallskip\noindent\emph{Razlog:} Razširenje Gaussova teorema na polinomy s algebraično-celočislenymi koeficientami zachovuje cělu hipotezu o koeficientah. Algebraično cěly koeficient ne možno zaměniti proizvoljnym algebraičnym koeficientom; takva skračenost by proměnila obseg tvrđenja. Razlog povezuje formulaciju s oddělom 9. On ne daje oddělno opravdanje za vse pravopisne i rodne varijanty teorem/teorema v širšem glosaru.\par
\noindent\emph{Reason:} Razširenje Gaussova teorema na polinomy s algebraično-celočislenymi koeficientami retains the full coefficient hypothesis. An algebraic-integer coefficient cannot be replaced with an arbitrary algebraic coefficient; that abbreviation would alter the assertion's scope. The reason links the formulation to Section 9. It does not separately justify every spelling and gender variant of teorem/teorema in the larger glossary.\par
\noindent{\footnotesize Izvorne zapisy / Original records: HIST-g054-013.}\par
\end{minipage}\par
\par\medskip\Needspace{7\baselineskip}
\noindent\begin{minipage}{\linewidth}
\hypertarget{TERM-0515}{}
\noindent\textbf{razširenje osnovnogo těla}\quad{\footnotesize\texttt{TERM-0515}}\par
\noindent\emph{Iztočnik / Source:} Erweiterung des Grundkörpers\par
\noindent\emph{Rusky / Russian:} расширение основного поля\par
\noindent\emph{Ukrajinsky / Ukrainian:} розширення основного поля\par
\smallskip\noindent\emph{Razlog:} Razširenje osnovnogo těla zachovuje tělo kako šire proektno ime za skalarni objekt, hoti konkretny osnovny objekt jest komutativno polje. Ruska i ukrajinska paralela uporabljajut polje. Razlog opira ograničenje na komutativnost v kontekstu, ne na samo ime tělo. Zato ne možno iz jednogo kandidata izvesti, že tělo i polje sut vsude sinonimy ili že komutativnost ne jest potrěbna.\par
\noindent\emph{Reason:} Razširenje osnovnogo těla retains tělo as the broader project name for a scalar object, although the particular ground object is a commutative field. Russian and Ukrainian use field. The reason relies on the context to impose commutativity, not on tělo alone. This candidate therefore does not establish that tělo and polje are interchangeable everywhere or that commutativity is dispensable.\par
\noindent{\footnotesize Izvorne zapisy / Original records: HIST-g199-001.}\par
\end{minipage}\par
\par\medskip\Needspace{7\baselineskip}
\noindent\begin{minipage}{\linewidth}
\hypertarget{TERM-0516}{}
\noindent\textbf{razširenje prvogo roda}\quad{\footnotesize\texttt{TERM-0516}}\par
\noindent\emph{Iztočnik / Source:} Erweiterung erster Art\par
\noindent\emph{Rusky / Russian:} расширение первого рода\par
\noindent\emph{Ukrajinsky / Ukrainian:} розширення першого роду\par
\smallskip\noindent\emph{Razlog:} Razširenje prvogo roda doslovno zachovuje historijno nazvanje Steinitzovoj klasifikacije. Dva razloga podpirajut kontinuitet izraza pri sovrěmennoj sintaksi. Oni ne davajut polnu definiciju razširenja prvogo roda ni razlog za razširenje prěd varijantoju razširjenje iz TERM-0517. Historijno nazvanje ostavaje svjazano s definicijeju v konkretnoj rabotě.\par
\noindent\emph{Reason:} Razširenje prvogo roda literally retains the historical name of the Steinitz classification. Two reasons support continuity of wording with modern syntax. They give neither the full definition of an extension of the first kind nor a reason to prefer razširenje to razširjenje in TERM-0517. The historical name remains tied to the definition in the particular paper.\par
\noindent{\footnotesize Izvorne zapisy / Original records: HIST-g095-018, HIST-g157-006.}\par
\end{minipage}\par
\par\medskip\Needspace{7\baselineskip}
\noindent\begin{minipage}{\linewidth}
\hypertarget{TERM-0517}{}
\noindent\textbf{razširjenje prvogo roda}\quad{\footnotesize\texttt{TERM-0517}}\par
\noindent\emph{Iztočnik / Source:} Erweiterung erster Art\par
\noindent\emph{Rusky / Russian:} расширение первого рода\par
\noindent\emph{Ukrajinsky / Ukrainian:} розширення першого роду\par
\smallskip\noindent\emph{Razlog:} Razširjenje prvogo roda jest druga zachovana forma istogo historijnogo nazvanja. Razlog povezuje ju s eksponentoju $(e)$ v teoremu IV i s razlicoju medžu doskonalymi i nedoskonalymi poljami. To jest matematična motivacija konteksta, ne dokumentovana fonetična prednost razširjenje prěd razširenje. Obe varijanty morajut ostati vidime pri rekonstrukciji izbora.\par
\noindent\emph{Reason:} Razširjenje prvogo roda is another preserved form of the historical name. The reason connects it with exponent $(e)$ in Theorem IV and the perfect/imperfect-field distinction. That supplies mathematical context, not a documented phonetic preference for razširjenje over razširenje. Both variants must remain visible in reconstructing the choice.\par
\noindent{\footnotesize Izvorne zapisy / Original records: HIST-g097-018.}\par
\end{minipage}\par
\par\medskip\Needspace{7\baselineskip}
\noindent\begin{minipage}{\linewidth}
\hypertarget{TERM-0518}{}
\noindent\textbf{razširjenje prvogo roda / vtorogo roda}\quad{\footnotesize\texttt{TERM-0518}}\par
\noindent\emph{Iztočnik / Source:} Erweiterung erster Art / zweiter Art\par
\noindent\emph{Rusky / Russian:} расширение первого рода / второго рода\par
\noindent\emph{Ukrajinsky / Ukrainian:} розширення першого роду / другого роду\par
\smallskip\noindent\emph{Razlog:} Razširjenje prvogo roda / vtorogo roda zachovuje iztočnu klasifikaciju vměsto tihoj zaměny terminami separabilny/neseparabilny. Razlog izrično ostavjaje odnos k sovrěmennoj terminologiji za daljše sravnjenje. To ne utvrđuje, že istorične i sovrěmenne kategorije sut v vsakom kontekstu različne; samo zabranjaje predpoložiti ih ravnost bez definicij. Takvo objasnjenje naleži k slovniku ili osobnomu naučnomu komentaru, ne k tihoj proměně avtorovogo teksta.\par
\noindent\emph{Reason:} Razširjenje prvogo roda / vtorogo roda preserves the source classification instead of silently substituting separable/inseparable. The rationale explicitly leaves the relation to modern terminology for further comparison. This does not establish that the historical and modern categories differ in every context; it prevents their equivalence being assumed without the definitions. Such explanation belongs in the dictionary or separate scholarly commentary, not a silent change to the author's text.\par
\noindent{\footnotesize Izvorne zapisy / Original records: HIST-g199-004.}\par
\end{minipage}\par
\par\medskip\Needspace{7\baselineskip}
\noindent\begin{minipage}{\linewidth}
\hypertarget{TERM-0519}{}
\noindent\textbf{razširjenja ... sovpadajut v vsih kolcah kvocientov}\quad{\footnotesize\texttt{TERM-0519}}\par
\noindent\emph{Iztočnik / Source:} Erweiterungen ... in allen Quotientenringen stimmen überein\par
\noindent\emph{Rusky / Russian:} расширения ... совпадают во всех кольцах частных\par
\noindent\emph{Ukrajinsky / Ukrainian:} розширення ... збігаються в усіх кільцях часток\par
\smallskip\noindent\emph{Razlog:} Razširjenja ... sovpadajut v vsih kolcah kvocientov zachovuje kvantifikator v vsih, potrěbny za kriterij prehoda od lokalnyh sravnjenij k obščej ravnosti. Razlog povezuje podnožnu primětku s koncem odděla 8, punkta 3. Německo Quotientenring i anglijsky quotient ring ne označajut avtomatično isto kako prsten ostatkov; odgovarjajuču konstrukciju trěba vzeti iz definicije. Historijno kolco ostavaje odvojeno od aktualnogo prsten v raněje pročitanyh pasažah.\par
\noindent\emph{Reason:} Razširjenja ... sovpadajut v vsih kolcah kvocientov retains in all, the quantifier needed for the local-to-global equality criterion. The reason connects the footnote with the end of Section 8, item 3. German Quotientenring and English quotient ring do not automatically mean a residue-class ring; the relevant construction must come from its definition. Historical kolco remains distinguished from current prsten in the previously examined passages.\par
\noindent{\footnotesize Izvorne zapisy / Original records: HIST-g162-008.}\par
\end{minipage}\par
\par\medskip\Needspace{7\baselineskip}
\noindent\begin{minipage}{\linewidth}
\hypertarget{TERM-0520}{}
\noindent\textbf{ideal razširenja}\quad{\footnotesize\texttt{TERM-0520}}\par
\noindent\emph{Iztočnik / Source:} Erweiterungsideal\par
\noindent\emph{Rusky / Russian:} идеал расширения\par
\noindent\emph{Ukrajinsky / Ukrainian:} ідеал розширення\par
\smallskip\noindent\emph{Razlog:} Ideal razširenja jest parny izraz k suženju ideala. Razlog govori o idealu generovanom iz centralnogo ideala v obhvačajučem prstenu, ne prosto o většem množestvě. Dva napravljanja trěbujut razne operacije, oprěděljene v konkretnoj definiciji. Prozračno slovotvorjenje samo ne dokazuje, že razširenje i suženje sut vzaimno obratne vo vseh slučajah.\par
\noindent\emph{Reason:} Ideal razširenja is paired with contraction of an ideal. The reason concerns the ideal generated from a central ideal in the ambient ring, not merely a larger set. The two directions require distinct operations supplied by the particular definition. Transparent word formation alone does not establish that extension and contraction are mutually inverse in all cases.\par
\noindent{\footnotesize Izvorne zapisy / Original records: HIST-g189-004.}\par
\end{minipage}\par
\par\medskip\Needspace{7\baselineskip}
\noindent\begin{minipage}{\linewidth}
\hypertarget{TERM-0521}{}
\noindent\textbf{razširjeny ideal}\quad{\footnotesize\texttt{TERM-0521}}\par
\noindent\emph{Iztočnik / Source:} Erweiterungsideal\par
\noindent\emph{Rusky / Russian:} идеал расширения\par
\noindent\emph{Ukrajinsky / Ukrainian:} ідеал розширення\par
\smallskip\noindent\emph{Razlog:} Razširjeny ideal prěnosi isti korenj kako razširjeno kolco i sohranjaje svez s suženjem ideala. Zapisy iz oddělov 1, 2 i 5 raboty 31 davajut prednost toj stalnoj parě prěd opisami prosto generovany ili uveličeny ideal. V osobnom kontekstu komponenty ideal možno utožsamiti s prstenom dobytym razširenjem; to ne znači, že ideal i prsten sut vsegda zaměnljive pojmy. Razlog zato trěbuje kontekstnu definiciju, ne samo doslovno razloženje německogo kompozita.\par
\noindent\emph{Reason:} Razširjeny ideal shares its root with razširjeno kolco and retains the connection with contraction of an ideal. Records from sections 1, 2 and 5 of Paper 31 favour this stable pair over merely describing the ideal as generated or enlarged. In the particular component construction, the ideal may be identified with the ring obtained by extension; this does not make ideal and ring interchangeable everywhere. The rationale therefore depends on the contextual definition, not just a literal division of the German compound.\par
\noindent{\footnotesize Izvorne zapisy / Original records: HIST-g120-001, HIST-g123-003, HIST-g140-003.}\par
\end{minipage}\par
\par\medskip\Needspace{7\baselineskip}
\noindent\begin{minipage}{\linewidth}
\hypertarget{TERM-0522}{}
\noindent\textbf{razširjeny ideal libovoljnogo ideala c}\quad{\footnotesize\texttt{TERM-0522}}\par
\noindent\emph{Iztočnik / Source:} Erweiterungsideal eines beliebigen Ideals c\par
\noindent\emph{Rusky / Russian:} идеал расширения произвольного идеала c\par
\noindent\emph{Ukrajinsky / Ukrainian:} ідеал розширення довільного ідеалу c\par
\smallskip\noindent\emph{Razlog:} Libovoljnogo ideala c govori, že razširenje priměnjaje se k proizvoljnomu idealu osnovnogo prstena množenja, ne samo k raněje izbranoj komponentě. Glavni termin razširjeny ideal ostavaje isti. Oznaka c identifikuje izhodny ideal; libovoljnogo sohranjaje obščnost tvrdženja. Zapis ne uvodi novu operaciju ni posebno ograničenje na prosty ili primarny ideal.\par
\noindent\emph{Reason:} Libovoljnogo ideala c makes the extension apply to an arbitrary ideal of the base multiplication ring, not only to an earlier selected component. The principal term razširjeny ideal stays unchanged. The symbol c identifies the starting ideal, while libovoljnogo preserves the statement's generality. The record introduces neither a new operation nor a restriction to prime or primary ideals.\par
\noindent{\footnotesize Izvorne zapisy / Original records: HIST-g160-007.}\par
\end{minipage}\par
\par\medskip\Needspace{7\baselineskip}
\noindent\begin{minipage}{\linewidth}
\hypertarget{TERM-0523}{}
\noindent\textbf{razširjene idealy}\quad{\footnotesize\texttt{TERM-0523}}\par
\noindent\emph{Iztočnik / Source:} Erweiterungsideale\par
\noindent\emph{Rusky / Russian:} идеалы расширения\par
\noindent\emph{Ukrajinsky / Ukrainian:} ідеали розширення\par
\smallskip\noindent\emph{Razlog:} Množinska forma razširjene idealy sohranjaje termin iz jednine pri prěhodu k lokalizaciji v oddělu 8 raboty 31. Zapis argumentuje kontinuitet, ne novy slovotvorny izbor. Da by takovy kontinuitet byl matematično opravdany, ideal mora byti razširjen do označenogo prstena; sam fakt, že se pojavi v většem prstenu, ne zastupaje definiciju razširenja.\par
\noindent\emph{Reason:} The plural razširjene idealy carries the established singular term into the localization discussion in section 8 of Paper 31. The recorded argument is continuity, not a new word-formation decision. That continuity is mathematically justified by extension to the specified ring; merely occurring in a larger ring does not replace the definition of extension.\par
\noindent{\footnotesize Izvorne zapisy / Original records: HIST-g159-007.}\par
\end{minipage}\par
\par\medskip\Needspace{7\baselineskip}
\noindent\begin{minipage}{\linewidth}
\hypertarget{TERM-0524}{}
\noindent\textbf{razširitelno polje}\quad{\footnotesize\texttt{TERM-0524}}\par
\noindent\emph{Iztočnik / Source:} Erweiterungskörper\par
\noindent\emph{Rusky / Russian:} поле расширения\par
\noindent\emph{Ukrajinsky / Ukrainian:} поле розширення\par
\smallskip\noindent\emph{Razlog:} Rusko pole razširenija i ukrajinsko pole rozšyrennja sut imenne konstrukcije, a zapis izbira pridavnik s imenem: razširitelno polje. Razlog jest ranějša upotrěba toj formy v prěvodu. To jest dokumentovany slovotvorny izbor, ne dokaz, že vsaka druga pridavnička forma jest nepravilna. Drugy zapis, TERM-0525, ima razširjeno polje; tute dvě historijne formy ne trěba skryti pod jednoj navodno jedinstvennoj rekomendacijej.\par
\noindent\emph{Reason:} The Russian and Ukrainian comparanda use a noun phrase, field of extension; the recorded choice uses adjective plus noun, razširitelno polje. Its reason is earlier use in the translation. This documents a word-formation choice, not the incorrectness of every other adjective. TERM-0525 records razširjeno polje instead; the two historical forms must not be concealed beneath a supposedly uniform recommendation.\par
\noindent{\footnotesize Izvorne zapisy / Original records: HIST-g083-011.}\par
\end{minipage}\par
\par\medskip\Needspace{7\baselineskip}
\noindent\begin{minipage}{\linewidth}
\hypertarget{TERM-0525}{}
\noindent\textbf{razširjeno polje / razširjeno kolco}\quad{\footnotesize\texttt{TERM-0525}}\par
\noindent\emph{Iztočnik / Source:} Erweiterungskörper / Erweiterungsring\par
\noindent\emph{Rusky / Russian:} поле расширения / кольцо расширения\par
\noindent\emph{Ukrajinsky / Ukrainian:} поле розширення / кільце розширення\par
\smallskip\noindent\emph{Razlog:} Para razširjeno polje / razširjeno kolco sohranjaje razlikovanje medžu poljem i prstenom, pri tom kako isti pridavnik vyražaje razširenje. Razlog odkazuje k prěhodu v točkě 2b od polj Q k prstenam R[Q]. Ty dva nazvy zato ne sut slobodne sinonimy. Oznaka polje ne može byti avtomatično prěnesena na vsaky prsten, dobyty prisojedinjenjem.\par
\noindent\emph{Reason:} The pair razširjeno polje / razširjeno kolco keeps field and ring distinct while using the same adjective for extension. The reason points to the passage in 2b from fields Q to rings R[Q]. The names are therefore not freely interchangeable. Field cannot automatically replace ring for every adjoining construction.\par
\noindent{\footnotesize Izvorne zapisy / Original records: HIST-g127-003.}\par
\end{minipage}\par
\par\medskip\Needspace{7\baselineskip}
\noindent\begin{minipage}{\linewidth}
\hypertarget{TERM-0526}{}
\noindent\textbf{razširjeno polje prvogo roda podpolja P}\quad{\footnotesize\texttt{TERM-0526}}\par
\noindent\emph{Iztočnik / Source:} Erweiterungskörper erster Art des Unterkörpers P\par
\noindent\emph{Rusky / Russian:} поле расширения первого рода подполя P\par
\noindent\emph{Ukrajinsky / Ukrainian:} поле розширення першого роду підполя P\par
\smallskip\noindent\emph{Razlog:} Prvogo roda jest klasifikacija razširenja, a podpolja P identifikuje polje, nad ktorym ono razsmatriva se. Zapis naměrno odlišaje tuto razširjeno polje od obščego razširenja Q. Zachovany razlog podkrepljaje kontinuitet polje, podpolje i razširjeno polje, ale ne dava samostatnu definiciju prvogo roda. Tu definiciju trěba čitati v rabotě, ne zaměniti domyslenym sovrěmennym ekvivalentom.\par
\noindent\emph{Reason:} Prvogo roda classifies the extension, and podpolja P identifies the field over which it is considered. The record deliberately distinguishes this extension from the general extension Q. It supports continuity among field, subfield and extension field but does not itself define first kind. That definition must be read in the work, not replaced by an assumed modern equivalent.\par
\noindent{\footnotesize Izvorne zapisy / Original records: HIST-g147-004.}\par
\end{minipage}\par
\par\medskip\Needspace{7\baselineskip}
\noindent\begin{minipage}{\linewidth}
\hypertarget{TERM-0527}{}
\noindent\textbf{razširjeno polje K polja kvocientov kolca $o_p$}\quad{\footnotesize\texttt{TERM-0527}}\par
\noindent\emph{Iztočnik / Source:} Erweiterungskörper K des Quotientenkörpers von $o_p$\par
\noindent\emph{Rusky / Russian:} поле расширения K поля частных кольца $o_p$\par
\noindent\emph{Ukrajinsky / Ukrainian:} поле розширення K поля часток кільця $o_p$\par
\smallskip\noindent\emph{Razlog:} Dobavjenje kolca čini vidimu vnutrenju strukturu frazy: najprvo jest prsten $o_p$, potom njegovo polje kvocientov, a potom razširjeno polje K. Razlog govori, že takovo objasnjenje bylo vključeno vo vseh trěh prěvodah. Ono ne dobavjaje novu algebraičnu hipotezu, ale pokazuje, k čemu prinadleži polje kvocientov. Tre nivoje konstrukcije ne trěba svesti k prostoj frazě razširenje prstena.\par
\noindent\emph{Reason:} Adding kolca makes the phrase's internal structure visible: first the ring $o_p$, then its quotient field, then the extension field K. The reason says that all three translations made this explicit. The addition does not supply a new algebraic hypothesis; it identifies whose quotient field is meant. The three levels of construction must not be collapsed into the simpler phrase extension of the ring.\par
\noindent{\footnotesize Izvorne zapisy / Original records: HIST-g161-006.}\par
\smallskip\noindent\emph{Sravnjenje s tekstom:} V točkě 4 odděla 8 raboty 31 aktualna pasaža ima v razširjenom polju K polja kvocientov prstena i potom oznaku osnovnogo prstena. Tuto potvrđaje objasnjenje vnutrenjej struktury frazy v historijnom razlogu. Aktualno prstena razlikuje se od ranějšego kolca. Sravnjenje ne sliva te formy i ne prenosi na tute slova kakvu-libo neprovedenu čitateljsku ocěnu.\par
\noindent\emph{Textual comparison:} In section 8, item 4 of Paper 31, the current passage reads v razširjenom polju K polja kvocientov prstena, followed by the base-ring notation. This supports the historical explanation of the phrase's nested structure. Current prstena differs from earlier kolca. The comparison keeps those forms distinct and makes no claim of an unperformed reader assessment.\par
\noindent{\footnotesize TEXT-COMP-18; pasaže / passages: 19003; \texttt{CORPUS\_FORM\_COMPARISONS.json}.}\par
\end{minipage}\par
\par\medskip\Needspace{7\baselineskip}
\noindent\begin{minipage}{\linewidth}
\hypertarget{TERM-0528}{}
\noindent\textbf{razširjeno kolco}\quad{\footnotesize\texttt{TERM-0528}}\par
\noindent\emph{Iztočnik / Source:} Erweiterungsring\par
\noindent\emph{Rusky / Russian:} кольцо расширения\par
\noindent\emph{Ukrajinsky / Ukrainian:} кільце розширення\par
\smallskip\noindent\emph{Razlog:} Četyri zapisy iz raboty 31 podderžujut razširjeno kolco kako parny izraz k razširjenomu polju. Oni uključajut prisojedinjenje elementa, razširenje koeficientnogo polja do algebraično zamknjenogo polja i prsten nad koeficientnoju oblastju Z. Pridavnička konstrukcija imaje sohraniti kratku, skladnu terminologiju. Jeden razlog odbacaje tęžky doslovny prěvod, ale ne navodi njegovu formu; slovnik zato ne može izmysliti konkretnogo konkurenta ili čislenu ocěnu za njega.\par
\noindent\emph{Reason:} Four records from Paper 31 support razširjeno kolco as the counterpart of extension field. Their contexts include adjoining an element, extending the coefficient field to an algebraically closed field, and a ring over the coefficient domain Z. The adjective construction is intended to keep the terminology compact and consistent. One reason rejects a cumbersome calque without naming it; the dictionary cannot invent a specific competitor or a numerical score for it.\par
\noindent{\footnotesize Izvorne zapisy / Original records: HIST-g119-001, HIST-g123-002, HIST-g126-003, HIST-g130-005.}\par
\end{minipage}\par
\par\medskip\Needspace{7\baselineskip}
\noindent\begin{minipage}{\linewidth}
\hypertarget{TERM-0529}{}
\noindent\textbf{razširjeno kolco $\overline{\mathfrak R}=\mathfrak R[\overline P]$}\quad{\footnotesize\texttt{TERM-0529}}\par
\noindent\emph{Iztočnik / Source:} Erweiterungsring $\overline{\mathfrak R}=\mathfrak R[\overline P]$\par
\noindent\emph{Rusky / Russian:} кольцо расширения $\overline{\mathfrak R}=\mathfrak R[\overline P]$\par
\noindent\emph{Ukrajinsky / Ukrainian:} кільце розширення $\overline{\mathfrak R}=\mathfrak R[\overline P]$\par
\smallskip\noindent\emph{Razlog:} Formula $\overline{\mathfrak R}=\mathfrak R[\overline P]$ jest sučstvena čest nazva: ona pokazuje, kako nastava označeny razširjeny prsten. Razlog dava prednost pridavniku razširjeno prěd abstraktnym imenem, da fraza ostane čitateljna v dokazě s razširenjem skalarov. Čitatelj mora moči razlikovati izhodny prsten, razširjene skalary i rezultat; odstranjenje formuly kako navodno lišnjego detalja by izgubilo tuto informaciju.\par
\noindent\emph{Reason:} The formula $\overline{\mathfrak R}=\mathfrak R[\overline P]$ is integral to the designation: it identifies how the extension ring is obtained. The reason favours adjectival razširjeno over an abstract noun for readability in a proof involving scalar extension. Readers need to distinguish the starting ring, the extended scalars and the result; dropping the formula as incidental detail would remove that information.\par
\noindent{\footnotesize Izvorne zapisy / Original records: HIST-g145-002.}\par
\end{minipage}\par
\par\medskip\Needspace{7\baselineskip}
\noindent\begin{minipage}{\linewidth}
\hypertarget{TERM-0530}{}
\noindent\textbf{razširjeno kolco polja / razširjeno polje}\quad{\footnotesize\texttt{TERM-0530}}\par
\noindent\emph{Iztočnik / Source:} Erweiterungsring eines Körpers / Erweiterungskörper\par
\noindent\emph{Rusky / Russian:} кольцо расширения поля / поле расширения\par
\noindent\emph{Ukrajinsky / Ukrainian:} кільце розширення поля / поле розширення\par
\smallskip\noindent\emph{Razlog:} Fraza razširjeno kolco polja imenuje prsten nad poljem, a razširjeno polje imenuje polje. Takovo razlikovanje jest važno v početku odděla 6 raboty 31, kde obe konstrukcije nastupajut zajedno. Podobnost pridavnika ne dovoljuje izpustiti osnovno ime: inače by prěvod pripisoval poljeve svojstva objektu, ktory v toj frazě jest samo prstenom.\par
\noindent\emph{Reason:} Razširjeno kolco polja names a ring over a field, whereas razširjeno polje names a field. This distinction matters at the opening of section 6 of Paper 31, where both constructions occur together. Sharing the adjective does not permit omission of the head noun: otherwise the translation would attribute field properties to an object designated there only as a ring.\par
\noindent{\footnotesize Izvorne zapisy / Original records: HIST-g148-003.}\par
\end{minipage}\par
\par\medskip\Needspace{7\baselineskip}
\noindent\begin{minipage}{\linewidth}
\hypertarget{TERM-0531}{}
\noindent\textbf{razširjeno kolco od Z bez nulovyh děliteljev}\quad{\footnotesize\texttt{TERM-0531}}\par
\noindent\emph{Iztočnik / Source:} Erweiterungsring ohne Nullteiler von Z\par
\noindent\emph{Rusky / Russian:} кольцо расширения кольца Z, не имеющее делителей нуля\par
\noindent\emph{Ukrajinsky / Ukrainian:} кільце розширення кільця Z, яке не має дільників нуля\par
\smallskip\noindent\emph{Razlog:} Bez nulovyh děliteljev ostavaje izrečeno uslovje razširjenogo prstena od Z. Zapis dava prednost toj polnoj konstrukciji prěd prěhodom k drugomu abstraktnomu nazvu v srědině odděla. Tym se ne sugeruje, že vsako razširenje sohranjaje otsutnost nulovyh děliteljev; fraza izrično pripisuje to svojstvo razsmatrivanomu objektu. Z mora ostati točno identifikovano.\par
\noindent\emph{Reason:} Bez nulovyh děliteljev retains the explicit no-zero-divisors condition on the extension ring of Z. The record prefers this full construction to switching to another abstract name mid-section. It does not suggest that every extension preserves the absence of zero divisors; the phrase expressly attributes that property to the object under discussion. Z must remain identifiable.\par
\noindent{\footnotesize Izvorne zapisy / Original records: HIST-g140-002.}\par
\end{minipage}\par
\par\medskip\Needspace{7\baselineskip}
\noindent\begin{minipage}{\linewidth}
\hypertarget{TERM-0532}{}
\noindent\textbf{kolco raz?irenja; modul raz?irenja}\quad{\footnotesize\texttt{TERM-0532}}\par
\noindent\emph{Iztočnik / Source:} Erweiterungsring; Erweiterungsmodul\par
\noindent\emph{Rusky / Russian:} [Bukvy ne sut zachovane / Lettering not recoverable]\par
\noindent\emph{Ukrajinsky / Ukrainian:} [Bukvy ne sut zachovane / Lettering not recoverable]\par
\smallskip\noindent\emph{Razlog:} Zapis para kolco raz?irenja s modul raz?irenja, čto pokazuje zajednu konstrukciju s imenem razširenja. Ale jedna bukva v medžuslovjanskoj formě i bukvy v ruskih i ukrajinskih sravnjenjah ne sut zachovane. Razlog opiraje se na kontinuitet s ranějšimi rabotami i tvrđenu medžuslovjansku prozračnost; ne predklada čitateljske testy. Oštečena forma zato ne jest gotova pravopisna rekomendacija, ni dokaz soglasnosti konkretnyh nacionalnyh form.\par
\noindent\emph{Reason:} The record pairs kolco raz?irenja with modul raz?irenja, showing a shared noun construction for extension. One letter in the Interslavic form and the lettering of the Russian and Ukrainian comparanda are not preserved. The rationale invokes continuity with earlier papers and claimed cross-Slavic transparency, but supplies no reader tests. The damaged form is therefore neither a finished spelling recommendation nor evidence of agreement among identifiable national forms.\par
\noindent{\footnotesize Izvorne zapisy / Original records: HIST-g210-004.}\par
\end{minipage}\par
\par\medskip\Needspace{7\baselineskip}
\noindent\begin{minipage}{\linewidth}
\hypertarget{TERM-0533}{}
\noindent\textbf{razširjene kolca polj}\quad{\footnotesize\texttt{TERM-0533}}\par
\noindent\emph{Iztočnik / Source:} Erweiterungsringe von Körpern\par
\noindent\emph{Rusky / Russian:} кольца расширения полей\par
\noindent\emph{Ukrajinsky / Ukrainian:} кільця розширення полів\par
\smallskip\noindent\emph{Razlog:} Množinsko razširjene kolca polj jest zaglavje o prstenah, ktore nastavajut prisojedinjenjem nad koeficientnymi poljami. Razlog izrično ohranjaje kolca: govoriti samo o razširjenyh poljah by suzilo sodržanje zaglavja. Polj jest zavisny član frazy, ne druga oznaka vseh dobytyh prstenov.\par
\noindent\emph{Reason:} The plural heading razširjene kolca polj concerns rings formed by adjoining over coefficient fields. The rationale expressly preserves rings: referring only to extension fields would narrow the heading's content. Polj is the dependent member of the phrase, not another designation for all the resulting rings.\par
\noindent{\footnotesize Izvorne zapisy / Original records: HIST-g119-002.}\par
\end{minipage}\par
\par\medskip\Needspace{7\baselineskip}
\noindent\begin{minipage}{\linewidth}
\hypertarget{TERM-0535}{}
\noindent\textbf{možno prisojedinjenje neopreděljenyh}\quad{\footnotesize\texttt{TERM-0535}}\par
\noindent\emph{Iztočnik / Source:} eventuelle Adjunktion von Unbestimmten\par
\noindent\emph{Rusky / Russian:} возможное присоединение неопределенных\par
\noindent\emph{Ukrajinsky / Ukrainian:} можливе приєднання невизначених\par
\smallskip\noindent\emph{Razlog:} Možno prisojedinjenje sohranjaje uslovnost německogo eventuelle: neoprěděljene mogu byti prisojedinjene, ale to ne jest bezuslovny povel. Razlog svęzuje tu možnost s pozdnějšej specializacijej, kogda P ima beskonečno mnogo elementov. Izpustiti možno by značilo proměniti dopuščenu konstrukciju v objazatelny korak; izpustiti kontekst specializacije by skrylo, začto konstrukcija jest upotrěbna.\par
\noindent\emph{Reason:} Možno prisojedinjenje preserves the conditional force of eventuelle: indeterminates may be adjoined, but this is not an unconditional instruction. The reason connects that possibility with later specialization when P has infinitely many elements. Omitting možno would turn an available construction into a mandatory step; omitting the specialization context would hide why the construction is useful.\par
\noindent{\footnotesize Izvorne zapisy / Original records: HIST-g098-008.}\par
\end{minipage}\par
\par\medskip\Needspace{7\baselineskip}
\noindent\begin{minipage}{\linewidth}
\hypertarget{TERM-0536}{}
\noindent\textbf{eksplicitno predstavjenje}\quad{\footnotesize\texttt{TERM-0536}}\par
\noindent\emph{Iztočnik / Source:} explizite Darstellung\par
\noindent\emph{Rusky / Russian:} явное представление\par
\noindent\emph{Ukrajinsky / Ukrainian:} явне подання\par
\smallskip\noindent\emph{Razlog:} Eksplicitno predstavjenje jest tehničny izraz, ne pohvala jasnosti izloženja. V zapisah k rabotě 4 ono znači predstavjenje posrědstvom samyh redov ili redov iz njih izvedenyh. Odděl 4 odkladaje pytanje, a odděl 5 dokazuje takovo predstavjenje identičnostij razloženja. Isti termin mora povezati definiciju, najavu i dokaz; slobodno zaměniti jego s podrobno objasnjenje by izgubilo točny smisyl.\par
\noindent\emph{Reason:} Eksplicitno predstavjenje is a technical expression, not praise for clear exposition. In the records for Paper 4 it denotes representation using the rows themselves or rows derived from them. Section 4 postpones the question, and section 5 proves such a representation of the decomposition identities. The same term links definition, announcement and proof; freely replacing it with detailed explanation would lose that specific meaning.\par
\noindent{\footnotesize Izvorne zapisy / Original records: HIST-g031-004, HIST-g033-009, HIST-g034-004.}\par
\end{minipage}\par
\par\medskip\Needspace{7\baselineskip}
\noindent\begin{minipage}{\linewidth}
\hypertarget{TERM-0537}{}
\noindent\textbf{eksplicitno predstavjenje invariantnyh tvorjenj}\quad{\footnotesize\texttt{TERM-0537}}\par
\noindent\emph{Iztočnik / Source:} explizite Darstellung der invarianten Bildungen\par
\noindent\emph{Rusky / Russian:} явное представление инвариантных образований\par
\noindent\emph{Ukrajinsky / Ukrainian:} явне подання інваріантних утворень\par
\smallskip\noindent\emph{Razlog:} Eksplicitno predstavjenje invariantnyh tvorjenj prodolžaje termin iz oddělov 2 i 5 raboty 4 v oddělu 6. Rusko predstavlenie i ukrajinsko podannja dajut sravniteljne nazvy predstavjenja, a medžuslovjansky zapis ostavaje pri predstavjenje. Invariantnyh tvorjenj oprěděljaje objekty, ne osobny stil objasnjenja. Razlog kontinuiteta ne dava osnovy zaměniti tvorjenja s libovoljnymi invariantami bez prověrky konteksta.\par
\noindent\emph{Reason:} Eksplicitno predstavjenje invariantnyh tvorjenj carries the term from sections 2 and 5 of Paper 4 into section 6. Russian predstavlenie and Ukrainian podannja provide comparison expressions for representation; the Interslavic record retains predstavjenje. Invariantnyh tvorjenj specifies the objects, not an expository style. Terminological continuity does not license replacing formations with arbitrary invariants without checking the context.\par
\noindent{\footnotesize Izvorne zapisy / Original records: HIST-g035-001.}\par
\end{minipage}\par
\par\medskip\Needspace{7\baselineskip}
\noindent\begin{minipage}{\linewidth}
\hypertarget{TERM-0538}{}
\noindent\textbf{eksplicitno razloženje}\quad{\footnotesize\texttt{TERM-0538}}\par
\noindent\emph{Iztočnik / Source:} explizite Zerlegung\par
\noindent\emph{Rusky / Russian:} явное разложение\par
\noindent\emph{Ukrajinsky / Ukrainian:} явний розклад\par
\smallskip\noindent\emph{Razlog:} Eksplicitno razloženje naznačaje, že produktna formula davaje samo razloženje, ne samo dokaz, že ono obstaja. Razlog k oddělu 6 raboty 22 preto sohranjaje oba slova. Eksplicitnost odnosi se k formě predstavjenja; ne garantuje sama po sobě jedinstvennost faktorov, kanoničnost ili legkost izčislenja.\par
\noindent\emph{Reason:} Eksplicitno razloženje indicates that the product formula supplies the decomposition itself, not merely a proof that one exists. The reason for section 6 of Paper 22 therefore retains both words. Explicitness concerns the representation supplied; by itself it guarantees neither uniqueness of factors, canonicity nor ease of computation.\par
\noindent{\footnotesize Izvorne zapisy / Original records: HIST-g090-014, HIST-g092-068.}\par
\end{minipage}\par
\par\medskip\Needspace{7\baselineskip}
\noindent\begin{minipage}{\linewidth}
\hypertarget{TERM-0539}{}
\noindent\textbf{eksponent}\quad{\footnotesize\texttt{TERM-0539}}\par
\noindent\emph{Iztočnik / Source:} Exponent\par
\noindent\emph{Rusky / Russian:} экспонента\par
\noindent\emph{Ukrajinsky / Ukrainian:} експонента\par
\smallskip\noindent\emph{Razlog:} Eksponent sohranjaje internacionalny nazv pri definiciji, ktora ispytuje najnižšu odgovarjajuču stepenj p ili P. Kratki historijni razlog govori o najmenšej stepeni, ale eksponent jest čislom, ktoro tu stepenj oprěděljaje, ne samym produktom ili idealom. Tuto razlikovanje trěba izrekti, da slovnik ne ponovi nepreciznost skračenogo razloga. Zapisna forma eksponent ne smi takože skryti drugu historijnu formu eksponenta.\par
\noindent\emph{Reason:} Eksponent retains the international name in a definition that considers the least suitable power of p or P. The brief historical reason speaks of the least power, but the exponent is the number specifying that power, not the product or ideal itself. This distinction needs to be stated so that the dictionary does not repeat the compressed rationale's imprecision. The recorded form eksponent must also remain distinguishable from historical eksponenta.\par
\noindent{\footnotesize Izvorne zapisy / Original records: HIST-g081-011, HIST-g082-011, HIST-g095-020.}\par
\smallskip\noindent\emph{Sravnjenje s tekstom:} V rabotě 19, teoremě V, aktualny tekst naziva najmenše čislo $\varrho$ s uslovjem $P^{\varrho}\equiv0\;(Q)$ eksponentom od $Q$. Podnožna primětka izrično odlišaje tuto uslovje od ravnosti $P^{\varrho}=Q$ i davaje priměr $Q=(x^2,y)$, $P=(x,y)$. Tuto odstraňaje dvosmyslnost kratkogo historijnogo razloga o najmenšej stepeni: vybrano jest čislo, a ne idealna stepenj. Sravnjenje jest novym čitanjem toj konkretnoj pasaže, ne dokazom ranějšej konsultacije i ne prověrkoju vseh smyslov eksponent v korpusu.\par
\noindent\emph{Textual comparison:} In Paper 19, Theorem V, the current text calls the least number $\varrho$ satisfying $P^{\varrho}\equiv0\;(Q)$ the exponent of $Q$. Its footnote expressly distinguishes this condition from $P^{\varrho}=Q$ and gives $Q=(x^2,y)$, $P=(x,y)$ as an example. This resolves the compressed historical rationale's ambiguity about the least power: the selected object is a number, not the ideal power. The comparison is a fresh reading of this particular passage, not evidence of earlier consultation or verification of every sense of exponent in the corpus.\par
\noindent{\footnotesize TEXT-COMP-15; pasaže / passages: 12834, 12834; \texttt{CORPUS\_FORM\_COMPARISONS.json}.}\par
\end{minipage}\par
\par\medskip\Needspace{7\baselineskip}
\noindent\begin{minipage}{\linewidth}
\hypertarget{TERM-0540}{}
\noindent\textbf{eksponent = indeks}\quad{\footnotesize\texttt{TERM-0540}}\par
\noindent\emph{Iztočnik / Source:} Exponent = Index\par
\noindent\emph{Rusky / Russian:} экспонента = индекс\par
\noindent\emph{Ukrajinsky / Ukrainian:} експонента = індекс\par
\smallskip\noindent\emph{Razlog:} Eksponent = indeks vyražaje ravnost dvoh invariantov v teoremě raboty 38 nad čislovymi poljami. Zapis sohranjaje znak ravnosti i oba imena, vměsto da by njih slil v jeden termin. Ravnost v toj teoremě ne jest slovnična sinonimija ni universalna definicija eksponenta črez indeks; uslovja tvrdženja ostavajut neotdělimi od jego upotrěby.\par
\noindent\emph{Reason:} Eksponent = indeks expresses the equality of two invariants in the theorem of Paper 38 over number fields. The record retains the equality sign and both names rather than collapsing them into one term. Equality in that theorem is neither lexical synonymy nor a universal definition of exponent through index; the theorem's hypotheses remain inseparable from its use.\par
\noindent{\footnotesize Izvorne zapisy / Original records: HIST-g208-010.}\par
\end{minipage}\par
\par\medskip\Needspace{7\baselineskip}
\noindent\begin{minipage}{\linewidth}
\hypertarget{TERM-0541}{}
\noindent\textbf{eksponent primarnogo ideala}\quad{\footnotesize\texttt{TERM-0541}}\par
\noindent\emph{Iztočnik / Source:} Exponent eines primären Ideals\par
\noindent\emph{Rusky / Russian:} экспонента примарного идеала\par
\noindent\emph{Ukrajinsky / Ukrainian:} експонента примарного ідеалу\par
\smallskip\noindent\emph{Razlog:} Primarnogo ideala privęzuje eksponent k oprěděljenomu idealu. Ako $P$ jest pridruženy prosty ideal k $Q$, išče se najmenše čislo $\varrho$ s uslovjem $P^{\varrho}\subseteq Q$; eksponent jest $\varrho$, ne $P^{\varrho}$. Zachovany razlog skračaje tuto razlikovanje slovami najmenša stepenj. Polno objasnjenje mora razlikovanje obnoviti, a ne prěnesti skračenost v definiciju. Iz uslovja takože ne slěduje ravnost $P^{\varrho}=Q$.\par
\noindent\emph{Reason:} Primarnogo ideala attaches the exponent to a specified primary ideal. If $P$ is the associated prime of $Q$, one seeks the least number $\varrho$ satisfying $P^{\varrho}\subseteq Q$; the exponent is $\varrho$, not $P^{\varrho}$. The surviving reason compresses this distinction into the words least power. A full explanation must restore it rather than transfer the compression into the definition. The condition also does not imply $P^{\varrho}=Q$.\par
\noindent{\footnotesize Izvorne zapisy / Original records: HIST-g081-043, HIST-g082-043.}\par
\smallskip\noindent\emph{Sravnjenje s tekstom:} V rabotě 19, teoremě V, aktualny tekst naziva najmenše čislo $\varrho$ s uslovjem $P^{\varrho}\equiv0\;(Q)$ eksponentom od $Q$. Podnožna primětka izrično odlišaje tuto uslovje od ravnosti $P^{\varrho}=Q$ i davaje priměr $Q=(x^2,y)$, $P=(x,y)$. Tuto odstraňaje dvosmyslnost kratkogo historijnogo razloga o najmenšej stepeni: vybrano jest čislo, a ne idealna stepenj. Sravnjenje jest novym čitanjem toj konkretnoj pasaže, ne dokazom ranějšej konsultacije i ne prověrkoju vseh smyslov eksponent v korpusu.\par
\noindent\emph{Textual comparison:} In Paper 19, Theorem V, the current text calls the least number $\varrho$ satisfying $P^{\varrho}\equiv0\;(Q)$ the exponent of $Q$. Its footnote expressly distinguishes this condition from $P^{\varrho}=Q$ and gives $Q=(x^2,y)$, $P=(x,y)$ as an example. This resolves the compressed historical rationale's ambiguity about the least power: the selected object is a number, not the ideal power. The comparison is a fresh reading of this particular passage, not evidence of earlier consultation or verification of every sense of exponent in the corpus.\par
\noindent{\footnotesize TEXT-COMP-15; pasaže / passages: 12834, 12834; \texttt{CORPUS\_FORM\_COMPARISONS.json}.}\par
\end{minipage}\par
\par\medskip\Needspace{7\baselineskip}
\noindent\begin{minipage}{\linewidth}
\hypertarget{TERM-0542}{}
\noindent\textbf{eksponent primarnogo ideala}\quad{\footnotesize\texttt{TERM-0542}}\par
\noindent\emph{Iztočnik / Source:} Exponent eines Primärideals\par
\noindent\emph{Rusky / Russian:} показатель первичного идеала\par
\noindent\emph{Ukrajinsky / Ukrainian:} показник первинного ідеалу\par
\smallskip\noindent\emph{Razlog:} V oddělu 8 raboty 30 eksponent primarnogo ideala prodolžaje termin iz oddělov 5 i 6. Zapisani sravniteljni izrazy pokazatelj i pokaznyk ne imajut isti korenj, ale imajut odigrati tu samu funkciju v tvrdženju. Prědnost eksponent motivovana jest kontinuitetom i potrebom odlišiti čislo od stepeni ideala. To ne jest glasovanje, v ktorom ruska i ukrajinska forma fonetično podderžujut eksponent.\par
\noindent\emph{Reason:} In section 8 of Paper 30, eksponent primarnogo ideala continues the terminology of sections 5 and 6. The recorded comparanda pokazatelj and pokaznyk have a different root but serve the corresponding role in the statement. Preference for eksponent is motivated by continuity and the need to distinguish the number from an ideal power. This is not a vote in which the Russian and Ukrainian forms phonetically support eksponent.\par
\noindent{\footnotesize Izvorne zapisy / Original records: HIST-g115-004.}\par
\end{minipage}\par
\par\medskip\Needspace{7\baselineskip}
\noindent\begin{minipage}{\linewidth}
\hypertarget{TERM-0543}{}
\noindent\textbf{ekstremaly}\quad{\footnotesize\texttt{TERM-0543}}\par
\noindent\emph{Iztočnik / Source:} Extremalen\par
\noindent\emph{Rusky / Russian:} экстремали\par
\noindent\emph{Ukrajinsky / Ukrainian:} екстремалі\par
\smallskip\noindent\emph{Razlog:} Ekstremaly imenuje krive iz variacijnoj teorije, ktore argument s normalnymi koordinatami prěobražaje v prjame. Zapis izbira internacionalny tehničny nazv, blizky ruskemu i ukrajinskomu sravnjenju, vměsto opisnoj zaměny. Predmetom sut krive, ne krajni čislene cěnnosti; zamešati ekstremaly s maksimumami ili minimumami by proměnilo objekt argumenta.\par
\noindent\emph{Reason:} Ekstremaly names the curves from variational calculus which the normal-coordinate argument maps to straight lines. The record chooses the international technical name, close to the Russian and Ukrainian comparanda, rather than a descriptive substitute. The objects are curves, not extreme numerical values; confusing extremals with maxima or minima would change the argument's subject.\par
\noindent{\footnotesize Izvorne zapisy / Original records: HIST-g074-015.}\par
\end{minipage}\par
\par\medskip\Needspace{7\baselineskip}
\noindent\begin{minipage}{\linewidth}
\hypertarget{TERM-0544}{}
\noindent\textbf{provodnik}\quad{\footnotesize\texttt{TERM-0544}}\par
\noindent\emph{Iztočnik / Source:} F?hrer\par
\noindent\emph{Rusky / Russian:} [Bukvy ne sut zachovane / Lettering not recoverable]\par
\noindent\emph{Ukrajinsky / Ukrainian:} [Bukvy ne sut zachovane / Lettering not recoverable]\par
\smallskip\noindent\emph{Razlog:} Provodnik jest izbrano za čislovoteoretičny conductor. Razlog se opiraje na slovjansky korenj i ustaljeny tehničny pojem, ne na obyčno značenje osoby, ktora vodi. Německo zaglavje F?hrer ima izgubjenu bukvu, a ruske i ukrajinske bukvy ne sut zachovane. Možno izložiti zabilěženy smisyl izbora, ale ne vydavati oštečene sravnjenja za nezavisne slovnične attestacije.\par
\noindent\emph{Reason:} Provodnik is selected for the number-theoretic conductor. The rationale appeals to a Slavic root and an established technical concept, not the ordinary sense of a person who guides. The German heading F?hrer has a lost letter, and the Russian and Ukrainian lettering is not preserved. The recorded intended sense can be explained without presenting damaged comparison fields as independent dictionary attestations.\par
\noindent{\footnotesize Izvorne zapisy / Original records: HIST-g209-013.}\par
\end{minipage}\par
\par\medskip\Needspace{7\baselineskip}
\noindent\begin{minipage}{\linewidth}
\hypertarget{TERM-0545}{}
\noindent\textbf{faktorna sistema}\quad{\footnotesize\texttt{TERM-0545}}\par
\noindent\emph{Iztočnik / Source:} Faktorensystem\par
\noindent\emph{Rusky / Russian:} [Bukvy ne sut zachovane / Lettering not recoverable]\par
\noindent\emph{Ukrajinsky / Ukrainian:} [Bukvy ne sut zachovane / Lettering not recoverable]\par
\smallskip\noindent\emph{Razlog:} Faktorna sistema jest pridavnik s imenem, prodolžajuči izbor iz raboty 39 v rabotě 41. Motivacija jest morfologična kontinuitet, ne novy smisyl faktora. Ruske i ukrajinske bukvy v tutom zapisě ne sut zachovane, zato sam zapis ne može potvrditi jejich točnu grafiju. Čitatelj može sravniti nepovrěđene formy v TERM-0547, bez pretvarjanja, že one sut obnovjene bukvy tutogo zapisa.\par
\noindent\emph{Reason:} Faktorna sistema uses adjective plus noun, continuing the choice from Paper 39 into Paper 41. The motivation is morphological continuity, not a new meaning for factor. The Russian and Ukrainian lettering in this record is lost, so the record cannot establish those exact spellings. Readers can compare the intact forms in TERM-0547 without treating them as recovered letters from this record.\par
\noindent{\footnotesize Izvorne zapisy / Original records: HIST-g211-007.}\par
\end{minipage}\par
\par\medskip\Needspace{7\baselineskip}
\noindent\begin{minipage}{\linewidth}
\hypertarget{TERM-0546}{}
\noindent\textbf{faktorna sistema jedinica}\quad{\footnotesize\texttt{TERM-0546}}\par
\noindent\emph{Iztočnik / Source:} Faktorensystem Eins\par
\noindent\emph{Rusky / Russian:} фактор-система единица\par
\noindent\emph{Ukrajinsky / Ukrainian:} фактор-система одиниця\par
\smallskip\noindent\emph{Razlog:} Faktorna sistema jedinica sohranjaje historijnu formulaciju Noether. Razlog izrično odbacaje zaměnu s trivialny kocikl, zato že formuly raboty upotrěbjajut jezyk faktornyh sistem. Sravnjenje s pozdnějšej terminologijej može byti polezno v objasnjenju, ale ne smi tiho zaměniti avtorske pojmenovanje v prěvodě. Jedinica jest čest toj formulacije, ne čislo sistemov.\par
\noindent\emph{Reason:} Faktorna sistema jedinica preserves Noether's historical formulation. The rationale explicitly rejects replacing it with trivial cocycle because the paper's formulas use factor-system terminology. Comparison with later terminology may help in an explanation, but must not silently replace the author's designation in the translation. Jedinica belongs to that formulation; it is not a count of systems.\par
\noindent{\footnotesize Izvorne zapisy / Original records: HIST-g212-003.}\par
\end{minipage}\par
\par\medskip\Needspace{7\baselineskip}
\noindent\begin{minipage}{\linewidth}
\hypertarget{TERM-0547}{}
\noindent\textbf{faktorne sistemy}\quad{\footnotesize\texttt{TERM-0547}}\par
\noindent\emph{Iztočnik / Source:} Faktorensysteme\par
\noindent\emph{Rusky / Russian:} [Bukvy ne sut zachovane / Lettering not recoverable]; фактор-системы\par
\noindent\emph{Ukrajinsky / Ukrainian:} [Bukvy ne sut zachovane / Lettering not recoverable]; фактор-системи\par
\smallskip\noindent\emph{Razlog:} Zapisy k rabotam 32 i 39 izbirajut faktorne sistemy v Brauerovoj konstrukciji. Pridavnik faktorne prěvodi složeny nazv, ne ostavjajuči faktor kako samostatno ime; historijni razlog dava tomu prednost prěd oblikom s defisom faktor-sistema. Vtoroj motiv jest kontinuitet, osoblivo v podnožnoj primětkě, kde novy sinonim byl by nepotrěbny. Opis formy s defisom kako rusko podobnoj jest stilistična ocěna zapisa, ne dokaz, že medžuslovjanska gramatika takovu formu zakazuje.\par
\noindent\emph{Reason:} Records for Papers 32 and 39 choose faktorne sistemy in Brauer's construction. Adjectival faktorne renders the compound without leaving faktor as a separate noun; the historical rationale prefers this to hyphenated faktor-sistema. A second motive is continuity, especially in a footnote where a new synonym would be unnecessary. Calling the hyphenated shape Russian-looking is the record's stylistic assessment, not evidence that Interslavic grammar prohibits it.\par
\noindent{\footnotesize Izvorne zapisy / Original records: HIST-g168-006, HIST-g176-005, HIST-g209-005.}\par
\end{minipage}\par
\par\medskip\Needspace{7\baselineskip}
\noindent\begin{minipage}{\linewidth}
\hypertarget{TERM-0548}{}
\noindent\textbf{faktorgrupa}\quad{\footnotesize\texttt{TERM-0548}}\par
\noindent\emph{Iztočnik / Source:} Faktorgruppe\par
\noindent\emph{Rusky / Russian:} фактор-группа\par
\noindent\emph{Ukrajinsky / Ukrainian:} фактор-група\par
\smallskip\noindent\emph{Razlog:} Faktorgrupa sohranjaje internacionalny složeny nazv, sravnjeny s ruskym i ukrajinskim faktor-grupa. Razlog jest kratky: on nazyvaje termin sovrěmennym i prozračnym, bez osobnogo sravnjenja s drugymi medžuslovjanskimi kandidatami. Pisanje jednym slovom razlikuje se od sravniteljnyh form s defisom. Zapis ne dokazuje ni obšču pravopisnu normu za vse kompozity s faktor-, ni čitateljski izměrenu razumělivost.\par
\noindent\emph{Reason:} Faktorgrupa retains the international compound, compared with Russian and Ukrainian hyphenated forms. The reason is brief: it calls the term modern and transparent without separately comparing other Interslavic candidates. Writing it as one word differs from the hyphenated comparanda. The record establishes neither a general spelling rule for every faktor- compound nor reader-tested intelligibility.\par
\noindent{\footnotesize Izvorne zapisy / Original records: HIST-g180-003.}\par
\end{minipage}\par
\par\medskip\Needspace{7\baselineskip}
\noindent\begin{minipage}{\linewidth}
\hypertarget{TERM-0549}{}
\noindent\textbf{svijanje}\quad{\footnotesize\texttt{TERM-0549}}\par
\noindent\emph{Iztočnik / Source:} Faltung\par
\noindent\emph{Rusky / Russian:} свертка\par
\noindent\emph{Ukrajinsky / Ukrainian:} згортання\par
\smallskip\noindent\emph{Razlog:} Sedm zapisov podderžuje svijanje za Faltung v rabotě 2. Glavni argument jest kontinuitet konstrukcije v tablicah, redukcijah, dokazah i završnoj listě form; kratko ime ne obtęžava ponavljane redy. Ruska i ukrajinska sravnjenja ukazujut k svęrtce ili kontrakciji, ale zapis iz odděla 26 ješče ostavjaje medžuslovjanski izbor k odbornomu razsmotrenju. Ponovjenje istogo izbora ne jest novym nezavisnym dokazom njegovoj priměrnosti. Aktualne pasaže s kontrakcija navodime odvojeno.\par
\noindent\emph{Reason:} Seven records support svijanje for Faltung in Paper 2. The main argument is continuity across construction tables, reductions, proofs and the concluding form list; a compact noun avoids burdening repeated lines. The Russian and Ukrainian comparisons point towards folding or contraction, but the section 26 record still leaves the Interslavic choice for specialist consideration. Repetition of the same choice is not fresh independent evidence of its suitability. Current passages using kontrakcija are reported separately.\par
\noindent{\footnotesize Izvorne zapisy / Original records: HIST-g012-002, HIST-g013-002, HIST-g014-003, HIST-g015-003, HIST-g016-003, HIST-g017-003, HIST-g027-001.}\par
\smallskip\noindent\emph{Sravnjenje s tekstom:} V oddělu 10 raboty 2 aktualna tablica ima kontrakcija, s grupoj $K,j,\mathfrak A$ s $\nu$ i odvojeno $f$ s $H$. Lokalna definicija TeX-komandy potvrđaje gotičnu $\mathfrak A$. V oddělu 19 opet stoi kontrakcija, ale spis argumentov počinaje latiničnym $v$, dokolě historijni glosar ima grečsku $\nu$. Tuto jest konkretna razlika, ne dopuščena tiha zaměna simbolov; bez sravnjenja s pečatnym iztočnikom ne odločujemo, ktory zapis jest pravilen. Tute dve pasaže potvrđajut pozdnějšu formu kontrakcija v tyh městah, ale ne dokazujut cělokorpusnu zaměnu svijanje ili skruženje.\par
\noindent\emph{Textual comparison:} In section 10 of Paper 2, the current table uses kontrakcija, grouping $K,j,\mathfrak A$ with $\nu$ and separately $f$ with $H$. The local TeX command definition confirms fraktur $\mathfrak A$. Section 19 again has kontrakcija, but its argument list begins with Latin $v$, whereas the historical glossary has Greek $\nu$. This is a concrete discrepancy, not permission for a silent symbol substitution; neither reading is adjudicated without comparison with the printed source. The two passages establish later kontrakcija at these locations, not a corpus-wide replacement of svijanje or skruženje.\par
\noindent{\footnotesize TEXT-COMP-16; pasaže / passages: 1277, 1286, 1287, 1288, 2653; \texttt{CORPUS\_FORM\_COMPARISONS.json}.}\par
\end{minipage}\par
\par\medskip\Needspace{7\baselineskip}
\noindent\begin{minipage}{\linewidth}
\hypertarget{TERM-0550}{}
\noindent\textbf{svijanje / skruženje}\quad{\footnotesize\texttt{TERM-0550}}\par
\noindent\emph{Iztočnik / Source:} Faltung\par
\noindent\emph{Rusky / Russian:} свёртка\par
\noindent\emph{Ukrajinsky / Ukrainian:} згортання\par
\smallskip\noindent\emph{Razlog:} Svijanje / skruženje jest nerazrěšena para kandidatov, ne dva dokazano ravnopravne termina. Zapis upominaje nizozemsko contractie kako čitateljski izbor i samo prědpoklada, že slovjansky iztočniki dopustet zgortannja ili svërtka. On ne navodi iztočnu pasažu ni razdělene ocěny dvoh medžuslovjanskih kandidatov. Slovnik sohranjaje tu nejistost; pozdnějša upotrěba svijanje ne može nazad dopisati otsutno sravnjenje.\par
\noindent\emph{Reason:} Svijanje / skruženje is an unresolved candidate pair, not two established interchangeable terms. The record mentions Dutch contractie as a readability choice and only conjectures that Slavic references would accommodate zgortannja or svërtka. It supplies neither an exact source passage nor separate assessments of the two Interslavic candidates. The dictionary preserves that uncertainty; later use of svijanje cannot retrospectively supply the missing comparison.\par
\noindent{\footnotesize Izvorne zapisy / Original records: HIST-g001-008.}\par
\smallskip\noindent\emph{Sravnjenje s tekstom:} V oddělu 10 raboty 2 aktualna tablica ima kontrakcija, s grupoj $K,j,\mathfrak A$ s $\nu$ i odvojeno $f$ s $H$. Lokalna definicija TeX-komandy potvrđaje gotičnu $\mathfrak A$. V oddělu 19 opet stoi kontrakcija, ale spis argumentov počinaje latiničnym $v$, dokolě historijni glosar ima grečsku $\nu$. Tuto jest konkretna razlika, ne dopuščena tiha zaměna simbolov; bez sravnjenja s pečatnym iztočnikom ne odločujemo, ktory zapis jest pravilen. Tute dve pasaže potvrđajut pozdnějšu formu kontrakcija v tyh městah, ale ne dokazujut cělokorpusnu zaměnu svijanje ili skruženje.\par
\noindent\emph{Textual comparison:} In section 10 of Paper 2, the current table uses kontrakcija, grouping $K,j,\mathfrak A$ with $\nu$ and separately $f$ with $H$. The local TeX command definition confirms fraktur $\mathfrak A$. Section 19 again has kontrakcija, but its argument list begins with Latin $v$, whereas the historical glossary has Greek $\nu$. This is a concrete discrepancy, not permission for a silent symbol substitution; neither reading is adjudicated without comparison with the printed source. The two passages establish later kontrakcija at these locations, not a corpus-wide replacement of svijanje or skruženje.\par
\noindent{\footnotesize TEXT-COMP-16; pasaže / passages: 1277, 1286, 1287, 1288, 2653; \texttt{CORPUS\_FORM\_COMPARISONS.json}.}\par
\end{minipage}\par
\par\medskip\Needspace{7\baselineskip}
\noindent\begin{minipage}{\linewidth}
\hypertarget{TERM-0551}{}
\noindent\textbf{svijanje s $s$ samo kako polarizacija}\quad{\footnotesize\texttt{TERM-0551}}\par
\noindent\emph{Iztočnik / Source:} Faltung mit $s$ nur als Polarisation\par
\noindent\emph{Rusky / Russian:} свертка с $s$ только как поляризация\par
\noindent\emph{Ukrajinsky / Ukrainian:} згортання з $s$ лише як поляризація\par
\smallskip\noindent\emph{Razlog:} Samo kako polarizacija ograničaje smysl konstrukcije so simbolom $s$. Razlog razlikuje stvarno tvorjenje form od konvencije njih uporadkovanja, sohranjajuči historijno svijanje i internacionalno polarizacija. Izraz samo ne jest stilističny dodatok: bez njega fraza mogla by pripisati upotrěbjenoj operaciji širšu funkciju, než toj konvenciji prinadleži.\par
\noindent\emph{Reason:} Samo kako polarizacija restricts the meaning of the construction involving $s$. The rationale distinguishes the actual construction of forms from the convention for ordering them, retaining historical svijanje and international polarizacija. Only is not a stylistic addition: without it, the phrase could assign the operation a broader role than the convention permits.\par
\noindent{\footnotesize Izvorne zapisy / Original records: HIST-g018-009.}\par
\end{minipage}\par
\par\medskip\Needspace{7\baselineskip}
\noindent\begin{minipage}{\linewidth}
\hypertarget{TERM-0552}{}
\noindent\textbf{svijanje $s$ s $\nu; H; L; H^2$}\quad{\footnotesize\texttt{TERM-0552}}\par
\noindent\emph{Iztočnik / Source:} Faltung von $s$ mit $\nu; H; L; H^2$\par
\noindent\emph{Rusky / Russian:} свертка $s$ с $\nu; H; L; H^2$\par
\noindent\emph{Ukrajinsky / Ukrainian:} згортання $s$ з $\nu; H; L; H^2$\par
\smallskip\noindent\emph{Razlog:} Zapis nosi istu historijnu upotrěbu svijanje k operacijam so $s$ i so členami $\nu$, $H$, $L$, $H^2$. Sam razlog argumentuje samo terminologičnu kontinuitet; ne dava osobnogo opravdanja vsakej oznaky. Zato formula jest čest zachovanogo zapisa, ale ne avtomatično prověrena transkripcija. Sravnjenje s aktualnym tekstom pokazuje takože pytanje $\nu$ protiv $v$, odvojeno od izbora slova za Faltung.\par
\noindent\emph{Reason:} The record carries historical svijanje into the operations involving $s$ and the listed $\nu$, $H$, $L$, $H^2$. Its reason argues only for terminological continuity, not separately for each symbol. The formula is therefore preserved evidence, not automatically a verified transcription. Comparison with the current text also reveals a $\nu$ versus $v$ question, distinct from the word chosen for Faltung.\par
\noindent{\footnotesize Izvorne zapisy / Original records: HIST-g020-002.}\par
\smallskip\noindent\emph{Sravnjenje s tekstom:} V oddělu 10 raboty 2 aktualna tablica ima kontrakcija, s grupoj $K,j,\mathfrak A$ s $\nu$ i odvojeno $f$ s $H$. Lokalna definicija TeX-komandy potvrđaje gotičnu $\mathfrak A$. V oddělu 19 opet stoi kontrakcija, ale spis argumentov počinaje latiničnym $v$, dokolě historijni glosar ima grečsku $\nu$. Tuto jest konkretna razlika, ne dopuščena tiha zaměna simbolov; bez sravnjenja s pečatnym iztočnikom ne odločujemo, ktory zapis jest pravilen. Tute dve pasaže potvrđajut pozdnějšu formu kontrakcija v tyh městah, ale ne dokazujut cělokorpusnu zaměnu svijanje ili skruženje.\par
\noindent\emph{Textual comparison:} In section 10 of Paper 2, the current table uses kontrakcija, grouping $K,j,\mathfrak A$ with $\nu$ and separately $f$ with $H$. The local TeX command definition confirms fraktur $\mathfrak A$. Section 19 again has kontrakcija, but its argument list begins with Latin $v$, whereas the historical glossary has Greek $\nu$. This is a concrete discrepancy, not permission for a silent symbol substitution; neither reading is adjudicated without comparison with the printed source. The two passages establish later kontrakcija at these locations, not a corpus-wide replacement of svijanje or skruženje.\par
\noindent{\footnotesize TEXT-COMP-16; pasaže / passages: 1277, 1286, 1287, 1288, 2653; \texttt{CORPUS\_FORM\_COMPARISONS.json}.}\par
\end{minipage}\par
\par\medskip\Needspace{7\baselineskip}
\noindent\begin{minipage}{\linewidth}
\hypertarget{TERM-0553}{}
\noindent\textbf{svijanje $s$ s ...}\quad{\footnotesize\texttt{TERM-0553}}\par
\noindent\emph{Iztočnik / Source:} Faltung von $s$ mit ...\par
\noindent\emph{Rusky / Russian:} свертка $s$ с ...\par
\noindent\emph{Ukrajinsky / Ukrainian:} згортання $s$ з ...\par
\smallskip\noindent\emph{Razlog:} Svijanje $s$ s ... prodolžaje ranějši termin raboty 2 v oddělu 20. Tri točky sut skračenjem v glosarnom zaglavju, ne dopolnjenym spisom argumentov. Zapis podkrepljaje kontinuitet konstrukcije, ale sam ne dovoljuje obnoviti izpuščene formy. Dlja konkretnoj matematičnoj upotrěby čitatelj mora iměti polnu pasažu, ne samo tuto zaglavje.\par
\noindent\emph{Reason:} Svijanje $s$ s ... continues the earlier Paper 2 term in section 20. The ellipsis abbreviates the glossary heading; it is not a complete list of arguments. The record supports continuity of the construction but does not recover the omitted forms. A mathematical application requires the full passage, not this heading alone.\par
\noindent{\footnotesize Izvorne zapisy / Original records: HIST-g021-002.}\par
\end{minipage}\par
\par\medskip\Needspace{7\baselineskip}
\noindent\begin{minipage}{\linewidth}
\hypertarget{TERM-0554}{}
\noindent\textbf{svijanje}\quad{\footnotesize\texttt{TERM-0554}}\par
\noindent\emph{Iztočnik / Source:} Faltung von ...\par
\noindent\emph{Rusky / Russian:} свертка\par
\noindent\emph{Ukrajinsky / Ukrainian:} згортання\par
\smallskip\noindent\emph{Razlog:} V oddělu 18 raboty 2 zapis opet dava prednost kratkomu svijanje, zato že isto ime može byti ponovjeno v sistemah i podsistemah. Ruske i ukrajinske formy sut označene kako idiomatične, ale bez navedenoj pasáže iz odbornogo iztočnika. Okrěm kontinuiteta i kratkosti zapis ne sohranjaje posebnu argumentaciju za novy kandidat. Zaglavje Faltung von ... ne jest polna specifikacija operacije.\par
\noindent\emph{Reason:} For section 18 of Paper 2, the record again favours compact svijanje because the same noun can recur across systems and subsystems. The Russian and Ukrainian forms are described as idiomatic without a cited passage from specialist writing. Beyond continuity and brevity, no separate argument for a new candidate is preserved. Faltung von ... is not a complete specification of the operation.\par
\noindent{\footnotesize Izvorne zapisy / Original records: HIST-g019-002.}\par
\end{minipage}\par
\par\medskip\Needspace{7\baselineskip}
\noindent\begin{minipage}{\linewidth}
\hypertarget{TERM-0555}{}
\noindent\textbf{svijanje $f$ s $\nu$}\quad{\footnotesize\texttt{TERM-0555}}\par
\noindent\emph{Iztočnik / Source:} Faltung von $f$ mit $\nu$\par
\noindent\emph{Rusky / Russian:} свертка $f$ с $\nu$\par
\noindent\emph{Ukrajinsky / Ukrainian:} згортання $f$ з $\nu$\par
\smallskip\noindent\emph{Razlog:} Svijanje $f$ s $\nu$ sohranjaje izbrany nazv i oba argumenta operacije. Zapis k oddělu 8 raboty 2 ne predstavjaje novu ocěnu alternativov: odkazuje k ranějšej parě ruske svërtka, ukrajinske zgortannja i medžuslovjanske svijanje. Bukva $f$ i grečska $\nu$ sut matematične oznaky, ne slova dlja prěvoda; njih identičnost mora ostati vidima i v kirilici.\par
\noindent\emph{Reason:} Svijanje $f$ s $\nu$ retains the selected noun and both arguments of the operation. The record for section 8 of Paper 2 does not supply a fresh assessment of alternatives; it refers back to the Russian, Ukrainian and Interslavic choices. Latin $f$ and Greek $\nu$ are mathematical symbols, not words to translate; their identity must remain visible in Cyrillic too.\par
\noindent{\footnotesize Izvorne zapisy / Original records: HIST-g009-002.}\par
\end{minipage}\par
\par\medskip\Needspace{7\baselineskip}
\noindent\begin{minipage}{\linewidth}
\hypertarget{TERM-0556}{}
\noindent\textbf{svijanje $K,j,\mathfrak{A}$ s $\nu$; $f$ s $H$}\quad{\footnotesize\texttt{TERM-0556}}\par
\noindent\emph{Iztočnik / Source:} Faltung von $K,j,\mathfrak{A}$ mit $\nu$; $f$ mit $H$\par
\noindent\emph{Rusky / Russian:} свертка $K,j,\mathfrak{A}$ с $\nu$; $f$ с $H$\par
\noindent\emph{Ukrajinsky / Ukrainian:} згортання $K,j,\mathfrak{A}$ з $\nu$; $f$ з $H$\par
\smallskip\noindent\emph{Razlog:} Razvětvena konstrukcija grupuje $K,j,\mathfrak A$ s $\nu$, a odvojeno $f$ s $H$. Razlog upominaje kratku formu s figurnoju skobkoju i podderžuje ranějši termin svijanje. Razděljenje na dve grupy jest sodržanje zapisa, ne samo tipografija. Komanda A v iztočnom TeX-u označaje gotičnu $\mathfrak A$, ne obyčnu latiničnu A; čitateljska forma mora sohraniti tuto razlikovanje, kako i pravym argumentam prinadležeče predlogy.\par
\noindent\emph{Reason:} The branched construction groups $K,j,\mathfrak A$ with $\nu$, and separately $f$ with $H$. The reason mentions the compact brace arrangement and supports the earlier term svijanje. Division into the two groups is part of the content, not merely typography. The source TeX command A denotes fraktur $\mathfrak A$, not ordinary Latin A; the readable form must preserve that distinction and the correct association of arguments.\par
\noindent{\footnotesize Izvorne zapisy / Original records: HIST-g011-002.}\par
\smallskip\noindent\emph{Sravnjenje s tekstom:} V oddělu 10 raboty 2 aktualna tablica ima kontrakcija, s grupoj $K,j,\mathfrak A$ s $\nu$ i odvojeno $f$ s $H$. Lokalna definicija TeX-komandy potvrđaje gotičnu $\mathfrak A$. V oddělu 19 opet stoi kontrakcija, ale spis argumentov počinaje latiničnym $v$, dokolě historijni glosar ima grečsku $\nu$. Tuto jest konkretna razlika, ne dopuščena tiha zaměna simbolov; bez sravnjenja s pečatnym iztočnikom ne odločujemo, ktory zapis jest pravilen. Tute dve pasaže potvrđajut pozdnějšu formu kontrakcija v tyh městah, ale ne dokazujut cělokorpusnu zaměnu svijanje ili skruženje.\par
\noindent\emph{Textual comparison:} In section 10 of Paper 2, the current table uses kontrakcija, grouping $K,j,\mathfrak A$ with $\nu$ and separately $f$ with $H$. The local TeX command definition confirms fraktur $\mathfrak A$. Section 19 again has kontrakcija, but its argument list begins with Latin $v$, whereas the historical glossary has Greek $\nu$. This is a concrete discrepancy, not permission for a silent symbol substitution; neither reading is adjudicated without comparison with the printed source. The two passages establish later kontrakcija at these locations, not a corpus-wide replacement of svijanje or skruženje.\par
\noindent{\footnotesize TEXT-COMP-16; pasaže / passages: 1277, 1286, 1287, 1288, 2653; \texttt{CORPUS\_FORM\_COMPARISONS.json}.}\par
\end{minipage}\par
\par\medskip\Needspace{7\baselineskip}
\noindent\begin{minipage}{\linewidth}
\hypertarget{TERM-0557}{}
\noindent\textbf{svijanje $s$ s formami 4. poręda}\quad{\footnotesize\texttt{TERM-0557}}\par
\noindent\emph{Iztočnik / Source:} Faltung von $s$ mit den Formen 4. Ordnung\par
\noindent\emph{Rusky / Russian:} свертка $s$ с формами 4-го порядка\par
\noindent\emph{Ukrajinsky / Ukrainian:} згортання $s$ з формами 4-го порядку\par
\smallskip\noindent\emph{Razlog:} V svijanje $s$ s formami 4. poręda prvy znak s jest matematična proměnna, a vtory znak s jest predlog. Zapis izrično sohranjaje oboje; odstranjati navodno ponovjeno slovo by zničilo argument ili gramatiku. Čislo 4 označaje pored form, ne količestvo form. Termin svijanje prodolžaje ranějši izbor, ale ne pretvarjaje operaciju v libovoljnu analitičnu konvoluciju.\par
\noindent\emph{Reason:} In svijanje $s$ s formami 4. poręda, the first s is a mathematical variable and the second is a preposition. The record explicitly retains both; deleting an apparent repeated word would remove either an argument or the grammar. The number 4 specifies the order of the forms, not their number. Svijanje continues the earlier choice without thereby identifying the operation with arbitrary analytic convolution.\par
\noindent{\footnotesize Izvorne zapisy / Original records: HIST-g022-002.}\par
\end{minipage}\par
\par\medskip\Needspace{7\baselineskip}
\noindent\begin{minipage}{\linewidth}
\hypertarget{TERM-0558}{}
\noindent\textbf{svijanje $s$ s formami 7. poręda}\quad{\footnotesize\texttt{TERM-0558}}\par
\noindent\emph{Iztočnik / Source:} Faltung von $s$ mit den Formen 7. Ordnung\par
\noindent\emph{Rusky / Russian:} свертка $s$ с формами 7-го порядка\par
\noindent\emph{Ukrajinsky / Ukrainian:} згортання $s$ з формами 7-го порядку\par
\smallskip\noindent\emph{Razlog:} Fraza sohranjaje operaciju so $s$ i izrično označaje formy sedmogo poreda. Razlog jest kontinuitet s ranějšimi zaglavjami, zajedno s točnym prěnosom čisla poreda. Sedmy pored ne znači sedm form; skračenje 7. jest čitateljska forma klasifikacije, ne proměna matematičnogo ograničenja.\par
\noindent\emph{Reason:} The phrase retains the operation involving $s$ and expressly identifies forms of seventh order. The reason is continuity with earlier headings together with faithful transmission of the order number. Seventh order does not mean seven forms; 7. abbreviates the classification without changing the mathematical restriction.\par
\noindent{\footnotesize Izvorne zapisy / Original records: HIST-g025-002.}\par
\end{minipage}\par
\par\medskip\Needspace{7\baselineskip}
\noindent\begin{minipage}{\linewidth}
\hypertarget{TERM-0559}{}
\noindent\textbf{svijanje $s$ s formami 5. poręda i s produktom}\quad{\footnotesize\texttt{TERM-0559}}\par
\noindent\emph{Iztočnik / Source:} Faltung von $s$ mit den Formen fünfter Ordnung und dem Produkt\par
\noindent\emph{Rusky / Russian:} свертка $s$ с формами 5-го порядка и с произведением\par
\noindent\emph{Ukrajinsky / Ukrainian:} згортання $s$ з формами 5-го порядку і з добутком\par
\smallskip\noindent\emph{Razlog:} S formami 5. poręda i s produktom sohranjaje dva roda objektov, s ktorymi provodi se operacija. Razlog izrično odlišaje produkt od dalšího svijanja: produkt jest argument, ne novo ime za operaciju. Ponovjeny predlog s pomaga pokazati koordinaciju, pri tom kako prvy znak $s$ ostavaje simbolom. Izpustiti i s produktom by suzilo konstrukciju na same formy petogo poreda.\par
\noindent\emph{Reason:} S formami 5. poręda i s produktom retains two kinds of objects with which the operation is performed. The rationale expressly distinguishes the product from another folding operation: the product is an argument, not a new name for the operation. Repeating the preposition s helps display the coordination, while the first $s$ remains a symbol. Omitting and with the product would restrict the construction to fifth-order forms alone.\par
\noindent{\footnotesize Izvorne zapisy / Original records: HIST-g023-002.}\par
\end{minipage}\par
\par\medskip\Needspace{7\baselineskip}
\noindent\begin{minipage}{\linewidth}
\hypertarget{TERM-0560}{}
\noindent\textbf{svijanje $s$ s formami 6. poręda}\quad{\footnotesize\texttt{TERM-0560}}\par
\noindent\emph{Iztočnik / Source:} Faltung von $s$ mit den Formen sechster Ordnung\par
\noindent\emph{Rusky / Russian:} свертка $s$ с формами 6-го порядка\par
\noindent\emph{Ukrajinsky / Ukrainian:} згортання $s$ з формами 6-го порядку\par
\smallskip\noindent\emph{Razlog:} Svijanje $s$ s formami 6. poręda nosi isti nazv operacije k formam šestogo poreda. Zapis podkrepljaje kontinuitet i tvrđenu prozračnost svijanje, pri sravnjenju s ruskym i ukrajinskim imenami. Tuto jest argument za stabilno pojmenovanje v seriji podobnyh konstrukcij, ne novy test razumělivosti. Oznaka $s$ i čislo poreda ostavajut samostatnymi čestjami matematičnoj informaciji.\par
\noindent\emph{Reason:} Svijanje $s$ s formami 6. poręda carries the same operation name to forms of sixth order. The record stresses continuity and the claimed transparency of svijanje, alongside the Russian and Ukrainian comparison nouns. This is an argument for stable naming across a series of related constructions, not a new intelligibility test. The symbol $s$ and the order number remain separate pieces of mathematical information.\par
\noindent{\footnotesize Izvorne zapisy / Original records: HIST-g024-002.}\par
\end{minipage}\par
\par\medskip\Needspace{7\baselineskip}
\noindent\begin{minipage}{\linewidth}
\hypertarget{TERM-0561}{}
\noindent\textbf{svijanje $\theta$ s $\nu$}\quad{\footnotesize\texttt{TERM-0561}}\par
\noindent\emph{Iztočnik / Source:} Faltung von $\theta$ mit $\nu$\par
\noindent\emph{Rusky / Russian:} свертка $\theta$ с $\nu$\par
\noindent\emph{Ukrajinsky / Ukrainian:} згортання $\theta$ з $\nu$\par
\smallskip\noindent\emph{Razlog:} Svijanje $\theta$ s $\nu$ sohranjaje historijni nazv operacije iz raboty 2 i oba označene argumenta. Zapis k oddělu 9 odkazuje k ranějšim ruskym, ukrajinskim i medžuslovjanskim izboram, ale ne dava novogo sravnjenja kandidatov. Grečske bukvy sut oznaky form; zaměna jednoj s podobnoju latiničnoju bukvoju ne jest dopustima jezyčna adaptacija.\par
\noindent\emph{Reason:} Svijanje $\theta$ s $\nu$ retains the historical operation name from Paper 2 and both specified arguments. The section 9 record refers to earlier Russian, Ukrainian and Interslavic choices but provides no fresh candidate comparison. The Greek letters identify forms; replacing either with a similar Latin letter would not be a permissible linguistic adaptation.\par
\noindent{\footnotesize Izvorne zapisy / Original records: HIST-g010-002.}\par
\end{minipage}\par
\par\medskip\Needspace{7\baselineskip}
\noindent\begin{minipage}{\linewidth}
\hypertarget{TERM-0562}{}
\noindent\textbf{svijanja s modulom}\quad{\footnotesize\texttt{TERM-0562}}\par
\noindent\emph{Iztočnik / Source:} Faltungen mit dem Modul\par
\noindent\emph{Rusky / Russian:} свёртки с модулем\par
\noindent\emph{Ukrajinsky / Ukrainian:} згортання з модулем\par
\smallskip\noindent\emph{Razlog:} Množinske svijanja s modulom prodolžaje termin iz oddělov 1-3 v oddělu 5 raboty 2. Zapis upominaje sposobnost imenja svijanje tvoriti take frazy, ale ostavjaje jeho priměrnost za daljše raboty otvorenoju. Modul tu jest čest terminologije invariantnoj teorije; bez definicije ne trěba pridavati mu vse svojstva modulov iz pozdnějšej abstraktnoj algebry.\par
\noindent\emph{Reason:} The plural svijanja s modulom carries the term from sections 1-3 into section 5 of Paper 2. The record invokes the noun's usefulness in such constructions but leaves its suitability for later papers open. Modul belongs here to the terminology of invariant theory; without the definition, it must not be assigned all the properties of modules in later abstract algebra.\par
\noindent{\footnotesize Izvorne zapisy / Original records: HIST-g006-013.}\par
\end{minipage}\par
\par\medskip\Needspace{7\baselineskip}
\noindent\begin{minipage}{\linewidth}
\hypertarget{TERM-0563}{}
\noindent\textbf{proces svijanja}\quad{\footnotesize\texttt{TERM-0563}}\par
\noindent\emph{Iztočnik / Source:} Faltungsprozeß\par
\noindent\emph{Rusky / Russian:} процесс свертки\par
\noindent\emph{Ukrajinsky / Ukrainian:} процес згортання\par
\smallskip\noindent\emph{Razlog:} Proces svijanja razklada německy kompozit na ime postupa i zavisno ime operacije. Razlog jest kontinuitet s ranějšim prěvodom Faltung kako svijanje. Zapis takože izrično ostavjaje izbor k odbornomu razsmotrenju. Nazvati formu prozračnoju bylo ocěnoju pri izborě, ne rezultatom izměrjenogo razuměnja čitateljev.\par
\noindent\emph{Reason:} Proces svijanja renders the German compound as a process noun with the operation's name dependent on it. The reason is continuity with earlier svijanje for Faltung. The record also explicitly leaves the choice for specialist consideration. Calling the form transparent was a selection judgment, not the result of a measured reader-comprehension study.\par
\noindent{\footnotesize Izvorne zapisy / Original records: HIST-g031-005.}\par
\end{minipage}\par
\par\medskip\Needspace{7\baselineskip}
\noindent\begin{minipage}{\linewidth}
\hypertarget{TERM-0564}{}
\noindent\textbf{proces svijanja / svijanje}\quad{\footnotesize\texttt{TERM-0564}}\par
\noindent\emph{Iztočnik / Source:} Faltungsprozeß / Faltung\par
\noindent\emph{Rusky / Russian:} процесс свертки / свертка; процесс свёртки / свёртка\par
\noindent\emph{Ukrajinsky / Ukrainian:} процес згортання / згортання\par
\smallskip\noindent\emph{Razlog:} Tri zapisy povezujut proces svijanja s svijanje. Formalna definicija v rabotě 2 dava prednost stalnomu imenu prěd stilističnym měnjanjem; zapisy k rabotě 4 prodolžajut isto imenuvanje. Oni ne predstavajut izbor kako dovršeny: osoblivo razlikujut simboličnu invariantnoteoretičnu kontrakciju od obyčnogo svivanja i od analitičnoj konvolucije. Preto kratkost i slovotvorna produktivnost ne sut dovoljne bez matematičnoj definicije.\par
\noindent\emph{Reason:} Three records connect proces svijanja with svijanje. The formal definition in Paper 2 favours a stable name over stylistic variation; the Paper 4 records continue it. They do not present the selection as settled: in particular, they distinguish symbolic invariant-theoretic contraction from ordinary folding and from analytic convolution. Brevity and productive word formation therefore remain insufficient without the mathematical definition.\par
\noindent{\footnotesize Izvorne zapisy / Original records: HIST-g003-001, HIST-g035-004, HIST-g037-002.}\par
\smallskip\noindent\emph{Sravnjenje s tekstom:} Aktualny odděl 1 raboty 2 naziva osnovny postup kontrakcijeju. Red form jest početna forma zajedno so vsimi formami dobytymi kontrakcijami samej so soboju; vyši formy razlikujut se po kontrakcijah. Pri veče redah simbolov tekst trěbuje izbrati specialne kontrakcije, da red byl jednoznačno normovany, i potom pokazuje premougolnu shemu. Tuto podkrepljaje izbor kratkogo imena za oprěděljenu generovanu strukturu, ne prosto množinu ili analitičny red. Aktualno red i kontrakcija razlikujut se od historijnyh ręd i svijanje; sravnjenje ne dokazuje cělokorpusno uskladnjenje ni oficijalny status kandidatov.\par
\noindent\emph{Textual comparison:} The current section 1 of Paper 2 calls the basic process contraction. A row of forms comprises an initial form and all forms obtained by self-contractions; higher forms are distinguished by the contractions performed. With more rows of symbols, the text requires selection of particular contractions to normalize the row uniquely, and then displays a rectangular scheme. This supports a compact name for a defined generated structure, not merely a set or an analytic series. Current red and kontrakcija differ from historical ręd and svijanje; the comparison does not establish corpus-wide harmonization or official status for the candidates.\par
\noindent{\footnotesize TEXT-COMP-19; pasaže / passages: 243, 295, 299, 301; \texttt{CORPUS\_FORM\_COMPARISONS.json}.}\par
\end{minipage}\par
\par\medskip\Needspace{7\baselineskip}
\noindent\begin{minipage}{\linewidth}
\hypertarget{TERM-0565}{}
\noindent\textbf{proces svijanja / svijanja}\quad{\footnotesize\texttt{TERM-0565}}\par
\noindent\emph{Iztočnik / Source:} Faltungsprozeß / Faltungen\par
\noindent\emph{Rusky / Russian:} процесс свертки / свертки\par
\noindent\emph{Ukrajinsky / Ukrainian:} процес згортання / згортання\par
\smallskip\noindent\emph{Razlog:} Proces svijanja / svijanja sohranjaje razlikovanje medžu postupom i množestvom pojedinyh operacij v rabotě 3. Zapis opiraje se na kontinuitet rabot 1-2, ne na novo sravnjenje slov. On ješče označaje medžuslovjansky izbor kako potrěbujuči odbornogo razsmotrenja. Množinska forma zato jest dokumentovana upotrěba kandidata, ne dodatno dokazanje jeho normativnogo statusa.\par
\noindent\emph{Reason:} Proces svijanja / svijanja distinguishes the process from multiple individual operations in Paper 3. The record relies on continuity with Papers 1-2, not a new comparison of words. It still marks the Interslavic selection for specialist consideration. The plural is therefore documented use of a candidate, not additional proof of its normative status.\par
\noindent{\footnotesize Izvorne zapisy / Original records: HIST-g028-006.}\par
\end{minipage}\par
\par\medskip\Needspace{7\baselineskip}
\noindent\begin{minipage}{\linewidth}
\hypertarget{TERM-0566}{}
\noindent\textbf{proces svijanja / osnovne svijanja}\quad{\footnotesize\texttt{TERM-0566}}\par
\noindent\emph{Iztočnik / Source:} Faltungsprozeß / Grundfaltungen\par
\noindent\emph{Rusky / Russian:} процесс свертки / основные свертки\par
\noindent\emph{Ukrajinsky / Ukrainian:} процес згортання / основні згортання\par
\smallskip\noindent\emph{Razlog:} Osnovne svijanja imenuje vyznačene osnovne operacije, a proces svijanja imenuje obšty postup. Razlog iz uvoda raboty 4 prodolžaje terminologiju rabot 1-3 i izrično sohranjaje jej provizorny status. Pridavnik osnovne ne trěba izpustiti: on odlišaje vybrane operacije od ostalyh, ktore možno iz njih sostaviti. Sam zapis ne ustanavjaje novo čislo ili novu klasifikaciju takyh operacij.\par
\noindent\emph{Reason:} Osnovne svijanja names the designated basic operations, whereas proces svijanja names the general process. The introductory Paper 4 record continues the terminology of Papers 1-3 and expressly retains its provisional status. Basic must not be omitted: it distinguishes selected operations from others constructed from them. The record itself establishes neither a new count nor a new classification of those operations.\par
\noindent{\footnotesize Izvorne zapisy / Original records: HIST-g029-012.}\par
\end{minipage}\par
\par\medskip\Needspace{7\baselineskip}
\noindent\begin{minipage}{\linewidth}
\hypertarget{TERM-0567}{}
\noindent\textbf{fiksovany klas}\quad{\footnotesize\texttt{TERM-0567}}\par
\noindent\emph{Iztočnik / Source:} feste Klasse\par
\noindent\emph{Rusky / Russian:} фиксированный класс\par
\noindent\emph{Ukrajinsky / Ukrainian:} фіксований клас\par
\smallskip\noindent\emph{Razlog:} Fiksovany klas ohranjaje zavisnost od izbranogo klasa. Razlog ukazuje na slědujuče sravnjenje so slědom i normoju v oblastjah, kde oni ne zavisjat od takogo izbora. Pridavnik zato nosi matematičnu informaciju, ne samo stilističny naglasok. Fiksovany v toj frazě znači izbrany i zadržany v razsmotrenju; bez drugogo konteksta ne znači invariantny pod vsem transformacijami.\par
\noindent\emph{Reason:} Fiksovany klas preserves dependence on a chosen class. The rationale points to the following contrast with trace and norm in domains, where they are independent of that choice. The adjective therefore carries mathematical information, not merely emphasis. Fixed here means selected and held in the discussion; without further context it does not mean invariant under every transformation.\par
\noindent{\footnotesize Izvorne zapisy / Original records: HIST-g136-001.}\par
\end{minipage}\par
\par\medskip\Needspace{7\baselineskip}
\noindent\begin{minipage}{\linewidth}
\hypertarget{TERM-0569}{}
\noindent\textbf{poslědok}\quad{\footnotesize\texttt{TERM-0569}}\par
\noindent\emph{Iztočnik / Source:} Folgerung\par
\noindent\emph{Rusky / Russian:} следствие; Следствие\par
\noindent\emph{Ukrajinsky / Ukrainian:} Наслідок; наслідок\par
\smallskip\noindent\emph{Razlog:} Poslědok imenuje matematično slědstvo i služi kako zaglavje corollarium. Zapisy iz raboty 2 i raboty 31 davajut prednost kontinuitetu so slovjanskim korenjem za rezultat ili posledovatelnost. Velika bukva v Poslědok jest razlika zaglavja, ne drugy pojem. Razlog ne dava čislenyh težin i ne dokazuje, že vsako obšče slovo za rezultat jest zaměnlivo s formalnym zaglavjem poslědka.\par
\noindent\emph{Reason:} Poslědok denotes a mathematical consequence and serves as a corollary heading. Records from Papers 2 and 31 favour continuity with a Slavic consequence/result family. Capitalization in Poslědok is a heading distinction, not a second concept. The rationale supplies no numerical weights and does not establish that every ordinary word for a result is interchangeable with the formal corollary label.\par
\noindent{\footnotesize Izvorne zapisy / Original records: HIST-g025-010, HIST-g026-011, HIST-g129-001.}\par
\end{minipage}\par
\par\medskip\Needspace{7\baselineskip}
\noindent\begin{minipage}{\linewidth}
\hypertarget{TERM-0570}{}
\noindent\textbf{poslědki}\quad{\footnotesize\texttt{TERM-0570}}\par
\noindent\emph{Iztočnik / Source:} Folgerungen\par
\noindent\emph{Rusky / Russian:} следствия\par
\noindent\emph{Ukrajinsky / Ukrainian:} висновки; наслідки\par
\smallskip\noindent\emph{Razlog:} Osam zapisov sohranjaje poslědki kako zaglavje zaključnyh punktov v rabotě 2. Šest zachovanyh formulacij razloga podkrepljaje kontinuitet, obyčny matematičny registr i kratkost. Ukrajinske sravnjenja imajut i vysnovky, i naslidky: ne sut jedinstvena forma. Najranějši razlog ostavjaje stilističnu prednost k odbornomu razsmotrenju; ne sohranjaje rezultata takogo razsmotrenja. Zato ne možno pripisati ljudskomu recenzentu potvrdženje izbora.\par
\noindent\emph{Reason:} Eight records retain poslědki as the heading for concluding items in Paper 2. Six preserved formulations of the rationale emphasize continuity, ordinary mathematical register and brevity. The Ukrainian comparanda include both vysnovky and naslidky, not one uniform form. The earliest reason leaves stylistic preference for specialist consideration; it contains no result of such consideration. The choice must therefore not be presented as approved by a human reviewer.\par
\noindent{\footnotesize Izvorne zapisy / Original records: HIST-g009-006, HIST-g013-013, HIST-g014-016, HIST-g015-013, HIST-g017-011, HIST-g018-012, HIST-g019-008, HIST-g024-012.}\par
\end{minipage}\par
\par\medskip\Needspace{7\baselineskip}
\noindent\begin{minipage}{\linewidth}
\hypertarget{TERM-0571}{}
\noindent\textbf{forma kriviny}\quad{\footnotesize\texttt{TERM-0571}}\par
\noindent\emph{Iztočnik / Source:} Form der Krümmung / Krümmungsform\par
\noindent\emph{Rusky / Russian:} форма кривизны\par
\noindent\emph{Ukrajinsky / Ukrainian:} форма кривини\par
\smallskip\noindent\emph{Razlog:} Forma kriviny razviva německy kompozit v jasnu konstrukciju s genitivom. Razlog svęzuje termin s konstrukcijeju Riemanna i Lipschitza i naměrno ne zaměnjaje jego nazvom konkretnogo pozdnějšego tenzora. Takovo zadržanje historijnoj širiny pojma ne jest zabranoju sovrěmennogo objasnjenja, ale ono mora ostati odvojeno od avtorskogo nazva. Sama fraza ne dovoljuje avtomatično pripisati vse definicije današnjej diferencialnoj geometrije.\par
\noindent\emph{Reason:} Forma kriviny expands the German compound into a clear genitive construction. The rationale connects the term with the Riemann-Lipschitz construction and deliberately avoids replacing it with the name of a particular later tensor. Retaining that historical breadth does not preclude modern explanation, but such explanation must remain distinct from the author's designation. The phrase alone does not license importing all present-day differential-geometric definitions.\par
\noindent{\footnotesize Izvorne zapisy / Original records: HIST-g084-008.}\par
\end{minipage}\par
\par\medskip\Needspace{7\baselineskip}
\noindent\begin{minipage}{\linewidth}
\hypertarget{TERM-0572}{}
\noindent\textbf{formalne linearne formy}\quad{\footnotesize\texttt{TERM-0572}}\par
\noindent\emph{Iztočnik / Source:} formale Linearformen\par
\noindent\emph{Rusky / Russian:} формальные линейные формы\par
\noindent\emph{Ukrajinsky / Ukrainian:} формальні лінійні форми\par
\smallskip\noindent\emph{Razlog:} Formalne linearne formy ohranjaje i formalnost, i linearnost objektov. Zapis objašnjava doslovnu blizkost potreboju razlikovati linearne transformacije od matric v primětkě pripisanoj van der Waerdenu. To ne jest razlog za zanemarenje rozdela medžu objektom i jego predstavjenjem v izbranoj bazě. Aktualna konstrukcija formalnyh form i proměna bazy sut sravnjene niže; to novo čitanje ne dokazuje samo po sobě historiju navedenej korekcijnoj primětky.\par
\noindent\emph{Reason:} Formalne linearne formy preserves both formality and linearity. The record explains its close wording by the need to distinguish linear transformations from matrices in a correction note attributed to van der Waerden. This is no reason to ignore the distinction between an object and its representation in a chosen basis. The current construction of formal forms and change of basis are compared below; that fresh reading does not by itself establish the history of the reported correction note.\par
\noindent{\footnotesize Izvorne zapisy / Original records: HIST-g193-004.}\par
\smallskip\noindent\emph{Sravnjenje s tekstom:} Aktualna konstrukcija modula predstavjenja bere formalne linearne formy $y=\sum_i x_i\alpha_i$ i izrično svęzuje matricu s izbranoju bazoju. Pri drugoj bazě, $y=xP$, matrica $B$ prěhodi v $P^{-1}BP$. To čini vidimym rozděl medžu formalnymi objektami, linearnoju akcijeju i jej matričnym predstavjenjem. Tute četyri pasaže podkrepljajut terminologično razlikovanje; one same ne sut novym sravnjenjem redakcijnoj primětky pripisanoj van der Waerdenu s pečatnym iztočnikom.\par
\noindent\emph{Textual comparison:} The current representation-module construction takes formal linear forms $y=\sum_i x_i\alpha_i$ and expressly ties a matrix to a chosen basis. With another basis, $y=xP$, matrix $B$ becomes $P^{-1}BP$. This makes visible the distinction between the formal objects, the linear action and its matrix representation. These four passages support the terminological distinction; they are not themselves a fresh collation of the editorial note attributed to van der Waerden against the printed source.\par
\noindent{\footnotesize TEXT-COMP-21; pasaže / passages: 21378, 21380, 21407, 21415; \texttt{CORPUS\_FORM\_COMPARISONS.json}.}\par
\end{minipage}\par
\par\medskip\Needspace{7\baselineskip}
\noindent\begin{minipage}{\linewidth}
\hypertarget{TERM-0573}{}
\noindent\textbf{formalny variacionny račun}\quad{\footnotesize\texttt{TERM-0573}}\par
\noindent\emph{Iztočnik / Source:} formale Variationsrechnung\par
\noindent\emph{Rusky / Russian:} формальное вариационное исчисление\par
\noindent\emph{Ukrajinsky / Ukrainian:} формальне варіаційне числення\par
\smallskip\noindent\emph{Razlog:} Formalny variacionny račun izbira račun prěd internacionalnym kalkulus. Razlog jest svez s ranějšej upotrěboj imenja račun za matematične postupy i s ruskym i ukrajinskim nazvami variacijnogo računa. Formalny ostavaje sučstvenny kvalifikator v toj frazě, ne sinonim neformalnogo ili netočnogo razsuždenja. Zapis dava kvalitativnu prednost domašnjemu korenju, ne čislenu uspěšnost v sravnjenju s kalkulus.\par
\noindent\emph{Reason:} Formalny variacionny račun chooses račun over international kalkulus. The rationale connects it with earlier use of račun for mathematical operations and with the Russian and Ukrainian names for variational calculus. Formal remains a substantive qualifier here, not a synonym for informal or imprecise reasoning. The record gives a qualitative preference for the Slavic-root term, not a numerical performance comparison with kalkulus.\par
\noindent{\footnotesize Izvorne zapisy / Original records: HIST-g084-001.}\par
\end{minipage}\par
\par\medskip\Needspace{7\baselineskip}
\noindent\begin{minipage}{\linewidth}
\hypertarget{TERM-0574}{}
\noindent\textbf{formy 11. poręda}\quad{\footnotesize\texttt{TERM-0574}}\par
\noindent\emph{Iztočnik / Source:} Formen 11. Ordnung\par
\noindent\emph{Rusky / Russian:} формы 11-го порядка\par
\noindent\emph{Ukrajinsky / Ukrainian:} форми 11-го порядку\par
\smallskip\noindent\emph{Razlog:} Formy 11. poręda jest kratka ordinalno-genitivna forma zaglavja, prodolžajuča konvenciju glavy IV raboty 2. Točka po čislu označaje ordinalnost: govorimo o jedinnadsetom poredu, ne o količestvě form. Zapis sohranjaje poręda dlja kontinuiteta i ješče ne predstavlja njegovu grafiju kako konečnu normativnu odluku.\par
\noindent\emph{Reason:} Formy 11. poręda is the compact ordinal-genitive heading continued in Chapter IV of Paper 2. The period marks ordinal use: the heading concerns eleventh order, not eleven forms. The record retains poręda for continuity and does not yet present its spelling as a final normative decision.\par
\noindent{\footnotesize Izvorne zapisy / Original records: HIST-g025-001.}\par
\end{minipage}\par
\par\medskip\Needspace{7\baselineskip}
\noindent\begin{minipage}{\linewidth}
\hypertarget{TERM-0575}{}
\noindent\textbf{formy 11., 12., 13. i 15. poręda}\quad{\footnotesize\texttt{TERM-0575}}\par
\noindent\emph{Iztočnik / Source:} Formen 11., 12., 13., 15. Ordnung\par
\noindent\emph{Rusky / Russian:} формы 11-го, 12-го, 13-го и 15-го порядков\par
\noindent\emph{Ukrajinsky / Ukrainian:} форми 11-го, 12-го, 13-го та 15-го порядків\par
\smallskip\noindent\emph{Razlog:} Spis 11., 12., 13. i 15. mora ostati točny: otsutnost 14 ne jest propust, ktory slovnik može avtomatično zapolniti. Rusko i ukrajinsko sravnjenje upotrěbjajut množinsky genitiv za několiko poredov. Medžuslovjanski zapis zadržuje poręda i izrično ostavjaje sintaksu čislovogo spisa k daljšej prověrkě. Kontinuitet zaglavij zato ne dokazuje, že zajedny jedinstvenny član pri vseh čislah jest jedina ili najlěpša gramatika.\par
\noindent\emph{Reason:} The list 11, 12, 13 and 15 must remain exact: absence of 14 is not a gap the dictionary may automatically fill. The Russian and Ukrainian comparanda use a plural genitive for several orders. The Interslavic record retains poręda and expressly leaves the syntax of the numeral list for further checking. Continuity among headings therefore does not establish that a shared singular expression is the only or best grammatical treatment.\par
\noindent{\footnotesize Izvorne zapisy / Original records: HIST-g017-001.}\par
\end{minipage}\par
\par\medskip\Needspace{7\baselineskip}
\noindent\begin{minipage}{\linewidth}
\hypertarget{TERM-0576}{}
\noindent\textbf{formy 12., 13., 14., 15. poręda}\quad{\footnotesize\texttt{TERM-0576}}\par
\noindent\emph{Iztočnik / Source:} Formen 12., 13., 14., 15. Ordnung\par
\noindent\emph{Rusky / Russian:} формы 12-го, 13-го, 14-го, 15-го порядков\par
\noindent\emph{Ukrajinsky / Ukrainian:} форми 12-го, 13-го, 14-го, 15-го порядків\par
\smallskip\noindent\emph{Razlog:} Formy 12., 13., 14., 15. poręda predstavlja drugy točny spis než TERM-0575. Razlog opet sohranjaje jedinstvennu konvenciju zaglavij pri tom, že ruske i ukrajinske formy imajut množinsky genitiv. Pytanje jednine ili množiny po spisku ordinalov ostavaje otvoreno v historijnom zapisě. Čitatelj mora viděti i spis, i tu gramatičnu nejistost, a ne samo obšče formy vysokogo poreda.\par
\noindent\emph{Reason:} Formy 12., 13., 14., 15. poręda represents a different exact list from TERM-0575. Again the reason retains a uniform heading convention while the Russian and Ukrainian comparanda use a plural genitive. Singular versus plural after a list of ordinals remains unresolved in the historical record. Readers need both the list and that grammatical uncertainty, not merely forms of high order.\par
\noindent{\footnotesize Izvorne zapisy / Original records: HIST-g026-001.}\par
\end{minipage}\par
\par\medskip\Needspace{7\baselineskip}
\noindent\begin{minipage}{\linewidth}
\hypertarget{TERM-0577}{}
\noindent\textbf{formy 7. poręda (sistem $s_{III}$)}\quad{\footnotesize\texttt{TERM-0577}}\par
\noindent\emph{Iztočnik / Source:} Formen 7. Ordnung (System $s_{III}$)\par
\noindent\emph{Rusky / Russian:} формы 7-го порядка (система $s_{III}$)\par
\noindent\emph{Ukrajinsky / Ukrainian:} форми 7-го порядку (система $s_{III}$)\par
\smallskip\noindent\emph{Razlog:} Sedmy pored i oznaka sistema $s_{III}$ nosut različnu informaciju. Fraza prodolžaje konvenciju zaglavij glavy III, sohranjajuči historijno poręda kako provizornu formu. Rimsko III i indeks pri $s$ identifikujut sistem; ne sut ordinalnym slovom, ktoro trěba prěvesti ili prěpisati v kirilicu.\par
\noindent\emph{Reason:} Seventh order and the system designation $s_{III}$ carry different information. The phrase continues the Chapter III heading convention while retaining historical poręda as provisional. Roman III and the subscript attached to $s$ identify the system; they are not an ordinal word to translate or transliterate into Cyrillic.\par
\noindent{\footnotesize Izvorne zapisy / Original records: HIST-g021-001.}\par
\end{minipage}\par
\par\medskip\Needspace{7\baselineskip}
\noindent\begin{minipage}{\linewidth}
\hypertarget{TERM-0578}{}
\noindent\textbf{formy 8. poręda}\quad{\footnotesize\texttt{TERM-0578}}\par
\noindent\emph{Iztočnik / Source:} Formen 8. Ordnung\par
\noindent\emph{Rusky / Russian:} формы 8-го порядка\par
\noindent\emph{Ukrajinsky / Ukrainian:} форми 8-го порядку\par
\smallskip\noindent\emph{Razlog:} Formy 8. poręda prodolžaje čislovo zaglavje iz oddělov 19-20. Ruska i ukrajinska ordinalno-genitivna konstrukcija služi sravnjenju, a medžuslovjanski izbor zadržuje ranějši nazv poreda dlja kontinuiteta. Cifra ne izměnjaje smisyl osmogo poreda i ne znači količestvo form. Zapis ne navodi posebno sravnjenje čitateljskoj prednosti cifry i polnogo ordinalnogo slova.\par
\noindent\emph{Reason:} Formy 8. poręda continues the numeral-heading practice of sections 19-20. Russian and Ukrainian ordinal-genitive constructions provide comparanda, while the Interslavic choice retains the earlier order term for continuity. The digit preserves eighth order and does not count the forms. The record supplies no separate reader comparison between the digit and the spelled-out ordinal.\par
\noindent{\footnotesize Izvorne zapisy / Original records: HIST-g022-001.}\par
\end{minipage}\par
\par\medskip\Needspace{7\baselineskip}
\noindent\begin{minipage}{\linewidth}
\hypertarget{TERM-0579}{}
\noindent\textbf{formy osmogo poręda}\quad{\footnotesize\texttt{TERM-0579}}\par
\noindent\emph{Iztočnik / Source:} Formen achter Ordnung\par
\noindent\emph{Rusky / Russian:} формы восьмого порядка\par
\noindent\emph{Ukrajinsky / Ukrainian:} форми восьмого порядку\par
\smallskip\noindent\emph{Razlog:} Osmogo poręda piše ordinal slovami i sohranjaje genitivnu konstrukciju iz ranějših oddělov raboty 2. Osmogo jest izbrano kako regularna forma od osmy; razlog takože podkrepljaje kontinuitet poręd/poręda. Cifrovany zapis v TERM-0578 i tuto polno slovo označajut isti čislovy pored, ale prinadležet različnym zachovanym pasajam, ne jednoj slučajno udvojenej slovničnej kartě.\par
\noindent\emph{Reason:} Osmogo poręda spells out the ordinal and retains the genitive construction from earlier sections of Paper 2. Osmogo is selected as the regular form of osmy, alongside continuity of poręd/poręda. The digit form in TERM-0578 and this spelled-out form designate the same order number, but belong to different preserved passages rather than an accidentally duplicated dictionary record.\par
\noindent{\footnotesize Izvorne zapisy / Original records: HIST-g014-001.}\par
\end{minipage}\par
\par\medskip\Needspace{7\baselineskip}
\noindent\begin{minipage}{\linewidth}
\hypertarget{TERM-0580}{}
\noindent\textbf{formy tretjego poręda}\quad{\footnotesize\texttt{TERM-0580}}\par
\noindent\emph{Iztočnik / Source:} Formen dritter Ordnung\par
\noindent\emph{Rusky / Russian:} формы третьего порядка\par
\noindent\emph{Ukrajinsky / Ukrainian:} форми третього порядку\par
\smallskip\noindent\emph{Razlog:} Tretjego poręda sohranjaje historijni pored v koeficientah, ne stepenj v proměnnyh. Zapis takože izrično govori o proměně tretjogo na tretjego. Tuto jest konkretna gramatična alternativa, ne novy matematičny pojem. Pozdnějše sravnjenje s gramatikoju podkrepljaje meku deklinaciju ordinala za tretji, ale ne dokazuje, že ona byla konsultovana pri prvom izborě. Točna ranějša grafija i aktualna forma ostavajut razlikovane.\par
\noindent\emph{Reason:} Tretjego poręda retains the historical order in the coefficients, not degree in the variables. The record also explicitly reports replacing tretjogo with tretjego. This is a concrete grammatical alternative, not a new mathematical concept. The later grammar comparison supports soft declension of the ordinal for third but does not establish that the grammar was consulted when the first choice was made. The earlier spelling and current form remain distinguished.\par
\noindent{\footnotesize Izvorne zapisy / Original records: HIST-g009-001.}\par
\smallskip\noindent\emph{Sravnjenje s tekstom:} Podnožna primětka v uvodě raboty 2 izrično razlikuje pored v koeficientah od stepena v proměnnyh; odděl 7 ponovno uporedkuje formy po poredu v koeficientah. Aktualne zaglavja oddělov 8–10 imajut tretjego, četvrtogo i petogo poreda. Pri novom sravnjenju gramatika čislovnikov dava meku deklinaciju za trětji, tvrdu za ostale ordinale i izrično navodi pety; gramatika pridavnikov objasnjaje alternaciju o/e po mekoj osnově. Tuto podkrepljaje morfologičny kontrast -ego/-ogo, ne ustanavjaje, že obě gramatiky byly raněje konsultovane. Ne potvrđaje takože vse grafijne detalje historijnogo poręda ili izključnost njegovej leksičnoj prednosti. Točne gramatične pasaže sut zachovane pod GRAM-01.\par
\noindent\emph{Textual comparison:} The introductory footnote to Paper 2 expressly distinguishes order in the coefficients from degree in the variables; section 7 again orders forms by coefficient order. Current headings in sections 8–10 have tretjego, četvrtogo and petogo poreda. In a fresh comparison, the numeral grammar assigns soft declension to trětji and hard declension to the other ordinals, expressly listing pety; the adjective grammar explains o/e alternation after a soft stem. This supports the morphological -ego/-ogo contrast, not earlier consultation of either grammar. It also does not validate every spelling detail of historical poręda or an exclusive lexical preference for it. The exact grammar passages are retained under GRAM-01.\par
\noindent{\footnotesize TEXT-COMP-20; pasaže / passages: 175, 1066, 1178, 1215, 1282; \texttt{CORPUS\_FORM\_COMPARISONS.json}.}\par
\noindent{\footnotesize\href{https://interslavic.fun/learn/grammar/numerals/}{Čislovniki / Numerals}: Ordinalne čisla / Ordinal numbers; GRAM-01.}\par
\noindent{\footnotesize\href{https://interslavic.fun/learn/grammar/adjectives/}{Pridavniki / Adjectives}: Sklanjanje / Inflection; GRAM-01.}\par
\end{minipage}\par
\par\medskip\Needspace{7\baselineskip}
\noindent\begin{minipage}{\linewidth}
\hypertarget{TERM-0581}{}
\noindent\textbf{formy petogo poręda}\quad{\footnotesize\texttt{TERM-0581}}\par
\noindent\emph{Iztočnik / Source:} Formen fünfter Ordnung\par
\noindent\emph{Rusky / Russian:} формы пятого порядка\par
\noindent\emph{Ukrajinsky / Ukrainian:} форми п'ятого порядку\par
\smallskip\noindent\emph{Razlog:} Petogo poręda prodolžaje zaglavja o poredu v koeficientah. Zapis vybira petogo od pety i izrično odbacaje dogadanu nazalizaciju; pri tom razlikuje tuto frazu od staršego petogo reda. Razlog tvrdi, že forma pety jest podkrepljena gramatikoju čislovnikov. Novo čitanje sohranjenoj gramatiky potvrđaje, že ona navodi pety i tvrdu deklinaciju, ale ne ustanavjaje, kto i kogda ju raněje konsultoval. Izbor poreda ostavaje odvojenym leksičnym pytanjem.\par
\noindent\emph{Reason:} Petogo poręda continues the headings about order in the coefficients. The record chooses petogo from pety and explicitly rejects a guessed nasalized form, while distinguishing this phrase from earlier petogo reda. The rationale claims support from numeral guidance. Fresh examination of the preserved grammar confirms that it lists pety and assigns hard declension, but does not establish who consulted it earlier or when. Selection of the order noun remains a separate lexical question.\par
\noindent{\footnotesize Izvorne zapisy / Original records: HIST-g011-001.}\par
\smallskip\noindent\emph{Sravnjenje s tekstom:} Podnožna primětka v uvodě raboty 2 izrično razlikuje pored v koeficientah od stepena v proměnnyh; odděl 7 ponovno uporedkuje formy po poredu v koeficientah. Aktualne zaglavja oddělov 8–10 imajut tretjego, četvrtogo i petogo poreda. Pri novom sravnjenju gramatika čislovnikov dava meku deklinaciju za trětji, tvrdu za ostale ordinale i izrično navodi pety; gramatika pridavnikov objasnjaje alternaciju o/e po mekoj osnově. Tuto podkrepljaje morfologičny kontrast -ego/-ogo, ne ustanavjaje, že obě gramatiky byly raněje konsultovane. Ne potvrđaje takože vse grafijne detalje historijnogo poręda ili izključnost njegovej leksičnoj prednosti. Točne gramatične pasaže sut zachovane pod GRAM-01.\par
\noindent\emph{Textual comparison:} The introductory footnote to Paper 2 expressly distinguishes order in the coefficients from degree in the variables; section 7 again orders forms by coefficient order. Current headings in sections 8–10 have tretjego, četvrtogo and petogo poreda. In a fresh comparison, the numeral grammar assigns soft declension to trětji and hard declension to the other ordinals, expressly listing pety; the adjective grammar explains o/e alternation after a soft stem. This supports the morphological -ego/-ogo contrast, not earlier consultation of either grammar. It also does not validate every spelling detail of historical poręda or an exclusive lexical preference for it. The exact grammar passages are retained under GRAM-01.\par
\noindent{\footnotesize TEXT-COMP-20; pasaže / passages: 175, 1066, 1178, 1215, 1282; \texttt{CORPUS\_FORM\_COMPARISONS.json}.}\par
\noindent{\footnotesize\href{https://interslavic.fun/learn/grammar/numerals/}{Čislovniki / Numerals}: Ordinalne čisla / Ordinal numbers; GRAM-01.}\par
\noindent{\footnotesize\href{https://interslavic.fun/learn/grammar/adjectives/}{Pridavniki / Adjectives}: Sklanjanje / Inflection; GRAM-01.}\par
\end{minipage}\par
\par\medskip\Needspace{7\baselineskip}
\noindent\begin{minipage}{\linewidth}
\hypertarget{TERM-0582}{}
\noindent\textbf{formy s vyšimi simbolami}\quad{\footnotesize\texttt{TERM-0582}}\par
\noindent\emph{Iztočnik / Source:} Formen mit höheren Symbolen\par
\noindent\emph{Rusky / Russian:} формы с более высокими символами\par
\noindent\emph{Ukrajinsky / Ukrainian:} форми з вищими символами\par
\smallskip\noindent\emph{Razlog:} S vyšimi simbolami razklada německu konstrukciju bez novogo složenogo termina. Zapis podkrepljaje komparativ i instrumental, ktore ostavajut čitateljne pri simboličnyh formulah. Vyši ne znači tu většy font ili vyše položenje na stranicě: odnos medžu simbolami trěba razuměti iz konstrukcije form. Razlog ne sohranjaje osobnu tablicu konkurentnyh pridavnikov.\par
\noindent\emph{Reason:} S vyšimi simbolami renders the German construction without coining a new compound. The rationale emphasizes a comparative adjective and instrumental construction readable alongside symbolic formulas. Higher does not here mean larger type or a higher position on the page: the relation between symbols must come from the construction of forms. No separate comparison table of competing adjectives survives in the reason.\par
\noindent{\footnotesize Izvorne zapisy / Original records: HIST-g012-007.}\par
\end{minipage}\par
\par\medskip\Needspace{7\baselineskip}
\noindent\begin{minipage}{\linewidth}
\hypertarget{TERM-0583}{}
\noindent\textbf{formy 9. poręda}\quad{\footnotesize\texttt{TERM-0583}}\par
\noindent\emph{Iztočnik / Source:} Formen neunter Ordnung\par
\noindent\emph{Rusky / Russian:} формы 9-го порядка\par
\noindent\emph{Ukrajinsky / Ukrainian:} форми 9-го порядку\par
\smallskip\noindent\emph{Razlog:} Formy 9. poręda zadržuje cifrovu ordinalno-genitivnu konstrukciju ranějših oddělov. Razlog dava prednost kontinuitetu prěd novym konkurentnym doslovnym prěvodom, ale ne imenuje togo konkurenta. Ne možno zato dopisati konkretny odbacany termin ili čislenu ocěnu. Čislo 9 ostavaje klasifikacijeju po poredu, ne čislom navedenyh form.\par
\noindent\emph{Reason:} Formy 9. poręda retains the digit-based ordinal-genitive construction of preceding sections. The reason prefers continuity over a competing calque but does not name that calque. A specific rejected term or numerical score therefore cannot be supplied. The number 9 remains a classification by order, not a count of the listed forms.\par
\noindent{\footnotesize Izvorne zapisy / Original records: HIST-g023-001.}\par
\end{minipage}\par
\par\medskip\Needspace{7\baselineskip}
\noindent\begin{minipage}{\linewidth}
\hypertarget{TERM-0584}{}
\noindent\textbf{formy devętogo poręda}\quad{\footnotesize\texttt{TERM-0584}}\par
\noindent\emph{Iztočnik / Source:} Formen neunter Ordnung\par
\noindent\emph{Rusky / Russian:} формы девятого порядка\par
\noindent\emph{Ukrajinsky / Ukrainian:} форми дев'ятого порядку\par
\smallskip\noindent\emph{Razlog:} Devętogo poręda sohranjaje polno ordinalno slovo i nazalizovanu grafiju kak zapisany predmet daljšego razsmotrenja. Razlog obračaje vnimanje na ukrajinski apostrof v dev’jatogo i na rusku genitivnu formu, ale te razlike ne sut samy po sobě opravdanjem medžuslovjanskogo ę. Zapis s cifroju 9. v TERM-0583 označaje isti čislovy pored; razlika form ostavaje vidima, bez tvrdženja, že vsaka grafija už jest potvrđena.\par
\noindent\emph{Reason:} Devętogo poręda preserves the spelled-out ordinal and its nasalized spelling as a recorded question for further consideration. The reason notes the Ukrainian apostrophe in dev’jatogo and the Russian genitive form, but those differences do not by themselves justify Interslavic ę. The digit form 9. in TERM-0583 expresses the same order number; the variants remain visible without claiming that every spelling has been validated.\par
\noindent{\footnotesize Izvorne zapisy / Original records: HIST-g015-001.}\par
\end{minipage}\par
\par\medskip\Needspace{7\baselineskip}
\noindent\begin{minipage}{\linewidth}
\hypertarget{TERM-0585}{}
\noindent\textbf{formy šestogo poręda}\quad{\footnotesize\texttt{TERM-0585}}\par
\noindent\emph{Iztočnik / Source:} Formen sechster Ordnung\par
\noindent\emph{Rusky / Russian:} формы шестого порядка\par
\noindent\emph{Ukrajinsky / Ukrainian:} форми шостого порядку\par
\smallskip\noindent\emph{Razlog:} Šestogo poręda nosi regularny ordinalny pridavnik v genitivě i prodolžaje zaglavja o poredu v koeficientah. Zapis sravnjaje rusku i ukrajinsku konstrukciju i ostavjaje cělu semju ordinalov i registra k odbornomu razsmotrenju. Regularnost pridavnika ne rěšaje avtomatično drugogo pytanja, koji medžuslovjanski nazv za Ordnung jest najpriměrnějši.\par
\noindent\emph{Reason:} Šestogo poręda uses the regular ordinal adjective in the genitive and continues the coefficient-order headings. The record compares Russian and Ukrainian constructions and leaves the wider ordinal/register family for specialist consideration. Regularity of the adjective does not automatically decide the separate question of the most suitable Interslavic noun for Ordnung.\par
\noindent{\footnotesize Izvorne zapisy / Original records: HIST-g012-001.}\par
\end{minipage}\par
\par\medskip\Needspace{7\baselineskip}
\noindent\begin{minipage}{\linewidth}
\hypertarget{TERM-0586}{}
\noindent\textbf{formy sedmogo poręda}\quad{\footnotesize\texttt{TERM-0586}}\par
\noindent\emph{Iztočnik / Source:} Formen siebenter Ordnung\par
\noindent\emph{Rusky / Russian:} формы седьмого порядка\par
\noindent\emph{Ukrajinsky / Ukrainian:} форми сьомого порядку\par
\smallskip\noindent\emph{Razlog:} Sedmogo poręda prodolžaje konvenciju iz oddělov 7-11 raboty 2: polno ordinalno slovo s genitivnym nazvom poreda. Ruska i ukrajinska sravnjenja pokazyvajut različne nacionalne formy ordinala, zato medžuslovjanski oblik ne jest prosto kopijeju jednoj iz njih. Razlog sohranjaje obšče petogo/šestogo/sedmogo sklanjanje i stalno imenuvanje matematičnoj veličiny.\par
\noindent\emph{Reason:} Sedmogo poręda continues the convention of sections 7-11 of Paper 2: a spelled-out ordinal with the genitive order noun. Russian and Ukrainian comparanda have different national ordinal forms, so the Interslavic shape is not simply a copy of either. The reason preserves the shared petogo/šestogo/sedmogo inflectional pattern and a stable name for the mathematical quantity.\par
\noindent{\footnotesize Izvorne zapisy / Original records: HIST-g013-001.}\par
\end{minipage}\par
\par\medskip\Needspace{7\baselineskip}
\noindent\begin{minipage}{\linewidth}
\hypertarget{TERM-0587}{}
\noindent\textbf{formy četvrtogo poręda}\quad{\footnotesize\texttt{TERM-0587}}\par
\noindent\emph{Iztočnik / Source:} Formen vierter Ordnung\par
\noindent\emph{Rusky / Russian:} формы четвертого порядка\par
\noindent\emph{Ukrajinsky / Ukrainian:} форми четвертого порядку\par
\smallskip\noindent\emph{Razlog:} Četvrtogo poręda jest tvrda pridavnička genitivna konstrukcija. Zapis izrično sravnjaje četvrtogo s mekim tretjego iz prědhodnogo odděla; razlika -ogo/-ego zato jest opravdana morfologijeju, ne nekonsistentnostju prěvoda. Pored ostavaje poredom v koeficientah. Novo sravnjenje s gramatikoju i aktualnymi zaglavjami podkrepljaje tu razliku, ne pretvarjajuči novy dokaz v staru konsultaciju.\par
\noindent\emph{Reason:} Četvrtogo poręda uses the hard-adjective genitive. The record explicitly contrasts četvrtogo with soft tretjego in the preceding section; the -ogo/-ego difference is therefore motivated by morphology, not inconsistency. Order remains order in the coefficients. Fresh comparison with the grammar and current headings supports that distinction without turning new evidence into an earlier consultation.\par
\noindent{\footnotesize Izvorne zapisy / Original records: HIST-g010-001.}\par
\smallskip\noindent\emph{Sravnjenje s tekstom:} Podnožna primětka v uvodě raboty 2 izrično razlikuje pored v koeficientah od stepena v proměnnyh; odděl 7 ponovno uporedkuje formy po poredu v koeficientah. Aktualne zaglavja oddělov 8–10 imajut tretjego, četvrtogo i petogo poreda. Pri novom sravnjenju gramatika čislovnikov dava meku deklinaciju za trětji, tvrdu za ostale ordinale i izrično navodi pety; gramatika pridavnikov objasnjaje alternaciju o/e po mekoj osnově. Tuto podkrepljaje morfologičny kontrast -ego/-ogo, ne ustanavjaje, že obě gramatiky byly raněje konsultovane. Ne potvrđaje takože vse grafijne detalje historijnogo poręda ili izključnost njegovej leksičnoj prednosti. Točne gramatične pasaže sut zachovane pod GRAM-01.\par
\noindent\emph{Textual comparison:} The introductory footnote to Paper 2 expressly distinguishes order in the coefficients from degree in the variables; section 7 again orders forms by coefficient order. Current headings in sections 8–10 have tretjego, četvrtogo and petogo poreda. In a fresh comparison, the numeral grammar assigns soft declension to trětji and hard declension to the other ordinals, expressly listing pety; the adjective grammar explains o/e alternation after a soft stem. This supports the morphological -ego/-ogo contrast, not earlier consultation of either grammar. It also does not validate every spelling detail of historical poręda or an exclusive lexical preference for it. The exact grammar passages are retained under GRAM-01.\par
\noindent{\footnotesize TEXT-COMP-20; pasaže / passages: 175, 1066, 1178, 1215, 1282; \texttt{CORPUS\_FORM\_COMPARISONS.json}.}\par
\noindent{\footnotesize\href{https://interslavic.fun/learn/grammar/numerals/}{Čislovniki / Numerals}: Ordinalne čisla / Ordinal numbers; GRAM-01.}\par
\noindent{\footnotesize\href{https://interslavic.fun/learn/grammar/adjectives/}{Pridavniki / Adjectives}: Sklanjanje / Inflection; GRAM-01.}\par
\end{minipage}\par
\par\medskip\Needspace{7\baselineskip}
\noindent\begin{minipage}{\linewidth}
\hypertarget{TERM-0588}{}
\noindent\textbf{formy v $n$ prěměnnyh}\quad{\footnotesize\texttt{TERM-0588}}\par
\noindent\emph{Iztočnik / Source:} Formen von $n$ Variabeln\par
\noindent\emph{Rusky / Russian:} формы от $n$ переменных\par
\noindent\emph{Ukrajinsky / Ukrainian:} форми від $n$ змінних\par
\smallskip\noindent\emph{Razlog:} Formy v $n$ prěměnnyh prodolžaje predmetny nazv iz raboty 3 v uvodě raboty 4. Predlog v čini frazu prirodnoju v medžuslovjanskom, bez doslovnogo kopirovanja německogo von ili nacionalnyh predlogov v sravnjenjah. $n$ označaje čislo proměnnyh, ne pored ili stepenj formy. Razlog ne podkrepljaje izpuščenje nijednoj proměnnoj ili suženje k binarnomu slučaju.\par
\noindent\emph{Reason:} Formy v $n$ prěměnnyh carries the subject designation from Paper 3 into Paper 4's introduction. The preposition v makes the phrase natural in Interslavic without literally copying German von or the national prepositions in the comparanda. $n$ counts variables, not the order or degree of a form. The rationale does not license omitting variables or restricting the statement to the binary case.\par
\noindent{\footnotesize Izvorne zapisy / Original records: HIST-g029-001.}\par
\end{minipage}\par
\par\medskip\Needspace{7\baselineskip}
\noindent\begin{minipage}{\linewidth}
\hypertarget{TERM-0589}{}
\noindent\textbf{formy 10. poręda}\quad{\footnotesize\texttt{TERM-0589}}\par
\noindent\emph{Iztočnik / Source:} Formen zehnter Ordnung\par
\noindent\emph{Rusky / Russian:} формы 10-го порядка\par
\noindent\emph{Ukrajinsky / Ukrainian:} форми 10-го порядку\par
\smallskip\noindent\emph{Razlog:} Formy 10. poręda zadržuje kratku cifrovu konstrukciju glavy IV. Motiv jest izběgnuti novogo konkurentnogo imena za pored pri povtorjenom tipu zaglavja. Konkretny drugy nazv ne jest navedeny. Cifrovo pisanje može byti sravnjeno s desętogo poręda v TERM-0590, ale razlikovanje dvou zapisov ne ustanavjaje novy matematičny smisyl.\par
\noindent\emph{Reason:} Formy 10. poręda retains Chapter IV's compact digit-based construction. The motive is to avoid introducing a competing order term into a repeated heading pattern. No particular alternative noun is named. The notation can be compared with spelled-out desętogo poręda in TERM-0590, but the difference between the records establishes no new mathematical sense.\par
\noindent{\footnotesize Izvorne zapisy / Original records: HIST-g024-001.}\par
\end{minipage}\par
\par\medskip\Needspace{7\baselineskip}
\noindent\begin{minipage}{\linewidth}
\hypertarget{TERM-0590}{}
\noindent\textbf{formy desętogo poręda}\quad{\footnotesize\texttt{TERM-0590}}\par
\noindent\emph{Iztočnik / Source:} Formen zehnter Ordnung\par
\noindent\emph{Rusky / Russian:} формы десятого порядка\par
\noindent\emph{Ukrajinsky / Ukrainian:} форми десятого порядку\par
\smallskip\noindent\emph{Razlog:} Desętogo poręda sohranjaje ordinal slovami, genitivnu konstrukciju i nazalizovanu grafiju vidimu dlja daljšego razsmotrenja. Rusko i ukrajinsko sravnjenje podkrepljajut gramatičny tip frazy, ne avtomatično točnu medžuslovjansku ortografiju. Tuto jest zapisana forma iz odděla 15 raboty 2; njej ne trěba ni skryti, ni predstavjati kako jedinu aktualnu rekomendaciju.\par
\noindent\emph{Reason:} Desętogo poręda retains the spelled-out ordinal, genitive construction and nasalized spelling for further consideration. Russian and Ukrainian comparanda support the grammatical type of phrase, not automatically the exact Interslavic orthography. This is the recorded form from section 15 of Paper 2; it must be neither concealed nor presented as the sole current recommendation.\par
\noindent{\footnotesize Izvorne zapisy / Original records: HIST-g016-001.}\par
\end{minipage}\par
\par\medskip\Needspace{7\baselineskip}
\noindent\begin{minipage}{\linewidth}
\hypertarget{TERM-0591}{}
\noindent\textbf{ręd form}\quad{\footnotesize\texttt{TERM-0591}}\par
\noindent\emph{Iztočnik / Source:} Formenreihe\par
\noindent\emph{Rusky / Russian:} ряд форм\par
\noindent\emph{Ukrajinsky / Ukrainian:} ряд форм\par
\smallskip\noindent\emph{Razlog:} Četrnadset zapisov davaje ręd form za historijni objekt Formenreihe. Glavni razlog ne jest samo podobnost ruskogo i ukrajinskogo rjad: Noether oprěděljaje generovanu semju form i razměščaje ju v premougolnoj shemě. Kratko ime sohranjaje svez medžu tabelami, redukcijami i pozdnějšimi konstrukcijami. Razlogy varujut prěd čitanjem kako analitičnogo stepenovogo reda, proste množiny ili objekta s dodatnoju modulnoju strukturoju. Něktore zapisy govoret o posledovatelnosti, druge o redu ili semji; ta razlika pokazuje potrěbu definicije, ne četiri slobodne sinonimy. Zapisani razlozi zadržujut semantičnu nejistost dlja odbornogo razsmotrenja, a ne dokazuju izměrjenu medžuslovjansku razumělivost.\par
\noindent\emph{Reason:} Fourteen records give ręd form for the historical object Formenreihe. The principal reason is not merely resemblance to Russian and Ukrainian rjad: Noether defines a generated family of forms and arranges it in a rectangular scheme. The compact name links tables, reductions and later constructions. The rationales warn against reading it as an analytic power series, a mere set or an object with additional module structure. Some say sequence, others row or family; that variation demonstrates the need for the definition, not four freely interchangeable synonyms. The recorded reasons retain semantic ambiguity for specialist consideration rather than proving measured cross-Slavic intelligibility.\par
\noindent{\footnotesize Izvorne zapisy / Original records: HIST-g001-009, HIST-g003-002, HIST-g011-004, HIST-g013-005, HIST-g014-008, HIST-g015-008, HIST-g021-006, HIST-g022-004, HIST-g023-005, HIST-g024-006, HIST-g025-007, HIST-g029-014, HIST-g036-008, HIST-g038-001.}\par
\smallskip\noindent\emph{Sravnjenje s tekstom:} Aktualny odděl 1 raboty 2 naziva osnovny postup kontrakcijeju. Red form jest početna forma zajedno so vsimi formami dobytymi kontrakcijami samej so soboju; vyši formy razlikujut se po kontrakcijah. Pri veče redah simbolov tekst trěbuje izbrati specialne kontrakcije, da red byl jednoznačno normovany, i potom pokazuje premougolnu shemu. Tuto podkrepljaje izbor kratkogo imena za oprěděljenu generovanu strukturu, ne prosto množinu ili analitičny red. Aktualno red i kontrakcija razlikujut se od historijnyh ręd i svijanje; sravnjenje ne dokazuje cělokorpusno uskladnjenje ni oficijalny status kandidatov.\par
\noindent\emph{Textual comparison:} The current section 1 of Paper 2 calls the basic process contraction. A row of forms comprises an initial form and all forms obtained by self-contractions; higher forms are distinguished by the contractions performed. With more rows of symbols, the text requires selection of particular contractions to normalize the row uniquely, and then displays a rectangular scheme. This supports a compact name for a defined generated structure, not merely a set or an analytic series. Current red and kontrakcija differ from historical ręd and svijanje; the comparison does not establish corpus-wide harmonization or official status for the candidates.\par
\noindent{\footnotesize TEXT-COMP-19; pasaže / passages: 243, 295, 299, 301; \texttt{CORPUS\_FORM\_COMPARISONS.json}.}\par
\end{minipage}\par
\par\medskip\Needspace{7\baselineskip}
\noindent\begin{minipage}{\linewidth}
\hypertarget{TERM-0592}{}
\noindent\textbf{ręd form / rędy form}\quad{\footnotesize\texttt{TERM-0592}}\par
\noindent\emph{Iztočnik / Source:} Formenreihe / Formenreihen\par
\noindent\emph{Rusky / Russian:} ряд форм / ряды форм\par
\noindent\emph{Ukrajinsky / Ukrainian:} ряд форм / ряди форм\par
\smallskip\noindent\emph{Razlog:} Ręd form / rędy form pokazuje jedninu i množinu istogo imena pri konstrukcijah odděla 18 raboty 2. Razlog želi zadržati obščeslovjansku čitateljnost bez suženja k jednoj nacionalnoj terminologiji. On ne dava mnogojezyčny eksperiment ili čislene težiny. Gramatična para ne može zaměniti definiciju reda form; ona govori, kako imenujemo jeden objekt i několiko takyh objektov.\par
\noindent\emph{Reason:} Ręd form / rędy form shows singular and plural of the same noun in section 18 of Paper 2. The rationale aims at broadly Slavic readability without narrowing the terminology to one national variety. It supplies no multilingual experiment or numerical weights. The grammatical pair does not replace the definition of a row of forms; it specifies how one such object and several are named.\par
\noindent{\footnotesize Izvorne zapisy / Original records: HIST-g019-005.}\par
\end{minipage}\par
\par\medskip\Needspace{7\baselineskip}
\noindent\begin{minipage}{\linewidth}
\hypertarget{TERM-0593}{}
\noindent\textbf{normovati ręd form}\quad{\footnotesize\texttt{TERM-0593}}\par
\noindent\emph{Iztočnik / Source:} Formenreihe normieren\par
\noindent\emph{Rusky / Russian:} нормировать ряд форм\par
\noindent\emph{Ukrajinsky / Ukrainian:} нормувати ряд форм\par
\smallskip\noindent\emph{Razlog:} Normovati ręd form prodolžaje definicijnu terminologiju odděla 1 pri razsmotrenju v oddělu 5 raboty 2. Jednoznačnost normovanja jest svezana s izbranymi kontrakcijami i uporadkovanjem, ne s libovoljnoju zaměnoju koeficientov. Sam kratky razlog ne opisuje vseh korakov; zato objasnjenje mora byti čitano s definicijeju. Nije oprěděljeno kako děljenje na euklidovu normu ili avtomatično postavjenje najvyššego koeficienta ravnym jedinici.\par
\noindent\emph{Reason:} Normovati ręd form carries the definitional terminology of section 1 into section 5 of Paper 2. Uniqueness of normalization concerns selected contractions and ordering, not arbitrary coefficient replacement. The brief reason does not describe every step, so the explanation must be read with the definition. It is not defined as division by a Euclidean norm or automatic setting of the leading coefficient to one.\par
\noindent{\footnotesize Izvorne zapisy / Original records: HIST-g006-003.}\par
\smallskip\noindent\emph{Sravnjenje s tekstom:} Aktualny odděl 1 raboty 2 naziva osnovny postup kontrakcijeju. Red form jest početna forma zajedno so vsimi formami dobytymi kontrakcijami samej so soboju; vyši formy razlikujut se po kontrakcijah. Pri veče redah simbolov tekst trěbuje izbrati specialne kontrakcije, da red byl jednoznačno normovany, i potom pokazuje premougolnu shemu. Tuto podkrepljaje izbor kratkogo imena za oprěděljenu generovanu strukturu, ne prosto množinu ili analitičny red. Aktualno red i kontrakcija razlikujut se od historijnyh ręd i svijanje; sravnjenje ne dokazuje cělokorpusno uskladnjenje ni oficijalny status kandidatov.\par
\noindent\emph{Textual comparison:} The current section 1 of Paper 2 calls the basic process contraction. A row of forms comprises an initial form and all forms obtained by self-contractions; higher forms are distinguished by the contractions performed. With more rows of symbols, the text requires selection of particular contractions to normalize the row uniquely, and then displays a rectangular scheme. This supports a compact name for a defined generated structure, not merely a set or an analytic series. Current red and kontrakcija differ from historical ręd and svijanje; the comparison does not establish corpus-wide harmonization or official status for the candidates.\par
\noindent{\footnotesize TEXT-COMP-19; pasaže / passages: 243, 295, 299, 301; \texttt{CORPUS\_FORM\_COMPARISONS.json}.}\par
\end{minipage}\par
\par\medskip\Needspace{7\baselineskip}
\noindent\begin{minipage}{\linewidth}
\hypertarget{TERM-0594}{}
\noindent\textbf{ręd form / rędy form}\quad{\footnotesize\texttt{TERM-0594}}\par
\noindent\emph{Iztočnik / Source:} Formenreihe/Formenreihen\par
\noindent\emph{Rusky / Russian:} ряд форм / ряды форм\par
\noindent\emph{Ukrajinsky / Ukrainian:} ряд форм / ряди форм\par
\smallskip\noindent\emph{Razlog:} Tuto para ręd form / rędy form jest zabilěžena v oddělu 19, odvojeno od zapisa dlja odděla 18. Razlog izrično ohranjaje razliku jednine i množiny za pozdnějše uskladnjenje. Dva povezana zapisy ne sut dva različne matematične termina, ale ni jeden ne trěba odstraniti iz dokazov o upotrěbě. Pravopisno ili leksično uskladnjenje mora sledovati konkretne pojavjenja, ne samo slovničnu zaglavnu formu.\par
\noindent\emph{Reason:} This ręd form / rędy form pair is recorded for section 19, separately from section 18. Its reason explicitly preserves singular and plural for later harmonization. The two linked records are not two different mathematical terms, but neither should disappear from the evidence of use. Spelling or lexical harmonization must follow actual occurrences, not only a dictionary headword.\par
\noindent{\footnotesize Izvorne zapisy / Original records: HIST-g020-004.}\par
\end{minipage}\par
\par\medskip\Needspace{7\baselineskip}
\noindent\begin{minipage}{\linewidth}
\hypertarget{TERM-0595}{}
\noindent\textbf{sistem form}\quad{\footnotesize\texttt{TERM-0595}}\par
\noindent\emph{Iztočnik / Source:} Formensystem\par
\noindent\emph{Rusky / Russian:} система форм\par
\noindent\emph{Ukrajinsky / Ukrainian:} система форм\par
\smallskip\noindent\emph{Razlog:} Sistem form prodolžaje nazv iz raboty 1 i sohranjaje historijni pojem Formensystem v uvodě raboty 2. Zapis izrično ne hoče prědčasno zaměniti jego sovrěmennym generujučim množestvom. Čitatelju ostavaje izložena definicija i konkretne svojstva polnoty, a ne svojstva, ktore by novy nazv mogl tiho dobaviti. To jest granica modernizacije pojmov, ne zabrana objasnjujučih komentarov v odvojenom slovniku.\par
\noindent\emph{Reason:} Sistem form continues the Paper 1 designation and preserves historical Formensystem in Paper 2's introduction. The record expressly avoids prematurely replacing it with a modern generating set. Readers receive the stated definition and particular completeness properties, not properties that a modern label might silently add. This limits conceptual modernization; it does not prohibit explanatory discussion in the separate dictionary.\par
\noindent{\footnotesize Izvorne zapisy / Original records: HIST-g002-001.}\par
\end{minipage}\par
\par\medskip\Needspace{7\baselineskip}
\noindent\begin{minipage}{\linewidth}
\hypertarget{TERM-0596}{}
\noindent\textbf{sistem form; [pri objasnjenju / when explanatory] sistem porodžujųćih form}\quad{\footnotesize\texttt{TERM-0596}}\par
\noindent\emph{Iztočnik / Source:} Formensystem\par
\noindent\emph{Rusky / Russian:} система форм; [pri objasnjenju / when explanatory] система порождающих форм\par
\noindent\emph{Ukrajinsky / Ukrainian:} система форм; [pri objasnjenju / when explanatory] система породжувальних форм\par
\smallskip\noindent\emph{Razlog:} Glavni kandidat jest sistem form; sistem porodžujųćih form jest zapisana objasnjujuča fraza, ne avtomatična zaměna vsakego Formensystem. Razlog odbacaje slepy doslovny kompozit formosistema i upominaje ukrajinsko породжуючі елементи za generatory. On ne imenuje konkretnu stranicu citovanogo iztočnika v tutom zapisě. Generujuču funkciju zato trěba ustanoviti iz konteksta, a ne iz navodne universalnoj sinonimije; grafija i sklanjanje dugego medžuslovjanskogo kandidata ostavajut predmetom daljšej prověrky.\par
\noindent\emph{Reason:} The main candidate is sistem form; sistem porodžujųćih form is a recorded explanatory phrase, not an automatic substitute for every Formensystem. The rationale rejects a blind compound calque formosistema and mentions Ukrainian породжуючі елементи for generators. This record does not identify an exact page in that reference. Generating function therefore has to be established contextually, not through supposed universal synonymy; the longer Interslavic candidate's spelling and inflection remain for further examination.\par
\noindent{\footnotesize Izvorne zapisy / Original records: HIST-g001-002.}\par
\end{minipage}\par
\par\medskip\Needspace{7\baselineskip}
\noindent\begin{minipage}{\linewidth}
\hypertarget{TERM-0597}{}
\noindent\textbf{sistema form}\quad{\footnotesize\texttt{TERM-0597}}\par
\noindent\emph{Iztočnik / Source:} Formensystem\par
\noindent\emph{Rusky / Russian:} система форм\par
\noindent\emph{Ukrajinsky / Ukrainian:} система форм\par
\smallskip\noindent\emph{Razlog:} Sistema form jest druga zapisana morfologična forma pri bibliografičnoj upotrěbě. Razlog svęzuje ju s ranějšej terminologijeju invariantnoj teorije i izběganiem tęžkyh kompozitov, ale ne argumentuje osobno za konečno -a v sravnjenju so sistem. Zato slovnik mora pokazati sistem form i sistema form odvojeno, ne pripisati zapisam nepostoječu jedinstvennu politiku roda imenja.\par
\noindent\emph{Reason:} Sistema form is another recorded morphological form, used bibliographically. The reason connects it with earlier invariant-theory terminology and avoidance of cumbersome compounds, but does not separately argue for final -a against sistem. The dictionary must therefore display sistem form and sistema form distinctly rather than attributing an unrecorded uniform gender policy to the records.\par
\noindent{\footnotesize Izvorne zapisy / Original records: HIST-g215-007.}\par
\end{minipage}\par
\par\medskip\Needspace{7\baselineskip}
\noindent\begin{minipage}{\linewidth}
\hypertarget{TERM-0598}{}
\noindent\textbf{svęz medžu formami}\quad{\footnotesize\texttt{TERM-0598}}\par
\noindent\emph{Iztočnik / Source:} Formenzusammenhang\par
\noindent\emph{Rusky / Russian:} связь между формами\par
\noindent\emph{Ukrajinsky / Ukrainian:} зв'язок між формами\par
\smallskip\noindent\emph{Razlog:} Svęz medžu formami zaměnjaje těžky německy kompozit jasnym odnošajnym izrazom. Razlog jest sintaktična prozračnost: predmety svęzi ostavajut izrečene. Konkretna odbacana medžuslovjanska složenica ne jest zabilěžena. V ukrajinskom sravnjenju apostrof jest zachovany kako kodovany znak; jeho čitateljsko pokazanje kako apostrofa ne dobavjaje novu bukvu ili novu attestaciju.\par
\noindent\emph{Reason:} Svęz medžu formami renders the dense German compound as a clear relational phrase. The reason is syntactic transparency: the related objects remain explicit. No particular rejected Interslavic compound is recorded. In the Ukrainian comparison, the apostrophe survives as an encoded character; displaying it as an apostrophe supplies neither a new letter nor a new attestation.\par
\noindent{\footnotesize Izvorne zapisy / Original records: HIST-g028-003.}\par
\end{minipage}\par
\par\medskip\Needspace{7\baselineskip}
\noindent\begin{minipage}{\linewidth}
\hypertarget{TERM-0601}{}
\noindent\textbf{svobodne od $x_1,\ldots,x_{i-1}$}\quad{\footnotesize\texttt{TERM-0601}}\par
\noindent\emph{Iztočnik / Source:} frei von $x_1,\ldots,x_{i-1}$\par
\noindent\emph{Rusky / Russian:} свободные от $x_1,\ldots,x_{i-1}$\par
\noindent\emph{Ukrajinsky / Ukrainian:} вільні від $x_1,\ldots,x_{i-1}$\par
\smallskip\noindent\emph{Razlog:} Svobodne od $x_1,\ldots,x_{i-1}$ govori, že navedene prěměnne ne nastupajut v polinomu. Predlog od jasno svęzuje odsutnost s točnym spiskom prěměnnyh; ne ide o svobodny modul ni o linearnu nezavisnost. Zapis podkrepljaje blizku konstrukciju ruskim i ukrajinskim sravnjenjem. Dva povęzane zapisy imajut jedin zachovany razlog, ne dva nezavisne dokazy jezyčnoj normy.\par
\noindent\emph{Reason:} Svobodne od $x_1,\ldots,x_{i-1}$ says that the listed variables do not occur in the polynomial. The preposition od ties absence to the exact list of variables; this is not freeness of a module or linear independence. The rationale supports close wording with Russian and Ukrainian comparanda. The two associated records contain one preserved rationale, not two independent attestations of a linguistic norm.\par
\noindent{\footnotesize Izvorne zapisy / Original records: HIST-g087-008, HIST-g092-028.}\par
\end{minipage}\par
\par\medskip\Needspace{7\baselineskip}
\noindent\begin{minipage}{\linewidth}
\hypertarget{TERM-0602}{}
\noindent\textbf{Frobeniusova teorija grupnyh karakterov}\quad{\footnotesize\texttt{TERM-0602}}\par
\noindent\emph{Iztočnik / Source:} Frobeniussche Theorie der Gruppencharaktere\par
\noindent\emph{Rusky / Russian:} фробениусова теория характеров групп\par
\noindent\emph{Ukrajinsky / Ukrainian:} фробеніусова теорія характерів груп\par
\smallskip\noindent\emph{Razlog:} Frobeniusova jest pridavničny oblik vlastnogo imena, soglasovany s teorija. Zapis razlikuje njegovo pisanje v običnoj proze od sohranjenja vlastnogo imena v bibliografiji: kirilična versija transliteruje pridavnik, a bibliografično ime ostavaje latinično. Grupnyh karakterov sohranjaje točny predmet teorije. To razlikovanje ne znači, že vse pojavjenja imena morajut iměti jedno i to samo grafično pravilo.\par
\noindent\emph{Reason:} Frobeniusova is an adjectival form of the proper name, agreeing with teorija. The rationale distinguishes its treatment in ordinary prose from preservation of a name in a bibliography: the Cyrillic version transliterates the adjective, while the bibliographical name remains Latin. Grupnyh karakterov retains the theory's precise subject. The distinction is contextual, not a rule that every occurrence of the name must receive identical graphic treatment.\par
\noindent{\footnotesize Izvorne zapisy / Original records: HIST-g105-002.}\par
\end{minipage}\par
\par\medskip\Needspace{7\baselineskip}
\noindent\begin{minipage}{\linewidth}
\hypertarget{TERM-0603}{}
\noindent\textbf{prědšedši paragrafy}\quad{\footnotesize\texttt{TERM-0603}}\par
\noindent\emph{Iztočnik / Source:} frühere Paragraphen\par
\noindent\emph{Rusky / Russian:} предыдущие параграфы\par
\noindent\emph{Ukrajinsky / Ukrainian:} попередні параграфи\par
\smallskip\noindent\emph{Razlog:} Prědšedši paragrafy sohranjaje Noetherino odsyłanje k ranějšim oddělam. Paragraf tu označaje čislovany odděl teksta, ne nužno jedin tekstovy odstup. Izbor pridavnika jest v zapise izrično ostavjen daljšemu pravopisnomu ili leksičnomu uskladnjenju. Zato jasnost odsyłanja i pravilnost samogo pridavnika ne trěba smatrati jednoj už rěšenoj věčju.\par
\noindent\emph{Reason:} Prědšedši paragrafy retains Noether's reference to earlier numbered sections. Here paragraf denotes a section, not necessarily a single typographical paragraph. The record expressly leaves the exact adjective open to later normalization. Clarity of the reference and correctness of that particular adjective are therefore separate questions.\par
\noindent{\footnotesize Izvorne zapisy / Original records: HIST-g017-009.}\par
\end{minipage}\par
\par\medskip\Needspace{7\baselineskip}
\noindent\begin{minipage}{\linewidth}
\hypertarget{TERM-0604}{}
\noindent\textbf{fundamentalno uravnenje}\quad{\footnotesize\texttt{TERM-0604}}\par
\noindent\emph{Iztočnik / Source:} Fundamentalgleichung\par
\noindent\emph{Rusky / Russian:} фундаментальное уравнение\par
\noindent\emph{Ukrajinsky / Ukrainian:} фундаментальне рівняння\par
\smallskip\noindent\emph{Razlog:} Fundamentalno uravnenje sohranjaje internacionalny pridavnik kako čest točnogo nazva. Razlog vezuje nazv k identifikaciji $P^{(i)}$ s fundamentalnym jednačenjem razširenja polja. On ne govori, že vsako važno jednačenje trěba nazvati fundamentalnym. Izbor uravnenje takože ostavaje različny od jednačenje v TERM-0605; sravnjenje tyh oblikov ne směje skryti jih različne konteksty.\par
\noindent\emph{Reason:} Fundamentalno uravnenje preserves the international adjective as part of a specific designation. The rationale ties it to the identification of $P^{(i)}$ with the fundamental equation of a field extension. It does not recommend calling every important equation fundamental. Uravnenje also differs from jednačenje in TERM-0605; comparing those forms must not obscure their different contexts.\par
\noindent{\footnotesize Izvorne zapisy / Original records: HIST-g097-016.}\par
\end{minipage}\par
\par\medskip\Needspace{7\baselineskip}
\noindent\begin{minipage}{\linewidth}
\hypertarget{TERM-0605}{}
\noindent\textbf{fundamentalno jednačenje / fundamentalna forma}\quad{\footnotesize\texttt{TERM-0605}}\par
\noindent\emph{Iztočnik / Source:} Fundamentalgleichung / Fundamentalform\par
\noindent\emph{Rusky / Russian:} фундаментальное уравнение / фундаментальная форма\par
\noindent\emph{Ukrajinsky / Ukrainian:} фундаментальне рівняння / фундаментальна форма\par
\smallskip\noindent\emph{Razlog:} Dva nazva, fundamentalno jednačenje i fundamentalna forma, sohranjajut historijnu terminologiju, ktora svęzuje definiciju diferenty s običnym proizvodnym polinoma. Jednačenje jest izbrano s naměroju južno- i zapadnoslovjanskoj čitateljnosti. To jest zapisana motivacija, ale priložena sravnjenja sut samo rusko i ukrajinsko; konkretne južne i zapadne pasaže ni čitateljski test tu ne sut navedene. Forma i jednačenje ostavajut dva nazva v argumentu, ne slobodno zaměnime slova.\par
\noindent\emph{Reason:} The two names, fundamentalno jednačenje and fundamentalna forma, retain historical terminology connecting the ideal-theoretic definition of the different with the usual polynomial derivative. Jednačenje is chosen with South and West Slavic recognizability in view. That is the recorded motivation, but the supplied comparanda are Russian and Ukrainian; no particular South or West Slavic passage or reader test accompanies it. Form and equation remain distinct names within the argument, not freely interchangeable words.\par
\noindent{\footnotesize Izvorne zapisy / Original records: HIST-g213-014.}\par
\end{minipage}\par
\par\medskip\Needspace{7\baselineskip}
\noindent\begin{minipage}{\linewidth}
\hypertarget{TERM-0606}{}
\noindent\textbf{fundamentalno uravnenje / fundamentalna forma}\quad{\footnotesize\texttt{TERM-0606}}\par
\noindent\emph{Iztočnik / Source:} Fundamentalgleichung / Fundamentalform\par
\noindent\emph{Rusky / Russian:} фундаментальное уравнение / фундаментальная форма\par
\noindent\emph{Ukrajinsky / Ukrainian:} фундаментальне рівняння / фундаментальна форма\par
\smallskip\noindent\emph{Razlog:} Fundamentalno uravnenje / fundamentalna forma sojedinjaje internacionalny korenj fundamental- so slovjanskim nazvom jednačenja. Razlog podkrepljaje historijnu kontinuitetnost i jasnost dlja sovrěmennogo čitatelja. On ne daje odvojenogo sravnjenja uravnenje s jednačenje. Obě zabilěžene formy zato sut vidime v slovniku; ne trěba iz različnyh zapisov tvoriti nezapisany konkurs s čislenym poběditeljem.\par
\noindent\emph{Reason:} Fundamentalno uravnenje / fundamentalna forma combines the international stem fundamental- with a Slavic term for equation. The rationale emphasizes historical continuity and clarity for modern readers. It offers no separate comparison of uravnenje with jednačenje. Both recorded forms therefore remain visible; differing records must not be turned into an undocumented competition with a numerical winner.\par
\noindent{\footnotesize Izvorne zapisy / Original records: HIST-g102-010.}\par
\end{minipage}\par
\par\medskip\Needspace{7\baselineskip}
\noindent\begin{minipage}{\linewidth}
\hypertarget{TERM-0607}{}
\noindent\textbf{fundamentalna točka}\quad{\footnotesize\texttt{TERM-0607}}\par
\noindent\emph{Iztočnik / Source:} Fundamentalpunkt\par
\noindent\emph{Rusky / Russian:} фундаментальная точка\par
\noindent\emph{Ukrajinsky / Ukrainian:} фундаментальна точка\par
\smallskip\noindent\emph{Razlog:} Fundamentalna točka prodolžaje geometričny nazv iz odděla 5 raboty 6. Internacionalny pridavnik jest sojedinjen s točka i soglasovan s njeju v rodě. Rusko i ukrajinsko sravnjenja podkrepljajut blizkost celej konstrukcije. Kratky razlog ne oprěděljaje sam geometričny objekt: fundamentalna točka ne može byti utožsamjena s proizvoljnoju osoblivoju točkoju bez čitanja definicije.\par
\noindent\emph{Reason:} Fundamentalna točka continues the geometrical designation from section 5 of Paper 6. The international adjective is combined with točka and agrees with it in gender. Russian and Ukrainian parallels support the construction's similarity across those languages. The short rationale does not itself define the geometrical object: a fundamental point cannot be identified with an arbitrary special point without the relevant definition.\par
\noindent{\footnotesize Izvorne zapisy / Original records: HIST-g052-006.}\par
\end{minipage}\par
\par\medskip\Needspace{7\baselineskip}
\noindent\begin{minipage}{\linewidth}
\hypertarget{TERM-0608}{}
\noindent\textbf{fundamentalne točke involucije}\quad{\footnotesize\texttt{TERM-0608}}\par
\noindent\emph{Iztočnik / Source:} Fundamentalpunkte der Involution\par
\noindent\emph{Rusky / Russian:} фундаментальные точки инволюции\par
\noindent\emph{Ukrajinsky / Ukrainian:} фундаментальні точки інволюції\par
\smallskip\noindent\emph{Razlog:} Fundamentalne točke involucije sohranjaje množinu i genitivnu svez s konkretnoju involucijeju. Razlog objašnjava kontinuitetnost nazva od uvoda raboty 6 do daljših razsmotrenij regularnyh sistemov i sistemov bez fundamentalnyh toček. Zato dopolnjenje involucije jest čestju točnogo predmeta, ne ukrasom. Zapis ne davaje osnovy zaměniti nazv modernym terminom bez dodatnogo matematičnogo sravnjenja.\par
\noindent\emph{Reason:} Fundamentalne točke involucije preserves the plural and the genitive relation to the particular involution. The rationale explains continuity from Paper 6's introduction to later discussion of regular systems and systems without fundamental points. Involucije specifies the object under discussion; it is not ornamental wording. The record does not justify substitution of a modern name without further mathematical comparison.\par
\noindent{\footnotesize Izvorne zapisy / Original records: HIST-g045-012.}\par
\end{minipage}\par
\par\medskip\Needspace{7\baselineskip}
\noindent\begin{minipage}{\linewidth}
\hypertarget{TERM-0609}{}
\noindent\textbf{fundamentalny teorem teorije eliminacije}\quad{\footnotesize\texttt{TERM-0609}}\par
\noindent\emph{Iztočnik / Source:} Fundamentalsatz der Eliminationstheorie\par
\noindent\emph{Rusky / Russian:} основная теорема теории элиминации\par
\noindent\emph{Ukrajinsky / Ukrainian:} фундаментальна теорема теорії елімінації\par
\smallskip\noindent\emph{Razlog:} Fundamentalny teorem teorije eliminacije sohranjaje historijno ime teorema. Razlog svęzuje jedinstvenost asociovanyh prostyh idealov s razložbenym teoremom teorije eliminacije v osoblivyh polinomnyh oblastjah. Pridavnik fundamentalny zato služi identifikaciji rezultata, ne ocěně njegovej važnosti. Rusko sravnjenje upotrěbjaje osnovnaja, ukrajinsko fundamentalna; ta razlika ne jest sama po sobě razlog prěimenovati avtorsky teorem.\par
\noindent\emph{Reason:} Fundamentalny teorem teorije eliminacije preserves a historical theorem name. The rationale connects uniqueness of associated prime ideals with the elimination-theoretic decomposition theorem in particular polynomial domains. Fundamentalny thus identifies the result rather than appraising its importance. The Russian comparator uses osnovnaja and the Ukrainian fundamentalna; that difference alone is not a reason to rename the author's theorem.\par
\noindent{\footnotesize Izvorne zapisy / Original records: HIST-g081-082.}\par
\end{minipage}\par
\par\medskip\Needspace{7\baselineskip}
\noindent\begin{minipage}{\linewidth}
\hypertarget{TERM-0610}{}
\noindent\textbf{funkcija prvogo roda / polje prvogo roda}\quad{\footnotesize\texttt{TERM-0610}}\par
\noindent\emph{Iztočnik / Source:} Funktion erster Art / Körper erster Art\par
\noindent\emph{Rusky / Russian:} функция первого рода / поле первого рода\par
\noindent\emph{Ukrajinsky / Ukrainian:} функція першого роду / поле першого роду\par
\smallskip\noindent\emph{Razlog:} Funkcija prvogo roda / polje prvogo roda sohranjaje označenje iz Steinitzovej doby, a ne zaměnjaje go novym nazvom klasa polj. Točno značenje zavisi od definicije v danoj rabotě. Novo sravnjenje niže pokazuje svez s različnymi linearnymi faktorami v rabotě 24, a takože drugu definiciju prvogo roda v rabotě 6. Jednaka fraza zato ne dokazuje jednakosti dvuh matematičnyh pojmov.\par
\noindent\emph{Reason:} Funkcija prvogo roda / polje prvogo roda preserves a designation from the Steinitz period instead of replacing it with a newer field classification. Its precise meaning depends on the definition in the particular paper. The fresh comparison below shows the connection with distinct linear factors in Paper 24 and a different first-kind definition in Paper 6. Identical phrasing does not establish identity of the two mathematical concepts.\par
\noindent{\footnotesize Izvorne zapisy / Original records: HIST-g100-007.}\par
\smallskip\noindent\emph{Sravnjenje s tekstom:} V odděle 1, punktah 14 i 15 raboty 24, prvogo roda jest svęzano s razpadom proste funkcije na različne linearne faktory i s eksponentom $f=0$; daljši punkt opisuje podpolje elementov prvogo roda. Naprotiv, definicija VII v rabotě 6 oprěděljaje relativno cělu oblast prvogo roda prěz oblast odobraženja i algebraičnu cělost polinomov nad njeju. To sut različne definicije dlja različnyh objektov. Slovnična fraza mora byti svęzana s točnym městom upotrěby; javna svez s sovrěmennoju terminologijeju može byti objasnjena odvojeno, ne tiho podstavjena v avtorsky tekst.\par
\noindent\emph{Textual comparison:} In Paper 24 section 1, items 14 and 15, first kind is tied to splitting of a simple function into distinct linear factors and to exponent $f=0$; the next item describes the subfield of first-kind elements. In contrast, definition VII in Paper 6 defines a relatively integral domain of the first kind by means of an image domain and the algebraic integrality of polynomials over it. These are different definitions for different objects. The dictionary phrase needs an exact context of use; a relationship to modern terminology can be explained separately, not silently substituted in the author's text.\par
\noindent{\footnotesize TEXT-COMP-23; pasaže / passages: 6065, 15363, 15365; \texttt{CORPUS\_FORM\_COMPARISONS.json}.}\par
\end{minipage}\par
\par\medskip\Needspace{7\baselineskip}
\noindent\begin{minipage}{\linewidth}
\hypertarget{TERM-0611}{}
\noindent\textbf{funkcionalna oblast}\quad{\footnotesize\texttt{TERM-0611}}\par
\noindent\emph{Iztočnik / Source:} Funktionalbereich\par
\noindent\emph{Rusky / Russian:} функциональная область\par
\noindent\emph{Ukrajinsky / Ukrainian:} функціональна область\par
\smallskip\noindent\emph{Razlog:} Funkcionalna oblast sohranjaje historijny nazv oblasti postrojenej iz racionalnyh funkcionalov. Tri razloga ukazujut na rědku sovrěmennu upotrěbu, kontinuitetnost medžu rabotami 9, 30 i 31, i neželatelnost užšego nazva koordinatno kolco. Prosba o odbornu proverku jest svedočstvom nejistosti, ne už dobytogo odobrenja. Oblast tu treba čitati s avtorskoj konstrukcijeju, ne kako proizvoljny prostor funkcij.\par
\noindent\emph{Reason:} Funkcionalna oblast retains the historical name of a domain constructed from rational functionals. Three rationales identify its uncommon modern usage, continuity across Papers 9, 30 and 31, and the danger of the narrower name coordinate ring. A request for specialist checking records uncertainty, not approval already obtained. Oblast must be read with the author's construction, not as an arbitrary space of functions.\par
\noindent{\footnotesize Izvorne zapisy / Original records: HIST-g060-007, HIST-g110-005, HIST-g130-010.}\par
\end{minipage}\par
\par\medskip\Needspace{7\baselineskip}
\noindent\begin{minipage}{\linewidth}
\hypertarget{TERM-0612}{}
\noindent\textbf{funkcionalny determinant}\quad{\footnotesize\texttt{TERM-0612}}\par
\noindent\emph{Iztočnik / Source:} Funktionaldeterminante\par
\noindent\emph{Rusky / Russian:} функциональный детерминант\par
\noindent\emph{Ukrajinsky / Ukrainian:} функціональний детермінант\par
\smallskip\noindent\emph{Razlog:} Funkcionalny determinant jest prodolženje nazva v několiko oddělah raboty 2. Tri zachovane razloga podkrepljajut internacionalny matematičny korenj i ustojčivost imenovanja prědměta v dokazu. Oni ne predlagajut novogo sinonima. Povtorjenje nazva pomaga čitatelju slěditi argument, ale samo čislo povtorjenij ne jest nezavisnym dokazom medžuslovjanskoj razumělivosti.\par
\noindent\emph{Reason:} Funkcionalny determinant continues the designation through several sections of Paper 2. The three preserved rationales emphasize an international mathematical stem and stable naming within the proof. They propose no new synonym. Repetition helps readers follow the argument, but its frequency is not independent evidence of cross-Slavic intelligibility.\par
\noindent{\footnotesize Izvorne zapisy / Original records: HIST-g013-007, HIST-g019-014, HIST-g020-010.}\par
\end{minipage}\par
\par\medskip\Needspace{7\baselineskip}
\noindent\begin{minipage}{\linewidth}
\hypertarget{TERM-0613}{}
\noindent\textbf{funkcionalny determinant / funkcionalna matrica}\quad{\footnotesize\texttt{TERM-0613}}\par
\noindent\emph{Iztočnik / Source:} Funktionaldeterminante / Funktionalmatrix\par
\noindent\emph{Rusky / Russian:} функциональный определитель / функциональная матрица\par
\noindent\emph{Ukrajinsky / Ukrainian:} функціональний визначник / функціональна матриця\par
\smallskip\noindent\emph{Razlog:} Funkcionalny determinant / funkcionalna matrica sohranjaje avtorsko razlikovanje determinanta i matrice. Razlog izrično odbija polnu zaměnu nazvom Jacobian, ukazujuči na formalny kontekst proizvodnyh u Noether. Rusko i ukrajinsko sravnjenja davajut domače nazvy determinanta, «определитель» i «визначник», zatym kako medžuslovjanski zapis izbiraje internacionalny determinant. Ne jest zabilěženo čisleno sravnjenje tyh korenov; važny jest i rozděl medžu samoj matriceju i jej determinantom.\par
\noindent\emph{Reason:} Funkcionalny determinant / funkcionalna matrica preserves the author's distinction between determinant and matrix. The rationale explicitly rejects full replacement by Jacobian, referring to Noether's formal derivative context. The Russian and Ukrainian comparanda supply native determinant terms, opredelitel and vyznačnyk, whereas the Interslavic record selects international determinant. No numerical comparison of those stems is recorded; the distinction between the matrix and its determinant matters in its own right.\par
\noindent{\footnotesize Izvorne zapisy / Original records: HIST-g205-005.}\par
\end{minipage}\par
\par\medskip\Needspace{7\baselineskip}
\noindent\begin{minipage}{\linewidth}
\hypertarget{TERM-0614}{}
\noindent\textbf{funkcionalne determinanty}\quad{\footnotesize\texttt{TERM-0614}}\par
\noindent\emph{Iztočnik / Source:} Funktionaldeterminanten\par
\noindent\emph{Rusky / Russian:} функциональные детерминанты\par
\noindent\emph{Ukrajinsky / Ukrainian:} функціональні детермінанти\par
\smallskip\noindent\emph{Razlog:} Funkcionalne determinanty jest množinna forma togo že internacionalnogo nazva. Tri razloga iz rabot 2 i 6 davajut prednost kratkomu poznatomu terminu prěd opisnym prěvodom, osoblivo v tvrdženju s mnogymi formulami. Razlog ne oprěděljaje konkretno odbijenu parafrazu, zato ju ne trěba izmisliti. Množina i pojavljanja v različnyh tvrdženjah ostavajut čestju dokazov upotrěby, ne novymi nezavisnymi slovničnymi korenjami.\par
\noindent\emph{Reason:} Funkcionalne determinanty is the plural of the same international designation. Three rationales from Papers 2 and 6 prefer a short recognizable term to descriptive paraphrase, especially in a statement containing many formulas. They name no particular rejected paraphrase, so none should be invented. Plural form and occurrence in different statements remain evidence of use, not new independent lexical stems.\par
\noindent{\footnotesize Izvorne zapisy / Original records: HIST-g005-008, HIST-g012-006, HIST-g040-018.}\par
\end{minipage}\par
\par\medskip\Needspace{7\baselineskip}
\noindent\begin{minipage}{\linewidth}
\hypertarget{TERM-0615}{}
\noindent\textbf{funkcionalne uravnenja}\quad{\footnotesize\texttt{TERM-0615}}\par
\noindent\emph{Iztočnik / Source:} Funktionalgleichungen\par
\noindent\emph{Rusky / Russian:} функциональные уравнения\par
\noindent\emph{Ukrajinsky / Ukrainian:} функціональні рівняння\par
\smallskip\noindent\emph{Razlog:} Funkcionalne uravnenja sojedinjaje internacionalny pridavnik s nazvom jednačenij na slovjanskoj osnově. Zapis stremi se k običnomu sovrěmennomu matematičnomu registru. Funkcionalne ne treba zameniti slovom funkcije: predmětem sut jednačenja, ne same funkcije. Razlog ne sravnjaje točno uravnenja s jednačenja i ne ustavjaje izključnu normu dlja vsih daljših rabot.\par
\noindent\emph{Reason:} Funkcionalne uravnenja combines an international adjective with a Slavic-root name for equations. The rationale aims at ordinary modern mathematical register. Funkcionalne should not collapse into functions: the objects named are equations, not the functions themselves. The record supplies no explicit comparison of uravnenja with jednačenja and establishes no exclusive norm for every later paper.\par
\noindent{\footnotesize Izvorne zapisy / Original records: HIST-g069-001.}\par
\end{minipage}\par
\par\medskip\Needspace{7\baselineskip}
\noindent\begin{minipage}{\linewidth}
\hypertarget{TERM-0616}{}
\noindent\textbf{funkcionalna matrica}\quad{\footnotesize\texttt{TERM-0616}}\par
\noindent\emph{Iztočnik / Source:} Funktionalmatrix\par
\noindent\emph{Rusky / Russian:} функциональная матрица\par
\noindent\emph{Ukrajinsky / Ukrainian:} функціональна матриця\par
\smallskip\noindent\emph{Razlog:} Funkcionalna matrica jest blizky prěvod historijnogo nazva. Zapis priznavaje svez izloženej matrice proizvodnyh s sovrěmennym Jacobian, ale ne želi tiho zaměniti avtorsko slovo. Pri novom čitanju odděla 2 raboty 6 možno viděti, kako rang matrice i nenulovy determinant vhodet v argument o algebraičnoj zavisnosti. To podkrepljaje predmetno jasno objasnjenje v slovniku, ne samovoljno prěimenovanje v prěvodu.\par
\noindent\emph{Reason:} Funkcionalna matrica closely translates the historical designation. The rationale recognizes the displayed derivative matrix's relation to a modern Jacobian but avoids silently replacing the author's word. A fresh reading of Paper 6 section 2 shows how matrix rank and a nonzero determinant enter the argument about algebraic dependence. This supports a mathematically explicit dictionary explanation, not unilateral renaming in the translation.\par
\noindent{\footnotesize Izvorne zapisy / Original records: HIST-g042-005.}\par
\smallskip\noindent\emph{Sravnjenje s tekstom:} Odděl 1 raboty 6 razsmatrjaje racionalne funkcije v redukovanom predstavjenju nad koeficientnym poljem $\Omega$ i izrično nazyvaje veličiny iz $\Omega$ funkcijami nultogo stepena. To lokalno oprěděljenje ne govori o nulovoj funkciji. V odděle 2 rang funkcionalnoj matrice jest $\rho$, a formula (2) izdvaja nenulovy determinant čestičnyh proizvodnyh. Sravnjenje čini točnym kontekst dvuh nazvov, ne predlagaje zaměnu historijnoj terminologije ni novu univerzalnu konvenciju o stepeni racionalnyh funkcij.\par
\noindent\emph{Textual comparison:} Section 1 of Paper 6 considers rational functions in reduced form over a coefficient field $\Omega$ and explicitly calls quantities in $\Omega$ functions of degree zero. This local convention does not mean the zero function. In section 2 the functional matrix has rank $\rho$, and equation (2) isolates a nonzero determinant of partial derivatives. The comparison specifies the context of the two names; it proposes neither replacement of the historical terminology nor a new universal convention for the degree of rational functions.\par
\noindent{\footnotesize TEXT-COMP-22; pasaže / passages: 5061, 5063, 5070, 5163, 5165, 5166; \texttt{CORPUS\_FORM\_COMPARISONS.json}.}\par
\end{minipage}\par
\par\medskip\Needspace{7\baselineskip}
\noindent\begin{minipage}{\linewidth}
\hypertarget{TERM-0617}{}
\noindent\textbf{funkcije nultogo stepena}\quad{\footnotesize\texttt{TERM-0617}}\par
\noindent\emph{Iztočnik / Source:} Funktionen nullten Grades\par
\noindent\emph{Rusky / Russian:} функции нулевой степени\par
\noindent\emph{Ukrajinsky / Ukrainian:} функції нульового степеня\par
\smallskip\noindent\emph{Razlog:} Funkcije nultogo stepena sohranjaje točnu formulaciju nulovogo stepena. Sam historijny razlog jest velmi kratky; njegov smysl možno utočniti čitanjem lokalnogo dogovora. Rabota 6 izrično utožsamnjaje te funkcije s veličinami iz koeficientnogo polja $\Omega$. Zato ne ide o funkcije, ktore imajut nulovu cěnnost; ne treba takože bezuslovno prěnesti to značenje na vsaku drugu konvenciju o stepeni racionalnoj funkcije.\par
\noindent\emph{Reason:} Funkcije nultogo stepena retains the exact degree-zero wording. The historical reason is very brief; its meaning can be clarified by reading the local convention. Paper 6 explicitly identifies these functions with quantities in the coefficient field $\Omega$. It therefore does not mean functions whose value is zero, nor should that meaning be transferred without qualification to every other convention for the degree of a rational function.\par
\noindent{\footnotesize Izvorne zapisy / Original records: HIST-g041-007.}\par
\smallskip\noindent\emph{Sravnjenje s tekstom:} Odděl 1 raboty 6 razsmatrjaje racionalne funkcije v redukovanom predstavjenju nad koeficientnym poljem $\Omega$ i izrično nazyvaje veličiny iz $\Omega$ funkcijami nultogo stepena. To lokalno oprěděljenje ne govori o nulovoj funkciji. V odděle 2 rang funkcionalnoj matrice jest $\rho$, a formula (2) izdvaja nenulovy determinant čestičnyh proizvodnyh. Sravnjenje čini točnym kontekst dvuh nazvov, ne predlagaje zaměnu historijnoj terminologije ni novu univerzalnu konvenciju o stepeni racionalnyh funkcij.\par
\noindent\emph{Textual comparison:} Section 1 of Paper 6 considers rational functions in reduced form over a coefficient field $\Omega$ and explicitly calls quantities in $\Omega$ functions of degree zero. This local convention does not mean the zero function. In section 2 the functional matrix has rank $\rho$, and equation (2) isolates a nonzero determinant of partial derivatives. The comparison specifies the context of the two names; it proposes neither replacement of the historical terminology nor a new universal convention for the degree of rational functions.\par
\noindent{\footnotesize TEXT-COMP-22; pasaže / passages: 5061, 5063, 5070, 5163, 5165, 5166; \texttt{CORPUS\_FORM\_COMPARISONS.json}.}\par
\end{minipage}\par
\par\medskip\Needspace{7\baselineskip}
\noindent\begin{minipage}{\linewidth}
\hypertarget{TERM-0618}{}
\noindent\textbf{polja funkcij}\quad{\footnotesize\texttt{TERM-0618}}\par
\noindent\emph{Iztočnik / Source:} Funktionenkörper\par
\noindent\emph{Rusky / Russian:} поля функций\par
\noindent\emph{Ukrajinsky / Ukrainian:} поля функцій\par
\smallskip\noindent\emph{Razlog:} Polja funkcij upotrěbjaje algebraičny nazv polje, a ne doslovny prěvod običnogo německogo slova Körper. Genitiv funkcij ukazuje, iz kakih objektov polja sostojat; množina polja ostavaje zachovana. Rusko i ukrajinsko sravnjenja imajut tu že konstrukciju. To jest izbor disciplinarnogo značenja, ne tvrdženje, že polje može avtomatično zaměniti Körper vo vsej matematikě ili obyčnom jezyku.\par
\noindent\emph{Reason:} Polja funkcij uses the algebraic term polje rather than a literal translation of the ordinary German word Körper. The genitive funkcij identifies the objects comprising the fields, and the plural polja is retained. Russian and Ukrainian comparanda use the same construction. This is a choice of disciplinary sense, not a claim that polje can automatically replace Körper throughout mathematics or ordinary language.\par
\noindent{\footnotesize Izvorne zapisy / Original records: HIST-g094-011.}\par
\end{minipage}\par
\par\medskip\Needspace{7\baselineskip}
\noindent\begin{minipage}{\linewidth}
\hypertarget{TERM-0619}{}
\noindent\textbf{provodnik}\quad{\footnotesize\texttt{TERM-0619}}\par
\noindent\emph{Iztočnik / Source:} Führer\par
\noindent\emph{Rusky / Russian:} проводник\par
\noindent\emph{Ukrajinsky / Ukrainian:} провідник\par
\smallskip\noindent\emph{Razlog:} Provodnik jest tu termin Artinovej teorije, ne obyčno ime osoby. Razlog upotrěbjaje blizkost ruskogo provodnik i ukrajinskogo providnyk kako osnovu medžuslovjanskogo oblika. Ta leksična svez ne davaje sama po sobě definicije provodnika i ne čini go totožnym s vsakym drugym algebraičnym objektom nazvanym provodnik. Konkretna definicija raboty 37 ostavaje potrěbna pri upotrěbě termina.\par
\noindent\emph{Reason:} Provodnik is an Artin-theoretic term here, not an ordinary name for a person. The rationale uses the resemblance between Russian provodnik and Ukrainian providnyk to support the Interslavic form. That lexical relation does not itself define the conductor or identify it with every other algebraic object bearing the same name. The specific definition in Paper 37 remains necessary when using the term.\par
\noindent{\footnotesize Izvorne zapisy / Original records: HIST-g207-007.}\par
\end{minipage}\par
\par\medskip\Needspace{7\baselineskip}
\noindent\begin{minipage}{\linewidth}
\hypertarget{TERM-0620}{}
\noindent\textbf{provodnik / porjadok / oblast cělosti}\quad{\footnotesize\texttt{TERM-0620}}\par
\noindent\emph{Iztočnik / Source:} Führer / Ordnung / Integritätsbereich\par
\noindent\emph{Rusky / Russian:} проводник / порядок / область целостности\par
\noindent\emph{Ukrajinsky / Ukrainian:} провідник / порядок / область цілісності\par
\smallskip\noindent\emph{Razlog:} Provodnik / porjadok / oblast cělosti jest skupina nazvov iz teorije prstenov i porjadkov, ne tri sinonimy. Zachovany razlog podrobněj govori samo o provodniku, svęzujuči go s ruskim i ukrajinskim obrazcem v algebraičnoj teoriji čisel. On ne davaje odvojenogo argumenta za vsaky iz ostalyh dvoh nazvov. Tu mežu dokumentacije treba sohraniti: povęzany spis ne jest polnym nezavisnym objasnjenjem vsih njegovih česti.\par
\noindent\emph{Reason:} Provodnik / porjadok / oblast cělosti is a group of terms from ring and order theory, not three synonyms. The preserved rationale discusses only the conductor in detail, linking it to the Russian and Ukrainian pattern in algebraic number theory. It supplies no separate argument for each of the other two names. This documentary limit matters: a grouped list is not a full independent justification of every member.\par
\noindent{\footnotesize Izvorne zapisy / Original records: HIST-g076-012.}\par
\end{minipage}\par
\par\medskip\Needspace{7\baselineskip}
\noindent\begin{minipage}{\linewidth}
\hypertarget{TERM-0621}{}
\noindent\textbf{Galoisova rezolventa}\quad{\footnotesize\texttt{TERM-0621}}\par
\noindent\emph{Iztočnik / Source:} Galoische Resolvente\par
\noindent\emph{Rusky / Russian:} резольвента Galois\par
\noindent\emph{Ukrajinsky / Ukrainian:} резольвента Galois\par
\smallskip\noindent\emph{Razlog:} Galoisova rezolventa razlikuje vlastno ime od nazva matematičnogo objekta. Rezolventa jest termin; Galoisova čini svez s imenem gramatično javnoju. Zapis za rabotu 23 sohranjaje latinično ime, ale ne davaje točnogo ranějšego i pozdnějšego teksta spomenutoj formulacije o sistemu koeficientov. Iz samogo slovničnogo zapisa zato ne možno rekonstruovati tu formulaciju; izbor nazva ostavaje dokumentovan.\par
\noindent\emph{Reason:} Galoisova rezolventa distinguishes the proper name from the mathematical object-name. Rezolventa is the term; Galoisova makes the relation to the name grammatical. The Paper 23 record preserves the Latin name but supplies no exact before-and-after text for the coefficient-system wording it mentions. That wording cannot be reconstructed from this dictionary record alone; the naming choice itself is documented.\par
\noindent{\footnotesize Izvorne zapisy / Original records: HIST-g094-019.}\par
\end{minipage}\par
\par\medskip\Needspace{7\baselineskip}
\noindent\begin{minipage}{\linewidth}
\hypertarget{TERM-0622}{}
\noindent\textbf{Galoisov modul}\quad{\footnotesize\texttt{TERM-0622}}\par
\noindent\emph{Iztočnik / Source:} Galoismodul\par
\noindent\emph{Rusky / Russian:} модуль Галуа\par
\noindent\emph{Ukrajinsky / Ukrainian:} модуль Галуа\par
\smallskip\noindent\emph{Razlog:} Galoisov modul sojedinjaje prěvojen pridavnik s internacionalnym modul. Rusko i ukrajinsko sravnjenja imajut konstrukciju s imenem Галуа posle imenika. Medžuslovjanski izbor ne trěbuje doslovno kopirovati taj porędok slov. On nazyvaje objekt teorije predstavjenij; kratky razlog ne davaje novoj definicije modula ni testa razumělivosti.\par
\noindent\emph{Reason:} Galoisov modul combines a possessive adjective with international modul. The Russian and Ukrainian comparanda place Галуа after the noun. The Interslavic construction need not copy that word order literally. It names a representation-theoretic object; the short rationale supplies neither a new definition of module nor a comprehension test.\par
\noindent{\footnotesize Izvorne zapisy / Original records: HIST-g207-009.}\par
\end{minipage}\par
\par\medskip\Needspace{7\baselineskip}
\noindent\begin{minipage}{\linewidth}
\hypertarget{TERM-0623}{}
\noindent\textbf{Galoisovy modul}\quad{\footnotesize\texttt{TERM-0623}}\par
\noindent\emph{Iztočnik / Source:} Galoismodul\par
\noindent\emph{Rusky / Russian:} [Bukvy ne sut zachovane / Lettering not recoverable]\par
\noindent\emph{Ukrajinsky / Ukrainian:} [Bukvy ne sut zachovane / Lettering not recoverable]\par
\smallskip\noindent\emph{Razlog:} Galoisovy modul ima tu že motivaciju vlastnogo imena i poznatogo algebraičnogo imenika, ale zapisany pridavnik razlikuje se od Galoisov v TERM-0622. Ruske i ukrajinske bukvy v tutom zapise ne sut zachovane. Ne trěba jih dopolnjati po susědnom zapise ni predstavljati dva različne oblika kako već uskladnjenu normu. Morfologična razlika ostavaje javna dlja sravnjenja.\par
\noindent\emph{Reason:} Galoisovy modul has the same motivation of a proper-name adjective and a recognizable algebraic noun, but its recorded adjective differs from Galoisov in TERM-0622. The Russian and Ukrainian lettering is not recoverable here. It must not be supplied from the neighbouring record, nor should the two forms be presented as an already harmonized norm. Their morphological difference remains available for comparison.\par
\noindent{\footnotesize Izvorne zapisy / Original records: HIST-g209-007.}\par
\end{minipage}\par
\par\medskip\Needspace{7\baselineskip}
\noindent\begin{minipage}{\linewidth}
\hypertarget{TERM-0624}{}
\noindent\textbf{Galoisova forma grupy}\quad{\footnotesize\texttt{TERM-0624}}\par
\noindent\emph{Iztočnik / Source:} Galoissche Form der Gruppe\par
\noindent\emph{Rusky / Russian:} форма Галуа группы\par
\noindent\emph{Ukrajinsky / Ukrainian:} форма Галуа групи\par
\smallskip\noindent\emph{Razlog:} Galoisova forma grupy označaje produktnu formu pripisanu grupě permutacij. Razlog za kirilično čitanje navodi konstrukciju s neizměnjenym imenem Галуа vo funkciji genitivnogo dopolnjenja. To jest gramatična i pismena konvencija, ne novy matematičny objekt. Ista motivacija ne dokazuje, že vsaka druga kirilična pojavjava imena v korpusu ima tu že formu.\par
\noindent\emph{Reason:} Galoisova forma grupy names the product form attached to a permutation group. For Cyrillic reading, the rationale specifies unchanged Галуа used as a genitive complement. This is a grammatical and script convention, not a different mathematical object. The same motivation does not prove that every other Cyrillic occurrence of the name in the corpus uses that form.\par
\noindent{\footnotesize Izvorne zapisy / Original records: HIST-g053-004.}\par
\end{minipage}\par
\par\medskip\Needspace{7\baselineskip}
\noindent\begin{minipage}{\linewidth}
\hypertarget{TERM-0625}{}
\noindent\textbf{Galoisova polja / Galoisovo razširjeno polje}\quad{\footnotesize\texttt{TERM-0625}}\par
\noindent\emph{Iztočnik / Source:} Galoissche Körper / Galoisscher Erweiterungskörper\par
\noindent\emph{Rusky / Russian:} поля Галуа / поле расширения Галуа\par
\noindent\emph{Ukrajinsky / Ukrainian:} поля Галуа / поле розширення Галуа\par
\smallskip\noindent\emph{Razlog:} Galoisova polja / Galoisovo razširjeno polje sohranjaje internacionalno vlastno ime i dodava slovjansku pridavničku morfologiju. Različne koncovky odpovědajut množině polja i jedině polje. Razširjeno ostavaje smyslovym dopolnjenjem; ne treba go izgubiti pri uskladnjenju imen. Zapis podkrepljaje kontinuitetnost s ranějšimi rabotami, ne novu definiciju Galoisovogo polja.\par
\noindent\emph{Reason:} Galoisova polja / Galoisovo razširjeno polje retains the international proper name with Slavic adjectival morphology. The endings correspond to plural polja and singular polje. Razširjeno remains a meaningful qualifier and must not disappear during name harmonization. The rationale supports continuity with earlier papers, not a new definition of a Galois field.\par
\noindent{\footnotesize Izvorne zapisy / Original records: HIST-g129-007.}\par
\end{minipage}\par
\par\medskip\Needspace{7\baselineskip}
\noindent\begin{minipage}{\linewidth}
\hypertarget{TERM-0626}{}
\noindent\textbf{Galoisova rezolventa}\quad{\footnotesize\texttt{TERM-0626}}\par
\noindent\emph{Iztočnik / Source:} Galoissche Resolvente\par
\noindent\emph{Rusky / Russian:} резольвента Galois; резольвента Галуа\par
\noindent\emph{Ukrajinsky / Ukrainian:} резольвента Galois; резольвента Галуа\par
\smallskip\noindent\emph{Razlog:} Dva zapisy za Galoisova rezolventa podkrepljajut ustojčivost nazva objekta i osobno pravilo za pismo vlastnogo imena. Ruske i ukrajinske sravnjenja zachovajut i latinično Galois, i kirilično Галуа. Ranějši razlog ostavjaje generovany kiriličny oblik dlja sravnjenja s konvencijeju raboty 6; pozdnějši opisuje pravilo automatičnogo prěpisa. To ne jest svedočstvo, že čitatelj už ocěnil oba oblika.\par
\noindent\emph{Reason:} Two records for Galoisova rezolventa emphasize a stable object-name and a separate script rule for the proper name. The Russian and Ukrainian comparanda preserve both Latin Galois and Cyrillic Галуа. The earlier rationale leaves the generated Cyrillic form for comparison with Paper 6's convention; the later one describes a deterministic transcription rule. Neither is evidence that readers have evaluated both forms.\par
\noindent{\footnotesize Izvorne zapisy / Original records: HIST-g056-006, HIST-g106-012.}\par
\end{minipage}\par
\par\medskip\Needspace{7\baselineskip}
\noindent\begin{minipage}{\linewidth}
\hypertarget{TERM-0627}{}
\noindent\textbf{Galoisova teorija nekomutativnyh těl}\quad{\footnotesize\texttt{TERM-0627}}\par
\noindent\emph{Iztočnik / Source:} Galoissche Theorie nichtkommutativer Körper\par
\noindent\emph{Rusky / Russian:} теория Галуа некоммутативных тел\par
\noindent\emph{Ukrajinsky / Ukrainian:} теорія Галуа некомутативних тіл\par
\smallskip\noindent\emph{Razlog:} Galoisova teorija nekomutativnyh těl povęzuje rabotu 33 s ranějšim razsmotrenjem razpadnyh polj i těl. Těl ne može byti avtomatično zaměnjeno s polj: kvalifikator nekomutativnyh jest sučstveny dlja predměta. Rusko i ukrajinsko sravnjenja podkrepljajut to razlikovanje. Kratky razlog motivuje kontinuitetnost matematičnogo programa, ne opću zaměnimost nazvov tělo i polje.\par
\noindent\emph{Reason:} Galoisova teorija nekomutativnyh těl links Paper 33 with the earlier discussion of splitting fields and division bodies. Těl cannot automatically become polj: the noncommutative qualifier is essential to the subject. Russian and Ukrainian comparanda support that distinction. The short rationale motivates continuity of the mathematical programme, not general interchangeability of tělo and polje.\par
\noindent{\footnotesize Izvorne zapisy / Original records: HIST-g177-008.}\par
\end{minipage}\par
\par\medskip\Needspace{7\baselineskip}
\noindent\begin{minipage}{\linewidth}
\hypertarget{TERM-0628}{}
\noindent\textbf{Galoisovo kolco razširenja}\quad{\footnotesize\texttt{TERM-0628}}\par
\noindent\emph{Iztočnik / Source:} Galoisscher Erweiterungsring\par
\noindent\emph{Rusky / Russian:} галуаово кольцо расширения\par
\noindent\emph{Ukrajinsky / Ukrainian:} галоаївське кільце розширення\par
\smallskip\noindent\emph{Razlog:} Galoisovo kolco razširenja jest nazv prstena dobytogo razširenjem koeficientov črez Galoisovo polje. Razlog izrično razlikuje polje, črez ktoro razširjenje ide, od rezultujučego prstena. Prisustvo Galoisovogo imena zato ne opravdava nazvati rezultat poljem. Kolco jest zachovany historijny kandidat; to ne jest samo po sobě izbor jego vměsto prsten vo vseh kontekstah.\par
\noindent\emph{Reason:} Galoisovo kolco razširenja names a ring obtained by extending coefficients through a Galois field. The rationale expressly distinguishes the field through which extension is made from the resulting ring. The presence of Galois's name therefore does not justify calling the result a field. Kolco is the preserved historical candidate, not by itself a decision to prefer it to prsten in every context.\par
\noindent{\footnotesize Izvorne zapisy / Original records: HIST-g213-009.}\par
\end{minipage}\par
\par\medskip\Needspace{7\baselineskip}
\noindent\begin{minipage}{\linewidth}
\hypertarget{TERM-0629}{}
\noindent\textbf{Galoisovo polje}\quad{\footnotesize\texttt{TERM-0629}}\par
\noindent\emph{Iztočnik / Source:} Galoisscher Körper\par
\noindent\emph{Rusky / Russian:} поле Галуа\par
\noindent\emph{Ukrajinsky / Ukrainian:} поле Галуа\par
\smallskip\noindent\emph{Razlog:} Galoisovo polje razvivaje vlastno ime v pridavnik, zatym kako rusko i ukrajinsko поле Галуа ostavjaje ime posle imenika. Razlog propisuje jedinoznačny prěpis pridavnika v kirilicu i razlikuje go od nezavisnogo leksičnogo izbora. Opis pravila sam ne garantuje, že vsaka pojavjava jest pravilno prěpisana. Matematično značenje polja treba brati iz danoj raboty, ne samo iz poznatogo eponima.\par
\noindent\emph{Reason:} Galoisovo polje develops the proper name into an adjective, whereas Russian and Ukrainian поле Галуа leave the name after the noun. The rationale specifies deterministic Cyrillic transcription of the adjective and distinguishes this from an independent lexical choice. A description of the rule does not guarantee correct conversion of every occurrence. The mathematical sense of field must come from the particular paper, not merely the familiar eponym.\par
\noindent{\footnotesize Izvorne zapisy / Original records: HIST-g097-012.}\par
\end{minipage}\par
\par\medskip\Needspace{7\baselineskip}
\noindent\begin{minipage}{\linewidth}
\hypertarget{TERM-0630}{}
\noindent\textbf{integralno zamknjeny / integralna zamknjenost}\quad{\footnotesize\texttt{TERM-0630}}\par
\noindent\emph{Iztočnik / Source:} ganz abgeschlossen\par
\noindent\emph{Rusky / Russian:} интегрально замкнутый / интегральная замкнутость\par
\noindent\emph{Ukrajinsky / Ukrainian:} інтегрально замкнений / інтегральна замкненість\par
\smallskip\noindent\emph{Razlog:} Integralno zamknjeny / integralna zamknjenost davaje pridavničny i imenny oblik jednoj algebraičnoj vlastnosti. Razlog jest kontinuitetnost s uvodom i prvim oddělom raboty 30. Zapis jednovrěmenno ostavjaje medžuslovjansku formulaciju dlja odbornogo sravnjenja. Vnųtrěšnja soglasnost nazvov zato ne jest svjedočstvom o ljudskoj proverke ni o izključnoj jezyčnoj normě.\par
\noindent\emph{Reason:} Integralno zamknjeny / integralna zamknjenost gives adjectival and nominal forms for one algebraic property. The reason is continuity with the introduction and first section of Paper 30. At the same time, the record leaves the Interslavic formulation for specialist consideration. Internal consistency is not evidence of human review or of an exclusive linguistic norm.\par
\noindent{\footnotesize Izvorne zapisy / Original records: HIST-g110-004.}\par
\end{minipage}\par
\par\medskip\Needspace{7\baselineskip}
\noindent\begin{minipage}{\linewidth}
\hypertarget{TERM-0631}{}
\noindent\textbf{cěly vzhodno k $R$}\quad{\footnotesize\texttt{TERM-0631}}\par
\noindent\emph{Iztočnik / Source:} ganz in bezug auf $R$\par
\noindent\emph{Rusky / Russian:} целый над $R$\par
\noindent\emph{Ukrajinsky / Ukrainian:} цілий над $R$\par
\smallskip\noindent\emph{Razlog:} Cěly vzhodno k $R$ pokušaje izložiti relaciju algebraičnoj cělosti. Rusko i ukrajinsko sravnjenja rabe nad, zatym kako medžuslovjanski zapis predlagaje vzhodno k i sam ostavjaje go dlja odbornogo razsmotrenja. Fraza zato jest kandidat, ne dokazana obća konvencija. Rasprava o po odnošenju k v DEC-0009 pokazuje potrěbu odvojiti leksičny izbor predloga od točnogo algebraičnogo značenja.\par
\noindent\emph{Reason:} Cěly vzhodno k $R$ attempts to express the relation of algebraic integrality. The Russian and Ukrainian comparanda use nad, whereas the Interslavic record proposes vzhodno k and itself leaves it for specialist consideration. This is a candidate, not a demonstrated general convention. The discussion of po odnošenju k in DEC-0009 illustrates why the choice of relational phrase must be distinguished from the precise algebraic meaning.\par
\noindent{\footnotesize Izvorne zapisy / Original records: HIST-g108-002.}\par
\end{minipage}\par
\par\medskip\Needspace{7\baselineskip}
\noindent\begin{minipage}{\linewidth}
\hypertarget{TERM-0632}{}
\noindent\textbf{cěly vzhodno k $R$ / elementy cěle vzhodno k $R$}\quad{\footnotesize\texttt{TERM-0632}}\par
\noindent\emph{Iztočnik / Source:} ganz in bezug auf $R$ / ganze Größen\par
\noindent\emph{Rusky / Russian:} целый над $R$ / элементы, целые над $R$\par
\noindent\emph{Ukrajinsky / Ukrainian:} цілий над $R$ / елементи, цілі над $R$\par
\smallskip\noindent\emph{Razlog:} Cěly vzhodno k $R$ / elementy cěle vzhodno k $R$ prodolžaje ranějši predložny izbor i prenosi periodno Größen kako elementy. Razlog želi algebraičnu jasnost vměsto doslovnyh veličin. To ne jest pravo izbrisati različnost avtorskyh objektov: elementy mora odgovarjati konkretnoj konstrukciji. Zapis ne dodava novogo dokaza za spornu frazu vzhodno k.\par
\noindent\emph{Reason:} Cěly vzhodno k $R$ / elementy cěle vzhodno k $R$ continues the earlier prepositional choice and renders period Größen as elementy. The rationale seeks algebraic clarity rather than literal quantities. That does not license erasing distinctions among the author's objects: elementy must fit the actual construction. The record supplies no new evidence for the contested phrase vzhodno k.\par
\noindent{\footnotesize Izvorne zapisy / Original records: HIST-g110-003.}\par
\end{minipage}\par
\par\medskip\Needspace{7\baselineskip}
\noindent\begin{minipage}{\linewidth}
\hypertarget{TERM-0633}{}
\noindent\textbf{cěly po neopreděljenyh}\quad{\footnotesize\texttt{TERM-0633}}\par
\noindent\emph{Iztočnik / Source:} ganz in den Unbestimmten\par
\noindent\emph{Rusky / Russian:} цело по неопределенным\par
\noindent\emph{Ukrajinsky / Ukrainian:} ціла за невизначеними\par
\smallskip\noindent\emph{Razlog:} Cěly po neopreděljenyh sohranjaje uslovje, že kvocient jest polinomom v danih neoprěděljenyh. Razlog izrično razlikuje tu vlastnost od samej algebraičnoj cělosti nad prstenom. Zato cěly ne može byti objasnjen jedinym bezkontekstnym sinonimom. Dopolnjenje po neopreděljenyh nosi matematičnu informaciju i mora ostati v prěvodu.\par
\noindent\emph{Reason:} Cěly po neopreděljenyh preserves the condition that the quotient is polynomial in the specified indeterminates. The rationale explicitly distinguishes that property from mere algebraic integrality over a ring. Cěly therefore cannot be explained by one context-free synonym. Po neopreděljenyh carries mathematical information and must remain in the translation.\par
\noindent{\footnotesize Izvorne zapisy / Original records: HIST-g049-003.}\par
\end{minipage}\par
\par\medskip\Needspace{7\baselineskip}
\noindent\begin{minipage}{\linewidth}
\hypertarget{TERM-0634}{}
\noindent\textbf{cělo-racionalny}\quad{\footnotesize\texttt{TERM-0634}}\par
\noindent\emph{Iztočnik / Source:} ganz rational\par
\noindent\emph{Rusky / Russian:} целорациональный\par
\noindent\emph{Ukrajinsky / Ukrainian:} ціло-раціональний\par
\smallskip\noindent\emph{Razlog:} Cělo-racionalny razlikuje polinomne izrazy od drobnyh racionalnyh funkcij. Razlog podkrepljaje to razlikovanje zato, že avtorska podnožna primětka pozdněje dopuščaje i drobne racionalne funkcije. Kratky blizky prěvod tako sohranjaje proměnu obsega. Cělo- tu ne može bez dodatnogo uslovja značiti, že vsi koeficienty sut cěle čisla.\par
\noindent\emph{Reason:} Cělo-racionalny distinguishes polynomial expressions from fractional rational functions. The rationale emphasizes this because the author's footnote later allows fractional rational functions as well. A close compact translation thus preserves the change in scope. Without an additional condition, cělo- cannot mean that all coefficients are integers.\par
\noindent{\footnotesize Izvorne zapisy / Original records: HIST-g214-008.}\par
\end{minipage}\par
\par\medskip\Needspace{7\baselineskip}
\noindent\begin{minipage}{\linewidth}
\hypertarget{TERM-0635}{}
\noindent\textbf{cěly i primitivny po sootvětnyh neopreděljenyh}\quad{\footnotesize\texttt{TERM-0635}}\par
\noindent\emph{Iztočnik / Source:} ganz und primitiv in den Variablen\par
\noindent\emph{Rusky / Russian:} целый и примитивный по соответствующим переменным\par
\noindent\emph{Ukrajinsky / Ukrainian:} цілий і примітивний за відповідними змінними\par
\smallskip\noindent\emph{Razlog:} Cěly i primitivny po sootvětnyh neopreděljenyh sohranjaje dva odvojene uslovja i spis prěměnnyh, k ktorym se odnosęt. Zachovany razlog govori o cěločislenyh koeficientah, ale novo čitanje odděla 5 raboty 22 takogo ograničenja ne podkrepljaje: razsmatranje idet nad poljem $P(u)$ i njegovymi razširenjami, a cěly znači polinomny po navedenej skupině prěměnnyh. To je konkretna meža historijnogo objasnjenja, ne razlog skryti go. Primitivny ostavaje dodatnym uslovjem, ne sinonimom cěly.\par
\noindent\emph{Reason:} Cěly i primitivny po sootvětnyh neopreděljenyh preserves two separate conditions and the variables to which they refer. The surviving rationale speaks of integer coefficients, but a fresh reading of Paper 22 section 5 does not support that restriction: the setting uses $P(u)$ and its extensions, and cěly means polynomial in the named group of variables. This is a concrete limitation of the historical explanation, not a reason to conceal it. Primitivny remains an additional condition, not a synonym for cěly.\par
\noindent{\footnotesize Izvorne zapisy / Original records: HIST-g089-004, HIST-g092-046.}\par
\smallskip\noindent\emph{Sravnjenje s tekstom:} Uvod raboty 22 dopuščaje libovoljno koeficientno polje. V odděle 5 prěhod k $P(u)$ s prisojedinjenymi prěměnnymi daje moduly linearnyh form, čije cěle veličiny sut polinomy v jednoj neoprěděljenoj. Slědujuči pasaž opisuje $R^{(i)}$ i $E^{(i)}$ kako cěle i primitivne po $x_{i+1},\ldots,x_n$. Tu cělost govori o polinomnom obliku v danih prěměnnyh, ne o neophodimo cěločislenyh koeficientah. Historijno objasnjenje, ktoro davaje poslědnje ograničenje, zato jest preuzko. Tuto sravnjenje proverjaje aktualny prěvod; ne zaměnjaje odvojeno sravnjenje s pečatnym originalom.\par
\noindent\emph{Textual comparison:} The introduction to Paper 22 allows an arbitrary coefficient field. In section 5, passing to $P(u)$ with adjoined variables gives modules of linear forms whose integral quantities are polynomials in one indeterminate. The next passage describes $R^{(i)}$ and $E^{(i)}$ as integral and primitive in $x_{i+1},\ldots,x_n$. Integrality here concerns polynomial form in the specified variables, not necessarily integer coefficients. The historical explanation imposing the latter restriction is therefore too narrow. This comparison checks the current translation; it does not replace separate collation with the printed original.\par
\noindent{\footnotesize TEXT-COMP-24; pasaže / passages: 14209, 14832, 14836; \texttt{CORPUS\_FORM\_COMPARISONS.json}.}\par
\end{minipage}\par
\par\medskip\Needspace{7\baselineskip}
\noindent\begin{minipage}{\linewidth}
\hypertarget{TERM-0636}{}
\noindent\textbf{integralno zamknjeno / integralna zamknjenost}\quad{\footnotesize\texttt{TERM-0636}}\par
\noindent\emph{Iztočnik / Source:} ganz-abgeschlossen / ganze Abgeschlossenheit\par
\noindent\emph{Rusky / Russian:} интегрально замкнутый / интегральная замкнутость\par
\noindent\emph{Ukrajinsky / Ukrainian:} інтегрально замкнений / інтегральна замкненість\par
\smallskip\noindent\emph{Razlog:} Integralno zamknjeno / integralna zamknjenost izbiraje korenj zamknjen- v običnom slovjanskom izražanju zatvorenosti. Razlog davaje prednost poznatoj konstrukciji prěd užšim novotvorom, ale ne imenuje konkretnogo odbijenogo novotvora. Integralno ograničaje značenje do algebraičnoj vlastnosti. Oblik zamknjeno i oblik zamknjeny iz TERM-0630 pokazuju gramatične varianty, ne dva različne matematične pojmy.\par
\noindent\emph{Reason:} Integralno zamknjeno / integralna zamknjenost selects zamknjen- from familiar Slavic closedness vocabulary. The rationale prefers a recognizable construction to a narrower neologism but names no particular rejected neologism. Integralno restricts the sense to the algebraic property. Zamknjeno and zamknjeny in TERM-0630 are grammatical variants, not two different mathematical concepts.\par
\noindent{\footnotesize Izvorne zapisy / Original records: HIST-g108-006.}\par
\end{minipage}\par
\par\medskip\Needspace{7\baselineskip}
\noindent\begin{minipage}{\linewidth}
\hypertarget{TERM-0637}{}
\noindent\textbf{integralna zamknjenost v polju kvocientov}\quad{\footnotesize\texttt{TERM-0637}}\par
\noindent\emph{Iztočnik / Source:} ganze Abgeschlossenheit im Quotientenkörper\par
\noindent\emph{Rusky / Russian:} интегральная замкнутость в поле частных\par
\noindent\emph{Ukrajinsky / Ukrainian:} інтегральна замкненість у полі часток\par
\smallskip\noindent\emph{Razlog:} Integralna zamknjenost v polju kvocientov sohranjaje i vlastnost, i polje, v ktorom ona jest razsmatrjana. Tri zapisy podkrepljajut algebraičnu cělost, ne topologično zatvorjenje, i povęzuju uvod raboty 30 s pozdnějšim oddělom 8. Dopolnjenje v polju kvocientov zato ne jest izlišne. Leksična blizkost korenja zamknjen- služi razumělivosti nazva, no ne zaměnjaje definiciju cělosti.\par
\noindent\emph{Reason:} Integralna zamknjenost v polju kvocientov preserves both the property and the field in which it is considered. Three records emphasize algebraic integrality rather than topological closure, linking Paper 30's introduction with section 8. V polju kvocientov is therefore not dispensable. Familiarity of the stem zamknjen- supports the wording but does not replace the definition of integrality.\par
\noindent{\footnotesize Izvorne zapisy / Original records: HIST-g107-005, HIST-g115-001, HIST-g115-006.}\par
\end{minipage}\par
\par\medskip\Needspace{7\baselineskip}
\noindent\begin{minipage}{\linewidth}
\hypertarget{TERM-0638}{}
\noindent\textbf{integralna zamknutost o v Omega}\quad{\footnotesize\texttt{TERM-0638}}\par
\noindent\emph{Iztočnik / Source:} ganze Abgeschlossenheit von o in Omega\par
\noindent\emph{Rusky / Russian:} целая замкнутость o в Omega\par
\noindent\emph{Ukrajinsky / Ukrainian:} ціла замкненість o в Omega\par
\smallskip\noindent\emph{Razlog:} Integralna zamknutost o v Omega izbiraje integralna, da bi se izběglo obyčno čitanje cěla kako polna ili nepodělena. Razlog jest semantično razlikovanje, ne čisleno izměrjena prednost. Zapis sohranjaje variantu zamknutost vměsto zamknjenost i bukveny zapis Omega; bez sravnjenja s konkretnym tekstom ne treba samovoljno zaměnjati niti jednu iz tyh česti.\par
\noindent\emph{Reason:} Integralna zamknutost o v Omega selects integralna to avoid an ordinary reading of cěla as complete or undivided. The reason is semantic disambiguation, not a measured numerical advantage. The record retains zamknutost rather than zamknjenost and the letter-spelled Omega; neither should be silently changed without comparison with the particular text.\par
\noindent{\footnotesize Izvorne zapisy / Original records: HIST-g157-009.}\par
\end{minipage}\par
\par\medskip\Needspace{7\baselineskip}
\noindent\begin{minipage}{\linewidth}
\hypertarget{TERM-0639}{}
\noindent\textbf{c?la algebrai?na funkcija}\quad{\footnotesize\texttt{TERM-0639}}\par
\noindent\emph{Iztočnik / Source:} ganze algebraische Funktion\par
\noindent\emph{Rusky / Russian:} [Bukvy ne sut zachovane / Lettering not recoverable]\par
\noindent\emph{Ukrajinsky / Ukrainian:} [Bukvy ne sut zachovane / Lettering not recoverable]\par
\smallskip\noindent\emph{Razlog:} Zapis c?la algebrai?na funkcija ima nezachovane bukvy. Jej razlog izrično sohranjaje značenje cěloj ili integralnoj algebraičnoj funkcije i razlikuje go od racionalnoj funkcije i od uslovja cěločislenyh koeficientov. To dovoljuje objasniti naměru, ale ne vosstanoviti izgubljene bukvy kako svedočstvo. I rusko, i ukrajinsko sravnjenje ostajut čitatelju označene kako nezachovane.\par
\noindent\emph{Reason:} The record c?la algebrai?na funkcija has missing lettering. Its rationale explicitly retains whole or integral algebraic-function wording, distinguishing it from rational functions and from the integer-coefficient condition. That allows the intention to be explained but not lost letters to be restored as evidence. Both the Russian and Ukrainian comparanda remain identified as unrecoverable.\par
\noindent{\footnotesize Izvorne zapisy / Original records: HIST-g063-006.}\par
\end{minipage}\par
\par\medskip\Needspace{7\baselineskip}
\noindent\begin{minipage}{\linewidth}
\hypertarget{TERM-0640}{}
\noindent\textbf{cěle algebraične funkcije}\quad{\footnotesize\texttt{TERM-0640}}\par
\noindent\emph{Iztočnik / Source:} ganze algebraische Funktionen\par
\noindent\emph{Rusky / Russian:} целые алгебраические функции\par
\noindent\emph{Ukrajinsky / Ukrainian:} цілі алгебраїчні функції\par
\smallskip\noindent\emph{Razlog:} Cěle algebraične funkcije jest zachovana množinska forma. Dva razloga ne davajut odinakovogo ograničenja: zapis dlja odděla 3 raboty 9 upominaje lokalnu podnožnu primětku o cěločislenyh koeficientah, a zapis dlja odděla 9 razlikuje opću integralnu funkciju od cěločislenej formulacije. Oba svědočstva trěba sohraniti s njihovymi městami. Iz lokalnogo uslovja ne slěduje, že vsaka upotrěba termina znači funkciju s cěločislenymi koeficientami.\par
\noindent\emph{Reason:} Cěle algebraične funkcije is the preserved plural. The two rationales do not impose the same restriction: the record for Paper 9 section 3 mentions a local footnote concerning integer coefficients, whereas the section 9 record distinguishes general integral-function wording from integer-coefficient phrasing. Both testimonies need their locations. A local condition does not imply that every use of the term means a function with integer coefficients.\par
\noindent{\footnotesize Izvorne zapisy / Original records: HIST-g061-002, HIST-g067-007.}\par
\end{minipage}\par
\par\medskip\Needspace{7\baselineskip}
\noindent\begin{minipage}{\linewidth}
\hypertarget{TERM-0641}{}
\noindent\textbf{cěle celočislene funkcije}\quad{\footnotesize\texttt{TERM-0641}}\par
\noindent\emph{Iztočnik / Source:} ganze ganzzahlige Funktionen\par
\noindent\emph{Rusky / Russian:} целые целочисленные функции\par
\noindent\emph{Ukrajinsky / Ukrainian:} цілі цілочисельні функції\par
\smallskip\noindent\emph{Razlog:} Cěle celočislene funkcije sohranjaje dva kvalifikatora avtorskoj frazy. V razloženju rezultanta prvogo iz njih ne treba utožsamniti s uslovjem na koeficienty: polinomny oblik i cěločislenost koeficientov sut različne tvrđenja. Historijny razlog dopuskal daljše sravnjenje s sovrěmennym izrazom polinomy s cěločislenymi koeficientami, ale ne zapisuje, že ta zaměna byla přijeta. Aktualny pasaž pri formuli (5) odděla 11 ukazuje, od kakih prěměnnyh i koeficientov funkcije zavisęt.\par
\noindent\emph{Reason:} Cěle celočislene funkcije preserves both qualifiers in the author's phrase. In the resultant expansion, polynomial form and integrality of coefficients are different assertions. The historical rationale allowed further comparison with the modern wording polynomials with integer coefficients, but does not record adoption of that replacement. The current passage following equation (5) in section 11 specifies the variables and coefficients on which these functions depend.\par
\noindent{\footnotesize Izvorne zapisy / Original records: HIST-g051-008.}\par
\smallskip\noindent\emph{Sravnjenje s tekstom:} V odděle 11 raboty 6 koeficienty $A_i(\lambda)$ sut polinomne izrazy v $\lambda_1,\ldots,\lambda_\rho$ i v koeficientah vhodnyh polinomov. Sostavny nazv zato zadržuje i polinomny oblik, i uslovje na čislove koeficienty takogo izraza. Bez ukazanja oběh grup argumentov fraza cěle funkcije byla by nedostatočna.\par
\noindent\emph{Textual comparison:} In Paper 6 section 11, the coefficients $A_i(\lambda)$ are polynomial expressions in $\lambda_1,\ldots,\lambda_\rho$ and in the coefficients of the input polynomials. The compound designation therefore retains both polynomial form and a condition on the numerical coefficients of that expression. Whole functions would be insufficient without specifying these two groups of arguments.\par
\noindent{\footnotesize TEXT-COMP-25; pasaže / passages: 6178, 6181; \texttt{CORPUS\_FORM\_COMPARISONS.json}.}\par
\end{minipage}\par
\par\medskip\Needspace{7\baselineskip}
\noindent\begin{minipage}{\linewidth}
\hypertarget{TERM-0642}{}
\noindent\textbf{cěle celočislene simetrične funkcije}\quad{\footnotesize\texttt{TERM-0642}}\par
\noindent\emph{Iztočnik / Source:} ganze ganzzahlige symmetrische Funktionen\par
\noindent\emph{Rusky / Russian:} целые целочисленные симметрические функции\par
\noindent\emph{Ukrajinsky / Ukrainian:} цілі цілочисельні симетричні функції\par
\smallskip\noindent\emph{Razlog:} Cěle celočislene simetrične funkcije sohranjaje polinomny oblik, uslovje na koeficienty i simetriju. Nijeden iz tyh kvalifikatorov ne jest izlišny. V odděle 14 raboty 6 cěločislenost jest izrično oprěděljena koeficientami iz prstena algebraično cělyh čisel $[\Omega]$, ne samo običnymi cělymi čislami. Točno čitanje definicije niže objašnjava, začto blizky doslovny nazv potrěbuje matematično objasnjenje.\par
\noindent\emph{Reason:} Cěle celočislene simetrične funkcije preserves polynomial form, a coefficient condition and symmetry. None of these qualifiers is redundant. In Paper 6 section 14, the coefficient condition is explicitly defined using the ring $[\Omega]$ of algebraic integers, not just ordinary integers. The exact definition compared below shows why a close literal name still requires mathematical explanation.\par
\noindent{\footnotesize Izvorne zapisy / Original records: HIST-g054-011.}\par
\smallskip\noindent\emph{Sravnjenje s tekstom:} Odděl 14 raboty 6 oprěděljaje $[\Omega]$ kak algebraično cěle čisla danogo konečnogo čislovogo polja. Pomočny teorem govori o simetričnyh funkcijah od $n$ redov veličin; završny priměr odděla 15 izrično vrača se k tomu teoremu. Taka svez podkrepljaje sostavno čitanje dvojnoj kvalifikacije pri TERM-0642 i TERM-0643.\par
\noindent\emph{Textual comparison:} Paper 6 section 14 defines $[\Omega]$ as the algebraic integers in the given finite number field. Its auxiliary theorem concerns symmetric functions of $n$ rows of quantities; the final example in section 15 expressly refers back to that theorem. This connection supports reading the double qualification compositionally in TERM-0642 and TERM-0643.\par
\noindent{\footnotesize TEXT-COMP-26; pasaže / passages: 6438, 6460, 6594; \texttt{CORPUS\_FORM\_COMPARISONS.json}.}\par
\end{minipage}\par
\par\medskip\Needspace{7\baselineskip}
\noindent\begin{minipage}{\linewidth}
\hypertarget{TERM-0643}{}
\noindent\textbf{cěle celočislene simetrične funkcije od $n$ redov veličin}\quad{\footnotesize\texttt{TERM-0643}}\par
\noindent\emph{Iztočnik / Source:} ganze ganzzahlige symmetrische Funktionen von $n$ Größenreihen\par
\noindent\emph{Rusky / Russian:} целые целочисленные симметрические функции от $n$ рядов величин\par
\noindent\emph{Ukrajinsky / Ukrainian:} цілі цілочисельні симетричні функції від $n$ рядів величин\par
\smallskip\noindent\emph{Razlog:} Fraza cěle celočislene simetrične funkcije od $n$ redov veličin dodava k TERM-0642 točny obseg argumentov. Čislo $n$ označaje redy, ne stepen funkcije. Razlog svęzuje završny priměr raboty 6 s pomočnym teoremom odděla 14; sohranjenje dvojnogo kvalifikatora čini tu svez vidimoju. Sravnjenje s definicijeju $[\Omega]$ pri TERM-0642 ostavaje potrěbno, da celočislene ne bylo bezosnovno suženo do racionalno cělyh koeficientov.\par
\noindent\emph{Reason:} The phrase cěle celočislene simetrične funkcije od $n$ redov veličin adds the precise range of arguments to TERM-0642. Here $n$ counts rows, not the function's degree. The rationale links Paper 6's final example to the auxiliary theorem of section 14; retaining both qualifiers makes that relation visible. The definition of $[\Omega]$ discussed at TERM-0642 is needed to avoid an unwarranted restriction to rational-integer coefficients.\par
\noindent{\footnotesize Izvorne zapisy / Original records: HIST-g055-010.}\par
\end{minipage}\par
\par\medskip\Needspace{7\baselineskip}
\noindent\begin{minipage}{\linewidth}
\hypertarget{TERM-0644}{}
\noindent\textbf{cěle veličiny}\quad{\footnotesize\texttt{TERM-0644}}\par
\noindent\emph{Iztočnik / Source:} ganze Größen\par
\noindent\emph{Rusky / Russian:} целые величины\par
\noindent\emph{Ukrajinsky / Ukrainian:} цілі величини\par
\smallskip\noindent\emph{Razlog:} Cěle veličiny sohranjaje obče ime dlja dvoh koeficientnyh oblastij v začetku raboty 22: cělyh racionalnyh čisel i polinomov jednoj prěměnnoj nad poljem. Veličiny zato ne znači samo čisla. Razlog ne želi prědčasno zaměniti avtorsku konstrukciju obćim modernym nazvom elementy prstena. Točna definicija niže pokazuje obě možnosti; prodolženje generalizacije v teoremu III ne treba pripisati samomu slovničnomu slovu.\par
\noindent\emph{Reason:} Cěle veličiny retains a common name for the two coefficient domains at the beginning of Paper 22: rational integers and one-variable polynomials over a field. Veličiny therefore does not mean numbers alone. The rationale avoids prematurely replacing the author's construction with the general modern name ring elements. The exact opening definition compared below specifies both possibilities; the later generalization in Theorem III cannot be inferred from the dictionary word alone.\par
\noindent{\footnotesize Izvorne zapisy / Original records: HIST-g085-001, HIST-g092-001.}\par
\smallskip\noindent\emph{Sravnjenje s tekstom:} V prvyh dvoh oddělah raboty 22 cěle veličiny sut ili racionalne cěle čisla, ili polinomy jednoj neoprěděljenej nad libovoljnym poljem $P$. Opis modula koristi tu istu oblast dlja množenja linearnyh form. Izraz cěle veličiny imaje zato širši obseg od čisla; koeficienty polinomov ne sut samym nazvom ograničene do cělyh čisel.\par
\noindent\emph{Textual comparison:} In the first two sections of Paper 22, integral quantities are either rational integers or polynomials in one indeterminate over an arbitrary field $P$. The module definition uses the same domain for scalar multiplication of linear forms. Integral quantities therefore has a wider range than numbers; the polynomial coefficients are not restricted to integers by that designation alone.\par
\noindent{\footnotesize TEXT-COMP-27; pasaže / passages: 14226, 14228; \texttt{CORPUS\_FORM\_COMPARISONS.json}.}\par
\end{minipage}\par
\par\medskip\Needspace{7\baselineskip}
\noindent\begin{minipage}{\linewidth}
\hypertarget{TERM-0645}{}
\noindent\textbf{cěle veličiny / cěle racionalne čisla}\quad{\footnotesize\texttt{TERM-0645}}\par
\noindent\emph{Iztočnik / Source:} ganze Größen / ganze rationale Zahlen\par
\noindent\emph{Rusky / Russian:} целые величины / целые рациональные числа\par
\noindent\emph{Ukrajinsky / Ukrainian:} цілі величини / цілі раціональні числа\par
\smallskip\noindent\emph{Razlog:} Cěle veličiny / cěle racionalne čisla sohranjaje dva povezana, ale ne opće totožna izraza. Zapis motivuje jih avtorskim jezykom koeficientov i sovrěmennoju čitateljnostju. Otvarjajuča definicija raboty 22, navedena pri TERM-0644, uključaje v cěle veličiny takože polinomy. Zato zajedni spis ne dopuščaje svesti obě frazy do običnyh cělyh čisel; pri konkretnom tvrđenju treba ustanoviti, ktora koeficientna oblast jest v sile.\par
\noindent\emph{Reason:} Cěle veličiny / cěle racionalne čisla preserves two related but not generally identical expressions. The record motivates them by the author's coefficient language and present-day readability. Paper 22's opening definition, discussed at TERM-0644, includes polynomials among cěle veličiny. The grouped list therefore does not reduce both phrases to ordinary integers; the coefficient domain in force must be identified in each statement.\par
\noindent{\footnotesize Izvorne zapisy / Original records: HIST-g088-007, HIST-g092-038.}\par
\end{minipage}\par
\par\medskip\Needspace{7\baselineskip}
\noindent\begin{minipage}{\linewidth}
\hypertarget{TERM-0646}{}
\noindent\textbf{cěle veličiny libovoljnogo polja}\quad{\footnotesize\texttt{TERM-0646}}\par
\noindent\emph{Iztočnik / Source:} ganze Größen eines beliebigen Körpers\par
\noindent\emph{Rusky / Russian:} целые величины произвольного поля\par
\noindent\emph{Ukrajinsky / Ukrainian:} цілі величини довільного поля\par
\smallskip\noindent\emph{Razlog:} Cěle veličiny libovoljnogo polja sohranjaje historičny imenik veličiny pri prěhodu od kompleksnyh čisel k libovoljnym abstraktnym poljam. Razlog daje prednost toj kontinuitetnosti prěd zaměnoju zaglavja na elementy. Libovoljnogo jest sučstveno dlja obsega razsmotrenja. Cělost tu mora byti čitana po definiciji raboty 9; sam pridavnik ne može značiti, že vsak element vsakego polja jest cěly nad libovoljnoju osnovoju.\par
\noindent\emph{Reason:} Cěle veličiny libovoljnogo polja preserves the historical noun veličiny when passing from complex numbers to arbitrary abstract fields. The rationale prefers that continuity to recasting the heading in terms of elements. Libovoljnogo is essential to the scope. Integrality must be read according to Paper 9's definition; the adjective alone cannot mean that every element of every field is integral over any base.\par
\noindent{\footnotesize Izvorne zapisy / Original records: HIST-g068-001.}\par
\end{minipage}\par
\par\medskip\Needspace{7\baselineskip}
\noindent\begin{minipage}{\linewidth}
\hypertarget{TERM-0647}{}
\noindent\textbf{cěle invariantne tvorjenja}\quad{\footnotesize\texttt{TERM-0647}}\par
\noindent\emph{Iztočnik / Source:} ganze invariante Bildungen\par
\noindent\emph{Rusky / Russian:} целые инвариантные образования\par
\noindent\emph{Ukrajinsky / Ukrainian:} цілі інваріантні утворення\par
\smallskip\noindent\emph{Razlog:} Cěle invariantne tvorjenja zadržuje široky smysl německogo Bildungen: invariantne konstrukcije ili objekty. Razlog ne želi bez dodatnogo osnovanja suziti jih do jednoj podklasy. Cěle sohranjaje algebraičny kvalifikator, ale ne govori o analitičnom integralu. Dlja točnogo čitanja cělosti potrěben jest lokalny matematičny kontekst, a ne jedina doslovna zaměna slova ganz.\par
\noindent\emph{Reason:} Cěle invariantne tvorjenja retains the breadth of German Bildungen: invariant constructions or objects. The rationale avoids narrowing them to one subclass without further justification. Cěle preserves the algebraic qualifier; it does not refer to an analytic integral. Its precise interpretation requires the local mathematical context, not a single literal replacement for ganz.\par
\noindent{\footnotesize Izvorne zapisy / Original records: HIST-g029-007.}\par
\end{minipage}\par
\par\medskip\Needspace{7\baselineskip}
\noindent\begin{minipage}{\linewidth}
\hypertarget{TERM-0648}{}
\noindent\textbf{cěly racionalny (absolutny) invariant}\quad{\footnotesize\texttt{TERM-0648}}\par
\noindent\emph{Iztočnik / Source:} ganze rationale (absolute) Invariante\par
\noindent\emph{Rusky / Russian:} целый рациональный (абсолютный) инвариант\par
\noindent\emph{Ukrajinsky / Ukrainian:} цілий раціональний (абсолютний) інваріант\par
\smallskip\noindent\emph{Razlog:} Cěly racionalny (absolutny) invariant sohranjaje cělu historijnu frazu i vstavku absolutny. Razlog izrično razlikuje polinomny smysl iz ranějšej invariantnoj teorije od uslovja na cěločislenost koeficientov v odděle 14 raboty 6. Skobky ne daju prava izpustiti kvalifikator, kogda on razlikuje razsmatrjany tip invarianta. Tuto jest kontinuitetnost pojma, ne novo ime za vsaky racionalny invariant.\par
\noindent\emph{Reason:} Cěly racionalny (absolutny) invariant preserves the historical phrase and its parenthetical absolute qualifier. The rationale expressly distinguishes the polynomial sense of the earlier invariant theory from section 14's coefficient-integrality condition in Paper 6. Parentheses do not license omission of a qualifier that distinguishes the relevant type of invariant. This continues a particular concept, not a new name for every rational invariant.\par
\noindent{\footnotesize Izvorne zapisy / Original records: HIST-g056-003.}\par
\end{minipage}\par
\par\medskip\Needspace{7\baselineskip}
\noindent\begin{minipage}{\linewidth}
\hypertarget{TERM-0649}{}
\noindent\textbf{cělo racionalno predstavjanje}\quad{\footnotesize\texttt{TERM-0649}}\par
\noindent\emph{Iztočnik / Source:} ganze rationale Darstellbarkeit\par
\noindent\emph{Rusky / Russian:} целая рациональная представимость\par
\noindent\emph{Ukrajinsky / Ukrainian:} ціла раціональна зображуваність\par
\smallskip\noindent\emph{Razlog:} Cělo racionalno predstavjanje imaje vyraziti možnost polinomnogo predstavjenja, odvojenu od obće racionalnoj predstavimosti. Německo Darstellbarkeit govori o možnosti, ne samo o činu predstavjanja. Historijny razlog objašnjava matematično razlikovanje, ale ne argumentuje osobno, kako izbrany imenik sohranjaje taj modalny smysl. Slovnik zato mora pokazati i naměru, i gramatyčno pytanje, a ne proglasiti ego už rěšenym.\par
\noindent\emph{Reason:} Cělo racionalno predstavjanje is intended to express polynomial representability, distinguished from general rational representability. German Darstellbarkeit denotes a possibility, not merely the act of presenting. The historical rationale explains the mathematical distinction but does not separately argue how its chosen noun preserves that modal sense. The dictionary therefore needs both the intention and the unresolved grammatical question.\par
\noindent{\footnotesize Izvorne zapisy / Original records: HIST-g040-011.}\par
\end{minipage}\par
\par\medskip\Needspace{7\baselineskip}
\noindent\begin{minipage}{\linewidth}
\hypertarget{TERM-0650}{}
\noindent\textbf{cělo racionalno predstavjenje}\quad{\footnotesize\texttt{TERM-0650}}\par
\noindent\emph{Iztočnik / Source:} ganze rationale Darstellung\par
\noindent\emph{Rusky / Russian:} целое рациональное представление\par
\noindent\emph{Ukrajinsky / Ukrainian:} ціле раціональне подання\par
\smallskip\noindent\emph{Razlog:} Cělo racionalno predstavjenje sohranjaje invariantno-teoretičny smysl izloženja prěz polinomny izraz. Razlog sravnjaje rusko представление s ukrajinskim подання i izbiraje predstavjenje. Cěly davaje slovjansku, racionalny internacionalnu čest frazy; te dvě česti ne sut sinonimy. Izbor imenika razlikuje se od pytanja o samoj možnosti predstavjenja v TERM-0649.\par
\noindent\emph{Reason:} Cělo racionalno predstavjenje retains the invariant-theoretic sense of expression by a polynomial. The rationale compares Russian представление with Ukrainian подання and selects predstavjenje. The phrase combines Slavic cěly with international racionalny; the two components are not synonyms. Choosing this noun is distinct from expressing the possibility of representation in TERM-0649.\par
\noindent{\footnotesize Izvorne zapisy / Original records: HIST-g057-001.}\par
\end{minipage}\par
\par\medskip\Needspace{7\baselineskip}
\noindent\begin{minipage}{\linewidth}
\hypertarget{TERM-0651}{}
\noindent\textbf{cěla racionalna funkcija}\quad{\footnotesize\texttt{TERM-0651}}\par
\noindent\emph{Iztočnik / Source:} ganze rationale Funktion\par
\noindent\emph{Rusky / Russian:} целая рациональная функция\par
\noindent\emph{Ukrajinsky / Ukrainian:} ціла раціональна функція\par
\smallskip\noindent\emph{Razlog:} Tri zapisy za cěla racionalna funkcija svęzujut frazu s rezultantami, invariantami i algebraičnoju cělostju. Jedin iz njih izrično odbija integralna funkcija, da ne směšati polinomny oblik s odvojenym pytanjem cělosti nad oblastju. Drugi zadržuje historijny jezyk pri razsmotrenju cělosti. To sut kontekstualne motivacije za jednaku frazu, ne dokaz, že vsako značenje slova integralny jest isto. Sravnjenje pri TERM-0653 davaje javny polinomny smysl.\par
\noindent\emph{Reason:} Three records for cěla racionalna funkcija connect the phrase with resultants, invariants and algebraic integrality. One expressly rejects integralna funkcija so that polynomial form is not confused with the separate question of integrality over a domain. Another retains the historical wording in a discussion of integrality. These are contextual motivations for the same phrase, not evidence that every sense of integral is identical. TERM-0653 gives the explicit polynomial reading.\par
\noindent{\footnotesize Izvorne zapisy / Original records: HIST-g051-007, HIST-g077-002, HIST-g094-008.}\par
\end{minipage}\par
\par\medskip\Needspace{7\baselineskip}
\noindent\begin{minipage}{\linewidth}
\hypertarget{TERM-0652}{}
\noindent\textbf{cěle racionalne funkcije}\quad{\footnotesize\texttt{TERM-0652}}\par
\noindent\emph{Iztočnik / Source:} ganze rationale Funktionen\par
\noindent\emph{Rusky / Russian:} целые рациональные функции\par
\noindent\emph{Ukrajinsky / Ukrainian:} цілі раціональні функції\par
\smallskip\noindent\emph{Razlog:} Cěle racionalne funkcije jest množinska forma historijnogo nazva. Razlogy iz rabot 6 i 10 sohranjajut razlikovanje medžu polinomami i racionalnymi izrazami s imenovateljami, vměsto avtomatičnoj zaměny vsih pojavjenij slovom polinomy. To ne skriva polinomny smysl: avtorsko objasnjenje ostavaje vidimo v TERM-0653 i TERM-0654. Povtorjena terminologija pomaga slěditi obsegy uslovij v različnyh oddělah.\par
\noindent\emph{Reason:} Cěle racionalne funkcije is the historical designation in the plural. The rationales from Papers 6 and 10 preserve the contrast between polynomials and rational expressions with denominators instead of automatically replacing every occurrence by polynomials. This does not conceal the polynomial meaning: the author's explicit explanation is retained in TERM-0653 and TERM-0654. Consistent terminology helps readers follow the scope of conditions across sections.\par
\noindent{\footnotesize Izvorne zapisy / Original records: HIST-g040-003, HIST-g047-005, HIST-g071-005.}\par
\end{minipage}\par
\par\medskip\Needspace{7\baselineskip}
\noindent\begin{minipage}{\linewidth}
\hypertarget{TERM-0653}{}
\noindent\textbf{cěle racionalne funkcije (polinomy)}\quad{\footnotesize\texttt{TERM-0653}}\par
\noindent\emph{Iztočnik / Source:} ganze rationale Funktionen (Polynome)\par
\noindent\emph{Rusky / Russian:} целые рациональные функции (полиномы)\par
\noindent\emph{Ukrajinsky / Ukrainian:} цілі раціональні функції (поліноми)\par
\smallskip\noindent\emph{Razlog:} Cěle racionalne funkcije (polinomy) sohranjaje i historijny nazv, i avtorsko objasnjenje v skobkah. Razlog navodi, že to ravenstvo značenij jest izrično v odděle 1 raboty 6. Aktualny tekst to podkrepljaje: polinomy sut različene od racionalnyh funkcij zapisanyh s čisliteljem i imenovateljem. Zato objasnjenje ne jest nezapisana moderna zaměna vvedena prevodnikom.\par
\noindent\emph{Reason:} Cěle racionalne funkcije (polinomy) retains both the historical name and the author's parenthetical explanation. The rationale says that section 1 of Paper 6 explicitly equates the two. The current text supports this: polynomials are distinguished from rational functions represented by numerator and denominator. The explanation is therefore not an undocumented modern substitution introduced by the translator.\par
\noindent{\footnotesize Izvorne zapisy / Original records: HIST-g041-005.}\par
\smallskip\noindent\emph{Sravnjenje s tekstom:} V odděle 1 raboty 6 fraza cěle racionalne funkcije jest bezposrědno dopolnjena slovom polinomy v skobkah. Odděl 8 povtorjaje to s to jest. Sravnjenje pokazuje, že dva slovnične zapisy TERM-0653 i TERM-0654 sohranjajut objasnjenje iz izloženja, ne vymysljaju novu redakcijnu definiciju.\par
\noindent\emph{Textual comparison:} In Paper 6 section 1, whole rational functions is immediately followed by polynomials in parentheses. Section 8 repeats the equivalence with that is. The comparison shows that TERM-0653 and TERM-0654 retain an explanation in the exposition rather than inventing a new editorial definition.\par
\noindent{\footnotesize TEXT-COMP-28; pasaže / passages: 5061, 5857; \texttt{CORPUS\_FORM\_COMPARISONS.json}.}\par
\end{minipage}\par
\par\medskip\Needspace{7\baselineskip}
\noindent\begin{minipage}{\linewidth}
\hypertarget{TERM-0654}{}
\noindent\textbf{cěle racionalne funkcije, to jest polinomy}\quad{\footnotesize\texttt{TERM-0654}}\par
\noindent\emph{Iztočnik / Source:} ganze rationale Funktionen (Polynome)\par
\noindent\emph{Rusky / Russian:} целые рациональные функции, то есть полиномы\par
\noindent\emph{Ukrajinsky / Ukrainian:} цілі раціональні функції, тобто поліноми\par
\smallskip\noindent\emph{Razlog:} Cěle racionalne funkcije, to jest polinomy razvivaje avtorsko objasnjenje v javnu frazu s to jest. Razlog govori o odděle 8, kde polinomny smysl jest čestju samogo izloženja. To jest rozděl v predstavjenju definicije, ne novy dopolnitelny matematičny smysl. Sravnjenje pri TERM-0653 vključaje i to město, tako že čitatelj može razlikovati avtorsko objasnjenje od osobnoj redakcijnoj primětky.\par
\noindent\emph{Reason:} Cěle racionalne funkcije, to jest polinomy expands the author's gloss into an explicit to jest phrase. The rationale concerns section 8, where the polynomial sense is part of the exposition itself. This changes how the definition is expressed, not its mathematical meaning. The comparison at TERM-0653 includes this passage so readers can distinguish an authorial gloss from a separate editorial note.\par
\noindent{\footnotesize Izvorne zapisy / Original records: HIST-g048-002.}\par
\end{minipage}\par
\par\medskip\Needspace{7\baselineskip}
\noindent\begin{minipage}{\linewidth}
\hypertarget{TERM-0655}{}
\noindent\textbf{cěle racionalne projektivne invarianty kvadratičnoj formy}\quad{\footnotesize\texttt{TERM-0655}}\par
\noindent\emph{Iztočnik / Source:} ganze rationale projektive Invarianten einer quadratischen Form\par
\noindent\emph{Rusky / Russian:} целые рациональные проективные инварианты квадратичной формы\par
\noindent\emph{Ukrajinsky / Ukrainian:} цілі раціональні проєктивні інваріанти квадратичної форми\par
\smallskip\noindent\emph{Razlog:} Cěle racionalne projektivne invarianty kvadratičnoj formy sohranjaje vse tri smyslove ograničenja: polinomny tip, projektivnu invariantnost i kvadratičnu formu. Genitiv kvadratičnoj formy izrično svęzuje invarianty s objektom. Kratky razlog ne predlagaje kračšego konkurentnogo nazva; zato slovnik ne treba prikazovati nezapisany izbor medžu alternativami. Dluga fraza jest smyslovo sostavjena, ne spis sinonimov.\par
\noindent\emph{Reason:} Cěle racionalne projektivne invarianty kvadratičnoj formy retains three restrictions: polynomial type, projective invariance and a quadratic form. The genitive kvadratičnoj formy explicitly connects the invariants to their object. The brief rationale proposes no competing shorter name, so the dictionary must not invent a comparison between alternatives. The long phrase is compositionally meaningful, not a list of synonyms.\par
\noindent{\footnotesize Izvorne zapisy / Original records: HIST-g053-008.}\par
\end{minipage}\par
\par\medskip\Needspace{7\baselineskip}
\noindent\begin{minipage}{\linewidth}
\hypertarget{TERM-0656}{}
\noindent\textbf{cěla racionalna kombinacija}\quad{\footnotesize\texttt{TERM-0656}}\par
\noindent\emph{Iztočnik / Source:} ganze rationale Verbindung\par
\noindent\emph{Rusky / Russian:} целая рациональная комбинация\par
\noindent\emph{Ukrajinsky / Ukrainian:} ціла раціональна комбінація\par
\smallskip\noindent\emph{Razlog:} Cěla racionalna kombinacija znači polinomnu kombinaciju, ne prosto linearnu kombinaciju i ne libovoljny racionalny izraz. Dva razloga svęzujut odděly 7 i 8 raboty 6 i podkrepljajut, že polinomny smysl stavaje izričny v daljšem odděle. Kombinacija zadržuje tu svez v dokazě. Izbor ne znači, že odstranjenje imenovateljev možno izpustiti v konkretnoj formuli.\par
\noindent\emph{Reason:} Cěla racionalna kombinacija means a polynomial combination, not merely a linear combination or an arbitrary rational expression. Two rationales link sections 7 and 8 of Paper 6, emphasizing the explicit polynomial reading in the later section. Kombinacija maintains that connection through the proof. The choice does not license omitting the question of denominators in an actual formula.\par
\noindent{\footnotesize Izvorne zapisy / Original records: HIST-g047-009, HIST-g048-015.}\par
\end{minipage}\par
\par\medskip\Needspace{7\baselineskip}
\noindent\begin{minipage}{\linewidth}
\hypertarget{TERM-0657}{}
\noindent\textbf{cělo racionalno čislo}\quad{\footnotesize\texttt{TERM-0657}}\par
\noindent\emph{Iztočnik / Source:} ganze rationale Zahl\par
\noindent\emph{Rusky / Russian:} целое рациональное число\par
\noindent\emph{Ukrajinsky / Ukrainian:} ціле раціональне число\par
\smallskip\noindent\emph{Razlog:} Cělo racionalno čislo zadržuje historijny nazv racionalnogo cělogo čisla. Racionalno tu razlikuje obične cěle čisla od širšego konteksta algebraično cělyh čisel; ne znači proizvoljny drobek. Razlog odbija skračenu modernu zaměnu, ktora by skryla tu avtorsku frazu. Čislo se pri tom ne smi směšati s cělym racionalnym funkcionalom, kotory imaje osobnu definiciju pri TERM-0662.\par
\noindent\emph{Reason:} Cělo racionalno čislo retains the historical designation rational integer. Racionalno distinguishes ordinary integers from the wider algebraic-integer context; it does not mean an arbitrary fraction. The rationale rejects a modern abbreviation that would hide the author's phrase. The number must also be distinguished from a whole rational functional, which has its own definition at TERM-0662.\par
\noindent{\footnotesize Izvorne zapisy / Original records: HIST-g060-004.}\par
\end{minipage}\par
\par\medskip\Needspace{7\baselineskip}
\noindent\begin{minipage}{\linewidth}
\hypertarget{TERM-0658}{}
\noindent\textbf{cěle racionalne čisla}\quad{\footnotesize\texttt{TERM-0658}}\par
\noindent\emph{Iztočnik / Source:} ganze rationale Zahlen\par
\noindent\emph{Rusky / Russian:} целые рациональные числа\par
\noindent\emph{Ukrajinsky / Ukrainian:} цілі раціональні числа\par
\smallskip\noindent\emph{Razlog:} Cěle racionalne čisla sohranjaje množinu i historijny imenik. Zachovany razlog svęzuje upotrěbu s cěločislenymi koeficientami v modulnoj identičnosti raboty 22. Dva povezana zapisy ne sut dva razne dokazy togo čitanja. Frazu trěba priměnjati k imenovanym koeficientam, ne prěnesti bez različenja na vsaku oblast veličin v celoj rabotě; sr. širšu definiciju pri TERM-0644.\par
\noindent\emph{Reason:} Cěle racionalne čisla preserves the plural and historical designation. The surviving rationale ties its use to integer coefficients in a module identity in Paper 22. Two associated records are not two separate proofs of that reading. The phrase must apply to the specified coefficients, not indiscriminately to every domain of quantities in the paper; see the broader definition at TERM-0644.\par
\noindent{\footnotesize Izvorne zapisy / Original records: HIST-g087-011, HIST-g092-031.}\par
\end{minipage}\par
\par\medskip\Needspace{7\baselineskip}
\noindent\begin{minipage}{\linewidth}
\hypertarget{TERM-0659}{}
\noindent\textbf{cěle racionalne celočislene kombinacije}\quad{\footnotesize\texttt{TERM-0659}}\par
\noindent\emph{Iztočnik / Source:} ganze rationale, ganzzahlige Verbindungen\par
\noindent\emph{Rusky / Russian:} целые рациональные целочисленные комбинации\par
\noindent\emph{Ukrajinsky / Ukrainian:} цілі раціональні цілочисельні комбінації\par
\smallskip\noindent\emph{Razlog:} Cěle racionalne celočislene kombinacije sohranjaje odvojeno polinomnu kombinaciju i ograničenje njej koeficientov. Razlog ukazuje na odděl 14 raboty 6, ktory dodava prsten $[\Omega]$ k ranějšim konstrukcijam. Aktualny pasaž razlikuje takve kombinacije od obćih racionalnyh kombinacij; obě formulacije zato ne možno sojediniti radi kračšego teksta. Definicija koeficientnogo prstena dana jest pri TERM-0642.\par
\noindent\emph{Reason:} Cěle racionalne celočislene kombinacije separately retains polynomial combination and the restriction on its coefficients. The rationale points to Paper 6 section 14, which introduces the ring $[\Omega]$ into earlier constructions. The current passage distinguishes these combinations from general rational combinations; the two formulations cannot be merged merely for brevity. The coefficient ring is defined at TERM-0642.\par
\noindent{\footnotesize Izvorne zapisy / Original records: HIST-g054-008.}\par
\smallskip\noindent\emph{Sravnjenje s tekstom:} Dva susědne pasaže odděla 14 razlikujut najmenšu sodrživajuču integralnu oblast od najmenšej relativno cěloj oblasti. Prva koristi polinomne kombinacije s cělymi koeficientami; druga dopuščaje racionalne kombinacije, čiji rezultat ostavaje cěločisleny polinom. Izpustiti cěle racionalne v prvom pasažu značilo by izgubiti to razlikovanje.\par
\noindent\emph{Textual comparison:} Two adjacent passages in section 14 distinguish the smallest containing integral domain from the smallest relatively integral domain. The first uses polynomial combinations with integral coefficients; the second permits rational combinations whose result is still an integral-coefficient polynomial. Omitting whole rational from the first passage would lose that distinction.\par
\noindent{\footnotesize TEXT-COMP-29; pasaže / passages: 6450, 6452; \texttt{CORPUS\_FORM\_COMPARISONS.json}.}\par
\end{minipage}\par
\par\medskip\Needspace{7\baselineskip}
\noindent\begin{minipage}{\linewidth}
\hypertarget{TERM-0660}{}
\noindent\textbf{cěla simetrična funkcija}\quad{\footnotesize\texttt{TERM-0660}}\par
\noindent\emph{Iztočnik / Source:} ganze symmetrische Funktion\par
\noindent\emph{Rusky / Russian:} целая симметрическая функция\par
\noindent\emph{Ukrajinsky / Ukrainian:} ціла симетрична функція\par
\smallskip\noindent\emph{Razlog:} Cěla simetrična funkcija zadržuje polinomny smysl v teoremu o simetričnyh funkcijah. Razlog povęzuje go s cěla racionalna kombinacija, a ne s analitičnoju cěloju funkcijeju. Simetrija ostavaje nezavisnym uslovjem na objekt. Kratky zapis ne opravdava dodavanje cěločislenosti koeficientov, ktora v drugyh městah jest označena dodatnym pridavnikom.\par
\noindent\emph{Reason:} Cěla simetrična funkcija retains the polynomial sense in the symmetric-function theorem. The rationale connects it with cěla racionalna kombinacija, not with an entire function of complex analysis. Symmetry remains a separate condition on the object. The brief record does not justify adding a coefficient-integrality condition that other passages mark with an additional adjective.\par
\noindent{\footnotesize Izvorne zapisy / Original records: HIST-g048-014.}\par
\end{minipage}\par
\par\medskip\Needspace{7\baselineskip}
\noindent\begin{minipage}{\linewidth}
\hypertarget{TERM-0661}{}
\noindent\textbf{cěle transcendentne čisla}\quad{\footnotesize\texttt{TERM-0661}}\par
\noindent\emph{Iztočnik / Source:} ganze transzendente Zahlen\par
\noindent\emph{Rusky / Russian:} целые трансцендентные числа\par
\noindent\emph{Ukrajinsky / Ukrainian:} цілі трансцендентні числа\par
\smallskip\noindent\emph{Razlog:} Cěle transcendentne čisla sohranjaje ime Zermelovoj konstrukcije, ne običnu klasifikaciju čisel na cěle i transcendentne. Doslovny anglijski izraz integer transcendental numbers byl by zavodlivy: racionalne cěle čisla ne sut transcendentne. V rabotě 9 idě o posebno oblast s oprěděljenymi svojstvami zamknjenosti, ktora sodrživaje algebraično cěle čisla i ne sodrživaje algebraične čisla, ktore ne sut cěle. Razlog daje prednost sohranjenju historijnogo nazva s objasnjenjem pred tihoju modernizacijeju.\par
\noindent\emph{Reason:} Cěle transcendentne čisla preserves the name of Zermelo's construction, not the ordinary classification of numbers as integers or transcendentals. The English rendering integer transcendental numbers would be misleading: rational integers are not transcendental. Paper 9 concerns a particular domain with specified closure properties, containing algebraic integers and excluding nonintegral algebraic numbers. The rationale favours retaining the historical name with an explanation over silently modernizing it.\par
\noindent{\footnotesize Izvorne zapisy / Original records: HIST-g058-002.}\par
\smallskip\noindent\emph{Sravnjenje s tekstom:} Uvod raboty 9 davaje četyri abstraktna svojstva oblasti $\mathfrak G$: zamknjenost pri sumě, razlici i produktu; predstavimost vsakogo kompleksnogo čisla kvocientom; sodržanje algebraično cělyh čisel; i izključenje necělyh algebraičnyh čisel. To objašnjavaje, začto cěle transcendentne čisla trěba čitati kak historijny nazv konstrukcije. Slovo cěle ne jest obična tvrdnja o racionalnoj cěločislenosti vsih jej elementov.\par
\noindent\emph{Textual comparison:} Paper 9's introduction gives four abstract properties of $\mathfrak G$: closure under addition, subtraction and multiplication; representation of every complex number as a quotient; inclusion of algebraic integers; and exclusion of nonintegral algebraic numbers. This explains why whole transcendental numbers must be read as the historical name of a construction. Whole is not an ordinary assertion that all its elements are rational integers.\par
\noindent{\footnotesize TEXT-COMP-30; pasaže / passages: 7134, 7136, 7137, 7138, 7139; \texttt{CORPUS\_FORM\_COMPARISONS.json}.}\par
\end{minipage}\par
\par\medskip\Needspace{7\baselineskip}
\noindent\begin{minipage}{\linewidth}
\hypertarget{TERM-0662}{}
\noindent\textbf{cěly racionalny funkcional}\quad{\footnotesize\texttt{TERM-0662}}\par
\noindent\emph{Iztočnik / Source:} ganzes rationales Funktional\par
\noindent\emph{Rusky / Russian:} целый рациональный функционал\par
\noindent\emph{Ukrajinsky / Ukrainian:} цілий раціональний функціонал\par
\smallskip\noindent\emph{Razlog:} Cěly racionalny funkcional mora byti čitany po definiciji raboty 9, ne po sovrěmennom značeniu funkcionala v funkcionalnoj analizi. Racionalna funkcija jest zapisana kako $A=aE_1/E_2$, kde polinomy sut primitivne i vzajemno proste. Cělost tu znači, že pozitivny racionalny faktor $a$ jest cělo čislo; ona ne potrěbuje, že imenovatelj $E_2$ jest jedinica. Historijny razlog pravilno razlikuje to svojstvo od cěločislenosti koeficientov, ale polna definicija jest neophodima dlja točnogo razuměnja.\par
\noindent\emph{Reason:} Cěly racionalny funkcional must be read according to Paper 9's definition, not the modern use of functional in functional analysis. A rational function is written as $A=aE_1/E_2$, with primitive, relatively prime polynomials. Here being whole means that the positive rational factor $a$ is an integer; it does not require the denominator $E_2$ to be one. The historical rationale rightly distinguishes this property from coefficient integrality, but the full definition is essential to understanding it.\par
\noindent{\footnotesize Izvorne zapisy / Original records: HIST-g060-002.}\par
\smallskip\noindent\emph{Sravnjenje s tekstom:} Definicija pripisana Veberu v odděle 2 raboty 9 koristi normalizovano predstavjenje $A=aE_1/E_2$. Polinomy $E_1,E_2$ sut primitivne, vzajemno proste i imajut cěle racionalne koeficienty. Cělost funkcionala jest uslovje na faktor $a$, ne uslovje $E_2=1$. Zato fraza ne može byti avtomatično zaměnjena slovom polinom.\par
\noindent\emph{Textual comparison:} The definition attributed to Weber in Paper 9 section 2 uses the normalized representation $A=aE_1/E_2$. The polynomials $E_1,E_2$ are primitive and relatively prime, with rational-integer coefficients. The functional is whole when the factor $a$ is an integer, not when $E_2=1$. The phrase therefore cannot automatically be replaced by polynomial.\par
\noindent{\footnotesize TEXT-COMP-31; pasaže / passages: 7223, 7226, 7228; \texttt{CORPUS\_FORM\_COMPARISONS.json}.}\par
\end{minipage}\par
\par\medskip\Needspace{7\baselineskip}
\noindent\begin{minipage}{\linewidth}
\hypertarget{TERM-0663}{}
\noindent\textbf{cělo-racionalne čislove koeficienty}\quad{\footnotesize\texttt{TERM-0663}}\par
\noindent\emph{Iztočnik / Source:} ganzrationale Zahlkoeffizienten\par
\noindent\emph{Rusky / Russian:} целые рациональные числовые коэффициенты\par
\noindent\emph{Ukrajinsky / Ukrainian:} цілі раціональні числові коефіцієнти\par
\smallskip\noindent\emph{Razlog:} Cělo-racionalne čislove koeficienty zadržuje složenu frazu za racionalno cěle koeficienty. Historijny razlog govori takože o racionalnyh čislovyh koeficientah, ale ne razlikuje dovolj jasno njihove nositelje. V aktualnom pasažu raboty 16 koeficienty polinomov $f_i$ sut cělo-racionalne; iz togo slěduje, že drugi koeficienty, $a_i$, sut racionalne čisla. Zato tuto ne smi byti utožsamnjenja dvuh uslovij. Defis pomaga čitati složeny nazv, ale sam ne rěšaje tu matematičnu razliku.\par
\noindent\emph{Reason:} Cělo-racionalne čislove koeficienty retains a compound designation for rational-integer coefficients. The historical rationale also mentions rational numerical coefficients without clearly distinguishing what they belong to. In the current Paper 16 passage, the polynomials $f_i$ have rational-integer coefficients; the other coefficients, $a_i$, are consequently rational numbers. These conditions must not be conflated. The hyphen helps readers parse the compound but does not itself settle the mathematical distinction.\par
\noindent{\footnotesize Izvorne zapisy / Original records: HIST-g078-013.}\par
\smallskip\noindent\emph{Sravnjenje s tekstom:} Rabota 16 izrično razlikuje cělo-racionalne koeficienty $f_i(\xi)$ od racionalnyh čisel $a_i$; pozdnějši pasaž podobno označaje $c_i$ kak racionalne čisla. Nejasnost kratkogo historijnogo razloga ne treba razrěšiti utožsamnjenjem tyh slov. Sravnjenje ukazuje točnu relaciju medžu nositeljami različnyh uslovij; proverka pečatnogo originala ostavaje odvojenym zadaniem.\par
\noindent\emph{Textual comparison:} Paper 16 expressly distinguishes the rational-integer coefficients of $f_i(\xi)$ from the rational numbers $a_i$; a later passage similarly identifies $c_i$ as rational numbers. The brief historical rationale must not be clarified by treating these conditions as equivalent. The comparison identifies the different expressions to which the conditions apply; collation with the printed original remains separate.\par
\noindent{\footnotesize TEXT-COMP-32; pasaže / passages: 10811, 10823; \texttt{CORPUS\_FORM\_COMPARISONS.json}.}\par
\end{minipage}\par
\par\medskip\Needspace{7\baselineskip}
\noindent\begin{minipage}{\linewidth}
\hypertarget{TERM-0664}{}
\noindent\textbf{racionalno-cěle koeficienty / cělo-racionalno-čislove invarianty}\quad{\footnotesize\texttt{TERM-0664}}\par
\noindent\emph{Iztočnik / Source:} ganzrationale Zahlkoeffizienten / ganzrationalzahlige Invarianten\par
\noindent\emph{Rusky / Russian:} рационально-целые коэффициенты / целорационально-числовые инварианты\par
\noindent\emph{Ukrajinsky / Ukrainian:} цілі раціональні числові коефіцієнти / ціло-раціонально-числові інваріанти\par
\smallskip\noindent\emph{Razlog:} Racionalno-cěle koeficienty / cělo-racionalno-čislove invarianty zadržuje razlikovanje medžu racionalno cělymi i algebraično cělymi koeficientami. Razlog sravnjaje zbitu rusku frazu s izričnejšeju ukrajinskoju i izbira analitičny medžuslovjansky izraz. To jest argument za vidimost sostavnyh smyslov, ne dokaz, že dolžši nazv jest legčejši dlja vsih čitateljev. Izbor byl prědložen dlja daljšego razsmotrenja; to ne znači, že ego odobril človečsky specialist.\par
\noindent\emph{Reason:} Racionalno-cěle koeficienty / cělo-racionalno-čislove invarianty preserves the distinction between rational-integer and algebraic-integer coefficients. The rationale compares a compact Russian expression with a more explicit Ukrainian one and chooses an analytic Interslavic phrase. This argues for making the components of meaning visible, not that a longer name has been shown easier for all readers. The choice was provisional; that does not imply approval by a human specialist.\par
\noindent{\footnotesize Izvorne zapisy / Original records: HIST-g077-015.}\par
\end{minipage}\par
\par\medskip\Needspace{7\baselineskip}
\noindent\begin{minipage}{\linewidth}
\hypertarget{TERM-0665}{}
\noindent\textbf{celočisleny}\quad{\footnotesize\texttt{TERM-0665}}\par
\noindent\emph{Iztočnik / Source:} ganzzahlig\par
\noindent\emph{Rusky / Russian:} целочисленный\par
\noindent\emph{Ukrajinsky / Ukrainian:} цілочисельний\par
\smallskip\noindent\emph{Razlog:} Celočisleny jest v uvodě raboty 9 izrično svęzany s koeficientami iz $[K]$, oblasti vsih algebraično cělyh čisel. Historijny razlog zato odbija suženje do običnyh racionalno cělyh koeficientov. Podnožna primětka razlikuje užšu konstrukciju, dlja ktorej te koeficienty vystarčajut, od širših oblastij. Slovjansky složeny pridavnik daje kratku formu, ale ego tehničny obseg nosi definicija, ne vsakodenno značenje slova čislo.\par
\noindent\emph{Reason:} Celočisleny is explicitly tied in Paper 9's introduction to coefficients in $[K]$, the domain of all algebraic integers. The historical rationale therefore rejects restriction to ordinary rational-integer coefficients. A footnote distinguishes the narrower construction, for which those coefficients suffice, from more general domains. The Slavic compound provides a short form, but its technical range comes from the definition, not the everyday meaning of number.\par
\noindent{\footnotesize Izvorne zapisy / Original records: HIST-g058-007.}\par
\smallskip\noindent\emph{Sravnjenje s tekstom:} Podnožna primětka uvoda raboty 9 izrično oprěděljaje cěločisleny prěz $[K]$, oblast vsih algebraično cělyh čisel. Ona objašnjavaje takože, začto užša upotrěba racionalno cělyh koeficientov ne jest izbrana dlja vsej teorije. To jest točno lokalno osnovanje dlja širokogo tehničnogo smysla pridavnika.\par
\noindent\emph{Textual comparison:} A footnote in Paper 9's introduction explicitly defines integral coefficients through $[K]$, the domain of all algebraic integers. It also explains why the narrower rational-integer restriction is not adopted throughout the theory. This is the precise local basis for the adjective's wider technical sense.\par
\noindent{\footnotesize TEXT-COMP-33; pasaže / passages: 7147; \texttt{CORPUS\_FORM\_COMPARISONS.json}.}\par
\end{minipage}\par
\par\medskip\Needspace{7\baselineskip}
\noindent\begin{minipage}{\linewidth}
\hypertarget{TERM-0666}{}
\noindent\textbf{celočisleno-regularny}\quad{\footnotesize\texttt{TERM-0666}}\par
\noindent\emph{Iztočnik / Source:} ganzzahlig-regulär\par
\noindent\emph{Rusky / Russian:} целочисленно-регулярный\par
\noindent\emph{Ukrajinsky / Ukrainian:} цілочисельно-регулярний\par
\smallskip\noindent\emph{Razlog:} Celočisleno-regularny dodava k regularnosti uslovje na rezultant: on mora byti obratimy element, jedinica prstena $[\Omega]$. To ne znači samo nenuljevy rezultant; i jedinica tu ne mora byti čislo $1$. Razlog svęzuje složeny pridavnik s novoju definicijeju v odděle 15 raboty 6. Sohranjenje oběh čestij nazva pomaga razlikovati to ostrějše uslovje od ranějšej regularnosti.\par
\noindent\emph{Reason:} Celočisleno-regularny adds a condition on the resultant to regularity: it must be an invertible element, a unit of $[\Omega]$. A nonzero resultant is not sufficient, and a unit need not be the number $1$. The rationale ties the compound adjective to the new definition in Paper 6 section 15. Retaining both components helps distinguish the stronger condition from the earlier notion of regularity.\par
\noindent{\footnotesize Izvorne zapisy / Original records: HIST-g055-007.}\par
\smallskip\noindent\emph{Sravnjenje s tekstom:} Definicija VIII' raboty 6 potrěbuje, že rezultant najvyšših homogenyh čestij jest jedinica iz $[\Omega]$. Sostavny pridavnik cěločisleno-regularny mora sohraniti to silnějše uslovje. Oznaka $R_0=\varepsilon$ ne daje osnovy bez dodatnogo uslovja zaměniti $\varepsilon$ čislom $1$.\par
\noindent\emph{Textual comparison:} Definition VIII' of Paper 6 requires the resultant of the highest homogeneous parts to be a unit of $[\Omega]$. The compound adjective must preserve this stronger condition. The notation $R_0=\varepsilon$ does not justify replacing $\varepsilon$ by $1$ without an additional assumption.\par
\noindent{\footnotesize TEXT-COMP-34; pasaže / passages: 6577; \texttt{CORPUS\_FORM\_COMPARISONS.json}.}\par
\end{minipage}\par
\par\medskip\Needspace{7\baselineskip}
\noindent\begin{minipage}{\linewidth}
\hypertarget{TERM-0667}{}
\noindent\textbf{cěločislne funkcije}\quad{\footnotesize\texttt{TERM-0667}}\par
\noindent\emph{Iztočnik / Source:} ganzzahlige Funktionen\par
\noindent\emph{Rusky / Russian:} целочисленные функции\par
\noindent\emph{Ukrajinsky / Ukrainian:} цілочисельні функції\par
\smallskip\noindent\emph{Razlog:} Cěločislne funkcije bylo izbrano, da ne směšati algebraično uslovje s cělymi funkcijami kompleksnoj analizy. Razlog govori o koeficientah i algebraičnoj cělosti v rabotě 35, ne o holomorfnosti na celoj kompleksnoj ravnině. On ne davaje polnu definiciju konkretnoj oblasti funkcij; slovnik zato ne dodavaje takvu definiciju iz samogo nazva. Pravopisna forma cěločislne jest zachovany prědlog, ne bez dodatnogo sravnjenja obvezna norma dlja vseh upotrěb.\par
\noindent\emph{Reason:} Cěločislne funkcije was chosen to avoid confusing an algebraic condition with entire functions of complex analysis. The rationale concerns coefficients and algebraic integrality in Paper 35, not holomorphy throughout the complex plane. It does not supply a full definition of the particular domain of functions; the dictionary cannot derive one from the name alone. The spelling cěločislne is a preserved proposal, not automatically a mandatory form for every occurrence.\par
\noindent{\footnotesize Izvorne zapisy / Original records: HIST-g205-003.}\par
\end{minipage}\par
\par\medskip\Needspace{7\baselineskip}
\noindent\begin{minipage}{\linewidth}
\hypertarget{TERM-0668}{}
\noindent\textbf{celočislena integralna baza}\quad{\footnotesize\texttt{TERM-0668}}\par
\noindent\emph{Iztočnik / Source:} ganzzahlige Integritätsbasis\par
\noindent\emph{Rusky / Russian:} целочисленный интегральный базис\par
\noindent\emph{Ukrajinsky / Ukrainian:} цілочисельна інтегральна база\par
\smallskip\noindent\emph{Razlog:} Celočislena integralna baza prěnosit nazv integralna baza iz ranějšego odděla raboty 6 i dodavaje uslovje na koeficienty iz $[\Omega]$. Razlog jest sostavny: sohraniti poznaty objekt i izrično pokazati nove ograničenje. Teorem VII' priměnjaje frazu k relativno cěloj oblasti prvogo roda. Slovo baza tu ne daje samo soboju prava zaključiti, že idě o svobodnu modulnu bazu nad običnymi cělymi čislami.\par
\noindent\emph{Reason:} Celočislena integralna baza carries integralna baza forward from an earlier section of Paper 6 and adds the coefficient condition in $[\Omega]$. The rationale is compositional: retain the established object and state the additional restriction. Theorem VII' applies the phrase to a relatively integral domain of the first kind. Baza alone does not warrant interpreting this as a free-module basis over ordinary integers.\par
\noindent{\footnotesize Izvorne zapisy / Original records: HIST-g055-001.}\par
\smallskip\noindent\emph{Sravnjenje s tekstom:} Teoremy VII' i IX' raboty 6 priměnjajut nazv cěločislena integralna baza k oblastjam prvogo roda i cěločisleno-regularnym sistemam. Te pojavy podkrepljajut svęzanu terminologiju novogo koeficientnogo uslovja. One ne ustanavljajut samy po sebi totožnost s libovoljnym sovrěmennym značenjem slova baza.\par
\noindent\emph{Textual comparison:} Theorems VII' and IX' of Paper 6 use integral-coefficient integral basis for first-kind domains and integrally regular systems. These occurrences support the related terminology under the new coefficient condition. They do not by themselves establish identity with every modern meaning of basis.\par
\noindent{\footnotesize TEXT-COMP-35; pasaže / passages: 6573, 6585; \texttt{CORPUS\_FORM\_COMPARISONS.json}.}\par
\end{minipage}\par
\par\medskip\Needspace{7\baselineskip}
\noindent\begin{minipage}{\linewidth}
\hypertarget{TERM-0669}{}
\noindent\textbf{cěločislovy invariant}\quad{\footnotesize\texttt{TERM-0669}}\par
\noindent\emph{Iztočnik / Source:} ganzzahlige Invariante\par
\noindent\emph{Rusky / Russian:} целочисленный инвариант\par
\noindent\emph{Ukrajinsky / Ukrainian:} цілочисельний інваріант\par
\smallskip\noindent\emph{Razlog:} Cěločislovy invariant označaje invariant s cělymi koeficientami, ne rezultat analitičnogo integrirovanja. Razlog opira se na zbližene rusky i ukrajinsky pridavniki i izbira prozračnu slovjansku složeninu. Pan-slovjanska razumlivost jest naměra izbora; v zapisu nema měrjenja razuměnja čitateljev. Dlatego čitatelj dostavaje i nazv, i matematično ograničenje, vměsto neopravdanoj obětnice, že sam nazv jest vsim očividny.\par
\noindent\emph{Reason:} Cěločislovy invariant denotes an invariant with integral coefficients, not the result of analytic integration. The rationale draws on the closely related Russian and Ukrainian adjectives and selects a transparent Slavic compound. Pan-Slavic intelligibility is the intention of the choice; the record contains no measurement of reader comprehension. The entry therefore supplies both the name and its mathematical restriction rather than claiming that the name is self-explanatory to everyone.\par
\noindent{\footnotesize Izvorne zapisy / Original records: HIST-g077-001.}\par
\end{minipage}\par
\par\medskip\Needspace{7\baselineskip}
\noindent\begin{minipage}{\linewidth}
\hypertarget{TERM-0670}{}
\noindent\textbf{celočislene koeficienty}\quad{\footnotesize\texttt{TERM-0670}}\par
\noindent\emph{Iztočnik / Source:} ganzzahlige Koeffizienten\par
\noindent\emph{Rusky / Russian:} целочисленные коэффициенты\par
\noindent\emph{Ukrajinsky / Ukrainian:} цілочисельні коефіцієнти\par
\smallskip\noindent\emph{Razlog:} Celočislene koeficienty izrično priradžuje cěločislenost koeficientam. Historijny razlog zadržuje razlikovanje medžu takym uslovjem i cělostju samoj funkcije ili drugogo objekta. Dlja raboty 9 oblast čisel jest oprěděljeno pri TERM-0665, a posebny racionalny funkcional pri TERM-0662. Kratka fraza ne zaměnjaje te definicije i ne dopuščaje bez konteksta svesti koeficienty samo do običnyh cělyh čisel.\par
\noindent\emph{Reason:} Celočislene koeficienty explicitly assigns integrality to the coefficients. The historical rationale retains the distinction between this condition and integrality of a function or another object. For Paper 9, TERM-0665 explains the domain of numbers and TERM-0662 the special rational functional. The short phrase does not replace those definitions or permit an automatic restriction to ordinary integers.\par
\noindent{\footnotesize Izvorne zapisy / Original records: HIST-g061-003.}\par
\end{minipage}\par
\par\medskip\Needspace{7\baselineskip}
\noindent\begin{minipage}{\linewidth}
\hypertarget{TERM-0671}{}
\noindent\textbf{celočislena linearna familija}\quad{\footnotesize\texttt{TERM-0671}}\par
\noindent\emph{Iztočnik / Source:} ganzzahlige lineare Schar\par
\noindent\emph{Rusky / Russian:} целочисленное линейное семейство\par
\noindent\emph{Ukrajinsky / Ukrainian:} цілочисельна лінійна родина\par
\smallskip\noindent\emph{Razlog:} Celočislena linearna familija nastavljaje ranějši nazv linearna familija, ale ograničaje skalary na $[\Omega]$. Definicija potrěbuje, že iz $F_1$ i $F_2$ slěduje sodržanje $c_1F_1+c_2F_2$ za libovoljne $c_1,c_2$ iz togo prstena. To jest više od spisa cěločislenyh polinomov i menje od izjavjenja, že dopustimy sut vse skalary polja $\Omega$. Razlog za familija jest terminologična kontinuitetnost, ne novo, nezapisano sravnjenje vseh mogućih imenikov.\par
\noindent\emph{Reason:} Celočislena linearna familija continues the earlier name linearna familija but restricts the scalars to $[\Omega]$. The definition requires $c_1F_1+c_2F_2$ to belong whenever $F_1,F_2$ do, for arbitrary $c_1,c_2$ in that ring. This is more than a list of integral-coefficient polynomials and does not assert closure under every scalar in the field $\Omega$. The recorded reason for familija is terminological continuity, not a new comparison of all possible nouns.\par
\noindent{\footnotesize Izvorne zapisy / Original records: HIST-g054-005.}\par
\smallskip\noindent\emph{Sravnjenje s tekstom:} Definicija v odděle 14 raboty 6 izrično ograničaje $c_1,c_2$ v kombinaciji $c_1F_1+c_2F_2$ na $[\Omega]$. Pridavnik cěločislena ne opisuje neoprěděljenu opću linearnost; on zadržuje točno koeficientno ograničenje.\par
\noindent\emph{Textual comparison:} The definition in Paper 6 section 14 explicitly restricts $c_1,c_2$ in $c_1F_1+c_2F_2$ to $[\Omega]$. The adjective does not express unspecified general linearity: it retains this precise coefficient restriction.\par
\noindent{\footnotesize TEXT-COMP-36; pasaže / passages: 6444; \texttt{CORPUS\_FORM\_COMPARISONS.json}.}\par
\end{minipage}\par
\par\medskip\Needspace{7\baselineskip}
\noindent\begin{minipage}{\linewidth}
\hypertarget{TERM-0672}{}
\noindent\textbf{cěločiselne matrice}\quad{\footnotesize\texttt{TERM-0672}}\par
\noindent\emph{Iztočnik / Source:} ganzzahlige Matrizen\par
\noindent\emph{Rusky / Russian:} целочисленные матрицы\par
\noindent\emph{Ukrajinsky / Ukrainian:} цілочислові матриці\par
\smallskip\noindent\emph{Razlog:} Cěločiselne matrice označaje matrice s cělymi elementami. Razlog zadržuje konkretny objekt v priměru o dělimosti i elementarnyh děliteljah. Pridavnik kvalifikuje zapisy matrice; ne govori o cěločislovoj veličině samogo determinanta kako dostatočnom uslovju. Spomenjenje počitanja po elementah ili klasah v razloge takože ne znači, že vse operacije matričnoj dělimosti možno zaměniti nezavisnym děljenjem jednotlivyh zapisov.\par
\noindent\emph{Reason:} Cěločiselne matrice denotes matrices with integer entries. The rationale keeps the object concrete in an example concerning divisibility and elementary divisors. The adjective qualifies the matrix entries; an integer determinant alone would not suffice. The rationale's reference to calculations by entries or classes also does not mean that matrix divisibility can always be replaced by independent division of individual entries.\par
\noindent{\footnotesize Izvorne zapisy / Original records: HIST-g081-096.}\par
\end{minipage}\par
\par\medskip\Needspace{7\baselineskip}
\noindent\begin{minipage}{\linewidth}
\hypertarget{TERM-0673}{}
\noindent\textbf{celočislno grupovo kolco}\quad{\footnotesize\texttt{TERM-0673}}\par
\noindent\emph{Iztočnik / Source:} ganzzahliger Gruppenring\par
\noindent\emph{Rusky / Russian:} целочисленное групповое кольцо\par
\noindent\emph{Ukrajinsky / Ukrainian:} цілочислове групове кільце\par
\smallskip\noindent\emph{Razlog:} Celočislno grupovo kolco zadržuje prstenovy objekt i koeficientno ograničenje, da razlikovati ego od racionalnoj grupovoj algebry. To razlikovanje jest cěl zapisanogo razloga. Forma kolco jest historijny prědlog; sama prisutnost v slovniku ne ustanavljaje jego prednost pred drugymi nazvami prstena v pozdnějšej terminologiji. Ni ruska, ni ukrajinska paralela ne zaměnjaje samostojno opravdanje konačnogo medžuslovjanskogo nazva.\par
\noindent\emph{Reason:} Celočislno grupovo kolco retains the ring-theoretic object and coefficient restriction to distinguish it from the rational group algebra. That is the purpose stated in the rationale. Kolco is the historical proposal; inclusion here does not establish its precedence over other ring terminology adopted later. Neither the Russian nor the Ukrainian parallel substitutes for a separate justification of the final Interslavic name.\par
\noindent{\footnotesize Izvorne zapisy / Original records: HIST-g207-012.}\par
\end{minipage}\par
\par\medskip\Needspace{7\baselineskip}
\noindent\begin{minipage}{\linewidth}
\hypertarget{TERM-0674}{}
\noindent\textbf{celočislena integralna oblast}\quad{\footnotesize\texttt{TERM-0674}}\par
\noindent\emph{Iztočnik / Source:} ganzzahliger Integritätsbereich\par
\noindent\emph{Rusky / Russian:} целочисленная область целостности\par
\noindent\emph{Ukrajinsky / Ukrainian:} цілочисельна інтегральна область\par
\smallskip\noindent\emph{Razlog:} Celočislena integralna oblast nastavljaje integralna oblast i dodavaje cěločisleno uslovje. V definiciji raboty 6 to jest linearna familija, ktora s dvoma polinomami sodrživaje i jih produkt. Taka sostavnost objasnjavaje izbor: ne vymyslja se novo kratko ime dlja objekta, čije svojstva sut dodavane po čestjah. Obče matematične asocijacije slova integralny ne smějut zaměniti to lokalno oprěděljenje.\par
\noindent\emph{Reason:} Celočislena integralna oblast continues integralna oblast and adds the integral-coefficient condition. Paper 6 defines it as a linear family containing the product whenever it contains two polynomials. This compositional structure explains the choice: no new short name is introduced for an object whose properties are being added step by step. General mathematical associations of integralny must not replace the local definition.\par
\noindent{\footnotesize Izvorne zapisy / Original records: HIST-g054-006.}\par
\smallskip\noindent\emph{Sravnjenje s tekstom:} Odmah poslě linearnoj familije definicija dodavaje zamknjenost pri produktu. Tako prěvod integralna oblast jest svęzany s konkretnym šagom konstrukcije. Nazvy v TERM-0671 i TERM-0674 ne sut dva različne imena za isti spis polinomov.\par
\noindent\emph{Textual comparison:} Immediately after the linear-family definition, closure under products is added. Integralna oblast is thus tied to a particular step of the construction. TERM-0671 and TERM-0674 are not merely different names for the same list of polynomials.\par
\noindent{\footnotesize TEXT-COMP-37; pasaže / passages: 6446; \texttt{CORPUS\_FORM\_COMPARISONS.json}.}\par
\end{minipage}\par
\par\medskip\Needspace{7\baselineskip}
\noindent\begin{minipage}{\linewidth}
\hypertarget{TERM-0675}{}
\noindent\textbf{celočislena relativno cěla oblast}\quad{\footnotesize\texttt{TERM-0675}}\par
\noindent\emph{Iztočnik / Source:} ganzzahliger relativ-ganzer Bereich\par
\noindent\emph{Rusky / Russian:} целочисленная относительно целая область\par
\noindent\emph{Ukrajinsky / Ukrainian:} цілочисельна відносно ціла область\par
\smallskip\noindent\emph{Razlog:} Celočislena relativno cěla oblast sojedinjaje ranějši nazv relativnoj cělosti s koeficientnym uslovjem. Definicija vključaje kvocient $F/G$ za $G\ne0$ togda i samo togda, kogda toj kvocient jest cěločisleny polinom. Zato oblast ne sodrživaje avtomatično vsaky drobek svojih elementov. Dolžina frazy sohranjaje dva različne uslovja, ktore by bez dodatnogo objasnjenja propali v kratkom nazvě polje.\par
\noindent\emph{Reason:} Celočislena relativno cěla oblast combines the earlier relative-integrality term with the coefficient condition. Its definition includes $F/G$, for $G\ne0$, if and only if that quotient is an integral-coefficient polynomial. The domain therefore does not automatically contain every fraction of its elements. The long phrase preserves two distinct conditions that would disappear if it were simply called a field.\par
\noindent{\footnotesize Izvorne zapisy / Original records: HIST-g054-007.}\par
\smallskip\noindent\emph{Sravnjenje s tekstom:} Definicija relativno cěloj oblasti v odděle 14 govori o kvocientě dvuh jej elementov samo pod uslovjem, že rezultat jest cěločisleny polinom. Točny uslovny oblik jest sučstveny i dlja nazva, i dlja njegovogo objasnjenja. Polje drobkov bylo by drugym objektom.\par
\noindent\emph{Textual comparison:} Section 14's definition of a relatively integral domain concerns a quotient of two of its elements only when the result is an integral-coefficient polynomial. The conditional form matters both to the name and to its explanation. A field of fractions would be a different object.\par
\noindent{\footnotesize TEXT-COMP-38; pasaže / passages: 6448; \texttt{CORPUS\_FORM\_COMPARISONS.json}.}\par
\end{minipage}\par
\par\medskip\Needspace{7\baselineskip}
\noindent\begin{minipage}{\linewidth}
\hypertarget{TERM-0676}{}
\noindent\textbf{celočislena relativno cěla oblast prvogo roda}\quad{\footnotesize\texttt{TERM-0676}}\par
\noindent\emph{Iztočnik / Source:} ganzzahliger relativ-ganzer Bereich erster Art\par
\noindent\emph{Rusky / Russian:} целочисленная относительно целая область первого рода\par
\noindent\emph{Ukrajinsky / Ukrainian:} цілочисельна відносно ціла область першого роду\par
\smallskip\noindent\emph{Razlog:} Celočislena relativno cěla oblast prvogo roda dodavaje klasifikaciju prvogo roda k objektu iz TERM-0675. Razlog opira se na kontinuitetnost oddělov 10, 14 i 15 raboty 6. Definicija VII' svęzuje tu klasifikaciju s oblastju odobraženja i algebraično cěloju zavisimostju polinomov. Ne treba prěnositi sem drugo značenje prvogo roda iz teorije razširjenij polj v rabotě 24; to različenje jest pokazano takože pri TERM-0610.\par
\noindent\emph{Reason:} Celočislena relativno cěla oblast prvogo roda adds the first-kind classification to TERM-0675. The rationale follows continuity between sections 10, 14 and 15 of Paper 6. Definition VII' connects this classification with an image domain and algebraically integral dependence of polynomials. The different first-kind terminology for field extensions in Paper 24 must not be imported here; TERM-0610 also explains that distinction.\par
\noindent{\footnotesize Izvorne zapisy / Original records: HIST-g055-004.}\par
\smallskip\noindent\emph{Sravnjenje s tekstom:} Definicija VII' razširjaje ranějšu klasifikaciju prěz oblast odobraženja, s novym uslovjem na cěločislene polinomy. To oprěděljenje treba čitati vměstě s definicijeju $[\Omega]$, ne s drugym upotrěbljenjem prvogo roda v razširjenjah polj.\par
\noindent\emph{Textual comparison:} Definition VII' extends the earlier image-domain classification with a condition on integral-coefficient polynomials. It must be read with the definition of $[\Omega]$, not with the different use of first kind for field extensions.\par
\noindent{\footnotesize TEXT-COMP-39; pasaže / passages: 6556; \texttt{CORPUS\_FORM\_COMPARISONS.json}.}\par
\end{minipage}\par
\par\medskip\Needspace{7\baselineskip}
\noindent\begin{minipage}{\linewidth}
\hypertarget{TERM-0677}{}
\noindent\textbf{celočisleny polinom}\quad{\footnotesize\texttt{TERM-0677}}\par
\noindent\emph{Iztočnik / Source:} ganzzahliges Polynom\par
\noindent\emph{Rusky / Russian:} целочисленный полином\par
\noindent\emph{Ukrajinsky / Ukrainian:} цілочисельний поліном\par
\smallskip\noindent\emph{Razlog:} Celočisleny polinom znači polinom s koeficientami iz $[\Omega]$, prstena algebraično cělyh čisel danogo polja. Razlog izbira pridavnik, da ostaviti isti oblik v nazvah polinoma, sistema i oblasti. To jest prosta gramatyčna strategija za svęzanu čest terminologije. Ona ne sužuje $[\Omega]$ do običnyh cělyh čisel i ne označaje samo svojstvo prijmati cěle vrijednosti pri cělyh argumentah.\par
\noindent\emph{Reason:} Celočisleny polinom means a polynomial with coefficients in $[\Omega]$, the ring of algebraic integers in the given field. The rationale chooses an adjective so the same construction can qualify polynomial, system and domain. This is a simple grammatical strategy for a related set of terms. It neither restricts $[\Omega]$ to ordinary integers nor denotes merely taking integer values at integer arguments.\par
\noindent{\footnotesize Izvorne zapisy / Original records: HIST-g054-003.}\par
\smallskip\noindent\emph{Sravnjenje s tekstom:} Definicija iz odděla 14 davaje polinomu koeficienty iz $[\Omega]$ i slědujuča definicija prenosi to uslovje na libovoljny sistem takyh polinomov. Zapis $\mathfrak S_{n\rho}$ pri TERM-0678 sohranjaje točnu matematičnu oznaku. Nazv sistema sam ne dodavaje ni linearnoj, ni multiplikativnoj zamknjenosti.\par
\noindent\emph{Textual comparison:} Section 14 first defines a polynomial with coefficients in $[\Omega]$, then extends the condition to an arbitrary system of such polynomials. The notation $\mathfrak S_{n\rho}$ at TERM-0678 preserves the mathematical identifier. The system designation alone adds neither linear nor multiplicative closure.\par
\noindent{\footnotesize TEXT-COMP-40; pasaže / passages: 6438, 6442; \texttt{CORPUS\_FORM\_COMPARISONS.json}.}\par
\end{minipage}\par
\par\medskip\Needspace{7\baselineskip}
\noindent\begin{minipage}{\linewidth}
\hypertarget{TERM-0678}{}
\noindent\textbf{celočisleny sistem $\mathfrak{S}_{n\rho}$}\quad{\footnotesize\texttt{TERM-0678}}\par
\noindent\emph{Iztočnik / Source:} ganzzahliges System $\mathfrak{S}_{n\rho}$\par
\noindent\emph{Rusky / Russian:} целочисленная система $\mathfrak{S}_{n\rho}$\par
\noindent\emph{Ukrajinsky / Ukrainian:} цілочисельна система $\mathfrak{S}_{n\rho}$\par
\smallskip\noindent\emph{Razlog:} Celočisleny sistem nastavljaje nazv libovoljnogo sistema iz odděla 7, nyně s polinomami, čije koeficienty sut cěle v smyslu odděla 14. Indeksy $n$ i $\rho$ ostavajut čestju matematičnoj oznaky i ne prevodet se kak slova. Razlog zato ne pripisyvaje sistemu zamknjenost pri linearnyh kombinacijah ili produktah: te dopolnitelne svojstva imajut svoje odvojene nazvy pri TERM-0671 i TERM-0674.\par
\noindent\emph{Reason:} Celočisleny sistem continues section 7's designation for an arbitrary system, now of polynomials whose coefficients are integral in section 14's sense. The indices $n$ and $\rho$ remain mathematical notation and are not translated as words. The rationale does not attribute closure under linear combinations or products to the system: those additional properties have separate names at TERM-0671 and TERM-0674.\par
\noindent{\footnotesize Izvorne zapisy / Original records: HIST-g054-004.}\par
\end{minipage}\par
\par\medskip\Needspace{7\baselineskip}
\noindent\begin{minipage}{\linewidth}
\hypertarget{TERM-0679}{}
\noindent\textbf{celočislovost}\quad{\footnotesize\texttt{TERM-0679}}\par
\noindent\emph{Iztočnik / Source:} Ganzzahligkeit\par
\noindent\emph{Rusky / Russian:} целочисленность\par
\noindent\emph{Ukrajinsky / Ukrainian:} цілочисельність\par
\smallskip\noindent\emph{Razlog:} Celočislovost imenuje uslovje, ktory odděl 14 dodavaje k ranějšemu izloženju: koeficienty polinomov naležet $[\Omega]$. Imenik dovolja kratko govoriti o toj svojstvě, dokolě pridavniki v TERM-0677 i TERM-0678 svęzujut ju s konkretnymi objektami. Historijny razlog ne kaže, že vse ranějše upotrěby cěly iměly točno isto značenje. Definicija koeficientnoj oblasti ostavaje obvezna dlja razuměnja.\par
\noindent\emph{Reason:} Celočislovost names the condition added in section 14: polynomial coefficients belong to $[\Omega]$. The noun allows concise discussion of that property, while the adjectives in TERM-0677 and TERM-0678 attach it to particular objects. The historical rationale does not say that all earlier uses of cěly had exactly this meaning. The coefficient-domain definition remains necessary.\par
\noindent{\footnotesize Izvorne zapisy / Original records: HIST-g054-001.}\par
\end{minipage}\par
\par\medskip\Needspace{7\baselineskip}
\noindent\begin{minipage}{\linewidth}
\hypertarget{TERM-0680}{}
\noindent\textbf{cěločislovost}\quad{\footnotesize\texttt{TERM-0680}}\par
\noindent\emph{Iztočnik / Source:} Ganzzahligkeit\par
\noindent\emph{Rusky / Russian:} целочисленность\par
\noindent\emph{Ukrajinsky / Ukrainian:} цілочисельність\par
\smallskip\noindent\emph{Razlog:} Cěločislovost jest kratky imenik za cěločisleno uslovje v pregledě invariantnoj teorije raboty 23. Historijny razlog daje dva anglijske objasnjenja, integrality i integer-valuedness, ale te pojmy ne sut opće totožne. Aktualny pasaž govori o polinomah s cělymi koeficientami i o ostrějšem teoremu konečnosti. Zato tu objašnjenje potrěbuje koeficientny smysl, ne samo cěle vrijednosti funkcije. Historijna oznaka dlja daljšego razsmotrenja ne jest dokaz izvršenoj ekspertnoj ocěnky.\par
\noindent\emph{Reason:} Cěločislovost is a compact noun for the integral-coefficient condition in Paper 23's survey of invariant theory. The historical rationale offers the English glosses integrality and integer-valuedness, but these are not generally equivalent. The current passage concerns polynomials with integer coefficients and a stronger finiteness theorem. The explanation therefore needs the coefficient sense, not merely integer values of a function. A historical flag for further consideration is not evidence of completed expert review.\par
\noindent{\footnotesize Izvorne zapisy / Original records: HIST-g094-012.}\par
\smallskip\noindent\emph{Sravnjenje s tekstom:} Punkt 5 raboty 23 razsmatrjaje ostrějši teorem konečnosti dlja polinomov s cělymi koeficientami. Cěločislovost tu ne jest prosto integer-valuedness, svojstvo davati cěle vrijednosti na cělyh argumentah. Tvrdnja o tada otvorjenom pytanju jest historijny obsah raboty, ne ocěnka nynějšego sostojanja matematiki.\par
\noindent\emph{Textual comparison:} Item 5 of Paper 23 considers a stronger finiteness theorem for polynomials with integer coefficients. Integrality here is not merely integer-valuedness, the property of taking integer values on integer arguments. The statement that a question was then open is part of the paper's historical content, not an assessment of present-day mathematics.\par
\noindent{\footnotesize TEXT-COMP-41; pasaže / passages: 15166; \texttt{CORPUS\_FORM\_COMPARISONS.json}.}\par
\end{minipage}\par
\par\medskip\Needspace{7\baselineskip}
\noindent\begin{minipage}{\linewidth}
\hypertarget{TERM-0681}{}
\noindent\textbf{integralnost predstavjenja}\quad{\footnotesize\texttt{TERM-0681}}\par
\noindent\emph{Iztočnik / Source:} Ganzzahligkeit der Darstellung\par
\noindent\emph{Rusky / Russian:} интегральность представления\par
\noindent\emph{Ukrajinsky / Ukrainian:} інтегральність подання\par
\smallskip\noindent\emph{Razlog:} Integralnost predstavjenja zadržuje algebraičny smysl cělosti, ne prosto svojstvo prijmati cěle čislove vrijednosti. Razlog odbija to užše aritmetično čitanje. V rabotě 5 integralna baza znači polinomno predstavjenje drugih polinomov, a posebno uslovje na rezultant osiguruje takože integralnost predstavjenja. To razlikovanje potrěbuje kontekst; sam nazv ne kaže, nad kakoju koeficientnoju oblastju jest predstavjenje vzeto.\par
\noindent\emph{Reason:} Integralnost predstavjenja retains an algebraic notion of integrality, not simply the property of taking integer values. The rationale rejects that narrower arithmetic reading. In Paper 5 an integral basis gives polynomial expressions for other polynomials, while a separate condition on the resultant also guarantees integrality of the representation. The distinction requires context; the name alone does not specify the coefficient domain.\par
\noindent{\footnotesize Izvorne zapisy / Original records: HIST-g039-019.}\par
\smallskip\noindent\emph{Sravnjenje s tekstom:} Rabota 5 najprvo oprěděljaje integralnu bazu pomoču polinomnogo predstavjenja s koeficientami iz danej oblasti. Potom razlikuje nenuljevu homogenu rezultantu od osobnogo uslovja, že ona jest jedinica koeficientnoj oblasti. Slovnično objasnjenje mora zadržati tu razliku; integralnost ne može byti svedena do cělyh vrijednostij funkcije.\par
\noindent\emph{Textual comparison:} Paper 5 first defines an integral basis through polynomial representation with coefficients in a specified domain. It then distinguishes a nonzero homogeneous resultant from the additional condition that it be a unit of the coefficient domain. The explanation must retain that distinction; integrality cannot be reduced to a function taking integer values.\par
\noindent{\footnotesize TEXT-COMP-42; pasaže / passages: 4966, 4968, 4973; \texttt{CORPUS\_FORM\_COMPARISONS.json}.}\par
\end{minipage}\par
\par\medskip\Needspace{7\baselineskip}
\noindent\begin{minipage}{\linewidth}
\hypertarget{TERM-0682}{}
\noindent\textbf{oblasti}\quad{\footnotesize\texttt{TERM-0682}}\par
\noindent\emph{Iztočnik / Source:} Gebiete\par
\noindent\emph{Rusky / Russian:} области\par
\noindent\emph{Ukrajinsky / Ukrainian:} області\par
\smallskip\noindent\emph{Razlog:} Oblasti v rabotě 42 sut klasy maksimalnyh poredkov, oprěděljene aritmetičnym uslovjem. Razlog zadržuje široky slovjansky imenik s izričnoju definicijeju, vměsto pripisanja jemu vsih svojstv sovrěmennogo algebraičnogo domena. V razloženom slučaju jedna oblast objedinjuje poredky s jednakym prěsěkom s poljem $k$; dalje klasifikacija potrěbuje utočnjenje na razvětvjenyh městah. Veliko čislo pojavjenij oblasti v korpusu ne dokazuje, že vsaka upotrěba imaje to značenje.\par
\noindent\emph{Reason:} Oblasti in Paper 42 are classes of maximal orders defined by an arithmetic condition. The rationale retains a broad Slavic noun with an explicit definition instead of attributing every property of a modern algebraic domain to it. In the split case, one class groups orders with the same intersection with the field $k$; the classification later requires refinement at ramified places. The many corpus occurrences of oblasti do not establish that every occurrence has this sense.\par
\noindent{\footnotesize Izvorne zapisy / Original records: HIST-g212-009.}\par
\smallskip\noindent\emph{Sravnjenje s tekstom:} Uvod raboty 42 izrično grupuje maksimalne poredky po jednakom prěsěku s $k$. Dlja občejših prostyh algebr potrěbuje se takože soglasje lokalnyh komponentov na razvětvjenyh městah. Zato oblasti tu označajut aritmetičnu klasifikaciju; ne sut bez daljšego totožne s integralnymi oblastjami ranějših rabot.\par
\noindent\emph{Textual comparison:} Paper 42's introduction explicitly groups maximal orders by equality of their intersection with $k$. For more general simple algebras, agreement of local components at ramified places is also required. Oblasti here therefore denotes an arithmetic classification, not automatically the integral domains of earlier papers.\par
\noindent{\footnotesize TEXT-COMP-43; pasaže / passages: 24451, 24453; \texttt{CORPUS\_FORM\_COMPARISONS.json}.}\par
\end{minipage}\par
\par\medskip\Needspace{7\baselineskip}
\noindent\begin{minipage}{\linewidth}
\hypertarget{TERM-0683}{}
\noindent\textbf{drobno algebraično čislo}\quad{\footnotesize\texttt{TERM-0683}}\par
\noindent\emph{Iztočnik / Source:} gebrochene algebraische Zahl\par
\noindent\emph{Rusky / Russian:} дробное алгебраическое число\par
\noindent\emph{Ukrajinsky / Ukrainian:} дробове алгебраїчне число\par
\smallskip\noindent\emph{Razlog:} Drobno algebraično čislo zadržuje protivopostavjenje algebraično cělym čislam. Drobno tu ne ograničaje čislo do racionalnogo drobka: čislo ostavaje algebraično, ale ne cělo v danom smyslu. Razlog jest sohranjenje toj pary v rabotě 9. On ne pridavaje slovu drobno značenje neveliko ili stepenj približenja.\par
\noindent\emph{Reason:} Drobno algebraično čislo preserves the contrast with algebraic integers. Drobno does not restrict the number to a rational fraction: it remains algebraic but is not integral in the relevant sense. The rationale is continuity of this contrast in Paper 9, not a suggestion that drobno means small or approximate.\par
\noindent{\footnotesize Izvorne zapisy / Original records: HIST-g061-007.}\par
\end{minipage}\par
\par\medskip\Needspace{7\baselineskip}
\noindent\begin{minipage}{\linewidth}
\hypertarget{TERM-0684}{}
\noindent\textbf{drobne funkcionaly}\quad{\footnotesize\texttt{TERM-0684}}\par
\noindent\emph{Iztočnik / Source:} gebrochene Funktionale\par
\noindent\emph{Rusky / Russian:} дробные функционалы\par
\noindent\emph{Ukrajinsky / Ukrainian:} дробові функціонали\par
\smallskip\noindent\emph{Razlog:} Drobne funkcionaly prodolžuje terminologiju odděla 2 v završnoj primětkě odděla 3 raboty 9. Forma funkcionaly povezuje racionalne i algebraične konstrukcije; ne znači avtomatično sovrěmenne linearne funkcionaly. Kvalifikator drobne treba čitati v protivopostavljenju cělym funkcionalam. Dlja racionalnogo slučaja definicija pri TERM-0662 pokazuje, začto prisutnost imenovatelja sama ne vystarčaje dlja razlikovanja.\par
\noindent\emph{Reason:} Drobne funkcionaly carries section 2's vocabulary into the closing remark of Paper 9 section 3. Funkcionaly links the rational and algebraic constructions; it does not automatically mean modern linear functionals. Drobne must be read in contrast to whole functionals. For the rational case, TERM-0662 explains why the presence of a denominator alone does not distinguish the classes.\par
\noindent{\footnotesize Izvorne zapisy / Original records: HIST-g061-012.}\par
\end{minipage}\par
\par\medskip\Needspace{7\baselineskip}
\noindent\begin{minipage}{\linewidth}
\hypertarget{TERM-0685}{}
\noindent\textbf{drobny ideal}\quad{\footnotesize\texttt{TERM-0685}}\par
\noindent\emph{Iztočnik / Source:} gebrochenes Ideal\par
\noindent\emph{Rusky / Russian:} дробный идеал\par
\noindent\emph{Ukrajinsky / Ukrainian:} дробовий ідеал\par
\smallskip\noindent\emph{Razlog:} Drobny ideal sohranjaje metaforu drobka dlja konečno generovanogo modula v polju kvocientov. Dva zapisy opirajut se na kontinuitetnost uvoda i odděla 9 raboty 30. Aktualny tekst ne ostavaje pri metaforě: on izrično opravdavaje nazv predstavjenjem modula kak modulnogo kvocienta dvoh idealov. Konečny tu ne treba bez daljšego čitati kak imějuči konečno čislo elementov.\par
\noindent\emph{Reason:} Drobny ideal retains the fraction metaphor for a finitely generated module in the quotient field. Two records appeal to continuity between Paper 30's introduction and section 9. The current text does more than use a metaphor: it explicitly justifies the name by representing such a module as a module quotient of two ideals. Finite here must not automatically be read as having finitely many elements.\par
\noindent{\footnotesize Izvorne zapisy / Original records: HIST-g116-002, HIST-g116-010.}\par
\smallskip\noindent\emph{Sravnjenje s tekstom:} Odděl 9 raboty 30 oprěděljaje drobny ideal kak konečny modul v polju kvocientov i potom daje osnovanje nazva: takvy modul možno predstaviti modulnym kvocientom dvoh idealov. Dokaz vybiraje konečnu bazu $t_1,\ldots,t_n$ i obči imenovatelj. Tute jest konečna generovanost, ne tvrđenje o konečnom čislu elementov modula.\par
\noindent\emph{Textual comparison:} Paper 30 section 9 defines a fractional ideal as a finite module in the quotient field and then justifies the name by expressing such a module as a module quotient of two ideals. The argument chooses a finite basis $t_1,\ldots,t_n$ and a common denominator. This is finite generation, not an assertion that the module has finitely many elements.\par
\noindent{\footnotesize TEXT-COMP-44; pasaže / passages: 17077, 17093; \texttt{CORPUS\_FORM\_COMPARISONS.json}.}\par
\end{minipage}\par
\par\medskip\Needspace{7\baselineskip}
\noindent\begin{minipage}{\linewidth}
\hypertarget{TERM-0686}{}
\noindent\textbf{drobny ideal / pravo drobny element}\quad{\footnotesize\texttt{TERM-0686}}\par
\noindent\emph{Iztočnik / Source:} gebrochenes Ideal / echt gebrochene Größe\par
\noindent\emph{Rusky / Russian:} дробный идеал / истинно дробный элемент\par
\noindent\emph{Ukrajinsky / Ukrainian:} дробовий ідеал / справді дробовий елемент\par
\smallskip\noindent\emph{Razlog:} Drobny ideal / pravo drobny element razlikuje algebraičny objekt od svojstva konkretnogo elementa. Pravo drobny podkrepljaje, že v Dedekindovoj lemmě idě o dejstvitelno drobny slučaj, ne samo o izraz s napisanym imenovateljem. Razlog želi zadržati tu razliku bez prěhoda k običnomu čislovomu drobku. Zapis ne ustanavljaje, že vsaky drobny ideal mora byti neintegralny: nazv klasa i ostro protivopostavjenje za element sut različne věči.\par
\noindent\emph{Reason:} Drobny ideal / pravo drobny element distinguishes an algebraic object from a property of a particular element. Pravo drobny marks the genuinely fractional case in Dedekind's lemma, not merely an expression written with a denominator. The rationale preserves that distinction without reducing the object to an ordinary numerical fraction. It does not establish that every fractional ideal is nonintegral: the class name and the strict contrast for an element are different matters.\par
\noindent{\footnotesize Izvorne zapisy / Original records: HIST-g107-009.}\par
\end{minipage}\par
\par\medskip\Needspace{7\baselineskip}
\noindent\begin{minipage}{\linewidth}
\hypertarget{TERM-0687}{}
\noindent\textbf{drobny racionalny funkcional}\quad{\footnotesize\texttt{TERM-0687}}\par
\noindent\emph{Iztočnik / Source:} gebrochenes rationales Funktional\par
\noindent\emph{Rusky / Russian:} дробный рациональный функционал\par
\noindent\emph{Ukrajinsky / Ukrainian:} дробовий раціональний функціонал\par
\smallskip\noindent\emph{Razlog:} Drobny racionalny funkcional jest parny nazv k cěly racionalny funkcional iz TERM-0662. Razlog sohranjaje Veberovo različenje v rabotě 9. V normalizovanom zapisu $A=aE_1/E_2$ jest važno, či racionalny faktor $a$ jest cěly, ne prosto či jest napisan drobek. Zato polinom s necělymi racionalnymi koeficientami može naležati k drobnym funkcionalam, hoti ne imaje prěměnnoj v imenovatelju.\par
\noindent\emph{Reason:} Drobny racionalny funkcional is the counterpart of cěly racionalny funkcional at TERM-0662. The rationale preserves the Weber distinction used in Paper 9. In the normalized expression $A=aE_1/E_2$, the relevant question is whether the rational factor $a$ is integral, not simply whether a fraction is written. Thus a polynomial with nonintegral rational coefficients can be a fractional functional despite having no variable in its denominator.\par
\noindent{\footnotesize Izvorne zapisy / Original records: HIST-g060-003.}\par
\end{minipage}\par
\par\medskip\Needspace{7\baselineskip}
\noindent\begin{minipage}{\linewidth}
\hypertarget{TERM-0688}{}
\noindent\textbf{vzaimno relativno proste}\quad{\footnotesize\texttt{TERM-0688}}\par
\noindent\emph{Iztočnik / Source:} gegenseitig relativprim\par
\noindent\emph{Rusky / Russian:} взаимно относительно просты; взаимно относительно простые\par
\noindent\emph{Ukrajinsky / Ukrainian:} взаємно відносно прості\par
\smallskip\noindent\emph{Razlog:} Vzaimno relativno proste zadržuje dva uslovja: relativnu prostost v jednom i v obratnom směru. Dva razloga izrično upozornjajut, že sama relacija v rabotě 19 ne jest simetrična. Zato vzaimno ne jest ukrašanje ili izlišny sinonim k relativno. Aktualna definicija s protivopriměrom pokazuje potrěbu oběh směrov; pravopisna razlika medžu historijnym vzaimno i aktualnym vzajemno ostavaje vidima.\par
\noindent\emph{Reason:} Vzaimno relativno proste retains two conditions: relative primeness in one direction and in the reverse direction. Both rationales explicitly note that Paper 19's relation is not symmetric. Vzaimno is therefore neither decoration nor a redundant synonym for relativno. The current definition and counterexample show why both directions matter; the spelling difference between historical vzaimno and current vzajemno remains visible.\par
\noindent{\footnotesize Izvorne zapisy / Original records: HIST-g081-052, HIST-g082-052.}\par
\smallskip\noindent\emph{Sravnjenje s tekstom:} Po definiciji V raboty 19 $\mathfrak R$ jest relativno prosty k $\mathfrak S$, ako iz $\mathfrak T\mathfrak R\equiv0\pmod{\mathfrak S}$ slěduje $\mathfrak T\equiv0\pmod{\mathfrak S}$. Podnožny priměr davaje $\mathfrak R=(x^2,y)$ i $\mathfrak S=(x)$: prvy směr jest pravdivy, obratny ne. Izraz vzajemno zato potrěbuje oba směry. Takvu definiciju ne možno tiho zaměniti običnym simetričnym pojmom vzajemnoj prostosti.\par
\noindent\emph{Textual comparison:} Definition V of Paper 19 says that $\mathfrak R$ is relatively prime to $\mathfrak S$ when $\mathfrak T\mathfrak R\equiv0\pmod{\mathfrak S}$ implies $\mathfrak T\equiv0\pmod{\mathfrak S}$. Its footnote gives $\mathfrak R=(x^2,y)$ and $\mathfrak S=(x)$: the first direction holds and the reverse fails. Mutual therefore requires both directions. The definition must not be silently replaced by the familiar symmetric notion of coprimality.\par
\noindent{\footnotesize TEXT-COMP-45; pasaže / passages: 13062, 13064, 13066; \texttt{CORPUS\_FORM\_COMPARISONS.json}.}\par
\end{minipage}\par
\par\medskip\Needspace{7\baselineskip}
\noindent\begin{minipage}{\linewidth}
\hypertarget{TERM-0689}{}
\noindent\textbf{vzaimna zavisimost form}\quad{\footnotesize\texttt{TERM-0689}}\par
\noindent\emph{Iztočnik / Source:} gegenseitige Abhängigkeit der Formen\par
\noindent\emph{Rusky / Russian:} взаимная зависимость форм\par
\noindent\emph{Ukrajinsky / Ukrainian:} взаємна залежність форм\par
\smallskip\noindent\emph{Razlog:} Vzaimna zavisimost form izražaje odnos medžu formami običnym pridavnikom i imenikom. Razlog daje prednost čitljivoj frazě pred umětnom složeninom, ale ne zapisuje konkretno odbijeno slovo. Zato takvu alternativu ne treba izmysliti. Fraza sama ne tvrdi, že zavisimost jest linearna, algebraična ili diferencijalna: to oprěděljaje razsmatrjano tvrđenje.\par
\noindent\emph{Reason:} Vzaimna zavisimost form expresses the relation between forms with an ordinary adjective and noun. The rationale favours a readable phrase over an artificial compound but records no particular rejected word. None should be invented. The phrase alone does not say whether the dependence is linear, algebraic or differential; that is determined by the statement under discussion.\par
\noindent{\footnotesize Izvorne zapisy / Original records: HIST-g029-004.}\par
\end{minipage}\par
\par\medskip\Needspace{7\baselineskip}
\noindent\begin{minipage}{\linewidth}
\hypertarget{TERM-0690}{}
\noindent\textbf{vhodi v diskriminantny ideal}\quad{\footnotesize\texttt{TERM-0690}}\par
\noindent\emph{Iztočnik / Source:} geht ... in dem Diskriminantenideal ... auf\par
\noindent\emph{Rusky / Russian:} входит в дискриминантный идеал\par
\noindent\emph{Ukrajinsky / Ukrainian:} входить до дискримінантного ідеалу\par
\smallskip\noindent\emph{Razlog:} Vhodi v diskriminantny ideal imaje tu vyraziti dělimost: prosty ideal nastupaje kak dělitelj diskriminantnogo ideala. Razlog izrično odbija tolkovanje kak običnu naležnost elementa množině ili prosty odnos sodržanja. Glagol vhodi bez konteksta može byti dvoznačny, zato slovnično objasnjenje mora imenovati točno relaciju. Aktualny dokaz svęzuje ju s lokalnym prostym elementom, ktory děli diskriminant.\par
\noindent\emph{Reason:} Vhodi v diskriminantny ideal here expresses divisibility: a prime ideal occurs as a divisor of the discriminant ideal. The rationale explicitly rejects reading it as ordinary membership or simply containment. Vhodi can be ambiguous without context, so the dictionary must name the relation precisely. The current proof links it to a local prime element dividing the discriminant.\par
\noindent{\footnotesize Izvorne zapisy / Original records: HIST-g163-003.}\par
\smallskip\noindent\emph{Sravnjenje s tekstom:} Teorem o diskriminantu v rabotě 31 i slědujuči lokalny argument svęzujut pojavjenje prostogo ideala v diskriminantnom ideale s pojavjenjem prostogo elementa v lokalnom diskriminantu. To podkrepljaje děliteljny smysl vhodi. Pasaže ne opravdavajut zaměnu na običnu naležnost elementa idealu.\par
\noindent\emph{Textual comparison:} The discriminant theorem in Paper 31 and its subsequent local argument connect a prime ideal's occurrence in the discriminant ideal with a prime element's occurrence in the local discriminant. This supports the divisibility sense of vhodi. The passages do not justify replacing it by ordinary membership of an element in an ideal.\par
\noindent{\footnotesize TEXT-COMP-46; pasaže / passages: 19071, 19105; \texttt{CORPUS\_FORM\_COMPARISONS.json}.}\par
\end{minipage}\par
\par\medskip\Needspace{7\baselineskip}
\noindent\begin{minipage}{\linewidth}
\hypertarget{TERM-0692}{}
\noindent\textbf{občna nuleva točka}\quad{\footnotesize\texttt{TERM-0692}}\par
\noindent\emph{Iztočnik / Source:} gemeinsame Nullstelle\par
\noindent\emph{Rusky / Russian:} общая нулевая точка\par
\noindent\emph{Ukrajinsky / Ukrainian:} спільний нуль / спільна нульова точка\par
\smallskip\noindent\emph{Razlog:} Občna nuleva točka sohranjaje geometričny smysl obče nuly plošnyh krivyh. Razlog izbira točka, da čitatelj vidi i algebraično zanulenje, i geometriju. To ne znači, že vsaka obča nula vo vseh rabotah mora byti nazvana točkoju; TERM-0693 pokazuje drugo kontekstualno rěšenje. Kratky zapis ne argumentuje osobno vsaku pravopisnu i sklonitvenu čest frazy.\par
\noindent\emph{Reason:} Občna nuleva točka preserves the geometric sense of a common zero of plane curves. The rationale chooses točka to retain both algebraic vanishing and geometry. It does not require every common zero in every paper to be called a point; TERM-0693 gives a different contextual choice. The short record does not separately justify every spelling and inflection in the phrase.\par
\noindent{\footnotesize Izvorne zapisy / Original records: HIST-g214-003.}\par
\end{minipage}\par
\par\medskip\Needspace{7\baselineskip}
\noindent\begin{minipage}{\linewidth}
\hypertarget{TERM-0693}{}
\noindent\textbf{obče nule}\quad{\footnotesize\texttt{TERM-0693}}\par
\noindent\emph{Iztočnik / Source:} gemeinsame Nullstellen\par
\noindent\emph{Rusky / Russian:} общие нули\par
\noindent\emph{Ukrajinsky / Ukrainian:} спільні нулі\par
\smallskip\noindent\emph{Razlog:} Obče nule koristi kratky algebraičny nazv nula vměsto obveznogo geometričnogo opisa točky. Razlog raboty 18 želi ne prědvarjati idealno-teoretično tolkovanje dodatnoju geometrijeju. Obče znači, že iste argumenty zanuljajut vse razsmatrjane polinomy, ne že vsak polinom imaje někaku svoju nulu. Historijna forma nule ostavaje odvojena od nuly pri TERM-0694.\par
\noindent\emph{Reason:} Obče nule uses the short algebraic term zero rather than requiring the geometric description point. The Paper 18 rationale avoids introducing extra geometry before the ideal-theoretic interpretation. Common means that the same arguments make all the polynomials vanish, not merely that each polynomial has a zero of its own. The historical form nule remains distinguished from nuly at TERM-0694.\par
\noindent{\footnotesize Izvorne zapisy / Original records: HIST-g080-002.}\par
\end{minipage}\par
\par\medskip\Needspace{7\baselineskip}
\noindent\begin{minipage}{\linewidth}
\hypertarget{TERM-0694}{}
\noindent\textbf{obče nuly}\quad{\footnotesize\texttt{TERM-0694}}\par
\noindent\emph{Iztočnik / Source:} gemeinsame Nullstellen\par
\noindent\emph{Rusky / Russian:} общие нули\par
\noindent\emph{Ukrajinsky / Ukrainian:} спільні нулі\par
\smallskip\noindent\emph{Razlog:} Obče nuly imaje isti osnovny algebraičny smysl kak TERM-0693, ale jest drugo zapisano sklonitveno rěšenje. Razlog raboty 22 daje prednost jezyku nul pred staršim opisom samo prěz točky. On ne sravnjaje nuly s nule i ne opravdavaje jih vzajemnu zaměnu. Aktualny uvod izrično govori o nulah vseh polinomov ideala nad algebraično zamknjenym nadpoljem; frazu ne treba suziti do realnyh toček.\par
\noindent\emph{Reason:} Obče nuly has the same basic algebraic meaning as TERM-0693 but records a different inflectional form. The Paper 22 rationale prefers the language of zeros to older point-only wording. It does not compare nuly with nule or justify treating them as interchangeable. The current introduction explicitly concerns zeros of all polynomials of an ideal over an algebraically closed extension; the phrase must not be restricted to real points.\par
\noindent{\footnotesize Izvorne zapisy / Original records: HIST-g093-004.}\par
\end{minipage}\par
\par\medskip\Needspace{7\baselineskip}
\noindent\begin{minipage}{\linewidth}
\hypertarget{TERM-0695}{}
\noindent\textbf{obći imenovatelj}\quad{\footnotesize\texttt{TERM-0695}}\par
\noindent\emph{Iztočnik / Source:} gemeinsamer Nenner\par
\noindent\emph{Rusky / Russian:} общий знаменатель\par
\noindent\emph{Ukrajinsky / Ukrainian:} спільний знаменник\par
\smallskip\noindent\emph{Razlog:} Obći imenovatelj zadržuje zajedni imenovatelj v dokazě i uslovje, že on ne izčezaje pri podstanovce. Razlog zato ne jest samo slovnična podobnost: bez togo uslovja racionalny izraz može postati neopreděljeny. Zapis sohranjaje formu obći; pri TERM-0696 jest zabilěžena obči. Te pravopisne formy ne predstavljajut dva različne matematične pojmy.\par
\noindent\emph{Reason:} Obći imenovatelj retains a common denominator in the proof and the requirement that it not vanish under substitution. The rationale is therefore more than lexical resemblance: without that condition, a rational expression can become undefined. This record preserves obći; TERM-0696 records obči. Those spellings do not represent different mathematical concepts.\par
\noindent{\footnotesize Izvorne zapisy / Original records: HIST-g042-006.}\par
\end{minipage}\par
\par\medskip\Needspace{7\baselineskip}
\noindent\begin{minipage}{\linewidth}
\hypertarget{TERM-0696}{}
\noindent\textbf{obči imenovatelj}\quad{\footnotesize\texttt{TERM-0696}}\par
\noindent\emph{Iztočnik / Source:} gemeinsamer Nenner\par
\noindent\emph{Rusky / Russian:} общий знаменатель\par
\noindent\emph{Ukrajinsky / Ukrainian:} спільний знаменник\par
\smallskip\noindent\emph{Razlog:} Obči imenovatelj nastavljaje istu terminologiju v odděle 5 raboty 6 i svęzuje ju s pozdnějšimi upotrěbami fiksovanogo imenovatelja. Obči i fiksovany ne sut sinonimy: prvy govori o zajednom prědstavjenju několikyh izrazov, vtory o tom, ktory imenovatelj ostavaje nezměnjeny v danom razsmotrenju. Razlog ne soderžuje čislovo sravnjenje alternativ, ale upira se na terminologičnu kontinuitetnost.\par
\noindent\emph{Reason:} Obči imenovatelj continues the same terminology in Paper 6 section 5 and links it to later fixed-denominator discussions. Common and fixed are not synonyms: one concerns a shared representation of several expressions, the other which denominator remains unchanged in a particular argument. The rationale rests on terminological continuity, not numerical comparison of alternatives.\par
\noindent{\footnotesize Izvorne zapisy / Original records: HIST-g045-015.}\par
\end{minipage}\par
\par\medskip\Needspace{7\baselineskip}
\noindent\begin{minipage}{\linewidth}
\hypertarget{TERM-0697}{}
\noindent\textbf{obči dělitelj}\quad{\footnotesize\texttt{TERM-0697}}\par
\noindent\emph{Iztočnik / Source:} gemeinsamer Teiler\par
\noindent\emph{Rusky / Russian:} общий делитель\par
\noindent\emph{Ukrajinsky / Ukrainian:} спільний дільник\par
\smallskip\noindent\emph{Razlog:} Obči dělitelj jest v odděle 8 raboty 6 oprěděljany vnutri danej oblasti. Ne vystarčaje, že polinom děli dva druge v někojem širšem polinomnom prstenu: k oblasti morajut naležati dělitelj i oba kvocienty. Razlog zadržuje obyčny nazv dělitelja s tym lokalnym uslovjem. Bez jeho izričnogo objasnjenja ta prosta fraza skryla by sučstvenu čest definicije.\par
\noindent\emph{Reason:} Obči dělitelj is defined internally to the given domain in Paper 6 section 8. It is not enough that a polynomial divides two others in some larger polynomial ring: the divisor and both quotients must belong to the domain. The rationale keeps familiar divisor vocabulary with this local condition. Without stating it explicitly, the simple phrase would conceal an essential part of the definition.\par
\noindent{\footnotesize Izvorne zapisy / Original records: HIST-g048-004.}\par
\smallskip\noindent\emph{Sravnjenje s tekstom:} V odděle 8 raboty 6 formula $F=LF_1$, $G=LG_1$ potrěbuje, da $L,F_1,G_1$ vse naležat k danej oblasti. V relativno cěloj oblasti odděla 9 potrěbno jest proveriti naležnost $L$; naležnost kvocientov togda slěduje iz uslovja 4a. Zato dodatek v oblasti jest matematično sučstveny, a kontinuitetnost nazva ne skriva proměnu uslovij.\par
\noindent\emph{Textual comparison:} In Paper 6 section 8, $F=LF_1$, $G=LG_1$ requires $L,F_1,G_1$ all to belong to the specified domain. In the relatively integral domain of section 9, membership of $L$ suffices: membership of the quotients then follows from condition 4a. In the domain is mathematically essential, and retaining the name must not conceal the changed conditions.\par
\noindent{\footnotesize TEXT-COMP-48; pasaže / passages: 5859, 5861, 5863, 6000, 6002; \texttt{CORPUS\_FORM\_COMPARISONS.json}.}\par
\end{minipage}\par
\par\medskip\Needspace{7\baselineskip}
\noindent\begin{minipage}{\linewidth}
\hypertarget{TERM-0698}{}
\noindent\textbf{obči dělitelj v $\mathfrak{G}$}\quad{\footnotesize\texttt{TERM-0698}}\par
\noindent\emph{Iztočnik / Source:} gemeinsamer Teiler in $\mathfrak{G}$\par
\noindent\emph{Rusky / Russian:} общий делитель в $\mathfrak{G}$\par
\noindent\emph{Ukrajinsky / Ukrainian:} спільний дільник у $\mathfrak{G}$\par
\smallskip\noindent\emph{Razlog:} Obči dělitelj v $\mathfrak G$ nastavljaje TERM-0697 v relativno cěloj oblasti. Odděl 9 može zjednodušiti proverku: k oblasti mora naležati sam dělitelj, a naležnost polinomnyh kvocientov už slěduje iz svojstva oblasti. To ne jest nova, nesvęzana definicija dělitelja. Bukva oblasti ostavaje matematičnym identifikatorom; ne zaměnjaje se obyčnym slovom grupa.\par
\noindent\emph{Reason:} Obči dělitelj v $\mathfrak G$ continues TERM-0697 in a relatively integral domain. Section 9 simplifies the check: the divisor must belong to the domain, while membership of the polynomial quotients follows from the domain's defining property. This is not an unrelated new definition of divisor. The domain letter remains a mathematical identifier, not an abbreviation to be translated as group.\par
\noindent{\footnotesize Izvorne zapisy / Original records: HIST-g049-008.}\par
\smallskip\noindent\emph{Sravnjenje s tekstom:} V odděle 8 raboty 6 formula $F=LF_1$, $G=LG_1$ potrěbuje, da $L,F_1,G_1$ vse naležat k danej oblasti. V relativno cěloj oblasti odděla 9 potrěbno jest proveriti naležnost $L$; naležnost kvocientov togda slěduje iz uslovja 4a. Zato dodatek v oblasti jest matematično sučstveny, a kontinuitetnost nazva ne skriva proměnu uslovij.\par
\noindent\emph{Textual comparison:} In Paper 6 section 8, $F=LF_1$, $G=LG_1$ requires $L,F_1,G_1$ all to belong to the specified domain. In the relatively integral domain of section 9, membership of $L$ suffices: membership of the quotients then follows from condition 4a. In the domain is mathematically essential, and retaining the name must not conceal the changed conditions.\par
\noindent{\footnotesize TEXT-COMP-48; pasaže / passages: 5859, 5861, 5863, 6000, 6002; \texttt{CORPUS\_FORM\_COMPARISONS.json}.}\par
\end{minipage}\par
\par\medskip\Needspace{7\baselineskip}
\noindent\begin{minipage}{\linewidth}
\hypertarget{TERM-0699}{}
\noindent\textbf{směšana grupa}\quad{\footnotesize\texttt{TERM-0699}}\par
\noindent\emph{Iztočnik / Source:} gemischte Gruppe\par
\noindent\emph{Rusky / Russian:} смешанная группа\par
\noindent\emph{Ukrajinsky / Ukrainian:} змішана група\par
\smallskip\noindent\emph{Razlog:} Směšana grupa označaje grupu transformacij zaviseču i od parametrov, i od libovoljnyh funkcij. Razlog daje prednost kratkomu strukturnomu nazvu, bez nove moderne parafrazy. V rabotě 13 te dva vida zavisimosti vedut k različnym relacijam, ktore v směšanom slučaju nastupajut zajedno. Pridavnik ne označaje proizvoljnu směs dvuh grup ili klasifikaciju abelovyh grup.\par
\noindent\emph{Reason:} Směšana grupa denotes a transformation group depending both on parameters and on arbitrary functions. The rationale favours a short structural name without adding a new modern paraphrase. In Paper 13 these kinds of dependence lead to different relations, which occur together in the mixed case. The adjective does not mean an arbitrary mixture of two groups or a classification of abelian groups.\par
\noindent{\footnotesize Izvorne zapisy / Original records: HIST-g075-006.}\par
\smallskip\noindent\emph{Sravnjenje s tekstom:} Uvod raboty 13 razlikuje parametry, libovoljne funkcije i beskonečno mnogo parametrov bez libovoljnyh funkcij. Směšana grupa imaje i parametry, i libovoljne funkcije. Pasaž poslě teoremov potvrđaje, že v tom slučaju nastupajut oba vida relacij. To jest lokalno značenje nazva, ne obča klasifikacija vseh grup.\par
\noindent\emph{Textual comparison:} Paper 13's introduction distinguishes parameters, arbitrary functions and infinitely many parameters without arbitrary functions. A mixed group has both parameters and arbitrary functions. The passage following the theorems confirms that both kinds of relations then arise. This is the local meaning of the term, not a general classification of all groups.\par
\noindent{\footnotesize TEXT-COMP-49; pasaže / passages: 9201, 9283; \texttt{CORPUS\_FORM\_COMPARISONS.json}.}\par
\end{minipage}\par
\par\medskip\Needspace{7\baselineskip}
\noindent\begin{minipage}{\linewidth}
\hypertarget{TERM-0700}{}
\noindent\textbf{cěly ręd form}\quad{\footnotesize\texttt{TERM-0700}}\par
\noindent\emph{Iztočnik / Source:} Gesamtformenreihe\par
\noindent\emph{Rusky / Russian:} полный ряд форм\par
\noindent\emph{Ukrajinsky / Ukrainian:} повний ряд форм\par
\smallskip\noindent\emph{Razlog:} Cěly ręd form razkladaje německu složeninu na čitlivu frazu: cěly označaje polny obseg, ręd form sam objekt. Razlog ne želi tvoriti neprozračno novo složeno slovo. Cěly tu ne znači algebraično integralny, kak v mnogih prethodnyh zapisah. Dlja točnogo smisla rěda form potrěbna jest konstrukcija raboty 2, objasnjena pri TERM-0564; ne ide avtomatično o stepenovy red.\par
\noindent\emph{Reason:} Cěly ręd form unpacks the German compound into a readable phrase: cěly marks the full extent and ręd form names the object. The rationale avoids an opaque new compound. Cěly does not mean algebraically integral here, as it does in many preceding entries. The exact meaning of the row of forms comes from Paper 2's construction, discussed at TERM-0564; it is not automatically a power series.\par
\noindent{\footnotesize Izvorne zapisy / Original records: HIST-g003-013.}\par
\end{minipage}\par
\par\medskip\Needspace{7\baselineskip}
\noindent\begin{minipage}{\linewidth}
\hypertarget{TERM-0701}{}
\noindent\textbf{obći ręd form}\quad{\footnotesize\texttt{TERM-0701}}\par
\noindent\emph{Iztočnik / Source:} Gesamtformenreihe\par
\noindent\emph{Rusky / Russian:} общий ряд форм\par
\noindent\emph{Ukrajinsky / Ukrainian:} загальний ряд форм\par
\smallskip\noindent\emph{Razlog:} Obći ręd form zadržuje imenik iz ranějšej terminologije i izražaje cěly obseg prěz obći. Razlog čita německo Gesamt- kak vkupny ili obči obseg, ne kak novo algebraično svojstvo. Sravnjenje s cěly ręd form pri TERM-0700 pokazuje dva historijne izbory. Zapis ne kaže, že oni byli izmerjeni po čitateljskoj razumlivosti, ni ne davaje osobny argument za ć protiv č.\par
\noindent\emph{Reason:} Obći ręd form keeps the established noun phrase and expresses total scope through obći. The rationale reads German Gesamt- as total or general extent, not a new algebraic property. Comparison with cěly ręd form at TERM-0700 reveals two historical choices. The record provides neither a reader-comprehension measurement nor a separate argument for ć rather than č.\par
\noindent{\footnotesize Izvorne zapisy / Original records: HIST-g012-008.}\par
\end{minipage}\par
\par\medskip\Needspace{7\baselineskip}
\noindent\begin{minipage}{\linewidth}
\hypertarget{TERM-0702}{}
\noindent\textbf{vkupna vaga formy}\quad{\footnotesize\texttt{TERM-0702}}\par
\noindent\emph{Iztočnik / Source:} Gesamtgewicht der Form\par
\noindent\emph{Rusky / Russian:} общий вес формы\par
\noindent\emph{Ukrajinsky / Ukrainian:} сумарна вага форми\par
\smallskip\noindent\emph{Razlog:} Vkupna vaga formy imenuje sumu matematičnyh vag pojedinyh redov prěměnnyh. Razlog svęzuje nazv s dokazom konečnosti poslě formuly (45). Vkupna zato ima točno aritmetično značenje, ne samo vpečatljenje važnosti. Ta vaga ne jest čislova ocěnka pri izborě slova v prevodničkoj metodě; ona jest veličina definirana v samoj matematičnoj rabotě. Definicija i slovnično sravnjenje sut navedene pri TERM-0709.\par
\noindent\emph{Reason:} Vkupna vaga formy names the sum of the mathematical weights of the individual rows of variables. The rationale connects the term to the finiteness argument following equation (45). Vkupna therefore has a precise arithmetic meaning, not a general impression of importance. This weight is not a numerical score used to choose a translation: it is a quantity defined in the mathematical work itself. The definition and dictionary comparison are given at TERM-0709.\par
\noindent{\footnotesize Izvorne zapisy / Original records: HIST-g035-010.}\par
\end{minipage}\par
\par\medskip\Needspace{7\baselineskip}
\noindent\begin{minipage}{\linewidth}
\hypertarget{TERM-0703}{}
\noindent\textbf{cělokupnost}\quad{\footnotesize\texttt{TERM-0703}}\par
\noindent\emph{Iztočnik / Source:} Gesamtheit aller / Gesamtheit der Restklassen\par
\noindent\emph{Rusky / Russian:} совокупность\par
\noindent\emph{Ukrajinsky / Ukrainian:} сукупність\par
\smallskip\noindent\emph{Razlog:} Cělokupnost znači kolekciju vseh razsmatrjanyh objektov, ne jih algebraičnu sumu. Razlog izrično odbija ranějše vsota po toj smyslovoj pričině. Razlika jest važna osoblivo za klasy ostatkov: kolekcija klasov ne jest samo operacija jih složenja. Odlučajuči kriterij jest razlika medžu kolekcijeju i operacijeju nad jej elementami.\par
\noindent\emph{Reason:} Cělokupnost denotes the collection of all the objects under consideration, not their algebraic sum. The rationale explicitly rejects the earlier vsota on that semantic ground. The distinction is especially important for residue classes: their collection is not merely the operation of adding them. The decisive criterion is the distinction between a collection and an operation on its members.\par
\noindent{\footnotesize Izvorne zapisy / Original records: HIST-g079-016.}\par
\end{minipage}\par
\par\medskip\Needspace{7\baselineskip}
\noindent\begin{minipage}{\linewidth}
\hypertarget{TERM-0704}{}
\noindent\textbf{cělost vseh funkcij}\quad{\footnotesize\texttt{TERM-0704}}\par
\noindent\emph{Iztočnik / Source:} Gesamtheit aller Funktionen\par
\noindent\emph{Rusky / Russian:} совокупность всех функций\par
\noindent\emph{Ukrajinsky / Ukrainian:} сукупність усіх функцій\par
\smallskip\noindent\emph{Razlog:} Cělost vseh funkcij izražaje totalnost kolekcije funkcij. Razlog nastavljaje ranějši široky jezyk množin, bez novogo simvola. Vseh nosi kvantifikaciju: ne samo někake funkcije. Cělost tu ne tvrdi algebraičnu cělost pojedinyh funkcij; lokalny kontekst mora razlikovati ta dva upotrěbljenja istogo korenja.\par
\noindent\emph{Reason:} Cělost vseh funkcij expresses the totality of a collection of functions. The rationale continues earlier collection language without introducing new notation. Vseh carries the quantification: all, not merely some functions. Cělost here does not assert that the individual functions are algebraically integral; context must distinguish these uses of the same stem.\par
\noindent{\footnotesize Izvorne zapisy / Original records: HIST-g043-006.}\par
\end{minipage}\par
\par\medskip\Needspace{7\baselineskip}
\noindent\begin{minipage}{\linewidth}
\hypertarget{TERM-0705}{}
\noindent\textbf{cělokupnost polinomov}\quad{\footnotesize\texttt{TERM-0705}}\par
\noindent\emph{Iztočnik / Source:} Gesamtheit der Polynome\par
\noindent\emph{Rusky / Russian:} совокупность полиномов\par
\noindent\emph{Ukrajinsky / Ukrainian:} сукупність поліномів\par
\smallskip\noindent\emph{Razlog:} Cělokupnost polinomov zadržuje smysl kolekcije v proze raboty 6. Razlog ne želi dodavati sovrěmennu množinovu notaciju tam, kde avtor jeju ne koristi. To ne jest argument protiv pojma množiny, ale za věrno predstavjenje izloženja. Ktore polinomy sodrživajut se v kolekciji, mora byti dany uslovjami pasaža, ne samym slovom cělokupnost.\par
\noindent\emph{Reason:} Cělokupnost polinomov retains the collection sense in Paper 6's prose. The rationale avoids adding modern set notation where the author does not use it. This is an argument for faithful presentation, not against the concept of a set. The passage's conditions must specify which polynomials belong; cělokupnost alone cannot do so.\par
\noindent{\footnotesize Izvorne zapisy / Original records: HIST-g049-004.}\par
\end{minipage}\par
\par\medskip\Needspace{7\baselineskip}
\noindent\begin{minipage}{\linewidth}
\hypertarget{TERM-0707}{}
\noindent\textbf{zakon neograničenogo i jednoznačnogo odčitanja}\quad{\footnotesize\texttt{TERM-0707}}\par
\noindent\emph{Iztočnik / Source:} Gesetz der unbeschränkten und eindeutigen Subtraktion\par
\noindent\emph{Rusky / Russian:} закон неограниченного и однозначного вычитания\par
\noindent\emph{Ukrajinsky / Ukrainian:} закон необмеженого й однозначного віднімання\par
\smallskip\noindent\emph{Razlog:} Zakon neograničenogo i jednoznačnogo odčitanja zadržuje oba zahtěva historijnoj aksiomy. Dlja vsakyh dopustimyh $a,b$ uravnenje $a+x=b$ mora iměti rěšenje, i to rěšenje mora byti jedino. Razlog ne zaměnjaje celu aksiomu sovrěmennym skračenjem o obratnyh elementah. Povtorjenje v definiciji modulnoj oblasti sohranjaje isti smysl, hoti elementy i oznaky sut druge.\par
\noindent\emph{Reason:} Zakon neograničenogo i jednoznačnogo odčitanja retains both requirements of the historical axiom. For every admissible $a,b$, the equation $a+x=b$ must have a solution, and that solution must be unique. The rationale does not replace the full axiom with a modern abbreviation about inverses. Its recurrence in the module-domain definition retains the same meaning with different elements and notation.\par
\noindent{\footnotesize Izvorne zapisy / Original records: HIST-g081-020, HIST-g082-020.}\par
\smallskip\noindent\emph{Sravnjenje s tekstom:} Obě pojavjenja aksiomy raboty 19 govoret o jedinom rěšenju aditivnogo uravnenja: v prstenu $a+x=b$, v modulnoj oblasti $\alpha+\xi=\beta$. Neograničeno znači, že rěšenje jest dostupno za vse navedene pary, ne čislovu neograničenost. Tekst odvojeno kaže, že prsten ne mora iměti multiplikativnu jedinicu.\par
\noindent\emph{Textual comparison:} The two occurrences of the axiom in Paper 19 require a unique solution of an additive equation: $a+x=b$ in the ring and $\alpha+\xi=\beta$ in the module domain. Unrestricted means solvability for all the specified pairs, not numerical unboundedness. The text separately notes that the ring need not have a multiplicative identity.\par
\noindent{\footnotesize TEXT-COMP-50; pasaže / passages: 12483, 12485, 13372; \texttt{CORPUS\_FORM\_COMPARISONS.json}.}\par
\end{minipage}\par
\par\medskip\Needspace{7\baselineskip}
\noindent\begin{minipage}{\linewidth}
\hypertarget{TERM-0708}{}
\noindent\textbf{obyčno linearno diferencialno uravnenje 2-go reda}\quad{\footnotesize\texttt{TERM-0708}}\par
\noindent\emph{Iztočnik / Source:} gew?hnliche lineare Differentialgleichung 2. Ordnung\par
\noindent\emph{Rusky / Russian:} обыкновенное линейное дифференциальное уравнение 2-го порядка\par
\noindent\emph{Ukrajinsky / Ukrainian:} звичайне лінійне диференціальне рівняння 2-го порядку\par
\smallskip\noindent\emph{Razlog:} Obyčno linearno diferencialno uravnenje 2-go reda sohranjaje tri nezavisne kvalifikacije: obyčno, linearno i drugogo reda. Razlog želi jasno vozvratiti čitatelja od obče modulnoj teorije k klasičnomu priměru. Red označaje red diferencialnogo uravnenja, ne stepen polinoma ni čislo uravnenij. V německom zaglavju zapisa jest povrěđena bukva; njej zmysl možno opisati, ale ne vydavati nezapisanu rekonstrukciju za zachovany original.\par
\noindent\emph{Reason:} Obyčno linearno diferencialno uravnenje 2-go reda retains three independent qualifications: ordinary, linear and second order. The rationale makes the return from general module theory to a classical example explicit. Order is the order of the differential equation, not polynomial degree or the number of equations. The record's German heading has a damaged letter; its intended sense can be explained without presenting an unrecorded restoration as the surviving original.\par
\noindent{\footnotesize Izvorne zapisy / Original records: HIST-g079-071.}\par
\end{minipage}\par
\par\medskip\Needspace{7\baselineskip}
\noindent\begin{minipage}{\linewidth}
\hypertarget{TERM-0709}{}
\noindent\textbf{vaga rěda / indeks vagy}\quad{\footnotesize\texttt{TERM-0709}}\par
\noindent\emph{Iztočnik / Source:} Gewicht der Reihe / Gewichtsindex\par
\noindent\emph{Rusky / Russian:} вес строки / весовой индекс\par
\noindent\emph{Ukrajinsky / Ukrainian:} вага рядка / ваговий індекс\par
\smallskip\noindent\emph{Razlog:} Vaga rěda / indeks vagy svęzuje dve veličiny iz odděla 1 raboty 4. Vaga reda jest $\rho(n-\rho)$; gorny i dolny indeksy sut $\rho$ i $n-\rho$, ne sama ta vaga. Historijny razlog predlagaje vaga radi slovjanskoj razumlivosti, ale ne soderžuje čitateljsky test. Nověje sravnjenje slovnika potvrđaje osnovno značenje vaga kak weight; to samo po sebi ješče ne dokazuje matematičnu upotrěbu.\par
\noindent\emph{Reason:} Vaga rěda / indeks vagy links two quantities from Paper 4 section 1. The row weight is $\rho(n-\rho)$; the upper and lower indices are $\rho$ and $n-\rho$, not the weight itself. The historical rationale proposes vaga for Slavic intelligibility but supplies no reader test. A fresh dictionary comparison confirms the ordinary weight sense of vaga; that alone does not attest its mathematical use.\par
\noindent{\footnotesize Izvorne zapisy / Original records: HIST-g030-014.}\par
\smallskip\noindent\emph{Sravnjenje s tekstom:} Rabota 4 oprěděljaje vagu reda kak $\rho(n-\rho)$, indeksa kak $\rho$ i $n-\rho$, i vkupnu vagu kak sumu vag redov. Formula za umenšenje podporaje dokaz konečnosti. Medžuslovjansky slovnik razlikuje vaga v značenju weight (zapis 5645) od vaga v značenju scales (22173); to jest leksična podpora, ne samostojna attestacija invariantnoj teorije. Gramatika imenikov opisuje tvrde i mekke osnovy i daje genitiv slugy od sluga. Pri regularnom sklonjenju vaga tomu odgovarja genitiv vagy, ne vagi. To jest novo gramatyčno sravnjenje; historijne formy pri TERM-0710 i TERM-0713 ne sut tiho proměnjene. Nema izměrjene čitateljskoj razumlivosti ni čislovyh težin prevodničkogo izbora.\par
\noindent\emph{Textual comparison:} Paper 4 defines row weight as $\rho(n-\rho)$, the indices as $\rho$ and $n-\rho$, and total weight as the sum of row weights. Its decrease formula supports the finiteness argument. The Interslavic Dictionary distinguishes vaga meaning weight (entry 5645) from vaga meaning scales (22173); this supports the ordinary lexical sense, not an independent attestation in invariant theory. The noun grammar describes hard and soft stems and gives genitive slugy from sluga. Under regular declension, vaga correspondingly has genitive vagy rather than vagi. This is a fresh grammatical comparison; the historical forms at TERM-0710 and TERM-0713 are not silently changed. No reader-comprehension measurement or numerical translation-selection weights are supplied.\par
\noindent{\footnotesize TEXT-COMP-51; pasaže / passages: 3784, 4433, 4725; \texttt{CORPUS\_FORM\_COMPARISONS.json}.}\par
\end{minipage}\par
\par\medskip\Needspace{7\baselineskip}
\noindent\begin{minipage}{\linewidth}
\hypertarget{TERM-0710}{}
\noindent\textbf{umenšenje vagi}\quad{\footnotesize\texttt{TERM-0710}}\par
\noindent\emph{Iztočnik / Source:} Gewichtsabnahme\par
\noindent\emph{Rusky / Russian:} убывание веса\par
\noindent\emph{Ukrajinsky / Ukrainian:} спадання ваги\par
\smallskip\noindent\emph{Razlog:} Umenšenje vagi znači pad matematičnoj vagy pri opisanoj operaciji. Razlog svęzuje odděl 9 s dokazom odděla 8, gde ta proměna služi počitanju kontrakcij. Sam nazv ne oprěděljaje ni velikost, ni znak proměny v vsakoj možnoj situaciji: to musi dokazati formula. Vagi jest historijny zapis; jejo odnos k regularnomu genitivu vagy jest razsmotren pri TERM-0709.\par
\noindent\emph{Reason:} Umenšenje vagi means a decrease of mathematical weight under the specified operation. The rationale links section 9 to section 8's use of that change to count contractions. The name alone determines neither the magnitude nor the sign in every possible situation; those depend on the formula. Vagi is the historical spelling; its relation to the regular genitive vagy is discussed at TERM-0709.\par
\noindent{\footnotesize Izvorne zapisy / Original records: HIST-g038-011.}\par
\end{minipage}\par
\par\medskip\Needspace{7\baselineskip}
\noindent\begin{minipage}{\linewidth}
\hypertarget{TERM-0711}{}
\noindent\textbf{umenšenje vagi / vkupna vaga}\quad{\footnotesize\texttt{TERM-0711}}\par
\noindent\emph{Iztočnik / Source:} Gewichtsabnahme / Gesamtgewicht\par
\noindent\emph{Rusky / Russian:} убывание веса / общий вес\par
\noindent\emph{Ukrajinsky / Ukrainian:} спадання ваги / сумарна вага\par
\smallskip\noindent\emph{Razlog:} Umenšenje vagi / vkupna vaga razlikuje proměnu veličiny od jejo ukupnogo značenja. Razlog svęzuje tu aritmetiku s čislom osnovnyh kontrakcij v odděle 8 raboty 4. Kratke frazy imajut pokazati strukturu argumenta, ne zavesti novu fizikalnu metaforu. Oznaka, že medžuslovjanske formy potrěbujut daljše razsmotrenje, ne jest svědočstvo, že specialist už jih proveril. Sravnjenje pri TERM-0709 uključaje definiciju i formulu umenšenja.\par
\noindent\emph{Reason:} Umenšenje vagi / vkupna vaga distinguishes a change in a quantity from its total value. The rationale connects this arithmetic to the number of elementary contractions in Paper 4 section 8. The short phrases are intended to expose the argument's structure, not introduce a new physical metaphor. The historical request for further consideration of the forms is not evidence that a specialist has checked them. TERM-0709 includes the definition and decrease formula.\par
\noindent{\footnotesize Izvorne zapisy / Original records: HIST-g037-006.}\par
\end{minipage}\par
\par\medskip\Needspace{7\baselineskip}
\noindent\begin{minipage}{\linewidth}
\hypertarget{TERM-0712}{}
\noindent\textbf{uslovja věsy / uslovja simetrije}\quad{\footnotesize\texttt{TERM-0712}}\par
\noindent\emph{Iztočnik / Source:} Gewichtsbedingungen / Symmetriebedingungen\par
\noindent\emph{Rusky / Russian:} условия веса / условия симметрии\par
\noindent\emph{Ukrajinsky / Ukrainian:} умови ваги / умови симетрії\par
\smallskip\noindent\emph{Razlog:} Uslovja věsy / uslovja simetrije označaje dve grupy ograničenij iz invariantnoj teorije. Razlog sohranjaje paru vaga i simetrija kak pojmovno čitlivu strukturu dokaza, ne kak opis fizikalnoj těžiny. On ne objašnjavaje osobno izbor korenja věs- naměsto vaga. Slovnik zato ne mora skrivati rozděl medžu tym historijnym prědlogom i formami pri TERM-0709, ni proglasiti jih bez dodatnogo osnovanja ravno doporučanymi.\par
\noindent\emph{Reason:} Uslovja věsy / uslovja simetrije denotes two sets of constraints in invariant theory. The rationale preserves weight and symmetry as a readable conceptual pairing, not a description of physical heaviness. It does not separately justify věs- instead of vaga. The dictionary must therefore retain the difference between this historical proposal and TERM-0709, rather than presenting them as equally recommended without further evidence.\par
\noindent{\footnotesize Izvorne zapisy / Original records: HIST-g077-006.}\par
\end{minipage}\par
\par\medskip\Needspace{7\baselineskip}
\noindent\begin{minipage}{\linewidth}
\hypertarget{TERM-0713}{}
\noindent\textbf{indeks vagi}\quad{\footnotesize\texttt{TERM-0713}}\par
\noindent\emph{Iztočnik / Source:} Gewichtsindex\par
\noindent\emph{Rusky / Russian:} весовой индекс\par
\noindent\emph{Ukrajinsky / Ukrainian:} ваговий індекс\par
\smallskip\noindent\emph{Razlog:} Indeks vagi zadržuje kratku frazu s imenikom indeks i genitivnym dopolnjenjem. Razlog izrično odbija težšu složeninu, ale ne daje jejo konkretno slovo. Smysl indeksa slěduje iz ranějšego oprěděljenja gornogo i dolnogo indeksa; on ne jest indeks čitateljskoj razumlivosti ni metodična ocěnka. Forma vagi ostavaje zachovana kak svědočstvo izbora, s gramatyčnym sravnjenjem pri TERM-0709.\par
\noindent\emph{Reason:} Indeks vagi retains a short phrase consisting of indeks and a genitive complement. The rationale rejects a heavier compound without recording a particular competing word. The index's meaning comes from the earlier upper/lower-index definition; it is neither a reader-intelligibility index nor a methodological score. Vagi remains preserved as evidence of the choice, with the grammatical comparison at TERM-0709.\par
\noindent{\footnotesize Izvorne zapisy / Original records: HIST-g037-007.}\par
\end{minipage}\par
\par\medskip\Needspace{7\baselineskip}
\noindent\begin{minipage}{\linewidth}
\hypertarget{TERM-0714}{}
\noindent\textbf{indeksy vagy}\quad{\footnotesize\texttt{TERM-0714}}\par
\noindent\emph{Iztočnik / Source:} Gewichtsindizes\par
\noindent\emph{Rusky / Russian:} весовые индексы\par
\noindent\emph{Ukrajinsky / Ukrainian:} вагові індекси\par
\smallskip\noindent\emph{Razlog:} Indeksy vagy jest množinska fraza, ktora prodolžuje terminologiju odděla 1 v odděle 3 raboty 4. Razlog argumentuje kontinuitetnost, ne novy nezavisny izbor slova. Vagy jest genitiv od vaga, a ne drugy nazv dlja pojedinogo indeksa. Sravnjenje pri TERM-0709 daje odvojeno leksično i gramatyčno osnovanje; ne treba pripisati to novějše čitanje originalnomu kratkomu zapisu.\par
\noindent\emph{Reason:} Indeksy vagy is the plural phrase continuing section 1's terminology in Paper 4 section 3. The rationale argues for continuity rather than a new independent lexical selection. Vagy is the genitive of vaga, not a different name for one index. TERM-0709 supplies separate lexical and grammatical grounds; that fresh consultation must not be attributed to the original brief record.\par
\noindent{\footnotesize Izvorne zapisy / Original records: HIST-g032-010.}\par
\end{minipage}\par
\par\medskip\Needspace{7\baselineskip}
\noindent\begin{minipage}{\linewidth}
\hypertarget{TERM-0715}{}
\noindent\textbf{ten sam ostatok / polny prěsěk}\quad{\footnotesize\texttt{TERM-0715}}\par
\noindent\emph{Iztočnik / Source:} gleiches Rest / vollständiger Schnitt\par
\noindent\emph{Rusky / Russian:} один и тот же остаток / полное пересечение\par
\noindent\emph{Ukrajinsky / Ukrainian:} той самий залишок / повний перетин\par
\smallskip\noindent\emph{Razlog:} Ten sam ostatok / polny prěsěk sohranjaje geometričny opis, ktory možno povezati s ekvivalentnymi idealami. Ten sam označaje identičnost ostatka, ne prosto podobnost dvuh objektov. Razlog ne daje prava zjednodušiti polny prěsěk do vsakego proizvoljnogo prěsěka. Oba izrazy potrěbujut geometričny kontekst raboty 14; zapis ne izlagaje osobnu sovrěmennu definiciju polnogo prěsěka.\par
\noindent\emph{Reason:} Ten sam ostatok / polny prěsěk retains geometric wording that can be related to equivalent ideals. Ten sam asserts identity of the residual object, not merely resemblance. The rationale does not license reducing complete intersection to any arbitrary intersection. Both expressions need Paper 14's geometric context; the record does not provide a separate modern definition of complete intersection.\par
\noindent{\footnotesize Izvorne zapisy / Original records: HIST-g076-022.}\par
\end{minipage}\par
\par\medskip\Needspace{7\baselineskip}
\noindent\begin{minipage}{\linewidth}
\hypertarget{TERM-0716}{}
\noindent\textbf{odnosi ravnosti i odnosi operacij}\quad{\footnotesize\texttt{TERM-0716}}\par
\noindent\emph{Iztočnik / Source:} Gleichheits- und Verknüpfungsbeziehungen\par
\noindent\emph{Rusky / Russian:} отношения равенства и операционные отношения\par
\noindent\emph{Ukrajinsky / Ukrainian:} відношення рівності та операційні відношення\par
\smallskip\noindent\emph{Razlog:} Odnosi ravnosti i odnosi operacij razlikuje, ktore elementy sut ravne i kak operacije medžu nimi sut oprěděljene. Razlog daje prednost operacij pred širokymi svezami ili kompozicijami. V aktualnom pasažu o $R[S]$ te odnosi naslědujut se iz okolnogo prstena $T$. Povtorjenje odnosi čini obe česti frazy javnymi; ne govori o dvuh neodvisnyh strukturah bez vzajemnoj svezi.\par
\noindent\emph{Reason:} Odnosi ravnosti i odnosi operacij distinguishes which elements are equal from how operations on them are determined. The rationale prefers operations to broad connections or compositions. In the current passage on $R[S]$, these relations are inherited from the ambient ring $T$. Repeating odnosi makes both parts explicit; it does not posit two unrelated structures.\par
\noindent{\footnotesize Izvorne zapisy / Original records: HIST-g119-006.}\par
\smallskip\noindent\emph{Sravnjenje s tekstom:} Pri konstrukciji $R[S]$ v rabotě 31 odnosi ravnosti i operacij sut dany tymi iz $T$. Pasaž takože kaže, začto někdy vystarčajut linearne formy vměsto eksplicitnogo tvorenja vsih polinomov. Zato operaciji jest konkretno algebraično slovo, ne opće ime dlja svezi medžu objektami.\par
\noindent\emph{Textual comparison:} In Paper 31's construction of $R[S]$, equality and operation relations are determined by those in $T$. The passage also explains why linear forms sometimes suffice instead of explicitly forming every polynomial. Operations is therefore a concrete algebraic designation, not a general word for connections between objects.\par
\noindent{\footnotesize TEXT-COMP-53; pasaže / passages: 17316; \texttt{CORPUS\_FORM\_COMPARISONS.json}.}\par
\end{minipage}\par
\par\medskip\Needspace{7\baselineskip}
\noindent\begin{minipage}{\linewidth}
\hypertarget{TERM-0717}{}
\noindent\textbf{uravnenje najnižšego stepena}\quad{\footnotesize\texttt{TERM-0717}}\par
\noindent\emph{Iztočnik / Source:} Gleichung niedrigsten Grades\par
\noindent\emph{Rusky / Russian:} уравнение наименьшей степени\par
\noindent\emph{Ukrajinsky / Ukrainian:} рівняння найменшого степеня\par
\smallskip\noindent\emph{Razlog:} Uravnenje najnižšego stepena ostavaje pri jeziku uravnenij i izrično pokazuje minimalnost stepena nad poljem $P$. Razlog ne želi zaměniti tu frazu samym pridavnikom nerazložimo. V konkretnoj upotrěbě minimalnost oprěděljaje, koliko poteg elementa može služiti kak linearno nezavisne elementy. Ne treba uspoređivati stepeny nad raznymi osnovnymi poljami bez ukazanja toj proměny.\par
\noindent\emph{Reason:} Uravnenje najnižšego stepena retains equation language and explicitly states minimal degree over the field $P$. The rationale avoids replacing the phrase with irreducible alone. In the particular passage, minimality determines how many powers of an element are linearly independent. Degrees over different base fields must not be compared without making that change explicit.\par
\noindent{\footnotesize Izvorne zapisy / Original records: HIST-g128-005.}\par
\smallskip\noindent\emph{Sravnjenje s tekstom:} Pasaž raboty 31 vybiraje element $\gamma$ drugogo roda i uravnenje najnižšego stepena nad $P$, pri harakteristice $p$. Daljši argument koristi linearno nezavisne potegy $\gamma$. Fraza najnižšego stepena zato nosi konkretno uslovje, ktoro ne treba izgubiti pri zaměně na opći nazv uravnenja.\par
\noindent\emph{Textual comparison:} The Paper 31 passage chooses an element $\gamma$ of the second kind and an equation of least degree over $P$, in characteristic $p$. The following argument uses linearly independent powers of $\gamma$. Least degree therefore carries a specific condition that must not disappear in a more general equation label.\par
\noindent{\footnotesize TEXT-COMP-54; pasaže / passages: 17737, 17739, 17741; \texttt{CORPUS\_FORM\_COMPARISONS.json}.}\par
\end{minipage}\par
\par\medskip\Needspace{7\baselineskip}
\noindent\begin{minipage}{\linewidth}
\hypertarget{TERM-0718}{}
\noindent\textbf{uravnenja s zadanoj grupoj}\quad{\footnotesize\texttt{TERM-0718}}\par
\noindent\emph{Iztočnik / Source:} Gleichungen mit vorgeschriebener Gruppe\par
\noindent\emph{Rusky / Russian:} уравнения с заданной группой\par
\noindent\emph{Ukrajinsky / Ukrainian:} рівняння із заданою групою\par
\smallskip\noindent\emph{Razlog:} Uravnenja s zadanoj grupoj izražaje konstrukcijnu zadaču: grupa jest zadana, a uravnenja trěba najdti s toju grupoju. Razlog bira obyčno sovrěmenno matematično zadany vměsto staršego doslovnogo predpisany. Ne idě o grupu samyh uravnenij, i ne dovoljno prosto ukazati uravnenja bez svojstva, ktoro ih grupa mora iměti. Glagolna naměra zadači ostavaje v naslovnoj frazě sohranjena.\par
\noindent\emph{Reason:} Uravnenja s zadanoj grupoj expresses a construction problem: the group is specified and equations with that group are sought. The rationale selects ordinary modern mathematical given rather than the older literal prescribed. This is not a group of equations, nor is it sufficient merely to exhibit equations without the required group. The title phrase preserves the direction of the problem.\par
\noindent{\footnotesize Izvorne zapisy / Original records: HIST-g073-001.}\par
\end{minipage}\par
\par\medskip\Needspace{7\baselineskip}
\noindent\begin{minipage}{\linewidth}
\hypertarget{TERM-0719}{}
\noindent\textbf{členy najvyššej dimenzije}\quad{\footnotesize\texttt{TERM-0719}}\par
\noindent\emph{Iztočnik / Source:} Glieder höchster Dimension\par
\noindent\emph{Rusky / Russian:} старшие однородные части\par
\noindent\emph{Ukrajinsky / Ukrainian:} старші однорідні частини\par
\smallskip\noindent\emph{Razlog:} Členy najvyššej dimenzije jest prěvod blizky iztočniku, dokolě ruska i ukrajinska paralely govoret izrično o vedučih homogenyh čestjah. Dva razloga nastavljajut to rěšenje medžu oddělami 11 i 12 raboty 6. Dimenzija tu ne mora byti čitana kak dimenzija prostora; govorit se o najvyššem stepenju polinomnyh členov. Slovnik mora pokazati tu razliku medžu doslovnym nazvom i sovrěmennym objasnjenjem, ne tiho zaměniti historijny tekst.\par
\noindent\emph{Reason:} Členy najvyššej dimenzije is a source-close rendering, whereas the Russian and Ukrainian comparanda explicitly refer to leading homogeneous parts. Two rationales continue this choice between sections 11 and 12 of Paper 6. Dimension here must not be read as the dimension of a space; it concerns the highest degree of polynomial terms. The dictionary needs to show the relation between the literal name and its modern explanation, not silently substitute for the historical text.\par
\noindent{\footnotesize Izvorne zapisy / Original records: HIST-g051-004, HIST-g052-008.}\par
\end{minipage}\par
\par\medskip\Needspace{7\baselineskip}
\noindent\begin{minipage}{\linewidth}
\hypertarget{TERM-0720}{}
\noindent\textbf{najve??i ob?i d?litelj / pr?s?k}\quad{\footnotesize\texttt{TERM-0720}}\par
\noindent\emph{Iztočnik / Source:} gr??ter gemeinsamer Teiler / Durchschnitt\par
\noindent\emph{Rusky / Russian:} [Bukvy ne sut zachovane / Lettering not recoverable]\par
\noindent\emph{Ukrajinsky / Ukrainian:} [Bukvy ne sut zachovane / Lettering not recoverable]\par
\smallskip\noindent\emph{Razlog:} Zapis največšego občego dělitelja i prěsěka ima povrěđene bukvy vo vseh jezykovyh poljah. Zachovany razlog vseže izrično sohranjaje oba pojmy: v rabotě 9 največši obči dělitelj razsmatrjanyh oblastij jest jih prěsěk. V odděle 5 toj prěsěk jest $[H]$, ktoro samo ne jest dopustima oblast $\mathfrak G$. Zato ne slěduje zamknjenost familije pri prěsěkah; to jest sučstvena razlika od odděla 4. Aktualny pasaž podkrepljaje pojmovny razlog, ale ne vosstanavlja izgubljeno pisanje starogo zapisa.\par
\noindent\emph{Reason:} The record for greatest common divisor and intersection has damaged lettering in every language field. Its surviving rationale nevertheless explicitly retains both concepts: the greatest common divisor of the domains under discussion in Paper 9 is their intersection. In section 5 that intersection is $[H]$, which is not itself an admissible domain $\mathfrak G$. Closure of the family under intersections therefore does not follow; this is an essential contrast with section 4. The current passage supports the conceptual reason but does not restore the lost spelling of the old record.\par
\noindent{\footnotesize Izvorne zapisy / Original records: HIST-g063-004.}\par
\smallskip\noindent\emph{Sravnjenje s tekstom:} V odděle 5 raboty 9 familija $\mathfrak M(\mathfrak G)$ objedinjuje oblasti, ktore sodrživajut vybranu algebraičnu bazu $H$. Jih največši obči dělitelj i prěsěk jest $[H]$, ale $[H]$ samo ne jest oblast $\mathfrak G$. Zaključenje izrično kaže, že familija ne jest zamknjena pri prěsěkah. Pasaž odděla 4 s inoj familijeju $\mathfrak M(H)$ ne može byti podstavjen vměsto togo rezultata.\par
\noindent\emph{Textual comparison:} In Paper 9 section 5, the family $\mathfrak M(\mathfrak G)$ groups domains containing the chosen algebraic basis $H$. Their greatest common divisor and intersection is $[H]$, but $[H]$ is not itself a domain $\mathfrak G$. The conclusion explicitly says that the family is not closed under intersections. Section 4's passage about the different family $\mathfrak M(H)$ cannot substitute for this result.\par
\noindent{\footnotesize TEXT-COMP-52; pasaže / passages: 7596, 7597, 7598, 7603, 7604, 7605, 7606, 7607, 7608, 7676, 7677, 7678; \texttt{CORPUS\_FORM\_COMPARISONS.json}.}\par
\end{minipage}\par
\par\medskip\Needspace{7\baselineskip}
\noindent\begin{minipage}{\linewidth}
\hypertarget{TERM-0721}{}
\noindent\textbf{najve??i ob?i d?litelj ili pr?s?k}\quad{\footnotesize\texttt{TERM-0721}}\par
\noindent\emph{Iztočnik / Source:} gr??ter gemeinsamer Teiler oder Durchschnitt\par
\noindent\emph{Rusky / Russian:} [Bukvy ne sut zachovane / Lettering not recoverable]\par
\noindent\emph{Ukrajinsky / Ukrainian:} [Bukvy ne sut zachovane / Lettering not recoverable]\par
\smallskip\noindent\emph{Razlog:} Zapis nazva največšego občego dělitelja ili prěsěka ima izgubljene bukvy. Zachovany razlog jest matematičny: v oddělě 6 raboty 9 oblast $[\Theta]$ sama izpolnjaje potrěbne uslovja i naleži k familiji $\mathfrak N(\mathfrak G)$. Zato prěsěk ostavaje v toj familiji, za razliku od slučaja pri TERM-0720. Dva nazvy vyražajut dva pogleda na isti objekt; ne sut dvě nezavisne konstrukcije. Aktualny tekst podkrěpjaje to razlikovanje, ale ne vrača izgubljenu grafiju starogo zapisa.\par
\noindent\emph{Reason:} The record of greatest common divisor or intersection has missing letters. Its surviving reason is mathematical: in Paper 9 section 6, the domain $[\Theta]$ itself satisfies the required conditions and belongs to $\mathfrak N(\mathfrak G)$. The intersection therefore remains within this family, unlike the case in TERM-0720. The two names describe the same object in different ways, not two independent constructions. The current passage supports that distinction but does not recover the damaged historical spelling.\par
\noindent{\footnotesize Izvorne zapisy / Original records: HIST-g064-008.}\par
\smallskip\noindent\emph{Sravnjenje s tekstom:} V rabotě 9, oddělě 6, $[\Theta]$ samo jest oblast $\mathfrak G$. Familija $\mathfrak N(\mathfrak G)$ sohranjaje zato prěsěky. To jest točna razlika proti oddělu 5 pri TERM-0720; ne možno ju zaměniti obščim tvrdženjem, že vse razsmatrjane familije sut zamknjene po prěsěku.\par
\noindent\emph{Textual comparison:} In Paper 9 section 6, $[\Theta]$ is itself a domain $\mathfrak G$. The family $\mathfrak N(\mathfrak G)$ therefore admits intersections internally. This is the precise contrast with section 5 in TERM-0720, not a general assertion that every family under consideration is closed under intersection.\par
\noindent{\footnotesize TEXT-COMP-55; pasaže / passages: 7766, 7767, 7773, 7774, 7775, 7776, 7777, 7778, 7779, 7780; \texttt{CORPUS\_FORM\_COMPARISONS.json}.}\par
\end{minipage}\par
\par\medskip\Needspace{7\baselineskip}
\noindent\begin{minipage}{\linewidth}
\hypertarget{TERM-0722}{}
\noindent\textbf{stepenj prostogo ideala}\quad{\footnotesize\texttt{TERM-0722}}\par
\noindent\emph{Iztočnik / Source:} Grad eines Primideals\par
\noindent\emph{Rusky / Russian:} степень простого идеала\par
\noindent\emph{Ukrajinsky / Ukrainian:} степінь простого ідеалу\par
\smallskip\noindent\emph{Razlog:} Stepenj prostogo ideala sohranjaje izrično avtorsko imenovanje veličiny $p^g$ v slučaju nedoskonalogo koeficientnogo polja. Stepenj tu ne znači sam eksponent $g$ ni stepenj proizvoljnogo polinoma. Rusko i ukrajinsko sravnjenje podkrěpjajut prěmy nazv stepenj. Razlog odnosi se k konkretnomu zaključku raboty 24; ne ustanavjaje jednu definiciju stepena dlja vsih prostyh idealov v vsih kontekstah.\par
\noindent\emph{Reason:} Stepenj prostogo ideala preserves the author's explicit name for $p^g$ in the imperfect-coefficient-field case. Degree here is not the exponent $g$ itself or the degree of an arbitrary polynomial. The Russian and Ukrainian comparanda support the direct degree term. The rationale concerns this particular conclusion of Paper 24, not a single definition of degree for every prime ideal in every context.\par
\noindent{\footnotesize Izvorne zapisy / Original records: HIST-g100-010.}\par
\smallskip\noindent\emph{Sravnjenje s tekstom:} Zaključny pasaž raboty 24 imenuje $p^g$ stepenjem prostogo ideala. Věta sravnjaje slučaj s algebraičnymi čislovymi poljami i oddělja jego od priměrov primarnyh idealov razvětvenja. Nazv zato ne znači eksponent $g$ i ne trěba zaměnjati prosty s primarny.\par
\noindent\emph{Textual comparison:} Paper 24's concluding passage calls $p^g$ the degree of a prime ideal. It compares the situation with algebraic number fields and distinguishes it from examples of primary ramification ideals. The name therefore does not mean the exponent $g$, and prime must not be replaced by primary.\par
\noindent{\footnotesize TEXT-COMP-56; pasaže / passages: 15714; \texttt{CORPUS\_FORM\_COMPARISONS.json}.}\par
\end{minipage}\par
\par\medskip\Needspace{7\baselineskip}
\noindent\begin{minipage}{\linewidth}
\hypertarget{TERM-0723}{}
\noindent\textbf{stepenj po prěměnnyh}\quad{\footnotesize\texttt{TERM-0723}}\par
\noindent\emph{Iztočnik / Source:} Grad in den Variablen\par
\noindent\emph{Rusky / Russian:} степень по переменным\par
\noindent\emph{Ukrajinsky / Ukrainian:} степінь за змінними\par
\smallskip\noindent\emph{Razlog:} Stepenj po prěměnnyh kaže, v ktoroj česti polinoma izměrjaje se stepenj. Razlog dava prednost toj izričnoj frazě, da by komentary k tablicam sohranjali ranějše razlikovanje reda i stepena. Doplnjenje po prěměnnyh ne jest nepotrěbno povtorjenje: ono oddělja zavisimost od prěměnnyh od zavisimosti od koeficientov. Zapis ne dava osobnogo čislovogo sravnjenja konkurentnyh nazvov.\par
\noindent\emph{Reason:} Stepenj po prěměnnyh states which part of the polynomial is being measured. The rationale favours an explicit phrase so that the table commentary preserves the earlier order/degree distinction. The qualification in the variables is not redundant: it separates dependence on variables from dependence on coefficients. The record supplies no separate numerical comparison of competing names.\par
\noindent{\footnotesize Izvorne zapisy / Original records: HIST-g002-010.}\par
\end{minipage}\par
\par\medskip\Needspace{7\baselineskip}
\noindent\begin{minipage}{\linewidth}
\hypertarget{TERM-0724}{}
\noindent\textbf{stepeni v rędah prěmennyh}\quad{\footnotesize\texttt{TERM-0724}}\par
\noindent\emph{Iztočnik / Source:} Gradzahlen\par
\noindent\emph{Rusky / Russian:} степени в строках переменных\par
\noindent\emph{Ukrajinsky / Ukrainian:} степені в рядках змінних\par
\smallskip\noindent\emph{Razlog:} Stepeni v rędah prěmennyh izjasnjaje funkciju německogo Gradzahlen, vměsto tvorenja osobnogo složenogo imena dlja čisel stepena. V formule (70) raboty 4 $m_\rho$ i $l_\rho$ sut stepeni v konkretnyh redah prěměnnyh prěd i poslě kontrakcije. Zato indeks reda jest sučstveny. Razlog ne odbija nijedno točno navedeno medžuslovjansko alternativno slovo; odbija težku doslovnu složeninu kako sposob vyraženja.\par
\noindent\emph{Reason:} Stepeni v rędah prěmennyh explains the function of German Gradzahlen rather than creating a separate compound noun for degree-numbers. In Paper 4 equation (70), $m_\rho$ and $l_\rho$ are degrees in particular rows of variables before and after contraction. The row index therefore matters. The rationale rejects a cumbersome literal compound as a mode of expression, but records no exact competing Interslavic word.\par
\noindent{\footnotesize Izvorne zapisy / Original records: HIST-g037-009.}\par
\smallskip\noindent\emph{Sravnjenje s tekstom:} V formuli (70) $m_\rho-l_\rho=-\alpha_{\rho-1}+2\alpha_\rho-\alpha_{\rho+1}$ stepeni $m_\rho,l_\rho$ odnoset se k redam $p_\rho$. Početkova i končna forma razlikujut se po kontrakcijah. To podkrěpjaje opisnu frazu stepeni v redah, ne samostalny novy pojem čisla stepena.\par
\noindent\emph{Textual comparison:} In equation (70), $m_\rho-l_\rho=-\alpha_{\rho-1}+2\alpha_\rho-\alpha_{\rho+1}$, the degrees $m_\rho,l_\rho$ refer to the rows $p_\rho$. Contractions relate the initial and resulting forms. This supports degrees in rows as an explanatory phrase, not a separate new degree-number concept.\par
\noindent{\footnotesize TEXT-COMP-57; pasaže / passages: 4864, 4866, 4867, 4870; \texttt{CORPUS\_FORM\_COMPARISONS.json}.}\par
\end{minipage}\par
\par\medskip\Needspace{7\baselineskip}
\noindent\begin{minipage}{\linewidth}
\hypertarget{TERM-0725}{}
\noindent\textbf{stepeni / ? kako najve?a med stepenami}\quad{\footnotesize\texttt{TERM-0725}}\par
\noindent\emph{Iztočnik / Source:} Gradzahlen / Dimension ?\par
\noindent\emph{Rusky / Russian:} [Bukvy ne sut zachovane / Lettering not recoverable]\par
\noindent\emph{Ukrajinsky / Ukrainian:} [Bukvy ne sut zachovane / Lettering not recoverable]\par
\smallskip\noindent\emph{Razlog:} Zapis stepeni / ? kako največa med stepenami ima povrědženy symbol i jezyčne formy. Zachovany razlog zadržuje največši stepenj med koeficientnymi polinomami. Aktualny tekst raboty 9 piše $\lambda$ i izrično svezuje to čislo k polinomam $a_1(\eta),\ldots,a_\chi(\eta)$. To ne jest dimenzija vektornogo prostora. Novějše čitanje može objasniti argument, ale ne smě byti predstavljeno kako nepovrědženy originalny zapis.\par
\noindent\emph{Reason:} The record stepeni / ? kako največa med stepenami has a damaged symbol and damaged language fields. Its surviving rationale retains the greatest degree among the coefficient polynomials. The current Paper 9 passage writes $\lambda$ and explicitly relates it to $a_1(\eta),\ldots,a_\chi(\eta)$. This is not vector-space dimension. The fresh reading can explain the argument but must not be presented as an undamaged historical record.\par
\noindent{\footnotesize Izvorne zapisy / Original records: HIST-g063-008.}\par
\smallskip\noindent\emph{Sravnjenje s tekstom:} Aktualna formula raboty 9 razlikuje stepenj uravnenja $\chi$ od $\lambda$, največšego stepena med koeficientnymi polinomami $a_i(\eta)$. Zapisane pytajniky starogo glosara ne dajut prava sjediniti te veličiny ili pripisati mu nepovrědženu grčku bukvu.\par
\noindent\emph{Textual comparison:} The current Paper 9 formula distinguishes the equation degree $\chi$ from $\lambda$, the greatest degree among coefficient polynomials $a_i(\eta)$. The old glossary's question marks do not license merging these quantities or attributing an undamaged Greek letter to its record.\par
\noindent{\footnotesize TEXT-COMP-58; pasaže / passages: 7630, 7633, 7634; \texttt{CORPUS\_FORM\_COMPARISONS.json}.}\par
\end{minipage}\par
\par\medskip\Needspace{7\baselineskip}
\noindent\begin{minipage}{\linewidth}
\hypertarget{TERM-0726}{}
\noindent\textbf{granične oblasti}\quad{\footnotesize\texttt{TERM-0726}}\par
\noindent\emph{Iztočnik / Source:} Grenzbereiche\par
\noindent\emph{Rusky / Russian:} граничные области\par
\noindent\emph{Ukrajinsky / Ukrainian:} межові області\par
\smallskip\noindent\emph{Razlog:} Granične oblasti sohranjaje oba konce obuhvačenogo razpona oblastij. Razlog daje prednost doslovnomu opisu, zato že zaključna věta raboty 9 prěnosi ranějše rezultaty na prostu karakteristiku. Aktualny pasaž izrično vključaje oblast izvedenu iz bazy i samo polje. Granične zato ne znači topologičnu granicu ni približenje limitu; važna jest vključennost oboju krajnyh oblastij.\par
\noindent\emph{Reason:} Granične oblasti preserves the two endpoints of the range of domains under consideration. The rationale prefers a literal description because Paper 9's final statement transfers earlier results to prime characteristic. The current passage explicitly includes both the domain generated from the basis and the field itself. Boundary here is not a topological boundary or a limiting approximation; the inclusion of both endpoint domains is the point.\par
\noindent{\footnotesize Izvorne zapisy / Original records: HIST-g068-013.}\par
\smallskip\noindent\emph{Sravnjenje s tekstom:} Zaključek za prostu karakteristiku dopušča oblast iz bazy, samo polje i medžuležeče oblasti. Pasaž izrično vključaje oba konca. Zato granične oblasti govori o koncah toga algebraičnogo razpona, ne o topologičnom kraju oblasti.\par
\noindent\emph{Textual comparison:} The prime-characteristic conclusion allows the basis-generated domain, the field itself and intermediate domains. It explicitly includes both endpoints. Boundary domains therefore concerns this algebraic range, not a topological boundary.\par
\noindent{\footnotesize TEXT-COMP-59; pasaže / passages: 8388, 8389, 8390, 8391, 8392, 8393, 8394; \texttt{CORPUS\_FORM\_COMPARISONS.json}.}\par
\end{minipage}\par
\par\medskip\Needspace{7\baselineskip}
\noindent\begin{minipage}{\linewidth}
\hypertarget{TERM-0727}{}
\noindent\textbf{največši obči dělitelj / najmene obče mnogokratno}\quad{\footnotesize\texttt{TERM-0727}}\par
\noindent\emph{Iztočnik / Source:} groesster gemeinsamer Teiler / kleinstes gemeinsames Vielfaches\par
\noindent\emph{Rusky / Russian:} наибольший общий делитель / наименьшее общее кратное\par
\noindent\emph{Ukrajinsky / Ukrainian:} найбільший спільний дільник / найменше спільне кратне\par
\smallskip\noindent\emph{Razlog:} Največši obči dělitelj / najmene obče mnogokratno prěnosi znanu aritmetičnu paru na jednostranne moduly. To ne znači, že moduly stavajut čislami: jih definicije oprěděljajut smysl. V rabotě 17 obči dělitelj tvoret sumy elementov dvoh modulov, a obče mnogokratno tvoret elementy, naležeče obom modulam. Razlog sohranjaje avtorsku analogiju, no ne dava prava zaměniti děliteljny poredok obyčnym sravnjenjem čislovoj velikosti.\par
\noindent\emph{Reason:} Največši obči dělitelj / najmene obče mnogokratno extends a familiar arithmetic pair to one-sided modules. This does not turn modules into numbers: the definitions determine the meaning. In Paper 17 the common divisor consists of sums of elements of the two modules, while the common multiple consists of elements belonging to both. The rationale preserves the author's analogy, not an ordering by ordinary numerical magnitude.\par
\noindent{\footnotesize Izvorne zapisy / Original records: HIST-g079-020.}\par
\smallskip\noindent\emph{Sravnjenje s tekstom:} Rabota 17, odděl 4, definira děliteljstvo čerez implikaciju kongruencij. Največši obči dělitelj sostoji se iz sum $M+N$, a najmenše obče mnogokratno iz polinomov, naležečih obom modulam. Vzajemna prostost znači, že obči dělitelj jest cěly jediničny modul. Tute definicije, ne čislova metafora sama po sobě, opravdyvajut prěvod.\par
\noindent\emph{Textual comparison:} Paper 17 section 4 defines divisibility through an implication between congruences. The greatest common divisor consists of sums $M+N$, and the least common multiple of polynomials belonging to both modules. Coprimeness means that the common divisor is the whole unit module. These definitions, not the numerical metaphor alone, justify the rendering.\par
\noindent{\footnotesize TEXT-COMP-60; pasaže / passages: 11157, 11159, 11161, 11163, 11165; \texttt{CORPUS\_FORM\_COMPARISONS.json}.}\par
\end{minipage}\par
\par\medskip\Needspace{7\baselineskip}
\noindent\begin{minipage}{\linewidth}
\hypertarget{TERM-0728}{}
\noindent\textbf{osnovne oblasti}\quad{\footnotesize\texttt{TERM-0728}}\par
\noindent\emph{Iztočnik / Source:} Grundbereiche\par
\noindent\emph{Rusky / Russian:} основные области\par
\noindent\emph{Ukrajinsky / Ukrainian:} основні області\par
\smallskip\noindent\emph{Razlog:} Osnovne oblasti imenuje izhodne algebraične oblasti, na ktore prěnoset se aksiomy. Razlog sohranjaje oblast vměsto avtomatičnogo prěimenovanja vseh objektov prstenami. V rabotě 30 trěba razlikovati cěle čisla, polinomy jednoj prěměnnoj i funkcionalnu oblast polinomov několiko prěměnnyh. Obyčny prsten polinomov několiko prěměnnyh ne izpolnjaje vsih tyh uslovij; imenno ta razlika objasnjaje potrěbu funkcionalnoj oblasti.\par
\noindent\emph{Reason:} Osnovne oblasti names the starting algebraic domains to which the axioms are transferred. The rationale retains domain rather than automatically renaming every object a ring. Paper 30 distinguishes integers, polynomials in one variable and the functional domain of polynomials in several variables. The ordinary multivariable polynomial ring does not satisfy all those conditions; this difference explains why the functional domain is needed.\par
\noindent{\footnotesize Izvorne zapisy / Original records: HIST-g110-008.}\par
\smallskip\noindent\emph{Sravnjenje s tekstom:} Rabota 30 razlikuje polinomy jednoj prěměnnoj od običnogo prstena polinomov několiko prěměnnyh: v drugom slučaju uslovje cěpov mnogokratnyh II ne imaje město. Prěhod k funkcionalnoj oblasti jest zatim sučstveny. Osnovne oblasti ne znači, že vsaky imenovany prsten bez dodatnogo strojenja izpolnjaje vse aksiomy.\par
\noindent\emph{Textual comparison:} Paper 30 distinguishes univariate polynomials from the ordinary multivariable polynomial ring: in the latter, the multiple-chain condition II fails. Passage to the functional domain is therefore substantive. Basic domains does not say that every named ring satisfies all the axioms without further construction.\par
\noindent{\footnotesize TEXT-COMP-61; pasaže / passages: 16401, 16403, 16405, 16407; \texttt{CORPUS\_FORM\_COMPARISONS.json}.}\par
\end{minipage}\par
\par\medskip\Needspace{7\baselineskip}
\noindent\begin{minipage}{\linewidth}
\hypertarget{TERM-0729}{}
\noindent\textbf{osnovno svijanje tipa $\rho$}\quad{\footnotesize\texttt{TERM-0729}}\par
\noindent\emph{Iztočnik / Source:} Grundfaltung vom Typus $\rho$\par
\noindent\emph{Rusky / Russian:} основная свертка типа $\rho$\par
\noindent\emph{Ukrajinsky / Ukrainian:} основне згортання типу $\rho$\par
\smallskip\noindent\emph{Razlog:} Osnovno svijanje tipa $\rho$ označaje osnovnu operaciju, iz ktoroj Teorema VIII sostavjaje obči proces. Razlog sohranjaje pridavnik osnovno i ranějši nazv svijanje. Aktualny tekst vměsto togo upotrěbjaje osnovna kontrakcija: to jest kogredientna kontrakcija defekta 1 dvoh redov istogo tipa $\rho$. Pridavnik osnovny zato ima konkretno strukturno opravdanje, ne samo značenje prosty ili legky.\par
\noindent\emph{Reason:} Osnovno svijanje tipa $\rho$ denotes a basic operation from which Theorem VIII constructs the general process. The rationale preserves the adjective osnovno and the earlier term svijanje. The current text instead uses osnovna kontrakcija: a cogredient contraction of defect 1 between two rows of the same type $\rho$. Basic therefore has a specific structural justification, not merely the meaning simple or easy.\par
\noindent{\footnotesize Izvorne zapisy / Original records: HIST-g037-004.}\par
\smallskip\noindent\emph{Sravnjenje s tekstom:} Aktualny odděl 8 raboty 4 definira osnovnu kontrakciju tipa $\rho$ kako kogredientnu kontrakciju defekta 1 dvoh redov istogo tipa. Teorema VIII tvori obči proces iz $n-1$ takyh operacij. Ta definicija objasnjaje osnovny; aktualno kontrakcija ne jest tožna bukvena forma historijnogo svijanje.\par
\noindent\emph{Textual comparison:} The current Paper 4 section 8 defines a basic contraction of type $\rho$ as a cogredient contraction of defect 1 between two rows of that type. Theorem VIII constructs the general process from $n-1$ such operations. This definition explains basic; current kontrakcija is not the same spelling as historical svijanje.\par
\noindent{\footnotesize TEXT-COMP-62; pasaže / passages: 4649, 4651, 4653; \texttt{CORPUS\_FORM\_COMPARISONS.json}.}\par
\end{minipage}\par
\par\medskip\Needspace{7\baselineskip}
\noindent\begin{minipage}{\linewidth}
\hypertarget{TERM-0730}{}
\noindent\textbf{osnovne svijanja}\quad{\footnotesize\texttt{TERM-0730}}\par
\noindent\emph{Iztočnik / Source:} Grundfaltungen\par
\noindent\emph{Rusky / Russian:} основные свертки; основные свёртки\par
\noindent\emph{Ukrajinsky / Ukrainian:} основні згортання\par
\smallskip\noindent\emph{Razlog:} Osnovne svijanja svezuje dva argumenty: formalny status v Teoremi I raboty 2 i prěhod k redukciji v oddělě 8 raboty 4. Oba razloga podkrěpjajut osnovny, ale samo povtorjenje nazva ne dokazuje, že vse definicije v oběh rabotah sut identične. Pri TERM-0729 jest dana točna pozdnějša definicija i aktualna forma osnovna kontrakcija. Historijna množinska forma ostavaje zapisom izbora, ne osobnym dokazom jej gramatičnoj pravilnosti.\par
\noindent\emph{Reason:} Osnovne svijanja connects two arguments: formal status in Paper 2's Theorem I and the transition to the reduction in Paper 4 section 8. Both rationales support basic, but recurrence of the name alone does not prove that all definitions in the two papers coincide. TERM-0729 gives the precise later definition and current wording osnovna kontrakcija. The historical plural remains evidence of the choice, not a separate demonstration of grammatical correctness.\par
\noindent{\footnotesize Izvorne zapisy / Original records: HIST-g003-007, HIST-g036-009.}\par
\smallskip\noindent\emph{Sravnjenje s tekstom:} Aktualny odděl 8 raboty 4 definira osnovnu kontrakciju tipa $\rho$ kako kogredientnu kontrakciju defekta 1 dvoh redov istogo tipa. Teorema VIII tvori obči proces iz $n-1$ takyh operacij. Ta definicija objasnjaje osnovny; aktualno kontrakcija ne jest tožna bukvena forma historijnogo svijanje.\par
\noindent\emph{Textual comparison:} The current Paper 4 section 8 defines a basic contraction of type $\rho$ as a cogredient contraction of defect 1 between two rows of that type. Theorem VIII constructs the general process from $n-1$ such operations. This definition explains basic; current kontrakcija is not the same spelling as historical svijanje.\par
\noindent{\footnotesize TEXT-COMP-62; pasaže / passages: 4649, 4651, 4653; \texttt{CORPUS\_FORM\_COMPARISONS.json}.}\par
\end{minipage}\par
\par\medskip\Needspace{7\baselineskip}
\noindent\begin{minipage}{\linewidth}
\hypertarget{TERM-0731}{}
\noindent\textbf{osnovna forma}\quad{\footnotesize\texttt{TERM-0731}}\par
\noindent\emph{Iztočnik / Source:} Grundform\par
\noindent\emph{Rusky / Russian:} основная форма\par
\noindent\emph{Ukrajinsky / Ukrainian:} основна форма\par
\smallskip\noindent\emph{Razlog:} Osnovna forma sohranjaje objekt, od ktorogo počinaje invariantno-teoretična konstrukcija. Tri razloga naglašajut prěmy matematičny izraz, kontinuitet med rabotami i izběganje užšego čitanja kako element bazy. V jednom zapisu rolju izhodnoj formy ima $\nu$; to ne čini slovo osnovna synonymom linearno nezavisna. Prednost jest pojmovna kontinuitetnost, ne čislena ocěnka razumělivosti.\par
\noindent\emph{Reason:} Osnovna forma preserves the object from which an invariant-theoretic construction starts. Three rationales emphasize direct mathematical wording, continuity across papers and avoidance of the narrower reading basis element. One record assigns the initial-form role to $\nu$; that does not make osnovna synonymous with linearly independent. The preference concerns conceptual continuity, not a measured intelligibility score.\par
\noindent{\footnotesize Izvorne zapisy / Original records: HIST-g006-010, HIST-g043-013, HIST-g057-002.}\par
\end{minipage}\par
\par\medskip\Needspace{7\baselineskip}
\noindent\begin{minipage}{\linewidth}
\hypertarget{TERM-0732}{}
\noindent\textbf{osnovna forma / sistema osnovnyh form}\quad{\footnotesize\texttt{TERM-0732}}\par
\noindent\emph{Iztočnik / Source:} Grundform / Grundformensystem\par
\noindent\emph{Rusky / Russian:} основная форма / система основных форм\par
\noindent\emph{Ukrajinsky / Ukrainian:} основна форма / система основних форм\par
\smallskip\noindent\emph{Razlog:} Osnovna forma / sistema osnovnyh form razděljuje pojedinu izhodnu formu i sistem takyh form. Razlog raboty 15 odkazuje k binarnoj formě, čije invarianty izslědujut se, ne k osnovnomu polju. Zato pridavnik osnovna vyražaje funkciju formy v konstrukciji. Sistem takože ne treba bez dokazanja zaměniti slovom baza: zapis ne tvrdi linearnoj nezavisnosti jego členov.\par
\noindent\emph{Reason:} Osnovna forma / sistema osnovnyh form distinguishes an initial form from a system of such forms. The Paper 15 rationale concerns the binary form whose invariants are studied, not a base field. The adjective describes the form's role in the construction. Nor should system silently become basis: the record does not assert linear independence of its members.\par
\noindent{\footnotesize Izvorne zapisy / Original records: HIST-g077-003.}\par
\end{minipage}\par
\par\medskip\Needspace{7\baselineskip}
\noindent\begin{minipage}{\linewidth}
\hypertarget{TERM-0733}{}
\noindent\textbf{osnovne formy}\quad{\footnotesize\texttt{TERM-0733}}\par
\noindent\emph{Iztočnik / Source:} Grundformen\par
\noindent\emph{Rusky / Russian:} основные формы\par
\noindent\emph{Ukrajinsky / Ukrainian:} основні форми\par
\smallskip\noindent\emph{Razlog:} Osnovne formy prodolžuje termin iz invariantnoj teorije. Jeden razlog govori o tom, že formy iz (31) razsmatrjajut se kako izhodne prěd substitucijeju (11); drugy sohranjaje tu terminologiju v rabotě 6. To jest razlog kontekstnogo kontinuiteta, ne dokaz osobnoj prednosti množinske formy prěd drugimi izrazami. Važna jest rolja form v substituciji, ne njih obična ili elementarna složnost.\par
\noindent\emph{Reason:} Osnovne formy continues the invariant-theoretic term. One rationale treats the forms in (31) as initial forms before substitution (11); the other carries the vocabulary into Paper 6. This is a reason for contextual continuity, not a separate demonstration that the plural outperforms competing expressions. The relevant feature is the forms' role in substitution, not their being ordinary or elementary.\par
\noindent{\footnotesize Izvorne zapisy / Original records: HIST-g034-005, HIST-g041-015.}\par
\end{minipage}\par
\par\medskip\Needspace{7\baselineskip}
\noindent\begin{minipage}{\linewidth}
\hypertarget{TERM-0734}{}
\noindent\textbf{osnovne funkcije}\quad{\footnotesize\texttt{TERM-0734}}\par
\noindent\emph{Iztočnik / Source:} Grundfunktionen\par
\noindent\emph{Rusky / Russian:} основные функции\par
\noindent\emph{Ukrajinsky / Ukrainian:} основні функції\par
\smallskip\noindent\emph{Razlog:} Osnovne funkcije imenuje funkcije, od ktoryh zavisi konstrukcija polnogo sistema v rabotě 12. Razlog sravnjaje anglijsku glosu basic/main functions (osnovne/glavne funkcije) s fundamental functions (fundamentalne funkcije). Vtora mogla by sugerovati drugy, samostalno definovany sovrěmenny objekt. Izbor ne umenšaje matematičnu važnost funkcij; on ograničaje tvrdženje do jih lokalnoj rolje. Zapis ne dava dodatnoj nezavisnoj definicije novogo klasa funkcij.\par
\noindent\emph{Reason:} Osnovne funkcije names the functions on which Paper 12's complete-system construction depends. The recorded comparison is between the English glosses basic/main functions and fundamental functions; the latter could suggest a different, independently defined modern object. The choice does not diminish their mathematical importance; it limits the claim to their local role. The record supplies no additional independent definition of a new class of functions.\par
\noindent{\footnotesize Izvorne zapisy / Original records: HIST-g074-011.}\par
\end{minipage}\par
\par\medskip\Needspace{7\baselineskip}
\noindent\begin{minipage}{\linewidth}
\hypertarget{TERM-0735}{}
\noindent\textbf{osnovny ideal}\quad{\footnotesize\texttt{TERM-0735}}\par
\noindent\emph{Iztočnik / Source:} Grundideal\par
\noindent\emph{Rusky / Russian:} основной идеал\par
\noindent\emph{Ukrajinsky / Ukrainian:} основний ідеал\par
\smallskip\noindent\emph{Razlog:} Osnovny ideal ima několiko kontekstov, ne jednu avtomatično obšču definiciju. Razlogi zadržujut svez s osnovny modul, kontinuitet stupenjevih idealov $g_i$ i razliku od fundamentalnogo uravnenja v rabotě 25. Ale rabota 38 izrično razděljuje svoj osnovny ideal, definovany prěz reduciranu normu reduciranoj diferenty, od Dedekindovogo imenovanja diferenty i diskriminanta. Edinstvo pridavnika zato ne smě zakryti različnost pojmov. Prednost osnovny prěd fundamentalny jest dokumentovana; edinstvo vseh smyslov ne jest.\par
\noindent\emph{Reason:} Osnovny ideal has several contexts, not one automatically shared definition. The rationales preserve its connection with osnovny modul, the sequence of basic ideals $g_i$, and the distinction from fundamental equation in Paper 25. Paper 38, however, explicitly distinguishes its basic ideal, defined through the reduced norm of the reduced different, from Dedekind's terminology for the different and discriminant. A common adjective must not conceal different concepts. The preference for osnovny over fundamentalny is documented; identity of every sense is not.\par
\noindent{\footnotesize Izvorne zapisy / Original records: HIST-g087-007, HIST-g092-027, HIST-g095-012, HIST-g099-002, HIST-g102-012, HIST-g208-013.}\par
\smallskip\noindent\emph{Sravnjenje s tekstom:} Primětka raboty 38 izrično zabranuje směšati njen osnovny ideal s Dedekindovym imenovanjem diferenty i diskriminanta. Lokalna definicija ide čerez reduciranu normu reduciranoj diferenty. Zadržanje istogo prěvodnogo nazva zato potrěbuje razděljenje smyslov i točne citaty, ne prosto zjedinjenje vseh pojavjenj.\par
\noindent\emph{Textual comparison:} Paper 38's footnote expressly warns against confusing its basic ideal with Dedekind's terminology for the different and discriminant. Its local definition uses the reduced norm of the reduced different. Retaining the same translated name therefore requires separated senses and exact references, not simply merging all occurrences.\par
\noindent{\footnotesize TEXT-COMP-63; pasaže / passages: 23088; \texttt{CORPUS\_FORM\_COMPARISONS.json}.}\par
\end{minipage}\par
\par\medskip\Needspace{7\baselineskip}
\noindent\begin{minipage}{\linewidth}
\hypertarget{TERM-0736}{}
\noindent\textbf{osnovny ideal $i$-go nivoja}\quad{\footnotesize\texttt{TERM-0736}}\par
\noindent\emph{Iztočnik / Source:} Grundideal $i$-ter Stufe\par
\noindent\emph{Rusky / Russian:} основной идеал $i$-й ступени\par
\noindent\emph{Ukrajinsky / Ukrainian:} основний ідеал $i$-го рівня\par
\smallskip\noindent\emph{Razlog:} Osnovny ideal $i$-go nivoja sohranjaje obči korenj s osnovny modul i osnovny ideal. Razlog izrično dava prednost osnovny prěd fundamentalny dlja togo terminologičnogo sveza. Doplnjenje $i$-go nivoja identifikuje město v stupenjevom strojenju; ono ne označaje stepenj ideala ni jeho mnogokratnost. Razlog ne sravnjaje nivoj s točno navedenym medžuslovjanskym synonymom, zato takovo sravnjenje ne treba domysljati.\par
\noindent\emph{Reason:} Osnovny ideal $i$-go nivoja preserves the shared stem of osnovny modul and osnovny ideal. The rationale explicitly favours osnovny over fundamentalny for that terminological connection. The qualification of stage $i$ locates the ideal in a staged construction; it is not an ideal power or its multiplicity. No exact Interslavic alternative to nivoj is compared in the record, so that comparison must not be invented.\par
\noindent{\footnotesize Izvorne zapisy / Original records: HIST-g096-014.}\par
\end{minipage}\par
\par\medskip\Needspace{7\baselineskip}
\noindent\begin{minipage}{\linewidth}
\hypertarget{TERM-0737}{}
\noindent\textbf{osnovno polje}\quad{\footnotesize\texttt{TERM-0737}}\par
\noindent\emph{Iztočnik / Source:} Grundkörper\par
\noindent\emph{Rusky / Russian:} основное поле\par
\noindent\emph{Ukrajinsky / Ukrainian:} основне поле\par
\smallskip\noindent\emph{Razlog:} Osnovno polje sohranjaje pridavnik osnovny, upotrěbjany takože pri idealah i modulah, ale imenik polje nosi samostalno matematično uslovje. V ciklotomičnom argumentu raboty 32 izhodny objekt jest komutativno koeficientno polje. Zato polje ne treba zaměniti širim tělo samo zaradi doslovnosti německogo Körper. Razlog o prozračnosti jest jezyčno sudženje; on ne daje izměrenogo rezultata čitatelskogo testa.\par
\noindent\emph{Reason:} Osnovno polje retains the adjective used with basic ideals and modules, but polje carries an independent mathematical condition. In Paper 32's cyclotomic argument the starting object is a commutative coefficient field. It should not become the broader tělo merely to follow German Körper word for word. The claim of transparency is a linguistic judgment, not a measured reader-test result.\par
\noindent{\footnotesize Izvorne zapisy / Original records: HIST-g097-005, HIST-g175-005, HIST-g203-006.}\par
\end{minipage}\par
\par\medskip\Needspace{7\baselineskip}
\noindent\begin{minipage}{\linewidth}
\hypertarget{TERM-0738}{}
\noindent\textbf{osnovno polje / nadpolje}\quad{\footnotesize\texttt{TERM-0738}}\par
\noindent\emph{Iztočnik / Source:} Grundkörper / Oberkörper\par
\noindent\emph{Rusky / Russian:} основное поле / надполе\par
\noindent\emph{Ukrajinsky / Ukrainian:} базове поле / надполе\par
\smallskip\noindent\emph{Razlog:} Osnovno polje / nadpolje vyražaje odnos polja i polja, ktore jego sodrživaje. Razlog odbija ukrajinsko «верхнє поле» i rusko «верхнее поле», t. j. doslovny prostranstveny opis vyššego polja: potrěbna jest algebraična relacija razširjenja, ne vysota. Nadpolje sohranjaje tu relaciju kratko. Sam nazv ne tvrdi, že razširjenje jest algebraično, konečno ili normalno; take uslovja morajut prihoditi iz konkretnoj pasaže.\par
\noindent\emph{Reason:} Osnovno polje / nadpolje expresses the relation between a field and a field containing it. The rationale rejects the literal Ukrainian верхнє поле and Russian верхнее поле, upper field, because the intended relation is algebraic extension, not spatial height. Nadpolje states that relation compactly. The name alone does not assert that the extension is algebraic, finite or normal; those conditions must come from the passage.\par
\noindent{\footnotesize Izvorne zapisy / Original records: HIST-g076-004.}\par
\end{minipage}\par
\par\medskip\Needspace{7\baselineskip}
\noindent\begin{minipage}{\linewidth}
\hypertarget{TERM-0739}{}
\noindent\textbf{osnovny modul}\quad{\footnotesize\texttt{TERM-0739}}\par
\noindent\emph{Iztočnik / Source:} Grundmodul\par
\noindent\emph{Rusky / Russian:} основной модуль\par
\noindent\emph{Ukrajinsky / Ukrainian:} основний модуль\par
\smallskip\noindent\emph{Razlog:} Osnovny modul zadržuje osnovnu rolju modula i paralelu s osnovny ideal. Četiri razloga sohranjajut prěmy izraz vměsto specialnogo neologizma; dva izrično ostavljajut fundamentalny modul kako alternativu dlja daljšego razsmotrenja. To jest dokumentovana otvorjena alternativa, ne svědočstvo, že ekspert ju už odbil. Prednost osnovny opiraje se na svez med pojmami, ne na novu tvrdnju o prirodnosti vsakoj upotrěby.\par
\noindent\emph{Reason:} Osnovny modul retains the module's basic role and the parallel with osnovny ideal. Four rationales preserve a direct expression instead of a special neologism; two explicitly leave fundamentalny modul for further consideration. This is a documented open alternative, not evidence that an expert rejected it. The preference for osnovny rests on the connection between concepts, not a new claim that every use is idiomatic.\par
\noindent{\footnotesize Izvorne zapisy / Original records: HIST-g080-008, HIST-g085-006, HIST-g092-006, HIST-g093-007, HIST-g095-013.}\par
\end{minipage}\par
\par\medskip\Needspace{7\baselineskip}
\noindent\begin{minipage}{\linewidth}
\hypertarget{TERM-0740}{}
\noindent\textbf{osnovny modul / osnovny ideal}\quad{\footnotesize\texttt{TERM-0740}}\par
\noindent\emph{Iztočnik / Source:} Grundmodul / Grundideal\par
\noindent\emph{Rusky / Russian:} основной модуль / основной идеал\par
\noindent\emph{Ukrajinsky / Ukrainian:} основний модуль / основний ідеал\par
\smallskip\noindent\emph{Razlog:} Osnovny modul / osnovny ideal zadržuje istu imennuju familiju pri prěnosu tvrđenja o polinomnyh idealah na moduly linearnyh form v rabotě 22. Razlog odbija prěhod k fundamentalny, zato že bi on oslabel vidimy svez konstrukcij. To ne tvrdi, že ideal i modul sut vsude isty objekt; obščij pridavnik pomaga slěditi konkretno prěformulovanje, dokolě imeniki sohranjajut svoje matematične funkcije.\par
\noindent\emph{Reason:} Osnovny modul / osnovny ideal retains the same naming family when Paper 22 transfers statements about polynomial ideals to modules of linear forms. The rationale avoids switching to fundamentalny because that would obscure the connection between the constructions. It does not claim that ideal and module are everywhere the same object: the shared adjective makes this reformulation visible while the nouns retain their mathematical functions.\par
\noindent{\footnotesize Izvorne zapisy / Original records: HIST-g088-006, HIST-g092-037.}\par
\end{minipage}\par
\par\medskip\Needspace{7\baselineskip}
\noindent\begin{minipage}{\linewidth}
\hypertarget{TERM-0741}{}
\noindent\textbf{grupa analitičnyh transformacij}\quad{\footnotesize\texttt{TERM-0741}}\par
\noindent\emph{Iztočnik / Source:} Gruppe analytischer Transformationen\par
\noindent\emph{Rusky / Russian:} группа аналитических преобразований\par
\noindent\emph{Ukrajinsky / Ukrainian:} група аналітичних перетворень\par
\smallskip\noindent\emph{Razlog:} Grupa analitičnyh transformacij sohranjaje tri česti pojma: grupu, analitičnost i transformacije. Razlog prodolžuje ranějšu terminologiju prěobraženij; on ne zaměnjaje grupu proizvoljnym soborom funkcij. Transformacija jest internacionalna forma, sravnena s ruskim preobrazovanie i ukrajinskim peretvorennja. Zapis ne kaže, že analitičnost možno izpustiti, i ne dava čislovoj ocěnky tih leksičnyh rěšenij.\par
\noindent\emph{Reason:} Grupa analitičnyh transformacij preserves three components: group, analyticity and transformations. The rationale continues the earlier transformation vocabulary; it does not replace a group with an arbitrary collection of functions. Transformacija is the international form compared with Russian preobrazovanie and Ukrainian peretvorennja. The record neither licenses omission of analyticity nor gives numerical scores for these lexical choices.\par
\noindent{\footnotesize Izvorne zapisy / Original records: HIST-g074-003.}\par
\end{minipage}\par
\par\medskip\Needspace{7\baselineskip}
\noindent\begin{minipage}{\linewidth}
\hypertarget{TERM-0742}{}
\noindent\textbf{grupa s oblastju operatorov}\quad{\footnotesize\texttt{TERM-0742}}\par
\noindent\emph{Iztočnik / Source:} Gruppe mit Operatorenbereich\par
\noindent\emph{Rusky / Russian:} группа с областью операторов\par
\noindent\emph{Ukrajinsky / Ukrainian:} група з областю операторів\par
\smallskip\noindent\emph{Razlog:} Grupa s oblastju operatorov zadržuje grupu i dodatny sistem operatorov kako različne česti konstrukcije. V definiciji raboty 33 operator posyla element grupy v element toj samoj grupy i izpolnjaje $\Theta(ab)=\Theta(a)\Theta(b)$. Razlog zato ne zaměnjaje oblast operatorov prosto grupnym dějanjem: tekst v toj definiciji ne trěbuje, aby sam sistem operatorov byl grupoju invertibilnyh odobraženij. Doslovna fraza sohranjaje točno tu obščnost.\par
\noindent\emph{Reason:} Grupa s oblastju operatorov retains the group and the additional operator system as distinct parts of the construction. In Paper 33's definition an operator maps a group element into the same group and satisfies $\Theta(ab)=\Theta(a)\Theta(b)$. The rationale therefore does not simply replace operator domain with group action: this definition does not require the operator system itself to be a group of invertible maps. The literal phrase preserves that generality.\par
\noindent{\footnotesize Izvorne zapisy / Original records: HIST-g177-002.}\par
\smallskip\noindent\emph{Sravnjenje s tekstom:} Definicija raboty 33 trěbuje jednoznačno odobraženje v istu grupu i zakon $\Theta(ab)=\Theta(a)\Theta(b)$. Ona v toj věti ne tvrdi invertibilnosti operatorov ni grupnogo zakona na njihovom sistemu. Zato grupa s oblastju operatorov ne možno avtomatično zaměniti sovrěmennym grupnym dějanjem.\par
\noindent\emph{Textual comparison:} Paper 33's definition requires a single-valued map into the same group and $\Theta(ab)=\Theta(a)\Theta(b)$. That sentence asserts neither invertibility of operators nor a group law on their system. Group with an operator domain therefore cannot automatically become a modern group action.\par
\noindent{\footnotesize TEXT-COMP-64; pasaže / passages: 19687; \texttt{CORPUS\_FORM\_COMPARISONS.json}.}\par
\end{minipage}\par
\par\medskip\Needspace{7\baselineskip}
\noindent\begin{minipage}{\linewidth}
\hypertarget{TERM-0743}{}
\noindent\textbf{grupne karaktere}\quad{\footnotesize\texttt{TERM-0743}}\par
\noindent\emph{Iztočnik / Source:} Gruppencharaktere\par
\noindent\emph{Rusky / Russian:} характеры групп\par
\noindent\emph{Ukrajinsky / Ukrainian:} характери груп\par
\smallskip\noindent\emph{Razlog:} Grupne karaktere sohranjaje specialny smysl karaktera v teoriji prědstavjenij. Razlog prodolžuje ranějšu terminologiju karakterov i ostatkovyh karakterov, vměsto obyčnogo opisa osobnosti ili svojstva. Takova smyslova kontinuitetnost ne dokazuje sama po sobě pravilnost istorijnoj množinske formy karaktere. Odborno značenje i gramatična forma sut dvě oddělne pytanja; zapis opravdyvaje prěd vsem prvo.\par
\noindent\emph{Reason:} Grupne karaktere retains the specialized sense of character in representation theory. The rationale continues the earlier character and residue-character vocabulary rather than using an ordinary description of personality or a feature. This semantic continuity does not by itself establish the correctness of the historical plural karaktere. Technical meaning and grammatical form are separate questions; the record chiefly argues for the former.\par
\noindent{\footnotesize Izvorne zapisy / Original records: HIST-g105-001.}\par
\end{minipage}\par
\par\medskip\Needspace{7\baselineskip}
\noindent\begin{minipage}{\linewidth}
\hypertarget{TERM-0744}{}
\noindent\textbf{grupny determinant}\quad{\footnotesize\texttt{TERM-0744}}\par
\noindent\emph{Iztočnik / Source:} Gruppendeterminante\par
\noindent\emph{Rusky / Russian:} групповой детерминант\par
\noindent\emph{Ukrajinsky / Ukrainian:} груповий детермінант\par
\smallskip\noindent\emph{Razlog:} Grupny determinant izbiraje prěmy nazv dlja Frobeniusovogo i Dedekindovogo objekta. Razlog odbija modernu opisnu zaměnu, ale ne imenuje konkretnu konkurentnu medžuslovjansku frazu. Determinant zato ostavaje prepoznavajemy matematičny korenj, a grupny označaje jego svez s grupoju. Zapis ne ustanavjaje, že vsaky determinant matrice iz prědstavjenja grupy možno bez dodatnogo uslovja nazvati grupnym determinantom.\par
\noindent\emph{Reason:} Grupny determinant selects a direct name for the Frobenius-Dedekind object. The rationale rejects a modern descriptive substitution but records no exact competing Interslavic phrase. Determinant retains a recognizable mathematical stem, while grupny expresses its relation to a group. The record does not establish that every determinant of a matrix in a group representation is thereby the group determinant.\par
\noindent{\footnotesize Izvorne zapisy / Original records: HIST-g178-002.}\par
\end{minipage}\par
\par\medskip\Needspace{7\baselineskip}
\noindent\begin{minipage}{\linewidth}
\hypertarget{TERM-0745}{}
\noindent\textbf{grupova determinanta}\quad{\footnotesize\texttt{TERM-0745}}\par
\noindent\emph{Iztočnik / Source:} Gruppendeterminante\par
\noindent\emph{Rusky / Russian:} групповой определитель\par
\noindent\emph{Ukrajinsky / Ukrainian:} груповий визначник\par
\smallskip\noindent\emph{Razlog:} Grupova determinanta prodolžuje internacionalny korenj determinant-, dokolě rusko opredelitel i ukrajinsko vyznačnyk pokazyvajut ine nacionalne leksične vybory. Dva razloga podkrěpjajut terminologičny kontinuitet. Oni ne davajut osobnogo sravnjenja mužskogo determinant pri TERM-0744 s ženskoju determinanta ni pridavnikov grupny i grupovy. Zato slovnik sohranjaje obě historijne formy, ale ne predstavlja ih gramatično i stilistično ravnocěnnymi bez dodatnogo razsmotrenja.\par
\noindent\emph{Reason:} Grupova determinanta continues the international determinant- stem, while Russian opredelitel and Ukrainian vyznačnyk exhibit different national lexical choices. Two rationales support terminological continuity. They do not separately compare masculine determinant in TERM-0744 with feminine determinanta, or grupny with grupovy. Both historical forms are therefore preserved without presenting them as grammatically and stylistically interchangeable without further examination.\par
\noindent{\footnotesize Izvorne zapisy / Original records: HIST-g201-009, HIST-g207-006.}\par
\end{minipage}\par
\par\medskip\Needspace{7\baselineskip}
\noindent\begin{minipage}{\linewidth}
\hypertarget{TERM-0746}{}
\noindent\textbf{grupno razděljenje pridruženih prostyh idealov}\quad{\footnotesize\texttt{TERM-0746}}\par
\noindent\emph{Iztočnik / Source:} Gruppeneinteilung der zugehörigen Primideale\par
\noindent\emph{Rusky / Russian:} разбиение ассоциированных простых идеалов на группы\par
\noindent\emph{Ukrajinsky / Ukrainian:} поділ належних простих ідеалів на групи\par
\smallskip\noindent\emph{Razlog:} Grupno razděljenje pridruženih prostyh idealov opisuje proceduru v dokazě jedinosti: idealy sobirajut se prěz poslědovateljne odnosi děliteljstva i mnogokratnosti. Grupa tu znači čest razděljenja, ne grupu s binarnoju operacijeju. Razlog sohranjaje procedurny opis vměsto jeho zaměny obščim nazvom relacije ekvivalentnosti. Da se dobyje čest razděljenja, ne dostavaje sravniti samo dva početkove idealy; trěba slěditi cělomu konstruovanomu rědu svezij.\par
\noindent\emph{Reason:} Grupno razděljenje pridruženih prostyh idealov describes the uniqueness proof's procedure: ideals are collected through successive divisor and multiple relations. Group here means a part of a partition, not a group with a binary operation. The rationale retains the procedural description rather than replacing it with a general equivalence-relation label. Comparing just two initial ideals is not enough to obtain a part; the constructed chain of relations matters.\par
\noindent{\footnotesize Izvorne zapisy / Original records: HIST-g081-056, HIST-g082-056.}\par
\smallskip\noindent\emph{Sravnjenje s tekstom:} Rabota 19 tvori grupy, dodavajuči dělitelje i mnogokratne med danymi pridruženymi prostymi idealami, dokolě postupok ne prestane davati nove člany. Konečnost sobora osiguruje zastavjenje. Grupa tu jest čest razděljenja; tekst ne nadaje jej grupnoj operacije.\par
\noindent\emph{Textual comparison:} Paper 19 forms groups by adding divisors and multiples among the specified associated prime ideals until no new members appear. Finiteness ensures termination. Group here is a part of a partition; no group operation is assigned to it.\par
\noindent{\footnotesize TEXT-COMP-65; pasaže / passages: 13150, 13152, 13156; \texttt{CORPUS\_FORM\_COMPARISONS.json}.}\par
\end{minipage}\par
\par\medskip\Needspace{7\baselineskip}
\noindent\begin{minipage}{\linewidth}
\hypertarget{TERM-0747}{}
\noindent\textbf{grupno razděljenje po nevzaimnoj prostosti}\quad{\footnotesize\texttt{TERM-0747}}\par
\noindent\emph{Iztočnik / Source:} Gruppeneinteilung nach Nicht-Teilerfremdheit\par
\noindent\emph{Rusky / Russian:} разбиение на группы по невзаимной простоте\par
\noindent\emph{Ukrajinsky / Ukrainian:} поділ на групи за невзаємною простотою\par
\smallskip\noindent\emph{Razlog:} Grupno razděljenje po nevzaimnoj prostosti opisuje povezanje idealov konečnymi cěpjami nevzajemnoj prostosti. To ne tvrdi, že vsake dva idealy jednoj grupy neposrědnje ne sut vzajemno proste. Razlog dava prednost avtorskomu postupku prěd zaměnoju nazvom sveznyh komponentov grafa. Aktualny pasaž raboty 19 izrično tvorit take cěpi i obosnovuje zastavjenje postupka konečnostju sobora idealov; starša forma nevzaimnoj ostavaje vidima kako historijny izbor.\par
\noindent\emph{Reason:} Grupno razděljenje po nevzaimnoj prostosti describes connection by finite chains of non-coprimeness. It does not assert that every pair within a group is directly non-coprime. The rationale prefers the author's procedure to substituting graph connected-component terminology. The current Paper 19 passage explicitly constructs these chains and grounds termination in the finite collection of ideals; the older form nevzaimnoj remains visible as a historical choice.\par
\noindent{\footnotesize Izvorne zapisy / Original records: HIST-g081-066, HIST-g082-066.}\par
\smallskip\noindent\emph{Sravnjenje s tekstom:} V oddělě 8 raboty 19 člany jednoj česti svezany sut cěpjami nevzajemnoj prostosti, dokolě člany različnyh čestij sut vzajemno proste. To razlikuje cěpnu svez od neposrědnogo odnosa vsakoj pary. Procedurny prěvod čuva imenno to razlikovanje.\par
\noindent\emph{Textual comparison:} In Paper 19 section 8, members of a part are connected through chains of non-coprimeness, while members of different parts are coprime. This distinguishes connection by a chain from a direct relation between every pair. The procedural translation preserves that distinction.\par
\noindent{\footnotesize TEXT-COMP-66; pasaže / passages: 13302, 13304, 13306; \texttt{CORPUS\_FORM\_COMPARISONS.json}.}\par
\end{minipage}\par
\par\medskip\Needspace{7\baselineskip}
\noindent\begin{minipage}{\linewidth}
\hypertarget{TERM-0748}{}
\noindent\textbf{porjadok grupy / porjadkov $h_i$}\quad{\footnotesize\texttt{TERM-0748}}\par
\noindent\emph{Iztočnik / Source:} Gruppenordnung / Ordnungen $h_i$\par
\noindent\emph{Rusky / Russian:} порядок группы / порядков $h_i$\par
\noindent\emph{Ukrajinsky / Ukrainian:} порядок групи / порядків $h_i$\par
\smallskip\noindent\emph{Razlog:} Porjadok grupy / porjadkov $h_i$ zadržuje tehničny korenj porjadok pri imenovanju poredkov grup. Oznaky $h_i$ razlikujut konkretne veličiny; oni ne sut čest jezyčnogo nazva. Razlog dava prednost kontinuitetu s ranějšim algebraičnym upotrěbjenjem, no sam isti korenj ne zjedinjava vse smysly slova porjadok. Osoblivo porjadok grupy ne treba směšati s poredkom koeficientov ili s uporjadkovanjem elementov.\par
\noindent\emph{Reason:} Porjadok grupy / porjadkov $h_i$ retains the technical stem porjadok for group orders. The indexed $h_i$ distinguish particular quantities; they are not part of the lexical name. The rationale favours continuity with earlier algebraic usage, but a shared stem does not make every sense of order identical. Group order must in particular remain distinct from coefficient order or an ordering of elements.\par
\noindent{\footnotesize Izvorne zapisy / Original records: HIST-g200-003.}\par
\end{minipage}\par
\par\medskip\Needspace{7\baselineskip}
\noindent\begin{minipage}{\linewidth}
\hypertarget{TERM-0749}{}
\noindent\textbf{grupno kolco}\quad{\footnotesize\texttt{TERM-0749}}\par
\noindent\emph{Iztočnik / Source:} Gruppenring\par
\noindent\emph{Rusky / Russian:} групповое кольцо\par
\noindent\emph{Ukrajinsky / Ukrainian:} групове кільце\par
\smallskip\noindent\emph{Razlog:} Grupno kolco jest ranějši prědlog dlja grupnogo prstena. Razlog zadržuje korenj kolco v terminologiji modulov nad prstenom raboty 34; jego slovo canonical vyražaje projektno rěšenje, ne dokaz oficialnoj slovničnoj normy. Aktualny tekst na nekotoryh městah ima grupno prsten. Tuto kombinaciju ne treba prinimati kako gramatično podkrěpljenje starogo prědloga: izbor korenja i soglasovanje roda potrěbujut oddělne razloge. Sravnjenje pri TERM-0750 obosnovuje tu razliku.\par
\noindent\emph{Reason:} Grupno kolco is an earlier proposal for group ring. The rationale retains kolco in Paper 34's module-over-ring vocabulary; its word canonical denotes a project choice, not evidence of an official dictionary norm. Some current passages contain grupno prsten. That combination should not be taken as grammatical support for the old proposal: stem selection and gender agreement need separate justification. TERM-0750 supplies the comparison.\par
\noindent{\footnotesize Izvorne zapisy / Original records: HIST-g184-007.}\par
\smallskip\noindent\emph{Sravnjenje s tekstom:} Rabota 32 definira grupne čisla kako linearne kombinacije grupnyh elementov. Rabota 34 daje i konečnu grupu, i priměr beskonečnoj grupy s konečnymi sumami. Aktualne pasaže pišut grupno prsten ili grupovo prsten; to ne jest gramatična rekomendacija. Slovnik ima grupovy (16935), koljce srednjego roda (8518) i pŕstėnj mužskogo roda (30033). Poslědnji zapis ne jest bukveno identičny projektnomu prsten i ne atestuje algebraičny smysl. Gramatika trěbuje soglasovanja roda, padeža i čisla; stary srednji rod pridavnika ne može ostati neopravdano pri mužskom imeniku. To jest novějše jezyčno sravnjenje, ne dokaz ranějšej konsultacije.\par
\noindent\emph{Textual comparison:} Paper 32 defines group numbers as linear combinations of group elements. Paper 34 gives both a finite group and an infinite-group example with finite sums. The current passages read grupno prsten or grupovo prsten; these are not grammatical recommendations. The dictionary has grupovy (16935), neuter koljce (8518) and masculine pŕstėnj (30033). The last is not letter-for-letter identical with the project's prsten and does not attest the algebraic sense. The grammar requires gender, case and number agreement; an old neuter adjective cannot remain unjustified with a masculine noun. This is a fresh linguistic comparison, not evidence of earlier consultation.\par
\noindent{\footnotesize TEXT-COMP-67; pasaže / passages: 20287, 21829, 19280; \texttt{CORPUS\_FORM\_COMPARISONS.json}.}\par
\end{minipage}\par
\par\medskip\Needspace{7\baselineskip}
\noindent\begin{minipage}{\linewidth}
\hypertarget{TERM-0750}{}
\noindent\textbf{grupovo kolco}\quad{\footnotesize\texttt{TERM-0750}}\par
\noindent\emph{Iztočnik / Source:} Gruppenring\par
\noindent\emph{Rusky / Russian:} групповое кольцо\par
\noindent\emph{Ukrajinsky / Ukrainian:} групове кільце\par
\smallskip\noindent\emph{Razlog:} Grupovo kolco sohranjaje prěmu paru pridavnika i imenika vměsto zaměny cělago nazva slovom algebra. Razlogi naglašajut srodnost s ruskim i ukrajinskim izrazom i potrebnu obščnost: v primětke k Burnsideovoj teoremě možno govoriti takože o beskonečnoj grupě, s konečnymi linearnymi kombinacijami jej elementov. Tvrdženje o največšej prozračnosti ne jest čitatelsky rezultat. Novějše sravnjenje s grupovy, koljce i pŕstėnj razlikuje slovničnu podkrěpu, matematično razširenje smysla i potrěbno gramatično soglasovanje.\par
\noindent\emph{Reason:} Grupovo kolco keeps a direct adjective-noun phrase instead of replacing the whole name with algebra. The rationales emphasize the Russian and Ukrainian parallels and the necessary generality: the Burnside-theorem footnote also allows an infinite group, using finite linear combinations of its elements. Maximum transparency is not a reader-test result. A fresh comparison with grupovy, koljce and pŕstėnj distinguishes dictionary support, extension to a mathematical sense and the necessary grammatical agreement.\par
\noindent{\footnotesize Izvorne zapisy / Original records: HIST-g198-008, HIST-g200-004, HIST-g204-001.}\par
\smallskip\noindent\emph{Sravnjenje s tekstom:} Rabota 32 definira grupne čisla kako linearne kombinacije grupnyh elementov. Rabota 34 daje i konečnu grupu, i priměr beskonečnoj grupy s konečnymi sumami. Aktualne pasaže pišut grupno prsten ili grupovo prsten; to ne jest gramatična rekomendacija. Slovnik ima grupovy (16935), koljce srednjego roda (8518) i pŕstėnj mužskogo roda (30033). Poslědnji zapis ne jest bukveno identičny projektnomu prsten i ne atestuje algebraičny smysl. Gramatika trěbuje soglasovanja roda, padeža i čisla; stary srednji rod pridavnika ne može ostati neopravdano pri mužskom imeniku. To jest novějše jezyčno sravnjenje, ne dokaz ranějšej konsultacije.\par
\noindent\emph{Textual comparison:} Paper 32 defines group numbers as linear combinations of group elements. Paper 34 gives both a finite group and an infinite-group example with finite sums. The current passages read grupno prsten or grupovo prsten; these are not grammatical recommendations. The dictionary has grupovy (16935), neuter koljce (8518) and masculine pŕstėnj (30033). The last is not letter-for-letter identical with the project's prsten and does not attest the algebraic sense. The grammar requires gender, case and number agreement; an old neuter adjective cannot remain unjustified with a masculine noun. This is a fresh linguistic comparison, not evidence of earlier consultation.\par
\noindent{\footnotesize TEXT-COMP-67; pasaže / passages: 20287, 21829, 19280; \texttt{CORPUS\_FORM\_COMPARISONS.json}.}\par
\end{minipage}\par
\par\medskip\Needspace{7\baselineskip}
\noindent\begin{minipage}{\linewidth}
\hypertarget{TERM-0751}{}
\noindent\textbf{grupno kolco / grupne čisla}\quad{\footnotesize\texttt{TERM-0751}}\par
\noindent\emph{Iztočnik / Source:} Gruppenring / Gruppenzahlen\par
\noindent\emph{Rusky / Russian:} групповое кольцо / групповые числа\par
\noindent\emph{Ukrajinsky / Ukrainian:} групове кільце / групові числа\par
\smallskip\noindent\emph{Razlog:} Grupno kolco / grupne čisla sohranjaje i nazv algebraičnoj struktury, i historijno imenovanje jej elementov. V rabotě 32 grupne čisla sut linearne kombinacije grupnyh elementov s koeficientami iz izhodnogo polja; oni ne značet čislo elementov grupy. Razlog zato ne svodi pasaž do samogo sovrěmennogo nazva grupova algebra. Čitatelj mora viděti, kako struktura nastavaje, a ne samo k kakomu klasu ona naleži.\par
\noindent\emph{Reason:} Grupno kolco / grupne čisla retains both the algebraic structure's name and the historical name for its elements. In Paper 32 group numbers are linear combinations of group elements with coefficients in the starting field; they do not mean the number of elements in the group. The rationale therefore does not reduce the passage to the modern label group algebra. The construction, not just its classification, must remain visible.\par
\noindent{\footnotesize Izvorne zapisy / Original records: HIST-g170-003.}\par
\smallskip\noindent\emph{Sravnjenje s tekstom:} Rabota 32 definira grupne čisla kako linearne kombinacije grupnyh elementov. Rabota 34 daje i konečnu grupu, i priměr beskonečnoj grupy s konečnymi sumami. Aktualne pasaže pišut grupno prsten ili grupovo prsten; to ne jest gramatična rekomendacija. Slovnik ima grupovy (16935), koljce srednjego roda (8518) i pŕstėnj mužskogo roda (30033). Poslědnji zapis ne jest bukveno identičny projektnomu prsten i ne atestuje algebraičny smysl. Gramatika trěbuje soglasovanja roda, padeža i čisla; stary srednji rod pridavnika ne može ostati neopravdano pri mužskom imeniku. To jest novějše jezyčno sravnjenje, ne dokaz ranějšej konsultacije.\par
\noindent\emph{Textual comparison:} Paper 32 defines group numbers as linear combinations of group elements. Paper 34 gives both a finite group and an infinite-group example with finite sums. The current passages read grupno prsten or grupovo prsten; these are not grammatical recommendations. The dictionary has grupovy (16935), neuter koljce (8518) and masculine pŕstėnj (30033). The last is not letter-for-letter identical with the project's prsten and does not attest the algebraic sense. The grammar requires gender, case and number agreement; an old neuter adjective cannot remain unjustified with a masculine noun. This is a fresh linguistic comparison, not evidence of earlier consultation.\par
\noindent{\footnotesize TEXT-COMP-67; pasaže / passages: 20287, 21829, 19280; \texttt{CORPUS\_FORM\_COMPARISONS.json}.}\par
\end{minipage}\par
\par\medskip\Needspace{7\baselineskip}
\noindent\begin{minipage}{\linewidth}
\hypertarget{TERM-0752}{}
\noindent\textbf{grupno kolco / kvaternionno tělo}\quad{\footnotesize\texttt{TERM-0752}}\par
\noindent\emph{Iztočnik / Source:} Gruppenring / Quaternionenkoerper\par
\noindent\emph{Rusky / Russian:} групповое кольцо / тело кватернионов\par
\noindent\emph{Ukrajinsky / Ukrainian:} групове кільце / кватерніонне тіло\par
\smallskip\noindent\emph{Razlog:} Grupno kolco / kvaternionno tělo sohranjaje razliku med prstenovoju konstrukcijeju i nekomutativnym tělom. Razlog prodolžuje izrično razděljenje ring/Körper v rabotě 32: prěvod ne smě iz samogo slova quaternion zaključiti komutativno polje. Takože ne treba tvrditi, že vsaka kvaternionna algebra jest tělo; svojstvo dělimosti mora byti danymi uslovjami osigurjeno. Starša fraza kolco i aktualna terminologija prstena ostavajut različnymi stupenjami jezyčnogo rěšenja.\par
\noindent\emph{Reason:} Grupno kolco / kvaternionno tělo preserves the distinction between a ring construction and a noncommutative division ring. The rationale continues Paper 32's explicit ring/Körper distinction: quaternion must not automatically become a commutative field. Nor does every quaternion algebra thereby become a division ring; the relevant hypotheses must guarantee that property. The earlier kolco phrase and current ring terminology remain distinct stages of the linguistic choice.\par
\noindent{\footnotesize Izvorne zapisy / Original records: HIST-g176-002.}\par
\end{minipage}\par
\par\medskip\Needspace{7\baselineskip}
\noindent\begin{minipage}{\linewidth}
\hypertarget{TERM-0753}{}
\noindent\textbf{elementy iz $\Omega$}\quad{\footnotesize\texttt{TERM-0753}}\par
\noindent\emph{Iztočnik / Source:} Größen aus $\Omega$\par
\noindent\emph{Rusky / Russian:} элементы из $\Omega$\par
\noindent\emph{Ukrajinsky / Ukrainian:} елементи з $\Omega$\par
\smallskip\noindent\emph{Razlog:} Elementy iz $\Omega$ zaměnjaje široke německe Größen izričnym označenjem elementov koeficientnogo polja. Razlog odnosi se k konstantam ili skaljaram, ne k izměrjenym fizičnym veličinam. Dodatok iz $\Omega$ sohranjaje članstvo v točnoj oblasti; on ne jest dekorativna oznaka. Iz togo ne slěduje, že vsako Größen v korpusu možno tako prěvesti: pri TERM-0754 ono služi drugomu razděljenju objektov.\par
\noindent\emph{Reason:} Elementy iz $\Omega$ replaces broad German Größen with explicit membership in the coefficient field. The rationale concerns constants or scalars, not measured physical quantities. The qualification from $\Omega$ preserves membership in a specified domain; it is not decorative notation. This does not license the same rendering of every Größen in the corpus: TERM-0754 uses it in a different classification.\par
\noindent{\footnotesize Izvorne zapisy / Original records: HIST-g047-006.}\par
\end{minipage}\par
\par\medskip\Needspace{7\baselineskip}
\noindent\begin{minipage}{\linewidth}
\hypertarget{TERM-0754}{}
\noindent\textbf{veličiny $y,z,\vartheta$}\quad{\footnotesize\texttt{TERM-0754}}\par
\noindent\emph{Iztočnik / Source:} Größen $y,z,\vartheta$\par
\noindent\emph{Rusky / Russian:} величины $y,z,\vartheta$\par
\noindent\emph{Ukrajinsky / Ukrainian:} величини $y,z,\vartheta$\par
\smallskip\noindent\emph{Razlog:} Veličiny $y,z,\vartheta$ zadržuje široky imenik i tri različne matematične oznaky. V rekursivnoj konstrukciji raboty 9 klasa $y$ pri adjunkciji vvodit algebraično-drobno čislo; $z$ ne ima togo svojstva i jest racionalno izrazimo prěz konečno mnogo poprědnyh čisel, različnyh od $y$; $\vartheta$ ne naleži nijednomu iz prvih dvoh slučajev. Zato oznaky ne možut byti jezyčno zamenjane ili sjedinjane. Razlog čuva strukturu razděljenja, ne samo tipografiju bukv.\par
\noindent\emph{Reason:} Veličiny $y,z,\vartheta$ retains a broad noun and three distinct mathematical labels. In Paper 9's recursive construction, adjoining a $y$ introduces an algebraically fractional number; a $z$ lacks that property and is rationally expressible through finitely many earlier numbers other than the $y$'s; a $\vartheta$ belongs to neither of the first two cases. The labels must therefore not be translated into one another or merged. The reason protects the classification, not merely the letters' typography.\par
\noindent{\footnotesize Izvorne zapisy / Original records: HIST-g065-004.}\par
\smallskip\noindent\emph{Sravnjenje s tekstom:} Rabota 9 razděljuje veličiny po tomu, dali jih adjunkcija vvodit algebraično-drobno čislo, dali sut racionalno izrazime prěz konečno mnogo poprědnyh čisel izvan klasa $y$, ili ne izpolnjajut nijedno iz tih uslovij. Oznaky $y,z,\vartheta$ sohranjajut točno tu trojnu razliku; njihova zaměna byla by proměnoju konstrukcije, ne stilističnym uproščenjem.\par
\noindent\emph{Textual comparison:} Paper 9 separates quantities according to whether adjoining them introduces an algebraically fractional number, whether they are rationally expressible through finitely many earlier numbers outside the $y$ class, or whether neither condition holds. The labels $y,z,\vartheta$ preserve that three-way distinction; merging them would change the construction, not merely simplify its style.\par
\noindent{\footnotesize TEXT-COMP-70; pasaže / passages: 7854, 7855, 7857, 7858, 7859, 7860, 7864, 7865, 7866, 7867, 7868, 7871, 7872; \texttt{CORPUS\_FORM\_COMPARISONS.json}.}\par
\end{minipage}\par
\par\medskip\Needspace{7\baselineskip}
\noindent\begin{minipage}{\linewidth}
\hypertarget{TERM-0755}{}
\noindent\textbf{red veličin / redy veličin}\quad{\footnotesize\texttt{TERM-0755}}\par
\noindent\emph{Iztočnik / Source:} Größenreihe / Größenreihen\par
\noindent\emph{Rusky / Russian:} ряд величин / ряды величин\par
\noindent\emph{Ukrajinsky / Ukrainian:} ряд величин / ряди величин\par
\smallskip\noindent\emph{Razlog:} Red veličin / redy veličin sohranjaje pojmovnu svez s ranějšimi rabotami i so simetričnymi funkcijami redov. Razlog podčrkava uporjadkovane redy, ne sumy članov beskonečnogo analitičnogo rěda. Razlika med jedinym redom i množinoju redov jest matematično važna: simetrija može odnositi se k zaměnjanju cělyh redov. Historijni zapis dava kontinuitet, ale ne zasebny test čitatelskoj razumělivosti.\par
\noindent\emph{Reason:} Red veličin / redy veličin retains the connection with earlier papers and with symmetric functions of rows. The rationale emphasizes ordered rows, not sums of terms of an infinite analytic series. The distinction between a row and a collection of rows matters mathematically: symmetry may concern permutation of whole rows. The historical record supports continuity, not a separate reader-comprehension test.\par
\noindent{\footnotesize Izvorne zapisy / Original records: HIST-g056-004.}\par
\end{minipage}\par
\par\medskip\Needspace{7\baselineskip}
\noindent\begin{minipage}{\linewidth}
\hypertarget{TERM-0756}{}
\noindent\textbf{redy veličin}\quad{\footnotesize\texttt{TERM-0756}}\par
\noindent\emph{Iztočnik / Source:} Größenreihen\par
\noindent\emph{Rusky / Russian:} ряды величин\par
\noindent\emph{Ukrajinsky / Ukrainian:} ряди величин\par
\smallskip\noindent\emph{Razlog:} Redy veličin sohranjaje možnost pozdnějšego geometričnogo tolkovanja redov kako toček. Razlog odkazuje takože k anglijskomu kontrolnomu izrazu rows of quantities. To jest zapisana pomoć pri izbore, ne nezavisna matematična avtoritetnost anglijskogo prěvoda. Red tu treba čitati po lokalnoj konstrukciji, ne avtomatično kako analitičny rěd ili konečnu sumu.\par
\noindent\emph{Reason:} Redy veličin preserves the later geometric interpretation of rows as points. The rationale also refers to the English control wording rows of quantities. That is a recorded aid to selection, not independent mathematical authority for the English translation. The local construction determines what a row is; it must not automatically be read as an analytic series or a finite sum.\par
\noindent{\footnotesize Izvorne zapisy / Original records: HIST-g045-008.}\par
\end{minipage}\par
\par\medskip\Needspace{7\baselineskip}
\noindent\begin{minipage}{\linewidth}
\hypertarget{TERM-0757}{}
\noindent\textbf{rędy veličin / rędy}\quad{\footnotesize\texttt{TERM-0757}}\par
\noindent\emph{Iztočnik / Source:} Größenreihen / Reihen\par
\noindent\emph{Rusky / Russian:} ряды величин / ряды\par
\noindent\emph{Ukrajinsky / Ukrainian:} ряди величин / ряди\par
\smallskip\noindent\emph{Razlog:} Rędy veličin / rędy sohranjaje redy, na ktore dějstvuje isto permutacijno dějanje. Razlog izrično odbija čitanje kak analitičnyh rędov. Bukva ę v historijnoj formě jest pravopisna čerta, ne znak drugogo matematičnogo pojma proti redy pri TERM-0755 i TERM-0756. V osnovnom pravopisu ę piše se kako e; identičnost smysla ipak trěba oprěděljati iz konteksta, ne samo iz takoje zaměny bukvy.\par
\noindent\emph{Reason:} Rędy veličin / rędy preserves rows subject to the same permutation action. The rationale expressly rejects the analytic-series reading. Historical ę is an orthographic feature, not evidence of a different mathematical concept from redy in TERM-0755 and TERM-0756. Basic spelling writes ę as e; contextual meaning still has to be established rather than inferred solely from that letter substitution.\par
\noindent{\footnotesize Izvorne zapisy / Original records: HIST-g077-007.}\par
\end{minipage}\par
\par\medskip\Needspace{7\baselineskip}
\noindent\begin{minipage}{\linewidth}
\hypertarget{TERM-0758}{}
\noindent\textbf{največa komutativna podpolja}\quad{\footnotesize\texttt{TERM-0758}}\par
\noindent\emph{Iztočnik / Source:} größte kommutative Unterkörper\par
\noindent\emph{Rusky / Russian:} максимальные коммутативные подтела\par
\noindent\emph{Ukrajinsky / Ukrainian:} найбільші комутативні підполя\par
\smallskip\noindent\emph{Razlog:} Največa komutativna podpolja sohranjaje komutativnost podoblastij v nekomutativnom tělu. Razlog dopušča rusko podtelo s pridavnikom komutativny, dokolě medžuslovjansko podpolje izrično kaže polje. Množinsky nazv ne tvrdi jedinogo največšego podpolja, ktoro sodrživaje vse ostale. Aktualny tekst sravnjaje stepeni takyh podpolj i razpadnyh polj; daljša jednakost stepenjev v matričnom slučaju ima dodatno uslovje regularnosti.\par
\noindent\emph{Reason:} Največa komutativna podpolja retains the commutativity of subobjects inside a noncommutative division ring. The rationale permits Russian podtelo qualified by commutative, while Interslavic podpolje explicitly marks a field. The plural name does not assert a unique greatest subfield containing all the others. The current text compares degrees of these subfields and splitting fields; the further equality of degrees in the matrix case has an additional regularity hypothesis.\par
\noindent{\footnotesize Izvorne zapisy / Original records: HIST-g169-003, HIST-g171-003.}\par
\smallskip\noindent\emph{Sravnjenje s tekstom:} V rabotě 32 največše komutativna podpolja od $\mathfrak A$ imajut stepenj $m$. Dlja matričnogo razširjenja $\mathfrak A_r$ tekst dava stepeni $ms$, kde $s$ děli $r$; rovnost s $mr$ jest tvrdžena pod dodatnym uslovjem regularnosti. Stepenjevi zaključki zato ne smějut byti prěnosene bez uslovij, ni množinsky nazv proměnen v tvrdženje jedinogo največšego podpolja.\par
\noindent\emph{Textual comparison:} In Paper 32 the greatest commutative subfields of $\mathfrak A$ have degree $m$. For the matrix extension $\mathfrak A_r$, the text gives degrees $ms$ with $s$ dividing $r$; equality with $mr$ is asserted under an additional regularity condition. These degree conclusions must not be transferred without their hypotheses, nor the plural name turned into a claim of a unique greatest subfield.\par
\noindent{\footnotesize TEXT-COMP-68; pasaže / passages: 19318, 19356; \texttt{CORPUS\_FORM\_COMPARISONS.json}.}\par
\end{minipage}\par
\par\medskip\Needspace{7\baselineskip}
\noindent\begin{minipage}{\linewidth}
\hypertarget{TERM-0759}{}
\noindent\textbf{najveće primarne idealy}\quad{\footnotesize\texttt{TERM-0759}}\par
\noindent\emph{Iztočnik / Source:} größte primäre Ideale\par
\noindent\emph{Rusky / Russian:} наибольшие примарные идеалы\par
\noindent\emph{Ukrajinsky / Ukrainian:} найбільші примарні ідеали\par
\smallskip\noindent\emph{Razlog:} Najveće primarne idealy sohranjaje avtorsky nazv v kontekstu razloženja, ne obšči nazv maksimalnogo ideala. Definicija IV raboty 19 trěbuje najkratše prědstavjenje s primarnymi komponentami, pri ktorom obče mnogokratno dvoh komponentov už ne jest primarno. Razlog opisuje sobiranje nerazložimyh idealov v takove komponenty. Slovo največši zato ne oprěděljaje prostě najveliky sobor elementov; ono zavisi od konkretnoj struktury prědstavjenja.\par
\noindent\emph{Reason:} Najveće primarne idealy retains the author's name within a decomposition, not the general term maximal ideal. Paper 19's Definition IV requires a shortest representation by primary components for which the common multiple of two components is no longer primary. The rationale describes collecting indecomposable ideals into such components. Greatest therefore does not simply mean the largest set of elements; its force depends on this representation structure.\par
\noindent{\footnotesize Izvorne zapisy / Original records: HIST-g081-046, HIST-g082-046.}\par
\smallskip\noindent\emph{Sravnjenje s tekstom:} Definicija IV raboty 19 svezuje največše primarne idealy k najkratšemu prědstavjenju i k tomu, že obče mnogokratno dvoh komponentov už ne jest primarno. V rabotě 31 kriterij diskriminanta potrebuje vlastno primarnu, ne prostu komponentu ideala $(p)$. Tute frazy zato ne možut byti skračene do maksimalny ideal ili prosty ideal.\par
\noindent\emph{Textual comparison:} Paper 19's Definition IV ties greatest primary ideals to a shortest representation in which the common multiple of two components is no longer primary. Paper 31's discriminant criterion requires a properly primary, non-prime component of $(p)$. Neither phrase can therefore be shortened to maximal ideal or prime ideal.\par
\noindent{\footnotesize TEXT-COMP-69; pasaže / passages: 12974, 12976, 17267; \texttt{CORPUS\_FORM\_COMPARISONS.json}.}\par
\end{minipage}\par
\par\medskip\Needspace{7\baselineskip}
\noindent\begin{minipage}{\linewidth}
\hypertarget{TERM-0760}{}
\noindent\textbf{največe primarne komponenty}\quad{\footnotesize\texttt{TERM-0760}}\par
\noindent\emph{Iztočnik / Source:} größte primäre Komponenten\par
\noindent\emph{Rusky / Russian:} наибольшие первичные компоненты\par
\noindent\emph{Ukrajinsky / Ukrainian:} найбільші первинні компоненти\par
\smallskip\noindent\emph{Razlog:} Največe primarne komponenty prodolžuje terminologiju primarnogo razloženja i sohranjaje rolju komponenty danogo ideala. V uvodu raboty 31 razsmatrjaje se razloženje glavnogo ideala $(p)$ v konkretnom konečnom poredku. Uslovje diskriminanta trěbuje, aby prinajmenje jedna komponenta byla vlastno primarna, t. j. ne prosta. Primarny zato ne jest synonym prosty, a komponenta ne jest proizvoljny podideal. Razlog zadržuje tute funkcije bez novogo nazva dlja vsakoj upotrěby.\par
\noindent\emph{Reason:} Največe primarne komponenty continues primary-decomposition terminology and preserves the role of a component of the specified ideal. Paper 31's introduction concerns the principal ideal $(p)$ in a particular finite order. The discriminant condition requires at least one component to be properly primary, that is, not prime. Primary is therefore not synonymous with prime, and a component is not an arbitrary subideal. The rationale retains these roles without inventing a new name for each occurrence.\par
\noindent{\footnotesize Izvorne zapisy / Original records: HIST-g118-004.}\par
\smallskip\noindent\emph{Sravnjenje s tekstom:} Definicija IV raboty 19 svezuje največše primarne idealy k najkratšemu prědstavjenju i k tomu, že obče mnogokratno dvoh komponentov už ne jest primarno. V rabotě 31 kriterij diskriminanta potrebuje vlastno primarnu, ne prostu komponentu ideala $(p)$. Tute frazy zato ne možut byti skračene do maksimalny ideal ili prosty ideal.\par
\noindent\emph{Textual comparison:} Paper 19's Definition IV ties greatest primary ideals to a shortest representation in which the common multiple of two components is no longer primary. Paper 31's discriminant criterion requires a properly primary, non-prime component of $(p)$. Neither phrase can therefore be shortened to maximal ideal or prime ideal.\par
\noindent{\footnotesize TEXT-COMP-69; pasaže / passages: 12974, 12976, 17267; \texttt{CORPUS\_FORM\_COMPARISONS.json}.}\par
\end{minipage}\par
\par\medskip\Needspace{7\baselineskip}
\noindent\begin{minipage}{\linewidth}
\hypertarget{TERM-0761}{}
\noindent\textbf{največše primarne komponenty}\quad{\footnotesize\texttt{TERM-0761}}\par
\noindent\emph{Iztočnik / Source:} größte primäre Komponenten\par
\noindent\emph{Rusky / Russian:} наибольшие первичные компоненты\par
\noindent\emph{Ukrajinsky / Ukrainian:} найбільші первинні компоненти\par
\smallskip\noindent\emph{Razlog:} Največše primarne komponenty zadržuje i tip ideala, i jeho rolju v razloženju. Razlog iz raboty 24 prodolžuje terminologiju prvogo paragrafa; ne predlagaje novy naziv dlja istogo objekta. Primarny ne znači prosty, a komponenta ne znači proizvoljny ideal. Sravniteljne ruske i ukrajinske formy podpirajut izbranu strukturu frazy, ale same po sobě ne dokazujut, že čitatelji svih slovjanskyh jezykov budu razuměti ju jednako lehko. Največši mora čitati se v smyslu danogo razloženja, ne kak fizična veličina.\par
\noindent\emph{Reason:} Največše primarne komponenty retains both the type of ideal and its role in a decomposition. The rationale for Paper 24 continues the terminology of its first section rather than naming the same object anew. Primary does not mean prime, and a component is not an arbitrary ideal. The Russian and Ukrainian comparanda support the phrase's structure but do not establish equally easy comprehension among speakers of every Slavic language. Greatest must be understood within the specified decomposition, not as physical size.\par
\noindent{\footnotesize Izvorne zapisy / Original records: HIST-g096-008.}\par
\end{minipage}\par
\par\medskip\Needspace{7\baselineskip}
\noindent\begin{minipage}{\linewidth}
\hypertarget{TERM-0762}{}
\noindent\textbf{največše primarne komponenty / največše primarne faktory}\quad{\footnotesize\texttt{TERM-0762}}\par
\noindent\emph{Iztočnik / Source:} größte primäre Komponenten / größte primäre Faktoren\par
\noindent\emph{Rusky / Russian:} наибольшие первичные компоненты / наибольшие первичные множители\par
\noindent\emph{Ukrajinsky / Ukrainian:} найбільші первинні компоненти / найбільші первинні множники\par
\smallskip\noindent\emph{Razlog:} Dvě formy v toj zapisi ne sut svobodno zaměnime synonymy. Komponenta pripadaje razloženju ideala; faktor pripadaje razloženju normy. Razlog zadržuje istu osnovu primarn- i superlativ, ale měnjaje glavno slovo po matematičnoj funkciji. Tako čitatelj može viděti svęz medžu konstrukcijami bez utraty razlike medžu objektami. Historijne pravopisne varianty superlativa sut sohranjene v podatkah; ne predstavljajut dva različne algebraične pojmy ni čisleno sravnjenje kandidatov.\par
\noindent\emph{Reason:} The two forms in this record are not freely interchangeable synonyms. Component belongs to an ideal decomposition; factor belongs to a norm decomposition. The rationale keeps the primary stem and superlative while changing the head noun to suit the mathematical role. This makes the relationship between the constructions visible without identifying their objects. Historical spelling variants of the superlative remain in the accompanying data; they are neither two algebraic concepts nor a numerical comparison of candidates.\par
\noindent{\footnotesize Izvorne zapisy / Original records: HIST-g102-006.}\par
\end{minipage}\par
\par\medskip\Needspace{7\baselineskip}
\noindent\begin{minipage}{\linewidth}
\hypertarget{TERM-0763}{}
\noindent\textbf{najveće primarne komponenty / idealy}\quad{\footnotesize\texttt{TERM-0763}}\par
\noindent\emph{Iztočnik / Source:} größte primäre Komponenten / Ideale\par
\noindent\emph{Rusky / Russian:} наибольшие примарные компоненты / идеалы\par
\noindent\emph{Ukrajinsky / Ukrainian:} найбільші примарні компоненти / ідеали\par
\smallskip\noindent\emph{Razlog:} Komponenty i idealy davajut dva opisa istih objektov v danom razloženju: prvo slovo nazyvaje jih funkciju, drugo jih matematičny tip. Razlog smatra kvalifikaciju največši sučstvenoju dlja tvrdženj o jednoznačnosti. Zato ne možno skratiti frazu do primarne komponenty všude, kde ona nosi tu dodatnu informaciju. Taka komponenta ne stavaje iz samogo nazva maksimalnym idealom. Točna definicija največšego primarnogo ideala mora byti vzeta iz danoj raboty, kako objasnjeno pri TERM-0759.\par
\noindent\emph{Reason:} Components and ideals describe the same objects from two directions within the specified decomposition: the first names their role, the second their mathematical type. The rationale treats greatest as essential to the uniqueness assertions. Primary components is therefore not an adequate abbreviation wherever the qualification carries additional information. Nor does the name make such a component a maximal ideal. The definition of a greatest primary ideal must be taken from the particular paper, as explained at TERM-0759.\par
\noindent{\footnotesize Izvorne zapisy / Original records: HIST-g081-007, HIST-g082-007.}\par
\end{minipage}\par
\par\medskip\Needspace{7\baselineskip}
\noindent\begin{minipage}{\linewidth}
\hypertarget{TERM-0764}{}
\noindent\textbf{največša primarna komponenta}\quad{\footnotesize\texttt{TERM-0764}}\par
\noindent\emph{Iztočnik / Source:} größte Primärkomponente\par
\noindent\emph{Rusky / Russian:} наибольшая первичная компонента\par
\noindent\emph{Ukrajinsky / Ukrainian:} найбільша первинна компонента\par
\smallskip\noindent\emph{Razlog:} Največša primarna komponenta sohranjaje děliteljsku terminologiju raboty 30. Največša ne jest ocěna važnosti ni počta elementov. V šestom paragrafu komponenty najkratšego prědstavjenja nazyvajut se največšimi zato, že najmenše obče mnogokratno kojihkoli dvuh iz njih už ne jest primarno. Izbor komponenta zadržuje rolju v razloženju, a primarna zadržuje nužnu razliku od prosta. Definicija daje frazě bolše preciznosti než samo sravnjenje jezyčnyh korenjev.\par
\noindent\emph{Reason:} Največša primarna komponenta preserves the divisibility vocabulary of Paper 30. Greatest is not a measure of importance or cardinality. In Section 6, the components of a shortest representation are called greatest because the least common multiple of any two of them is no longer primary. Component retains their role in the decomposition, while primary preserves the distinction from prime. The definition gives the phrase greater precision than a comparison of word roots alone could provide.\par
\noindent{\footnotesize Izvorne zapisy / Original records: HIST-g113-011.}\par
\smallskip\noindent\emph{Sravnjenje s tekstom:} Rabota 30 nazyvaje primarne komponenty največšimi v točno oprěděljenom smyslu najkratšego prědstavjenja. Teorema IV izključaje iz tvrdženja o jednoznačnosti komponenty naležeče k $\mathfrak{o}$. To ne jest tvrdženje, že vsaka taka komponenta jest jednaka $\mathfrak{o}$: dokaz togo rovnostva pozdněj upotrěbjaje egzistenciju jediničnogo elementa. Pasaže podpirajut potrěbu sohraniti oba različenja.\par
\noindent\emph{Textual comparison:} Paper 30 calls primary components greatest in the specified sense of a shortest representation. Theorem IV excludes components belonging to $\mathfrak{o}$ from its uniqueness assertion. This does not say that every such component equals $\mathfrak{o}$: the subsequent proof of that equality uses the existence of an identity element. The passages support retaining both distinctions.\par
\noindent{\footnotesize TEXT-COMP-71; pasaže / passages: 16717, 16719, 16749, 16780, 16788, 16792, 16794, 16800; \texttt{CORPUS\_FORM\_COMPARISONS.json}.}\par
\end{minipage}\par
\par\medskip\Needspace{7\baselineskip}
\noindent\begin{minipage}{\linewidth}
\hypertarget{TERM-0765}{}
\noindent\textbf{največše primarne komponenty / izolovane komponenty}\quad{\footnotesize\texttt{TERM-0765}}\par
\noindent\emph{Iztočnik / Source:} größte Primärkomponenten / isolierte Komponenten\par
\noindent\emph{Rusky / Russian:} наибольшие первичные компоненты / изолированные компоненты\par
\noindent\emph{Ukrajinsky / Ukrainian:} найбільші первинні компоненти / ізольовані компоненти\par
\smallskip\noindent\emph{Razlog:} Največše primarne komponenty i izolovane komponenty ne označajut dve proizvoljne stilistične možnosti. Prva fraza opisuje vrstu komponent v razloženju; druga izděljuje posebnu klasu, za koju tekst potvrđuje jednoznačnost. Razlog čuva obě kvalifikacije i obču osnovu komponent-. Aktualna pasaža o najkratšem prědstavjenju otdělno upominaje asociovane proste idealy i izolovane komponenty. Ona ne daje osnovy zaměniti vsak pojav največši slovom izolovany ili obratno.\par
\noindent\emph{Reason:} Največše primarne komponenty and izolovane komponenty are not arbitrary stylistic alternatives. The first describes components in the decomposition; the second distinguishes a particular class whose uniqueness the text asserts. The rationale preserves both qualifications and the common component stem. The current passage on a shortest representation separately mentions associated prime ideals and isolated components. It does not justify replacing every occurrence of greatest with isolated, or conversely.\par
\noindent{\footnotesize Izvorne zapisy / Original records: HIST-g113-005.}\par
\smallskip\noindent\emph{Sravnjenje s tekstom:} Rabota 30 nazyvaje primarne komponenty največšimi v točno oprěděljenom smyslu najkratšego prědstavjenja. Teorema IV izključaje iz tvrdženja o jednoznačnosti komponenty naležeče k $\mathfrak{o}$. To ne jest tvrdženje, že vsaka taka komponenta jest jednaka $\mathfrak{o}$: dokaz togo rovnostva pozdněj upotrěbjaje egzistenciju jediničnogo elementa. Pasaže podpirajut potrěbu sohraniti oba različenja.\par
\noindent\emph{Textual comparison:} Paper 30 calls primary components greatest in the specified sense of a shortest representation. Theorem IV excludes components belonging to $\mathfrak{o}$ from its uniqueness assertion. This does not say that every such component equals $\mathfrak{o}$: the subsequent proof of that equality uses the existence of an identity element. The passages support retaining both distinctions.\par
\noindent{\footnotesize TEXT-COMP-71; pasaže / passages: 16717, 16719, 16749, 16780, 16788, 16792, 16794, 16800; \texttt{CORPUS\_FORM\_COMPARISONS.json}.}\par
\end{minipage}\par
\par\medskip\Needspace{7\baselineskip}
\noindent\begin{minipage}{\linewidth}
\hypertarget{TERM-0766}{}
\noindent\textbf{največše primarne komponenty / primarna komponenta naležeča k $\mathfrak{o}$}\quad{\footnotesize\texttt{TERM-0766}}\par
\noindent\emph{Iztočnik / Source:} größte Primärkomponenten / zu $\mathfrak{o}$ gehörige Primärkomponente\par
\noindent\emph{Rusky / Russian:} наибольшие первичные компоненты / первичная компонента, принадлежащая $\mathfrak{o}$\par
\noindent\emph{Ukrajinsky / Ukrainian:} найбільші первинні компоненти / первинна компонента, що належить до $\mathfrak{o}$\par
\smallskip\noindent\emph{Razlog:} Fraza naležeča k $\mathfrak{o}$ sohranjaje odnos komponenty k asociovanomu idealu; ne tvrdži, že sama komponenta jest jednaka $\mathfrak{o}$. To jest nužno v Teoremě IV raboty 30: jednoznačnost pri dvojnom uslovju cěpov dana jest komponentam, ktore ne naležat k jediničnomu idealu. Dalši tekst razděljuje obči slučaj od slučaja s jediničnym elementom. V tom drugom slučaju dokazuje se, že komponenta naležeča k $\mathfrak{o}$ mora byti jednaka $\mathfrak{o}$. Vynešenje togo dodatnogo uslovja by změnilo smysl izjatyja.\par
\noindent\emph{Reason:} The phrase naležeča k $\mathfrak{o}$ preserves a component's relation to its associated ideal; it does not say that the component itself equals $\mathfrak{o}$. This matters in Theorem IV of Paper 30, where uniqueness under the double chain condition is asserted for components not belonging to the unit ideal. The subsequent discussion distinguishes the general case from the case with an identity element. Only in the latter does it prove that a component belonging to $\mathfrak{o}$ must equal $\mathfrak{o}$. Dropping that additional hypothesis would change the exception's meaning.\par
\noindent{\footnotesize Izvorne zapisy / Original records: HIST-g114-004.}\par
\smallskip\noindent\emph{Sravnjenje s tekstom:} Rabota 30 nazyvaje primarne komponenty največšimi v točno oprěděljenom smyslu najkratšego prědstavjenja. Teorema IV izključaje iz tvrdženja o jednoznačnosti komponenty naležeče k $\mathfrak{o}$. To ne jest tvrdženje, že vsaka taka komponenta jest jednaka $\mathfrak{o}$: dokaz togo rovnostva pozdněj upotrěbjaje egzistenciju jediničnogo elementa. Pasaže podpirajut potrěbu sohraniti oba različenja.\par
\noindent\emph{Textual comparison:} Paper 30 calls primary components greatest in the specified sense of a shortest representation. Theorem IV excludes components belonging to $\mathfrak{o}$ from its uniqueness assertion. This does not say that every such component equals $\mathfrak{o}$: the subsequent proof of that equality uses the existence of an identity element. The passages support retaining both distinctions.\par
\noindent{\footnotesize TEXT-COMP-71; pasaže / passages: 16717, 16719, 16749, 16780, 16788, 16792, 16794, 16800; \texttt{CORPUS\_FORM\_COMPARISONS.json}.}\par
\end{minipage}\par
\par\medskip\Needspace{7\baselineskip}
\noindent\begin{minipage}{\linewidth}
\hypertarget{TERM-0767}{}
\noindent\textbf{največi obči dělitelj}\quad{\footnotesize\texttt{TERM-0767}}\par
\noindent\emph{Iztočnik / Source:} größter gemeinsamer Teiler\par
\noindent\emph{Rusky / Russian:} наибольший общий делитель\par
\noindent\emph{Ukrajinsky / Ukrainian:} найбільший спільний дільник\par
\smallskip\noindent\emph{Razlog:} Največi obči dělitelj sohranjaje staršu děliteljsku terminologiju tam, kde Noether objasnjaje modulnu sumu. Razlog za rabotu 31 ne rekomenduje običny čislovy algoritm: objašnjenje dolžno uputiti čitatelja na operaciju nad modulami. Suma dava obči dělitelj v avtorovom poredku děliteljnosti. Zadržanje nazva i izrično objasnjenje sumy služat različnym cěljam: prvo sohranjaje jezyk argumenta, drugo pomogaje sovrěmennomu čitatelju razpoznati operaciju.\par
\noindent\emph{Reason:} Največi obči dělitelj retains the older divisibility vocabulary where Noether explains module addition. The rationale for Paper 31 does not recommend an ordinary numerical algorithm: the explanation must direct the reader to an operation on modules. Their sum supplies the common divisor in the author's divisibility order. Retaining the name and explicitly explaining the sum serve different purposes: one preserves the argument's language, the other helps a modern reader recognize the operation.\par
\noindent{\footnotesize Izvorne zapisy / Original records: HIST-g120-006.}\par
\end{minipage}\par
\par\medskip\Needspace{7\baselineskip}
\noindent\begin{minipage}{\linewidth}
\hypertarget{TERM-0768}{}
\noindent\textbf{največši obči dělitelj}\quad{\footnotesize\texttt{TERM-0768}}\par
\noindent\emph{Iztočnik / Source:} größter gemeinsamer Teiler\par
\noindent\emph{Rusky / Russian:} наибольший общий делитель\par
\noindent\emph{Ukrajinsky / Ukrainian:} найбільший спільний дільник\par
\smallskip\noindent\emph{Razlog:} Največši obči dělitelj v modulnom smyslu jest motivovany definicijeju črez sumu elementov. Slovo največši ne označaje maksimum po počtu elementov. Sravniteljne formy podpirajut obči naziv dělitelja; matematična definicija oprěděljaje, koja operacija jest nazyvana. V rabotě 22 razlikovanje modulnogo i polinomnogo smysla jest važno. Čestota istogo izraza v korpusu ne dokazuje, že vsaky jeho pojav imaje istu priměnu.\par
\noindent\emph{Reason:} Največši obči dělitelj in the module sense is motivated by a definition through sums of elements. Greatest must therefore not be read as a maximum of cardinalities. The comparanda support the common divisor terminology; the mathematical definition determines which operation it names. Paper 22 makes the distinction between module and polynomial senses important. The phrase's frequency in the corpus does not establish that every occurrence has the same application.\par
\noindent{\footnotesize Izvorne zapisy / Original records: HIST-g086-007, HIST-g092-017.}\par
\end{minipage}\par
\par\medskip\Needspace{7\baselineskip}
\noindent\begin{minipage}{\linewidth}
\hypertarget{TERM-0769}{}
\noindent\textbf{največši obči dělitelj / prěsěk}\quad{\footnotesize\texttt{TERM-0769}}\par
\noindent\emph{Iztočnik / Source:} größter gemeinsamer Teiler / Durchschnitt\par
\noindent\emph{Rusky / Russian:} наибольший общий делитель / пересечение\par
\noindent\emph{Ukrajinsky / Ukrainian:} найбільший спільний дільник / перетин\par
\smallskip\noindent\emph{Razlog:} Največši obči dělitelj i prěsěk ostavajut oboje, zato že avtorova děliteljska terminologija i opis operacije sut važni dlja dokaza. V četvrtom paragrafu raboty 9 familija oblastij $\mathfrak M(H)$ jest zatvorjena po odnošenju k prěsěkam; prěsěk vseh členov sam pripadaje toj familiji. To ne jest pravilo, že svaki obči dělitelj v vseh rabotah jest prěsěk. Pri modulah sravniteljne zapisi govoręt o sumě, a dalši paragraf raboty 9 trěbuje osobno razvažanje za drugu familiju oblastij.\par
\noindent\emph{Reason:} Največši obči dělitelj and prěsěk are both retained because the author's divisibility terminology and the description of the operation each matter to the proof. In Section 4 of Paper 9, the family of domains $\mathfrak M(H)$ is closed under intersections; the intersection of all its members is itself a member. This is not a rule that every common divisor throughout the corpus is an intersection. The module records refer to sums, while the next section of Paper 9 requires a separate argument for a different family of domains.\par
\noindent{\footnotesize Izvorne zapisy / Original records: HIST-g062-006.}\par
\smallskip\noindent\emph{Sravnjenje s tekstom:} V zaključku četvrtogo paragrafa raboty 9 prěsěk vseh oblastij iz $\mathfrak M(H)$ jest sam oblastju danogo tipa. Zatvorenost familije po odnošenju k prěsěkam jest izrična čest tvrdženja. Ona ne perenosi se avtomatično na druge familije ili na modulnu operaciju, ktora drugde nosi naziv občego dělitelja.\par
\noindent\emph{Textual comparison:} At the end of Section 4 of Paper 9, the intersection of all domains in $\mathfrak M(H)$ is itself a domain of the specified type. Closure under intersections is an explicit part of the assertion. It does not transfer automatically to other families or to the module operation elsewhere called a common divisor.\par
\noindent{\footnotesize TEXT-COMP-72; pasaže / passages: 7560, 7561, 7562; \texttt{CORPUS\_FORM\_COMPARISONS.json}.}\par
\end{minipage}\par
\par\medskip\Needspace{7\baselineskip}
\noindent\begin{minipage}{\linewidth}
\hypertarget{TERM-0770}{}
\noindent\textbf{največši obči dělitelj / najmenše obče mnogokratno}\quad{\footnotesize\texttt{TERM-0770}}\par
\noindent\emph{Iztočnik / Source:} größter gemeinsamer Teiler / kleinstes gemeinsames Vielfaches\par
\noindent\emph{Rusky / Russian:} наибольший общий делитель / наименьшее общее кратное\par
\noindent\emph{Ukrajinsky / Ukrainian:} найбільший спільний дільник / найменше спільне кратне\par
\smallskip\noindent\emph{Razlog:} Par dělitelj i mnogokratno zadržuje dva vzajemno svęzane pojmy avtorovoj teorije idealov. Razlog iz raboty 30 daje prednost tomu paru prěd potihom zaměnoju nazvami suma i prěsěk. Objasnjenje može pokazati odgovarjajuče operacije, ale ne smije skryti poredok děliteljnosti, ktory upravjaje argumentom. Največši i najmenši dolžno čitati v tom poredku. Jezyčna sličnost s obyčnym čislovym največšim občim děliteljem ne jest dostatočna definicija idealnogo slučaja.\par
\noindent\emph{Reason:} The pair divisor and multiple retains two related concepts in the author's ideal theory. The rationale for Paper 30 prefers this pair to silently substituting sum and intersection. An explanation may identify the corresponding operations, but it must not conceal the divisibility order governing the argument. Greatest and least are read in that order. Linguistic resemblance to the ordinary numerical greatest common divisor does not suffice to define the ideal case.\par
\noindent{\footnotesize Izvorne zapisy / Original records: HIST-g111-004.}\par
\end{minipage}\par
\par\medskip\Needspace{7\baselineskip}
\noindent\begin{minipage}{\linewidth}
\hypertarget{TERM-0771}{}
\noindent\textbf{največši obči dělitelj jest jednaky jediničnomu idealu $\mathfrak{D}=(\varepsilon)$}\quad{\footnotesize\texttt{TERM-0771}}\par
\noindent\emph{Iztočnik / Source:} größter gemeinsamer Teiler gleich dem Einheitsideal $\mathfrak{D}=(\varepsilon)$\par
\noindent\emph{Rusky / Russian:} наибольший общий делитель равен единичному идеалу $\mathfrak{D}=(\varepsilon)$\par
\noindent\emph{Ukrajinsky / Ukrainian:} найбільший спільний дільник дорівнює одиничному ідеалу $\mathfrak{D}=(\varepsilon)$\par
\smallskip\noindent\emph{Razlog:} Cěla fraza vyražaje definiciju vzajemno prostyh idealov, ne samo naziv. V osmom paragrafu raboty 19 ih najveći obči dělitelj mora byti jednaky jediničnomu idealu $\mathfrak D=(\varepsilon)$, ktory sodrživaje vse elementy podloženoj oblasti. Razlog sohranjaje avtorov izraz vměsto kratkogo sovrěmennogo nazva komaksimalny. Poslědujuči tekst razlikuje tu definiciju od vzajemnoj relativnoj prostosti i daje priměr, kde obratno tvrdženje ne imaje město. Zato toj dugi opis ne možno bez utraty informaciji skratiti do proizvoljnogo značenja slova prosty.\par
\noindent\emph{Reason:} The whole phrase states a definition of coprime ideals, not merely a name. In Section 8 of Paper 19 their greatest common divisor must equal the unit ideal $\mathfrak D=(\varepsilon)$, which contains every element of the underlying domain. The rationale preserves the author's expression instead of the compact modern term comaximal. The subsequent text distinguishes this definition from mutual relative primality and gives a counterexample to the converse. The longer description therefore cannot be abbreviated to an unspecified meaning of prime without losing information.\par
\noindent{\footnotesize Izvorne zapisy / Original records: HIST-g081-065, HIST-g082-065.}\par
\smallskip\noindent\emph{Sravnjenje s tekstom:} Definicija VIII raboty 19 oprěděljaje vzajemno proste idealy črez rovnost ih občego dělitelja s $\mathfrak D=(\varepsilon)$. Slědujuča pasaža razlikuje tu vlastnost od vzajemnoj relativnoj prostosti i daje $(x)$, $(y)$ kako priměr nevaženja obratnogo tvrdženja. Zato celosna definicija daje bolše informaciji než izolovany naziv. Lokalna makrodefinicija potvrđuje frakturny zapis $\mathfrak D$.\par
\noindent\emph{Textual comparison:} Definition VIII of Paper 19 defines coprime ideals through equality of their common divisor with $\mathfrak D=(\varepsilon)$. The following passage distinguishes this property from mutual relative primality, using $(x)$ and $(y)$ to show failure of the converse. The full definition therefore carries more information than an isolated label. The local macro definition confirms the fraktur $\mathfrak D$.\par
\noindent{\footnotesize TEXT-COMP-73; pasaže / passages: 13269, 13276, 13278, 13280; \texttt{CORPUS\_FORM\_COMPARISONS.json}.}\par
\end{minipage}\par
\par\medskip\Needspace{7\baselineskip}
\noindent\begin{minipage}{\linewidth}
\hypertarget{TERM-0772}{}
\noindent\textbf{največši obči dělitelj v polinomnom smyslu}\quad{\footnotesize\texttt{TERM-0772}}\par
\noindent\emph{Iztočnik / Source:} größter gemeinsamer Teiler im Polynomsinn\par
\noindent\emph{Rusky / Russian:} наибольший общий делитель в полиномиальном смысле\par
\noindent\emph{Ukrajinsky / Ukrainian:} найбільший спільний дільник у поліномному сенсі; найбільший спільний дільник у поліноміальному сенсі\par
\smallskip\noindent\emph{Razlog:} Dodatek v polinomnom smyslu ograničaje značenje občego dělitelja. Razlogi iz raboty 22 i absolutnoj prostosti v rabotě 24 nastajut na razlikovanju děliteljnosti polinomov od děliteljnosti idealov. Izpuščenje dodatka moglo by sugerovati drugu operaciju, iako korenj děl- ostaje isti. Sravnjenje sohranjaje takože historijne pravopisne varianty superlativa; njihova unifikacija ne smije zaměniti matematično razlikovanje. Izbor jest obrazložen smyslom i kontinuitetom nazva, ne novym čislenym rangiranjem slov.\par
\noindent\emph{Reason:} The qualification v polinomnom smyslu restricts the sense of common divisor. The rationales from Paper 22 and the discussion of absolute primality in Paper 24 insist on distinguishing polynomial divisibility from ideal divisibility. Omitting the qualification could suggest a different operation even though the divisor stem remains unchanged. The comparison also retains historical spelling variants of the superlative; standardizing them cannot replace the mathematical distinction. The choice is justified by sense and terminological continuity, not by a new numerical ranking of words.\par
\noindent{\footnotesize Izvorne zapisy / Original records: HIST-g089-006, HIST-g090-009, HIST-g092-048, HIST-g092-063, HIST-g101-008.}\par
\end{minipage}\par
\par\medskip\Needspace{7\baselineskip}
\noindent\begin{minipage}{\linewidth}
\hypertarget{TERM-0773}{}
\noindent\textbf{najveći obšty dělitelj idealov}\quad{\footnotesize\texttt{TERM-0773}}\par
\noindent\emph{Iztočnik / Source:} größter gemeinsamer Teiler von Idealen\par
\noindent\emph{Rusky / Russian:} наибольший общий делитель идеалов\par
\noindent\emph{Ukrajinsky / Ukrainian:} найбільший спільний дільник ідеалів\par
\smallskip\noindent\emph{Razlog:} Najveći obšty dělitelj idealov jest historijna forma dlja operacije, ktora v danoj teoriji odpovědaje sumě idealov. Sohranjena motivacija zadržuje i děliteljsky naziv, i izričnu definiciju. Tablica pokazuje obšty, dokolje druge zapisi upotrěbjajut obči. Ta variacija sama po sobě ne dokazuje, že oba pridavnika byli osobno sravnjeni i odobreni. Matematična vloga frazy jest objasnjena; polna jezyčna ocěna različnyh form trěbuje dopolniteljno svědočstvo.\par
\noindent\emph{Reason:} Najveći obšty dělitelj idealov is a historical form for the operation that corresponds to the sum of ideals in this theory. The preserved motivation retains both the divisor name and its explicit definition. The table records obšty, whereas other records use obči. That variation does not itself show that the two adjectives were separately compared and endorsed. The phrase's mathematical role is explained; a full linguistic assessment of the variants requires additional evidence.\par
\noindent{\footnotesize Izvorne zapisy / Original records: HIST-g081-025, HIST-g082-025.}\par
\end{minipage}\par
\par\medskip\Needspace{7\baselineskip}
\noindent\begin{minipage}{\linewidth}
\hypertarget{TERM-0774}{}
\noindent\textbf{najveći primarny faktor normy}\quad{\footnotesize\texttt{TERM-0774}}\par
\noindent\emph{Iztočnik / Source:} größter primärer Faktor der Norm\par
\noindent\emph{Rusky / Russian:} наибольший первичный множитель нормы\par
\noindent\emph{Ukrajinsky / Ukrainian:} найбільший первинний множник норми\par
\smallskip\noindent\emph{Razlog:} Najveći primarny faktor normy sohranjaje faktor tam, kde predmetom jest norma, a ne prosto razloženje ideala. Razlog svęzuje frazu s Teoremami X i XI raboty 24 i vybira medžunarodny korenj faktor- radi očekivanoj razpoznatelnosti. To očekivanje ne jest rezultat čitatelskogo eksperimenta. Naziv trěbuje dodatka normy: bez njega by tablica mogla skryti razliku od primarnoj komponenty. Primarny i najveći nadalej noset svoj specifičny algebraičny smysl.\par
\noindent\emph{Reason:} Najveći primarny faktor normy retains factor where the object is a norm, rather than simply an ideal decomposition. The rationale connects the phrase to Theorems X and XI of Paper 24 and selects the international factor stem for anticipated recognizability. That expectation is not the result of a reader experiment. The qualification of the norm matters: without it, the table could obscure the distinction from a primary component. Primary and greatest continue to carry their specific algebraic meanings.\par
\noindent{\footnotesize Izvorne zapisy / Original records: HIST-g099-007.}\par
\end{minipage}\par
\par\medskip\Needspace{7\baselineskip}
\noindent\begin{minipage}{\linewidth}
\hypertarget{TERM-0775}{}
\noindent\textbf{najveći primarny dělitelj / primarny dělitelj}\quad{\footnotesize\texttt{TERM-0775}}\par
\noindent\emph{Iztočnik / Source:} größter primärer Teiler / primärer Teiler\par
\noindent\emph{Rusky / Russian:} наибольший примарный делитель / примарный делитель\par
\noindent\emph{Ukrajinsky / Ukrainian:} найбільший примарний дільник / примарний дільник\par
\smallskip\noindent\emph{Razlog:} Primarny dělitelj opisuje komponentu s pomočju avtorovoj děliteljskoj relacije. Dodatek najveći zadržuje osobnu kvalifikaciju, kogda ju tekst trěbuje. Razlog ne svede oba izrazy do obyčnogo podmnožestva: idealna děliteljnost i uključanje množestv ne možut byti zaměnene bez objasnjenja poredka. Obči korenj děl- svezuje frazu s ostaloju terminologijeju raboty 19. Sohranjena motivacija podpiraje tu sistematičnost, ale ne dava čislenyh težin za vsaku možnu slovnu alternativu.\par
\noindent\emph{Reason:} Primarny dělitelj describes a component through the author's divisibility relation. Greatest retains a further qualification wherever the text requires it. The rationale does not reduce either expression to ordinary subset language: ideal divisibility and set inclusion cannot be interchanged without explaining the order. The divisor stem connects the phrase to the rest of Paper 19's terminology. The preserved record supports this consistency but supplies no numerical weights for every possible lexical alternative.\par
\noindent{\footnotesize Izvorne zapisy / Original records: HIST-g081-050, HIST-g082-050.}\par
\end{minipage}\par
\par\medskip\Needspace{7\baselineskip}
\noindent\begin{minipage}{\linewidth}
\hypertarget{TERM-0776}{}
\noindent\textbf{največši popolno reducibilny modul}\quad{\footnotesize\texttt{TERM-0776}}\par
\noindent\emph{Iztočnik / Source:} größter vollständig reduzibler Modul\par
\noindent\emph{Rusky / Russian:} наибольший вполне приводимый модуль\par
\noindent\emph{Ukrajinsky / Ukrainian:} найбільший повністю звідний модуль\par
\smallskip\noindent\emph{Razlog:} Največši popolno reducibilny modul označaje oprěděljenu konstrukciju, ne proizvoljny maksimalny podmodul. V rabotě 17 najprvo razsmatrjaje se najmenše obče mnogokratno vseh prostyh děliteljev danogo modula. Naziv največši polno reducibilny modul priměnjaje se, kogda ta konstrukcija može byti prědstavjena črez konečno mnogo prostyh modulov. Tekst izrično razsmatrjaje takože slučaj, kde takogo modula nema. Zato ne možno prědložiti bezuslovnu egzistenciju ni potihom zaměniti največši s maksimalny. Popolno v staršoj zapisi i polno v pasaži ne měnjajut togo uslovja.\par
\noindent\emph{Reason:} Največši popolno reducibilny modul denotes a specified construction, not an arbitrary maximal submodule. Paper 17 first considers the least common multiple of all simple divisors of the given module. The name greatest completely reducible module applies when that construction can be represented using finitely many simple modules. The text also explicitly considers the case in which no such module exists. The wording must therefore neither imply unconditional existence nor silently substitute maximal for greatest. Historical popolno and the passage's polno do not alter that condition.\par
\noindent{\footnotesize Izvorne zapisy / Original records: HIST-g079-059.}\par
\smallskip\noindent\emph{Sravnjenje s tekstom:} Rabota 17 najprvo definira polnu reducibilnost črez konečnu sumu prostyh grup i odgovarjajuči modul. Pozdnějšemu nazvu največši polno reducibilny modul odpovědaje prědstavimost najmenšego občego mnogokratnogo vseh prostyh děliteljev črez konečno mnogo prostyh modulov. Teorem X izrično dopušča, že taki modul ne obstavaje. To jest sučstveno uslovje nazva, ne fakultativno objasnjenje.\par
\noindent\emph{Textual comparison:} Paper 17 first defines complete reducibility using a finite sum of simple groups and the corresponding module. The later name greatest completely reducible module requires the least common multiple of all simple divisors to be representable using finitely many simple modules. Theorem X explicitly allows that such a module may not exist. This is an essential condition of the name, not an optional explanatory gloss.\par
\noindent{\footnotesize TEXT-COMP-74; pasaže / passages: 11576, 11949, 11951; \texttt{CORPUS\_FORM\_COMPARISONS.json}.}\par
\end{minipage}\par
\par\medskip\Needspace{7\baselineskip}
\noindent\begin{minipage}{\linewidth}
\hypertarget{TERM-0777}{}
\noindent\textbf{najveće možno invariantno ograničenje}\quad{\footnotesize\texttt{TERM-0777}}\par
\noindent\emph{Iztočnik / Source:} größtmögliche invariante Einschränkung\par
\noindent\emph{Rusky / Russian:} наибольшее возможное инвариантное ограничение\par
\noindent\emph{Ukrajinsky / Ukrainian:} найбільше можливе інваріантне обмеження\par
\smallskip\noindent\emph{Razlog:} Najveće možno invariantno ograničenje vyražaje kvalifikovanu krajnju možnost, na koju se poziva rabota 4 pri Teoremě VI. Vsaka čest frazy imaje funkciju: ograničenje jest predmet, invariantno oprěděljaje njegov vid, a najveće možno zadržuje silu tvrdženja. Kratky historijny razlog ne opisuva dodatnu čislenu optimizaciju ni novy formalny kriterij. Zato objasnjenje čuva frazu, ne dopisuje nezasvědočeny algoritm dlja njezinogo izbora.\par
\noindent\emph{Reason:} Najveće možno invariantno ograničenje expresses the qualified extremal possibility invoked in Paper 4 in connection with Theorem VI. Each part has a function: restriction names the object, invariant specifies its kind, and greatest possible preserves the strength of the assertion. The brief historical rationale describes neither an additional numerical optimization nor a new formal criterion. The explanation therefore retains the phrase without supplying an unattested algorithm for its selection.\par
\noindent{\footnotesize Izvorne zapisy / Original records: HIST-g031-015.}\par
\end{minipage}\par
\par\medskip\Needspace{7\baselineskip}
\noindent\begin{minipage}{\linewidth}
\hypertarget{TERM-0779}{}
\noindent\textbf{poluprosto kolco}\quad{\footnotesize\texttt{TERM-0779}}\par
\noindent\emph{Iztočnik / Source:} Halbeinfacher Ring\par
\noindent\emph{Rusky / Russian:} полупростое кольцо\par
\noindent\emph{Ukrajinsky / Ukrainian:} напівпросте кільце\par
\smallskip\noindent\emph{Razlog:} Poluprosto kolco jest sohranjena fraza dlja nazva, ktory rabota 34 dava v skobkah pri opisu prstena bez radikala. V toj pasaži radikal oprěděljen jest kak maksimalny nilpotentny ideal pri uslovju maksimalnosti za prave idealy. Pozdněj tekst svęzuje odsutstvo radikala s polnoju reducibilnostju pri dopolniteljnom uslovju minimalnosti. Proto parentetičny synonym ne dava pravo zanedbati te uslovja ili proglasiti vsaky prsten bez radikala poluprostym v vsakom sovrěmennom smyslu. Pojmovna vloga jasna jest iz definicij; sohranjena forma kolco ne potvrđuje gramatiku vseh pozdnějših fraz s prsten.\par
\noindent\emph{Reason:} Poluprosto kolco is the preserved phrase for the parenthetical name accompanying ring without radical in Paper 34. There the radical is defined as the maximal nilpotent ideal under the maximum condition for right ideals. The text later connects absence of the radical with complete reducibility under an additional minimum condition. The parenthetical synonym therefore does not license omitting those conditions or declaring every radical-free ring semisimple in every modern sense. The definitions clarify the conceptual role; the recorded kolco form does not certify the grammar of every later phrase using prsten.\par
\noindent{\footnotesize Izvorne zapisy / Original records: HIST-g190-004.}\par
\smallskip\noindent\emph{Sravnjenje s tekstom:} V rabotě 34 radikal jest maksimalny nilpotentny ideal pri uslovju maksimalnosti za prave idealy. Poluprosty jest parentetičny naziv prstena s nulovym radikalom. Povezanje s polnoju reducibilnostju trěbuje dalših hipotez, osoblivo uslovja minimalnosti v obratnom směru. Sravnjenje ne dava osnovy perenesti tu lokalnu terminologiju na vse sovrěmenne teorije radikala bez proverky definicij.\par
\noindent\emph{Textual comparison:} In Paper 34 the radical is the maximal nilpotent ideal under the maximum condition for right ideals. Semisimple is the parenthetical name for a ring with zero radical. The connection with complete reducibility requires further hypotheses, notably the minimum condition in the converse direction. This comparison does not justify transferring the local terminology to every modern radical theory without checking definitions.\par
\noindent{\footnotesize TEXT-COMP-75; pasaže / passages: 20860, 20864, 20916, 20924; \texttt{CORPUS\_FORM\_COMPARISONS.json}.}\par
\end{minipage}\par
\par\medskip\Needspace{7\baselineskip}
\noindent\begin{minipage}{\linewidth}
\hypertarget{TERM-0780}{}
\noindent\textbf{poluprosty sistem}\quad{\footnotesize\texttt{TERM-0780}}\par
\noindent\emph{Iztočnik / Source:} halbeinfaches System\par
\noindent\emph{Rusky / Russian:} полупростая система\par
\noindent\emph{Ukrajinsky / Ukrainian:} напівпроста система\par
\smallskip\noindent\emph{Razlog:} Poluprosty sistem zadržuje algebraičny pridavnik i obči naziv objekta. Historijna motivacija ograničaje se na tvrdženje o běžnom algebraičnom registru, uz rusku i ukrajinsku sravniteljnicu. To ne jest polna dokumentacija oficialnogo medžuslovjanskogo termina. Konkretna pasaža raboty 37 dava bolše matematičnogo objasnjenja: po razširjenju koeficientov nastavaje suma izomorfnyh polj, zato poluprosty sistem. Ta konstrukcija, ne opče tvrđenje o svakom sistemu, podpiraje smysl pridavnika. Jezyčna normativnost i matematična priměnimost ostavajut različne pytanja.\par
\noindent\emph{Reason:} Poluprosty sistem retains the algebraic adjective and the general name of the object. The historical motivation consists only of an assertion about standard algebraic register, accompanied by Russian and Ukrainian comparanda. That is not full documentation of an official Interslavic term. The particular passage in Paper 37 provides a more informative mathematical explanation: extending coefficients yields a sum of isomorphic fields, hence a semisimple system. This construction, not a claim about every system, supports the adjective's sense. Linguistic normativity and mathematical applicability remain different questions.\par
\noindent{\footnotesize Izvorne zapisy / Original records: HIST-g207-011.}\par
\smallskip\noindent\emph{Sravnjenje s tekstom:} Pasaža raboty 37 opisuva konkretno razširjenje koeficientov: iz njega nastavaje suma izomorfnyh polj, i tekst zato govori o poluprostom sistemu. To podpiraje matematičny smysl pridavnika v toj pasaži, ne tvrdženje, že ista fraza jest uznata jezyčna norma za vsaky kontekst.\par
\noindent\emph{Textual comparison:} The passage in Paper 37 describes a particular extension of coefficients: it yields a sum of isomorphic fields, and the text therefore calls it a semisimple system. This supports the adjective's mathematical sense in that passage, not a claim that the phrase is an established linguistic norm for every context.\par
\noindent{\footnotesize TEXT-COMP-76; pasaže / passages: 22846; \texttt{CORPUS\_FORM\_COMPARISONS.json}.}\par
\end{minipage}\par
\par\medskip\Needspace{7\baselineskip}
\noindent\begin{minipage}{\linewidth}
\hypertarget{TERM-0781}{}
\noindent\textbf{teorem obstajanja Hassea}\quad{\footnotesize\texttt{TERM-0781}}\par
\noindent\emph{Iztočnik / Source:} Hasse existence theorem\par
\noindent\emph{Rusky / Russian:} теорема существования Хассе\par
\noindent\emph{Ukrajinsky / Ukrainian:} теорема існування Гассе\par
\smallskip\noindent\emph{Razlog:} Teorem obstajanja Hassea označaje imenovany teorem o obstajanju, ne vlastnost samogo Hassea. Razlog vybira prěmy naziv matematičnogo tvrđenja vměsto parafrazy; ruska i ukrajinska sravniteljnica davajut odgovarjajuče nazvy. Sam razlog ne obrazlagaje podrobno koncovku Hassea. Zadržanje identiteta avtorovogo imena i izbor gramatičnoj formy sut zato dva različne pytanja. Niti naziv teorema zaměnjaje njegovo cělo tvrđenje s hipotezami.\par
\noindent\emph{Reason:} Teorem obstajanja Hassea names a theorem asserting existence, not a property of Hasse himself. The rationale chooses a direct mathematical designation instead of paraphrase; the Russian and Ukrainian comparanda supply corresponding names. It does not, however, explain the ending Hassea in detail. Preserving the person's identity and choosing a grammatical form of the name are therefore separate questions. Nor does naming the theorem replace its full statement and hypotheses.\par
\noindent{\footnotesize Izvorne zapisy / Original records: HIST-g172-003.}\par
\end{minipage}\par
\par\medskip\Needspace{7\baselineskip}
\noindent\begin{minipage}{\linewidth}
\hypertarget{TERM-0782}{}
\noindent\textbf{normna teorema Hasse}\quad{\footnotesize\texttt{TERM-0782}}\par
\noindent\emph{Iztočnik / Source:} Hassescher Normensatz\par
\noindent\emph{Rusky / Russian:} норменная теорема Гассе\par
\noindent\emph{Ukrajinsky / Ukrainian:} нормова теорема Гассе\par
\smallskip\noindent\emph{Razlog:} Normna teorema Hasse sohranjaje vlastno ime i povezanje teoremy s algebraičnoju normoju. Kratka motivacija govori o prilagodženju pridavnika matematičnoj proze, ale ne dava podrobne morfologične sravnjenje. Zato tablica dokumentuje izbranu formu, ne proglasuje ju edinoju možnoju. Norma tu nosi tehničny algebraičny smysl; ne označaje društveno pravilo. Različny zapis imena pri TERM-0781 ostaje vidimy kak gramatično pytanje, ne razlika medžu dvoma ljudami.\par
\noindent\emph{Reason:} Normna teorema Hasse preserves the name and the theorem's association with an algebraic norm. The brief motivation speaks of adapting the adjective to mathematical prose but supplies no detailed morphological comparison. The table therefore documents the selected form without declaring it the only possible one. Norm has its technical algebraic sense here, not that of a social rule. The different treatment of the name at TERM-0781 remains visible as a grammatical question, not a distinction between two people.\par
\noindent{\footnotesize Izvorne zapisy / Original records: HIST-g208-012.}\par
\end{minipage}\par
\par\medskip\Needspace{7\baselineskip}
\noindent\begin{minipage}{\linewidth}
\hypertarget{TERM-0783}{}
\noindent\textbf{glavna diagonala}\quad{\footnotesize\texttt{TERM-0783}}\par
\noindent\emph{Iztočnik / Source:} Hauptdiagonale\par
\noindent\emph{Rusky / Russian:} главная диагональ\par
\noindent\emph{Ukrajinsky / Ukrainian:} головна діагональ\par
\smallskip\noindent\emph{Razlog:} Glavna diagonala jest prozračna kombinacija pridavnika i medžunarodnogo matematičnogo imenika. Ruska i ukrajinska sravniteljnica podpirajut tu strukturu. V matričnom kontekstu ide o diagonalu s jednakymi indeksami reda i stolbca, ne o proizvoljnu geometričnu diagonalu. Historijna ocěna, že termin jest nesporny, ne jest svědočstvo nezavisnoj ekspertnoj proverky. Zapis dava kratku motivaciju iz běžnogo matričnogo jezyka, ne čislene rezultate razumělivosti.\par
\noindent\emph{Reason:} Glavna diagonala combines an intelligible adjective with an international mathematical noun. The Russian and Ukrainian comparanda support that structure. In the matrix context it is the diagonal with equal row and column indices, not an arbitrary geometric diagonal. The historical assessment that the term is uncontroversial is not evidence of independent expert review. The record offers a short justification from ordinary matrix language, not numerical comprehension results.\par
\noindent{\footnotesize Izvorne zapisy / Original records: HIST-g196-010.}\par
\end{minipage}\par
\par\medskip\Needspace{7\baselineskip}
\noindent\begin{minipage}{\linewidth}
\hypertarget{TERM-0784}{}
\noindent\textbf{glavna jedinica}\quad{\footnotesize\texttt{TERM-0784}}\par
\noindent\emph{Iztočnik / Source:} Haupteinheit\par
\noindent\emph{Rusky / Russian:} главная единица\par
\noindent\emph{Ukrajinsky / Ukrainian:} головна одиниця\par
\smallskip\noindent\emph{Razlog:} Glavna jedinica sohranjaje historijno imenovanje iz teorije hiperkompleksnyh čisel. Razlog izrično ne hoče svesti to imenovanje do pozdnějšego občego nazva jediničny element. To ne dokazuje, že v vsakoj pasaži ide o različne matematične objekty; ono zadržuje avtorov izbor jezyka. Glavna zato ne jest slobodna pohvala objekta, ale čest sohranjenogo tehničnogo nazva. Identitet funkciji v konkretnoj konstrukciji mora byti ustanovjen iz definicije.\par
\noindent\emph{Reason:} Glavna jedinica preserves historical terminology from the theory of hypercomplex numbers. The rationale expressly avoids reducing that wording to the later general name identity element. This does not establish that every passage concerns different mathematical objects; it preserves the author's choice of vocabulary. Principal is therefore not a free expression of importance, but part of the retained technical name. Identity of function in a particular construction must be established from the definition.\par
\noindent{\footnotesize Izvorne zapisy / Original records: HIST-g118-014.}\par
\end{minipage}\par
\par\medskip\Needspace{7\baselineskip}
\noindent\begin{minipage}{\linewidth}
\hypertarget{TERM-0785}{}
\noindent\textbf{teorem o glavnom rodu}\quad{\footnotesize\texttt{TERM-0785}}\par
\noindent\emph{Iztočnik / Source:} Hauptgeschlechtssatz\par
\noindent\emph{Rusky / Russian:} [Bukvy ne sut zachovane / Lettering not recoverable]\par
\noindent\emph{Ukrajinsky / Ukrainian:} [Bukvy ne sut zachovane / Lettering not recoverable]\par
\smallskip\noindent\emph{Razlog:} Teorem o glavnom rodu upotrěbjaje rod v aritmetičnom smyslu teorije klasnyh polj. Oba razloga odbivajut doslovny prěvod, ktory by sugeroval pol ili obyčnu vrstu naměsto tehničnogo roda. Kontinuitet medžu rabotami 39 i 41 jest vtorym razlogom za isti naziv. Ruske i ukrajinske bukvy v zapisi ne sut sohranjene; rod ne može služiti za nagađanje izgubljenyh form. Razpoznatelnost občego slovjanskogo korenja jest obrazloženje izbora, ne izměrena čitatelska ocěna.\par
\noindent\emph{Reason:} Teorem o glavnom rodu uses genus in the arithmetic sense of class-field theory. Both rationales reject a literal rendering that would suggest sex or an ordinary kind instead of the technical genus. Continuity between Papers 39 and 41 is a second reason for retaining the name. The Russian and Ukrainian lettering has not survived in this record; rod cannot be used to guess the missing forms. Recognizability of the shared Slavic root is a reason for the choice, not a measured reader assessment.\par
\noindent{\footnotesize Izvorne zapisy / Original records: HIST-g209-002, HIST-g211-001.}\par
\end{minipage}\par
\par\medskip\Needspace{7\baselineskip}
\noindent\begin{minipage}{\linewidth}
\hypertarget{TERM-0786}{}
\noindent\textbf{glavny ideal}\quad{\footnotesize\texttt{TERM-0786}}\par
\noindent\emph{Iztočnik / Source:} Hauptideal\par
\noindent\emph{Rusky / Russian:} главный идеал\par
\noindent\emph{Ukrajinsky / Ukrainian:} головний ідеал\par
\smallskip\noindent\emph{Razlog:} Glavny ideal jest komponovany tehničny naziv, ne ocěna važnosti ideala. Razlogi svęzujut jego s jednim porodžujučim elementom, polinomnoju bazu, funkcionalnoju oblastju i idealom porodženym prostym čislom. Osoblivo važen jest priměr parnyh cělyh čisel v rabotě 19: vsaky ideal jest glavny, ale oblast ne ima jedinice i ne dava jednoznačnogo produktnogo prědstavjenja čisel. Zato naziv ne dopušča potihom dodati hipotezy sovrěmennogo teorema o faktorizaciji. Obči termin mora čuvati različne uslovja svih teh priměnjenj.\par
\noindent\emph{Reason:} Glavny ideal is a compositional technical name, not an assessment of an ideal's importance. The rationales connect it with a single generator, a polynomial basis, a functional domain, and an ideal generated by a prime number. The even-integer example in Paper 19 is particularly important: every ideal is principal, yet the domain lacks an identity and does not give unique product representations of its numbers. The name therefore does not license silently adding the hypotheses of a modern factorization theorem. A common term must preserve the different conditions of these applications.\par
\noindent{\footnotesize Izvorne zapisy / Original records: HIST-g081-089, HIST-g085-009, HIST-g092-009, HIST-g110-010, HIST-g118-003.}\par
\smallskip\noindent\emph{Sravnjenje s tekstom:} Rabota 19 dava oblast parnyh cělyh čisel: vsaky ideal jest glavny, ale jedinica ne obstaja. Tekst razděljuje idealno svojstvo od jednoznačnogo produktnogo prědstavjenja čisel. Pri prěvodu ne možno potihom vključiti hipotezu jedinice samo zato, že obče slovo glavny pripomina sovrěmenny specialny naziv.\par
\noindent\emph{Textual comparison:} Paper 19 uses the domain of even integers: every ideal is principal, but there is no identity. The text separates the ideal property from unique product representation of numbers. A translation must not silently insert an identity hypothesis merely because principal resembles a modern specialized designation.\par
\noindent{\footnotesize TEXT-COMP-77; pasaže / passages: 13643, 13658; \texttt{CORPUS\_FORM\_COMPARISONS.json}.}\par
\end{minipage}\par
\par\medskip\Needspace{7\baselineskip}
\noindent\begin{minipage}{\linewidth}
\hypertarget{TERM-0787}{}
\noindent\textbf{glavny ideal / bazny polinom}\quad{\footnotesize\texttt{TERM-0787}}\par
\noindent\emph{Iztočnik / Source:} Hauptideal / Basispolynom\par
\noindent\emph{Rusky / Russian:} главный идеал / базисный полином\par
\noindent\emph{Ukrajinsky / Ukrainian:} головний ідеал / базисний поліном\par
\smallskip\noindent\emph{Razlog:} Glavny ideal i bazny polinom opisuju svęzanu dvojicu: ideal i polinom, ktory daje jeho bazu v argumentu o jednoj prěměnnoj. Razlog pri Teoremě XI raboty 22 zadržuje te dve funkcije. Bazny ne znamenaje, že polinom jest jediny bez dodatnoj normalizacije, ni že vsak mnogoprěměnny ideal imaje jednogo porodžujučego elementa. Prědnost prostoj konstrukciji iz poznanyh korenjev jest kvalitativna motivacija; zaměna svakego slova izolovano ne by sohranila vezu medžu nimi.\par
\noindent\emph{Reason:} Glavny ideal and bazny polinom describe a related pair: the ideal and the polynomial supplying its basis in a one-variable argument. The rationale for Theorem XI of Paper 22 retains these two roles. Basic does not mean that the polynomial is unique without further normalization, or that every multivariable ideal has a single generator. Preference for a straightforward construction from familiar stems is a qualitative motivation; translating each word in isolation would not preserve their relationship.\par
\noindent{\footnotesize Izvorne zapisy / Original records: HIST-g090-004, HIST-g092-058.}\par
\end{minipage}\par
\par\medskip\Needspace{7\baselineskip}
\noindent\begin{minipage}{\linewidth}
\hypertarget{TERM-0788}{}
\noindent\textbf{glavny ideal $\mathfrak{T}p$ / po parah vzajemno proste primarne komponenty}\quad{\footnotesize\texttt{TERM-0788}}\par
\noindent\emph{Iztočnik / Source:} Hauptideal $\mathfrak{T}p$ / paarweise teilerfremde Primärkomponenten\par
\noindent\emph{Rusky / Russian:} главный идеал $\mathfrak{T}p$ / попарно взаимно простые первичные компоненты\par
\noindent\emph{Ukrajinsky / Ukrainian:} головний ідеал $\mathfrak{T}p$ / попарно взаємно прості первинні компоненти\par
\smallskip\noindent\emph{Razlog:} Glavny ideal $\mathfrak T p$ i po parah vzajemno proste primarne komponenty sohranjajut vse sučstvene česti frazy. Po parah zahtěvaje uslovje za vsaku dvojicu različnyh komponent, ne samo za vybranu dvojicu. Primarny ne možno zaměniti s prosty: teorem o diskriminantu v rabotě 31 izrično razlikuje vlastno primarnu komponentu od prostogo ideala drugogo roda. Zapis prodolžuje ustaljenu frazu o vzajemnoj prostosti; njegov kratky razlog ne jest nova ocěna vseh jezyčnyh alternativ.\par
\noindent\emph{Reason:} Glavny ideal $\mathfrak T p$ and po parah vzajemno proste primarne komponenty retain all the phrase's essential parts. Pairwise requires the condition for every pair of distinct components, not only a selected pair. Primary cannot be replaced by prime: Paper 31's discriminant theorem expressly distinguishes a properly primary component from a prime ideal of the second kind. The record continues the established coprimeness phrase; its short rationale is not a new assessment of every linguistic alternative.\par
\noindent{\footnotesize Izvorne zapisy / Original records: HIST-g153-003.}\par
\smallskip\noindent\emph{Sravnjenje s tekstom:} Diskriminantny teorem raboty 31 govori o razloženju glavnogo ideala $\mathfrak T p$ na po parah vzajemno proste primarne komponenty. On potom razlikuje vlastno primarnu komponentu i prosty ideal drugogo roda. Te odvojene kvalifikacije sut nužne dlja točnogo tvrđenja; lokalna makrodefinicija potvrđuje frakturnu bukvu $\mathfrak T$.\par
\noindent\emph{Textual comparison:} Paper 31's discriminant theorem concerns decomposition of the principal ideal $\mathfrak T p$ into pairwise coprime primary components. It then distinguishes a properly primary component from a prime ideal of the second kind. Those separate qualifications are necessary to the exact statement; the local macro definition confirms fraktur $\mathfrak T$.\par
\noindent{\footnotesize TEXT-COMP-78; pasaže / passages: 18695, 18715; \texttt{CORPUS\_FORM\_COMPARISONS.json}.}\par
\end{minipage}\par
\par\medskip\Needspace{7\baselineskip}
\noindent\begin{minipage}{\linewidth}
\hypertarget{TERM-0789}{}
\noindent\textbf{kolco glavnyh idealov}\quad{\footnotesize\texttt{TERM-0789}}\par
\noindent\emph{Iztočnik / Source:} Hauptidealring\par
\noindent\emph{Rusky / Russian:} кольцо главных идеалов\par
\noindent\emph{Ukrajinsky / Ukrainian:} кільце головних ідеалів\par
\smallskip\noindent\emph{Razlog:} Kolco glavnyh idealov razkriva svojstvo, že vsaky ideal jest glavny. Tri razloga podpirajut kontinuitet nazva glavny ideal i polnu frazu vměsto neobjasnjenoj kratice. V rabotě 30 izraz priměnjaje se takože k prstenam klasov ostatkov. Zato ne možno avtomatično prěvoditi vsaky takovy naziv kak oblast bez děliteljev nuly: ta hipoteza jest osobna. Starša forma kolco ostaje v sravnjenju kak dokumentovana izbira, ne kak potvrda vseh njezinyh aktualnyh gramatičnyh priměnjenj.\par
\noindent\emph{Reason:} Kolco glavnyh idealov spells out the property that every ideal is principal. Three rationales support continuity with glavny ideal and a full phrase instead of an unexplained abbreviation. Paper 30 also applies the expression to residue-class rings. It must therefore not automatically become a domain without zero divisors: that is a separate hypothesis. The earlier kolco form remains in the comparison as a documented choice, not as certification of all its current grammatical applications.\par
\noindent{\footnotesize Izvorne zapisy / Original records: HIST-g116-003, HIST-g116-009, HIST-g160-004.}\par
\end{minipage}\par
\par\medskip\Needspace{7\baselineskip}
\noindent\begin{minipage}{\linewidth}
\hypertarget{TERM-0790}{}
\noindent\textbf{kolco glavnyh idealov / izhodno kolco}\quad{\footnotesize\texttt{TERM-0790}}\par
\noindent\emph{Iztočnik / Source:} Hauptidealring / Ausgangsring\par
\noindent\emph{Rusky / Russian:} кольцо главных идеалов / исходное кольцо\par
\noindent\emph{Ukrajinsky / Ukrainian:} кільце головних ідеалів / вихідне кільце\par
\smallskip\noindent\emph{Razlog:} Kontrast medžu kolco glavnyh idealov i izhodno kolco jest sučstven dlja razširjenja argumenta v osmom paragrafu raboty 31. Pervy naziv dava algebraično svojstvo; vtory dava rolju prstena v konstrukciji. Novy početkovy prsten ne mora byti istogo specialnogo tipa kak ranějši. Razlog zato zadržuje polny naziv glavnyh idealov i ne vvodit kratice, ktora by skryla tę razliku. Aktualna pasaža izrično zaměnjaje ranějšu osnovu obščim multiplikacijskim prstenom.\par
\noindent\emph{Reason:} The contrast between kolco glavnyh idealov and izhodno kolco matters to the extension of the argument in Section 8 of Paper 31. The first names an algebraic property; the second names a ring's role in the construction. The new starting ring need not have the same special type as the earlier one. The rationale accordingly retains the full principal-ideal expression rather than an abbreviation that could obscure the contrast. The current passage expressly replaces the earlier base with a more general multiplication ring.\par
\noindent{\footnotesize Izvorne zapisy / Original records: HIST-g155-002.}\par
\smallskip\noindent\emph{Sravnjenje s tekstom:} Početok osmogo paragrafa raboty 31 zaměnjaje ranějši prsten glavnyh idealov glavnym poredkom čislovogo ili funkcionalnogo polja i potom obščim multiplikacijskim prstenom. Njegove hipotezy izrično uključajut odsutstvo děliteljev nuly i egzistenciju jediničnogo elementa. Imena prstena, poredka i polja zato ne sut zaměnime opisy jednoj istej hipotezy.\par
\noindent\emph{Textual comparison:} The opening of Section 8 of Paper 31 replaces the earlier principal ideal ring by a principal order of a number or function field, then by a more general multiplication ring. Its hypotheses expressly include absence of zero divisors and existence of an identity. The names of the ring, order, and field are therefore not interchangeable descriptions of one hypothesis.\par
\noindent{\footnotesize TEXT-COMP-79; pasaže / passages: 18787, 18789; \texttt{CORPUS\_FORM\_COMPARISONS.json}.}\par
\end{minipage}\par
\par\medskip\Needspace{7\baselineskip}
\noindent\begin{minipage}{\linewidth}
\hypertarget{TERM-0791}{}
\noindent\textbf{kolco glavnyh idealov $\mathfrak{o}_{\mathfrak{p}}$}\quad{\footnotesize\texttt{TERM-0791}}\par
\noindent\emph{Iztočnik / Source:} Hauptidealring $\mathfrak{o}_{\mathfrak{p}}$\par
\noindent\emph{Rusky / Russian:} кольцо главных идеалов $\mathfrak{o}_{\mathfrak{p}}$\par
\noindent\emph{Ukrajinsky / Ukrainian:} кільце головних ідеалів $\mathfrak{o}_{\mathfrak{p}}$\par
\smallskip\noindent\emph{Razlog:} Kolco glavnyh idealov s oznakoju $\mathfrak{o}_{\mathfrak p}$ prodolžuje polny naziv prstena, ne kratice. Indeks identificuje konkretno lokalno razširjenje v rabotě 31; ne jest čestjo jezyčnogo nazva. Aktualna pasaža svęzuje toj prsten s modulnoju bazoju drugogo prstena kvocientov i priměnjenjem diskriminantnogo teorema. Lokalni makrodefiniciji oprěděljajut frakturne bukvy; ih tipografija mora byti sohranjena osobno od izbora slov.\par
\noindent\emph{Reason:} Kolco glavnyh idealov with the label $\mathfrak{o}_{\mathfrak p}$ continues the full name instead of an abbreviation. The subscript identifies a particular local extension in Paper 31; it is not part of the lexical expression. The current passage connects this ring with a module basis of another quotient ring and application of the discriminant theorem. Local macro definitions specify the fraktur letters, whose typography must be preserved independently of word choice.\par
\noindent{\footnotesize Izvorne zapisy / Original records: HIST-g161-004.}\par
\smallskip\noindent\emph{Sravnjenje s tekstom:} Pasaža raboty 31 piše $\mathfrak{o}_{\mathfrak p}$ pri opisu modulnoj bazy i priměnjenja diskriminantnogo teorema. Makrodefiniciji oprěděljajut obě frakturne bukvy. To podpiraje čitlivy matematičny zapis historijnogo zaglavja o\_p bez proměny izvornyh podatkov.\par
\noindent\emph{Textual comparison:} The Paper 31 passage writes $\mathfrak{o}_{\mathfrak p}$ in describing the module basis and application of the discriminant theorem. Macro definitions specify both fraktur letters. This supports a readable mathematical presentation of historical heading o\_p without altering the original data.\par
\noindent{\footnotesize TEXT-COMP-80; pasaže / passages: 18982, 18985, 19003; \texttt{CORPUS\_FORM\_COMPARISONS.json}.}\par
\end{minipage}\par
\par\medskip\Needspace{7\baselineskip}
\noindent\begin{minipage}{\linewidth}
\hypertarget{TERM-0792}{}
\noindent\textbf{kolco glavnyh idealov bez děliteljev nuly i s jediniceju}\quad{\footnotesize\texttt{TERM-0792}}\par
\noindent\emph{Iztočnik / Source:} Hauptidealring ohne Nullteiler und mit Einheitselement\par
\noindent\emph{Rusky / Russian:} кольцо главных идеалов без делителей нуля и с единицей\par
\noindent\emph{Ukrajinsky / Ukrainian:} кільце головних ідеалів без дільників нуля і з одиницею\par
\smallskip\noindent\emph{Razlog:} Cela fraza zadržuje tri odvojene uslovja: vsi idealy sut glavne, dělitelji nuly ne nastupajut i postoji jediničny element. Razlog izrično predpočita sohraniti avtorove hipotezy, iako sovrěmenna kratica često vključaje něktore iz njih po definiciji. Izpuščenje dodatkov bi učinilo prěvod zavisnym od definicije, ktoru čitatelj sam prinese. Zato opisna dolžina služi preciznosti. Zapis ne rešaje sam po sobě vse morfologične varianty instrumentalne formy jedinice.\par
\noindent\emph{Reason:} The full phrase retains three separate conditions: every ideal is principal, there are no zero divisors, and an identity element exists. The rationale expressly favours preserving the author's hypotheses even though a modern abbreviation often incorporates some of them by definition. Omitting the qualifications would make the translation depend on whichever definition the reader supplies. Its descriptive length therefore serves precision. The record does not itself settle every morphological variant of the instrumental form of unit.\par
\noindent{\footnotesize Izvorne zapisy / Original records: HIST-g149-002.}\par
\end{minipage}\par
\par\medskip\Needspace{7\baselineskip}
\noindent\begin{minipage}{\linewidth}
\hypertarget{TERM-0793}{}
\noindent\textbf{glavny klas}\quad{\footnotesize\texttt{TERM-0793}}\par
\noindent\emph{Iztočnik / Source:} Hauptklasse\par
\noindent\emph{Rusky / Russian:} главный класс\par
\noindent\emph{Ukrajinsky / Ukrainian:} головний клас\par
\smallskip\noindent\emph{Razlog:} Glavny klas svęzuje pridavnik s idealnym klasom, ne s proizvoljnoju glavnoju skupinoju objektov. V rabotě 31 ide o klas jediničnogo ideala; druhy razlog povezuje jego s klasom prědstavjenja danym bazoju. Prvy razlog ostavja otkrytu alternativu specialnogo medžunarodnogo pridavnika za principal, ale ne dava konkretno medžuslovjansko slovo ni tvrdženje o pozdnějšem odbacanju. To jest dokumentovana možnost dalšego sravnjenja, ne izmyšljeny poraženy kandidat.\par
\noindent\emph{Reason:} Glavny klas connects the adjective with an ideal class, not an arbitrary leading group of objects. In Paper 31 it is the class of the unit ideal; a second rationale connects it with the representation class supplied by a basis. The first leaves open the alternative of a specialized international adjective for principal but gives neither an exact Interslavic form nor evidence of its subsequent rejection. This is a documented possibility for further comparison, not an invented losing candidate.\par
\noindent{\footnotesize Izvorne zapisy / Original records: HIST-g134-004, HIST-g140-004.}\par
\end{minipage}\par
\par\medskip\Needspace{7\baselineskip}
\noindent\begin{minipage}{\linewidth}
\hypertarget{TERM-0794}{}
\noindent\textbf{glavny modul}\quad{\footnotesize\texttt{TERM-0794}}\par
\noindent\emph{Iztočnik / Source:} Hauptmodul\par
\noindent\emph{Rusky / Russian:} главный модуль\par
\noindent\emph{Ukrajinsky / Ukrainian:} головний модуль\par
\smallskip\noindent\emph{Razlog:} Glavny modul zadržuje kratky termin dlja modula s jednim porodžujučim elementom. Razlog izrično sravnjaje njegovu ponovnu upotrěblivost v dokazu s dolžšim opisom črez jedin generator i dava prednost kratkomu nazvu. To ne zabranjaje opis pri prvom definovanju. Razlika jest medžu imenom i objasnjenjem, ne medžu dvoma matematičnyma objektama. Prednost kompaktnosti jest zabilěžena kvalitativno; čislena težina za nju ne jest dana.\par
\noindent\emph{Reason:} Glavny modul retains a compact term for a module with one generator. The rationale explicitly compares its reusability in the proof with the longer one-generator description and favours the short name. That does not prohibit the description at the first definition. The distinction is between a name and an explanation, not between two mathematical objects. Compactness is a recorded qualitative preference; no numerical weight is supplied for it.\par
\noindent{\footnotesize Izvorne zapisy / Original records: HIST-g108-007.}\par
\end{minipage}\par
\par\medskip\Needspace{7\baselineskip}
\noindent\begin{minipage}{\linewidth}
\hypertarget{TERM-0795}{}
\noindent\textbf{glavny porjadok}\quad{\footnotesize\texttt{TERM-0795}}\par
\noindent\emph{Iztočnik / Source:} Hauptordnung\par
\noindent\emph{Rusky / Russian:} главный порядок\par
\noindent\emph{Ukrajinsky / Ukrainian:} головний порядок\par
\smallskip\noindent\emph{Razlog:} Glavny porjadok v rabotě 43 prodolžuje staršu terminologiju algebraičnyh poredkov. Aktualny uvod identificuje jeho s sistemom cělyh čisel čislovogo polja. Zato porjadok tu ne jest poredok grupy ni uporjadkovanje elementov. Historijno nazvanje termina stabilnym znači kontinuitet proektnogo jezyka, ne dokaz obščej jezyčnoj normy. Sravnjenje s TERM-0796 nužno jest, zato že isto německo slovo v drugoj rabotě dobilo naziv maksimalny porjadok.\par
\noindent\emph{Reason:} Glavny porjadok in Paper 43 continues earlier terminology for algebraic orders. The current introduction identifies it with the system of integers of the number field. Order here is therefore neither a group's order nor an ordering of elements. The historical description of the term as stable means continuity of project usage, not proof of a general linguistic norm. Comparison with TERM-0796 is necessary because the same German word received maksimalny porjadok in another paper.\par
\noindent{\footnotesize Izvorne zapisy / Original records: HIST-g213-011.}\par
\smallskip\noindent\emph{Sravnjenje s tekstom:} Rabota 43 nazyvaje sistem cělyh čisel čislovogo polja glavnym poredkom. Rabota 37 govori o maksimalnyh poredkah i v konkretnom dokazu objasnjaje maksimalnost nemožnostju dalšego razširjenja. Dve historijne motivacije zato podpirajut različne nazvy za svęzane aritmetične objekty. Konsistentna rekomendacija trěbuje sravnjenja definicij, ne avtomatičnoj zaměny vseh pojavjenj Haupt-.\par
\noindent\emph{Textual comparison:} Paper 43 calls the system of integers of a number field its principal order. Paper 37 speaks of maximal orders and explains maximality in a particular proof through the impossibility of further enlargement. The two historical motivations therefore support different names for related arithmetic objects. A consistent recommendation requires comparison of definitions, not automatic replacement of every Haupt- occurrence.\par
\noindent{\footnotesize TEXT-COMP-81; pasaže / passages: 22785, 22804, 24741, 24769; \texttt{CORPUS\_FORM\_COMPARISONS.json}.}\par
\end{minipage}\par
\par\medskip\Needspace{7\baselineskip}
\noindent\begin{minipage}{\linewidth}
\hypertarget{TERM-0796}{}
\noindent\textbf{maksimalny porjadok}\quad{\footnotesize\texttt{TERM-0796}}\par
\noindent\emph{Iztočnik / Source:} Hauptordnung\par
\noindent\emph{Rusky / Russian:} максимальный порядок\par
\noindent\emph{Ukrajinsky / Ukrainian:} максимальний порядок\par
\smallskip\noindent\emph{Razlog:} Maksimalny porjadok vybira smyslovy prěvod vměsto doslovnogo glavny. Razlog pri rabotě 37 čita Hauptordnung kako maksimalny poredok, v čislovom polju prsten cělyh čisel. Aktualny dokaz izrično govori o nemožnosti razširiti takovy poredok do obširnějšego pri danoj osnově. Ta motivacija različna jest od kontinuiteta frazy glavny porjadok pri TERM-0795. Oba razloga trěba pokazati vměstě; prosty spis ne može bez kontekstnogo sravnjenja proglasiti vse upotrěby jednoj formy pravilnymi, a drugoj nepravilnymi.\par
\noindent\emph{Reason:} Maksimalny porjadok chooses a sense-based rendering rather than literal principal. The rationale for Paper 37 reads Hauptordnung as maximal order, the ring of integers in a number field. The current proof expressly discusses the impossibility of enlarging such an order at the specified base. This differs from the continuity argument for glavny porjadok at TERM-0795. Both reasons must be visible: a word list cannot, without contextual comparison, declare every use of one form correct and every use of the other incorrect.\par
\noindent{\footnotesize Izvorne zapisy / Original records: HIST-g207-003.}\par
\smallskip\noindent\emph{Sravnjenje s tekstom:} Rabota 43 nazyvaje sistem cělyh čisel čislovogo polja glavnym poredkom. Rabota 37 govori o maksimalnyh poredkah i v konkretnom dokazu objasnjaje maksimalnost nemožnostju dalšego razširjenja. Dve historijne motivacije zato podpirajut različne nazvy za svęzane aritmetične objekty. Konsistentna rekomendacija trěbuje sravnjenja definicij, ne avtomatičnoj zaměny vseh pojavjenj Haupt-.\par
\noindent\emph{Textual comparison:} Paper 43 calls the system of integers of a number field its principal order. Paper 37 speaks of maximal orders and explains maximality in a particular proof through the impossibility of further enlargement. The two historical motivations therefore support different names for related arithmetic objects. A consistent recommendation requires comparison of definitions, not automatic replacement of every Haupt- occurrence.\par
\noindent{\footnotesize TEXT-COMP-81; pasaže / passages: 22785, 22804, 24741, 24769; \texttt{CORPUS\_FORM\_COMPARISONS.json}.}\par
\end{minipage}\par
\par\medskip\Needspace{7\baselineskip}
\noindent\begin{minipage}{\linewidth}
\hypertarget{TERM-0797}{}
\noindent\textbf{glavny porjadok / glavna oblast}\quad{\footnotesize\texttt{TERM-0797}}\par
\noindent\emph{Iztočnik / Source:} Hauptordnung / Hauptgebiet\par
\noindent\emph{Rusky / Russian:} главный порядок / главная область\par
\noindent\emph{Ukrajinsky / Ukrainian:} головний порядок / головна область\par
\smallskip\noindent\emph{Razlog:} Glavny porjadok i glavna oblast zadržujut obču čest Haupt-, ale ne označajut isti tip objekta. V rabotě 42 oblast grupuje maksimalne poredky po točnyh aritmetičnyh uslovjah; glavna oblast oprěděljaje se črez prěsěk s podpoljem. Proto oblast ne dolžno avtomatično čitati kak oblast cělosti. Historijna motivacija priznaje potrěbu togo kontekstnogo razlikovanja. Samo sosedstvo slova porjadok ne jest dostatočna definicija; čitatelju potrěbne sut uslovja klasifikacije.\par
\noindent\emph{Reason:} Glavny porjadok and glavna oblast retain the shared Haupt- element but do not name the same type of object. In Paper 42 an oblast groups maximal orders by specified arithmetic conditions; the principal oblast is characterized through intersection with a subfield. Oblast must therefore not automatically be read as integral domain. The historical rationale acknowledges this need for contextual distinction. Merely placing order nearby does not supply a definition; the reader needs the classification conditions.\par
\noindent{\footnotesize Izvorne zapisy / Original records: HIST-g212-008.}\par
\smallskip\noindent\emph{Sravnjenje s tekstom:} V rabotě 42 oblast po danomu podpolju obuhvata maksimalne poredky s jednakym prěsěkom i dodatnymi lokalnymi uslovjami. Ako prěsěk jest glavny poredok, tekst govori o glavnoj oblasti. To jest klasifikacija poredkov, ne opča definicija oblasti cělosti.\par
\noindent\emph{Textual comparison:} In Paper 42 an oblast relative to a subfield comprises maximal orders with the same intersection and additional local conditions. If the intersection is a principal order, the text calls the oblast principal. This classifies orders; it is not a general definition of an integral domain.\par
\noindent{\footnotesize TEXT-COMP-82; pasaže / passages: 24657; \texttt{CORPUS\_FORM\_COMPARISONS.json}.}\par
\end{minipage}\par
\par\medskip\Needspace{7\baselineskip}
\noindent\begin{minipage}{\linewidth}
\hypertarget{TERM-0798}{}
\noindent\textbf{glavny porjadok / kvadrat prostogo ideala}\quad{\footnotesize\texttt{TERM-0798}}\par
\noindent\emph{Iztočnik / Source:} Hauptordnung / Quadrat eines Primideals\par
\noindent\emph{Rusky / Russian:} главный порядок / квадрат простого идеала\par
\noindent\emph{Ukrajinsky / Ukrainian:} головний порядок / квадрат простого ідеалу\par
\smallskip\noindent\emph{Razlog:} Glavny porjadok i kvadrat prostogo ideala povezuju terminologiju poredkov s idealnym umnoženjem. Razlog prodolžuje uvod i sravnjenje s Dedekindom. Kvadrat tu znači produkt ideala so soboju, ne geometričnu figuru ni samo množestvo kvadratov njegovih elementov. Prosty kvalifikuje ideal, ne eksponent. Taka composicionna jasnost jest razlog zadržati vse česti frazy; očekivana razpoznatelnost ne byla podkrepljena čitatelskym testom.\par
\noindent\emph{Reason:} Glavny porjadok and kvadrat prostogo ideala connect order terminology with ideal multiplication. The rationale continues the introduction and its comparison with Dedekind. Square means the ideal's product with itself, not a geometric figure or merely the set of squares of its elements. Prime qualifies the ideal, not the exponent. This compositional clarity supports retaining every part of the phrase; anticipated recognizability was not backed by a reader test.\par
\noindent{\footnotesize Izvorne zapisy / Original records: HIST-g154-003.}\par
\end{minipage}\par
\par\medskip\Needspace{7\baselineskip}
\noindent\begin{minipage}{\linewidth}
\hypertarget{TERM-0799}{}
\noindent\textbf{glavny porjadok relativnogo polja}\quad{\footnotesize\texttt{TERM-0799}}\par
\noindent\emph{Iztočnik / Source:} Hauptordnung des Relativkörpers\par
\noindent\emph{Rusky / Russian:} главный порядок относительного поля\par
\noindent\emph{Ukrajinsky / Ukrainian:} головний порядок відносного поля\par
\smallskip\noindent\emph{Razlog:} Glavny porjadok relativnogo polja kombinuje dva ranějše nazvy, ne tvori neobrazloženy novy pojem. Relativny sohranjaje zavisimost razsmotrenja od podloženogo polja, a glavny porjadok nosi aritmetičnu funkciju. Razlog daje prednost terminologičnoj kontinuitetnosti v osmom paragrafu raboty 31. On ne dava samostojnyh težin za vse sinonimy slova relativny. Točno osnovno polje i ostale hipotezy morajut ostati v tekstu; pridavnik ne zaměnjaje ih.\par
\noindent\emph{Reason:} Glavny porjadok relativnogo polja combines two earlier terms rather than creating an unexplained new concept. Relative preserves dependence on an underlying field, while principal order carries the arithmetic role. The rationale favours continuity within Section 8 of Paper 31. It does not supply separate weights for every synonym of relative. The precise base field and remaining hypotheses must remain in the text; the adjective does not replace them.\par
\noindent{\footnotesize Izvorne zapisy / Original records: HIST-g162-006.}\par
\end{minipage}\par
\par\medskip\Needspace{7\baselineskip}
\noindent\begin{minipage}{\linewidth}
\hypertarget{TERM-0800}{}
\noindent\textbf{glavny porjadok čislovogo abo funkcionalnogo polja}\quad{\footnotesize\texttt{TERM-0800}}\par
\noindent\emph{Iztočnik / Source:} Hauptordnung eines Zahl- oder Funktionenkörpers\par
\noindent\emph{Rusky / Russian:} главный порядок числового или функционального поля\par
\noindent\emph{Ukrajinsky / Ukrainian:} головний порядок числового або функціонального поля\par
\smallskip\noindent\emph{Razlog:} Glavny porjadok čislovogo abo funkcionalnogo polja zadržuje alternativu medžu dvoma vidami polja. Funkcionalno polje označaje tehničny objekt, ne proizvoljny prostor funkcij. Razlog izbira tę razpoznatelnu konstrukciju prěd užšim nacionalnym kalkom, ale ne nazyvaje konkretny kalk: jeho točnu formu ne trěba izmyšljati. Aktualny uvod osmogo paragrafa daje toj frazě osobnu rolju pri prěhodu od prstena glavnyh idealov k obščej koeficientnoj osnově. Toj prěhod ne znači, že vsaky glavny poredok imaje vse idealy glavne.\par
\noindent\emph{Reason:} Glavny porjadok čislovogo abo funkcionalnogo polja retains the alternative between two kinds of field. Function field names a technical object, not an arbitrary space of functions. The rationale prefers this recognizable construction to a narrower national calque but does not identify that calque; its exact form must not be invented. The current opening of Section 8 gives the phrase a particular role in passing from a principal ideal ring to a more general coefficient base. That passage does not mean that every principal order has every ideal principal.\par
\noindent{\footnotesize Izvorne zapisy / Original records: HIST-g155-003.}\par
\smallskip\noindent\emph{Sravnjenje s tekstom:} Početok osmogo paragrafa raboty 31 zaměnjaje ranějši prsten glavnyh idealov glavnym poredkom čislovogo ili funkcionalnogo polja i potom obščim multiplikacijskim prstenom. Njegove hipotezy izrično uključajut odsutstvo děliteljev nuly i egzistenciju jediničnogo elementa. Imena prstena, poredka i polja zato ne sut zaměnime opisy jednoj istej hipotezy.\par
\noindent\emph{Textual comparison:} The opening of Section 8 of Paper 31 replaces the earlier principal ideal ring by a principal order of a number or function field, then by a more general multiplication ring. Its hypotheses expressly include absence of zero divisors and existence of an identity. The names of the ring, order, and field are therefore not interchangeable descriptions of one hypothesis.\par
\noindent{\footnotesize TEXT-COMP-79; pasaže / passages: 18787, 18789; \texttt{CORPUS\_FORM\_COMPARISONS.json}.}\par
\end{minipage}\par
\par\medskip\Needspace{7\baselineskip}
\noindent\begin{minipage}{\linewidth}
\hypertarget{TERM-0801}{}
\noindent\textbf{glavny porjadok čislovogo polja}\quad{\footnotesize\texttt{TERM-0801}}\par
\noindent\emph{Iztočnik / Source:} Hauptordnung eines Zahlkörpers\par
\noindent\emph{Rusky / Russian:} главный порядок числового поля\par
\noindent\emph{Ukrajinsky / Ukrainian:} головний порядок числового поля\par
\smallskip\noindent\emph{Razlog:} Glavny porjadok čislovogo polja prodolžuje imenovanje iz ranějših česti raboty 31. Razlog jest terminologična kontinuitetnost, ne novy dokaz, že glavny jest najlučši prěvod vsih pojavjenj Hauptordnung. Kvalifikacija čislovogo polja čuva aritmetičny kontekst: ne ide o poredok grupy ili obyčno uporjadkovanje. Sravnjenje pri TERM-0795 i TERM-0796 pokazuje, začto tu frazu trěba svęzati s definicijeju poredka, a ne oceniti samo po slovnoj sličnosti.\par
\noindent\emph{Reason:} Glavny porjadok čislovogo polja continues the wording of earlier parts of Paper 31. The reason is terminological continuity, not a new demonstration that principal is the best translation of every Hauptordnung. The number-field qualification preserves the arithmetic context: this is neither group order nor ordinary ordering. The comparison at TERM-0795 and TERM-0796 shows why the phrase needs a definition of the order, not assessment by verbal resemblance alone.\par
\noindent{\footnotesize Izvorne zapisy / Original records: HIST-g164-008.}\par
\end{minipage}\par
\par\medskip\Needspace{7\baselineskip}
\noindent\begin{minipage}{\linewidth}
\hypertarget{TERM-0802}{}
\noindent\textbf{glavny rjad}\quad{\footnotesize\texttt{TERM-0802}}\par
\noindent\emph{Iztočnik / Source:} Hauptreihe\par
\noindent\emph{Rusky / Russian:} главный ряд\par
\noindent\emph{Ukrajinsky / Ukrainian:} головний ряд\par
\smallskip\noindent\emph{Razlog:} Glavny rjad sohranjaje specialny smysl glavnoj serije podgrup. Oba razloga odbivajut zaměnu obščim nazvom normalnogo reda. V rabotě 34 dopustimost zavisi od operatorov; pri vključanju vseh vnutrnyh avtomorfizmov členi sut normalne podgrupy, a red ne dopušča dopolniteljnogo člena togo tipa medžu sosednymi členami. Ta nemožnost dopolnjenja nosi informaciju, koju prosty naziv normalny red ne izražaje. Historijno rjad i aktualno red ostavajut različnymi zapisami, ne dokazom dvuh različnih matematičnyh konstrukcij.\par
\noindent\emph{Reason:} Glavny rjad preserves the special sense of a chief series of subgroups. Both rationales reject the more general label normal series. In Paper 34 admissibility depends on the operators; when all inner automorphisms are included, the members are normal subgroups and no further member of the required kind can be inserted between adjacent members. This inability to refine the series carries information absent from normal series alone. Historical rjad and current red remain different wordings, not evidence of two different mathematical constructions.\par
\noindent{\footnotesize Izvorne zapisy / Original records: HIST-g181-003, HIST-g182-007.}\par
\smallskip\noindent\emph{Sravnjenje s tekstom:} V rabotě 34 kompozicijny red ne dopušča vstavjenje dalšego dopustimogo normalnogo člena medžu dvoma sosednymi členami. Pri operatorah, vključajučih vse vnutrne avtomorfizmy, tekst nazyvaje jego glavnym redom. Zato istorijne slovo rjad ne možno objasniti samo kak proizvoljny normalny red; potrěbna jest i nemožnost dopolnjenja v danej kategoriji.\par
\noindent\emph{Textual comparison:} In Paper 34 a composition series admits no further admissible normal member between two adjacent members. With operators including all inner automorphisms, the text calls it a chief series. Historical rjad therefore cannot be explained merely as an arbitrary normal series: the absence of refinement in the specified category also matters.\par
\noindent{\footnotesize TEXT-COMP-83; pasaže / passages: 19965, 19967, 19969, 19971, 19973; \texttt{CORPUS\_FORM\_COMPARISONS.json}.}\par
\end{minipage}\par
\par\medskip\Needspace{7\baselineskip}
\noindent\begin{minipage}{\linewidth}
\hypertarget{TERM-0803}{}
\noindent\textbf{glavna teorema}\quad{\footnotesize\texttt{TERM-0803}}\par
\noindent\emph{Iztočnik / Source:} Hauptsatz\par
\noindent\emph{Rusky / Russian:} главная теорема\par
\noindent\emph{Ukrajinsky / Ukrainian:} головна теорема\par
\smallskip\noindent\emph{Razlog:} Glavna teorema prevodi Hauptsatz kako centralno matematično tvrđenje. Razlog izrično odbiva doslovnu interpretaciju Satz kak rečenice. Pridavnik govori o rolji rezultata v rabotě, ne o dodatnoj hipotezě ili osobnoj matematičnoj klasifikaciji. Ruska i ukrajinska sravniteljnica podpirajut konstrukciju pridavnik plus naziv teoremy. Sohranjena motivacija ne rešaje odnos femininoj teorema k maskulinomu teorem iz drugih zapisov; tu morfologičnu variaciju trěba razlikovati od matematičnogo smysla.\par
\noindent\emph{Reason:} Glavna teorema renders Hauptsatz as a central mathematical assertion. The rationale expressly rejects treating Satz literally as sentence. The adjective concerns the result's role in the paper, not an additional hypothesis or a separate mathematical classification. Russian and Ukrainian comparanda support the adjective-plus-theorem construction. The preserved motivation does not settle the relation between feminine teorema and masculine teorem elsewhere; that morphological variation must be distinguished from mathematical meaning.\par
\noindent{\footnotesize Izvorne zapisy / Original records: HIST-g208-001.}\par
\end{minipage}\par
\par\medskip\Needspace{7\baselineskip}
\noindent\begin{minipage}{\linewidth}
\hypertarget{TERM-0804}{}
\noindent\textbf{glavny teorem / pomočny teorem / dodatok / zamětka}\quad{\footnotesize\texttt{TERM-0804}}\par
\noindent\emph{Iztočnik / Source:} Hauptsatz / Hilfssatz / Zusatz / Bemerkung\par
\noindent\emph{Rusky / Russian:} главная теорема / лемма / дополнение / замечание\par
\noindent\emph{Ukrajinsky / Ukrainian:} головна теорема / допоміжна теорема / додаток / зауваження\par
\smallskip\noindent\emph{Razlog:} Četyri nazvy čuvajut hierarhiju izloženja: glavny rezultat, pomočny rezultat, dodatok i zamětku. Pri Hilfssatz ruska sravniteljnica izbira lemmu, a ukrajinska pomočnu teoremu. Medžuslovjansky razlog daje prednost pomočny teorem, ibo lema v toj česti prěvoda ješče nebyla ustaljena. To jest lokalno obrazloženje, ne tvrdženje, že lema ne može byti medžuslovjansky termin: TERM-0816 i TERM-0818 dokumentujut upravo tu alternativu. Dodatok i zamětka takože ne morajut byti avtomatično preimenovane v teoremy.\par
\noindent\emph{Reason:} The four names preserve the exposition's hierarchy: principal result, auxiliary result, addition, and remark. For Hilfssatz the Russian comparandum chooses lemma and the Ukrainian auxiliary theorem. The Interslavic rationale favours pomočny teorem because lemma had not been settled in that part of the translation. This is a local justification, not a claim that lema cannot be an Interslavic term: TERM-0816 and TERM-0818 document that very alternative. Addition and remark likewise must not automatically become theorems.\par
\noindent{\footnotesize Izvorne zapisy / Original records: HIST-g214-009.}\par
\end{minipage}\par
\par\medskip\Needspace{7\baselineskip}
\noindent\begin{minipage}{\linewidth}
\hypertarget{TERM-0805}{}
\noindent\textbf{glavny slěd}\quad{\footnotesize\texttt{TERM-0805}}\par
\noindent\emph{Iztočnik / Source:} Hauptspur\par
\noindent\emph{Rusky / Russian:} главный след\par
\noindent\emph{Ukrajinsky / Ukrainian:} головний слід\par
\smallskip\noindent\emph{Razlog:} Glavny slěd jest definovany črez regularno prědstavjenje. Tri razloga zadržujut taj naziv i njegov kontinuitet v rabotě 34, vměsto slobodnoj modernizacije. Smysl se opira na definiciju, ne na predstavu, že jeden slěd jest važnějši od drugogo. Tekst osobno definira redukovany slěd črez različne ireducibilne prědstavjenja; pri razširjenju polja dvě veličiny ne vedut se jednako. Zato glavny i redukovany ne sut zaměnime stilistične dodatky.\par
\noindent\emph{Reason:} Glavny slěd is defined through the regular representation. Three rationales retain that name and its continuity through Paper 34 rather than modernizing it freely. Its meaning rests on the definition, not on the idea that one trace is more important than another. The text separately defines reduced trace through distinct irreducible representations; the two quantities do not behave identically under field extension. Principal and reduced are therefore not interchangeable stylistic qualifications.\par
\noindent{\footnotesize Izvorne zapisy / Original records: HIST-g202-002, HIST-g203-004, HIST-g204-005.}\par
\smallskip\noindent\emph{Sravnjenje s tekstom:} Definicija glavnogo slěda v rabotě 34 odkazuje na regularno prědstavjenje; redukovany slěd definira se odvojeno. Pozdněj pasaža govori, že prvy ne měnjaje se pri prěhodu k algebraično zatvorjenomu polju, a to samo ne jest tvrdženo za vtory. Ti dva nazvy zato ne označajut avtomatično istu veličinu.\par
\noindent\emph{Textual comparison:} The definition of principal trace in Paper 34 refers to the regular representation; reduced trace is defined separately. A later passage says that the former is unchanged on passing to an algebraically closed field, while the same is not asserted for the latter. The two names therefore do not automatically denote the same quantity.\par
\noindent{\footnotesize TEXT-COMP-84; pasaže / passages: 22203, 22234; \texttt{CORPUS\_FORM\_COMPARISONS.json}.}\par
\end{minipage}\par
\par\medskip\Needspace{7\baselineskip}
\noindent\begin{minipage}{\linewidth}
\hypertarget{TERM-0806}{}
\noindent\textbf{najvyšši elementarny dělitelj}\quad{\footnotesize\texttt{TERM-0806}}\par
\noindent\emph{Iztočnik / Source:} Highest elementary divisor\par
\noindent\emph{Rusky / Russian:} наивысший элементарный делитель\par
\noindent\emph{Ukrajinsky / Ukrainian:} найвищий елементарний дільник\par
\smallskip\noindent\emph{Razlog:} Najvyšši elementarny dělitelj prodolžuje determinantnu i elementarno-děliteljsku terminologiju raboty 22. Sohranjena motivacija jest kratka: upira se na kontinuitet s ranějšimi paragrafami, ne dava novogo sravnjenja superlativov. Najvyšši dolžno razuměti po poredku, danomu v toj teoriji, ne kak čislenu veličinu proizvoljnogo polinoma. Sam naziv ne može zaměniti definiciju dělitelja i način njegovogo uporjadkovanja. Tablica sohranjaje izbor, ale ne pripisuje jemu nezasvědočenu čislenu ocěnu.\par
\noindent\emph{Reason:} Najvyšši elementarny dělitelj continues the determinant and elementary-divisor vocabulary of Paper 22. The preserved motivation is brief: it relies on continuity with earlier sections rather than offering a new comparison of superlatives. Highest must be understood in the order specified by that theory, not as the numerical magnitude of an arbitrary polynomial. The name cannot replace the divisor's definition or its ordering. The table preserves the choice without attributing an undocumented numerical assessment to it.\par
\noindent{\footnotesize Izvorne zapisy / Original records: HIST-g091-004, HIST-g092-073.}\par
\end{minipage}\par
\par\medskip\Needspace{7\baselineskip}
\noindent\begin{minipage}{\linewidth}
\hypertarget{TERM-0807}{}
\noindent\textbf{Hilbertove funkcije}\quad{\footnotesize\texttt{TERM-0807}}\par
\noindent\emph{Iztočnik / Source:} Hilbertsche Anzahlen\par
\noindent\emph{Rusky / Russian:} гильбертовы функции\par
\noindent\emph{Ukrajinsky / Ukrainian:} Гільбертові функції\par
\smallskip\noindent\emph{Razlog:} Hilbertove funkcije jest naměrna modernizacija historijnogo zaglavja Hilbertsche Anzahlen, ne doslovny prěvod slova Anzahlen. Sohranjeny razlog tumači matematičny predmet prěz sovrěmennu terminologiju Hilbertovoj funkcije. Ale rabota 27 sama razlikuje Hilbertove čisla, karakterističnu funkciju i druge čisla Macaulaya. Zato modernizacija može zamagliti razliku, ktora jest predmetom teksta. V diplomatičnom izdanju ne trěba predstavljati taki preimenovany naslov kak věrno sohranjeno zaglavje; moderny pojasniteljny naziv može byti odvojeno označeny. Tuto jest redakcijna odluka i jejo omeženje, ne dokaz, že vse nazvane veličiny sut jedna funkcija.\par
\noindent\emph{Reason:} Hilbertove funkcije is an intentional modernization of historical title Hilbertsche Anzahlen, not a literal translation of Anzahlen. The preserved rationale interprets the mathematical subject through modern Hilbert-function terminology. Yet Paper 27 itself distinguishes Hilbert numbers, the characteristic function, and other numbers introduced by Macaulay. Modernization can therefore obscure a distinction that the text investigates. A diplomatic edition should not present such a renamed title as faithfully retained; a modern explanatory designation can be separately identified. This is an editorial choice and its limitation, not proof that all the named quantities form one function.\par
\noindent{\footnotesize Izvorne zapisy / Original records: HIST-g215-006.}\par
\smallskip\noindent\emph{Sravnjenje s tekstom:} Rabota 27 ima zaglavje Hilbertove čisla i razlikuje několiko svęzanyh počtov i funkcij. Bibliografija v dvoh městah piše Hilbertove funkcije. Sravnjenje zato potvrđuje konkretnu razliku medžu tekstom i bibliografičnym opisom; ne opravdava ih potihu identifikaciju. Starša modernizacija mora byti obrazložena kak redakcijna interpretacija, ne predstavljena kak jediny možny věrny prěvod zaglavja.\par
\noindent\emph{Textual comparison:} Paper 27 is headed Hilbertove čisla and distinguishes several related counts and functions. The bibliography reads Hilbertove funkcije in two places. The comparison therefore establishes a concrete difference between text and bibliographical description, not justification for silently identifying them. The earlier modernization needs explanation as an editorial interpretation, not presentation as the only possible faithful title translation.\par
\noindent{\footnotesize TEXT-COMP-85; pasaže / passages: 15911, 15917, bib-isv.tex:111, bib-isv.tex:189; \texttt{CORPUS\_FORM\_COMPARISONS.json}.}\par
\end{minipage}\par
\par\medskip\Needspace{7\baselineskip}
\noindent\begin{minipage}{\linewidth}
\hypertarget{TERM-0808}{}
\noindent\textbf{Hilbertove čisla}\quad{\footnotesize\texttt{TERM-0808}}\par
\noindent\emph{Iztočnik / Source:} Hilbertsche Anzahlen\par
\noindent\emph{Rusky / Russian:} гильбертовы числа\par
\noindent\emph{Ukrajinsky / Ukrainian:} Гільбертові числа\par
\smallskip\noindent\emph{Razlog:} Hilbertove čisla zadržuje historijno imenovanje počtov. Razlog izrično razlikuje ih od Hilbertovoj karakterističnoj funkcije i od pozdnějšeho skračenogo nazva Hilbertova funkcija. Aktualna rabota 27 objasnjaje nekoliko svęzanyh počtov ostatkovyh klasov i kompozicijnyh členov; ne svodi jih vseh do jednogo bezuslovnogo funkcionalnogo nazva. Zato čisla jest smyslovo važny izbor, ne prosta neochota upotrěbiti sovrěmenny termin. Sravnjenje s modernizovanym bibliografičnym zaglavjem pri TERM-0807 ostaje izrično.\par
\noindent\emph{Reason:} Hilbertove čisla retains the historical name for counts. The rationale expressly distinguishes them from Hilbert's characteristic function and the later shorthand Hilbert function. The current Paper 27 explains several related counts of residue classes and composition factors; it does not reduce them all to one unqualified functional label. Numbers is therefore a consequential semantic choice, not merely reluctance to use a modern term. The comparison with the modernized bibliographical title at TERM-0807 remains explicit.\par
\noindent{\footnotesize Izvorne zapisy / Original records: HIST-g104-001.}\par
\smallskip\noindent\emph{Sravnjenje s tekstom:} Rabota 27 ima zaglavje Hilbertove čisla i razlikuje několiko svęzanyh počtov i funkcij. Bibliografija v dvoh městah piše Hilbertove funkcije. Sravnjenje zato potvrđuje konkretnu razliku medžu tekstom i bibliografičnym opisom; ne opravdava ih potihu identifikaciju. Starša modernizacija mora byti obrazložena kak redakcijna interpretacija, ne predstavljena kak jediny možny věrny prěvod zaglavja.\par
\noindent\emph{Textual comparison:} Paper 27 is headed Hilbertove čisla and distinguishes several related counts and functions. The bibliography reads Hilbertove funkcije in two places. The comparison therefore establishes a concrete difference between text and bibliographical description, not justification for silently identifying them. The earlier modernization needs explanation as an editorial interpretation, not presentation as the only possible faithful title translation.\par
\noindent{\footnotesize TEXT-COMP-85; pasaže / passages: 15911, 15917, bib-isv.tex:111, bib-isv.tex:189; \texttt{CORPUS\_FORM\_COMPARISONS.json}.}\par
\end{minipage}\par
\par\medskip\Needspace{7\baselineskip}
\noindent\begin{minipage}{\linewidth}
\hypertarget{TERM-0809}{}
\noindent\textbf{Hilbertovo tvrđenje}\quad{\footnotesize\texttt{TERM-0809}}\par
\noindent\emph{Iztočnik / Source:} Hilbertsche Behauptung\par
\noindent\emph{Rusky / Russian:} утверждение Гильберта\par
\noindent\emph{Ukrajinsky / Ukrainian:} твердження Гільберта\par
\smallskip\noindent\emph{Razlog:} Hilbertovo tvrđenje prevodi Behauptung bez povyšenja do teoremy. Motivacija poštuje označenje, ktore Noether sama vybira v zaglavju odděla. Tvrđenje ne znači neformalnu ili nepravdivu mysl; naziv sohranjaje avtorovu strukturu izloženja i sam po sobě ne govori, dali dokaz slěduje niže. Ruska i ukrajinska sravniteljnica podpirajut tu razliku. Zaměna iz ljubvi k sovrěmennoj terminologiji by byla redakcijnoju proměnoju, ne neutralnym prěvodom.\par
\noindent\emph{Reason:} Hilbertovo tvrđenje translates Behauptung without promoting it to theorem. The motivation respects the designation Noether herself uses in the section heading. Assertion does not mean an informal or false idea; the label preserves the author's expository structure and does not itself say whether a proof follows. The Russian and Ukrainian comparanda support that distinction. Replacing it merely out of preference for modern terminology would be an editorial change, not a neutral translation.\par
\noindent{\footnotesize Izvorne zapisy / Original records: HIST-g075-015.}\par
\end{minipage}\par
\par\medskip\Needspace{7\baselineskip}
\noindent\begin{minipage}{\linewidth}
\hypertarget{TERM-0810}{}
\noindent\textbf{Hilbertov teorem o nerazložimosti}\quad{\footnotesize\texttt{TERM-0810}}\par
\noindent\emph{Iztočnik / Source:} Hilbertscher Irreduzibilitätssatz\par
\noindent\emph{Rusky / Russian:} теорема Гильберта о неприводимости\par
\noindent\emph{Ukrajinsky / Ukrainian:} теорема Гільберта про незвідність\par
\smallskip\noindent\emph{Razlog:} Hilbertov teorem o nerazložimosti kombinuje razpoznatelno ime s opisom matematičnogo svojstva. Razlog predpočita nerazložimost radi očekivanoj prozračnosti v slovjanskoj algebri; ruska i ukrajinska sravniteljnica imajut druge korenjeve konstrukcije za isti pojem. To ne jest izměrena razumělivost. Naziv teorema takože ne oprěděljaje samostojno vsih njegovih hipotez. Avtorova pozdnějša dopolnitelna zamětka trěbuje konečno algebraično čislovo polje pri formulaciji v rabotě 11; rěč o jezyčnom nazvu ne daje pravo izpustiti to uslovje.\par
\noindent\emph{Reason:} Hilbertov teorem o nerazložimosti combines a recognizable name with a description of the mathematical property. The rationale favours nerazložimost for anticipated transparency in Slavic algebra; the Russian and Ukrainian comparanda use other root constructions for the same concept. This is not measured comprehension. Nor does the theorem's name determine all its hypotheses. The author's later supplementary note requires a finite algebraic number field in the formulation in Paper 11; discussion of the linguistic label does not permit omitting that condition.\par
\noindent{\footnotesize Izvorne zapisy / Original records: HIST-g073-006.}\par
\smallskip\noindent\emph{Sravnjenje s tekstom:} Avtorova dopolnitelna zamětka trěbuje pri formulaciji Hilbertovogo teorema v rabotě 11 konečno algebraično čislovo polje, a objasnjaje takože smysl ranějšej frazy o oblasti racionalnosti. To podčrknuje razliku medžu izborom nazva teorema i sohranjenjem njegovih uslovij. Sravnjenje ne jest nova proverka originalnoj pečati; opira se na konkretnu pasažu aktualnogo prěvoda.\par
\noindent\emph{Textual comparison:} The author's supplementary note requires a finite algebraic number field in the formulation of Hilbert's theorem in Paper 11 and also explains the earlier phrase concerning a domain of rationality. This emphasizes the difference between choosing the theorem's name and preserving its conditions. The comparison is not a fresh inspection of the original print; it uses a particular passage of the current translation.\par
\noindent{\footnotesize TEXT-COMP-86; pasaže / passages: 10898; \texttt{CORPUS\_FORM\_COMPARISONS.json}.}\par
\end{minipage}\par
\par\medskip\Needspace{7\baselineskip}
\noindent\begin{minipage}{\linewidth}
\hypertarget{TERM-0811}{}
\noindent\textbf{Hilbertova teorema o nerazložimosti}\quad{\footnotesize\texttt{TERM-0811}}\par
\noindent\emph{Iztočnik / Source:} Hilbertscher Irreduzibilitätssatz\par
\noindent\emph{Rusky / Russian:} теорема Гильберта о неприводимости\par
\noindent\emph{Ukrajinsky / Ukrainian:} теорема Гільберта про незвідність\par
\smallskip\noindent\emph{Razlog:} Hilbertova teorema o nerazložimosti jest feminina konstrukcija istogo imenovanogo teorema, ktory TERM-0810 zapisuje s maskulinym teorem. Sohranjena motivacija govori o medžunarodnom nazvu i slovjanskoj morfologiji, ale ne daje sravniteljnogo dokaza za prednost jednogo roda. Korenj nerazložim- nosi matematičny smysl; rod imenika upravjaje formoju Hilbertova ili Hilbertov. Te dve pytanja ne trěba sliti v jednu ocěnu. Povtorjenje nazva v različnyh rabotah samo po sobě ne opravdava vsaku njegovu formu.\par
\noindent\emph{Reason:} Hilbertova teorema o nerazložimosti is a feminine construction for the same named theorem that TERM-0810 records with masculine teorem. The preserved motivation mentions an international name and Slavic morphology, but offers no comparative demonstration favouring one gender. The irreducibility stem carries the mathematical meaning; the noun's gender determines Hilbertova or Hilbertov. These are separate questions. Recurrence of the name in different papers does not itself justify every form.\par
\noindent{\footnotesize Izvorne zapisy / Original records: HIST-g039-014.}\par
\end{minipage}\par
\par\medskip\Needspace{7\baselineskip}
\noindent\begin{minipage}{\linewidth}
\hypertarget{TERM-0812}{}
\noindent\textbf{Hilbertov teorem o bazisu modula}\quad{\footnotesize\texttt{TERM-0812}}\par
\noindent\emph{Iztočnik / Source:} Hilbertscher Satz von der Modulbasis\par
\noindent\emph{Rusky / Russian:} теорема Гильберта о базисе модуля\par
\noindent\emph{Ukrajinsky / Ukrainian:} теорема Гільберта про базис модуля\par
\smallskip\noindent\emph{Razlog:} Hilbertov teorem o bazisu modula sohranjaje lokalny način, kako Noether prenaša Hilbertovo i Gordanovo razvažanje o konečnoj bazi. Razlog ne hoče zaměniti to imenovanje pozdnějšim obščim jezykom noetherovosti. Baza tu služi obrazloženju konečnoj porodženosti; ne možno bez proverky pridati linearnu nezavisnost, koju slovo baza nosi v nekotoryh drugih kontekstah. Sohranjenje modula i bazy pomogaje čitatelju slěditi samomu dokazu, ne samo razpoznati sovrěmennu klasifikaciju.\par
\noindent\emph{Reason:} Hilbertov teorem o bazisu modula preserves the local way in which Noether transfers Hilbert's and Gordan's finite-basis reasoning. The rationale does not replace this name with the later general language of Noetherianity. Basis here serves the finite-generation argument; one must not add linear independence without checking, merely because basis carries that condition in some other contexts. Retaining module and basis helps the reader follow the proof itself, not only recognize a modern classification.\par
\noindent{\footnotesize Izvorne zapisy / Original records: HIST-g079-032.}\par
\end{minipage}\par
\par\medskip\Needspace{7\baselineskip}
\noindent\begin{minipage}{\linewidth}
\hypertarget{TERM-0813}{}
\noindent\textbf{Hilbertovy teorem I o modulnoj bazi}\quad{\footnotesize\texttt{TERM-0813}}\par
\noindent\emph{Iztočnik / Source:} Hilbertsches Theorem I von der Modulbasis\par
\noindent\emph{Rusky / Russian:} гильбертова теорема I о модульном базисе\par
\noindent\emph{Ukrajinsky / Ukrainian:} Гільбертова теорема I про модульну базу\par
\smallskip\noindent\emph{Razlog:} Hilbertovy teorem I o modulnoj bazi čuva vlastno ime, rimsko čislo i mechanizm modulnoj bazy. I jest identifikator konkretno citovanogo teorema, ne slovo, ktore potrěbuje prěvoda. Aktualny paragraf raboty 6 izrično priměnjaje teorem I k sistemu odgovarjajučih funkcij. Motivacija ne podava polnu gramatičnu obranu historijnoj formy Hilbertovy; njezina razpoznatelnost ne jest dostatočno opravdanje koncovky. Razlikovanje I i II pri TERM-0814 mora ostati i pri pravopisnom uredženju.\par
\noindent\emph{Reason:} Hilbertovy teorem I o modulnoj bazi preserves the name, Roman numeral, and module-basis mechanism. I identifies the particular theorem cited; it is not a word to translate. The current section of Paper 6 expressly applies Theorem I to the system of corresponding functions. The motivation supplies no full grammatical defence of historical Hilbertovy; recognizability is not sufficient justification of the ending. The distinction between I and II at TERM-0814 must survive orthographic editing.\par
\noindent{\footnotesize Izvorne zapisy / Original records: HIST-g052-004.}\par
\smallskip\noindent\emph{Sravnjenje s tekstom:} Teorem I priměnjaje se k sistemu funkcij, aby dobyti konečno prědstavjenje črez vybrane polinomy. Pri teoremě II tekst izrično nazyvaje cěle racionalne čislove koeficienty i pokazuje prěnos na algebraično-cěločislene polinomy. Rimsko čislo i cěločisleny pridavnik zato nesut različnu, sučstvenu informaciju; njih ne možno sliti v obči naziv bez objasnjenja.\par
\noindent\emph{Textual comparison:} Theorem I is applied to a system of functions to obtain a finite representation through selected polynomials. For Theorem II, the text explicitly names rational-integer coefficients and explains the extension to algebraic-integer polynomials. The Roman numeral and integer adjective therefore carry distinct, essential information; they cannot be merged into a general name without explanation.\par
\noindent{\footnotesize TEXT-COMP-87; pasaže / passages: 6252, 6254, 6256, 6258, 6261, 6264, 6265, 6564, 6568, 6571, 6573; \texttt{CORPUS\_FORM\_COMPARISONS.json}.}\par
\end{minipage}\par
\par\medskip\Needspace{7\baselineskip}
\noindent\begin{minipage}{\linewidth}
\hypertarget{TERM-0814}{}
\noindent\textbf{Hilbertov teorem II o celočislenej modulnoj bazi}\quad{\footnotesize\texttt{TERM-0814}}\par
\noindent\emph{Iztočnik / Source:} Hilbertsches Theorem II von der ganzzahligen Modulbasis\par
\noindent\emph{Rusky / Russian:} гильбертова теорема II о целочисленном модульном базисе\par
\noindent\emph{Ukrajinsky / Ukrainian:} Гільбертова теорема II про цілочисельну модульну базу\par
\smallskip\noindent\emph{Razlog:} Hilbertov teorem II o celočislenej modulnoj bazi zadržuje dve razlike od ranějšego nazva: čislo II i cěločislenost. Pasaža raboty 6 prvo priměnjaje teorem za cěle racionalne čislove koeficienty, potom objasnjaje jeho razširjenje na algebraično-cěločislene polinomy pomoču fundamentalnogo sistema algebraičnyh cělyh čisel. Zato cěločisleny ne jest prazdny pridevek pri bazi. Koeficientna oblast i razširjenje sut čestju matematičnogo sodržanja, ktoru kratky naziv ne može zaměniti.\par
\noindent\emph{Reason:} Hilbertov teorem II o celočislenej modulnoj bazi retains two differences from the earlier name: II and the integer qualification. The Paper 6 passage first applies the theorem for rational-integer coefficients, then explains its extension to algebraic-integer polynomials using a fundamental system of algebraic integers. Integer is therefore not an empty adjective attached to basis. The coefficient domain and extension belong to the mathematical content, which the short name cannot replace.\par
\noindent{\footnotesize Izvorne zapisy / Original records: HIST-g055-005.}\par
\smallskip\noindent\emph{Sravnjenje s tekstom:} Teorem I priměnjaje se k sistemu funkcij, aby dobyti konečno prědstavjenje črez vybrane polinomy. Pri teoremě II tekst izrično nazyvaje cěle racionalne čislove koeficienty i pokazuje prěnos na algebraično-cěločislene polinomy. Rimsko čislo i cěločisleny pridavnik zato nesut različnu, sučstvenu informaciju; njih ne možno sliti v obči naziv bez objasnjenja.\par
\noindent\emph{Textual comparison:} Theorem I is applied to a system of functions to obtain a finite representation through selected polynomials. For Theorem II, the text explicitly names rational-integer coefficients and explains the extension to algebraic-integer polynomials. The Roman numeral and integer adjective therefore carry distinct, essential information; they cannot be merged into a general name without explanation.\par
\noindent{\footnotesize TEXT-COMP-87; pasaže / passages: 6252, 6254, 6256, 6258, 6261, 6264, 6265, 6564, 6568, 6571, 6573; \texttt{CORPUS\_FORM\_COMPARISONS.json}.}\par
\end{minipage}\par
\par\medskip\Needspace{7\baselineskip}
\noindent\begin{minipage}{\linewidth}
\hypertarget{TERM-0815}{}
\noindent\textbf{Hilbertov teorem o modulnom bazisu}\quad{\footnotesize\texttt{TERM-0815}}\par
\noindent\emph{Iztočnik / Source:} Hilbertsches Theorem von der Modulbasis\par
\noindent\emph{Rusky / Russian:} теорема Гильберта о базисе модулей\par
\noindent\emph{Ukrajinsky / Ukrainian:} теорема Гільберта про базис модулів\par
\smallskip\noindent\emph{Razlog:} Hilbertov teorem o modulnom bazisu zadržuje Hilbertovo ime i historijnu modulno-baznu formulaciju. Razlog iz raboty 19 daje tomu prednost prěd pozdnějšim jezykom noetherovosti. To ne govori, že sovrěmenny opis jest nespravedlivy; govori, začto on ne zaměnjaje izvorno imenovanje v prěvodu. Bazis i baza v različnyh zapisah ostavajut različnymi jezyčnymi formami. Sohranjena motivacija obrazlagaje matematičnu kontinuitetnost, ale ne daje osobnu morfologičnu ocěnu vseh tih form.\par
\noindent\emph{Reason:} Hilbertov teorem o modulnom bazisu retains Hilbert's name and the historical module-basis formulation. The Paper 19 rationale prefers this to the later language of Noetherianity. That does not declare the modern description wrong; it explains why it does not replace the original designation in translation. Bazis and baza in different records remain different linguistic forms. The motivation justifies mathematical continuity but does not separately assess the morphology of all those forms.\par
\noindent{\footnotesize Izvorne zapisy / Original records: HIST-g081-079.}\par
\end{minipage}\par
\par\medskip\Needspace{7\baselineskip}
\noindent\begin{minipage}{\linewidth}
\hypertarget{TERM-0816}{}
\noindent\textbf{pomocna lema}\quad{\footnotesize\texttt{TERM-0816}}\par
\noindent\emph{Iztočnik / Source:} Hilfssatz\par
\noindent\emph{Rusky / Russian:} вспомогательная лемма\par
\noindent\emph{Ukrajinsky / Ukrainian:} допоміжна лема\par
\smallskip\noindent\emph{Razlog:} Pomocna lema vybira medžunarodny naziv lema i pridavnik označajuči pomočnu rolju. Razlog smatra to obyčnoju matematičnoju prozoju, ne ponizovanjem formalnosti tvrđenja. To se razlikuje od argumenta pri TERM-0804, ktory zadržuje pomočny teorem zaradi lokalno neustaljenogo nazva lema. Obe motivacije sut sohranjene, bez izmyšljanja obščego glasovanja medžu nimi. Forma pomocna v staršej zapisi takože ne dokazuje, že c i č sut svobodno zaměnime; pravopis jest osobno pytanje.\par
\noindent\emph{Reason:} Pomocna lema selects the international noun lema and an adjective marking an auxiliary role. The rationale treats this as ordinary mathematical prose, not a reduction in the statement's formality. This differs from TERM-0804's argument for pomočny teorem because lemma was locally unsettled. Both motivations are preserved without inventing a general vote between them. Historical pomocna also does not establish that c and č are freely interchangeable; spelling is a separate question.\par
\noindent{\footnotesize Izvorne zapisy / Original records: HIST-g057-010.}\par
\end{minipage}\par
\par\medskip\Needspace{7\baselineskip}
\noindent\begin{minipage}{\linewidth}
\hypertarget{TERM-0817}{}
\noindent\textbf{pomoćny teorem}\quad{\footnotesize\texttt{TERM-0817}}\par
\noindent\emph{Iztočnik / Source:} Hilfssatz\par
\noindent\emph{Rusky / Russian:} вспомогательная теорема\par
\noindent\emph{Ukrajinsky / Ukrainian:} допоміжна теорема\par
\smallskip\noindent\emph{Razlog:} Pomoćny teorem zadržuje naziv teorema dlja Hilbertovogo pomočnogo rezultata. Sohranjena motivacija upira se na matematičny register, ne na dokaz, že lema byla sravnjena i odbačena. Historijno ć jest vidimo v tablici; po osnovnom pravopisu čitatelska forma ima č. Taka pravopisna zaměna ne rěšaje leksičnu alternativu teorem ili lema. Sravnjenje medžu zapisami trěbuje oba vida informaciji: kako se slovo piše i začto jest izbrano.\par
\noindent\emph{Reason:} Pomoćny teorem retains theorem for Hilbert's auxiliary result. The preserved motivation appeals to mathematical register, not to evidence that lemma was compared and rejected. Historical ć remains visible in the table; under basic orthography the reading form has č. That spelling correspondence does not settle the lexical alternative theorem or lemma. Comparison between records needs both kinds of information: how the word is written and why it was chosen.\par
\noindent{\footnotesize Izvorne zapisy / Original records: HIST-g050-005.}\par
\end{minipage}\par
\par\medskip\Needspace{7\baselineskip}
\noindent\begin{minipage}{\linewidth}
\hypertarget{TERM-0818}{}
\noindent\textbf{pomočna lema}\quad{\footnotesize\texttt{TERM-0818}}\par
\noindent\emph{Iztočnik / Source:} Hilfssatz\par
\noindent\emph{Rusky / Russian:} вспомогательная лемма\par
\noindent\emph{Ukrajinsky / Ukrainian:} допоміжна лема\par
\smallskip\noindent\emph{Razlog:} Pomočna lema opisuje pomočny rezultat v lokalnom dokazu. Jeden razlog naglaša jeho funkciju, a drugy prodolžuje ranějše imenovanje. To podpiraje kontinuitet lemy, ale ne dokazuje, že vse raboty izbrali isti naziv: TERM-0817 i TERM-0819 zadržujut teorem. Izbor ne dolžno objasnjati kak razliku v istinitosti ili rigoroznosti tvrđenj. Jeho predmet jest matematičny register i organizacija izloženja, s pravopisnoju razlikoju od pomocna pri TERM-0816 sohranjenoj v podatkah.\par
\noindent\emph{Reason:} Pomočna lema describes an auxiliary result in a local proof. One rationale emphasizes its function; the other continues earlier terminology. This supports continuity of lemma but does not establish that every paper chose the same label: TERM-0817 and TERM-0819 retain theorem. The choice should not be explained as a difference in truth or rigour between statements. It concerns mathematical register and expository organization; the spelling difference from pomocna at TERM-0816 remains in the data.\par
\noindent{\footnotesize Izvorne zapisy / Original records: HIST-g042-001, HIST-g060-009.}\par
\end{minipage}\par
\par\medskip\Needspace{7\baselineskip}
\noindent\begin{minipage}{\linewidth}
\hypertarget{TERM-0819}{}
\noindent\textbf{pomočny teorem}\quad{\footnotesize\texttt{TERM-0819}}\par
\noindent\emph{Iztočnik / Source:} Hilfssatz\par
\noindent\emph{Rusky / Russian:} вспомогательное утверждение; Лемма\par
\noindent\emph{Ukrajinsky / Ukrainian:} допоміжне твердження; Лема\par
\smallskip\noindent\emph{Razlog:} Pomočny teorem jest zadrženy kak prozračny opis lokalnogo pomočnogo rezultata. Razlogi upominajut moderny naziv lema v ruskom i ukrajinskom i ostavjajut medžuslovjansky izbor otkrytym dlja odbornogo sravnjenja. To ne jest dokaz, že specialist jeho odobril ili že lemma objektivno jest menje ustaljena v cělom jezyku. Velika početna bukva pokazuje funkciju zaglavja, ne drugy leksem. Uporědženje raznyh variantov mora sohraniti izvorne motivacije i jasnu rolju rezultata v dokazu.\par
\noindent\emph{Reason:} Pomočny teorem is retained as a transparent description of a local auxiliary result. The reasons mention modern lemma terminology in Russian and Ukrainian and leave the Interslavic choice open to specialist comparison. This is not evidence that a specialist approved it, or that lemma is objectively less established throughout the language. Initial capitalization marks a heading function, not a separate lexeme. Reconciling the variants must preserve the original motivations and the result's clear role in the proof.\par
\noindent{\footnotesize Izvorne zapisy / Original records: HIST-g079-060, HIST-g115-005.}\par
\end{minipage}\par
\par\medskip\Needspace{7\baselineskip}
\noindent\begin{minipage}{\linewidth}
\hypertarget{TERM-0820}{}
\noindent\textbf{pomoćny teorem o algebraično cěloj zavisnosti}\quad{\footnotesize\texttt{TERM-0820}}\par
\noindent\emph{Iztočnik / Source:} Hilfssatz über algebraisch-ganze Abhängigkeit\par
\noindent\emph{Rusky / Russian:} вспомогательная теорема об алгебраически целой зависимости\par
\noindent\emph{Ukrajinsky / Ukrainian:} допоміжна теорема про алгебраїчно цілу залежність\par
\smallskip\noindent\emph{Razlog:} Pomoćny teorem o algebraično cěloj zavisnosti zadržuje pojem zavisnosti v imenovanju odděla raboty 6. Sohranjena motivacija razlikuje zaglavje od rečeničnoj konstrukcije v těle dokaza, kde odnos možno vyraziti kak byti cěly nad danoju oblastju. To ne sut dva različne matematične pojmy: nazov i skladnja služe različnym funkcijam. Cěla zavisnost ne znači proizvoljnu algebraičnu ili linearnu zavisnost. Po osnovnom pravopisu pomoćny odgovarja pomočny; to ne měnjaje uslovja integralnosti.\par
\noindent\emph{Reason:} Pomoćny teorem o algebraično cěloj zavisnosti preserves dependence in the name of the section in Paper 6. The motivation distinguishes the heading from sentence construction in the proof, where the relation can be expressed as being integral over a specified domain. These are not two mathematical concepts: name and syntax serve different functions. Integral dependence does not mean arbitrary algebraic or linear dependence. Basic spelling changes pomoćny to pomočny without changing the integrality condition.\par
\noindent{\footnotesize Izvorne zapisy / Original records: HIST-g051-001.}\par
\end{minipage}\par
\par\medskip\Needspace{7\baselineskip}
\noindent\begin{minipage}{\linewidth}
\hypertarget{TERM-0821}{}
\noindent\textbf{pomocne prěměnne}\quad{\footnotesize\texttt{TERM-0821}}\par
\noindent\emph{Iztočnik / Source:} Hilfsvariablen\par
\noindent\emph{Rusky / Russian:} вспомогательные переменные\par
\noindent\emph{Ukrajinsky / Ukrainian:} допоміжні змінні\par
\smallskip\noindent\emph{Razlog:} Pomocne prěměnne imenuje proměnne uvedene za posrědny račun, ne novu klasu matematičnyh objektov. Rusky i ukrajinsky prěvod tak samo vyražajut pomočnu funkciju. Kratky historijny razlog naziva izraz prozračnym, ale ne podava sravniteljne ocěnjenje čitateljev. Aktualny pasaž raboty 2 govori o pomočnyh prěměnnyh, ktore se identifikujut na koncu računa. Razlika pomocne/pomočne zato ostavaje slovotvorno-pravopisno pytanje, odvojeno od dobro oprěděljenoj funkcije proměnnyh.\par
\noindent\emph{Reason:} Pomocne prěměnne names variables introduced for an intermediate calculation, not a new class of mathematical objects. The Russian and Ukrainian versions likewise express their auxiliary role. The brief historical reason calls the expression transparent but supplies no comparative reader assessment. The present passage in Paper 2 describes auxiliary variables identified at the end of the calculation. The distinction between pomocne and pomočne therefore remains a word-formation and spelling question, separate from the variables’ clearly specified role.\par
\noindent{\footnotesize Izvorne zapisy / Original records: HIST-g004-008.}\par
\smallskip\noindent\emph{Sravnjenje s tekstom:} Dva pasaža raboty 2 ukazyvajut funkciju pomočnyh proměnnyh: posrědny račun i pozdnějša identifikacija. Aktualno pomočnyh ma č; ranějša pomocne i pomoćne ostavajut različne zapisane formy, ne tri različne računske pojmy.\par
\noindent\emph{Textual comparison:} Two passages in Paper 2 identify the variables’ role in intermediate calculation and later identification. Present pomočnyh has č; earlier pomocne and pomoćne remain distinct recorded forms, not three different computational concepts.\par
\noindent{\footnotesize TEXT-COMP-88; pasaže / passages: 435, 974; \texttt{CORPUS\_FORM\_COMPARISONS.json}.}\par
\end{minipage}\par
\par\medskip\Needspace{7\baselineskip}
\noindent\begin{minipage}{\linewidth}
\hypertarget{TERM-0822}{}
\noindent\textbf{pomoćne prěměnne}\quad{\footnotesize\texttt{TERM-0822}}\par
\noindent\emph{Iztočnik / Source:} Hilfsvariablen\par
\noindent\emph{Rusky / Russian:} вспомогательные переменные\par
\noindent\emph{Ukrajinsky / Ukrainian:} допоміжні змінні\par
\smallskip\noindent\emph{Razlog:} Pomoćne prěměnne prodolžuje termin iz §2 v §6 toj samoj raboty. Kontinuitetnost jest zapisany razlog: čitatelj ma razpoznati to samo računsko srědstvo pri njegovom daljšem upotrěbljenju. Osnovny pravopis može zaměniti ć s č, ale to pravilo samo ne izjednačava ranějše pomocne s pomočne. Fraza ne trěbuje nove definicije pri vsakom pojavjenju; trěbuje prověrku, že proměnne imajut tu samu pomočnu funkciju.\par
\noindent\emph{Reason:} Pomoćne prěměnne carries the term from §2 into §6 of the same paper. Continuity is the recorded reason: readers should recognize the same computational device when it is used again. Basic spelling can replace ć with č, but that rule alone does not equate the earlier pomocne with pomočne. The phrase does not need a new definition at every occurrence; it needs confirmation that the variables perform the same auxiliary role.\par
\noindent{\footnotesize Izvorne zapisy / Original records: HIST-g007-008.}\par
\smallskip\noindent\emph{Sravnjenje s tekstom:} Dva pasaža raboty 2 ukazyvajut funkciju pomočnyh proměnnyh: posrědny račun i pozdnějša identifikacija. Aktualno pomočnyh ma č; ranějša pomocne i pomoćne ostavajut različne zapisane formy, ne tri različne računske pojmy.\par
\noindent\emph{Textual comparison:} Two passages in Paper 2 identify the variables’ role in intermediate calculation and later identification. Present pomočnyh has č; earlier pomocne and pomoćne remain distinct recorded forms, not three different computational concepts.\par
\noindent{\footnotesize TEXT-COMP-88; pasaže / passages: 435, 974; \texttt{CORPUS\_FORM\_COMPARISONS.json}.}\par
\end{minipage}\par
\par\medskip\Needspace{7\baselineskip}
\noindent\begin{minipage}{\linewidth}
\hypertarget{TERM-0823}{}
\noindent\textbf{homogena forma}\quad{\footnotesize\texttt{TERM-0823}}\par
\noindent\emph{Iztočnik / Source:} homogene Form\par
\noindent\emph{Rusky / Russian:} однородная форма\par
\noindent\emph{Ukrajinsky / Ukrainian:} однорідна форма\par
\smallskip\noindent\emph{Razlog:} Homogena forma sohranjaje internacionalno imenovanje iz teorije invariantov, dokolě rusko i ukrajinsko sravnjenje upotrěbjajut domače pridavniki. V toj odluci jest ješče drugy krok: historijno i-ter Dimension vyražaje se kako stepen i. To objasnjava matematičny smysl, ne proměnjaje prostorovu dimenziju v stepen. Slovnik potvrđaje osnovu homogenny, s dvoma n, i takože alternativu jednorodny. Zapis homogena zato ne može byti opravdany samo obščim tvrdženjem o internacionalnosti; leksičny izbor i točna forma pridavnika sut dva različne pytanja.\par
\noindent\emph{Reason:} Homogena forma retains an international invariant-theory term where the Russian and Ukrainian comparanda use native adjectives. The decision also makes a separate interpretive move: historical i-ter Dimension is rendered as degree i. This explains the mathematical sense rather than turning a spatial dimension into a degree. The dictionary attests homogenny, with double n, as well as the alternative jednorodny. Homogena therefore cannot be justified solely by an appeal to international vocabulary: lexical selection and the adjective’s exact form are distinct questions.\par
\noindent{\footnotesize Izvorne zapisy / Original records: HIST-g045-005.}\par
\smallskip\noindent\emph{Sravnjenje s tekstom:} Medžuslovjansky slovnik podava homogenny (19718) i jednorodny (21312), homogennosť (19717) i jednorodnosť (21311), i homogenizovati (23412). To podporuje obe leksične možnosti, ne jedinu obveznu formu. Čislove polja slovnika ne sut težiny historijnogo izbora v tom prěvodu. Tekst raboty 6 objasnjava uvedenje $x_0$ pri homogenizaciji; rabota 21 govori o formě f(dx) s nenulovym determinantom drugih izvodnyh. Te različne konstrukcije potrěbujut konkretnu matematičnu definiciju, ne samo sličnu osnovu. Slovnikovo dvojno n trěba odlišiti od padežnogo i brojovogo soglasovanja. To jest nyněšnje sravnjenje s točnoj kopijeju slovnika, ne tvrdženje o ranějšem konsultovanju.\par
\noindent\emph{Textual comparison:} The Interslavic Dictionary lists homogenny (19718) and jednorodny (21312), homogennosť (19717) and jednorodnosť (21311), and homogenizovati (23412). This supports both lexical possibilities, not a uniquely compulsory form. The dictionary’s numerical fields are not historical selection weights for this translation. Paper 6 explains the introduction of $x_0$ in homogenization; Paper 21 treats f(dx) with a nonzero determinant of second derivatives. These constructions need their particular mathematical definitions, not merely a related stem. The dictionary’s double n must also be distinguished from agreement in case and number. This is a present comparison with a fixed dictionary copy, not a claim of earlier consultation.\par
\noindent{\footnotesize TEXT-COMP-89; pasaže / passages: 6165, 14063, 14065, 14067; \texttt{CORPUS\_FORM\_COMPARISONS.json}.}\par
\end{minipage}\par
\par\medskip\Needspace{7\baselineskip}
\noindent\begin{minipage}{\linewidth}
\hypertarget{TERM-0824}{}
\noindent\textbf{homogenna forma}\quad{\footnotesize\texttt{TERM-0824}}\par
\noindent\emph{Iztočnik / Source:} homogene Form\par
\noindent\emph{Rusky / Russian:} однородная форма\par
\noindent\emph{Ukrajinsky / Ukrainian:} однорідна форма\par
\smallskip\noindent\emph{Razlog:} Homogenna forma zadržuje internacionalnu osnovu homogenny, potvrđenu v slovniku. V korennom argumentu raboty 9 tekst sravnjava stepen $\nu\chi$ formy po vybranyh η s ukupnym stepenem λ polinoma po η. Koeficienty samy mogut zavisěti od η; zato homogennost po vybranyh argumentah ne trěba směšati s tvrdženjem, že cěly izraz ima točno stepen $\nu\chi$ po vsih proměnnyh. Ta razlika objasnjava, začto dokaz dava neravnost $\lambda\geq\nu\chi$.\par
\noindent\emph{Reason:} Homogenna forma retains the dictionary-attested international stem homogenny. In Paper 9’s root argument, the text compares the form’s degree $\nu\chi$ in selected η arguments with the polynomial’s total degree λ in η. The coefficients may themselves depend on η; homogeneity in the selected arguments must not be confused with a claim that the whole expression has exactly degree $\nu\chi$ in all variables. This distinction explains why the argument gives the inequality $\lambda\geq\nu\chi$.\par
\noindent{\footnotesize Izvorne zapisy / Original records: HIST-g061-008.}\par
\smallskip\noindent\emph{Sravnjenje s tekstom:} V §3 raboty 9 $F_\nu$ i $G_{\nu\chi}$ sut formy po vybranyh η, ale jih koeficienty takože zaviset od η; tekst zato zaključaje $\lambda\geq\nu\chi$. V §9 ta sama rabota upotrěbjaje historijno dimenzija za stepen homogenyh česti $a_i(\eta)$. Ta dva pasaža pokazyvajut, začto trěba izrično ukazati argumenty i smysl stepena. Modernizacija slova Dimension ne jest opča zaměna vsakego pojavjenja slova dimenzija.\par
\noindent\emph{Textual comparison:} In Paper 9, §3, $F_\nu$ and $G_{\nu\chi}$ are forms in selected η arguments, but their coefficients also depend on η; the text consequently concludes $\lambda\geq\nu\chi$. In §9, the same paper uses historical dimension for the degree of the homogeneous components $a_i(\eta)$. These passages show why the variables and the meaning of degree must be explicit. Modernizing Dimension is not a general replacement of every occurrence of dimension.\par
\noindent{\footnotesize TEXT-COMP-90; pasaže / passages: 7345, 7347, 7348, 7361, 7362, 7363, 7364, 8190, 8191, 8192, 8193, 8194; \texttt{CORPUS\_FORM\_COMPARISONS.json}.}\par
\end{minipage}\par
\par\medskip\Needspace{7\baselineskip}
\noindent\begin{minipage}{\linewidth}
\hypertarget{TERM-0825}{}
\noindent\textbf{homogenna forma $f(dx)$}\quad{\footnotesize\texttt{TERM-0825}}\par
\noindent\emph{Iztočnik / Source:} homogene Form $f(dx)$\par
\noindent\emph{Rusky / Russian:} однородная форма $f(dx)$\par
\noindent\emph{Ukrajinsky / Ukrainian:} однорідна форма $f(dx)$\par
\smallskip\noindent\emph{Razlog:} Homogenna forma f(dx) nazyvaje formu po diferencialah dx; nejde o proizvoljno vyrovnanje teksta ili o identičnost vsih koeficientov. Historijny izbor daje prednost internacionalnoj osnově radi očekovanoj poznatosti v slovjanskom matematičnom registru. Slovnik nezavisno potvrđaje homogenny, ale takože jednorodny; ne dokazuje, že prva možnost jest vsim čitateljam razumlivějša. Aktualny pasaž spojuje formu s variacionnoju zadačeju i izričnym nenulovym determinantom drugih izvodnyh. Kratko imenovanje ne zaměnjaje togo uslovja.\par
\noindent\emph{Reason:} Homogenna forma f(dx) names a form in the differentials dx, not general uniformity or equality of all coefficients. The historical choice favours the international stem for its expected familiarity in Slavic mathematical registers. The dictionary independently attests homogenny, but also jednorodny; it does not establish that every reader understands the former more readily. The present passage connects the form with a variational problem and an explicit nonzero determinant of second derivatives. The short name does not replace that hypothesis.\par
\noindent{\footnotesize Izvorne zapisy / Original records: HIST-g084-004.}\par
\smallskip\noindent\emph{Sravnjenje s tekstom:} Medžuslovjansky slovnik podava homogenny (19718) i jednorodny (21312), homogennosť (19717) i jednorodnosť (21311), i homogenizovati (23412). To podporuje obe leksične možnosti, ne jedinu obveznu formu. Čislove polja slovnika ne sut težiny historijnogo izbora v tom prěvodu. Tekst raboty 6 objasnjava uvedenje $x_0$ pri homogenizaciji; rabota 21 govori o formě f(dx) s nenulovym determinantom drugih izvodnyh. Te različne konstrukcije potrěbujut konkretnu matematičnu definiciju, ne samo sličnu osnovu. Slovnikovo dvojno n trěba odlišiti od padežnogo i brojovogo soglasovanja. To jest nyněšnje sravnjenje s točnoj kopijeju slovnika, ne tvrdženje o ranějšem konsultovanju.\par
\noindent\emph{Textual comparison:} The Interslavic Dictionary lists homogenny (19718) and jednorodny (21312), homogennosť (19717) and jednorodnosť (21311), and homogenizovati (23412). This supports both lexical possibilities, not a uniquely compulsory form. The dictionary’s numerical fields are not historical selection weights for this translation. Paper 6 explains the introduction of $x_0$ in homogenization; Paper 21 treats f(dx) with a nonzero determinant of second derivatives. These constructions need their particular mathematical definitions, not merely a related stem. The dictionary’s double n must also be distinguished from agreement in case and number. This is a present comparison with a fixed dictionary copy, not a claim of earlier consultation.\par
\noindent{\footnotesize TEXT-COMP-89; pasaže / passages: 6165, 14063, 14065, 14067; \texttt{CORPUS\_FORM\_COMPARISONS.json}.}\par
\end{minipage}\par
\par\medskip\Needspace{7\baselineskip}
\noindent\begin{minipage}{\linewidth}
\hypertarget{TERM-0826}{}
\noindent\textbf{homogenny formy}\quad{\footnotesize\texttt{TERM-0826}}\par
\noindent\emph{Iztočnik / Source:} homogene Formen\par
\noindent\emph{Rusky / Russian:} однородные формы\par
\noindent\emph{Ukrajinsky / Ukrainian:} однорідні форми\par
\smallskip\noindent\emph{Razlog:} Homogenny formy bylo izbrano radi kontinuitetnosti s rabotami 3–5. To jest razlog za sohranjanje poznatoj terminologije, ale ne polno gramatično opravdanje zapisa. Slovnik podava homogenny kako osnovnu formu pridavnika; imenovanje form v množině potrěbuje so toju množinoju soglasovanu formu. Tu se zato čuva historijny zapis, ne proglašaje se samostojnym novym terminom. Sravnjenje s homogenna forma i homogene formy razděljuje leksičnu osnovu, broj i pravopis.\par
\noindent\emph{Reason:} Homogenny formy was chosen for continuity with Papers 3–5. That explains retention of familiar terminology but does not fully justify the recorded grammar. The dictionary gives homogenny as the adjective’s citation form; a plural noun phrase needs the appropriate agreeing form. The historical spelling is therefore preserved here, not promoted as a separate new term. Comparison with homogenna forma and homogene formy distinguishes lexical stem, number, and spelling.\par
\noindent{\footnotesize Izvorne zapisy / Original records: HIST-g042-013.}\par
\smallskip\noindent\emph{Sravnjenje s tekstom:} Medžuslovjansky slovnik podava homogenny (19718) i jednorodny (21312), homogennosť (19717) i jednorodnosť (21311), i homogenizovati (23412). To podporuje obe leksične možnosti, ne jedinu obveznu formu. Čislove polja slovnika ne sut težiny historijnogo izbora v tom prěvodu. Tekst raboty 6 objasnjava uvedenje $x_0$ pri homogenizaciji; rabota 21 govori o formě f(dx) s nenulovym determinantom drugih izvodnyh. Te različne konstrukcije potrěbujut konkretnu matematičnu definiciju, ne samo sličnu osnovu. Slovnikovo dvojno n trěba odlišiti od padežnogo i brojovogo soglasovanja. To jest nyněšnje sravnjenje s točnoj kopijeju slovnika, ne tvrdženje o ranějšem konsultovanju.\par
\noindent\emph{Textual comparison:} The Interslavic Dictionary lists homogenny (19718) and jednorodny (21312), homogennosť (19717) and jednorodnosť (21311), and homogenizovati (23412). This supports both lexical possibilities, not a uniquely compulsory form. The dictionary’s numerical fields are not historical selection weights for this translation. Paper 6 explains the introduction of $x_0$ in homogenization; Paper 21 treats f(dx) with a nonzero determinant of second derivatives. These constructions need their particular mathematical definitions, not merely a related stem. The dictionary’s double n must also be distinguished from agreement in case and number. This is a present comparison with a fixed dictionary copy, not a claim of earlier consultation.\par
\noindent{\footnotesize TEXT-COMP-89; pasaže / passages: 6165, 14063, 14065, 14067; \texttt{CORPUS\_FORM\_COMPARISONS.json}.}\par
\end{minipage}\par
\par\medskip\Needspace{7\baselineskip}
\noindent\begin{minipage}{\linewidth}
\hypertarget{TERM-0827}{}
\noindent\textbf{homogene formy / nehomogene polinomy}\quad{\footnotesize\texttt{TERM-0827}}\par
\noindent\emph{Iztočnik / Source:} homogene Formen / inhomogene Polynome\par
\noindent\emph{Rusky / Russian:} однородные формы / неоднородные полиномы\par
\noindent\emph{Ukrajinsky / Ukrainian:} однорідні форми / неоднорідні поліноми\par
\smallskip\noindent\emph{Razlog:} Par homogene formy / nehomogene polinomy sohranjaje vidimu svez medžu homogennostju i homogenizacijeju v dokazu. Prednost internacionalnoj osnově jest izrično navedena; razlog ne imenuje konkurenčny domači izraz, ktory byl faktično sravnjen i odbrazen. Slovnikovo jednorodny daje nyně dokumentovanu alternativu, ne dokaz, že byla razgledana raněje. U prěvoda raboty 6 sam tekst objasnjava homogenizaciju uvedenjem dodatnogo neoprěděljenogo $x_0$. Razlika medžu formami i nehomogennymi polinomami zato ima strukturnu funkciju v argumentu.\par
\noindent\emph{Reason:} The pair homogene formy / nehomogene polinomy keeps the relationship between homogeneity and homogenization visible in the proof. Preference for an international stem is explicit, but the reason names no native competitor actually compared and rejected. Dictionary-attested jednorodny supplies a documented alternative now, not evidence that it was considered earlier. In the translation of Paper 6, the text itself explains homogenization by introducing an additional indeterminate $x_0$. The distinction between forms and inhomogeneous polynomials therefore has a structural role in the argument.\par
\noindent{\footnotesize Izvorne zapisy / Original records: HIST-g051-006.}\par
\smallskip\noindent\emph{Sravnjenje s tekstom:} Medžuslovjansky slovnik podava homogenny (19718) i jednorodny (21312), homogennosť (19717) i jednorodnosť (21311), i homogenizovati (23412). To podporuje obe leksične možnosti, ne jedinu obveznu formu. Čislove polja slovnika ne sut težiny historijnogo izbora v tom prěvodu. Tekst raboty 6 objasnjava uvedenje $x_0$ pri homogenizaciji; rabota 21 govori o formě f(dx) s nenulovym determinantom drugih izvodnyh. Te različne konstrukcije potrěbujut konkretnu matematičnu definiciju, ne samo sličnu osnovu. Slovnikovo dvojno n trěba odlišiti od padežnogo i brojovogo soglasovanja. To jest nyněšnje sravnjenje s točnoj kopijeju slovnika, ne tvrdženje o ranějšem konsultovanju.\par
\noindent\emph{Textual comparison:} The Interslavic Dictionary lists homogenny (19718) and jednorodny (21312), homogennosť (19717) and jednorodnosť (21311), and homogenizovati (23412). This supports both lexical possibilities, not a uniquely compulsory form. The dictionary’s numerical fields are not historical selection weights for this translation. Paper 6 explains the introduction of $x_0$ in homogenization; Paper 21 treats f(dx) with a nonzero determinant of second derivatives. These constructions need their particular mathematical definitions, not merely a related stem. The dictionary’s double n must also be distinguished from agreement in case and number. This is a present comparison with a fixed dictionary copy, not a claim of earlier consultation.\par
\noindent{\footnotesize TEXT-COMP-89; pasaže / passages: 6165, 14063, 14065, 14067; \texttt{CORPUS\_FORM\_COMPARISONS.json}.}\par
\end{minipage}\par
\par\medskip\Needspace{7\baselineskip}
\noindent\begin{minipage}{\linewidth}
\hypertarget{TERM-0828}{}
\noindent\textbf{homogene formy / nehomogenny polinomy}\quad{\footnotesize\texttt{TERM-0828}}\par
\noindent\emph{Iztočnik / Source:} homogene Formen / inhomogene Polynome\par
\noindent\emph{Rusky / Russian:} однородные формы / неоднородные многочлены\par
\noindent\emph{Ukrajinsky / Ukrainian:} однорідні форми / неоднорідні многочлени\par
\smallskip\noindent\emph{Razlog:} Homogene formy / nehomogenny polinomy iz raboty 20 sleduje isti terminologičny princip kak TERM-0827: povezana imenovanja imajut pomagati čitatelju slediti prěhody medžu homogennym i nehomogennym slučajem. Rusko i ukrajinsko mnogočlen sut sravniteljne imenovanja polinoma, ne dokaz da je jedno slovo izključeno iz medžuslovjanskogo. Različne zapisy pridavnikov v samoj parě ne sut obosnovane tym semantičnym razlogom. Jej gramatičnu skladnost trěba ocěniti odvojeno od izbora internacionalnoj leksičnoj osnovy.\par
\noindent\emph{Reason:} Homogene formy / nehomogenny polinomy in Paper 20 follows the same terminological principle as TERM-0827: related names should help readers follow transitions between homogeneous and inhomogeneous cases. The Russian and Ukrainian mnogočlen forms are comparanda for polynomial, not evidence that a word is excluded from Interslavic. The adjective spellings within the pair are not justified by that semantic rationale. Their grammatical consistency must be assessed separately from the choice of an international lexical stem.\par
\noindent{\footnotesize Izvorne zapisy / Original records: HIST-g083-007.}\par
\smallskip\noindent\emph{Sravnjenje s tekstom:} Medžuslovjansky slovnik podava homogenny (19718) i jednorodny (21312), homogennosť (19717) i jednorodnosť (21311), i homogenizovati (23412). To podporuje obe leksične možnosti, ne jedinu obveznu formu. Čislove polja slovnika ne sut težiny historijnogo izbora v tom prěvodu. Tekst raboty 6 objasnjava uvedenje $x_0$ pri homogenizaciji; rabota 21 govori o formě f(dx) s nenulovym determinantom drugih izvodnyh. Te različne konstrukcije potrěbujut konkretnu matematičnu definiciju, ne samo sličnu osnovu. Slovnikovo dvojno n trěba odlišiti od padežnogo i brojovogo soglasovanja. To jest nyněšnje sravnjenje s točnoj kopijeju slovnika, ne tvrdženje o ranějšem konsultovanju.\par
\noindent\emph{Textual comparison:} The Interslavic Dictionary lists homogenny (19718) and jednorodny (21312), homogennosť (19717) and jednorodnosť (21311), and homogenizovati (23412). This supports both lexical possibilities, not a uniquely compulsory form. The dictionary’s numerical fields are not historical selection weights for this translation. Paper 6 explains the introduction of $x_0$ in homogenization; Paper 21 treats f(dx) with a nonzero determinant of second derivatives. These constructions need their particular mathematical definitions, not merely a related stem. The dictionary’s double n must also be distinguished from agreement in case and number. This is a present comparison with a fixed dictionary copy, not a claim of earlier consultation.\par
\noindent{\footnotesize TEXT-COMP-89; pasaže / passages: 6165, 14063, 14065, 14067; \texttt{CORPUS\_FORM\_COMPARISONS.json}.}\par
\end{minipage}\par
\par\medskip\Needspace{7\baselineskip}
\noindent\begin{minipage}{\linewidth}
\hypertarget{TERM-0829}{}
\noindent\textbf{homogene formy i-toj dimenzije}\quad{\footnotesize\texttt{TERM-0829}}\par
\noindent\emph{Iztočnik / Source:} homogene Formen i-ter Dimension\par
\noindent\emph{Rusky / Russian:} однородные формы i-го измерения\par
\noindent\emph{Ukrajinsky / Ukrainian:} однорідні форми i-го виміру\par
\smallskip\noindent\emph{Razlog:} I-toj dimenzije naměrno sohranjaje historijno Dimension tam, kde ranějša politika často upotrěbjala stepen. Sam zapis priznava to napetost. V rabotě 9, §9, $a_i(\eta)$ označa homogene formy v razloženju po stepenjah, a modul sodrživaje produkty σ-tego stepena. Zato i ne jest dimenzija vektornogo prostora. Slovnik ma pokazati i historijno slovo, i jego lokalny smysl; ne smy prikryti razliku medžu věrnym prěnosom imenovanja i modernizovanym objasnjenjem.\par
\noindent\emph{Reason:} I-toj dimenzije deliberately retains historical Dimension where earlier practice often used degree. The record itself acknowledges the tension. In Paper 9, §9, $a_i(\eta)$ denotes homogeneous forms in a degree expansion, and the module contains products of degree σ. Thus i is not the dimension of a vector space. A dictionary should give both the historical word and its local meaning, without concealing the distinction between retaining the author’s terminology and explaining it in modern terms.\par
\noindent{\footnotesize Izvorne zapisy / Original records: HIST-g067-003.}\par
\smallskip\noindent\emph{Sravnjenje s tekstom:} V §3 raboty 9 $F_\nu$ i $G_{\nu\chi}$ sut formy po vybranyh η, ale jih koeficienty takože zaviset od η; tekst zato zaključaje $\lambda\geq\nu\chi$. V §9 ta sama rabota upotrěbjaje historijno dimenzija za stepen homogenyh česti $a_i(\eta)$. Ta dva pasaža pokazyvajut, začto trěba izrično ukazati argumenty i smysl stepena. Modernizacija slova Dimension ne jest opča zaměna vsakego pojavjenja slova dimenzija.\par
\noindent\emph{Textual comparison:} In Paper 9, §3, $F_\nu$ and $G_{\nu\chi}$ are forms in selected η arguments, but their coefficients also depend on η; the text consequently concludes $\lambda\geq\nu\chi$. In §9, the same paper uses historical dimension for the degree of the homogeneous components $a_i(\eta)$. These passages show why the variables and the meaning of degree must be explicit. Modernizing Dimension is not a general replacement of every occurrence of dimension.\par
\noindent{\footnotesize TEXT-COMP-90; pasaže / passages: 7345, 7347, 7348, 7361, 7362, 7363, 7364, 8190, 8191, 8192, 8193, 8194; \texttt{CORPUS\_FORM\_COMPARISONS.json}.}\par
\end{minipage}\par
\par\medskip\Needspace{7\baselineskip}
\noindent\begin{minipage}{\linewidth}
\hypertarget{TERM-0830}{}
\noindent\textbf{homogene polinomy}\quad{\footnotesize\texttt{TERM-0830}}\par
\noindent\emph{Iztočnik / Source:} homogene Polynome\par
\noindent\emph{Rusky / Russian:} однородные полиномы\par
\noindent\emph{Ukrajinsky / Ukrainian:} однорідні поліноми\par
\smallskip\noindent\emph{Razlog:} Homogene polinomy bylo zadržano radi poznatosti internacionalnoj osnovy i blizosti k citatom. To objasnjava leksičny smer, ne dokazuje obščeslovjansku razumlivost. Slovnik ima i homogenny, i jednorodny; nacionalne sravniteljne formy takože ne ukazyvajut jedinu mogućnost. Pojavjenje izraza v aktualnom korpusu potvrđaje upotrěbljenje, ale ne pravilnost vsakego n, padeža ili kontekstualnogo značenja. Izbor osnovy i proverka gramatičnoj formy zato ostavajut različne česti ocěnjenja.\par
\noindent\emph{Reason:} Homogene polinomy was retained for the international stem’s expected familiarity and its proximity to citations. That explains the lexical direction, not demonstrated pan-Slavic intelligibility. The dictionary has both homogenny and jednorodny, and the national-language comparanda likewise do not establish a unique option. Occurrence in the present corpus establishes use, not the correctness of every n, inflection, or contextual meaning. Choosing the stem and checking its grammatical realization are separate parts of the assessment.\par
\noindent{\footnotesize Izvorne zapisy / Original records: HIST-g214-013.}\par
\smallskip\noindent\emph{Sravnjenje s tekstom:} Medžuslovjansky slovnik podava homogenny (19718) i jednorodny (21312), homogennosť (19717) i jednorodnosť (21311), i homogenizovati (23412). To podporuje obe leksične možnosti, ne jedinu obveznu formu. Čislove polja slovnika ne sut težiny historijnogo izbora v tom prěvodu. Tekst raboty 6 objasnjava uvedenje $x_0$ pri homogenizaciji; rabota 21 govori o formě f(dx) s nenulovym determinantom drugih izvodnyh. Te različne konstrukcije potrěbujut konkretnu matematičnu definiciju, ne samo sličnu osnovu. Slovnikovo dvojno n trěba odlišiti od padežnogo i brojovogo soglasovanja. To jest nyněšnje sravnjenje s točnoj kopijeju slovnika, ne tvrdženje o ranějšem konsultovanju.\par
\noindent\emph{Textual comparison:} The Interslavic Dictionary lists homogenny (19718) and jednorodny (21312), homogennosť (19717) and jednorodnosť (21311), and homogenizovati (23412). This supports both lexical possibilities, not a uniquely compulsory form. The dictionary’s numerical fields are not historical selection weights for this translation. Paper 6 explains the introduction of $x_0$ in homogenization; Paper 21 treats f(dx) with a nonzero determinant of second derivatives. These constructions need their particular mathematical definitions, not merely a related stem. The dictionary’s double n must also be distinguished from agreement in case and number. This is a present comparison with a fixed dictionary copy, not a claim of earlier consultation.\par
\noindent{\footnotesize TEXT-COMP-89; pasaže / passages: 6165, 14063, 14065, 14067; \texttt{CORPUS\_FORM\_COMPARISONS.json}.}\par
\end{minipage}\par
\par\medskip\Needspace{7\baselineskip}
\noindent\begin{minipage}{\linewidth}
\hypertarget{TERM-0831}{}
\noindent\textbf{homogena rezultanta členov najvyššego stepena}\quad{\footnotesize\texttt{TERM-0831}}\par
\noindent\emph{Iztočnik / Source:} homogene Resultante der Glieder höchster Dimension\par
\noindent\emph{Rusky / Russian:} однородная результанта членов наивысшей степени\par
\noindent\emph{Ukrajinsky / Ukrainian:} однорідна результанта членів найвищого степеня\par
\smallskip\noindent\emph{Razlog:} Homogena rezultanta členov najvyššego stepena prěvodit Dimension kak stepen v konkretnom polinomnom kontekstu; rezultanta ostavaje tehničnym imenovanjem. Ne ide o rezultant proizvoljnyh členov, ale o najvyšše homogene česti polinomov. Pasaž raboty 5 razlikuje nenulovu rezultantu, dostatočnu za obstojanje integralnoj bazy, od silnějšego uslovja, že rezultanta jest jedinica koeficientnoj oblasti, potrěbnogo v navedenom argumentu za integralnost prědstavjenja. Prěvod ne može sliti ne-nula i jedinica.\par
\noindent\emph{Reason:} Homogena rezultanta členov najvyššego stepena renders Dimension as degree in this polynomial context while retaining resultant as a technical term. It concerns the polynomials’ highest homogeneous parts, not arbitrary terms. The passage in Paper 5 distinguishes a nonzero resultant, sufficient there for an integral basis, from the stronger unit condition used in the stated argument for integrality of the representation. Nonzero and unit must not be merged in translation.\par
\noindent{\footnotesize Izvorne zapisy / Original records: HIST-g039-018.}\par
\smallskip\noindent\emph{Sravnjenje s tekstom:} V rabotě 5 nenulova rezultanta najvyšših členov daje dostatočno uslovje za integralnu bazu. Slědujuči pasaž osoblivo zahtěvaje jedinicu koeficientnoj oblasti za integralnost prědstavjenja; još dalje avtor naglašaje, že prvo uslovje ne jest potrěbno. Prěvod mora sohranjati vse tri razlikovanja: najvyšše česti, ne-nula protiv jedinice, i dostatočnost protiv potrěbnosti.\par
\noindent\emph{Textual comparison:} In Paper 5 a nonzero resultant of the highest-degree parts gives a sufficient condition for an integral basis. The next passage separately requires a unit of the coefficient domain for integrality of the representation, and the author then stresses that the first condition is not necessary. Translation must preserve all three distinctions: highest-degree parts, nonzero versus unit, and sufficiency versus necessity.\par
\noindent{\footnotesize TEXT-COMP-91; pasaže / passages: 4968, 4973, 4975; \texttt{CORPUS\_FORM\_COMPARISONS.json}.}\par
\end{minipage}\par
\par\medskip\Needspace{7\baselineskip}
\noindent\begin{minipage}{\linewidth}
\hypertarget{TERM-0832}{}
\noindent\textbf{homogenny izobarny invarianty}\quad{\footnotesize\texttt{TERM-0832}}\par
\noindent\emph{Iztočnik / Source:} homogene, isobare Invarianten\par
\noindent\emph{Rusky / Russian:} однородные изобарные инварианты\par
\noindent\emph{Ukrajinsky / Ukrainian:} однорідні ізобарні інваріанти\par
\smallskip\noindent\emph{Razlog:} Homogenny izobarny invarianty čuva dva pridavnika, zato že oni ne nesut jednu i tu samu informaciju. V invariantnoj teoriji homogennost odnosi se k stepenu, a izobarnost k jednakoj važe členov po danoj konvenciji. Historijny razlog naglašaje kratkost internacionalnyh imenovanj; ne dava test njihove razumlivosti. Pasaž raboty 6 razlikuje sistem takyh invariantov od polja, ktoro dobivaje se pri dopuščanju racionalnyh kombinacij i odgovarjajučem ograničenju transformacij. Samo izobara v obščem slovniku ne bi dokazala matematičny smysl izobarny.\par
\noindent\emph{Reason:} Homogenny izobarny invarianty preserves two adjectives because they do not convey the same information. In invariant theory, homogeneity concerns degree and isobaricity concerns equal term weight under the specified convention. The historical reason emphasizes the international terms’ compactness, without an intelligibility test. The passage in Paper 6 distinguishes a system of such invariants from the field obtained with rational combinations and the appropriate restriction on transformations. A general dictionary entry for isobar would not by itself attest the mathematical sense of isobaric.\par
\noindent{\footnotesize Izvorne zapisy / Original records: HIST-g041-016.}\par
\smallskip\noindent\emph{Sravnjenje s tekstom:} Rabota 6 razlikuje sistem homogennyh izobarnyh invariantov od polja racionalnyh invariantov. Rabota 23 definira homogenu oblast cělosti tym, že ona sodrživaje vsaku homogennu čest polinoma, ktory k nej prinaleži. Drugy izraz zato ne označa oblast, v kojoj vsi elementy imajut jedin stepen. Sohranjanje pridavnika bez definicije moglo by zakryti tu razliku.\par
\noindent\emph{Textual comparison:} Paper 6 distinguishes a system of homogeneous isobaric invariants from the field of rational invariants. Paper 23 defines a homogeneous integral domain by its containing each homogeneous component of any polynomial belonging to it. The second expression therefore does not mean a domain whose elements all have a single degree. Retaining the adjective without its definition could conceal the distinction.\par
\noindent{\footnotesize TEXT-COMP-92; pasaže / passages: 5130, 15099; \texttt{CORPUS\_FORM\_COMPARISONS.json}.}\par
\end{minipage}\par
\par\medskip\Needspace{7\baselineskip}
\noindent\begin{minipage}{\linewidth}
\hypertarget{TERM-0833}{}
\noindent\textbf{homogena oblast cělosti}\quad{\footnotesize\texttt{TERM-0833}}\par
\noindent\emph{Iztočnik / Source:} homogener Integritätsbereich\par
\noindent\emph{Rusky / Russian:} однородная область целостности\par
\noindent\emph{Ukrajinsky / Ukrainian:} однорідна область цілісності\par
\smallskip\noindent\emph{Razlog:} Homogena oblast cělosti prodolžuje ranějše imenovanje oblasti, ale pridavnik potrěbuje definiciju iz raboty 23. Tu on znači: ako polinom prinaleži k oblasti, njegove homogene česti takože prinaležet k nej, vsaka odvojeno. To ne znači, že vsaky element oblasti sam jest homogenny polinom jednogo stepena. Ta razlika čuva smysl razloženja i dopušča sumy različnyh stepenjev. Historijna motivacija kontinuitetnosti ne izgovarja tu definiciju; zato ona jest dana v kontekstualnom sravnjenju.\par
\noindent\emph{Reason:} Homogena oblast cělosti continues the earlier name for the domain, but its adjective needs Paper 23’s definition: if a polynomial belongs to the domain, each of its homogeneous components belongs to it separately. This does not mean that every element of the domain is itself homogeneous of a single degree. The distinction preserves decomposition and permits sums of different degrees. The historical continuity rationale does not state that definition, so it is supplied in the contextual comparison.\par
\noindent{\footnotesize Izvorne zapisy / Original records: HIST-g094-005.}\par
\smallskip\noindent\emph{Sravnjenje s tekstom:} Rabota 6 razlikuje sistem homogennyh izobarnyh invariantov od polja racionalnyh invariantov. Rabota 23 definira homogenu oblast cělosti tym, že ona sodrživaje vsaku homogennu čest polinoma, ktory k nej prinaleži. Drugy izraz zato ne označa oblast, v kojoj vsi elementy imajut jedin stepen. Sohranjanje pridavnika bez definicije moglo by zakryti tu razliku.\par
\noindent\emph{Textual comparison:} Paper 6 distinguishes a system of homogeneous isobaric invariants from the field of rational invariants. Paper 23 defines a homogeneous integral domain by its containing each homogeneous component of any polynomial belonging to it. The second expression therefore does not mean a domain whose elements all have a single degree. Retaining the adjective without its definition could conceal the distinction.\par
\noindent{\footnotesize TEXT-COMP-92; pasaže / passages: 5130, 15099; \texttt{CORPUS\_FORM\_COMPARISONS.json}.}\par
\end{minipage}\par
\par\medskip\Needspace{7\baselineskip}
\noindent\begin{minipage}{\linewidth}
\hypertarget{TERM-0834}{}
\noindent\textbf{uslovje homogennosti / homogenizovanje}\quad{\footnotesize\texttt{TERM-0834}}\par
\noindent\emph{Iztočnik / Source:} Homogenitätsvoraussetzung / Homogenisieren\par
\noindent\emph{Rusky / Russian:} условие однородности / гомогенизация\par
\noindent\emph{Ukrajinsky / Ukrainian:} умова однорідності / гомогенізація\par
\smallskip\noindent\emph{Razlog:} Uslovje homogennosti i homogenizovanje razlikujut prědpoloženje od operacije: prvo opisuje svojstvo vhoda, drugo jest postupok, ktory dovoljava razgledati inoj slučaj pomoču homogennogo. Internacionalna osnova zbližava ta imenovanja. Arhivny razlog tvrdi, že ona jest menje dvosmyslena než domači kalk, ale ne nazyvaje konkretno sravnjanu formu ni měru dvosmyslenosti. Slovnik nyně potvrđaje homogennosť i homogenizovati, i takože jednorodnosť. To podkrepljaje tvorjenje povezanoj terminologije, ne dokazuje historijno odbacanje vseh domačih alternativ.\par
\noindent\emph{Reason:} Uslovje homogennosti and homogenizovanje distinguish a hypothesis from an operation: the former describes an input property, the latter a procedure for treating another case through a homogeneous one. The international stem connects the two names. The archival reason claims less ambiguity than a native calque but names neither a specific competitor nor a measure of ambiguity. The dictionary now attests homogennosť and homogenizovati, as well as jednorodnosť. This supports a related terminological family, not a historical rejection of every native alternative.\par
\noindent{\footnotesize Izvorne zapisy / Original records: HIST-g074-016.}\par
\smallskip\noindent\emph{Sravnjenje s tekstom:} Medžuslovjansky slovnik podava homogenny (19718) i jednorodny (21312), homogennosť (19717) i jednorodnosť (21311), i homogenizovati (23412). To podporuje obe leksične možnosti, ne jedinu obveznu formu. Čislove polja slovnika ne sut težiny historijnogo izbora v tom prěvodu. Tekst raboty 6 objasnjava uvedenje $x_0$ pri homogenizaciji; rabota 21 govori o formě f(dx) s nenulovym determinantom drugih izvodnyh. Te različne konstrukcije potrěbujut konkretnu matematičnu definiciju, ne samo sličnu osnovu. Slovnikovo dvojno n trěba odlišiti od padežnogo i brojovogo soglasovanja. To jest nyněšnje sravnjenje s točnoj kopijeju slovnika, ne tvrdženje o ranějšem konsultovanju.\par
\noindent\emph{Textual comparison:} The Interslavic Dictionary lists homogenny (19718) and jednorodny (21312), homogennosť (19717) and jednorodnosť (21311), and homogenizovati (23412). This supports both lexical possibilities, not a uniquely compulsory form. The dictionary’s numerical fields are not historical selection weights for this translation. Paper 6 explains the introduction of $x_0$ in homogenization; Paper 21 treats f(dx) with a nonzero determinant of second derivatives. These constructions need their particular mathematical definitions, not merely a related stem. The dictionary’s double n must also be distinguished from agreement in case and number. This is a present comparison with a fixed dictionary copy, not a claim of earlier consultation.\par
\noindent{\footnotesize TEXT-COMP-89; pasaže / passages: 6165, 14063, 14065, 14067; \texttt{CORPUS\_FORM\_COMPARISONS.json}.}\par
\end{minipage}\par
\par\medskip\Needspace{7\baselineskip}
\noindent\begin{minipage}{\linewidth}
\hypertarget{TERM-0835}{}
\noindent\textbf{homomorfny / modulno-homomorfny / kolcevo-homomorfny}\quad{\footnotesize\texttt{TERM-0835}}\par
\noindent\emph{Iztočnik / Source:} homomorph / modul-homomorph / ring-homomorph\par
\noindent\emph{Rusky / Russian:} гомоморфный / модульно-гомоморфный / кольцево-гомоморфный\par
\noindent\emph{Ukrajinsky / Ukrainian:} гомоморфний / модульно-гомоморфний / кільцево-гомоморфний\par
\smallskip\noindent\emph{Razlog:} Homomorfny zadržuje internacionalnu osnovu; modulno- i kolcevo- ukazyvajut, ktoru strukturu zobraženje čuva. Najvažnějša informacija jest v definiciji raboty 30: cilj jest izčerpany obrazom, i čuvajut se razlika i množenje operatorom. To jest surjektivno zobraženje v današnjem vyraženju, ne libovoljny nesurjektivny homomorfizm v zadany cilj. Prstenovy slučaj dodatno čuva produkt. Historijno kolcevo- ne smy avtomatično prenesti v aktualnu terminologiju prstena; i izbor korenja, i uslovja definicije potrěbujut svoji odvojene opravdanja.\par
\noindent\emph{Reason:} Homomorfny retains the international stem; modulno- and kolcevo- specify which structure the map preserves. The crucial information lies in Paper 30’s definition: the image exhausts the target, and subtraction and multiplication by an operator are preserved. In present terminology this is a surjective map, not an arbitrary nonsurjective homomorphism into a specified target. The ring case also preserves products. Historical kolcevo- cannot automatically be carried into the present prsten terminology: root selection and the definition’s conditions each need their own justification.\par
\noindent{\footnotesize Izvorne zapisy / Original records: HIST-g111-002.}\par
\smallskip\noindent\emph{Sravnjenje s tekstom:} Definicija raboty 30 zahtěvaje, že obraz izčerpava cilj, i izpisuje sohranjanje razlike i operatornogo množenja. Pri prstenah dodava se produkt. Rabota 34 izrično svęzuje obraz s faktorprstenom po dvustrannom idealu i čuva dodatne operatorne uslovja. Nazov homomorfizm sam ne može zaměniti tu informaciju o zobraženju i njegovom cilju.\par
\noindent\emph{Textual comparison:} Paper 30’s definition requires the image to exhaust the target and writes out preservation of subtraction and operator multiplication. Products enter in the ring case. Paper 34 explicitly relates the image to a quotient by a two-sided ideal and retains additional operator conditions. The name homomorphism cannot substitute for this information about the map and its target.\par
\noindent{\footnotesize TEXT-COMP-93; pasaže / passages: 16455, 16457, 16462, 16466, 16482, 20404, 20406, 20412; \texttt{CORPUS\_FORM\_COMPARISONS.json}.}\par
\end{minipage}\par
\par\medskip\Needspace{7\baselineskip}
\noindent\begin{minipage}{\linewidth}
\hypertarget{TERM-0836}{}
\noindent\textbf{homomorfno matrično kolco ... stvorjeno idealom}\quad{\footnotesize\texttt{TERM-0836}}\par
\noindent\emph{Iztočnik / Source:} homomorphe Matrizenring ... durch ein Ideal ... erzeugt\par
\noindent\emph{Rusky / Russian:} гомоморфное матричное кольцо ... порожденное идеалом\par
\noindent\emph{Ukrajinsky / Ukrainian:} гомоморфне матричне кільце ... породжене ідеалом\par
\smallskip\noindent\emph{Razlog:} Fraza homomorfno matrično kolco … stvorjeno idealom sohranjaje ograničenje v avtorskoj primětkě. Klasy prědstavjenij tu vključajut samo matrične prsteny, dobyte iz idealov po navedenoj konstrukciji, ne vse myslime homomorfne matrične obrazy. Pytanje o vsih takyh obrazah avtor izrično ostavlja izvan razgledanja. Aktualno homomorfno matrično prsten pokazuje dodatny problem: pri zaměně kolco na prsten ostalo je srednorodno soglasovanje. To jest gramatičny nedostatok, ne nova razlika medžu matematičnymi objektami.\par
\noindent\emph{Reason:} The phrase homomorfno matrično kolco … stvorjeno idealom preserves a restriction in the author’s footnote. Representation classes here include only matrix rings obtained from ideals by the specified construction, not every conceivable homomorphic matrix image. The author expressly leaves the more general question outside the discussion. Present homomorfno matrično prsten reveals a further problem: replacement of kolco by prsten has left neuter agreement behind. That is a grammatical defect, not a new distinction between mathematical objects.\par
\noindent{\footnotesize Izvorne zapisy / Original records: HIST-g133-006.}\par
\smallskip\noindent\emph{Sravnjenje s tekstom:} Avtorska primětka ne razrěšaje obče pytanje, može li vsak homomorfny matričny obraz byti stvorjen idealom; ona omežuje definiciju klasov do obrazov stvorjenyh po danoj konstrukciji. Aktualny zapis homomorfno matrično prsten ima srednorodne pridavniki pri imenovanju prstena. Slovnikova rodova informacija i pravilo soglasovanja razgledane pri TEXT-COMP-67 pokazujut, začto tu zaměnu imenovanja potrěbno dopolniti gramatičnoju korekcijeju; matematično ograničenje ostavaje odvojeno od njej.\par
\noindent\emph{Textual comparison:} The author’s footnote does not settle whether every homomorphic matrix image can be obtained from an ideal; it restricts the definition of representation classes to images obtained by the specified construction. Present homomorfno matrično prsten leaves neuter adjectives with the ring noun. The dictionary gender evidence and agreement rule discussed at TEXT-COMP-67 show why the noun replacement needs a grammatical correction as well; the mathematical restriction is a separate issue.\par
\noindent{\footnotesize TEXT-COMP-94; pasaže / passages: 17956; \texttt{CORPUS\_FORM\_COMPARISONS.json}.}\par
\end{minipage}\par
\par\medskip\Needspace{7\baselineskip}
\noindent\begin{minipage}{\linewidth}
\hypertarget{TERM-0837}{}
\noindent\textbf{teorem o homomorfizmu za kolca}\quad{\footnotesize\texttt{TERM-0837}}\par
\noindent\emph{Iztočnik / Source:} Homomorphiesatz für Ringe\par
\noindent\emph{Rusky / Russian:} теорема о гомоморфизме для колец\par
\noindent\emph{Ukrajinsky / Ukrainian:} теорема про гомоморфізм для кілець\par
\smallskip\noindent\emph{Razlog:} Teorem o homomorfizmu za kolca prodolžuje imenovanje iz §2, ale omežuje predmet do prstenov. Aktualna teorema govori, že homomorfny obraz jest izomorfny faktorprstenu po dvustrannom idealu; operatorne uslovja ostavajut važny, kogda sut zadane. Nazov zato ne trěba razuměti kak tvrdženje o vsakom zobraženju ili vsakom jednostrannom idealu. Kolca jest historijna leksična forma, dokolě aktualno zaglavje ima prsteny. Ta zaměna ne smy zatemniti neproměnjenu matematičnu strukturu teoremy.\par
\noindent\emph{Reason:} Teorem o homomorfizmu za kolca continues the name used in §2 while restricting its subject to rings. The present theorem states that a homomorphic image is isomorphic to a quotient ring by a two-sided ideal; specified operator conditions remain relevant. Its name must not be read as a claim about every map or every one-sided ideal. Kolca is the historical lexical form, whereas the present heading has prsteny. That replacement must not obscure the theorem’s unchanged mathematical structure.\par
\noindent{\footnotesize Izvorne zapisy / Original records: HIST-g186-001.}\par
\smallskip\noindent\emph{Sravnjenje s tekstom:} Definicija raboty 30 zahtěvaje, že obraz izčerpava cilj, i izpisuje sohranjanje razlike i operatornogo množenja. Pri prstenah dodava se produkt. Rabota 34 izrično svęzuje obraz s faktorprstenom po dvustrannom idealu i čuva dodatne operatorne uslovja. Nazov homomorfizm sam ne može zaměniti tu informaciju o zobraženju i njegovom cilju.\par
\noindent\emph{Textual comparison:} Paper 30’s definition requires the image to exhaust the target and writes out preservation of subtraction and operator multiplication. Products enter in the ring case. Paper 34 explicitly relates the image to a quotient by a two-sided ideal and retains additional operator conditions. The name homomorphism cannot substitute for this information about the map and its target.\par
\noindent{\footnotesize TEXT-COMP-93; pasaže / passages: 16455, 16457, 16462, 16466, 16482, 20404, 20406, 20412; \texttt{CORPUS\_FORM\_COMPARISONS.json}.}\par
\end{minipage}\par
\par\medskip\Needspace{7\baselineskip}
\noindent\begin{minipage}{\linewidth}
\hypertarget{TERM-0839}{}
\noindent\textbf{homomorfizmy iz centra v $\Omega$}\quad{\footnotesize\texttt{TERM-0839}}\par
\noindent\emph{Iztočnik / Source:} Homomorphismen von Z in $\Omega$\par
\noindent\emph{Rusky / Russian:} гомоморфизмы из центра в $\Omega$\par
\noindent\emph{Ukrajinsky / Ukrainian:} гомоморфізми з центра в $\Omega$\par
\smallskip\noindent\emph{Razlog:} Homomorfizmy iz centra v Omega čuva smer zobraženja: iz centra Z v polje Ω. Zaměna slova homomorfizm slovom karakter v toj točkě byla by prědčasna. Pri prědstavjenju centra matricami Eξ korespondencija z ↦ ξ ima stepen jeden; v daljšej formuli karakter na centru jest $n_\nu$-kratny toj korespondencije. Zato internacionalno imenovanje ne jest samo stil: ono pomaga razlikovati zobraženje od slěda prědstavjenja. Pismeno Omega v starom zapisu označaje simbol Ω, ne novu leksičnu jedinicu.\par
\noindent\emph{Reason:} Homomorfizmy iz centra v Omega preserves the map’s direction, from the centre Z to the field Ω. Replacing homomorphism with character at this point would be premature. When the centre is represented by matrices Eξ, the correspondence z ↦ ξ has degree one; in the later formula, the character restricted to the centre is $n_\nu$ times that correspondence. The international term therefore helps distinguish a map from the trace of a representation. Written-out Omega in the earlier record denotes Ω, not a new lexical item.\par
\noindent{\footnotesize Izvorne zapisy / Original records: HIST-g199-005.}\par
\smallskip\noindent\emph{Sravnjenje s tekstom:} §21 raboty 34 povezuje prědstavjenje centra matricami Eξ s korespondencijeju z ↦ ξ. §24 razděljuje tu korespondenciju θ od karaktera χ i dava faktor $n_\nu$: $\chi^{(\nu)}(z)=n_\nu\theta^{(\nu)}(z)$. Pri $n_\nu$ = 1 ony soglašajut se; za děljenje prěz $n_\nu$ daljši tekst prědpoklada karakteristiku nula. Slovnik mora sohranjati te uslovja, ne podavati karakter kak neograničeny sinonim homomorfizma centra.\par
\noindent\emph{Textual comparison:} Paper 34, §21, connects representation of the centre by matrices Eξ with the correspondence z ↦ ξ. Section 24 distinguishes this correspondence θ from the character χ and gives the factor $n_\nu$: $\chi^{(\nu)}(z)=n_\nu\theta^{(\nu)}(z)$. They coincide when $n_\nu$ = 1; for division by $n_\nu$ the following text assumes characteristic zero. The dictionary must retain these conditions instead of offering character as an unrestricted synonym for centre homomorphism.\par
\noindent{\footnotesize TEXT-COMP-96; pasaže / passages: 21952, 21964, 22236, 22238, 22240, 22243, 22245; \texttt{CORPUS\_FORM\_COMPARISONS.json}.}\par
\end{minipage}\par
\par\medskip\Needspace{7\baselineskip}
\noindent\begin{minipage}{\linewidth}
\hypertarget{TERM-0840}{}
\noindent\textbf{homomorfizm centra}\quad{\footnotesize\texttt{TERM-0840}}\par
\noindent\emph{Iztočnik / Source:} Homomorphismus des Zentrums\par
\noindent\emph{Rusky / Russian:} гомоморфизм центра\par
\noindent\emph{Ukrajinsky / Ukrainian:} гомоморфізм центру\par
\smallskip\noindent\emph{Razlog:} Homomorfizm centra zadržuje vezu medžu algebričnoju strukturoju i karakterami, ale ne izjednačava ih. V rabotě 34, §24, formula jest $\chi^{(\nu)}(z)=n_\nu\theta^{(\nu)}(z)$: θ jest homomorfizm centra, χ jest slěd prědstavjenja. V komutativnom slučaju $n_\nu$ = 1, i karakter sam dava homomorfizm. Čislovy faktor i domen centra zato sut sučstveni za imenovanje. Historijny razlog stabilnosti internacionalnogo termina podkrepljaje kontinuitetnost, ne opravdava izpuščenje tih uslovij.\par
\noindent\emph{Reason:} Homomorfizm centra maintains the connection between algebraic structure and characters without identifying them. In Paper 34, §24, the formula is $\chi^{(\nu)}(z)=n_\nu\theta^{(\nu)}(z)$: θ is the centre homomorphism and χ the representation’s trace. In the commutative case $n_\nu$ = 1, the character itself gives a homomorphism. The numerical factor and the domain of the centre are therefore essential. The historical appeal to terminological stability supports continuity, not omission of these conditions.\par
\noindent{\footnotesize Izvorne zapisy / Original records: HIST-g202-006.}\par
\smallskip\noindent\emph{Sravnjenje s tekstom:} §21 raboty 34 povezuje prědstavjenje centra matricami Eξ s korespondencijeju z ↦ ξ. §24 razděljuje tu korespondenciju θ od karaktera χ i dava faktor $n_\nu$: $\chi^{(\nu)}(z)=n_\nu\theta^{(\nu)}(z)$. Pri $n_\nu$ = 1 ony soglašajut se; za děljenje prěz $n_\nu$ daljši tekst prědpoklada karakteristiku nula. Slovnik mora sohranjati te uslovja, ne podavati karakter kak neograničeny sinonim homomorfizma centra.\par
\noindent\emph{Textual comparison:} Paper 34, §21, connects representation of the centre by matrices Eξ with the correspondence z ↦ ξ. Section 24 distinguishes this correspondence θ from the character χ and gives the factor $n_\nu$: $\chi^{(\nu)}(z)=n_\nu\theta^{(\nu)}(z)$. They coincide when $n_\nu$ = 1; for division by $n_\nu$ the following text assumes characteristic zero. The dictionary must retain these conditions instead of offering character as an unrestricted synonym for centre homomorphism.\par
\noindent{\footnotesize TEXT-COMP-96; pasaže / passages: 21952, 21964, 22236, 22238, 22240, 22243, 22245; \texttt{CORPUS\_FORM\_COMPARISONS.json}.}\par
\end{minipage}\par
\par\medskip\Needspace{7\baselineskip}
\noindent\begin{minipage}{\linewidth}
\hypertarget{TERM-0841}{}
\noindent\textbf{hiperkompleksne veličiny}\quad{\footnotesize\texttt{TERM-0841}}\par
\noindent\emph{Iztočnik / Source:} hyperkomplexe Größen\par
\noindent\emph{Rusky / Russian:} гиперкомплексные величины\par
\noindent\emph{Ukrajinsky / Ukrainian:} гіперкомплексні величини\par
\smallskip\noindent\emph{Razlog:} Izraz sohranja historičny jezyk abstrakta iz 1925 goda: avtorica govori o hiperkompleksnyh veličinah, ne prosto o modulah. Rusky i ukrajinsky ekvivalenty podpirajut medžunarodny koren hiperkompleks-, a veličiny oddavaje německe Größen. Glavny argument jest věrnost pojmovomu i historičnomu kontekstu. Zaměna slovom algebry mogla by pomagati dnešnjemu čitatelju, ale ona by izbrisala razliku medžu nazvanjem elementov i nazvanjem struktury. Zato objašnjenje može povezati stary izraz s teorijeju algebr, ne podměnjajuči avtorin tekst. Zapis ne ustanovjava, že izraz jest nova tvorba ili občepriznany medžuslovjansky termin.\par
\noindent\emph{Reason:} The expression preserves the historical language of the 1925 abstract: Noether speaks of hypercomplex quantities, not simply modules. The Russian and Ukrainian equivalents support the international stem hypercomplex-, while veličiny renders Größen. The governing consideration is fidelity to the conceptual and historical setting. Substituting algebry might help a present-day reader, but would erase the distinction between naming elements and naming their structure. An explanation can therefore connect the older expression with algebra theory without replacing the author's words. The record does not establish either a new coinage or general Interslavic acceptance.\par
\noindent{\footnotesize Izvorne zapisy / Original records: HIST-g105-009.}\par
\end{minipage}\par
\par\medskip\Needspace{7\baselineskip}
\noindent\begin{minipage}{\linewidth}
\hypertarget{TERM-0842}{}
\noindent\textbf{hiperkompleksne sistemy, algebry}\quad{\footnotesize\texttt{TERM-0842}}\par
\noindent\emph{Iztočnik / Source:} hyperkomplexe Systeme, Algebren\par
\noindent\emph{Rusky / Russian:} [Bukvy ne sut zachovane / Lettering not recoverable]\par
\noindent\emph{Ukrajinsky / Ukrainian:} [Bukvy ne sut zachovane / Lettering not recoverable]\par
\smallskip\noindent\emph{Razlog:} Dvojica hiperkompleksne sistemy, algebry sohranja starše nazvanje i jegovo bezposrědnje objašnjenje. Čitatelj ne mora vybirati medžu dvoma konkurenčnymi prěvodami: v danom kontekstu vtore nazvanje objasnjaje prvo. To jest razlog sohraniti obě česti, a ne ostaviti samo algebry. Izraz ne znači, že vsaka upotrěba slova sistema v korpusu označaje algebru. Ruske i ukrajinske polja v tom zapisu izgubili sut bukvy i ne mogut služiti kak čitateljne attestacije. Argument o historičnoj kontinuitetě možno rekonstruovati, ale sravnjenje tyh dvou jezykov ne možno dopolniti iz domysla.\par
\noindent\emph{Reason:} The pair hiperkompleksne sistemy, algebry retains an older designation together with its immediate explanatory equivalent. Readers are not being offered two rival translations: in this context the second term explains the first. That is the reason for keeping both instead of retaining only algebry. It does not mean that every occurrence of sistema elsewhere in the corpus denotes an algebra. The Russian and Ukrainian fields in this record have lost their lettering and cannot serve as readable attestations. The historical-continuity argument survives; the missing language comparison cannot be supplied by conjecture.\par
\noindent{\footnotesize Izvorne zapisy / Original records: HIST-g209-001.}\par
\end{minipage}\par
\par\medskip\Needspace{7\baselineskip}
\noindent\begin{minipage}{\linewidth}
\hypertarget{TERM-0843}{}
\noindent\textbf{hiperkompleksna sistema}\quad{\footnotesize\texttt{TERM-0843}}\par
\noindent\emph{Iztočnik / Source:} hyperkomplexes System\par
\noindent\emph{Rusky / Russian:} гиперкомплексная система\par
\noindent\emph{Ukrajinsky / Ukrainian:} гіперкомплексна система\par
\smallskip\noindent\emph{Razlog:} Hiperkompleksna sistema sohranja avtorin tehničny izraz iz raboty 32. Zaměna prostym prsten ili algebrnoju parafrazoju by mogla skryti način, na ktorym avtorica uvodi objekt i jegove prědstavjenja. Vybor opiraje se na kontinuitet termina, ne na tvrdženje, že sistema i prsten vsegda sut ravnoznačne. Žensky rod tu jest čest zachovanogo prěvoda; drugy zapis dava mužsky hiperkompleksny sistem. Te dvě formy trěba sravniti kak gramatične varianty jednogo nazvanja, ne kak dva različne matematične pojmy. Iz zapisanoj motivacije ne slěduje, ktory rod ima prednost v občej medžuslovjanskoj normě.\par
\noindent\emph{Reason:} Hiperkompleksna sistema retains the technical expression used in Paper 32. Replacing it with prsten alone or an algebraic paraphrase could obscure the way Noether introduces the object and its representations. The choice rests on terminological continuity, not on a claim that system and ring are always synonymous. Feminine gender is part of the preserved translation here; another record has masculine hiperkompleksny sistem. These are grammatical variants of one designation, not two mathematical concepts. The recorded justification does not decide which gender should prevail in general Interslavic usage.\par
\noindent{\footnotesize Izvorne zapisy / Original records: HIST-g170-002.}\par
\smallskip\noindent\emph{Sravnjenje s tekstom:} §6 raboty 34 sam dava nazvanje algebra nad poljem $K$, ale to ne oslobodžaje od čitanja definicije. Ona potrěbuje konečny rang, soglasnost prstenovogo množenja s operatorami iz $K$ i jedinično dějstvo elementa 1 iz polja. Odvojena slědujuča věta prědpoklada jediničny element samogo sistema. Te dve jedinične uslovja ne sut identične; sovrěmenno objašnjenje ne smi sliti jih.\par
\noindent\emph{Textual comparison:} Paper 34, §6, itself gives algebra over $K$ as an alternative designation, but that does not remove the need to read the definition. It requires finite rank, compatibility of ring multiplication with scalar operators and identity action by the field's 1. A separate following sentence assumes an identity element in the system itself. These two identity conditions are not identical and must not be merged in a modern explanation.\par
\noindent{\footnotesize TEXT-COMP-97; pasaže / passages: 20266, 20268, 20270, 20271, 20272, 20285; \texttt{CORPUS\_FORM\_COMPARISONS.json}.}\par
\end{minipage}\par
\par\medskip\Needspace{7\baselineskip}
\noindent\begin{minipage}{\linewidth}
\hypertarget{TERM-0844}{}
\noindent\textbf{hiperkompleksny sistem}\quad{\footnotesize\texttt{TERM-0844}}\par
\noindent\emph{Iztočnik / Source:} hyperkomplexes System\par
\noindent\emph{Rusky / Russian:} гиперкомплексная система\par
\noindent\emph{Ukrajinsky / Ukrainian:} гіперкомплексна система\par
\smallskip\noindent\emph{Razlog:} Mužska forma hiperkompleksny sistem imaje osoblivo jasny kontekst v §6 raboty 34: tekst sam pridava nazvanje algebra nad poljem. Definicija potrěbuje konečny rang i soglasnost množenja s dějstvom polja; jediničny operator iz polja ne jest samo soboju tvrdženje o jediničnom elementu same algebry. Slovnik zato sohranja nazvanje i objasnjaje hipotezy, vměsto da avtomatično prěpiše dnešnju definiciju algebry s jedinicoju. Historičny razlog za sistem ne rěšaje rodovuju variantnost sistema/sistem. Ta variantnost ostaje odvojena jezykova zadacha.\par
\noindent\emph{Reason:} The masculine hiperkompleksny sistem has an especially clear setting in Paper 34, §6: the text itself supplies algebra over a field as an alternative designation. Its definition requires finite rank and compatibility of multiplication with scalar action; the field's identity acting as the identity operator does not by itself assert that the algebra has a multiplicative identity. The dictionary therefore retains the designation and explains its hypotheses rather than automatically importing a modern unital definition. The historical reason for system does not settle the gender variation sistema/sistem. That remains a separate linguistic question.\par
\noindent{\footnotesize Izvorne zapisy / Original records: HIST-g184-005, HIST-g188-004.}\par
\smallskip\noindent\emph{Sravnjenje s tekstom:} §6 raboty 34 sam dava nazvanje algebra nad poljem $K$, ale to ne oslobodžaje od čitanja definicije. Ona potrěbuje konečny rang, soglasnost prstenovogo množenja s operatorami iz $K$ i jedinično dějstvo elementa 1 iz polja. Odvojena slědujuča věta prědpoklada jediničny element samogo sistema. Te dve jedinične uslovja ne sut identične; sovrěmenno objašnjenje ne smi sliti jih.\par
\noindent\emph{Textual comparison:} Paper 34, §6, itself gives algebra over $K$ as an alternative designation, but that does not remove the need to read the definition. It requires finite rank, compatibility of ring multiplication with scalar operators and identity action by the field's 1. A separate following sentence assumes an identity element in the system itself. These two identity conditions are not identical and must not be merged in a modern explanation.\par
\noindent{\footnotesize TEXT-COMP-97; pasaže / passages: 20266, 20268, 20270, 20271, 20272, 20285; \texttt{CORPUS\_FORM\_COMPARISONS.json}.}\par
\end{minipage}\par
\par\medskip\Needspace{7\baselineskip}
\noindent\begin{minipage}{\linewidth}
\hypertarget{TERM-0845}{}
\noindent\textbf{najvyšša dimenzija}\quad{\footnotesize\texttt{TERM-0845}}\par
\noindent\emph{Iztočnik / Source:} Höchstdimension\par
\noindent\emph{Rusky / Russian:} наивысшая размерность\par
\noindent\emph{Ukrajinsky / Ukrainian:} найвища розмірність\par
\smallskip\noindent\emph{Razlog:} Najvyšša dimenzija oddavaje Höchstdimension kak termin iz teorije idealov raboty 24. Prědnost ima prozračna konstrukcija so superlativom, ne novo jednoslovno nazvanje. V toj rabotě dimenzija opiraje se na osnovne idealy i individualne normy; za libovoljny ideal jest važno, ktora dimenzija jest največša sredi komponentov. Zato najvyšša tu ne znači proizvoljno največšu dimenziju objekta v kakom-koli drugom smyslu. Jediničny ideal imaje v avtorinoj konvenciji dimenziju −1. To uslovje prinaleži matematičnomu objašnjenju i ne može byti tiho zaměnjeno dnešnjeju konvencijeju.\par
\noindent\emph{Reason:} Najvyšša dimenzija renders Höchstdimension as a term in Paper 24's ideal theory. A transparent superlative construction is preferred to a new single-word designation. In that paper dimension is developed through basic ideals and individual norms; for an arbitrary ideal, the largest dimension among its components matters. Thus highest does not mean whichever notion of dimension happens to be largest in some other setting. The unit ideal has dimension −1 under Noether's convention. That convention belongs in the mathematical explanation and must not be silently replaced with a present-day one.\par
\noindent{\footnotesize Izvorne zapisy / Original records: HIST-g095-016.}\par
\smallskip\noindent\emph{Sravnjenje s tekstom:} Rabota 24 povezuje najvyššu dimenziju s prvym nejediničnym individualnym normom i potom s dimenzijami primarnyh komponentov. Za prosty ili primarny ideal najvyšša dimenzija jest jego dimenzija; za libovoljny ideal važny jest maksimum dimenzij komponentov. Konvencija o jediničnom idealu s dimenzijeju −1 ostaje čest avtorinoj teorije. Prěvod ne trěbuje tiho podměniti tu teoriju inym obščim oprěděljenjem dimenzije.\par
\noindent\emph{Textual comparison:} Paper 24 relates highest dimension to the first nonunit individual norm and then to the dimensions of primary components. For a prime or primary ideal, highest dimension is its dimension; for an arbitrary ideal the maximum of component dimensions matters. The convention giving the unit ideal dimension −1 remains part of Noether's theory. Translation does not call for silently replacing that theory with another general definition of dimension.\par
\noindent{\footnotesize TEXT-COMP-98; pasaže / passages: 15339, 15392, 15509; \texttt{CORPUS\_FORM\_COMPARISONS.json}.}\par
\end{minipage}\par
\par\medskip\Needspace{7\baselineskip}
\noindent\begin{minipage}{\linewidth}
\hypertarget{TERM-0846}{}
\noindent\textbf{najvyšša dimenzija libovoljnogo ideala}\quad{\footnotesize\texttt{TERM-0846}}\par
\noindent\emph{Iztočnik / Source:} Höchstdimension eines beliebigen Ideals\par
\noindent\emph{Rusky / Russian:} наивысшая размерность произвольного идеала\par
\noindent\emph{Ukrajinsky / Ukrainian:} найвища розмірність довільного ідеалу\par
\smallskip\noindent\emph{Razlog:} Doplnjenje libovoljnogo ideala jest sučstveno: ono razširjaje izraz za proste i primarne idealy na ideal, ktory ne mora byti takovy. V rabotě 24 najvyšša dimenzija takogo ideala polučaje se kak maksimum dimenzij primarnyh komponentov, odnosno jih asociovanyh prostyh idealov. Prěvod zato sohranja libovoljnost i razlikuje maksimum od pojedinoj dimenzije prosteho komponenta. Rusky i ukrajinsky paralely podpirajut strukturu izraza, ale same ne dokazujut to matematično pravilo. Pravilo slěduje iz konkrětnogo odlomka, ktory slovnik dava k sravnjenju.\par
\noindent\emph{Reason:} The qualification libovoljnogo ideala is substantive: it extends the expression from prime and primary ideals to an ideal that need be neither. In Paper 24 the highest dimension of such an ideal is the maximum of the dimensions of its primary components, equivalently their associated prime ideals. The translation therefore preserves arbitrary and distinguishes the maximum from the dimension of an individual prime component. The Russian and Ukrainian parallels support the phrase's construction but do not establish that mathematical rule. The rule comes from the specific passage supplied for comparison.\par
\noindent{\footnotesize Izvorne zapisy / Original records: HIST-g098-004.}\par
\smallskip\noindent\emph{Sravnjenje s tekstom:} Rabota 24 povezuje najvyššu dimenziju s prvym nejediničnym individualnym normom i potom s dimenzijami primarnyh komponentov. Za prosty ili primarny ideal najvyšša dimenzija jest jego dimenzija; za libovoljny ideal važny jest maksimum dimenzij komponentov. Konvencija o jediničnom idealu s dimenzijeju −1 ostaje čest avtorinoj teorije. Prěvod ne trěbuje tiho podměniti tu teoriju inym obščim oprěděljenjem dimenzije.\par
\noindent\emph{Textual comparison:} Paper 24 relates highest dimension to the first nonunit individual norm and then to the dimensions of primary components. For a prime or primary ideal, highest dimension is its dimension; for an arbitrary ideal the maximum of component dimensions matters. The convention giving the unit ideal dimension −1 remains part of Noether's theory. Translation does not call for silently replacing that theory with another general definition of dimension.\par
\noindent{\footnotesize TEXT-COMP-98; pasaže / passages: 15339, 15392, 15509; \texttt{CORPUS\_FORM\_COMPARISONS.json}.}\par
\end{minipage}\par
\par\medskip\Needspace{7\baselineskip}
\noindent\begin{minipage}{\linewidth}
\hypertarget{TERM-0847}{}
\noindent\textbf{najvyšši elementarny dělitelj}\quad{\footnotesize\texttt{TERM-0847}}\par
\noindent\emph{Iztočnik / Source:} höchster Elementarteiler\par
\noindent\emph{Rusky / Russian:} высший элементарный делитель\par
\noindent\emph{Ukrajinsky / Ukrainian:} найвищий елементарний дільник\par
\smallskip\noindent\emph{Razlog:} Najvyšši elementarny dělitelj sohranja razliku medžu normoju i vyznačenym elementarnym děliteljem v rabotě 22. Najvyšši označaje položenje v strukturnom poredku, ne največšu čislovu cěnnost. Odlomok povezuje normu s děliteljnostju: norma jest dělima tym děliteljem, a něktora jego stepenj jest dělima normoju. Taka relacija objasnjaje, začto ne trěba zaměnjati izraz prostym največši koeficient ili norma. Slovjansky dělitelj i medžunarodny elementarny sohranjajut razpoznavajemu strukturu. Anglijska iztočna variantnost v TERM-0806 odnosi se k toj samoj terminologičnoj rodině.\par
\noindent\emph{Reason:} Najvyšši elementarny dělitelj preserves the distinction between the norm and the distinguished elementary divisor in Paper 22. Highest denotes a position in the structural ordering, not the largest numerical value. The passage relates norm and divisor through divisibility: the norm is divisible by the divisor, while a power of the divisor is divisible by the norm. This explains why neither largest coefficient nor norm is an adequate substitution. The Slavic dělitelj and international elementarny retain a recognizable construction. The English-source variant in TERM-0806 belongs to the same terminological family.\par
\noindent{\footnotesize Izvorne zapisy / Original records: HIST-g089-005, HIST-g092-047.}\par
\smallskip\noindent\emph{Sravnjenje s tekstom:} V rabotě 22 norma i najvyšši elementarny dělitelj ostajut različne objekty. Norma jest dělima prěz toj dělitelj, a něktora stepenj dělitelja jest dělima normoju. To jest konkreten strukturny svez, ne sravnjenje čislovoj velikosti; on objasnjaje, začto pridevnik najvyšši ne može byti čitany kak največši koeficient.\par
\noindent\emph{Textual comparison:} In Paper 22 the norm and highest elementary divisor remain different objects. The norm is divisible by this divisor, while a power of the divisor is divisible by the norm. This is a specific structural relation, not a comparison of numerical magnitudes; it explains why highest cannot be read as largest coefficient.\par
\noindent{\footnotesize TEXT-COMP-99; pasaže / passages: 14836, 14850; \texttt{CORPUS\_FORM\_COMPARISONS.json}.}\par
\end{minipage}\par
\par\medskip\Needspace{7\baselineskip}
\noindent\begin{minipage}{\linewidth}
\hypertarget{TERM-0848}{}
\noindent\textbf{najvyšši koeficient}\quad{\footnotesize\texttt{TERM-0848}}\par
\noindent\emph{Iztočnik / Source:} höchster Koeffizient\par
\noindent\emph{Rusky / Russian:} старший коэффициент\par
\noindent\emph{Ukrajinsky / Ukrainian:} старший коефіцієнт\par
\smallskip\noindent\emph{Razlog:} Najvyšši koeficient označaje koeficient najvyššej stepeni glavnoj proměnnoj involucijskej formy. Ne ide o koeficient s največšeju absolutnoju cěnnostju. Zachovany razlog povezuje toj izraz s jegovoju upotrěboju v rabotě 6, uključajuči rolju v imeniteljah; kontinuitet važny jest, zato že čitatelj mora slěditi isti objekt skozi dokaz. Nově sravnjene odlomky pokazyvajut i normalizaciju: děljenje najvyššim koeficientom može dati koeficient 1. To ne dopušča bez uslovja proglasiti izvornu formu moničnoju; trěba razlučiti formu prěd normalizacijeju i poslě jej.\par
\noindent\emph{Reason:} Najvyšši koeficient denotes the coefficient of the highest power of the principal variable in the involution form. It is not the coefficient with greatest absolute value. The preserved reason links the expression to its use throughout Paper 6, including its role in denominators; continuity matters because readers must follow the same object through the argument. The passages newly compared also show normalization: division by the leading coefficient can make it 1. This does not permit calling the original form monic without qualification; the form before normalization and the form after it must be distinguished.\par
\noindent{\footnotesize Izvorne zapisy / Original records: HIST-g048-017, HIST-g054-014.}\par
\smallskip\noindent\emph{Sravnjenje s tekstom:} V rabotě 6 koeficient najvyššej stepeni $z$ involucijskej formy jest nezavisny od pomočnih proměnnyh; děljenjem nim možno normalizovati formu. Definicije V i V′ daljše upotrěbjajut uslovja o jediničnom koeficientu. Razlika medžu koeficientom 1 i obratimym elementom koeficientnoj oblasti mora byti izrečena pri objašnjenju, a ne skryta za obščim nazvanjem moničny. Priložene definicije davajut točne oblasti, k ktorym se odnoset.\par
\noindent\emph{Textual comparison:} In Paper 6 the coefficient of the highest power of $z$ in the involution form is independent of the auxiliary variables; dividing by it normalizes the form. Definitions V and V′ later impose identity or unit conditions on the coefficient. The distinction between coefficient 1 and an invertible element of the coefficient domain needs to be explained rather than hidden behind the general label monic. The attached definitions specify their respective domains.\par
\noindent{\footnotesize TEXT-COMP-100; pasaže / passages: 5542, 6006, 6550; \texttt{CORPUS\_FORM\_COMPARISONS.json}.}\par
\end{minipage}\par
\par\medskip\Needspace{7\baselineskip}
\noindent\begin{minipage}{\linewidth}
\hypertarget{TERM-0849}{}
\noindent\textbf{najvyšši koeficient jest jedinica}\quad{\footnotesize\texttt{TERM-0849}}\par
\noindent\emph{Iztočnik / Source:} höchster Koeffizient die Einheit\par
\noindent\emph{Rusky / Russian:} старший коэффициент является единицей\par
\noindent\emph{Ukrajinsky / Ukrainian:} старший коефіцієнт є одиницею\par
\smallskip\noindent\emph{Razlog:} Cěla věta najvyšši koeficient jest jedinica sohranja definicijny uslov, ne samo nazvanje moničny. V definiciji cělosti v rabotě 6 element izpolnjaje uravnenje s vedučim koeficientom 1 i s ostalymi koeficientami iz ukazanoj oblasti. V drugyh odlomkah korpusa koeficient može byti jedinicoju v prstenu, to jest obratimym elementom; to ne mora byti doslovno element 1. Slovnik mora pokazati razliku medžu tymi kontekstami. Historična motivacija za frazu jest kontinuitet s najvyšši koeficient, ale sama fraza bez definicije ne odstranja dvosmyslnost slova jedinica.\par
\noindent\emph{Reason:} The full clause najvyšši koeficient jest jedinica preserves a defining condition rather than replacing it with the label monic. In Paper 6's definition of integrality, an element satisfies an equation with leading coefficient 1 and the remaining coefficients in the specified domain. Elsewhere in the corpus a coefficient may be a unit of a ring, meaning an invertible element; this need not literally be 1. The dictionary must distinguish those settings. The historical reason for the clause is continuity with najvyšši koeficient, but the wording alone does not remove the ambiguity of jedinica.\par
\noindent{\footnotesize Izvorne zapisy / Original records: HIST-g049-011.}\par
\smallskip\noindent\emph{Sravnjenje s tekstom:} V rabotě 6 koeficient najvyššej stepeni $z$ involucijskej formy jest nezavisny od pomočnih proměnnyh; děljenjem nim možno normalizovati formu. Definicije V i V′ daljše upotrěbjajut uslovja o jediničnom koeficientu. Razlika medžu koeficientom 1 i obratimym elementom koeficientnoj oblasti mora byti izrečena pri objašnjenju, a ne skryta za obščim nazvanjem moničny. Priložene definicije davajut točne oblasti, k ktorym se odnoset.\par
\noindent\emph{Textual comparison:} In Paper 6 the coefficient of the highest power of $z$ in the involution form is independent of the auxiliary variables; dividing by it normalizes the form. Definitions V and V′ later impose identity or unit conditions on the coefficient. The distinction between coefficient 1 and an invertible element of the coefficient domain needs to be explained rather than hidden behind the general label monic. The attached definitions specify their respective domains.\par
\noindent{\footnotesize TEXT-COMP-100; pasaže / passages: 5542, 6006, 6550; \texttt{CORPUS\_FORM\_COMPARISONS.json}.}\par
\end{minipage}\par
\par\medskip\Needspace{7\baselineskip}
\noindent\begin{minipage}{\linewidth}
\hypertarget{TERM-0850}{}
\noindent\textbf{vyša forma}\quad{\footnotesize\texttt{TERM-0850}}\par
\noindent\emph{Iztočnik / Source:} höhere Form\par
\noindent\emph{Rusky / Russian:} более высокая форма\par
\noindent\emph{Ukrajinsky / Ukrainian:} вища форма\par
\smallskip\noindent\emph{Razlog:} Vyša forma prěvodi vyšši položenje v historičnom poredku form iz raboty 2. To ne jest svobodny sinonim formy višego stepena: avtorin poredok i polinomny stepen mogu davati različne informacije. Zapis prědpoklada, že definicija togo poredka bude čitatelju dostupna. V sravnjenom odlomku avtorica izrično govori o vyšših formah pri tom samom redu i stepenu; time ograniča možlivo čitanje izraza. Pravopisna forma vyša ostaje zachovana kak historična variantnost, ne kak rěšenje odnošenja k formam vyše i vyšše v srodnyh zapisah.\par
\noindent\emph{Reason:} Vyša forma renders a higher position in Paper 2's historical ordering of forms. It is not a free synonym for a form of higher degree: Noether's ordering and polynomial degree can convey different information. The record assumes that the definition of this ordering remains available to the reader. In the passage compared, Noether explicitly discusses higher forms at the same order and degree, constraining the expression's interpretation. The spelling vyša is retained as a historical variant, not as a resolution of its relation to vyše and vyšše in related records.\par
\noindent{\footnotesize Izvorne zapisy / Original records: HIST-g003-010.}\par
\smallskip\noindent\emph{Sravnjenje s tekstom:} Odlomok iz §17 raboty 2 izrično govori o vyšših formah pri tom samom redu i stepenu. To ograniča čitanje pridevnika: on prinaleži avtorinomu dodatnomu poredku, a ne prosto povećanju polinomnogo stepena. Uvodne definicije reducibilnosti v rabotah 2 i 4 povezuju redukciju s invariantnym faktorom ili s vyššimi formami. Matematičny svez ne dokazuje pravilnost vsih zapisanyh koncovok; za pridevnično-prislovnično razlikovanje relevantna jest gramatika pridevnikov, navedena v GRAM-03.\par
\noindent\emph{Textual comparison:} The passage from Paper 2, §17, explicitly discusses higher forms at the same order and degree. This constrains the modifier: it belongs to Noether's additional ordering, not simply to an increase in polynomial degree. The opening definitions of reducibility in Papers 2 and 4 connect reduction with an invariant factor or higher forms. This mathematical connection does not establish the correctness of every recorded ending; the adjective grammar cited in GRAM-03 is relevant to distinguishing adjective from adverb.\par
\noindent{\footnotesize TEXT-COMP-101; pasaže / passages: 547, 2466, 4910; \texttt{CORPUS\_FORM\_COMPARISONS.json}.}\par
\end{minipage}\par
\par\medskip\Needspace{7\baselineskip}
\noindent\begin{minipage}{\linewidth}
\hypertarget{TERM-0851}{}
\noindent\textbf{vyše formy}\quad{\footnotesize\texttt{TERM-0851}}\par
\noindent\emph{Iztočnik / Source:} höhere Formen\par
\noindent\emph{Rusky / Russian:} высшие формы\par
\noindent\emph{Ukrajinsky / Ukrainian:} вищі форми\par
\smallskip\noindent\emph{Razlog:} Plural vyše formy sohranja nazvanje iz §17 raboty 2, ale jegova gramatična forma ne byla prědstavjena kak definitivna. Izvorna motivacija izrično ostavila pytanje koncovky i pravopisa za razgledanje. Važno jest odděliti dve věči: matematičny smysl vyššego položenja v avtorinom poredku i jezykovu pravilnost pridevnika pri formy. Obča gramatika razlikuje pridevnik vysši od prislovnika vyše; sama nazočnost formy v rabotom tekstu ne jest jezykovo dokazateljstvo. Slovnik zato dava i zapisanu formu, i razlog jejnoj nejistoty, ne vyrabjajuči iz istorije avtomatičnu rekomendaciju.\par
\noindent\emph{Reason:} The plural vyše formy preserves the designation from Paper 2, §17, but its grammatical form was not presented as final. The original justification explicitly left the ending and spelling for review. Two issues must be separated: the mathematical meaning of higher position in Noether's ordering and the linguistic correctness of the modifier of formy. The general grammar distinguishes adjectival vysši from adverbial vyše; occurrence in a working translation is not itself linguistic evidence. The dictionary therefore gives both the recorded form and the reason for uncertainty rather than turning historical usage into an automatic recommendation.\par
\noindent{\footnotesize Izvorne zapisy / Original records: HIST-g018-008.}\par
\smallskip\noindent\emph{Sravnjenje s tekstom:} Odlomok iz §17 raboty 2 izrično govori o vyšših formah pri tom samom redu i stepenu. To ograniča čitanje pridevnika: on prinaleži avtorinomu dodatnomu poredku, a ne prosto povećanju polinomnogo stepena. Uvodne definicije reducibilnosti v rabotah 2 i 4 povezuju redukciju s invariantnym faktorom ili s vyššimi formami. Matematičny svez ne dokazuje pravilnost vsih zapisanyh koncovok; za pridevnično-prislovnično razlikovanje relevantna jest gramatika pridevnikov, navedena v GRAM-03.\par
\noindent\emph{Textual comparison:} The passage from Paper 2, §17, explicitly discusses higher forms at the same order and degree. This constrains the modifier: it belongs to Noether's additional ordering, not simply to an increase in polynomial degree. The opening definitions of reducibility in Papers 2 and 4 connect reduction with an invariant factor or higher forms. This mathematical connection does not establish the correctness of every recorded ending; the adjective grammar cited in GRAM-03 is relevant to distinguishing adjective from adverb.\par
\noindent{\footnotesize TEXT-COMP-101; pasaže / passages: 547, 2466, 4910; \texttt{CORPUS\_FORM\_COMPARISONS.json}.}\par
\end{minipage}\par
\par\medskip\Needspace{7\baselineskip}
\noindent\begin{minipage}{\linewidth}
\hypertarget{TERM-0852}{}
\noindent\textbf{vyšše formy / vyšše samosvijanje}\quad{\footnotesize\texttt{TERM-0852}}\par
\noindent\emph{Iztočnik / Source:} höhere Formen / höhere Faltung in sich\par
\noindent\emph{Rusky / Russian:} более высокие формы / более высокая свертка с самой собой\par
\noindent\emph{Ukrajinsky / Ukrainian:} вищі форми / вище самозгортання\par
\smallskip\noindent\emph{Razlog:} Vyšše formy i vyšše samosvijanje sut svezane s avtorinym poredkom pri simboličnom svijanji v rabotě 4. Zapis objasnjaje, že položenje zavisi od čisel osnovnyh svijanij i od dodatnogo pravila pri ravnosti; ne ide prosto o to, že operacija jest težša ili vrši se pozdněje. Samosvijanje čini vozvratnost operacije vidimoju, vměsto da zamenja ju obščim svijanjem. Srodne izrazy v rabotě 2 podpirajut kontinuitet pojma. Iz toga ne slěduje, že različne zapisane formy vyša, vyše i vyšše sut vse gramatično ravnopravne; to pytanje ostaje odvojeno od matematičnoj definicije.\par
\noindent\emph{Reason:} Vyšše formy and vyšše samosvijanje are tied to Noether's ordering in symbolic contraction in Paper 4. The record explains that position depends on counts of basic contractions and a further rule for ties, not simply on an operation being more difficult or occurring later. Samosvijanje makes self-contraction explicit instead of replacing it with contraction generally. Related expressions in Paper 2 support conceptual continuity. This does not establish that the recorded forms vyša, vyše and vyšše are grammatically interchangeable; that question is separate from the mathematical definition.\par
\noindent{\footnotesize Izvorne zapisy / Original records: HIST-g038-004.}\par
\smallskip\noindent\emph{Sravnjenje s tekstom:} Odlomok iz §17 raboty 2 izrično govori o vyšših formah pri tom samom redu i stepenu. To ograniča čitanje pridevnika: on prinaleži avtorinomu dodatnomu poredku, a ne prosto povećanju polinomnogo stepena. Uvodne definicije reducibilnosti v rabotah 2 i 4 povezuju redukciju s invariantnym faktorom ili s vyššimi formami. Matematičny svez ne dokazuje pravilnost vsih zapisanyh koncovok; za pridevnično-prislovnično razlikovanje relevantna jest gramatika pridevnikov, navedena v GRAM-03.\par
\noindent\emph{Textual comparison:} The passage from Paper 2, §17, explicitly discusses higher forms at the same order and degree. This constrains the modifier: it belongs to Noether's additional ordering, not simply to an increase in polynomial degree. The opening definitions of reducibility in Papers 2 and 4 connect reduction with an invariant factor or higher forms. This mathematical connection does not establish the correctness of every recorded ending; the adjective grammar cited in GRAM-03 is relevant to distinguishing adjective from adverb.\par
\noindent{\footnotesize TEXT-COMP-101; pasaže / passages: 547, 2466, 4910; \texttt{CORPUS\_FORM\_COMPARISONS.json}.}\par
\end{minipage}\par
\par\medskip\Needspace{7\baselineskip}
\noindent\begin{minipage}{\linewidth}
\hypertarget{TERM-0853}{}
\noindent\textbf{vyše razvětvjenje}\quad{\footnotesize\texttt{TERM-0853}}\par
\noindent\emph{Iztočnik / Source:} höhere Verzweigung\par
\noindent\emph{Rusky / Russian:} высшее ветвление\par
\noindent\emph{Ukrajinsky / Ukrainian:} вища розгалуженість\par
\smallskip\noindent\emph{Razlog:} Vyše razvětvjenje jest zapisany prěvod tehničnogo izraza višego razvětvjenja iz raboty 37. Motivacija dava prednost slovjanskomu korenu razvětv-, zamišljenomu kak razumlivy napreč jezykam, prěd užšim nacionalnym kalkom. Konkretna zavržena forma ni čislene vagi ne sut navedene; ne možemo jih dopisati. Vyše mora sohraniti tehnično ograničenje, ne značiti samo bolju količinu običnogo razvětvjenja. Razumlivost korena jest naměr vyboru, ne rezultat izpytanja čitateljev. I gramatično razlučenje pridevnika i prislovnika potrěbuje svoj dokaz, nezavisno od smysla termina.\par
\noindent\emph{Reason:} Vyše razvětvjenje is the recorded rendering of the technical expression higher ramification in Paper 37. The rationale favours the Slavic stem razvětv-, intended to be recognizable across languages, over a narrower national calque. Neither a specific rejected form nor numerical weights are supplied; they cannot be added retrospectively. Higher must retain its technical qualification, not merely mean a greater amount of ordinary ramification. Cross-language intelligibility is the intention behind the choice, not a result of reader testing. The adjective–adverb distinction also needs its own linguistic justification, independently of the term's mathematical meaning.\par
\noindent{\footnotesize Izvorne zapisy / Original records: HIST-g207-002.}\par
\end{minipage}\par
\par\medskip\Needspace{7\baselineskip}
\noindent\begin{minipage}{\linewidth}
\hypertarget{TERM-0854}{}
\noindent\textbf{vyšši modul}\quad{\footnotesize\texttt{TERM-0854}}\par
\noindent\emph{Iztočnik / Source:} höherer Modul\par
\noindent\emph{Rusky / Russian:} более высокий модуль\par
\noindent\emph{Ukrajinsky / Ukrainian:} вищий модуль\par
\smallskip\noindent\emph{Razlog:} Vyšši modul pripada hierarhiji modulov pri redukciji v rabotě 2. Zapisani argument jest kratky: sohraniti nazvanje togo poredka, vměsto da ga raztvorimo v dolgu parafrazu. Ale hierarhija sama ne oprědělja směr vključenja modulov. Iz slova vyšši ne slěduje ni podmodul ni nadmodul bez avtorinoj definicije. Rusky i ukrajinsky paralely pokazyvajut podobnu konstrukciju; one ne ustanovjajut tu definiciju i ne rěšajut pravopis viššego pridevnika. Korpusny slovnik mora zato povezati izraz s konkrětnym redukcijnym kontekstom, a ne dati jegovu univerzalnu zaměnu.\par
\noindent\emph{Reason:} Vyšši modul belongs to the hierarchy of modules used in reduction in Paper 2. The recorded argument is brief: preserve the designation of that ordering instead of dissolving it into a long paraphrase. A hierarchy alone, however, does not determine the direction of module inclusion. Neither submodule nor supermodule follows from higher without Noether's definition. The Russian and Ukrainian parallels show a similar construction; they neither establish that definition nor settle the comparative's spelling. A corpus dictionary must therefore connect the phrase with its particular reduction context rather than offer a universal substitution.\par
\noindent{\footnotesize Izvorne zapisy / Original records: HIST-g006-012.}\par
\end{minipage}\par
\par\medskip\Needspace{7\baselineskip}
\noindent\begin{minipage}{\linewidth}
\hypertarget{TERM-0855}{}
\noindent\textbf{ideal / glavny ideal / prosty ideal / primarny ideal}\quad{\footnotesize\texttt{TERM-0855}}\par
\noindent\emph{Iztočnik / Source:} Ideal / Hauptideal / Primideal / primäres Ideal\par
\noindent\emph{Rusky / Russian:} идеал / главный идеал / простой идеал / первичный идеал\par
\noindent\emph{Ukrajinsky / Ukrainian:} ідеал / головний ідеал / простий ідеал / первинний ідеал\par
\smallskip\noindent\emph{Razlog:} Ta četvorka tvorji svezanu terminologičnu rodinu: ideal, glavny ideal, prosty ideal i primarny ideal. Izbor medžunarodnogo ideal olakšava vez s algebraičnoju literaturoju; glavny imaje sravnjene ruske i ukrajinske paralely. Sučstveno jest ne sliti prosty s primarny. Prosty i primarny odgovarjajut različnym definicijam i mogut igrati različne rolje v razloženju. Proto sohranjanje dvou različnyh pridevnikov ima bolju važnost než površna sličnost slova primarny so slovom prime. Rodina ne jest dokaz nove tvorby vsih četyreh izrazov; ona jest zapis soglasovanogo vyboru v danom prěvodu.\par
\noindent\emph{Reason:} These four expressions form a coordinated terminological family: ideal, glavny ideal, prosty ideal and primarny ideal. International ideal connects the text with algebraic literature; glavny has the recorded Russian and Ukrainian parallels. The crucial requirement is not to merge prosty with primarny. Prime and primary have different definitions and can play different roles in a decomposition. Keeping two distinct modifiers therefore matters more than the superficial resemblance between primarny and English prime. The family is not evidence that all four expressions were newly coined; it records a coordinated choice for this translation.\par
\noindent{\footnotesize Izvorne zapisy / Original records: HIST-g076-009.}\par
\end{minipage}\par
\par\medskip\Needspace{7\baselineskip}
\noindent\begin{minipage}{\linewidth}
\hypertarget{TERM-0856}{}
\noindent\textbf{ideal / modul / modul form}\quad{\footnotesize\texttt{TERM-0856}}\par
\noindent\emph{Iztočnik / Source:} Ideal / Modul / Formenmodul\par
\noindent\emph{Rusky / Russian:} идеал / модуль / модуль форм\par
\noindent\emph{Ukrajinsky / Ukrainian:} ідеал / модуль / модуль форм\par
\smallskip\noindent\emph{Razlog:} Ideal, modul i modul form ne sut tri bezuslovno zaměnjajeme nazvanja. V rabotě 19 primětka povezuje ideal polinomov s staršim obyčajem govoriti o modulah; ona takože daje zatvorenost pri razlikě i množenju polinomom. Prěvod sohranja tu historičnu vez, aby čitatelj razuměl i starše nazvanje, i matematičny objekt. Slovnik mora razlikovati takovu kontekstnu ravnoznačnost od dnešnjego obščego pojma modula. Opis staršego upotrěbljenja ne opravdava globalnu zaměnu vsih modulov idealami. Izraz modul form sohranja imenovanu kategoriju, ne skryvajuči jejne historične granice.\par
\noindent\emph{Reason:} Ideal, modul and modul form are not three unconditionally interchangeable labels. A footnote in Paper 19 connects a polynomial ideal with the older practice of speaking of modules; it also supplies closure under subtraction and multiplication by a polynomial. The translation retains this historical connection so that readers can understand both the older designation and the mathematical object. The dictionary must distinguish such contextual equivalence from the modern general notion of a module. Reporting the older usage does not justify replacing every module with an ideal. Modul form preserves the named category while making its historical limits explicit.\par
\noindent{\footnotesize Izvorne zapisy / Original records: HIST-g081-022, HIST-g082-022.}\par
\smallskip\noindent\emph{Sravnjenje s tekstom:} Primětka raboty 19 objasnjaje starše nazvanje modul i odrazu dava idealno svojstvo. V rabotě 22 definicija idealov iz polinomov sohranja zatvorenost pri razlikě i množenju libovoljnym polinomom; primětka razlikuje jih od modulov linearnyh form. Ta dva odlomky opravdavajut historično objašnjenje, ne globalnu zaměnu slova modul slovom ideal.\par
\noindent\emph{Textual comparison:} The Paper 19 footnote explains the older designation module and immediately supplies the ideal property. Paper 22's definition of polynomial ideals retains closure under subtraction and multiplication by an arbitrary polynomial; its footnote distinguishes them from modules of linear forms. These passages support a historical explanation, not a global replacement of module with ideal.\par
\noindent{\footnotesize TEXT-COMP-102; pasaže / passages: 13859, 14501; \texttt{CORPUS\_FORM\_COMPARISONS.json}.}\par
\end{minipage}\par
\par\medskip\Needspace{7\baselineskip}
\noindent\begin{minipage}{\linewidth}
\hypertarget{TERM-0857}{}
\noindent\textbf{ideal iz polinomov}\quad{\footnotesize\texttt{TERM-0857}}\par
\noindent\emph{Iztočnik / Source:} Ideal aus Polynomen\par
\noindent\emph{Rusky / Russian:} идеал из полиномов\par
\noindent\emph{Ukrajinsky / Ukrainian:} ідеал з поліномів\par
\smallskip\noindent\emph{Razlog:} Ideal iz polinomov byl izbrany vměsto samogo pridevničnogo skračenja, zato že odlomok najprvo oprědělja objekt. Fraza pokazuje, iz čego ideal sostoji se, a slědujuče uslovja davajut zatvorenost pri razlikě i množenju libovoljnym polinomom. V rabotě 22 primětka izrično razlikuje to upotrěbljenje od modulov linearnyh form. V prěvodu ne dostača prepoznati koren ideal: trěba sohraniti polje koeficientov, neoprěděljene i kvantifikaciju v definiciji. Rusky i ukrajinsky izrazy takože podpirajut konstrukciju iz polinomov; to jest jezykova opora, ne zaměna matematičnogo teksta.\par
\noindent\emph{Reason:} Ideal iz polinomov was preferred to an adjectival shortening because the passage first defines the object. The phrase identifies what the ideal consists of, while the following conditions give closure under subtraction and multiplication by an arbitrary polynomial. In Paper 22 a footnote expressly distinguishes this usage from modules of linear forms. Recognizing the stem ideal is not enough: the translation must preserve the coefficient field, indeterminates and quantification in the definition. The Russian and Ukrainian expressions also support a noun construction involving polynomials; that is linguistic support, not a substitute for the mathematics.\par
\noindent{\footnotesize Izvorne zapisy / Original records: HIST-g087-001, HIST-g092-021.}\par
\smallskip\noindent\emph{Sravnjenje s tekstom:} Primětka raboty 19 objasnjaje starše nazvanje modul i odrazu dava idealno svojstvo. V rabotě 22 definicija idealov iz polinomov sohranja zatvorenost pri razlikě i množenju libovoljnym polinomom; primětka razlikuje jih od modulov linearnyh form. Ta dva odlomky opravdavajut historično objašnjenje, ne globalnu zaměnu slova modul slovom ideal.\par
\noindent\emph{Textual comparison:} The Paper 19 footnote explains the older designation module and immediately supplies the ideal property. Paper 22's definition of polynomial ideals retains closure under subtraction and multiplication by an arbitrary polynomial; its footnote distinguishes them from modules of linear forms. These passages support a historical explanation, not a global replacement of module with ideal.\par
\noindent{\footnotesize TEXT-COMP-102; pasaže / passages: 13859, 14501; \texttt{CORPUS\_FORM\_COMPARISONS.json}.}\par
\end{minipage}\par
\par\medskip\Needspace{7\baselineskip}
\noindent\begin{minipage}{\linewidth}
\hypertarget{TERM-0858}{}
\noindent\textbf{baza ideala}\quad{\footnotesize\texttt{TERM-0858}}\par
\noindent\emph{Iztočnik / Source:} Idealbasis\par
\noindent\emph{Rusky / Russian:} базис идеала\par
\noindent\emph{Ukrajinsky / Ukrainian:} база ідеалу\par
\smallskip\noindent\emph{Razlog:} Baza ideala sohranja istorijsko upotrěbljenje bazy za konečnu tvorjaču sovkupnost ideala v rabotě 22. Tu baza ne potrěbuje linearnoj nezavisnosti kak baza vektornogo prostora. Genitiv ideala čini nositelja bazy izričnym, a rusky bazis i ukrajinska baza davajut sravnitelnu oporu. Taka prozračnost ne odstranja polisemičnosti: ekspert mora viděti definiciju, a čitatelj mora byti prěduprědženy o različnyh smyslah bazy. Izraz tvorjača sovkupnost možno upotrěbiti v objašnjenju; ne trěba tiho podměnjati im avtorino historično nazvanje.\par
\noindent\emph{Reason:} Baza ideala preserves the historical use of basis for a finite generating set of an ideal in Paper 22. Here basis need not carry the linear-independence condition of a vector-space basis. The genitive ideala makes the object explicit, while Russian bazis and Ukrainian baza supply comparative support. This transparency does not remove the polysemy: an expert needs the definition, and a reader needs to see the different senses of basis. Generating set can be used in the explanation without silently replacing Noether's historical designation.\par
\noindent{\footnotesize Izvorne zapisy / Original records: HIST-g089-007, HIST-g092-049.}\par
\smallskip\noindent\emph{Sravnjenje s tekstom:} §8 raboty 31 dopušča modulnu bazu s $m\ge n$ elementami. Potrěbne determinanty linearno izražajut se prěz konečno mnogo determinantov iz te bazy; ony tvorjat bazu diskriminantnogo ideala. Slučaj linearno nezavisnoj modulnoj bazy iz $n$ elementov jest razgledany odvojeno. Ta izrična razlika daje čitatelju možlivosť razuměti istoričnu polisemičnost bazy bez izměnjenja avtorinogo nazvanja.\par
\noindent\emph{Textual comparison:} Paper 31, §8, allows a module basis of $m\ge n$ elements. The required determinants are linear combinations of finitely many determinants drawn from that basis, which generate the discriminant ideal. A linearly independent module basis of $n$ elements is considered separately. This explicit distinction lets readers understand the historical polysemy of basis without replacing Noether's designation.\par
\noindent{\footnotesize TEXT-COMP-103; pasaže / passages: 18897; \texttt{CORPUS\_FORM\_COMPARISONS.json}.}\par
\end{minipage}\par
\par\medskip\Needspace{7\baselineskip}
\noindent\begin{minipage}{\linewidth}
\hypertarget{TERM-0859}{}
\noindent\textbf{idealna baza}\quad{\footnotesize\texttt{TERM-0859}}\par
\noindent\emph{Iztočnik / Source:} Idealbasis\par
\noindent\emph{Rusky / Russian:} идеальный базис\par
\noindent\emph{Ukrajinsky / Ukrainian:} ідеальний базис\par
\smallskip\noindent\emph{Razlog:} Idealna baza jest pridevnična variantnost srodnogo izraza baza ideala. Zapis iz raboty 23 podpira kontinuitet medžunarodnyh korenjev ideal i baza; ne prědstavjaje novy matematičny pojem. Pridevnik idealna tu znači naležeča idealu, ne sovršena ili najlepsa. Jeho kompaktnost možno sravniti s jasnějšim genitivom baza ideala, ale zachovany razlog ne dava formalno sravnjenje s čislenymi vagami i ne dokazuje, že jedna forma jest vsegda prednostna. Smysl bazy mora byti vzety iz odlomka; sam pridevnik ne dodava linearnoj nezavisnosti.\par
\noindent\emph{Reason:} Idealna baza is an adjectival variant of baza ideala. The Paper 23 record supports continuity of the international stems ideal and baza; it does not introduce a new mathematical concept. Idealna here means belonging to an ideal, not perfect or best. Its compactness can be compared with the more explicit genitive baza ideala, but the preserved reason provides neither a numerical comparison nor a demonstration that one form should always be preferred. The sense of basis must come from the passage; the adjective does not add linear independence.\par
\noindent{\footnotesize Izvorne zapisy / Original records: HIST-g094-007.}\par
\smallskip\noindent\emph{Sravnjenje s tekstom:} §8 raboty 31 dopušča modulnu bazu s $m\ge n$ elementami. Potrěbne determinanty linearno izražajut se prěz konečno mnogo determinantov iz te bazy; ony tvorjat bazu diskriminantnogo ideala. Slučaj linearno nezavisnoj modulnoj bazy iz $n$ elementov jest razgledany odvojeno. Ta izrična razlika daje čitatelju možlivosť razuměti istoričnu polisemičnost bazy bez izměnjenja avtorinogo nazvanja.\par
\noindent\emph{Textual comparison:} Paper 31, §8, allows a module basis of $m\ge n$ elements. The required determinants are linear combinations of finitely many determinants drawn from that basis, which generate the discriminant ideal. A linearly independent module basis of $n$ elements is considered separately. This explicit distinction lets readers understand the historical polysemy of basis without replacing Noether's designation.\par
\noindent{\footnotesize TEXT-COMP-103; pasaže / passages: 18897; \texttt{CORPUS\_FORM\_COMPARISONS.json}.}\par
\end{minipage}\par
\par\medskip\Needspace{7\baselineskip}
\noindent\begin{minipage}{\linewidth}
\hypertarget{TERM-0860}{}
\noindent\textbf{baza diskriminantnogo ideala}\quad{\footnotesize\texttt{TERM-0860}}\par
\noindent\emph{Iztočnik / Source:} Idealbasis des Diskriminantenideals\par
\noindent\emph{Rusky / Russian:} базис дискриминантного идеала\par
\noindent\emph{Ukrajinsky / Ukrainian:} базис дискримінантного ідеалу\par
\smallskip\noindent\emph{Razlog:} Baza diskriminantnogo ideala sohranja bazu kak avtorino nazvanje i genitivom točno ukazuje objekt. V §8 raboty 31 konečno mnogo determinantov tvoret taku bazu, zato že ostale potrebne determinanty izražajut se jih kombinacijami. Tekst dopušča modulnu bazu s količstvom elementov $m\ge n$ i posebno razgleduje slučaj linearno nezavisnoj bazy iz $n$ elementov. To razlikovanje pokazyvaje, začto obča baza ideala ne mora byti linearno nezavisna. Vybor ne jest samo doslovny kalk: on sohranja termin i istočasno potrěbuje objašnjenja jegovogo istoričnogo smysla.\par
\noindent\emph{Reason:} Baza diskriminantnogo ideala retains basis as Noether's designation and identifies its object precisely through the genitive. In Paper 31, §8, finitely many determinants form such a basis because the other required determinants can be expressed as their combinations. The text allows a module basis with $m\ge n$ elements and separately treats the case of a linearly independent basis of $n$ elements. That distinction shows why an ideal basis need not generally be linearly independent. The choice is not merely a literal calque: it preserves the term while requiring an explanation of its historical sense.\par
\noindent{\footnotesize Izvorne zapisy / Original records: HIST-g158-007.}\par
\smallskip\noindent\emph{Sravnjenje s tekstom:} §8 raboty 31 dopušča modulnu bazu s $m\ge n$ elementami. Potrěbne determinanty linearno izražajut se prěz konečno mnogo determinantov iz te bazy; ony tvorjat bazu diskriminantnogo ideala. Slučaj linearno nezavisnoj modulnoj bazy iz $n$ elementov jest razgledany odvojeno. Ta izrična razlika daje čitatelju možlivosť razuměti istoričnu polisemičnost bazy bez izměnjenja avtorinogo nazvanja.\par
\noindent\emph{Textual comparison:} Paper 31, §8, allows a module basis of $m\ge n$ elements. The required determinants are linear combinations of finitely many determinants drawn from that basis, which generate the discriminant ideal. A linearly independent module basis of $n$ elements is considered separately. This explicit distinction lets readers understand the historical polysemy of basis without replacing Noether's designation.\par
\noindent{\footnotesize TEXT-COMP-103; pasaže / passages: 18897; \texttt{CORPUS\_FORM\_COMPARISONS.json}.}\par
\end{minipage}\par
\par\medskip\Needspace{7\baselineskip}
\noindent\begin{minipage}{\linewidth}
\hypertarget{TERM-0861}{}
\noindent\textbf{diferenciranje idealov}\quad{\footnotesize\texttt{TERM-0861}}\par
\noindent\emph{Iztočnik / Source:} Idealdifferentiation\par
\noindent\emph{Rusky / Russian:} дифференцирование идеалов\par
\noindent\emph{Ukrajinsky / Ukrainian:} диференціювання ідеалів\par
\smallskip\noindent\emph{Razlog:} Diferenciranje idealov nazvaje operaciju, a ne samo formalnu sličnost s diferenciranjem polinomov. Izbor procesnogo imennika so korenom diferenc- i doplnjenja idealov dava kratky, razpoznavajemy izraz. Rabota 36 objavjaje vez s diferentoju; rabota 43 razviva definiciju s razlikovym idealom i idealnym kvocientom, poslě čego sledi specializacija. Zato ne trěba objašnjati termin prosto kak priměnjenje običnogo izvodnogo k svakomu tvorjačemu elementu. Tekst pokazuje, pri kakom specialnom slučaju konstrukcija sběgaje se s polinomnym diferencialnym kvocientom. Sohranjanje toj svezi i jejnyh uslovij jest važnějše než sličnost nazvanij.\par
\noindent\emph{Reason:} Diferenciranje idealov names an operation, not merely a formal resemblance to differentiating polynomials. A process noun with the stem diferenc- and the complement idealov gives a short, recognizable expression. Paper 36 announces the connection with the different; Paper 43 develops the definition through a difference ideal and an ideal quotient, followed by specialization. The term should therefore not be explained simply as applying the ordinary derivative to each generator. The text identifies the special case in which the construction coincides with a polynomial differential quotient. Preserving this connection and its conditions matters more than the resemblance between the names.\par
\noindent{\footnotesize Izvorne zapisy / Original records: HIST-g206-001, HIST-g213-001.}\par
\smallskip\noindent\emph{Sravnjenje s tekstom:} Kratka rabota 36 objavjaje interpretaciju diferenty kak diferencialnogo kvocienta definujučego ideala. V rabotě 43 ide o razlikovy ideal, idealny kvocient v ukazanoj polinomnoj oblasti i specializaciju. §3 izrično upotrěbja $(0):\mathfrak B$. Vez s obyčnym diferencialnym kvocientom polinoma nastupaje pri specialnom uslovju, že baza sostoji se iz stepenjev jednogo elementa. Slovnik zato ne smi zaměniti obšču konstrukciju proizvoljnym diferenciranjem tvorjačih polinomov.\par
\noindent\emph{Textual comparison:} The short Paper 36 announces an interpretation of the different as the differential quotient of the defining ideal. Paper 43 uses a difference ideal, an ideal quotient in a specified polynomial domain and specialization. Section 3 explicitly uses $(0):\mathfrak B$. The connection with an ordinary polynomial differential quotient appears under the special condition that the basis consists of powers of one element. The dictionary must not replace the general construction with arbitrary differentiation of generating polynomials.\par
\noindent{\footnotesize TEXT-COMP-106; pasaže / passages: 22743, 24741, 24751, 24755, 24759, 24861, 24863; \texttt{CORPUS\_FORM\_COMPARISONS.json}.}\par
\end{minipage}\par
\par\medskip\Needspace{7\baselineskip}
\noindent\begin{minipage}{\linewidth}
\hypertarget{TERM-0862}{}
\noindent\textbf{klas idealov / klas predstavjenj}\quad{\footnotesize\texttt{TERM-0862}}\par
\noindent\emph{Iztočnik / Source:} Idealklasse / Darstellungsklasse\par
\noindent\emph{Rusky / Russian:} класс идеалов / класс представлений\par
\noindent\emph{Ukrajinsky / Ukrainian:} клас ідеалів / клас представлень\par
\smallskip\noindent\emph{Razlog:} Klas idealov i klas predstavjenj tvoret paru, zato že tekst povezuje dve klasifikacije. Genitivne konstrukcije izrično kažu, čto jest klasifikovano, vměsto da pridevnik idealny mogl zvučati kak obča ocěna. V sravnjenom odlomku raboty 31 idealy naležat istomu klasu, kogda sut izomorfne kak moduly nad danom prstenom; zatim ustanovjava se korespondencija s klasami prědstavjenij. To jest konkreten odnos ekvivalentnosti, ne svobodno grupovanje podobnyh objektov. Zapisana kontinuitet s ranějšimi izrazami podpiraje vybor, ale ne rěšaje variantnost klas/klasa ni pravopis predstavjenj/prědstavjenij v celom korpusu.\par
\noindent\emph{Reason:} Klas idealov and klas predstavjenj form a pair because the text connects two classifications. The genitive constructions make explicit what is classified, avoiding an adjectival idealny that could sound evaluative outside mathematics. In the compared passage from Paper 31, ideals belong to the same class when they are isomorphic as modules over the specified ring; a correspondence with representation classes is then established. This is a particular equivalence relation, not an informal grouping of similar objects. Recorded continuity supports the choice but does not settle klas/klasa or predstavjenj/prědstavjenij throughout the corpus.\par
\noindent{\footnotesize Izvorne zapisy / Original records: HIST-g138-002.}\par
\smallskip\noindent\emph{Sravnjenje s tekstom:} Definicija v §4 raboty 31 povezuje naležnost dvoh idealov istomu klasu s izomorfizmom kak modulov nad osnovnym prstenom. Poslě togo tekst ustanovjava korespondenciju s klasami prědstavjenij. Zato klas ne jest prosto sovkupnost podobnyh idealov; genitivna konstrukcija sohranja objekt klasifikacije, a priložena definicija sohranja odnos ekvivalentnosti.\par
\noindent\emph{Textual comparison:} The definition in Paper 31, §4, connects two ideals belonging to the same class with their isomorphism as modules over the underlying ring. The text then establishes a correspondence with representation classes. Thus class is not merely a collection of similar ideals: the genitive construction preserves the object classified, and the attached definition preserves the equivalence relation.\par
\noindent{\footnotesize TEXT-COMP-104; pasaže / passages: 17950, 17985; \texttt{CORPUS\_FORM\_COMPARISONS.json}.}\par
\end{minipage}\par
\par\medskip\Needspace{7\baselineskip}
\noindent\begin{minipage}{\linewidth}
\hypertarget{TERM-0863}{}
\noindent\textbf{idealna klasa; glavna klasa; glavne idealy}\quad{\footnotesize\texttt{TERM-0863}}\par
\noindent\emph{Iztočnik / Source:} Idealklasse; Hauptklasse; Hauptideale\par
\noindent\emph{Rusky / Russian:} [Bukvy ne sut zachovane / Lettering not recoverable]\par
\noindent\emph{Ukrajinsky / Ukrainian:} [Bukvy ne sut zachovane / Lettering not recoverable]\par
\smallskip\noindent\emph{Razlog:} Idealna klasa, glavna klasa i glavne idealy sohranjajut rodinu nazvanij iz raboty 41. Pridevnik glavny povezuje glavnu klasu s glavnimi idealami i s terminom glavny rod, ne s proizvoljno največšoju ili najvažnějšoju klasoju. To jest matematična motivacija za soglasovanost vyboru. Pridevnična idealna klasa razlikuje se gramatično od genitivnogo klasy idealov; razlika ne ustanovjava dva različne pojmy. Ruske i ukrajinske bukvy v tom zapisu ne sut zachovane. Zato njim ne možno pripisati konkretne sličnosti ni vagovany dokaz, hoti obči razlog kontinuiteta ostaje čitateljny.\par
\noindent\emph{Reason:} Idealna klasa, glavna klasa and glavne idealy preserve the family of expressions in Paper 41. Glavny connects principal class with principal ideals and with principal genus, not with an arbitrarily largest or most important class. That is the mathematical reason for coordinating the choices. Adjectival idealna klasa differs grammatically from genitive klasy idealov; the difference does not establish two mathematical concepts. The Russian and Ukrainian lettering in this record has not survived. Specific similarities or weighted support cannot therefore be attributed to those fields, although the general continuity argument remains readable.\par
\noindent{\footnotesize Izvorne zapisy / Original records: HIST-g211-009.}\par
\end{minipage}\par
\par\medskip\Needspace{7\baselineskip}
\noindent\begin{minipage}{\linewidth}
\hypertarget{TERM-0864}{}
\noindent\textbf{klasy idealov}\quad{\footnotesize\texttt{TERM-0864}}\par
\noindent\emph{Iztočnik / Source:} Idealklassen\par
\noindent\emph{Rusky / Russian:} классы идеалов\par
\noindent\emph{Ukrajinsky / Ukrainian:} класи ідеалів\par
\smallskip\noindent\emph{Razlog:} Klasy idealov daje prozračny plural: klasifikovane objekty sut idealy. Zapis dava toj konstrukciji prednost prěd težšoju pridevničnoju kalkoju, sohranjajuči vez s glavnoju klasoju i korespondencijeju klasov. Izbor jest stilističny i terminologičny, ne rezultat izměrenoj čitateljskoj razumlivosti. V rabotě 31 smysl klasa opiraje se na modulny izomorfizm; ne treba avtomatično prěnesti sem vsaku definiciju idealnogo klasa iz drugogo razděla teorije čisel. Srodny zapis s jedinicoju klas idealov i variant klasa ostajut vidime, aby ekspert mogl posuditi i smyslovu kontinuitet, i jezykovu normu.\par
\noindent\emph{Reason:} Klasy idealov gives a transparent plural: the classified objects are ideals. The record prefers this construction to a heavier adjectival calque while preserving links with the principal class and class correspondence. This is a stylistic and terminological judgment, not measured reader comprehension. In Paper 31 the sense of class rests on module isomorphism; not every definition of ideal class from another part of number theory should be imported automatically. The related singular klas idealov and the variant klasa remain visible so that an expert can assess both continuity of meaning and linguistic usage.\par
\noindent{\footnotesize Izvorne zapisy / Original records: HIST-g133-001.}\par
\smallskip\noindent\emph{Sravnjenje s tekstom:} Definicija v §4 raboty 31 povezuje naležnost dvoh idealov istomu klasu s izomorfizmom kak modulov nad osnovnym prstenom. Poslě togo tekst ustanovjava korespondenciju s klasami prědstavjenij. Zato klas ne jest prosto sovkupnost podobnyh idealov; genitivna konstrukcija sohranja objekt klasifikacije, a priložena definicija sohranja odnos ekvivalentnosti.\par
\noindent\emph{Textual comparison:} The definition in Paper 31, §4, connects two ideals belonging to the same class with their isomorphism as modules over the underlying ring. The text then establishes a correspondence with representation classes. Thus class is not merely a collection of similar ideals: the genitive construction preserves the object classified, and the attached definition preserves the equivalence relation.\par
\noindent{\footnotesize TEXT-COMP-104; pasaže / passages: 17950, 17985; \texttt{CORPUS\_FORM\_COMPARISONS.json}.}\par
\end{minipage}\par
\par\medskip\Needspace{7\baselineskip}
\noindent\begin{minipage}{\linewidth}
\hypertarget{TERM-0865}{}
\noindent\textbf{idealny kvocient}\quad{\footnotesize\texttt{TERM-0865}}\par
\noindent\emph{Iztočnik / Source:} Idealquotient\par
\noindent\emph{Rusky / Russian:} частное идеалов\par
\noindent\emph{Ukrajinsky / Ukrainian:} частка ідеалів\par
\smallskip\noindent\emph{Razlog:} Idealny kvocient nazvaje kolonnu operaciju $\mathfrak a:\mathfrak b$, ne prsten klasov ostatkov i ne polje kvocientov. Pri obyčnom čitanju kolonnej operacije ishodi se od elementov, kotoryh produkty s $\mathfrak b$ ležat v $\mathfrak a$; kontekst oprědělja osnovnu oblast. Zapis dava prednost kratkoj konstrukciji pridevnik plus imennik, sposobnoj povtorjati se v dokazah. Ruske i ukrajinske paralely upotrěbjajut različne srodne imenniky; oni podpirajut pojem, ale sami ne ustanovjavajut medžuslovjansky kvocient. Osoblivo važno jest ne sputati avtorin děliteljny poredok s obyčnym poredkom vključenja množstv.\par
\noindent\emph{Reason:} Idealny kvocient names the colon operation $\mathfrak a:\mathfrak b$, not a residue-class ring or a field of fractions. Under the usual reading of the colon operation, one considers elements whose products with $\mathfrak b$ lie in $\mathfrak a$; the context specifies the ambient domain. The record favours a short adjective–noun construction usable repeatedly in proofs. Russian and Ukrainian render quotient by related but distinct nouns; these support the concept without themselves establishing Interslavic kvocient. It is especially important not to confuse Noether's divisibility ordering with ordinary set inclusion.\par
\noindent{\footnotesize Izvorne zapisy / Original records: HIST-g107-008.}\par
\smallskip\noindent\emph{Sravnjenje s tekstom:} Definicija raboty 30 upotrěbja děliteljny jezyk: jejino najmenši modul ne možno bez objašnjenja prěnesti v jezyk najmenej sovkupnosti po vključenju. Uslovje kolonnej operacije oprědělja elementy, ktoryh produkty popadajut v pervy ideal. Pri $(0):\mathfrak a$ ono dava ideal elementov, ktore uništožajut $\mathfrak a$ pri množenju; rabota 31 izrično razlikuje toj ideal od ostatkovogo prstena, v ktorogo imenitelju on stoji. Takže $\mathfrak b$ prost po odnošenju k $\mathfrak a$ jest orientovano uslovje; vzajemna prostost potrěbuje obě orientacije.\par
\noindent\emph{Textual comparison:} Paper 30's definition uses divisibility language: its smallest module cannot be transferred without explanation into the language of smallest sets under inclusion. The colon condition determines elements whose products fall in the first ideal. For $(0):\mathfrak a$ it gives the annihilator of the ideal; Paper 31 explicitly distinguishes this ideal from the quotient ring in whose denominator it occurs. Likewise, $\mathfrak b$ prime relative to $\mathfrak a$ is a directed condition; mutual relative primality requires both directions.\par
\noindent{\footnotesize TEXT-COMP-105; pasaže / passages: 16628, 16634, 16636, 17874, 18152; \texttt{CORPUS\_FORM\_COMPARISONS.json}.}\par
\end{minipage}\par
\par\medskip\Needspace{7\baselineskip}
\noindent\begin{minipage}{\linewidth}
\hypertarget{TERM-0866}{}
\noindent\textbf{idealny kvocient nulovogo ideala črez $\mathfrak{a}$}\quad{\footnotesize\texttt{TERM-0866}}\par
\noindent\emph{Iztočnik / Source:} Idealquotient des Nullideals durch $\mathfrak{a}$\par
\noindent\emph{Rusky / Russian:} частное нулевого идеала по $\mathfrak{a}$\par
\noindent\emph{Ukrajinsky / Ukrainian:} частка нульового ідеалу за $\mathfrak{a}$\par
\smallskip\noindent\emph{Razlog:} V izrazu idealny kvocient nulovogo ideala črez $\mathfrak a$ nulovy ideal jest prvy operand, a $\mathfrak a$ vtory. Konstrukcija $(0):\mathfrak a$ vyznača elementy, kotore uništožajut $\mathfrak a$ pri množenju. V rabotě 31 ona pojavjaje se kak ideal v imenitelju prstena klasov ostatkov; sam kolonny ideal ne jest tym ostatkovym prstenom. Proto prěvod sohranja směr operacije, ne skračaje frazu do kvocienta po idealu. Istorično črez ne mora ostati rekomendovanym prědlogom; srodny korpusny odlomok upotrěbja prěz. Smysl operacije i izbor prědloga sut odvojene pytanja.\par
\noindent\emph{Reason:} In idealny kvocient nulovogo ideala črez $\mathfrak a$, the zero ideal is the first operand and $\mathfrak a$ the second. The construction $(0):\mathfrak a$ singles out elements that annihilate $\mathfrak a$ under multiplication. In Paper 31 it occurs as the ideal in the denominator of a residue-class ring; the colon ideal itself is not that quotient ring. The translation therefore retains the direction of the operation instead of shortening the phrase to quotient by an ideal. Historical črez need not remain the recommended preposition; the related corpus passage uses prěz. The operation's meaning and the prepositional choice are separate questions.\par
\noindent{\footnotesize Izvorne zapisy / Original records: HIST-g131-006.}\par
\smallskip\noindent\emph{Sravnjenje s tekstom:} Definicija raboty 30 upotrěbja děliteljny jezyk: jejino najmenši modul ne možno bez objašnjenja prěnesti v jezyk najmenej sovkupnosti po vključenju. Uslovje kolonnej operacije oprědělja elementy, ktoryh produkty popadajut v pervy ideal. Pri $(0):\mathfrak a$ ono dava ideal elementov, ktore uništožajut $\mathfrak a$ pri množenju; rabota 31 izrično razlikuje toj ideal od ostatkovogo prstena, v ktorogo imenitelju on stoji. Takže $\mathfrak b$ prost po odnošenju k $\mathfrak a$ jest orientovano uslovje; vzajemna prostost potrěbuje obě orientacije.\par
\noindent\emph{Textual comparison:} Paper 30's definition uses divisibility language: its smallest module cannot be transferred without explanation into the language of smallest sets under inclusion. The colon condition determines elements whose products fall in the first ideal. For $(0):\mathfrak a$ it gives the annihilator of the ideal; Paper 31 explicitly distinguishes this ideal from the quotient ring in whose denominator it occurs. Likewise, $\mathfrak b$ prime relative to $\mathfrak a$ is a directed condition; mutual relative primality requires both directions.\par
\noindent{\footnotesize TEXT-COMP-105; pasaže / passages: 16628, 16634, 16636, 17874, 18152; \texttt{CORPUS\_FORM\_COMPARISONS.json}.}\par
\end{minipage}\par
\par\medskip\Needspace{7\baselineskip}
\noindent\begin{minipage}{\linewidth}
\hypertarget{TERM-0867}{}
\noindent\textbf{kvocienty idealov}\quad{\footnotesize\texttt{TERM-0867}}\par
\noindent\emph{Iztočnik / Source:} Idealquotienten\par
\noindent\emph{Rusky / Russian:} частные идеалов\par
\noindent\emph{Ukrajinsky / Ukrainian:} частки ідеалів\par
\smallskip\noindent\emph{Razlog:} Kvocienty idealov jest pluralna genitivna variantnost idealny kvocient. Pozdny zapis povezuje ju s definicijnoju primětkoju pri rabotě 45 i s ranějšeju rabotoju 43. To objasnjaje kontinuitet vyboru: idealov čini argumenty operacije jasnymi, a kvocient sohranja kratky tehničny koren. Prěvod ne smi prětvoriti kolonnu operaciju v děljenje čisel ili v tvorjenje ostatkovyh klasov. Motivacija nazvaje izraz rabotnym tehničnym rěšenjem; ona ne dava dokaz prijetja od medžuslovjanskoj obščiny. Nově sravnjene definicije v rabotah 30 i 43 razjasnjajut smysl, ale ne zaměnjajut potrěbnu prověrku točnogo odlomka raboty 45.\par
\noindent\emph{Reason:} Kvocienty idealov is a plural genitive variant of idealny kvocient. The later record connects it with a defining footnote in Paper 45 and with Paper 43. This explains the continuity of the choice: idealov makes the operands explicit, while kvocient retains a compact technical stem. The translation must not turn the colon operation into numerical division or the formation of residue classes. The rationale treats it as a working technical solution, not evidence of acceptance by the Interslavic community. Newly compared definitions in Papers 30 and 43 clarify the meaning but do not replace checking the exact Paper 45 passage.\par
\noindent{\footnotesize Izvorne zapisy / Original records: HIST-g214-012.}\par
\smallskip\noindent\emph{Sravnjenje s tekstom:} Definicija raboty 30 upotrěbja děliteljny jezyk: jejino najmenši modul ne možno bez objašnjenja prěnesti v jezyk najmenej sovkupnosti po vključenju. Uslovje kolonnej operacije oprědělja elementy, ktoryh produkty popadajut v pervy ideal. Pri $(0):\mathfrak a$ ono dava ideal elementov, ktore uništožajut $\mathfrak a$ pri množenju; rabota 31 izrično razlikuje toj ideal od ostatkovogo prstena, v ktorogo imenitelju on stoji. Takže $\mathfrak b$ prost po odnošenju k $\mathfrak a$ jest orientovano uslovje; vzajemna prostost potrěbuje obě orientacije.\par
\noindent\emph{Textual comparison:} Paper 30's definition uses divisibility language: its smallest module cannot be transferred without explanation into the language of smallest sets under inclusion. The colon condition determines elements whose products fall in the first ideal. For $(0):\mathfrak a$ it gives the annihilator of the ideal; Paper 31 explicitly distinguishes this ideal from the quotient ring in whose denominator it occurs. Likewise, $\mathfrak b$ prime relative to $\mathfrak a$ is a directed condition; mutual relative primality requires both directions.\par
\noindent{\footnotesize TEXT-COMP-105; pasaže / passages: 16628, 16634, 16636, 17874, 18152; \texttt{CORPUS\_FORM\_COMPARISONS.json}.}\par
\end{minipage}\par
\par\medskip\Needspace{7\baselineskip}
\noindent\begin{minipage}{\linewidth}
\hypertarget{TERM-0868}{}
\noindent\textbf{idealne kvocienty / glavny ideal / prost vzhodno k $t$}\quad{\footnotesize\texttt{TERM-0868}}\par
\noindent\emph{Iztočnik / Source:} Idealquotienten / Hauptideal / prim zu $t$\par
\noindent\emph{Rusky / Russian:} частные идеалов / главный идеал / прост относительно $t$\par
\noindent\emph{Ukrajinsky / Ukrainian:} частки ідеалів / головний ідеал / простий щодо $t$\par
\smallskip\noindent\emph{Razlog:} Ta skupina povezuje idealny kvocient, glavny ideal i orientovanu relaciju prim zu. V rabotě 30 uslovje $\mathfrak a:\mathfrak b=\mathfrak a$ znači, že $\mathfrak b$ jest prost po odnošenju k $\mathfrak a$; obratny směr ne jest dan avtomatično. Preto izraz ne možno zaměniti simetričnym vzajemno prosty bez dodatnogo uslovja. Zapisana konstrukcija prost vzhodno k jest istorična: srodny prědložny vybor razgledan jest v DEC-0009, a sravnjena definicija upotrěbja po odnošenju k. Slovnik sohranja stary zapis za razuměnje vyboru, ne kak bezuslovnu rekomendaciju njegovej formulacije.\par
\noindent\emph{Reason:} This group connects ideal quotient, principal ideal and the directed relation prim zu. In Paper 30 the condition $\mathfrak a:\mathfrak b=\mathfrak a$ means that $\mathfrak b$ is prime relative to $\mathfrak a$; the reverse direction is not automatic. It therefore cannot be replaced with symmetric relatively prime without an additional condition. The recorded prost vzhodno k is historical: the related prepositional choice is discussed in DEC-0009, and the compared definition uses po odnošenju k. The dictionary retains the earlier wording to explain the choice, not as an unconditional recommendation of its formulation.\par
\noindent{\footnotesize Izvorne zapisy / Original records: HIST-g113-004.}\par
\smallskip\noindent\emph{Sravnjenje s tekstom:} Definicija raboty 30 upotrěbja děliteljny jezyk: jejino najmenši modul ne možno bez objašnjenja prěnesti v jezyk najmenej sovkupnosti po vključenju. Uslovje kolonnej operacije oprědělja elementy, ktoryh produkty popadajut v pervy ideal. Pri $(0):\mathfrak a$ ono dava ideal elementov, ktore uništožajut $\mathfrak a$ pri množenju; rabota 31 izrično razlikuje toj ideal od ostatkovogo prstena, v ktorogo imenitelju on stoji. Takže $\mathfrak b$ prost po odnošenju k $\mathfrak a$ jest orientovano uslovje; vzajemna prostost potrěbuje obě orientacije.\par
\noindent\emph{Textual comparison:} Paper 30's definition uses divisibility language: its smallest module cannot be transferred without explanation into the language of smallest sets under inclusion. The colon condition determines elements whose products fall in the first ideal. For $(0):\mathfrak a$ it gives the annihilator of the ideal; Paper 31 explicitly distinguishes this ideal from the quotient ring in whose denominator it occurs. Likewise, $\mathfrak b$ prime relative to $\mathfrak a$ is a directed condition; mutual relative primality requires both directions.\par
\noindent{\footnotesize TEXT-COMP-105; pasaže / passages: 16628, 16634, 16636, 17874, 18152; \texttt{CORPUS\_FORM\_COMPARISONS.json}.}\par
\end{minipage}\par
\par\medskip\Needspace{7\baselineskip}
\noindent\begin{minipage}{\linewidth}
\hypertarget{TERM-0869}{}
\noindent\textbf{teorija idealov v $\mathfrak{T}$}\quad{\footnotesize\texttt{TERM-0869}}\par
\noindent\emph{Iztočnik / Source:} Idealtheorie in $\mathfrak{T}$\par
\noindent\emph{Rusky / Russian:} теория идеалов в $\mathfrak{T}$\par
\noindent\emph{Ukrajinsky / Ukrainian:} теорія ідеалів у $\mathfrak{T}$\par
\smallskip\noindent\emph{Razlog:} Teorija idealov v $\mathfrak T$ označaje teoriju konkretno danogo poredka, ne obči ekskurs o idealah. V §7 raboty 31 za tymi slovami slěduje strukturny teorem o jednoznačnom prědstavjenju nenulovogo ideala produktom primarnyh idealov. Zato sohraniti $\mathfrak T$ i jegovo město v frazě jest čest matematičnoj věrnosti. Slovjansky genitiv idealov i prědlog v činjat konstrukciju prozračnoju; nije potrěbno tvoriti novy jednoslovny termin. Istorijski zapis T jest skračenym pokazateljem simbola, ktory v odlomku jest nabrany frakturom.\par
\noindent\emph{Reason:} Teorija idealov v $\mathfrak T$ concerns the theory of a specified order, not a general digression about ideals. In Paper 31, §7, the phrase introduces a structural theorem on uniquely representing a nonzero ideal as a product of primary ideals. Preserving $\mathfrak T$ and its place in the phrase is therefore part of mathematical fidelity. The Slavic genitive idealov and the preposition v make the construction explicit; no new single-word term is needed. Historical plain T is an abbreviated record of a symbol set in Fraktur in the passage.\par
\noindent{\footnotesize Izvorne zapisy / Original records: HIST-g150-001.}\par
\smallskip\noindent\emph{Sravnjenje s tekstom:} Tri primětky raboty 31 imajut konkretne funkcije: §7, 3 pri teoriji idealov poredka; §9, 7 pri integralnoj zamknutosti; §9 i §3, 1 pri egzistenciji konečnoj modulnoj bazy. Německe kratko nazvanje Idealtheorie identifikuje iztočnik. Prěvod veznika v tretjoj primětkě ne měnja nazvanje. Točne odlomky potrvrdžajut funkciju citacij v nynějšem korpusu; ne potrvrdžajut samo soboju prijašnje tvrdženja o pečatnom svědku.\par
\noindent\emph{Textual comparison:} The three Paper 31 footnotes have specific functions: §7, 3 accompanies the ideal theory of the order; §9, 7 accompanies integral closure; §9 and §3, 1 accompany existence of a finite module basis. The shortened German title Idealtheorie identifies the source. Translating the conjunction in the third note does not change the title. These passages establish the citations' functions in the present corpus; they do not by themselves verify earlier claims about the printed witness.\par
\noindent{\footnotesize TEXT-COMP-107; pasaže / passages: 18593, 18861, 18897; \texttt{CORPUS\_FORM\_COMPARISONS.json}.}\par
\end{minipage}\par
\par\medskip\Needspace{7\baselineskip}
\noindent\begin{minipage}{\linewidth}
\hypertarget{TERM-0873}{}
\noindent\textbf{razloženje ideala / razloženje nulovogo ideala}\quad{\footnotesize\texttt{TERM-0873}}\par
\noindent\emph{Iztočnik / Source:} Idealzerlegung / Zerlegung des Nullideals\par
\noindent\emph{Rusky / Russian:} разложение идеала / разложение нулевого идеала\par
\noindent\emph{Ukrajinsky / Ukrainian:} розклад ідеалу / розклад нульового ідеалу\par
\smallskip\noindent\emph{Razlog:} Razloženje ideala i razloženje nulovogo ideala upotrěbjajut isti koren za svezane konstrukcije. Doplnjenje nulovogo ne jest izlišno: pri prěhodu k ostatkovomu prstenu razloženje vybranogo ideala odgovarja razloženju nulovogo ideala v novoj strukturě. Zapis dava prednost kontinuitetu s terminom primarnogo razloženja i razpoznavajemomu slovjanskomu korenu. Ale jedno slovo razloženje ne odstranja razliku medžu produktom idealov, jih prěsěčenjem i direktnoju sumoju; konkretny teorem oprědělja operaciju. Slovnik sohranja svez, ne spaja različne teoremy v jednu neograničenu definiciju.\par
\noindent\emph{Reason:} Razloženje ideala and razloženje nulovogo ideala use the same stem for related constructions. Nulovogo is not redundant: on passing to a residue-class ring, a decomposition of the chosen ideal corresponds to a decomposition of the zero ideal in the new structure. The record favours continuity with primary-decomposition terminology and a recognizable Slavic stem. But decomposition alone does not erase the distinctions between a product of ideals, their intersection and a direct sum; the particular theorem determines the operation. The dictionary preserves the connection without merging different theorems into one unrestricted definition.\par
\noindent{\footnotesize Izvorne zapisy / Original records: HIST-g152-002.}\par
\smallskip\noindent\emph{Sravnjenje s tekstom:} V §7 raboty 31 homomorfizm prěnosi razloženje ideala $\mathfrak T p$ na razloženje nulovogo ideala v $\mathfrak R$. Odlomok takože razlučuje prědstavjenje slěda matricami od prědstavjenja prěz konjugovane elementy: poslednje ne mora byti priměnimo k $\mathfrak R$. Izbor slova ideal i sohranjanje obou simbolov jest nužno, aby slovnik ne izravnal iztočny objekt s jegovym obrazom.\par
\noindent\emph{Textual comparison:} In Paper 31, §7, a homomorphism takes decomposition of the ideal $\mathfrak T p$ to decomposition of the zero ideal in $\mathfrak R$. The passage also distinguishes the matrix description of trace from a description through conjugate elements: the latter need not apply to $\mathfrak R$. The word ideal and both symbols must be preserved so that the dictionary does not identify the original object with its image.\par
\noindent{\footnotesize TEXT-COMP-108; pasaže / passages: 18671; \texttt{CORPUS\_FORM\_COMPARISONS.json}.}\par
\end{minipage}\par
\par\medskip\Needspace{7\baselineskip}
\noindent\begin{minipage}{\linewidth}
\hypertarget{TERM-0874}{}
\noindent\textbf{razloženje ideala $\mathfrak{T}p$}\quad{\footnotesize\texttt{TERM-0874}}\par
\noindent\emph{Iztočnik / Source:} Idealzerlegung von $\mathfrak{T}p$\par
\noindent\emph{Rusky / Russian:} разложение идеала $\mathfrak{T}p$\par
\noindent\emph{Ukrajinsky / Ukrainian:} розклад ідеалу $\mathfrak{T}p$\par
\smallskip\noindent\emph{Razlog:} Razloženje ideala $\mathfrak T p$ izrično nazyvaje ideal, o ktorym ide govor. Sohranjanje simbola jest važno: tekst povezuje razloženje togo ideala s razloženjem nulovogo ideala pri prěhodu k ostatkovomu prstenu. Prěvod ne smi nazvati sam objekt ostatkovym prstenom ili izpustiti množenje v $\mathfrak T p$. Motivacija sohranja prědmet teorema, a ne prosto německy poredok slov. V istorijskom spisku T p jest ploška tekstova forma; matematično objašnjenje razlikuje frakturny poredok $\mathfrak T$ od obyčnogo simbola $p$.\par
\noindent\emph{Reason:} Razloženje ideala $\mathfrak T p$ explicitly names the ideal under discussion. Preserving the symbol matters: the text connects decomposition of this ideal with decomposition of the zero ideal upon passage to a residue-class ring. The translation must not call the object itself a quotient ring or lose the multiplication in $\mathfrak T p$. The rationale preserves the theorem's subject, not merely German word order. In the historical list T p is a plain-text rendering; the mathematical explanation distinguishes the Fraktur order $\mathfrak T$ from the ordinary symbol $p$.\par
\noindent{\footnotesize Izvorne zapisy / Original records: HIST-g163-005.}\par
\smallskip\noindent\emph{Sravnjenje s tekstom:} V §7 raboty 31 homomorfizm prěnosi razloženje ideala $\mathfrak T p$ na razloženje nulovogo ideala v $\mathfrak R$. Odlomok takože razlučuje prědstavjenje slěda matricami od prědstavjenja prěz konjugovane elementy: poslednje ne mora byti priměnimo k $\mathfrak R$. Izbor slova ideal i sohranjanje obou simbolov jest nužno, aby slovnik ne izravnal iztočny objekt s jegovym obrazom.\par
\noindent\emph{Textual comparison:} In Paper 31, §7, a homomorphism takes decomposition of the ideal $\mathfrak T p$ to decomposition of the zero ideal in $\mathfrak R$. The passage also distinguishes the matrix description of trace from a description through conjugate elements: the latter need not apply to $\mathfrak R$. The word ideal and both symbols must be preserved so that the dictionary does not identify the original object with its image.\par
\noindent{\footnotesize TEXT-COMP-108; pasaže / passages: 18671; \texttt{CORPUS\_FORM\_COMPARISONS.json}.}\par
\end{minipage}\par
\par\medskip\Needspace{7\baselineskip}
\noindent\begin{minipage}{\linewidth}
\hypertarget{TERM-0875}{}
\noindent\textbf{idempotentne elementy}\quad{\footnotesize\texttt{TERM-0875}}\par
\noindent\emph{Iztočnik / Source:} idempotente Elemente\par
\noindent\emph{Rusky / Russian:} идемпотентные элементы\par
\noindent\emph{Ukrajinsky / Ukrainian:} ідемпотентні елементи\par
\smallskip\noindent\emph{Razlog:} Idempotentne elementy sohranja medžunarodny tehničny koren; opisujuče zaměny ne sut navedene kak konkretny konkurenty. Motivacija povezuje termin s rabotoju 32 i spominaje češku i poljsku upotrěbu. No zapis ne dava točnyh odlomkov tyh dvoh jezykov, zato oni ne mogut byti prědstavjene kak neposrědnje citovane attestacije. Matematicky idempotentny element izpolnjaje $e^2=e$; v rabotě 34 on učestvuje v razloženjah idealov i prědstavjenijah. Nazvanje samo ne pomeni primitivnost, centralnost ni nenulovost. Takove dodatne uslovja trěba sohraniti tam, kde avtorica jih potrěbuje.\par
\noindent\emph{Reason:} Idempotentne elementy retains the international technical stem; no specific descriptive replacements are recorded as competitors. The rationale links the term with Paper 32 and mentions Czech and Polish usage. It supplies no exact passages from those languages, however, so they cannot be presented as directly cited attestations. Mathematically an idempotent satisfies $e^2=e$; in Paper 34 it participates in ideal decompositions and representations. The name alone does not imply primitivity, centrality or nonzeroness. Such additional conditions must be preserved wherever Noether requires them.\par
\noindent{\footnotesize Izvorne zapisy / Original records: HIST-g187-001.}\par
\smallskip\noindent\emph{Sravnjenje s tekstom:} Rabota 32 povezuje idempotentny element $r$ s razloženjem $S=rS+(e-r)S$, kde $e$ jest jedinica. Primitivne idempotenty vyděljajut proste idealy. V kvaternionnom slučaju tekst dodatno izključuje 0 i 1 i razgleduje razpad. Tym sučstveno jest odlišiti osnovne svojstvo idempotentnosti od primitivnosti i od dodatnyh hipotez; sam koren idempotent- ne nosi vsih tyh uslovij.\par
\noindent\emph{Textual comparison:} Paper 32 connects an idempotent $r$ with a decomposition $S=rS+(e-r)S$, where $e$ is the identity. Primitive idempotents single out simple ideals. In the quaternion case the text additionally excludes 0 and 1 and discusses splitting. It is therefore essential to distinguish idempotency itself from primitivity and from additional hypotheses; the stem idempotent- does not carry all those conditions.\par
\noindent{\footnotesize TEXT-COMP-109; pasaže / passages: 19445, 19447, 19457; \texttt{CORPUS\_FORM\_COMPARISONS.json}.}\par
\end{minipage}\par
\par\medskip\Needspace{7\baselineskip}
\noindent\begin{minipage}{\linewidth}
\hypertarget{TERM-0876}{}
\noindent\textbf{idempotenty; komplementarne bazy}\quad{\footnotesize\texttt{TERM-0876}}\par
\noindent\emph{Iztočnik / Source:} Idempotente; komplement?re Basen\par
\noindent\emph{Rusky / Russian:} [Bukvy ne sut zachovane / Lettering not recoverable]\par
\noindent\emph{Ukrajinsky / Ukrainian:} [Bukvy ne sut zachovane / Lettering not recoverable]\par
\smallskip\noindent\emph{Razlog:} Idempotenty i komplementarne bazy sut v rabotě 40 matematično svezane, ne samo susědne tehnične slova. Odlomok povezuje koeficienty prědstavjenja idempotenta s obrazom komplementarnoj bazy. Zapis dava prednost medžunarodnym nazvanjam vměsto umětnyh novotvorb; ne uvodi za vsako novo město novy sinonim. Komplementarna baza tu ne znači proizvoljnu bazu doplňajučego podprostora: smysl oprědělja konkreten teorem. Ruske i ukrajinske bukvy izgubljene sut, a německy tekst imaje komplement?re. Tyto povrědy ostajut vidime; slovnik ne izdava domysljenu rekonstrukciju za zachovane dokazateljstvo.\par
\noindent\emph{Reason:} Idempotenty and komplementarne bazy are mathematically connected in Paper 40, not merely adjacent technical words. The passage links the coefficients in an idempotent's expression with the image of a complementary basis. The record favours international designations over artificial coinages rather than introducing a new synonym at each occurrence. Complementary basis here does not mean an arbitrary basis of a complementary subspace: the particular theorem determines its sense. The Russian and Ukrainian lettering is lost, and the German field has komplement?re. These defects remain visible; conjectural reconstruction is not presented as surviving evidence.\par
\noindent{\footnotesize Izvorne zapisy / Original records: HIST-g210-010.}\par
\smallskip\noindent\emph{Sravnjenje s tekstom:} Teorem v rabotě 40 izražaje idempotent $e$ po bazi $a_1,\ldots,a_n$; pri odgovarajučem odobraženju koeficienty davajut obraz komplementarnoj bazy. Zato komplementarna ne jest tu svobodno opisujuči pridevnik. Slovnik sohranja vez s teoremom i ne zaměnja jeju pojmom libovoljnogo doplňajučego podprostora.\par
\noindent\emph{Textual comparison:} The theorem in Paper 40 expresses an idempotent $e$ in a basis $a_1,\ldots,a_n$; under the corresponding mapping its coefficients give the image of the complementary basis. Complementary is therefore not a free descriptive adjective here. The dictionary retains the connection with the theorem rather than replacing it with the notion of an arbitrary complementary subspace.\par
\noindent{\footnotesize TEXT-COMP-110; pasaže / passages: 24055, 24057, 24061, 24065, 24067; \texttt{CORPUS\_FORM\_COMPARISONS.json}.}\par
\end{minipage}\par
\par\medskip\Needspace{7\baselineskip}
\noindent\begin{minipage}{\linewidth}
\hypertarget{TERM-0877}{}
\noindent\textbf{idempotentny element}\quad{\footnotesize\texttt{TERM-0877}}\par
\noindent\emph{Iztočnik / Source:} idempotentes Element\par
\noindent\emph{Rusky / Russian:} идемпотентный элемент\par
\noindent\emph{Ukrajinsky / Ukrainian:} ідемпотентний елемент\par
\smallskip\noindent\emph{Razlog:} Izraz idempotentny element označaje jeden element, a idempotentne elementy jest odgovarjajuča pluralna forma. Zapis pripisuje medžunarodnomu korenu vysoku razpoznavajemost; to jest kvalitativna ocěna, ne izměreny procent razuměnja. V kvaternionnom kontekstu raboty 32 sučstveno jest, že element ne jest ni 0 ni 1; tekst govori o razpadu poslě odgovarajučego razširjenja. Ne možno iz togo vyvesti, že izvorno tělo s děljenjem samo imaje netrivialne idempotenty. Podrobny slovnik sohranja i kratko nazvanje, i oblast upotrěby, vključno s ograničenjem na konkretno razgledany razširjeny sistem.\par
\noindent\emph{Reason:} Idempotentny element is the singular counterpart of the plural term. The record attributes high recognizability to its international stem; this is a qualitative judgment, not a measured comprehension percentage. In the quaternion setting of Paper 32 it matters that the element is neither 0 nor 1; the text concerns splitting after the relevant extension. One cannot infer that the original division algebra itself contains a nontrivial idempotent. A detailed dictionary preserves both the compact designation and its domain of use, including the restriction to the extended system under discussion.\par
\noindent{\footnotesize Izvorne zapisy / Original records: HIST-g174-002.}\par
\smallskip\noindent\emph{Sravnjenje s tekstom:} Rabota 32 povezuje idempotentny element $r$ s razloženjem $S=rS+(e-r)S$, kde $e$ jest jedinica. Primitivne idempotenty vyděljajut proste idealy. V kvaternionnom slučaju tekst dodatno izključuje 0 i 1 i razgleduje razpad. Tym sučstveno jest odlišiti osnovne svojstvo idempotentnosti od primitivnosti i od dodatnyh hipotez; sam koren idempotent- ne nosi vsih tyh uslovij.\par
\noindent\emph{Textual comparison:} Paper 32 connects an idempotent $r$ with a decomposition $S=rS+(e-r)S$, where $e$ is the identity. Primitive idempotents single out simple ideals. In the quaternion case the text additionally excludes 0 and 1 and discusses splitting. It is therefore essential to distinguish idempotency itself from primitivity and from additional hypotheses; the stem idempotent- does not carry all those conditions.\par
\noindent{\footnotesize TEXT-COMP-109; pasaže / passages: 19445, 19447, 19457; \texttt{CORPUS\_FORM\_COMPARISONS.json}.}\par
\end{minipage}\par
\par\medskip\Needspace{7\baselineskip}
\noindent\begin{minipage}{\linewidth}
\hypertarget{TERM-0878}{}
\noindent\textbf{identično v $Y$}\quad{\footnotesize\texttt{TERM-0878}}\par
\noindent\emph{Iztočnik / Source:} identisch in $Y$\par
\noindent\emph{Rusky / Russian:} тождественно по $Y$\par
\noindent\emph{Ukrajinsky / Ukrainian:} тотожно щодо $Y$\par
\smallskip\noindent\emph{Razlog:} Identično v $Y$ sohranja točno matematično ograničenje: relacija važi kak identičnost po proměnnoj $Y$, ne samo pri jednoj specialnoj cěnnosti. Zapis povezuje frazu s relacijeju (7) v §7 raboty 9. Rusko toždestvenno i ukrajinsko totožno podpirajut potrebu izraziti identičnost, ale ne daju osnovu tiho zaměniti proměnnu ili oslabiti kvantifikaciju. Prědlog v jest čest zapisanoj konstrukcije; dlja čitatelja objašnjenje mora pokazati, v ktoroj proměnnoj identičnost razgleduje se. To jest izbor sintakse s matematičnym slědstvom, ne samo slovnična zaměna identisch.\par
\noindent\emph{Reason:} Identično v $Y$ preserves a precise mathematical qualification: the relation holds as an identity in $Y$, not merely at one special value. The record connects the phrase with relation (7) in Paper 9, §7. Russian toždestvenno and Ukrainian totožno support the need to express identity, but do not justify silently changing the variable or weakening the quantification. The preposition v is part of the recorded construction; the explanation must make clear to readers in which variable the identity is considered. This is a syntactic choice with mathematical consequences, not just a word-for-word substitution for identisch.\par
\noindent{\footnotesize Izvorne zapisy / Original records: HIST-g065-016.}\par
\end{minipage}\par
\par\medskip\Needspace{7\baselineskip}
\noindent\begin{minipage}{\linewidth}
\hypertarget{TERM-0879}{}
\noindent\textbf{identične razklady}\quad{\footnotesize\texttt{TERM-0879}}\par
\noindent\emph{Iztočnik / Source:} identische Zerlegungen\par
\noindent\emph{Rusky / Russian:} тождественные разложения\par
\noindent\emph{Ukrajinsky / Ukrainian:} тотожні розклади\par
\smallskip\noindent\emph{Razlog:} Identične razklady razlikuje istost samyh razloženij od izomorfnosti jih česti. Motivacija iz raboty 17 naglašaje, že ponavljanje istogo izomorfizmova klasa v različnyh razloženjah ne čini te razloženja identičnymi. Zato ne trěba zaměnjati identične slovom ekvivalentne bez definicije, ni razklady s izomorfnymi komponentami. Rusky i ukrajinsky paralely podpirajut silny smysl istosti. Izbor razklady koegzistuje s razloženje v drugyh zapisah; slovnik mora pokazati variantnost, ale ne izdumati smyslovu razliku samo iz različnyh korenjev prěvoda.\par
\noindent\emph{Reason:} Identične razklady distinguishes identity of the decompositions themselves from isomorphism of their parts. The Paper 17 rationale stresses that recurrence of the same isomorphism class in different decompositions does not make the decompositions identical. Identical should therefore not become equivalent without a definition, nor should decompositions become isomorphic components. The Russian and Ukrainian parallels support the stronger sense of identity. Razklady coexists with razloženje in other records; the dictionary must show that variation without inventing a conceptual distinction solely from different translated stems.\par
\noindent{\footnotesize Izvorne zapisy / Original records: HIST-g079-049.}\par
\end{minipage}\par
\par\medskip\Needspace{7\baselineskip}
\noindent\begin{minipage}{\linewidth}
\hypertarget{TERM-0880}{}
\noindent\textbf{identičnost polj}\quad{\footnotesize\texttt{TERM-0880}}\par
\noindent\emph{Iztočnik / Source:} Identität der Körper\par
\noindent\emph{Rusky / Russian:} тождество полей\par
\noindent\emph{Ukrajinsky / Ukrainian:} тотожність полів\par
\smallskip\noindent\emph{Razlog:} Identičnost polj sohranja ravnost konkretno razgledanyh polj, ne samo jih abstraktnu izomorfnost. Zapis iz §6 raboty 6 upiraje se na sravnjenje stepenjev razširjenja v dokazu pripisanom Steinitzu. V sravnjenom odlomku važno jest takože vključenje polj v obče nadpolje. Jednaky konečny stepen sám po sobě ne dokazuje, že dva proizvoljne polja sut identične; zaključanje rabota zaradi togo vključenja. Zato slovnik objasnjaje i termin, i logičnu osnovu frazy. Izbor ne dodava novy teorem: on brani silu ravnosti, ktoru argument už potrěbuje.\par
\noindent\emph{Reason:} Identičnost polj preserves equality of the particular fields under consideration, not merely their abstract isomorphism. The Paper 6, §6, record relies on a comparison of extension degrees in an argument attributed to Steinitz. In the passage compared, inclusion of the fields within a common overfield also matters. Equal finite degrees alone do not prove that two arbitrary fields are identical; the conclusion works because of that inclusion. The dictionary therefore explains both the term and the logical basis of the wording. The choice adds no new theorem: it protects the strength of equality already required by the argument.\par
\noindent{\footnotesize Izvorne zapisy / Original records: HIST-g046-015.}\par
\smallskip\noindent\emph{Sravnjenje s tekstom:} V rabotě 6 argument upotrěbja vključenje podpolj v isto polje racionalnyh funkcij i ravnost odgovarajučih konečnyh stepenjev. Zato identičnost polj označaje ravnost vybranyh podpolj v danoj situaciji. Bez vključenja ravnost stepenjev ne dostava. To jest uslovje samogo argumenta, a ne dodatno opravdanje vybrano samo zaradi prěvoda.\par
\noindent\emph{Textual comparison:} Paper 6's argument uses inclusion of subfields within the same rational-function field and equality of the corresponding finite degrees. Identity of fields therefore means equality of the selected subfields in that setting. Without inclusion, equal degrees are insufficient. This is a condition of the argument itself, not an extra justification invented solely for the translation.\par
\noindent{\footnotesize TEXT-COMP-111; pasaže / passages: 5574, 5704; \texttt{CORPUS\_FORM\_COMPARISONS.json}.}\par
\end{minipage}\par
\par\medskip\Needspace{7\baselineskip}
\noindent\begin{minipage}{\linewidth}
\hypertarget{TERM-0881}{}
\noindent\textbf{identičnost medžu polinomami}\quad{\footnotesize\texttt{TERM-0881}}\par
\noindent\emph{Iztočnik / Source:} Identität zwischen Polynomen\par
\noindent\emph{Rusky / Russian:} тождество между полиномами\par
\noindent\emph{Ukrajinsky / Ukrainian:} тотожність між поліномами\par
\smallskip\noindent\emph{Razlog:} Identičnost medžu polinomami čini vidimym konkreten krok dokaza: od racionalnyh funkcij prěhodi se k polinomnoj identičnosti. Doplnjenje medžu polinomami zato ne jest samo stilistična izbytočnost. Ono kaže, v kakih objektah relacija nyně važi, i ne dopušča zaměniti identičnost ravnostju pri jednoj specializaciji. Rusky i ukrajinsky paralely podpirajut istu konstrukciju. Zapis daje to funkcionalno opravdanje; ne dava sravnjenu listu konkurentnyh nazvanij ni dokaz, že fraza byla novotvorba.\par
\noindent\emph{Reason:} Identičnost medžu polinomami makes a particular proof step visible: the passage moves from rational functions to a polynomial identity. The complement medžu polinomami is therefore not merely stylistic redundancy. It identifies the objects in which the relation now holds and prevents identity from becoming equality at one specialization. The Russian and Ukrainian parallels support the same construction. The record supplies this functional justification, not a comparison of competing names or evidence that the phrase was newly coined.\par
\noindent{\footnotesize Izvorne zapisy / Original records: HIST-g042-010.}\par
\end{minipage}\par
\par\medskip\Needspace{7\baselineskip}
\noindent\begin{minipage}{\linewidth}
\hypertarget{TERM-0883}{}
\noindent\textbf{teorema identičnosti}\quad{\footnotesize\texttt{TERM-0883}}\par
\noindent\emph{Iztočnik / Source:} Identitätssatz\par
\noindent\emph{Rusky / Russian:} теорема тождества\par
\noindent\emph{Ukrajinsky / Ukrainian:} теорема тотожності\par
\smallskip\noindent\emph{Razlog:} Teorema identičnosti sohranja prozračnu dvuslovnu strukturu německogo nazvanja. Zapis dava medžunarodnomu korenu identičnost prednost zaradi tehničnoj točnosti. Ale samo nazvanje ne identificiruje jednu univerzalnu teoremu: v rabotě 2 ono prinaleži invariantno-teoretičnomu kontekstu. Čitatelj ne trěbuje prědpolagati, že ide o poznanu jemu teoremu identičnosti iz kompleksnoj analizy. Točna formulacija i odkaz oprěděljajut smysl. Zapisani paralely pokazyvajut srodnu konstrukciju, ne merenje točnosti ni polno dokazateljstvo medžuslovjanskoj normy.\par
\noindent\emph{Reason:} Teorema identičnosti retains the transparent two-word structure of the German designation. The record favours the international stem identičnost for technical precision. The name alone, however, does not identify one universal theorem: in Paper 2 it belongs to an invariant-theoretic setting. Readers should not assume that it denotes the familiar identity theorem of complex analysis. The statement and reference determine its meaning. The recorded parallels show a related construction, not a measurement of precision or complete evidence of an Interslavic norm.\par
\noindent{\footnotesize Izvorne zapisy / Original records: HIST-g003-005.}\par
\smallskip\noindent\emph{Sravnjenje s tekstom:} V rabotě 2 teorema identičnosti pojavjaje se pri relacijah medžu polarami, vměstě s teoremoju produkta za determinanty. To lokalizuje nazvanje v invariantno-teoretičnom argumentu. Samo obče nazvanje ne jest dostatočnym odkazom na točnu formulaciju ni dovoljenjem zaměniti ju drugoju poznanoju teoremoju identičnosti.\par
\noindent\emph{Textual comparison:} In Paper 2 the identity theorem appears in the discussion of relations between polars, together with the product theorem for determinants. This locates the name within the invariant-theoretic argument. The general designation alone is neither an adequate reference to the exact statement nor permission to substitute another familiar identity theorem.\par
\noindent{\footnotesize TEXT-COMP-118; pasaže / passages: 719; \texttt{CORPUS\_FORM\_COMPARISONS.json}.}\par
\end{minipage}\par
\par\medskip\Needspace{7\baselineskip}
\noindent\begin{minipage}{\linewidth}
\hypertarget{TERM-0884}{}
\noindent\textbf{v imenovatelju}\quad{\footnotesize\texttt{TERM-0884}}\par
\noindent\emph{Iztočnik / Source:} im Nenner\par
\noindent\emph{Rusky / Russian:} в знаменателе\par
\noindent\emph{Ukrajinsky / Ukrainian:} у знаменнику\par
\smallskip\noindent\emph{Razlog:} V imenovatelju sohranja město uslovja v definiciji prstena kvocientov. Dopuščene imenovatelje, ne proizvoljne čislitelje, oprěděljajut lokalizaciju iz raboty 31. V sravnjenom odlomku imenovatelj mora byti elementom osnovnogo poredka prostym k ukazanomu idealu; primětka razjasnjaje smysl toj prostosti. Zato fraza ne smi skratiti se do neoprěděljenogo s takim elementom. Imennik ima konkretno matematično dejstvo v definiciji. Varijant imenitelj v drugyh odlomkah mora ostati predmetom terminologičnogo sravnjenja, ne pričinoju izpustiti samo uslovje.\par
\noindent\emph{Reason:} V imenovatelju preserves the location of a condition in the definition of a quotient ring. Permitted denominators, not arbitrary numerators, determine the localization in Paper 31. In the compared passage the denominator must be an element of the underlying order prime to the specified ideal; a footnote clarifies that relation. The phrase must therefore not be shortened to an unspecified with such an element. The noun performs a precise mathematical function in the definition. The variant imenitelj elsewhere calls for terminological comparison, not omission of the condition.\par
\noindent{\footnotesize Izvorne zapisy / Original records: HIST-g159-004.}\par
\smallskip\noindent\emph{Sravnjenje s tekstom:} Točka o prěhodu k prstenu kvocientov v rabotě 31 oprědělja, ktore elementy mogut byti v imenovatelju. Primětka pri toj točkě odkazyvaje na Grella i sohranja avtorin lokator primětka 6, str. 83. Definicija i odkaz imajut različne funkcije: prva oprědělja objekt, vtory ukazuje iztočnik dokazov. Ni jedno ne trěba izpustiti pri prěvodu.\par
\noindent\emph{Textual comparison:} Paper 31's paragraph on passage to the quotient ring specifies which elements may occur in the denominator. Its footnote refers to Grell and retains Noether's locator note 6, p. 83. Definition and reference have different functions: one determines the object, the other identifies the source of proofs. Neither should disappear in translation.\par
\noindent{\footnotesize TEXT-COMP-115; pasaže / passages: 18932; \texttt{CORPUS\_FORM\_COMPARISONS.json}.}\par
\end{minipage}\par
\par\medskip\Needspace{7\baselineskip}
\noindent\begin{minipage}{\linewidth}
\hypertarget{TERM-0885}{}
\noindent\textbf{nerazložimo vzhodno k $[H]$}\quad{\footnotesize\texttt{TERM-0885}}\par
\noindent\emph{Iztočnik / Source:} in bezug auf $[H]$ irreduzibel\par
\noindent\emph{Rusky / Russian:} неприводимое относительно $[H]$\par
\noindent\emph{Ukrajinsky / Ukrainian:} незвідне відносно $[H]$\par
\smallskip\noindent\emph{Razlog:} Nerazložimo vzhodno k $[H]$ sohranja zavisnost nerazložimosti od oblasti koeficientov. Izpustiti $[H]$ značilo by oslabiti informaciju potrěbnu pri minimalnom uravnenju v rabotě 9. Skobky sut čest matematičnogo označenja, ne ukazanje na tekst, ktory možno izpustiti. Forma vzhodno k jest istorijska: nynějšne sravnjene odlomky upotrěbjajut po odnošenju k, a odvojena odluka DEC-0009 objasnjaje prědložnu politiku. Smyslova opravdanost relativnogo doplnjenja ne opravdava avtomatično vsaku ranějšu formulaciju.\par
\noindent\emph{Reason:} Nerazložimo vzhodno k $[H]$ retains the dependence of irreducibility on the coefficient domain. Omitting $[H]$ would remove information required at the minimal-equation step in Paper 9. The brackets belong to the mathematical identifier; they do not mark dispensable text. Vzhodno k is historical: the present passages compared use po odnošenju k, and DEC-0009 separately explains the prepositional policy. The semantic justification for the relative qualification does not automatically justify every earlier formulation.\par
\noindent{\footnotesize Izvorne zapisy / Original records: HIST-g061-010.}\par
\smallskip\noindent\emph{Sravnjenje s tekstom:} Rabota 9 razlikuje bazu $H$, oblast $[H]$ i drugu bazu $H'$. Primětka pokazuje, že bazove svojstva mogut važiti po odnošenju k $H'$, ale ne k $H$. Odlomky o nerazložimom uravnenju izrično sohranjajut $[H]$. Zato zapisane relativne frazy nesut matematičnu informaciju; nynějšny izraz po odnošenju k zaměnja formulaciju, ne izpustja objekt odnosa.\par
\noindent\emph{Textual comparison:} Paper 9 distinguishes the basis $H$, the domain $[H]$ and another basis $H'$. A remark shows that the basis properties can hold relative to $H'$ but not $H$. The passages on an irreducible equation expressly retain $[H]$. The recorded relative phrases therefore carry mathematical information: the present po odnošenju k changes the wording without dropping the object of the relation.\par
\noindent{\footnotesize TEXT-COMP-112; pasaže / passages: 7199, 7201, 7350, 7628; \texttt{CORPUS\_FORM\_COMPARISONS.json}.}\par
\end{minipage}\par
\par\medskip\Needspace{7\baselineskip}
\noindent\begin{minipage}{\linewidth}
\hypertarget{TERM-0886}{}
\noindent\textbf{uravnenje nerazlo?imo vzhodno k $[H]$}\quad{\footnotesize\texttt{TERM-0886}}\par
\noindent\emph{Iztočnik / Source:} in bezug auf $[H]$ irreduzible Gleichung\par
\noindent\emph{Rusky / Russian:} [Bukvy ne sut zachovane / Lettering not recoverable]; $[H]$\par
\noindent\emph{Ukrajinsky / Ukrainian:} [Bukvy ne sut zachovane / Lettering not recoverable]; $[H]$\par
\smallskip\noindent\emph{Razlog:} Cěly izraz povezuje uravnenje, nerazložimost i oblast $[H]$. Zapisana motivacija želi obyčny sovrěmenny jezyk nerazložimosti, pri sohranjanju toj trojnoj svezi. No latiničny tekst imaje povrědu nerazlo?imo, a ruske i ukrajinske slova sut izgubjene; v njih ostaje čitateljny simbol $[H]$. Slovnik ne smi predstaviti domysljene bukvy kak zachovany prěvod. Nově čitany korpusny odlomok dava nerazložimo uravnenje po odnošenju k $[H]$; to jest odvojene dokazateljstvo aktualnoj formy, ne povratno dopolnjenje povrědženogo zapisa.\par
\noindent\emph{Reason:} The full expression connects an equation, irreducibility and the domain $[H]$. The recorded motivation favours ordinary contemporary irreducibility language while preserving that three-part relation. The Latin text is damaged at nerazlo?imo, however, and the Russian and Ukrainian words are lost; their symbol $[H]$ remains readable. The dictionary must not present conjectural letters as a surviving translation. The newly read corpus passage supplies nerazložimo uravnenje po odnošenju k $[H]$; this is separate evidence of the current form, not a retrospective filling of the damaged record.\par
\noindent{\footnotesize Izvorne zapisy / Original records: HIST-g063-007.}\par
\smallskip\noindent\emph{Sravnjenje s tekstom:} Rabota 9 razlikuje bazu $H$, oblast $[H]$ i drugu bazu $H'$. Primětka pokazuje, že bazove svojstva mogut važiti po odnošenju k $H'$, ale ne k $H$. Odlomky o nerazložimom uravnenju izrično sohranjajut $[H]$. Zato zapisane relativne frazy nesut matematičnu informaciju; nynějšny izraz po odnošenju k zaměnja formulaciju, ne izpustja objekt odnosa.\par
\noindent\emph{Textual comparison:} Paper 9 distinguishes the basis $H$, the domain $[H]$ and another basis $H'$. A remark shows that the basis properties can hold relative to $H'$ but not $H$. The passages on an irreducible equation expressly retain $[H]$. The recorded relative phrases therefore carry mathematical information: the present po odnošenju k changes the wording without dropping the object of the relation.\par
\noindent{\footnotesize TEXT-COMP-112; pasaže / passages: 7199, 7201, 7350, 7628; \texttt{CORPUS\_FORM\_COMPARISONS.json}.}\par
\end{minipage}\par
\par\medskip\Needspace{7\baselineskip}
\noindent\begin{minipage}{\linewidth}
\hypertarget{TERM-0887}{}
\noindent\textbf{vzhodno k glavnomu klasu}\quad{\footnotesize\texttt{TERM-0887}}\par
\noindent\emph{Iztočnik / Source:} in bezug auf die Hauptklasse\par
\noindent\emph{Rusky / Russian:} относительно главного класса\par
\noindent\emph{Ukrajinsky / Ukrainian:} відносно головного класу\par
\smallskip\noindent\emph{Razlog:} Vzhodno k glavnomu klasu sohranja referenciju k vybranomu klasu idealov, od ktorogo zaviset definicije slěda, normy i diskriminanta v rabotě 31. Glavny klas tu ne jest prosto najvažnějšij klas v obyčnom smyslu. Zapis objasnjaje kontinuitet s §4 i §6; to jest razlog, začto relativne doplnjenje ostaje v tekstu. Sam prědlog ne jest definitivny: v sravnjennyh odlomkah stojit po odnošenju k. Polna dokumentacija oddělja stabilny matematičny odnos od historijskej prědložnej formulacije, ktora byla predmetom pozdnějšego vyboru.\par
\noindent\emph{Reason:} Vzhodno k glavnomu klasu preserves reference to the selected ideal class on which the definitions of trace, norm and discriminant depend in Paper 31. Principal class here does not simply mean the most important class in ordinary speech. The record explains continuity with §§4 and 6, which is why the relative qualification remains in the text. The preposition itself is not final: the compared passages have po odnošenju k. Full documentation separates the stable mathematical relation from the historical prepositional wording considered in a later choice.\par
\noindent{\footnotesize Izvorne zapisy / Original records: HIST-g146-008.}\par
\smallskip\noindent\emph{Sravnjenje s tekstom:} Rabota 31 izrično kvalifikuje slěd po odnošenju k glavnomu klasu pri izčezanju na asociovanom prostom idealu. Pri formulah prěz konjugovane elementy tekst ponovno kaže, že slěd i norma sut vzete po odnošenju k tomu klasu. Točni odlomky podpirajut potrebu sohraniti kvalifikaciju; ne dopuščajut prevesti frazu kak neograničeno nazvanje slěda.\par
\noindent\emph{Textual comparison:} Paper 31 expressly qualifies trace relative to the principal class when discussing its vanishing on the associated prime ideal. In the formulas through conjugate elements it again specifies that trace and norm are taken relative to that class. These passages support retaining the qualification, not translating the phrase as an unrestricted name for trace.\par
\noindent{\footnotesize TEXT-COMP-113; pasaže / passages: 18310, 18383, 18474, 18513; \texttt{CORPUS\_FORM\_COMPARISONS.json}.}\par
\end{minipage}\par
\par\medskip\Needspace{7\baselineskip}
\noindent\begin{minipage}{\linewidth}
\hypertarget{TERM-0888}{}
\noindent\textbf{slěd, vzety vzhodno k glavnomu klasu}\quad{\footnotesize\texttt{TERM-0888}}\par
\noindent\emph{Iztočnik / Source:} in bezug auf die Hauptklasse genommene Spur\par
\noindent\emph{Rusky / Russian:} след, взятый относительно главного класса\par
\noindent\emph{Ukrajinsky / Ukrainian:} слід, узятий відносно головного класу\par
\smallskip\noindent\emph{Razlog:} Slěd, vzety vzhodno k glavnomu klasu sohranja dvě svezane věči: nazvanje slěd iz §5 i vybrany klas iz §4. Dva zachovane razloga podčrkivajut, že bez klasnogo doplnjenja čitatelj mogl by razuměti obyčny slěd bez ukazanja prědstavjenja. Jeden zapis takože govori o pečatnej formulaciji; to tvrdženje ne jest novo čitanje skana. Nynějšne odlomky raboty 31 izrično upotrěbjajut po odnošenju k glavnomu klasu, naprimer pri izčezanju slěda elementov asociovanogo prostogo ideala. Matematično ograničenje sohranja se, iako prědložny izraz menja se.\par
\noindent\emph{Reason:} Slěd, vzety vzhodno k glavnomu klasu retains two connected elements: the trace terminology from §5 and the selected class from §4. Both surviving rationales stress that without the class qualification a reader could understand an ordinary trace with no representation specified. One record also invokes the printed wording; that claim is not a new scan inspection. Present passages in Paper 31 expressly use po odnošenju k glavnomu klasu, for example when stating vanishing of the trace of elements of the associated prime ideal. The mathematical qualification survives even though the prepositional expression changes.\par
\noindent{\footnotesize Izvorne zapisy / Original records: HIST-g141-004, HIST-g143-005.}\par
\smallskip\noindent\emph{Sravnjenje s tekstom:} Rabota 31 izrično kvalifikuje slěd po odnošenju k glavnomu klasu pri izčezanju na asociovanom prostom idealu. Pri formulah prěz konjugovane elementy tekst ponovno kaže, že slěd i norma sut vzete po odnošenju k tomu klasu. Točni odlomky podpirajut potrebu sohraniti kvalifikaciju; ne dopuščajut prevesti frazu kak neograničeno nazvanje slěda.\par
\noindent\emph{Textual comparison:} Paper 31 expressly qualifies trace relative to the principal class when discussing its vanishing on the associated prime ideal. In the formulas through conjugate elements it again specifies that trace and norm are taken relative to that class. These passages support retaining the qualification, not translating the phrase as an unrestricted name for trace.\par
\noindent{\footnotesize TEXT-COMP-113; pasaže / passages: 18310, 18383, 18474, 18513; \texttt{CORPUS\_FORM\_COMPARISONS.json}.}\par
\end{minipage}\par
\par\medskip\Needspace{7\baselineskip}
\noindent\begin{minipage}{\linewidth}
\hypertarget{TERM-0889}{}
\noindent\textbf{vzhodno k $H$}\quad{\footnotesize\texttt{TERM-0889}}\par
\noindent\emph{Iztočnik / Source:} in bezug auf $H$\par
\noindent\emph{Rusky / Russian:} относительно $H$\par
\noindent\emph{Ukrajinsky / Ukrainian:} щодо $H$\par
\smallskip\noindent\emph{Razlog:} Vzhodno k $H$ označaje zavisnost bazovyh svojstv od izbranoj bazy. Zapis zavrgaje nejasnu formulaciju s $H$, zato že ona ne kaže, kak jest baza svazana s tvrdženjem. V rabotě 9 sravnjena primětka pokazyvaje oblast, ktora ne imaje danih svojstv po odnošenju k $H$, ale imaje jih po odnošenju k inoj bazi $H'$. Zato baza ne jest proizvoljna dekoracija frazy. Sohraniti toj odnos jest nužno; sohraniti staru formulaciju vzhodno k ne jest nužno, kako pokazyvaje pozdnějša prědložna odluka.\par
\noindent\emph{Reason:} Vzhodno k $H$ marks the dependence of basis properties on the chosen basis. The record rejects the vague wording with $H$ because it does not say how the basis relates to the assertion. A compared remark in Paper 9 gives a domain lacking the specified properties relative to $H$ but possessing them relative to another basis $H'$. The basis is therefore not incidental decoration. The relation must be retained; the earlier wording vzhodno k need not be, as the later prepositional choice shows.\par
\noindent{\footnotesize Izvorne zapisy / Original records: HIST-g059-005.}\par
\smallskip\noindent\emph{Sravnjenje s tekstom:} Rabota 9 razlikuje bazu $H$, oblast $[H]$ i drugu bazu $H'$. Primětka pokazuje, že bazove svojstva mogut važiti po odnošenju k $H'$, ale ne k $H$. Odlomky o nerazložimom uravnenju izrično sohranjajut $[H]$. Zato zapisane relativne frazy nesut matematičnu informaciju; nynějšny izraz po odnošenju k zaměnja formulaciju, ne izpustja objekt odnosa.\par
\noindent\emph{Textual comparison:} Paper 9 distinguishes the basis $H$, the domain $[H]$ and another basis $H'$. A remark shows that the basis properties can hold relative to $H'$ but not $H$. The passages on an irreducible equation expressly retain $[H]$. The recorded relative phrases therefore carry mathematical information: the present po odnošenju k changes the wording without dropping the object of the relation.\par
\noindent{\footnotesize TEXT-COMP-112; pasaže / passages: 7199, 7201, 7350, 7628; \texttt{CORPUS\_FORM\_COMPARISONS.json}.}\par
\end{minipage}\par
\par\medskip\Needspace{7\baselineskip}
\noindent\begin{minipage}{\linewidth}
\hypertarget{TERM-0890}{}
\noindent\textbf{elementy cěle vzhodno k $\mathfrak{o}$}\quad{\footnotesize\texttt{TERM-0890}}\par
\noindent\emph{Iztočnik / Source:} in bezug auf $\mathfrak{o}$ ganze Größen\par
\noindent\emph{Rusky / Russian:} элементы, целые относительно $\mathfrak{o}$\par
\noindent\emph{Ukrajinsky / Ukrainian:} елементи, цілі відносно $\mathfrak{o}$\par
\smallskip\noindent\emph{Razlog:} Elementy cěle vzhodno k $\mathfrak o$ sohranja osnovny prsten, po odnošenju k ktoromu cělost jest oprěděljena. Rabota 31 iz takyh elementov v razširjenom polju tvorji $\mathfrak O$ i potom oprědělja poredky. Cěly tu ne znači proizvoljno cělo čislo; smysl jest algebraična cělost nad ukazanym prstenom. V istorijskom spisku mo jest skračenym zapisom makra, ne slovom prěvoda. Nynějšny odlomok dava cělyh po odnošenju k $\mathfrak o$. Slovnik mora sohraniti i razlikovanje dvou prstenov, i relativny smysl cělosti.\par
\noindent\emph{Reason:} Elementy cěle vzhodno k $\mathfrak o$ retains the base ring relative to which integrality is defined. Paper 31 forms $\mathfrak O$ from such elements in the extension field and then defines orders. Integral here does not mean an arbitrary integer: it means algebraic integrality over the specified ring. In the historical list mo abbreviates a macro rather than a translated word. The present passage has cělyh po odnošenju k $\mathfrak o$. The dictionary must retain both the distinction between the rings and the relative sense of integrality.\par
\noindent{\footnotesize Izvorne zapisy / Original records: HIST-g149-004.}\par
\smallskip\noindent\emph{Sravnjenje s tekstom:} §7 raboty 31 oprědělja $\mathfrak O$ kak prsten elementov iz $K$, cělyh po odnošenju k $\mathfrak o$, pri $\Omega$ kak polju kvocientov $\mathfrak o$. §8 potom govori o linearnoj nezavisnosti nad $\Omega$. Trěba sohraniti i različne simbole, i različne odnose: cělost nad osnovnym prstenom ne jest isto, čto linearna nezavisnost nad poljem.\par
\noindent\emph{Textual comparison:} Paper 31, §7, defines $\mathfrak O$ as the ring of elements of $K$ integral relative to $\mathfrak o$, with $\Omega$ the fraction field of $\mathfrak o$. Section 8 then discusses linear independence over $\Omega$. Both the distinct symbols and the distinct relations must be retained: integrality over the base ring is not linear independence over the field.\par
\noindent{\footnotesize TEXT-COMP-114; pasaže / passages: 18559, 18857; \texttt{CORPUS\_FORM\_COMPARISONS.json}.}\par
\end{minipage}\par
\par\medskip\Needspace{7\baselineskip}
\noindent\begin{minipage}{\linewidth}
\hypertarget{TERM-0891}{}
\noindent\textbf{linearno nezavisnyh nad $\Omega$ elementov}\quad{\footnotesize\texttt{TERM-0891}}\par
\noindent\emph{Iztočnik / Source:} in bezug auf $\Omega$ linear unabhängigen Elementen\par
\noindent\emph{Rusky / Russian:} линейно независимых над $\Omega$ элементов\par
\noindent\emph{Ukrajinsky / Ukrainian:} лінійно незалежних над $\Omega$ елементів\par
\smallskip\noindent\emph{Razlog:} Linearno nezavisnyh nad $\Omega$ elementov prěvodit německe in bezug auf kratkym matematičnym nad. Oblast koeficientov linearnej relacije ostaje izrična; ničto v slově nezavisny ne daje prava ju izpustiti. Motivacija nazyva nad obyčnym sovrěmennym rěšenjem, aby izběgnuti težšu doslovnu konstrukciju. V §8 raboty 31 odlomok točno upotrěbja sistemu iz $n$ elementov, linearno nezavisnyh nad $\Omega$, za tvorjenje slěda. Soglasnost s odlomkom podpiraje smysl vyboru; sama ne jest sravniteljnym empirnym izslědovanjem stilja.\par
\noindent\emph{Reason:} Linearno nezavisnyh nad $\Omega$ elementov renders German in bezug auf with the compact mathematical over. The coefficient domain of the linear relation remains explicit; nothing in independent licenses its omission. The rationale calls nad a contemporary solution that avoids a cumbersome literal construction. In Paper 31, §8, the passage uses a system of $n$ elements linearly independent over $\Omega$ to form the trace. Agreement with the passage supports the choice's meaning; it is not itself an empirical comparative study of style.\par
\noindent{\footnotesize Izvorne zapisy / Original records: HIST-g157-004.}\par
\smallskip\noindent\emph{Sravnjenje s tekstom:} §7 raboty 31 oprědělja $\mathfrak O$ kak prsten elementov iz $K$, cělyh po odnošenju k $\mathfrak o$, pri $\Omega$ kak polju kvocientov $\mathfrak o$. §8 potom govori o linearnoj nezavisnosti nad $\Omega$. Trěba sohraniti i različne simbole, i različne odnose: cělost nad osnovnym prstenom ne jest isto, čto linearna nezavisnost nad poljem.\par
\noindent\emph{Textual comparison:} Paper 31, §7, defines $\mathfrak O$ as the ring of elements of $K$ integral relative to $\mathfrak o$, with $\Omega$ the fraction field of $\mathfrak o$. Section 8 then discusses linear independence over $\Omega$. Both the distinct symbols and the distinct relations must be retained: integrality over the base ring is not linear independence over the field.\par
\noindent{\footnotesize TEXT-COMP-114; pasaže / passages: 18559, 18857; \texttt{CORPUS\_FORM\_COMPARISONS.json}.}\par
\end{minipage}\par
\par\medskip\Needspace{7\baselineskip}
\noindent\begin{minipage}{\linewidth}
\hypertarget{TERM-0892}{}
\noindent\textbf{najnižše po koeficientah formy}\quad{\footnotesize\texttt{TERM-0892}}\par
\noindent\emph{Iztočnik / Source:} in den Koeffizienten niedrigsten Formen\par
\noindent\emph{Rusky / Russian:} формы, самые низкие по коэффициентам\par
\noindent\emph{Ukrajinsky / Ukrainian:} форми, найнижчі за коефіцієнтами\par
\smallskip\noindent\emph{Razlog:} Najnižše po koeficientah formy jest izrično provizorno rěšenje. Zapis govori, že točny invariantno-teoretičny poredok zavisi od konteksta i že pozdnějša tablica mora objasniti upotrěbu. Zato slovnik ne može iz samej frazy zaključiti, že ide o najmanjše absolutne cěnnosti koeficientov ili o najmanjši polinomny stepen. Rusky i ukrajinsky paralely sohranjajut podobnu doslovnost; tri slične formulacije ne rěšajut odloženo pytanje. Polna dokumentacija uključaje tu nejistotu i potrěbu identifikovati točny kriterij poredka.\par
\noindent\emph{Reason:} Najnižše po koeficientah formy is expressly provisional. The record says that the precise invariant-theoretic ordering depends on context and that a later table should clarify the usage. The phrase alone therefore cannot establish whether smallest absolute coefficient values or smallest polynomial degree are intended. The Russian and Ukrainian parallels retain similar literalness; three similar formulations do not resolve the deferred question. Complete documentation includes this uncertainty and the need to identify the exact ordering criterion.\par
\noindent{\footnotesize Izvorne zapisy / Original records: HIST-g007-003.}\par
\end{minipage}\par
\par\medskip\Needspace{7\baselineskip}
\noindent\begin{minipage}{\linewidth}
\hypertarget{TERM-0894}{}
\noindent\textbf{predstavjenje racionalno nad poljem $K$}\quad{\footnotesize\texttt{TERM-0894}}\par
\noindent\emph{Iztočnik / Source:} in einem Koerper $K$ rationale Darstellung\par
\noindent\emph{Rusky / Russian:} представление, рациональное над полем $K$\par
\noindent\emph{Ukrajinsky / Ukrainian:} представлення, раціональне над полем $K$\par
\smallskip\noindent\emph{Razlog:} Predstavjenje racionalno nad poljem $K$ sohranja $K$ kak komutativno osnovno polje. Zapis izrično vybiraje polje, ne tělo, zato že v tom odlomku ne ide o nekomutativnu strukturu s děljenjem. To razlikovanje jest osoblivo važno v kvaternionnom kontekstu raboty 32, kde različne oblasti mogut stojati blizko jedna drugoj. Racionalno nad $K$ takože ne znači prosto iměti racionalne čislene koeficienty: nazvane polje mora ostati čest definicije prědstavjenja. Jezykova politika ne opravdava globalnu zaměnu vsih německyh Körper slovom polje.\par
\noindent\emph{Reason:} Predstavjenje racionalno nad poljem $K$ retains $K$ as a commutative base field. The record expressly selects polje rather than tělo because this passage does not concern a noncommutative division structure. That distinction is particularly important in the quaternion setting of Paper 32, where different domains can occur close together. Rational over $K$ also does not simply mean having rational-number coefficients: the named field must remain part of the representation's definition. This linguistic policy does not justify replacing every German Körper with polje.\par
\noindent{\footnotesize Izvorne zapisy / Original records: HIST-g176-003.}\par
\end{minipage}\par
\par\medskip\Needspace{7\baselineskip}
\noindent\begin{minipage}{\linewidth}
\hypertarget{TERM-0895}{}
\noindent\textbf{prehoditi v sebe / sootvětovati samoj sobě}\quad{\footnotesize\texttt{TERM-0895}}\par
\noindent\emph{Iztočnik / Source:} in sich übergehen / entspricht sich selbst\par
\noindent\emph{Rusky / Russian:} переходить в себя / соответствовать самому себе\par
\noindent\emph{Ukrajinsky / Ukrainian:} переходити у себе / відповідати самому собі\par
\smallskip\noindent\emph{Razlog:} Prehoditi v sebe i sootvětovati samoj sobě sohranjajut avtorin jezyk samokorespondencije realnyh cěnnostij pri odobraženju $f$. V takom kontekstu čitatelj mora razlikovati element, ktory ostaje fiksovany, od množstva, ktoro samo prěhodi v sebe bez fiksovanja každogo elementa. Zapis nedava prava izravnati te dve tvrdženja; konkretna věta oprědělja predmet korespondencije. Dva glagolne izrazy sohranjajut stilističnu razliku originala. Motivacija ne soderži posebno normativno opravdanje korena sootvětovati, zato jeho zachovano upotrěbljenje ne jest samo po sobě jezykova rekomendacija.\par
\noindent\emph{Reason:} Prehoditi v sebe and sootvětovati samoj sobě preserve Noether's language of self-correspondence for real values under the map $f$. In this setting a reader must distinguish an element being fixed from a set being taken into itself without every element being fixed. The record does not license identifying those assertions; the particular sentence determines the object of the correspondence. The two verbal expressions retain a stylistic distinction in the original. The rationale supplies no separate normative justification for sootvětovati, so its recorded use is not by itself a linguistic recommendation.\par
\noindent{\footnotesize Izvorne zapisy / Original records: HIST-g072-011.}\par
\end{minipage}\par
\par\medskip\Needspace{7\baselineskip}
\noindent\begin{minipage}{\linewidth}
\hypertarget{TERM-0896}{}
\noindent\textbf{indukcija}\quad{\footnotesize\texttt{TERM-0896}}\par
\noindent\emph{Iztočnik / Source:} Induktionsschluss\par
\noindent\emph{Rusky / Russian:} индукция\par
\noindent\emph{Ukrajinsky / Ukrainian:} індукція\par
\smallskip\noindent\emph{Razlog:} Indukcija izbrana jest kak kratky izraz obyčnogo dokaznogo registra. Zapis izrično odbira težšu doslovnu parafrazu induktivny zaključok, zato že v odlomku dostavaje nazvanja metody. To skračenje ne znači skratiti sam dokaz: baza indukcije, parametr i prěhod ostajut potrěbne tam, kde argument jih potrěbuje. Ruska i ukrajinska paralela imaju isti medžunarodny koren. Smysl v rabotě 10 jest matematična indukcija, ne empirno obobčenje iz neskolikyh priměrov.\par
\noindent\emph{Reason:} Indukcija is selected as a compact expression in ordinary proof language. The record explicitly rejects the heavier literal paraphrase inductive conclusion because naming the method suffices in this passage. Shortening the designation does not shorten the proof: the induction base, parameter and step remain necessary where the argument requires them. The Russian and Ukrainian parallels share the international stem. In Paper 10 the intended sense is mathematical induction, not empirical generalization from a few examples.\par
\noindent{\footnotesize Izvorne zapisy / Original records: HIST-g071-007.}\par
\end{minipage}\par
\par\medskip\Needspace{7\baselineskip}
\noindent\begin{minipage}{\linewidth}
\hypertarget{TERM-0897}{}
\noindent\textbf{indukovano prěobraženje}\quad{\footnotesize\texttt{TERM-0897}}\par
\noindent\emph{Iztočnik / Source:} induzierte Transformation\par
\noindent\emph{Rusky / Russian:} индуцированное преобразование\par
\noindent\emph{Ukrajinsky / Ukrainian:} індуковане перетворення\par
\smallskip\noindent\emph{Razlog:} Indukovano prěobraženje označaje transformaciju oprěděljenu drugoju transformacijeju ili konstrukcijeju. Zapis želi sohraniti tu zavisnost v sovrěmennom matematičnom izloženju, ne vrativši se k staršemu nazvanju substitucija. To ne jest razlog izbrisati pravilo, ktorym odobraženje indukuje se. Prěobraženje daje slovjansky imennik pri medžunarodnom pridevniku; drugy zapis ima transformacije. Ta leksična variantnost ne mora značiti različny matematičny smysl, ale to trěba proveriti v odlomku. Konkrětna sravnitelna statistika za prěobraženje i transformacija ne jest dana.\par
\noindent\emph{Reason:} Indukovano prěobraženje denotes a transformation determined by another transformation or construction. The record seeks to preserve that dependence in contemporary mathematical exposition rather than revert to the older designation substitution. This is not a reason to omit the rule inducing the map. Prěobraženje supplies a Slavic noun with an international modifier; another record uses transformacije. This lexical variation need not signal a different mathematical sense, but that must be checked in context. No comparative statistics for prěobraženje and transformacija are supplied.\par
\noindent{\footnotesize Izvorne zapisy / Original records: HIST-g094-003.}\par
\end{minipage}\par
\par\medskip\Needspace{7\baselineskip}
\noindent\begin{minipage}{\linewidth}
\hypertarget{TERM-0898}{}
\noindent\textbf{indukovane transformacije / totalna diferencialna forma}\quad{\footnotesize\texttt{TERM-0898}}\par
\noindent\emph{Iztočnik / Source:} induzierte Transformationen / totale Differentialform\par
\noindent\emph{Rusky / Russian:} индуцированные преобразования / тотальная дифференциальная форма\par
\noindent\emph{Ukrajinsky / Ukrainian:} індуковані перетворення / тотальна диференціальна форма\par
\smallskip\noindent\emph{Razlog:} Indukovane transformacije i totalna diferencialna forma svezane sut s transformacijami koeficientov diferencialnyh form na koncu §6 raboty 13. Zapis sohranja medžunarodne pridevniki zaradi tehničnoj jasnosti i kontinuiteta s sravnjenymi jezykami. Indukovane ukazuje, že koeficienty ne transformujut se proizvoljno, ale po pravilah izvornogo prěobraženja. Totalna ne mora avtomatično značiti točna forma: takove razlučenje mora rešiti avtorina definicija, ne sličnost nazvanij. Opravdanje izbora pridevnikov ne jest dokaz, že vse podrobnosti formy i jej transformacije už sut proverjene.\par
\noindent\emph{Reason:} Indukovane transformacije and totalna diferencialna forma concern transformations of differential-form coefficients at the end of Paper 13, §6. The record retains international modifiers for technical clarity and continuity with the compared languages. Induced indicates that the coefficients do not transform arbitrarily but according to the original transformation's rules. Total must not automatically be read as exact: Noether's definition, not similarity between labels, must settle that distinction. Justifying the modifiers is not evidence that every detail of the form and its transformation has already been checked.\par
\noindent{\footnotesize Izvorne zapisy / Original records: HIST-g075-019.}\par
\end{minipage}\par
\par\medskip\Needspace{7\baselineskip}
\noindent\begin{minipage}{\linewidth}
\hypertarget{TERM-0899}{}
\noindent\textbf{vnutrišnje avtomorfizmy}\quad{\footnotesize\texttt{TERM-0899}}\par
\noindent\emph{Iztočnik / Source:} innere Automorphismen\par
\noindent\emph{Rusky / Russian:} внутренние автоморфизмы\par
\noindent\emph{Ukrajinsky / Ukrainian:} внутрішні автоморфізми\par
\smallskip\noindent\emph{Razlog:} Vnutrišnje avtomorfizmy sočeta razpoznavajemy slovjansky pridevnik s medžunarodnym nazvanjem odobraženja. V rabotě 34 perehod medžu dvoma razloženjami jest dan konjugacijeju; slovo vnutrišnje sohranja imenno to ograničenje, ne prosto fakt, že oba objekty ležat v jednoj strukturě. Odlomok izrično piše odobraženje s $x^{-1}$ i $x$. Ta formula objasnjaje smysl pridevnika. Tvrdžena razumlivost slovjanskogo korena jest kvalitativna motivacija vyboru; čitateljski test ili čislene vagi ne sut navedene.\par
\noindent\emph{Reason:} Vnutrišnje avtomorfizmy combines a recognizable Slavic modifier with the international name of the map. In Paper 34 the passage between two decompositions is given by conjugation; inner preserves that restriction, not simply the fact that both objects belong to one structure. The passage explicitly writes a map using $x^{-1}$ and $x$. This formula explains the adjective's meaning. The claimed intelligibility of the Slavic stem is a qualitative reason for the choice; no reader test or numerical weights are recorded.\par
\noindent{\footnotesize Izvorne zapisy / Original records: HIST-g192-006.}\par
\smallskip\noindent\emph{Sravnjenje s tekstom:} Teorem raboty 34 dava konjugaciju $\alpha'\mapsto x^{-1}\alpha'x$ pri prěhodu medžu dvoma razloženjami. To jest matematičny smysl pridevnika vnutrišnji v danom vyborě. Nije dostatočno rekti, že avtomorfizm jest něgde vnutri struktury; potrěbno sohraniti točno odobraženje.\par
\noindent\emph{Textual comparison:} The theorem in Paper 34 gives conjugation $\alpha'\mapsto x^{-1}\alpha'x$ in passing between two decompositions. This is the mathematical meaning of inner in the choice at issue. Saying that the automorphism is somewhere inside the structure is not enough; the specified map must be retained.\par
\noindent{\footnotesize TEXT-COMP-116; pasaže / passages: 21264; \texttt{CORPUS\_FORM\_COMPARISONS.json}.}\par
\end{minipage}\par
\par\medskip\Needspace{7\baselineskip}
\noindent\begin{minipage}{\linewidth}
\hypertarget{TERM-0900}{}
\noindent\textbf{uslovja integrabilnosti}\quad{\footnotesize\texttt{TERM-0900}}\par
\noindent\emph{Iztočnik / Source:} Integrabilitaetsbedingungen\par
\noindent\emph{Rusky / Russian:} условия интегрируемости\par
\noindent\emph{Ukrajinsky / Ukrainian:} умови інтегровності\par
\smallskip\noindent\emph{Razlog:} Uslovja integrabilnosti sohranja jezyk primětky, ktora objasnjaje redukcijny postupok prěz neizpolnjenje potrěbnyh uslovij za sistem diferencialnyh uravnenij. Potrěbnyh jest sučstveno: ne znači, že samo izpolnjenje bez ostalyh hipotez vsegda dava rěšenje. V sravnjenom tekstě racionalnost koeficienta izključuje izraz s $\log x$; slědujuči priměr potom adjunguje $\log x$ oblasti racionalnosti i opisuje drugu situaciju. Proto ne trěba čitati frazu kak bezuslovnu nerěšimost togo samogo diferencialnogo zapisa v vseh klasah funkcij. Izbor termina podpiraje smyslovu vez s sovrěmennom kompatibilnostju uravnenij.\par
\noindent\emph{Reason:} Uslovja integrabilnosti retains the footnote's language explaining the reduction procedure through failure of necessary conditions for a differential-equation system. Necessary matters: it does not say that satisfying these conditions alone always supplies a solution without further hypotheses. In the compared text, rationality of a coefficient excludes an expression involving $\log x$; the next example adjoins $\log x$ to the rationality domain and describes a different situation. The phrase must therefore not be read as unconditional insolubility of the same differential expressions in every class of functions. The choice connects the terminology with contemporary equation-compatibility language.\par
\noindent{\footnotesize Izvorne zapisy / Original records: HIST-g079-076.}\par
\smallskip\noindent\emph{Sravnjenje s tekstom:} Primětka raboty 17 opisuje redukciju kak neizpolnjenje potrěbnyh uslovij integrabilnosti. Okolny tekst izključuje koeficient s $\log x$ zaradi racionalnosti. Slědujuči priměr razširjaje oblast racionalnosti adjunkcijeju toj funkcije i govori o izpolnjenom uslovju. Ta razlika pokazuje, začto tvrdženje mora ostati svazano s ukazanoju oblastju funkcij.\par
\noindent\emph{Textual comparison:} Paper 17's footnote describes the reduction as failure of necessary integrability conditions. The surrounding text excludes a coefficient involving $\log x$ because it must be rational. The following example enlarges the rationality domain by adjoining that function and speaks of the condition being satisfied. This difference shows why the assertion must remain tied to the specified function domain.\par
\noindent{\footnotesize TEXT-COMP-117; pasaže / passages: 12283, 12285, 12287, 12289; \texttt{CORPUS\_FORM\_COMPARISONS.json}.}\par
\end{minipage}\par
\par\medskip\Needspace{7\baselineskip}
\noindent\begin{minipage}{\linewidth}
\hypertarget{TERM-0901}{}
\noindent\textbf{integraly sistema diferencialnyh uravnenj}\quad{\footnotesize\texttt{TERM-0901}}\par
\noindent\emph{Iztočnik / Source:} Integrale eines Differentialgleichungssystems\par
\noindent\emph{Rusky / Russian:} интегралы системы дифференциальных уравнений\par
\noindent\emph{Ukrajinsky / Ukrainian:} інтеграли системи диференціальних рівнянь\par
\smallskip\noindent\emph{Razlog:} Integraly sistema diferencialnyh uravnenj sohranja istoričnu upotrěbu slova integral za funkcije, ktore izpolnjajut uravnenja. V rabotě 17, §11, modul diferencialnyh izrazov svazuje se s konečnym sistemom uravnenij; rěč jest o njihovih simultannyh rěšenjah, ne o zapisu oprěděljenogo integrala. Sohranjenje nazvanja dopušča čitati avtorovu vez medžu algebroj i diferencialnymi uravnenjami bez tihoj modernizacije. Pojasnjenje rěšenja jest zato korisno, ale ne zaměnjaje iztočnogo termina. Zapis ne dava čislene vagi ni dokaz, že ta konstrukcija jest najprirodnějša dlja vseh slovjanskyh čitateljev.\par
\noindent\emph{Reason:} Integraly sistema diferencialnyh uravnenj preserves the historical use of integral for functions satisfying the equations. In Paper 17, §11, a module of differential expressions is related to a finite equation system; the objects are its simultaneous solutions, not expressions denoting definite integrals. Retaining the name lets readers follow the author's connection between algebra and differential equations without silent modernization. The explanatory gloss solutions is therefore useful, but does not replace the source term. The record supplies neither numerical weights nor evidence that this construction is the most idiomatic one for all Slavic readers.\par
\noindent{\footnotesize Izvorne zapisy / Original records: HIST-g079-066.}\par
\smallskip\noindent\emph{Sravnjenje s tekstom:} V §11 raboty 17 integral označaje rěšenje diferencialnyh uravnenij. Sravnjene odlomky govoret o simultannyh integralah i o čislu linearno nezavisnyh integralov, svazanom s redom grupy ostatkov. To ustanavjaje lokalny smysl, ne dokaz polnoj analitičnoj teorije egzistencije za libovoljne koeficienty.\par
\noindent\emph{Textual comparison:} In §11 of Paper 17, integral denotes a solution of differential equations. The compared passages discuss simultaneous integrals and the number of linearly independent integrals in relation to the order of the residue group. They establish the local sense, not a full analytic existence theorem for arbitrary coefficients.\par
\noindent{\footnotesize TEXT-COMP-119; pasaže / passages: 11978, 11986; \texttt{CORPUS\_FORM\_COMPARISONS.json}.}\par
\end{minipage}\par
\par\medskip\Needspace{7\baselineskip}
\noindent\begin{minipage}{\linewidth}
\hypertarget{TERM-0902}{}
\noindent\textbf{oblast c?losti $[H]$}\quad{\footnotesize\texttt{TERM-0902}}\par
\noindent\emph{Iztočnik / Source:} Integrit?tsbereich $[H]$\par
\noindent\emph{Rusky / Russian:} [Bukvy ne sut zachovane / Lettering not recoverable]; $[H]$\par
\noindent\emph{Ukrajinsky / Ukrainian:} [Bukvy ne sut zachovane / Lettering not recoverable]; $[H]$\par
\smallskip\noindent\emph{Razlog:} Oblast cělosti $[H]$ imenuje polinomovu oblast, ne polje $(H)$. To jest smysl zachovanoj motivacije za rabotu 9, §5. Něktore bukvy v německoj i medžuslovjanskoj natuknici sut izgubjene; v ruskoj i ukrajinskoj koloně ostalo jest čitljivo samo označenje $[H]$. Ty prazdniny ne sut jezykovymi variantami. Sravnjeni tekst objasnjaje $[H]$ kak oblast polinomov od bazovyh veličin s cělymi algebraičnymi koeficientami. Ta kontekstna podpora ne vozvrača izgubjene bukvy sravnjajučih prevodov. Kvadratne skobky sut čest razlikovanja matematičnyh objektov.\par
\noindent\emph{Reason:} Oblast cělosti $[H]$ names the polynomial domain, not the field $(H)$. That is the surviving rationale for Paper 9, §5. Some letters in the German and Interslavic headwords are lost; in the Russian and Ukrainian columns only the symbol $[H]$ remains legible. These gaps are not linguistic variants. The compared text describes $[H]$ as the domain of polynomials in the basis quantities with algebraic-integer coefficients. That contextual support does not recover the lost wording of the comparator translations. The square brackets contribute to the distinction between mathematical objects.\par
\noindent{\footnotesize Izvorne zapisy / Original records: HIST-g063-003.}\par
\smallskip\noindent\emph{Sravnjenje s tekstom:} Tekst raboty 9 opreděljaje $[H]$ polinomami od bazovyh veličin s cělymi algebraičnymi koeficientami. Odlomok s $[\Theta]$ izrično govori o cěločislenyh polinomah od $\vartheta$. Zato kvadratne skobky sohranjajut tip objekta. Čitljive matematične odlomky podpirajut smysl poškodženyh slovarnyh zapisov, ale ne vozstanavjajut njihove izgubjene ruske i ukrajinske bukvy.\par
\noindent\emph{Textual comparison:} Paper 9 defines $[H]$ through polynomials in the basis quantities with algebraic-integer coefficients. The passage with $[\Theta]$ explicitly refers to integer polynomials in the $\vartheta$. The square brackets therefore preserve the type of object. The legible mathematical passages support the sense of damaged dictionary records, but do not recover their lost Russian and Ukrainian letters.\par
\noindent{\footnotesize TEXT-COMP-122; pasaže / passages: 7199, 7749, 7750, 7751, 7752, 7753, 7754, 7755; \texttt{CORPUS\_FORM\_COMPARISONS.json}.}\par
\end{minipage}\par
\par\medskip\Needspace{7\baselineskip}
\noindent\begin{minipage}{\linewidth}
\hypertarget{TERM-0903}{}
\noindent\textbf{oblast c?losti $[\Theta]$}\quad{\footnotesize\texttt{TERM-0903}}\par
\noindent\emph{Iztočnik / Source:} Integrit?tsbereich $[\Theta]$\par
\noindent\emph{Rusky / Russian:} [Bukvy ne sut zachovane / Lettering not recoverable]; $[\Theta]$\par
\noindent\emph{Ukrajinsky / Ukrainian:} [Bukvy ne sut zachovane / Lettering not recoverable]; $[\Theta]$\par
\smallskip\noindent\emph{Razlog:} Za rabotu 9, §6, motivacija vybiraje oblast cělosti $[\Theta]$, zato že objekt sostoji se iz cěločislenyh polinomov od bazovyh veličin. Polje porodženo bazoju jest drugy objekt. Arhivny zapis dava Theta slovami; sravnjeni matematičny tekst piše grčsko $\Theta$. Bukvy v někotoryh slovah natuknice i sravnjajučih prevodov sut poškodžene, i ne trěba jih dopolnjati kak by byly zachovane. Smyslova rekonstrukcija jest vozmožna iz motivacije i konteksta, ale polna leksična rekonstrukcija ruskoj i ukrajinskoj formulacije iz togo zapisa ne jest vozmožna.\par
\noindent\emph{Reason:} For Paper 9, §6, the rationale selects oblast cělosti $[\Theta]$ because the object consists of integer polynomials in basis quantities. The field generated by the basis is a different object. The archival record spells out Theta; the compared mathematical text uses Greek $\Theta$. Letters in some headword and comparator words are damaged, and must not be filled in as though they had survived. The rationale and context permit semantic reconstruction, but not full recovery of the Russian and Ukrainian wording from this record.\par
\noindent{\footnotesize Izvorne zapisy / Original records: HIST-g064-006.}\par
\smallskip\noindent\emph{Sravnjenje s tekstom:} Tekst raboty 9 opreděljaje $[H]$ polinomami od bazovyh veličin s cělymi algebraičnymi koeficientami. Odlomok s $[\Theta]$ izrično govori o cěločislenyh polinomah od $\vartheta$. Zato kvadratne skobky sohranjajut tip objekta. Čitljive matematične odlomky podpirajut smysl poškodženyh slovarnyh zapisov, ale ne vozstanavjajut njihove izgubjene ruske i ukrajinske bukvy.\par
\noindent\emph{Textual comparison:} Paper 9 defines $[H]$ through polynomials in the basis quantities with algebraic-integer coefficients. The passage with $[\Theta]$ explicitly refers to integer polynomials in the $\vartheta$. The square brackets therefore preserve the type of object. The legible mathematical passages support the sense of damaged dictionary records, but do not recover their lost Russian and Ukrainian letters.\par
\noindent{\footnotesize TEXT-COMP-122; pasaže / passages: 7199, 7749, 7750, 7751, 7752, 7753, 7754, 7755; \texttt{CORPUS\_FORM\_COMPARISONS.json}.}\par
\end{minipage}\par
\par\medskip\Needspace{7\baselineskip}
\noindent\begin{minipage}{\linewidth}
\hypertarget{TERM-0904}{}
\noindent\textbf{oblast cělosti}\quad{\footnotesize\texttt{TERM-0904}}\par
\noindent\emph{Iztočnik / Source:} Integritaetsbereich\par
\noindent\emph{Rusky / Russian:} область целостности\par
\noindent\emph{Ukrajinsky / Ukrainian:} область цілісності\par
\smallskip\noindent\emph{Razlog:} Oblast cělosti v rabotě 17 označaje polinomovu oblast $J$, pri kotoroj ne prědpokladaje se komutativny zakon množenja. Proto sovrěmenno naviknuty smysl komutativnoj integralnoj oblasti ne směje tiho dobaviti novu hipotezu. Motivacija sohranja istoričny Bereich i ne zaměnjaje jego automatično slovom prsten; naglas na otsutstvu nulovyh děliteljev ostaje svazany s konkretnym objektom. V izloženju $J$ nastava posrědstvom adjungovanja neizvěstnyh k polju $P$. Neprědpoložena komutativnost jest matematično važne ograničenje analogije, ne pravopisna različnost.\par
\noindent\emph{Reason:} Oblast cělosti in Paper 17 denotes a polynomial domain $J$ for which the commutative law of multiplication is not assumed. The familiar modern reading of integral domain as commutative must therefore not silently add a hypothesis. The rationale preserves historical Bereich instead of automatically replacing it with ring; its emphasis on absence of zero divisors remains tied to the particular object. In the exposition, $J$ is formed by adjoining indeterminates to a field $P$. Non-assumed commutativity is a mathematically important limit on the analogy, not a spelling difference.\par
\noindent{\footnotesize Izvorne zapisy / Original records: HIST-g079-009.}\par
\smallskip\noindent\emph{Sravnjenje s tekstom:} Rabota 6 stroji integralnu oblast iz sistemy racionalnyh funkcij posrědstvom zatvorenosti pod skalarnym množenjem, dodavanjem i množenjem; dopuščene kvocienty dajut polje. Rabota 17 vvodži oblast cělosti $J$ bez prědpoloženja komutativnogo množenja. Rabota 19 razlikuje obči prsten od vlastnoj oblasti cělosti s jediniceju i bez nulovyh děliteljev. Te sravnjenja ne razrěšajut ustaliti vse pojavy jedinoju slovarnoju definicijeju.\par
\noindent\emph{Textual comparison:} Paper 6 constructs an integralna oblast from a system of rational functions by closure under scalar multiplication, addition and multiplication; admitting quotients gives a field. Paper 17 introduces oblast cělosti $J$ without assuming commutative multiplication. Paper 19 distinguishes a general ring from a proper domain of integrity with identity and without zero divisors. These comparisons do not license imposing one dictionary definition on every occurrence.\par
\noindent{\footnotesize TEXT-COMP-120; pasaže / passages: 5740, 5744, 5746, 5748, 5750, 5752, 5754, 11044, 12485; \texttt{CORPUS\_FORM\_COMPARISONS.json}.}\par
\end{minipage}\par
\par\medskip\Needspace{7\baselineskip}
\noindent\begin{minipage}{\linewidth}
\hypertarget{TERM-0905}{}
\noindent\textbf{integralna baza}\quad{\footnotesize\texttt{TERM-0905}}\par
\noindent\emph{Iztočnik / Source:} Integritätsbasis\par
\noindent\emph{Rusky / Russian:} интегральная база; интегральный базис\par
\noindent\emph{Ukrajinsky / Ukrainian:} інтегральна база; інтегральний базис\par
\smallskip\noindent\emph{Razlog:} Integralna baza svazuje Integritätsbasis s algebraičnoju cělostju, ne s moralnoju ili pravnoju integritetnostju. Šest motivacij pokryvajut rabotu 5, uvod raboty 6, jej §§10 i 12, i raboty 23 i 29. V rabotě 6 definicija izrično potrěbuje konečno mnogo polinomov, prěz ktore vsaky polinom oblasti jest cěla racionalna kombinacija s koeficientami iz $\Omega$. Tu baza znači polinomove generatory: ne trěba dobaviti linearnu nezavisnost ni čitati termin kak svobodnu cěločislenu bazu. §12 svazuje to nazvanje s teoremami VIII i IX. Ruske zapisy kolebajut se medžu baza i bazis, pri čem jedna motivacija preferuje bazis kak idiomatičny; ukrajinske zapisy takože imajut oba vida. Medžuslovjanska odluka daje prědnost kontinuitetu integralna baza. To jest dokumentovany kvalitativny izbor, ne izčisljeny rezultat sravnjenja.\par
\noindent\emph{Reason:} Integralna baza links Integritätsbasis to algebraic integrality, not moral or legal integrity. Six rationales cover Paper 5, the introduction to Paper 6, its §§10 and 12, and Papers 23 and 29. In Paper 6 the definition explicitly requires finitely many polynomials through which every polynomial of the domain is a whole rational combination with coefficients in $\Omega$. Here basis means polynomial generators: neither linear independence nor a free integral-module basis should be added. Section 12 connects the name to Theorems VIII and IX. The Russian records alternate between baza and bazis, with one rationale preferring bazis as idiomatic; Ukrainian records also contain both forms. The Interslavic decision prioritizes continuity of integralna baza. This is a documented qualitative choice, not a calculated comparison result.\par
\noindent{\footnotesize Izvorne zapisy / Original records: HIST-g039-009, HIST-g040-007, HIST-g050-002, HIST-g052-002, HIST-g094-006, HIST-g106-005.}\par
\smallskip\noindent\emph{Sravnjenje s tekstom:} Definicija VI raboty 6 potrěbuje polinomove kombinacije konečno mnogih generatorov s koeficientami iz $\Omega$, ne linearne kombinacije nezavisnyh vektorov. Teorema VII utočnjaje uslovja, pri ktoryh involucijska baza jest integralna baza. Rabota 35 upotrěbjaje integritetna baza pri porodženju cěločislenymi polinomami. Blizkost nazvanj ne izključuje razlikovanja koeficientnyh oblastij i operacij.\par
\noindent\emph{Textual comparison:} Definition VI in Paper 6 requires polynomial combinations of finitely many generators with coefficients in $\Omega$, not linear combinations of independent vectors. Theorem VII specifies conditions under which an involution basis is an integral basis. Paper 35 uses integritetna baza for generation by integer polynomials. The related names do not remove distinctions between coefficient domains and operations.\par
\noindent{\footnotesize TEXT-COMP-121; pasaže / passages: 6061, 6069, 22583; \texttt{CORPUS\_FORM\_COMPARISONS.json}.}\par
\end{minipage}\par
\par\medskip\Needspace{7\baselineskip}
\noindent\begin{minipage}{\linewidth}
\hypertarget{TERM-0906}{}
\noindent\textbf{integralna oblast}\quad{\footnotesize\texttt{TERM-0906}}\par
\noindent\emph{Iztočnik / Source:} Integritätsbereich\par
\noindent\emph{Rusky / Russian:} область целостности\par
\noindent\emph{Ukrajinsky / Ukrainian:} інтегральна область\par
\smallskip\noindent\emph{Razlog:} Integralna oblast jest izbor raboty 6, §7, povezany s uvodnym nazvanjem i s algebraičnym kontekstom. Iztočny predmet jest sistem racionalnyh funkcij zatvorjeny pod skalarnym množenjem, dodavanjem i množenjem funkcij; daljše dopuščenje kvocientov dava polje. Zato pridevnik integralna tu ne govori o integralnom računu. Ruski sravnjajuči izraz imaje konstrukciju oblast celostnosti, ukrajinski integralna oblast jest bližši medžuslovjanskomu vidu. Motivacija nazyvaje termin standardnym, ale sama ne dava odlomok iz nezavisnoj medžuslovjanskoj matematičnoj publikacije.\par
\noindent\emph{Reason:} Integralna oblast is the choice in Paper 6, §7, linked to the introductory terminology and algebraic context. The object is a system of rational functions closed under scalar multiplication, addition and multiplication of functions; further admitting quotients gives a field. The adjective integralna therefore does not refer to integral calculus here. The Russian comparator has a domain-of-integrity construction, while Ukrainian integralna oblast is closer to the Interslavic form. The rationale calls the terminology standard, but does not itself supply a passage from an independent Interslavic mathematical publication.\par
\noindent{\footnotesize Izvorne zapisy / Original records: HIST-g047-003.}\par
\smallskip\noindent\emph{Sravnjenje s tekstom:} Rabota 6 stroji integralnu oblast iz sistemy racionalnyh funkcij posrědstvom zatvorenosti pod skalarnym množenjem, dodavanjem i množenjem; dopuščene kvocienty dajut polje. Rabota 17 vvodži oblast cělosti $J$ bez prědpoloženja komutativnogo množenja. Rabota 19 razlikuje obči prsten od vlastnoj oblasti cělosti s jediniceju i bez nulovyh děliteljev. Te sravnjenja ne razrěšajut ustaliti vse pojavy jedinoju slovarnoju definicijeju.\par
\noindent\emph{Textual comparison:} Paper 6 constructs an integralna oblast from a system of rational functions by closure under scalar multiplication, addition and multiplication; admitting quotients gives a field. Paper 17 introduces oblast cělosti $J$ without assuming commutative multiplication. Paper 19 distinguishes a general ring from a proper domain of integrity with identity and without zero divisors. These comparisons do not license imposing one dictionary definition on every occurrence.\par
\noindent{\footnotesize TEXT-COMP-120; pasaže / passages: 5740, 5744, 5746, 5748, 5750, 5752, 5754, 11044, 12485; \texttt{CORPUS\_FORM\_COMPARISONS.json}.}\par
\end{minipage}\par
\par\medskip\Needspace{7\baselineskip}
\noindent\begin{minipage}{\linewidth}
\hypertarget{TERM-0907}{}
\noindent\textbf{oblast cělosti}\quad{\footnotesize\texttt{TERM-0907}}\par
\noindent\emph{Iztočnik / Source:} Integritätsbereich\par
\noindent\emph{Rusky / Russian:} область целостности\par
\noindent\emph{Ukrajinsky / Ukrainian:} область цілісності\par
\smallskip\noindent\emph{Razlog:} Oblast cělosti imaje šest zachovanyh obosnovanj. V uvodu i §4 raboty 9 ona sohranja slovjansku konstrukciju oblast cělosti vměsto latinizovanoj integralna oblast; v §10 jest svazana s izčerpavanjem obkladajučego polja kvocientami. Uvod raboty 10 i raboty 23 i 29 podpirajut kontinuitet uporěby. Jedna motivacija izrično sohranja literalny, ješče potrěbujuči jezykovu ocěnku vid, vměsto nezasvědčenogo novotvora. Odkazy na ranějše raboty ne dokazujut, že vsude byla ta sama forma: rabota 6 upotrěbjaje takože integralna oblast. Slovnik mora pokazati oba vybory i jih konteksty, ne prědstaviti retrospektivnu soglasnost, ktorej zapisy ne dokazujut.\par
\noindent\emph{Reason:} Oblast cělosti has six preserved rationales. In the introduction and §4 of Paper 9 it retains a Slavic domain-of-integrity construction instead of Latin-derived integralna oblast; in §10 it is linked to exhausting the ambient field by quotients. The opening of Paper 10 and Papers 23 and 29 support continuity of usage. One rationale expressly retains a literal form still requiring linguistic appraisal rather than an unattested coinage. References to earlier papers do not prove that they all used this same form: Paper 6 also uses integralna oblast. The dictionary must show both choices and their contexts, not impose retrospective agreement unsupported by the records.\par
\noindent{\footnotesize Izvorne zapisy / Original records: HIST-g058-006, HIST-g062-003, HIST-g068-008, HIST-g069-011, HIST-g094-004, HIST-g106-004.}\par
\smallskip\noindent\emph{Sravnjenje s tekstom:} Rabota 6 stroji integralnu oblast iz sistemy racionalnyh funkcij posrědstvom zatvorenosti pod skalarnym množenjem, dodavanjem i množenjem; dopuščene kvocienty dajut polje. Rabota 17 vvodži oblast cělosti $J$ bez prědpoloženja komutativnogo množenja. Rabota 19 razlikuje obči prsten od vlastnoj oblasti cělosti s jediniceju i bez nulovyh děliteljev. Te sravnjenja ne razrěšajut ustaliti vse pojavy jedinoju slovarnoju definicijeju.\par
\noindent\emph{Textual comparison:} Paper 6 constructs an integralna oblast from a system of rational functions by closure under scalar multiplication, addition and multiplication; admitting quotients gives a field. Paper 17 introduces oblast cělosti $J$ without assuming commutative multiplication. Paper 19 distinguishes a general ring from a proper domain of integrity with identity and without zero divisors. These comparisons do not license imposing one dictionary definition on every occurrence.\par
\noindent{\footnotesize TEXT-COMP-120; pasaže / passages: 5740, 5744, 5746, 5748, 5750, 5752, 5754, 11044, 12485; \texttt{CORPUS\_FORM\_COMPARISONS.json}.}\par
\end{minipage}\par
\par\medskip\Needspace{7\baselineskip}
\noindent\begin{minipage}{\linewidth}
\hypertarget{TERM-0908}{}
\noindent\textbf{integritetna oblast / kolco bez nulovyh děliteljev}\quad{\footnotesize\texttt{TERM-0908}}\par
\noindent\emph{Iztočnik / Source:} Integritätsbereich (Ring ohne Nullteiler)\par
\noindent\emph{Rusky / Russian:} область целостности / кольцо без делителей нуля\par
\noindent\emph{Ukrajinsky / Ukrainian:} область цілісності / кільце без дільників нуля\par
\smallskip\noindent\emph{Razlog:} Integritetna oblast dopolnjena pojasnjenjem kolco bez nulovyh děliteljev jest dvojny zapis za rabotu 35: medžunarodny kalek i opis matematičnogo svojstva. Pojasnjenje ne jest nova hipoteza dodana prevoditeljem; motivacija svazuje jego s iztočnym Ring ohne Nullteiler. V sravnjenom tekstě jest prsten bez nulovyh děliteljev, a ne arhivno kolco. To sravnjenje ustanavjaje leksičnu različnost medžu zapisom izbora i tekstom, ne proměnu matematičnogo svojstva. Sama podobnost integritetna s německym slovom ne dokazuje oficialnu medžuslovjansku upotrěbu.\par
\noindent\emph{Reason:} Integritetna oblast with the explanation kolco bez nulovyh děliteljev is a two-part record for Paper 35: an international calque and a description of the mathematical property. The explanation is not a new hypothesis supplied by the translator; the rationale connects it to the source phrase Ring ohne Nullteiler. The compared text reads prsten bez nulovyh děliteljev, rather than archival kolco. This establishes a lexical difference between the choice record and the text, not a change in the mathematical property. Resemblance of integritetna to the German word does not establish official Interslavic usage.\par
\noindent{\footnotesize Izvorne zapisy / Original records: HIST-g205-002.}\par
\smallskip\noindent\emph{Sravnjenje s tekstom:} V rabotě 35 tekst dava integritetnu oblast i tutčas pojasnjaje ju kak prsten bez nulovyh děliteljev. Slovarny zapis imaje kolco v tom pojasnjenju. Ustanovjena različnost otnosi se k imeniku za prsten; ne jest osnovom za tvrdženje, že matematično uslovje bylo proměnjeno.\par
\noindent\emph{Textual comparison:} In Paper 35 the text gives integritetna oblast and immediately explains it as prsten bez nulovyh děliteljev. The dictionary record has kolco in that explanation. The established difference concerns the noun for ring; it is not grounds for claiming that the mathematical condition changed.\par
\noindent{\footnotesize TEXT-COMP-123; pasaže / passages: 22576; \texttt{CORPUS\_FORM\_COMPARISONS.json}.}\par
\end{minipage}\par
\par\medskip\Needspace{7\baselineskip}
\noindent\begin{minipage}{\linewidth}
\hypertarget{TERM-0909}{}
\noindent\textbf{oblast cělosti}\quad{\footnotesize\texttt{TERM-0909}}\par
\noindent\emph{Iztočnik / Source:} Integritätsbereich / Bereich der Integrität\par
\noindent\emph{Rusky / Russian:} область целостности\par
\noindent\emph{Ukrajinsky / Ukrainian:} область цілісності\par
\smallskip\noindent\emph{Razlog:} Oblast cělosti obhvata dvě německe formulacije, Integritätsbereich i Bereich der Integrität, v rabotě 19. Motivacija obhvača zakres do §§6 i 12 i razlikuje oblasti cělosti od občejših prstenovyh oblastij. V uvodnom oprěděljenju vlastna oblast cělosti imaje jedinicu i ne imaje nulovyh děliteljev, dokolě obči prsten može ne iměti jedinicu i može iměti nulove dělitelje. Tak pojasnjen izbor sohranja avtorovo razlikovanje; samo leksično zjednoobraznjenje vsih pojavjenij Bereich zamazalo by jeho. Dvě německe formulacije ne potrěbujut dvě nezasvědčene medžuslovjanske novotvorne nazvanja.\par
\noindent\emph{Reason:} Oblast cělosti covers two German formulations, Integritätsbereich and Bereich der Integrität, in Paper 19. The rationale spans the material through §§6 and 12 and distinguishes domains of integrity from the more general ring domains. In the introductory definition, a proper domain of integrity has an identity and no zero divisors, whereas a general ring may lack an identity and have zero divisors. The contextualized choice preserves the author's distinction; lexical uniformity imposed on every occurrence of Bereich would obscure it. The two German formulations do not require two unattested Interslavic coinages.\par
\noindent{\footnotesize Izvorne zapisy / Original records: HIST-g081-002, HIST-g082-002.}\par
\smallskip\noindent\emph{Sravnjenje s tekstom:} Rabota 6 stroji integralnu oblast iz sistemy racionalnyh funkcij posrědstvom zatvorenosti pod skalarnym množenjem, dodavanjem i množenjem; dopuščene kvocienty dajut polje. Rabota 17 vvodži oblast cělosti $J$ bez prědpoloženja komutativnogo množenja. Rabota 19 razlikuje obči prsten od vlastnoj oblasti cělosti s jediniceju i bez nulovyh děliteljev. Te sravnjenja ne razrěšajut ustaliti vse pojavy jedinoju slovarnoju definicijeju.\par
\noindent\emph{Textual comparison:} Paper 6 constructs an integralna oblast from a system of rational functions by closure under scalar multiplication, addition and multiplication; admitting quotients gives a field. Paper 17 introduces oblast cělosti $J$ without assuming commutative multiplication. Paper 19 distinguishes a general ring from a proper domain of integrity with identity and without zero divisors. These comparisons do not license imposing one dictionary definition on every occurrence.\par
\noindent{\footnotesize TEXT-COMP-120; pasaže / passages: 5740, 5744, 5746, 5748, 5750, 5752, 5754, 11044, 12485; \texttt{CORPUS\_FORM\_COMPARISONS.json}.}\par
\end{minipage}\par
\par\medskip\Needspace{7\baselineskip}
\noindent\begin{minipage}{\linewidth}
\hypertarget{TERM-0910}{}
\noindent\textbf{oblast cělosti $[H]$}\quad{\footnotesize\texttt{TERM-0910}}\par
\noindent\emph{Iztočnik / Source:} Integritätsbereich $[H]$\par
\noindent\emph{Rusky / Russian:} область целостности $[H]$\par
\noindent\emph{Ukrajinsky / Ukrainian:} область цілісності $[H]$\par
\smallskip\noindent\emph{Razlog:} Oblast cělosti $[H]$ prodolžuje terminologiju Zermelovogo konteksta v rabotě 9, §1. Kvadratne skobky sohranjajut oblast, otličnu od pridruženogo polja. Motivacija takože izrično potrěbuje sravniti terminologiju raboty 6 prěd občim ustaljenjem. To potrěbovanje ne jest dokazom, že sravnjenje bylo togda zakončeno ili že vse varianty sut ravnoznačne vo vseh kontekstah. Sravnjene definicije nyně pokazyvajut motivaciju dvoh vidov, integralna oblast i oblast cělosti; one ne ustanavjajut jedinu nezavisno zasvědčenu medžuslovjansku normu.\par
\noindent\emph{Reason:} Oblast cělosti $[H]$ continues the terminology of the Zermelo context in Paper 9, §1. Square brackets preserve the domain, distinct from the associated field. The rationale also explicitly calls for comparison with Paper 6 before general standardization. That request is not evidence that the comparison was completed at the time, or that every variant is interchangeable in every context. The definitions compared here explain the motivation of integralna oblast and oblast cělosti; they do not establish a single independently attested Interslavic norm.\par
\noindent{\footnotesize Izvorne zapisy / Original records: HIST-g059-003.}\par
\smallskip\noindent\emph{Sravnjenje s tekstom:} Rabota 6 stroji integralnu oblast iz sistemy racionalnyh funkcij posrědstvom zatvorenosti pod skalarnym množenjem, dodavanjem i množenjem; dopuščene kvocienty dajut polje. Rabota 17 vvodži oblast cělosti $J$ bez prědpoloženja komutativnogo množenja. Rabota 19 razlikuje obči prsten od vlastnoj oblasti cělosti s jediniceju i bez nulovyh děliteljev. Te sravnjenja ne razrěšajut ustaliti vse pojavy jedinoju slovarnoju definicijeju.\par
\noindent\emph{Textual comparison:} Paper 6 constructs an integralna oblast from a system of rational functions by closure under scalar multiplication, addition and multiplication; admitting quotients gives a field. Paper 17 introduces oblast cělosti $J$ without assuming commutative multiplication. Paper 19 distinguishes a general ring from a proper domain of integrity with identity and without zero divisors. These comparisons do not license imposing one dictionary definition on every occurrence.\par
\noindent{\footnotesize TEXT-COMP-120; pasaže / passages: 5740, 5744, 5746, 5748, 5750, 5752, 5754, 11044, 12485; \texttt{CORPUS\_FORM\_COMPARISONS.json}.}\par
\smallskip\noindent\emph{Sravnjenje s tekstom:} Tekst raboty 9 opreděljaje $[H]$ polinomami od bazovyh veličin s cělymi algebraičnymi koeficientami. Odlomok s $[\Theta]$ izrično govori o cěločislenyh polinomah od $\vartheta$. Zato kvadratne skobky sohranjajut tip objekta. Čitljive matematične odlomky podpirajut smysl poškodženyh slovarnyh zapisov, ale ne vozstanavjajut njihove izgubjene ruske i ukrajinske bukvy.\par
\noindent\emph{Textual comparison:} Paper 9 defines $[H]$ through polynomials in the basis quantities with algebraic-integer coefficients. The passage with $[\Theta]$ explicitly refers to integer polynomials in the $\vartheta$. The square brackets therefore preserve the type of object. The legible mathematical passages support the sense of damaged dictionary records, but do not recover their lost Russian and Ukrainian letters.\par
\noindent{\footnotesize TEXT-COMP-122; pasaže / passages: 7199, 7749, 7750, 7751, 7752, 7753, 7754, 7755; \texttt{CORPUS\_FORM\_COMPARISONS.json}.}\par
\end{minipage}\par
\par\medskip\Needspace{7\baselineskip}
\noindent\begin{minipage}{\linewidth}
\hypertarget{TERM-0911}{}
\noindent\textbf{oblast cělosti $[\Theta]$}\quad{\footnotesize\texttt{TERM-0911}}\par
\noindent\emph{Iztočnik / Source:} Integritätsbereich $[\Theta]$\par
\noindent\emph{Rusky / Russian:} область целостности $[\Theta]$\par
\noindent\emph{Ukrajinsky / Ukrainian:} область цілісності $[\Theta]$\par
\smallskip\noindent\emph{Razlog:} Oblast cělosti $[\Theta]$ prodolžuje razlikovanje polinomovoj oblasti i porodženogo polja v rabotě 9, §7. Medžu tym zapisom i poškodženym zapisom iz §6 jest kontinuitet matematičnogo objekta, ale ne dokaz, že izgubjene bukvy možut byti vozstanovjene iz pozdnějšej kolonki. Ruski i ukrajinski sravnjajuči vidy zdes sut čitljive i podpirajut konstrukciju s oblastju cělosti. Zapis Theta slovami v sravnjenoj tablici služi za grčsko $\Theta$ v tekstě. Sohranjenje skobok jest ne menje važno od izbora pridevnika.\par
\noindent\emph{Reason:} Oblast cělosti $[\Theta]$ continues the distinction between the polynomial domain and generated field in Paper 9, §7. Its relation to the damaged §6 record establishes continuity of the mathematical object, not proof that missing letters can be recovered from a later column. Here the Russian and Ukrainian comparators are legible and support the domain-of-integrity construction. Spelled-out Theta in the comparison table stands for Greek $\Theta$ in the text. Preserving the brackets matters as much as choosing the adjective.\par
\noindent{\footnotesize Izvorne zapisy / Original records: HIST-g065-007.}\par
\smallskip\noindent\emph{Sravnjenje s tekstom:} Tekst raboty 9 opreděljaje $[H]$ polinomami od bazovyh veličin s cělymi algebraičnymi koeficientami. Odlomok s $[\Theta]$ izrično govori o cěločislenyh polinomah od $\vartheta$. Zato kvadratne skobky sohranjajut tip objekta. Čitljive matematične odlomky podpirajut smysl poškodženyh slovarnyh zapisov, ale ne vozstanavjajut njihove izgubjene ruske i ukrajinske bukvy.\par
\noindent\emph{Textual comparison:} Paper 9 defines $[H]$ through polynomials in the basis quantities with algebraic-integer coefficients. The passage with $[\Theta]$ explicitly refers to integer polynomials in the $\vartheta$. The square brackets therefore preserve the type of object. The legible mathematical passages support the sense of damaged dictionary records, but do not recover their lost Russian and Ukrainian letters.\par
\noindent{\footnotesize TEXT-COMP-122; pasaže / passages: 7199, 7749, 7750, 7751, 7752, 7753, 7754, 7755; \texttt{CORPUS\_FORM\_COMPARISONS.json}.}\par
\end{minipage}\par
\par\medskip\Needspace{7\baselineskip}
\noindent\begin{minipage}{\linewidth}
\hypertarget{TERM-0912}{}
\noindent\textbf{integralna oblast iz polinomov}\quad{\footnotesize\texttt{TERM-0912}}\par
\noindent\emph{Iztočnik / Source:} Integritätsbereich aus Polynomen\par
\noindent\emph{Rusky / Russian:} область целостности из полиномов\par
\noindent\emph{Ukrajinsky / Ukrainian:} інтегральна область з поліномів\par
\smallskip\noindent\emph{Razlog:} Integralna oblast iz polinomov sohranja termin integralna oblast iz §§7 i 8 raboty 6 v §9, ale ne izpušča ograničenja na polinomy. Oblast racionalnyh funkcij i oblast iz polinomov ne sut prosto dvě nazvanja togo samogo klasa objektov. Konstrukcija iz polinomov zato nosi matematičnu informaciju, ne samo stilističnu podrobnost. Motivacija opiraje se na kontinuitet vnutri raboty; ne dava samostojnu novu ocěnku vsih slovjanskyh alternativ.\par
\noindent\emph{Reason:} Integralna oblast iz polinomov carries integralna oblast from §§7 and 8 of Paper 6 into §9 without dropping the restriction to polynomials. A domain of rational functions and a domain consisting of polynomials are not simply two names for the same class of objects. Iz polinomov therefore conveys mathematical information, not merely stylistic detail. The rationale relies on continuity within the paper; it does not report a new independent comparison of all Slavic alternatives.\par
\noindent{\footnotesize Izvorne zapisy / Original records: HIST-g049-002.}\par
\smallskip\noindent\emph{Sravnjenje s tekstom:} Rabota 6 stroji integralnu oblast iz sistemy racionalnyh funkcij posrědstvom zatvorenosti pod skalarnym množenjem, dodavanjem i množenjem; dopuščene kvocienty dajut polje. Rabota 17 vvodži oblast cělosti $J$ bez prědpoloženja komutativnogo množenja. Rabota 19 razlikuje obči prsten od vlastnoj oblasti cělosti s jediniceju i bez nulovyh děliteljev. Te sravnjenja ne razrěšajut ustaliti vse pojavy jedinoju slovarnoju definicijeju.\par
\noindent\emph{Textual comparison:} Paper 6 constructs an integralna oblast from a system of rational functions by closure under scalar multiplication, addition and multiplication; admitting quotients gives a field. Paper 17 introduces oblast cělosti $J$ without assuming commutative multiplication. Paper 19 distinguishes a general ring from a proper domain of integrity with identity and without zero divisors. These comparisons do not license imposing one dictionary definition on every occurrence.\par
\noindent{\footnotesize TEXT-COMP-120; pasaže / passages: 5740, 5744, 5746, 5748, 5750, 5752, 5754, 11044, 12485; \texttt{CORPUS\_FORM\_COMPARISONS.json}.}\par
\end{minipage}\par
\par\medskip\Needspace{7\baselineskip}
\noindent\begin{minipage}{\linewidth}
\hypertarget{TERM-0913}{}
\noindent\textbf{integralne oblasti}\quad{\footnotesize\texttt{TERM-0913}}\par
\noindent\emph{Iztočnik / Source:} Integritätsbereiche\par
\noindent\emph{Rusky / Russian:} области целостности\par
\noindent\emph{Ukrajinsky / Ukrainian:} інтегральні області\par
\smallskip\noindent\emph{Razlog:} Integralne oblasti jest množstveny vid uvodnogo termina v rabotě 6. Ruski sravnjajuči prevod imaje genitivnu konstrukciju s celostnostju, dokolě ukrajinski i medžuslovjanski upotrěbjajut pridevnik integralny. Vybor mnogočislennogo vida ne zavodi drugy matematičny pojam. Uvod sohranja vez s pozdnějšimi definicijami i ne trěbuje doslovno kopirovati rusku sintaksu samo zato, že ruski jest jeden iz sravnjajučih jezykov.\par
\noindent\emph{Reason:} Integralne oblasti is the plural of the introductory term in Paper 6. The Russian comparator has a genitive construction referring to integrity, whereas Ukrainian and Interslavic use the adjective integralny. The plural introduces no different mathematical concept. The introduction retains its connection to the later definitions; Russian syntax need not be copied merely because Russian is one of the comparison languages.\par
\noindent{\footnotesize Izvorne zapisy / Original records: HIST-g040-009.}\par
\smallskip\noindent\emph{Sravnjenje s tekstom:} Rabota 6 stroji integralnu oblast iz sistemy racionalnyh funkcij posrědstvom zatvorenosti pod skalarnym množenjem, dodavanjem i množenjem; dopuščene kvocienty dajut polje. Rabota 17 vvodži oblast cělosti $J$ bez prědpoloženja komutativnogo množenja. Rabota 19 razlikuje obči prsten od vlastnoj oblasti cělosti s jediniceju i bez nulovyh děliteljev. Te sravnjenja ne razrěšajut ustaliti vse pojavy jedinoju slovarnoju definicijeju.\par
\noindent\emph{Textual comparison:} Paper 6 constructs an integralna oblast from a system of rational functions by closure under scalar multiplication, addition and multiplication; admitting quotients gives a field. Paper 17 introduces oblast cělosti $J$ without assuming commutative multiplication. Paper 19 distinguishes a general ring from a proper domain of integrity with identity and without zero divisors. These comparisons do not license imposing one dictionary definition on every occurrence.\par
\noindent{\footnotesize TEXT-COMP-120; pasaže / passages: 5740, 5744, 5746, 5748, 5750, 5752, 5754, 11044, 12485; \texttt{CORPUS\_FORM\_COMPARISONS.json}.}\par
\end{minipage}\par
\par\medskip\Needspace{7\baselineskip}
\noindent\begin{minipage}{\linewidth}
\hypertarget{TERM-0914}{}
\noindent\textbf{integralne oblasti iz polinomov}\quad{\footnotesize\texttt{TERM-0914}}\par
\noindent\emph{Iztočnik / Source:} Integritätsbereiche aus Polynomen\par
\noindent\emph{Rusky / Russian:} области целостности из полиномов\par
\noindent\emph{Ukrajinsky / Ukrainian:} інтегральні області з поліномів\par
\smallskip\noindent\emph{Razlog:} Integralne oblasti iz polinomov jest nazvanje klasa objektov raboty 6, §8. Motivacija sohranja i ranějši termin za oblast, i izrično ograničenje na polinomy. Poslědnje stava osoblivo važnym pri prěnosu involucijskoj bazy iz polja na polinomovu oblast: koeficienty morajut pripadati toj oblasti. Skračenje do integralne oblasti izgubilo by izričnu kvalifikaciju iztočnogo izraza. Množstveny vid jest gramatična prilagodba, ne poseben neologizm.\par
\noindent\emph{Reason:} Integralne oblasti iz polinomov names the class of objects in Paper 6, §8. The rationale retains both the established domain term and the explicit polynomial restriction. The latter becomes especially important when transferring an involution basis from a field to a polynomial domain: the coefficients must belong to that domain. Shortening the phrase to integralne oblasti would lose an express qualification of the source expression. The plural is a grammatical adjustment, not a separate neologism.\par
\noindent{\footnotesize Izvorne zapisy / Original records: HIST-g048-001.}\par
\smallskip\noindent\emph{Sravnjenje s tekstom:} Rabota 6 stroji integralnu oblast iz sistemy racionalnyh funkcij posrědstvom zatvorenosti pod skalarnym množenjem, dodavanjem i množenjem; dopuščene kvocienty dajut polje. Rabota 17 vvodži oblast cělosti $J$ bez prědpoloženja komutativnogo množenja. Rabota 19 razlikuje obči prsten od vlastnoj oblasti cělosti s jediniceju i bez nulovyh děliteljev. Te sravnjenja ne razrěšajut ustaliti vse pojavy jedinoju slovarnoju definicijeju.\par
\noindent\emph{Textual comparison:} Paper 6 constructs an integralna oblast from a system of rational functions by closure under scalar multiplication, addition and multiplication; admitting quotients gives a field. Paper 17 introduces oblast cělosti $J$ without assuming commutative multiplication. Paper 19 distinguishes a general ring from a proper domain of integrity with identity and without zero divisors. These comparisons do not license imposing one dictionary definition on every occurrence.\par
\noindent{\footnotesize TEXT-COMP-120; pasaže / passages: 5740, 5744, 5746, 5748, 5750, 5752, 5754, 11044, 12485; \texttt{CORPUS\_FORM\_COMPARISONS.json}.}\par
\end{minipage}\par
\par\medskip\Needspace{7\baselineskip}
\noindent\begin{minipage}{\linewidth}
\hypertarget{TERM-0915}{}
\noindent\textbf{integralne oblasti iz relativno cělyh funkcij}\quad{\footnotesize\texttt{TERM-0915}}\par
\noindent\emph{Iztočnik / Source:} Integritätsbereiche aus relativ-ganzen Funktionen\par
\noindent\emph{Rusky / Russian:} области целостности из относительно целых функций\par
\noindent\emph{Ukrajinsky / Ukrainian:} інтегральні області з відносно цілих функцій\par
\smallskip\noindent\emph{Razlog:} Integralne oblasti iz relativno cělyh funkcij sohranja Hilbertovo nazvanje v prose raboty 6, §9; bibliografičny naslov ostaje německy. To razlikovanje dopušča čitatelju razuměti argument i istovrěmenno najti citovano dělo. Relativno cěly tu ne znači prosto analitično cěly: lokalno uslovje potrěbuje, že kvocient dvoh funkcij oblasti pripada jej togda i samo togda, kogda jest polinom. Slědujuči argument svazuje to oprěděljenje s Hilbertovym nazvanjem. Proto glosu trěba čitati s toj definicijeju, ne kak sovrěmenny neograničeny termin za cele funkcije kompleksnoj proměnnoj.\par
\noindent\emph{Reason:} Integralne oblasti iz relativno cělyh funkcij preserves Hilbert's designation in the prose of Paper 6, §9; the bibliographic title remains German. This distinction lets readers understand the argument while still locating the cited work. Relativno cěly here does not simply mean analytically entire: the local condition requires that a quotient of two functions of the domain belong to it if and only if it is a polynomial. The ensuing argument relates that definition to Hilbert's designation. The gloss must therefore be read with this definition, not as an unrestricted modern term for entire functions of a complex variable.\par
\noindent{\footnotesize Izvorne zapisy / Original records: HIST-g049-006.}\par
\smallskip\noindent\emph{Sravnjenje s tekstom:} Lokalno uslovje 4a v rabotě 6 potrěbuje prinaležnost kvocienta oblasti togda i samo togda, kogda kvocient jest polinom. Odlomok potom svazuje takovu oblast s Hilbertovym nazvanjem. To jest konkretna opora interpretacije relativno cěly; ne trěba iz samogo slova izvoditi ni analitičnu cělost, ni vsako pozdnějše svojstvo integralnoj zamknjenosti.\par
\noindent\emph{Textual comparison:} Local condition 4a in Paper 6 requires that a quotient belong to the domain if and only if it is a polynomial. The subsequent passage relates such a domain to Hilbert's designation. This is concrete support for interpreting relativno cěly; the word alone does not establish analytic entirety or every later property of integral closure.\par
\noindent{\footnotesize TEXT-COMP-124; pasaže / passages: 5975, 5977, 5989; \texttt{CORPUS\_FORM\_COMPARISONS.json}.}\par
\end{minipage}\par
\par\medskip\Needspace{7\baselineskip}
\noindent\begin{minipage}{\linewidth}
\hypertarget{TERM-0916}{}
\noindent\textbf{invariantne tvorjenja}\quad{\footnotesize\texttt{TERM-0916}}\par
\noindent\emph{Iztočnik / Source:} invariante Bildungen\par
\noindent\emph{Rusky / Russian:} инвариантные образования\par
\noindent\emph{Ukrajinsky / Ukrainian:} інваріантні утворення\par
\smallskip\noindent\emph{Razlog:} Invariantne tvorjenja vybrano jest za široko německo Bildungen. V rabotě 3 nazivaje formovane invariantne objekty, ne samo formy ili kovarianty; v rabotě 12 naglas jest na konstrukcijah, prěz ktore dobyvajut se vse invarianty. Oboje podpiraje široky obhvat imenika. Tvorjenja ostaje izrično označeno kak forma, ktorej prirodnost potrěbuje jezykovu ocěnku. Zato motivaciju ne trěba prědstaviti kak potvrždženu idiomatičnost; i ni užše formy, ni samo kovarianty ne možut bez konteksta zaměniti izbrany izraz.\par
\noindent\emph{Reason:} Invariantne tvorjenja was chosen for the broad German Bildungen. In Paper 3 it names formed invariant objects, not just forms or covariants; in Paper 12 the emphasis is on constructions through which all invariants are obtained. Both favour a broadly inclusive noun. Tvorjenja is expressly marked as requiring linguistic appraisal for naturalness. The rationale must not be presented as confirmed idiomatic usage; neither the narrower forms nor covariants alone can replace it without reference to context.\par
\noindent{\footnotesize Izvorne zapisy / Original records: HIST-g028-005, HIST-g074-007.}\par
\end{minipage}\par
\par\medskip\Needspace{7\baselineskip}
\noindent\begin{minipage}{\linewidth}
\hypertarget{TERM-0917}{}
\noindent\textbf{invarianta grupy / invariantny integral}\quad{\footnotesize\texttt{TERM-0917}}\par
\noindent\emph{Iztočnik / Source:} Invariante der Gruppe / invariantes Integral\par
\noindent\emph{Rusky / Russian:} инвариант группы / инвариантный интеграл\par
\noindent\emph{Ukrajinsky / Ukrainian:} інваріант групи / інваріантний інтеграл\par
\smallskip\noindent\emph{Razlog:} Invarianta grupy i invariantny integral v rabotě 13 razlikujut imenik ot pridevnika. Žensky imenik invarianta naziva objekt invariantny po odnošenju k grupě; invariantny pri integral govori o svojstvu integrala. Ta lokalna posledovatelnost jest smysl motivacije. Ona ne znači, že slovnik už izključil množstveno invarianty ili druge gramatične varianty vo vseh rabotah. Nazvanje standardnym v ranějšem obosnovanju ne zaměnjaje nezavisnu attestaciju medžuslovjanskoj matematičnoj upotrěby.\par
\noindent\emph{Reason:} Invarianta grupy and invariantny integral in Paper 13 distinguish a noun from an adjective. The feminine noun invarianta names an object invariant under the group; invariantny modifying integral states a property of the integral. This local consistency is the point of the rationale. It does not mean the dictionary has excluded plural invarianty or other grammatical variants throughout the corpus. Calling the wording standard in the earlier rationale does not replace independent attestation of Interslavic mathematical usage.\par
\noindent{\footnotesize Izvorne zapisy / Original records: HIST-g075-007.}\par
\end{minipage}\par
\par\medskip\Needspace{7\baselineskip}
\noindent\begin{minipage}{\linewidth}
\hypertarget{TERM-0918}{}
\noindent\textbf{invariantne diferencialne operacije}\quad{\footnotesize\texttt{TERM-0918}}\par
\noindent\emph{Iztočnik / Source:} invariante Differentialoperationen\par
\noindent\emph{Rusky / Russian:} инвариантные дифференциальные операции\par
\noindent\emph{Ukrajinsky / Ukrainian:} інваріантні диференціальні операції\par
\smallskip\noindent\emph{Razlog:} Invariantne diferencialne operacije v rabotě 3 sohranja medžunarodnu slovotvornu semju diferencialnogo računa. Motivacija preferuje terminologično sočetanje vměsto opisa každoj operacije inymi slovami. Operacije ostavajut operacijami; pridevnik invariantne ne prěobrazuje jih v invariantne veličiny. Tak izbira se stabilno nazvanje dlja opakovanej konstrukcije, ale zapis ne dava čislene ocěnky kandidata ni sravnjene čitateljske rezultaty.\par
\noindent\emph{Reason:} Invariantne diferencialne operacije in Paper 3 retains the international word family of differential calculus. The rationale favours a terminological combination over paraphrasing each operation. The operations remain operations; the modifier invariantne does not turn them into invariant quantities. The choice supplies a stable name for a recurring construction, but the record gives neither numerical candidate scores nor comparative reader results.\par
\noindent{\footnotesize Izvorne zapisy / Original records: HIST-g028-017.}\par
\end{minipage}\par
\par\medskip\Needspace{7\baselineskip}
\noindent\begin{minipage}{\linewidth}
\hypertarget{TERM-0919}{}
\noindent\textbf{invariantne procesy diferenciranja}\quad{\footnotesize\texttt{TERM-0919}}\par
\noindent\emph{Iztočnik / Source:} invariante Differentiationsprozesse\par
\noindent\emph{Rusky / Russian:} инвариантные процессы дифференцирования\par
\noindent\emph{Ukrajinsky / Ukrainian:} інваріантні процеси диференціювання\par
\smallskip\noindent\emph{Razlog:} Invariantne procesy diferenciranja prědava Differentiationsprozesse v rabotě 16. Motivacija svazuje izraz s Fischerovym operatornym postupkom i razvitjem v redy; imenik procesy sohranja to, že rěč jest o provedenju diferenciranja, ne samo o gotovoj derivatě. Internacionalny koren diferenc- svazuje termin s ostaloj terminologijeju računa. Taka smyslova vez jest zachovana motivacija, ne dokaz, že procesy jest jedina vozmožna ili najlěpša medžuslovjanska kolokacija.\par
\noindent\emph{Reason:} Invariantne procesy diferenciranja renders Differentiationsprozesse in Paper 16. The rationale connects it to Fischer's operator-based procedure and series expansions; procesy retains the sense of carrying out differentiation rather than naming only the resulting derivative. The international diferenc- stem connects the term with the surrounding calculus vocabulary. This semantic connection is the preserved reason, not proof that procesy is the sole possible or best Interslavic collocation.\par
\noindent{\footnotesize Izvorne zapisy / Original records: HIST-g078-010.}\par
\end{minipage}\par
\par\medskip\Needspace{7\baselineskip}
\noindent\begin{minipage}{\linewidth}
\hypertarget{TERM-0920}{}
\noindent\textbf{invariantne variacijske problemy}\quad{\footnotesize\texttt{TERM-0920}}\par
\noindent\emph{Iztočnik / Source:} Invariante Variationsprobleme\par
\noindent\emph{Rusky / Russian:} инвариантные вариационные задачи\par
\noindent\emph{Ukrajinsky / Ukrainian:} інваріантні варіаційні задачі\par
\smallskip\noindent\emph{Razlog:} Invariantne variacijske problemy jest naslov raboty 13. Vybor sohranja medžunarodno razpoznavajemy sočetanje teorije variacij, vměsto doslovnogo periodnogo oborota problema variacije. To ne znamenaje proměnu obhvata raboty: invariantnost otnosi se k variacijnomu problemu. Motivacija predpolaga odbornu razpoznavajemost korenov, ale ne dokumentuje test razumlivosti. V rabotě 21 jest zachovan takože vybor variacijska zadača; naslovny problem i to nazvanje trěba pokazati kak dokumentovane varianty, ne vydavati jih za dokaz davno ustaljenoj jedinoj normy.\par
\noindent\emph{Reason:} Invariantne variacijske problemy is the title of Paper 13. The choice retains an internationally recognizable combination from the calculus of variations rather than a literal period-style problem of variation. This does not change the paper's scope: invariance concerns the variational problem. The rationale assumes disciplinary recognition of the stems, but documents no intelligibility test. Paper 21 also records variacijska zadača; the title's problem and that designation should be shown as documented variants, not evidence of a long-established single norm.\par
\noindent{\footnotesize Izvorne zapisy / Original records: HIST-g075-001.}\par
\end{minipage}\par
\par\medskip\Needspace{7\baselineskip}
\noindent\begin{minipage}{\linewidth}
\hypertarget{TERM-0921}{}
\noindent\textbf{invarianty libovoljnyh diferencialnyh izrazov}\quad{\footnotesize\texttt{TERM-0921}}\par
\noindent\emph{Iztočnik / Source:} Invarianten beliebiger Differentialausdrücke\par
\noindent\emph{Rusky / Russian:} инварианты произвольных дифференциальных выражений\par
\noindent\emph{Ukrajinsky / Ukrainian:} інваріанти довільних диференціальних виразів\par
\smallskip\noindent\emph{Razlog:} Invarianty libovoljnyh diferencialnyh izrazov jest naslovny izraz raboty 12. Libovoljny prědava beliebig kak proizvoljno izbrany v razgledajomom klasu, ne kak slučajny v probabilističnom smyslu. Diferencialny izraz ostaje izrazom; naslov ne oběščaje teoriju samo uravnenij ili samo operatorov v pozdnějšem smyslu. Motivacija podpiraje medžunarodny invariant- i točno sohranjenje širiny iztočnogo naslovu. Ne dava porědok kandidatov s čislennymi vagami.\par
\noindent\emph{Reason:} Invarianty libovoljnyh diferencialnyh izrazov is the title phrase in Paper 12. Libovoljny renders beliebig as arbitrarily chosen within the class considered, not random in a probabilistic sense. A differential expression remains an expression; the title does not restrict the theory to equations or to operators in a later sense. The rationale supports international invariant- and preservation of the source title's breadth. It provides no numerically weighted candidate ranking.\par
\noindent{\footnotesize Izvorne zapisy / Original records: HIST-g074-001.}\par
\end{minipage}\par
\par\medskip\Needspace{7\baselineskip}
\noindent\begin{minipage}{\linewidth}
\hypertarget{TERM-0922}{}
\noindent\textbf{invariantno polje}\quad{\footnotesize\texttt{TERM-0922}}\par
\noindent\emph{Iztočnik / Source:} Invariantenkörper\par
\noindent\emph{Rusky / Russian:} инвариантное поле; поле инвариантов\par
\noindent\emph{Ukrajinsky / Ukrainian:} поле інваріантів; інваріантне поле\par
\smallskip\noindent\emph{Razlog:} Invariantno polje sohranja matematičny tip objekta, zato že Körper v rabotah 5 i 11 znači polje, ne tělo v obyčnom smyslu. Motivacija raboty 5 svazuje izbor s invariantno-teoretičnoju terminologijeju raboty 4; motivacija raboty 11 sohranja sovrěmenno nazvanje polja pored istoričnogo Lagrangevogo nazvanja. Ruske i ukrajinske sravnjenja imajut i pridevnik, i konstrukciju polje invariantov. Ti vidy objasnjajut smyslovu blizkost, ale ne dokazujut, že pridevnična forma vsegda jest lěpša od genitivnoj. Polje trěba razlikovati od prstena invariantov.\par
\noindent\emph{Reason:} Invariantno polje preserves the mathematical type of the object, since Körper in Papers 5 and 11 means field, not body in the ordinary sense. Paper 5's rationale connects the choice to the invariant-theoretic terminology of Paper 4; Paper 11's retains the contemporary field name beside a historical Lagrangian designation. The Russian and Ukrainian comparisons include both adjectival forms and field-of-invariants constructions. These explain semantic proximity but do not establish that the adjective is always preferable to the genitive. The field must remain distinct from a ring of invariants.\par
\noindent{\footnotesize Izvorne zapisy / Original records: HIST-g039-006, HIST-g073-009.}\par
\end{minipage}\par
\par\medskip\Needspace{7\baselineskip}
\noindent\begin{minipage}{\linewidth}
\hypertarget{TERM-0923}{}
\noindent\textbf{invariantna variacijska zadača}\quad{\footnotesize\texttt{TERM-0923}}\par
\noindent\emph{Iztočnik / Source:} invariantes Variationsproblem\par
\noindent\emph{Rusky / Russian:} инвариантная вариационная задача\par
\noindent\emph{Ukrajinsky / Ukrainian:} інваріантна варіаційна задача\par
\smallskip\noindent\emph{Razlog:} Invariantna variacijska zadača v rabotě 21 imenuje konkretny variacijny problem, iz kotorogo proizhodet Lagrangeve izrazy. Izbor zadača sohranja smysl problemu ili zadanja; zaměna na obči princip byla by širšoju interpretacijeju, ne ekvivalentnym prevodom. Motivacija otnosi se k toj strukturnoj roli. Nazvanie raboty 13 upotrěbjaje problemy, takže zadača jest dokumentovana alternativa vnutri korpusa, ne povsudno provedeny izbor. Ne trěba iz te različnosti izvoditi različnost matematičnogo pojma bez dodatnogo konteksta.\par
\noindent\emph{Reason:} Invariantna variacijska zadača in Paper 21 names the specific variational problem from which the Lagrangian expressions arise. Zadača preserves the sense of a problem or task; replacing it with a general principle would broaden the interpretation rather than translate it equivalently. The rationale concerns that structural role. The title of Paper 13 uses problemy, so zadača is a documented alternative within the corpus, not a universally implemented choice. The lexical difference alone does not establish a difference in mathematical concept.\par
\noindent{\footnotesize Izvorne zapisy / Original records: HIST-g084-003.}\par
\end{minipage}\par
\par\medskip\Needspace{7\baselineskip}
\noindent\begin{minipage}{\linewidth}
\hypertarget{TERM-0924}{}
\noindent\textbf{involucija v linearnom prostoru}\quad{\footnotesize\texttt{TERM-0924}}\par
\noindent\emph{Iztočnik / Source:} Involution in einem linearen Raum\par
\noindent\emph{Rusky / Russian:} инволюция в линейном пространстве\par
\noindent\emph{Ukrajinsky / Ukrainian:} інволюція в лінійному просторі\par
\smallskip\noindent\emph{Razlog:} Involucija v linearnom prostoru sohranja geometrično nazvanje iz raboty 6, §5. V odlomku vsaky red $x$ razuměje se kak točka, a sistem rěšenij dava grupy po $i$ toček; vsaka točka oprěděljaje cělu grupu, ktorej pripada. To jest nazvana involucija. Čitatelj ne směje avtomatično podstaviti samo sovrěmenno odobraženje s kvadratom ravnym identičnosti. Medžunarodny imenik sohranja istoričnu vez s involucijskoju formoju i bazoju; smysl daje lokalno geometrično oprěděljenje, ne samo podobnost slov.\par
\noindent\emph{Reason:} Involucija v linearnom prostoru preserves the geometric designation in Paper 6, §5. In the passage, each row $x$ is understood as a point, and the solution system gives groups of $i$ points; each point determines the whole group to which it belongs. This is what the text calls an involution. Readers must not automatically substitute only the modern notion of a map whose square is the identity. The international noun preserves the historical connection with the involution form and basis; the local geometric definition, not word resemblance alone, supplies the meaning.\par
\noindent{\footnotesize Izvorne zapisy / Original records: HIST-g045-011.}\par
\smallskip\noindent\emph{Sravnjenje s tekstom:} Geometrično oprěděljenje govori o grupah po $i$ toček. Koeficienty involucijskoj formy daju racionalnu bazu s izrično ograničenoju jednoznačnostju; v dokazu Lurotovoj teoremy izbrany koeficient stvarno zavisi od neoprěděljenoj. Formulaciju o libovoljnoj funkciji bazy zato trěba čitati s tym kontekstom. Pri prěnosu na polinomovu oblast uslovja na koeficienty i njih vzajemnu prostotu sut čest definicije.\par
\noindent\emph{Textual comparison:} The geometric definition concerns groups of $i$ points. Coefficients of the involution form give a rational basis with explicitly qualified uniqueness; in the proof of Lüroth's theorem the chosen coefficient genuinely depends on the indeterminate. The formulation about any function of the basis must be read in that context. In the transfer to a polynomial domain, conditions on the coefficients and their relative primality form part of the definition.\par
\noindent{\footnotesize TEXT-COMP-125; pasaže / passages: 5658, 5576, 5704, 5708, 5901; \texttt{CORPUS\_FORM\_COMPARISONS.json}.}\par
\end{minipage}\par
\par\medskip\Needspace{7\baselineskip}
\noindent\begin{minipage}{\linewidth}
\hypertarget{TERM-0925}{}
\noindent\textbf{involucijska baza}\quad{\footnotesize\texttt{TERM-0925}}\par
\noindent\emph{Iztočnik / Source:} Involutionsbasis\par
\noindent\emph{Rusky / Russian:} инволюционная база; инволюционный базис\par
\noindent\emph{Ukrajinsky / Ukrainian:} інволюційна база\par
\smallskip\noindent\emph{Razlog:} Involucijska baza jest čest sočetanoj terminologije raboty 6: uvod vvodži nazvanje, §5 dava koeficientnu konstrukciju, §6 svazuje ju s minimalnoju bazoju, §8 prěnosi ju na polinomove oblasti, a §10 v Teoremě VII svazuje ju s integralnoju bazoju. Ruske zapisy kolebajut se medžu baza i bazis; medžuslovjanski sohranja baza. V §6 formulacija o libovoljnoj funkciji bazy ne směje byti odtrgnuta od predhodnogo dokaza, ktory izrično izbira koeficient stvarno zaviseči od neoprěděljenoj. Zato ne trěba utvrdžati, že libovoljna konstanta generuje isto polje. Kvalifikacije o jednoznačnosti v §5 i uslovje prvogo roda v §10 sut matematičny obhvat rezultata, ne izlišne podrobnosti slovnika.\par
\noindent\emph{Reason:} Involucijska baza belongs to a connected terminology in Paper 6: the introduction names it, §5 gives the coefficient construction, §6 relates it to a minimal basis, §8 transfers it to polynomial domains, and §10 relates it to an integral basis in Theorem VII. Russian records alternate between baza and bazis; Interslavic retains baza. In §6 the statement about any function of the basis must not be detached from the preceding proof, which explicitly chooses a coefficient genuinely depending on the indeterminate. It must not be expanded into a claim that an arbitrary constant generates the same field. The qualifications on uniqueness in §5 and the first-kind hypothesis in §10 delimit the mathematical result; they are not dispensable dictionary detail.\par
\noindent{\footnotesize Izvorne zapisy / Original records: HIST-g040-008, HIST-g045-002, HIST-g046-011, HIST-g048-010, HIST-g050-007.}\par
\smallskip\noindent\emph{Sravnjenje s tekstom:} Geometrično oprěděljenje govori o grupah po $i$ toček. Koeficienty involucijskoj formy daju racionalnu bazu s izrično ograničenoju jednoznačnostju; v dokazu Lurotovoj teoremy izbrany koeficient stvarno zavisi od neoprěděljenoj. Formulaciju o libovoljnoj funkciji bazy zato trěba čitati s tym kontekstom. Pri prěnosu na polinomovu oblast uslovja na koeficienty i njih vzajemnu prostotu sut čest definicije.\par
\noindent\emph{Textual comparison:} The geometric definition concerns groups of $i$ points. Coefficients of the involution form give a rational basis with explicitly qualified uniqueness; in the proof of Lüroth's theorem the chosen coefficient genuinely depends on the indeterminate. The formulation about any function of the basis must be read in that context. In the transfer to a polynomial domain, conditions on the coefficients and their relative primality form part of the definition.\par
\noindent{\footnotesize TEXT-COMP-125; pasaže / passages: 5658, 5576, 5704, 5708, 5901; \texttt{CORPUS\_FORM\_COMPARISONS.json}.}\par
\end{minipage}\par
\par\medskip\Needspace{7\baselineskip}
\noindent\begin{minipage}{\linewidth}
\hypertarget{TERM-0926}{}
\noindent\textbf{involucijska forma}\quad{\footnotesize\texttt{TERM-0926}}\par
\noindent\emph{Iztočnik / Source:} Involutionsform\par
\noindent\emph{Rusky / Russian:} инволюционная форма\par
\noindent\emph{Ukrajinsky / Ukrainian:} інволюційна форма\par
\smallskip\noindent\emph{Razlog:} Involucijska forma naziva izbranu homogennu formu v rabotě 6, §5, i sohranja geometričnu vez s involucijeju. V §6 jej koeficient služi pri dokazu Lurotovoj teoremy. V §8 forma polja množi se polinomom tako, že jej koeficienty naležet polinomovoj oblasti i sut v nej vzajemno proste; ti koeficienty potom tvoret odgovarjajuču bazu. Nazvanje forma ne dava osnovu zaměniti tu konstrukciju definicijeju sovrěmennogo operatora. Tri motivacije podpirajut kontinuitet nazvanja pri proměně uslovij na koeficienty, ne tvrdženje, že vse te objekty sut bez uslovij identične.\par
\noindent\emph{Reason:} Involucijska forma names the distinguished homogeneous form in Paper 6, §5, preserving its geometric connection with involution. In §6 a coefficient of the form is used in the proof of Lüroth's theorem. In §8 the field's form is multiplied by a polynomial so that its coefficients belong to the polynomial domain and are relatively prime there; those coefficients then form the associated basis. Choosing forma must not replace this construction with a modern operator definition. The three rationales support continuity of the name as coefficient conditions change, not unconditional identity of all the objects involved.\par
\noindent{\footnotesize Izvorne zapisy / Original records: HIST-g045-001, HIST-g046-010, HIST-g048-009.}\par
\smallskip\noindent\emph{Sravnjenje s tekstom:} Geometrično oprěděljenje govori o grupah po $i$ toček. Koeficienty involucijskoj formy daju racionalnu bazu s izrično ograničenoju jednoznačnostju; v dokazu Lurotovoj teoremy izbrany koeficient stvarno zavisi od neoprěděljenoj. Formulaciju o libovoljnoj funkciji bazy zato trěba čitati s tym kontekstom. Pri prěnosu na polinomovu oblast uslovja na koeficienty i njih vzajemnu prostotu sut čest definicije.\par
\noindent\emph{Textual comparison:} The geometric definition concerns groups of $i$ points. Coefficients of the involution form give a rational basis with explicitly qualified uniqueness; in the proof of Lüroth's theorem the chosen coefficient genuinely depends on the indeterminate. The formulation about any function of the basis must be read in that context. In the transfer to a polynomial domain, conditions on the coefficients and their relative primality form part of the definition.\par
\noindent{\footnotesize TEXT-COMP-125; pasaže / passages: 5658, 5576, 5704, 5708, 5901; \texttt{CORPUS\_FORM\_COMPARISONS.json}.}\par
\end{minipage}\par
\par\medskip\Needspace{7\baselineskip}
\noindent\begin{minipage}{\linewidth}
\hypertarget{TERM-0927}{}
\noindent\textbf{«iracionalny» / «racionalny» směr}\quad{\footnotesize\texttt{TERM-0927}}\par
\noindent\emph{Iztočnik / Source:} irrationale / rationale Richtung\par
\noindent\emph{Rusky / Russian:} «иррациональное» / «рациональное» направление\par
\noindent\emph{Ukrajinsky / Ukrainian:} «ірраціональний» / «раціональний» напрям\par
\smallskip\noindent\emph{Razlog:} «Iracionalny» i «racionalny» směr v rabotě 11 sut nazvanja dvaju sposobov opisa: prvy harakterizuje korene, drugy koeficienty. Kavyčky sohranjajut status istoričnyh nazvanj. Iracionalny tu ne znači nelogičny ili nerazumny; racionalny ne jest pohvala lěpšego razmišljanja. Vybor internacionalnyh pridevnikov sohranja vez s algebraičnym razlikovanjem, a směr naziva sposob izslědovanja. Ta natuknica dokumentuje takože izložitelnu konvenciju: bez kavyček i objasnjenja čitatelj mog by čitati ocěnku razuma vměsto klasa metoda.\par
\noindent\emph{Reason:} “Iracionalny” and “racionalny” směr in Paper 11 name two modes of description: the first characterizes roots, the second coefficients. Quotation marks retain their status as historical designations. Iracionalny does not mean illogical or unreasonable, and racionalny is not praise for better reasoning. The international adjectives preserve the algebraic connection, while směr names an approach to investigation. This entry also documents an expository convention: without the quotation marks and explanation, readers could mistake a classification of methods for an assessment of rational thought.\par
\noindent{\footnotesize Izvorne zapisy / Original records: HIST-g073-002.}\par
\end{minipage}\par
\par\medskip\Needspace{7\baselineskip}
\noindent\begin{minipage}{\linewidth}
\hypertarget{TERM-0928}{}
\noindent\textbf{nerazložimy polinomy}\quad{\footnotesize\texttt{TERM-0928}}\par
\noindent\emph{Iztočnik / Source:} Irreducible polynomials\par
\noindent\emph{Rusky / Russian:} неприводимые полиномы\par
\noindent\emph{Ukrajinsky / Ukrainian:} незвідні поліноми\par
\smallskip\noindent\emph{Razlog:} Nerazložimy polinomy odpovědaje anglijskomu Irreducible polynomials v rabotě 22, §7. Vybor transparentnogo slovjanskogo složenja govori o nemožnosti razloženja; rusko неприводимые i ukrajinsko незвідні imajut različne ustaljene korene. Zapis ne tvrdi, že medžuslovjanski doslovno prěuzima jedan iz njih. Za polinom jest važno koeficientno polje ili prsten, v ktorom razloženje razgledaje se; slovo samo ne dava bezuslovnu nerazloživost po vsakom razširjenju koeficientov. Iztočna natuknica jest anglijsky citat, ne poškodžena německa forma.\par
\noindent\emph{Reason:} Nerazložimy polinomy corresponds to English Irreducible polynomials in Paper 22, §7. The transparent Slavic compound refers to inability to decompose; the Russian and Ukrainian comparators use different established stems. The record does not claim that Interslavic directly adopts either one. For a polynomial, the coefficient field or ring in which factorization is considered matters; the word alone does not assert irreducibility under every coefficient extension. The source headword is an English quotation, not a damaged German form.\par
\noindent{\footnotesize Izvorne zapisy / Original records: HIST-g091-010, HIST-g092-079.}\par
\end{minipage}\par
\par\medskip\Needspace{7\baselineskip}
\noindent\begin{minipage}{\linewidth}
\hypertarget{TERM-0929}{}
\noindent\textbf{nerazložimy / nerazložimo}\quad{\footnotesize\texttt{TERM-0929}}\par
\noindent\emph{Iztočnik / Source:} irreduzibel\par
\noindent\emph{Rusky / Russian:} неразложимый\par
\noindent\emph{Ukrajinsky / Ukrainian:} нерозкладний\par
\smallskip\noindent\emph{Razlog:} Nerazložimy i nerazložimo v rabotě 19 sohranjajut svojstvo idealov i drugih objektov po odnošenju k oprěděljenomu razloženju. V uvodnom priměru $q$ ne može byti najmenšim občim mnogokratnym dvoh vlastnyh děliteljev. Zato ne trěba prěvesti tu upotrěbu kak obyčnu nemožnost produkta dvaju nejediničnyh polinomov. Ruski i ukrajinski sravnjajuči vidy ukazujut na nerazloživost, a medžuslovjanski sohranja tu čitljivu vez. Alternacija konca -y i -o služi gramatičnomu kontekstu; ne označaje dve različne teorije.\par
\noindent\emph{Reason:} Nerazložimy and nerazložimo in Paper 19 preserve a property of ideals and other objects relative to a specified decomposition. In the introductory example, $q$ cannot be the least common multiple of two proper divisors. This usage must not be translated as ordinary inability to factor into two nonunit polynomials. The Russian and Ukrainian comparators point toward indecomposability, and Interslavic retains that transparent connection. The -y/-o alternation serves the grammatical context; it does not denote two different theories.\par
\noindent{\footnotesize Izvorne zapisy / Original records: HIST-g081-008, HIST-g082-008.}\par
\smallskip\noindent\emph{Sravnjenje s tekstom:} Raboty 19 i 30 objasnjajut nerazloživost pomoču najmenšego občego mnogokratnogo dvoh vlastnyh děliteljev. Pri uslovju cěpov nerazložimy ideal jest primarny, ale kontekst s nulovymi děliteljami razgledaje takože razložime primarne idealy. V rabotě 17 ireducibilne sostavne česti govoret o aditivnom razloženju grupy ostatkov. Te upotrěby ne trěba spojiti s polinomovym razloženjem na faktory bez ukazanija operacije.\par
\noindent\emph{Textual comparison:} Papers 19 and 30 explain irreducibility by the least common multiple of two proper divisors. Under the chain condition an irreducible ideal is primary, but the zero-divisor context also considers reducible primary ideals. In Paper 17 irreducible constituents concern additive decomposition of a residue group. These usages must not be merged with polynomial factorization without specifying the operation.\par
\noindent{\footnotesize TEXT-COMP-126; pasaže / passages: 12424, 16680, 16686, 17147, 17149, 11446; \texttt{CORPUS\_FORM\_COMPARISONS.json}.}\par
\end{minipage}\par
\par\medskip\Needspace{7\baselineskip}
\noindent\begin{minipage}{\linewidth}
\hypertarget{TERM-0930}{}
\noindent\textbf{nerazložimy / razložimy}\quad{\footnotesize\texttt{TERM-0930}}\par
\noindent\emph{Iztočnik / Source:} irreduzibel / reduzibel\par
\noindent\emph{Rusky / Russian:} неразложимый / разложимый\par
\noindent\emph{Ukrajinsky / Ukrainian:} нерозкладний / розкладний\par
\smallskip\noindent\emph{Razlog:} Nerazložimy / razložimy prědava irreduzibel / reduzibel v rabotě 30, §6, kak protivopoložne svojstva idealov. Teorema I izrično definuje prvy vid pomoču nemožnosti prědstaviti ideal kak najmenše obče mnogokratno dvoh vlastnyh děliteljev. V Teoremě II pod uslovjem cěpov děliteljev iz nerazloživosti slěduje primarnost. To uslovje i směr implikacije morajut ostati vidime; termin ne čini vsaky primarny ideal avtomatično nerazložimym. Vybory sut svazany s operacijeju razloženja, ne samo s obštrim značenjem německogo pridevnika.\par
\noindent\emph{Reason:} Nerazložimy / razložimy renders irreduzibel / reduzibel in Paper 30, §6, as opposed properties of ideals. Theorem I explicitly defines the first through inability to express the ideal as the least common multiple of two proper divisors. Under the divisor-chain condition, Theorem II derives primaryness from irreducibility. Both the hypothesis and direction of implication must remain visible; the terminology does not automatically make every primary ideal irreducible. These choices depend on the decomposition operation, not simply on the general meaning of the German adjective.\par
\noindent{\footnotesize Izvorne zapisy / Original records: HIST-g113-002.}\par
\smallskip\noindent\emph{Sravnjenje s tekstom:} Raboty 19 i 30 objasnjajut nerazloživost pomoču najmenšego občego mnogokratnogo dvoh vlastnyh děliteljev. Pri uslovju cěpov nerazložimy ideal jest primarny, ale kontekst s nulovymi děliteljami razgledaje takože razložime primarne idealy. V rabotě 17 ireducibilne sostavne česti govoret o aditivnom razloženju grupy ostatkov. Te upotrěby ne trěba spojiti s polinomovym razloženjem na faktory bez ukazanija operacije.\par
\noindent\emph{Textual comparison:} Papers 19 and 30 explain irreducibility by the least common multiple of two proper divisors. Under the chain condition an irreducible ideal is primary, but the zero-divisor context also considers reducible primary ideals. In Paper 17 irreducible constituents concern additive decomposition of a residue group. These usages must not be merged with polynomial factorization without specifying the operation.\par
\noindent{\footnotesize TEXT-COMP-126; pasaže / passages: 12424, 16680, 16686, 17147, 17149, 11446; \texttt{CORPUS\_FORM\_COMPARISONS.json}.}\par
\end{minipage}\par
\par\medskip\Needspace{7\baselineskip}
\noindent\begin{minipage}{\linewidth}
\hypertarget{TERM-0931}{}
\noindent\textbf{nerazložimy / razložimy}\quad{\footnotesize\texttt{TERM-0931}}\par
\noindent\emph{Iztočnik / Source:} irreduzibel / reduzibel (Primärideal)\par
\noindent\emph{Rusky / Russian:} неразложимый / разложимый\par
\noindent\emph{Ukrajinsky / Ukrainian:} нерозкладний / розкладний\par
\smallskip\noindent\emph{Razlog:} Nerazložimy / razložimy pri primarnom idealu prodolžuje to samo oprěděljenje raboty 30 v §9. Dva zachovana zapisy ne vvodžut novy smysl slova: ide o razloženje kak najmenše obče mnogokratno. Kontekst s nulovymi děliteljami izrično razgledaje razložime primarne idealy. Zato fraza primarny ideal i fraza nerazložimy ideal ne sut samostojno zaměnljive. Sohranjenje toj različnosti važnějše jest od kosmetičnogo jedinstva s terminom za ireducibilny polinom.\par
\noindent\emph{Reason:} Nerazložimy / razložimy for a primary ideal continues the same definition in Paper 30, §9. The two preserved records introduce no new sense of the word: the decomposition is as a least common multiple. The context with zero divisors explicitly considers reducible primary ideals. Primary ideal and irreducible ideal are therefore not independently interchangeable expressions. Preserving this distinction matters more than cosmetic uniformity with the term for an irreducible polynomial.\par
\noindent{\footnotesize Izvorne zapisy / Original records: HIST-g116-005, HIST-g116-013.}\par
\smallskip\noindent\emph{Sravnjenje s tekstom:} Raboty 19 i 30 objasnjajut nerazloživost pomoču najmenšego občego mnogokratnogo dvoh vlastnyh děliteljev. Pri uslovju cěpov nerazložimy ideal jest primarny, ale kontekst s nulovymi děliteljami razgledaje takože razložime primarne idealy. V rabotě 17 ireducibilne sostavne česti govoret o aditivnom razloženju grupy ostatkov. Te upotrěby ne trěba spojiti s polinomovym razloženjem na faktory bez ukazanija operacije.\par
\noindent\emph{Textual comparison:} Papers 19 and 30 explain irreducibility by the least common multiple of two proper divisors. Under the chain condition an irreducible ideal is primary, but the zero-divisor context also considers reducible primary ideals. In Paper 17 irreducible constituents concern additive decomposition of a residue group. These usages must not be merged with polynomial factorization without specifying the operation.\par
\noindent{\footnotesize TEXT-COMP-126; pasaže / passages: 12424, 16680, 16686, 17147, 17149, 11446; \texttt{CORPUS\_FORM\_COMPARISONS.json}.}\par
\end{minipage}\par
\par\medskip\Needspace{7\baselineskip}
\noindent\begin{minipage}{\linewidth}
\hypertarget{TERM-0932}{}
\noindent\textbf{ireducibilne sostavne česti}\quad{\footnotesize\texttt{TERM-0932}}\par
\noindent\emph{Iztočnik / Source:} irreduzible Bestandteile\par
\noindent\emph{Rusky / Russian:} неприводимые составные части\par
\noindent\emph{Ukrajinsky / Ukrainian:} незвідні складові частини\par
\smallskip\noindent\emph{Razlog:} Ireducibilne sostavne česti v rabotě 17 sohranja medžunarodny pridevnik i prozračno slovjansko nazvanje komponentov. Motivacija ne zaměnjaje Bestandteile užšim podgrupy: v teoremě rěč jest o prědstavjenju grupy ostatkov kak sumy takyh čestij, povezanym s prědstavjenjem modula kak najmenšego občego mnogokratnogo. Lokalno značenje ireducibilnosti ustanavjaje argument o daljšem aditivnom razloženju. Iz te frazy samoj ne možno izvoditi, že avtor govori o ireducibilnych prědstavjenjah v točno tom samom smyslu kak v pozdnějših rabotah.\par
\noindent\emph{Reason:} Ireducibilne sostavne česti in Paper 17 retains an international adjective and a transparent Slavic name for constituents. The rationale does not replace Bestandteile with the narrower subgroups: the theorem concerns expressing a residue group as a sum of such constituents, linked to expressing a module as a least common multiple. The argument about further additive decomposition establishes the local meaning of irreducibility. The phrase alone cannot establish that the author means irreducible representations in exactly the same sense as in the later papers.\par
\noindent{\footnotesize Izvorne zapisy / Original records: HIST-g079-036.}\par
\smallskip\noindent\emph{Sravnjenje s tekstom:} Raboty 19 i 30 objasnjajut nerazloživost pomoču najmenšego občego mnogokratnogo dvoh vlastnyh děliteljev. Pri uslovju cěpov nerazložimy ideal jest primarny, ale kontekst s nulovymi děliteljami razgledaje takože razložime primarne idealy. V rabotě 17 ireducibilne sostavne česti govoret o aditivnom razloženju grupy ostatkov. Te upotrěby ne trěba spojiti s polinomovym razloženjem na faktory bez ukazanija operacije.\par
\noindent\emph{Textual comparison:} Papers 19 and 30 explain irreducibility by the least common multiple of two proper divisors. Under the chain condition an irreducible ideal is primary, but the zero-divisor context also considers reducible primary ideals. In Paper 17 irreducible constituents concern additive decomposition of a residue group. These usages must not be merged with polynomial factorization without specifying the operation.\par
\noindent{\footnotesize TEXT-COMP-126; pasaže / passages: 12424, 16680, 16686, 17147, 17149, 11446; \texttt{CORPUS\_FORM\_COMPARISONS.json}.}\par
\end{minipage}\par
\par\medskip\Needspace{7\baselineskip}
\noindent\begin{minipage}{\linewidth}
\hypertarget{TERM-0933}{}
\noindent\textbf{nerazložime predstavjenja}\quad{\footnotesize\texttt{TERM-0933}}\par
\noindent\emph{Iztočnik / Source:} irreduzible Darstellungen\par
\noindent\emph{Rusky / Russian:} неприводимые представления\par
\noindent\emph{Ukrajinsky / Ukrainian:} незвідні зображення\par
\smallskip\noindent\emph{Razlog:} Nerazložime predstavjenja jest zachovany izbor za irreduzible Darstellungen v uvodu raboty 32. Motivacija daje prědnost transparentnomu pridevniku i medžuslovjanskomu predstavjenje pored ruskogo predstavlenie i ukrajinskogo zobrazhennia. Ta lokalna odluka ne soglasuje se s pravilom raboty 28, ktore izrično rezervovalo ireducibilny za prědstavjenja, a nerazložimy za idealy. Sravnjeni tekst raboty 32 stvarno upotrěbjaje nerazložime prědstavjenja. Zato slovnik ne směje prikryti različnost tvrdženjem o posledovatelno provedenoj trojnoj shemě. Razlikovanje ireducibilnogo i samo direktno nerazložimogo prědstavjenja potrěbuje konkretnu definiciju i hipotezy, ne samo zaměnu pridevnika.\par
\noindent\emph{Reason:} Nerazložime predstavjenja is the recorded choice for irreduzible Darstellungen at the opening of Paper 32. The rationale favours a transparent adjective and Interslavic predstavjenje alongside Russian predstavlenie and Ukrainian zobrazhennia. This local decision does not match Paper 28's explicit rule reserving ireducibilny for representations and nerazložimy for ideals. The compared Paper 32 text does indeed use nerazložime prědstavjenja. The dictionary must not conceal this difference by claiming a consistently implemented three-way scheme. Distinguishing an irreducible representation from a merely directly indecomposable one requires the particular definition and hypotheses, not just an adjective replacement.\par
\noindent{\footnotesize Izvorne zapisy / Original records: HIST-g166-002.}\par
\smallskip\noindent\emph{Sravnjenje s tekstom:} Rabota 28 v jednom odlomku razlikuje nerazložime idealy i ireducibilne klasy prědstavjenij, soglasno izričnoj trojnoj konvenciji. Naslov i uvod raboty 32 upotrěbjajut nerazložime prědstavjenja i nerazložime sostavne česti. To jest konkretna nesoglasnost leksičnyh pravil, ktora ostaje predmetom ocěnky; čitateljsko izloženje ne směje izmysliti jedinu doslědno provedenu politiku.\par
\noindent\emph{Textual comparison:} Paper 28 distinguishes indecomposable ideals from irreducible representation classes in one passage, matching its explicit three-way convention. Paper 32's title and opening use nerazložime prědstavjenja and nerazložime sostavne česti. This is a concrete inconsistency between lexical rules, still requiring appraisal; a reader-facing account must not invent a single consistently implemented policy.\par
\noindent{\footnotesize TEXT-COMP-127; pasaže / passages: 15950, 19171, 19210; \texttt{CORPUS\_FORM\_COMPARISONS.json}.}\par
\end{minipage}\par
\par\medskip\Needspace{7\baselineskip}
\noindent\begin{minipage}{\linewidth}
\hypertarget{TERM-0934}{}
\noindent\textbf{ireducibilne klasy predstavjenj}\quad{\footnotesize\texttt{TERM-0934}}\par
\noindent\emph{Iztočnik / Source:} irreduzible Darstellungsklassen\par
\noindent\emph{Rusky / Russian:} неприводимые классы представлений\par
\noindent\emph{Ukrajinsky / Ukrainian:} незвідні класи представлень\par
\smallskip\noindent\emph{Razlog:} Ireducibilne klasy predstavjenj v rabotě 28 jest čest izričnoj trojnoj konvencije: popolno reducibilny za polno reducibilne prsteny, nerazložimy za nerazložime idealy, ireducibilny za prědstavjenja. V sravnjenom odlomku klasy nerazložimyh idealov odpovědajut klasam ireducibilnyh prědstavjenij pri specializaciji na polno reducibilne hiperkompleksne sistemy. Klas jest klasom prědstavjenij; ne trěba iz pridevničnogo poredka izvoditi novu operaciju razloženja na samyh klasah. Taka lokalna shema objasnjaje izbor, ale variant raboty 32 pokazyvaje, že ona ne byla jedinoju politikoj cělogo prevoda.\par
\noindent\emph{Reason:} Ireducibilne klasy predstavjenj in Paper 28 belongs to an explicit three-way convention: popolno reducibilny for completely reducible rings, nerazložimy for indecomposable ideals, and ireducibilny for representations. In the compared passage, classes of indecomposable ideals correspond to classes of irreducible representations after specialization to completely reducible hypercomplex systems. A class is a class of representations; adjective order must not be read as introducing a new decomposition operation on the classes themselves. The local scheme explains the choice, but Paper 32's alternative shows that it was not the translation's only policy.\par
\noindent{\footnotesize Izvorne zapisy / Original records: HIST-g105-010.}\par
\smallskip\noindent\emph{Sravnjenje s tekstom:} Rabota 28 v jednom odlomku razlikuje nerazložime idealy i ireducibilne klasy prědstavjenij, soglasno izričnoj trojnoj konvenciji. Naslov i uvod raboty 32 upotrěbjajut nerazložime prědstavjenja i nerazložime sostavne česti. To jest konkretna nesoglasnost leksičnyh pravil, ktora ostaje predmetom ocěnky; čitateljsko izloženje ne směje izmysliti jedinu doslědno provedenu politiku.\par
\noindent\emph{Textual comparison:} Paper 28 distinguishes indecomposable ideals from irreducible representation classes in one passage, matching its explicit three-way convention. Paper 32's title and opening use nerazložime prědstavjenja and nerazložime sostavne česti. This is a concrete inconsistency between lexical rules, still requiring appraisal; a reader-facing account must not invent a single consistently implemented policy.\par
\noindent{\footnotesize TEXT-COMP-127; pasaže / passages: 15950, 19171, 19210; \texttt{CORPUS\_FORM\_COMPARISONS.json}.}\par
\end{minipage}\par
\par\medskip\Needspace{7\baselineskip}
\noindent\begin{minipage}{\linewidth}
\hypertarget{TERM-0935}{}
\noindent\textbf{nereducibilne formy}\quad{\footnotesize\texttt{TERM-0935}}\par
\noindent\emph{Iztočnik / Source:} Irreduzible Formen\par
\noindent\emph{Rusky / Russian:} неприводимые формы; нередуцируемые формы\par
\noindent\emph{Ukrajinsky / Ukrainian:} незвідні форми; нередуковні форми\par
\smallskip\noindent\emph{Razlog:} Nereducibilne formy jest dokumentovano jedinnadsetju motivacijami raboty 2: §§8, 10, 11, 13, 14, 15, 18, 20, 23, 24 i 26; §16 ima posebnju natuknicu za jedinstveny i množstveny vid. Glavna motivacija jest sohranjenje semji redukcija, reducent i nereducibilny. V §§8 i 10 svaz jest izrično s ranějšimi §§3, 7 i 9; v §11 medžunarodny koren motivovan jest želanjem ne privilegovati jeden nacionalny jezyk. Pozdnějše zapisy podpirajut kontinuitet, osoblivo v gustoj redukcijnoj terminologiji §§23--26. Ruske i ukrajinske varianty koristajut različne korene; od §13 nadalje někotore motivacije ostavjajut nerazložime kak slovjansku alternativu za ocěnku. To jest realna alternativa, ne odbačeny kandidat s izmysljenoju nižšoju čislenoju ocěnkoju.\par
\noindent\emph{Reason:} Nereducibilne formy is supported by eleven rationales in Paper 2, covering §§8, 10, 11, 13, 14, 15, 18, 20, 23, 24 and 26; §16 has a separate singular/plural entry. The main reason is to retain the family redukcija, reducent and nereducibilny. In §§8 and 10 the connection is explicitly to earlier §§3, 7 and 9; in §11 the international stem is motivated by avoiding preference for one national language. Later records favour continuity, particularly in the dense reduction terminology of §§23--26. Russian and Ukrainian variants use different stems; from §13 onward some rationales retain nerazložime as a Slavic alternative for appraisal. This is a real alternative, not a rejected candidate assigned an invented lower numerical score.\par
\noindent{\footnotesize Izvorne zapisy / Original records: HIST-g009-003, HIST-g011-003, HIST-g012-003, HIST-g014-004, HIST-g015-004, HIST-g016-004, HIST-g019-003, HIST-g021-003, HIST-g024-003, HIST-g025-003, HIST-g027-004.}\par
\end{minipage}\par
\par\medskip\Needspace{7\baselineskip}
\noindent\begin{minipage}{\linewidth}
\hypertarget{TERM-0936}{}
\noindent\textbf{nereducibilne formy / nereducibilna forma}\quad{\footnotesize\texttt{TERM-0936}}\par
\noindent\emph{Iztočnik / Source:} Irreduzible Formen / Irreduzible Form\par
\noindent\emph{Rusky / Russian:} неприводимые формы / неприводимая форма\par
\noindent\emph{Ukrajinsky / Ukrainian:} незвідні форми / незвідна форма\par
\smallskip\noindent\emph{Razlog:} Nereducibilne formy i nereducibilna forma v rabotě 2, §16, sohranjajut istu semju s redukcijeju i reducentom kak ranějše paragrafy. Jedinstveny i množstveny vid pokazujut gramatično soglasovanje, ne dve različne leksične odluki. Ruski i ukrajinski sravnjajuči vidy podpirajut smysl nereducibilnosti, ale njih različne korene ne možno sbrojiti kak glasove bez naznačenogo pravila. Arhivna motivacija jest kvalitativna: kontinuitet odbornogo izraza, ne izčisljeny rezultat.\par
\noindent\emph{Reason:} Nereducibilne formy and nereducibilna forma in Paper 2, §16, retain the same family with redukcija and reducent as the earlier sections. The singular and plural show grammatical agreement, not two different lexical decisions. The Russian and Ukrainian comparators support the intended irreducibility sense, but their different stems cannot be counted as votes without a stated rule. The archival rationale is qualitative: continuity of technical expression, not a calculated result.\par
\noindent{\footnotesize Izvorne zapisy / Original records: HIST-g017-004.}\par
\end{minipage}\par
\par\medskip\Needspace{7\baselineskip}
\noindent\begin{minipage}{\linewidth}
\hypertarget{TERM-0937}{}
\noindent\textbf{nereducibilne formy / poręd}\quad{\footnotesize\texttt{TERM-0937}}\par
\noindent\emph{Iztočnik / Source:} Irreduzible Formen / Ordnung\par
\noindent\emph{Rusky / Russian:} неприводимые формы / порядок\par
\noindent\emph{Ukrajinsky / Ukrainian:} незвідні форми / порядок\par
\smallskip\noindent\emph{Razlog:} Nereducibilne formy / poręd v rabotě 2, §19, svazuje dve odluki v kontekstě vyčisljenja Sistema III. Prva prodolžuje semju redukcije iz §18. Druga prědava Ordnung, matematičny parametr formy, vidom poręd. Ranějši zapis izrično ostavjaje toj vid za jezykovu ocěnku; slovnik ne směje vydavati jeho za potvrždženu normu. Poręd ne trěba čitati kak porědok stroček v spisku, ni bez konteksta zaměnjati stepenjem: točny pojam zadava lokalna teorija form.\par
\noindent\emph{Reason:} Nereducibilne formy / poręd in Paper 2, §19, links two choices in the enumeration of System III. The first continues the reduction family from §18. The second renders Ordnung, a mathematical parameter of the form, as poręd. The earlier record expressly leaves this form open to linguistic appraisal; the dictionary must not present it as an established norm. Poręd should not be read as the order of rows in a list, or replaced without context by degree: the local theory of forms determines the underlying concept.\par
\noindent{\footnotesize Izvorne zapisy / Original records: HIST-g020-003.}\par
\end{minipage}\par
\par\medskip\Needspace{7\baselineskip}
\noindent\begin{minipage}{\linewidth}
\hypertarget{TERM-0938}{}
\noindent\textbf{nerazložime funkcije}\quad{\footnotesize\texttt{TERM-0938}}\par
\noindent\emph{Iztočnik / Source:} irreduzible Funktionen\par
\noindent\emph{Rusky / Russian:} неприводимые функции\par
\noindent\emph{Ukrajinsky / Ukrainian:} незвідні функції\par
\smallskip\noindent\emph{Razlog:} Nerazložime funkcije v rabotě 6, §8, sohranja iztočny imenik Funktionen v kontekstě polinomovoj oblasti. Prěvod ne zaměnjaje jego tiho sovrěmennym elementi, hotě by se mogli opisati podobni algebraični objekty. Smysl nerazloživosti potrěbuje danu oblast i jej ponimanje děliteljnosti. Transparentny pridevnik ne ustanavjaje ni analitičnu nerozložimost libovoljnoj funkcije, ni univerzalno svojstvo po vsakom razširjenju oblasti.\par
\noindent\emph{Reason:} Nerazložime funkcije in Paper 6, §8, retains the source noun Funktionen in the context of a polynomial domain. It does not silently substitute modern elements, even where similar algebraic objects could be described that way. Irreducibility depends on the stated domain and its notion of divisibility. The transparent adjective establishes neither analytic indecomposability of an arbitrary function nor a universal property under every domain extension.\par
\noindent{\footnotesize Izvorne zapisy / Original records: HIST-g048-008.}\par
\end{minipage}\par
\par\medskip\Needspace{7\baselineskip}
\noindent\begin{minipage}{\linewidth}
\hypertarget{TERM-0939}{}
\noindent\textbf{nerazložima cěla funkcija}\quad{\footnotesize\texttt{TERM-0939}}\par
\noindent\emph{Iztočnik / Source:} irreduzible ganze Funktion\par
\noindent\emph{Rusky / Russian:} неприводимая целая функция\par
\noindent\emph{Ukrajinsky / Ukrainian:} незвідна ціла функція\par
\smallskip\noindent\emph{Razlog:} Nerazložima cěla funkcija v rabotě 6, §§5 i 9, sohranja istorično sočetanje ganze Funktion. V tom algebraičnom kontekstě cěla funkcija jest polinom, ne libovoljna analitična funkcija bez konečnyh singularnostij. Obe motivacije sohranjajut funkcija vměsto tihoj zaměny slovom polinom ili uravnenje. Tak možno objasniti sovrěmenny smysl, ne izbrisavajuči način, kak avtor imenuje objekt. Nerazloživost, kak i prije, jest relativna k razgledajomoj koeficientnoj oblasti.\par
\noindent\emph{Reason:} Nerazložima cěla funkcija in Paper 6, §§5 and 9, retains historical ganze Funktion. In this algebraic context, cěla funkcija is a polynomial, not an arbitrary analytic function without finite singularities. Both rationales keep funkcija rather than silently substituting polynomial or equation. This permits a modern explanatory gloss without erasing the author's designation. Irreducibility remains relative to the coefficient domain under consideration.\par
\noindent{\footnotesize Izvorne zapisy / Original records: HIST-g045-004, HIST-g049-012.}\par
\end{minipage}\par
\par\medskip\Needspace{7\baselineskip}
\noindent\begin{minipage}{\linewidth}
\hypertarget{TERM-0940}{}
\noindent\textbf{nerazložimo jednačenje}\quad{\footnotesize\texttt{TERM-0940}}\par
\noindent\emph{Iztočnik / Source:} irreduzible Gleichung\par
\noindent\emph{Rusky / Russian:} неприводимое уравнение\par
\noindent\emph{Ukrajinsky / Ukrainian:} незвідне рівняння\par
\smallskip\noindent\emph{Razlog:} Nerazložimo jednačenje v rabotě 6, §2, jest izrično svazano s izborom raboty 5. Motivacija opisuje retrospektivno soglasovanje i odklanja ranějše neděljivo, nazvano menje kanoničnym. To poslednje jest sud v zachovanom obosnovanju, ne dokaz nezavisnoj normativnoj publikacije. Ruski i ukrajinski sravnjajuči izrazy podpirajut algebraičnu irreducibilnost; prosto neděljivo moglo by zamučiti, čto imenno ne može byti razloženo. Sravnjeni tekst stvarno imaje nerazložimo jednačenje pri uravnenju za $h(x)$ po odnošenju k danym funkcijam. Relativny kontekst jest sučstven, i sama leksična soglasnost ne proverjaje vsaky matematičny znak uravnenja.\par
\noindent\emph{Reason:} Nerazložimo jednačenje in Paper 6, §2, is explicitly tied to the choice in Paper 5. The rationale describes retrospective alignment and rejects earlier neděljivo as less canonical. That last description is a judgment in the surviving rationale, not evidence of an independent normative publication. The Russian and Ukrainian comparators support algebraic irreducibility; simply indivisible could obscure what is unable to factor. The compared text does use nerazložimo jednačenje for the equation of $h(x)$ relative to the given functions. The relative context is essential, and agreement in wording does not verify every mathematical sign in the equation.\par
\noindent{\footnotesize Izvorne zapisy / Original records: HIST-g042-009.}\par
\smallskip\noindent\emph{Sravnjenje s tekstom:} Sravnjeni red raboty 6 imaje nerazložimo jednačenje za $h(x)$ po odnošenju k ranějše danym funkcijam. To potvrždžaje konkretnu upotrěbu zapisanogo sočetanja i sohranja relativny smysl nerazloživosti. Ne potvrždžaje nezavisnu normativnu attestaciju ni polnu proverku slědujučeho uravnenja.\par
\noindent\emph{Textual comparison:} The compared line in Paper 6 has nerazložimo jednačenje for $h(x)$ relative to the previously given functions. It confirms a concrete occurrence of the recorded combination and preserves the relative sense of irreducibility. It establishes neither independent normative attestation nor full verification of the ensuing equation.\par
\noindent{\footnotesize TEXT-COMP-128; pasaže / passages: 5178; \texttt{CORPUS\_FORM\_COMPARISONS.json}.}\par
\end{minipage}\par
\par\medskip\Needspace{7\baselineskip}
\noindent\begin{minipage}{\linewidth}
\hypertarget{TERM-0941}{}
\noindent\textbf{nerazložimo uravnenje vzhodno k $\Omega_i$}\quad{\footnotesize\texttt{TERM-0941}}\par
\noindent\emph{Iztočnik / Source:} irreduzible Gleichung in bezug auf $\Omega_i$\par
\noindent\emph{Rusky / Russian:} неприводимое относительно $\Omega_i$ уравнение\par
\noindent\emph{Ukrajinsky / Ukrainian:} незвідне відносно $\Omega_i$ рівняння\par
\smallskip\noindent\emph{Razlog:} Nerazložimo uravnenje vzhodno k $\Omega_i$ sohranja relativno uslovje, ot ktorogo zavisi dokaz v rabotě 9, §7. Motivacija svazuje to uravnenje s primitivnymi funkcijami, integralnoju zavisnostju i pozdnějšim argumentom o stepenu. Izpuščenje $\Omega_i$ sdelalo by svojstvo neoprěděljenym. V sravnjenom tekstě jest po odnošenju k $\Omega_i$, pri tom koeficienty uravnenja pripadajut $[\Theta]$. Izmenjenje predložnoj frazy ne odměnjaje matematičnogo uslovja. Arhivno vzhodno k ostaje čest istorije izbora, ne rekomendacija bez sravnjenja s drugimi konstrukcijami.\par
\noindent\emph{Reason:} Nerazložimo uravnenje vzhodno k $\Omega_i$ preserves the relative condition on which the argument in Paper 9, §7 depends. The rationale links the equation to primitive functions, integral dependence and the later degree argument. Omitting $\Omega_i$ would leave the property unspecified. The compared text has po odnošenju k $\Omega_i$, while the equation's coefficients belong to $[\Theta]$. The changed prepositional phrase does not remove the mathematical condition. Archival vzhodno k remains part of the history of the choice, not an unqualified recommendation.\par
\noindent{\footnotesize Izvorne zapisy / Original records: HIST-g065-013.}\par
\smallskip\noindent\emph{Sravnjenje s tekstom:} Odlomok §7 raboty 9 prědpisuje uravnenje po odnošenju k $\Omega_i$ s koeficientami iz $[\Theta]$, potom množi $y$ vedučim koeficientom. Tak ustanavjaje se kontekst relativnogo izraza. V tekstě jest po odnošenju k, ne arhivno vzhodno k.\par
\noindent\emph{Textual comparison:} The passage in Paper 9, §7 specifies an equation relative to $\Omega_i$ with coefficients in $[\Theta]$, then multiplies $y$ by the leading coefficient. This establishes the context of the relative expression. The text has po odnošenju k, not archival vzhodno k.\par
\noindent{\footnotesize TEXT-COMP-129; pasaže / passages: 7939, 7945, 7946, 7947; \texttt{CORPUS\_FORM\_COMPARISONS.json}.}\par
\end{minipage}\par
\par\medskip\Needspace{7\baselineskip}
\noindent\begin{minipage}{\linewidth}
\hypertarget{TERM-0942}{}
\noindent\textbf{nerazložime komponenty / podsistemy}\quad{\footnotesize\texttt{TERM-0942}}\par
\noindent\emph{Iztočnik / Source:} irreduzible Komponenten / Teilsysteme\par
\noindent\emph{Rusky / Russian:} неприводимые компоненты / подсистемы\par
\noindent\emph{Ukrajinsky / Ukrainian:} незвідні компоненти / підсистеми\par
\smallskip\noindent\emph{Razlog:} Nerazložime komponenty / podsistemy v rabotě 27 sohranja razložnu terminologiju i razděljenje sistemov na česti. Odlomok o Hilbertovyh čislah govori o podsistemah danyh čislami nerazložimyh komponentov idealov i o čislah členov kompozicijnogo reda. Proto komponenta i podsistem ne sut dve proizvoljne sinonimne nazvanja: imajut različne roli v odlomku. Izbor ne trěbuje zaměniti vsak podsistem odmahu podmodulom; avtorovo konkretno grupovanje čisel mora ostati vidimo.\par
\noindent\emph{Reason:} Nerazložime komponenty / podsistemy in Paper 27 retains the decomposition vocabulary and the division of systems into parts. The passage on Hilbert numbers discusses subsystems given by counts of indecomposable components of ideals and by lengths of a composition series. Component and subsystem are therefore not arbitrarily interchangeable names: they serve different roles in the passage. The choice does not require replacing every subsystem with submodule; the author's particular grouping of numbers must remain visible.\par
\noindent{\footnotesize Izvorne zapisy / Original records: HIST-g104-007.}\par
\smallskip\noindent\emph{Sravnjenje s tekstom:} V rabotě 27 podsistemy Hilbertovyh čisel dany sut čislami nerazložimyh komponentov idealov i čislami členov kompozicijne serije. Zato dva imenika imajut svazane, ale različne funkcije; ne trěba jih svesti k jednoj univerzalnoj komponentě.\par
\noindent\emph{Textual comparison:} In Paper 27, subsystems of Hilbert numbers are given by counts of indecomposable ideal components and by lengths of a composition series. The two nouns thus have related but different functions; they should not be collapsed into a universal component.\par
\noindent{\footnotesize TEXT-COMP-144; pasaže / passages: 15917; \texttt{CORPUS\_FORM\_COMPARISONS.json}.}\par
\end{minipage}\par
\par\medskip\Needspace{7\baselineskip}
\noindent\begin{minipage}{\linewidth}
\hypertarget{TERM-0944}{}
\noindent\textbf{nereduktibilne nerazložime identičnosti}\quad{\footnotesize\texttt{TERM-0944}}\par
\noindent\emph{Iztočnik / Source:} irreduzible unzerlegbare Identitäten\par
\noindent\emph{Rusky / Russian:} неприводимые неразложимые тождества\par
\noindent\emph{Ukrajinsky / Ukrainian:} нередуковні нерозкладні тотожності\par
\smallskip\noindent\emph{Razlog:} Nereduktibilne nerazložime identičnosti prědava oba pridevnika v irreduzible unzerlegbare Identitäten raboty 4, §3. Motivacija izrično odklanja jih slitje v jedno sovrěmenno slovo. Nereduktibilne upotrěbjaje medžunarodny koren redukcije, nerazložime slovjansky koren razloženja; sravnjajuči prevody takože sohranjajut paru. To jest naměrno terminologično razlikovanje, ne stilično povtorjenje. Za operativne definicije oběh svojstv potrěbny jest kontekst identičnostij; slovnik ne směje sam dobaviti jednu iz modernyh definicij polinomovej irreducibilnosti.\par
\noindent\emph{Reason:} Nereduktibilne nerazložime identičnosti renders both adjectives in irreduzible unzerlegbare Identitäten in Paper 4, §3. The rationale explicitly rejects merging them into one modern word. Nereduktibilne uses the international reduction stem, nerazložime the Slavic decomposition stem; the comparator translations likewise retain a pair. This is an intentional terminological distinction, not stylistic repetition. Operational definitions of the two properties require the context of the identities; the dictionary must not supply a modern definition of polynomial irreducibility in their place.\par
\noindent{\footnotesize Izvorne zapisy / Original records: HIST-g032-002.}\par
\end{minipage}\par
\par\medskip\Needspace{7\baselineskip}
\noindent\begin{minipage}{\linewidth}
\hypertarget{TERM-0945}{}
\noindent\textbf{nerazložime dvustranne klasy}\quad{\footnotesize\texttt{TERM-0945}}\par
\noindent\emph{Iztočnik / Source:} irreduzible zweiseitige Klassen\par
\noindent\emph{Rusky / Russian:} неразложимые двусторонние классы\par
\noindent\emph{Ukrajinsky / Ukrainian:} нерозкладні двобічні класи\par
\smallskip\noindent\emph{Razlog:} Nerazložime dvustranne klasy jest formulacija zavěršajučej teoremy raboty 19, §12. Dvustranne otnosi se k klasam i ne směje odpadnuti: odlomok svazuje jih s ekvivalentnymi matricami i elementarnymi děliteljami, a primětka izrično protivostavlja jednostranne klasy, za ktore takova jedina razložimost ne slěduje. Zato medžu pridevnikami i imenikom sohranja se blizka sintaktična vez. Obšča fraza nerazložime klasy bez ukazanija storonovosti proširila by tvrdženje izvan dokazanogo slučaja.\par
\noindent\emph{Reason:} Nerazložime dvustranne klasy formulates the closing theorem of Paper 19, §12. Dvustranne modifies the classes and cannot be dropped: the passage relates them to equivalent matrices and elementary divisors, while a footnote explicitly contrasts one-sided classes, for which the same unique decomposition does not follow. The adjectives are therefore kept syntactically close to the noun. Bare indecomposable classes would broaden the assertion beyond the case established in the text.\par
\noindent{\footnotesize Izvorne zapisy / Original records: HIST-g081-104.}\par
\smallskip\noindent\emph{Sravnjenje s tekstom:} V rabotě 19 svaz s elementarnymi děliteljami dany jest za dvustranne idealy i klasy ekvivalentnyh matric. Primětka otměčuje, že jednostranne klasy ne imajut togo samogo jednoznačnogo razloženja. Zato dvustranny jest uslovje tvrdženja, ne odstranima rečnikova podrobnost.\par
\noindent\emph{Textual comparison:} Paper 19 supplies the connection with elementary divisors for two-sided ideals and classes of equivalent matrices. The footnote states that one-sided classes lack the same unique decomposition. Two-sided is therefore a condition of the assertion, not a dispensable verbal detail.\par
\noindent{\footnotesize TEXT-COMP-131; pasaže / passages: 13707, 13754; \texttt{CORPUS\_FORM\_COMPARISONS.json}.}\par
\end{minipage}\par
\par\medskip\Needspace{7\baselineskip}
\noindent\begin{minipage}{\linewidth}
\hypertarget{TERM-0946}{}
\noindent\textbf{nerazložima algebraična mnogovidnost}\quad{\footnotesize\texttt{TERM-0946}}\par
\noindent\emph{Iztočnik / Source:} irreduzibles algebraisches Gebilde\par
\noindent\emph{Rusky / Russian:} неприводимое алгебраическое многообразие\par
\noindent\emph{Ukrajinsky / Ukrainian:} незвідний алгебричний многовид\par
\smallskip\noindent\emph{Razlog:} Nerazložima algebraična mnogovidnost interpretuje Gebilde geometrično. Motivacija ukazuje na §10 raboty 19, kde proste idealy i nerazložime algebraične tvorby izrično svazany sut prěz nule polinomov. Rusko многообразие i ukrajinsko многовид podpirajut izbrany odborny registar; medžuslovjansko mnogovidnost jest izrično naměrny tvoreny termin, otvorjeny jezykovoj ocěnce. To ne jest dokaz ranějšej medžuslovjanskoj attestacije. Kontekst takože ne dava prava razuměti mnogovidnost kak gladku mnogovidnost diferencialnoj geometrije: ide o algebraičny objekt oprěděljeny polinomami.\par
\noindent\emph{Reason:} Nerazložima algebraična mnogovidnost interprets Gebilde geometrically. The rationale points to §10 of Paper 19, where prime ideals and irreducible algebraic objects are expressly related through polynomial zeros. Russian многообразие and Ukrainian многовид support the disciplinary register; Interslavic mnogovidnost is explicitly a deliberately constructed term open to linguistic appraisal. This is not evidence of prior Interslavic attestation. Nor does the context license reading mnogovidnost as a smooth manifold of differential geometry: the object is algebraic and specified by polynomials.\par
\noindent{\footnotesize Izvorne zapisy / Original records: HIST-g081-081.}\par
\smallskip\noindent\emph{Sravnjenje s tekstom:} Kontekst raboty 19 svazuje proste idealy s nerazložimymi mnogovidnostjami prěz izčezanje polinomov. To podpiraje geometrično čitanje Gebilde i odlišaje njegovo značenje od gladkoj mnogovidnosti. Sam odlomok ne zasvědčaje mnogovidnost kak ranějši medžuslovjansky termin; to ostaje projektna leksična odluka.\par
\noindent\emph{Textual comparison:} Paper 19 relates prime ideals to irreducible varieties through vanishing polynomials. This supports the geometric reading of Gebilde and distinguishes the sense from a smooth manifold. The passage does not attest mnogovidnost as a pre-existing Interslavic term; that remains the project's lexical decision.\par
\noindent{\footnotesize TEXT-COMP-132; pasaže / passages: 13470, 13471; \texttt{CORPUS\_FORM\_COMPARISONS.json}.}\par
\end{minipage}\par
\par\medskip\Needspace{7\baselineskip}
\noindent\begin{minipage}{\linewidth}
\hypertarget{TERM-0947}{}
\noindent\textbf{nerazložimy ideal}\quad{\footnotesize\texttt{TERM-0947}}\par
\noindent\emph{Iztočnik / Source:} irreduzibles Ideal\par
\noindent\emph{Rusky / Russian:} неразложимый идеал\par
\noindent\emph{Ukrajinsky / Ukrainian:} нерозкладний ідеал\par
\smallskip\noindent\emph{Razlog:} Nerazložimy ideal v rabotě 30, §6, znači ideal, ktory ne možno prědstaviti kak najmenše obče mnogokratno dvoh vlastnyh děliteljev. Ta definicija jest neposrědnja pričina izbora korena razlož-. Ona sohranja razlikovanje od irreducibilnosti elementa pri produktnom razloženju. Rusko неразложимый i ukrajinsko нерозкладний pokazyvajut tu samu smyslovu orientaciju. Slovo ne dava samostojno informaciju o faktorovoj algebrě ili o prostosti ideala.\par
\noindent\emph{Reason:} Nerazložimy ideal in Paper 30, §6 means an ideal that cannot be expressed as the least common multiple of two proper divisors. That definition directly motivates the razlož- stem. It preserves the distinction from irreducibility of an element under product factorization. Russian неразложимый and Ukrainian нерозкладний show the same semantic orientation. The word alone supplies no further assertion about a quotient algebra or primality of the ideal.\par
\noindent{\footnotesize Izvorne zapisy / Original records: HIST-g113-009.}\par
\smallskip\noindent\emph{Sravnjenje s tekstom:} Raboty 19 i 30 objasnjajut nerazloživost pomoču najmenšego občego mnogokratnogo dvoh vlastnyh děliteljev. Pri uslovju cěpov nerazložimy ideal jest primarny, ale kontekst s nulovymi děliteljami razgledaje takože razložime primarne idealy. V rabotě 17 ireducibilne sostavne česti govoret o aditivnom razloženju grupy ostatkov. Te upotrěby ne trěba spojiti s polinomovym razloženjem na faktory bez ukazanija operacije.\par
\noindent\emph{Textual comparison:} Papers 19 and 30 explain irreducibility by the least common multiple of two proper divisors. Under the chain condition an irreducible ideal is primary, but the zero-divisor context also considers reducible primary ideals. In Paper 17 irreducible constituents concern additive decomposition of a residue group. These usages must not be merged with polynomial factorization without specifying the operation.\par
\noindent{\footnotesize TEXT-COMP-126; pasaže / passages: 12424, 16680, 16686, 17147, 17149, 11446; \texttt{CORPUS\_FORM\_COMPARISONS.json}.}\par
\end{minipage}\par
\par\medskip\Needspace{7\baselineskip}
\noindent\begin{minipage}{\linewidth}
\hypertarget{TERM-0948}{}
\noindent\textbf{izolovane najveće primarne idealy}\quad{\footnotesize\texttt{TERM-0948}}\par
\noindent\emph{Iztočnik / Source:} isolierte größte primäre Ideale\par
\noindent\emph{Rusky / Russian:} изолированные наибольшие примарные идеалы\par
\noindent\emph{Ukrajinsky / Ukrainian:} ізольовані найбільші примарні ідеали\par
\smallskip\noindent\emph{Razlog:} Izolovane najveće primarne idealy sohranja vse tri kvalifikacije iz dodatka k rabotě 19. Zapis svazuje jedinost iz §5 s pozdnějšim rezultatom §7. Izolovane ne jest zaměnjeno sovrěmennym nazvanjem nekogo klasa pridruženyh prostyh idealov; največši primarny ideal i pridruženy prosty ideal ostajut različnymi objektami. Teorema XIII govori o jednoznačnosti izolovanyh komponentov, ne vseh primarnyh komponentov bez različenija. Nynějši tekst piše največših vměsto arhivnogo najveće; ta razlika formy takože ostaje vidima v sravnjenju.\par
\noindent\emph{Reason:} Izolovane najveće primarne idealy retains all three qualifications in the addendum to Paper 19. The record links the uniqueness discussion in §5 to the later result in §7. Isolated is not replaced by a modern name for a category of associated primes; a greatest primary ideal and its associated prime remain different objects. Theorem XIII concerns uniqueness of isolated components, not all primary components indiscriminately. The present text has največših rather than archival najveće; the comparison also keeps that difference in form visible.\par
\noindent{\footnotesize Izvorne zapisy / Original records: HIST-g081-047, HIST-g082-047.}\par
\smallskip\noindent\emph{Sravnjenje s tekstom:} Prědbežne odkazy raboty 19 vodat k Definiciji VIa i Teoremě XIII. Definicija sohranja uslovja na pridružene proste idealy i prědstavjenje, ale ne zavisi od specialnogo komplementa. Jedinost ne proširjaje se avtomatično na neizolovane komponenty. V rabotě 24 osnovny ideal jest izolovana komponenta oprěděljena cěloju sovokupnostju prostyh idealov; zato komponenta ne znači vsegda jeden primarny ideal. Rabota 30 povtorjaje tvrdženje jedinosti v podnožnoj primětke.\par
\noindent\emph{Textual comparison:} Forward references in Paper 19 lead to Definition VIa and Theorem XIII. The definition retains conditions on associated primes and the representation while becoming independent of a particular complement. Uniqueness does not automatically extend to non-isolated components. In Paper 24 a basic ideal is an isolated component determined by a whole collection of primes; component therefore does not always mean a single primary ideal. Paper 30 repeats the uniqueness assertion in a footnote.\par
\noindent{\footnotesize TEXT-COMP-133; pasaže / passages: 12950, 13026, 13222, 13244, 15544, 16719; \texttt{CORPUS\_FORM\_COMPARISONS.json}.}\par
\end{minipage}\par
\par\medskip\Needspace{7\baselineskip}
\noindent\begin{minipage}{\linewidth}
\hypertarget{TERM-0949}{}
\noindent\textbf{izolovane idealy}\quad{\footnotesize\texttt{TERM-0949}}\par
\noindent\emph{Iztočnik / Source:} isolierte Ideale\par
\noindent\emph{Rusky / Russian:} изолированные идеалы\par
\noindent\emph{Ukrajinsky / Ukrainian:} ізольовані ідеали\par
\smallskip\noindent\emph{Razlog:} Izolovane idealy jest kratko nazvanje v prědbežnom odkazě raboty 19, §4, na teoremu jedinosti v §7. Motivacija sohranja izrično nazvanje i ne dodava polnu definiciju v podnožnu primětku. Slovnik može dati objasnjenje i odkaz, ale ne směje z takogo kratkogo pojavljanja zaključiti, že lokalny argument už dokazal pozdnějšu teoremu. Izolacija jest svojstvo v kontekstě idealnogo razloženja, ne prostorna odděljenost množiny bodov.\par
\noindent\emph{Reason:} Izolovane idealy is the short designation in the forward reference from Paper 19, §4 to the uniqueness theorem in §7. The rationale keeps the name without inserting its full definition into the footnote. The dictionary may supply an explanation and reference, but must not infer from the short occurrence that the local argument already proves the later theorem. Isolation is a property within ideal decomposition, not spatial separation of a set of points.\par
\noindent{\footnotesize Izvorne zapisy / Original records: HIST-g081-013, HIST-g082-013.}\par
\smallskip\noindent\emph{Sravnjenje s tekstom:} Prědbežne odkazy raboty 19 vodat k Definiciji VIa i Teoremě XIII. Definicija sohranja uslovja na pridružene proste idealy i prědstavjenje, ale ne zavisi od specialnogo komplementa. Jedinost ne proširjaje se avtomatično na neizolovane komponenty. V rabotě 24 osnovny ideal jest izolovana komponenta oprěděljena cěloju sovokupnostju prostyh idealov; zato komponenta ne znači vsegda jeden primarny ideal. Rabota 30 povtorjaje tvrdženje jedinosti v podnožnoj primětke.\par
\noindent\emph{Textual comparison:} Forward references in Paper 19 lead to Definition VIa and Theorem XIII. The definition retains conditions on associated primes and the representation while becoming independent of a particular complement. Uniqueness does not automatically extend to non-isolated components. In Paper 24 a basic ideal is an isolated component determined by a whole collection of primes; component therefore does not always mean a single primary ideal. Paper 30 repeats the uniqueness assertion in a footnote.\par
\noindent{\footnotesize TEXT-COMP-133; pasaže / passages: 12950, 13026, 13222, 13244, 15544, 16719; \texttt{CORPUS\_FORM\_COMPARISONS.json}.}\par
\end{minipage}\par
\par\medskip\Needspace{7\baselineskip}
\noindent\begin{minipage}{\linewidth}
\hypertarget{TERM-0950}{}
\noindent\textbf{izolovana komponenta}\quad{\footnotesize\texttt{TERM-0950}}\par
\noindent\emph{Iztočnik / Source:} isolierte Komponente\par
\noindent\emph{Rusky / Russian:} изолированная компонента\par
\noindent\emph{Ukrajinsky / Ukrainian:} ізольована компонента\par
\smallskip\noindent\emph{Razlog:} Izolovana komponenta sohranja komponentu kak čest razloženja v rabotě 30, §6. V sravnjenoj primětke pridružene proste idealy i izolovane komponenty sut jednoznačno oprěděljene. To ne pravi vsaku komponentu izolovanoju. Ranějša motivacija nazyvaje izraz standardnym i daje paralele vo vseh treh prevodah; te paralele podpirajut izbor, ale samy po sobě ne ustanavjajut nezavisnu medžuslovjansku normu. Izolovana jest tehničny kvalifikator, ne zaměna polnoj definicije.\par
\noindent\emph{Reason:} Izolovana komponenta preserves component as a constituent of the decomposition in Paper 30, §6. In the compared footnote, associated prime ideals and isolated components are uniquely determined. That does not make every component isolated. The earlier rationale calls the phrase standard and supplies parallels in all three translations; these support the choice but do not by themselves establish an independent Interslavic norm. Isolated is a technical qualifier, not a substitute for the full definition.\par
\noindent{\footnotesize Izvorne zapisy / Original records: HIST-g113-012.}\par
\smallskip\noindent\emph{Sravnjenje s tekstom:} Prědbežne odkazy raboty 19 vodat k Definiciji VIa i Teoremě XIII. Definicija sohranja uslovja na pridružene proste idealy i prědstavjenje, ale ne zavisi od specialnogo komplementa. Jedinost ne proširjaje se avtomatično na neizolovane komponenty. V rabotě 24 osnovny ideal jest izolovana komponenta oprěděljena cěloju sovokupnostju prostyh idealov; zato komponenta ne znači vsegda jeden primarny ideal. Rabota 30 povtorjaje tvrdženje jedinosti v podnožnoj primětke.\par
\noindent\emph{Textual comparison:} Forward references in Paper 19 lead to Definition VIa and Theorem XIII. The definition retains conditions on associated primes and the representation while becoming independent of a particular complement. Uniqueness does not automatically extend to non-isolated components. In Paper 24 a basic ideal is an isolated component determined by a whole collection of primes; component therefore does not always mean a single primary ideal. Paper 30 repeats the uniqueness assertion in a footnote.\par
\noindent{\footnotesize TEXT-COMP-133; pasaže / passages: 12950, 13026, 13222, 13244, 15544, 16719; \texttt{CORPUS\_FORM\_COMPARISONS.json}.}\par
\end{minipage}\par
\par\medskip\Needspace{7\baselineskip}
\noindent\begin{minipage}{\linewidth}
\hypertarget{TERM-0951}{}
\noindent\textbf{izolovana komponenta razloženja}\quad{\footnotesize\texttt{TERM-0951}}\par
\noindent\emph{Iztočnik / Source:} isolierte Komponente der Zerlegung\par
\noindent\emph{Rusky / Russian:} изолированная компонента разложения\par
\noindent\emph{Ukrajinsky / Ukrainian:} ізольована компонента розкладу\par
\smallskip\noindent\emph{Razlog:} Izolovana komponenta razloženja v rabotě 24 imenuje objekt razložnoj teorije, ne pridruženy prosty ideal sam. Motivacija želela sohraniti isti tip objekta v ruskom, ukrajinskom i medžuslovjanskom tekstě. Teorema VIII tak imenuje osnovne idealy, oprěděljene cěloju sovokupnostju pridruženyh prostyh idealov ukazanyh dimenzij. Zato tu komponenta ne mora značiti samo jeden primarny ideal. Genitiv razloženja pomaže čitatelju viděti, čego jest ona čest, i ne trěba go zamenjati sovrěmennoju skračenoju etiketkoju, ktora bi suzila avtorov smysl.\par
\noindent\emph{Reason:} Izolovana komponenta razloženja in Paper 24 names an object of decomposition theory, not the associated prime ideal itself. The rationale aims to preserve the same type of object in Russian, Ukrainian and Interslavic. Theorem VIII uses the designation for basic ideals determined by a whole collection of associated primes of specified dimensions. Component here therefore need not mean just one primary ideal. The genitive razloženja helps readers see what it is a constituent of; a modern abbreviated label should not narrow the author's sense.\par
\noindent{\footnotesize Izvorne zapisy / Original records: HIST-g099-003.}\par
\smallskip\noindent\emph{Sravnjenje s tekstom:} Prědbežne odkazy raboty 19 vodat k Definiciji VIa i Teoremě XIII. Definicija sohranja uslovja na pridružene proste idealy i prědstavjenje, ale ne zavisi od specialnogo komplementa. Jedinost ne proširjaje se avtomatično na neizolovane komponenty. V rabotě 24 osnovny ideal jest izolovana komponenta oprěděljena cěloju sovokupnostju prostyh idealov; zato komponenta ne znači vsegda jeden primarny ideal. Rabota 30 povtorjaje tvrdženje jedinosti v podnožnoj primětke.\par
\noindent\emph{Textual comparison:} Forward references in Paper 19 lead to Definition VIa and Theorem XIII. The definition retains conditions on associated primes and the representation while becoming independent of a particular complement. Uniqueness does not automatically extend to non-isolated components. In Paper 24 a basic ideal is an isolated component determined by a whole collection of primes; component therefore does not always mean a single primary ideal. Paper 30 repeats the uniqueness assertion in a footnote.\par
\noindent{\footnotesize TEXT-COMP-133; pasaže / passages: 12950, 13026, 13222, 13244, 15544, 16719; \texttt{CORPUS\_FORM\_COMPARISONS.json}.}\par
\end{minipage}\par
\par\medskip\Needspace{7\baselineskip}
\noindent\begin{minipage}{\linewidth}
\hypertarget{TERM-0952}{}
\noindent\textbf{izolovany ideal / neizolovany ideal}\quad{\footnotesize\texttt{TERM-0952}}\par
\noindent\emph{Iztočnik / Source:} isoliertes Ideal / nicht-isoliertes Ideal\par
\noindent\emph{Rusky / Russian:} изолированный идеал / неизолированный идеал\par
\noindent\emph{Ukrajinsky / Ukrainian:} ізольований ідеал / неізольований ідеал\par
\smallskip\noindent\emph{Razlog:} Izolovany ideal / neizolovany ideal v rabotě 19, §7, sohranja dvojicu, na kotoroj osnovana jest teorema jedinosti. Definicija VIa čini izolaciju nezavisnoju od izbranogo komplementa, pri izričnyh uslovjah na prědstavjenje i pridružene proste idealy. Teorema XIII potom ograničuje mnogoznačnost na neizolovane komponenty. Ne- zato nosi matematično značenje. Macaulayevo istorično imbedded jest citovano v svojem kontekstě; ono ne jest dovoljenjem povsudu tiho zaměniti avtorovu paru slov modernym embedded/non-embedded.\par
\noindent\emph{Reason:} Izolovany ideal / neizolovany ideal in Paper 19, §7 preserves the opposition underlying the uniqueness theorem. Definition VIa makes isolation independent of the chosen complement, under explicit conditions on the representation and associated primes. Theorem XIII then confines possible nonuniqueness to the non-isolated components. The negative prefix therefore carries mathematical content. Macaulay's historical imbedded is cited in its own context; it does not license silently replacing the author's pair throughout by modern embedded/non-embedded.\par
\noindent{\footnotesize Izvorne zapisy / Original records: HIST-g081-057, HIST-g082-057.}\par
\smallskip\noindent\emph{Sravnjenje s tekstom:} Prědbežne odkazy raboty 19 vodat k Definiciji VIa i Teoremě XIII. Definicija sohranja uslovja na pridružene proste idealy i prědstavjenje, ale ne zavisi od specialnogo komplementa. Jedinost ne proširjaje se avtomatično na neizolovane komponenty. V rabotě 24 osnovny ideal jest izolovana komponenta oprěděljena cěloju sovokupnostju prostyh idealov; zato komponenta ne znači vsegda jeden primarny ideal. Rabota 30 povtorjaje tvrdženje jedinosti v podnožnoj primětke.\par
\noindent\emph{Textual comparison:} Forward references in Paper 19 lead to Definition VIa and Theorem XIII. The definition retains conditions on associated primes and the representation while becoming independent of a particular complement. Uniqueness does not automatically extend to non-isolated components. In Paper 24 a basic ideal is an isolated component determined by a whole collection of primes; component therefore does not always mean a single primary ideal. Paper 30 repeats the uniqueness assertion in a footnote.\par
\noindent{\footnotesize TEXT-COMP-133; pasaže / passages: 12950, 13026, 13222, 13244, 15544, 16719; \texttt{CORPUS\_FORM\_COMPARISONS.json}.}\par
\end{minipage}\par
\par\medskip\Needspace{7\baselineskip}
\noindent\begin{minipage}{\linewidth}
\hypertarget{TERM-0953}{}
\noindent\textbf{izomorfne / izomorfna sootvětnost / izomorfno odobraženje}\quad{\footnotesize\texttt{TERM-0953}}\par
\noindent\emph{Iztočnik / Source:} isomorph / isomorphe Beziehung / isomorphe Abbildung\par
\noindent\emph{Rusky / Russian:} изоморфные / изоморфное соответствие / изоморфное отображение\par
\noindent\emph{Ukrajinsky / Ukrainian:} ізоморфні / ізоморфне відношення / ізоморфне відображення\par
\smallskip\noindent\emph{Razlog:} Izomorfne / izomorfna sootvětnost / izomorfno odobraženje v rabotě 9, §8, sohranja medžunarodny koren i razlikuje sootvětnost od imenovanogo odobraženja. Ruski par соответствие i отображение daje prěmu sravnjajuču oporu; ukrajinski odnos i odobraženje pokazyvajut podobno razlikovanje, ne identičnu sintaksu. Tak prěvod može nazvati i relaciju medžu oblastjami, i konkretno odobraženje. Izomorfija ne znači, že dve oblasti sut bukvalno jedna i ta sama podmnožina.\par
\noindent\emph{Reason:} Izomorfne / izomorfna sootvětnost / izomorfno odobraženje in Paper 9, §8 retains the international stem and distinguishes correspondence from a designated map. Russian соответствие and отображение provide direct comparators; Ukrainian relation and mapping show a similar distinction, not identical syntax. The translation can thus name both a relation between domains and a particular mapping. Isomorphism does not mean that the two domains are literally the same subset.\par
\noindent{\footnotesize Izvorne zapisy / Original records: HIST-g066-003.}\par
\end{minipage}\par
\par\medskip\Needspace{7\baselineskip}
\noindent\begin{minipage}{\linewidth}
\hypertarget{TERM-0954}{}
\noindent\textbf{izomorfny / izomorfija}\quad{\footnotesize\texttt{TERM-0954}}\par
\noindent\emph{Iztočnik / Source:} isomorph / Isomorphie\par
\noindent\emph{Rusky / Russian:} изоморфный / изоморфизм\par
\noindent\emph{Ukrajinsky / Ukrainian:} ізоморфний / ізоморфізм\par
\smallskip\noindent\emph{Razlog:} Izomorfny / izomorfija v rabotě 22, §4, sohranja medžunarodny matematičny koren. Motivacija svazuje jego s definicijeju prěz sistemy prědstaviteljev i s nazvanjem teoremy: volny opis mogl by zakryti tu vez. Pridevnik naziva svojstvo ili odnošenje, imenik izomorfija samo odnošenje ili odobraženje po kontekstu. Rusko i ukrajinsko nazvanja na -izm ne objazujut medžuslovjanski avtomatično vybrati isto zakončenje; v dokumentah jest izbrano -ija.\par
\noindent\emph{Reason:} Izomorfny / izomorfija in Paper 22, §4 retains the international mathematical stem. The rationale links it to the definition through systems of representatives and to the theorem's name: a free paraphrase could obscure that connection. The adjective names a property or relation, while izomorfija names the relation or map as the context requires. The Russian and Ukrainian nouns ending in -izm do not automatically require the same Interslavic ending; the documented choice is -ija.\par
\noindent{\footnotesize Izvorne zapisy / Original records: HIST-g088-003, HIST-g092-034.}\par
\end{minipage}\par
\par\medskip\Needspace{7\baselineskip}
\noindent\begin{minipage}{\linewidth}
\hypertarget{TERM-0955}{}
\noindent\textbf{izomorfno odobraženje}\quad{\footnotesize\texttt{TERM-0955}}\par
\noindent\emph{Iztočnik / Source:} isomorphe Abbildung\par
\noindent\emph{Rusky / Russian:} изоморфное отображение\par
\noindent\emph{Ukrajinsky / Ukrainian:} ізоморфне відображення\par
\smallskip\noindent\emph{Razlog:} Izomorfno odobraženje v uvodu raboty 10 prodolžuje ranějšu terminologiju odobraženij. Rusko отображение i ukrajinsko відображення podpirajut izbranu semju; pridevnik sohranja medžunarodno izomorf-. Motivacija jest i smyslova, i tekstova: odobraženje dolžno byti razpoznavajemo kak toj samy tip konstrukcije v različnyh čestjah prevoda. To ne otměnjaje potrěbu pri konkretnom odobraženju ukazati objekty i strukturu, ktora sohranjaje se.\par
\noindent\emph{Reason:} Izomorfno odobraženje at the opening of Paper 10 continues the earlier mapping vocabulary. Russian отображение and Ukrainian відображення support the chosen family; the adjective retains international izomorf-. The reason is semantic and textual: readers should recognize the same type of construction across the translation. This does not remove the need to specify the objects and the preserved structure for a particular map.\par
\noindent{\footnotesize Izvorne zapisy / Original records: HIST-g069-002.}\par
\end{minipage}\par
\par\medskip\Needspace{7\baselineskip}
\noindent\begin{minipage}{\linewidth}
\hypertarget{TERM-0956}{}
\noindent\textbf{izomorfna prirědba / nulovy klas / dělitelji nuly}\quad{\footnotesize\texttt{TERM-0956}}\par
\noindent\emph{Iztočnik / Source:} isomorphe Zuordnung / Nullklasse / Nullteiler\par
\noindent\emph{Rusky / Russian:} изоморфное сопоставление / нулевой класс / делители нуля\par
\noindent\emph{Ukrajinsky / Ukrainian:} ізоморфне співвіднесення / нульовий клас / дільники нуля\par
\smallskip\noindent\emph{Razlog:} Izomorfna prirědba / nulovy klas / dělitelji nuly sut sočetane termini dokaza Teoremy VI v rabotě 24. Prirědba izbrana jest za formalno prisvojenje, različeno od volnějšego sootvętstvovati v prose. Nulovy klas jest klas ostatkov, ktory predstavjaje nulu; dělitelj nuly jest element s drugim svojstvom v prstenu. Frazy ne smějut sliti te dva pojma. Motivacija ne govori, že samo porěd nazvanj už ustanavjaje izomorfiju; sootvětnost i operacije morajut byti dany argumentom.\par
\noindent\emph{Reason:} Izomorfna prirědba / nulovy klas / dělitelji nuly are related terms in the proof of Theorem VI in Paper 24. Prirědba is selected for a formal assignment, distinct from looser sootvętstvovati in prose. A zero class is the residue class representing zero; a zero divisor has a different property in the ring. The wording must not merge those concepts. The rationale does not imply that the list of names establishes an isomorphism: the correspondence and operations must be supplied by the argument.\par
\noindent{\footnotesize Izvorne zapisy / Original records: HIST-g098-010.}\par
\end{minipage}\par
\par\medskip\Needspace{7\baselineskip}
\noindent\begin{minipage}{\linewidth}
\hypertarget{TERM-0957}{}
\noindent\textbf{matrično kolco, izomorfno ... pri tom ekvivalentno}\quad{\footnotesize\texttt{TERM-0957}}\par
\noindent\emph{Iztočnik / Source:} isomorpher, und zwar äquivalenter, Matrizenring\par
\noindent\emph{Rusky / Russian:} матричное кольцо, изоморфное ... причем эквивалентное\par
\noindent\emph{Ukrajinsky / Ukrainian:} матричне кільце, ізоморфне ... причому еквівалентне\par
\smallskip\noindent\emph{Razlog:} Matrično kolco, izomorfno ... pri tom ekvivalentno sohranja dva nivoja tvrdženja raboty 31, §4. Ekvivalentnost jest konkretna konjugacija matričnyh prstenov pri zaměně bazy, ne prosto stilističny sinonim izomorfije. Sravnjeni tekst dava $R_{\bar{\mathfrak a}}=P^{-1}R_{\mathfrak a}P$ i uslovje, že determinant matrice $P$ jest jedinica koeficientnoj oblasti. Nyně tam jest matrično prsten vměsto arhivnogo matrično kolco. Takova zaměna imenika potrěbuje osobnu proverku roda i soglasovanja; ona ne mjenja potrebu sohraniti oba matematična odnošenja.\par
\noindent\emph{Reason:} Matrično kolco, izomorfno ... pri tom ekvivalentno preserves two levels of assertion in Paper 31, §4. Equivalence is the specific conjugacy of matrix rings under a change of basis, not a stylistic synonym for isomorphism. The compared text gives $R_{\bar{\mathfrak a}}=P^{-1}R_{\mathfrak a}P$ and requires the determinant of $P$ to be a unit of the coefficient domain. It now has matrično prsten rather than archival matrično kolco. This noun substitution requires a separate gender-and-agreement check; it does not change the need to preserve both mathematical relations.\par
\noindent{\footnotesize Izvorne zapisy / Original records: HIST-g132-002.}\par
\smallskip\noindent\emph{Sravnjenje s tekstom:} Sravnjeni odlomok raboty 31 dava konkretno konjugacijsko ravnenstvo s $P^{-1}$ i $P$; determinant $P$ jest jedinica. To objasnjaje, počemu izomorfno i ekvivalentno sohranjajut se odvojeno. Matrično prsten v nynějšem prevodu takože ukazuje na potrěbu proveriti soglasovanje po zaměně imenika kolco.\par
\noindent\emph{Textual comparison:} The compared passage in Paper 31 gives a specific conjugacy equation using $P^{-1}$ and $P$; the determinant of $P$ is a unit. This explains why isomorphic and equivalent remain separate. Matrično prsten in the present translation also points to the need to check agreement after replacing the noun kolco.\par
\noindent{\footnotesize TEXT-COMP-134; pasaže / passages: 17909, 17911, 17913; \texttt{CORPUS\_FORM\_COMPARISONS.json}.}\par
\end{minipage}\par
\par\medskip\Needspace{7\baselineskip}
\noindent\begin{minipage}{\linewidth}
\hypertarget{TERM-0958}{}
\noindent\textbf{teoremy o izomorfiji pri aditivnom razkladu}\quad{\footnotesize\texttt{TERM-0958}}\par
\noindent\emph{Iztočnik / Source:} Isomorphiesätze bei der additiven Zerlegung\par
\noindent\emph{Rusky / Russian:} теоремы об изоморфизме при аддитивном разложении\par
\noindent\emph{Ukrajinsky / Ukrainian:} теореми про ізоморфізм при адитивному розкладі\par
\smallskip\noindent\emph{Razlog:} Teoremy o izomorfiji pri aditivnom razkladu v rabotě 28 sohranja strukturny obhvat iztočnogo izraza. Pri aditivnom razkladu ne jest opustimo dopolnjenje: ukazuje, v kakoj konstrukciji teoremy upotrěbjajut se. Motivacija svazuje izraz s ranějšimi rabotami 17 i 19, ne s proizvoljno vybranym sovrěmennym spiskom teoremov o izomorfizmah. Internacionalna izomorfija i aditivny ostajut razpoznavajemy odborny semji; razklad daje nazvanje operacije razloženja.\par
\noindent\emph{Reason:} Teoremy o izomorfiji pri aditivnom razkladu in Paper 28 retains the source phrase's structural scope. Pri aditivnom razkladu is not dispensable: it specifies the construction in which the theorems are used. The rationale connects the phrase to Papers 17 and 19, not to an arbitrarily chosen modern list of isomorphism theorems. International izomorfija and aditivny remain recognizable technical families, while razklad names the decomposition operation.\par
\noindent{\footnotesize Izvorne zapisy / Original records: HIST-g105-005.}\par
\end{minipage}\par
\par\medskip\Needspace{7\baselineskip}
\noindent\begin{minipage}{\linewidth}
\hypertarget{TERM-0959}{}
\noindent\textbf{teoremy o izomorfiji; direktne sumy}\quad{\footnotesize\texttt{TERM-0959}}\par
\noindent\emph{Iztočnik / Source:} Isomorphiesätze. Direkte Summen\par
\noindent\emph{Rusky / Russian:} теоремы об изоморфизме; прямые суммы\par
\noindent\emph{Ukrajinsky / Ukrainian:} теореми про ізоморфізм; прямі суми\par
\smallskip\noindent\emph{Razlog:} Teoremy o izomorfiji; direktne sumy prědava dvočestny naslov §4 raboty 30. Pervy člen sohranja ranějše izomorfija, drugy ranějše direktna suma iz raboty 28. To dava kontinuitet čitatelju, ale ne svodi dve temy k jednoj. Interpunkcija razděljaje naslovne temy; ne znamenaje, že direktna suma jest synonym izomorfije. Prěvod ne zaměnjaje avtorov poredok sovrěmennoju učebničnoju klasifikacijeju.\par
\noindent\emph{Reason:} Teoremy o izomorfiji; direktne sumy renders the two-part title of §4 in Paper 30. The first part retains earlier izomorfija, the second direktna suma from Paper 28. This gives readers continuity without reducing the two topics to one. Punctuation separates the topics in the heading; it does not make direct sum synonymous with isomorphism. The translation does not replace the author's organization with a modern textbook classification.\par
\noindent{\footnotesize Izvorne zapisy / Original records: HIST-g111-001.}\par
\end{minipage}\par
\par\medskip\Needspace{7\baselineskip}
\noindent\begin{minipage}{\linewidth}
\hypertarget{TERM-0960}{}
\noindent\textbf{izomorfija / izomorfno prirěditi}\quad{\footnotesize\texttt{TERM-0960}}\par
\noindent\emph{Iztočnik / Source:} Isomorphismus / isomorph zuordnen\par
\noindent\emph{Rusky / Russian:} изоморфизм / изоморфно поставить в соответствие\par
\noindent\emph{Ukrajinsky / Ukrainian:} ізоморфізм / ізоморфно поставити у відповідність\par
\smallskip\noindent\emph{Razlog:} Izomorfija / izomorfno prirěditi v rabotě 24 sohranja medžunarodny imenik i daje glagol za formalno prirědženje. Motivacija izrično razlikuje prirěditi od volnějšego sootvětstvovati. Tak tekst govori ne samo, že dva objekty kažut se sličnymi, ale že konstruirova se oprěděljena sootvětnost. Koren izomorf- jest obči vsim trema zapisanym prevodam; izbor konca -ija i glagola ostaje konkretnym medžuslovjanskim rěšenjem.\par
\noindent\emph{Reason:} Izomorfija / izomorfno prirěditi in Paper 24 keeps the international noun and supplies a verb for a formal assignment. The rationale expressly distinguishes prirěditi from looser sootvětstvovati. The text thus says more than that two objects seem similar: it constructs a specified correspondence. The izomorf- stem is shared by all three recorded translations; the -ija ending and the verb remain particular Interslavic choices.\par
\noindent{\footnotesize Izvorne zapisy / Original records: HIST-g097-008.}\par
\end{minipage}\par
\par\medskip\Needspace{7\baselineskip}
\noindent\begin{minipage}{\linewidth}
\hypertarget{TERM-0961}{}
\noindent\textbf{izomorfizm kolc klasov ostatkov}\quad{\footnotesize\texttt{TERM-0961}}\par
\noindent\emph{Iztočnik / Source:} Isomorphismus der Restklassenringe\par
\noindent\emph{Rusky / Russian:} изоморфизм колец классов вычетов\par
\noindent\emph{Ukrajinsky / Ukrainian:} ізоморфізм кілець класів лишків\par
\smallskip\noindent\emph{Razlog:} Izomorfizm kolc klasov ostatkov v rabotě 31 sohranja polnu konstrukciju Restklassenringe. Rěč jest o prstenovoj strukturě, ne samo o bijekciji množin klasov. Motivacija odklanja skračenje do samogo nazvanja faktor-prstenov, zato že ono moglo by zakryti izričny odkaz na ostatkove klasy v dokazě. To ne ustanavjaje različnost medžu prstenom klasov ostatkov i odgovarjajučim faktor-prstenom; ide o način predstavjenja. Zdes imenik jest izomorfizm, dokolě druge zapisy izbirajut izomorfija. Oba istoricky zasvědčena vybora ostajut vidimy.\par
\noindent\emph{Reason:} Izomorfizm kolc klasov ostatkov in Paper 31 preserves the full Restklassenringe construction. The assertion concerns ring structure, not merely a bijection between sets of classes. The rationale rejects shortening the expression to the name quotient rings because that might conceal the proof's express reference to residue classes. This does not establish a mathematical difference between a residue-class ring and the corresponding quotient ring; it concerns presentation. Here the noun is izomorfizm, while other records select izomorfija. Both historically recorded choices remain visible.\par
\noindent{\footnotesize Izvorne zapisy / Original records: HIST-g159-009.}\par
\end{minipage}\par
\par\medskip\Needspace{7\baselineskip}
\noindent\begin{minipage}{\linewidth}
\hypertarget{TERM-0963}{}
\noindent\textbf{iteracija}\quad{\footnotesize\texttt{TERM-0963}}\par
\noindent\emph{Iztočnik / Source:} Iteration\par
\noindent\emph{Rusky / Russian:} итерация\par
\noindent\emph{Ukrajinsky / Ukrainian:} ітерація\par
\smallskip\noindent\emph{Razlog:} Iteracija v rabotě 2, §2, izbrana jest kak kratky medžunarodny matematičny termin s prěmymi ruskymi i ukrajinskymi paralelami. Sravnjeni odlomok govori o opakovanoj priměně postupka, iz kotoroj nastavaje razvitje; koeficienty računajut se rekursivno. Zato iteracija tu ne naziva avtomatično čislenu približajemu metodu s dokazanoju shodnostju. Izbor sohranja obči pojam povtorjenja, a konkretne svojstva dava matematičny kontekst.\par
\noindent\emph{Reason:} Iteracija in Paper 2, §2 was chosen as a compact international mathematical term with direct Russian and Ukrainian parallels. The compared passage concerns repeated application of a procedure yielding an expansion; its coefficients are calculated recursively. Iteration here therefore does not automatically name a numerical approximation method with proved convergence. The choice preserves the general notion of repetition, while the mathematical context supplies its particular properties.\par
\noindent{\footnotesize Izvorne zapisy / Original records: HIST-g004-012.}\par
\smallskip\noindent\emph{Sravnjenje s tekstom:} Razvitje dobyvaje se iteracijeju, a koeficienty izračunajut se rekursivno. Sravnjeni odlomok daje konkretno algebraično ponimanje povtorjenja; sam imenik ne govori o čislenoj shodnosti.\par
\noindent\emph{Textual comparison:} An expansion is obtained by iteration, and the coefficients are calculated recursively. The compared passage gives a concrete algebraic meaning to repetition; the noun itself makes no claim about numerical convergence.\par
\noindent{\footnotesize TEXT-COMP-143; pasaže / passages: 515, 523; \texttt{CORPUS\_FORM\_COMPARISONS.json}.}\par
\end{minipage}\par
\par\medskip\Needspace{7\baselineskip}
\noindent\begin{minipage}{\linewidth}
\hypertarget{TERM-0966}{}
\noindent\textbf{vsaka potenca indeksa / vsaky parny stepen}\quad{\footnotesize\texttt{TERM-0966}}\par
\noindent\emph{Iztočnik / Source:} jede Potenz des Index / jeder gerade Grad\par
\noindent\emph{Rusky / Russian:} каждая степень индекса / любая четная степень\par
\noindent\emph{Ukrajinsky / Ukrainian:} кожний степінь індексу / кожний парний степінь\par
\smallskip\noindent\emph{Razlog:} Vsaka potenca indeksa / vsaky parny stepen razlikuje čislenu potencu i stepen poljovogo razširjenja v rabotě 32, §2. Potenz ne jest neoprěděljena sila ili intensivnost. Glavny tekst pokazuje neograničenost stepenov minimalnyh razpadnyh polj i pojavjenje vsakoj potence indeksa. Podnožna primětka pripisuje R. Braueru silnějše tvrdženje o minimalnyh razpadnyh poljah kvaternionnogo těla vsakogo parnogo stepena. Ne trěba pripisati silnějši rezultat ranějšemu argumentu ni poměšati minimalno polje s poljem najmenšego čislenogo stepena.\par
\noindent\emph{Reason:} Vsaka potenca indeksa / vsaky parny stepen distinguishes a numerical power from the degree of a field extension in Paper 32, §2. Potenz is not vague strength or intensity. The main text establishes unbounded degrees of minimal splitting fields and the occurrence of every power of the index. A footnote attributes to R. Brauer the stronger statement about minimal splitting fields of a quaternion division algebra of every even degree. The stronger result must not be attributed to the preceding argument, nor should a minimal field be confused with one of numerically least degree.\par
\noindent{\footnotesize Izvorne zapisy / Original records: HIST-g172-004.}\par
\smallskip\noindent\emph{Sravnjenje s tekstom:} Glavny tekst raboty 32 govori o potencah indeksa kak stepenah minimalnyh razpadnyh polj. Primětka potom pripisuje Braueru rezultat za vsaky parny stepen v kvaternionnom slučaju. Sohranjenje granicy medžu tekstom i primětkoju sohranja i različny obhvat tvrdženij.\par
\noindent\emph{Textual comparison:} The main text of Paper 32 discusses powers of the index as degrees of minimal splitting fields. The footnote then attributes to Brauer the result for every even degree in the quaternion case. Preserving the distinction between text and footnote also preserves the different scope of the assertions.\par
\noindent{\footnotesize TEXT-COMP-136; pasaže / passages: 19490; \texttt{CORPUS\_FORM\_COMPARISONS.json}.}\par
\end{minipage}\par
\par\medskip\Needspace{7\baselineskip}
\noindent\begin{minipage}{\linewidth}
\hypertarget{TERM-0967}{}
\noindent\textbf{teorem Jordan--Hölder}\quad{\footnotesize\texttt{TERM-0967}}\par
\noindent\emph{Iztočnik / Source:} Jordan-Hölderscher Satz\par
\noindent\emph{Rusky / Russian:} теорема Jordan--Hölder\par
\noindent\emph{Ukrajinsky / Ukrainian:} теорема Jordan--Hölder\par
\smallskip\noindent\emph{Razlog:} Teorem Jordan--Hölder sohranja medžunarodno razpoznavajemo ime v rabotě 30, §10. Vsě tri zapisane formy ostavjajut imena v latinici. Motivacija takože govori, že imenny izraz jest ochranjen pri avtomatičnom prěpisu okolnogo teksta v kirilicu. To jest pravilo za vlastna imena, ne prepověd kiriličnogo zapisa matematičnoj prozy. Zapis pravila sam po sobě ne proverjaje každogo kiriličnogo pojavjenja. Nazvanje teoremy identifikuje rezultat; jego točny obhvat ostaje zavisny od formulacije v tekstě.\par
\noindent\emph{Reason:} Teorem Jordan--Hölder preserves the internationally recognizable name in Paper 30, §10. All three recorded forms retain the names in Latin script. The rationale also says the name is protected during automatic conversion of surrounding prose to Cyrillic. This is a proper-name convention, not a prohibition on Cyrillic mathematical prose. The recorded rule alone does not verify every Cyrillic occurrence. The theorem's name identifies the result; its exact scope still depends on the statement in the text.\par
\noindent{\footnotesize Izvorne zapisy / Original records: HIST-g117-010.}\par
\end{minipage}\par
\par\medskip\Needspace{7\baselineskip}
\noindent\begin{minipage}{\linewidth}
\hypertarget{TERM-0969}{}
\noindent\textbf{pravy $K$-modul}\quad{\footnotesize\texttt{TERM-0969}}\par
\noindent\emph{Iztočnik / Source:} $K$-Rechtsmodul\par
\noindent\emph{Rusky / Russian:} правый $K$-модуль\par
\noindent\emph{Ukrajinsky / Ukrainian:} правий $K$-модуль\par
\smallskip\noindent\emph{Razlog:} Pravy $K$-modul sohranja storonu skalarnogo dějstva v rabotě 34, §6. Ona jest matematično sučstvena, zato že $K$ ne mora byti komutativno: zapis $aK$ kaže skalary sprava. Odlomok takože izrično trěbuje, že jedinica $e$ těla dějstvuje kak jediničny operator, $ae=a$. Pravost ne jest izbranym grafičnym poredkom slov, ktory možno bezuslovno obratiti. Ruski i ukrajinski sravnjajuči vidy sohranjajut to samo ograničenje.\par
\noindent\emph{Reason:} Pravy $K$-modul preserves the side of scalar action in Paper 34, §6. This is mathematically substantive because $K$ need not be commutative: $aK$ places scalars on the right. The passage also expressly requires the division ring's identity $e$ to act as the identity operator, $ae=a$. Right-sidedness is not an arbitrary typographic arrangement that may be reversed unconditionally. The Russian and Ukrainian comparators preserve the same restriction.\par
\noindent{\footnotesize Izvorne zapisy / Original records: HIST-g184-003.}\par
\smallskip\noindent\emph{Sravnjenje s tekstom:} V rabotě 31 $K$ razgledaje se kak $\Omega$-modul pri oprěděljenju slěda. V rabotě 34 tělo ne mora byti komutativno, a pravy $K$-modul imaje izrično uslovje $ae=a$. Fraza razgledati kak i pridevnik pravy nosat različne, oba važné strukturne informacije.\par
\noindent\emph{Textual comparison:} In Paper 31, $K$ is regarded as an $\Omega$-module when defining trace. In Paper 34 the division ring need not be commutative, and the right $K$-module has the explicit condition $ae=a$. Regarded as and right carry different, equally important structural information.\par
\noindent{\footnotesize TEXT-COMP-137; pasaže / passages: 18857, 20221, 20223; \texttt{CORPUS\_FORM\_COMPARISONS.json}.}\par
\end{minipage}\par
\par\medskip\Needspace{7\baselineskip}
\noindent\begin{minipage}{\linewidth}
\hypertarget{TERM-0970}{}
\noindent\textbf{t?lo; komutativno polje}\quad{\footnotesize\texttt{TERM-0970}}\par
\noindent\emph{Iztočnik / Source:} K?rper; kommutativer K?rper\par
\noindent\emph{Rusky / Russian:} [Bukvy ne sut zachovane / Lettering not recoverable]\par
\noindent\emph{Ukrajinsky / Ukrainian:} [Bukvy ne sut zachovane / Lettering not recoverable]\par
\smallskip\noindent\emph{Razlog:} Poškodženy zapis t?lo; komutativno polje naměrno razlikuje dva smysly německogo Körper v rabotě 40. Tělo dolžno označati vozmožno nekomutativny dělitveny objekt, polje komutativny. Ruska i ukrajinska slovna sravnjenja v tom zapisě ne sut čitljive, zato ne možno iz njih napraviti tabelu realnyh alternativ. Sravnjeni odlomok o tělu $A$ s centrom $P$ i komutativnom razširitelnom polju $Z$ podpiraje potrěbu razlikovanja. Ne dokazuje, že vse pojavjenja Körper v cěloj rabotě už byli kontekstno proverjene.\par
\noindent\emph{Reason:} The damaged record t?lo; komutativno polje deliberately distinguishes two senses of German Körper in Paper 40. Tělo is intended for a possibly noncommutative division object, polje for a commutative field. The Russian and Ukrainian lexical comparators in this record are illegible, so they cannot supply a table of actual alternatives. The compared passage concerning a division object $A$ with centre $P$ and a commutative extension field $Z$ supports the need for the distinction. It does not prove that every occurrence of Körper throughout the paper has been contextually checked.\par
\noindent{\footnotesize Izvorne zapisy / Original records: HIST-g210-001.}\par
\smallskip\noindent\emph{Sravnjenje s tekstom:} Rabota 40 razlikuje dělitvene algebry s centrom $P$ od komutativnyh razširitelnyh polj. Klasy podobnyh algeber imajut umnoženje, jediničny klas i obratne klasy, izrično opisane v §7. To podpiraje smyslovu rekonstrukciju, ale ne vozvrača izgubjene bukvy ruskoj i ukrajinskoj kolonki.\par
\noindent\emph{Textual comparison:} Paper 40 distinguishes division algebras with centre $P$ from commutative extension fields. Classes of similar algebras have multiplication, an identity class and inverse classes, explicitly described in §7. This supports semantic reconstruction but does not restore the lost Russian and Ukrainian lettering.\par
\noindent{\footnotesize TEXT-COMP-141; pasaže / passages: 23895, 23900, 23902, 23905, 23907, 23913; \texttt{CORPUS\_FORM\_COMPARISONS.json}.}\par
\end{minipage}\par
\par\medskip\Needspace{7\baselineskip}
\noindent\begin{minipage}{\linewidth}
\hypertarget{TERM-0971}{}
\noindent\textbf{$K_i$ konjugovane polja v Galoisovom razširjenom polju $Q$}\quad{\footnotesize\texttt{TERM-0971}}\par
\noindent\emph{Iztočnik / Source:} $K_i$ konjugierte Körper im Galoisschen Erweiterungskörper $Q$\par
\noindent\emph{Rusky / Russian:} $K_i$ сопряженные поля в поле расширения Галуа $Q$\par
\noindent\emph{Ukrajinsky / Ukrainian:} $K_i$ спряжені поля в полі розширення Галуа $Q$\par
\smallskip\noindent\emph{Razlog:} $K_i$ konjugovane polja v Galoisovom razširjenom polju $Q$ sohranja i odnos konjugacije, i obče obkladajuče polje v rabotě 31. Motivacija preferuje razpoznavajemy medžunarodny koren konjug- i izrično nazvanje Galoisovogo razširjenja. Konjugovane polja ne sut avtomatično bukvalno to samo podpolje; ni izomorfija sama po sobě ne ustanavjaje jednakost podmnožin. Ocěna, že sočetanje jest dostatočno razumljivo medžu slovjanskymi jezykami, jest zachovany jezykovy sud, ne rezultat čitateljskogo izpytanja.\par
\noindent\emph{Reason:} $K_i$ konjugovane polja v Galoisovom razširjenom polju $Q$ preserves both conjugacy and the common ambient field in Paper 31. The rationale favours the recognizable international conjug- stem and explicit naming of the Galois extension. Conjugate fields are not automatically literally the same subfield; isomorphism alone likewise does not establish equality as subsets. The assessment that the combination is sufficiently transparent across Slavic languages is a recorded linguistic judgment, not the result of a reader experiment.\par
\noindent{\footnotesize Izvorne zapisy / Original records: HIST-g147-006.}\par
\end{minipage}\par
\par\medskip\Needspace{7\baselineskip}
\noindent\begin{minipage}{\linewidth}
\hypertarget{TERM-0972}{}
\noindent\textbf{$k_i$-rędne determinanty}\quad{\footnotesize\texttt{TERM-0972}}\par
\noindent\emph{Iztočnik / Source:} $k_i$-reihige Determinanten\par
\noindent\emph{Rusky / Russian:} $k_i$-рядные детерминанты\par
\noindent\emph{Ukrajinsky / Ukrainian:} $k_i$-рядні детермінанти\par
\smallskip\noindent\emph{Razlog:} $k_i$-rędne determinanty sohranja indeksovany poręd determinanta v rabotě 22, §6, i prodolžuje terminologiju minorov. $i$ jest indeks, ne faktor, a $k_i$ ne naziva čislo determinantov. Sravnjeni tekst upotrěbjaje izričnu konstrukciju determinanty reda $k_i$. Oba načina predstavjenja sohranjajut parametr, ale složeny pridevnik i genitivna fraza sut različne jezykove odluki. Matematičnu ekvivalentnost dvaju sistemov determinantov dava argument o modulah, ne samo pravilno prěpisany indeks.\par
\noindent\emph{Reason:} $k_i$-rędne determinanty retains the indexed order of the determinant in Paper 22, §6 and continues the terminology of minors. The $i$ is an index, not a factor, and $k_i$ does not count the determinants. The compared text uses the explicit construction determinanty reda $k_i$. Both presentations preserve the parameter, but the compound adjective and genitive phrase are different linguistic choices. The module argument, not merely a correctly copied index, establishes the equivalence of the two determinant systems.\par
\noindent{\footnotesize Izvorne zapisy / Original records: HIST-g090-008, HIST-g092-062.}\par
\smallskip\noindent\emph{Sravnjenje s tekstom:} Sravnjeni odlomok imaje determinanty reda $k_i$ i reda $\bar k_i$, povezane argumentom o modulah linearnyh form. To ustanavjaje matematičnu rolju indeksa i nynějšu opisnu frazu pored arhivnogo $k_i$-rędne.\par
\noindent\emph{Textual comparison:} The compared passage has determinants of order $k_i$ and order $\bar k_i$, related by an argument about modules of linear forms. This establishes the mathematical role of the index and the present descriptive phrase beside archival $k_i$-rędne.\par
\noindent{\footnotesize TEXT-COMP-142; pasaže / passages: 14945; \texttt{CORPUS\_FORM\_COMPARISONS.json}.}\par
\end{minipage}\par
\par\medskip\Needspace{7\baselineskip}
\noindent\begin{minipage}{\linewidth}
\hypertarget{TERM-0975}{}
\noindent\textbf{cepi prostyh idealov / cepova definicija}\quad{\footnotesize\texttt{TERM-0975}}\par
\noindent\emph{Iztočnik / Source:} Ketten von Primidealen / Kettendefinition\par
\noindent\emph{Rusky / Russian:} цепочки простых идеалов / цепочное определение\par
\noindent\emph{Ukrajinsky / Ukrainian:} ланцюги простих ідеалів / ланцюгове означення\par
\smallskip\noindent\emph{Razlog:} Cepi prostyh idealov / cepova definicija prěnosi terminologiju konečnyh cepij iz raboty 24, §1, 10, na definiciju dimenzije. Cepova definicija jest prozračno sočetanje, ne nova neoprěděljena etiketa. Odlomok računa člany cepi i izrično vključaje jediničny ideal; pri avtorovoj konvenciji njegova dimenzija jest $-1$. Zato nazvanje ne opravyvaje tiho zaměniti ni definiciju prostogo ideala, ni način računanja členov sovrěmennoju verzijeju. Sravnjeni tekst takože govori, kogda možno obojti se bez adjunkcije novyh neoprěděljenyh.\par
\noindent\emph{Reason:} Cepi prostyh idealov / cepova definicija extends the finite-chain vocabulary of Paper 24, §1, item 10 to the dimension definition. Cepova definicija is a transparent phrase, not a new unexplained label. The passage counts members of a chain and explicitly includes the unit ideal; under the author's convention its dimension is $-1$. The name therefore does not justify silently replacing either the definition of prime ideal or the counting convention with a modern version. The compared text also states when adjoining new indeterminates can be avoided.\par
\noindent{\footnotesize Izvorne zapisy / Original records: HIST-g098-006.}\par
\smallskip\noindent\emph{Sravnjenje s tekstom:} V rabotě 24 definicija dimenzije računa člany cepi s izrično vključonym jediničnym idealom dimenzije $-1$ i s uslovjem o adjunkciji neoprěděljenyh. V rabotě 30 naslov govori o uslovjah, a uvod o teoremě, ktora jih prěnosi na moduly. Sravnjenje sohranja i avtorovu konvenciju računanja, i razlikovanje svojstva od teoremy.\par
\noindent\emph{Textual comparison:} Paper 24's dimension definition counts chain members with the unit ideal explicitly included at dimension $-1$, and with a qualification concerning adjoining indeterminates. Paper 30's title names conditions, while its opening names the theorem transferring them to modules. The comparison preserves both the author's counting convention and the distinction between a property and a theorem.\par
\noindent{\footnotesize TEXT-COMP-138; pasaže / passages: 15531, 15535, 15537, 16305, 16307; \texttt{CORPUS\_FORM\_COMPARISONS.json}.}\par
\end{minipage}\par
\par\medskip\Needspace{7\baselineskip}
\noindent\begin{minipage}{\linewidth}
\hypertarget{TERM-0976}{}
\noindent\textbf{uslovja cěpov v konečnyh modulnyh oblastjah}\quad{\footnotesize\texttt{TERM-0976}}\par
\noindent\emph{Iztočnik / Source:} Kettensätze in endlichen Modulbereichen\par
\noindent\emph{Rusky / Russian:} условия цепей в конечных модульных областях\par
\noindent\emph{Ukrajinsky / Ukrainian:} умови ланцюгів у скінченних модульних областях\par
\smallskip\noindent\emph{Razlog:} Uslovja cěpov v konečnyh modulnyh oblastjah v rabotě 30, §2, prodolžuje ranějši izbor prěvoditi Kettensatz jezykom uslovij. Cěp jest zachovany koren za konečny red svazanyh objektov, modulna oblast sohranja avtorov širi pojam. Uvod paragrafa govori o teoremě, ktora prěnosi cěpove uslovja na moduly. Zato uslovje i teorema ne sut sinonimy: prěvod naslovu naglašaje svojstvo, a izloženje daje rezultat o jeho prěnosu. Ta razlika dolžna byti jasna v metodologičnom obosnovanju.\par
\noindent\emph{Reason:} Uslovja cěpov v konečnyh modulnyh oblastjah in Paper 30, §2 continues the earlier decision to render Kettensatz in the language of conditions. Cěp supplies the established chain stem, while modulna oblast retains the author's broader notion. The section's opening describes a theorem transferring chain conditions to modules. Condition and theorem are therefore not synonyms: the translated heading emphasizes the property, while the exposition supplies the result about transferring it. That difference belongs in the methodological explanation.\par
\noindent{\footnotesize Izvorne zapisy / Original records: HIST-g109-001.}\par
\smallskip\noindent\emph{Sravnjenje s tekstom:} V rabotě 24 definicija dimenzije računa člany cepi s izrično vključonym jediničnym idealom dimenzije $-1$ i s uslovjem o adjunkciji neoprěděljenyh. V rabotě 30 naslov govori o uslovjah, a uvod o teoremě, ktora jih prěnosi na moduly. Sravnjenje sohranja i avtorovu konvenciju računanja, i razlikovanje svojstva od teoremy.\par
\noindent\emph{Textual comparison:} Paper 24's dimension definition counts chain members with the unit ideal explicitly included at dimension $-1$, and with a qualification concerning adjoining indeterminates. Paper 30's title names conditions, while its opening names the theorem transferring them to modules. The comparison preserves both the author's counting convention and the distinction between a property and a theorem.\par
\noindent{\footnotesize TEXT-COMP-138; pasaže / passages: 15531, 15535, 15537, 16305, 16307; \texttt{CORPUS\_FORM\_COMPARISONS.json}.}\par
\end{minipage}\par
\par\medskip\Needspace{7\baselineskip}
\noindent\begin{minipage}{\linewidth}
\hypertarget{TERM-0977}{}
\noindent\textbf{skobkovy faktor}\quad{\footnotesize\texttt{TERM-0977}}\par
\noindent\emph{Iztočnik / Source:} Klammerfaktor\par
\noindent\emph{Rusky / Russian:} скобочный множитель\par
\noindent\emph{Ukrajinsky / Ukrainian:} дужковий множник\par
\smallskip\noindent\emph{Razlog:} Skobkovy faktor v rabotě 2, §5, jest doslovna fraza svazana s notacijeju. Formula $(abc)$ dava čitatelju konkretno značenje skobkovogo zapisa; faktor sohranja algebraičnu rolju. Ne ide prosto o každý izraz, zapisany v skobkah, ni o sam znak skobky kak samostojny množitelj. Motivacija sočetaje razpoznavajemy medžunarodny imenik s opisnym slovjanskim pridevnikom; jeje dostatočnost zavisi od togo, že formula ostaje pored objasnjenja.\par
\noindent\emph{Reason:} Skobkovy faktor in Paper 2, §5 is a literal notation-facing phrase. The formula $(abc)$ gives the bracket notation a concrete referent; faktor preserves its algebraic role. It does not simply mean every expression in parentheses, or the bracket character itself as an independent multiplier. The rationale combines a recognizable international noun with a descriptive Slavic adjective; its adequacy depends on keeping the formula alongside the explanation.\par
\noindent{\footnotesize Izvorne zapisy / Original records: HIST-g006-002.}\par
\end{minipage}\par
\par\medskip\Needspace{7\baselineskip}
\noindent\begin{minipage}{\linewidth}
\hypertarget{TERM-0978}{}
\noindent\textbf{klasa grupovogo elementa}\quad{\footnotesize\texttt{TERM-0978}}\par
\noindent\emph{Iztočnik / Source:} Klasse eines Gruppenelements\par
\noindent\emph{Rusky / Russian:} класс элемента группы\par
\noindent\emph{Ukrajinsky / Ukrainian:} клас елемента групи\par
\smallskip\noindent\emph{Razlog:} Klasa grupovogo elementa izbira kratky pridevničny izraz vměsto genitivnoj konstrukcije ruskogo i ukrajinskogo sravnjenja. Odlomok raboty 34, §26, dava sučstveno utočnjenje: klasoju nazvana jest suma $K_a=\sum_s s^{-1}as$ v grupovom prstenu, pri čem sumuje se samo po različnyh konjugatnyh elementah. Zato čitatelj ne směje razuměti frazu samo kak množinu konjugatov, ostavjajuči sumu nevyrečenoju. Klasy v tom smyslu sut centralne elementy prstena. Izbor sintaksy i pojasnjenje lokalnogo značenja sut različne, ale obě potrěbne česti prevoda.\par
\noindent\emph{Reason:} Klasa grupovogo elementa selects a compact adjectival phrase instead of the genitive constructions in the Russian and Ukrainian comparisons. Paper 34, §26 adds an essential qualification: what it calls a class is the group-ring sum $K_a=\sum_s s^{-1}as$, with summation only over distinct conjugates. Readers must therefore not understand the phrase merely as a set of conjugates while leaving the sum unexpressed. Classes in this sense are central ring elements. Choosing the syntax and explaining the local sense are distinct but both necessary tasks.\par
\noindent{\footnotesize Izvorne zapisy / Original records: HIST-g204-006.}\par
\smallskip\noindent\emph{Sravnjenje s tekstom:} V §26 raboty 34 klasa jest izrično suma različnyh elementov $s^{-1}as$ v grupovom prstenu. Daljše tvrdženje, že te klasy sut centralne elementy, zavisi od takogo ponimanja. Glosa samo množina konjugatov ostavila by matematičnu operaciju nevyrečenoju.\par
\noindent\emph{Textual comparison:} In Paper 34, §26, class is explicitly a sum of the distinct elements $s^{-1}as$ in the group ring. The subsequent statement that these classes are central elements depends on that reading. The gloss set of conjugates alone would leave a mathematical operation unexpressed.\par
\noindent{\footnotesize TEXT-COMP-139; pasaže / passages: 22475, 22477, 22479; \texttt{CORPUS\_FORM\_COMPARISONS.json}.}\par
\end{minipage}\par
\par\medskip\Needspace{7\baselineskip}
\noindent\begin{minipage}{\linewidth}
\hypertarget{TERM-0979}{}
\noindent\textbf{klasa $\mathfrak{R}_i$ / podmnožstvo $\mathfrak{R}_i^*$}\quad{\footnotesize\texttt{TERM-0979}}\par
\noindent\emph{Iztočnik / Source:} Klasse $\mathfrak{R}_i$ / $\mathfrak{R}_i^*$\par
\noindent\emph{Rusky / Russian:} класс $\mathfrak{R}_i$ / подмножество $\mathfrak{R}_i^*$\par
\noindent\emph{Ukrajinsky / Ukrainian:} клас $\mathfrak{R}_i$ / підмножина $\mathfrak{R}_i^*$\par
\smallskip\noindent\emph{Razlog:} Klasa $\mathfrak R_i$ / podmnožstvo $\mathfrak R_i^*$ v rabotě 9, §8, razlikuje zbirku s vozmožnymi povtorjenjami od izbora bez povtorjenja. Iz oblasti $\mathfrak G_i$ stroji se klasa vseh izomorfnyh k nej oblastij, pri čem ista oblast može pojavjati se mnogo razov. Podmnožstvo so zvezdoju sodrživaje vsaku takovu oblast jedin raz. To ne znači izbor samo jednogo prědstavitelja za cěly izomorfny klas: ostavajut vse različne oblasti, izomorfne k $\mathfrak G_i$. Arhivni latinski R i G predstavjajut gotične označenja tekstovogo odlomka.\par
\noindent\emph{Reason:} Klasa $\mathfrak R_i$ / podmnožstvo $\mathfrak R_i^*$ in Paper 9, §8 distinguishes a collection with possible repetition from a selection without repetition. Starting from $\mathfrak G_i$, the construction gives a class of all domains isomorphic to it, in which the same domain may occur many times. The starred subset contains each such domain once. This does not mean choosing only one representative for the entire isomorphism class: all distinct domains isomorphic to $\mathfrak G_i$ remain. Archival Latin R and G represent the passage's fraktur identifiers.\par
\noindent{\footnotesize Izvorne zapisy / Original records: HIST-g066-010.}\par
\smallskip\noindent\emph{Sravnjenje s tekstom:} $\mathfrak R_i$ može sodrživati istu oblast mnogo razov; $\mathfrak R_i^*$ sodrživaje vsaku oblast izomorfnu k $\mathfrak G_i$ jedin raz. Ne ide o izbor jednogo jedinogo objekta vměsto vseh izomorfnyh objektov. Razlikovanje klasa i podmnožstva sohranja imenno točnu izbirovu operaciju.\par
\noindent\emph{Textual comparison:} $\mathfrak R_i$ may contain the same domain many times; $\mathfrak R_i^*$ contains each domain isomorphic to $\mathfrak G_i$ once. This is not selection of one single object in place of all isomorphic objects. The distinction between class and subset preserves the particular selection operation.\par
\noindent{\footnotesize TEXT-COMP-140; pasaže / passages: 8122, 8123, 8124, 8125, 8126, 8127, 8128; \texttt{CORPUS\_FORM\_COMPARISONS.json}.}\par
\end{minipage}\par
\par\medskip\Needspace{7\baselineskip}
\noindent\begin{minipage}{\linewidth}
\hypertarget{TERM-0980}{}
\noindent\textbf{klasy podobnyh algeber; grupa klasov}\quad{\footnotesize\texttt{TERM-0980}}\par
\noindent\emph{Iztočnik / Source:} Klassen ?hnlicher Algebren; Gruppe der Klassen\par
\noindent\emph{Rusky / Russian:} [Bukvy ne sut zachovane / Lettering not recoverable]\par
\noindent\emph{Ukrajinsky / Ukrainian:} [Bukvy ne sut zachovane / Lettering not recoverable]\par
\smallskip\noindent\emph{Razlog:} Klasy podobnyh algeber; grupa klasov sohranja opisny jezyk raboty 40 vměsto tihoj zaměny pozdnějšim eponymnym nazvanjem. V tekstě jeden klas sobiraje matrične algebry s tym samym pridruženym dělitvenym objektom, s izomorfnymi objektami identificirovanymi. Produkt nad obščim centrom definuje umnoženje klasov; jediničny klas davajut matrične prsteny nad centrom, obratny klas recipročna algebra. Tak podprta jest vez s pozdnějšim jezykom Brauerovoj grupy, ale podobny ne směje byti pročitan kak libovoljna matrična podobnost. Ruska i ukrajinska polja v zachovanom zapisě poškodžena sut; njih leksika ne rekonstruuje se iz paměti.\par
\noindent\emph{Reason:} Klasy podobnyh algeber; grupa klasov retains Paper 40's descriptive language instead of silently replacing it with a later eponym. In the text, a class collects matrix algebras with the same associated division object, identifying isomorphic division objects. Product over the common centre defines multiplication of classes; matrix rings over the centre give the identity class, and the reciprocal algebra gives the inverse. This supports the connection with later Brauer-group language, but similar must not be read as arbitrary matrix similarity. The Russian and Ukrainian fields in the surviving record are damaged; their wording is not reconstructed from memory.\par
\noindent{\footnotesize Izvorne zapisy / Original records: HIST-g210-009.}\par
\smallskip\noindent\emph{Sravnjenje s tekstom:} Rabota 40 razlikuje dělitvene algebry s centrom $P$ od komutativnyh razširitelnyh polj. Klasy podobnyh algeber imajut umnoženje, jediničny klas i obratne klasy, izrično opisane v §7. To podpiraje smyslovu rekonstrukciju, ale ne vozvrača izgubjene bukvy ruskoj i ukrajinskoj kolonki.\par
\noindent\emph{Textual comparison:} Paper 40 distinguishes division algebras with centre $P$ from commutative extension fields. Classes of similar algebras have multiplication, an identity class and inverse classes, explicitly described in §7. This supports semantic reconstruction but does not restore the lost Russian and Ukrainian lettering.\par
\noindent{\footnotesize TEXT-COMP-141; pasaže / passages: 23895, 23900, 23902, 23905, 23907, 23913; \texttt{CORPUS\_FORM\_COMPARISONS.json}.}\par
\end{minipage}\par
\par\medskip\Needspace{7\baselineskip}
\noindent\begin{minipage}{\linewidth}
\hypertarget{TERM-0981}{}
\noindent\textbf{klasovy invariant}\quad{\footnotesize\texttt{TERM-0981}}\par
\noindent\emph{Iztočnik / Source:} Klasseninvariante\par
\noindent\emph{Rusky / Russian:} классовый инвариант\par
\noindent\emph{Ukrajinsky / Ukrainian:} класовий інваріант\par
\smallskip\noindent\emph{Razlog:} Klasovy invariant v rabotě 31, §4, sojedinja medžunarodny invariant s pridevnikom od klasa. Rusky классовый инвариант i ukrajinsky класовий інваріант davajut blizke sravnjenje; motivacija predpočitaje kratku frazu vměsto težšego opisa invariant po odnošenju k klasu. Konkretno ide o invariant klasa prědstavjenij i, po predhodnom rezultatu, klasa idealov. Odlomok ustanavjaje invariantnost $|tE-C|$ i njegovih koeficientov. Ne ide o proizvoljnu klasu objektov. Tvrdženje o čitateljskoj prozračnosti jest jezykova ocěna, ne izměreny rezultat.\par
\noindent\emph{Reason:} Klasovy invariant in Paper 31, §4 combines the international noun with the adjective formed from klasa. The Russian and Ukrainian adjective–noun phrases provide close comparators; the rationale favours this compact construction over the heavier invariant with respect to a class. Here the class is a representation class and, by the preceding result, an ideal class: the passage establishes invariance of $|tE-C|$ and its coefficients. It does not mean an arbitrary class of objects. Claimed readability is a linguistic judgment, not a measured result.\par
\noindent{\footnotesize Izvorne zapisy / Original records: HIST-g135-003.}\par
\smallskip\noindent\emph{Sravnjenje s tekstom:} V rabotě 31 koeficienty $|tE-C|$ sut invarianty klasa. Slěd jest koeficient pri $(-t^{s-1})$ i ravny sumě diagonalnyh elementov. Taka definicija svazuje leksičny izbor s přesnym znakom i s konkretnym smyslom klasa.\par
\noindent\emph{Textual comparison:} In Paper 31 the coefficients of $|tE-C|$ are class invariants. The trace is the coefficient of $(-t^{s-1})$ and equals the sum of diagonal entries. The definition links the lexical choice to an exact sign and a particular meaning of class.\par
\noindent{\footnotesize TEXT-COMP-145; pasaže / passages: 18024, 18026, 18028; \texttt{CORPUS\_FORM\_COMPARISONS.json}.}\par
\end{minipage}\par
\par\medskip\Needspace{7\baselineskip}
\noindent\begin{minipage}{\linewidth}
\hypertarget{TERM-0982}{}
\noindent\textbf{klasno polje / prim-formy}\quad{\footnotesize\texttt{TERM-0982}}\par
\noindent\emph{Iztočnik / Source:} Klassenkörper / Primformen\par
\noindent\emph{Rusky / Russian:} поле классов / прайм-формы\par
\noindent\emph{Ukrajinsky / Ukrainian:} поле класів / прайм-форми\par
\smallskip\noindent\emph{Razlog:} Klasno polje / prim-formy čuvaje dva različna termina raboty 14. Za Primformen zapis izbira zaimstvovani koren prim-, sravnjujuči rusko прайм-формы i ukrajinsko прайм-форми, da by ne poměšal specialny pojem s obyčnoju prostoju formoju. V tekstě oblast transcendentnyh prim-form jest samo v někoj mery transcendentny analog klasnogo polja. Zato slovnik ne směje udělati iz analogije identičnost dvoh teorij. Zapis nazyvaje termine standardnymi, ale ne dava nezavisnu medžuslovjansku atestaciju.\par
\noindent\emph{Reason:} Klasno polje / prim-formy preserves two distinct terms in Paper 14. For Primformen the record chooses the borrowed prim- stem, compared with Russian прайм-формы and Ukrainian прайм-форми, to avoid confusing the specialized notion with an ordinary simple form. The text describes the domain of transcendental prime forms only as an analogue, in a certain respect, of a class field. The dictionary must not turn that analogy into an identity of theories. The record calls the terminology standard but supplies no independent Interslavic attestation.\par
\noindent{\footnotesize Izvorne zapisy / Original records: HIST-g076-014.}\par
\smallskip\noindent\emph{Sravnjenje s tekstom:} Oblast transcendentnyh prim-form opisana jest samo v někoj mery kak analog klasnogo polja. Upotrěba različnyh terminov ne trěbuje prěvratiti tu ograničenu analogiju v identičnost.\par
\noindent\emph{Textual comparison:} The domain of transcendental prime forms is described only in a certain respect as an analogue of a class field. The terminology should preserve this qualified analogy, not turn it into an identity.\par
\noindent{\footnotesize TEXT-COMP-146; pasaže / passages: 9922; \texttt{CORPUS\_FORM\_COMPARISONS.json}.}\par
\end{minipage}\par
\par\medskip\Needspace{7\baselineskip}
\noindent\begin{minipage}{\linewidth}
\hypertarget{TERM-0983}{}
\noindent\textbf{najmenša soderžeča integralna oblast}\quad{\footnotesize\texttt{TERM-0983}}\par
\noindent\emph{Iztočnik / Source:} kleinster enthaltender Integritätsbereich\par
\noindent\emph{Rusky / Russian:} наименьшая содержащая область целостности\par
\noindent\emph{Ukrajinsky / Ukrainian:} найменша інтегральна область, що містить систему\par
\smallskip\noindent\emph{Razlog:} Najmenša soderžeča integralna oblast naziva objekt potrěbny pri obobčenju poslě Teoremy VIII v rabotě 6, §12. Najmenša nosi uslovje vključenja, ne ocěnu čislovoj veličiny. Sravnjeni tekst govori o polinomah najmenšej oblasti, ktore možut byti prědstavjene cělymi racionalnymi kombinacijami konečno mnogyh polinomov početnogo sistema. Čislo tyh polinomov može zavisěti od izbranogo elementa. To jest sučstveno za prěnos argumenta o bazi; zaměna samym slovom zamknutost skryla by, kakoj objekt jest skonstruovany.\par
\noindent\emph{Reason:} Najmenša soderžeča integralna oblast names the object needed in the generalization following Theorem VIII in Paper 6, §12. Least expresses an inclusion condition, not numerical size. The compared text describes polynomials in the least domain as integral rational combinations of finitely many polynomials from the initial system; the number may depend on the element represented. This matters to the transfer of the basis argument. Replacing the phrase by closure alone would conceal which object has been constructed.\par
\noindent{\footnotesize Izvorne zapisy / Original records: HIST-g052-005.}\par
\smallskip\noindent\emph{Sravnjenje s tekstom:} Rabota 6 razlikuje cěle racionalne kombinacije od kvocientov i vsih racionalnyh kombinacij. Poslě Teoremy VIII jasno jest, že element najmenšej oblasti koristi konečno čislo polinomov iz početnogo sistema, zavisno od elementa. Zato sodržanje i minimalnost ne sut samo ukrasne pridevniky.\par
\noindent\emph{Textual comparison:} Paper 6 distinguishes integral rational combinations from quotients and arbitrary rational combinations. Following Theorem VIII, an element of the least domain uses finitely many initial-system polynomials, with the number depending on that element. Containment and leastness are therefore substantive qualifications.\par
\noindent{\footnotesize TEXT-COMP-147; pasaže / passages: 5770, 5772, 5774, 5778, 5780, 6271, 6273; \texttt{CORPUS\_FORM\_COMPARISONS.json}.}\par
\end{minipage}\par
\par\medskip\Needspace{7\baselineskip}
\noindent\begin{minipage}{\linewidth}
\hypertarget{TERM-0984}{}
\noindent\textbf{najmenjše polje, ktore soderži sistem}\quad{\footnotesize\texttt{TERM-0984}}\par
\noindent\emph{Iztočnik / Source:} kleinster enthaltender Körper\par
\noindent\emph{Rusky / Russian:} наименьшее поле, содержащее систему\par
\noindent\emph{Ukrajinsky / Ukrainian:} найменше поле, що містить систему\par
\smallskip\noindent\emph{Razlog:} Najmenjše polje, ktore soderži sistem v uvodě raboty 6 izrično imenuje i polje, i sodržany sistem. Motivacija otklanja prěkratky izraz za zamknutost bez nazvanja polja i sodržanogo sistema. Odlomok §7 pozdněje definira to polje prěz vse racionalne kombinacije konečno mnogyh funkcij sistema s koeficientami iz $\Omega$. Zato najmenjše ne znači najmanje po čislu elementov. To relativno prědloženje i krače pridevnične sočetanje v slědujučem zapisu sut dva načina prědati isto uslovje, ne dva nova vida polj.\par
\noindent\emph{Reason:} Najmenjše polje, ktore soderži sistem in the introduction to Paper 6 names both the field and the system it contains. The rationale rejects over-compression to a bare hull or closure. Section 7 subsequently defines the field through all rational combinations of finitely many system functions with coefficients in $\Omega$. Least therefore does not mean smallest cardinality. This relative clause and the shorter adjectival construction in the next record are two ways of expressing the containing condition, not two new kinds of field.\par
\noindent{\footnotesize Izvorne zapisy / Original records: HIST-g040-016.}\par
\smallskip\noindent\emph{Sravnjenje s tekstom:} Rabota 6 razlikuje cěle racionalne kombinacije od kvocientov i vsih racionalnyh kombinacij. Poslě Teoremy VIII jasno jest, že element najmenšej oblasti koristi konečno čislo polinomov iz početnogo sistema, zavisno od elementa. Zato sodržanje i minimalnost ne sut samo ukrasne pridevniky.\par
\noindent\emph{Textual comparison:} Paper 6 distinguishes integral rational combinations from quotients and arbitrary rational combinations. Following Theorem VIII, an element of the least domain uses finitely many initial-system polynomials, with the number depending on that element. Containment and leastness are therefore substantive qualifications.\par
\noindent{\footnotesize TEXT-COMP-147; pasaže / passages: 5770, 5772, 5774, 5778, 5780, 6271, 6273; \texttt{CORPUS\_FORM\_COMPARISONS.json}.}\par
\end{minipage}\par
\par\medskip\Needspace{7\baselineskip}
\noindent\begin{minipage}{\linewidth}
\hypertarget{TERM-0985}{}
\noindent\textbf{najmenše soderžeče polje}\quad{\footnotesize\texttt{TERM-0985}}\par
\noindent\emph{Iztočnik / Source:} kleinster enthaltender Körper\par
\noindent\emph{Rusky / Russian:} наименьшее содержащее поле\par
\noindent\emph{Ukrajinsky / Ukrainian:} найменше поле, що містить\par
\smallskip\noindent\emph{Razlog:} Najmenše soderžeče polje jest kratka forma konstrukcije iz raboty 6, §7. Čitatelj mora uviděti sodržanje sistema i minimalnost polja, zato že argument prěnosi racionalne bazy polj na libovoljne sistemy. Sravnjeni tekst razlikuje linearne kombinacije, cěle racionalne kombinacije i kvocienty, prěhodeči poslědovateljno k linearnoj familiji, integralnoj oblasti i polju. Najmenše polje ne jest samo drugoje ime za jednu iz predhodnyh oblastij. Arhivno soderžeče i nynešnje sodrživajuče ostajut različnymi formami v dokumentaciji.\par
\noindent\emph{Reason:} Najmenše soderžeče polje is the compact construction in Paper 6, §7. Both containment and leastness must remain visible because the argument transfers rational bases of fields to arbitrary systems. The compared text distinguishes linear combinations, integral rational combinations and quotients, passing successively to a linear family, a domain and a field. The least field is not merely another name for either preceding object. Archival soderžeče and present sodrživajuče remain distinct documented forms.\par
\noindent{\footnotesize Izvorne zapisy / Original records: HIST-g047-010.}\par
\smallskip\noindent\emph{Sravnjenje s tekstom:} Rabota 6 razlikuje cěle racionalne kombinacije od kvocientov i vsih racionalnyh kombinacij. Poslě Teoremy VIII jasno jest, že element najmenšej oblasti koristi konečno čislo polinomov iz početnogo sistema, zavisno od elementa. Zato sodržanje i minimalnost ne sut samo ukrasne pridevniky.\par
\noindent\emph{Textual comparison:} Paper 6 distinguishes integral rational combinations from quotients and arbitrary rational combinations. Following Theorem VIII, an element of the least domain uses finitely many initial-system polynomials, with the number depending on that element. Containment and leastness are therefore substantive qualifications.\par
\noindent{\footnotesize TEXT-COMP-147; pasaže / passages: 5770, 5772, 5774, 5778, 5780, 6271, 6273; \texttt{CORPUS\_FORM\_COMPARISONS.json}.}\par
\end{minipage}\par
\par\medskip\Needspace{7\baselineskip}
\noindent\begin{minipage}{\linewidth}
\hypertarget{TERM-0986}{}
\noindent\textbf{najmenša soderžeča relativno cěla oblast}\quad{\footnotesize\texttt{TERM-0986}}\par
\noindent\emph{Iztočnik / Source:} kleinster enthaltender relativ-ganzer Bereich\par
\noindent\emph{Rusky / Russian:} наименьшая содержащая относительно целая область\par
\noindent\emph{Ukrajinsky / Ukrainian:} найменша відносно ціла область, що містить\par
\smallskip\noindent\emph{Razlog:} Najmenša soderžeča relativno cěla oblast v rabotě 6, §9, naměrno povtarjaje konstrukciju najmenšego sodrživajučego polja. Analogia ne odstranjaje ograničenje na polinomy: nova oblast sodrživaje vse polinomy, ktore možut byti kvocientami dvoh polinomov početnoj oblasti. Ne vključaje proizvoljne racionalne funkcije bez togo uslovja. Zato relativno cěla jest tehnična kvalifikacija v lokalnoj definiciji; ne trěba jeje bez objasnjenja zaměniti sovrěmennym nazvanjem integralnogo zamknutja.\par
\noindent\emph{Reason:} Najmenša soderžeča relativno cěla oblast in Paper 6, §9 deliberately parallels the least-containing-field construction. The analogy does not remove the restriction to polynomials: the new domain contains all polynomials expressible as quotients of two polynomials from the initial domain. It does not admit arbitrary rational functions without that condition. Relativno cěla is therefore a technical qualification under the local definition, not a licence to substitute the modern label integral closure without explanation.\par
\noindent{\footnotesize Izvorne zapisy / Original records: HIST-g049-014.}\par
\smallskip\noindent\emph{Sravnjenje s tekstom:} Definicija dopušča kvocienty polinomov samo togda, kogda rezultat jest polinom. Predhodna i nova oblast imajut isto najmenše sodrživajuče polje; iz togo ne slěduje ravnost samyh oblastij.\par
\noindent\emph{Textual comparison:} The definition admits quotients of polynomials only when the result is a polynomial. The initial and new domains have the same least containing field; this does not make the domains themselves equal.\par
\noindent{\footnotesize TEXT-COMP-148; pasaže / passages: 6018, 6020, 6024; \texttt{CORPUS\_FORM\_COMPARISONS.json}.}\par
\end{minipage}\par
\par\medskip\Needspace{7\baselineskip}
\noindent\begin{minipage}{\linewidth}
\hypertarget{TERM-0987}{}
\noindent\textbf{najmenše razširjeno polje}\quad{\footnotesize\texttt{TERM-0987}}\par
\noindent\emph{Iztočnik / Source:} kleinster Erweiterungskörper\par
\noindent\emph{Rusky / Russian:} наименьшее поле расширения\par
\noindent\emph{Ukrajinsky / Ukrainian:} найменше поле розширення\par
\smallskip\noindent\emph{Razlog:} Najmenše razširjeno polje sohranja minimalnost $Q$ v rabotě 31, §6, i prodolžuje uže izbranu frazu razširjeno polje. V sravnjenom odlomku to jest najmenše polje, nad ktorym prsten dopušča ukazano prědstavjenje kako direktna suma prstenov ranga jedan, izomorfnyh k $Q$. Ne ide o libovoljno razširjenje ni o polje izbrano samo po nejnižšemu čislovomu stepenu. Razlog o stabilnosti frazy dokazuje kontinuitet lokalnogo izbora, ne samostojnu jezykovu normu.\par
\noindent\emph{Reason:} Najmenše razširjeno polje keeps the minimality of $Q$ explicit in Paper 31, §6 and continues the previously chosen extension-field phrase. In the compared passage it is the least field over which the stated decomposition into rank-one rings, each isomorphic to $Q$, becomes possible. It is not an arbitrary extension or a field selected only by numerically smallest degree. The record’s claim of a stable phrase establishes continuity of the project choice, not an independent linguistic norm.\par
\noindent{\footnotesize Izvorne zapisy / Original records: HIST-g146-006.}\par
\smallskip\noindent\emph{Sravnjenje s tekstom:} Uslovje na $Q$ jest vezano s razloženjem prstena na komponenty ranga jedan. Ranějša primětka opisuje odgovarjajuče najmenše polje prěz kompozitum Galoisovyh polj za razpadanje oprědělenyh polinomov. Sam pridevnik najmenše ne može zaměniti toto svojstvo.\par
\noindent\emph{Textual comparison:} The condition on $Q$ is tied to decomposition of the ring into rank-one components. An earlier note describes the corresponding least field by the compositum of Galois fields splitting the specified polynomials. Least alone cannot replace this property.\par
\noindent{\footnotesize TEXT-COMP-149; pasaže / passages: 18474, 17794; \texttt{CORPUS\_FORM\_COMPARISONS.json}.}\par
\end{minipage}\par
\par\medskip\Needspace{7\baselineskip}
\noindent\begin{minipage}{\linewidth}
\hypertarget{TERM-0988}{}
\noindent\textbf{najmenši obči imenovatelj}\quad{\footnotesize\texttt{TERM-0988}}\par
\noindent\emph{Iztočnik / Source:} kleinster gemeinsamer Nenner\par
\noindent\emph{Rusky / Russian:} наименьший общий знаменатель\par
\noindent\emph{Ukrajinsky / Ukrainian:} найменший спільний знаменник\par
\smallskip\noindent\emph{Razlog:} Najmenši obči imenovatelj v rabotě 6, §8, svazuje konstrukciju s ranějšim objasnjenjem fiksovanogo imenovatelja. Lokalna definicija trěbuje, da polinomy $F,F_1,\ldots,F_k$ budut vzajemno proste v danoj oblasti, kogda $f_i=F_i/F$. Slědujuči odlomok dopušča različne take prědstavjenja, zato že jednoznačno razloženje v oblasti obče ne dějstvuje. Izbor znanogo aritmetičnogo nazvanja ne daje prava predpolagati čislovy poredok ili jednoznačnost imenovatelja, ktoryh tekst ne utvrdžaje.\par
\noindent\emph{Reason:} Najmenši obči imenovatelj in Paper 6, §8 connects the construction to the earlier fixed-denominator explanation. The local definition requires $F,F_1,\ldots,F_k$ to be relatively prime in the specified domain when $f_i=F_i/F$. The next paragraph permits different such representations because unique factorization does not generally hold there. Familiar arithmetic terminology does not license importing a numerical ordering or uniqueness of the denominator that the text has not established.\par
\noindent{\footnotesize Izvorne zapisy / Original records: HIST-g048-006.}\par
\smallskip\noindent\emph{Sravnjenje s tekstom:} Uslovje za najmenši obči imenovatelj govori o vzajemnoj prostosti $F,F_1,\ldots,F_k$ v izbranoj oblasti. Tekst odmah dopušča različne prědstavjenja i opozarja, že jednoznačno razloženje tam obče ne dějstvuje.\par
\noindent\emph{Textual comparison:} The least-common-denominator condition concerns relative primality of $F,F_1,\ldots,F_k$ in the chosen domain. The text immediately allows different representations and warns that unique factorization does not generally hold there.\par
\noindent{\footnotesize TEXT-COMP-150; pasaže / passages: 5876, 5878, 5880, 5882, 5891; \texttt{CORPUS\_FORM\_COMPARISONS.json}.}\par
\end{minipage}\par
\par\medskip\Needspace{7\baselineskip}
\noindent\begin{minipage}{\linewidth}
\hypertarget{TERM-0989}{}
\noindent\textbf{najmene obče mnogokratno}\quad{\footnotesize\texttt{TERM-0989}}\par
\noindent\emph{Iztočnik / Source:} kleinstes gemeinsames Vielfaches\par
\noindent\emph{Rusky / Russian:} наименьшее общее кратное\par
\noindent\emph{Ukrajinsky / Ukrainian:} найменше спільне кратне\par
\smallskip\noindent\emph{Razlog:} Najmene obče mnogokratno prěnosi obyčno aritmetično nazvanje na moduly v rabotě 17. Zapis ostavja otvorjeno pytanje razumljivosti v abstraktnom izloženju. Nyněšnji tekst imaje najmenše obče mnogokratno i definira objekt kako cělost polinomov, dělimyh prěz oba moduly; definicija razširjaje se takože na semeju proizvoljnoj močnosti. Zato čitateljska podpora trěbuje definiciju operacije, ne samo razpoznavajemy koren. Arhivno najmene ne jest predstavljeno kak bezuslovno rekomendovana forma.\par
\noindent\emph{Reason:} Najmene obče mnogokratno transfers ordinary arithmetic terminology to modules in Paper 17. The record expressly leaves its readability in abstract exposition open. The present text uses najmenše obče mnogokratno and defines the object as all polynomials divisible by both modules, extending the definition to families of arbitrary cardinality. Readers therefore need the operation’s definition, not merely a recognizable stem. Archival najmene is not presented as an unqualified recommended form.\par
\noindent{\footnotesize Izvorne zapisy / Original records: HIST-g079-008.}\par
\smallskip\noindent\emph{Sravnjenje s tekstom:} Raboty 17 i 19 definirajut mnogokratno prěz obču děliteljnost; rabota 31 upotrěbjaje kratnik pri razloženju nulovogo ideala. V rabotě 30 ravnost s produktom ima izrično uslovje poparnoj vzajemnoj prostosti v komutativnom prstenu s jediniceju. Sam naziv ne izjednja operacije.\par
\noindent\emph{Textual comparison:} Papers 17 and 19 define the multiple through common divisibility; Paper 31 uses kratnik in decomposing the zero ideal. In Paper 30, equality with the product has the explicit hypothesis of pairwise coprimality in a commutative ring with identity. The name alone does not identify the operations.\par
\noindent{\footnotesize TEXT-COMP-151; pasaže / passages: 11163, 12562, 17273, 17665, 16493, 16507; \texttt{CORPUS\_FORM\_COMPARISONS.json}.}\par
\end{minipage}\par
\par\medskip\Needspace{7\baselineskip}
\noindent\begin{minipage}{\linewidth}
\hypertarget{TERM-0990}{}
\noindent\textbf{najmenše obče mnogokratno}\quad{\footnotesize\texttt{TERM-0990}}\par
\noindent\emph{Iztočnik / Source:} kleinstes gemeinsames Vielfaches\par
\noindent\emph{Rusky / Russian:} наименьшее общее кратное\par
\noindent\emph{Ukrajinsky / Ukrainian:} найменше спільне кратне\par
\smallskip\noindent\emph{Razlog:} Najmenše obče mnogokratno sbira devet zapisov iz rabot 19, 22, 24 i 25. Razlogy sohranjajut Noetherin jezyk děliteljnosti, vezu s modulami i kontinuitet medžu definicijeju, faktorizacijeju i povratom k prědstavjenju ideala. Mnogokratno izbrano jest kak sostavny slovjansky izraz vměsto latinskoj skračenosti ili užšej nacionalnoj formy. V rabotě 19 objekt definiran jest vsimi elementami, dělimymi prěz vsaky dany ideal. Sravnjenje s prěsěkom može objasniti konstrukciju, ale ne směje izbrisati avtorov izraz iz prevoda. Opis prozračnosti ostaje kvalitativny.\par
\noindent\emph{Reason:} Najmenše obče mnogokratno brings together nine records from Papers 19, 22, 24 and 25. Their reasons preserve Noether’s divisibility language, its module interpretation, and continuity from the definition through factorization back to ideal representation. Mnogokratno is chosen as a compositional Slavic expression rather than a Latin abbreviation or a narrower national form. Paper 19 defines the object through elements divisible by every specified ideal. An intersection gloss can explain the construction, but should not erase the author’s wording from the translation. Claims of transparency remain qualitative.\par
\noindent{\footnotesize Izvorne zapisy / Original records: HIST-g081-009, HIST-g082-009, HIST-g086-006, HIST-g089-012, HIST-g092-016, HIST-g092-054, HIST-g096-007, HIST-g099-004, HIST-g102-007.}\par
\smallskip\noindent\emph{Sravnjenje s tekstom:} Raboty 17 i 19 definirajut mnogokratno prěz obču děliteljnost; rabota 31 upotrěbjaje kratnik pri razloženju nulovogo ideala. V rabotě 30 ravnost s produktom ima izrično uslovje poparnoj vzajemnoj prostosti v komutativnom prstenu s jediniceju. Sam naziv ne izjednja operacije.\par
\noindent\emph{Textual comparison:} Papers 17 and 19 define the multiple through common divisibility; Paper 31 uses kratnik in decomposing the zero ideal. In Paper 30, equality with the product has the explicit hypothesis of pairwise coprimality in a commutative ring with identity. The name alone does not identify the operations.\par
\noindent{\footnotesize TEXT-COMP-151; pasaže / passages: 11163, 12562, 17273, 17665, 16493, 16507; \texttt{CORPUS\_FORM\_COMPARISONS.json}.}\par
\end{minipage}\par
\par\medskip\Needspace{7\baselineskip}
\noindent\begin{minipage}{\linewidth}
\hypertarget{TERM-0991}{}
\noindent\textbf{najmenši obči kratnik}\quad{\footnotesize\texttt{TERM-0991}}\par
\noindent\emph{Iztočnik / Source:} kleinstes gemeinsames Vielfaches\par
\noindent\emph{Rusky / Russian:} наименьшее общее кратное\par
\noindent\emph{Ukrajinsky / Ukrainian:} найменше спільне кратне\par
\smallskip\noindent\emph{Razlog:} Najmenši obči kratnik v rabotě 31 jest druga dokumentovana forma dlja najmenšego občego mnogokratnogo. Motivacija sohranja starši jezyk děliteljnosti idealov i jego paru s največšim občim děliteljem. Sravnjene odlomky upotrěbjajut kratnik pri prědstavjenju nulovogo ideala prěz proste idealy. To ne jest nova operacija samo zato, že druge raboty imajut mnogokratno. Ne trěba tiho zaměniti izraz ni prěsěkom, ni produktom: ravnost s produktom trěbuje svoje osobne uslovja.\par
\noindent\emph{Reason:} Najmenši obči kratnik in Paper 31 is another recorded form for least common multiple. The rationale retains older ideal-divisibility language and its pairing with greatest common divisor. The compared passages use kratnik when expressing the zero ideal through prime ideals. The different noun does not itself introduce a new operation beside mnogokratno elsewhere. Neither intersection nor product should silently replace the term; equality with the product requires its own hypotheses.\par
\noindent{\footnotesize Izvorne zapisy / Original records: HIST-g121-009, HIST-g126-004.}\par
\smallskip\noindent\emph{Sravnjenje s tekstom:} Raboty 17 i 19 definirajut mnogokratno prěz obču děliteljnost; rabota 31 upotrěbjaje kratnik pri razloženju nulovogo ideala. V rabotě 30 ravnost s produktom ima izrično uslovje poparnoj vzajemnoj prostosti v komutativnom prstenu s jediniceju. Sam naziv ne izjednja operacije.\par
\noindent\emph{Textual comparison:} Papers 17 and 19 define the multiple through common divisibility; Paper 31 uses kratnik in decomposing the zero ideal. In Paper 30, equality with the product has the explicit hypothesis of pairwise coprimality in a commutative ring with identity. The name alone does not identify the operations.\par
\noindent{\footnotesize TEXT-COMP-151; pasaže / passages: 11163, 12562, 17273, 17665, 16493, 16507; \texttt{CORPUS\_FORM\_COMPARISONS.json}.}\par
\end{minipage}\par
\par\medskip\Needspace{7\baselineskip}
\noindent\begin{minipage}{\linewidth}
\hypertarget{TERM-0992}{}
\noindent\textbf{najmenše obče mnogokratno / največši obči dělitelj}\quad{\footnotesize\texttt{TERM-0992}}\par
\noindent\emph{Iztočnik / Source:} kleinstes gemeinsames Vielfaches / größter gemeinsamer Teiler\par
\noindent\emph{Rusky / Russian:} наименьшее общее кратное / наибольший общий делитель\par
\noindent\emph{Ukrajinsky / Ukrainian:} найменше спільне кратне / найбільший спільний дільник\par
\smallskip\noindent\emph{Razlog:} Najmenše obče mnogokratno / največši obči dělitelj v rabotě 30, §5, sohranja paru iz avtorovoj aritmetiky idealov. Razlog upotrěbljenja jest kontinuitet s §4, ne odkaz od mogućego sovrěmennogo objasnjenja. Razlika medžu děliteljem i mnogokratnym upravjaje směr zaključkov, zato oba izrazy dolžat ostati različny. Pojasnjenje čitatelju može sravniti jih s sumoju i prěsěkom po lokalnyh definicijah; prevod ne směje avtomatično podměniti tu dvojicu drugoju terminologijeju.\par
\noindent\emph{Reason:} Najmenše obče mnogokratno / največši obči dělitelj in Paper 30, §5 preserves the paired vocabulary of the author’s ideal arithmetic. The reason is continuity with §4, not a prohibition on modern explanatory glosses. Divisor and multiple govern the direction of the arguments and must remain distinct. A reader’s explanation may compare them with sum and intersection under the local definitions; the translation should not automatically replace the pair with that different terminology.\par
\noindent{\footnotesize Izvorne zapisy / Original records: HIST-g112-011.}\par
\smallskip\noindent\emph{Sravnjenje s tekstom:} Raboty 17 i 19 definirajut mnogokratno prěz obču děliteljnost; rabota 31 upotrěbjaje kratnik pri razloženju nulovogo ideala. V rabotě 30 ravnost s produktom ima izrično uslovje poparnoj vzajemnoj prostosti v komutativnom prstenu s jediniceju. Sam naziv ne izjednja operacije.\par
\noindent\emph{Textual comparison:} Papers 17 and 19 define the multiple through common divisibility; Paper 31 uses kratnik in decomposing the zero ideal. In Paper 30, equality with the product has the explicit hypothesis of pairwise coprimality in a commutative ring with identity. The name alone does not identify the operations.\par
\noindent{\footnotesize TEXT-COMP-151; pasaže / passages: 11163, 12562, 17273, 17665, 16493, 16507; \texttt{CORPUS\_FORM\_COMPARISONS.json}.}\par
\end{minipage}\par
\par\medskip\Needspace{7\baselineskip}
\noindent\begin{minipage}{\linewidth}
\hypertarget{TERM-0995}{}
\noindent\textbf{koeficienty kako simetrične funkcije integralov}\quad{\footnotesize\texttt{TERM-0995}}\par
\noindent\emph{Iztočnik / Source:} Koeffizienten als symmetrische Funktionen der Integrale\par
\noindent\emph{Rusky / Russian:} коэффициенты как симметрические функции интегралов\par
\noindent\emph{Ukrajinsky / Ukrainian:} коефіцієнти як симетричні функції інтегралів\par
\smallskip\noindent\emph{Razlog:} Koeficienty kako simetrične funkcije integralov v rabotě 17 čuvaje avtorovo nazvanje integralov iz §11. V priměru ide o dva rěšenja diferencialnogo uravnenja drugogo reda; oblast racionalnosti sostoji se iz racionalnyh funkcij tih integralov i jih derivatov v dostatočno maloj okolini regularnoj točky. Slovnik može objasniti rěšiteljny smysl integralov, ale sam prevod ne zaměnjaje ih tiho slovom rěšenja. Taka historična věrnost ne dava prava čitati ih kak proizvoljne oprěděljene integraly.\par
\noindent\emph{Reason:} Koeficienty kako simetrične funkcije integralov in Paper 17 retains the author’s integral terminology from §11. The example concerns two solutions of a second-order differential equation; its rationality domain consists of rational functions of those integrals and their derivatives in a sufficiently small neighbourhood of a regular point. The dictionary can explain the solution sense, while the translation does not silently replace integrals by solutions. Historical fidelity does not justify reading them as arbitrary definite integrals.\par
\noindent{\footnotesize Izvorne zapisy / Original records: HIST-g079-072.}\par
\smallskip\noindent\emph{Sravnjenje s tekstom:} Priměr raboty 17 govori o uravnenju drugogo reda, dvoh integralah i jih derivatah v maloj okolini regularnoj točky. To jest rěšiteljny smysl integralov, različny od proizvoljnoj čislovoj integracije.\par
\noindent\emph{Textual comparison:} The Paper 17 example concerns a second-order equation, two integrals and their derivatives near a regular point. This is the solution sense of integrals, not arbitrary numerical integration.\par
\noindent{\footnotesize TEXT-COMP-153; pasaže / passages: 12079; \texttt{CORPUS\_FORM\_COMPARISONS.json}.}\par
\end{minipage}\par
\par\medskip\Needspace{7\baselineskip}
\noindent\begin{minipage}{\linewidth}
\hypertarget{TERM-0996}{}
\noindent\textbf{koeficientno podružje}\quad{\footnotesize\texttt{TERM-0996}}\par
\noindent\emph{Iztočnik / Source:} Koeffizientenbereich\par
\noindent\emph{Rusky / Russian:} область коэффициентов\par
\noindent\emph{Ukrajinsky / Ukrainian:} область коефіцієнтів\par
\smallskip\noindent\emph{Razlog:} Koeficientno podružje v rabotě 5 sohranja širinu německogo Bereich. Motivacija naměrno ne izbira jedno sovrěmenno strukturalno nazvanje tam, kde značenje zavisi od konteksta. To ne znači, že svojstva oblasti sut proizvoljne: ih trěbuje uviděti iz matematičnyh prědpoloženij. Rusko i ukrajinsko область davaje oblastnu metaforu, ale ne dokazuje točnu medžuslovjansku formu podružje. Slovnik sohranja toj izbor vedle kasnějšego oblast koeficientov, ne predstavlja obě formy kak jedno bezsporno pravilo.\par
\noindent\emph{Reason:} Koeficientno podružje in Paper 5 retains the breadth of German Bereich. The rationale deliberately avoids imposing one modern structural name where the sense depends on context. That does not make the domain’s properties arbitrary: they must be read from the mathematical assumptions. Russian and Ukrainian область support the domain metaphor, but do not establish the exact Interslavic form podružje. The dictionary preserves this choice alongside later oblast koeficientov rather than presenting both as one uncontested rule.\par
\noindent{\footnotesize Izvorne zapisy / Original records: HIST-g039-004.}\par
\end{minipage}\par
\par\medskip\Needspace{7\baselineskip}
\noindent\begin{minipage}{\linewidth}
\hypertarget{TERM-0997}{}
\noindent\textbf{oblast koeficientov}\quad{\footnotesize\texttt{TERM-0997}}\par
\noindent\emph{Iztočnik / Source:} Koeffizientenbereich\par
\noindent\emph{Rusky / Russian:} область коэффициентов\par
\noindent\emph{Ukrajinsky / Ukrainian:} область коефіцієнтів\par
\smallskip\noindent\emph{Razlog:} Oblast koeficientov sohranja oblastny izraz v rabotah 18, 20 i 43. Zapis o rabotě 20 izrično kaže, že ranějše koeficientno podružje bylo zaměnjeno po obrazu raboty 18; to jest svědectvo o obrazu izbora, ne nezavisna norma. Pri prisojedinjenju proměnnyh i pri konstrukciji direktnogo produkta Bereich ne dolžen avtomatično stati poljem. Genitiv koeficientov pojasnjaje, kakvu rolu oblast igraje. Poslědnji razlog ostavja oblast na jezykovu ocěnu zaradi obyčnojezykovoj širiny slova.\par
\noindent\emph{Reason:} Oblast koeficientov retains domain wording in Papers 18, 20 and 43. The Paper 20 record explicitly reports replacing koeficientno podružje on the model of Paper 18; this is evidence about the choice’s basis, not an independent norm. In adjoining variables and constructing a direct product, Bereich must not automatically become field. The genitive specifies the domain’s role. The last rationale leaves oblast open to linguistic appraisal because of the ordinary word’s breadth.\par
\noindent{\footnotesize Izvorne zapisy / Original records: HIST-g080-007, HIST-g083-003, HIST-g213-005.}\par
\end{minipage}\par
\par\medskip\Needspace{7\baselineskip}
\noindent\begin{minipage}{\linewidth}
\hypertarget{TERM-0998}{}
\noindent\textbf{oblast koeficientov / polje $P(u)$}\quad{\footnotesize\texttt{TERM-0998}}\par
\noindent\emph{Iztočnik / Source:} Koeffizientenbereich / Körper $P(u)$\par
\noindent\emph{Rusky / Russian:} область коэффициентов / поле $P(u)$\par
\noindent\emph{Ukrajinsky / Ukrainian:} область коефіцієнтів / поле $P(u)$\par
\smallskip\noindent\emph{Razlog:} Oblast koeficientov / polje $P(u)$ naměrno razlikuje dva stupnje konstrukcije v rabotě 22, §3. K oblasti prisojedinjajut se neoprěděljene $u_{\mu\nu}$, a daljši koeficienty $\alpha_\lambda(u)$ sut racionalne funkcije tih neoprěděljenyh s koeficientami iz $P$. Zato $P(u)$ ne označuje samo polinomny prsten $P[u]$. Izbor polje v drugoj polovici sočetanja podpiraje konkretna konstrukcija, ne obče pravilo zaměnjati vsako Bereich poljem.\par
\noindent\emph{Reason:} Oblast koeficientov / polje $P(u)$ deliberately distinguishes two stages in Paper 22, §3. The indeterminates $u_{\mu\nu}$ are adjoined, and the subsequent coefficients $\alpha_\lambda(u)$ are rational functions of them with coefficients in $P$. Thus $P(u)$ does not denote merely the polynomial ring $P[u]$. Field wording in the second half is justified by this construction, not by a general rule replacing every Bereich with field.\par
\noindent{\footnotesize Izvorne zapisy / Original records: HIST-g087-002, HIST-g092-022.}\par
\smallskip\noindent\emph{Sravnjenje s tekstom:} V rabotě 22 koeficienty $\alpha_\lambda(u)$ izrično sut racionalne funkcije. Polje $P(u)$ zato ne treba zaměniti polinomnym prstenom, hotě za konkretny konečny izraz možno odstraniti imenovatelje vhodnym množenjem.\par
\noindent\emph{Textual comparison:} Paper 22 explicitly makes $\alpha_\lambda(u)$ rational functions. The field $P(u)$ must therefore not be confused with a polynomial ring, even though a suitable multiplication can clear denominators in a particular finite expression.\par
\noindent{\footnotesize TEXT-COMP-155; pasaže / passages: 14514; \texttt{CORPUS\_FORM\_COMPARISONS.json}.}\par
\end{minipage}\par
\par\medskip\Needspace{7\baselineskip}
\noindent\begin{minipage}{\linewidth}
\hypertarget{TERM-0999}{}
\noindent\textbf{koeficientna oblast $\Omega$}\quad{\footnotesize\texttt{TERM-0999}}\par
\noindent\emph{Iztočnik / Source:} Koeffizientenbereich $\Omega$\par
\noindent\emph{Rusky / Russian:} область коэффициентов $\Omega$\par
\noindent\emph{Ukrajinsky / Ukrainian:} область коефіцієнтів $\Omega$\par
\smallskip\noindent\emph{Razlog:} Koeficientna oblast $\Omega$ v rabotě 6, §1, sohranja Bereich kako oblast, dokolě sam tekst ne utočni strukturu. Odlomok stvarno prědpoklada čislovo polje i izrično dopušča vse kompleksne čisla i ješče konečno čislo parametrov. Zato tu ne trěba avtomatično čitati čislovo polje v užšem sovrěmennom smyslu konečnogo razširjenja racionalnyh čisel. Pridevnična konstrukcija čuva označenje $\Omega$ i oddělja nazvanje roli od polnyh prědpoloženij o objektu.\par
\noindent\emph{Reason:} Koeficientna oblast $\Omega$ in Paper 6, §1 retains Bereich as domain until the text specifies the structure. The passage assumes a Zahlkörper but expressly allows all complex numbers and finitely many further parameters. It should therefore not automatically be read as number field in the narrower modern sense of a finite extension of the rationals. The adjectival construction preserves $\Omega$ while distinguishing the name of its role from the full assumptions on the object.\par
\noindent{\footnotesize Izvorne zapisy / Original records: HIST-g041-006.}\par
\smallskip\noindent\emph{Sravnjenje s tekstom:} Na početku raboty 6 $\Omega$ možno vzeti kako vse kompleksne čisla i vključiti parametry. §6 govori o algebraičnom razširjenju koeficientnogo polja, a primětka §9 zadržava oblast koeficientov pri $\rho=0$. Taka kontekstualna proměna ne dolžna byti skryta za jedinym avtomatičnym prevodom Bereich.\par
\noindent\emph{Textual comparison:} At the start of Paper 6, $\Omega$ may include all complex numbers and parameters. Section 6 speaks of an algebraic extension of the coefficient field, while the §9 footnote retains coefficient-domain wording at $\rho=0$. This contextual variation should not be hidden behind one automatic translation of Bereich.\par
\noindent{\footnotesize TEXT-COMP-154; pasaže / passages: 5061, 5698, 5987; \texttt{CORPUS\_FORM\_COMPARISONS.json}.}\par
\end{minipage}\par
\par\medskip\Needspace{7\baselineskip}
\noindent\begin{minipage}{\linewidth}
\hypertarget{TERM-1000}{}
\noindent\textbf{koeficientno polje $\Omega$}\quad{\footnotesize\texttt{TERM-1000}}\par
\noindent\emph{Iztočnik / Source:} Koeffizientenbereich $\Omega$\par
\noindent\emph{Rusky / Russian:} поле коэффициентов $\Omega$\par
\noindent\emph{Ukrajinsky / Ukrainian:} поле коефіцієнтів $\Omega$\par
\smallskip\noindent\emph{Razlog:} Koeficientno polje $\Omega$ v rabotě 6, §6, jest kontekstualny izbor, različny od oblast v drugyh odlomkah. Motivacija opiraje se na argumenty o stepeni razširjenja, v ktoryh $\Omega$ služi kak osnovno polje. Nyněšnji odlomok o minimalnoj bazi takože govori o dopuščenju algebraičnogo razširjenja koeficientnogo polja. Zapis ostavja redakcijnu ostorožnost zaradi širšego Bereich. Slovnik zato ne udělaje iz togo jednogo izbora pravilo, že oblast i polje sut povsudu zaměnljive.\par
\noindent\emph{Reason:} Koeficientno polje $\Omega$ in Paper 6, §6 is a contextual choice, distinct from oblast in other passages. The rationale rests on extension-degree arguments in which $\Omega$ serves as the base field. The present minimal-basis passage likewise speaks of allowing an algebraic extension of the coefficient field. The record retains an editorial reservation because Bereich is broader. This one choice therefore does not make domain and field interchangeable throughout the corpus.\par
\noindent{\footnotesize Izvorne zapisy / Original records: HIST-g046-007.}\par
\smallskip\noindent\emph{Sravnjenje s tekstom:} Na početku raboty 6 $\Omega$ možno vzeti kako vse kompleksne čisla i vključiti parametry. §6 govori o algebraičnom razširjenju koeficientnogo polja, a primětka §9 zadržava oblast koeficientov pri $\rho=0$. Taka kontekstualna proměna ne dolžna byti skryta za jedinym avtomatičnym prevodom Bereich.\par
\noindent\emph{Textual comparison:} At the start of Paper 6, $\Omega$ may include all complex numbers and parameters. Section 6 speaks of an algebraic extension of the coefficient field, while the §9 footnote retains coefficient-domain wording at $\rho=0$. This contextual variation should not be hidden behind one automatic translation of Bereich.\par
\noindent{\footnotesize TEXT-COMP-154; pasaže / passages: 5061, 5698, 5987; \texttt{CORPUS\_FORM\_COMPARISONS.json}.}\par
\end{minipage}\par
\par\medskip\Needspace{7\baselineskip}
\noindent\begin{minipage}{\linewidth}
\hypertarget{TERM-1001}{}
\noindent\textbf{oblast koeficientov $\Omega$}\quad{\footnotesize\texttt{TERM-1001}}\par
\noindent\emph{Iztočnik / Source:} Koeffizientenbereich $\Omega$\par
\noindent\emph{Rusky / Russian:} область коэффициентов $\Omega$\par
\noindent\emph{Ukrajinsky / Ukrainian:} область коефіцієнтів $\Omega$\par
\smallskip\noindent\emph{Razlog:} Oblast koeficientov $\Omega$ v podnožnoj primětkě raboty 6, §9, naměrno ostaje bližše k Bereich, než ranějše polje. Primětka govori o slučaju $\rho=0$, kogda cělost polinomov v razgledanom polju sostoji se iz elementov $\Omega$. Genitivna konstrukcija dozvaljaje sohraniti avtorovo nazvanje bez zaměny samogo objekta. Tuto i koeficientna oblast sut različne sintaktične formy; razlika oblast/polje potrěbuje strukturny kontekst, ne samo soglasovanje slov.\par
\noindent\emph{Reason:} Oblast koeficientov $\Omega$ in the footnote to Paper 6, §9 deliberately stays closer to Bereich than the field wording used earlier. The note concerns $\rho=0$, when the polynomials in the field under discussion consist of elements of $\Omega$. The genitive construction preserves the author’s designation without changing the object. It and koeficientna oblast are different syntactic forms; the domain/field distinction requires structural context, not merely consistent vocabulary.\par
\noindent{\footnotesize Izvorne zapisy / Original records: HIST-g049-005.}\par
\smallskip\noindent\emph{Sravnjenje s tekstom:} Na početku raboty 6 $\Omega$ možno vzeti kako vse kompleksne čisla i vključiti parametry. §6 govori o algebraičnom razširjenju koeficientnogo polja, a primětka §9 zadržava oblast koeficientov pri $\rho=0$. Taka kontekstualna proměna ne dolžna byti skryta za jedinym avtomatičnym prevodom Bereich.\par
\noindent\emph{Textual comparison:} At the start of Paper 6, $\Omega$ may include all complex numbers and parameters. Section 6 speaks of an algebraic extension of the coefficient field, while the §9 footnote retains coefficient-domain wording at $\rho=0$. This contextual variation should not be hidden behind one automatic translation of Bereich.\par
\noindent{\footnotesize TEXT-COMP-154; pasaže / passages: 5061, 5698, 5987; \texttt{CORPUS\_FORM\_COMPARISONS.json}.}\par
\end{minipage}\par
\par\medskip\Needspace{7\baselineskip}
\noindent\begin{minipage}{\linewidth}
\hypertarget{TERM-1002}{}
\noindent\textbf{ravnost koeficientov}\quad{\footnotesize\texttt{TERM-1002}}\par
\noindent\emph{Iztočnik / Source:} Koeffizientengleichheit\par
\noindent\emph{Rusky / Russian:} равенство коэффициентов\par
\noindent\emph{Ukrajinsky / Ukrainian:} рівність коефіцієнтів\par
\smallskip\noindent\emph{Razlog:} Ravnost koeficientov izbira prostu genitivnu konstrukciju dlja kriterija jednoznačnosti v rabotě 31, §1. Sravnjeni odlomok prvo trěbuje izraziti elementy prěz linearno nezavisne elementy; potom ravnost určuje se koeficientami v razširjenom polju. Bez toho prědpoloženja ravnost libovoljnyh linearnih izrazov ne trěbuje člennogo sběganja koeficientov. Kratkost frazy zato oslanjaje se na jasno izloženu bazu, ne zaměnjaje ju. Ruska i ukrajinska konstrukcija davajut prěmo sravnjenje, a izbor prigodny jest za povtorno upotrěbljenje.\par
\noindent\emph{Reason:} Ravnost koeficientov selects a simple genitive construction for the uniqueness criterion in Paper 31, §1. The compared passage first requires expressing elements in terms of linearly independent elements; equality is then determined by coefficients in the extended field. Without that qualification, equality of arbitrary linear expressions need not mean termwise equality of coefficients. The phrase’s brevity therefore depends on the basis being stated. Russian and Ukrainian provide direct comparators, and the rationale favours a reusable formulation.\par
\noindent{\footnotesize Izvorne zapisy / Original records: HIST-g119-011.}\par
\smallskip\noindent\emph{Sravnjenje s tekstom:} Tekst raboty 31 svazuje ravnost koeficientov s prědstavjenjem po linearno nezavisnyh elementah. Podnožna primětka razgledaje, kako razširjenje koeficientov može vvesti nove zavisnosti i nulove dělitelje. Kratky termin ne odstranjaje ty ograničenja.\par
\noindent\emph{Textual comparison:} Paper 31 ties coefficient equality to expression in linearly independent elements. Its footnote considers how extending coefficients can introduce new dependencies and zero divisors. The compact term does not remove those qualifications.\par
\noindent{\footnotesize TEXT-COMP-156; pasaže / passages: 17363; \texttt{CORPUS\_FORM\_COMPARISONS.json}.}\par
\end{minipage}\par
\par\medskip\Needspace{7\baselineskip}
\noindent\begin{minipage}{\linewidth}
\hypertarget{TERM-1003}{}
\noindent\textbf{kogredientny}\quad{\footnotesize\texttt{TERM-1003}}\par
\noindent\emph{Iztočnik / Source:} kogredient\par
\noindent\emph{Rusky / Russian:} когредиентный\par
\noindent\emph{Ukrajinsky / Ukrainian:} когредієнтний\par
\smallskip\noindent\emph{Razlog:} Kogredientny v rabotě 12 sohranja rědky historijny termin teorije invariantov. Tekst sam objasnjaje jego rolju: veličiny, ktore zaměnjajut vyšše diferencialy, podležat tym samym linearnym transformacijam kak prve diferencialy. Zato ne dostatočno skazati, že one odgovarjajut jedna drugoj; sučstvena jest transformacijna relacija. Motivacija otklanja tihu zaměnu s kovariantny ili kontragredientny. Sravnjeni odlomok podpiraje to različenje, ale ne ustanavjaje, že svaki kovariant v vsakoj sovrěmennoj konvenciji imaje toto nazvanje.\par
\noindent\emph{Reason:} Kogredientny in Paper 12 preserves a rare historical invariant-theory term. The text explains its role: the quantities replacing higher differentials undergo the same linear transformations as the first differentials. Merely saying that they correspond would lose the transformation relation. The rationale rejects silent replacement by covariant or contragredient. The compared passage supports that distinction, not a claim that every covariant under every modern convention has this designation.\par
\noindent{\footnotesize Izvorne zapisy / Original records: HIST-g074-013.}\par
\smallskip\noindent\emph{Sravnjenje s tekstom:} Rabota 12 pojasnjaje kogredientnost tymi samymi linearnymi transformacijami. Rabota 2 razlikuje kogredientne i kontragredientne redy v formulah, a rabota 23 koristi kogredientnost za oprěděljenje grupy. Poslědnja pri tom ne tvrdi, že obči dokaz navedenej redukcije už jest znany.\par
\noindent\emph{Textual comparison:} Paper 12 explains cogredience through the same linear transformations. Paper 2 distinguishes cogredient and contragredient rows in formulas, and Paper 23 uses cogredience to define the group. The latter does not claim that a general proof of the stated reduction is already known.\par
\noindent{\footnotesize TEXT-COMP-157; pasaže / passages: 400, 402, 9074, 9146, 15205; \texttt{CORPUS\_FORM\_COMPARISONS.json}.}\par
\end{minipage}\par
\par\medskip\Needspace{7\baselineskip}
\noindent\begin{minipage}{\linewidth}
\hypertarget{TERM-1004}{}
\noindent\textbf{kogredientny / kontragredientny}\quad{\footnotesize\texttt{TERM-1004}}\par
\noindent\emph{Iztočnik / Source:} kogredient / kontragradient\par
\noindent\emph{Rusky / Russian:} коградиентный / контраградиентный\par
\noindent\emph{Ukrajinsky / Ukrainian:} коградієнтний / контраградієнтний\par
\smallskip\noindent\emph{Razlog:} Kogredientny / kontragredientny v rabotě 21 sohranja medžunarodnu semeju transformacijnyh terminov. Razlog predpočitaje kratke nazvanja dolgomu opisnomu prevodu, zato že dokaz sledi različne sposoby transformacije. To ne jest dovoljenje čitati predponu kontra- kak obyčnu protivnost bez formul. Napisanje iztočnogo zapisa kontragradient i ruskih/ukrajinskih sravnjenij s -градиент- različaje se od izbranogo -gredient-. Tute formy sut dokumentovane; sam zapis ne ustanavjaje, kaky variant byl na pečatnoj stranici.\par
\noindent\emph{Reason:} Kogredientny / kontragredientny in Paper 21 retains the international family of transformation terms. The rationale prefers compact names to a long paraphrase because the proof tracks different transformation behaviours. This does not permit reading contra- as an unspecified opposition without the formulas. The archived source label kontragradient and the Russian/Ukrainian comparators with -градиент- differ from the chosen -gredient- spelling. Those forms are documented; the record alone does not establish which spelling appeared on the printed page.\par
\noindent{\footnotesize Izvorne zapisy / Original records: HIST-g084-006.}\par
\end{minipage}\par
\par\medskip\Needspace{7\baselineskip}
\noindent\begin{minipage}{\linewidth}
\hypertarget{TERM-1005}{}
\noindent\textbf{kogredientny / kontragredientny}\quad{\footnotesize\texttt{TERM-1005}}\par
\noindent\emph{Iztočnik / Source:} kogredient / kontragredient\par
\noindent\emph{Rusky / Russian:} когредиентный / контрагредиентный\par
\noindent\emph{Ukrajinsky / Ukrainian:} когредієнтний / контрагредієнтний\par
\smallskip\noindent\emph{Razlog:} Kogredientny / kontragredientny v zapisah rabot 3 i 4 izbrano jest s -gredient-, po vzoru iztočnogo termina i ruskogo/ukrajinskogo sravnjenja. Razlog izrično opozarja na variant kogradient iz §26 i trěbuje jeho različenje od izbranogo pravopisa. V sravnjenom korpusu jest kogradientne pri pozdnějšem odlomku raboty 2, ale kogredientnyh pri objasnjenju dvoh kontrakcij v rabotě 3. Zapisano pravilo zato ne jest dokaz jego vsestrannogo priměnjenja. Razlika pisanja ne směje skryti matematičnu protivnost cogredient/contragredient.\par
\noindent\emph{Reason:} Kogredientny / kontragredientny in the records for Papers 3 and 4 chooses -gredient- on the model of the source term and the Russian/Ukrainian comparators. The rationale expressly flags kogradient from §26 as a spelling variant to distinguish from that choice. The compared corpus has kogradientne in a later Paper 2 passage but kogredientnyh in Paper 3’s explanation of two contractions. A recorded rule is therefore not proof of uniform implementation. The spelling issue must not obscure the mathematical cogredient/contragredient distinction.\par
\noindent{\footnotesize Izvorne zapisy / Original records: HIST-g028-007, HIST-g030-010.}\par
\smallskip\noindent\emph{Sravnjenje s tekstom:} V rabotě 2 čita se kogradientne, dokolě rabota 3 imaje kogredientnyh pri dvoh imenovanyh kontrakcijah. To potvrždaje, že arhivno pravopisno pravilo ne ustanavjaje avtomatično doslědnost vsego korpusa.\par
\noindent\emph{Textual comparison:} Paper 2 reads kogradientne, whereas Paper 3 has kogredientnyh for the two named contractions. This confirms that an archived spelling rule does not automatically establish consistency throughout the corpus.\par
\noindent{\footnotesize TEXT-COMP-158; pasaže / passages: 3344, 3649; \texttt{CORPUS\_FORM\_COMPARISONS.json}.}\par
\end{minipage}\par
\par\medskip\Needspace{7\baselineskip}
\noindent\begin{minipage}{\linewidth}
\hypertarget{TERM-1006}{}
\noindent\textbf{kogredientno svijanje}\quad{\footnotesize\texttt{TERM-1006}}\par
\noindent\emph{Iztočnik / Source:} kogrediente Faltung\par
\noindent\emph{Rusky / Russian:} когредиентная свертка\par
\noindent\emph{Ukrajinsky / Ukrainian:} когредієнтне згортання\par
\smallskip\noindent\emph{Razlog:} Kogredientno svijanje v rabotě 4, §8, naziva operaciju medžu kogredientnymi redami $S^\rho,T^\rho$. Pridevnik ne znači samo sposobnost sojedinjenja ili obyčno odgovaranje. Sravnjeni tekst upotrěbjaje kogredientna kontrakcija i razlikuje osnovnu kontrakciju defekta 1 od obče kogredientnoj kontrakcije. Slovnik sohranja arhivno svijanje i nynešnje kontrakcija kak dva izbora dlja toj semeji operacij. Ne směje izbrisati razliku medžu občeju i osnovnoju operacijeju, na ktoroj opiraje se teorema zvedenja.\par
\noindent\emph{Reason:} Kogredientno svijanje in Paper 4, §8 names an operation between cogredient rows $S^\rho,T^\rho$. The adjective means more than compatibility or ordinary correspondence. The compared text uses kogredientna kontrakcija and distinguishes a basic contraction of defect 1 from a general cogredient contraction. The dictionary preserves archival svijanje and present kontrakcija as two choices for this family of operations. It must also preserve the general/basic distinction on which the reduction theorem depends.\par
\noindent{\footnotesize Izvorne zapisy / Original records: HIST-g037-003.}\par
\smallskip\noindent\emph{Sravnjenje s tekstom:} Sravnjeni tekst raboty 4 imaje kontrakcija na městě arhivnogo svijanje. On razlikuje obče kogredientne kontrakcije od osnovnyh kontrakcij defekta 1; zvedenje na osnovne jest silnějše tvrdženje od samogo zvedenja na kogredientne.\par
\noindent\emph{Textual comparison:} The compared Paper 4 text has kontrakcija where the record has svijanje. It distinguishes general cogredient contractions from basic contractions of defect 1; reduction to basic contractions is stronger than reduction merely to cogredient ones.\par
\noindent{\footnotesize TEXT-COMP-159; pasaže / passages: 4321, 4649, 4653; \texttt{CORPUS\_FORM\_COMPARISONS.json}.}\par
\end{minipage}\par
\par\medskip\Needspace{7\baselineskip}
\noindent\begin{minipage}{\linewidth}
\hypertarget{TERM-1007}{}
\noindent\textbf{kogredientno svijanje s defektom $\rho$}\quad{\footnotesize\texttt{TERM-1007}}\par
\noindent\emph{Iztočnik / Source:} kogrediente Faltung vom Defekt $\rho$\par
\noindent\emph{Rusky / Russian:} когредиентная свертка с дефектом $\rho$\par
\noindent\emph{Ukrajinsky / Ukrainian:} когредієнтне згортання з дефектом $\rho$\par
\smallskip\noindent\emph{Razlog:} Kogredientno svijanje s defektom $\rho$ prodolžuje semeju terminov iz §2 pri specialnoj konstrukciji raboty 4, §5. Defekt jest parametr operacije, ne ocěna nesovršenosti. Sravnjeni tekst nazyvaje elementarne kovarianty rezultatami kogredientnoj kontrakcije s defektom $\rho$. Zato potrěbno čuvati i pridevnik, i simbol parametra. Svijanje jest arhivny imenik, kontrakcija nyněšnji; ni jedna forma ne směje byti upotrěbjena kak opravdanje za gubitok ograničenja na konkretny vid operacije.\par
\noindent\emph{Reason:} Kogredientno svijanje s defektom $\rho$ continues the folding/defect family from §2 for the special construction in Paper 4, §5. Defect is an operation parameter, not an assessment of imperfection. The compared text calls the elementary covariants results of cogredient contraction with defect $\rho$. Both adjective and parameter must therefore remain. Svijanje is the archival noun and kontrakcija the present one; neither licenses dropping the restriction to the specified operation.\par
\noindent{\footnotesize Izvorne zapisy / Original records: HIST-g034-011.}\par
\smallskip\noindent\emph{Sravnjenje s tekstom:} Sravnjeni tekst raboty 4 imaje kontrakcija na městě arhivnogo svijanje. On razlikuje obče kogredientne kontrakcije od osnovnyh kontrakcij defekta 1; zvedenje na osnovne jest silnějše tvrdženje od samogo zvedenja na kogredientne.\par
\noindent\emph{Textual comparison:} The compared Paper 4 text has kontrakcija where the record has svijanje. It distinguishes general cogredient contractions from basic contractions of defect 1; reduction to basic contractions is stronger than reduction merely to cogredient ones.\par
\noindent{\footnotesize TEXT-COMP-159; pasaže / passages: 4321, 4649, 4653; \texttt{CORPUS\_FORM\_COMPARISONS.json}.}\par
\end{minipage}\par
\par\medskip\Needspace{7\baselineskip}
\noindent\begin{minipage}{\linewidth}
\hypertarget{TERM-1008}{}
\noindent\textbf{kogredientna transformacija}\quad{\footnotesize\texttt{TERM-1008}}\par
\noindent\emph{Iztočnik / Source:} kogrediente Transformation\par
\noindent\emph{Rusky / Russian:} когредиентное преобразование\par
\noindent\emph{Ukrajinsky / Ukrainian:} когредієнтне перетворення\par
\smallskip\noindent\emph{Razlog:} Kogredientna transformacija v rabotě 31, §4, prodolžuje semeju iz ranějšej teorije invariantov. Konkretno, prěhod medžu bazami měnja diskriminantny determinant kvadratom determinanta transformacije. Sravnjeni odlomok kaže, že koeficienty slěda imajut te same linearne relacije, i zato determinant dobyvaje toj samy faktor. Obča fraza odgovarjajuča proměna ne by izrekla toj transformacijny zakon. Izbor ostaje otvorjen jezykovoj ocěně; kontinuitet s ranějšim slovom ne jest samostojna atestacija.\par
\noindent\emph{Reason:} Kogredientna transformacija in Paper 31, §4 continues the earlier invariant-theory family. Here a basis change alters the discriminant determinant by the square of the transformation determinant. The compared passage says that the trace coefficients satisfy the same linear relations, so their determinant acquires that same factor. A vague corresponding change would not express this transformation law. The choice remains open to linguistic appraisal; continuity with an earlier term is not independent attestation.\par
\noindent{\footnotesize Izvorne zapisy / Original records: HIST-g137-005.}\par
\smallskip\noindent\emph{Sravnjenje s tekstom:} Pri prěhodu medžu bazami determinant měnjaje se kvadratom determinanta transformacije. Sravnjeni odlomok prěnosi te same linearne relacije na koeficienty slěda; zato kogredientna transformacija znači opreděleny zakon, ne obče odgovaranje.\par
\noindent\emph{Textual comparison:} On changing basis the determinant changes by the square of the transformation determinant. The compared passage transfers the same linear relations to the trace coefficients; cogredient transformation names a particular law, not a general correspondence.\par
\noindent{\footnotesize TEXT-COMP-160; pasaže / passages: 18102; \texttt{CORPUS\_FORM\_COMPARISONS.json}.}\par
\end{minipage}\par
\par\medskip\Needspace{7\baselineskip}
\noindent\begin{minipage}{\linewidth}
\hypertarget{TERM-1009}{}
\noindent\textbf{kogredientne varijable}\quad{\footnotesize\texttt{TERM-1009}}\par
\noindent\emph{Iztočnik / Source:} kogrediente Variabeln\par
\noindent\emph{Rusky / Russian:} когредиентные переменные\par
\noindent\emph{Ukrajinsky / Ukrainian:} когредієнтні змінні\par
\smallskip\noindent\emph{Razlog:} Kogredientne varijable v rabotě 8 čuvaje klasičny pridevnik teorije invariantov i njegove veze s transformacijami. Motivacija naměrno ne raztvarjaje jeho v obče pojasnjenje, ktore by skrylo disciplinarny pojem. Varijable jest medžunarodny imenik; v drugom zapisu za tu semeju izbrano jest prěměnne. Razlika imenikov ne jest razlika matematičnoj relacije. Za čitatelja bez znanja staršej terminologije treba podati transformacijno objasnjenje, kakovo korpus dava v rabotě 12; sam zaimstvovany pridevnik ne jest polna definicija.\par
\noindent\emph{Reason:} Kogredientne varijable in Paper 8 retains the classical invariant-theory adjective and its connection with transformations. The rationale deliberately avoids dissolving it into a loose paraphrase that would conceal the disciplinary notion. Varijable is the international noun; another record in the same family chooses prěměnne. That noun variation is not a difference in the mathematical relation. Readers unfamiliar with the older terminology need a transformation-based explanation such as the corpus gives in Paper 12; the borrowed adjective alone is not a definition.\par
\noindent{\footnotesize Izvorne zapisy / Original records: HIST-g057-012.}\par
\smallskip\noindent\emph{Sravnjenje s tekstom:} Rabota 12 pojasnjaje kogredientnost tymi samymi linearnymi transformacijami. Rabota 2 razlikuje kogredientne i kontragredientne redy v formulah, a rabota 23 koristi kogredientnost za oprěděljenje grupy. Poslědnja pri tom ne tvrdi, že obči dokaz navedenej redukcije už jest znany.\par
\noindent\emph{Textual comparison:} Paper 12 explains cogredience through the same linear transformations. Paper 2 distinguishes cogredient and contragredient rows in formulas, and Paper 23 uses cogredience to define the group. The latter does not claim that a general proof of the stated reduction is already known.\par
\noindent{\footnotesize TEXT-COMP-157; pasaže / passages: 400, 402, 9074, 9146, 15205; \texttt{CORPUS\_FORM\_COMPARISONS.json}.}\par
\end{minipage}\par
\par\medskip\Needspace{7\baselineskip}
\noindent\begin{minipage}{\linewidth}
\hypertarget{TERM-1010}{}
\noindent\textbf{kogredientne prěměnne}\quad{\footnotesize\texttt{TERM-1010}}\par
\noindent\emph{Iztočnik / Source:} kogrediente Variablen\par
\noindent\emph{Rusky / Russian:} когредиентные переменные\par
\noindent\emph{Ukrajinsky / Ukrainian:} когредієнтні змінні\par
\smallskip\noindent\emph{Razlog:} Kogredientne prěměnne v rabotě 2, §2, sohranja paru s kontragredientny. Odlomok razlikuje formy s dvoma kogredientnymi redami $s_x^m t_x^n$ od form s kontragredientnymi redami. Zato ne možno zaměniti pridevnik samym sočetanim povezane prěměnne: protivnost upravjaje različnymi razvitjami. Prěměnne dava slovjansky imenik, dokolě varijable v drugom zapisu jest zaimstvovana alternativa. Dokumentacija ne pretvarjaje dve imenni formy v dva različno definovana vida kogredientnosti.\par
\noindent\emph{Reason:} Kogredientne prěměnne in Paper 2, §2 preserves the pairing with kontragredientny. The passage distinguishes forms with two cogredient rows $s_x^m t_x^n$ from forms with contragredient rows. Related variables would lose an opposition governing different expansions. Prěměnne supplies the Slavic noun, whereas varijable elsewhere is a borrowed alternative. The two noun forms do not establish two differently defined types of cogredience.\par
\noindent{\footnotesize Izvorne zapisy / Original records: HIST-g004-004.}\par
\smallskip\noindent\emph{Sravnjenje s tekstom:} Rabota 12 pojasnjaje kogredientnost tymi samymi linearnymi transformacijami. Rabota 2 razlikuje kogredientne i kontragredientne redy v formulah, a rabota 23 koristi kogredientnost za oprěděljenje grupy. Poslědnja pri tom ne tvrdi, že obči dokaz navedenej redukcije už jest znany.\par
\noindent\emph{Textual comparison:} Paper 12 explains cogredience through the same linear transformations. Paper 2 distinguishes cogredient and contragredient rows in formulas, and Paper 23 uses cogredience to define the group. The latter does not claim that a general proof of the stated reduction is already known.\par
\noindent{\footnotesize TEXT-COMP-157; pasaže / passages: 400, 402, 9074, 9146, 15205; \texttt{CORPUS\_FORM\_COMPARISONS.json}.}\par
\end{minipage}\par
\par\medskip\Needspace{7\baselineskip}
\noindent\begin{minipage}{\linewidth}
\hypertarget{TERM-1011}{}
\noindent\textbf{kogredientnost}\quad{\footnotesize\texttt{TERM-1011}}\par
\noindent\emph{Iztočnik / Source:} Kogredienz\par
\noindent\emph{Rusky / Russian:} когредиентность\par
\noindent\emph{Ukrajinsky / Ukrainian:} когредієнтність\par
\smallskip\noindent\emph{Razlog:} Kogredientnost tvori abstraktny imenik od kogredientny s koncom -nost. Razlog v rabotě 23 predpočitaje razpoznavajemy i legko iskomy termin vměsto proměnnyh opisov. V sravnjenom odlomku trěbovanje kogredientnosti veličin $\psi$ k $dx$ inducuje transformaciju jih koeficientov i tym oprědělja grupu. To je konkretna rolja uslovja. Ona ne daje prava odstraniti priznanie teksta, že občy pristup k opisanomu teoremu redukcije ješče nedostaje; stabilno nazvanje ne jest dokaz rezultata.\par
\noindent\emph{Reason:} Kogredientnost forms the abstract noun from kogredientny with -nost. The Paper 23 rationale favours a recognizable, searchable term over variable paraphrases. In the compared passage, requiring the quantities $\psi$ to be cogredient with $dx$ induces transformation of their coefficients and thereby defines the group. That is the condition’s specific role. It does not remove the text’s statement that a general approach to the described reduction theorem is still lacking: stable terminology is not a proof.\par
\noindent{\footnotesize Izvorne zapisy / Original records: HIST-g094-018.}\par
\smallskip\noindent\emph{Sravnjenje s tekstom:} Rabota 12 pojasnjaje kogredientnost tymi samymi linearnymi transformacijami. Rabota 2 razlikuje kogredientne i kontragredientne redy v formulah, a rabota 23 koristi kogredientnost za oprěděljenje grupy. Poslědnja pri tom ne tvrdi, že obči dokaz navedenej redukcije už jest znany.\par
\noindent\emph{Textual comparison:} Paper 12 explains cogredience through the same linear transformations. Paper 2 distinguishes cogredient and contragredient rows in formulas, and Paper 23 uses cogredience to define the group. The latter does not claim that a general proof of the stated reduction is already known.\par
\noindent{\footnotesize TEXT-COMP-157; pasaže / passages: 400, 402, 9074, 9146, 15205; \texttt{CORPUS\_FORM\_COMPARISONS.json}.}\par
\end{minipage}\par
\par\medskip\Needspace{7\baselineskip}
\noindent\begin{minipage}{\linewidth}
\hypertarget{TERM-1012}{}
\noindent\textbf{kombinatoričny/progresivny/regresivny/směšany produkt}\quad{\footnotesize\texttt{TERM-1012}}\par
\noindent\emph{Iztočnik / Source:} kombinatorisches/progressives/regressives/gemischtes Produkt\par
\noindent\emph{Rusky / Russian:} комбинаторное/прогрессивное/регрессивное/смешанное произведение\par
\noindent\emph{Ukrajinsky / Ukrainian:} комбінаторний/прогресивний/регресивний/мішаний добуток\par
\smallskip\noindent\emph{Razlog:} Kombinatoričny, progresivny, regresivny i směšany produkt sohranjajut Grassmannovo razděljenje produktov v uvodě raboty 4. Medžunarodne koreni prvih trěh pridevnikov sojedinjajut se s medžuslovjanskim směšany i občim imenikom produkt. Tekst svazuje progresivny produkt s matricom samyh redov, regresivny s prirědženymi redami, a směšany s konstrukcijeju novyh redov prěz substituciju. Zato ne ide o stilistične variante jednogo umnoženja. Osoblivo matričny produkt opisan jest kak směšany produkt stupnja 0, ne kak cela kategorija bez ograničenja.\par
\noindent\emph{Reason:} Kombinatoričny, progresivny, regresivny and směšany produkt preserve Grassmann’s product taxonomy in the introduction to Paper 4. The international stems of the first three adjectives combine with Slavic směšany and the shared noun produkt. The text relates progressive product to the matrix of the rows themselves, regressive product to associated rows, and mixed product to forming new rows by substitution. These are not stylistic variants for one multiplication. In particular, matrix product is identified with mixed product of degree 0, not the whole category without qualification.\par
\noindent{\footnotesize Izvorne zapisy / Original records: HIST-g029-016.}\par
\smallskip\noindent\emph{Sravnjenje s tekstom:} Odlomok svazuje progresivny, regresivny i směšany produkt s različnymi konstrukcijami matric i redov. Vlastny matričny produkt raboty sravnjaje se samo so směšanym produktom stupnja 0.\par
\noindent\emph{Textual comparison:} The passage associates progressive, regressive and mixed products with different matrix and row constructions. The paper’s matrix product is compared specifically with mixed product of degree 0.\par
\noindent{\footnotesize TEXT-COMP-161; pasaže / passages: 3709; \texttt{CORPUS\_FORM\_COMPARISONS.json}.}\par
\end{minipage}\par
\par\medskip\Needspace{7\baselineskip}
\noindent\begin{minipage}{\linewidth}
\hypertarget{TERM-1013}{}
\noindent\textbf{komutativno povezano s ...}\quad{\footnotesize\texttt{TERM-1013}}\par
\noindent\emph{Iztočnik / Source:} kommutativ mit ... verbunden\par
\noindent\emph{Rusky / Russian:} коммутативно связанное с ...\par
\noindent\emph{Ukrajinsky / Ukrainian:} комутативно пов'язане з ...\par
\smallskip\noindent\emph{Razlog:} Komutativno povezano s v rabotě 34, §§18–19, prodolžuje tehničnu frazu iz §15. Nyněšnji odlomok trěbuje bimoduly s dějstvom $P$ sprava i prstena slěva, ravnost $(am)\rho=a(m\rho)=(a\rho)m$ i jediničny operator iz $P$. Podmoduly takože morajut dopuščati množenje s $P$. Obča reč povezano ne by sama prědala ty uslovja; ona potrěbuje definovanu tehničnu upotrěbu. Sočetanje ne tvrdi, že sam razgledany prsten jest komutativny.\par
\noindent\emph{Reason:} Komutativno povezano s in Paper 34, §§18–19 continues the technical phrase from §15. The present passage requires bimodules with $P$ acting on the right and the ring on the left, the equality $(am)\rho=a(m\rho)=(a\rho)m$, and the unit of $P$ acting as the identity. Submodules must also admit multiplication by $P$. Connected alone would not convey these conditions; the phrase needs its defined technical use. It does not assert that the ring itself is commutative.\par
\noindent{\footnotesize Izvorne zapisy / Original records: HIST-g196-009.}\par
\smallskip\noindent\emph{Sravnjenje s tekstom:} §15 i pozdnějši odlomok raboty 34 izrično svazujut operatorno sohranjenje s $(am)\rho=a(m\rho)=(a\rho)m$. Druhy odlomok dodaje uslovje jediničnogo operatora i stabilnost podmodulov po $P$. To ne jest tvrdženje komutativnosti vsego prstena.\par
\noindent\emph{Textual comparison:} Section 15 and a later Paper 34 passage explicitly relate preservation of the operator action to $(am)\rho=a(m\rho)=(a\rho)m$. The latter adds the identity-operator condition and stability of submodules under $P$. This is not a statement that the entire ring is commutative.\par
\noindent{\footnotesize TEXT-COMP-162; pasaže / passages: 21441, 21443, 21447, 21449, 21745, 21747, 21749; \texttt{CORPUS\_FORM\_COMPARISONS.json}.}\par
\end{minipage}\par
\par\medskip\Needspace{7\baselineskip}
\noindent\begin{minipage}{\linewidth}
\hypertarget{TERM-1014}{}
\noindent\textbf{komutativno povezany s $P$}\quad{\footnotesize\texttt{TERM-1014}}\par
\noindent\emph{Iztočnik / Source:} kommutativ mit $P$ verbunden\par
\noindent\emph{Rusky / Russian:} коммутативно связанный с $P$\par
\noindent\emph{Ukrajinsky / Ukrainian:} комутативно пов'язаний з $P$\par
\smallskip\noindent\emph{Razlog:} Komutativno povezany s $P$ v rabotě 34, §15, zadržava kompaktnu konstrukciju istorijnogo iztočnika. Povezanost znači konkretno soglasje s komutativnoju oblastju multiplikatorov: prsten-nositelj prědstavjenja imaje $P$ v centru, a prědstavjenje mora sohranjati takože množenje s $\rho\in P$. Zapis zato otklanja neoprěděljen govor o skalarnoj linejnosti bez opisa dějstva. Frazu možno objasniti formulami, ale ne razmyti do običnoj svezi medžu objektami.\par
\noindent\emph{Reason:} Komutativno povezany s $P$ in Paper 34, §15 retains the compact historical construction. The connection means specific compatibility with a commutative multiplier domain: the representation’s carrier division ring has $P$ in its centre, and the representation must preserve multiplication by $\rho\in P$. The rationale therefore rejects vague scalar-linearity language without specifying the action. Formulas can explain the phrase; an ordinary connection between objects cannot replace it.\par
\noindent{\footnotesize Izvorne zapisy / Original records: HIST-g193-006.}\par
\smallskip\noindent\emph{Sravnjenje s tekstom:} §15 i pozdnějši odlomok raboty 34 izrično svazujut operatorno sohranjenje s $(am)\rho=a(m\rho)=(a\rho)m$. Druhy odlomok dodaje uslovje jediničnogo operatora i stabilnost podmodulov po $P$. To ne jest tvrdženje komutativnosti vsego prstena.\par
\noindent\emph{Textual comparison:} Section 15 and a later Paper 34 passage explicitly relate preservation of the operator action to $(am)\rho=a(m\rho)=(a\rho)m$. The latter adds the identity-operator condition and stability of submodules under $P$. This is not a statement that the entire ring is commutative.\par
\noindent{\footnotesize TEXT-COMP-162; pasaže / passages: 21441, 21443, 21447, 21449, 21745, 21747, 21749; \texttt{CORPUS\_FORM\_COMPARISONS.json}.}\par
\end{minipage}\par
\par\medskip\Needspace{7\baselineskip}
\noindent\begin{minipage}{\linewidth}
\hypertarget{TERM-1015}{}
\noindent\textbf{komutativno polje $K$}\quad{\footnotesize\texttt{TERM-1015}}\par
\noindent\emph{Iztočnik / Source:} kommutativer Körper $K$\par
\noindent\emph{Rusky / Russian:} коммутативное поле $K$\par
\noindent\emph{Ukrajinsky / Ukrainian:} комутативне поле $K$\par
\smallskip\noindent\emph{Razlog:} Komutativno polje $K$ v rabotě 34, §6, izbira polje, kogda komutativnost izrično postavjena. V tom samom paragrafu ranějši moduly razgledajut se nad ne potrěbno komutativnym tělom, a definicija hiperkompleksnogo sistema uže trěbuje komutativno polje. To različenje podpiraje lokalnu motivaciju. Ale zapis komutativno tělo $P$ iz §20 pokazuje, že razděljenje polje/tělo ne bylo vsestranno provedeno. Slovnik mora pokazati oba izbora i ne pripisyvati korpusu polnu doslědnost, ktorej sravnjenje ne ustanavjaje.\par
\noindent\emph{Reason:} Komutativno polje $K$ in Paper 34, §6 chooses polje when commutativity is explicit. Earlier in the same section, modules are considered over a division ring that need not be commutative, whereas the hypercomplex-system definition requires a commutative field. This supports the local rationale. Yet komutativno tělo $P$ in §20 shows that the polje/tělo allocation was not uniform. The dictionary must expose both choices rather than attribute complete consistency to the corpus.\par
\noindent{\footnotesize Izvorne zapisy / Original records: HIST-g184-002.}\par
\smallskip\noindent\emph{Sravnjenje s tekstom:} V §6 raboty 34 nekomutativnost dopušča se pri tělu, a definicija hiperkompleksnogo sistema upotrěbjaje komutativno polje. Ale §20 govori o komutativnom tělu. Taka konkurencija form jest v samom sravnjenom tekstě, ne samo v spisku terminov.\par
\noindent\emph{Textual comparison:} In Paper 34, §6 permits a noncommutative division ring under tělo, then uses komutativno polje in defining a hypercomplex system. Section 20, however, says komutativno tělo. The competition occurs in the compared text itself, not merely in the term list.\par
\noindent{\footnotesize TEXT-COMP-163; pasaže / passages: 20221, 20268, 20271, 20272, 21794; \texttt{CORPUS\_FORM\_COMPARISONS.json}.}\par
\end{minipage}\par
\par\medskip\Needspace{7\baselineskip}
\noindent\begin{minipage}{\linewidth}
\hypertarget{TERM-1016}{}
\noindent\textbf{komutativno tělo $P$}\quad{\footnotesize\texttt{TERM-1016}}\par
\noindent\emph{Iztočnik / Source:} kommutativer Körper $P$\par
\noindent\emph{Rusky / Russian:} коммутативное поле $P$\par
\noindent\emph{Ukrajinsky / Ukrainian:} комутативне поле $P$\par
\smallskip\noindent\emph{Razlog:} Komutativno tělo $P$ v rabotě 34, §20, jest sohranjena alternativa k polje. Razlog opiraje se na kontinuitet projektnogo tělo i prědpolagaje največšu medžuslovjansku razpoznavajemost; čitateljsko izpytanje za tu ocěnu ne jest podano. Rusko i ukrajinsko sravnjenje upotrěbjajut pole/polje v komutativnom smyslu. Nyněšnji tekst stvarno imaje komutativno tělo, dokolě §6 imaje komutativno polje. Zato pridevnik nosi javno uslovje, a konkurencija imenikov ostaje pytanjem doslědnosti i jezykovej ocěny, ne različije dvoh komutativnyh struktur samo po nazvanju.\par
\noindent\emph{Reason:} Komutativno tělo $P$ in Paper 34, §20 is a retained alternative to polje. The rationale invokes continuity with the project’s tělo and claims maximum cross-Slavic recognizability, but supplies no reader experiment for that assessment. Russian and Ukrainian use their field nouns in the commutative sense. The present text indeed has komutativno tělo, while §6 has komutativno polje. The adjective makes the condition explicit; the competing nouns remain an issue of consistency and linguistic appraisal, not two commutative structures distinguished merely by name.\par
\noindent{\footnotesize Izvorne zapisy / Original records: HIST-g198-002.}\par
\smallskip\noindent\emph{Sravnjenje s tekstom:} V §6 raboty 34 nekomutativnost dopušča se pri tělu, a definicija hiperkompleksnogo sistema upotrěbjaje komutativno polje. Ale §20 govori o komutativnom tělu. Taka konkurencija form jest v samom sravnjenom tekstě, ne samo v spisku terminov.\par
\noindent\emph{Textual comparison:} In Paper 34, §6 permits a noncommutative division ring under tělo, then uses komutativno polje in defining a hypercomplex system. Section 20, however, says komutativno tělo. The competition occurs in the compared text itself, not merely in the term list.\par
\noindent{\footnotesize TEXT-COMP-163; pasaže / passages: 20221, 20268, 20271, 20272, 21794; \texttt{CORPUS\_FORM\_COMPARISONS.json}.}\par
\end{minipage}\par
\par\medskip\Needspace{7\baselineskip}
\noindent\begin{minipage}{\linewidth}
\hypertarget{TERM-1017}{}
\noindent\textbf{komutativny zakon / zam?nljivost}\quad{\footnotesize\texttt{TERM-1017}}\par
\noindent\emph{Iztočnik / Source:} kommutatives Gesetz / Vertauschbarkeit\par
\noindent\emph{Rusky / Russian:} [Bukvy ne sut zachovane / Lettering not recoverable]\par
\noindent\emph{Ukrajinsky / Ukrainian:} [Bukvy ne sut zachovane / Lettering not recoverable]\par
\smallskip\noindent\emph{Razlog:} Komutativny zakon / zam?nljivost v rabotě 17 razlikuje algebraično uslovje od obyčnogo prěnumerovanja. Sravnjeni odlomok koristi ravnost produktov klasov ostatkov, da by potrěbna grupa mogla byti prědstavjena kak aditivna suma. V tekstě čita se zaměnljivost, ale to ne vrača poškodžene bukvy samogo arhivnogo zapisa. Ruska i ukrajinska polja tam sut nečitateljne i ne mogut služiti kak dokazane leksične alternativy.\par
\noindent\emph{Reason:} Komutativny zakon / zam?nljivost in Paper 17 distinguishes an algebraic condition from ordinary reindexing. The compared passage uses equality of products of residue classes so that the relevant constituents lie in the required group and give its additive decomposition. The text reads zaměnljivost, but this does not restore the damaged lettering of the archived record itself. Its Russian and Ukrainian fields are illegible and cannot serve as evidenced lexical alternatives.\par
\noindent{\footnotesize Izvorne zapisy / Original records: HIST-g079-041.}\par
\smallskip\noindent\emph{Sravnjenje s tekstom:} Odlomok raboty 17 koristi zaměnljivost konkretnyh klasov ostatkov za aditivno razloženje. Čitatelj može najdti čitateljnu formu zaměnljivost v korpusu; poškodženy arhivny zapis i njegove nečitateljne sravnjenja ostajut oddělno označene.\par
\noindent\emph{Textual comparison:} The Paper 17 passage uses commutation of specified residue classes to obtain an additive decomposition. The corpus supplies a legible occurrence of zaměnljivost; the damaged archived form and its illegible comparators remain separately marked.\par
\noindent{\footnotesize TEXT-COMP-164; pasaže / passages: 11527; \texttt{CORPUS\_FORM\_COMPARISONS.json}.}\par
\end{minipage}\par
\par\medskip\Needspace{7\baselineskip}
\noindent\begin{minipage}{\linewidth}
\hypertarget{TERM-1018}{}
\noindent\textbf{komplement}\quad{\footnotesize\texttt{TERM-1018}}\par
\noindent\emph{Iztočnik / Source:} Komplement\par
\noindent\emph{Rusky / Russian:} дополнение\par
\noindent\emph{Ukrajinsky / Ukrainian:} комплемент\par
\smallskip\noindent\emph{Razlog:} Komplement v rabotě 19 jest najmenše obče mnogokratno vsih ostalyh komponentov danogo prědstavjenja. Nyněšnja Definicija I izrično dava tu konstrukciju. Medžunarodny imenik izbrany jest vměsto ruskogo dopolnenije; ukrajinsky zapis takože imaje komplement. Kontinuitet slova važen jest za pozdnějše kriterije, ale pojem ostaje zavisny od razloženja: ne znači množinovo dopolnjenje vsih elementov mimo jednogo ideala. Definicija, ne sam razpoznavajemy zvuk, predotvrača takvo pogrěšno čitanje.\par
\noindent\emph{Reason:} Komplement in Paper 19 is the least common multiple of all the other components of a given representation, as the present Definition I explicitly states. The international noun is chosen instead of the Russian supplement word; the Ukrainian record also has komplement. Continuity matters to the later criteria, but the object remains relative to the decomposition. It is not the set-theoretic complement of all elements outside one ideal. The definition, not merely the recognizable noun, prevents that misreading.\par
\noindent{\footnotesize Izvorne zapisy / Original records: HIST-g081-032, HIST-g082-032.}\par
\smallskip\noindent\emph{Sravnjenje s tekstom:} Komplement danogo komponenta jest najmenše obče mnogokratno ostalyh, ne množina elementov mimo njega. V rabotě 30 ortogonalnost opisana jest produktami idealov i idempotentnostju, a direktny prěsěk jest imenovany v odkazu na Prüfera. Te svezane upotrěby ne možno svesti k geometričnomu ugolu ili množinovomu dopolnjenju.\par
\noindent\emph{Textual comparison:} A component’s complement is the least common multiple of the others, not the set of elements outside it. Paper 30 describes orthogonality by ideal products and idempotency, while direct intersection appears in the reference to Prüfer. These related uses cannot be reduced to geometric angle or set complement.\par
\noindent{\footnotesize TEXT-COMP-165; pasaže / passages: 12570, 12358, 12359, 12360, 12362, 12363, 12364, 16528, 16536, 16553; \texttt{CORPUS\_FORM\_COMPARISONS.json}.}\par
\end{minipage}\par
\par\medskip\Needspace{7\baselineskip}
\noindent\begin{minipage}{\linewidth}
\hypertarget{TERM-1019}{}
\noindent\textbf{komplement / ortogonalne relacije / direktny prěsěk}\quad{\footnotesize\texttt{TERM-1019}}\par
\noindent\emph{Iztočnik / Source:} Komplement / Orthogonalitätsrelationen / direkter Durchschnitt\par
\noindent\emph{Rusky / Russian:} комплемент / ортогональные соотношения / прямое пересечение\par
\noindent\emph{Ukrajinsky / Ukrainian:} комплемент / ортогональні співвідношення / прямий перетин\par
\smallskip\noindent\emph{Razlog:} Komplement / ortogonalne relacije / direktny prěsěk sohranja tri svazane pojmy raboty 30, §4. Komplement jest izrično imenovany ideal; „ortogonalne relacije” sut ravnosti $\mathfrak a_i\mathfrak a_k=(0)$ za $i\ne k$ i $\mathfrak a_i^2=\mathfrak a_i$. Ne vveden jest skalarni produkt ili geometričny ugol. Direktny prěsěk ostaje v odkazu na Prüferovu formulaciju svezi s direktnoju sumoju. Zato trěba čuvati navodky, formuly i bibliografičny kontekst, ne zaměnjati istorijnu trojicu jednym sovrěmennym sloganom.\par
\noindent\emph{Reason:} Komplement / ortogonalne relacije / direktny prěsěk preserves three linked notions in Paper 30, §4. The complement is a named ideal; the “orthogonality relations” are $\mathfrak a_i\mathfrak a_k=(0)$ for $i\ne k$ and $\mathfrak a_i^2=\mathfrak a_i$. No inner product or geometric angle is introduced. Direct intersection remains in the reference to Prüfer’s formulation of its connection with direct sum. Quotation marks, formulas and bibliographic context must remain, rather than replacing the historical trio with a modern slogan.\par
\noindent{\footnotesize Izvorne zapisy / Original records: HIST-g111-006.}\par
\smallskip\noindent\emph{Sravnjenje s tekstom:} Komplement danogo komponenta jest najmenše obče mnogokratno ostalyh, ne množina elementov mimo njega. V rabotě 30 ortogonalnost opisana jest produktami idealov i idempotentnostju, a direktny prěsěk jest imenovany v odkazu na Prüfera. Te svezane upotrěby ne možno svesti k geometričnomu ugolu ili množinovomu dopolnjenju.\par
\noindent\emph{Textual comparison:} A component’s complement is the least common multiple of the others, not the set of elements outside it. Paper 30 describes orthogonality by ideal products and idempotency, while direct intersection appears in the reference to Prüfer. These related uses cannot be reduced to geometric angle or set complement.\par
\noindent{\footnotesize TEXT-COMP-165; pasaže / passages: 12570, 12358, 12359, 12360, 12362, 12363, 12364, 16528, 16536, 16553; \texttt{CORPUS\_FORM\_COMPARISONS.json}.}\par
\end{minipage}\par
\par\medskip\Needspace{7\baselineskip}
\noindent\begin{minipage}{\linewidth}
\hypertarget{TERM-1020}{}
\noindent\textbf{komplement $Q$}\quad{\footnotesize\texttt{TERM-1020}}\par
\noindent\emph{Iztočnik / Source:} Komplement von $Q$\par
\noindent\emph{Rusky / Russian:} дополнение к $Q$\par
\noindent\emph{Ukrajinsky / Ukrainian:} доповнення до $Q$\par
\smallskip\noindent\emph{Razlog:} Komplement $Q$ v rabotě 18 čuvaje Noetherino idealno nazvanje i medžunarodny imenik. V sravnjenom odlomku $\mathfrak M=[\mathfrak Q,\mathfrak L]$, a $\mathfrak L$ jest komplement $\mathfrak Q$ v primarnom razloženju; potom nastupaje norma $\mathfrak L$ po $\mathfrak Q$. Zato kratka fraza ne znači vsih elementov mimo $Q$. Arhivno latinsko Q predstavlja matematičnu oznaku, v sravnjenom tekstě gotičnu $\mathfrak Q$. Čitateljsko objasnjenje može razkryti konstrukciju, ale ne može zaměniti imenovany ideal množinovym dopolnjenjem.\par
\noindent\emph{Reason:} Komplement $Q$ in Paper 18 retains Noether’s ideal-theoretic designation and the international noun. In the compared passage $\mathfrak M=[\mathfrak Q,\mathfrak L]$, with $\mathfrak L$ the complement of $\mathfrak Q$ in a primary decomposition; the norm of $\mathfrak L$ relative to $\mathfrak Q$ then occurs. The compact phrase does not mean all elements outside $Q$. Archival Latin Q represents a mathematical identifier, printed in fraktur in the compared text. An explanatory gloss may unpack the construction, but cannot replace the named ideal by a set complement.\par
\noindent{\footnotesize Izvorne zapisy / Original records: HIST-g080-012.}\par
\smallskip\noindent\emph{Sravnjenje s tekstom:} Komplement danogo komponenta jest najmenše obče mnogokratno ostalyh, ne množina elementov mimo njega. V rabotě 30 ortogonalnost opisana jest produktami idealov i idempotentnostju, a direktny prěsěk jest imenovany v odkazu na Prüfera. Te svezane upotrěby ne možno svesti k geometričnomu ugolu ili množinovomu dopolnjenju.\par
\noindent\emph{Textual comparison:} A component’s complement is the least common multiple of the others, not the set of elements outside it. Paper 30 describes orthogonality by ideal products and idempotency, while direct intersection appears in the reference to Prüfer. These related uses cannot be reduced to geometric angle or set complement.\par
\noindent{\footnotesize TEXT-COMP-165; pasaže / passages: 12570, 12358, 12359, 12360, 12362, 12363, 12364, 16528, 16536, 16553; \texttt{CORPUS\_FORM\_COMPARISONS.json}.}\par
\end{minipage}\par
\par\medskip\Needspace{7\baselineskip}
\noindent\begin{minipage}{\linewidth}
\hypertarget{TERM-1021}{}
\noindent\textbf{stepen komplementa / dělimy črez primarny ideal}\quad{\footnotesize\texttt{TERM-1021}}\par
\noindent\emph{Iztočnik / Source:} Komplementpotenz / durch ein Primärideal teilbar\par
\noindent\emph{Rusky / Russian:} степень дополнения / делима на примарный идеал\par
\noindent\emph{Ukrajinsky / Ukrainian:} степінь комплемента / подільний на первинний ідеал\par
\smallskip\noindent\emph{Razlog:} Stepen komplementa / dělimy črez primarny ideal sohranja i potenciju, i směr děliteljnosti v rabotě 19. V zaključnoj harakterizaciji neizolovanyh primarnyh idealov někoja potencija vsakogo komplementa jest dělima prěz dany ideal; za izolovane take uslovje ne može nastati. Zato vsaky a někoja ne sut izmenljive s jednoj fiksovanoj potencijeju dlja vsih. Stepen tu jest algebraična potencija komplementa, ne jego dimenzija ili stepenj razširjenja.\par
\noindent\emph{Reason:} Stepen komplementa / dělimy črez primarny ideal retains both the power and the direction of divisibility in Paper 19. The concluding characterization says that some power of every complement is divisible by the non-isolated primary ideal; for isolated ideals this cannot occur. Every and some do not imply one fixed power for all complements. Stepen here means an algebraic power of the complement, not its dimension or an extension degree.\par
\noindent{\footnotesize Izvorne zapisy / Original records: HIST-g081-061, HIST-g082-061.}\par
\smallskip\noindent\emph{Sravnjenje s tekstom:} Harakterizacija v rabotě 19 ohranja razliku medžu někoju potencijeju vsakogo komplementa i nemožnostju take děliteljnosti za izolovane komponente. Izjava ne izbira jeden obči eksponent za vse komplementy.\par
\noindent\emph{Textual comparison:} The Paper 19 characterization distinguishes some power of every complement from the impossibility of such divisibility for isolated components. It does not select a single common exponent for all complements.\par
\noindent{\footnotesize TEXT-COMP-166; pasaže / passages: 13246, 13248, 13250, 13252; \texttt{CORPUS\_FORM\_COMPARISONS.json}.}\par
\end{minipage}\par
\par\medskip\Needspace{7\baselineskip}
\noindent\begin{minipage}{\linewidth}
\hypertarget{TERM-1022}{}
\noindent\textbf{komplementarne bazy / komplementarne moduly}\quad{\footnotesize\texttt{TERM-1022}}\par
\noindent\emph{Iztočnik / Source:} Komplementärbasen / Komplementärmoduln\par
\noindent\emph{Rusky / Russian:} комплементарные базы / комплементарные модули\par
\noindent\emph{Ukrajinsky / Ukrainian:} комплементарні бази / комплементарні модулі\par
\smallskip\noindent\emph{Razlog:} Komplementarne bazy / komplementarne moduly v rabotě 42 zadržava medžunarodny pridevnik, zajedničny sravniteljnym prevodam. Jegova rolja oprěděljena jest slědovoju dualnostju: $\operatorname{Sp}(a_j\bar a_h)=\delta_{jh}$. Formula daje pravilno umnoženje matričnyh jedinic v argumentu. Zato komplementarny ne znači samo dodatny i ne označuje avtomatično direktny sumand. Razpoznavajemost jest razlog izbora, ale točno značenje davaje svaz s formulami.\par
\noindent\emph{Reason:} Komplementarne bazy / komplementarne moduly in Paper 42 retains the international adjective shared by the comparator translations. Its role is specified by trace duality: $\operatorname{Sp}(a_j\bar a_h)=\delta_{jh}$. That identity yields the matrix-unit multiplication in the argument. Complementary therefore means more than additional and does not automatically mean a direct summand. Recognizability motivates the choice; the formulas determine its precise sense.\par
\noindent{\footnotesize Izvorne zapisy / Original records: HIST-g212-006.}\par
\smallskip\noindent\emph{Sravnjenje s tekstom:} Rabota 42 dava $\operatorname{Sp}(a_j\bar a_h)=\delta_{jh}$ kak definicijne ravnosti komplementarnyh baz i koristi ih za matrične jedinice. Takvo komplementarno uslovje ne jest samo tvrdženje o direktnoj sumi.\par
\noindent\emph{Textual comparison:} Paper 42 gives $\operatorname{Sp}(a_j\bar a_h)=\delta_{jh}$ as the defining equalities of complementary bases and uses them for matrix units. This complementary condition is not merely a direct-sum assertion.\par
\noindent{\footnotesize TEXT-COMP-167; pasaže / passages: 24500, 24501, 24503, 24505, 24507, 24555, 24621; \texttt{CORPUS\_FORM\_COMPARISONS.json}.}\par
\end{minipage}\par
\par\medskip\Needspace{7\baselineskip}
\noindent\begin{minipage}{\linewidth}
\hypertarget{TERM-1023}{}
\noindent\textbf{komplementarna baza / komplementarny modul}\quad{\footnotesize\texttt{TERM-1023}}\par
\noindent\emph{Iztočnik / Source:} Komplementärbasis / Komplementärmodul\par
\noindent\emph{Rusky / Russian:} комплементарная база / комплементарный модуль\par
\noindent\emph{Ukrajinsky / Ukrainian:} комплементарна база / комплементарний модуль\par
\smallskip\noindent\emph{Razlog:} Komplementarna baza / komplementarny modul v rabotě 43 čuvaje tehnično nazvanje centralno dlja sravnjenja s Dedekindom. Sravnjeni tekst daje komplementarny modul elementami, čije slědy produktov s elementami poredka sut cěle v osnovnom prstenu. To jest konkretno uslovje, ne obyčno dopolnjenje spiska elementov. Medžunarodny pridevnik izbrany jest radi toj preciznosti; njegova jedinstvena forma soglasuje se s baza i modul, ne vvedaje dva nepovezana smysla.\par
\noindent\emph{Reason:} Komplementarna baza / komplementarny modul in Paper 43 preserves the technical designation central to the Dedekind comparison. The compared text specifies the complementary module through elements whose products with order elements have integral traces in the base ring. This is a definite condition, not merely adding more elements to a list. The international adjective is chosen for that precision; its singular forms agree with baza and modul rather than introduce unrelated senses.\par
\noindent{\footnotesize Izvorne zapisy / Original records: HIST-g213-010.}\par
\smallskip\noindent\emph{Sravnjenje s tekstom:} Rabota 43 povezuje koeficienty komplementarnyh baz so slědom i karakterizuje modul $\mathfrak C$ cělostju $\operatorname{Sp}(co)$ dlja vsih $o$ iz poredka. To daje sodržanje pridevnika i svaz s Dedekindovoju diferentoju.\par
\noindent\emph{Textual comparison:} Paper 43 connects the coefficients of complementary bases with trace and characterizes $\mathfrak C$ by integrality of $\operatorname{Sp}(co)$ for all $o$ in the order. This supplies the adjective’s content and its connection with Dedekind’s different.\par
\noindent{\footnotesize TEXT-COMP-168; pasaže / passages: 24769, 25047, 25049, 25051, 25083, 25101, 25103, 25105; \texttt{CORPUS\_FORM\_COMPARISONS.json}.}\par
\end{minipage}\par
\par\medskip\Needspace{7\baselineskip}
\noindent\begin{minipage}{\linewidth}
\hypertarget{TERM-1024}{}
\noindent\textbf{komplementarny modul}\quad{\footnotesize\texttt{TERM-1024}}\par
\noindent\emph{Iztočnik / Source:} komplementärer Modul\par
\noindent\emph{Rusky / Russian:} комплементарный модуль\par
\noindent\emph{Ukrajinsky / Ukrainian:} комплементарний модуль\par
\smallskip\noindent\emph{Razlog:} Komplementarny modul v rabotě 14 jest izbor vměsto užših ili običnyh dopolnjajučih pridevnikov, navedenyh v ruskom i ukrajinskom sravnjenju. Tekst odkazuje na Henselov tehničny pojem, iz ktorego izvodi ideal razvětvenja, i pozdněje svazuje komplementarny modul s normalnoju bazoju i teoremoju Riemanna–Rocha. Zato čitatelj potrěbuje termin, ktory možno svezati s tými odkazami. Sama reč komplementarny ne znači, že konstrukcija sběgaje se s idealnym komplementom iz drugoj raboty.\par
\noindent\emph{Reason:} Komplementarny modul in Paper 14 is preferred to the looser supplementary adjectives named in the Russian and Ukrainian comparison. The text invokes Hensel’s technical notion to obtain the ramification ideal, then relates the complementary module to a normal basis and the Riemann–Roch theorem. Readers need a term they can connect with those references. The shared word complementary does not make this construction identical to the ideal complement in another paper.\par
\noindent{\footnotesize Izvorne zapisy / Original records: HIST-g076-010.}\par
\smallskip\noindent\emph{Sravnjenje s tekstom:} Rabota 14 odkazuje na Hensela, potom svazuje komplementarny modul s idealom razvětvenja, normalnoju bazoju i Riemannom–Rochom. Te odkazy opravdavajut tehnično, ne obyčno dopolnjajuče čitanje pridevnika.\par
\noindent\emph{Textual comparison:} Paper 14 refers to Hensel, then links the complementary module to the ramification ideal, a normal basis and Riemann–Roch. These references support a technical rather than an ordinary supplementary reading.\par
\noindent{\footnotesize TEXT-COMP-169; pasaže / passages: 9912, 10028; \texttt{CORPUS\_FORM\_COMPARISONS.json}.}\par
\end{minipage}\par
\par\medskip\Needspace{7\baselineskip}
\noindent\begin{minipage}{\linewidth}
\hypertarget{TERM-1025}{}
\noindent\textbf{kompleksne koordinaty / linijne koordinaty}\quad{\footnotesize\texttt{TERM-1025}}\par
\noindent\emph{Iztočnik / Source:} Komplexkoordinaten / Linienkoordinaten\par
\noindent\emph{Rusky / Russian:} комплексные координаты / линейные координаты\par
\noindent\emph{Ukrajinsky / Ukrainian:} комплексні координати / лінійні координати\par
\smallskip\noindent\emph{Razlog:} Kompleksne koordinaty / linijne koordinaty v rabotě 16 sut nazvanja iz simboličnoj geometrije. Sravnjeni tekst zaměnjaje $t_{ik}$ veličinami $p_{ik}=x_i y_k-x_k y_i$ i razgledaje relacije medžu nimi. Zato linijne znači koordinaty linij v toj konstrukciji, ne proizvoljno linearnu zaměnu proměnnyh. Kompleksne ne dolžno bez konteksta byti pročteno samo kak kompleksnočislene. Arhivny razlog predpočitaje prěmo nazvanje koordinat; slovnik mora dodati lokalnu definiciju, da privyčny pridevnik ne skryje istorijny smysl.\par
\noindent\emph{Reason:} Kompleksne koordinaty / linijne koordinaty in Paper 16 names coordinate systems in symbolic geometry. The compared text replaces $t_{ik}$ by $p_{ik}=x_i y_k-x_k y_i$ and studies their relations. Linijne therefore refers to line coordinates in this construction, not an arbitrary linear change of variables. Kompleksne should not be read without context as merely complex-number-valued. The record favours direct coordinate terminology; the dictionary must supply the local construction so familiar adjectives do not conceal the historical sense.\par
\noindent{\footnotesize Izvorne zapisy / Original records: HIST-g078-005.}\par
\smallskip\noindent\emph{Sravnjenje s tekstom:} Formula $p_{ik}=x_i y_k-x_k y_i$ daje konkretno značenje linijnyh koordinat; relacije medžu nimi trěbujut normalnu formu. Nazvanja ne možno razuměti samo kak linearny i kompleksnočislovy v obyčnom sovrěmennom smyslu.\par
\noindent\emph{Textual comparison:} The formula $p_{ik}=x_i y_k-x_k y_i$ gives the concrete meaning of line coordinates; their relations require a normal form. The names cannot be read merely as linear and complex-number-valued in an ordinary modern sense.\par
\noindent{\footnotesize TEXT-COMP-170; pasaže / passages: 10735, 10737, 10739, 10747; \texttt{CORPUS\_FORM\_COMPARISONS.json}.}\par
\end{minipage}\par
\par\medskip\Needspace{7\baselineskip}
\noindent\begin{minipage}{\linewidth}
\hypertarget{TERM-1026}{}
\noindent\textbf{komponenta}\quad{\footnotesize\texttt{TERM-1026}}\par
\noindent\emph{Iztočnik / Source:} Komponente\par
\noindent\emph{Rusky / Russian:} компонента\par
\noindent\emph{Ukrajinsky / Ukrainian:} компонента\par
\smallskip\noindent\emph{Razlog:} Komponenta v rabotě 31 sohranja medžunarodny imenik, blizky ruskomu i ukrajinskomu sravnjenju. Jeden razlog podpiraje projekcijno i razloženo čitanje, drugy prodolžuje je v dokazu, ktory svodi pytanje k jednoj česti direktnoj sumy. Zato komponenta naziva čest konkretnoj algebraičnoj struktury, ne proizvoljny člen spiska. Kontinuitet imenika pomaga čitatelju slěditi redukciju, ale njegov precizny objekt vsegda zavisi od nazvanogo razloženja.\par
\noindent\emph{Reason:} Komponenta in Paper 31 preserves the international noun shared closely with Russian and Ukrainian. One rationale supports the projection/decomposition reading; the other continues it in a proof reducing the problem to one direct-sum component. Thus component denotes a part of a specified algebraic structure, not an arbitrary list item. Lexical continuity helps readers follow the reduction, but the exact object still depends on the stated decomposition.\par
\noindent{\footnotesize Izvorne zapisy / Original records: HIST-g122-003, HIST-g127-002.}\par
\end{minipage}\par
\par\medskip\Needspace{7\baselineskip}
\noindent\begin{minipage}{\linewidth}
\hypertarget{TERM-1029}{}
\noindent\textbf{komponentny klas}\quad{\footnotesize\texttt{TERM-1029}}\par
\noindent\emph{Iztočnik / Source:} Komponentenklasse\par
\noindent\emph{Rusky / Russian:} компонентный класс\par
\noindent\emph{Ukrajinsky / Ukrainian:} компонентний клас\par
\smallskip\noindent\emph{Razlog:} Komponentny klas v rabotě 31, §5, prodolžuje termin iz predhodnogo odlomka i čuva objekt klase. Nyněšnji tekst razlikuje diskriminant komponentnogo ideala po odnošenju k celomu klasu i po odnošenju k komponentnomu klasu. Zato odbrositi klas i govoriti samo o idealu značilo by skryti relativnost konstrukcije. Jedinstvena forma ne jest novi pojem vedle množstvenyh komponentnyh klasov; ona naziva jeden iz njih v argumentu.\par
\noindent\emph{Reason:} Komponentny klas in Paper 31, §5 continues the preceding terminology and preserves the class object. The present text distinguishes the discriminant of a component ideal relative to the whole class from that relative to the component class. Dropping class and speaking only of the ideal would conceal this relativity. The singular does not introduce a concept different from the plural component classes; it names one of them in the argument.\par
\noindent{\footnotesize Izvorne zapisy / Original records: HIST-g139-004.}\par
\smallskip\noindent\emph{Sravnjenje s tekstom:} V rabotě 31 komponentnym klasam idealov odgovarjajut komponentne klasy prstena. Slědy i diskriminanty potom berut se po odnošenju k tym klasam. Vynehanje slova klas skrylo by osnovu sravnjenja.\par
\noindent\emph{Textual comparison:} In Paper 31 component classes of ideals correspond to component classes of the ring. Traces and discriminants are then taken relative to those classes. Omitting class would conceal the basis of comparison.\par
\noindent{\footnotesize TEXT-COMP-173; pasaže / passages: 18152, 18185, 18193, 18195; \texttt{CORPUS\_FORM\_COMPARISONS.json}.}\par
\end{minipage}\par
\par\medskip\Needspace{7\baselineskip}
\noindent\begin{minipage}{\linewidth}
\hypertarget{TERM-1030}{}
\noindent\textbf{komponentne klasy}\quad{\footnotesize\texttt{TERM-1030}}\par
\noindent\emph{Iztočnik / Source:} Komponentenklassen\par
\noindent\emph{Rusky / Russian:} компонентные классы\par
\noindent\emph{Ukrajinsky / Ukrainian:} компонентні класи\par
\smallskip\noindent\emph{Razlog:} Komponentne klasy zadržava dva imenni nivoje: komponente i jih klasy. V rabotě 31 blokovo prědstavjenje prstena daje odgovaranje medžu komponentnymi klasami idealnogo klasa i komponentnymi klasami matričnogo prstena. Zato to ne jest samo spisok sumandov. Kratka pridevnično-imenna konstrukcija imaje prěme ruske i ukrajinske sravnjenja, ale njeje opravdanje v tekstě jest zachovanje različenogo matematičnogo objekta.\par
\noindent\emph{Reason:} Komponentne klasy retains two levels of designation: components and their classes. In Paper 31 the ring’s block representation gives a correspondence between component classes of an ideal class and those of the matrix ring. This is more than a list of summands. The compact adjective–noun construction has close Russian and Ukrainian comparators; its textual justification is preservation of the distinct mathematical object.\par
\noindent{\footnotesize Izvorne zapisy / Original records: HIST-g138-003.}\par
\smallskip\noindent\emph{Sravnjenje s tekstom:} V rabotě 31 komponentnym klasam idealov odgovarjajut komponentne klasy prstena. Slědy i diskriminanty potom berut se po odnošenju k tym klasam. Vynehanje slova klas skrylo by osnovu sravnjenja.\par
\noindent\emph{Textual comparison:} In Paper 31 component classes of ideals correspond to component classes of the ring. Traces and discriminants are then taken relative to those classes. Omitting class would conceal the basis of comparison.\par
\noindent{\footnotesize TEXT-COMP-173; pasaže / passages: 18152, 18185, 18193, 18195; \texttt{CORPUS\_FORM\_COMPARISONS.json}.}\par
\end{minipage}\par
\par\medskip\Needspace{7\baselineskip}
\noindent\begin{minipage}{\linewidth}
\hypertarget{TERM-1031}{}
\noindent\textbf{kompozicijne faktory}\quad{\footnotesize\texttt{TERM-1031}}\par
\noindent\emph{Iztočnik / Source:} Kompositionsfaktoren\par
\noindent\emph{Rusky / Russian:} композиционные факторы\par
\noindent\emph{Ukrajinsky / Ukrainian:} композиційні фактори\par
\smallskip\noindent\emph{Razlog:} Kompozicijne faktory v rabotě 34, §3, prodolžuje faktorgrupnu terminologiju potrěbnu za teoremu Jordana–Höldera. Tekst izrično definira faktory kak kvocienty poslědovateljnyh členov kompozicijnogo reda. Ih prostota razuměje se po odnošenju k dopustimym normalnym děliteljam, pri danyh operatorah. Zato faktor ne znači prosty faktor čisla, a jednakost multiskupa faktorov ne znači, že samy podgrupy v dvoh redah morajut byti identične.\par
\noindent\emph{Reason:} Kompozicijne faktory in Paper 34, §3 continues factor-group terminology needed for Jordan–Hölder. The text explicitly defines the factors as quotients of successive terms of a composition series. Their simplicity is relative to admissible normal subgroups under the specified operators. Factor therefore does not mean prime factor of a number, and agreement of the factors does not mean the subgroups in two series must themselves be identical.\par
\noindent{\footnotesize Izvorne zapisy / Original records: HIST-g181-002.}\par
\smallskip\noindent\emph{Sravnjenje s tekstom:} §3 raboty 34 definira konečny nevstavimy red dopustimyh podgrup i njegove kvocientne faktory. On upotrěbjaje red, dokolě §23 pri matricah prědstavjenja imaje serija. Arhivno rjad jest treti variant istej terminologije; promena imenika ne dava novu konstrukciju.\par
\noindent\emph{Textual comparison:} Paper 34, §3 defines a finite unrefinable series of admissible subgroups and its quotient factors. It uses red, whereas §23 uses serija for representation matrices. Archival rjad is a third form of the same terminology; the noun variation does not introduce a new construction.\par
\noindent{\footnotesize TEXT-COMP-174; pasaže / passages: 19965, 19967, 19969, 19971, 19981, 19985, 22114, 22137; \texttt{CORPUS\_FORM\_COMPARISONS.json}.}\par
\end{minipage}\par
\par\medskip\Needspace{7\baselineskip}
\noindent\begin{minipage}{\linewidth}
\hypertarget{TERM-1032}{}
\noindent\textbf{kompozicijny red}\quad{\footnotesize\texttt{TERM-1032}}\par
\noindent\emph{Iztočnik / Source:} Kompositionsreihe\par
\noindent\emph{Rusky / Russian:} композиционный ряд\par
\noindent\emph{Ukrajinsky / Ukrainian:} композиційний ряд\par
\smallskip\noindent\emph{Razlog:} Kompozicijny red v rabotah 27 i 30 izbrany jest vměsto novoizmysljenogo domačego sočetanja. Oba razloga cěnjat kratkost i znany medžunarodny koren kompozicij-, svazany s grupami i modulami. Nazvanje red samo po sobě ne ustanavjaje kompozicijnost: potrěbno jest ograničenje na dopustime vstavjene členy. Opis izbora kak bole kanoničnogo jest ocěna arhivnogo zapisa; bez navedenogo jezykovogo iztočnika ne može byti predstavjen kak dokaz nezavisnoj medžuslovjanskoj normy.\par
\noindent\emph{Reason:} Kompozicijny red in Papers 27 and 30 is chosen instead of inventing a native compound. Both rationales value concision and the familiar international composition stem in group and module theory. Series alone does not establish a composition series; the restriction on admissible intermediate terms is essential. The record’s description of the choice as more canonical is an editorial assessment, not proof of an independent Interslavic norm without a cited language source.\par
\noindent{\footnotesize Izvorne zapisy / Original records: HIST-g104-008, HIST-g117-002.}\par
\smallskip\noindent\emph{Sravnjenje s tekstom:} §3 raboty 34 definira konečny nevstavimy red dopustimyh podgrup i njegove kvocientne faktory. On upotrěbjaje red, dokolě §23 pri matricah prědstavjenja imaje serija. Arhivno rjad jest treti variant istej terminologije; promena imenika ne dava novu konstrukciju.\par
\noindent\emph{Textual comparison:} Paper 34, §3 defines a finite unrefinable series of admissible subgroups and its quotient factors. It uses red, whereas §23 uses serija for representation matrices. Archival rjad is a third form of the same terminology; the noun variation does not introduce a new construction.\par
\noindent{\footnotesize TEXT-COMP-174; pasaže / passages: 19965, 19967, 19969, 19971, 19981, 19985, 22114, 22137; \texttt{CORPUS\_FORM\_COMPARISONS.json}.}\par
\end{minipage}\par
\par\medskip\Needspace{7\baselineskip}
\noindent\begin{minipage}{\linewidth}
\hypertarget{TERM-1033}{}
\noindent\textbf{kompozicijny rjad}\quad{\footnotesize\texttt{TERM-1033}}\par
\noindent\emph{Iztočnik / Source:} Kompositionsreihe\par
\noindent\emph{Rusky / Russian:} композиционный ряд\par
\noindent\emph{Ukrajinsky / Ukrainian:} композиційний ряд\par
\smallskip\noindent\emph{Razlog:} Kompozicijny rjad v zapisu raboty 34, §3, spojuje medžunarodny pridevnik so slovjanskim imenikom, blizkym ruskomu ряд i ukrajinskomu ряд. Nyněšnji tekst togo paragrafa ipak imaje kompozicijny red. Tuto jest formna alternativa, ne razlika medžu dvoma definicijama kompozicijnogo reda. Metodologično važno jest čuvati i prvonačalny razlog, i formu v tekstě; stabilnost korena ne dokazuje, že vsaka časť termina ostala neproměnena.\par
\noindent\emph{Reason:} Kompozicijny rjad in the Paper 34, §3 record combines the international adjective with a Slavic noun close to Russian and Ukrainian ряд. The present section, however, has kompozicijny red. This is a formal alternative, not two definitions of composition series. Methodologically, both the original reason and the text’s form must remain visible: stability of the adjective does not prove that every part of the phrase stayed unchanged.\par
\noindent{\footnotesize Izvorne zapisy / Original records: HIST-g181-001.}\par
\smallskip\noindent\emph{Sravnjenje s tekstom:} §3 raboty 34 definira konečny nevstavimy red dopustimyh podgrup i njegove kvocientne faktory. On upotrěbjaje red, dokolě §23 pri matricah prědstavjenja imaje serija. Arhivno rjad jest treti variant istej terminologije; promena imenika ne dava novu konstrukciju.\par
\noindent\emph{Textual comparison:} Paper 34, §3 defines a finite unrefinable series of admissible subgroups and its quotient factors. It uses red, whereas §23 uses serija for representation matrices. Archival rjad is a third form of the same terminology; the noun variation does not introduce a new construction.\par
\noindent{\footnotesize TEXT-COMP-174; pasaže / passages: 19965, 19967, 19969, 19971, 19981, 19985, 22114, 22137; \texttt{CORPUS\_FORM\_COMPARISONS.json}.}\par
\end{minipage}\par
\par\medskip\Needspace{7\baselineskip}
\noindent\begin{minipage}{\linewidth}
\hypertarget{TERM-1034}{}
\noindent\textbf{kompozicijny red / kompozicijny faktor}\quad{\footnotesize\texttt{TERM-1034}}\par
\noindent\emph{Iztočnik / Source:} Kompositionsreihe / Kompositionsfaktor\par
\noindent\emph{Rusky / Russian:} композиционный ряд / композиционный фактор\par
\noindent\emph{Ukrajinsky / Ukrainian:} композиційний ряд / композиційний фактор\par
\smallskip\noindent\emph{Razlog:} Kompozicijny red / kompozicijny faktor v pozdnějših paragrafah raboty 34 sohranja svaz s grupami s operatorami iz glavy I. Zapis izbira red radi soglasovanja s izloženjem prědstavjenij i otvorjeno pomina ranějše rjad. Sravnjeni §3 uže govori red, ale §23 imaje serija. Zato dokumentovana želja doslědnosti ne jest dokaz jednotnogo imenika v celoj rabotě. Razlika red/faktor ostaje strukturna: red sostoji se iz podobjektov, faktor jest kvocient medžu poslědovateljnymi členami.\par
\noindent\emph{Reason:} Kompozicijny red / kompozicijny faktor in later sections of Paper 34 retains the connection with Chapter I’s operator groups. The record selects red for consistency with the representation-theory exposition and openly flags earlier rjad. The compared §3 already has red, but §23 uses serija. An intention of consistency is therefore not evidence of a uniform noun throughout the paper. Series and factor remain structurally distinct: the former consists of subobjects, the latter is a quotient between successive terms.\par
\noindent{\footnotesize Izvorne zapisy / Original records: HIST-g194-003, HIST-g196-006.}\par
\smallskip\noindent\emph{Sravnjenje s tekstom:} §3 raboty 34 definira konečny nevstavimy red dopustimyh podgrup i njegove kvocientne faktory. On upotrěbjaje red, dokolě §23 pri matricah prědstavjenja imaje serija. Arhivno rjad jest treti variant istej terminologije; promena imenika ne dava novu konstrukciju.\par
\noindent\emph{Textual comparison:} Paper 34, §3 defines a finite unrefinable series of admissible subgroups and its quotient factors. It uses red, whereas §23 uses serija for representation matrices. Archival rjad is a third form of the same terminology; the noun variation does not introduce a new construction.\par
\noindent{\footnotesize TEXT-COMP-174; pasaže / passages: 19965, 19967, 19969, 19971, 19981, 19985, 22114, 22137; \texttt{CORPUS\_FORM\_COMPARISONS.json}.}\par
\end{minipage}\par
\par\medskip\Needspace{7\baselineskip}
\noindent\begin{minipage}{\linewidth}
\hypertarget{TERM-1035}{}
\noindent\textbf{kompozicijny rjad / kompozicijny faktor}\quad{\footnotesize\texttt{TERM-1035}}\par
\noindent\emph{Iztočnik / Source:} Kompositionsreihe / Kompositionsfaktor\par
\noindent\emph{Rusky / Russian:} композиционный ряд / композиционный фактор\par
\noindent\emph{Ukrajinsky / Ukrainian:} композиційний ряд / композиційний фактор\par
\smallskip\noindent\emph{Razlog:} Kompozicijny rjad / kompozicijny faktor v rabotě 34, §5, naměrno prodolžuje razlog iz §3 i formulaciju Jordana–Höldera. To opravdava kontinuitet značenja, ne novi samostojny leksičny vybor za vsako pojavljanje. Arhivna forma rjad ostaje sravniteljnym variantom k red i serija. Pri objašnjenju trěba jasno razlikovati dolgost reda od njegovih faktorov i ne prěvratiti formnu variantnost v razlikovanje trěh matematičnyh konstrukcij.\par
\noindent\emph{Reason:} Kompozicijny rjad / kompozicijny faktor in Paper 34, §5 deliberately continues the §3 rationale and Jordan–Hölder framing. This supports continuity of meaning, not a new independent lexical selection for every occurrence. Archival rjad remains a comparator to red and serija. The exposition must distinguish series length from its factors and must not turn noun variation into three different mathematical constructions.\par
\noindent{\footnotesize Izvorne zapisy / Original records: HIST-g183-006.}\par
\smallskip\noindent\emph{Sravnjenje s tekstom:} §3 raboty 34 definira konečny nevstavimy red dopustimyh podgrup i njegove kvocientne faktory. On upotrěbjaje red, dokolě §23 pri matricah prědstavjenja imaje serija. Arhivno rjad jest treti variant istej terminologije; promena imenika ne dava novu konstrukciju.\par
\noindent\emph{Textual comparison:} Paper 34, §3 defines a finite unrefinable series of admissible subgroups and its quotient factors. It uses red, whereas §23 uses serija for representation matrices. Archival rjad is a third form of the same terminology; the noun variation does not introduce a new construction.\par
\noindent{\footnotesize TEXT-COMP-174; pasaže / passages: 19965, 19967, 19969, 19971, 19981, 19985, 22114, 22137; \texttt{CORPUS\_FORM\_COMPARISONS.json}.}\par
\end{minipage}\par
\par\medskip\Needspace{7\baselineskip}
\noindent\begin{minipage}{\linewidth}
\hypertarget{TERM-1036}{}
\noindent\textbf{kompozicijna serija modula predstavjenja}\quad{\footnotesize\texttt{TERM-1036}}\par
\noindent\emph{Iztočnik / Source:} Kompositionsreihe des Darstellungsmoduls\par
\noindent\emph{Rusky / Russian:} композиционный ряд модуля представления\par
\noindent\emph{Ukrajinsky / Ukrainian:} композиційний ряд модуля представлення\par
\smallskip\noindent\emph{Razlog:} Kompozicijna serija modula predstavjenja v rabotě 34, §23, naměrno izbira serija vměsto rjad radi širšej medžuslovjanskoj razumljivosti. To jest zabilěžena ocěna, ne rezultat ispytanja. Nyněšnji odlomok dava kompozicijna serija modula prědstavjenja pri izboru bazy, ktora vodi k blokovoj formě matrice $W$. Smysl zato ostaje kompozicijny red modula prědstavjenja, ne osobny novy vid serije. Slovnik mora javno pokazati konkurenciju serija/red/rjad i razumljivost ne proglasiti dokazanom.\par
\noindent\emph{Reason:} Kompozicijna serija modula predstavjenja in Paper 34, §23 deliberately chooses serija over rjad for broader cross-Slavic readability. That is a recorded judgment, not an experimental result. The present passage has kompozicijna serija modula prědstavjenja when choosing a basis that gives the matrix $W$ a block form. The meaning remains a composition series of the representation module, not a new kind of series. The dictionary must expose serija/red/rjad competition without declaring intelligibility proved.\par
\noindent{\footnotesize Izvorne zapisy / Original records: HIST-g201-005.}\par
\smallskip\noindent\emph{Sravnjenje s tekstom:} §3 raboty 34 definira konečny nevstavimy red dopustimyh podgrup i njegove kvocientne faktory. On upotrěbjaje red, dokolě §23 pri matricah prědstavjenja imaje serija. Arhivno rjad jest treti variant istej terminologije; promena imenika ne dava novu konstrukciju.\par
\noindent\emph{Textual comparison:} Paper 34, §3 defines a finite unrefinable series of admissible subgroups and its quotient factors. It uses red, whereas §23 uses serija for representation matrices. Archival rjad is a third form of the same terminology; the noun variation does not introduce a new construction.\par
\noindent{\footnotesize TEXT-COMP-174; pasaže / passages: 19965, 19967, 19969, 19971, 19981, 19985, 22114, 22137; \texttt{CORPUS\_FORM\_COMPARISONS.json}.}\par
\end{minipage}\par
\par\medskip\Needspace{7\baselineskip}
\noindent\begin{minipage}{\linewidth}
\hypertarget{TERM-1037}{}
\noindent\textbf{kongruentny po modulu}\quad{\footnotesize\texttt{TERM-1037}}\par
\noindent\emph{Iztočnik / Source:} kongruent mod $M$\par
\noindent\emph{Rusky / Russian:} конгруэнтные по модулю\par
\noindent\emph{Ukrajinsky / Ukrainian:} конгруентні за модулем\par
\smallskip\noindent\emph{Razlog:} Kongruentny po modulu čuva algebraično nazvanje v rabotě 17. Formula $F\equiv G$ po modulu $\mathfrak M$ znači, že $F-G$ naleži k $\mathfrak M$. To uslovje oprědělja klasy ostatkov, ale ne ustanavjaje avtomatično množenje klasov: tekst potom opozarja, že v nekomutativnom slučaju produkt dvoh proizvoljnyh klasov ne mora byti dobro oprěděljen. Zato znano nazvanje kongruencije ne može zaměniti prověrku operacij, ktore konkretna struktura dopušča.\par
\noindent\emph{Reason:} Kongruentny po modulu retains the algebraic congruence term in Paper 17. The relation $F\equiv G$ modulo $\mathfrak M$ means that $F-G$ belongs to $\mathfrak M$. This defines residue classes but does not automatically give them a multiplication: the text warns that, in the noncommutative setting, the product of arbitrary classes need not be well defined. Familiar congruence terminology cannot replace checking which operations the particular structure admits.\par
\noindent{\footnotesize Izvorne zapisy / Original records: HIST-g079-015.}\par
\smallskip\noindent\emph{Sravnjenje s tekstom:} Razlika $F-G$ dolžna naležati k modulu. Slědujuči odlomok pokazyvaje, že produkt dvoh proizvoljnyh klasov ne mora byti dobro oprěděljen, hotě lěvo umnoženje polinomom jest dopustimo.\par
\noindent\emph{Textual comparison:} The difference $F-G$ must belong to the module. The following passage shows that multiplication of two arbitrary classes need not be well defined, although left multiplication by a polynomial is admissible.\par
\noindent{\footnotesize TEXT-COMP-175; pasaže / passages: 11098, 11110; \texttt{CORPUS\_FORM\_COMPARISONS.json}.}\par
\end{minipage}\par
\par\medskip\Needspace{7\baselineskip}
\noindent\begin{minipage}{\linewidth}
\hypertarget{TERM-1039}{}
\noindent\textbf{konjugaty}\quad{\footnotesize\texttt{TERM-1039}}\par
\noindent\emph{Iztočnik / Source:} Konjugierte\par
\noindent\emph{Rusky / Russian:} сопряжённые\par
\noindent\emph{Ukrajinsky / Ukrainian:} спряжені\par
\smallskip\noindent\emph{Razlog:} Konjugaty v zapisah raboty 9, §§3 i 9, i raboty 10, §2, zadržava medžunarodny koren dlja algebraičnoj konjugacije. Razlogy svezujut tri odlomky s ranějšim terminom, dokolě rusky i ukrajinsky sravniteljne formy imajut slovjanske koreni. To jest izbor kratkogo tehničnogo nazvanja, ne tvrdženje, že vse vidy konjugacije sut jedna operacija. Razgledane objekty i osnovna oblast morajut byti nazvane v konkretnom argumentu; na primer, konjugaty oblastij ne sut obyčno kompleksno sprjažene čisla.\par
\noindent\emph{Reason:} Konjugaty in the records for Paper 9, §§3 and 9 and Paper 10, §2 retains the international root for algebraic conjugacy. The rationales connect all three passages with earlier terminology, while Russian and Ukrainian use Slavic-root comparators. This selects a compact technical designation; it does not identify every kind of conjugation as one operation. The objects and base domain must be specified in the argument: conjugate domains, for example, are not merely complex-conjugate numbers.\par
\noindent{\footnotesize Izvorne zapisy / Original records: HIST-g061-011, HIST-g067-010, HIST-g070-010.}\par
\end{minipage}\par
\par\medskip\Needspace{7\baselineskip}
\noindent\begin{minipage}{\linewidth}
\hypertarget{TERM-1040}{}
\noindent\textbf{konjugat / konjugaty vzhodno k $(H)$}\quad{\footnotesize\texttt{TERM-1040}}\par
\noindent\emph{Iztočnik / Source:} Konjugierte / in bezug auf $(H)$ Konjugierten\par
\noindent\emph{Rusky / Russian:} сопряжённое / сопряжённые относительно $(H)$\par
\noindent\emph{Ukrajinsky / Ukrainian:} спряжене / спряжені відносно $(H)$\par
\smallskip\noindent\emph{Razlog:} Konjugat / konjugaty vzhodno k $(H)$ čuvaje relativnost konstrukcije klasov v rabotě 9, §8. Zapis naměrno zadržava konjugat, ale sam označaje vzhodno k kak konstrukciju za daljšu ocěnu. Nyněšnji tekst imaje po odnošenju k $(H)$; primětka kaže, že elementy konjugovanyh oblastij možut byti racionalno izrazimy prěz početnu oblast, ale ne potrěbno cělo i racionalno. Zato oblasti možut byti različne. Ispuščenje osnovy $(H)$ ili tiho tvrdženje identičnosti skrylo by matematično uslovje.\par
\noindent\emph{Reason:} Konjugat / konjugaty vzhodno k $(H)$ preserves the relative character of the class construction in Paper 9, §8. The record deliberately keeps konjugat but itself flags vzhodno k for further appraisal. The present text has po odnošenju k $(H)$; its note says elements of the conjugate domains can be expressed rationally through the initial domain, but not necessarily integrally and rationally. The domains may therefore differ. Removing $(H)$ or silently asserting identity would lose a mathematical qualification.\par
\noindent{\footnotesize Izvorne zapisy / Original records: HIST-g066-009.}\par
\smallskip\noindent\emph{Sravnjenje s tekstom:} Sravnjeni odlomok raboty 9 imaje po odnošenju k $(H)$ i izrično dopušča različnost konjugovanyh oblastij. Racionalna izrazimost ne mora byti cělo racionalna izrazimost. To jest matematična granica, ktoru predložna fraza mora čuvati.\par
\noindent\emph{Textual comparison:} The compared Paper 9 passage has po odnošenju k $(H)$ and expressly allows the conjugate domains to differ. Rational expressibility need not be integral rational expressibility. The prepositional phrase must preserve that mathematical qualification.\par
\noindent{\footnotesize TEXT-COMP-177; pasaže / passages: 8115, 8116, 8117, 8118, 8119, 8120; \texttt{CORPUS\_FORM\_COMPARISONS.json}.}\par
\end{minipage}\par
\par\medskip\Needspace{7\baselineskip}
\noindent\begin{minipage}{\linewidth}
\hypertarget{TERM-1041}{}
\noindent\textbf{konjugaty imenovatelja}\quad{\footnotesize\texttt{TERM-1041}}\par
\noindent\emph{Iztočnik / Source:} Konjugierte des Nenners\par
\noindent\emph{Rusky / Russian:} сопряженные знаменателя\par
\noindent\emph{Ukrajinsky / Ukrainian:} спряжені знаменника\par
\smallskip\noindent\emph{Razlog:} Konjugaty imenovatelja v rabotě 7 naziva elementy potrěbne za odstranjenje imenovatelja pri racionalnom absolutnom invariantu. Genitiv imenovatelja čuva, čije konjugaty se berut; tekst dodatno oprěděljaje odnos k grupě. Množenje čislitelja i imenovatelja dotičnymi konjugatami dava količnik cělyh racionalnyh absolutnyh invariantov. Tekst izrično ne zahtěvaje, že oba sut vzajemno proste. Zato obično navodjenje na racionalizaciju kompleksnogo imenovatelja ne objašnjaje toj konkretny postupok.\par
\noindent\emph{Reason:} Konjugaty imenovatelja in Paper 7 names the conjugates needed to clear the denominator of a rational absolute invariant. The genitive identifies what is conjugated; the text further specifies conjugacy with respect to the group. Multiplying numerator and denominator by the relevant conjugates produces a quotient of integral rational absolute invariants. The text expressly allows these to have common factors. Ordinary complex-denominator rationalization would therefore not explain the construction.\par
\noindent{\footnotesize Izvorne zapisy / Original records: HIST-g056-011.}\par
\smallskip\noindent\emph{Sravnjenje s tekstom:} Odlomok raboty 7 izrično svjazuje konjugaty imenovatelja s grupoju. Dobijeny čislitelj i imenovatelj sut cěle racionalne absolutne invarianty, ale ne morajut byti vzajemno proste.\par
\noindent\emph{Textual comparison:} Paper 7 specifies conjugates relative to the group. The resulting numerator and denominator are integral rational absolute invariants, but need not be relatively prime.\par
\noindent{\footnotesize TEXT-COMP-178; pasaže / passages: 6706; \texttt{CORPUS\_FORM\_COMPARISONS.json}.}\par
\end{minipage}\par
\par\medskip\Needspace{7\baselineskip}
\noindent\begin{minipage}{\linewidth}
\hypertarget{TERM-1042}{}
\noindent\textbf{konjugovane elementy}\quad{\footnotesize\texttt{TERM-1042}}\par
\noindent\emph{Iztočnik / Source:} konjugierte Elemente\par
\noindent\emph{Rusky / Russian:} сопряженные элементы\par
\noindent\emph{Ukrajinsky / Ukrainian:} спряжені елементи\par
\smallskip\noindent\emph{Razlog:} Konjugovane elementy v rabotě 31, §6, prodolžuje medžunarodny korenj ranějših konjugovanyh faktorov i polj. Zapis sravnjaje rusko сопряженные i ukrajinsko спряжені, ale predpočita konjugovane prěd prinudnym kalkom. V tekstě elementy pojavjajut se v prědstavjenju slěda, normy i diskriminanta pri polnoj razložimosti prvogo roda; v poljevom slučaju odpovědajut izomorfizmam, ktore fiksujut $P$. Nijedna opča podobnost elementov ni kompleksno konjugovanje bez ukazanja osnovy ne daje togo smysla.\par
\noindent\emph{Reason:} Konjugovane elementy in Paper 31, §6 continues the international stem used for conjugate factors and fields. The record compares Russian сопряженные and Ukrainian спряжені but prefers konjugovane to a forced calque. In the text, the elements occur in formulas for trace, norm and discriminant under complete decomposability of the first kind; in the field case they correspond to isomorphisms fixing $P$. Neither unspecified similarity nor complex conjugation without a base-field qualification supplies that meaning.\par
\noindent{\footnotesize Izvorne zapisy / Original records: HIST-g146-002.}\par
\smallskip\noindent\emph{Sravnjenje s tekstom:} V poljevom slučaju raboty 31 izomorfizmy fiksujut osnovno polje. Tekst dodatno razlikuje izomorfny objekt od polja, už vključenogo v Galoisovo razširjenje.\par
\noindent\emph{Textual comparison:} In Paper 31's field case, the isomorphisms fix the base field. The text further distinguishes an isomorphic object from a field already included in the Galois extension.\par
\noindent{\footnotesize TEXT-COMP-179; pasaže / passages: 18470, 18474, 18525; \texttt{CORPUS\_FORM\_COMPARISONS.json}.}\par
\end{minipage}\par
\par\medskip\Needspace{7\baselineskip}
\noindent\begin{minipage}{\linewidth}
\hypertarget{TERM-1043}{}
\noindent\textbf{konjugovane veličiny}\quad{\footnotesize\texttt{TERM-1043}}\par
\noindent\emph{Iztočnik / Source:} konjugierte Größen\par
\noindent\emph{Rusky / Russian:} сопряженные величины\par
\noindent\emph{Ukrajinsky / Ukrainian:} спряжені величини\par
\smallskip\noindent\emph{Razlog:} Konjugovane veličiny v rabotě 6, §5, odnosi se k linearnym formam v algebraičnom razširjenju. Tekst razglědaje $z-u(x)$ i njegove konjugaty pri tvorjenju involucijskoj formy. Veličina ne znači samo realno izměrimy razměr: tuto sut algebraične izrazy. Zapis kaže, že formulacija slědovala anglijskoj kontrolnoj versiji; to jest podatok o posrědnom iztočniku izbora, ne dokaz samostojnoj medžuslovjanskoj normy. Osnovno polje i relacija konjugovanosti ostajut čestju objašnjenja.\par
\noindent\emph{Reason:} Konjugovane veličiny in Paper 6, §5 refers to linear forms in an algebraic extension. The text uses $z-u(x)$ and its conjugates to construct the involution form. Veličina here denotes an algebraic expression, not necessarily a measurable real magnitude. The record says that the phrasing followed an English control version: this identifies an intermediary in the choice, not independent Interslavic authority. The base field and the conjugacy relation remain part of the explanation.\par
\noindent{\footnotesize Izvorne zapisy / Original records: HIST-g045-006.}\par
\smallskip\noindent\emph{Sravnjenje s tekstom:} Rabota 6 povezuje faktory $z-u(x)$ s redami konjugovanyh veličin i jejich elementarnymi simetričnymi funkcijami. Odstranjenje redov iz nazvanja skrylo by koordiniranu strukturu, na ktoroj stoji argument.\par
\noindent\emph{Textual comparison:} Paper 6 connects the factors $z-u(x)$ to conjugate tuples and their elementary symmetric functions. Omitting the tuples from the name would conceal the coordinated structure used by the argument.\par
\noindent{\footnotesize TEXT-COMP-180; pasaže / passages: 5533, 5538, 5540, 5542, 5580, 5583, 5584, 5586, 5595, 5916, 5918, 5922, 5923, 5925; \texttt{CORPUS\_FORM\_COMPARISONS.json}.}\par
\end{minipage}\par
\par\medskip\Needspace{7\baselineskip}
\noindent\begin{minipage}{\linewidth}
\hypertarget{TERM-1044}{}
\noindent\textbf{konjugovane redy veličin}\quad{\footnotesize\texttt{TERM-1044}}\par
\noindent\emph{Iztočnik / Source:} konjugierte Größenreihen\par
\noindent\emph{Rusky / Russian:} сопряженные ряды величин\par
\noindent\emph{Ukrajinsky / Ukrainian:} спряжені ряди величин\par
\smallskip\noindent\emph{Razlog:} Konjugovane redy veličin v rabotě 6, §8, sohranja redy iz §5, zato že dokaz ponovno koristi elementarne simetrične funkcije tih redov. Red znači zajedno vzety sistem koordinat; ne trěba čitati go kak beskonečnu čislovu seriju. Sravnjeny odlomok izrično ukazuje polje, po odnošenju k ktoromu redy sut konjugovane, i piše početny red s jego konjugatami. Ponavjanje termina čuva tu vezu v dokazu; ne opravdava izpuščenje slova redy ili njegovo zaměnjenje neprirědženym spiskom elementov.\par
\noindent\emph{Reason:} Konjugovane redy veličin in Paper 6, §8 retains the rows-of-quantities wording from §5 because the proof again uses their elementary symmetric functions. A row is a coordinated tuple, not an infinite numerical series. The compared passage specifies the field relative to which the rows are conjugate and displays the initial row with its conjugates. Reusing the phrase preserves that proof connection; it does not justify dropping rows or replacing coordinated tuples with an unstructured list.\par
\noindent{\footnotesize Izvorne zapisy / Original records: HIST-g048-012.}\par
\smallskip\noindent\emph{Sravnjenje s tekstom:} Rabota 6 povezuje faktory $z-u(x)$ s redami konjugovanyh veličin i jejich elementarnymi simetričnymi funkcijami. Odstranjenje redov iz nazvanja skrylo by koordiniranu strukturu, na ktoroj stoji argument.\par
\noindent\emph{Textual comparison:} Paper 6 connects the factors $z-u(x)$ to conjugate tuples and their elementary symmetric functions. Omitting the tuples from the name would conceal the coordinated structure used by the argument.\par
\noindent{\footnotesize TEXT-COMP-180; pasaže / passages: 5533, 5538, 5540, 5542, 5580, 5583, 5584, 5586, 5595, 5916, 5918, 5922, 5923, 5925; \texttt{CORPUS\_FORM\_COMPARISONS.json}.}\par
\end{minipage}\par
\par\medskip\Needspace{7\baselineskip}
\noindent\begin{minipage}{\linewidth}
\hypertarget{TERM-1045}{}
\noindent\textbf{konjugacijna klasa}\quad{\footnotesize\texttt{TERM-1045}}\par
\noindent\emph{Iztočnik / Source:} Konjugiertenklasse / Klassen konjugierter Gruppenelemente\par
\noindent\emph{Rusky / Russian:} класс сопряженности\par
\noindent\emph{Ukrajinsky / Ukrainian:} клас спряженості\par
\smallskip\noindent\emph{Razlog:} Konjugacijna klasa naziva klasu grupovyh elementov, srodnu ruskomu класс сопряженности i ukrajinskomu клас спряженості. Zapis predpočita medžunarodny pridavnik vměsto novogo domačego kalka, ale ocěna jasnosti ne jest izměrena. V §26 raboty 34 trěba razlikovati klasu konjugovanosti od sume jej elementov v grupovom prstenu. Koeficienty centralnogo elementa sut konstantne na klasi, dokolě $K_a$ jest suma različnyh konjugatov. Dvě svjazane upotrěby slova klasa ne sut doslovno jedin objekt.\par
\noindent\emph{Reason:} Konjugacijna klasa names the conjugacy class corresponding to Russian класс сопряженности and Ukrainian клас спряженості. The record favours the international adjective over inventing a native calque, but its clarity assessment is not measured. Paper 34, §26 requires a distinction between the class of group elements and its sum in the group ring. A central element has coefficients constant on a class, whereas $K_a$ is the sum of its distinct conjugates. These related uses of class are not literally one object.\par
\noindent{\footnotesize Izvorne zapisy / Original records: HIST-g204-007.}\par
\smallskip\noindent\emph{Sravnjenje s tekstom:} Koeficienty centralnogo elementa sut konstantne na konjugacijnoj klasi. Ranějše $K_a$ jest suma različnyh konjugatov; klasa elementov i suma v grupovom prstenu trěbujut različne objašnjenja.\par
\noindent\emph{Textual comparison:} Coefficients of a central element are constant on a conjugacy class. The preceding $K_a$ is a sum of distinct conjugates; the class of elements and its group-ring sum need separate explanations.\par
\noindent{\footnotesize TEXT-COMP-181; pasaže / passages: 22475, 22477, 22479, 22484, 22486; \texttt{CORPUS\_FORM\_COMPARISONS.json}.}\par
\end{minipage}\par
\par\medskip\Needspace{7\baselineskip}
\noindent\begin{minipage}{\linewidth}
\hypertarget{TERM-1046}{}
\noindent\textbf{konstanty}\quad{\footnotesize\texttt{TERM-1046}}\par
\noindent\emph{Iztočnik / Source:} Konstanten ($\xi_i a=a\xi_i$)\par
\noindent\emph{Rusky / Russian:} константы\par
\noindent\emph{Ukrajinsky / Ukrainian:} константи\par
\smallskip\noindent\emph{Razlog:} Konstanty jest obično medžunarodno slovo, no rabota 17, §10, dava jemu točno lokalno uslovje: veličiny $a$ iz oblasti racionalnosti trěbujut izpolnjati $\xi_i a=a\xi_i$. Zapis naměrno sohranja tu definiciju, da by slovo ne bylo čitano samo kak neproměnny čislovy parametr. Teorem VIII predpolagaje beskonečno mnogo takyh konstant. Izpuščenje relacije komutovanja ili te beskonečnosti proměnilo by hipotezu, iako samo slovo ostalo by poznano.\par
\noindent\emph{Reason:} Konstanty is an ordinary international term, but Paper 17, §10 imposes a precise local condition: quantities $a$ in the rationality domain must satisfy $\xi_i a=a\xi_i$. The record preserves this definition so that constant is not read merely as a fixed numerical parameter. Theorem VIII requires infinitely many such constants. Omitting the commutation condition or that infinitude would change the hypothesis even though the familiar word remained.\par
\noindent{\footnotesize Izvorne zapisy / Original records: HIST-g079-057.}\par
\smallskip\noindent\emph{Sravnjenje s tekstom:} Teorem VIII raboty 17 sojedinja uslovje $\xi_i a=a\xi_i$ s potrěbovanjem beskonečno mnogyh konstant. Obe česti sut hipotezy, ne osobnosti nazvanja.\par
\noindent\emph{Textual comparison:} Paper 17, Theorem VIII combines $\xi_i a=a\xi_i$ with the requirement of infinitely many constants. Both are hypotheses, not optional features of the terminology.\par
\noindent{\footnotesize TEXT-COMP-182; pasaže / passages: 11737; \texttt{CORPUS\_FORM\_COMPARISONS.json}.}\par
\end{minipage}\par
\par\medskip\Needspace{7\baselineskip}
\noindent\begin{minipage}{\linewidth}
\hypertarget{TERM-1048}{}
\noindent\textbf{konstrukcija $f(z)$}\quad{\footnotesize\texttt{TERM-1048}}\par
\noindent\emph{Iztočnik / Source:} Konstruktion von $f(z)$\par
\noindent\emph{Rusky / Russian:} построение $f(z)$\par
\noindent\emph{Ukrajinsky / Ukrainian:} побудова $f(z)$\par
\smallskip\noindent\emph{Razlog:} Konstrukcija $f(z)$ v rabotě 10, §3, izbira medžunarodny imennik vměsto ruskogo построение i ukrajinskogo побудова. Sravnjeny početok paragrafa govori, že najdene potrěbne uslovja davajut takože konstrukciju: zadavaje drugo dobro uporjadkovanje i postupno prirědžaje ješče neupotrěbjene bazove elementy. Zato konstrukcija ima točno dokazno sodržanje. Moderno nazvanje ne daje prava prědložiti konvergentnu čislovu proceduru ili konečnu implementaciju, ktoryh fraza sama ne obeščava.\par
\noindent\emph{Reason:} Konstrukcija $f(z)$ in Paper 10, §3 chooses an international noun beside Russian построение and Ukrainian побудова. The opening says that the necessary conditions also yield a construction: it specifies a second well-ordering and successively assigns previously unassigned basis elements. Construction therefore has definite proof content. Modern terminology does not authorize presenting it as a convergent numerical procedure or a finite implementation, neither of which the phrase promises.\par
\noindent{\footnotesize Izvorne zapisy / Original records: HIST-g071-001.}\par
\smallskip\noindent\emph{Sravnjenje s tekstom:} §7 raboty 9 razlikuje dokaz obstajanja od konstrukcije pri zadannom dobrom uporjadkovanju. §3 raboty 10 izrično prirědžaje bazove elementy s pomočju drugogo dobrogo uporjadkovanja. Reč konstrukcija ne obeščava konečny algoritm.\par
\noindent\emph{Textual comparison:} Paper 9, §7 distinguishes existence from construction under a specified well-order. Paper 10, §3 explicitly assigns basis elements using another well-order. The word construction does not promise a finite algorithm.\par
\noindent{\footnotesize TEXT-COMP-183; pasaže / passages: 7839, 7840, 7841, 7842, 8518; \texttt{CORPUS\_FORM\_COMPARISONS.json}.}\par
\end{minipage}\par
\par\medskip\Needspace{7\baselineskip}
\noindent\begin{minipage}{\linewidth}
\hypertarget{TERM-1049}{}
\noindent\textbf{kontinuirna grupa v smislu Lija}\quad{\footnotesize\texttt{TERM-1049}}\par
\noindent\emph{Iztočnik / Source:} kontinuierliche Gruppe im Lieschen Sinne\par
\noindent\emph{Rusky / Russian:} непрерывная группа в смысле Ли\par
\noindent\emph{Ukrajinsky / Ukrainian:} неперервна група у сенсі Лі\par
\smallskip\noindent\emph{Razlog:} Kontinuirna grupa v smislu Lija sohranja istorijno ograničenje raboty 13. Pridavnik kontinuirna ne nosi vsu informaciju: dodatok v smislu Lija svjazuje variacijske problemy s konkretnymi grupami transformacij. Vvedenje samogo teksta razlikuje uravnenja, nastavajuče iz variacijskyh problemov, od libovoljnyh diferencialnyh uravnenij, dopuščajučih grupu. Zato ne trěba skračati frazu do obče kontinuirnoj simetrije i potratiti toj kontrast. Zapis opira izbor na teoriju Lija, ne na tvrdženje o novom izslědovanji razumljivosti.\par
\noindent\emph{Reason:} Kontinuirna grupa v smislu Lija retains the historical qualification in Paper 13. The adjective continuous does not carry all the information: in Lie's sense connects the variational problems to the specified transformation-group setting. The introduction distinguishes equations arising from variational problems from arbitrary differential equations admitting a group. Reducing the phrase to generic continuous symmetry would lose that contrast. The record grounds the choice in Lie-theoretic framing, not in a new comprehension study.\par
\noindent{\footnotesize Izvorne zapisy / Original records: HIST-g075-002.}\par
\smallskip\noindent\emph{Sravnjenje s tekstom:} Vvedenje raboty 13 izrično razlikuje diferencialne uravnenja variacijskyh problemov od libovoljnyh uravnenij, dopuščajučih grupu. Liejev kontekst i to ograničenje trěbujut ostati vidime.\par
\noindent\emph{Textual comparison:} Paper 13's introduction distinguishes differential equations arising from variational problems from arbitrary equations admitting a group. The Lie-theoretic context and that restriction must remain visible.\par
\noindent{\footnotesize TEXT-COMP-184; pasaže / passages: 9195; \texttt{CORPUS\_FORM\_COMPARISONS.json}.}\par
\end{minipage}\par
\par\medskip\Needspace{7\baselineskip}
\noindent\begin{minipage}{\linewidth}
\hypertarget{TERM-1050}{}
\noindent\textbf{kontragradientny / kogradientny}\quad{\footnotesize\texttt{TERM-1050}}\par
\noindent\emph{Iztočnik / Source:} kontragredient / kogredient\par
\noindent\emph{Rusky / Russian:} контрагредиентный / когредиентный\par
\noindent\emph{Ukrajinsky / Ukrainian:} контрагредієнтний / когредієнтний\par
\smallskip\noindent\emph{Razlog:} Kontragradientny / kogradientny jest arhivna dvojica iz §26 raboty 2. Jej razlog jest sohranitje medžunarodnyh terminov različno usměrenyh transformacij. Ale zapis ima -gradient-, dokolě ruska i ukrajinska sravnjenja i ranějše medžuslovjanske zapisy imajut -gredient-. Nyněšnji odlomok sohranja kontragradientnym i kogradientne. To jest dokumentovana pravopisna nedoslědnost, ne dokaz dvuh novyh matematičnyh smyslov. Slovnik pokazyvaje ju zajedno s ranějšim pravilom, ne izdavaje neprovedeno soglasovanje za završeno.\par
\noindent\emph{Reason:} Kontragradientny / kogradientny is the archived pair from Paper 2, §26. Its rationale is retention of international terms for different transformation behaviours. Yet the record has -gradient-, whereas its Russian and Ukrainian comparators and earlier Interslavic records have -gredient-. The current passage retains kontragradientnym and kogradientne. This is a documented spelling inconsistency, not evidence of two additional mathematical meanings. It must remain visible beside the earlier rule rather than being described as resolved.\par
\noindent{\footnotesize Izvorne zapisy / Original records: HIST-g027-003.}\par
\smallskip\noindent\emph{Sravnjenje s tekstom:} Odlomok §26 raboty 2 imaje kontragradientnym i kogradientne. To jest točno svědčenje upotrěby -gradient- v sravnjenom tekstě, a ne dokaz završenoj normalizacije k -gredient-.\par
\noindent\emph{Textual comparison:} Paper 2, §26 has kontragradientnym and kogradientne. These attest -gradient- in the compared text, rather than a completed normalization to -gredient-.\par
\noindent{\footnotesize TEXT-COMP-185; pasaže / passages: 3340, 3344; \texttt{CORPUS\_FORM\_COMPARISONS.json}.}\par
\end{minipage}\par
\par\medskip\Needspace{7\baselineskip}
\noindent\begin{minipage}{\linewidth}
\hypertarget{TERM-1051}{}
\noindent\textbf{kontragredientne formy}\quad{\footnotesize\texttt{TERM-1051}}\par
\noindent\emph{Iztočnik / Source:} kontragrediente Formen\par
\noindent\emph{Rusky / Russian:} контрагредиентные формы\par
\noindent\emph{Ukrajinsky / Ukrainian:} контрагредієнтні форми\par
\smallskip\noindent\emph{Razlog:} Kontragredientne formy sojedinja šest zapisov iz raboty 2. Razlog počinaje od terminologije simbolov v vvedenju i §1, potom prodolžuje ju za formy v §8 i daljših redukcijah. Bitna jest razlika od prosto pridruženoj formy: kontragredientnost oprěděljaje, ktore produkty vhodet v dotične kontrakcije. Zapis o prěnosah slov pri prelomu reda v ruskoj i ukrajinskoj versijah označaje tipografiju, ne razliku termina. Povtorjena motivacija čuva kontinuum pojma, no ne stvara šest nezavisnyh jezykovyh dokazov.\par
\noindent\emph{Reason:} Kontragredientne formy brings together six Paper 2 records. The rationale extends the symbol terminology of the introduction and §1 to forms in §8 and later reductions. Its distinction from merely associated forms matters: contragredience governs which products enter the stated contractions. A record about discretionary hyphenation in the Russian and Ukrainian versions concerns typesetting, not a different term. Repeated reasons establish continuity of the concept, not six independent linguistic attestations.\par
\noindent{\footnotesize Izvorne zapisy / Original records: HIST-g009-007, HIST-g010-007, HIST-g011-014, HIST-g013-014, HIST-g019-010, HIST-g020-007.}\par
\smallskip\noindent\emph{Sravnjenje s tekstom:} Početna konstrukcija razlikuje čislo redov kontragredientnyh simbolov. Daljše redukcije ograničajut vhodenje odredženyh form v kontrakcije produktami s kontragredientnymi formami. To jest matematična relacija, ne opča asociacija.\par
\noindent\emph{Textual comparison:} The initial construction distinguishes the number of contragredient symbol rows. Later reductions restrict specified forms' participation in contractions to products with contragredient forms. This is a mathematical relation, not generic association.\par
\noindent{\footnotesize TEXT-COMP-186; pasaže / passages: 299, 424, 1197, 1260, 2004; \texttt{CORPUS\_FORM\_COMPARISONS.json}.}\par
\end{minipage}\par
\par\medskip\Needspace{7\baselineskip}
\noindent\begin{minipage}{\linewidth}
\hypertarget{TERM-1052}{}
\noindent\textbf{kontragredientne simboly}\quad{\footnotesize\texttt{TERM-1052}}\par
\noindent\emph{Iztočnik / Source:} kontragrediente Symbole\par
\noindent\emph{Rusky / Russian:} контрагредиентные символы\par
\noindent\emph{Ukrajinsky / Ukrainian:} контрагредієнтні символи\par
\smallskip\noindent\emph{Razlog:} Kontragredientne simboly zadržava pridavnik teorije invariantov, už upotrěbjeny v primětkě vvedenja raboty 2. V §1 čislo redov simbolov vlijaje na jednoznačnost reda form: slučaj dvuh redov razlikuje se od slučaja s bolje redami. Simboly zato sut čestju točnoj simboličnoj konstrukcije, ne prosto drugoje slovo za proizvoljne znaky. Kratkost termina ne zaměnjaje objašnjenje transformacijnogo odnosa, no sohranja vezu s daljšimi kontragredientnymi formami.\par
\noindent\emph{Reason:} Kontragredientne simboly retains the invariant-theory adjective already used in the introductory footnote of Paper 2. In §1, the number of symbol rows affects uniqueness of the series of forms: two rows are distinguished from a greater number. Symbols thus belong to a specific symbolic construction, not merely an arbitrary collection of marks. The concise phrase does not replace explanation of the transformation relation, but preserves its connection with the later contragredient forms.\par
\noindent{\footnotesize Izvorne zapisy / Original records: HIST-g003-011.}\par
\smallskip\noindent\emph{Sravnjenje s tekstom:} Početna konstrukcija razlikuje čislo redov kontragredientnyh simbolov. Daljše redukcije ograničajut vhodenje odredženyh form v kontrakcije produktami s kontragredientnymi formami. To jest matematična relacija, ne opča asociacija.\par
\noindent\emph{Textual comparison:} The initial construction distinguishes the number of contragredient symbol rows. Later reductions restrict specified forms' participation in contractions to products with contragredient forms. This is a mathematical relation, not generic association.\par
\noindent{\footnotesize TEXT-COMP-186; pasaže / passages: 299, 424, 1197, 1260, 2004; \texttt{CORPUS\_FORM\_COMPARISONS.json}.}\par
\end{minipage}\par
\par\medskip\Needspace{7\baselineskip}
\noindent\begin{minipage}{\linewidth}
\hypertarget{TERM-1053}{}
\noindent\textbf{kontravariant/kovariant}\quad{\footnotesize\texttt{TERM-1053}}\par
\noindent\emph{Iztočnik / Source:} Kontravariante/Kovariante\par
\noindent\emph{Rusky / Russian:} контравариант/ковариант\par
\noindent\emph{Ukrajinsky / Ukrainian:} контраваріант/коваріант\par
\smallskip\noindent\emph{Razlog:} Kontravariant/kovariant sohranja dva medžunarodna imennika v §11 raboty 2. Ruska i ukrajinska sravnjenja majut iste različne koreni; ne sut predložene dva konkurentne prěvody jednogo termina. Sravnjeny tekst razglědaje odredžene formy različno, kak kontravariant i kak kovariant, pri uslovjah kontrakcije. Zato ne trěba sjediniti oba pod nazvanjem invarianta ili podobna forma. Slovnično vysvětlenje mora povezati termin s konkretnym objektom i njegovym transformovanjem, a ne samo pokazati sličnost napisa.\par
\noindent\emph{Reason:} Kontravariant/kovariant preserves two international nouns in Paper 2, §11. The Russian and Ukrainian comparisons maintain the same distinction; these are not competing translations of one term. The compared text treats specified forms differently, as contravariant and covariant, in contraction conditions. They must therefore not be merged under invariant or similar form. Explanation must connect each term to its particular object and transformation behaviour, rather than merely displaying cognate spellings.\par
\noindent{\footnotesize Izvorne zapisy / Original records: HIST-g012-013.}\par
\smallskip\noindent\emph{Sravnjenje s tekstom:} Rabota 2 razglědaje $\theta$ kak kontravariant, a $H$ kak kovariant v uslovjah kontrakcij. Dvojica nazvanij označaje različne roli, ne dva slobodne sinonimy.\par
\noindent\emph{Textual comparison:} Paper 2 treats $\theta$ as contravariant and $H$ as covariant in contraction conditions. The names mark different roles, not interchangeable synonyms.\par
\noindent{\footnotesize TEXT-COMP-187; pasaže / passages: 1802, 2126; \texttt{CORPUS\_FORM\_COMPARISONS.json}.}\par
\end{minipage}\par
\par\medskip\Needspace{7\baselineskip}
\noindent\begin{minipage}{\linewidth}
\hypertarget{TERM-1054}{}
\noindent\textbf{kontravarianty}\quad{\footnotesize\texttt{TERM-1054}}\par
\noindent\emph{Iztočnik / Source:} Kontravarianten\par
\noindent\emph{Rusky / Russian:} контраварианты\par
\noindent\emph{Ukrajinsky / Ukrainian:} контраваріанти\par
\smallskip\noindent\emph{Razlog:} Kontravarianty prodolžuje medžunarodny korenj v §13 raboty 2. Ovdě zapis dava konkretny jezykovy iztočnik: Bedratyukovu rabotu math/0702732. Nova konsultacija jej §1, str. 1, potvrdžaje ukrajinske kontravarianty v pisemnoj formě контраваріанти; avtor razlikuje kovarianty i kontravarianty prěz razne kolca invariantov s prěměnnymi i dualnymi prěměnnymi. To podpiraje ukrajinsku tehničnu upotrěbu i smyslovu razliku, ne samo medžuslovjansko slovo ni jeho obšču razumljivost. Nova konsultacija ostaje oddělena od istorijnogo navodjenja iztočnika.\par
\noindent\emph{Reason:} Kontravarianty continues the international stem in Paper 2, §13. Here the record identifies a specific linguistic source: Bedratyuk, arXiv:math/0702732v1 (2007). Fresh examination of §1, p. 1 confirms Ukrainian контраваріанти and a distinction between covariants and contravariants through different invariant rings involving variables and dual variables. This supports the Ukrainian technical use and semantic distinction, not independent Interslavic adoption or general intelligibility. The new consultation is distinct from the earlier record's citation.\par
\noindent{\footnotesize Izvorne zapisy / Original records: HIST-g014-012.}\par
\smallskip\noindent\emph{Sravnjenje s tekstom:} Rabota 2 izrično izključaje produkty s kontravariantami v konkretnom koraku redukcije. Tehničny termin odnosi se k toj ulozi, ne samo k podobnosti s ukrajinskim slovom.\par
\noindent\emph{Textual comparison:} Paper 2 explicitly excludes products with contravariants at a particular reduction step. The technical term has this role, beyond its resemblance to the Ukrainian word.\par
\noindent{\footnotesize TEXT-COMP-191; pasaže / passages: 2093, 2126; \texttt{CORPUS\_FORM\_COMPARISONS.json}.}\par
\end{minipage}\par
\par\medskip\Needspace{7\baselineskip}
\noindent\begin{minipage}{\linewidth}
\hypertarget{TERM-1055}{}
\noindent\textbf{odgovarjajuče matrice}\quad{\footnotesize\texttt{TERM-1055}}\par
\noindent\emph{Iztočnik / Source:} korrespondierende Matrizen\par
\noindent\emph{Rusky / Russian:} соответствующие матрицы\par
\noindent\emph{Ukrajinsky / Ukrainian:} відповідні матриці\par
\smallskip\noindent\emph{Razlog:} Odgovarjajuče matrice jest opisna fraza, ne novy izmysljeny termin. Rabota 4 prodolžuje upotrěbu iz raboty 3. Sravnjeny teorem prirědžaje redu $p_\rho$ red $q_{n-\rho}$, svjazany s minorami i uslovjem $(u\mid x)=0$. Zato odgovarjanje tuto ne znači jednakost matric niti njih podobnost bez dodatnyh uslovij. Obyčna konstrukcija ostaje prikladna, zato že sam teorem dava točnu relaciju; slovnik ne trěbuje izmysljati osobny internacionalizm za vsako opisno sočetanje.\par
\noindent\emph{Reason:} Odgovarjajuče matrice is descriptive correspondence language, not a new coinage; Paper 4 continues its use in Paper 3. The compared theorem associates a row $p_\rho$ with a row $q_{n-\rho}$ through minors and the condition $(u\mid x)=0$. Correspondence here does not mean equality of matrices or similarity without further conditions. An ordinary phrase can suffice because the theorem specifies the relation; a separate internationalism need not be invented for every descriptive expression.\par
\noindent{\footnotesize Izvorne zapisy / Original records: HIST-g028-011, HIST-g030-003.}\par
\smallskip\noindent\emph{Sravnjenje s tekstom:} Teorem v rabotah 3 i 4 sojedinja redy $p_\rho$ i $q_{n-\rho}$ prěz minory i relaciju $(u\mid x)=0$. Obična fraza odgovarjajuče matrice opira se na to oprěděljeno odgovarjanje.\par
\noindent\emph{Textual comparison:} The theorem in Papers 3 and 4 relates rows $p_\rho$ and $q_{n-\rho}$ through minors and $(u\mid x)=0$. The ordinary phrase corresponding matrices relies on that specified correspondence.\par
\noindent{\footnotesize TEXT-COMP-188; pasaže / passages: 3653, 3742, 3782; \texttt{CORPUS\_FORM\_COMPARISONS.json}.}\par
\end{minipage}\par
\par\medskip\Needspace{7\baselineskip}
\noindent\begin{minipage}{\linewidth}
\hypertarget{TERM-1056}{}
\noindent\textbf{kovariantna derivacija}\quad{\footnotesize\texttt{TERM-1056}}\par
\noindent\emph{Iztočnik / Source:} kovariante Ableitung\par
\noindent\emph{Rusky / Russian:} ковариантная производная; ковариантная производная / ковариантное дифференцирование\par
\noindent\emph{Ukrajinsky / Ukrainian:} коваріантна похідна; коваріантна похідна / коваріантне диференціювання\par
\smallskip\noindent\emph{Razlog:} Kovariantna derivacija byla izbrana v rabotě 12 i prodolžena v rabotě 21. Prvy zapis upominaje domaču alternativu izvodna i ostavlja ju za jezykovu ocěnu; vtory opravdava kontinuum, ne novu nezavisnu selekciju. Tekst raboty 21 dava formulu za $h^{(1)}$, ktora odstranjaje druge diferencialy i ostavlja invariantnu formu. Zato ne trěba izjedniti termin s običnym parcialnym diferenciranjem. Ruska i ukrajinska dvojica derivata/diferenciranja pomaga razlikovati rezultat i proces, no sama ne izbira medžuslovjansky korenj.\par
\noindent\emph{Reason:} Kovariantna derivacija was selected in Paper 12 and continued in Paper 21. The first record explicitly leaves native izvodna open for linguistic appraisal; the second invokes continuity, not a new independent selection. Paper 21 supplies a formula for $h^{(1)}$ that eliminates second differentials while retaining an invariant form. The term therefore cannot be equated with ordinary partial differentiation. The Russian and Ukrainian derivative/differentiation pair helps distinguish result and process without itself settling the Interslavic stem.\par
\noindent{\footnotesize Izvorne zapisy / Original records: HIST-g074-012, HIST-g084-007.}\par
\smallskip\noindent\emph{Sravnjenje s tekstom:} Rabota 12 naziva proces kovariantnoju derivacijeju. Formula (143) raboty 21 izrično odstranjaje druge diferencialy; rabota 23 prěnosi konstrukciju na formy $[\Omega_i]$ i ih poslědovateljne derivacije. Kontinuum termina ne rěšaje jezykovy izbor medžu derivacija i izvodna.\par
\noindent\emph{Textual comparison:} Paper 12 names the process covariant differentiation. Paper 21's formula (143) explicitly eliminates second differentials; Paper 23 uses the construction for $[\Omega_i]$ and successive derivatives. Terminological continuity does not settle the linguistic choice between derivacija and izvodna.\par
\noindent{\footnotesize TEXT-COMP-189; pasaže / passages: 9061, 14117, 14119, 14120, 14121, 14124, 15197, 15199, 15201; \texttt{CORPUS\_FORM\_COMPARISONS.json}.}\par
\end{minipage}\par
\par\medskip\Needspace{7\baselineskip}
\noindent\begin{minipage}{\linewidth}
\hypertarget{TERM-1057}{}
\noindent\textbf{kovariantne derivacije}\quad{\footnotesize\texttt{TERM-1057}}\par
\noindent\emph{Iztočnik / Source:} kovariante Ableitungen\par
\noindent\emph{Rusky / Russian:} ковариантные производные\par
\noindent\emph{Ukrajinsky / Ukrainian:} коваріантні похідні\par
\smallskip\noindent\emph{Razlog:} Kovariantne derivacije v rabotě 23 jest množina ranějšego termina, ne nova operacija izbrana radi drugogo konca. Zapis izrično prěnosi diferencialno-geometričny registr iz raboty 21. V sravnjenom odlomku po odstranjenju vyšših diferencialov pojavjajut se formy $[\Omega_i]$ i ih poslědovateljne kovariantne derivacije. Te formy svedut pytanje k algebraičnym invariantam. Razlog čuva vezu medžu rabotami, ale ostavjaje otvorjeno ranějše jezykovědne pytanje o derivacija i izvodna.\par
\noindent\emph{Reason:} Kovariantne derivacije in Paper 23 is the plural of the earlier term, not a new operation selected because of an ending. The record explicitly carries over Paper 21's differential-geometric register. In the compared passage, elimination of higher differentials yields forms $[\Omega_i]$ and their successive covariant derivatives, reducing the problem to algebraic invariants. The reason preserves a connection between papers while leaving the earlier linguistic question about derivacija versus izvodna open.\par
\noindent{\footnotesize Izvorne zapisy / Original records: HIST-g094-015.}\par
\smallskip\noindent\emph{Sravnjenje s tekstom:} Rabota 12 naziva proces kovariantnoju derivacijeju. Formula (143) raboty 21 izrično odstranjaje druge diferencialy; rabota 23 prěnosi konstrukciju na formy $[\Omega_i]$ i ih poslědovateljne derivacije. Kontinuum termina ne rěšaje jezykovy izbor medžu derivacija i izvodna.\par
\noindent\emph{Textual comparison:} Paper 12 names the process covariant differentiation. Paper 21's formula (143) explicitly eliminates second differentials; Paper 23 uses the construction for $[\Omega_i]$ and successive derivatives. Terminological continuity does not settle the linguistic choice between derivacija and izvodna.\par
\noindent{\footnotesize TEXT-COMP-189; pasaže / passages: 9061, 14117, 14119, 14120, 14121, 14124, 15197, 15199, 15201; \texttt{CORPUS\_FORM\_COMPARISONS.json}.}\par
\end{minipage}\par
\par\medskip\Needspace{7\baselineskip}
\noindent\begin{minipage}{\linewidth}
\hypertarget{TERM-1058}{}
\noindent\textbf{ciklotomične polja}\quad{\footnotesize\texttt{TERM-1058}}\par
\noindent\emph{Iztočnik / Source:} Kreiskörper\par
\noindent\emph{Rusky / Russian:} круговые поля\par
\noindent\emph{Ukrajinsky / Ukrainian:} кругові поля\par
\smallskip\noindent\emph{Razlog:} Forma ciklotomične polja v rabotě 11 jest medžunarodna alternativa bukvalnomu kružno polje, sravnjenomu s ruskim круговое поле i ukrajinskim кругове поле. Zapis predpočita ciklotomične radi izběganja obyčnogo geometričnogo smysla kruga. Nyněšnji tekst upotrěbjaje toj korenj pri Abelovyh poljah i pri koeficientnom polju, izvedenom iz grupovyh harakterov. To jest argument o čitatelskoj orientaciji, ne izměreny rezultat. Drugy zapis naměrno izbira kružno; slovnik mora pokazati oba i razjasniti smysl v odlomku, ne izmisliti obče bezisključno pravilo.\par
\noindent\emph{Reason:} Ciklotomične polja in Paper 11 is the international alternative to the literal circle-field wording compared with Russian круговое поле and Ukrainian кругове поле. The record favours ciklotomične to avoid the ordinary geometric sense of circle. The current text uses it in discussion of Abelian fields and a coefficient field derived from group characters. This is a reader-orientation argument, not a measured result. Another record deliberately chooses kružno; both choices and their local senses must remain visible rather than being forced into an invented universal rule.\par
\noindent{\footnotesize Izvorne zapisy / Original records: HIST-g073-005.}\par
\smallskip\noindent\emph{Sravnjenje s tekstom:} Tekst ima ciklotomično v rabotě 11, kružno v 38 i kružno-dělitelno v 32. V poslednjem slučaji ciklično $K$ jest podpolje $P(\varepsilon)$, kde $\varepsilon$ jest primitivny korenj jedinice. Formy nazvanij i različne relacije vključenja ne trěba směšivati.\par
\noindent\emph{Textual comparison:} The text has ciklotomično in Paper 11, kružno in 38 and kružno-dělitelno in 32. In the last case the cyclic field $K$ is a subfield of $P(\varepsilon)$, with $\varepsilon$ a primitive root of unity. Variations in naming must not obscure the different inclusion relations.\par
\noindent{\footnotesize TEXT-COMP-190; pasaže / passages: 8724, 8904, 19616, 19620, 19638, 23185, 23188, 23192; \texttt{CORPUS\_FORM\_COMPARISONS.json}.}\par
\end{minipage}\par
\par\medskip\Needspace{7\baselineskip}
\noindent\begin{minipage}{\linewidth}
\hypertarget{TERM-1059}{}
\noindent\textbf{kružno polje}\quad{\footnotesize\texttt{TERM-1059}}\par
\noindent\emph{Iztočnik / Source:} Kreiskörper\par
\noindent\emph{Rusky / Russian:} круговое поле\par
\noindent\emph{Ukrajinsky / Ukrainian:} кругове поле\par
\smallskip\noindent\emph{Razlog:} Kružno polje v rabotě 38 jest naměrny slovjansky kalk. Tekst govori o kružnyh poljah i potom konkretno o polju $n^h$-tyh korenjev jednoty za dostatočno veliko $h$; pri poljah harakterov ukazuje takože na podpolja polja $n$-tyh korenjev jednoty. Zato objašnjenje mora čuvati različne uslovja o polju i podpolju, ne izvoditi jih iz metafory kruga. Izbor protivostavi se ranějšej preferenciji ciklotomične, ale ne označaje novu matematiku samo radi drugogo korenja.\par
\noindent\emph{Reason:} Kružno polje in Paper 38 is a deliberate Slavic calque. The text speaks of circle fields, then specifically of the field of $n^h$-th roots of unity for sufficiently large $h$; for character fields it also invokes subfields of the field of $n$-th roots. An explanation must preserve those field/subfield qualifications instead of deriving them from the circle metaphor. The choice contrasts with the earlier preference for ciklotomične; a different stem does not itself introduce different mathematics.\par
\noindent{\footnotesize Izvorne zapisy / Original records: HIST-g208-018.}\par
\smallskip\noindent\emph{Sravnjenje s tekstom:} Tekst ima ciklotomično v rabotě 11, kružno v 38 i kružno-dělitelno v 32. V poslednjem slučaji ciklično $K$ jest podpolje $P(\varepsilon)$, kde $\varepsilon$ jest primitivny korenj jedinice. Formy nazvanij i različne relacije vključenja ne trěba směšivati.\par
\noindent\emph{Textual comparison:} The text has ciklotomično in Paper 11, kružno in 38 and kružno-dělitelno in 32. In the last case the cyclic field $K$ is a subfield of $P(\varepsilon)$, with $\varepsilon$ a primitive root of unity. Variations in naming must not obscure the different inclusion relations.\par
\noindent{\footnotesize TEXT-COMP-190; pasaže / passages: 8724, 8904, 19616, 19620, 19638, 23185, 23188, 23192; \texttt{CORPUS\_FORM\_COMPARISONS.json}.}\par
\end{minipage}\par
\par\medskip\Needspace{7\baselineskip}
\noindent\begin{minipage}{\linewidth}
\hypertarget{TERM-1060}{}
\noindent\textbf{kružno-dělitelno polje}\quad{\footnotesize\texttt{TERM-1060}}\par
\noindent\emph{Iztočnik / Source:} Kreisteilungskörper\par
\noindent\emph{Rusky / Russian:} поле деления круга\par
\noindent\emph{Ukrajinsky / Ukrainian:} поле поділу кола\par
\smallskip\noindent\emph{Razlog:} Kružno-dělitelno polje v rabotě 32 sohranja složenu istorijnu formu německogo Kreisteilungskörper. Zapis izrično predpočita ju plnomu modernomu zaimstvovanju. V dokazě $\varepsilon$ jest primitivny $p$-ty korenj jedinice, a ciklično polje $K$ izbrano jest kak podpolje $P(\varepsilon)$. Trěba odlišiti kružno-dělitelno polje od libovoljnogo cikličnogo polja: tekst sojedinja je prěz vključenje, ne prěz sinonimiju. To jest tretja dokumentovana strategija v toj terminologičnoj familiji.\par
\noindent\emph{Reason:} Kružno-dělitelno polje in Paper 32 preserves the compound historical wording of Kreisteilungskörper. The record expressly prefers this to a fully modern borrowing. In the proof, $\varepsilon$ is a primitive $p$-th root of unity and the cyclic field $K$ is selected as a subfield of $P(\varepsilon)$. A cyclotomic field must be distinguished from an arbitrary cyclic field: the text connects them by inclusion, not synonymy. This is a third documented strategy within the terminology family.\par
\noindent{\footnotesize Izvorne zapisy / Original records: HIST-g175-007.}\par
\smallskip\noindent\emph{Sravnjenje s tekstom:} Tekst ima ciklotomično v rabotě 11, kružno v 38 i kružno-dělitelno v 32. V poslednjem slučaji ciklično $K$ jest podpolje $P(\varepsilon)$, kde $\varepsilon$ jest primitivny korenj jedinice. Formy nazvanij i različne relacije vključenja ne trěba směšivati.\par
\noindent\emph{Textual comparison:} The text has ciklotomično in Paper 11, kružno in 38 and kružno-dělitelno in 32. In the last case the cyclic field $K$ is a subfield of $P(\varepsilon)$, with $\varepsilon$ a primitive root of unity. Variations in naming must not obscure the different inclusion relations.\par
\noindent{\footnotesize TEXT-COMP-190; pasaže / passages: 8724, 8904, 19616, 19620, 19638, 23185, 23188, 23192; \texttt{CORPUS\_FORM\_COMPARISONS.json}.}\par
\end{minipage}\par
\par\medskip\Needspace{7\baselineskip}
\noindent\begin{minipage}{\linewidth}
\hypertarget{TERM-1061}{}
\noindent\textbf{Kroneckerova substitucija}\quad{\footnotesize\texttt{TERM-1061}}\par
\noindent\emph{Iztočnik / Source:} Kroneckersche Substitution\par
\noindent\emph{Rusky / Russian:} подстановка Kronecker'а\par
\noindent\emph{Ukrajinsky / Ukrainian:} підстановка Kronecker'а\par
\smallskip\noindent\emph{Razlog:} Kroneckerova substitucija spaja pridavnik od vlastnogo imena s internacionalnym nazvanjem operacije. V rabotě 20 ona zaměnjaje mnogo prěměnnyh stepenjami jednoj prěměnnoj; formula, ne samo ime, oprědělja postupok. Ruskij i ukrajinskij sravnjeni izrazy koristajut imenicu podstavjenja s imenom v zavisnom padežu. Oni podpirajut identifikaciju operacije, ale sami ne dokazujut, že substitucija jest najrazumělivějša vsim slovjanskim čitateljam. Arhivny argument o razumělivosti ostaje nepoměrenym očekivanjom.\par
\noindent\emph{Reason:} Kroneckerova substitucija combines a possessive adjective with an international name for the operation. In Paper 20, several variables are replaced by powers of one variable; the formula, not the name alone, specifies the procedure. The recorded Russian and Ukrainian expressions use a noun for substitution followed by the name in a dependent case. They help identify the operation but do not establish that substitucija is the most intelligible form for all Slavic readers. The archival appeal to intelligibility remains an unmeasured expectation.\par
\noindent{\footnotesize Izvorne zapisy / Original records: HIST-g083-006.}\par
\smallskip\noindent\emph{Sravnjenje s tekstom:} Formula v rabotě 20 zaměnjaje prěměnne stepenjami jednoj prěměnnoj. Pridavnik od imena ne jest zaměna za tu matematičnu specifikaciju.\par
\noindent\emph{Textual comparison:} Paper 20 replaces the variables by powers of one variable. The possessive naming construction is not a substitute for that mathematical specification.\par
\noindent{\footnotesize TEXT-COMP-192; pasaže / passages: 13863, 13865, 13867, 13872, 13873, 13874, 13875; \texttt{CORPUS\_FORM\_COMPARISONS.json}.}\par
\end{minipage}\par
\par\medskip\Needspace{7\baselineskip}
\noindent\begin{minipage}{\linewidth}
\hypertarget{TERM-1062}{}
\noindent\textbf{teorem Kronekera / relativno-Abelove polja}\quad{\footnotesize\texttt{TERM-1062}}\par
\noindent\emph{Iztočnik / Source:} Kroneckerscher Satz / relativ-Abelsche Körper\par
\noindent\emph{Rusky / Russian:} теорема Кронекера / относительно абелевы поля\par
\noindent\emph{Ukrajinsky / Ukrainian:} теорема Кронекера / відносно абелеві поля\par
\smallskip\noindent\emph{Razlog:} Teorem Kronekera i relativno-Abelove polja čuvajut vlastne imena i matematično ograničenje relativno. V rabotě 11 razlika medžu osnovnym poljem racionalnyh čisel i kvadratnymi čislovymi poljami jest bitna: vsak teorem ima svoj obseg. Kronekera jest prilagodžena imenska forma, dokolě drugi zapisy koristajut Kroneckerov. To sut različne jezykove realizacije imena, ne različni matematici ni nezavisne terminologične attestacije. Vidimost vlastnogo imena ne opravdyvaje izpuščenje uslovij teorema.\par
\noindent\emph{Reason:} Teorem Kronekera and relativno-Abelove polja retain the names and the mathematical qualification relative. Paper 11 distinguishes the rational base field from quadratic number fields; each theorem has its own scope. Kronekera adapts the name, whereas other records use Kroneckerov. These are different linguistic realizations of one name, not different mathematicians or independent terminological attestations. Recognizability of the name does not justify omitting the theorem’s conditions.\par
\noindent{\footnotesize Izvorne zapisy / Original records: HIST-g073-004.}\par
\smallskip\noindent\emph{Sravnjenje s tekstom:} Rabota 11 prilagodžaje ime kako Kronekera; rabota 6 koristi Kroneckerov. Oba slučaja trěbujut čuvati specialny sadržanje nazvanogo rezultata.\par
\noindent\emph{Textual comparison:} Paper 11 adapts the name as Kronekera; Paper 6 uses Kroneckerov. Both require preservation of the particular content of the named result.\par
\noindent{\footnotesize TEXT-COMP-193; pasaže / passages: 8724, 6486; \texttt{CORPUS\_FORM\_COMPARISONS.json}.}\par
\end{minipage}\par
\par\medskip\Needspace{7\baselineskip}
\noindent\begin{minipage}{\linewidth}
\hypertarget{TERM-1063}{}
\noindent\textbf{Kroneckerov fundamentalny sistem}\quad{\footnotesize\texttt{TERM-1063}}\par
\noindent\emph{Iztočnik / Source:} Kroneckersches Fundamentalsystem\par
\noindent\emph{Rusky / Russian:} кронекерова фундаментальная система\par
\noindent\emph{Ukrajinsky / Ukrainian:} Кронекерова фундаментальна система\par
\smallskip\noindent\emph{Razlog:} Kroneckerov fundamentalny sistem označuje nazvany sistem algebraično cělyh veličin. Arhivny razlog jest internacionalna imenska konstrukcija, priměnjena v dokazu pomoćnogo teorema. Kontekst §14 raboty 6 veže taj sistem s elementarnymi simetričnymi funkcijami i algebraično cěloju zavisimostju. Zato fundamentalny ne znači prosto važny, a sistem ne směje byti bez objašnjenja zaměnjen bazoju v libovoljnom smyslu. Zachovanje imena i konstrukcije podpiraje svez s citovanym rezultatom, ne tvrdženje o novoj medžuslovjanskoj kovanici.\par
\noindent\emph{Reason:} Kroneckerov fundamentalny sistem names a system of algebraically integral quantities. The recorded reason is an international naming construction used in an auxiliary theorem. Paper 6, §14 connects the system with elementary symmetric functions and algebraic integrality. Fundamental therefore does not simply mean important, nor may system be replaced without explanation by basis in an arbitrary sense. Retaining the name and construction connects the passage with the cited result; it is not evidence of a newly coined Interslavic expression.\par
\noindent{\footnotesize Izvorne zapisy / Original records: HIST-g054-012.}\par
\smallskip\noindent\emph{Sravnjenje s tekstom:} Rabota 11 prilagodžaje ime kako Kronekera; rabota 6 koristi Kroneckerov. Oba slučaja trěbujut čuvati specialny sadržanje nazvanogo rezultata.\par
\noindent\emph{Textual comparison:} Paper 11 adapts the name as Kronekera; Paper 6 uses Kroneckerov. Both require preservation of the particular content of the named result.\par
\noindent{\footnotesize TEXT-COMP-193; pasaže / passages: 8724, 6486; \texttt{CORPUS\_FORM\_COMPARISONS.json}.}\par
\end{minipage}\par
\par\medskip\Needspace{7\baselineskip}
\noindent\begin{minipage}{\linewidth}
\hypertarget{TERM-1064}{}
\noindent\textbf{forma kriviny}\quad{\footnotesize\texttt{TERM-1064}}\par
\noindent\emph{Iztočnik / Source:} Krümmungsform\par
\noindent\emph{Rusky / Russian:} форма кривизны\par
\noindent\emph{Ukrajinsky / Ukrainian:} форма кривини\par
\smallskip\noindent\emph{Razlog:} Forma kriviny jest izbrana vměsto modernizujučej tensornoj parafrazy zato, že termin izvornika govori o formě. V rabotě 12 forma Riemanna ima specialno město pri kvadratičnyh formah v teoremu redukcije; tekst raboty 23 ponovno veže ju s istim slučajem. Ruskij izraz s kriviznoju i ukrajinskij s krivinoju podpirajut geometričny smysl kriviny, ale ne izjednačajut vsaku formu s tensorom. Izbor čuvaje istoriju matematičnogo izloženja, ne zastarělu pravopisnu osoblivost.\par
\noindent\emph{Reason:} Forma kriviny was preferred to a modernizing tensor paraphrase because the source names a form. In Paper 12, Riemann’s form has a special role for quadratic forms in the reduction theorem; Paper 23 connects it with the same case. The Russian and Ukrainian comparanda support the meaning curvature, not an unrestricted equivalence between form and tensor. The choice preserves the mathematical exposition’s historical vocabulary, rather than merely an old spelling.\par
\noindent{\footnotesize Izvorne zapisy / Original records: HIST-g074-018.}\par
\smallskip\noindent\emph{Sravnjenje s tekstom:} Raboty 12 i 23 vežut formu Riemanna s kvadratičnym slučajem. Moderna tensorna parafraza ne jest čisto jezykova zaměna bez nužnosti objašnjenja.\par
\noindent\emph{Textual comparison:} Papers 12 and 23 connect Riemann’s form with the quadratic case. A modern tensor paraphrase is not a purely linguistic substitution needing no explanation.\par
\noindent{\footnotesize TEXT-COMP-194; pasaže / passages: 9152, 15203; \texttt{CORPUS\_FORM\_COMPARISONS.json}.}\par
\end{minipage}\par
\par\medskip\Needspace{7\baselineskip}
\noindent\begin{minipage}{\linewidth}
\hypertarget{TERM-1066}{}
\noindent\textbf{polje}\quad{\footnotesize\texttt{TERM-1066}}\par
\noindent\emph{Iztočnik / Source:} Körper\par
\noindent\emph{Rusky / Russian:} поле\par
\noindent\emph{Ukrajinsky / Ukrainian:} поле\par
\smallskip\noindent\emph{Razlog:} Polje prodolžaje ustanovjene nazvanje komutativnogo algebraičnogo objekta. Dva zapisa davajut različne razloge: kontinuitet raboty 6 i matematičny dokaz v rabotě 30, §7. V drugom slučaju odsutnost děliteljev nuly daje jedinost rěšenja jednačenja $ax=b$ pri nenulovom $a$; uslovje cěpov jest takože potrěbno v navedenom dokazu obstajanja. Zato kratko slovo polje ne směje prikryti uslovja, od ktoryh vyhodi zaključok. Samo podobnost k ruskomu i ukrajinskomu pole ne izpolnjaje tu matematičnu kontrolju.\par
\noindent\emph{Reason:} Polje continues the established name for the commutative algebraic object. The two records give different reasons: continuity within Paper 6 and the proof in Paper 30, §7. In the latter, absence of zero divisors gives uniqueness of a solution to $ax=b$ for nonzero $a$; the chain condition is also needed in the stated existence argument. The compact word field must not obscure the assumptions supporting the conclusion. Resemblance to Russian and Ukrainian pole cannot replace that mathematical check.\par
\noindent{\footnotesize Izvorne zapisy / Original records: HIST-g047-004, HIST-g114-011.}\par
\smallskip\noindent\emph{Sravnjenje s tekstom:} Rabota 30, §7 koristi uslovje cěpov zajedno s odsutnostju děliteljev nuly. Poslědnje uslovje samo ne směje zastupiti cěly dokaz; nynějšni imennik jest prsten, ne arhivny kolco.\par
\noindent\emph{Textual comparison:} Paper 30, §7 uses the chain condition together with absence of zero divisors. The latter must not stand in for the whole proof; the current noun is prsten, not archival kolco.\par
\noindent{\footnotesize TEXT-COMP-196; pasaže / passages: 16761, 16763, 16772, 16774, 16778; \texttt{CORPUS\_FORM\_COMPARISONS.json}.}\par
\end{minipage}\par
\par\medskip\Needspace{7\baselineskip}
\noindent\begin{minipage}{\linewidth}
\hypertarget{TERM-1067}{}
\noindent\textbf{tělo}\quad{\footnotesize\texttt{TERM-1067}}\par
\noindent\emph{Iztočnik / Source:} Körper\par
\noindent\emph{Rusky / Russian:} тело\par
\noindent\emph{Ukrajinsky / Ukrainian:} тіло\par
\smallskip\noindent\emph{Razlog:} Tělo v rabotě 34 označuje strukturu s obratnymi nenulovymi elementami bez občeho prědpoloženja komutativnosti. Zapisy iz §§2 i 14 zato ne forsujut polje. To ne znači, že vsaky objekt nazvany tělom jest strogo nekomutativny: definicija dopušča takože komutativny slučaj. Matematičny znak jest obratimost nenulovyh elementov, ne metafora těla. Pozdnějši izraz komutativno tělo v §20 potvrdžaje, že prosty bezuslovny slovnik Körper=polje byl by nedostatočny.\par
\noindent\emph{Reason:} Tělo in Paper 34 covers a division structure without a general assumption of commutativity. The records for §§2 and 14 therefore resist forcing field terminology. This does not mean that every object called tělo is strictly noncommutative: the definition permits the commutative case. The mathematical criterion is invertibility of nonzero elements, not the metaphor of a body. The later phrase komutativno tělo in §20 shows why an unconditional lookup rule Körper=polje would be inadequate.\par
\noindent{\footnotesize Izvorne zapisy / Original records: HIST-g180-006, HIST-g192-001.}\par
\smallskip\noindent\emph{Sravnjenje s tekstom:} Rabota 34 dopušča těla bez prědpoloženja komutativnosti, a pri definiciji hiperkompleksnogo sistema dodaje komutativnost. Izraz komutativno tělo v §20 jest protivdokaz beziznimkovomu leksičnomu razděljenju, ne protivdokaz matematičnoj definiciji.\par
\noindent\emph{Textual comparison:} Paper 34 permits division rings without assuming commutativity, then adds commutativity in the definition of a hypercomplex system. Komutativno tělo in §20 contradicts an exceptionless lexical allocation, not the mathematical definition.\par
\noindent{\footnotesize TEXT-COMP-197; pasaže / passages: 19821, 19916, 19918, 20221, 20268, 21228, 21794, 19744, 19746, 19756; \texttt{CORPUS\_FORM\_COMPARISONS.json}.}\par
\end{minipage}\par
\par\medskip\Needspace{7\baselineskip}
\noindent\begin{minipage}{\linewidth}
\hypertarget{TERM-1068}{}
\noindent\textbf{tělo}\quad{\footnotesize\texttt{TERM-1068}}\par
\noindent\emph{Iztočnik / Source:} Körper (nicht notwendig kommutativ)\par
\noindent\emph{Rusky / Russian:} тело\par
\noindent\emph{Ukrajinsky / Ukrainian:} тіло\par
\smallskip\noindent\emph{Razlog:} Tělo v početnom razgledanju modulov §6 raboty 34 ne namečuje komutativnost koeficientov. Pozdnějša definicija hiperkompleksnogo sistema v tom samom paragrafu izrično trěbuje komutativno polje $K$. Razlika zato jest lokalna i dokazno važna: obče modulno razgledanje i specialna definicija ne imajut iste uslovja. Prevod mora čuvati město, na ktorom nove uslovje vvodit se; ne može dodati jego už k prvomu pojavjenju nazvanja.\par
\noindent\emph{Reason:} Tělo in the opening module discussion of Paper 34, §6 does not impose commutativity on the coefficients. The later definition of a hypercomplex system in that section explicitly requires a commutative field $K$. The distinction is local and mathematically consequential: the general module discussion and the special definition have different assumptions. Translation must retain the point at which the additional condition is introduced, rather than importing it at the first occurrence of the noun.\par
\noindent{\footnotesize Izvorne zapisy / Original records: HIST-g184-001.}\par
\smallskip\noindent\emph{Sravnjenje s tekstom:} Rabota 34 dopušča těla bez prědpoloženja komutativnosti, a pri definiciji hiperkompleksnogo sistema dodaje komutativnost. Izraz komutativno tělo v §20 jest protivdokaz beziznimkovomu leksičnomu razděljenju, ne protivdokaz matematičnoj definiciji.\par
\noindent\emph{Textual comparison:} Paper 34 permits division rings without assuming commutativity, then adds commutativity in the definition of a hypercomplex system. Komutativno tělo in §20 contradicts an exceptionless lexical allocation, not the mathematical definition.\par
\noindent{\footnotesize TEXT-COMP-197; pasaže / passages: 19821, 19916, 19918, 20221, 20268, 21228, 21794, 19744, 19746, 19756; \texttt{CORPUS\_FORM\_COMPARISONS.json}.}\par
\end{minipage}\par
\par\medskip\Needspace{7\baselineskip}
\noindent\begin{minipage}{\linewidth}
\hypertarget{TERM-1069}{}
\noindent\textbf{tělo}\quad{\footnotesize\texttt{TERM-1069}}\par
\noindent\emph{Iztočnik / Source:} Körper (noncommutative)\par
\noindent\emph{Rusky / Russian:} тело\par
\noindent\emph{Ukrajinsky / Ukrainian:} тіло\par
\smallskip\noindent\emph{Razlog:} Arhivny naslovič Körper (noncommutative) jest kratša orientacija, ne doslovna izjava uslovja. Podnožna primětka §1 raboty 34 kaže, že komutativny zakon ne prědpoklada se. To jest slabějše od izjavy, že zakon ne jest izpolnjen: komutativne priměry ne sut izključene. Tělo zato čuvaje obči smysl struktury s obratnymi nenulovymi elementami. Razumna terminologična razlika ne směje sama stvoriti novo matematično ograničenje.\par
\noindent\emph{Reason:} The archival heading Körper (noncommutative) is shorthand, not a literal statement of the hypothesis. The footnote in Paper 34, §1 says that commutativity is not assumed. That is weaker than saying it fails: commutative examples remain allowed. Tělo therefore retains the general division-ring sense. A useful terminological distinction must not itself create an additional mathematical restriction.\par
\noindent{\footnotesize Izvorne zapisy / Original records: HIST-g179-004.}\par
\smallskip\noindent\emph{Sravnjenje s tekstom:} Rabota 34 dopušča těla bez prědpoloženja komutativnosti, a pri definiciji hiperkompleksnogo sistema dodaje komutativnost. Izraz komutativno tělo v §20 jest protivdokaz beziznimkovomu leksičnomu razděljenju, ne protivdokaz matematičnoj definiciji.\par
\noindent\emph{Textual comparison:} Paper 34 permits division rings without assuming commutativity, then adds commutativity in the definition of a hypercomplex system. Komutativno tělo in §20 contradicts an exceptionless lexical allocation, not the mathematical definition.\par
\noindent{\footnotesize TEXT-COMP-197; pasaže / passages: 19821, 19916, 19918, 20221, 20268, 21228, 21794, 19744, 19746, 19756; \texttt{CORPUS\_FORM\_COMPARISONS.json}.}\par
\end{minipage}\par
\par\medskip\Needspace{7\baselineskip}
\noindent\begin{minipage}{\linewidth}
\hypertarget{TERM-1070}{}
\noindent\textbf{polje / polje funkcij / čislovo polje}\quad{\footnotesize\texttt{TERM-1070}}\par
\noindent\emph{Iztočnik / Source:} Körper / Funktionenkörper / Zahlkörper\par
\noindent\emph{Rusky / Russian:} поле / поле функций / числовое поле\par
\noindent\emph{Ukrajinsky / Ukrainian:} поле / поле функцій / числове поле\par
\smallskip\noindent\emph{Razlog:} Polje, polje funkcij i čislovo polje tvoret svezanu terminologičnu rodinu. Bibliografičny zapis govori o naměri razděljati komutativne polja od obče pojatyh těl. To jest razumna smernica, ale ne točny opis vsego sravnjenogo teksta: rabota 34, §20 ima komutativno tělo. Zato pravilo trěba čitati s kontekstom i iznimkami, ne kak avtomatičnu zaměnu. Bibliografična povtornost sama po sobě ne jest nova nezavisna attestacija.\par
\noindent\emph{Reason:} Polje, polje funkcij and čislovo polje form a related terminological family. The bibliographical record states an intention to distinguish commutative fields from division rings in the broader sense. This is a useful guideline, but not an exact description of the entire compared text: Paper 34, §20 has komutativno tělo. The rule must therefore be read with its contexts and exceptions, not as an automatic substitution. Bibliographical repetition is not a new independent attestation.\par
\noindent{\footnotesize Izvorne zapisy / Original records: HIST-g215-008.}\par
\smallskip\noindent\emph{Sravnjenje s tekstom:} Rabota 34 dopušča těla bez prědpoloženja komutativnosti, a pri definiciji hiperkompleksnogo sistema dodaje komutativnost. Izraz komutativno tělo v §20 jest protivdokaz beziznimkovomu leksičnomu razděljenju, ne protivdokaz matematičnoj definiciji.\par
\noindent\emph{Textual comparison:} Paper 34 permits division rings without assuming commutativity, then adds commutativity in the definition of a hypercomplex system. Komutativno tělo in §20 contradicts an exceptionless lexical allocation, not the mathematical definition.\par
\noindent{\footnotesize TEXT-COMP-197; pasaže / passages: 19821, 19916, 19918, 20221, 20268, 21228, 21794, 19744, 19746, 19756; \texttt{CORPUS\_FORM\_COMPARISONS.json}.}\par
\end{minipage}\par
\par\medskip\Needspace{7\baselineskip}
\noindent\begin{minipage}{\linewidth}
\hypertarget{TERM-1071}{}
\noindent\textbf{polje / kolco klasov ostatkov bez děliteljev nuly}\quad{\footnotesize\texttt{TERM-1071}}\par
\noindent\emph{Iztočnik / Source:} Körper / Restklassenring ohne Nullteiler\par
\noindent\emph{Rusky / Russian:} поле / кольцо классов вычетов без делителей нуля\par
\noindent\emph{Ukrajinsky / Ukrainian:} поле / кільце класів лишків без дільників нуля\par
\smallskip\noindent\emph{Razlog:} Polje i prsten klasov ostatkov bez děliteljev nuly ne sut bezuslovne sinonimy. V §7 raboty 30 zaključok, že prsten klasov ostatkov jest polje, koristi takože uslovje cěpov kratnyh. Arhivna forma kolco jest v sravnjenom tekstu prsten; ta leksična razlika ne měnja logičnu potrěbu zachovati dopolnitelno uslovje. Izbor terminov mora pokazati prěhod v dokazu, ne izbrisati jego kak že dokončenu definiciju.\par
\noindent\emph{Reason:} Field and a residue-class ring without zero divisors are not unconditional synonyms. In Paper 30, §7 the conclusion that the residue-class ring is a field also uses the multiple-chain condition. The archival kolco appears as prsten in the compared text; that lexical difference does not remove the need to retain the additional assumption. Terminology must express the transition made by the proof, not erase it as though it were already a definition.\par
\noindent{\footnotesize Izvorne zapisy / Original records: HIST-g114-005.}\par
\smallskip\noindent\emph{Sravnjenje s tekstom:} Rabota 30, §7 koristi uslovje cěpov zajedno s odsutnostju děliteljev nuly. Poslědnje uslovje samo ne směje zastupiti cěly dokaz; nynějšni imennik jest prsten, ne arhivny kolco.\par
\noindent\emph{Textual comparison:} Paper 30, §7 uses the chain condition together with absence of zero divisors. The latter must not stand in for the whole proof; the current noun is prsten, not archival kolco.\par
\noindent{\footnotesize TEXT-COMP-196; pasaže / passages: 16761, 16763, 16772, 16774, 16778; \texttt{CORPUS\_FORM\_COMPARISONS.json}.}\par
\end{minipage}\par
\par\medskip\Needspace{7\baselineskip}
\noindent\begin{minipage}{\linewidth}
\hypertarget{TERM-1072}{}
\noindent\textbf{polje vsih kompleksnyh čisel}\quad{\footnotesize\texttt{TERM-1072}}\par
\noindent\emph{Iztočnik / Source:} Körper aller komplexen Zahlen\par
\noindent\emph{Rusky / Russian:} поле всех комплексных чисел\par
\noindent\emph{Ukrajinsky / Ukrainian:} поле всіх комплексних чисел\par
\smallskip\noindent\emph{Razlog:} Polje vsih kompleksnyh čisel imenuje obseg konstrukcije v §7 raboty 9. Vsih jest važno: ne ide prosto o někoje podpolje kompleksnyh čisel. Izbor polje govori v poznanoj algebraičnoj terminologiji, a ne doslovno prenosi metaforu německogo Körper. Sam naziv ješče ne kaže, kako racionalna baza konstruuje se; taj postupok i zadany dobry poredok morajut ostati v izloženju.\par
\noindent\emph{Reason:} Polje vsih kompleksnyh čisel specifies the domain of the construction in Paper 9, §7. All matters: the passage is not merely about some subfield of the complex numbers. Polje uses familiar algebraic terminology rather than literally reproducing the German metaphor Körper. The name alone does not explain the construction of the rational basis; that argument and its specified well-ordering must remain in the exposition.\par
\noindent{\footnotesize Izvorne zapisy / Original records: HIST-g065-003.}\par
\smallskip\noindent\emph{Sravnjenje s tekstom:} Rabota 9, §7 odlišaje obči predmet konstrukcije od specialnyh baz v primětce. Kvantifikacija vsih i zadany poredok ne sut ukrasne česti izraza.\par
\noindent\emph{Textual comparison:} Paper 9, §7 distinguishes the general construction from the special bases discussed in the footnote. All and the specified ordering are substantive parts of the statement.\par
\noindent{\footnotesize TEXT-COMP-198; pasaže / passages: 7843, 7844, 7845, 7846, 7847, 7848, 7849, 7850, 7851, 7852, 7854, 7855; \texttt{CORPUS\_FORM\_COMPARISONS.json}.}\par
\end{minipage}\par
\par\medskip\Needspace{7\baselineskip}
\noindent\begin{minipage}{\linewidth}
\hypertarget{TERM-1073}{}
\noindent\textbf{polje vsih kompleksnyh čisel / čislovo polje / tělo}\quad{\footnotesize\texttt{TERM-1073}}\par
\noindent\emph{Iztočnik / Source:} Körper aller komplexen Zahlen / Zahlkörper / nicht notwendig kommutativer Körper\par
\noindent\emph{Rusky / Russian:} поле всех комплексных чисел / числовое поле / тело\par
\noindent\emph{Ukrajinsky / Ukrainian:} поле всіх комплексних чисел / числове поле / тіло\par
\smallskip\noindent\emph{Razlog:} Tri izrazy iz vvedenja raboty 34 razlikujut polje vsih kompleksnyh čisel, čislove polja i těla, za ktore komutativnost ne jest potrěbna. To ne jest prosta pravopisna varijacija: različne oblasti koeficientov igrat različne roli v teoriji prědstavjenij. Izbor polje/tělo jest opravdany tekstovym značenjem, ale ne daje licenciju na slepu globalnu zaměnu. Obči smysl tělo može uključati komutativne slučaje, kako pokazyvajut drugde izrične uslovja.\par
\noindent\emph{Reason:} The three expressions in the introduction to Paper 34 distinguish the field of all complex numbers, number fields, and division rings not required to be commutative. This is not mere spelling variation: different coefficient domains play different roles in representation theory. The polje/tělo distinction is justified by the passage’s meaning, but does not license blind global replacement. The broader sense of tělo can include commutative cases, as explicit qualifications elsewhere show.\par
\noindent{\footnotesize Izvorne zapisy / Original records: HIST-g178-003.}\par
\smallskip\noindent\emph{Sravnjenje s tekstom:} Rabota 34 dopušča těla bez prědpoloženja komutativnosti, a pri definiciji hiperkompleksnogo sistema dodaje komutativnost. Izraz komutativno tělo v §20 jest protivdokaz beziznimkovomu leksičnomu razděljenju, ne protivdokaz matematičnoj definiciji.\par
\noindent\emph{Textual comparison:} Paper 34 permits division rings without assuming commutativity, then adds commutativity in the definition of a hypercomplex system. Komutativno tělo in §20 contradicts an exceptionless lexical allocation, not the mathematical definition.\par
\noindent{\footnotesize TEXT-COMP-197; pasaže / passages: 19821, 19916, 19918, 20221, 20268, 21228, 21794, 19744, 19746, 19756; \texttt{CORPUS\_FORM\_COMPARISONS.json}.}\par
\end{minipage}\par
\par\medskip\Needspace{7\baselineskip}
\noindent\begin{minipage}{\linewidth}
\hypertarget{TERM-1074}{}
\noindent\textbf{polja $\mathfrak{K}_{n\rho}$ [rekonstrukcija indeksa / reconstructed subscript]}\quad{\footnotesize\texttt{TERM-1074}}\par
\noindent\emph{Iztočnik / Source:} Körper $\mathfrak{K}_{n\rho}$ [rekonstrukcija indeksa / reconstructed subscript]\par
\noindent\emph{Rusky / Russian:} поля $\mathfrak{K}_{n\rho}$ [rekonstrukcija indeksa / reconstructed subscript]\par
\noindent\emph{Ukrajinsky / Ukrainian:} поля $\mathfrak{K}_{n\rho}$ [rekonstrukcija indeksa / reconstructed subscript]\par
\smallskip\noindent\emph{Razlog:} Nazvanje polja jest čitljivo, ale arhivny indeks K\_\{n ho\} jest poškodženy. V §1 raboty 6 polje jest označeno frakturnym K s indeksami n i rho; n čisli neoprěděljene, a rho jest algebraičny rang. Čitateljsko zapisanje indeksa jest zato označeno kako rekonstrukcija iz sravnjenogo teksta, ne kak točna kopija arhivnogo niza. Arhivna izjava o zachovanoj notaciji sama po sobě ne dokazuje, že notacija preživěla v vsakoj izvedenoj tabeli.\par
\noindent\emph{Reason:} The word for field is readable, but the archival subscript K\_\{n ho\} is damaged. Paper 6, §1 denotes the field by a fraktur K with subscripts n and rho: n counts the indeterminates and rho is the algebraic rank. The reader’s typeset subscript is therefore marked as reconstructed from the compared text, not as a literal copy of the archival string. An archival assertion that notation was preserved does not establish that it survived in every derived table.\par
\noindent{\footnotesize Izvorne zapisy / Original records: HIST-g041-001.}\par
\smallskip\noindent\emph{Sravnjenje s tekstom:} V rabotě 6 sistemy i polja ne sut sinonimy. Notacija ima frakturny K s n i rho; rho jest algebraičny rang. Arhivny niz K\_\{n ho\} jest zachovany v danych, a rekonstrukcija indeksa iz teksta jest izrično označena.\par
\noindent\emph{Textual comparison:} In Paper 6, systems and fields are not synonyms. The notation is fraktur K with n and rho, with rho the algebraic rank. The damaged archival string K\_\{n ho\} remains in the data; the subscript reconstructed from the text is explicitly identified.\par
\noindent{\footnotesize TEXT-COMP-199; pasaže / passages: 4998, 5051, 5057, 5067, 5070, 5099, 5115, 5117; \texttt{CORPUS\_FORM\_COMPARISONS.json}.}\par
\end{minipage}\par
\par\medskip\Needspace{7\baselineskip}
\noindent\begin{minipage}{\linewidth}
\hypertarget{TERM-1075}{}
\noindent\textbf{polje racionalnyh funkcij}\quad{\footnotesize\texttt{TERM-1075}}\par
\noindent\emph{Iztočnik / Source:} Körper rationaler Funktionen\par
\noindent\emph{Rusky / Russian:} поле рациональных функций\par
\noindent\emph{Ukrajinsky / Ukrainian:} поле раціональних функцій\par
\smallskip\noindent\emph{Razlog:} Polje racionalnyh funkcij v rabotě 23 jest oblast, iz ktorej berut se vsi polinomy, ne prosto množstvo vsih polinomov bez ograničenja. Relacija pripadnosti jest matematično sadržanje izbora. Arhivny zapis opisuje ranějšu preširoku formulaciju; sravnjena pasaža nyně izrično govori o polinomah danogo polja racionalnyh funkcij. Sama ta pasaža potvrdžaje nynějšu relaciju, ale bez prědhodnogo teksta ne dokazuje točny cěly prěběg ranějšej proměny.\par
\noindent\emph{Reason:} Polje racionalnyh funkcij in Paper 23 is the field from which all the polynomials under consideration are taken, not an unrestricted set of all polynomials. Membership is the mathematical substance of the choice. The archival record describes an earlier overbroad wording; the compared passage explicitly refers to the polynomials of a field of rational functions. This establishes the present relation, but without the predecessor text it does not establish the complete sequence of the earlier revision.\par
\noindent{\footnotesize Izvorne zapisy / Original records: HIST-g094-021.}\par
\smallskip\noindent\emph{Sravnjenje s tekstom:} Rabota 23 govori o vsih polinomah prinaležečih danomu polju racionalnyh funkcij. To ograničenje ne možno zaměniti neograničenoju množinoju polinomov.\par
\noindent\emph{Textual comparison:} Paper 23 concerns all the polynomials belonging to a given field of rational functions. The restriction cannot be replaced by an unrestricted collection of polynomials.\par
\noindent{\footnotesize TEXT-COMP-200; pasaže / passages: 15164; \texttt{CORPUS\_FORM\_COMPARISONS.json}.}\par
\end{minipage}\par
\par\medskip\Needspace{7\baselineskip}
\noindent\begin{minipage}{\linewidth}
\hypertarget{TERM-1076}{}
\noindent\textbf{Polja i sistemy racionalnyh funkcij}\quad{\footnotesize\texttt{TERM-1076}}\par
\noindent\emph{Iztočnik / Source:} Körper und Systeme rationaler Funktionen\par
\noindent\emph{Rusky / Russian:} Поля и системы рациональных функций\par
\noindent\emph{Ukrajinsky / Ukrainian:} Поля та системи раціональних функцій\par
\smallskip\noindent\emph{Razlog:} Polja i sistemy racionalnyh funkcij čuvaje oba člena naslova raboty 6. Polje jest specialny sistem s oprěděljenymi algebraičnymi svojstvami; sistemy ne smějut bez razloga izčeznuti iz naslova ni stati sinonimom polj. Izbor polje vměsto doslovnogo tělo prědaje komutativny funkcijny kontekst. Razumělivost naslova zavisi od različenja objektov, ne samo od podobnosti posamečnyh slov.\par
\noindent\emph{Reason:} Polja i sistemy racionalnyh funkcij retains both components of Paper 6’s title. A field is a system with specified algebraic properties; systems must neither disappear from the title nor become synonyms of fields. Polje rather than a literal body expresses the commutative function-field context. The title’s clarity depends on distinguishing the objects, not merely on resemblance between individual words.\par
\noindent{\footnotesize Izvorne zapisy / Original records: HIST-g040-001.}\par
\smallskip\noindent\emph{Sravnjenje s tekstom:} V rabotě 6 sistemy i polja ne sut sinonimy. Notacija ima frakturny K s n i rho; rho jest algebraičny rang. Arhivny niz K\_\{n ho\} jest zachovany v danych, a rekonstrukcija indeksa iz teksta jest izrično označena.\par
\noindent\emph{Textual comparison:} In Paper 6, systems and fields are not synonyms. The notation is fraktur K with n and rho, with rho the algebraic rank. The damaged archival string K\_\{n ho\} remains in the data; the subscript reconstructed from the text is explicitly identified.\par
\noindent{\footnotesize TEXT-COMP-199; pasaže / passages: 4998, 5051, 5057, 5067, 5070, 5099, 5115, 5117; \texttt{CORPUS\_FORM\_COMPARISONS.json}.}\par
\end{minipage}\par
\par\medskip\Needspace{7\baselineskip}
\noindent\begin{minipage}{\linewidth}
\hypertarget{TERM-1077}{}
\noindent\textbf{najkrajše predstavjenje}\quad{\footnotesize\texttt{TERM-1077}}\par
\noindent\emph{Iztočnik / Source:} kürzeste Darstellung\par
\noindent\emph{Rusky / Russian:} кратчайшее представление\par
\noindent\emph{Ukrajinsky / Ukrainian:} найкоротше подання\par
\smallskip\noindent\emph{Razlog:} Najkrajše predstavjenje jest arhivna forma za irredundantno prědstavjenje ideala, ne za stilistično kratku rečenicu. Definicija §5 raboty 30 kaže, že nijedna komponenta ne može byti izpuščena. Sravnjeni tekst ima najkratše prědstavjenje, zato slovnik mora razlikovati arhivnu oblikovnu varijantu od formy v korpusu. Izmenjenje pisanja ne izmenjaje definicionny test; nijedna od tih form ne znači samo najmenše čislo pisanih symbolov.\par
\noindent\emph{Reason:} Najkrajše predstavjenje is the archival form for an irredundant ideal representation, not a stylistically short sentence. Paper 30, §5 defines it by the impossibility of omitting a component. The compared text has najkratše prědstavjenje, so the dictionary must distinguish the recorded variant from the corpus form. The spelling difference does not change the defining test; neither form simply means the smallest number of written symbols.\par
\noindent{\footnotesize Izvorne zapisy / Original records: HIST-g112-010.}\par
\smallskip\noindent\emph{Sravnjenje s tekstom:} Rabota 19 razlikuje dva uslovja redukovanosti; najkratše koristi prvo. Rabota 30 ponovno objašnjaje nemožnost izpuščenja komponenty. Arhivna varijanta najkrajše ne jest forma v toj sravnjenoj pasaži.\par
\noindent\emph{Textual comparison:} Paper 19 distinguishes two reducedness conditions; shortest uses the first. Paper 30 again explains the impossibility of omitting a component. Archival najkrajše is not the form in that compared passage.\par
\noindent{\footnotesize TEXT-COMP-201; pasaže / passages: 12570, 16617; \texttt{CORPUS\_FORM\_COMPARISONS.json}.}\par
\end{minipage}\par
\par\medskip\Needspace{7\baselineskip}
\noindent\begin{minipage}{\linewidth}
\hypertarget{TERM-1078}{}
\noindent\textbf{najkratše predstavjenja}\quad{\footnotesize\texttt{TERM-1078}}\par
\noindent\emph{Iztočnik / Source:} kürzeste Darstellungen\par
\noindent\emph{Rusky / Russian:} кратчайшие представления\par
\noindent\emph{Ukrajinsky / Ukrainian:} найкоротші подання\par
\smallskip\noindent\emph{Razlog:} Najkratše predstavjenja čuvaje razliku medžu prědstavjenjami, iz ktoryh ne može odpasti komponenta, i polno redukovanymi prědstavjenjami. Definicija I v §6 raboty 19 davaje dva uslovja za redukovanost; za najkratše dostatočno jest prvo. Arhivny razlog izrično zabranjaje sliti oba pojmy. To jest semantično ograničenje termina, ne prědpočitavanje jednoj sinonimičnoj frazy. Dva zapisa priměnjenja toj samoj zasady ne sut dva nezavisne dokazy jezykovej normy.\par
\noindent\emph{Reason:} Najkratše predstavjenja preserves the distinction between representations from which no component can be omitted and fully reduced representations. Definition I in Paper 19, §6 gives two conditions for reducedness; the first alone defines shortest representations. The recorded rationale explicitly prevents conflating the notions. This is a semantic restriction, not a preference between synonymous phrases. Two records applying the same principle are not two independent demonstrations of a linguistic norm.\par
\noindent{\footnotesize Izvorne zapisy / Original records: HIST-g081-031, HIST-g082-031.}\par
\smallskip\noindent\emph{Sravnjenje s tekstom:} Rabota 19 razlikuje dva uslovja redukovanosti; najkratše koristi prvo. Rabota 30 ponovno objašnjaje nemožnost izpuščenja komponenty. Arhivna varijanta najkrajše ne jest forma v toj sravnjenoj pasaži.\par
\noindent\emph{Textual comparison:} Paper 19 distinguishes two reducedness conditions; shortest uses the first. Paper 30 again explains the impossibility of omitting a component. Archival najkrajše is not the form in that compared passage.\par
\noindent{\footnotesize TEXT-COMP-201; pasaže / passages: 12570, 16617; \texttt{CORPUS\_FORM\_COMPARISONS.json}.}\par
\end{minipage}\par
\par\medskip\Needspace{7\baselineskip}
\noindent\begin{minipage}{\linewidth}
\hypertarget{TERM-1079}{}
\noindent\textbf{dolžina ideala / kompozicijna serija}\quad{\footnotesize\texttt{TERM-1079}}\par
\noindent\emph{Iztočnik / Source:} Laenge des Ideals / Kompositionsreihe\par
\noindent\emph{Rusky / Russian:} длина идеала / композиционный ряд\par
\noindent\emph{Ukrajinsky / Ukrainian:} довжина ідеалу / композиційний ряд\par
\smallskip\noindent\emph{Razlog:} Dolžina ideala i kompozicijna serija v dodatku 45 sut svezane s rezultanta-mnogokratnostju, ne s geometričnoju dolžinoju. Tekst identifikuje čislo s dolžinoju ideala i s rangom prstena klasov ostatkov; podnožna primětka objašnjaje dolžinu črez kompozicijnu seriju. Zato trěba zachovati konečne točky toj serije i prědpoloženje algebraično zatvorenogo osnovnogo polja v navedenom obsegu. Arhivna ocěna prozračnosti slova dolžina jest argument za izbor, ne rezultat čitateljskogo ispytanja.\par
\noindent\emph{Reason:} Dolžina ideala and kompozicijna serija in Supplement 45 concern resultant multiplicity, not geometric length. The text identifies the number with ideal length and with the rank of the residue-class ring; a footnote explains length through a composition series. The series’ endpoints and the stated assumption of an algebraically closed base field must therefore be retained within the passage’s scope. The recorded appeal to the transparency of dolžina is a reason for the choice, not a reader-test result.\par
\noindent{\footnotesize Izvorne zapisy / Original records: HIST-g214-011.}\par
\smallskip\noindent\emph{Sravnjenje s tekstom:} Dodatok 45 veže dolžinu s kompozicijnoju serijeju i rangom prstena klasov ostatkov. Primětka 14 izrično postavlja algebraično zatvorjeno osnovno polje v označenom obsegu; definicionne uslovja ne smějut odpasti.\par
\noindent\emph{Textual comparison:} Supplement 45 connects length with a composition series and the rank of the residue-class ring. Footnote 14 explicitly assumes an algebraically closed base field in the stated scope; the defining conditions must not disappear.\par
\noindent{\footnotesize TEXT-COMP-202; pasaže / passages: 45-isv.tex:399, 45-isv.tex:400, 45-isv.tex:410, 45-isv.tex:411, 45-isv.tex:412, 45-isv.tex:413, 45-isv.tex:414, 45-isv.tex:415, 45-isv.tex:416, 45-isv.tex:417, 45-isv.tex:425, 45-isv.tex:426; \texttt{CORPUS\_FORM\_COMPARISONS.json}.}\par
\end{minipage}\par
\par\medskip\Needspace{7\baselineskip}
\noindent\begin{minipage}{\linewidth}
\hypertarget{TERM-1080}{}
\noindent\textbf{Lagrangeve derivacije}\quad{\footnotesize\texttt{TERM-1080}}\par
\noindent\emph{Iztočnik / Source:} Lagrangesche Ableitungen\par
\noindent\emph{Rusky / Russian:} лагранжевы производные\par
\noindent\emph{Ukrajinsky / Ukrainian:} лагранжеві похідні\par
\smallskip\noindent\emph{Razlog:} Lagrangeve derivacije v rabotě 23 ne sut obči naziv za vse derivacije, svezane s imenom Lagrange. Tekst veže jih s variacionnoju zadačeju i zaměnoju vyšših diferencialov; daljše polarizovane kombinacije imajut vlastne transformacijske svojstva. Arhivny zapis govori takože o zachovanju latinskogo imena v kyriličnom tekstu radi izběganja mehaničnoj pogrěšky. To jest oddělna pisemna konvencija. Tu sravnjeny latinsky tekst podpiraje termin i kontekst, ale ne jest kontroljeju togo kyriličnogo pravila.\par
\noindent\emph{Reason:} Lagrangeve derivacije in Paper 23 does not mean any derivative associated with Lagrange’s name. The text connects them with the variational problem and the replacement of higher differentials; further polarized combinations have specified transformation properties. The archival record also states a policy of retaining the Latin name within Cyrillic text to avoid a mechanical misspelling. That is a separate script convention. The Latin passage compared here supports the term and context, but is not a check of that Cyrillic policy.\par
\noindent{\footnotesize Izvorne zapisy / Original records: HIST-g094-014.}\par
\smallskip\noindent\emph{Sravnjenje s tekstom:} Rabota 23 opisuje zaměnu vyšših diferencialov i polarizovane kombinacije. Imenska konstrukcija ne opravdyvaje izpuščenje togo konteksta niti sama potvrdžaje pravopis kyriličnogo prevoda.\par
\noindent\emph{Textual comparison:} Paper 23 describes replacing higher differentials and forming polarized combinations. The naming construction neither licenses omission of that context nor independently verifies the Cyrillic translation’s spelling.\par
\noindent{\footnotesize TEXT-COMP-203; pasaže / passages: 15197, 15205; \texttt{CORPUS\_FORM\_COMPARISONS.json}.}\par
\end{minipage}\par
\par\medskip\Needspace{7\baselineskip}
\noindent\begin{minipage}{\linewidth}
\hypertarget{TERM-1081}{}
\noindent\textbf{Lagrangeove izrazy}\quad{\footnotesize\texttt{TERM-1081}}\par
\noindent\emph{Iztočnik / Source:} Lagrangesche Ausdrücke\par
\noindent\emph{Rusky / Russian:} лагранжевы выражения\par
\noindent\emph{Ukrajinsky / Ukrainian:} лагранжеві вирази\par
\smallskip\noindent\emph{Razlog:} Naziv Lagrangeove izrazy zachovuje ime i razlikuje izraz od uravnenja. V rabotě 21 veličiny $\psi_i$ nastavajut pri integraciji po čestjah; uravnenja nastavajut, kogda te izrazy položimo ravnymi nuli. Zapis daje prědnost poznavajemosti imena, ne novomu nazivu objekta. Rusky i ukrajinsky sravniteljny izraz imajut podobnu strukturu. Forma Lagrangeove jest jedna iz několiko eponimnyh grafij v korpusu, ne dokaz jedinoj prějete medžuslovjanskoj normy; v rabotě 13 stoji Lagranževe. Različna grafija sama ne znači različny matematičny pojem.\par
\noindent\emph{Reason:} The name Lagrangeove izrazy preserves the eponym and distinguishes an expression from an equation. In Paper 21 the quantities $\psi_i$ arise through integration by parts; equations result when those expressions are set equal to zero. The recorded preference is for a recognizable name rather than a newly coined object name. The Russian and Ukrainian comparison forms have a similar structure. Lagrangeove is one of several eponymic spellings in the corpus, not evidence of a single accepted Interslavic norm: Paper 13 uses Lagranževe. The spelling difference alone does not identify a different mathematical concept.\par
\noindent{\footnotesize Izvorne zapisy / Original records: HIST-g084-005.}\par
\smallskip\noindent\emph{Sravnjenje s tekstom:} Raboty 13 i 21 vežut izrazy s variacionnym računom, ale upotrěbjajut različne grafije imena. Formuly razlikujut sam izraz od uravnenja nastavajučego položenjem togo izraza ravnym nuli.\par
\noindent\emph{Textual comparison:} Papers 13 and 21 connect the expressions with variational calculus but use different spellings of the eponym. The formulas distinguish the expression itself from the equation obtained by setting it equal to zero.\par
\noindent{\footnotesize TEXT-COMP-204; pasaže / passages: 9225, 9231, 14075, 14077, 14079; \texttt{CORPUS\_FORM\_COMPARISONS.json}.}\par
\end{minipage}\par
\par\medskip\Needspace{7\baselineskip}
\noindent\begin{minipage}{\linewidth}
\hypertarget{TERM-1082}{}
\noindent\textbf{Lagranževe izrazy}\quad{\footnotesize\texttt{TERM-1082}}\par
\noindent\emph{Iztočnik / Source:} Lagrangesche Ausdrücke\par
\noindent\emph{Rusky / Russian:} лагранжевы выражения\par
\noindent\emph{Ukrajinsky / Ukrainian:} лагранжеві вирази\par
\smallskip\noindent\emph{Razlog:} Izrazy jest tu sučstveny izbor: rabota 13 izrično nazyvaje $\psi_i$ lěvymi stranami Lagrangevyh uravnenij. Slovo uravnenja by utožnilo izraz s tvrđenjem, že on jest ravny nuli; operatori by vneslo drugy način opisa. Arhivny razlog zato zachovuje smysl německogo Ausdrücke. Lagranževe prilagodžuje ime slavjanskomu pisanju, v različju od Lagrangeove v rabotě 21. Tu ne jest ustanovjena prěvaha jednoj grafije za vse konteksty. Važna jest razlika medžu variacionnym izrazom, uravnenjem i operacijeju, ne samo podobnost slov.\par
\noindent\emph{Reason:} Izrazy is a substantive choice here: Paper 13 explicitly identifies $\psi_i$ as the left-hand sides of the Lagrange equations. Calling them equations would identify an expression with the assertion that it vanishes; operators would introduce a different description. The archived reason therefore preserves the sense of German Ausdrücke. Lagranževe adapts the name to Slavic spelling, unlike Lagrangeove in Paper 21. No general superiority of one spelling has been established here. What matters is the distinction between the variational expression, the equation and the operation, not simply resemblance between words.\par
\noindent{\footnotesize Izvorne zapisy / Original records: HIST-g075-008.}\par
\smallskip\noindent\emph{Sravnjenje s tekstom:} Raboty 13 i 21 vežut izrazy s variacionnym računom, ale upotrěbjajut različne grafije imena. Formuly razlikujut sam izraz od uravnenja nastavajučego položenjem togo izraza ravnym nuli.\par
\noindent\emph{Textual comparison:} Papers 13 and 21 connect the expressions with variational calculus but use different spellings of the eponym. The formulas distinguish the expression itself from the equation obtained by setting it equal to zero.\par
\noindent{\footnotesize TEXT-COMP-204; pasaže / passages: 9225, 9231, 14075, 14077, 14079; \texttt{CORPUS\_FORM\_COMPARISONS.json}.}\par
\end{minipage}\par
\par\medskip\Needspace{7\baselineskip}
\noindent\begin{minipage}{\linewidth}
\hypertarget{TERM-1083}{}
\noindent\textbf{Lagrangeve rodove oblasti}\quad{\footnotesize\texttt{TERM-1083}}\par
\noindent\emph{Iztočnik / Source:} Lagrangesche Gattungsbereiche\par
\noindent\emph{Rusky / Russian:} лагранжевы области родов\par
\noindent\emph{Ukrajinsky / Ukrainian:} Лагранжеві області родів\par
\smallskip\noindent\emph{Razlog:} Rodove oblasti zachovuje historijny naziv Gattungsbereiche, ale rod tu ne znači rod algebraičnoj krivoj. V rabotah 5 i 6 ide o polja racionalnyh funkcij, ktore ostavajut neizměnjene pri permutacijah zadanoj grupy. Rabota 6 potom razgleduje jih cěle racionalne funkcije i integralnost nad elementarnymi simetričnymi funkcijami. Tri zachovane razloga podčrkivajut historijnu identičnost i kontinuitet termina. To opravdyvaje zachovanje naziva s objasnjenjem objekta, ne prenos v sovrěmennu terminologiju roda bez opory v definiciji. Ruske i ukrajinske formy sut zabilěžene sravnjenja, ne samostatny dokaz ustaljenoj upotrěby.\par
\noindent\emph{Reason:} Rodove oblasti preserves the historical designation Gattungsbereiche, but rod does not mean the genus of an algebraic curve here. Papers 5 and 6 concern fields of rational functions unchanged by the permutations of a specified group. Paper 6 subsequently studies their polynomial elements and integrality over the elementary symmetric functions. All three preserved explanations emphasize historical identity and continuity of terminology. That supports retaining the name with an explanation of its object, not transferring it into modern genus terminology without support from the definition. The Russian and Ukrainian forms are recorded comparanda, not independent evidence of established usage.\par
\noindent{\footnotesize Izvorne zapisy / Original records: HIST-g039-005, HIST-g041-013, HIST-g053-002.}\par
\smallskip\noindent\emph{Sravnjenje s tekstom:} Definicije v rabotah 5, 6 i 11 svęzujut naziv s racionalnymi funkcijami invariantnymi pri permutacijah. Razgovor o cělyh racionalnyh funkcijah v §13 i §15 ne znači, že polje i polinomny podprsten sut to samo.\par
\noindent\emph{Textual comparison:} The definitions in Papers 5, 6 and 11 connect the name with rational functions invariant under permutations. The discussion of polynomial elements in §§13 and 15 does not identify the field with its polynomial subring.\par
\noindent{\footnotesize TEXT-COMP-205; pasaže / passages: 4950, 4952, 4956, 5124, 5126, 5128, 6367, 6369, 6373, 6587, 8728, 8763; \texttt{CORPUS\_FORM\_COMPARISONS.json}.}\par
\end{minipage}\par
\par\medskip\Needspace{7\baselineskip}
\noindent\begin{minipage}{\linewidth}
\hypertarget{TERM-1084}{}
\noindent\textbf{Lagrangeva interpolacijska formula}\quad{\footnotesize\texttt{TERM-1084}}\par
\noindent\emph{Iztočnik / Source:} Lagrangesche Interpolationsformel\par
\noindent\emph{Rusky / Russian:} интерполяционная формула Лагранжа\par
\noindent\emph{Ukrajinsky / Ukrainian:} інтерполяційна формула Лагранжа\par
\smallskip\noindent\emph{Razlog:} Lagrangeva interpolacijska formula sostoji se iz poznavajemogo imena i internacionalnogo korenja interpolacija. Sravniteljne ruske i ukrajinske konstrukcije vyražajut ime genitivom; medžuslovjanska forma upotrěbjaje pridavnik. Matematična uloga formuly zavisi od města: v rabotě 7 podnožna primětka govori o propustu v Weberovom dokazu i zaměnjuje interpolaciju diferenciacijeju; v rabotě 36 ona služi kako sravnjenje pri obobčenju. Naziv zato ne govori, že oba dokaza upotrěbjajut tu samu formulu. Ranějša nedouměnica o kyriličnom pisanju imena ne može byti rěšena samym latinskym sravnjenjem.\par
\noindent\emph{Reason:} Lagrangeva interpolacijska formula combines a recognizable eponym with the international interpolation stem. The Russian and Ukrainian comparanda express the name through a genitive construction; the Interslavic form uses an adjective. The formula’s mathematical role varies: a footnote in Paper 7 identifies a gap in Weber’s proof and replaces interpolation by differentiation, whereas Paper 36 invokes it as a comparison in a generalization. The name does not imply that both arguments apply the same formula. The earlier uncertainty about the Cyrillic spelling of the eponym cannot be resolved from the Latin comparison alone.\par
\noindent{\footnotesize Izvorne zapisy / Original records: HIST-g056-013, HIST-g206-006.}\par
\smallskip\noindent\emph{Sravnjenje s tekstom:} Podnožna primětka v rabotě 7 zaměnjuje interpolaciju diferenciacijeju radi propusta v citovanom dokazu. Rabota 36 upotrěbjaje interpolaciju kako sravnjenje pri teoriji diferenta. Jedin naziv zato ne označuje jeden i toj samy korak oběh argumentov.\par
\noindent\emph{Textual comparison:} The footnote in Paper 7 replaces interpolation by differentiation because of a gap in the cited proof. Paper 36 uses interpolation as a comparison in the theory of the different. One name therefore does not designate the same step in both arguments.\par
\noindent{\footnotesize TEXT-COMP-206; pasaže / passages: 6706, 22743, 22745; \texttt{CORPUS\_FORM\_COMPARISONS.json}.}\par
\end{minipage}\par
\par\medskip\Needspace{7\baselineskip}
\noindent\begin{minipage}{\linewidth}
\hypertarget{TERM-1085}{}
\noindent\textbf{Lagrangeva rodova oblast}\quad{\footnotesize\texttt{TERM-1085}}\par
\noindent\emph{Iztočnik / Source:} Lagrangescher Gattungsbereich\par
\noindent\emph{Rusky / Russian:} лагранжева родовая область\par
\noindent\emph{Ukrajinsky / Ukrainian:} Лагранжева родова область\par
\smallskip\noindent\emph{Razlog:} Jednina Lagrangeva rodova oblast v §15 raboty 6 prodolžuje naziv iz §13. Zadača tu jest najti integralnu bazu za racionalne funkcije s cělymi čislovymi koeficientami v toj oblasti. Kontinuitet naziva čuva vez s ranějšoju konstrukcijeju, ale ne izravnjuje polje vsih invariantnyh racionalnyh funkcij s jego polinomnym podprstenom. Arhivny razlog jest kratky i ne prinosi novogo jezykovogo dokaza za rodova. Smysl trěba čitati iz definicije i konteksta; historijny termin ne směje sugerovati sovrěmenny rod krivoj.\par
\noindent\emph{Reason:} The singular Lagrangeva rodova oblast in §15 of Paper 6 continues the name used in §13. The problem here is to find an integral basis for the rational functions with integral numerical coefficients in that domain. Continuity of naming maintains the link with the earlier construction, but does not identify the field of all invariant rational functions with its polynomial subring. The archived explanation is brief and supplies no new linguistic evidence for rodova. Its meaning must be read from the definition and context; the historical term must not suggest the modern genus of a curve.\par
\noindent{\footnotesize Izvorne zapisy / Original records: HIST-g055-009.}\par
\smallskip\noindent\emph{Sravnjenje s tekstom:} Definicije v rabotah 5, 6 i 11 svęzujut naziv s racionalnymi funkcijami invariantnymi pri permutacijah. Razgovor o cělyh racionalnyh funkcijah v §13 i §15 ne znači, že polje i polinomny podprsten sut to samo.\par
\noindent\emph{Textual comparison:} The definitions in Papers 5, 6 and 11 connect the name with rational functions invariant under permutations. The discussion of polynomial elements in §§13 and 15 does not identify the field with its polynomial subring.\par
\noindent{\footnotesize TEXT-COMP-205; pasaže / passages: 4950, 4952, 4956, 5124, 5126, 5128, 6367, 6369, 6373, 6587, 8728, 8763; \texttt{CORPUS\_FORM\_COMPARISONS.json}.}\par
\end{minipage}\par
\par\medskip\Needspace{7\baselineskip}
\noindent\begin{minipage}{\linewidth}
\hypertarget{TERM-1086}{}
\noindent\textbf{Lagranžev rodovy domen}\quad{\footnotesize\texttt{TERM-1086}}\par
\noindent\emph{Iztočnik / Source:} Lagrangescher Gattungsbereich\par
\noindent\emph{Rusky / Russian:} лагранжев родовой домен\par
\noindent\emph{Ukrajinsky / Ukrainian:} лагранжів родовий домен\par
\smallskip\noindent\emph{Razlog:} V rabotě 11 Lagranžev rodovy domen jest objasnjen kako polje funkcij neizměnjajemyh pri zadanoj permutacijskoj grupě. Domen i ranějše oblast sut dvě prěvodne formy historijnogo naziva, ne dva dokazano različna objekta. Arhivny zapis daje takože pisemny razlog: Lagranžev jest prilagodženy oblik imena, ktory ne potrěbuje neprilagodženogo francuzskogo korenja v kyriličnom pridavniku. To jest argument o pisanju, ne o matematičnoj definiciji. V toj rabotě vopros o racionalnosti polja ostavaje odděljeny od samogo naziva; slovo domen ne směje sugerovati, že racionalnost už dokazana.\par
\noindent\emph{Reason:} In Paper 11 Lagranžev rodovy domen is explained as the field of functions fixed by the specified permutation group. Domen and the earlier oblast are two renderings of a historical name, not two demonstrably distinct objects. The archival record also gives a script-related reason: Lagranžev adapts the name without requiring an unadapted French stem inside a Cyrillic adjective. That is an argument about spelling, not mathematical definition. The rationality question remains separate from naming the field; domen must not suggest that rationality has already been proved.\par
\noindent{\footnotesize Izvorne zapisy / Original records: HIST-g073-008.}\par
\smallskip\noindent\emph{Sravnjenje s tekstom:} Definicije v rabotah 5, 6 i 11 svęzujut naziv s racionalnymi funkcijami invariantnymi pri permutacijah. Razgovor o cělyh racionalnyh funkcijah v §13 i §15 ne znači, že polje i polinomny podprsten sut to samo.\par
\noindent\emph{Textual comparison:} The definitions in Papers 5, 6 and 11 connect the name with rational functions invariant under permutations. The discussion of polynomial elements in §§13 and 15 does not identify the field with its polynomial subring.\par
\noindent{\footnotesize TEXT-COMP-205; pasaže / passages: 4950, 4952, 4956, 5124, 5126, 5128, 6367, 6369, 6373, 6587, 8728, 8763; \texttt{CORPUS\_FORM\_COMPARISONS.json}.}\par
\end{minipage}\par
\par\medskip\Needspace{7\baselineskip}
\noindent\begin{minipage}{\linewidth}
\hypertarget{TERM-1087}{}
\noindent\textbf{Laplaceovo razloženje}\quad{\footnotesize\texttt{TERM-1087}}\par
\noindent\emph{Iztočnik / Source:} Laplacesche Entwicklung\par
\noindent\emph{Rusky / Russian:} разложение Лапласа\par
\noindent\emph{Ukrajinsky / Ukrainian:} розклад Лапласа\par
\smallskip\noindent\emph{Razlog:} Laplaceovo razloženje označuje razvinutje determinanta po minorah. V rabotě 4 ta operacija daje sumu produktov odgovarjajučih determinantov pri definiciji matričnogo produkta. Razloženje jest srodno ruskomu razloženiju i ukrajinskomu rozkladu v zabilěženom sravnjenju; ime Laplace ostavaje vidimo. Tu ne ide o libovoljno razloženje polinoma na faktory. Izbor zato trěba vezati s operacijeju v dokazu. Arhivny razlog govori o obyčajnom nazivu, ale ne daje konkretnogo jezykovogo iztočnika, ktory by sam potvrdil normativnost medžuslovjanskoj formy.\par
\noindent\emph{Reason:} Laplaceovo razloženje denotes expansion of a determinant by minors. In Paper 4 that operation gives a sum of products of corresponding determinants in the definition of a matrix product. Razloženje resembles the Russian razloženie and Ukrainian rozklad in the recorded comparison, while keeping Laplace’s name visible. This is not arbitrary factorization of a polynomial. The choice must therefore be tied to the operation in the proof. The archived explanation calls the name standard, but provides no specific linguistic source that would independently establish the Interslavic form as normative.\par
\noindent{\footnotesize Izvorne zapisy / Original records: HIST-g030-012.}\par
\smallskip\noindent\emph{Sravnjenje s tekstom:} V rabotě 4 determinanty i jih minory daju konkretno značenje slovu razloženje. Ta vez čuva matematičnu operaciju, koju samo slovo bez konteksta ne oprěděljaje.\par
\noindent\emph{Textual comparison:} In Paper 4, determinants and their minors give razloženje its specific meaning. That connection preserves the mathematical operation, which the word alone does not determine.\par
\noindent{\footnotesize TEXT-COMP-207; pasaže / passages: 3782, 3805, 3807, 3809; \texttt{CORPUS\_FORM\_COMPARISONS.json}.}\par
\end{minipage}\par
\par\medskip\Needspace{7\baselineskip}
\noindent\begin{minipage}{\linewidth}
\hypertarget{TERM-1088}{}
\noindent\textbf{leksikografičny red / najvyšše členy}\quad{\footnotesize\texttt{TERM-1088}}\par
\noindent\emph{Iztočnik / Source:} lexikographische Anordnung / höchste Glieder\par
\noindent\emph{Rusky / Russian:} лексикографический порядок / старшие члены\par
\noindent\emph{Ukrajinsky / Ukrainian:} лексикографічний порядок / найстарші члени\par
\smallskip\noindent\emph{Razlog:} Najvyšše členy bylo izbrano vměsto kalka najstarše členy: porjadok členov jest formalny, ne vrěmenny. Rusky i ukrajinsky sravniteljny izraz upotrěbjajut metaforu staršinstva; to samo po sobě ne čini jih matematično netočnymi. Medžuslovjansky izbor daje prědnost vysotnoj metaforě. V rabotě 24 dokaz vybira najvyšši član s koeficientom ne dělimym zadanym prostym idealom, po leksikografičnom redu. To ne znači prostě najvyššu obču stepenj. Arhivny zapis ostavlja vopros prirodnosti frazy otkrytym; tu ne imajemo rezultata čitateljskogo izpytanja.\par
\noindent\emph{Reason:} Najvyšše členy was preferred to the calque najstarše členy because the ordering is formal, not temporal. The Russian and Ukrainian comparanda use a seniority metaphor; that does not itself make them mathematically inaccurate. The Interslavic choice instead prefers a height metaphor. In Paper 24 the proof selects the highest term whose coefficient is not divisible by a specified prime ideal, using lexicographic order. This does not simply mean highest total degree. The original record leaves the naturalness of the phrase open; no reader-testing result is available here.\par
\noindent{\footnotesize Izvorne zapisy / Original records: HIST-g096-004.}\par
\smallskip\noindent\emph{Sravnjenje s tekstom:} Rabota 17 upotrěbjaje leksikografičnu redukciju, a rabota 24 izbira najvyšši člen s koeficientom ne dělimym zadanym prostym idealom. V oběh slučajah red mora byti znany; najvyšši ne znači avtomatično najvyššu obču stepenj.\par
\noindent\emph{Textual comparison:} Paper 17 uses lexicographic reduction; Paper 24 selects the highest term whose coefficient is not divisible by the specified prime ideal. In both cases the order must be known: highest does not automatically mean highest total degree.\par
\noindent{\footnotesize TEXT-COMP-208; pasaže / passages: 11417, 11419, 11420, 11421, 11422, 11423, 11424, 11426, 15374, 15378; \texttt{CORPUS\_FORM\_COMPARISONS.json}.}\par
\end{minipage}\par
\par\medskip\Needspace{7\baselineskip}
\noindent\begin{minipage}{\linewidth}
\hypertarget{TERM-1089}{}
\noindent\textbf{leksikografičny red}\quad{\footnotesize\texttt{TERM-1089}}\par
\noindent\emph{Iztočnik / Source:} lexikographische Ordnung\par
\noindent\emph{Rusky / Russian:} лексикографический порядок\par
\noindent\emph{Ukrajinsky / Ukrainian:} лексикографічний порядок\par
\smallskip\noindent\emph{Razlog:} Leksikografičny red označuje dogovor, po ktorom se sravnjavajut eksponenty monomov v ustanovjenom porjadku prěměnnyh. V redukciji raboty 17 vybrany najvyšši član može byti odstranjen, tako že ostanu nižše členy po tomu redu. Arhivny razlog daje prědnost kontinuitetu slova red i internacionalnomu pridavniku, podobnomu ruskomu i ukrajinskomu sravnjenju. Red tu ne znači ni čislovy rjad ni stepenj polinoma. Matematično objasnjenje izbora ne zaměnjuje dokaz, že konkretna redukcija jest pravilna; ono govori, koju informaciju termin mora zachovati.\par
\noindent\emph{Reason:} Leksikografičny red denotes the convention by which monomial exponents are compared in a specified variable order. In the reduction in Paper 17, the selected leading term can be removed, leaving terms lower in that order. The recorded preference is for continuity in red and an international adjective resembling the Russian and Ukrainian comparanda. Red here means neither a numerical series nor polynomial degree. Explaining the terminology does not replace a proof that the particular reduction is valid; it identifies the information that the term must retain.\par
\noindent{\footnotesize Izvorne zapisy / Original records: HIST-g079-035.}\par
\smallskip\noindent\emph{Sravnjenje s tekstom:} Rabota 17 upotrěbjaje leksikografičnu redukciju, a rabota 24 izbira najvyšši člen s koeficientom ne dělimym zadanym prostym idealom. V oběh slučajah red mora byti znany; najvyšši ne znači avtomatično najvyššu obču stepenj.\par
\noindent\emph{Textual comparison:} Paper 17 uses lexicographic reduction; Paper 24 selects the highest term whose coefficient is not divisible by the specified prime ideal. In both cases the order must be known: highest does not automatically mean highest total degree.\par
\noindent{\footnotesize TEXT-COMP-208; pasaže / passages: 11417, 11419, 11420, 11421, 11422, 11423, 11424, 11426, 15374, 15378; \texttt{CORPUS\_FORM\_COMPARISONS.json}.}\par
\end{minipage}\par
\par\medskip\Needspace{7\baselineskip}
\noindent\begin{minipage}{\linewidth}
\hypertarget{TERM-1090}{}
\noindent\textbf{linearno-drobna zavisnost}\quad{\footnotesize\texttt{TERM-1090}}\par
\noindent\emph{Iztočnik / Source:} linear gebrochene Abhängigkeit\par
\noindent\emph{Rusky / Russian:} дробно-линейная зависимость\par
\noindent\emph{Ukrajinsky / Ukrainian:} дробово-лінійна залежність\par
\smallskip\noindent\emph{Razlog:} Linearno-drobna zavisnost govori o zavisnosti, vyražajemej drobju s linearnym čisliteljem i imenovateljem. Arhivny razlog veže naziv s ranějšim terminom za linearno-drobno prěobraženje. V §6 raboty 6 takova zavisnost slěduje za funkcije, ktore porodžajut to samo racionalno polje, v okviru razgledanogo voprosa o Lürothovoj funkciji. Ona ne slěduje iz každoj racionalnoj zavisnosti dvuh funkcij. Slovo zavisnost zachovuje ulogu tvrđenja, dokolě prěobraženje imenuje način prěhoda medžu izrazami. Sravniteljne formy podpirajut složeny naziv, ne čislovu ocěnku jego razumlivosti.\par
\noindent\emph{Reason:} Linearno-drobna zavisnost describes dependence expressed as a fraction with linear numerator and denominator. The archived reason links it to the earlier term for a fractional-linear transformation. In §6 of Paper 6 this dependence arises for functions generating the same rational function field, within the discussion of a Lüroth function. It does not follow from every rational dependence between two functions. Zavisnost preserves the role of the assertion, whereas transformation names the passage between expressions. The comparison forms support the compound terminology, not a numerical estimate of comprehensibility.\par
\noindent{\footnotesize Izvorne zapisy / Original records: HIST-g046-006.}\par
\smallskip\noindent\emph{Sravnjenje s tekstom:} §6 raboty 6 veže linearno-drobnu zavisnost s ravnostju odgovarjajučih racionalnyh polj. Tuto ograničenje brani čitanje tvrđenja kako pravila za libovoljnu racionalnu zavisnost.\par
\noindent\emph{Textual comparison:} §6 of Paper 6 connects fractional-linear dependence with equality of the associated rational function fields. This qualification prevents reading the assertion as a rule for arbitrary rational dependence.\par
\noindent{\footnotesize TEXT-COMP-209; pasaže / passages: 5692, 5696, 5704, 5708; \texttt{CORPUS\_FORM\_COMPARISONS.json}.}\par
\end{minipage}\par
\par\medskip\Needspace{7\baselineskip}
\noindent\begin{minipage}{\linewidth}
\hypertarget{TERM-1091}{}
\noindent\textbf{linearno nezavisny vzhodno k $P$}\quad{\footnotesize\texttt{TERM-1091}}\par
\noindent\emph{Iztočnik / Source:} linear unabhängig in bezug auf $P$\par
\noindent\emph{Rusky / Russian:} линейно независимый относительно $P$\par
\noindent\emph{Ukrajinsky / Ukrainian:} лінійно незалежний відносно $P$\par
\smallskip\noindent\emph{Razlog:} Dopolnjenje po odnošenju k $P$ ne može se izpustiti: linearna nezavisnost zavisi od polja dopustimyh koeficientov. Ranějši zapis měl vzhodno k $P$, a sravnjeny tekst §1 raboty 31 imaje po odnošenju k $P$. Nezavisnost nad $\overline P$ tam jest čestju specialnoj konstrukcije razširjenja; podnožna primětka objasnjaje, začto ona ne jest avtomatična pri libovoljnom vključenju v veče polje. Racionalna cělj ranějšej frazy jest jasna, ale tvrđenje, že vzhodno k bylo ustaljenym jezykovym obratom, ne jest tym dokazano.\par
\noindent\emph{Reason:} The complement po odnošenju k $P$ is mathematically indispensable: linear independence depends on the field of allowed coefficients. The earlier record has vzhodno k $P$, while the compared text of §1 in Paper 31 has po odnošenju k $P$. Independence over $\overline P$ belongs there to a particular construction of an extension; a footnote explains why it is not automatic under an arbitrary embedding into a larger field. The intended function of the earlier phrase is clear, but that does not establish the claim that vzhodno k was an accepted linguistic idiom.\par
\noindent{\footnotesize Izvorne zapisy / Original records: HIST-g119-010.}\par
\smallskip\noindent\emph{Sravnjenje s tekstom:} V §1 raboty 31 tekst upotrěbjaje po odnošenju k, ne ranějše vzhodno k. Konstrukcija razširjenja čuva nezavisnost; podnožna primětka izrično upozornjaje na inačny rezultat bez jej prědpoloženij.\par
\noindent\emph{Textual comparison:} The text of §1 in Paper 31 uses po odnošenju k, not the earlier vzhodno k. The extension is constructed to preserve independence; the footnote explicitly warns of a different result without its assumptions.\par
\noindent{\footnotesize TEXT-COMP-210; pasaže / passages: 17324, 17344, 17350, 17363; \texttt{CORPUS\_FORM\_COMPARISONS.json}.}\par
\end{minipage}\par
\par\medskip\Needspace{7\baselineskip}
\noindent\begin{minipage}{\linewidth}
\hypertarget{TERM-1092}{}
\noindent\textbf{linearno nezavisne medžu soboju i od ...}\quad{\footnotesize\texttt{TERM-1092}}\par
\noindent\emph{Iztočnik / Source:} linear unabhängig voneinander und von ...\par
\noindent\emph{Rusky / Russian:} линейно независимы между собой и от ...\par
\noindent\emph{Ukrajinsky / Ukrainian:} лінійно незалежні між собою і від ...\par
\smallskip\noindent\emph{Razlog:} Konstrukcija medžu soboju i od dopolnjaje, katere veličiny trěba razgledati zajedno pri izpytanju linearnoj nezavisnosti. Samo parna nezavisnost ne by dostala. Arhivny razlog za §4 raboty 22 različa semejstva $\xi$ i $\omega$. Sravnjeny tekst, nasuprot tomu, upotrěbjaje $\xi$ za obe semejstva, vključajuči tvrđenje o nezavisnosti $\xi$ od $\xi$. To jest sučstvena nejasnost oznak, ne razlog izbrisati vtoru polovinu frazy. Bez sravnjenja s izvornym svědkom ne pripisujemo točnu zaměnu pojedinym simbolam; razlikujemo naměru konstrukcije od nepodtvrđenogo zapisa dokaza.\par
\noindent\emph{Reason:} The construction medžu soboju i od specifies which quantities must be considered together when testing linear independence. Pairwise independence alone would not suffice. The archived explanation for §4 of Paper 22 distinguishes $\xi$ and $\omega$ families. The compared text, however, uses $\xi$ for both, including a statement of independence of $\xi$ from $\xi$. This is a substantive ambiguity of notation, not a reason to omit the second part of the phrase. Without comparison against the original witness, no exact replacement is assigned to individual symbols; the intended construction is distinguished from the unconfirmed notation of the proof.\par
\noindent{\footnotesize Izvorne zapisy / Original records: HIST-g088-010, HIST-g092-041.}\par
\smallskip\noindent\emph{Sravnjenje s tekstom:} Sravnjeny tekst §4 raboty 22 označuje konečno i konstruovano semejstvo črez $\xi$, dokolě zachovany razlog govori o $\xi$ i $\omega$. To sravnjenje ustanovjaje nejasnost, ale samo ne daje izvornogo raspoloženja vsih simbolov. Dokaz ne jest tu proglašen izpravnym.\par
\noindent\emph{Textual comparison:} The compared text of §4 in Paper 22 denotes both the finite and the constructed family by $\xi$, while the preserved rationale distinguishes $\xi$ and $\omega$. The comparison establishes an ambiguity but does not itself supply the original placement of every symbol. The proof is not certified as correct here.\par
\noindent{\footnotesize TEXT-COMP-211; pasaže / passages: 14652, 14654, 14656, 14715, 14716, 14717, 14721; \texttt{CORPUS\_FORM\_COMPARISONS.json}.}\par
\end{minipage}\par
\par\medskip\Needspace{7\baselineskip}
\noindent\begin{minipage}{\linewidth}
\hypertarget{TERM-1093}{}
\noindent\textbf{linearno nezavisna modulna baza / klasy ostatkov / kongruencija}\quad{\footnotesize\texttt{TERM-1093}}\par
\noindent\emph{Iztočnik / Source:} linear unabhängige Modulbasis / Restklassen / Kongruenz\par
\noindent\emph{Rusky / Russian:} линейно независимый модульный базис / классы вычетов / конгруенция\par
\noindent\emph{Ukrajinsky / Ukrainian:} лінійно незалежна модульна база / класи лишків / конгруенція\par
\smallskip\noindent\emph{Razlog:} Tute tri frazy rabotajut zajedno v dokazu: modulna baza daje jednoznačne koeficienty, klasy ostatkov sut obrazy elementov po redukciji, a kongruencija vyražaje ravnost po modulu ideala. V §7 raboty 31 iz relacije medžu klasami bazovyh elementov dobyva se ravnost s kratnikom $p$; nezavisnost izvornoj bazy togda prinudžaje vse koeficienty byti dělimymi prěz $p$. Internacionalne korenja linear-, modul- i kongruenc- byly izbrany radi prědpoloženej medžujezykovoj poznavajemosti. To jest zabilěžen razlog izbora, ne izměrena razumlivost. Bazu ne možno zaměniti libovoljnym porodžajučim sistemom bez utraty argumenta.\par
\noindent\emph{Reason:} These three phrases work together in the proof: a module basis gives unique coefficients, residue classes are the images of elements under reduction, and a congruence expresses equality modulo the ideal. In §7 of Paper 31, a relation between the classes of basis elements becomes an equality with a multiple of $p$; independence of the original basis then forces every coefficient to be divisible by $p$. The international stems linear-, modul- and kongruenc- were chosen for presumed cross-language recognizability. That is a recorded reason, not measured comprehension. Replacing the basis by an arbitrary generating system would lose the argument.\par
\noindent{\footnotesize Izvorne zapisy / Original records: HIST-g151-007.}\par
\smallskip\noindent\emph{Sravnjenje s tekstom:} §7 raboty 31 upotrěbjaje nezavisnost bazovyh elementov, aby iz kongruencije dobyti dělitelnost koeficientov. To jest razlog za zachovanje vsih treh različenij: baza, klasa i kongruencija.\par
\noindent\emph{Textual comparison:} §7 of Paper 31 uses independence of the basis elements to deduce divisibility of coefficients from a congruence. That is why all three distinctions matter: basis, class and congruence.\par
\noindent{\footnotesize TEXT-COMP-212; pasaže / passages: 18629, 18631, 18633, 18635, 18637; \texttt{CORPUS\_FORM\_COMPARISONS.json}.}\par
\end{minipage}\par
\par\medskip\Needspace{7\baselineskip}
\noindent\begin{minipage}{\linewidth}
\hypertarget{TERM-1094}{}
\noindent\textbf{linearno nezavisne relacije}\quad{\footnotesize\texttt{TERM-1094}}\par
\noindent\emph{Iztočnik / Source:} linear unabhängige Relationen\par
\noindent\emph{Rusky / Russian:} линейно независимые соотношения\par
\noindent\emph{Ukrajinsky / Ukrainian:} лінійно незалежні співвідношення\par
\smallskip\noindent\emph{Razlog:} Linearno nezavisne relacije imajut linearno nezavisne vektory koeficientov; nijedna iz relacij ne jest linearna kombinacija ostalyh. V §4 raboty 10 relacije s realnymi koeficientami ograničajut veličiny $x,y,X,Y$; rang četiri, tri, dva ili jedin odgovarja nuli, jednoj, dvěma ili trem nezavisnym relacijam. Pridavnik ne govori, že same veličiny sut nezavisne. Relacije prodolžuje ranějšu terminologiju algebraičnoj zavisnosti; rusky i ukrajinsky sravniteljny izraz vyražajut soodnos. Presna priroda koeficientov i čislo nezavisnyh relacij sut važnejše od samogo izbora internacionalnogo korenja.\par
\noindent\emph{Reason:} Linearno nezavisne relacije have linearly independent coefficient vectors: no relation is a linear combination of the others. In §4 of Paper 10, relations with real coefficients constrain $x,y,X,Y$; ranks four, three, two and one correspond to zero, one, two and three independent relations. The adjective does not assert that the quantities themselves are independent. Relacije continues the earlier terminology of algebraic dependence; the Russian and Ukrainian comparanda express a relation or interrelation. The coefficient domain and the number of independent relations matter more than the choice of an international stem alone.\par
\noindent{\footnotesize Izvorne zapisy / Original records: HIST-g072-005.}\par
\smallskip\noindent\emph{Sravnjenje s tekstom:} §4 raboty 10 svęzuje rang s čislom nezavisnyh relacij medžu četyroma veličinami. Koeficienty sut realne; izpuščenje togo uslovja ili zaměna nezavisnosti relacij nezavisnostju veličin by změnilo tvrđenje.\par
\noindent\emph{Textual comparison:} §4 of Paper 10 connects rank with the number of independent relations between four quantities. The coefficients are real; omitting that qualification or replacing independence of relations by independence of the quantities would change the assertion.\par
\noindent{\footnotesize TEXT-COMP-213; pasaže / passages: 8611, 8613, 8615, 8617; \texttt{CORPUS\_FORM\_COMPARISONS.json}.}\par
\end{minipage}\par
\par\medskip\Needspace{7\baselineskip}
\noindent\begin{minipage}{\linewidth}
\hypertarget{TERM-1095}{}
\noindent\textbf{linearno nezavisna $Z$-modulna baza}\quad{\footnotesize\texttt{TERM-1095}}\par
\noindent\emph{Iztočnik / Source:} linear unabhängige $Z$-Modulbasis\par
\noindent\emph{Rusky / Russian:} линейно независимая $Z$-модульная база\par
\noindent\emph{Ukrajinsky / Ukrainian:} лінійно незалежна $Z$-модульна база\par
\smallskip\noindent\emph{Razlog:} Oznaka $\mathfrak Z$ imenuje osnovnu oblast, ne obavezno prsten cělyh čisel. §4 raboty 31 izrično dopušćaje polje, prsten racionalnyh cělyh čisel ili odgovarjajuču funkcionalnu oblast. Baza jest linearno nezavisna nad toj osnovoj; egzistencija konečnoj takovoj bazy jest prědpoloženje, ne svojstvo vsakogo modula. Arhivny razlog čuva osnovu v složenom nazivu Z-modulna baza. V objasnjenju vračajemo frakturnu oznaku iz matematičnogo teksta, ale ne pretvarjajemo arhivno prosto Z v oznaku $\mathbb Z$. Time ostavaje vidima i osnovna oblast i sila prědpoloženja.\par
\noindent\emph{Reason:} The symbol $\mathfrak Z$ names the base domain, not necessarily the ring of integers. §4 of Paper 31 explicitly allows a field, the rational integers or the relevant function domain. The basis is linearly independent over that base; the existence of a finite such basis is an assumption, not a property of every module. The archived reason preserves the base in the compound Z-modulna baza. The explanation restores the fraktur symbol from the mathematical text, but does not turn the archive’s plain Z into $\mathbb Z$. Both the base domain and the strength of the assumption remain visible.\par
\noindent{\footnotesize Izvorne zapisy / Original records: HIST-g131-002.}\par
\smallskip\noindent\emph{Sravnjenje s tekstom:} Prědpoloženje §4 raboty 31 zahtěvaje konečne nezavisne modulne bazy. Osnovna oblast jest označena frakturnym Z i v tekstu ima několiko dopustimyh tipov; ona ne jest izključitelno prstenom cělyh čisel.\par
\noindent\emph{Textual comparison:} The assumption in §4 of Paper 31 requires finite independent module bases. The base domain is denoted by fraktur Z and has several allowed types in the text; it is not exclusively the ring of integers.\par
\noindent{\footnotesize TEXT-COMP-214; pasaže / passages: 17827, 17829, 17858, 17860; \texttt{CORPUS\_FORM\_COMPARISONS.json}.}\par
\end{minipage}\par
\par\medskip\Needspace{7\baselineskip}
\noindent\begin{minipage}{\linewidth}
\hypertarget{TERM-1096}{}
\noindent\textbf{linearna / racionalna / algebraična baza}\quad{\footnotesize\texttt{TERM-1096}}\par
\noindent\emph{Iztočnik / Source:} lineare / rationale / algebraische Basis\par
\noindent\emph{Rusky / Russian:} линейный / рациональный / алгебраический базис\par
\noindent\emph{Ukrajinsky / Ukrainian:} лінійна / раціональна / алгебраїчна база\par
\smallskip\noindent\emph{Razlog:} Tri pridavnika ne sut zaměnjajeme nazvy jednoj bazy. Uvod raboty 10 sravnjava linearno prědstavjenje s racionalnymi i algebraičnymi odnosami, osoblivo pri funkcionalnyh uravnenjah. Baza jest obči imenny člen, ale dopustime operacije i zahtěvana nezavisnost zaviset od pridavnika i definicije. Zapis daje prědnost medžuslovjanskomu baza radi kontinuiteta s ranějšimi razdělami, uz rusko bazis i ukrajinsko baza. Njegovo tvrđenje o sovrěmennoj ukrajinskoj upotrěbě nema tu konkretnogo bibliografičnogo dokaza; sravniteljne formy ne ustanovjajut, že druge forme sut nepravilne.\par
\noindent\emph{Reason:} The three adjectives are not interchangeable names for one kind of basis. The introduction to Paper 10 contrasts linear representation with rational and algebraic relations, particularly in functional equations. Baza is the shared noun, but the permitted operations and required independence depend on the adjective and definition. The record prefers Interslavic baza for continuity with earlier sections, alongside Russian bazis and Ukrainian baza. Its claim about contemporary Ukrainian usage has no specific bibliographical support here; the comparison forms do not establish that alternatives are incorrect.\par
\noindent{\footnotesize Izvorne zapisy / Original records: HIST-g069-008.}\par
\smallskip\noindent\emph{Sravnjenje s tekstom:} Uvod raboty 10 razlikuje linearnu, racionalnu i algebraičnu bazu. Podnožna primětka v §4 daje uslovja za Hamelovu bazu: racionalne koeficienty, konečne sumy i nezavisnost vsakogo konečnogo podsistema.\par
\noindent\emph{Textual comparison:} The introduction to Paper 10 distinguishes linear, rational and algebraic bases. The footnote in §4 states the conditions for a Hamel basis: rational coefficients, finite sums and independence of every finite subsystem.\par
\noindent{\footnotesize TEXT-COMP-215; pasaže / passages: 8443, 8445, 8601, 8603; \texttt{CORPUS\_FORM\_COMPARISONS.json}.}\par
\end{minipage}\par
\par\medskip\Needspace{7\baselineskip}
\noindent\begin{minipage}{\linewidth}
\hypertarget{TERM-1097}{}
\noindent\textbf{linearna zavisnost / relacija}\quad{\footnotesize\texttt{TERM-1097}}\par
\noindent\emph{Iztočnik / Source:} lineare Abhängigkeit / Relation\par
\noindent\emph{Rusky / Russian:} линейная зависимость / соотношение\par
\noindent\emph{Ukrajinsky / Ukrainian:} лінійна залежність / співвідношення\par
\smallskip\noindent\emph{Razlog:} Zavisnost označuje matematično svojstvo elementov; relacija označuje ravnost, posrědstvom ktore se to svojstvo vyražaje. V §2 raboty 31 ta razlika pojavjaje se pri sravnjenju koeficientov iz $P$ i $P_i=P\varepsilon_i$: množenje s $\varepsilon_i$ prěnosi zavisnost v jedno napravljanje, a oblik $P_i$ daje obratno napravljanje. Zato sut odgovarjajuče rangi ravne. Arhivny izbor relacija vměsto širšego svakodennogo slova čuva tehničnu ulogu ravnosti. Relacija tu ne znači proizvoljny binarny odnos na množině; kontekst koeficientov ostavaje neobhodny.\par
\noindent\emph{Reason:} Zavisnost names a mathematical property of the elements; relacija names the equality through which that property is expressed. In §2 of Paper 31 this distinction appears when comparing coefficients in $P$ and $P_i=P\varepsilon_i$: multiplication by $\varepsilon_i$ transfers dependence in one direction, while the form of $P_i$ supplies the converse. The corresponding ranks are therefore equal. The recorded preference for relacija over a broader everyday word preserves the technical role of the equality. Relacija here does not mean an arbitrary binary relation on a set; the coefficient context remains essential.\par
\noindent{\footnotesize Izvorne zapisy / Original records: HIST-g122-008.}\par
\smallskip\noindent\emph{Sravnjenje s tekstom:} §2 raboty 31 daje konkretno sravnjenje zavisnosti nad $P$ i nad $P_i=P\varepsilon_i$. Prěnos relacije v oba napravljanja jest matematičny razlog ravnosti rangov, ne prosta zaměna sinonimov.\par
\noindent\emph{Textual comparison:} §2 of Paper 31 supplies a specific comparison of dependence over $P$ and over $P_i=P\varepsilon_i$. Transferring the relation in both directions is the mathematical reason for equality of ranks, not a substitution of synonyms.\par
\noindent{\footnotesize TEXT-COMP-216; pasaže / passages: 17499, 17524; \texttt{CORPUS\_FORM\_COMPARISONS.json}.}\par
\end{minipage}\par
\par\medskip\Needspace{7\baselineskip}
\noindent\begin{minipage}{\linewidth}
\hypertarget{TERM-1098}{}
\noindent\textbf{linearne kombinacije}\quad{\footnotesize\texttt{TERM-1098}}\par
\noindent\emph{Iztočnik / Source:} lineare Aggregate\par
\noindent\emph{Rusky / Russian:} линейные комбинации\par
\noindent\emph{Ukrajinsky / Ukrainian:} лінійні комбінації\par
\smallskip\noindent\emph{Razlog:} Linearne kombinacije prědaje historijne lineare Aggregate po matematičnoj operaciji, ne po vněšnej podobnosti slov. V §5 raboty 2 ide o sumy umnoženyh izrazov; obrat linearnogo sistema daje relacije medžu nimi. Samo agregaty by ne skazalo, či ide o sumu, množinu ili slučajno sobranje. Zachovany razlog daje prědnost jasnomu sovrěmennomu nazivu pri neizměnjenej operaciji. Ruska i ukrajinska forma v sravnjenju sut soglasne s tym izborom. To jest modernizacija terminologije, ne dozvoljenje modernizovati formuly ili izpuščati jih koeficienty.\par
\noindent\emph{Reason:} Linearne kombinacije renders historical lineare Aggregate by its mathematical operation rather than surface resemblance. In §5 of Paper 2 the objects are sums of multiplied expressions; inversion of the linear system produces relations between them. Bare agregaty would not say whether a sum, set or collection was intended. The preserved reason prefers a clear modern name while retaining the operation. The Russian and Ukrainian comparison forms agree with that choice. This modernizes terminology, not the formulas, and does not license omission of their coefficients.\par
\noindent{\footnotesize Izvorne zapisy / Original records: HIST-g006-005.}\par
\smallskip\noindent\emph{Sravnjenje s tekstom:} V §5 raboty 2 operacija jest sravniva s rěšenjem linearnogo sistema. Izraz linearne kombinacije čini operaciju jasnoju bez proměny izvornoj formuly.\par
\noindent\emph{Textual comparison:} In §5 of Paper 2 the operation is connected with solving a linear system. Linearne kombinacije makes the operation clear without changing the source formula.\par
\noindent{\footnotesize TEXT-COMP-217; pasaže / passages: 713, 715, 716, 717, 719; \texttt{CORPUS\_FORM\_COMPARISONS.json}.}\par
\end{minipage}\par
\par\medskip\Needspace{7\baselineskip}
\noindent\begin{minipage}{\linewidth}
\hypertarget{TERM-1099}{}
\noindent\textbf{linearna baza}\quad{\footnotesize\texttt{TERM-1099}}\par
\noindent\emph{Iztočnik / Source:} lineare Basis\par
\noindent\emph{Rusky / Russian:} линейный базис\par
\noindent\emph{Ukrajinsky / Ukrainian:} лінійна база\par
\smallskip\noindent\emph{Razlog:} V §4 raboty 10 linearna baza jest Hamelova baza realnyh čisel nad racionalnymi čislami: vsako realno čislo ima konečno linearno prědstavjenje s racionalnymi koeficientami, a nijeden konečny podsistem ne ima netrivialnoj takovoj relacije. Podnožna primětka čini oba uslovja vidimymi. Baza zato ne jest samo sistem porodžajučih elementov. Arhivny izbor prodolžuje racionalna i algebraična baza v toj samoj rabotě, uz rusko bazis. To terminologično srodstvo ne izbrisuje razliku medžu linearnym razponom i prědstavjenjem posrědstvom racionalnyh ili algebraičnyh operacij.\par
\noindent\emph{Reason:} In §4 of Paper 10, linearna baza is a Hamel basis of the real numbers over the rationals: every real number has a finite linear representation with rational coefficients, and no finite subset has a nontrivial such relation. A footnote makes both conditions explicit. A basis is therefore more than a generating system. The recorded choice continues racionalna and algebraična baza in the same paper, alongside Russian bazis. That terminological relationship does not erase the distinction between linear span and representation through rational or algebraic operations.\par
\noindent{\footnotesize Izvorne zapisy / Original records: HIST-g072-003.}\par
\smallskip\noindent\emph{Sravnjenje s tekstom:} Uvod raboty 10 razlikuje linearnu, racionalnu i algebraičnu bazu. Podnožna primětka v §4 daje uslovja za Hamelovu bazu: racionalne koeficienty, konečne sumy i nezavisnost vsakogo konečnogo podsistema.\par
\noindent\emph{Textual comparison:} The introduction to Paper 10 distinguishes linear, rational and algebraic bases. The footnote in §4 states the conditions for a Hamel basis: rational coefficients, finite sums and independence of every finite subsystem.\par
\noindent{\footnotesize TEXT-COMP-215; pasaže / passages: 8443, 8445, 8601, 8603; \texttt{CORPUS\_FORM\_COMPARISONS.json}.}\par
\end{minipage}\par
\par\medskip\Needspace{7\baselineskip}
\noindent\begin{minipage}{\linewidth}
\hypertarget{TERM-1100}{}
\noindent\textbf{linearna baza grupy ostatkov}\quad{\footnotesize\texttt{TERM-1100}}\par
\noindent\emph{Iztočnik / Source:} lineare Basis der Restgruppe\par
\noindent\emph{Rusky / Russian:} линейный базис группы остатков\par
\noindent\emph{Ukrajinsky / Ukrainian:} лінійний базис групи лишків\par
\smallskip\noindent\emph{Razlog:} Baza grupy ostatkov v §11 raboty 17 znači linearno nezavisne klasy, ktore porodžajut vse klasy s koeficientami iz $P$. Konečna grupa tu ima konečnu razměrnost nad $P$, ne nužno konečno čislo elementov. Teorem XI veže tu razměrnost s čislom nezavisnyh integralov odgovarjajučego sistema diferencialnyh uravnenij. Zato arhivny razlog zachovuje baza, ne samo porodžajuči sistem. Primer s klasami $1,\xi_1,\xi_2,\xi_1\xi_2$ pokazyvaje, kako se to čislo čita. Naziv čuva historijnu Restgruppe, dokolě objasnjenje čini linearnu strukturu javnoju.\par
\noindent\emph{Reason:} In §11 of Paper 17, baza grupy ostatkov means linearly independent classes spanning all classes with coefficients in $P$. A finite group here has finite dimension over $P$, not necessarily finitely many elements. Theorem XI connects that dimension with the number of independent integrals of the associated differential-equation system. The archived reason therefore retains basis rather than merely generating system. The example with classes $1,\xi_1,\xi_2,\xi_1\xi_2$ shows how the number is read. The name preserves historical Restgruppe, while the explanation makes its linear structure explicit.\par
\noindent{\footnotesize Izvorne zapisy / Original records: HIST-g079-064.}\par
\smallskip\noindent\emph{Sravnjenje s tekstom:} Definicija v §11 raboty 17 ustanovjaje linearnu nezavisnost i red nad poljem P. Teorem XI veže toj red s čislom nezavisnyh integralov. Čitanje konečna kako konečnogo čisla elementov by utratilo tu definiciju.\par
\noindent\emph{Textual comparison:} The definition in §11 of Paper 17 specifies linear independence and order over the field P. Theorem XI connects that order with the number of independent integrals. Reading finite as finite cardinality would lose the definition.\par
\noindent{\footnotesize TEXT-COMP-218; pasaže / passages: 11980, 11982, 11986, 11988, 11989, 11991, 11993, 11994; \texttt{CORPUS\_FORM\_COMPARISONS.json}.}\par
\end{minipage}\par
\par\medskip\Needspace{7\baselineskip}
\noindent\begin{minipage}{\linewidth}
\hypertarget{TERM-1101}{}
\noindent\textbf{linearna familija form}\quad{\footnotesize\texttt{TERM-1101}}\par
\noindent\emph{Iztočnik / Source:} lineare Formenschar\par
\noindent\emph{Rusky / Russian:} линейное семейство форм\par
\noindent\emph{Ukrajinsky / Ukrainian:} лінійна сім’я форм\par
\smallskip\noindent\emph{Razlog:} Linearna familija form v rabotě 8 jest definirana zamknjenostju pri tvorjenju $c_1\theta_1+c_2\theta_2$ nad zadanym poljem $K$. Formy imajut jednaky stepenj, a $K$ sodrživaje jih koeficienty; podnožna primětka izrično govori o občejših slučajah, v ktoryh rang može byti beskonečny. Familija zato ne znači proizvoljno sobranje niti samo jednoparametričnu familiju. Arhivny razlog upotrěbljaje metaforu semejstva, podobnu ruskomu i ukrajinskomu sravnjenju. On ne dava konkretnogo iztočnika za svoje tvrđenje o sovrěmennoj normě; podrobna definicija jest silnějša opora za smysl nego to obče tvrđenje.\par
\noindent\emph{Reason:} In Paper 8, linearna familija form is defined by closure under $c_1\theta_1+c_2\theta_2$ over a specified field $K$. The forms have the same degree and $K$ contains their coefficients; a footnote explicitly discusses broader cases in which rank can be infinite. Familija therefore means neither an arbitrary collection nor only a one-parameter family. The archived explanation adopts a family metaphor resembling the Russian and Ukrainian comparanda. It gives no specific source for its claim about contemporary usage; the detailed definition is stronger evidence for the meaning than that general claim.\par
\noindent{\footnotesize Izvorne zapisy / Original records: HIST-g057-006.}\par
\smallskip\noindent\emph{Sravnjenje s tekstom:} Raboty 6 i 8 daju izrične definicije linearnoj familije bez ograničenja na jednoparametričny slučaj. Uvod raboty 16 govori o linearnyh snopeh form, ale toj odlomok ne dokazuje užše tolkovanje pencil iz glosara.\par
\noindent\emph{Textual comparison:} Papers 6 and 8 give explicit definitions of linear family without restriction to one parameter. The introduction to Paper 16 speaks of linearne snope form, but that passage does not establish the narrower pencil interpretation in the glossary.\par
\noindent{\footnotesize TEXT-COMP-219; pasaže / passages: 5740, 5742, 5744, 5746, 5748, 5750, 6755, 6806, 6812, 10729, 10731, 10733, 10759; \texttt{CORPUS\_FORM\_COMPARISONS.json}.}\par
\end{minipage}\par
\par\medskip\Needspace{7\baselineskip}
\noindent\begin{minipage}{\linewidth}
\hypertarget{TERM-1102}{}
\noindent\textbf{linearne snope form}\quad{\footnotesize\texttt{TERM-1102}}\par
\noindent\emph{Iztočnik / Source:} lineare Formenscharen\par
\noindent\emph{Rusky / Russian:} линейные пучки форм\par
\noindent\emph{Ukrajinsky / Ukrainian:} лінійні пучки форм\par
\smallskip\noindent\emph{Razlog:} Linearne snope form jest druga prěvodna forma německogo Formenscharen. Arhivny razlog nazval objekt linearnym snopom v smyslu pencil, ale sravnjeny uvod raboty 16 ne ustanovjaje ograničenje na jednoparametričnu familiju: on govori o razkladah form, polarnyh procesah i modulnyh sistemah. Zato užše tolkovanje ne može byti prějeto samo na osnově glosara. Ruske i ukrajinske pučki dajut sravnjenje metafory, ne dokaz razměrnosti. Snop tu ne znači avtomatično snop v sovrěmennoj teoriji snopov. Forma i jej smyslovy obseg ostavajut prědlogom, ktory potrěbuje točnoj disciplinarnoj opory.\par
\noindent\emph{Reason:} Linearne snope form is another rendering of German Formenscharen. The archived rationale calls the object a linear pencil, but the compared introduction to Paper 16 establishes no restriction to a one-parameter family: it discusses expansions of forms, polar processes and module systems. The narrower interpretation cannot therefore be accepted on the glossary’s authority alone. Russian and Ukrainian pučki provide a comparison of metaphors, not evidence of dimension. Nor does snop automatically mean a sheaf in modern sheaf theory. The form and its semantic range remain proposals requiring precise disciplinary support.\par
\noindent{\footnotesize Izvorne zapisy / Original records: HIST-g078-009.}\par
\smallskip\noindent\emph{Sravnjenje s tekstom:} Raboty 6 i 8 daju izrične definicije linearnoj familije bez ograničenja na jednoparametričny slučaj. Uvod raboty 16 govori o linearnyh snopeh form, ale toj odlomok ne dokazuje užše tolkovanje pencil iz glosara.\par
\noindent\emph{Textual comparison:} Papers 6 and 8 give explicit definitions of linear family without restriction to one parameter. The introduction to Paper 16 speaks of linearne snope form, but that passage does not establish the narrower pencil interpretation in the glossary.\par
\noindent{\footnotesize TEXT-COMP-219; pasaže / passages: 5740, 5742, 5744, 5746, 5748, 5750, 6755, 6806, 6812, 10729, 10731, 10733, 10759; \texttt{CORPUS\_FORM\_COMPARISONS.json}.}\par
\end{minipage}\par
\par\medskip\Needspace{7\baselineskip}
\noindent\begin{minipage}{\linewidth}
\hypertarget{TERM-1103}{}
\noindent\textbf{linearne sistemy jednačenj}\quad{\footnotesize\texttt{TERM-1103}}\par
\noindent\emph{Iztočnik / Source:} lineare Gleichungssysteme\par
\noindent\emph{Rusky / Russian:} линейные системы уравнений\par
\noindent\emph{Ukrajinsky / Ukrainian:} лінійні системи рівнянь\par
\smallskip\noindent\emph{Razlog:} Linearne sistemy jednačenj označuje uslovja, ktore trěba rěšiti za koeficienty prědstavjenja elementa. V §3 raboty 30 prěhod k konjugovanym elementam daje takovy sistem; nenulovost diskriminantnogo izraza i integralna zamknjenost imajut svojo město v argumentu. Naziv ne govori, že vsaky linearny sistem ima rěšenje ili že koeficienty mogut byti izbrany proizvoljno. Arhivny razlog prědpočitaje sovrěmennu matematičnu frazu. Sravnjeny tekst imaje linearny sistem jednačenij; ta morfologična forma ne jest identična glosarnomu množnomu zaglavju. Matematičny pojem ostavaje toj samy.\par
\noindent\emph{Reason:} Linearne sistemy jednačenj names the conditions to be solved for the coefficients representing an element. In §3 of Paper 30, passing to the conjugate elements gives such a system; nonvanishing of the discriminant expression and integral closedness have their own roles in the argument. The term does not assert that every linear system is solvable or that its coefficients may be chosen arbitrarily. The archived reason prefers modern mathematical phrasing. The compared passage has singular linearny sistem jednačenij rather than the glossary’s plural heading. The mathematical concept is the same.\par
\noindent{\footnotesize Izvorne zapisy / Original records: HIST-g110-011.}\par
\smallskip\noindent\emph{Sravnjenje s tekstom:} §3 raboty 30 upotrěbjaje konjugovane elementy i nenulovy diskriminantny izraz. Naziv linearny sistem ne izpušča ty uslovja niti ne zaměnjuje dokaz integralnosti koeficientov.\par
\noindent\emph{Textual comparison:} §3 of Paper 30 uses conjugate elements and a nonzero discriminant expression. Calling the equations a linear system neither removes those hypotheses nor replaces the argument about integrality of coefficients.\par
\noindent{\footnotesize TEXT-COMP-220; pasaže / passages: 16391, 16395; \texttt{CORPUS\_FORM\_COMPARISONS.json}.}\par
\end{minipage}\par
\par\medskip\Needspace{7\baselineskip}
\noindent\begin{minipage}{\linewidth}
\hypertarget{TERM-1104}{}
\noindent\textbf{linearna mnogovidnost}\quad{\footnotesize\texttt{TERM-1104}}\par
\noindent\emph{Iztočnik / Source:} lineare Mannigfaltigkeit\par
\noindent\emph{Rusky / Russian:} линейное многообразие\par
\noindent\emph{Ukrajinsky / Ukrainian:} лінійний многовид\par
\smallskip\noindent\emph{Razlog:} Mnogovidnost v rabotě 10 označuje geometričnu množinu cěnnostij, ktorym se možno pribliziti v okolině zadanoj točky. Pri fiksovanyh $a,b$ relacija $c_1a+c_2b+c_3A+c_4B=0$ oprěděljaje prěmu v ravnině cěnnostij $A,B$; ona ne mora prohoditi črez početok. Zato historijno linearna tu dopušćaje afinnu geometriju. Ne ide o novo prědpoloženje gladkosti niti o ranějše množiny rešetkovyh toček. Arhivny zapis izrično tretuje mnogovidnost kak konstruovany medžuslovjansky termin, ne kak potvrđenu normu. Matematična definicija objasnjaje pojem, ale ne rěšaje sama jezykovu prirodnost naziva.\par
\noindent\emph{Reason:} In Paper 10, mnogovidnost denotes the geometric set of values that can be approached near a specified point. With $a,b$ fixed, $c_1a+c_2b+c_3A+c_4B=0$ defines a line in the $A,B$ value plane; it need not pass through the origin. Historical linear thus allows an affine interpretation here. This is neither an added smoothness assumption nor the earlier collection of lattice points. The archival record explicitly treats mnogovidnost as a constructed Interslavic mathematical term, not an established norm. The mathematical definition clarifies the concept but does not by itself settle the name’s linguistic naturalness.\par
\noindent{\footnotesize Izvorne zapisy / Original records: HIST-g072-007.}\par
\smallskip\noindent\emph{Sravnjenje s tekstom:} V §4 raboty 10 fiksovane veličiny $a,b$ i linearna relacija v $A,B$ dopuščajut afinnu prěmu; drugy slučaj daje dvurazměrnu oblast cěnnostij. Mnogovidnost tu ne prinosi novogo prědpoloženja gladkosti.\par
\noindent\emph{Textual comparison:} In §4 of Paper 10, fixed $a,b$ and a linear relation in $A,B$ allow an affine line; the other case gives a two-dimensional value set. Mnogovidnost introduces no new smoothness assumption here.\par
\noindent{\footnotesize TEXT-COMP-221; pasaže / passages: 8635, 8637, 8638, 8640, 8642, 8644, 8646; \texttt{CORPUS\_FORM\_COMPARISONS.json}.}\par
\end{minipage}\par
\par\medskip\Needspace{7\baselineskip}
\noindent\begin{minipage}{\linewidth}
\hypertarget{TERM-1105}{}
\noindent\textbf{linearna familija}\quad{\footnotesize\texttt{TERM-1105}}\par
\noindent\emph{Iztočnik / Source:} lineare Schar\par
\noindent\emph{Rusky / Russian:} линейное семейство\par
\noindent\emph{Ukrajinsky / Ukrainian:} лінійна сім'я\par
\smallskip\noindent\emph{Razlog:} Linearna familija v §7 raboty 6 sodrživaje s vsakoju funkcijeju $f$ takože $cf$, i s funkcijami $f,g$ takože $f+g$. Množenje dvuh funkcij jest dodatno uslovje slědujuče definicije, ne čest samoj linearnoj familije. Arhivny razlog zato odbira familija, ne užši naziv za jednorazměrny projektivny snop. Internacionalno korenje bylo izbrano radi očekyvane prozračnosti, ale zapis ostavlja termin jezykovo nepodtvrđenym. Oba uslovja i oblast koeficientov sut nužne za čitatelja; bez njih by naziv skryl rozděl medžu linearnym sistemom i prstenovoju strukturoju.\par
\noindent\emph{Reason:} In §7 of Paper 6, linearna familija contains $cf$ with each function $f$, and $f+g$ with functions $f,g$. Multiplication of two functions is an additional condition in the next definition, not part of linearna familija itself. The archived reason therefore chooses family rather than a narrower name for a projective pencil. The international stem was selected for anticipated transparency, but the record leaves the term linguistically unconfirmed. Both closure conditions and the coefficient domain matter: without them the name would conceal the distinction between a linear system and a ring structure.\par
\noindent{\footnotesize Izvorne zapisy / Original records: HIST-g047-002.}\par
\smallskip\noindent\emph{Sravnjenje s tekstom:} Raboty 6 i 8 daju izrične definicije linearnoj familije bez ograničenja na jednoparametričny slučaj. Uvod raboty 16 govori o linearnyh snopeh form, ale toj odlomok ne dokazuje užše tolkovanje pencil iz glosara.\par
\noindent\emph{Textual comparison:} Papers 6 and 8 give explicit definitions of linear family without restriction to one parameter. The introduction to Paper 16 speaks of linearne snope form, but that passage does not establish the narrower pencil interpretation in the glossary.\par
\noindent{\footnotesize TEXT-COMP-219; pasaže / passages: 5740, 5742, 5744, 5746, 5748, 5750, 6755, 6806, 6812, 10729, 10731, 10733, 10759; \texttt{CORPUS\_FORM\_COMPARISONS.json}.}\par
\end{minipage}\par
\par\medskip\Needspace{7\baselineskip}
\noindent\begin{minipage}{\linewidth}
\hypertarget{TERM-1106}{}
\noindent\textbf{linearna transformacija}\quad{\footnotesize\texttt{TERM-1106}}\par
\noindent\emph{Iztočnik / Source:} lineare Transformation\par
\noindent\emph{Rusky / Russian:} линейное преобразование\par
\noindent\emph{Ukrajinsky / Ukrainian:} лінійне перетворення\par
\smallskip\noindent\emph{Razlog:} Linearna transformacija jest v §3 raboty 6 operacija, ktoroju možno specializovati něktore koeficienty osnovnoj formy, zachovajuči razgledane relacije invariantov. Transformacija zachovuje internacionalno korenje, dokolě rusky i ukrajinsky sravniteljny izraz upotrěbjajut domače nazvy prěobraženja. Zachovany razlog dava samo kratku ocěnku obyčajnosti termina; on ne dokumentuje porjadkovanje alternativ po čislovyh kriterijah. Naziv ne směje byti zaměnjen s linearno-drobnoju transformacijeju iz §6: vtoroj frazě prinaleži děljenje, ktoro samo linearna ne vyražaje.\par
\noindent\emph{Reason:} In §3 of Paper 6, linearna transformacija is the operation used to specialize certain coefficients of the basic form while retaining the invariant relations under discussion. Transformacija retains an international stem, whereas the Russian and Ukrainian comparanda use native transformation terms. The preserved explanation only gives a brief judgment that the vocabulary is standard; it documents no numerical ranking of alternatives. The term must not be conflated with the fractional-linear transformation in §6: the latter involves division, which linear alone does not express.\par
\noindent{\footnotesize Izvorne zapisy / Original records: HIST-g043-012.}\par
\smallskip\noindent\emph{Sravnjenje s tekstom:} §3 raboty 6 opisuje čislovo fiksovanje někojih koeficientov osnovnoj formy s zachovanjem invariantnyh relacij. Ta upotrěba ne jest identična linearno-drobnoj transformaciji iz §6.\par
\noindent\emph{Textual comparison:} §3 of Paper 6 describes numerically fixing some coefficients of the basic form while preserving the invariant relations. This use is not identical to the fractional-linear transformation in §6.\par
\noindent{\footnotesize TEXT-COMP-222; pasaže / passages: 5384, 5696; \texttt{CORPUS\_FORM\_COMPARISONS.json}.}\par
\end{minipage}\par
\par\medskip\Needspace{7\baselineskip}
\noindent\begin{minipage}{\linewidth}
\hypertarget{TERM-1107}{}
\noindent\textbf{linearna nezavisimost / definicija ravnosti}\quad{\footnotesize\texttt{TERM-1107}}\par
\noindent\emph{Iztočnik / Source:} lineare Unabhängigkeit / Gleichheitsdefinition\par
\noindent\emph{Rusky / Russian:} линейная независимость / определение равенства\par
\noindent\emph{Ukrajinsky / Ukrainian:} лінійна незалежність / означення рівності\par
\smallskip\noindent\emph{Razlog:} Definicija ravnosti jest tu čestju konstrukcije razširjenogo prstena, ne obyčajna primětka o znaku ravnosti. V §3 raboty 31 nezavisne stepeni elementa $\gamma$ ostavajut nezavisnymi nad razširjenym poljem po toj definiciji. To daje nenulovy element $G(\gamma)$, čija $p$-ta stepenj jest nula. Izpuščenje definicije by skrylo, začto dva tvrđenja ne protivorěčat jedno drugomu. Arhivny razlog zato veže nezavisimost i definiciju ravnosti. Oblik nezavisimost i drugde upotrěbljany nezavisnost ne dokazuju dva različna matematična pojma; razlika jest v tvorjenju slova.\par
\noindent\emph{Reason:} Definicija ravnosti is part of the construction of the extended ring, not a routine comment on the equals sign. In §3 of Paper 31, independent powers of $\gamma$ remain independent over the extended field by that definition. This produces a nonzero element $G(\gamma)$ whose $p$th power is zero. Omitting the definition would hide why those assertions are compatible. The archived explanation therefore connects independence with the definition of equality. Nezavisimost and the alternative nezavisnost used elsewhere do not establish two different mathematical concepts; they differ in word formation.\par
\noindent{\footnotesize Izvorne zapisy / Original records: HIST-g128-007.}\par
\smallskip\noindent\emph{Sravnjenje s tekstom:} Definicija ravnosti razširjenogo prstena čuva nezavisnost nižših stepenjev $\gamma$. Zato $G(\gamma)$ ne izčezaje, dokolě jego $p$-ta stepenj izčezaje. Ta razlika jest ključ k argumentu, ne jezykovy ukras.\par
\noindent\emph{Textual comparison:} The equality definition in the extended ring preserves independence of the lower powers of $\gamma$. Thus $G(\gamma)$ is nonzero although its $p$th power vanishes. That distinction is central to the argument, not verbal decoration.\par
\noindent{\footnotesize TEXT-COMP-223; pasaže / passages: 17735, 17737, 17739, 17741, 17743, 17745, 17747, 17748, 17750; \texttt{CORPUS\_FORM\_COMPARISONS.json}.}\par
\end{minipage}\par
\par\medskip\Needspace{7\baselineskip}
\noindent\begin{minipage}{\linewidth}
\hypertarget{TERM-1108}{}
\noindent\textbf{linearno sočinenje / linearna forma}\quad{\footnotesize\texttt{TERM-1108}}\par
\noindent\emph{Iztočnik / Source:} lineare Verbindung / Linearform\par
\noindent\emph{Rusky / Russian:} линейное соединение / линейная форма\par
\noindent\emph{Ukrajinsky / Ukrainian:} лінійна сполука / лінійна форма\par
\smallskip\noindent\emph{Razlog:} Linearno sočinenje podčrkivaje formalno sostavljanje novyh elementov pri razširjenju prstena v §1 raboty 31. Tekst najprvo piše sumy $\sum \bar c_i\gamma_i$, a potom definovaje jih ravnost, sumu i produkt. Arhivny razlog zachovuje bližšu vez s německym Verbindung, vměsto avtomatičnoj zaměny s kombinacija. No v toj samoj rabotě kombinacije takože nastupajut; razlika slov sama ne ustanovjaje razliku operacij. Linearna forma označuje izraz s koeficientami, a ne nužno uravnenje. Izbor sočinenje jest obosnovany naměroju formalnoj konstrukcije, ne dokazanoju občeju jezykovoj prěvahoj.\par
\noindent\emph{Reason:} Linearno sočinenje emphasizes the formal assembly of new elements in the ring extension in §1 of Paper 31. The text first writes sums $\sum \bar c_i\gamma_i$ and then defines their equality, addition and multiplication. The archived reason retains a closer link to German Verbindung instead of automatically substituting combination. Yet kombinacije also occurs in the same paper; the wording alone does not establish different operations. Linearna forma names an expression with coefficients, not necessarily an equation. Sočinenje is motivated by the formal construction, not by demonstrated general linguistic superiority.\par
\noindent{\footnotesize Izvorne zapisy / Original records: HIST-g119-009.}\par
\smallskip\noindent\emph{Sravnjenje s tekstom:} Rabota 31 upotrěbjaje kombinacije i sočinenja pri svojem opisu razširjenja, a rabota 22 izrično piše konečne sumy s racionalnymi funkcijami novyh neoprěděljenyh kako koeficientami. Razlika nazvanij ne dokazuje razliku matematičnyh operacij.\par
\noindent\emph{Textual comparison:} Paper 31 uses both kombinacije and sočinenja in its account of extension, while Paper 22 explicitly writes finite sums with rational functions of the new indeterminates as coefficients. Different names do not by themselves establish different mathematical operations.\par
\noindent{\footnotesize TEXT-COMP-224; pasaže / passages: 17316, 17326, 17328, 17330, 17332, 17334, 17335, 17339, 17341, 17342, 14514, 14516, 14518; \texttt{CORPUS\_FORM\_COMPARISONS.json}.}\par
\end{minipage}\par
\par\medskip\Needspace{7\baselineskip}
\noindent\begin{minipage}{\linewidth}
\hypertarget{TERM-1109}{}
\noindent\textbf{linearne kombinacije}\quad{\footnotesize\texttt{TERM-1109}}\par
\noindent\emph{Iztočnik / Source:} lineare Verbindungen\par
\noindent\emph{Rusky / Russian:} линейные комбинации\par
\noindent\emph{Ukrajinsky / Ukrainian:} лінійні комбінації\par
\smallskip\noindent\emph{Razlog:} Linearne kombinacije prědaje lineare Verbindungen v §3 raboty 22 kako sumy s koeficientami. Pri prisojedinjenju neoprěděljenyh $u$ k polju koeficientov tekst izrično dopušča racionalne funkcije od $u$ kako koeficienty konečnyh sum. Glosar daje prědnost sovrěmennomu nazivu vměsto doslovnogo spojanja. To jest inačny jezykovy izbor nego sočinenje v rabotě 31, ne samo po sobě inačny vid linearnosti. Konečnost sumy i oblast koeficientov ne mogut byti izpuščene pod občim nazvom kombinacija.\par
\noindent\emph{Reason:} Linearne kombinacije renders lineare Verbindungen in §3 of Paper 22 as sums with coefficients. When the indeterminates $u$ are adjoined to the coefficient field, the text explicitly allows rational functions of $u$ as coefficients in finite sums. The glossary prefers modern combination terminology to a literal connection. This differs linguistically from sočinenje in Paper 31, but does not itself identify a different kind of linearity. Finiteness of the sum and the coefficient domain must not disappear beneath the general name combination.\par
\noindent{\footnotesize Izvorne zapisy / Original records: HIST-g087-004, HIST-g092-024.}\par
\smallskip\noindent\emph{Sravnjenje s tekstom:} Rabota 31 upotrěbjaje kombinacije i sočinenja pri svojem opisu razširjenja, a rabota 22 izrično piše konečne sumy s racionalnymi funkcijami novyh neoprěděljenyh kako koeficientami. Razlika nazvanij ne dokazuje razliku matematičnyh operacij.\par
\noindent\emph{Textual comparison:} Paper 31 uses both kombinacije and sočinenja in its account of extension, while Paper 22 explicitly writes finite sums with rational functions of the new indeterminates as coefficients. Different names do not by themselves establish different mathematical operations.\par
\noindent{\footnotesize TEXT-COMP-224; pasaže / passages: 17316, 17326, 17328, 17330, 17332, 17334, 17335, 17339, 17341, 17342, 14514, 14516, 14518; \texttt{CORPUS\_FORM\_COMPARISONS.json}.}\par
\end{minipage}\par
\par\medskip\Needspace{7\baselineskip}
\noindent\begin{minipage}{\linewidth}
\hypertarget{TERM-1110}{}
\noindent\textbf{polna linearna sistema}\quad{\footnotesize\texttt{TERM-1110}}\par
\noindent\emph{Iztočnik / Source:} lineare Vollschar\par
\noindent\emph{Rusky / Russian:} полная линейная система\par
\noindent\emph{Ukrajinsky / Ukrainian:} повна лінійна система\par
\smallskip\noindent\emph{Razlog:} Polna linearna sistema v rabotě 14 obhvača vse koresidualne grupy toček s tym samym čislom toček kako zadana grupa. Rěč polna jest zato sučstvena: ne ide samo o proizvoljny podsistem. Arhivny razlog upotrěbjaje naziv poznavajemy iz algebraičnoj geometrije, ne doslovno polna familija ili snop. Sam tekst ale upozornjaje, že koresidualne grupy obče ne morajut iměti jednako čislo toček. To ograničenje čuva vez s klasami divizorov. Historijnu konvenciju čislenja dimenzije trěba čitati iz teksta, ne tiho zaměnjati sovrěmennoju konvencijeju.\par
\noindent\emph{Reason:} In Paper 14, polna linearna sistema contains all co-residual point groups with the same number of points as the given group. Complete is therefore substantive: this is not an arbitrary subsystem. The archived reason uses the name recognizable from algebraic geometry rather than literally full family or sheaf. The text itself warns that co-residual point groups need not generally have the same number of points. That restriction maintains the link with divisor classes. The historical convention for counting dimension must be read from the passage, not silently replaced by a modern convention.\par
\noindent{\footnotesize Izvorne zapisy / Original records: HIST-g076-021.}\par
\smallskip\noindent\emph{Sravnjenje s tekstom:} Rabota 14 razlikuje koresidualnost bez ograničenja čisla toček od polnoj linearnoj sistemy s fiksovanym čislom. Definicija dimenzije takože mora ostati vezana s historijnym tekstom.\par
\noindent\emph{Textual comparison:} Paper 14 distinguishes co-residuality without a restriction on point count from the complete linear system with fixed point count. Its dimension convention likewise has to remain tied to the historical text.\par
\noindent{\footnotesize TEXT-COMP-225; pasaže / passages: 10022, 10024, 10026, 10028; \texttt{CORPUS\_FORM\_COMPARISONS.json}.}\par
\end{minipage}\par
\par\medskip\Needspace{7\baselineskip}
\noindent\begin{minipage}{\linewidth}
\hypertarget{TERM-1111}{}
\noindent\textbf{linearny faktor}\quad{\footnotesize\texttt{TERM-1111}}\par
\noindent\emph{Iztočnik / Source:} Linearfaktor\par
\noindent\emph{Rusky / Russian:} линейный множитель\par
\noindent\emph{Ukrajinsky / Ukrainian:} лінійний множник\par
\smallskip\noindent\emph{Razlog:} Linearny faktor v §6 raboty 22 ima vid $x_i-\bar x_i$ v odgovarjajučem razširjenom polju. Internacionalno faktor bylo izbrano radi kratkosti; rusko i ukrajinsko sravnjenje upotrěbjajut korenje množiti. Tu ne ide o faktor-grupu. Tekst veže faktor s nulju ideala i potom zadržuje mnogokratnosti v produktnom zapisu. Pri tom sravnjenje formul otkryvaje oddělnu nejasnost: substitucija (35) ima minusy, a slědujuči zapis transformovanoj nuly plusy. Točny znak potrěbuje izvornogo svědka; jasny naziv faktora ne potvrdžaje avtomatično ty detali.\par
\noindent\emph{Reason:} In §6 of Paper 22, linearny faktor has the form $x_i-\bar x_i$ over the relevant extension field. The international factor stem was preferred for compactness; the Russian and Ukrainian comparanda use the multiplication stem. This is not a quotient-group sense of factor. The text connects a factor with a zero of the ideal and then retains multiplicities in the product. Comparing the formulas also exposes a separate ambiguity: substitution (35) uses minus signs, while the subsequent transformed zero uses plus signs. The exact signs require the original witness; clear factor terminology does not certify those details.\par
\noindent{\footnotesize Izvorne zapisy / Original records: HIST-g090-010, HIST-g092-064.}\par
\smallskip\noindent\emph{Sravnjenje s tekstom:} Teorem XII raboty 22 veže faktor s nulju ideala. Slědujuča substitucija (35) ima minusy, a zapis transformovanoj nuly plusy. To sravnjenje ne opravdyvaje tihe izpravjenje; točny izvor mora razrěšiti nesoglasje.\par
\noindent\emph{Textual comparison:} Theorem XII in Paper 22 connects a factor with a zero of the ideal. The following substitution (35) has minus signs, but the expression for the transformed zero has plus signs. This comparison does not warrant a silent correction; the original witness must resolve the disagreement.\par
\noindent{\footnotesize TEXT-COMP-226; pasaže / passages: 14952, 14956, 14958, 14961, 14963, 14965, 14967, 14969; \texttt{CORPUS\_FORM\_COMPARISONS.json}.}\par
\end{minipage}\par
\par\medskip\Needspace{7\baselineskip}
\noindent\begin{minipage}{\linewidth}
\hypertarget{TERM-1112}{}
\noindent\textbf{linearne faktory}\quad{\footnotesize\texttt{TERM-1112}}\par
\noindent\emph{Iztočnik / Source:} Linearfaktoren\par
\noindent\emph{Rusky / Russian:} линейные множители\par
\noindent\emph{Ukrajinsky / Ukrainian:} лінійні множники\par
\smallskip\noindent\emph{Razlog:} Linearne faktory ima različne ulogi v trěh zabilěženih kontekstah. Rabota 4 zaměnjuje stepenj faktora produktom linearnyh form, ktore potom možno položiti ravnymi; rabota 24 veže faktory s nuljami i dopušča višekratnost pri nedoskonalom polju; rabota 31 upotrěbjaje razpadanje na linearne faktory za opis polja razloženja. Zapis daje prědnost kontinuitetu internacionalnogo faktor. Sam naziv ne govori, že faktory sut različne ili že jih koeficienty ležet v početkovom polju. Ty uslovja prinadležet konkretnomu tvrđenju i morajut byti zachovane.\par
\noindent\emph{Reason:} Linearne faktory has distinct roles in its three recorded contexts. Paper 4 replaces a power of a factor by a product of linear forms that can subsequently be set equal; Paper 24 connects factors with zeros and allows multiplicity over an imperfect field; Paper 31 uses splitting into linear factors to describe a splitting field. The recorded preference is for continuity in the international factor stem. The name alone does not assert that the factors are distinct or that their coefficients lie in the original field. Those qualifications belong to the particular assertion and must be retained.\par
\noindent{\footnotesize Izvorne zapisy / Original records: HIST-g035-011, HIST-g097-014, HIST-g129-008.}\par
\smallskip\noindent\emph{Sravnjenje s tekstom:} Rabota 4 vvodit linearne faktory kako pomočne formy; rabota 24 zadržuje mnogokratnosti; rabota 31 harakterizuje polje razloženja. Termin ne daje sam ni različnost faktorov ni jih prinadležnost početkovomu polju.\par
\noindent\emph{Textual comparison:} Paper 4 introduces linear factors as auxiliary forms; Paper 24 retains multiplicities; Paper 31 characterizes a splitting field. The term alone guarantees neither distinctness of the factors nor their belonging to the original field.\par
\noindent{\footnotesize TEXT-COMP-227; pasaže / passages: 4435, 15484, 15496, 15498, 17781, 17790, 17794; \texttt{CORPUS\_FORM\_COMPARISONS.json}.}\par
\end{minipage}\par
\par\medskip\Needspace{7\baselineskip}
\noindent\begin{minipage}{\linewidth}
\hypertarget{TERM-1113}{}
\noindent\textbf{linearne faktory / produktno predstavjenje}\quad{\footnotesize\texttt{TERM-1113}}\par
\noindent\emph{Iztočnik / Source:} Linearfaktoren / Produktdarstellung\par
\noindent\emph{Rusky / Russian:} линейные множители / представление произведением\par
\noindent\emph{Ukrajinsky / Ukrainian:} лінійні множники / подання добутком\par
\smallskip\noindent\emph{Razlog:} Produktno prědstavjenje govori, kako objekt jest zapisan; linearne faktory govori, iz čego toj produkt sostoji. Teorem XI raboty 24 dava faktory normy s eksponentami $\delta_\lambda\rho_\lambda$. Obe česti eksponenta sut važne: naziv linearny ne odstranjava mnogokratnost, a pri nedoskonalyh poljah može nastupati dodatna stepenj karakteristiki. Arhivny razlog daje prědnost internacionalnym korenjam produkt i faktor radi matematičnoj poznavajemosti. To ne jest izměrena razumlivost. Objasnjenje mora zachovati nulje, osnovno polje i eksponenty, ne samo skazati, že norma razpadaje se.\par
\noindent\emph{Reason:} Produktno prědstavjenje says how the object is written; linearne faktory says what the product consists of. Theorem XI in Paper 24 gives factors of the norm with exponents $\delta_\lambda\rho_\lambda$. Both parts of the exponent matter: linear does not remove multiplicity, and imperfect fields can introduce a further power of the characteristic. The archived reason prefers the international product and factor stems for mathematical recognizability. That is not measured comprehension. An explanation must retain the zeros, base field and exponents, not merely say that the norm splits.\par
\noindent{\footnotesize Izvorne zapisy / Original records: HIST-g099-010.}\par
\smallskip\noindent\emph{Sravnjenje s tekstom:} Teorem XI raboty 24 čuva oba množitelja eksponenta. Odgovarjajuče nulje i svojstva osnovnogo polja ostavajut čestju smysla, iako vsaky faktor jest linearny.\par
\noindent\emph{Textual comparison:} Theorem XI in Paper 24 retains both factors of the exponent. The corresponding zeros and properties of the base field remain part of the assertion even though each factor is linear.\par
\noindent{\footnotesize TEXT-COMP-228; pasaže / passages: 15631, 15633, 15634, 15635, 15637, 15639; \texttt{CORPUS\_FORM\_COMPARISONS.json}.}\par
\end{minipage}\par
\par\medskip\Needspace{7\baselineskip}
\noindent\begin{minipage}{\linewidth}
\hypertarget{TERM-1114}{}
\noindent\textbf{linearne faktory; odnorodnjene korenji}\quad{\footnotesize\texttt{TERM-1114}}\par
\noindent\emph{Iztočnik / Source:} Linearfaktoren; homogen gemachte Wurzeln\par
\noindent\emph{Rusky / Russian:} линейные множители; однородненные корни\par
\noindent\emph{Ukrajinsky / Ukrainian:} лінійні множники; одноріднені корені\par
\smallskip\noindent\emph{Razlog:} Odnorodnjene korenji svezany sut s razloženjem binarnoj formy na linearne faktory i s odnorodnym zapisom korenjev. Ne ide o izvlečenje čislovogo korenja niti samo o pereimenovanje faktora. V rabotě 15 funkcija takyh korenjev mora izpolnjati tri razne uslovja: invariantnost, věsu i simetriju; poslědnje dva čuvajut cělu racionalnu zavisnost od koeficientov osnovnyh form. Arhivny izbor kombinuje internacionalno faktor so slavjanskymi odnorodnjeny i korenj. Njegova ocěnka prozračnosti jest razlog za prědlog, ne rezultat jezykovogo izpytanja.\par
\noindent\emph{Reason:} Odnorodnjene korenji connects the decomposition of a binary form into linear factors with homogeneous coordinates for its roots. This is neither numerical root extraction nor merely renaming a factor. In Paper 15, a function of such roots must satisfy three separate conditions: invariance, weight and symmetry; the latter two preserve polynomial dependence on the coefficients of the original forms. The recorded choice combines international faktor with Slavic odnorodnjeny and korenj. Its judgment of transparency motivates the proposal; it is not a language-testing result.\par
\noindent{\footnotesize Izvorne zapisy / Original records: HIST-g077-005.}\par
\smallskip\noindent\emph{Sravnjenje s tekstom:} Uvod raboty 15 veže linearne faktory s korenjami i razlikuje uslovja invariantnosti, věsy i simetrije. Korenj tu ne jest nazvom operacije čislovogo izvlečenja korenja.\par
\noindent\emph{Textual comparison:} The introduction to Paper 15 connects linear factors with roots and distinguishes invariance, weight and symmetry conditions. Root here does not name numerical root extraction.\par
\noindent{\footnotesize TEXT-COMP-229; pasaže / passages: 10091, 10094, 10096, 10098; \texttt{CORPUS\_FORM\_COMPARISONS.json}.}\par
\end{minipage}\par
\par\medskip\Needspace{7\baselineskip}
\noindent\begin{minipage}{\linewidth}
\hypertarget{TERM-1115}{}
\noindent\textbf{linearne formy}\quad{\footnotesize\texttt{TERM-1115}}\par
\noindent\emph{Iztočnik / Source:} Linearformen\par
\noindent\emph{Rusky / Russian:} линейные формы\par
\noindent\emph{Ukrajinsky / Ukrainian:} лінійні форми\par
\smallskip\noindent\emph{Razlog:} Linearne formy v §2 raboty 30 sut izrazy $r_1\xi_1+\cdots+r_k\xi_k$ s koeficientami iz prstena. Na tym osnovana jest definicija konečno porodženoj modulnoj oblasti; ona ne zahtěvaje, že porodžajuče elementy sut nezavisne. Tekst takože razlikuje slučaj prstena s jedinicoju od slučaja bez njej, kde prihodet dodatne cěločislene kratnosti. Arhivny razlog čuva prosty paralelny naziv v trěh jezykah. Ta podobnost ne směje voditi k zaměně modula vektornym prostorom niti k dodanju nezavisnosti, ktorej definicija ne ima.\par
\noindent\emph{Reason:} In §2 of Paper 30, linearne formy are expressions $r_1\xi_1+\cdots+r_k\xi_k$ with coefficients in the ring. They define a finitely generated module domain; the generators are not required to be independent. The text also distinguishes a ring with identity from one without it, where additional integer multiples occur. The archived reason retains a simple parallel name in the three languages. That resemblance must not lead to replacing the module by a vector space or adding independence absent from the definition.\par
\noindent{\footnotesize Izvorne zapisy / Original records: HIST-g109-005.}\par
\smallskip\noindent\emph{Sravnjenje s tekstom:} §2 raboty 30 definovaje konečnu modulnu oblast prěz vse linearne kombinacije konečno mnogyh elementov. Nezavisnost ne jest zahtěvana, a prsten bez jedinicy potrěbuje dodatne členy.\par
\noindent\emph{Textual comparison:} §2 of Paper 30 defines a finite module domain through the linear span of finitely many elements. Independence is not required, and a ring without identity needs additional terms.\par
\noindent{\footnotesize TEXT-COMP-230; pasaže / passages: 16309, 16311; \texttt{CORPUS\_FORM\_COMPARISONS.json}.}\par
\end{minipage}\par
\par\medskip\Needspace{7\baselineskip}
\noindent\begin{minipage}{\linewidth}
\hypertarget{TERM-1116}{}
\noindent\textbf{modul linearnyh form}\quad{\footnotesize\texttt{TERM-1116}}\par
\noindent\emph{Iztočnik / Source:} Linearformenmodul\par
\noindent\emph{Rusky / Russian:} модуль линейных форм\par
\noindent\emph{Ukrajinsky / Ukrainian:} модуль лінійних форм\par
\smallskip\noindent\emph{Razlog:} Modul linearnyh form razkriva strukturu objekta v imennom sočetaniu, vměsto stiskanja vsego německogo složenogo slova v jedin pridavnik. Rabota 18 razgleduje polinomny ideal kako modul, linearny v stepenovyh produktah vybranyh prěměnnyh; ostale prěměnne imajut ulogu pri oblasti koeficientov. To ne znači, že polinomy sut stepena jedin po početkovyh prěměnnyh. Zachovany razlog daje prědnost čitateljnosti te struktury. Zato mora byti jasno, v čem formy sut linearne i nad čem modul jest razgledany, osoblivo pri prěhodu k normam.\par
\noindent\emph{Reason:} Modul linearnyh form exposes the object’s structure in a noun phrase instead of compressing the German compound into one derived adjective. Paper 18 regards a polynomial ideal as a module linear in monomials of selected variables; the remaining variables have a role in the coefficient domain. This does not make the polynomials degree one in the original variables. The preserved reason prefers readability of that structure. The explanation must therefore say in what the forms are linear and over what domain the module is considered, particularly when passing to norms.\par
\noindent{\footnotesize Izvorne zapisy / Original records: HIST-g080-005.}\par
\smallskip\noindent\emph{Sravnjenje s tekstom:} Raboty 18 i 22 razgledujut stepenove produkty vybranyh prěměnnyh kako linearne argumenty. Polinomy v ostalyh prěměnnyh služat kako koeficienty; to ne čini početkove polinomy linearnymi po vsih prěměnnyh.\par
\noindent\emph{Textual comparison:} Papers 18 and 22 treat monomials in selected variables as linear arguments. Polynomials in the remaining variables serve as coefficients; this does not make the original polynomials linear in all variables.\par
\noindent{\footnotesize TEXT-COMP-231; pasaže / passages: 12353, 12354, 12355, 12356, 14501, 14617, 14619, 14621, 14626; \texttt{CORPUS\_FORM\_COMPARISONS.json}.}\par
\end{minipage}\par
\par\medskip\Needspace{7\baselineskip}
\noindent\begin{minipage}{\linewidth}
\hypertarget{TERM-1117}{}
\noindent\textbf{modul linearnyh form / modul iz linearnyh form}\quad{\footnotesize\texttt{TERM-1117}}\par
\noindent\emph{Iztočnik / Source:} Linearformenmodul / Modul aus Linearformen\par
\noindent\emph{Rusky / Russian:} модуль линейных форм / модуль из линейных форм\par
\noindent\emph{Ukrajinsky / Ukrainian:} модуль лінійних форм / модуль з лінійних форм\par
\smallskip\noindent\emph{Razlog:} Modul linearnyh form i modul iz linearnyh form sut dve konstrukcije za toj samy strukturny opis. Rěč iz podčrkivaje, iz kakovih elementov modul sostoji; ona ne tvori novy vid modula. V rabotě 22 podnožna primětka izrično opravdyvaje naziv ideal za polinomne objekty, da by se oni razlikovali od modulov linearnyh form. Potom se polinom piše kako suma koeficientov pomnoženyh s monomami, tretovanymi linearno. Arhivny razlog čuva to različenje. Bez ukazanja koeficientov i monomov formalna podobnost nazvanij by skryla proměnu točky viděnja.\par
\noindent\emph{Reason:} Modul linearnyh form and modul iz linearnyh form are two constructions for the same structural description. Iz emphasizes what elements the module consists of; it does not create a new type of module. In Paper 22 a footnote explicitly justifies ideal for the polynomial objects to distinguish them from modules of linear forms. A polynomial is then written as a sum of coefficients times monomials treated linearly. The archived reason preserves that distinction. Without identifying the coefficients and monomials, the similarity of the names would obscure the change of viewpoint.\par
\noindent{\footnotesize Izvorne zapisy / Original records: HIST-g093-006.}\par
\smallskip\noindent\emph{Sravnjenje s tekstom:} Raboty 18 i 22 razgledujut stepenove produkty vybranyh prěměnnyh kako linearne argumenty. Polinomy v ostalyh prěměnnyh služat kako koeficienty; to ne čini početkove polinomy linearnymi po vsih prěměnnyh.\par
\noindent\emph{Textual comparison:} Papers 18 and 22 treat monomials in selected variables as linear arguments. Polynomials in the remaining variables serve as coefficients; this does not make the original polynomials linear in all variables.\par
\noindent{\footnotesize TEXT-COMP-231; pasaže / passages: 12353, 12354, 12355, 12356, 14501, 14617, 14619, 14621, 14626; \texttt{CORPUS\_FORM\_COMPARISONS.json}.}\par
\end{minipage}\par
\par\medskip\Needspace{7\baselineskip}
\noindent\begin{minipage}{\linewidth}
\hypertarget{TERM-1118}{}
\noindent\textbf{linearno-drobno}\quad{\footnotesize\texttt{TERM-1118}}\par
\noindent\emph{Iztočnik / Source:} lineargebrochen\par
\noindent\emph{Rusky / Russian:} линейно-дробно\par
\noindent\emph{Ukrajinsky / Ukrainian:} лінійно-дробово\par
\smallskip\noindent\emph{Razlog:} Linearno-drobno jest prilovičny oblik za zavisnost od funkcije prěz kvocient linearnyh izrazov. V §13 raboty 6 toj opis ne jest koncem argumenta: za razgledane polinomy zavisnost jest takože algebraično cěla, i tekst iz togo vyvodi, že ona jest v istině cěla i linearna. Izpuščenje togo poslednjego koraka by utratilo uslovje, ktoro daje integralnu bazu iz jednogo polinoma. Kratky arhivny razlog govori o sovrěmennom nazivu, ale ne opravdyvaje zaměnu vsego argumenta samym nazvom Möbiusova transformacija.\par
\noindent\emph{Reason:} Linearno-drobno is the adverbial form for dependence through a quotient of linear expressions. In §13 of Paper 6 that description is not the end of the argument: for the polynomials under consideration, the dependence is also integral, from which the text concludes that it is actually polynomial and linear. Omitting that last step would lose the condition yielding an integral basis consisting of one polynomial. The brief archived reason invokes modern terminology, but does not justify replacing the argument by the label Möbius transformation.\par
\noindent{\footnotesize Izvorne zapisy / Original records: HIST-g053-007.}\par
\smallskip\noindent\emph{Sravnjenje s tekstom:} §13 raboty 6 priměnjuje dodatno uslovje integralnosti k linearno-drobnoj zavisnosti. Zaključenje o integralnoj bazi iz jednogo polinoma zavisi od togo koraka.\par
\noindent\emph{Textual comparison:} §13 of Paper 6 applies an additional integrality condition to fractional-linear dependence. The conclusion that one polynomial forms an integral basis depends on that step.\par
\noindent{\footnotesize TEXT-COMP-232; pasaže / passages: 6402, 6404, 6406, 6408; \texttt{CORPUS\_FORM\_COMPARISONS.json}.}\par
\end{minipage}\par
\par\medskip\Needspace{7\baselineskip}
\noindent\begin{minipage}{\linewidth}
\hypertarget{TERM-1119}{}
\noindent\textbf{drobno-linearne preobraženje}\quad{\footnotesize\texttt{TERM-1119}}\par
\noindent\emph{Iztočnik / Source:} lineargebrochene Transformation\par
\noindent\emph{Rusky / Russian:} дробно-линейное преобразование\par
\noindent\emph{Ukrajinsky / Ukrainian:} дробово-лінійне перетворення\par
\smallskip\noindent\emph{Razlog:} Drobno-linearne preobraženje opisuje nejedinstvenost izbora Lurotovej funkcije v rabotě 5: razne dopustime porodžajuče funkcije svęzany sut takym prěobraženjem. Obě česti pridavnika morajut ostati, zato že samo linearne by izključilo děljenje. Forma jest blizka zabilěženym ruskomu i ukrajinskomu sravnjenjam. Rabota 6 ima obratny porjadok korenjev i internacionalno transformacija; to jest dokumentovana terminologična variacija, ne dokaz različnogo vida prěobraženja. Arhivny razlog ne daje čislovyh težin za izbor porjadka korenjev.\par
\noindent\emph{Reason:} Drobno-linearne preobraženje describes the nonuniqueness of the Lüroth function in Paper 5: different allowed generators are related by such a transformation. Both parts of the adjective matter, because linear alone would exclude division. The form is close to the recorded Russian and Ukrainian comparanda. Paper 6 reverses the stem order and uses international transformacija; this is documented terminological variation, not evidence of a different transformation. The archived explanation supplies no numerical weights for the choice of stem order.\par
\noindent{\footnotesize Izvorne zapisy / Original records: HIST-g039-011.}\par
\smallskip\noindent\emph{Sravnjenje s tekstom:} Rabota 5 piše drobno-linearne prěobraženje, a rabota 6 linearno-drobna transformacija. Obe frazy govoret o nejedinstvenosti Lurotovej funkcije pri tom samom porodženom polju.\par
\noindent\emph{Textual comparison:} Paper 5 uses drobno-linearne prěobraženje, while Paper 6 uses linearno-drobna transformacija. Both concern the nonuniqueness of the Lüroth function for the same generated field.\par
\noindent{\footnotesize TEXT-COMP-233; pasaže / passages: 4958, 5696, 5704, 5708; \texttt{CORPUS\_FORM\_COMPARISONS.json}.}\par
\end{minipage}\par
\par\medskip\Needspace{7\baselineskip}
\noindent\begin{minipage}{\linewidth}
\hypertarget{TERM-1120}{}
\noindent\textbf{linearno-drobna transformacija}\quad{\footnotesize\texttt{TERM-1120}}\par
\noindent\emph{Iztočnik / Source:} lineargebrochene Transformation\par
\noindent\emph{Rusky / Russian:} дробно-линейное преобразование\par
\noindent\emph{Ukrajinsky / Ukrainian:} дробово-лінійна трансформація\par
\smallskip\noindent\emph{Razlog:} Linearno-drobna transformacija v §6 raboty 6 čuva smysl německogo složenogo termina bez zaměny širšoju projektivnoju terminologijeju. To jest prěobraženje tipa Möbius, ale matematično tvrđenje tu govori o funkciji, jednoznačnoj do takogo prěobraženja, i o ravnosti porodžajemyh polj. Arhivny razlog daje prědnost kratkomu opisu. Drobno-linearne preobraženje iz raboty 5 i linearno-drobna transformacija ne predstavljajut dva ustanovjene termina za dva različna objekta. Slovnik čuva obe formy, dokolě objašnjaje obči pojem i dokumentovany rozděl v pisanju.\par
\noindent\emph{Reason:} In §6 of Paper 6, linearno-drobna transformacija preserves the German compound’s meaning without replacing it by broader projective terminology. It is a Möbius-type transformation, but the assertion here concerns a function unique up to that transformation and equality of the generated fields. The archived reason prefers a concise description. Drobno-linearne preobraženje in Paper 5 and linearno-drobna transformacija are not two established terms for different objects. The dictionary preserves both forms while explaining their common concept and documented difference in wording.\par
\noindent{\footnotesize Izvorne zapisy / Original records: HIST-g046-005.}\par
\smallskip\noindent\emph{Sravnjenje s tekstom:} Rabota 5 piše drobno-linearne prěobraženje, a rabota 6 linearno-drobna transformacija. Obe frazy govoret o nejedinstvenosti Lurotovej funkcije pri tom samom porodženom polju.\par
\noindent\emph{Textual comparison:} Paper 5 uses drobno-linearne prěobraženje, while Paper 6 uses linearno-drobna transformacija. Both concern the nonuniqueness of the Lüroth function for the same generated field.\par
\noindent{\footnotesize TEXT-COMP-233; pasaže / passages: 4958, 5696, 5704, 5708; \texttt{CORPUS\_FORM\_COMPARISONS.json}.}\par
\end{minipage}\par
\par\medskip\Needspace{7\baselineskip}
\noindent\begin{minipage}{\linewidth}
\hypertarget{TERM-1121}{}
\noindent\textbf{linearna nezavisnost $(\alpha_i)$}\quad{\footnotesize\texttt{TERM-1121}}\par
\noindent\emph{Iztočnik / Source:} Linearunabhängigkeit der $(\alpha_i)$\par
\noindent\emph{Rusky / Russian:} линейная независимость $(\alpha_i)$\par
\noindent\emph{Ukrajinsky / Ukrainian:} лінійна незалежність $(\alpha_i)$\par
\smallskip\noindent\emph{Razlog:} Linearna nezavisnost tu odnosi se k klasam ostatkov $(\alpha_i)$, ne samo k početkovym elementam $\alpha_i$. V §7 raboty 31 prědhodny dokaz pokazuje, že kongruencija linearnoj kombinacije modulo $p$ prisiljaje vsaky koeficient byti dělivym prěz $p$. Zato klasy tvorjat bazu nad poljem ostatkov. Arhivny razlog čuva svezu s tym dokazom; slědujuči prěhod diskriminanta opiraje se na nju. Kratka fraza blizka ruskomu i ukrajinskomu sravnjenju, ale sama podobnost slov ne zaměnjuje ukazanja oblasti koeficientov i značenja skobok.\par
\noindent\emph{Reason:} Here linearna nezavisnost concerns the residue classes $(\alpha_i)$, not merely the original elements $\alpha_i$. The preceding proof in §7 of Paper 31 shows that a linear combination congruent to zero modulo $p$ forces every coefficient to be divisible by $p$. The classes therefore form a basis over the residue field. The recorded reason preserves the connection with this proof, on which the subsequent passage of the discriminant depends. The short phrase resembles the Russian and Ukrainian comparanda, but verbal similarity cannot replace the coefficient domain or the meaning of the parentheses.\par
\noindent{\footnotesize Izvorne zapisy / Original records: HIST-g152-003.}\par
\smallskip\noindent\emph{Sravnjenje s tekstom:} V §7 raboty 31 skobki označajut klasy ostatkov. Nezavisnost nad poljem ostatkov slěduje iz dělivosti vseh koeficientov prěz $p$; tek potom po toj bazi tvori se diskriminant.\par
\noindent\emph{Textual comparison:} In §7 of Paper 31, parentheses mark residue classes. Independence over the residue field follows from divisibility of every coefficient by $p$; the discriminant is then formed using that basis.\par
\noindent{\footnotesize TEXT-COMP-234; pasaže / passages: 18627, 18631, 18633, 18635, 18637, 18671; \texttt{CORPUS\_FORM\_COMPARISONS.json}.}\par
\end{minipage}\par
\par\medskip\Needspace{7\baselineskip}
\noindent\begin{minipage}{\linewidth}
\hypertarget{TERM-1122}{}
\noindent\textbf{lěvy dělitelj nulja}\quad{\footnotesize\texttt{TERM-1122}}\par
\noindent\emph{Iztočnik / Source:} linker Nullteiler\par
\noindent\emph{Rusky / Russian:} левый делитель нуля\par
\noindent\emph{Ukrajinsky / Ukrainian:} лівий дільник нуля\par
\smallskip\noindent\emph{Razlog:} Lěvy dělitelj nulja jest doslovno složeny naziv s domačim korenjem děliti, podobno ruskomu i ukrajinskomu sravnjenju. O smyslu lěvy odlučuje formula: v §7 raboty 34 $P$ stoji lěvo v $PA=0$. Obstajanje lěvoj obratnoj matrice dava $A=P^{-1}PA=0$, zato $P$ ne može anulirati nenulovu matricu s toj strany. Strana ne směje byti izpuščena niti zaměnjena po inoj konvenciji nazvanja. Arhivny zapis naziva termin standardnym, ale ne dava odělne slovnične potvrđenje; izričny dokaz jest tu točna opora za matematičny smysl.\par
\noindent\emph{Reason:} Lěvy dělitelj nulja is a compositional calque using the native division stem, parallel to the Russian and Ukrainian comparanda. The formula determines the sense of left: in §7 of Paper 34, $P$ stands on the left in $PA=0$. A left inverse gives $A=P^{-1}PA=0$, so $P$ cannot annihilate a nonzero matrix on that side. The side must neither be omitted nor silently reversed to follow another naming convention. The archived record calls the term standard but supplies no separate dictionary attestation; the explicit proof establishes its mathematical sense here.\par
\noindent{\footnotesize Izvorne zapisy / Original records: HIST-g185-004.}\par
\smallskip\noindent\emph{Sravnjenje s tekstom:} V §7 raboty 34 formula $PA=0$ pokazuje, že lěvy odnosi se k městu $P$ v produktu. Lěva obratna matrica poništaje tuty produkt i dava $A=0$.\par
\noindent\emph{Textual comparison:} In §7 of Paper 34, $PA=0$ fixes the convention: left refers to the position of $P$ in the product. Multiplication by a left inverse gives $A=0$.\par
\noindent{\footnotesize TEXT-COMP-235; pasaže / passages: 20329, 20331, 20335, 20347, 20348, 20350; \texttt{CORPUS\_FORM\_COMPARISONS.json}.}\par
\end{minipage}\par
\par\medskip\Needspace{7\baselineskip}
\noindent\begin{minipage}{\linewidth}
\hypertarget{TERM-1124}{}
\noindent\textbf{lěva jedinica}\quad{\footnotesize\texttt{TERM-1124}}\par
\noindent\emph{Iztočnik / Source:} Linkseinheit\par
\noindent\emph{Rusky / Russian:} левая единица\par
\noindent\emph{Ukrajinsky / Ukrainian:} ліва одиниця\par
\smallskip\noindent\emph{Razlog:} Lěva jedinica jest element, ktory pri množenju s lěva ostavlja elementy razgledanogo prstena nepoměnjenymi. V §9 raboty 34 ona nastavaje kako suma odgovarjajučih idempotentov i služi prstenu, sostavljenomu iz pravyh idealov. Naziv zato mora čuvati i stranu množenja, i točny prsten, za ktory jedinica dějstvuje. Arhivny razlog svezuje frazu s ranějšim nazvom prstena s jedinicoju; rusky i ukrajinsky sravnitelje imajut tu samu strukturu. Ne tvrdi se, že lěva jedinica už jest dvustranna ili že svaki idempotent jest jedinica cělogo prstena.\par
\noindent\emph{Reason:} A lěva jedinica leaves elements of the ring under consideration unchanged when multiplied on the left. In §9 of Paper 34 it arises as a sum of suitable idempotents and serves the ring formed from the indicated right ideals. The name must therefore retain both the multiplication side and the ring on which the identity acts. The recorded reason connects the phrase with the preceding ring-with-identity terminology; Russian and Ukrainian have the same construction. It does not assert that a left identity is already two-sided or that every idempotent is the identity of the whole ring.\par
\noindent{\footnotesize Izvorne zapisy / Original records: HIST-g187-006.}\par
\smallskip\noindent\emph{Sravnjenje s tekstom:} §9 raboty 34 dava lěvu jedinicu na sumě pravyh idealov. §13 najprvo dobiva lěvu jedinicu za cěly prsten i potom izključenjem nilpotentnogo ostatka dokazuje pravo dějstvovanje. Te koraky ne sut zaměnime.\par
\noindent\emph{Textual comparison:} Section 9 of Paper 34 gives a left identity on the sum of right ideals. Section 13 first obtains a left identity for the whole ring and then proves its right-identity action by ruling out a nilpotent remainder. These steps are not interchangeable.\par
\noindent{\footnotesize TEXT-COMP-237; pasaže / passages: 20509, 20513, 20551, 20962, 20964, 20979, 20981, 20985, 20987, 20989; \texttt{CORPUS\_FORM\_COMPARISONS.json}.}\par
\end{minipage}\par
\par\medskip\Needspace{7\baselineskip}
\noindent\begin{minipage}{\linewidth}
\hypertarget{TERM-1125}{}
\noindent\textbf{lěva jedinica / prava jedinica}\quad{\footnotesize\texttt{TERM-1125}}\par
\noindent\emph{Iztočnik / Source:} Linkseinheit / Rechtseinheit\par
\noindent\emph{Rusky / Russian:} левая единица / правая единица\par
\noindent\emph{Ukrajinsky / Ukrainian:} ліва одиниця / права одиниця\par
\smallskip\noindent\emph{Razlog:} Par lěva jedinica / prava jedinica čuva dva koraka dokaza v §13 raboty 34. Najprvo se konstruuje lěva jedinica. Potom Peirceovo razloženje daje dodatny lěvy ideal s kvadratom nula; pod danym prědpoloženjem toj ideal mora izčeznuti, i jedinica jest takože prava. Arhivny razlog daje prědnost konkretnomu nazvanju identičnogo dějstvovanja s každoj strany. Tut ne ide o libovoljny obratimy element: ravnosti $ex=x$ i $xe=x$ sut različne uslovja, čije sjedinenje jest rezultat dokaza.\par
\noindent\emph{Reason:} The pair lěva jedinica / prava jedinica preserves two steps in the proof in §13 of Paper 34. A left identity is first constructed. A Peirce decomposition then gives a remaining left ideal whose square is zero; under the stated hypothesis this ideal vanishes, and the identity is also a right identity. The recorded preference is for concrete names for identity action on each side. This is not an arbitrary invertible element: $ex=x$ and $xe=x$ are distinct conditions, and proving both is the point of the argument.\par
\noindent{\footnotesize Izvorne zapisy / Original records: HIST-g191-004.}\par
\smallskip\noindent\emph{Sravnjenje s tekstom:} §9 raboty 34 dava lěvu jedinicu na sumě pravyh idealov. §13 najprvo dobiva lěvu jedinicu za cěly prsten i potom izključenjem nilpotentnogo ostatka dokazuje pravo dějstvovanje. Te koraky ne sut zaměnime.\par
\noindent\emph{Textual comparison:} Section 9 of Paper 34 gives a left identity on the sum of right ideals. Section 13 first obtains a left identity for the whole ring and then proves its right-identity action by ruling out a nilpotent remainder. These steps are not interchangeable.\par
\noindent{\footnotesize TEXT-COMP-237; pasaže / passages: 20509, 20513, 20551, 20962, 20964, 20979, 20981, 20985, 20987, 20989; \texttt{CORPUS\_FORM\_COMPARISONS.json}.}\par
\end{minipage}\par
\par\medskip\Needspace{7\baselineskip}
\noindent\begin{minipage}{\linewidth}
\hypertarget{TERM-1126}{}
\noindent\textbf{lěvy ideal / pravy ideal / recipročny}\quad{\footnotesize\texttt{TERM-1126}}\par
\noindent\emph{Iztočnik / Source:} Linksideal / Rechtsideal / reziprok\par
\noindent\emph{Rusky / Russian:} левый идеал / правый идеал / взаимно обратный\par
\noindent\emph{Ukrajinsky / Ukrainian:} лівий ідеал / правий ідеал / взаємно обернений\par
\smallskip\noindent\emph{Razlog:} Lěvy ideal i pravy ideal govoret, s ktorej strany dopuščeno jest množenje elementami poredka. Recipročny v teoremu 2 raboty 42 ima užši smysl než neoprěděljeno vzajemno odnošenje: idealy $\mathcal L$ i $\mathcal R$ sut utvorjene iz modula $a$ i jego slědovo komplementarnogo modula $\bar a$. Arhivny razlog zadržuje internacionalny naziv, dokolě rusky i ukrajinsky sravnitelje upotrěbjajut frazy za vzajemnu obratnost. To ne jest děljenje na ideal kako na čislo; čitatelj mora viděti i maksimalne poredky, i komplementarnost, ktora obosnovuje to sparjenje.\par
\noindent\emph{Reason:} Lěvy ideal and pravy ideal specify the side on which multiplication by elements of the order is allowed. In Theorem 2 of Paper 42, recipročny has a more precise meaning than an unspecified mutual relation: the ideals $\mathcal L$ and $\mathcal R$ are constructed from a module $a$ and its trace-complementary module $\bar a$. The recorded reason retains the international term, while the Russian and Ukrainian comparanda use phrases for mutual inverseness. This is not division by an ideal as though it were a number; the orders and the complementarity underlying the pairing must remain explicit.\par
\noindent{\footnotesize Izvorne zapisy / Original records: HIST-g212-011.}\par
\smallskip\noindent\emph{Sravnjenje s tekstom:} Rabota 34 izrično definira stranu množenja bez prědpoloženja komutativnosti. Teorem 2 raboty 42 konstruuje $\mathcal L,\mathcal R$ iz $a,\bar a$; slědova komplementarnost dava matematično sodržanje recipročnoj svezi.\par
\noindent\emph{Textual comparison:} Paper 34 explicitly defines the side of multiplication without assuming commutativity. Theorem 2 of Paper 42 constructs $\mathcal L,\mathcal R$ from $a,\bar a$; trace complementarity gives the reciprocal relation its mathematical content.\par
\noindent{\footnotesize TEXT-COMP-238; pasaže / passages: 19821, 19833, 19835, 24555, 24560, 24569, 24571, 24573, 24575, 24580; \texttt{CORPUS\_FORM\_COMPARISONS.json}.}\par
\end{minipage}\par
\par\medskip\Needspace{7\baselineskip}
\noindent\begin{minipage}{\linewidth}
\hypertarget{TERM-1127}{}
\noindent\textbf{lěve idealy / pravy ideal}\quad{\footnotesize\texttt{TERM-1127}}\par
\noindent\emph{Iztočnik / Source:} Linksideale / Rechtsideal\par
\noindent\emph{Rusky / Russian:} левые идеалы / правый идеал\par
\noindent\emph{Ukrajinsky / Ukrainian:} ліві ідеали / правий ідеал\par
\smallskip\noindent\emph{Razlog:} Arhivny zapis dlja §23 raboty 34 upotrěbjaje obyčne nazvy lěvyh i pravyh idealov. Njih matematičny sadržaj dan jest ranějšim oprěděljenjem: $r\mathfrak a\subseteq\mathfrak a$ dlja lěvogo ideala, $\mathfrak a r\subseteq\mathfrak a$ dlja pravogo. Komutativnost ne jest prědpoložena. Razlika medžu množnym lěve idealy i jedinstvenym pravy ideal v zaglavju ne označa razliku v definiciji. Razlog o stabilnosti nazvov podkrepljen jest zapisanymi ruskimi i ukrajinskimi formami, ne dokazom o vsem slavjanskom upotrěbjenju.\par
\noindent\emph{Reason:} The archived entry associated with §23 of Paper 34 uses ordinary left- and right-ideal names. Their mathematical content is supplied by the earlier definition: $r\mathfrak a\subseteq\mathfrak a$ for a left ideal and $\mathfrak a r\subseteq\mathfrak a$ for a right ideal. Commutativity is not assumed. The plural left ideals and singular right ideal in the heading do not signal different kinds of definition. The claim that the names are stable is supported here by recorded Russian and Ukrainian forms, not by a survey of all Slavic usage.\par
\noindent{\footnotesize Izvorne zapisy / Original records: HIST-g201-007.}\par
\smallskip\noindent\emph{Sravnjenje s tekstom:} Rabota 34 izrično definira stranu množenja bez prědpoloženja komutativnosti. Teorem 2 raboty 42 konstruuje $\mathcal L,\mathcal R$ iz $a,\bar a$; slědova komplementarnost dava matematično sodržanje recipročnoj svezi.\par
\noindent\emph{Textual comparison:} Paper 34 explicitly defines the side of multiplication without assuming commutativity. Theorem 2 of Paper 42 constructs $\mathcal L,\mathcal R$ from $a,\bar a$; trace complementarity gives the reciprocal relation its mathematical content.\par
\noindent{\footnotesize TEXT-COMP-238; pasaže / passages: 19821, 19833, 19835, 24555, 24560, 24569, 24571, 24573, 24575, 24580; \texttt{CORPUS\_FORM\_COMPARISONS.json}.}\par
\end{minipage}\par
\par\medskip\Needspace{7\baselineskip}
\noindent\begin{minipage}{\linewidth}
\hypertarget{TERM-1128}{}
\noindent\textbf{dlžina $i$}\quad{\footnotesize\texttt{TERM-1128}}\par
\noindent\emph{Iztočnik / Source:} Länge $i$\par
\noindent\emph{Rusky / Russian:} длина $i$\par
\noindent\emph{Ukrajinsky / Ukrainian:} довжина $i$\par
\smallskip\noindent\emph{Razlog:} Dlžina $i$ v §2 raboty 30 jest položaj elementa v postupnom sistemu porodžajučih elementov $\xi_1,\ldots,\xi_k$. Element možno izraziti prěz prve $i$ elementov, ale ne prěz prve $i-1$. To ne jest geometrična dlžina niti bez dodatnogo dokaza dlžina kompozicijnogo reda modula. Arhivny razlog izbira prěmy naziv za to postupno razširjanje. Poredok porodžajučih elementov jest čest definicije; oni ne muset tvoriti nezavisnu bazu. Zato tu ne dostatočno govoriti samo o čislu koeficientov v jednom izbranom zapisu.\par
\noindent\emph{Reason:} In §2 of Paper 30, dlžina $i$ gives an element's position in the successive spans of the ordered generators $\xi_1,\ldots,\xi_k$. It can be expressed using the first $i$ generators, but not the first $i-1$. This is neither geometric length nor, without further argument, the composition length of a module. The recorded reason chooses a direct name for this successive construction. The ordering belongs to the definition, and the generators need not form an independent basis. Counting the coefficients in one chosen expression would therefore be insufficient.\par
\noindent{\footnotesize Izvorne zapisy / Original records: HIST-g109-007.}\par
\smallskip\noindent\emph{Sravnjenje s tekstom:} §2 raboty 30 definira dlžinu prěz najmanjši početkovy spisok $\xi_1,\ldots,\xi_i$, dostatočny dlja izraženja elementa. Definicija ne požada nezavisnost generatorov i ne jest definicija kompozicijnoj dlžiny.\par
\noindent\emph{Textual comparison:} Section 2 of Paper 30 defines length through the shortest initial list $\xi_1,\ldots,\xi_i$ sufficient to express the element. The definition does not require independent generators and is not the definition of composition length.\par
\noindent{\footnotesize TEXT-COMP-239; pasaže / passages: 16311, 16314, 16316, 16318, 16320; \texttt{CORPUS\_FORM\_COMPARISONS.json}.}\par
\end{minipage}\par
\par\medskip\Needspace{7\baselineskip}
\noindent\begin{minipage}{\linewidth}
\hypertarget{TERM-1129}{}
\noindent\textbf{sistemy rěšenj}\quad{\footnotesize\texttt{TERM-1129}}\par
\noindent\emph{Iztočnik / Source:} Lösungssysteme\par
\noindent\emph{Rusky / Russian:} системы решений\par
\noindent\emph{Ukrajinsky / Ukrainian:} системи розв'язків\par
\smallskip\noindent\emph{Razlog:} Sistemy rěšenj čuva istorijny način govorenja o cělyh redah cěnnostij neizvěstnyh. V §5 raboty 6 rěšenje jest red $\xi$, ktory izpolnjaje relacije (7) i uslovje $H^*(\xi)\ne0$. Čislo takyh dopustimyh redov svezano jest so stepenom polja. Arhivny razlog zadržuje sistem, da by ne podměniti red jedinoju koordinatoju. To ne znači, že vsaka moderna fraza rěšenje sistema jest nepravilna; važno jest, da ona označa cěly red i čuva izključenie nulovogo imenovatelja.\par
\noindent\emph{Reason:} Sistemy rěšenj retains the historical way of referring to complete tuples of values. In §5 of Paper 6, a solution is a tuple $\xi$ satisfying relations (7) and $H^*(\xi)\ne0$. The number of admissible tuples is related to the field degree. The recorded reason retains system to avoid replacing a tuple by an individual coordinate. This does not make the modern expression solution of the system wrong: what matters is that it denotes the whole tuple and preserves the exclusion of a zero denominator.\par
\noindent{\footnotesize Izvorne zapisy / Original records: HIST-g045-010.}\par
\smallskip\noindent\emph{Sravnjenje s tekstom:} Relacije (7) raboty 6 govoret o cělom redu $\xi$ s $H^*(\xi)\ne0$. Geometrično tolkovanje i tvrdnja o čislu sistemov čuvajut to ograničenje.\par
\noindent\emph{Textual comparison:} Relations (7) in Paper 6 concern the whole tuple $\xi$ with $H^*(\xi)\ne0$. Both the geometric interpretation and the count of solution systems retain that restriction.\par
\noindent{\footnotesize TEXT-COMP-240; pasaže / passages: 5628, 5631, 5632, 5636, 5646, 5651, 5653, 5654, 5656, 5658, 5662; \texttt{CORPUS\_FORM\_COMPARISONS.json}.}\par
\end{minipage}\par
\par\medskip\Needspace{7\baselineskip}
\noindent\begin{minipage}{\linewidth}
\hypertarget{TERM-1130}{}
\noindent\textbf{Lürothova funkcija}\quad{\footnotesize\texttt{TERM-1130}}\par
\noindent\emph{Iztočnik / Source:} Lürothsche Funktion\par
\noindent\emph{Rusky / Russian:} функция Люрота\par
\noindent\emph{Ukrajinsky / Ukrainian:} функція Люрота\par
\smallskip\noindent\emph{Razlog:} Lürothova funkcija čuva ime, čerez ktore možno najdti teoremu i bibliografiju. Dva arhivne zapisy svezujut §§6 i 13 raboty 6 i opisujut različno pisanje v latinici i kirilici. V nynějšem latiničnom tekstu stoi Lurotova funkcija, dokolě točna bibliografična citacija čuva Lüroth. Tuto razliku slovnik ne skryvaje. Matematično ide o funkciju, ktora porodžaje razgledano polje i jest oprěděljena do linearno-drobnoj transformacije; v §13 dodatne uslovja dajut polinomialnu integralnu bazu. Arhivno objašnjenje kiriličnogo pisanja ne jest nova prověrka vseh kiriličnyh pojavjenij.\par
\noindent\emph{Reason:} Lürothova funkcija preserves the name by which readers can find the theorem and bibliography. Two archived reasons connect §§6 and 13 of Paper 6 and describe different Latin and Cyrillic spellings. The current Latin prose reads Lurotova funkcija, while the exact bibliographical citation retains Lüroth; the dictionary makes that distinction visible. Mathematically, the function generates the field under discussion and is determined up to a fractional-linear transformation. In §13 further hypotheses yield a polynomial integral basis. The archived account of Cyrillic spelling is not a fresh verification of every Cyrillic occurrence.\par
\noindent{\footnotesize Izvorne zapisy / Original records: HIST-g046-004, HIST-g053-006.}\par
\smallskip\noindent\emph{Sravnjenje s tekstom:} §6 raboty 6 piše Lurotova funkcija v prozě, ale Lüroth v citaciji. §13 upotrěbjaje to samo funkcionalno nazvanje pri dodatnom argumentu o integralnosti. Arhivno Lürothova i nynějšno Lurotova ne smějut byti nerazličeno prikazane.\par
\noindent\emph{Textual comparison:} Section 6 of Paper 6 uses Lurotova funkcija in prose but Lüroth in the citation. Section 13 uses the same functional name with an additional integrality argument. The archived Lürothova and current Lurotova should not be presented as identical spellings.\par
\noindent{\footnotesize TEXT-COMP-241; pasaže / passages: 5694, 5696, 5702, 6402, 6408, 6410; \texttt{CORPUS\_FORM\_COMPARISONS.json}.}\par
\end{minipage}\par
\par\medskip\Needspace{7\baselineskip}
\noindent\begin{minipage}{\linewidth}
\hypertarget{TERM-1131}{}
\noindent\textbf{definujuči ideal}\quad{\footnotesize\texttt{TERM-1131}}\par
\noindent\emph{Iztočnik / Source:} $m$ / definierendes Ideal (implicit in scan context)\par
\noindent\emph{Rusky / Russian:} определяющий идеал\par
\noindent\emph{Ukrajinsky / Ukrainian:} визначальний ідеал\par
\smallskip\noindent\emph{Razlog:} Definujuči ideal opisuje ideal relacij, prěz ktory polinomialny prsten dava potrěbny prsten klasov ostatkov. V §22 raboty 34 $\mathfrak m$ jest porodženy relacijami $x_i^{h_i}-e$, i faktorprsten prědstavjaje grupovy prsten. Arhivny zapis sam govori, že slovny opis byl izvedeny iz konteksta simbola $m$; to ne jest dokaz, že německy tekst tu upotrěbil ravno te slova. Razlog dlja definujuči jest internacionalna prozračnost i očekyvana pravilnost tvorjenja slova. To jest objašnjenje predloga, ne samostojna jezykova potvrda.\par
\noindent\emph{Reason:} Definujuči ideal describes the ideal of relations through which a polynomial ring gives the required quotient ring. In §22 of Paper 34, $\mathfrak m$ is generated by the relations $x_i^{h_i}-e$, and the quotient represents the group ring. The archived source field itself says that the verbal description was inferred from the context of the symbol $m$; it does not establish that the German passage used those exact words. The stated reasons for definujuči are international transparency and expected regular word formation. They explain the proposal but do not independently attest it.\par
\noindent{\footnotesize Izvorne zapisy / Original records: HIST-g200-006.}\par
\smallskip\noindent\emph{Sravnjenje s tekstom:} §22 raboty 34 prikazuje $\mathfrak m=(x_1^{h_1}-e,\ldots,x_r^{h_r}-e)$ i odgovarjajuči faktorprsten. To podkrepljaje opisnu funkciju nazva definujuči ideal, ale ne dokazuje prisutnost ravno toj slovnoj frazy v německom svědku.\par
\noindent\emph{Textual comparison:} Section 22 of Paper 34 displays $\mathfrak m=(x_1^{h_1}-e,\ldots,x_r^{h_r}-e)$ and the associated quotient ring. This supports the descriptive function of defining ideal, but does not establish that the German witness contained that exact phrase.\par
\noindent{\footnotesize TEXT-COMP-242; pasaže / passages: 22006, 22010, 22012, 22021, 22023, 22025, 22027; \texttt{CORPUS\_FORM\_COMPARISONS.json}.}\par
\end{minipage}\par
\par\medskip\Needspace{7\baselineskip}
\noindent\begin{minipage}{\linewidth}
\hypertarget{TERM-1132}{}
\noindent\textbf{mnogovidnost}\quad{\footnotesize\texttt{TERM-1132}}\par
\noindent\emph{Iztočnik / Source:} Mannigfaltigkeit\par
\noindent\emph{Rusky / Russian:} многообразие\par
\noindent\emph{Ukrajinsky / Ukrainian:} многовид\par
\smallskip\noindent\emph{Razlog:} Mnogovidnost v rabotě 11 opisuje množstvo parametrov uravnenij s zadanoju grupoju. Tvrđenje govori, že izključenie cěnnostij, ktore ponižajut grupu, ne ponižaje toj mnogovidnosti. Arhivny zapis daje rusko mnogoobrazije i ukrajinsko mnohovyd kako sravnitelje, ale izrično ostavlja ukrajinsku i medžuslovjansku formu dlja daljšego jezykovogo razgledanja. Sam zapis ne dokazuje, že čelověčny ekspert jego prověril. V objašnjenju trěba čuvati istorijny parametrny smysl; ne možno bez argumenta dodati suvremenno uslovje gladkoj topologičnoj mnogovidnosti.\par
\noindent\emph{Reason:} In Paper 11, mnogovidnost describes the extent of a parametrized family of equations with a prescribed group. The assertion is that excluding parameter values which reduce the group does not diminish this extent. The archived record supplies Russian mnogoobrazie and Ukrainian mnohovyd as comparanda but explicitly leaves the Ukrainian and Interslavic choices open to further linguistic assessment. Its review flag does not establish that a human expert reviewed them. The explanation must retain the historical parameter-counting sense, without importing the modern requirements of a smooth topological manifold.\par
\noindent{\footnotesize Izvorne zapisy / Original records: HIST-g073-016.}\par
\smallskip\noindent\emph{Sravnjenje s tekstom:} V rabotě 11 ide o parametrnu mnogovidnost uravnenij; v rabotě 4 o uporedženy sklop rěšetkovyh toček. Jedin německy naziv zato ne potrěbuje jedino slovo bez smyslovogo različenja v prěvodu.\par
\noindent\emph{Textual comparison:} Paper 11 concerns the parameter extent of a family of equations, whereas Paper 4 concerns an ordered arrangement of lattice points. One German word therefore does not require one undifferentiated translation.\par
\noindent{\footnotesize TEXT-COMP-243; pasaže / passages: 4906, 8830, 8832, 8834; \texttt{CORPUS\_FORM\_COMPARISONS.json}.}\par
\end{minipage}\par
\par\medskip\Needspace{7\baselineskip}
\noindent\begin{minipage}{\linewidth}
\hypertarget{TERM-1133}{}
\noindent\textbf{sklop rěšetkovyh toček}\quad{\footnotesize\texttt{TERM-1133}}\par
\noindent\emph{Iztočnik / Source:} Mannigfaltigkeit von Gitterpunkten\par
\noindent\emph{Rusky / Russian:} множество решеточных точек\par
\noindent\emph{Ukrajinsky / Ukrainian:} сукупність ґраткових точок\par
\smallskip\noindent\emph{Razlog:} Sklop rěšetkovyh toček jest smyslovy prěvod Mannigfaltigkeit v drugom kontekstu: §9 raboty 4 razměšča formy po cěločislnyh koordinatah v $n-1$ pozitivnyh směrjah. To ne jest ta sama situacija kako parametrna mnogovidnost raboty 11. Arhivny razlog zato izbira opis raspoloženja toček, podobno ruskomu množestvo i ukrajinskomu sukupnist, vměsto slova, ktore by sugerovalo topologičnu mnogovidnost. Nekoliko form može iměti te same koordinaty; dodatno pravilo oprěděljaje jih poredok. Čitatelj ne směje iz nazva zaključiti, že vsaka točka označa jedinu formu.\par
\noindent\emph{Reason:} Sklop rěšetkovyh toček translates Mannigfaltigkeit by its function in a different setting: §9 of Paper 4 arranges forms by integer coordinates in $n-1$ positive directions. This is not the parameter-family setting of Paper 11. The recorded reason therefore favors an arrangement of points, alongside Russian set and Ukrainian collection terminology, instead of suggesting a topological manifold. Several forms may share the same coordinates; an additional rule orders them. The name must not imply that every point identifies only one form.\par
\noindent{\footnotesize Izvorne zapisy / Original records: HIST-g038-006.}\par
\smallskip\noindent\emph{Sravnjenje s tekstom:} V rabotě 11 ide o parametrnu mnogovidnost uravnenij; v rabotě 4 o uporedženy sklop rěšetkovyh toček. Jedin německy naziv zato ne potrěbuje jedino slovo bez smyslovogo različenja v prěvodu.\par
\noindent\emph{Textual comparison:} Paper 11 concerns the parameter extent of a family of equations, whereas Paper 4 concerns an ordered arrangement of lattice points. One German word therefore does not require one undifferentiated translation.\par
\noindent{\footnotesize TEXT-COMP-243; pasaže / passages: 4906, 8830, 8832, 8834; \texttt{CORPUS\_FORM\_COMPARISONS.json}.}\par
\end{minipage}\par
\par\medskip\Needspace{7\baselineskip}
\noindent\begin{minipage}{\linewidth}
\hypertarget{TERM-1134}{}
\noindent\textbf{matrično predstavjenje / matrične predstavjenja}\quad{\footnotesize\texttt{TERM-1134}}\par
\noindent\emph{Iztočnik / Source:} Matrizendarstellung\par
\noindent\emph{Rusky / Russian:} матричное представление / матричные представления\par
\noindent\emph{Ukrajinsky / Ukrainian:} матричне зображення / матричні зображення\par
\smallskip\noindent\emph{Razlog:} Matrično predstavjenje označa algebraično prědstavjenje, ne samo vizualny zapis v tablici. V §4 raboty 31 množenje elementom prstena na idealu, s izbranoju bazoju, dava koeficientne matrice; nuljevoj matrici odpovědajut elementy, ktore anulirajut ideal. Arhivny razlog sojedinjaje internacionalno matričny so slavjanskim predstavjenje, sravnivajuči rusko predstavlenije i ukrajinsko zobrazhennia. Nynějšni tekst piše prědstavjenje. Pri sravnjenju formul odkriva se i osobna težkost: indeksy v skalarnoj relaciji ne soglašajut se s pisanym matričnym produktom. Naziv ostaje objašnjivym, ale točny indeksny zapis potrěbuje izvornogo svědka.\par
\noindent\emph{Reason:} Matrično predstavjenje denotes an algebraic representation, not merely a visual arrangement in a table. In §4 of Paper 31, multiplication on an ideal with a chosen basis gives coefficient matrices; the zero matrix corresponds to elements annihilating the ideal. The recorded reason combines international matričny with Slavic predstavjenje, comparing Russian predstavlenie and Ukrainian zobrazhennia. The current prose spells it prědstavjenje. Comparing the formulas exposes a separate difficulty: the indices in the scalar relation do not agree with the displayed matrix product. The term can be explained, but the exact index arrangement requires the original witness.\par
\noindent{\footnotesize Izvorne zapisy / Original records: HIST-g130-001.}\par
\smallskip\noindent\emph{Sravnjenje s tekstom:} §4 raboty 31 svezuje prědstavjenje s množenjem na idealu. Nynějšny tekst dava $\gamma\alpha_i=\sum_jc_{ij}\alpha_j$, ale potom piše red $(\alpha_1,\ldots,\alpha_s)$ pomnoženy na matricu $(c_{ij})$, čija $i$-ta komponenta upotrěbjaje $c_{ji}$. To jest konkretno nesoglasje zapisa; točnu izvornu formu trěba prověriti po pečatnom svědku, ne izbrati prěz tiho transponovanje.\par
\noindent\emph{Textual comparison:} Section 4 of Paper 31 constructs a representation from multiplication on an ideal. The current text gives $\gamma\alpha_i=\sum_jc_{ij}\alpha_j$, then writes the row $(\alpha_1,\ldots,\alpha_s)$ multiplied by $(c_{ij})$, whose $i$th component instead uses $c_{ji}$. This is a concrete inconsistency in the displayed notation. The printed witness must establish the original form; a silent transpose would not do so.\par
\noindent{\footnotesize TEXT-COMP-244; pasaže / passages: 17823, 17827, 17858, 17860, 17863, 17865, 17866, 17868, 17870, 17872, 17874, 17876, 17877, 17878, 17880; \texttt{CORPUS\_FORM\_COMPARISONS.json}.}\par
\end{minipage}\par
\par\medskip\Needspace{7\baselineskip}
\noindent\begin{minipage}{\linewidth}
\hypertarget{TERM-1135}{}
\noindent\textbf{matrične jedinice}\quad{\footnotesize\texttt{TERM-1135}}\par
\noindent\emph{Iztočnik / Source:} Matrizeneinheiten\par
\noindent\emph{Rusky / Russian:} матричные единицы\par
\noindent\emph{Ukrajinsky / Ukrainian:} матричні одиниці\par
\smallskip\noindent\emph{Razlog:} Matrične jedinice sut elementy $c_{ik}$ s pravilami $c_{ij}c_{jk}=c_{ik}$ i $c_{ij}c_{hk}=0$ pri $j\ne h$. Jedinice tu ne znači, že vsaky $c_{ik}$ jest jediničny ili obratimy element cělogo prstena. Tri arhivne zapisy čuvajut obči naziv; jeden svezuje jego s ranějšim slovom jedinica, drugi tvrdi široku slavjansku prozračnost, tretji ukazuje izrične produkty. V rabotě 42 ty pravila proizhodet iz slědovo komplementarnyh baz. Pravilo množenja jest opora smysla; tvrdnja o prozračnosti ostaje kvalitativnoju ocěnkoju, ne izměrenym razuměnjem.\par
\noindent\emph{Reason:} Matrične jedinice are elements $c_{ik}$ satisfying $c_{ij}c_{jk}=c_{ik}$ and $c_{ij}c_{hk}=0$ when $j\ne h$. Units here does not mean that every $c_{ik}$ is an identity or an invertible element of the whole ring. Three preserved reasons retain the shared name: one connects it with earlier identity terminology, another claims broad Slavic transparency, and the third points to the explicit products. In Paper 42 these rules follow from trace-complementary bases. The multiplication rule establishes the sense; the transparency claim remains a qualitative judgment, not measured comprehension.\par
\noindent{\footnotesize Izvorne zapisy / Original records: HIST-g197-007, HIST-g198-004, HIST-g212-005.}\par
\smallskip\noindent\emph{Sravnjenje s tekstom:} Rabota 34 dava relacije množenja matričnyh jedinic i prědstavjaje prave idealy prěz jih redy. V rabotě 42 relacije proizhodet iz $\operatorname{Sp}(a_j\bar a_h)=\delta_{jh}$. Slovo jedinica tu ima svojstvo, različno od dvustrannoj jedinice cělogo prstena.\par
\noindent\emph{Textual comparison:} Paper 34 gives the multiplication relations for matrix units and describes right ideals using their rows. In Paper 42 the relations follow from $\operatorname{Sp}(a_j\bar a_h)=\delta_{jh}$. Unit here has a role distinct from the two-sided identity of the whole ring.\par
\noindent{\footnotesize TEXT-COMP-245; pasaže / passages: 20352, 20354, 20356, 20358, 20360, 21130, 21136, 21138, 21140, 21194, 21196, 21198, 21200, 24486, 24488, 24498, 24499, 24500, 24501, 24503, 24505; \texttt{CORPUS\_FORM\_COMPARISONS.json}.}\par
\end{minipage}\par
\par\medskip\Needspace{7\baselineskip}
\noindent\begin{minipage}{\linewidth}
\hypertarget{TERM-1136}{}
\noindent\textbf{matričny produkt}\quad{\footnotesize\texttt{TERM-1136}}\par
\noindent\emph{Iztočnik / Source:} Matrizenprodukt\par
\noindent\emph{Rusky / Russian:} матричное произведение\par
\noindent\emph{Ukrajinsky / Ukrainian:} матричний добуток\par
\smallskip\noindent\emph{Razlog:} Matričny produkt v rabotah 3 i 4 potrěbuje istorijno objašnjenje. Tekst pod nim razuměje sumu produktov odgovarjajučih determinantov dvoh matric s ravnym čislom redov; to ne jest samo obyčno množenje matric s matričnym rezultatom. Arhivne zapisy daju prědnost internacionalnomu produkt, dokolě rusky sravnitelj piše proizvedenije a ukrajinsky dobutok; pozdnějše zapisy čuvajut kontinuitet s prvim izborom i teoremom IV. Nazvati slovo standardnym ne ustanovjaje standardnogo sovrěmennogo smysla. Točna lokalna definicija mora byti dana pored nazva.\par
\noindent\emph{Reason:} Matričny produkt in Papers 3 and 4 requires a historical explanation. The text defines it as a sum of products of corresponding determinants from two matrices with the same number of rows; it is not simply ordinary matrix multiplication with a matrix as result. The preserved reasons prefer international produkt, against Russian proizvedenie and Ukrainian dobutok, and later records maintain continuity with the first choice and Theorem IV. Calling the vocabulary standard does not establish the ordinary modern meaning. The local definition must accompany the name.\par
\noindent{\footnotesize Izvorne zapisy / Original records: HIST-g028-013, HIST-g031-001, HIST-g035-003.}\par
\smallskip\noindent\emph{Sravnjenje s tekstom:} Raboty 3 i 4 izrično definujut matričny produkt kako sumu produktov odgovarjajučih determinantov. Produktna matrica jest iny objekt, čij determinant dava tu cěnnost. Tretji vid v §2 raboty 4 znači tretju lokalnu možnost sjedinjanja redov simbolov, ne tretji razměr ili tretju mocninu.\par
\noindent\emph{Textual comparison:} Papers 3 and 4 explicitly define the matrix product as a sum of products of corresponding determinants. The product matrix is a different object whose determinant gives that value. Third kind in §2 of Paper 4 means the third local arrangement of symbolic rows, not a third dimension or third power.\par
\noindent{\footnotesize TEXT-COMP-246; pasaže / passages: 3655, 3703, 3709, 3782, 3805, 3807, 3809, 3811, 4361, 4366, 4372, 4386; \texttt{CORPUS\_FORM\_COMPARISONS.json}.}\par
\end{minipage}\par
\par\medskip\Needspace{7\baselineskip}
\noindent\begin{minipage}{\linewidth}
\hypertarget{TERM-1137}{}
\noindent\textbf{matričny produkt / produktna matrica}\quad{\footnotesize\texttt{TERM-1137}}\par
\noindent\emph{Iztočnik / Source:} Matrizenprodukt / Produktmatrix\par
\noindent\emph{Rusky / Russian:} матричное произведение / матрица произведения\par
\noindent\emph{Ukrajinsky / Ukrainian:} матричний добуток / матриця добутку\par
\smallskip\noindent\emph{Razlog:} Matričny produkt i produktna matrica ne sut tu zaměnime nazvy jednogo objekta. V §1 raboty 4 suma produktov odgovarjajučih determinantov jest cěnnost determinanta produktnoj matrice. Prvy naziv označa toj skalarno vyraženy produkt, drugy matricu, iz ktorej se bere determinant. Arhivny zapis zadržuje internacionalne korenje, s paralelnymi ruskimi i ukrajinskimi frazami. Razlika v směru tvorjenja slov ne směje skryti razliku medžu matricoju i jej determinantom.\par
\noindent\emph{Reason:} Matričny produkt and produktna matrica are not interchangeable names for one object here. In §1 of Paper 4, the sum of products of corresponding determinants is the value of the determinant of the product matrix. The first expression names that scalar-valued product; the second names the matrix whose determinant is taken. The archived record retains international stems alongside parallel Russian and Ukrainian phrases. The change in the compound's structure must not conceal the distinction between a matrix and its determinant.\par
\noindent{\footnotesize Izvorne zapisy / Original records: HIST-g030-013.}\par
\smallskip\noindent\emph{Sravnjenje s tekstom:} Raboty 3 i 4 izrično definujut matričny produkt kako sumu produktov odgovarjajučih determinantov. Produktna matrica jest iny objekt, čij determinant dava tu cěnnost. Tretji vid v §2 raboty 4 znači tretju lokalnu možnost sjedinjanja redov simbolov, ne tretji razměr ili tretju mocninu.\par
\noindent\emph{Textual comparison:} Papers 3 and 4 explicitly define the matrix product as a sum of products of corresponding determinants. The product matrix is a different object whose determinant gives that value. Third kind in §2 of Paper 4 means the third local arrangement of symbolic rows, not a third dimension or third power.\par
\noindent{\footnotesize TEXT-COMP-246; pasaže / passages: 3655, 3703, 3709, 3782, 3805, 3807, 3809, 3811, 4361, 4366, 4372, 4386; \texttt{CORPUS\_FORM\_COMPARISONS.json}.}\par
\end{minipage}\par
\par\medskip\Needspace{7\baselineskip}
\noindent\begin{minipage}{\linewidth}
\hypertarget{TERM-1138}{}
\noindent\textbf{matričny produkt tretjega vida}\quad{\footnotesize\texttt{TERM-1138}}\par
\noindent\emph{Iztočnik / Source:} Matrizenprodukt dritter Art\par
\noindent\emph{Rusky / Russian:} матричное произведение третьего типа\par
\noindent\emph{Ukrajinsky / Ukrainian:} матричний добуток третього типу\par
\smallskip\noindent\emph{Razlog:} Tretjego vida označa tretju iz trěh lokalno navedenyh možnostij v §2 raboty 4: redy matrice mogut byti dane oddělno, sjedinjene v jedin red simbolov, ili sjedinjene v nekoliko redov simbolov. Iz dvoh redov $S^\sigma,T^\tau$ s pomočju redov prěměnnyh tvori se produkt togo tretjego vida. Arhivny razlog izbira vid za německo Art bez tvrdnje o obče prijetoj klasifikaciji. To zato ne jest tretji stepen matrice, tretja mocnina produkta ili novy obči način množenja matric.\par
\noindent\emph{Reason:} Tretjego vida identifies the third of three possibilities listed locally in §2 of Paper 4: matrix rows may be given separately, assembled into one symbolic row, or assembled into several symbolic rows. Two rows $S^\sigma,T^\tau$, together with variable rows, give a product of the third kind. The archived reason chooses vid for German Art without claiming a generally established classification. It does not mean a matrix of order three, the third power of a product, or a new general operation of matrix multiplication.\par
\noindent{\footnotesize Izvorne zapisy / Original records: HIST-g031-006.}\par
\smallskip\noindent\emph{Sravnjenje s tekstom:} Raboty 3 i 4 izrično definujut matričny produkt kako sumu produktov odgovarjajučih determinantov. Produktna matrica jest iny objekt, čij determinant dava tu cěnnost. Tretji vid v §2 raboty 4 znači tretju lokalnu možnost sjedinjanja redov simbolov, ne tretji razměr ili tretju mocninu.\par
\noindent\emph{Textual comparison:} Papers 3 and 4 explicitly define the matrix product as a sum of products of corresponding determinants. The product matrix is a different object whose determinant gives that value. Third kind in §2 of Paper 4 means the third local arrangement of symbolic rows, not a third dimension or third power.\par
\noindent{\footnotesize TEXT-COMP-246; pasaže / passages: 3655, 3703, 3709, 3782, 3805, 3807, 3809, 3811, 4361, 4366, 4372, 4386; \texttt{CORPUS\_FORM\_COMPARISONS.json}.}\par
\end{minipage}\par
\par\medskip\Needspace{7\baselineskip}
\noindent\begin{minipage}{\linewidth}
\hypertarget{TERM-1139}{}
\noindent\textbf{matrične produkty}\quad{\footnotesize\texttt{TERM-1139}}\par
\noindent\emph{Iztočnik / Source:} Matrizenprodukte\par
\noindent\emph{Rusky / Russian:} матричные произведения\par
\noindent\emph{Ukrajinsky / Ukrainian:} матричні добутки\par
\smallskip\noindent\emph{Razlog:} Množny naziv matrične produkty v uvodu raboty 4 čuva to samo nazvanje kako odělne konstrukcije. Uvod najprvo najavjaje, že invariantne tvorjenja možno prědstaviti prěz te produkty, a dodatna zamětka sravnivaje jih s Grassmannovym směšanym produktom stupnja nula. Arhivny razlog govori o prozračnosti produkt i o inoj ruskoj i ukrajinskoj leksice. Gramatično čislo imenika ne tvori novy matematičny smysl. Ale obča razumlivost nazva ne oslobađa objašnjenje od lokalnoj definicije, osoblivo dlja čitatelja, ktory znajet samo obyčno množenje matric.\par
\noindent\emph{Reason:} The plural matrične produkty in the introduction to Paper 4 continues the name used for the individual constructions. The introduction announces representation of invariant expressions by these products; an added note compares them with Grassmann's mixed product of degree zero. The recorded reason invokes the transparency of produkt and the different Russian and Ukrainian vocabulary. Grammatical number does not create another mathematical sense. Yet a familiar name does not remove the need for the local definition, especially for readers who know only ordinary matrix multiplication.\par
\noindent{\footnotesize Izvorne zapisy / Original records: HIST-g029-009.}\par
\smallskip\noindent\emph{Sravnjenje s tekstom:} Raboty 3 i 4 izrično definujut matričny produkt kako sumu produktov odgovarjajučih determinantov. Produktna matrica jest iny objekt, čij determinant dava tu cěnnost. Tretji vid v §2 raboty 4 znači tretju lokalnu možnost sjedinjanja redov simbolov, ne tretji razměr ili tretju mocninu.\par
\noindent\emph{Textual comparison:} Papers 3 and 4 explicitly define the matrix product as a sum of products of corresponding determinants. The product matrix is a different object whose determinant gives that value. Third kind in §2 of Paper 4 means the third local arrangement of symbolic rows, not a third dimension or third power.\par
\noindent{\footnotesize TEXT-COMP-246; pasaže / passages: 3655, 3703, 3709, 3782, 3805, 3807, 3809, 3811, 4361, 4366, 4372, 4386; \texttt{CORPUS\_FORM\_COMPARISONS.json}.}\par
\end{minipage}\par
\par\medskip\Needspace{7\baselineskip}
\noindent\begin{minipage}{\linewidth}
\hypertarget{TERM-1140}{}
\noindent\textbf{matrično kolco}\quad{\footnotesize\texttt{TERM-1140}}\par
\noindent\emph{Iztočnik / Source:} Matrizenring\par
\noindent\emph{Rusky / Russian:} матричное кольцо\par
\noindent\emph{Ukrajinsky / Ukrainian:} матричне кільце\par
\smallskip\noindent\emph{Razlog:} Matrično kolco jest arhivna forma nazva; nynějšni tekst §25 raboty 34 govori o matričnom prstenu. Obe formy tu odnoset se k prstenu matric, v ktorom računaje se diskriminant posrědstvom matričnyh jedinic. Arhivny razlog daje kolco radi očekyvanoj razumlivosti pored ruskogo kolco i ukrajinskogo kilce. Zapis ne dava čitateljskogo testa ni opravdanja zaměny za prsten vo vseh kontekstah. Slovnik mora pokazati obě istorijne formy i ih isto značenje, ne predstaviti staru formu kako nynějšnu normu.\par
\noindent\emph{Reason:} Matrično kolco is the archived form; the current text of §25 of Paper 34 speaks of a matričny prsten. Both refer here to a matrix ring whose discriminant is calculated using matrix units. The archived reason prefers kolco for anticipated intelligibility alongside Russian koltso and Ukrainian kiltse. It supplies neither a reader test nor a justification for replacement by prsten in every context. The dictionary must show both historical forms and their shared meaning, rather than present the earlier form as the current norm.\par
\noindent{\footnotesize Izvorne zapisy / Original records: HIST-g203-009.}\par
\smallskip\noindent\emph{Sravnjenje s tekstom:} §25 raboty 34 upotrěbjaje matričny prsten i računaje jego diskriminant s pomočju dvoh uporedženij matričnyh jedinic. Posebno indeksno pytanje pojavlja se raněje: pri zapisanoj proměně bazy $(b)=(a)P$ slědujuče formuly stavljajut matricu Grama bazy $(a)$ ravnoju $P^t$ puta matrici Grama bazy $(b)$ puta $P$. Napravljenje proměny i te ravnosti ne soglašajut se; potrěbno jest sravnjenje s izvornym svědkom.\par
\noindent\emph{Textual comparison:} Section 25 of Paper 34 uses matričny prsten and computes its discriminant using two orderings of the matrix units. A separate notation question occurs earlier: after the displayed basis change $(b)=(a)P$, the following formulas equate the Gram matrix of $(a)$ with $P^t$ times the Gram matrix of $(b)$ times $P$. The direction of the basis change and these equalities disagree; the original witness needs to be checked.\par
\noindent{\footnotesize TEXT-COMP-247; pasaže / passages: 22307, 22309, 22313, 22317, 22321, 22325, 22329, 22374, 22396, 22397; \texttt{CORPUS\_FORM\_COMPARISONS.json}.}\par
\end{minipage}\par
\par\medskip\Needspace{7\baselineskip}
\noindent\begin{minipage}{\linewidth}
\hypertarget{TERM-1141}{}
\noindent\textbf{prěobraženje matric}\quad{\footnotesize\texttt{TERM-1141}}\par
\noindent\emph{Iztočnik / Source:} Matrizenumformung\par
\noindent\emph{Rusky / Russian:} преобразование матриц\par
\noindent\emph{Ukrajinsky / Ukrainian:} перетворення матриць\par
\smallskip\noindent\emph{Razlog:} Prěobraženje matric razklada německy složeny naziv na imenik dějanja i genitiv objekta. Arhivny razlog izbiraje sovrěmenny linearno-algebraičny izraz vměsto doslovnogo prěnosa složenogo Umformung; rusky i ukrajinsky sravniteljny izraz imajut tu samu stavbu. To opravdava konstrukciju nazva, ale samo po sobě ne oprěděljaje, kakovo prěobraženje jest dopuščeno. V zachovanom zapisu jest odkaz na rabotu 23, ale točno tako nazvana operacija ne byla ustanovjena v nynějšnom tekstu toj raboty. Zato ne priravnjajemo nazva bez daljšego svědka k podobnosti matric, elementarnym rjadovym operacijam ili proměně bazy.\par
\noindent\emph{Reason:} Prěobraženje matric renders the German compound as a noun of action followed by the object in the genitive. The archived reason favours modern linear-algebraic wording over a literal reshaping compound; the Russian and Ukrainian comparanda have the same construction. This supports the wording but does not determine which transformations are allowed. The record points to Paper 23, but an operation with this precise designation has not been established in its current text. Without a further witness, the expression should therefore not be equated with matrix similarity, elementary row operations or change of basis.\par
\noindent{\footnotesize Izvorne zapisy / Original records: HIST-g094-024.}\par
\end{minipage}\par
\par\medskip\Needspace{7\baselineskip}
\noindent\begin{minipage}{\linewidth}
\hypertarget{TERM-1142}{}
\noindent\textbf{matrice}\quad{\footnotesize\texttt{TERM-1142}}\par
\noindent\emph{Iztočnik / Source:} Matrizes\par
\noindent\emph{Rusky / Russian:} матрицы\par
\noindent\emph{Ukrajinsky / Ukrainian:} матриці\par
\smallskip\noindent\emph{Razlog:} Matrice jest množna forma od matrica. Arhivny razlog spojuje medžunarodny korenj, rusku i ukrajinsku formu i kontinuitet s ranějšim v matricah. V §7 raboty 34 ide o kvadratnyh matricah nad prstenom, s dodavanjem po elementah i obyčnym matričnym množenjem. Tu ne jest potrěbny novy naziv za stary německy plural Matrizes. Ale obyčno množenje v toj rabotě trěba razlikovati od specialno definovanyh matričnyh produktov ranih rabot 3 i 4; sam korenj matrica ne garantuje isto značenje vsakoj složenej frazy.\par
\noindent\emph{Reason:} Matrice is the plural of matrica. The archived reason combines the international stem, Russian and Ukrainian forms, and continuity with the earlier v matricah. Section 7 of Paper 34 concerns square matrices over a ring, with entrywise addition and ordinary matrix multiplication. The older German plural Matrizes does not require a different object name. Ordinary multiplication here must nevertheless be distinguished from the specially defined matrix products of Papers 3 and 4: a shared matrix stem does not guarantee that every compound has the same meaning.\par
\noindent{\footnotesize Izvorne zapisy / Original records: HIST-g185-001.}\par
\smallskip\noindent\emph{Sravnjenje s tekstom:} §7 raboty 34 definira obyčno matrično množenje i potom razlikuje tělo $K$ od dodatnogo komutativnogo slučaja. Matrice nad $K$ zato imajut elementy iz $K$ i ne prědpokladajut avtomatično komutativnost.\par
\noindent\emph{Textual comparison:} Section 7 of Paper 34 defines ordinary matrix multiplication and then distinguishes a division ring $K$ from the additional commutative case. Matrices over $K$ thus have their entries in $K$ without automatically assuming commutativity.\par
\noindent{\footnotesize TEXT-COMP-248; pasaže / passages: 20327, 20329, 20331, 20333, 20335, 20337, 20352, 20354, 20356, 20362; \texttt{CORPUS\_FORM\_COMPARISONS.json}.}\par
\end{minipage}\par
\par\medskip\Needspace{7\baselineskip}
\noindent\begin{minipage}{\linewidth}
\hypertarget{TERM-1143}{}
\noindent\textbf{matrice nad $K$}\quad{\footnotesize\texttt{TERM-1143}}\par
\noindent\emph{Iztočnik / Source:} Matrizes in $K$\par
\noindent\emph{Rusky / Russian:} матрицы над $K$\par
\noindent\emph{Ukrajinsky / Ukrainian:} матриці над $K$\par
\smallskip\noindent\emph{Razlog:} Nad v matricah nad $K$ označaje, že matrične elementy naležet $K$, ne že matrice same sut elementami $K$. Arhivny razlog izbiraje tu konstrukciju kako sovrěmenny algebraičny izraz, soglasny s ruskim i ukrajinskim nad. §7 raboty 34 izrično dopušča tělo $K$ bez prědpoloženja komutativnosti; komutativnost pojavlja se kako dodatno uslovje pri nazivu hiperkompleksny sistem. Zato slovnično objašnjenje ne směje zaměniti $K$ komutativnym poljem bez ogovorki. Formulacija obhvata i nekomutativny slučaj, pri ktorom poredok množenja elementov ostaje sučstvenny.\par
\noindent\emph{Reason:} Nad in matrice nad $K$ says that the matrix entries belong to $K$, not that the matrices themselves are elements of $K$. The archived reason selects the modern algebraic construction, parallel to Russian and Ukrainian nad. Section 7 of Paper 34 explicitly allows a division ring $K$ without assuming commutativity; commutativity is added when the matrix ring is described as a hypercomplex system. A dictionary explanation must therefore not silently replace $K$ by a commutative field. The phrase includes the noncommutative case, where the order of coefficient multiplication matters.\par
\noindent{\footnotesize Izvorne zapisy / Original records: HIST-g185-005.}\par
\smallskip\noindent\emph{Sravnjenje s tekstom:} §7 raboty 34 definira obyčno matrično množenje i potom razlikuje tělo $K$ od dodatnogo komutativnogo slučaja. Matrice nad $K$ zato imajut elementy iz $K$ i ne prědpokladajut avtomatično komutativnost.\par
\noindent\emph{Textual comparison:} Section 7 of Paper 34 defines ordinary matrix multiplication and then distinguishes a division ring $K$ from the additional commutative case. Matrices over $K$ thus have their entries in $K$ without automatically assuming commutativity.\par
\noindent{\footnotesize TEXT-COMP-248; pasaže / passages: 20327, 20329, 20331, 20333, 20335, 20337, 20352, 20354, 20356, 20362; \texttt{CORPUS\_FORM\_COMPARISONS.json}.}\par
\end{minipage}\par
\par\medskip\Needspace{7\baselineskip}
\noindent\begin{minipage}{\linewidth}
\hypertarget{TERM-1144}{}
\noindent\textbf{uslovje maksimalnosti}\quad{\footnotesize\texttt{TERM-1144}}\par
\noindent\emph{Iztočnik / Source:} Maximalbedingung\par
\noindent\emph{Rusky / Russian:} условие максимальности\par
\noindent\emph{Ukrajinsky / Ukrainian:} умова максимальності\par
\smallskip\noindent\emph{Razlog:} Uslovje maksimalnosti čuva istorijny naziv i vezu s uslovjem prekinjenja vzrastajučih verig. V §3 raboty 34 maksimalna podgrupa znači maksimalnu v razgledanom množstvu dopustimyh podgrup, ne najveliku podgrupu vsěh grup. Smysl jest, že v vsakom nepustom takim množstvu obstaja maksimalny element; ekvivalentno, strogo vzrastajuča veriga ne može prodolžati se bez konca. Slovo nepusto tu objasnjava logičny obseg definicije: nynějšna formulacija teksta jego ne izgovarja. Arhivny razlog daje prědnost istorijnomu uslovju, a ne neposrědnoj zaměně nazvom neterovost, ktory by zakryl razgledany vid podgrup.\par
\noindent\emph{Reason:} Uslovje maksimalnosti preserves the historical name and its connection with termination of ascending chains. In §3 of Paper 34, maximal means maximal within the specified collection of admissible subgroups, not a greatest subgroup of all groups. The intended condition is that every nonempty such collection has a maximal member; equivalently, a strictly ascending chain cannot continue indefinitely. Nonempty here makes the logical scope explicit: the current wording of the text does not say it. The archived reason favours the historical condition rather than directly replacing it with Noetherianity, which could obscure the class of subgroups under consideration.\par
\noindent{\footnotesize Izvorne zapisy / Original records: HIST-g181-005.}\par
\smallskip\noindent\emph{Sravnjenje s tekstom:} Definicije v §3 povezuju maksimalnost so vzrastajučimi verigami i minimalnost so spadajučimi. Slovnično objasnjenje dodava nepusto k razgledanom množstvu; to jest logično objasnjenje, ne citat iz toj formulacije. §4 razlikuje prědpoloženje za konečno razloženje od silnějših prědpoloženij dlja jednoznačnosti.\par
\noindent\emph{Textual comparison:} The definitions in §3 connect maximality with ascending chains and minimality with descending chains. The dictionary explanation supplies nonempty for the collection considered; this is a logical clarification, not a quotation from that wording. Section 4 distinguishes the hypothesis for a finite decomposition from the stronger hypotheses used for uniqueness.\par
\noindent{\footnotesize TEXT-COMP-249; pasaže / passages: 19993, 19995, 19997, 19999, 20071, 20073; \texttt{CORPUS\_FORM\_COMPARISONS.json}.}\par
\end{minipage}\par
\par\medskip\Needspace{7\baselineskip}
\noindent\begin{minipage}{\linewidth}
\hypertarget{TERM-1145}{}
\noindent\textbf{uslovje maksimalnosti / minimalnosti za prave idealy}\quad{\footnotesize\texttt{TERM-1145}}\par
\noindent\emph{Iztočnik / Source:} Maximalbedingung / Minimalbedingung für Rechtsideale\par
\noindent\emph{Rusky / Russian:} условие максимальности / минимальности для правых идеалов\par
\noindent\emph{Ukrajinsky / Ukrainian:} умова максимальності / мінімальності для правих ідеалів\par
\smallskip\noindent\emph{Razlog:} Dopolnjenje za prave idealy jest částju uslovja, ne nepovinna stilistična podrobnost. Arhivny razlog prěnosi nazvy maksimalnosti i minimalnosti iz §3 v §8 raboty 34. Tu dopustimy pravy ideal jest jednočasno $K$-modul; konečny rang daje ograničenje dlja strogo vzrastajučih i spadajučih verig. Prave ne možno izpustiti, osoblivo v nekomutativnom prstenu: uslovje za prave idealy samo svojom nazvoju ne ustanavjaje analogičnogo uslovja za lěve. Parna formulacija maksimalnosti / minimalnosti pokazuje dva uslovja, ne jednu nerazličenu vlastnost.\par
\noindent\emph{Reason:} The qualifier za prave idealy is part of the condition, not optional stylistic detail. The recorded reason carries the maximality and minimality terminology from §3 into §8 of Paper 34. Here an admissible right ideal is also a $K$-module, and finite rank bounds both strictly ascending and strictly descending chains. Right cannot be omitted, especially in a noncommutative ring: a condition on right ideals does not establish the corresponding left condition merely by being named. The paired wording maksimalnosti / minimalnosti denotes two conditions, not one undifferentiated property.\par
\noindent{\footnotesize Izvorne zapisy / Original records: HIST-g186-007.}\par
\smallskip\noindent\emph{Sravnjenje s tekstom:} V §8 raboty 34 konečny rang $K$-modula ograničaje verigi dopustimyh pravyh idealov. Ograničenje prave jest izrično; njegovo izpuščenje by razširilo tvrdnju izvan citovanogo argumenta.\par
\noindent\emph{Textual comparison:} In §8 of Paper 34, finite $K$-module rank bounds chains of admissible right ideals. The right-sided qualification is explicit; removing it would extend the statement beyond the cited argument.\par
\noindent{\footnotesize TEXT-COMP-250; pasaže / passages: 20429; \texttt{CORPUS\_FORM\_COMPARISONS.json}.}\par
\end{minipage}\par
\par\medskip\Needspace{7\baselineskip}
\noindent\begin{minipage}{\linewidth}
\hypertarget{TERM-1146}{}
\noindent\textbf{maksimalna oblast}\quad{\footnotesize\texttt{TERM-1146}}\par
\noindent\emph{Iztočnik / Source:} Maximalbereich\par
\noindent\emph{Rusky / Russian:} максимальная область\par
\noindent\emph{Ukrajinsky / Ukrainian:} максимальна область\par
\smallskip\noindent\emph{Razlog:} Maksimalna oblast v rabotě 35 imaje osobnu definiciju. V danoj okolnoj oblasti $\mathfrak Z$ ona ne može byti razširjena dodavanjem cěločislnyh imenovateljev: ako $ag$ naleži jej, $g\in\mathfrak Z$ i $a$ jest nenulovo cělo čislo, togda jej naleži i $g$. To jest zamknjenost pod takim děljenjem, ne maksimalnost sredi vsih sobstvennyh podprstenov. Arhivny razlog čuva oblast za Bereich, medžutym kako prstenova struktura jest oprěděljena odělno. Razlikovanje jest važnějše od same doslovnosti: naziv ne znači proizvoljno integralno zamknjenje ni normalizaciju, a ograničenje na cěločislne imenovatelje mora ostati vidimo.\par
\noindent\emph{Reason:} Maksimalna oblast has a particular definition in Paper 35. Within a fixed ambient domain $\mathfrak Z$, it cannot be enlarged by adjoining integer denominators: if it contains $ag$, where $g\in\mathfrak Z$ and $a$ is a nonzero integer, it also contains $g$. This is closure under that division operation, not maximality among all proper subrings. The recorded reason preserves oblast for Bereich while the ring structure is specified separately. The distinction matters more than literal wording: the name does not mean arbitrary integral closure or normalization, and the restriction to integer denominators must remain visible.\par
\noindent{\footnotesize Izvorne zapisy / Original records: HIST-g205-001.}\par
\smallskip\noindent\emph{Sravnjenje s tekstom:} Rabota 35 definira maksimalnost vnutri danoj $\mathfrak Z$ posrědstvom uslovja, že $ag\in M$ vleče $g\in M$ za nenulove cěle $a$ i $g\in\mathfrak Z$. Opisana oblast $M(\mathfrak I)$ jest takym načinom zamknuta oblast porodžena prěz $\mathfrak I$; naziv ne označaje maksimalny sobstvenny podprsten.\par
\noindent\emph{Textual comparison:} Paper 35 defines maximality inside a fixed $\mathfrak Z$ by the condition that $ag\in M$ implies $g\in M$ for nonzero integers $a$ and $g\in\mathfrak Z$. The described $M(\mathfrak I)$ is the domain generated from $\mathfrak I$ by this closure operation; the term does not denote a maximal proper subring.\par
\noindent{\footnotesize TEXT-COMP-251; pasaže / passages: 22525, 22527, 22576, 22578; \texttt{CORPUS\_FORM\_COMPARISONS.json}.}\par
\end{minipage}\par
\par\medskip\Needspace{7\baselineskip}
\noindent\begin{minipage}{\linewidth}
\hypertarget{TERM-1147}{}
\noindent\textbf{maksimalny nilpotentny ideal}\quad{\footnotesize\texttt{TERM-1147}}\par
\noindent\emph{Iztočnik / Source:} maximales nilpotentes Ideal\par
\noindent\emph{Rusky / Russian:} максимальный нильпотентный идеал\par
\noindent\emph{Ukrajinsky / Ukrainian:} максимальний нільпотентний ідеал\par
\smallskip\noindent\emph{Razlog:} Nilpotentny jest zachovany medžunarodny matematičny korenj, ktory podpirajut i oba zapisane istočnoslovjanske sravnitve. Nilpotentnost ideala znači zanuljenje někojej potegy cělogo ideala, ne samo nilpotentnost vsakogo elementa odělno. V rabotě 34 uslovje maksimalnosti za prave idealy i teorem o sumě dvoh nilpotentnyh pravyh idealov vedut k jednomu takomu idealu, ktory sodrživaje vse ostale. Tekst potom pokazuje, že on jest dvustranny, i nazyvaje jego radikalom. Zato maksimalny tu stavaje takože najveliky v toj klasě, ale to slěduje iz argumenta; ne jest opča definicija slova maksimalny.\par
\noindent\emph{Reason:} Nilpotentny retains the international mathematical stem also represented in both recorded East Slavic comparanda. Nilpotence of an ideal means that some power of the whole ideal is zero, not merely that each element is separately nilpotent. In Paper 34, the ascending-chain condition on right ideals and the result about sums of two nilpotent right ideals lead to one such ideal containing all the others. The text then shows that it is two-sided and calls it the radical. Thus maximal is also greatest in this particular class, but that follows from the argument; it is not the general meaning of maximal.\par
\noindent{\footnotesize Izvorne zapisy / Original records: HIST-g196-004.}\par
\smallskip\noindent\emph{Sravnjenje s tekstom:} Argument o sumě nilpotentnyh idealov vedet k najvelikomu takomu idealu. Posebne pytanje pojavlja se v §19: formulacija teorema davaje anulovanje od radikala kako jednu iz dvoh alternativ, ale dokaz ustanavjaje to anulovanje dlja vsakogo prostogo modula. Dalje dokaz razlikuje anulovanje cělym kvocientnym prstenom od dejstvija jego jedinice kako jediničnogo operatora. Formulaciju i poslědok trěba sravniti s izvornym svědkom; ne prěnosimo ih kakoby už soglasovanu dihotomiju.\par
\noindent\emph{Textual comparison:} The sum argument gives a greatest nilpotent ideal. A separate question arises in §19: the theorem states annihilation by the radical as one of two alternatives, whereas the proof establishes this annihilation for every simple module. The proof then distinguishes annihilation by the entire quotient ring from its identity acting as the identity operator. The statement and corollary need comparison with the original witness; they are not adopted here as an already reconciled dichotomy.\par
\noindent{\footnotesize TEXT-COMP-252; pasaže / passages: 20857, 20860, 20862, 21735, 21737, 21739, 21741, 21743; \texttt{CORPUS\_FORM\_COMPARISONS.json}.}\par
\end{minipage}\par
\par\medskip\Needspace{7\baselineskip}
\noindent\begin{minipage}{\linewidth}
\hypertarget{TERM-1148}{}
\noindent\textbf{maksimalne porjadky}\quad{\footnotesize\texttt{TERM-1148}}\par
\noindent\emph{Iztočnik / Source:} Maximalordnungen\par
\noindent\emph{Rusky / Russian:} максимальные порядки\par
\noindent\emph{Ukrajinsky / Ukrainian:} максимальні порядки\par
\smallskip\noindent\emph{Razlog:} Maksimalne porjadky jest arhivny naziv; nynějšny tekst raboty 42 imaje maksimalne poredky. Arhivny razlog davaje prědnost kontinuitetu s ranějšim nazivom i istočnoslovjanskimi sravnitvami. V obojih formah ide o algebraičnyh poredkah v algebraji nad čislovym poljem, ne o relacijah uporjadkovanja. §II tvori take poredky iz komplementarnyh konečnyh modulov i opisuje jih prěz lěve i prave idealy. Zato anglijsky izraz order theory v starom opravdanju ne trěba razuměti kako teoriju čestičnyh poredkov. Razlika porjadok / poredok ostaje slovničnym pytanjem; matematičny objekt od nje ne mjenja se.\par
\noindent\emph{Reason:} Maksimalne porjadky is the archived wording; the current text of Paper 42 has maksimalne poredky. The recorded reason favours continuity with the earlier term and the East Slavic comparanda. Both forms denote algebraic orders in an algebra over a number field, not ordering relations. Section II constructs such orders from complementary finite modules and describes them through left and right ideals. The phrase order theory in the older rationale should therefore not be understood as the theory of partially ordered sets. The porjadok / poredok variation remains a lexical question, not a change of mathematical object.\par
\noindent{\footnotesize Izvorne zapisy / Original records: HIST-g212-002.}\par
\smallskip\noindent\emph{Sravnjenje s tekstom:} §II raboty 42 imaje poredky i konstruuje jih v razloženoj algebraji iz komplementarnyh modulov. To kontekstualno oprěděljaje smysl arhivnogo porjadky; ne govori se o relaciji čestičnogo uporjadkovanja.\par
\noindent\emph{Textual comparison:} Section II of Paper 42 uses poredky and constructs them in a split algebra from complementary modules. This establishes the contextual sense of the archived porjadky; the subject is not a partial-order relation.\par
\noindent{\footnotesize TEXT-COMP-253; pasaže / passages: 24553, 24555, 24557, 24558, 24560, 24562, 24564, 24566; \texttt{CORPUS\_FORM\_COMPARISONS.json}.}\par
\end{minipage}\par
\par\medskip\Needspace{7\baselineskip}
\noindent\begin{minipage}{\linewidth}
\hypertarget{TERM-1149}{}
\noindent\textbf{množinno-teoretičny prěsěk}\quad{\footnotesize\texttt{TERM-1149}}\par
\noindent\emph{Iztočnik / Source:} mengentheoretischer Durchschnitt\par
\noindent\emph{Rusky / Russian:} теоретико-множественное пересечение\par
\noindent\emph{Ukrajinsky / Ukrainian:} теоретико-множинний перетин\par
\smallskip\noindent\emph{Razlog:} Množinno-teoretičny prěsěk izrično govori o občih elementah dvoh množstv. V začetku raboty 31 množstva prstena $\mathfrak R$ i razširjenogo polja $\overline P$ imajut obči prěsěk ravny $P$; dodatno uslovje govori o odnošenjah operacij izvan $P$. Arhivny razlog hoče razlikovati tu konstrukciju od razgovora o idealah i modulah. To ne znači, že prěsěk idealov jest druga množinstvena operacija. Pridevnik označaje pogled na osnovne množstva i ne može sam vyraziti dodatno algebraično uslovje. Uslovja zato ne slivajemo v neoprěděljenu tvrdnju o nezavisnosti.\par
\noindent\emph{Reason:} Množinno-teoretičny prěsěk explicitly concerns the common elements of two sets. At the beginning of Paper 31, the underlying sets of the ring $\mathfrak R$ and the extension field $\overline P$ intersect in $P$; a further condition concerns relations between operations outside $P$. The archived reason seeks to distinguish this construction from discussion of ideals and modules. This does not make intersection of ideals a different set operation. The adjective identifies the underlying-set perspective and cannot by itself express the additional algebraic condition. The two conditions should not be collapsed into an unspecified claim of independence.\par
\noindent{\footnotesize Izvorne zapisy / Original records: HIST-g119-007.}\par
\smallskip\noindent\emph{Sravnjenje s tekstom:} Prěsěk prstena $\mathfrak R$ i polja $\overline P$ jest $P$, a slědujuče uslovje o operacijah jest odělno. Primětka opisuje ekvivalentne razširjenja kako izomorfne razširjenja, čij izomorfizm fiksuje $P$ element po elementu.\par
\noindent\emph{Textual comparison:} The ring $\mathfrak R$ and field $\overline P$ intersect in $P$, while the following condition about operations is separate. The footnote describes equivalent extensions as isomorphic extensions whose isomorphism fixes $P$ element by element.\par
\noindent{\footnotesize TEXT-COMP-254; pasaže / passages: 17322, 17324, 17326; \texttt{CORPUS\_FORM\_COMPARISONS.json}.}\par
\end{minipage}\par
\par\medskip\Needspace{7\baselineskip}
\noindent\begin{minipage}{\linewidth}
\hypertarget{TERM-1150}{}
\noindent\textbf{determinanta reda najmanje dva}\quad{\footnotesize\texttt{TERM-1150}}\par
\noindent\emph{Iztočnik / Source:} mindestens zweireihige Determinante\par
\noindent\emph{Rusky / Russian:} определитель порядка не менее двух\par
\noindent\emph{Ukrajinsky / Ukrainian:} визначник порядку принаймні два\par
\smallskip\noindent\emph{Razlog:} Determinanta reda najmanje dva zaměnjaje doslovno dvurjadnu determinantu nazivom za razměr kvadratnoj matrice. Arhivny razlog izbiraje medžunarodno determinanta i red, a ruska i ukrajinska sravnitva takože označajut red. V §6 raboty 31 govor ide o matričnoj razměrnosti prinajmenje dva: vse elementy diskriminantnoj matrice kromě vozmožno jednoj pozicije izčezajut, zato njejna determinanta jest nula. Ne trěba čitati najmanje dva kako točno dva, ni kako tvrdnju samo o grafičnom čislu rjadov. Nynějšny tekst imaje prinajmenje dva; to jest ta sama dolna granica, ne nova matematična definicija.\par
\noindent\emph{Reason:} Determinanta reda najmanje dva replaces the literal two-rowed determinant with a description of the size of a square matrix. The archived reason selects international determinanta and red, while the Russian and Ukrainian comparanda likewise specify the order. In §6 of Paper 31, the size is at least two: every entry of the discriminant matrix except possibly one vanishes, so its determinant is zero. Najmanje dva must not be read as exactly two or as a statement merely about the visual number of rows. The current text uses prinajmenje dva, expressing the same lower bound rather than another mathematical definition.\par
\noindent{\footnotesize Izvorne zapisy / Original records: HIST-g143-006.}\par
\smallskip\noindent\emph{Sravnjenje s tekstom:} Diskriminantna matrica imaje prinajmenje dva rjada i stolpca, a samo jedna pozicija može byti nenulova. To jest lokalny razlog zanuljenja determinanty; najmanje i prinajmenje vyražajut tu samu neostru dolnu granicu.\par
\noindent\emph{Textual comparison:} The discriminant matrix has at least two rows and columns, but only one position can be nonzero. This is the local reason its determinant vanishes; najmanje and prinajmenje express the same inclusive lower bound.\par
\noindent{\footnotesize TEXT-COMP-255; pasaže / passages: 18383; \texttt{CORPUS\_FORM\_COMPARISONS.json}.}\par
\end{minipage}\par
\par\medskip\Needspace{7\baselineskip}
\noindent\begin{minipage}{\linewidth}
\hypertarget{TERM-1151}{}
\noindent\textbf{minimalna baza}\quad{\footnotesize\texttt{TERM-1151}}\par
\noindent\emph{Iztočnik / Source:} Minimalbasis\par
\noindent\emph{Rusky / Russian:} минимальная база; минимальный базис\par
\noindent\emph{Ukrajinsky / Ukrainian:} мінімальна база\par
\smallskip\noindent\emph{Razlog:} Četyri zapisy za minimalna baza podčrkivajut, že točny smysl davaje definicija, ne sama obča slova minimalna i baza. V §6 raboty 6 racionalna baza jest minimalna, ako sodrživaje samo $\rho$ funkcij; one togda sut algebraično nezavisne. Ide o racionalnom generovanju polja funkcij, ne o linearnoj bazi vektornogo prostora ni samo o neizbytočnom spisku generatorov. Teorem o egzistenciji takej bazy imaje svoja ograničenja: tekst odělno govori o razširjenju koeficientnogo polja i o neuspehu v občem pri $\rho\ge3$. Arhivno opravdanje zato podpiraje kratky naziv vměstě s definicijeju, ne bezuslovnu oběčanu egzistenciju.\par
\noindent\emph{Reason:} The four records for minimalna baza stress that the definition, not the ordinary words minimal and basis alone, supplies the precise sense. In §6 of Paper 6, a rational basis is minimal if it contains only $\rho$ functions; these are then algebraically independent. This concerns rational generation of a function field, not a vector-space basis or merely an irredundant generating list. Existence has restrictions: the text separately discusses extension of the coefficient field and failure in general for $\rho\ge3$. The recorded rationale therefore supports a concise name together with its definition, not an unconditional promise that such a basis exists.\par
\noindent{\footnotesize Izvorne zapisy / Original records: HIST-g039-008, HIST-g040-006, HIST-g046-001, HIST-g073-010.}\par
\smallskip\noindent\emph{Sravnjenje s tekstom:} Definicija IV v rabotě 6 žada racionalnu bazu iz $\rho$ funkcij, ktore sut algebraično nezavisne. Tri slědujuče tvrdnje razlikujut slučaj odnogo parametra, slučaj dvoh s koeficientnym razširjenjem i obči neuspeh pri vyššem $\rho$. Slovnik ne prěpisyvaje tu istorijnu informaciju kako bezuslovny sovrěmenny teorem.\par
\noindent\emph{Textual comparison:} Definition IV in Paper 6 requires a rational basis of $\rho$ algebraically independent functions. The following three statements distinguish one parameter, two with coefficient extension, and general failure for higher $\rho$. The dictionary does not recast this historical discussion as an unconditional modern theorem.\par
\noindent{\footnotesize TEXT-COMP-256; pasaže / passages: 5692, 5694, 5696, 5698, 5700; \texttt{CORPUS\_FORM\_COMPARISONS.json}.}\par
\end{minipage}\par
\par\medskip\Needspace{7\baselineskip}
\noindent\begin{minipage}{\linewidth}
\hypertarget{TERM-1152}{}
\noindent\textbf{uslovje minimalnosti}\quad{\footnotesize\texttt{TERM-1152}}\par
\noindent\emph{Iztočnik / Source:} Minimalbedingung\par
\noindent\emph{Rusky / Russian:} условие минимальности\par
\noindent\emph{Ukrajinsky / Ukrainian:} умова мінімальності\par
\smallskip\noindent\emph{Razlog:} Uslovje minimalnosti jest parny naziv k uslovju maksimalnosti. Arhivna dva razloga čuvajut kontinuitet medžu §3 i §4 raboty 34 i vezu s nynějšim govorom o spadajučih verigah. Smysl jest konečnost strogo spadajučih verig dopustimyh podgrup, ili obstajanje minimalnogo elementa v vsakom nepustom ih množstvu. Minimalny ne znači nužno jeden najmenši element dlja cělogo množstva. V §4 to uslovje daje konečno razloženje na direktno nerazložime faktory; primětka o jednoznačnosti žada oba uslovja. Ta razlika v prědpoloženjah ne směje izčeznuti pri objasnjenju modernym nazvom artinovost.\par
\noindent\emph{Reason:} Uslovje minimalnosti is paired with uslovje maksimalnosti. The two archived reasons preserve continuity between §§3 and 4 of Paper 34 and the connection with present-day descending-chain language. The condition requires termination of strictly descending chains of admissible subgroups, or a minimal member in every nonempty collection of them. Minimal need not mean one least member of the whole collection. In §4 the condition yields a finite decomposition into directly indecomposable factors; the note about uniqueness assumes both chain conditions. That difference in hypotheses must not disappear when explaining the term through modern Artinian terminology.\par
\noindent{\footnotesize Izvorne zapisy / Original records: HIST-g181-006, HIST-g182-006.}\par
\smallskip\noindent\emph{Sravnjenje s tekstom:} Definicije v §3 povezuju maksimalnost so vzrastajučimi verigami i minimalnost so spadajučimi. Slovnično objasnjenje dodava nepusto k razgledanom množstvu; to jest logično objasnjenje, ne citat iz toj formulacije. §4 razlikuje prědpoloženje za konečno razloženje od silnějših prědpoloženij dlja jednoznačnosti.\par
\noindent\emph{Textual comparison:} The definitions in §3 connect maximality with ascending chains and minimality with descending chains. The dictionary explanation supplies nonempty for the collection considered; this is a logical clarification, not a quotation from that wording. Section 4 distinguishes the hypothesis for a finite decomposition from the stronger hypotheses used for uniqueness.\par
\noindent{\footnotesize TEXT-COMP-249; pasaže / passages: 19993, 19995, 19997, 19999, 20071, 20073; \texttt{CORPUS\_FORM\_COMPARISONS.json}.}\par
\end{minipage}\par
\par\medskip\Needspace{7\baselineskip}
\noindent\begin{minipage}{\linewidth}
\hypertarget{TERM-1153}{}
\noindent\textbf{minimalne razpadne polja kvaternionnogo těla}\quad{\footnotesize\texttt{TERM-1153}}\par
\noindent\emph{Iztočnik / Source:} minimale Zerfaellungskorper des Quaternionenkorpers\par
\noindent\emph{Rusky / Russian:} минимальные поля расщепления тела кватернионов\par
\noindent\emph{Ukrajinsky / Ukrainian:} мінімальні поля розщеплення кватерніонного тіла\par
\smallskip\noindent\emph{Razlog:} Minimalno razpadno polje jest tu polje, ktoro razklada kvaternionno tělo, ale ne sodrživaje sobstvennogo podpolja s toju samoju vlastnostju. Minimalnost jest po vključanju, ne po najnižšem stepenu. Rabota 32 izrično davaje taka polja neograničenyh stepenjev. Arhivny razlog čuva razpadno polje radi kontinuiteta, a tělo radi nekomutativnogo kvaternionnogo objekta. On takože spominaje češsky i poljsky korenj rozklad- kako alternativu za razgledanje, ale ne dava točnogo slovničnogo ili knižnogo města; to ne jest dokaz čitanja tih izvorov. Izbor korenja i definicija minimalnosti sut zato dva odělna pytanja.\par
\noindent\emph{Reason:} A minimal splitting field here splits the quaternion division algebra but has no proper subfield with that property. Minimality is by inclusion, not by the smallest degree. Paper 32 explicitly gives such fields of unbounded degrees. The archived reason keeps razpadno polje for continuity and tělo for the noncommutative quaternion object. It also mentions the Czech and Polish rozklad- stem as an alternative for consideration, but supplies no exact dictionary or book passage; this is not evidence that those sources were read. The choice of stem and the definition of minimality are therefore separate questions.\par
\noindent{\footnotesize Izvorne zapisy / Original records: HIST-g172-002.}\par
\smallskip\noindent\emph{Sravnjenje s tekstom:} Rabota 32 razgledaje ciklične polja stepena $2^n$: prědstavimost $-1$ kako sumy dvoh kvadratov davaje razpad, a realnost sobstvennyh podpolj izključaje razpad v nih. Minimalnost po vključanju soglašaje se s neograničenymi stepenjami.\par
\noindent\emph{Textual comparison:} Paper 32 considers cyclic fields of degree $2^n$: expressing $-1$ as a sum of two squares gives splitting, while the reality of proper subfields excludes splitting there. Minimality by inclusion is compatible with unbounded degrees.\par
\noindent{\footnotesize TEXT-COMP-257; pasaže / passages: 19486, 19490; \texttt{CORPUS\_FORM\_COMPARISONS.json}.}\par
\end{minipage}\par
\par\medskip\Needspace{7\baselineskip}
\noindent\begin{minipage}{\linewidth}
\hypertarget{TERM-1154}{}
\noindent\textbf{minimalna rezolventa}\quad{\footnotesize\texttt{TERM-1154}}\par
\noindent\emph{Iztočnik / Source:} Minimalresolvente\par
\noindent\emph{Rusky / Russian:} минимальная резольвента\par
\noindent\emph{Ukrajinsky / Ukrainian:} мінімальна резольвента\par
\smallskip\noindent\emph{Razlog:} Minimalna rezolventa čuva medžunarodny naziv s istorijnym algebraičnym smyslom. Arhivny razlog ne zaměnjaje jego občim pomočno uravnenje, zato že to by izbrisalo vezu s teorijeju eliminacije. V §12 raboty 19 nerazložime mnogovidnosti izolovanyh primarnyh komponentov sut spojene s minimalnoju rezolventoju; primětka odkazuje na Königovu knigu iz 1903, stranicu 235. Tu jest prověren odkaz v tekstu, ne nezavisno čitanje samoj Königovoj stranicy. Zato slovnik može pokazati lokalnu funkciju nazva i jego bibliografičnu vezu, ale ne vydavati odkaz za polnu, izvorně prověrenu definiciju rezolventy.\par
\noindent\emph{Reason:} Minimalna rezolventa preserves an international term in its historical algebraic sense. The archived reason avoids the generic auxiliary equation because that would erase the connection with elimination theory. In §12 of Paper 19, the irreducible varieties associated with isolated primary components are linked to the minimal resolvent; the footnote cites page 235 of König's 1903 book. What has been checked here is the citation in the text, not an independent reading of that page in König. The dictionary can therefore show the term's local role and bibliographical connection without presenting the citation as a fully source-checked definition.\par
\noindent{\footnotesize Izvorne zapisy / Original records: HIST-g081-083.}\par
\smallskip\noindent\emph{Sravnjenje s tekstom:} §12 raboty 19 spojuje minimalnu rezolventu s nerazložimymi mnogovidnostjami izolovanyh primarnyh komponentov. Odkaz na Königa ustanavjaje bibliografičnu vezu; čitanje samoj citovanoj stranicy ne jest tu zasvědčeno.\par
\noindent\emph{Textual comparison:} Section 12 of Paper 19 connects the minimal resolvent with the irreducible varieties of isolated primary components. The König reference establishes a bibliographical connection; reading the cited page itself is not attested here.\par
\noindent{\footnotesize TEXT-COMP-258; pasaže / passages: 13471, 13473, 13475; \texttt{CORPUS\_FORM\_COMPARISONS.json}.}\par
\end{minipage}\par
\par\medskip\Needspace{7\baselineskip}
\noindent\begin{minipage}{\linewidth}
\hypertarget{TERM-1157}{}
\noindent\textbf{s koeficientami iz $K$}\quad{\footnotesize\texttt{TERM-1157}}\par
\noindent\emph{Iztočnik / Source:} mit Koeffizienten aus $K$\par
\noindent\emph{Rusky / Russian:} [Bukvy ne sut zachovane / Lettering not recoverable]; $K$\par
\noindent\emph{Ukrajinsky / Ukrainian:} [Bukvy ne sut zachovane / Lettering not recoverable]; $K$\par
\smallskip\noindent\emph{Razlog:} S koeficientami iz $K$ ostavjaje koeficientno polje vidimo. V §6 raboty 9 $K$ znači polje vsih algebraičnyh čisel; primětka objasnjava, kako by formulacije ostali važeče pri zaměně njegovoj poljem racionalnyh čisel. Zato racionalna funkcija ne znači tu nužno funkciju s racionalnymi čislovymi koeficientami. Arhivny razlog čuva imenno tu razliku. Bukvy v ruskoj i ukrajinskoj sravnitvě sut povređene; zachovany jest samo simbol $K$. Ih ne dopolnjujemo dogadkoju i ne račitajemo kako prověrjene jezyčne attestacije.\par
\noindent\emph{Reason:} S koeficientami iz $K$ keeps the coefficient field visible. In §6 of Paper 9, $K$ is the field of all algebraic numbers; a footnote explains how the statements would remain valid if it were replaced by the rational number field. Rational function here therefore does not necessarily mean a function with rational numerical coefficients. This is the distinction the archived reason preserves. The lettering in the Russian and Ukrainian comparison fields is damaged; only the symbol $K$ survives. Those fields are neither completed by guesswork nor counted as verified linguistic attestations.\par
\noindent{\footnotesize Izvorne zapisy / Original records: HIST-g064-002.}\par
\smallskip\noindent\emph{Sravnjenje s tekstom:} Definicija racionalnoj bazy v rabotě 9 sodrživaje koeficientno polje $K$. Primětka izrično razlikuje algebraične i racionalne čislove koeficienty i opisuje dovolenu zaměnu polja; tu razliku ne možno odčitati iz samogo pridevnika racionalna.\par
\noindent\emph{Textual comparison:} The definition of rational basis in Paper 9 includes the coefficient field $K$. Its footnote expressly distinguishes algebraic from rational numerical coefficients and describes the permitted field replacement; that distinction cannot be recovered from rational alone.\par
\noindent{\footnotesize TEXT-COMP-261; pasaže / passages: 7718, 7719, 7720, 7721, 7722, 7723, 7724, 7725, 7726, 7727, 7728, 7729, 7730, 7731, 7732; \texttt{CORPUS\_FORM\_COMPARISONS.json}.}\par
\end{minipage}\par
\par\medskip\Needspace{7\baselineskip}
\noindent\begin{minipage}{\linewidth}
\hypertarget{TERM-1158}{}
\noindent\textbf{predstavime s koeficientami iz $o$}\quad{\footnotesize\texttt{TERM-1158}}\par
\noindent\emph{Iztočnik / Source:} mit Koeffizienten aus $o$ darstellbar\par
\noindent\emph{Rusky / Russian:} выражаются с коэффициентами из $o$\par
\noindent\emph{Ukrajinsky / Ukrainian:} виражаються з коефіцієнтами з $o$\par
\smallskip\noindent\emph{Razlog:} Predstavime s koeficientami iz o čuva ne samo možnost prědstavjenja, ale i oblast dopuščenyh koeficientov. Arhivny zapis upotrěbjaje proste o; v §8 raboty 31 toj oblastju jest prsten $\mathfrak o$. Ne ide o čislu nula. Tekst nazyvaje vybrane determinanty bazoju diskriminantnogo ideala i tvrdi, že druge determinanty sut ih linearnymi kombinacijami s koeficientami iz $\mathfrak o$. Koeficienty iz bolšego polja ne byli by dostatočne dlja takej tvrdnje o idealu v $\mathfrak o$. Arhivny razlog čuva tu nužnost; samo jezyčno sravnjenje frazy ješče ne jest polnym dokazom tvrdnje o vybranyh determinantah.\par
\noindent\emph{Reason:} Predstavime s koeficientami iz o preserves not just expressibility but the allowed coefficient domain. The archive uses plain o; in §8 of Paper 31 the domain is the ring $\mathfrak o$, not the number zero. The text calls selected determinants a basis of the discriminant ideal and asserts that the other determinants are their linear combinations with coefficients in $\mathfrak o$. Coefficients in a larger field would not suffice for that assertion about an ideal in $\mathfrak o$. The archived reason preserves this requirement; comparison of the wording alone is not a full proof of the claim about the selected determinants.\par
\noindent{\footnotesize Izvorne zapisy / Original records: HIST-g158-009.}\par
\smallskip\noindent\emph{Sravnjenje s tekstom:} V §8 raboty 31 determinanty naležet $\mathfrak o$, i tvrdnja o konečnom generovanju žada linearne kombinacije nad tym samym prstenom. To jest točna informacija, ktoru čuva fraza s koeficientami iz o; stary proste o ne jest nula.\par
\noindent\emph{Textual comparison:} In §8 of Paper 31, the determinants belong to $\mathfrak o$, and the finite-generation assertion requires linear combinations over that same ring. This is the precise information preserved by s koeficientami iz o; the archived plain o is not zero.\par
\noindent{\footnotesize TEXT-COMP-262; pasaže / passages: 18895, 18897; \texttt{CORPUS\_FORM\_COMPARISONS.json}.}\par
\end{minipage}\par
\par\medskip\Needspace{7\baselineskip}
\noindent\begin{minipage}{\linewidth}
\hypertarget{TERM-1159}{}
\noindent\textbf{pomnoženyh konstantami}\quad{\footnotesize\texttt{TERM-1159}}\par
\noindent\emph{Iztočnik / Source:} mit Konstanten multipliziert\par
\noindent\emph{Rusky / Russian:} умноженных на постоянные\par
\noindent\emph{Ukrajinsky / Ukrainian:} помножених на сталі\par
\smallskip\noindent\emph{Razlog:} Pomnoženyh konstantami v §4 raboty 4 govori o skalarnyh množiteljah polarnyh operacij. Arhivny razlog sravnivaje medžunarodnu konstantu s ruskim postojannye i ukrajinskim stali; on ne tvrdi, že tam sut nove matematične konstanty s osobnymi nazvami. Tekst označaje s $\Pi$ agregat takyh operacij, čto dozvoljaje zapisati ranějše sumy jednoj formuloju. Množny oblik pomnoženyh mora ostati vezany k operacijam, a konstante ne možut byti zamenjene funkcijami proměnnyh, po ktoryh vrši se diferencijacija: to by izměnilo samy operator.\par
\noindent\emph{Reason:} Pomnoženyh konstantami in §4 of Paper 4 describes scalar multipliers of polar operations. The archived reason compares international konstanta with Russian postoyannye and Ukrainian stali; it does not claim that new named mathematical constants are introduced. The text denotes an aggregate of these operations by $\Pi$, allowing the earlier sums to be written in one formula. The plural participle must remain attached to the operations, and the constants cannot be replaced by functions of the differentiation variables: that would change the operator itself.\par
\noindent{\footnotesize Izvorne zapisy / Original records: HIST-g033-008.}\par
\smallskip\noindent\emph{Sravnjenje s tekstom:} Poslě izpisanja polarnyh operacij tekst raboty 4 naziva $\Pi$ jih agregatom s konstantnymi množiteljami. Formulacija govori o kombinaciji operacij, ne o zaměně koeficientov proizvoljnymi funkcijami.\par
\noindent\emph{Textual comparison:} After displaying the polar operations, Paper 4 names $\Pi$ as their aggregate with constant multipliers. The wording concerns a combination of operations, not replacement of the coefficients by arbitrary functions.\par
\noindent{\footnotesize TEXT-COMP-263; pasaže / passages: 4202, 4204, 4205, 4206, 4207, 4209, 4212, 4213, 4214, 4215, 4217; \texttt{CORPUS\_FORM\_COMPARISONS.json}.}\par
\end{minipage}\par
\par\medskip\Needspace{7\baselineskip}
\noindent\begin{minipage}{\linewidth}
\hypertarget{TERM-1162}{}
\noindent\textbf{modul}\quad{\footnotesize\texttt{TERM-1162}}\par
\noindent\emph{Iztočnik / Source:} Modul\par
\noindent\emph{Rusky / Russian:} модуль\par
\noindent\emph{Ukrajinsky / Ukrainian:} модуль\par
\smallskip\noindent\emph{Razlog:} Modul čuva obščy medžunarodny korenj, prisutny i v obojih slovjanskyh sravnjenjah. Dva zachovane razloga ne sut jednako podrobne: prvy prěduprědžaje prěd avtomatičnym modernym čitanjem, vtory samo naziva slovo algebraičnym terminom. V rabotě 1 moduly upravljajut postupnoju konstrukcijeju odnosno polnyh sistemov form; tekst měnjaje $(abc)$ na druge označene formy. V §9 raboty 4 konečno mnogo form zjedinjaje se v modul pri izslědovanju viših form. Eti konteksty opravdyvajut čuvanje slova, ale ne davajut osnovu pripisati vsěm pojavjenjam jednu pozdnějšu aksiomatičnu definiciju. Različenje smyslov potrěbno jest v slovniku, ne posrědstvom tihoj zaměny avtorovogo nazva.\par
\noindent\emph{Reason:} Modul retains the international stem shared by both Slavic comparanda. The two preserved reasons differ in detail: the first warns against an automatic modern reading, while the second merely calls it an algebraic term. In Paper 1 the moduli govern the successive construction of relatively complete systems of forms; the text replaces $(abc)$ by other specified forms. In §9 of Paper 4 finitely many forms are collected into a module in studying higher forms. These contexts justify retaining the word but do not justify assigning every occurrence a single later axiomatic definition. The dictionary must distinguish the senses rather than silently replace the author's terminology.\par
\noindent{\footnotesize Izvorne zapisy / Original records: HIST-g001-006, HIST-g038-007.}\par
\smallskip\noindent\emph{Sravnjenje s tekstom:} Raboty 1 i 2 opisyvajut postupnu zaměnu modulov pri odnosno polnyh sistemah. V §17 raboty 2 višina pripada etoj postupnosti, a ne dimenziji modula v sovrěmennoj algebraičnoj definiciji.\par
\noindent\emph{Textual comparison:} Papers 1 and 2 describe successive changes of moduli for relatively complete systems. In §17 of Paper 2, higher belongs to this succession, not to module dimension in a modern algebraic definition.\par
\noindent{\footnotesize TEXT-COMP-266; pasaže / passages: 71, 73, 77, 84, 182, 184, 188, 195, 2387, 2389, 2391, 2394, 2395; \texttt{CORPUS\_FORM\_COMPARISONS.json}.}\par
\smallskip\noindent\emph{Sravnjenje s tekstom:} V §9 konečno mnogo form zjedinjaje se v modul za daljše izslědovanje. Nazvanie prěmy algebraičny termin ne zaměnjaje matematičny opis etoj konstrukcije.\par
\noindent\emph{Textual comparison:} In §9 finitely many forms are collected into a module for the subsequent investigation. Calling it a direct algebraic term cannot replace the mathematical account of this construction.\par
\noindent{\footnotesize TEXT-COMP-267; pasaže / passages: 4910, 4912; \texttt{CORPUS\_FORM\_COMPARISONS.json}.}\par
\end{minipage}\par
\par\medskip\Needspace{7\baselineskip}
\noindent\begin{minipage}{\linewidth}
\hypertarget{TERM-1163}{}
\noindent\textbf{modul}\quad{\footnotesize\texttt{TERM-1163}}\par
\noindent\emph{Iztočnik / Source:} Modul (historical generated polynomial set)\par
\noindent\emph{Rusky / Russian:} модуль\par
\noindent\emph{Ukrajinsky / Ukrainian:} модуль\par
\smallskip\noindent\emph{Razlog:} V rabotě 45 modul označaje sistem, zadany polinomami; pojavjajut se izrazy $(\varphi,\psi,p^\rho)$ i $(\varphi,\psi,p^\sigma)$. Arhivny razlog prědpočita modul prěd avtomatičnoju zaměnoju na ideal, bo uvod izrično izbira izlaganje bez aparata teorije idealov. To ne znači, že sodržina nema idealno-teoretičnogo tumačenja. Avtorska granica takože jest važna: osnovny članok jest od H. Kapferera, a dodatok napisan jest vměstě s E. Noether; jezyk osnovnogo uvoda ne trěba pripisyvati lično Noether. Anglijsky opis v skobkah pri izvornom slově jest objašnjenje iz terminologičnogo zapisa, ne němečska citata.\par
\noindent\emph{Reason:} In Work 45, modul denotes a system specified by polynomials; expressions include $(\varphi,\psi,p^\rho)$ and $(\varphi,\psi,p^\sigma)$. The archived reason prefers modul to automatic replacement by ideal because the introduction explicitly chooses an exposition without the apparatus of ideal theory. This does not mean that the content lacks an ideal-theoretic interpretation. The authorship boundary also matters: the main article is by H. Kapferer, with a supplement written jointly with E. Noether; the language of its main introduction should not be attributed personally to Noether. The parenthetical English description in the source field is a terminology-record gloss, not a German quotation.\par
\noindent{\footnotesize Izvorne zapisy / Original records: HIST-g214-005.}\par
\smallskip\noindent\emph{Sravnjenje s tekstom:} Uvod izbira elementarno izlaganje bez idealno-teoretičnogo aparata. Nazv modul čuva etu prezentaciju; zaglavje pripisyvaje osnovnu rabotu H. Kapfereru, a dodatok spoločnomu avtorstvu s E. Noether.\par
\noindent\emph{Textual comparison:} The introduction chooses an elementary exposition without ideal-theoretic apparatus. Modul preserves that presentation; the title attributes the main work to H. Kapferer and the supplement to joint authorship with E. Noether.\par
\noindent{\footnotesize TEXT-COMP-268; pasaže / passages: 45-isv.tex:37, 45-isv.tex:38, 45-isv.tex:44, 45-isv.tex:45, 45-isv.tex:46, 45-isv.tex:47, 45-isv.tex:48, 45-isv.tex:49, 45-isv.tex:50, 45-isv.tex:53, 45-isv.tex:54, 45-isv.tex:55, 45-isv.tex:56, 45-isv.tex:57, 45-isv.tex:58, 45-isv.tex:59, 45-isv.tex:60; \texttt{CORPUS\_FORM\_COMPARISONS.json}.}\par
\end{minipage}\par
\par\medskip\Needspace{7\baselineskip}
\noindent\begin{minipage}{\linewidth}
\hypertarget{TERM-1164}{}
\noindent\textbf{modul / osnovna oblast cělosti}\quad{\footnotesize\texttt{TERM-1164}}\par
\noindent\emph{Iztočnik / Source:} Modul / Grundintegritätsbereich\par
\noindent\emph{Rusky / Russian:} модуль / основная область целостности\par
\noindent\emph{Ukrajinsky / Ukrainian:} модуль / базова область цілісності\par
\smallskip\noindent\emph{Razlog:} Para modul / osnovna oblast cělosti razděljuje objekt i oblast jego koeficientov. V rabotě 14 modul jest zamknuty pri odčitanju i pri množenju elementami osnovnoj oblasti $\mathfrak{S}$. Ako $\mathfrak{S}$ jest oblast cělyh čisel, druga potrěba slěduje iz prvoj; pri obščem $\mathfrak{S}$ ona neset dodatnu informaciju. Arhivny razlog upomina Dedekindovu tradiciju i slovjansky idiom oblasti cělosti. Jednak razlika od vektornogo prostora ne jest razlika od vsakego modernogo modula: koeficienty ne morajut tvoriti polje. Tekst posebno razlikuje moduly od idealov, čije elementy ležet v samoj oblasti $\mathfrak{S}$. Čuvanje obojih nazvov čini tu razliku vidimoju.\par
\noindent\emph{Reason:} The pair modul / osnovna oblast cělosti separates the object from its coefficient domain. In Paper 14 a module is closed under subtraction and multiplication by elements of the ground domain $\mathfrak{S}$. If $\mathfrak{S}$ is the integers, the second requirement follows from the first; for general $\mathfrak{S}$ it adds information. The archived reason invokes the Dedekind tradition and the Slavic expression for an integral domain. However, a distinction from vector spaces is not a distinction from every modern module: the coefficients need not form a field. The text separately identifies ideals, whose elements belong to $\mathfrak{S}$ itself. Retaining both names makes that distinction visible.\par
\noindent{\footnotesize Izvorne zapisy / Original records: HIST-g076-008.}\par
\smallskip\noindent\emph{Sravnjenje s tekstom:} Definicija imenuje odčitanje i množenje elementami $\mathfrak{S}$. Uslovje, že elementy samogo modula ležet v $\mathfrak{S}$, razlikuje idealny slučaj; vektorna nezavisnost ne jest častju toj definicije.\par
\noindent\emph{Textual comparison:} The definition names subtraction and multiplication by elements of $\mathfrak{S}$. Requiring the module's own elements to lie in $\mathfrak{S}$ distinguishes the ideal case; vector-space independence is not part of the definition.\par
\noindent{\footnotesize TEXT-COMP-269; pasaže / passages: 9883, 9885; \texttt{CORPUS\_FORM\_COMPARISONS.json}.}\par
\end{minipage}\par
\par\medskip\Needspace{7\baselineskip}
\noindent\begin{minipage}{\linewidth}
\hypertarget{TERM-1165}{}
\noindent\textbf{modul / vyši modul}\quad{\footnotesize\texttt{TERM-1165}}\par
\noindent\emph{Iztočnik / Source:} Modul / höherer Modul\par
\noindent\emph{Rusky / Russian:} модуль / более высокий модуль\par
\noindent\emph{Ukrajinsky / Ukrainian:} модуль / вищий модуль\par
\smallskip\noindent\emph{Razlog:} Vyši modul jest sravnitelny izraz iz postupka konstrukcije form, ne tvrdženje o višej dimenziji ili večšej kardinalnosti. V §17 raboty 2 iz sistema form $s$ prěhodi se k slědujučemu odnosno polnomu sistemu; redukcija izraza $\nu(s)=(ss'u)^4$ privodi k modulam $\varrho$ i $t$. Arhivny izbor čuva medžunarodny modul i sravnitelny oblik vyši, pri čem sam zapis ostavjaje morfologiju za daljše izslědovanje. Rusko i ukrajinsko sravnjenje podpirajut semantičnu strukturu, ale ne sut dokazom, že pravopis vyši byl jezykově potvrđeny. Slovnik zato razděljuje matematičny smysl sravnjenja i jezykovu ocenku oblika.\par
\noindent\emph{Reason:} Vyši modul is a comparative expression within the construction of forms, not a claim of higher dimension or greater cardinality. In §17 of Paper 2 the system of forms $s$ leads to the next relatively complete system; reduction of $\nu(s)=(ss'u)^4$ leads to the moduli $\varrho$ and $t$. The archived choice retains international modul and renders the comparative with vyši, while expressly leaving its morphology open to further examination. The Russian and Ukrainian comparison supports the semantic structure but does not establish that the spelling vyši received linguistic validation. The dictionary therefore separates the mathematical meaning of the comparison from assessment of the word form.\par
\noindent{\footnotesize Izvorne zapisy / Original records: HIST-g018-004.}\par
\smallskip\noindent\emph{Sravnjenje s tekstom:} Raboty 1 i 2 opisyvajut postupnu zaměnu modulov pri odnosno polnyh sistemah. V §17 raboty 2 višina pripada etoj postupnosti, a ne dimenziji modula v sovrěmennoj algebraičnoj definiciji.\par
\noindent\emph{Textual comparison:} Papers 1 and 2 describe successive changes of moduli for relatively complete systems. In §17 of Paper 2, higher belongs to this succession, not to module dimension in a modern algebraic definition.\par
\noindent{\footnotesize TEXT-COMP-266; pasaže / passages: 71, 73, 77, 84, 182, 184, 188, 195, 2387, 2389, 2391, 2394, 2395; \texttt{CORPUS\_FORM\_COMPARISONS.json}.}\par
\end{minipage}\par
\par\medskip\Needspace{7\baselineskip}
\noindent\begin{minipage}{\linewidth}
\hypertarget{TERM-1166}{}
\noindent\textbf{modul / moduly}\quad{\footnotesize\texttt{TERM-1166}}\par
\noindent\emph{Iztočnik / Source:} Modul / Moduln\par
\noindent\emph{Rusky / Russian:} модуль / модули\par
\noindent\emph{Ukrajinsky / Ukrainian:} модуль / модулі\par
\smallskip\noindent\emph{Razlog:} Modul / moduly čuva jedinstveno i množestveno čislo němečskogo Modul / Moduln. V uvodu raboty 2 različne formy postupno služet modulami pri konstrukciji odnosno polnyh sistemov; množina ne jest nazvom nove vrste algebraičnogo objekta. Arhivny razlog izbira prěmy kognat, no jego slovo o neizběžnosti togo izbora jest ocěna, ne dokaz nedostatka alternativ. Smyslova granica jest konkretna: odnosi medžu modulami v etom postupku ne smějut byti zamenjene obščim opisom modernoj teorije modulov. Oba čislove obliki zato ostajut pod jednoj lemoju s kontekstualnoju definicijeju.\par
\noindent\emph{Reason:} Modul / moduly preserves the singular and plural of German Modul / Moduln. In the introduction to Paper 2, different forms successively serve as moduli in constructing relatively complete systems; the plural does not name a new kind of algebraic object. The archived reason selects the direct cognate, but its claim that this choice is unavoidable is an assessment, not evidence that no alternatives exist. The semantic boundary is specific: relations between moduli in this procedure must not be replaced by a generic account of modern module theory. The two number forms therefore belong under one headword with a contextual definition.\par
\noindent{\footnotesize Izvorne zapisy / Original records: HIST-g002-005.}\par
\smallskip\noindent\emph{Sravnjenje s tekstom:} Raboty 1 i 2 opisyvajut postupnu zaměnu modulov pri odnosno polnyh sistemah. V §17 raboty 2 višina pripada etoj postupnosti, a ne dimenziji modula v sovrěmennoj algebraičnoj definiciji.\par
\noindent\emph{Textual comparison:} Papers 1 and 2 describe successive changes of moduli for relatively complete systems. In §17 of Paper 2, higher belongs to this succession, not to module dimension in a modern algebraic definition.\par
\noindent{\footnotesize TEXT-COMP-266; pasaže / passages: 71, 73, 77, 84, 182, 184, 188, 195, 2387, 2389, 2391, 2394, 2395; \texttt{CORPUS\_FORM\_COMPARISONS.json}.}\par
\end{minipage}\par
\par\medskip\Needspace{7\baselineskip}
\noindent\begin{minipage}{\linewidth}
\hypertarget{TERM-1167}{}
\noindent\textbf{modul vseh identičnyh relacij}\quad{\footnotesize\texttt{TERM-1167}}\par
\noindent\emph{Iztočnik / Source:} Modul aller identischen Relationen\par
\noindent\emph{Rusky / Russian:} модуль всех тождественных соотношений\par
\noindent\emph{Ukrajinsky / Ukrainian:} модуль усіх тотожних співвідношень\par
\smallskip\noindent\emph{Razlog:} Modul vseh identičnyh relacij čuva i avtorovo modul i točno potrěbu identiteta. V rabotě 16 polinomy $G(t)$ pripadajut tomu modulu, ako poslě zaměny $t$ vyrazami $p$ rezultat $G(p)$ jest identično nula v prěměnnyh $x,y$. Identične zato ne znači, že vsě relacije sut jedna i ta sama relacija. Arhivny razlog prědpočita ta slova prěd modernizacijeju na ideal, v soglasnosti s izlaganjem posrědstvom sistemov modulov. Dnešnje objašnjenje kako jadro podstanovky može pomoči čitatelju, ale jest objašnjenjem, ne zamenoj avtorovogo termina. V prověrenom tekstu stojí oblik vsih; slovnik čuva razliku od arhivnogo vseh.\par
\noindent\emph{Reason:} Modul vseh identičnyh relacij retains both the author's modul and the requirement that a relation hold identically. In Paper 16 the polynomials $G(t)$ belong to this module when substitution of the expressions $p$ for $t$ makes $G(p)$ identically zero in the variables $x,y$. Identične therefore does not mean that every relation is the very same relation. The archived reason prefers this wording to modernization as ideal, consistent with the exposition through systems of modules. A present-day explanation as the kernel of substitution can help the reader, but is an explanation rather than a replacement for the author's term. The checked text has vsih; the dictionary preserves its distinction from archived vseh.\par
\noindent{\footnotesize Izvorne zapisy / Original records: HIST-g078-007.}\par
\smallskip\noindent\emph{Sravnjenje s tekstom:} Uslovje $G(p)=0$ jest identitetom poslě podstanovky, ne tvrdženjem, že različne polinomialne relacije sut identične jedna drugoj. Tekst čuva modul kako nazv sistema takyh relacij.\par
\noindent\emph{Textual comparison:} The condition $G(p)=0$ is an identity after substitution, not a claim that different polynomial relations are identical to one another. The text retains module as the name of this system of relations.\par
\noindent{\footnotesize TEXT-COMP-270; pasaže / passages: 10737, 10742, 10743, 10744, 10745, 10746, 10747, 10767; \texttt{CORPUS\_FORM\_COMPARISONS.json}.}\par
\end{minipage}\par
\par\medskip\Needspace{7\baselineskip}
\noindent\begin{minipage}{\linewidth}
\hypertarget{TERM-1168}{}
\noindent\textbf{modul iz linearnyh form}\quad{\footnotesize\texttt{TERM-1168}}\par
\noindent\emph{Iztočnik / Source:} Modul aus Linearformen\par
\noindent\emph{Rusky / Russian:} модуль из линейных форм\par
\noindent\emph{Ukrajinsky / Ukrainian:} модуль з лінійних форм\par
\smallskip\noindent\emph{Razlog:} Modul iz linearnyh form jest prěmy opis elementov modula. V §1 raboty 22 formy sut linearne v konečnom spisku prěměnnyh $\xi_i$, dokoliko jih koeficienty mogut byti cělymi čislami ili polinomami jednoj drugoj prěměnnoj nad poljem $P$. Slovo linearnyh zato ne ograničuje stepen koeficientnyh polinomov. Arhivny razlog želi razlikovati etu konstrukciju od polinomialnyh idealov uvedenyh v §3; to jest razlika ulog v izlaganju, ne zabrana opisati někoje ideale kako moduly. Prědlog iz i genitiv form prědavaju sastav objekta bez dodatnogo teoretičnogo nazva.\par
\noindent\emph{Reason:} Modul iz linearnyh form directly describes the elements of the module. In §1 of Paper 22 the forms are linear in a finite list of variables $\xi_i$, while their coefficients may be integers or polynomials in a different single variable over a field $P$. Linearnyh therefore does not bound the degree of the coefficient polynomials. The archived reason seeks to distinguish this construction from polynomial ideals introduced in §3; this is a distinction between roles in the exposition, not a prohibition on describing some ideals as modules. The preposition iz and genitive form describe the object's constituents without adding another theoretical label.\par
\noindent{\footnotesize Izvorne zapisy / Original records: HIST-g085-002, HIST-g092-002.}\par
\smallskip\noindent\emph{Sravnjenje s tekstom:} Formy sut linearne v $\xi_i$, a koeficienty imajut svoj odvojeno oprěděljany obseg. Tekst dozvoljaje polinomialne koeficienty; linearny zato ne znači stepen najvyše jeden vo vsěh prisutnyh prěměnnyh.\par
\noindent\emph{Textual comparison:} The forms are linear in $\xi_i$, while the coefficients have a separately specified domain. Polynomial coefficients are allowed; linear therefore does not mean degree at most one in every variable present.\par
\noindent{\footnotesize TEXT-COMP-271; pasaže / passages: 14224, 14226, 14228, 14230, 14231; \texttt{CORPUS\_FORM\_COMPARISONS.json}.}\par
\end{minipage}\par
\par\medskip\Needspace{7\baselineskip}
\noindent\begin{minipage}{\linewidth}
\hypertarget{TERM-1169}{}
\noindent\textbf{modul v $(\Sigma,T)$}\quad{\footnotesize\texttt{TERM-1169}}\par
\noindent\emph{Iztočnik / Source:} Modul in $(\Sigma,T)$\par
\noindent\emph{Rusky / Russian:} модуль в $(\Sigma,T)$\par
\noindent\emph{Ukrajinsky / Ukrainian:} модуль у $(\Sigma,T)$\par
\smallskip\noindent\emph{Razlog:} Modul v $(\Sigma,T)$ čuva dve razne ulogi: $\Sigma$ dava skalary, a $T$ jest aditivna grupa, v kojoj ležet elementy modula. V §9 raboty 19 podmnožstvo $T$ mora byti zamknuto pri odčitanju i pri dejstvu skalarov. Arhivny razlog pravilno podčrkuje obobčenje idealov; to ne jest para dvoh modulov. Poradok množenja jest suščestveny, bo $\Sigma$ ne mora byti komutativny. Prověreny tekst piše $c\alpha$, ale v primětci naziva toj slučaj pravym modulom. Slovnik čuva to dejstvo i opisuje konvenciju, ne opravjaje je po navikah sovrěmennogo čitatelja. Sravnjenje s pečatnym originalom togo označenja ostaje potrebno.\par
\noindent\emph{Reason:} Modul v $(\Sigma,T)$ preserves two different roles: $\Sigma$ supplies scalars, while $T$ is the additive group containing the module elements. In §9 of Paper 19 a subset of $T$ must be closed under subtraction and scalar action. The archived reason correctly stresses the generalization of ideals; the notation is not a pair of modules. Multiplication order matters because $\Sigma$ need not be commutative. The checked text writes $c\alpha$ but calls this the right-module case in a footnote. The dictionary preserves the action and describes that convention rather than correcting it to a present-day reader's expectations. Comparison of this designation with the printed original remains necessary.\par
\noindent{\footnotesize Izvorne zapisy / Original records: HIST-g081-070.}\par
\smallskip\noindent\emph{Sravnjenje s tekstom:} Definicija jest v §9. Tekst piše $c\alpha$ i v primětci naziva toj slučaj pravym; etu kombinaciju ne možno tiho obrnuti. Ako nema jedinice, zapis generatorov vključaje takože cěločislne množitelji. Pečatny svědok nazvanija strany ne byl prověren v etom sravnjenju.\par
\noindent\emph{Textual comparison:} The definition is in §9. The text writes $c\alpha$ and calls this the right-hand case in a footnote; that combination must not be silently reversed. Without an identity, the expression in generators also includes integer multipliers. The printed witness for the side designation was not checked in this comparison.\par
\noindent{\footnotesize TEXT-COMP-272; pasaže / passages: 13360, 13362, 13364, 13366, 13370, 13372, 13374, 13376, 13383, 13385, 13386, 13387, 13388, 13390, 13398, 13399, 13400, 13401; \texttt{CORPUS\_FORM\_COMPARISONS.json}.}\par
\end{minipage}\par
\par\medskip\Needspace{7\baselineskip}
\noindent\begin{minipage}{\linewidth}
\hypertarget{TERM-1170}{}
\noindent\textbf{modul $M_\sigma$ / modulno svojstvo}\quad{\footnotesize\texttt{TERM-1170}}\par
\noindent\emph{Iztočnik / Source:} Modul $M_\sigma$ / Moduleigenschaft\par
\noindent\emph{Rusky / Russian:} модуль $M_\sigma$ / модульное свойство\par
\noindent\emph{Ukrajinsky / Ukrainian:} модуль $M_\sigma$ / модульна властивість\par
\smallskip\noindent\emph{Razlog:} Modul $M_\sigma$ i modulno svojstvo sut vezane k konkretnomu argumentu v §9 raboty 9. V tekste modul, zapisan frakturnym $\mathfrak{M}_\sigma$, porađajut monomy stepena $\sigma$, s dozvoljenymi množiteljami iz oblasti cělyh algebraičnyh funkcij. Dokaz potrěbuje, že izomorfizm oblastij prěnosi toj modul v sootvětny modul drugoj oblasti; sama bijekcija elementov ne vyjadřaje vsego potrěbnogo. Arhivny razlog prědpočita prěmu terminologiju modula i čuvanje jeje veze s §3. Zapis $M_\sigma$ v lematě jest prostějši arhivny zapis, ne utvrđženje, že v avtorovoj notaciji stojí obyčno latinske M.\par
\noindent\emph{Reason:} Modul $M_\sigma$ and modulno svojstvo belong to a particular argument in §9 of Paper 9. In the text the module, written with fraktur $\mathfrak{M}_\sigma$, is generated by monomials of degree $\sigma$, with permitted multipliers from the domain of integral algebraic functions. The proof requires an isomorphism of the domains to carry this module to its counterpart in the other domain; a mere bijection of elements does not express the whole requirement. The archived reason favours direct module terminology and continuity with §3. The headword's $M_\sigma$ is the simpler archived notation, not a claim that the author's symbol is an ordinary Latin M.\par
\noindent{\footnotesize Izvorne zapisy / Original records: HIST-g067-004.}\par
\smallskip\noindent\emph{Sravnjenje s tekstom:} Modul $\mathfrak{M}_\sigma$ porađajut označene monomy. Dokaz upotrěbjaje jego zamknutost pri množenju dozvoljenymi funkcijami, da izvesti sootvětstvo modulov. To svojstvo jest bolše než samo aditivno sravnjenje dvuh oblastij.\par
\noindent\emph{Textual comparison:} The module $\mathfrak{M}_\sigma$ is generated by the specified monomials. The proof uses closure under multiplication by the permitted functions to derive correspondence of the modules. This property is more than an additive comparison of the two domains.\par
\noindent{\footnotesize TEXT-COMP-273; pasaže / passages: 8186, 8188, 8190, 8192, 8193, 8195, 8198, 8207, 8211, 8213; \texttt{CORPUS\_FORM\_COMPARISONS.json}.}\par
\end{minipage}\par
\par\medskip\Needspace{7\baselineskip}
\noindent\begin{minipage}{\linewidth}
\hypertarget{TERM-1171}{}
\noindent\textbf{modulne bazy}\quad{\footnotesize\texttt{TERM-1171}}\par
\noindent\emph{Iztočnik / Source:} Modulbasen\par
\noindent\emph{Rusky / Russian:} модульные базы\par
\noindent\emph{Ukrajinsky / Ukrainian:} модульні бази\par
\smallskip\noindent\emph{Razlog:} Modulne bazy jest množina izraza modulna baza; arhivny razlog prědpočita ju težšemu opisu kako bazisny sistem. V §4,2 raboty 31 različne bazy jednogo i togo že ideala davajut različne matricne zapisy, ale istu klasu ekvivalentnyh prědstavjenij. Sootvětny izbor baz pri izomorfnyh idealah može dati identične matrice. Zato slovnik ne směje skratiti rezultat do tvrdženja, že matrice ne zaviset od bazy. Nezavisna od izbora jest klasa, a ne vsaka matrica. Slovjansko soglasje pridevnika s množinoju bazy čuva običnu gramatičnu strukturu prěmogo kalka.\par
\noindent\emph{Reason:} Modulne bazy is the plural of modulna baza; the archived reason prefers it to a heavier expression such as basis system. In §4,2 of Paper 31 different bases of the same ideal give different matrix expressions but the same class of equivalent representations. Corresponding choices of bases for isomorphic ideals can give identical matrices. The dictionary must therefore not shorten the result to a claim that the matrices are independent of the basis. The class is independent of that choice, not each matrix. Agreement of the adjective with plural bazy preserves the ordinary grammatical structure of the direct calque.\par
\noindent{\footnotesize Izvorne zapisy / Original records: HIST-g134-002.}\par
\smallskip\noindent\emph{Sravnjenje s tekstom:} Zaměna bazy ne ostavjaje každu matricu neizměnnoju; ona ostavjaje klasu ekvivalentnosti. Sootvětstvo baz izomorfnyh idealov dava točniji rezultat o identičnyh matricah.\par
\noindent\emph{Textual comparison:} A change of basis does not leave each matrix unchanged; it preserves the equivalence class. Corresponding bases of isomorphic ideals give the more specific result about identical matrices.\par
\noindent{\footnotesize TEXT-COMP-274; pasaže / passages: 17950, 17985, 17989, 17993, 17995; \texttt{CORPUS\_FORM\_COMPARISONS.json}.}\par
\end{minipage}\par
\par\medskip\Needspace{7\baselineskip}
\noindent\begin{minipage}{\linewidth}
\hypertarget{TERM-1172}{}
\noindent\textbf{modulna baza}\quad{\footnotesize\texttt{TERM-1172}}\par
\noindent\emph{Iztočnik / Source:} Modulbasis\par
\noindent\emph{Rusky / Russian:} модульная база; модульный базис\par
\noindent\emph{Ukrajinsky / Ukrainian:} модульна база\par
\smallskip\noindent\emph{Razlog:} Modulna baza čuva němečske složenje kako pridevnik s imenikom; arhivne razloge združajut Hilbertovu teoriju konečnyh generatorov, konstrukciju iz §3 raboty 9 i bazy iz raboty 31. To ne dava jednu potrěbu nezavisnosti za vsě slučaje. V rabotě 6 važna jest možnost vyraziti formy kombinacijami konečnogo spiska; v rabotě 9 slova o roli prěměnnyh kako bazy ne tvrđet samy po sebě jih linearnoj nezavisnosti. Naprotiv v §2 raboty 31 nezavisnost nad poljem jest izrična, a v §3 posebne komponenty jedinice tvoret nezavisnu bazu razširenogo prstena. Zato baza mora iměti kontekstualny opis: generatory, ili generatory s dokazanoju nezavisnostju.\par
\noindent\emph{Reason:} Modulna baza renders the German compound as an adjective and noun. The archived reasons bring together Hilbert's finite-generation theory, the construction in §3 of Paper 9, and the bases in Paper 31. This does not establish one independence requirement for every case. In Paper 6 the important property is expression of forms as combinations of a finite list; in Paper 9 the statement that indeterminates play the role of a basis does not itself assert their linear independence. By contrast, §2 of Paper 31 expressly requires independence over a field, and in §3 particular components of the identity form an independent basis of the extended ring. Baza therefore needs a contextual description: generators, or generators with established independence.\par
\noindent{\footnotesize Izvorne zapisy / Original records: HIST-g040-004, HIST-g061-004, HIST-g121-002, HIST-g129-003.}\par
\smallskip\noindent\emph{Sravnjenje s tekstom:} V uvodu raboty 6 konečne formy služe za vyraženje ostalyh form kombinacijami s koeficientami oprěděljenogo vida. Sama upotrěba baza ne ustanavjaje nezavisnost generatorov.\par
\noindent\emph{Textual comparison:} In the introduction to Paper 6 finitely many forms express the remaining forms through combinations with coefficients of a specified kind. The use of basis alone does not establish independence of the generators.\par
\noindent{\footnotesize TEXT-COMP-275; pasaže / passages: 5007, 5013, 5015, 5017, 5019, 5021; \texttt{CORPUS\_FORM\_COMPARISONS.json}.}\par
\smallskip\noindent\emph{Sravnjenje s tekstom:} Rabota 9 prěnosi opis s prěměnnymi k cělym algebraičnym funkcijam. Zapis $z=\alpha+\sum g_i\eta_i$ i primětka o svojstvu $z-\alpha$ davajut konkretny smysl slovam modulna oblast i baza; sam nazv ne dokazyvaje nezavisnosti $\eta_i$.\par
\noindent\emph{Textual comparison:} Paper 9 passes from a description with indeterminates to integral algebraic functions. The expression $z=\alpha+\sum g_i\eta_i$ and the footnote on the property of $z-\alpha$ give concrete content to module domain and basis; the name alone does not prove independence of the $\eta_i$.\par
\noindent{\footnotesize TEXT-COMP-276; pasaže / passages: 7264, 7266, 7268, 7270, 7273, 7277, 7284; \texttt{CORPUS\_FORM\_COMPARISONS.json}.}\par
\smallskip\noindent\emph{Sravnjenje s tekstom:} V rabotě 31 nezavisnost v etyh městah jest izrična. To oprěděljaje silnějši smysl bazy než samo porađanje iz konečnogo spiska v drugih rabotah.\par
\noindent\emph{Textual comparison:} In these passages of Paper 31, independence is explicit. It establishes a stronger meaning of basis than generation by a finite list in other papers.\par
\noindent{\footnotesize TEXT-COMP-277; pasaže / passages: 17434, 17785, 17786, 17787, 17788; \texttt{CORPUS\_FORM\_COMPARISONS.json}.}\par
\end{minipage}\par
\par\medskip\Needspace{7\baselineskip}
\noindent\begin{minipage}{\linewidth}
\hypertarget{TERM-1173}{}
\noindent\textbf{bazis modula / konečny bazis modula}\quad{\footnotesize\texttt{TERM-1173}}\par
\noindent\emph{Iztočnik / Source:} Modulbasis / endliche Modulbasis\par
\noindent\emph{Rusky / Russian:} базис модуля / конечный базис модуля\par
\noindent\emph{Ukrajinsky / Ukrainian:} базис модуля / скінченний базис модуля\par
\smallskip\noindent\emph{Razlog:} Bazis modula jest imenna konstrukcija s genitivom, v sravnjenju s pridevnikovym modulna baza. Arhivny razlog ne dava čislove ocěnky dvoh oblikov; on objašnjaje, začto čuvati slovo bazis: v §6 raboty 17 konečno mnogo izbranyh elementov dozvoljaje predstaviti cěly sistem linearnymi kombinacijami. Konečny odnosi se k spisku generatorov, ne k čislu elementov v oblasti. V prověrenom tekste stojí modulny bazis. Množiteli $A_{ij}$ v zapise $F_i=\sum A_{ij}G_j$ stojat s lěvoj strany; uslovja na nekomutativne koeficientne zakony ne smějut byti izgubjene. Toj konečny bazis ne jest avtomatično nezavisnym bazisom proizvoljnogo modula.\par
\noindent\emph{Reason:} Bazis modula is a noun with a genitive complement, as opposed to adjectival modulna baza. The archived reason supplies no numerical ranking of these forms; it explains why basis is retained: in §6 of Paper 17 finitely many selected elements express the whole system by linear combinations. Konečny concerns the list of generators, not the number of elements in the ambient domain. The checked text has modulny bazis. The multipliers $A_{ij}$ in $F_i=\sum A_{ij}G_j$ stand on the left; the hypotheses governing the noncommutative coefficient laws must not disappear. This finite basis is not automatically an independent basis of an arbitrary module.\par
\noindent{\footnotesize Izvorne zapisy / Original records: HIST-g079-033.}\par
\smallskip\noindent\emph{Sravnjenje s tekstom:} V §6 raboty 17 zapis $F_i=\sum A_{ij}G_j$ oprěděljaje ulogu konečnogo bazisa. Lokalni zakony za koeficienty sut častju uslovij; to ne jest bezuslovny teorem o proizvoljnom nekomutativnom prstenu.\par
\noindent\emph{Textual comparison:} In §6 of Paper 17, $F_i=\sum A_{ij}G_j$ specifies the role of a finite basis. The local coefficient laws are part of the hypotheses; this is not an unconditional theorem about an arbitrary noncommutative ring.\par
\noindent{\footnotesize TEXT-COMP-278; pasaže / passages: 11391, 11393, 11395, 11397, 11399, 11401, 11403; \texttt{CORPUS\_FORM\_COMPARISONS.json}.}\par
\end{minipage}\par
\par\medskip\Needspace{7\baselineskip}
\noindent\begin{minipage}{\linewidth}
\hypertarget{TERM-1174}{}
\noindent\textbf{modulna baza / fundamentalna forma / fundamentalno uravnenje}\quad{\footnotesize\texttt{TERM-1174}}\par
\noindent\emph{Iztočnik / Source:} Modulbasis / Fundamentalform / Fundamentalgleichung\par
\noindent\emph{Rusky / Russian:} модульная база / фундаментальная форма / фундаментальное уравнение\par
\noindent\emph{Ukrajinsky / Ukrainian:} модульна база / фундаментальна форма / фундаментальне рівняння\par
\smallskip\noindent\emph{Razlog:} Modulna baza, fundamentalna forma i fundamentalno uravnenje sut tri povezane, ale ne sinonimne nazvy. Uvod raboty 31 objašnjaje prěhod od stepenej primitivnogo elementa i jego oprěděljajučego uravnenja k bazě s relacijami; analogično govori o stepeneh fundamentalnoj formy i jeje uravnenju. Arhivny razlog bira prěmy modulny naziv i medžunarodne fundamentalna forma / fundamentalno uravnenje, čuvajuči strukturu argumenta. Motivacija jest matematična: rabota s bazami dozvoljaje redukciju po prostom modulu v uslovjah, v ktoryh izbor primitivnogo elementa može byti neprigodny. Iz etogo ne slěduje, že vsak algebraičny poredok ima bazu iz stepenej jednogo elementa.\par
\noindent\emph{Reason:} Modulna baza, fundamentalna forma and fundamentalno uravnenje are three related names, not synonyms. The introduction to Paper 31 explains the move from powers of a primitive element and its defining equation to a basis with relations; it also speaks of powers of the fundamental form and its equation. The archived reason chooses direct module terminology and the international expressions fundamental form / fundamental equation, preserving the argument's structure. The motivation is mathematical: working with bases permits reduction modulo a prime under circumstances in which a primitive-element choice can be unsuitable. This does not imply that every algebraic order has a basis consisting of powers of one element.\par
\noindent{\footnotesize Izvorne zapisy / Original records: HIST-g118-013.}\par
\smallskip\noindent\emph{Sravnjenje s tekstom:} Uvod raboty 31 sravnivaje različne sposoby zadanja algebraičnogo objekta i objašnjaje prědnost bazy pri redukciji. Fundamentalna forma, jeje uravnenje i modulna baza imajut različne ulogi v tom sravnjenju.\par
\noindent\emph{Textual comparison:} The introduction to Paper 31 compares ways of specifying the algebraic object and explains the advantage of a basis for reduction. The fundamental form, its equation and the module basis play different roles in that comparison.\par
\noindent{\footnotesize TEXT-COMP-279; pasaže / passages: 17279, 17281; \texttt{CORPUS\_FORM\_COMPARISONS.json}.}\par
\end{minipage}\par
\par\medskip\Needspace{7\baselineskip}
\noindent\begin{minipage}{\linewidth}
\hypertarget{TERM-1175}{}
\noindent\textbf{modulna baza / linearno nezavisna modulna baza}\quad{\footnotesize\texttt{TERM-1175}}\par
\noindent\emph{Iztočnik / Source:} Modulbasis / linear unabhängige Modulbasis\par
\noindent\emph{Rusky / Russian:} модульная база / линейно независимая модульная база\par
\noindent\emph{Ukrajinsky / Ukrainian:} модульна база / лінійно незалежна модульна база\par
\smallskip\noindent\emph{Razlog:} Dodatok linearno nezavisna ne jest stilistična izbytočnost. Na početku §4 raboty 31 izrično potrěbuje se, že vsak $\mathfrak{Z}$-modul v oblasti $\mathfrak{R}$, uključajuči samu $\mathfrak{R}$, ima konečnu nezavisnu bazu. Zapis $\mathfrak{Z}$ označaje osnovnu oblast, ne avtomatično cěle čisla. Arhivny razlog čuva etu formulaciju vměsto vvedenja novogo nazva svobodny modul. Tak čitatelj vidi točno uslovje, ktoro nosi daljši argument o zaměně baz i determinantu. Bez lokalnyh uslovij ne možno prěnesti toj rezultat na vsě moduly nad vsěmi prstenami.\par
\noindent\emph{Reason:} The addition linearno nezavisna is not stylistic redundancy. At the beginning of §4 of Paper 31 the explicit requirement is that every $\mathfrak{Z}$-module in $\mathfrak{R}$, including $\mathfrak{R}$ itself, have a finite independent basis. The symbol $\mathfrak{Z}$ denotes the ground domain, not automatically the integers. The archived reason retains this formulation rather than introduce the additional name free module. The reader can thus see the exact hypothesis supporting the subsequent argument about changes of basis and their determinant. Without the local hypotheses, the result cannot be transferred to all modules over all rings.\par
\noindent{\footnotesize Izvorne zapisy / Original records: HIST-g130-007.}\par
\smallskip\noindent\emph{Sravnjenje s tekstom:} Uslovje govori o vsěh $\mathfrak{Z}$-modulah v danoj $\mathfrak{R}$, ne o vsěh modulah bez ograničenja. Bazovy symbol ne označaje avtomatično prsten cělyh čisel.\par
\noindent\emph{Textual comparison:} The hypothesis concerns all $\mathfrak{Z}$-modules in the given $\mathfrak{R}$, not every module without restriction. The ground-domain symbol does not automatically denote the ring of integers.\par
\noindent{\footnotesize TEXT-COMP-280; pasaže / passages: 17823, 17827, 17829; \texttt{CORPUS\_FORM\_COMPARISONS.json}.}\par
\end{minipage}\par
\par\medskip\Needspace{7\baselineskip}
\noindent\begin{minipage}{\linewidth}
\hypertarget{TERM-1176}{}
\noindent\textbf{modulna baza iz linearno nezavisnyh elementov}\quad{\footnotesize\texttt{TERM-1176}}\par
\noindent\emph{Iztočnik / Source:} Modulbasis aus linear unabhängigen Elementen\par
\noindent\emph{Rusky / Russian:} модульный базис из линейно независимых элементов\par
\noindent\emph{Ukrajinsky / Ukrainian:} модульний базис із лінійно незалежних елементів\par
\smallskip\noindent\emph{Razlog:} Modulna baza iz linearno nezavisnyh elementov čuva polny uslovny izraz: slova iz linearno nezavisnyh elementov ne možno izpustiti pod prědlogom, že vsaka baza už po definiciji jest nezavisna. V §8,4 raboty 31 tvrđženje govori o lokalizovanom prstenu $\mathfrak{T}_{\mathfrak{P}}$ kako konečnom poredku nad glavnym prstenom $\mathfrak{o}_{\mathfrak{p}}$, s vezju $\mathfrak{P}=\mathfrak{T}\mathfrak{p}$. Arhivny razlog izrično čuva nezavisnost. To jest baza nad označenym koeficientnym prstenom, ne proizvoljny spisok elementov. Primětka citira H. Grella; konsultacija njegovoj raboty ne jest zasvědčena tym, že citata prisutstvuje v perevodě.\par
\noindent\emph{Reason:} Modulna baza iz linearno nezavisnyh elementov retains the full condition: from linearly independent elements cannot be omitted on the assumption that every historical use of basis already entails independence. In §8,4 of Paper 31 the assertion concerns the localized ring $\mathfrak{T}_{\mathfrak{P}}$ as a finite order over the principal ideal ring $\mathfrak{o}_{\mathfrak{p}}$, with $\mathfrak{P}=\mathfrak{T}\mathfrak{p}$. The archived reason explicitly retains independence. This is a basis over the specified coefficient ring, not an arbitrary list of elements. A footnote cites H. Grell; the presence of that citation in the translation does not establish consultation of his work.\par
\noindent{\footnotesize Izvorne zapisy / Original records: HIST-g161-003.}\par
\smallskip\noindent\emph{Sravnjenje s tekstom:} Definicija prstena kvocientov upotrěbjaje dozvoljene imenovatelje, ne klasy ostatkov. Baza iz nezavisnyh elementov pripada konkretnoj lokalnoj konstrukciji nad $\mathfrak{o}_{\mathfrak{p}}$. Glavny prsten i konečnost poredka sut matematične uslovja.\par
\noindent\emph{Textual comparison:} The definition of the ring of quotients uses permitted denominators, not residue classes. The basis of independent elements belongs to a particular local construction over $\mathfrak{o}_{\mathfrak{p}}$. The principal ideal ring and finite-order conditions are mathematical hypotheses.\par
\noindent{\footnotesize TEXT-COMP-281; pasaže / passages: 18932, 19003; \texttt{CORPUS\_FORM\_COMPARISONS.json}.}\par
\end{minipage}\par
\par\medskip\Needspace{7\baselineskip}
\noindent\begin{minipage}{\linewidth}
\hypertarget{TERM-1177}{}
\noindent\textbf{modulna oblast}\quad{\footnotesize\texttt{TERM-1177}}\par
\noindent\emph{Iztočnik / Source:} Modulbereich\par
\noindent\emph{Rusky / Russian:} модульная область\par
\noindent\emph{Ukrajinsky / Ukrainian:} модульна область\par
\smallskip\noindent\emph{Razlog:} Modulna oblast čuva istoricky obraz izraza Modulbereich i soglasnost medžu rabotami 9 i 30. Arhivni razloge protivjajut se neobjašnjenoj modernizaciji, ale ne trěba iz njih izvesti, že objekt nikogda ne jest modulom v dnešnjem smyslu. V §2 raboty 30 oblast jest aditivna Abelova grupa s dejstvom elementov prstena, a tekst izrično dava distributivne i asociativne zakony; v §10 analizuje jeje lance i kompozicionne redy. To ne jest prosto kolekcija modulov. V §3 raboty 9 nazva pripada specifičnoj konstrukciji s algebraičnymi funkcijami. Jednaky korenj podpiraje jezykovu kontinuitetnost, no vsaka definicija mora byti čitana v svojem městu.\par
\noindent\emph{Reason:} Modulna oblast preserves the historical expression Modulbereich and continuity between Papers 9 and 30. The archived reasons oppose unexplained modernization, but they should not be taken to establish that the object is never a module in the present-day sense. In §2 of Paper 30 the domain is an additive Abelian group acted on by ring elements, with explicit distributive and associative laws; §10 studies its chains and composition series. It is not simply a collection of modules. In §3 of Paper 9 the name belongs to a particular construction with algebraic functions. A shared stem supports linguistic continuity, but each definition must be read in its own setting.\par
\noindent{\footnotesize Izvorne zapisy / Original records: HIST-g061-001, HIST-g109-002, HIST-g117-003.}\par
\smallskip\noindent\emph{Sravnjenje s tekstom:} Rabota 9 prěnosi opis s prěměnnymi k cělym algebraičnym funkcijam. Zapis $z=\alpha+\sum g_i\eta_i$ i primětka o svojstvu $z-\alpha$ davajut konkretny smysl slovam modulna oblast i baza; sam nazv ne dokazyvaje nezavisnosti $\eta_i$.\par
\noindent\emph{Textual comparison:} Paper 9 passes from a description with indeterminates to integral algebraic functions. The expression $z=\alpha+\sum g_i\eta_i$ and the footnote on the property of $z-\alpha$ give concrete content to module domain and basis; the name alone does not prove independence of the $\eta_i$.\par
\noindent{\footnotesize TEXT-COMP-276; pasaže / passages: 7264, 7266, 7268, 7270, 7273, 7277, 7284; \texttt{CORPUS\_FORM\_COMPARISONS.json}.}\par
\smallskip\noindent\emph{Sravnjenje s tekstom:} Tekst dava aditivnu grupu, skalarnu akciju i jeje zakony, potom lance i kompozicionny red. Arhivno prěduprědženje protiv modernizacije ne opravyvaje negaciju etoj modulnoj struktury.\par
\noindent\emph{Textual comparison:} The text supplies an additive group, scalar action and its laws, followed by chains and a composition series. The archived warning against modernization does not justify denying this module structure.\par
\noindent{\footnotesize TEXT-COMP-282; pasaže / passages: 16309, 16311, 16314, 17189, 17190, 17193, 17195; \texttt{CORPUS\_FORM\_COMPARISONS.json}.}\par
\end{minipage}\par
\par\medskip\Needspace{7\baselineskip}
\noindent\begin{minipage}{\linewidth}
\hypertarget{TERM-1178}{}
\noindent\textbf{modulno predstavjenje}\quad{\footnotesize\texttt{TERM-1178}}\par
\noindent\emph{Iztočnik / Source:} Moduldarstellung\par
\noindent\emph{Rusky / Russian:} модульное представление\par
\noindent\emph{Ukrajinsky / Ukrainian:} модульне подання\par
\smallskip\noindent\emph{Razlog:} Modulno predstavjenje v rabotě 24 znači zapis polinomialnogo ideala kako modula s oprěděljenymi generatorami i koeficientami. Formula (3) dava diagonalny oblik s faktorami $E,E_1,\ldots$; pozdnějši primjer upotrěbjaje toj oblik. Arhivny razlog bira prosty kalk, da slovo predstavjenje ne sugeruje teoriju matricnyh prědstavjenij prstenov ili grup. Eti smysly istogo němečskogo Darstellung trěba razlikovati, ne avtomatično sjediniti. V prověrenom tekste upotrěbljaje se prědstavjenje s ě; arhivna lema čuvaje svoj zapis. Nijedny od dvoh zapisov ne može sam zaměniti opis togo, čto jest predstavjeno i nad kakymi koeficientami.\par
\noindent\emph{Reason:} Modulno predstavjenje in Paper 24 means expression of a polynomial ideal as a module with specified generators and coefficients. Formula (3) gives a diagonal form involving $E,E_1,\ldots$, used in a later example. The archived reason chooses a plain calque so that predstavjenje does not suggest the theory of matrix representations of rings or groups. These senses of the same German Darstellung must be distinguished, not automatically merged. The checked text uses prědstavjenje with ě; the archived headword retains its own spelling. Neither spelling can replace an account of what is being represented and over which coefficients.\par
\noindent{\footnotesize Izvorne zapisy / Original records: HIST-g100-012.}\par
\smallskip\noindent\emph{Sravnjenje s tekstom:} Formula (3) i pozdnějši primjer govoret o predstavjenju ideala posrědstvom modulnyh generatorov. To ne jest to samoe značenje prědstavjenija kako matrice, ktore dejstvajut na prostorě.\par
\noindent\emph{Textual comparison:} Formula (3) and the later example concern expression of an ideal through module generators. This is not the same sense of representation as matrices acting on a space.\par
\noindent{\footnotesize TEXT-COMP-283; pasaže / passages: 15307, 15309, 15311, 15317, 15681, 15683, 15685, 15687, 15689, 15691; \texttt{CORPUS\_FORM\_COMPARISONS.json}.}\par
\end{minipage}\par
\par\medskip\Needspace{7\baselineskip}
\noindent\begin{minipage}{\linewidth}
\hypertarget{TERM-1179}{}
\noindent\textbf{modulne svojstvo}\quad{\footnotesize\texttt{TERM-1179}}\par
\noindent\emph{Iztočnik / Source:} Moduleigenschaft der Restgruppe\par
\noindent\emph{Rusky / Russian:} модульное свойство\par
\noindent\emph{Ukrajinsky / Ukrainian:} модульна властивість\par
\smallskip\noindent\emph{Razlog:} Modulne svojstvo v §3 raboty 17 odnosi se k grupě ostatkov: klasi možno skladyvati i množiti elementom $F$ s lěvoj strany. Arhivny razlog naziva to modulopodobnym svojstvom i izběgaje dodatnu modernu etiketu. Matematična granica jest točnije vidima v dokazě: množenje dvoh klasov ne jest občno dobro oprěděljeno, dokoliko dejstvo $F$ na klas ostaje dobro oprěděljeno. Zato ne trěba tvrdžiti, že nema modulnoj struktury; ne ustanovjena jest struktura kvocientnogo prstena. Klas jedinice služi cikličnym generatorom pri tom dejstvu, ne sam po sebě vnutrišnoju jedinicoju neoprěděljenogo množenja klasov.\par
\noindent\emph{Reason:} Modulne svojstvo in §3 of Paper 17 concerns the residue group: its classes can be added and multiplied on the left by an element $F$. The archived reason calls this a module-like property and avoids an additional modern label. The mathematical boundary is clearer in the proof: multiplying two classes is not generally well defined, whereas the action of $F$ on a class is. One should therefore not claim that there is no module structure; what has not been established is a quotient-ring structure. The class of the identity serves as a cyclic generator for this action, not by itself as an internal identity for an undefined multiplication of classes.\par
\noindent{\footnotesize Izvorne zapisy / Original records: HIST-g079-018.}\par
\smallskip\noindent\emph{Sravnjenje s tekstom:} Dokaz razděljuje dobre oprěděljeno lěvo dejstvo i občno nedozvoljeno množenje dvuh klasov. Zato izraz modulno svojstvo ne trěba zaměniti kvocientnym prstenom, ni tumačiti kako negaciju modula.\par
\noindent\emph{Textual comparison:} The proof distinguishes a well-defined left action from generally unavailable multiplication of two classes. Module property should therefore neither be replaced by quotient ring nor interpreted as a denial of a module structure.\par
\noindent{\footnotesize TEXT-COMP-284; pasaže / passages: 11110, 11112, 11114, 11122; \texttt{CORPUS\_FORM\_COMPARISONS.json}.}\par
\end{minipage}\par
\par\medskip\Needspace{7\baselineskip}
\noindent\begin{minipage}{\linewidth}
\hypertarget{TERM-1180}{}
\noindent\textbf{modulna homomorfnost ... k sistemě slědov (§4, 3)}\quad{\footnotesize\texttt{TERM-1180}}\par
\noindent\emph{Iztočnik / Source:} Modulhomomorphie ... zum System der Spuren (§4, 3)\par
\noindent\emph{Rusky / Russian:} модульная гомоморфность ... к системе следов (§4, 3)\par
\noindent\emph{Ukrajinsky / Ukrainian:} модульна гомоморфність ... до системи слідів (§4, 3)\par
\smallskip\noindent\emph{Razlog:} Modulna homomorfnost k sistemě slědov čuva naziv konkretnoj mapy i odkaz k §4,3 raboty 31. Pri fiksovanom klasu idealov slědy tvoret modul. Pišemo $S$ za toj slěd; formula $S(\beta-\gamma)=S(\beta)-S(\gamma)$ vměstě s $S(c\beta)=cS(\beta)$ vyjadřaje aditivnost i soglasnost s koeficientami. V §8,2 to upotrěbjaje se v argumentu o konečnoj bazě. Arhivny razlog prědpočita polny opis skraćenju po linearnosti: čitatelj mora moči najdti, ktora mapa i ktory zakony sut potrěbne. Homomorfnost ne znači injektivnost i ne tvrdi zakona $S(\beta\gamma)=S(\beta)S(\gamma)$; to byl by drugy, prstenovy uslov.\par
\noindent\emph{Reason:} Modulna homomorfnost k sistemě slědov retains the name of the particular map and the reference to §4,3 of Paper 31. For a fixed ideal class, the traces form a module. Writing $S$ for this trace, $S(\beta-\gamma)=S(\beta)-S(\gamma)$ together with $S(c\beta)=cS(\beta)$ expresses additivity and compatibility with coefficients. In §8,2 this enters an argument about a finite basis. The archived reason prefers the full description to the abbreviation by linearity: the reader must be able to locate the map and the required laws. Homomorfnost does not mean injectivity and does not assert $S(\beta\gamma)=S(\beta)S(\gamma)$; that would be a different, ring-theoretic condition.\par
\noindent{\footnotesize Izvorne zapisy / Original records: HIST-g158-008.}\par
\smallskip\noindent\emph{Sravnjenje s tekstom:} Zakony za razliku i skalarny množitelj oprěděljajut potrěbnu homomorfnost. Odkaz iz §8,2 k §4,3 imenuje etu konstrukciju, bez tvrdženja o injektivnosti ili množenju dvuh slědov.\par
\noindent\emph{Textual comparison:} The laws for a difference and a scalar multiple specify the homomorphism required. The reference from §8,2 to §4,3 identifies this construction without asserting injectivity or a law for multiplying two traces.\par
\noindent{\footnotesize TEXT-COMP-285; pasaže / passages: 18024, 18057, 18059, 18061, 18063, 18857, 18897; \texttt{CORPUS\_FORM\_COMPARISONS.json}.}\par
\end{minipage}\par
\par\medskip\Needspace{7\baselineskip}
\noindent\begin{minipage}{\linewidth}
\hypertarget{TERM-1181}{}
\noindent\textbf{moduly}\quad{\footnotesize\texttt{TERM-1181}}\par
\noindent\emph{Iztočnik / Source:} Moduln\par
\noindent\emph{Rusky / Russian:} модули\par
\noindent\emph{Ukrajinsky / Ukrainian:} модулі\par
\smallskip\noindent\emph{Razlog:} Moduly jest nazvanje algebraičnoj struktury, ne obče slovo za sistemy. V rabotě 17 polinomy mogut ne komutovati. Klasy ostatkov možno dodavati i množiti polinomom s lěvoj strany, ale produkt dvuh klasov ne musi byti dobro oprěděljeny. Modulny jezyk zato čuva potrěbnu razliku medžu dejstvom polinomov i množenjem v prstenu ostatkov. Arhivny razlog izbira prěmy algebraičny termin vměsto opisov sistemy ili idealy; to ne znači, že idealna interpretacija nikogda ne jest možna. Konečnost v uvodu znači konečno čislo linearno nezavisnyh klasov, ne konečno čislo vsih klasov. Rusko i ukrajinsko sravnjena podpirajut izbor toj slovnoj rodiny, ale sama ne dokazujut ponimanje medžu čitateljami vsih slovjanskyh jezykov. Dostupny tekst nazyva zatvorenost na $FM$ pravostrannom, a na $MF$ lěvostrannom. To označenje strany trěba sravniti s pečatnym izvornikom; ne zaměnjujemo go tiho današnym obyčajem.\par
\noindent\emph{Reason:} Moduly names an algebraic structure, not systems in general. In Paper 17 the polynomials need not commute. Residue classes can be added and multiplied by a polynomial on the left, but multiplication of two classes need not be well defined. Module language therefore preserves the distinction between a polynomial action and multiplication in a residue-class ring. The archived reason chooses direct algebraic terminology over the paraphrases systems or ideals; it does not rule out an ideal interpretation where appropriate. Finiteness in the introduction means finitely many linearly independent classes, not finitely many classes altogether. The Russian and Ukrainian comparanda support the lexical family, but do not themselves establish comprehension across Slavic readerships. The available text calls closure under $FM$ right-sided and under $MF$ left-sided. That side designation requires comparison with the printed witness; it is not silently replaced by modern convention.\par
\noindent{\footnotesize Izvorne zapisy / Original records: HIST-g079-001.}\par
\smallskip\noindent\emph{Sravnjenje s tekstom:} Korpus odlišaje množenje klasa polinomom od množenja dvuh klasov. To objasnjaje, počemu obecny opis sistema ne zaměnjaje modulny jezyk. Označenja stran pri FM i MF ostavajut posebnym pitanjem za pečatny izvornik.\par
\noindent\emph{Textual comparison:} The corpus distinguishes multiplying a class by a polynomial from multiplying two classes. This explains why the generic description system cannot replace module language. The side designations attached to FM and MF remain a separate printed-witness question.\par
\noindent{\footnotesize TEXT-COMP-286; pasaže / passages: 10957, 11088, 11090, 11092, 11094, 11096, 11102, 11110, 11122; \texttt{CORPUS\_FORM\_COMPARISONS.json}.}\par
\end{minipage}\par
\par\medskip\Needspace{7\baselineskip}
\noindent\begin{minipage}{\linewidth}
\hypertarget{TERM-1182}{}
\noindent\textbf{modulny produkt}\quad{\footnotesize\texttt{TERM-1182}}\par
\noindent\emph{Iztočnik / Source:} Modulprodukt\par
\noindent\emph{Rusky / Russian:} модульное произведение; произведение модулей\par
\noindent\emph{Ukrajinsky / Ukrainian:} добуток модулів; модульний добуток\par
\smallskip\noindent\emph{Razlog:} Modulny produkt označa sumy produktov elementov, ne množinu uporjadkovanyh par. V §8 raboty 34 produkt $AB$ sostoji se iz konečnyh sum $ab$, kde $a$ naleži k $A$ i $b$ k $B$. To jest važno takože za konstrukcije poredkov v rabotě 42. Arhiv imaje dva razloge za izbor: čuvati modulnu konstrukciju, ne nazyvati ju prědčasno idealnym produktom; i upotrěbjati produkt soglasno s terminom skrěženy produkt, vměsto alternativy dobytok. To jest kvalitativny izbor doslědnosti, ne dokaz, že dobytok jest loš slovo. Termin ne znači tenzorny produkt ni prěmy produkt modulov; konkretny zakon množenja davaje rabota.\par
\noindent\emph{Reason:} Modulny produkt denotes sums of products of elements, not a set of ordered pairs. In §8 of Paper 34, $AB$ consists of finite sums of products $ab$, with $a$ in $A$ and $b$ in $B$. This also matters for the constructions of orders in Paper 42. The archive gives two reasons: preserve the module-level construction without prematurely calling it an ideal product; and use produkt consistently with skrěženy produkt, rather than the alternative dobytok. This is a qualitative preference for consistency, not evidence that dobytok is an unsuitable word. The term does not mean a tensor product or a direct product of modules; the work supplies the particular multiplication law.\par
\noindent{\footnotesize Izvorne zapisy / Original records: HIST-g120-003, HIST-g186-006, HIST-g212-007.}\par
\smallskip\noindent\emph{Sravnjenje s tekstom:} Definicija v rabotě 34 zahtěvaje sumy produktov. Konstrukcije poredkov v rabotah 31 i 42 čuvajut toj smysl; ne ide o tenzorny produkt.\par
\noindent\emph{Textual comparison:} The definition in Paper 34 requires sums of products. The constructions of orders in Papers 31 and 42 retain this meaning; they are not tensor products.\par
\noindent{\footnotesize TEXT-COMP-287; pasaže / passages: 17394, 20418, 20420, 20423, 24486, 24487, 24488, 24489, 24490, 24491, 24492, 24493, 24494, 24495, 24616, 24617, 24618, 24619; \texttt{CORPUS\_FORM\_COMPARISONS.json}.}\par
\end{minipage}\par
\par\medskip\Needspace{7\baselineskip}
\noindent\begin{minipage}{\linewidth}
\hypertarget{TERM-1183}{}
\noindent\textbf{modulny kvocient}\quad{\footnotesize\texttt{TERM-1183}}\par
\noindent\emph{Iztočnik / Source:} Modulquotient\par
\noindent\emph{Rusky / Russian:} частное модулей\par
\noindent\emph{Ukrajinsky / Ukrainian:} частка модулів\par
\smallskip\noindent\emph{Razlog:} Modulny kvocient v §9 raboty 30 označa konstrukciju s dvojtočkoju. Ako jednačenje $\mathfrak A T=\mathfrak B$ imaje jedino rěšenje, tekst identifikuje go s $\mathfrak B:\mathfrak A$. Uslovje jedinstvenosti ne možno izpustiti i tvrditi, že vsak takovy kvocient bez daljših prědposylok rěšaje produktno jednačenje. Arhivny razlog oddělja toj pojem od polja kvocientov, prstena klasov ostatkov i obyčnyh drobjev. Rusko častnoe modulej i ukrajinsko častka moduliv sut sravnjena za toj specialny smysl, ne instrukcija čitati dvojtočku kako obyčno děljenje čisel.\par
\noindent\emph{Reason:} Modulny kvocient in §9 of Paper 30 denotes the colon construction. If the equation $\mathfrak A T=\mathfrak B$ has a unique solution, the text identifies it with $\mathfrak B:\mathfrak A$. The uniqueness condition must not be dropped to suggest that every such quotient solves the product equation without further hypotheses. The archived reason distinguishes this object from a quotient field, a residue-class ring and ordinary fractions. The Russian частное модулей and Ukrainian частка модулів are comparanda for this specialised sense, not instructions to read the colon as ordinary numerical division.\par
\noindent{\footnotesize Izvorne zapisy / Original records: HIST-g116-011.}\par
\smallskip\noindent\emph{Sravnjenje s tekstom:} Definicija kvocienta koristi obratny porjadok děliteljnosti. Uvjet jedinstvenosti v §9 ne trěba izpustiti pri prěhodu od produktnogo jednačenja k kolon-kvocientu.\par
\noindent\emph{Textual comparison:} The quotient definition uses the reverse divisibility ordering. The uniqueness condition in §9 must be retained when passing from a product equation to a colon quotient.\par
\noindent{\footnotesize TEXT-COMP-288; pasaže / passages: 16628, 16630, 16632, 16634, 17087, 17089, 17091, 17093; \texttt{CORPUS\_FORM\_COMPARISONS.json}.}\par
\end{minipage}\par
\par\medskip\Needspace{7\baselineskip}
\noindent\begin{minipage}{\linewidth}
\hypertarget{TERM-1184}{}
\noindent\textbf{modulny kvocient / idealny kvocient}\quad{\footnotesize\texttt{TERM-1184}}\par
\noindent\emph{Iztočnik / Source:} Modulquotient / Idealquotient\par
\noindent\emph{Rusky / Russian:} частное модулей / частное идеалов\par
\noindent\emph{Ukrajinsky / Ukrainian:} частка модулів / частка ідеалів\par
\smallskip\noindent\emph{Razlog:} Modulny kvocient i idealny kvocient tvorjat jednu terminologičnu paru, zato že §5,4 raboty 30 prvo oprěděljaje konstrukciju za moduly v razširjajučem prstenu, a potom specializuje ju k idealam. $C=A:B$ izpolnjaje $CB\subseteq A$, a vsaky modul $D$ s $DB\subseteq A$ jest sodržany v $C$. Po vključanju to jest največši takovy modul. Slovo najmenši v istorijskom opisu odnosi se k obratnomu porjadku děliteljnosti; ne trěba čitati go kako najmenši po vključanju. Arhivny razlog zahtěvaje odlišati etot kvocient od prstena ostatkov $R/I$: elementy kolon-kvocienta sut množitelje, ne klasy po ekvivalenciji. Pari odgovarja para modul/ideal, ne dva slučajno izbrana sinonima.\par
\noindent\emph{Reason:} Modulny kvocient and idealny kvocient form a terminological pair because §5,4 of Paper 30 first defines the construction for modules in an extension ring and then specialises it to ideals. $C=A:B$ satisfies $CB\subseteq A$, and every module $D$ with $DB\subseteq A$ is contained in $C$. It is therefore the largest such module by inclusion. Smallest in the historical description refers to the reverse divisibility ordering, not to the smallest set by inclusion. The archived reason requires a distinction from the residue-class ring $R/I$: elements of a colon quotient are multipliers, not equivalence classes. The two names reflect the module/ideal distinction rather than being arbitrary synonyms.\par
\noindent{\footnotesize Izvorne zapisy / Original records: HIST-g112-003, HIST-g116-004.}\par
\smallskip\noindent\emph{Sravnjenje s tekstom:} Definicija kvocienta koristi obratny porjadok děliteljnosti. Uvjet jedinstvenosti v §9 ne trěba izpustiti pri prěhodu od produktnogo jednačenja k kolon-kvocientu.\par
\noindent\emph{Textual comparison:} The quotient definition uses the reverse divisibility ordering. The uniqueness condition in §9 must be retained when passing from a product equation to a colon quotient.\par
\noindent{\footnotesize TEXT-COMP-288; pasaže / passages: 16628, 16630, 16632, 16634, 17087, 17089, 17091, 17093; \texttt{CORPUS\_FORM\_COMPARISONS.json}.}\par
\end{minipage}\par
\par\medskip\Needspace{7\baselineskip}
\noindent\begin{minipage}{\linewidth}
\hypertarget{TERM-1185}{}
\noindent\textbf{modulny teorem}\quad{\footnotesize\texttt{TERM-1185}}\par
\noindent\emph{Iztočnik / Source:} Modulsatz\par
\noindent\emph{Rusky / Russian:} модульная теорема\par
\noindent\emph{Ukrajinsky / Ukrainian:} модульна теорема\par
\smallskip\noindent\emph{Razlog:} Modulny teorem jest kratko nazvanje točno ukazanogo teorema prěnosa uslovij cěpov. V §2 raboty 30 $M$ jest konečno poraždany modul nad komutativnym prstenom $R$ s jediniceju. Uslovje rastučih odnosno padajučih cěpov idealov v $R$ togda prěhodi na podmoduly $M$. Konečny tu ne znači, že $M$ imaje konečno mnogo elementov. Arhivny razlog čuva nazvanje zato, že §1 odkazuje k tomu teoremu v lemě o cělyh elementah. To ne jest nazvanje vsih teoremov o modulah ni zakon modularnoj rešetky. V rěčniku trěba čitati nazvanje vměstě s jego prědposylkami.\par
\noindent\emph{Reason:} Modulny teorem is a short name for the specified theorem transferring chain conditions. In §2 of Paper 30, $M$ is a finitely generated module over a commutative ring $R$ with identity. The ascending or descending chain condition on ideals of $R$ then transfers to submodules of $M$. Finite here does not mean that $M$ has finitely many elements. The archived reason retains the name because §1 invokes this theorem in the integral-elements lemma. It is neither a name for all module theorems nor the modular law for lattices. The dictionary name must be read with its hypotheses.\par
\noindent{\footnotesize Izvorne zapisy / Original records: HIST-g109-003.}\par
\smallskip\noindent\emph{Sravnjenje s tekstom:} Teorem imaje eksplicitne prědposylky o prstenu i modulě. Konečnost jest oprěděljena prěz konečnu poraždajuču sistemu, ne prěz čislo vsih elementov.\par
\noindent\emph{Textual comparison:} The theorem has explicit hypotheses on the ring and module. Finiteness is defined through a finite generating system, not by counting all elements.\par
\noindent{\footnotesize TEXT-COMP-289; pasaže / passages: 16307, 16309, 16311, 16314, 16316, 16318; \texttt{CORPUS\_FORM\_COMPARISONS.json}.}\par
\end{minipage}\par
\par\medskip\Needspace{7\baselineskip}
\noindent\begin{minipage}{\linewidth}
\hypertarget{TERM-1186}{}
\noindent\textbf{teorija modulov / bazis modula / prosty modul / nulovy modul}\quad{\footnotesize\texttt{TERM-1186}}\par
\noindent\emph{Iztočnik / Source:} Modultheorie / Modulbasis / Primmodul / Nullmodul\par
\noindent\emph{Rusky / Russian:} теория модулей / базис модуля / простой модуль / нулевой модуль\par
\noindent\emph{Ukrajinsky / Ukrainian:} теорія модулів / базис модуля / простий модуль / нульовий модуль\par
\smallskip\noindent\emph{Razlog:} Teorija modulov, bazis modula, prosty modul i nulovy modul sut povezane nazvanja v rabotě 15, ale ne vse imajut današny najočekivany smysl. Modul relacij sostoji se iz cěločislovyh polinomov, ktore izčezajut po podstavjenju determinantov; zatvoren jest na dodavanje i množenje polinomami. Bazis tu označa poraždajuču sistemu, bez tvrdženja linearnoj nezavisnosti. Nulovy modul sodrži samo nulu. Arhivny razlog dopušča, že prosty pokryva i prostotu v idealnom smyslu, i prostotu modula. To ne opravdava matematično izravnjenje: prostota v etom kontekstu relacij i redukcije modulo $p$ ne znači, že modul ne imaje vlastnyh nenulovyh podmodulov. Dodatno, identifikacija $M=M_p$ v dostupnom tekstu zahtěvaje sravnjenje s pečatnym izvornikom; po tam danoj cěločislovoj definiciji konstanta $p$ naleži k $M_p$, ale ne k $M$. Ne predlagajemo tiho popravljanje etoj ravnosti.\par
\noindent\emph{Reason:} Teorija modulov, bazis modula, prosty modul and nulovy modul are related names in Paper 15, but not all have the most immediately expected modern meaning. The relation module consists of integer polynomials that vanish after substitution of the determinants; it is closed under addition and polynomial multiplication. Basis here means a generating system, without an assertion of linear independence. The zero module contains only zero. The archived rationale allows a prime/simple ambiguity in prosty. This does not justify mathematical equivalence: prime in this relation-theoretic, modulo-$p$ context does not mean a module with no proper nonzero submodules. Moreover, the identification $M=M_p$ in the available text requires comparison with the printed witness: under the integer-polynomial definition given there, the constant $p$ belongs to $M_p$ but not $M$. No silent correction of that equality is proposed.\par
\noindent{\footnotesize Izvorne zapisy / Original records: HIST-g077-009.}\par
\smallskip\noindent\emph{Sravnjenje s tekstom:} Množenje relacije libovoljnym polinomom ostavja relaciju. Kvadratične relacije poraždajut N; pri menje než četyroh indeksah N jest nulovy modul. Bazis jest poraždajuča sistema.\par
\noindent\emph{Textual comparison:} Multiplying a relation by any polynomial preserves a relation. The quadratic relations generate N; with fewer than four indices N is the zero module. Basis means a generating system.\par
\noindent{\footnotesize TEXT-COMP-290; pasaže / passages: 10319, 10321, 10325, 10371, 10374, 10376, 10378, 10394; \texttt{CORPUS\_FORM\_COMPARISONS.json}.}\par
\smallskip\noindent\emph{Sravnjenje s tekstom:} Dostupny tekst poistovnjaje $M$ i $M_p$, ale opisuje $M_p$ cěločislovymi polinomami, ktore izčezajut modulo $p$. Konstanta $p$ zato razděljuje dvě dane definicije. Trěba sravniti pečatny izvornik i njegovu konvenciju redukcije; ne imajemo ovdě dokazano izvorno popravljanje.\par
\noindent\emph{Textual comparison:} The available text identifies $M$ with $M_p$ but describes $M_p$ by integer polynomials vanishing modulo $p$. The constant $p$ therefore separates the two definitions as stated. The printed witness and its reduction convention require comparison; no source correction is established here.\par
\noindent{\footnotesize TEXT-COMP-291; pasaže / passages: 10319, 10331, 10335, 10337, 10390; \texttt{CORPUS\_FORM\_COMPARISONS.json}.}\par
\end{minipage}\par
\par\medskip\Needspace{7\baselineskip}
\noindent\begin{minipage}{\linewidth}
\hypertarget{TERM-1187}{}
\noindent\textbf{čisla mnogokratnosti}\quad{\footnotesize\texttt{TERM-1187}}\par
\noindent\emph{Iztočnik / Source:} Multiplicity numbers\par
\noindent\emph{Rusky / Russian:} числа кратности\par
\noindent\emph{Ukrajinsky / Ukrainian:} числа кратності\par
\smallskip\noindent\emph{Razlog:} Čisla mnogokratnosti imenujut čisla, ktore čuvajut informaciju o povtorjenju faktorov. V §7 raboty 22 eksponenty linearnyh faktorov i čisla pri nerazložimyh faktorah morajut byti odlišene. Ako $\gamma_\mu$ jest stepen nerazložimogo faktora i $\beta_\mu$ jego eksponent, tekst davaje mnogokratnost $\beta_\mu\gamma_\mu$ kako čislo linearno nezavisnyh ostatkovyh klasov nad ukazanym poljem. Arhivny opis čislo ostatkovyh klasov zato trěba dopolniti slovami linearno nezavisnyh. Anglejsko Multiplicity numbers v arhivu jest opis pojma, ne německy citat. Izbor prodolžaje rodinu mnogokratnost; prirodnost izraza ostala jest otvorjena za jezykovu ocěnku. Čislove vagi kandidatov ne sut zapisane.\par
\noindent\emph{Reason:} Čisla mnogokratnosti names numbers retaining information about repeated factors. In §7 of Paper 22, exponents of linear factors must be distinguished from the numbers attached to irreducible factors. If $\gamma_\mu$ is the degree of an irreducible factor and $\beta_\mu$ its exponent, the text gives the multiplicity $\beta_\mu\gamma_\mu$ as the number of linearly independent residue classes over the specified field. The archival shorthand residue-class count therefore needs the qualification linearly independent. The archived English Multiplicity numbers is a concept description, not a German quotation. The choice continues the mnogokratnost family; its naturalness remains a linguistic review question. No numerical candidate weights are recorded.\par
\noindent{\footnotesize Izvorne zapisy / Original records: HIST-g091-006, HIST-g092-075.}\par
\smallskip\noindent\emph{Sravnjenje s tekstom:} Eksponent i stepen faktora imajut odlišne uloge. Njih produkt daje navedenu mnogokratnost; izpustiti linearnu nezavisnost značilo by poměniti čislo, ktoro opisuje tekst.\par
\noindent\emph{Textual comparison:} Exponent and factor degree have distinct roles. Their product gives the stated multiplicity; omitting linear independence would change the quantity described.\par
\noindent{\footnotesize TEXT-COMP-292; pasaže / passages: 15000, 15005, 15012, 15014, 15016, 15631, 15634, 15635, 15637, 15642; \texttt{CORPUS\_FORM\_COMPARISONS.json}.}\par
\end{minipage}\par
\par\medskip\Needspace{7\baselineskip}
\noindent\begin{minipage}{\linewidth}
\hypertarget{TERM-1188}{}
\noindent\textbf{multiplikacijsko kolco / regularne multiplikacijske kolca}\quad{\footnotesize\texttt{TERM-1188}}\par
\noindent\emph{Iztočnik / Source:} Multiplikationsring / reguläre Multiplikationsringe\par
\noindent\emph{Rusky / Russian:} мультипликационное кольцо / регулярные мультипликационные кольца\par
\noindent\emph{Ukrajinsky / Ukrainian:} мультиплікаційне кільце / регулярні мультиплікаційні кільця\par
\smallskip\noindent\emph{Razlog:} Multiplikacijsko kolco jest istorijno nazvanje specialnoj klasy prstenov, ne opis libovoljnogo prstena ili oblasti operatorov. Rabota 31, §8, citira Krulla i objasnjaje, že on upotrěbjaje nazvanje regularne multiplikacijske prsteny za razgledany slučaj bez děliteljev nuly. Uslovja vključajut jedinicu i jednoznačno razloženje nenulovyh vlastnyh idealov na potegy prostyh idealov, ktore sut maksimalne. Arhivny izbor čuva poznavajemo medžunarodno osnovu zaradi etogo konkretnogo nazvanja. Poznavatelnost medžu slovjanskymi čitateljami jest naměr, ne rezultat testa. V dostupnom korpusu kolco odgovarja prsten, ale fraza multiplikacijsko prsten ne imaje soglasovanja roda; očekivana mužska forma jest multiplikacijsky prsten. Eto jest jezykovo pitanje za popravljanje, ne nova matematična definicija. Citiranje Krulla v rabotě ne dokazuje, že njegov članok byl pročitany pri izboru termina.\par
\noindent\emph{Reason:} Multiplikacijsko kolco is the historical name of a particular class of rings, not a description of any ring or operator domain. Paper 31, §8 cites Krull and explains that he uses regular multiplication rings for the case without zero divisors considered here. The hypotheses include an identity and unique factorisation of nonzero proper ideals into powers of prime ideals that are maximal. The archived choice retains a recognisable international stem because this is a specific historical name. Cross-Slavic recognisability is an aim, not a measured result. In the available corpus kolco corresponds to prsten, but multiplikacijsko prsten lacks gender agreement; the expected masculine form is multiplikacijsky prsten. This is a linguistic correction question, not a new mathematical definition. The work's citation of Krull does not establish that his article was read when the term was chosen.\par
\noindent{\footnotesize Izvorne zapisy / Original records: HIST-g155-004.}\par
\smallskip\noindent\emph{Sravnjenje s tekstom:} Citacija razlikuje istorijno nazvanje od obče frazy o množenju. Korpusna fraza multiplikacijsko prsten imaje nezavisno jezykovo pitanje: neutralna forma pridavnika ne soglašaje se s mužskym prsten.\par
\noindent\emph{Textual comparison:} The citation distinguishes a historical name from a generic phrase about multiplication. The corpus phrase multiplikacijsko prsten also presents a separate linguistic issue: the neuter adjective does not agree with masculine prsten.\par
\noindent{\footnotesize TEXT-COMP-293; pasaže / passages: 18787, 18789, 18823; \texttt{CORPUS\_FORM\_COMPARISONS.json}.}\par
\end{minipage}\par
\par\medskip\Needspace{7\baselineskip}
\noindent\begin{minipage}{\linewidth}
\hypertarget{TERM-1189}{}
\noindent\textbf{tablica množenja}\quad{\footnotesize\texttt{TERM-1189}}\par
\noindent\emph{Iztočnik / Source:} Multiplikationstafel\par
\noindent\emph{Rusky / Russian:} таблица умножения\par
\noindent\emph{Ukrajinsky / Ukrainian:} таблиця множення\par
\smallskip\noindent\emph{Razlog:} Tablica množenja v §6 raboty 34 opisuje produkty baznyh elementov hiperkompleksnogo sistema. Formula $u_i u_k=\sum_j u_j\gamma^j_{ik}$ davaje strukturne konstanty; njim sut naložene uslovja asociativnosti. To ne jest prosto tablica cěločislovogo množenja i ne proizvoljna zbirka koeficientov. Arhivny razlog daje prědnost poznatoj slovjanskoj konstrukciji tablica množenja, blizkoj ruskomu i ukrajinskomu sravnjenju, vměsto opisnogo hibrida. Konkretny odbačeny hibrid ne jest zapisan; zato ne izmysljajemo kandidata ni čislovu ocěnku. Značenje tablicy jest oprěděljeno njenoju funkcijeju v algebraičnom sistemu.\par
\noindent\emph{Reason:} Tablica množenja in §6 of Paper 34 specifies products of basis elements in a hypercomplex system. The formula $u_i u_k=\sum_j u_j\gamma^j_{ik}$ gives structure constants, which are constrained by associativity. This is not merely an elementary integer multiplication table or an arbitrary collection of coefficients. The archived reason prefers the familiar Slavic construction tablica množenja, close to the Russian and Ukrainian comparanda, over a descriptive hybrid. No particular rejected hybrid is recorded, so neither a candidate wording nor a numerical score is invented. The table's meaning is determined by its function in the algebraic system.\par
\noindent{\footnotesize Izvorne zapisy / Original records: HIST-g184-006.}\par
\smallskip\noindent\emph{Sravnjenje s tekstom:} Tablica davaje produkty baznyh elementov s koeficientami. Asociativnost naložuje relacije tym koeficientam; tablica ne jest proizvoljna.\par
\noindent\emph{Textual comparison:} The table gives products of basis elements with coefficients. Associativity constrains those coefficients; the table is not arbitrary.\par
\noindent{\footnotesize TEXT-COMP-294; pasaže / passages: 20279, 20281, 20283, 20285, 20287, 20289; \texttt{CORPUS\_FORM\_COMPARISONS.json}.}\par
\end{minipage}\par
\par\medskip\Needspace{7\baselineskip}
\noindent\begin{minipage}{\linewidth}
\hypertarget{TERM-1190}{}
\noindent\textbf{množitelje}\quad{\footnotesize\texttt{TERM-1190}}\par
\noindent\emph{Iztočnik / Source:} Multiplikatoren\par
\noindent\emph{Rusky / Russian:} множители\par
\noindent\emph{Ukrajinsky / Ukrainian:} множники\par
\smallskip\noindent\emph{Razlog:} Množitelje v §9 raboty 9 sut elementy, s ktorimi tvorjat se linearne kombinacije generatorov modula. Tekst izrično dopušča vse cěle algebraične funkcije od $\eta$ kako množitelje; ne ograničuje jih na konstanty i ne zaměnjaje jih libovoljnymi racionalnymi funkcijami. Arhivny razlog bira obče algebraično slovo blizko ruskomu množiteĺ i ukrajinskomu množnyk. Jego primětka, že tu ne ide o operatornu oblast množiteljev, razjasnjaje toj konkretny pasus; ne zabranjaje upotrěbu iste slovnoj rodiny za dejstvo prstena v rabotě 34. Kontekst, a ne razne koreni za vsaku upotrěbu, razlikuje te smysly.\par
\noindent\emph{Reason:} Množitelje in §9 of Paper 9 are the elements used to form linear combinations of module generators. The text expressly allows all integral algebraic functions of $\eta$ as multipliers; it neither restricts them to constants nor replaces them by arbitrary rational functions. The archived reason chooses the general algebraic word corresponding to Russian множитель and Ukrainian множник. Its observation that this is not an operator-theoretic multiplier domain explains this particular passage; it does not prohibit the same lexical family for a ring action in Paper 34. Context, rather than a separate root for every use, distinguishes the senses.\par
\noindent{\footnotesize Izvorne zapisy / Original records: HIST-g067-006.}\par
\smallskip\noindent\emph{Sravnjenje s tekstom:} Modul tvori se s cělymi algebraičnymi funkcijami kak množiteljami. Točny obseg množiteljev jest čestju definicije, ne dodatna sinonimična glosa.\par
\noindent\emph{Textual comparison:} The module is formed using integral algebraic functions as multipliers. The specified range of multipliers is part of the definition, not an optional synonymous gloss.\par
\noindent{\footnotesize TEXT-COMP-295; pasaže / passages: 8190, 8191, 8192, 8193, 8194; \texttt{CORPUS\_FORM\_COMPARISONS.json}.}\par
\end{minipage}\par
\par\medskip\Needspace{7\baselineskip}
\noindent\begin{minipage}{\linewidth}
\hypertarget{TERM-1191}{}
\noindent\textbf{oblast množiteljev}\quad{\footnotesize\texttt{TERM-1191}}\par
\noindent\emph{Iztočnik / Source:} Multiplikatorenbereich\par
\noindent\emph{Rusky / Russian:} область множителей\par
\noindent\emph{Ukrajinsky / Ukrainian:} область множників\par
\smallskip\noindent\emph{Razlog:} Oblast množiteljev označa prsten, ktory dejstvuje na aditivnoj grupě modula. V §1, priměru 4 raboty 34 množitelje stojat lěvo; tri zakony izražajut distributivnost i soglasnost dejstva s množenjem prstena. Absolutna oblast množiteljev znači, že različne elementy davajut različne operacije, to jest dejstvo jest věrno. §8 poistovnja elementy s jednakym dejstvom i dobyvaje prsten ostatkov po idealu, ktory anuluje cěly modul. Arhivny razlog čuva paralelu s oblastju operatorov i točny odkaz nazad k §1. Oblast tu ne znači oblast cělosti: odsutnost děliteljev nuly ne jest utvrđena. Takože ne mora vključati vse endomorfizmy aditivnoj grupy.\par
\noindent\emph{Reason:} Oblast množiteljev denotes a ring acting on the additive group of a module. In §1, example 4 of Paper 34 the multipliers stand on the left; three laws express distributivity and compatibility with ring multiplication. An absolute multiplier domain means that distinct elements give distinct operations: the action is faithful. Section 8 identifies elements with the same action and obtains the residue-class ring modulo the ideal annihilating the whole module. The archived reason preserves the parallel with operator domain and the precise reference back to §1. Domain here does not mean integral domain: absence of zero divisors is not asserted. Nor need it include every endomorphism of the additive group.\par
\noindent{\footnotesize Izvorne zapisy / Original records: HIST-g179-005, HIST-g186-005.}\par
\smallskip\noindent\emph{Sravnjenje s tekstom:} Različne množitelje v absolutnoj oblasti davajut različne operacije. Prěhod k njoj jest kvocient po idealu, ktory anuluje cěly modul; ne jest tvrdženje o odsutnosti děliteljev nuly.\par
\noindent\emph{Textual comparison:} Distinct multipliers in the absolute domain induce distinct operations. Passing to it takes the quotient by the ideal annihilating the whole module; it does not assert absence of zero divisors.\par
\noindent{\footnotesize TEXT-COMP-296; pasaže / passages: 19823, 19827, 19828, 19829, 19833, 19835, 20414; \texttt{CORPUS\_FORM\_COMPARISONS.json}.}\par
\end{minipage}\par
\par\medskip\Needspace{7\baselineskip}
\noindent\begin{minipage}{\linewidth}
\hypertarget{TERM-1192}{}
\noindent\textbf{uslovja mnogokratnosti / mnogokratnost}\quad{\footnotesize\texttt{TERM-1192}}\par
\noindent\emph{Iztočnik / Source:} Multiplizitaetsbedingungen / Multiplizitaet\par
\noindent\emph{Rusky / Russian:} условия кратности / кратность\par
\noindent\emph{Ukrajinsky / Ukrainian:} умови кратності / кратність\par
\smallskip\noindent\emph{Razlog:} Uslovja mnogokratnosti i mnogokratnost upotrěbjajut jednu slovjansku osnovu za povtorjenje s točnym matematičnym čislom. V rabotě 45 Kapferer oprěděljaje lokalne uslovja prěz eksponenty izčezanja, označene $(K_i)$ i $(g_i)$; suma čisel $(g_i)$ jest mnogokratnost prěsěčenja $\mu$. V Teoremu III, pri ego prědposylkah o prostyh točkah na $\psi$ i vzajemnoj prostotě $K$ s $\psi$, pripadnost $K$ idealu $(\varphi,\psi)$ odgovarja neravnostjam mnogokratnostij prěsěčenja v vsakom občem punktu. To ne jest samo uslovje, že kriva prohodi prěz toj punkt. Arhiv bira mnogokratnost vměsto medžunarodnoj zaimky zaradi prozračnosti i soglasnosti s jezykom o punktah i prěsěčenjah. To jest razlog izbora, ne dokaz, že vse lokalne pojmy mnogokratnosti sut identične.\par
\noindent\emph{Reason:} Uslovja mnogokratnosti and mnogokratnost use one Slavic stem for repetition measured by a precise mathematical number. In Work 45, Kapferer defines local conditions through vanishing orders denoted by $(K_i)$ and $(g_i)$; the sum of the numbers $(g_i)$ is the intersection multiplicity $\mu$. In Theorem III, under its hypotheses concerning simple points on $\psi$ and coprimality of $K$ with $\psi$, membership of $K$ in $(\varphi,\psi)$ corresponds to inequalities between intersection multiplicities at every common point. This is not merely the condition that a curve passes through the point. The archive chooses mnogokratnost over an international loan for transparency and compatibility with point and intersection language. This explains the choice; it does not establish that all local notions of multiplicity coincide.\par
\noindent{\footnotesize Izvorne zapisy / Original records: HIST-g214-001.}\par
\smallskip\noindent\emph{Sravnjenje s tekstom:} Uslovja $(K_i)$ i $(g_i)$ ne znači samo prohod prěz punkt. Teorem III imaje osobne prědposylky, ktore trěba čuvati pri tumačenju mnogokratnosti.\par
\noindent\emph{Textual comparison:} The conditions on $(K_i)$ and $(g_i)$ do not merely assert passage through a point. Theorem III has particular hypotheses that must be retained when interpreting multiplicity.\par
\noindent{\footnotesize TEXT-COMP-297; pasaže / passages: 45-isv.tex:158, 45-isv.tex:159, 45-isv.tex:160, 45-isv.tex:161, 45-isv.tex:162, 45-isv.tex:163, 45-isv.tex:164, 45-isv.tex:165, 45-isv.tex:166, 45-isv.tex:167, 45-isv.tex:168, 45-isv.tex:169, 45-isv.tex:170, 45-isv.tex:172, 45-isv.tex:173, 45-isv.tex:174, 45-isv.tex:175, 45-isv.tex:176, 45-isv.tex:177, 45-isv.tex:330, 45-isv.tex:331, 45-isv.tex:332, 45-isv.tex:333, 45-isv.tex:334, 45-isv.tex:335, 45-isv.tex:336, 45-isv.tex:337, 45-isv.tex:338; \texttt{CORPUS\_FORM\_COMPARISONS.json}.}\par
\end{minipage}\par
\par\medskip\Needspace{7\baselineskip}
\noindent\begin{minipage}{\linewidth}
\hypertarget{TERM-1193}{}
\noindent\textbf{mnogokratnost}\quad{\footnotesize\texttt{TERM-1193}}\par
\noindent\emph{Iztočnik / Source:} Multiplizität\par
\noindent\emph{Rusky / Russian:} кратность\par
\noindent\emph{Ukrajinsky / Ukrainian:} кратність\par
\smallskip\noindent\emph{Razlog:} Mnogokratnost povezuje terminologiju rezultantov v rabotah 22 i 24 s jezykom primarnyh idealov. Arhivny razlog daje prědnost prozračnoj slovjanskoj tvorbě i zahtěvaje doslědnost medžu etymi smyslami, pri otvorjenoj jezykovoj ocěnci. Matematična doslědnost ne znači, že vsekde trěba vzeti samo eksponent faktora. V rabotě 22 stepen nerazložimogo faktora množi se s jego eksponentom; rezultat měri čislo linearno nezavisnyh ostatkovyh klasov nad ukazanym poljem. Rabota 24, Teorem XI i prinaležna primětka H.-N., povezuje stepen primarnogo faktora normy s odgovarjajučim prostorom klasov. Rěčnik mora pokazati, ktory objekt i ktoro polje sut zadane, vměsto obče glosy čislo povtorjenij.\par
\noindent\emph{Reason:} Mnogokratnost links the resultant terminology of Papers 22 and 24 with primary-ideal language. The archived reason favours a transparent Slavic formation and consistency between these uses, while leaving its linguistic assessment open. Mathematical consistency does not mean taking the factor exponent alone everywhere. In Paper 22, the degree of the irreducible factor multiplies its exponent; the result measures the number of linearly independent residue classes over the specified field. Paper 24, Theorem XI and the accompanying H.-N. footnote, relates the degree of a primary factor of the norm to the corresponding space of classes. A dictionary must identify the object and field rather than offer only the generic gloss number of repetitions.\par
\noindent{\footnotesize Izvorne zapisy / Original records: HIST-g093-005, HIST-g099-011.}\par
\smallskip\noindent\emph{Sravnjenje s tekstom:} Eksponent i stepen faktora imajut odlišne uloge. Njih produkt daje navedenu mnogokratnost; izpustiti linearnu nezavisnost značilo by poměniti čislo, ktoro opisuje tekst.\par
\noindent\emph{Textual comparison:} Exponent and factor degree have distinct roles. Their product gives the stated multiplicity; omitting linear independence would change the quantity described.\par
\noindent{\footnotesize TEXT-COMP-292; pasaže / passages: 15000, 15005, 15012, 15014, 15016, 15631, 15634, 15635, 15637, 15642; \texttt{CORPUS\_FORM\_COMPARISONS.json}.}\par
\end{minipage}\par
\par\medskip\Needspace{7\baselineskip}
\noindent\begin{minipage}{\linewidth}
\hypertarget{TERM-1194}{}
\noindent\textbf{močnosť}\quad{\footnotesize\texttt{TERM-1194}}\par
\noindent\emph{Iztočnik / Source:} Mächtigkeit\par
\noindent\emph{Rusky / Russian:} мощность\par
\noindent\emph{Ukrajinsky / Ukrainian:} потужність\par
\smallskip\noindent\emph{Razlog:} Močnost označa kardinalnost množiny, ne fyzičnu silu, algebraičnu potenciju ili stepen poljnogo razširjenja. V §8 raboty 9 upotrěbjena jest za postavjenje vzajemno jednoznačnogo odgovaranja medžu osnovami; argument citira Steinitza. Arhivny izbor močnosť prodolžaje už uživanu proektnu terminologiju, a dostupny korpus piše močnost. Razlika v označenju končnogo t jest pravopisna, ne smyslova. Opis mora čuvati etu razliku i pri čitanju latiničnoj i kiriličnoj formy. Jednakost kardinalnostij razgledanyh beskonečnyh množin ne opravdava opču tvrdnju, že vsako algebraično razširjenje, uključajuči razširjenja konečnyh polj, imaje tu samu kardinalnost. Steinitzova citacija jest svedočstvo izvora argumenta v rabotě, ne dokaz ranějšego pročitanja njegovogo teksta pri prěvodě.\par
\noindent\emph{Reason:} Močnost denotes set cardinality, not physical power, an algebraic power or field-extension degree. In §8 of Paper 9 it is used to establish a bijective correspondence between bases; the argument cites Steinitz. The archived choice močnosť continues existing project terminology, while the available corpus writes močnost. The treatment of final t is an orthographic distinction, not a difference of meaning, and should remain recognisable in the Latin and Cyrillic forms. Equality of cardinalities for the infinite sets considered does not justify the unrestricted claim that every algebraic extension, including an extension of finite fields, has the same cardinality. The Steinitz citation identifies the work's source for the argument, not an earlier reading of his text during translation.\par
\noindent{\footnotesize Izvorne zapisy / Original records: HIST-g066-005.}\par
\smallskip\noindent\emph{Sravnjenje s tekstom:} Močnost v etyh pasusah označa kardinalnost. Citacije Steinitza i upotrěba baz svězujut formulaciju s točnym argumentom, ne s fyzičnoju ili algebraičnoju potencijeju.\par
\noindent\emph{Textual comparison:} Močnost in these passages means cardinality. The Steinitz citations and use of bases link the wording to a particular argument, not to physical or algebraic powers.\par
\noindent{\footnotesize TEXT-COMP-298; pasaže / passages: 8061, 8062, 8063, 8064, 8065, 8066, 8067, 8068, 8069, 8070, 8071, 8072, 8073, 8074, 8487, 8489, 8491; \texttt{CORPUS\_FORM\_COMPARISONS.json}.}\par
\end{minipage}\par
\par\medskip\Needspace{7\baselineskip}
\noindent\begin{minipage}{\linewidth}
\hypertarget{TERM-1195}{}
\noindent\textbf{močnosť / kontinuum / dobra uporjadkovanja}\quad{\footnotesize\texttt{TERM-1195}}\par
\noindent\emph{Iztočnik / Source:} Mächtigkeit / Kontinuum / Wohlordnungen\par
\noindent\emph{Rusky / Russian:} мощность / континуум / вполне упорядочения\par
\noindent\emph{Ukrajinsky / Ukrainian:} потужність / континуум / добрі упорядкування\par
\smallskip\noindent\emph{Razlog:} Močnost, kontinuum i dobro uporjadkovanje imajut različne funkcije v §9 raboty 9. Močnost sravnjava veličinu množin; kontinuum davaje kardinal $c$; dobro uporjadkovanje jest porjadok, v ktorom vsaka neprazdna podmnožina imaje najmenši element. Poslědnje svojstvo jest silnějše než obyčny linearny porjadok. Arhivny razlog prodolžaje terminologiju §8 i ostavjaje dobro uporjadkovanje za daljšu jezykovu ocěnku. Tekst daje $2^c$ za čislo razgledanyh konstrukcij, ne avtomatično za čislo različnyh domen ili klasov ekvivalencije: sám ukazuje slučaj s samo odčislimo mnogo neekvivalentnymi domenami. Oznaka kontinuuma ne zahtěvaje hipotezu kontinuuma.\par
\noindent\emph{Reason:} Močnost, kontinuum and dobro uporjadkovanje have different functions in §9 of Paper 9. Močnost compares sizes of sets; kontinuum supplies the cardinal $c$; a well-order is an order in which every nonempty subset has a least element, a stronger property than an ordinary linear order. The archived reason continues the vocabulary of §8 and leaves dobro uporjadkovanje open to further linguistic assessment. The text gives $2^c$ for the number of constructions considered, not automatically for the number of distinct domains or equivalence classes: it itself notes a case with only countably many inequivalent domains. The continuum notation does not require the continuum hypothesis.\par
\noindent{\footnotesize Izvorne zapisy / Original records: HIST-g067-002.}\par
\smallskip\noindent\emph{Sravnjenje s tekstom:} Tekst izrično razlikuje mnohostvo konstrukcij i čislo neekvivalentnyh domen. Rěčnik ne može izbrisati etu razliku pri objasnjenju kontinuuma i dobryh uporjadkovanij.\par
\noindent\emph{Textual comparison:} The text expressly distinguishes the multitude of constructions from the number of inequivalent domains. Explaining continuum and well-orderings must preserve this distinction.\par
\noindent{\footnotesize TEXT-COMP-299; pasaže / passages: 8172, 8173, 8174, 8175, 8176, 8177, 8178, 8179, 8180, 8181; \texttt{CORPUS\_FORM\_COMPARISONS.json}.}\par
\end{minipage}\par
\par\medskip\Needspace{7\baselineskip}
\noindent\begin{minipage}{\linewidth}
\hypertarget{TERM-1196}{}
\noindent\textbf{močnosť kontinuuma}\quad{\footnotesize\texttt{TERM-1196}}\par
\noindent\emph{Iztočnik / Source:} Mächtigkeit des Kontinuums\par
\noindent\emph{Rusky / Russian:} мощность континуума\par
\noindent\emph{Ukrajinsky / Ukrainian:} потужність континууму\par
\smallskip\noindent\emph{Razlog:} Močnost kontinuuma jest rodilnikova konstrukcija, ktora imenuje kardinalnost kontinuuma. V §2 raboty 10 ona nastupaje pri sravnjenju racionalnyh baz za polje vsih realnyh i kompleksnyh čisel i pri postavjenju vzajemno jednoznačnogo odgovaranja. Arhivny razlog jest doslědnost s rabotoju 9; ne sodrži novo sravnjenje alternativ ni čislove vagi. Arhiv piše močnosť kontinuuma, a dostupny tekst močnost kontinuuma. Dobro uporjadkovanje i porjadok elementov v argumentu ne davajut sami po sobě konečnogo algoritma za izpisanje cěloj takoj bazy. Termin imenuje teoretičnu kardinalnost, ne praktičnu proceduru enumeracije.\par
\noindent\emph{Reason:} Močnost kontinuuma is a genitive construction naming the cardinality of the continuum. In §2 of Paper 10 it occurs when comparing rational bases for the field of all real and complex numbers and establishing a bijective correspondence. The archived reason is consistency with Paper 9; it supplies neither a new comparison of alternatives nor numerical weights. The archive writes močnosť kontinuuma and the available text močnost kontinuuma. The well-ordering and ordering of elements in the argument do not by themselves supply a finite algorithm for listing such an entire basis. The term names a theoretical cardinality, not a practical enumeration procedure.\par
\noindent{\footnotesize Izvorne zapisy / Original records: HIST-g070-009.}\par
\smallskip\noindent\emph{Sravnjenje s tekstom:} Močnost v etyh pasusah označa kardinalnost. Citacije Steinitza i upotrěba baz svězujut formulaciju s točnym argumentom, ne s fyzičnoju ili algebraičnoju potencijeju.\par
\noindent\emph{Textual comparison:} Močnost in these passages means cardinality. The Steinitz citations and use of bases link the wording to a particular argument, not to physical or algebraic powers.\par
\noindent{\footnotesize TEXT-COMP-298; pasaže / passages: 8061, 8062, 8063, 8064, 8065, 8066, 8067, 8068, 8069, 8070, 8071, 8072, 8073, 8074, 8487, 8489, 8491; \texttt{CORPUS\_FORM\_COMPARISONS.json}.}\par
\end{minipage}\par
\par\medskip\Needspace{7\baselineskip}
\noindent\begin{minipage}{\linewidth}
\hypertarget{TERM-1204}{}
\noindent\textbf{prirodne iracionalnosti prinadležeče grupě}\quad{\footnotesize\texttt{TERM-1204}}\par
\noindent\emph{Iztočnik / Source:} natürliche zur Gruppe gehörige Irrationalitäten\par
\noindent\emph{Rusky / Russian:} естественные иррациональности, принадлежащие группе\par
\noindent\emph{Ukrajinsky / Ukrainian:} природні ірраціональності, що належать групі\par
\smallskip\noindent\emph{Razlog:} Prirodne iracionalnosti prinaležeče grupě čuva historičny izraz, za ktory rabota 11 davaje posebno oprěděljenje. Arhiv piše prinadležeče; dostupny tekst ima prinaležeče. Izraz ne znači libovoljne generatori. V §1, č. 2 parametry zaměnjajut se racionalnymi funkcijami korenj smotrenyh kako neoprěděljene; te funkcije morajut byti algebraično nezavisne, a polinom nad nastaly poljem mora iměti ukazanu grupu. To, že funkcije sut racionalne v korenjah, ne protivrěči historičnomu nazvanju iracionalnosti: važno jest, po odnošenju k ktoromu polju se govori. Sravnjajeme ruska i ukrajinska fraza zachovajut isto povezanje s grupoj. Razlog vernosti definiciji jest silnějši od želanja zaměniti izraz modernym kratkym terminom. Avtor izrično postavlja silnějše požadanje i potom dostatočne uslovja, ne obče potrěbnu karakterizaciju.\par
\noindent\emph{Reason:} Prirodne iracionalnosti prinaležeče grupě preserves a historical expression given a specific definition in Paper 11. The archive writes prinadležeče; the available text has prinaležeče. It does not mean arbitrary generators. In §1, item 2, the parameters are replaced by rational functions of roots treated as indeterminates; these functions must be algebraically independent, and the polynomial over the resulting field must have the specified group. Rationality in the roots does not contradict the historical name irrationalities: the reference field matters. The Russian and Ukrainian comparanda preserve the same relation to the group. Fidelity to the definition takes precedence over replacing it with a shorter modern term. The author expressly imposes a stronger requirement and obtains sufficient conditions, not a universally necessary characterization.\par
\noindent{\footnotesize Izvorne zapisy / Original records: HIST-g073-013.}\par
\smallskip\noindent\emph{Sravnjenje s tekstom:} Prirodne iracionalnosti sut definovane točnymi uslovjami na parametry, funkcije korenj i grupu. Dostatočnost uslovij ne směje stati potrěbnostju v prěvodu.\par
\noindent\emph{Textual comparison:} Natural irrationalities are defined through exact conditions on parameters, root functions, and the group. Sufficiency of the conditions must not become necessity in translation.\par
\noindent{\footnotesize TEXT-COMP-305; pasaže / passages: 8756, 8758, 8761, 8763, 8765; \texttt{CORPUS\_FORM\_COMPARISONS.json}.}\par
\end{minipage}\par
\par\medskip\Needspace{7\baselineskip}
\noindent\begin{minipage}{\linewidth}
\hypertarget{TERM-1205}{}
\noindent\textbf{dodatno uslovje}\quad{\footnotesize\texttt{TERM-1205}}\par
\noindent\emph{Iztočnik / Source:} Nebenbedingung\par
\noindent\emph{Rusky / Russian:} дополнительное условие\par
\noindent\emph{Ukrajinsky / Ukrainian:} додаткова умова\par
\smallskip\noindent\emph{Razlog:} Dodatno uslovje prědava Nebenbedingung bez doslovnogo obraza bočnogo uslovja. Arhivny razlog vključaje i kontinuitet: v rabotě 9 to jest uslovje konstrukcije racionalnoj bazy, ne novy nazov za hipotezu. Ruska i ukrajinska fraza takože zachovajut značenje dodatnogo uslovja. Po §6 oblast $[\Theta]$, vyvedena iz bazy, ne směje sodrživati algebraično-drobnogo čisla. Dodatno zato ne znači dobrovoljno. V §10 avtor razděljuje karakteristiku nula i karakteristiku $p$; v drugom slučaju uslovje otpadaje iz navedenogo svojstva prvotnogo polja. Ne trěba prěnesti to izključenje na slučaj karakteristiki nula. Izbor jest kvalitativny: jasnost uslovja i jeho svez s ranějšim oprěděljenjem imajut prědnost prěd doslovnym kopirovanjem německoj struktury.\par
\noindent\emph{Reason:} Dodatno uslovje renders Nebenbedingung without the literal image of a side condition. The recorded reason also concerns continuity: in Paper 9 it is a condition on the construction of a rational basis, not a new synonym for hypothesis. The Russian and Ukrainian phrases likewise retain the additional-condition sense. In §6 the domain $[\Theta]$ generated from the basis must contain no algebraically fractional number. Additional therefore does not mean optional. In §10 the author distinguishes characteristic zero from characteristic $p$; in the latter case the condition falls away because of the stated property of the prime field. That exception must not be transferred to characteristic zero. The choice is qualitative: clarity of the condition and its connection to the earlier definition outweigh literal reproduction of the German structure.\par
\noindent{\footnotesize Izvorne zapisy / Original records: HIST-g058-012, HIST-g068-012.}\par
\smallskip\noindent\emph{Sravnjenje s tekstom:} Uslovje izključaje algebraično-drobne čisla iz vyvedenej oblasti. Njegovo otpadenje v karakteristikě $p$ jest odvojeno od slučaja karakteristiki nula.\par
\noindent\emph{Textual comparison:} The condition excludes algebraically fractional numbers from the generated domain. Its disappearance in characteristic $p$ is distinguished from characteristic zero.\par
\noindent{\footnotesize TEXT-COMP-306; pasaže / passages: 7751, 7752, 7753, 7754, 7755, 7756, 7757, 7758, 8382, 8383, 8384, 8385, 8386, 8387; \texttt{CORPUS\_FORM\_COMPARISONS.json}.}\par
\end{minipage}\par
\par\medskip\Needspace{7\baselineskip}
\noindent\begin{minipage}{\linewidth}
\hypertarget{TERM-1206}{}
\noindent\textbf{imenovatelj}\quad{\footnotesize\texttt{TERM-1206}}\par
\noindent\emph{Iztočnik / Source:} Nenner\par
\noindent\emph{Rusky / Russian:} знаменатель\par
\noindent\emph{Ukrajinsky / Ukrainian:} знаменник\par
\smallskip\noindent\emph{Razlog:} Imenovatelj označa imenovatelj racionalnogo izraza, ne nazov ili indeks. Arhivny razlog jest vezany k točnomu dokazu v rabotě 10, §3. Pri zaměně bazovyh veličin jih obrazami muset se zachovati i relacije $H=0$, i uslovja $H\ne0$. Druge potrěbovanje postoji zato, že polinom može nastupati v imenovatelju: nule ne smějut nastati tam, kde dolžen ostati oprěděljen količnik. Rusko znamenatelj i ukrajinsky znamennyk sut sravnjajeme nazvy iste funkcije; sam zapis ne dokazuje, že imenovatelj jest jedin možny vseslovjansky izbor. Formula s čisliteljem i imenovateljem pokazuje, začto kratky nazov potrěbuje točno uslovje nenulovosti. Zaměna frazoj vse funkcije sut nenulove byla by nepravdiva: govori se o zachovanju togo, bylo li dano značenje nula ili ne.\par
\noindent\emph{Reason:} Imenovatelj denotes the denominator of a rational expression, not a name or an index. The recorded reason concerns a precise step in Paper 10, §3. Replacing basis quantities by their images must preserve both relations $H=0$ and conditions $H\ne0$. The latter requirement matters because a polynomial may occur as a denominator: a zero must not appear where the quotient needs to remain defined. Russian znamenatelʹ and Ukrainian znamennyk name the same function; the record alone does not establish imenovatelj as the only possible pan-Slavic choice. The numerator-denominator formula shows why the short term requires the exact nonzero condition. Replacing the statement by all functions are nonzero would be wrong: what is preserved is whether the given value is zero or not.\par
\noindent{\footnotesize Izvorne zapisy / Original records: HIST-g071-006.}\par
\smallskip\noindent\emph{Sravnjenje s tekstom:} Tekst potrěbuje zachovanje nulovyh i nenulovyh relacij. Nenulovost imenovatelja čuva oprěděljenost racionalnogo izraza.\par
\noindent\emph{Textual comparison:} The text requires preservation of both zero and nonzero relations. Nonvanishing of the denominator preserves the definedness of the rational expression.\par
\noindent{\footnotesize TEXT-COMP-307; pasaže / passages: 8533, 8535, 8539, 8540, 8544, 8545, 8547; \texttt{CORPUS\_FORM\_COMPARISONS.json}.}\par
\end{minipage}\par
\par\medskip\Needspace{7\baselineskip}
\noindent\begin{minipage}{\linewidth}
\hypertarget{TERM-1209}{}
\noindent\textbf{ne jest zatvorjena vzhodno k tvorenju pr?s?kov}\quad{\footnotesize\texttt{TERM-1209}}\par
\noindent\emph{Iztočnik / Source:} nicht abgeschlossen gegen?ber Durchschnittsbildung\par
\noindent\emph{Rusky / Russian:} [Bukvy ne sut zachovane / Lettering not recoverable]\par
\noindent\emph{Ukrajinsky / Ukrainian:} [Bukvy ne sut zachovane / Lettering not recoverable]\par
\smallskip\noindent\emph{Razlog:} Arhivny zapis ne zachoval vse bukvy: znaky pytanja v pr?s?kov i cělo poškodžene ruska i ukrajinska polja ne dovoljut obnoviti izvorne formy bez dodatnogo iztočnika. Dostupny tekst raboty 9, §5 odvojeno potvrđuje frazu ne jest zatvorjena po odnošenju k tvorenju prěsěkov. To jest tekstovy dokaz za tekuču formulaciju, ne obnova vseh arhivnyh bukv. Smysl jest algebraična zamknjenost klasa pri operaciji: prěsěk vsih oblastij iz $\mathfrak M(\mathfrak G)$ jest $[H]$, a $[H]$ sam ne naleži tomu klasu. Arhivny razlog požadaje jasnost togo strukturnogo zaključka. Ne ide o topologičnu zamknjenost; iz navedenogo prěsěka vsih členov ne trěba bez daljšego dokaza zaključiti nezamknjenost pri prěsěku dvou členov.\par
\noindent\emph{Reason:} The archival record does not preserve all lettering: question marks in pr?s?kov and wholly damaged Russian and Ukrainian fields prevent certain recovery of the original forms. The available text of Paper 9, §5 separately attests ne jest zatvorjena po odnošenju k tvorenju prěsěkov. This is evidence for the current wording, not recovery of all archival letters. The meaning is algebraic closure of a class under an operation: the intersection of all domains in $\mathfrak M(\mathfrak G)$ is $[H]$, which is not itself a member of that class. The recorded reason calls for an explicit statement of this structural conclusion. Topological closedness is not meant; the stated intersection of all members does not by itself establish failure of closure under binary intersections.\par
\noindent{\footnotesize Izvorne zapisy / Original records: HIST-g063-009.}\par
\smallskip\noindent\emph{Sravnjenje s tekstom:} Prěsěk vsih členov klasa jest $[H]$, ktory ne jest členom togo klasa. To jest izričny tekstovy základ za čtenje poškodženogo arhivnogo izraza, ne obnova jeho izgubljenyh bukv.\par
\noindent\emph{Textual comparison:} The intersection of all class members is $[H]$, which is not itself a member. This supplies explicit textual context for reading the damaged archival phrase, not recovery of its lost lettering.\par
\noindent{\footnotesize TEXT-COMP-309; pasaže / passages: 7672, 7673, 7674, 7676, 7677, 7678; \texttt{CORPUS\_FORM\_COMPARISONS.json}.}\par
\end{minipage}\par
\par\medskip\Needspace{7\baselineskip}
\noindent\begin{minipage}{\linewidth}
\hypertarget{TERM-1210}{}
\noindent\textbf{ne izčezaje identično / identično izčezaje}\quad{\footnotesize\texttt{TERM-1210}}\par
\noindent\emph{Iztočnik / Source:} nicht identisch verschwinden / identisch verschwinden\par
\noindent\emph{Rusky / Russian:} не обращается тождественно в нуль / тождественно обращается в нуль\par
\noindent\emph{Ukrajinsky / Ukrainian:} не зникає тотожно / тотожно зникає\par
\smallskip\noindent\emph{Razlog:} Identično izčezati znači byti nulovym polinomom v ukazanyh neoprěděljenyh, ne izčezati samo pri jednoj izbranoj cěnnosti. Po substituciji koeficientov rabota 20 ispytuje identično izčezanje formy reducibilnosti v ostalyh neoprěděljenyh. Ne izčezati identično zato ne znači nikogda ne byti nula. Arhivny razlog prědava rusku razširjenu frazu obračati se v nulj i kratše ukrajinsko i medžuslovjansko izčezati kako odgovarjajuče izrazy za isto svojstvo. Matematična točnost ima prědnost prěd doslovnym izravnanjem glagolov. Potrěbno jest zachovati takože prědpoloženje, že specializacija ne snižaje stepen polinoma; bez njega iskazana ekvivalencija razložimosti i izčezanja ne jest ta sama. Zapis ne daje čislove vahy ni dokaz čitateljskoj razumlivosti.\par
\noindent\emph{Reason:} To vanish identically means to be the zero polynomial in the specified indeterminates, not merely to vanish at one selected value. After coefficient substitution, Paper 20 tests whether the reducibility form vanishes identically in the remaining indeterminates. Not vanishing identically therefore does not mean never being zero. The recorded reason treats the expanded Russian turn into zero wording and the shorter Ukrainian and Interslavic vanish wording as expressions of the same property. Mathematical precision takes precedence over matching the verbs literally. The condition that specialization does not lower the polynomial's degree must also be retained; without it the stated equivalence between reducibility and vanishing is not the same assertion. The record supplies neither numerical weights nor evidence from reader testing.\par
\noindent{\footnotesize Izvorne zapisy / Original records: HIST-g083-005.}\par
\smallskip\noindent\emph{Sravnjenje s tekstom:} Forma reducibilnosti ispytuje se po koeficientnoj substituciji. V glavnom tvrđenju ostavaje uslovje nesniženja stepena.\par
\noindent\emph{Textual comparison:} The reducibility form is tested after coefficient substitution. The main assertion retains the condition that the degree does not drop.\par
\noindent{\footnotesize TEXT-COMP-310; pasaže / passages: 13805, 13807, 13998, 14000, 14004, 14006, 14009; \texttt{CORPUS\_FORM\_COMPARISONS.json}.}\par
\end{minipage}\par
\par\medskip\Needspace{7\baselineskip}
\noindent\begin{minipage}{\linewidth}
\hypertarget{TERM-1211}{}
\noindent\textbf{ne nužno celočislene funkcije}\quad{\footnotesize\texttt{TERM-1211}}\par
\noindent\emph{Iztočnik / Source:} nicht notwendig ganzzahlige Funktionen\par
\noindent\emph{Rusky / Russian:} не обязательно целочисленные функции\par
\noindent\emph{Ukrajinsky / Ukrainian:} не обов'язково цілочисельні функції\par
\smallskip\noindent\emph{Razlog:} Ne nužno celočislene funkcije zachovaje obseg negacije v německom nicht notwendig: uslovje cěločislennosti ne jest potrěbno, ale něktore dopustime funkcije mogut jego izpolnjati. Prěvod necěločislene funkcije by naprotiv izključil te slučaji. Rusky i ukrajinsky ekvivalent takože izrično neguje potrěbnost. V koncu §3 raboty 9 dostupny tekst piše ne potrěbno cěločislene. Sravnjenje s ranějšeju primětkoju o cělyh algebraičnyh funkcijah pokazyvaje, že ide o koeficientne uslovja v algebraičnom kontekstě. Cěločisleny tu ne možno bez objasnjenja čitati kako prijmajuči samo cěle čislove cěnnosti; takože algebraično cěly ne jest synonym racionalnogo cělogo čisla. Razlog izbora jest zachovati razširjenje dopustimoj klasy i razlikovati njegove uslovja, ne samo prěvesti jedino prislovje.\par
\noindent\emph{Reason:} Ne nužno celočislene funkcije preserves the scope of negation in German nicht notwendig: the integer-coefficient condition is not required, although some permitted functions may satisfy it. Rendering it as non-integer-coefficient functions would instead exclude those cases. The Russian and Ukrainian equivalents likewise negate necessity explicitly. At the end of §3 of Paper 9 the available text reads ne potrěbno cěločislene. Comparison with the earlier footnote on integral algebraic functions shows that coefficient conditions in an algebraic setting are at issue. Cěločisleny must not be read without explanation as integer-valued; nor is algebraically integral synonymous with rational integer. The reason for the choice is to preserve the enlargement of the allowed class and distinguish its conditions, not merely to translate an adverb.\par
\noindent{\footnotesize Izvorne zapisy / Original records: HIST-g061-013.}\par
\smallskip\noindent\emph{Sravnjenje s tekstom:} Ranějša primětka oprědělja koeficientno uslovje; pozdnějši odstavec jego oslabjaje, ne neguje cělost vsih dopustimyh funkcij.\par
\noindent\emph{Textual comparison:} The earlier footnote defines the coefficient condition; the later paragraph relaxes it rather than denying integrality of all permitted functions.\par
\noindent{\footnotesize TEXT-COMP-311; pasaže / passages: 7147, 7266, 7267, 7268, 7269, 7270, 7271, 7272, 7373, 7374, 7376, 7377, 7378, 7379, 7380, 7381; \texttt{CORPUS\_FORM\_COMPARISONS.json}.}\par
\end{minipage}\par
\par\medskip\Needspace{7\baselineskip}
\noindent\begin{minipage}{\linewidth}
\hypertarget{TERM-1213}{}
\noindent\textbf{ne jest popolno reducibilno prvogo roda}\quad{\footnotesize\texttt{TERM-1213}}\par
\noindent\emph{Iztočnik / Source:} nicht vollständig reduzibel erster Art\par
\noindent\emph{Rusky / Russian:} не является вполне приводимым первого рода\par
\noindent\emph{Ukrajinsky / Ukrainian:} не є цілком редуцибельним першого роду\par
\smallskip\noindent\emph{Razlog:} Popolno reducibilny prvogo roda jest sostavny historičny termin. Arhivny izbor reducibil- měl zachovati ranějšu terminologičnu familiju i ne sugerovati moderny pojem redukovanogo prstena. Negacija se odnosi k cělomu svojstvu, uključajuči prvogo roda. V rabotě 31, §§7-8 ono nastupaje pri prstenu klasov ostatkov i izčezanju diskriminanta. §6 davaje kontrolu smysla: poslě ukazanogo razširjenja koeficientnogo polja nerazložime komponenty morajut byti poljami izomorfnymi algebraično zatvorjenomu polju. Dostupny tekst ne jest terminologično jednorodny: §6 piše polna razložimost, a pozdnějše města polno reducibilno. Ta dva stemna rěšenja ne smějut byti smolčana kako už provedena standardizacija. Ruska i ukrajinska fraza takože čuvajut ograničenje prvogo roda; ne jestdna kratka zaměna redukovany ne davaje vse to značenje.\par
\noindent\emph{Reason:} Popolno reducibilny prvogo roda is a compound historical term. The recorded choice of reducibil- was intended to preserve an earlier terminological family without suggesting the modern notion of a reduced ring. The negation applies to the whole property, including of the first kind. In Paper 31, §§7-8 it concerns a residue-class ring and the vanishing of its discriminant. Section 6 controls the meaning: after the specified extension of the coefficient field, the indecomposable components must be fields isomorphic to the algebraically closed field. The available text is not terminologically uniform: §6 has polna razložimost, whereas later passages use polno reducibilno. These two stem choices must not be represented as an already accomplished standardization. The Russian and Ukrainian phrases also retain of the first kind; the short replacement reduced does not convey the entire property.\par
\noindent{\footnotesize Izvorne zapisy / Original records: HIST-g154-005, HIST-g164-005.}\par
\smallskip\noindent\emph{Sravnjenje s tekstom:} §6 i pozdnějše paragrafy uživajut različne osnovy za isto povezane svojstvo. Izrična definicija prědnostno oprědělja smysl; negacija cělosti svojstva ne znači samo redukovanost prstena.\par
\noindent\emph{Textual comparison:} Section 6 and later sections use different stems for the connected property. The explicit definition controls the meaning; negating the whole property is not simply a statement about a ring being reduced.\par
\noindent{\footnotesize TEXT-COMP-313; pasaže / passages: 18412, 18441, 18751, 19105; \texttt{CORPUS\_FORM\_COMPARISONS.json}.}\par
\end{minipage}\par
\par\medskip\Needspace{7\baselineskip}
\noindent\begin{minipage}{\linewidth}
\hypertarget{TERM-1214}{}
\noindent\textbf{neizolovane (imbedded) primarne idealy}\quad{\footnotesize\texttt{TERM-1214}}\par
\noindent\emph{Iztočnik / Source:} nicht-isolierte (imbedded) primäre Ideale\par
\noindent\emph{Rusky / Russian:} неизолированные (imbedded) примарные идеалы\par
\noindent\emph{Ukrajinsky / Ukrainian:} неізольовані (imbedded) первинні ідеали\par
\smallskip\noindent\emph{Razlog:} Neizolovane primarne idealy označaje komponenty v teoriji razloženja idealov, ne neizolovane točky topologičnogo prostora. Pri tomu ostavaje historično anglijsko imbedded v skobkah. Arhivny razlog čuva citovanu formu, ne zaměnjaje jej tiho modernym pravopisom embedded i ne prědstavjaje jej kako novo medžuslovjansko slovo. V dostupnom tekstě izraz stoi v primětce k teoremu XIII raboty 19, §7; arhivne oznaky jedinic ne sut same po sobě točny naslov paragrafa. Teorem pripisuje možno mnogoznačno prědstavjenje neizolovanym komponentam, togda kak izolovane sut jednoznačno oprěděljene. To ne znači, že vsaka neizolovana komponenta mora byti mnogoznačna: priměry v §11 izrično dopuščajut jednoznačnost i tutyh komponent. Macaulay jest citovany avtorom; ta dokumentacija ne znači, že jeho original byl odvojeno pročitan.\par
\noindent\emph{Reason:} Neizolovane primarne idealy denotes components in ideal-decomposition theory, not non-isolated points of a topological space. The historical English imbedded remains in parentheses. The recorded reason preserves the quoted form rather than silently modernizing it to embedded or presenting it as a new Interslavic word. In the available text the expression occurs in the footnote to Theorem XIII of Paper 19, §7; archival unit labels are not by themselves exact section headings. The theorem locates possible ambiguity in the non-isolated components, while isolated components are uniquely determined. It does not say that every non-isolated component must be ambiguous: the examples in §11 expressly allow uniqueness for those components too. Macaulay is cited by the author; this documentation does not imply a separate reading of his original.\par
\noindent{\footnotesize Izvorne zapisy / Original records: HIST-g081-060, HIST-g082-060.}\par
\smallskip\noindent\emph{Sravnjenje s tekstom:} Teorem XIII izločuje izolovane komponenty iz možné mnogoznačnosti. Priměry v §11 pokazyvajut, že i neizolovane komponenty mogut byti jednoznačno oprěděljene.\par
\noindent\emph{Textual comparison:} Theorem XIII excludes isolated components from possible ambiguity. Examples in §11 show that non-isolated components can also be uniquely determined.\par
\noindent{\footnotesize TEXT-COMP-314; pasaže / passages: 13212, 13242, 13244, 13246, 13252, 13651, 13656, 13664; \texttt{CORPUS\_FORM\_COMPARISONS.json}.}\par
\end{minipage}\par
\par\medskip\Needspace{7\baselineskip}
\noindent\begin{minipage}{\linewidth}
\hypertarget{TERM-1215}{}
\noindent\textbf{neizomorfne ili suštno različne}\quad{\footnotesize\texttt{TERM-1215}}\par
\noindent\emph{Iztočnik / Source:} nicht-isomorph oder wesentlich verschieden\par
\noindent\emph{Rusky / Russian:} неизоморфны, или существенно отличны\par
\noindent\emph{Ukrajinsky / Ukrainian:} неізоморфні, або суттєво відмінні\par
\smallskip\noindent\emph{Razlog:} Neizomorfne ili suštno različne svęzuje formalny termin s objasnjujučim nazvanjem iz prědhodečego paragrafa. Dostupny tekst piše sučstveno različne. V rabotě 9, §§8-9 ili v toj frazě davaje druge nazvanje istogo odnosa, ne slabu alternativu medžu neizomorfijoju i obyčnoju neravnostju množstv. Različne zapisy ili različne elementy ješče ne dokazyvajut neizomorfiju. Dokaz sravnjaje algebraične svojstva oblastij s različnymi indeksami; zato zaključek opisuje različne klasy po izomorfiji. Slědujuča primětka ješče različa izomorfiju cělyh oblastij od izomorfije s čestju drugoj oblasti: vsaka iz dvou neizomorfnyh oblastij može byti izomorfna česti drugoj. Arhivny razlog jest kontinuitet togo pojma, ne nova empirčna ocěnka razumlivosti slova suštno.\par
\noindent\emph{Reason:} Neizomorfne ili suštno različne links a formal term to the explanatory name used in the preceding section. The available text writes sučstveno različne. In Paper 9, §§8-9, ili in this phrase supplies another name for the same relation, not a weaker alternative between non-isomorphism and ordinary inequality of sets. Different descriptions or different elements do not alone establish non-isomorphism. The proof compares algebraic properties of domains with different indices, so its conclusion concerns distinct isomorphism classes. The following note further distinguishes isomorphism of the whole domains from isomorphism with part of the other domain: each of two non-isomorphic domains may be isomorphic to part of the other. The recorded reason is continuity of this concept, not a new empirical assessment of the intelligibility of suštno.\par
\noindent{\footnotesize Izvorne zapisy / Original records: HIST-g067-011.}\par
\smallskip\noindent\emph{Sravnjenje s tekstom:} Sučstveno različne objasnjaje neizomorfne. Slědujuča primětka ne dopuščaje směšati izomorfiju cělosti s izomorfijoju s čestju drugoj oblasti.\par
\noindent\emph{Textual comparison:} Essentially different explains non-isomorphic. The following note prevents confusion between isomorphism of the wholes and isomorphism with part of the other domain.\par
\noindent{\footnotesize TEXT-COMP-315; pasaže / passages: 8075, 8076, 8077, 8078, 8079, 8080, 8294, 8295, 8296, 8297, 8298, 8299, 8301, 8302, 8303, 8304; \texttt{CORPUS\_FORM\_COMPARISONS.json}.}\par
\end{minipage}\par
\par\medskip\Needspace{7\baselineskip}
\noindent\begin{minipage}{\linewidth}
\hypertarget{TERM-1216}{}
\noindent\textbf{nekomutativna kolceva oblast / nekomutativno kolco}\quad{\footnotesize\texttt{TERM-1216}}\par
\noindent\emph{Iztočnik / Source:} nicht-kommutativer Ringbereich / nicht-kommutativer Ring\par
\noindent\emph{Rusky / Russian:} некоммутативная кольцевая область / некоммутативное кольцо\par
\noindent\emph{Ukrajinsky / Ukrainian:} некомутативна кільцева область / некомутативне кільце\par
\smallskip\noindent\emph{Razlog:} Nekomutativna kolceva oblast i nekomutativno kolco zachovajut razliku medžu německym Ringbereich i Ring; dostupny tekst uživa prstenova oblast i prsten. Arhivny razlog jest, že ne možno tiho prěnesti komutativne prědpoloženja v nekomutativny kontekst. V §9 raboty 19 definicija vypuščaje zakon komutativnosti množenja; to ne znači, že vsake dva elementy muset nekomutovati. Slovo oblast samo ne ustanovjaje odsutnost nulovyh děliteljev. V §12 priměrom jest oblast cěločislenyh matric, s jednostrannymi idealami i dvustrannymi idealami kako specialnym slučajem. Čislove vahy za kolco protiv prsten v tutom zapisu ne sut dane. Vybor imennoj osnovy i soglasovanje po rodu trěba razlikovati od matematičnogo svojstva; nekomutativnost ne směje zakryti ni strannost idealov, ni ostale izrične prědpoloženja.\par
\noindent\emph{Reason:} Nekomutativna kolceva oblast and nekomutativno kolco retain the distinction between German Ringbereich and Ring; the available text uses prstenova oblast and prsten. The recorded reason is that commutative assumptions must not be silently imported into a noncommutative setting. In §9 of Paper 19 the definition omits the commutative law of multiplication; this does not mean that every two elements must fail to commute. Oblast alone does not establish the absence of zero divisors. In §12 the example is the domain of integer matrices, with one-sided ideals and two-sided ideals as a special case. This record supplies no numerical weights for kolco against prsten. Choice of noun stem and grammatical gender agreement must be distinguished from the mathematical property; noncommutativity must obscure neither sidedness of ideals nor other explicit hypotheses.\par
\noindent{\footnotesize Izvorne zapisy / Original records: HIST-g081-073.}\par
\smallskip\noindent\emph{Sravnjenje s tekstom:} Rabota 19 vključaje matričnu oblast i jednostranne idealy. Rabota 17 odvojeno postavlja zakony polinomnoj oblasti, iz ktoryh vyvodi nenulovost produkta. Slovo oblast samo ne zaměnjaje te zakony.\par
\noindent\emph{Textual comparison:} Paper 19 includes a matrix domain and one-sided ideals. Paper 17 separately imposes the laws of its polynomial domain from which nonvanishing of a product follows. Oblast alone does not replace those laws.\par
\noindent{\footnotesize TEXT-COMP-316; pasaže / passages: 13360, 13362, 13364, 13390, 13409, 13691, 13695, 10955, 10957, 11044, 11048, 11050, 11055, 11057, 11059, 11062; \texttt{CORPUS\_FORM\_COMPARISONS.json}.}\par
\end{minipage}\par
\par\medskip\Needspace{7\baselineskip}
\noindent\begin{minipage}{\linewidth}
\hypertarget{TERM-1217}{}
\noindent\textbf{neracionalna involucija}\quad{\footnotesize\texttt{TERM-1217}}\par
\noindent\emph{Iztočnik / Source:} nicht-rationale Involution\par
\noindent\emph{Rusky / Russian:} нерациональная инволюция\par
\noindent\emph{Ukrajinsky / Ukrainian:} нераціональна інволюція\par
\smallskip\noindent\emph{Razlog:} Neracionalna involucija čuva termin geometričnogo argumenta, ne obyčno značenje nerazumnosti. Arhivna motivacija odsyla k Enriquesovomu priměru, ktory Noether navodi v §6 raboty 6. V tom kontekstě minimalna baza polja $\mathfrak{K}_{n\rho}$ imaje $\rho$ algebraično nezavisnyh racionalnyh generatorov; citovany priměr daje polje $\mathfrak{K}_{33}$ bez takoj bazy. Odgovarjajuče ruska i ukrajinska forma zachovajut internacionalny korenj racional-. Ne jest korektno zaměniti historičnu geometričnu involuciju bez dodatnogo objasnjenja prostym modernym opisom funkcija, čije dvojno priměnjenje jest identičnost: tuto město svęzuje involuciju s poljem funkcij i jego racionalnostju. Takože priměr ne govori, že vsako polje ranga tri ne imaje minimalnoj bazy. Zapis dokumentuje upotrěbu i citaciju v Noetherovom tekstě, ne nezavisnu prověrku Enriquesovogo dokaza.\par
\noindent\emph{Reason:} Neracionalna involucija preserves the term in the geometric argument, not the everyday sense of irrational judgment. The archived motivation points to the example by Enriques cited in §6 of Paper 6. Here a minimal basis of $\mathfrak{K}_{n\rho}$ has $\rho$ algebraically independent rational generators; the cited example supplies a field $\mathfrak{K}_{33}$ without such a basis. The corresponding Russian and Ukrainian forms retain the international racional- root. It would be inappropriate to replace the historical geometric involution, without further explanation, by the simple modern description of a function whose square is the identity: this passage connects involution with a function field and its rationality. Nor does the example say that every field of rank three lacks a minimal basis. The record documents usage and citation in Noether's text, not an independent verification of Enriques's proof.\par
\noindent{\footnotesize Izvorne zapisy / Original records: HIST-g046-012.}\par
\smallskip\noindent\emph{Sravnjenje s tekstom:} Noether svęzuje Enriquesov priměr s poljem bez minimalnoj bazy. Opis jest ograničeny tym priměrom, ne vsimi poljami istogo ranga.\par
\noindent\emph{Textual comparison:} Noether connects the example by Enriques to a field without a minimal basis. The description is restricted to that example, not all fields of the same rank.\par
\noindent{\footnotesize TEXT-COMP-317; pasaže / passages: 5688, 5692, 5694, 5700, 5712; \texttt{CORPUS\_FORM\_COMPARISONS.json}.}\par
\end{minipage}\par
\par\medskip\Needspace{7\baselineskip}
\noindent\begin{minipage}{\linewidth}
\hypertarget{TERM-1218}{}
\noindent\textbf{neobstajanje eksplicitnogo predstavjenja}\quad{\footnotesize\texttt{TERM-1218}}\par
\noindent\emph{Iztočnik / Source:} Nichtexistenz der expliziten Darstellung\par
\noindent\emph{Rusky / Russian:} несуществование явного представления\par
\noindent\emph{Ukrajinsky / Ukrainian:} неіснування явного подання\par
\smallskip\noindent\emph{Razlog:} Neobstajanje eksplicitnogo prědstavjenja jest tvrđenje o točno definovanom vidu prědstavjenja, ne izvěstje o neuspešnom iskanju formuly. Arhivny zapis veže termin k §2 raboty 4: eksplicitno prědstavjenje uživa cěle redy simbolov ili jednoznačno iz njih vyvedene redy, ne razložene častne redy. V §3 teorema II ješče razlikuje prědstavjenje bez substitucije (11). Zaključna negacija ograničena jest razgledanymi identičnostjami i uslovjem, že medžu redami prěměnnyh ne jest red $x$ ili $u$. Ne možno zato prěvesti frazu kako ne obstaje nikako prědstavjenje. Rusky javny i ukrajinsky javny izraz sut v arhivu sravnjajemi s internacionalnym eksplicitny; razlog medžuslovjanskogo izbora jest kontinuitet definicije. Dostupny tekst dopolnjaje pravopis prědstavjenje, kde arhiv imaje predstavjenje.\par
\noindent\emph{Reason:} Neobstajanje eksplicitnogo prědstavjenja asserts nonexistence of a precisely defined kind of representation, not failure to find a formula. The archival record connects the term to §2 of Paper 4: explicit representation uses whole symbol rows, or rows uniquely derived from them, rather than their decomposed constituent rows. In §3 Theorem II further distinguishes representation without substitution (11). The concluding negative statement is restricted to the identities considered and the condition that no variable row is an $x$ or $u$ row. It therefore cannot be rendered as no representation exists. The archive compares the Russian and Ukrainian explicit wordings with international eksplicitny; continuity of the definition is the stated reason for the Interslavic choice. The available text writes prědstavjenje where the archive has predstavjenje.\par
\noindent{\footnotesize Izvorne zapisy / Original records: HIST-g032-013.}\par
\smallskip\noindent\emph{Sravnjenje s tekstom:} Oprěděljenje v §2 i uslovja teoremy v §3 ograničavajut neobstajanje. Ne ide o neobstajanje vsakoj formuly ili vsake reprezentacije.\par
\noindent\emph{Textual comparison:} The definition in §2 and the conditions of the theorem in §3 delimit the nonexistence claim. It is not nonexistence of every formula or every representation.\par
\noindent{\footnotesize TEXT-COMP-318; pasaže / passages: 3809, 4026, 4055; \texttt{CORPUS\_FORM\_COMPARISONS.json}.}\par
\end{minipage}\par
\par\medskip\Needspace{7\baselineskip}
\noindent\begin{minipage}{\linewidth}
\hypertarget{TERM-1219}{}
\noindent\textbf{necěly element}\quad{\footnotesize\texttt{TERM-1219}}\par
\noindent\emph{Iztočnik / Source:} nichtganzes Element\par
\noindent\emph{Rusky / Russian:} нецелый элемент\par
\noindent\emph{Ukrajinsky / Ukrainian:} нецілий елемент\par
\smallskip\noindent\emph{Razlog:} Necěly element potrěbuje ukazanje osnovnogo prstena i smysla cělosti. Arhivna motivacija kaže, že izraz prodolžaje familiju cěly iz §1 raboty 30. Tam cělost po odnošenju k $R$ jest oprěděljena konečnostju vyvedenogo poredka i ekvivalentno moničnym polinomnym uravnenjem s koeficientami v $R$. Ale §9, č. 7 dokazyvaje samu integralnu zamknjenost. Početkovo necěly v tom dokazu mora označati element izvan osnovnogo prstena, dlja ktorogo se potom izključa cělost po odnošenju k němu; inače by zaključek byl už prědpoloženy. Tuto čtenje podpiraje uvod raboty, govoreči o pravo drobnom elementu. Kvocient $a/c$ s nepodělimymi stepenjami čislitelja služi k proturěčju s moničnym uravnenjem. Bez prověrky německogo města ne trěba proglasiti razliku slov izvornoju pogrěškoju; slovnik mora jasno razděljati dva smysla, ne směšati jih s necělym racionalnym čislom.\par
\noindent\emph{Reason:} Necěly element requires the base ring and the relevant sense of integrality to be specified. The archival motivation says that it continues the cěly family from §1 of Paper 30. There, integrality over $R$ is defined by finiteness of the generated order and equivalently by a monic polynomial equation with coefficients in $R$. But §9, item 7 is proving integral closedness itself. At the beginning of that argument necěly must refer to an element outside the base ring, for which integrality over that ring is then excluded; otherwise the conclusion would already be assumed. The introduction supports this reading by referring to a properly fractional element. A quotient $a/c$ whose numerator powers are not divisible by the denominator supplies the contradiction with a monic equation. Without checking the German passage, the wording difference should not be declared an original error; the dictionary must distinguish these two senses rather than conflate them with a non-integer rational number.\par
\noindent{\footnotesize Izvorne zapisy / Original records: HIST-g116-012.}\par
\smallskip\noindent\emph{Sravnjenje s tekstom:} §1 oprědělja cělost po odnošenju k $R$. Uvod i dokaz aksioma V razlikujut pravo drobny element od možno cělogo elementa nad $R$; prěvod dolžen čuvati to razlikovanje, da ne prědpoloži zaključek.\par
\noindent\emph{Textual comparison:} Section 1 defines integrality over $R$. The introduction and proof of axiom V distinguish a properly fractional element from possible integrality over $R$; the translation must preserve that distinction rather than assume the conclusion.\par
\noindent{\footnotesize TEXT-COMP-319; pasaže / passages: 16124, 16137, 16200, 16208, 16210, 16212, 17109, 17119, 17121, 17123, 17125, 17127; \texttt{CORPUS\_FORM\_COMPARISONS.json}.}\par
\end{minipage}\par
\par\medskip\Needspace{7\baselineskip}
\noindent\begin{minipage}{\linewidth}
\hypertarget{TERM-1220}{}
\noindent\textbf{nekomutativne oblasti}\quad{\footnotesize\texttt{TERM-1220}}\par
\noindent\emph{Iztočnik / Source:} nichtkommutative Bereiche\par
\noindent\emph{Rusky / Russian:} некоммутативные области\par
\noindent\emph{Ukrajinsky / Ukrainian:} некомутативні області\par
\smallskip\noindent\emph{Razlog:} Nekomutativne oblasti jest šire nazvanje algebraičnyh oblastij prěd oprěděljenjem konkretnoj polinomnoj oblasti $J$ v rabotě 17. To jest izričny arhivny razlog za oblast, sravnjajemu s ruskym i ukrajinskym oblastj. Slovo samo ne ustanovjaje vse aksiomy. V §2 koeficientno polje jest komutativno; nekomutativnost može nastati v množenju neizvěstnyh s koeficientami. Pozdnějše prědpoloženja govorjat o vzajemnoj komutativnosti neizvěstnyh i o osobom zakoně množenja, davajučem aditivnost stepenov. Tek te zakony podpirajut zaključek, že produkt dvou nenulovyh polinomov ne jest nula. Ne trěba ni prěložiti vse te prědpoloženja v samo slovo oblast, ni prenesti vse teoremy na libovoljny nekomutativny prsten. To jest svez terminologije s točnym obsegom definicije, ne obča unifikacija vseh historičnyh upotrěbjenij Bereich.\par
\noindent\emph{Reason:} Nekomutativne oblasti is the broader designation for algebraic domains before the particular polynomial domain $J$ is defined in Paper 17. That is the express archival reason for oblast, compared with Russian and Ukrainian oblastʹ. The word alone does not establish all the axioms. In §2 the coefficient field is commutative; noncommutativity can arise when indeterminates multiply coefficients. Later assumptions concern mutual commutativity of the indeterminates and a special multiplication law yielding additivity of degree. Those laws support the conclusion that a product of two nonzero polynomials is nonzero. The translator should neither pack all these assumptions into oblast itself nor extend every theorem to an arbitrary noncommutative ring. This is a connection between terminology and the exact scope of a definition, not a universal unification of historical uses of Bereich.\par
\noindent{\footnotesize Izvorne zapisy / Original records: HIST-g079-002.}\par
\smallskip\noindent\emph{Sravnjenje s tekstom:} Rabota 19 vključaje matričnu oblast i jednostranne idealy. Rabota 17 odvojeno postavlja zakony polinomnoj oblasti, iz ktoryh vyvodi nenulovost produkta. Slovo oblast samo ne zaměnjaje te zakony.\par
\noindent\emph{Textual comparison:} Paper 19 includes a matrix domain and one-sided ideals. Paper 17 separately imposes the laws of its polynomial domain from which nonvanishing of a product follows. Oblast alone does not replace those laws.\par
\noindent{\footnotesize TEXT-COMP-316; pasaže / passages: 13360, 13362, 13364, 13390, 13409, 13691, 13695, 10955, 10957, 11044, 11048, 11050, 11055, 11057, 11059, 11062; \texttt{CORPUS\_FORM\_COMPARISONS.json}.}\par
\end{minipage}\par
\par\medskip\Needspace{7\baselineskip}
\noindent\begin{minipage}{\linewidth}
\hypertarget{TERM-1221}{}
\noindent\textbf{nekomutativno vstavjenje}\quad{\footnotesize\texttt{TERM-1221}}\par
\noindent\emph{Iztočnik / Source:} nichtkommutative Einbettung\par
\noindent\emph{Rusky / Russian:} некоммутативное вложение\par
\noindent\emph{Ukrajinsky / Ukrainian:} некомутативне вкладення\par
\smallskip\noindent\emph{Razlog:} Nekomutativno vstavjenje označaje vstavjenje polja v nekomutativnu algebraičnu strukturu; ne znači, že samo polje prestaje byti komutativno. Arhivny razlog jest kontinuitet s ranějšim vstavjenje i prěvod Einbettung v modernom smyslu. Rusky vloženije i ukrajinsky vkladennja sut zapisane paralely, ne rezultaty testa razumlivosti. V rabotě 32 konstrukcija prědstavjaje razpadno polje kako podpolje matričnogo prstena nad tělom $A$. To objasnjaje, čto dolžen čitatelj razuměti pod vstavjenjem: izomorfno prědstavjenje vnutri veče struktury. Ista osnova v drugih městah korpusa može označati vstavjenje cěnnostij v formulu. Zato kontinuitet slova ne osvobadžaje od razlikovanja pojmov. Tut ne trěba ni prěvoditi vstavjenje kako obyčnu substituciju, ni dodavati tvrdženje, že vsako razpadno polje vhodi neposrědnje v samo $A$.\par
\noindent\emph{Reason:} Nekomutativno vstavjenje denotes embedding a field in a noncommutative algebraic structure; it does not make the embedded field noncommutative. The preserved reason is continuity with earlier vstavjenje and a modern reading of Einbettung. Russian vlozhenie and Ukrainian vkladennia are recorded comparanda, not comprehension-test results. In Paper 32 the construction realizes a splitting field as a subfield of a matrix ring over a division ring $A$. This explains what embedding must convey: an isomorphic realization inside a larger structure. Elsewhere in the corpus the same stem can denote substitution of values into a formula. Lexical continuity therefore does not remove the need to distinguish concepts. Here the reader must neither understand ordinary substitution nor infer that every splitting field embeds directly in $A$ itself.\par
\noindent{\footnotesize Izvorne zapisy / Original records: HIST-g168-007.}\par
\smallskip\noindent\emph{Sravnjenje s tekstom:} Razpadno polje, pridruženo tělo i cěly matričny prsten imajut različne ulohy. Izomorfno vstavjenje čuva tu razliku.\par
\noindent\emph{Textual comparison:} The splitting field, associated division ring, and whole matrix ring have different roles. An isomorphic embedding preserves that distinction.\par
\noindent{\footnotesize TEXT-COMP-320; pasaže / passages: 19244, 19314, 19352, 19356; \texttt{CORPUS\_FORM\_COMPARISONS.json}.}\par
\end{minipage}\par
\par\medskip\Needspace{7\baselineskip}
\noindent\begin{minipage}{\linewidth}
\hypertarget{TERM-1222}{}
\noindent\textbf{nekomutativne těla}\quad{\footnotesize\texttt{TERM-1222}}\par
\noindent\emph{Iztočnik / Source:} nichtkommutative Körper\par
\noindent\emph{Rusky / Russian:} некоммутативные тела\par
\noindent\emph{Ukrajinsky / Ukrainian:} некомутативні тіла\par
\smallskip\noindent\emph{Razlog:} Nekomutativne těla jest izrična terminologična odluka za historično německo Körper v rabotě 32. Zapis različa tělo, v ktorom nenulove elementy imajut multiplikativne obraty, od komutativnogo polja. Rusky tělo i ukrajinsky tilo služe kako sravnjeni priměry; ta dva zapisa sama ne dokazyvajut obče prijetu medžuslovjansku normu. Zato osnova tělo označena jest kako projektna tvorba, potrěbujuča kontekstne prověrky. Pridavnik nekomutativny čuva važnu razliku v temě raboty, ale ne znači, že vsaka para elementov ne komutuje. Centar ostava komutativnym poljem. Odluka dolžna čuvati različne algebraične objekty, ne mehanično zaměnjati vsako Körper slovom tělo, vključajuči města, kde iztočnik govori o komutativnom polju.\par
\noindent\emph{Reason:} Nekomutativne těla is an explicit terminological decision for historical German Körper in Paper 32. The record distinguishes tělo, whose nonzero elements have multiplicative inverses, from a commutative polje. Russian telo and Ukrainian tilo provide the recorded parallels; those two entries do not establish a generally accepted Interslavic norm. The choice is therefore a project proposal requiring contextual consistency checks. The adjective noncommutative preserves an important distinction in the paper, but does not mean that every pair of elements fails to commute. The center remains a commutative field. The policy must preserve the different algebraic objects, not mechanically replace every occurrence of Körper by tělo, including passages about commutative fields.\par
\noindent{\footnotesize Izvorne zapisy / Original records: HIST-g168-002.}\par
\smallskip\noindent\emph{Sravnjenje s tekstom:} Razpadno polje, pridruženo tělo i cěly matričny prsten imajut različne ulohy. Izomorfno vstavjenje čuva tu razliku.\par
\noindent\emph{Textual comparison:} The splitting field, associated division ring, and whole matrix ring have different roles. An isomorphic embedding preserves that distinction.\par
\noindent{\footnotesize TEXT-COMP-320; pasaže / passages: 19244, 19314, 19352, 19356; \texttt{CORPUS\_FORM\_COMPARISONS.json}.}\par
\end{minipage}\par
\par\medskip\Needspace{7\baselineskip}
\noindent\begin{minipage}{\linewidth}
\hypertarget{TERM-1224}{}
\noindent\textbf{nekomutativno tělo}\quad{\footnotesize\texttt{TERM-1224}}\par
\noindent\emph{Iztočnik / Source:} nichtkommutativer Körper\par
\noindent\emph{Rusky / Russian:} некоммутативное тело\par
\noindent\emph{Ukrajinsky / Ukrainian:} некомутативне тіло\par
\smallskip\noindent\emph{Razlog:} Nekomutativno tělo jest jednična forma togo samogo izbora kako nekomutativne těla. Dva arhivne zapisy svęzujut ju s občimi razpadnymi poljami i redukcijnym teoremom raboty 32. Njih razlog jest točna razlika medžu tělom i komutativnym poljem; ocěnka, že tělo jest najjasnější obšteslovjansky kalk, ostava uredničnym sudom, ne izměrenym rezultatom. Dostupny tekst različa tělo $A$ od prstena $A_r$ vseh matric reda $r$ nad $A$. Za $r>1$ ne vsaka nenulova matrica jest obratima, tako že pridruženo tělo ne jest drugo nazvanje cělogo matričnogo prstena. Jedna terminologična osnova dolžna sledovati isto značenje prěz padeže i čisla, ale ne zakryvati tutu strukturnu razliku.\par
\noindent\emph{Reason:} Nekomutativno tělo is the singular of the same choice recorded as nekomutativne těla. Two historical records connect it with the general splitting fields and reduction theorem of Paper 32. Their reason is the precise distinction between a division ring and a commutative field; the description of tělo as the clearest pan-Slavic calque remains editorial judgment, not a measured result. The available text distinguishes the division ring $A$ from the ring $A_r$ of all matrices of order $r$ over $A$. For $r>1$ not every nonzero matrix is invertible, so associated division ring is not another name for the entire matrix ring. A common term should retain its sense across cases and numbers without concealing this structural distinction.\par
\noindent{\footnotesize Izvorne zapisy / Original records: HIST-g169-001, HIST-g171-001.}\par
\smallskip\noindent\emph{Sravnjenje s tekstom:} Razpadno polje, pridruženo tělo i cěly matričny prsten imajut različne ulohy. Izomorfno vstavjenje čuva tu razliku.\par
\noindent\emph{Textual comparison:} The splitting field, associated division ring, and whole matrix ring have different roles. An isomorphic embedding preserves that distinction.\par
\noindent{\footnotesize TEXT-COMP-320; pasaže / passages: 19244, 19314, 19352, 19356; \texttt{CORPUS\_FORM\_COMPARISONS.json}.}\par
\end{minipage}\par
\par\medskip\Needspace{7\baselineskip}
\noindent\begin{minipage}{\linewidth}
\hypertarget{TERM-1225}{}
\noindent\textbf{nekomutativny dvustranny slu?aj}\quad{\footnotesize\texttt{TERM-1225}}\par
\noindent\emph{Iztočnik / Source:} nichtkommutativer zweiseitiger Fall\par
\noindent\emph{Rusky / Russian:} [Bukvy ne sut zachovane / Lettering not recoverable]\par
\noindent\emph{Ukrajinsky / Ukrainian:} [Bukvy ne sut zachovane / Lettering not recoverable]\par
\smallskip\noindent\emph{Razlog:} Dvustranny jest suštna čest frazy: podčrtna primětka k rabotě 17 govori o nekomutativnom dvustrannom slučaju, svęzanom s citovanym rezultatom Schmeidlera. Ona ne prěnosi bez uslovij komutativny dokaz na libovoljny jednostranny slučaj. Pokazany produkt elementov različnih podgrup izčezaje, i ta vlastnost podpiraje dokaz jednoznačnosti razklada. Arhivny razlog izrično zahtěvaje sačuvanje dvustrannosti. Arhivna forma slu?aj i ošteđene kirilične paralely ne mogut služiti kako polno sačuvane attestacije. Dostupna fraza v nekomutativnom dvustrannom slučaju potvrđaje čitljivu formu slučaj v prěvodu, ale ne vosstanavjaje izgubljene bukvy iz drugyh zapisov. Schmeidlerova rabota tu jest citovana v korpusu; to ne znači, že ona byla nezavisno pročitana za tuto slovarnu statju.\par
\noindent\emph{Reason:} Dvustranny is essential: the footnote in Paper 17 concerns a noncommutative two-sided case associated with the cited result of Schmeidler. It does not transfer the commutative argument without conditions to an arbitrary one-sided setting. The displayed products of elements from different subgroups vanish, supporting the uniqueness argument for the decomposition. Preserving two-sidedness is the express archival reason. The damaged archival slu?aj and Cyrillic comparanda cannot count as fully preserved attestations. The available phrase v nekomutativnom dvustrannom slučaju attests the readable form slučaj in the translation, but does not recover the lost letters in other records. Schmeidler's work is cited by the corpus; this is not a claim to have independently read it for this entry.\par
\noindent{\footnotesize Izvorne zapisy / Original records: HIST-g079-042.}\par
\smallskip\noindent\emph{Sravnjenje s tekstom:} Dvustrannost i izčezanje ukazanyh produktov sut česti konkretnogo dokaza, ne stilistične dopolnjenja.\par
\noindent\emph{Textual comparison:} Two-sidedness and the specified vanishing products belong to the particular argument, not to stylistic embellishment.\par
\noindent{\footnotesize TEXT-COMP-322; pasaže / passages: 11527, 11529, 11531; \texttt{CORPUS\_FORM\_COMPARISONS.json}.}\par
\end{minipage}\par
\par\medskip\Needspace{7\baselineskip}
\noindent\begin{minipage}{\linewidth}
\hypertarget{TERM-1226}{}
\noindent\textbf{nelinearne identične relacije}\quad{\footnotesize\texttt{TERM-1226}}\par
\noindent\emph{Iztočnik / Source:} nichtlineare identische Beziehungen\par
\noindent\emph{Rusky / Russian:} нелинейные тождественные соотношения\par
\noindent\emph{Ukrajinsky / Ukrainian:} нелінійні тотожні співвідношення\par
\smallskip\noindent\emph{Razlog:} Nelinearne identične relacije označajut identičnosti medžu komponentami, ne tvrđenje, že te komponenty sut identične jedna drugoj. Arhivny razlog prědpočita terminologiju relacije i identičnosti bez dodavanja statusa samostojnoj teoremy. V §1 raboty 4 relacije medžu elementami $p_\rho$ omogućajut normovati symbole $S^\rho$ tako, da oni izpolnjajut te same identičnosti kako determinanty ukazanoj matrice. Nelinearny opisuje formu tyh relacij; on ne neguje vse linearne transformacijske zakony, prisutne v istom razvažanju. Ruska i ukrajinska paralela takože sodrživajut termin za identičnost. Za razuměnje potrěbno jest čitati te slova vměstě: relacija, ktora drži identično, ne prosto podobnost ili jednako imenovanje objektov.\par
\noindent\emph{Reason:} Nelinearne identične relacije denotes identities among components, not a claim that the components are identical to one another. The recorded reason favors the terminology of relation and identity without assigning an additional theorem-level status. In §1 of Paper 4, relations among the elements $p_\rho$ permit normalization of the symbols $S^\rho$ so that they satisfy the same identities as determinants of the specified matrix. Nonlinear describes these relations; it does not negate all the linear transformation laws used in the argument. The Russian and Ukrainian comparanda likewise express identity. The words must be read together: a relation that holds identically, not merely resemblance or giving objects the same name.\par
\noindent{\footnotesize Izvorne zapisy / Original records: HIST-g030-009.}\par
\smallskip\noindent\emph{Sravnjenje s tekstom:} Relacije ogrančivajut komponenty i podpirajut normalizaciju simbolov; ne poistovětjajut vse komponenty.\par
\noindent\emph{Textual comparison:} The relations constrain components and support normalization of symbols; they do not identify all components with one another.\par
\noindent{\footnotesize TEXT-COMP-323; pasaže / passages: 3762, 3764, 3766, 3769, 3771; \texttt{CORPUS\_FORM\_COMPARISONS.json}.}\par
\end{minipage}\par
\par\medskip\Needspace{7\baselineskip}
\noindent\begin{minipage}{\linewidth}
\hypertarget{TERM-1227}{}
\noindent\textbf{neměrime funkcije}\quad{\footnotesize\texttt{TERM-1227}}\par
\noindent\emph{Iztočnik / Source:} nichtmeßbare Funktionen\par
\noindent\emph{Rusky / Russian:} неизмеримые функции\par
\noindent\emph{Ukrajinsky / Ukrainian:} невимірні функції\par
\smallskip\noindent\emph{Razlog:} Neměrime funkcije upotrěbjajut měrimost v smyslu teorije měry, ne kako opis veličin, ktore nikto ješče ne izměril. To razlikovanje jest cěly izričny razlog arhivnogo izbora. V §4 raboty 10 podčrtna primětka pripisuje Lebesgue rezultat o neměrimosti realnoj i imaginarnoj česti funkcije dvou realnyh prěmennyh. Avtor osobno upozornjaje, že iz samoj toj neměrimosti ješče ne slěduje ekstremna diskontinuirnost v kompleksnoj oblasti. Ta poslědnja vlastnost znači možnost približiti libovoljnu zadanu cěnnost slikami toček v vsakoj okolině. Zato neměrimy, diskontinuirny i ekstremno diskontinuirny ne sut tri zaměnljive etikety. Tut prověrena jest formulacija dostupnogo korpusa, ne nezavisny prěgled Lebesgueovoj citovanoj raboty.\par
\noindent\emph{Reason:} Neměrime funkcije uses measurability in the sense of measure theory, not quantities nobody has yet measured. This distinction is the entire explicit reason preserved for the choice. In §4 of Paper 10 a footnote attributes to Lebesgue the nonmeasurability of the real and imaginary parts as functions of two real variables. The author expressly warns that this alone does not imply extreme discontinuity in the complex domain. That latter property means being able to approximate any prescribed value by images of points in every neighborhood. Nonmeasurable, discontinuous, and extremely discontinuous are therefore not interchangeable labels. The evidence checked here is the wording of the available corpus, not an independent examination of Lebesgue's cited work.\par
\noindent{\footnotesize Izvorne zapisy / Original records: HIST-g072-009.}\par
\smallskip\noindent\emph{Sravnjenje s tekstom:} Avtorova definicija i podčrtna primětka odděljajut dvě svojstva, ktore ne možno spojiti v jednu etiketu.\par
\noindent\emph{Textual comparison:} The author's definition and footnote distinguish two properties that cannot be merged into a single label.\par
\noindent{\footnotesize TEXT-COMP-324; pasaže / passages: 8601, 8605, 8609, 8690; \texttt{CORPUS\_FORM\_COMPARISONS.json}.}\par
\end{minipage}\par
\par\medskip\Needspace{7\baselineskip}
\noindent\begin{minipage}{\linewidth}
\hypertarget{TERM-1228}{}
\noindent\textbf{nenegativne eksponenty}\quad{\footnotesize\texttt{TERM-1228}}\par
\noindent\emph{Iztočnik / Source:} nichtnegative Exponenten\par
\noindent\emph{Rusky / Russian:} неотрицательные показатели\par
\noindent\emph{Ukrajinsky / Ukrainian:} невід'ємні показники\par
\smallskip\noindent\emph{Razlog:} Nenegativne eksponenty čuva razliku medžu čislami većimi od nuly i čislami većimi ili ravnymi nuli. Arhiv prědpočita internacionalnu osnovu eksponent, dokolě ruska i ukrajinska paralela uživajut pokazatel odnosno pokaznyk; to jest kvalitativny izbor, ne čislovy rezultat sravnjanja. V polinomnoj relaciji (6) iz §2 raboty 6 eksponenty $\sigma_k$ sut nenegativne cěle čisla, i prinajmenje jeden mora byti nula. Zaměna s pozitivne zato ne prosto skrati stil: ona izključi slučaj, ktory tekst izrično zahtěvaje. Padežna forma nenegativnyh eksponentov v rečenju i nominativna slovarna forma označajut isti pojem. Tut ne postoji novo prědpoloženje o libovoljnyh realnyh eksponentah.\par
\noindent\emph{Reason:} Nenegativne eksponenty preserves the difference between numbers greater than zero and numbers greater than or equal to zero. The archive favors the international stem eksponent, whereas its Russian and Ukrainian comparanda use pokazatelʹ and pokaznyk; this is a qualitative choice, not a numerical comparison result. In polynomial relation (6) of Paper 6, §2, the exponents $\sigma_k$ are nonnegative integers and at least one must be zero. Substituting positive would therefore not merely shorten the style: it would exclude a case expressly required by the text. The inflected nenegativnyh eksponentov in the sentence and the dictionary's nominative form designate the same concept. No new assumption about arbitrary real exponents is involved.\par
\noindent{\footnotesize Izvorne zapisy / Original records: HIST-g042-014.}\par
\smallskip\noindent\emph{Sravnjenje s tekstom:} Relacija i slědujuče uslovje zahtěvajut prinajmenje jeden nulovy eksponent.\par
\noindent\emph{Textual comparison:} The relation and its following condition require at least one zero exponent.\par
\noindent{\footnotesize TEXT-COMP-325; pasaže / passages: 5196, 5197, 5198, 5202; \texttt{CORPUS\_FORM\_COMPARISONS.json}.}\par
\end{minipage}\par
\par\medskip\Needspace{7\baselineskip}
\noindent\begin{minipage}{\linewidth}
\hypertarget{TERM-1229}{}
\noindent\textbf{neprimarny}\quad{\footnotesize\texttt{TERM-1229}}\par
\noindent\emph{Iztočnik / Source:} nichtprimär\par
\noindent\emph{Rusky / Russian:} непримарный\par
\noindent\emph{Ukrajinsky / Ukrainian:} непримарний\par
\smallskip\noindent\emph{Razlog:} Neprimarny jest neposrědna negacija primarny. Arhivny razlog čuva vidlivu opoziciju v dokazě teoremy VI raboty 19, bez tvorjenja drugogo žargonnogo nazvanja. Primarny tu ne znači prosty: definicija govori o produktu $ab$ i o stepenah faktora. Ako ideal ne jest primarny, tekst izbira $ab$ v idealu, ale $a$ izvan njega i vse pozitivne stepeny $b$ izvan njega. Teorema svęzuje neprimarnost s razložimostju; ona ne ustanovjaje ekvivalenciju primarnosti i nerazložimosti. Podčrtna primětka davaje primarny, ale razložimy ideal $(x^2,xy,y^\lambda)$ pri $\lambda\ge2$. To jest važna ograda za prěvod negacije: slovarna blizkost ne směje obračati směr logičnogo zaključka.\par
\noindent\emph{Reason:} Neprimarny directly negates primarny. The recorded reason preserves the visible opposition in the proof of Theorem VI of Paper 19 without inventing a separate technical label. Primary does not mean prime here: the definition concerns a product $ab$ and powers of a factor. For a nonprimary ideal the text chooses $ab$ inside it, but $a$ outside it and every positive power of $b$ outside it. The theorem connects nonprimaryness with reducibility; it does not establish equivalence between primaryness and irreducibility. Its footnote gives the primary but reducible ideal $(x^2,xy,y^\lambda)$ for $\lambda\ge2$. This is an important safeguard in translating negation: lexical proximity must not reverse the direction of the implication.\par
\noindent{\footnotesize Izvorne zapisy / Original records: HIST-g081-044, HIST-g082-044.}\par
\smallskip\noindent\emph{Sravnjenje s tekstom:} Rabotě 19 prinaleži izričny protivpriměr k obratu; rabota 30 objasnjaje, pod kojim prědpoloženjem slaba i silna primarnost sovpadajut.\par
\noindent\emph{Textual comparison:} Paper 19 supplies an explicit counterexample to the converse; Paper 30 explains the hypothesis under which weak and strong primaryness coincide.\par
\noindent{\footnotesize TEXT-COMP-326; pasaže / passages: 12838, 12840, 12842, 12843, 12844, 16607, 16609, 16611, 16613, 16615, 16678, 16684, 16686; \texttt{CORPUS\_FORM\_COMPARISONS.json}.}\par
\end{minipage}\par
\par\medskip\Needspace{7\baselineskip}
\noindent\begin{minipage}{\linewidth}
\hypertarget{TERM-1230}{}
\noindent\textbf{neprimarny / primarny}\quad{\footnotesize\texttt{TERM-1230}}\par
\noindent\emph{Iztočnik / Source:} nichtprimär / primär\par
\noindent\emph{Rusky / Russian:} непервичный / первичный\par
\noindent\emph{Ukrajinsky / Ukrainian:} непервинний / первинний\par
\smallskip\noindent\emph{Razlog:} Par neprimarny / primarny prodolžaje terminologiju §5 raboty 30. Zapis izrično opravdava kratku formu primarny tym, že slaba i silna primarnost sovpadajut pri aksiomě I. To ne jest odstranjenje nepotrěbnyh pridavnikov povsudu: bez togo uslovja tekst davaje priměr, kde oni ne sovpadajut. Slaba forma govori o stepenah elementov-děliteljev nuly, silna o stepenah idealov-děliteljev nuly. V §6 uslovje cěpov děliteljev ostava prědpoloženjem teoremy o nerazložimyh idealah. Ruska i ukrajinska paralela v tom zapisě imajut drugu osnovu za primarnost než paralely k rabotě 19; kontinuitet medžuslovjanskogo termina zato jest svesna leksična odluka, ne bukvalno kopiranje vseh paralel.\par
\noindent\emph{Reason:} The pair neprimarny / primarny continues the terminology of §5 of Paper 30. The record expressly justifies the shorter primary by the coincidence of weak and strong primaryness under axiom I. This is not permission to remove the qualifiers everywhere: without that hypothesis the text gives a counterexample to their coincidence. The weak form concerns powers of zero-divisor elements, the strong form powers of zero-divisor ideals. In §6 the divisor-chain condition remains a hypothesis of the theorem about irreducible ideals. The Russian and Ukrainian comparanda in this record use a different stem for primaryness from those for Paper 19; continuity of the Interslavic term is thus a deliberate lexical decision, not literal copying of every parallel.\par
\noindent{\footnotesize Izvorne zapisy / Original records: HIST-g113-010.}\par
\smallskip\noindent\emph{Sravnjenje s tekstom:} Rabotě 19 prinaleži izričny protivpriměr k obratu; rabota 30 objasnjaje, pod kojim prědpoloženjem slaba i silna primarnost sovpadajut.\par
\noindent\emph{Textual comparison:} Paper 19 supplies an explicit counterexample to the converse; Paper 30 explains the hypothesis under which weak and strong primaryness coincide.\par
\noindent{\footnotesize TEXT-COMP-326; pasaže / passages: 12838, 12840, 12842, 12843, 12844, 16607, 16609, 16611, 16613, 16615, 16678, 16684, 16686; \texttt{CORPUS\_FORM\_COMPARISONS.json}.}\par
\end{minipage}\par
\par\medskip\Needspace{7\baselineskip}
\noindent\begin{minipage}{\linewidth}
\hypertarget{TERM-1231}{}
\noindent\textbf{neregularne sistemy}\quad{\footnotesize\texttt{TERM-1231}}\par
\noindent\emph{Iztočnik / Source:} nichtreguläre Systeme\par
\noindent\emph{Rusky / Russian:} нерегулярные системы\par
\noindent\emph{Ukrajinsky / Ukrainian:} нерегулярні системи\par
\smallskip\noindent\emph{Razlog:} Neregularne sistemy jest negacija definiranoj regularnosti, ne ocěnka, že sistemy sut neobyčne ili loše uredžene. Arhivny razlog jest kontrast potrěbny za Hilbertov protivpriměr i slědujuči paragraf raboty 6. Definicija VIII svęzuje regularnost s nenulovym rezultantom členov najvyššej dimenzije v ukazanoj oblasti. Teorema IX davaje integralnu bazu za regularny sistem; podčrtna primětka pokazuje, že obstojat neregularne sistemy bez takoj bazy. Ona ne tvrdi, že nijeden neregularny sistem ne ima bazu. Prikazany sistem $x_1,x_1x_2,x_1x_2^2,\ldots$ ne može byti generovany konečnym početkom toga rěda kako polinomnymi kombinacijami. Zato integralna baza tu dolžna byti čitana v avtorovom kontekstě, ne avtomatično kako baza cělyh čisel čislovogo polja.\par
\noindent\emph{Reason:} Neregularne sistemy negates a defined regularity condition; it does not describe systems as unusual or badly arranged. The recorded reason is the contrast needed for Hilbert's counterexample and the next section of Paper 6. Definition VIII ties regularity to a nonzero resultant of highest-degree terms in the specified domain. Theorem IX supplies an integral basis for a regular system; the footnote shows that some nonregular systems have no such basis. It does not assert that no nonregular system can have a basis. The displayed system $x_1,x_1x_2,x_1x_2^2,\ldots$ cannot be generated polynomially by a finite initial segment of this sequence. Integral basis must therefore be read in the author's context, not automatically as an integral basis of a number field.\par
\noindent{\footnotesize Izvorne zapisy / Original records: HIST-g052-007.}\par
\smallskip\noindent\emph{Sravnjenje s tekstom:} Definicija regularnosti, teorema i protivpriměr imajut različne logične funkcii. Protivpriměr dokazyvaje obstojanje, ne tvrđenje o vseh neregularnyh sistemah.\par
\noindent\emph{Textual comparison:} The definition of regularity, theorem, and counterexample have different logical functions. The counterexample proves existence, not a statement about all nonregular systems.\par
\noindent{\footnotesize TEXT-COMP-327; pasaže / passages: 6278, 6280, 6284, 6286, 6287, 6288, 6289, 6290, 6291, 6293, 6295, 6296, 6298, 6300, 6302, 6304; \texttt{CORPUS\_FORM\_COMPARISONS.json}.}\par
\end{minipage}\par
\par\medskip\Needspace{7\baselineskip}
\noindent\begin{minipage}{\linewidth}
\hypertarget{TERM-1232}{}
\noindent\textbf{nenulovost diskriminanta}\quad{\footnotesize\texttt{TERM-1232}}\par
\noindent\emph{Iztočnik / Source:} Nichtverschwinden der Diskriminante\par
\noindent\emph{Rusky / Russian:} ненулевость дискриминанта\par
\noindent\emph{Ukrajinsky / Ukrainian:} ненульовість дискримінанта\par
\smallskip\noindent\emph{Razlog:} Nenulovost diskriminanta zaměnjaje dolgu frazu diskriminant ne jest nula kratkym imennym nazvanjem svojstva. To jest zapisany stilističny razlog, podpirany ruskoju i ukrajinskoju paralelami, ne zahtěvanje, aby vsaka polna rečenija byla skračena. V §6 raboty 31 svojstvo jest uslovje teoremy o polnoj razložimosti prvogo roda. Diskriminant jest oprěděljen do ukazanogo nenulovogo faktora, a njegova nenulovost od izbora ne zavisi. Nenulovy ne znači pozitivny. Isto tako, ekvivalencija s jediničnym diskriminantnym idealom v danoj teoremě zavisi od koeficientnogo polja; nad libovoljnym prstenom nenulovy element ne mora byti jedinica. Kratkost termina ne směje izgubiti ty uslovja.\par
\noindent\emph{Reason:} Nenulovost diskriminanta replaces the longer phrase the discriminant is not zero by a compact nominal designation of the property. That is the recorded stylistic reason, supported by the Russian and Ukrainian comparanda, not a prohibition on the full sentence when it is clearer. In §6 of Paper 31 the property is the criterion in the theorem on complete reducibility of the first kind. The discriminant is determined up to the specified nonzero factor, while its nonvanishing is independent of that choice. Nonzero does not mean positive. Likewise, equivalence with the discriminant ideal being the unit ideal in that theorem depends on the coefficient field; over an arbitrary ring a nonzero element need not be a unit. The compact term must not erase these hypotheses.\par
\noindent{\footnotesize Izvorne zapisy / Original records: HIST-g148-002.}\par
\smallskip\noindent\emph{Sravnjenje s tekstom:} Slovesno nazvanje uslovja odgovarja točnoj teoremě s koeficientnym poljem i konečnym rangom.\par
\noindent\emph{Textual comparison:} The verbal name of the criterion corresponds to the precise theorem with a coefficient field and finite rank.\par
\noindent{\footnotesize TEXT-COMP-328; pasaže / passages: 18279, 18412; \texttt{CORPUS\_FORM\_COMPARISONS.json}.}\par
\end{minipage}\par
\par\medskip\Needspace{7\baselineskip}
\noindent\begin{minipage}{\linewidth}
\hypertarget{TERM-1233}{}
\noindent\textbf{nižše formy rěda form}\quad{\footnotesize\texttt{TERM-1233}}\par
\noindent\emph{Iztočnik / Source:} niedrigere Formen der Formenreihe\par
\noindent\emph{Rusky / Russian:} низшие формы ряда форм\par
\noindent\emph{Ukrajinsky / Ukrainian:} нижчі форми ряду форм\par
\smallskip\noindent\emph{Razlog:} Nižše formy rěda form jest blizky prěvod německoj frazy. Arhiv prědpočita obšteslovjanske osnovy nižši i ręd, sravnjene s ruskym i ukrajinskym vyraženjami, vměsto posebno nacionalnoj konstrukcije. Ne davaje ale čislove ocěnky razumlivosti ni konkretno imenovanoj odvrgnutoj alternativy. V §14 raboty 2 fraza ukazuje zaměnu $j$ ili $A$ nižšimi formami danogo rěda pri tvorjenju reducenta; ona odsyla do ranějše konstrukcije rěda, ne do položaja formuly niže na stranicě. Zato samo slovo nižši ne ustanovjaje, po kojemu parametru sut formy uporedžene. Dostupno reda i arhivno rěda ili ręd ne opravdavajut tvorenje različnih matematičnyh pojmov zaradi samoj ortografije.\par
\noindent\emph{Reason:} Nižše formy rěda form closely translates the German phrase. The archive favors the broadly Slavic stems nižši and ręd, compared with the Russian and Ukrainian expressions, over a specifically national construction. It supplies neither a numerical intelligibility score nor a named rejected alternative. In §14 of Paper 2 the phrase instructs replacement of $j$ or $A$ by lower forms of the relevant series in constructing a reducent; it refers to the earlier construction of that series, not to formulas lower on the page. Lower alone therefore does not determine the parameter by which the forms are ordered. Current reda and archival rěda or ręd do not justify creating different mathematical concepts merely from the spelling variation.\par
\noindent{\footnotesize Izvorne zapisy / Original records: HIST-g015-011.}\par
\smallskip\noindent\emph{Sravnjenje s tekstom:} Instrukcije ukazyvajut, ktoru formu trěba zaměniti, i odsylajut do konstrukcije rěda.\par
\noindent\emph{Textual comparison:} The instructions specify the form to be replaced and refer to the construction of the series.\par
\noindent{\footnotesize TEXT-COMP-329; pasaže / passages: 2204, 2206, 2207, 2210, 2212, 2214; \texttt{CORPUS\_FORM\_COMPARISONS.json}.}\par
\end{minipage}\par
\par\medskip\Needspace{7\baselineskip}
\noindent\begin{minipage}{\linewidth}
\hypertarget{TERM-1234}{}
\noindent\textbf{ni??a mnogovidnost}\quad{\footnotesize\texttt{TERM-1234}}\par
\noindent\emph{Iztočnik / Source:} niedrigere Mannigfaltigkeit\par
\noindent\emph{Rusky / Russian:} меньшая многообразность\par
\noindent\emph{Ukrajinsky / Ukrainian:} нижча многовидність\par
\smallskip\noindent\emph{Razlog:} Arhivno ni??a mnogovidnost jest ošteđeno; dostupny prěvod ima nižšej mnogovidnosti. Zapis objašnjaje izbor bukvalnoj osnovy mnogovidnost kako način sačuvanja historičnogo kontrasta medžu avtorovym prostym modulom i komutativnym upotrěbjenjem toga nazvanja. Podčrtna primětka raboty 17 govori o děliteljah nižšej mnogovidnosti i o izključenju, kogda mnogovidnost jest nula. Tu smysl svęzan jest s dimenzijeju v algebraičnom razvažanju; samo slovo ne uvodi gladku ili topologičnu mnogovidnost v modernom smyslu. Bukvalnost čuva vidlivu historičnu terminologiju, ale možna nejasnost ostava predmetom stručnogo razsmotrenja. Ona ne daje pravo poistovětiti avtorov prosty modul s vsakoj pozdnějšej definicijeju prostogo ideala.\par
\noindent\emph{Reason:} The archival ni??a mnogovidnost is damaged; the available translation reads nižšej mnogovidnosti. The record explains the literal stem mnogovidnost as a way to preserve the historical contrast between the author's prime module and commutative usage of that name. The footnote in Paper 17 concerns divisors of lower Mannigfaltigkeit and the exception when it is zero. The sense is dimensional within an algebraic argument; the word alone introduces neither a smooth nor a topological manifold in the modern sense. Literalness keeps the historical terminology visible, but possible ambiguity remains a subject for expert examination. It does not license identifying the author's prime module with every later definition of a prime ideal.\par
\noindent{\footnotesize Izvorne zapisy / Original records: HIST-g079-078.}\par
\smallskip\noindent\emph{Sravnjenje s tekstom:} Podčrtna primětka izrično različa dva upotrěbjenja prostogo modula; sam termin ne dovoljaje utočniti modernu geometričnu kategoriju.\par
\noindent\emph{Textual comparison:} The footnote explicitly distinguishes two uses of prime module; the term alone does not determine a modern geometric category.\par
\noindent{\footnotesize TEXT-COMP-330; pasaže / passages: 12233, 12235, 12237; \texttt{CORPUS\_FORM\_COMPARISONS.json}.}\par
\end{minipage}\par
\par\medskip\Needspace{7\baselineskip}
\noindent\begin{minipage}{\linewidth}
\hypertarget{TERM-1235}{}
\noindent\textbf{nižšego stepena}\quad{\footnotesize\texttt{TERM-1235}}\par
\noindent\emph{Iztočnik / Source:} niedrigeren Grades\par
\noindent\emph{Rusky / Russian:} более низкой степени\par
\noindent\emph{Ukrajinsky / Ukrainian:} нижчого степеня\par
\smallskip\noindent\emph{Razlog:} Nižšego stepena jest izbor za niedrigeren Grades, ne za nizši red v smyslu naslova paragrafa. Ta razlika medžu stepenom i redom jest izričny arhivny razlog, s ruskym stepenj i ukrajinskym stepinj kako paralelami. V §23 raboty 2 avtorova podčrtna primětka objašnjaje znak $\equiv$: produkty form nižšego stepena sut izpuščene. Zato pokazane formule ne trěba čitati kako obyčne točne ravnosti bez tih produktov. Tut terminologija i notacija služat jednej informaciji: redukcijny račun čuva odgovarjajuču ekvivalenciju, a ne tvrdi, že zanemarjene produkty sut nula. Prěvoditelj dolžen sačuvati i frazu i avtorovo objasnjenje znaka; iz samogo $\equiv$ ne trěba domysljati čislovy modul.\par
\noindent\emph{Reason:} Nižšego stepena is chosen for niedrigeren Grades, not lower order in the sense used by the section heading. This distinction between degree and order is the express archival reason, with Russian stepenʹ and Ukrainian stepinʹ as comparanda. In §23 of Paper 2 the author's footnote explains $\equiv$: products of forms of lower degree have been omitted. The displayed formulas must therefore not be read as ordinary exact equalities with those products simply absent. Terminology and notation carry the same information: the reduction calculation preserves the relevant equivalence, not a claim that the omitted products are zero. Both the phrase and the author's explanation of the sign must remain; the symbol $\equiv$ alone does not imply an integer modulus.\par
\noindent{\footnotesize Izvorne zapisy / Original records: HIST-g024-009.}\par
\smallskip\noindent\emph{Sravnjenje s tekstom:} Znak $\equiv$ imaje tut izrično avtorovo objasnjenje o izpuščenyh produktah nižšego stepena.\par
\noindent\emph{Textual comparison:} Here $\equiv$ has an explicit authorial explanation concerning omitted products of lower degree.\par
\noindent{\footnotesize TEXT-COMP-331; pasaže / passages: 3094, 3097, 3098, 3099, 3100; \texttt{CORPUS\_FORM\_COMPARISONS.json}.}\par
\end{minipage}\par
\par\medskip\Needspace{7\baselineskip}
\noindent\begin{minipage}{\linewidth}
\hypertarget{TERM-1236}{}
\noindent\textbf{nilpotentny ideal}\quad{\footnotesize\texttt{TERM-1236}}\par
\noindent\emph{Iztočnik / Source:} nilpotentes Ideal\par
\noindent\emph{Rusky / Russian:} нильпотентный идеал\par
\noindent\emph{Ukrajinsky / Ukrainian:} нільпотентний ідеал\par
\smallskip\noindent\emph{Razlog:} Nilpotentny ideal čuva internacionalny algebraičny termin. Oba arhivne razloga prědpočitajut tu osnovu, s ruskym i ukrajinskym paralelami, prěd tvorjenjem nezasvědčenogo slovjanskogo neologizma; nijeden takovy neologizm ne jest imenovany kako stvarno razsmotrěna alternativa. Definicija §12 raboty 34 zahtěvaje jednu pozitivnu stepen cělogo ideala, $\mathfrak c^\rho=0$. To jest silnějše uslovje než samo zahtěvanje, aby vsaky element iměl někoj stepen ravny nuli. Definicija vključaje nulovy ideal. Zato tvrđenje §25, že diskriminant jest nula pri obstojanju nilpotentnogo ideala, potrěbuje kontekstnu ogradu: pokazany dokaz davaje nulovy red v matricě samo pri nenulovom idealu. Dostupna rečenija ne piše nenulovy. Ostava otvorjeno, jest li to izpuščeno uslovje v prěvodu ili historična konvencija iztočnika; tu ne jest učinjena tiha ispravka.\par
\noindent\emph{Reason:} Nilpotentny ideal retains an international algebraic term. Both preserved reasons favor this stem, with Russian and Ukrainian parallels, over an unattested Slavic neologism; neither names such a neologism as an actually considered alternative. The definition in §12 of Paper 34 requires one positive power of the whole ideal to vanish, $\mathfrak c^\rho=0$. This is stronger than merely requiring some power of each individual element to vanish. The definition includes the zero ideal. Consequently, the §25 assertion that the discriminant vanishes when a nilpotent ideal exists needs a contextual qualification: the displayed proof supplies a zero matrix row only when the ideal is nonzero. The available sentence does not say nonzero. Whether this is an omitted translation hypothesis or a historical source convention remains unresolved; no silent emendation is made here.\par
\noindent{\footnotesize Izvorne zapisy / Original records: HIST-g190-001, HIST-g203-007.}\par
\smallskip\noindent\emph{Sravnjenje s tekstom:} Definicija uključaje nulovy ideal. Dokaz tvrđenja o diskriminantu potrěbuje nenulovy ideal; dostupna formulacija ostava iztočnym pytanjem.\par
\noindent\emph{Textual comparison:} The definition includes the zero ideal. The discriminant argument needs a nonzero ideal; the available wording remains a source question.\par
\noindent{\footnotesize TEXT-COMP-332; pasaže / passages: 20833, 20837, 20839, 20845, 20847, 20860, 22331, 22333, 22335, 22337, 22340, 22341, 22345, 22346, 22349; \texttt{CORPUS\_FORM\_COMPARISONS.json}.}\par
\end{minipage}\par
\par\medskip\Needspace{7\baselineskip}
\noindent\begin{minipage}{\linewidth}
\hypertarget{TERM-1237}{}
\noindent\textbf{Noetherov fundamentalny teorem algebraičnyh funkcij}\quad{\footnotesize\texttt{TERM-1237}}\par
\noindent\emph{Iztočnik / Source:} Noetherscher Fundamentalsatz der algebraischen Funktionen\par
\noindent\emph{Rusky / Russian:} основная теорема Нётер об алгебраических функциях\par
\noindent\emph{Ukrajinsky / Ukrainian:} Нетерова основна теорема алгебраїчних функцій\par
\smallskip\noindent\emph{Razlog:} Noetherov fundamentalny teorem algebraičnyh funkcij jest historično ime teoremy, ne obče pohvala njezinoj važnosti. Zapis prědpočita fundamentalny teorem za medžuslovjansky naslov; ruska i ukrajinska paralela upotrěbjajut prirodnu konstrukciju s osnovnoj teoremoju. V dostupnom naslově stoji Neterovomu, dokolě bibliografičny citat čuva Max Noether. Tekst i citat jasno pokazyvajut, že imenovany teorem jest Maksov, ne Emmyjin teorem o simetrijah. Rabota jest podpisana H. Kapfererom i ima dodatok sovměstny s E. Noether. Prinaležnost polinoma idealu, generovanomu dvoma polinomami, svęzuje se tu s konečnymi lokalnymi uslovjami. Takovo kratko objasnjenje pomaga identifikovati imenovany rezultat bez zaměny avtorstva i bez tvrdženja, že sam naslov davaje polny algoritm.\par
\noindent\emph{Reason:} Noetherov fundamentalny teorem algebraičnyh funkcij is the historical name of a theorem, not general praise for its importance. The record favors fundamentalny teorem in the Interslavic title, while the Russian and Ukrainian comparanda use natural constructions with fundamental or basic theorem. The available title has Neterovomu, whereas the bibliographic citation preserves Max Noether. Text and citation make clear that the named theorem is Max's, not Emmy's theorem concerning symmetries. The paper is signed by H. Kapferer and includes a supplement jointly with E. Noether. Membership of a polynomial in an ideal generated by two polynomials is related here to finite local conditions. This short identification distinguishes the named result and its authorship without claiming that the title itself supplies a complete algorithm.\par
\noindent{\footnotesize Izvorne zapisy / Original records: HIST-g214-002.}\par
\smallskip\noindent\emph{Sravnjenje s tekstom:} Naslov, podpis, historično objasnjenje i bibliografičny citat vměstě identifikujut teorem i razlikujut avtory.\par
\noindent\emph{Textual comparison:} The title, byline, historical explanation, and bibliographic citation together identify the theorem and distinguish the authors.\par
\noindent{\footnotesize TEXT-COMP-333; pasaže / passages: 45-isv.tex:34, 45-isv.tex:35, 45-isv.tex:36, 45-isv.tex:37, 45-isv.tex:38, 45-isv.tex:39, 45-isv.tex:63, 45-isv.tex:64, 45-isv.tex:65, 45-isv.tex:66, 45-isv.tex:68, 45-isv.tex:69, 45-isv.tex:70, 45-isv.tex:72, 45-isv.tex:73, 45-isv.tex:76, 45-isv.tex:77, 45-isv.tex:78, 45-isv.tex:79, 45-isv.tex:80, 45-isv.tex:81, 45-isv.tex:83, 45-isv.tex:85, 45-isv.tex:134, 45-isv.tex:135; \texttt{CORPUS\_FORM\_COMPARISONS.json}.}\par
\end{minipage}\par
\par\medskip\Needspace{7\baselineskip}
\noindent\begin{minipage}{\linewidth}
\hypertarget{TERM-1238}{}
\noindent\textbf{norma}\quad{\footnotesize\texttt{TERM-1238}}\par
\noindent\emph{Iztočnik / Source:} Norm\par
\noindent\emph{Rusky / Russian:} норма\par
\noindent\emph{Ukrajinsky / Ukrainian:} норма\par
\smallskip\noindent\emph{Razlog:} Norma jest v zapisě prijeta kako internacionalna osnova, s ruskym i ukrajinskym norma i s padežnoju formoju normoju v medžuslovjanskom tekstu. V rabotě 31, §4, avtor imenuje normoju koeficient svobodny od $t$ v determinantu $|tE-C|$. Dlja matrice reda $s$ to jest $(-1)^s\det C$, v tekstu zapisano $\pm|C|$. To znamenje ne trěba tiho odstraniti zaradi drugoj moderne konvencije normy. Vertikalne črty tu označajut determinant, ne absolutnu cěnnost; norma ne jest po definiciji nenegativna měra dliny. Sačuvane objašnjenja izbora dlja vseh trěh jezykov napisane sut po-anglijsky. Nynějšnje medžuslovjansko objasnjenje rekonstruuje te razloge; ne pripisuje arhivu tekst na jezyku, na ktorom on ne byl napisan.\par
\noindent\emph{Reason:} Norma is adopted in the record as an international stem, with Russian and Ukrainian norma and the inflected Interslavic normoju in the running text. In Paper 31, §4, the author calls the coefficient independent of $t$ in $|tE-C|$ the norm. For a matrix of order $s$ this is $(-1)^s\det C$, written $\pm|C|$ in the text. That sign must not be silently discarded to match another modern norm convention. The vertical bars denote a determinant, not an absolute value; the norm is not by definition a nonnegative measure of length. The preserved explanations of the choice for all three languages are written in English. The present Interslavic explanation reconstructs those reasons; it does not attribute to the archive a text in a language in which it was not written.\par
\noindent{\footnotesize Izvorne zapisy / Original records: HIST-g135-006.}\par
\smallskip\noindent\emph{Sravnjenje s tekstom:} Avtor oprědělja normu prěz ukazany harakterističny determinant. To fiksuje takože konvenciju znamenja.\par
\noindent\emph{Textual comparison:} The author defines the norm through the stated characteristic determinant. This also fixes the sign convention.\par
\noindent{\footnotesize TEXT-COMP-334; pasaže / passages: 18024, 18026, 18028; \texttt{CORPUS\_FORM\_COMPARISONS.json}.}\par
\end{minipage}\par
\par\medskip\Needspace{7\baselineskip}
\noindent\begin{minipage}{\linewidth}
\hypertarget{TERM-1239}{}
\noindent\textbf{norma / elementarny dělitelj}\quad{\footnotesize\texttt{TERM-1239}}\par
\noindent\emph{Iztočnik / Source:} Norm / Elementarteiler\par
\noindent\emph{Rusky / Russian:} норма / элементарный делитель\par
\noindent\emph{Ukrajinsky / Ukrainian:} норма / елементарний дільник\par
\smallskip\noindent\emph{Razlog:} Norma i elementarny dělitelj sut različne, povezane pojmy, ne sinonimy. Arhivny razlog za dělitelj jest semantična jasnost i kontinuitet s ranějšej terminologijeju korpusa; ruska i ukrajinska paralela imaju odgovarjajuče osnovy. V rabotě 25 norma ideala jest oprěděljena posrědstvom produktov elementarnyh děliteljev modulov, pri čem samo konečno mnogo faktorov razlikujut se od jedinicy. To ne jest avtomatično norma elementa iz drugyh rabot. Pri prostom idealu tekst različa normu od najvyššego elementarnogo dělitelja: nad nedoskonalym poljem norma može byti stepen odgovarjajučej proste funkcije. Zato zaměna normy jednym elementarnym děliteljem bez ukazanogo uslovja izgubila by informaciju o mnogokratnosti. Opis izboru ne pripisuje slovu norma jedno isto oprěděljenje prěz cěly korpus.\par
\noindent\emph{Reason:} Norma and elementarny dělitelj are distinct, related concepts, not synonyms. The archival reason for dělitelj is semantic transparency and continuity with earlier corpus terminology; the Russian and Ukrainian comparanda have corresponding stems. In Paper 25 the norm of an ideal is defined through products of elementary divisors of modules, with only finitely many factors different from the unit. It is not automatically the norm of an element used in other papers. For a prime ideal the text distinguishes its norm from its highest elementary divisor: over an imperfect field the norm may be a power of the associated irreducible polynomial. Replacing the norm by a single elementary divisor without the specified qualification would lose multiplicity information. The explanation of the lexical choice therefore does not impose one definition of norm throughout the corpus.\par
\noindent{\footnotesize Izvorne zapisy / Original records: HIST-g102-011.}\par
\smallskip\noindent\emph{Sravnjenje s tekstom:} Polje koeficientov, konečnost nejediničnyh faktorov i možno pojavjenje stepenov sut česti definicije i priměnjenja normy.\par
\noindent\emph{Textual comparison:} The coefficient field, finiteness of nonunit factors, and possible appearance of powers belong to the definition and use of the norm.\par
\noindent{\footnotesize TEXT-COMP-335; pasaže / passages: 15841, 15843, 15845, 15847; \texttt{CORPUS\_FORM\_COMPARISONS.json}.}\par
\end{minipage}\par
\par\medskip\Needspace{7\baselineskip}
\noindent\begin{minipage}{\linewidth}
\hypertarget{TERM-1240}{}
\noindent\textbf{norma / obrazovanje normy}\quad{\footnotesize\texttt{TERM-1240}}\par
\noindent\emph{Iztočnik / Source:} Norm / Normbildung\par
\noindent\emph{Rusky / Russian:} норма / построение нормы\par
\noindent\emph{Ukrajinsky / Ukrainian:} норма / побудова норми\par
\smallskip\noindent\emph{Razlog:} Obrazovanje normy označaje operaciju, norma jej rezultat. Arhivny zapis čuva norma v vseh trěh prěvodah i izbira obrazovanje dlja konstrukcijnogo kroka v rabotě 20. Razlika od normalizacije ili normovanja slěduje iz konkretnogo matematičnogo postupa; zapis ne davaje spisok odvrgnutyh sinonimov. V punktu 4 avtor tvori linearnu formu $\zeta(u)$ i množi ju s različnymi konjugatami, dobytymi permutacijami ukazanyh redov prěmennyh. Dobytaja simetrična funkcija vyražaje se prěz koeficienty produkta dvou form. Konstrukcija potom davaje polinomnu identičnost, upotrěbjanu pri specializaciji. Slovarny zapis dolžen čuvati ty koeficientne polja i način tvorjenja konjugatov: opis konkretnogo produkta ne jest neograničena definicija normy dlja vsakego algebraičnogo razširjenja. Ruske postrojenije i ukrajinske pobudova sut paralely dlja konstrukcijnoj funkcii frazy, ne novi sinonimy dlja samoj normy.\par
\noindent\emph{Reason:} Obrazovanje normy denotes the operation, norma its result. The archival record retains norma in all three translations and chooses obrazovanje for the construction step in Paper 20. The distinction from normalization follows from the actual mathematical operation; the record does not supply a list of rejected synonyms. In item 4 the author forms a linear expression $\zeta(u)$ and multiplies it by the distinct conjugates obtained by permuting the specified rows of variables. The resulting symmetric function is expressed through the coefficients of the product of two forms. The construction then supplies a polynomial identity used under specialization. The dictionary must preserve the coefficient fields and the specified formation of conjugates: this particular product is not an unrestricted definition of norm for every algebraic extension. Russian postroenie and Ukrainian pobudova are parallels for the constructive function of the phrase, not new synonyms for the norm itself.\par
\noindent{\footnotesize Izvorne zapisy / Original records: HIST-g083-010.}\par
\smallskip\noindent\emph{Sravnjenje s tekstom:} Konjugaty nastavajut konkretno zadanymi permutacijami; dobytaja identičnost prodolžaje se do specializovanyh cěnnostij.\par
\noindent\emph{Textual comparison:} The conjugates arise from specified permutations; the resulting identity continues to specialized values.\par
\noindent{\footnotesize TEXT-COMP-336; pasaže / passages: 13902, 13909, 13910, 13911, 13912, 13913, 13914, 13918, 13922, 13925, 13927, 13932, 13935, 13937, 13938, 13942; \texttt{CORPUS\_FORM\_COMPARISONS.json}.}\par
\end{minipage}\par
\par\medskip\Needspace{7\baselineskip}
\noindent\begin{minipage}{\linewidth}
\hypertarget{TERM-1241}{}
\noindent\textbf{norma osnovnogo ideala vzhodno k $m$}\quad{\footnotesize\texttt{TERM-1241}}\par
\noindent\emph{Iztočnik / Source:} Norm des Grundideals nach $m$\par
\noindent\emph{Rusky / Russian:} норма основного идеала относительно $m$\par
\noindent\emph{Ukrajinsky / Ukrainian:} норма основного ідеалу відносно $m$\par
\smallskip\noindent\emph{Razlog:} Norma osnovnogo ideala v rabotě 22, §5, jest norma osnovnogo modula po odnošenju k modulu danogo ideala, poslě prisojedinjenja ukazanyh prěmennyh k polju koeficientov. Arhivny razlog čuva tu vez medžu idealami i modulami: to ne jest nova, nezavisna definicija čislovoj normy. Definicija VI nazyvaje rezultat rezultantom odgovarjajučego nivoja; produkt rezultantov tvori rezultantnu formu. Ruske otnositeljno i ukrajinske vidnosno izražajut relaciju, ne operaciju děljenja. Arhivno vzhodno k jest ranějši izbor, označeny dlja ješče jednogo razgledanja; nynějšnji tekst ima po odnošenju k. Zato stary zapis ostaje dokazom istorije izbora, ne preporukom oživiti napuščenu konstrukciju. Simbol $m$ jest imenom ideala, ne měrnoju jedinicoju.\par
\noindent\emph{Reason:} The norm of the basic ideal in Paper 22, §5, is the norm of the basic module relative to the module of the given ideal, after adjoining the specified variables to the coefficient field. The recorded reason preserves this link between ideals and modules; it does not introduce an independent numerical norm. Definition VI calls the result the resultant at the corresponding level, and the product of these resultants forms the resultant form. Russian otnositelno and Ukrainian vidnosno express a relation, not division. The archived vzhodno k was an earlier choice explicitly awaiting further review; the current text uses po odnošenju k. Its preservation documents the earlier decision, not a recommendation to revive the superseded construction. The symbol $m$ names an ideal, not a unit of measurement.\par
\noindent{\footnotesize Izvorne zapisy / Original records: HIST-g089-003, HIST-g092-045.}\par
\smallskip\noindent\emph{Sravnjenje s tekstom:} Definicija VI povězuje rezultant $i$-go nivoja s normoju modulov poslě oprěděljenogo prisojedinjenja. Idealni i modulni jazyk zato opisujut tu samu konstrukciju v dvou oblikah. Čitanje po odnošenju k ohranjaje poredok argumentov; rezultant i najvyšši elementarny dělitelj ne sut zamenime.\par
\noindent\emph{Textual comparison:} Definition VI links the level-$i$ resultant to a module norm after a specified adjunction. Ideal and module language thus describe the same construction in two forms. Reading po odnošenju k preserves the order of the arguments; the resultant and greatest elementary divisor are not interchangeable.\par
\noindent{\footnotesize TEXT-COMP-337; pasaže / passages: 14830, 14832, 14834, 14839, 14841, 14842, 14848, 14850; \texttt{CORPUS\_FORM\_COMPARISONS.json}.}\par
\end{minipage}\par
\par\medskip\Needspace{7\baselineskip}
\noindent\begin{minipage}{\linewidth}
\hypertarget{TERM-1242}{}
\noindent\textbf{norma po tom modulu / po $Q$}\quad{\footnotesize\texttt{TERM-1242}}\par
\noindent\emph{Iztočnik / Source:} Norm nach diesem Modul / nach $Q$\par
\noindent\emph{Rusky / Russian:} норма по этому модулю / по $Q$\par
\noindent\emph{Ukrajinsky / Ukrainian:} норма за цим модулем / за $Q$\par
\smallskip\noindent\emph{Razlog:} Arhivny zapis izbira po, da by ostaviti vidimoj relaciju izraženu německym nach; on priznavaje, že sam predlog jest široky. V rabotě 18 norma najprvo odnoši se k danomu modulu, a potom, pri razloženju ideala na primarnu komponentu i dopolnjujuču čest, k komponentě $Q$. Ne ide o normu samogo $Q$ bez drugogo argumenta, ni o obyčny ostatok pri děljenju. Takovo po dolžno se čitati vměstě s definicijoju modulov i prisojedinennyh koeficientov. Ruske po i ukrajinske za pokazyvajut različne sposoby izraženja relacije, ale ne sut samostalnymi dokazami odobrenja širokoj konstrukcije v medžuslovjanskom. Pozdnějše po odnošenju k daje jasnějši opis te iste relacije; zaměna potrěbuje prověrku argumentov normy, ne samo zaměnu predloga.\par
\noindent\emph{Reason:} The archive chooses po to retain the relation expressed by German nach, while acknowledging that the preposition alone is broad. In Paper 18 the norm is first relative to the given module and subsequently, when an ideal is decomposed into a primary component and its complement, relative to the component $Q$. This is neither the norm of $Q$ alone nor an ordinary remainder on division. The preposition must be read with the definitions of the modules and the adjoined coefficients. Russian po and Ukrainian za illustrate different ways of expressing the relation, but do not independently establish the broad construction as accepted Interslavic usage. The later po odnošenju k makes the relation more explicit; replacing it requires checking the arguments of the norm, not merely replacing a preposition.\par
\noindent{\footnotesize Izvorne zapisy / Original records: HIST-g080-009.}\par
\smallskip\noindent\emph{Sravnjenje s tekstom:} Tekst najprvo davaje normu osnovnogo modula po danomu modulu, potom tvori primarny faktor iz dopolnjujučej česti ideala po $Q$. Sam predlog ne ustanovjaje, ktory objekt jest prvym ili drugim argumentom. Pojasnjenje dolžno imenovati te objekty i izběgati smysl obyčnogo ostatka pri děljenju.\par
\noindent\emph{Textual comparison:} The text first gives the norm of a basic module relative to the given module, then forms a primary factor from the complementary part of the ideal relative to $Q$. The preposition alone does not determine the first and second arguments. The explanation must name them and avoid suggesting an ordinary remainder on division.\par
\noindent{\footnotesize TEXT-COMP-338; pasaže / passages: 12343, 12353, 12354, 12355, 12356, 12357, 12358, 12359, 12360, 12361, 12362, 12363, 12364; \texttt{CORPUS\_FORM\_COMPARISONS.json}.}\par
\end{minipage}\par
\par\medskip\Needspace{7\baselineskip}
\noindent\begin{minipage}{\linewidth}
\hypertarget{TERM-1243}{}
\noindent\textbf{norma $G$ vzhodno k $A$}\quad{\footnotesize\texttt{TERM-1243}}\par
\noindent\emph{Iztočnik / Source:} Norm von $G$ nach $A$\par
\noindent\emph{Rusky / Russian:} норма $G$ относительно $A$\par
\noindent\emph{Ukrajinsky / Ukrainian:} норма $G$ відносно $A$\par
\smallskip\noindent\emph{Razlog:} V definiciji III raboty 22 norma $G$ po odnošenju k $A$ jest determinant prěobraženja, ktoro vyražaje bazu $A$ prěz bazu $G$. Směr jest suščestvenny: obratiti argumenty znači obratiti matematičnu konstrukciju. Formula (9) izražaje normu kako produkt elementarnyh děliteljev. V cěločislnom slučaju ona izražaje čislo ostatkovyh klasov; v polinomnom slučaju s dimenzijeju prostora klasov sravnjaje se jej stepenj, ne sam polinom. Arhivny razlog čuva kontinuitet s vvodom i privrěmenno vzhodno k, izrično označeno dlja daljšego razgledanja. Nynějšnja formulacija po odnošenju k jest jasna relacijna konstrukcija; istoricky zapis ne davaje razloga vrnuti ranějši predlog. Oba simbola ostajut matematičnymi imenami.\par
\noindent\emph{Reason:} In Definition III of Paper 22 the norm of $G$ relative to $A$ is the determinant of the transformation expressing a basis of $A$ in a basis of $G$. The direction matters: reversing the arguments reverses the construction. Formula (9) expresses the norm as the product of the elementary divisors. In the integer case it gives the number of residue classes; in the polynomial case its degree, not the polynomial itself, is compared with the dimension of the space of classes. The recorded reason preserves continuity with the introduction and provisionally retains vzhodno k, expressly marked for further review. The current po odnošenju k is an explicit relational construction; the historical entry supplies no reason to restore its predecessor. Both symbols remain mathematical names.\par
\noindent{\footnotesize Izvorne zapisy / Original records: HIST-g086-003, HIST-g092-013.}\par
\smallskip\noindent\emph{Sravnjenje s tekstom:} Definicija vyražaje bazu $A$ prěz bazu $G$, ne obratno. Čislo ostatkovyh klasov v cěločislnom slučaju i dimenzija v polinomnom slučaju ne sut totožne formulacije: vo vtorom slučaju trěba vzeti stepenj polinomnoj normy. Relacijna fraza ne dolžna tu razliku skryti.\par
\noindent\emph{Textual comparison:} The definition expresses a basis of $A$ in a basis of $G$, not conversely. The integer-case class count and polynomial-case dimension are not identical formulations: the latter takes the degree of the polynomial norm. The relational phrase must not conceal this distinction.\par
\noindent{\footnotesize TEXT-COMP-339; pasaže / passages: 14404, 14408, 14414; \texttt{CORPUS\_FORM\_COMPARISONS.json}.}\par
\end{minipage}\par
\par\medskip\Needspace{7\baselineskip}
\noindent\begin{minipage}{\linewidth}
\hypertarget{TERM-1244}{}
\noindent\textbf{normalna baza}\quad{\footnotesize\texttt{TERM-1244}}\par
\noindent\emph{Iztočnik / Source:} Normalbasis\par
\noindent\emph{Rusky / Russian:} нормальный базис\par
\noindent\emph{Ukrajinsky / Ukrainian:} нормальний базис\par
\smallskip\noindent\emph{Razlog:} Normalna baza čuva poznaty medžunarodny korenj i kratku konstrukciju pridavnika s imennikom baza. Arhivny razlog jest prozračnost toj konstrukcije, ne čislovo sravnjenje kandidatov. V rabotě 37 potrěbno razlikovati bazu polja nad osnovnym poljem od cěloj bazy maksimalnogo poredka nad osnovnym poredkom: v obojih slučajah elementy sut konjugaty, ale uslovja byta ne sut jednake. V rabotě 40 ta ista fraza označaje drugu konstrukciju, normalnu bazu podmodula po zadanoj bazě polnogo modula. Zato sama fraza bez definicije ne ustanovjaje jediny pojem. Tvrdženje poslě teoremy 5 v nynějšnjem prěvodu raboty 37 o dělimosti stepena potrěbuje sravnjenje s pečatnym iztočnikom; ono ne dolžno stati občeju slovarnoju definicijeju ili pravilom dlja drugogo teksta.\par
\noindent\emph{Reason:} Normalna baza retains the familiar international stem and a short adjective–noun construction. The recorded reason is its transparency, not a numerical comparison of candidates. Paper 37 requires a distinction between a field basis over the ground field and an integral basis of the maximal order over the ground order: both consist of conjugates, but their existence conditions differ. In Paper 40 the same phrase denotes another construction, a normal basis of a submodule relative to a prescribed basis of the full module. The phrase alone therefore does not identify one uniform concept. The current translation's assertion following Theorem 5 of Paper 37 about divisibility of the degree requires comparison with the printed source; it must not become a general dictionary definition or a rule for another text.\par
\noindent{\footnotesize Izvorne zapisy / Original records: HIST-g207-001.}\par
\smallskip\noindent\emph{Sravnjenje s tekstom:} Poljna baza, cěla baza maksimalnogo poredka i normalna baza podmodula po zadanoj bazě imajut različne definicije. Prve dvě upotrěbjajut konjugaty, tretja ukazanu formu koeficientov. Rečenica poslě teoremy 5 tvrdi, že dělimost stepena prěz $p$ izključaje cělu normalnu bazu. Dělimost stepena ne jest isto kako vyšše razvětvjenje, zato takovo čitanje potrěbuje sravnjenje s originalom; ne izvodi se iz njega obča norma terminologije.\par
\noindent\emph{Textual comparison:} A field basis, an integral basis of the maximal order and a normal submodule basis relative to a prescribed basis have different definitions. The first two use conjugates, the third the displayed coefficient form. The sentence after Theorem 5 says divisibility of the degree by $p$ excludes an integral normal basis. Divisibility of the degree is not the same as higher ramification, so that reading requires comparison with the original; no general terminology rule is inferred from it.\par
\noindent{\footnotesize TEXT-COMP-340; pasaže / passages: 22785, 22787, 22828, 22831, 22833, 22849, 22851, 22853, 23606, 23607, 23609, 23611, 23613, 23615, 23617, 23618, 23620; \texttt{CORPUS\_FORM\_COMPARISONS.json}.}\par
\end{minipage}\par
\par\medskip\Needspace{7\baselineskip}
\noindent\begin{minipage}{\linewidth}
\hypertarget{TERM-1245}{}
\noindent\textbf{normalna dělitvena algebra}\quad{\footnotesize\texttt{TERM-1245}}\par
\noindent\emph{Iztočnik / Source:} normale Divisionsalgebra\par
\noindent\emph{Rusky / Russian:} нормальная алгебра с делением\par
\noindent\emph{Ukrajinsky / Ukrainian:} нормальна алгебра з діленням\par
\smallskip\noindent\emph{Razlog:} Arhivny razlog izbira dělitvena, da by svojstvo děljenja bylo vidimo v samom nazvanju, a ne skryto v specialnom nacionalnom terminu. Konkretne odvrgnute slova ne sut zapisane, zato jih ne možno dopolniti kako istoricky sravnjane kandidaty. Ruska algebra s děljenjem i ukrajinska algebra z dilennjam sut opisne paralely, ne nezavisna medžuslovjanska atestacija. V rabotě 38 normalna takože ima točny smysl: centr algebry jest osnovno polje. Dělitvena označaje obratimost vsakego nenulovogo elementa; ne slěduje, že algebra jest komutativna. Glavny teorem govori o cikličnosti nad algebraičnym čislovym poljem, ne nad libovoljnym poljem. Pri daljšem upotrěbjenju trěba čuvati oba uslovja, centralnost i děljenje.\par
\noindent\emph{Reason:} The recorded reason chooses dělitvena to make the division property visible in the name rather than conceal it in a specialized national term. No particular rejected words are recorded, so they cannot be supplied as historical candidates. Russian algebra s deleniem and Ukrainian algebra z dilenniam provide descriptive parallels, not independent Interslavic attestation. In Paper 38 normalna also has a precise meaning: the centre is the ground field. Dělitvena means that every nonzero element is invertible; it does not imply commutativity. The main cyclicity theorem concerns algebraic number fields, not arbitrary ground fields. Subsequent use must preserve both requirements, centrality and division.\par
\noindent{\footnotesize Izvorne zapisy / Original records: HIST-g208-002.}\par
\smallskip\noindent\emph{Sravnjenje s tekstom:} Avtor ustanovjaje, že algebra nad osnovnym poljem ima centr ravny tomu polju, i upotrěbjaje polnu matričnu algebru kako razščepjeny slučaj. To davaje konkretny smysl slovu normalna. Dělitvena jest užši uslovje než prosta; te dva pridavniki ne smut se izgubiti pri obobčenju hesla.\par
\noindent\emph{Textual comparison:} The author stipulates that an algebra over the ground field has precisely that field as its centre and uses a full matrix algebra as the split case. This fixes the sense of normalna. Division is a stronger condition than simplicity; neither adjective may disappear when the entry is generalized.\par
\noindent{\footnotesize TEXT-COMP-341; pasaže / passages: 22999, 23013, 23081; \texttt{CORPUS\_FORM\_COMPARISONS.json}.}\par
\end{minipage}\par
\par\medskip\Needspace{7\baselineskip}
\noindent\begin{minipage}{\linewidth}
\hypertarget{TERM-1246}{}
\noindent\textbf{normalna prosta algebra}\quad{\footnotesize\texttt{TERM-1246}}\par
\noindent\emph{Iztočnik / Source:} normale einfache Algebra\par
\noindent\emph{Rusky / Russian:} нормальная простая алгебра\par
\noindent\emph{Ukrajinsky / Ukrainian:} нормальна проста алгебра\par
\smallskip\noindent\emph{Razlog:} Razlog dlja normalna prosta algebra v arhivu jest kratok: neposrědny matematičny termin. Obširno sravnjenje ili čislove vagi ne sut zachovane. Prosta ne znači lahka: ide o algebru bez netrivialnyh dvustrannyh idealov. V rabotě 38 normalna dopolnjaje uslovje, že centr jest osnovno polje. To jest smysl, ktory nyně obyčno izraža nazvanje centralna prosta algebra; ta pojasnjujuča paralela ne dokazuje, že byla istorickym kandidatom prěvoda. Prosta algebra ne musi byti dělitvena: polna matrična algebra jest suščestvenny priměr v samom argumentu. Zato normalna prosta i normalna dělitvena ne smut se sjediniti v jedno heslo bez različenja. Ruske i ukrajinske prostaja/prosta podpirajut disciplinarny smysl pridavnika, ne jeho obyčno značenje.\par
\noindent\emph{Reason:} The archival reason for normalna prosta algebra is brief: a direct mathematical term. No extended comparison or numerical weights survive. Prosta does not mean easy: the algebra has no nontrivial two-sided ideals. In Paper 38 normalna adds the requirement that the centre be the ground field. This is the sense now usually expressed by central simple algebra; that explanatory parallel is not evidence that it was a historical translation candidate. A simple algebra need not be a division algebra: full matrix algebras are essential to the argument itself. Normalna prosta and normalna dělitvena must therefore not be merged without distinction. Russian prostaya and Ukrainian prosta support the disciplinary sense of the adjective, not its everyday meaning.\par
\noindent{\footnotesize Izvorne zapisy / Original records: HIST-g208-005.}\par
\smallskip\noindent\emph{Sravnjenje s tekstom:} Avtor ustanovjaje, že algebra nad osnovnym poljem ima centr ravny tomu polju, i upotrěbjaje polnu matričnu algebru kako razščepjeny slučaj. To davaje konkretny smysl slovu normalna. Dělitvena jest užši uslovje než prosta; te dva pridavniki ne smut se izgubiti pri obobčenju hesla.\par
\noindent\emph{Textual comparison:} The author stipulates that an algebra over the ground field has precisely that field as its centre and uses a full matrix algebra as the split case. This fixes the sense of normalna. Division is a stronger condition than simplicity; neither adjective may disappear when the entry is generalized.\par
\noindent{\footnotesize TEXT-COMP-341; pasaže / passages: 22999, 23013, 23081; \texttt{CORPUS\_FORM\_COMPARISONS.json}.}\par
\end{minipage}\par
\par\medskip\Needspace{7\baselineskip}
\noindent\begin{minipage}{\linewidth}
\hypertarget{TERM-1247}{}
\noindent\textbf{normalna forma}\quad{\footnotesize\texttt{TERM-1247}}\par
\noindent\emph{Iztočnik / Source:} Normalform\par
\noindent\emph{Rusky / Russian:} нормальная форма\par
\noindent\emph{Ukrajinsky / Ukrainian:} нормальна форма\par
\smallskip\noindent\emph{Razlog:} Četyri zachovane razlogy dlja normalna forma naglašajut prozračnost kognatov i čuvanje kratkogo termina v formulah. Jednako isto nazvanje ne jest dovoljnym razlogom dlja jednake definicije v vseh rabotah. V rabotah 2 i 4 normalnost jest vezana s izčeznenjem ukazanyh invariantnyh diferencialnyh tvorjenij; v rabotě 4 osnovna forma jest izrično izključena. V rabotě 16 ide takože o izbor prědstavitelja po idealu relacij medžu koordinatami, pri čem normalna forma zadovoljava dopolniteljne uslovja na koeficienty. Ne slěduje obča jedinstvenost iz samogo slova normalna. Dlja logičnogo teksta možno zadržati nazvanje, ale konjunktivna, disjunktivna ili druga normalna forma potrěbuje vlastnu definiciju i dokaz točno tvrdžene ekvivalencije.\par
\noindent\emph{Reason:} Four preserved reasons for normalna forma emphasize transparent cognates and retaining a compact term in formulas. Identical wording nevertheless does not justify one definition across all the papers. In Papers 2 and 4 normality is tied to the vanishing of specified invariant differential formations; Paper 4 expressly excepts the ground form. Paper 16 also concerns selection of a representative modulo the ideal of relations between coordinates, with additional conditions on the coefficients. The adjective normalna alone does not imply a general uniqueness theorem. A logic text may retain the name, but conjunctive, disjunctive or other normal forms require their own definitions and proofs of the precise equivalence asserted.\par
\noindent{\footnotesize Izvorne zapisy / Original records: HIST-g003-003, HIST-g029-013, HIST-g036-002, HIST-g078-006.}\par
\smallskip\noindent\emph{Sravnjenje s tekstom:} Raboty 2, 3 i 4 vezujut normalne formy s invariantnymi diferencialnymi uslovjami. Rabota 16 uključaje relacije koordinatov i kongruenciju prědstaviteljev po idealu. Sličnost imen ne znači jedinstvenost procedury ili identičnost s normalnymi formami logiky. Prěnos termina potrěbuje novu definiciju, ne prěnos celogo istorickogo smysla.\par
\noindent\emph{Textual comparison:} Papers 2, 3 and 4 connect normal forms to invariant differential conditions. Paper 16 includes coordinate relations and congruence of representatives modulo an ideal. Similar names imply neither a unique procedure nor identity with logical normal forms. Transferring the term requires a new definition, not transplantation of the entire historical sense.\par
\noindent{\footnotesize TEXT-COMP-342; pasaže / passages: 267, 272, 274, 3675, 3707, 4456, 4457, 4458, 4459, 4460, 4461, 4462, 4463, 4464, 10729, 10735, 10736, 10737, 10738, 10739, 10740, 10741, 10742, 10743, 10744, 10745, 10746, 10747; \texttt{CORPUS\_FORM\_COMPARISONS.json}.}\par
\end{minipage}\par
\par\medskip\Needspace{7\baselineskip}
\noindent\begin{minipage}{\linewidth}
\hypertarget{TERM-1248}{}
\noindent\textbf{normalna forma $\rho$-tej variacije}\quad{\footnotesize\texttt{TERM-1248}}\par
\noindent\emph{Iztočnik / Source:} Normalform der $\rho$-ten Variation\par
\noindent\emph{Rusky / Russian:} нормальная форма $\rho$-й вариации\par
\noindent\emph{Ukrajinsky / Ukrainian:} нормальна форма $\rho$-ї варіації\par
\smallskip\noindent\emph{Razlog:} Normalna forma $\rho$-tej variacije jest specialno heslo, ne prosta povtorka normalna forma. Arhiv čuva grčske $\rho$ kako indeks reda variacije, ne kako bukvu, ktoru trěba prěvesti. V rabotě 12 linearna kombinacija invariantov odstranjaje směšane diferencialy reda $\rho+1$; ne možno zato zaměniti indeks v nazvanju s $\rho+1$. Koeficienty, znaky i faktorijaly v opreděljajučoj formulě ostajut čestju matematičnogo sadržanja. Poslědovna redukcija do prvih diferencialov jest daljši krok, ne už sama definicija toj normalnoj formy. Ruska i ukrajinska fraza imajut isti indeksny vzor. Razlog naziva termin neposrědnym; ne sodrživa dokaz jedinstvenosti ni čislovo ocěnjenje alternativ.\par
\noindent\emph{Reason:} Normalna forma of the $\rho$th variation is a specific entry, not a repetition of normalna forma. The archive preserves Greek $\rho$ as the variation-order index, not a letter to be translated. In Paper 12 a linear combination of invariants eliminates mixed differentials of order $\rho+1$; this does not permit replacing the index in the name by $\rho+1$. Coefficients, signs and factorials in the defining formula remain mathematical content. The subsequent reduction to first differentials is a further step, not already the definition of this normal form. The Russian and Ukrainian phrases have the same index pattern. The reason calls the terminology direct; it contains neither a uniqueness proof nor numerical assessments of alternatives.\par
\noindent{\footnotesize Izvorne zapisy / Original records: HIST-g074-010.}\par
\smallskip\noindent\emph{Sravnjenje s tekstom:} Indeks $\rho$ označaje variaciju, a $\rho+1$ red směšanyh diferencialov, ktore odstranjaje konstrukcija. Faktorijaly i alternujuče znaky ne sut tipografične dodatky. Slědujuči prěhod do prvih diferencialov jest oddělny krok, zato ne trěba vkladati jego rezultat v samu definiciju normalnoj formy.\par
\noindent\emph{Textual comparison:} The index $\rho$ denotes the variation, while $\rho+1$ is the order of the mixed differentials eliminated by the construction. Factorials and alternating signs are not typographical embellishments. The subsequent passage to first differentials is a separate step, whose result should not be read into the definition of the normal form.\par
\noindent{\footnotesize TEXT-COMP-343; pasaže / passages: 9017, 9018, 9019, 9020, 9021, 9022, 9023, 9024, 9025, 9031; \texttt{CORPUS\_FORM\_COMPARISONS.json}.}\par
\end{minipage}\par
\par\medskip\Needspace{7\baselineskip}
\noindent\begin{minipage}{\linewidth}
\hypertarget{TERM-1249}{}
\noindent\textbf{normalne formy}\quad{\footnotesize\texttt{TERM-1249}}\par
\noindent\emph{Iztočnik / Source:} Normalformen\par
\noindent\emph{Rusky / Russian:} нормальные формы\par
\noindent\emph{Ukrajinsky / Ukrainian:} нормальні форми\par
\smallskip\noindent\emph{Razlog:} Normalne formy jest množina od normalna forma, ne novo matematično svojstvo. Arhiv raboty 3 čuva samo obči razlog, že termin jest obyčny v vseh trěh prěvodah. Matematičny smysl davaje zaključečny odstavec raboty: formy izčezajut pri ukazanyh invariantnyh diferencialnyh operacijah, a njih upotrěbjenje podpiraje redukciju kontrakcij do kogredientnyh kontrakcij defekta 1. Mertensov rezultat v kvaternarnoj oblasti služi v tekstu kako izhodišče dlja prěnosa na $n$-arnu oblast. Citacija ne znači, že iztočnik Mertensa byl nanovo pročitan pri rekonstrukciji. Termin dolžen ostati povezan s tym kriterijem, ne s prosto udobnym zapisom ili s logičnymi normalnymi formami. Kratky istoricky razlog ne trěba predstavjati kako obširny dokumentovany izbor.\par
\noindent\emph{Reason:} Normalne formy is the plural of normalna forma, not a new mathematical property. The Paper 3 archive preserves only the general reason that the term is customary in all three translations. The mathematical sense comes from the paper's closing paragraph: the forms vanish under the specified invariant differential operations, and their use supports reduction to cogredient contractions of defect 1. The text uses Mertens's quaternary result as the starting point for an extension to the $n$-ary case. The citation does not mean that Mertens's original was newly read during reconstruction. The term must remain tied to that criterion, not merely a convenient notation or logical normal forms. A brief historical reason must not be presented as an extensively documented selection.\par
\noindent{\footnotesize Izvorne zapisy / Original records: HIST-g028-016.}\par
\smallskip\noindent\emph{Sravnjenje s tekstom:} Raboty 2, 3 i 4 vezujut normalne formy s invariantnymi diferencialnymi uslovjami. Rabota 16 uključaje relacije koordinatov i kongruenciju prědstaviteljev po idealu. Sličnost imen ne znači jedinstvenost procedury ili identičnost s normalnymi formami logiky. Prěnos termina potrěbuje novu definiciju, ne prěnos celogo istorickogo smysla.\par
\noindent\emph{Textual comparison:} Papers 2, 3 and 4 connect normal forms to invariant differential conditions. Paper 16 includes coordinate relations and congruence of representatives modulo an ideal. Similar names imply neither a unique procedure nor identity with logical normal forms. Transferring the term requires a new definition, not transplantation of the entire historical sense.\par
\noindent{\footnotesize TEXT-COMP-342; pasaže / passages: 267, 272, 274, 3675, 3707, 4456, 4457, 4458, 4459, 4460, 4461, 4462, 4463, 4464, 10729, 10735, 10736, 10737, 10738, 10739, 10740, 10741, 10742, 10743, 10744, 10745, 10746, 10747; \texttt{CORPUS\_FORM\_COMPARISONS.json}.}\par
\end{minipage}\par
\par\medskip\Needspace{7\baselineskip}
\noindent\begin{minipage}{\linewidth}
\hypertarget{TERM-1250}{}
\noindent\textbf{normalny dělitelj}\quad{\footnotesize\texttt{TERM-1250}}\par
\noindent\emph{Iztočnik / Source:} Normalteiler\par
\noindent\emph{Rusky / Russian:} нормальный делитель\par
\noindent\emph{Ukrajinsky / Ukrainian:} нормальний дільник\par
\smallskip\noindent\emph{Razlog:} Normalny dělitelj jest naměrno istoricko nazvanje normalnoj podgrupy. Petj arhivnyh razlogov čuva blizost k německomu Normalteiler i kontinuitet medžu rabotami; oni ne utvrdžajut, že ide o dělitelj čisla. Rabota 34 razjasnjaje smysl prěz vnutrne avtomorfizmy: podgrupa jest invariantna pri konjugaciji. Pri operatornyh homomorfizmah potrěbna jest takože dopustimost po zadanoj oblasti operatorov. V kompozicijnom redu normalnost v poprědnom členu ne smije se zaměniti normalnostju v cěloj grupě. Zapis 1251 prěvodi to samo německo slovo modernizujučim nazvanjem; to jest razlika medžu dvoma prěvodnymi politikami, ne dva različne grupove pojmy. Slovar može pojasniti ekvivalenciju, ne prikryti razliku v istorickom izboru.\par
\noindent\emph{Reason:} Normalny dělitelj deliberately retains a historical name for a normal subgroup. Five archival reasons prioritize closeness to German Normalteiler and continuity across papers; none identifies an integer divisor. Paper 34 explains the sense through inner automorphisms: the subgroup is invariant under conjugation. For operator homomorphisms, admissibility under the prescribed operator domain is an additional requirement. In a composition series, normality in the preceding member must not be replaced by normality in the whole group. Entry 1251 translates the same German word with a modernizing name. This is a difference between two translation policies, not two group-theoretic concepts. The dictionary can explain their equivalence without concealing the historical difference in choice.\par
\noindent{\footnotesize Izvorne zapisy / Original records: HIST-g117-008, HIST-g179-003, HIST-g181-007, HIST-g182-005, HIST-g183-005.}\par
\smallskip\noindent\emph{Sravnjenje s tekstom:} Konjugacija i tvorjenje faktorgrupy ustanovjajut skupny matematičny smysl dvou nazvanij. Normalnost po vsim vnutrnym avtomorfizmam, karakterističnost po vsim avtomorfizmam i dopustimost po zadanym operatoram ne sut jedno uslovje. Mužsko podgrup v rabotě 41 i žensko podgrupa v rabotě 34 ostajut dokumentovanoju gramatičnoju različnostju, ne novymi grupovymi pojmami.\par
\noindent\emph{Textual comparison:} Conjugation and formation of the quotient establish the common mathematical sense of the two names. Normality under all inner automorphisms, characteristicity under all automorphisms and admissibility under specified operators are not one condition. Masculine podgrup in Paper 41 and feminine podgrupa in Paper 34 remain a documented grammatical difference, not new group-theoretic concepts.\par
\noindent{\footnotesize TEXT-COMP-344; pasaže / passages: 17190, 19811, 19813, 19815, 19893, 19895, 19965, 19967, 19969, 19971, 19973, 24253, 24255, 24257; \texttt{CORPUS\_FORM\_COMPARISONS.json}.}\par
\end{minipage}\par
\par\medskip\Needspace{7\baselineskip}
\noindent\begin{minipage}{\linewidth}
\hypertarget{TERM-1251}{}
\noindent\textbf{normalny podgrup}\quad{\footnotesize\texttt{TERM-1251}}\par
\noindent\emph{Iztočnik / Source:} Normalteiler\par
\noindent\emph{Rusky / Russian:} [Bukvy ne sut zachovane / Lettering not recoverable]\par
\noindent\emph{Ukrajinsky / Ukrainian:} [Bukvy ne sut zachovane / Lettering not recoverable]\par
\smallskip\noindent\emph{Razlog:} Arhivny razlog dlja normalny podgrup izrično modernizuje istoricko Normalteiler. To se razlikuje od razloga dlja normalny dělitelj, ktory čuva istoricko nazvanje. V rabotě 41 $K^*$ jest abelova normalna podgrupa razširjenej grupy, a faktorgrupa jest izomorfna zadanoj grupě; to ustanovjaje matematičny smysl. Nynějšnji tekst pri tom ima mužsku formu podgrup, dok rabota 34 upotrěbjaje žensku podgrupa. Gramatična razlika potrěbuje samostalno usoglasnjenje: podgrupa potrěbuje normalna, ne mehanično prěnesti normalny iz drugogo nazvanja. Ruske i ukrajinske bukvy v tom arhivnom zapisu ne sut zachovane, zato ih ne rekonstruujemo iz očekovanogo značenja. Modernizacija jest dokumentovana naměra; obča jezična pravilnost konkretnoj mužskoj formy ne jest tym dokazana.\par
\noindent\emph{Reason:} The archival reason for normalny podgrup explicitly modernizes historical Normalteiler. It differs from the reason for normalny dělitelj, which retains the historical designation. In Paper 41, $K^*$ is an abelian normal subgroup of the extension group and the quotient is isomorphic to the prescribed group, establishing the mathematical sense. The current text uses masculine podgrup, whereas Paper 34 uses feminine podgrupa. This grammatical difference needs separate harmonization: podgrupa requires normalna, not mechanical transfer of normalny from the other designation. The Russian and Ukrainian lettering in this archival entry is lost and is not reconstructed from the expected meaning. Modernization is a documented intention; that intention does not establish the general linguistic correctness of this particular masculine form.\par
\noindent{\footnotesize Izvorne zapisy / Original records: HIST-g211-011.}\par
\smallskip\noindent\emph{Sravnjenje s tekstom:} Konjugacija i tvorjenje faktorgrupy ustanovjajut skupny matematičny smysl dvou nazvanij. Normalnost po vsim vnutrnym avtomorfizmam, karakterističnost po vsim avtomorfizmam i dopustimost po zadanym operatoram ne sut jedno uslovje. Mužsko podgrup v rabotě 41 i žensko podgrupa v rabotě 34 ostajut dokumentovanoju gramatičnoju različnostju, ne novymi grupovymi pojmami.\par
\noindent\emph{Textual comparison:} Conjugation and formation of the quotient establish the common mathematical sense of the two names. Normality under all inner automorphisms, characteristicity under all automorphisms and admissibility under specified operators are not one condition. Masculine podgrup in Paper 41 and feminine podgrupa in Paper 34 remain a documented grammatical difference, not new group-theoretic concepts.\par
\noindent{\footnotesize TEXT-COMP-344; pasaže / passages: 17190, 19811, 19813, 19815, 19893, 19895, 19965, 19967, 19969, 19971, 19973, 24253, 24255, 24257; \texttt{CORPUS\_FORM\_COMPARISONS.json}.}\par
\end{minipage}\par
\par\medskip\Needspace{7\baselineskip}
\noindent\begin{minipage}{\linewidth}
\hypertarget{TERM-1252}{}
\noindent\textbf{normno-ostatkovy symbol}\quad{\footnotesize\texttt{TERM-1252}}\par
\noindent\emph{Iztočnik / Source:} Normenrestsymbol\par
\noindent\emph{Rusky / Russian:} [Bukvy ne sut zachovane / Lettering not recoverable]; символ норменного вычета\par
\noindent\emph{Ukrajinsky / Ukrainian:} [Bukvy ne sut zachovane / Lettering not recoverable]; символ норменного лишку\par
\smallskip\noindent\emph{Razlog:} Normno-ostatkovy symbol sveziva medžunarodny norm- so slovjanskim ostatok. Dva zachovane razlogy naglašajut tu konstrukciju i kontinuitet medžu rabotami 38 i 39. Rusky vyčet i ukrajinsky lyšok pokazyvajut različne korenje dlja ostatka; oni ne davajut čislovu vagoju automatično izbrati medžuslovjansko slovo. V cikličnom slučaju symbol po městu ima jediničnu vrědnost ravno pri odgovarjajučem uslovju normy. To ne jest obyčny ostatok pri děljenju ni razlika medžu normoju i ostatkom. Potrěbno čuvati argument, razširjenje i město v simboličnom zapisu. Něktore arhivne paralely sut nečitljive, druge sut zachovane; čitljivy zapis ne popravjaje potajno bukvy drugogo. Odkazy na Hasseho i Chevalleyja sut citacije v korpusu, ne novy dokaz konsultacije ih izvornyh tekstov.\par
\noindent\emph{Reason:} Normno-ostatkovy symbol joins international norm- to Slavic ostatok. Two preserved reasons emphasize that construction and continuity between Papers 38 and 39. Russian vychet and Ukrainian lyshok show different roots for a residue; they do not provide numerical weights that automatically select the Interslavic word. In the cyclic case the symbol at a place has identity value precisely under the corresponding norm condition. It is neither an ordinary remainder on division nor a difference between norm and residue. The argument, extension and place must remain explicit in the notation. Some archived comparanda are unreadable and others survive; a surviving record does not silently restore another's lost letters. References to Hasse and Chevalley are corpus citations, not fresh evidence that their originals were consulted.\par
\noindent{\footnotesize Izvorne zapisy / Original records: HIST-g208-011, HIST-g209-006.}\par
\smallskip\noindent\emph{Sravnjenje s tekstom:} Normno-ostatkovy symbol patrí k zadanoj cikličnoj situaciji i městu. Jegova jedinična vrědnost izražaje uslovje normy; text potom povězuje ty uslovja s razščepjenjem i produktnoju relacijeju. Termin ne trěba čitati kako prosty čislovy ostatok ili kako samostalnu definiciju normy.\par
\noindent\emph{Textual comparison:} The norm-residue symbol belongs to the specified cyclic setting and place. Its identity value expresses a norm condition, subsequently related to splitting and a product relation. The term should not be read as a simple numerical remainder or an independent definition of norm.\par
\noindent{\footnotesize TEXT-COMP-345; pasaže / passages: 23065, 23066, 23067, 23068, 23069, 23070, 23071, 23072, 23073, 23152, 23153, 23154, 23155, 23156, 23157, 23158, 23159, 23160, 23161, 23162, 23163, 23164, 23165, 23238; \texttt{CORPUS\_FORM\_COMPARISONS.json}.}\par
\end{minipage}\par
\par\medskip\Needspace{7\baselineskip}
\noindent\begin{minipage}{\linewidth}
\hypertarget{TERM-1253}{}
\noindent\textbf{normny teorem}\quad{\footnotesize\texttt{TERM-1253}}\par
\noindent\emph{Iztočnik / Source:} Normensatz\par
\noindent\emph{Rusky / Russian:} [Bukvy ne sut zachovane / Lettering not recoverable]; [Bukvy ne sut zachovane / Lettering not recoverable]\par
\noindent\emph{Ukrajinsky / Ukrainian:} [Bukvy ne sut zachovane / Lettering not recoverable]; [Bukvy ne sut zachovane / Lettering not recoverable]\par
\smallskip\noindent\emph{Razlog:} Normny teorem jest kompaktno nazvanje, izbrano radi kontinuiteta normnoj terminologije. Arhivny opis spominaje lokalnu teoriju polj klasov, ale konkretny teorem v rabotě 39 govori o prěhodu od vseh lokalnyh uslovij k globalnomu. V cikličnom slučaju element osnovnogo čislovogo polja jest norma iz razširjenja, ako jest norma pri vsakom konečnom i beskonečnom městu. Algebrična formulacija govori o razščepjenju algebry i ima svoj obseg. Ne možno prěnesti tvrdženje bez uslovij na libovoljno razširjenje; rabota 38 sama upozornjaje na granicu neabelovogo obobčenja. Ruske i ukrajinske bukvy v arhivnyh paralelah ne sut zachovane. Slovar zato opira objasnjenje na čitlivy zapis razloga i konkretne tvrdženja, ne na izmysljene obnovjenje paralel.\par
\noindent\emph{Reason:} Normny teorem is a compact name chosen for continuity in norm terminology. The archival description mentions local class field theory, but the particular theorem in Paper 39 passes from all local conditions to a global conclusion. In the cyclic case an element of the ground number field is a norm from the extension if it is a norm at every finite and infinite place. The algebraic formulation concerns splitting and has its own scope. The statement cannot be transferred without its hypotheses to an arbitrary extension; Paper 38 itself marks the limit of a nonabelian generalization. The Russian and Ukrainian lettering in the archived comparanda is lost. The dictionary therefore grounds its explanation in the readable reason and actual assertions, not invented restorations of those parallels.\par
\noindent{\footnotesize Izvorne zapisy / Original records: HIST-g209-003, HIST-g211-003.}\par
\smallskip\noindent\emph{Sravnjenje s tekstom:} V cikličnoj česti raboty 39 vsi konečne i beskonečne města učastvujut v uslovju normy. Algebrično razščepjenje ima širšu formulaciju, no ne dopuščaje neograničene zaměny cikličnogo razširjenja libovoljnym. Arhivny kratky opis lokalnoj teorije ne dolžen skryti globalny vyvod samogo teorema.\par
\noindent\emph{Textual comparison:} In the cyclic part of Paper 39 all finite and infinite places enter the norm condition. Algebraic splitting has a broader formulation but does not license replacing a cyclic extension by an arbitrary one without restriction. The archive's brief reference to local theory must not hide the theorem's global conclusion.\par
\noindent{\footnotesize TEXT-COMP-346; pasaže / passages: 23075, 23231, 23238, 23243, 23278, 23294, 23296, 23317, 24190; \texttt{CORPUS\_FORM\_COMPARISONS.json}.}\par
\end{minipage}\par
\par\medskip\Needspace{7\baselineskip}
\noindent\begin{minipage}{\linewidth}
\hypertarget{TERM-1254}{}
\noindent\textbf{potrěbne i dostatočne uslovja}\quad{\footnotesize\texttt{TERM-1254}}\par
\noindent\emph{Iztočnik / Source:} notwendige und hinreichende Bedingungen\par
\noindent\emph{Rusky / Russian:} необходимые и достаточные условия\par
\noindent\emph{Ukrajinsky / Ukrainian:} необхідні й достатні умови\par
\smallskip\noindent\emph{Razlog:} Potrěbne i dostatočne uslovja čuva oboje směry implikacije. Arhivny razlog upotrěbjaje prozračnu slovjansku paru, sravnimu s ruskim neobhodimye i dostatočnye i ukrajinskim neobhidni i dostatni. V rabotě 4, §3, uslovja karakterizujut pojavu cělogo reda kako nerazložimogo argumenta; ne ide samo o obstajanje jeho jednotlivyh komponentov. Primětka razlikuje ranějše dostatočno uslovje od nynějšnjej potrěbnosti i dostatočnosti. Spojka i zato ne jest zamenima s ili. Ako uslovja tvore sistem, oboje směry odnošet se k cělomu sistemu; ne slěduje, že vsako pojedino uslovje jest samo dostatočno. To jest važno takože dlja logičnogo prěvoda: kratka fraza dolžna čuvati obseg uslovij i tvrdženja, ne oslabiti potrěbnost do praktičnoj poleznosti.\par
\noindent\emph{Reason:} Potrěbne i dostatočne uslovja preserves both directions of implication. The recorded reason uses a transparent Slavic pair comparable to Russian neobkhodimye i dostatochnye and Ukrainian neobkhidni i dostatni. In Paper 4, §3, the conditions characterize occurrence of an entire row as an indecomposable argument, not merely occurrence of its separate components. A note distinguishes an earlier sufficient condition from the present necessity and sufficiency. The conjunction i is therefore not interchangeable with ili. Where the conditions form a system, both implications concern the whole system; it does not follow that each condition alone is sufficient. This also matters for translation of logic: the compact phrase must retain the scope of the conditions and assertion, not weaken necessity into practical usefulness.\par
\noindent{\footnotesize Izvorne zapisy / Original records: HIST-g032-008.}\par
\smallskip\noindent\emph{Sravnjenje s tekstom:} Primětka razlikuje dostatočno uslovje ranějšego rezultata od obostrannej charakterizacije. Diferencialne uslovja odnošet se k cělomu nerazložimomu redu. Logična sila spojky i i skupny obseg sistema dolžny ostati vidime; to ne jest tvrdženje, že vsako posamezno jednačenje samo daje cěly rezultat.\par
\noindent\emph{Textual comparison:} The note distinguishes a sufficient condition in an earlier result from a characterization in both directions. The differential conditions concern the complete indecomposable row. The force of and and the joint scope of the system must remain visible; this does not assert that each separate equation alone gives the whole result.\par
\noindent{\footnotesize TEXT-COMP-347; pasaže / passages: 3987, 3989, 3990, 3991, 3992, 3993, 3994, 3995, 3996, 3997, 3998, 3999, 4000, 4001, 4002, 4003, 4004; \texttt{CORPUS\_FORM\_COMPARISONS.json}.}\par
\end{minipage}\par
\par\medskip\Needspace{7\baselineskip}
\noindent\begin{minipage}{\linewidth}
\hypertarget{TERM-1255}{}
\noindent\textbf{nula / jedinica / vlastne oblasti cělosti}\quad{\footnotesize\texttt{TERM-1255}}\par
\noindent\emph{Iztočnik / Source:} Null / Einheit / eigentliche Integritätsbereiche\par
\noindent\emph{Rusky / Russian:} нуль / единица / собственные области целостности\par
\noindent\emph{Ukrajinsky / Ukrainian:} нуль / одиниця / власні області цілісності\par
\smallskip\noindent\emph{Razlog:} Nula, jedinica i vlastna oblast cělosti ne sut tri zamenime nazvanja. Arhivny razlog čuva razliku medžu osnovnym pojmom prstena i silnějšimi uslovjami. V rabotě 19, §1, zakony davajut nulu, ale ne potrěbujut jedinicu ni izključenje nulovyh děliteljev. Vlastne oblasti cělosti imajut jedinicu i svojstvo, že iz nule produkta slěduje nula jednogo faktora. Jedinica tu jest neutralny element množenja, ne libovoljny obratimy element. Vlastny v istorickom nazvanju ne označaje vlastnu podmnožinu někoj drugoj oblasti. Nynějšnja definicija inogo teksta može izrično dodati $0\ne1$; ne trěba tvrditi, že takovo dopolnjenje bylo izraženo v toj rečenici Noether. Čisla v zapisu $na$ tu označajut povtoreno složenje, ne už zadano vloženje cělyh čisel v prsten.\par
\noindent\emph{Reason:} Nula, jedinica and vlastna oblast cělosti are not interchangeable names. The recorded reason preserves the distinction between the underlying notion of ring and stronger requirements. In Paper 19, §1, the laws supply zero but require neither an identity nor exclusion of zero divisors. Vlastne oblasti cělosti have an identity and the property that a zero product forces one factor to be zero. Jedinica here means the multiplicative identity, not an arbitrary invertible element. Vlastny in the historical name does not mean a proper subset of another domain. Another text's modern definition may expressly add $0\ne1$; this must not be represented as an explicit requirement in this sentence of Noether's. Integers in $na$ here denote repeated addition, not an already specified embedding of the integers into the ring.\par
\noindent{\footnotesize Izvorne zapisy / Original records: HIST-g081-021, HIST-g082-021.}\par
\smallskip\noindent\emph{Sravnjenje s tekstom:} Zakony osnovnogo prstena ne vključajut jedinicu. Tekst dodavaje ju vměstě s uslovjem o nulovom produktu, kogda nazyvaje vlastne oblasti cělosti. Čislovy množitelj $n$ v $na$ označaje povtoreno složenje; on ne potvrdžaje, že osnovny prsten sodrživa kopiju prstena cělyh čisel.\par
\noindent\emph{Textual comparison:} The underlying ring laws do not include an identity. The text adds one together with the zero-product condition when naming eigentliche Integritätsbereiche. The numerical multiplier $n$ in $na$ denotes repeated addition; it does not establish that the ring contains a copy of the integers.\par
\noindent{\footnotesize TEXT-COMP-348; pasaže / passages: 12476, 12478, 12479, 12480, 12481, 12482, 12483, 12485; \texttt{CORPUS\_FORM\_COMPARISONS.json}.}\par
\end{minipage}\par
\par\medskip\Needspace{7\baselineskip}
\noindent\begin{minipage}{\linewidth}
\hypertarget{TERM-1256}{}
\noindent\textbf{nulovy ili jediničny ideal}\quad{\footnotesize\texttt{TERM-1256}}\par
\noindent\emph{Iztočnik / Source:} Null- bzw. Einheitsideal\par
\noindent\emph{Rusky / Russian:} нулевой или единичный идеал\par
\noindent\emph{Ukrajinsky / Ukrainian:} нульовий або одиничний ідеал\par
\smallskip\noindent\emph{Razlog:} Nulovy ili jediničny ideal izražaje dvě vozmožne vrědnosti ideala, ne dva imena jednogo objekta. Arhiv čuva tu disjunkciju i nazvanje jediničny ideal kako termin projekta; takovo vnutrne upotrěbjenje ne jest nezavisna jezična autorizacija. V rabotě 31, §6, diskriminantny ideal nad odgovarjajučim poljem jest nulovy ili cěle polje. Jediničny ideal jest ideal porodženy jedinicoju, ne jedinoelementna množina, sostoječa samo iz jedinice. Dokaz tvori produkt diskriminantnyh idealov komponentov: vsi jedinične faktory davajut jediničny produkt, a nula jednogo davaje nulu. Ta dichotomija ne prenosi se bez uslovij na idealy libovoljnogo prstena. Ruska i ukrajinska para čuva tu samu razliku medžu nulovym i jediničnym.\par
\noindent\emph{Reason:} Nulovy ili jediničny ideal expresses two possible values of an ideal, not two names for one object. The archive preserves this disjunction and treats jediničny ideal as the project's term; internal usage is not independent linguistic authorization. In Paper 31, §6, the discriminant ideal over the relevant field is zero or the whole field. The unit ideal is generated by the identity, not the singleton containing only the identity. The proof takes the product of the component discriminant ideals: all unit factors give the unit ideal, while a zero factor gives zero. This dichotomy does not extend without hypotheses to ideals of an arbitrary ring. The Russian and Ukrainian pairs retain the same distinction between zero and unit.\par
\noindent{\footnotesize Izvorne zapisy / Original records: HIST-g145-004.}\par
\smallskip\noindent\emph{Sravnjenje s tekstom:} Jediničny ideal jest cěly osnovny prsten, nulovy ima samo nulu. V razgledanom poljnom kontekstu diskriminantny ideal ima jednu iz tyh dvou vrědnostij; argument upotrěbjaje produkt komponentnyh diskriminantnyh idealov. Ili izražaje alternativu rezultata, ne nejasnost čitanja.\par
\noindent\emph{Textual comparison:} The unit ideal is the whole base ring, while the zero ideal contains only zero. In the field setting considered, the discriminant ideal has one of these two values; the argument uses a product of component discriminant ideals. Or expresses alternative outcomes, not uncertainty in reading.\par
\noindent{\footnotesize TEXT-COMP-349; pasaže / passages: 18350, 18352, 18354, 18441; \texttt{CORPUS\_FORM\_COMPARISONS.json}.}\par
\end{minipage}\par
\par\medskip\Needspace{7\baselineskip}
\noindent\begin{minipage}{\linewidth}
\hypertarget{TERM-1257}{}
\noindent\textbf{nulovy i jediničny ideal}\quad{\footnotesize\texttt{TERM-1257}}\par
\noindent\emph{Iztočnik / Source:} Null- und Einheitsideal\par
\noindent\emph{Rusky / Russian:} нулевой и единичный идеалы\par
\noindent\emph{Ukrajinsky / Ukrainian:} нульовий та одиничний ідеали\par
\smallskip\noindent\emph{Razlog:} Nulovy i jediničny ideal nabrajaje dva izključene slučaja. Za razliku od ili v heslu 1256, i tu sojedinjaje dva objekta, ktore treba sledovati oddělno pri korespondenciji idealov. V rabotě 31, §8, prěhod k prstenu kvocientov jest lokalizacija s oprěděljenymi imenovateljami, ne faktor-prsten po idealu. Tekst najprvo opisuje korespondenciju dlja ostalyh idealov i potom izrično priradžuje nulovy i jediničny ideal. Arhivny razlog spominaje razliku medžu skračenym opisom netrivialnyh idealov i polnoju frazoju pečatnogo iztočnika; to jest zachovano tvrdženje ranějšego zapisa, ne nova prověrka skana. Pri prěvodu ne trěba ni skryti te izključene slučaje pod kratkym pridavnikom, ni nenaměrno prěmeniti istoricku konvenciju o prostyh idealah.\par
\noindent\emph{Reason:} Nulovy i jediničny ideal lists two exceptional cases. Unlike ili in entry 1256, i joins two objects to be tracked separately in an ideal correspondence. In Paper 31, §8, passage to the ring of quotients is localization with specified denominators, not quotienting by an ideal. The text first describes the correspondence for other ideals and then expressly matches the zero and unit ideals. The recorded reason mentions a difference between an abbreviated reference to nontrivial ideals and the printed source's fuller phrase. This remains a claim in the earlier record, not a fresh inspection of the scan. Translation must neither hide the exceptions under a short adjective nor inadvertently change the historical convention concerning prime ideals.\par
\noindent{\footnotesize Izvorne zapisy / Original records: HIST-g159-006.}\par
\smallskip\noindent\emph{Sravnjenje s tekstom:} Prsten kvocientov tvori se dopuščanjem zadanyh imenovateljev. Korespondencija nenulovyh nejediničnyh idealov dopolnjaje se dvoma izričnymi priradženjami. Nulovy i jediničny ideal zato dolžny ostati različenymi objektami; nabrajanje ne znači, že vsaky ideal jest jednym iz njih.\par
\noindent\emph{Textual comparison:} The ring of quotients is formed by allowing specified denominators. The correspondence for nonzero nonunit ideals is supplemented by two explicit matches. Zero and unit must therefore remain distinct objects; listing them does not imply that every ideal is one of the two.\par
\noindent{\footnotesize TEXT-COMP-350; pasaže / passages: 18932; \texttt{CORPUS\_FORM\_COMPARISONS.json}.}\par
\end{minipage}\par
\par\medskip\Needspace{7\baselineskip}
\noindent\begin{minipage}{\linewidth}
\hypertarget{TERM-1258}{}
\noindent\textbf{nulovy ideal}\quad{\footnotesize\texttt{TERM-1258}}\par
\noindent\emph{Iztočnik / Source:} Nullideal\par
\noindent\emph{Rusky / Russian:} нулевой идеал\par
\noindent\emph{Ukrajinsky / Ukrainian:} нульовий ідеал\par
\smallskip\noindent\emph{Razlog:} Nulovy ideal jest ideal, sostoječi samo iz nuly. Arhivny razlog čuva to nazvanje v cěloj teoriji razloženja raboty 31; on ne predlagaje novy pojem. V §2 nulovy ideal jest vyraženy prěz produkt i prěsěčenje primarnyh idealov pod uslovjami razgledanogo prstena konečnogo ranga. Iz odgovarjajučego razloženja dobyvajut se direktne komponenty prstena i izomorfizmy s prstenami ostatkovyh klasov. To ne znači, že sam prsten jest nulovy: nula jest idealom v prsteně, ktory može iměti mnoho elementov. Takože ne možno izrěkti ravnost produkta i prěsěčenja dlja libovoljnyh idealov bez uslovij. Ruske nulevoj i ukrajinske nuljovyj davajut blizke paralely, ale osnovny razlog čuvanja termina jest jego jasno i stalno značenje v dokazu.\par
\noindent\emph{Reason:} Nulovy ideal is the ideal containing only zero. The recorded reason retains the name throughout the decomposition theory of Paper 31; it does not propose a new concept. In §2 the zero ideal is expressed as both a product and an intersection of primary ideals under the assumptions on the finite-rank ring. The corresponding decomposition yields direct components of the ring and isomorphisms with residue-class rings. This does not mean that the ring itself is zero: zero is an ideal in a ring that may have many elements. Nor may equality of product and intersection be asserted for arbitrary ideals without the relevant hypotheses. Russian nulevoy and Ukrainian nulovyi provide close parallels, but the principal reason for retaining the term is its clear and stable object reference in the proof.\par
\noindent{\footnotesize Izvorne zapisy / Original records: HIST-g122-002.}\par
\smallskip\noindent\emph{Sravnjenje s tekstom:} Produkt i prěsěčenje primarnyh idealov v toj situaciji davajut nulu, a odgovarjajuče komponenty davajut direktnu sumu prstena. Nulovy ideal ne jest isto kako nulovy prsten. Formuly i uslovja razloženja, ne samo slovo nulovy, oprěděljajut, kogda ty ravnosti imajut město.\par
\noindent\emph{Textual comparison:} In this setting the product and intersection of the primary ideals give zero, while the corresponding components give a direct sum of the ring. A zero ideal is not a zero ring. The decomposition formulas and hypotheses, not the adjective zero alone, determine when these equalities hold.\par
\noindent{\footnotesize TEXT-COMP-351; pasaže / passages: 17469, 17471, 17472, 17474, 17476, 17478, 17480, 17482, 17484, 17486, 17487, 17488, 17491, 17493, 17499; \texttt{CORPUS\_FORM\_COMPARISONS.json}.}\par
\end{minipage}\par
\par\medskip\Needspace{7\baselineskip}
\noindent\begin{minipage}{\linewidth}
\hypertarget{TERM-1259}{}
\noindent\textbf{nulovy ideal / jediničny ideal}\quad{\footnotesize\texttt{TERM-1259}}\par
\noindent\emph{Iztočnik / Source:} Nullideal / Einheitsideal\par
\noindent\emph{Rusky / Russian:} нулевой идеал / единичный идеал\par
\noindent\emph{Ukrajinsky / Ukrainian:} нульовий ідеал / одиничний ідеал\par
\smallskip\noindent\emph{Razlog:} Nulovy ideal i jediničny ideal v rabotě 22 označajut konce poslědovateljnoj eliminacije. Arhivny zapis čuva te konce, ne predlagaje novu slovesnu konkurenciju. Poslě odstranjenja vseh polinomnyh prěmennyh poslednji ideal leži v polju koeficientov i zato jest nulovy ili jediničny. Postup može skončiti nulovym idealom už raněje; iz togo ne slěduje, že izhodny ideal byl nulovy. Naopak, vyvod o jediničnom slučaju v tekstu upotrěbjaje ustanovjene relacije dělimosti. Simboly nuly i jedinice ne zaměnjajut rozličenje medžu elementom, idealom i cělym poljem. Ruska i ukrajinska para imajut tu samu terminologičnu strukturu. Objasnjenje trěba vezati k tomu postupu, a ne učiniti iz vseh idealov prosto dvě kategorije.\par
\noindent\emph{Reason:} In Paper 22 the zero and unit ideals name endpoints of successive elimination. The archival entry preserves those endpoints rather than proposing a new lexical competition. Once all polynomial variables have been removed, the final ideal lies in the coefficient field and is therefore zero or the unit ideal. The procedure can terminate at zero earlier; this does not imply that the original ideal was zero. Conversely, the text's conclusion about the unit case uses the established divisibility relations. Zero and one symbols do not remove the distinction between an element, an ideal and the entire field. The Russian and Ukrainian pair has the same terminological structure. The explanation belongs to this procedure, not a classification of all ideals into just two kinds.\par
\noindent{\footnotesize Izvorne zapisy / Original records: HIST-g090-007, HIST-g092-061.}\par
\smallskip\noindent\emph{Sravnjenje s tekstom:} Poslednji ideal bez polinomnyh prěmennyh jest idealom polja, zato ima samo nulovu ili jediničnu vozmožnost. Ranějši nulovy konec procedury ne govori, že izhodny ideal byl nulovy. Ta razlika dolžna ostati v slovarnom objasnjenju i pri každym prěnosu termina v inoj kontekst.\par
\noindent\emph{Textual comparison:} The final ideal without polynomial variables is an ideal of a field, leaving only zero or unit. An earlier zero endpoint does not say that the initial ideal was zero. That distinction belongs in the dictionary explanation and in any transfer of the term to another context.\par
\noindent{\footnotesize TEXT-COMP-352; pasaže / passages: 14935, 14937, 14939, 14941, 14943, 14945, 14947, 14949, 14950; \texttt{CORPUS\_FORM\_COMPARISONS.json}.}\par
\end{minipage}\par
\par\medskip\Needspace{7\baselineskip}
\noindent\begin{minipage}{\linewidth}
\hypertarget{TERM-1260}{}
\noindent\textbf{nulovy ideal / prosty ideal}\quad{\footnotesize\texttt{TERM-1260}}\par
\noindent\emph{Iztočnik / Source:} Nullideal / Primideal\par
\noindent\emph{Rusky / Russian:} нулевой идеал / простой идеал\par
\noindent\emph{Ukrajinsky / Ukrainian:} нульовий ідеал / простий ідеал\par
\smallskip\noindent\emph{Razlog:} Nulovy ideal i prosty ideal tu ne tvore disjunkciju. Uslovje govori, že sam nulovy ideal jest prosty. Arhivny razlog izrično odvrgaje dodavanje pojma priradženogo prostogo ideala: to by bylo ino tvrdženje. V rabotě 31, §6, toj slučaj razgleda se nad algebraično zamknjenym osnovnym poljem i v kontekstu konečnogo ranga; otud slěduje, že razgledany prsten jest samo osnovno polje. Bez tyh uslovij ne možno zaključiti, že vsaka oblast cělosti jest osnovnym poljem. Pridavnik prosty ima idealno-teoretičny smysl, podprěty ruskoju i ukrajinskoju paraleloju, ne značenje lahky. Slovarny zapis čuva točnost predikata, bo proměna od jest prosty do ima prosty ideal proměnila by hipotezu dokaza.\par
\noindent\emph{Reason:} Here nulovy ideal and prosty ideal do not form a disjunction. The hypothesis says that the zero ideal itself is prime. The recorded reason expressly rejects introducing an associated prime, which would assert something different. Paper 31, §6, treats this case over an algebraically closed ground field and in the finite-rank setting; those assumptions imply that the ring is the ground field itself. Without them, one cannot conclude that every integral domain is its ground field. Prosty has its ideal-theoretic sense, supported by the Russian and Ukrainian parallels, not the meaning easy. The dictionary preserves the exact predicate: changing is prime to has a prime ideal would change the proof's hypothesis.\par
\noindent{\footnotesize Izvorne zapisy / Original records: HIST-g142-003.}\par
\smallskip\noindent\emph{Sravnjenje s tekstom:} Uslovje jest prostota samogo nulovogo ideala, ne bytje priradženogo prostogo ideala. Vyvod o rangu 1 upotrěbjaje konečny rang i algebraičnu zamknjenost osnovnogo polja. Vytrgnuta iz tyh hipotez fraza dala by mnogo silnějše i nepravdivo tvrdženje.\par
\noindent\emph{Textual comparison:} The condition is primality of the zero ideal itself, not the existence of an associated prime. The rank-one conclusion uses finite rank and algebraic closedness of the ground field. Detached from those hypotheses, the phrase would assert something much stronger and false.\par
\noindent{\footnotesize TEXT-COMP-353; pasaže / passages: 18350, 18352, 18354; \texttt{CORPUS\_FORM\_COMPARISONS.json}.}\par
\end{minipage}\par
\par\medskip\Needspace{7\baselineskip}
\noindent\begin{minipage}{\linewidth}
\hypertarget{TERM-1261}{}
\noindent\textbf{nuljevy operator}\quad{\footnotesize\texttt{TERM-1261}}\par
\noindent\emph{Iztočnik / Source:} Nulloperator\par
\noindent\emph{Rusky / Russian:} нулевой оператор\par
\noindent\emph{Ukrajinsky / Ukrainian:} нульовий оператор\par
\smallskip\noindent\emph{Razlog:} Nuljevy operator jest operator, ktory posylaje vsaky element modula v nulu. Arhivny razlog daje prědnost paraleli s jediničnym operatorom i dotadašnjemu projektovomu korenju nulj-; samo imenno nula moglo by skryti, že ide o dějstvo. V rabotě 34, §17, jediničny element prstena dějstvuje jedinično na jednom sčlenu razloženja modula i nuljevo na drugom. To jest važno uslovje konteksta: tu ne prědpolagaje se, že jediničny element prstena mora dějstvovati jedinično na cělom modulu. Nuljevy operator zato ne znači nulovy modul. Ruska i ukrajinska paralela podpirajut sostav frazy; same po sobě ne ustanavljajut medžuslovjanske normy ani prědnosti nuljevy nad nulovy.\par
\noindent\emph{Reason:} A nuljevy operator sends every module element to zero. The recorded reason prefers a construction parallel to identity operator and the project's established nulj- stem; the noun zero alone could conceal that an action is meant. In Paper 34, §17, the ring identity acts as the identity on one summand and as zero on the other. This context matters: modules are not assumed here to have the ring identity acting identically throughout. A zero operator therefore does not mean a zero module. The Russian and Ukrainian parallels support the phrase's composition; they do not independently establish an Interslavic norm or a preference for nuljevy over nulovy.\par
\noindent{\footnotesize Izvorne zapisy / Original records: HIST-g195-003.}\par
\smallskip\noindent\emph{Sravnjenje s tekstom:} Rabota 34, §17, razlagaje modul na čest s jediničnym dějstvom elementa $e$ i čest s nuljevym dějstvom. Formula s razlikoju $m-em$ pokazuje, začto vsi elementy prstena poničujut drugu čest. Nuljevy operator i jediničny operator zato sut različne dějstva na module, ne dva nazvanja togo samogo nulovogo objekta.\par
\noindent\emph{Textual comparison:} Paper 34, §17, decomposes the module into the part where $e$ acts identically and the part where it acts as zero. The formula using $m-em$ explains why every ring element annihilates the second part. Zero operator and identity operator are therefore different actions on a module, not two names for the same zero object.\par
\noindent{\footnotesize TEXT-COMP-354; pasaže / passages: 21587, 21589, 21591, 21593, 21594, 21595, 21596, 21597, 21598, 21599, 21600, 21601, 21602, 21603, 21604, 21605, 21607; \texttt{CORPUS\_FORM\_COMPARISONS.json}.}\par
\end{minipage}\par
\par\medskip\Needspace{7\baselineskip}
\noindent\begin{minipage}{\linewidth}
\hypertarget{TERM-1262}{}
\noindent\textbf{nuljeva točka}\quad{\footnotesize\texttt{TERM-1262}}\par
\noindent\emph{Iztočnik / Source:} Nullstelle\par
\noindent\emph{Rusky / Russian:} нуль\par
\noindent\emph{Ukrajinsky / Ukrainian:} нуль\par
\smallskip\noindent\emph{Razlog:} Dodavanje slova točka k nule jest izrično obrazložena odluka. V rabotě 34, §22, nula jest sistem koordinatov, ktory zadovoljavaje vsim uravnenjam razgledanogo ideala, ne skalar 0. Koordinaty tu sut koreni jedinice i ne morajut byti nulami. Rusky i ukrajinsky zapis dopuščajut kratko нуль; medžuslovjanska fraza nuljeva točka iměje učiniti tu razliku vidimoju. Dostupny tekst daje sistemy koordinatov pri razloženju grupovogo prstena konečnoj abeljevoj grupy. Uslovje, že karakteristika ne děli poredka grupy, i prěhod k razlagajučemu polju sut česti argumenta o različnyh korenjah. Termin ne smije skryti tyh uslovij ni zaměniti točku idealom, ktory ona oprěděljaje.\par
\noindent\emph{Reason:} Adding točka to the word for zero is an expressly reasoned choice. In Paper 34, §22, a zero is a coordinate tuple satisfying every equation of the ideal, not the scalar 0. Here its coordinates are roots of unity and need not be zero. The Russian and Ukrainian records allow the shorter нуль; nuljeva točka is intended to make the distinction explicit. The available passage gives these tuples in the decomposition of a finite abelian group's group ring. The characteristic not dividing the group order, and passage to a splitting field, belong to the argument for distinct roots. The term must not conceal those hypotheses or identify the point with its defining ideal.\par
\noindent{\footnotesize Izvorne zapisy / Original records: HIST-g200-008.}\par
\smallskip\noindent\emph{Sravnjenje s tekstom:} V rabotě 34, §22, idealy s linearnymi generatorami zadavajut sistemy koordinatov iz korenjev jedinice. Nuljeva točka govori o společnom rěšenju, ne o sistemu s nulami vo vsih koordinatah. Različenost korenjev zavisi od navedenej karakteristiky i polja; terminologično objašnjenje ne izključaje tyh hipotez.\par
\noindent\emph{Textual comparison:} In Paper 34, §22, ideals with linear generators specify tuples of roots of unity. A zero point is a common solution, not a tuple with every coordinate zero. Distinctness of the roots depends on the stated characteristic and field; the terminology explanation retains these hypotheses.\par
\noindent{\footnotesize TEXT-COMP-355; pasaže / passages: 22006, 22010, 22021, 22022, 22023, 22024, 22025, 22027, 22029, 22031, 22032, 22033, 22039, 22041, 22043, 22045, 22047, 22049; \texttt{CORPUS\_FORM\_COMPARISONS.json}.}\par
\end{minipage}\par
\par\medskip\Needspace{7\baselineskip}
\noindent\begin{minipage}{\linewidth}
\hypertarget{TERM-1263}{}
\noindent\textbf{nula / polje nuljev}\quad{\footnotesize\texttt{TERM-1263}}\par
\noindent\emph{Iztočnik / Source:} Nullstelle / Nullstellenkörper\par
\noindent\emph{Rusky / Russian:} нуль / поле нулей\par
\noindent\emph{Ukrajinsky / Ukrainian:} нуль / поле нулів\par
\smallskip\noindent\emph{Razlog:} Par nula / polje nuljev čuva svezu medžu společnym rěšenjem polinomnyh uravnenij i poljem, generovanym njegovymi koordinatami. Arhivny razlog odvrgaje prěvod, ktory by vse slučaje svel samo k polju korenja jednogo uravnenja; on takože prědpolagaje širšu razumlivost korenja nulj-, bez čitateljskogo izslědovanja. V rabotě 25 prvo nastavaje polje od korenja ireducibilnogo polinoma. Potom tekst razširjaje pojem k prostym idealam v mnogo proměnnyh: iz oblasti klasov ostatkov nastavaje polje drobej. To ne znači, že vsaky kvocient po prostom idealu uže jest poljem. Razlagajuče polje vseh korenjev v tekstu ima oddělnu konstrukciju. Pozvanje na Steinitza jest avtorova citacija, ne dokaz, že njegov original byl nanovo pročitan pri toj rekonstrukciji.\par
\noindent\emph{Reason:} The pair nula / polje nuljev retains the link between a common zero of polynomials and the field generated by its coordinates. The recorded reason rejects reducing every instance to the field of a root of one equation; it also anticipates wider recognition of nulj-, without reader testing. Paper 25 first constructs a field from a root of an irreducible polynomial, then extends the account to prime ideals in several variables: the residue-class domain is followed by its fraction field. This does not make every quotient by a prime ideal a field already. The splitting field of all roots has a separate construction. The reference to Steinitz is the author's citation, not evidence of a fresh reading of his original during this reconstruction.\par
\noindent{\footnotesize Izvorne zapisy / Original records: HIST-g102-002.}\par
\smallskip\noindent\emph{Sravnjenje s tekstom:} Raboty 24, §3, i 25 razlikujut oblast klasov ostatkov od njejnogo polja drobej i od polja generovanogo koordinatami nuly. Polje vseh korenjev ima oddělnu ulogu. Rečenica raboty 24 o generatorah do i i nezavisnyh koordinatah od i+1 do n ima indeksnu nesoglasnost; ona ostavaje svědčenjem otvorenogo voprosa, ne podtvrđenjem pravilnosti vseh navedenyh formul.\par
\noindent\emph{Textual comparison:} Papers 24, §3, and 25 distinguish the residue-class domain, its fraction field and the field generated by a zero's coordinates. The field of all roots has a separate role. Paper 24's sentence about generators through i and independent coordinates from i+1 through n has an index inconsistency; it remains evidence of an open question, not confirmation that every quoted formula is correct.\par
\noindent{\footnotesize TEXT-COMP-356; pasaže / passages: 15822, 15824, 15828, 15839, 15441, 15447, 15449, 15453, 15459, 15461; \texttt{CORPUS\_FORM\_COMPARISONS.json}.}\par
\end{minipage}\par
\par\medskip\Needspace{7\baselineskip}
\noindent\begin{minipage}{\linewidth}
\hypertarget{TERM-1264}{}
\noindent\textbf{korenj uravnenja}\quad{\footnotesize\texttt{TERM-1264}}\par
\noindent\emph{Iztočnik / Source:} Nullstelle / Wurzel einer Gleichung\par
\noindent\emph{Rusky / Russian:} корень уравнения\par
\noindent\emph{Ukrajinsky / Ukrainian:} корінь рівняння / нуль функції\par
\smallskip\noindent\emph{Razlog:} Korenj uravnenja jest vybrano za konkretny krok konstrukcije v rabotě 10, §3. Zapis daje prědnost čitateljno neposrědnoj frazě pred nulovym nazvanjem v toj rečenici, ne ustanavljaje občego zakaza slova nula. Rusko korenjno nazvanje i ukrajinske korenjno i nulovo nazvanja pokazyvajut oba sposoby vyraženja. Poslě izbora obrazov rannějših bazovyh veličin, obraz slědujučej algebraičnoj veličiny izbiraje se kako korenj uravnenja s uže prěobraženymi koeficientami. Nije dosta izbrati libovoljny korenj izvornogo uravnenja. Fraza takože ne tvrdži, že korenj jest jedinstveny ili že konstrukcija s dobrym uporedkovanjem jest konečnym vyčisliteljnym algoritmom.\par
\noindent\emph{Reason:} Korenj uravnenja was chosen for a specific construction step in Paper 10, §3. The record prefers the immediately readable root-of-an-equation phrase over zero terminology in that sentence; it does not prohibit nula generally. The Russian root expression and the Ukrainian root and zero expressions show both approaches. After images of earlier basis quantities have been chosen, the next algebraic quantity's image is selected as a root of the equation with coefficients already transformed. An arbitrary root of the original equation would not suffice. Nor does the phrase assert uniqueness of the root or turn the well-ordering construction into a finite computational algorithm.\par
\noindent{\footnotesize Izvorne zapisy / Original records: HIST-g071-004.}\par
\smallskip\noindent\emph{Sravnjenje s tekstom:} Rabota 10, §3, najprvo zadavaje obrazy rannějših veličin. Uravnenje za slědujuči obraz koristi te obrazy v koeficientah. Korenj uravnenja tu čuva ulogu izbranej veličiny v prěobraženom uravnenju; ne označaje prosto lubovoljne rěšenje izvornogo polinoma.\par
\noindent\emph{Textual comparison:} Paper 10, §3, first specifies images of earlier quantities. The equation for the next image uses those images in its coefficients. Root of the equation here retains the selected quantity's role in the transformed equation, rather than denoting an arbitrary solution of the original polynomial.\par
\noindent{\footnotesize TEXT-COMP-357; pasaže / passages: 8522, 8523, 8524, 8525; \texttt{CORPUS\_FORM\_COMPARISONS.json}.}\par
\end{minipage}\par
\par\medskip\Needspace{7\baselineskip}
\noindent\begin{minipage}{\linewidth}
\hypertarget{TERM-1265}{}
\noindent\textbf{nula prostogo ideala}\quad{\footnotesize\texttt{TERM-1265}}\par
\noindent\emph{Iztočnik / Source:} Nullstelle eines Primideals\par
\noindent\emph{Rusky / Russian:} нулевая точка простого идеала\par
\noindent\emph{Ukrajinsky / Ukrainian:} нульова точка простого ідеалу\par
\smallskip\noindent\emph{Razlog:} Nula prostogo ideala jest kratko nazvanje společnogo rěšenja polinomnyh uravnenij: vsaky polinom ideala ima na tom rěšenju vrijednost 0. Ruska i ukrajinska paralela kažu točka jasněje; arhiv izbral medžuslovjansko nula i ostavil razliku medžu točkoju i množstvom toček kontekstu. Ta kratkost potrěbuje opatrnost. V rabotě 19, §12, polinom mora izčezati na vsih konečnyh nuljah ideala, ne samo na jednoj. Takože, ako produkt izčezaje na cělom nulovom množstvu prostogo ideala, jediny od faktorov izčezaje na cělom množstvu; to nije samo različny izbor faktora pri vsakoj točkě. Tekst rabotaje nad kompleksnymi koeficientami i izrično pomina algebraičnu zamknjenost. Kratky termin ne zaměnjaje tyh kvantifikatorov i uslovij.\par
\noindent\emph{Reason:} Nula prostogo ideala is a compact name for a common zero: every polynomial of the ideal has value 0 at that tuple. The Russian and Ukrainian parallels say point more explicitly; the archive chose Interslavic nula and left the distinction between a point and a set of points to context. That economy requires care. In Paper 19, §12, a polynomial must vanish at all finite zeros of the ideal, not just one. Likewise, if a product vanishes on the whole zero set of a prime ideal, one factor vanishes on that whole set; this is not merely a different choice of factor at each point. The passage uses complex coefficients and expressly mentions algebraic closure. The short term cannot replace these quantifiers and hypotheses.\par
\noindent{\footnotesize Izvorne zapisy / Original records: HIST-g081-080.}\par
\smallskip\noindent\emph{Sravnjenje s tekstom:} Rabota 19, §12, izrično govori o vseh konečnyh nuljah i o algebraično zamknjenoj oblasti vrijednostij. Teoremny zaključok o pripadnosti idealu ne slěduje iz izčezanja na jednoj točkě. Pri produktu, jediny faktor mora izčezati na cělom nulovom množstvu; izbor raznogo faktora pro razne točky jest slabše tvrdženje.\par
\noindent\emph{Textual comparison:} Paper 19, §12, explicitly refers to all finite zeros and an algebraically closed range of values. Ideal membership does not follow from vanishing at one point. For a product, one factor must vanish on the entire zero set; choosing different factors at different points is a weaker assertion.\par
\noindent{\footnotesize TEXT-COMP-358; pasaže / passages: 13461, 13462, 13466, 13467, 13468, 13470; \texttt{CORPUS\_FORM\_COMPARISONS.json}.}\par
\end{minipage}\par
\par\medskip\Needspace{7\baselineskip}
\noindent\begin{minipage}{\linewidth}
\hypertarget{TERM-1266}{}
\noindent\textbf{nulje}\quad{\footnotesize\texttt{TERM-1266}}\par
\noindent\emph{Iztočnik / Source:} Nullstellen\par
\noindent\emph{Rusky / Russian:} нули\par
\noindent\emph{Ukrajinsky / Ukrainian:} нулі\par
\smallskip\noindent\emph{Razlog:} Množina nulje čuva algebraično nazvanje, naměsto systematičnogo prěhoda k geometričnym točkam. Zapis povězuje toj izbor s povtorjajučim se terminom polje nuljev v rabotě 24. V §1 nuly nastavajut nad poljem s parametrami; njihove koordinaty ne morajut ležati v prvonačalnom koeficientnom polju. Faktory rezultanta pokazyvajut razgledane nuly i njihove kratnosti. Zato nulje ne znači ni vse koordinaty ravne nule, ni samo čislove točky, ni nužno množstvo bez povtorjenij. Rusko нули i ukrajinsko нулі sut zachovane paralele. Oblik nulje jest istorijsky izbor zapisa; sovpadanje s korpusom samo po sobě ne dokazuje, že on jest jediny pravilny množinski oblik v medžuslovjanskom.\par
\noindent\emph{Reason:} The plural nulje retains algebraic zero terminology instead of systematically replacing it with geometric points. The record links that choice to the recurrent polje nuljev in Paper 24. In §1 the zeros are taken over a field containing parameters; their coordinates need not lie in the original coefficient field. Resultant factors describe these zeros with their multiplicities. Nulje therefore means neither that all coordinates equal zero, nor only numerical points, nor necessarily a collection without repetition. Russian нули and Ukrainian нулі are the preserved parallels. Nulje is the historical recorded choice; occurrence in the corpus does not itself establish it as the uniquely correct Interslavic plural.\par
\noindent{\footnotesize Izvorne zapisy / Original records: HIST-g095-003.}\par
\smallskip\noindent\emph{Sravnjenje s tekstom:} Rabota 24, §1, razgleda algebraične vrijednosti nad poljem s ostalymi proměnnymi kako parametrami. Produkt linearnyh faktorov rezultanta čuva takože kratnosti. Zapis o pismenosti v §3 prodolžaje termin nulj-; sam latinicky kontekst ne jest novym testom kiriličnogo konvertera ni nezavisnym svědčenjem jezykovoj upotrěby.\par
\noindent\emph{Textual comparison:} Paper 24, §1, considers algebraic values over a field retaining the other variables as parameters. The resultant's product of linear factors also carries multiplicities. The §3 script record continues nulj- terminology; the Latin context is not itself a new test of the Cyrillic converter or independent language attestation.\par
\noindent{\footnotesize TEXT-COMP-359; pasaže / passages: 15263, 15341, 15343, 15344, 15346; \texttt{CORPUS\_FORM\_COMPARISONS.json}.}\par
\end{minipage}\par
\par\medskip\Needspace{7\baselineskip}
\noindent\begin{minipage}{\linewidth}
\hypertarget{TERM-1268}{}
\noindent\textbf{polje nuljev}\quad{\footnotesize\texttt{TERM-1268}}\par
\noindent\emph{Iztočnik / Source:} Nullstellenkörper\par
\noindent\emph{Rusky / Russian:} поле нулей\par
\noindent\emph{Ukrajinsky / Ukrainian:} поле нулів\par
\smallskip\noindent\emph{Razlog:} Polje nuljev jest doslovny, sostavny termin, izbrany za povtorjajuči pojem v dokazě raboty 24. To nije polje, čiji vse elementy sut nule, ni avtomatično razlagajuče polje. V §3 prosty ideal daje oblast klasov ostatkov; poslě dodavanja drobej dobyvaje se polje, izomorfno polju generovanomu koordinatami společnoj nuly nad danym osnovnym poljem. Sravnenje mora čuvati parametry i osnovno polje. Zapis ne davaje čislovyh ocěnok kandidatom. Pri čitanju dostupnogo prěvoda odkryva se nerazrěšena nesoglasnost indeksov: rečenica najprvo govori o generatorah do $i$, potom vybiraje medžu njimi $n-i$ nezavisnyh i imenuje ih indeksami od $i+1$ do $n$. To potrěbuje sravnenje s izvornym svědkom; slovar ne popravjaje citatu molčalivo.\par
\noindent\emph{Reason:} Polje nuljev is a literal, compositional term selected for a recurring concept in Paper 24's argument. It is not a field whose elements are all zero, nor automatically a splitting field. In §3 a prime ideal gives a residue-class domain; adjoining fractions yields a field isomorphic to that generated by a common zero's coordinates over the specified ground field. The parameters and ground field must remain explicit. The record supplies no numerical candidate scores. Reading the available translation reveals an unresolved index discrepancy: a sentence first gives generators only through $i$, then selects $n-i$ independent ones among them and names them with indices from $i+1$ through $n$. Comparison with the original witness is needed; the dictionary does not silently repair the quotation.\par
\noindent{\footnotesize Izvorne zapisy / Original records: HIST-g095-004.}\par
\smallskip\noindent\emph{Sravnjenje s tekstom:} Raboty 24, §3, i 25 razlikujut oblast klasov ostatkov od njejnogo polja drobej i od polja generovanogo koordinatami nuly. Polje vseh korenjev ima oddělnu ulogu. Rečenica raboty 24 o generatorah do i i nezavisnyh koordinatah od i+1 do n ima indeksnu nesoglasnost; ona ostavaje svědčenjem otvorenogo voprosa, ne podtvrđenjem pravilnosti vseh navedenyh formul.\par
\noindent\emph{Textual comparison:} Papers 24, §3, and 25 distinguish the residue-class domain, its fraction field and the field generated by a zero's coordinates. The field of all roots has a separate role. Paper 24's sentence about generators through i and independent coordinates from i+1 through n has an index inconsistency; it remains evidence of an open question, not confirmation that every quoted formula is correct.\par
\noindent{\footnotesize TEXT-COMP-356; pasaže / passages: 15822, 15824, 15828, 15839, 15441, 15447, 15449, 15453, 15459, 15461; \texttt{CORPUS\_FORM\_COMPARISONS.json}.}\par
\end{minipage}\par
\par\medskip\Needspace{7\baselineskip}
\noindent\begin{minipage}{\linewidth}
\hypertarget{TERM-1269}{}
\noindent\textbf{nulovy dělitelj}\quad{\footnotesize\texttt{TERM-1269}}\par
\noindent\emph{Iztočnik / Source:} Nullteiler\par
\noindent\emph{Rusky / Russian:} делитель нуля\par
\noindent\emph{Ukrajinsky / Ukrainian:} дільник нуля\par
\smallskip\noindent\emph{Razlog:} Nulovy dělitelj čuva pojem nenulovogo elementa, ktory s drugym nenulovym elementom daje produkt 0. Arhivny razlog ukazuje na tradicionalnu terminologiju i na važno različenje dvou prěhodov v rabotě 31, §1: dodavanje kopije polja i razširenje koeficientnogo polja. V priměrě $\mathbb{Q}[x]/(x^2+1)$ jest polje, ale poslě razširenja koeficientov do $\mathbb{Q}(i)$ isti kvocientny zapis ima faktory $x-i$ i $x+i$, čiji produkt jest nula. To ne znači, že samo razširjeno polje ima dělitelje nuly. Sostav prstena i njegova relacija ostavajut čestju objašnjenja. Oblik nulovy razlikuje se od nuljevy v drugih zapisah; podoben smysl ne daje osnovu za slepu zaměnu vsih pojavov.\par
\noindent\emph{Reason:} Nulovy dělitelj retains the notion of a nonzero element whose product with another nonzero element is 0. The recorded reason invokes traditional terminology and the important distinction in Paper 31, §1, between adjoining a field copy and enlarging the coefficient field. In the example, $\mathbb{Q}[x]/(x^2+1)$ is a field, whereas after extending coefficients to $\mathbb{Q}(i)$ the same quotient presentation has factors $x-i$ and $x+i$ whose product is zero. This does not say that the extended coefficient field itself has zero divisors. The algebra and its defining relation remain part of the explanation. Nulovy differs from nuljevy elsewhere; related meanings do not justify a blind replacement of every occurrence.\par
\noindent{\footnotesize Izvorne zapisy / Original records: HIST-g119-012.}\par
\smallskip\noindent\emph{Sravnjenje s tekstom:} Priměr iz primětky raboty 31, §1, čuva i prsten polinomov i relaciju $x^2+1=0$ pri proměně koeficientov. Nenulove klasy $x-i$ i $x+i$ potom imajut nulovy produkt. To objašnjaje razliku medžu razširenjem koeficientov danej algebry i tvrdženjem, že samo polje ima dělitelje nuly.\par
\noindent\emph{Textual comparison:} The footnote in Paper 31, §1, retains both the polynomial algebra and its relation $x^2+1=0$ while changing coefficients. The nonzero classes $x-i$ and $x+i$ then have zero product. This explains the distinction between extending an algebra's coefficients and claiming that the field itself has zero divisors.\par
\noindent{\footnotesize TEXT-COMP-360; pasaže / passages: 17316, 17324, 17326, 17363; \texttt{CORPUS\_FORM\_COMPARISONS.json}.}\par
\end{minipage}\par
\par\medskip\Needspace{7\baselineskip}
\noindent\begin{minipage}{\linewidth}
\hypertarget{TERM-1270}{}
\noindent\textbf{ideal-dělitelj nuly}\quad{\footnotesize\texttt{TERM-1270}}\par
\noindent\emph{Iztočnik / Source:} Nullteilerideal\par
\noindent\emph{Rusky / Russian:} идеал-делитель нуля\par
\noindent\emph{Ukrajinsky / Ukrainian:} ідеал-дільник нуля\par
\smallskip\noindent\emph{Razlog:} Ideal-dělitelj nuly prenosi podmět děljenja nuly od elementa k idealu. Zapis daje prědnost prozračnomu složenju, zato že rabota 30, §5, oprěděljaje pojem na městě črez prsten klasov ostatkov. Razgledany nenulovy ideal može iměti nenulovy ideal, čiji produkt s njim jest nulovy ideal. Fraza ne znači ideal sostavjen iz vsih děliteljev nuly, ni prosto nulovy ideal. To različenje jest nužno za slabo i silno primarny uslovja: nilpotentnost vsakogo elementa ne daje bez dodatnyh uslovij jedinu stepen, ktora poničuje cěly ideal. Tekst davaje protivpriměr s nekonečno mnogo proměnnymi i osobno koristi konečnu generovanost pri drugom zaključku. Ruska i ukrajinska sostavne paralele ne smijut izbrisati tu razliku objekta.\par
\noindent\emph{Reason:} Ideal-dělitelj nuly moves the subject of zero division from an element to an ideal. The record prefers a transparent compound because Paper 30, §5, defines the notion locally through the residue-class ring. The nonzero ideal in question can have a nonzero ideal whose product with it is the zero ideal. The phrase does not mean the ideal consisting of all zero divisors, or simply the zero ideal. This distinction is necessary for weak and strong primary conditions: nilpotence of each element does not, without further assumptions, give a single exponent annihilating the whole ideal. The passage provides an infinitely generated counterexample and separately uses finite generation in the other direction. The Russian and Ukrainian compounds must not erase the distinction of objects.\par
\noindent{\footnotesize Izvorne zapisy / Original records: HIST-g112-002.}\par
\smallskip\noindent\emph{Sravnjenje s tekstom:} Rabota 30, §5, različaje nilpotentnost elementov i nilpotentnost idealov, s različnymi slabymi i silnymi uslovjami. Njih ekvivalenciju ne možno podrazuměvati bez dodatnyh hipotez. V dokaznom tekstu jest takože nerazrěšena jednakost $a^x b^\lambda=(ab)^x$ bez navedenogo uslovja $\lambda=x$; zachovany odlomok ne služi kako podtvrđenje pravilnosti toj formule.\par
\noindent\emph{Textual comparison:} Paper 30, §5, distinguishes nilpotence of elements from nilpotence of ideals, with different weak and strong conditions. Their equivalence cannot be assumed without further hypotheses. The proof text also has the unresolved equality $a^x b^\lambda=(ab)^x$ without a stated assumption $\lambda=x$; preserving the passage does not endorse that formula.\par
\noindent{\footnotesize TEXT-COMP-361; pasaže / passages: 16601, 16603, 16605, 16607, 16609, 16611, 16613, 16615; \texttt{CORPUS\_FORM\_COMPARISONS.json}.}\par
\end{minipage}\par
\par\medskip\Needspace{7\baselineskip}
\noindent\begin{minipage}{\linewidth}
\hypertarget{TERM-1271}{}
\noindent\textbf{čislove koeficienty}\quad{\footnotesize\texttt{TERM-1271}}\par
\noindent\emph{Iztočnik / Source:} numerische Koeffizienten\par
\noindent\emph{Rusky / Russian:} числовые коэффициенты\par
\noindent\emph{Ukrajinsky / Ukrainian:} числові коефіцієнти\par
\smallskip\noindent\emph{Razlog:} Čislove koeficienty jest izbrano kako obyčajno matematično nazvanje, v paraleli s ruskoju i ukrajinskoju frazoju. Arhivny razlog daje tomu prědnost pred doslovnym internacionalizmom za německo numerisch, ale ne zachovaje točnogo alternativnogo medžuslovjanskogo oblika. Zato takogo kandidata ne dopisujemo. V rabotě 2, §2, ide o koeficienty v formulah kontrakciji; nekotore sut vyražene odnošenjami binomialnyh čisel, druge opreděljajut se iz prědhodnyh koeficientov. Čislovy ne znači približny, desetičny ili nutno cěločislovy. Indeksy mogut upravljati vrijednostju koeficienta, ne prěměnjajuči njegove ulohy. Slovar tako odlišuje koeficientnu konstantu od formy, ktora sama sodrživaje proměnne.\par
\noindent\emph{Reason:} Čislove koeficienty was chosen as ordinary mathematical wording, parallel to the Russian and Ukrainian phrases. The archived reason prefers it to a literal internationalism for German numerisch, but does not preserve the exact alternative Interslavic form; none is supplied retrospectively. In Paper 2, §2, these are coefficients in contraction formulas: some are ratios of binomial coefficients, while others are determined from preceding coefficients. Numerical does not mean approximate, decimal or necessarily integral. Indices can determine a coefficient's value without changing its role. The dictionary thus distinguishes a coefficient constant from a form that itself contains variables.\par
\noindent{\footnotesize Izvorne zapisy / Original records: HIST-g004-011.}\par
\smallskip\noindent\emph{Sravnjenje s tekstom:} V rabotě 2, §2, koeficienty dana sut točnymi formulami i rekurentnym odnošenjem k rannějšim koeficientam. Zato čislove označaje njih ulogu pri formah, ne približnost vyčislenja. Citacija drugih avtorov v primětkě jest svědčenjem bibliografičnoj svezy, ne tvrdženjem o novom pročitanju ih rabot.\par
\noindent\emph{Textual comparison:} In Paper 2, §2, coefficients are given by exact formulas and relations to earlier coefficients. Numerical describes their role alongside the forms, not approximate computation. The footnote's references establish a bibliographic connection, not a claim that those other works were freshly read.\par
\noindent{\footnotesize TEXT-COMP-362; pasaže / passages: 402, 410, 411, 412, 413, 414, 415, 515, 516, 517, 518, 519, 520, 521, 522, 523; \texttt{CORPUS\_FORM\_COMPARISONS.json}.}\par
\end{minipage}\par
\par\medskip\Needspace{7\baselineskip}
\noindent\begin{minipage}{\linewidth}
\hypertarget{TERM-1273}{}
\noindent\textbf{najbližši vyši modul}\quad{\footnotesize\texttt{TERM-1273}}\par
\noindent\emph{Iztočnik / Source:} nächsthöherer Modul\par
\noindent\emph{Rusky / Russian:} следующий более высокий модуль\par
\noindent\emph{Ukrajinsky / Ukrainian:} наступний вищий модуль\par
\smallskip\noindent\emph{Razlog:} Najbližši vyši modul čuva poredok korakov konstrukcije v rabotě 2, §7. Arhivny razlog pomina hierarhiju modulov; ruska i ukrajinska paralela jasno vyražajut slědujuči vyši korak. Važno jest istorijsko značenje slova modul. V toj rabotě ono služi pri relativnoj polnotě sistema form: dopušče se ostatny izrazy, dobyvane kontrakcijami so sistemom modula. To ne jest samo današnji abstraktny modul s višim rangom. §7 prěhodi od modula k nazvanomu vyšemu modulu i potom k grupě vyših modulov. Pridavniki označajut toj konstrukcijny poredok, ne fizičnu vysotu ni avtomatično vključenje jednogo abstraktnogo modula v drugy.\par
\noindent\emph{Reason:} Najbližši vyši modul preserves the sequence of construction steps in Paper 2, §7. The archived reason mentions a hierarchy of moduli; the Russian and Ukrainian parallels expressly indicate the next higher step. The historical sense of modul matters. In this paper it belongs to relative completeness of a system of forms: remainder expressions obtained by contractions with the modulus system are allowed. It is not simply a modern abstract module of higher rank. Section 7 passes from a modulus to the designated higher one and then to a collection of higher moduli. The adjectives mark that construction order, not physical height or automatic inclusion between abstract modules.\par
\noindent{\footnotesize Izvorne zapisy / Original records: HIST-g008-004.}\par
\smallskip\noindent\emph{Sravnjenje s tekstom:} Oprěděljenje v rabotě 2, §3, govori o relativnoj polnotě do dodatnyh kontrakcijnyh izrazov. Prěhody v §7 koristajut toj smysl modula. Vyši zato potrěbuje konkretny red form i pravilo redukcije; ne možno go čitati kako obču teoremu o rangah abstraktnyh modulov.\par
\noindent\emph{Textual comparison:} The definition in Paper 2, §3, states relative completeness up to additional contraction expressions. The steps in §7 use that sense of modulus. Higher consequently requires the specified series of forms and reduction rule; it cannot be read as a general theorem about ranks of abstract modules.\par
\noindent{\footnotesize TEXT-COMP-364; pasaže / passages: 545, 547, 700, 1046, 1064; \texttt{CORPUS\_FORM\_COMPARISONS.json}.}\par
\end{minipage}\par
\par\medskip\Needspace{7\baselineskip}
\noindent\begin{minipage}{\linewidth}
\hypertarget{TERM-1274}{}
\noindent\textbf{slědujući vyši relativno polny sistem}\quad{\footnotesize\texttt{TERM-1274}}\par
\noindent\emph{Iztočnik / Source:} nächsthöheres relativ vollständiges System\par
\noindent\emph{Rusky / Russian:} следующая более высокая относительно полная система\par
\noindent\emph{Ukrajinsky / Ukrainian:} наступна вища відносно повна система\par
\smallskip\noindent\emph{Razlog:} Slědujući vyši relativno polny sistem razlagaje německy složeny opis v čitateljnu frazu. Arhiv ocěnjuje rusku i ukrajinsku formulaciju kako gladšu, ale pri medžuslovjanskoj izrično ostavjaje voprosy soglasovanja pridavnikov i ih poredka za prověrku. To ograničenje ne možno proměniti v tvrdženje o uže ustanovjenoj normě. V rabotě 2, §17, sistem ostavaje polnym po danym modulam; vyše i slědujući označajut prěhod konstrukcije, ne bezuslovnu polnotu. Dostupna fraza slědujučego vyšego relativno polnogo sistema pokazyvaje i sklonjeny oblik i pisanje č naměsto arhivnogo ć. Različenje formulacij trěba razgledati v rečenici, bez slepogo prěnosa starogo pravopisa v nove prěvody.\par
\noindent\emph{Reason:} Slědujući vyši relativno polny sistem expands the German compound description into a readable phrase. The archive judges the Russian and Ukrainian wording smoother, but expressly leaves Interslavic adjective agreement and comparative word order for review. That qualification must not become a claim of an already established norm. In Paper 2, §17, completeness remains relative to specified moduli; higher and next indicate the construction step, not unconditional completeness. The available phrase slědujučego vyšego relativno polnogo sistema shows both inflection and č in place of archival ć. The formulations need comparison in their sentences, without blindly transferring older spelling to new translations.\par
\noindent{\footnotesize Izvorne zapisy / Original records: HIST-g018-003.}\par
\smallskip\noindent\emph{Sravnjenje s tekstom:} Dostupny §17 raboty 2 koristi frazu slědujučego vyšego relativno polnogo sistema, potom davaje slědujuči modul i sistemy form. Tekst podpiraje relativny, korakovy smysl frazy. On ne razrěšaje sam po sobě arhivny vopros o najlepšem slovoporedku i ne opravdavaje prěnosa pisma ć bez sravnenja s nynějšim č.\par
\noindent\emph{Textual comparison:} The available §17 of Paper 2 uses slědujučego vyšego relativno polnogo sistema, then supplies the next modulus and systems of forms. It supports the relative, stepwise sense of the phrase. By itself it does not settle the archived word-order question or justify transferring ć without comparison with current č.\par
\noindent{\footnotesize TEXT-COMP-365; pasaže / passages: 2387, 2389, 2391, 2394, 2395; \texttt{CORPUS\_FORM\_COMPARISONS.json}.}\par
\end{minipage}\par
\par\medskip\Needspace{7\baselineskip}
\noindent\begin{minipage}{\linewidth}
\hypertarget{TERM-1276}{}
\noindent\textbf{gorny / dolny indeks vagy}\quad{\footnotesize\texttt{TERM-1276}}\par
\noindent\emph{Iztočnik / Source:} oberer / unterer Gewichtsindex\par
\noindent\emph{Rusky / Russian:} верхний / нижний весовой индекс\par
\noindent\emph{Ukrajinsky / Ukrainian:} верхній / нижній ваговий індекс\par
\smallskip\noindent\emph{Razlog:} Gorny i dolny v tom paru označajut položenje indeksa, ne najvyššu i najnižšu vrijednost vagy. V rabotě 4, §1, gorny indeks $\rho$ broji symboly jednogo roda, a dolny $n-\rho$ drugogo; vaga reda jest $\rho(n-\rho)$. Isti red može zato byti označeny črez gorny ili dolny indeks. Arhivny razlog povězuje imenovanje neposrědno s položenjem vo formulah; ruska i ukrajinska paralela podpirajut toj sostav. Slovo vaga tu ima avtorovo kombinatorno značenje. Ono ne označaje ocěnku terminologičnogo kandidata, statističnu važnost ni čislovu ocěnku dověrja. Ta razlika jest osobno važna v slovaru, ktory takože objašnjaje kvalitativno odvažanje jezykovyh izborov.\par
\noindent\emph{Reason:} Gorny and dolny here denote index position, not the greatest and least values of the weight. In Paper 4, §1, the upper index $\rho$ counts symbols of one kind and the lower index $n-\rho$ the other; the row's weight is $\rho(n-\rho)$. The same row can consequently be designated by an upper or a lower index. The archived reason connects the names directly with their placement in formulas; the Russian and Ukrainian parallels support that composition. Weight has the author's combinatorial meaning. It is not a terminology candidate's score, statistical importance or a numerical confidence rating. This distinction matters especially in a dictionary that also discusses qualitative weighing of language choices.\par
\noindent{\footnotesize Izvorne zapisy / Original records: HIST-g030-015.}\par
\smallskip\noindent\emph{Sravnjenje s tekstom:} Rabota 4, §1, daje gorny indeks $\rho$, dolny indeks $n-\rho$ i vagu $\rho(n-\rho)$. Točna oznaka povězuje dva sposoby zapisa togo samogo reda. Ulohu vsakogo indeksa trěba čuvati; zaměna gorny/dolny s veči/menši proměnila by smysl.\par
\noindent\emph{Textual comparison:} Paper 4, §1, gives the upper index $\rho$, lower index $n-\rho$ and weight $\rho(n-\rho)$. The notation connects two ways of designating the same row. Each index's role must be retained; replacing upper/lower with greater/smaller would change the meaning.\par
\noindent{\footnotesize TEXT-COMP-366; pasaže / passages: 3768, 3769, 3770, 3771, 3772, 3773, 3774, 3775, 3776, 3777, 3778, 3779, 3780, 3781, 3782, 3784; \texttt{CORPUS\_FORM\_COMPARISONS.json}.}\par
\end{minipage}\par
\par\medskip\Needspace{7\baselineskip}
\noindent\begin{minipage}{\linewidth}
\hypertarget{TERM-1278}{}
\noindent\textbf{$\Omega$-proces}\quad{\footnotesize\texttt{TERM-1278}}\par
\noindent\emph{Iztočnik / Source:} Omega-Prozeß\par
\noindent\emph{Rusky / Russian:} $\Omega$-процесс\par
\noindent\emph{Ukrajinsky / Ukrainian:} $\Omega$-процес\par
\smallskip\noindent\emph{Razlog:} Symbol $\Omega$ ostavaje čestju nazvanja operacije. Arhivny razlog izrično daje prědnost čuvanju Noetherinoj oznaki pred opisnym prěimenovanjem operatora. Ruska i ukrajinska paralela takože čuvajut grčku bukvu, s prilagodjenym slovom proces. V rabotě 2, §1, normalna forma opisana jest uslovjem, že $\Omega$-proces po ukazanyh parah proměnnyh davaje nulu. Nazvanje zato nosi konkretnu ulogu v oprěděljenju, a ne samo ukrašajuču simboliku. Sama bukva ne dokazuje svezu s vsakoj drugoju $\Omega$ v korpusu: v rabotě 34 ona, na priměr, označaje oblast operatorov. To sut različne ulohy iste oznaki, ne jeden termin za avtomatično objedinjenje.\par
\noindent\emph{Reason:} The symbol $\Omega$ remains part of the operation's name. The recorded reason explicitly prefers preserving Noether's notation to replacing the operator name with a descriptive paraphrase. Russian and Ukrainian likewise retain the Greek letter with adapted forms of process. In Paper 2, §1, normal form is characterized by the $\Omega$-process giving zero for specified pairs of variables. The name therefore has a concrete defining role rather than decorative symbolism. The letter alone does not establish a connection with every other $\Omega$ in the corpus: in Paper 34 it denotes an operator domain, for example. These are different roles of the same symbol, not one term to merge automatically.\par
\noindent{\footnotesize Izvorne zapisy / Original records: HIST-g003-004.}\par
\smallskip\noindent\emph{Sravnjenje s tekstom:} V rabotě 2, §1, priměna $\Omega$-procesa k ukazanym proměnnym poničuje normalnu formu. Zachovanje symbola daje neposrědnu svezu s tym kriterijem. To ne jest ista uloga $\Omega$, ktora v rabotě 34 imenuje množstvo operatorov; sovpadanje bukvy ne jest terminologičnym dokazom.\par
\noindent\emph{Textual comparison:} In Paper 2, §1, applying the $\Omega$-process to the specified variables annihilates a normal form. Preserving the symbol keeps a direct link to that criterion. This is not the role of $\Omega$ naming a collection of operators in Paper 34; sharing a letter is not terminological evidence.\par
\noindent{\footnotesize TEXT-COMP-368; pasaže / passages: 267, 270, 272, 274; \texttt{CORPUS\_FORM\_COMPARISONS.json}.}\par
\end{minipage}\par
\par\medskip\Needspace{7\baselineskip}
\noindent\begin{minipage}{\linewidth}
\hypertarget{TERM-1280}{}
\noindent\textbf{oblast operatorov}\quad{\footnotesize\texttt{TERM-1280}}\par
\noindent\emph{Iztočnik / Source:} Operatorenbereich\par
\noindent\emph{Rusky / Russian:} область операторов\par
\noindent\emph{Ukrajinsky / Ukrainian:} область операторів\par
\smallskip\noindent\emph{Razlog:} Oblast operatorov jest cělost danyh operatorov, ne oblast oprěděljenja jednogo operatora. Zapis čuva doslovno nazvanje, zato že ta cělost upravljaje dopustimymi podgrupami i operatornymi izomorfizmami. V rabotě 34, §1, vsaky operator davaje endomorfizm grupy; obratimost vsakogo dějstva ne jest potrěbna. Različne oznaki mogut iměti isto dějstvo, dokolko ne dodano osobno uslovje razlikovanja. Dopustima podgrupa mora byti stabilna pod vsimi danymi operatorami. Iz toga ne slěduje, že vsaka jest normalna ili karakteristična: te sveze vyžadajut osobno izbrane oblasti operatorov. Ruska i ukrajinska paralela podpirajut nazvanje. Načalno oprěděljenje samo ne naměta prstenovu strukturu oblasti.\par
\noindent\emph{Reason:} Oblast operatorov is the collection of specified operators, not the domain of definition of one operator. The record retains the literal name because this collection governs admissible subgroups and operator isomorphisms. In Paper 34, §1, every operator induces a group endomorphism; invertibility of every action is not required. Distinct operator labels may have the same action unless a separate distinguishing condition is added. An admissible subgroup must be stable under all specified operators. It is not therefore automatically normal or characteristic: those connections require particular operator collections. The Russian and Ukrainian parallels support the name. The initial definition does not itself impose a ring structure on the collection.\par
\noindent{\footnotesize Izvorne zapisy / Original records: HIST-g179-001.}\par
\smallskip\noindent\emph{Sravnjenje s tekstom:} V rabotě 34, §1, dopustimost podgrupy znači stabilnost pod vsimi operatorami danej oblasti. Vnutrne automorfizmy i vsi automorfizmy davajut osobne slučaje normalnosti i karakterističnosti. Pozdějši operatorny homomorfizm mora usoglasiti grupovo zobraženje s operatornym dějstvom, ne samo čuvati grupovy produkt.\par
\noindent\emph{Textual comparison:} In Paper 34, §1, subgroup admissibility means stability under every operator in the specified collection. Inner automorphisms and all automorphisms give the special normal and characteristic cases. A later operator homomorphism must respect the operator action as well as the group product.\par
\noindent{\footnotesize TEXT-COMP-370; pasaže / passages: 19801, 19803, 19805, 19807, 19809, 19811, 19813, 19815, 19817, 19891; \texttt{CORPUS\_FORM\_COMPARISONS.json}.}\par
\end{minipage}\par
\par\medskip\Needspace{7\baselineskip}
\noindent\begin{minipage}{\linewidth}
\hypertarget{TERM-1281}{}
\noindent\textbf{operatorny homomorfizm / operatorno-homomorfny}\quad{\footnotesize\texttt{TERM-1281}}\par
\noindent\emph{Iztočnik / Source:} operatorhomomorph\par
\noindent\emph{Rusky / Russian:} операторный гомоморфизм / операторно-гомоморфный\par
\noindent\emph{Ukrajinsky / Ukrainian:} операторний гомоморфізм / операторно-гомоморфний\par
\smallskip\noindent\emph{Razlog:} Operatorny homomorfizm i pridavnik operatorno-homomorfny čuvajut ne samo grupovo zobraženje, ale i jego soglasje s danymi operatorami. Zapis raboty 34, §§18-19, izrično prodolžaje terminologiju §15 i ostavjaje spojnik v pridavniku kako kratku tehničnu konstrukciju. To jest zapisany pravopisny izbor, ne dokaz, že ina pisanja ne sut dopustime. V §15 operatorne homomorfizmy povězuju dejstvo množiteljev s modulom prědstavjenja; v §18 različajut se operatorny homomorfizm i izomorfizm. Ne možno izostaviti operatorny samo zato, že bližši kontekst govori o module. Ruska i ukrajinska paralela zachovajut istu čest značenja; čislove važnosti sravnenyh form ne sut zabilěžene.\par
\noindent\emph{Reason:} Operatorny homomorfizm and the adjective operatorno-homomorfny preserve both the group map and its compatibility with the specified operators. The record for Paper 34, §§18-19, expressly continues §15 and retains the hyphenated adjective as a compact technical construction. This is a recorded spelling choice, not evidence that other spellings are inadmissible. In §15 operator homomorphisms connect multiplier action with a representation module; §18 distinguishes operator homomorphism from isomorphism. Operatorny cannot be dropped merely because the surrounding discussion concerns modules. The Russian and Ukrainian parallels retain the same qualification; no numerical weights for the compared forms were recorded.\par
\noindent{\footnotesize Izvorne zapisy / Original records: HIST-g196-011.}\par
\smallskip\noindent\emph{Sravnjenje s tekstom:} Rabota 34, §§2, 15 i 18, čuva razliku medžu operatornym homomorfizmom i izomorfizmom pri prědstavjenjah. Spojnik v arhivnom pridavniku jest pravopisna odluka; formula o soglasju dějstvij jest matematično uslovje. Jedno ne može zaměniti drugo.\par
\noindent\emph{Textual comparison:} Paper 34, §§2, 15 and 18, retains the distinction between operator homomorphism and isomorphism in representation theory. The archival adjective's hyphen is a spelling decision; compatibility of the actions is a mathematical condition. Neither substitutes for the other.\par
\noindent{\footnotesize TEXT-COMP-371; pasaže / passages: 19891, 19893, 21422, 21424, 21426, 21428, 21430, 21669, 21729; \texttt{CORPUS\_FORM\_COMPARISONS.json}.}\par
\end{minipage}\par
\par\medskip\Needspace{7\baselineskip}
\noindent\begin{minipage}{\linewidth}
\hypertarget{TERM-1282}{}
\noindent\textbf{operatorny homomorfizm}\quad{\footnotesize\texttt{TERM-1282}}\par
\noindent\emph{Iztočnik / Source:} Operatorhomomorphismus\par
\noindent\emph{Rusky / Russian:} операторный гомоморфизм\par
\noindent\emph{Ukrajinsky / Ukrainian:} операторний гомоморфізм\par
\smallskip\noindent\emph{Razlog:} Pridavnik operatorny jest sučstvena čest oprěděljenja. V rabotě 34, §2, homomorfizm čuva grupovy produkt i mora soglasovati se s dějstvom vsakogo operatora. Ako obraz od $a$ jest $\bar a$, obraz od $\Theta a$ mora byti $\Theta\bar a$. Arhiv daje tomu uslovju prědnost pred kratšim, neoprěděljenym homomorfizmom. Podčrtna primětka razširjaje pojem na različne oblasti operatorov, ale togda potrěbno jest takože odgovarjajuče zobraženje samyh operatorov. Slovar čuva oba slučaje razděljeno. Ruska i ukrajinska fraza imajut isto eksplicitno ograničenje; sama podobnost nazvanja ne ustanavljaje, ktory operatory dane zobraženje v konkretnom priměrě čuva.\par
\noindent\emph{Reason:} Operatorny is an essential part of the definition. In Paper 34, §2, a homomorphism preserves the group product and must commute with each specified operator action. If $a$ maps to $\bar a$, then $\Theta a$ must map to $\Theta\bar a$. The archive prioritizes this condition over the shorter, unspecified homomorphism. A footnote extends the concept to different operator collections, but then a corresponding map of the operators is also required. The dictionary keeps these cases separate. Russian and Ukrainian express the same qualification; similarity of names alone does not determine which operators a particular map preserves.\par
\noindent{\footnotesize Izvorne zapisy / Original records: HIST-g180-001.}\par
\smallskip\noindent\emph{Sravnjenje s tekstom:} Oprěděljenje v rabotě 34, §2, vyžadaje čuvanje grupovogo produkta i operatornogo dějstva. Primětka razširjaje pojem na različne oblasti operatorov, s dodatnym zobraženjem medžu njimi. Bijektivnost stavaje osobnym uslovjem operatornogo izomorfizma, ne čestju vsakogo homomorfizma.\par
\noindent\emph{Textual comparison:} Paper 34, §2, requires preservation of both the group product and operator action. Its footnote extends the notion to different operator collections with an additional map between them. Bijectivity is the separate condition for operator isomorphism, not part of every homomorphism.\par
\noindent{\footnotesize TEXT-COMP-372; pasaže / passages: 19891, 19893, 19914, 19920, 19927; \texttt{CORPUS\_FORM\_COMPARISONS.json}.}\par
\end{minipage}\par
\par\medskip\Needspace{7\baselineskip}
\noindent\begin{minipage}{\linewidth}
\hypertarget{TERM-1283}{}
\noindent\textbf{operatorno izomorfny}\quad{\footnotesize\texttt{TERM-1283}}\par
\noindent\emph{Iztočnik / Source:} operatorisomorph\par
\noindent\emph{Rusky / Russian:} операторно изоморфный\par
\noindent\emph{Ukrajinsky / Ukrainian:} операторно ізоморфний\par
\smallskip\noindent\emph{Razlog:} Dva zapisy čuvajut operatorno izomorfny pri modulah i prědstavjenjah v rabotě 34 i pri Galoisovyh modulah v rabotě 37. Oba razloga odvrgajut skračenje k izomorfizmu samyh aditivnyh grup. V prvom slučaju zobraženje mora čuvati množenje s elementami danej oblasti; v drugom takože substitucije Galoisovej grupy. Polje razgledaje se kako modul nad osnovnym poljem, ne kako kvocient dvou polj. Ne mora iz toga slědovati izomorfizm polja s grupnym prstenom kako prstenov: predmetom jest modul s operatorami. Različenje izomorfizma i homomorfizma ostavaje osobno uslovje bijektivnosti. Pisanje bez spojnika jest zachovany oblik tyh zapisov, ne oddělny matematičny pojem.\par
\noindent\emph{Reason:} Two records retain operatorno izomorfny for modules and representations in Paper 34 and Galois modules in Paper 37. Both reasons reject shortening this to an isomorphism of the underlying additive groups. In the former setting the map respects multiplication by the specified operators; in the latter it also respects Galois substitutions. The field is treated as a module over its base field, not as a quotient of two fields. This need not make the field and group ring isomorphic as rings: the object is a module with operators. Bijectivity remains the additional distinction from homomorphism. The unhyphenated spelling is the form preserved in these records, not a separate mathematical concept.\par
\noindent{\footnotesize Izvorne zapisy / Original records: HIST-g197-003, HIST-g207-010.}\par
\smallskip\noindent\emph{Sravnjenje s tekstom:} Rabota 34, §18, razgleda proste modulne komponenty; rabota 37, §2, razgleda polje s Galoisovym dějstvom. V obojih slučajah izomorfizm čuva imenovanu operatornu strukturu. Ravnost rangov sama po sobě ne davaje izomorfizm proizvoljnogo zobraženja: važna jest dana konstrukcija i njejna baza.\par
\noindent\emph{Textual comparison:} Paper 34, §18, concerns simple module components; Paper 37, §2, concerns a field with Galois action. In both settings the isomorphism preserves the specified operator structure. Equality of ranks alone does not make an arbitrary map an isomorphism: the actual construction and its basis matter.\par
\noindent{\footnotesize TEXT-COMP-373; pasaže / passages: 21669, 21673, 21675, 21679, 21681, 21683, 21685, 22821, 22823, 22824, 22828, 22831, 22833, 22835, 22837, 22838, 22839, 22841, 22844; \texttt{CORPUS\_FORM\_COMPARISONS.json}.}\par
\end{minipage}\par
\par\medskip\Needspace{7\baselineskip}
\noindent\begin{minipage}{\linewidth}
\hypertarget{TERM-1284}{}
\noindent\textbf{operatorno-izomorfny}\quad{\footnotesize\texttt{TERM-1284}}\par
\noindent\emph{Iztočnik / Source:} operatorisomorph\par
\noindent\emph{Rusky / Russian:} операторно изоморфный\par
\noindent\emph{Ukrajinsky / Ukrainian:} операторно ізоморфний\par
\smallskip\noindent\emph{Razlog:} Operatorno-izomorfny prodolžaje terminologiju raboty 33 i §2 raboty 34. Arhivny razlog vyžadaje zachovanje operatornoj struktury, ne samo složenja. V §10 raboty 34 toj uslovje razlučuje nenulove prave idealy iz različnyh dvustrannyh komponentov: operator, ktory poničuje jeden ideal, ne poničuje drugy. Aditivna podobnost zato ne jest dosta. Zachovany spojnik razlikuje pisanje od operatorno izomorfny v drugih zapisah, ale ne davaje osnovu za dva različne matematične značenja. Ruska i ukrajinska paralela pišut odnos razděljeno. Izbor jedinoobraznogo pravopisa potrěbuje jezykovu ocěnku; argument s poničenjem stanovjaje semantično ograničenje, ne pravopisnu normu.\par
\noindent\emph{Reason:} Operatorno-izomorfny continues Paper 33 and §2 of Paper 34. The archived reason requires preservation of the operator structure, not addition alone. In §10 of Paper 34 this distinguishes nonzero right ideals in different two-sided components: an operator annihilating one ideal does not annihilate the other. Additive similarity is therefore insufficient. The retained hyphen differs from operatorno izomorfny elsewhere, but does not establish two mathematical meanings. Russian and Ukrainian write the relation as separate words. Selecting a uniform spelling requires linguistic assessment; the annihilator argument establishes a semantic constraint, not an orthographic norm.\par
\noindent{\footnotesize Izvorne zapisy / Original records: HIST-g188-003.}\par
\smallskip\noindent\emph{Sravnjenje s tekstom:} Rabota 34, §10, različaje idealy prěz dějstva, ktore ih poničujut. Argument potrěbuje nenulove idealy: nulovy ideal ne možno razlučiti takim neponičujučim dějstvom. To objašnjaje, začto operatorny uslovje jest silnějši od samotnogo aditivnogo izomorfizma, bez ustanovjanja prědnosti jednogo pisanja pridavnika.\par
\noindent\emph{Textual comparison:} Paper 34, §10, distinguishes ideals by their annihilating actions. The argument requires nonzero ideals: the zero ideal cannot be distinguished by a non-annihilating action. This explains why the operator condition is stronger than additive isomorphism alone, without establishing a preference for one spelling of the adjective.\par
\noindent{\footnotesize TEXT-COMP-374; pasaže / passages: 20623, 20626, 20629, 20633, 20634, 20636, 20663; \texttt{CORPUS\_FORM\_COMPARISONS.json}.}\par
\end{minipage}\par
\par\medskip\Needspace{7\baselineskip}
\noindent\begin{minipage}{\linewidth}
\hypertarget{TERM-1285}{}
\noindent\textbf{operatorno izomorfny / operatorny homomorfizm}\quad{\footnotesize\texttt{TERM-1285}}\par
\noindent\emph{Iztočnik / Source:} operatorisomorph / Operatorhomomorphismus\par
\noindent\emph{Rusky / Russian:} операторно изоморфный / операторный гомоморфизм\par
\noindent\emph{Ukrajinsky / Ukrainian:} операторно ізоморфний / операторний гомоморфізм\par
\smallskip\noindent\emph{Razlog:} V rabotě 34, §14, zapis naměrno ostavjaje operatornu terminologiju vměsto zaměny s občim opisom modulnogo zobraženja. To čuva jezyk dokaza. Množenje s vhodnym elementom davaje operatorny homomorfizm medžu prostymi pravymi idealami; nenulovost zobraženja, prostota ideala i izčrpanje obraza davajut izomorfizm. Sama jedina slova homomorfizm i izomorfny ne sut vzajemno zaměnime. Operatorno tu vyražaje soglasje s danym množenjem; ono ne znači libovoljnu funkciju medžu množstvami. Arhivna alternativa modulno zobraženje jest zachovana kako odvrgany opis, ne kako dokaz, že današnji modulny jezyk byl by matematično nepravilny.\par
\noindent\emph{Reason:} For Paper 34, §14, the record deliberately retains operator terminology instead of replacing it with a general module-map description. This preserves the proof's language. Multiplication by a suitable element gives an operator homomorphism between simple right ideals; nonzeroness, simplicity and coverage of the image give an isomorphism. Homomorphism and isomorphic are not interchangeable labels. Operatorno expresses compatibility with the specified multiplication, not an arbitrary function between sets. The module-map alternative is preserved as the rejected wording, not as evidence that modern module language would be mathematically wrong.\par
\noindent{\footnotesize Izvorne zapisy / Original records: HIST-g192-004.}\par
\smallskip\noindent\emph{Sravnjenje s tekstom:} V rabotě 34, §14, element medzi idealami zadavaje homomorfizm prěz množenje. Prosty ideal ne ima netrivialnogo jadra, a nenulovy obraz izčrpava ciljevy prosty ideal. Te uslovja, ne samo slovo homomorfizm, opravdavajut slědujuče izomorfny.\par
\noindent\emph{Textual comparison:} In Paper 34, §14, multiplication by an element defines a homomorphism between the ideals. Simplicity excludes a nontrivial kernel, and the nonzero image fills the simple target ideal. These conditions, rather than the word homomorphism alone, justify the subsequent isomorphic.\par
\noindent{\footnotesize TEXT-COMP-375; pasaže / passages: 21091, 21095, 21097, 21099, 21101, 21103, 21105, 21107; \texttt{CORPUS\_FORM\_COMPARISONS.json}.}\par
\end{minipage}\par
\par\medskip\Needspace{7\baselineskip}
\noindent\begin{minipage}{\linewidth}
\hypertarget{TERM-1286}{}
\noindent\textbf{operatorny izomorfizm}\quad{\footnotesize\texttt{TERM-1286}}\par
\noindent\emph{Iztočnik / Source:} Operatorisomorphie\par
\noindent\emph{Rusky / Russian:} [Bukvy ne sut zachovane / Lettering not recoverable]; операторный изоморфизм\par
\noindent\emph{Ukrajinsky / Ukrainian:} [Bukvy ne sut zachovane / Lettering not recoverable]; операторний ізоморфізм\par
\smallskip\noindent\emph{Razlog:} Operatorny izomorfizm sjedinjaje tri zapisy iz rabot 33, 34 i 39. Razlogi sut zachovanje operatornogo uslovja, paralela s operatornym homomorfizmom i upotrěba internacionalnogo korenja s pridavnikom, a ne s německym složenjem. Jeden rusky i jeden ukrajinsky zapis imajut izgubjene bukvy; pozdějše čitateljne zapisy ne razrěšajut rekonstrukciju izgubljenyh tekstov. Rabota 39 izrično objašnjaje izomorfizm polja i grupovoj algebry kako zobraženje, ktoro čuva linearne kombinacije i odgovarjanje Galoisovyh substitucij grupovomu množenju. To ne jest tvrdženje o izomorfizmu prstenov. Jej pozvanje na Deuringov dokaz i prazdninu rannějšego dokaza ostavaje avtorovo svědčenje, ne novy nezavisny pregled tyh dokazov.\par
\noindent\emph{Reason:} Operatorny izomorfizm brings together three records from Papers 33, 34 and 39. Their reasons are preservation of the operator condition, parallelism with operator homomorphism, and use of an international stem with an agreeing adjective rather than a German compound. One Russian and one Ukrainian record have lost lettering; later readable records do not license reconstruction of those lost strings. Paper 39 expressly describes the field-to-group-algebra isomorphism as preserving linear combinations and the correspondence between Galois substitutions and group multiplication. This is not a ring-isomorphism claim. Its reference to Deuring's proof and a gap in the earlier proof remains the author's testimony, not a fresh independent review of those proofs.\par
\noindent{\footnotesize Izvorne zapisy / Original records: HIST-g177-003, HIST-g180-002, HIST-g209-008.}\par
\smallskip\noindent\emph{Sravnjenje s tekstom:} Raboty 33 i 34 oprěděljajut operatorny izomorfizm črez soglasje s vsim danym dějstvom. Rabota 39 govori o linearnyh kombinacijah i Galoisovyh substitucijah pri izomorfizmu polja i grupovoj algebry. Ne utvrđaje se čuvanje proizvoljnogo množenja elementov polja. Rannějši dokazny propust imaje avtorovo priznanje v primětkě.\par
\noindent\emph{Textual comparison:} Papers 33 and 34 define operator isomorphism by compatibility with the specified action. Paper 39 describes linear combinations and Galois substitutions in the field-to-group-algebra isomorphism. It does not assert preservation of arbitrary field multiplication. The earlier proof gap is acknowledged by the author in a footnote.\par
\noindent{\footnotesize TEXT-COMP-376; pasaže / passages: 19687, 19689, 19891, 23240; \texttt{CORPUS\_FORM\_COMPARISONS.json}.}\par
\end{minipage}\par
\par\medskip\Needspace{7\baselineskip}
\noindent\begin{minipage}{\linewidth}
\hypertarget{TERM-1287}{}
\noindent\textbf{porjadok}\quad{\footnotesize\texttt{TERM-1287}}\par
\noindent\emph{Iztočnik / Source:} Ordnung\par
\noindent\emph{Rusky / Russian:} порядок\par
\noindent\emph{Ukrajinsky / Ukrainian:} порядок\par
\smallskip\noindent\emph{Razlog:} Porjadok byl izbrany pro algebraično-čislovoteoretičny smysl Ordnung, bez tvorjenja novogo purističnogo nazvanja. Dva razloga ukazujut na rusko i ukrajinsko порядок i prědpolagajut širšu slovjansku razumlivost; ne davajut test čitanja v zapadnyh i južnyh jezykah. V rabotě 30, §1, oprěděljenje jest širše než obyčajny konečny čislovy poredok: prsten medžu fiksovanym osnovnym prstenom i danym nadprstenom nazyvaje se poredkom. Konečnost jest dodatno uslovje konečnoj generovanosti modula. Dostupny prěvod piše poredok. Slovar čuva arhivno porjadok kako dokumentovany izbor, ne kako razlog vrnuti staru formu vo vsaky nynějšnji pojav.\par
\noindent\emph{Reason:} Porjadok was chosen for the algebraic-number-theoretic sense of Ordnung, without coining a purist replacement. Two reasons invoke Russian and Ukrainian порядок and anticipate wider Slavic recognition; they supply no West- or South-Slavic reading test. Paper 30, §1, defines the term more broadly than a usual finite number-theoretic order: a ring between the fixed base ring and the given overring is called an order. Finiteness is the additional condition of finite generation as a module. The available translation spells it poredok. The dictionary preserves archival porjadok as a documented choice, not a reason to restore it at every current occurrence.\par
\noindent{\footnotesize Izvorne zapisy / Original records: HIST-g107-010, HIST-g108-004.}\par
\smallskip\noindent\emph{Sravnjenje s tekstom:} Rabota 30, §1, oddělja oprěděljenje poredka od dopolnitelnoj konečnosti modula. V rabotě 31 uvodna primětka zatem ograničavaje vse poredky na konečne. To jest razlika medžu opreděljenjem i lokalnoju dogovorenoju kratkostju; konečny tu ne broji vse elementy prstena.\par
\noindent\emph{Textual comparison:} Paper 30, §1, separates the definition of an order from the additional finiteness of its module. Paper 31's introductory footnote then restricts all orders to finite ones. This distinguishes a definition from a local abbreviation; finite here does not count all the ring's elements.\par
\noindent{\footnotesize TEXT-COMP-377; pasaže / passages: 16133, 16174, 16176, 16190, 16192, 16194, 16196, 16198, 17267, 17269; \texttt{CORPUS\_FORM\_COMPARISONS.json}.}\par
\end{minipage}\par
\par\medskip\Needspace{7\baselineskip}
\noindent\begin{minipage}{\linewidth}
\hypertarget{TERM-1288}{}
\noindent\textbf{porjadok / konečny porjadok}\quad{\footnotesize\texttt{TERM-1288}}\par
\noindent\emph{Iztočnik / Source:} Ordnung / endliche Ordnung\par
\noindent\emph{Rusky / Russian:} порядок / конечный порядок\par
\noindent\emph{Ukrajinsky / Ukrainian:} порядок / скінченний порядок\par
\smallskip\noindent\emph{Razlog:} Konečny porjadok v toj rabotě ne znači prsten s konečno mnogo elementami. Rabota 30 oprěděljaje konečnost črez konečnu modulnu bazu nad osnovnym prstenom. Rabota 31 potom v uvodnoj primětkě izrično govori, že pod poredkom vsegda razuměje konečny poredok. Arhivny razlog čuva tu lokalnu dogovorenost i prodolžaje rano izbrano nazvanje. Kratše slovo poredok možet zato byti dovoljno vnutri toj raboty, ale v samostojnom slovarnom zapisu potrěbno jest ukazati vynešene uslovje. Ruska i ukrajinska paralela ne odstraňajut razliku medžu konečnym čislom generatorov i konečnym čislom elementov. Ne možno prěnesti uvodnu dogovorenost na vsaky drugy tekst.\par
\noindent\emph{Reason:} Konečny porjadok here does not mean a ring with finitely many elements. Paper 30 defines finiteness through a finite module generating set over the base ring. Paper 31 then expressly stipulates in its introductory footnote that order always means finite order. The recorded reason preserves this local convention and continues the earlier terminology. The shorter poredok can therefore suffice within that paper, but a standalone dictionary entry must state the suppressed hypothesis. The Russian and Ukrainian parallels do not remove the distinction between finitely many generators and finitely many elements. The introductory convention cannot be transferred to every other text.\par
\noindent{\footnotesize Izvorne zapisy / Original records: HIST-g118-002.}\par
\smallskip\noindent\emph{Sravnjenje s tekstom:} Rabota 30, §1, oddělja oprěděljenje poredka od dopolnitelnoj konečnosti modula. V rabotě 31 uvodna primětka zatem ograničavaje vse poredky na konečne. To jest razlika medžu opreděljenjem i lokalnoju dogovorenoju kratkostju; konečny tu ne broji vse elementy prstena.\par
\noindent\emph{Textual comparison:} Paper 30, §1, separates the definition of an order from the additional finiteness of its module. Paper 31's introductory footnote then restricts all orders to finite ones. This distinguishes a definition from a local abbreviation; finite here does not count all the ring's elements.\par
\noindent{\footnotesize TEXT-COMP-377; pasaže / passages: 16133, 16174, 16176, 16190, 16192, 16194, 16196, 16198, 17267, 17269; \texttt{CORPUS\_FORM\_COMPARISONS.json}.}\par
\end{minipage}\par
\par\medskip\Needspace{7\baselineskip}
\noindent\begin{minipage}{\linewidth}
\hypertarget{TERM-1289}{}
\noindent\textbf{porjadok / glavny porjadok}\quad{\footnotesize\texttt{TERM-1289}}\par
\noindent\emph{Iztočnik / Source:} Ordnung / Hauptordnung\par
\noindent\emph{Rusky / Russian:} порядок / главный порядок\par
\noindent\emph{Ukrajinsky / Ukrainian:} порядок / головний порядок\par
\smallskip\noindent\emph{Razlog:} Arhiv odděljaje Ordnung kako prstenovy objekt od poredka ili stepena form. Glavny porjadok byl izbrany kako prozračny opis, s jezykovoj paraleloju k glavnomu idealu. Ta paralela ne ustanavljaje isto matematično svojstvo. V rabotě 31, §7, glavny poredok jest cělost elementov konečnogo razširjenja polja, cělyh nad danym osnovnym prstenom. On ne jest tim samym ideal porodženy jedinym elementom. Vsaky razgledany poredok leži medžu osnovnym prstenom i tym prstenom cělyh elementov. Dostupny tekst daje formy poredok i glavny poredok; arhivne porjadok i glavny porjadok ostavajut istorijskymi oblikami. Jezykovu srodnost trěba objasniti, ne proměniti v nepravilnu definiciju.\par
\noindent\emph{Reason:} The archive distinguishes Ordnung as a ring object from the order or degree of forms. Glavny porjadok was selected as a transparent phrase, with a linguistic parallel to principal ideal. That parallel does not establish the same mathematical property. In Paper 31, §7, the main order is the collection of extension-field elements integral over the specified base ring. It is not thereby an ideal generated by one element. Every order under consideration lies between the base ring and that ring of integral elements. The available text has poredok and glavny poredok; archival porjadok and glavny porjadok remain historical forms. Linguistic resemblance needs explanation, not conversion into a false definition.\par
\noindent{\footnotesize Izvorne zapisy / Original records: HIST-g149-005.}\par
\smallskip\noindent\emph{Sravnjenje s tekstom:} Rabota 31, §7, imenuje cěly prsten integralnyh elementov glavnym poredkom. Rang inogo poredka měri se nad navedenym osnovnym prstenom; nižši rang vodi k medžupolju. Glavny v toj frazě ne znači, že vsaky ideal poredka jest glavny, ni že rang jest velkostju determinanta.\par
\noindent\emph{Textual comparison:} Paper 31, §7, calls the whole ring of integral elements the main order. Another order's rank is measured over the specified base ring; lower rank leads to an intermediate field. Main here does not mean every ideal is principal, nor that rank is the size of a determinant.\par
\noindent{\footnotesize TEXT-COMP-378; pasaže / passages: 18559, 18561; \texttt{CORPUS\_FORM\_COMPARISONS.json}.}\par
\end{minipage}\par
\par\medskip\Needspace{7\baselineskip}
\noindent\begin{minipage}{\linewidth}
\hypertarget{TERM-1290}{}
\noindent\textbf{red grupy}\quad{\footnotesize\texttt{TERM-1290}}\par
\noindent\emph{Iztočnik / Source:} Ordnung der Gruppe\par
\noindent\emph{Rusky / Russian:} порядок группы\par
\noindent\emph{Ukrajinsky / Ukrainian:} порядок групи\par
\smallskip\noindent\emph{Razlog:} Dva zapisy koristajut red grupy pro različne veličiny. V rabotě 7 $h$ jest čislo elementov konečnoj grupy i davaje granicu stepena invariantov. V rabotě 17, §11, red grupy ostatkov izrično oprěděljen jest čislom elementov linearnoj bazy nad $P$. To jest dimenzija, ne kardinalnost aditivnoj grupy; pri beskonečnom $P$ cělo množstvo može byti beskonečno. Rannějši razlog označaje red kako kratku formu, ale ostavjaje alternativu porjadok za jezykovu ocěnku. Pozdějši zapis čuva doslědnost i istorijsko bazovno značenje. Sovpadanje ruskyh i ukrajinskyh nazvanij ne dopušča sliti te smysly. Oba morajut byti vidime pro čitatelja.\par
\noindent\emph{Reason:} Two records use red grupy for different quantities. In Paper 7, $h$ counts the elements of a finite group and bounds invariant degree. In Paper 17, §11, the order of a residue group is expressly defined as the number of elements in a linear basis over $P$. It is dimension, not the cardinality of the additive group; over infinite $P$ the underlying set can be infinite. The earlier reason calls red compact but leaves porjadok for linguistic assessment. The later record preserves consistency and the historical basis-count meaning. Matching Russian and Ukrainian names do not license merging these senses. Both must be visible to the reader.\par
\noindent{\footnotesize Izvorne zapisy / Original records: HIST-g056-009, HIST-g079-063.}\par
\smallskip\noindent\emph{Sravnjenje s tekstom:} V rabotě 7 red $h$ upravljaje čislom grupnyh substitucij i graniceju stepena invariantov. Rabota 17, §11, davaje ino izrično oprěděljenje: red $\mu$ jest dimenzija grupy ostatkov nad poljem $P$. Čitatelj potrěbuje oba smysly; klasifikacija svih pojavov kako obyčajnoj kardinalnosti grupy byla by nepravilna.\par
\noindent\emph{Textual comparison:} In Paper 7, order $h$ governs the number of group substitutions and the invariant-degree bound. Paper 17, §11, gives a different explicit definition: order $\mu$ is the residue group's dimension over $P$. Readers need both senses; classifying every occurrence as ordinary group cardinality would be wrong.\par
\noindent{\footnotesize TEXT-COMP-379; pasaže / passages: 6681, 6683, 6684, 6686, 6688, 6689, 6691, 6693, 6695, 6697, 6701, 6703, 11978, 11980, 11982, 11986, 11988, 11989; \texttt{CORPUS\_FORM\_COMPARISONS.json}.}\par
\end{minipage}\par
\par\medskip\Needspace{7\baselineskip}
\noindent\begin{minipage}{\linewidth}
\hypertarget{TERM-1291}{}
\noindent\textbf{porjadok vzhodno k $o$}\quad{\footnotesize\texttt{TERM-1291}}\par
\noindent\emph{Iztočnik / Source:} Ordnung in bezug auf $o$\par
\noindent\emph{Rusky / Russian:} порядок относительно $o$\par
\noindent\emph{Ukrajinsky / Ukrainian:} порядок відносно $o$\par
\smallskip\noindent\emph{Razlog:} Vzhodno k v arhivnom nazvanju označava odnos poredka k osnovnomu prstenu $o$. Zapis daje prědnost povtorjanju togo odnosa pred novym opisom zaměny osnovy. Dostupna rečenica občeho diskriminantnogo teorema nyně piše po odnošenju k i koristi frakturno $o$; slovo ne govori o vostočnom napravjenju. To jest relativny prstenovy pojem: osnovny prsten upravljaje cělostju, modulom i diskriminantnym idealom. Pravilo čuvanja odnosa ne vyžadaje vozvrata arhivnoj formulaciji. Ni podobno slovo otnosno v sravniteljnyh jezykah, ni kratka oznaka osnovy ne dopuščajut izostaviti, nad ktorym prstenom stanovjajut se uslovja. Slovo $o$ jest matematična oznaka, ne čislo 0.\par
\noindent\emph{Reason:} Vzhodno k in the archived phrase marks the order's relation to the base ring $o$. The record prefers repeating that relation to a new base-change paraphrase. The available general discriminant theorem now says po odnošenju k and uses fraktur $o$; the phrase does not express an eastward direction. This is a relative ring-theoretic notion: the base ring governs integrality, module structure and the discriminant ideal. Preserving the relation does not require restoring the archived wording. Neither a similar relational word in the comparison languages nor an abbreviated base symbol licenses omitting the ring over which the conditions are imposed. The letter $o$ is a mathematical identifier, not the numeral 0.\par
\noindent{\footnotesize Izvorne zapisy / Original records: HIST-g163-004.}\par
\smallskip\noindent\emph{Sravnjenje s tekstom:} Rabota 31, §8, povtorjaje osnovny prsten pri poredku i diskriminantnom idealu. Dostupno po odnošenju k odpovědaje funkciji arhivnogo vzhodno k, bez vozvrata staroj frazě. Prěhod od osnovnogo prstena k njegovoj lokalizaciji menja město, v ktorom čitajut se uslovja dělitelnosti.\par
\noindent\emph{Textual comparison:} Paper 31, §8, repeats the base ring for the order and discriminant ideal. Current po odnošenju k performs the role of archival vzhodno k without restoring the older phrase. Passing from the base ring to its localization changes where the divisibility conditions are interpreted.\par
\noindent{\footnotesize TEXT-COMP-380; pasaže / passages: 19003, 19037, 19071, 19105; \texttt{CORPUS\_FORM\_COMPARISONS.json}.}\par
\end{minipage}\par
\par\medskip\Needspace{7\baselineskip}
\noindent\begin{minipage}{\linewidth}
\hypertarget{TERM-1292}{}
\noindent\textbf{poręd v koeficientah}\quad{\footnotesize\texttt{TERM-1292}}\par
\noindent\emph{Iztočnik / Source:} Ordnung in den Koeffizienten\par
\noindent\emph{Rusky / Russian:} порядок в коэффициентах\par
\noindent\emph{Ukrajinsky / Ukrainian:} порядок у коефіцієнтах\par
\smallskip\noindent\emph{Razlog:} Poręd v koeficientah čuva razliku medžu stepenom po koeficientah osnovnoj formy i stepenom po jeje proměnnyh. Arhivny razlog prodolžaje tu razliku i nazyvaje poręd prozračnym, ale izrično ne predstavlja to kako okončanu jezykovu ocěnku. Rabota 2, §7, uporedžuje formy po tom čislu i davaje tabelu kontrakcij: priměry različnih poredkov ne sut različne moči množstva form. Dostupny tekst piše pored, bez ę, i jasno čuva dopolnjenje v koeficientah. Izostavjenje togo dopolnjenja moglo by skryti, po ktoryh veličinah ide račun. Poręd, red i porjadok zato ne smijut byti avtomatično objedinjene bez različenja matematičnyh značenij.\par
\noindent\emph{Reason:} Poręd v koeficientah preserves the distinction between degree in the original form's coefficients and degree in its variables. The recorded reason continues that distinction and calls poręd transparent, while expressly leaving linguistic assessment open. Paper 2, §7, orders forms by this number and provides a contraction table: different orders are not different cardinalities of collections of forms. The available text writes pored without ę and retains v koeficientah. Dropping that complement could obscure which quantities are being counted. Poręd, red and porjadok must therefore not be automatically merged without distinguishing their mathematical meanings.\par
\noindent{\footnotesize Izvorne zapisy / Original records: HIST-g008-006.}\par
\smallskip\noindent\emph{Sravnjenje s tekstom:} Tabela v rabotě 2, §7, grupuje formy po poredku v koeficientah. Dopolnjenje v koeficientah jest čestju smysla, zato že stepen po proměnnyh jest drugo čislo. Dostupno pored ne imaje nosovogo znaka iz arhivnogo poręd; sravnenje čuva oba zapisy bez avtomatičnoj zaměny.\par
\noindent\emph{Textual comparison:} The table in Paper 2, §7, groups forms by order in the coefficients. The complement in the coefficients is part of the meaning because variable degree is a different number. Current pored lacks the nasal mark of archival poręd; the comparison preserves both spellings without automatic replacement.\par
\noindent{\footnotesize TEXT-COMP-381; pasaže / passages: 1066, 1068, 1070, 1072, 1074, 1076, 1078, 1080, 1082, 1096, 1097, 1098, 1099, 1100, 1101; \texttt{CORPUS\_FORM\_COMPARISONS.json}.}\par
\end{minipage}\par
\par\medskip\Needspace{7\baselineskip}
\noindent\begin{minipage}{\linewidth}
\hypertarget{TERM-1293}{}
\noindent\textbf{porjadky v čislovyh i funkcionalnyh poljah}\quad{\footnotesize\texttt{TERM-1293}}\par
\noindent\emph{Iztočnik / Source:} Ordnungen in Zahl- und Funktionenkörpern\par
\noindent\emph{Rusky / Russian:} порядки в числовых и функциональных полях\par
\noindent\emph{Ukrajinsky / Ukrainian:} порядки в числових і функціональних полях\par
\smallskip\noindent\emph{Razlog:} Porjadky v čislovyh i funkcionalnyh poljah prodolžaje terminologiju teorije poredkov, ne teorije usporědženij ili dobryh usporědženij. To jest izričny razlog arhivnogo izbora. Rabota 31, §4, razširjaje prědpoloženja o koeficientnom prstenu, da bi objediniti poredky i ranějše konečno-razměrne algebry nad poljem. V tom odlomku obstajanje konečnoj linearno nezavisnoj modulnoj bazy jest prědpoloženo; ono ne vyhodi iz samogo nazvanja poredok. Dostupny tekst čuva i čislove i funkcionalne polja. Relativny poredok nad konkretnym osnovnym prstenom ne jest prosto libovoljna kolekcija funkcij ani relacija menše nego. Ruska i ukrajinska paralela podpirajut terminologičnu svezu, ne zaměnjajut navedene definicije.\par
\noindent\emph{Reason:} Porjadky v čislovyh i funkcionalnyh poljah continues the terminology of orders, not order relations or well-orderings. That is the archive's express reason. Paper 31, §4, broadens assumptions on the coefficient ring to treat orders together with the preceding finite-dimensional algebras over a field. Existence of a finite linearly independent module basis is an assumption in that passage; it does not follow from the name order alone. The available text retains both number fields and function fields. A relative order over a specified base ring is neither an arbitrary collection of functions nor a less-than relation. The Russian and Ukrainian parallels support the terminological connection without replacing the definitions.\par
\noindent{\footnotesize Izvorne zapisy / Original records: HIST-g130-009.}\par
\smallskip\noindent\emph{Sravnjenje s tekstom:} Rabota 31, §4, davaje konkretne prědpoloženja o oblasti cělosti i modulnyh bazah, aby objediniti čislove i funkcionalne poredky s ranějšimi algebrajami. To oprěděljaje predmet frazy, ne pravilo uporedženja elementov. Glavny zadatok prěvoda jest čuvanje toj struktury i njejnyh hipotez.\par
\noindent\emph{Textual comparison:} Paper 31, §4, gives explicit domain and module-basis assumptions to treat number- and function-field orders together with the earlier algebras. These define the objects denoted by the phrase, not an ordering rule on elements. Translation must preserve that structure and its hypotheses.\par
\noindent{\footnotesize TEXT-COMP-382; pasaže / passages: 17825, 17827, 17829; \texttt{CORPUS\_FORM\_COMPARISONS.json}.}\par
\end{minipage}\par
\par\medskip\Needspace{7\baselineskip}
\noindent\begin{minipage}{\linewidth}
\hypertarget{TERM-1294}{}
\noindent\textbf{porjadky ranga $n$}\quad{\footnotesize\texttt{TERM-1294}}\par
\noindent\emph{Iztočnik / Source:} Ordnungen vom Range $n$\par
\noindent\emph{Rusky / Russian:} порядки ранга $n$\par
\noindent\emph{Ukrajinsky / Ukrainian:} порядки рангу $n$\par
\smallskip\noindent\emph{Razlog:} Rang $n$ označaje modulny rang poredka nad osnovnym prstenom. Arhivny razlog prodolžaje rang iz §2 i odlišuje go od poredka determinanta i stepena. V rabotě 31, §7, $n$ jest takože stepen okolnogo polja nad osnovnym poljem. To sovpadanje nastavaje iz ukazanyh uslovij; ne znači, že rang i stepen vsegda sut sinonimy. Poredok nižšego ranga poraždaje svoje medžupolje črez tvorjenje kvocientov, zato tekst može ograničiti razgledanje na polny rang. Ruska i ukrajinska paralela čuvajut rang i symbol $n$. Dostupny oblik poredky ranga $n$ razlikuje se pravopisno od arhivnogo porjadky, ne značenjem togo čisla.\par
\noindent\emph{Reason:} Rank $n$ is the module rank of the order over the base ring. The archived reason continues rang from §2 and distinguishes it from determinant order and degree. In Paper 31, §7, $n$ is also the ambient field's degree over the base field. That agreement follows from the stated setting; rank and degree are not universally synonyms. An order of lower rank generates its own intermediate field by taking fractions, so the text can restrict attention to full rank. Russian and Ukrainian preserve rank and the identifier $n$. Current poredky ranga $n$ differs orthographically from archival porjadky, not in what the number denotes.\par
\noindent{\footnotesize Izvorne zapisy / Original records: HIST-g149-006.}\par
\smallskip\noindent\emph{Sravnjenje s tekstom:} Rabota 31, §7, imenuje cěly prsten integralnyh elementov glavnym poredkom. Rang inogo poredka měri se nad navedenym osnovnym prstenom; nižši rang vodi k medžupolju. Glavny v toj frazě ne znači, že vsaky ideal poredka jest glavny, ni že rang jest velkostju determinanta.\par
\noindent\emph{Textual comparison:} Paper 31, §7, calls the whole ring of integral elements the main order. Another order's rank is measured over the specified base ring; lower rank leads to an intermediate field. Main here does not mean every ideal is principal, nor that rank is the size of a determinant.\par
\noindent{\footnotesize TEXT-COMP-378; pasaže / passages: 18559, 18561; \texttt{CORPUS\_FORM\_COMPARISONS.json}.}\par
\end{minipage}\par
\par\medskip\Needspace{7\baselineskip}
\noindent\begin{minipage}{\linewidth}
\hypertarget{TERM-1295}{}
\noindent\textbf{porjadky relativnyh polj}\quad{\footnotesize\texttt{TERM-1295}}\par
\noindent\emph{Iztočnik / Source:} Ordnungen von Relativkörpern\par
\noindent\emph{Rusky / Russian:} порядки относительных полей\par
\noindent\emph{Ukrajinsky / Ukrainian:} порядки відносних полів\par
\smallskip\noindent\emph{Razlog:} Relativny v nazvanju porjadky relativnyh polj čuva zavisnost od izbranogo osnovnogo polja i prstena. Arhiv daje prědnost internacionalno-slovjanskoj formě, so znanym matematičnym značenjem; ruska i ukrajinska paralela koristajut svoje odnošenjske pridavniki. V rabotě 31, §8, osnovny prsten uže ne mora byti prstenom glavnyh idealov. Razgledaje se širša klasa multiplikacionnyh prstenov, a vměsto jednogo diskriminantnogo elementa nastavaje diskriminantny ideal. Zato relativny ne znači približny ni uslovno pravilny. I polje, i prsten, po odnošenju k ktoromu oprěděljajut se cěle elementy, sut česti značenja. Relativny poredok ne možno bez proverky zaměniti s apsolutnym.\par
\noindent\emph{Reason:} Relativny in porjadky relativnyh polj retains dependence on a selected base field and ring. The archive prefers an international-Slavic form with mathematical relational meaning; Russian and Ukrainian use their own relational adjectives. In Paper 31, §8, the base ring need no longer be a principal ideal ring. A broader multiplication-ring class is considered, and a discriminant ideal replaces a single discriminant element. Relative therefore means neither approximate nor conditionally correct. Both the field and the ring relative to which integrality is defined belong to the meaning. A relative order cannot be replaced by an absolute one without checking the context.\par
\noindent{\footnotesize Izvorne zapisy / Original records: HIST-g155-001.}\par
\smallskip\noindent\emph{Sravnjenje s tekstom:} V rabotě 31, §8, obči osnovny prsten vyžadaje diskriminantny ideal. Čislovopoljny specialny slučaj govori o relativnom diskriminantu. Relativny ne jest samo stilističny pridavnik: osnovny prsten i razširjenje polja ostavajut neizostavime.\par
\noindent\emph{Textual comparison:} In Paper 31, §8, the more general base ring calls for a discriminant ideal. The number-field specialization concerns the relative discriminant. Relative is not stylistic: the base ring and field extension remain indispensable.\par
\noindent{\footnotesize TEXT-COMP-383; pasaže / passages: 18785, 18787, 18789, 19071, 19107; \texttt{CORPUS\_FORM\_COMPARISONS.json}.}\par
\end{minipage}\par
\par\medskip\Needspace{7\baselineskip}
\noindent\begin{minipage}{\linewidth}
\hypertarget{TERM-1296}{}
\noindent\textbf{ortogonalnostne uslovja}\quad{\footnotesize\texttt{TERM-1296}}\par
\noindent\emph{Iztočnik / Source:} Orthogonalitaetsbedingungen\par
\noindent\emph{Rusky / Russian:} [Bukvy ne sut zachovane / Lettering not recoverable]\par
\noindent\emph{Ukrajinsky / Ukrainian:} [Bukvy ne sut zachovane / Lettering not recoverable]\par
\smallskip\noindent\emph{Razlog:} Ortogonalnostne uslovja pomenovajut relacije idempotentov: $a_i^2=a_i$ i $a_i a_j=0$ za $i\ne j$. Arhivny razlog smatra ortogonalnost poznanym tehničnym opisom; zachovane ruske i ukrajinske polja však ne imajut čitateljnyh bukv, zato jih formy ne možemo podati kako točne sravnjenja. V rabotě 17, §5, to sut osobne klasy ostatkov, za ktore potrěbne produkty obstajut; ne trěba podrazuměvati libovoljno množenje vseh klasov po module. Relacije so sumoju ravnoju jedinici objašnjajut razloženje i jeho jedinstvenost. Tu ne zadano vnutrěšnje množenje vektorov ili měrenje ugla. Algebraična ortogonalnost ne mora izvirati iz geometrije.\par
\noindent\emph{Reason:} Ortogonalnostne uslovja names the idempotent relations $a_i^2=a_i$ and $a_i a_j=0$ for $i\ne j$. The archived reason treats orthogonality as familiar technical language, but the preserved Russian and Ukrainian fields have no recoverable lettering, so their exact forms cannot be supplied as comparanda. In Paper 17, §5, these are special residue classes for which the required products exist; arbitrary multiplication of all classes modulo a module must not be assumed. Together with the sum being the identity, the relations explain the decomposition and its uniqueness. No vector inner product or angle measurement is specified. Algebraic orthogonality need not arise from geometry.\par
\noindent{\footnotesize Izvorne zapisy / Original records: HIST-g079-027.}\par
\smallskip\noindent\emph{Sravnjenje s tekstom:} Rabota 17, §5, ustanavljaje, za ktore klase ostatkov potrěbne produkty sut oprěděljene. V rabotah 31 i 34 relacije idempotentov razlučujut komponenty prstena. Sveza jest algebraična: kvadrat čuva element, različne komponenty davajut nulovy produkt. Iz toga ne slěduje geometričny ugol ili centralnost vsih idempotentov.\par
\noindent\emph{Textual comparison:} Paper 17, §5, establishes which residue classes admit the required products. In Papers 31 and 34 the idempotent relations separate ring components. The connection is algebraic: squaring preserves an element and different components have zero product. Neither a geometric angle nor centrality of every idempotent follows.\par
\noindent{\footnotesize TEXT-COMP-384; pasaže / passages: 11245, 11247, 11249, 11251, 11253, 11255, 11257, 11259, 11261, 11265, 11267, 11269, 11274, 11276, 11277, 11278, 11281, 11283, 11286, 17503, 17505, 17507, 17509, 17511, 17513, 17514, 17516, 17518, 17519, 17520, 17522, 20471, 20473, 20475, 20476, 20480, 20481, 20482, 20485; \texttt{CORPUS\_FORM\_COMPARISONS.json}.}\par
\end{minipage}\par
\par\medskip\Needspace{7\baselineskip}
\noindent\begin{minipage}{\linewidth}
\hypertarget{TERM-1297}{}
\noindent\textbf{relacije ortogonalnosti}\quad{\footnotesize\texttt{TERM-1297}}\par
\noindent\emph{Iztočnik / Source:} Orthogonalitätsrelationen\par
\noindent\emph{Rusky / Russian:} соотношения ортогональности\par
\noindent\emph{Ukrajinsky / Ukrainian:} співвідношення ортогональності\par
\smallskip\noindent\emph{Razlog:} Relacije ortogonalnosti čuva internacionalny korenj vměsto novogo purističnogo opisa. Dva zapisy povězujut toj izbor s razloženjem na idempotentne komponenty v rabotah 31 i 34; ruska i ukrajinska paralela sut čitateljne i tak samo imenujut ortogonalnost. V rabotě 31 produkt elementov razloženogo prstena računa se po komponentah. V rabotě 34 iz razloženja na prave idealy dobyvajut se idempotenty, čiji mješane produkty sut nula. Ne trěba iz samej relacije zaključati, že vsaky idempotent jest centralny. Takože to nije tvrdženje o pravom uglu. Glavny kriterij izbora jest neizgubljenje postoječej algebraičnoj svezy, a širša razumlivost ostavaje jezykovym očekivanjem, ne izměrjenym rezultatom.\par
\noindent\emph{Reason:} Relacije ortogonalnosti retains the international stem instead of introducing a purist description. Two records connect the choice with idempotent decompositions in Papers 31 and 34; their Russian and Ukrainian parallels are readable and likewise name orthogonality. In Paper 31 products in the decomposed ring are computed componentwise. In Paper 34 decomposition into right ideals supplies idempotents whose mixed products are zero. The relations alone do not imply that every idempotent is central. Nor is this a claim about right angles. Preserving the established algebraic connection is the principal selection criterion; wider intelligibility remains a linguistic expectation, not a measured result.\par
\noindent{\footnotesize Izvorne zapisy / Original records: HIST-g122-005, HIST-g187-003.}\par
\smallskip\noindent\emph{Sravnjenje s tekstom:} Rabota 17, §5, ustanavljaje, za ktore klase ostatkov potrěbne produkty sut oprěděljene. V rabotah 31 i 34 relacije idempotentov razlučujut komponenty prstena. Sveza jest algebraična: kvadrat čuva element, različne komponenty davajut nulovy produkt. Iz toga ne slěduje geometričny ugol ili centralnost vsih idempotentov.\par
\noindent\emph{Textual comparison:} Paper 17, §5, establishes which residue classes admit the required products. In Papers 31 and 34 the idempotent relations separate ring components. The connection is algebraic: squaring preserves an element and different components have zero product. Neither a geometric angle nor centrality of every idempotent follows.\par
\noindent{\footnotesize TEXT-COMP-384; pasaže / passages: 11245, 11247, 11249, 11251, 11253, 11255, 11257, 11259, 11261, 11265, 11267, 11269, 11274, 11276, 11277, 11278, 11281, 11283, 11286, 17503, 17505, 17507, 17509, 17511, 17513, 17514, 17516, 17518, 17519, 17520, 17522, 20471, 20473, 20475, 20476, 20480, 20481, 20482, 20485; \texttt{CORPUS\_FORM\_COMPARISONS.json}.}\par
\end{minipage}\par
\par\medskip\Needspace{7\baselineskip}
\noindent\begin{minipage}{\linewidth}
\hypertarget{TERM-1301}{}
\noindent\textbf{$p$-adično razširjenje}\quad{\footnotesize\texttt{TERM-1301}}\par
\noindent\emph{Iztočnik / Source:} $p$-adische Erweiterung\par
\noindent\emph{Rusky / Russian:} $p$-адическое расширение\par
\noindent\emph{Ukrajinsky / Ukrainian:} $p$-адичне розширення\par
\smallskip\noindent\emph{Razlog:} $p$-adično razširjenje čuva internacionalny pridavnik i običny algebraičny pojem razširjenja. Arhivny razlog povězuje jego razumlivost s ruskoju i ukrajinskoju matematičnoju terminologijeju; konkretna čitateljska prověrka ne jest navedena. V rabotě 37 osnovno polje prěhodi k polju pri izbranom prostom idealu, a koeficienty grupovogo prstena razširjajut se odgovarjajuče. To ne znači prostě dodati jeden algebraičny korenj ni zaměniti polje njegovym ostatkovym poljem. Dostupny tekst označaje ideal frakturou; prosto p v starom glosaru ne opravdava zameny tyh oznak. Primětka razlikuje polno p-adično razširjenje od slabšego prěhoda k prstenu kvocientov v osobnom komutativnom slučaju.\par
\noindent\emph{Reason:} $p$-adično razširjenje retains the international adjective and the ordinary algebraic word for extension. The recorded reason links intelligibility to Russian and Ukrainian mathematical terminology; no reader test is reported. In Paper 37 the base field passes to the field at a specified prime ideal, and group-ring coefficients are extended accordingly. This is not simply adjunction of one algebraic root or replacement by the residue field. The available text uses a fraktur ideal symbol; the early glossary’s plain p does not justify interchanging those identifiers. A footnote distinguishes the full p-adic extension from the weaker passage to a ring of quotients in a special commutative case.\par
\noindent{\footnotesize Izvorne zapisy / Original records: HIST-g207-005.}\par
\smallskip\noindent\emph{Sravnjenje s tekstom:} Rabota 37 povězuje razširjenje osnovnogo polja s odgovarjajučim razširjenjem koeficientov grupovogo prstena. Primětka izrično razlikuje toj postupok od slabšego prěhoda k prstenu kvocientov v osobnom slučaju; ostatkovo polje ne može zaměniti ni jedno.\par
\noindent\emph{Textual comparison:} Paper 37 relates extension of the base field to corresponding extension of group-ring coefficients. A footnote expressly distinguishes this from a weaker passage to a ring of quotients in a special case; a residue field is not a substitute for either.\par
\noindent{\footnotesize TEXT-COMP-388; pasaže / passages: 22785, 22791, 22797, 22799; \texttt{CORPUS\_FORM\_COMPARISONS.json}.}\par
\end{minipage}\par
\par\medskip\Needspace{7\baselineskip}
\noindent\begin{minipage}{\linewidth}
\hypertarget{TERM-1302}{}
\noindent\textbf{$p$-adične komponenty / razvětvjene města}\quad{\footnotesize\texttt{TERM-1302}}\par
\noindent\emph{Iztočnik / Source:} $p$-adische Komponenten / Verzweigungsstellen\par
\noindent\emph{Rusky / Russian:} $p$-адические компоненты / места ветвления\par
\noindent\emph{Ukrajinsky / Ukrainian:} $p$-адичні компоненти / розгалужені місця\par
\smallskip\noindent\emph{Razlog:} $p$-adične komponenty i razvětvjene města prodolžajut ranějši lokalno-globalny slovnik: město imenuje matematičny objekt, a razvětvjenje njegovo svojstvo. Rabota 42 ne razděljuje poredky samo po jih prěsěku s podpoljem: v občej prostoj algebrě trěbuje takože ravnost lokalnyh komponentov na razvětvjenyh městah. Izostavjenje poslednogo uslovja změnilo by razdělbu. Komponenta tu ne jest proizvoljny člen direktnoj sumy; ona jest p-adična komponenta danogo poredka. Ruska i ukrajinska paralela podpirajut oba nazvanja, ale ne zaměnjajut tu definiciju. Arhivny razlog jest doslědnost, ne dokaz, že vsako razvětvjenje v drugoj oblasti ima istu matematičnu prirodu.\par
\noindent\emph{Reason:} $p$-adične komponenty and razvětvjene města continue the earlier local-global vocabulary: město names the mathematical place, and razvětvjenje its ramification property. Paper 42 does not classify orders solely by their intersection with a subfield: for a general simple algebra it also requires equality of local components at ramified places. Omitting that condition would change the classification. A component here is the p-adic component of the given order, not an arbitrary direct-sum summand. Russian and Ukrainian parallels support both names but do not replace the definition. The recorded reason is continuity, not evidence that every use of ramification in another subject denotes the same structure.\par
\noindent{\footnotesize Izvorne zapisy / Original records: HIST-g212-010.}\par
\smallskip\noindent\emph{Sravnjenje s tekstom:} V rabotě 42, §3, ravnost prěsěkov s podpoljem ne dostačuje: definicija oblasti uključaje ravnost p-adičnyh komponentov na razvětvjenyh městah. To jest dopolnitelno uslovje na poredky, ne svojstvo, ktore možno izostaviti iz nazvanja.\par
\noindent\emph{Textual comparison:} In Paper 42, §3, equality of intersections with the subfield is insufficient: the definition of a region includes equality of p-adic components at ramified places. This is an additional condition on orders, not a qualification that can be dropped from the wording.\par
\noindent{\footnotesize TEXT-COMP-389; pasaže / passages: 24451, 24453, 24654, 24657, 24662, 24665, 24668, 24670, 24675, 24677, 24679; \texttt{CORPUS\_FORM\_COMPARISONS.json}.}\par
\end{minipage}\par
\par\medskip\Needspace{7\baselineskip}
\noindent\begin{minipage}{\linewidth}
\hypertarget{TERM-1303}{}
\noindent\textbf{$p$-togo razmjera}\quad{\footnotesize\texttt{TERM-1303}}\par
\noindent\emph{Iztočnik / Source:} $p$-ter Dimension / $p$-te Dimension\par
\noindent\emph{Rusky / Russian:} $p$-го измерения\par
\noindent\emph{Ukrajinsky / Ukrainian:} $p$-го виміру\par
\smallskip\noindent\emph{Razlog:} $p$-togo razmjera jest zachovany prěvod historijnogo Dimension pri homogennyh formah. Sam arhiv ostavjaje razmjer otvorenym k jezykovoj ocěnki i nazyvaje dimenzija kako alternativu. Kontekst ale ne dopuščaje avtomatičnu zaměnu: v rabotě 12 $p$ označaje stepen homogenej formy, s kvadratičnym slučajem $p=2$, ne čislo nezavisnyh koordinat. Ranějše razvitje po $u$ oddělja homogene česti po takomu stepenju. Slovnik mora zato pokazati matematičny smysl prěd izborom jednoobraznogo nazvanja. Dostupny tekst piše razměr, ne arhivno razmjer. Ni podoba s ruskoju i ukrajinskoju formou, ni prědložena alternativa sama ne rěšaje razliku medžu stepenjem i dimenzijeju.\par
\noindent\emph{Reason:} $p$-togo razmjera is the recorded rendering of historical Dimension for homogeneous forms. The archive itself leaves razmjer open to linguistic assessment and names dimenzija as an alternative. Context prevents an automatic replacement: in Paper 12, $p$ denotes homogeneous degree, with the quadratic case $p=2$, not the number of independent coordinates. The earlier expansion in $u$ separates homogeneous parts by that degree. The dictionary must therefore explain the mathematical sense before choosing uniform wording. The available text spells razměr rather than archival razmjer. Neither resemblance to the Russian and Ukrainian forms nor the proposed alternative itself resolves the distinction between degree and dimension.\par
\noindent{\footnotesize Izvorne zapisy / Original records: HIST-g074-019.}\par
\smallskip\noindent\emph{Sravnjenje s tekstom:} V rabotě 12 razvitje po $u$ razděljuje homogene česti po stepenju; koncovy odlomok razgleda kvadratične formy kako slučaj $p=2$. Historijno Dimension zato potrěbuje kontekstualno objasnjenje, ne mehaničnu zaměnu s dimenzija.\par
\noindent\emph{Textual comparison:} Paper 12’s expansion in $u$ separates homogeneous parts by degree; its concluding passage treats quadratic forms as the case $p=2$. Historical Dimension therefore requires contextual explanation, not mechanical replacement by dimenzija.\par
\noindent{\footnotesize TEXT-COMP-390; pasaže / passages: 8976, 8978, 8979, 9097, 9099, 9100, 9120, 9122, 9125, 9127, 9152; \texttt{CORPUS\_FORM\_COMPARISONS.json}.}\par
\end{minipage}\par
\par\medskip\Needspace{7\baselineskip}
\noindent\begin{minipage}{\linewidth}
\hypertarget{TERM-1304}{}
\noindent\textbf{konjugacija črez $P$}\quad{\footnotesize\texttt{TERM-1304}}\par
\noindent\emph{Iztočnik / Source:} $P^{-1} R_a P$ / $P^{-1} C P$\par
\noindent\emph{Rusky / Russian:} сопряжение посредством $P$\par
\noindent\emph{Ukrajinsky / Ukrainian:} спряження через $P$\par
\smallskip\noindent\emph{Razlog:} Konjugacija črez $P$ jest pojasnjujuče nazvanje operacije, zapisanej formulami $P^{-1}R_aP$ i $P^{-1}CP$. Arhiv davaje prědnost zachovanju formul; nazvanje ma služiti daljšoj terminologiji klasov prědstavjenij. V rabotě 31, §4, zaměna modulnoj bazy jest dana matricou $P$ s determinantem, ktory jest jedinica v osnovnom prstenu. To jest silnějše než nenulovost determinanta nad proizvoljnym prstenom. Red faktorov slěduje iz izbranoj konvencije bazovyh redov i ne može byti obratny samo radi stilja. Ne ide o kompleksno konjugovanje ni o konjugovane korenje polinoma; ruska i ukrajinska fraza ne uklanjajut potrěbu ukazati upravo matričnu operaciju.\par
\noindent\emph{Reason:} Konjugacija črez $P$ is an explanatory name for the operations written as $P^{-1}R_aP$ and $P^{-1}CP$. The record prioritizes preservation of the formulas and supplies the name for later representation-class terminology. In Paper 31, §4, the module-basis change uses a matrix $P$ whose determinant is a unit in the base ring. This is stronger than merely being nonzero over an arbitrary ring. Factor order follows the chosen row-basis convention and cannot be reversed for stylistic reasons. The operation is neither complex conjugation nor passage to conjugate polynomial roots; the Russian and Ukrainian phrases do not remove the need to specify the matrix operation.\par
\noindent{\footnotesize Izvorne zapisy / Original records: HIST-g132-004.}\par
\smallskip\noindent\emph{Sravnjenje s tekstom:} Rabota 31, §4, izvodi $P^{-1}CP$ iz konkretnoj konvencije bazovyh redov. Determinant $P$ jest jedinica osnovnogo prstena. Red faktorov i to uslovje čuvajut smisao; nenulovy determinant sam ne daje isti rezultat nad libovoljnym prstenom.\par
\noindent\emph{Textual comparison:} Paper 31, §4, derives $P^{-1}CP$ from the particular row-basis convention. The determinant of $P$ is a unit of the base ring. Factor order and that condition carry the meaning; a nonzero determinant alone does not give the same result over an arbitrary ring.\par
\noindent{\footnotesize TEXT-COMP-391; pasaže / passages: 17909, 17911, 17913, 17916, 17917, 17918, 17921; \texttt{CORPUS\_FORM\_COMPARISONS.json}.}\par
\end{minipage}\par
\par\medskip\Needspace{7\baselineskip}
\noindent\begin{minipage}{\linewidth}
\hypertarget{TERM-1305}{}
\noindent\textbf{po parah vzajemno proste}\quad{\footnotesize\texttt{TERM-1305}}\par
\noindent\emph{Iztočnik / Source:} paarweise teilerfremd\par
\noindent\emph{Rusky / Russian:} попарно взаимно просты; попарно взаимно простые\par
\noindent\emph{Ukrajinsky / Ukrainian:} попарно взаємно прості\par
\smallskip\noindent\emph{Razlog:} Po parah vzajemno proste čuva uslovje za vsaku dvojicu, ne samo za cěly zbir faktorov. Tri zapisy daju tri svęzane razloge: prozračnu frazu vměsto zgustnjenogo pridavnika, matematičnu ulogu pri direktnom razloženju i doslědnost medžu paragrafami. V rabotě 31 ista fraza nastupaje pri idealah, pri nerazložimyh polinomah i pri primarnyh komponentah. To ne čini te objekte identičnymi: pri polinomah nad poljem prostota garantuje odgovarjajuče razloženje ostatkovogo prstena, a pri idealah trěba čitati ih srodnu definiciju. Dodatok po parah ne možno izostaviti. Arhiv ne dava čislove težiny niti nazyvaje konkretny kratši pridavnik, ktory byl odvrgnuty.\par
\noindent\emph{Reason:} Po parah vzajemno proste preserves a condition on every pair, not merely on the whole collection of factors. Three records give related reasons: a transparent phrase rather than a dense adjective, its mathematical role in direct decomposition, and consistency across sections. Paper 31 uses the phrase for ideals, irreducible polynomials and primary components. That does not identify the objects: coprimality of polynomials over a field supports the corresponding residue-ring decomposition, while ideals require their own definition. Po parah must not disappear. The archive gives no numerical weights and does not name the particular shorter adjective it rejected.\par
\noindent{\footnotesize Izvorne zapisy / Original records: HIST-g123-006, HIST-g127-008, HIST-g152-008.}\par
\smallskip\noindent\emph{Sravnjenje s tekstom:} Rabota 31 koristi istu frazu pri idealah i pri polinomah nad poljem. Uslovje prvogo roda v §3 čuva razpad na različne nerazložime faktory; razloženje ostatkovogo prstena ne slěduje samo iz slova nerazložimy. V §7 komponenty ponovno sut idealy.\par
\noindent\emph{Textual comparison:} Paper 31 uses the same phrase for ideals and for polynomials over a field. The first-kind condition in §3 preserves factorization into distinct irreducible factors; the residue-ring decomposition does not follow from the word irreducible alone. In §7 the components are again ideals.\par
\noindent{\footnotesize TEXT-COMP-392; pasaže / passages: 17553, 17555, 17557, 17559, 17561, 17702, 17704, 18715; \texttt{CORPUS\_FORM\_COMPARISONS.json}.}\par
\end{minipage}\par
\par\medskip\Needspace{7\baselineskip}
\noindent\begin{minipage}{\linewidth}
\hypertarget{TERM-1306}{}
\noindent\textbf{poparno vzaimno proste}\quad{\footnotesize\texttt{TERM-1306}}\par
\noindent\emph{Iztočnik / Source:} paarweise teilerfremd\par
\noindent\emph{Rusky / Russian:} попарно взаимно просты\par
\noindent\emph{Ukrajinsky / Ukrainian:} попарно взаємно прості\par
\smallskip\noindent\emph{Razlog:} Poparno vzaimno proste jest ranějša forma frazy o vzajemnoj prostosti za vsaku dvojicu. Zapis upotrěbja prozračno slovjansko izrazjenje i prodolžaje ranějše razloženja idealov. Dostupna rabota 19 piše poparno vzajemno proste. Važnějša od toj razlike pisanja jest razlika medžu vzajemno prostymi i vzajemno relativno prostymi idealami: §6 definira relativnu prostost jednostrannym uslovjem; §8 izrično govori, že obratna implikacija ne drži. Zato podobne frazy ne sut slobodne sinonimy. Pri vzajemnoj prostosti suma dvoh idealov jest cěly prsten; upravo to podpiraje ravnost produkta i najmenšego občego mnogokratnogo. Glosarova doslědnost ne opravdava gubitok slova relativno tam, kde ono pripada definiciji.\par
\noindent\emph{Reason:} Poparno vzaimno proste is an earlier form of the phrase requiring coprimality for every pair. Its record favors a transparent Slavic expression continuous with earlier ideal decompositions. The available Paper 19 spells it poparno vzajemno proste. More important than that spelling difference is the distinction between coprime and mutually relatively prime ideals: §6 defines relative primality by a directional condition; §8 explicitly says the converse implication fails. The similar phrases are not interchangeable synonyms. For coprime ideals their sum is the whole ring, which supports equality of product and least common multiple. Glossary consistency does not justify losing relativno where it belongs to the definition.\par
\noindent{\footnotesize Izvorne zapisy / Original records: HIST-g081-004, HIST-g082-004.}\par
\smallskip\noindent\emph{Sravnjenje s tekstom:} Rabota 19, §6, definira jednostrannu relativnu prostost. §8 oddělja od njej vzajemnu prostotu s jediničnym največšim občim děliteljem i daje priměr, gde obratna implikacija ne drži. To različenje mora ostati vidimo i pri ponovjenju frazy v nekomutativnom priměru §12.\par
\noindent\emph{Textual comparison:} Paper 19, §6, defines directional relative primality. Section 8 distinguishes coprimality with unit greatest common divisor and gives an example where the converse implication fails. That distinction must remain visible when the phrase recurs in the noncommutative example of §12.\par
\noindent{\footnotesize TEXT-COMP-393; pasaže / passages: 13062, 13064, 13066, 13278, 13280, 13316, 13318, 13320, 13322, 13324, 13326, 13328, 13330, 13332, 13333, 13337, 13666; \texttt{CORPUS\_FORM\_COMPARISONS.json}.}\par
\end{minipage}\par
\par\medskip\Needspace{7\baselineskip}
\noindent\begin{minipage}{\linewidth}
\hypertarget{TERM-1307}{}
\noindent\textbf{po parah vzajemno proste primarne ideale}\quad{\footnotesize\texttt{TERM-1307}}\par
\noindent\emph{Iztočnik / Source:} paarweise teilerfremde Primärideale\par
\noindent\emph{Rusky / Russian:} попарно взаимно простые первичные идеалы\par
\noindent\emph{Ukrajinsky / Ukrainian:} попарно взаємно прості первинні ідеали\par
\smallskip\noindent\emph{Razlog:} Po parah vzajemno proste primarne ideale sodrživaje dva nezavisne uslovja: vsaky faktor jest primarny, a vsaka dvojica jest vzajemno prosta. Arhiv izrično čuva vtory uslovje, ibo on dopuščaje zaměniti najmenše obče mnogokratno produktom. V rabotě 31, §2, konečny rang dava dvojno uslovje cěpov; s jediničnym elementom ono vodi k takomu razloženju. Ne tvrdži se, že vsako primarno razloženje v libovoljnom prstenu jest po parah vzajemno prosto. Arhivna forma ideale ne jest osnovou za mehaničnu pravopisnu zaměnu nyněšnjih padežnyh form. Slovnik čuva i primarnost, i dvojicovost, i uslovja prstena.\par
\noindent\emph{Reason:} Po parah vzajemno proste primarne ideale contains two separate conditions: each factor is primary, and each pair is coprime. The record explicitly preserves the latter because it permits replacing the least common multiple by the product. In Paper 31, §2, finite rank gives the double chain condition, which, with an identity element, leads to this decomposition. It does not say every primary decomposition in an arbitrary ring is pairwise coprime. The archival ideale spelling is not a basis for mechanically replacing current inflected forms. The dictionary preserves primaryness, the pairwise quantifier and the ring hypotheses.\par
\noindent{\footnotesize Izvorne zapisy / Original records: HIST-g121-008.}\par
\smallskip\noindent\emph{Sravnjenje s tekstom:} Rabota 30 povězuje produktnu formu s dvojicovoju vzajemnoju prostotoju; teorema V takože trěbuje komutativnost, jedinicu i dvojno uslovje cěpov. Rabota 31 dobyva to uslovje iz konečnogo ranga. Nijedny iz tyh korakov ne možno zaměniti samym pridavnikom primarny.\par
\noindent\emph{Textual comparison:} Paper 30 connects product form to pairwise coprimality; Theorem V also requires commutativity, an identity and the double chain condition. Paper 31 obtains the chain condition from finite rank. None of these steps can be replaced by the adjective primary alone.\par
\noindent{\footnotesize TEXT-COMP-394; pasaže / passages: 16507, 16509, 16511, 16794, 16796, 16798, 16800, 17434, 17436, 17438, 17440; \texttt{CORPUS\_FORM\_COMPARISONS.json}.}\par
\end{minipage}\par
\par\medskip\Needspace{7\baselineskip}
\noindent\begin{minipage}{\linewidth}
\hypertarget{TERM-1308}{}
\noindent\textbf{poparno vzaimno proste primarne idealy}\quad{\footnotesize\texttt{TERM-1308}}\par
\noindent\emph{Iztočnik / Source:} paarweise teilerfremde Primärideale\par
\noindent\emph{Rusky / Russian:} попарно взаимно простые первичные идеалы\par
\noindent\emph{Ukrajinsky / Ukrainian:} попарно взаємно прості первинні ідеали\par
\smallskip\noindent\emph{Razlog:} Poparno vzaimno proste primarne idealy povězuje tri svojstva bez ih slivanja. Arhivny razlog odkazuje na teoremu V raboty 30: pod dvojnym uslovjem cěpov i pri jedinici različne pridružene proste idealy sut vzajemno proste, zato i primarne komponenty imajut to svojstvo. Togda prědstavjenje najmenšim občim mnogokratnym stava produktnym. Sam fakt, že dva proste idealy sut različne, ne dava takogo zaključka bez ukazanyh prědpoloženij. Dostupny tekst piše vzajemno. Ta pravopisna razlika ostavaje oddělena od matematičnogo razloga, kako i alternativa po parah vměsto poparno. Nikakve čislove važnosti medžu tymi formami ne sut zachovane.\par
\noindent\emph{Reason:} Poparno vzaimno proste primarne idealy joins three properties without conflating them. The recorded reason invokes Theorem V of Paper 30: under the double chain condition and with an identity, distinct associated prime ideals are coprime, and the primary components inherit that property. A least-common-multiple representation then becomes a product representation. Distinctness of two prime ideals alone does not establish this conclusion without the stated hypotheses. The available text spells vzajemno. That spelling difference remains separate from the mathematical reason, as does the choice of po parah rather than poparno. No numerical weighting between these forms is preserved.\par
\noindent{\footnotesize Izvorne zapisy / Original records: HIST-g114-010.}\par
\smallskip\noindent\emph{Sravnjenje s tekstom:} Rabota 30 povězuje produktnu formu s dvojicovoju vzajemnoju prostotoju; teorema V takože trěbuje komutativnost, jedinicu i dvojno uslovje cěpov. Rabota 31 dobyva to uslovje iz konečnogo ranga. Nijedny iz tyh korakov ne možno zaměniti samym pridavnikom primarny.\par
\noindent\emph{Textual comparison:} Paper 30 connects product form to pairwise coprimality; Theorem V also requires commutativity, an identity and the double chain condition. Paper 31 obtains the chain condition from finite rank. None of these steps can be replaced by the adjective primary alone.\par
\noindent{\footnotesize TEXT-COMP-394; pasaže / passages: 16507, 16509, 16511, 16794, 16796, 16798, 16800, 17434, 17436, 17438, 17440; \texttt{CORPUS\_FORM\_COMPARISONS.json}.}\par
\end{minipage}\par
\par\medskip\Needspace{7\baselineskip}
\noindent\begin{minipage}{\linewidth}
\hypertarget{TERM-1309}{}
\noindent\textbf{po parah vzajemno proste primarne komponenty}\quad{\footnotesize\texttt{TERM-1309}}\par
\noindent\emph{Iztočnik / Source:} paarweise teilerfremde Primärkomponenten\par
\noindent\emph{Rusky / Russian:} попарно взаимно простые первичные компоненты\par
\noindent\emph{Ukrajinsky / Ukrainian:} попарно взаємно прості первинні компоненти\par
\smallskip\noindent\emph{Razlog:} Po parah vzajemno proste primarne komponenty prodolžaje terminologiju raboty 31, §7, v občem diskriminantnom teoremu §8. Razlog davaje prědnost prozračnoj slovjanskoj frazě radi očěkuvanoj razumlivosti za češske, slovačske i poljske čitatelje; ne prědstavjaje izměreny čitateljski rezultat i ne nazyvaje točny kratši internacionalny konkurent. V teoremu komponenty sut idealy v nadprstenu. Uslovje, že prinajmenje jedna jest vlastno primarna ili prosty ideal drugogo roda, mora ostati različne od jih po-parnoj vzajemnoj prostosti. Nazvanje komponenty ne može samo zaměniti dokaz takogo razloženja. Ni primarny ne znači vsegda prosty.\par
\noindent\emph{Reason:} Po parah vzajemno proste primarne komponenty carries the terminology of Paper 31, §7, into the general discriminant theorem of §8. The reason prefers a transparent Slavic phrase for anticipated Czech, Slovak and Polish comprehension; it reports no measured reader result and names no exact shorter international competitor. The theorem’s components are ideals in the overring. Its condition that at least one is properly primary or a prime ideal of the second kind must remain distinct from their pairwise coprimality. The word component cannot itself replace proof of that decomposition. Nor does primary always mean prime.\par
\noindent{\footnotesize Izvorne zapisy / Original records: HIST-g163-006.}\par
\smallskip\noindent\emph{Sravnjenje s tekstom:} Odlomky raboty 31, §§7-8, čuvajut po-parnu prostotu komponentov i odděleno uslovje, že pojavjaje se prinajmenje jedna vlastno primarna komponenta ili prosty ideal drugogo roda. Svez ili i kvantifikator prinajmenje jedna sut česti tvrdženja.\par
\noindent\emph{Textual comparison:} The passages in Paper 31, §§7-8, retain pairwise coprimality of components and the separate condition that at least one properly primary component or prime ideal of the second kind occurs. The disjunction and at-least-one quantifier are parts of the assertion.\par
\noindent{\footnotesize TEXT-COMP-395; pasaže / passages: 18715, 19071; \texttt{CORPUS\_FORM\_COMPARISONS.json}.}\par
\end{minipage}\par
\par\medskip\Needspace{7\baselineskip}
\noindent\begin{minipage}{\linewidth}
\hypertarget{TERM-1311}{}
\noindent\textbf{paralelno prěnešenje vektora}\quad{\footnotesize\texttt{TERM-1311}}\par
\noindent\emph{Iztočnik / Source:} parallele Vektorübertragung\par
\noindent\emph{Rusky / Russian:} параллельное перенесение вектора\par
\noindent\emph{Ukrajinsky / Ukrainian:} паралельне перенесення вектора\par
\smallskip\noindent\emph{Razlog:} Paralelno prěnešenje vektora daje slovjansko nazvanje tehničnomu ponjatju paralelnogo transporta. Ruska i ukrajinska paralela imajut srodny glagolni korenj; arhiv ipak izrično ostavjaje izbor medžu prěnešenje i internacionalnoju semjeju transport otvorenym k jezykovoj ocěnki. V rabotě 21 nazvanje povězano jest s uslovjem $p(d,\delta)=0$ i s geometričnym tolkovanjem Levi-Civity. Ne ide o proizvoljno poměščenje vektora ni prosto o pravu liniju v izbranyh koordinatah. Tekst razlikuje variacionnu osnovu i Weylovo obobčenje, kde odgovarjajuča geodezična uslovja ne izhode več iz istej variacionnoj zadači. Slovnik ne smije činiti vse te konstrukcije identičnymi samo radi odnogo nazvanja.\par
\noindent\emph{Reason:} Paralelno prěnešenje vektora gives a Slavic name to the technical notion of parallel transport. Russian and Ukrainian use a related transfer stem; the archive nevertheless explicitly leaves prěnešenje versus the international transport family open to linguistic assessment. In Paper 21 the name is tied to $p(d,\delta)=0$ and Levi-Civita’s geometric interpretation. It is not arbitrary relocation of a vector or simply a straight line in chosen coordinates. The text distinguishes the variational setting from Weyl’s generalization, where the corresponding geodesic condition no longer follows from the same variational problem. A shared name must not identify all these constructions.\par
\noindent{\footnotesize Izvorne zapisy / Original records: HIST-g084-012.}\par
\smallskip\noindent\emph{Sravnjenje s tekstom:} Rabota 21 povězuje prěnešenje s uslovjem $p(d,\delta)=0$. Vyvod trěbuje diferenciranja prěd eliminacijeju, a Weylovo obobčenje ne daje istu variacionnu osnovu. Nazvanje transporta ne uklanjaje te razlike konstrukcije.\par
\noindent\emph{Textual comparison:} Paper 21 connects transport to $p(d,\delta)=0$. The derivation requires differentiation before elimination, and Weyl’s generalization does not supply the same variational foundation. The transport label does not remove these differences of construction.\par
\noindent{\footnotesize TEXT-COMP-397; pasaže / passages: 14147, 14149, 14151, 15205; \texttt{CORPUS\_FORM\_COMPARISONS.json}.}\par
\end{minipage}\par
\par\medskip\Needspace{7\baselineskip}
\noindent\begin{minipage}{\linewidth}
\hypertarget{TERM-1312}{}
\noindent\textbf{parametrično predstavjenje}\quad{\footnotesize\texttt{TERM-1312}}\par
\noindent\emph{Iztočnik / Source:} Parameterdarstellung\par
\noindent\emph{Rusky / Russian:} параметрическое представление\par
\noindent\emph{Ukrajinsky / Ukrainian:} параметричне подання\par
\smallskip\noindent\emph{Razlog:} Parametrično predstavjenje označaje racionalnu konstrukciju jednačenij pomoču parametrov, ne prilagodženje čislovyh koeficientov k danym. To jest izričny arhivny razlog. V rabotě 5 konstrukcija zavisi od byta minimalnoj bazy pripadajučeho invariantnogo polja i od koeficientnoj oblasti. Najprostejši priměr jest obče jednačenje s neoprěděljenymi koeficientami. Zato fraza ne obeščaje bezuslovno rěšenje za vsaku grupu. Rusko predstavlenie i ukrajinsko podannja pokazyvajut različne prirodne nazvanja istej uloge; ne ustanavjajut jediny dopustimy medžuslovjansky korenj. Dostupny tekst piše prědstavjenje, a rabota 11 koristi pridavnik parametrno; pravopisny i tvorbeny izbori ostavajut oddělene od algebraičnogo smysla.\par
\noindent\emph{Reason:} Parametrično predstavjenje denotes rational construction of equations through parameters, not numerical fitting of coefficients to data. That is the explicit archived motivation. In Paper 5 the construction depends on a minimal basis of the associated invariant field and on the coefficient domain. The simplest example is the general equation with indeterminate coefficients. The phrase therefore promises no unconditional solution for every group. Russian predstavlenie and Ukrainian podannja illustrate different natural names for the same role; they do not establish one uniquely admissible Interslavic stem. The available text spells prědstavjenje, while Paper 11 uses the adjective parametrno; orthographic and derivational choices remain separate from the algebraic sense.\par
\noindent{\footnotesize Izvorne zapisy / Original records: HIST-g039-012.}\par
\smallskip\noindent\emph{Sravnjenje s tekstom:} Raboty 5 i 11 povězujut parametre s minimalnoju bazoju invariantnogo polja i s prědstavjenjem jednačenij. Množstvo primerov, polna familija i singularne izključky ne sut isto. Drugo popravjenje raboty 16 dalje ograničaje osnovnu oblast; slovnično objasnjenje mora ukazati tu granicu bez molčalivoj proměny ranějšego izdanja.\par
\noindent\emph{Textual comparison:} Papers 5 and 11 connect parameters with a minimal basis of the invariant field and a presentation of equations. Many examples, a whole family and singular exceptions are different claims. Paper 16’s second correction further restricts the base domain; the dictionary must explain that boundary without silently changing the earlier edition.\par
\noindent{\footnotesize TEXT-COMP-398; pasaže / passages: 4956, 4958, 4960, 4962, 4964, 8726, 8728, 8730, 8734, 8737, 8739, 8742, 8744, 8747, 8749, 8751, 8754, 8756, 10898; \texttt{CORPUS\_FORM\_COMPARISONS.json}.}\par
\end{minipage}\par
\par\medskip\Needspace{7\baselineskip}
\noindent\begin{minipage}{\linewidth}
\hypertarget{TERM-1313}{}
\noindent\textbf{parametrno predstavjenje}\quad{\footnotesize\texttt{TERM-1313}}\par
\noindent\emph{Iztočnik / Source:} Parameterdarstellung\par
\noindent\emph{Rusky / Russian:} параметрическое представление\par
\noindent\emph{Ukrajinsky / Ukrainian:} параметричне подання\par
\smallskip\noindent\emph{Razlog:} Parametrno predstavjenje ma izražati koeficienty ili korenje kako funkcije parametrov, ne samo davati jim nove oznaky. Arhivny razlog čuva tu konstrukcijnu ulogu i paralelnost prěvodov. Rabota 11 razlikuje obstojanje mnogo jednačenij od obhvačenja cěloj zadanoj familije; ona takože ograničaje prědstavjenje radi singularnyh cěnnostij. Pozdějše avtorovo popravjenje v rabotě 16 izrično sužuje uslovja na osnovnu oblast pri Hilbertovom teoremu: konečno algebraično čislovo polje, s navedenym razširjenjem obsega uvoda. V uvodu dopusčaje takože adjunkciju konečno mnogyh parametrov i medžupolja takovyh razširjenij. To ograničenje ne smije izgubiti se v slovničnom objasnjenju. Izbor parametrno proti parametrično jest tvorbena tema; on ne rěšaje ni polnotu konstrukcije, ni potrěbne uslovja na polje.\par
\noindent\emph{Reason:} Parametrno predstavjenje expresses coefficients or roots as functions of parameters, rather than merely assigning them new symbols. The recorded reason preserves that constructive role and parallel wording across translations. Paper 11 distinguishes existence of many equations from coverage of the whole specified family and also restricts the presentation because of singular values. The later authorial correction in Paper 16 expressly narrows the base-domain conditions for Hilbert’s theorem to a finite algebraic number field. For the introduction, it also permits adjoining finitely many parameters and taking intermediate fields of such extensions. This restriction must not disappear from the dictionary’s explanation. Parametrno versus parametrično is a derivational choice; it settles neither completeness of the construction nor the necessary field hypotheses.\par
\noindent{\footnotesize Izvorne zapisy / Original records: HIST-g073-007.}\par
\smallskip\noindent\emph{Sravnjenje s tekstom:} Raboty 5 i 11 povězujut parametre s minimalnoju bazoju invariantnogo polja i s prědstavjenjem jednačenij. Množstvo primerov, polna familija i singularne izključky ne sut isto. Drugo popravjenje raboty 16 dalje ograničaje osnovnu oblast; slovnično objasnjenje mora ukazati tu granicu bez molčalivoj proměny ranějšego izdanja.\par
\noindent\emph{Textual comparison:} Papers 5 and 11 connect parameters with a minimal basis of the invariant field and a presentation of equations. Many examples, a whole family and singular exceptions are different claims. Paper 16’s second correction further restricts the base domain; the dictionary must explain that boundary without silently changing the earlier edition.\par
\noindent{\footnotesize TEXT-COMP-398; pasaže / passages: 4956, 4958, 4960, 4962, 4964, 8726, 8728, 8730, 8734, 8737, 8739, 8742, 8744, 8747, 8749, 8751, 8754, 8756, 10898; \texttt{CORPUS\_FORM\_COMPARISONS.json}.}\par
\end{minipage}\par
\par\medskip\Needspace{7\baselineskip}
\noindent\begin{minipage}{\linewidth}
\hypertarget{TERM-1314}{}
\noindent\textbf{Peirceovo razloženje}\quad{\footnotesize\texttt{TERM-1314}}\par
\noindent\emph{Iztočnik / Source:} Peircesche Zerlegung\par
\noindent\emph{Rusky / Russian:} разложение Пирса\par
\noindent\emph{Ukrajinsky / Ukrainian:} розклад Пірса\par
\smallskip\noindent\emph{Razlog:} Peirceovo razloženje čuva ime pri nazvanom algebraičnom razloženju. Arhiv izbiraje prisvojny pridavnik vměsto transliterovanej imennej frazy; v ruskom i ukrajinskom zapisu ime stoji v genitivu. V rabotě 34, §13, ide o jednostrannom razloženju $r=cr+(r-cr)$, kde $c^2=c$. Jedna čest leži v pravom idealu porodženom $c$, a drugu $c$ anuliraje. To ne jest utvrždenje, že oba sumandy sut dvustranne idealy; ni ne potrěbuje molčalivo napisati $1-c$ v prstenu, ktoremu jedinica ješče ne jest dokazana. Slovnik mora čuvati jednostrannost i ulogu idempotenta, ne samo eponim.\par
\noindent\emph{Reason:} Peirceovo razloženje preserves the name of the algebraic decomposition. The archive chooses a possessive adjective instead of a transliterated noun phrase; Russian and Ukrainian use the name in a genitive construction. In Paper 34, §13, the passage concerns the one-sided decomposition $r=cr+(r-cr)$ with $c^2=c$. One part lies in the right ideal generated by $c$, and the other is annihilated by $c$. This does not assert that both summands are two-sided ideals, nor does it license silently writing $1-c$ in a ring whose identity has not yet been established. The dictionary must preserve one-sidedness and the idempotent’s role, not just the eponym.\par
\noindent{\footnotesize Izvorne zapisy / Original records: HIST-g191-005.}\par
\smallskip\noindent\emph{Sravnjenje s tekstom:} V rabotě 34, §13, idempotent $c$ dava razloženje $r=cr+(r-cr)$ na dva prave idealy. Drugu čest anuliraje $c$, a suma jest direktna. Tekst ne trěbuje uvođenja $1-c$ prěd dokazom byta jedinice i ne daje dvustrannost oběma sumandam.\par
\noindent\emph{Textual comparison:} In Paper 34, §13, the idempotent $c$ gives $r=cr+(r-cr)$ as a decomposition into two right ideals. The second part is annihilated by $c$, and the sum is direct. The text neither requires introducing $1-c$ before existence of an identity is established nor makes both summands two-sided.\par
\noindent{\footnotesize TEXT-COMP-399; pasaže / passages: 20926, 20928, 20930, 20932, 20934, 20935, 20936, 20937, 20939, 20941, 20943, 20945, 20947, 20949; \texttt{CORPUS\_FORM\_COMPARISONS.json}.}\par
\end{minipage}\par
\par\medskip\Needspace{7\baselineskip}
\noindent\begin{minipage}{\linewidth}
\hypertarget{TERM-1315}{}
\noindent\textbf{proces Poissonovyh skobok}\quad{\footnotesize\texttt{TERM-1315}}\par
\noindent\emph{Iztočnik / Source:} Poissonscher Klammerprozeß\par
\noindent\emph{Rusky / Russian:} процесс скобок Пуассона\par
\noindent\emph{Ukrajinsky / Ukrainian:} процес дужок Пуассона\par
\smallskip\noindent\emph{Razlog:} Proces Poissonovyh skobok čuva razpoznavajemo nazvanje tehničnogo postupka. Arhivny razlog jest dostupnost pojma dnešnjim čitateljam, s ostavljenjem německogo bibliografičnogo naslova. V rabotě 4, §3, fraza nahodi se v primětkě: njeno upotrěbjenje pri Wellsteinu povězuje dostatočne uslovja Capellija s potrěbnymi i dostatočnymi uslovjami za diferencialne procesy. To jest konkretna matematična svez, ne ime tipografskoj skobky. Ruska forma upotrěbja skobki, ukrajinska dužky; zapis ne dava čislovo sravnjenje tih korenjev. Sama primětka ne zadava novu obču formulu Poissonovej skobky i ne dokazuje, že citovana rabota byla znova prověrena pri sadašnjoj rekonstrukciji.\par
\noindent\emph{Reason:} Proces Poissonovyh skobok retains a recognizable name for the technical procedure. The recorded reason is accessibility to modern readers while retaining the German bibliographic title. In Paper 4, §3, the phrase occurs in a footnote: its use in Wellstein links Capelli’s sufficient conditions to necessary and sufficient conditions for differential processes. This is a particular mathematical connection, not a name for a typographic bracket. Russian uses skobki and Ukrainian dužky; the record supplies no numerical comparison of those stems. The footnote neither defines a new general Poisson-bracket formula nor establishes that the cited work was freshly checked during this reconstruction.\par
\noindent{\footnotesize Izvorne zapisy / Original records: HIST-g032-009.}\par
\smallskip\noindent\emph{Sravnjenje s tekstom:} Primětka raboty 4, §3, imenuje proces pri prěhodu od dostatočnyh k potrěbnym i dostatočnym uslovjam. Německy naslov citovanej raboty ostaje neproměnjeny. To jest navedena priměna, ne dokaz novej prověrky citovanogo izvoda.\par
\noindent\emph{Textual comparison:} A footnote in Paper 4, §3, names the process when passing from sufficient to necessary and sufficient conditions. The cited German title remains unchanged. This is a reported application, not evidence of a fresh check of the cited derivation.\par
\noindent{\footnotesize TEXT-COMP-400; pasaže / passages: 3987, 3989, 3990, 3991; \texttt{CORPUS\_FORM\_COMPARISONS.json}.}\par
\end{minipage}\par
\par\medskip\Needspace{7\baselineskip}
\noindent\begin{minipage}{\linewidth}
\hypertarget{TERM-1316}{}
\noindent\textbf{polary}\quad{\footnotesize\texttt{TERM-1316}}\par
\noindent\emph{Iztočnik / Source:} Polaren\par
\noindent\emph{Rusky / Russian:} поляры\par
\noindent\emph{Ukrajinsky / Ukrainian:} поляри\par
\smallskip\noindent\emph{Razlog:} Polary čuva kratky internacionalny termin iz teorije invariantov. Tri zapisy opirajut izbor na ruske i ukrajinske paralele i na kontinuitet tehničnogo pojma; zapis raboty 12 ipak ostavjaje oblik množiny otvorenym k ocěnki, bez točnogo imenovanja alternative. V rabotě 2 polary sut formy, ktore nastavajut iz danej formy pri ukazanyh operacijah i redah prěměnnyh; rabota 4 koristi takve formy v razvitjah. V rabotě 12 one nastavajut iz razvitja $f(x,dx+\lambda\delta x)$. Zajedničny korenj ne čini vse te formule identičnymi. Polara jest rezultat operacije, ne sama operacija; ni naziv ne govori o polarnyh koordinatah ili fizičnoj polarizaciji.\par
\noindent\emph{Reason:} Polary retains the compact international invariant-theory term. Three records support it through Russian and Ukrainian parallels and technical continuity; the Paper 12 record nevertheless leaves the plural form open to assessment without naming a precise alternative. In Paper 2, polars are forms obtained from a given form through specified operations and variable rows; Paper 4 uses such forms in expansions. In Paper 12 they arise from expanding $f(x,dx+\lambda\delta x)$. A shared stem does not identify all these formulas. A polar is an operation’s result, not the operation itself; the name says nothing about polar coordinates or physical polarization.\par
\noindent{\footnotesize Izvorne zapisy / Original records: HIST-g004-002, HIST-g033-006, HIST-g074-009.}\par
\smallskip\noindent\emph{Sravnjenje s tekstom:} Raboty 2 i 4 povězujut polary s ukazanymi redami prěměnnyh i polarnymi procesami. Rabota 12 daje osobnu konstrukciju iz razvitja diferencialnogo izraza. Različenje medžu rezultatom i operacijeju jest obče; konkretne formule i koeficienty ostavajut zavisne od konteksta.\par
\noindent\emph{Textual comparison:} Papers 2 and 4 relate polars to specified variable rows and polar processes. Paper 12 gives a particular construction by expanding a differential expression. The distinction between result and operation is shared; the actual formulas and coefficients remain context-dependent.\par
\noindent{\footnotesize TEXT-COMP-401; pasaže / passages: 398, 400, 402, 404, 405, 406, 408, 411, 412, 413, 414, 415, 418, 420, 422, 443, 445, 446, 447, 448, 450, 452, 453, 454, 455, 4202, 4204, 4205, 4206, 4207, 4209, 8985, 8986, 8989, 8990, 8991, 8993, 8994, 8995, 8996; \texttt{CORPUS\_FORM\_COMPARISONS.json}.}\par
\end{minipage}\par
\par\medskip\Needspace{7\baselineskip}
\noindent\begin{minipage}{\linewidth}
\hypertarget{TERM-1317}{}
\noindent\textbf{polary / tvorjenje polarov / polarny procesy}\quad{\footnotesize\texttt{TERM-1317}}\par
\noindent\emph{Iztočnik / Source:} Polaren / Polarenbildung / Polarprozesse\par
\noindent\emph{Rusky / Russian:} поляры / образование поляр / полярные процессы\par
\noindent\emph{Ukrajinsky / Ukrainian:} поляри / утворення поляр / полярні процеси\par
\smallskip\noindent\emph{Razlog:} Polary, tvorjenje polarov i polarny procesy tvoret svęzanu terminologičnu semju, ale imajut različne gramatične i matematične ulogi. Arhiv čuva tehničny internacionalny korenj vměsto domačej kalki; konkretna odvrgnuta kalka ne jest navedena. V rabotě 16 polary sut koeficientne formy razvitja, tvorjenje označaje ih dobivanje, a procesy sut operacije. Tekst razlikuje te operacije od $\Omega$-procesa i od množenja determinantem, hoti v konstrukciji nastupajut vměstě. Pozdějše avtorovo popravjenje ograničaje jednu ranějše upotrěbljenu identičnost na binarny slučaj. Zato ponovjenje istej terminologije ne smije prěnositi nepravilnu obču formulu. Arhivno polarny procesy takože potrěbuje oddělnu prověrku soglasovanja pridavnika s množinoju.\par
\noindent\emph{Reason:} Polary, tvorjenje polarov and polarny procesy form a related terminological family but have different grammatical and mathematical roles. The archive retains the technical international stem rather than a native calque; it records no exact rejected calque. In Paper 16, polars are coefficient forms in an expansion, formation names their production, and processes are operations. The text distinguishes those operations from the $\Omega$-process and multiplication by a determinant, although they occur together in the construction. A later authorial correction restricts one previously used identity to the binary case. Terminological continuity must not carry an invalid general formula with it. Archival polarny procesy also requires a separate adjective-plural agreement check.\par
\noindent{\footnotesize Izvorne zapisy / Original records: HIST-g078-004.}\par
\smallskip\noindent\emph{Sravnjenje s tekstom:} Rabota 16 razlikuje polary koeficientnyh form, tvorjenje polarov i operacije s determinantom i $\Omega$. Avtorovo popravjenje izrično ograničaje jednu identičnost na binarnu oblast. Jednaky slovny korenj ne opravdava prěnos formule na šire uslovja.\par
\noindent\emph{Textual comparison:} Paper 16 distinguishes polars of coefficient forms, formation of polars, and operations involving the determinant and $\Omega$. The authorial correction explicitly restricts one identity to the binary setting. A common word stem does not justify extending a formula beyond its hypotheses.\par
\noindent{\footnotesize TEXT-COMP-402; pasaže / passages: 10729, 10731, 10733, 10749, 10757, 10771, 10773, 10775, 10777; \texttt{CORPUS\_FORM\_COMPARISONS.json}.}\par
\end{minipage}\par
\par\medskip\Needspace{7\baselineskip}
\noindent\begin{minipage}{\linewidth}
\hypertarget{TERM-1318}{}
\noindent\textbf{polary / polarne operacije / polarny proces}\quad{\footnotesize\texttt{TERM-1318}}\par
\noindent\emph{Iztočnik / Source:} Polaren / Polaroperationen / Polarprozeß\par
\noindent\emph{Rusky / Russian:} поляры / полярные операции / полярный процесс\par
\noindent\emph{Ukrajinsky / Ukrainian:} поляри / полярні операції / полярний процес\par
\smallskip\noindent\emph{Razlog:} Polary, polarne operacije i polarny proces naměrno razlikujut formy od operacij nad formami. To jest izričny razlog zapisa raboty 4, §7. Tam $\Pi$ označaje agregat polarnyh operacij, v kotoryh se diferencira po jednom redu prěměnnyh i kombinira s drugim. Rezultat priměnjenja jest forma, ne novy sinonim za operator. Ruska i ukrajinska paralela čuvajut istu razliku medžu imenami objektov i dějstvij. Pri prěnosu termina potrěbno jest zachovati argumenty, směru operacije i ukazane indekse; ne dostaje samo podoba korenja polar-. Arhiv ne dava osobne čislove težiny za operacija proti proces.\par
\noindent\emph{Reason:} Polary, polarne operacije and polarny proces deliberately distinguish forms from operations on forms. That is the explicit reason recorded for Paper 4, §7. There $\Pi$ denotes a combination of polar operations that differentiate with respect to one variable row and combine with another. Applying it produces a form, not a new synonym for the operator. Russian and Ukrainian preserve the same distinction between objects and actions. Transferring the term requires preservation of arguments, direction and specified indices; resemblance of the polar- stem is insufficient. The archive gives no separate numerical weights for operacija versus proces.\par
\noindent{\footnotesize Izvorne zapisy / Original records: HIST-g036-006.}\par
\smallskip\noindent\emph{Sravnjenje s tekstom:} V rabotě 4, §7, $\Pi$ jest agregat operacij s različnymi redami prěměnnyh; formy v razloženju sut jih rezultaty. Indekse, redy i granice sumovanja trěba čuvati vměstě s nazvanjem polarny.\par
\noindent\emph{Textual comparison:} In Paper 4, §7, $\Pi$ is a combination of operations on different variable rows; the forms in the expansion are their results. Indices, rows and summation bounds must be preserved along with the label polar.\par
\noindent{\footnotesize TEXT-COMP-403; pasaže / passages: 4531, 4548, 4550, 4551, 4554, 4556, 4557, 4560, 4562, 4565; \texttt{CORPUS\_FORM\_COMPARISONS.json}.}\par
\end{minipage}\par
\par\medskip\Needspace{7\baselineskip}
\noindent\begin{minipage}{\linewidth}
\hypertarget{TERM-1319}{}
\noindent\textbf{polarizacija $x$ po $\nu_1$}\quad{\footnotesize\texttt{TERM-1319}}\par
\noindent\emph{Iztočnik / Source:} Polarisation von $x$ nach $\nu_1$\par
\noindent\emph{Rusky / Russian:} поляризация $x$ по $\nu_1$\par
\noindent\emph{Ukrajinsky / Ukrainian:} поляризація $x$ за $\nu_1$\par
\smallskip\noindent\emph{Razlog:} Polarizacija $x$ po $\nu_1$ čuva internacionalny termin i eksplicitny odnos medžu dvoma argumentami. Arhivni rusky po i ukrajinsky za pokazyvajut različne prirodne prědložne konstrukcije; medžuslovjansky zapis izbiraje po. V rabotě 2, §17, ta instrukcija vodi od ukazanej identičnosti k formuli (21). Zato $x$ i $\nu_1$ ne smut byti zaměnjene, a dolny indeks jest čest matematičnoj oznaky, ne padežna koncovka. Slovny vyraz sam ne zadava normalizaciju ni čislove koeficienty polarnogo procesa: ty trěba zachovati iz konkretnyh formul. Sravnjenje tyh dvoh odlomkov čuva směr operacije; ne jest novy nezavisny dokaz vsej redukcije.\par
\noindent\emph{Reason:} Polarizacija $x$ po $\nu_1$ retains the international term and an explicit relation between its two arguments. Recorded Russian po and Ukrainian za show different natural prepositional constructions; Interslavic selects po. In Paper 2, §17, this instruction leads from the displayed identity to formula (21). Thus $x$ and $\nu_1$ must not be interchanged, and the subscript is part of a mathematical identifier, not an inflectional ending. The wording alone defines neither normalization nor numerical coefficients of the polar process: these must be preserved from the actual formulas. Comparing the two passages preserves operation direction; it is not a new independent proof of the whole reduction.\par
\noindent{\footnotesize Izvorne zapisy / Original records: HIST-g018-006.}\par
\smallskip\noindent\emph{Sravnjenje s tekstom:} Rabota 2, §17, izrično naznačaje polarizaciju $x$ po $\nu_1$ medžu počatnoju identičnostju i formulou (21). Sravnjenje čuva red argumentov i indeks; same slova ne ustanavjajut normalizaciju operacije ni prověrku vseh koeficientov.\par
\noindent\emph{Textual comparison:} Paper 2, §17, explicitly specifies polarization of $x$ with respect to $\nu_1$ between the starting identity and formula (21). The comparison preserves argument order and subscript; the words alone establish neither normalization of the operation nor verification of every coefficient.\par
\noindent{\footnotesize TEXT-COMP-404; pasaže / passages: 2399, 2401, 2402, 2403, 2404, 2405, 2407, 2408, 2409, 2410, 2411, 2412, 2413, 2414; \texttt{CORPUS\_FORM\_COMPARISONS.json}.}\par
\end{minipage}\par
\par\medskip\Needspace{7\baselineskip}
\noindent\begin{minipage}{\linewidth}
\hypertarget{TERM-1320}{}
\noindent\textbf{polarizacija / polarny členy}\quad{\footnotesize\texttt{TERM-1320}}\par
\noindent\emph{Iztočnik / Source:} Polarisierung / Polarenglieder\par
\noindent\emph{Rusky / Russian:} поляризация / полярные члены\par
\noindent\emph{Ukrajinsky / Ukrainian:} поляризація / полярні члени\par
\smallskip\noindent\emph{Razlog:} Polarizacija i polarny členy prodolžajut terminologiju polarov iz §2. Arhivny razlog jest kratka podoba s ruskoju i ukrajinskoju matematičnoju formou. V §5 polarizacija jest postupok, a polarny členy sut izrazy razvivane kako linearne kombinacije s danym simboličnym faktorom. Člen znači član formule, ne element proizvoljnoj grupy i ne korenj jednačenja. Različne kombinacije prěměnnyh ostavajut važne pri povězivanju polarov. Staru formu polarny členy trěba odděliti od nyněšnjego padeža polarnyh členov i od občej prověrky soglasovanja; ona ne ustanavjaje posebnogo matematičnogo smysla. Izbor internacionalnogo korenja ne razrěšaje izostaviti operaciju ili njene rezultate.\par
\noindent\emph{Reason:} Polarizacija and polarny členy continue the vocabulary of polars from §2. The archived reason is concise resemblance to Russian and Ukrainian mathematical forms. In §5, polarization is the procedure, while polar terms are expressions expanded as linear combinations with a specified symbolic factor. Člen denotes a formula term, not an arbitrary group element or an equation’s root. Different combinations of variables remain important when relating the polars. Archival polarny členy must be distinguished from the current inflected polarnyh členov and from a general agreement check; it establishes no separate mathematical sense. Choosing the international stem does not permit omission of either the operation or its results.\par
\noindent{\footnotesize Izvorne zapisy / Original records: HIST-g006-004.}\par
\smallskip\noindent\emph{Sravnjenje s tekstom:} V rabotě 2, §5, polarne členy razvijajut se v linearne kombinacije, a različne kombinacije prěměnnyh daju odnose medžu polarami. Postupok, izrazy i njene prěměnne ne smut byti slite v jedno neoprěděljeno nazvanje.\par
\noindent\emph{Textual comparison:} In Paper 2, §5, polar terms are expanded as linear combinations, and different variable combinations give relations among polars. The procedure, expressions and their variables must not collapse into one unspecified label.\par
\noindent{\footnotesize TEXT-COMP-405; pasaže / passages: 711, 713, 715, 716, 717, 719; \texttt{CORPUS\_FORM\_COMPARISONS.json}.}\par
\end{minipage}\par
\par\medskip\Needspace{7\baselineskip}
\noindent\begin{minipage}{\linewidth}
\hypertarget{TERM-1321}{}
\noindent\textbf{polarne procesy / polary}\quad{\footnotesize\texttt{TERM-1321}}\par
\noindent\emph{Iztočnik / Source:} Polarprozesse / Polare\par
\noindent\emph{Rusky / Russian:} полярные процессы / поляры\par
\noindent\emph{Ukrajinsky / Ukrainian:} полярні процеси / поляри\par
\smallskip\noindent\emph{Razlog:} Polarne procesy / polary čuvajut različje medžu operacijeju i formami, ktore ona davaje. Arhiv izrično odbira kratko polary vměsto razširjenogo nazvanja polarne formy, zato že tekst upotrěbljaje to kako nazvanje rezultata. Ruska i ukrajinska paralela čuvajut isti internacionalny korenj. V rabotě 8 proces $P_{hk}$ jest zapisan kako diferencialny operator; produkt procesov znači poslědovateljno priměnjenje, s pravym operatorom prvim. Redukcijny teorem izražaje formu kako sumu polarov. Zato proces i polara ne sut zaměnjajeme slova, ni to jest terminologija polarnyh koordinat. Kratkost jest zapisany razlog, ne dokaz izměrenoj razumlivosti.\par
\noindent\emph{Reason:} Polarne procesy / polary distinguishes an operation from the forms it produces. The record expressly chooses the short polary instead of expanding it to polarne formy, because the text names the outputs directly. Russian and Ukrainian retain the same international stem. In Paper 8, $P_{hk}$ is displayed as a differential operator; a product of processes means successive application, with the right-hand operator acting first. The reduction theorem expresses a form as a sum of polars. Process and polar are therefore not interchangeable, nor is this polar-coordinate terminology. Brevity is the recorded reason, not evidence of measured intelligibility.\par
\noindent{\footnotesize Izvorne zapisy / Original records: HIST-g057-003.}\par
\smallskip\noindent\emph{Sravnjenje s tekstom:} Rabota 8 definira diferencialny proces i izrično oprědělja porjadok njegovogo složenja. Redukcija dava sumu polarov; proizvodeča operacija i proizvedena forma morajut ostati različene.\par
\noindent\emph{Textual comparison:} Paper 8 defines a differential process and expressly specifies composition order. Reduction yields a sum of polars; the producing operation and resulting form must remain distinct.\par
\noindent{\footnotesize TEXT-COMP-406; pasaže / passages: 6751, 6753, 6760, 6764, 6766, 6768, 6770, 6771, 6772, 6773, 6774, 6776, 6778, 6780, 6783, 6785, 6787, 6803; \texttt{CORPUS\_FORM\_COMPARISONS.json}.}\par
\end{minipage}\par
\par\medskip\Needspace{7\baselineskip}
\noindent\begin{minipage}{\linewidth}
\hypertarget{TERM-1322}{}
\noindent\textbf{polarne procesy / polarne operacije}\quad{\footnotesize\texttt{TERM-1322}}\par
\noindent\emph{Iztočnik / Source:} Polarprozesse / Polaroperationen\par
\noindent\emph{Rusky / Russian:} полярные процессы / полярные операции\par
\noindent\emph{Ukrajinsky / Ukrainian:} полярні процеси / полярні операції\par
\smallskip\noindent\emph{Razlog:} Polarne procesy / polarne operacije čuva dejstvenu ulogu v konstrukciji formuly (31) raboty 4. Zapis navodi soglasnost s modernym matematičnym registrom, ale ne imenuje osobnu novu čitateljsku prověrku. Ruska i ukrajinska paralela imajut te iste internacionalne nazvanja. Dostupny odlomok najprvo dava pojedine diferencialne procesy, potom označaje s $\Pi$ agregat operacij s konstantnymi množiteljami. Polara jest rezultat, a $\Pi$ ne jest proizvoljna jedina derivacija. Indeksy, porjadok argumentov i dolna i gornja granica sumovanja ostavajut čestju konstrukcije; obče slovo operacija jih ne može zaměniti.\par
\noindent\emph{Reason:} Polarne procesy / polarne operacije retains the operational role in the construction of formula (31) in Paper 4. The record cites modern mathematical register, but no new reader assessment. Russian and Ukrainian use the same international names. The passage first displays individual differential processes, then denotes an aggregate of operations with constant multipliers by $\Pi$. A polar is an output, while $\Pi$ is not an arbitrary single derivative. Indices, argument order and both summation bounds remain part of the construction; the general word operation cannot replace them.\par
\noindent{\footnotesize Izvorne zapisy / Original records: HIST-g033-007.}\par
\smallskip\noindent\emph{Sravnjenje s tekstom:} V rabotě 4 formula (31) upotrěbljaje agregat polarnyh operacij s konstantnymi množiteljami. Točno ukazane argumenty i granice sumy ne slědujut samo iz nazvanja proces.\par
\noindent\emph{Textual comparison:} In Paper 4, formula (31) uses an aggregate of polar operations with constant multipliers. Its specified arguments and summation bounds do not follow from the word process alone.\par
\noindent{\footnotesize TEXT-COMP-407; pasaže / passages: 4202, 4204, 4205, 4206, 4207, 4209, 4212, 4213, 4214, 4215, 4217; \texttt{CORPUS\_FORM\_COMPARISONS.json}.}\par
\end{minipage}\par
\par\medskip\Needspace{7\baselineskip}
\noindent\begin{minipage}{\linewidth}
\hypertarget{TERM-1323}{}
\noindent\textbf{poligon (odnosno divizor) / klas poligonov}\quad{\footnotesize\texttt{TERM-1323}}\par
\noindent\emph{Iztočnik / Source:} Polygon bzw. Divisor / Polygonklasse\par
\noindent\emph{Rusky / Russian:} полигон (соответственно дивизор) / класс полигонов\par
\noindent\emph{Ukrajinsky / Ukrainian:} полігон (відповідно дивізор) / клас полігонів\par
\smallskip\noindent\emph{Razlog:} Poligon (odnosno divizor) / klas poligonov čuva historično nazvanje, prědano v izloženju Dedekindove i Weberove teorije. Izrična alternativa jest povsudna modernizacija na divizor; zapis ju ne izbira, ale dodava divizor tam, kde to děla iztočnik. Ruska i ukrajinska paralela čuvajut istu dvojicu. V rabotě 14 poligon ne jest mnogougoljnik v ravnině: odlomok govori o točkah Riemannovoj površiny, nulah i beskonečnyh městah funkcije. Klasy idealov i klasy poligonov ne sut prosto identične; tekst potrebuje jednak stepen za odgovarjajuču ekvivalentnost i objašnjava ulogu imenovateljnogo poligona. Slovnik mora čuvati tu razliku.\par
\noindent\emph{Reason:} Poligon (odnosno divizor) / klas poligonov preserves the historical name in the account of Dedekind-Weber theory. The explicit alternative is uniform modernization to divizor; the record rejects that policy and adds divisor where the source does. Russian and Ukrainian preserve the same pair. In Paper 14, polygon does not mean a plane polygon: the passage concerns points of a Riemann surface, zeros and infinite places of a function. Ideal classes and polygon classes are not simply identical; the text requires equal degree for the corresponding equivalence and explains the denominator polygon's role. The dictionary must retain that distinction.\par
\noindent{\footnotesize Izvorne zapisy / Original records: HIST-g076-006.}\par
\smallskip\noindent\emph{Sravnjenje s tekstom:} Rabota 14 čuva historično poligon i objašnjava jego odnos k divizoru. Prěhod medžu klasami idealov i poligonov zavisi od stepena i imenovateljnyh toček; ne jest prosta zaměna imen.\par
\noindent\emph{Textual comparison:} Paper 14 retains historical polygon and explains its relation to divisor. Passage between ideal and polygon classes depends on degree and denominator points; it is not a mere renaming.\par
\noindent{\footnotesize TEXT-COMP-408; pasaže / passages: 10022, 10024, 10026, 10028; \texttt{CORPUS\_FORM\_COMPARISONS.json}.}\par
\end{minipage}\par
\par\medskip\Needspace{7\baselineskip}
\noindent\begin{minipage}{\linewidth}
\hypertarget{TERM-1324}{}
\noindent\textbf{polinomialno kolco}\quad{\footnotesize\texttt{TERM-1324}}\par
\noindent\emph{Iztočnik / Source:} Polynombereich\par
\noindent\emph{Rusky / Russian:} кольцо многочленов\par
\noindent\emph{Ukrajinsky / Ukrainian:} кільце многочленів\par
\smallskip\noindent\emph{Razlog:} Polinomialno kolco jest arhivna forma iz raboty 34, §22. Zapis bira internacionalny polinom- vměsto nacionalno oblikovanogo mnogochlen- i tolkuje Bereich kako prsten zaradi konkretnoj algebraičnoj struktury. Ruska i ukrajinska paralela čuvajut imenno nazvanje prstena mnogočlenov. Dostupny kumulativny tekst už piše polinomialny prsten: to jest važna razlika medžu arhivnym kandidatom i nyněšnju formoju. Mapa $x_i\mapsto a_i$ dava grupovo prsten kako homomorfny obraz polinomnogo prstena, ne kako sam polinomny prsten bez relacij. To opravdava strukturalno čitanje nazvanja, ale ne obče pravilo zaměniti vsako německo Bereich slovom prsten.\par
\noindent\emph{Reason:} Polinomialno kolco is the archived form from Paper 34, §22. The record chooses international polinom- over a nationally shaped mnogochlen- form and interprets Bereich as ring because of the particular algebraic structure. Russian and Ukrainian explicitly name a polynomial ring. The cumulative text now uses polinomialny prsten: the historical candidate must be distinguished from the current form. The map $x_i\mapsto a_i$ makes the group ring a homomorphic image of the polynomial ring, not the unrestricted polynomial ring itself. This supports a structural interpretation here, not a rule replacing every German Bereich by ring.\par
\noindent{\footnotesize Izvorne zapisy / Original records: HIST-g200-005.}\par
\smallskip\noindent\emph{Sravnjenje s tekstom:} V rabotě 34 grupovo prsten nastavaje kako homomorfny obraz polinomnogo prstena pri ukazanyh relacijah. Dostupna forma prsten razlikuje se od arhivnogo kolco; nazvanje ne izbrisava relacije.\par
\noindent\emph{Textual comparison:} In Paper 34 the group ring arises as a homomorphic image of a polynomial ring under specified relations. Current prsten differs from archival kolco; the label does not remove those relations.\par
\noindent{\footnotesize TEXT-COMP-409; pasaže / passages: 22004, 22006, 22008, 22010, 22012, 22015, 22017, 22018, 22019, 22021, 22023, 22025, 22027, 22029; \texttt{CORPUS\_FORM\_COMPARISONS.json}.}\par
\end{minipage}\par
\par\medskip\Needspace{7\baselineskip}
\noindent\begin{minipage}{\linewidth}
\hypertarget{TERM-1325}{}
\noindent\textbf{polinomna oblast}\quad{\footnotesize\texttt{TERM-1325}}\par
\noindent\emph{Iztočnik / Source:} Polynombereich\par
\noindent\emph{Rusky / Russian:} полиномиальная область\par
\noindent\emph{Ukrajinsky / Ukrainian:} поліномна область; поліноміальна область\par
\smallskip\noindent\emph{Razlog:} Polinomna oblast upotrěbljaje internacionalny korenj i šire slovo oblast. Tri zapisy davajut kako razlog svez s už upotrěbljanoju terminologijeju idealov i eliminacije; ne davajut čislove ocěny kandidatov. Ruska i ukrajinska paralela takože imajut oblast. V rabotě 19, §10, osnovu tvoret polinomy nad kompleksnymi čislami; rabota 27 govori o Hilbertovyh čislah, a rabota 31 izrično razlikuje $P[x]$ od njegovogo prstena klasov ostatkov po $f(x)$. Jednak korenj ne čini te objekte jednakymi. Takože rabota 19 posebno razgledaje homogene formy, kde složenje ograničeno jest na jednak stepen. Zato oblast ne može bez konteksta obeščati vse operacije obyčnogo prstena.\par
\noindent\emph{Reason:} Polinomna oblast combines the international stem with the broader word oblast. Three records justify it by continuity with existing ideal- and elimination-theory terminology; they provide no numerical candidate scores. Russian and Ukrainian likewise use oblast. Paper 19, §10, starts with polynomials over the complex numbers; Paper 27 discusses Hilbert numbers, while Paper 31 explicitly distinguishes $P[x]$ from its residue-class ring modulo $f(x)$. Sharing a stem does not identify these objects. Paper 19 also separately considers homogeneous forms, where addition is restricted to equal degree. Without context, oblast cannot promise all operations of an ordinary ring.\par
\noindent{\footnotesize Izvorne zapisy / Original records: HIST-g081-078, HIST-g104-004, HIST-g127-007.}\par
\smallskip\noindent\emph{Sravnjenje s tekstom:} Raboty 19, 27 i 31 upotrěbljajut isti korenj v različnyh konstrukcijah. Polinomny prsten, homogena oblast s ograničenym složenjem i prsten klasov ostatkov ne sut jediny objekt.\par
\noindent\emph{Textual comparison:} Papers 19, 27 and 31 use the same stem in different constructions. A polynomial ring, a homogeneous domain with restricted addition, and a residue-class ring are not one object.\par
\noindent{\footnotesize TEXT-COMP-410; pasaže / passages: 13459, 13461, 13462, 13466, 13467, 13468, 13477, 15917, 17702, 17704; \texttt{CORPUS\_FORM\_COMPARISONS.json}.}\par
\end{minipage}\par
\par\medskip\Needspace{7\baselineskip}
\noindent\begin{minipage}{\linewidth}
\hypertarget{TERM-1326}{}
\noindent\textbf{polinomna oblast / polinomy}\quad{\footnotesize\texttt{TERM-1326}}\par
\noindent\emph{Iztočnik / Source:} Polynombereich / Polynome\par
\noindent\emph{Rusky / Russian:} полиномиальная область / многочлены\par
\noindent\emph{Ukrajinsky / Ukrainian:} поліноміальна область / многочлени\par
\smallskip\noindent\emph{Razlog:} Polinomna oblast / polinomy čuva razliku medžu cěloju oblastju $J$ i jej elementami. To jest izričny razlog zapisa; ruska i ukrajinska paralela razlikujut oblast od mnogočlenov. V rabotě 17 polinomy sut konečne sumy uporjadkovanyh produktov s koeficientami. Važno jest, že obča konstrukcija ne prědpoklada komutativnost koeficientov s neizvěstnymi: tekst dava osobne pravila za $\xi_i a$. Komutativnost samyh neizvěstnyh uvodi se oddělno, s ukazanym městom priměnjenja. Prěvod polinom ne smije tiho zaměniti tu oblast običnym komutativnym prstenom nad centralnym poljem. Razlika izmeždu struktury i elementa jest neodvisna od izbora medžu polinom- i mnogočlen-.\par
\noindent\emph{Reason:} Polinomna oblast / polinomy distinguishes the whole domain $J$ from its elements. This is the record's explicit reason; the Russian and Ukrainian comparanda distinguish the domain from its polynomials. In Paper 17, polynomials are finite sums of ordered products with coefficients. Crucially, the general construction does not assume that coefficients commute with the indeterminates: the text gives special rules for $\xi_i a$. Commutativity among the indeterminates is imposed separately, with a stated point of use. The word polynomial must not silently turn this into an ordinary commutative ring over a central field. Structure versus element is a separate question from polinom- versus mnogočlen-.\par
\noindent{\footnotesize Izvorne zapisy / Original records: HIST-g079-011.}\par
\smallskip\noindent\emph{Sravnjenje s tekstom:} Rabota 17 definira polinomy kako konečne sumy i osobno reguliše množenje neizvěstnyh s koeficientami. Pozdnějša komutativnost neizvěstnyh ne dozvalja prědpokladati centralnost vsih koeficientov.\par
\noindent\emph{Textual comparison:} Paper 17 defines polynomials as finite sums and separately regulates multiplication of indeterminates with coefficients. Later commutativity among the indeterminates does not license centrality of all coefficients.\par
\noindent{\footnotesize TEXT-COMP-411; pasaže / passages: 11042, 11044, 11046, 11048, 11050, 11053, 11055, 11057, 11059, 11062; \texttt{CORPUS\_FORM\_COMPARISONS.json}.}\par
\end{minipage}\par
\par\medskip\Needspace{7\baselineskip}
\noindent\begin{minipage}{\linewidth}
\hypertarget{TERM-1327}{}
\noindent\textbf{polinomna oblast iz izčetno mnogih neopreděljenyh}\quad{\footnotesize\texttt{TERM-1327}}\par
\noindent\emph{Iztočnik / Source:} Polynombereich von abzählbar vielen Unbestimmten\par
\noindent\emph{Rusky / Russian:} полиномиальная область от счетно многих неопределенных\par
\noindent\emph{Ukrajinsky / Ukrainian:} поліномна область від зліченно багатьох невизначених\par
\smallskip\noindent\emph{Razlog:} Polinomna oblast iz izčetno mnogih neopreděljenyh čuva izričnu beskonečnost i izčetnost generatorov v Hölzerovom kontrapriměru. Ruska i ukrajinska paralela davajut tu istu količinstvenu ogradu; zapis ne predlaga krači konkurentni termin. V rabotě 30, §5, polinomy imajut koeficienty iz polja, a ideal sodrživaje rastuče stepeni pojedinyh $x_i$ i směšane produkty. Vsaky odgovarjajuči element imaje někoj eksponent, ale ne obstaje jedna stepenj za cěle idealne podsistemy, upotrěbljene v kontrapriměru. Zato konečnost vsakego polinoma ne znači konečno čislo generatorov oblasti. Zaměna izčetno mnogih s prosto mnogimi izgubila by uslovje konstrukcije.\par
\noindent\emph{Reason:} Polinomna oblast iz izčetno mnogih neopreděljenyh preserves the explicit countable infinity of generators in Hölzer's counterexample. Russian and Ukrainian give the same cardinality qualification; the record names no shorter competing term. In Paper 30, §5, coefficients lie in a field and the ideal contains increasing powers of individual $x_i$ together with mixed products. Each relevant element has some exponent, but no single power works for the entire ideal subsystems used in the counterexample. Finiteness of each polynomial therefore does not imply finitely many domain generators. Replacing countably many by merely many would lose a construction hypothesis.\par
\noindent{\footnotesize Izvorne zapisy / Original records: HIST-g112-012.}\par
\smallskip\noindent\emph{Sravnjenje s tekstom:} Hölzerov priměr v rabotě 30 upotrěbljaje beskonečnu poslědovateljnost generatorov i rastuče eksponenty. Konečnost pojedinogo polinoma ne dava jednu obču stepenj za cěly ideal.\par
\noindent\emph{Textual comparison:} Hölzer's example in Paper 30 uses an infinite sequence of generators and increasing exponents. Finiteness of an individual polynomial does not give a uniform power for the entire ideal.\par
\noindent{\footnotesize TEXT-COMP-412; pasaže / passages: 16611, 16613, 16615; \texttt{CORPUS\_FORM\_COMPARISONS.json}.}\par
\end{minipage}\par
\par\medskip\Needspace{7\baselineskip}
\noindent\begin{minipage}{\linewidth}
\hypertarget{TERM-1329}{}
\noindent\textbf{polinomne idealy / idealy iz polinomov}\quad{\footnotesize\texttt{TERM-1329}}\par
\noindent\emph{Iztočnik / Source:} Polynomideale\par
\noindent\emph{Rusky / Russian:} полиномиальные идеалы / идеалы из полиномов\par
\noindent\emph{Ukrajinsky / Ukrainian:} поліномні ідеали / ідеали з поліномів\par
\smallskip\noindent\emph{Razlog:} Polinomne idealy / idealy iz polinomov predstavlja dva načina izraženja istogo tipa objekta. Zapis izbira kratky pridavnik za naslov i razloženu frazu za prozu, kde potrěbno jest razlikovati idealy od modulov linearnyh form. Ruska i ukrajinska paralela dopuščajut tu istu izměnu registra. Rabota 22 ne ostavlja razliky na nazvanju: definira zatvorenost pri razlici i pri množenju libovoljnym polinomom, a primětka izrično objašnjava potrebu nazvanja ideal vměsto staršego modul. Zato slovnik ne smije klasifikovati oba slovosloženja kako dva različne ponjatja, ni smije izbrisati razliku medžu idealom i modulom v danoj konstrukciji.\par
\noindent\emph{Reason:} Polinomne idealy / idealy iz polinomov offers two expressions for the same object type. The record chooses the compact adjective for the title and the expanded phrase in prose where ideals must be distinguished from modules of linear forms. Russian and Ukrainian allow the same register variation. Paper 22 does not leave the distinction to a name: it defines closure under difference and multiplication by an arbitrary polynomial, and a footnote explicitly explains why ideal is needed in place of the older module. The dictionary should neither classify the two phrases as different concepts nor erase the ideal/module distinction in this construction.\par
\noindent{\footnotesize Izvorne zapisy / Original records: HIST-g093-001.}\par
\smallskip\noindent\emph{Sravnjenje s tekstom:} Rabota 22 dava obe nazvanja idealov i izrično objašnjava, začem jih razlikuje od modulov linearnyh form. Zatvorenost pri polinomnom množenju jest čest definicije.\par
\noindent\emph{Textual comparison:} Paper 22 gives both ideal expressions and explicitly explains why they must be distinguished from modules of linear forms. Closure under polynomial multiplication is part of the definition.\par
\noindent{\footnotesize TEXT-COMP-414; pasaže / passages: 14204, 14209, 14220, 14222, 14226, 14228, 14501, 14521; \texttt{CORPUS\_FORM\_COMPARISONS.json}.}\par
\end{minipage}\par
\par\medskip\Needspace{7\baselineskip}
\noindent\begin{minipage}{\linewidth}
\hypertarget{TERM-1330}{}
\noindent\textbf{stepenj svojego asociovanogo prostogo ideala}\quad{\footnotesize\texttt{TERM-1330}}\par
\noindent\emph{Iztočnik / Source:} Potenz seines zugehörigen Primideals\par
\noindent\emph{Rusky / Russian:} степень своего соответствующего простого идеала\par
\noindent\emph{Ukrajinsky / Ukrainian:} степінь свого асоційованого простого ідеалу\par
\smallskip\noindent\emph{Razlog:} Stepenj svojego asociovanogo prostogo ideala čuva silnějše tvrdženje, ne samo obče prědstavjenje prěz primarne komponenty. Zapis upomina imenno prěhod v rabotě 30, §8. Rusky odgovarjajuči i ukrajinsky asociovany označajut odnos k tomu prostomu idealu, ktory prinaleži danomu primarnomu idealu; ne k proizvoljnomu prostomu idealu. Dostupny tekst dodava integralnu zamknjenost k ranějšim aksiomam i izključaje nulovy ideal. Bez tyh uslovij primarny ideal ne mora byti stepenju svojego prostogo ideala. Slovnik mora izraziti jednakost idealov, ne samo děliteljnost někojoj stepeni; povratno svojego čuva taj odnos.\par
\noindent\emph{Reason:} Stepenj svojego asociovanogo prostogo ideala preserves a stronger assertion than general decomposition into primary components. The record specifically identifies the transition in Paper 30, §8. Russian corresponding and Ukrainian associated refer to the prime ideal belonging to the particular primary ideal, not an arbitrary prime. The passage adds integral closedness to the earlier axioms and excludes the zero ideal. Without those assumptions, a primary ideal need not be a power of its prime ideal. The dictionary must express equality of ideals, not merely divisibility involving some power; svojego preserves that association.\par
\noindent{\footnotesize Izvorne zapisy / Original records: HIST-g115-008.}\par
\smallskip\noindent\emph{Sravnjenje s tekstom:} Rabota 30 pridava integralnu zamknjenost i pod vsimi aksiomami dava jednoznačny produkt stepenjev prostyh idealov. To jest silnějše od občego primarnogo razloženja i čuva izključenija nulovogo i jediničnogo idealov.\par
\noindent\emph{Textual comparison:} Paper 30 adds integral closedness and, under all the axioms, gives a unique product of prime-ideal powers. This is stronger than general primary decomposition and retains the zero/unit exclusions.\par
\noindent{\footnotesize TEXT-COMP-415; pasaže / passages: 16866, 16868, 16870, 16872, 16968, 16970; \texttt{CORPUS\_FORM\_COMPARISONS.json}.}\par
\end{minipage}\par
\par\medskip\Needspace{7\baselineskip}
\noindent\begin{minipage}{\linewidth}
\hypertarget{TERM-1331}{}
\noindent\textbf{stepenj svojego asociovanogo prostogo ideala / produkt stepenjev}\quad{\footnotesize\texttt{TERM-1331}}\par
\noindent\emph{Iztočnik / Source:} Potenz seines zugehörigen Primideals / Potenzprodukt\par
\noindent\emph{Rusky / Russian:} степень своего соответствующего простого идеала / произведение степеней\par
\noindent\emph{Ukrajinsky / Ukrainian:} степінь свого асоційованого простого ідеалу / добуток степенів\par
\smallskip\noindent\emph{Razlog:} Stepenj svojego asociovanogo prostogo ideala / produkt stepenjev povězuje jednotlivu primarnu komponentu s produktnym prědstavjenjem cělogo ideala. To jest zapisany razlog za izrične slova stepenj i produkt v Teoremě VI raboty 30. Ruska i ukrajinska paralela čuvajut oba nivoja. Formula s produktom stepenjev prostyh idealov ne jest suma ni obče prěsěčenje primarnyh komponentov. Tekst daje jednoznačnost pri aksiomah I-V, s izključenjem nulovogo i jediničnogo idealov. Tvrditi samo, že idealy možno razložiti, bylo by nedostatočno: izgubili by se typ faktorov, eksponenty i sila uslovij.\par
\noindent\emph{Reason:} Stepenj svojego asociovanogo prostogo ideala / produkt stepenjev connects an individual primary component with the product representation of the whole ideal. This is the recorded reason for making power and product explicit in Paper 30's Theorem VI. Russian and Ukrainian preserve both levels. The displayed product of powers of prime ideals is neither a sum nor an unrestricted intersection of primary components. The text gives uniqueness under axioms I-V, excluding the zero and unit ideals. Saying only that ideals can be decomposed would lose the factor type, exponents and strength of the hypotheses.\par
\noindent{\footnotesize Izvorne zapisy / Original records: HIST-g115-003.}\par
\smallskip\noindent\emph{Sravnjenje s tekstom:} Rabota 30 pridava integralnu zamknjenost i pod vsimi aksiomami dava jednoznačny produkt stepenjev prostyh idealov. To jest silnějše od občego primarnogo razloženja i čuva izključenija nulovogo i jediničnogo idealov.\par
\noindent\emph{Textual comparison:} Paper 30 adds integral closedness and, under all the axioms, gives a unique product of prime-ideal powers. This is stronger than general primary decomposition and retains the zero/unit exclusions.\par
\noindent{\footnotesize TEXT-COMP-415; pasaže / passages: 16866, 16868, 16870, 16872, 16968, 16970; \texttt{CORPUS\_FORM\_COMPARISONS.json}.}\par
\end{minipage}\par
\par\medskip\Needspace{7\baselineskip}
\noindent\begin{minipage}{\linewidth}
\hypertarget{TERM-1332}{}
\noindent\textbf{produkt stepenjev}\quad{\footnotesize\texttt{TERM-1332}}\par
\noindent\emph{Iztočnik / Source:} Potenzprodukt\par
\noindent\emph{Rusky / Russian:} произведение степеней\par
\noindent\emph{Ukrajinsky / Ukrainian:} добуток степенів\par
\smallskip\noindent\emph{Razlog:} Produkt stepenjev razklada německo složeno slovo na razumlivu genitivnu konstrukciju. Izričny razlog jest razlika medžu produktom stepenjev prostyh idealov i prosto produktom primarnyh idealov. Ruska i ukrajinska paralela upotrěbljajut analogičnu razloženu frazu. V Teoremě VI raboty 30 konečno mnogo prostyh faktorov stoji pod konkretnymi aksiomami; to ne jest obča definicija primarnogo razloženja. V innych odlomkah isti izraz može označati produkt stepenjev proměnnyh. Tymy upotrěbami možno děliti tvorbenu formu, ale ne matematičny typ faktorov; slovnik dolžen jest razlikovati te smysly.\par
\noindent\emph{Reason:} Produkt stepenjev expands the German compound into a transparent genitive construction. The explicit reason is to distinguish a product of powers of prime ideals from a mere product of primary ideals. Russian and Ukrainian use corresponding expanded phrases. In Paper 30's Theorem VI, finitely many prime factors occur under specific axioms; this is not a general definition of primary decomposition. Elsewhere the same expression can denote a product of powers of variables. The wording can be shared, but the mathematical factor type cannot: the dictionary must distinguish these senses.\par
\noindent{\footnotesize Izvorne zapisy / Original records: HIST-g115-009.}\par
\smallskip\noindent\emph{Sravnjenje s tekstom:} Rabota 30 pridava integralnu zamknjenost i pod vsimi aksiomami dava jednoznačny produkt stepenjev prostyh idealov. To jest silnějše od občego primarnogo razloženja i čuva izključenija nulovogo i jediničnogo idealov.\par
\noindent\emph{Textual comparison:} Paper 30 adds integral closedness and, under all the axioms, gives a unique product of prime-ideal powers. This is stronger than general primary decomposition and retains the zero/unit exclusions.\par
\noindent{\footnotesize TEXT-COMP-415; pasaže / passages: 16866, 16868, 16870, 16872, 16968, 16970; \texttt{CORPUS\_FORM\_COMPARISONS.json}.}\par
\end{minipage}\par
\par\medskip\Needspace{7\baselineskip}
\noindent\begin{minipage}{\linewidth}
\hypertarget{TERM-1333}{}
\noindent\textbf{stepenovy produkt}\quad{\footnotesize\texttt{TERM-1333}}\par
\noindent\emph{Iztočnik / Source:} Potenzprodukt\par
\noindent\emph{Rusky / Russian:} степенное произведение\par
\noindent\emph{Ukrajinsky / Ukrainian:} степеневий добуток\par
\smallskip\noindent\emph{Razlog:} Stepenovy produkt čuva doslovnu strukturu německogo nazvanja vměsto modernizacije na monom. Zapis navodi ranějši priměr, na ktory se primětka vrati; dve arhivne svezky ne sut dve nezavisne čitateljske prověrky. V dostupnoj primětkě raboty 19 ideal $(x^2,xy,y^\lambda)$ sodrživaje vse produkty stepenjev danogo občego stepena, ale ostava primarny i razložimy. Izraz označaje komutativne produkty od $x,y$, ne produkty prostyh idealov. Njegova uloga jest v dokazu nilpotentnosti ostatkov polinomov bez konstantnogo člena. Čuvanje historičnogo nazvanja ne opravdava mněnija, že každý primarny ideal jest nerazložimy.\par
\noindent\emph{Reason:} Stepenovy produkt retains the literal structure of the German name instead of modernizing it to monomial. The record refers back to an earlier example; its two archival links are not two independent reader assessments. In Paper 19's available footnote, $(x^2,xy,y^\lambda)$ contains all power products of a specified total degree but is primary and reducible. The expression denotes commutative products in $x,y$, not products of prime ideals. Its role is in showing nilpotence of residues of polynomials without a constant term. Preserving the historical name does not justify treating every primary ideal as irreducible.\par
\noindent{\footnotesize Izvorne zapisy / Original records: HIST-g081-045, HIST-g082-045.}\par
\smallskip\noindent\emph{Sravnjenje s tekstom:} Primětka raboty 19 upotrěbljaje produkty stepenjev od x i y v kontrapriměru k obratu teorema. Ne govori o produktah prostyh idealov ni o nerazložimosti vsakego primarnogo ideala.\par
\noindent\emph{Textual comparison:} Paper 19's footnote uses power products in x and y in a counterexample to the converse theorem. It concerns neither products of prime ideals nor irreducibility of every primary ideal.\par
\noindent{\footnotesize TEXT-COMP-416; pasaže / passages: 12836, 12838; \texttt{CORPUS\_FORM\_COMPARISONS.json}.}\par
\end{minipage}\par
\par\medskip\Needspace{7\baselineskip}
\noindent\begin{minipage}{\linewidth}
\hypertarget{TERM-1334}{}
\noindent\textbf{produkt stepenjev prostyh idealov / vlastny dělitelj}\quad{\footnotesize\texttt{TERM-1334}}\par
\noindent\emph{Iztočnik / Source:} Potenzprodukt von Primidealen / echter Teiler\par
\noindent\emph{Rusky / Russian:} произведение степеней простых идеалов / собственный делитель\par
\noindent\emph{Ukrajinsky / Ukrainian:} добуток степенів простих ідеалів / власний дільник\par
\smallskip\noindent\emph{Razlog:} Produkt stepenjev prostyh idealov / vlastny dělitelj bira razloženu frazu vměsto stisnutogo složenogo slova. Zapis očekuje razumlivost za južne, istočne i zapadne slovjanske čitatelje, ale ne navodi izměreny rezultat ni točnu kraču alternativu. V rabotě 31, §8, ta fraza jest čestju aksiomatičnogo opisa prstena bez děliteljev nuly i s jedinicoju. Jednoznačno razloženje izključaje nulovy i jediničny ideal; prosty faktory ne imajut vlastnogo dělitelja različnogo od jediničnogo ideala. Vlastny tu znači strogo različny v danom odnosu děliteljnosti, ne vlastništvo. Niti tu ogradu niti osnovne uslovja ne možno izpustiti radi kratkosti.\par
\noindent\emph{Reason:} Produkt stepenjev prostyh idealov / vlastny dělitelj chooses an expanded phrase over a compressed compound. The record anticipates comprehension across South, East and West Slavic readers, but reports no measurement or exact shorter alternative. In Paper 31, §8, it belongs to the axiomatic description of a ring without zero divisors and with an identity. Unique factorization excludes the zero and unit ideals; the prime factors have no proper divisor other than the unit ideal. Vlastny means strict difference in the specified divisibility relation, not ownership. Neither that qualification nor the ring assumptions may be omitted for brevity.\par
\noindent{\footnotesize Izvorne zapisy / Original records: HIST-g155-008.}\par
\smallskip\noindent\emph{Sravnjenje s tekstom:} Rabota 31 daje prstenove uslovja, jednoznačnost, izključenija i ogradu vlastnogo dělitelja v jednom opisu. Smysl vlastny jest relativen k děliteljnosti, ne k posědanju.\par
\noindent\emph{Textual comparison:} Paper 31 gives ring hypotheses, uniqueness, exclusions and the proper-divisor qualification in one description. Proper is relative to divisibility, not possession.\par
\noindent{\footnotesize TEXT-COMP-417; pasaže / passages: 18789; \texttt{CORPUS\_FORM\_COMPARISONS.json}.}\par
\end{minipage}\par
\par\medskip\Needspace{7\baselineskip}
\noindent\begin{minipage}{\linewidth}
\hypertarget{TERM-1335}{}
\noindent\textbf{potencijalne produkty}\quad{\footnotesize\texttt{TERM-1335}}\par
\noindent\emph{Iztočnik / Source:} Potenzprodukte\par
\noindent\emph{Rusky / Russian:} степенные произведения\par
\noindent\emph{Ukrajinsky / Ukrainian:} степеневі добутки\par
\smallskip\noindent\emph{Razlog:} Potencijalne produkty jest stvarno zapisana i v dostupnom tekstu raboty 9 ješče upotrěbljena forma. Arhiv opravdava ju kratkostju internacionalnogo složenja; odlomok ale jasno govori o monomnyh produktah stepenjev proměnnyh $\eta$. Nazvanje zato ne znači možlive ili ješče neizpolnjene produkty. To jest slabo město zapisanogo izbora: forma sama može skryti značenje stepeni. Drugde upotrěbljane produkt stepenjev i stepenovy produkt davajut izričnje sravnjenje; to jest nyněšnje sravniteljno objasnenje, ne izmysljena ranějša kandidatska ocěna. Svaky pojediny element sodrživaje samo konečno mnogo takyh produktov, čeho nazvanje potencialny ne objašnjava. Arhivnu formu čuvajemo kako svědectvo, ne kako potvrđenu normu.\par
\noindent\emph{Reason:} Potencijalne produkty is genuinely recorded and still occurs in the available Paper 9 text. The archive defends a compact international compound, but the passage clearly concerns monomial products of powers of the variables $\eta$, not possible or unrealized products. This is a weakness in the recorded choice: the wording can obscure exponentiation. Produkt stepenjev and stepenovy produkt elsewhere provide more explicit comparanda; this is a retrospective assessment, not an invented earlier candidate score. Each individual element contains only finitely many such products, a fact not explained by potential. The archival form is preserved as evidence, not endorsed as an established norm.\par
\noindent{\footnotesize Izvorne zapisy / Original records: HIST-g067-005.}\par
\smallskip\noindent\emph{Sravnjenje s tekstom:} Rabota 9 čuva potencijalne produkty, ale konstrukcija zahtěvaje monomne produkty fiksovanogo občego stepena i konečne sumy. Razuměnje mora slědovati formule; arhivna kratkost sama ne opravdava dvosmyslno nazvanje.\par
\noindent\emph{Textual comparison:} Paper 9 retains potencijalne produkty, but the construction requires monomial products of fixed total degree and finite sums. Meaning must follow the formula; archival brevity alone does not justify the ambiguous wording.\par
\noindent{\footnotesize TEXT-COMP-418; pasaže / passages: 8183, 8184, 8186, 8187, 8190, 8191, 8192, 8193, 8194, 8195, 8196, 8197, 8198; \texttt{CORPUS\_FORM\_COMPARISONS.json}.}\par
\end{minipage}\par
\par\medskip\Needspace{7\baselineskip}
\noindent\begin{minipage}{\linewidth}
\hypertarget{TERM-1336}{}
\noindent\textbf{produkty stepenev}\quad{\footnotesize\texttt{TERM-1336}}\par
\noindent\emph{Iztočnik / Source:} Potenzprodukte\par
\noindent\emph{Rusky / Russian:} произведения степеней; степенные произведения\par
\noindent\emph{Ukrajinsky / Ukrainian:} добутки степенів; степеневі добутки\par
\smallskip\noindent\emph{Razlog:} Produkty stepenev spojuje kratko matematično produkt s slovjanskim korenjem stepenj. Dva zapisy podčrkajut kompozicijnu razumlivost i čuvanje děliteljnosti vměsto zaměny s monomy; jeden izrično ostavlja tehničny stil k daljšej ocěně. Ruska i ukrajinska paralela imajut i genitivnu i pridavničnu konstrukciju. V rabotě 15 ide o nenegativne eksponenty i uslovje zatvorenosti pri kvocientu, kogda kvocient jest polinomny. V rabotě 17 dokaz upotrěbljaje komutativnost neizvěstnyh i děliteljnost produktov v konečnom čislu proměnnyh. Konečny bazis ne znači konečnu sistemu vsih produktov. Arhivno stepenev i drugde stepenjev potrěbujut gramatično sravnjenje, ne množenje ponjatij.\par
\noindent\emph{Reason:} Produkty stepenev combines compact mathematical produkt with the Slavic stepenj stem. Two records emphasize compositional intelligibility and retention of divisibility language instead of monomials; one expressly leaves technical style for further assessment. Russian and Ukrainian offer both genitive and adjectival constructions. Paper 15 uses nonnegative exponents and closure under a quotient when that quotient is polynomial. Paper 17 uses commutativity of the indeterminates and divisibility of products in finitely many variables. A finite basis does not mean a finite collection of all products. Archival stepenev and stepenjev elsewhere require grammatical comparison, not multiplication of concepts.\par
\noindent{\footnotesize Izvorne zapisy / Original records: HIST-g077-008, HIST-g079-034.}\par
\smallskip\noindent\emph{Sravnjenje s tekstom:} Raboty 15 i 17 čuvajut nenegativne eksponenty, uslovja kvocienta ili děliteljnosti i konečno čislo proměnnyh. Konečny bazis ne znači konečno mnogo vsih produktov.\par
\noindent\emph{Textual comparison:} Papers 15 and 17 retain nonnegative exponents, quotient or divisibility conditions, and finitely many variables. A finite basis does not mean finitely many products altogether.\par
\noindent{\footnotesize TEXT-COMP-419; pasaže / passages: 10108, 10110, 10112, 11405, 11407, 11409, 11411, 11413, 11415, 11417; \texttt{CORPUS\_FORM\_COMPARISONS.json}.}\par
\end{minipage}\par
\par\medskip\Needspace{7\baselineskip}
\noindent\begin{minipage}{\linewidth}
\hypertarget{TERM-1337}{}
\noindent\textbf{stepenove produkty}\quad{\footnotesize\texttt{TERM-1337}}\par
\noindent\emph{Iztočnik / Source:} Potenzprodukte\par
\noindent\emph{Rusky / Russian:} степенные произведения\par
\noindent\emph{Ukrajinsky / Ukrainian:} степеневі добутки\par
\smallskip\noindent\emph{Razlog:} Stepenove produkty čuva pridavničnu konstrukciju, paralelnu ruskoj i ukrajinskoj. Zapis daje prědnost doslovnomu power products pred monomial basis, zato že Noether opisuje modul v tyh produktah. Njegova historična formulacija ale sama ne dokazuje, kogda nastalo slovo monom. V rabotě 18 ideal pojma se kako modul linearnyh form v produktah stepenjev ranějših proměnnyh, poslě adjunkcije pozdnějših k oblasti koeficientov. To jest izbor načina prědstavjenja, ne tvrdženje, že ideal jest generovany samo monomami. Zameniti frazu slovami monomna baza bez dodatnogo dokaza moglo by změniti matematičny smysl, a ne samo modernizovati register.\par
\noindent\emph{Reason:} Stepenove produkty retains an adjectival construction parallel to Russian and Ukrainian. The record prefers literal power products to monomial basis because Noether describes a module in those products. Its historical wording does not by itself establish when the word monomial arose. In Paper 18, an ideal is viewed as a module of linear forms in powers of the earlier variables after the later variables are adjoined to the coefficient domain. This is a choice of representation, not an assertion that the ideal is generated by monomials alone. Substituting monomial basis without further justification could change the mathematics rather than merely modernize the register.\par
\noindent{\footnotesize Izvorne zapisy / Original records: HIST-g080-006.}\par
\smallskip\noindent\emph{Sravnjenje s tekstom:} Rabota 18 pojma ideal kako modul v produktah ranějših proměnnyh po razširjenju koeficientov. Ta reprezentacija ne tvrdi, že sam ideal jest generovany monomami.\par
\noindent\emph{Textual comparison:} Paper 18 views an ideal as a module in products of earlier variables after coefficient extension. This representation does not assert that the ideal itself is generated by monomials.\par
\noindent{\footnotesize TEXT-COMP-420; pasaže / passages: 12343, 12345, 12347, 12348, 12349, 12351, 12353, 12354, 12355, 12356; \texttt{CORPUS\_FORM\_COMPARISONS.json}.}\par
\end{minipage}\par
\par\medskip\Needspace{7\baselineskip}
\noindent\begin{minipage}{\linewidth}
\hypertarget{TERM-1338}{}
\noindent\textbf{produkty stepenev v imenovatelju / častne derivaty}\quad{\footnotesize\texttt{TERM-1338}}\par
\noindent\emph{Iztočnik / Source:} Potenzprodukte im Nenner / partielle Ableitungen\par
\noindent\emph{Rusky / Russian:} степенные произведения в знаменателе / частные производные\par
\noindent\emph{Ukrajinsky / Ukrainian:} степеневі добутки у знаменнику / частинні похідні\par
\smallskip\noindent\emph{Razlog:} Produkty stepenev v imenovatelju / častne derivaty povězuje algebraičnu redukciju s diferencialnoju interpretacijeju. Zapis čuva terminologiju raboty 17, §6, i izrične nazvanja imenovatelja i častnoj derivaty; ruska i ukrajinska paralela takože razlikujut oba objekta. Dokaz ne utvrdžuje, že vsaky imenovatelj jest jeden iz konečno mnogyh elementov. Dopušča produkty jih stepenjev, s eksponentami zavisečimi od redukcije. Pri prěhodu k derivatam tekst izbira regularnu točku, kde dany koeficienty ne izčezajut. V formule red derivacije jest suma indeksov, a imenovatelj označaje pojedine proměnne i jih redy. Čuvanje frazy ne dozvalja izbrisati te uslovja.\par
\noindent\emph{Reason:} Produkty stepenev v imenovatelju / častne derivaty connects algebraic reduction with its differential interpretation. The record continues Paper 17, §6, terminology and explicitly names denominators and partial derivatives; Russian and Ukrainian likewise distinguish the objects. The proof does not say that every denominator is one of finitely many elements. It permits products of their powers, with exponents depending on the reduction. The differential interpretation chooses a regular point where the specified coefficients do not vanish. In the formula, derivative order is the sum of the indices, while the denominator specifies each variable and its order. Retaining the phrase does not permit those conditions to disappear.\par
\noindent{\footnotesize Izvorne zapisy / Original records: HIST-g079-069.}\par
\smallskip\noindent\emph{Sravnjenje s tekstom:} Rabota 17 ograničaje osnovne množitelje imenovateljev, ne vse jih stepeni. Diferencialna interpretacija izrično zahtěvaje regularnu točku i neizčezajuče koeficienty.\par
\noindent\emph{Textual comparison:} Paper 17 restricts the basic denominator factors, not all their powers. The differential interpretation explicitly requires a regular point and nonvanishing coefficients.\par
\noindent{\footnotesize TEXT-COMP-421; pasaže / passages: 11991, 11993, 11994, 11996, 11997, 11998, 11999, 12001, 12003, 12005, 12010, 12011, 12012, 12015, 12016, 12018, 12019, 12020, 12021, 12022, 12023, 12024, 12026; \texttt{CORPUS\_FORM\_COMPARISONS.json}.}\par
\end{minipage}\par
\par\medskip\Needspace{7\baselineskip}
\noindent\begin{minipage}{\linewidth}
\hypertarget{TERM-1339}{}
\noindent\textbf{stepenove produkty $q$-oj stepeni}\quad{\footnotesize\texttt{TERM-1339}}\par
\noindent\emph{Iztočnik / Source:} Potenzprodukte $q$-ter Dimension\par
\noindent\emph{Rusky / Russian:} степенные произведения $q$-й степени\par
\noindent\emph{Ukrajinsky / Ukrainian:} степеневі добутки $q$-го ступеня\par
\smallskip\noindent\emph{Razlog:} Stepenove produkty $q$-oj stepeni tolkuje historično Dimension kako obči stepen, ne dimenziju prostora. To jest izričny razlog zapisa; ruska i ukrajinska paralela upotrěbljajut nazvanje stepeni. V rabotě 22, §3, produkty nastupajut v modulnoj identičnosti Dedekinda-Mertensa. $A^q$ sodrživaje produkty stepenjev generatorov občego stepena $q$ i jih cěločislene linearne kombinacije. To ne znači, že vsak pojediny generator imaje eksponent $q$, ni že prostor ima dimenziju $q$. Gornji indeks v oznakě modula i stepen monomnogo produkta povezani sut konstrukcijeju, ale ne sut prosto dve oznaky za čislo proměnnyh.\par
\noindent\emph{Reason:} Stepenove produkty $q$-oj stepeni interprets historical Dimension as total degree, not dimension of a space. This is the record's explicit reason; Russian and Ukrainian likewise use degree terminology. In Paper 22, §3, these products occur in the Dedekind-Mertens module identity. $A^q$ contains power products in the generators of total degree $q$ and their integer linear combinations. This does not give every individual generator exponent $q$, nor does it assert that a space has dimension $q$. The module superscript and the product's degree are linked by the construction, not alternative names for the number of variables.\par
\noindent{\footnotesize Izvorne zapisy / Original records: HIST-g087-010, HIST-g092-030.}\par
\smallskip\noindent\emph{Sravnjenje s tekstom:} Rabota 22 opisuje produkty občego stepena q i cěločislene linearne kombinacije v A na stepenj q. Stepenj ne znači dimenziju prostora ni jednak eksponent vsakego generatora.\par
\noindent\emph{Textual comparison:} Paper 22 describes products of total degree q and integer linear combinations in the qth power of A. Degree means neither vector-space dimension nor the same exponent on every generator.\par
\noindent{\footnotesize TEXT-COMP-422; pasaže / passages: 14560, 14562, 14563, 14565, 14568, 14570, 14572, 14574, 14576, 14579; \texttt{CORPUS\_FORM\_COMPARISONS.json}.}\par
\end{minipage}\par
\par\medskip\Needspace{7\baselineskip}
\noindent\begin{minipage}{\linewidth}
\hypertarget{TERM-1340}{}
\noindent\textbf{stepenna suma}\quad{\footnotesize\texttt{TERM-1340}}\par
\noindent\emph{Iztočnik / Source:} Potenzsumme\par
\noindent\emph{Rusky / Russian:} степенная сумма\par
\noindent\emph{Ukrajinsky / Ukrainian:} сума степенів\par
\smallskip\noindent\emph{Razlog:} Stepenna suma jest kompaktny termin za sumu jednakogo stepena različnyh argumentov. Arhiv povezuje jego s modernym jezykom simetričnyh funkcij i s drugim polnym sistemom invariantov; ruska paralela jest pridavnična, ukrajinska imaje genitivnu konstrukciju. V rabotě 7 formula izrično dava $S_\mu$ kako sumu $\mu$-tyh stepenjev $h$ linearnyh form. To ne jest $\mu$-ta stepenj jih sumy. Pri izvlačenju invariantov $J$ tekst čuva čislovy faktor: koeficient pri danom produktu od $u$ ne možno prosto identifikovati s $J$ bez toga faktora. Nazvanje samo ne oprědělja tyh normalizacij ni oblasti koeficientov.\par
\noindent\emph{Reason:} Stepenna suma is a compact term for a sum of the same power of different arguments. The archive relates it to modern symmetric-function terminology and the second complete invariant system; Russian is adjectival and Ukrainian uses a genitive construction. In Paper 7, the formula explicitly defines $S_\mu$ as the sum of the $\mu$th powers of $h$ linear forms, not the $\mu$th power of their sum. When extracting the invariants $J$, the text preserves a numerical factor: the coefficient of a given product in the $u$ variables cannot simply be identified with $J$ without that factor. The name alone does not fix these normalizations or the coefficient domain.\par
\noindent{\footnotesize Izvorne zapisy / Original records: HIST-g056-008.}\par
\smallskip\noindent\emph{Sravnjenje s tekstom:} Rabota 7 dava sumu jednakogo stepena linearnyh form, ne stepenj jih sumy. Svez s invariantami J čuva čislovy faktor pri izvlačenju koeficienta.\par
\noindent\emph{Textual comparison:} Paper 7 gives the sum of equal powers of linear forms, not the power of their sum. The relation to invariants J preserves a numerical factor when extracting the coefficient.\par
\noindent{\footnotesize TEXT-COMP-423; pasaže / passages: 6676, 6678, 6679, 6681, 6683, 6684, 6686, 6688, 6689, 6691, 6693, 6695, 6697; \texttt{CORPUS\_FORM\_COMPARISONS.json}.}\par
\end{minipage}\par
\par\medskip\Needspace{7\baselineskip}
\noindent\begin{minipage}{\linewidth}
\hypertarget{TERM-1341}{}
\noindent\textbf{prost vzhodno k / vzaimno prosty}\quad{\footnotesize\texttt{TERM-1341}}\par
\noindent\emph{Iztočnik / Source:} prim zu / teilerfremd\par
\noindent\emph{Rusky / Russian:} прост относительно / взаимно прост\par
\noindent\emph{Ukrajinsky / Ukrainian:} простий щодо / взаємно простий\par
\smallskip\noindent\emph{Razlog:} Prost vzhodno k / vzaimno prosty zapisuje dva odnosa, ne dva ravnoznačna izraza. V rabotě 30 prostost po odnošenju k idealu znači, že idealny kvocient ne měnja toj ideal; vzajemna prostost kako děliteljev jest druga definicija. Arhiv izrično čuva razlikovanje do teorema V, kde pri jediničnom elementu i dvojnym uslovju cěpov pojmy sběgajut se. Ruska i ukrajinska paralela takože čuvajut odnositelnu i vzajemnu konstrukciju. Zato kratši obči izraz vzaimno prosty ne može zaměniti vse pojavjenja. Starše vzhodno k jest arhivna forma; dostupny tekst ima po odnošenju k, čiji osobny razlog jest opisan v heslu o toj predložnoj konstrukciji. Novy slovnik mora pokazati i smysl, i razliku medžu arhivnoju i aktualnoju formoju.\par
\noindent\emph{Reason:} Prost vzhodno k / vzaimno prosty records two relations, not two equivalent expressions. In Paper 30, being prime with respect to an ideal means that taking the ideal quotient leaves that ideal unchanged; being coprime as divisors has a different definition. The record expressly preserves the distinction until Theorem V, where identity and the double chain condition make the notions coincide. Russian and Ukrainian likewise retain relational and mutual constructions. A shorter blanket translation as vzaimno prosty would therefore lose information. Vzhodno k is the older recorded wording; the available text uses po odnošenju k, whose separate rationale is explained in the entry on that prepositional construction. The dictionary must preserve both the meaning and the difference between recorded and current wording.\par
\noindent{\footnotesize Izvorne zapisy / Original records: HIST-g114-003, HIST-g114-009.}\par
\smallskip\noindent\emph{Sravnjenje s tekstom:} Idealny kvocient oprědělja napravljenje odnosa. Sběženje s vzajemnoju prostostoju jest pozdnějši rezultat pod silnějšimi uslovjami, ne početkova ravnoznačnost slov.\par
\noindent\emph{Textual comparison:} The ideal quotient fixes the direction of the relation. Coincidence with coprimeness is a later result under stronger assumptions, not an initial equivalence of words.\par
\noindent{\footnotesize TEXT-COMP-424; pasaže / passages: 16628, 16630, 16632, 16634, 16636, 16794, 16796, 16798, 16800, 16802, 16804; \texttt{CORPUS\_FORM\_COMPARISONS.json}.}\par
\end{minipage}\par
\par\medskip\Needspace{7\baselineskip}
\noindent\begin{minipage}{\linewidth}
\hypertarget{TERM-1342}{}
\noindent\textbf{prost vzhodno k $a$ / vzaimno proste v tom smyslu}\quad{\footnotesize\texttt{TERM-1342}}\par
\noindent\emph{Iztočnik / Source:} prim zu $a$ / gegenseitig prim\par
\noindent\emph{Rusky / Russian:} прост относительно $a$ / взаимно просты в этом смысле\par
\noindent\emph{Ukrajinsky / Ukrainian:} простий щодо $a$ / взаємно прості в цьому сенсі\par
\smallskip\noindent\emph{Razlog:} Prost vzhodno k $a$ / vzaimno proste v tom smyslu čuva napravljenje odnosa. V točnoj definiciji raboty 30, §5, ravnost $\mathfrak a:\mathfrak b=\mathfrak a$ kaže, že $\mathfrak b$ jest prost po odnošenju k $\mathfrak a$, ne obratno. Vzajemnost trěbuje oba napravljenja. Dodatok v tom smyslu oddělja toj pojem od raněje definiranoj vzajemnoj prostosti kako děliteljev; v obščem slučaju tekst dava samo implikaciju od poslednje k prvoj. Zapisany izbor zato plača dolžinoju frazy za točnost odnosa. Arhivno obične $a$ ostaje tak v zaglavju, dokole sravnjena formula upotrěbljaje frakturu. Aktualno po odnošenju k ne opravdava tiho prěpisanje staršego vzhodno k v historijnoj tabeli.\par
\noindent\emph{Reason:} Prost vzhodno k $a$ / vzaimno proste v tom smyslu preserves the direction of the relation. In the exact definition in Paper 30, §5, $\mathfrak a:\mathfrak b=\mathfrak a$ says that $\mathfrak b$ is prime with respect to $\mathfrak a$, not the reverse. Mutuality requires both directions. The qualification v tom smyslu distinguishes this notion from coprimeness as divisors, defined earlier; in general the text gives only the implication from the latter to the former. The recorded choice thus accepts a longer phrase to preserve relational precision. The archive's plain $a$ remains in the headword, while the compared formula uses fraktur. Current po odnošenju k does not justify silently rewriting historical vzhodno k in the comparison table.\par
\noindent{\footnotesize Izvorne zapisy / Original records: HIST-g112-004.}\par
\smallskip\noindent\emph{Sravnjenje s tekstom:} Idealny kvocient oprědělja napravljenje odnosa. Sběženje s vzajemnoju prostostoju jest pozdnějši rezultat pod silnějšimi uslovjami, ne početkova ravnoznačnost slov.\par
\noindent\emph{Textual comparison:} The ideal quotient fixes the direction of the relation. Coincidence with coprimeness is a later result under stronger assumptions, not an initial equivalence of words.\par
\noindent{\footnotesize TEXT-COMP-424; pasaže / passages: 16628, 16630, 16632, 16634, 16636, 16794, 16796, 16798, 16800, 16802, 16804; \texttt{CORPUS\_FORM\_COMPARISONS.json}.}\par
\end{minipage}\par
\par\medskip\Needspace{7\baselineskip}
\noindent\begin{minipage}{\linewidth}
\hypertarget{TERM-1343}{}
\noindent\textbf{nerazložimy polinom}\quad{\footnotesize\texttt{TERM-1343}}\par
\noindent\emph{Iztočnik / Source:} Primfunktion\par
\noindent\emph{Rusky / Russian:} неприводимый многочлен\par
\noindent\emph{Ukrajinsky / Ukrainian:} незвідний многочлен\par
\smallskip\noindent\emph{Razlog:} Nerazložimy polinom jest vědomo osuvremenjenje německogo Primfunktion v Štejnicovoj definiciji citovanoj v rabotě 31. Ruska i ukrajinska paralela takože imenujut polinom i jego nerazložimost, ne obču funkciju. Zapis dava prednost razumlivosti dlja sovremennogo algebraista, ale ne navodi čitateljsky test. Definicija razširjenja prvogo roda trěbuje, že vsaky element jest nulja takogo polinoma i že v vhodnom razširjenju polinom ima različne linearne faktory. Sama nerazložimost ne vključaje poslednje uslovje. Literalno prosta funkcija jest zachovano v drugyh rabotah; to različje politika po kontekstu mora objasniti, ne skryti pod jedinom obščim prěvodom.\par
\noindent\emph{Reason:} Nerazložimy polinom deliberately modernizes German Primfunktion in the Steinitz definition cited in Paper 31. Russian and Ukrainian likewise name a polynomial and its irreducibility rather than a function in general. The record prioritizes readability for present-day algebraists but supplies no reader test. The definition of an extension of the first kind requires every element to be a root of such a polynomial, which must split into distinct linear factors in a suitable extension. Irreducibility alone does not include the latter condition. The literal prosta funkcija is retained elsewhere in the corpus; a context-sensitive policy must explain that difference rather than conceal it beneath one universal translation.\par
\noindent{\footnotesize Izvorne zapisy / Original records: HIST-g124-008.}\par
\smallskip\noindent\emph{Sravnjenje s tekstom:} Štejnicova definicija vyžaduje nerazložimy polinom i različne linearne faktory po vhodnom razširjenju. Ostatkovo polje prostogo ideala daje kontekst definicije.\par
\noindent\emph{Textual comparison:} The Steinitz definition requires an irreducible polynomial and distinct linear factors after suitable extension. The residue field of a prime ideal supplies the context.\par
\noindent{\footnotesize TEXT-COMP-425; pasaže / passages: 17595, 17597, 17599; \texttt{CORPUS\_FORM\_COMPARISONS.json}.}\par
\end{minipage}\par
\par\medskip\Needspace{7\baselineskip}
\noindent\begin{minipage}{\linewidth}
\hypertarget{TERM-1344}{}
\noindent\textbf{prosta funkcija}\quad{\footnotesize\texttt{TERM-1344}}\par
\noindent\emph{Iztočnik / Source:} Primfunktion\par
\noindent\emph{Rusky / Russian:} простая функция\par
\noindent\emph{Ukrajinsky / Ukrainian:} проста функція\par
\smallskip\noindent\emph{Razlog:} Prosta funkcija čuva historijnu analogiju medžu polinomom i prostym idealom. To jest izričny razlog v zapisu raboty 24; ruska i ukrajinska paralela slědujut toj samoj literalnoj formě. Tekst definira izraz kako polinom nerazložimy nad $P(u)$, takže prosta tu ne znači lehka i ne označaje prostyh funkcij iz teorije měry. Prednost jest vidima paralela s idealami; cena jest potrěba čitati definiciju, zato že samo zaglavje ne kaže polinom. Alternativa nerazložimy polinom, izbrana v rabotě 31, čini objekt neposrědněje vidimym. Obe strategije sut zapisane v korpusu; slovnik ne predstavlja jih kako dva različna matematična objekta v tom kontekstu.\par
\noindent\emph{Reason:} Prosta funkcija retains the historical analogy between a polynomial and a prime ideal. That is the explicit reason in the Paper 24 record; Russian and Ukrainian follow the same literal wording. The text defines the expression as a polynomial irreducible over $P(u)$, so prosta means neither easy nor a simple function in measure theory. The advantage is a visible parallel with ideals; the cost is reliance on the definition, because the headword alone does not identify a polynomial. Nerazložimy polinom, chosen in Paper 31, makes the object more immediately explicit. Both strategies are recorded in the corpus; the dictionary does not treat them as different mathematical objects in this context.\par
\noindent{\footnotesize Izvorne zapisy / Original records: HIST-g095-010.}\par
\smallskip\noindent\emph{Sravnjenje s tekstom:} Literalna funkcija tu jest polinom nad ukazanym poljem. Primarna funkcija jest njegova stepenj; vlastna primarna funkcija ne jest prosta. Množinnost korenja i stepenj polinoma ostavajut različne veličiny.\par
\noindent\emph{Textual comparison:} Literal function here denotes a polynomial over the specified field. A primary function is a power of a prime function; a properly primary function is not prime. Root multiplicity and polynomial degree remain different quantities.\par
\noindent{\footnotesize TEXT-COMP-426; pasaže / passages: 15333, 15814, 15816, 15818, 15820, 15822, 15824, 15826; \texttt{CORPUS\_FORM\_COMPARISONS.json}.}\par
\end{minipage}\par
\par\medskip\Needspace{7\baselineskip}
\noindent\begin{minipage}{\linewidth}
\hypertarget{TERM-1345}{}
\noindent\textbf{ireducibilna funkcija / ireducibilne faktory}\quad{\footnotesize\texttt{TERM-1345}}\par
\noindent\emph{Iztočnik / Source:} Primfunktion / Primfaktoren\par
\noindent\emph{Rusky / Russian:} неприводимая функция / неприводимые множители\par
\noindent\emph{Ukrajinsky / Ukrainian:} незвідна функція / незвідні множники\par
\smallskip\noindent\emph{Razlog:} Ireducibilna funkcija / ireducibilne faktory osuvremenjuje Primfunktion i Primfaktoren v determinantnom odlomku raboty 34. Zapis dava internacionalnomu ireducibilny prednost zaradi kontinuiteta s jezykom teorije prědstavjenj; ruska i ukrajinska paralela upotrěbljajut svoje izrazy za nerazložimost. To ne jest tvrdženje, že polinomna nerazložimost i ireducibilnost prědstavjenja sut jedna definicija. Tekst povezuje je prěz determinanty blokov i pri algebraično zamknjenom osnovnom polju pokazuje nerazložime, medžu soboju različne faktory za neekvivalentne prědstavjenja. Uslovje o polju ostaje potrěbno. Arhiv ne dava čislenogo sravnjenja s nerazložimy; kontinuitet jest kvalitativny razlog, ne izměrena jezyčna prevaga.\par
\noindent\emph{Reason:} Ireducibilna funkcija / ireducibilne faktory modernizes Primfunktion and Primfaktoren in the determinant passage of Paper 34. The record favours international ireducibilny for continuity with representation-theoretic vocabulary; Russian and Ukrainian use their expressions for irreducibility. This does not mean that polynomial irreducibility and irreducibility of a representation have one definition. The text connects them through block determinants and, over an algebraically closed base field, obtains irreducible, mutually distinct factors for inequivalent representations. The field hypothesis remains essential. The record gives no numerical comparison with nerazložimy; continuity is a qualitative reason, not a measured linguistic advantage.\par
\noindent{\footnotesize Izvorne zapisy / Original records: HIST-g201-006.}\par
\smallskip\noindent\emph{Sravnjenje s tekstom:} Razloženje prědstavjenja dava determinantne faktory. Tvrdženje o jih nerazložimosti i različnosti izrično koristi algebraično zamknjeno osnovno polje; slovny korenj sam to uslovje ne nese.\par
\noindent\emph{Textual comparison:} Decomposing the representation gives determinant factors. Their irreducibility and distinctness expressly use an algebraically closed base field; the shared lexical stem does not carry that hypothesis by itself.\par
\noindent{\footnotesize TEXT-COMP-427; pasaže / passages: 22104, 22106, 22108, 22109, 22110, 22111, 22112, 22114, 22115, 22116, 22117, 22118, 22119, 22120, 22121, 22122, 22123, 22125, 22127, 22128, 22129, 22130, 22131, 22132, 22133, 22134, 22135; \texttt{CORPUS\_FORM\_COMPARISONS.json}.}\par
\end{minipage}\par
\par\medskip\Needspace{7\baselineskip}
\noindent\begin{minipage}{\linewidth}
\hypertarget{TERM-1346}{}
\noindent\textbf{prosta funkcija / primarna funkcija}\quad{\footnotesize\texttt{TERM-1346}}\par
\noindent\emph{Iztočnik / Source:} Primfunktion / Primärfunktion\par
\noindent\emph{Rusky / Russian:} простая функция / первичная функция\par
\noindent\emph{Ukrajinsky / Ukrainian:} проста функція / первинна функція\par
\smallskip\noindent\emph{Razlog:} Prosta funkcija / primarna funkcija čuva paru, ktora v rabotah 24 i 25 prěnosi se od polinomov k idealam. Prva jest nerazložimy polinom nad danom poljem, druga stepenj takogo polinoma; vlastna primarna funkcija ima stepenj vyšu od prve. Arhiv bira internacionalno primarna vměsto nove domašnej kovanice, ale ne imenuje tu kovanicu. Zato ne možno dodati izmišljenu tabelu kandidatov. Ruska i ukrajinska paralela čuvajut prost-/pervičn- odnosno prost-/pervyn- razlikovanje. Teoremy XII i XIII raboty 24 činjut razlikovanje nositelem argumenta: norma prostogo ideala jest prosta funkcija pod dodatnym uslovjem doskonalosti polja, ne bez njega. Nad nedoskonalym poljem norma može byti viša stepenj. Terminologična paralela ne dozvaljaje izbrisati to uslovje.\par
\noindent\emph{Reason:} Prosta funkcija / primarna funkcija retains a pair carried from polynomials to ideals in Papers 24 and 25. The first is an irreducible polynomial over the specified field; the second is a power of one; a properly primary function has exponent greater than one. The record chooses international primarna instead of a new native coinage but does not name that coinage. An invented candidate table would therefore be inappropriate. Russian and Ukrainian preserve their prime/primary contrasts. The distinction is load-bearing in Theorems XII and XIII of Paper 24: a prime ideal's norm is a prime function under the additional perfect-field hypothesis, not without it. Over an imperfect field the norm can be a higher power. The terminological parallel does not permit deletion of that hypothesis.\par
\noindent{\footnotesize Izvorne zapisy / Original records: HIST-g096-010, HIST-g100-003, HIST-g102-004.}\par
\smallskip\noindent\emph{Sravnjenje s tekstom:} Literalna funkcija tu jest polinom nad ukazanym poljem. Primarna funkcija jest njegova stepenj; vlastna primarna funkcija ne jest prosta. Množinnost korenja i stepenj polinoma ostavajut različne veličiny.\par
\noindent\emph{Textual comparison:} Literal function here denotes a polynomial over the specified field. A primary function is a power of a prime function; a properly primary function is not prime. Root multiplicity and polynomial degree remain different quantities.\par
\noindent{\footnotesize TEXT-COMP-426; pasaže / passages: 15333, 15814, 15816, 15818, 15820, 15822, 15824, 15826; \texttt{CORPUS\_FORM\_COMPARISONS.json}.}\par
\smallskip\noindent\emph{Sravnjenje s tekstom:} Teorem XIII daje prime/primary paralelu za normy nad doskonalym poljem. Teorem XIV i slědujuče priměry čuvajut izključenja nad nedoskonalym poljem.\par
\noindent\emph{Textual comparison:} Theorem XIII gives the prime/primary parallel for norms over a perfect field. Theorem XIV and the following examples retain the exceptions over an imperfect field.\par
\noindent{\footnotesize TEXT-COMP-428; pasaže / passages: 15661, 15663, 15665, 15669, 15679; \texttt{CORPUS\_FORM\_COMPARISONS.json}.}\par
\end{minipage}\par
\par\medskip\Needspace{7\baselineskip}
\noindent\begin{minipage}{\linewidth}
\hypertarget{TERM-1347}{}
\noindent\textbf{nerazložimy polinom prvogo roda}\quad{\footnotesize\texttt{TERM-1347}}\par
\noindent\emph{Iztočnik / Source:} Primfunktion erster Art\par
\noindent\emph{Rusky / Russian:} неприводимый многочлен первого рода\par
\noindent\emph{Ukrajinsky / Ukrainian:} незвідний многочлен першого роду\par
\smallskip\noindent\emph{Razlog:} Nerazložimy polinom prvogo roda prodolžaje izrično zapisano osuvremenjenje Primfunktion iz raboty 31, §2. Prvogo roda znači v tom odlomku separabilnost: po razširjenju polja polinom razpada se bez povtorenyh korenjev. V §3 tekst najprvo piše polje kako $P[x]/f(x)$, potom razširjaje koeficienty do $Q$. Ireducibilnost $f$ ne musi ostati nad $Q$; jego nove nerazložime faktory ostavajut parno vzajemno proste. Zato prvogo roda ne znači prvogo stepena i ne obešča, že polinom ostaje nerazložimy nad vsakym poljem. Isty modernizujuči izbor v ruskoj i ukrajinskoj paraleli pomaga vyjavniti objekt, ale ne zaměnjaje definiciju prědanogo historijnogo kvalifikatora.\par
\noindent\emph{Reason:} Nerazložimy polinom prvogo roda continues the explicitly recorded modernization of Primfunktion in Paper 31, §2. First kind here expresses separability: after field extension the polynomial splits without repeated roots. In §3 the text first writes a field as $P[x]/f(x)$, then extends coefficients to $Q$. Irreducibility of $f$ need not survive over $Q$; its new irreducible factors remain pairwise coprime. Thus prvogo roda means neither first degree nor irreducibility over every field. The same modernizing choice in the Russian and Ukrainian parallels helps identify the object but cannot replace the definition of the retained historical qualifier.\par
\noindent{\footnotesize Izvorne zapisy / Original records: HIST-g127-006.}\par
\smallskip\noindent\emph{Sravnjenje s tekstom:} Polinom, nerazložimy nad $P$, može razpasti se nad $Q$. Prvi rod garantuje različne faktory, ne obšču neproměnnost nerazložimosti.\par
\noindent\emph{Textual comparison:} A polynomial irreducible over $P$ may factor over $Q$. First kind guarantees distinct factors, not universal preservation of irreducibility.\par
\noindent{\footnotesize TEXT-COMP-429; pasaže / passages: 17702, 17704; \texttt{CORPUS\_FORM\_COMPARISONS.json}.}\par
\end{minipage}\par
\par\medskip\Needspace{7\baselineskip}
\noindent\begin{minipage}{\linewidth}
\hypertarget{TERM-1348}{}
\noindent\textbf{prosta funkcija prvogo roda}\quad{\footnotesize\texttt{TERM-1348}}\par
\noindent\emph{Iztočnik / Source:} Primfunktion erster Art\par
\noindent\emph{Rusky / Russian:} простая функция первого рода\par
\noindent\emph{Ukrajinsky / Ukrainian:} проста функція першого роду\par
\smallskip\noindent\emph{Razlog:} Prosta funkcija prvogo roda čuva jezyk raboty 24 i njegovu paralelu s prostymi idealami. Prvogo roda dodava uslovje o različnyh korenjah; ne jest samo ukras k prostoj funkciji. Lema V dokazuje to svojstvo dlja formy elementarnyh děliteljev prostogo ideala pri doskonalom osnovnom polju. Slědujuči odstavek izrično upozornja, že dodanje neoprěděljenyh može dati nedoskonalo racionalno funkcionalno polje v pozitivnoj karakteristice. Zato rezultat ne jest neposrědny poslědok tvrdženja, že osnovno polje jest doskonalo; tekst trěbuje osobny argument. Zapis bira kontinuitet s ranějšimi definicijami, dokole alternativa nerazložimy polinom prvogo roda v rabotě 31 vyjavnja polinom v samom nazvanju.\par
\noindent\emph{Reason:} Prosta funkcija prvogo roda preserves Paper 24's vocabulary and its parallel with prime ideals. First kind adds the distinct-root condition; it is not an ornamental addition to prime function. Lemma V establishes this property for the elementary-divisor form of a prime ideal over a perfect base field. The following paragraph expressly warns that adjoining indeterminates can produce an imperfect rational function field in positive characteristic. The result therefore does not follow immediately from perfection of the base field; the text needs a separate argument. The recorded choice favours continuity with earlier definitions, whereas nerazložimy polinom prvogo roda in Paper 31 identifies the polynomial directly in its name.\par
\noindent{\footnotesize Izvorne zapisy / Original records: HIST-g097-011.}\par
\smallskip\noindent\emph{Sravnjenje s tekstom:} Dostupny dokaz izrično oddělja doskonalo osnovno polje od racionalnogo funkcionalnogo razširjenja, ktoro može byti nedoskonalo. Lema jest potrěbna dlja danej formy.\par
\noindent\emph{Textual comparison:} The available proof expressly distinguishes the perfect base field from the rational function extension, which can be imperfect. The lemma is needed for the particular form.\par
\noindent{\footnotesize TEXT-COMP-430; pasaže / passages: 15467, 15469, 15471, 15474; \texttt{CORPUS\_FORM\_COMPARISONS.json}.}\par
\end{minipage}\par
\par\medskip\Needspace{7\baselineskip}
\noindent\begin{minipage}{\linewidth}
\hypertarget{TERM-1349}{}
\noindent\textbf{nerazložimy polinom drugogo roda}\quad{\footnotesize\texttt{TERM-1349}}\par
\noindent\emph{Iztočnik / Source:} Primfunktion zweiter Art\par
\noindent\emph{Rusky / Russian:} неприводимый многочлен второго рода\par
\noindent\emph{Ukrajinsky / Ukrainian:} незвідний многочлен другого роду\par
\smallskip\noindent\emph{Razlog:} Nerazložimy polinom drugogo roda prodolžaje modernizaciju Primfunktion, ale čuva historijnu opoziciju prvogo i drugogo roda. V priměru raboty 31 polinom $t^2-x$ nad ukazanym poljem karakteristiki $2$ jest nerazložimy, no neseparabilny. Zato ostatkovy prsten jest polje, ale razširjenje jest drugogo roda. To ne znači polinom drugogo stepena: stepenj $2$ pripada priměru, ne definiciji drugogo roda. Arhiv izrično dava toj formě prednost prěd literalnoju prostoju funkcijeju zaradi jasnosti ostatkovogo polinomnogo priměra. Teorem o diskriminantu odliča takovu prostu komponentu od vlastno primarnoj komponenty; ni jedna ne směje byti tiho zaměnjena drugoju.\par
\noindent\emph{Reason:} Nerazložimy polinom drugogo roda continues the modernization of Primfunktion while retaining the historical first-kind/second-kind opposition. In Paper 31's example, $t^2-x$ over the specified characteristic-$2$ field is irreducible but inseparable. The quotient ring is therefore a field, yet the extension is of the second kind. This does not mean a polynomial of second degree: degree two belongs to the example, not the definition of second kind. The record expressly prefers this wording to literal prosta funkcija for clarity in the residue-polynomial example. The discriminant theorem distinguishes such a prime component from a properly primary component; neither may silently replace the other.\par
\noindent{\footnotesize Izvorne zapisy / Original records: HIST-g153-007.}\par
\smallskip\noindent\emph{Sravnjenje s tekstom:} Priměr $t^2-x$ v karakteristice $2$ oddělja nerazložimost od separabilnosti. Teorem rozlišuje prostu komponentu drugogo roda od vlastno primarnoj komponenty.\par
\noindent\emph{Textual comparison:} The characteristic-$2$ example $t^2-x$ separates irreducibility from separability. The theorem distinguishes a prime component of the second kind from a properly primary component.\par
\noindent{\footnotesize TEXT-COMP-431; pasaže / passages: 18715; \texttt{CORPUS\_FORM\_COMPARISONS.json}.}\par
\end{minipage}\par
\par\medskip\Needspace{7\baselineskip}
\noindent\begin{minipage}{\linewidth}
\hypertarget{TERM-1350}{}
\noindent\textbf{prosta grupa}\quad{\footnotesize\texttt{TERM-1350}}\par
\noindent\emph{Iztočnik / Source:} Primgruppe\par
\noindent\emph{Rusky / Russian:} простая группа\par
\noindent\emph{Ukrajinsky / Ukrainian:} проста група\par
\smallskip\noindent\emph{Razlog:} Prosta grupa prěnosi Primgruppe, ale jej smysl trěbuje definiciju raboty 17. Tu jest to grupa ostatkov modula, ktory ne ima drugyh děliteljev kromě samogo sebe i polnogo jediničnogo modula. Arhiv izrično upozornja na riziko anahroničnogo čitanja kako sovremennoj abstraktnoj proste grupy bez netrivialnyh normalnyh podgrup. Ruska i ukrajinska paralela čuvajut isto kratko nazvanje; cena jest potrěba držati modulnu strukturu i dopustime operacije vidimymi. Tekst dodava, že takove proste grupy mogut byti beskonečne i že sut ireducibilne v jeho smyslu. Slovnik ne prěnosi ty tvrdženja na libovoljnu modernu kategoriju grup. Kontinuitet historijnogo nazvanja jest opravdanje izbora, ne dokaz ravnoznačnosti vsego slova prosty.\par
\noindent\emph{Reason:} Prosta grupa translates Primgruppe, but its meaning requires Paper 17's definition. Here it is the residue group of a module having no divisors other than itself and the full unit module. The record expressly warns against an anachronistic reading as a modern abstract simple group with no nontrivial normal subgroups. Russian and Ukrainian retain the same short name; the cost is that the module structure and admissible operations must remain visible. The text adds that these simple groups can be infinite and are irreducible in its sense. The dictionary does not transfer those statements to an arbitrary modern category of groups. Continuity of the historical name motivates the choice, not equivalence of every use of prosty.\par
\noindent{\footnotesize Izvorne zapisy / Original records: HIST-g079-044.}\par
\smallskip\noindent\emph{Sravnjenje s tekstom:} Definicija počina modulnymi děliteljami i potom nazyvaje grupu ostatkov prostoju. To jest kontrola historijnogo smysla, ne dovoljenje uvoziti lubovoljnu modernu definiciju proste grupy.\par
\noindent\emph{Textual comparison:} The definition begins with module divisors and then calls the residue group simple. This controls the historical sense; it does not license importing an arbitrary modern definition of simple group.\par
\noindent{\footnotesize TEXT-COMP-432; pasaže / passages: 11568, 11570, 11572, 11574, 11576; \texttt{CORPUS\_FORM\_COMPARISONS.json}.}\par
\end{minipage}\par
\par\medskip\Needspace{7\baselineskip}
\noindent\begin{minipage}{\linewidth}
\hypertarget{TERM-1351}{}
\noindent\textbf{prosty element}\quad{\footnotesize\texttt{TERM-1351}}\par
\noindent\emph{Iztočnik / Source:} Primgröße\par
\noindent\emph{Rusky / Russian:} простой элемент\par
\noindent\emph{Ukrajinsky / Ukrainian:} простий елемент\par
\smallskip\noindent\emph{Razlog:} Prosty element zaměnja literalnu prostu veličinu konkretnym algebraičnym objektom. Oba zapisa raboty 31 davajut toj modernizaciji prednost; ruska i ukrajinska paralela takože imajut element. §7 definira $p$ takim, že glavny ideal, ktory on poraždaje v osnovnoj oblasti, jest prosty i različny od nulovogo i jediničnogo ideala. To ne kaže, že njegov razširjeny ideal ostaje prosty. V §8 po lokalizaciji element $p$ poraždaje lokalny prosty ideal; element i ideal ne sut ista věč. Različny šrift i okolny prsten čuvajut razliku. Nerazložimy element ne jest avtomatična zaměna: pri obščem prstenu nerazložimost i prostost ne sut ista vlastnost.\par
\noindent\emph{Reason:} Prosty element replaces literal prime magnitude with a concrete algebraic object. Both Paper 31 records favour that modernization; Russian and Ukrainian likewise use element. Section 7 defines $p$ by requiring the principal ideal it generates in the base domain to be prime and different from the zero and unit ideals. This does not say that its extended ideal remains prime. In §8, after localization, an element $p$ generates the local prime ideal; element and ideal are not the same object. Typography and the ambient ring preserve the distinction. Irreducible element is not an automatic substitute: in a general ring, irreducibility and primality are different properties.\par
\noindent{\footnotesize Izvorne zapisy / Original records: HIST-g151-002, HIST-g160-005.}\par
\smallskip\noindent\emph{Sravnjenje s tekstom:} Element $p$ jest odličen od ideala, ktory poraždaje. Prostost v osnovnoj oblasti i ponašanje razširjenogo ideala sut različne tvrdžbe; lokalizacija pozdněje dava osobny generator.\par
\noindent\emph{Textual comparison:} The element $p$ is distinct from the ideal it generates. Primality in the base domain and the behaviour of the extended ideal are different assertions; localization later supplies a particular generator.\par
\noindent{\footnotesize TEXT-COMP-433; pasaže / passages: 18627, 18629, 18631, 18633, 18637, 18970; \texttt{CORPUS\_FORM\_COMPARISONS.json}.}\par
\end{minipage}\par
\par\medskip\Needspace{7\baselineskip}
\noindent\begin{minipage}{\linewidth}
\hypertarget{TERM-1352}{}
\noindent\textbf{prosty element osnovnoj oblasti}\quad{\footnotesize\texttt{TERM-1352}}\par
\noindent\emph{Iztočnik / Source:} Primgröße des Grundbereichs\par
\noindent\emph{Rusky / Russian:} простой элемент основной области\par
\noindent\emph{Ukrajinsky / Ukrainian:} простий елемент основної області\par
\smallskip\noindent\emph{Razlog:} Prosty element osnovnoj oblasti čuva ne samo izbor slova element za Größe, ale i město, kde prostost jest utvrdžena. Rabota 31 sravnjava osnovnu oblast s poredkom v razširjenom polju; element prosty v osnovnoj oblasti može v razširjenom poredku dati ideal s několiko primarnymi komponentami. Zato skračenje do prosty element bez konteksta može skryti potrěbnu razliku. Uvod povezuje dělitelje diskriminanta s vlastno primarnymi komponentami abo prostymi komponentami drugogo roda. Ruska i ukrajinska paralela čuvajut kvalifikator osnovnoj oblasti. Zapisany razlog jest semantična točnost i jasny osnovny objekt, ne tvrdženje, že každa osnovna oblast jest polje ili že vsaka komponenta jest prosta.\par
\noindent\emph{Reason:} Prosty element osnovnoj oblasti preserves not only element for Größe but also where primality is asserted. Paper 31 compares the base domain with an order in an extension field; an element prime in the base domain can generate an ideal with several primary components in the extended order. Shortening the phrase to prime element without context can therefore hide a necessary distinction. The introduction relates discriminant divisors to properly primary components or prime components of the second kind. Russian and Ukrainian retain the base-domain qualifier. The recorded reason is semantic precision and an explicit base object, not a claim that every base domain is a field or every component prime.\par
\noindent{\footnotesize Izvorne zapisy / Original records: HIST-g118-010.}\par
\smallskip\noindent\emph{Sravnjenje s tekstom:} Uvod i diskriminantny teorem označajut osnovnu oblast i razširjeny poredok. Kvalifikator osnovnoj oblasti čuva to napravljenje i razliku medžu prostym elementom i komponentami njegovogo razširjenogo ideala.\par
\noindent\emph{Textual comparison:} The introduction and discriminant theorem identify the base domain and extended order. The base-domain qualifier preserves that direction and distinguishes a prime element from the components of its extended ideal.\par
\noindent{\footnotesize TEXT-COMP-434; pasaže / passages: 17275, 17277, 17279, 17281, 18627, 18715; \texttt{CORPUS\_FORM\_COMPARISONS.json}.}\par
\end{minipage}\par
\par\medskip\Needspace{7\baselineskip}
\noindent\begin{minipage}{\linewidth}
\hypertarget{TERM-1353}{}
\noindent\textbf{prost ideal}\quad{\footnotesize\texttt{TERM-1353}}\par
\noindent\emph{Iztočnik / Source:} Primideal\par
\noindent\emph{Rusky / Russian:} простой идеал\par
\noindent\emph{Ukrajinsky / Ukrainian:} простий ідеал\par
\smallskip\noindent\emph{Razlog:} Prost ideal jest kratka arhivna forma za Primideal v rabotě 24. Razlog jest stabilny algebraičny smysl i razlikovanje od primarnogo ideala; ruska i ukrajinska paralela imajut odgovarajuči prost- korenj. Definicija kaže, že ako produkt dvoh elementov leži v idealu, prinajmenje jeden faktor leži v njem. Pri primarnom idealu vměsto togo može byti potrěbna stepenj drugogo faktora. Zato primarny ne jest stilska razširjena forma od prosty. Arhivno prost i drugde zapisano prosty ne dokazuju dva koncepta; slovnik čuva razliku form bez izmišljanja matematičnoj razlike. Obča normativna tvrdžba zapisa sama po sobě ne dokumentuje nezavisnu konsultaciju vsih slovjanskyh tradicij.\par
\noindent\emph{Reason:} Prost ideal is the short recorded form for Primideal in Paper 24. The reason is a stable algebraic meaning distinct from primary ideal; Russian and Ukrainian have the corresponding prost- stem. The definition says that if a product of two elements lies in the ideal, at least one factor lies in it. For a primary ideal, a power of the other factor may instead be required. Primarny is therefore not a stylistically expanded form of prosty. Recorded prost and the prosty found elsewhere do not establish two concepts; the dictionary preserves the difference in form without inventing a mathematical difference. The record's general appeal to standard usage does not itself document independent consultation of all Slavic traditions.\par
\noindent{\footnotesize Izvorne zapisy / Original records: HIST-g095-006.}\par
\smallskip\noindent\emph{Sravnjenje s tekstom:} Definicija prostogo ideala koristi same faktory; definicija primarnogo ideala dopušča stepenj. Kratka forma prost ne směje izbrisati razlikovanje od primarny.\par
\noindent\emph{Textual comparison:} The prime-ideal definition uses the factors themselves; the primary-ideal definition allows a power. The short form prost must not erase its distinction from primarny.\par
\noindent{\footnotesize TEXT-COMP-435; pasaže / passages: 15348, 15350, 15352, 15355; \texttt{CORPUS\_FORM\_COMPARISONS.json}.}\par
\end{minipage}\par
\par\medskip\Needspace{7\baselineskip}
\noindent\begin{minipage}{\linewidth}
\hypertarget{TERM-1354}{}
\noindent\textbf{prosty ideal}\quad{\footnotesize\texttt{TERM-1354}}\par
\noindent\emph{Iztočnik / Source:} Primideal\par
\noindent\emph{Rusky / Russian:} простой идеал\par
\noindent\emph{Ukrajinsky / Ukrainian:} простий ідеал\par
\smallskip\noindent\emph{Razlog:} Prosty ideal spojuje tri zapisa iz rabot 20, 31 i 35. Vse čuvajut isti prost- korenj kako ruska i ukrajinska paralela, ale matematična uloga zavisi od odlomka. Ostrovskogo poslědok dopušča konečno mnogo izključenij pri redukciji modulo prosty ideal. V rabotě 31 prostost kontrakcije slěduje iz vloženja ostatkovogo prstena; to ne jest tvrdženje, že razširjenje vsečego prostogo ideala ostaje prosto. V rabotě 35 ideal relacij jest prost zato, že homomorfny obraz jest bez nulovyh děliteljev. Kontinuitet nazvanja jest zapisany razlog, no on ne zaměnjaje tri različne argumenty i jih uslovja. Kratka varianta prost ideal iz raboty 24 ostaje formovym, ne konceptualnym razlikovanjem.\par
\noindent\emph{Reason:} Prosty ideal joins three records from Papers 20, 31 and 35. All retain the same prost- stem as Russian and Ukrainian, but the mathematical role depends on the passage. Ostrowski's consequence permits finitely many exceptions when reducing modulo a prime ideal. In Paper 31, primality of contraction follows from embedding the residue-class ring; this does not assert that extending every prime ideal preserves primality. In Paper 35, the relation ideal is prime because the homomorphic image has no zero divisors. Continuity of naming is the recorded reason, but it cannot replace those three different arguments and their hypotheses. The short prost ideal from Paper 24 remains a difference of form, not of concept.\par
\noindent{\footnotesize Izvorne zapisy / Original records: HIST-g083-008, HIST-g120-008, HIST-g205-007.}\par
\smallskip\noindent\emph{Sravnjenje s tekstom:} Tri odlomky koriste isti termin, ale različne argumenty: konečna izključenja pri redukciji, vloženje pri kontrakciji i otsutnost nulovyh děliteljev v homomorfnom obrazu. Njih uslovja ostavajut osobne.\par
\noindent\emph{Textual comparison:} Three passages use the same term but different arguments: finitely many exceptions under reduction, an embedding under contraction, and absence of zero divisors in a homomorphic image. Their hypotheses remain separate.\par
\noindent{\footnotesize TEXT-COMP-436; pasaže / passages: 13809, 13811, 13813, 17394, 17396, 17398, 22594, 22621, 22623; \texttt{CORPUS\_FORM\_COMPARISONS.json}.}\par
\end{minipage}\par
\par\medskip\Needspace{7\baselineskip}
\noindent\begin{minipage}{\linewidth}
\hypertarget{TERM-1355}{}
\noindent\textbf{prost ideal / vlastny primarny ideal}\quad{\footnotesize\texttt{TERM-1355}}\par
\noindent\emph{Iztočnik / Source:} Primideal / eigentliches Primärideal\par
\noindent\emph{Rusky / Russian:} простой идеал / собственный первичный идеал\par
\noindent\emph{Ukrajinsky / Ukrainian:} простий ідеал / власний первинний ідеал\par
\smallskip\noindent\emph{Razlog:} Prost ideal / vlastny primarny ideal čuva razlikovanje medžu prostym i stvarno ne-prostym primarnym slučajem v rabotě 24. Vlastny tu prěnosi eigentlich, ne samo uslovje, že ideal jest različny od cělogo prstena. To jest važno: prosty ideal takože može byti vlastny podideal, ale ne jest zato vlastno primarny v tom smyslu. Zapis bira primarny radi kontinuiteta i očekivanoj medžuslovjanskoj poznatosti; ruska i ukrajinska paralela imajut svoje izrazy za primary. Teorem XIII prěnosi razlikovanje na normy i formy elementarnyh děliteljev pri doskonalom osnovnom polju. Bez doskonalosti dostupny tekst dava posebne pojavy, takže prosta norma i prost ideal ne smějut byti bezuslovno izravnjene.\par
\noindent\emph{Reason:} Prost ideal / vlastny primarny ideal preserves the distinction between the prime case and the genuinely non-prime primary case in Paper 24. Vlastny here translates eigentlich, not merely the condition that an ideal differs from the whole ring. This matters: a prime ideal can also be a proper subideal, but is not thereby properly primary in this sense. The record chooses primarny for continuity and anticipated cross-Slavic familiarity; Russian and Ukrainian use their primary terms. Theorem XIII carries the distinction to norms and elementary-divisor forms over a perfect base field. Without perfection, the available text supplies exceptional behaviour, so a prime norm and a prime ideal cannot be equated unconditionally.\par
\noindent{\footnotesize Izvorne zapisy / Original records: HIST-g100-001.}\par
\smallskip\noindent\emph{Sravnjenje s tekstom:} Vlastno primarny v toj opoziciji znači primarny, ale ne prosty. Uslovje doskonalosti v normnom kriteriju ne směje propasti pri skračenju terminologije.\par
\noindent\emph{Textual comparison:} Properly primary in this opposition means primary but not prime. The perfect-field hypothesis in the norm criterion must not disappear when terminology is shortened.\par
\noindent{\footnotesize TEXT-COMP-437; pasaže / passages: 15333, 15596, 15598, 15600, 15602, 15604, 15661, 15663, 15665, 15669, 15679; \texttt{CORPUS\_FORM\_COMPARISONS.json}.}\par
\end{minipage}\par
\par\medskip\Needspace{7\baselineskip}
\noindent\begin{minipage}{\linewidth}
\hypertarget{TERM-1356}{}
\noindent\textbf{prosty ideal / primarny ideal / primarna komponenta}\quad{\footnotesize\texttt{TERM-1356}}\par
\noindent\emph{Iztočnik / Source:} Primideal / Primaerideal / Primaerkomponente\par
\noindent\emph{Rusky / Russian:} простой идеал / примарный идеал / примарная компонента\par
\noindent\emph{Ukrajinsky / Ukrainian:} простий ідеал / первинний ідеал / первинна компонента\par
\smallskip\noindent\emph{Razlog:} Prosty ideal / primarny ideal / primarna komponenta prěnosi slovnik teorije idealov v Noetherin dodatok k rabotě 45. Zapis izrično čuva ranějše primarny; ruska paralela tu ima primarnyj, dokole ukrajinska ima pervynnyj. Razlika medžu paralelami ne jest dokumentovany konkurs s čislenymi težinami. Tekst povezuje rezultat o multiplicitetu s dolžinoju ideala i izolyraje primarnu komponentu pri danom nuljovom punktu. Viša dostatočna stepenj prostogo ideala pomaže vyděliti komponentu, ale to ne znači, že sama komponenta jest samo ta stepenj. Jasna opozicija prosty/primarny/komponenta čuva i tip objekta, i jego ulogu v razloženju.\par
\noindent\emph{Reason:} Prosty ideal / primarny ideal / primarna komponenta carries ideal-theoretic vocabulary into Noether's supplement to Paper 45. The record expressly retains earlier primarny; the Russian parallel here has primarnyj while Ukrainian has pervynnyj. That difference between parallels is not a documented numerical competition. The text connects a multiplicity result with ideal length and isolates a primary component at a specified zero. A sufficiently high power of the prime ideal helps isolate the component, but this does not mean that the component itself is just that power. The explicit prime/primary/component contrast preserves both object type and role in the decomposition.\par
\noindent{\footnotesize Izvorne zapisy / Original records: HIST-g214-006.}\par
\smallskip\noindent\emph{Sravnjenje s tekstom:} Noetherin dodatok koristi teoriju idealov dlja multiplicitetov. Dodanje dostatočno vysokoj stepeni prostogo ideala vydělja primarnu komponentu; komponenta ne jest tym definovana kako gola stepenj.\par
\noindent\emph{Textual comparison:} Noether's supplement uses ideal theory for multiplicities. Adding a sufficiently high power of the prime ideal isolates a primary component; this does not define the component as a bare power.\par
\noindent{\footnotesize TEXT-COMP-438; pasaže / passages: 45-isv.tex:391, 45-isv.tex:393, 45-isv.tex:394, 45-isv.tex:395, 45-isv.tex:396, 45-isv.tex:397, 45-isv.tex:398, 45-isv.tex:399, 45-isv.tex:400, 45-isv.tex:407, 45-isv.tex:408, 45-isv.tex:410, 45-isv.tex:411, 45-isv.tex:412, 45-isv.tex:413, 45-isv.tex:414, 45-isv.tex:415, 45-isv.tex:416, 45-isv.tex:417, 45-isv.tex:425, 45-isv.tex:426, 45-isv.tex:479, 45-isv.tex:480, 45-isv.tex:481, 45-isv.tex:482, 45-isv.tex:483, 45-isv.tex:484, 45-isv.tex:485, 45-isv.tex:486, 45-isv.tex:487, 45-isv.tex:488, 45-isv.tex:489, 45-isv.tex:490, 45-isv.tex:492, 45-isv.tex:493, 45-isv.tex:494, 45-isv.tex:495; \texttt{CORPUS\_FORM\_COMPARISONS.json}.}\par
\end{minipage}\par
\par\medskip\Needspace{7\baselineskip}
\noindent\begin{minipage}{\linewidth}
\hypertarget{TERM-1357}{}
\noindent\textbf{prosty ideal / proste idealy}\quad{\footnotesize\texttt{TERM-1357}}\par
\noindent\emph{Iztočnik / Source:} Primideal / Primideale\par
\noindent\emph{Rusky / Russian:} простой идеал / простые идеалы\par
\noindent\emph{Ukrajinsky / Ukrainian:} простий ідеал / прості ідеали\par
\smallskip\noindent\emph{Razlog:} Prosty ideal / proste idealy zapisuje jednotnu i množnu formu istogo pojma v rabotě 19. Arhiv čuva prost- za ideal pridruženy primarnomu idealu i odběgaje zaměšanje s primarny. Definicija ukazuje razliku v eksponentah: prosty slučaj ne trěbuje vyšej stepeni, dokole primarny može ju trěbovati. Pri konečnoj bazě tekst prěhodi od uslovja na elementy k uslovju na idealy i polučaje jedinu dostatočnu stepenj. Množina elementov, ktoryh neka stepenj leži v primarnom idealu, dava njegov pridruženy prosty ideal. To ne tvrdi, že vsaky primarny ideal jest stepenj prostogo ideala. Gramatična množina proste idealy ne dodava osobnu kategoriju.\par
\noindent\emph{Reason:} Prosty ideal / proste idealy records singular and plural forms of one concept in Paper 19. The archive keeps prost- for the prime attached to a primary ideal and avoids conflating it with primarny. The definition distinguishes exponents: the prime case needs no higher power, whereas the primary case may. With a finite basis, the text passes from a condition on elements to a condition on ideals and obtains a uniform sufficient power. The set of elements some power of which lies in a primary ideal gives its associated prime. This does not assert that every primary ideal is a power of a prime ideal. The grammatical plural proste idealy adds no separate category.\par
\noindent{\footnotesize Izvorne zapisy / Original records: HIST-g081-040, HIST-g082-040.}\par
\smallskip\noindent\emph{Sravnjenje s tekstom:} Elementne definicije prěhodet pri konečnoj bazě k idealnym uslovjam. Pridruženy prosty ideal jest oprěděljen stepenjami elementov; nět obščego utožsamjenja primarnyh idealov s prostymi stepenjami.\par
\noindent\emph{Textual comparison:} With a finite basis, element-based definitions pass to ideal conditions. The associated prime is determined by powers of elements; primary ideals are not generally identified with prime powers.\par
\noindent{\footnotesize TEXT-COMP-439; pasaže / passages: 12785, 12787, 12789, 12791, 12793, 12795, 12797, 12799, 12801, 12803; \texttt{CORPUS\_FORM\_COMPARISONS.json}.}\par
\end{minipage}\par
\par\medskip\Needspace{7\baselineskip}
\noindent\begin{minipage}{\linewidth}
\hypertarget{TERM-1358}{}
\noindent\textbf{prosty ideal / primarny ideal / primarna komponenta}\quad{\footnotesize\texttt{TERM-1358}}\par
\noindent\emph{Iztočnik / Source:} Primideal / Primärideal / Primärkomponente\par
\noindent\emph{Rusky / Russian:} простой идеал / первичный идеал / первичная компонента\par
\noindent\emph{Ukrajinsky / Ukrainian:} простий ідеал / первинний ідеал / первинна компонента\par
\smallskip\noindent\emph{Razlog:} Prosty ideal / primarny ideal / primarna komponenta jest izrično razděljena trojica v uvodu raboty 30. Prosty označaje prime, internacionalno primarny označaje primary, a komponenta jeho ulogu v razloženju. Ruska i ukrajinska paralela čuvajut to razlikovanje svojimi formami. Uvod sam razdělja stupenje rezultata: uslovje vzrastajučih cěpov dava primarno razloženje, silnějše uslovja daju jednoznačnost i produkt vzajemno prostyh komponentov, a integralna zamknjenost s ostalymi aksiomami dovodi k stepenjam prostyh idealov. Zato čitatelj ne směje zameniti komponentu stepenju samo zaradi poznatogo specialnogo slučaja. Arhivny razlog jest terminologična razlučivost, ne dokaz, že vse tri slova imajut jednaku razprostranjenost v každom slovjanskom jezyku.\par
\noindent\emph{Reason:} Prosty ideal / primarny ideal / primarna komponenta is an expressly differentiated triad in Paper 30's introduction. Prosty denotes prime, international primarny denotes primary, and component identifies a role in a decomposition. Russian and Ukrainian retain the contrast in their own forms. The introduction itself separates levels of result: the ascending chain condition gives primary decomposition, stronger assumptions give uniqueness and a product of coprime components, and integral closure with the other axioms leads to powers of prime ideals. A reader must therefore not replace component by power merely because of that familiar special case. The archived reason is terminological differentiation, not proof of equal currency of all three words in every Slavic language.\par
\noindent{\footnotesize Izvorne zapisy / Original records: HIST-g107-007.}\par
\smallskip\noindent\emph{Sravnjenje s tekstom:} Uvod stupenjevano usiljava prědpoloženja. Primarno razloženje, jednoznačne komponenty, produkt i produkt prostyh stepenjev ne sut obščo zaměnjajeme opisy.\par
\noindent\emph{Textual comparison:} The introduction strengthens the assumptions in stages. Primary decomposition, unique components, a product and a product of prime powers are not generally interchangeable descriptions.\par
\noindent{\footnotesize TEXT-COMP-440; pasaže / passages: 16115, 16117, 16120, 16121, 16122, 16123, 16124, 16127, 16129, 16131, 16133, 16135, 16139, 16143; \texttt{CORPUS\_FORM\_COMPARISONS.json}.}\par
\end{minipage}\par
\par\medskip\Needspace{7\baselineskip}
\noindent\begin{minipage}{\linewidth}
\hypertarget{TERM-1359}{}
\noindent\textbf{prost ideal / primarny ideal / asociovany prost ideal}\quad{\footnotesize\texttt{TERM-1359}}\par
\noindent\emph{Iztočnik / Source:} Primideal / Primärideal / zugehöriges Primideal\par
\noindent\emph{Rusky / Russian:} простой идеал / первичный идеал / ассоциированный простой идеал\par
\noindent\emph{Ukrajinsky / Ukrainian:} простий ідеал / первинний ідеал / асоційований простий ідеал\par
\smallskip\noindent\emph{Razlog:} Prost ideal / primarny ideal / asociovany prost ideal modernizuje literalno naležeči izraz v rabotě 25. Arhiv bira asociovany za kratkost i sovremenny algebraičny idiom; ruska i ukrajinska paralela takože upotrěbljajut asociacijny korenj. Tvrdženje zapisa o razumlivosti jest jezyčna ocěna, ne rezultat testovanja čitateljev. Odlomok oprědělja pridruženy prosty ideal prěz dva děliteljske uslovja: on děli primarny ideal i neka njegova stepenj jest děljena tym primarnym idealom. Potom razlikuje jednoznačno oprěděljene proste idealy od ne vždy jednoznačnyh primarnyh komponentov. Asociovany ne znači proizvoljno povezany i ne čini vse komponenty jedinstvenymi.\par
\noindent\emph{Reason:} Prost ideal / primarny ideal / asociovany prost ideal modernizes the literal belonging expression in Paper 25. The record selects asociovany for brevity and modern algebraic idiom; Russian and Ukrainian likewise use an association stem. Its claim of intelligibility is a linguistic judgement, not a reader-test result. The passage specifies the associated prime through two divisibility conditions: it divides the primary ideal, and some power of it is divisible by that primary ideal. It then distinguishes uniquely determined primes from primary components that are not always unique. Associated therefore means neither arbitrarily connected nor that every component is unique.\par
\noindent{\footnotesize Izvorne zapisy / Original records: HIST-g102-005.}\par
\smallskip\noindent\emph{Sravnjenje s tekstom:} Tekst definira pridruženy prosty ideal i potom oddělja jednoznačnost prostyh idealov od nejednoznačnosti něktoryh primarnyh komponentov. Kratko asociovany nosi tu konkretnu relaciju.\par
\noindent\emph{Textual comparison:} The text defines the associated prime, then separates uniqueness of the primes from nonuniqueness of some primary components. The concise asociovany carries that specific relation.\par
\noindent{\footnotesize TEXT-COMP-441; pasaže / passages: 15828, 15832, 15834, 15835, 15837, 15839, 15841, 15843, 15845; \texttt{CORPUS\_FORM\_COMPARISONS.json}.}\par
\end{minipage}\par
\par\medskip\Needspace{7\baselineskip}
\noindent\begin{minipage}{\linewidth}
\hypertarget{TERM-1360}{}
\noindent\textbf{prosty ideal / primarny ideal / asociovany prosty ideal}\quad{\footnotesize\texttt{TERM-1360}}\par
\noindent\emph{Iztočnik / Source:} Primideal / Primärideal / zugehöriges Primideal\par
\noindent\emph{Rusky / Russian:} простой идеал / первичный идеал / ассоциированный простой идеал\par
\noindent\emph{Ukrajinsky / Ukrainian:} простий ідеал / первинний ідеал / асоційований простий ідеал\par
\smallskip\noindent\emph{Razlog:} Prosty ideal / primarny ideal / asociovany prosty ideal prodolžaje ranějše razlikovanje rabot 24 i 30. Zapis ne navodi novu sravnitelnu procěnu; razlog jest kontinuitet už oprěděljenyh pojmov. V §5 raboty 30 slaba definicija prostogo ideala koristi ostatkovy prsten bez nulovyh děliteljev; primarny slučaj dopušča nilpotentne nulove dělitelje i vodi k pridruženomu prostomu idealu. Osobno se razgleda silna idealna formulacija. Dostupny tekst ima v jednom od tamošnjih dokazov nerazrěšeno pytanje o eksponentě; to ne ukryvajem tvrdženjem, že cěly dokaz jest prověreny. Različje prosty/primarny ostaje jasne, ale ocěna termina ne jest potvrđenoju pravilnostju vsakogo matematičnogo kroka.\par
\noindent\emph{Reason:} Prosty ideal / primarny ideal / asociovany prosty ideal continues the distinction established in Papers 24 and 30. The record gives no new comparative evaluation; its reason is continuity of already defined concepts. In Paper 30, §5, the weak definition of prime ideal uses a residue ring without zero divisors; the primary case permits nilpotent zero divisors and leads to its associated prime. A stronger ideal-based formulation is considered separately. The available text has an unresolved exponent question in one of these proofs; this must not be hidden by claiming that the whole proof has been verified. The prime/primary distinction remains clear, but evaluating a term does not certify every mathematical step.\par
\noindent{\footnotesize Izvorne zapisy / Original records: HIST-g112-001.}\par
\smallskip\noindent\emph{Sravnjenje s tekstom:} Ostatkovo prsten i produkt idealov davajut dve formulacije. Nerazrěšeno pytanje o eksponentě v dostupnom dokaznom tekstu ostaje odděljeno od semantičnogo razlikovanja prosty/primarny.\par
\noindent\emph{Textual comparison:} The residue ring and products of ideals provide two formulations. The unresolved exponent question in the available proof text remains separate from the semantic prime/primary distinction.\par
\noindent{\footnotesize TEXT-COMP-442; pasaže / passages: 16601, 16603, 16605, 16607, 16609; \texttt{CORPUS\_FORM\_COMPARISONS.json}.}\par
\end{minipage}\par
\par\medskip\Needspace{7\baselineskip}
\noindent\begin{minipage}{\linewidth}
\hypertarget{TERM-1361}{}
\noindent\textbf{prosty ideal prvogo / drugogo roda}\quad{\footnotesize\texttt{TERM-1361}}\par
\noindent\emph{Iztočnik / Source:} Primideal erster / zweiter Art\par
\noindent\emph{Rusky / Russian:} простой идеал первого / второго рода\par
\noindent\emph{Ukrajinsky / Ukrainian:} простий ідеал першого / другого роду\par
\smallskip\noindent\emph{Razlog:} Prosty ideal prvogo / drugogo roda čuva Štejnicovu klasifikaciju, a ne zaměnjaje rod obščim slovom tip. To jest izrično zapisany izbor; ruska i ukrajinska paralela takože imajut rod. V uvodu raboty 31 klasifikacija ideala zavisi od razširjenja, ktoro tvori njegovo ostatkovo polje nad obrazom osnovnogo polja. Prvi rod označaje separabilny slučaj, drugy rod njegov protivny slučaj; ne ide o dva stepena ideala. Razlikovanje stava potrěbno nad nedoskonalym poljem. Teorem o diskriminantu zato ima dve alternative: vlastno primarna komponenta ili prosty ideal drugogo roda. Samo slovo prosty ne izključaje drugu alternativu.\par
\noindent\emph{Reason:} Prosty ideal prvogo / drugogo roda retains the Steinitz classification rather than replacing kind with a generic type word. That is the expressly recorded choice; Russian and Ukrainian likewise use rod. In Paper 31's introduction, the ideal's classification depends on the extension formed by its residue field over the image of the base field. First kind denotes the separable case and second kind the contrary case, not two powers of an ideal. The distinction becomes necessary over an imperfect field. The discriminant theorem therefore has two alternatives: a properly primary component or a prime ideal of the second kind. Prime alone does not exclude the second alternative.\par
\noindent{\footnotesize Izvorne zapisy / Original records: HIST-g118-009.}\par
\smallskip\noindent\emph{Sravnjenje s tekstom:} Rod ideala opira se na algebraičnom razširjenju ostatkovogo polja. Dve kategorije ne sut dva stepena, i zaglavno i ne pravi ih jednovremennymi vlastnostjami jednogo ideala.\par
\noindent\emph{Textual comparison:} The ideal's kind is determined through the residue-field extension. The two categories are not powers, and the heading's and does not make them simultaneous properties of one ideal.\par
\noindent{\footnotesize TEXT-COMP-443; pasaže / passages: 17275, 17595, 17597, 17599; \texttt{CORPUS\_FORM\_COMPARISONS.json}.}\par
\end{minipage}\par
\par\medskip\Needspace{7\baselineskip}
\noindent\begin{minipage}{\linewidth}
\hypertarget{TERM-1362}{}
\noindent\textbf{prosty ideal prvogo i drugogo roda}\quad{\footnotesize\texttt{TERM-1362}}\par
\noindent\emph{Iztočnik / Source:} Primideal erster und zweiter Art\par
\noindent\emph{Rusky / Russian:} простой идеал первого и второго рода\par
\noindent\emph{Ukrajinsky / Ukrainian:} простий ідеал першого і другого роду\par
\smallskip\noindent\emph{Razlog:} Prosty ideal prvogo i drugogo roda čuva prost- za prime i ne směšiva jego s primarny. Zapis ukazuje tu razlikovateljnu ulogu, a ne novu kovanicu. Definicija raboty 31, §2, razgleda prsten konečnogo ranga i izključaje jediničny ideal; ostatkovo polje togda jest konečno algebraično razširjenje obraza $P$. Vrsta togo razširjenja oprědělja rod ideala. Sojuz i v zaglavju nabraja dve kategorije: ne znači, že jedin ideal jest jednovremenno prvogo i drugogo roda. Takože konečny rang nad poljem ne znači konečno čislo elementov. Te uslovja sut čestju objasnjenja, hoti ne vse vidime v kratkom nazvanju.\par
\noindent\emph{Reason:} Prosty ideal prvogo i drugogo roda reserves prost- for prime and does not confuse it with primarny. The record identifies that discriminating role rather than a new coinage. Paper 31, §2, considers a ring of finite rank and excludes the unit ideal; the residue field is then a finite algebraic extension of the image of $P$. The kind of this extension determines the ideal's kind. And in the heading lists two categories; it does not mean that one ideal is simultaneously of both kinds. Finite rank over a field also does not mean finitely many elements. These conditions belong in the explanation even though the short name does not display them.\par
\noindent{\footnotesize Izvorne zapisy / Original records: HIST-g124-001.}\par
\smallskip\noindent\emph{Sravnjenje s tekstom:} Rod ideala opira se na algebraičnom razširjenju ostatkovogo polja. Dve kategorije ne sut dva stepena, i zaglavno i ne pravi ih jednovremennymi vlastnostjami jednogo ideala.\par
\noindent\emph{Textual comparison:} The ideal's kind is determined through the residue-field extension. The two categories are not powers, and the heading's and does not make them simultaneous properties of one ideal.\par
\noindent{\footnotesize TEXT-COMP-443; pasaže / passages: 17275, 17595, 17597, 17599; \texttt{CORPUS\_FORM\_COMPARISONS.json}.}\par
\end{minipage}\par
\par\medskip\Needspace{7\baselineskip}
\noindent\begin{minipage}{\linewidth}
\hypertarget{TERM-1363}{}
\noindent\textbf{prosty ideal $p$ iz multiplikacijskogo kolca $o$}\quad{\footnotesize\texttt{TERM-1363}}\par
\noindent\emph{Iztočnik / Source:} Primideal $p$ aus einem Multiplikationsring $o$\par
\noindent\emph{Rusky / Russian:} простой идеал $p$ из мультипликационного кольца $o$\par
\noindent\emph{Ukrajinsky / Ukrainian:} простий ідеал $p$ з мультиплікаційного кільця $o$\par
\smallskip\noindent\emph{Razlog:} Prosty ideal $p$ iz multiplikacijskogo kolca $o$ prodolžaje ranějši slovnik raboty 31. Arhivni $p$ i $o$ sut obične bukvy; aktualny teorem upotrěbljaje simvoly prostogo ideala i osnovnogo prstena v frakturě i slovo prsten. Zato $p$ ne směje byti prěčitano kako prosto čislo samo zaradi običnogo šrifta. Teorem govori o idealu v osnovnom multiplikacijskom prstenu i o njegovom razširjenju v poredok. Arhiv dava kontinuitet kako razlog, ale njegov izraz kanoničny opisuje projektov izbor, ne obšču jezyčnu sankciju. Izměna kolco na prsten trěbuje takože kontekstualnu prověrku soglasovanja; samo zaměnjenje imenika ne ustanavlja pravilnost cele frazy.\par
\noindent\emph{Reason:} Prosty ideal $p$ iz multiplikacijskogo kolca $o$ continues Paper 31's earlier vocabulary. The archive uses plain $p$ and $o$; the current theorem uses fraktur symbols for the prime ideal and base ring, and the noun prsten. Plain typography must therefore not turn $p$ into a prime number. The theorem concerns an ideal in the base multiplication ring and its extension to an order. The record cites continuity, but its word canonical describes a project choice, not general linguistic endorsement. Replacing kolco with prsten also requires contextual agreement checks; substituting the noun alone does not establish the correctness of the whole phrase.\par
\noindent{\footnotesize Izvorne zapisy / Original records: HIST-g163-002.}\par
\smallskip\noindent\emph{Sravnjenje s tekstom:} Obči teorem čuva osnovny prsten, razširjeny ideal, prinajmenje jednu komponentu i disjunkciju. Prosta komponenta drugogo roda ostaje prosta, ne vlastno primarna.\par
\noindent\emph{Textual comparison:} The general theorem retains the base ring, extended ideal, at least one component and the disjunction. A prime component of the second kind remains prime, not properly primary.\par
\noindent{\footnotesize TEXT-COMP-444; pasaže / passages: 17275, 19071; \texttt{CORPUS\_FORM\_COMPARISONS.json}.}\par
\end{minipage}\par
\par\medskip\Needspace{7\baselineskip}
\noindent\begin{minipage}{\linewidth}
\hypertarget{TERM-1364}{}
\noindent\textbf{prosty ideal drugogo roda}\quad{\footnotesize\texttt{TERM-1364}}\par
\noindent\emph{Iztočnik / Source:} Primideal zweiter Art\par
\noindent\emph{Rusky / Russian:} простой идеал второго рода\par
\noindent\emph{Ukrajinsky / Ukrainian:} простий ідеал другого роду\par
\smallskip\noindent\emph{Razlog:} Prosty ideal drugogo roda jest druga alternativa v obščem diskriminantnom teoremu raboty 31. Zapis izrično čuva paralelu s prvim rodom; ruska i ukrajinska paralela imajut isto razlikovanje. Ideal ostaje prosty, no njegovo ostatkovo polje jest neseparabilno razširjenje danogo osnovnogo polja. To ne jest novo nazvanje vlastno primarnogo ideala. Teorem koristi ili i prinajmenje jeden: jedna vlastno primarna komponenta abo jedna prosta komponenta drugogo roda dostačuje za navedenu preprěku. Prěvod, ktory by skratil obe alternative do ne-prosty komponent, izměnil by matematičny smysl.\par
\noindent\emph{Reason:} Prosty ideal drugogo roda is the second alternative in Paper 31's general discriminant theorem. The record expressly preserves its parallel with first kind; Russian and Ukrainian make the same distinction. The ideal remains prime, but its residue field is an inseparable extension of the specified base field. This is not another name for a properly primary ideal. The theorem uses or and at least one: one properly primary component or one prime component of the second kind suffices for the stated obstruction. Compressing both alternatives to a non-prime component would change the mathematics.\par
\noindent{\footnotesize Izvorne zapisy / Original records: HIST-g163-008.}\par
\smallskip\noindent\emph{Sravnjenje s tekstom:} Obči teorem čuva osnovny prsten, razširjeny ideal, prinajmenje jednu komponentu i disjunkciju. Prosta komponenta drugogo roda ostaje prosta, ne vlastno primarna.\par
\noindent\emph{Textual comparison:} The general theorem retains the base ring, extended ideal, at least one component and the disjunction. A prime component of the second kind remains prime, not properly primary.\par
\noindent{\footnotesize TEXT-COMP-444; pasaže / passages: 17275, 19071; \texttt{CORPUS\_FORM\_COMPARISONS.json}.}\par
\end{minipage}\par
\par\medskip\Needspace{7\baselineskip}
\noindent\begin{minipage}{\linewidth}
\hypertarget{TERM-1365}{}
\noindent\textbf{proste idealy}\quad{\footnotesize\texttt{TERM-1365}}\par
\noindent\emph{Iztočnik / Source:} Primideale\par
\noindent\emph{Rusky / Russian:} простые идеалы\par
\noindent\emph{Ukrajinsky / Ukrainian:} прості ідеали\par
\smallskip\noindent\emph{Razlog:} Proste idealy čuva množnu formu raněje oprěděljenogo prosty ideal. V zapisu raboty 31, §3, razlog jest kontinuitet, osoblivo pri kontrakciji idealov. Dostupny dokaz prěhodit od prědstavjenja nulovogo ideala v razširjenom prstenu k prěsěkam prostyh idealov s izhodnym prstenom. Te prěsěki ostavajut proste idealy. Zato argument prěnosi polnu razložimost od razširjenogo prstena k izhodnomu; obratno napravljenje ne jest dano tym samym korakom. Množina ne trěbuje nove jezyčnoj politike, ale čitatelj mora viděti, ktera konstrukcija čuva prostost.\par
\noindent\emph{Reason:} Proste idealy preserves the plural of the previously defined prosty ideal. The Paper 31, §3 record gives continuity as its reason, particularly under contraction. The available proof passes from a representation of the zero ideal in an extended ring to intersections of its prime ideals with the original ring. Those intersections remain prime. The argument thus transfers complete reducibility from the extended ring to the original ring; the reverse direction is not supplied by the same step. The plural needs no new lexical policy, but the reader must see which construction preserves primality.\par
\noindent{\footnotesize Izvorne zapisy / Original records: HIST-g126-006.}\par
\smallskip\noindent\emph{Sravnjenje s tekstom:} Prěsěki idealov razširjenogo prstena s izhodnym prstenom sut kontrakcije. Dokaz koristi to napravljenje čuvanja prostosti.\par
\noindent\emph{Textual comparison:} Intersections of ideals of the extended ring with the original ring are contractions. The proof uses that direction of preservation of primality.\par
\noindent{\footnotesize TEXT-COMP-445; pasaže / passages: 17665, 17666, 17667, 17668, 17669, 17670, 17671, 17672, 17673; \texttt{CORPUS\_FORM\_COMPARISONS.json}.}\par
\end{minipage}\par
\par\medskip\Needspace{7\baselineskip}
\noindent\begin{minipage}{\linewidth}
\hypertarget{TERM-1366}{}
\noindent\textbf{proste idealy}\quad{\footnotesize\texttt{TERM-1366}}\par
\noindent\emph{Iztočnik / Source:} Primidealen\par
\noindent\emph{Rusky / Russian:} простые идеалы\par
\noindent\emph{Ukrajinsky / Ukrainian:} прості ідеали\par
\smallskip\noindent\emph{Razlog:} Proste idealy jest slovjansky oblik izbrany v rabotě 34, §22, s tvrdženjem o neposrědnoj razumlivosti. Ruska i ukrajinska paralela pokazujut soglasnost korenja, ale ne predstavljajut čitateljsky test ni polny pregled slovjanske terminologije. Matematičny odlomok koristi konečnu Abelovu grupu i karakteristiku polja, ktora ne děli red grupy. Poslě razširjenja koeficientov polinomy $x_i^{h_i}-e$ imajut različne linearne faktory, a definirajuči ideal razpada se na različne proste idealy. Različnost i uslovje o karakteristice sut nositeli argumenta. Pridavnik prosty jih ne vključaje avtomatično.\par
\noindent\emph{Reason:} Proste idealy is the Slavic form chosen in Paper 34, §22, with a claim of immediate intelligibility. Russian and Ukrainian show agreement of stem, but are neither a reader test nor a full survey of Slavic terminology. The mathematical passage uses a finite abelian group and a field characteristic not dividing its order. After coefficient extension, the polynomials $x_i^{h_i}-e$ have distinct linear factors, and the defining ideal decomposes into distinct prime ideals. Distinctness and the characteristic hypothesis carry the argument; the adjective prime does not include them automatically.\par
\noindent{\footnotesize Izvorne zapisy / Original records: HIST-g200-007.}\par
\smallskip\noindent\emph{Sravnjenje s tekstom:} Uslovje na karakteristiku polja pozvoljaje različne linearne faktory i odgovarjajuče različne proste idealy. Termin ne zaměnjaje to uslovje.\par
\noindent\emph{Textual comparison:} The field-characteristic condition permits distinct linear factors and corresponding distinct prime ideals. The term cannot substitute for that condition.\par
\noindent{\footnotesize TEXT-COMP-446; pasaže / passages: 22006, 22007, 22008, 22009, 22010, 22011, 22012, 22013, 22015, 22016, 22017, 22018, 22019, 22020, 22021, 22022, 22023, 22024, 22025, 22026, 22027, 22028, 22029, 22030, 22031, 22032, 22033, 22034; \texttt{CORPUS\_FORM\_COMPARISONS.json}.}\par
\end{minipage}\par
\par\medskip\Needspace{7\baselineskip}
\noindent\begin{minipage}{\linewidth}
\hypertarget{TERM-1367}{}
\noindent\textbf{primitivne idempotentne elementy}\quad{\footnotesize\texttt{TERM-1367}}\par
\noindent\emph{Iztočnik / Source:} primitive idempotente Elemente\par
\noindent\emph{Rusky / Russian:} примитивные идемпотентные элементы\par
\noindent\emph{Ukrajinsky / Ukrainian:} примітивні ідемпотентні елементи\par
\smallskip\noindent\emph{Razlog:} Primitivne idempotentne elementy čuva dva internacionalne termina s različnymi ulogami. Idempotentny govori o množenju elementa so soboju; primitivny tu označaje komponenty jedinice, povezane s prostymi jednostrannymi idealami pri direktnom razloženju. To jest zapisany razlog, ne značenje primitivny kako nerazvity. Rabota 32 dava razloženje $S=rS+(e-r)S$, kde $e$ jest jedinica; strana idealov ostaje važna. Uslovje $r\ne0,1$ ne garantuje primitivnost v vsakom proizvoljnom prstenu: posebny zaključok teksta odnosi se k kvaternionnomu slučaju s indeksom dva. Ruska i ukrajinska paralela zadrživajut isti internacionalny par; izvorna literatura citovana v odlomku ne byla tu nanovo nezavisno prověrjena.\par
\noindent\emph{Reason:} Primitivne idempotentne elementy retains two international terms with different roles. Idempotent describes multiplication of an element by itself; primitive here marks the components of the identity associated with simple one-sided ideals in a direct decomposition. That is the recorded reason, not primitive in the sense of undeveloped. Paper 32 gives $S=rS+(e-r)S$, with $e$ the identity; the side of the ideals remains important. The condition $r\ne0,1$ does not guarantee primitivity in every arbitrary ring: the text's special conclusion concerns the quaternion case with index two. Russian and Ukrainian retain the same international pair; the literature cited in the passage was not independently rechecked here.\par
\noindent{\footnotesize Izvorne zapisy / Original records: HIST-g173-002.}\par
\smallskip\noindent\emph{Sravnjenje s tekstom:} Komponenty jedinice svazany sut s jednostrannym razloženjem. Specialny kvaternionny slučaj ne jest obča definicija primitivnosti.\par
\noindent\emph{Textual comparison:} Components of the identity are linked to a one-sided decomposition. The special quaternion case is not a general definition of primitivity.\par
\noindent{\footnotesize TEXT-COMP-447; pasaže / passages: 19445, 19446, 19447, 19448, 19449, 19450, 19451, 19452, 19453, 19457; \texttt{CORPUS\_FORM\_COMPARISONS.json}.}\par
\end{minipage}\par
\par\medskip\Needspace{7\baselineskip}
\noindent\begin{minipage}{\linewidth}
\hypertarget{TERM-1368}{}
\noindent\textbf{primitivny $p$-ty korenj iz jedinice}\quad{\footnotesize\texttt{TERM-1368}}\par
\noindent\emph{Iztočnik / Source:} primitive $p$-te Einheitswurzel\par
\noindent\emph{Rusky / Russian:} первообразный $p$-й корень из единицы\par
\noindent\emph{Ukrajinsky / Ukrainian:} примітивний $p$-й корінь з одиниці\par
\smallskip\noindent\emph{Razlog:} Primitivny $p$-ty korenj iz jedinice prodolžaje ranějšu frazu korenje iz jedinice; sam zapis označuje to kako izbor strojenogo matematičnogo registra. Ruska paralela koristi pervoobraznyj, ukrajinska primityvnyj: ne možno predstaviti je kako tri identične izbory. Primitivny označaje točny multiplikativny red $p$, ne prosto proizvoljny korenj jednačenja $z^p=1$. Rabota 32 izbira prosto čislo $p$ s kongruencijnym uslovjem na $2^r+1$ i potom koristi korenj $\varepsilon$ v normnom argumentu za razpadno polje kvaternionnoj grupy. Indeks $p$ i stepenj $2^n$ razgledanogo podpolja imajut različne ulogi; ne sut dva zapisa iste veličiny.\par
\noindent\emph{Reason:} Primitivny $p$-ty korenj iz jedinice continues the earlier root-of-unity phrase; the record itself labels this a constructed-register choice. Russian uses pervoobraznyj, Ukrainian primityvnyj: they cannot be presented as three identical choices. Primitive specifies exact multiplicative order $p$, not an arbitrary solution of $z^p=1$. Paper 32 selects a prime $p$ satisfying a congruence condition on $2^r+1$, then uses a root $\varepsilon$ in a norm argument for a splitting field of the quaternion group. The index $p$ and the degree $2^n$ of the subfield have different roles; they are not two names for the same quantity.\par
\noindent{\footnotesize Izvorne zapisy / Original records: HIST-g175-003.}\par
\smallskip\noindent\emph{Sravnjenje s tekstom:} Korenj ima točny red $p$. Normny argument i izbor cikličnogo podpolja čuvajut osobne uslovja na $p$, $r$ i $n$.\par
\noindent\emph{Textual comparison:} The root has exact order $p$. The norm argument and choice of cyclic subfield retain separate conditions on $p$, $r$ and $n$.\par
\noindent{\footnotesize TEXT-COMP-448; pasaže / passages: 19604, 19608, 19610, 19612, 19614, 19616, 19617, 19618, 19619, 19620, 19622, 19624, 19625, 19626, 19628, 19630, 19632, 19634, 19636, 19638, 19640; \texttt{CORPUS\_FORM\_COMPARISONS.json}.}\par
\end{minipage}\par
\par\medskip\Needspace{7\baselineskip}
\noindent\begin{minipage}{\linewidth}
\hypertarget{TERM-1369}{}
\noindent\textbf{primitivne polinomy}\quad{\footnotesize\texttt{TERM-1369}}\par
\noindent\emph{Iztočnik / Source:} primitive Polynome\par
\noindent\emph{Rusky / Russian:} примитивные многочлены\par
\noindent\emph{Ukrajinsky / Ukrainian:} примітивні поліноми\par
\smallskip\noindent\emph{Razlog:} Primitivne polinomy jest kratky termin s izričnoju lokalnoju definicijeju: največši obči dělitelj koeficientov jest jedinica v osnovnom prstenu. Arhiv zato ne dodava parafrazu do samogo nazvanja; ruska i ukrajinska paralela takože koristajut primitiv- korenj. Rabota 30 izrično govori o dělitelju v polinomnom smyslu, ne idealnom, i dava osnovny prsten s jednoznačnym razloženjem. Primitivnost zato ne znači nerazložimost, moničnost ni minimalny stepenj. Primitivne imenitelji potom definirajut funkcionalnu oblast. Vynesti izraz iz takogo konteksta bez osnovnogo prstena moglo by izměniti njegov smysl.\par
\noindent\emph{Reason:} Primitivne polinomy is a short term with an explicit local definition: the greatest common divisor of the coefficients is a unit in the base ring. The record therefore adds no paraphrase to the name itself; Russian and Ukrainian also use the primitiv- stem. Paper 30 expressly means divisibility in the polynomial sense, not ideal divisibility, and supplies a base ring with unique factorization. Primitive therefore means neither irreducible, monic nor of least degree. Primitive denominators then define the functional domain. Removing the expression from this context without specifying the base ring could change its meaning.\par
\noindent{\footnotesize Izvorne zapisy / Original records: HIST-g110-009.}\par
\smallskip\noindent\emph{Sravnjenje s tekstom:} Največši obči dělitelj koeficientov jest jedinica v ukazanom osnovnom prstenu. To jest druga vlastnost od nerazložimosti i služi dlja izbora imeniteljev.\par
\noindent\emph{Textual comparison:} The coefficients' greatest common divisor is a unit in the specified base ring. This differs from irreducibility and supplies the denominator condition.\par
\noindent{\footnotesize TEXT-COMP-449; pasaže / passages: 16401, 16403, 16405, 16407, 16409, 16411; \texttt{CORPUS\_FORM\_COMPARISONS.json}.}\par
\end{minipage}\par
\par\medskip\Needspace{7\baselineskip}
\noindent\begin{minipage}{\linewidth}
\hypertarget{TERM-1370}{}
\noindent\textbf{primitivny element / primitivna funkcija}\quad{\footnotesize\texttt{TERM-1370}}\par
\noindent\emph{Iztočnik / Source:} primitives Element / primitive Funktion\par
\noindent\emph{Rusky / Russian:} примитивный элемент / примитивная функция\par
\noindent\emph{Ukrajinsky / Ukrainian:} примітивний елемент / примітивна функція\par
\smallskip\noindent\emph{Razlog:} Primitivny element / primitivna funkcija čuva nazvanje generatora polja, zato že dokaz raboty 9 svede pytanje k takim generatoram. Funkcija tu jest element funkcionalnogo polja, ne prvobrazna funkcija v diferencialnom računu. Tekst izbira medžupolja, ktoryh primitivny element jest racionalno izrazimy prěz $\vartheta$, i potom gradi polinomne izrazy. Dostupna pasaža o konečno mnogih medžupoljah vključaje takože konjugovane elementy; toj obseg ne směje byti skračen do proizvoljnogo konečnogo razširjenja. Zapis naglašaje primitivno-elementnu redukciju; novějše sravnjenje ne dokazuje, že vse potrebne uslovja byli prověrjene v momentu prvogo izbora.\par
\noindent\emph{Reason:} Primitivny element / primitivna funkcija preserves the name for a field generator because Paper 9 reduces its problem to such generators. Function here is an element of a function field, not an antiderivative. The text selects intermediate fields whose primitive element is rationally expressible in $\vartheta$, then constructs polynomial expressions. The available finite-intermediate-field passage includes conjugate elements; that scope must not be shortened to an arbitrary finite extension. The record emphasizes the primitive-element reduction; the present comparison does not prove that all necessary hypotheses were checked when the wording was first selected.\par
\noindent{\footnotesize Izvorne zapisy / Original records: HIST-g065-011.}\par
\smallskip\noindent\emph{Sravnjenje s tekstom:} Primitivny element generuje ukazano polje i mora byti racionalno izrazimy prěz dane veličiny. Dostupny odlomok s konečno mnogimi medžupoljami vključaje konjugovane elementy.\par
\noindent\emph{Textual comparison:} The primitive element generates the specified field and must be rationally expressible in the given quantities. The available finite-intermediate-field passage includes conjugate elements.\par
\noindent{\footnotesize TEXT-COMP-450; pasaže / passages: 7903, 7904, 7905, 7906, 7907, 7908, 7909, 7910, 7911, 7913, 7914, 7915, 7916, 7917, 7918, 7919, 7920, 7921, 7922, 7923, 7924, 7925, 7926, 7927, 7928, 7929, 7930, 7931, 7932, 7933, 7934, 7935, 7936; \texttt{CORPUS\_FORM\_COMPARISONS.json}.}\par
\end{minipage}\par
\par\medskip\Needspace{7\baselineskip}
\noindent\begin{minipage}{\linewidth}
\hypertarget{TERM-1371}{}
\noindent\textbf{primitivny element kone?nogo polja raz?irjenja [Čestično nezachovane bukvy / Partly unrecoverable lettering]}\quad{\footnotesize\texttt{TERM-1371}}\par
\noindent\emph{Iztočnik / Source:} primitives Element eines endlichen Erweiterungsk?rpers [Čestično nezachovane bukvy / Partly unrecoverable lettering]\par
\noindent\emph{Rusky / Russian:} [Bukvy ne sut zachovane / Lettering not recoverable]\par
\noindent\emph{Ukrajinsky / Ukrainian:} [Bukvy ne sut zachovane / Lettering not recoverable]\par
\smallskip\noindent\emph{Razlog:} Arhivno zaglavje čestično jest poškodženo: ono zachovava primitivny element i polja, ale ne vse bukvy; ruska i ukrajinska polja ne zachovajut čitljive bukvy. Dostupny tekst raboty 9 dozvaljaje nezavisno objasniti koncept, ne obnoviti izginuly zapis. Za vsako polje $\mathfrak K$, konečno nad $(H)$, trěbuje se konečno razširjenje $\mathfrak L$ nad $\mathfrak K$, čiji primitivny element leži v $[H,\mathfrak S]$. To poslednje vključenje jest imenno zapisany razlog za čuvanje kvalifikacije. Izbor ne kaže samo, že nekde obstaja primitivny element; on mora byti v ukazanoj oblasti. Vlastnost II ostaje osobno uslovje konstrukcije, a ne avtomatičny poslědok adjunkcije proizvoljnogo sistema.\par
\noindent\emph{Reason:} The archived headword is partly damaged: it preserves primitivny element and polja but not every letter; the Russian and Ukrainian fields retain no readable lettering. The available Paper 9 text permits an independent explanation of the concept, not recovery of the lost record. For every field $\mathfrak K$ finite over $(H)$, the text requires a finite extension $\mathfrak L$ of $\mathfrak K$ whose primitive element lies in $[H,\mathfrak S]$. This inclusion is precisely the recorded reason for retaining the qualification. The choice does not merely say that a primitive element exists somewhere: it must belong to the specified domain. Property II remains a separate condition of the construction, not an automatic consequence of adjoining an arbitrary system.\par
\noindent{\footnotesize Izvorne zapisy / Original records: HIST-g063-011.}\par
\smallskip\noindent\emph{Sravnjenje s tekstom:} Egzistencija primitivnogo elementa ne jest cělo uslovje: tekst trěbuje njegovo vključenje v $[H,\mathfrak S]$ i pravilny red kvantifikatorov nad razširjenjami.\par
\noindent\emph{Textual comparison:} Existence of a primitive element is not the whole condition: the text requires membership in $[H,\mathfrak S]$ and the correct quantifier order over extensions.\par
\noindent{\footnotesize TEXT-COMP-451; pasaže / passages: 7680, 7681, 7682, 7683, 7684, 7685, 7686, 7687, 7688, 7689, 7690, 7691, 7692, 7693, 7694, 7695, 7696, 7697; \texttt{CORPUS\_FORM\_COMPARISONS.json}.}\par
\end{minipage}\par
\par\medskip\Needspace{7\baselineskip}
\noindent\begin{minipage}{\linewidth}
\hypertarget{TERM-1372}{}
\noindent\textbf{prvotno polje}\quad{\footnotesize\texttt{TERM-1372}}\par
\noindent\emph{Iztočnik / Source:} Primkörper\par
\noindent\emph{Rusky / Russian:} простое поле\par
\noindent\emph{Ukrajinsky / Ukrainian:} первинне поле\par
\smallskip\noindent\emph{Razlog:} Prvotno polje jest projektov izbor za polje porodženo jediničnym elementom. Arhiv izrično dava jemu prednost prěd prosto polje, aby razlikovati strukturu polja od izraza prosta karakteristika. Rusky zapis ima prostoe pole, ukrajinsky pervynne pole; ukrajinsky i medžuslovjansky izbor ostali označeni za autoritativnu jezyčnu prověrku. Tekst rabot 9 i 10 definira objekt operacijami nad jediniceju: nastava tip racionalnyh čisel ili ostatkovyh klasov modulo prosto $p$. Ni starost polja, ni prostost algebraičnogo razširjenja ne sut smyslom prvotno. Novy slovnik čuva predběžnost izbora; ne predstavlja internu prevenciju dvosmysla kako dokazan obči slovjansky standard. Rabota 22 v dostupnom tekstu vsejedno ima prosto polje za isti objekt; to jest konkretno nesoglasje korpusa, ne dokaz, že jedno iz imen jest oficialno.\par
\noindent\emph{Reason:} Prvotno polje is a project choice for the field generated by the multiplicative identity. The archive explicitly prefers it to prosto polje to distinguish the field structure from prosta karakteristika. Russian records prostoe pole, Ukrainian pervynne pole; the Ukrainian and Interslavic choices remained marked for authoritative language review. Papers 9 and 10 define the object through operations on the identity, producing the rational-number type or residue classes modulo a prime $p$. Neither antiquity of the field nor simplicity of an algebraic extension is meant by prvotno. The dictionary retains the provisional status: an internal attempt to avoid ambiguity is not established pan-Slavic standard usage. The available Paper 22 text nevertheless uses prosto polje for the same object: a concrete corpus inconsistency, not proof that either name is official.\par
\noindent{\footnotesize Izvorne zapisy / Original records: HIST-g068-002, HIST-g071-011.}\par
\smallskip\noindent\emph{Sravnjenje s tekstom:} Raboty 9 i 10 definirajut prvotno polje iz jedinice. Rabota 22 koristi prosto polje za isti objekt; to konkretno različje ostaje vidimo, dokole jezyčna prednost ne jest ustanovljena novym dokazom.\par
\noindent\emph{Textual comparison:} Papers 9 and 10 define prvotno polje from the identity. Paper 22 uses prosto polje for the same object; that concrete variation remains visible without claiming new evidence of linguistic superiority.\par
\noindent{\footnotesize TEXT-COMP-452; pasaže / passages: 8338, 8339, 8340, 8341, 8342, 8343, 8344, 8345, 8346, 8347, 8348, 8349, 8350, 8351, 8353, 8354, 8355, 8356, 8357, 8358, 8359, 8360, 8361, 8362, 8363, 8572, 14827; \texttt{CORPUS\_FORM\_COMPARISONS.json}.}\par
\end{minipage}\par
\par\medskip\Needspace{7\baselineskip}
\noindent\begin{minipage}{\linewidth}
\hypertarget{TERM-1373}{}
\noindent\textbf{prosty modul}\quad{\footnotesize\texttt{TERM-1373}}\par
\noindent\emph{Iztočnik / Source:} Primmodul\par
\noindent\emph{Rusky / Russian:} простой модуль\par
\noindent\emph{Ukrajinsky / Ukrainian:} простий модуль\par
\smallskip\noindent\emph{Razlog:} Prosty modul prěnosi Primmodul podle izričnoj definicije raboty 17: jedini dělitelji sut sam modul i jediničny modul vsih polinomov. Arhiv bira prost- paralelno s ruskim i ukrajinskim, ale privjazuje slovo k toj definiciji, ne k libovoljnomu pozdnějšemu pojmu prostogo ideala. V historijnoj děliteljskoj terminologiji trěba čitati napravljenje vključenja; dělitelj ne jest prosto synonym slova podmodul bez objasnjenja. Grupa ostatkov potom nazyvaje se prosta grupa. Čuvanje kratkogo nazvanja zato trěbuje jasne definicije obojih objektov; ne opravdava preměnu modula v ideal ni odebranje operatornoj struktury.\par
\noindent\emph{Reason:} Prosty modul translates Primmodul according to Paper 17's explicit definition: its only divisors are itself and the unit module of all polynomials. The record chooses prost- in parallel with Russian and Ukrainian but ties the word to that definition, not an arbitrary later notion of prime ideal. Historical divisor terminology requires attention to inclusion direction; divisor cannot simply be treated as a synonym for submodule without explanation. The residue group is then called a simple group. Retaining the short name therefore requires clear definitions of both objects; it does not authorize turning the module into an ideal or removing its operator structure.\par
\noindent{\footnotesize Izvorne zapisy / Original records: HIST-g079-043.}\par
\smallskip\noindent\emph{Sravnjenje s tekstom:} Prosty modul jest oprěděljen děliteljami, a prosta grupa jeho ostatkovym objektom. Historijno nazvanje ostaje svazano s toj definicijeju.\par
\noindent\emph{Textual comparison:} The prime module is defined through divisors, and the simple group as its residue object. The historical name remains tied to that definition.\par
\noindent{\footnotesize TEXT-COMP-453; pasaže / passages: 11572, 11574; \texttt{CORPUS\_FORM\_COMPARISONS.json}.}\par
\end{minipage}\par
\par\medskip\Needspace{7\baselineskip}
\noindent\begin{minipage}{\linewidth}
\hypertarget{TERM-1374}{}
\noindent\textbf{primarny}\quad{\footnotesize\texttt{TERM-1374}}\par
\noindent\emph{Iztočnik / Source:} primär\par
\noindent\emph{Rusky / Russian:} примарный\par
\noindent\emph{Ukrajinsky / Ukrainian:} примарний\par
\smallskip\noindent\emph{Razlog:} Primarny čuva internacionalny primar- korenj vo vseh trěh arhivnyh paralelah raboty 19. Izričny razlog jest razlikovanje od prost- pri modernom algebraičnom izlaganju. Definicija III dopušča konečnu stepenj drugogo faktora, kogda produkt leži v idealu, a prvy faktor ne leži v njem; prosty slučaj ima stepenj jednu. Definicija IIIa prěnosi uslovje na produkty idealov pri danej konečnoj bazě. Teorem o nerazložimyh idealah dava primarnost, ale ne obrat bez dodatnyh uslovij. Zato primarny ne znači samo fundamentalny, prvy ili nerazložimy. Soglasnost zapisanyh paralel ne jest sama po sobě dokazom nezavisnoj konsultacije vseh jezyčnyh autoritetov.\par
\noindent\emph{Reason:} Primarny retains the international primar- stem in all three archived parallels for Paper 19. The explicit reason is differentiation from prost- in modern algebraic exposition. Definition III allows a finite power of the other factor when a product lies in the ideal but the first factor does not; the prime case has exponent one. Definition IIIa transfers the condition to ideal products with the available finite basis. The theorem on irreducible ideals yields primaryness, but not the converse without additional conditions. Primarny therefore means neither merely fundamental, first nor irreducible. Agreement among the recorded parallels is not by itself evidence of independent consultation of all linguistic authorities.\par
\noindent{\footnotesize Izvorne zapisy / Original records: HIST-g081-006, HIST-g082-006.}\par
\smallskip\noindent\emph{Sravnjenje s tekstom:} Primarnost dopušča stepenj; prostost jest osobny slučaj. Nerazložimost implicira primarnost pri danyh uslovjah, ale obrat jest izrično oprovrgnuty priměrom.\par
\noindent\emph{Textual comparison:} Primaryness allows a power; primality is a special case. Irreducibility implies primaryness under the stated conditions, but an explicit example refutes the converse.\par
\noindent{\footnotesize TEXT-COMP-454; pasaže / passages: 12785, 12787, 12789, 12791, 12834, 12838; \texttt{CORPUS\_FORM\_COMPARISONS.json}.}\par
\end{minipage}\par
\par\medskip\Needspace{7\baselineskip}
\noindent\begin{minipage}{\linewidth}
\hypertarget{TERM-1375}{}
\noindent\textbf{primarne idealy / primarny faktor}\quad{\footnotesize\texttt{TERM-1375}}\par
\noindent\emph{Iztočnik / Source:} primäre Ideale / primärer Faktor\par
\noindent\emph{Rusky / Russian:} примарные идеалы / примарный множитель\par
\noindent\emph{Ukrajinsky / Ukrainian:} примарні ідеали / примарний множник\par
\smallskip\noindent\emph{Razlog:} Primarne idealy / primarny faktor čuva primar- v rabotě 18 i jeho polinomnoj interpretaciji. Arhiv dava prednost internacionalnomu korenju prěd neimenovanym čisto slovjanskim izrazom, zaradi očekivanoj slabšej kanoničnosti poslednjego i rizika zaměšanja s prostym modulom. To jest kvalitativny sud; ne zachovano čisleno sravnjenje ni konkretna odbacena kovanica. V pasaži primarnomu idealu $Q$ odpovědaje polinomny faktor $Q^{(i)}$, s interpretacijeju normy komplementa po $Q$. Podobna bukva i isto pridavno imę ne činjut ideal i polinom jednym objektom. Matematična svez jest razlog za terminologičnu paralelu, ne za zaměnjenje jednoj kategorije drugoju.\par
\noindent\emph{Reason:} Primarne idealy / primarny faktor retains primar- in Paper 18 and its polynomial interpretation. The archive favours the international stem over an unnamed purely Slavic expression, anticipating weaker established usage and possible confusion with simple-module vocabulary. This is a qualitative judgement: neither numerical comparison nor a specific rejected coinage survives. In the passage, a primary ideal $Q$ corresponds to a polynomial factor $Q^{(i)}$, interpreted through the norm of the complement relative to $Q$. Similar lettering and the shared adjective do not make ideal and polynomial one object. The mathematical connection motivates terminological parallelism, not substitution of one category for the other.\par
\noindent{\footnotesize Izvorne zapisy / Original records: HIST-g080-011.}\par
\smallskip\noindent\emph{Sravnjenje s tekstom:} Normna interpretacija povezuje ideal s faktorom rezultantnoj formy. Ona ne utožsamjaje ideal s polinomom i čuva multiplicitet.\par
\noindent\emph{Textual comparison:} The norm interpretation links an ideal with a factor of the resultant form. It does not identify ideal and polynomial, and preserves multiplicity.\par
\noindent{\footnotesize TEXT-COMP-455; pasaže / passages: 12343, 12345, 12346, 12347, 12348, 12349, 12350, 12351, 12353, 12354, 12355, 12356, 12357, 12358, 12359, 12360, 12361, 12362, 12363, 12364; \texttt{CORPUS\_FORM\_COMPARISONS.json}.}\par
\end{minipage}\par
\par\medskip\Needspace{7\baselineskip}
\noindent\begin{minipage}{\linewidth}
\hypertarget{TERM-1376}{}
\noindent\textbf{primarne kolca}\quad{\footnotesize\texttt{TERM-1376}}\par
\noindent\emph{Iztočnik / Source:} primäre Ringe\par
\noindent\emph{Rusky / Russian:} первичные кольца\par
\noindent\emph{Ukrajinsky / Ukrainian:} первинні кільця\par
\smallskip\noindent\emph{Razlog:} Primarne kolca jest arhivna množna forma; aktualna pasaža raboty 31, §6, ima primarne prsteny. Zapis bira kontinuitet s teorijeju primarnyh prstenov i pripravjaje slednje izslědovanje, ne uvodi novy koncept. Uvod §6 svodi računanje diskriminanta na primarne komponenty razširjenogo prstena. To ne znači, že proizvoljno razširjenje primarnogo prstena ostaje primarno: ranějši §2 izrično dopušča protivny slučaj. Slědujuča konstrukcija specialnoj bazy koristi stepeni nilpotentnogo prostogo ideala pridruženogo nule. Isty pridavnik primarny povezuje ty odlomky, ale ne odstranja uslovja konečnogo ranga i ukazanogo osnovnogo polja.\par
\noindent\emph{Reason:} Primarne kolca is the archived plural form; the current passage in Paper 31, §6 uses primarne prsteny. The record chooses continuity with primary-ring theory and prepares the following discussion, rather than introducing a new concept. The §6 opening reduces the discriminant calculation to primary components of the extended ring. This does not mean that an arbitrary extension of a primary ring remains primary: §2 expressly allows the contrary. The following special-basis construction uses powers of the nilpotent prime associated with zero. The same adjective connects these passages but does not remove the finite-rank and base-field conditions.\par
\noindent{\footnotesize Izvorne zapisy / Original records: HIST-g148-005.}\par
\smallskip\noindent\emph{Sravnjenje s tekstom:} Prsten konečnogo ranga dopušča opis nulovyh děliteljev i razloženje. Razširjena komponenta ne musi ostati primarna; obrazovanje slov ne garantuje čuvanja vlastnosti.\par
\noindent\emph{Textual comparison:} Finite rank permits the zero-divisor description and decomposition. An extended component need not remain primary; word formation does not guarantee preservation of the property.\par
\noindent{\footnotesize TEXT-COMP-456; pasaže / passages: 17434, 17436, 17438, 17440, 17553, 17555, 17557, 17559, 17561, 17563, 17564, 17566, 18279, 18308, 18310; \texttt{CORPUS\_FORM\_COMPARISONS.json}.}\par
\end{minipage}\par
\par\medskip\Needspace{7\baselineskip}
\noindent\begin{minipage}{\linewidth}
\hypertarget{TERM-1377}{}
\noindent\textbf{primarno kolco}\quad{\footnotesize\texttt{TERM-1377}}\par
\noindent\emph{Iztočnik / Source:} primärer Ring\par
\noindent\emph{Rusky / Russian:} первичное кольцо\par
\noindent\emph{Ukrajinsky / Ukrainian:} первинне кільце\par
\smallskip\noindent\emph{Razlog:} Primarno kolco razširjaje slovnik primarnogo ideala na ostatkovy prsten, v ktorom nulove dělitelje sut nilpotentne. Oba arhivna razloga naglašajut kontinuitet i tehničnu stabilnost internacionalnogo korenja; ruska i ukrajinska paralela imajut pervičn-/pervyn- formy. Rabota 31 pracuje s konečnym rangom nad poljem i izrično dobiva takože, že vsaky element ostatkovogo prstena jest jedinica ili nulovy dělitelj. Ta dodatna vlastnost ne smeje byti vnesena v definiciju libovoljnogo prstena samo iz imena. Dostupny tekst ima primarno prsten, dokole arhiv ima primarno kolco; slovnik zato razdělja matematičny smysl, historijny izbor imenika i potrěbu osobnoj gramatičnoj prověrki aktualnogo soglasovanja.\par
\noindent\emph{Reason:} Primarno kolco extends primary-ideal vocabulary to a residue ring whose zero divisors are nilpotent. Both archived rationales emphasize continuity and the international stem's technical stability; Russian and Ukrainian use their pervichn-/pervyn- forms. Paper 31 works with finite rank over a field and also explicitly obtains that every element of the residue ring is a unit or a zero divisor. That additional property must not be inserted into a definition of an arbitrary ring merely from its name. The available text has primarno prsten whereas the archive has primarno kolco; the dictionary therefore separates mathematical meaning, historical noun choice and the need for a distinct grammatical check of current agreement.\par
\noindent{\footnotesize Izvorne zapisy / Original records: HIST-g121-006, HIST-g123-008.}\par
\smallskip\noindent\emph{Sravnjenje s tekstom:} Prsten konečnogo ranga dopušča opis nulovyh děliteljev i razloženje. Razširjena komponenta ne musi ostati primarna; obrazovanje slov ne garantuje čuvanja vlastnosti.\par
\noindent\emph{Textual comparison:} Finite rank permits the zero-divisor description and decomposition. An extended component need not remain primary; word formation does not guarantee preservation of the property.\par
\noindent{\footnotesize TEXT-COMP-456; pasaže / passages: 17434, 17436, 17438, 17440, 17553, 17555, 17557, 17559, 17561, 17563, 17564, 17566, 18279, 18308, 18310; \texttt{CORPUS\_FORM\_COMPARISONS.json}.}\par
\end{minipage}\par
\par\medskip\Needspace{7\baselineskip}
\noindent\begin{minipage}{\linewidth}
\hypertarget{TERM-1378}{}
\noindent\textbf{primarne faktory}\quad{\footnotesize\texttt{TERM-1378}}\par
\noindent\emph{Iztočnik / Source:} Primärfaktoren\par
\noindent\emph{Rusky / Russian:} первичные множители\par
\noindent\emph{Ukrajinsky / Ukrainian:} первинні множники\par
\smallskip\noindent\emph{Razlog:} Primarne faktory čuva konzervativnu polinomnu frazu raboty 22. Zapis izrično ne zaměnjaje ju asociovanymi prostymi idealami. Na koncu §5 tekst sam pojasnjaje, že faktorizaciju rezultantnogo faktora možno prodolžiti do stepenjev nerazložimyh polinomov, imenno primarnyh faktorov. Njih razloženju odpovědaje prědstavjenje osnovnogo ideala kako najmenšego občego mnogokratnogo. Faktor, ideal i pridruženy prosty ideal imaju različne ulogi; osuvremenjenje ne směje skryti stepenj ni multiplicitet. Ruska i ukrajinska paralela takože čuvajut faktor/primary izraz. Prednost věrnosti frazi jest zapisany razlog, ne dokaz, že vsaka pojavjenost slova faktor v korpusu ima toj smysl.\par
\noindent\emph{Reason:} Primarne faktory conservatively retains Paper 22's polynomial phrase. The record expressly does not replace it by associated prime ideals. At the end of §5, the text itself explains that factorization of a resultant factor can be continued to powers of irreducible polynomials, namely primary factors. Their decomposition corresponds to a representation of the fundamental ideal as a least common multiple. Factor, ideal and associated prime have different roles; modernization must not conceal the exponent or multiplicity. Russian and Ukrainian also retain the primary-factor expression. Fidelity to the phrase is the recorded reason, not proof that every occurrence of factor in the corpus has this sense.\par
\noindent{\footnotesize Izvorne zapisy / Original records: HIST-g089-011, HIST-g092-053.}\par
\smallskip\noindent\emph{Sravnjenje s tekstom:} Primarne faktory sut tu stepeni nerazložimyh polinomov. Njih idealna interpretacija ne zaměnjaje faktory asociovanymi prostymi idealami.\par
\noindent\emph{Textual comparison:} Primary factors here are powers of irreducible polynomials. Their ideal-theoretic interpretation does not replace the factors by associated primes.\par
\noindent{\footnotesize TEXT-COMP-457; pasaže / passages: 14913, 14914, 14915, 14916, 14917, 14918, 14919, 14920, 14921, 14922, 14923, 14925; \texttt{CORPUS\_FORM\_COMPARISONS.json}.}\par
\end{minipage}\par
\par\medskip\Needspace{7\baselineskip}
\noindent\begin{minipage}{\linewidth}
\hypertarget{TERM-1379}{}
\noindent\textbf{primarna funkcija}\quad{\footnotesize\texttt{TERM-1379}}\par
\noindent\emph{Iztočnik / Source:} Primärfunktion\par
\noindent\emph{Rusky / Russian:} первичная функция\par
\noindent\emph{Ukrajinsky / Ukrainian:} первинна функція\par
\smallskip\noindent\emph{Razlog:} Primarna funkcija jest polinomny analog primarnogo ideala v rabotě 24. Arhiv upira izbor na toj analogiji; on ne navodi samostojnu čisleno ocěnjenu konkurenciju kandidatov. Definicija kaže, že to jest stepenj proste, to jest nerazložimoj, funkcije nad ukazanym poljem. Stepenj jedna jest uključena; samo vlastna primarna funkcija izključaje prostu funkciju. Tako največši primarny faktory grupujut stepeni istogo nerazložimogo faktora, a produkt različnyh takyh faktorov više ne jest primarny. Prosta i primarna ne sut zaměnjajeme stilske varianty; primarna takože ne znači primitivna v smyslu koeficientov ili generatora polja.\par
\noindent\emph{Reason:} Primarna funkcija is the polynomial analogue of primary ideal in Paper 24. The archive bases the choice on that analogy; it gives no separate numerically scored competition. The definition makes it a power of a prime, meaning irreducible, function over the specified field. Exponent one is included; only properly primary function excludes the prime case. Maximal primary factors thus collect powers of the same irreducible factor, whereas a product of different such factors is no longer primary. Prime and primary are not interchangeable stylistic variants; primary also does not mean primitive in the coefficient or field-generator sense.\par
\noindent{\footnotesize Izvorne zapisy / Original records: HIST-g095-011.}\par
\smallskip\noindent\emph{Sravnjenje s tekstom:} Teorem IX razlikuje prostu funkciju od njezinoj stepeni i izrično ne izključaje eksponent jednu. Inkluzija specialnogo slučaja ne znači identične definicije.\par
\noindent\emph{Textual comparison:} Theorem IX distinguishes a prime function from its power and expressly leaves exponent one included. Inclusion of the special case does not mean identical definitions.\par
\noindent{\footnotesize TEXT-COMP-458; pasaže / passages: 15333, 15594, 15596, 15598, 15599, 15600, 15601, 15602, 15603, 15604, 15606; \texttt{CORPUS\_FORM\_COMPARISONS.json}.}\par
\end{minipage}\par
\par\medskip\Needspace{7\baselineskip}
\noindent\begin{minipage}{\linewidth}
\hypertarget{TERM-1380}{}
\noindent\textbf{primarna funkcija / prosta funkcija}\quad{\footnotesize\texttt{TERM-1380}}\par
\noindent\emph{Iztočnik / Source:} Primärfunktion / Primfunktion\par
\noindent\emph{Rusky / Russian:} первичная функция / простая функция\par
\noindent\emph{Ukrajinsky / Ukrainian:} первинна функція / проста функція\par
\smallskip\noindent\emph{Razlog:} Primarna funkcija / prosta funkcija čuva opoziciju potrěbnu v teoremu IX raboty 24. Zapis ne dopušča složenje obojih pojmov v jedin. Teorem veže primarnost ideala s primarnostju njegove formy elementarnyh děliteljev, i zato normy; dokaz piše $Q^{(i)}=(P^{(i)})^\rho$ s prostoju funkcijeju $P^{(i)}$. Zapis eksponenta jest sučestveny, ne ukras. Tekst izrično ostavlja $\rho=1$ i ne izključaje prosty ideal kako specialny slučaj. Silnějši normny kriterij za prostost trěbuje dodatne uslovja razgledane pozdněje. Zato slovnik čuva inkluziju prostogo slučaja v primarny bez utožsamjenja definicij.\par
\noindent\emph{Reason:} Primarna funkcija / prosta funkcija preserves the opposition needed in Theorem IX of Paper 24. The record does not permit merging the two notions. The theorem relates primaryness of an ideal to that of its elementary-divisor form, and therefore its norm; the proof writes $Q^{(i)}=(P^{(i)})^\rho$ with a prime function $P^{(i)}$. The exponent is substantive, not decoration. The text expressly retains $\rho=1$ and does not exclude a prime ideal as a special case. The stronger norm criterion for primality requires additional conditions considered later. The dictionary therefore preserves inclusion of the prime case within the primary case without identifying their definitions.\par
\noindent{\footnotesize Izvorne zapisy / Original records: HIST-g099-006.}\par
\smallskip\noindent\emph{Sravnjenje s tekstom:} Teorem IX razlikuje prostu funkciju od njezinoj stepeni i izrično ne izključaje eksponent jednu. Inkluzija specialnogo slučaja ne znači identične definicije.\par
\noindent\emph{Textual comparison:} Theorem IX distinguishes a prime function from its power and expressly leaves exponent one included. Inclusion of the special case does not mean identical definitions.\par
\noindent{\footnotesize TEXT-COMP-458; pasaže / passages: 15333, 15594, 15596, 15598, 15599, 15600, 15601, 15602, 15603, 15604, 15606; \texttt{CORPUS\_FORM\_COMPARISONS.json}.}\par
\end{minipage}\par
\par\medskip\Needspace{7\baselineskip}
\noindent\begin{minipage}{\linewidth}
\hypertarget{TERM-1381}{}
\noindent\textbf{primarny ideal}\quad{\footnotesize\texttt{TERM-1381}}\par
\noindent\emph{Iztočnik / Source:} Primärideal\par
\noindent\emph{Rusky / Russian:} первичный идеал\par
\noindent\emph{Ukrajinsky / Ukrainian:} первинний ідеал\par
\smallskip\noindent\emph{Razlog:} Primarny ideal bira internacionalny korenj primar-, dokolě ruska i ukrajinska paralela koristajut pervičn-/pervyn-. Arhiv pripisuje izboru točnost i medžuslovjansku razumlivost; toj zapis ne jest dokaz čitateljskogo testovanja ni slovničnogo priznanja. Rabota 24, §1, 11, razlikuje primarnost od prostoty pomoču stepeni faktora: iz $ab\in Q$ i $a\notin Q$ slěduje $b^k\in Q$ za někoj $k$. Prosty ideal jest specialny primarny slučaj, ne protivna kategorija. Slovo vlastny ograđuje one primarne idealy, ktore ne sut proste. Izbor pridavnogo imena mora čuvati tuto trojnu razliku.\par
\noindent\emph{Reason:} Primarny ideal selects the international stem primar-, whereas the recorded Russian and Ukrainian comparanda use pervichn-/pervyn-. The archive attributes precision and cross-Slavic intelligibility to this choice; that statement is neither a reader test nor evidence of dictionary recognition. Paper 24, §1, 11, distinguishes primary from prime through a power of a factor: from $ab\in Q$ and $a\notin Q$, some $b^k\in Q$ follows. Prime is a special primary case, not its opposite. Properly primary excludes that special case. The adjective must preserve all three distinctions.\par
\noindent{\footnotesize Izvorne zapisy / Original records: HIST-g095-007.}\par
\smallskip\noindent\emph{Sravnjenje s tekstom:} Definicija čuva stepen faktora i vključaje prosty ideal kako specialny slučaj. Vlastny izključaje tuty slučaj.\par
\noindent\emph{Textual comparison:} The definition retains the power of a factor and includes prime ideals as a special case. Properly primary excludes that case.\par
\noindent{\footnotesize TEXT-COMP-459; pasaže / passages: 15348, 15350, 15351, 15352, 15353, 15354, 15355, 15356, 15357; \texttt{CORPUS\_FORM\_COMPARISONS.json}.}\par
\end{minipage}\par
\par\medskip\Needspace{7\baselineskip}
\noindent\begin{minipage}{\linewidth}
\hypertarget{TERM-1382}{}
\noindent\textbf{primarny ideal / asociovany prosty ideal}\quad{\footnotesize\texttt{TERM-1382}}\par
\noindent\emph{Iztočnik / Source:} Primärideal / zugehöriges Primideal\par
\noindent\emph{Rusky / Russian:} первичный идеал / ассоциированный простой идеал\par
\noindent\emph{Ukrajinsky / Ukrainian:} первинний ідеал / асоційований простий ідеал\par
\smallskip\noindent\emph{Razlog:} Primarny ideal / asociovany prosty ideal čuva formalny odnos, ne samo obče značenje něčego povezanogo. To jest izričny razlog arhivnogo izbora asociovany vměsto neopreděljenogo opisa pripadanja. V rabotě 27 polje ostatkovyh klasov asociovanogo prostogo ideala davaje osnovu za čislo linearno nezavisnyh klasov. Zapis potom razgleđaje red kvocientov idealov i kompozicijny red; ne govori samo o čislu elementov konečnoj množiny. Paralelne ruske i ukrajinske formy takože označajut asociaciju. Čuvanje formalnogo nazvanja ne daje pravo zaměniti primarny ideal njegovym asociovanym prostym idealom.\par
\noindent\emph{Reason:} Primarny ideal / asociovany prosty ideal preserves a formal association, not merely a generic relationship or belonging. That is the archive's explicit reason for asociovany. In Paper 27 the residue field of the associated prime supplies the scalar field for counting linearly independent residue classes. The passage then considers ideal quotients and a composition series; it is not simply counting the elements of a finite set. Russian and Ukrainian comparanda likewise mark association. Preserving the formal name does not authorize replacing the primary ideal by its associated prime.\par
\noindent{\footnotesize Izvorne zapisy / Original records: HIST-g104-005.}\par
\smallskip\noindent\emph{Sravnjenje s tekstom:} Asociovany prosty ideal oprěděljaje polje ostatkovyh klasov. Čislo linearno nezavisnyh klasov ne jest obče čislo elementov.\par
\noindent\emph{Textual comparison:} The associated prime supplies the residue field. The number of linearly independent classes is not an unrestricted count of elements.\par
\noindent{\footnotesize TEXT-COMP-460; pasaže / passages: 15917; \texttt{CORPUS\_FORM\_COMPARISONS.json}.}\par
\end{minipage}\par
\par\medskip\Needspace{7\baselineskip}
\noindent\begin{minipage}{\linewidth}
\hypertarget{TERM-1384}{}
\noindent\textbf{produkt diskriminantnyh idealov komponentov}\quad{\footnotesize\texttt{TERM-1384}}\par
\noindent\emph{Iztočnik / Source:} Produkt der Diskriminantenideale der Komponenten\par
\noindent\emph{Rusky / Russian:} произведение дискриминантных идеалов компонент\par
\noindent\emph{Ukrajinsky / Ukrainian:} добуток дискримінантних ідеалів компонент\par
\smallskip\noindent\emph{Razlog:} Produkt diskriminantnyh idealov komponentov čuva množenje idealov; ne zaměnjaje je direktnoju sumoju, hoti dokaz počinaje od direktnoj sumy komponentov. Oba arhivne razloga naglašajut tuty razny operacije i birajut produkt kako internacionalny matematičny izraz. V rabotě 31, §5, blokova determinantna forma daje produkt, pri čem vsaky diskriminantny ideal jest vzety po odnošenju k svojemu komponentnomu klasu. §6 koristi to razlikovanje pri jediničnyh i nulovyh diskriminantnyh idealah. Izpuščenje odnosnogo klasa ili zaměna produkta sumoju by proměnili tvrdženje. Prozirnost slova ostavaje arhivno razsudženje, ne izměrena vlastnost.\par
\noindent\emph{Reason:} Produkt diskriminantnyh idealov komponentov preserves ideal multiplication, not a direct sum, although the proof begins with a direct-sum decomposition. Both archived reasons stress this distinction and choose produkt as a mathematical internationalism. In Paper 31, §5, the block determinant yields a product, with each component's discriminant ideal taken relative to its own component class. Section 6 uses that distinction for unit and zero discriminant ideals. Omitting the reference class or replacing the product by a sum would change the assertion. Transparency remains the archive's judgment, not a measured property.\par
\noindent{\footnotesize Izvorne zapisy / Original records: HIST-g139-005, HIST-g145-005.}\par
\smallskip\noindent\emph{Sravnjenje s tekstom:} Blokovy determinant daje produkt, ne sumu diskriminantnyh idealov. Odnos k komponentnym klasam jest čest tvrdženja.\par
\noindent\emph{Textual comparison:} The block determinant gives a product, not a sum of discriminant ideals. Reference to the component classes is part of the assertion.\par
\noindent{\footnotesize TEXT-COMP-462; pasaže / passages: 18187, 18189, 18190, 18191, 18192, 18193, 18194, 18195, 18196, 18197, 18198, 18199, 18200, 18201, 18202, 18203, 18204, 18205, 18206, 18207, 18208, 18209, 18210, 18211, 18212, 18213, 18214, 18215, 18441; \texttt{CORPUS\_FORM\_COMPARISONS.json}.}\par
\end{minipage}\par
\par\medskip\Needspace{7\baselineskip}
\noindent\begin{minipage}{\linewidth}
\hypertarget{TERM-1385}{}
\noindent\textbf{produkt konečno mnogih poparno vzaimno prostyh idealov}\quad{\footnotesize\texttt{TERM-1385}}\par
\noindent\emph{Iztočnik / Source:} Produkt von endlich vielen paarweise teilerfremden Idealen\par
\noindent\emph{Rusky / Russian:} произведение конечного числа попарно взаимно простых идеалов\par
\noindent\emph{Ukrajinsky / Ukrainian:} добуток скінченно багатьох попарно взаємно простих ідеалів\par
\smallskip\noindent\emph{Razlog:} Produkt konečno mnogih poparno vzaimno prostyh idealov čuva dva bitna uslovja: konečnost čisla faktorov i prostotu vsakoj pary. Arhiv naglašaje, že teorem XV raboty 19 jest teorem o produktu, ne samo daljši teorem o najmenšem občem mnogokratnom. Aktualny tekst piše konečno mnogyh i vzajemno, a polno tvrdženje dodava, že faktory sut vzajemno-prosto-nerazložime. To poslědnje uslovje ne směje izčeznuti zato, že kratko slovnično zaglavje go ne sodrživaje. Produkt i najmenše obče mnogokratno tu sut ravne po dokazanoj poparnoj prostotě, ne po značenju samyh slov.\par
\noindent\emph{Reason:} Produkt konečno mnogih poparno vzaimno prostyh idealov retains finiteness and coprimeness of every pair. The archive emphasizes that Paper 19's Theorem XV concerns a product, not merely another least-common-multiple theorem. The current text spells the forms konečno mnogyh and vzajemno, and the full theorem also requires the factors to be coprime-indecomposable. That last qualification must not disappear because the abbreviated dictionary headword omits it. Product and ideal least common multiple coincide here by the proved coprimeness condition, not because the words mean the same thing.\par
\noindent{\footnotesize Izvorne zapisy / Original records: HIST-g081-068, HIST-g082-068.}\par
\smallskip\noindent\emph{Sravnjenje s tekstom:} Polny teorem dodava vzajemno-prosto-nerazložimost faktorov. Kratko zaglavje ne smeje zameniti vse uslovja teorema.\par
\noindent\emph{Textual comparison:} The full theorem also requires coprime-indecomposable factors. The short headword cannot replace the theorem's complete hypotheses.\par
\noindent{\footnotesize TEXT-COMP-463; pasaže / passages: 13312, 13316, 13330, 13331, 13332, 13333, 13334, 13335, 13336, 13337; \texttt{CORPUS\_FORM\_COMPARISONS.json}.}\par
\end{minipage}\par
\par\medskip\Needspace{7\baselineskip}
\noindent\begin{minipage}{\linewidth}
\hypertarget{TERM-1386}{}
\noindent\textbf{produkt konečno mnogih po parah vzajemno prostyh primarnyh idealov}\quad{\footnotesize\texttt{TERM-1386}}\par
\noindent\emph{Iztočnik / Source:} Produkt von endlich vielen paarweise teilerfremden Primäridealen\par
\noindent\emph{Rusky / Russian:} произведение конечного числа попарно взаимно простых первичных идеалов\par
\noindent\emph{Ukrajinsky / Ukrainian:} добуток скінченно багатьох попарно взаємно простих первинних ідеалів\par
\smallskip\noindent\emph{Razlog:} Produkt konečno mnogih po parah vzajemno prostyh primarnyh idealov prodolžaje terminologiju raboty 31, §2; to jest zapisany razlog za po parah i primarnyh. V §7 toj izraz priměnjaje se k idealam poredka $T$, različnym od nulovogo ideala. To ograničenje jest bitno: rabota už prědtym razgleđa produktno razloženje nulovogo ideala v drugom, konečnorangovom kontekstu. Pridružene proste idealy faktorov v §7 nemajut vlastnogo dělitelja različnogo od jediničnogo. Konečnost odnosi se k čislu faktorov, ne k čislu elementov ideala. Paralela s §2 opravdava kontinuitet nazvanja, ne izpuščenje novih uslovij.\par
\noindent\emph{Reason:} Produkt konečno mnogih po parah vzajemno prostyh primarnyh idealov continues Paper 31, §2 terminology; that is the recorded reason for po parah and primarnyh. Section 7 applies the phrase to ideals of an order $T$ other than the zero ideal. This exclusion matters: the paper previously decomposes the zero ideal in a different, finite-rank setting. The associated primes in §7 have no proper divisor other than the unit ideal. Finiteness concerns the number of factors, not the number of elements of an ideal. Continuity of vocabulary does not license omission of the new hypotheses.\par
\noindent{\footnotesize Izvorne zapisy / Original records: HIST-g150-004.}\par
\smallskip\noindent\emph{Sravnjenje s tekstom:} V §7 izključeny jest nulovy ideal. Ranějše razloženje nulovogo ideala v konečnorangovom prstenu ne odstranjaje tuto ograničenje.\par
\noindent\emph{Textual comparison:} Section 7 excludes the zero ideal. Its earlier decomposition in a finite-rank ring does not remove this restriction.\par
\noindent{\footnotesize TEXT-COMP-464; pasaže / passages: 17440, 18559, 18561, 18593; \texttt{CORPUS\_FORM\_COMPARISONS.json}.}\par
\end{minipage}\par
\par\medskip\Needspace{7\baselineskip}
\noindent\begin{minipage}{\linewidth}
\hypertarget{TERM-1387}{}
\noindent\textbf{produkt idealov}\quad{\footnotesize\texttt{TERM-1387}}\par
\noindent\emph{Iztočnik / Source:} Produkt von Idealen / Definition IIIa\par
\noindent\emph{Rusky / Russian:} произведение идеалов\par
\noindent\emph{Ukrajinsky / Ukrainian:} добуток ідеалів\par
\smallskip\noindent\emph{Razlog:} Produkt idealov jest izbran, da by čitatelj viděl točno tvorjenje ideala iz produktov elementov. Arhivna odkazka k definiciji IIIa raboty 19 vodi k formulaciji primarnosti prěz idealy; sama konstrukcija produkta jest objasnjena v prědhodečej podnožnoj primětke. Ako $A$ i $B$ imajut baze $a_i$ i $b_k$, produkty $a_i b_k$ porodžajut $AB$. Zato ide o ideal, ktory oni porodžajut, ne samo o množinu pojedinyh produktov bez zamknjenja. Primětka izrično koristi komutativnost. Ruska i ukrajinska paralela ostavajut nazvanjami idealnogo produkta; nikaky konkurentny termin ni čislena vaga ne sut zapisane.\par
\noindent\emph{Reason:} Produkt idealov was chosen to retain the construction of an ideal from products of elements. The archival reference to Paper 19's Definition IIIa points to primaryness formulated with ideals; the product itself is explained in the preceding footnote. If $A$ and $B$ have bases $a_i$ and $b_k$, the products $a_i b_k$ generate $AB$. This is the ideal they generate, not merely the set of individual products without closure. The footnote explicitly uses commutativity. Russian and Ukrainian comparanda likewise name ideal multiplication; no competing term or numerical weight is recorded.\par
\noindent{\footnotesize Izvorne zapisy / Original records: HIST-g081-042, HIST-g082-042.}\par
\smallskip\noindent\emph{Sravnjenje s tekstom:} Podnožna primětka daje produkty bazovyh elementov kako bazu ideala. Definicija IIIa potom koristi takovo množenje.\par
\noindent\emph{Textual comparison:} The footnote gives products of basis elements as generators of the ideal. Definition IIIa then uses that multiplication.\par
\noindent{\footnotesize TEXT-COMP-465; pasaže / passages: 12785, 12787, 12789, 12791, 12793, 12795; \texttt{CORPUS\_FORM\_COMPARISONS.json}.}\par
\end{minipage}\par
\par\medskip\Needspace{7\baselineskip}
\noindent\begin{minipage}{\linewidth}
\hypertarget{TERM-1389}{}
\noindent\textbf{produktno predstavjenje}\quad{\footnotesize\texttt{TERM-1389}}\par
\noindent\emph{Iztočnik / Source:} Produktdarstellung\par
\noindent\emph{Rusky / Russian:} представление в виде произведения\par
\noindent\emph{Ukrajinsky / Ukrainian:} добуткове подання / подання у вигляді добутку; подання у вигляді добутку\par
\smallskip\noindent\emph{Razlog:} Produktno predstavjenje bira kratky internacionalny pridavnik za prědstavjenje kako produkt. Arhiv takože sodrživaje važnu rusku alternativu: proizvedennoe predstavlenie bylo odbačeno v korist predstavlenie v vide proizvedenija; ranějša pojavnost ostala je kandidat za prověrku, ne dokazano popravjena. Toj kontrast pokazuje razliku medžu prevedenym matematičnym smyslom i mehaničnoju kalkoju. Aktualny medžuslovjansky tekst piše prědstavjenje. V rabotě 19 pojavjaje se i neobstajanje takogo prědstavjenja; v rabotě 31 razgleđa se produktno razloženje nulovogo ideala. Termin sam ne tvrdži egzistenciju, jedinost ni faktorizačny teorem bez uslovij.\par
\noindent\emph{Reason:} Produktno predstavjenje selects a compact international adjective for representation as a product. The archive also preserves a meaningful Russian alternative: proizvedennoe predstavlenie was rejected for predstavlenie v vide proizvedenija; an earlier occurrence remained a consistency candidate, not an established correction. This contrast illustrates semantic translation rather than mechanical calquing. The current Interslavic text writes prědstavjenje. Paper 19 also discusses failure of such a representation; Paper 31 decomposes the zero ideal as a product. The term itself asserts neither existence nor uniqueness nor an unconditional factorization theorem.\par
\noindent{\footnotesize Izvorne zapisy / Original records: HIST-g081-091, HIST-g123-005.}\par
\smallskip\noindent\emph{Sravnjenje s tekstom:} Tekst razlikuje obstajanje i neobstajanje produktnogo prědstavjenja. Konečnorangovo razloženje nulovogo ideala ima svoje uslovja.\par
\noindent\emph{Textual comparison:} The text distinguishes existence from failure of product representation. The finite-rank decomposition of the zero ideal has its own hypotheses.\par
\noindent{\footnotesize TEXT-COMP-467; pasaže / passages: 13543, 13643, 13658, 17469, 17470, 17471, 17472, 17473, 17474, 17475, 17476, 17477, 17478; \texttt{CORPUS\_FORM\_COMPARISONS.json}.}\par
\end{minipage}\par
\par\medskip\Needspace{7\baselineskip}
\noindent\begin{minipage}{\linewidth}
\hypertarget{TERM-1390}{}
\noindent\textbf{produktno predstavjenje jego nulovogo ideala}\quad{\footnotesize\texttt{TERM-1390}}\par
\noindent\emph{Iztočnik / Source:} Produktdarstellung seines Nullideals\par
\noindent\emph{Rusky / Russian:} представление его нулевого идеала в виде произведения\par
\noindent\emph{Ukrajinsky / Ukrainian:} подання його нульового ідеалу у вигляді добутку\par
\smallskip\noindent\emph{Razlog:} Produktno predstavjenje jego nulovogo ideala prodolžaje izraz iz §2 pri definiciji polnoj razložimosti v rabotě 31, §3. To jest cěly zapisany razlog; ne imajemo nezavisnu novu ocěnku drugih kandidatov. Jego odnosi se k razgleđanomu prstenu. Produktno prědstavjenje nulovogo ideala $ (0)=q_1\cdots q_r $ ne jest samo druga forma direktnoj sumy prstena: dve razloženja sut svezane konstrukcijoju, ale imaju razne objekty i operacije. Polna razložimost dalje trěbuje, že vse primarne komponenty budut proste. Slovničny opis čuva taj uslovny prěhod vměsto izjednačenja primarnosti s prostotoju.\par
\noindent\emph{Reason:} Produktno predstavjenje jego nulovogo ideala carries §2 vocabulary into Paper 31, §3's definition of complete reducibility. That is the complete recorded rationale; there is no separate new assessment of alternative candidates. Jego refers to the ring under discussion. The zero-ideal product $ (0)=q_1\cdots q_r $ is not simply another spelling of the ring's direct sum: a construction relates the two decompositions, but their objects and operations differ. Complete reducibility further requires all primary components to be prime. The explanation preserves that conditional step instead of equating primaryness with primality.\par
\noindent{\footnotesize Izvorne zapisy / Original records: HIST-g125-002.}\par
\smallskip\noindent\emph{Sravnjenje s tekstom:} Produkt idealov i direktna suma prstena imajut razne objekty. Polna razložimost trěbuje prostotu vsih primarnyh komponentov.\par
\noindent\emph{Textual comparison:} The product of ideals and the direct sum of the ring concern different objects. Complete reducibility requires every primary component to be prime.\par
\noindent{\footnotesize TEXT-COMP-461; pasaže / passages: 17434, 17436, 17438, 17440, 17469, 17470, 17471, 17472, 17473, 17474, 17475, 17476, 17477, 17478, 17479, 17480, 17630, 17632; \texttt{CORPUS\_FORM\_COMPARISONS.json}.}\par
\end{minipage}\par
\par\medskip\Needspace{7\baselineskip}
\noindent\begin{minipage}{\linewidth}
\hypertarget{TERM-1391}{}
\noindent\textbf{produkty}\quad{\footnotesize\texttt{TERM-1391}}\par
\noindent\emph{Iztočnik / Source:} Produkte\par
\noindent\emph{Rusky / Russian:} произведения\par
\noindent\emph{Ukrajinsky / Ukrainian:} добутки\par
\smallskip\noindent\emph{Razlog:} Produkty jest mnoga forma internacionalnogo produkt, izbrana v arhivu kako prozirna matematična leksika pored ruskogo proizvedenija i ukrajinskogo dobutky. Ne jest zapisana čislena ocěnka razumlivosti. V rabotě 2, §§25–26, ide o produkty form, s kotorymi se poslě činjat kontrakcije i redukcije. Množenje formy i izvršenje kontrakcije ostavajut razne operacije. Oznaky II–III i III–III imenovajut sistemy form, ne odčitanje čisel; sravnenje produkta kako vyššego zavisi od ukazanogo poreda v simbolah. Kratko obče slovo zato ne može zaměniti strukturne oznaky ni njih lokalno značenje.\par
\noindent\emph{Reason:} Produkty is the plural of the international produkt, chosen in the archive as transparent mathematical vocabulary alongside Russian proizvedenija and Ukrainian dobutky. No numerical intelligibility assessment is recorded. In Paper 2, §§25–26, these are products of forms subsequently subjected to contractions and reductions. Multiplication of forms and contraction remain different operations. II–III and III–III identify systems of forms, not numerical subtraction; treating one product as higher depends on the stated order in the symbols. The short general word cannot replace those structural labels or their local interpretation.\par
\noindent{\footnotesize Izvorne zapisy / Original records: HIST-g026-008, HIST-g027-006.}\par
\smallskip\noindent\emph{Sravnjenje s tekstom:} Množenje form ne jest sama kontrakcija. Oznaky sistemov i uslovje poreda v simbolah nesut smysl sravnenja.\par
\noindent\emph{Textual comparison:} Multiplication of forms is not the contraction itself. System labels and the condition on symbolic order carry the comparison's meaning.\par
\noindent{\footnotesize TEXT-COMP-468; pasaže / passages: 3258, 3271, 3272, 3273, 3274, 3275, 3350, 3355, 3356, 3357, 3383, 3384, 3385, 3386, 3387, 3388, 3389, 3414; \texttt{CORPUS\_FORM\_COMPARISONS.json}.}\par
\end{minipage}\par
\par\medskip\Needspace{7\baselineskip}
\noindent\begin{minipage}{\linewidth}
\hypertarget{TERM-1392}{}
\noindent\textbf{produktna matrica}\quad{\footnotesize\texttt{TERM-1392}}\par
\noindent\emph{Iztočnik / Source:} Produktenmatrix\par
\noindent\emph{Rusky / Russian:} продуктовая матрица\par
\noindent\emph{Ukrajinsky / Ukrainian:} продуктова матриця\par
\smallskip\noindent\emph{Razlog:} Produktna matrica ostavlja svez medžu matriceju, množenjem i determinantom, na ktory opira se Laplaceovo razloženje v rabotě 4, §5. Arhivni razlog jest kontinuitet s §2, ne tvrdženje, že vsaky moderny smysl matričnogo produkta jest toj samy. §2 izrično definira svoj matričny produkt kako sumu produktov odgovarjajučih determinantov; §5 koristi determinant produktnoj matrice za formulu (32). Matrica ne jest svoj determinant, a suma minorov ne směje bez objasnjenja stati običnym zapisom množenja dvuh pravokutnyh matric. Ruska i ukrajinska arhivna pridavna imena sut sravniteljne formy, ne tu novo prověrene norme disciplinarnogo registra.\par
\noindent\emph{Reason:} Produktna matrica retains the connection between matrix, multiplication and determinant used by the Laplace expansion in Paper 4, §5. The recorded reason is continuity with §2, not an assertion that every modern meaning of matrix product is identical. Section 2 expressly defines its matrix product as a sum of products of corresponding determinants; §5 uses the product matrix's determinant for formula (32). A matrix is not its determinant, and a sum of minors must not silently become ordinary rectangular-matrix multiplication. The archived Russian and Ukrainian adjectives are comparanda, not newly verified disciplinary norms.\par
\noindent{\footnotesize Izvorne zapisy / Original records: HIST-g034-003.}\par
\smallskip\noindent\emph{Sravnjenje s tekstom:} Lokalna definicija matričnogo produkta i Laplaceovo razloženje trěbujut različenja matrice od jej determinanta.\par
\noindent\emph{Textual comparison:} The local matrix-product definition and Laplace expansion require distinguishing a matrix from its determinant.\par
\noindent{\footnotesize TEXT-COMP-469; pasaže / passages: 3782, 3805, 4239, 4240, 4241, 4242, 4243, 4244, 4245, 4246, 4247, 4248, 4249, 4250, 4251, 4252, 4253, 4254, 4255; \texttt{CORPUS\_FORM\_COMPARISONS.json}.}\par
\end{minipage}\par
\par\medskip\Needspace{7\baselineskip}
\noindent\begin{minipage}{\linewidth}
\hypertarget{TERM-1393}{}
\noindent\textbf{kolco produkta}\quad{\footnotesize\texttt{TERM-1393}}\par
\noindent\emph{Iztočnik / Source:} Produktenring\par
\noindent\emph{Rusky / Russian:} [Bukvy ne sut zachovane / Lettering not recoverable]\par
\noindent\emph{Ukrajinsky / Ukrainian:} [Bukvy ne sut zachovane / Lettering not recoverable]\par
\smallskip\noindent\emph{Razlog:} Kolco produkta jest starše nazvanje; aktualna rabota 40 piše prsten produkta. Arhiv povezuje internacionalny produkt s kolco i spominaje češsku, poljsku i vzhodnoslovjansku triangulaciju, ale sam zapis ne daje konkretne citate; poškodžene ruske i ukrajinske bukvy ne mogut byti domysljene. Izrična podnožna primětka definira objekt kako prsten porodženy iz dvoh prstenov v občem nadprstenu; produkt ne mora byti direktny. Teorem trěbuje elementnu komutativnost i egzistenciju odgovarjajučego prstena. Priměr pokazuje, že zadany prěsěk ne vsegda možno zachovati. Zato ni moderno direktno tvorjenje produkta ni tensorna konstrukcija ne mogut byti podstavjene samo po nazvanju.\par
\noindent\emph{Reason:} Kolco produkta is the older label; current Paper 40 writes prsten produkta. The archive combines international produkt with kolco and mentions Czech, Polish and East Slavic triangulation, but supplies no specific citations in this record; damaged Russian and Ukrainian lettering cannot be reconstructed by guesswork. An explicit footnote defines a ring generated by two rings in a common overring, not necessarily a direct product. The theorem requires elementwise commutation and existence of the relevant ring. A counterexample shows that a prescribed intersection need not be realizable. Neither a modern direct-product nor a tensor construction can therefore be substituted on the strength of the name.\par
\noindent{\footnotesize Izvorne zapisy / Original records: HIST-g210-003.}\par
\smallskip\noindent\emph{Sravnjenje s tekstom:} Izrična definicija ograničuje smysl produkta. Egzistencija i zadany prěsěk sut dodatne uslovja, ne avtomatične slědstva nazvanja.\par
\noindent\emph{Textual comparison:} The explicit definition fixes the product sense. Existence and a prescribed intersection are additional conditions, not automatic consequences of the name.\par
\noindent{\footnotesize TEXT-COMP-470; pasaže / passages: 23476, 23480, 23482, 23484, 23485, 23486, 23487, 23488, 23489, 23490, 23491, 23492, 23493, 23494, 23495, 23496; \texttt{CORPUS\_FORM\_COMPARISONS.json}.}\par
\end{minipage}\par
\par\medskip\Needspace{7\baselineskip}
\noindent\begin{minipage}{\linewidth}
\hypertarget{TERM-1394}{}
\noindent\textbf{Landauove produktne teoremy}\quad{\footnotesize\texttt{TERM-1394}}\par
\noindent\emph{Iztočnik / Source:} Produktsaetze von Landau\par
\noindent\emph{Rusky / Russian:} теоремы Ландау о произведениях\par
\noindent\emph{Ukrajinsky / Ukrainian:} добуткові теореми Ландау\par
\smallskip\noindent\emph{Razlog:} Landauove produktne teoremy čuva historijno imenovanu grupu teoremov. Arhiv izrično odklanja obči opis kako pravila faktorizacije, zato že rabota 17, §12, pokazuje granicu jednoproměnnovogo rezultata. V priměru s diferencialnymi operatorami od dvoh proměnnyh modul nema polinoma drugogo stepena, ale ima dva linearno nezavisne polinomy tretjego stepena; tekst iz toga zaključaje, že modul ne može iměti bazu iz jednogo polinoma. To jest ukazany razlog neuspěha prěnosa. Ne tvrdži se, že Landauove teoremy ne sut pravdive pri svojem jednoproměnnovom uslovju, ni že nikako drugo razloženje ne jest možno.\par
\noindent\emph{Reason:} Landauove produktne teoremy preserves a historically identified theorem family. The archive explicitly rejects the generic paraphrase factorization rules because Paper 17, §12, exposes the boundary of the one-variable result. In the two-variable differential-operator example, the module has no degree-two polynomial but has two linearly independent degree-three polynomials; the text concludes that it cannot have a one-polynomial basis. That is the stated obstruction to transfer. The passage does not say that Landau's theorems fail under their one-variable hypotheses or that every other decomposition is impossible.\par
\noindent{\footnotesize Izvorne zapisy / Original records: HIST-g079-075.}\par
\smallskip\noindent\emph{Sravnjenje s tekstom:} Dva nezavisne polinomy tretjego stepena i nedostajanje drugogo stepena podpirajut protivopriměr prěnosu. Oni ne opovrgajut izvorno jednoproměnnovo tvrdženje.\par
\noindent\emph{Textual comparison:} Two independent degree-three polynomials and the absence of degree two support the counterexample to transfer. They do not refute the original one-variable assertion.\par
\noindent{\footnotesize TEXT-COMP-471; pasaže / passages: 12209, 12210, 12211, 12212, 12213, 12214, 12215, 12216, 12217, 12218, 12219, 12220, 12221, 12222, 12223, 12224, 12225, 12226, 12227, 12228, 12229, 12230, 12231; \texttt{CORPUS\_FORM\_COMPARISONS.json}.}\par
\end{minipage}\par
\par\medskip\Needspace{7\baselineskip}
\noindent\begin{minipage}{\linewidth}
\hypertarget{TERM-1395}{}
\noindent\textbf{teorema produkta}\quad{\footnotesize\texttt{TERM-1395}}\par
\noindent\emph{Iztočnik / Source:} Produktsatz\par
\noindent\emph{Rusky / Russian:} теорема о произведении\par
\noindent\emph{Ukrajinsky / Ukrainian:} теорема про добуток\par
\smallskip\noindent\emph{Razlog:} Teorema produkta čuva nazvanje rezultata priměnjenogo v rabotě 4, §2; arhiv veže izbor k upotrěbljenju poslě formule (5). V čitanom tekstu rezultat davaje nulovy determinant reda $n-\lambda$ i vodi k relacijam (12). Zato teorema ne směje byti zaměnjena s neformalnym pravilom ili s pravidlom diferenciranja produkta. Kratko zaglavje samo ne identifikuje rezultat: potrěbne sut lokalne formule i značenje matričnogo produkta. Ruska i ukrajinska paralela takože čuvajut teoremu o produktu. Zapis ne sodrživaje daljšu konkurenciju kandidatov ni čislene vagi.\par
\noindent\emph{Reason:} Teorema produkta retains the name of the result invoked in Paper 4, §2; the archive ties it to the application following formula (5). The passage uses the result to obtain a zero determinant of order $n-\lambda$ and the relations (12). It must not become an informal rule or the differentiation rule for a product. The short headword alone does not identify the result: its local formulas and the specified matrix-product meaning are necessary. Russian and Ukrainian comparanda likewise retain a product theorem. No further candidate contest or numerical weighting survives.\par
\noindent{\footnotesize Izvorne zapisy / Original records: HIST-g031-013.}\par
\smallskip\noindent\emph{Sravnjenje s tekstom:} Priměnjenje zavisi od relacij (5) i vodi k (12). Kratko nazvanje ne jest samostojna identifikacija teoremy.\par
\noindent\emph{Textual comparison:} The application depends on relations (5) and leads to (12). The short name does not identify the theorem by itself.\par
\noindent{\footnotesize TEXT-COMP-472; pasaže / passages: 3859, 3860, 3861, 3862, 3863, 3864, 3865, 3866, 3867, 3868, 3869, 3870, 3871, 3872, 3873, 3874; \texttt{CORPUS\_FORM\_COMPARISONS.json}.}\par
\end{minipage}\par
\par\medskip\Needspace{7\baselineskip}
\noindent\begin{minipage}{\linewidth}
\hypertarget{TERM-1396}{}
\noindent\textbf{teorema produkta za determinanty}\quad{\footnotesize\texttt{TERM-1396}}\par
\noindent\emph{Iztočnik / Source:} Produktsatz für Determinanten\par
\noindent\emph{Rusky / Russian:} теорема произведения для детерминантов\par
\noindent\emph{Ukrajinsky / Ukrainian:} теорема добутку для детермінантів\par
\smallskip\noindent\emph{Razlog:} Teorema produkta za determinanty izbrana jest po analogiji s ranějšim nazvanjem teoremy produkta za matrice. Arhivny razlog jest terminologičny kontinuitet; ne sodrživaje osobno dokazanu ocěnku vsih možnyh prevodov. V rabotě 2, §5, teorema pomaga vyraziti polarne členy i izvěsti relacije prěz obrat sistema uravnenij. Za determinanty čuva objekt, k ktoromu se nazvanje odnosi; odstranjenje tyh slov skrylo by razliku od matričnoj formulacije. Izraz teorema ostavja dokazny rezultat, ne instrukciju samo umnožiti vidime faktory. Aktualna upotrěba jest doložena, ale ne daje sama po sebě nadjezyčnu normu.\par
\noindent\emph{Reason:} Teorema produkta za determinanty was selected by analogy with the earlier name for the matrix product theorem. The recorded reason is terminological continuity, not a separately established comparison of every possible translation. In Paper 2, §5, the theorem helps express polar terms and derive relations by reversing a system of equations. Za determinanty preserves the object of the theorem; omitting it would conceal the distinction from the matrix formulation. Teorema retains a proved result rather than an instruction merely to multiply visible factors. Current occurrence is established, but occurrence alone does not establish a language-wide norm.\par
\noindent{\footnotesize Izvorne zapisy / Original records: HIST-g006-007.}\par
\smallskip\noindent\emph{Sravnjenje s tekstom:} Uslovje (a), kontrakcije i relacije polarov davajut točny kontekst. Slovo teorema ne trěba svesti do neformalnogo pravila.\par
\noindent\emph{Textual comparison:} Condition (a), contractions and polar relations supply the precise context. Theorem should not be reduced to an informal rule.\par
\noindent{\footnotesize TEXT-COMP-473; pasaže / passages: 267, 268, 269, 270, 271, 272, 273, 274, 276, 278, 280, 281, 282, 283, 284, 285, 286, 287, 288, 289, 290, 291, 292, 700, 702, 703, 704, 705, 706, 707, 708, 709, 710, 711, 712, 713, 714, 715, 716, 717, 718, 719; \texttt{CORPUS\_FORM\_COMPARISONS.json}.}\par
\end{minipage}\par
\par\medskip\Needspace{7\baselineskip}
\noindent\begin{minipage}{\linewidth}
\hypertarget{TERM-1397}{}
\noindent\textbf{teorema produkta za determinanty / matrice}\quad{\footnotesize\texttt{TERM-1397}}\par
\noindent\emph{Iztočnik / Source:} Produktsatz für Determinanten / Matrizen\par
\noindent\emph{Rusky / Russian:} теорема о произведении для детерминантов / матриц\par
\noindent\emph{Ukrajinsky / Ukrainian:} теорема про добуток для детермінантів / матриць\par
\smallskip\noindent\emph{Razlog:} Teorema produkta za determinanty / matrice čuva dva svezane, ale različne opisy v rabotě 4, §3. Arhiv prodolžaje teorema produkta iz §2 i dodava točny objekt. Formula (14) davaje determinantnu identičnost; sistem (16) razširjaje matričnu formulaciju s uslovjem $2\le\rho\le n$ i s alternativnymi zapisami istogo izraza. Svez medžu formulami ne znači, že matrica i determinant sut sinonimy. Aktualny tekst takože sodrživaje redakcijnu podnožnu primětku o dopolnjenyh točkah; to jest osobny redakcijny sloj, ne avtorova terminologična odluka. Tuto sravnenje ne zaměnjaje prověrku originalnogo otiska.\par
\noindent\emph{Reason:} Teorema produkta za determinanty / matrice preserves related but distinct descriptions in Paper 4, §3. The archive continues teorema produkta from §2 and makes the object explicit. Formula (14) gives the determinant identity; system (16) develops the matrix formulation for $2\le\rho\le n$, with alternative expressions of the same quantity. Their relationship does not make matrix and determinant synonyms. The current text also contains an editorial footnote about supplied ellipsis dots; that is a separate editorial layer, not an authorial terminology decision. This comparison does not replace examination of the original printing.\par
\noindent{\footnotesize Izvorne zapisy / Original records: HIST-g032-006.}\par
\smallskip\noindent\emph{Sravnjenje s tekstom:} Determinantna identičnost i matričny sistem imajut točne formy i granice. Redakcijna primětka o točkah ostavaje osobnym slojem.\par
\noindent\emph{Textual comparison:} The determinant identity and matrix system have exact forms and bounds. The editorial note about ellipsis dots remains a separate layer.\par
\noindent{\footnotesize TEXT-COMP-474; pasaže / passages: 3907, 3908, 3909, 3910, 3911, 3912, 3913, 3914, 3924, 3925, 3926, 3927, 3928, 3929, 3930, 3931, 3932, 3933, 3934, 3935, 3936, 3937, 3938, 3939, 3940, 3941, 3942, 3943, 3944, 3945, 3946, 3947, 3948, 3949, 3950, 3951; \texttt{CORPUS\_FORM\_COMPARISONS.json}.}\par
\end{minipage}\par
\par\medskip\Needspace{7\baselineskip}
\noindent\begin{minipage}{\linewidth}
\hypertarget{TERM-1398}{}
\noindent\textbf{teorema produkta za matrice}\quad{\footnotesize\texttt{TERM-1398}}\par
\noindent\emph{Iztočnik / Source:} Produktsatz für Matrizen\par
\noindent\emph{Rusky / Russian:} теорема произведения для матриц\par
\noindent\emph{Ukrajinsky / Ukrainian:} теорема добутку для матриць\par
\smallskip\noindent\emph{Razlog:} Teorema produkta za matrice ostavlja rezultat vidimo teoremom, ne neformalnym pravilom produkta; to jest izričny arhivny motiv. Rabota 2, §1, priměnjaje go vměstě s teoremoju identičnosti i uslovjem (a), aby svesti kontrakcije III i IV k osnovnym I i II. Uslovje i identičnosti sut čest razsudženja, ne fakultativno okruženje nazvanja. V izrazah nastupajut simbolične faktory i kontrakcije, zato nazvanje ne směje izazvati zaměnu njihovogo smysla s libovoljnoju modernoj matričnoj operacijoju. Ruski i ukrajinski zapisy takože čuvajut imenovanje teoremy; razumlivost ne byla izměrena.\par
\noindent\emph{Reason:} Teorema produkta za matrice keeps the result visibly a theorem rather than an informal product rule; this is the explicit archival motivation. Paper 2, §1, uses it with the identity theorem and condition (a) to reduce contractions III and IV to the fundamental I and II. The condition and identities belong to the reasoning, not optional surroundings of the name. The expressions involve symbolic factors and contractions, so the label must not prompt substitution of an arbitrary modern matrix operation. Russian and Ukrainian records likewise name a theorem; intelligibility was not measured.\par
\noindent{\footnotesize Izvorne zapisy / Original records: HIST-g003-006.}\par
\smallskip\noindent\emph{Sravnjenje s tekstom:} Uslovje (a), kontrakcije i relacije polarov davajut točny kontekst. Slovo teorema ne trěba svesti do neformalnogo pravila.\par
\noindent\emph{Textual comparison:} Condition (a), contractions and polar relations supply the precise context. Theorem should not be reduced to an informal rule.\par
\noindent{\footnotesize TEXT-COMP-473; pasaže / passages: 267, 268, 269, 270, 271, 272, 273, 274, 276, 278, 280, 281, 282, 283, 284, 285, 286, 287, 288, 289, 290, 291, 292, 700, 702, 703, 704, 705, 706, 707, 708, 709, 710, 711, 712, 713, 714, 715, 716, 717, 718, 719; \texttt{CORPUS\_FORM\_COMPARISONS.json}.}\par
\end{minipage}\par
\par\medskip\Needspace{7\baselineskip}
\noindent\begin{minipage}{\linewidth}
\hypertarget{TERM-1399}{}
\noindent\textbf{teorema produkta v najobćejšej formě}\quad{\footnotesize\texttt{TERM-1399}}\par
\noindent\emph{Iztočnik / Source:} Produktsatz in allgemeinster Form\par
\noindent\emph{Rusky / Russian:} теорема о произведении в самой общей форме\par
\noindent\emph{Ukrajinsky / Ukrainian:} теорема про добуток у найзагальнішій формі\par
\smallskip\noindent\emph{Razlog:} Teorema produkta v najobćejšej formě čuva obobčenje, izrično svězane v arhivu s formuloju (32) raboty 4, §5. Aktualny tekst piše najobčejšej; slovnik razlikuje taj pravopisny kontrast od matematičnogo značenja. Obobčenje dopušča daljše razloženje redov simbolov i libovoljno mnogo redov, s točno danimi granicami sumiranja v (36). Najobčejša znači najobčejša v opisanoj klasi identičnostij, ne teorema bez kakihkoli uslovij. Cěly razsah, ograničenja indeksov i različne redy morajut ostati vidime. Arhiv ne predlagaje novy sinonim za produkt ni novy rezultat.\par
\noindent\emph{Reason:} Teorema produkta v najobćejšej formě retains the generalization explicitly linked in the archive to Paper 4, §5, formula (32). The current text spells najobčejšej; the dictionary separates that spelling contrast from mathematical meaning. The generalization permits further decomposition of symbol rows and arbitrarily many rows, with the exact summation bounds in (36). Most general means most general within the described class of identities, not a theorem without hypotheses. Scope, index restrictions and distinct rows must remain visible. The archive proposes neither a new synonym for product nor a new result.\par
\noindent{\footnotesize Izvorne zapisy / Original records: HIST-g034-002.}\par
\smallskip\noindent\emph{Sravnjenje s tekstom:} Obobčenje čuva redy simbolov i granice sumiranja. Ono ne izkazuje teoremu bez uslovij ili mimo zadanoj klasy.\par
\noindent\emph{Textual comparison:} The generalization retains symbol rows and summation bounds. It does not assert a theorem without hypotheses or beyond the specified class.\par
\noindent{\footnotesize TEXT-COMP-475; pasaže / passages: 4239, 4240, 4241, 4242, 4243, 4244, 4245, 4246, 4247, 4248, 4249, 4250, 4251, 4252, 4253, 4254, 4255, 4278, 4279, 4280, 4281, 4282, 4283, 4284, 4285, 4286, 4287, 4288, 4289, 4290, 4291, 4292, 4293, 4294, 4295, 4296, 4297, 4298, 4299, 4300; \texttt{CORPUS\_FORM\_COMPARISONS.json}.}\par
\end{minipage}\par
\par\medskip\Needspace{7\baselineskip}
\noindent\begin{minipage}{\linewidth}
\hypertarget{TERM-1400}{}
\noindent\textbf{projektivna invarianta}\quad{\footnotesize\texttt{TERM-1400}}\par
\noindent\emph{Iztočnik / Source:} projektive Invariante\par
\noindent\emph{Rusky / Russian:} проективный инвариант\par
\noindent\emph{Ukrajinsky / Ukrainian:} проєктивний інваріант\par
\smallskip\noindent\emph{Razlog:} Projektivna invarianta čuva internacionalny pridavnik i kratko matematično nazvanje. Arhiv opisuje paralelu s ruskim i ukrajinskim kako standardnu, ale sam toj zapis ne ustanavlja status v nezavisnom medžuslovjanskom kanonu. V rabotě 12 invarianta pri fiksovanom městě $x=(x)_0$ jest prěvedena na projektivnu invariantu sistema $F_\rho$, pri čem vyšše diferencialy zaměnjajut se kogredientnymi proměnnymi. Točno prěobraženje i fiksovano město noset značenje pridavnika. Ne ide o projektivny modul ni o čislo, ktoro ostavaje neproměnjeno pod vsimi možnymi prěobraženjami. Nazvanje mora byti čitano vměstě s transformacionnym zakonom.\par
\noindent\emph{Reason:} Projektivna invarianta retains an international adjective and a compact mathematical name. The archive describes the Russian and Ukrainian parallels as standard, but that record alone does not establish independent Interslavic canon status. In Paper 12 an invariant at a fixed point $x=(x)_0$ is related to a projective invariant of the system $F_\rho$, with higher differentials replaced by cogredient variables. The transformation and fixed point carry the adjective's meaning. This is neither a projective module nor a number unchanged under every possible transformation. The name must be read with its transformation law.\par
\noindent{\footnotesize Izvorne zapisy / Original records: HIST-g074-005.}\par
\smallskip\noindent\emph{Sravnjenje s tekstom:} Fiksovano město, kogredientne proměnne i linearno prěobraženje konkretizujut smysl pridavnika. Ne ide o projektivny modul.\par
\noindent\emph{Textual comparison:} The fixed point, cogredient variables and linear transformation specify the adjective's meaning. This is not a projective module.\par
\noindent{\footnotesize TEXT-COMP-476; pasaže / passages: 9130, 9131, 9132, 9133, 9134, 9135, 9136, 9137, 9138, 9139, 9140, 9141, 9142, 9143, 9144, 9145, 9146, 9147, 9148; \texttt{CORPUS\_FORM\_COMPARISONS.json}.}\par
\end{minipage}\par
\par\medskip\Needspace{7\baselineskip}
\noindent\begin{minipage}{\linewidth}
\hypertarget{TERM-1401}{}
\noindent\textbf{projektivna invariantna polja}\quad{\footnotesize\texttt{TERM-1401}}\par
\noindent\emph{Iztočnik / Source:} projektive Invariantenkörper\par
\noindent\emph{Rusky / Russian:} проективные поля инвариантов\par
\noindent\emph{Ukrajinsky / Ukrainian:} проєктивні поля інваріантів\par
\smallskip\noindent\emph{Razlog:} Projektivna invariantna polja označaje polja invariantov, ne polja, ktore samy sut nepodměnene pri někoj operaciji. Arhiv čuva projektivny kvalifikator i paralelne ruske i ukrajinske izrazy za polja invariantov; ne davaje čislove ocenky alternativ. Rabota 6, §1, sama objasnjaje razliku: pri razgledanyh transformacijah s determinantom 1 racionalne kombinacije invariantov ostavajut invariantami, tako že nastaje polje. Samy homogene izobarne invarianty, nasuprot tomu, tvorjat sistem, ne polje. Izbor polje zato nosi tvrđenje o algebraičnoj strukturě. Izpustiti pridavno ime invariantny ili zaměniti polje slovom sistem izgubilo by razliku, na kojoj stojit izloženje.\par
\noindent\emph{Reason:} Projektivna invariantna polja names fields of invariants, rather than fields merely described as unchanged by some operation. The preserved rationale retains the projective qualifier and Russian and Ukrainian expressions for invariant fields; it supplies no numerical ranking of alternatives. Paper 6, §1, makes the distinction explicit: for the determinant-one transformations considered there, rational combinations remain invariants and constitute a field. Homogeneous isobaric invariants alone instead form a system, not a field. Polje thus carries a claim about algebraic structure. Omitting invariantny or replacing polje with system would erase a distinction on which the exposition depends.\par
\noindent{\footnotesize Izvorne zapisy / Original records: HIST-g041-014.}\par
\smallskip\noindent\emph{Sravnjenje s tekstom:} Zamknjenost pri racionalnyh operacijah razlikuje polje od sistema homogennyh izobarnyh invariantov.\par
\noindent\emph{Textual comparison:} Closure under rational operations distinguishes the field from the system of homogeneous isobaric invariants.\par
\noindent{\footnotesize TEXT-COMP-477; pasaže / passages: 5130; \texttt{CORPUS\_FORM\_COMPARISONS.json}.}\par
\end{minipage}\par
\par\medskip\Needspace{7\baselineskip}
\noindent\begin{minipage}{\linewidth}
\hypertarget{TERM-1402}{}
\noindent\textbf{projektivna invariantna teorija}\quad{\footnotesize\texttt{TERM-1402}}\par
\noindent\emph{Iztočnik / Source:} projektive Invariantentheorie\par
\noindent\emph{Rusky / Russian:} проективная инвариантная теория\par
\noindent\emph{Ukrajinsky / Ukrainian:} проективна інваріантна теорія\par
\smallskip\noindent\emph{Razlog:} Projektivna invariantna teorija jest nazvanje discipliny, ne opis jednoho svojstva teorije. Oba zachovane razloga, iz rabot 3 i 4, voljajut prěmy matematičny izraz vměsto opisnoj parafrazy. Korpus potvrdžaje jeden i toj samy smysl: konečnost sistema form služi kako vyhodna točka, a pytanja o relacijah medžu formami i simboličnyh metodah sut daljši predmet. Ruska i ukrajinska paralela podpirajut takovo razčlenjenje nazvanja. Tuty povtorjeny zapis davaje kontinuitet upotrěby, ne dvě nezavisne lingvistične recenzije. Arhiv ne ustanovjaje, že red pridavnyh imen jest jediny dopustimy v Medžuslovjanskom.\par
\noindent\emph{Reason:} Projektivna invariantna teorija is a disciplinary name, not a description of a property possessed by a theory. Both preserved reasons, from Papers 3 and 4, prefer the direct mathematical expression to a descriptive paraphrase. The corpus confirms the same sense in both: finiteness of the system of forms is the starting point, followed by relations among forms and symbolic methods. The Russian and Ukrainian comparanda support this parsing of the name. The two records establish continuity of use, not two independent linguistic reviews. They do not establish this adjective order as the only permissible Interslavic formulation.\par
\noindent{\footnotesize Izvorne zapisy / Original records: HIST-g028-001, HIST-g029-002.}\par
\smallskip\noindent\emph{Sravnjenje s tekstom:} Oba uvoda koriste isto nazvanje za teoriju form i relacij medžu nimi.\par
\noindent\emph{Textual comparison:} Both introductions use the same name for the theory of forms and their relations.\par
\noindent{\footnotesize TEXT-COMP-478; pasaže / passages: 3647, 3699; \texttt{CORPUS\_FORM\_COMPARISONS.json}.}\par
\end{minipage}\par
\par\medskip\Needspace{7\baselineskip}
\noindent\begin{minipage}{\linewidth}
\hypertarget{TERM-1405}{}
\noindent\textbf{faktor proporcionalnosti}\quad{\footnotesize\texttt{TERM-1405}}\par
\noindent\emph{Iztočnik / Source:} Proportionalitätsfaktor\par
\noindent\emph{Rusky / Russian:} множитель пропорциональности\par
\noindent\emph{Ukrajinsky / Ukrainian:} множник пропорційності\par
\smallskip\noindent\emph{Razlog:} Faktor proporcionalnosti jest kratko sostavjeno nazvanje množitelja, ktory ustanavjaje proporciju. Arhiv nazyvaje go standardnym skalarnym izrazom i navodi rusku i ukrajinsku konstrukciju s množiteljem; ne dokumentuje sravniteljno testovanje. V rabotě 4, §1, redy determinantov sut normovane tako, že tuty faktor v formuli (1) stavaje ravnym 1. To ne jest proizvoljno izpuščenje množitelja, ale izrično vyslovjeno normovanje. Izbor mora čuvati razliku medžu občim faktorom proporcionalnosti i njegovoj cěnnostju po normovanju. Slovo koeficient moglo by opisati rolu v formuli, ale arhiv ne pokazuje, že ono bylo istorijskim kandidatom.\par
\noindent\emph{Reason:} Faktor proporcionalnosti is a compact compositional name for a proportionality multiplier. The record calls it standard scalar terminology and supplies Russian and Ukrainian constructions with 'multiplier'; it documents no comparative testing. In Paper 4, §1, the determinant rows are normalized so that this factor in formula (1) equals 1. The multiplier is not casually omitted: the normalization is expressly stated. The terminology must retain the difference between a general proportionality factor and its value after normalization. 'Coefficient' could describe its role in a formula, but the record does not show it was a historical candidate.\par
\noindent{\footnotesize Izvorne zapisy / Original records: HIST-g030-007.}\par
\smallskip\noindent\emph{Sravnjenje s tekstom:} Ravnost faktora jedinici jest poslědok normovanja, ne izpuščenogo uslovja.\par
\noindent\emph{Textual comparison:} The factor equals 1 by normalization, not by omission of a condition.\par
\noindent{\footnotesize TEXT-COMP-481; pasaže / passages: 3742, 3743, 3744, 3745, 3746, 3747, 3748, 3749, 3750, 3751, 3752, 3753; \texttt{CORPUS\_FORM\_COMPARISONS.json}.}\par
\end{minipage}\par
\par\medskip\Needspace{7\baselineskip}
\noindent\begin{minipage}{\linewidth}
\hypertarget{TERM-1406}{}
\noindent\textbf{pojem točke / grupa toček}\quad{\footnotesize\texttt{TERM-1406}}\par
\noindent\emph{Iztočnik / Source:} Punktbegriff / Punktgruppe\par
\noindent\emph{Rusky / Russian:} понятие точки / группа точек\par
\noindent\emph{Ukrajinsky / Ukrainian:} поняття точки / група точок\par
\smallskip\noindent\emph{Razlog:} Pojem točke / grupa toček čuva točku kako obči predmet aritmetičnogo i geometričnogo izloženja. To jest zachovany razlog prěmogo prěvoda. Rabota 14 sravnjuje aritmetično vvedenje točky s pojmom vzetym iz teorije funkcij, a potom razgleđaje koresidualne grupy toček. Tutaj grupa ne trěbuje značiti algebraičnu grupu s operacijeju: riječ označaje geometričny skup toček, ktory se sravnjuje s idealom i divizorom. Sam tekst govori, že koresidualne grupy ne muset iměti jednako mnogo toček. Pri prěnosu terminologije v drugu rabotu zato trěba ustanoviti smysl, ne avtomatično prěvesti vsako Punktgruppe kakoby šlo o abstraktnu grupu.\par
\noindent\emph{Reason:} Pojem točke / grupa toček retains the point as the common object of arithmetic and geometric exposition; that is the preserved reason for a direct rendering. Paper 14 compares an arithmetically introduced point with the point concept taken from function theory, then discusses coresidual groups of points. Grupa here need not denote an algebraic group equipped with an operation: it names a geometric collection compared with ideals and divisors. The passage expressly allows coresidual point groups to have unequal numbers of points. Terminology transfer must therefore establish the sense, not automatically treat every Punktgruppe as an abstract group.\par
\noindent{\footnotesize Izvorne zapisy / Original records: HIST-g076-015.}\par
\smallskip\noindent\emph{Sravnjenje s tekstom:} Geometrična grupa toček ne mora byti abstraktna grupa. Tekst sravnjuje točke, idealy i divizory.\par
\noindent\emph{Textual comparison:} A geometric group of points need not be an abstract group. The text compares points, ideals and divisors.\par
\noindent{\footnotesize TEXT-COMP-482; pasaže / passages: 9876, 9937, 10022, 10024; \texttt{CORPUS\_FORM\_COMPARISONS.json}.}\par
\end{minipage}\par
\par\medskip\Needspace{7\baselineskip}
\noindent\begin{minipage}{\linewidth}
\hypertarget{TERM-1407}{}
\noindent\textbf{kvadrat jedinice}\quad{\footnotesize\texttt{TERM-1407}}\par
\noindent\emph{Iztočnik / Source:} Quadrat einer Einheit\par
\noindent\emph{Rusky / Russian:} квадрат единицы\par
\noindent\emph{Ukrajinsky / Ukrainian:} квадрат одиниці\par
\smallskip\noindent\emph{Razlog:} Kvadrat jedinice čuva točno dopustimo različenje dvoh diskriminantov. Jedinica tutaj znači obratimy element osnovnogo prstena, ne obyčno čislo 1. V rabotě 31, §4, proměna bazy davaje kvadrat determinanta transformacije kako faktor; diskriminant jest zato oprěděljen do kvadrata jedinice. Izraz do jedinice byl by slabši i stratil by informaciju o kvadratu. Izraz ravny diskriminant byl by silnějši, než govori tekst. Ruska i ukrajinska paralela čuvajut tož ograničenje. Arhiv opravdava izbor točnostju relacije ekvivalentnosti, ne novostju korenja kvadrat.\par
\noindent\emph{Reason:} Kvadrat jedinice preserves the exact permitted discrepancy between two discriminants. Jedinica here means an invertible element of the base ring, not simply the number 1. In Paper 31, §4, a basis change contributes the square of its transformation determinant; the discriminant is accordingly determined up to a unit square. 'Up to a unit' would discard the square restriction, whereas equality would assert more than the passage does. The Russian and Ukrainian comparanda preserve the same qualification. The rationale concerns the precise equivalence relation, not novelty of the stem kvadrat.\par
\noindent{\footnotesize Izvorne zapisy / Original records: HIST-g137-003.}\par
\smallskip\noindent\emph{Sravnjenje s tekstom:} Proměna bazy množi diskriminant kvadratom jedinice; samo do jedinice ne čuva polnu informaciju.\par
\noindent\emph{Textual comparison:} A basis change multiplies the discriminant by a unit square; up to a unit loses information.\par
\noindent{\footnotesize TEXT-COMP-483; pasaže / passages: 18102, 18104; \texttt{CORPUS\_FORM\_COMPARISONS.json}.}\par
\end{minipage}\par
\par\medskip\Needspace{7\baselineskip}
\noindent\begin{minipage}{\linewidth}
\hypertarget{TERM-1408}{}
\noindent\textbf{kvadratične relacije}\quad{\footnotesize\texttt{TERM-1408}}\par
\noindent\emph{Iztočnik / Source:} quadratische Relationen\par
\noindent\emph{Rusky / Russian:} квадратичные соотношения\par
\noindent\emph{Ukrajinsky / Ukrainian:} квадратичні співвідношення\par
\smallskip\noindent\emph{Razlog:} Kvadratične relacije voljaje internacionalny korenj relacij- za algebraične identičnosti medžu elementami reda. Arhiv to naznačuje kako moderny algebraičny registr; rusky i ukrajinsky upotrěbjajut svoje izrazy za sootnošenija odnosno spivvidnošennja. V rabotě 4, §2, te relacije možno zapisati invariantno, i one odgovarjajut relacijam linearnym v koeficientah osnovnoj formy. Kvadratičny i linearny zato ne sut protivrěčne opisy: odnoset se k različnym argumentam. Prěvod mora čuvati, v čem jest relacija kvadratična i v čem linearna; ne može zaměniti oba opisa jednym neograničenym izrazom jednačenje drugogo stepena.\par
\noindent\emph{Reason:} Kvadratične relacije selects the international stem relacij- for algebraic identities among row elements. The record presents this as modern algebraic register, alongside the native Russian and Ukrainian relation terms. In Paper 4, §2, these relations admit an invariant expression and correspond to relations linear in the basic form's coefficients. Quadratic and linear are therefore not contradictory descriptions: they concern different arguments. Translation must retain what the relation is quadratic in and what it is linear in, rather than replacing both with an unqualified 'second-degree equation'.\par
\noindent{\footnotesize Izvorne zapisy / Original records: HIST-g031-010.}\par
\smallskip\noindent\emph{Sravnjenje s tekstom:} Relacije medžu elementami reda odgovarjajut relacijam linearnym v koeficientah osnovnoj formy.\par
\noindent\emph{Textual comparison:} Relations among row elements correspond to relations linear in the basic form's coefficients.\par
\noindent{\footnotesize TEXT-COMP-484; pasaže / passages: 3859, 3860, 3861, 3862, 3863, 3864, 3865, 3874; \texttt{CORPUS\_FORM\_COMPARISONS.json}.}\par
\end{minipage}\par
\par\medskip\Needspace{7\baselineskip}
\noindent\begin{minipage}{\linewidth}
\hypertarget{TERM-1409}{}
\noindent\textbf{kvadratične relacije / kvadratične formy}\quad{\footnotesize\texttt{TERM-1409}}\par
\noindent\emph{Iztočnik / Source:} quadratische Relationen / quadratische Formen\par
\noindent\emph{Rusky / Russian:} квадратичные соотношения / квадратичные формы\par
\noindent\emph{Ukrajinsky / Ukrainian:} квадратичні співвідношення / квадратичні форми\par
\smallskip\noindent\emph{Razlog:} Kvadratične relacije / kvadratične formy razlikuje uslovje ravnosti od polinoma, ktory go izražaje. V rabotě 15 formy $f_{ik,rs}$ v neizvěstnyh $\xi_{ik}$ sut polinomy drugogo stepena; po zaměně $\xi_{ik}$ determinantami $(ik)$ odgovarjajuče izrazy izčezajut identično. Te formy tvorjat modul, ktory služi dokazu. Arhiv izbral relacija za takovo algebraično uslovje, kako čest terminologije modula relacij. Formy i relacije zato ne sut dva nazvanja, ktore možno svobodno zaměnjati vo vsih větah. Zachovany razlog podpiraje razlikovanje objektov; ne davaje čislovoj oceny drugih korenjev.\par
\noindent\emph{Reason:} Kvadratične relacije / kvadratične formy distinguishes an equality condition from the polynomial expressing it. In Paper 15 the forms $f_{ik,rs}$ in the indeterminates $\xi_{ik}$ are quadratic polynomials; substituting the determinants $(ik)$ yields identically vanishing expressions. Those forms generate the module used in the proof. The preserved rationale selects relacija for this algebraic condition within relation-module terminology. Forms and relations are thus not freely interchangeable names in every sentence. The recorded reason supports an object distinction; it supplies no numerical evaluation of competing stems.\par
\noindent{\footnotesize Izvorne zapisy / Original records: HIST-g077-011.}\par
\smallskip\noindent\emph{Sravnjenje s tekstom:} Polinomy v neizvěstnyh poraždajut modul; po substituciji determinantov nastajut identičnosti.\par
\noindent\emph{Textual comparison:} Polynomials in indeterminates generate a module; substitution of determinants yields identities.\par
\noindent{\footnotesize TEXT-COMP-485; pasaže / passages: 10369, 10370, 10371, 10372, 10373, 10374, 10375, 10376, 10377, 10378, 10379, 10380, 10381, 10382, 10383; \texttt{CORPUS\_FORM\_COMPARISONS.json}.}\par
\end{minipage}\par
\par\medskip\Needspace{7\baselineskip}
\noindent\begin{minipage}{\linewidth}
\hypertarget{TERM-1410}{}
\noindent\textbf{kvazi-leksikografičny porjadok}\quad{\footnotesize\texttt{TERM-1410}}\par
\noindent\emph{Iztočnik / Source:} quasi-lexikographische Anordnung\par
\noindent\emph{Rusky / Russian:} квазилексикографический порядок\par
\noindent\emph{Ukrajinsky / Ukrainian:} квазілексикографічний порядок\par
\smallskip\noindent\emph{Razlog:} Kvazi-leksikografičny porjadok označaje osobny poredok konečnyh podmnožstv, ne samo obyčny slovničny red. Arhiv čuva internacionalny kvalifikator vměsto dlgoj parafrazy; nynějšny tekst ima poredok. Rabota 30, §6, najprvo sravnjuje čislo elementov; samo pri jednakej moči odlučuje prvi različny element v danym dobrom poredku. Po tom poredku bira se prva konečna baza vsakogo ideala. Prědpoloženje dobrogo poredka elementov i uslovje konečnoj bazy sut sučstvene: sam adjektiv ne dokazuje dobry poredok. Slovo kvazi- čuva razliku od neograničenogo leksikografičnogo sravnjenja; podrobna definicija mora ostati čitljiva pri vsakom prěnosu togo pojma.\par
\noindent\emph{Reason:} Kvazi-leksikografičny porjadok names a particular ordering of finite subsets, not merely ordinary dictionary order. The archive retains the international qualifier instead of a long paraphrase; the current text uses poredok. Paper 30, §6, first compares cardinalities and, only when these agree, compares the first differing element in the given well-order. The first finite basis of each ideal is then selected in this ordering. The assumed well-order of elements and finite generation matter: the adjective alone establishes no well-ordering theorem. Kvazi- preserves the distinction from unrestricted lexicographic comparison; the explicit definition must accompany any transfer of this concept.\par
\noindent{\footnotesize Izvorne zapisy / Original records: HIST-g113-006.}\par
\smallskip\noindent\emph{Sravnjenje s tekstom:} Konečne množiny najprvo sravnjujut se po moči. Definicija potom koristi dany dobry poredok elementov.\par
\noindent\emph{Textual comparison:} Finite sets are compared first by cardinality; the definition then uses the given well-order of elements.\par
\noindent{\footnotesize TEXT-COMP-486; pasaže / passages: 16678; \texttt{CORPUS\_FORM\_COMPARISONS.json}.}\par
\end{minipage}\par
\par\medskip\Needspace{7\baselineskip}
\noindent\begin{minipage}{\linewidth}
\hypertarget{TERM-1411}{}
\noindent\textbf{kvaternionno tělo}\quad{\footnotesize\texttt{TERM-1411}}\par
\noindent\emph{Iztočnik / Source:} Quaternionenkörper\par
\noindent\emph{Rusky / Russian:} тело кватернионов\par
\noindent\emph{Ukrajinsky / Ukrainian:} кватерніонне тіло\par
\smallskip\noindent\emph{Razlog:} Kvaternionno tělo čuva nekomutativny smysl německogo Körper. Arhiv izrično razlikuje tuty objekt od komutativnogo polja. V rabotě 32, §2, bazu nad racionalnym poljem tvorjat $1,i,j,k$, s $ij=k$ i $ji=-k$; zaměna tělo slovom polje izgubila by matematičnu razliku. Daljše prisojedinjenje razpadnogo polja može proizvesti netrivialny idempotent, tako že povědě o razširenoj algebrě ne trěba čitati kakoby ona ostavala tělo. Rusko tělo kvaternionov i ukrajinsko kvaternionne tilo v spisku podpirajut poimenovanje, ale ne sut nezavisno proverjena medžuslovjanska norma.\par
\noindent\emph{Reason:} Kvaternionno tělo preserves the noncommutative sense of German Körper. The rationale explicitly distinguishes it from a commutative field. In Paper 32, §2, the rational basis is $1,i,j,k$, with $ij=k$ and $ji=-k$; replacing tělo with polje would erase a mathematical distinction. Subsequent extension to a splitting field can produce a nontrivial idempotent, so statements about the extended algebra must not be read as asserting it remains a division ring. The recorded Russian and Ukrainian body terms support the naming choice, but do not constitute an independently established Interslavic norm.\par
\noindent{\footnotesize Izvorne zapisy / Original records: HIST-g174-001.}\par
\smallskip\noindent\emph{Sravnjenje s tekstom:} Zajedny korenj ne sbližaje definicije těla i grupy. Dva pridavne varianty ne imajut zachovanu sravniteljnu ocenu.\par
\noindent\emph{Textual comparison:} A common stem does not identify division ring and group. The two adjectival variants have no preserved comparative ranking.\par
\noindent{\footnotesize TEXT-COMP-487; pasaže / passages: 19244, 19391, 19392, 19393, 19394, 19395, 24683, 24684, 24685, 24686, 24687, 24688, 24689, 24690, 24691; \texttt{CORPUS\_FORM\_COMPARISONS.json}.}\par
\end{minipage}\par
\par\medskip\Needspace{7\baselineskip}
\noindent\begin{minipage}{\linewidth}
\hypertarget{TERM-1412}{}
\noindent\textbf{kvaternionsko tělo}\quad{\footnotesize\texttt{TERM-1412}}\par
\noindent\emph{Iztočnik / Source:} Quaternionenkörper\par
\noindent\emph{Rusky / Russian:} кватернионное тело\par
\noindent\emph{Ukrajinsky / Ukrainian:} кватерніонне тіло\par
\smallskip\noindent\emph{Razlog:} Kvaternionsko tělo čuva tož strukturno razlikovanje tělo/polje kako kvaternionno tělo, ale ima drugi tvorbeny sufiks. Arhiv povezaje izbor s politikoj raboty 40 za nekomutativny Körper; ne pokazuje sravniteljne oceny sufiksov -n- i -sk-. Rabota 42 razgleđaje maksimalny poredok kvaternionskogo těla, ktory sodrživaje poredky dvoh različnyh polj. To potvrdžaje potrebu razlikovati obhvatnu nekomutativnu strukturu od komutativnyh podpolj. Dvě zachovane pridavne formy zato trěba pokazati kako varianty istej terminologičnoj semji, ne izjaviti jednu pobědnicom bez razloga. Pozdnějše ujednačenje sufiksa trěbuje izrično lingvistično opravdanje.\par
\noindent\emph{Reason:} Kvaternionsko tělo preserves the same body/field distinction as kvaternionno tělo, but uses a different adjectival suffix. The record links this choice to Paper 40's policy for noncommutative Körper; it provides no comparative ranking of -n- and -sk-. Paper 42 discusses a maximal order of the quaternion division ring containing orders of two different fields, confirming the need to distinguish the surrounding noncommutative structure from its commutative subfields. The two recorded adjectives should therefore appear as variants in one terminological family, not as a winner and loser without evidence. A later suffix standardization requires an explicit linguistic justification.\par
\noindent{\footnotesize Izvorne zapisy / Original records: HIST-g212-014.}\par
\smallskip\noindent\emph{Sravnjenje s tekstom:} Zajedny korenj ne sbližaje definicije těla i grupy. Dva pridavne varianty ne imajut zachovanu sravniteljnu ocenu.\par
\noindent\emph{Textual comparison:} A common stem does not identify division ring and group. The two adjectival variants have no preserved comparative ranking.\par
\noindent{\footnotesize TEXT-COMP-487; pasaže / passages: 19244, 19391, 19392, 19393, 19394, 19395, 24683, 24684, 24685, 24686, 24687, 24688, 24689, 24690, 24691; \texttt{CORPUS\_FORM\_COMPARISONS.json}.}\par
\end{minipage}\par
\par\medskip\Needspace{7\baselineskip}
\noindent\begin{minipage}{\linewidth}
\hypertarget{TERM-1413}{}
\noindent\textbf{kvaternionno tělo / kvaternionna grupa}\quad{\footnotesize\texttt{TERM-1413}}\par
\noindent\emph{Iztočnik / Source:} Quaternionenkörper / Quaternionengruppe\par
\noindent\emph{Rusky / Russian:} тело кватернионов / группа кватернионов\par
\noindent\emph{Ukrajinsky / Ukrainian:} тіло кватерніонів / група кватерніонів\par
\smallskip\noindent\emph{Razlog:} Kvaternionno tělo / kvaternionna grupa čuva zajedny korenj za svezane, ale različne objekty. Arhiv pripisuje internacionalnomu korenju medžuslovjansku razpoznavajemost; to jest motivacija, ne rezultat měrjenja razumlivosti. V drugom uvodu raboty 32 razpadne polja grupy sut svezane s prědstavjenjami kvaternionnogo těla. Teorija prědstavjenij objasnjaje vez, ale ne stvara identičnost medžu grupoj i tělom. Zato imenica, ne samo korenj kvaternion-, mora ostati eksplicitna. Takovo obhvačenje semji terminov jest korisno pri širjenju slovnika, dokolě definicije objektov ne sut směšane.\par
\noindent\emph{Reason:} Kvaternionno tělo / kvaternionna grupa retains one stem for related but different objects. The archive attributes cross-Slavic recognizability to the international stem; that is a motivation, not a measured comprehension result. Paper 32's second introduction connects the group's splitting fields with representations of the quaternion division ring. Representation theory explains the connection, but does not identify the group with the division ring. The noun, not merely the shared quaternion stem, must remain explicit. Recording a terminological family is useful for dictionary expansion provided that its object definitions remain distinct.\par
\noindent{\footnotesize Izvorne zapisy / Original records: HIST-g168-004.}\par
\smallskip\noindent\emph{Sravnjenje s tekstom:} Zajedny korenj ne sbližaje definicije těla i grupy. Dva pridavne varianty ne imajut zachovanu sravniteljnu ocenu.\par
\noindent\emph{Textual comparison:} A common stem does not identify division ring and group. The two adjectival variants have no preserved comparative ranking.\par
\noindent{\footnotesize TEXT-COMP-487; pasaže / passages: 19244, 19391, 19392, 19393, 19394, 19395, 24683, 24684, 24685, 24686, 24687, 24688, 24689, 24690, 24691; \texttt{CORPUS\_FORM\_COMPARISONS.json}.}\par
\end{minipage}\par
\par\medskip\Needspace{7\baselineskip}
\noindent\begin{minipage}{\linewidth}
\hypertarget{TERM-1415}{}
\noindent\textbf{kvocient $a/D^{(i)}$}\quad{\footnotesize\texttt{TERM-1415}}\par
\noindent\emph{Iztočnik / Source:} Quotient $a/D^{(i)}$\par
\noindent\emph{Rusky / Russian:} частное $a/D^{(i)}$\par
\noindent\emph{Ukrajinsky / Ukrainian:} частка $a/D^{(i)}$\par
\smallskip\noindent\emph{Razlog:} Kvocient $a/D^{(i)}$ ohranjaje izraz iz zaključnej primětky raboty 22, §6. Arhiv zahtěvaje kontinuitet s terminologijeju Hentzelta i Noether, ne prěpisanje v pozdnějši obči izraz faktor-modul. V nynějšem tekstu $a$ označaje ideal, a $D^{(i)}$ jest polinom povezan s sukcesivnoj eliminacijeju. Slědujuči paragraf izrično razgleđaje kvocient ideala prěz osnovny ideal; definicija v §1 vysvětljaje kvocient pomoču množenja i prinaležnosti. Sama kosa črta zato ne znači, že elementy ideala sut prosto razděljene na ostatkove klasy. Symboly v iztočnom spisku sut skračeny zapis; slovnična glava mora jih prikazati matematično, bez zaměny smysla operacije.\par
\noindent\emph{Reason:} Kvocient $a/D^{(i)}$ retains the expression in Paper 22, §6's closing note. The record calls for continuity with Hentzelt and Noether's quotient terminology, not recasting it as a later generic quotient module. In the current text $a$ denotes an ideal and $D^{(i)}$ a polynomial associated with successive elimination. The next section expressly considers an ideal quotient by a basic ideal; §1 defines quotients through multiplication and membership conditions. A slash alone therefore does not imply passage to residue classes of the ideal. The archived headword is a compressed notation that needs mathematical typesetting without changing the operation's sense.\par
\noindent{\footnotesize Izvorne zapisy / Original records: HIST-g090-015, HIST-g092-069.}\par
\smallskip\noindent\emph{Sravnjenje s tekstom:} Zapis s kosoj črtoj nosi lokalny idealny smysl; ne ustanavjaje avtomatično faktor-modul.\par
\noindent\emph{Textual comparison:} The slash carries the local ideal-quotient sense; it does not automatically denote a quotient module.\par
\noindent{\footnotesize TEXT-COMP-489; pasaže / passages: 14971, 14976, 14977, 14978, 14979, 14980, 14981, 14982, 14983; \texttt{CORPUS\_FORM\_COMPARISONS.json}.}\par
\end{minipage}\par
\par\medskip\Needspace{7\baselineskip}
\noindent\begin{minipage}{\linewidth}
\hypertarget{TERM-1416}{}
\noindent\textbf{kvocient ... črez osnovny ideal}\quad{\footnotesize\texttt{TERM-1416}}\par
\noindent\emph{Iztočnik / Source:} Quotient by the basic ideal\par
\noindent\emph{Rusky / Russian:} частное ... по основному идеалу\par
\noindent\emph{Ukrajinsky / Ukrainian:} частка ... за основним ідеалом\par
\smallskip\noindent\emph{Razlog:} Kvocient … črez osnovny ideal povezuje kvocient s už uvedenym osnovnym idealom. Arhiv to opravdava kontinuitetom terminologije osnovnogo modula i idealnoj dělimosti; anglijsky iztočny naslov v spisku jest opis pojma, ne citat německogo originala. Rabota 22, §7, piše idealny kvocient $\mathfrak a=\mathfrak m/\mathfrak g_{i-1}$ i koristi prědlog prěz. Zato črez i prěz sut dokumentovane povrhnostne varianty, ne dokaz različnyh algebraičnyh operacij. Kontekst govori o transformovanom idealu, jeho osnovnom idealu i uslovju množenja, ne o množině ostatkovyh klasov. Terminologična kontinuitet ne znači pravo odstraniti takovo kontekstno razlikovanje.\par
\noindent\emph{Reason:} Kvocient … črez osnovny ideal connects the quotient with the already introduced basic ideal. The rationale is continuity with basic-module terminology and ideal divisibility; the English source-side heading is a concept description, not a quotation from the German original. Paper 22, §7, writes the ideal quotient $\mathfrak a=\mathfrak m/\mathfrak g_{i-1}$ and uses prěz. Thus črez and prěz are documented surface variants, not evidence of different algebraic operations. The context concerns a transformed ideal, its basic ideal and a multiplication condition, not a set of residue classes. Terminological continuity must not erase that contextual distinction.\par
\noindent{\footnotesize Izvorne zapisy / Original records: HIST-g091-002, HIST-g092-071.}\par
\smallskip\noindent\emph{Sravnjenje s tekstom:} Prěz v tekstu odgovarjaje črez v spisku. Smysl ustanavjajut idealy i uslovje množenja.\par
\noindent\emph{Textual comparison:} The corpus's prěz corresponds to the list's črez. The ideals and multiplication condition establish the sense.\par
\noindent{\footnotesize TEXT-COMP-490; pasaže / passages: 14976, 14977, 14978, 14979, 14980, 14981, 14982, 14983; \texttt{CORPUS\_FORM\_COMPARISONS.json}.}\par
\end{minipage}\par
\par\medskip\Needspace{7\baselineskip}
\noindent\begin{minipage}{\linewidth}
\hypertarget{TERM-1417}{}
\noindent\textbf{kvocient dvoh modulov}\quad{\footnotesize\texttt{TERM-1417}}\par
\noindent\emph{Iztočnik / Source:} Quotient zweier Moduln\par
\noindent\emph{Rusky / Russian:} частное двух модулей\par
\noindent\emph{Ukrajinsky / Ukrainian:} частка двох модулів\par
\smallskip\noindent\emph{Razlog:} Kvocient dvoh modulov ohranjaje istorijsko nazvanje, ale njegov smysl trěba vzeti iz definicije. Rabota 22, §1, oprěděljaje kvocient modulov linearnyh form kako ideal vsih skalarov $c$, za ktore $cB$ leži v $A$. To ne jest obyčny faktor-modul $A/B$. Primětka takože izrično govori, že tuty kvocient ne sběgaje se s Dedekindovym kvocientom čislovyh modulov. Arhivna charakterizacija kako Dedekindovy typ zato označaje istorijsku vez, ne polnu identičnost definicij. Racionalny izbor slova kvocient može byti ohranjen samo zajedno s tymi ograničenjami; tri povtorjene zapisy prěvoda ne zaměnjajut matematičnu definiciju.\par
\noindent\emph{Reason:} Kvocient dvoh modulov preserves the historical name, but its sense must come from the definition. Paper 22, §1, defines the quotient of modules of linear forms as the ideal of scalars $c$ for which $cB$ lies in $A$. This is not the ordinary quotient module $A/B$. The footnote also expressly says that the quotient defined here does not coincide with Dedekind's quotient of numerical modules. The archive's 'Dedekind-style' characterization therefore records a historical connection, not identical definitions. Kvocient is defensible only with these qualifications; three repeated translation records cannot replace the mathematical definition.\par
\noindent{\footnotesize Izvorne zapisy / Original records: HIST-g085-008, HIST-g092-008, HIST-g093-010.}\par
\smallskip\noindent\emph{Sravnjenje s tekstom:} Definicija govori o skalarah, čiji produkty s drugim modulom ležat v prvom. Primětka čuva razliku od Dedekinda.\par
\noindent\emph{Textual comparison:} The definition concerns scalars whose products with the second module lie in the first; the note retains the distinction from Dedekind.\par
\noindent{\footnotesize TEXT-COMP-491; pasaže / passages: 14294, 14295, 14296, 14297, 14298, 14299, 14300, 14301, 14302, 14303, 14304, 14305; \texttt{CORPUS\_FORM\_COMPARISONS.json}.}\par
\end{minipage}\par
\par\medskip\Needspace{7\baselineskip}
\noindent\begin{minipage}{\linewidth}
\hypertarget{TERM-1418}{}
\noindent\textbf{tvorjenje kvocientov}\quad{\footnotesize\texttt{TERM-1418}}\par
\noindent\emph{Iztočnik / Source:} Quotientenbildung\par
\noindent\emph{Rusky / Russian:} образование частных\par
\noindent\emph{Ukrajinsky / Ukrainian:} утворення часток\par
\smallskip\noindent\emph{Razlog:} Tvorjenje kvocientov sostavjaje slovjanske ime procesa s internacionalnym algebraičnym korenjem. Tri zachovane razloga odnoset se k rabotam 9, 24 i 25; oni čuvajut terminologičnu kontinuitet i kratkost, bez dokumentovanogo čislovogo vaganja. V rabotě 9 tvorjenje kvocientov prěnosi izomorfnu korespondenciju. V rabotah 24 i 25 prsten ostatkovyh klasov po vlastnom prostom idealu najprvo nema nulovyh děliteljev, a potom razširjaje se do polja tvorjenjem drobov. To sut dva koraka: obrazovanje ostatkovyh klasov i tvorjenje kvocientov v dobytoj oblasti. Izraz ne směje skryti drugy korak ni sugerovati, že kvocient po libovoljnym idealu vsegda jest polje.\par
\noindent\emph{Reason:} Tvorjenje kvocientov combines a Slavic process noun with the international algebraic stem. The three preserved reasons concern Papers 9, 24 and 25 and emphasize continuity and compactness without a numerical weighting exercise. In Paper 9, formation of quotients extends an isomorphic correspondence. In Papers 24 and 25, residue classes modulo a proper prime ideal first form a ring without zero divisors, which is then extended to a field by forming fractions. These are two operations: taking residue classes and forming fractions in the resulting domain. The expression must not conceal the second step or suggest that a quotient by an arbitrary ideal is already a field.\par
\noindent{\footnotesize Izvorne zapisy / Original records: HIST-g067-009, HIST-g097-003, HIST-g102-008.}\par
\smallskip\noindent\emph{Sravnjenje s tekstom:} Tvorjenje ostatkovyh klasov i razširenje do polja sut različne operacije, eksplicitno povezane v tekstu.\par
\noindent\emph{Textual comparison:} Taking residue classes and extending to a field are distinct operations explicitly linked in the text.\par
\noindent{\footnotesize TEXT-COMP-492; pasaže / passages: 8200, 8201, 8202, 8203, 8204, 8205, 15441, 15443, 15839; \texttt{CORPUS\_FORM\_COMPARISONS.json}.}\par
\end{minipage}\par
\par\medskip\Needspace{7\baselineskip}
\noindent\begin{minipage}{\linewidth}
\hypertarget{TERM-1419}{}
\noindent\textbf{idealny kvocient nulovogo ideala črez $c$}\quad{\footnotesize\texttt{TERM-1419}}\par
\noindent\emph{Iztočnik / Source:} Quotientenideal $(0):c$\par
\noindent\emph{Rusky / Russian:} частное нулевого идеала по $c$\par
\noindent\emph{Ukrajinsky / Ukrainian:} частка нульового ідеалу за $c$\par
\smallskip\noindent\emph{Razlog:} Idealny kvocient nulovogo ideala črez $c$ razkladaje zapis $(0):c$ v čitljivu rečenicu. Arhiv izrično pokračuje v terminologiji iz §4, ne tvori drugo ime za tož operaciju. V rabotě 31, §5, $c$ jest ideal: relevantne elementy $\delta$ izpolnjajut $\delta c=0$. Kvocient nulovogo ideala zato označaje anihilator togo ideala, ne prsten ostatkovyh klasov po njem. Nynějšny tekst ima prěz i frakturno označenje ideala; te razlike trěba odděliti od smysla. Čuvanje dvotočky i obuhvata uslovja množenja jest važnějše od mehaničnogo ohranjanja skračenogo tekstovogo zapisa v staroj tabeli.\par
\noindent\emph{Reason:} Idealny kvocient nulovogo ideala črez $c$ spells out $(0):c$ in readable prose. The rationale explicitly continues §4's vocabulary rather than inventing a second name for the same operation. In Paper 31, §5, $c$ is an ideal: the relevant elements $\delta$ satisfy $\delta c=0$. This quotient of the zero ideal is thus the ideal's annihilator, not a residue-class ring modulo it. The current text uses prěz and a Fraktur ideal symbol; those surface differences must be separated from meaning. Preserving the colon and the scope of the multiplication condition matters more than mechanically reproducing the old table's abbreviated plain-text notation.\par
\noindent{\footnotesize Izvorne zapisy / Original records: HIST-g138-005.}\par
\smallskip\noindent\emph{Sravnjenje s tekstom:} Uslovje $\delta c=0$ označaje anihilator ideala, ne kvocientny prsten.\par
\noindent\emph{Textual comparison:} The condition $\delta c=0$ defines an ideal annihilator, not a quotient ring.\par
\noindent{\footnotesize TEXT-COMP-493; pasaže / passages: 18152; \texttt{CORPUS\_FORM\_COMPARISONS.json}.}\par
\end{minipage}\par
\par\medskip\Needspace{7\baselineskip}
\noindent\begin{minipage}{\linewidth}
\hypertarget{TERM-1420}{}
\noindent\textbf{kvocientno polje}\quad{\footnotesize\texttt{TERM-1420}}\par
\noindent\emph{Iztočnik / Source:} Quotientenkörper\par
\noindent\emph{Rusky / Russian:} поле частных\par
\noindent\emph{Ukrajinsky / Ukrainian:} поле часток\par
\smallskip\noindent\emph{Razlog:} Kvocientno polje zachovuje medžunarodny korenj kvocient- za polje drobov oblasti cělostnosti. Arhiv povezuje rabotu 9 s terminologijeju jej poprědnyh paragrafov i rabotu 35 s točnostju internacionalnogo izraza; ruska i ukrajinska paralela koristajut nazvanja polja častnyh odnosno častok. Rabota 9 izrično opisuje elementy polja kako kvocienty elementov oblasti; pri smyslovom vysvětlenju imenitelj mora byti nenulovy. Rabota 35 ustanavjaje polje kvocientov $Q$ oblasti $I$ i medžu nimi oblast $Z$. To ne jest polje ostatkovyh klasov po idealu. Slovnik zato čuva zajednu semju kvocient-, ale ne spojuje različne konstrukcije v jednu definiciju.\par
\noindent\emph{Reason:} Kvocientno polje uses the international quotient stem for the fraction field of an integral domain. The record links Paper 9 to its preceding sections and Paper 35 to the precision of the international expression; Russian and Ukrainian comparanda use their established quotient-field names. Paper 9 explicitly describes field elements as quotients of elements of the domain; an explanatory definition must require a nonzero denominator. Paper 35 specifies the fraction field $Q$ of $I$ and an intermediate domain $Z$. This is not a residue-class field modulo an ideal. The dictionary therefore retains the common kvocient- family without collapsing its different constructions into one definition.\par
\noindent{\footnotesize Izvorne zapisy / Original records: HIST-g068-009, HIST-g205-010.}\par
\smallskip\noindent\emph{Sravnjenje s tekstom:} Kvocientno polje nastaje iz oblasti drobami s nenulovym imeniteljem. Tuta konstrukcija ne jest tvorjenje ostatkovyh klasov.\par
\noindent\emph{Textual comparison:} The fraction field is formed from a domain using fractions with nonzero denominators, not by taking residue classes.\par
\noindent{\footnotesize TEXT-COMP-494; pasaže / passages: 8370, 8371, 8372, 8373, 8374, 8376, 8377, 8378, 8379, 8380, 22576, 22578; \texttt{CORPUS\_FORM\_COMPARISONS.json}.}\par
\end{minipage}\par
\par\medskip\Needspace{7\baselineskip}
\noindent\begin{minipage}{\linewidth}
\hypertarget{TERM-1421}{}
\noindent\textbf{polje kvocientov}\quad{\footnotesize\texttt{TERM-1421}}\par
\noindent\emph{Iztočnik / Source:} Quotientenkörper\par
\noindent\emph{Rusky / Russian:} поле частных\par
\noindent\emph{Ukrajinsky / Ukrainian:} поле часток\par
\smallskip\noindent\emph{Razlog:} Polje kvocientov jest imenična konstrukcija s genitivom, odgovarjajuča pridavnoj formě kvocientno polje. Sedm zachovanyh zapisov povezuje jeje upotrěbu v rabotah 24, 29, 30 i 31: koeficientno polje po prisojedinjenju neoprěděljenyh, razširenja polj, integralna zamknjenost, slěd i norma. Razlog za korenj kvocient- jest kontinuitet i izběganje obyčnogo čestičnogo ili drobnogo smysla; razpoznavajemost jest tvrđenje arhiva, ne izměrjena vlastnost. Tekst razlikuje polje kvocientov od prstena kvocientov i od idealnyh/modulnyh kvocientov. Povtorjenje izbora ne dokazuje, že genitivna forma jest bolja od pridavnoj; arhiv ne davaje takej sravniteljnoj oceny. Obě formy morajut byti povezane jednym smyslom, dokolě lokalne prědpoloženja ostavajut vidime.\par
\noindent\emph{Reason:} Polje kvocientov is a noun-plus-genitive construction corresponding to the adjectival kvocientno polje. Seven preserved records connect its use in Papers 24, 29, 30 and 31: coefficient fields after adjoining indeterminates, field extensions, integral closure, trace and norm. The reason for kvocient- is continuity and avoidance of ordinary partitive or fractional senses; recognizability is asserted by the archive, not measured. The corpus distinguishes fraction fields from rings of fractions and ideal/module quotients. Repeated selection does not establish that the genitive construction is preferable to the adjective; no such comparative ranking survives. Both forms need one sense connection while local hypotheses remain visible.\par
\noindent{\footnotesize Izvorne zapisy / Original records: HIST-g100-005, HIST-g106-010, HIST-g108-010, HIST-g110-002, HIST-g115-002, HIST-g136-005, HIST-g149-003.}\par
\smallskip\noindent\emph{Sravnjenje s tekstom:} Konstrukcija so genitivom ostavaje v različnih kontekstah. Povtorjenje ne jest sravniteljne jezyčno testovanje.\par
\noindent\emph{Textual comparison:} The genitive construction recurs in different contexts. Recurrence is not a comparative language test.\par
\noindent{\footnotesize TEXT-COMP-495; pasaže / passages: 15655, 15983, 16021, 16135, 16391, 16845, 18065, 18559; \texttt{CORPUS\_FORM\_COMPARISONS.json}.}\par
\end{minipage}\par
\par\medskip\Needspace{7\baselineskip}
\noindent\begin{minipage}{\linewidth}
\hypertarget{TERM-1422}{}
\noindent\textbf{polje kvocientov / kolco kvocientov}\quad{\footnotesize\texttt{TERM-1422}}\par
\noindent\emph{Iztočnik / Source:} Quotientenkörper / Quotientenring\par
\noindent\emph{Rusky / Russian:} поле частных / кольцо частных\par
\noindent\emph{Ukrajinsky / Ukrainian:} поле часток / кільце часток\par
\smallskip\noindent\emph{Razlog:} Polje kvocientov / kolco kvocientov čuva razliku polje/prsten pri zajednom korenju kvocient-. V nynějšem tekstu staromu kolco odgovarjaje prsten. Rabota 30 pri odsutnosti nulovyh děliteljev razgleđaje polje drobov; pri dopuščenju nulovyh děliteljev imenitelj mora byti regularny element. To jest element, ktory ne jest nulovy dělitelj, ne prosto libovoljny nenulovy element. V §9 ostatkovy prsten služi protivpriměru i jest jednak svojemu prstenu kvocientov, zato že vse regularne elementy sut jedinice. Ostatkovy prsten i prsten kvocientov zato imajut različne konstrukcije, hoti v tom priměru objekty sběgajut se. Arhiv izrično zahtěvaje čuvati tutu razliku, ne izravniti nazvanja.\par
\noindent\emph{Reason:} Polje kvocientov / kolco kvocientov preserves the field/ring distinction within the common kvocient- family. The current text uses prsten where the older list has kolco. Paper 30 treats a fraction field in the absence of zero divisors; when zero divisors are allowed, the denominator must be regular, meaning a non-zero-divisor rather than merely a nonzero element. In §9 a residue-class ring supplies a counterexample and equals its own ring of fractions because every regular element is a unit. Residue-class ring and ring of fractions are different constructions even when their resulting objects coincide in that example. The archive expressly calls for preserving this distinction.\par
\noindent{\footnotesize Izvorne zapisy / Original records: HIST-g107-006, HIST-g116-006, HIST-g116-014.}\par
\smallskip\noindent\emph{Sravnjenje s tekstom:} Pri nulovyh děliteljah regularnost imenitelja jest važna. Ostatkovy prsten ne jest po definiciji prsten kvocientov.\par
\noindent\emph{Textual comparison:} With zero divisors, denominator regularity matters. A residue-class ring is not by definition a ring of fractions.\par
\noindent{\footnotesize TEXT-COMP-496; pasaže / passages: 16115, 16124, 16139, 16261, 17147, 17149, 17277; \texttt{CORPUS\_FORM\_COMPARISONS.json}.}\par
\end{minipage}\par
\par\medskip\Needspace{7\baselineskip}
\noindent\begin{minipage}{\linewidth}
\hypertarget{TERM-1423}{}
\noindent\textbf{polje kvocientov integralnoj oblasti}\quad{\footnotesize\texttt{TERM-1423}}\par
\noindent\emph{Iztočnik / Source:} Quotientenkörper eines Integritätsbereichs\par
\noindent\emph{Rusky / Russian:} поле частных области целостности\par
\noindent\emph{Ukrajinsky / Ukrainian:} поле часток інтегральної області\par
\smallskip\noindent\emph{Razlog:} Polje kvocientov integralnoj oblasti zadržuje osnovnu oblast v samom nazvanju. Arhiv opravdava internacionalny korenj prozračnostju standardnogo pojma, ne predlagaje nove konstrukcije. Uvod raboty 6 opisuje najmenše polje, ktore sodrživaje libovoljny razgledany sistem, kako obobčenje polja kvocientov integralnoj oblasti. Fraza zato mora čuvati razliku medžu danoj oblastju i poljem, ktore iz njej nastaje; ne govori, že sama oblast už jest polje. Integralny tutaj znači algebraičnu oblast bez nulovyh děliteljev, ne integral iz analizy. To konkretne sravnjenje potvrđaje upotrěbu pojma v uvodu; ono ne jest nezavisna istorijska atestacija medžuslovjanskogo nazvanja.\par
\noindent\emph{Reason:} Polje kvocientov integralnoj oblasti keeps the originating domain in the name. The preserved motivation attributes transparency to the international stem for a standard concept; it proposes no new construction. Paper 6's introduction describes the smallest field containing an arbitrary system under consideration as a generalization of the fraction field of an integral domain. The wording must distinguish the given domain from the field constructed from it, not imply that the domain is already a field. Integral concerns the algebraic domain without zero divisors, not an integral in analysis. This comparison confirms the concept in the introduction, not an independent historical attestation of the Interslavic name.\par
\noindent{\footnotesize Izvorne zapisy / Original records: HIST-g040-017.}\par
\smallskip\noindent\emph{Sravnjenje s tekstom:} Uvod razlikuje integralnu oblast od najmenšego polja, ktore ju sodrživaje.\par
\noindent\emph{Textual comparison:} The introduction distinguishes the integral domain from its smallest containing field.\par
\noindent{\footnotesize TEXT-COMP-497; pasaže / passages: 5033; \texttt{CORPUS\_FORM\_COMPARISONS.json}.}\par
\end{minipage}\par
\par\medskip\Needspace{7\baselineskip}
\noindent\begin{minipage}{\linewidth}
\hypertarget{TERM-1424}{}
\noindent\textbf{polje kvocientov od $T$}\quad{\footnotesize\texttt{TERM-1424}}\par
\noindent\emph{Iztočnik / Source:} Quotientenkörper von $T$\par
\noindent\emph{Rusky / Russian:} поле частных $T$\par
\noindent\emph{Ukrajinsky / Ukrainian:} поле часток $T$\par
\smallskip\noindent\emph{Razlog:} Polje kvocientov od $T$ izrično imenuje prsten, iz ktorogo polje nastaje. Zachovany razlog jest soglasnost s ranějšim slovnikom polj kvocientov; nema osobnogo vaganja prědloga od. V rabotě 31, §8, polje kvocientov od $T$ jest obhvatna oblast za prsten $T_a$, ktory dopuščaje samo imenitelje proste k danomu idealu. Zato polje kvocientov i takovy prsten ne sut synonimy. Bukva $T$ jest matematičny identifikator, ktory ostavaje latinskoju bukvoju i v kiriličnom izloženju; typografsko prikazanje ne směje zaměniti ju podobnoju kiriličnoju bukvoj. Proverka tutoj upotrěby ne dokazuje, že v vsakom drugom kontekstu T znači toj samy objekt.\par
\noindent\emph{Reason:} Polje kvocientov od $T$ explicitly names the ring from which the field is formed. The surviving reason is consistency with the earlier quotient-field glossary; it supplies no separate ranking of the preposition od. In Paper 31, §8, the fraction field of $T$ contains the ring $T_a$, whose denominators are restricted by coprimality to the given ideal. The fraction field and that ring are not synonyms. $T$ is a mathematical identifier and remains Latin in Cyrillic prose; typesetting must not replace it with a similar Cyrillic letter. Checking this occurrence does not establish that every T elsewhere denotes the same object.\par
\noindent{\footnotesize Izvorne zapisy / Original records: HIST-g159-003.}\par
\smallskip\noindent\emph{Sravnjenje s tekstom:} Prsten $T_a$ jest oprěděljen imeniteljami v polju kvocientov. Ideal $a$ ne jest v definiciji ograničeny na prosty ideal.\par
\noindent\emph{Textual comparison:} The ring $T_a$ is specified by denominators within the fraction field. The definition does not restrict $a$ to prime ideals.\par
\noindent{\footnotesize TEXT-COMP-498; pasaže / passages: 18932, 18970; \texttt{CORPUS\_FORM\_COMPARISONS.json}.}\par
\end{minipage}\par
\par\medskip\Needspace{7\baselineskip}
\noindent\begin{minipage}{\linewidth}
\hypertarget{TERM-1425}{}
\noindent\textbf{kolco kvocientov}\quad{\footnotesize\texttt{TERM-1425}}\par
\noindent\emph{Iztočnik / Source:} Quotientenring\par
\noindent\emph{Rusky / Russian:} кольцо частных\par
\noindent\emph{Ukrajinsky / Ukrainian:} кільце часток\par
\smallskip\noindent\emph{Razlog:} Kolco kvocientov bylo vybrano za prsten drobov, ne za faktor-prsten. Dva zachovane razloga imajut različne lokalne uslovja: rabota 30 dopuščaje nulove dělitelje i zato uporablja regularne imenitelje; uvod raboty 31 razgleđaje prěhod po odnošenju k prostomu idealu. Nynějšny tekst v oboju slučajah piše prsten kvocientov. Razlika medžu poljem i prstenom jest strukturna, dokolě kolco/prsten jest leksična variacija. Sama prisutnost slova kvocient ne oprěděljaje ni množinu dopustimyh imeniteljev ni kvocient po idealu. Slovnik zato povezuje formy, ale ne izpuščaje definiciju konkretnoj lokalizacije.\par
\noindent\emph{Reason:} Kolco kvocientov was selected for a ring of fractions, not a factor ring. The two preserved reasons have different local conditions: Paper 30 permits zero divisors and therefore uses regular denominators, while Paper 31's introduction concerns passage relative to a prime ideal. The current text calls both prsten kvocientov. Field versus ring is a structural distinction, whereas kolco versus prsten is a lexical variation. The word quotient alone specifies neither the allowed denominators nor passage modulo an ideal. The dictionary connects the forms without suppressing the definition of the particular localization.\par
\noindent{\footnotesize Izvorne zapisy / Original records: HIST-g108-011, HIST-g118-011.}\par
\smallskip\noindent\emph{Sravnjenje s tekstom:} Pri nulovyh děliteljah regularnost imenitelja jest važna. Ostatkovy prsten ne jest po definiciji prsten kvocientov.\par
\noindent\emph{Textual comparison:} With zero divisors, denominator regularity matters. A residue-class ring is not by definition a ring of fractions.\par
\noindent{\footnotesize TEXT-COMP-496; pasaže / passages: 16115, 16124, 16139, 16261, 17147, 17149, 17277; \texttt{CORPUS\_FORM\_COMPARISONS.json}.}\par
\end{minipage}\par
\par\medskip\Needspace{7\baselineskip}
\noindent\begin{minipage}{\linewidth}
\hypertarget{TERM-1426}{}
\noindent\textbf{kolco kvocientov $T_a$}\quad{\footnotesize\texttt{TERM-1426}}\par
\noindent\emph{Iztočnik / Source:} Quotientenring $T_a$\par
\noindent\emph{Rusky / Russian:} кольцо частных $T_a$\par
\noindent\emph{Ukrajinsky / Ukrainian:} кільце часток $T_a$\par
\smallskip\noindent\emph{Razlog:} Kolco kvocientov $T_a$ daje kratko nazvanje lokalizovanomu prstenu, da by daljše formule ostali čitljive. To jest izričny razlog arhiva. Nynějšna definicija v rabotě 31, §8, piše prsten kvocientov i koristi frakturne symboly: elementy ležat v polju kvocientov od $T$, a jih imenitelje sut proste k idealu $a$. Tutaj $a$ jest libovoljny ideal razgledanoj osnovnoj oblasti, ne avtomatično prosty ideal. Sledovateljno ne možno iz samogo indeksa zaključiti, že prsten mora byti lokalny s jednym maksimalnym idealom. Podindeks jest čest definicije dopustimyh imeniteljev; ne jest nomer prstena ni pokazatelj faktor-prstena.\par
\noindent\emph{Reason:} Kolco kvocientov $T_a$ gives the localized ring a short name so that subsequent formulas remain readable; this is the archive's explicit rationale. The current definition in Paper 31, §8, uses prsten kvocientov and Fraktur symbols: its elements lie in the fraction field of $T$ and have denominators coprime to the ideal $a$. Here $a$ is an arbitrary ideal of the specified base domain, not automatically a prime ideal. The subscript alone therefore does not establish that the ring is local with a single maximal ideal. It identifies the denominator restriction, not a serial number or a factor-ring construction.\par
\noindent{\footnotesize Izvorne zapisy / Original records: HIST-g159-002.}\par
\smallskip\noindent\emph{Sravnjenje s tekstom:} Prsten $T_a$ jest oprěděljen imeniteljami v polju kvocientov. Ideal $a$ ne jest v definiciji ograničeny na prosty ideal.\par
\noindent\emph{Textual comparison:} The ring $T_a$ is specified by denominators within the fraction field. The definition does not restrict $a$ to prime ideals.\par
\noindent{\footnotesize TEXT-COMP-498; pasaže / passages: 18932, 18970; \texttt{CORPUS\_FORM\_COMPARISONS.json}.}\par
\end{minipage}\par
\par\medskip\Needspace{7\baselineskip}
\noindent\begin{minipage}{\linewidth}
\hypertarget{TERM-1429}{}
\noindent\textbf{$R$-porjadok}\quad{\footnotesize\texttt{TERM-1429}}\par
\noindent\emph{Iztočnik / Source:} $R$-Ordnung\par
\noindent\emph{Rusky / Russian:} $R$-порядок\par
\noindent\emph{Ukrajinsky / Ukrainian:} $R$-порядок\par
\smallskip\noindent\emph{Razlog:} $R$-porjadok čuva osnovny prsten $R$ pri imenici. Arhiv izrično zahtěvaje ohraniti latinsku bukvu kako matematičny symbol vo vseh pismovyh versijah. Nynějšny tekst ima $R$-poredok i v rabotě 30, §1, definira go kako prsten medžu $R$ i obhvatnym $T$. Samo nazvanje tutaj ne vključaje konečnoj porodženosti: ta vlastnost objašnjaje se oddělno za moduly. Takože ne ide o poredok kakoby binarnu relaciju. Terminologična kontinuitet zato trěbuje čuvati osnovny prsten i lokalnu definiciju, ne samo povtoriti korenj porjad-/pored-. Izpustiti $R$ možno jedino kogda on jest nedvosmysleno fiksovany v okolnom izloženju.\par
\noindent\emph{Reason:} $R$-porjadok keeps the base ring $R$ attached to the noun. The archive expressly requires the Latin letter to remain a mathematical symbol in every script version. The current text uses $R$-poredok, defined in Paper 30, §1, as a ring between $R$ and the ambient $T$. The name alone does not require finite generation here; that property is introduced separately for modules. Nor is this order a binary ordering relation. Terminological continuity must retain the base ring and local definition, not merely repeat a porjad-/pored- stem. $R$ can be omitted only when the exposition fixes it unambiguously.\par
\noindent{\footnotesize Izvorne zapisy / Original records: HIST-g108-005.}\par
\smallskip\noindent\emph{Sravnjenje s tekstom:} $R$ jest fiksovana osnovna oblast, ne oznaka binarnoj relacije poredka.\par
\noindent\emph{Textual comparison:} $R$ is the fixed base ring, not a label for a binary ordering relation.\par
\noindent{\footnotesize TEXT-COMP-501; pasaže / passages: 16174, 16196; \texttt{CORPUS\_FORM\_COMPARISONS.json}.}\par
\end{minipage}\par
\par\medskip\Needspace{7\baselineskip}
\noindent\begin{minipage}{\linewidth}
\hypertarget{TERM-1430}{}
\noindent\textbf{Brauerova grupa}\quad{\footnotesize\texttt{TERM-1430}}\par
\noindent\emph{Iztočnik / Source:} R. Brauersche Gruppe\par
\noindent\emph{Rusky / Russian:} группа Брауэра\par
\noindent\emph{Ukrajinsky / Ukrainian:} група Брауера\par
\smallskip\noindent\emph{Razlog:} Brauerova grupa jest prisvojna konstrukcija iz imena Brauer i imenice grupa. Arhiv nazyvaje ju sovrěmennym nazvanjem za klasy centralnyh prostyh algebr; to ne jest tvrđenje, že AI izumila tuty termin. Rabota 38 sama utočnjaje, že elementy grupy sut klasy podobnyh algebr, ne vse algebry kako različne pojedične objekty. V tej rabotě algebra znači normalnu prostu algebru nad algebraičnym čislovym poljem, s tym poljem kako centrom. Pridavno ime Brauerova zadržuje identifikaciju matematika, ale ne zaměnjaje definiciju klasa ni osnovnogo polja. Kratka glava slovnika ne směje razširiti lokalnu teoremu na sve vozmožne algebry.\par
\noindent\emph{Reason:} Brauerova grupa is a possessive construction from Brauer and grupa. The archive describes it as the modern name for classes of central simple algebras, not as a term invented by the AI. Paper 38 itself specifies that the group's elements are similarity classes of algebras, not all individual algebras counted separately. There algebra means a normal simple algebra over an algebraic number field, with that field as its centre. Brauerova retains the mathematician's identity but does not replace the class definition or base-field specification. A short dictionary heading must not extend the local theorem to every possible algebra.\par
\noindent{\footnotesize Izvorne zapisy / Original records: HIST-g208-015.}\par
\smallskip\noindent\emph{Sravnjenje s tekstom:} Normalna prosta algebra i osnovno polje sut izrično oprěděljene; grupa ne sobira algebry bez ekvivalencije.\par
\noindent\emph{Textual comparison:} Normal simple algebra and the base field are specified; the group is not a collection of algebras without equivalence.\par
\noindent{\footnotesize TEXT-COMP-502; pasaže / passages: 23013, 23140; \texttt{CORPUS\_FORM\_COMPARISONS.json}.}\par
\end{minipage}\par
\par\medskip\Needspace{7\baselineskip}
\noindent\begin{minipage}{\linewidth}
\hypertarget{TERM-1431}{}
\noindent\textbf{sistemy blokovyh matric $R^b$ i $R^c$}\quad{\footnotesize\texttt{TERM-1431}}\par
\noindent\emph{Iztočnik / Source:} $R^b$ and $R^c$ block systems\par
\noindent\emph{Rusky / Russian:} системы блочных матриц $R^b$ и $R^c$\par
\noindent\emph{Ukrajinsky / Ukrainian:} системи блокових матриць $R^b$ та $R^c$\par
\smallskip\noindent\emph{Razlog:} Sistemy blokovyh matric $R^b$ i $R^c$ čuvajut prědstavjenje, na ktorom opiraje se dokaz direktnoj sumy. Arhiv pripisuje polnějšemu skanovanomu svědku tute blokove sistemy; tuto tvrđenje ostavaje istorijskim zapisom, ne novoju proverkoju skana. V nynějšem tekstu raboty 31, §5, jedna sistema ima blok $D^b$ i nulovy dopolniteljny blok; druga jest oprěděljena odgovarjajuče. Treba potom pokazati, že one sut podprsteny i idealy odgovarjajučego matričnogo prstena. Vrhni indeks $b$ ne znači potegu: označaje komponentu. Forma matričnogo bloka ne jest dekoracija, a nazvanje ne daje pravo preskočiti dokaz zamknjenosti i direktnosti sumy.\par
\noindent\emph{Reason:} Sistemy blokovyh matric $R^b$ i $R^c$ retains the representation on which the direct-sum argument depends. The archive attributes these block systems to a fuller scan witness; that remains a historical report, not a new inspection of the scan. In the current Paper 31, §5, one system contains a $D^b$ block with a complementary zero block; the other is defined correspondingly. The text then has to establish that these are subrings and ideals of the relevant matrix ring. The superscript b labels a component, not a power. Block form is substantive, and its name does not dispense with closure or directness of the sum.\par
\noindent{\footnotesize Izvorne zapisy / Original records: HIST-g138-006.}\par
\smallskip\noindent\emph{Sravnjenje s tekstom:} Vrhni indeks imenuje komponentu; blokova konstrukcija jest čest dokaza.\par
\noindent\emph{Textual comparison:} The superscript labels a component; the block construction is part of the proof.\par
\noindent{\footnotesize TEXT-COMP-503; pasaže / passages: 18142, 18143, 18144, 18145, 18146, 18147, 18148, 18149, 18150, 18151, 18152; \texttt{CORPUS\_FORM\_COMPARISONS.json}.}\par
\end{minipage}\par
\par\medskip\Needspace{7\baselineskip}
\noindent\begin{minipage}{\linewidth}
\hypertarget{TERM-1432}{}
\noindent\textbf{radikal}\quad{\footnotesize\texttt{TERM-1432}}\par
\noindent\emph{Iztočnik / Source:} Radikal\par
\noindent\emph{Rusky / Russian:} радикал\par
\noindent\emph{Ukrajinsky / Ukrainian:} радикал\par
\smallskip\noindent\emph{Razlog:} Radikal koristi medžunarodny algebraičny korenj, jednaky v zachovanyh ruskih i ukrajinskih paralelah. Arhiv opravdava izbor kontinuitetom raboty 34; ne predlagaje novu rěč ni čislovo vaganje. V §12 radikal jest maksimalny dvustranny nilpotentny ideal, dobyty pri uslovju maksimalnosti za prave idealy. Zato definiciju i jeje prědpoloženja trěba dati zajedno s imenem. Ne možno prosto prenesti tuty opis na libovoljno prsten ili změšati go s korenjem čisla, radikalom ideala ili jinym radikalom iz občej teorije prstenov. Stabilna kratka forma jest korisna, ale jeje jednoznačnost vyhodi iz konkretnoj definicije, ne iz slovničnogo korenja.\par
\noindent\emph{Reason:} Radikal uses an international algebraic stem shared by the recorded Russian and Ukrainian comparanda. The rationale is continuity within Paper 34; it proposes neither a new word nor a numerical ranking. In §12 the radical is the largest two-sided nilpotent ideal, obtained under the maximum condition on right ideals. The definition and its hypotheses therefore belong with the name. This description cannot simply be transferred to arbitrary rings or conflated with a numerical root, an ideal radical or another radical in general ring theory. A stable short form is useful, but its precision comes from the local definition rather than the stem alone.\par
\noindent{\footnotesize Izvorne zapisy / Original records: HIST-g190-002.}\par
\smallskip\noindent\emph{Sravnjenje s tekstom:} Lokalna definicija koristi maksimalny nilpotentny ideal. Prědpoloženje ne směje izčeznuti iz vysvětlenja.\par
\noindent\emph{Textual comparison:} The local definition uses the largest nilpotent ideal. Its hypothesis must not disappear from the explanation.\par
\noindent{\footnotesize TEXT-COMP-504; pasaže / passages: 20833, 20860, 20864; \texttt{CORPUS\_FORM\_COMPARISONS.json}.}\par
\end{minipage}\par
\par\medskip\Needspace{7\baselineskip}
\noindent\begin{minipage}{\linewidth}
\hypertarget{TERM-1433}{}
\noindent\textbf{radikal / kolco bez radikala}\quad{\footnotesize\texttt{TERM-1433}}\par
\noindent\emph{Iztočnik / Source:} Radikal / Ring ohne Radikal\par
\noindent\emph{Rusky / Russian:} радикал / кольцо без радикала\par
\noindent\emph{Ukrajinsky / Ukrainian:} радикал / кільце без радикала\par
\smallskip\noindent\emph{Razlog:} Radikal / kolco bez radikala čuva svez medžu radikalom i ostatkovym prstenom po njem. Nynějšna forma jest prsten bez radikala, i §12 izrično govori, že radikal togda jest nulovy ideal: ne ide o odsutnost definicije. V §19 dokaz upotrěbljaje nilpotentnost, da pokaže anulovanje prostogo modula radikalom i prěhod k ostatkovomu prstenu. Zachovana formulacija teoremy imaje nejasnost: prva alternativa govori o anulovanju radikalom, čo dokaz ustanavjaje za vse razgledane proste moduly; daljše slučaje dokaz razlikuje po dějstvu cělogo ostatkovogo prstena. To trěbuje sravnjenja s iztočnym svědkom prěd redakcijnoj popravkoju. Kontinuitet termina ne jest dokaz točnosti vsakej formulacije teoremy.\par
\noindent\emph{Reason:} Radikal / kolco bez radikala retains the relationship between the radical and the residue-class ring modulo it. The current form is prsten bez radikala, and §12 explicitly means that the radical is the zero ideal, not that it is undefined. In §19 nilpotence shows that the radical annihilates a simple module, allowing passage to the residue-class ring. The preserved theorem wording is unclear: its first alternative says the radical annihilates the module, which the proof establishes for all simple modules considered; the proof distinguishes the subsequent cases by the action of the entire residue-class ring. This needs source-witness comparison before editorial correction. Terminological continuity does not certify every theorem formulation.\par
\noindent{\footnotesize Izvorne zapisy / Original records: HIST-g196-005.}\par
\smallskip\noindent\emph{Sravnjenje s tekstom:} Nulovy radikal imaje definiciju. V §19 prva alternativa teoremy ne razlikuje slučaje, ktore oddělja dokaz; iztočny svědok potrěbuje sravnjenje.\par
\noindent\emph{Textual comparison:} A zero radical is defined. In §19 the theorem's first alternative does not distinguish the cases separated by the proof; source-witness comparison is needed.\par
\noindent{\footnotesize TEXT-COMP-505; pasaže / passages: 20864, 21735, 21737, 21739, 21741, 21743, 21751; \texttt{CORPUS\_FORM\_COMPARISONS.json}.}\par
\end{minipage}\par
\par\medskip\Needspace{7\baselineskip}
\noindent\begin{minipage}{\linewidth}
\hypertarget{TERM-1434}{}
\noindent\textbf{granične členy / graničny integral}\quad{\footnotesize\texttt{TERM-1434}}\par
\noindent\emph{Iztočnik / Source:} Randglieder / Randintegral\par
\noindent\emph{Rusky / Russian:} граничные члены / граничный интеграл\par
\noindent\emph{Ukrajinsky / Ukrainian:} крайові члени / крайовий інтеграл\par
\smallskip\noindent\emph{Razlog:} Granične členy / graničny integral čuva konkretnu granicu oblasti pri parcialnom integrirovanju. Arhiv podčrkava, že dokaz razlikuje graničny integral od identičnosti s divergencijeju. Rabota 13 najprvo zapisuje granične členy; pozdněje izbira funkcije i dostatočno mnogo izvodnyh nulovymi na granici, tako že graničny integral izčezaje. To izčezanje zato jest poslědok uslovij, ne opravdanje bezuslovno vymazati člen. Ruska paralela ima graničn-, ukrajinska krajov-; medžuslovjansky izbor čuva prvu semju. Arhiv ne davaje oceny, ktora by dokazovala, že druga semja jest nepravilna ili menje razumljiva.\par
\noindent\emph{Reason:} Granične členy / graničny integral keeps the boundary of the integration region explicit. The archive emphasizes the proof's distinction between a boundary integral and a divergence identity. Paper 13 first records the boundary terms; later it chooses functions and sufficiently many derivatives to vanish on the boundary, making the boundary integral vanish. That disappearance follows from conditions, not permission to discard the term unconditionally. The Russian comparandum has graničn- and the Ukrainian krajov-; the Interslavic choice retains the former family. No preserved ranking establishes the latter as incorrect or less intelligible.\par
\noindent{\footnotesize Izvorne zapisy / Original records: HIST-g075-011.}\par
\smallskip\noindent\emph{Sravnjenje s tekstom:} Funkcije i potrěbne izvodne izčezajut na granici. Divergencija i graničny integral imajut različne roli.\par
\noindent\emph{Textual comparison:} The functions and required derivatives vanish at the boundary. Divergence and boundary integral have different roles.\par
\noindent{\footnotesize TEXT-COMP-506; pasaže / passages: 9225, 9401, 9457, 9743; \texttt{CORPUS\_FORM\_COMPARISONS.json}.}\par
\end{minipage}\par
\par\medskip\Needspace{7\baselineskip}
\noindent\begin{minipage}{\linewidth}
\hypertarget{TERM-1435}{}
\noindent\textbf{krajeva zadacha / vsude konečne integraly}\quad{\footnotesize\texttt{TERM-1435}}\par
\noindent\emph{Iztočnik / Source:} Randwertproblem / überall endliche Integrale\par
\noindent\emph{Rusky / Russian:} краевая задача / всюду конечные интегралы\par
\noindent\emph{Ukrajinsky / Ukrainian:} крайова задача / всюди скінченні інтеграли\par
\smallskip\noindent\emph{Razlog:} Krajeva zadacha / vsude konečne integraly zadržuje istorijsku analitičnu formulaciju. Arhiv izrično odbacaje jeje avtomatično zaměnjenje občim izrazom holomorfne integraly, zato že sravnjenje Noether jest istorijsko i formalno. Rabota 14 povezuje dokaz obstajanja vsude konečnyh integralov s krajevoj zadacheju, a potom postavjaje analogiju s aritmetičnym teoremom obstajanja prostyh idealov s predpisanymi harakterami. To ne jest tvrđenje, že dva matematična predmeta sut identične. Vsude ograničuje města, gde integral ostavaje konečny; ne znači, že jest samo konečno mnogo integralov. Forma zadacha jest zachovano pisanje zapisa i teksta, ne samostalno potvrdženje jezyčnoj normy.\par
\noindent\emph{Reason:} Krajeva zadacha / vsude konečne integraly retains the historical analytic formulation. The record expressly rejects automatically replacing it with 'holomorphic integrals', because Noether's comparison is historical and formal. Paper 14 connects existence of everywhere-finite integrals with a boundary-value problem, then draws an analogy with an arithmetic existence theorem for prime ideals with prescribed characters. It does not identify the two mathematical objects. Everywhere specifies where an integral remains finite; it does not assert that there are only finitely many integrals. Zadacha is the spelling preserved in the record and corpus, not independent confirmation of a language norm.\par
\noindent{\footnotesize Izvorne zapisy / Original records: HIST-g076-025.}\par
\smallskip\noindent\emph{Sravnjenje s tekstom:} Analitičny opis i aritmetična analogija ostavajut različne. Pytanje jest o svojstvu integralov, ne o jih čislu.\par
\noindent\emph{Textual comparison:} The analytic description and arithmetic analogy remain distinct. The property concerns integrals, not their number.\par
\noindent{\footnotesize TEXT-COMP-507; pasaže / passages: 10033; \texttt{CORPUS\_FORM\_COMPARISONS.json}.}\par
\end{minipage}\par
\par\medskip\Needspace{7\baselineskip}
\noindent\begin{minipage}{\linewidth}
\hypertarget{TERM-1436}{}
\noindent\textbf{rang}\quad{\footnotesize\texttt{TERM-1436}}\par
\noindent\emph{Iztočnik / Source:} Rang\par
\noindent\emph{Rusky / Russian:} ранг\par
\noindent\emph{Ukrajinsky / Ukrainian:} ранг\par
\smallskip\noindent\emph{Razlog:} Rang jest zajedny internacionalny termin, ktory arhiv veže k modulam, matricam i elementarnym děliteljam. V rabotě 22 on jest maksimalno čislo linearno nezavisnyh form; koeficientna oblast jest izrično zadana. V koncu raboty 19 nerazložima klasa, odgovarjajuča dělitelju 0, davaje rang po prědstavjenju elementarnyh děliteljev. Ne govori se, že rang sam mora byti 0: nuly v konečnej česti sistema označajut granicu nenulovyh děliteljev. Izbor kratkogo korenja podpiraje kontinuitet, ale ne ovlastnjaje prenesti tuty modulny smysl na vsak drugy pojem, nazvany rangom. Tri svezane zapisy prěvoda ne sut tri nezavisne definicije.\par
\noindent\emph{Reason:} Rang is the shared international term linked in the records to modules, matrices and elementary divisors. In Paper 22 it is the maximum number of linearly independent forms, with the coefficient domain expressly specified. At the end of Paper 19, the irreducible class corresponding to divisor 0 records the rank through the elementary-divisor representation. This does not say that the rank itself is zero: the terminal zeros mark the limit of the nonzero divisors. The short stem supports continuity but does not license transferring this module sense to every other notion called rank. Three connected translation records are not three independent definitions.\par
\noindent{\footnotesize Izvorne zapisy / Original records: HIST-g081-103, HIST-g085-003, HIST-g092-003.}\par
\smallskip\noindent\emph{Sravnjenje s tekstom:} Rang mjeri linearnu nezavisnost. Klasa, odgovarjajuča dělitelju 0, koduje rang, ne zahtěvaje nulovy rang.\par
\noindent\emph{Textual comparison:} Rank measures linear independence. The class associated with divisor 0 encodes rank; it does not require rank zero.\par
\noindent{\footnotesize TEXT-COMP-508; pasaže / passages: 13719, 13720, 13721, 13722, 13723, 13724, 13732, 13735, 13766, 13768, 14226, 14230, 14231; \texttt{CORPUS\_FORM\_COMPARISONS.json}.}\par
\end{minipage}\par
\par\medskip\Needspace{7\baselineskip}
\noindent\begin{minipage}{\linewidth}
\hypertarget{TERM-1437}{}
\noindent\textbf{rang vzhodno k $K$}\quad{\footnotesize\texttt{TERM-1437}}\par
\noindent\emph{Iztočnik / Source:} Rang in bezug auf $K$\par
\noindent\emph{Rusky / Russian:} ранг относительно $K$\par
\noindent\emph{Ukrajinsky / Ukrainian:} ранг відносно $K$\par
\smallskip\noindent\emph{Razlog:} Rang vzhodno k $K$ nosi izričnu otnosnost k koeficientnomu polju. To jest zachovany razlog izbora; prědlog vzhodno k v nynějšem tekstu ustupil po odnošenju k. V rabotě 34, §10, dva prave idealy v trěhrazměrnom sistemu imajut različny rang nad $K$: jeden poraždajut $e_1,u$, drugy $e_2$. To dokazuje, že oni ne mogut byti operatorno izomorfne v razgledanom kontekstu. Izraz ne govori samo o različnom čislu imenovanyh generatorov; relevantna jest linearna nezavisnost nad istom poljem. Prěnos opravdanoj predložnoj konstrukcije mora čuvati $K$, stranu ideala i konkretno dějstvo operatorov.\par
\noindent\emph{Reason:} Rang vzhodno k $K$ makes relativity to the coefficient field explicit; that is the preserved rationale. The current text replaces vzhodno k with po odnošenju k. In Paper 34, §10, two right ideals in a three-dimensional system have different ranks over $K$: one is spanned by $e_1,u$, the other by $e_2$. This prevents operator isomorphism in the specified context. The point is not merely that their displayed generating lists have different lengths; it is linear independence over the same field. Transferring the justified prepositional construction must retain $K$, the side of the ideals and the operator action.\par
\noindent{\footnotesize Izvorne zapisy / Original records: HIST-g188-005.}\par
\smallskip\noindent\emph{Sravnjenje s tekstom:} Različny rang po odnošenju k $K$ izključaje navedenu operatornu izomorfiju.\par
\noindent\emph{Textual comparison:} Different ranks over $K$ exclude the stated operator isomorphism.\par
\noindent{\footnotesize TEXT-COMP-509; pasaže / passages: 20638, 20639, 20640, 20641, 20642, 20643, 20644, 20645, 20646, 20647, 20648, 20649, 20650, 20651, 20652, 20653, 20654, 20655, 20656, 20657, 20658, 20659, 20660, 20661, 20662, 20663; \texttt{CORPUS\_FORM\_COMPARISONS.json}.}\par
\end{minipage}\par
\par\medskip\Needspace{7\baselineskip}
\noindent\begin{minipage}{\linewidth}
\hypertarget{TERM-1438}{}
\noindent\textbf{rang četyri/tri/dva/jedin}\quad{\footnotesize\texttt{TERM-1438}}\par
\noindent\emph{Iztočnik / Source:} Rang vier/drei/zwei/eins\par
\noindent\emph{Rusky / Russian:} ранг четыре/три/два/один\par
\noindent\emph{Ukrajinsky / Ukrainian:} ранг чотири/три/два/один\par
\smallskip\noindent\emph{Razlog:} Rang četyri/tri/dva/jedin zadržuje čislovu klasifikaciju iz raboty 10, §4. Arhiv nazyvaje izraz prěmym i centralnym za argument; sam toj opis ne davaje jednu universalnu definiciju ranga. Tutaj pri $z=x+iy$ i $\varphi(z)=X+iY$ rang jest četyri minus čislo nezavisnyh realnyh linearnyh relacij medžu $x,y,X,Y$, ktore važat za funkciju. Ne ide o rang kompleksnogo Jacobiana ni o stepen poljnogo razširenja. Prěvod mora ohraniti obratny odnos: nula relacij davaje rang četyri, a tri relacije rang jedin. Prosto navedenje čisel bez togo pravila izgubilo by definiciju, potrěbnu dokazu.\par
\noindent\emph{Reason:} Rang četyri/tri/dva/jedin retains the numerical classification in Paper 10, §4. The archive calls the expression direct and central to the argument, but that description supplies no universal definition of rank. Here, for $z=x+iy$ and $\varphi(z)=X+iY$, rank is four minus the number of independent real linear relations among $x,y,X,Y$ that hold for the function. It is not a complex Jacobian rank or field-extension degree. Translation must preserve the reverse relationship: no relation gives rank four, whereas three relations give rank one. Merely listing the numbers would lose the definition needed by the proof.\par
\noindent{\footnotesize Izvorne zapisy / Original records: HIST-g072-004.}\par
\smallskip\noindent\emph{Sravnjenje s tekstom:} Definicija odčityvaje čislo realnyh linearnyh relacij od četyri; ne govori o kompleksnom Jacobianu.\par
\noindent\emph{Textual comparison:} The definition subtracts the number of real linear relations from four; it is not a complex Jacobian rank.\par
\noindent{\footnotesize TEXT-COMP-510; pasaže / passages: 8611, 8612, 8613, 8614, 8615; \texttt{CORPUS\_FORM\_COMPARISONS.json}.}\par
\end{minipage}\par
\par\medskip\Needspace{7\baselineskip}
\noindent\begin{minipage}{\linewidth}
\hypertarget{TERM-1439}{}
\noindent\textbf{rangi}\quad{\footnotesize\texttt{TERM-1439}}\par
\noindent\emph{Iztočnik / Source:} Rangzahlen\par
\noindent\emph{Rusky / Russian:} ранги\par
\noindent\emph{Ukrajinsky / Ukrainian:} ранги\par
\smallskip\noindent\emph{Razlog:} Rangi jest skračenie německogo Rangzahlen, izbrano radi sovrěmennogo stručnogo izloženja. Arhiv izrično zachovuje izvorno složeno slovo za sravnjenje, ne tvori novu čislovu veličinu. V rabotě 31, §2, sravnjujut se rang komponenty po odnošenju k $P$ i rang po odnošenju k izomorfnomu polju $P_i$. Ravnanje rangov slěduje iz prěnosa linearnyh relacij pomoču idempotenta $\varepsilon_i$. Kratši plural ne směje prikryti, že ide o dvě početno različne osnovne polja i izričny argument jih sravnjenja. Ujednačenje jezika zato čuva vse parametry, hoti ne povtorjuje doslovno čest -zahlen.\par
\noindent\emph{Reason:} Rangi shortens German Rangzahlen for concise modern exposition. The archive explicitly preserves the source compound for comparison; it introduces no new numerical quantity. Paper 31, §2, compares a component's rank over $P$ with its rank over the isomorphic field $P_i$. Equality follows by transporting linear relations through the idempotent $\varepsilon_i$. The shorter plural must not conceal the two initially distinct scalar fields or the argument comparing them. Linguistic modernization retains all parameters even though it does not reproduce the -zahlen component literally.\par
\noindent{\footnotesize Izvorne zapisy / Original records: HIST-g122-007.}\par
\smallskip\noindent\emph{Sravnjenje s tekstom:} Ravnanje rangov jest opravdano prěnosom relacij pomoču idempotenta, ne kratšim nazvanjem.\par
\noindent\emph{Textual comparison:} Equality of ranks is justified by transporting relations through an idempotent, not by the shortened name.\par
\noindent{\footnotesize TEXT-COMP-511; pasaže / passages: 17434, 17524; \texttt{CORPUS\_FORM\_COMPARISONS.json}.}\par
\end{minipage}\par
\par\medskip\Needspace{7\baselineskip}
\noindent\begin{minipage}{\linewidth}
\hypertarget{TERM-1440}{}
\noindent\textbf{racionalno zavisny / nezavisny}\quad{\footnotesize\texttt{TERM-1440}}\par
\noindent\emph{Iztočnik / Source:} rational abhängig/unabhängig\par
\noindent\emph{Rusky / Russian:} рационально зависимый / независимый\par
\noindent\emph{Ukrajinsky / Ukrainian:} раціонально залежний / незалежний\par
\smallskip\noindent\emph{Razlog:} Racionalno zavisny / nezavisny čuva razlikovanje, na ktorom osnovana jest racionalna baza v rabotě 10, §2. Pri danym dobrom poredku čislo ili jest racionalno izrazimo prěz konečno mnogo poprědnyh, ili ne jest; koeficienty racionalnyh funkcij takože sut racionalne čisla. Potom tekst oddělno razlikuje algebraičnu zavisimost od algebraičnoj nezavisimosti po odnošenju k polju početnogo odrězka. Arhiv govori, že para zrcali pozdnějšu algebraičnu paru; to znači paralelnu jezyčnu konstrukciju, ne ravnost dvoh matematičnyh relacij. Poredok, konečnost upotrěbljenyh argumentov i oblast koeficientov morajut ostati eksplicitne. Nezavisny ne znači prosto neočekivany ili nesvezany.\par
\noindent\emph{Reason:} Racionalno zavisny / nezavisny preserves the distinction underlying the rational basis in Paper 10, §2. Given a well-order, a number either is rationally expressible in finitely many predecessors or is not; the rational functions' coefficients are themselves rational numbers. The text then separately distinguishes algebraic dependence and independence over an initial-segment field. The archive says that the pair mirrors the later algebraic pair: this is parallel linguistic construction, not identity of the two mathematical relations. The order, finite number of arguments used and coefficient domain must remain explicit. Independent does not merely mean unexpected or unrelated.\par
\noindent{\footnotesize Izvorne zapisy / Original records: HIST-g070-006.}\par
\smallskip\noindent\emph{Sravnjenje s tekstom:} Racionalna izrazimost koristi konečno mnogo poprědnyh elementov i racionalne koeficienty. Algebraična zavisimost jest daljše oddělno pytanje.\par
\noindent\emph{Textual comparison:} Rational expressibility uses finitely many predecessors and rational coefficients. Algebraic dependence is a further, separate question.\par
\noindent{\footnotesize TEXT-COMP-512; pasaže / passages: 8487, 8489, 8491; \texttt{CORPUS\_FORM\_COMPARISONS.json}.}\par
\end{minipage}\par
\par\medskip\Needspace{7\baselineskip}
\noindent\begin{minipage}{\linewidth}
\hypertarget{TERM-1441}{}
\noindent\textbf{racionalno nezavisimy}\quad{\footnotesize\texttt{TERM-1441}}\par
\noindent\emph{Iztočnik / Source:} rational unabh?ngig [Čestično nezachovane bukvy / Partly unrecoverable lettering]\par
\noindent\emph{Rusky / Russian:} [Bukvy ne sut zachovane / Lettering not recoverable]\par
\noindent\emph{Ukrajinsky / Ukrainian:} [Bukvy ne sut zachovane / Lettering not recoverable]\par
\smallskip\noindent\emph{Razlog:} Racionalno nezavisimy tu ne znači algebraično nezavisimy. V rabotě 9, §6, čislo ne može byti racionalno vyraženo prěz konečno mnogo poprědnyh čisel v danom dobrom uporjadkovanju. Ono ipak može izpolnjati algebraično uravnenje nad poljem poprědnyh čisel. Navedeni priměr s korenjem pokazuje, že poredok elementov jest važny; kriterij ne jest nezavisnost libovoljnogo neuporjadkovanogo mnošestva. Arhiv izrično čuva razlikovanje racionalnoj i algebraičnoj zavisimosti. Německy naslov imaje nezachovanu bukvu, a ruske i ukrajinske bukvy ne sut zachovane; iz njih zato nemožno izvesti jezyčno sravnjenje. Nynějša kontekstna interpretacija ne vosstanavja te nezachovane bukvy.\par
\noindent\emph{Reason:} Racionalno nezavisimy here does not mean algebraically independent. In Paper 9, §6, a number cannot be expressed rationally through finitely many predecessors in a specified well-order. It may nevertheless satisfy an algebraic equation over the field generated by those predecessors. The example involving a root shows that the order matters; this is not independence of an arbitrary unordered set. The archive expressly preserves rational versus algebraic dependence. One letter of the German heading is unrecoverable, and the Russian and Ukrainian lettering is lost, so those fields cannot support a linguistic comparison. The present contextual interpretation does not restore their missing letters.\par
\noindent{\footnotesize Izvorne zapisy / Original records: HIST-g064-004.}\par
\smallskip\noindent\emph{Sravnjenje s tekstom:} Dobro uporjadkovanje i koeficientno polje ograničajut smysl nezavisnosti.\par
\noindent\emph{Textual comparison:} The well-order and coefficient field delimit the sense of independence.\par
\noindent{\footnotesize TEXT-COMP-513; pasaže / passages: 7727, 7728, 7729, 7730, 7731, 7732, 7734, 7735, 7736, 7737, 7738, 7739, 7740, 7741, 7742, 7743, 7744, 7745, 7746, 7747, 7748, 7749, 8491; \texttt{CORPUS\_FORM\_COMPARISONS.json}.}\par
\end{minipage}\par
\par\medskip\Needspace{7\baselineskip}
\noindent\begin{minipage}{\linewidth}
\hypertarget{TERM-1442}{}
\noindent\textbf{racionalno- i algebraično-drobne funkcije}\quad{\footnotesize\texttt{TERM-1442}}\par
\noindent\emph{Iztočnik / Source:} rational- und algebraisch-gebrochene Funktionen\par
\noindent\emph{Rusky / Russian:} рационально- и алгебраически-дробные функции\par
\noindent\emph{Ukrajinsky / Ukrainian:} раціонально- та алгебраїчно-дробові функції\par
\smallskip\noindent\emph{Razlog:} V racionalno- i algebraično-drobne funkcije zajedny končny člen drobne funkcije nosi dva različne opisy. Skračenje prvogo složenogo pridavnogo slova s veznikom čuva oba pojma; ne smě ih stopiti v jednu klasu racionalnyh funkcij. V rabotě 9, §2, svojstva 4) i 5) sut različne zahtěvy izključenja. Konstrukcija s racionalnymi funkcionalami pokazuje, že abstraktne svojstva oblasti ne potrěbujut ni jednogo iz njih. Arhiv zato davaje matematičnomu razlikovanju prědnost nad kratkostju naslova. Drobny tu označaje protivnost odgovarjajučej cělosti, ne prosto pojavjenje znaka děljenja v slučajno izbranom zapisu.\par
\noindent\emph{Reason:} In racionalno- i algebraično-drobne funkcije, two distinct modifiers share the ending drobne funkcije. Suspending the first compound adjective at the hyphen preserves both concepts; it must not collapse them into a single class of rational functions. In Paper 9, §2, properties 4 and 5 are different exclusion requirements. The construction using rational functionals shows that neither is required by the abstract properties of the domain. The archive thus gives mathematical differentiation priority over a shorter label. Drobny contrasts with the relevant notion of integrality, not simply the appearance of a division sign in an arbitrarily chosen expression.\par
\noindent{\footnotesize Izvorne zapisy / Original records: HIST-g060-010.}\par
\smallskip\noindent\emph{Sravnjenje s tekstom:} Drobne racionalne i drobne algebraične funkcije ostavajut različne kategorije.\par
\noindent\emph{Textual comparison:} Fractional rational and fractional algebraic functions remain distinct categories.\par
\noindent{\footnotesize TEXT-COMP-514; pasaže / passages: 7147, 7148, 7152, 7153, 7234, 7242, 7244; \texttt{CORPUS\_FORM\_COMPARISONS.json}.}\par
\end{minipage}\par
\par\medskip\Needspace{7\baselineskip}
\noindent\begin{minipage}{\linewidth}
\hypertarget{TERM-1443}{}
\noindent\textbf{racionalno-c?lo adjungovati [Čestično nezachovane bukvy / Partly unrecoverable lettering]}\quad{\footnotesize\texttt{TERM-1443}}\par
\noindent\emph{Iztočnik / Source:} rational-ganz adjungieren\par
\noindent\emph{Rusky / Russian:} [Bukvy ne sut zachovane / Lettering not recoverable]\par
\noindent\emph{Ukrajinsky / Ukrainian:} [Bukvy ne sut zachovane / Lettering not recoverable]\par
\smallskip\noindent\emph{Razlog:} Adjungovati jest glagol za operaciju, ktorej imenno nazvanje jest adjunkcija. Arhiv bira jego kako kratky tehničny kalk i povezuje jego s terminologijeju §4; tvrđenje o obyčnoj ruskoj i ukrajinskoj upotrěbě ne možno prověriti iz izgubljenyh bukv tih dvoh polej. V §5 racionalno-cělo adjungovanje davaje polinomy od prisojedinjanyh elementov nad početkovoju oblastju. Ono samo ne označaje prisojedinjenje vsih drobov niti algebraičnu integralnu zamknjenost. Zapis racionalno-c?lo čestično poškodženy ostavaje takym v arhivnom sravnjenju; čitljiva forma racionalno-cělo pochodi iz konkretnoj pasaže prěvoda. Kontinuitet nazvanja ne odpravja dodatnyh uslovij potrěbnyh za konstruovanu oblast.\par
\noindent\emph{Reason:} Adjungovati is the verb for the operation named adjunkcija. The archive selects it as a compact technical calque continuous with §4; its claim about ordinary Russian and Ukrainian usage cannot be checked from the lost lettering in those fields. In §5, rational-integral adjoining forms polynomials in the adjoined elements over the initial domain. It does not itself mean adjoining all fractions or taking algebraic integral closure. The partly damaged archival form racionalno-c?lo remains unchanged in the comparison; the readable racionalno-cělo comes from the actual translated passage. Terminological continuity does not remove the additional conditions required of the resulting domain.\par
\noindent{\footnotesize Izvorne zapisy / Original records: HIST-g063-010.}\par
\smallskip\noindent\emph{Sravnjenje s tekstom:} Definicija operacije daje kriterij za glagol i imenno nazvanje; jezyčno odobrenje ne slěduje iz definicije.\par
\noindent\emph{Textual comparison:} The definition supplies the criterion for the verb and noun; linguistic approval does not follow from that definition.\par
\noindent{\footnotesize TEXT-COMP-515; pasaže / passages: 7461, 7462, 7463, 7464, 7465, 7466, 7467, 7468, 7469, 7470, 7471, 7472, 7520, 7521, 7522, 7523, 7524, 7525, 7680, 7681, 7682, 7683, 7684, 7685, 7686, 7687, 7688, 7689, 7690, 7691, 7692, 7693, 7694, 7695, 7696, 7697, 7785, 7786, 7787, 7788, 7789, 7790, 7791, 7792, 7793, 7794, 7810, 7811, 7812, 7813, 7814, 7815, 7816, 7817, 7818; \texttt{CORPUS\_FORM\_COMPARISONS.json}.}\par
\end{minipage}\par
\par\medskip\Needspace{7\baselineskip}
\noindent\begin{minipage}{\linewidth}
\hypertarget{TERM-1444}{}
\noindent\textbf{racionalno-c?la adjunkcija [Čestično nezachovane bukvy / Partly unrecoverable lettering]}\quad{\footnotesize\texttt{TERM-1444}}\par
\noindent\emph{Iztočnik / Source:} rational-ganze Adjunktion\par
\noindent\emph{Rusky / Russian:} [Bukvy ne sut zachovane / Lettering not recoverable]\par
\noindent\emph{Ukrajinsky / Ukrainian:} [Bukvy ne sut zachovane / Lettering not recoverable]\par
\smallskip\noindent\emph{Razlog:} Jediny zachovany razlog za tutu formu adjunkcija jest kontinuitet s §§4–5 i kratkost internacionalizma; zapis izrično ostavja jezyčno potvrdženje otvorjenym. Zato ne možno predstavljati izbor kako už priznanu normu. Čestično poškodženy pridavny oblik v arhivu ne zaměnjaje se tiho drugim. §6 davaje čitljivu upotrěbu: racionalno-cěla adjunkcija vodi od početkovoj oblasti do prstena polinomnyh izrazov v dodatnom sistemu, s izključenjem nedopustimyh algebraično-drobnyh čisel. To jest smyslovy kriterij za upotrěbu termina. Arhiv ne sodrživaje čislove oceny kandidatov niti dokumentovano rěšenje jezyčnogo avtoriteta.\par
\noindent\emph{Reason:} The sole preserved rationale for this form of adjunkcija is continuity with §§4–5 and the compactness of the international term; the record expressly leaves linguistic authority review pending. The choice therefore cannot be presented as an already recognized norm. Its partly damaged archival adjective is not silently replaced. Section 6 provides a readable occurrence: rational-integral adjoining takes the initial domain to the ring of polynomial expressions in an additional system, while excluding inadmissible algebraically fractional numbers. That supplies the semantic condition for using the term. The archive contains neither numerical candidate scores nor a documented ruling by a language authority.\par
\noindent{\footnotesize Izvorne zapisy / Original records: HIST-g064-010.}\par
\smallskip\noindent\emph{Sravnjenje s tekstom:} Definicija operacije daje kriterij za glagol i imenno nazvanje; jezyčno odobrenje ne slěduje iz definicije.\par
\noindent\emph{Textual comparison:} The definition supplies the criterion for the verb and noun; linguistic approval does not follow from that definition.\par
\noindent{\footnotesize TEXT-COMP-515; pasaže / passages: 7461, 7462, 7463, 7464, 7465, 7466, 7467, 7468, 7469, 7470, 7471, 7472, 7520, 7521, 7522, 7523, 7524, 7525, 7680, 7681, 7682, 7683, 7684, 7685, 7686, 7687, 7688, 7689, 7690, 7691, 7692, 7693, 7694, 7695, 7696, 7697, 7785, 7786, 7787, 7788, 7789, 7790, 7791, 7792, 7793, 7794, 7810, 7811, 7812, 7813, 7814, 7815, 7816, 7817, 7818; \texttt{CORPUS\_FORM\_COMPARISONS.json}.}\par
\end{minipage}\par
\par\medskip\Needspace{7\baselineskip}
\noindent\begin{minipage}{\linewidth}
\hypertarget{TERM-1445}{}
\noindent\textbf{racionalno-cěla adjunkcija}\quad{\footnotesize\texttt{TERM-1445}}\par
\noindent\emph{Iztočnik / Source:} rational-ganze Adjunktion\par
\noindent\emph{Rusky / Russian:} рационально-целое присоединение\par
\noindent\emph{Ukrajinsky / Ukrainian:} раціонально-ціле приєднання\par
\smallskip\noindent\emph{Razlog:} Racionalno-cěla adjunkcija jest čitljiva forma zachovana v zapisah za §§4 i 7 raboty 9. Rusko prisojedinenije i ukrajinsko pryjednannja predstavljajut rodne izrazove rodiny; ISV zapis bira internacionalno adjunkcija, zato že smatraje go kratšim i jasnějšim. Ne navodi konkretno odbrošenu rodnu ISV frazu niti test razuměvanja, zato toj razlog ostavaje kvalitativny. §4 definira operaciju kako tvorjenje polinomov v dodatnyh elementah; slědujuča algebraično-cěla adjunkcija jest druga operacija. Pri sravnjenju s poškodženymi zapisami iz §§5–6 kontinuitet smysla jest podprěny tekstom, ale različna sostojanja arhivnyh zapisov ostavajut vidime.\par
\noindent\emph{Reason:} Racionalno-cěla adjunkcija is the readable form preserved for §§4 and 7 of Paper 9. Russian prisoyedineniye and Ukrainian pryiednannia represent native expression families; the ISV record chooses international adjunkcija as shorter and clearer. It supplies neither a particular rejected native ISV phrase nor a comprehension test, so that judgment remains qualitative. Section 4 defines the operation by forming polynomials in the additional elements; the subsequent algebraic-integral adjoining is a different operation. The text supports semantic continuity with the damaged records from §§5–6, but their different states of preservation remain visible.\par
\noindent{\footnotesize Izvorne zapisy / Original records: HIST-g062-007, HIST-g065-005.}\par
\smallskip\noindent\emph{Sravnjenje s tekstom:} Definicija operacije daje kriterij za glagol i imenno nazvanje; jezyčno odobrenje ne slěduje iz definicije.\par
\noindent\emph{Textual comparison:} The definition supplies the criterion for the verb and noun; linguistic approval does not follow from that definition.\par
\noindent{\footnotesize TEXT-COMP-515; pasaže / passages: 7461, 7462, 7463, 7464, 7465, 7466, 7467, 7468, 7469, 7470, 7471, 7472, 7520, 7521, 7522, 7523, 7524, 7525, 7680, 7681, 7682, 7683, 7684, 7685, 7686, 7687, 7688, 7689, 7690, 7691, 7692, 7693, 7694, 7695, 7696, 7697, 7785, 7786, 7787, 7788, 7789, 7790, 7791, 7792, 7793, 7794, 7810, 7811, 7812, 7813, 7814, 7815, 7816, 7817, 7818; \texttt{CORPUS\_FORM\_COMPARISONS.json}.}\par
\end{minipage}\par
\par\medskip\Needspace{7\baselineskip}
\noindent\begin{minipage}{\linewidth}
\hypertarget{TERM-1446}{}
\noindent\textbf{racionalno-drobny / algebraično-drobny}\quad{\footnotesize\texttt{TERM-1446}}\par
\noindent\emph{Iztočnik / Source:} rational-gebrochen / algebraisch-gebrochen\par
\noindent\emph{Rusky / Russian:} рационально-дробный / алгебраически-дробный\par
\noindent\emph{Ukrajinsky / Ukrainian:} раціонально-дробовий / алгебраїчно-дробовий\par
\smallskip\noindent\emph{Razlog:} Dvojica racionalno-drobny / algebraično-drobny zachovaje protivnost cěly/drobny pri dvoh tipah funkcij. Uvod raboty 9 oddělno navodi cěle racionalne i cěle algebraične funkcije, potom izključaje odgovarjajuče drobne funkcije kako svojstva 4) i 5). Korenj drob- zato ne jest izložiteljska ozdoba: on oprěděljaje, ktore elementy naležet oblasti. Arhiv davaje prědnost vernosti toj protivnosti, s paralelnymi ruskimi i ukrajinskimi složenymi pridavnymi slovami. Ne kaže, že slovnična dolga forma jest vsegda potrěbna; može se skratiti samo ako oba opisy i oblast koeficientov ostavajut jasne.\par
\noindent\emph{Reason:} The pair racionalno-drobny / algebraično-drobny preserves the integral/fractional contrast for two types of functions. Paper 9's introduction separately lists integral rational and integral algebraic functions, then excludes the corresponding fractional functions as properties 4 and 5. The stem drob- is therefore not ornamental exposition: it determines which elements belong to the domain. The archive prioritizes fidelity to this contrast, with parallel Russian and Ukrainian compound adjectives. It does not require the full long forms in every sentence; they may be abbreviated only while both qualifications and the coefficient domain remain clear.\par
\noindent{\footnotesize Izvorne zapisy / Original records: HIST-g058-009.}\par
\smallskip\noindent\emph{Sravnjenje s tekstom:} Drobne racionalne i drobne algebraične funkcije ostavajut različne kategorije.\par
\noindent\emph{Textual comparison:} Fractional rational and fractional algebraic functions remain distinct categories.\par
\noindent{\footnotesize TEXT-COMP-514; pasaže / passages: 7147, 7148, 7152, 7153, 7234, 7242, 7244; \texttt{CORPUS\_FORM\_COMPARISONS.json}.}\par
\end{minipage}\par
\par\medskip\Needspace{7\baselineskip}
\noindent\begin{minipage}{\linewidth}
\hypertarget{TERM-1447}{}
\noindent\textbf{racionalna baza}\quad{\footnotesize\texttt{TERM-1447}}\par
\noindent\emph{Iztočnik / Source:} Rationalbasis\par
\noindent\emph{Rusky / Russian:} рациональная база; рациональный базис\par
\noindent\emph{Ukrajinsky / Ukrainian:} раціональна база; раціональний базис\par
\smallskip\noindent\emph{Razlog:} Racionalna baza v rabotah 5, 6 i 29 označaje konečnu porodžajuču familiju za racionalne izrazy s koeficientami iz zadannogo polja. Ne znači avtomatično algebraično nezavisnu familiju: rabota 5 oddělno definira minimalnu bazu, a rabota 6 razširjaje nezavisny sistem do racionalnoj bazy. Osm zapisov podprěnajut kratkost i kontinuitet nazvanja, ne osm nezavisnyh jezyčnyh prověrjenij. Ruske i ukrajinske pole sodrživajut obě formy baza i bazis; sam arhiv ne davaje novega izvažanja med nimi. Smyslovy kriterij jest racionalna prědstavimost vsih funkcij. Pri prěnosu v drugy predmet ne smě kratko slovo baza prikryvati razliku medžu porodženjem, linearnom nezavisnostju i algebraičnom nezavisnostju.\par
\noindent\emph{Reason:} In Papers 5, 6 and 29, racionalna baza denotes a finite generating family for rational expressions with coefficients in the specified field. It does not automatically mean an algebraically independent family: Paper 5 separately defines a minimal basis, while Paper 6 extends an independent system to a rational basis. Eight records support brevity and continuity, not eight independent linguistic assessments. The Russian and Ukrainian fields preserve both baza and bazis; the archive offers no new weighing between them. The semantic criterion is rational expressibility of all the functions. Transfer to another subject must not let the short word baza conceal the differences among generation, linear independence and algebraic independence.\par
\noindent{\footnotesize Izvorne zapisy / Original records: HIST-g039-007, HIST-g040-005, HIST-g044-001, HIST-g046-002, HIST-g047-011, HIST-g049-007, HIST-g054-009, HIST-g106-007.}\par
\smallskip\noindent\emph{Sravnjenje s tekstom:} Konečna racionalna baza ne mora byti minimalna baza.\par
\noindent\emph{Textual comparison:} A finite rational basis need not be a minimal basis.\par
\noindent{\footnotesize TEXT-COMP-516; pasaže / passages: 4954, 4956, 4958, 5410, 5412, 5479, 5481, 5482, 5576, 5690, 6454, 6456, 15995, 15997, 15999; \texttt{CORPUS\_FORM\_COMPARISONS.json}.}\par
\end{minipage}\par
\par\medskip\Needspace{7\baselineskip}
\noindent\begin{minipage}{\linewidth}
\hypertarget{TERM-1448}{}
\noindent\textbf{racionalna baza}\quad{\footnotesize\texttt{TERM-1448}}\par
\noindent\emph{Iztočnik / Source:} rationale Basis\par
\noindent\emph{Rusky / Russian:} [Bukvy ne sut zachovane / Lettering not recoverable]\par
\noindent\emph{Ukrajinsky / Ukrainian:} [Bukvy ne sut zachovane / Lettering not recoverable]\par
\smallskip\noindent\emph{Razlog:} V §6 raboty 9 racionalna baza jest sistem čisel s vlastnym uslovjem dobrogo uporjadkovanja, ne samo isto, čo konečna racionalna baza polja funkcij v ranějših rabotah. Vsako čislo mora byti racionalno izrazimo prěz konečno mnogo bazovyh čisel, a nijedno bazovo čislo ne smi byti tako vyraženo prěz svoje poprědniki pri prinajmenje jednom dobrom uporjadkovanju. Početkova definicija dopušćaje algebraične čislove koeficienty; primětka objasnjaje variant s racionalnymi koeficientami. Arhiv bira doslovny i paralelny termin, ale ruski i ukrajinski zapisy ne sut čitljive. Rozhodujuče za upotrěbu jest cela lokalna definicija, ne to, že fraza racionalna baza povtorjaje se drugde.\par
\noindent\emph{Reason:} In Paper 9, §6, racionalna baza is a system of numbers with a particular well-ordering condition, not simply the same object as the finite rational basis of a function field in the earlier papers. Every number must be rationally expressible through finitely many basis numbers, while no basis number is so expressible through its predecessors in at least one well-order. The initial definition allows algebraic-number coefficients; a footnote explains the variant with rational coefficients. The archive chooses a literal parallel term, but its Russian and Ukrainian records are unreadable. The complete local definition governs usage, not merely recurrence of the phrase racionalna baza elsewhere.\par
\noindent{\footnotesize Izvorne zapisy / Original records: HIST-g064-001.}\par
\smallskip\noindent\emph{Sravnjenje s tekstom:} Uslovje poredka i koeficienty morajut ostati izrične.\par
\noindent\emph{Textual comparison:} The ordering condition and coefficients must remain explicit.\par
\noindent{\footnotesize TEXT-COMP-517; pasaže / passages: 7716, 7718, 7719, 7720, 7721, 7722, 7723, 7724, 7725, 7726, 7727, 7728, 7729, 7730, 7731, 7732; \texttt{CORPUS\_FORM\_COMPARISONS.json}.}\par
\end{minipage}\par
\par\medskip\Needspace{7\baselineskip}
\noindent\begin{minipage}{\linewidth}
\hypertarget{TERM-1450}{}
\noindent\textbf{racionalna baza / algebraična baza}\quad{\footnotesize\texttt{TERM-1450}}\par
\noindent\emph{Iztočnik / Source:} rationale Basis / algebraische Basis\par
\noindent\emph{Rusky / Russian:} рациональная база / алгебраическая база\par
\noindent\emph{Ukrajinsky / Ukrainian:} раціональна база / алгебраїчна база\par
\smallskip\noindent\emph{Razlog:} Racionalna baza / algebraična baza ostavajut paralelne konstrukcije v završnom obobčenju raboty 9. Paralelna forma ne znači jednake definicije: racionalno porodženje i algebraično porodženje imajut različne uslovja. V karakteristiki prostej ostaje potrěbno čitati, kogo polja baza jest baza, i pri algebraično-cělyh veličinah izrično prědpoloženje algebraičnoj zamknjenosti. Tekst odpravja dodatno uslovje zato že prvotno polje ne imaje drobnyh veličin v tutom smyslu; ne odpravja vse prědpoloženja konstrukcije. Arhiv bira kontinuitet s ranějšimi paragrafami, ne novu terminologiju za specialnu karakteristiku.\par
\noindent\emph{Reason:} Racionalna baza / algebraična baza remain parallel constructions in Paper 9's concluding generalization. Parallel wording does not make their definitions identical: rational and algebraic generation have different conditions. In prime characteristic, the reader must still identify the field to which the basis belongs and retain the explicit algebraic-closure hypothesis for algebraic-integral elements. The text removes the additional condition because the prime field has no fractional elements in this sense; it does not remove every hypothesis of the construction. The archive chooses continuity with earlier sections, not a new terminology for the special characteristic.\par
\noindent{\footnotesize Izvorne zapisy / Original records: HIST-g068-010.}\par
\smallskip\noindent\emph{Sravnjenje s tekstom:} Otpadenje dodatnogo uslovja ne odpravja ostatnyh hipotez.\par
\noindent\emph{Textual comparison:} Removal of the additional condition does not remove the remaining hypotheses.\par
\noindent{\footnotesize TEXT-COMP-519; pasaže / passages: 8370, 8371, 8372, 8373, 8374, 8376, 8377, 8378, 8379, 8380, 8382, 8383, 8384, 8385, 8386, 8387, 8388, 8389, 8390, 8391, 8392, 8393, 8394; \texttt{CORPUS\_FORM\_COMPARISONS.json}.}\par
\end{minipage}\par
\par\medskip\Needspace{7\baselineskip}
\noindent\begin{minipage}{\linewidth}
\hypertarget{TERM-1451}{}
\noindent\textbf{racionalna baza vsih čisel}\quad{\footnotesize\texttt{TERM-1451}}\par
\noindent\emph{Iztočnik / Source:} rationale Basis aller Zahlen\par
\noindent\emph{Rusky / Russian:} рациональный базис всех чисел\par
\noindent\emph{Ukrajinsky / Ukrainian:} раціональна база всіх чисел\par
\smallskip\noindent\emph{Razlog:} Racionalna baza vsih čisel v rabotě 10, č. 2, ima izrično ograničenu oblast: vsa realna i kompleksna čisla, s racionalnymi čislovymi koeficientami. Vsih ne davaje tvrđenje o vsevozmožnyh čislovyh strukturah. Definicija dalej zahtěvaje dobro uporjadkovanje, pri ktorom bazovo čislo ne može byti racionalno vyraženo prěz konečno mnogo poprědnyh. Arhiv prodolžaje termin iz raboty 9, ale koeficientne konvencije morajut byti čitane nanovo. Zapis smatraje rusko bazis idiomatičnym, a ukrajinsko baza i ISV baza doslědnymi s ranějšim prěvodom; ne dokazuje tym obču višestepennost jednogo nazvanja. Genitiv vsih čisel čuva obseg tvrđenja, ne vnosit novy tip bazy.\par
\noindent\emph{Reason:} In Paper 10, part 2, racionalna baza vsih čisel has an explicit domain: all real and complex numbers, with rational-number coefficients. Vsih does not assert a claim about every conceivable number structure. The definition further requires a well-order in which a basis number cannot be expressed rationally through finitely many predecessors. The archive continues Paper 9's term, but the coefficient convention must be read afresh. It describes Russian bazis as idiomatic and Ukrainian and ISV baza as consistent with the earlier translation; this does not prove a general superiority of either noun. The genitive vsih čisel preserves the claim's scope without introducing a new type of basis.\par
\noindent{\footnotesize Izvorne zapisy / Original records: HIST-g070-003.}\par
\smallskip\noindent\emph{Sravnjenje s tekstom:} Čisla i dopustime koeficienty sut izrično oprěděljene.\par
\noindent\emph{Textual comparison:} The numbers and admissible coefficients are explicitly specified.\par
\noindent{\footnotesize TEXT-COMP-520; pasaže / passages: 8487, 8489, 8491; \texttt{CORPUS\_FORM\_COMPARISONS.json}.}\par
\end{minipage}\par
\par\medskip\Needspace{7\baselineskip}
\noindent\begin{minipage}{\linewidth}
\hypertarget{TERM-1452}{}
\noindent\textbf{racionalna baza s dodatnym uslovjem}\quad{\footnotesize\texttt{TERM-1452}}\par
\noindent\emph{Iztočnik / Source:} rationale Basis mit Nebenbedingung\par
\noindent\emph{Rusky / Russian:} [Bukvy ne sut zachovane / Lettering not recoverable]; рациональная база с дополнительным условием\par
\noindent\emph{Ukrajinsky / Ukrainian:} [Bukvy ne sut zachovane / Lettering not recoverable]; раціональна база з додатковою умовою\par
\smallskip\noindent\emph{Razlog:} Dodatno uslovje jest hipoteza, ne sporedna primětka. §6 zahtěvaje, da polinomy od bazovyh čisel s dopustimymi cělymi koeficientami ne porodet algebraično-drobno čislo; pri daljšem adjungovanju isto izključenje ograničaje sistem dodatnyh elementov. Navedeny priměr s polovinoju obratnogo kvadratnogo korenja pokazuje, že uslovje ne jest prazdno. Tri zapisy povezuju §§6–8 i čuvaju kvantifikator vsaka pri dopolnjenju algebraičnoj bazy do racionalnoj. Arhiv sodrživaje i poškodžene, i čitljive ruske i ukrajinske formy; čitljive ne dokazyvajut točne bukvy poškodženyh. Izbor s dodatnym uslovjem daje čitatelju jasnu hipotezu za prověrjenje.\par
\noindent\emph{Reason:} Dodatno uslovje is a hypothesis, not a passing aside. Section 6 requires that polynomials in the basis numbers with the admissible integral coefficients do not generate an algebraically fractional number; in subsequent adjoining, the same exclusion constrains the additional system. The example involving half the reciprocal square root shows that this condition is substantive. Three records connect §§6–8 and preserve the quantifier every when completing an algebraic basis to a rational one. The archive contains both damaged and readable Russian and Ukrainian forms; the readable forms do not establish the exact lettering of the damaged ones. S dodatnym uslovjem gives the reader an explicit hypothesis to examine.\par
\noindent{\footnotesize Izvorne zapisy / Original records: HIST-g064-005, HIST-g065-001, HIST-g066-008.}\par
\smallskip\noindent\emph{Sravnjenje s tekstom:} Izključenje algebraično-drobnyh čisel jest hipoteza; vsaka baza ne znači samo jeden priměr.\par
\noindent\emph{Textual comparison:} Excluding algebraically fractional numbers is a hypothesis; every basis does not mean just one example.\par
\noindent{\footnotesize TEXT-COMP-521; pasaže / passages: 7171, 7174, 7179, 7750, 7751, 7752, 7753, 7754, 7755, 7756, 7757, 7758, 7759, 7785, 7786, 7787, 7788, 7789, 7790, 7791, 7802, 7803, 7804, 8014, 8015, 8016, 8017, 8018, 8019, 8020, 8021, 8022, 8023; \texttt{CORPUS\_FORM\_COMPARISONS.json}.}\par
\end{minipage}\par
\par\medskip\Needspace{7\baselineskip}
\noindent\begin{minipage}{\linewidth}
\hypertarget{TERM-1453}{}
\noindent\textbf{racionalno predstavjenje}\quad{\footnotesize\texttt{TERM-1453}}\par
\noindent\emph{Iztočnik / Source:} rationale Darstellung\par
\noindent\emph{Rusky / Russian:} рациональное представление\par
\noindent\emph{Ukrajinsky / Ukrainian:} раціональне подання\par
\smallskip\noindent\emph{Razlog:} Racionalno predstavjenje v arhivu odgovarjaje nynějšemu racionalno prědstavjenje. Znači vyraženje invarianta kako racionalnej funkcije zadannyh invariantov, ne prědstavjenje grupy matricami. Rabota 7, §3, prědstavjaje racionalny absolutny invariant kako kvocient dvoh polinomnyh invariantov i potom vyražaje go prěz koeficienty Galoisovej rezolventy. Koncovka raboty odlišaje takovo prědstavjenje od teorema konečnosti. Razlog zapisa jest smyslova i terminologična kontinuitet s rabotoju 6. Nynějšna bukva ě jest svedočena tekstom; ne trěba prěpisyvati staru formu, da by ona izględala bezproměnna. Racionalna prědstavimost sama po sobě ne jest tvrđenje o polinomnom porodženju.\par
\noindent\emph{Reason:} The archival racionalno predstavjenje corresponds to current racionalno prědstavjenje. It means expressing an invariant as a rational function of specified invariants, not representing a group by matrices. Paper 7, §3, writes a rational absolute invariant as a quotient of two polynomial invariants and then expresses it through coefficients of the Galois resolvent. The paper's conclusion distinguishes such representation from the finiteness theorem. Semantic and terminological continuity with Paper 6 is the recorded rationale. The current ě is attested by the text; the earlier form should not be rewritten to make the wording appear unchanged. Rational expressibility alone is not a claim of polynomial generation.\par
\noindent{\footnotesize Izvorne zapisy / Original records: HIST-g056-010.}\par
\smallskip\noindent\emph{Sravnjenje s tekstom:} Tuto prědstavjenje ne jest matrično prědstavjenje grupy.\par
\noindent\emph{Textual comparison:} This representation is not a matrix representation of a group.\par
\noindent{\footnotesize TEXT-COMP-522; pasaže / passages: 6705, 6706, 6711, 6712, 6713, 6714, 6715, 6716, 6717, 6718, 6719, 6720, 6723; \texttt{CORPUS\_FORM\_COMPARISONS.json}.}\par
\end{minipage}\par
\par\medskip\Needspace{7\baselineskip}
\noindent\begin{minipage}{\linewidth}
\hypertarget{TERM-1454}{}
\noindent\textbf{polja racionalnyh funkcij / polje racionalnyh funkcij}\quad{\footnotesize\texttt{TERM-1454}}\par
\noindent\emph{Iztočnik / Source:} Rationale Funktionenkörper / rationaler Funktionenkörper\par
\noindent\emph{Rusky / Russian:} поля рациональных функций / поле рациональных функций\par
\noindent\emph{Ukrajinsky / Ukrainian:} поля раціональних функцій / поле раціональних функцій\par
\smallskip\noindent\emph{Razlog:} Polje racionalnyh funkcij imenuje matematičny objekt vměsto doslovnogo běžnogo značenja německogo Körper. Zapis upotrěbljaje polje, ne tělo, i daje jedinstveno i množno čislo. Definicija raboty 5 jest važna: elementy polja sut racionalne funkcije od zadannyh neoprěděljenyh, ale polje ne mora byti polnym poljem vsih takyh funkcij niti samo čistym transcendentnym razširjenjem. Invariantne podpolja takože naležet k opisu. Tvrđenje o racionalnoj bazě zato ne smi byti prěvedeno kako tvrđenje o nezavisnoj minimalnoj bazě. Arhiv bira neposrědno nazvanje objekta; matematičny obseg davaje definicija, ne sovrěmenno volno čitanje pridavnogo slova racionalny.\par
\noindent\emph{Reason:} Polje racionalnyh funkcij names the mathematical object instead of translating the everyday meaning of German Körper literally. The record uses polje rather than tělo and supplies singular and plural. Paper 5's definition matters: the field's elements are rational functions in specified indeterminates, but the field need not be the full field of all such functions or itself a purely transcendental extension. Invariant subfields also fit the description. A claim about a rational basis must therefore not become a claim about an independent minimal basis. The archive chooses a direct object name; the definition, not a loose modern reading of rational, determines its mathematical scope.\par
\noindent{\footnotesize Izvorne zapisy / Original records: HIST-g039-001.}\par
\smallskip\noindent\emph{Sravnjenje s tekstom:} Definicija ne zahtěvaje čistej transcendentnosti samogo podpolja.\par
\noindent\emph{Textual comparison:} The definition does not require the subfield itself to be purely transcendental.\par
\noindent{\footnotesize TEXT-COMP-523; pasaže / passages: 4941, 4950, 4952, 4954, 4956, 4958; \texttt{CORPUS\_FORM\_COMPARISONS.json}.}\par
\end{minipage}\par
\par\medskip\Needspace{7\baselineskip}
\noindent\begin{minipage}{\linewidth}
\hypertarget{TERM-1455}{}
\noindent\textbf{racionalno prosto čislo}\quad{\footnotesize\texttt{TERM-1455}}\par
\noindent\emph{Iztočnik / Source:} rationale Primzahl\par
\noindent\emph{Rusky / Russian:} рациональное простое число\par
\noindent\emph{Ukrajinsky / Ukrainian:} раціональне просте число\par
\smallskip\noindent\emph{Razlog:} Racionalno prosto čislo označaje obyčno prosto čislo v racionalnoj čislovoj oblasti, ne prosty ideal ili prosty element proizvoljnogo razširjenja. V rabotě 32, §3, ono učestvuje v kongruenciji $2^r+1\equiv0\pmod p$, a stepenj dvojky, děleči $p-1$, oprěděljaje daljšu konstrukciju. Takovy aritmetičny kontekst opravdyvaje čislo i čuva funkciju slova racionalno. Arhiv izrično bira razlikovanje obyčnyh prostyh čisel od objektov v razširenjah; ruski i ukrajinski zapisy imajut isto složenje. Uslovje na $p$ i kvantifikator o beskonečno mnogih stepenjah ostavajut v teoremu; ne trěba ih zaměnjati občeju rečeniceju pro proste idealy.\par
\noindent\emph{Reason:} Racionalno prosto čislo denotes an ordinary prime in the rational-number setting, not a prime ideal or a prime element of an arbitrary extension. In Paper 32, §3, it occurs in the congruence $2^r+1\equiv0\pmod p$, and the power of two dividing $p-1$ determines the subsequent construction. This arithmetic context supports čislo and preserves the function of racionalno. The archive expressly distinguishes ordinary rational primes from objects in extensions; the Russian and Ukrainian records have the same construction. The condition on $p$ and the quantification over infinitely many exponents belong to the theorem and cannot be replaced by a general statement about prime ideals.\par
\noindent{\footnotesize Izvorne zapisy / Original records: HIST-g175-002.}\par
\smallskip\noindent\emph{Sravnjenje s tekstom:} Kontekst različaje čislo od ideala i čuva uslovje na prosty modul kongruencije.\par
\noindent\emph{Textual comparison:} The context distinguishes a number from an ideal and retains the condition on the prime modulus.\par
\noindent{\footnotesize TEXT-COMP-524; pasaže / passages: 19604, 19608, 19610, 19612, 19614; \texttt{CORPUS\_FORM\_COMPARISONS.json}.}\par
\end{minipage}\par
\par\medskip\Needspace{7\baselineskip}
\noindent\begin{minipage}{\linewidth}
\hypertarget{TERM-1456}{}
\noindent\textbf{racionalne relacije}\quad{\footnotesize\texttt{TERM-1456}}\par
\noindent\emph{Iztočnik / Source:} rationale Relationen\par
\noindent\emph{Rusky / Russian:} рациональные соотношения\par
\noindent\emph{Ukrajinsky / Ukrainian:} раціональні співвідношення\par
\smallskip\noindent\emph{Razlog:} Racionalne relacije sut zavisimosti ohranjane izomorfnym odobraženjem polj. Rabota 10, č. 1, opisuje racionalne operacije kako konečne sostavjenja dodavanja, množenja i obratnyh operacij, potom piše relaciju $F(x,y,\ldots,t)=0$. Slovo relacija tu ne jest sinonim proizvoljnogo binarnogo odnošenja iz logiki. Arhiv bira kratky internacionalizm i tvrdi razpoznavajemost, ale ne davaje izslědovanje čitateljev; ruske sootnošenija i ukrajinske spivvidnošennja sut sravniteljne formy. Pri prěnosu v OLP trěba zato razlikovati smysly matematičnoj relacije, zavisimosti i formalnogo relacijnogo simvola, ne avtomatično napraviti jeden prěvod vseh pojavjenij.\par
\noindent\emph{Reason:} Racionalne relacije are dependencies preserved by a field isomorphism. Paper 10, part 1, describes rational operations as finite compositions of addition, multiplication and their inverses, then displays a relation $F(x,y,\ldots,t)=0$. Relacija here is not a synonym for an arbitrary binary relation in logic. The archive selects a short international term and asserts recognizability, but supplies no reader study; Russian sootnosheniya and Ukrainian spivvidnoshennia are the recorded comparanda. Transfer to OLP must therefore distinguish mathematical dependence, relations and formal relation symbols rather than automatically imposing one translation on every occurrence.\par
\noindent{\footnotesize Izvorne zapisy / Original records: HIST-g069-010.}\par
\smallskip\noindent\emph{Sravnjenje s tekstom:} Racionalna relacija tu proizhodi iz operacij polja; ne jest proizvoljny relacijny simvol.\par
\noindent\emph{Textual comparison:} A rational relation here arises from field operations; it is not an arbitrary relation symbol.\par
\noindent{\footnotesize TEXT-COMP-525; pasaže / passages: 8434, 8436, 8439, 8440, 8452, 8454, 8456; \texttt{CORPUS\_FORM\_COMPARISONS.json}.}\par
\end{minipage}\par
\par\medskip\Needspace{7\baselineskip}
\noindent\begin{minipage}{\linewidth}
\hypertarget{TERM-1457}{}
\noindent\textbf{racionalna kombinacija}\quad{\footnotesize\texttt{TERM-1457}}\par
\noindent\emph{Iztočnik / Source:} rationale Verbindung\par
\noindent\emph{Rusky / Russian:} рациональная комбинация\par
\noindent\emph{Ukrajinsky / Ukrainian:} раціональна комбінація\par
\smallskip\noindent\emph{Razlog:} Racionalna kombinacija znači racionalny izraz v zadannyh funkcijah, ne obyčnu svez i ne samo linearnu kombinaciju. To razlikovanje opravdyvaje zaměnu doslovnogo značenja Verbindung slovom kombinacija. Pet zapisov čuvaju isto nazvanje od raboty 5 do različnyh paragrafov raboty 6. Jediny ranějšny opis spomina podobnost linearnomu izrazu; to ne može služiti definicijeju, zato že tuty operacije dopušćajut produkty i kvocienty. Definicija racionalnoj bazy vyžaduje upravo racionalne, ne samo linearne izrazy. Kvalitativna prědnost dana jest prirodnomu algebraičnomu čitanju i kontinuitetu; arhiv ne sodrživaje čislovogo sravnjenja.\par
\noindent\emph{Reason:} Racionalna kombinacija means a rational expression in specified functions, neither an ordinary connection nor merely a linear combination. That distinction motivates replacing the everyday sense of Verbindung with kombinacija. Five records maintain the term from Paper 5 through several sections of Paper 6. One early description calls the expression linear-like; that cannot serve as a definition because the operations allow products and quotients. The rational-basis definition requires rational, not merely linear, expressions. Natural algebraic reading and continuity receive qualitative priority; no numerical comparison is preserved.\par
\noindent{\footnotesize Izvorne zapisy / Original records: HIST-g039-010, HIST-g043-008, HIST-g044-004, HIST-g045-009, HIST-g047-008.}\par
\smallskip\noindent\emph{Sravnjenje s tekstom:} Dokumentovane upotrěby dopušćajut obče racionalne izrazy.\par
\noindent\emph{Textual comparison:} The documented uses admit general rational expressions.\par
\noindent{\footnotesize TEXT-COMP-526; pasaže / passages: 4956, 5130, 5356, 5410, 5463, 5465, 5466, 5468, 5469, 5470, 5472, 5473, 5605, 15999; \texttt{CORPUS\_FORM\_COMPARISONS.json}.}\par
\end{minipage}\par
\par\medskip\Needspace{7\baselineskip}
\noindent\begin{minipage}{\linewidth}
\hypertarget{TERM-1458}{}
\noindent\textbf{racionalno / algebraično / transcendentno čislo}\quad{\footnotesize\texttt{TERM-1458}}\par
\noindent\emph{Iztočnik / Source:} rationale Zahl / algebraische Zahl / transzendente Zahl\par
\noindent\emph{Rusky / Russian:} рациональное / алгебраическое / трансцендентное число\par
\noindent\emph{Ukrajinsky / Ukrainian:} раціональне / алгебраїчне / трансцендентне число\par
\smallskip\noindent\emph{Razlog:} Racionalno, algebraično i transcendentno čislo čuvaju internacionalne korenje zajedno s rodnym čislo. Arhiv opisuje ih kako standardnu terminologiju; to jest zapis o razlogu izbora, ne samostojno potvrdženje vseh ciljnyh form od jezyčnogo avtoriteta. V rabotě 10 racionalne čisla pri odobraženju ostavajut fiksovane, algebraične mogut prějdti v konjugate, a za transcendentne trěba izbrati bazu. Tri nazvanja ne označajut tri razděljene klasy: racionalne čisla sut algebraične. Spoločny končny član čislo možno deliti med pridavnymi slovami v slovničnom naslovi, ale razlike v tvrdženjah i uključanje jednoj klasy v drugu morajut ostati izrične.\par
\noindent\emph{Reason:} Racionalno, algebraično and transcendentno čislo retain international stems with native čislo. The archive calls them standard terminology; this records the reason for selection, not independent approval of every target form by a language authority. In Paper 10, rational numbers are fixed by the map, algebraic numbers may move to conjugates, and transcendental numbers require a basis construction. These are not three disjoint classes: rational numbers are algebraic. A dictionary heading can share the final noun čislo among the adjectives, but the different assertions and the inclusion of one class in another must remain explicit.\par
\noindent{\footnotesize Izvorne zapisy / Original records: HIST-g070-002.}\par
\smallskip\noindent\emph{Sravnjenje s tekstom:} Racionalne čisla sut algebraične; tri nazvanja ne dajut tri razděljene klasy.\par
\noindent\emph{Textual comparison:} Rational numbers are algebraic; the three labels do not give three disjoint classes.\par
\noindent{\footnotesize TEXT-COMP-527; pasaže / passages: 8487, 8489, 8491; \texttt{CORPUS\_FORM\_COMPARISONS.json}.}\par
\end{minipage}\par
\par\medskip\Needspace{7\baselineskip}
\noindent\begin{minipage}{\linewidth}
\hypertarget{TERM-1459}{}
\noindent\textbf{racionalny / cěly Galoisov modul}\quad{\footnotesize\texttt{TERM-1459}}\par
\noindent\emph{Iztočnik / Source:} rationaler / ganzer Galoismodul\par
\noindent\emph{Rusky / Russian:} рациональный / целый модуль Галуа\par
\noindent\emph{Ukrajinsky / Ukrainian:} раціональний / цілий модуль Галуа\par
\smallskip\noindent\emph{Razlog:} Racionalny / cěly Galoisov modul razlikuje oblast operatorov, ne samo stilističny opis modula. V rabotě 37 racionalny slučaj upotrěbljaje grupovu algebru nad poljem $k$, a cěly slučaj grupovy prsten nad oblastju cělyh elementov. Teoremy o poljnoj i cěloj normalnoj bazě zato imajut različne hipotezy i različne zaključky. Arhiv izrično zahtěvaje čuvati protivnost racionalny/cěly vo vseh prěvodah. To ne opravdyvaje tvrđenje, že vsaky cěly Galoisov modul jest svobodny ili imaje normalnu bazu. Pytanje o uslovju na stepenj v pozdnějšem tvrdženju togo teksta ostavaje oddělno zaznamenano; lexikalna soglasnost jego ne rěšaje.\par
\noindent\emph{Reason:} Racionalny / cěly Galoisov modul distinguishes the operator domain, not merely a descriptive style. In Paper 37 the rational case uses the group algebra over the field $k$, while the integral case uses the group ring over the domain of integral elements. The field and integral normal-basis theorems consequently have different hypotheses and conclusions. The archive expressly requires the rational/integral contrast in every translation. It does not justify saying that every integral Galois module is free or has a normal basis. The separately recorded question about the degree condition in a later assertion of this text remains unresolved; lexical consistency does not settle it.\par
\noindent{\footnotesize Izvorne zapisy / Original records: HIST-g207-013.}\par
\smallskip\noindent\emph{Sravnjenje s tekstom:} Poljny i cěly slučaj ne imajut jednakogo tvrđenja o normalnoj bazě.\par
\noindent\emph{Textual comparison:} The field and integral cases do not carry the same normal-basis assertion.\par
\noindent{\footnotesize TEXT-COMP-528; pasaže / passages: 22821, 22823, 22824, 22828, 22831, 22833, 22841, 22844, 22846, 22849, 22851, 22853, 22856; \texttt{CORPUS\_FORM\_COMPARISONS.json}.}\par
\end{minipage}\par
\par\medskip\Needspace{7\baselineskip}
\noindent\begin{minipage}{\linewidth}
\hypertarget{TERM-1460}{}
\noindent\textbf{racionalny funkcional}\quad{\footnotesize\texttt{TERM-1460}}\par
\noindent\emph{Iztočnik / Source:} rationales Funktional\par
\noindent\emph{Rusky / Russian:} рациональный функционал\par
\noindent\emph{Ukrajinsky / Ukrainian:} раціональний функціонал\par
\smallskip\noindent\emph{Razlog:} Racionalny funkcional jest historičny termin definovany prěz normalizovano prědstavjenje racionalnoj funkcije, ne obyčny sovrěmenny funkcional na prostoru funkcij. V rabotě 9, §2, zapis ima formu $A=aE_1/E_2$, kde polinomy sut primitivne i vzajemno proste s cělymi racionalnymi koeficientami, a $a$ jest pozitivno racionalno čislo. Cělost funkcionala zavisi od cělosti $a$, ne od odsutnosti imenitelja. Arhiv čuva nazvanje funkcional, zato že Noether izrično upuča na Webera; ona ne govori samo o proizvoljnoj funkciji. Ta upučenost jest svedočena pasažeju Noetherovej raboty. Weberova navedena kniga ne byla tu nezavisno pročitana, zato novo neposrědno potvrdženje od Webera ne tvrdi se.\par
\noindent\emph{Reason:} Racionalny funkcional is a historical term defined through a normalized presentation of a rational function, not an ordinary modern functional on a function space. In Paper 9, §2, the expression has the form $A=aE_1/E_2$, with relatively prime primitive polynomials having integral rational coefficients and positive rational $a$. Integrality of the functional depends on integrality of $a$, not on the absence of a denominator. The archive keeps funkcional because Noether explicitly refers to Weber rather than speaking only of an arbitrary function. That reference is attested by Noether's passage. Weber's cited book has not been independently read here, so no fresh direct verification from Weber is claimed.\par
\noindent{\footnotesize Izvorne zapisy / Original records: HIST-g060-001.}\par
\smallskip\noindent\emph{Sravnjenje s tekstom:} Cělost zavisi od racionalnogo množitelja v normalizovanoj formě, ne od odsutnosti imenitelja.\par
\noindent\emph{Textual comparison:} Integrality depends on the rational multiplier in the normalized form, not on the absence of a denominator.\par
\noindent{\footnotesize TEXT-COMP-529; pasaže / passages: 7221, 7223, 7225, 7226, 7228, 7230, 7232; \texttt{CORPUS\_FORM\_COMPARISONS.json}.}\par
\end{minipage}\par
\par\medskip\Needspace{7\baselineskip}
\noindent\begin{minipage}{\linewidth}
\hypertarget{TERM-1461}{}
\noindent\textbf{oblast racionalnosti}\quad{\footnotesize\texttt{TERM-1461}}\par
\noindent\emph{Iztočnik / Source:} Rationalitätsbereich\par
\noindent\emph{Rusky / Russian:} область рациональности\par
\noindent\emph{Ukrajinsky / Ukrainian:} область раціональності\par
\smallskip\noindent\emph{Razlog:} Oblast racionalnosti čuva historično nazvanje za koeficientnu ili osnovnu oblast, vměsto avtomatičnoj zaměny slovom polje. Tri zapisy povezuju raboty 8, 11 i 17. V rabotě 8 pripadnost koeficientov zadannoj oblasti ograničaje linearne kombinacije; rabota 11 izrično govori i o poljah, takže slovo oblast ne odrěka poljnu strukturu. V rabotě 17 osobno uslovje jest beskonečno mnogo konstant, ktore komutujut s oprěděljenymi operatorami. Arhiv davaje prědnost vernosti avtorovoj terminologiji i ne naměravaje vsude modernizovati Bereich. Tutu leksičnu rezervu ne trěba prěvratiti v zabranu objasniti stvarnu strukturu, kogda ju lokalna definicija ustanavja.\par
\noindent\emph{Reason:} Oblast racionalnosti retains the historical name for a coefficient or base domain instead of automatically substituting polje. Three records connect Papers 8, 11 and 17. In Paper 8, membership in the specified coefficient domain constrains linear combinations; Paper 11 also speaks explicitly of fields, so oblast does not deny field structure. In Paper 17 the particular hypothesis is the existence of infinitely many constants commuting with specified operators. The archive prioritizes fidelity to the author's terminology and declines to modernize Bereich everywhere. This lexical restraint does not prohibit explaining the actual structure when the local definition establishes it.\par
\noindent{\footnotesize Izvorne zapisy / Original records: HIST-g057-007, HIST-g073-003, HIST-g079-056.}\par
\smallskip\noindent\emph{Sravnjenje s tekstom:} Tri priměny razlikajut historično nazvanje oblasti od jeje dodatnyh strukturnyh uslovij.\par
\noindent\emph{Textual comparison:} Three uses distinguish the historical domain label from its additional structural conditions.\par
\noindent{\footnotesize TEXT-COMP-530; pasaže / passages: 6768, 6806, 8724, 8726, 8728, 8742, 8763, 11735, 11737; \texttt{CORPUS\_FORM\_COMPARISONS.json}.}\par
\end{minipage}\par
\par\medskip\Needspace{7\baselineskip}
\noindent\begin{minipage}{\linewidth}
\hypertarget{TERM-1462}{}
\noindent\textbf{racionalno-čislove koeficienty / funkcije}\quad{\footnotesize\texttt{TERM-1462}}\par
\noindent\emph{Iztočnik / Source:} rationalzahlige Koeffizienten / Funktionen\par
\noindent\emph{Rusky / Russian:} рационально-числовые коэффициенты / функции\par
\noindent\emph{Ukrajinsky / Ukrainian:} раціонально-числові коефіцієнти / функції\par
\smallskip\noindent\emph{Razlog:} Racionalno-čislove koeficienty izrično ograničajut koeficienty na racionalne čisla. To jest drugo svojstvo než racionalnost funkcije svojih argumentov. V rabotě 11 invariantno polje sestojit se iz racionalnyh funkcij s racionalno-čislovymi koeficientami; pri razširjenju oblasti koeficientov prěhodit se k drugomu nazvanomu polju. Arhiv bira složeno pridavno slovo, da razlika ne izčezne pri skračenju. Forma racionalno-čislove funkcije zato potrěbuje kontekst, ktory kaže, že čislova kvalifikacija govori o koeficientah, ne o tom, že vse cěnnosti funkcije sut racionalne. Čislove oceny alternativ ne sut zachovane.\par
\noindent\emph{Reason:} Racionalno-čislove koeficienty explicitly restricts the coefficients to rational numbers. This is different from a function being rational in its arguments. In Paper 11 the invariant field consists of rational functions with rational-number coefficients; extending the coefficient domain produces another specified field. The archive chooses the compound adjective to keep this distinction visible in abbreviated prose. Racionalno-čislove funkcije accordingly needs a context making clear that the numerical qualification concerns coefficients, not a claim that every function value is rational. No numerical assessment of alternatives is preserved.\par
\noindent{\footnotesize Izvorne zapisy / Original records: HIST-g073-011.}\par
\smallskip\noindent\emph{Sravnjenje s tekstom:} Kvalifikacija govori o koeficientah, ne o vseh cěnnostjah funkcije.\par
\noindent\emph{Textual comparison:} The qualification concerns the coefficients, not all values of the function.\par
\noindent{\footnotesize TEXT-COMP-531; pasaže / passages: 8763; \texttt{CORPUS\_FORM\_COMPARISONS.json}.}\par
\end{minipage}\par
\par\medskip\Needspace{7\baselineskip}
\noindent\begin{minipage}{\linewidth}
\hypertarget{TERM-1463}{}
\noindent\textbf{računno sredstvo}\quad{\footnotesize\texttt{TERM-1463}}\par
\noindent\emph{Iztočnik / Source:} rechnerisches Mittel\par
\noindent\emph{Rusky / Russian:} вычислительное средство\par
\noindent\emph{Ukrajinsky / Ukrainian:} обчислювальний засіб\par
\smallskip\noindent\emph{Razlog:} Računno sredstvo imenuje konkretny sposob izvedenja redukcij, ne računalnik ili program. Uvod raboty 2 povezuje jego s rekursivnym tvorjenjem specialnyh razvitij v redy: teoremy davajut osnovu, a razvitja omogućajut skutečny račun. Arhiv izrično odbira računno vměsto komputacijny, ktoro ocěnjuje kako prěviše anglijskopodobno za toj tekst, dokolě drugy kontekst ne pokaže potrěbu. To jest stilistična ocěna togo zapisa, ne obča zabrana internacionalizmov v ISV. Ruski i ukrajinski izrazy takože označajut sredstvo vyčislenja. Pri prěnosu v logiku trěba nanovo razlikovati konkretny račun, algoritm i teoriju vyčislimosti.\par
\noindent\emph{Reason:} Računno sredstvo names a concrete means of carrying out reductions, not a computer or program. Paper 2's introduction links it to recursive construction of particular series expansions: the theorems supply the basis, while the expansions enable the actual calculation. The archive explicitly prefers računno to komputacijny, judging the latter too English-like for this passage unless another context requires it. That is the record's stylistic judgment, not a general prohibition on internationalisms in ISV. The Russian and Ukrainian expressions likewise denote a means of calculation. Transfer to logic requires renewed distinctions among concrete calculation, an algorithm and computability theory.\par
\noindent{\footnotesize Izvorne zapisy / Original records: HIST-g002-009.}\par
\smallskip\noindent\emph{Sravnjenje s tekstom:} Rekursivne razvitja dopuščajut izvesti redukcije, ktoryh osnovu davajut teoremy.\par
\noindent\emph{Textual comparison:} Recursive expansions carry out reductions whose basis the theorems supply.\par
\noindent{\footnotesize TEXT-COMP-532; pasaže / passages: 180, 182; \texttt{CORPUS\_FORM\_COMPARISONS.json}.}\par
\end{minipage}\par
\par\medskip\Needspace{7\baselineskip}
\noindent\begin{minipage}{\linewidth}
\hypertarget{TERM-1464}{}
\noindent\textbf{pręmougolna shema}\quad{\footnotesize\texttt{TERM-1464}}\par
\noindent\emph{Iztočnik / Source:} rechteckiges Schema\par
\noindent\emph{Rusky / Russian:} прямоугольная схема\par
\noindent\emph{Ukrajinsky / Ukrainian:} прямокутна схема\par
\smallskip\noindent\emph{Razlog:} Pręmougolna shema jest doslovny opis razměščenja izrazov po redah i stolbcah. Zachovany razlog upuča na vid tablice, ale sam po sobě ne jest dokaz, že vsaka strana izdankov byla vidovo prověrjena. Nynějšny tekst raboty 2, §1, imaje oblik premougolna i eksplicitnu tablicu s po sobě idučimi kontrakcijami. Shema tuto organizuje red form; ne trěba bez dopolniteljnogo obosnovanja nazvati ju matricom s matričnym množenjem. Bukvova razlika medžu arhivnym ę i nynějšim e ostavaje zabilěžena. Sravnjenje objasnjaje izbrany smysl, ale ne nadoměnjaje pravopisnogo obosnovanja ni vidovoj prověrky konečnogo čitatelja.\par
\noindent\emph{Reason:} Pręmougolna shema literally describes expressions arranged in rows and columns. The preserved rationale appeals to the array's appearance, but is not itself evidence that every edition page was visually checked. The current text of Paper 2, §1, uses premougolna and an explicit array of successive contractions. Here shema organizes a series of forms; it should not be called a matrix with matrix multiplication without further justification. The archival ę and current e remain recorded as different spellings. This comparison explains the chosen sense but does not replace orthographic justification or visual inspection of the final reader.\par
\noindent{\footnotesize Izvorne zapisy / Original records: HIST-g003-012.}\par
\smallskip\noindent\emph{Sravnjenje s tekstom:} Premougolna opisuje oblik shemy; formula ne uvodi samostalno matričnu operaciju.\par
\noindent\emph{Textual comparison:} Premougolna describes the array's shape; the formula does not by itself introduce a matrix operation.\par
\noindent{\footnotesize TEXT-COMP-533; pasaže / passages: 294, 296, 298, 299, 300, 301, 304, 305, 306, 307, 308, 309, 310, 311, 312, 313, 314, 315, 316, 317, 318, 319, 320, 321, 322, 323, 324, 325, 326, 327, 328, 329, 330, 331, 332, 333, 334, 335, 336, 337, 338; \texttt{CORPUS\_FORM\_COMPARISONS.json}.}\par
\end{minipage}\par
\par\medskip\Needspace{7\baselineskip}
\noindent\begin{minipage}{\linewidth}
\hypertarget{TERM-1465}{}
\noindent\textbf{prave odnosno lěve klasy}\quad{\footnotesize\texttt{TERM-1465}}\par
\noindent\emph{Iztočnik / Source:} Rechts- bzw. Linksklassen\par
\noindent\emph{Rusky / Russian:} правые, соответственно левые, классы\par
\noindent\emph{Ukrajinsky / Ukrainian:} праві, відповідно ліві, класи\par
\smallskip\noindent\emph{Razlog:} Prave odnosno lěve klasy čuva orientaciju jednostrannoj ekvivalencije matric. V rabotě 19 bazisne matrice za odgovarjajuči ideal imajut formu $UA$ ili $AV$, kde množitelje sut unimodularne; dvustranna klasa dopuščaje $UAV$. Arhiv zato odbira jasno nazvane obě strany, ne obče jednostranne klasy. Odnosno povezuje nazvanje s odgovarjajučim slučajem, ne dopuščaje libovoljno zaměnjati strany. Tuty klasy ne sut prosto klasy ostatkov po idealu; opreděljaju ih dopustime proměny bazisa. Za novo priměnjenje trěba čuvati konkretne operacije i historiju strannyh nazvanij, ne zaključati jih samo iz obyčnogo značenja pravy.\par
\noindent\emph{Reason:} Prave odnosno lěve klasy preserves the orientation of one-sided matrix equivalence. In Paper 19, basis matrices for the corresponding ideal have the form $UA$ or $AV$, with unimodular multipliers; a two-sided class allows $UAV$. The archive therefore names both sides explicitly rather than saying only one-sided classes. Odnosno connects each name with its corresponding case, not an arbitrary interchange of sides. These are not simply residue classes modulo an ideal: the admissible changes of basis define them. A new application must retain the actual operations and historical side-naming convention rather than infer them from the ordinary meaning of right.\par
\noindent{\footnotesize Izvorne zapisy / Original records: HIST-g081-099.}\par
\smallskip\noindent\emph{Sravnjenje s tekstom:} Operacije UA, AV i UAV ustanavjajut, ktory tip ekvivalencije imenovany jest.\par
\noindent\emph{Textual comparison:} The operations UA, AV and UAV establish which equivalence is being named.\par
\noindent{\footnotesize TEXT-COMP-534; pasaže / passages: 13695, 13697, 13703, 13705, 13707, 13709; \texttt{CORPUS\_FORM\_COMPARISONS.json}.}\par
\end{minipage}\par
\par\medskip\Needspace{7\baselineskip}
\noindent\begin{minipage}{\linewidth}
\hypertarget{TERM-1466}{}
\noindent\textbf{pravo predstavjenje / lěvo predstavjenje}\quad{\footnotesize\texttt{TERM-1466}}\par
\noindent\emph{Iztočnik / Source:} Rechtsdarstellung / Linksdarstellung\par
\noindent\emph{Rusky / Russian:} правое представление / левое представление\par
\noindent\emph{Ukrajinsky / Ukrainian:} праве подання / ліве подання\par
\smallskip\noindent\emph{Razlog:} Pravo predstavjenje / lěvo predstavjenje zadržuje zajedny termin Darstellung, ale v rabotě 34, §9, znači razloženje prstena na direktnu sumu pravyh ili lěvyh idealov. Pravomu zapisu s $e_iR$ odgovarjaje lěvy zapis s $Re_i$; ortogonalne idempotenty daju svez med nimi. To ne jest samo nazvanje matričnogo prědstavjenja. Arhiv izrično čuva toj smysl pomoču naslova i formul, ne izmišljaje novego korenja. Nynějšny tekst imaje prědstavjenje s ě; razlika od starogo zapisa ostavaje vidima. Slovnična lokalizacija mora povezati obě slovesne formy s operacijeju razloženja i ne prěnesti samostalno smysl teorije prědstavjenij.\par
\noindent\emph{Reason:} Pravo predstavjenje / lěvo predstavjenje retains the shared word Darstellung, but in Paper 34, §9, it means decomposition of the ring into a direct sum of right or left ideals. A right-hand expression with $e_iR$ corresponds to a left-hand expression with $Re_i$ through orthogonal idempotents. It is not merely the name of a matrix representation. The archive explicitly relies on the title and formulas to preserve that sense rather than inventing another stem. The current text has prědstavjenje with ě; its difference from the earlier spelling remains visible. Lexicographic description must connect both forms to decomposition rather than import the separate representation-theoretic sense.\par
\noindent{\footnotesize Izvorne zapisy / Original records: HIST-g187-005.}\par
\smallskip\noindent\emph{Sravnjenje s tekstom:} Prědstavjenje govori tuto o direktnoj sumě idealov; idempotenty čuvajut svez med oběma stranami.\par
\noindent\emph{Textual comparison:} Representation here means a direct sum of ideals; the idempotents relate the two sides.\par
\noindent{\footnotesize TEXT-COMP-535; pasaže / passages: 20509, 20511, 20513, 20514, 20515, 20516, 20518, 20519, 20520, 20521, 20522, 20523, 20524, 20525, 20528, 20530, 20532, 20534, 20536, 20538, 20540, 20542, 20544, 20545; \texttt{CORPUS\_FORM\_COMPARISONS.json}.}\par
\end{minipage}\par
\par\medskip\Needspace{7\baselineskip}
\noindent\begin{minipage}{\linewidth}
\hypertarget{TERM-1467}{}
\noindent\textbf{prava / lěva obratna matrica}\quad{\footnotesize\texttt{TERM-1467}}\par
\noindent\emph{Iztočnik / Source:} Rechtsinverse / Linksinverse\par
\noindent\emph{Rusky / Russian:} правая / левая обратная матрица\par
\noindent\emph{Ukrajinsky / Ukrainian:} права / ліва обернена матриця\par
\smallskip\noindent\emph{Razlog:} Prava / lěva obratna matrica čuva, na ktorej straně množenje davaje jediničnu matricu. §7 raboty 34 rabotaje s kvadratnymi matricami nad možno nekomutativnym tělom. Definicija regularnosti najprvo zahtěvaje obě obratne matrice; slědujuči dokaz pokazuje, že jedna strana dostaje i že obě obratne matrice sut jednake. To jest dokazane tvrđenje pod danymi prědpoloženjami, ne pričina vymazati pravy/lěvy iz definicije. Arhiv davaje upravo toj dokazny razlikě prědnost. Jednaka terminologija ne opravdyvaje prěnos zaključka na proizvoljne nekvadratne matrice ili na proizvoljny prsten.\par
\noindent\emph{Reason:} Prava / lěva obratna matrica preserves which side of multiplication yields the identity matrix. Paper 34, §7, treats square matrices over a possibly noncommutative division ring. The definition of regularity initially requires both inverses; the following proof shows that one side suffices and that the inverses agree. This is a proved assertion under the stated assumptions, not a reason to erase right/left from the definition. The archive gives precisely that proof-bearing distinction priority. Shared terminology does not justify transferring the conclusion to arbitrary rectangular matrices or arbitrary rings.\par
\noindent{\footnotesize Izvorne zapisy / Original records: HIST-g185-003.}\par
\smallskip\noindent\emph{Sravnjenje s tekstom:} Dokaz za kvadratne matrice nad tělom ne jest obča definicija obratnosti.\par
\noindent\emph{Textual comparison:} A proof for square matrices over a division ring is not a general definition of invertibility.\par
\noindent{\footnotesize TEXT-COMP-536; pasaže / passages: 20329, 20331, 20333, 20335, 20336, 20337, 20338, 20339, 20340, 20341, 20342, 20343, 20344, 20345, 20346, 20347, 20348, 20349, 20350; \texttt{CORPUS\_FORM\_COMPARISONS.json}.}\par
\end{minipage}\par
\par\medskip\Needspace{7\baselineskip}
\noindent\begin{minipage}{\linewidth}
\hypertarget{TERM-1468}{}
\noindent\textbf{pravostranny / lěvostranny / dvustranny modul}\quad{\footnotesize\texttt{TERM-1468}}\par
\noindent\emph{Iztočnik / Source:} rechtsseitig / linksseitig / zweiseitig\par
\noindent\emph{Rusky / Russian:} правосторонний / левосторонний / двусторонний модуль\par
\noindent\emph{Ukrajinsky / Ukrainian:} правобічний / лівобічний / двобічний модуль\par
\smallskip\noindent\emph{Razlog:} Pravostranny, lěvostranny i dvustranny konkretno označajut različne zamknjenosti modula. Rabota 17, §3, nazyvaje FM-zamknjenost pravostrannoju, a MF-zamknjenost lěvostrannoju. Točny poredok faktorov jest zato važnějši než domysl iz sovrěmennogo nazvanja. Arhiv izrično zachovaje tuty směr i protivstavi go neoprěděljenomu jednostranny; dvustranny znači oba uslovja. Otvorjeno pytanje jest, je li takovo nazvanje v cělosti iztočno-věrno historičnoj konvenciji; bez čitanja iztočnika ne trěba go ni opravdati, ni obratiti. Za prěnos v novy tekst potrěbny sut formula operacije i jasno ukazana konvencija.\par
\noindent\emph{Reason:} Pravostranny, lěvostranny and dvustranny explicitly name different module closure conditions. Paper 17, §3, calls FM-closure right-sided and MF-closure left-sided. The exact order of factors therefore matters more than an inference from a present-day label. The archive expressly preserves this direction rather than using only the vague one-sided; two-sided means both conditions. Whether the naming faithfully reproduces the historical source convention remains an open source question; it should neither be certified nor reversed without reading that source. Transfer to a new text requires the operation formula and an explicit convention.\par
\noindent{\footnotesize Izvorne zapisy / Original records: HIST-g079-014.}\par
\smallskip\noindent\emph{Sravnjenje s tekstom:} Formule čuvajut točnu orientaciju nynějšego teksta; iztočna historična konvencija potrěbuje odvojenogo prověrjenja.\par
\noindent\emph{Textual comparison:} The formulas preserve the current text's precise orientation; the historical source convention requires separate verification.\par
\noindent{\footnotesize TEXT-COMP-537; pasaže / passages: 11090, 11092, 11094, 11096, 11098, 11102, 13360, 13362, 13364, 13366, 13370, 13372, 13374, 13376, 13378, 13380, 13382, 13384, 13386, 13387, 13388, 13390; \texttt{CORPUS\_FORM\_COMPARISONS.json}.}\par
\end{minipage}\par
\par\medskip\Needspace{7\baselineskip}
\noindent\begin{minipage}{\linewidth}
\hypertarget{TERM-1469}{}
\noindent\textbf{pravostranne / lěvostranne / dvostranne moduly i idealy}\quad{\footnotesize\texttt{TERM-1469}}\par
\noindent\emph{Iztočnik / Source:} rechtsseitige / linksseitige / zweiseitige Moduln und Ideale\par
\noindent\emph{Rusky / Russian:} правосторонние / левосторонние / двусторонние модули и идеалы\par
\noindent\emph{Ukrajinsky / Ukrainian:} правобічні / лівобічні / двобічні модулі та ідеали\par
\smallskip\noindent\emph{Razlog:} Rodina stran- davaje prozirnu tvorbu nazvanij za obě orientacije i za obě zajedno. Zapis imaje dvostranne, dokolě drugy zapis i česti nynějšego teksta imajut dvustranne; arhiv ne davaje sravniteljne oceny dvoh samoglasnyh variantov. Smysl ustanavjajut operacije: v definiciji raboty 19, §9, množenje $c\alpha$ nazvano jest pravostrannym, a $\alpha c$ lěvostrannym. Toj konkretni dokaz naleži definiciji, hoti arhivna oznaka jedinice upuča na §12. Nazvanje lokacije samo po sobě zato ne jest dovoljne svedočstvo. Latinski tekst ostavaje osnovoju projekcije; kirilična forma ne davaje drugogo nezavisnogo dokaza za strannu konvenciju.\par
\noindent\emph{Reason:} The stran- family makes the names for the two orientations and their combination morphologically explicit. This record has dvostranne, while another record and parts of the current text have dvustranne; no comparative assessment of the vowel variants is preserved. Operations establish the sense: the definition in Paper 19, §9, calls multiplication $c\alpha$ right-sided and $\alpha c$ left-sided. That concrete evidence belongs to the definition, although the archival unit label points to §12. A location label is therefore not sufficient evidence by itself. The Latin text remains the basis for projection; the Cyrillic form supplies no second independent witness for the side convention.\par
\noindent{\footnotesize Izvorne zapisy / Original records: HIST-g081-071.}\par
\smallskip\noindent\emph{Sravnjenje s tekstom:} Formule čuvajut točnu orientaciju nynějšego teksta; iztočna historična konvencija potrěbuje odvojenogo prověrjenja.\par
\noindent\emph{Textual comparison:} The formulas preserve the current text's precise orientation; the historical source convention requires separate verification.\par
\noindent{\footnotesize TEXT-COMP-537; pasaže / passages: 11090, 11092, 11094, 11096, 11098, 11102, 13360, 13362, 13364, 13366, 13370, 13372, 13374, 13376, 13378, 13380, 13382, 13384, 13386, 13387, 13388, 13390; \texttt{CORPUS\_FORM\_COMPARISONS.json}.}\par
\end{minipage}\par
\par\medskip\Needspace{7\baselineskip}
\noindent\begin{minipage}{\linewidth}
\hypertarget{TERM-1470}{}
\noindent\textbf{redukcija}\quad{\footnotesize\texttt{TERM-1470}}\par
\noindent\emph{Iztočnik / Source:} Reduktion\par
\noindent\emph{Rusky / Russian:} редукция\par
\noindent\emph{Ukrajinsky / Ukrainian:} редукція\par
\smallskip\noindent\emph{Razlog:} Redukcija jest činnost vyraženja formy pomoču dopustimyh drugih form, v lokalno oprěděljenom redukcijnom sistemu. V §8 raboty 2 kratak naslov v jedinstvenom čislu prědhodi jednoj eksplicitnoj relaciji; potom poslědki opreděljajut reducent. Zachovany razlog jest kontinuitet terminologije raboty, ne poseben izum za toj paragraf. Vyše položenje v porjadku form ili pojavjenje invariantnogo faktora naleži definiciji reducibilnosti; redukcija ne mora značiti čislovo umenšenje stepenja v vsakom argumentu. Pri daljšej upotrěbě v logikě forma slova može ostati podobna, ale lokalna relacija redukcije mora byti oprěděljena nanovo.\par
\noindent\emph{Reason:} Redukcija is the act of expressing a form through admissible other forms within the locally defined reduction system. In Paper 2, §8, a short singular heading precedes one explicit relation; the consequences then identify a reducer. The preserved reason is continuity with the paper's terminology, not a separate invention for this section. A higher position in the ordering of forms or an invariant factor belongs to the definition of reducibility; reduction need not mean a numerical decrease of degree in every argument. Later use in logic may retain a similar word form, but must define its own reduction relation.\par
\noindent{\footnotesize Izvorne zapisy / Original records: HIST-g009-004.}\par
\smallskip\noindent\emph{Sravnjenje s tekstom:} Jedinstveny naslov uvodi konkretnu relaciju v raněje oprěděljenom porjadku form.\par
\noindent\emph{Textual comparison:} The singular heading introduces a particular relation within the previously defined ordering of forms.\par
\noindent{\footnotesize TEXT-COMP-538; pasaže / passages: 545, 547, 549, 551, 1193, 1195, 1197; \texttt{CORPUS\_FORM\_COMPARISONS.json}.}\par
\end{minipage}\par
\par\medskip\Needspace{7\baselineskip}
\noindent\begin{minipage}{\linewidth}
\hypertarget{TERM-1472}{}
\noindent\textbf{redukcija abo razklad}\quad{\footnotesize\texttt{TERM-1472}}\par
\noindent\emph{Iztočnik / Source:} Reduktion oder Zerfallen\par
\noindent\emph{Rusky / Russian:} редукция или разложение\par
\noindent\emph{Ukrajinsky / Ukrainian:} редукція або розклад\par
\smallskip\noindent\emph{Razlog:} Redukcija abo razklad zachovaje dva matematične opisy, ne dva proizvoljne sinonima. V §24 raboty 2 služi jednokratna kontrakcija za redukciju ili razklad odgovarjajučih form. Tekst potom zasebno upotrěbljaje i formulaciju redukcija i razklad. Abo i i zato ne smějut byti izravnjene: ony različajut slučaje i zajedne poslědky. Arhiv čuva razklad radi svezi s razkladna forma i razlikuje proceduru redukovanja od razlaganja. To ne jest tvrđenje, že dva procesa nikogda ne mogut nastupati zajedno; sama sosedna formulacija pokazuje, zašto ih trěba razuměti po konkretnoj formě i relaciji.\par
\noindent\emph{Reason:} Redukcija abo razklad preserves two mathematical descriptions, not two arbitrary synonyms. In Paper 2, §24, a single contraction produces reduction or decomposition of the relevant forms. The text then separately uses redukcija i razklad. Or and and must not be levelled: they distinguish alternative cases from joint consequences. The archive retains razklad for continuity with razkladna forma and distinguishes reducing from decomposing. This does not assert that the two processes never occur together; the neighbouring formulation itself shows why they must be interpreted through the particular form and relation.\par
\noindent{\footnotesize Izvorne zapisy / Original records: HIST-g025-006.}\par
\smallskip\noindent\emph{Sravnjenje s tekstom:} Sosedne pasaže upotrěbljajut dva sveza radi različnyh poslědkov kontrakcije.\par
\noindent\emph{Textual comparison:} Neighbouring passages use the two conjunctions for different consequences of contraction.\par
\noindent{\footnotesize TEXT-COMP-540; pasaže / passages: 3174, 3176, 3178, 3180, 3182, 3184, 3186, 3188, 3190, 3192, 3194; \texttt{CORPUS\_FORM\_COMPARISONS.json}.}\par
\end{minipage}\par
\par\medskip\Needspace{7\baselineskip}
\noindent\begin{minipage}{\linewidth}
\hypertarget{TERM-1473}{}
\noindent\textbf{redukcije}\quad{\footnotesize\texttt{TERM-1473}}\par
\noindent\emph{Iztočnik / Source:} Reduktionen\par
\noindent\emph{Rusky / Russian:} редукции\par
\noindent\emph{Ukrajinsky / Ukrainian:} редукції\par
\smallskip\noindent\emph{Razlog:} Množny naslov redukcije služi opakovanym skupinam vyvodov v §§9 i 13–16 raboty 2. Pet zapisov opravdyvaje jeho upotrěbu čislom redukcij, kontinuitetom naslovov i internacionalnym korenjem. Tvrđenje o razumljivosti ne podprěno jest testom čitateljev. Povtorjenje naslova ne čini vsako pojavjenje novym terminom: jednotlivne formule ostavajut različne priměny jednogo opravdanogo nazvanja. Sravnjenje mora čuvati i poredok vyvodov, i to, ktory reducent ili ranějša relacija podprěnaje vsaky z nih. Leksična doslědnost sama ne potvrdžaje točnost svih simbolov v formulačnom bloku.\par
\noindent\emph{Reason:} The plural heading redukcije introduces recurring groups of derivations in Paper 2, §§9 and 13–16. Five records justify it by the multiple reductions, consistent headings and international stem. The claim of comprehensibility is not supported by a reader test. Repeating the heading does not make each occurrence a newly coined term: the individual formulas remain separate applications of one justified label. Comparison must preserve the order of derivation and identify the reducer or earlier relation supporting each application. Lexical consistency alone does not verify every symbol in the formula block.\par
\noindent{\footnotesize Izvorne zapisy / Original records: HIST-g010-003, HIST-g014-005, HIST-g015-005, HIST-g016-005, HIST-g017-005.}\par
\smallskip\noindent\emph{Sravnjenje s tekstom:} Pet lokalnyh skupin redukcij imaje zajedny naslov; jeho doslědnost sama ne prověrjaje njihove formule.\par
\noindent\emph{Textual comparison:} Five local groups of reductions share a heading; its consistency does not verify their formulas.\par
\noindent{\footnotesize TEXT-COMP-541; pasaže / passages: 1229, 1232, 1233, 1234, 1235, 1236, 1237, 1238, 1239, 1240, 1241, 1242, 1243, 1244, 2057, 2059, 2061, 2176, 2178, 2180, 2182, 2184, 2263, 2265, 2267, 2269, 2271, 2349, 2351, 2353, 2355; \texttt{CORPUS\_FORM\_COMPARISONS.json}.}\par
\end{minipage}\par
\par\medskip\Needspace{7\baselineskip}
\noindent\begin{minipage}{\linewidth}
\hypertarget{TERM-1474}{}
\noindent\textbf{redukcije / reducent}\quad{\footnotesize\texttt{TERM-1474}}\par
\noindent\emph{Iztočnik / Source:} Reduktionen / Reduzent\par
\noindent\emph{Rusky / Russian:} редукции / редуцент\par
\noindent\emph{Ukrajinsky / Ukrainian:} редукції / редуцент\par
\smallskip\noindent\emph{Razlog:} Redukcije označajut izvedenja, a reducent imenuje tehničnu rolju v njih. Po definiciji raboty 2 reducent jest reducibilny red form; ne jest čověk, ktory izvodi račun, ni obči nazov računalnogo orudja. §18 izrično upotrěbljaje takovy red kako faktor, iz ktorego nastavajut različne redukcije. Arhiv čuva reducent skupaj s rodinoju reducibilny i navodi odgovarjajuče novotvorene ruske i ukrajinske termine. Prozirnost togo internacionalnogo korenja jest razlog zapisa, ne nezavisno izměrjen rezultat. Smysl roditeljskogo termina i imenovana rolja morajut ostati povezane, da čitatelj ne směša sredstvo s njegovymi poslědkami.\par
\noindent\emph{Reason:} Redukcije names the derivations, whereas reducent names a technical role within them. Paper 2 defines a reducer as a reducible series of forms; it is not the person carrying out a calculation or a generic computing tool. Section 18 explicitly uses such a series as a factor from which several reductions follow. The archive retains reducent alongside reducibilny and records the corresponding coined Russian and Ukrainian technical terms. Transparency of the international stem is the recorded reason, not an independently measured outcome. The underlying definition and named role must remain connected so the reader does not confuse the means with its consequences.\par
\noindent{\footnotesize Izvorne zapisy / Original records: HIST-g019-006.}\par
\smallskip\noindent\emph{Sravnjenje s tekstom:} Definicija reducenta daje tehničny smysl imenovanomu sredstvu redukcij.\par
\noindent\emph{Textual comparison:} The definition of reducer supplies the technical sense of the named means of reduction.\par
\noindent{\footnotesize TEXT-COMP-542; pasaže / passages: 547, 551, 2592, 2594, 2595, 2596, 2597, 2598, 2599, 2600, 2601, 2602, 2603, 2604, 2606, 2613; \texttt{CORPUS\_FORM\_COMPARISONS.json}.}\par
\end{minipage}\par
\par\medskip\Needspace{7\baselineskip}
\noindent\begin{minipage}{\linewidth}
\hypertarget{TERM-1476}{}
\noindent\textbf{redukcijne formuly}\quad{\footnotesize\texttt{TERM-1476}}\par
\noindent\emph{Iztočnik / Source:} Reduktionsformeln\par
\noindent\emph{Rusky / Russian:} редукционные формулы; формулы редукции\par
\noindent\emph{Ukrajinsky / Ukrainian:} редукційні формули; формули редукції\par
\smallskip\noindent\emph{Razlog:} Redukcijne formuly imenuje zapisane relacije, ne vse rěšenja redukcijnoj metody. Dva zapisy iz §§5 i 10 raboty 2 birajut kratak pridavny naslov radi gustyh formulačnyh blokov i kontinuiteta s §3. Ruski i ukrajinski primeri sodrživajut i pridavnu konstrukciju, i konstrukciju s genitivom; arhiv ne davaje čislovogo sravnjenja med nimi. §5 izrično govori, ktore formule nastavajut polarizacijeju i izključaje relacije za razložime formy. Tuty izključek ne smě izčeznuti pri skračenju termina. Kratkost naslova jest izložiteljska prědnost, ne pravo skratiti sam matematičny obseg formul.\par
\noindent\emph{Reason:} Redukcijne formuly names the written relations, not every result obtainable from the reduction method. Two records from Paper 2, §§5 and 10, choose a short adjectival heading before dense formula blocks and for continuity with §3. The Russian and Ukrainian comparanda include both adjectival and genitive constructions; the archive supplies no numerical ranking between them. Section 5 expressly identifies which formulas arise by polarization and excepts relations for decomposable forms. That exception must survive any abbreviation of the term. Brevity of the heading is an expository benefit, not permission to reduce the mathematical scope of the formulas.\par
\noindent{\footnotesize Izvorne zapisy / Original records: HIST-g006-009, HIST-g011-010.}\par
\smallskip\noindent\emph{Sravnjenje s tekstom:} Naslov može byti kratak, ale uslovje polarizacije i njegov izključek morajut ostati polne.\par
\noindent\emph{Textual comparison:} The heading may be brief, but the polarization condition and its exception must remain complete.\par
\noindent{\footnotesize TEXT-COMP-544; pasaže / passages: 775, 778, 779, 780, 781, 816, 831, 832, 1316, 1318, 1320; \texttt{CORPUS\_FORM\_COMPARISONS.json}.}\par
\end{minipage}\par
\par\medskip\Needspace{7\baselineskip}
\noindent\begin{minipage}{\linewidth}
\hypertarget{TERM-1477}{}
\noindent\textbf{redukcijny teorem}\quad{\footnotesize\texttt{TERM-1477}}\par
\noindent\emph{Iztočnik / Source:} Reduktionssatz\par
\noindent\emph{Rusky / Russian:} редукционная теорема\par
\noindent\emph{Ukrajinsky / Ukrainian:} редукційна теорема\par
\smallskip\noindent\emph{Razlog:} Redukcijny teorem označaje konkretno tvrđenje s dokazom, ne samo neformalnu tehniku. Za rabotu 8 arhiv naglašaje rolju teorema v cělom dokazu; v rabotě 16 razlog upuča na specialny poslědok formuly (2). Nynějša pasaža raboty 16 čuva avtorovo ograničenje: ranějši postup davaje redukcijny teorem, ale bez izraza za razliku ne davaje polnogo razklada v redy. Arhivna izjava o kontinuitetu s rabotoju 12 ne ustanavja, že ide o toj samy teorem; tvrdženja trěba identifikovati po formulah i upučenjah. Takože prědnost nazvanja teorem nad tvrdženje jest lokalna urednička ocěna, ne universalny prěklad německogo Satz.\par
\noindent\emph{Reason:} Redukcijny teorem denotes a particular proved assertion, not merely an informal technique. For Paper 8 the archive emphasizes its role in the whole argument; for Paper 16 it identifies a special consequence of formula (2). The current Paper 16 passage preserves the author's limitation: the earlier procedure yields the reduction theorem, but without an expression for the difference it does not establish the full series expansion. The archival claim of continuity with Paper 12 does not establish that these are the same theorem; statements must be identified by their formulas and references. Preference for theorem over proposition is likewise a local editorial judgment, not a universal translation of German Satz.\par
\noindent{\footnotesize Izvorne zapisy / Original records: HIST-g057-004, HIST-g078-012.}\par
\smallskip\noindent\emph{Sravnjenje s tekstom:} Razlikovanje redukcijnogo teorema i polnogo razvitja ograničaje, čo izvedenje dokazuje.\par
\noindent\emph{Textual comparison:} Distinguishing the reduction theorem from the full expansion limits what the derivation establishes.\par
\noindent{\footnotesize TEXT-COMP-545; pasaže / passages: 6751, 6753, 6755, 6778, 6806, 6808, 6810, 6812, 10757, 10759, 10760, 10761, 10762, 10763, 10764, 10765, 10766, 10767, 10768, 10769, 10770, 10771, 10772, 10773, 10774, 10775, 10776, 10777, 10778, 10779; \texttt{CORPUS\_FORM\_COMPARISONS.json}.}\par
\end{minipage}\par
\par\medskip\Needspace{7\baselineskip}
\noindent\begin{minipage}{\linewidth}
\hypertarget{TERM-1478}{}
\noindent\textbf{teorem redukcije}\quad{\footnotesize\texttt{TERM-1478}}\par
\noindent\emph{Iztočnik / Source:} Reduktionssatz\par
\noindent\emph{Rusky / Russian:} теорема редукции\par
\noindent\emph{Ukrajinsky / Ukrainian:} редукційна теорема; теорема редукції\par
\smallskip\noindent\emph{Razlog:} Teorem redukcije jest genitivna alternativa k redukcijny teorem. Zapisy za raboty 12 i 23 čuvaju nazvanje strukturnogo tvrđenja, ktore svodi pytanja diferencialnyh invariantov k algebraičnoj, linearnoj teoriji invariantov. Oni odbirajut redukciju vměsto neoprěděljenogo sniženja i teorem vměsto širšego principa. Arhiv odkazuje na kontinuitet s rabotoju 21; konkretne pasaže pokazyvajut, čo imenovano tvrđenje čini. Ne davajut dokaz, že genitivna konstrukcija jest jezyčno bolja od pridavnoj. Važno jest takože ograničenje: utvrdženje za variacionny kontekst ne prěnosi se avtomatično na libovoljnu novuju strukturu prěnosa.\par
\noindent\emph{Reason:} Teorem redukcije is the genitive alternative to redukcijny teorem. The records for Papers 12 and 23 preserve the name of a structural assertion reducing questions about differential invariants to algebraic, linear invariant theory. They choose reduction rather than vague lowering, and theorem rather than the broader principle. The archive points to continuity with Paper 21; the actual passages show what the named assertion does. They do not establish that the genitive construction is linguistically better than the adjective. The scope restriction also matters: a result in the variational setting does not automatically extend to an arbitrary new transport structure.\par
\noindent{\footnotesize Izvorne zapisy / Original records: HIST-g074-006, HIST-g094-013.}\par
\smallskip\noindent\emph{Sravnjenje s tekstom:} Genitivno nazvanje ne ustanavja novo tvrđenje; smysl čuvajut konkretne definicije i ograničenja.\par
\noindent\emph{Textual comparison:} The genitive label does not establish a new theorem; the definitions and restrictions establish its sense.\par
\noindent{\footnotesize TEXT-COMP-546; pasaže / passages: 8974, 8976, 8978, 8979, 9083, 9086, 9087, 9088, 9089, 9090, 9091, 14147, 14149, 14151, 14153, 15201, 15203, 15205; \texttt{CORPUS\_FORM\_COMPARISONS.json}.}\par
\end{minipage}\par
\par\medskip\Needspace{7\baselineskip}
\noindent\begin{minipage}{\linewidth}
\hypertarget{TERM-1479}{}
\noindent\textbf{redukcijne teoremy}\quad{\footnotesize\texttt{TERM-1479}}\par
\noindent\emph{Iztočnik / Source:} Reduktionssätze\par
\noindent\emph{Rusky / Russian:} редукционные теоремы\par
\noindent\emph{Ukrajinsky / Ukrainian:} редукційні теореми\par
\smallskip\noindent\emph{Razlog:} Redukcijne teoremy jest množna forma naslova, ktory uvodi definicije i tvrđenja o redukciji v rabotě 2, §3. Internacionalny korenj redukcij- s pridavnoju morfologijeju čuva svez s redukcija i reducent; arhiv ocěnjuje takovy naslov kako kratak i razumlivy. No to poslednje ostavaje ocěnoju, ne rezultatom jezyčnogo eksperimenta. Slovo teoremy važe se k eksplicitnym uslovjam, osoblivo k roli reducenta i početkovoj i konečnoj formy reda. Nynějšny uvod raboty ima takože variant redukcijske. Arhiv ne davaje argument, ktory by kvantitativno rangoval te sufiksalne možnosti.\par
\noindent\emph{Reason:} Redukcijne teoremy is the plural heading introducing definitions and assertions about reduction in Paper 2, §3. The international stem redukcij- with adjectival morphology connects redukcija and reducent; the archive judges this heading compact and comprehensible. The latter remains an assessment, not a language experiment. Teoremy is tied to explicit conditions, notably the reducer's role and the initial and final forms of the series. The current introduction also uses redukcijske. No preserved argument quantitatively ranks these suffixal alternatives.\par
\noindent{\footnotesize Izvorne zapisy / Original records: HIST-g005-001.}\par
\smallskip\noindent\emph{Sravnjenje s tekstom:} Redukcijne i redukcijske ostavajut zachovane varianty; doslědna budučna norma ne smě izbrisati razliku med matematičnymi uslovjami.\par
\noindent\emph{Textual comparison:} Redukcijne and redukcijske remain recorded variants; a consistent future norm must not erase differences in mathematical conditions.\par
\noindent{\footnotesize TEXT-COMP-547; pasaže / passages: 180, 540, 543, 545, 547, 549, 551, 553, 4891, 4893, 4895, 4896, 4897, 4898, 4899, 4900, 4901, 4902, 4903, 4904, 4905, 4906, 4908, 4910, 4912, 4914, 4916, 4918; \texttt{CORPUS\_FORM\_COMPARISONS.json}.}\par
\end{minipage}\par
\par\medskip\Needspace{7\baselineskip}
\noindent\begin{minipage}{\linewidth}
\hypertarget{TERM-1480}{}
\noindent\textbf{redukcijske teoremy}\quad{\footnotesize\texttt{TERM-1480}}\par
\noindent\emph{Iztočnik / Source:} Reduktionssätze\par
\noindent\emph{Rusky / Russian:} редукционные теоремы\par
\noindent\emph{Ukrajinsky / Ukrainian:} редукційні теореми\par
\smallskip\noindent\emph{Razlog:} Redukcijske teoremy čuva variant pridavnogo nazvanja v rabotě 4, §9. Arhivny razlog jest algebraičny registar i vernost teoremnomu naslovu. Tekst prostranno oprěděljaje, čo znači vyša forma, i potom formuluje redukcijno tvrđenje za cěly red form. To ne jest prosto metoda za zaměnu složnogo izraza kratšim. Sravnjenje s redukcijne pokazuje opravdane terminologične svezi, ale ne zachovane rěšenje o tom, ktory sufiks imaje byti obča norma. Za novy prěvod mogoče izbrati doslědnu formu po jezyčnom kanonu; historične zapisy i različne matematične uslovja ne smějut byti odstranene, da bi doslědnost izgledala starša než jest.\par
\noindent\emph{Reason:} Redukcijske teoremy retains the adjectival variant in Paper 4, §9. Its archival rationale is algebraic register and fidelity to a theorem-oriented heading. The text carefully defines higher form before stating the reduction assertion for a whole series of forms. This is not simply a method for replacing a complicated expression with a shorter one. Comparison with redukcijne establishes terminological connections, but no preserved decision selecting one suffix as the general norm. A new translation may choose a consistent form through the language canon; historical records and differing mathematical conditions must not be removed to make that consistency appear older than it is.\par
\noindent{\footnotesize Izvorne zapisy / Original records: HIST-g038-002.}\par
\smallskip\noindent\emph{Sravnjenje s tekstom:} Redukcijne i redukcijske ostavajut zachovane varianty; doslědna budučna norma ne smě izbrisati razliku med matematičnymi uslovjami.\par
\noindent\emph{Textual comparison:} Redukcijne and redukcijske remain recorded variants; a consistent future norm must not erase differences in mathematical conditions.\par
\noindent{\footnotesize TEXT-COMP-547; pasaže / passages: 180, 540, 543, 545, 547, 549, 551, 553, 4891, 4893, 4895, 4896, 4897, 4898, 4899, 4900, 4901, 4902, 4903, 4904, 4905, 4906, 4908, 4910, 4912, 4914, 4916, 4918; \texttt{CORPUS\_FORM\_COMPARISONS.json}.}\par
\end{minipage}\par
\par\medskip\Needspace{7\baselineskip}
\noindent\begin{minipage}{\linewidth}
\hypertarget{TERM-1481}{}
\noindent\textbf{teorem redukcije}\quad{\footnotesize\texttt{TERM-1481}}\par
\noindent\emph{Iztočnik / Source:} Reduktionstheorem\par
\noindent\emph{Rusky / Russian:} теорема редукции\par
\noindent\emph{Ukrajinsky / Ukrainian:} теорема редукції\par
\smallskip\noindent\emph{Razlog:} Teorem redukcije v rabotě 21 imenuje prěvod problema diferencialnyh invariantov na problem projektivnyh invariantov krivosti i jeje kovariantnyh derivacij. Arhiv dava prednost kratkomu internacionalnomu korenju, zajedničnomu s ruskoj i ukrajinskoj formoj, i hoče razlikovati toj teorem od algebraičnoj reducibilnosti raboty 20. Sam korenj redukc- ne garantuje to razlikovanje: potrěbno jest izjasniti, čo se na čo svodi. Nynějšny tekst ukazuje takože granicu teorema pri drugoj definiciji prěnosa. Zato nazvanje ne znači neograničeno obečanje redukcije za vsaku geometriju. Arhivna ocěna občeho terminologičnogo statusa ne jest nova prověrka vsih slovjanskyh kanonov.\par
\noindent\emph{Reason:} In Paper 21, teorem redukcije names the passage from a problem about differential invariants to one about projective invariants of curvature and its covariant derivatives. The record favours a compact international stem shared with its Russian and Ukrainian comparanda and seeks to distinguish this theorem from the algebraic reducibility of Paper 20. The stem alone cannot make that distinction: the entry must explain what is reduced to what. The current text also marks the theorem's limits under a different definition of transport. Its title therefore does not promise reduction for every geometry. The archival assessment of established terminology is not a fresh survey of all Slavic scholarly canons.\par
\noindent{\footnotesize Izvorne zapisy / Original records: HIST-g084-010.}\par
\smallskip\noindent\emph{Sravnjenje s tekstom:} Teorem svodi diferencialne invariante na izrično opisane projektivne podatky; daljši odstup označuje granicu pri inoj definiciji prěnosa.\par
\noindent\emph{Textual comparison:} The theorem reduces differential invariants to explicitly described projective data; the subsequent qualification marks a limit under a different definition of transport.\par
\noindent{\footnotesize TEXT-COMP-548; pasaže / passages: 14147, 14149, 14151, 14153; \texttt{CORPUS\_FORM\_COMPARISONS.json}.}\par
\end{minipage}\par
\par\medskip\Needspace{7\baselineskip}
\noindent\begin{minipage}{\linewidth}
\hypertarget{TERM-1482}{}
\noindent\textbf{reducent}\quad{\footnotesize\texttt{TERM-1482}}\par
\noindent\emph{Iztočnik / Source:} Reduzent\par
\noindent\emph{Rusky / Russian:} редуцент\par
\noindent\emph{Ukrajinsky / Ukrainian:} редуцент\par
\smallskip\noindent\emph{Razlog:} Reducent jest historično tehnično nazvanje v teoriji invariantov. Četrnaest zapisov davaje prednost zachovanomu nazvanju prěd mehaničnym reduktorom, občim faktorom ili opisom, ktory samo govori o uproščenju. Priznana cena jest menša idiomatičnost než pri parafrazi; korist jest prepoznavajema svez medžu povtorjenymi definicijami i priměnami. V rabotě 2 reducent ne označuje libovoljnu pomočnu veličinu: definicija važe jego k reducibilnomu redu form. Tekst pri tom citira Gordana. Zato arhivno označenje „coined/historical“ ne dokazuje, že slovo izmyslila Noether ili AI. Ostavaje obvezok čitati točnu definiciju v danom kontekstu; povtorjenje nazvanja samo ne dokazuje jezyčnu pravilnost.\par
\noindent\emph{Reason:} Reducent retains a historical technical noun of invariant theory. Fourteen records favour that identifiable term over a mechanical reducer, a generic factor, or a paraphrase referring only to simplification. They acknowledge reduced idiomatic ease as the cost of keeping a repeated role recognizable across definitions and applications. In Paper 2 the word does not mean any auxiliary quantity: its definition ties it to a reducible series of forms. The passage also cites Gordan. Thus the archive's label “coined/historical” does not establish that either Noether or the AI invented the term. Its precise definition must still be read in each setting; recurrence alone is not evidence of linguistic correctness.\par
\noindent{\footnotesize Izvorne zapisy / Original records: HIST-g001-010, HIST-g005-003, HIST-g009-005, HIST-g011-005, HIST-g012-014, HIST-g013-012, HIST-g014-006, HIST-g015-006, HIST-g017-006, HIST-g020-005, HIST-g023-006, HIST-g024-004, HIST-g025-008, HIST-g026-003.}\par
\smallskip\noindent\emph{Sravnjenje s tekstom:} Nazvanje reducent jest svezano s reducibilnym redom form. Citacija Gordana ne jest dokaz novoj konsultacije jeho raboty.\par
\noindent\emph{Textual comparison:} Reducent is tied to a reducible series of forms. The citation of Gordan is not evidence that his work was newly consulted.\par
\noindent{\footnotesize TEXT-COMP-549; pasaže / passages: 121, 543, 545, 547, 549, 551, 553, 567; \texttt{CORPUS\_FORM\_COMPARISONS.json}.}\par
\end{minipage}\par
\par\medskip\Needspace{7\baselineskip}
\noindent\begin{minipage}{\linewidth}
\hypertarget{TERM-1483}{}
\noindent\textbf{reducenty}\quad{\footnotesize\texttt{TERM-1483}}\par
\noindent\emph{Iztočnik / Source:} Reduzenten\par
\noindent\emph{Rusky / Russian:} редуценты\par
\noindent\emph{Ukrajinsky / Ukrainian:} редуценти\par
\smallskip\noindent\emph{Razlog:} Reducenty jest množna forma od reducent, ne nova teorija redukcije. Zapis za rabotu 2, §26, izrično prodolžuje termin iz §§22–25. Ta svez jest vidima i v samom tekstu: pod naslovom „Reducenty“ sut navedene formy s odkazami na §23 i §24, ktore zajedno s ranějšimi formulami daju potrěbnu redukciju. Tako možno ukazati i razlog terminologičnoj doslědnosti, i konkretno město priměny. Ranějšne čislo pojavjenj v cělom korpusu samo ne by dokazovalo pojavjenje v §26. Arhiv ne čuva novy sravnitelny izbor množnogo nazvanja; to jest prodolženje už uvedenoj morfologije.\par
\noindent\emph{Reason:} Reducenty is the plural of reducent, not a newly proposed theory of reduction. The record for Paper 2, §26, explicitly continues the term from §§22–25. The connection is also visible in the text: a heading “Reducenty” introduces forms referring back to §§23 and 24 which, together with earlier formulas, give the required reduction. This establishes both the reason for terminological continuity and a specific application. A whole-corpus occurrence count would not by itself establish occurrence in §26. The archive records no new comparative selection of a plural term; it extends the morphology already introduced.\par
\noindent{\footnotesize Izvorne zapisy / Original records: HIST-g027-008.}\par
\smallskip\noindent\emph{Sravnjenje s tekstom:} §26 priamo imaje naslov Reducenty i odkazy na dva ranějšne paragrafy. Takovo pojavjenje potvrđaje konkretno prodolženje termina.\par
\noindent\emph{Textual comparison:} Section 26 explicitly has the heading Reducenty and references two earlier sections, establishing a particular continuation of the term.\par
\noindent{\footnotesize TEXT-COMP-550; pasaže / passages: 3334, 3492, 3494, 3497, 3498, 3499, 3501; \texttt{CORPUS\_FORM\_COMPARISONS.json}.}\par
\end{minipage}\par
\par\medskip\Needspace{7\baselineskip}
\noindent\begin{minipage}{\linewidth}
\hypertarget{TERM-1484}{}
\noindent\textbf{reducenty / redukcija}\quad{\footnotesize\texttt{TERM-1484}}\par
\noindent\emph{Iztočnik / Source:} Reduzenten / Reduktion\par
\noindent\emph{Rusky / Russian:} редуценты / редукция\par
\noindent\emph{Ukrajinsky / Ukrainian:} редуценти / редукція\par
\smallskip\noindent\emph{Razlog:} Para reducenty / redukcija prěnosi nazvanje objekta i nazvanje postupka iz raboty 2 v rabotu 3. Zachovana motivacija jest morfologična doslědnost, ne novo rangovanje konkurentov. Uvod raboty 3 to podpira sodržatelno: govori o teoriji reducentov, redukciji poslědovatelnyh osnovnyh kontrakcij i ulozi konstrukcije, ktora ranějše rezultaty razširjaje. No prěnos slova ne smije skryti uslovnost ranějših tvrđenij za obči slučaj. Reducent i redukcija sut svezane slova s različnymi funkcijami, ne vzaimno zaměnljive sinonimy. Ničto v tom zapisu ne oprěděljaje redukciju dokazov ili izračunjenj v drugoj disciplině.\par
\noindent\emph{Reason:} The pair reducenty / redukcija carries the object-name and operation-name from Paper 2 into Paper 3. The preserved motivation is morphological continuity, not a new ranking of competitors. Paper 3's introduction gives that continuity substantive content: it discusses the theory of reducers, reduction of successive basic contractions, and a construction extending earlier results. Terminological continuity must not conceal the conditional status of earlier statements in the general case. Reducent and redukcija are related words with different functions, not interchangeable synonyms. This record establishes no terminology for proof reduction or computation in another discipline.\par
\noindent{\footnotesize Izvorne zapisy / Original records: HIST-g028-008.}\par
\smallskip\noindent\emph{Sravnjenje s tekstom:} Uvod svezava reducenty s redukcijeju i s ranějšimi uslovnymi rezultatami; kontinuitet nazvanja ne odstranjaje uslovja.\par
\noindent\emph{Textual comparison:} The introduction links reducers with reduction and earlier conditional results; continuity of the name does not remove the conditions.\par
\noindent{\footnotesize TEXT-COMP-551; pasaže / passages: 3649, 3651; \texttt{CORPUS\_FORM\_COMPARISONS.json}.}\par
\end{minipage}\par
\par\medskip\Needspace{7\baselineskip}
\noindent\begin{minipage}{\linewidth}
\hypertarget{TERM-1485}{}
\noindent\textbf{reducibilny}\quad{\footnotesize\texttt{TERM-1485}}\par
\noindent\emph{Iztočnik / Source:} reduzibel\par
\noindent\emph{Rusky / Russian:} редуцируемый\par
\noindent\emph{Ukrajinsky / Ukrainian:} редуковний\par
\smallskip\noindent\emph{Razlog:} Reducibilny jest izbrano v rabotě 2 radi internacionalnogo korenja i svezanosti s reducent. Arhiv sravnjuje rusko reducirujemoje i ukrajinsko redukovnyj, pričem nacionalny status tych form jest ocěna ranějšego prevoditelja, ne nova kanonska prověrka. Definicija govori o vyraženju formy prěz druge invariante i o reducibilnyh redah form. Ne govori prosto, že element jest produktom dvoh netrivialnyh faktorov. Zato kratka para „reducibilny = razložimy“ bez definicije byla by zavodna. Težina svezanosti slovnoj familije jest zachovana kvalitativno; čisleny bal ili sravnitelny test čitatelskoj razumělivosti ne jest zachovany.\par
\noindent\emph{Reason:} Paper 2 selects reducibilny for its international stem and connection with reducent. The archive compares Russian редуцируемый and Ukrainian редуковний; its assessments of their national disciplinary status belong to the earlier translator, not a fresh canon check. The definition concerns expressing a form through other invariants and defining reducible series of forms. It does not simply assert that an element factors into two nontrivial factors. A bare equivalence “reducibilny = decomposable” would therefore mislead without the definition. The importance assigned to word-family coherence is recorded qualitatively; no numerical score or comparative reader-comprehension test survives.\par
\noindent{\footnotesize Izvorne zapisy / Original records: HIST-g005-004.}\par
\smallskip\noindent\emph{Sravnjenje s tekstom:} Nazvanje reducent jest svezano s reducibilnym redom form. Citacija Gordana ne jest dokaz novoj konsultacije jeho raboty.\par
\noindent\emph{Textual comparison:} Reducent is tied to a reducible series of forms. The citation of Gordan is not evidence that his work was newly consulted.\par
\noindent{\footnotesize TEXT-COMP-549; pasaže / passages: 121, 543, 545, 547, 549, 551, 553, 567; \texttt{CORPUS\_FORM\_COMPARISONS.json}.}\par
\end{minipage}\par
\par\medskip\Needspace{7\baselineskip}
\noindent\begin{minipage}{\linewidth}
\hypertarget{TERM-1486}{}
\noindent\textbf{razložimy / nerazložimy}\quad{\footnotesize\texttt{TERM-1486}}\par
\noindent\emph{Iztočnik / Source:} reduzibel / irreduzibel\par
\noindent\emph{Rusky / Russian:} разложимый / неразложимый\par
\noindent\emph{Ukrajinsky / Ukrainian:} розкладний / нерозкладний\par
\smallskip\noindent\emph{Razlog:} Razložimy / nerazložimy jest izbrano radi samej definicije: ideal jest prědstavimy kako najmenše obče mnogokratno vlastnyh děliteljev. V upotrěbjenoj idealnoj terminologiji to odgovarjaje prěsěku; ne ide o faktorizaciju jednogo elementa. Ruska i ukrajinska para daju arhivne sravnitelne formy za razloženje, ale ne nezavisne ocěny istogo izbora. Dva zapisa sodržajut jednu zachovanu motivaciju. Nynějšna definicija nalazi se v §2 raboty 19, iako oznaky rabotnyh jedinic upominajut §6 i §12. Oboje trěba zachovati: historičny identitet zapisa i točno město definicije. To značenje ne smije byti složeno v jedno s reducibilnostju polinomov ili prědstavjenij.\par
\noindent\emph{Reason:} Razložimy / nerazložimy is chosen because of the definition itself: an ideal is represented as the least common multiple of proper divisors. In the ideal terminology used here, this corresponds to intersection, not factorization of an element. The Russian and Ukrainian pairs provide historical decomposition comparanda, not independent evaluations of the same choice. Two records preserve one rationale. In the current text the definition is in Paper 19, §2, although the working-unit labels mention §§6 and 12. Both identities matter: the historical record and the definition's actual location. This sense must not be merged with reducibility of polynomials or representations.\par
\noindent{\footnotesize Izvorne zapisy / Original records: HIST-g081-034, HIST-g082-034.}\par
\smallskip\noindent\emph{Sravnjenje s tekstom:} Smysl jest oprěděljen najmenšim občim mnogokratnym vlastnyh děliteljev. Definicija nalazi se v §2, različnom od oznak dvoh historičnyh rabotnyh jedinic.\par
\noindent\emph{Textual comparison:} The sense is defined by the least common multiple of proper divisors. The definition is in §2, distinct from the labels of the two historical working units.\par
\noindent{\footnotesize TEXT-COMP-552; pasaže / passages: 12568, 12600, 12604, 12613; \texttt{CORPUS\_FORM\_COMPARISONS.json}.}\par
\end{minipage}\par
\par\medskip\Needspace{7\baselineskip}
\noindent\begin{minipage}{\linewidth}
\hypertarget{TERM-1487}{}
\noindent\textbf{reducibilny / ireducibilny}\quad{\footnotesize\texttt{TERM-1487}}\par
\noindent\emph{Iztočnik / Source:} reduzibel / irreduzibel\par
\noindent\emph{Rusky / Russian:} приводимый / неприводимый\par
\noindent\emph{Ukrajinsky / Ukrainian:} звідний / незвідний\par
\smallskip\noindent\emph{Razlog:} V rabotě 17 reducibilny / ireducibilny odnosi se k grupě ostatkov s operatorami. Redukcijno uslovje trěbuje razloženje vsakogo elementa na jednoznačno oprěděljene sostavne česti iz dvoh nenulovyh podgrup danogo tipa. Izomorfizm v okolnom tekstu čuva takože djelovanje operatorov; ne dosta jest libovoljno aditivno razloženje. Jednoznačnost sostavnyh čestij pri danom razloženju ne znači jednoznačnost izbora samyh podgrup. Arhiv dava prednost internacionalnomu korenju v medžuslovjanskom, različnomu od zachovanyh ruskyh i ukrajinskyh sravnenij. Čest ukrajinskogo teksta v motivaciji jest stracena; čitatelj može viděti zachovane forme v sravnitelnoj tablici, no ne rekonstruovanu citaciju izgubljenogo razsudka.\par
\noindent\emph{Reason:} In Paper 17, reducibilny / ireducibilny concerns a residue group equipped with operators. The condition requires each element to decompose uniquely into components from two nonzero subgroups of the specified kind. The surrounding definition of isomorphism also preserves operator action; an arbitrary additive decomposition is not enough. Uniqueness of components for a fixed decomposition does not imply uniqueness of the subgroups themselves. The archive favours the international stem in Interslavic, unlike its preserved Russian and Ukrainian comparanda. Part of the Ukrainian wording inside the rationale is lost. The surviving forms can be displayed in the comparison table, but they do not license an invented quotation of the missing argument.\par
\noindent{\footnotesize Izvorne zapisy / Original records: HIST-g079-025.}\par
\smallskip\noindent\emph{Sravnjenje s tekstom:} Jednoznačnost komponentov mora byti čitana zajedno s definicijeju operatornogo izomorfizma i dopustimoj podgrupy.\par
\noindent\emph{Textual comparison:} Uniqueness of components must be read with the definitions of operator-preserving isomorphism and admissible subgroup.\par
\noindent{\footnotesize TEXT-COMP-553; pasaže / passages: 11191, 11195, 11198, 11206, 11208, 11212, 11214; \texttt{CORPUS\_FORM\_COMPARISONS.json}.}\par
\end{minipage}\par
\par\medskip\Needspace{7\baselineskip}
\noindent\begin{minipage}{\linewidth}
\hypertarget{TERM-1488}{}
\noindent\textbf{redukovalny / reducent}\quad{\footnotesize\texttt{TERM-1488}}\par
\noindent\emph{Iztočnik / Source:} reduzibel / Reduzent\par
\noindent\emph{Rusky / Russian:} редуцируемый / редуцент\par
\noindent\emph{Ukrajinsky / Ukrainian:} звідний / редуцент\par
\smallskip\noindent\emph{Razlog:} Redukovalny / reducent v rabotě 4 čuva tehnično nazvanje reducenta i drugu pridavnu formu za redukcijnu možnost. Definicija prědhodi teoremu IX i važe red form k točnomu uslovju, pri ktorom forma jest vyša od drugoj. Takova vyšina jest oprěděljena kontrakcijami, ne obyčnym rangovanjem po veličině. Arhiv ocěnjuje izbor kako prikladny matematičnomu registru, ale ne sravnjuje redukovalny s ranějšim reducibilny dovolno podrobno, da bi dokazyval prednost jednoj občeje norme. Slovnik čuva oba historična izbora i jih uslovja. Doslědnost v novom tekstu trěbuje jezyčno opravdanje, ne prosto zaměnu vseh različnyh pridavnyh form.\par
\noindent\emph{Reason:} Redukovalny / reducent in Paper 4 preserves the reducer noun while using another adjective for susceptibility to reduction. The definition precedes Theorem IX and ties a series of forms to a precise condition under which one form is higher than another. That ordering is defined through contractions, not an ordinary comparison of size. The archive judges the wording suitable for mathematical register, but does not compare redukovalny with the earlier reducibilny in enough detail to establish a preferred general norm. The dictionary preserves both historical choices and their conditions. Consistency in a new text requires linguistic justification, not indiscriminate replacement of every differing adjective.\par
\noindent{\footnotesize Izvorne zapisy / Original records: HIST-g038-008.}\par
\smallskip\noindent\emph{Sravnjenje s tekstom:} Adjektivny variant naleži k redu form s oprěděljenym porjadkom. Nema zachovanogo rangovanja dvaju medžuslovjanskyh adjektivov.\par
\noindent\emph{Textual comparison:} The adjectival variant belongs to a series of forms with a defined ordering. No ranking of the two Interslavic adjectives is preserved.\par
\noindent{\footnotesize TEXT-COMP-554; pasaže / passages: 4895, 4906, 4910, 4912; \texttt{CORPUS\_FORM\_COMPARISONS.json}.}\par
\end{minipage}\par
\par\medskip\Needspace{7\baselineskip}
\noindent\begin{minipage}{\linewidth}
\hypertarget{TERM-1489}{}
\noindent\textbf{stavaje reducibilny / reducent}\quad{\footnotesize\texttt{TERM-1489}}\par
\noindent\emph{Iztočnik / Source:} reduzibel wird / Reduzent\par
\noindent\emph{Rusky / Russian:} становится приводимым / редуцент\par
\noindent\emph{Ukrajinsky / Ukrainian:} стає звідним / редуцент\par
\smallskip\noindent\emph{Razlog:} Stavaje reducibilny čuva proměnu, vyraženu německym wird, a ne samo statičnu vlastnost. V §17 raboty 2 ta proměna zavisi od vvedenogo modula: tekst govori, kogda forma stavaje reducibilna po tom modulu. Zato izpuščenje dopolnjenja moglo by prevratiti relativno tvrđenje v bezuslovno. Reducent ostavaje nazvanjem drugogo, svezanogo objekta. Arhiv dava prednost doslědnosti redukcijnoj familije, no izrično ostavjaje otvorjeno sravnjenje s možnymi domácimi alternativami. On ne imenuje takovu alternativu i ne dokumentuje završeny ekspertni pregled. Slovnik ne dopolnjaje tu prazninu izmysljenym konkurentom ili izmysljenym glasovanjem.\par
\noindent\emph{Reason:} Stavaje reducibilny preserves the change expressed by German wird, rather than merely naming a static property. In Paper 2, §17, that change is relative to the modulus introduced in the argument: the text specifies when the form becomes reducible modulo it. Omitting that qualification could turn a relative assertion into an unconditional one. Reducent continues to name a different, related object. The record prioritizes continuity of the reduction family but explicitly leaves comparison with possible native alternatives open. It names no such alternative and documents no completed expert review. The dictionary does not fill that gap with an invented competitor or vote.\par
\noindent{\footnotesize Izvorne zapisy / Original records: HIST-g018-005.}\par
\smallskip\noindent\emph{Sravnjenje s tekstom:} Slovesna fraza čuva proměnu po odnošenju k zadannomu modulu, ne bezuslovnu proměnu objekta.\par
\noindent\emph{Textual comparison:} The verbal phrase preserves change relative to the specified modulus, not an unconditional change of the object.\par
\noindent{\footnotesize TEXT-COMP-555; pasaže / passages: 2391, 2394, 2395, 2399; \texttt{CORPUS\_FORM\_COMPARISONS.json}.}\par
\end{minipage}\par
\par\medskip\Needspace{7\baselineskip}
\noindent\begin{minipage}{\linewidth}
\hypertarget{TERM-1490}{}
\noindent\textbf{forma reducibilnosti}\quad{\footnotesize\texttt{TERM-1490}}\par
\noindent\emph{Iztočnik / Source:} Reduzibilitätsform\par
\noindent\emph{Rusky / Russian:} форма приводимости\par
\noindent\emph{Ukrajinsky / Ukrainian:} форма звідності\par
\smallskip\noindent\emph{Razlog:} Forma reducibilnosti imenuje konkretnu pomočnu polinomnu konstrukciju raboty 20. Arhiv izbira internacionalny korenj radi jasnogo tehničnogo nazvanja, pri čem ruska i ukrajinska forma upotrěbljajut svoje slovne familije. Matematična uslovja sut suščestvena: za polinom od najmenje dvoh prěměnnyh, pri specializaciji bez poniženja stepena, neizčezanje toj formy identično v pomočnih neoprěděljenyh jest kriterij absolutnoj nerazložimosti. Ne ide o neizčezanje pri jednom libovoljnom izboru pomočnih čisel. Identično izčezanje odgovarjaje razložimosti. Nazvanje zato ne znači ni reducent iz teorije invariantov, ni prostě opisanije možnosti uproščenja. Izvorne matematične uslovja, ne samo zajedničny korenj, upravjajut priměnoju termina.\par
\noindent\emph{Reason:} Forma reducibilnosti names a specific auxiliary polynomial construction in Paper 20. The archive chooses the international stem to make the technical object recognizable, while its Russian and Ukrainian comparanda use their own word families. The mathematical qualifications are essential: for a polynomial in at least two variables, under specialization without degree loss, being nonzero as a polynomial in the auxiliary indeterminates characterizes absolute irreducibility. This does not mean nonvanishing at one arbitrary numerical choice of the auxiliary variables. Identical vanishing corresponds to reducibility. The term therefore means neither the reducer of invariant theory nor a general possibility of simplification. Its mathematical conditions, not merely a shared stem, govern its use.\par
\noindent{\footnotesize Izvorne zapisy / Original records: HIST-g083-004.}\par
\smallskip\noindent\emph{Sravnjenje s tekstom:} Kriterij svezuje absolutnu nerazložimost s neizčezanjem polinoma v pomočnih neoprěděljenyh. Uslovje o stepenju i směr kriterija sut neodstranime.\par
\noindent\emph{Textual comparison:} The criterion relates absolute irreducibility to a nonzero polynomial in the auxiliary indeterminates. Its degree qualification and direction are essential.\par
\noindent{\footnotesize TEXT-COMP-556; pasaže / passages: 13803, 13805, 13807, 14004, 14006, 14009; \texttt{CORPUS\_FORM\_COMPARISONS.json}.}\par
\end{minipage}\par
\par\medskip\Needspace{7\baselineskip}
\noindent\begin{minipage}{\linewidth}
\hypertarget{TERM-1491}{}
\noindent\textbf{reducibilne predstavjenja}\quad{\footnotesize\texttt{TERM-1491}}\par
\noindent\emph{Iztočnik / Source:} reduzible Darstellungen\par
\noindent\emph{Rusky / Russian:} приводимые представления\par
\noindent\emph{Ukrajinsky / Ukrainian:} звідні зображення\par
\smallskip\noindent\emph{Razlog:} Reducibilne prědstavjenja v rabotě 34 sut prědstavjenja s netrivialnym invariantnym podmodulom. Baza, prilagodžena podmodulu, dava blokovo-triangularnu matricu; blok vniz pod diagonaloju ne mora byti nula. Zato reducibilnost ne znači avtomatično razloženje na direktnu sumu. Tekst oddělno razgledava blokovo-diagonalny slučaj i polnu reducibilnost. Nad občim koeficientnym prstenom takože čuva uslovje razloživosti po koeficientnom modulu, ktoro nad tělom jest izpolnjeno. Arhiv motivuje medžuslovjansky izbor kontinuitetom familije reducibilny / ireducibilny i ocěnjenoju internacionalnoju prozračnostju. To poslednje ne jest izměreny rezultat. Historično predstavjenje i nynějšno prědstavjenje ostavajut različenymi pisanymi formami, ne novymi matematičnymi smyslami.\par
\noindent\emph{Reason:} In Paper 34, reducibilne prědstavjenja are representations with a nontrivial invariant submodule. An adapted basis gives a block-triangular matrix; the block below the diagonal need not vanish. Reducibility therefore does not automatically mean a direct-sum decomposition. The text treats the block-diagonal case and complete reducibility separately. Over a general coefficient ring it also retains a splitting condition on the coefficient module, satisfied over a division ring. The archive motivates the Interslavic choice by continuity with reducibilny / ireducibilny and presumed international transparency. That last claim is not a measured result. Historical predstavjenje and current prědstavjenje remain distinguishable spellings, not different mathematical senses.\par
\noindent{\footnotesize Izvorne zapisy / Original records: HIST-g194-001.}\par
\smallskip\noindent\emph{Sravnjenje s tekstom:} Podmodul daje triangularnu formu pri ukazanyh bazovyh uslovjah. Direktna suma jest oddělno razgledany, silnějši slučaj.\par
\noindent\emph{Textual comparison:} A submodule gives triangular form under the stated basis conditions. A direct sum is a separately treated, stronger case.\par
\noindent{\footnotesize TEXT-COMP-557; pasaže / passages: 21486, 21488, 21490, 21494, 21496, 21513, 21515, 21517, 21519, 21543, 21545, 21547, 21549; \texttt{CORPUS\_FORM\_COMPARISONS.json}.}\par
\end{minipage}\par
\par\medskip\Needspace{7\baselineskip}
\noindent\begin{minipage}{\linewidth}
\hypertarget{TERM-1492}{}
\noindent\textbf{redukovano v odnosu k komplementu / k $R$}\quad{\footnotesize\texttt{TERM-1492}}\par
\noindent\emph{Iztočnik / Source:} reduziert in bezug auf das Komplement / in bezug auf $R$\par
\noindent\emph{Rusky / Russian:} редуцированное относительно дополнения / относительно $R$\par
\noindent\emph{Ukrajinsky / Ukrainian:} зредуковане щодо комплемента / щодо $R$\par
\smallskip\noindent\emph{Razlog:} Redukovano v odnosu k čuva točny objekt, po odnošenju k ktoromu prědstavjenje jest redukovano. V lemi VI raboty 19 najkratše prědstavjenje jest najprvo redukovano po odnošenju k komplementu; izolovanost ideala dopuščaje izvesti redukovanost takože po odnošenju k njemu. Odbaciti dopolnjenje ili prědstaviti odnos kako avtomatično simetričny značilo by odstraniti čest tvrđenja. Arhivny razlog jest imenno čuvanje toj napravjenoj zavisimosti, ne prosto obča fraza o uproščenju. Nynějšna lema jest v §7. Historičnu konstrukciju v odnosu k ne podměnjajemo pozdnějšoj domovoj prednostju po odnošenju k; njih odnos jest izjasnjen oddělno. Matematično $R$ jest identifikator, ne slovo za prěvod.\par
\noindent\emph{Reason:} Redukovano v odnosu k preserves the object relative to which the representation is reduced. In Paper 19, Lemma VI, a shortest representation is initially reduced with respect to the complement; isolation of the ideal permits reducedness with respect to that ideal to be deduced as well. Dropping the complement of the phrase, or treating the relation as automatically symmetric, would remove part of the assertion. Preserving this directed dependence is the recorded reason, not a generic reference to simplification. The current lemma is in §7. The historical construction v odnosu k is not silently replaced by the later house preference po odnošenju k; their relationship is discussed separately. Mathematical $R$ is an identifier, not a word to translate.\par
\noindent{\footnotesize Izvorne zapisy / Original records: HIST-g081-058, HIST-g082-058.}\par
\smallskip\noindent\emph{Sravnjenje s tekstom:} Lema prenosi vlastnost od komplementa k izolovanomu idealu pod konkretnymi uslovjami. Slovo komplement i identifikator $R$ ne možno izpustiti.\par
\noindent\emph{Textual comparison:} The lemma transfers the property from the complement to the isolated ideal under specific hypotheses. Neither complement nor identifier $R$ may be dropped.\par
\noindent{\footnotesize TEXT-COMP-558; pasaže / passages: 12570, 13212, 13214, 13216, 13218, 13222; \texttt{CORPUS\_FORM\_COMPARISONS.json}.}\par
\end{minipage}\par
\par\medskip\Needspace{7\baselineskip}
\noindent\begin{minipage}{\linewidth}
\hypertarget{TERM-1493}{}
\noindent\textbf{redukovano predstavjenje}\quad{\footnotesize\texttt{TERM-1493}}\par
\noindent\emph{Iztočnik / Source:} reduzierte Darstellung\par
\noindent\emph{Rusky / Russian:} редуцированное представление; сокращенный вид\par
\noindent\emph{Ukrajinsky / Ukrainian:} редуковане подання; скорочене подання; скорочений вигляд\par
\smallskip\noindent\emph{Razlog:} Redukovano predstavjenje sodrživaje dva smysla, ktore ne smějut byti slite. V rabotě 6 ide o zapis racionalnoj funkcije kako drobi so vzaimno prostym čisliteljem i imenovateljem. Vzaimna prostost polinomov ne znači, že pri vsakoj specializaciji oni ne mogut zajedno izčeznuti; okolny tekst čuva uslovja dopustimogo podstavjenja. V rabotě 19 ide o prědstavjenje ideala najmenšim občim mnogokratnym komponentov. Redukovanost tam trěbuje dve uslovja, najkratšest samo prvo. Arhiv zato odbacaje poistovětjenje redukovanogo s najkratšim, a ne samo predlaga drugu formu slovesa „skratiti“. Jedno slovnično zaglavje može sodržavati oba smysla, ale svaki trěbuje svoju definiciju, svoje priměry i svoje omeženja.\par
\noindent\emph{Reason:} Redukovano predstavjenje contains two senses which must not be collapsed. In Paper 6 it denotes a rational function written as a fraction with coprime numerator and denominator. Coprime multivariate polynomials can still vanish together at a specialization; the surrounding text retains the conditions on admissible substitution. In Paper 19 it denotes an ideal represented as the least common multiple of its components. Reducedness there requires two conditions, whereas shortest representation requires only the first. The archive therefore rejects identification of reduced with shortest; it is not merely choosing another word for cancelling a fraction. One dictionary heading may contain both senses, but each needs its own definition, examples, and qualifications.\par
\noindent{\footnotesize Izvorne zapisy / Original records: HIST-g041-004, HIST-g042-011, HIST-g081-030, HIST-g082-030.}\par
\smallskip\noindent\emph{Sravnjenje s tekstom:} Drob s vzaimno prostymi polinomami i redukovano idealno razloženje sut različne konstrukcije. Vo vtorom slučaji najkratše i redukovano ne sut ravnoznačne.\par
\noindent\emph{Textual comparison:} A fraction with coprime polynomials and a reduced ideal decomposition are different constructions. In the second setting, shortest and reduced are not equivalent.\par
\noindent{\footnotesize TEXT-COMP-559; pasaže / passages: 5061, 5193, 5212, 5213, 5214, 5215, 5216, 5217, 12570; \texttt{CORPUS\_FORM\_COMPARISONS.json}.}\par
\end{minipage}\par
\par\medskip\Needspace{7\baselineskip}
\noindent\begin{minipage}{\linewidth}
\hypertarget{TERM-1494}{}
\noindent\textbf{redukovany diskriminant}\quad{\footnotesize\texttt{TERM-1494}}\par
\noindent\emph{Iztočnik / Source:} reduzierte Diskriminante\par
\noindent\emph{Rusky / Russian:} приведенный дискриминант\par
\noindent\emph{Ukrajinsky / Ukrainian:} зведений дискримінант\par
\smallskip\noindent\emph{Razlog:} Redukovany diskriminant prodolžuje izbor redukovany slěd iz prědhodnogo paragrafa. V rabotě 34 on jest determinantom matrice redukovanyh slědov produktov bazovyh elementov, pri čem ty slědy sut utvorjene v algebraično zatvorjenom polju. Matrica i jej determinant sut različne objekty; pridavno redukovany ne označuje čislo, ktoro jest prosto menše. Arhiv dava prednost internacionalnoj formě radi doslědnosti i izběganja neprězrěčnoj nacionalnoj kalki. Konkretna odrěčena kalka ne jest imenovana, zato ju ne izmysljajemo. Jasnost izbora proizhodi takože iz definicije i jeje poljovyh uslovij; sam internacionalny vid slova ne dokazuje, že vsakomu čitatelju ono bude neposrědno razumělive.\par
\noindent\emph{Reason:} Redukovany diskriminant continues the preceding section's choice redukovany slěd. In Paper 34 it is the determinant of the matrix of reduced traces of products of basis elements, with those traces formed over an algebraically closed field. Matrix and determinant are distinct objects; reduced does not simply mean a smaller number. The archive favours the international adjective for consistency and to avoid an opaque national calque. It names no particular rejected calque, so none is invented here. Clarity also depends on the definition and its field conditions: the international appearance of a word does not establish immediate comprehension by every reader.\par
\noindent{\footnotesize Izvorne zapisy / Original records: HIST-g203-003.}\par
\smallskip\noindent\emph{Sravnjenje s tekstom:} Redukovany diskriminant jest determinant, ne sama matrica. Definicija ukazuje, nad kakym poljem sut obrazovane redukovane slědy.\par
\noindent\emph{Textual comparison:} The reduced discriminant is a determinant, not the matrix itself. The definition specifies the field over which the reduced traces are formed.\par
\noindent{\footnotesize TEXT-COMP-560; pasaže / passages: 22301, 22305, 22307, 22309; \texttt{CORPUS\_FORM\_COMPARISONS.json}.}\par
\end{minipage}\par
\par\medskip\Needspace{7\baselineskip}
\noindent\begin{minipage}{\linewidth}
\hypertarget{TERM-1495}{}
\noindent\textbf{redukovany slěd}\quad{\footnotesize\texttt{TERM-1495}}\par
\noindent\emph{Iztočnik / Source:} reduzierte Spur\par
\noindent\emph{Rusky / Russian:} приведенный след; приведённый след\par
\noindent\emph{Ukrajinsky / Ukrainian:} зведений слід\par
\smallskip\noindent\emph{Razlog:} Redukovany slěd zachovuje internacionalno pridavno i slavjansko nazvanje slěda. Arhiv hoče izběgti privilegovanje jednoj nacionalnoj kalki i ocěnjuje kombinaciju kako razumělivu; test čitateljev ne jest zachovany. Definicija raboty 34 jest suma slědov različnyh ireducibilnyh prědstavjenij, v odličenju od glavnogo slěda regularnogo prědstavjenja. Tekst izrično upozornjaje, že redukovany slěd ne ostavaje neizměnjeny pri prěhodu k algebraično zatvorjenomu polju. Zato jeho historično značenje ne možno bez prověrky poistovětiti so vsakoju sovrěmennoju upotrěboju togo nazvanja. Formula v razščepjenom matričnom slučaji ne opravdava obču definiciju prěz děljenje glavnogo slěda čislom vo vsakoj karakteristikě.\par
\noindent\emph{Reason:} Redukovany slěd combines an international adjective with the Slavic trace noun. The archive seeks to avoid privileging one national calque and judges the combination intelligible; no reader test is preserved. Paper 34 defines it as the sum of traces of distinct irreducible representations, unlike the principal trace of the regular representation. The text explicitly warns that reduced trace need not remain unchanged on passage to an algebraically closed field. Its historical meaning therefore cannot be equated without checking with every modern use of the name. The formula in the split matrix case does not justify defining reduced trace generally by dividing principal trace by an integer in every characteristic.\par
\noindent{\footnotesize Izvorne zapisy / Original records: HIST-g202-003, HIST-g203-005.}\par
\smallskip\noindent\emph{Sravnjenje s tekstom:} Suma slědov različnyh ireducibilnyh prědstavjenij jest razlikovana od slěda regularnogo prědstavjenja. Tekst čuva upozornjenje o razširjenju polja.\par
\noindent\emph{Textual comparison:} The sum of traces of distinct irreducible representations is distinguished from the trace of the regular representation. The text retains its warning about field extension.\par
\noindent{\footnotesize TEXT-COMP-561; pasaže / passages: 22203, 22209, 22217, 22219, 22221, 22234; \texttt{CORPUS\_FORM\_COMPARISONS.json}.}\par
\end{minipage}\par
\par\medskip\Needspace{7\baselineskip}
\noindent\begin{minipage}{\linewidth}
\hypertarget{TERM-1496}{}
\noindent\textbf{redukovany stepen}\quad{\footnotesize\texttt{TERM-1496}}\par
\noindent\emph{Iztočnik / Source:} reduzierter Grad\par
\noindent\emph{Rusky / Russian:} приведенная степень\par
\noindent\emph{Ukrajinsky / Ukrainian:} зведений степінь\par
\smallskip\noindent\emph{Razlog:} Redukovany stepen v rabotě 24 označuje stepenj podpolja prvogo roda v konečnom razširjenju nedoskonalogo polja. To jest konkretna algebraična charakteristika, ne rezultat proizvoljnogo poniženja stepena polinoma. Arhivni razlog za redukovany jest imenno izběganje togo obyčnogo čitanja. Oznaka $r$ važe nazvanje k definiciji; eksponent $f$ jest drugy podatok o razširjenju. Ranějšny zapis stepen i pozdnějšny stepenj ostavajut vidime, zato že arhiv ne dava sravnitelno obrazloženje njihove razlike. Slovnik objasnjaje matematičny smysl bez tvrđenja, že vsaka jezyčna ili matematična podrobnost okolnogo teksta už jest opravdana. V svezanom zapisu o eksponentě jest označeno otvorjeno pitanje o priměru.\par
\noindent\emph{Reason:} In Paper 24, redukovany stepen denotes the degree of the first-kind subfield in a finite extension of an imperfect field. It is a particular algebraic invariant, not an arbitrary lowering of a polynomial's degree. Avoiding that ordinary reading is the recorded reason for redukovany. The symbol $r$ connects the term to the definition; the exponent $f$ is a different datum of the extension. Earlier stepen and later stepenj remain visible because the archive provides no comparative explanation of their difference. The dictionary clarifies the mathematical sense without certifying every linguistic or mathematical detail of the surrounding text. The related exponent entry identifies an unresolved question about the example.\par
\noindent{\footnotesize Izvorne zapisy / Original records: HIST-g095-019.}\par
\smallskip\noindent\emph{Sravnjenje s tekstom:} Definicija razlikuje $r$ i $f$; teorem dalje uvodi $g$. Primě$r$ s dvoma kvadratnymi korenjami v karakteristikě dva imaje neobjasnjeno nesoglasje s ranějšim opisom $f$; to ne jest dozvoljenje za tihu popravu originala.\par
\noindent\emph{Textual comparison:} The definition distinguishes $r$ and $f$; the theorem further introduces $g$. The characteristic-two example with two square roots has an unexplained discrepancy against the earlier description of $f$; this does not authorize silent emendation of the original.\par
\noindent{\footnotesize TEXT-COMP-562; pasaže / passages: 15365, 15669, 15673, 15675, 15677, 15681, 15683, 15685; \texttt{CORPUS\_FORM\_COMPARISONS.json}.}\par
\end{minipage}\par
\par\medskip\Needspace{7\baselineskip}
\noindent\begin{minipage}{\linewidth}
\hypertarget{TERM-1497}{}
\noindent\textbf{redukovany stepenj / eksponenta}\quad{\footnotesize\texttt{TERM-1497}}\par
\noindent\emph{Iztočnik / Source:} reduzierter Grad / Exponent\par
\noindent\emph{Rusky / Russian:} редуцированная степень / экспонента\par
\noindent\emph{Ukrajinsky / Ukrainian:} редукований степінь / експонента\par
\smallskip\noindent\emph{Razlog:} Redukovany stepenj / eksponenta čuva dva različne podatka v razsudku raboty 24. Pomočna lema VIII govori o konečnom razširjenju, podpolju prvogo roda i medžupolju; teorem XIV dalje razlikuje eksponent $f$ od eksponenta $g$ v stepeni normy. Arhiv ocěnjuje redukovany stepenj kako prozračnu tehničnu konstrukciju i predlaga ekspertno razgledanje, ale ne dokumentuje, že ono bylo izvedeno. Nynějšny priměr polja karakteristiki dva naziva eksponentom čislo dva, iako kvadrat vsakogo elementa leži v osnovnom polju. To vytvarjaje pitanje po odnošenju k ranějšej definiciji. Bez sravnjenja s německym originalom ne pripisujemo nesoglasje avtoru i ne zaměnjajemo čislo v prěvodě.\par
\noindent\emph{Reason:} Redukovany stepenj / eksponenta preserves two distinct data in Paper 24's argument. Auxiliary Lemma VIII concerns a finite extension, its first-kind subfield, and an intermediate field; Theorem XIV further distinguishes exponent $f$ from exponent $g$ in the power defining the norm. The archive judges redukovany stepenj a transparent technical construction and proposes expert consideration, but records no completed review. The current characteristic-two example calls the exponent two although every element's square lies in the base field. This raises a question against the earlier definition. Without comparison with the German original, the discrepancy is not attributed to the author and the number is not silently changed in the translation.\par
\noindent{\footnotesize Izvorne zapisy / Original records: HIST-g100-008.}\par
\smallskip\noindent\emph{Sravnjenje s tekstom:} Definicija razlikuje $r$ i $f$; teorem dalje uvodi $g$. Primě$r$ s dvoma kvadratnymi korenjami v karakteristikě dva imaje neobjasnjeno nesoglasje s ranějšim opisom $f$; to ne jest dozvoljenje za tihu popravu originala.\par
\noindent\emph{Textual comparison:} The definition distinguishes $r$ and $f$; the theorem further introduces $g$. The characteristic-two example with two square roots has an unexplained discrepancy against the earlier description of $f$; this does not authorize silent emendation of the original.\par
\noindent{\footnotesize TEXT-COMP-562; pasaže / passages: 15365, 15669, 15673, 15675, 15677, 15681, 15683, 15685; \texttt{CORPUS\_FORM\_COMPARISONS.json}.}\par
\end{minipage}\par
\par\medskip\Needspace{7\baselineskip}
\noindent\begin{minipage}{\linewidth}
\hypertarget{TERM-1498}{}
\noindent\textbf{regularny sistem vrědnostij / singularnosti}\quad{\footnotesize\texttt{TERM-1498}}\par
\noindent\emph{Iztočnik / Source:} regulaeres Wertsystem / Singularitaeten\par
\noindent\emph{Rusky / Russian:} регулярная система значений / сингулярности\par
\noindent\emph{Ukrajinsky / Ukrainian:} регулярна система значень / сингулярності\par
\smallskip\noindent\emph{Razlog:} Regularny sistem vrědnostij čuva sistem podstavjenyh cěnnostij prěměnnyh, ne samo neoprěděljenu „dobru točku“. Arhiv izbira sistem radi vidimej struktury n-torky. V analitičnom dokazu raboty 17 potrěbno jest vybrati cěnnosti, pri ktoryh veduče koeficienty ne izčezajut i ostale nastupajuče koeficienty sut regularne v okolini. Singularnosti sut ono, čemu se pri tom izboru izběgaje. Takova regularnost ne jest regularnost matrice ili prědstavjenja. Arhivno vrědnostij i nynějšno cěnnostij ne treba tajno poistovětiti kako jeden zapis. Ni toj izbor, ni jeho matematika ne rozhoduje o nazvanju regularnyh jezykov v logikě i teoriji avtomatov.\par
\noindent\emph{Reason:} Regularny sistem vrědnostij preserves a system of assigned variable values, not merely an unspecified “good point”. The record prefers system to retain the tuple's explicit structure. Paper 17's analytic proof requires values where the leading coefficients are nonzero and the other coefficients involved are regular in a neighbourhood. Singularities are what this selection avoids. This regularity is neither matrix regularity nor the regular representation. Historical vrědnostij and current cěnnostij must not silently be treated as one spelling. Neither this lexical decision nor its mathematics settles the terminology of regular languages in logic and automata theory.\par
\noindent{\footnotesize Izvorne zapisy / Original records: HIST-g079-068.}\par
\smallskip\noindent\emph{Sravnjenje s tekstom:} Sistem cěnnostij jest vybrany s obzirom na koeficienty i singularnosti. To ne jest matrična ili reprezentacijna regularnost.\par
\noindent\emph{Textual comparison:} The value tuple is selected with regard to coefficients and singularities. This is not matrix or representation-theoretic regularity.\par
\noindent{\footnotesize TEXT-COMP-563; pasaže / passages: 11988, 11989, 12001, 12015, 12016, 12018, 12019, 12020, 12021, 12022, 12023, 12024; \texttt{CORPUS\_FORM\_COMPARISONS.json}.}\par
\end{minipage}\par
\par\medskip\Needspace{7\baselineskip}
\noindent\begin{minipage}{\linewidth}
\hypertarget{TERM-1499}{}
\noindent\textbf{regularna}\quad{\footnotesize\texttt{TERM-1499}}\par
\noindent\emph{Iztočnik / Source:} regulär\par
\noindent\emph{Rusky / Russian:} регулярная\par
\noindent\emph{Ukrajinsky / Ukrainian:} регулярна\par
\smallskip\noindent\emph{Razlog:} Regularna čuva izrično oprěděljeno nazvanje matrice v rabotě 34. Definicija trěbuje pravu i lěvu obratnu matricu nad tělom; slědujuči dokaz pokazuje, že v tom okviru dostatočna jest jedna strana. Arhiv zato ne zaměnjaje avtorsko nazvanje občim opisom obratimosti. To ne izključuje izjasnjenje smysla, ale čuva razlikovanje medžu definicijeju i teoremom o njej. Konečno -a jest ženskogo roda radi slova matrica, ne osobna nova lema. Pri novoj priměně treba prověriti osnovnu strukturu: argument za konečne kvadratne matrice nad tělom ne jest bezuslovno tvrđenje o všech odnostrannyh obratnyh elementah libovoljnogo prstena.\par
\noindent\emph{Reason:} Regularna preserves an explicitly defined matrix term in Paper 34. The definition requires right and left inverse matrices over a division ring; the following argument establishes that one side suffices in this setting. The archive therefore retains the author's term instead of replacing it with a general description of invertibility. Explaining the meaning is still appropriate, but definition and theorem about the definition remain distinct. The ending -a agrees with feminine matrica; it is not a separate new lemma. Reuse requires checking the underlying structure: the argument for finite square matrices over a division ring is not an unrestricted assertion about one-sided inverses in arbitrary rings.\par
\noindent{\footnotesize Izvorne zapisy / Original records: HIST-g185-002.}\par
\smallskip\noindent\emph{Sravnjenje s tekstom:} Dve obratne matrice sut v definiciji; dostatočnost jednoj strany jest slědujuče tvrđenje v zadannom algebraičnom okviru.\par
\noindent\emph{Textual comparison:} Two inverse matrices occur in the definition; sufficiency of one side is a subsequent assertion in the specified algebraic setting.\par
\noindent{\footnotesize TEXT-COMP-564; pasaže / passages: 20327, 20329, 20331, 20333, 20335, 20343, 20345; \texttt{CORPUS\_FORM\_COMPARISONS.json}.}\par
\end{minipage}\par
\par\medskip\Needspace{7\baselineskip}
\noindent\begin{minipage}{\linewidth}
\hypertarget{TERM-1500}{}
\noindent\textbf{regularno predstavjenje}\quad{\footnotesize\texttt{TERM-1500}}\par
\noindent\emph{Iztočnik / Source:} reguläre Darstellung\par
\noindent\emph{Rusky / Russian:} регулярное представление\par
\noindent\emph{Ukrajinsky / Ukrainian:} регулярне зображення; регулярне представлення\par
\smallskip\noindent\emph{Razlog:} Regularno predstavjenje jest nazvanje prědstavjenja algebraičnogo sistema jego vlastnym jediničnym idealom. V grupovom prstenu baza iz grupovyh elementov dava konkretno permutacijno prědstavjenje. Arhivni razlogy sut internacionalno pridavno, disciplinarno nazvanje i doslědnost s ranějšim predstavjenje. To jest kvalitativna kontinuitetna odluka, ne tri nezavisne jezyčne prověrky. Regularno tu ne znači, že vsaka prědstavjajuča matrica jest obratima: algebra sodrživaje nulu i može sodržavati takože nenulove neobratime elementy. Pri izjasnjenju grupovogo priměra treba čuvati uslovje o karakteristikě iz teksta. Pisanje prědstavjenje v nynějšnom korpusu takože ne treba tajno pripisovati staršemu zapisu bez ě.\par
\noindent\emph{Reason:} Regularno predstavjenje names the representation of an algebraic system on its own unit ideal. In a group ring, the group-element basis gives a concrete permutation representation. The recorded reasons are an international adjective, disciplinary terminology, and continuity with the earlier predstavjenje. This is a qualitative continuity decision, not three independent linguistic evaluations. Regular does not mean that every representing matrix is invertible: an algebra includes zero and may include nonzero nonunits. Explanation of the group example must retain the text's characteristic hypothesis. The current corpus's spelling prědstavjenje must not silently be attributed to the older record without ě.\par
\noindent{\footnotesize Izvorne zapisy / Original records: HIST-g198-007, HIST-g201-004, HIST-g204-004.}\par
\smallskip\noindent\emph{Sravnjenje s tekstom:} Regularno prědstavjenje jest oprěděljeno nositeljem i dějstvom. Grupovy priměr i njegovo uslovje o karakteristikě ne měnjajut smysla nazvanja v tvrđenje, že vse matrične obrazy sut obratime.\par
\noindent\emph{Textual comparison:} The regular representation is defined by its carrier and action. The group example and its characteristic hypothesis do not turn the name into an assertion that all matrix images are invertible.\par
\noindent{\footnotesize TEXT-COMP-565; pasaže / passages: 21861, 21863, 22123, 22125, 22449, 22451, 22453, 22455, 22457; \texttt{CORPUS\_FORM\_COMPARISONS.json}.}\par
\end{minipage}\par
\par\medskip\Needspace{7\baselineskip}
\noindent\begin{minipage}{\linewidth}
\hypertarget{TERM-1501}{}
\noindent\textbf{regularne sistemy iz polinomov}\quad{\footnotesize\texttt{TERM-1501}}\par
\noindent\emph{Iztočnik / Source:} reguläre Systeme aus Polynomen\par
\noindent\emph{Rusky / Russian:} регулярные системы из полиномов\par
\noindent\emph{Ukrajinsky / Ukrainian:} регулярні системи з поліномів\par
\smallskip\noindent\emph{Razlog:} Regularne sistemy iz polinomov označajut klasu, oprěděljenu v rabotě 6 definicijeju VIII, ne prosto „dobre“ sistemy. Uslovje trěbuje, aby v oblasti odobraženja najmenšej sodrživajučej integralnoj oblasti obstajalo potrěbno čislo polinomov, čije najvyšše homogene česti imajut nenulovy rezultant. Arhiv izbira prěmu internacionalnu formu radi prepoznavajemosti definicije. Nova konsultacija slovnika potvrđaje samo regularny kako obče pridavno slovo, ne ovu specialnu definiciju. Teorem o integralnoj bazě jest poslědok uslovja, ne zaměna za njega. Isto slovo regularny v drugyh rabotah može značiti drugu vlastnost; slovnična doslědnost trěbuje zachovati te razlike.\par
\noindent\emph{Reason:} Regularne sistemy iz polinomov denotes the class specified by Definition VIII in Paper 6, not simply “well-behaved” systems. The condition requires the appropriate number of polynomials in the image domain of the smallest containing integral domain, whose highest homogeneous parts have nonzero resultant. The archive chooses a direct international expression to keep that definition recognizable. Fresh dictionary consultation attests regularny as a general adjective, not this particular definition. Existence of an integral basis is a consequence of the condition, not a substitute for it. The same adjective elsewhere may name another property; terminological consistency requires preserving those differences.\par
\noindent{\footnotesize Izvorne zapisy / Original records: HIST-g052-001.}\par
\smallskip\noindent\emph{Sravnjenje s tekstom:} Definicija VIII dava algebraično uslovje, a geometrične tvrđenja razlikujut oblast odobraženja od samogo homogenogo sistema. Slovničny zapis 5671 potvrđaje regularny, ne specialnu definiciju regularnosti.\par
\noindent\emph{Textual comparison:} Definition VIII gives the algebraic condition, while the geometric statements distinguish the image domain from the homogeneous system itself. Dictionary row 5671 attests regularny, not the particular regularity definition.\par
\noindent{\footnotesize TEXT-COMP-566; pasaže / passages: 5025, 6273, 6280, 6284, 6322, 6324, 6328, 6329, 6332, 6334, 6338; \texttt{CORPUS\_FORM\_COMPARISONS.json}.}\par
\end{minipage}\par
\par\medskip\Needspace{7\baselineskip}
\noindent\begin{minipage}{\linewidth}
\hypertarget{TERM-1502}{}
\noindent\textbf{regularne sistemy; bez fundamentalnyh toček}\quad{\footnotesize\texttt{TERM-1502}}\par
\noindent\emph{Iztočnik / Source:} reguläre Systeme; fundamentalpunktlos\par
\noindent\emph{Rusky / Russian:} регулярные системы; без фундаментальных точек\par
\noindent\emph{Ukrajinsky / Ukrainian:} регулярні системи; без фундаментальних точок\par
\smallskip\noindent\emph{Razlog:} Bez fundamentalnyh toček razkriva složeno německo pridavno fundamentalpunktlos kak jasnu predložnu frazu. Arhivni razlog jest čuvanje geometričnogo smysla i prepoznavajemoj svezi s iztočnikom. Fundamentalna točka tu jest oprěděljena posrědstvom homogenizovanyh uravnenij; ona ne znači libovoljnu „važnu točku“ ili prosto singularnost. Uvod dava kratak geometričny opis regularnosti, a §12 točnejše razlikuje oblast odobraženja najmenšej sodrživajučej oblasti od samogo sistema homogenyh form. Frazu „bez toček“ zato ne treba iztrgnuti iz njeje oprěděljene oblasti. Izbor internacionalnogo nazvanja i razvitogo negacijnogo izraza jest kvalitativny; čislene težiny konkurentov ne sut zachovane.\par
\noindent\emph{Reason:} Bez fundamentalnyh toček unpacks the German compound fundamentalpunktlos into a clear prepositional phrase. The recorded reason is preservation of the geometric sense and recognizability of the source terminology. A fundamental point is defined through the homogenized equations, not as any “important point” or simply a singularity. The introduction gives a short geometric account of regularity; §12 distinguishes the image domain of the smallest containing domain from the system of homogeneous forms itself. The phrase “without points” must therefore retain its specified setting. Choosing an international noun phrase with explicit negation is a qualitative decision; no numerical comparison of competitors survives.\par
\noindent{\footnotesize Izvorne zapisy / Original records: HIST-g040-014.}\par
\smallskip\noindent\emph{Sravnjenje s tekstom:} Definicija VIII dava algebraično uslovje, a geometrične tvrđenja razlikujut oblast odobraženja od samogo homogenogo sistema. Slovničny zapis 5671 potvrđaje regularny, ne specialnu definiciju regularnosti.\par
\noindent\emph{Textual comparison:} Definition VIII gives the algebraic condition, while the geometric statements distinguish the image domain from the homogeneous system itself. Dictionary row 5671 attests regularny, not the particular regularity definition.\par
\noindent{\footnotesize TEXT-COMP-566; pasaže / passages: 5025, 6273, 6280, 6284, 6322, 6324, 6328, 6329, 6332, 6334, 6338; \texttt{CORPUS\_FORM\_COMPARISONS.json}.}\par
\end{minipage}\par
\par\medskip\Needspace{7\baselineskip}
\noindent\begin{minipage}{\linewidth}
\hypertarget{TERM-1503}{}
\noindent\textbf{regularny element / nedělitelj nuly}\quad{\footnotesize\texttt{TERM-1503}}\par
\noindent\emph{Iztočnik / Source:} reguläres Element / Nichtnullteiler\par
\noindent\emph{Rusky / Russian:} регулярный элемент / неделитель нуля\par
\noindent\emph{Ukrajinsky / Ukrainian:} регулярний елемент / недільник нуля\par
\smallskip\noindent\emph{Razlog:} Regularny element jest tu dopolnjen smyslovym izjasnjenjem nedělitelj nuly. V rabotě 30 samo nenulovost ne dosta, kogda prsten dopuščaje dělitelje nuly: uslovje se odnosi k generatoru glavnogo modula i k imeniteljam v prstenu kvocientov. Arhiv hoče zachovati to jačše uslovje i ocěnjuje regularny kako stabilny algebraičny termin. Obči slovnik potvrđaje slovo, ale ne njegovu isključnu matematičnu priměnu. Nynějšna podnožna primětka ima frazu ne-nulovy dělitelj. Ona može sugerovati nenulovogo dělitelja vměsto elementa, ktory ne jest děliteljem nuly; tu nejasnost treba izjasniti, ne prědstaviti ju kak už rěšenu. Regularnost takože ne znači avtomatično obratimost v izhodnom prstenu.\par
\noindent\emph{Reason:} Regularny element is here accompanied by the sense-defining gloss nedělitelj nuly. In Paper 30, being nonzero alone is insufficient when the ring permits zero divisors: the stronger condition concerns the generator of a principal module and denominators in the quotient ring. The archive seeks to retain that condition and regards regularny as a stable algebraic term. The general dictionary attests the word, not an exclusive mathematical application. The current footnote says ne-nulovy dělitelj. That wording can suggest a nonzero divisor rather than an element which is not a zero divisor; this ambiguity needs clarification and is not represented as already resolved. Regularity also does not automatically mean invertibility in the original ring.\par
\noindent{\footnotesize Izvorne zapisy / Original records: HIST-g108-008.}\par
\smallskip\noindent\emph{Sravnjenje s tekstom:} Arhivni nedělitelj nuly i nynějšno ne-nulovy dělitelj ne sut bezuslovno jednoznačne frazy. Okolny tekst trěbuje jačše uslovje nego samu nenulovost.\par
\noindent\emph{Textual comparison:} Archived nedělitelj nuly and current ne-nulovy dělitelj are not unambiguously equivalent expressions. The surrounding argument requires more than being nonzero.\par
\noindent{\footnotesize TEXT-COMP-567; pasaže / passages: 16206, 16244, 16261; \texttt{CORPUS\_FORM\_COMPARISONS.json}.}\par
\end{minipage}\par
\par\medskip\Needspace{7\baselineskip}
\noindent\begin{minipage}{\linewidth}
\hypertarget{TERM-1504}{}
\noindent\textbf{regularny polinom vzhodno k $x_i$}\quad{\footnotesize\texttt{TERM-1504}}\par
\noindent\emph{Iztočnik / Source:} reguläres Polynom in bezug auf $x_i$\par
\noindent\emph{Rusky / Russian:} регулярный полином относительно $x_i$\par
\noindent\emph{Ukrajinsky / Ukrainian:} регулярний поліном відносно $x_i$\par
\smallskip\noindent\emph{Razlog:} Regularny polinom vzhodno k $x_i$ sodrživaje neodstranimy argument: regularnost jest oprěděljena po odnošenju k izbranoj prěměnnoj. V rabotě 22 polinom stepena $r$ jest regularny po odnošenju k $x_i$, ako sodrživaje člen $x_i^r$ s nenulovym koeficientom. Definicija ne trěbuje, aby toj koeficient byl $1$, i ne znači prosto „polinom bez singularnostij“. Arhiv dava prednost internacionalnomu pridavnomu i čuvanju definicije; dva zapisa ne tvorjat dva nezavisne jezyčne testa. Starše vzhodno k i nynějšno po odnošenju k treba razlikovati. Argument $x_i$ ostavaje točnym matematičnym identifikatorom pri vsakoj jezyčnoj redakciji.\par
\noindent\emph{Reason:} Regularny polinom vzhodno k $x_i$ includes an indispensable argument: regularity is relative to a chosen variable. Paper 22 calls a degree-$r$ polynomial regular with respect to $x_i$ when it contains the term $x_i^r$ with nonzero coefficient. The definition does not require that coefficient to be one, nor does it simply mean “a polynomial without singularities”. The archive favours the international adjective and retention of the explicit definition; its two records are not two independent language tests. Earlier vzhodno k and current po odnošenju k must remain distinguishable. The argument $x_i$ remains an exact mathematical identifier through linguistic editing.\par
\noindent{\footnotesize Izvorne zapisy / Original records: HIST-g087-006, HIST-g092-026.}\par
\smallskip\noindent\emph{Sravnjenje s tekstom:} Definicija navodi najvyššu stepenj izbranoj prěměnnoj s nenulovym koeficientom. Predložna čest nazvanja zato sodrživaje neobhodimy argument.\par
\noindent\emph{Textual comparison:} The definition specifies the highest power of the chosen variable with nonzero coefficient. The phrase identifying that variable carries an indispensable argument.\par
\noindent{\footnotesize TEXT-COMP-568; pasaže / passages: 14539, 14541, 14544; \texttt{CORPUS\_FORM\_COMPARISONS.json}.}\par
\end{minipage}\par
\par\medskip\Needspace{7\baselineskip}
\noindent\begin{minipage}{\linewidth}
\hypertarget{TERM-1505}{}
\noindent\textbf{redy po $n$ varijablov}\quad{\footnotesize\texttt{TERM-1505}}\par
\noindent\emph{Iztočnik / Source:} Reihen von je $n$ Variabeln\par
\noindent\emph{Rusky / Russian:} ряды по $n$ переменных\par
\noindent\emph{Ukrajinsky / Ukrainian:} ряди по $n$ змінних\par
\smallskip\noindent\emph{Razlog:} Redy po $n$ varijablov čuva raspored v grupe s jednakym čislom elementov. V rabotě 8 forma može zavisěti od $N$ redov, a vsaky red sodrživaje $n$ prěměnnyh; ta dva čisla ne sut zaměnljive. Arhiv izbira red radi kontinuiteta s rabotami 6–7 i zato, že ne ide o beskonečny analitičny red. Nyně izčitany slovničny zapis ręd dava obči smysl „row, queue“, drugi zapis „order, arrangement, sequence“. Oni podpirajut leksičnu osnovu, ale sami ne oprěděljajut matematične bloky. Kontekst, vključno s polarnymi procesami medžu redami, ustanovjaje tuto tehnično priměnjenje. Čislo redov i dolžina vsakogo reda morajut ostati eksplicitne.\par
\noindent\emph{Reason:} Redy po $n$ varijablov preserves organization into groups of equal size. In Paper 8 a form may depend on $N$ rows, each containing $n$ variables; these two counts are not interchangeable. The archive chooses red for continuity with Papers 6–7 and because an infinite analytic series is not intended. The newly consulted dictionary gives ręd with ordinary senses “row, queue” and, in another record, “order, arrangement, sequence”. These support the lexical base but do not themselves define the mathematical blocks. Context, including the polar operators between rows, establishes the technical use. Both the number of rows and the length of each must remain explicit.\par
\noindent{\footnotesize Izvorne zapisy / Original records: HIST-g057-005.}\par
\smallskip\noindent\emph{Sravnjenje s tekstom:} Vsaky iz $N$ redov imaje $n$ prěměnnyh. Slovnične zapisy 2283 i 3798 podpirajut obče smysly ręd, ale specialny raspored daje matematičny tekst.\par
\noindent\emph{Textual comparison:} Each of the $N$ rows has $n$ variables. Dictionary rows 2283 and 3798 support ordinary senses of ręd, while the mathematical text supplies the particular arrangement.\par
\noindent{\footnotesize TEXT-COMP-569; pasaže / passages: 6753, 6760, 6762, 6764, 6766; \texttt{CORPUS\_FORM\_COMPARISONS.json}.}\par
\end{minipage}\par
\par\medskip\Needspace{7\baselineskip}
\noindent\begin{minipage}{\linewidth}
\hypertarget{TERM-1506}{}
\noindent\textbf{razvitje v rěd}\quad{\footnotesize\texttt{TERM-1506}}\par
\noindent\emph{Iztočnik / Source:} Reihenentwicklung\par
\noindent\emph{Rusky / Russian:} разложение в ряд\par
\noindent\emph{Ukrajinsky / Ukrainian:} розклад у ряд\par
\smallskip\noindent\emph{Razlog:} Razvitje v rěd jest zachovana fraza za algebraično razloženje v rabotě 2. Arhivni razlog jest svez s §2 i čitlivost konstrukcije reda form; on ne dava sravnenja vsih sinonimov. V §6 fraza uvodi izračunjenje konkretnoj formuly s pomoču ranějših razloženij i rekursivnogo sistema. To ne jest tvrđenje o konvergenciji beskonečnogo analitičnogo reda. Nynějšny tekst piše razvitje v red. Slovničny zapis ręd o redah i poredku ne opravdava avtomatično ni staršu bukvu ě, ni svaki specialny smysl slova. Ti različni dokazy treba držati oddělno: historično pisanje, sovrěmenna leksična osnova i matematična priměna.\par
\noindent\emph{Reason:} Razvitje v rěd is the preserved expression for an algebraic expansion in Paper 2. Its recorded reason is connection with §2 and readability of the form-series construction, not a comparison of every synonym. In §6 it introduces the calculation of a particular formula using earlier expansions and a recurrent system. It makes no assertion about convergence of an infinite analytic series. The current text writes razvitje v red. The dictionary's ręd entries for rows and order do not automatically justify either the historical letter ě or every technical sense. Historical spelling, present lexical support, and mathematical application are distinct kinds of evidence.\par
\noindent{\footnotesize Izvorne zapisy / Original records: HIST-g007-005.}\par
\smallskip\noindent\emph{Sravnjenje s tekstom:} Tekst objasnjaje, čto ukazujut dvě linije koeficientov, i izrično započinaje račun iz poslědnjego uravnenja. Slovo konečne odnosi se k rezultatu računa.\par
\noindent\emph{Textual comparison:} The passage explains the two lines of coefficients and explicitly starts the calculation from the last equation. Final characterizes the outcome of that calculation.\par
\noindent{\footnotesize TEXT-COMP-570; pasaže / passages: 922, 974, 978, 979; \texttt{CORPUS\_FORM\_COMPARISONS.json}.}\par
\end{minipage}\par
\par\medskip\Needspace{7\baselineskip}
\noindent\begin{minipage}{\linewidth}
\hypertarget{TERM-1507}{}
\noindent\textbf{rędovo razloženje}\quad{\footnotesize\texttt{TERM-1507}}\par
\noindent\emph{Iztočnik / Source:} Reihenentwicklung\par
\noindent\emph{Rusky / Russian:} рядовое разложение\par
\noindent\emph{Ukrajinsky / Ukrainian:} рядовий розклад\par
\smallskip\noindent\emph{Razlog:} Rędovo razloženje jest drugy historičny izbor za Reihenentwicklung. Zapis za §7 raboty 4 hoče zachovati Mertensovo nazvanje razloženja, ne poistovětiti je s prostym seznamom redov prěměnnyh, i izrično dopuščaje možnost razvoj. Zapis za §9 prodolžuje toj izbor, ale ne rěšaje sravnenje. Kontekst govori o razloženju form po polarah normalnyh form s pomoču invariantnyh procesov diferenciacije. Polary i transformovane prěměnne sut česti konstrukcije, ne dokazy analitičnoj konvergencije. Nynějšny tekst upotrěbljaje razloženje v redy. Arhiv ne čuva zaključny argument, ktory by daval obču prednost tomu izrazu, rędovo razloženje ili razvoj.\par
\noindent\emph{Reason:} Rędovo razloženje is another historical choice for Reihenentwicklung. The Paper 4, §7, record seeks to retain Mertens's expansion terminology without reducing it to a list of variable rows, and explicitly leaves razvoj as a possible alternative. The §9 record continues the choice but does not settle the comparison. The context concerns expansion of forms in polars of normal forms through invariant differentiation operators. Polars and transformed variables are parts of the construction, not evidence of analytic convergence. The current text uses razloženje v redy. No final comparative argument survives establishing a general preference between that expression, rędovo razloženje, and razvoj.\par
\noindent{\footnotesize Izvorne zapisy / Original records: HIST-g036-005, HIST-g038-010.}\par
\smallskip\noindent\emph{Sravnjenje s tekstom:} Razloženje jest svezano s invariantnymi procesami diferenciacije i polarami normalnyh form. Arhivna alternativa razvoj ne imaje zachovanogo zaključnogo sravnenja.\par
\noindent\emph{Textual comparison:} The expansion is connected with invariant differentiation operators and polars of normal forms. No final comparison survives settling the archived alternative razvoj.\par
\noindent{\footnotesize TEXT-COMP-571; pasaže / passages: 4454, 4456, 4464, 4597, 4606, 4609, 4908; \texttt{CORPUS\_FORM\_COMPARISONS.json}.}\par
\end{minipage}\par
\par\medskip\Needspace{7\baselineskip}
\noindent\begin{minipage}{\linewidth}
\hypertarget{TERM-1508}{}
\noindent\textbf{razklad v rědy v teoriji form}\quad{\footnotesize\texttt{TERM-1508}}\par
\noindent\emph{Iztočnik / Source:} Reihenentwicklung der Formentheorie\par
\noindent\emph{Rusky / Russian:} разложение в ряды в теории форм\par
\noindent\emph{Ukrajinsky / Ukrainian:} розклад у ряди в теорії форм\par
\smallskip\noindent\emph{Razlog:} Razklad v rědy v teoriji form jest titulna fraza, ktora čuva disciplinarny okvir. Uvod raboty 16 razlikuje dva vida razloženja: po stepenjah determinanta i po normalnyh formah za veličiny različnyh stupnjev. Pozdnějša formula pokazuje konečno razloženje do stepena $p$. Arhiv zato izbira doslovnu konstrukciju s redami, bez dodavanja analitičnoj konvergencije. To ne znači, že red prěměnnyh, red form i razloženje po stepenjah sut jeden objekt. Nynějšny naslov piše razkladam v redy; slovnik zachovuje staršu formu i jeje razlog, ne izmysljaje novu motivaciju za vsaku razliku v bukvah. Citacije Mertensa i Fišera v korpusu ne značat, že njihove raboty byly nyně izčitane.\par
\noindent\emph{Reason:} Razklad v rědy v teoriji form is a title phrase retaining the disciplinary setting. Paper 16's introduction distinguishes expansions in powers of a determinant from expansions in normal forms for quantities of different grades. A later formula displays a finite expansion through degree $p$. The archive therefore chooses literal row/series wording without adding analytic convergence. This does not make a variable row, a series of forms, and a power expansion the same object. The current title has razkladam v redy; the dictionary preserves the earlier form and its recorded reason without inventing a new motivation for every spelling difference. Citations of Mertens and Fischer in the corpus do not mean their works were newly read.\par
\noindent{\footnotesize Izvorne zapisy / Original records: HIST-g078-001.}\par
\smallskip\noindent\emph{Sravnjenje s tekstom:} Uvod razlikuje dva vida razloženja; izpisana formula po stepenjah determinanta zaključaje se pri $p$. Ona ne dodava tvrđenje o analitičnoj konvergenciji.\par
\noindent\emph{Textual comparison:} The introduction distinguishes two expansion types; the displayed determinant-power expansion ends at $p$. It adds no claim of analytic convergence.\par
\noindent{\footnotesize TEXT-COMP-572; pasaže / passages: 10724, 10729, 10731, 10733, 10854, 10856; \texttt{CORPUS\_FORM\_COMPARISONS.json}.}\par
\end{minipage}\par
\par\medskip\Needspace{7\baselineskip}
\noindent\begin{minipage}{\linewidth}
\hypertarget{TERM-1509}{}
\noindent\textbf{razvitje v ręd po polarah}\quad{\footnotesize\texttt{TERM-1509}}\par
\noindent\emph{Iztočnik / Source:} Reihenentwicklung nach Polaren\par
\noindent\emph{Rusky / Russian:} разложение в ряд по полярам\par
\noindent\emph{Ukrajinsky / Ukrainian:} розклад у ряд за полярами\par
\smallskip\noindent\emph{Razlog:} Razvitje v ręd po polarah čuva polaru kako nazvany objekt invariantnoj teorije. Arhiv izrično odbacaje parafrazu „derivata“ ili „projekcija“: takie slova mogli by skryti, po čemu jest razloženje organizovano. Konec §1 raboty 2 predvidi konstrukciju §2 i ukazuje specialny slučaj, v ktorom ostavaje samo početny člen. Predložna čest po polarah zato nosi matematičnu informaciju, ne jest ukrasna dopolnitelna fraza. Novi slovničny zapis ręd može podpirati obči korenj, ale ne zaměnjaje definiciju polary ili objašnjenje posebnosti togo razloženja. Zachovano ręd i nynějšno red sut dokumentovane različno.\par
\noindent\emph{Reason:} Razvitje v ręd po polarah keeps the polar visible as a named object of invariant theory. The archive explicitly rejects paraphrase as “derivative” or “projection”: those words could conceal what organizes the expansion. The end of Paper 2, §1, anticipates §2's construction and identifies a special case in which only the initial term remains. The phrase po polarah therefore carries mathematical information rather than decorative elaboration. The newly consulted ręd entry may support the general lexical base, but cannot replace the definition of a polar or explanation of this particular expansion. Preserved ręd and current red are documented separately.\par
\noindent{\footnotesize Izvorne zapisy / Original records: HIST-g003-017.}\par
\smallskip\noindent\emph{Sravnjenje s tekstom:} Specialny slučaj iz §1 predhodi dvom tipam razloženja i ih kombinaciji v §2. Polara nosi oprěděljenu funkciju; parafrazy derivata ili projekcija ne sut ravnocěnny dokaz.\par
\noindent\emph{Textual comparison:} The special case in §1 precedes two expansion types and their combination in §2. A polar has a specified role; paraphrase as derivative or projection does not supply equivalent information.\par
\noindent{\footnotesize TEXT-COMP-573; pasaže / passages: 378, 395, 398, 400, 402, 418, 420, 422, 424, 437, 439, 441; \texttt{CORPUS\_FORM\_COMPARISONS.json}.}\par
\end{minipage}\par
\par\medskip\Needspace{7\baselineskip}
\noindent\begin{minipage}{\linewidth}
\hypertarget{TERM-1510}{}
\noindent\textbf{razvitja v ręd po polarah}\quad{\footnotesize\texttt{TERM-1510}}\par
\noindent\emph{Iztočnik / Source:} Reihenentwicklungen nach Polaren\par
\noindent\emph{Rusky / Russian:} разложения в ряд по полярам\par
\noindent\emph{Ukrajinsky / Ukrainian:} розклади у ряди за полярами\par
\smallskip\noindent\emph{Razlog:} Razvitja v ręd po polarah jest množna forma nazvanja v naslově §2. Arhiv čuva klasicky jezyk razloženij i prepoznavajemo polary, prodolžujuči frazu iz §1. Nynějšny tekst najprvo razlikuje dva tipa razloženij, potom jih kombinuje za formy s dvoma kontragredientnymi redami. Množina zato imaje konkretno opravdanje i ne znači, že jeden obči postupok proizvoljno razdělen jest na neodvisne metody. Zachovana motivacija ne daje čisleny test razumělivosti i ne ustanovjaje prednost nad vsimi drugymi slovami za razloženje. Organizacionna rolja polarov i razlika medžu koeficientami razloženja i redami prěměnnyh sut suščestvene.\par
\noindent\emph{Reason:} Razvitja v ręd po polarah is the plural heading of §2. The archive retains classical expansion language and recognizable polars, continuing the expression from §1. The current text first distinguishes two expansion types and then combines them for forms with two contragredient rows. The plural thus has a specific basis; it does not mean that one general procedure has arbitrarily been divided into independent methods. The preserved rationale provides no numerical intelligibility test or preference over every competing expansion term. The organizing role of polars and the distinction between expansion coefficients and variable rows remain essential.\par
\noindent{\footnotesize Izvorne zapisy / Original records: HIST-g004-001.}\par
\smallskip\noindent\emph{Sravnjenje s tekstom:} Specialny slučaj iz §1 predhodi dvom tipam razloženja i ih kombinaciji v §2. Polara nosi oprěděljenu funkciju; parafrazy derivata ili projekcija ne sut ravnocěnny dokaz.\par
\noindent\emph{Textual comparison:} The special case in §1 precedes two expansion types and their combination in §2. A polar has a specified role; paraphrase as derivative or projection does not supply equivalent information.\par
\noindent{\footnotesize TEXT-COMP-573; pasaže / passages: 378, 395, 398, 400, 402, 418, 420, 422, 424, 437, 439, 441; \texttt{CORPUS\_FORM\_COMPARISONS.json}.}\par
\end{minipage}\par
\par\medskip\Needspace{7\baselineskip}
\noindent\begin{minipage}{\linewidth}
\hypertarget{TERM-1511}{}
\noindent\textbf{čislo rędov}\quad{\footnotesize\texttt{TERM-1511}}\par
\noindent\emph{Iztočnik / Source:} Reihenzahl\par
\noindent\emph{Rusky / Russian:} число строк\par
\noindent\emph{Ukrajinsky / Ukrainian:} кількість рядків\par
\smallskip\noindent\emph{Razlog:} Čislo rędov znači počet redov matric, ne čislo členov analitičnogo reda i ne čislo različnih razloženij. V §2 raboty 4 matričny produkt izpolzovany jest v historičnom smyslu sumy produktov odgovarjajučih determinantov dvoh matric s jednakym čislom redov. Predzadanno čislo redov i ograničenja na njega sut zato česti tvrđenja. Arhiv upozornjaje, že ręd ostavaje otvorjeny jezyčnomu pregledanju; to ne jest potvrdženje, že pregled byl izvedeny. Nova slovnična konsultacija podpira obči smysl „row“, ne nužno cělu specialnu frazu. Nynějšny tekst imaje čislo redov, čto zachovujemo oddělno od staršego pisanja.\par
\noindent\emph{Reason:} Čislo rędov means the number of matrix rows, not the number of terms in an analytic series or the number of distinct expansions. Paper 4, §2, uses matrix product in the historical sense of a sum of products of corresponding determinants of two matrices with equal row counts. The prescribed row count and its bounds are therefore part of the assertion. The archive leaves ręd open to linguistic review; that is not evidence that review occurred. Fresh dictionary consultation supports the ordinary row sense, not necessarily the complete technical phrase. The current text has čislo redov, preserved distinctly from the older spelling.\par
\noindent{\footnotesize Izvorne zapisy / Original records: HIST-g031-002.}\par
\smallskip\noindent\emph{Sravnjenje s tekstom:} Čislo redov ograničaje tvrđenje o sumě produktov odgovarjajučih determinantov. Obči slovničny smysl ręd ne zaměnjaje ovu historičnu definiciju.\par
\noindent\emph{Textual comparison:} The row count bounds the assertion about a sum of products of corresponding determinants. The ordinary dictionary sense of ręd does not replace this historical definition.\par
\noindent{\footnotesize TEXT-COMP-574; pasaže / passages: 3782, 3784, 3803, 3805; \texttt{CORPUS\_FORM\_COMPARISONS.json}.}\par
\end{minipage}\par
\par\medskip\Needspace{7\baselineskip}
\noindent\begin{minipage}{\linewidth}
\hypertarget{TERM-1512}{}
\noindent\textbf{čisto imaginarny}\quad{\footnotesize\texttt{TERM-1512}}\par
\noindent\emph{Iztočnik / Source:} rein imaginär\par
\noindent\emph{Rusky / Russian:} чисто мнимый\par
\noindent\emph{Ukrajinsky / Ukrainian:} чисто уявний\par
\smallskip\noindent\emph{Razlog:} Čisto imaginarny čuva oba česti matematičnogo nazvanja: čisto znači izčezanje realnoj česti, a ne samo to, že čislo ne jest realno. V dokazu raboty 10 slučaj $X-x=0$ vodi od čisto imaginarnogo argumenta k čisto imaginarnomu obrazu. Fraza zato ne može byti skračena do občeho imaginarny bez izgubljenja uslovja. Arhiv dava prednost internacionalnomu korenju, dokolě ruska i ukrajinska para upotrěbljajut nacionalne nazvanja. To jest zachovana ocěna jasnosti, ne izměreny sravnitelny rezultat. Pri sravnjenju formula $f(iy)$ i znak $i$ morajut ostati matematičnymi podatkami, ne jezyčnymi metaforami o nečemu izmysljenom.\par
\noindent\emph{Reason:} Čisto imaginarny preserves both parts of the mathematical expression: purely means that the real part vanishes, not merely that the number is nonreal. In Paper 10's argument, the case $X-x=0$ sends a purely imaginary argument to a purely imaginary image. Shortening the phrase to imaginary could therefore lose a condition. The archive favours an international stem while the Russian and Ukrainian comparanda use national terms. This is a recorded judgment about clarity, not a measured comparative result. In the accompanying argument, $f(iy)$ and the symbol $i$ must remain mathematical data rather than metaphors for something invented.\par
\noindent{\footnotesize Izvorne zapisy / Original records: HIST-g072-008.}\par
\smallskip\noindent\emph{Sravnjenje s tekstom:} Dokaz izděljuje slučaj nulovoj realnoj česti i izpisuje podobu $f(iy)$. Uslovje čisto ne možno odstraniti.\par
\noindent\emph{Textual comparison:} The proof isolates the zero-real-part case and displays the form of $f(iy)$. The qualification purely cannot be dropped.\par
\noindent{\footnotesize TEXT-COMP-575; pasaže / passages: 8664, 8666, 8668, 8670, 8672, 8674, 8676; \texttt{CORPUS\_FORM\_COMPARISONS.json}.}\par
\end{minipage}\par
\par\medskip\Needspace{7\baselineskip}
\noindent\begin{minipage}{\linewidth}
\hypertarget{TERM-1513}{}
\noindent\textbf{rekurzivno označene bazove veličiny}\quad{\footnotesize\texttt{TERM-1513}}\par
\noindent\emph{Iztočnik / Source:} rekurrierend definierte Basisgrößen\par
\noindent\emph{Rusky / Russian:} рекурсивно определенные базовые величины\par
\noindent\emph{Ukrajinsky / Ukrainian:} рекурсивно означені базові величини\par
\smallskip\noindent\emph{Razlog:} Rekurzivno označene bazove veličiny upuča na zavisimost definicije od ranějše izbranyh veličin. V §7 raboty 9 izbor jest proveden po zadannom dobrom uporjadkovanju vsih kompleksnyh čisel. Tip 1 prověrjaje adjunkciju po odnošenju k poprědnim bazovym veličinam, a tip 3 opreděljaje te bazove veličiny črez izključenje prvyh dvoh tipov. Arhiv hoče učiniti tutu zavisimost vidimoju. Ne daje tem dokaz, že konstrukcija jest efektivny računny algoritm. To razlikovanje jest osoblivo važno pri budučoj priměně v OLP. Nynějšno rekursivno definiranym jest jasnějše svezano s definicijeju nego historično označene; oba zapisa ostavajut dokumentovane, bez pripisovanja pozdnějšej formulacije staršej motivaciji.\par
\noindent\emph{Reason:} Rekurzivno označene bazove veličiny points to the dependence of the definition on previously selected quantities. Paper 9, §7, proceeds along a given well-ordering of all complex numbers. Type 1 tests adjunction relative to preceding basis quantities, while type 3 defines those basis quantities by excluding the first two types. The archive seeks to make that dependence visible. It does not thereby establish an effective computational algorithm. This distinction is particularly important for later use in OLP. Current rekursivno definiranym expresses the definitional role more explicitly than historical označene; both forms remain documented without attributing the later wording to the earlier rationale.\par
\noindent{\footnotesize Izvorne zapisy / Original records: HIST-g065-008.}\par
\smallskip\noindent\emph{Sravnjenje s tekstom:} Bazove veličiny oprěděljajut se s pomoču poprědnih veličin pri zadannom dobrom uporjadkovanju. To ne jest samo po sobě dokaz efektivnoj vyčislimosti.\par
\noindent\emph{Textual comparison:} The basis quantities are defined using preceding quantities in a given well-ordering. This alone does not establish effective computability.\par
\noindent{\footnotesize TEXT-COMP-576; pasaže / passages: 7839, 7840, 7841, 7842, 7843, 7854, 7855, 7856, 7857, 7858, 7859, 7860, 7861, 7862, 7863, 7864, 7865, 7866, 7867, 7868, 7869, 7870, 7871, 7872, 7873, 7874, 7875; \texttt{CORPUS\_FORM\_COMPARISONS.json}.}\par
\end{minipage}\par
\par\medskip\Needspace{7\baselineskip}
\noindent\begin{minipage}{\linewidth}
\hypertarget{TERM-1514}{}
\noindent\textbf{konečne vrědnosti najdene rekursivno}\quad{\footnotesize\texttt{TERM-1514}}\par
\noindent\emph{Iztočnik / Source:} rekurrierend gefundenen Schlußwerte\par
\noindent\emph{Rusky / Russian:} конечные значения, найденные рекурсивно\par
\noindent\emph{Ukrajinsky / Ukrainian:} кінцеві значення, знайдені рекурсивно\par
\smallskip\noindent\emph{Razlog:} Konečne vrědnosti najdene rekursivno govori o zaključnyh rezultatah računa, ne o konečnosti čisel i ne o limitě beskonečnogo procesa. Pasaža §6 raboty 2 sama objasnjaje poredok: prva linija daje podstavjene cěnnosti, druga daje rezultaty, izračunane počinajuči s poslědnjim uravnenjem sistema. Arhiv čuva tuto proceduru i internacionalny korenj rekursiv-. Nynějšny tekst imaje konečne cěnnosti; ta leksična razlika ne zaměnjaje uslovje o směru računa. Poznanje zaključnyh formul ne opravdava izpuščenje razsudka, ako cělj jest objasniti, kako ony sut dobyte. To jest konkretna rekursija, ne obča deklaracija o algoritmičnoj složenosti.\par
\noindent\emph{Reason:} Konečne vrědnosti najdene rekursivno refers to the final results of a calculation, not finite-valued numbers or a limit of an infinite process. Paper 2, §6, specifies the order: the first line gives substituted values, the second the results calculated starting from the system's last equation. The archive retains this procedure and the international recursive stem. The current text uses konečne cěnnosti; that lexical difference does not replace the qualification about direction of calculation. Knowing the final formulas does not justify omitting their derivation when the purpose is to explain how they were obtained. This is a specific recurrence, not a general assertion about algorithmic complexity.\par
\noindent{\footnotesize Izvorne zapisy / Original records: HIST-g007-007.}\par
\smallskip\noindent\emph{Sravnjenje s tekstom:} Tekst objasnjaje, čto ukazujut dvě linije koeficientov, i izrično započinaje račun iz poslědnjego uravnenja. Slovo konečne odnosi se k rezultatu računa.\par
\noindent\emph{Textual comparison:} The passage explains the two lines of coefficients and explicitly starts the calculation from the last equation. Final characterizes the outcome of that calculation.\par
\noindent{\footnotesize TEXT-COMP-570; pasaže / passages: 922, 974, 978, 979; \texttt{CORPUS\_FORM\_COMPARISONS.json}.}\par
\end{minipage}\par
\par\medskip\Needspace{7\baselineskip}
\noindent\begin{minipage}{\linewidth}
\hypertarget{TERM-1516}{}
\noindent\textbf{rekursivny sistem ravnanj}\quad{\footnotesize\texttt{TERM-1516}}\par
\noindent\emph{Iztočnik / Source:} rekurrierendes Gleichungssystem\par
\noindent\emph{Rusky / Russian:} рекурсивная система уравнений\par
\noindent\emph{Ukrajinsky / Ukrainian:} рекурсивна система рівнянь\par
\smallskip\noindent\emph{Razlog:} Rekursivny sistem ravnanj imenuje sistem, ktorogo odnosy dopuščajut poslědovateljne izračunjenje iskanyh koeficientov. V §2 raboty 2 želano razloženje jest dobyto obratom nastavajučego sistema; tekst potom govori o jednoznačnom rekursivnom izračunjenju čislovyh koeficientov. Arhiv ocěnjuje sovrěmenny korenj rekursiv- kako čitlivy i doslědny s uvodom. To ne znači, že obyčna opakovljivost uravnenja dostatočna jest za nazvanje rekursivny. Treba pokazati zavisimost, početne podatky i směr vyvoda v konkretnoj konstrukciji. Starše ravnanj i nynějšno ravnanij sut zachovane formy, ne dokaz dvaju različnih matematičnyh pojmov.\par
\noindent\emph{Reason:} Rekursivny sistem ravnanj names a system whose relations permit successive calculation of the required coefficients. In Paper 2, §2, the desired expansion is obtained by reversing the resulting system; the text then describes unique recursive calculation of numerical coefficients. The archive regards the modern recursive stem as readable and consistent with the introduction. This does not mean that mere repetition of an equation makes it recursive. The dependency, initial data, and direction of derivation must be shown in the particular construction. Historical ravnanj and current ravnanij are preserved forms, not evidence of two different mathematical concepts.\par
\noindent{\footnotesize Izvorne zapisy / Original records: HIST-g004-010.}\par
\smallskip\noindent\emph{Sravnjenje s tekstom:} Uvodna metodna fraza imaje konkretno prodolženje v sistemu za koeficienty. Jeje dokumentovanje ne trěbuje pripisati historičnomu tekstu sovrěmennu teoriju algoritmov.\par
\noindent\emph{Textual comparison:} The introductory methodological phrase has a concrete continuation in the coefficient system. Documenting it does not require attributing a modern theory of algorithms to the historical text.\par
\noindent{\footnotesize TEXT-COMP-577; pasaže / passages: 180, 437, 439, 441, 515, 518, 519, 520, 521, 523; \texttt{CORPUS\_FORM\_COMPARISONS.json}.}\par
\end{minipage}\par
\par\medskip\Needspace{7\baselineskip}
\noindent\begin{minipage}{\linewidth}
\hypertarget{TERM-1517}{}
\noindent\textbf{relacija}\quad{\footnotesize\texttt{TERM-1517}}\par
\noindent\emph{Iztočnik / Source:} Relation\par
\noindent\emph{Rusky / Russian:} соотношение\par
\noindent\emph{Ukrajinsky / Ukrainian:} співвідношення\par
\smallskip\noindent\emph{Razlog:} Relacija jest kratko internacionalno nazvanje, ktoro uvodi matematičny odnos v zapisě raboty 2, §16. Arhiv sravnjuje je s ruskim соотношение i ukrajinskim співвідношення, davajuči prednost kratkosti prěd formulom. Pri aktualnom pojavjenju relacija jest kongruencija po zadannom modulu. Zato slovo ne ograničava se na strogu jednakost i samo ne označuje binarnu relaciju kako množstvo uporjadkovanyh parov. To drugo značenje mora byti objasnjeno oddělno pri OLP. Pasaža ima takože ranějše registrovano povtorjenje nazvanja formy; terminologično izjasnjenje ne jest potvrdženje točnosti vseh znakov toj pasaži.\par
\noindent\emph{Reason:} Relacija is the compact international term introducing a mathematical relation in the Paper 2, §16, record. The archive compares it with Russian соотношение and Ukrainian співвідношення, preferring compact wording before a formula. In the current occurrence the relation is a congruence modulo a specified modulus. The term is therefore not restricted to strict equality, nor does it by itself designate a binary relation as a set of ordered pairs. That other sense requires separate explanation in OLP. The passage also contains the previously recorded repeated-form question; clarifying its terminology does not certify every symbol in the passage.\par
\noindent{\footnotesize Izvorne zapisy / Original records: HIST-g017-007.}\par
\smallskip\noindent\emph{Sravnjenje s tekstom:} Priměr relacije jest kongruencija, ne prosta jednakost. Ranějše registrovano povtorjenje nazvanja formy ostavaje otvorjeno iztočnikovomu sravnjenju.\par
\noindent\emph{Textual comparison:} The example of a relation is a congruence, not simple equality. The previously recorded repeated-form question remains open to source comparison.\par
\noindent{\footnotesize TEXT-COMP-578; pasaže / passages: 2351, 2353, 2355; \texttt{CORPUS\_FORM\_COMPARISONS.json}.}\par
\end{minipage}\par
\par\medskip\Needspace{7\baselineskip}
\noindent\begin{minipage}{\linewidth}
\hypertarget{TERM-1518}{}
\noindent\textbf{relacija medžu razkladnymi formami}\quad{\footnotesize\texttt{TERM-1518}}\par
\noindent\emph{Iztočnik / Source:} Relation zwischen zerfallenden Formen\par
\noindent\emph{Rusky / Russian:} соотношение между разлагающимися формами\par
\noindent\emph{Ukrajinsky / Ukrainian:} співвідношення між формами, що розкладаються\par
\smallskip\noindent\emph{Razlog:} Relacija medžu razkladnymi formami upotrěbljena jest v rabotě 2 pri produktoh s funkcionalnym determinantom. Formula izpisuje jedan produkt kak razliku dvoh produktov; kompozicionalnost form jest zato vidima, ne samo oběčana slovom razkladny. Arhiv prodolžuje ranějšnu familiju razkladny / razkladajut se i ostavjaje jeje jezyčnu ocěnu otvorjenu. On ne dava osnovy, da bi izraz poistovětiti s idealnym razloženjem ili s reducibilnostju prědstavjenja. Pri čitanju frazy treba čuvati smysl forma iz invariantnoj teorije i podane uslovja na množitelj. Sravnitelne ruske i ukrajinske frazy sut zachovane dokazy izbora, ne rezultaty novogo testa čitatelskoj razumělivosti.\par
\noindent\emph{Reason:} Relacija medžu razkladnymi formami occurs in Paper 2 for products involving a functional determinant. The displayed formula writes one product as a difference of two products, so the forms' compositional structure is shown rather than merely asserted by razkladny. The archive continues the earlier razkladny / razkladajut se family and leaves its linguistic assessment open. It gives no basis for equating this with ideal decomposition or reducibility of a representation. Reading the phrase requires the invariant-theoretic meaning of form and the stated qualifications on the factor. Its Russian and Ukrainian comparanda are preserved evidence of the choice, not results of a new reader-comprehension test.\par
\noindent{\footnotesize Izvorne zapisy / Original records: HIST-g019-015.}\par
\smallskip\noindent\emph{Sravnjenje s tekstom:} Izpisana identičnost pokazuje produkt funkcionalnogo determinanta i razliku dvoh produktov; treba zachovati uslovja na množitelj.\par
\noindent\emph{Textual comparison:} The displayed identity shows a product involving a functional determinant and a difference of two products; the qualifications on the factor must be retained.\par
\noindent{\footnotesize TEXT-COMP-579; pasaže / passages: 2625, 2627; \texttt{CORPUS\_FORM\_COMPARISONS.json}.}\par
\end{minipage}\par
\par\medskip\Needspace{7\baselineskip}
\noindent\begin{minipage}{\linewidth}
\hypertarget{TERM-1519}{}
\noindent\textbf{relacije}\quad{\footnotesize\texttt{TERM-1519}}\par
\noindent\emph{Iztočnik / Source:} Relationen\par
\noindent\emph{Rusky / Russian:} соотношения\par
\noindent\emph{Ukrajinsky / Ukrainian:} співвідношення\par
\smallskip\noindent\emph{Razlog:} Relacije v rabotě 6 jesut smyslovo jače nego neoprěděljene „svezi“ medžu funkcijami. V §3 formula vyražaje identično izčezanje polinomnogo izraza od funkcij; princip prěnosa čuva takove relacije pri vzajemno jednoznačnom odobraženju i obratno. Arhiv hoče dovoljno široko nazvanje za uravnenja i identičnosti, bez izpuščenja ih logičnoj roli. Čislo relacij može byti konečno ili beskonečno, ale vsaka navedena formula imaje svoje argumenty i oblast identičnosti. Zato širina slova ne znači dozvoljenje za volnu parafrazu matematičnogo uslovja. To takože ne ustanovjaje terminologiju svih relacij v logikě bez prověrky konkretnogo značenja.\par
\noindent\emph{Reason:} Relacije in Paper 6 means more than unspecified “connections” between functions. In §3 the formula asserts identical vanishing of a polynomial expression in functions; the transfer principle preserves such relations under a one-to-one correspondence and conversely. The archive seeks a term broad enough for equations and identities without losing their logical role. There may be finitely or infinitely many relations, but each displayed formula has its arguments and domain of identity. Breadth of vocabulary is not permission to paraphrase away the mathematical condition. It also does not settle terminology for every kind of relation in logic without checking the sense.\par
\noindent{\footnotesize Izvorne zapisy / Original records: HIST-g043-007.}\par
\smallskip\noindent\emph{Sravnjenje s tekstom:} Prěnos i obratno tvrđenje čuvajut polinomne relacije. Široko nazvanje ne uklanjaje točne argumenty ni oblast, v ktoroj izraz identično izčezaje.\par
\noindent\emph{Textual comparison:} Transfer and the converse assertion preserve polynomial relations. The broad term removes neither the exact arguments nor the domain on which the expression vanishes identically.\par
\noindent{\footnotesize TEXT-COMP-580; pasaže / passages: 5345, 5347, 5348, 5350, 5352, 5353, 5354, 5374, 5376, 5378; \texttt{CORPUS\_FORM\_COMPARISONS.json}.}\par
\end{minipage}\par
\par\medskip\Needspace{7\baselineskip}
\noindent\begin{minipage}{\linewidth}
\hypertarget{TERM-1520}{}
\noindent\textbf{relacije za razložime formy}\quad{\footnotesize\texttt{TERM-1520}}\par
\noindent\emph{Iztočnik / Source:} Relationen für die zerfallenden Formen\par
\noindent\emph{Rusky / Russian:} соотношения для распадающихся форм\par
\noindent\emph{Ukrajinsky / Ukrainian:} співвідношення для форм, що розкладаються на множники\par
\smallskip\noindent\emph{Razlog:} Relacije za razložime formy označajut izključenu grupu slučajev v rabotě 2, §5: obče tvrđenje o vyvodženju pozdnějših redukcijnyh formul polarizacijeju izrično dela ovo izključenje. Zato fraza ne može byti izpuščena kak periferny primer. Arhiv upominaje, že ukrajinsky prěvod razkriva razloženje na množitelje, rusky prodolžuje ranějšnu formulaciju, a medžuslovjansky ponovjaje razložime formy. Ta ponovjena upotrěba ne znači, že termin jest nova tvorba toga prěvoda ili že jeho jezyčny pregled byl završen. Sravnjenje s razkladnymi formami v §18 čuva oba zapisa i izjasnjaje matematiku, ale ne dodava nezachovano glasovanje o prednosti jednoj formy.\par
\noindent\emph{Reason:} Relacije za razložime formy identifies an expressly excepted class in Paper 2, §5: the general assertion that later reduction formulas follow by polarization explicitly makes this exception. The phrase cannot be omitted as a peripheral example. The archive notes that Ukrainian makes factorization explicit, Russian continues earlier wording, and Interslavic repeats razložime formy. Repetition establishes neither that the translation coined the term nor that its linguistic review was completed. Comparison with razkladnymi formami in §18 preserves both forms and clarifies the mathematics, but adds no unrecorded vote for one adjective over the other.\par
\noindent{\footnotesize Izvorne zapisy / Original records: HIST-g006-011.}\par
\smallskip\noindent\emph{Sravnjenje s tekstom:} Tvrđenje o polarizaciji izrično izključaje relacije za razložime formy. Fraza ograničaje obseg tvrđenja, ne dava samo dopolnitelny primer.\par
\noindent\emph{Textual comparison:} The polarization assertion expressly excepts relations for decomposable forms. The phrase limits the assertion's scope rather than merely adding an example.\par
\noindent{\footnotesize TEXT-COMP-581; pasaže / passages: 811, 813, 816; \texttt{CORPUS\_FORM\_COMPARISONS.json}.}\par
\end{minipage}\par
\par\medskip\Needspace{7\baselineskip}
\noindent\begin{minipage}{\linewidth}
\hypertarget{TERM-1521}{}
\noindent\textbf{relativno cěle funkcije}\quad{\footnotesize\texttt{TERM-1521}}\par
\noindent\emph{Iztočnik / Source:} relativ ganze Funktionen\par
\noindent\emph{Rusky / Russian:} относительно целые функции\par
\noindent\emph{Ukrajinsky / Ukrainian:} відносно цілі функції\par
\smallskip\noindent\emph{Razlog:} Relativno cěle funkcije zachovuje historično nazvanje Hilbertovoj problemy 14. V rabotě 5 ide o vse polinomy, sodrživane v danom polju racionalnyh funkcij, i o možnost vyraziti jih črez konečno mnogo takyh polinomov. Cěly tu ne znači ni cěločisleny koeficient, ni moralnu vlastnost, ni obču analitičnu funkciju bez singularnostij. Arhiv izbira doslovnu frazu radi svezi s problemom, ne zato, že vsaky sovrěmenny smysl integralnosti jest zaměnljiv. Izrično podane uslovje na polinomy čuva smysl bolje nego sama internacionalna čest relativno. Citacija Hilberta ostavaje citacijeju iz korpusa; jeho izvorni tekst ne byl v tutom sravnjenju nanovo izčitany.\par
\noindent\emph{Reason:} Relativno cěle funkcije preserves the historical name associated with Hilbert's fourteenth problem. Paper 5 considers all polynomials contained in a given field of rational functions and whether finitely many of those polynomials suffice to express them. Here whole or integral does not mean integer coefficients, a moral quality, or a general analytic function without singularities. The archive chooses literal wording to retain the connection with the problem, not because every modern sense of integrality is interchangeable. The explicit polynomial condition carries more meaning than relativno alone. Hilbert remains a citation within the corpus; his original text was not newly consulted for this comparison.\par
\noindent{\footnotesize Izvorne zapisy / Original records: HIST-g039-017.}\par
\smallskip\noindent\emph{Sravnjenje s tekstom:} Integralna baza jest tu konečna polinomna sistema poroditeljev. Historično nazvanje relativno cěle funkcije mora byti čitano v tom okviru.\par
\noindent\emph{Textual comparison:} An integral basis here is a finite polynomial generating system. The historical label relatively integral functions must be read in that setting.\par
\noindent{\footnotesize TEXT-COMP-582; pasaže / passages: 4950, 4954, 4956, 4966, 4968, 4973; \texttt{CORPUS\_FORM\_COMPARISONS.json}.}\par
\end{minipage}\par
\par\medskip\Needspace{7\baselineskip}
\noindent\begin{minipage}{\linewidth}
\hypertarget{TERM-1522}{}
\noindent\textbf{vzaimno proste v polinomnom smyslu}\quad{\footnotesize\texttt{TERM-1522}}\par
\noindent\emph{Iztočnik / Source:} relativ prim im Polynomsinn\par
\noindent\emph{Rusky / Russian:} взаимно простые в полиномиальном смысле\par
\noindent\emph{Ukrajinsky / Ukrainian:} взаємно прості у поліноміальному сенсі\par
\smallskip\noindent\emph{Razlog:} Vzaimno proste v polinomnom smyslu čuva ograničenje, ktoro lahko izgubiti pri skračenju nazvanja. Teorem X raboty 22 najprvo razklada polinom na dva vzajemno proste faktory, a potom konstruuje idealne dělitelje s oprěděljenymi normami. Polinomnogo uslovja ne treba zaměnjati neizrečennym tvrđenjem, že dva proizvoljne idealy dajut jediničny ideal. Arhiv izrično navodi tu blizost polinomnyh i idealnyh pojmov kako razlog za polnu frazu. Starše vzaimno i nynějšno vzajemno ostavajut različne zapisy istogo naměra; dva historične zapisy ne dokazujut dva nezavisne sravnenja. Smysl pridavnogo prosty ustanovjaje njegov odnos k drugomu polinomu, ne tvrdženje o nerazložimosti vsakogo faktora.\par
\noindent\emph{Reason:} Vzaimno proste v polinomnom smyslu retains a qualification easily lost by shortening the term. Paper 22's Theorem X first factors a polynomial into two relatively prime factors and then constructs ideal divisors with specified norms. The polynomial condition must not be replaced by an unstated assertion that two arbitrary ideals generate the unit ideal. The archive expressly gives this proximity of polynomial and ideal concepts as its reason for retaining the full phrase. Historical vzaimno and current vzajemno are different spellings of the intended relation; two records do not establish two independent comparisons. Prime here describes the relation between the polynomials, not irreducibility of each factor.\par
\noindent{\footnotesize Izvorne zapisy / Original records: HIST-g089-010, HIST-g092-052.}\par
\smallskip\noindent\emph{Sravnjenje s tekstom:} Teorem razlišuje razloženje polinoma i poslědujuču konstrukciju dvoh idealnyh děliteljev. Dodatek v polinomnom smyslu čuva tu granicu.\par
\noindent\emph{Textual comparison:} The theorem distinguishes polynomial factorization from the subsequent construction of two ideal divisors. The qualifier in the polynomial sense preserves that boundary.\par
\noindent{\footnotesize TEXT-COMP-583; pasaže / passages: 14910, 14913, 14915, 14918, 14920, 14922, 14925; \texttt{CORPUS\_FORM\_COMPARISONS.json}.}\par
\end{minipage}\par
\par\medskip\Needspace{7\baselineskip}
\noindent\begin{minipage}{\linewidth}
\hypertarget{TERM-1523}{}
\noindent\textbf{relativno polny sistem}\quad{\footnotesize\texttt{TERM-1523}}\par
\noindent\emph{Iztočnik / Source:} relativ vollständige System\par
\noindent\emph{Rusky / Russian:} относительно полная система\par
\noindent\emph{Ukrajinsky / Ukrainian:} відносно повна система\par
\smallskip\noindent\emph{Razlog:} Relativno polny sistem v zaključku raboty 2 vrača čitatelja k definiciji §3, a ne proglašaje polnost bez uslovij. Kontrakcije produktov sistemnyh form vyražajut se črez formy sistema do dodatnyh izrazov, nastavajučih s pomoču sistema modula. §26 sobira 331 form i podsistemov; ta suma sama ne jest definicija polnosti niti dokaz polnoty prěvoda. Arhiv predpočitaje doslovno i čitljivo nazvanje relativnoj polnosti. Pojem jest iz invariantnoj teorije. Pri OLP isto pridavno polny može nositi drugu definiciju, napriměr odnos medžu dokazivostju i istinnostju; tu definiciju ne možno prěnesti iz Noether samo po podobnosti slova.\par
\noindent\emph{Reason:} Relativno polny sistem in Paper 2's conclusion refers back to §3's definition rather than asserting unqualified completeness. Contractions of products of system forms are expressible through the system forms up to additional expressions involving the modulus system. Section 26 collects 331 forms and subsystems; that count is neither the definition of completeness nor evidence that the translation is complete. The archive favours a literal, readable expression for relative completeness. This is an invariant-theoretic concept. In OLP the same adjective may have another definition, such as a relation between provability and truth; that definition cannot be imported from Noether by word matching.\par
\noindent{\footnotesize Izvorne zapisy / Original records: HIST-g027-010.}\par
\smallskip\noindent\emph{Sravnjenje s tekstom:} Uvod upotrěbljaje odnosno polny, a §3 relativno polny. Definicija sodrživaje dodatne izrazy s pomoču modula; zaključna suma 331 form ne zaměnjaje to uslovje.\par
\noindent\emph{Textual comparison:} The introduction uses odnosno polny and §3 relativno polny. The definition includes additional expressions involving the modulus; the concluding count of 331 forms does not replace that condition.\par
\noindent{\footnotesize TEXT-COMP-584; pasaže / passages: 69, 71, 178, 182, 196, 545, 3503, 3505, 3517, 3518; \texttt{CORPUS\_FORM\_COMPARISONS.json}.}\par
\end{minipage}\par
\par\medskip\Needspace{7\baselineskip}
\noindent\begin{minipage}{\linewidth}
\hypertarget{TERM-1524}{}
\noindent\textbf{odnosno polny sistem}\quad{\footnotesize\texttt{TERM-1524}}\par
\noindent\emph{Iztočnik / Source:} relativ vollständiges System\par
\noindent\emph{Rusky / Russian:} относительно полная система\par
\noindent\emph{Ukrajinsky / Ukrainian:} відносно повна система\par
\smallskip\noindent\emph{Razlog:} Odnosno polny sistem jest ranějša fraza iz kratkoj raboty 1 i uvoda raboty 2. Arhivni razlog jest jeje status kak nazvanogo tehničnogo objekta: ne treba měnjati nazvanje radi stilističnogo raznoobrazja. Vvodny tekst razlikuje relativnu polnost od absolutnoj i opisuje uslovja prěhoda medžu nimi. Vnutri raboty 2 definicija i pozdnějše razděly upotrěbljajut relativno polny. Slovnik zachovuje oba izbora i točnu svez medžu nimi, ne izmysljaje matematičnu razliku radi raznogo priloga. Arhivno tvrđenje o razumljivosti vo vsih trěh jezykah jest ocěna autora zapisa, ne rezultat izčitannogo čitatelskogo testa.\par
\noindent\emph{Reason:} Odnosno polny sistem is the earlier expression in the short Paper 1 and Paper 2's introduction. The recorded reason is its status as a named technical object: the label should not be varied merely for stylistic diversity. The introduction distinguishes relative from absolute completeness and describes the conditions for passing between them. Paper 2's definition and later sections use relativno polny. The dictionary preserves both choices and their connection rather than inventing a mathematical distinction from the different adverb. The archive's claim of intelligibility in all three languages is the recorder's assessment, not a documented reader-comprehension test.\par
\noindent{\footnotesize Izvorne zapisy / Original records: HIST-g001-004, HIST-g002-003.}\par
\smallskip\noindent\emph{Sravnjenje s tekstom:} Uvod upotrěbljaje odnosno polny, a §3 relativno polny. Definicija sodrživaje dodatne izrazy s pomoču modula; zaključna suma 331 form ne zaměnjaje to uslovje.\par
\noindent\emph{Textual comparison:} The introduction uses odnosno polny and §3 relativno polny. The definition includes additional expressions involving the modulus; the concluding count of 331 forms does not replace that condition.\par
\noindent{\footnotesize TEXT-COMP-584; pasaže / passages: 69, 71, 178, 182, 196, 545, 3503, 3505, 3517, 3518; \texttt{CORPUS\_FORM\_COMPARISONS.json}.}\par
\end{minipage}\par
\par\medskip\Needspace{7\baselineskip}
\noindent\begin{minipage}{\linewidth}
\hypertarget{TERM-1525}{}
\noindent\textbf{relativno polny sistem mod $(\varrho,t)$}\quad{\footnotesize\texttt{TERM-1525}}\par
\noindent\emph{Iztočnik / Source:} relativ vollständiges System mod $(\varrho,t)$\par
\noindent\emph{Rusky / Russian:} относительно полная система по модулю $(\varrho,t)$\par
\noindent\emph{Ukrajinsky / Ukrainian:} відносно повна система за модулем $(\varrho,t)$\par
\smallskip\noindent\emph{Razlog:} Relativno polny sistem mod $(\varrho,t)$ imenuje relativnu polnost s dvoma točno označenymi modulami. Arhiv čuva kratko internacionalno mod pokraj formule, dokolě ruska i ukrajinska para davajut polne predložne izrazy. Tuto različno jezyčno obličje ne měnjaje matematičnogo argumenta. V §17 trěba oddělno pokazati redukciju modula i svođenje sistema od $s$; rezultat potom služi za konstrukciju absolutno polnogo sistema. To ne jest dozvoljenje izpustiti par $(\varrho,t)$ ili čitati mod kako neobveznu primětku. Zapis upominaje buduču jezyčnu ocěnu, no ne dava dokaza, že ona už byla provedena.\par
\noindent\emph{Reason:} Relativno polny sistem mod $(\varrho,t)$ names relative completeness with respect to two precisely identified moduli. The archive retains compact international mod beside a formula, whereas the Russian and Ukrainian comparanda use fuller prepositional expressions. This linguistic difference leaves the mathematical argument unchanged. Section 17 separately requires reduction of the modulus and reduction of the system of $s$; the result then serves the construction of an absolutely complete system. Neither the pair $(\varrho,t)$ nor the qualification mod is dispensable. The record anticipates linguistic review but does not show that such review occurred.\par
\noindent{\footnotesize Izvorne zapisy / Original records: HIST-g018-001.}\par
\smallskip\noindent\emph{Sravnjenje s tekstom:} Naslov čuva točny par argumentov. Razsudok potom razlikuje redukciju modula i sistema od $s$.\par
\noindent\emph{Textual comparison:} The heading retains the exact argument pair. The argument then distinguishes reduction of the modulus from reduction of the system of $s$.\par
\noindent{\footnotesize TEXT-COMP-585; pasaže / passages: 2380, 2384, 2387, 2389, 2391, 2394, 2395, 2451; \texttt{CORPUS\_FORM\_COMPARISONS.json}.}\par
\end{minipage}\par
\par\medskip\Needspace{7\baselineskip}
\noindent\begin{minipage}{\linewidth}
\hypertarget{TERM-1526}{}
\noindent\textbf{relativno polny sistem mod $(A,\nu)$}\quad{\footnotesize\texttt{TERM-1526}}\par
\noindent\emph{Iztočnik / Source:} relativ vollständiges System mod $(A,\nu)$\par
\noindent\emph{Rusky / Russian:} относительно полная система по модулю $(A,\nu)$\par
\noindent\emph{Ukrajinsky / Ukrainian:} відносно повна система за модулем $(A,\nu)$\par
\smallskip\noindent\emph{Razlog:} Relativno polny sistem mod $(A,\nu)$ prodolžuje definiciju §3 v konkretnom rezultatu §5. Nynějšna pasaža imenuje četyri formy $f,\theta,K,j$ i navodi, kak se ostale potrěbne tvorjenja vyražajut po modulu. Arhiv čuva nazvanje i algebraičnu konstrukciju mod, vměsto stylističnoj zaměny. Bukva $A$ jest v historičnom slovničnom polju prosty znak; korpus upotrěbljaje frakturnu bukvu $\mathfrak{A}$ za tu formu. Jeje tipografično čitanje ne směje skryti razliku medžu zachovanym zapisom i nynějšnim tekstom. Grupa četyreh form jest podana kako rezultat v danom kontekstu, ne kako obča definicija vsih relativno polnyh sistemov.\par
\noindent\emph{Reason:} Relativno polny sistem mod $(A,\nu)$ carries §3's definition into the specific result of §5. The current passage names four forms, $f,\theta,K,j$, and explains the relevant expressions modulo the specified forms. The archive preserves the term and algebraic mod construction rather than making a stylistic substitution. Historical dictionary $A$ is a plain letter; the corpus uses the Fraktur letter $\mathfrak{A}$ for this form. Typesetting must not conceal the difference between the preserved record and current text. The four-form collection is a result in this setting, not a definition of every relatively complete system.\par
\noindent{\footnotesize Izvorne zapisy / Original records: HIST-g006-014.}\par
\smallskip\noindent\emph{Sravnjenje s tekstom:} Četyri formy i modulus sut izrično dane; fraza prodolžuje definiciju relativnoj polnosti, ne novu obču definiciju četyreh form.\par
\noindent\emph{Textual comparison:} The four forms and modulus are explicit; the phrase applies the relative-completeness definition rather than defining an arbitrary four-form collection.\par
\noindent{\footnotesize TEXT-COMP-586; pasaže / passages: 545, 652, 773, 775; \texttt{CORPUS\_FORM\_COMPARISONS.json}.}\par
\end{minipage}\par
\par\medskip\Needspace{7\baselineskip}
\noindent\begin{minipage}{\linewidth}
\hypertarget{TERM-1527}{}
\noindent\textbf{relativno polny sistem mod $(s,\rho,t)$}\quad{\footnotesize\texttt{TERM-1527}}\par
\noindent\emph{Iztočnik / Source:} relativ vollständiges System mod $(s,\rho,t)$\par
\noindent\emph{Rusky / Russian:} относительно полная система по модулю $(s,\rho,t)$\par
\noindent\emph{Ukrajinsky / Ukrainian:} відносно повна система за модулем $(s,\rho,t)$\par
\smallskip\noindent\emph{Razlog:} Relativno polny sistem mod $(s,\rho,t)$ čuva vse tri argumenty, ktore sut potrěbne za postupnost konstrukcije. §7 najprvo tvori sistem po modulu $s$, potom dodava dvě kvadratične formy; par $(\varrho,t)$ razgleda se kak vyši modul. Arhiv motivuje izraz kontinuitetom s §3, §5 i §6. Taky kontinuitet jest svezano objašnjenje izbora, ne tri nove nezavisne jezyčne prověrky. V historičnom polju stoji znak $\rho$, a v nynějšnem tekstu njegova tipografična varianta $\varrho$; oba zapisa treba činiti čitlivymi bez tihogo prěpisanja arhiva. Izraz ne znači, že tri argumenty možno libovoljno proměniti ili izpustiti.\par
\noindent\emph{Reason:} Relativno polny sistem mod $(s,\rho,t)$ retains all three arguments required by the construction. Section 7 first forms a system modulo $s$ and then adjoins two quadratic forms; the pair $(\varrho,t)$ is regarded as a higher modulus. The archived reason is continuity with §§3, 5, and 6. That is a connected explanation of the choice, not three new independent linguistic checks. The historical field uses $\rho$ and the current text its typographical variant $\varrho$; both must be legible without silently rewriting the archive. The phrase does not make its three arguments optional.\par
\noindent{\footnotesize Izvorne zapisy / Original records: HIST-g008-003.}\par
\smallskip\noindent\emph{Sravnjenje s tekstom:} Dvě kvadratične formy pridajut se k ranějšemu modulu; ih par jest razgledan kak vyši modul.\par
\noindent\emph{Textual comparison:} Two quadratic forms are adjoined to the earlier modulus; their pair is treated as a higher modulus.\par
\noindent{\footnotesize TEXT-COMP-587; pasaže / passages: 1039, 1046, 1048, 1051, 1052, 1053, 1056, 1058, 1060, 1062, 1064; \texttt{CORPUS\_FORM\_COMPARISONS.json}.}\par
\end{minipage}\par
\par\medskip\Needspace{7\baselineskip}
\noindent\begin{minipage}{\linewidth}
\hypertarget{TERM-1528}{}
\noindent\textbf{relativno polny sistem po modulu ... / mod ...}\quad{\footnotesize\texttt{TERM-1528}}\par
\noindent\emph{Iztočnik / Source:} relativ vollständiges System modulo ...\par
\noindent\emph{Rusky / Russian:} относительно полная система по модулю ...\par
\noindent\emph{Ukrajinsky / Ukrainian:} відносно повна система за модулем ...\par
\smallskip\noindent\emph{Razlog:} Relativno polny sistem po modulu jest polna slovna konstrukcija; mod jest jeje kompaktna forma pokraj formule. Arhiv izrično dopuščaje oba vida i zato ne trěbuje od redaktora odstraniti jedno iz njih radi vněšnej jednakosti. Trěbuje se smyslova doslědnost: pri vsakoj upotrěbě mora ostati jasno, po kojem modulu i v kojem smislu sistem jest polny. Tri točky v historičnom zaglavju sut město za konkretny argument, ne povolenje izpustiti argument v matematičnom tekstu. Definicija §3 uključaje dodatny člen iz sistemov modula, čto razlikuje ovu relativnu polnost od absolutnoj.\par
\noindent\emph{Reason:} Relativno polny sistem po modulu is the full verbal construction; mod is its compact counterpart beside a formula. The archive expressly permits both and therefore does not require an editor to eliminate one for superficial uniformity. Consistency concerns meaning: each use must identify the modulus and the sense in which the system is complete. The ellipsis in the historical headword is a placeholder for a specific argument, not permission to omit that argument from mathematical prose. Section 3's definition includes an additional term from the modulus systems, distinguishing this relative completeness from absolute completeness.\par
\noindent{\footnotesize Izvorne zapisy / Original records: HIST-g005-002.}\par
\smallskip\noindent\emph{Sravnjenje s tekstom:} Uvod upotrěbljaje odnosno polny, a §3 relativno polny. Definicija sodrživaje dodatne izrazy s pomoču modula; zaključna suma 331 form ne zaměnjaje to uslovje.\par
\noindent\emph{Textual comparison:} The introduction uses odnosno polny and §3 relativno polny. The definition includes additional expressions involving the modulus; the concluding count of 331 forms does not replace that condition.\par
\noindent{\footnotesize TEXT-COMP-584; pasaže / passages: 69, 71, 178, 182, 196, 545, 3503, 3505, 3517, 3518; \texttt{CORPUS\_FORM\_COMPARISONS.json}.}\par
\end{minipage}\par
\par\medskip\Needspace{7\baselineskip}
\noindent\begin{minipage}{\linewidth}
\hypertarget{TERM-1529}{}
\noindent\textbf{relativno algebraično cěla}\quad{\footnotesize\texttt{TERM-1529}}\par
\noindent\emph{Iztočnik / Source:} relativ-algebraisch-ganz\par
\noindent\emph{Rusky / Russian:} относительно алгебраически целая\par
\noindent\emph{Ukrajinsky / Ukrainian:} відносно алгебраїчно ціла\par
\smallskip\noindent\emph{Razlog:} Relativno algebraično cěla jest ženskaja forma pridavnogo izraza, soglasovana s veličina. Definicija V raboty 6 izrično nazyvaje oblast, nad ktoroju veličina izpolnjaje uravnenje s najvyššim koeficientom $1$ i s ostalymi koeficientami iz toj oblasti. Arhiv razkriva německy složeny izraz, aby čitatelj ne izgubil relativnu osnovu uslovja. Relativno ne dodava tajnu novu vlastnost k vyrazu algebraično cěla po odnošenju k oblasti: sam tekst predstavjaje ty formulacije kak dva nazvanja. No ovu moničnu algebraičnu zavisimost ne treba směšati s ranějše definovanymi relativno cělymi oblastjami, čije uslovje govori o kvocientah, ktore sut polinomy.\par
\noindent\emph{Reason:} Relativno algebraično cěla is a feminine adjectival form agreeing with veličina. Paper 6's Definition V specifies the domain over which the quantity satisfies a monic equation with its remaining coefficients in that domain. The archive unpacks the German compound so that the relative base of the condition stays visible. Relativno does not add an undisclosed property to algebraically integral with respect to the domain: the text itself presents these as two names. But this monic algebraic dependence must not be conflated with the earlier definition of relatively integral domains, whose condition concerns quotients that are polynomials.\par
\noindent{\footnotesize Izvorne zapisy / Original records: HIST-g049-010.}\par
\smallskip\noindent\emph{Sravnjenje s tekstom:} Definicija V trěbuje monično uravnenje s koeficientami iz označenej oblasti i dava dva nazvanja toga uslovja.\par
\noindent\emph{Textual comparison:} Definition V requires a monic equation with coefficients in the specified domain and gives two names for that condition.\par
\noindent{\footnotesize TEXT-COMP-588; pasaže / passages: 6004, 6006, 6008, 6010, 6012, 6014, 6016; \texttt{CORPUS\_FORM\_COMPARISONS.json}.}\par
\end{minipage}\par
\par\medskip\Needspace{7\baselineskip}
\noindent\begin{minipage}{\linewidth}
\hypertarget{TERM-1530}{}
\noindent\textbf{relativno-Galoisove ?islove polja [Čestično nezachovane bukvy / Partly unrecoverable lettering]}\quad{\footnotesize\texttt{TERM-1530}}\par
\noindent\emph{Iztočnik / Source:} relativ-galoissche Zahlk?rper [Čestično nezachovane bukvy / Partly unrecoverable lettering]\par
\noindent\emph{Rusky / Russian:} [Bukvy ne sut zachovane / Lettering not recoverable]\par
\noindent\emph{Ukrajinsky / Ukrainian:} [Bukvy ne sut zachovane / Lettering not recoverable]\par
\smallskip\noindent\emph{Razlog:} Relativno-Galoisove čislove polja jest čitljiva nynějšna forma nazvanja, ale něktore bukvy historičnogo zapisa ne sut zachovane. Arhiv dava konkretno smyslovo opravdanje za polje: čislove polja sut komutativne, hoti dokaz upotrěbljaje nekomutativne algebry. §1 raboty 41 potom izrično razširjaje okvir do libovoljnogo komutativnogo polja $k$ i separabilnogo Galoisovogo razširjenja $K/k$. Nazvanje ne govori, že samo polje jest nekomutativno. Poškodžene ruske i ukrajinske bukvy ne možno rekonstruovati samo iz toga, čto model znajet očekivany termin; oni sut označene kako nezachovane, a nynějšna forma čislove jest citovana iz korpusa.\par
\noindent\emph{Reason:} Relativno-Galoisove čislove polja is the readable current form, but some letters in the historical record are lost. The archive gives a specific semantic reason for polje: the number fields are commutative even though the proof uses noncommutative algebras. Paper 41, §1, then explicitly broadens the setting to an arbitrary commutative field $k$ and a separable Galois extension $K/k$. The label does not say that the field itself is noncommutative. Lost Russian and Ukrainian lettering cannot be reconstructed merely because a model knows an expected term; it is marked unrecoverable, while current čislove is quoted from the corpus.\par
\noindent{\footnotesize Izvorne zapisy / Original records: HIST-g211-002.}\par
\smallskip\noindent\emph{Sravnjenje s tekstom:} Naslov imenuje čislove polja; §1 izrično dopuščaje libovoljno komutativno osnovno polje. Skrěženy produkt jest algebroju nad tim poljem, ne dokaz nekomutativnosti polja.\par
\noindent\emph{Textual comparison:} The title names number fields; §1 expressly permits an arbitrary commutative base field. The crossed product is an algebra over it, not evidence that the field is noncommutative.\par
\noindent{\footnotesize TEXT-COMP-589; pasaže / passages: 24184, 24190, 24192, 24198, 24200, 24202, 24204; \texttt{CORPUS\_FORM\_COMPARISONS.json}.}\par
\end{minipage}\par
\par\medskip\Needspace{7\baselineskip}
\noindent\begin{minipage}{\linewidth}
\hypertarget{TERM-1531}{}
\noindent\textbf{relativno cěle oblasti prvogo roda}\quad{\footnotesize\texttt{TERM-1531}}\par
\noindent\emph{Iztočnik / Source:} relativ-ganze Bereiche erster Art\par
\noindent\emph{Rusky / Russian:} относительно целые области первого рода\par
\noindent\emph{Ukrajinsky / Ukrainian:} відносно цілі області першого роду\par
\smallskip\noindent\emph{Razlog:} Relativno cěle oblasti prvogo roda imenuje klasu iz Definicije VII raboty 6. Ona trěbuje oblast odobraženja, nad ktoroju vsaky polinom od navedenyh prěměnnyh jest algebraično cěly. Konečnost i involucijska integralna baza sut daljši teorem, ne zaměna definicije. Arhiv izbira rod kako klasifikacijno slovo i odbacaje slabějše čitanje s tip; to jest zachovana terminologična preferencija. Nynějšny uvod ipak ima prvogo tipa, a §10 prvogo roda, takže doslědnost ješče trěbuje sravnjenje kontekstov. Prvy rod ne znači prvy poredok v logikě i ne nese vsak drugy matematičny smysl slova rod.\par
\noindent\emph{Reason:} Relativno cěle oblasti prvogo roda names the class in Paper 6's Definition VII. It requires an image domain over which every polynomial in the specified variables is algebraically integral. Finiteness and the involution integral basis are a subsequent theorem, not replacements for the definition. The archive chooses rod as a classification word and rejects a weaker reading with tip; this is a preserved terminological preference. The current introduction nevertheless has prvogo tipa, while §10 has prvogo roda, so consistency still requires contextual comparison. First kind does not mean first order in logic or carry every other mathematical sense of rod.\par
\noindent{\footnotesize Izvorne zapisy / Original records: HIST-g050-001.}\par
\smallskip\noindent\emph{Sravnjenje s tekstom:} Definicija VII govori o oblasti odobraženja i algebraičnoj cělosti svih polinomov. Konečnost jest daljši teorem; tip v uvodu i rod v sekciji ostavajut dokumentovane variante.\par
\noindent\emph{Textual comparison:} Definition VII concerns the image domain and algebraic integrality of all polynomials. Finiteness is a subsequent theorem; introductory tip and sectional rod remain documented variants.\par
\noindent{\footnotesize TEXT-COMP-590; pasaže / passages: 5025, 6057, 6059, 6061, 6063, 6065, 6067, 6069, 6093; \texttt{CORPUS\_FORM\_COMPARISONS.json}.}\par
\end{minipage}\par
\par\medskip\Needspace{7\baselineskip}
\noindent\begin{minipage}{\linewidth}
\hypertarget{TERM-1532}{}
\noindent\textbf{relativno-cěle oblasti prvogo roda}\quad{\footnotesize\texttt{TERM-1532}}\par
\noindent\emph{Iztočnik / Source:} relativ-ganze Bereiche erster Art\par
\noindent\emph{Rusky / Russian:} относительно целые области первого рода\par
\noindent\emph{Ukrajinsky / Ukrainian:} відносно цілі області першого роду\par
\smallskip\noindent\emph{Razlog:} Relativno-cěle oblasti prvogo roda prodolžuje nazvanje iz §10 pri priměrah §13. Spojnik v historičnom zapisu grupuje relativno-cěle, ale ne sodrživaje dokaz noveho matematičnogo uslovja. Prvy priměr govori o polinomah invariantnyh pri danoj permutacijnoj grupě; uravnenje s elementarnymi simetričnymi funkcijami ustanovjaje potrěbnu algebraičnu cělost. Arhivni razlog jest kontinuitet, s izrično zachovanoju relativnostju. Tu konkretnu priměnu treba sravniti s definicijeju, a ne nazvati vsaky konečno porodženy prsten oblastju prvogo roda bez prověrky jeje uslovij. Razlika medžu zapisom so spojnikom i bez njega ostavaje vidima v slovniku.\par
\noindent\emph{Reason:} Relativno-cěle oblasti prvogo roda continues §10's term in §13's examples. The historical hyphen groups relativno-cěle but supplies no evidence of a new mathematical condition. The first example concerns polynomials invariant under a specified permutation group; the equation using elementary symmetric functions establishes the required algebraic integrality. The archived reason is continuity with explicit relativity. This application must be checked against the definition rather than calling every finitely generated ring a domain of the first kind without checking its conditions. Hyphenated and unhyphenated records remain distinguishable in the dictionary.\par
\noindent{\footnotesize Izvorne zapisy / Original records: HIST-g053-001.}\par
\smallskip\noindent\emph{Sravnjenje s tekstom:} Uravnenje s elementarnymi simetričnymi funkcijami opravdava priměnjenje definicije prvogo roda v konkretnoj oblasti.\par
\noindent\emph{Textual comparison:} The equation using elementary symmetric functions justifies applying the first-kind definition to the particular domain.\par
\noindent{\footnotesize TEXT-COMP-591; pasaže / passages: 6365, 6367, 6369, 6370, 6373; \texttt{CORPUS\_FORM\_COMPARISONS.json}.}\par
\end{minipage}\par
\par\medskip\Needspace{7\baselineskip}
\noindent\begin{minipage}{\linewidth}
\hypertarget{TERM-1533}{}
\noindent\textbf{relativno cěle funkcije}\quad{\footnotesize\texttt{TERM-1533}}\par
\noindent\emph{Iztočnik / Source:} relativ-ganze Funktionen\par
\noindent\emph{Rusky / Russian:} относительно целые функции\par
\noindent\emph{Ukrajinsky / Ukrainian:} відносно цілі функції\par
\smallskip\noindent\emph{Razlog:} Fraza relativno cěle funkcije vračaje se v rabotah 23 i 35, ale s različnym okolnym zadaniem. Rabota 23 govori o konečnosti oblastij vsih polinomov iz polja racionalnyh funkcij. Rabota 35 razgleduje maksimalnu oblast cěločislnyh funkcij i razlikuje konečno čislo prvočisel v imenovateljah od ograničenosti ih eksponentov. Arhiv čuva historično nazvanje Hilbertovoj problemy, vměsto tihoj zaměny novym algebraičnym opisom. To ne opravdava izpuščenje dodatnyh cěločislnyh uslovij. V nynějšnoj pasaži raboty 35 jest protivrečivo zapisano uslovje s $a g$ i $a^{\varrho}g$; to jest otvorjeno iztočnikovo pytanje, ne potvrđena čest definicije.\par
\noindent\emph{Reason:} Relativno cěle funkcije recurs in Papers 23 and 35 with different surrounding questions. Paper 23 concerns finite generation of the domains of all polynomials in a field of rational functions. Paper 35 examines a maximal domain of integer-coefficient functions and distinguishes finitely many denominator primes from bounded exponents of those primes. The archive retains the historical Hilbert-problem label instead of silently replacing it with a modern algebraic description. This does not justify dropping the additional arithmetic conditions. The current Paper 35 passage contains a contradictory condition involving $a g$ and $a^{\varrho}g$; that is an open source question, not an endorsed part of the definition.\par
\noindent{\footnotesize Izvorne zapisy / Original records: HIST-g094-009, HIST-g205-004.}\par
\smallskip\noindent\emph{Sravnjenje s tekstom:} Rabota 23 govori o polinomah polja racionalnyh funkcij; rabota 35 dodava cěločislne uslovja i ograničenje imenovateljev. Protivrěčivo uslovje v redku 22531 ostavaje otvorjeno primarnomu sravnjenju.\par
\noindent\emph{Textual comparison:} Paper 23 concerns polynomials in a rational-function field; Paper 35 adds arithmetic conditions and restrictions on denominators. The contradictory condition at line 22531 remains open to primary-source comparison.\par
\noindent{\footnotesize TEXT-COMP-592; pasaže / passages: 15162, 15164, 15166, 22523, 22525, 22527, 22529, 22531, 22533; \texttt{CORPUS\_FORM\_COMPARISONS.json}.}\par
\end{minipage}\par
\par\medskip\Needspace{7\baselineskip}
\noindent\begin{minipage}{\linewidth}
\hypertarget{TERM-1534}{}
\noindent\textbf{relativno cěle funkcije / oblasti prvogo roda}\quad{\footnotesize\texttt{TERM-1534}}\par
\noindent\emph{Iztočnik / Source:} relativ-ganze Funktionen / Bereiche erster Art\par
\noindent\emph{Rusky / Russian:} относительно целые функции / области первого рода\par
\noindent\emph{Ukrajinsky / Ukrainian:} відносно цілі функції / області першого роду\par
\smallskip\noindent\emph{Razlog:} Relativno cěle funkcije / oblasti prvogo roda sobira dva svezane nazvanja iz uvoda raboty 6. Kosa črta ne znači, že funkcija i oblast jesut jeden tip objekta. Prvy izraz upuča k Hilbertovoj problemě, drugy k specialnoj klasě oblastij, za ktoru rabota dava pozitiven rezultat. Arhiv prodolžuje jezyk raboty 5 i čuva historično nazvanje klasy. Uvod kratko objasnjaje dosah rezultata, a §9 i §10 davajut različne definicije. Pri redakciji i budučej priměně trěba razděljati tuto motivacijnu svez od tvrđenja o ekvivalentnosti pojmov; kontinuitet terminov ne čini ih sinonimami.\par
\noindent\emph{Reason:} Relativno cěle funkcije / oblasti prvogo roda collects two connected labels from Paper 6's introduction. The slash does not make a function and a domain the same type of object. The first expression refers to Hilbert's problem, the second to the special class of domains for which the paper obtains a positive result. The archive continues Paper 5's language while retaining the historical class name. The introduction summarizes the result's reach, whereas §§9 and 10 supply distinct definitions. Editing and later reuse must distinguish that motivational connection from an assertion of conceptual equivalence; terminological continuity does not make the terms synonyms.\par
\noindent{\footnotesize Izvorne zapisy / Original records: HIST-g040-013.}\par
\smallskip\noindent\emph{Sravnjenje s tekstom:} Motivacijny svez iz uvoda ne poistovětjaje funkciju s oblastju. Definicije razlikujut relativnu cělost od dodatnogo uslovja prvogo roda.\par
\noindent\emph{Textual comparison:} The introductory motivation does not identify a function with a domain. The definitions distinguish relative integrality from the first-kind condition.\par
\noindent{\footnotesize TEXT-COMP-593; pasaže / passages: 5025, 5975, 5977, 5996, 6065, 6069; \texttt{CORPUS\_FORM\_COMPARISONS.json}.}\par
\end{minipage}\par
\par\medskip\Needspace{7\baselineskip}
\noindent\begin{minipage}{\linewidth}
\hypertarget{TERM-1535}{}
\noindent\textbf{relativno cěla oblast}\quad{\footnotesize\texttt{TERM-1535}}\par
\noindent\emph{Iztočnik / Source:} relativ-ganzer Bereich\par
\noindent\emph{Rusky / Russian:} относительно целая область\par
\noindent\emph{Ukrajinsky / Ukrainian:} відносно ціла область\par
\smallskip\noindent\emph{Razlog:} Relativno cěla oblast prodolžuje jezyk relativno cělyh funkcij i zachovuje oblast kak algebraičny objekt. Uslovje 4a raboty 6 jest točno: ako $F$ i $G$ sut v oblasti i $G\ne0$, kvocient $F/G$ naleži oblasti togda i samo togda, kogda on jest polinom. Iz togo tekst vyvodi opis posrědstvom vseh polinomov najmenšego sodrživajučego polja. Arhiv čuva tuto leksičnu familiju, ne navrhuje novy geometričny smysl oblasti. Izraz ne možno avtomatično zaměniti integralno zamknjena oblast: takova zaměna trěbovala by točno sravnjenje definicij i osnovnyh prstenov. Relativna cělost i algebraična cělost veličiny v §9 jesut svezane, ale različne uslovja.\par
\noindent\emph{Reason:} Relativno cěla oblast continues the language of relatively integral functions and retains oblast as an algebraic object. Paper 6's condition 4a is precise: for $F$ and $G$ in the domain with $G\ne0$, the quotient $F/G$ belongs to the domain exactly when it is a polynomial. The text derives the description using all polynomials in the smallest containing field. The archive maintains this lexical family rather than proposing a geometric sense of domain. The phrase must not automatically be replaced by integrally closed domain without comparing definitions and ambient rings. Relative integrality of the domain and algebraic integrality of a quantity in §9 are connected but distinct conditions.\par
\noindent{\footnotesize Izvorne zapisy / Original records: HIST-g049-001.}\par
\smallskip\noindent\emph{Sravnjenje s tekstom:} Kvocient prinaleži oblasti togda i samo togda, kogda jest polinom. Tekst potom daje ekvivalentny opis pomoču najmenšego sodrživajučego polja.\par
\noindent\emph{Textual comparison:} The quotient belongs to the domain exactly when it is a polynomial. The text then gives the equivalent description through the smallest containing field.\par
\noindent{\footnotesize TEXT-COMP-594; pasaže / passages: 5973, 5975, 5977, 5979, 5981, 5983, 5985, 5987, 5989, 5993, 5996, 6000, 6002; \texttt{CORPUS\_FORM\_COMPARISONS.json}.}\par
\end{minipage}\par
\par\medskip\Needspace{7\baselineskip}
\noindent\begin{minipage}{\linewidth}
\hypertarget{TERM-1536}{}
\noindent\textbf{relativny diskriminant}\quad{\footnotesize\texttt{TERM-1536}}\par
\noindent\emph{Iztočnik / Source:} Relativdiskriminante\par
\noindent\emph{Rusky / Russian:} относительный дискриминант\par
\noindent\emph{Ukrajinsky / Ukrainian:} відносний дискримінант\par
\smallskip\noindent\emph{Razlog:} Relativny diskriminant čuva literaturno nazvanje, ktoro podnožna primětka raboty 31 navodi v navodnikah. V tutom kontekstu ide o diskriminantny ideal glavnogo poredka po odnošenju k osnovnomu poredku, ne prosto o jeden proizvoljny determinant. Pozdnějša primětka sravnjuje konstrukcije Hilberta i Hecke posrědstvom lokalizacij i determinantov konjugovanyh elementov. Arhiv radi toj svezi čuva kompaktno nazvanje vměsto parafrazy. Takovo čitanje primětky ne jest nova nezavisna prověrka knig Hilberta i Hecke. Takože ne rěšaje ranějše registrovano pytanje o susědnom upučenju po punktu 1.\par
\noindent\emph{Reason:} Relativny diskriminant retains the literature label quoted in Paper 31's footnote. Here it is the discriminant ideal of the principal order relative to the base order, not simply an arbitrary single determinant. A later footnote compares the Hilbert and Hecke constructions through localization and determinants of conjugate elements. The archive keeps the compact label to preserve this connection rather than replacing it with a paraphrase. Reading that footnote is not a new independent verification of Hilbert's or Hecke's books. Nor does it resolve the previously recorded question about the adjacent reference to point 1.\par
\noindent{\footnotesize Izvorne zapisy / Original records: HIST-g155-006, HIST-g162-005.}\par
\smallskip\noindent\emph{Sravnjenje s tekstom:} Primětky zachovujut literaturno nazvanje i sravnjenje s konstrukcijami Hilberta i Hecke. Diskriminantny ideal ne jest prosto zaměnjen jednim determinantom.\par
\noindent\emph{Textual comparison:} The footnotes retain the literature label and comparison with the Hilbert and Hecke constructions. The discriminant ideal is not simply replaced by a single determinant.\par
\noindent{\footnotesize TEXT-COMP-595; pasaže / passages: 18785, 18787, 18789, 19037; \texttt{CORPUS\_FORM\_COMPARISONS.json}.}\par
\end{minipage}\par
\par\medskip\Needspace{7\baselineskip}
\noindent\begin{minipage}{\linewidth}
\hypertarget{TERM-1537}{}
\noindent\textbf{relativny diskriminant razširjenogo polja}\quad{\footnotesize\texttt{TERM-1537}}\par
\noindent\emph{Iztočnik / Source:} Relativdiskriminante eines Erweiterungskörpers\par
\noindent\emph{Rusky / Russian:} относительный дискриминант поля расширения\par
\noindent\emph{Ukrajinsky / Ukrainian:} відносний дискримінант поля розширення\par
\smallskip\noindent\emph{Razlog:} Relativny diskriminant razširjenogo polja razkriva, kojemu polju prinaleži relativno-aritmetičny podatok. Rabota 31 zaključaje, že prosty ideal osnovnogo glavnogo poredka vhodi v toj diskriminant togda i samo togda, kogda v razširjenom glavnom poredku ima kvadrat prostogo ideala kako dělitelj. Arhiv izbira už upotrěbljeno razširjeno polje i prodolžuje nazvanje iz Hilbertovej i Heckeovej primětky. Prednost jest kontinuitet i jasno označenje odnosa, ne nova kovanica. Polje razširjenja ne jest zaměna za osnovno polje: pri priměně frazy obě strani relativity morajut ostati poznane iz teksta.\par
\noindent\emph{Reason:} Relativny diskriminant razširjenogo polja identifies the extension field to which the relative arithmetic datum belongs. Paper 31 concludes that a prime ideal of the base principal order occurs in this discriminant exactly when it has a prime-ideal square as a divisor in the extended principal order. The archive uses the established razširjeno polje and continues the label from the Hilbert/Hecke footnote. Its reasons are continuity and explicit relational structure, not a new coinage. The extension field cannot replace the base field: both sides of the relativity must remain recoverable from the text.\par
\noindent{\footnotesize Izvorne zapisy / Original records: HIST-g164-009.}\par
\smallskip\noindent\emph{Sravnjenje s tekstom:} Specializacija teorema povezuje relativny diskriminant s kvadratom prostogo ideala v razširjenom glavnom poredku.\par
\noindent\emph{Textual comparison:} The theorem's specialization connects the relative discriminant with a prime-ideal square in the extended principal order.\par
\noindent{\footnotesize TEXT-COMP-596; pasaže / passages: 19105, 19107; \texttt{CORPUS\_FORM\_COMPARISONS.json}.}\par
\end{minipage}\par
\par\medskip\Needspace{7\baselineskip}
\noindent\begin{minipage}{\linewidth}
\hypertarget{TERM-1538}{}
\noindent\textbf{relativna invariantnost}\quad{\footnotesize\texttt{TERM-1538}}\par
\noindent\emph{Iztočnik / Source:} relative Invarianz\par
\noindent\emph{Rusky / Russian:} относительная инвариантность\par
\noindent\emph{Ukrajinsky / Ukrainian:} відносна інваріантність\par
\smallskip\noindent\emph{Razlog:} Relativna invariantnost v rabotě 13 ne znači približnu invariantnost. Izraz može pri transformaciji dobiti točno oprěděljen faktor; v §5 toj faktor jest Jakobijev determinant transformacije nezavisnyh prěměnnyh. Arhiv izrično razlikuje relativnu od absolutnoj invariantnosti, zato čuva pridavno. Uslovja sekcije sut suščestvene: transformacije ne sodrživajut izvodne zavisnyh veličin, a nove nezavisne prěměnne zaviset samo od staryh nezavisnyh prěměnnyh. Podnožna primětka upozornjaje na nevalidnost toga zaključka bez tyh uslovij. Imenovanje jednotlivyh čestij relacij invariantnymi ne dava pravo odstraniti faktor ili prědpoloženja.\par
\noindent\emph{Reason:} Relativna invariantnost in Paper 13 does not mean approximate invariance. The expression may acquire a precisely specified factor; in §5 that factor is the Jacobian determinant of the independent-variable transformation. The archive expressly distinguishes relative from absolute invariance and therefore retains the adjective. The section's hypotheses matter: transformations contain no derivatives of dependent quantities, and the new independent variables depend only on the old independent variables. The footnote warns that the argument does not hold without these conditions. Calling individual parts of the relations invariant does not authorize dropping the factor or hypotheses.\par
\noindent{\footnotesize Izvorne zapisy / Original records: HIST-g075-013.}\par
\smallskip\noindent\emph{Sravnjenje s tekstom:} Jakobijev faktor čuva relativnu invariantnost izraza i invariantnost integrala. Podnožna primětka izrično ograničaje argument.\par
\noindent\emph{Textual comparison:} The Jacobian factor preserves relative invariance of the expression and invariance of the integral. The footnote expressly limits the argument.\par
\noindent{\footnotesize TEXT-COMP-597; pasaže / passages: 9293, 9600, 9602, 9604, 9606, 9634, 9638, 9640, 9641, 9642, 9643, 9645, 9647, 9649, 9683, 9685; \texttt{CORPUS\_FORM\_COMPARISONS.json}.}\par
\end{minipage}\par
\par\medskip\Needspace{7\baselineskip}
\noindent\begin{minipage}{\linewidth}
\hypertarget{TERM-1539}{}
\noindent\textbf{relativny / modularny / formalny invariant}\quad{\footnotesize\texttt{TERM-1539}}\par
\noindent\emph{Iztočnik / Source:} Relativinvariante / Modularinvariante / formale Invariante\par
\noindent\emph{Rusky / Russian:} относительный / модулярный / формальный инвариант\par
\noindent\emph{Ukrajinsky / Ukrainian:} відносний / модулярний / формальний інваріант\par
\smallskip\noindent\emph{Razlog:} Relativny / modularny / formalny invariant čuva tri blizke, no ne zaměnljive nazvanja. V rabotě 29 relativny invariant pri grupovoj substituciji měnjaje se samo nenulovym faktorom iz koeficientnogo polja. Modularny invariant mora davati iste cěnnosti pri vsih podstavjenjah iz toga polja; nad konečnym poljem to ne trěbuje identičnosti polinomov. Uvod prědstavjaje razliku medžu modularnymi i obče formalnymi invariantami, a §2 dava identičnost kak uslovje za absolutny polinomny invariant. Arhiv izbira internacionalne koreni, da by razlike ostali vidime i otvorjene obosnovanoj jezyčnoj reviziji. On ne tvrdi, že te tri nazvanja razděljajut vse invarianty na tri neperesěkajuče se klasy.\par
\noindent\emph{Reason:} Relativny / modularny / formalny invariant retains three related but non-interchangeable labels. In Paper 29 a relative invariant changes under a group substitution only by a nonzero factor from the coefficient field. A modular invariant must give equal values at every substitution from that field; over a finite field this need not be a polynomial identity. The introduction contrasts modular with general formal invariants, while §2 gives identity as the condition for an absolute polynomial invariant. The archive chooses international stems to keep the distinctions visible and open to reasoned linguistic revision. It does not claim that these three labels partition all invariants into disjoint classes.\par
\noindent{\footnotesize Izvorne zapisy / Original records: HIST-g106-013.}\par
\smallskip\noindent\emph{Sravnjenje s tekstom:} Absolutna polinomna invariantnost jest identičnost; modularna invariantnost trěbuje ravnosti pri cěnnostjah iz polja. Nad konečnym poljem te uslovja ne morajut byti ekvivalentne.\par
\noindent\emph{Textual comparison:} Absolute polynomial invariance is an identity; modular invariance requires equality at values from the field. Over a finite field the conditions need not be equivalent.\par
\noindent{\footnotesize TEXT-COMP-598; pasaže / passages: 15983, 15989, 15991, 16045, 16053, 16055, 16057, 16082, 16084; \texttt{CORPUS\_FORM\_COMPARISONS.json}.}\par
\end{minipage}\par
\par\medskip\Needspace{7\baselineskip}
\noindent\begin{minipage}{\linewidth}
\hypertarget{TERM-1540}{}
\noindent\textbf{relativna norma}\quad{\footnotesize\texttt{TERM-1540}}\par
\noindent\emph{Iztočnik / Source:} Relativnorm\par
\noindent\emph{Rusky / Russian:} относительная норма\par
\noindent\emph{Ukrajinsky / Ukrainian:} відносна норма\par
\smallskip\noindent\emph{Razlog:} Relativna norma jest algebraična norma s ukazanym osnovnym poljem, ne dolžina vektora. Rabota 32 izračunava normu produktem dvoh konjugovanyh faktorov v razširjenju $P(\varepsilon,i)$ nad $P(\varepsilon)$ i dosega cěnnost $-1$. Ta negativna cěnnost sama upozornjaje, že ne ide o metriku. Arhiv dava prednost už poznatomu internacionalnomu nazvanju i izrično ne predlaga novoj tvorby; tvrđenje o jeho stabilnosti v slovjanskoj algebraičnoj literaturě ostavaje historičnym razlogom, ne rezultatom novogo obširnogo izslědovanja. Pri daljšej priměně treba vsegda zachovati polje razširjenja, osnovno polje i smysl normy.\par
\noindent\emph{Reason:} Relativna norma is an algebraic norm with a specified base field, not a vector length. Paper 32 calculates it as a product of two conjugate factors in $P(\varepsilon,i)$ over $P(\varepsilon)$ and obtains the value $-1$. That negative value itself warns against a metric interpretation. The archive favours a familiar international term and expressly avoids a new coinage; its claim of stability across Slavic algebra literature remains a historical rationale, not the result of a new comprehensive survey. Later reuse must retain the extension field, base field, and sense of norm.\par
\noindent{\footnotesize Izvorne zapisy / Original records: HIST-g175-004.}\par
\smallskip\noindent\emph{Sravnjenje s tekstom:} Produkt konjugovanyh faktorov daje relativnu normu, a daljši račun daje cěnnost $-1$. Smysl jest algebraičny, ne metričny.\par
\noindent\emph{Textual comparison:} The product of conjugate factors gives the relative norm, and the subsequent calculation gives the value $-1$. The sense is algebraic, not metric.\par
\noindent{\footnotesize TEXT-COMP-599; pasaže / passages: 19608, 19610, 19612, 19616, 19618, 19620, 19622, 19624, 19625, 19626, 19628, 19630, 19632, 19634, 19636; \texttt{CORPUS\_FORM\_COMPARISONS.json}.}\par
\end{minipage}\par
\par\medskip\Needspace{7\baselineskip}
\noindent\begin{minipage}{\linewidth}
\hypertarget{TERM-1541}{}
\noindent\textbf{relativno proste}\quad{\footnotesize\texttt{TERM-1541}}\par
\noindent\emph{Iztočnik / Source:} relativprim / relativ prim\par
\noindent\emph{Rusky / Russian:} относительно просты\par
\noindent\emph{Ukrajinsky / Ukrainian:} відносно прості\par
\smallskip\noindent\emph{Razlog:} Relativno proste v uvodu raboty 19 imenuje lokalny dělitelnostny odnos: ako je produkt $bq_i$ dělimy prěz $q_k$, trěba, aby takože $b$ byl dělimy prěz $q_k$. Arhiv zato ne zaměnjaje tutu frazu izrazom za poparnu vzajemnu prostost bez dalnogo uslovja. §7 pozdněje razlikuje vzajemno proste idealy od vzajemno relativno prostyh; obrat prvoj implikacije ne imaje město. To ne jest samo stilistična razlika medžu relativno i vzajemno. Slovnik čuva definiciju i smer odnosa, aby znano slovo prosty ne uneslo silnějše uslovje iz aritmetike cělyh čisel. Dva povezana arhivna zapisa ne dokazyvajut dva nezavisne jezyčne testa.\par
\noindent\emph{Reason:} Relativno proste in Paper 19's introduction denotes a local divisibility relation: divisibility of $bq_i$ by $q_k$ must imply divisibility of $b$ by $q_k$. The archive therefore does not substitute an unqualified expression for pairwise coprimality. Section 7 later distinguishes coprime ideals from mutually relatively prime ideals and explicitly denies the converse implication. This is not merely stylistic variation between relativno and vzajemno. The definition and direction must remain visible so that the familiar adjective does not import the stronger condition from integer arithmetic. Two linked archival records do not establish two independent linguistic tests.\par
\noindent{\footnotesize Izvorne zapisy / Original records: HIST-g081-005, HIST-g082-005.}\par
\smallskip\noindent\emph{Sravnjenje s tekstom:} Dělitelnostny kriterij i pozdnějša definicija vzajemnoj prostosti sut različne. Obrat implikacije jest izrično izključen.\par
\noindent\emph{Textual comparison:} The divisibility criterion differs from the later coprimality definition; the converse implication is expressly excluded.\par
\noindent{\footnotesize TEXT-COMP-600; pasaže / passages: 12420, 13278, 13280; \texttt{CORPUS\_FORM\_COMPARISONS.json}.}\par
\end{minipage}\par
\par\medskip\Needspace{7\baselineskip}
\noindent\begin{minipage}{\linewidth}
\hypertarget{TERM-1542}{}
\noindent\textbf{relativno prosty k}\quad{\footnotesize\texttt{TERM-1542}}\par
\noindent\emph{Iztočnik / Source:} relativprim / relativprim zu\par
\noindent\emph{Rusky / Russian:} относительно прост к\par
\noindent\emph{Ukrajinsky / Ukrainian:} відносно простий до\par
\smallskip\noindent\emph{Razlog:} Relativno prosty k jest usměrjena konstrukcija. Definicija V trěbuje, da iz $TR\equiv0\pmod S$ slěduje $T\equiv0\pmod S$; vzajemnost trěbuje i obraten odnos. Podnožna primětka izrično daje priměr nesimetrije s idealami $(x^2,y)$ i $(x)$. Arhiv zato čuva prědlog k, vměsto da bi frazu tiho prěměnil v simetrično svojstvo para. Ruska konstrukcija s k i ukrajinska s do služat sravnjenju upravjanja, ne rěšaju matematičnu definiciju. Pri prěnosu v drugu rabotu trěba prověriti oboje: kto jest relativno prosty k komu i kak se ta relacija tamo definira.\par
\noindent\emph{Reason:} Relativno prosty k is a directed construction. Definition V requires $TR\equiv0\pmod S$ to imply $T\equiv0\pmod S$; mutuality requires the reverse relation as well. Its footnote explicitly demonstrates asymmetry using $(x^2,y)$ and $(x)$. The archive retains k rather than silently turning the expression into a symmetric property of a pair. Russian k and Ukrainian do provide comparanda for grammatical government, not the mathematical definition. Reuse in another work must check both the direction and the local definition of the relation.\par
\noindent{\footnotesize Izvorne zapisy / Original records: HIST-g081-051, HIST-g082-051.}\par
\smallskip\noindent\emph{Sravnjenje s tekstom:} Definicija i priměr čuvajut nesimetriju; vzajemnost trěbuje oba směry.\par
\noindent\emph{Textual comparison:} The definition and example preserve asymmetry; mutuality requires both directions.\par
\noindent{\footnotesize TEXT-COMP-601; pasaže / passages: 13062, 13064, 13066; \texttt{CORPUS\_FORM\_COMPARISONS.json}.}\par
\end{minipage}\par
\par\medskip\Needspace{7\baselineskip}
\noindent\begin{minipage}{\linewidth}
\hypertarget{TERM-1543}{}
\noindent\textbf{relativno-prosto-nerazlozimy}\quad{\footnotesize\texttt{TERM-1543}}\par
\noindent\emph{Iztočnik / Source:} relativprim-irreduzibel\par
\noindent\emph{Rusky / Russian:} относительно-просто-неразложимый\par
\noindent\emph{Ukrajinsky / Ukrainian:} відносно-просто-нерозкладний\par
\smallskip\noindent\emph{Razlog:} Zapis relativno-prosto-nerazlozimy čuva složeno nazvanje razloženja iz §6 raboty 19. Ideal jest takovy, ako se ne može prědstaviti kako najmenše obče mnogokratno vzajemno relativno prostyh vlastnyh děliteljev. Dodatok relativno-prosto zato ograničuje dopustimo razloženje; ne smi otpasti pri kračenju do nerazložimy. Arhivny zapis bez ž ostavaje čitliv kak zapis, ale nynějšny korpus i sosedny termin imaju nerazložimy. Ta bukvena razlika sama ne ustanovjava novu klasu idealov. Nema čislovogo sravnjenja kandidatov; zachovany razlog jest potreba držati uslovje relativnoj prostosti skupaj s nerazložimostju.\par
\noindent\emph{Reason:} The record relativno-prosto-nerazlozimy preserves the compound naming Paper 19, §6's decomposition class. An ideal belongs to this class when it cannot be expressed as a least common multiple of mutually relatively prime proper divisors. Relativno-prosto restricts the permissible decomposition and cannot be dropped when shortening the label to indecomposable. The archive's spelling without ž is retained as recorded, whereas the current corpus and adjacent entry have nerazložimy. That letter difference does not define a new class of ideals. No numerical candidate comparison survives; the recorded reason is to keep relative primeness attached to indecomposability.\par
\noindent{\footnotesize Izvorne zapisy / Original records: HIST-g082-053.}\par
\smallskip\noindent\emph{Sravnjenje s tekstom:} Definicija trěbuje vlastne děliteli i vzajemnu relativnu prostost. Jača vzajemna prostost jest poseben pojem.\par
\noindent\emph{Textual comparison:} The definition requires proper divisors and mutual relative primeness. Stronger coprimality is a separate notion.\par
\noindent{\footnotesize TEXT-COMP-602; pasaže / passages: 13060, 13068, 13278, 13280; \texttt{CORPUS\_FORM\_COMPARISONS.json}.}\par
\end{minipage}\par
\par\medskip\Needspace{7\baselineskip}
\noindent\begin{minipage}{\linewidth}
\hypertarget{TERM-1544}{}
\noindent\textbf{relativno-prosto-nerazložimy}\quad{\footnotesize\texttt{TERM-1544}}\par
\noindent\emph{Iztočnik / Source:} relativprim-irreduzibel\par
\noindent\emph{Rusky / Russian:} относительно-просто-неразложим\par
\noindent\emph{Ukrajinsky / Ukrainian:} відносно-просто-нерозкладний\par
\smallskip\noindent\emph{Razlog:} Relativno-prosto-nerazložimy jest izrično konstruovany složeny termin, ne novo matematično svojstvo dodano prěvodom. Jeho tri česti čuvajut definiciju iz §6: nerazložimost razgledaje se samo po odnošenju k vzajemno relativno prostym vlastnym děliteljam. Arhiv označuje termin kako izbor kontrolovanogo registra. To znači projektnu terminologičnu odluku, ne potvrdženje od jezyčnoj obščiny. Smysl ne jest totožny s pozdnějšim vzajemno-prosto-nerazložimy, ktory se opira na silnějšu prostost idealov. Potrěbna jest kratka definicija pri prvom vvedenju; sam složeny izraz ne zaměnjaje jeje.\par
\noindent\emph{Reason:} Relativno-prosto-nerazložimy is an explicitly constructed compound, not a new mathematical property introduced by translation. Its components preserve §6's definition: indecomposability is relative to mutually relatively prime proper divisors. The archive marks it as a controlled-register choice. That is a project terminology decision, not approval by the language community. It is distinct from the later vzajemno-prosto-nerazložimy, which uses the stronger coprimality condition. A short definition is needed at first introduction; the compound alone does not replace it.\par
\noindent{\footnotesize Izvorne zapisy / Original records: HIST-g081-053.}\par
\smallskip\noindent\emph{Sravnjenje s tekstom:} Definicija trěbuje vlastne děliteli i vzajemnu relativnu prostost. Jača vzajemna prostost jest poseben pojem.\par
\noindent\emph{Textual comparison:} The definition requires proper divisors and mutual relative primeness. Stronger coprimality is a separate notion.\par
\noindent{\footnotesize TEXT-COMP-602; pasaže / passages: 13060, 13068, 13278, 13280; \texttt{CORPUS\_FORM\_COMPARISONS.json}.}\par
\end{minipage}\par
\par\medskip\Needspace{7\baselineskip}
\noindent\begin{minipage}{\linewidth}
\hypertarget{TERM-1545}{}
\noindent\textbf{sistem predstaviteljev}\quad{\footnotesize\texttt{TERM-1545}}\par
\noindent\emph{Iztočnik / Source:} Repräsentantensystem\par
\noindent\emph{Rusky / Russian:} система представителей\par
\noindent\emph{Ukrajinsky / Ukrainian:} система представників\par
\smallskip\noindent\emph{Razlog:} Sistem predstaviteljev nazyvaje izbrane prědstavitelje klasov, ne prědstavjenje grupy matricami. Arhiv izrično predpočitaje tuto konstrukciju doslovnomu sistemu prědstavjenja. V §4 raboty 22 polinomno děljenje davaje prědstavitelje, ktoryh stepenj po $x_1$ jest menši než $k$; to uslovje rabotaje v dokazu izomorfizma sistemov klasov. Sam izbor prědstaviteljev ješče ne jest dokaz izomorfizma: definicija trěbuje čuvanje razlik i množenja toju samoju cěloju veličinoju. Arhivno predstaviteljev i nynějšno prědstaviteljev sut odlišene v zapisu, bez izmyšljanja dvuh matematičnyh pojmov.\par
\noindent\emph{Reason:} Sistem predstaviteljev names chosen representatives of classes, not a matrix representation of a group. The archive expressly prefers this construction to a literal representation system. In Paper 22, §4, polynomial division gives representatives of degree less than $k$ in $x_1$; that bound is used in the residue-system isomorphism argument. Choosing representatives does not itself prove an isomorphism: the definition also requires preservation of differences and multiplication by the same integral quantity. Archival predstaviteljev and current prědstaviteljev remain distinguishable spellings without inventing two mathematical concepts.\par
\noindent{\footnotesize Izvorne zapisy / Original records: HIST-g088-004, HIST-g092-035.}\par
\smallskip\noindent\emph{Sravnjenje s tekstom:} Prědstavitelji imajut stepenj menšu než k. Definicija izomorfizma takože čuva operacije, ne samo izbor prědstaviteljev.\par
\noindent\emph{Textual comparison:} The representatives have degree less than k. The isomorphism definition also preserves operations, not merely a choice of representatives.\par
\noindent{\footnotesize TEXT-COMP-603; pasaže / passages: 14631, 14635, 14637, 14639, 14642; \texttt{CORPUS\_FORM\_COMPARISONS.json}.}\par
\end{minipage}\par
\par\medskip\Needspace{7\baselineskip}
\noindent\begin{minipage}{\linewidth}
\hypertarget{TERM-1546}{}
\noindent\textbf{residualny modul / kvocient dvoh modulov}\quad{\footnotesize\texttt{TERM-1546}}\par
\noindent\emph{Iztočnik / Source:} Residualmodul / Quotient zweier Moduln\par
\noindent\emph{Rusky / Russian:} остаточный модуль / частное двух модулей\par
\noindent\emph{Ukrajinsky / Ukrainian:} залишковий модуль / частка двох модулів\par
\smallskip\noindent\emph{Razlog:} Residualny modul i kvocient dvoh modulov sut tu dva historična nazvanja citovanyh konstrukcij. Podnožna primětka identificira $T_0$ s Laskerovym residualnym modulom i, pri razširjenju na čislove moduly, s Dedekindovym kvocientom. $T_0$ jest definovany črez uslovje $T_0R\equiv0\pmod S$. Ne treba jeho zaměnjati s modulom klasov ostatkov, v ktorom elementy sut kongruencijne klasy. Arhiv čuva internacionalne koreni radi točnoj svezi s citovanoj literaturoj; to ne jest dokaz, že vse sovrěmenne izrazy s kvocient imajut isti smysl. Izvorne raboty Laskera i Dedekinda ne byly za tuto sravnjenje nanovo pročitane.\par
\noindent\emph{Reason:} Residualny modul and kvocient dvoh modulov are historical names for the cited constructions. The footnote identifies $T_0$ with Lasker's residual module and, after extension to number modules, with Dedekind's quotient. The defining condition is $T_0R\equiv0\pmod S$. This must not be confused with a residue-class module whose elements are congruence classes. The archive retains international stems to preserve the connection with the cited literature, not because all modern uses of quotient have the same sense. Lasker's and Dedekind's original works were not newly read for this comparison.\par
\noindent{\footnotesize Izvorne zapisy / Original records: HIST-g081-055, HIST-g082-055.}\par
\smallskip\noindent\emph{Sravnjenje s tekstom:} Definicija T0 jest dana črez produktno uslovje; citacija ju povezuje s Laskerom i Dedekindom, ne s prostu kolekcijeju ostatkovyh klasov.\par
\noindent\emph{Textual comparison:} T0 is defined by a product condition; the citation connects it with Lasker and Dedekind, not with a mere collection of residue classes.\par
\noindent{\footnotesize TEXT-COMP-604; pasaže / passages: 13070; \texttt{CORPUS\_FORM\_COMPARISONS.json}.}\par
\end{minipage}\par
\par\medskip\Needspace{7\baselineskip}
\noindent\begin{minipage}{\linewidth}
\hypertarget{TERM-1547}{}
\noindent\textbf{faktor-grupy}\quad{\footnotesize\texttt{TERM-1547}}\par
\noindent\emph{Iztočnik / Source:} Restgruppen\par
\noindent\emph{Rusky / Russian:} фактор-группы\par
\noindent\emph{Ukrajinsky / Ukrainian:} фактор-групи\par
\smallskip\noindent\emph{Razlog:} Faktor-grupy jest arhivny izbor za Restgruppen v primětce k Teoremu III raboty 19. Primětka povezuje ideal $C$ s faktor-grupoju, a dva označene idealy s podgrupami, v ktore se ona razlagaje. Arhiv upotrěbljaje sovrěmenno faktor-grupno nazvanje zajedno s ruskoju i ukrajinskoju formoju, vměsto novogo samostatnogo slova. To jest objasnjen izbor registra. Ne daje oprěděliti vsaku upotrěbu Restgruppe kako prostu abstraktnu grupu bez dopolnitelnyh operacij: osoblivo rabota 17 izrično čuva dejstvo polinomov. Jeje uslovja trěba čitati oddělno, ne ujednačati samym slovnikom.\par
\noindent\emph{Reason:} Faktor-grupy is the recorded rendering of Restgruppen in Paper 19's footnote to Theorem III. That note associates the ideal $C$ with a factor group and two specified ideals with the subgroups into which it decomposes. The archive uses contemporary factor-group terminology alongside the Russian and Ukrainian forms instead of coining a separate word. This is an explained register choice. It does not establish that every occurrence of Restgruppe means a bare abstract group: Paper 17 expressly retains polynomial actions. Those conditions must be read separately, not flattened by glossary matching.\par
\noindent{\footnotesize Izvorne zapisy / Original records: HIST-g081-036, HIST-g082-036.}\par
\smallskip\noindent\emph{Sravnjenje s tekstom:} Primětka izrično prirědžuje idealam faktor-grupu i podgrupy. Ona jest lokalna terminologična svez, ne obča definicija vseh ostatkovyh objektov.\par
\noindent\emph{Textual comparison:} The footnote explicitly associates ideals with a factor group and subgroups. This is a local terminological connection, not a definition of every residue object.\par
\noindent{\footnotesize TEXT-COMP-605; pasaže / passages: 12641; \texttt{CORPUS\_FORM\_COMPARISONS.json}.}\par
\end{minipage}\par
\par\medskip\Needspace{7\baselineskip}
\noindent\begin{minipage}{\linewidth}
\hypertarget{TERM-1548}{}
\noindent\textbf{grupy ostatkov sut izomorfne}\quad{\footnotesize\texttt{TERM-1548}}\par
\noindent\emph{Iztočnik / Source:} Restgruppen isomorph\par
\noindent\emph{Rusky / Russian:} группы остатков изоморфны\par
\noindent\emph{Ukrajinsky / Ukrainian:} групи лишків ізоморфні\par
\smallskip\noindent\emph{Razlog:} Grupy ostatkov sut izomorfne čuva objekty Teorema VI raboty 17: grupy, pridružene modulam jednakogo vida. Ale grupovy nivo ne znači, že trěba čuvati samo složenje. Ranějša definicija izomorfizma trěbuje takože čuvanje množenja klasa vsakym istim polinomom. Tuto uslovje jest sučestveno; svobodna zaměna s izomorfizmoju goloj aditivnoj grupy oslabila by tvrđenje. Arhivny razlog isključa zaměnu samih ostatkovyh grup modulami ili jih polinomnymi prědstaviteljami. Nynějšno sravnjenje zato utočnjaje, kakva struktura na tyh grupah ostavaje čuvana.\par
\noindent\emph{Reason:} Grupy ostatkov sut izomorfne preserves the objects of Paper 17's Theorem VI: the residue groups attached to modules of the same kind. Group-level wording does not, however, mean that only addition is preserved. The earlier definition of isomorphism also requires preservation of multiplication of a class by each same polynomial. That condition matters: replacing it with an isomorphism of bare additive groups would weaken the assertion. The archived rationale excludes substituting the original modules or their polynomial representatives for the residue groups. The present comparison clarifies the structure that those groups retain.\par
\noindent{\footnotesize Izvorne zapisy / Original records: HIST-g079-050.}\par
\smallskip\noindent\emph{Sravnjenje s tekstom:} Pored složenja definicija izrično čuva množenje vsakym istim polinomom. Teorem VI zato ne govori samo o golyh aditivnyh grupah.\par
\noindent\emph{Textual comparison:} The definition preserves multiplication by each same polynomial as well as addition. Theorem VI therefore does not concern bare additive groups alone.\par
\noindent{\footnotesize TEXT-COMP-606; pasaže / passages: 11206, 11650, 11652, 11653, 11655, 11657, 11658, 11664; \texttt{CORPUS\_FORM\_COMPARISONS.json}.}\par
\end{minipage}\par
\par\medskip\Needspace{7\baselineskip}
\noindent\begin{minipage}{\linewidth}
\hypertarget{TERM-1549}{}
\noindent\textbf{klasa ostatkov / ostatkova klasa}\quad{\footnotesize\texttt{TERM-1549}}\par
\noindent\emph{Iztočnik / Source:} Restklasse\par
\noindent\emph{Rusky / Russian:} класс вычетов\par
\noindent\emph{Ukrajinsky / Ukrainian:} клас лишків\par
\smallskip\noindent\emph{Razlog:} Klasa ostatkov i ostatkova klasa sut dvě gramatične konstrukcije za isti osnovny pojem, ne nazvanja dvuh raznyh faktornih objektov. Arhivny počatkovy slovnik raboty 22 čuva analitičnu formu s genitivom, ale odkazuje na pridavnu formu v §2. Tam se kongruentne linearne formy sobirajut v jednu klasu, a sistem klasov imaje označeno početkovo prostranstvo i modul. Ne treba čitati klasu kako jedinogo izbranogo prědstavitelja. Pri novom tekstu možno izbrati doslědnu formu, ale konkretnu gramatičnu pravilnost i istoriju izmeny trěba prověriti, ne izvoditi iz samoj podobnosti slov.\par
\noindent\emph{Reason:} Klasa ostatkov and ostatkova klasa are two grammatical constructions for the same basic notion, not names for two different quotient objects. Paper 22's terminology seed preserves the analytic genitive expression while pointing to the adjectival form in §2. There, congruent linear forms are grouped into a class, and the system of classes has a specified ambient object and modulus. A class is not a single chosen representative. A new text may adopt consistent phrasing, but grammatical correctness and the history of a change require evidence rather than inference from similar words.\par
\noindent{\footnotesize Izvorne zapisy / Original records: HIST-g093-011.}\par
\smallskip\noindent\emph{Sravnjenje s tekstom:} Definicija razlikuje klasu, sistem klasov i izbranogo prědstavitelja. Dalši tekst razlikuje kardinalnost od linearnoj dimenzije.\par
\noindent\emph{Textual comparison:} The definition distinguishes a class, a system of classes and a representative. The subsequent text distinguishes cardinality from linear dimension.\par
\noindent{\footnotesize TEXT-COMP-607; pasaže / passages: 14404, 14414; \texttt{CORPUS\_FORM\_COMPARISONS.json}.}\par
\end{minipage}\par
\par\medskip\Needspace{7\baselineskip}
\noindent\begin{minipage}{\linewidth}
\hypertarget{TERM-1550}{}
\noindent\textbf{ostatkova klasa}\quad{\footnotesize\texttt{TERM-1550}}\par
\noindent\emph{Iztočnik / Source:} Restklasse\par
\noindent\emph{Rusky / Russian:} класс вычетов\par
\noindent\emph{Ukrajinsky / Ukrainian:} клас лишків\par
\smallskip\noindent\emph{Razlog:} Ostatkova klasa v §6 raboty 22 prodolžuje ranějšu pridavnu konstrukciju. Teorem XI razgledaje ideal $\mathfrak b$ modulo polinoma $h(t)$; vybrane klasy, prědstavjene polinomami $f(t),tf(t),\ldots$, tvoret linearno nezavisny sistem nad $P$. Čislo $k-p$ jest rang, a ne avtomatično čislo vseh klasov. Zachovana odluka daje kontinuitet nazvanja, ne novu kovanciju ili novu kvantitativnu ocenu. Pri priměnjenju važno jest odděliti klasu, polinom-prědstavitelja i bazu sistema klasov. Samo jednotno pravopisanje ne garantuje te razlike.\par
\noindent\emph{Reason:} Ostatkova klasa in Paper 22, §6 continues the earlier adjectival construction. Theorem XI considers the ideal $\mathfrak b$ modulo a polynomial $h(t)$; the classes represented by $f(t),tf(t),\ldots$ form a linearly independent system over $P$. The number $k-p$ is a rank, not automatically the cardinality of all classes. The recorded decision supplies terminological continuity, not a new coinage or a new quantitative assessment. Application must distinguish a class, a polynomial representative and a basis for the system of classes. Uniform spelling alone does not secure those distinctions.\par
\noindent{\footnotesize Izvorne zapisy / Original records: HIST-g090-005, HIST-g092-059.}\par
\smallskip\noindent\emph{Sravnjenje s tekstom:} Teorem i dokaz davajut bazu iz k-p klasov. To jest linearna nezavisnost nad P, ne avtomatično konečno čislo vseh klasov.\par
\noindent\emph{Textual comparison:} The theorem and proof give a basis of k-p classes. This concerns linear independence over P, not automatically a finite total number of classes.\par
\noindent{\footnotesize TEXT-COMP-608; pasaže / passages: 14935, 14937; \texttt{CORPUS\_FORM\_COMPARISONS.json}.}\par
\end{minipage}\par
\par\medskip\Needspace{7\baselineskip}
\noindent\begin{minipage}{\linewidth}
\hypertarget{TERM-1551}{}
\noindent\textbf{ostatkova klasa / sistem ostatkovyh klasov}\quad{\footnotesize\texttt{TERM-1551}}\par
\noindent\emph{Iztočnik / Source:} Restklasse / System der Restklassen\par
\noindent\emph{Rusky / Russian:} класс вычетов / система классов вычетов\par
\noindent\emph{Ukrajinsky / Ukrainian:} клас лишків / система класів лишків\par
\smallskip\noindent\emph{Razlog:} Ostatkova klasa imenuje jednu klasu, a sistem ostatkovyh klasov cěly označeny sistem. Definicija III raboty 22 davaje symbolu $G\mid A$ točny smysl: berut se klasy po $A$, sostojęče iz elementov $G$. Ne ide o proizvoljny seznam vybranyh polinomov. Arhiv predpočitaje pridavnu konstrukciju radi jeje prirodnosti v definiciji i svezanosti s ranějšim klasy ostatkov; to jest kvalitativna jezyčna ocena, ne čitateljsky eksperiment. Dalje tekst izrično razlikuje čislo klasov nad cělymi čislami od čisla linearno nezavisnyh klasov v polinomnom slučaju. Slovnik mora so slovami prěnesti i tuto razliku.\par
\noindent\emph{Reason:} Ostatkova klasa names one class, whereas sistem ostatkovyh klasov names the specified system. Paper 22's Definition III assigns $G\mid A$ the precise meaning of the classes modulo $A$ consisting of elements of $G$, not an arbitrary list of selected polynomials. The archive favours the adjectival construction as natural in the definition and continuous with klasy ostatkov; that is a qualitative linguistic assessment, not a reader experiment. The following text expressly distinguishes the number of classes over the integers from the number of linearly independent classes in the polynomial case. The dictionary must carry that distinction with the vocabulary.\par
\noindent{\footnotesize Izvorne zapisy / Original records: HIST-g086-002, HIST-g092-012.}\par
\smallskip\noindent\emph{Sravnjenje s tekstom:} Definicija razlikuje klasu, sistem klasov i izbranogo prědstavitelja. Dalši tekst razlikuje kardinalnost od linearnoj dimenzije.\par
\noindent\emph{Textual comparison:} The definition distinguishes a class, a system of classes and a representative. The subsequent text distinguishes cardinality from linear dimension.\par
\noindent{\footnotesize TEXT-COMP-607; pasaže / passages: 14404, 14414; \texttt{CORPUS\_FORM\_COMPARISONS.json}.}\par
\end{minipage}\par
\par\medskip\Needspace{7\baselineskip}
\noindent\begin{minipage}{\linewidth}
\hypertarget{TERM-1552}{}
\noindent\textbf{klasy ostatkov}\quad{\footnotesize\texttt{TERM-1552}}\par
\noindent\emph{Iztočnik / Source:} Restklassen\par
\noindent\emph{Rusky / Russian:} классы вычетов\par
\noindent\emph{Ukrajinsky / Ukrainian:} класи лишків\par
\smallskip\noindent\emph{Razlog:} Klasy ostatkov čuva korenj ostatok v dvuh arhivnyh kontekstah: modulnom tolkovanju raboty 19 i teoremu o redukciji po prostom idealu v rabotě 20. Oba razloga upominajut meždujezyčnu prozračnost, ale ne davajut čitateljsky test ili čislove vagi. Nynějšna pasaža raboty 20 upotrěbljaje polje klasov modulo označenogo prostogo ideala pri izključenju konečno mnogyh idealov. To ustanovjava konkretnu lokalnu upotrěbu, ne teorem, že faktor po libovoljnom prostom idealu libovoljnogo prstena vsegda jest polje. Ruske formy vyčet i ostatok i ukrajinsky lišok ne trěba svoditi k jednomu korenju radi vněšnej jednakosti.\par
\noindent\emph{Reason:} Klasy ostatkov retains ostatok in two recorded settings: Paper 19's module interpretation and Paper 20's reduction modulo a prime ideal. Both rationales invoke cross-language transparency but provide no reader test or numerical weights. The current Paper 20 passage uses the residue field modulo the specified prime ideal while excluding finitely many ideals. This documents that local use, not a theorem that quotienting any ring by any prime ideal always produces a field. Russian vyčet and ostatok and Ukrainian lišok need not be reduced to one common stem for superficial uniformity.\par
\noindent{\footnotesize Izvorne zapisy / Original records: HIST-g081-094, HIST-g083-009.}\par
\smallskip\noindent\emph{Sravnjenje s tekstom:} Polje klasov modulo prostogo ideala služi v konkretnom teoremu s konečnym čislom izključenij.\par
\noindent\emph{Textual comparison:} The residue field modulo a prime ideal is used within a particular theorem with finitely many exceptions.\par
\noindent{\footnotesize TEXT-COMP-609; pasaže / passages: 14022; \texttt{CORPUS\_FORM\_COMPARISONS.json}.}\par
\end{minipage}\par
\par\medskip\Needspace{7\baselineskip}
\noindent\begin{minipage}{\linewidth}
\hypertarget{TERM-1553}{}
\noindent\textbf{klasy ostatkov / grupy ostatkov}\quad{\footnotesize\texttt{TERM-1553}}\par
\noindent\emph{Iztočnik / Source:} Restklassen / Restgruppen\par
\noindent\emph{Rusky / Russian:} классы остатков / группы остатков\par
\noindent\emph{Ukrajinsky / Ukrainian:} класи лишків / групи лишків\par
\smallskip\noindent\emph{Razlog:} Klasy ostatkov i grupy ostatkov morajut ostati različne po nivoju objekta. V rabotě 17 klasa jest cělost polinomov kongruentnyh danomu polinomu; grupa sobira te klasy s dodavanjem i dejstvom polinomov. Jednostrannost modula znači, že produkt dvoh proizvoljnyh klasov ne mora byti dobro oprěděljen. Zato grupa ne može bez dalnogo uslovja stati prstenom. Arhivny razlog čuva ruski i medžuslovjansky korenj ostatok, ukrajinsky lišok, i ne zatajuje različnyh matematičnyh struktur. Ta razlika jest potrěbna takože v poučnom tekstu: čitatelj mora viděti, katere operacije sut v definiciji dopuščene.\par
\noindent\emph{Reason:} Klasy ostatkov and grupy ostatkov must remain distinct levels of object. In Paper 17 a class is the collection of polynomials congruent to a given polynomial; the group assembles those classes with addition and polynomial action. Because the module is one-sided, multiplication of two arbitrary classes need not be well-defined. The group therefore cannot become a ring without an additional condition. The archival rationale preserves Russian and Interslavic ostatok beside Ukrainian lišok without suppressing the different mathematical structures. An educational text likewise needs to identify which operations its definition permits.\par
\noindent{\footnotesize Izvorne zapisy / Original records: HIST-g079-005.}\par
\smallskip\noindent\emph{Sravnjenje s tekstom:} Jednostrannost dopušča dodavanje klasov i dejstvo polinomov, ale ne obči produkt dvoh klasov.\par
\noindent\emph{Textual comparison:} One-sidedness permits class addition and polynomial action, but not a general product of two classes.\par
\noindent{\footnotesize TEXT-COMP-610; pasaže / passages: 10957, 11098, 11102, 11108, 11110, 11112, 11122; \texttt{CORPUS\_FORM\_COMPARISONS.json}.}\par
\end{minipage}\par
\par\medskip\Needspace{7\baselineskip}
\noindent\begin{minipage}{\linewidth}
\hypertarget{TERM-1554}{}
\noindent\textbf{ostatkove klase modulo $p$}\quad{\footnotesize\texttt{TERM-1554}}\par
\noindent\emph{Iztočnik / Source:} Restklassen modulo $p$\par
\noindent\emph{Rusky / Russian:} классы вычетов по модулю $p$\par
\noindent\emph{Ukrajinsky / Ukrainian:} класи лишків за модулем $p$\par
\smallskip\noindent\emph{Razlog:} Ostatkove klase modulo $p$ spojuje slovjansku osnovu s internacionalnym modulo. Arhiv sravnjaje rusku i ukrajinsku prědložnu formu, ne nudi novu matematičnu operaciju. V rabotě 35 prěhod jest zaměna cěločislnyh koeficientov jih klasami po prostom čislu $p$. Čisla v imenovateljah davajut izključenja: bez jih prověrky redukcija ne mora byti definovana. To uslovje i dalne prědpoloženje teorema morajut ostati v tekstu. Pojem ne jest proizvoljno čislovo zaokrugljenje. Prosty symbol $p$ iz historičnogo zaglavja ostavaje istim identifikatorom pri matematičnom tipografičnom prikazu.\par
\noindent\emph{Reason:} Ostatkove klase modulo $p$ combines the Slavic base with international modulo. The archive compares Russian and Ukrainian prepositional forms; it does not propose a new mathematical operation. In Paper 35 the passage replaces integer coefficients by their classes modulo a prime number $p$. Primes occurring in denominators are excluded: reduction need not be defined without that check. Those exclusions and the theorem's further hypothesis must remain in the text. This is not numerical rounding. The plain p of the historical heading remains the same identifier when typeset mathematically.\par
\noindent{\footnotesize Izvorne zapisy / Original records: HIST-g205-006.}\par
\smallskip\noindent\emph{Sravnjenje s tekstom:} Redukcija zaměnjaje koeficienty klasami po p; izključenja iz imenovateljev i dodatne uslovja ostavajut potrěbne.\par
\noindent\emph{Textual comparison:} Reduction replaces coefficients by classes modulo p; denominator exclusions and further hypotheses remain necessary.\par
\noindent{\footnotesize TEXT-COMP-611; pasaže / passages: 22535, 22569, 22599; \texttt{CORPUS\_FORM\_COMPARISONS.json}.}\par
\end{minipage}\par
\par\medskip\Needspace{7\baselineskip}
\noindent\begin{minipage}{\linewidth}
\hypertarget{TERM-1555}{}
\noindent\textbf{klasy ostatkov; normovanje klasov ostatkov; normovany polinom}\quad{\footnotesize\texttt{TERM-1555}}\par
\noindent\emph{Iztočnik / Source:} Restklassen; Normierung der Restklassen; normiertes Polynom\par
\noindent\emph{Rusky / Russian:} классы вычетов; нормирование классов вычетов; нормированный многочлен\par
\noindent\emph{Ukrajinsky / Ukrainian:} класи лишків; нормування класів лишків; нормований поліном\par
\smallskip\noindent\emph{Razlog:} Klasy ostatkov, normovanje i normovany polinom sobirajut konstrukciju iz raboty 15, ne tri slobodno zaměnljive synonimy. Dla četyreh redov normovana forma ne sodrživaje produkt $\xi_{14}\xi_{23}$; kongruencija zaměnjaje jeho drugymi členami. V občem slučaju tekst davaje ograničenje vloženyh parov indeksov i dokazuje jedinoznačnost normovanogo prědstavitelja. Normovany tu zato ne znači avtomatično moničny polinom ili vektor jednotkovej dolžiny. Arhiv opiraje se na znanu modularnu slovnu familiju i internacionalne normovanje/normovany, bez kvantitativnogo rangovanja. Pri OLP nazvanje normalnoj formy trěbuje svoju logičnu definiciju i svoj dokaz, ne prěnos tutogo redukcijnogo pravila.\par
\noindent\emph{Reason:} Klasy ostatkov, normovanje and normovany polinom collect a construction from Paper 15, not three interchangeable synonyms. For four rows, a normalized form excludes the product $\xi_{14}\xi_{23}$, replacing it by other terms through a congruence. In general the text imposes a restriction on nested index pairs and proves uniqueness of the normalized representative. Normalized therefore does not automatically mean monic or of unit vector length. The archive relies on familiar residue terminology and international normovanje/normovany without quantitative candidate ranking. An OLP normal form requires its own logical definition and proof, not transfer of this reduction rule.\par
\noindent{\footnotesize Izvorne zapisy / Original records: HIST-g077-010.}\par
\smallskip\noindent\emph{Sravnjenje s tekstom:} Normovanje izključuje označene produkty indeksov i daje jedinogo prědstavitelja; ne znači prosto moničnost.\par
\noindent\emph{Textual comparison:} Normalization excludes specified indexed products and yields a unique representative; it does not simply mean monicity.\par
\noindent{\footnotesize TEXT-COMP-612; pasaže / passages: 10440, 10442, 10444, 10446, 10495, 10514; \texttt{CORPUS\_FORM\_COMPARISONS.json}.}\par
\end{minipage}\par
\par\medskip\Needspace{7\baselineskip}
\noindent\begin{minipage}{\linewidth}
\hypertarget{TERM-1556}{}
\noindent\textbf{oblast/kolco klasov ostatkov}\quad{\footnotesize\texttt{TERM-1556}}\par
\noindent\emph{Iztočnik / Source:} Restklassenbereich / Restklassenring\par
\noindent\emph{Rusky / Russian:} область/кольцо классов вычетов\par
\noindent\emph{Ukrajinsky / Ukrainian:} область/кільце класів лишків\par
\smallskip\noindent\emph{Razlog:} Oblast/kolco klasov ostatkov zachovuje historičnu slovnu familiju, vměsto kračej sovrěmennej faktor-prsten. U raboty 34, §8, oblast $\mathfrak o/\mathfrak a$ jest prsten zato, že $\mathfrak a$ jest dvustranny ideal; samo modulno svojstvo by ne dostačalo. Arhiv čuva nazvanje konstrukcije, ale ne tvrdi, že oblast vsegda znači oblast bez nulovyh děliteljev. Nynějšny korpus piše prsten, ne arhivno kolco. Slovnik ukazuje oba zapisy i jih odnos k tomu samomu kvocientnomu objektu; to ne jest dovoljenje slepo vrnuti ranějšny korenj do novogo teksta. Prěd izborem trěba sravniti nynějšny kanon i kontekst.\par
\noindent\emph{Reason:} Oblast/kolco klasov ostatkov preserves the historical lexical family rather than compressing it into the modern factor-ring label. In Paper 34, §8, $\mathfrak o/\mathfrak a$ is a ring because $\mathfrak a$ is a two-sided ideal; the module property alone would not suffice. Retaining oblast does not imply an integral domain. The current corpus uses prsten rather than archival kolco. The dictionary shows both records and their relation to the quotient construction; it does not authorize blindly reintroducing the older stem. Current canon and context must be checked before reuse.\par
\noindent{\footnotesize Izvorne zapisy / Original records: HIST-g186-002.}\par
\smallskip\noindent\emph{Sravnjenje s tekstom:} Dvustrannost ideala jest izrično uslovje prstenovoj struktury. Nynějšny termin jest prsten.\par
\noindent\emph{Textual comparison:} The ideal's two-sidedness is an explicit condition for ring structure. The current term is prsten.\par
\noindent{\footnotesize TEXT-COMP-613; pasaže / passages: 20402, 20404; \texttt{CORPUS\_FORM\_COMPARISONS.json}.}\par
\end{minipage}\par
\par\medskip\Needspace{7\baselineskip}
\noindent\begin{minipage}{\linewidth}
\hypertarget{TERM-1557}{}
\noindent\textbf{polje ostatkovyh klasov}\quad{\footnotesize\texttt{TERM-1557}}\par
\noindent\emph{Iztočnik / Source:} Restklassenkörper\par
\noindent\emph{Rusky / Russian:} поле классов вычетов\par
\noindent\emph{Ukrajinsky / Ukrainian:} поле класів лишків\par
\smallskip\noindent\emph{Razlog:} Polje ostatkovyh klasov jest povezana konstrukcija, vračajuča se v rabotah 24, 25 i 27. Zachovane razloženja predpočitaju prozračnu frazu novomu jednoslovnomu neologizmu i čuvaju kontinuitet medžu sekcijami. Ale pojavenje slova polje trěbuje matematično uslovje: lema v rabotě 24 trěbuje konečnu dimenziju sistemu klasov po prostom idealu nad koeficientnym poljem. Rabota 25 konstruuje polje s jednim korenjem nerazložimogo polinoma; polje vseh korenjev nastaje povtorjenjem konstrukcije. Rabota 27 broji linearno nezavisne klasy po odnošenju k označenomu polju. Te priměry ne davajut poistovětiti klasu, ostatkovy prsten, polje s jednim korenjem i razpadno polje.\par
\noindent\emph{Reason:} Polje ostatkovyh klasov is a connected construction recurring in Papers 24, 25 and 27. The preserved rationales prefer an explicit phrase to an opaque one-word neologism and maintain continuity between sections. But field requires mathematical hypotheses: Paper 24's lemma assumes finite dimension of the residue system modulo a prime ideal over the coefficient field. Paper 25 constructs a field containing one root of an irreducible polynomial; repeated construction produces the field containing all roots. Paper 27 counts linearly independent classes relative to the specified field. These uses do not identify a class, a residue ring, a one-root extension and a splitting field.\par
\noindent{\footnotesize Izvorne zapisy / Original records: HIST-g095-005, HIST-g097-002, HIST-g101-006, HIST-g102-003, HIST-g104-006.}\par
\smallskip\noindent\emph{Sravnjenje s tekstom:} Konečna linearna dimenzija, izključenja pri redukciji i različje medžu jednim korenjem i vsemi korenjami morajut ostati točne.\par
\noindent\emph{Textual comparison:} Finite linear dimension, exclusions in reduction and the distinction between one root and all roots must remain explicit.\par
\noindent{\footnotesize TEXT-COMP-614; pasaže / passages: 15414, 15416, 15420, 15739, 15745, 15822, 15824, 15917; \texttt{CORPUS\_FORM\_COMPARISONS.json}.}\par
\end{minipage}\par
\par\medskip\Needspace{7\baselineskip}
\noindent\begin{minipage}{\linewidth}
\hypertarget{TERM-1558}{}
\noindent\textbf{polje klasov ostatkov / sovršeno polje}\quad{\footnotesize\texttt{TERM-1558}}\par
\noindent\emph{Iztočnik / Source:} Restklassenkörper / vollkommener Körper\par
\noindent\emph{Rusky / Russian:} поле классов вычетов / совершенное поле\par
\noindent\emph{Ukrajinsky / Ukrainian:} поле класів лишків / досконале поле\par
\smallskip\noindent\emph{Razlog:} Polje klasov ostatkov i sovršeno polje sut dva različne opisy: prvy govori o konstrukciji, drugi o dodatnoj vlastnosti polja. Arhiv prodolžuje ranějši izbor sovršeno, ale nynějšna pasaža raboty 31 piše doskonalo. Ta razlika ne směje byti skryta kakoby vybrana forma vsegda ostala nezměnjena. V razgledanom poslědku doskonalost osnovnogo ostatkovogo polja odstranjaje slučaj prostogo ideala drugogo roda iz kriterija diskriminanta; ne zaměnjaje sam kriterij ani ne znači neformalnu pohvalu polja. Zachovana je kontinuitetna motivacija, ne ocena novogo kalka ili čislovo hlasovanje o dvuh pridavnyh imenah.\par
\noindent\emph{Reason:} Polje klasov ostatkov and sovršeno polje describe different things: the first a construction, the second an additional property of the field. The archive continues the earlier choice sovršeno, but the current Paper 31 passage reads doskonalo. The dictionary must not conceal that difference by implying that the recorded choice remained unchanged. In the inspected consequence, perfectness of the base residue field removes the second-kind prime-ideal case from the discriminant criterion; it neither replaces the criterion nor expresses informal praise. The surviving reason is continuity, not a fresh evaluation of a calque or a numerical vote between adjectives.\par
\noindent{\footnotesize Izvorne zapisy / Original records: HIST-g154-001.}\par
\smallskip\noindent\emph{Sravnjenje s tekstom:} Pasaža upotrěbljaje doskonalo i izrično ograničuje diskriminantny poslědok; arhivny izbor sovršeno ostavaje odlišeny.\par
\noindent\emph{Textual comparison:} The passage uses doskonalo and explicitly qualifies the discriminant consequence; archival sovršeno remains distinguishable.\par
\noindent{\footnotesize TEXT-COMP-615; pasaže / passages: 18749, 18751; \texttt{CORPUS\_FORM\_COMPARISONS.json}.}\par
\end{minipage}\par
\par\medskip\Needspace{7\baselineskip}
\noindent\begin{minipage}{\linewidth}
\hypertarget{TERM-1559}{}
\noindent\textbf{modul ostatkovyh klasov}\quad{\footnotesize\texttt{TERM-1559}}\par
\noindent\emph{Iztočnik / Source:} Restklassenmodul\par
\noindent\emph{Rusky / Russian:} модуль классов вычетов\par
\noindent\emph{Ukrajinsky / Ukrainian:} модуль класів лишків\par
\smallskip\noindent\emph{Razlog:} Modul ostatkovyh klasov prodolžuje familiju ostatk- v zaključenju raboty 30. Tamo prěhod k podmodulu i prěhod k modulu klasov ostatkov pomogajut povezati kompozicijny red s dvojnym uslovjem cěpov. Trěba čuvati oba směry i uslovja, ne objaviti konečnu kardinalnost modula zato, že red imaje konečnu dolgost. Arhiv daje jezyčny razlog paralelnosti s ranějšimi terminami za prsten i modul klasov. Ne davaje nezavisny čitateljsky test prozračnosti. Nazvanje objasnjava kvocientnu konstrukciju, ale samo po sobě ne dokazuje prěryvanja cěpov u proizvoljnogo modula.\par
\noindent\emph{Reason:} Modul ostatkovyh klasov continues the ostatk- family in Paper 30's conclusion. Passage to a submodule and to the residue-class module is used there to connect a composition series with the double chain condition. Both directions and their hypotheses matter; finite length of the series does not imply finite cardinality of the module. The archive gives a linguistic reason for keeping module and ring terminology parallel, not an independent reader test of transparency. The phrase identifies a quotient construction but does not prove termination of chains for an arbitrary module.\par
\noindent{\footnotesize Izvorne zapisy / Original records: HIST-g117-005.}\par
\smallskip\noindent\emph{Sravnjenje s tekstom:} Zaključenje čuva prěhody k podmodulu i k modulu klasov; konečna dolgost ne jest konečna kardinalnost.\par
\noindent\emph{Textual comparison:} The conclusion retains passage to a submodule and to the residue-class module; finite length is not finite cardinality.\par
\noindent{\footnotesize TEXT-COMP-616; pasaže / passages: 17193, 17197, 17199, 17203, 17207, 17230; \texttt{CORPUS\_FORM\_COMPARISONS.json}.}\par
\end{minipage}\par
\par\medskip\Needspace{7\baselineskip}
\noindent\begin{minipage}{\linewidth}
\hypertarget{TERM-1560}{}
\noindent\textbf{modul klasov ostatkov / kolco klasov ostatkov / ostatkova klasa}\quad{\footnotesize\texttt{TERM-1560}}\par
\noindent\emph{Iztočnik / Source:} Restklassenmodul / Restklassenring / Restklasse\par
\noindent\emph{Rusky / Russian:} модуль классов вычетов / кольцо классов вычетов / класс вычетов\par
\noindent\emph{Ukrajinsky / Ukrainian:} модуль класів лишків / кільце класів лишків / клас лишків\par
\smallskip\noindent\emph{Razlog:} Trojica modul klasov ostatkov, kolco klasov ostatkov i ostatkova klasa čuva tri nivoja opisa. §4 raboty 30 najprvo sobira kongruentne elementy v jednu klasu i daje modulnu strukturu sistemu klasov. Prstenova struktura dodatno trěbuje čuvanje produktov; pri nekomutativnyh prstenah primětka trěbuje dvustranne idealy. Zato izomorfizm modulov ne smi avtomatično stati izomorfizmom prstenov. Arhivny korenj kolco nynějšny tekst zaměnjaje formoju prsten, medžutym ostatk- ostavaje obča familija. Pridavna forma za jednu klasu i genitivna za sistem mogut služiti različnym gramatičnym roljam; ta razlika ne jest matematična nepodudarnost.\par
\noindent\emph{Reason:} The trio modul klasov ostatkov, kolco klasov ostatkov and ostatkova klasa preserves three descriptive levels. Paper 30, §4 first groups congruent elements into a class and equips the system with module structure. Ring structure additionally requires preservation of products; for noncommutative rings the footnote requires two-sided ideals. A module isomorphism must therefore not automatically become a ring isomorphism. Current prsten replaces archival kolco while the ostatk- family persists. An adjective for a single class and a genitive construction for the system can serve different grammatical roles; that difference is not a mathematical disagreement.\par
\noindent{\footnotesize Izvorne zapisy / Original records: HIST-g111-003.}\par
\smallskip\noindent\emph{Sravnjenje s tekstom:} Modulny izomorfizm čuva modulne operacije; prstenovy dodatno čuva produkty. Dvustranne idealy sut potrěbne v nekomutativnom slučaju.\par
\noindent\emph{Textual comparison:} Module isomorphism preserves module operations; ring isomorphism also preserves products. Two-sided ideals are required in the noncommutative case.\par
\noindent{\footnotesize TEXT-COMP-617; pasaže / passages: 16460, 16462, 16464, 16466, 16480, 16482; \texttt{CORPUS\_FORM\_COMPARISONS.json}.}\par
\end{minipage}\par
\par\medskip\Needspace{7\baselineskip}
\noindent\begin{minipage}{\linewidth}
\hypertarget{TERM-1561}{}
\noindent\textbf{kolco klasov ostatkov}\quad{\footnotesize\texttt{TERM-1561}}\par
\noindent\emph{Iztočnik / Source:} Restklassenring\par
\noindent\emph{Rusky / Russian:} кольцо классов вычетов\par
\noindent\emph{Ukrajinsky / Ukrainian:} кільце класів лишків\par
\smallskip\noindent\emph{Razlog:} Kolco klasov ostatkov jest arhivny naziv za prsten, čijimi elementami sut kongruentne klasy, ne izbrane prědstavitelje. Šest zapisov iz rabot 30 i 31 daje prědnost prozrěčnoj frazi i kontinuitetu: čitatelj vidi i klasy i prstenovu strukturu. Kratše faktor-kolco ne bylo izključeno iz cělogo korpusa; razvažanje tu odbivaje nemotivovanu zaměnu naziva v srědině dokaza. Razlikovanje od polja kvocientov i od kvocienta idealov čuva različne konstrukcije. Nynějšny tekst už upotrěbjaje prsten klasov ostatkov, zato arhivno kolco ne jest nova rekomendacija. Slovnik pokazuje oba stanja i matematicky smysl, po ktorom trěba prověrjati kontinuitet. Šest zapisov jednej tradicije ne sut šest nezavisnyh potvrđenij jezyčnoj pravilnosti.\par
\noindent\emph{Reason:} Kolco klasov ostatkov is the archived name for a ring whose elements are congruence classes, not selected representatives. Six records from Papers 30 and 31 favour a transparent phrase and continuity: both the classes and the ring structure remain visible. The shorter faktor-kolco was not excluded from the entire corpus; the stated objection here is an unmotivated change of name within an argument. Distinguishing this construction from a field of fractions and an ideal quotient protects different mathematical objects. The current text uses prsten klasov ostatkov, so the archived kolco is not a new recommendation. Both forms are documented against their mathematical use. Six records within one translation tradition are not six independent confirmations of linguistic correctness.\par
\noindent{\footnotesize Izvorne zapisy / Original records: HIST-g107-011, HIST-g120-005, HIST-g124-002, HIST-g127-005, HIST-g131-005, HIST-g151-004.}\par
\smallskip\noindent\emph{Sravnjenje s tekstom:} Različne raboty čuvajut prstenovu strukturu klasov. Arhivno kolco ne zaměnjuje nynějšno prsten.\par
\noindent\emph{Textual comparison:} Different papers retain the ring structure of classes. Archived kolco does not displace current prsten.\par
\noindent{\footnotesize TEXT-COMP-618; pasaže / passages: 16141, 17363, 17398, 17595, 17704, 17874, 18627, 18635; \texttt{CORPUS\_FORM\_COMPARISONS.json}.}\par
\end{minipage}\par
\par\medskip\Needspace{7\baselineskip}
\noindent\begin{minipage}{\linewidth}
\hypertarget{TERM-1562}{}
\noindent\textbf{kolco klasov ostatkov / polje klasov ostatkov}\quad{\footnotesize\texttt{TERM-1562}}\par
\noindent\emph{Iztočnik / Source:} Restklassenring / Restklassenkörper\par
\noindent\emph{Rusky / Russian:} кольцо классов вычетов / поле классов вычетов\par
\noindent\emph{Ukrajinsky / Ukrainian:} кільце класів лишків / поле класів лишків\par
\smallskip\noindent\emph{Razlog:} V uvodu raboty 31 zajedny izraz klasov ostatkov svyzuje prsten i polje, ale ne izravnjuje jih. Polje imaje dodatnu strukturu: nenuleve elementy dopuščajut děljenje. Tekst razlikuje polja prvogo i drugogo roda pri nedoskonalom osnovnom polju; ta razlika nestavaje zaměnoju naziva prstena nazvom polja. Dalje je izrično rečeno, že prsten iz polinomnogo prědstavjenja ne mora byti izomorfny prstenu ostatkovyh klasov poredka. Arhivna motivacija zato jest čitlivost zajednoj terminologičnoj rodiny, ne tvrdženje, že vsi takovi prsteni sut polja ili medžu soboju izomorfni. Kratši kvocientny naziv ne směl by skryti, po ktorom idealu i v ktoroj oblasti se tvori kvocient. Nynějšny prsten zaměnjuje arhivno kolco.\par
\noindent\emph{Reason:} In Paper 31's introduction, the shared phrase klasov ostatkov links ring and field without identifying them. A field has additional structure: its nonzero elements admit division. The text distinguishes extensions of the first and second kind over an imperfect base field; replacing “ring” with “field” would not express that distinction. It also explicitly warns that a ring obtained from a polynomial presentation need not be isomorphic to the order's residue ring. The recorded motive is therefore a readable terminology family, not a claim that all these rings are fields or mutually isomorphic. A shorter quotient label must still identify the ideal and ambient structure. Current prsten replaces archived kolco.\par
\noindent{\footnotesize Izvorne zapisy / Original records: HIST-g118-006.}\par
\smallskip\noindent\emph{Sravnjenje s tekstom:} Nedoskonalo polje i neizomorfne polinomne prědstavjenja trěbujut različnih hipotez.\par
\noindent\emph{Textual comparison:} Imperfect fields and nonisomorphic polynomial presentations require distinct hypotheses.\par
\noindent{\footnotesize TEXT-COMP-619; pasaže / passages: 17275, 17281; \texttt{CORPUS\_FORM\_COMPARISONS.json}.}\par
\end{minipage}\par
\par\medskip\Needspace{7\baselineskip}
\noindent\begin{minipage}{\linewidth}
\hypertarget{TERM-1563}{}
\noindent\textbf{kolco klasov ostatkov modulo $p^t$}\quad{\footnotesize\texttt{TERM-1563}}\par
\noindent\emph{Iztočnik / Source:} Restklassenring modulo $p^t$\par
\noindent\emph{Rusky / Russian:} кольцо классов вычетов modulo $p^t$\par
\noindent\emph{Ukrajinsky / Ukrainian:} кільце класів лишків modulo $p^t$\par
\smallskip\noindent\emph{Razlog:} Modifikator modulo $p^t$ jest čestju ponjatja: rabota 43 razgledaje prsten po stepenu, ne samo po prostom elementu. Arhiv čuva međunarodny izraz modulo i cěly eksponent, vměsto skračenja do polja po $p$. V §7 diferenta glavnogo poredka prěhodi v diferentu prstena klasov ostatkov; pri razloženju $p$ na proste idealne faktory pojavjajut se eksponenty $te_i$. Zato zaměna $p^t$ na $p$ izměnila by samu konstrukciju. Za točno oprěděljenje eksponentov uvod dodatno potrěbuje dostatočno veliko $t$. To, že $p$ jest prosty element osnovnoj oblasti, samo po sobě ne daje polje pri vzjatju klasov po jego stepenu. Arhivno kolco ostavaje istorijnoju formoju, nynějšno prsten jest vidimo v tekstu. Matematične bukvy i eksponenty se pokazyvajut kako formulu, ne kako obyčna jezyčna slova.\par
\noindent\emph{Reason:} The modifier modulo $p^t$ is part of the concept: Paper 43 considers a quotient by a power, not merely by a prime element. The archive retains international modulo and the whole exponent instead of abbreviating the construction to a field modulo $p$. In Section 7 the different of the principal order passes to the different of the residue ring; factorization of $p$ into prime ideals introduces exponents $te_i$. Replacing $p^t$ by $p$ would therefore alter the construction. For exact determination of the exponents, the introduction additionally requires sufficiently large $t$. Primality of the base element alone does not make the quotient by its power a field. Kolco remains the historical form and prsten the current one; mathematical letters and exponents are typeset as notation rather than prose.\par
\noindent{\footnotesize Izvorne zapisy / Original records: HIST-g213-012.}\par
\smallskip\noindent\emph{Sravnjenje s tekstom:} Eksponenty ostavajut čestju konstrukcije i odgovarjajučih komponentov.\par
\noindent\emph{Textual comparison:} Exponents remain part of the construction and its corresponding components.\par
\noindent{\footnotesize TEXT-COMP-620; pasaže / passages: 24773, 25122, 25124, 25126, 25128, 25132, 25133, 25135, 25137, 25138, 25140; \texttt{CORPUS\_FORM\_COMPARISONS.json}.}\par
\end{minipage}\par
\par\medskip\Needspace{7\baselineskip}
\noindent\begin{minipage}{\linewidth}
\hypertarget{TERM-1564}{}
\noindent\textbf{kolco klasov ostatkov po radikalu}\quad{\footnotesize\texttt{TERM-1564}}\par
\noindent\emph{Iztočnik / Source:} Restklassenring nach dem Radikal\par
\noindent\emph{Rusky / Russian:} кольцо классов вычетов по радикалу\par
\noindent\emph{Ukrajinsky / Ukrainian:} кільце класів лишків за радикалом\par
\smallskip\noindent\emph{Razlog:} Po radikalu označaje ideal, po ktorom se tvori prsten klasov ostatkov. Arhiv daje prědnost razvitoj frazi prěd novym kratkym faktor-kolco, aby ostala jasna svyza s rabotami 30 i 31. V §12 raboty 34 se pod uslovjem maksimuma za prave idealy konstruuje največši nilpotentny dvustranny ideal; kvocient po njem jest bez radikala v tutom smyslu. Dvustrannost potrěbna jest za prstenovu strukturu kvocienta. Racionala zato ne podupiraje neograničeno prěnosenje teorema na libovoljny radikal ili libovoljny nekomutativny prsten. Čitatelj potrěbuje i naziv konstrukcije i hipotezu, pod ktoroj avtor dobiva svoj rezultat. Istorijno kolco i nynějšno prsten sut različne jezyčne formy toho samogo lokalno oprěděljenogo objekta.\par
\noindent\emph{Reason:} Po radikalu identifies the ideal by which the residue ring is formed. The archive prefers the expanded phrase to a newly introduced faktor-kolco to retain continuity with Papers 30 and 31. In Section 12 of Paper 34, a maximum condition on right ideals yields a largest nilpotent two-sided ideal; the quotient is radical-free in the stated sense. Two-sidedness matters to the quotient's ring structure. This rationale does not support transferring the theorem to every notion of radical or every noncommutative ring. Readers need both the construction's name and the hypothesis under which the result is established. Historical kolco and current prsten denote this same locally defined object.\par
\noindent{\footnotesize Izvorne zapisy / Original records: HIST-g190-005.}\par
\smallskip\noindent\emph{Sravnjenje s tekstom:} Uslovje maksimuma i dvustrannost ideala ne mogut izčeznuti iz objašnjenja.\par
\noindent\emph{Textual comparison:} The maximum condition and the ideal's two-sidedness cannot disappear from the explanation.\par
\noindent{\footnotesize TEXT-COMP-621; pasaže / passages: 20860, 20862, 20864, 20866; \texttt{CORPUS\_FORM\_COMPARISONS.json}.}\par
\end{minipage}\par
\par\medskip\Needspace{7\baselineskip}
\noindent\begin{minipage}{\linewidth}
\hypertarget{TERM-1565}{}
\noindent\textbf{kolco klasov ostatkov bez děliteljev nuly / dělitelj nuly}\quad{\footnotesize\texttt{TERM-1565}}\par
\noindent\emph{Iztočnik / Source:} Restklassenring ohne Nullteiler / Nullteiler\par
\noindent\emph{Rusky / Russian:} кольцо классов вычетов без делителей нуля / делитель нуля\par
\noindent\emph{Ukrajinsky / Ukrainian:} кільце класів лишків без дільників нуля / дільник нуля\par
\smallskip\noindent\emph{Razlog:} Dělitelj nuly jest konkretna sostavna fraza; bez děliteljev nuly čuva negaciju cělogo svojstva. Rabota 30, §5, razvažaje komutativny prsten bez dodatnoj prědpoložki o jedinici i razlikuje dělitelje nuly medžu elementami od idealov-děliteljev nuly. Ona potom dokazyvaje ekvivalenciju dvuh definicij prostogo ideala. Termin ne směje prědstaviti tutu dokazanu ekvivalenciju kako početnu zaměnimost elementa i ideala. V nynějšnem zapisu obe prědpoložki upotrěbjajut te same frakturne simvoly, hoti tekst razlikuje tipy objektov; toj detalj trěbuje sravnjanja s prvotnym svědkom, ne tihu popravku. Arhivni izbor slov vysvětljaje smysl bez izmišljanja takogo sravnjanja. Kolco klasov ostatkov jest istorijna forma nynějšnogo prsten klasov ostatkov.\par
\noindent\emph{Reason:} Dělitelj nuly is a concrete compound expression; bez děliteljev nuly preserves negation of the whole property. Paper 30, Section 5, considers a commutative ring without additionally assuming an identity and distinguishes zero divisors among elements from zero-divisor ideals. It then proves equivalence of two definitions of prime ideal. Terminology must not present that proved equivalence as an initial interchangeability of elements and ideals. The current formulas use the same fraktur symbols in both conditions, although the prose distinguishes object types; this detail requires comparison with the primary witness, not silent repair. The lexical account does not pretend that such comparison has occurred. Kolco klasov ostatkov is the historical form of current prsten klasov ostatkov.\par
\noindent{\footnotesize Izvorne zapisy / Original records: HIST-g112-007.}\par
\smallskip\noindent\emph{Sravnjenje s tekstom:} Tekst razlikuje dva tipa objektov i dokazyvaje ekvivalenciju. Jednaka frakturna notacija trěbuje prověrky po prvotnom svědku.\par
\noindent\emph{Textual comparison:} The prose distinguishes two object types and proves equivalence. Identical fraktur notation needs comparison with the primary witness.\par
\noindent{\footnotesize TEXT-COMP-622; pasaže / passages: 16585, 16586, 16601, 16603, 16605; \texttt{CORPUS\_FORM\_COMPARISONS.json}.}\par
\end{minipage}\par
\par\medskip\Needspace{7\baselineskip}
\noindent\begin{minipage}{\linewidth}
\hypertarget{TERM-1566}{}
\noindent\textbf{faktor-kolco $R/p$ ranga jedan}\quad{\footnotesize\texttt{TERM-1566}}\par
\noindent\emph{Iztočnik / Source:} Restklassenring $R/p$ vom Rang eins\par
\noindent\emph{Rusky / Russian:} факторкольцо $R/p$ ранга один\par
\noindent\emph{Ukrajinsky / Ukrainian:} фактор-кільце $R/p$ рангу один\par
\smallskip\noindent\emph{Razlog:} Faktor-kolco jest tu kratka alternativa k razvitoj frazi kolco klasov ostatkov. Arhivna racionala svyzuje izbor s gustym tehničnym kontekstom i tvrdženjem o rangu; nynějšny tekst upotrěbjaje faktor-prsten. Pri algebraično zamknjenom osnovnom polju kvocient vlastno primarnogo prstena po asociovanom idealu imaje rang jedan, ale sam primarny prsten imaje prinajmenje dva bazne elementy. Rang jedan zato ne znači, že kvocient imaje samo jeden element. Argument o diskriminantu zavisi od baznyh produktov i izčezanja jih slědov. Arhiv takože tvrdí, že objašnjenje kvocienta bylo obnovjeno iz pečatnogo teksta; ovde se toj istorijni zapis čuva, ale ne vydavaje za novo nezavisno sravnjanje s pečatnym originalom.\par
\noindent\emph{Reason:} Faktor-kolco is the compact alternative here to kolco klasov ostatkov. The recorded rationale connects that choice to a dense technical passage and its rank statement; the current text uses faktor-prsten. Over the algebraically closed base field, the quotient of a properly primary ring by its associated ideal has rank one, while the primary ring itself has at least two basis elements. Rank one therefore does not mean that the quotient has only one element. The discriminant argument depends on basis products and vanishing traces. The historical note also says that the quotient explanation was restored from the printed text. That claim is preserved as an archival claim, not presented as a new independent comparison with the printed original.\par
\noindent{\footnotesize Izvorne zapisy / Original records: HIST-g143-004.}\par
\smallskip\noindent\emph{Sravnjenje s tekstom:} Kvocient imaje rang jedan, medžutym vlastno primarny prsten imaje veću bazu.\par
\noindent\emph{Textual comparison:} The quotient has rank one while the properly primary ring has a larger basis.\par
\noindent{\footnotesize TEXT-COMP-623; pasaže / passages: 18383; \texttt{CORPUS\_FORM\_COMPARISONS.json}.}\par
\end{minipage}\par
\par\medskip\Needspace{7\baselineskip}
\noindent\begin{minipage}{\linewidth}
\hypertarget{TERM-1567}{}
\noindent\textbf{kolco klasov ostatkov $T_P/T_P p$}\quad{\footnotesize\texttt{TERM-1567}}\par
\noindent\emph{Iztočnik / Source:} Restklassenring $T_P/T_P p$\par
\noindent\emph{Rusky / Russian:} кольцо классов вычетов $T_P/T_P p$\par
\noindent\emph{Ukrajinsky / Ukrainian:} кільце класів лишків $T_P/T_P p$\par
\smallskip\noindent\emph{Razlog:} Fraza s formulom $T_P/T_P p$ označaje konkretny kvocient poslje lokalizacije. Imenovatelj jest ideal $T_P p$, ne vněšny množitelj poslě děljenja $T_P$ samo soboju. V §8 raboty 31 toj prsten jest izomorfny drugomu prstenu ostatkovyh klasov; izomorfija jest svyza medžu dvoma prědstavjenjami, ne njih bukvalna identičnost. Izbor razvitoj frazy svyzuje lokalizaciju, dělitelnost diskriminanta i razloženje bez zaměny kolca za polje. Nynějšny tekst čuva prsten, a frakturne bukvy pripominajut točne tipy idealov i poredkov. Slovnik pokazuje arhivnu formulu čitljivo, no ne zaměnjaje arhivnu običnu bukvu novym simvolom i ne izmišlja drugu matematičnu operaciju.\par
\noindent\emph{Reason:} The phrase containing $T_P/T_P p$ names a particular quotient after localization. Its denominator is the ideal $T_P p$, not an external multiplier applied after dividing $T_P$ by itself. In Section 8 of Paper 31 this ring is isomorphic to another residue ring: isomorphism relates two presentations rather than asserting literal identity. The expanded phrase connects localization, discriminant divisibility and decomposition without replacing ring by field. The current text uses prsten and fraktur letters marking the exact ideals and orders. The dictionary typesets the archived formula legibly without silently substituting a different symbol or mathematical operation for the archived notation.\par
\noindent{\footnotesize Izvorne zapisy / Original records: HIST-g164-004.}\par
\smallskip\noindent\emph{Sravnjenje s tekstom:} Izomorfija dvuh prstenov klasov svyzuje prědstavjenja; ne jest bukvalna jednakost zapisov.\par
\noindent\emph{Textual comparison:} An isomorphism of residue rings relates presentations; it is not literal equality of expressions.\par
\noindent{\footnotesize TEXT-COMP-624; pasaže / passages: 19105, 19107; \texttt{CORPUS\_FORM\_COMPARISONS.json}.}\par
\end{minipage}\par
\par\medskip\Needspace{7\baselineskip}
\noindent\begin{minipage}{\linewidth}
\hypertarget{TERM-1568}{}
\noindent\textbf{sistem ostatkovyh klasov}\quad{\footnotesize\texttt{TERM-1568}}\par
\noindent\emph{Iztočnik / Source:} Restklassensystem\par
\noindent\emph{Rusky / Russian:} система классов вычетов\par
\noindent\emph{Ukrajinsky / Ukrainian:} система класів лишків\par
\smallskip\noindent\emph{Razlog:} Sistem ostatkovyh klasov čuva sistem kako nadredny izraz, kogda argument ješče ne tvrdi, že objekt jest polje. V rabotě 22, §4, beskonečny modulny opis sravnjaje se s konečnym prědstavjenjem; izomorfija čuva razlike i umnoženje koeficientami. Iz činjenice, že samo konečno mnogo elementarnyh děliteljev jest različno od jedinice, ne slěduje, že obstaje samo konečno mnogo klasov. Arhiv daje prědnost kontinuitetu s ranějšimi paragrafami i adjektivnomu ostatkovyh, ne novomu nezavisnomu nazivu. Dva zachovane zapisa dokumentujut ponovjenje izbora, ne dva nezavisne pokusy ponimanja. Pri pozdnějšem prěnosenju na logiku analogija s klasami može byti polezna, ale modulne operacije ne smějut byti pripisane libovoljnomu sistemu ekvivalencijnyh klasov.\par
\noindent\emph{Reason:} Sistem ostatkovyh klasov retains system as the broader expression where the argument has not established a field. In Paper 22, Section 4, an infinite module description is compared with a finite presentation; the isomorphism preserves differences and scalar multiplication. Having only finitely many elementary divisors different from one does not imply having finitely many classes. The archive favours continuity with earlier sections and the adjective ostatkovyh rather than an unrelated new name. Two surviving records document recurrence of the choice, not two independent comprehension experiments. For later reuse in logic, the analogy with classes may help, but module operations must not be attributed to an arbitrary system of equivalence classes.\par
\noindent{\footnotesize Izvorne zapisy / Original records: HIST-g088-002, HIST-g092-033.}\par
\smallskip\noindent\emph{Sravnjenje s tekstom:} Konečnost relevantnyh elementarnyh děliteljev ne daje konečnost množenstva klasov.\par
\noindent\emph{Textual comparison:} Finitely many relevant elementary divisors do not imply finitely many classes.\par
\noindent{\footnotesize TEXT-COMP-625; pasaže / passages: 14621, 14623, 14626, 14628, 14631, 14633; \texttt{CORPUS\_FORM\_COMPARISONS.json}.}\par
\end{minipage}\par
\par\medskip\Needspace{7\baselineskip}
\noindent\begin{minipage}{\linewidth}
\hypertarget{TERM-1569}{}
\noindent\textbf{sistem ostatkovyh klasov / polje ostatkovyh klasov}\quad{\footnotesize\texttt{TERM-1569}}\par
\noindent\emph{Iztočnik / Source:} Restklassensystem / Restklassenkörper\par
\noindent\emph{Rusky / Russian:} система классов вычетов / поле классов вычетов\par
\noindent\emph{Ukrajinsky / Ukrainian:} система класів лишків / поле класів лишків\par
\smallskip\noindent\emph{Razlog:} Sistem i polje ostatkovyh klasov tvoret svyazanu, ale ne zaměnimu paru. V rabotě 24, §6, lema o izomorfiji govori o podsistemu, a pozdnějši teorem svyzuje tvrđenja pri doskonalom polju koeficientov. Slovo polje zato nosi matematičnu obvezku, koju samo obstajanje kongruentnyh klasov ne izpolnjaje. Arhivna racionala čuva jednak korenj ostatkov- za čitlivy prěhod medžu lemoju VII i teoremom XIII; ona ne daje pravo skratiti prědpoložki ili zaměniti podsistem cělym objektom. Ruska i ukrajinska para služat sravnjajučimi formami v zachovanom zapisu. Bez novogo svědočstva one ne dokazujut ni nezavisnogo izslědovanja vsih slovjanskyh tradicij, ni toho, že svaka klasa učastvuje v tvrđenoj izomorfiji.\par
\noindent\emph{Reason:} System and field of residue classes form a related but non-interchangeable pair. In Paper 24, Section 6, an isomorphism lemma concerns a subsystem, while the later theorem relates assertions over a perfect coefficient field. Field therefore carries a mathematical commitment that the mere existence of congruence classes does not satisfy. The recorded rationale keeps the same ostatkov- stem across Lemma VII and Theorem XIII; it does not license dropping hypotheses or replacing the subsystem with the entire object. Russian and Ukrainian expressions are comparanda in the preserved record. Without further evidence they establish neither an independent survey of all Slavic traditions nor that every class participates in the stated isomorphism.\par
\noindent{\footnotesize Izvorne zapisy / Original records: HIST-g100-004.}\par
\smallskip\noindent\emph{Sravnjenje s tekstom:} Trěba čuvati podsistem, uslovje doskonalosti i smysl izomorfije.\par
\noindent\emph{Textual comparison:} The subsystem, perfectness assumption and type of isomorphism must remain explicit.\par
\noindent{\footnotesize TEXT-COMP-626; pasaže / passages: 15649, 15651, 15655, 15657, 15659, 15661, 15663, 15665; \texttt{CORPUS\_FORM\_COMPARISONS.json}.}\par
\end{minipage}\par
\par\medskip\Needspace{7\baselineskip}
\noindent\begin{minipage}{\linewidth}
\hypertarget{TERM-1570}{}
\noindent\textbf{sistem ostatkovyh klasov / kolco bez nulovyh děliteljev}\quad{\footnotesize\texttt{TERM-1570}}\par
\noindent\emph{Iztočnik / Source:} Restklassensystem / Ring ohne Nullteiler\par
\noindent\emph{Rusky / Russian:} система классов вычетов / кольцо без делителей нуля\par
\noindent\emph{Ukrajinsky / Ukrainian:} система класів лишків / кільце без дільників нуля\par
\smallskip\noindent\emph{Razlog:} Rabota 24 veže sistem ostatkovyh klasov s svojstvom prstena bez děliteljev nuly. Složenje nulovy dělitelj jest arhivna alternativa k dělitelj nuly; negacija ostavaje pri svojstvu, ne pri samom nazivu prstena. Lema v začetnoj česti čuva otsutnost děliteljev nuly pri prisojedinjenju neoprěděljenyh. Čitatelj zato potrěbuje razlikovati te elementy, polinomy nad njimi i prěhod k klasam. Otsutnost děliteljev nuly sama ne jest vsegda dostatočna za polje; kjer rabota dobiva polje iz prostogo ideala, dodatne hipotezy o konečnoj razměrnosti ne možno odstraniti. Arhiv daje prědnost kontinuitetu s rabotoju 22 i jasnomu sostavlenju slov, ne tvrđenju o ravnosti vsih tutyh struktur. Nynějšny prsten ostavaje odličen od arhivnogo kolco.\par
\noindent\emph{Reason:} Paper 24 links a system of residue classes to the property of a ring without zero divisors. Nulovy dělitelj is the archived adjectival alternative to dělitelj nuly; negation qualifies the property, not the ring's name. An early lemma preserves absence of zero divisors when indeterminates are adjoined. Readers must therefore distinguish the elements, polynomials over them and passage to classes. Absence of zero divisors alone does not always give a field; where the paper obtains a field from a prime ideal, additional finite-dimensional hypotheses cannot be discarded. The archive favours continuity with Paper 22 and transparent composition, not identity of all these structures. Current prsten remains distinct from historical kolco.\par
\noindent{\footnotesize Izvorne zapisy / Original records: HIST-g096-005.}\par
\smallskip\noindent\emph{Sravnjenje s tekstom:} Bez děliteljev nuly ne znači bez dodatnyh uslovij polje.\par
\noindent\emph{Textual comparison:} Without zero divisors does not mean field without further hypotheses.\par
\noindent{\footnotesize TEXT-COMP-627; pasaže / passages: 15374, 15376, 15378, 15380, 15382, 15414, 15416, 15420; \texttt{CORPUS\_FORM\_COMPARISONS.json}.}\par
\end{minipage}\par
\par\medskip\Needspace{7\baselineskip}
\noindent\begin{minipage}{\linewidth}
\hypertarget{TERM-1571}{}
\noindent\textbf{sistem ostatkovyh klasov prostogo čisla $p$}\quad{\footnotesize\texttt{TERM-1571}}\par
\noindent\emph{Iztočnik / Source:} Restklassensystem einer Primzahl $p$\par
\noindent\emph{Rusky / Russian:} система классов вычетов по простому числу $p$\par
\noindent\emph{Ukrajinsky / Ukrainian:} система класів лишків за простим числом $p$\par
\smallskip\noindent\emph{Razlog:} Prostogo čisla $p$ v rabotě 9, §10, oprěděljaje osnovny aritmetičny slučaj. Tekst razlikuje prosto polje tipa racionalnyh čisel od prostogo polja klasov po $p$; v drugom slučaju mnogokratne jedinice s cělymi koeficientami už daju cělo prosto polje. Izbor sistem ostatkovyh klasov ostavaje blizky imenovanoj konstrukciji, a ne nejasnomu skračenju do prostogo sistema. Prostota čisla jest sučestvena: libovoljny složeny modul ne dava tu samu poljevu strukturu. Bukva $p$ jest matematičny parametr i ne podlěgaje transliteraciji. Tej priměr objašnjuje konkretny smysl rodiny ostatkovyh klasov, ale ne opravdyvaje avtomatičnu zaměnu prsten na polje v ostalyh článkah slovnika.\par
\noindent\emph{Reason:} Prostogo čisla $p$ specifies the basic arithmetic case in Paper 9, Section 10. The text distinguishes the prime field of rational-number type from the prime field of classes modulo $p$; in the latter case integer multiples of the identity already give the whole prime field. Sistem ostatkovyh klasov stays close to that construction rather than shortening it to an ambiguous “prime system.” Primality matters: an arbitrary composite modulus does not give the same field structure. The letter $p$ is a mathematical parameter and is not transliterated. This example clarifies one instance of the residue-class family but does not justify replacing ring with field throughout the other entries.\par
\noindent{\footnotesize Izvorne zapisy / Original records: HIST-g068-005.}\par
\smallskip\noindent\emph{Sravnjenje s tekstom:} Čislo p jest prosty modul; aritmetičny slučaj ne oprěděljaje vse ostatkove prsteni.\par
\noindent\emph{Textual comparison:} The modulus p is prime; this arithmetic case does not define all residue rings.\par
\noindent{\footnotesize TEXT-COMP-628; pasaže / passages: 8350, 8351, 8353, 8354, 8355, 8356, 8357, 8358, 8359, 8360, 8361, 8362, 8363; \texttt{CORPUS\_FORM\_COMPARISONS.json}.}\par
\end{minipage}\par
\par\medskip\Needspace{7\baselineskip}
\noindent\begin{minipage}{\linewidth}
\hypertarget{TERM-1572}{}
\noindent\textbf{podpolje ostatkovyh klasov}\quad{\footnotesize\texttt{TERM-1572}}\par
\noindent\emph{Iztočnik / Source:} Restklassenteilkörper\par
\noindent\emph{Rusky / Russian:} подполе классов вычетов\par
\noindent\emph{Ukrajinsky / Ukrainian:} підполе класів лишків\par
\smallskip\noindent\emph{Razlog:} Podpolje ostatkovyh klasov čuva v jednom nazivu i inkluziju i poljevu strukturu. V rabotě 24, §5, tekst izrično govori o podpolju vyvedenom tvorjenjem kvocientov i měri po njem stepen podprstena, ktory jest prědstavjen klasami iz osnovnogo ideala. Zato podpolje ne jest libovoljna podmnožina ostatkovyh klasov, a stepen ne jest prosto čislo elementov. Arhivna racionala čuva tę svyzu medžu kvocientami, podpoljem i prstenom, vměsto opuščenja pod- ili zaměny polja prstenom. Pri primarnom idealu tekst osoblivo prěhodi k prstenu vsih klasov; to jest kontekstualny rezultat, ne slovnična definicija. Različne matematične tipy ostavajut vidime črěz kratak, sostavny naziv.\par
\noindent\emph{Reason:} Podpolje ostatkovyh klasov preserves both inclusion and field structure in one expression. In Paper 24, Section 5, the text explicitly describes a subfield obtained by forming quotients and measures over it the degree of a subring represented by classes from the base ideal. Thus the subfield is not an arbitrary subset of residue classes, and degree is not simply a count of elements. The preserved rationale retains this connection among quotients, subfield and ring rather than omitting pod- or substituting ring for field. For a primary ideal the text specifically passes to the ring of all classes; that is a contextual result, not the dictionary definition. A short compound keeps the mathematical types visible.\par
\noindent{\footnotesize Izvorne zapisy / Original records: HIST-g099-009.}\par
\smallskip\noindent\emph{Sravnjenje s tekstom:} Podpolje nastava tvorjenjem kvocientov; stepen zavisi od osnovy.\par
\noindent\emph{Textual comparison:} The subfield is formed by quotients; degree is relative to a base.\par
\noindent{\footnotesize TEXT-COMP-629; pasaže / passages: 15637, 15639, 15642; \texttt{CORPUS\_FORM\_COMPARISONS.json}.}\par
\end{minipage}\par
\par\medskip\Needspace{7\baselineskip}
\noindent\begin{minipage}{\linewidth}
\hypertarget{TERM-1573}{}
\noindent\textbf{teorem o ostatku / ostatkova grupa / koresidualny}\quad{\footnotesize\texttt{TERM-1573}}\par
\noindent\emph{Iztočnik / Source:} Restsatz / Restgruppe / korresidual\par
\noindent\emph{Rusky / Russian:} теорема об остатке / остаточная группа / корезидуальный\par
\noindent\emph{Ukrajinsky / Ukrainian:} теорема про залишок / залишкова група / корезидуальний\par
\smallskip\noindent\emph{Razlog:} Teorem o ostatku v rabotě 14 jest teorem algebraičnoj geometrije, ne teorem o reziduah kompleksnoj analize. Ostatkova grupa označaje grupu toček na krivoj, ne abstraktnu grupu po kvocientu. Dvě grupy sut koresidualne, kogda se mogut dopolniti toju samoju grupoju do odgovarjajučih prěsěkov; tekst objašnjuje svyzu s ekvivalencijoju idealov i s adjungovanymi krivymi. Koresidualne grupy ne morajut iměti jednako mnogo toček. Arhiv zato čuva ostatkov- za geometričny ostatok i međunarodny koresidualny za definovany odnos. Srodnost s nazvami ostatkovyh klasov ne znači identičnost konstrukcij. Takova razlika jest potrěbna i v slovniku: jadro slova pomaže paměti, ale definiciju dava matematičny kontekst.\par
\noindent\emph{Reason:} Teorem o ostatku in Paper 14 is an algebraic-geometric theorem, not the residue theorem of complex analysis. Ostatkova grupa denotes a group of points on a curve, not an abstract quotient group. Two groups are co-residual when the same group can complete them to the corresponding intersections; the text relates this to ideal equivalence and adjoint curves. Co-residual groups need not contain equally many points. The archive therefore keeps ostatkov- for the geometric remainder and international koresidualny for the defined relation. Resemblance to residue-class terminology does not identify the constructions. The dictionary must retain that distinction: a shared stem helps recognition, but the mathematical context supplies the definition.\par
\noindent{\footnotesize Izvorne zapisy / Original records: HIST-g076-020.}\par
\smallskip\noindent\emph{Sravnjenje s tekstom:} Grupe toček i koresidualnost ne sut kvocientna grupa ili teorem kompleksnoj analize.\par
\noindent\emph{Textual comparison:} Point groups and co-residuality are neither an abstract quotient group nor the theorem of complex analysis.\par
\noindent{\footnotesize TEXT-COMP-630; pasaže / passages: 10008, 10010, 10012, 10014, 10016, 10018, 10020, 10022, 10024, 10026; \texttt{CORPUS\_FORM\_COMPARISONS.json}.}\par
\end{minipage}\par
\par\medskip\Needspace{7\baselineskip}
\noindent\begin{minipage}{\linewidth}
\hypertarget{TERM-1574}{}
\noindent\textbf{sistem ostatkov / kolco ostatkovyh klasov / ostatkove klasy}\quad{\footnotesize\texttt{TERM-1574}}\par
\noindent\emph{Iztočnik / Source:} Restsystem / Restklassenring / Restklassen\par
\noindent\emph{Rusky / Russian:} система вычетов / кольцо классов вычетов / классы вычетов\par
\noindent\emph{Ukrajinsky / Ukrainian:} система лишків / кільце класів лишків / класи лишків\par
\smallskip\noindent\emph{Razlog:} Dodatok raboty 45 izrično oprěděljaje polny sistem ostatkov kako sistem prědstaviteljev konečno mnogyh linearno nezavisnyh klasov nad oblastju koeficientov. Zato sistem, prsten i posamezne klasy ne možno sbiti do jednogo objekta. V teoremu IV polinomy davajut prědstavitelje; relacije pokazyvajut, že ih linearne kombinacije izčerpavajut odgovarjajuče klasy i že prědstavitelje sut linearno nezavisne. Dolžina ideala svyzuje se s rangom prstena, ne s kardinalnostju množenstva vsih klasov. Arhiv izbira sostavnu rodinu ostatok/ostatkovy i čuva kolco; nynějšny tekst imaje prsten. Tej izbor trěbuje čitanja definicije: sama fraza polny sistem ostatkov mogla by v drugom kontekstu značiti po jednom prědstavitelju každoj kongruentnoj klasy.\par
\noindent\emph{Reason:} The supplement to Paper 45 explicitly defines a complete residue system as representatives of finitely many linearly independent classes over the coefficient domain. System, ring and individual classes must therefore not be collapsed into one object. In Theorem IV polynomials provide representatives; the relations show that their linear combinations exhaust the corresponding classes and that they are linearly independent. Ideal length is related to ring rank, not to the cardinality of all classes. The archive chooses the compositional ostatok/ostatkovy family and retains kolco; the current text has prsten. Reading the definition is essential: in another context “complete residue system” could mean one representative of every congruence class.\par
\noindent{\footnotesize Izvorne zapisy / Original records: HIST-g214-010.}\par
\smallskip\noindent\emph{Sravnjenje s tekstom:} Dodatok izrično daje linearnu nezavisnost i izčerpanje klasov kombinacijami.\par
\noindent\emph{Textual comparison:} The supplement explicitly requires linear independence and exhaustion of classes by combinations.\par
\noindent{\footnotesize TEXT-COMP-631; pasaže / passages: 45-isv.tex:393, 45-isv.tex:394, 45-isv.tex:395, 45-isv.tex:396, 45-isv.tex:397, 45-isv.tex:398, 45-isv.tex:399, 45-isv.tex:400, 45-isv.tex:425, 45-isv.tex:426, 45-isv.tex:428, 45-isv.tex:429, 45-isv.tex:431, 45-isv.tex:432, 45-isv.tex:433, 45-isv.tex:434, 45-isv.tex:435, 45-isv.tex:436, 45-isv.tex:437, 45-isv.tex:438, 45-isv.tex:439, 45-isv.tex:440, 45-isv.tex:441, 45-isv.tex:442, 45-isv.tex:444, 45-isv.tex:445, 45-isv.tex:446, 45-isv.tex:447, 45-isv.tex:448, 45-isv.tex:449, 45-isv.tex:450, 45-isv.tex:451, 45-isv.tex:452, 45-isv.tex:453, 45-isv.tex:454, 45-isv.tex:455, 45-isv.tex:456, 45-isv.tex:458, 45-isv.tex:459, 45-isv.tex:460, 45-isv.tex:461, 45-isv.tex:462; \texttt{CORPUS\_FORM\_COMPARISONS.json}.}\par
\end{minipage}\par
\par\medskip\Needspace{7\baselineskip}
\noindent\begin{minipage}{\linewidth}
\hypertarget{TERM-1575}{}
\noindent\textbf{ostatny sistem $G-[\Theta]$}\quad{\footnotesize\texttt{TERM-1575}}\par
\noindent\emph{Iztočnik / Source:} Restsystem $G-[\Theta]$\par
\noindent\emph{Rusky / Russian:} [Bukvy ne sut zachovane / Lettering not recoverable] $G-[\Theta]$\par
\noindent\emph{Ukrajinsky / Ukrainian:} [Bukvy ne sut zachovane / Lettering not recoverable] $G-[\Theta]$\par
\smallskip\noindent\emph{Razlog:} Ostatny sistem v rabotě 9, §6, imenuje ostavajuču čest oblasti, zapisanu $G-[\Theta]$, ktoru možno vzeti za sistem prisojedinjanyh funkcij. V nynějšnom tekstu oblast označena jest frakturnym $G$; konstrukcija daje prědstavjenje oblasti prěz $[\Theta,\mathfrak S]$ pri uslovju dopustimosti. Znak minus tu ne oprěděljaje kvocientny prsten, i blizkost k Restsystem v rabotě 45 ne opravdyvaje taku zaměnu. Arhivni rusky i ukrajinsky slovny tekst jest izgubljen v znakach pytanja; zachovana jest samo formula. Slovnik ne domysljaje izgubjene slova iz sadašnjego ponimanja. On pokazuje granicu zachovanogo svědočstva i razlikuje istorijno zapisan izbor od novogo možnogo prěvoda.\par
\noindent\emph{Reason:} Ostatny sistem in Paper 9, Section 6, names the remaining part of a domain, written $G-[\Theta]$, which can serve as the system of functions to adjoin. The current text uses a fraktur $G$ for the domain and represents it as $[\Theta,\mathfrak S]$ under the admissibility condition. The minus sign here does not define a quotient ring, and resemblance to Restsystem in Paper 45 does not justify that substitution. The Russian and Ukrainian wording in the archive has been lost to question marks; only the formula survives. The dictionary does not infer the lost words from present understanding. It marks the evidence boundary and distinguishes the recorded choice from a possible new translation.\par
\noindent{\footnotesize Izvorne zapisy / Original records: HIST-g064-012.}\par
\smallskip\noindent\emph{Sravnjenje s tekstom:} Minus i dopustimy sistem davajut prědstavjenje oblasti, ne faktor-prsten.\par
\noindent\emph{Textual comparison:} The subtraction expression and admissible system provide a representation of the domain, not a quotient ring.\par
\noindent{\footnotesize TEXT-COMP-632; pasaže / passages: 7785, 7786, 7787, 7788, 7789, 7790, 7791, 7792, 7793, 7794, 7795, 7801, 7802, 7803, 7804, 7805; \texttt{CORPUS\_FORM\_COMPARISONS.json}.}\par
\end{minipage}\par
\par\medskip\Needspace{7\baselineskip}
\noindent\begin{minipage}{\linewidth}
\hypertarget{TERM-1576}{}
\noindent\textbf{rezultant}\quad{\footnotesize\texttt{TERM-1576}}\par
\noindent\emph{Iztočnik / Source:} Resultante\par
\noindent\emph{Rusky / Russian:} результант\par
\noindent\emph{Ukrajinsky / Ukrainian:} результант\par
\smallskip\noindent\emph{Razlog:} Rezultant jest međunarodny algebraičny termin, čij korenj čitatelj može poznavati iz ruskoj i ukrajinskoj matematičnoj tradicije. Arhivny izbor čuva registr eliminacionnoj teorije, svyazanoj v tekstu s Mertensom i Königom; to ne znači, že jih raboty byly nyně nezavisno pročitane. V §11 raboty 6 nenulevy rezultant homogennyh form isključaje zajedno nenulevo rěšenje. Pri nehomogennom sistemu uslovje o členah najvyššego stepena daje dostatočny kriterij v obsudženom argumentu; ne trěba zaměnjati dostatočnost nužnostju. Rezultant zato ne jest prosto konečny rezultat izčislenja. Racionala daje prědnost poznavajemomu tehničnomu korenju prěd opisnym novotvorom, pri sohranjenju uslovij, od ktoryh zavisi jego matematično značenje.\par
\noindent\emph{Reason:} Rezultant is an international algebraic term whose stem is recognizable in the Russian and Ukrainian comparanda. The recorded choice preserves the register of elimination theory associated in the text with Mertens and König; that does not imply their works have now been independently read. In Section 11 of Paper 6, a nonzero resultant of homogeneous forms excludes a common nonzero solution. For a nonhomogeneous system the condition on highest-degree parts supplies a sufficient criterion in the argument; sufficiency must not be silently changed to necessity. A resultant is therefore not simply the final result of a calculation. The rationale favours a recognizable technical stem over a descriptive coinage while retaining the hypotheses governing its mathematical use.\par
\noindent{\footnotesize Izvorne zapisy / Original records: HIST-g051-003.}\par
\smallskip\noindent\emph{Sravnjenje s tekstom:} Homogenny kriterij i dostatočno uslovje za nehomogenny sistem trěbujut razlikovanja.\par
\noindent\emph{Textual comparison:} The homogeneous criterion and a sufficient condition for a nonhomogeneous system must be distinguished.\par
\noindent{\footnotesize TEXT-COMP-633; pasaže / passages: 6119, 6121, 6123, 6125, 6139, 6141, 6165, 6173; \texttt{CORPUS\_FORM\_COMPARISONS.json}.}\par
\end{minipage}\par
\par\medskip\Needspace{7\baselineskip}
\noindent\begin{minipage}{\linewidth}
\hypertarget{TERM-1577}{}
\noindent\textbf{rezultant členov najvyššej dimenzije}\quad{\footnotesize\texttt{TERM-1577}}\par
\noindent\emph{Iztočnik / Source:} Resultante der Glieder höchster Dimension\par
\noindent\emph{Rusky / Russian:} результант членов наивысшей размерности\par
\noindent\emph{Ukrajinsky / Ukrainian:} результант членів найвищої вимірності\par
\smallskip\noindent\emph{Razlog:} Členy najvyššej dimenzije v arhivnom nazivu znače homogennye česti najvyššego stepena; dimenzija tu ne jest dimenzija algebraičnoj množiny. V §15 raboty 6 potrěbno jest, aby jih rezultant byl jedinicoju v oblasti koeficientov, ne samo nenulevym. Ta silnějša prědpoložka dava odgovarjajuče tvrdženje nad oblastju, v ktoroj ne vsako nenulevo čislo imaje obratny element. Tekst uvodi rezultant $R_0=\varepsilon$ i ilustruje uslovje elementarnymi simetričnymi funkcijami s rezultantom $+1$ ili $-1$. Arhiv čuva frazu raděje nego uvodi nov naziv pri prěhodu od ranějšego kriterija. Slovnik zato objašnjuje istorijni smysl dimenzije i rozděl medžu nenulevostju i obratimostju; ne vyhladžuje jih pod odnym poznatym slovom.\par
\noindent\emph{Reason:} The archived členov najvyššej dimenzije means the homogeneous parts of highest degree, not the dimension of an algebraic set. In Section 15 of Paper 6 their resultant must be a unit in the coefficient domain, not merely nonzero. That stronger hypothesis supports the corresponding assertion over a domain where nonzero elements need not be invertible. The text writes $R_0=\varepsilon$ and illustrates the condition using elementary symmetric functions with resultant $+1$ or $-1$. The archive retains the phrase rather than introducing a new name when moving beyond the earlier criterion. The dictionary therefore explains the historical degree sense of dimenzija and the difference between nonvanishing and invertibility instead of smoothing both away under familiar vocabulary.\par
\noindent{\footnotesize Izvorne zapisy / Original records: HIST-g055-008.}\par
\smallskip\noindent\emph{Sravnjenje s tekstom:} Najvyšša dimenzija označaje stepen; jedinica jest silnějše uslovje nego nenulevost.\par
\noindent\emph{Textual comparison:} Highest dimension means degree; being a unit is stronger than being nonzero.\par
\noindent{\footnotesize TEXT-COMP-634; pasaže / passages: 6573, 6575, 6577, 6579, 6581, 6583, 6585, 6587, 6588, 6589, 6590, 6591, 6592; \texttt{CORPUS\_FORM\_COMPARISONS.json}.}\par
\end{minipage}\par
\par\medskip\Needspace{7\baselineskip}
\noindent\begin{minipage}{\linewidth}
\hypertarget{TERM-1578}{}
\noindent\textbf{rezultant $i$-go nivoja / elementarny dělitelj $i$-go nivoja}\quad{\footnotesize\texttt{TERM-1578}}\par
\noindent\emph{Iztočnik / Source:} Resultante $i$-ter Stufe / Elementarteiler $i$-ter Stufe\par
\noindent\emph{Rusky / Russian:} результант $i$-го уровня / элементарный делитель $i$-го уровня\par
\noindent\emph{Ukrajinsky / Ukrainian:} результант $i$-го рівня / елементарний дільник $i$-го рівня\par
\smallskip\noindent\emph{Razlog:} Nivoj prěvodí Stufe kako indeksovanu vrstvu konstrukcije, ne prosto vremenny krok algoritma. Definicija VI v rabotě 22, §5, nazyvaje $R^{(i)}(x)$ rezultantom $i$-go nivoja i $E^{(i)}(x)$ elementarnym děliteljem toho nivoja. Oni nastavajut poslě prisojedinjenja ostalyh neoprěděljenyh; prvi jest odgovarjajuča norma, vtory najvyšši elementarny dělitelj. Produkty po vseh indeksah daju dvě formy, medžu ktorymi tekst dokazyvaje dělitelnost. Arhivni izbor nivoj čuva takvu indeksaciju i izběgaje pretežno proceduralno čitanje etapy. Tej razlog ne jest obči zakaz slova etapa v drugih kontekstah. Dva arhivna zapisa čuvajut to samo pravilo; bukva $i$ ostavaje matematičnym indeksom i pokazyvaje se odlično od običnogo slova.\par
\noindent\emph{Reason:} Nivoj renders Stufe as an indexed level of the construction, not simply a chronological algorithmic step. Definition VI in Paper 22, Section 5, calls $R^{(i)}(x)$ the resultant of level $i$ and $E^{(i)}(x)$ the elementary divisor of that level. They arise after the remaining indeterminates are adjoined; the former is the corresponding norm and the latter the highest elementary divisor. Products over all indices give two forms whose divisibility relation is proved. The recorded choice nivoj preserves this indexing and avoids a predominantly procedural reading of “stage.” It is not a general ban on etapa elsewhere. Two archival records preserve the same rule; $i$ remains a mathematical index, typeset distinctly from ordinary wording.\par
\noindent{\footnotesize Izvorne zapisy / Original records: HIST-g089-002, HIST-g092-044.}\par
\smallskip\noindent\emph{Sravnjenje s tekstom:} Definicija VI veže indeks s normoju, elementarnym děliteljem i produktami po nivojah.\par
\noindent\emph{Textual comparison:} Definition VI links the index to a norm, an elementary divisor and products over levels.\par
\noindent{\footnotesize TEXT-COMP-635; pasaže / passages: 14830, 14832, 14834, 14836, 14839, 14841, 14842, 14848, 14850; \texttt{CORPUS\_FORM\_COMPARISONS.json}.}\par
\end{minipage}\par
\par\medskip\Needspace{7\baselineskip}
\noindent\begin{minipage}{\linewidth}
\hypertarget{TERM-1579}{}
\noindent\textbf{rezultantno značenje}\quad{\footnotesize\texttt{TERM-1579}}\par
\noindent\emph{Iztočnik / Source:} Resultantenbedeutung\par
\noindent\emph{Rusky / Russian:} значение результанта\par
\noindent\emph{Ukrajinsky / Ukrainian:} значення результанта\par
\smallskip\noindent\emph{Razlog:} Rezultantno značenje označaje funkciju, ktoru norma izpolnjaje v eliminacionnoj teoriji. V rabotě 25 svyza medžu razloženjem ideala i razloženjem normy daje rezultantom ih ulogu; slovo značenje tu ne označaje samo numeričnu cěnnost. Kratky adjektiv rezultantno veže vysvětlenje s poznatym terminom bez tvorjenja novogo složenogo substantiva. Arhiv tak objašnjuje izbor, ale ne dava měrene skory za vse možne sinonimy. Pri čitanju potrěbno jest odlišiti normu v tutoj algebraičnoj konstrukciji od metrijnoj normy, a svyzu razloženij od jejich bukvalnoj identičnosti. Takože mnogokratnost, vyražena stepenom, zavisi od relativnoj bazy; ona ne jest neobjasnjeno opakovanje slova rezultant.\par
\noindent\emph{Reason:} Rezultantno značenje denotes the role played by the norm in elimination theory. In Paper 25 the relation between ideal decomposition and factorization of the norm gives resultants their significance; značenje does not mean merely a numerical value. The short adjective rezultantno connects the explanation to the familiar technical term without creating a new compound noun. That is the preserved rationale, not a measured ranking of every possible synonym. Readers must distinguish this algebraic norm from a metric norm and a relation between factorizations from their literal identity. The multiplicity expressed by degree also depends on the relative base; it is not an unexplained repetition of the word resultant.\par
\noindent{\footnotesize Izvorne zapisy / Original records: HIST-g102-013.}\par
\smallskip\noindent\emph{Sravnjenje s tekstom:} Rezultantno značenje označaje matematičnu ulogu, ne samo numeričnu cěnnost.\par
\noindent\emph{Textual comparison:} Resultant significance denotes a mathematical role, not merely a numerical value.\par
\noindent{\footnotesize TEXT-COMP-636; pasaže / passages: 15847, 15849, 15851; \texttt{CORPUS\_FORM\_COMPARISONS.json}.}\par
\end{minipage}\par
\par\medskip\Needspace{7\baselineskip}
\noindent\begin{minipage}{\linewidth}
\hypertarget{TERM-1580}{}
\noindent\textbf{rezultanta forma}\quad{\footnotesize\texttt{TERM-1580}}\par
\noindent\emph{Iztočnik / Source:} Resultantenform\par
\noindent\emph{Rusky / Russian:} результирующая форма\par
\noindent\emph{Ukrajinsky / Ukrainian:} результантна форма\par
\smallskip\noindent\emph{Razlog:} Rezultanta forma v arhivu jest kratka međunarodna fraza. Ukrajinsky sravnjajuči izraz imaje jasny adjektivny korenj rezultant-, pokym rusky izraz rezultirujuščaja možet v običnom jezyku značiti prosto výslednu formu. Matematičny kontekst zato jest važnějši než mehanično kopiranje približno sličnogo slova. V rabotě 18 idealu odgovarja forma složena iz rezultantnyh faktorov; ona izčezaje na zajednyh nuljah, no argument čuva takože faktory i mnogokratnosti. Jednakost forma pri dělitelnosti idealov jest silnějše tvrđenje nego jednakost golyh množenstv nuljev. Slovnik čuva istorijno zapisanu formu i jej razlog, ne proglašaje neizslědovanu morfologičnu variantu jedinym pravilnym standardom. V OLP rezultant takože ne směje byti zaměnjen logičnym pojmom rezolucije po sličnosti zvuka.\par
\noindent\emph{Reason:} Rezultanta forma is the compact international phrase preserved in the archive. The Ukrainian comparandum has an explicit rezultant- adjective, while the Russian rezultirujuščaja can mean simply “resulting” in ordinary language. Mathematical context therefore matters more than mechanical copying of a superficially similar word. In Paper 18 an ideal is associated with a form made from resultant factors; it vanishes at the common zeros, but the argument retains factors and multiplicities as well. Equality of the forms under ideal divisibility is stronger than equality of the bare zero sets. The dictionary preserves the historical form and its rationale without declaring an unexamined morphological variant uniquely correct. Nor should OLP identify resultant with logical resolution through resemblance of sound.\par
\noindent{\footnotesize Izvorne zapisy / Original records: HIST-g080-003.}\par
\smallskip\noindent\emph{Sravnjenje s tekstom:} Jednakost form nosi vyše informacije nego samo jednakost množenstv nuljev.\par
\noindent\emph{Textual comparison:} Equality of forms carries more information than equality of their zero sets alone.\par
\noindent{\footnotesize TEXT-COMP-637; pasaže / passages: 12339, 12341, 12343, 12345, 12347, 12349, 12351, 12353, 12355, 12357; \texttt{CORPUS\_FORM\_COMPARISONS.json}.}\par
\end{minipage}\par
\par\medskip\Needspace{7\baselineskip}
\noindent\begin{minipage}{\linewidth}
\hypertarget{TERM-1581}{}
\noindent\textbf{rezultantna forma / forma elementarnyh děliteljev}\quad{\footnotesize\texttt{TERM-1581}}\par
\noindent\emph{Iztočnik / Source:} Resultantenform / Elementarteilerform\par
\noindent\emph{Rusky / Russian:} результантная форма / форма элементарных делителей\par
\noindent\emph{Ukrajinsky / Ukrainian:} результантна форма / форма елементарних дільників\par
\smallskip\noindent\emph{Razlog:} Rezultantna forma i forma elementarnyh děliteljev nazyvajut dvě konstrukcije, ktore v §5 raboty 22 stavajut invariantami. Arhiv davaje prědnost stalym složenym nazivam prěd svobodnym opisom: internationalny korenj rezultant- ohranjaje prepoznavajemy naziv, a dělitelj prodolžaje už upotrěbjenu terminologičnu rodinu. Rusky i ukrajinsky zapisy dajut slične složene formy, ale ne sut nezavisny test razumljivosti. V tekste rezultantna forma $R$ i forma $E$ ne sut obče identične: navedeno jest $E\mid R$ i $R\mid E^k$. Zato podobnost nazivov ne dozvaljaje jih sjediniti. Prědnost tu imajut razlikovanje invariantov i kontinuitet nazivov; čislove ocěnky kandidatov ne sut zachovane.\par
\noindent\emph{Reason:} Rezultantna forma and forma elementarnyh děliteljev name two constructions that become invariants in Paper 22, §5. The record favours stable compound names over free paraphrase: the international resultant stem remains recognizable, while dělitelj continues an established terminological family. The Russian and Ukrainian records supply similar compounds, not an independent comprehension test. The text does not generally identify the resultant form $R$ with the elementary-divisor form $E$: it states $E\mid R$ and $R\mid E^k$. Similar names therefore do not justify merging the objects. Distinction between the invariants and continuity of naming are the documented priorities; numerical candidate scores are not preserved.\par
\noindent{\footnotesize Izvorne zapisy / Original records: HIST-g089-001, HIST-g092-043.}\par
\smallskip\noindent\emph{Sravnjenje s tekstom:} Nazivy razlikujut invarianty; relacija dělimosti ne jest identičnost.\par
\noindent\emph{Textual comparison:} The names distinguish invariants; divisibility is not identity.\par
\noindent{\footnotesize TEXT-COMP-638; pasaže / passages: 14832, 14834, 14836, 14839, 14841, 14842, 14848, 14850; \texttt{CORPUS\_FORM\_COMPARISONS.json}.}\par
\end{minipage}\par
\par\medskip\Needspace{7\baselineskip}
\noindent\begin{minipage}{\linewidth}
\hypertarget{TERM-1582}{}
\noindent\textbf{rezultantna forma / rezultant}\quad{\footnotesize\texttt{TERM-1582}}\par
\noindent\emph{Iztočnik / Source:} Resultantenform / Resultante\par
\noindent\emph{Rusky / Russian:} результантная форма / результант\par
\noindent\emph{Ukrajinsky / Ukrainian:} результантна форма / результант\par
\smallskip\noindent\emph{Razlog:} Rezultantna forma / rezultant čuva svez medžu nazvom formy i jeji algebraičnoju funkcijeju. Arhiv opiraje izbor na internationalnu rodinu rezultant-, s blizkimi ruskymi i ukrajinskymi formami, i potrěbuje razlikovati ju od determinant i norma. To potrěbovanje odnosi se na terminologične role, ne na tvrdženje, že objekty nikdy ne mogut byti ravne. V razgledanoj konstrukciji raboty 22 tekst upravo svezuje normu, determinant i rezultantnu formu. Slovnik zato ne može povtoriti kratku arhivnu formulaciju kako bezuslovnu zabranu takoj identifikacije. Smysl prěklada jest ohranjati naziv, pod ktorym autor uvodi objekt, i izjaviti matematičnu ravnost tam, kde ona jest dokazana.\par
\noindent\emph{Reason:} Rezultantna forma / rezultant preserves the relation between the name of the form and its algebraic function. The record favours the international resultant family, with close Russian and Ukrainian forms, and requires a distinction from determinant and norm. That requirement concerns terminological roles, not a claim that the objects can never coincide. In the construction examined in Paper 22, the text expressly connects the norm, determinant, and resultant form. The dictionary must not turn the abbreviated historical rationale into an unconditional prohibition on that identification. Translation should retain the name under which the author introduces an object and state the mathematical equality where the argument establishes it.\par
\noindent{\footnotesize Izvorne zapisy / Original records: HIST-g093-002.}\par
\smallskip\noindent\emph{Sravnjenje s tekstom:} Terminologične role sut različne i togda, kogda konstrukcija davaje ravne objekty.\par
\noindent\emph{Textual comparison:} The terms have different roles even when a construction identifies their objects.\par
\noindent{\footnotesize TEXT-COMP-639; pasaže / passages: 14832, 14834, 14836; \texttt{CORPUS\_FORM\_COMPARISONS.json}.}\par
\end{minipage}\par
\par\medskip\Needspace{7\baselineskip}
\noindent\begin{minipage}{\linewidth}
\hypertarget{TERM-1583}{}
\noindent\textbf{rezultantna forma i teorija eliminacije}\quad{\footnotesize\texttt{TERM-1583}}\par
\noindent\emph{Iztočnik / Source:} Resultantenform und Eliminationstheorie\par
\noindent\emph{Rusky / Russian:} результантная форма и теория исключения\par
\noindent\emph{Ukrajinsky / Ukrainian:} результантна форма і теорія елімінування\par
\smallskip\noindent\emph{Razlog:} Teorija eliminacije opisuje v rabotě 22 odstranjenje proměnnyh v teoriji polinomnyh idealov. Arhiv izbral internationalny korenj zaradi prepoznavajemosti, ale sravnjenje ne pokazuje jednaky korenj vo vseh jezykah: rusky zapis uživaje isključenije, ukrajinsky eliminuvannja. To jest izbor disciplinarnogo naziva, ne proste glasovanje po bukvah. §7 svezuje rezultantnu formu s postopnym eliminovanjem, dimenzijeju i mnogokratnostju; pri specializaciji ostavajut važni izrečene uslovja. Naziv zato ne davaje dozvoljenja vypustiti te uslovja. Za OLP trěba oddělno obosnovati eliminaciju kvantifikatorov: sličny korenj sam ne dokazyvaje identičnost algebraičnogo i logičnogo pojma.\par
\noindent\emph{Reason:} Teorija eliminacije concerns the elimination of variables in polynomial-ideal theory in Paper 22. Recognizability motivates the international stem, but the comparanda do not all use that stem: the Russian record uses exclusion terminology, whereas the Ukrainian uses an elimination form. This is a choice of disciplinary expression, not a simple vote on matching letters. Section 7 relates the resultant form to successive elimination, dimension, and multiplicity; its stated specialization conditions remain essential. The term does not permit those conditions to disappear. Quantifier elimination in OLP requires its own contextual justification: a shared stem alone does not identify the algebraic and logical concepts.\par
\noindent{\footnotesize Izvorne zapisy / Original records: HIST-g091-001, HIST-g092-070.}\par
\smallskip\noindent\emph{Sravnjenje s tekstom:} Naziv glavy ne izpuščaje uslovij o dimenziji, mnogokratnosti i specializaciji.\par
\noindent\emph{Textual comparison:} The section title does not dispense with conditions on dimension, multiplicity, and specialization.\par
\noindent{\footnotesize TEXT-COMP-640; pasaže / passages: 14974, 14976, 14994, 14998, 15003, 15008, 15010, 15016, 15018; \texttt{CORPUS\_FORM\_COMPARISONS.json}.}\par
\end{minipage}\par
\par\medskip\Needspace{7\baselineskip}
\noindent\begin{minipage}{\linewidth}
\hypertarget{TERM-1584}{}
\noindent\textbf{recipro?no predstavjenje; direktno predstavjenje [Čestično nezachovane bukvy / Partly unrecoverable lettering]}\quad{\footnotesize\texttt{TERM-1584}}\par
\noindent\emph{Iztočnik / Source:} reziproke Darstellung; direkte Darstellung\par
\noindent\emph{Rusky / Russian:} [Bukvy ne sut zachovane / Lettering not recoverable]\par
\noindent\emph{Ukrajinsky / Ukrainian:} [Bukvy ne sut zachovane / Lettering not recoverable]\par
\smallskip\noindent\emph{Razlog:} Recipročno prědstavjenje v nynějšnom tekste raboty 40 obracaje poredok množenja; direktno prědstavjenje go ohranjaje. Razlika vyhodi iz lěvyh i pravyh operatorov a ne iz pojetja obratnoj matrice. Arhivna motivacija davaje prědnost internationalnomu nazivu prěd nenazvanoju domačnoju kalkoju, ktora je ocenjena kako menje prepoznavajema; ta ocěnka ne jest podkrepljena testom čitateljev. Sam zapis jest poškodženy: medžuslovjansko recipro?no ne čuva vseh bukv, a rusky i ukrajinsky redy ne dozvaljajut obnoviti slova. Tekst korpusa zato služi kako nynějšny kontekst, ne kako dokaz, čto je stalo v vseh izgubljenyh poljah. Slovo recipročno ne trěba mechanicky zaměnjati s obratno.\par
\noindent\emph{Reason:} In the current text of Paper 40, recipročno prědstavjenje reverses multiplication order, whereas direktno prědstavjenje preserves it. The distinction comes from left and right operators, not from the concept of an inverse matrix. The historical rationale prefers international terminology to an unnamed native calque judged less recognizable; no reader test supports that judgement. The record itself is damaged: Interslavic recipro?no lacks a letter, and the Russian and Ukrainian lines do not preserve recoverable words. The corpus provides a current contextual witness, not proof of what every lost field originally contained. Recipročno must therefore not be mechanically replaced by obratno.\par
\noindent{\footnotesize Izvorne zapisy / Original records: HIST-g210-005.}\par
\smallskip\noindent\emph{Sravnjenje s tekstom:} Obračenje poredka operatorov ne jest isto čto obračenje matrice. Poškodžene arhivne bukvy ne sut obnovjene bez dokazov.\par
\noindent\emph{Textual comparison:} Reversing the order of operators is not matrix inversion. Damaged historical lettering has not been reconstructed without evidence.\par
\noindent{\footnotesize TEXT-COMP-641; pasaže / passages: 23375, 23377, 23379, 23381, 23421, 23548, 23556, 23566; \texttt{CORPUS\_FORM\_COMPARISONS.json}.}\par
\end{minipage}\par
\par\medskip\Needspace{7\baselineskip}
\noindent\begin{minipage}{\linewidth}
\hypertarget{TERM-1585}{}
\noindent\textbf{obratna matrica}\quad{\footnotesize\texttt{TERM-1585}}\par
\noindent\emph{Iztočnik / Source:} Reziproke der Matrix\par
\noindent\emph{Rusky / Russian:} обратная матрица\par
\noindent\emph{Ukrajinsky / Ukrainian:} обернена матриця\par
\smallskip\noindent\emph{Razlog:} Obratna matrica zaměnjuje historijsko označenje Reziproke der Matrix modernym nazivom. Rusky obratnaja i ukrajinsky obernena podkrepljajut to pojetjno čitanje. Arhiv tvrdži, že taka modernizacija ne měnjaje algebraičny smysl; to ale ne zaměnjuje kontrolu indeksov. V nynějšnom (43) raboty 17 stojut $\alpha^{12}=-\alpha_{21}/A$ i $\alpha^{21}=-\alpha_{12}/A$, pri $A=\det(\alpha_{ik})\ne0$. Pri obyčnom poredku redov i stolpcov tako označena matrica jest transponovana obratna matrica. Možna jest ina konvencija indeksov ili pogrěška prěpisa; bez primarnogo svědka ne trěba pripisyvati pogrěšku autoru. Leksičny izbor obratna jest razumlivy, ale presna notacija ostavaje otvoreno pitanje.\par
\noindent\emph{Reason:} Obratna matrica modernizes the historical expression Reziproke der Matrix. Russian obratnaja and Ukrainian obernena support that conceptual reading. The record says that modernization leaves the algebra unchanged, but this does not replace an index check. In the current equation (43) of Paper 17, $\alpha^{12}=-\alpha_{21}/A$ and $\alpha^{21}=-\alpha_{12}/A$, with $A=\det(\alpha_{ik})\ne0$. Under ordinary row-and-column conventions, the matrix so labelled is the transpose of the inverse. A different indexing convention or a transcription defect may be involved; no error should be attributed to the author without the primary witness. The lexical choice obratna is intelligible, but the exact notation remains an open question.\par
\noindent{\footnotesize Izvorne zapisy / Original records: HIST-g079-074.}\par
\smallskip\noindent\emph{Sravnjenje s tekstom:} Relacije v (43) potrěbujut sravnjenje s primarnym svědkom; leksična modernizacija sama ne potvrđuje notaciju.\par
\noindent\emph{Textual comparison:} The relations in (43) require comparison with the primary witness; lexical modernization alone does not validate the notation.\par
\noindent{\footnotesize TEXT-COMP-642; pasaže / passages: 12110, 12112, 12114, 12150, 12151, 12152, 12153, 12154, 12155, 12156, 12157, 12158, 12159, 12163; \texttt{CORPUS\_FORM\_COMPARISONS.json}.}\par
\end{minipage}\par
\par\medskip\Needspace{7\baselineskip}
\noindent\begin{minipage}{\linewidth}
\hypertarget{TERM-1586}{}
\noindent\textbf{vzaimnosti}\quad{\footnotesize\texttt{TERM-1586}}\par
\noindent\emph{Iztočnik / Source:} Reziprozitäten\par
\noindent\emph{Rusky / Russian:} соотношения взаимности\par
\noindent\emph{Ukrajinsky / Ukrainian:} співвідношення взаємності\par
\smallskip\noindent\emph{Razlog:} Vzaimnosti jest kratok arhivny naziv za recipročne relacije medžu modulnymi kvocientami. Nynějšny tekst uživaje vzajemnosti; ta razlika form ostavaje vidima v slovniku. Arhiv ne potrěboval dolgy opis, zato že jednačenja noset presny sodržaj relacije. To jest opravdanje kratkosti v konkretnom izloženju, ne pravilo izpuščati definicije. V §2 raboty 22 čitatelj može sravniti oba naprěvenja relacije i prědpoloženja o faktorah. Naziv ne označaje obče aritmetičny zakon vzajemnosti, niti vsaku simetričnu relaciju. Pri prěnosu do drugogo teksta trěba vozstanoviti objekt, operaciju i uslovja, ne samo izbrati obče slovo za vzajemnost.\par
\noindent\emph{Reason:} Vzaimnosti is the compact archived name for reciprocal relations between module quotients. The current text uses vzajemnosti; the difference remains visible in the dictionary. The record considers a short name sufficient because the displayed equations supply the precise content of the relation. That justifies brevity in this particular exposition, not the omission of definitions as a general rule. Paper 22, §2 lets the reader compare the two directions and the assumptions on the factors. The name denotes neither an unspecified arithmetic reciprocity law nor every symmetric relation. Transfer to another text requires recovering the objects, operation, and conditions, not merely selecting a generic word for reciprocity.\par
\noindent{\footnotesize Izvorne zapisy / Original records: HIST-g086-008, HIST-g092-018.}\par
\smallskip\noindent\emph{Sravnjenje s tekstom:} Jednačenja određajut, v kakom smyslu relacije sut vzajemne.\par
\noindent\emph{Textual comparison:} The equations determine the sense in which the relations are reciprocal.\par
\noindent{\footnotesize TEXT-COMP-643; pasaže / passages: 14436, 14438, 14440, 14444, 14446, 14466; \texttt{CORPUS\_FORM\_COMPARISONS.json}.}\par
\end{minipage}\par
\par\medskip\Needspace{7\baselineskip}
\noindent\begin{minipage}{\linewidth}
\hypertarget{TERM-1587}{}
\noindent\textbf{determinant porjadka $\rho$}\quad{\footnotesize\texttt{TERM-1587}}\par
\noindent\emph{Iztočnik / Source:} $\rho$-reihige Determinante\par
\noindent\emph{Rusky / Russian:} детерминант порядка $\rho$\par
\noindent\emph{Ukrajinsky / Ukrainian:} детермінант порядку $\rho$\par
\smallskip\noindent\emph{Razlog:} Determinant porjadka $\rho$ jest arhivny popraveny izbor: prvy zapis mešal poredok determinanta s rangom modula. V §7 raboty 24 oba čisla sut označene $\rho$, ale imajut različne role. Poredok govori o veličině kvadratnoj matrice, čij determinant jest vzety; rang govori o modulnoj strukturě. Nynějšny tekst imaje formu poredka, ne arhivno porjadka. Slovnik ohranjaje i razliku pojmov, i sviděteljstvo o jezyčnoj formě; ne prědstavjaje staru formu kako novu normu. Odločujuči kriterij byl matematična točnost, ne kratkost. Ta razlika ostavaje potrěbna i togda, kogda vrednosti ranga i poredka slučajno sut ravne.\par
\noindent\emph{Reason:} Determinant porjadka $\rho$ is the archived corrected choice: an earlier wording confused determinant order with module rank. Both numbers are denoted by $\rho$ in Paper 24, §7, but they have different roles. Order describes the size of the square matrix whose determinant is taken; rank describes module structure. The current text has poredka rather than archived porjadka. The dictionary preserves both the conceptual distinction and the evidence of the linguistic form, without presenting the older spelling as a new norm. Mathematical accuracy, not brevity, was decisive. The distinction remains necessary when rank and determinant order happen to have equal values.\par
\noindent{\footnotesize Izvorne zapisy / Original records: HIST-g101-007.}\par
\smallskip\noindent\emph{Sravnjenje s tekstom:} Jednaky simbol za dvě čisla ne sjedinjaje njihove matematične role.\par
\noindent\emph{Textual comparison:} The same symbol for two numbers does not merge their mathematical roles.\par
\noindent{\footnotesize TEXT-COMP-644; pasaže / passages: 15727, 15739, 15743, 15745, 15747, 15755; \texttt{CORPUS\_FORM\_COMPARISONS.json}.}\par
\end{minipage}\par
\par\medskip\Needspace{7\baselineskip}
\noindent\begin{minipage}{\linewidth}
\hypertarget{TERM-1588}{}
\noindent\textbf{teorem Riemanna--Rocha / redukcijny teorem}\quad{\footnotesize\texttt{TERM-1588}}\par
\noindent\emph{Iztočnik / Source:} Riemann-Rochscher Satz / Reduktionssatz\par
\noindent\emph{Rusky / Russian:} теорема Римана--Роха / теорема редукции\par
\noindent\emph{Ukrajinsky / Ukrainian:} теорема Рімана--Роха / теорема редукції\par
\smallskip\noindent\emph{Razlog:} Teorem Riemanna--Rocha i redukcijny teorem ohranjajut dva imena, ne dva zaměnljive opisy jednogo tvrđenja. Arhiv se opiraje na disciplinarny upotrěb i ohranjaje autorovo Reduktionssatz, s objasněnjem o fiksovanyh točkah v proze. V rabotě 14 geometricky pristup vodi od rezidualnogo teorema črez redukciju do teorema Riemanna--Rocha; aritmetično izloženje měnjaje poredok dokaza. Z toga ne slěduje, že jeden teorem možno izpustiti kako sinonim drugogo. Fiksovane točky tu trěba čitati v danom geometrickom kontekstu, ne kako automaticky odkaz na teorem o kontrakcijah. Eponim ohranjaje naučnu identifikaciju, ale sam naziv ne potvrđuje věrnost vsego dokaza.\par
\noindent\emph{Reason:} Teorem Riemanna--Rocha and redukcijny teorem retain two names, not interchangeable descriptions of one assertion. The record appeals to disciplinary usage and preserves the author's Reduktionssatz with an explanation involving fixed points in the prose. In Paper 14, the geometric route proceeds from the residual theorem through reduction to Riemann--Roch, while the arithmetic exposition changes the order of proof. That does not make one theorem dispensable as a synonym of another. Fixed points must be read in this geometric context, not treated as an automatic reference to a contraction theorem. The eponym preserves scholarly identification; the name alone does not certify the entire proof.\par
\noindent{\footnotesize Izvorne zapisy / Original records: HIST-g076-023.}\par
\smallskip\noindent\emph{Sravnjenje s tekstom:} Geometricky i aritmetičny pristupy ne činet nazivy redukcije i Riemanna--Rocha zaměnljivymi.\par
\noindent\emph{Textual comparison:} Geometric and arithmetic approaches do not make the reduction and Riemann--Roch theorem names interchangeable.\par
\noindent{\footnotesize TEXT-COMP-645; pasaže / passages: 9864, 9866, 10024, 10026, 10028, 10030; \texttt{CORPUS\_FORM\_COMPARISONS.json}.}\par
\end{minipage}\par
\par\medskip\Needspace{7\baselineskip}
\noindent\begin{minipage}{\linewidth}
\hypertarget{TERM-1589}{}
\noindent\textbf{Riemannova površina / absolutna Riemannova površina}\quad{\footnotesize\texttt{TERM-1589}}\par
\noindent\emph{Iztočnik / Source:} Riemannsche Fläche / absolute Riemannsche Fläche\par
\noindent\emph{Rusky / Russian:} риманова поверхность / абсолютная риманова поверхность\par
\noindent\emph{Ukrajinsky / Ukrainian:} ріманова поверхня / абсолютна ріманова поверхня\par
\smallskip\noindent\emph{Razlog:} Riemannova površina spojuje ohranjeno vlašćne ime s široko slovjanskim nazivom površina; tak jest formulovana arhivna motivacija. Dodatek absolutna ne jest ukras. V rabotě 14 označaje on konstrukciju toček nezavisnu od izbranoj funkcije $z$; tekst sravnjaje konečne vrednosti $z$ s upotrěboju $1/z$ za ostale točky. Izbor zato trěba ocěnjati po toj definiciji, ne samo po sodobnom zvučenju naziva. Ruske i ukrajinske formy podkrepljajut svez s terminologijeju površin, ale ne stanovjajut samostojno medžuslovjansku normu. Slovnik ne prěnosi kvalifikator absolutna na drugy objekt bez novogo čitanja definicije i ne zaměnjuje historijsko izloženje modernym učebničnym.\par
\noindent\emph{Reason:} Riemannova površina combines the retained proper name with the broadly Slavic noun površina, as the historical rationale puts it. Absolutna is not decorative. In Paper 14 it identifies a construction of points independent of the selected function $z$; the text contrasts finite values of $z$ with the use of $1/z$ for the remaining points. The choice must therefore be assessed against that definition, not merely the modern familiarity of the name. The Russian and Ukrainian forms support a connection with surface terminology but do not independently establish an Interslavic norm. The dictionary does not transfer absolutna to another object without rereading its definition or replace the historical exposition with a modern textbook account.\par
\noindent{\footnotesize Izvorne zapisy / Original records: HIST-g076-005.}\par
\smallskip\noindent\emph{Sravnjenje s tekstom:} Kvalifikator absolutna imaje lokalnu definicijsku funkciju.\par
\noindent\emph{Textual comparison:} The qualifier absolute has a specific role in the local definition.\par
\noindent{\footnotesize TEXT-COMP-646; pasaže / passages: 9922, 9924, 9926, 9928, 9930; \texttt{CORPUS\_FORM\_COMPARISONS.json}.}\par
\end{minipage}\par
\par\medskip\Needspace{7\baselineskip}
\noindent\begin{minipage}{\linewidth}
\hypertarget{TERM-1590}{}
\noindent\textbf{Rimanove normalne koordinaty}\quad{\footnotesize\texttt{TERM-1590}}\par
\noindent\emph{Iztočnik / Source:} Riemannsche Normalkoordinaten\par
\noindent\emph{Rusky / Russian:} римановы нормальные координаты\par
\noindent\emph{Ukrajinsky / Ukrainian:} ріманові нормальні координати\par
\smallskip\noindent\emph{Razlog:} Rimanove normalne koordinaty ohranjajut geometricky naziv i eponim. Tri arhivne zapisy svezujut izbor s upotrěboju v rabotah 12, 21 i 23, osoblivo s obosnovanjem redukcije invariantov. V rabotě 12 normalne koordinaty prevodet ekstremaly iz fiksovanoj točky v prave linije; v tom leži smysl naziva, ne v normalizaciji dlžiny vektora. Zapisy o kirilici govorjat i o transliteraciji, i o ohranjanju vlašćnyh imen. Bez sravnjenja konkretnyh kiriličnyh měst te skraćene opisy ne dokazyvajut jednu detaljnu politiku. Pri OLP ne trěba iz geometričnogo normalne vyvoditi izbor za normalnu formu formula: kvalifikator imaje svoju funkciju v vsakom disciplinarnom kontekstu.\par
\noindent\emph{Reason:} Rimanove normalne koordinaty retains both the geometric term and the eponym. Three records connect the choice with Papers 12, 21, and 23, particularly the justification for reducing invariants. In Paper 12, normal coordinates turn extremals through a fixed point into straight lines; that is the relevant meaning, not normalization of a vector's length. The Cyrillic notes refer both to transliteration and to preserving proper names. Without comparison of the actual Cyrillic passages, those abbreviated descriptions do not establish one detailed policy. Nor should geometric normalne determine the term for a normal form of formulas in OLP: the qualifier has its own function in each disciplinary context.\par
\noindent{\footnotesize Izvorne zapisy / Original records: HIST-g074-014, HIST-g084-009, HIST-g094-016.}\par
\smallskip\noindent\emph{Sravnjenje s tekstom:} Smysl normalnosti ide iz konstrukcije ekstremalov, ne iz občeho normalizovanja.\par
\noindent\emph{Textual comparison:} Normality derives from the construction using extremals, not generic normalization.\par
\noindent{\footnotesize TEXT-COMP-647; pasaže / passages: 9105, 9116, 9118, 9120, 9122, 14145, 14149, 14151, 15201, 15203, 15205; \texttt{CORPUS\_FORM\_COMPARISONS.json}.}\par
\end{minipage}\par
\par\medskip\Needspace{7\baselineskip}
\noindent\begin{minipage}{\linewidth}
\hypertarget{TERM-1591}{}
\noindent\textbf{kolco vsih cělyh algebraičnyh čisel}\quad{\footnotesize\texttt{TERM-1591}}\par
\noindent\emph{Iztočnik / Source:} Ring aller ganzen algebraischen Zahlen\par
\noindent\emph{Rusky / Russian:} кольцо всех целых алгебраических чисел\par
\noindent\emph{Ukrajinsky / Ukrainian:} кільце всіх цілих алгебраїчних чисел\par
\smallskip\noindent\emph{Razlog:} Kolco vsih cělyh algebraičnyh čisel čuva razvity opis oblasti koeficientov prěd kratkym nazivom prsten cělyh čisel. Arhiv davaje tomu opisu prědnost, zato že autor izjavjaje, ktore čisla prinadležet k prstenu. Slovo vsih jest ale ograničeno kontekstom: v §7 raboty 24 ide o cěle algebraične čisla danogo konečnogo algebraičnogo čislovogo polja $K$, ne o vse take čisla bez ograničenja polja. Nynějšny korpus uživaje prsten na městu kolco. Važno jest ohranjati i algebraičny uslov, i polje, v ktorom on jest razgledany. Kratšejši naziv može byti užitečny poslě definicije, ale ne směje skryti drugu oblast koeficientov.\par
\noindent\emph{Reason:} Kolco vsih cělyh algebraičnyh čisel preserves an expanded description of the coefficient domain instead of the shorter ring-of-integers name. The record favours that description because the author specifies which numbers belong to the ring. Vsih, however, is bounded by context: Paper 24, §7 concerns the algebraic integers of the specified finite algebraic number field $K$, not all algebraic integers without a field restriction. The current corpus uses prsten where the record has kolco. Both the integrality condition and its ambient field must remain explicit. A shorter name may be useful after definition, but must not conceal a different coefficient domain.\par
\noindent{\footnotesize Izvorne zapisy / Original records: HIST-g101-004.}\par
\smallskip\noindent\emph{Sravnjenje s tekstom:} Slovo vsih ostavaje ograničeno danym čislovym poljem.\par
\noindent\emph{Textual comparison:} All remains restricted to the specified number field.\par
\noindent{\footnotesize TEXT-COMP-648; pasaže / passages: 15725, 15727, 15729, 15739, 15745, 15755; \texttt{CORPUS\_FORM\_COMPARISONS.json}.}\par
\end{minipage}\par
\par\medskip\Needspace{7\baselineskip}
\noindent\begin{minipage}{\linewidth}
\hypertarget{TERM-1592}{}
\noindent\textbf{kolco matric / matrične jedinice}\quad{\footnotesize\texttt{TERM-1592}}\par
\noindent\emph{Iztočnik / Source:} Ring der Matrizes / Matrizeneinheiten\par
\noindent\emph{Rusky / Russian:} кольцо матриц / матричные единицы\par
\noindent\emph{Ukrajinsky / Ukrainian:} кільце матриць / матричні одиниці\par
\smallskip\noindent\emph{Razlog:} Kolco matric / matrične jedinice svezuje naziv prstena s jeho koordinatnym opisom. Arhiv podkrepljaje soglasnost treh jezyčnyh zapisov i kontinuitet algebraičnoj konstrukcije; nynějšny medžuslovjansky tekst uživaje prsten. V §14 raboty 34 množenje matričnyh jedinic zavisi od srednjih indeksov: pri jih soglasnosti nastavaje odgovarjajuča jedinica, pri nesoglasnosti nula. Zato jedinice ne znači, že vsi ti elementi sut obratljive; to ne jest sinonim grupy jedinic prstena. Takože koeficientny objekt v tom opisu može byti nekomutativno tělo. Prěklad dolžen ohranjati te razlike, a ne neprimětno zaměniti teoriju matric nad tělom teorijeju nad komutativnym poljem.\par
\noindent\emph{Reason:} Kolco matric / matrične jedinice connects the name of the ring with its coordinate description. The record emphasizes agreement among the three language records and continuity of the algebraic construction; the current Interslavic text uses prsten. In Paper 34, §14, multiplication of matrix units depends on the middle indices: matching indices give the corresponding unit, mismatching indices give zero. Jedinice therefore does not mean that all these elements are invertible; it is not another name for the ring's group of units. The coefficient object may also be a noncommutative division ring. Translation must preserve these distinctions, not silently substitute matrix theory over a commutative field.\par
\noindent{\footnotesize Izvorne zapisy / Original records: HIST-g192-002.}\par
\smallskip\noindent\emph{Sravnjenje s tekstom:} Matrična jedinica ne jest obče obratlivy element prstena.\par
\noindent\emph{Textual comparison:} A matrix unit is not generally an invertible ring element.\par
\noindent{\footnotesize TEXT-COMP-649; pasaže / passages: 21116, 21118, 21120, 21130, 21132, 21136, 21138, 21140, 21142; \texttt{CORPUS\_FORM\_COMPARISONS.json}.}\par
\end{minipage}\par
\par\medskip\Needspace{7\baselineskip}
\noindent\begin{minipage}{\linewidth}
\hypertarget{TERM-1593}{}
\noindent\textbf{kolco v abstraktnoj definiciji}\quad{\footnotesize\texttt{TERM-1593}}\par
\noindent\emph{Iztočnik / Source:} Ring in abstrakter Definition\par
\noindent\emph{Rusky / Russian:} кольцо в абстрактном определении\par
\noindent\emph{Ukrajinsky / Ukrainian:} кільце в абстрактному означенні\par
\smallskip\noindent\emph{Razlog:} Kolco v abstraktnoj definiciji bylo izbrano skoro doslovno, zato že autor ne tolko imenuje prsten, ale daje jeho aksiomy. V rabotě 19 definicija potrěbuje komutativnost množenja i specifičnu aditivnu strukturu; jediničny element ne jest potrěbny a nulove dělitelji sut dozvoljene. Kratky moderny naziv s inačejšimi konvencijami ne zaměnjuje toj definicije. Nynějšny prsten ne měnjaje etih uslovij. Iztočny tekst takože pripisuje prědhodnu definiciju Fraenkelu a objasnjaje odstranjenje něktoryh ograničenij; to jest avtorska pohođenost, ne dostignjenje prěklada. Dva historijske zapisy dokumentujut kontinuitet naziva, ne dva nezavisne dokazy pravilnosti vsakoj aksiomy.\par
\noindent\emph{Reason:} Kolco v abstraktnoj definiciji was kept close to the source because the author supplies ring axioms rather than merely naming the structure. In Paper 19, the definition requires commutative multiplication and a specified additive structure; an identity element is not required and zero divisors are permitted. A compact modern term with different conventions cannot replace that definition. The current name prsten does not change these conditions. The source also credits Fraenkel's earlier definition and explains the removal of restrictions; that is inherited mathematical provenance, not an achievement of the translation. Two historical records document continuity of naming, not two independent checks of every axiom.\par
\noindent{\footnotesize Izvorne zapisy / Original records: HIST-g081-018, HIST-g082-018.}\par
\smallskip\noindent\emph{Sravnjenje s tekstom:} Definiciju određajut izrečene aksiomy, ne konvencije dodane kratkym modernym nazivom.\par
\noindent\emph{Textual comparison:} The stated axioms govern the definition, not conventions imported by a shorter modern label.\par
\noindent{\footnotesize TEXT-COMP-650; pasaže / passages: 12476, 12478, 12479, 12480, 12481, 12482, 12483, 12485; \texttt{CORPUS\_FORM\_COMPARISONS.json}.}\par
\end{minipage}\par
\par\medskip\Needspace{7\baselineskip}
\noindent\begin{minipage}{\linewidth}
\hypertarget{TERM-1594}{}
\noindent\textbf{kolco bez jediničnogo elementa}\quad{\footnotesize\texttt{TERM-1594}}\par
\noindent\emph{Iztočnik / Source:} Ring ohne Einheitselement\par
\noindent\emph{Rusky / Russian:} кольцо без единичного элемента\par
\noindent\emph{Ukrajinsky / Ukrainian:} кільце без одиничного елемента\par
\smallskip\noindent\emph{Razlog:} Kolco bez jediničnogo elementa ohranjaje vidimu ograničenost prědpoloženja. V dodatku k §2 raboty 30 izpuščenje potrěbovanja o jedinici měnjaje opis poradženogo modula: kromě prstenovyh množiteljev trěba dopustiti cěločiselne mnogokratnosti poradžajučih elementov. Arhiv zato ne zaměnil tu frazu jednoslovnym nazivom. Važno jest razlikovati prsten, za ktory jedinica ne jest prědpoložena, od prstena, za ktory se tvrdži, že jedinicy ne imaje. Oboje ne jest identično ni s tvrdženjem, že nijeden element ne jest obratlivy. Nynějšny tekst imenuje prsten, ne kolco; nasledovanie naziva trěba provoditi zajedno s pravilnym dosegom aksiomy.\par
\noindent\emph{Reason:} Kolco bez jediničnogo elementa keeps the hypothesis visible. In the addendum to Paper 30, §2, dropping the identity assumption changes the description of a generated module: integer multiples of the generators are needed in addition to multiples by ring elements. The record consequently retains the explicit phrase rather than replacing it with a one-word label. A ring in which an identity is not assumed must be distinguished from a ring asserted to have no identity. Neither expression simply says that no individual element is invertible. The current text says prsten, not kolco; continuity of naming must be checked together with the exact scope of the axiom.\par
\noindent{\footnotesize Izvorne zapisy / Original records: HIST-g109-009.}\par
\smallskip\noindent\emph{Sravnjenje s tekstom:} Izpuščena aksioma měnjaje opis poradženogo modula.\par
\noindent\emph{Textual comparison:} Dropping the axiom changes the description of the generated module.\par
\noindent{\footnotesize TEXT-COMP-651; pasaže / passages: 16346, 16348, 16350; \texttt{CORPUS\_FORM\_COMPARISONS.json}.}\par
\end{minipage}\par
\par\medskip\Needspace{7\baselineskip}
\noindent\begin{minipage}{\linewidth}
\hypertarget{TERM-1595}{}
\noindent\textbf{kolco bez nulovyh děliteljev}\quad{\footnotesize\texttt{TERM-1595}}\par
\noindent\emph{Iztočnik / Source:} Ring ohne Nullteiler\par
\noindent\emph{Rusky / Russian:} кольцо без делителей нуля\par
\noindent\emph{Ukrajinsky / Ukrainian:} кільце без дільників нуля\par
\smallskip\noindent\emph{Razlog:} Kolco bez nulovyh děliteljev ponehava izrečeny uslov, vměsto da avtomatično poda kratko oblast cělosti. Arhiv opiraje to na kontinuitet terminologije v rabotah 29–31. V §4, punktu 7 raboty 31 taj uslov vystupaje pri sravnjenju ranga ideala s rangom prstena i pri opisah slěda i normy. Slovnik čuva motivaciju za eksplicitnu frazu; ona ne dokazyvaje, že moderny naziv nikde ne jest prikladny. Drugy zapis za taj sam německy izraz, oblast cělosti, sviděteljstvuje o drugoj redakcijnoj strategiji. Oba izbory trěba sravniti po kontekstu, osoblivo po komutativnosti i jedinici, prěd ustanovjenjem jedinoj politiki za OLP.\par
\noindent\emph{Reason:} Kolco bez nulovyh děliteljev retains the stated condition instead of automatically shortening it to oblast cělosti. The record bases this on continuity of zero-divisor terminology across Papers 29–31. In Paper 31, §4, item 7, the condition accompanies comparisons of ideal and ring rank and the treatment of trace and norm. The dictionary preserves the reason for the explicit phrase without treating it as proof that modern domain terminology is never appropriate. Another record for the same German expression chooses oblast cělosti and thus represents a different editorial strategy. Both choices require contextual comparison, especially concerning commutativity and identity, before one policy is established for OLP.\par
\noindent{\footnotesize Izvorne zapisy / Original records: HIST-g136-004.}\par
\smallskip\noindent\emph{Sravnjenje s tekstom:} Literalna fraza i modernizovany naziv vyžadujut kontekstualno sravnjenje.\par
\noindent\emph{Textual comparison:} The literal phrase and modernized name require contextual comparison.\par
\noindent{\footnotesize TEXT-COMP-652; pasaže / passages: 18057, 18059, 18061, 18063, 18065, 18067; \texttt{CORPUS\_FORM\_COMPARISONS.json}.}\par
\end{minipage}\par
\par\medskip\Needspace{7\baselineskip}
\noindent\begin{minipage}{\linewidth}
\hypertarget{TERM-1596}{}
\noindent\textbf{oblast cělosti}\quad{\footnotesize\texttt{TERM-1596}}\par
\noindent\emph{Iztočnik / Source:} Ring ohne Nullteiler\par
\noindent\emph{Rusky / Russian:} область целостности\par
\noindent\emph{Ukrajinsky / Ukrainian:} область цілісності\par
\smallskip\noindent\emph{Razlog:} Oblast cělosti jest svěstna modernizacija německogo Ring ohne Nullteiler za koeficientny prsten $Z$. Arhiv ju opravdava sodobnym matematičnym registrom a ohranjanjem iztočnogo izraza v aparatu. V danom početku §4 raboty 31 jediničny element osnovnogo prstena prinadleži k $Z$; tekst potom postavjaje dopolniteljne uslovja o konečnyh slobodnyh modulah. Te daljše uslovja ne sut čest definicije vsakoj oblasti cělosti. Izbor jest zato prikladny tolko s ustanovjenymi konvencijami. On ne odpravjaje susědny literalny izbor bez nulovyh děliteljev; razlika medžu modernizacijeju i věrnym ohranjanjem razvitogo uslova ostavaje predmetom redakcijnogo rěšenja.\par
\noindent\emph{Reason:} Oblast cělosti is a deliberate modernization of Ring ohne Nullteiler for the coefficient ring $Z$. The record appeals to current mathematical register while retaining the source expression in the apparatus. At this opening of Paper 31, §4, the ambient ring's identity belongs to $Z$; further assumptions on finite free modules then follow. Those additional conditions are not part of the definition of every integral domain. The choice is therefore appropriate only with the conventions established. It does not erase the neighbouring literal choice bez nulovyh děliteljev: the difference between modernization and preserving an explicit hypothesis remains an editorial decision.\par
\noindent{\footnotesize Izvorne zapisy / Original records: HIST-g130-006.}\par
\smallskip\noindent\emph{Sravnjenje s tekstom:} Uslovja o konečnyh slobodnyh modulah ne prinadležet avtomatično definiciji oblasti cělosti.\par
\noindent\emph{Textual comparison:} Finite-free module assumptions do not automatically belong to the definition of integral domain.\par
\noindent{\footnotesize TEXT-COMP-653; pasaže / passages: 17823, 17825, 17827, 17829; \texttt{CORPUS\_FORM\_COMPARISONS.json}.}\par
\end{minipage}\par
\par\medskip\Needspace{7\baselineskip}
\noindent\begin{minipage}{\linewidth}
\hypertarget{TERM-1597}{}
\noindent\textbf{kolco bez děliteljev nula i s jediničnym elementom}\quad{\footnotesize\texttt{TERM-1597}}\par
\noindent\emph{Iztočnik / Source:} Ring ohne Nullteiler und mit Einheitselement\par
\noindent\emph{Rusky / Russian:} кольцо без делителей нуля и с единицей\par
\noindent\emph{Ukrajinsky / Ukrainian:} кільце без дільників нуля і з одиничним елементом\par
\smallskip\noindent\emph{Razlog:} Kolco bez děliteljev nula i s jediničnym elementom izjavjaje dva uslova oddělno: odsutnost nulovyh děliteljev i obstajanje jedinicy. Arhiv im pripisyvaje visoku razumljivost medžu slovjanskimi matematičnymi registrami, ale ne podava empiricky test ili čislove ocěnky. V §8 raboty 31 te dva uslova ješče ne izčerpajut definiciju razgledanogo prstena: slědujut i potrěbovanja o razloženju idealov. Nynějšny tekst uživaje prsten, ale čuva izraz děliteljev nula. To ne rěšaje padežnu razliku prěd formami nuly i nulovyh drugde v korpusu; obosnovanje korenja ne jest obosnovanje vsakoj fleksije. Slovnik zapisuje razliku bez izmyslenogo historijskogo gramatičnogo obosnovanja.\par
\noindent\emph{Reason:} Kolco bez děliteljev nula i s jediničnym elementom states two conditions separately: absence of zero divisors and existence of an identity. The record attributes high cross-Slavic transparency to them, but supplies neither an empirical test nor numerical scores. In Paper 31, §8, those two conditions do not exhaust the definition of the ring under discussion: requirements on ideal factorization follow. The current text uses prsten but retains děliteljev nula. That does not settle its case-form difference from nuly and nulovyh elsewhere in the corpus; justification of a lexical stem does not justify every inflection. The dictionary records the difference without inventing a historical grammatical rationale.\par
\noindent{\footnotesize Izvorne zapisy / Original records: HIST-g155-007.}\par
\smallskip\noindent\emph{Sravnjenje s tekstom:} Prsten bez nulovyh děliteljev i s jediniceju ješče ne izpolnjaje avtomatično vsih slědujučih aksiomov.\par
\noindent\emph{Textual comparison:} A ring with identity and no zero divisors does not automatically satisfy all the further axioms.\par
\noindent{\footnotesize TEXT-COMP-654; pasaže / passages: 18785, 18787, 18789, 16603; \texttt{CORPUS\_FORM\_COMPARISONS.json}.}\par
\end{minipage}\par
\par\medskip\Needspace{7\baselineskip}
\noindent\begin{minipage}{\linewidth}
\hypertarget{TERM-1598}{}
\noindent\textbf{kolco bez radikala}\quad{\footnotesize\texttt{TERM-1598}}\par
\noindent\emph{Iztočnik / Source:} Ring ohne Radikal\par
\noindent\emph{Rusky / Russian:} кольцо без радикала\par
\noindent\emph{Ukrajinsky / Ukrainian:} кільце без радикала\par
\smallskip\noindent\emph{Razlog:} Kolco bez radikala ohranjaje uslov o idealu, ne zaměnjaje go opisom vsih elementov. V §12 raboty 34 radikal jest maksimalny nilpotentny dvustranny ideal pri danom uslovju cěpov. Prsten bez radikala imaje taj ideal ravny nuly; to ne znači, že v njem ne može byti nijednogo nenulovogo nilpotentnogo elementa. Zato naziv ne jest avtomatično sinonim modernogo redukovanogo prstena. Drugy arhivny zapis prodolžaje uslov pri priměnjenju k grupovym prstenam, kde važna jest karakteristika koeficientnogo polja. Nynějšny naziv prsten ne měnjaje nijednu od tih granic. Svez s polnoju razložimostju treba čitati po konkretnyh teoremah, ne izvoditi samo iz naziva.\par
\noindent\emph{Reason:} Kolco bez radikala preserves a condition on an ideal, not a description of every element. In Paper 34, §12, the radical is the maximal nilpotent two-sided ideal under the stated chain condition. A radical-free ring has that ideal equal to zero; this does not say that it cannot contain any nonzero nilpotent element. The name is therefore not automatically synonymous with a modern reduced ring. A second record carries the condition into the group-ring application, where the characteristic of the coefficient field matters. The current name prsten changes none of these boundaries. Relations with complete reducibility must be read from the particular theorems, not inferred from the name alone.\par
\noindent{\footnotesize Izvorne zapisy / Original records: HIST-g190-003, HIST-g204-003.}\par
\smallskip\noindent\emph{Sravnjenje s tekstom:} Nulovy radikal i odsutnost vsih nenulovyh nilpotentnyh elementov ne sut isto potrěbovanje.\par
\noindent\emph{Textual comparison:} Zero radical and absence of all nonzero nilpotent elements are not the same requirement.\par
\noindent{\footnotesize TEXT-COMP-655; pasaže / passages: 20860, 20862, 20864, 20866, 21138, 22447, 22449, 22451, 22453, 22455, 22457; \texttt{CORPUS\_FORM\_COMPARISONS.json}.}\par
\end{minipage}\par
\par\medskip\Needspace{7\baselineskip}
\noindent\begin{minipage}{\linewidth}
\hypertarget{TERM-1599}{}
\noindent\textbf{kolco konečnogo ranga vzhodno k podpolju $P$}\quad{\footnotesize\texttt{TERM-1599}}\par
\noindent\emph{Iztočnik / Source:} Ring von endlichem Rang in bezug auf einen Unterkörper $P$\par
\noindent\emph{Rusky / Russian:} кольцо конечного ранга относительно подполя $P$\par
\noindent\emph{Ukrajinsky / Ukrainian:} кільце скінченного рангу відносно підполя $P$\par
\smallskip\noindent\emph{Razlog:} Kolco konečnogo ranga vzhodno k podpolju $P$ čuva i konečnost ranga, i polje, po odnošenju k ktoromu rang jest izčisljeny. Arhiv svezuje izbor s ranějšeju terminologijeju ranga i s prozrěčnym složenjem pod-polje, za ktoro rusky i ukrajinsky zapisy dajut slične nazivy. To ne znači konečno čislo elementov prstena. Teorem v §6, punktu 4 raboty 31 razgledaje taj relativny rang pri uslovjah o diskriminantu i razložimosti. Nynějšny tekst imaje prsten i po odnošenju k, ne arhivno kolco i vzhodno k. Slovnik ohranjaje istoriju izbora, ale pri novoj prěkladnoj rabotě ne rekomenduje vozvračanje k staroj prědložnoj konstrukciji.\par
\noindent\emph{Reason:} Kolco konečnogo ranga vzhodno k podpolju $P$ retains both finite rank and the field relative to which rank is measured. The record connects the choice with earlier rank terminology and the transparent compound pod-polje, supported by similar Russian and Ukrainian names. It does not mean that the ring has finitely many elements. The theorem in Paper 31, §6, item 4 concerns this relative rank together with discriminant and reducibility conditions. The current text uses prsten and po odnošenju k rather than archived kolco and vzhodno k. The dictionary preserves the history of the choice without recommending a return to the older prepositional construction in new translation.\par
\noindent{\footnotesize Izvorne zapisy / Original records: HIST-g144-002.}\par
\smallskip\noindent\emph{Sravnjenje s tekstom:} Rang zavisi od danej koeficientnoj oblasti; ne jest čislo vsih elementov.\par
\noindent\emph{Textual comparison:} Rank is relative to the specified coefficient field; it is not the cardinality of the ring.\par
\noindent{\footnotesize TEXT-COMP-656; pasaže / passages: 18408, 18410, 18412, 18414; \texttt{CORPUS\_FORM\_COMPARISONS.json}.}\par
\end{minipage}\par
\par\medskip\Needspace{7\baselineskip}
\noindent\begin{minipage}{\linewidth}
\hypertarget{TERM-1600}{}
\noindent\textbf{kolcova adjunkcija}\quad{\footnotesize\texttt{TERM-1600}}\par
\noindent\emph{Iztočnik / Source:} Ring-Adjunktion\par
\noindent\emph{Rusky / Russian:} присоединение как кольцу\par
\noindent\emph{Ukrajinsky / Ukrainian:} приєднання як кільця\par
\smallskip\noindent\emph{Razlog:} Kolcova adjunkcija izpolzovaje internationalny korenj s kvalifikatorom, ktory kaže, čto jest konstruovano. Arhivny rusky i ukrajinsky zapisy izbirajut obično opisno prisojedinjenje; toj razhod ne jest skryt za izmyslenym edinym slovjanskim nazivom. V §1 raboty 31 prstenova adjunkcija davaje poradženy podprsten: polinomy v prisojeđenih elementah s ohranjanymi relacijami okolnogo prstena. Ne davaje avtomatično polje niti obratne elementy vsih nenulovyh elementov. Takože to ne jest kategorična adjunkcija funktorov. Odločujuči jest smyslovy kvalifikator prstenova v nynějšnom tekste; staro kolcova ostavaje arhivna forma. Prosta zaměna slova adjunkcija v drugoj disciplině izgubila by razliku konstrukcij.\par
\noindent\emph{Reason:} Kolcova adjunkcija uses an international stem with a qualifier identifying what is constructed. The Russian and Ukrainian records instead choose ordinary descriptive adjoining language; that divergence is not concealed behind an invented single Slavic label. In Paper 31, §1, ring adjunction gives the generated subring: polynomials in the adjoined elements, subject to the relations of the ambient ring. It does not automatically give a field or inverses for all nonzero elements. Nor is it a categorical adjunction of functors. The meaningful qualifier is decisive: the current text has prstenova, while kolcova remains the archived form. Replacing every occurrence of adjunction in another discipline would lose the distinction between constructions.\par
\noindent{\footnotesize Izvorne zapisy / Original records: HIST-g119-004.}\par
\smallskip\noindent\emph{Sravnjenje s tekstom:} Poradženy prsten ohranjaje relacije okolnogo prstena; ne jest avtomatično polje.\par
\noindent\emph{Textual comparison:} The generated ring retains ambient relations; it is not automatically a field.\par
\noindent{\footnotesize TEXT-COMP-657; pasaže / passages: 17302, 17304, 17310, 17312, 17314, 17316, 17318, 17320, 17322, 17324; \texttt{CORPUS\_FORM\_COMPARISONS.json}.}\par
\end{minipage}\par
\par\medskip\Needspace{7\baselineskip}
\noindent\begin{minipage}{\linewidth}
\hypertarget{TERM-1601}{}
\noindent\textbf{kolceva oblast}\quad{\footnotesize\texttt{TERM-1601}}\par
\noindent\emph{Iztočnik / Source:} Ringbereich\par
\noindent\emph{Rusky / Russian:} кольцевая область\par
\noindent\emph{Ukrajinsky / Ukrainian:} кільцева область\par
\smallskip\noindent\emph{Razlog:} Kolceva oblast doslovno ohranjaje dvočlenny německy naziv Ringbereich. Arhiv tomu dava prednost, aby ne suzil pojem prěd pozdnějšim razširjenjem k modulam i nekomutativnym oblastjam. Oblast tu ne znači automatno oblast cělosti: v §1 raboty 19 nulove dělitelje sut dozvoljene i jedinica ne jest potrěbna; §9 razděljuje prstenovu oblast i oblast, na ktoro ona dějstvuje. Rusko i ukrajinsko sravnjenje podpirajut formu frazy, ne dodavaju nove aksiomy. Nynějšny prěvod ima prstenova oblast. Slovnik ohranjaje staru formu i jeji razlog, ale pri novoj upotrěbě trěba čitati definiciju oblasti, ne izvoditi odsutnost nulovyh děliteljev iz jedinogo slova oblast.\par
\noindent\emph{Reason:} Kolceva oblast closely preserves the two-part name Ringbereich. The recorded reason is to avoid narrowing the concept before the later treatment of modules and noncommutative structures. Oblast does not automatically mean integral domain here: Paper 19, §1 permits zero divisors and does not require an identity; §9 distinguishes the ring from the structure on which it acts. The Russian and Ukrainian comparanda support the construction of the phrase, not additional axioms. The current translation has prstenova oblast. The dictionary preserves the earlier form and its reason, but future use must follow the local definition, not infer absence of zero divisors from oblast alone.\par
\noindent{\footnotesize Izvorne zapisy / Original records: HIST-g081-001, HIST-g082-001.}\par
\smallskip\noindent\emph{Sravnjenje s tekstom:} Naziv oblasti ne dodava odsutnost nulovyh děliteljev ni komutativnost v pozdnějšem obobčenju.\par
\noindent\emph{Textual comparison:} The domain label adds neither absence of zero divisors nor commutativity in the later generalization.\par
\noindent{\footnotesize TEXT-COMP-658; pasaže / passages: 12432, 12440, 12474, 12476, 12485, 13383, 13385, 13387, 13388, 13390; \texttt{CORPUS\_FORM\_COMPARISONS.json}.}\par
\end{minipage}\par
\par\medskip\Needspace{7\baselineskip}
\noindent\begin{minipage}{\linewidth}
\hypertarget{TERM-1602}{}
\noindent\textbf{kolca konečnogo ranga}\quad{\footnotesize\texttt{TERM-1602}}\par
\noindent\emph{Iztočnik / Source:} Ringe endlichen Ranges\par
\noindent\emph{Rusky / Russian:} кольца конечного ранга\par
\noindent\emph{Ukrajinsky / Ukrainian:} кільця скінченного рангу\par
\smallskip\noindent\emph{Razlog:} Kolca konečnogo ranga prodolžuje ranějši naziv i izrično ohranjaje rang kako modulny pojem. V §2 raboty 31 rang jest maksimalno čislo linearno nezavisnyh elementov nad danym poljem; ne jest čislo vseh elementov prstena. §4 potom razširjaje prědpoloženja, aby objel i poredky: baza i nezavisnost sut tam odnosne k oblasti cělosti. Arhiv čuva jednu frazu prěz toj prěhod, ne tvrdži, že vsi osnovni objekty sut polja. Ruska i ukrajinska forma davajut paralelu konečnosti i ranga. Nynějšny korpus piše prsteny, ne kolca; pri standardizaciji trěba ohraniti osnovnu oblast i vsako dodatno uslovje o bazi.\par
\noindent\emph{Reason:} Kolca konečnogo ranga continues the earlier terminology while retaining rank as a module-theoretic condition. In Paper 31, §2, rank is the largest number of linearly independent elements over the specified field, not the number of all ring elements. Section 4 broadens the assumptions to include orders: basis and independence are then relative to an integral domain. The records preserve the phrase across that transition without saying that every base is a field. The Russian and Ukrainian forms parallel both finiteness and rank. Current Interslavic uses prsteny rather than kolca; standardization must retain the coefficient domain and every additional basis hypothesis.\par
\noindent{\footnotesize Izvorne zapisy / Original records: HIST-g121-001, HIST-g130-004.}\par
\smallskip\noindent\emph{Sravnjenje s tekstom:} Konečnost ranga jest različna od konečnosti množstva elementov.\par
\noindent\emph{Textual comparison:} Finite rank differs from finiteness of the set of elements.\par
\noindent{\footnotesize TEXT-COMP-659; pasaže / passages: 17432, 17434, 17436, 17823, 17825, 17827, 17829; \texttt{CORPUS\_FORM\_COMPARISONS.json}.}\par
\end{minipage}\par
\par\medskip\Needspace{7\baselineskip}
\noindent\begin{minipage}{\linewidth}
\hypertarget{TERM-1603}{}
\noindent\textbf{kolca s jediničnym elementom / kolco s jedinico}\quad{\footnotesize\texttt{TERM-1603}}\par
\noindent\emph{Iztočnik / Source:} Ringe mit Einheitselement / Ring mit Einheit\par
\noindent\emph{Rusky / Russian:} кольца с единицей / кольцо с единицей\par
\noindent\emph{Ukrajinsky / Ukrainian:} кільця з одиничним елементом / кільце з одиницею\par
\smallskip\noindent\emph{Razlog:} Polny izraz s jediničnym elementom v naslovu i kratky s jedinico v tekučoj prozě sut dvě formy togo samogo uslovja. Arhiv objašnjaje kratkost ruskoj formy i polnějše ukrajinske i medžuslovjanske zaglavje; toj izbor jest registr i ekonomija izloženja, ne različna definicija jedinice. §17 raboty 34 razděljuje modul na čest, kde jedinica dějstvuje kako identičny operator, i čest, kde ona dějstvuje kako nula. Zato prsten s jedinico sam ne znači, že vsak razgledany modul už ima identično dějstvovanje jedinice. Ranějša prednost kolco prěd formoju ring jest zachovana kako historijno obrazloženje; aktualny naziv jest prsten.\par
\noindent\emph{Reason:} The full heading s jediničnym elementom and the shorter prose form s jedinico express the same assumption. The record explains the short Russian wording and the fuller Ukrainian and Interslavic headings as choices of register and economy, not different definitions of identity. Paper 34, §17 decomposes a module into a part where the ring identity acts as the identity operator and a part where it acts as zero. Thus a ring with identity does not by itself mean that every module being considered already has unital action. The earlier preference for kolco over a borrowed ring is preserved as historical reasoning; current terminology is prsten.\par
\noindent{\footnotesize Izvorne zapisy / Original records: HIST-g195-001.}\par
\smallskip\noindent\emph{Sravnjenje s tekstom:} Jedinica prstena ne dějstvuje potrěbno identično na cělom modulu.\par
\noindent\emph{Textual comparison:} The ring identity need not act identically on the entire module.\par
\noindent{\footnotesize TEXT-COMP-660; pasaže / passages: 21585, 21587, 21589, 21591, 21595, 21605, 21607; \texttt{CORPUS\_FORM\_COMPARISONS.json}.}\par
\end{minipage}\par
\par\medskip\Needspace{7\baselineskip}
\noindent\begin{minipage}{\linewidth}
\hypertarget{TERM-1604}{}
\noindent\textbf{kolca s jediničnym elementom i bez nuljevyh děliteljev}\quad{\footnotesize\texttt{TERM-1604}}\par
\noindent\emph{Iztočnik / Source:} Ringe mit Einheitselement und ohne Nullteiler\par
\noindent\emph{Rusky / Russian:} кольца с единицей и без делителей нуля\par
\noindent\emph{Ukrajinsky / Ukrainian:} кільця з одиницею і без дільників нуля\par
\smallskip\noindent\emph{Razlog:} Kolca s jediničnym elementom i bez nuljevyh děliteljev jest rezultat zabilěženoj zaměny provizornogo ringy na tada upotrěbljano kolca. Arhiv se opira na doslědnost s rabotami 19 i 24, ne na čislenu ocěnku kandidatov. Dva dodatka morajut ostati odvojeno vidime: egzistencija jedinice i odsutnost nulovyh děliteljev. Kratka rabota 26 k nim dodava cěpove uslovja i integralnu zamknjenost; prva fraza ne sodržuje cělu karakterizaciju. Nynějšny korpus ima prsteny i ohranjaje nuljevyh na tom městu. Razlika s nulovyh v drugyh pasažah jest stvarna pravopisna ili tvorbena razlika za sravnjenje, ne dokaz različnyh matematičnyh pojmov.\par
\noindent\emph{Reason:} Kolca s jediničnym elementom i bez nuljevyh děliteljev records a change from provisional ringy to the then-used kolca. Its stated basis is consistency with Papers 19 and 24, not numerical candidate scores. The two qualifications must remain separately visible: existence of an identity and absence of zero divisors. The brief Paper 26 adds chain conditions and integral closedness, so the opening phrase does not contain the whole characterization. The current text has prsteny and retains nuljevyh here. Its difference from nulovyh elsewhere is a spelling or formation difference to examine, not evidence of distinct mathematical concepts.\par
\noindent{\footnotesize Izvorne zapisy / Original records: HIST-g103-003.}\par
\smallskip\noindent\emph{Sravnjenje s tekstom:} Odsutnost nulovyh děliteljev i jedinica ne zaměnjajut cěpovyh uslovij i integralnoj zamknjenosti.\par
\noindent\emph{Textual comparison:} Absence of zero divisors and an identity do not replace the chain and integral-closedness conditions.\par
\noindent{\footnotesize TEXT-COMP-661; pasaže / passages: 15878, 15884, 15886, 15888; \texttt{CORPUS\_FORM\_COMPARISONS.json}.}\par
\end{minipage}\par
\par\medskip\Needspace{7\baselineskip}
\noindent\begin{minipage}{\linewidth}
\hypertarget{TERM-1605}{}
\noindent\textbf{kolca ranga jedan}\quad{\footnotesize\texttt{TERM-1605}}\par
\noindent\emph{Iztočnik / Source:} Ringe vom Rang eins\par
\noindent\emph{Rusky / Russian:} кольца ранга один\par
\noindent\emph{Ukrajinsky / Ukrainian:} кільця рангу один\par
\smallskip\noindent\emph{Razlog:} Kolca ranga jedan ohranjaje točnu veličinu ranga. Poslědok v §3 raboty 31 govori o razloženju po razširjenju osnovnogo polja: komponenty sut jednodimenzionalne nad odgovarjajučym poljem i jemu izomorfne. On ne pravi, že početny prsten imaje jedin element, ni že on už jest polje. Prěhod k algebraično zatvorjenomu polju i potom k dostatočnomu konečnomu razširjenju jest čestju argumenta. Rusky i ukrajinsky zapisy takože izgovarjajut čislo, ne nový tip konečnosti. Aktualno prsteny ranga jedan čuva toj smysl; slovo rang trěba čitati s jeho osnovnym poljem.\par
\noindent\emph{Reason:} Kolca ranga jedan retains the precise value of the rank. The corollary in Paper 31, §3 concerns decomposition after extension of the base field: the components are one-dimensional over, and isomorphic to, the appropriate field. It says neither that the initial ring has one element nor that it is already a field. Passage to an algebraically closed field and then to a sufficient finite extension is part of the argument. Russian and Ukrainian likewise state the number rather than a new kind of finiteness. Current prsteny ranga jedan preserves this meaning; rank must be read with its base field.\par
\noindent{\footnotesize Izvorne zapisy / Original records: HIST-g129-002.}\par
\smallskip\noindent\emph{Sravnjenje s tekstom:} Komponenty sut jednodimenzionalne nad odgovarjajučym poljem, ne množstva iz jednogo elementa.\par
\noindent\emph{Textual comparison:} The components are one-dimensional over the appropriate field, not singleton sets.\par
\noindent{\footnotesize TEXT-COMP-662; pasaže / passages: 17781, 17783, 17785, 17786, 17788, 17790, 17792, 17794; \texttt{CORPUS\_FORM\_COMPARISONS.json}.}\par
\end{minipage}\par
\par\medskip\Needspace{7\baselineskip}
\noindent\begin{minipage}{\linewidth}
\hypertarget{TERM-1606}{}
\noindent\textbf{kolcovy homomorfizm od $R$ k $R_a$}\quad{\footnotesize\texttt{TERM-1606}}\par
\noindent\emph{Iztočnik / Source:} Ringhomomorphismus von $R$ zu $R_a$\par
\noindent\emph{Rusky / Russian:} гомоморфизм колец из $R$ в $R_a$\par
\noindent\emph{Ukrajinsky / Ukrainian:} гомоморфізм кілець із $R$ до $R_a$\par
\smallskip\noindent\emph{Razlog:} Kolcovy homomorfizm od $R$ k $R_a$ spojuje ustaljeny naziv preslikanja s izričnym iztočnikom i cěljem. Predložky zato noset směr, ne samo stil. V §4 raboty 31 korpus govori o prstenovom homomorfizmu v matrično prsten, a potom iz njega izvodi aditivne i skalarne relacije slěda. Homomorfizm sam ne jest tvrdženje obratnoj jednoznačnosti; takože matrično preslikanje ne trěba směšati s posljedujučym slědom. Arhivny podindeks a jest ohranjen v zaglavju; tekst upotrěbjaje idealny podindeks. Pri novoj upotrěbě každe preslikanje trěbuje svoje operacije i směr, ne samo povtorjenje slova homomorfizm.\par
\noindent\emph{Reason:} Kolcovy homomorfizm od $R$ k $R_a$ combines the established map name with an explicit source and target. The prepositions therefore carry direction, not merely style. In Paper 31, §4, the current text gives a ring homomorphism into a matrix ring and then derives additive and scalar relations for trace. A homomorphism is not by itself a claim of bijectivity, and the matrix map must not be confused with the subsequent trace. The archived subscript a is preserved in the heading; the text uses an ideal subscript. In a new context, each map needs its own preserved operations and direction, not just reuse of the label homomorfizm.\par
\noindent{\footnotesize Izvorne zapisy / Original records: HIST-g136-003.}\par
\smallskip\noindent\emph{Sravnjenje s tekstom:} Prstenovo preslikanje i iz njega vyvedeny slěd imajut različne role.\par
\noindent\emph{Textual comparison:} The ring map and the trace derived from it have distinct roles.\par
\noindent{\footnotesize TEXT-COMP-663; pasaže / passages: 18057, 18059, 18061, 18063, 18065; \texttt{CORPUS\_FORM\_COMPARISONS.json}.}\par
\end{minipage}\par
\par\medskip\Needspace{7\baselineskip}
\noindent\begin{minipage}{\linewidth}
\hypertarget{TERM-1607}{}
\noindent\textbf{kolcove operacije iz elementov kolca}\quad{\footnotesize\texttt{TERM-1607}}\par
\noindent\emph{Iztočnik / Source:} Ringverknüpfungen aus Ringelementen\par
\noindent\emph{Rusky / Russian:} кольцевые операции из элементов кольца\par
\noindent\emph{Ukrajinsky / Ukrainian:} кільцеві операції з елементів кільця\par
\smallskip\noindent\emph{Razlog:} Kolcove operacije iz elementov kolca prěvodi Verknüpfungen nazvom operacij, aby konstrukcija ostala konkretna. Arhiv odpovědno ne zavodi kategorijny jezyk, ktory ne jest v iztočnoj frazě. V §7 raboty 31 sravnjavajuća pasaža razlikuje matričnu definiciju slěda od definicije prěz konjugovane elementy: druga ne mora byti priměnima k rezidualnomu prstenu. Slěd i diskriminant tam nastavajut prěz prstenove operacije, što objašnjaje jih prěnos homomorfizmom. Kratko operacije ne daje pravo dodati děljenje ili proizvoljne nove operacije. Nynějšna forma jest prstenove operacije; formulacija objašnjaje konstrukciju, ne zavodi novy kategoričny objekt.\par
\noindent\emph{Reason:} Kolcove operacije iz elementov kolca renders Verknüpfungen as operations to keep the construction concrete. The record deliberately avoids introducing categorical language absent from the phrase. In Paper 31, §7, the comparison passage distinguishes the matrix definition of trace from a definition through conjugate elements: the latter need not apply to the residue ring. Trace and discriminant are formed by ring operations, explaining their passage through the homomorphism. The short word operacije does not license adding division or arbitrary new operations. Current wording is prstenove operacije; the expression explains a construction rather than naming a new categorical object.\par
\noindent{\footnotesize Izvorne zapisy / Original records: HIST-g152-005.}\par
\smallskip\noindent\emph{Sravnjenje s tekstom:} Matrična definicija ostavaje relevantna, kogda konjugovane elementy ne davajut potrěbno prědstavjenje.\par
\noindent\emph{Textual comparison:} The matrix definition remains relevant when conjugate elements need not provide the required representation.\par
\noindent{\footnotesize TEXT-COMP-664; pasaže / passages: 18671, 18675, 18677, 18679, 18681; \texttt{CORPUS\_FORM\_COMPARISONS.json}.}\par
\end{minipage}\par
\par\medskip\Needspace{7\baselineskip}
\noindent\begin{minipage}{\linewidth}
\hypertarget{TERM-1608}{}
\noindent\textbf{obrubljenje}\quad{\footnotesize\texttt{TERM-1608}}\par
\noindent\emph{Iztočnik / Source:} Ränderung\par
\noindent\emph{Rusky / Russian:} окаймление\par
\noindent\emph{Ukrajinsky / Ukrainian:} облямування\par
\smallskip\noindent\emph{Razlog:} Obrubljenje upotrěbjaje sliku okraja, podobno ruskomu okajmlenie i ukrajinskomu obljamuvannja. Arhiv ocěnjuje tu sliku kako razumlivu i prosi daljšu jezyčnu prověrku; ne zapisyvaje izvršeny čitatelsky test. V §3 raboty 2 to jest konkretna operacija Clebschovogo prěnosa, ktora veže binarnu relaciju s ternarnoju, ne proizvoljno dodanje okraja objektu. Novo konsultovane zapisy 21388 i 21390 oficialnogo slovnika podpirajut korenj obrųb v smyslah okraja i podšitego kraja oděže. Oni ne zasvědčajut matematičny smysl ni točno odvođeno obrubljenje; to zaglavje v konzultovanoj kopiji ne bylo nalezeno. Obyčny korenj i tehnično rěšenje zato ostavajut različne vrstvy dokaza.\par
\noindent\emph{Reason:} Obrubljenje uses an edging image, like Russian okajmlenie and Ukrainian obljamuvannja. The record judges that image intelligible and requests further linguistic checking; it does not report a completed reader test. In Paper 2, §3, this is a specific operation in Clebsch's transfer from a binary relation to a ternary one, not arbitrary addition of a border. Newly consulted official-dictionary rows 21388 and 21390 support the stem obrųb in the ordinary senses of border and clothing hem. They attest neither the mathematical sense nor the precise derivative obrubljenje, which was not found as a headword in the consulted copy. Ordinary lexical support and the technical proposal therefore remain distinct kinds of evidence.\par
\noindent{\footnotesize Izvorne zapisy / Original records: HIST-g005-010.}\par
\smallskip\noindent\emph{Sravnjenje s tekstom:} Tehnično značenje vyhodi iz sravnjenja binarnoj i ternarnoj relacije.\par
\noindent\emph{Textual comparison:} The technical meaning comes from the binary and ternary relations being compared.\par
\noindent{\footnotesize TEXT-COMP-665; pasaže / passages: 614, 615, 618, 620, 622, 624, 626; \texttt{CORPUS\_FORM\_COMPARISONS.json}.}\par
\end{minipage}\par
\par\medskip\Needspace{7\baselineskip}
\noindent\begin{minipage}{\linewidth}
\hypertarget{TERM-1609}{}
\noindent\textbf{$S(\gamma)$ slěd / $N(\gamma)$ norma vzhodno k glavnomu klasu}\quad{\footnotesize\texttt{TERM-1609}}\par
\noindent\emph{Iztočnik / Source:} $S(\gamma)$ die Spur / $N(\gamma)$ die Norm in bezug auf die Hauptklasse\par
\noindent\emph{Rusky / Russian:} $S(\gamma)$ след / $N(\gamma)$ норма относительно главного класса\par
\noindent\emph{Ukrajinsky / Ukrainian:} $S(\gamma)$ слід / $N(\gamma)$ норма відносно головного класу\par
\smallskip\noindent\emph{Razlog:} Slěd i norma ohranjajut dva matematična invarianta i jih zavisnost od glavnogo klasa. Arhiv dava prednost kontinuitetu s ranějšimi sekcijami, a ne proglasu jednoj frazy za vse upotrěby. V nynějšnom §6 raboty 31 diagonalno prědstavjenje dava $S(\gamma)=c_1+\cdots+c_n$ i $N(\gamma)=c_1\cdots c_n$. Suma i produkt ne sut zaměnjive, a norma tu ne znači metriku ili dolžinu vektora. Korpus už piše po odnošenju k vměsto arhivnogo vzhodno k. Pri skračenju zaglavja trěba ohraniti obě oznaky i relativny klas; bez nego izpustil by se parametr definicije.\par
\noindent\emph{Reason:} Slěd and norma preserve two mathematical invariants and their dependence on the principal class. The recorded preference is continuity with earlier sections, not one phrase for every use. In the current Paper 31, §6, a diagonal representation gives $S(\gamma)=c_1+\cdots+c_n$ and $N(\gamma)=c_1\cdots c_n$. Sum and product are not interchangeable, and norm here does not mean a metric or vector length. The corpus now uses po odnošenju k instead of archived vzhodno k. A shortened heading must retain both symbols and the relative class; otherwise it loses a parameter of the definition.\par
\noindent{\footnotesize Izvorne zapisy / Original records: HIST-g147-002.}\par
\smallskip\noindent\emph{Sravnjenje s tekstom:} Suma, produkt i relativny klas ohranjajut različne česti definicije.\par
\noindent\emph{Textual comparison:} Sum, product, and the relative class preserve different parts of the definition.\par
\noindent{\footnotesize TEXT-COMP-666; pasaže / passages: 18503, 18505, 18507, 18509, 18511, 18513, 18525; \texttt{CORPUS\_FORM\_COMPARISONS.json}.}\par
\end{minipage}\par
\par\medskip\Needspace{7\baselineskip}
\noindent\begin{minipage}{\linewidth}
\hypertarget{TERM-1611}{}
\noindent\textbf{teorem}\quad{\footnotesize\texttt{TERM-1611}}\par
\noindent\emph{Iztočnik / Source:} Satz\par
\noindent\emph{Rusky / Russian:} теорема\par
\noindent\emph{Ukrajinsky / Ukrainian:} теорема\par
\smallskip\noindent\emph{Razlog:} Teorem prěvodi Satz kako oznaku dokazovanogo tvrdženja, ne kako obyčnu větu. Dva zapisy raboty 31 davajut prednost kontinuitetu mužskogo internacionalnogo naziva i tipografskomu razděljenju oznaky od tvrdženja. Čislo, tučny tisk i kursiv pomagajut najti granicu teorema; oni ne dokazujut jeho matematičnu pravilnost. Ta mužska forma ne jest edinstvena historijna forma korpusa: v rabotě 6 stojit teorema. V novo prověrenoj kopiji oficialnogo slovnika ni teorem ni teorema ne sut samostojne zaglavja. Taka odsutnost ne čini jih nepravilnymi, ale ne dava nezavisny slovnikovy dokaz za izbor roda. Rod i soglasovanje trěba obrazložiti posebno prěd občeju standardizacijeju.\par
\noindent\emph{Reason:} Teorem renders Satz as the label of a mathematical assertion, not an ordinary sentence. Two Paper 31 records favour continuity of the masculine international form and typographic separation of label from statement. Numbering, bold type, and italics help readers locate the theorem; they do not establish its mathematical correctness. The masculine form is not the corpus's only historical choice: Paper 6 has teorema. Neither teorem nor teorema occurs as an independent headword in the newly checked official-dictionary snapshot. Absence does not make them invalid, but it supplies no independent dictionary authority for the gender choice. Gender and agreement need separate justification before corpus-wide standardization.\par
\noindent{\footnotesize Izvorne zapisy / Original records: HIST-g126-001, HIST-g144-001.}\par
\smallskip\noindent\emph{Sravnjenje s tekstom:} Tipografija ohranjaje granicu tvrdženja, ne dokazuje njegovu pravilnost.\par
\noindent\emph{Textual comparison:} Typography preserves the assertion boundary, not its correctness.\par
\noindent{\footnotesize TEXT-COMP-668; pasaže / passages: 17661, 18412; \texttt{CORPUS\_FORM\_COMPARISONS.json}.}\par
\end{minipage}\par
\par\medskip\Needspace{7\baselineskip}
\noindent\begin{minipage}{\linewidth}
\hypertarget{TERM-1612}{}
\noindent\textbf{Teorema I}\quad{\footnotesize\texttt{TERM-1612}}\par
\noindent\emph{Iztočnik / Source:} Satz I\par
\noindent\emph{Rusky / Russian:} Теорема I\par
\noindent\emph{Ukrajinsky / Ukrainian:} Теорема I\par
\smallskip\noindent\emph{Razlog:} Teorema I označaje formalno numerovano tvrdženje raboty 6 o vzajemno jednoznačnom priredženju sistemov funkcij s ohranjenjem relacij v oběh směrah. Arhivny razlog jest kratky: Satz jest tu teorem, ne proza bez osobnoj matematičnoj roli. Rimsky I identifikuje tvrdženje; ne jest promenna ni bukva za transliteraciju. Forma teorema i genitiv teoremy v nynějšnoj pasaži stojut vedle teorem v drugyh rabotah. Slovnik tu razliku ne skryvaje i ne pripisyvaje arhivu nepostojajuče čislene sravnjenje roda. Pri jedinoj budučej normě trěba spolu kontrolovati pádne formy i vse odkazy, ne samo zaglavje.\par
\noindent\emph{Reason:} Teorema I labels the formal numbered assertion in Paper 6 about a one-to-one correspondence between function systems preserving relations in both directions. The archived reason is brief: Satz denotes a theorem here, not prose without a distinct mathematical role. Roman I identifies the assertion; it is neither a variable nor a letter for transliteration. Teorema and the genitive teoremy in this passage coexist with teorem elsewhere. The dictionary neither conceals that difference nor attributes an unrecorded numerical gender comparison to the archive. A future common convention must check inflection and references together, not just replace the heading.\par
\noindent{\footnotesize Izvorne zapisy / Original records: HIST-g043-009.}\par
\smallskip\noindent\emph{Sravnjenje s tekstom:} Forma teorema jest zasvědočena s odkazom v genitivu; tvrdženje ohranjaje obě smeri.\par
\noindent\emph{Textual comparison:} Teorema is attested with a genitive reference; the assertion retains both directions.\par
\noindent{\footnotesize TEXT-COMP-669; pasaže / passages: 5372, 5374, 5376, 5378; \texttt{CORPUS\_FORM\_COMPARISONS.json}.}\par
\end{minipage}\par
\par\medskip\Needspace{7\baselineskip}
\noindent\begin{minipage}{\linewidth}
\hypertarget{TERM-1614}{}
\noindent\textbf{teorem o konečnoj cěpi}\quad{\footnotesize\texttt{TERM-1614}}\par
\noindent\emph{Iztočnik / Source:} Satz von der endlichen Kette\par
\noindent\emph{Rusky / Russian:} теорема о конечной цепи\par
\noindent\emph{Ukrajinsky / Ukrainian:} теорема про скінченний ланцюг\par
\smallskip\noindent\emph{Razlog:} Teorem o konečnoj cěpi ohranjaje lokalny opis i historijno nazvanje, ktore tekst pripisyvaje Dedekindu. Ruska cěpna metafora i ukrajinsky lancjug dajut različne korenje za isti porjadkovy obraz; arhiv ne dava čislene sravnjenje teh korenjev. Sam teorem izgovarjaje stabilizaciju ideala od někojego konečnogo indeksa, odnosno nemožnost beskonečnoj strogoj cěpi danogo směra. Ne tvrdži, že cěly prsten ima konečno mnogo idealov ili elementov. Izbor slova konečny mora zato byti svęzany s cěpju i uslovjem konečnosti bazy. Dedekindov i Laskerov prispevek v izvornoj primětce ostavajut avtorsko iztočnikovo dědictvo.\par
\noindent\emph{Reason:} Teorem o konečnoj cěpi preserves the local description and the historical chain terminology attributed in the text to Dedekind. Russian and Ukrainian use different chain roots for the same ordering image; the archive provides no numerical comparison of those roots. The theorem asserts stabilization from some finite index, equivalently the impossibility of an infinite strict chain in the specified direction. It does not say that the ring has finitely many ideals or elements. Konečny must therefore remain attached to the chain and the finite-basis assumption. Dedekind's and Lasker's contributions in the source footnote remain inherited scholarly provenance.\par
\noindent{\footnotesize Izvorne zapisy / Original records: HIST-g081-027, HIST-g082-027.}\par
\smallskip\noindent\emph{Sravnjenje s tekstom:} Konečny indeks stabilizacije ne znači konečno mnogo idealov v cělom prstenu.\par
\noindent\emph{Textual comparison:} A finite stabilization index does not mean finitely many ideals in the whole ring.\par
\noindent{\footnotesize TEXT-COMP-671; pasaže / passages: 12438, 12515, 12516, 12518, 12520, 12526; \texttt{CORPUS\_FORM\_COMPARISONS.json}.}\par
\end{minipage}\par
\par\medskip\Needspace{7\baselineskip}
\noindent\begin{minipage}{\linewidth}
\hypertarget{TERM-1615}{}
\noindent\textbf{teorem o izomorfiji}\quad{\footnotesize\texttt{TERM-1615}}\par
\noindent\emph{Iztočnik / Source:} Satz von der Isomorphie\par
\noindent\emph{Rusky / Russian:} теорема об изоморфии\par
\noindent\emph{Ukrajinsky / Ukrainian:} теорема про ізоморфію\par
\smallskip\noindent\emph{Razlog:} Teorem o izomorfiji imenuje rezultat v §8 raboty 17 opisno, bez zaměny novym slavnym imenom. Arhiv izrično dava prednost toj věrnosti iztočnomu naslovu. Tekst sravnjava dva konečna razklady polno reducibilnoj grupy na proste grupy: každa komponenta ima izomorfnu odgovarjajuču komponentu, a osobne uslovja odlučajut identičnost razkladov. To ne jest toliko formula o kvocientu i jadru homomorfizma, ktora v drugom kontekstu nosi podobno nazvanje. Pri ponovnoj upotrěbě treba identifikovati tvrdženje i prědpoloženja, ne izenačiti vse teoremy o izomorfiji po frazě.\par
\noindent\emph{Reason:} Teorem o izomorfiji names the result in Paper 17, §8 descriptively rather than substituting a later famous name. The record expressly prefers this fidelity to the source heading. The text compares two finite decompositions of a completely reducible group into simple groups: each component has an isomorphic counterpart, and particular conditions govern identity of the decompositions. This is not merely the quotient-and-kernel formula called an isomorphism theorem in another setting. Reuse requires identifying the assertion and its hypotheses, not equating all results sharing the phrase.\par
\noindent{\footnotesize Izvorne zapisy / Original records: HIST-g079-047.}\par
\smallskip\noindent\emph{Sravnjenje s tekstom:} Opisny naslov jest svęzany s konkretnym tvrdženjem o komponentah.\par
\noindent\emph{Textual comparison:} The descriptive heading is tied to a specific assertion about components.\par
\noindent{\footnotesize TEXT-COMP-672; pasaže / passages: 11568, 11570, 11572, 11576, 11612, 11614, 11616; \texttt{CORPUS\_FORM\_COMPARISONS.json}.}\par
\end{minipage}\par
\par\medskip\Needspace{7\baselineskip}
\noindent\begin{minipage}{\linewidth}
\hypertarget{TERM-1621}{}
\noindent\textbf{ostrej?e relacije (27), (28)}\quad{\footnotesize\texttt{TERM-1621}}\par
\noindent\emph{Iztočnik / Source:} schaerfere Relationen (27), (28)\par
\noindent\emph{Rusky / Russian:} более сильные соотношения (27), (28)\par
\noindent\emph{Ukrajinsky / Ukrainian:} сильніші відношення (27), (28)\par
\smallskip\noindent\emph{Razlog:} Ostrějše relacije označajut silnějši rezultat dokaza: ne samo izomorfiju grup ostatkov, ale konkretne kongruencije (27), (28) s $X=Q$ i $Y=P$. Arhivny razlog izrično čuva tu razliku; rusko bolee sil'nye i ukrajinsko syl'niši davajut sravniteljne konstrukcije s tym samym logičnym smyslom. Nynějšny tekst raboty 17 potvrđuje i komparativ, i oba svědka. Ostrějši tu ne znači bolje točny čislovy odhad. Zaglavje starogo zapisa ima nečitljivu bukvu; nynějšnje ostrějše jest oddělno potvrđena forma, ne obnovjena bukva arhiva. Izbor jest kvalitativny: čitljiva sveza s siloj tvrdženja, ne izměreny rezultat sravnjenja jezykov.\par
\noindent\emph{Reason:} Ostrějše relacije marks a stronger result: not merely an isomorphism of residue groups, but the particular congruences (27), (28) with $X=Q$ and $Y=P$. The recorded reason expressly preserves this distinction; Russian bolee sil'nye and Ukrainian syl'niši supply comparative constructions with the same logical force. The current text of paper 17 confirms both the comparative and the two witnesses. Sharper here does not mean a more accurate numerical estimate. One letter in the old heading is unreadable; current ostrějše is separately attested, not a recovered archival letter. The choice is qualitative, favouring a clear relation to strength of assertion, not based on measured language-comparison scores.\par
\noindent{\footnotesize Izvorne zapisy / Original records: HIST-g079-053.}\par
\smallskip\noindent\emph{Sravnjenje s tekstom:} Kongruencije s $X=Q$ i $Y=P$ davajut vyše informaciji než samo izomorfija.\par
\noindent\emph{Textual comparison:} Congruences with $X=Q$ and $Y=P$ give more information than the isomorphism alone.\par
\noindent{\footnotesize TEXT-COMP-678; pasaže / passages: 11664, 11681, 11695; \texttt{CORPUS\_FORM\_COMPARISONS.json}.}\par
\end{minipage}\par
\par\medskip\Needspace{7\baselineskip}
\noindent\begin{minipage}{\linewidth}
\hypertarget{TERM-1622}{}
\noindent\textbf{shema svijanja}\quad{\footnotesize\texttt{TERM-1622}}\par
\noindent\emph{Iztočnik / Source:} Schema zur Faltung\par
\noindent\emph{Rusky / Russian:} схема свертки\par
\noindent\emph{Ukrajinsky / Ukrainian:} схема згортання\par
\smallskip\noindent\emph{Razlog:} Shema svijanja byla kratka povratna upominka o konstrukciji iz §7. Dva zapisy iz raboty 2 čuvajut toj razlog, ale ostavljajut shema/schema za pravopisno razsmotrenje. Novo čitanje slovnika podava shema, ženskogo roda; to podpiraje pisanje imenika, ne samo po sebě termin svijanje. Nynějšny §15 imaje shema kontrakcije, a §16 takože označaje operacije slovom kontrakcija. Zato trěba razděliti dva pytanja: kako pisati shema i kako nazvati tu konkretnu operaciju teorije invariantov. Jej zaměna na obče slovo za skladanje ili svivanje bez prověrky formul mogla by izgubiti tehničny smysl.\par
\noindent\emph{Reason:} Shema svijanja was a short backward reference to the construction in §7. Two records from paper 2 preserve that reason while leaving shema/schema for spelling review. Fresh dictionary consultation gives feminine shema: evidence for the noun's spelling, not in itself for the technical term svijanje. Current §15 has shema kontrakcije, and §16 also labels the operations kontrakcija. Two questions must therefore be separated: how to spell shema, and how to name this particular operation of invariant theory. Replacing it with a general word for folding without checking the formulas could lose the technical sense.\par
\noindent{\footnotesize Izvorne zapisy / Original records: HIST-g016-011, HIST-g017-010.}\par
\smallskip\noindent\emph{Sravnjenje s tekstom:} Nynějšny tekst čuva shema, ale nazva operacije jest kontrakcija.\par
\noindent\emph{Textual comparison:} The current text retains shema, but names the operation kontrakcija.\par
\noindent{\footnotesize TEXT-COMP-679; pasaže / passages: 2271, 2297, 2298, 2299, 2300, 2301, 2302; \texttt{CORPUS\_FORM\_COMPARISONS.json}.}\par
\end{minipage}\par
\par\medskip\Needspace{7\baselineskip}
\noindent\begin{minipage}{\linewidth}
\hypertarget{TERM-1623}{}
\noindent\textbf{zaključny izraz}\quad{\footnotesize\texttt{TERM-1623}}\par
\noindent\emph{Iztočnik / Source:} Schlussausdruck\par
\noindent\emph{Rusky / Russian:} заключительное выражение\par
\noindent\emph{Ukrajinsky / Ukrainian:} підсумковий вираз\par
\smallskip\noindent\emph{Razlog:} Zaključny izraz imenuje izraz dostany na koncu opisanogo postupka, ne novy matematičny objekt s osobnoju definicijeju. Zapisany razlog daje prednost tomu funkcionalnomu objasnjenju prěd doslovnym obrazom zatvarjanja. V §4 raboty 4 simetrična forma pokazuje, že dva razkladanja vedut k tomu samomu izrazu; nezavisnost od prvotnogo razkladanja jest sučstveny rezultat. Rusko zaključitel'noe i ukrajinsko pidsumkovyj podpirajut čitanje zaključny kako završny. Ni zaglavje, ni ta motivacija ne utvrđajut, že izraz jest zatvorjena formula v posebnom modernom smyslu. Takova dodatna vlastnost ne směje byti vložena prěz izbor slova.\par
\noindent\emph{Reason:} Zaključny izraz names the expression reached at the end of the procedure, not a new mathematical object with a separate definition. The recorded reason favours this functional description over a literal image of closing. In paper 4, §4 the symmetric form shows that two decompositions lead to the same expression; independence of the original decomposition is the substantive result. Russian zaključitel'noe and Ukrainian pidsumkovyj support the concluding or final sense. Neither the heading nor its rationale establishes that this is a closed-form expression in a particular modern technical sense. The wording must not introduce that extra property.\par
\noindent{\footnotesize Izvorne zapisy / Original records: HIST-g033-002.}\par
\smallskip\noindent\emph{Sravnjenje s tekstom:} Simetrična forma zaključnogo izraza čuva smysl rezultata.\par
\noindent\emph{Textual comparison:} The symmetric final expression preserves the substantive conclusion.\par
\noindent{\footnotesize TEXT-COMP-680; pasaže / passages: 4278; \texttt{CORPUS\_FORM\_COMPARISONS.json}.}\par
\end{minipage}\par
\par\medskip\Needspace{7\baselineskip}
\noindent\begin{minipage}{\linewidth}
\hypertarget{TERM-1624}{}
\noindent\textbf{konečne izrazy / početne izrazy}\quad{\footnotesize\texttt{TERM-1624}}\par
\noindent\emph{Iztočnik / Source:} Schlußausdrücke / Anfangsausdrücke\par
\noindent\emph{Rusky / Russian:} конечные выражения / исходные выражения\par
\noindent\emph{Ukrajinsky / Ukrainian:} остаточні вирази / початкові вирази\par
\smallskip\noindent\emph{Razlog:} Konečne i početne izrazy sut dva konca redukcije. Arhiv posebno zabeležuje idiomatično rusko ishodnye za početne izrazy; ukrajinsko počatkovi čuva to samo značenje. Nynějšny §8 raboty 4 imaje važno ograničenje: bez relacij (62) ne slěduje, že konečne izrazy sut ekvivalentne početnym. Oprěděljenje ekvivalentnosti uključaje razliku formami vyššej kontrakcije, odnosno s vyššim umenšenjem vagi. Zato konečne tu znači završne, ne izrazy s konečnym čislom členov. Jasna para početne/konečne čuva poredok postupka, ale ne zaměnjaje uslovje, pod ktorym rezultat ostava ekvivalentny izhodu.\par
\noindent\emph{Reason:} Konečne and početne izrazy are the two ends of a reduction. The archive specifically notes idiomatic Russian ishodnye for the starting expressions; Ukrainian počatkovi carries the same sense. Current paper 4, §8 contains a crucial restriction: without relations (62), the final expressions need not be equivalent to the initial ones. Equivalence involves a difference consisting of forms of higher contraction, hence greater reduction in weight. Konečne therefore means final here, not expressions with finitely many terms. The initial/final pair preserves the procedure's order but does not replace the condition under which its result remains equivalent to its starting point.\par
\noindent{\footnotesize Izvorne zapisy / Original records: HIST-g037-010.}\par
\smallskip\noindent\emph{Sravnjenje s tekstom:} Ekvivalentnost ostava ograničena navedenymi relacijami.\par
\noindent\emph{Textual comparison:} Equivalence remains conditional on the specified relations.\par
\noindent{\footnotesize TEXT-COMP-681; pasaže / passages: 4872; \texttt{CORPUS\_FORM\_COMPARISONS.json}.}\par
\end{minipage}\par
\par\medskip\Needspace{7\baselineskip}
\noindent\begin{minipage}{\linewidth}
\hypertarget{TERM-1625}{}
\noindent\textbf{zaključenje}\quad{\footnotesize\texttt{TERM-1625}}\par
\noindent\emph{Iztočnik / Source:} Schlußfolgerung\par
\noindent\emph{Rusky / Russian:} заключение\par
\noindent\emph{Ukrajinsky / Ukrainian:} висновок\par
\smallskip\noindent\emph{Razlog:} Zaključenje uvodi sobranje rezultatov dotedašnjih izčisljenj. Zapisany razlog dava prednost obyčnomu matematičnomu prose i nazyvaje formu medžuslovjanski prozračnoju; to jest redakcijna ocěna, ne čitatelsky eksperiment. V §26 raboty 2 za tym slovom slěduje relativno polny sistem iz 331 form i podsistemov modulo zadanyh objektov. Čislo 331 pripada avtorskym konstrukcijam, ne slovniku ili izměrenju kačestva prěvoda. Rusko zaključenie i ukrajinsko vysnovok pokazujut različne leksične sredstva za tu samu izložiteljnu funkciju. Izbor ne trěbuje, aby vse tri jezyky iměli jednaky korenj.\par
\noindent\emph{Reason:} Zaključenje introduces a synthesis of the preceding calculations. The recorded reason favours ordinary mathematical prose and calls the form pan-Slavically transparent; that is an editorial judgement, not a reader experiment. In paper 2, §26 it introduces a relatively complete system of 331 forms and subsystems modulo the specified objects. The number 331 belongs to the author's constructions, not to the dictionary or a measure of translation quality. Russian zaključenie and Ukrainian vysnovok illustrate different lexical resources for the same expository function. The choice does not require all three languages to use the same root.\par
\noindent{\footnotesize Izvorne zapisy / Original records: HIST-g027-009.}\par
\smallskip\noindent\emph{Sravnjenje s tekstom:} Naslov uvodi rezultat o 331 formah i podsistemah.\par
\noindent\emph{Textual comparison:} The wording introduces a result about 331 forms and subsystems.\par
\noindent{\footnotesize TEXT-COMP-682; pasaže / passages: 3503, 3505; \texttt{CORPUS\_FORM\_COMPARISONS.json}.}\par
\end{minipage}\par
\par\medskip\Needspace{7\baselineskip}
\noindent\begin{minipage}{\linewidth}
\hypertarget{TERM-1626}{}
\noindent\textbf{konečna formula}\quad{\footnotesize\texttt{TERM-1626}}\par
\noindent\emph{Iztočnik / Source:} Schlußformel\par
\noindent\emph{Rusky / Russian:} конечная формула\par
\noindent\emph{Ukrajinsky / Ukrainian:} кінцева формула\par
\smallskip\noindent\emph{Razlog:} Konečna formula jest formula dostana na koncu računa i potom neposrědnje upotrěbljana. Arhivny razlog jest kratky: obyčno moderno matematično vyraženje bylo bolje vhodno než historijna kalka. V §6 raboty 2 avtorka najprvo podrobno izvodi priměr iz razvitja II, a za daljše računy upominaje formulu III. Čislovy identifikator trěba čuvati, zato že svęzuje objasnjenje s konkretnym rezultatom. Konečna ne znači, že formula ima konečnu oblast priměnenja, ni že vsaka slična formula jest zatvorjena forma. Nynějšna upotrěba potvrđuje smysl završnogo rezultata, ne širu terminologičnu normu.\par
\noindent\emph{Reason:} Konečna formula is the formula obtained at the end of a calculation and then used directly. The archival reason is brief: ordinary modern mathematical wording was preferable to a historical calque. In paper 2, §6 the author first works through an example using expansion II, then refers to formula III for later calculations. That numerical identifier must remain because it links the explanation to a particular result. Konečna does not mean that the formula has a finite domain of application, or that every such formula is a closed form. Current usage confirms the final-result sense, not a broader terminological standard.\par
\noindent{\footnotesize Izvorne zapisy / Original records: HIST-g007-006.}\par
\smallskip\noindent\emph{Sravnjenje s tekstom:} Oznaka III svęzuje prose s konkretnym rezultatom računa.\par
\noindent\emph{Textual comparison:} Identifier III connects the prose to a particular result of the calculation.\par
\noindent{\footnotesize TEXT-COMP-683; pasaže / passages: 922; \texttt{CORPUS\_FORM\_COMPARISONS.json}.}\par
\end{minipage}\par
\par\medskip\Needspace{7\baselineskip}
\noindent\begin{minipage}{\linewidth}
\hypertarget{TERM-1627}{}
\noindent\textbf{oblast prěsěka / zatvorjena vzhodno k tvorenju prěsěkov}\quad{\footnotesize\texttt{TERM-1627}}\par
\noindent\emph{Iztočnik / Source:} Schnittbereich / Durchschnitte abgeschlossen\par
\noindent\emph{Rusky / Russian:} область пересечения / замкнута относительно образования пересечений\par
\noindent\emph{Ukrajinsky / Ukrainian:} область перетину / замкнена щодо творення перетинів\par
\smallskip\noindent\emph{Razlog:} Oblast prěsěka i zatvorjenost po odnošenju k prěsěkam trěba čitati vměstě. V rabotě 9 prěsěk oblastij kojegokoli podsistema iz $\mathfrak M(H)$ ponovno naleži k $\mathfrak M(H)$. To jest zatvorjenost klasa pod operacijeju, ne topologična zatvorjenost ili cělost elementov. Starši zapis čuva vzhodno k kako tadašnju domovu konvenciju i izrično ostavljaje ju za prověrku. Nynějšny tekst ima po odnošenju k, čemu jest posvěčeno osobno objasnjenje. Ruskemu peresečenie i ukrajinskomu peretyn odgovarja prěsěk; zajedničny smysl jest ostavanje v tom samom klasu poslě operacije, ne podobnost korenjev sama po sebě.\par
\noindent\emph{Reason:} Oblast prěsěka and closure with respect to intersections must be read together. In paper 9, the intersection of the regions in any subsystem of $\mathfrak M(H)$ again belongs to $\mathfrak M(H)$. This is closure of a class under an operation, not topological closedness or integrality of elements. The earlier record preserves vzhodno k as the house convention of its time and expressly leaves it for review. The current text has po odnošenju k, discussed separately. Russian peresečenie and Ukrainian peretyn correspond to prěsěk; the shared meaning is remaining in the same class after the operation, not root resemblance alone.\par
\noindent{\footnotesize Izvorne zapisy / Original records: HIST-g062-012.}\par
\smallskip\noindent\emph{Sravnjenje s tekstom:} Prěsěk podsistema ponovno naleži k razgledanoj sovkupnosti oblastij.\par
\noindent\emph{Textual comparison:} The intersection of a subsystem again belongs to the class of regions under discussion.\par
\noindent{\footnotesize TEXT-COMP-684; pasaže / passages: 7445, 7446, 7447, 7448, 7449, 7450, 7451, 7452, 7453, 7454, 7455, 7456, 7562; \texttt{CORPUS\_FORM\_COMPARISONS.json}.}\par
\end{minipage}\par
\par\medskip\Needspace{7\baselineskip}
\noindent\begin{minipage}{\linewidth}
\hypertarget{TERM-1628}{}
\noindent\textbf{mnogokratnost prěseka / rezultanta-mnogokratnost}\quad{\footnotesize\texttt{TERM-1628}}\par
\noindent\emph{Iztočnik / Source:} Schnittpunktmultiplizitaet / Resultantenmultiplizitaet\par
\noindent\emph{Rusky / Russian:} кратность пересечения / результантная кратность\par
\noindent\emph{Ukrajinsky / Ukrainian:} кратність перетину / результантна кратність\par
\smallskip\noindent\emph{Razlog:} Mnogokratnost prěseka imenuje geometričnu veličinu, a složeno rezultanta-mnogokratnost izrično pokazuje jej svezu s rezultantoju. Arhiv dava prednost takoj prozračnoj konstrukciji prěd kratkim nejasnym novym imenikom. V rabotě 45 $\mu$ jest mnogokratnost prěseka oprěděljena rezultantoju; suma čisel $(g_i)$ jest ravna toj samoj $\mu$, a dokaz upotrěbjaje produktny teorem rezultantov. To ne sut dvě nezavisne mnogokratnosti proizvoljno pridružene točce. Različne nazvy osvětljajut geometričnu interpretaciju i sposob izčisljenja. Smysl složenogo termina možno zato objasniti bez tvrdženja, že forma s defisom jest už obče prijeta v medžuslovjanskoj matematičnoj literaturě.\par
\noindent\emph{Reason:} Mnogokratnost prěseka names the geometric quantity, while the compound rezultanta-mnogokratnost explicitly signals its relation to a resultant. The archive favours that transparent construction over a short but opaque new noun. In paper 45, $\mu$ is the intersection multiplicity defined by the resultant; the numbers $(g_i)$ sum to that same $\mu$, and the proof uses the product theorem for resultants. These are not two independent multiplicities arbitrarily attached to a point. The names emphasize the geometric interpretation and the method of calculation. The compound can therefore be explained without claiming that its hyphenated form is already established in Interslavic mathematical literature.\par
\noindent{\footnotesize Izvorne zapisy / Original records: HIST-g214-007.}\par
\smallskip\noindent\emph{Sravnjenje s tekstom:} Geometrična mnogokratnost $\mu$ jest oprěděljena rezultantoju.\par
\noindent\emph{Textual comparison:} The geometric multiplicity $\mu$ is defined by the resultant.\par
\noindent{\footnotesize TEXT-COMP-685; pasaže / passages: 45-isv.tex:164, 45-isv.tex:165, 45-isv.tex:166, 45-isv.tex:167, 45-isv.tex:168, 45-isv.tex:169, 45-isv.tex:206, 45-isv.tex:207, 45-isv.tex:208, 45-isv.tex:209, 45-isv.tex:210, 45-isv.tex:211, 45-isv.tex:213, 45-isv.tex:214, 45-isv.tex:215, 45-isv.tex:216, 45-isv.tex:217, 45-isv.tex:218, 45-isv.tex:219, 45-isv.tex:220, 45-isv.tex:221, 45-isv.tex:222, 45-isv.tex:223, 45-isv.tex:224, 45-isv.tex:225, 45-isv.tex:226; \texttt{CORPUS\_FORM\_COMPARISONS.json}.}\par
\end{minipage}\par
\par\medskip\Needspace{7\baselineskip}
\noindent\begin{minipage}{\linewidth}
\hypertarget{TERM-1629}{}
\noindent\textbf{Šurov indeks}\quad{\footnotesize\texttt{TERM-1629}}\par
\noindent\emph{Iztočnik / Source:} Schurscher Index\par
\noindent\emph{Rusky / Russian:} индекс Шура\par
\noindent\emph{Ukrajinsky / Ukrainian:} індекс Шура\par
\smallskip\noindent\emph{Razlog:} Šurov indeks čuva ime Schura i označaje konkretno čislo $m$. Starši razlog sravnjava kratku pridavničnu formu Šurov s analitičnoju konstrukcijeju indeks Schura; rusky i ukrajinsky zapisy imajut drugy tip. V rabotě 32 pridruženo nekomutativno tělo ima dimenziju $m^2$ nad centrom, ale indeks jest $m$, ne $m^2$. To samo $m$ jest stepen največših komutativnyh podpolj v razgledanom slučaju i čislo prostyh jednostrannyh idealov poslě razširjenja do razpadnogo polja. Takove konteksty utvrđajut identičnost veličiny. Pridavnična adaptacija imena ne odstranjava bibliografičnu pripisanost Schurovi i ne dava novogo avtora pojma.\par
\noindent\emph{Reason:} Šurov indeks preserves Schur's name and denotes the particular integer $m$. The earlier rationale compares compact adjectival Šurov with an analytic index-of-Schur construction; the Russian and Ukrainian records use the latter pattern. In paper 32 the associated division algebra has dimension $m^2$ over its centre, but the index is $m$, not $m^2$. In the setting discussed, the same $m$ is the degree of maximal commutative subfields and the number of simple one-sided ideals after extension to a splitting field. These contexts establish the quantity's identity. Adjectival adaptation does not remove Schur's scholarly credit or assign a new originator to the concept.\par
\noindent{\footnotesize Izvorne zapisy / Original records: HIST-g169-004, HIST-g171-004, HIST-g173-004.}\par
\smallskip\noindent\emph{Sravnjenje s tekstom:} Indeks jest $m$, medžu tym kako dimenzija těla nad centrom jest $m^2$.\par
\noindent\emph{Textual comparison:} The index is $m$, whereas the division algebra's dimension over its centre is $m^2$.\par
\noindent{\footnotesize TEXT-COMP-686; pasaže / passages: 19314, 19318, 19453; \texttt{CORPUS\_FORM\_COMPARISONS.json}.}\par
\end{minipage}\par
\par\medskip\Needspace{7\baselineskip}
\noindent\begin{minipage}{\linewidth}
\hypertarget{TERM-1630}{}
\noindent\textbf{slaby / silny prosty ideal}\quad{\footnotesize\texttt{TERM-1630}}\par
\noindent\emph{Iztočnik / Source:} schwaches / starkes Primideal\par
\noindent\emph{Rusky / Russian:} слабый / сильный простой идеал\par
\noindent\emph{Ukrajinsky / Ukrainian:} слабкий / сильний простий ідеал\par
\smallskip\noindent\emph{Razlog:} Slaby i silny prosty ideal sut začasne definicijne oznaky, ne ocěny kačestva idealov. V §5 raboty 30 jedna definicija govori o elementah-děliteljah nuly v prstenu ostatkov, druga o idealah-děliteljah nuly. Prsten jest komutativno, ale jediničny element ne jest prědpoloženy. Dokaz potom ustanovjaje ekvivalentnost oběh definicij i dozvaljaje dalje govoriti prosto o prostyh idealah. Imenno toj dokazny krok jest zapisany razlog za čuvanje oběh pridavnikov. Nynějšny prěpis ima posebno zabeleženo pytanje o tipah simbolov v definiciji; razjasnjenje značenja slov ne rěšava togo pytanja i ne jest potvrđjenjem věrnosti vsakej formuly.\par
\noindent\emph{Reason:} Slaby and silny prosty ideal are temporary definitional labels, not assessments of the quality of ideals. In paper 30, §5 one definition concerns element zero divisors in the residue ring, the other ideal zero divisors. The ring is commutative, but an identity is not assumed. The proof then establishes equivalence and permits the subsequent unqualified term prime ideal. Preserving this proof step is the recorded reason for retaining both adjectives. The current transcription has a separately recorded question about the types of symbols used in the definition; explaining the words does not resolve that question or certify every formula.\par
\noindent{\footnotesize Izvorne zapisy / Original records: HIST-g112-005.}\par
\smallskip\noindent\emph{Sravnjenje s tekstom:} Ekvivalentnost oběh definicij jest rezultat dokaza, ne pravopisna konvencija.\par
\noindent\emph{Textual comparison:} Equivalence of the definitions is a proved result, not a spelling convention.\par
\noindent{\footnotesize TEXT-COMP-687; pasaže / passages: 16601, 16603, 16605; \texttt{CORPUS\_FORM\_COMPARISONS.json}.}\par
\end{minipage}\par
\par\medskip\Needspace{7\baselineskip}
\noindent\begin{minipage}{\linewidth}
\hypertarget{TERM-1631}{}
\noindent\textbf{slaby / silny primarny ideal}\quad{\footnotesize\texttt{TERM-1631}}\par
\noindent\emph{Iztočnik / Source:} schwaches / starkes Primärideal\par
\noindent\emph{Rusky / Russian:} слабый / сильный первичный идеал\par
\noindent\emph{Ukrajinsky / Ukrainian:} слабкий / сильний первинний ідеал\par
\smallskip\noindent\emph{Razlog:} Slaby i silny primarny ideal ne možno sliti samo zato, že susědny dokaz sliva dva pojma prostogo ideala. V tom samom §5 slaba definicija trěbuje nilpotentnost elementa-dělitelja nuly, silna nilpotentnost ideala-dělitelja nuly. Bez aksiomy I obrat od slabogo k silnomu ne imaje město; priměr R. Hölzera s beskonečno mnogymi prěměnnymi pokazuje razliku. Pod uslovjem cěpov konečna baza dava obči eksponent, čto opravdava sjedinenje definicij. Arhivny razlog izrično čuva tu zavisimost od aksiomy. Zato prěvod ne směje modernizovati razlikovanje do jednogo neograničenogo pojma. Zabeleženo pytanje o eksponentah v odlomku ostava oddělno od toj terminologičnoj motivacije.\par
\noindent\emph{Reason:} Slaby and silny primarny ideal cannot be merged merely because the neighbouring proof merges the two notions of prime ideal. In the same §5, the weak definition requires nilpotence of an element zero divisor, the strong definition nilpotence of an ideal zero divisor. Without Axiom I, the weak-to-strong implication fails; R. Hölzer's example in infinitely many variables shows the difference. Under the chain condition, a finite basis supplies a common exponent and justifies identifying the definitions. The recorded reason expressly preserves this dependence on the axiom. Translation must not modernize the distinction into one unrestricted concept. The recorded question about exponents in the passage remains separate from this terminological rationale.\par
\noindent{\footnotesize Izvorne zapisy / Original records: HIST-g112-006.}\par
\smallskip\noindent\emph{Sravnjenje s tekstom:} Bez uslovja cěpov razlika medžu elementarnymi i idealnymi eksponentami ostava sučstvena.\par
\noindent\emph{Textual comparison:} Without the chain condition, the distinction between elementwise and ideal nilpotence remains substantive.\par
\noindent{\footnotesize TEXT-COMP-688; pasaže / passages: 16607, 16609, 16611, 16613, 16615; \texttt{CORPUS\_FORM\_COMPARISONS.json}.}\par
\end{minipage}\par
\par\medskip\Needspace{7\baselineskip}
\noindent\begin{minipage}{\linewidth}
\hypertarget{TERM-1632}{}
\noindent\textbf{simultanna sistema}\quad{\footnotesize\texttt{TERM-1632}}\par
\noindent\emph{Iztočnik / Source:} simultanes System\par
\noindent\emph{Rusky / Russian:} симультанная система\par
\noindent\emph{Ukrajinsky / Ukrainian:} симультанна система\par
\smallskip\noindent\emph{Razlog:} Simultanna sistema znači sistemu objektov, ktore razgledajut se vměstě pod istoju proměnoju prěměnnyh. Ne ide o dva sobytja v jednom času. Arhivny razlog daje prednost internacionalnomu pridavniku za teoriju invariantov; v tom zapisu ruska i ukrajinska forma takože sut simultanne. V rabotě 12 ide o diferencialne izrazy i jih projektivne invarianty. Polna sistema može byti sčetno beskonečna, pri konečnom čislu form potrěbnyh za vsaky fiksovany red izvodnyh. Novo konsultovany slovnik potvrđuje žensky imenik sistema, ne attestaciju točnoj frazy simultanna sistema. To rodovo svědectvo treba sravniti s mužskym sistem v drugyh častjah korpusa.\par
\noindent\emph{Reason:} Simultanna sistema means a system of objects considered together under the same change of variables, not two events occurring at one time. The recorded rationale favours the international adjective for invariant theory; this record's Russian and Ukrainian forms also use it. Paper 12 concerns differential expressions and their projective invariants. The full system may be countably infinite, with only finitely many forms required for each fixed derivative order. Fresh dictionary consultation supports the feminine noun sistema, not an attestation of the exact phrase simultanna sistema. This gender evidence must be considered alongside masculine sistem elsewhere in the corpus.\par
\noindent{\footnotesize Izvorne zapisy / Original records: HIST-g074-004.}\par
\smallskip\noindent\emph{Sravnjenje s tekstom:} Objekty razgledajut se pod zajedničnoju transformacijeju.\par
\noindent\emph{Textual comparison:} The objects are considered under a common transformation.\par
\noindent{\footnotesize TEXT-COMP-689; pasaže / passages: 8974, 8976; \texttt{CORPUS\_FORM\_COMPARISONS.json}.}\par
\end{minipage}\par
\par\medskip\Needspace{7\baselineskip}
\noindent\begin{minipage}{\linewidth}
\hypertarget{TERM-1633}{}
\noindent\textbf{simultanny sistem dvuh kvadratičnyh form}\quad{\footnotesize\texttt{TERM-1633}}\par
\noindent\emph{Iztočnik / Source:} simultanes System zweier quadratischer Formen\par
\noindent\emph{Rusky / Russian:} одновременная система двух квадратичных форм\par
\noindent\emph{Ukrajinsky / Ukrainian:} одночасна система двох квадратичних форм\par
\smallskip\noindent\emph{Razlog:} Simultanny sistem dvuh kvadratičnyh form jest starša forma za dvě kvadratične formy razgledane vměstě. Zapis uvodi internacionalny pridavnik kako tehnično jasnějši, medžu tym kako rusky i ukrajinsky sravniteljne zapisy upotrěbjajut slova za odnovrěmennost. Ta ocěna jasnosti nebyla podprěta čitatelskym testom. Nynějšny uvod raboty 2 ima simultanny sistem dvoh kvadratičnyh form; slovnik, konsultovany nyně, ima žensko sistema. Zato razlikovanje sistem/sistema vključaje i soglasovanje pridavnika, ne samo jednu koncovu bukvu. Smysl jest sovměstna teorija invariantov oběh form, ne nezavisno ponavljanje procedury za vsaku formu. Obě historijne formy ostavajut vidljive do opravdanogo sjednoobraznjenja.\par
\noindent\emph{Reason:} Simultanny sistem dvuh kvadratičnyh form is an earlier wording for two quadratic forms considered together. The record presents the international adjective as technically clearer, whereas its Russian and Ukrainian comparanda use words for simultaneity. That clarity judgement was not supported by reader testing. The current introduction to paper 2 has simultanny sistem dvoh kvadratičnyh form; the dictionary consulted now lists feminine sistema. The sistem/sistema variation therefore affects adjective agreement, not just a final letter. The meaning is joint invariant theory of the two forms, not an independent repetition of a procedure for each. Both historical forms remain visible pending justified standardization.\par
\noindent{\footnotesize Izvorne zapisy / Original records: HIST-g002-007.}\par
\smallskip\noindent\emph{Sravnjenje s tekstom:} Aktualno simultanny sistem trěbuje sravnjenja s slovnikovym ženskym sistema.\par
\noindent\emph{Textual comparison:} Current simultanny sistem needs comparison with dictionary feminine sistema.\par
\noindent{\footnotesize TEXT-COMP-690; pasaže / passages: 84; \texttt{CORPUS\_FORM\_COMPARISONS.json}.}\par
\end{minipage}\par
\par\medskip\Needspace{7\baselineskip}
\noindent\begin{minipage}{\linewidth}
\hypertarget{TERM-1634}{}
\noindent\textbf{simultanne invarianty}\quad{\footnotesize\texttt{TERM-1634}}\par
\noindent\emph{Iztočnik / Source:} Simultaninvarianten\par
\noindent\emph{Rusky / Russian:} симультанные инварианты\par
\noindent\emph{Ukrajinsky / Ukrainian:} симультанні інваріанти\par
\smallskip\noindent\emph{Razlog:} Simultanne invarianty sut invarianty několiko objektov pod zajedničnoju transformacijeju. Arhiv čuva internacionalny pridavnik vměsto opisnogo sovměstne, ale označaje medžuslovjanski izbor kako potrěbujuči daljšego razsmotrenja. V §7 raboty 4 formy $\Phi$ proizhodet iz $\varphi$ i $H$ invariantnymi diferencialnymi operacijami; zato nazva soderžatelno svęzana jest s oběma rodinami form. Ona ne znači, že vsaka rodina može prěobražati se nezavisno, ni že svojstva nastavajut v jednom času. Ruska i ukrajinska sravniteljne formy sut leksična podpora internacionalnomu obrascu, ne nezavisny dokaz opštej prijatosti v medžuslovjanskom.\par
\noindent\emph{Reason:} Simultanne invarianty are invariants of several objects under a common transformation. The archive retains the international adjective rather than paraphrasing it as joint, but marks the Interslavic choice for further review. In paper 4, §7 the forms $\Phi$ arise from $\varphi$ and $H$ by invariant differential operations, so the name is substantively tied to both families of forms. It does not mean that each family may be transformed independently, or that the properties arise at the same moment. The Russian and Ukrainian comparanda support the international lexical pattern, not an independent claim of general acceptance in Interslavic.\par
\noindent{\footnotesize Izvorne zapisy / Original records: HIST-g036-007.}\par
\smallskip\noindent\emph{Sravnjenje s tekstom:} Formy proizhodet iz oběh rodin invariantnymi operacijami.\par
\noindent\emph{Textual comparison:} The forms arise from both families by invariant operations.\par
\noindent{\footnotesize TEXT-COMP-691; pasaže / passages: 4600, 4609; \texttt{CORPUS\_FORM\_COMPARISONS.json}.}\par
\end{minipage}\par
\par\medskip\Needspace{7\baselineskip}
\noindent\begin{minipage}{\linewidth}
\hypertarget{TERM-1635}{}
\noindent\textbf{singularne primarne čisla / ostatkove karaktere / stepenne ostatky}\quad{\footnotesize\texttt{TERM-1635}}\par
\noindent\emph{Iztočnik / Source:} singuläre Primärzahlen / Restcharaktere / Potenzreste\par
\noindent\emph{Rusky / Russian:} особые первичные числа / вычетные характеры / степенные вычеты\par
\noindent\emph{Ukrajinsky / Ukrainian:} особливі первинні числа / лишкові характери / степеневі лишки\par
\smallskip\noindent\emph{Razlog:} Singularne primarne čisla, ostatkove karaktere i stepenne ostatky pripadajut aritmetičnoj teoriji čislovyh polj. Zapisany razlog posebno razděljaje toj smysl ostatka od geometričnogo teorema o ostatku, pri zachovanji etimologičnoj svezy. Rusko vyčet i ukrajinsko lyšok sut nacionalne sravniteljne formy, ne vyžadovanje jednogo korenja za vse jezyky. Rabota 14 govori o relativno nerazvětvenyh razširjenjah, predpisanyh karakterah i zakonah recipročnosti. Zato singularny tu ne možno tumačiti iz geometričnoj singularnosti točke. Kratky iztočny pregled ne dava polnu definiciju vsakej specialnoj aritmetičnoj kategorije; slovnik čuva to ograničenje vměsto domyslanja definicij iz podobnosti slov.\par
\noindent\emph{Reason:} Singularne primarne čisla, ostatkove karaktere and stepenne ostatky belong to the arithmetic of number fields. The recorded rationale expressly separates this residue vocabulary from the geometric residue theorem while retaining the etymological connection. Russian vyčet and Ukrainian lyšok are national comparanda, not a requirement that all languages share a root. Paper 14 discusses relatively unramified extensions, prescribed characters and reciprocity laws. Singular here must therefore not be interpreted through the geometric singularity of a point. The brief source survey does not give a full definition of every specialized arithmetic category; the dictionary retains that limitation rather than inferring definitions from similar words.\par
\noindent{\footnotesize Izvorne zapisy / Original records: HIST-g076-024.}\par
\smallskip\noindent\emph{Sravnjenje s tekstom:} Kontekst čislovyh polj ograničava smysl terminov.\par
\noindent\emph{Textual comparison:} The number-field context constrains the terms' meaning.\par
\noindent{\footnotesize TEXT-COMP-692; pasaže / passages: 10033, 10035, 10037; \texttt{CORPUS\_FORM\_COMPARISONS.json}.}\par
\end{minipage}\par
\par\medskip\Needspace{7\baselineskip}
\noindent\begin{minipage}{\linewidth}
\hypertarget{TERM-1636}{}
\noindent\textbf{osoblive točke / točka vozvrata}\quad{\footnotesize\texttt{TERM-1636}}\par
\noindent\emph{Iztočnik / Source:} singuläre Punkte / Rückkehrpunkt\par
\noindent\emph{Rusky / Russian:} особые точки / точка возврата\par
\noindent\emph{Ukrajinsky / Ukrainian:} особливі точки / точка повернення\par
\smallskip\noindent\emph{Razlog:} Osobliva točka jest obča nazva za singularnu točku krivoj; točka vozvrata čuva historijnu nazvu osobnogo tipa. Arhiv podava rusko osobaja točka i ukrajinsko osoblyva točka, a Rückkehrpunkt tumači doslovnym obrazom vozvrata. V rabotě 14 ta razlika jest sučstvena: jedna sistema koordinat može odgovarjati několiko točkam normalizovanoj krivoj ili byti prijeta v vyšem redu. Poslědnji slučaj svęzuje se s vyššoju točkoju vozvrata. Zato vse osoblive točke ne sut točky vozvrata, a historijny izraz ne trěba bez dodatnogo dokazanja sužati do obyčnoj točky vozvrata. Razuměnje konkretnoj geometrije ima prednost prěd samym obrazom dviganja nazad.\par
\noindent\emph{Reason:} Osobliva točka is a general name for a singular point of a curve; točka vozvrata retains the historical name of a particular type. The archive gives Russian osobaja točka and Ukrainian osoblyva točka, rendering Rückkehrpunkt through the literal image of return. The distinction matters in paper 14: one coordinate tuple can correspond to several points of the normalized curve or be attained with higher order. The latter case is connected with a higher return point. Thus not every singular point is a cusp, and the historical expression should not be narrowed to an ordinary cusp without further evidence. The particular geometry takes priority over the image of moving backwards.\par
\noindent{\footnotesize Izvorne zapisy / Original records: HIST-g076-017.}\par
\smallskip\noindent\emph{Sravnjenje s tekstom:} Obča singularnost ne znači vsegda točku vozvrata.\par
\noindent\emph{Textual comparison:} A general singularity is not necessarily a cusp.\par
\noindent{\footnotesize TEXT-COMP-693; pasaže / passages: 9958, 9965, 9967; \texttt{CORPUS\_FORM\_COMPARISONS.json}.}\par
\end{minipage}\par
\par\medskip\Needspace{7\baselineskip}
\noindent\begin{minipage}{\linewidth}
\hypertarget{TERM-1637}{}
\noindent\textbf{singularne sistemy vrědnostij}\quad{\footnotesize\texttt{TERM-1637}}\par
\noindent\emph{Iztočnik / Source:} singuläre Wertsysteme\par
\noindent\emph{Rusky / Russian:} сингулярные системы значений\par
\noindent\emph{Ukrajinsky / Ukrainian:} сингулярні системи значень\par
\smallskip\noindent\emph{Razlog:} Singularne sistemy vrědnostij označajut izključne n-torky parametrov ili koeficientov, ne patologične čisla vzete oddělno. Arhivny razlog čuva slovo sistema, zato že neuspěh parametrnogo prědstavjenja zavisi od sovkupnosti cěnnostij. Nynějšna rabota 11 ima sistemy cěnnostij. Argument s imenovateljami pokazuje, že pri nenulovom produktu $H_1\cdots H_n$ prědstavjenje uspěvaje; možne izključne sistemy ležet tam, kde produkt izčezaje. To jest potrěbno uslovje izključnosti, ne bezuslovna ekvivalentnost. Daljše osobno izslědovanje oprěděljaje, ktore od tyh sistem stvarno odgovarjajut iskanym uravnenjam. Zaměna vrědnostij na cěnnostij jest oddělna leksična historija i ne smije změniti tu logičnu ogradu.\par
\noindent\emph{Reason:} Singularne sistemy vrědnostij denotes exceptional tuples of parameters or coefficients, not pathological numbers taken individually. The recorded reason retains system because failure of the parametrization depends on the values collectively. Current paper 11 has sistemy cěnnostij. The denominator argument shows that the representation succeeds when $H_1\cdots H_n$ is nonzero; potentially exceptional systems lie where it vanishes. This is a necessary condition for exceptionality, not an unconditional equivalence. Further investigation determines which such systems actually correspond to the equations sought. The change from vrědnostij to cěnnostij is a separate lexical history and must not alter that logical qualification.\par
\noindent{\footnotesize Izvorne zapisy / Original records: HIST-g073-014.}\par
\smallskip\noindent\emph{Sravnjenje s tekstom:} Izčezanje produkta imenovateljev ostavljaje osobno pytanje o rěšenijah.\par
\noindent\emph{Textual comparison:} Vanishing of the denominator product leaves a separate question about solutions.\par
\noindent{\footnotesize TEXT-COMP-694; pasaže / passages: 8754, 8803, 8805, 8809, 8812, 8814, 8816, 8818, 8821, 8826, 8828, 8830, 8832; \texttt{CORPUS\_FORM\_COMPARISONS.json}.}\par
\end{minipage}\par
\par\medskip\Needspace{7\baselineskip}
\noindent\begin{minipage}{\linewidth}
\hypertarget{TERM-1639}{}
\noindent\textbf{togda $o$ jest kolco glavnyh idealov}\quad{\footnotesize\texttt{TERM-1639}}\par
\noindent\emph{Iztočnik / Source:} so wird $o$ Hauptidealring\par
\noindent\emph{Rusky / Russian:} то $o$ является кольцом главных идеалов\par
\noindent\emph{Ukrajinsky / Ukrainian:} то $o$ є кільцем головних ідеалів\par
\smallskip\noindent\emph{Razlog:} Togda $o$ jest kolco glavnyh idealov čuva polnu nazvu rezultata vměsto skračenja PID. Arhiv opiraje toj izbor na ranějšnju terminologičnu konvenciju. Nynějšny §8 raboty 31 ima prsten glavnyh idealov. Tvrdženje govori o multiplikacijskom prstenu s jedinym prostym idealom, različnym od nulovogo i jediničnogo; dokaz vybira element, ktory poraždaje toj ideal, i upotrěbjaje ranějše razloženje idealov v produkt stepenjev. Ne možno iztrgnuti frazu i pripisati toj rezultat vsakomu proizvoljnomu prstenu s jedinym prostym idealom. Polna nazva pomaga čitatelju slediti logiku, ale ne zaměnjaje hipotezy i ne dodava novyh aksiom modernoj oblasti glavnyh idealov.\par
\noindent\emph{Reason:} Togda $o$ jest kolco glavnyh idealov retains the full name of the conclusion rather than switching to PID. The archive bases this choice on an earlier terminology convention. Current paper 31, §8 has prsten glavnyh idealov. The assertion concerns a multiplication ring with just one prime ideal distinct from the zero and unit ideals; the proof selects an element generating that ideal and uses the earlier factorization of ideals into products of powers. The phrase cannot be detached and asserted of every arbitrary ring with one prime ideal. Spelling out the name helps the reader follow the argument but does not replace its hypotheses or add the axioms of a modern principal ideal domain.\par
\noindent{\footnotesize Izvorne zapisy / Original records: HIST-g156-003.}\par
\smallskip\noindent\emph{Sravnjenje s tekstom:} Dokaz upotrěbjaje ranějše razloženje idealov v produkt stepenjev.\par
\noindent\emph{Textual comparison:} The proof uses the earlier factorization of ideals into products of powers.\par
\noindent{\footnotesize TEXT-COMP-696; pasaže / passages: 18823; \texttt{CORPUS\_FORM\_COMPARISONS.json}.}\par
\end{minipage}\par
\par\medskip\Needspace{7\baselineskip}
\noindent\begin{minipage}{\linewidth}
\hypertarget{TERM-1643}{}
\noindent\textbf{specializacija}\quad{\footnotesize\texttt{TERM-1643}}\par
\noindent\emph{Iztočnik / Source:} Spezialisierung\par
\noindent\emph{Rusky / Russian:} специализация\par
\noindent\emph{Ukrajinsky / Ukrainian:} спеціалізація\par
\smallskip\noindent\emph{Razlog:} Specializacija čuva internacionalny matematičny korenj, razpoznavajemy takože v ruskom i ukrajinskom zapisu. Četyri arhivne zapisy davajut kontinuitet rabotam 2 i 24; oni ne vyžadajut novogo domašnjego slova samo zato, že ciljny jezyk jest skonstruovany. V rabotě 2 konkretne zaměny a), a'), b), b') pokazyvajut, čto specializacija děla. V rabotě 24 slovo opisuje zaměnu neoprěděljenyh veličin. Vyraženje ne uključaje samo po sebě ohranjenje vseh svojstv: vsako takovo tvrdženje trěbuje svoje uslovja. Jasnost zato proizhodi iz formul vměstě s ustaljenym imenikom, ne iz imenika oddělno.\par
\noindent\emph{Reason:} Specializacija retains an international mathematical stem also recognizable in the recorded Russian and Ukrainian forms. Four historical records establish continuity across papers 2 and 24; they do not demand a new native coinage merely because the target language is constructed. In paper 2 the explicit substitutions a), a'), b), b') show what specialization does. In paper 24 the word describes replacing indeterminates. It does not itself imply preservation of every property: each such assertion requires its own conditions. Clarity therefore comes from the formulas together with the established noun, not from the noun alone.\par
\noindent{\footnotesize Izvorne zapisy / Original records: HIST-g004-006, HIST-g013-008, HIST-g014-011, HIST-g097-017.}\par
\smallskip\noindent\emph{Sravnjenje s tekstom:} Konkretne zaměny oprěděljajut smysl imenika.\par
\noindent\emph{Textual comparison:} The actual substitutions determine the noun's meaning.\par
\noindent{\footnotesize TEXT-COMP-700; pasaže / passages: 429, 430, 433, 1913, 2079, 15496, 15498; \texttt{CORPUS\_FORM\_COMPARISONS.json}.}\par
\end{minipage}\par
\par\medskip\Needspace{7\baselineskip}
\noindent\begin{minipage}{\linewidth}
\hypertarget{TERM-1644}{}
\noindent\textbf{specializacija / specializovati}\quad{\footnotesize\texttt{TERM-1644}}\par
\noindent\emph{Iztočnik / Source:} Spezialisierung\par
\noindent\emph{Rusky / Russian:} специализация / специализировать\par
\noindent\emph{Ukrajinsky / Ukrainian:} спеціалізація / спеціалізувати\par
\smallskip\noindent\emph{Razlog:} Para specializacija/specializovati dava imenik za operaciju i glagol za jej izvršenje. Arhivny razlog predpočitaje toj internacionalny obrazec i konsistentnost s prědhodečimi sekcijami; tvrdženje o medžuslovjanskoj prozračnosti jest kvalitativna ocěna. V §4 raboty 24 pri beskonečnom osnovnom polju veličiny možno izbrati tako, že norma, forma elementarnyh děliteljev i dimenzija ostavajut shranjene. Ipak ta forma ne mora ostati prosta. Zato opis operacije mora čuvati razliku medžu svojstvami, ktore izbor ohranja, i svojstvami, za ktore takovo tvrdženje ne jest dano.\par
\noindent\emph{Reason:} The pair specializacija/specializovati supplies a noun for the operation and a verb for performing it. The recorded reason prefers the international pattern and continuity with preceding sections; its judgement of pan-Slavic transparency is qualitative. In paper 24, §4, over an infinite base field the values can be selected so that norm, elementary-divisor form and dimension are preserved. Yet that form need not remain irreducible. Describing the operation must therefore retain the distinction between properties the choice preserves and properties for which no such claim is made.\par
\noindent{\footnotesize Izvorne zapisy / Original records: HIST-g098-011.}\par
\smallskip\noindent\emph{Sravnjenje s tekstom:} Dimenzija može byti ohranjena, kogda nerazložimost ne jest garantovana.\par
\noindent\emph{Textual comparison:} Dimension can be preserved while irreducibility is not guaranteed.\par
\noindent{\footnotesize TEXT-COMP-701; pasaže / passages: 15527; \texttt{CORPUS\_FORM\_COMPARISONS.json}.}\par
\end{minipage}\par
\par\medskip\Needspace{7\baselineskip}
\noindent\begin{minipage}{\linewidth}
\hypertarget{TERM-1645}{}
\noindent\textbf{specializacija na relativne diskriminanty čislovyh polj}\quad{\footnotesize\texttt{TERM-1645}}\par
\noindent\emph{Iztočnik / Source:} Spezialisierung auf Relativdiskriminanten von Zahlkörpern\par
\noindent\emph{Rusky / Russian:} специализация к относительным дискриминантам числовых полей\par
\noindent\emph{Ukrajinsky / Ukrainian:} спеціалізація до відносних дискримінантів числових полів\par
\smallskip\noindent\emph{Razlog:} Specializacija na relativne diskriminanty čislovyh polj označaje prěhod od občej teorije k konkretnomu čislovo-teoretičnomu slučaju. Arhiv čuva ranějše relativny diskriminant i izrično slovo specializacija, zato že rezultat jest vyvedeny kak sledstvo, ne vvedeny kako nova nezavisna teorema. V §8 raboty 31 uslovje govori o prostom ideale, njegovom vhode v relativny diskriminant i dělimosti kvadratom prostogo ideala v razširjenom polju. Relativny i oba poredky sut sučstvene kvalifikacije. Opisna zaměna na prosto diskriminant polja by skryla osnovno polje i směr razširjenja.\par
\noindent\emph{Reason:} Specializacija na relativne diskriminanty čislovyh polj marks passage from the general theory to a particular number-theoretic case. The archive retains the earlier relative-discriminant term and the explicit word specialization because the result is derived as a corollary, not introduced as an unrelated theorem. In paper 31, §8 the condition concerns a prime ideal, its occurrence in the relative discriminant, and divisibility by a prime-ideal square in the extension. Relative and the two orders are substantive qualifications. Replacing the phrase with merely field discriminant would hide the base field and the direction of extension.\par
\noindent{\footnotesize Izvorne zapisy / Original records: HIST-g164-007.}\par
\smallskip\noindent\emph{Sravnjenje s tekstom:} Relativny diskriminant čuva svezu s osnovnym poljem i razširjenjem.\par
\noindent\emph{Textual comparison:} The relative discriminant retains the relation to the base field and extension.\par
\noindent{\footnotesize TEXT-COMP-702; pasaže / passages: 19107; \texttt{CORPUS\_FORM\_COMPARISONS.json}.}\par
\end{minipage}\par
\par\medskip\Needspace{7\baselineskip}
\noindent\begin{minipage}{\linewidth}
\hypertarget{TERM-1646}{}
\noindent\textbf{specializacija někojih neopreděljenyh}\quad{\footnotesize\texttt{TERM-1646}}\par
\noindent\emph{Iztočnik / Source:} Spezialisierung einiger Unbestimmten\par
\noindent\emph{Rusky / Russian:} специализация некоторых неопределенных\par
\noindent\emph{Ukrajinsky / Ukrainian:} спеціалізація деяких невизначених\par
\smallskip\noindent\emph{Razlog:} Specializacija někojih neoprěděljenyh označaje zaměnu vybranyh neoprěděljenyh, ne izčezanje neizvěstnosti v obyčnom jezyku. Zapisany razlog čuva objasnjenje, začemu oblast odobraženja ne mora izpolnjati uslovje 4a). V §9 raboty 6 kvocient, ktory prvotno ne byl polinom, može po specializaciji stati polinomom; priměr upotrěbjaje $x_3=0$. Zato slovo někojih i konkretna zaměna sut važnejše od samej podobnosti nacionalnyh form. Specializacija može ohraniti jedno potrěbno svojstvo i utratiti drugo; prěvod ne směje dodati bezuslovnu svezu medžu iztočnoju i obraznoju oblastju.\par
\noindent\emph{Reason:} Specializacija někojih neoprěděljenyh means substitution for selected indeterminates, not the removal of uncertainty in ordinary language. The recorded reason preserves the explanation of why the image domain need not satisfy condition 4a). In paper 6, §9 a quotient that was not originally a polynomial can become one after specialization; the example uses $x_3=0$. The qualifier some and the actual substitution therefore matter more than resemblance between national forms. Specialization may preserve one required property and lose another; translation must not add an unconditional equivalence between source and image domains.\par
\noindent{\footnotesize Izvorne zapisy / Original records: HIST-g049-017.}\par
\smallskip\noindent\emph{Sravnjenje s tekstom:} Kvocient može stati polinomom poslě zaměny neoprěděljenoj.\par
\noindent\emph{Textual comparison:} A quotient can become a polynomial after substitution for an indeterminate.\par
\noindent{\footnotesize TEXT-COMP-703; pasaže / passages: 6033; \texttt{CORPUS\_FORM\_COMPARISONS.json}.}\par
\end{minipage}\par
\par\medskip\Needspace{7\baselineskip}
\noindent\begin{minipage}{\linewidth}
\hypertarget{TERM-1647}{}
\noindent\textbf{specializacija $u_{\mu\nu}=\bar u_{\mu\nu}$}\quad{\footnotesize\texttt{TERM-1647}}\par
\noindent\emph{Iztočnik / Source:} Spezialisierung $u_{\mu\nu}=\bar u_{\mu\nu}$\par
\noindent\emph{Rusky / Russian:} специализация $u_{\mu\nu}=\bar u_{\mu\nu}$\par
\noindent\emph{Ukrajinsky / Ukrainian:} спеціалізація $u_{\mu\nu}=\bar u_{\mu\nu}$\par
\smallskip\noindent\emph{Razlog:} Specializacija $u_{\mu\nu}=\bar u_{\mu\nu}$ izrično kaže, ktore neoprěděljene veličiny sut zaměnjene i kak označeny sut jih nove cěnnosti. Arhiv čuva internacionalny korenj zato, že zaključny odstavek raboty 22 razgledaje ohranjenje regularnosti. Tekst trěbuje podobno izbranu specializaciju nad poljem karakteristiky nul; ne govori, že vsaka proizvoljna zaměna uspěvaje. Nynějšna pasaža takože razlikuje obstajanje teorema od avtomatičnogo specializovanja osnovnogo modula i elementarnyh děliteljev. Ohranjenje indeksov i crty nad $u$ zato čini čest smysla, ne samo tipografiju.\par
\noindent\emph{Reason:} Specializacija $u_{\mu\nu}=\bar u_{\mu\nu}$ specifies which indeterminates are replaced and how their new values are denoted. The archive retains the international stem because paper 22's closing discussion concerns preservation of regularity. The text requires a suitable specialization over a characteristic-zero field; it does not say that every arbitrary substitution succeeds. The current passage also distinguishes survival of the theorem from automatic specialization of the underlying module and elementary divisors. Preserving both indices and the bar on $u$ is consequently part of the meaning, not mere typography.\par
\noindent{\footnotesize Izvorne zapisy / Original records: HIST-g088-011, HIST-g092-042.}\par
\smallskip\noindent\emph{Sravnjenje s tekstom:} Trěba podobny izbor cěnnostij, ne proizvoljnu zaměnu.\par
\noindent\emph{Textual comparison:} A suitable choice of values is required, not an arbitrary substitution.\par
\noindent{\footnotesize TEXT-COMP-704; pasaže / passages: 14827; \texttt{CORPUS\_FORM\_COMPARISONS.json}.}\par
\end{minipage}\par
\par\medskip\Needspace{7\baselineskip}
\noindent\begin{minipage}{\linewidth}
\hypertarget{TERM-1648}{}
\noindent\textbf{specialna modulna baza / dopolniti libovoljnu bazu od $p$}\quad{\footnotesize\texttt{TERM-1648}}\par
\noindent\emph{Iztočnik / Source:} spezielle Modulbasis / beliebige Basis von $p$ ergänzen\par
\noindent\emph{Rusky / Russian:} специальный модульный базис / дополнить произвольный базис $p$\par
\noindent\emph{Ukrajinsky / Ukrainian:} спеціальна модульна база / доповнити довільну базу $p$\par
\smallskip\noindent\emph{Razlog:} Specialna modulna baza i dopolniti libovoljnu bazu opisujut različne, ale svezane vybory v dokazu. Arhiv čuva libovoljnu v smyslu proizvoljno vybranoj, ne slučajno izbranoj bazy; od $p$ vyražaje pripadnost idealu. V §6 raboty 31 možno upotrěbiti ranějše skonstruovanu specialnu bazu ili dopolniti bazu prostogo ideala do bazy cělogo prstena. To daje argument za izčezanje diskriminanta pri danyh prědpoloženjah. Slovo specialna ne znači numerično optimizovana, a libovoljna ne zaměnjaje dokaz, že dopolnjenje obstaja.\par
\noindent\emph{Reason:} Specialna modulna baza and dopolniti libovoljnu bazu describe different but related choices in a proof. The archive keeps libovoljnu in the sense of arbitrarily chosen, not randomly sampled; od $p$ expresses belonging to the ideal. In paper 31, §6 one can use the specially constructed earlier basis or extend a basis of the prime ideal to one of the whole ring. This supplies the argument for vanishing of the discriminant under the stated assumptions. Special does not mean numerically optimized, and arbitrary does not replace the proof that an extension exists.\par
\noindent{\footnotesize Izvorne zapisy / Original records: HIST-g143-003.}\par
\smallskip\noindent\emph{Sravnjenje s tekstom:} Proizvoljna baza ideala dopolnjaje se pri navedenyh hipotezah.\par
\noindent\emph{Textual comparison:} An arbitrary ideal basis is extended under the stated hypotheses.\par
\noindent{\footnotesize TEXT-COMP-705; pasaže / passages: 18308, 18383; \texttt{CORPUS\_FORM\_COMPARISONS.json}.}\par
\end{minipage}\par
\par\medskip\Needspace{7\baselineskip}
\noindent\begin{minipage}{\linewidth}
\hypertarget{TERM-1649}{}
\noindent\textbf{specialne $P$-modulne bazy}\quad{\footnotesize\texttt{TERM-1649}}\par
\noindent\emph{Iztočnik / Source:} Spezielle $P$-Modulbasen\par
\noindent\emph{Rusky / Russian:} специальные $P$-модульные базисы\par
\noindent\emph{Ukrajinsky / Ukrainian:} спеціальні $P$-модульні бази\par
\smallskip\noindent\emph{Razlog:} Specialne $P$-modulne bazy čuva internacionalny pridavnik i ranějše modulna baza. Arhiv opravdava specialne tehničnoju razpoznavajemošću, ale ne dava izměrenje razumělivosti. V rabotě 31 $P$ jest polje, nad ktorym bazy sut linearno nezavisne; dokaz dopolnja bazu jednogo ideala do bazy drugogo. To jest osobno prilagodžena baza dlja razgledanogo vključenja. Ne možna obča zaměna na libovoljny modul nad proizvoljnym prstenom, gde podobno dopolnjenje ne musi obstajati. Bukva $P$ i konačnorazměrno prědpoloženje sut zato nositeli matematičnoj informaciji.\par
\noindent\emph{Reason:} Specialne $P$-modulne bazy preserves the international adjective and the earlier module-basis term. The archive justifies specialne by technical recognizability, not a measurement of comprehension. In paper 31, $P$ is the field over which the bases are linearly independent; the proof extends a basis of one ideal to a basis of another. This is a basis adapted to the inclusion under discussion. The statement cannot be generalized to arbitrary modules over arbitrary rings, where such extension need not exist. The letter $P$ and the finite-dimensional hypothesis therefore carry mathematical information.\par
\noindent{\footnotesize Izvorne zapisy / Original records: HIST-g141-001.}\par
\smallskip\noindent\emph{Sravnjenje s tekstom:} Linearno nezavisna modulna baza trěbuje ukazanu oblast koeficientov.\par
\noindent\emph{Textual comparison:} The linearly independent module basis requires the stated coefficient domain.\par
\noindent{\footnotesize TEXT-COMP-706; pasaže / passages: 18308; \texttt{CORPUS\_FORM\_COMPARISONS.json}.}\par
\end{minipage}\par
\par\medskip\Needspace{7\baselineskip}
\noindent\begin{minipage}{\linewidth}
\hypertarget{TERM-1650}{}
\noindent\textbf{specialno dvojnorazovo svijanje}\quad{\footnotesize\texttt{TERM-1650}}\par
\noindent\emph{Iztočnik / Source:} spezielle zweimalige Faltung\par
\noindent\emph{Rusky / Russian:} специальная двукратная свертка\par
\noindent\emph{Ukrajinsky / Ukrainian:} спеціальне дворазове згортання\par
\smallskip\noindent\emph{Razlog:} Dvojnorazovo svijanje bylo izbrano, da izrično kaže dvokratnost kontrakcije. Arhiv priznava cenu dolžiny i podava dvojno kako prostějši kandidat za daljše razsmotrenje; čislena sravnitelna vaga ne jest zapisana. Nynějšny §13 raboty 2 ima specialnu dvojnorazovu kontrakciju s zadanoju razložimoju formoju. Slovo kontrakcija jest pozdnějša leksična forma, ale čislova kvalifikacija ostava. Pri možnom skračenji do dvojno trěba čuvati razliku medžu dvokratnoju kontrakcijeju i dvojnym objektom. Neobhodima jest matematična prozračnost, ne maksimalna dolžina frazy sama po sebě.\par
\noindent\emph{Reason:} Dvojnorazovo svijanje was chosen to express explicitly the twofold character of the contraction. The archive acknowledges the cost in length and offers simpler dvojno for later consideration; it records no numerical comparative weight. Current paper 2, §13 has a special twofold contraction with the specified decomposable form. Kontrakcija is the later lexical form, but the numerical qualification remains. If the wording is shortened to dvojno, the distinction between a twofold contraction and a double object must survive. Mathematical transparency matters, not maximal explicit length for its own sake.\par
\noindent{\footnotesize Izvorne zapisy / Original records: HIST-g014-010.}\par
\smallskip\noindent\emph{Sravnjenje s tekstom:} Aktualna fraza čuva čislovu kvalifikaciju operacije.\par
\noindent\emph{Textual comparison:} The current phrase preserves the operation's numerical qualification.\par
\noindent{\footnotesize TEXT-COMP-707; pasaže / passages: 2079, 2083; \texttt{CORPUS\_FORM\_COMPARISONS.json}.}\par
\end{minipage}\par
\par\medskip\Needspace{7\baselineskip}
\noindent\begin{minipage}{\linewidth}
\hypertarget{TERM-1651}{}
\noindent\textbf{sl?d}\quad{\footnotesize\texttt{TERM-1651}}\par
\noindent\emph{Iztočnik / Source:} Spur\par
\noindent\emph{Rusky / Russian:} [Bukvy ne sut zachovane / Lettering not recoverable]\par
\noindent\emph{Ukrajinsky / Ukrainian:} [Bukvy ne sut zachovane / Lettering not recoverable]\par
\smallskip\noindent\emph{Razlog:} Poškodženo sl?d s nečitljivymi ruskim i ukrajinskim poljami ne dozvaljaje vosstanoviti točnu staršu ortografiju. Čitlivy razlog ipak izrično svęzuje izbor s ranějšim slědom v rabotah 34 i 37 i daje prednost slovjanskomu korenju prěd internacionalnym zaimanjem. Nynějšna rabota 40 upotrěbjaje reducirany slěd pri konstrukciji separabilnogo razpadnogo polja. Nenulovy slěd tam vybira element izvan osnovnogo polja; slěd skalara izčezaje v danej pozitivnoj karakteristice. To potvrđuje tehničny smysl, ale ne izginule bukvy. Jezyčna motivacija i sostojanje arhiva morajut ostati različne svědectva.\par
\noindent\emph{Reason:} Damaged sl?d, with unreadable Russian and Ukrainian fields, does not permit recovery of the exact earlier spelling. The surviving rationale nevertheless explicitly links the choice to earlier trace terminology in papers 34 and 37 and prefers a Slavic root to an international loan. Current paper 40 uses reduced trace in constructing a separable splitting field. A nonzero trace selects an element outside the base field there; scalar traces vanish in the positive characteristic under discussion. This confirms the technical sense, not the lost lettering. Linguistic motivation and the archive's condition remain different kinds of evidence.\par
\noindent{\footnotesize Izvorne zapisy / Original records: HIST-g210-011.}\par
\smallskip\noindent\emph{Sravnjenje s tekstom:} Nenulovy redukovany slěd razlikuje izbrany element od skalarov.\par
\noindent\emph{Textual comparison:} A nonzero reduced trace distinguishes the selected element from scalars.\par
\noindent{\footnotesize TEXT-COMP-708; pasaže / passages: 23975; \texttt{CORPUS\_FORM\_COMPARISONS.json}.}\par
\end{minipage}\par
\par\medskip\Needspace{7\baselineskip}
\noindent\begin{minipage}{\linewidth}
\hypertarget{TERM-1652}{}
\noindent\textbf{slěd}\quad{\footnotesize\texttt{TERM-1652}}\par
\noindent\emph{Iztočnik / Source:} Spur\par
\noindent\emph{Rusky / Russian:} след\par
\noindent\emph{Ukrajinsky / Ukrainian:} слід\par
\smallskip\noindent\emph{Razlog:} Slěd prěvodi Spur ustaljenym slovjanskym korenjem, v sravnjenju s ruskim sled i ukrajinskim slid. Arhiv čuva takože množstveno slědy i povtorjanje imenika pri definiciji, da sveza s němčinskim tekstom ostane jasna. V §4 raboty 31 definicija jest suma diagonalnyh elementov prirědženej matrice; §8 čita polje razširjenja kako modul nad osnovnym poljem. Zato slěd ne jest sled dviganja ili ostatok objekta. Pri korpusnom poisku niz slěd takože nastupaje v slěduje; prosty počet podnizov ne daje počet vseh terminologičnyh pojavjenj. Svědectvom sut konkretne definicije i jih kontekst.\par
\noindent\emph{Reason:} Slěd renders Spur with an established Slavic root, compared with Russian sled and Ukrainian slid. The archive also preserves plural slědy and repetition of the noun in the definition to keep the relation to the German clear. In paper 31, §4 the definition is the sum of diagonal entries of the assigned matrix; §8 treats the extension field as a module over its base field. This is not a track left by movement or a remnant of an object. In corpus searches, the string slěd also occurs in slěduje; a substring count is not a count of technical occurrences. The evidence is the located definitions in context.\par
\noindent{\footnotesize Izvorne zapisy / Original records: HIST-g130-002, HIST-g157-001.}\par
\smallskip\noindent\emph{Sravnjenje s tekstom:} Suma diagonalnyh elementov čuva konkretnu definiciju.\par
\noindent\emph{Textual comparison:} The diagonal sum preserves the concrete definition.\par
\noindent{\footnotesize TEXT-COMP-709; pasaže / passages: 18024, 18026, 18857, 18861; \texttt{CORPUS\_FORM\_COMPARISONS.json}.}\par
\end{minipage}\par
\par\medskip\Needspace{7\baselineskip}
\noindent\begin{minipage}{\linewidth}
\hypertarget{TERM-1653}{}
\noindent\textbf{slěd / matrična predstavjenost / konjugovane elementy}\quad{\footnotesize\texttt{TERM-1653}}\par
\noindent\emph{Iztočnik / Source:} Spur / Matrizendarstellung / konjugierte Elemente\par
\noindent\emph{Rusky / Russian:} след / матричное представление / сопряженные элементы\par
\noindent\emph{Ukrajinsky / Ukrainian:} слід / матричне подання / спряжені елементи\par
\smallskip\noindent\emph{Razlog:} Slěd, matrična prědstavjenost i konjugovane elementy svęzany sut jednoj važnoju razlikoj definicij. Arhiv opisuje dopolnjenje teksta po skanu i dava prednost internacionalnym matrična i konjugovane prěd novymi kalkami. Nynějšny §7 raboty 31 govori, že definicija prěz konjugovane elementy jest pravilna za izhodny poredok, ale ne potrěbno za jego prsten ostatkov. Matrična definicija dovoljava prěnesti slěd i diskriminant homomorfizmom. Zato ti slova ne sut svobodna stilistična zaměna: oni označajut sposob, ktory opravdava prěhod. Novo čitanje pasaže ne zaměnjaje zabeleženu staršu prověrku skana.\par
\noindent\emph{Reason:} Slěd, matrična prědstavjenost and konjugovane elementy are linked by an important distinction between definitions. The archive reports expansion from a scan and prefers international matrična and konjugovane to new calques. Current paper 31, §7 says that a definition through conjugate elements is valid for the original order but not necessarily for its residue ring. The matrix definition permits trace and discriminant to pass through the homomorphism. These words are therefore not freely interchangeable stylistic variants: they identify the construction justifying the passage. Fresh reading of the paragraph is not a repeat of the recorded earlier scan examination.\par
\noindent{\footnotesize Izvorne zapisy / Original records: HIST-g152-004.}\par
\smallskip\noindent\emph{Sravnjenje s tekstom:} Matrična konstrukcija ne trěbuje prědstavjenja prěz konjugovane elementy.\par
\noindent\emph{Textual comparison:} The matrix construction does not require a description through conjugate elements.\par
\noindent{\footnotesize TEXT-COMP-710; pasaže / passages: 18671; \texttt{CORPUS\_FORM\_COMPARISONS.json}.}\par
\end{minipage}\par
\par\medskip\Needspace{7\baselineskip}
\noindent\begin{minipage}{\linewidth}
\hypertarget{TERM-1654}{}
\noindent\textbf{slěd / slědy}\quad{\footnotesize\texttt{TERM-1654}}\par
\noindent\emph{Iztočnik / Source:} Spur / Spuren\par
\noindent\emph{Rusky / Russian:} след / следы\par
\noindent\emph{Ukrajinsky / Ukrainian:} слід / сліди\par
\smallskip\noindent\emph{Razlog:} Slěd/slědy podava edinstveno i množstveno čislo togo samogo termina. Arhiv izrično sravnjava rusky sled, ukrajinsky slid i poljsky ślad; to jest kvalitativna etimologična i disciplinarna podpora, ne sčitanje nezavisnyh glasov vseh slovjanskyh jezykov. V §24 raboty 34 slěd zavisi od prědstavjenja, jest linearny i ne měnjaje se pri ekvivalentnom prědstavjenju. Tekst dalje razlikuje glavny i redukovany slěd. Taka kvalifikacija ne směje izginuti, kogda imenik ostava jednaky: v različnih prědstavjenjah čisla ne muset byti ravne. Neutratena jest i svez s karakterami.\par
\noindent\emph{Reason:} Slěd/slědy gives singular and plural forms of the same term. The archive expressly compares Russian sled, Ukrainian slid and Polish ślad; this is qualitative etymological and disciplinary support, not a tally of independent votes from every Slavic language. In paper 34, §24 trace depends on the representation, is linear, and is unchanged by equivalent representations. The text further distinguishes principal and reduced trace. Such qualifications cannot disappear merely because the noun is unchanged: different representations need not give the same number. The connection with characters also remains.\par
\noindent{\footnotesize Izvorne zapisy / Original records: HIST-g202-001.}\par
\smallskip\noindent\emph{Sravnjenje s tekstom:} Klas prědstavjenja i kvalifikacije glavny/redukovany ostavajut važni.\par
\noindent\emph{Textual comparison:} The representation class and the principal/reduced qualifiers remain important.\par
\noindent{\footnotesize TEXT-COMP-711; pasaže / passages: 22175, 22177, 22179, 22181, 22183, 22185, 22187, 22203; \texttt{CORPUS\_FORM\_COMPARISONS.json}.}\par
\end{minipage}\par
\par\medskip\Needspace{7\baselineskip}
\noindent\begin{minipage}{\linewidth}
\hypertarget{TERM-1655}{}
\noindent\textbf{slěd elementa vzhodno k klasu}\quad{\footnotesize\texttt{TERM-1655}}\par
\noindent\emph{Iztočnik / Source:} Spur eines Elementes bezüglich einer Klasse\par
\noindent\emph{Rusky / Russian:} след элемента относительно класса\par
\noindent\emph{Ukrajinsky / Ukrainian:} слід елемента відносно класу\par
\smallskip\noindent\emph{Razlog:} Slěd elementa vzhodno k klasu izrično čuva klas kako parametr definicije. Tri jezyčne objasnjenja arhiva soglašajut se v toj motivaciji: ne govori se o neograničenom slědu, ale o slědu v danom klasu prědstavjenij. V nynějšnem §4 raboty 31 pojavjaje se po odnošenju k klasu i oznaka $S_{(\mathfrak a)}(\gamma)$. Zaměna bazy daje ekvivalentnu matricu i ne měnjaje sumu diagonalnyh elementov. To opravdava nezavisnost od izbranoj bazy, ne od svakogo možnogo klasa. Stara predložna konstrukcija jest ohranjena kako historijny izbor, a ne kako nynějšna preporuka.\par
\noindent\emph{Reason:} Slěd elementa vzhodno k klasu expressly preserves the class as a parameter of the definition. The archive's three language-specific rationales agree on that motivation: this is trace in a given representation class, not an unqualified trace. Current paper 31, §4 has po odnošenju k klasu and notation $S_{(\mathfrak a)}(\gamma)$. A change of basis gives an equivalent matrix and leaves the sum of diagonal entries unchanged. This justifies independence of the chosen basis, not independence of every possible class. The earlier prepositional construction is retained as a historical choice, not as the current recommendation.\par
\noindent{\footnotesize Izvorne zapisy / Original records: HIST-g135-001.}\par
\smallskip\noindent\emph{Sravnjenje s tekstom:} Nezavisnost od bazy ne jest bezuslovna nezavisnost od klasa.\par
\noindent\emph{Textual comparison:} Basis independence is not unconditional class independence.\par
\noindent{\footnotesize TEXT-COMP-712; pasaže / passages: 18024, 18026; \texttt{CORPUS\_FORM\_COMPARISONS.json}.}\par
\end{minipage}\par
\par\medskip\Needspace{7\baselineskip}
\noindent\begin{minipage}{\linewidth}
\hypertarget{TERM-1656}{}
\noindent\textbf{slěd elementa, vzety vzhodno k klasu}\quad{\footnotesize\texttt{TERM-1656}}\par
\noindent\emph{Iztočnik / Source:} Spur eines Elementes, genommen in bezug auf eine Klasse\par
\noindent\emph{Rusky / Russian:} след элемента, взятый относительно класса\par
\noindent\emph{Ukrajinsky / Ukrainian:} слід елемента, взятий відносно класу\par
\smallskip\noindent\emph{Razlog:} Vzety vzhodno k klasu prodolžuje klasovo zavisnu definiciju slěda i čuva objekt, ktory měnjaje se v slědujučem rezultatu. Arhivny razlog jest imenno prěhod medžu klasom i njegovymi komponentami. V §5 raboty 31 blokovo diagonalna matrica pokazuje, že slěd za direktnu sumu jest suma slědov komponentnyh klasov. To ne jest povtorjena definicija togo samogo neparametrizovanogo čisla. Nynějšny tekst ima vzety po odnošenju k klasu; proměna predloga ne měnjaje aditivnost ni to, čego slěd se tvori. Pri skračenji frazy trěba oba odnosa ostaviti vidlivymi.\par
\noindent\emph{Reason:} Vzety vzhodno k klasu continues the class-dependent definition of trace and preserves the parameter that changes in the following result. The recorded reason is precisely the passage between a class and its components. In paper 31, §5 a block-diagonal matrix shows that the trace for a direct sum is the sum of the traces for the component classes. This is not a repeated definition of the same unparameterized number. Current wording is vzety po odnošenju k klasu; changing the preposition alters neither additivity nor what the trace is taken of. A shorter phrase must still make both relationships clear.\par
\noindent{\footnotesize Izvorne zapisy / Original records: HIST-g139-002.}\par
\smallskip\noindent\emph{Sravnjenje s tekstom:} Blokova forma daje sumu komponentnyh slědov.\par
\noindent\emph{Textual comparison:} The block form yields the sum of the component traces.\par
\noindent{\footnotesize TEXT-COMP-713; pasaže / passages: 18181, 18183, 18185; \texttt{CORPUS\_FORM\_COMPARISONS.json}.}\par
\end{minipage}\par
\par\medskip\Needspace{7\baselineskip}
\noindent\begin{minipage}{\linewidth}
\hypertarget{TERM-1657}{}
\noindent\textbf{prosto slěd / nezavisno od klasa}\quad{\footnotesize\texttt{TERM-1657}}\par
\noindent\emph{Iztočnik / Source:} Spur schlechthin / unabhängig von der Klasse\par
\noindent\emph{Rusky / Russian:} просто след / независимо от класса\par
\noindent\emph{Ukrajinsky / Ukrainian:} просто слід / незалежно від класу\par
\smallskip\noindent\emph{Razlog:} Prosto slěd prěvodi schlechthin po jego funkciji v argumentu, ne doslovnym slovom za prostotu. Arhiv dava razlog: pri prstenah bez nulovyh děliteljev klasovy parametr može byti izpuščeny. Nynějšna primětka §4 raboty 31 objasnjaje, že baza nenulovogo ideala jest takože baza jediničnogo ideala v polju kvocientov; slědy i normy zato sodržajut tu samu konstrukciju. Odsutnost zavisimosti od klasa jest rezultat s hipotezami, ne svojstvo každogo slěda v vsakom prstenu. Prosto ne znači, že slěd jest elementarny, jediničny ili matematicky prosty objekt.\par
\noindent\emph{Reason:} Prosto slěd translates schlechthin by its argumentative function, not by a literal word for simplicity. The archive gives the reason: in rings without zero divisors the class parameter can be dropped. The current remark in paper 31, §4 explains that a basis of a nonzero ideal is also a basis of the unit ideal over the quotient field, so traces and norms come from the same construction. Class independence is a result with hypotheses, not a property of every trace in every ring. Prosto does not say that the trace is elementary, a unit, or a mathematically simple object.\par
\noindent{\footnotesize Izvorne zapisy / Original records: HIST-g136-006.}\par
\smallskip\noindent\emph{Sravnjenje s tekstom:} Izpuščenje klasa vyhodi iz navedenyh hipotez.\par
\noindent\emph{Textual comparison:} Omission of the class follows from the stated hypotheses.\par
\noindent{\footnotesize TEXT-COMP-714; pasaže / passages: 18065; \texttt{CORPUS\_FORM\_COMPARISONS.json}.}\par
\end{minipage}\par
\par\medskip\Needspace{7\baselineskip}
\noindent\begin{minipage}{\linewidth}
\hypertarget{TERM-1658}{}
\noindent\textbf{Štejnicova terminologija / razširjenje prvogo ili drugogo roda}\quad{\footnotesize\texttt{TERM-1658}}\par
\noindent\emph{Iztočnik / Source:} Steinitzsche Bezeichnung / Erweiterung erster oder zweiter Art\par
\noindent\emph{Rusky / Russian:} терминология Штейница / расширение первого или второго рода\par
\noindent\emph{Ukrajinsky / Ukrainian:} термінологія Штейніца / розширення першого або другого роду\par
\smallskip\noindent\emph{Razlog:} Štejnicova terminologija pripisyvaje razlikovanje prvogo i drugogo roda Steinitzu. Arhiv trěbuje čuvati historijne nazvy, zato že oni vhodet v kriterij diskriminanta i v daljše razděljenje slučajev. Nynějšna podnožna primětka raboty 31 definira prvy rod tym, že vsaky element jest korenj nerazložimogo polinoma s različnymi linearnymi faktorami v vhodnom razširjenju. To objasnjaje svez so separabilnošću. Drugy rod jest protivny slučaj, ne samo drugy stepen razširjenja. Moderno objasnjenje može pomoći čitatelju, ale ne směje ticho zaměniti avtorske nazvy ili odstraniti imenovanu historijnu pripisanost.\par
\noindent\emph{Reason:} Štejnicova terminologija attributes the first-kind/second-kind distinction to Steinitz. The archive requires retaining the historical names because they enter the discriminant criterion and the subsequent division into cases. The current footnote in paper 31 defines the first kind by every element being a root of an irreducible polynomial with distinct linear factors in a suitable extension. This explains the connection with separability. The second kind is the contrary case, not merely extension degree two. A modern explanation may help readers but must not silently replace the author's terminology or remove the named historical attribution.\par
\noindent{\footnotesize Izvorne zapisy / Original records: HIST-g124-007.}\par
\smallskip\noindent\emph{Sravnjenje s tekstom:} Definicija prědpokladaje različne korenje, ne stepen dva.\par
\noindent\emph{Textual comparison:} The definition concerns distinct roots, not degree two.\par
\noindent{\footnotesize TEXT-COMP-715; pasaže / passages: 17599; \texttt{CORPUS\_FORM\_COMPARISONS.json}.}\par
\end{minipage}\par
\par\medskip\Needspace{7\baselineskip}
\noindent\begin{minipage}{\linewidth}
\hypertarget{TERM-1659}{}
\noindent\textbf{m?sto}\quad{\footnotesize\texttt{TERM-1659}}\par
\noindent\emph{Iztočnik / Source:} Stelle\par
\noindent\emph{Rusky / Russian:} [Bukvy ne sut zachovane / Lettering not recoverable]\par
\noindent\emph{Ukrajinsky / Ukrainian:} [Bukvy ne sut zachovane / Lettering not recoverable]\par
\smallskip\noindent\emph{Razlog:} Poškodženo m?sto čuva samo čest staršej formy; rusky i ukrajinsky zapisy nemajut čitlivyh bukv. Ostaly razlog jest jasny: Stelle tu jest termin lokalnoj teorije čislovyh polj, ne obyčna pozicija. Nynějšna rabota 39 imaje město i izrično objasnjaje prěhod k p-adičnym dopolnjenjam, a takože vključaje beskonečne města. To dava nezavisno čitanje nynešnjego termina, ne vosstanovjenje arhiva. Pri upotrěbě lokalno-globalnogo teorema vse města ne možno bez opravdanja sužati do konečnyh. Slovo dolžno byti tumačeno vměstě s lokalnym poljem i uslovjami teorema.\par
\noindent\emph{Reason:} Damaged m?sto preserves only part of the earlier form; the Russian and Ukrainian fields have no readable lettering. The surviving reason is clear: Stelle is a term of local number-field theory, not an ordinary location. Current paper 39 has město and explicitly explains passage to p-adic completions while also including infinite places. This supplies an independent reading of the current term, not a restoration of the archive. In applying the local-global theorem, all places must not be narrowed without justification to finite places. The word must be explained together with its local field and the theorem's conditions.\par
\noindent{\footnotesize Izvorne zapisy / Original records: HIST-g209-011.}\par
\smallskip\noindent\emph{Sravnjenje s tekstom:} Lokalno-globalno tvrdženje vključaje oba vida měst.\par
\noindent\emph{Textual comparison:} The local-global assertion includes both kinds of place.\par
\noindent{\footnotesize TEXT-COMP-716; pasaže / passages: 23294, 23296; \texttt{CORPUS\_FORM\_COMPARISONS.json}.}\par
\end{minipage}\par
\par\medskip\Needspace{7\baselineskip}
\noindent\begin{minipage}{\linewidth}
\hypertarget{TERM-1660}{}
\noindent\textbf{město}\quad{\footnotesize\texttt{TERM-1660}}\par
\noindent\emph{Iztočnik / Source:} Stelle\par
\noindent\emph{Rusky / Russian:} место\par
\noindent\emph{Ukrajinsky / Ukrainian:} місце\par
\smallskip\noindent\emph{Razlog:} Město vyražaje aritmetično Stelle s pomoču obyčnogo slovjanskogo imenika, kotory tu ima specialny smysl. Arhiv sravnjava rusko mesto i ukrajinsko misce, a p-adičny kontekst nazyvaje razlikujučim priznakom. V rabotě 37 podindeks označaje lokalnu konstrukciju pri prostom ideale; teorem o normalnoj bazi uključaje uslovje, že toj ideal ne děli stepen razširjenja. Obyčna fraza imaje město v tom samom korpusu ne jest priměr togo termina, zato sčitanje samogo slova bylo by obmančivo. Čitatelj mora viděti aritmetičny objekt i hipotezu, ne samo kratku slovnikovu paru.\par
\noindent\emph{Reason:} Město renders arithmetic Stelle using an ordinary Slavic noun with a specialized sense here. The archive compares Russian mesto and Ukrainian misce, identifying the p-adic context as the disambiguating feature. In paper 37 the subscript denotes the local construction at a prime ideal; the normal-basis theorem includes the condition that the ideal not divide the extension degree. The ordinary phrase imaje město elsewhere in the corpus is not an occurrence of this technical term, so counting the word alone would be misleading. Readers need the arithmetic object and the hypothesis, not just a short dictionary pair.\par
\noindent{\footnotesize Izvorne zapisy / Original records: HIST-g207-004.}\par
\smallskip\noindent\emph{Sravnjenje s tekstom:} Uslovje o dělimosti stepena ostava čest teorema.\par
\noindent\emph{Textual comparison:} The condition involving divisibility of the degree remains part of the theorem.\par
\noindent{\footnotesize TEXT-COMP-717; pasaže / passages: 22849; \texttt{CORPUS\_FORM\_COMPARISONS.json}.}\par
\end{minipage}\par
\par\medskip\Needspace{7\baselineskip}
\noindent\begin{minipage}{\linewidth}
\hypertarget{TERM-1661}{}
\noindent\textbf{vsegda obstojna $o$-modulna baza}\quad{\footnotesize\texttt{TERM-1661}}\par
\noindent\emph{Iztočnik / Source:} stets existierende $o$-Modulbasis\par
\noindent\emph{Rusky / Russian:} всегда существующий $o$-модульный базис\par
\noindent\emph{Ukrajinsky / Ukrainian:} завжди існуючий $o$-модульний базис\par
\smallskip\noindent\emph{Razlog:} Vsegda obstojna modulna baza vyražaje egzistenciju v danoj oblasti, ne tvrdženje o libovoljnom modulu. Osoblivo važno jest, že baza tu može značiti konečnu porodžajuču sistemu: pasaža dopušča $m\ge n$, a linearno nezavisnu bazu razgleduje potom kako specialny slučaj. Zato ne možno bez dodatnogo uslovja tumačiti termin kako svobodnu bazu s točno $n$ elementami. Nynějšny prěvod sodrživaje vsegda obstojna v skobkah i piše prsten matematičnym simbolom; otsutstvo cěloj neprekinutoj frazy iz staršego slovnika ne znači, že smysl izčezl. Arhiv govori o obnovjenju po otisku, ale novyj pregled toj pasaže potvrđaje tekst prěvoda, ne samo po sobě věrnost otisku.\par
\noindent\emph{Reason:} The phrase asserts the existence of a module basis in the setting under discussion, not for every module. Crucially, basis here can mean a finite generating set: the passage allows $m\ge n$ and subsequently treats a linearly independent basis as a special case. It must therefore not be interpreted without an additional hypothesis as a free basis of exactly $n$ elements. The current translation places vsegda obstojna in parentheses and writes the ring symbol mathematically; failure to find the whole uninterrupted glossary phrase does not mean that its meaning was lost. The archive describes a restoration from the printed source, but this fresh examination confirms the translated passage, not by itself its fidelity to that print.\par
\noindent{\footnotesize Izvorne zapisy / Original records: HIST-g158-006.}\par
\smallskip\noindent\emph{Sravnjenje s tekstom:} Rozděl medžu porodžajučoju sistemoju i linearno nezavisnoju bazoju jest izričny.\par
\noindent\emph{Textual comparison:} The distinction between a generating set and a linearly independent basis is explicit.\par
\noindent{\footnotesize TEXT-COMP-718; pasaže / passages: 18897, 19003; \texttt{CORPUS\_FORM\_COMPARISONS.json}.}\par
\end{minipage}\par
\par\medskip\Needspace{7\baselineskip}
\noindent\begin{minipage}{\linewidth}
\hypertarget{TERM-1662}{}
\noindent\textbf{lu?}\quad{\footnotesize\texttt{TERM-1662}}\par
\noindent\emph{Iztočnik / Source:} Strahl\par
\noindent\emph{Rusky / Russian:} [Bukvy ne sut zachovane / Lettering not recoverable]\par
\noindent\emph{Ukrajinsky / Ukrainian:} [Bukvy ne sut zachovane / Lettering not recoverable]\par
\smallskip\noindent\emph{Razlog:} Luč jest aritmetičny termin v razsuždenju o provodniku i normovanju faktornyh sistemov, ne geometrična poluprěma. Starši zapis lu? ima poškodženu poslědnju bukvu, a rusko i ukrajinsko pole ne sut čitlive; jih ne možno upotrěbiti kako dvojnogo nezavisnogo potvrdženja. Nynějšna pasaža izrično piše luč i povezuje go s elementami jediničnogo klasa. To podkrepljaje kontinuitet pojma, ne dokazuje obstajanje oficijalnogo medžuslovjanskogo termina teorije klasnyh polj. Definiciju i kongruentne uslovja trěba čitati v matematičnom kontekstu; sam obraz luča jih ne zaměnjaje.\par
\noindent\emph{Reason:} Luč is an arithmetic term in the discussion of conductors and norms of factor systems, not a geometric half-line. The older record lu? has a damaged final letter, while its Russian and Ukrainian fields are unreadable; these cannot serve as two independent confirmations. The current passage explicitly writes luč and connects it with elements of the identity class. This supports continuity of the concept, not the existence of an official Interslavic class-field-theory term. The definition and congruence conditions must be read in their mathematical context; the image of a ray does not supply them.\par
\noindent{\footnotesize Izvorne zapisy / Original records: HIST-g209-014.}\par
\smallskip\noindent\emph{Sravnjenje s tekstom:} Aritmetičny kontekst ne jest geometrična definicija.\par
\noindent\emph{Textual comparison:} The arithmetic context is not a geometric definition.\par
\noindent{\footnotesize TEXT-COMP-719; pasaže / passages: 23319; \texttt{CORPUS\_FORM\_COMPARISONS.json}.}\par
\end{minipage}\par
\par\medskip\Needspace{7\baselineskip}
\noindent\begin{minipage}{\linewidth}
\hypertarget{TERM-1663}{}
\noindent\textbf{lu?ne klasy}\quad{\footnotesize\texttt{TERM-1663}}\par
\noindent\emph{Iztočnik / Source:} Strahlklassen\par
\noindent\emph{Rusky / Russian:} [Bukvy ne sut zachovane / Lettering not recoverable]\par
\noindent\emph{Ukrajinsky / Ukrainian:} [Bukvy ne sut zachovane / Lettering not recoverable]\par
\smallskip\noindent\emph{Razlog:} Lučne klasy prodolžajut paru luč / lučny, aby jedno aritmetično ponjatje imělo prepoznavajemu slovotvornu vez. Arhivny oblik lu?ne čuva poškodženje, ale nynějšna rabota 41 piše lučne klasy. Tam sut različene absolutne idealne klasy v gornjem polju i lučne klasy v dolnjem polju; izpustiti lučne značilo by izgubiti imenno razliku, na ktoroj stoji sravnjenje. Rusky i ukrajinsky zapisy v arhivu ne možno obnoviti iz znakov pytanja. Tvrdženje staršego razloga o prozračnosti jest jezyčna ocěnka, ne rezultat testovanja čitateljev.\par
\noindent\emph{Reason:} Lučne klasy continues the luč / lučny word family so that the arithmetic concept has a recognizable derivational connection. The archived spelling lu?ne retains damage, whereas current paper 41 writes lučne klasy. It distinguishes absolute ideal classes in the upper field from ray classes in the lower field; dropping lučne would remove precisely the distinction on which the comparison depends. The Russian and Ukrainian archival forms cannot be recovered from question marks. The older rationale's claim of transparency is a linguistic judgment, not a reader-test result.\par
\noindent{\footnotesize Izvorne zapisy / Original records: HIST-g211-008.}\par
\smallskip\noindent\emph{Sravnjenje s tekstom:} Gornje i dolnje polje imajut tu različne klasne razdělby.\par
\noindent\emph{Textual comparison:} Upper and lower fields have different class partitions here.\par
\noindent{\footnotesize TEXT-COMP-720; pasaže / passages: 24295; \texttt{CORPUS\_FORM\_COMPARISONS.json}.}\par
\end{minipage}\par
\par\medskip\Needspace{7\baselineskip}
\noindent\begin{minipage}{\linewidth}
\hypertarget{TERM-1664}{}
\noindent\textbf{strukturne teoremy}\quad{\footnotesize\texttt{TERM-1664}}\par
\noindent\emph{Iztočnik / Source:} Struktursätze\par
\noindent\emph{Rusky / Russian:} структурные теоремы\par
\noindent\emph{Ukrajinsky / Ukrainian:} структурні теореми\par
\smallskip\noindent\emph{Razlog:} Strukturne teoremy zbližaje se s oběma zapisanima vostočnoslovjanskima formama i imenuje rezultaty o strukturě algebraičnyh objektov. Razlog odbacaje prěvod Sätze kako propozicije v toj frazě; to jest izbor registra, ne pravilo, že vsako německo Satz v korpusu mora stati teoremoju. Konkretny nynějšny priměr govori o matričnyh prstenah nad nekomutativnymi tělami i o recipročnyh prědstavjenjah. Pri takoj terminologiji dolžno ostati javno, ktore rezultaty trěbujut konečny rang nad centrom: sama fraza strukturne teoremy ne dava toj hipotezy. Arhiv ne priložil čislene težiny za teoremy protiv propozicij.\par
\noindent\emph{Reason:} Strukturne teoremy closely follows both recorded East Slavic comparanda and names results about the structure of algebraic objects. The rationale rejects propositions for Sätze in this phrase; that is a register choice, not a rule requiring every German Satz in the corpus to become a theorem. A current example concerns matrix rings over noncommutative division rings and reciprocal representations. The terminology must leave clear which results require finite rank over the centre: the phrase structure theorems does not itself supply that hypothesis. No numerical weighting of theorems against propositions was recorded.\par
\noindent{\footnotesize Izvorne zapisy / Original records: HIST-g206-005.}\par
\smallskip\noindent\emph{Sravnjenje s tekstom:} Konečny rang nad centrom vključaje se tamo, kde jego argument trěbuje.\par
\noindent\emph{Textual comparison:} Finite rank over the centre is introduced where the argument requires it.\par
\noindent{\footnotesize TEXT-COMP-721; pasaže / passages: 23381; \texttt{CORPUS\_FORM\_COMPARISONS.json}.}\par
\end{minipage}\par
\par\medskip\Needspace{7\baselineskip}
\noindent\begin{minipage}{\linewidth}
\hypertarget{TERM-1665}{}
\noindent\textbf{substitucija (35)}\quad{\footnotesize\texttt{TERM-1665}}\par
\noindent\emph{Iztočnik / Source:} Substitution (35)\par
\noindent\emph{Rusky / Russian:} подстановка (35)\par
\noindent\emph{Ukrajinsky / Ukrainian:} підстановка (35)\par
\smallskip\noindent\emph{Razlog:} Substitucija tu jest opisana zaměna prěměnnyh, s točnym odkazom na formulu (35). Internacionalny korenj izbran jest pri ruskom podstanovka i ukrajinskom pidstanovka; zapis ne dava sravniteljne ocěny drugih medžuslovjanskyh kandidatov. V pasaži složenje (12) s (35) ponovno dava substituciju tipa (12), i posrědstvom nje nastavajut odgovarjajuče idealy. Zato čislo (35) jest čest funkcionalnogo odkaza, ne dekoracija slovnikovogo zaglavja. Jednočasna nulovost pary idealov jest matematičny poslědok toj konstrukcije, ne definicija slova substitucija.\par
\noindent\emph{Reason:} Here substitucija names the specified change of variables, with an exact reference to formula (35). The international stem was chosen alongside Russian podstanovka and Ukrainian pidstanovka; the record supplies no comparative scoring of other Interslavic candidates. In the passage, composing (12) with (35) again gives a substitution of type (12), from which corresponding ideals are constructed. The number (35) is consequently a functional reference, not decorative wording in a dictionary heading. The simultaneous vanishing of the paired ideals is a mathematical consequence of the construction, not the definition of substitution.\par
\noindent{\footnotesize Izvorne zapisy / Original records: HIST-g090-013, HIST-g092-067.}\par
\smallskip\noindent\emph{Sravnjenje s tekstom:} Odkazy (12) i (35) določajut konkretne zaměny.\par
\noindent\emph{Textual comparison:} References (12) and (35) specify the actual changes of variables.\par
\noindent{\footnotesize TEXT-COMP-722; pasaže / passages: 14961; \texttt{CORPUS\_FORM\_COMPARISONS.json}.}\par
\end{minipage}\par
\par\medskip\Needspace{7\baselineskip}
\noindent\begin{minipage}{\linewidth}
\hypertarget{TERM-1666}{}
\noindent\textbf{substitucija novih rędov}\quad{\footnotesize\texttt{TERM-1666}}\par
\noindent\emph{Iztočnik / Source:} Substitution neuer Reihen\par
\noindent\emph{Rusky / Russian:} подстановка новых строк\par
\noindent\emph{Ukrajinsky / Ukrainian:} підстановка нових рядків\par
\smallskip\noindent\emph{Razlog:} Substitucija novih redov razlikuje zaměnu redov od kontrakcije. Arhivny ręd jest starši pravopis, a nynějšna pasaža piše redov. Red tu znači red simboličnyh veličin v matričnom prědstavjenju, ne beskonečny čislovy red; rusko stroka i ukrajinsko rjadok pomogajut viděti toj smysl. Formuly (9) i (11) uvodet nove redy, poslě čego možno izvesti daljšu kontrakciju. Těsna svez dvoh postupkov ne daje prava sliti jih nazvy. Za to razlikovanje jest točny pasažny razlog, a ne izmysljene čislene rangovanje.\par
\noindent\emph{Reason:} Substitucija novih redov distinguishes replacement of rows from contraction. Archived ręd is an earlier spelling; the present passage writes redov. Here a row is a row of symbolic quantities in a matrix representation, not an infinite numerical series; Russian stroka and Ukrainian rjadok help identify that sense. Formulas (9) and (11) introduce new rows, after which a further contraction may be performed. The close connection between the two operations does not justify merging their names. This distinction has a specific textual basis, not an invented numerical ranking.\par
\noindent{\footnotesize Izvorne zapisy / Original records: HIST-g031-009.}\par
\smallskip\noindent\emph{Sravnjenje s tekstom:} Dvě operacije sustavno sut povezane, ale ne identične.\par
\noindent\emph{Textual comparison:} The two operations are systematically connected but not identical.\par
\noindent{\footnotesize TEXT-COMP-723; pasaže / passages: 3851, 3853, 3854, 3857; \texttt{CORPUS\_FORM\_COMPARISONS.json}.}\par
\end{minipage}\par
\par\medskip\Needspace{7\baselineskip}
\noindent\begin{minipage}{\linewidth}
\hypertarget{TERM-1667}{}
\noindent\textbf{poslědovateljno jednorazovo svijanje}\quad{\footnotesize\texttt{TERM-1667}}\par
\noindent\emph{Iztočnik / Source:} sukzessive einmalige Faltung\par
\noindent\emph{Rusky / Russian:} последовательная однократная свертка; последовательные однократные свертки\par
\noindent\emph{Ukrajinsky / Ukrainian:} послідовне одноразове згортання; послідовні одноразові згортання\par
\smallskip\noindent\emph{Razlog:} Poslědovateljno jednorazovo svijanje opisuje povtorjany postupok: vsaka pojedina operacija jest jednorazova. To ne jest to samo kako jedna kontrakcija viššego reda. Starši zapis izrično ostavil pytanje, bylo li by množstveno jednorazove svijanja prirodnjejše od jedinčnogo nazvanja procesa. Nynějšne pasaže čuvajut konstrukciju s jedinčnym imenikom, ale svijanje zaměnjeno jest terminom kontrakcija. Korpus tako pokazuje realizovanu formu, ne dokazuje, že množstvena alternativa byla jezyčno vyvrgnuta. Razlog izbora ostaje prozračnost reda operacij i kontinuitet terminologije; težiny ili test razumělivosti ne zapisany.\par
\noindent\emph{Reason:} Poslědovateljno jednorazovo svijanje describes a repeated procedure in which each individual operation is onefold. It is not the same as a single contraction of higher order. The older record explicitly leaves open whether plural jednorazove svijanja would be more natural than a singular process noun. Current passages retain the singular construction but replace svijanje with kontrakcija. They show which form was used, not that the plural alternative was linguistically disproved. The recorded reasons are clarity about the sequence of operations and terminological continuity; no weights or intelligibility test were recorded.\par
\noindent{\footnotesize Izvorne zapisy / Original records: HIST-g014-009, HIST-g015-009.}\par
\smallskip\noindent\emph{Sravnjenje s tekstom:} Nynějšny tekst čuva jedinčny procesny imenik pri povtorjanyh operacijah.\par
\noindent\emph{Textual comparison:} The current text retains a singular process noun for repeated operations.\par
\noindent{\footnotesize TEXT-COMP-724; pasaže / passages: 2072, 2178; \texttt{CORPUS\_FORM\_COMPARISONS.json}.}\par
\end{minipage}\par
\par\medskip\Needspace{7\baselineskip}
\noindent\begin{minipage}{\linewidth}
\hypertarget{TERM-1668}{}
\noindent\textbf{sukcesivna eliminacija}\quad{\footnotesize\texttt{TERM-1668}}\par
\noindent\emph{Iztočnik / Source:} sukzessive Elimination\par
\noindent\emph{Rusky / Russian:} последовательное исключение\par
\noindent\emph{Ukrajinsky / Ukrainian:} послідовне елімінування\par
\smallskip\noindent\emph{Razlog:} Sukcesivna eliminacija nazyvaje postupno izključanje prěměnnyh, ne jedno nerazčlenjeno izčislenje rezultanta. V rabotě 22 povezuje se s redom idealov; starši razlog naglašaje prodolženje korenja eliminacija iz ranějšej terminologije. Rusky zapis upotrěbljaje isključenije, ukrajinsky eliminuvannja, tako že sravnjenje ne jest glasovanje dvoh jednakyh form. Nynějšna primětka pripisana E. Noether naglašaje nezavisnost postupka od izbranoj idealnoj bazy. I to svojstvo dolžno ostati v prěvodu: kratko ime postupka go ne vključaje avtomatično.\par
\noindent\emph{Reason:} Sukcesivna eliminacija names elimination of variables in stages, not one undifferentiated resultant calculation. In paper 22 it is connected with a sequence of ideals; the older rationale stresses continuity with the previously used eliminacija stem. Russian uses isključenije and Ukrainian eliminuvannja, so the comparison is not a vote by two identical forms. A current note attributed to E. Noether emphasizes the procedure's independence from the chosen ideal basis. That property must remain in the translation too: the short name of the procedure does not automatically express it.\par
\noindent{\footnotesize Izvorne zapisy / Original records: HIST-g090-002, HIST-g092-056.}\par
\smallskip\noindent\emph{Sravnjenje s tekstom:} Postupnost i nezavisnost od bazy sut različne svojstva.\par
\noindent\emph{Textual comparison:} Successiveness and basis independence are distinct properties.\par
\noindent{\footnotesize TEXT-COMP-725; pasaže / passages: 14932, 14971; \texttt{CORPUS\_FORM\_COMPARISONS.json}.}\par
\end{minipage}\par
\par\medskip\Needspace{7\baselineskip}
\noindent\begin{minipage}{\linewidth}
\hypertarget{TERM-1669}{}
\noindent\textbf{sukcesivna substitucija}\quad{\footnotesize\texttt{TERM-1669}}\par
\noindent\emph{Iztočnik / Source:} sukzessive Substitution\par
\noindent\emph{Rusky / Russian:} последовательная подстановка\par
\noindent\emph{Ukrajinsky / Ukrainian:} послідовна підстановка\par
\smallskip\noindent\emph{Razlog:} Sukcesivna substitucija i sukcesivna specializacija tut opisujut razne aspekty togo samogo postupka. Substitucija govori o zaměně prěměnnoj, specializacija o prěhodu k danym čislovym cěnnostjam. Formula izrično zaměnjaje najprvo $x_n$, potom $x_{n-1}$ i tako dalje. Vsaka racionalna funkcija mora pri tom byti privedena k redukovanomu prědstavjenju; bez togo ogradženja prěvod iztratil by uslovje, ktore izključaje identično nulovy čislitelj, imenovatelj ili izraz $0/0$. Internacionalno sukcesivna jest zapisany izbor, ne nova nazva za inoj matematičnoj operaciji.\par
\noindent\emph{Reason:} Sukcesivna substitucija and successive specialization describe different aspects of the same procedure here. Substitution names replacement of a variable; specialization names passage to specified numerical values. The formula explicitly replaces $x_n$, then $x_{n-1}$, and so on. Each rational function must be put into reduced form at the relevant step; omitting that qualification would remove a condition used to exclude an identically zero numerator or denominator, or an expression $0/0$. International sukcesivna is the recorded wording, not the name of a different mathematical operation.\par
\noindent{\footnotesize Izvorne zapisy / Original records: HIST-g042-016.}\par
\smallskip\noindent\emph{Sravnjenje s tekstom:} Redukovane racionalne izrazy sut čest uslovij.\par
\noindent\emph{Textual comparison:} Reduced rational expressions are part of the assumptions.\par
\noindent{\footnotesize TEXT-COMP-726; pasaže / passages: 5242, 5245, 5246, 5247, 5250; \texttt{CORPUS\_FORM\_COMPARISONS.json}.}\par
\end{minipage}\par
\par\medskip\Needspace{7\baselineskip}
\noindent\begin{minipage}{\linewidth}
\hypertarget{TERM-1670}{}
\noindent\textbf{sukcesivna čislova specializacija}\quad{\footnotesize\texttt{TERM-1670}}\par
\noindent\emph{Iztočnik / Source:} sukzessive zahlenmäßige Spezialisierung\par
\noindent\emph{Rusky / Russian:} последовательная числовая специализация\par
\noindent\emph{Ukrajinsky / Ukrainian:} послідовна числова спеціалізація\par
\smallskip\noindent\emph{Razlog:} Čislova specializacija razjasnjaje, čto prěměnny dostavajut čislove cěnnosti, a sukcesivna, že to děje se jedna za drugoju. Arhiv motivuje izbor imenno toju strukturoju dokaza, s ruskoju i ukrajinskoju paralelami. Lema ne govori, že libovoljne cěnnosti sut dozvoljene: ostajuče funkcije trěba shraniti algebraično nezavisnymi. Pod tym prědpoloženjem zavisna funkcija ne stava identično nuloju ili neoprěděljenoju. Zato specializovati ne smije v slovniku stati bezuslovnym zamenjati čislami. Nijedna ocěnka v procenteh ne byla zapisana.\par
\noindent\emph{Reason:} Čislova specializacija makes explicit that variables receive numerical values, while sukcesivna indicates that this happens one variable after another. The archive motivates the choice by this structure of the proof, with Russian and Ukrainian parallels. The lemma does not permit arbitrary values: the remaining functions must retain their algebraic independence. Under that hypothesis a dependent function does not become identically zero or indeterminate. Specialize must therefore not be glossed as unconditional replacement by numbers. No percentage score was recorded.\par
\noindent{\footnotesize Izvorne zapisy / Original records: HIST-g042-003.}\par
\smallskip\noindent\emph{Sravnjenje s tekstom:} Ne vse čislove cěnnosti sut dozvoljene.\par
\noindent\emph{Textual comparison:} Not every numerical value is admissible.\par
\noindent{\footnotesize TEXT-COMP-727; pasaže / passages: 5156; \texttt{CORPUS\_FORM\_COMPARISONS.json}.}\par
\end{minipage}\par
\par\medskip\Needspace{7\baselineskip}
\noindent\begin{minipage}{\linewidth}
\hypertarget{TERM-1671}{}
\noindent\textbf{sukcesivno razloženje}\quad{\footnotesize\texttt{TERM-1671}}\par
\noindent\emph{Iztočnik / Source:} sukzessive Zerlegung\par
\noindent\emph{Rusky / Russian:} последовательное разложение\par
\noindent\emph{Ukrajinsky / Ukrainian:} послідовне розкладання\par
\smallskip\noindent\emph{Razlog:} Sukcesivno razloženje znači konstruovanje razloženja korak po koraku. V rabotě 19 redukovano prědstavjenje ideala nastava iz redov dvočlennyh prědstavjenij; dokaz prověrjaje, že prědstavjenje s většim čislom komponentov ostaje redukovano. Ne govori se samo o egzistenciji gotovogo razloženja. Slovjanske razloženije i rozkladannja podkrepljajut imenik postupka, ale ne opreděljajut, kakym matematičnym spojanjem komponenty sostavljajut cělo. Tu trěba čuvati avtorsko najmenše obče mnogokratno idealov, ne zaměnjati go bez objasnjenja obyčnoju sumoju.\par
\noindent\emph{Reason:} Sukcesivno razloženje means constructing a decomposition step by step. In paper 19 a reduced representation of an ideal is obtained from a sequence of two-component representations; the proof checks that the representation with more components remains reduced. This is more than the existence of an already completed decomposition. Russian razloženije and Ukrainian rozkladannja support a process noun, but do not specify the mathematical operation combining the components. Here the author's least common multiple of ideals must be retained, not silently replaced by an ordinary sum.\par
\noindent{\footnotesize Izvorne zapisy / Original records: HIST-g081-033, HIST-g082-033.}\par
\smallskip\noindent\emph{Sravnjenje s tekstom:} Nastavajuče prědstavjenje dolžno ostati redukovano.\par
\noindent\emph{Textual comparison:} The resulting representation must remain reduced.\par
\noindent{\footnotesize TEXT-COMP-728; pasaže / passages: 12574, 12588, 12590, 12591, 12592, 12594, 12598; \texttt{CORPUS\_FORM\_COMPARISONS.json}.}\par
\end{minipage}\par
\par\medskip\Needspace{7\baselineskip}
\noindent\begin{minipage}{\linewidth}
\hypertarget{TERM-1672}{}
\noindent\textbf{sukcesivny razklad v dv? sostavne ?esti}\quad{\footnotesize\texttt{TERM-1672}}\par
\noindent\emph{Iztočnik / Source:} sukzessive Zerlegung in zwei Bestandteile\par
\noindent\emph{Rusky / Russian:} [Bukvy ne sut zachovane / Lettering not recoverable]\par
\noindent\emph{Ukrajinsky / Ukrainian:} [Bukvy ne sut zachovane / Lettering not recoverable]\par
\smallskip\noindent\emph{Razlog:} Starši izraz imaje nečitlive bukvy v dv? i ?esti; nynějšna pasaža nezavisno dava razklad v dvě sostavne česti. Dvě jest suštna čest konstrukcije: mnogočlenno prědstavjenje dobyvaje se povtorjanymi dvočlennymi razkladami. V tom razsuždenju avtor takože dopušća beskonečno prodolženje, ako nijedna iz pojavjajučih se čestij ne jest nula. Zato množstveno sostavne česti ne dolžno prikryvati uslovje nenulovosti. Poškodžene rusko i ukrajinsko pole ostajut poškodžene; čitlivy nynějšny tekst ne vosstanovjaje ih starše bukvy.\par
\noindent\emph{Reason:} The older phrase has unreadable letters in dv? and ?esti; the current passage independently supplies razklad v dvě sostavne česti. Two is an essential part of the construction: a many-component representation is obtained through repeated two-component decompositions. The discussion also allows indefinite continuation if none of the resulting components is zero. Thus the phrase constituent parts must not obscure the nonvanishing condition. The damaged Russian and Ukrainian fields remain damaged; the readable current translation does not restore their earlier lettering.\par
\noindent{\footnotesize Izvorne zapisy / Original records: HIST-g079-029.}\par
\smallskip\noindent\emph{Sravnjenje s tekstom:} Bez konečnogo ograničenja ne znači bez uslovja o komponentah.\par
\noindent\emph{Textual comparison:} Indefinite continuation does not remove the condition on the components.\par
\noindent{\footnotesize TEXT-COMP-729; pasaže / passages: 11267, 11269, 11272; \texttt{CORPUS\_FORM\_COMPARISONS.json}.}\par
\end{minipage}\par
\par\medskip\Needspace{7\baselineskip}
\noindent\begin{minipage}{\linewidth}
\hypertarget{TERM-1673}{}
\noindent\textbf{predstavjenje jediničnogo elementa kako sumy}\quad{\footnotesize\texttt{TERM-1673}}\par
\noindent\emph{Iztočnik / Source:} Summendarstellung des Einheitselementes\par
\noindent\emph{Rusky / Russian:} представление единичного элемента в виде суммы\par
\noindent\emph{Ukrajinsky / Ukrainian:} подання одиничного елемента як суми\par
\smallskip\noindent\emph{Razlog:} Predstavjenje jediničnogo elementa kako sumy jest razložena, razumna fraza vměsto kalkovanja cělogo německogo složenogo imenika. Ona zachovuje objekt, jedinicu prstena, i operaciju, sumu. Zapisany razlog govori o idempotentnom razloženju: formula dodavaje, že različne komponenty imajut nulovy produkt, a vsaka komponenta jest idempotentna. To razlikuje konstrukciju od proizvoljnogo zapisa jedinice kako sumy. Nynějšny prěvod piše prědstavjenje; razlika pravopisa ne znači razliku ponjatja. Bližšie ruska i ukrajinska opisna frazy pokazyvajut model izraženja, ne cifrovo dokazanu razumělivost.\par
\noindent\emph{Reason:} Predstavjenje jediničnogo elementa kako sumy is an explicit phrase rather than a word-for-word reproduction of the entire German compound. It retains both the object, the ring identity, and the operation, summation. The recorded rationale identifies an idempotent decomposition: the formula additionally says that distinct components have zero product and each component is idempotent. This distinguishes the construction from an arbitrary expression of the identity as a sum. The current translation writes prědstavjenje; that spelling difference does not change the concept. The comparable Russian and Ukrainian phrases support the mode of expression, not a numerically established intelligibility claim.\par
\noindent{\footnotesize Izvorne zapisy / Original records: HIST-g122-004.}\par
\smallskip\noindent\emph{Sravnjenje s tekstom:} Suma jest dopolnjena nulovymi vzajemnymi produktami i idempotentnostju.\par
\noindent\emph{Textual comparison:} The sum is accompanied by zero cross-products and idempotence.\par
\noindent{\footnotesize TEXT-COMP-730; pasaže / passages: 17501, 17503, 17505, 17507; \texttt{CORPUS\_FORM\_COMPARISONS.json}.}\par
\end{minipage}\par
\par\medskip\Needspace{7\baselineskip}
\noindent\begin{minipage}{\linewidth}
\hypertarget{TERM-1674}{}
\noindent\textbf{suplementne čisla}\quad{\footnotesize\texttt{TERM-1674}}\par
\noindent\emph{Iztočnik / Source:} Supplementzahlen\par
\noindent\emph{Rusky / Russian:} суплементные числа\par
\noindent\emph{Ukrajinsky / Ukrainian:} суплементні числа\par
\smallskip\noindent\emph{Razlog:} Suplementne čisla čuvajut istorijno tehnično ime, povezano s O. Ore i izčislenjem eksponentov diferenty. Arhiv izrično predpočita blizost k iztočniku prěd neobjasnjenym zaměnjenjem suvremennym terminom. Ruska i ukrajinska forma takože imajut internacionalny korenj. Nynějšna pasaža imenuje te čisla, ale ne dava tamo samostojnu definiciju; iz samogo zaglavja ne možno izvesti, s ktorym pozdnějšim invariantom sut točno identične. Zato slovnik čuva istorijno ime i njegovu funkciju v pasusu, bez izmysljenogo suvremennogo sinonima.\par
\noindent\emph{Reason:} Suplementne čisla retains the historical technical name associated with O. Ore and the calculation of different exponents. The archive explicitly prefers proximity to the source over an unexplained replacement by a modern term. The recorded Russian and Ukrainian forms likewise use the international stem. The current passage names these numbers but does not supply a self-contained definition there; the heading alone does not establish an exact identification with a later invariant. The dictionary therefore preserves the historical name and its role in the passage without inventing a modern synonym.\par
\noindent{\footnotesize Izvorne zapisy / Original records: HIST-g213-015.}\par
\smallskip\noindent\emph{Sravnjenje s tekstom:} Ime ostaje istorijno; ta pasaža ne dava polnu definiciju.\par
\noindent\emph{Textual comparison:} The name remains historical; this passage does not give a full definition.\par
\noindent{\footnotesize TEXT-COMP-731; pasaže / passages: 25169; \texttt{CORPUS\_FORM\_COMPARISONS.json}.}\par
\end{minipage}\par
\par\medskip\Needspace{7\baselineskip}
\noindent\begin{minipage}{\linewidth}
\hypertarget{TERM-1675}{}
\noindent\textbf{Sylowovo polje}\quad{\footnotesize\texttt{TERM-1675}}\par
\noindent\emph{Iztočnik / Source:} Sylowkörper\par
\noindent\emph{Rusky / Russian:} силовское поле\par
\noindent\emph{Ukrajinsky / Ukrainian:} силовське поле\par
\smallskip\noindent\emph{Razlog:} Sylowovo polje jest historijna pisana forma; nynějšna pasaža ima Silovovyh polj i dava izričnu definiciju: invariantno polje Silovovoj $p$-podgrupy Galoisovoj grupy. Tako polje ne jest samo podgrupa, ani libovoljno polje stepena $p$. Razlika Sylow / Silov dotykaje prilagodženje vlastnogo imena, ne raznyh matematičnyh objektov. V dokazě stepen toga polja nad osnovnym poljem jest vzajemno prosty s $p$; to umožnjaje argument o indeksu algebry. Kratky arhivny razlog govori o shranjenju grupno-teoretičnoj redukcije, ale ne opravdava nynějšnu formu imena oficijalnym slovnikom.\par
\noindent\emph{Reason:} Sylowovo polje is the historical recorded spelling; the current passage has Silovovyh polj and explicitly defines the object as the fixed field of a Sylow $p$-subgroup of the Galois group. Such a field is neither the subgroup itself nor an arbitrary field of degree $p$. Sylow versus Silov concerns adaptation of the proper name, not different mathematical objects. In the proof the field's degree over the ground field is prime to $p$, enabling the argument about the algebra's index. The short archival rationale invokes preservation of this group-theoretic reduction, but does not establish the name's spelling from an official dictionary.\par
\noindent{\footnotesize Izvorne zapisy / Original records: HIST-g208-009.}\par
\smallskip\noindent\emph{Sravnjenje s tekstom:} Pasaža sama dava definiciju i nužny odnos stepena k prostomu čislu.\par
\noindent\emph{Textual comparison:} The passage itself gives the definition and required relation of the degree to the prime.\par
\noindent{\footnotesize TEXT-COMP-732; pasaže / passages: 23031; \texttt{CORPUS\_FORM\_COMPARISONS.json}.}\par
\end{minipage}\par
\par\medskip\Needspace{7\baselineskip}
\noindent\begin{minipage}{\linewidth}
\hypertarget{TERM-1677}{}
\noindent\textbf{simbolična predstavimost}\quad{\footnotesize\texttt{TERM-1677}}\par
\noindent\emph{Iztočnik / Source:} symbolische Darstellbarkeit\par
\noindent\emph{Rusky / Russian:} символическая представимость\par
\noindent\emph{Ukrajinsky / Ukrainian:} символічна зображуваність; символічна поданність\par
\smallskip\noindent\emph{Razlog:} Simbolična predstavimost znači možnost prědstaviti objekt simbolično, ne pojedino prědstavjenje. Ta razlika medžu svojstvom i konkretnym vyrazom jest obči razlog v trěh arhivnyh zapisah; ukrajinske alternativy zobrazhuvanist i podannist ne sut tuta čislovno rangovane. Nynějšny tekst piše prědstavimost i formuluje teoremu, že invariantne tvorjenja možno prědstaviti matričnymi produktami. Pozdnějše razsuždenje dodavaje eksplicitnu prědstavimost. Zato vozmožnost i podany eksplicitny vyraz ne trěba zaměšati; isto tako tut ne ide o teoriju matričnyh prědstavjenij grup.\par
\noindent\emph{Reason:} Simbolična predstavimost means the possibility of representing an object symbolically, not an individual representation. This distinction between a property and a particular expression is the common rationale in three archival records; the Ukrainian alternatives zobrazhuvanist and podannist are not numerically ranked there. The current text writes prědstavimost and states the theorem that invariant formations can be represented by matrix products. A later argument adds explicit representability. Possibility and an exhibited explicit expression must therefore not be conflated; nor is this passage about the theory of matrix representations of groups.\par
\noindent{\footnotesize Izvorne zapisy / Original records: HIST-g029-006, HIST-g031-003, HIST-g035-002.}\par
\smallskip\noindent\emph{Sravnjenje s tekstom:} Teorema o možnosti ne jest samo pokaz konkretnogo izraza.\par
\noindent\emph{Textual comparison:} A theorem about possibility is not merely the display of a particular expression.\par
\noindent{\footnotesize TEXT-COMP-734; pasaže / passages: 3807, 4359; \texttt{CORPUS\_FORM\_COMPARISONS.json}.}\par
\end{minipage}\par
\par\medskip\Needspace{7\baselineskip}
\noindent\begin{minipage}{\linewidth}
\hypertarget{TERM-1678}{}
\noindent\textbf{simbolične identičnosti}\quad{\footnotesize\texttt{TERM-1678}}\par
\noindent\emph{Iztočnik / Source:} symbolische Identitäten\par
\noindent\emph{Rusky / Russian:} символические тождества\par
\noindent\emph{Ukrajinsky / Ukrainian:} символічні тотожності\par
\smallskip\noindent\emph{Razlog:} Simbolične identičnosti nazyvajut algebraične ravnosti v simboličnom metodu. Smysl německogo Identität tut ne jest filosofska totožnost objekta ani jedinična matrica. Rusko toždestvo i ukrajinsko totožnist objasnjajut smysl, no ne dokazyvajut, že internacionalna identičnost jest jedina možna medžuslovjanska forma. Arhiv sam ostavil jezyčnu ocěnku otkrytu. Nynějšny paragraf izrično govori o generovanju relacij iz identičnostij pri pravilah o redah simbolov i prěměnnyh; ty pravila dolžno čuvati vměstě s nazvanjem, ne shoditi cělu teoremu do slova ravnost.\par
\noindent\emph{Reason:} Simbolične identičnosti names algebraic identities in the symbolic method. German Identität here means neither philosophical identity of an object nor an identity matrix. Russian toždestvo and Ukrainian totožnist clarify the sense, but do not prove that international identičnost is the only possible Interslavic form. The archive itself leaves the linguistic assessment open. The current section explicitly discusses generating relations from identities under rules governing rows of symbols and variables; those rules must be retained alongside the name, not collapsed into the word equality.\par
\noindent{\footnotesize Izvorne zapisy / Original records: HIST-g032-001.}\par
\smallskip\noindent\emph{Sravnjenje s tekstom:} Pravila o redah simbolov i prěměnnyh ostajut čest tvrdženja.\par
\noindent\emph{Textual comparison:} Rules for rows of symbols and variables remain part of the statement.\par
\noindent{\footnotesize TEXT-COMP-735; pasaže / passages: 3895, 3897, 3901, 3902, 3903; \texttt{CORPUS\_FORM\_COMPARISONS.json}.}\par
\end{minipage}\par
\par\medskip\Needspace{7\baselineskip}
\noindent\begin{minipage}{\linewidth}
\hypertarget{TERM-1679}{}
\noindent\textbf{simboličny metod}\quad{\footnotesize\texttt{TERM-1679}}\par
\noindent\emph{Iztočnik / Source:} symbolische Methode\par
\noindent\emph{Rusky / Russian:} символический метод\par
\noindent\emph{Ukrajinsky / Ukrainian:} символічний метод\par
\smallskip\noindent\emph{Razlog:} Simboličny metod jest nazva metoda teorije invariantov, ne obči opis vsakoj matematiky s simbolami. Arhiv daje blizke rusku i ukrajinsku formu i opiraje se na internacionalnu disciplinarnu terminologiju, bez citovane tada konkretnoj slovnikovej normy. Pasaža naznačaje dvě fundamentalne teoremy: simbolično prědstavjenje tvorjenij i dobyvanje relacij posrědstvom malogo čisla identičnostij. Tymy funkcijami metod dostava točny obseg. Točno upominanje Studyjeve raboty jest naučny odkaz avtora; ne možno prěpisati jego konsultaciju prevoditeljskomu AI bez osobnogo dokaza.\par
\noindent\emph{Reason:} Simboličny metod names the method of invariant theory, not mathematics using symbols in general. The archive supplies closely related Russian and Ukrainian forms and invokes international disciplinary terminology without citing a particular dictionary authority consulted at the time. The passage identifies two fundamental theorems: symbolic representation of formations and generation of relations by means of a small number of identities. Those functions determine the method's scope. Its precise reference to Study's work is the author's scholarly citation; consultation of that work cannot be attributed to the translating AI without separate evidence.\par
\noindent{\footnotesize Izvorne zapisy / Original records: HIST-g028-004.}\par
\smallskip\noindent\emph{Sravnjenje s tekstom:} Avtorsky naučny odkaz ne dokazyvaje konsultaciju od AI.\par
\noindent\emph{Textual comparison:} The author's scholarly citation does not establish consultation by the AI.\par
\noindent{\footnotesize TEXT-COMP-736; pasaže / passages: 3649; \texttt{CORPUS\_FORM\_COMPARISONS.json}.}\par
\end{minipage}\par
\par\medskip\Needspace{7\baselineskip}
\noindent\begin{minipage}{\linewidth}
\hypertarget{TERM-1680}{}
\noindent\textbf{simboličny ręd}\quad{\footnotesize\texttt{TERM-1680}}\par
\noindent\emph{Iztočnik / Source:} Symbolreihe\par
\noindent\emph{Rusky / Russian:} символическая строка\par
\noindent\emph{Ukrajinsky / Ukrainian:} символічний рядок\par
\smallskip\noindent\emph{Razlog:} Simboličny red označaje red simbolov, a ne čislovy red ili tekstovu stroku. Starši zapis ręd i nynějšny redov razlikujut se pravopisom i padežom. Zaključna ograda v rabotě 4 jest osoblivo važna: redy prěměnnyh i redy, izvedene substitucijeju iz simboličnyh redov, ne imajut avtomatično te iste svojstva pri redukciji. Vyražajnost jednogo izraza prěz drugy zato ješče ne znači ekvivalentnost v avtorskom smyslu. Razlog izbora ne sostojit samo v sličnosti s ruskim stroka ili ukrajinskim rjadok; on takože čuva razliku, potrěbnu za pravilno čitanje dokaza.\par
\noindent\emph{Reason:} Simboličny red denotes a row of symbols, not a numerical series or a line of text. The older spelling ręd and current redov differ in spelling and case. The closing qualification in paper 4 is particularly important: rows of variables and rows obtained by substitution from symbolic rows do not automatically have the same reduction properties. Expressibility of one expression in terms of another therefore does not yet establish equivalence in the author's sense. The reason for the choice is not merely resemblance to Russian stroka or Ukrainian rjadok; it also preserves a distinction needed to read the proof correctly.\par
\noindent{\footnotesize Izvorne zapisy / Original records: HIST-g037-008.}\par
\smallskip\noindent\emph{Sravnjenje s tekstom:} Rozděl medžu redami prěměnnyh i simboličnymi redami jest potrěbny.\par
\noindent\emph{Textual comparison:} The distinction between variable rows and symbolic rows is necessary.\par
\noindent{\footnotesize TEXT-COMP-737; pasaže / passages: 4872; \texttt{CORPUS\_FORM\_COMPARISONS.json}.}\par
\end{minipage}\par
\par\medskip\Needspace{7\baselineskip}
\noindent\begin{minipage}{\linewidth}
\hypertarget{TERM-1681}{}
\noindent\textbf{rędy simbolov}\quad{\footnotesize\texttt{TERM-1681}}\par
\noindent\emph{Iztočnik / Source:} Symbolreihen\par
\noindent\emph{Rusky / Russian:} строки символов\par
\noindent\emph{Ukrajinsky / Ukrainian:} символьні рядки\par
\smallskip\noindent\emph{Razlog:} Redy simbolov jest množstvena forma za Symbolreihen. Starši ręd byl označeny kako provisionalny, a nynějšny tekst piše redy; ni jedno z toho ne jest dokaz samostojnej jezyčnoj certifikacije. V definiciji redy simbolov izpolnjajut relacije analogične relacijam determinantov, a jih elementy noset indekse. Zato red ne treba čitati kako beskonečnu sumu ili kakykoli uporedženy spis. Rusko stroki simvolov i ukrajinsko symvolni rjadky pomogajut razlikovati smysl. Vybrany opisny genitiv simbolov čuva, iz čego se red sostoji.\par
\noindent\emph{Reason:} Redy simbolov is the plural expression for Symbolreihen. Earlier ręd was marked provisional; the present text writes redy, neither fact amounting to independent linguistic certification. In the definition the symbolic rows satisfy relations analogous to determinant relations and their elements carry indices. Row must therefore not be read as an infinite sum or an arbitrary ordered list. Russian stroki simvolov and Ukrainian symvolni rjadky help disambiguate the sense. The chosen descriptive genitive simbolov retains what the row consists of.\par
\noindent{\footnotesize Izvorne zapisy / Original records: HIST-g030-008.}\par
\smallskip\noindent\emph{Sravnjenje s tekstom:} Red ima konkretnu algebraičnu strukturu.\par
\noindent\emph{Textual comparison:} The row has a specified algebraic structure.\par
\noindent{\footnotesize TEXT-COMP-738; pasaže / passages: 3762, 3764, 3766, 3769; \texttt{CORPUS\_FORM\_COMPARISONS.json}.}\par
\end{minipage}\par
\par\medskip\Needspace{7\baselineskip}
\noindent\begin{minipage}{\linewidth}
\hypertarget{TERM-1682}{}
\noindent\textbf{rędy simbolov / rędy prěměnnyh}\quad{\footnotesize\texttt{TERM-1682}}\par
\noindent\emph{Iztočnik / Source:} Symbolreihen / Variabelnreihen\par
\noindent\emph{Rusky / Russian:} строки символов / строки переменных\par
\noindent\emph{Ukrajinsky / Ukrainian:} символьні рядки / рядки змінних\par
\smallskip\noindent\emph{Razlog:} Para redy simbolov / redy prěměnnyh čuva jednu obšču glavu redy, ale ne zjednačuje dva objekty. V dokazu o identičnostjah red prěměnnyh može byti specializovany, a čislo redov simbolov vključeno jest v pravilo o nerazložimosti identičnosti. Oba prilastka zato nesut matematičnu informaciju. Ruske i ukrajinske pary sledujut tu samu opisnu strukturu. Starši znak ę prědstavjaje pravopisnu razliku; poznamka o recenziji v arhivu označuje otvorjeno jezyčno pytanje, ne zabranjenje čitati ili dalje razvivati dokaz.\par
\noindent\emph{Reason:} The pair redy simbolov / redy prěměnnyh keeps the shared head noun rows without identifying the two objects. In the argument about identities a variable row can be specialized, while the number of symbolic rows enters the rule for treating an identity as indecomposable. Both modifiers therefore carry mathematical information. The Russian and Ukrainian pairs follow the same descriptive structure. The older ę is a spelling difference; the archive's review note identifies an open linguistic question, not a prohibition on reading or developing the argument.\par
\noindent{\footnotesize Izvorne zapisy / Original records: HIST-g032-004.}\par
\smallskip\noindent\emph{Sravnjenje s tekstom:} Različne funkcii redov ne sut izbrisane obščim imenikom.\par
\noindent\emph{Textual comparison:} The shared noun does not erase the different roles of the rows.\par
\noindent{\footnotesize TEXT-COMP-739; pasaže / passages: 3897; \texttt{CORPUS\_FORM\_COMPARISONS.json}.}\par
\end{minipage}\par
\par\medskip\Needspace{7\baselineskip}
\noindent\begin{minipage}{\linewidth}
\hypertarget{TERM-1683}{}
\noindent\textbf{elementarne simetrične funkcije}\quad{\footnotesize\texttt{TERM-1683}}\par
\noindent\emph{Iztočnik / Source:} symmetrische Elementarfunktionen\par
\noindent\emph{Rusky / Russian:} элементарные симметрические функции\par
\noindent\emph{Ukrajinsky / Ukrainian:} елементарні симетричні функції\par
\smallskip\noindent\emph{Razlog:} Elementarne simetrične funkcije jest ustojiva internacionalna fraza v trěh zapisanih prěvodah. Elementarne tu znači oprěděljenu klasu simetričnyh funkcij, ne proste ili legke funkcije v pedagogičnom smyslu. V korpusu one vystupajut takože za redy veličin; opreděljenje mora sledovati konkretnu konstrukciju, a ne avtomatično prěnesti vsako svojstvo iz najprostějšego odnoprěměnnogo priměra. Starši zapisy opravdavajut povtorjanje frazy konsistentnostju medžu paragrafami. Oni ne davajut čislene težiny, i tri pojavjenja ne sut tri nezavisne jezyčne autoritety.\par
\noindent\emph{Reason:} Elementarne simetrične funkcije is a stable international phrase across the three recorded translations. Elementary here denotes a specified class of symmetric functions, not simple or easy functions in a pedagogical sense. In the corpus they also occur for rows of quantities; the definition must follow the construction in question rather than automatically importing every property from a simpler example. The historical records justify repeated wording by consistency between sections. They give no numerical weights, and three occurrences are not three independent linguistic authorities.\par
\noindent{\footnotesize Izvorne zapisy / Original records: HIST-g045-007, HIST-g048-013, HIST-g053-003.}\par
\smallskip\noindent\emph{Sravnjenje s tekstom:} Elementarne jest tehnično ograničenje, ne ocěnka legkosti.\par
\noindent\emph{Textual comparison:} Elementary is a technical qualification, not a judgment of ease.\par
\noindent{\footnotesize TEXT-COMP-740; pasaže / passages: 5595, 5605, 5916; \texttt{CORPUS\_FORM\_COMPARISONS.json}.}\par
\end{minipage}\par
\par\medskip\Needspace{7\baselineskip}
\noindent\begin{minipage}{\linewidth}
\hypertarget{TERM-1684}{}
\noindent\textbf{simetrične funkcije redov veličin}\quad{\footnotesize\texttt{TERM-1684}}\par
\noindent\emph{Iztočnik / Source:} symmetrische Funktionen von Größenreihen\par
\noindent\emph{Rusky / Russian:} симметрические функции рядов величин\par
\noindent\emph{Ukrajinsky / Ukrainian:} симетричні функції рядів величин\par
\smallskip\noindent\emph{Razlog:} Simetrične funkcije redov veličin uvodet simetriju medžu redami, prikazanymi posrědstvom $y,z,\ldots,t$. Redy tut ne sut stepenove redy. Arhiv ostavil za sravnjenje red i riad, ale ne zapisal razultat čislennogo ocěnjanja. Nynějšna lema govori o cělyh cěločislenyh funkcijah i o konečnom čislu porodžajučih funkcij; oba ogradženja dolžno čuvati. Slovnik može pokazati, kako vybrana fraza razčlenjaje objekt, bez neosnovanogo tvrdženja, že vse narodne tradicije imajut istu terminologiju za mnogo-redny slučaj.\par
\noindent\emph{Reason:} Simetrične funkcije redov veličin introduces symmetry among the rows displayed using $y,z,\ldots,t$. These rows are not power series. The archive leaves red and riad for comparison but records no numerical assessment. The present lemma concerns integral functions with integer coefficients and a finite number of generators; both qualifications must remain. The dictionary can show how the chosen phrase identifies its object without claiming that every national tradition uses identical terminology for the many-row case.\par
\noindent{\footnotesize Izvorne zapisy / Original records: HIST-g054-010.}\par
\smallskip\noindent\emph{Sravnjenje s tekstom:} Cělost i cěločislenost sut ohranjene v lemi.\par
\noindent\emph{Textual comparison:} Integrality and integer coefficients are retained in the lemma.\par
\noindent{\footnotesize TEXT-COMP-741; pasaže / passages: 6460, 6463, 6464, 6466, 6470, 6486; \texttt{CORPUS\_FORM\_COMPARISONS.json}.}\par
\end{minipage}\par
\par\medskip\Needspace{7\baselineskip}
\noindent\begin{minipage}{\linewidth}
\hypertarget{TERM-1685}{}
\noindent\textbf{sistem $s_I$}\quad{\footnotesize\texttt{TERM-1685}}\par
\noindent\emph{Iztočnik / Source:} System $s_I$\par
\noindent\emph{Rusky / Russian:} система $s_I$\par
\noindent\emph{Ukrajinsky / Ukrainian:} система $s_I$\par
\smallskip\noindent\emph{Razlog:} V sistem $s_I$ samo imenik jest jezyčny izbor; $s$ i rimski podindeks $I$ sut matematične identifikatory. Arhiv povezuje sistem s kontrakcijami $s$ i njegovyh stepenjev s formami Sistema I. Pasaža iz §18 konkretno vypisyvaje te vhody, a zaključny pregled broji 35 form. Taj broj jest avtorska struktura raboty, ne broj slovnikovih zapisov. Pri kiriličnom prěpise ne smije se zaměniti latinsko $s$ ili $I$ vizualno podobnoju bukvoju. Izbor staršego muškogo sistem ne izključa slovnikovu žensku formu sistema, razmotrenu oddělno.\par
\noindent\emph{Reason:} In sistem $s_I$, only the noun is a language choice; $s$ and the Roman subscript $I$ are mathematical identifiers. The archive connects the system with contractions of $s$ and its powers with forms from System I. Section 18 explicitly lists those inputs, and the closing synopsis counts 35 forms. That number belongs to the author's organization of the work, not to the dictionary's entry count. A Cyrillic rendering must not replace Latin $s$ or $I$ with a visually similar letter. Historical masculine sistem does not exclude dictionary feminine sistema, discussed separately.\par
\noindent{\footnotesize Izvorne zapisy / Original records: HIST-g019-001.}\par
\smallskip\noindent\emph{Sravnjenje s tekstom:} Zaglavje i zaključny broj pripadajut istomu podsistemu.\par
\noindent\emph{Textual comparison:} The heading and closing count belong to the same subsystem.\par
\noindent{\footnotesize TEXT-COMP-742; pasaže / passages: 2530, 2534, 2535, 2536, 3525; \texttt{CORPUS\_FORM\_COMPARISONS.json}.}\par
\end{minipage}\par
\par\medskip\Needspace{7\baselineskip}
\noindent\begin{minipage}{\linewidth}
\hypertarget{TERM-1686}{}
\noindent\textbf{sistem $s_{II}$}\quad{\footnotesize\texttt{TERM-1686}}\par
\noindent\emph{Iztočnik / Source:} System $s_{II}$\par
\noindent\emph{Rusky / Russian:} система $s_{II}$\par
\noindent\emph{Ukrajinsky / Ukrainian:} система $s_{II}$\par
\smallskip\noindent\emph{Razlog:} Sistem $s_{II}$ prodolžaje konvenciju podsistemov, pri čem podindeks $II$ ne prevodi se. Tekst §19 nazyvaje kontrakciju $s$ s vypisanymi formami, a zaključny pregled navodi devet form togo podsistema. Jednakost obščego imenika sistem s poprědnim zaglavjem ne označuje jednakost dvou podsistemov. Starši razlog točno naglašaje zaštitu rimskogo podindeksa pri kiriličnom prěpise. To jest pravilo věrnogo prěnosa notacije, ne novy termin, dobity jezyčnym glasovanjem.\par
\noindent\emph{Reason:} Sistem $s_{II}$ continues the subsystem convention, leaving the subscript $II$ untranslated. Section 19 names the contraction of $s$ with the listed forms, and the closing synopsis gives nine forms for the subsystem. Sharing the noun sistem with the preceding heading does not identify the two subsystems. The older rationale specifically calls for preservation of the Roman subscript in Cyrillic. This is a rule for faithful transmission of notation, not a new term obtained by a linguistic vote.\par
\noindent{\footnotesize Izvorne zapisy / Original records: HIST-g020-001.}\par
\smallskip\noindent\emph{Sravnjenje s tekstom:} Kontrakcije i broj form opreděljajut konkretny podsistem.\par
\noindent\emph{Textual comparison:} Contractions and the count of forms identify the particular subsystem.\par
\noindent{\footnotesize TEXT-COMP-743; pasaže / passages: 2650, 2653, 3526; \texttt{CORPUS\_FORM\_COMPARISONS.json}.}\par
\end{minipage}\par
\par\medskip\Needspace{7\baselineskip}
\noindent\begin{minipage}{\linewidth}
\hypertarget{TERM-1687}{}
\noindent\textbf{sistem $s_{IV}$}\quad{\footnotesize\texttt{TERM-1687}}\par
\noindent\emph{Iztočnik / Source:} System $s_{IV}$\par
\noindent\emph{Rusky / Russian:} система $s_{IV}$\par
\noindent\emph{Ukrajinsky / Ukrainian:} система $s_{IV}$\par
\smallskip\noindent\emph{Razlog:} Sistem $s_{IV}$ razčlenjaje se na lokalizovany obči imenik i nezměnjene matematične ime. V zaključnom spisu podsistem imaje 32 formy, zato izguba rimskogo podindeksa uničila by vez s konkretnym mestom v konstrukciji. Arhiv ne navodi alternativnu nazvu ni čislenu ocěnku; njegov razlog jest čuvanje stabilnogo označenja. Slične bukvy iz drugogo pisma ne sut ravnocěnna zaměna za izvorne matematične znaky. Pravopis i rod imenika možno razgledati oddělno od togo zaštitenogo imena.\par
\noindent\emph{Reason:} Sistem $s_{IV}$ separates a localized common noun from an unchanged mathematical name. The closing list assigns the subsystem 32 forms, so losing the Roman subscript would break the connection with its particular place in the construction. The archive gives neither an alternative name nor a numerical score; its rationale is preservation of a stable designation. Similar-looking letters from another script are not interchangeable with the original mathematical symbols. The noun's spelling and gender can be considered separately from the protected name.\par
\noindent{\footnotesize Izvorne zapisy / Original records: HIST-g027-002.}\par
\smallskip\noindent\emph{Sravnjenje s tekstom:} Označenje jest čest matematičnogo spiska.\par
\noindent\emph{Textual comparison:} The designation is part of the mathematical inventory.\par
\noindent{\footnotesize TEXT-COMP-744; pasaže / passages: 3528; \texttt{CORPUS\_FORM\_COMPARISONS.json}.}\par
\end{minipage}\par
\par\medskip\Needspace{7\baselineskip}
\noindent\begin{minipage}{\linewidth}
\hypertarget{TERM-1688}{}
\noindent\textbf{sistem klasov ostatkov / komplement}\quad{\footnotesize\texttt{TERM-1688}}\par
\noindent\emph{Iztočnik / Source:} System der Restklassen / Komplement\par
\noindent\emph{Rusky / Russian:} система классов вычетов / дополнение\par
\noindent\emph{Ukrajinsky / Ukrainian:} система класів лишків / комплемент\par
\smallskip\noindent\emph{Razlog:} Sistem klasov ostatkov i komplement ne sut sinonimy. Slovnikova kosa črta združuje dva povezana ponjatja iz jednogo razsuždenja: prěhod k klasam ostatkov povezuje razloženje idealov s aditivnym razloženjem, a svojstvo komplementa služi drugomu opisu razložimosti. Komplement tu ne znači prosto dopolnjenje množiny. Rusky zapis imaje dopolnenije, ukrajinsky komplement; medžuslovjanska forma jest internacionalny izbor, ne dokaz jedinoglasnosti. Potrěbno jest obě ponjatja čitati v algebraičnom svezu, kotory arhivny razlog imenno hotěl shraniti.\par
\noindent\emph{Reason:} Sistem klasov ostatkov and komplement are not synonyms. The glossary's slash groups two connected concepts from one argument: passage to residue classes relates ideal decomposition to additive decomposition, while the complement property gives another description of decomposability. Complement here does not simply mean set-theoretic complement. Russian records dopolnenije and Ukrainian komplement; the Interslavic form is an international choice, not evidence of unanimity. Both concepts need the algebraic connection that the archival rationale sought to preserve.\par
\noindent{\footnotesize Izvorne zapisy / Original records: HIST-g081-016, HIST-g082-016.}\par
\smallskip\noindent\emph{Sravnjenje s tekstom:} Povezane ponjatja ne sut sinonimy.\par
\noindent\emph{Textual comparison:} The related concepts are not synonyms.\par
\noindent{\footnotesize TEXT-COMP-745; pasaže / passages: 12448; \texttt{CORPUS\_FORM\_COMPARISONS.json}.}\par
\end{minipage}\par
\par\medskip\Needspace{7\baselineskip}
\noindent\begin{minipage}{\linewidth}
\hypertarget{TERM-1690}{}
\noindent\textbf{sistem koeficientov Galoisovoj rezolventy}\quad{\footnotesize\texttt{TERM-1690}}\par
\noindent\emph{Iztočnik / Source:} System von Koeffizienten der Galoischen Resolvente\par
\noindent\emph{Rusky / Russian:} система коэффициентов резольвенты Galois\par
\noindent\emph{Ukrajinsky / Ukrainian:} система коефіцієнтів резольвенти Galois\par
\smallskip\noindent\emph{Razlog:} Sistem koeficientov Galoisovoj rezolventy mora iměti oba genitivne odnosa: čto tvori sistem i čijimi koeficientami ty veličiny sut. Izpustiti koeficientov pretvorilo by frazu v drugo tvrdženje. Nynějšna pasaža nazyvaje upravo toj sistem integralnoju bazoju za invarianty konečnyh grup. Starši zapis čuva razlog, že ranějša fraza byla prěkratka; važny naučny rezultat jest polny odnos koeficientov k rezolventě, ne povědanje o redakciji. Vlastno ime Galois ostaje prepoznavajemy v prilastku, ne zaměnjaje se obščim slovom.\par
\noindent\emph{Reason:} Sistem koeficientov Galoisovoj rezolventy must preserve both genitive relations: what constitutes the system and what those quantities are coefficients of. Omitting koeficientov would turn the phrase into a different assertion. The current passage names precisely this coefficient system as an integral basis for invariants of finite groups. The older record explains that an earlier phrase was overcompressed; the scholarly point is the complete coefficient–resolvent relation, not the editing history. The proper name Galois remains recognizable in the modifier rather than becoming a generic word.\par
\noindent{\footnotesize Izvorne zapisy / Original records: HIST-g094-020.}\par
\smallskip\noindent\emph{Sravnjenje s tekstom:} Oba genitivne odnosa sut nužne.\par
\noindent\emph{Textual comparison:} Both genitive relations are required.\par
\noindent{\footnotesize TEXT-COMP-747; pasaže / passages: 15164; \texttt{CORPUS\_FORM\_COMPARISONS.json}.}\par
\end{minipage}\par
\par\medskip\Needspace{7\baselineskip}
\noindent\begin{minipage}{\linewidth}
\hypertarget{TERM-1691}{}
\noindent\textbf{sistemna determinanta}\quad{\footnotesize\texttt{TERM-1691}}\par
\noindent\emph{Iztočnik / Source:} Systemdeterminante\par
\noindent\emph{Rusky / Russian:} системный определитель\par
\noindent\emph{Ukrajinsky / Ukrainian:} системний визначник\par
\smallskip\noindent\emph{Razlog:} Sistemna determinanta jest determinant konkretnoj sistemnoj matrice, ne proizvoljno imenujemy determinant. V definiciji elementy matrice sut linearne formy v neoprěděljenyh, a stepen determinanta odpovědaje stepenu prědstavjenja. Prilastok sistemna svazuje termin s toj konstrukcijeju. Rusky opredelitel i ukrajinsky vyznačnyk pokazyvajut različne narodne imeniki; medžuslovjanska determinanta izbira internacionalny korenj. Arhivny opis, že objekt ne jest samo obyčny determinant, treba čitati kako ograničenje konteksta, ne kako izmysljenje nove operacije determinantovanja.\par
\noindent\emph{Reason:} Sistemna determinanta is the determinant of the specified system matrix, not an arbitrarily named determinant. In the definition the matrix entries are linear forms in indeterminates, and the determinant's degree matches the representation's degree. Sistemna connects the term with that construction. Russian opredelitel and Ukrainian vyznačnyk illustrate different national nouns; Interslavic determinanta chooses the international stem. The archival claim that this is more than an ordinary determinant is best read as a contextual restriction, not as the invention of a different determinant operation.\par
\noindent{\footnotesize Izvorne zapisy / Original records: HIST-g201-002.}\par
\smallskip\noindent\emph{Sravnjenje s tekstom:} Sistemna ukazuje specifikovanu matricu; operacija ostaje determinant.\par
\noindent\emph{Textual comparison:} System identifies the specified matrix; the operation remains a determinant.\par
\noindent{\footnotesize TEXT-COMP-748; pasaže / passages: 22098, 22100, 22102, 22104; \texttt{CORPUS\_FORM\_COMPARISONS.json}.}\par
\end{minipage}\par
\par\medskip\Needspace{7\baselineskip}
\noindent\begin{minipage}{\linewidth}
\hypertarget{TERM-1692}{}
\noindent\textbf{sistemy $\mathfrak{S}_{n\rho}$ / $\mathfrak{S}^{*}_{\rho\rho}$}\quad{\footnotesize\texttt{TERM-1692}}\par
\noindent\emph{Iztočnik / Source:} Systeme $\mathfrak{S}_{n\rho}$ / $\mathfrak{S}^{*}_{\rho\rho}$\par
\noindent\emph{Rusky / Russian:} системы $\mathfrak{S}_{n\rho}$ / $\mathfrak{S}^{*}_{\rho\rho}$\par
\noindent\emph{Ukrajinsky / Ukrainian:} системи $\mathfrak{S}_{n\rho}$ / $\mathfrak{S}^{*}_{\rho\rho}$\par
\smallskip\noindent\emph{Razlog:} Označenja $\mathfrak{S}_{n\rho}$ i $\mathfrak{S}^{*}_{\rho\rho}$ razlikujut početny sistem i njegov obraz po specializaciji. Zvězdička i oba podindeksy sut nužne: formula ne opisuje libovoljnu zaměnu imena. Nynějšny tekst govori o vzajemno jednoznačnom odobraženju i identifikuje algebraičny rang $\rho$. Arhiv zapisuje TeX makro Ssys, ktorogo definicija jest frakturna bukva S; čitateljska forma pokazyvaje matematičny znak, ne ime makra. Samo množstveno sistemy jest jezyčno lokalizovano.\par
\noindent\emph{Reason:} The designations $\mathfrak{S}_{n\rho}$ and $\mathfrak{S}^{*}_{\rho\rho}$ distinguish the initial system from its image after specialization. The star and both subscripts matter: the formula is not an arbitrary renaming. The current text describes a one-to-one correspondence and identifies the algebraic rank $\rho$. The archive records the TeX macro Ssys, defined as a Fraktur S; the reading form displays the mathematical symbol, not the macro name. Only the plural noun sistemy is localized.\par
\noindent{\footnotesize Izvorne zapisy / Original records: HIST-g043-002.}\par
\smallskip\noindent\emph{Sravnjenje s tekstom:} Zvězdička i indekse razlikujut objekte.\par
\noindent\emph{Textual comparison:} The star and subscripts distinguish the objects.\par
\noindent{\footnotesize TEXT-COMP-749; pasaže / passages: 5052, 5276, 5278; \texttt{CORPUS\_FORM\_COMPARISONS.json}.}\par
\end{minipage}\par
\par\medskip\Needspace{7\baselineskip}
\noindent\begin{minipage}{\linewidth}
\hypertarget{TERM-1693}{}
\noindent\textbf{sistemy linearnyh diferencialnyh uravnenj v častnyh derivatah}\quad{\footnotesize\texttt{TERM-1693}}\par
\noindent\emph{Iztočnik / Source:} Systeme linearer partieller Differentialgleichungen\par
\noindent\emph{Rusky / Russian:} системы линейных уравнений в частных производных\par
\noindent\emph{Ukrajinsky / Ukrainian:} системи лінійних диференціальних рівнянь у частинних похідних\par
\smallskip\noindent\emph{Razlog:} Sistemy linearnyh diferencialnyh uravnenj v častnyh derivatah rastvarjaje skompresovanu německu frazu v jasny opis. Linearne modifikuje uravnenja, a v častnyh derivatah razlikuje ih od obyčnyh diferencialnyh uravnenj. Izrazi i uravnenja ipak ne sut isto: pasaža polaga diferencialne izrazy ravnymi nuli, aby dobiti uravnenja. To jest nužna svez, ktoru slovnik ne smije prećutěti. Arhiv imenuje derivata provisionalnym terminom za ekspertnu ocěnku; novy pasažny pregled objasnjaje njegovu funkciju, ne izmysljaje už izvršenu ekspertnu recenziju.\par
\noindent\emph{Reason:} Sistemy linearnyh diferencialnyh uravnenj v častnyh derivatah expands the compressed German phrase into an explicit description. Linear modifies the equations, while v častnyh derivatah distinguishes partial from ordinary differential equations. Expressions and equations are nevertheless different: the passage sets differential expressions equal to zero to obtain equations. That necessary connection must not disappear from the dictionary. The archive identifies derivata as provisional wording for specialist assessment; this fresh passage examination explains its function without inventing a completed expert review.\par
\noindent{\footnotesize Izvorne zapisy / Original records: HIST-g079-065.}\par
\smallskip\noindent\emph{Sravnjenje s tekstom:} Položenje izraza ravnym nuli jest izrično navedena operacija.\par
\noindent\emph{Textual comparison:} Setting an expression equal to zero is an explicitly stated operation.\par
\noindent{\footnotesize TEXT-COMP-750; pasaže / passages: 11978, 11986; \texttt{CORPUS\_FORM\_COMPARISONS.json}.}\par
\end{minipage}\par
\par\medskip\Needspace{7\baselineskip}
\noindent\begin{minipage}{\linewidth}
\hypertarget{TERM-1694}{}
\noindent\textbf{sistemy parcialnyh ili obyčnyh linearnyh diferencialnyh izrazov}\quad{\footnotesize\texttt{TERM-1694}}\par
\noindent\emph{Iztočnik / Source:} Systeme partieller oder gewöhnlicher linearer Differentialausdrücke\par
\noindent\emph{Rusky / Russian:} системы частных или обыкновенных линейных дифференциальных выражений\par
\noindent\emph{Ukrajinsky / Ukrainian:} системи частинних або звичайних лінійних диференціальних виразів\par
\smallskip\noindent\emph{Razlog:} Sistemy parcialnyh ili obyčnyh linearnyh diferencialnyh izrazov čuvajut obě vetvi avtorskoj formulacije: parcialne i obyčne. Izrazov jest naměrno, ne nebrěžna zaměna za uravnenj. V kontekstu jde o razloženje idealov v nekomutativnoj oblasti i o čislo komponentov, s priměnjenjem k diferencialnym izrazam. Zato kratko sistemy diferencialnyh uravnenj ne bylo by polnym prěvodom toj frazy. Arhiv pri različnyh narodnyh formah izbira internacionalne parcialny i diferencialny; pravopis ne dokazuje oficijalno prijetje tehničnogo slovnika.\par
\noindent\emph{Reason:} Sistemy parcialnyh ili obyčnyh linearnyh diferencialnyh izrazov preserves both branches of the author's formulation: partial and ordinary. Izrazov is intentional, not a careless substitute for equations. The context is decomposition of ideals in a noncommutative domain and the number of components, with an application to differential expressions. Simply writing systems of differential equations would not translate the whole phrase. Alongside differing national forms the archive chooses international parcialny and diferencialny; the spelling does not establish official adoption of this technical vocabulary.\par
\noindent{\footnotesize Izvorne zapisy / Original records: HIST-g081-093.}\par
\smallskip\noindent\emph{Sravnjenje s tekstom:} Obě vetvi frazy i imenik izrazy ostajut čitlive.\par
\noindent\emph{Textual comparison:} Both branches of the phrase and the noun expressions remain explicit.\par
\noindent{\footnotesize TEXT-COMP-751; pasaže / passages: 13666; \texttt{CORPUS\_FORM\_COMPARISONS.json}.}\par
\end{minipage}\par
\par\medskip\Needspace{7\baselineskip}
\noindent\begin{minipage}{\linewidth}
\hypertarget{TERM-1695}{}
\noindent\textbf{sistemy $\mathfrak{S}_{n\rho}$ [rekonstrukcija indeksa / reconstructed subscript]}\quad{\footnotesize\texttt{TERM-1695}}\par
\noindent\emph{Iztočnik / Source:} Systeme $\mathfrak{S}_{n\rho}$ [rekonstrukcija indeksa / reconstructed subscript]\par
\noindent\emph{Rusky / Russian:} системы $\mathfrak{S}_{n\rho}$ [rekonstrukcija indeksa / reconstructed subscript]\par
\noindent\emph{Ukrajinsky / Ukrainian:} системи $\mathfrak{S}_{n\rho}$ [rekonstrukcija indeksa / reconstructed subscript]\par
\smallskip\noindent\emph{Razlog:} Zaglavje iz ranějšego spiska izgledaje kako S\_\{n ho\}, ale polny zachovany zapis ima indeks $n\rho$, a definicija makra v nynějšnom tekstu dava $\mathfrak{S}_{n\rho}$. Zato čitateljska rekonstrukcija možna jest bez nagađanja nepoznatoj bukvy; ona ipak dolžna byti izrično označena. $n$ jest čislo neoprěděljenyh, a $\rho$ algebraičny rang, ne prosto drugy redny numer. Imenik sistemy obščo imenuje skupinu funkcij; ne može se bez prědpoloženij zaměniti poljami. Oštečena skračena grafija i polny arhivny zapis ostajut oba dohledljive.\par
\noindent\emph{Reason:} The earlier heading appears as S\_\{n ho\}, but the complete preserved record has subscript $n\rho$, and the macro definition in the current text supplies $\mathfrak{S}_{n\rho}$. A reading reconstruction is therefore possible without guessing an unknown letter, although it must be identified explicitly. Here $n$ is the number of indeterminates and $\rho$ the algebraic rank, not merely another serial number. Sistemy names the collection of functions generally; it cannot be replaced by fields without the requisite hypotheses. Both the damaged abbreviated form and the complete archival record remain traceable.\par
\noindent{\footnotesize Izvorne zapisy / Original records: HIST-g041-002.}\par
\smallskip\noindent\emph{Sravnjenje s tekstom:} Polna grafija makra i kontekst podkrepljajut označenu rekonstrukciju.\par
\noindent\emph{Textual comparison:} The full macro spelling and context support the labelled reconstruction.\par
\noindent{\footnotesize TEXT-COMP-752; pasaže / passages: 5052, 5057, 5059, 5278; \texttt{CORPUS\_FORM\_COMPARISONS.json}.}\par
\end{minipage}\par
\par\medskip\Needspace{7\baselineskip}
\noindent\begin{minipage}{\linewidth}
\hypertarget{TERM-1696}{}
\noindent\textbf{sistemy $\mathfrak{S}_{n1}$}\quad{\footnotesize\texttt{TERM-1696}}\par
\noindent\emph{Iztočnik / Source:} Systeme $\mathfrak{S}_{n1}$\par
\noindent\emph{Rusky / Russian:} системы $\mathfrak{S}_{n1}$\par
\noindent\emph{Ukrajinsky / Ukrainian:} системи $\mathfrak{S}_{n1}$\par
\smallskip\noindent\emph{Razlog:} Sistemy $\mathfrak{S}_{n1}$ označajut sistemy polinomov algebraičnogo ranga jedan. Tekstova fraza o jedinom algebraično nezavisnom polinomu ne znači, že sistem sodrživaje samo jeden polinom: odtud izrično slěduje tvrdženje takože za sistemy s beskonečno mnogymi polinomami jednoj neoprěděljenej. Dolžno čuvati i imenik i indeks, aby ne zaměšati čislo elementov s rangom. Arhivny razlog naglašaje imeno to objasnjenje. Notacija ostaje matematična pri oběh pismennyh verzijah prěvoda.\par
\noindent\emph{Reason:} Sistemy $\mathfrak{S}_{n1}$ denotes polynomial systems of algebraic rank one. The wording about a single algebraically independent polynomial does not mean that the system contains just one polynomial: the text explicitly includes systems with infinitely many polynomials in one indeterminate. Both the noun and the subscript must be preserved to avoid confusing cardinality with rank. The archival rationale emphasizes exactly this clarification. The notation remains mathematical in both scripts of the translation.\par
\noindent{\footnotesize Izvorne zapisy / Original records: HIST-g053-005.}\par
\smallskip\noindent\emph{Sravnjenje s tekstom:} Takože beskonečne sistemy vključene sut v tvrdženje.\par
\noindent\emph{Textual comparison:} Infinite systems are also included in the assertion.\par
\noindent{\footnotesize TEXT-COMP-753; pasaže / passages: 6387, 6400; \texttt{CORPUS\_FORM\_COMPARISONS.json}.}\par
\end{minipage}\par
\par\medskip\Needspace{7\baselineskip}
\noindent\begin{minipage}{\linewidth}
\hypertarget{TERM-1697}{}
\noindent\textbf{sistemna matrica / grupova matrica}\quad{\footnotesize\texttt{TERM-1697}}\par
\noindent\emph{Iztočnik / Source:} Systemmatrix / Gruppenmatrix\par
\noindent\emph{Rusky / Russian:} системная матрица / групповая матрица\par
\noindent\emph{Ukrajinsky / Ukrainian:} системна матриця / групова матриця\par
\smallskip\noindent\emph{Razlog:} Sistemna matrica i grupova matrica sut povezane, ale ne oběma směrama ravnoznačne nazvy. Pasaža definije $W=A_1x_1+\cdots+A_hx_h$ iz danogo prědstavjenja, a grupovoj matriceju nazyvaje specialny slučaj, kogda osnovne elementy tvoret grupu. Prilastok grupova zato ne možno odstraniti kako nepotrěbny sinonim. Internacionalna matrica daje kontinuitet s drugimi algebraičnymi zapisami. Kosá črta v arhivu krati sravnjenje, ne sostavljaje matematičnu identičnost dvou ponjatij.\par
\noindent\emph{Reason:} Sistemna matrica and grupova matrica are related names, but not interchangeable in both directions. The passage defines $W=A_1x_1+\cdots+A_hx_h$ from the given representation and calls it a group matrix in the special case where the underlying elements form a group. Grupova therefore cannot be discarded as a redundant synonym. International matrica provides continuity with the other algebraic terminology. The archive's slash abbreviates the comparison; it does not assert identity of the two concepts.\par
\noindent{\footnotesize Izvorne zapisy / Original records: HIST-g201-003.}\par
\smallskip\noindent\emph{Sravnjenje s tekstom:} Definicija ne dava obratnu ekvivalentnost dvou nazv.\par
\noindent\emph{Textual comparison:} The definition does not make the two names interchangeable in both directions.\par
\noindent{\footnotesize TEXT-COMP-754; pasaže / passages: 22094, 22096, 22098, 22100; \texttt{CORPUS\_FORM\_COMPARISONS.json}.}\par
\end{minipage}\par
\par\medskip\Needspace{7\baselineskip}
\noindent\begin{minipage}{\linewidth}
\hypertarget{TERM-1698}{}
\noindent\textbf{syzygije}\quad{\footnotesize\texttt{TERM-1698}}\par
\noindent\emph{Iztočnik / Source:} Syzygien\par
\noindent\emph{Rusky / Russian:} сизигии\par
\noindent\emph{Ukrajinsky / Ukrainian:} сизигії\par
\smallskip\noindent\emph{Razlog:} Syzygije čuvaje internacionalno ime za relacije, ktore v toj pasaži vystupajut medžu produktami. Arhiv predpočita zajęty termin opisnej frazě radi menšej dvojsmyslnosti, no ne dava dokaz, že medžuslovjanska forma už byla oficijalno normovana. Ruske i ukrajinske formy imajut drugu prilagodženost pisanija. Nynějšny tekst sam objasnjaje smysl posrědstvom skobok; taka kratka glosa pomogaje čitatelju bez zaměny termina. Ne slěduje iz toga, že vsaka relacija v cělom korpusu dolžna nazyvati se syzygijeju.\par
\noindent\emph{Reason:} Syzygije preserves the international name for relations, here relations between products. The archive prefers the borrowing to a paraphrase on grounds of reduced ambiguity, but supplies no evidence that the Interslavic form was already officially standardized. The Russian and Ukrainian forms adapt the spelling differently. The current text explains the meaning parenthetically; such a short gloss helps readers without replacing the term. It does not follow that every relation throughout the corpus should be called a syzygy.\par
\noindent{\footnotesize Izvorne zapisy / Original records: HIST-g027-007.}\par
\smallskip\noindent\emph{Sravnjenje s tekstom:} Kratko objasnjenje v skobkah razjasnjaje tehničny termin.\par
\noindent\emph{Textual comparison:} A short parenthetical gloss clarifies the technical term.\par
\noindent{\footnotesize TEXT-COMP-755; pasaže / passages: 3416; \texttt{CORPUS\_FORM\_COMPARISONS.json}.}\par
\end{minipage}\par
\par\medskip\Needspace{7\baselineskip}
\noindent\begin{minipage}{\linewidth}
\hypertarget{TERM-1700}{}
\noindent\textbf{tabela}\quad{\footnotesize\texttt{TERM-1700}}\par
\noindent\emph{Iztočnik / Source:} Tabelle\par
\noindent\emph{Rusky / Russian:} таблица\par
\noindent\emph{Ukrajinsky / Ukrainian:} таблиця\par
\smallskip\noindent\emph{Razlog:} Tabela jest obči imenik za organizovany prikaz form. Arhiv dopušča lokalizaciju nadpisov, no ne zaměnu matematičnyh těl tabely; v zaključnom prikazu 331 form i podsistemov čisla, indekse i razměščenje imajut naučnu funkciju. Rusko i ukrajinsko tablica ne sut istovětna s vybranym tabela, zato izraz ne opisujemo kako edinu obšču slovjansku formu. Važno jest razlikovati prěvod nazvy i redakciju tabelarnogo sodržanja. Opravdanje termina samo po sobě ne ustanovjaje věrnosti vseh jeho bunek.\par
\noindent\emph{Reason:} Tabela is the common noun for the organized display of forms. The archive permits localization of captions but not replacement of the mathematical table bodies; in the concluding presentation of 331 forms and subsystems, numbers, indices, and arrangement have scholarly functions. Russian and Ukrainian tablica differ from the chosen tabela, so it should not be described as the single common Slavic form. Translating the name and editing the table's contents are distinct operations. Justifying the term does not by itself establish the fidelity of every cell.\par
\noindent{\footnotesize Izvorne zapisy / Original records: HIST-g027-012.}\par
\smallskip\noindent\emph{Sravnjenje s tekstom:} Ime tabely ne zaměnjaje njeje polno matematično sodržanje.\par
\noindent\emph{Textual comparison:} The table's name does not replace its complete mathematical contents.\par
\noindent{\footnotesize TEXT-COMP-757; pasaže / passages: 3505, 3512, 3525, 3526, 3528; \texttt{CORPUS\_FORM\_COMPARISONS.json}.}\par
\end{minipage}\par
\par\medskip\Needspace{7\baselineskip}
\noindent\begin{minipage}{\linewidth}
\hypertarget{TERM-1702}{}
\noindent\textbf{dělimy nerazložimym uravnenjem}\quad{\footnotesize\texttt{TERM-1702}}\par
\noindent\emph{Iztočnik / Source:} teilbar durch eine irreduzible Gleichung\par
\noindent\emph{Rusky / Russian:} делится на неприводимое уравнение\par
\noindent\emph{Ukrajinsky / Ukrainian:} ділиться на незвідне рівняння\par
\smallskip\noindent\emph{Razlog:} Dělimy nerazložimym uravnenjem čuvaje historijno govorenje o dělimosti uravnenjem. Matematično děli se polinomny izraz, ne istinnost tvrdženja: pasaža izrično piše $H=FG$ i ukazuje polje, nad ktorym jest $F$ nerazložimy. Zato uravnenje ne trěba čitati kako proizvoljnu logičnu formulu. Arhiv predpočita věrnost avtorskemu izrazu modernizaciji na polinom; slovnik može objasniti smysl bez tichej zaměny v diplomatskom prěvodu. Imenovano polje jest neizpustnym uslovjem nerazložimosti.\par
\noindent\emph{Reason:} Dělimy nerazložimym uravnenjem preserves the historical language of divisibility by an equation. Mathematically, the divisible object is a polynomial expression, not the truth of a proposition: the passage explicitly writes $H=FG$ and specifies the field over which $F$ is irreducible. Uravnenje must therefore not be read as an arbitrary logical formula. The record favours fidelity to the author's expression over modernization to polynomial; the dictionary can explain that usage without silently replacing it in a diplomatic translation. The named field is essential to irreducibility.\par
\noindent{\footnotesize Izvorne zapisy / Original records: HIST-g071-009.}\par
\smallskip\noindent\emph{Sravnjenje s tekstom:} Formula objasnjaje historijno govorenje o uravnenju.\par
\noindent\emph{Textual comparison:} The formula explains the historical equation terminology.\par
\noindent{\footnotesize TEXT-COMP-759; pasaže / passages: 8558, 8560; \texttt{CORPUS\_FORM\_COMPARISONS.json}.}\par
\end{minipage}\par
\par\medskip\Needspace{7\baselineskip}
\noindent\begin{minipage}{\linewidth}
\hypertarget{TERM-1703}{}
\noindent\textbf{dělimost / dělimy čerez modul}\quad{\footnotesize\texttt{TERM-1703}}\par
\noindent\emph{Iztočnik / Source:} Teilbarkeit / durch einen Modul teilbar\par
\noindent\emph{Rusky / Russian:} делимость / делимый на модуль\par
\noindent\emph{Ukrajinsky / Ukrainian:} подільність / подільний на модуль\par
\smallskip\noindent\emph{Razlog:} Dělimost označaje uslovje medžu modulami, vyraženo kongruencijami: vsaky element prvogo modula mora byti nuleju po drugom modulu. Arhivna fraza dělimy čerez modul čuvaje tu relaciju; nynějšna definicija piše dělimy prěz. Razlika predloga ne daje pravo obrnuti směr relacije. Zaměna modul s ideal byla by osoblivo riskovna v jednostrannom kontekstu. Sravnjenje s obyčnoju dělimostju polinomov pomogaje intuiciji, ale ne zaměnjaje danu definiciju.\par
\noindent\emph{Reason:} Dělimost denotes a relation between modules expressed through congruences: every element of the first module must be zero modulo the second. The recorded phrase dělimy čerez modul retains that relation; the current definition uses dělimy prěz. This prepositional difference does not license reversing its direction. Replacing module with ideal would be especially risky in the one-sided setting. Ordinary polynomial divisibility provides an analogy, not a substitute for the stated definition.\par
\noindent{\footnotesize Izvorne zapisy / Original records: HIST-g079-019.}\par
\smallskip\noindent\emph{Sravnjenje s tekstom:} Definicije čuvajut relaciju medžu modulami.\par
\noindent\emph{Textual comparison:} The definitions preserve a relation between modules.\par
\noindent{\footnotesize TEXT-COMP-760; pasaže / passages: 11157, 11159, 11161, 11163; \texttt{CORPUS\_FORM\_COMPARISONS.json}.}\par
\end{minipage}\par
\par\medskip\Needspace{7\baselineskip}
\noindent\begin{minipage}{\linewidth}
\hypertarget{TERM-1704}{}
\noindent\textbf{dělimost $D_T$ črez $p$}\quad{\footnotesize\texttt{TERM-1704}}\par
\noindent\emph{Iztočnik / Source:} Teilbarkeit von $D_T$ durch $p$\par
\noindent\emph{Rusky / Russian:} делимость $D_T$ на $p$\par
\noindent\emph{Ukrajinsky / Ukrainian:} подільність $D_T$ на $p$\par
\smallskip\noindent\emph{Razlog:} Dělimost $D_T$ črez $p$ jest imenska fraza za to, že $p$ děli diskriminant. Ona sohranjaje pored dělitelja i dělimogo, ktory ne možno obrnuti pri sravnjenju s ruskym na ili ukrajinskym na. V danom teoremu uslovje ima dve možnosti: istinno primarnu komponentu ili prosty ideal vtorogo vida. Kratko objasnjenje samo kako razvětvenje moglo by zakryti vtoru možnost. Arhivna latinica $T$ jest kratky zapis; v pasaži ona prědstavjena osobym matematičnym pismom, ne drugym objektom.\par
\noindent\emph{Reason:} Dělimost $D_T$ črez $p$ is the nominal expression for $p$ dividing the discriminant. It preserves which object divides which; the Russian and Ukrainian use of na must not lead to reversal. The theorem gives two possibilities: a genuinely primary component or a prime ideal of the second kind. Explaining the condition simply as ramification could obscure the second alternative. The archive's plain $T$ is shorthand for the specially styled mathematical symbol in the passage, not a different object.\par
\noindent{\footnotesize Izvorne zapisy / Original records: HIST-g154-007.}\par
\smallskip\noindent\emph{Sravnjenje s tekstom:} Teorem ne trěba skracati do jednogo specialnogo slučaja.\par
\noindent\emph{Textual comparison:} The theorem must not be reduced to one special case.\par
\noindent{\footnotesize TEXT-COMP-761; pasaže / passages: 18715; \texttt{CORPUS\_FORM\_COMPARISONS.json}.}\par
\end{minipage}\par
\par\medskip\Needspace{7\baselineskip}
\noindent\begin{minipage}{\linewidth}
\hypertarget{TERM-1705}{}
\noindent\textbf{pojmy děliteljnosti}\quad{\footnotesize\texttt{TERM-1705}}\par
\noindent\emph{Iztočnik / Source:} Teilbarkeitsbegriffe\par
\noindent\emph{Rusky / Russian:} понятия делимости\par
\noindent\emph{Ukrajinsky / Ukrainian:} поняття подільності\par
\smallskip\noindent\emph{Razlog:} Pojmy děliteljnosti jest starši izbor za pojmy dělimosti. Děliteljnost i dělimost ne naznačajut tut dvě različne teorije: razlika jest v tvorjenju imenika. Arhiv nastojaje na dělimosti v zadanoj oblasti, kde obči dělitelj i oba količnika morajut ostati v toj oblasti. Zato analogija s faktorovanjem v polinomnom kolcu ne jest dostatočna. Slovnik čuvaje obě formy za nahodženje pasaž, no ne dava izmyšljenu čislenu prednost jednoj z njih; za konsistentny terminologičny ukazatelj vhodny jest obči naslov dělimost.\par
\noindent\emph{Reason:} Pojmy děliteljnosti is an earlier wording for concepts of divisibility. Děliteljnost and dělimost do not name two theories here; they are competing noun formations. The record insists on divisibility in the specified domain, where the common divisor and both quotients must remain in that domain. Factorization in a polynomial ring is therefore not an adequate substitute. Both forms are retained for finding passages, without inventing a numerical preference; dělimost is a suitable shared index heading.\par
\noindent{\footnotesize Izvorne zapisy / Original records: HIST-g048-003.}\par
\smallskip\noindent\emph{Sravnjenje s tekstom:} Pripadanje dělitelja i količnikov oblasti jest čest definicije.\par
\noindent\emph{Textual comparison:} Membership of the divisor and quotients in the domain is part of the definition.\par
\noindent{\footnotesize TEXT-COMP-762; pasaže / passages: 5857, 5859, 5861, 5863, 5865; \texttt{CORPUS\_FORM\_COMPARISONS.json}.}\par
\end{minipage}\par
\par\medskip\Needspace{7\baselineskip}
\noindent\begin{minipage}{\linewidth}
\hypertarget{TERM-1706}{}
\noindent\textbf{dělitelj}\quad{\footnotesize\texttt{TERM-1706}}\par
\noindent\emph{Iztočnik / Source:} Teiler\par
\noindent\emph{Rusky / Russian:} делитель\par
\noindent\emph{Ukrajinsky / Ukrainian:} дільник\par
\smallskip\noindent\emph{Razlog:} Dělitelj věrno prěnosi Teiler, no jeho smysl zavisi od objekta. V devetom članku oblast sodržaje drugu oblast kako dělitelj; ne govori se o čislovom faktorě. Arhiv zadržaje historijno ime vměsto svobodnoj modernizacije. Čitatelju zato trěba objasniti relaciju sodržanja, ne ostaviti samo ednoslovno sravnjenje. Izbor se razlikuje od podpolje za Teiler des Körpers v drugom članku: to jest razlika redakčnej politiki, ne dokaz, že dělitelj i podpolje sut obče sinonimy.\par
\noindent\emph{Reason:} Dělitelj closely renders Teiler, but its meaning depends on the object. In the ninth paper one domain contains another as a divisor; this is not a numerical factor. The record retains the historical name rather than freely modernizing it. Readers therefore need the inclusion relation explained, not merely a one-word equivalence. This differs from the choice podpolje for Teiler des Körpers elsewhere: it is a difference of editorial policy, not evidence that divisor and subfield are general synonyms.\par
\noindent{\footnotesize Izvorne zapisy / Original records: HIST-g059-004, HIST-g060-012.}\par
\smallskip\noindent\emph{Sravnjenje s tekstom:} Dělitelj tu ne označaje čislovy faktor.\par
\noindent\emph{Textual comparison:} Divisor does not mean a numerical factor here.\par
\noindent{\footnotesize TEXT-COMP-763; pasaže / passages: 7199, 7232; \texttt{CORPUS\_FORM\_COMPARISONS.json}.}\par
\end{minipage}\par
\par\medskip\Needspace{7\baselineskip}
\noindent\begin{minipage}{\linewidth}
\hypertarget{TERM-1707}{}
\noindent\textbf{podpolje}\quad{\footnotesize\texttt{TERM-1707}}\par
\noindent\emph{Iztočnik / Source:} Teiler des Körpers\par
\noindent\emph{Rusky / Russian:} подполе\par
\noindent\emph{Ukrajinsky / Ukrainian:} підполе\par
\smallskip\noindent\emph{Razlog:} Podpolje jest smyslovy, ne doslovny prěvod Teiler des Körpers. Pasaža razgledaje vložene polja i jih prěsěčenje; značenje ne jest količnično polje niti faktor polinoma. Moderno ime zato odstraňaje dvojsmyslnost staršego Teiler. Rusko podpole i ukrajinsko pidpole davajut prěhlednu sravniteljnu osnovu. Věrnosť smyslu mora byti oddělena od věrnosti historijnomu slovu: drugde zachovany dělitelj trěbuje istogo objasnjenja, a ne avtomatičnej zaměny po cělom korpusu.\par
\noindent\emph{Reason:} Podpolje translates the sense, rather than the literal wording, of Teiler des Körpers. The passage concerns embedded fields and their intersection, not a quotient field or a polynomial factor. The modern name therefore removes an ambiguity of the older Teiler. Russian podpole and Ukrainian pidpole provide a transparent comparison. Fidelity to the concept must be distinguished from fidelity to historical wording: retained dělitelj elsewhere needs the same explanation, not an automatic corpus-wide replacement.\par
\noindent{\footnotesize Izvorne zapisy / Original records: HIST-g073-012.}\par
\smallskip\noindent\emph{Sravnjenje s tekstom:} Sodržanje i prěsěčenje poljev objasnjajut vybrany termin.\par
\noindent\emph{Textual comparison:} Field inclusion and intersection explain the selected term.\par
\noindent{\footnotesize TEXT-COMP-764; pasaže / passages: 8763; \texttt{CORPUS\_FORM\_COMPARISONS.json}.}\par
\end{minipage}\par
\par\medskip\Needspace{7\baselineskip}
\noindent\begin{minipage}{\linewidth}
\hypertarget{TERM-1708}{}
\noindent\textbf{dělitelj jednakogo ranga}\quad{\footnotesize\texttt{TERM-1708}}\par
\noindent\emph{Iztočnik / Source:} Teiler gleichen Ranges\par
\noindent\emph{Rusky / Russian:} делитель того же ранга\par
\noindent\emph{Ukrajinsky / Ukrainian:} дільник того самого рангу\par
\smallskip\noindent\emph{Razlog:} Dělitelj jednakogo ranga zadržaje dve česti uslovja: dělitelj v avtorskome poredku modulov i jednakost ranga. Definicija osnovnogo modula ukazuje kongruenciju $bg\equiv0$ s nenulovym $b$; ne govori, že proizvoljne moduly jednakogo ranga sut ravne. Slovo jednakogo samo ograničaje dopustime dělitelje. Arhiv odbacaje svobodnu zaměnu s podmodulom ili nadmodulom bez objasnjenja směra. Taka smyslova glosa jest koristna, ale ne smi izgubiti avtorskeho poredka.\par
\noindent\emph{Reason:} Dělitelj jednakogo ranga retains two conditions: divisor in the author's ordering of modules, and equality of rank. The definition of the basic module uses the congruence $bg\equiv0$ with nonzero $b$; it does not say that arbitrary modules of equal rank are equal. Jednakogo restricts the admissible divisors. The record resists replacing the expression with submodule or supermodule without explaining the direction. Such a gloss can help, but must preserve the author's ordering.\par
\noindent{\footnotesize Izvorne zapisy / Original records: HIST-g086-005, HIST-g092-015.}\par
\smallskip\noindent\emph{Sravnjenje s tekstom:} Jednakost ranga ne jest samostatnym tvrdženjem jednakosti modulov.\par
\noindent\emph{Textual comparison:} Equality of rank is not a standalone assertion of equality of modules.\par
\noindent{\footnotesize TEXT-COMP-765; pasaže / passages: 14265, 14267, 14269; \texttt{CORPUS\_FORM\_COMPARISONS.json}.}\par
\end{minipage}\par
\par\medskip\Needspace{7\baselineskip}
\noindent\begin{minipage}{\linewidth}
\hypertarget{TERM-1709}{}
\noindent\textbf{dělitelj v idealnom smyslu}\quad{\footnotesize\texttt{TERM-1709}}\par
\noindent\emph{Iztočnik / Source:} Teiler im Idealsinn\par
\noindent\emph{Rusky / Russian:} делитель в идеальном смысле\par
\noindent\emph{Ukrajinsky / Ukrainian:} дільник в ідеальному сенсі\par
\smallskip\noindent\emph{Razlog:} V idealnom smyslu jest potrěbno ograničenje slova dělitelj. Teorem sravnjuje dva idealy i jih rezultantne formy; on ne prěnosi bez dodatka obyčno polinomno děljenje na te formy. Zadržanje kvalifikatora razlučaje relaciju medžu idealami od faktorovanja izraza. Zapis ima přednost prěd kračim dělitelj zato, že kratša forma zakryvaje imenno razliku, na ktorej zavisi tvrdženje. To ne jest novy vid ideala, ale utočnjenje relacije.\par
\noindent\emph{Reason:} V idealnom smyslu is a necessary qualification of dělitelj. The theorem compares two ideals and their resultant forms; it does not simply transfer ordinary polynomial division to those forms. Keeping the qualifier distinguishes the relation between ideals from factorization of an expression. The longer wording is preferable because shortening it hides the distinction on which the statement depends. It specifies a relation, not a new kind of ideal.\par
\noindent{\footnotesize Izvorne zapisy / Original records: HIST-g089-008, HIST-g092-050.}\par
\smallskip\noindent\emph{Sravnjenje s tekstom:} Kvalifikator razlikuje idealy od jih rezultantnyh form.\par
\noindent\emph{Textual comparison:} The qualification distinguishes ideals from their resultant forms.\par
\noindent{\footnotesize TEXT-COMP-766; pasaže / passages: 14895; \texttt{CORPUS\_FORM\_COMPARISONS.json}.}\par
\end{minipage}\par
\par\medskip\Needspace{7\baselineskip}
\noindent\begin{minipage}{\linewidth}
\hypertarget{TERM-1710}{}
\noindent\textbf{dělitelj polja $K$}\quad{\footnotesize\texttt{TERM-1710}}\par
\noindent\emph{Iztočnik / Source:} Teiler von $K$\par
\noindent\emph{Rusky / Russian:} делитель поля $K$\par
\noindent\emph{Ukrajinsky / Ukrainian:} дільник поля $K$\par
\smallskip\noindent\emph{Razlog:} Dělitelj polja $K$ zadržaje historijno Teiler, v odlišnosti od modernizovanogo podpolje v drugom zapisu. Kontekst svazuje polja s sistemami uravnenj i govori o jih děliteljah; ne tvrdi, že polje ima nenulove čislove faktory. Arhivna prednost doslovnomu imeniku jest obnovima, ale ostavjaje potrěbu objasniti sodržanje. Za čitatelsky ukazatelj trěba svazati zapis s podpolje i javno pokazati tu razliku politiki. Ne možno prikazati vse historijne izbory kako už edinstveno normovane.\par
\noindent\emph{Reason:} Dělitelj polja $K$ retains historical Teiler, unlike the modernized podpolje in another record. The context connects fields with systems of equations and discusses their divisors; it is not asserting that a field has nonzero numerical factors. The archive's preference for the literal noun is recoverable, but leaves the inclusion relation in need of explanation. A reader's index should connect this entry with podpolje and make the policy difference explicit. The historical choices cannot all be presented as already standardized.\par
\noindent{\footnotesize Izvorne zapisy / Original records: HIST-g045-013.}\par
\smallskip\noindent\emph{Sravnjenje s tekstom:} Termin trěbuje smyslovu glosu, ne čislovo tumačenje.\par
\noindent\emph{Textual comparison:} The term needs a conceptual gloss, not a numerical interpretation.\par
\noindent{\footnotesize TEXT-COMP-767; pasaže / passages: 5664, 5666; \texttt{CORPUS\_FORM\_COMPARISONS.json}.}\par
\end{minipage}\par
\par\medskip\Needspace{7\baselineskip}
\noindent\begin{minipage}{\linewidth}
\hypertarget{TERM-1711}{}
\noindent\textbf{cěp děliteljev ili kratnikov}\quad{\footnotesize\texttt{TERM-1711}}\par
\noindent\emph{Iztočnik / Source:} Teiler- oder Vielfachenkette\par
\noindent\emph{Rusky / Russian:} цепь делителей или кратных\par
\noindent\emph{Ukrajinsky / Ukrainian:} ланцюг дільників або кратних\par
\smallskip\noindent\emph{Razlog:} Cěp děliteljev ili kratnikov imenovaje dva směra jednoj poredkovej struktury. V pasaži konečnost cěpi slěduje iz strogoj proměny ranga v konečnom intervalu; samo slovo cěp ne ustanovjaje zakona konečnosti dlja vseh idealov. Pravi dělitelj ili kratnik izključaje ravnost na danom stupenju. Arhivno kratnik jest formotvorny izbor blizky mnogokratno v drugih pasažah. Slovnik mora pokazati obě nazvy, ne proměniti smer cěpi radi povrhovoj jezyčnoj konsistencije.\par
\noindent\emph{Reason:} Cěp děliteljev ili kratnikov names the two directions of an ordering. In the passage, termination follows from strict changes of rank within a finite interval; the word chain does not establish a termination law for all ideals. A proper divisor or multiple excludes equality at that stage. The recorded kratnik is a morphological alternative to mnogokratno elsewhere. Both names belong in the dictionary; surface consistency must not reverse the chain's direction.\par
\noindent{\footnotesize Izvorne zapisy / Original records: HIST-g121-003.}\par
\smallskip\noindent\emph{Sravnjenje s tekstom:} Konečnost ima konkretny razlog v rangovom argumentu.\par
\noindent\emph{Textual comparison:} Termination has a specific reason in the rank argument.\par
\noindent{\footnotesize TEXT-COMP-768; pasaže / passages: 17436; \texttt{CORPUS\_FORM\_COMPARISONS.json}.}\par
\end{minipage}\par
\par\medskip\Needspace{7\baselineskip}
\noindent\begin{minipage}{\linewidth}
\hypertarget{TERM-1712}{}
\noindent\textbf{medžusobno prosty}\quad{\footnotesize\texttt{TERM-1712}}\par
\noindent\emph{Iztočnik / Source:} teilerfremd\par
\noindent\emph{Rusky / Russian:} взаимно просты\par
\noindent\emph{Ukrajinsky / Ukrainian:} взаємно прості\par
\smallskip\noindent\emph{Razlog:} Medžusobno prosty opisuje čislitelj i znamenatelj racionalnoj funkcije po skračenju. Medžusobno činje jasnym, že svojstvo prinaleži paru, ne každome elementu oddělno. Rusko vzaimno i ukrajinsko vzajemno podpirajut tu vzajemnu relaciju, ale ne prisiljajut jednu medžuslovjansku prědponu. Nynějšny tekst uživa vzajemno prosty. Razlika medžu medžusobno i vzajemno jest jezyčna; smysl običnoj vzaimnej prostosti ne trěba směšati s osoboj, směrnoj definicijeju relativnoj prostosti v teoriji idealov.\par
\noindent\emph{Reason:} Medžusobno prosty describes the numerator and denominator of a rational function after reduction. Medžusobno makes clear that the property belongs to the pair, not to each element separately. Russian vzaimno and Ukrainian vzajemno support the reciprocal relation without forcing one Interslavic form. The current text uses vzajemno prosty. The difference is linguistic; ordinary coprimality must not be confused with the special directed definition of relative primality in the ideal-theoretic paper.\par
\noindent{\footnotesize Izvorne zapisy / Original records: HIST-g042-012.}\par
\smallskip\noindent\emph{Sravnjenje s tekstom:} Vzaimna prostost odnosi se k čislitelju i znamenatelju.\par
\noindent\emph{Textual comparison:} Coprimality relates the numerator and denominator.\par
\noindent{\footnotesize TEXT-COMP-769; pasaže / passages: 5061; \texttt{CORPUS\_FORM\_COMPARISONS.json}.}\par
\end{minipage}\par
\par\medskip\Needspace{7\baselineskip}
\noindent\begin{minipage}{\linewidth}
\hypertarget{TERM-1713}{}
\noindent\textbf{vzaimno prosty}\quad{\footnotesize\texttt{TERM-1713}}\par
\noindent\emph{Iztočnik / Source:} teilerfremd\par
\noindent\emph{Rusky / Russian:} взаимно простой\par
\noindent\emph{Ukrajinsky / Ukrainian:} взаємно простий\par
\smallskip\noindent\emph{Razlog:} Vzaimno prosty označaje tu idealy, čiji največši obči dělitelj jest jediničny ideal. Arhiv naměrno razlikuje to od relativno prosty, čija definicija jest směrna. Pasaža dokazuje implikaciju od vzaimnej k oboustrannej relativnoj prostosti, ale dava $(x)$ i $(y)$ kako protivpriměr obratnomu tvrdženju. Zato sovrěmenna obyčna sinonimija slov coprime i relatively prime ne može upravljati tym prěvodom. Relacija ima prednost prěd jezyčnoju ekonomijeju.\par
\noindent\emph{Reason:} Vzaimno prosty here means ideals whose greatest common divisor is the unit ideal. The record deliberately distinguishes it from relativno prosty, whose definition is directed. The passage proves that coprimality implies relative primality in both directions, but gives $(x)$ and $(y)$ as a counterexample to the converse. The ordinary modern synonymy of coprime and relatively prime must therefore not govern this translation. Preserving the relation takes priority over verbal economy.\par
\noindent{\footnotesize Izvorne zapisy / Original records: HIST-g081-063, HIST-g082-063.}\par
\smallskip\noindent\emph{Sravnjenje s tekstom:} Protivpriměr čuva razliku medžu definicijami.\par
\noindent\emph{Textual comparison:} The counterexample preserves the distinction between definitions.\par
\noindent{\footnotesize TEXT-COMP-770; pasaže / passages: 13278, 13280; \texttt{CORPUS\_FORM\_COMPARISONS.json}.}\par
\end{minipage}\par
\par\medskip\Needspace{7\baselineskip}
\noindent\begin{minipage}{\linewidth}
\hypertarget{TERM-1714}{}
\noindent\textbf{vzaimno prosty / pravi dělitelji}\quad{\footnotesize\texttt{TERM-1714}}\par
\noindent\emph{Iztočnik / Source:} teilerfremd / echte Teiler\par
\noindent\emph{Rusky / Russian:} взаимно простые / собственные делители\par
\noindent\emph{Ukrajinsky / Ukrainian:} взаємно прості / справжні дільники\par
\smallskip\noindent\emph{Razlog:} Vzaimno prosty i pravi dělitelji vyražajut nezavisne uslovja: dva dělitelja imajut jediničny obči dělitelj i ani jedin ne jest cělym danym modulom. Pravi znači proper, ne pravostranny; v tekstu o jednostrannyh modulah ta dvojsmyslnost jest osoblivo važna. Vlastny dělitelj može služiti objasnjujučeju alternativoju bez proměny matematičnogo uslovja. Obyčna množenstvena glosa dolžna čuvati obratny směr poredka pri dělimosti, a ne tiho zaměniti relaciju inkluzijeju.\par
\noindent\emph{Reason:} Vzaimno prosty and pravi dělitelji express independent conditions: the two divisors have unit greatest common divisor, and neither equals the whole specified module. Pravi means proper, not right-sided; that ambiguity particularly matters in a discussion of one-sided modules. Vlastny dělitelj can serve as a clarifying alternative without changing the mathematical condition. Any set-theoretic gloss must preserve the reversed ordering involved in divisibility rather than silently substitute inclusion.\par
\noindent{\footnotesize Izvorne zapisy / Original records: HIST-g079-021.}\par
\smallskip\noindent\emph{Sravnjenje s tekstom:} Pravi i vzaimno prosty ne naznačajut isto svojstvo.\par
\noindent\emph{Textual comparison:} Proper and coprime do not name the same property.\par
\noindent{\footnotesize TEXT-COMP-771; pasaže / passages: 11165, 11167; \texttt{CORPUS\_FORM\_COMPARISONS.json}.}\par
\end{minipage}\par
\par\medskip\Needspace{7\baselineskip}
\noindent\begin{minipage}{\linewidth}
\hypertarget{TERM-1715}{}
\noindent\textbf{vzajemno prosty v $J$}\quad{\footnotesize\texttt{TERM-1715}}\par
\noindent\emph{Iztočnik / Source:} teilerfremd in $J$\par
\noindent\emph{Rusky / Russian:} взаимно простые в $J$\par
\noindent\emph{Ukrajinsky / Ukrainian:} взаємно прості в $J$\par
\smallskip\noindent\emph{Razlog:} Vzajemno prosty v $J$ ograničaje relaciju k imenovanoj oblasti. Obči dělitelj i količniki dolžny byti elementami $J$; rezultat v většem polinomnom kolcu ne rěšaje avtomatično vopros v $J$. Zato v $J$ ne jest dekorativny dodatok. Vybrano vzajemno prosty svazuje zapis s drugimi izrazami za coprimality, no ne dava prava pridati Bézoutovu identičnost v proizvoljnoj oblasti. Ta silnějša charakterizacija potrěbuje svoje uslovja.\par
\noindent\emph{Reason:} Vzajemno prosty v $J$ restricts the relation to the named domain. The common divisor and quotients must belong to $J$; a result in a larger polynomial ring does not automatically settle the question in $J$. The qualifier is therefore essential. Vzajemno prosty connects the entry with other expressions for coprimality, but does not authorize adding a Bézout identity in an arbitrary domain. That stronger characterization requires its own hypotheses.\par
\noindent{\footnotesize Izvorne zapisy / Original records: HIST-g048-005.}\par
\smallskip\noindent\emph{Sravnjenje s tekstom:} Izpustiti oblast značilo by proměniti vopros.\par
\noindent\emph{Textual comparison:} Omitting the domain would change the question.\par
\noindent{\footnotesize TEXT-COMP-772; pasaže / passages: 5859, 5863, 5865, 5871; \texttt{CORPUS\_FORM\_COMPARISONS.json}.}\par
\end{minipage}\par
\par\medskip\Needspace{7\baselineskip}
\noindent\begin{minipage}{\linewidth}
\hypertarget{TERM-1716}{}
\noindent\textbf{vzaimno prosty s modulom}\quad{\footnotesize\texttt{TERM-1716}}\par
\noindent\emph{Iztočnik / Source:} teilerfremd mit einem Modul\par
\noindent\emph{Rusky / Russian:} взаимно простой с модулем\par
\noindent\emph{Ukrajinsky / Ukrainian:} взаємно простий з модулем\par
\smallskip\noindent\emph{Razlog:} Vzaimno prosty s modulom govori o polinomu i modulu, ne o paru polinomov. V definiciji jednakogo vida formulě $YQ\equiv1$ i $XP\equiv1$ po različnyh modulah davajut točno značenje. Zadržanje s modulom čuvaje tip drugogo argumenta. $P$ i $Q$ ne sut tim samym oglašeny jedinicami cělogo kolca; obratimy sut odgovarajuče klasy ostatkov. Predlog s jest prirodny izraz relacije, ale jeje algebraičnu sodržinu ustanovjajut kongruencije.\par
\noindent\emph{Reason:} Vzaimno prosty s modulom relates a polynomial to a module, not two polynomials. In the definition of modules of the same type, $YQ\equiv1$ and $XP\equiv1$ modulo different modules specify the meaning. Keeping s modulom preserves the type of the second argument. This does not declare $P$ and $Q$ units of the whole ring: the corresponding residue classes are invertible. The preposition expresses the relation naturally, but the congruences determine its algebraic content.\par
\noindent{\footnotesize Izvorne zapisy / Original records: HIST-g079-051.}\par
\smallskip\noindent\emph{Sravnjenje s tekstom:} Kongruencije utočnjajut smysl relacije s modulom.\par
\noindent\emph{Textual comparison:} The congruences specify the relation to a module.\par
\noindent{\footnotesize TEXT-COMP-773; pasaže / passages: 11650, 11655, 11657, 11658, 11664; \texttt{CORPUS\_FORM\_COMPARISONS.json}.}\par
\end{minipage}\par
\par\medskip\Needspace{7\baselineskip}
\noindent\begin{minipage}{\linewidth}
\hypertarget{TERM-1717}{}
\noindent\textbf{vzaimno-prosto-nerazlozimy}\quad{\footnotesize\texttt{TERM-1717}}\par
\noindent\emph{Iztočnik / Source:} teilerfremd-irreduzibel\par
\noindent\emph{Rusky / Russian:} взаимно-просто-неразложимый\par
\noindent\emph{Ukrajinsky / Ukrainian:} взаємно-просто-нерозкладний\par
\smallskip\noindent\emph{Razlog:} Vzaimno-prosto-nerazlozimy jest historijna forma složenogo prilastka; neznačené z v nerazlozimy ne ustanovjaje odlišny termin od nerazložimy. Pasaža definije nerazložimost po odnošenju k prědstavjenju kako najmenše obče mnogokratno vzaimno prostyh idealov. Ne ide prosto o neobstavanje produktnogo razloženja. Spojniki pomogajut svazati vzaimno-prosto s konkretnym vidom nerazložimosti. Pravopisna variacija može byti objedinjena v ukazatelju, ale smyslova razlika od relativno-prosto-nerazložimy mora ostati.\par
\noindent\emph{Reason:} Vzaimno-prosto-nerazlozimy is a historical spelling of the compound adjective; its unaccented z does not establish a term distinct from nerazložimy. The passage defines indecomposability with respect to expression as a least common multiple of coprime ideals, not merely the absence of a product factorization. Hyphens help attach vzaimno-prosto to the particular kind of indecomposability. The spelling variant can be grouped in the index, but the distinction from relativno-prosto-nerazložimy must remain.\par
\noindent{\footnotesize Izvorne zapisy / Original records: HIST-g082-064.}\par
\smallskip\noindent\emph{Sravnjenje s tekstom:} Diakritična varianta ne ustanovjaje novu matematičnu kategoriju.\par
\noindent\emph{Textual comparison:} The diacritic variant does not establish a new mathematical category.\par
\noindent{\footnotesize TEXT-COMP-774; pasaže / passages: 13278, 13280; \texttt{CORPUS\_FORM\_COMPARISONS.json}.}\par
\end{minipage}\par
\par\medskip\Needspace{7\baselineskip}
\noindent\begin{minipage}{\linewidth}
\hypertarget{TERM-1718}{}
\noindent\textbf{vzaimno-prosto-nerazložimy}\quad{\footnotesize\texttt{TERM-1718}}\par
\noindent\emph{Iztočnik / Source:} teilerfremd-irreduzibel\par
\noindent\emph{Rusky / Russian:} взаимно-просто-неразложимый\par
\noindent\emph{Ukrajinsky / Ukrainian:} взаємно-просто-нерозкладний\par
\smallskip\noindent\emph{Razlog:} Vzaimno-prosto-nerazložimy zachovuje vid nerazložimosti prěd vseobčim imenikom. Vybor složenogo prilastka naměrno ne skracuje jeho do nerazložimy: cěly članok sravnjuje različne razloženja. Definicija VIII govori o najmenšem občem mnogokratnom vzaimno prostyh idealov. Dana forma jest diakritično polnějša než arhivno nerazlozimy, ne nova matematična kategorija. Uporjadžanje pravopisa ne smi odstraniti uslovje vzaimnej prostosti.\par
\noindent\emph{Reason:} Vzaimno-prosto-nerazložimy places the kind of indecomposability within the adjective. It is deliberately not shortened to nerazložimy, since the paper compares different decompositions. Definition VIII concerns least common multiples of coprime ideals. This form supplies the diacritic absent from archived nerazlozimy; it does not introduce a new mathematical category. Regularizing spelling must not erase the coprimality condition.\par
\noindent{\footnotesize Izvorne zapisy / Original records: HIST-g081-064.}\par
\smallskip\noindent\emph{Sravnjenje s tekstom:} Nerazložimost ograničena jest imenovanym vidom razloženja.\par
\noindent\emph{Textual comparison:} Indecomposability is restricted to the named kind of decomposition.\par
\noindent{\footnotesize TEXT-COMP-775; pasaže / passages: 13278; \texttt{CORPUS\_FORM\_COMPARISONS.json}.}\par
\end{minipage}\par
\par\medskip\Needspace{7\baselineskip}
\noindent\begin{minipage}{\linewidth}
\hypertarget{TERM-1719}{}
\noindent\textbf{vzaimno-prosto nerazložimy / relativno-prosto nerazložimy}\quad{\footnotesize\texttt{TERM-1719}}\par
\noindent\emph{Iztočnik / Source:} teilerfremd-irreduzibel / relativprim-irreduzibel\par
\noindent\emph{Rusky / Russian:} взаимно-просто неразложимый / относительно-просто неразложимый\par
\noindent\emph{Ukrajinsky / Ukrainian:} взаємно-просто нерозкладний / відносно-просто нерозкладний\par
\smallskip\noindent\emph{Razlog:} Vzaimno-prosto nerazložimy i relativno-prosto nerazložimy čuvajut dve kategorije, ktore Noether razlučaje. Medžuslovjanska konstrukcija označaje, k kakomu dopustimomu razloženju odnosi se nerazložimost. Pasaže o vzaimnej i relativnoj prostosti pokazujut, začto te nazvy ne možno objediniti: implikacija ima samo jedin směr. Spojnik ili razděl ne měnjaje togo. Trěba sravniti definicije, ne izabrati kračši izraz po zvučanji i odstraniti jednu kategoriju.\par
\noindent\emph{Reason:} Vzaimno-prosto nerazložimy and relativno-prosto nerazložimy preserve two categories distinguished by Noether. The construction specifies the admissible decomposition relative to which indecomposability is asserted. The passages on coprimality and relative primality explain why the names cannot be merged: the implication goes only one way. Hyphenation or spacing does not change that fact. The definitions, not the sound of the shorter expression, determine the distinction.\par
\noindent{\footnotesize Izvorne zapisy / Original records: HIST-g081-015, HIST-g082-015.}\par
\smallskip\noindent\emph{Sravnjenje s tekstom:} Različne definicije čuvajut različne nazvy.\par
\noindent\emph{Textual comparison:} Different definitions retain different names.\par
\noindent{\footnotesize TEXT-COMP-776; pasaže / passages: 12420, 12422, 12424, 13278, 13280; \texttt{CORPUS\_FORM\_COMPARISONS.json}.}\par
\end{minipage}\par
\par\medskip\Needspace{7\baselineskip}
\noindent\begin{minipage}{\linewidth}
\hypertarget{TERM-1720}{}
\noindent\textbf{vzaimno-prosto-nerazlozime idealy v nekomutativnom kolcu s jedinicoju}\quad{\footnotesize\texttt{TERM-1720}}\par
\noindent\emph{Iztočnik / Source:} teilerfremd-irreduzible Ideale in einem nicht-kommutativen Ring mit Einheit\par
\noindent\emph{Rusky / Russian:} взаимно-просто-неразложимые идеалы в некоммутативном кольце с единицей\par
\noindent\emph{Ukrajinsky / Ukrainian:} взаємно-просто-нерозкладні ідеали в некомутативному кільці з одиницею\par
\smallskip\noindent\emph{Razlog:} Dolga fraza čuvaje i vid idealov, i nekomutativnost, i obstavanje jedinice. Pasaža prěnosi teorem obstavanja konečnogo razloženja na vzaimno-prosto-nerazložime idealy v to okruženje, no ne prěnosi bez uslovij cělu teoriju primarnogo razloženja i jedinosti. Izpustiti s jedinicoju značilo by oslabiti hipotezu. Arhivno nerazlozime bez diakritiki jest pravopisna varianta, ne opravdanje dlja gubitka smyslovogo prilastka. Jasnost dostigaje se razloženjem frazy v objasnjenju, ne skracenjem samogo tvrdženja.\par
\noindent\emph{Reason:} The long phrase preserves the kind of ideals, noncommutativity, and the existence of an identity. The passage extends the existence of a finite decomposition into coprime-indecomposable ideals to this setting, not the entire theory of primary decomposition and uniqueness without conditions. Omitting s jedinicoju would weaken a hypothesis. Archived nerazlozime without the diacritic is a spelling variant, not a reason to lose the semantic qualifier. Clarity comes from unpacking the phrase in the explanation, not abridging the assertion.\par
\noindent{\footnotesize Izvorne zapisy / Original records: HIST-g081-076.}\par
\smallskip\noindent\emph{Sravnjenje s tekstom:} Obstavanje razloženja ne označaje neograničenu jedinost.\par
\noindent\emph{Textual comparison:} Existence of a decomposition does not assert unrestricted uniqueness.\par
\noindent{\footnotesize TEXT-COMP-777; pasaže / passages: 13409; \texttt{CORPUS\_FORM\_COMPARISONS.json}.}\par
\end{minipage}\par
\par\medskip\Needspace{7\baselineskip}
\noindent\begin{minipage}{\linewidth}
\hypertarget{TERM-1721}{}
\noindent\textbf{vzaimno proste idealy / jediničny ideal}\quad{\footnotesize\texttt{TERM-1721}}\par
\noindent\emph{Iztočnik / Source:} teilerfremde Ideale / Einheitsideal\par
\noindent\emph{Rusky / Russian:} взаимно простые идеалы / единичный идеал\par
\noindent\emph{Ukrajinsky / Ukrainian:} взаємно прості ідеали / одиничний ідеал\par
\smallskip\noindent\emph{Razlog:} Vzaimno proste idealy i jediničny ideal trěba čitati vměstě: v toj definiciji največši obči dělitelj dvoh idealov jest ideal sodržajuči vse elementy prstena. Jediničny ne znači ideal iz jedinogo elementa. Pravila za produkty i najmenše obče mnogokratne dalje predpokladajut komutativnost i jediničny element množenja. Arhiv daje prednost obyčnomu algebraičnomu izrazu, ale njegov smysl ustanovjaje ta definicija. Nynějšno vzajemno jest jezyčna varianta arhivnogo vzaimno, ne proměna uslovja.\par
\noindent\emph{Reason:} Vzaimno proste idealy and jediničny ideal must be read together: the greatest common divisor of the ideals is the ideal containing every ring element. Unit ideal does not mean an ideal with only one element. The subsequent rules for products and least common multiples also assume commutativity and a multiplicative identity. The record prefers ordinary algebraic terminology, but this definition determines its meaning. Current vzajemno varies the archived vzaimno linguistically, not mathematically.\par
\noindent{\footnotesize Izvorne zapisy / Original records: HIST-g111-005.}\par
\smallskip\noindent\emph{Sravnjenje s tekstom:} Jediničny ideal ustanovjaje smysl definicije.\par
\noindent\emph{Textual comparison:} The unit ideal determines the sense of the definition.\par
\noindent{\footnotesize TEXT-COMP-778; pasaže / passages: 16493, 16503, 16507; \texttt{CORPUS\_FORM\_COMPARISONS.json}.}\par
\end{minipage}\par
\par\medskip\Needspace{7\baselineskip}
\noindent\begin{minipage}{\linewidth}
\hypertarget{TERM-1722}{}
\noindent\textbf{vzaimno proste polinomy}\quad{\footnotesize\texttt{TERM-1722}}\par
\noindent\emph{Iztočnik / Source:} teilerfremde Polynome\par
\noindent\emph{Rusky / Russian:} взаимно простые полиномы\par
\noindent\emph{Ukrajinsky / Ukrainian:} взаємно прості поліноми\par
\smallskip\noindent\emph{Razlog:} Vzaimno proste polinomy zadržaje poznany algebraičny izraz bez doslovnogo kalkovanja njemeckogo fremd. V članku o krivyh dva taka polinomy mogut iměti obču nulovu točku: prvy teorem baš počinaje s takoj točkoju. Zato slovnik ne smi objasniti vzaimnu prostost kako neobstavanje občih nuljev ili, bez dodatnyh uslovij, kako tvorjenje jediničnogo ideala. Kontekst jest mnogoprěměnny. Ustanoviti točny smysl važněje jest než uskladniti izraz so vsakoju drugoju upotrěboju prosty.\par
\noindent\emph{Reason:} Vzaimno proste polinomy retains familiar algebraic language without literally calquing German fremd. In the paper on curves, two such polynomials may have a common zero: the first theorem starts with one. The dictionary must therefore not explain coprimality as absence of common zeros or, without further hypotheses, generation of the unit ideal. This is a multivariable setting. Identifying the precise sense matters more than making every use of prosty sound alike.\par
\noindent{\footnotesize Izvorne zapisy / Original records: HIST-g214-004.}\par
\smallskip\noindent\emph{Sravnjenje s tekstom:} V mnogoprěměnnom slučaju prostost ne izključaje obče nulove točky.\par
\noindent\emph{Textual comparison:} In several variables, coprimality does not exclude common zeros.\par
\noindent{\footnotesize TEXT-COMP-779; pasaže / passages: 45-isv.tex:53, 45-isv.tex:54, 45-isv.tex:55, 45-isv.tex:56, 45-isv.tex:57, 45-isv.tex:58, 45-isv.tex:59, 45-isv.tex:60, 45-isv.tex:348, 45-isv.tex:349, 45-isv.tex:350, 45-isv.tex:351, 45-isv.tex:352, 45-isv.tex:354, 45-isv.tex:355; \texttt{CORPUS\_FORM\_COMPARISONS.json}.}\par
\end{minipage}\par
\par\medskip\Needspace{7\baselineskip}
\noindent\begin{minipage}{\linewidth}
\hypertarget{TERM-1723}{}
\noindent\textbf{vzajemno proste primitivne polinomy}\quad{\footnotesize\texttt{TERM-1723}}\par
\noindent\emph{Iztočnik / Source:} teilerfremde primitive Polynome\par
\noindent\emph{Rusky / Russian:} взаимно простые примитивные многочлены\par
\noindent\emph{Ukrajinsky / Ukrainian:} взаємно прості примітивні поліноми\par
\smallskip\noindent\emph{Razlog:} Vzajemno proste primitivne polinomy spojuje dva uslovja normalizacije. V prědstavjenju racionalnogo funkcionala dva polinomy ne imajut obči netrivialny faktor, a primitivne odnosi se k jih cěločiselnym koeficientam. Pozitivny racionalny faktor izneseny jest oddělno. Zato primitivne ne znači prosto nerazložime, a vzaimna prostost ne može zaměniti koeficientnu normalizaciju. Arhiv prěnosi tu konstrukciju iz citovanogo Vebera; avtorskaja citacija sama po sobě ne dokazuje, že prěvoditelj pročital njegovo cělo izloženje.\par
\noindent\emph{Reason:} Vzajemno proste primitivne polinomy combines two normalization conditions. In the rational functional's representation, the two polynomials have no nontrivial common factor, while primitive concerns their integer coefficients. A positive rational factor is taken out separately. Primitive therefore does not merely mean irreducible, and coprimality cannot replace coefficient normalization. The archive carries over the construction attributed to Weber; that citation alone does not establish that the translator read Weber's full treatment.\par
\noindent{\footnotesize Izvorne zapisy / Original records: HIST-g060-005.}\par
\smallskip\noindent\emph{Sravnjenje s tekstom:} Primitivnost i vzaimna prostost imajut različne ulogi.\par
\noindent\emph{Textual comparison:} Primitivity and coprimality have different roles.\par
\noindent{\footnotesize TEXT-COMP-780; pasaže / passages: 7223, 7226, 7228; \texttt{CORPUS\_FORM\_COMPARISONS.json}.}\par
\end{minipage}\par
\par\medskip\Needspace{7\baselineskip}
\noindent\begin{minipage}{\linewidth}
\hypertarget{TERM-1724}{}
\noindent\textbf{vzaimno proste i vzaimno-prosto-nerazložime klasy}\quad{\footnotesize\texttt{TERM-1724}}\par
\noindent\emph{Iztočnik / Source:} teilerfremde und teilerfremd-irreduzible Klassen\par
\noindent\emph{Rusky / Russian:} взаимно простые и взаимно-просто-неразложимые классы\par
\noindent\emph{Ukrajinsky / Ukrainian:} взаємно прості та взаємно-просто-нерозкладні класи\par
\smallskip\noindent\emph{Razlog:} Klasy čuvaje tip objekta, na ktory prěnesena jest terminologija vzaimnej prostosti i odgovarjajučej nerazložimosti. Pasaža izrično prěhodi od idealov do dvustrannyh klasov matric i oprěděljaje dělimost prěz $A=PBQ$. Pri polnom rangu formulovano jest jednoznačno razloženje. Zato ne možno zaměniti klasy idealami radi známogo slovosloženja niti prěnesti tvrdženje o jedinosti na proizvoljne jednostranne klasy. Dolgy prilastok zachovuje vid razloženja, a imenik zachovuje jego oblast priměnjenja.\par
\noindent\emph{Reason:} Klasy preserves the type of object to which coprimality and the corresponding indecomposability are transferred. The passage explicitly moves from ideals to two-sided matrix classes and defines divisibility through $A=PBQ$. It states uniqueness of the decomposition at full rank. Replacing classes with ideals for verbal familiarity, or extending the uniqueness claim to arbitrary one-sided classes, would be misleading. The long adjective preserves the kind of decomposition; the noun preserves its domain of application.\par
\noindent{\footnotesize Izvorne zapisy / Original records: HIST-g081-101.}\par
\smallskip\noindent\emph{Sravnjenje s tekstom:} Imenik čuva oblast priměnjenja tvrdženja.\par
\noindent\emph{Textual comparison:} The noun preserves the statement's domain of application.\par
\noindent{\footnotesize TEXT-COMP-781; pasaže / passages: 13709, 13711, 13724; \texttt{CORPUS\_FORM\_COMPARISONS.json}.}\par
\end{minipage}\par
\par\medskip\Needspace{7\baselineskip}
\noindent\begin{minipage}{\linewidth}
\hypertarget{TERM-1725}{}
\noindent\textbf{vzaimna prostota}\quad{\footnotesize\texttt{TERM-1725}}\par
\noindent\emph{Iztočnik / Source:} Teilerfremdheit\par
\noindent\emph{Rusky / Russian:} взаимная простота\par
\noindent\emph{Ukrajinsky / Ukrainian:} взаємна простота\par
\smallskip\noindent\emph{Razlog:} Vzaimna prostota jest imenska forma relacije, v toj pasaži zapisanej kako $(r,s)=1$. Ona govori o dvou faktorah normy, ne o tom, že vsaky faktor jest prostym čislom. Rusko prostota i ukrajinsko prostota podpirajut izbor imenika, ale sravnjenje ne zaměnjaje proverku matematičnogo konteksta. Prostota bez vzaimna mogla by obrnuti relaciju paru v svojstvo jedinogo objekta. Osebne idealne definicije iz članka o razloženjah ne smějut byti prěnesene sem samo radi jednakoj njemeckoj nazvy.\par
\noindent\emph{Reason:} Vzaimna prostota is the noun for the relation written here as $(r,s)=1$. It concerns two factors of the norm, not the assertion that each factor is a prime number. Russian and Ukrainian prostota support the noun choice, but comparison does not replace checking the mathematical context. Omitting vzaimna could turn a relation between two objects into a property of one. The special ideal-theoretic definitions in the decomposition paper must not be imported merely because the German name is the same.\par
\noindent{\footnotesize Izvorne zapisy / Original records: HIST-g086-009, HIST-g092-019.}\par
\smallskip\noindent\emph{Sravnjenje s tekstom:} Formula utočnjaje smysl imenika.\par
\noindent\emph{Textual comparison:} The formula specifies the noun's meaning.\par
\noindent{\footnotesize TEXT-COMP-782; pasaže / passages: 14424; \texttt{CORPUS\_FORM\_COMPARISONS.json}.}\par
\end{minipage}\par
\par\medskip\Needspace{7\baselineskip}
\noindent\begin{minipage}{\linewidth}
\hypertarget{TERM-1726}{}
\noindent\textbf{cěp děliteljev}\quad{\footnotesize\texttt{TERM-1726}}\par
\noindent\emph{Iztočnik / Source:} Teilerkette\par
\noindent\emph{Rusky / Russian:} цепь делителей\par
\noindent\emph{Ukrajinsky / Ukrainian:} ланцюг дільників\par
\smallskip\noindent\emph{Razlog:} Cěp děliteljev pojmenovuje poslědovatelnost po dělimosti. Arhiv objasnjaje jeju kako descending chain, no to slovo samo ne ustanovjaje směr inkluzije: dělitelj ideala sodržaje ideal, ktory děli. V danom izloženju cěpi děliteljev i kratnyh oddělno upotrěbljajut se pri konstrukciji kompozicijnogo reda. Zato glosa rastuča ili padajuča mora vsegda imenovati poredok. Slovjanska nazva čuvaje avtorske različenje bez skrytoj zaměny jednoju iz sovrěmennyh konvencij.\par
\noindent\emph{Reason:} Cěp děliteljev names a sequence ordered by divisibility. The archival rationale calls it a descending chain, but that phrase alone does not specify inclusion direction: an ideal divisor contains the ideal it divides. The exposition uses divisor and multiple chains separately in constructing a composition series. An ascending or descending gloss must therefore state which ordering it means. The Slavic expression preserves the author's distinction without silently substituting a modern convention.\par
\noindent{\footnotesize Izvorne zapisy / Original records: HIST-g117-006.}\par
\smallskip\noindent\emph{Sravnjenje s tekstom:} Rastuča i padajuča potrěbujut imenovany poredok.\par
\noindent\emph{Textual comparison:} Ascending and descending require a specified ordering.\par
\noindent{\footnotesize TEXT-COMP-783; pasaže / passages: 17189, 17193, 17199, 17203, 17207; \texttt{CORPUS\_FORM\_COMPARISONS.json}.}\par
\end{minipage}\par
\par\medskip\Needspace{7\baselineskip}
\noindent\begin{minipage}{\linewidth}
\hypertarget{TERM-1727}{}
\noindent\textbf{cěp idealov po dělimosti / cěp mnogokratnyh modulo ideal}\quad{\footnotesize\texttt{TERM-1727}}\par
\noindent\emph{Iztočnik / Source:} Teilerkette aus Idealen / Vielfachenkette modulo einem Ideal\par
\noindent\emph{Rusky / Russian:} цепь идеалов по делимости / цепь кратных по модулю идеала\par
\noindent\emph{Ukrajinsky / Ukrainian:} ланцюг ідеалів за діленням / ланцюг кратних за модулем ідеалу\par
\smallskip\noindent\emph{Razlog:} Cěp idealov po dělimosti i cěp mnogokratnyh modulo ideal čuvajut dve različne hipotezy. V pasaži vtora ograničena jest fiksovanym nenulovym idealom, a ne vsimi padajučimi cěpami bez ograničenja. Dalje potrěbne sut jedinica, odsutnost nulovyh děliteljev i integralna zamknjenost. Zato kratko ime modernogo klasa prstenov ne jest dovoljnym prěvodom celogo tvrdženja. Arhiv daje prednost izričnomu opisu uslovij; slovnik objasnjaje jih poredok, ne udaljaje jih.\par
\noindent\emph{Reason:} Cěp idealov po dělimosti and cěp mnogokratnyh modulo ideal preserve two distinct hypotheses. The second is restricted to a fixed nonzero ideal, not all descending chains without restriction. The passage also requires an identity, absence of zero divisors, and integral closure. A modern class name for rings is therefore not an adequate translation of the whole statement. The archive favours explicit conditions; the dictionary explains their ordering rather than removing them.\par
\noindent{\footnotesize Izvorne zapisy / Original records: HIST-g103-005.}\par
\smallskip\noindent\emph{Sravnjenje s tekstom:} Nenulovy ideal ograničaje vtoru cěpovu hipotezu.\par
\noindent\emph{Textual comparison:} The nonzero ideal restricts the second chain hypothesis.\par
\noindent{\footnotesize TEXT-COMP-784; pasaže / passages: 15886; \texttt{CORPUS\_FORM\_COMPARISONS.json}.}\par
\end{minipage}\par
\par\medskip\Needspace{7\baselineskip}
\noindent\begin{minipage}{\linewidth}
\hypertarget{TERM-1728}{}
\noindent\textbf{teorem o cěpah děliteljev}\quad{\footnotesize\texttt{TERM-1728}}\par
\noindent\emph{Iztočnik / Source:} Teilerkettensatz\par
\noindent\emph{Rusky / Russian:} теорема о цепях делителей\par
\noindent\emph{Ukrajinsky / Ukrainian:} теорема про ланцюги дільників\par
\smallskip\noindent\emph{Razlog:} Teorem o cěpah děliteljev čuva dokazny smysl slova Satz v toj pasaži. Članok razgledaje, kogda svojstvo ostaje validno pri konečnom integralnom razširjenju, i izrično govori, začto ne može byti priměnjeno neposrědnje bez integralnoj zamknjenosti. Zaměna s neograničenym sloganom o konečnosti izgubila by tu ogradu. To že drugde vybrano jest uslovje cěpov děliteljev ne jest potrěbno protivrěčje: jedno ime opisuje teorem prěnosa, drugo prědpoloženu aksiomu.\par
\noindent\emph{Reason:} Teorem o cěpah děliteljev retains the theorem sense of Satz in this passage. The paper considers when the property persists under a finite integral extension and expressly explains why it cannot be applied directly without integral closure. An unrestricted slogan about finiteness would lose that qualification. Choosing uslovje cěpov děliteljev elsewhere need not be inconsistent: one expression describes a transfer theorem, the other an assumed axiom.\par
\noindent{\footnotesize Izvorne zapisy / Original records: HIST-g106-009.}\par
\smallskip\noindent\emph{Sravnjenje s tekstom:} Pasaža čuva uslovja priměnjenja teorema.\par
\noindent\emph{Textual comparison:} The passage preserves the theorem's conditions of application.\par
\noindent{\footnotesize TEXT-COMP-785; pasaže / passages: 15987, 16021, 16033; \texttt{CORPUS\_FORM\_COMPARISONS.json}.}\par
\end{minipage}\par
\par\medskip\Needspace{7\baselineskip}
\noindent\begin{minipage}{\linewidth}
\hypertarget{TERM-1729}{}
\noindent\textbf{uslovje cěpov děliteljev}\quad{\footnotesize\texttt{TERM-1729}}\par
\noindent\emph{Iztočnik / Source:} Teilerkettensatz\par
\noindent\emph{Rusky / Russian:} условие цепей делителей\par
\noindent\emph{Ukrajinsky / Ukrainian:} умова ланцюгів дільників\par
\smallskip\noindent\emph{Razlog:} Uslovje cěpov děliteljev daje prednost funkcionu izraza v uvodu: ono jest prědpoložena aksioma I. Definicija govori o prěryvanju strogoj cěpi i, ravnoznačno, o stabilizaciji cěpi. Nazvati to tut samo teorem moglo by zakryti, čto jest prědpoloženo, a čto bude dokazano. Čislo I ostaje stabilnym označenjem. Izbor uslovje ne modernizuje sodržinu, ale činje logičnu ulogu čitljivu.\par
\noindent\emph{Reason:} Uslovje cěpov děliteljev follows the expression's role in the introduction: it is assumed as axiom I. The definition describes termination of a strict chain and, equivalently, eventual stabilization of a chain. Calling it simply a theorem here could obscure what is assumed and what will be proved. The label I remains unchanged. Uslovje makes the logical role readable without modernizing the substance.\par
\noindent{\footnotesize Izvorne zapisy / Original records: HIST-g107-002.}\par
\smallskip\noindent\emph{Sravnjenje s tekstom:} Nazva pokazuje, čto jest prědpoloženo.\par
\noindent\emph{Textual comparison:} The name identifies what is assumed.\par
\noindent{\footnotesize TEXT-COMP-786; pasaže / passages: 16120, 16121; \texttt{CORPUS\_FORM\_COMPARISONS.json}.}\par
\end{minipage}\par
\par\medskip\Needspace{7\baselineskip}
\noindent\begin{minipage}{\linewidth}
\hypertarget{TERM-1730}{}
\noindent\textbf{uslovje cěpov děliteljev / fiksovany dobry porjadok / konečna baza ideala}\quad{\footnotesize\texttt{TERM-1730}}\par
\noindent\emph{Iztočnik / Source:} Teilerkettensatz / feste Wohlordnung / endliche Idealbasis\par
\noindent\emph{Rusky / Russian:} условие цепей делителей / фиксированный хороший порядок / конечный базис идеала\par
\noindent\emph{Ukrajinsky / Ukrainian:} умова ланцюгів дільників / фіксований добрий порядок / скінченний базис ідеалу\par
\smallskip\noindent\emph{Razlog:} Uslovje cěpov děliteljev, fiksovany dobry porjadok i konečna baza ideala ne sut tri sinonimy. Pasaža izrično povezuje dobyvanje konečnoj bazy s fiksovanym dobrym poredkom; to jest zavisimost argumenta. Baza ideala tut znači konečny sistem tvorečih elementov, ne nepreměnno linearno nezavisnu bazu. Arhiv želi sohraniti množinsko-teoretičnu ogradu. Jasno izloženje mora razděliti uslovje, pomoćnu strukturu i zaključok, a ne ostaviti jih kako nerazlikovan spisok terminov.\par
\noindent\emph{Reason:} The chain condition, a fixed well-ordering, and a finite ideal basis are not three synonyms. The passage explicitly connects obtaining the finite basis with a fixed well-ordering; this is a dependency in the argument. Ideal basis here means a finite generating set, not necessarily a linearly independent basis. The archive seeks to preserve the set-theoretic qualification. A clear explanation separates the condition, the auxiliary structure, and the conclusion rather than leaving them as an undifferentiated list.\par
\noindent{\footnotesize Izvorne zapisy / Original records: HIST-g112-009.}\par
\smallskip\noindent\emph{Sravnjenje s tekstom:} Svez uslovij ne jest sinonimija terminov.\par
\noindent\emph{Textual comparison:} The dependency between conditions is not synonymy.\par
\noindent{\footnotesize TEXT-COMP-787; pasaže / passages: 16615, 16678; \texttt{CORPUS\_FORM\_COMPARISONS.json}.}\par
\end{minipage}\par
\par\medskip\Needspace{7\baselineskip}
\noindent\begin{minipage}{\linewidth}
\hypertarget{TERM-1731}{}
\noindent\textbf{uslovje cěpov děliteljev / uslovje cěpov kratnyh}\quad{\footnotesize\texttt{TERM-1731}}\par
\noindent\emph{Iztočnik / Source:} Teilerkettensatz / Vielfachenkettensatz\par
\noindent\emph{Rusky / Russian:} условие цепей делителей / условие цепей кратных\par
\noindent\emph{Ukrajinsky / Ukrainian:} умова ланцюгів дільників / умова ланцюгів кратних\par
\smallskip\noindent\emph{Razlog:} Dva izraza s uslovje naznačajut dve cěpove hipotezy, čija protivnost po dělimosti mora ostati vidima. Modulny teorem prěnosi odgovarjajuče uslovje iz prstena na konečnu modulnu oblast; ne tvrdi, že oba sut avtomatično ravnoznačne. Nazvy kratnyh, kratnikov i mnogokratnyh sut jezyčne varianty, ktore slovnik može svazati. Izbirati jednu formu v ukazatelju možno, ale spojiti dve matematične uslovja v jedno bez dodatka ne možno.\par
\noindent\emph{Reason:} The two expressions with uslovje name separate chain hypotheses whose opposite divisibility directions must remain visible. The module theorem transfers the corresponding condition from the ring to a finite module domain; it does not say the two conditions are automatically equivalent. Kratnyh, kratnikov, and mnogokratnyh are linguistic variants the dictionary can link. A shared index form is possible; conflating the mathematical conditions is not.\par
\noindent{\footnotesize Izvorne zapisy / Original records: HIST-g109-004.}\par
\smallskip\noindent\emph{Sravnjenje s tekstom:} Odgovarjajuče uslovje prěnosi se bez objedinjanja dvou směrov.\par
\noindent\emph{Textual comparison:} The corresponding condition transfers without merging the two directions.\par
\noindent{\footnotesize TEXT-COMP-788; pasaže / passages: 16314, 16330, 16334, 16338; \texttt{CORPUS\_FORM\_COMPARISONS.json}.}\par
\end{minipage}\par
\par\medskip\Needspace{7\baselineskip}
\noindent\begin{minipage}{\linewidth}
\hypertarget{TERM-1732}{}
\noindent\textbf{uslovje cěpov děliteljev / dobry porjadok / postulat izbora}\quad{\footnotesize\texttt{TERM-1732}}\par
\noindent\emph{Iztočnik / Source:} Teilerkettensatz / Wohlordnung / Auswahlpostulat\par
\noindent\emph{Rusky / Russian:} условие цепей делителей / хороший порядок / постулат выбора\par
\noindent\emph{Ukrajinsky / Ukrainian:} умова ланцюгів дільників / добрий порядок / постулат вибору\par
\smallskip\noindent\emph{Razlog:} Dobry porjadok jest tehničny well-order, ne pohvala kakovosti uporjadkovanja. Noether fiksuje takovy poredok, iz njega izbiraje prvu dopustimu bazu i objasnjaje, že pri už danom poredku ne trěba novogo postulata izbora. Arhiv zachovuje to razlikovanje vměsto tichej modernizacije osnov dokaza. Postulat izbora ne znači izbiranje slova prěvoditeljem. Terminologija mora činit čitljivu baš matematiku izbora, vključajuči ogradu za izčetny slučaj.\par
\noindent\emph{Reason:} Dobry porjadok means a technical well-order, not a favourable assessment of an ordering. Noether fixes such an ordering, uses it to choose the first admissible basis, and explains that an already given ordering requires no further choice postulate. The record preserves that distinction rather than silently modernizing the proof's foundations. Postulat izbora does not refer to a translator's choice of words. The terminology must make the mathematics of selection readable, including the qualification for the countable case.\par
\noindent{\footnotesize Izvorne zapisy / Original records: HIST-g113-001.}\par
\smallskip\noindent\emph{Sravnjenje s tekstom:} Avtorske ogrady o izboru ostajut čitljive.\par
\noindent\emph{Textual comparison:} The author's qualifications about choice remain visible.\par
\noindent{\footnotesize TEXT-COMP-789; pasaže / passages: 16678, 16682; \texttt{CORPUS\_FORM\_COMPARISONS.json}.}\par
\end{minipage}\par
\par\medskip\Needspace{7\baselineskip}
\noindent\begin{minipage}{\linewidth}
\hypertarget{TERM-1733}{}
\noindent\textbf{podsistemy}\quad{\footnotesize\texttt{TERM-1733}}\par
\noindent\emph{Iztočnik / Source:} Teilsysteme\par
\noindent\emph{Rusky / Russian:} подсистемы\par
\noindent\emph{Ukrajinsky / Ukrainian:} підсистеми\par
\smallskip\noindent\emph{Razlog:} Podsistemy spojuje prozračnu slovjansku prědponu pod- s internacionalnym sistema. Ono označaje sostavny sistem vnutri většego sistema form, ne bezuslovno proizvoljno podmnožstvo. Zaključna tabela razlučaje pojedine formy i podsistemy, zato imeniki ne sut izměnljivi pri sčitanju. Arhiv daje kratko sravnjenje s ruskim podsistemy i ukrajinskim pidsystemy; tvrditi na toj osnově izměrenu razumělivost dlja vseh Slovjan bylo by neopravdano.\par
\noindent\emph{Reason:} Podsistemy combines transparent Slavic pod- with international sistema. It denotes a constituent system within the larger system of forms, not automatically an arbitrary subset. The concluding table distinguishes individual forms from subsystems, so the nouns cannot be exchanged when counting. The archive gives a brief comparison with Russian podsistemy and Ukrainian pidsystemy; that does not establish measured intelligibility for all Slavic readers.\par
\noindent{\footnotesize Izvorne zapisy / Original records: HIST-g027-011.}\par
\smallskip\noindent\emph{Sravnjenje s tekstom:} Sistemy ne sut isto, čto pojedine formy.\par
\noindent\emph{Textual comparison:} Systems are not the same as individual forms.\par
\noindent{\footnotesize TEXT-COMP-790; pasaže / passages: 3505, 3509, 3510, 3511, 3512; \texttt{CORPUS\_FORM\_COMPARISONS.json}.}\par
\end{minipage}\par
\par\medskip\Needspace{7\baselineskip}
\noindent\begin{minipage}{\linewidth}
\hypertarget{TERM-1735}{}
\noindent\textbf{ternarna bikvadratična forma; [objasnjujuče prvo nazvanje / explanatory first naming]: ternarna forma četvrtogo stepena}\quad{\footnotesize\texttt{TERM-1735}}\par
\noindent\emph{Iztočnik / Source:} ternäre biquadratische Form\par
\noindent\emph{Rusky / Russian:} тернарная биквадратичная форма; [objasnjujuče prvo nazvanje / explanatory first naming]: тернарная форма четвёртой степени\par
\noindent\emph{Ukrajinsky / Ukrainian:} тернарна біквадратична форма; [objasnjujuče prvo nazvanje / explanatory first naming]: тернарна форма четвертого степеня\par
\smallskip\noindent\emph{Razlog:} Ternarna bikvadratična forma čuvaje historijno nazvanje, a ternarna forma četvrtogo stepena objasnjaje njegovo značenje. Ternarna upučaje na tri prěměnne; četvrty stepenj na homogennu formu stepena četyri. Bikvadratična ne znači tut, že vsaka prěměnna ima zasebno stepenj dva. Arhiv dopušča objasnjujuče prvo nazvanje i kračšu historijnu formu v proze. Nynějšny uvod prvogo članka daje obe nazvy vměstě, tako že čitatelj ne mora ugadyvati starši termin.\par
\noindent\emph{Reason:} Ternarna bikvadratična forma retains the historical name, while ternarna forma četvrtogo stepena explains its meaning. Ternary refers to three variables; degree four refers to a homogeneous quartic form. Biquadratic here does not mean degree two separately in every variable. The record allows an explanatory first naming and the shorter historical expression in the prose. The current first-paper introduction gives both names together, so readers need not guess the older usage.\par
\noindent{\footnotesize Izvorne zapisy / Original records: HIST-g002-002.}\par
\smallskip\noindent\emph{Sravnjenje s tekstom:} Dve nazvy izrično povezane sut v uvodu.\par
\noindent\emph{Textual comparison:} The introduction explicitly connects the two names.\par
\noindent{\footnotesize TEXT-COMP-792; pasaže / passages: 43, 53, 150, 161, 178; \texttt{CORPUS\_FORM\_COMPARISONS.json}.}\par
\end{minipage}\par
\par\medskip\Needspace{7\baselineskip}
\noindent\begin{minipage}{\linewidth}
\hypertarget{TERM-1736}{}
\noindent\textbf{ternarna forma četvrtogo stepena; kratko: ternarna kvartika}\quad{\footnotesize\texttt{TERM-1736}}\par
\noindent\emph{Iztočnik / Source:} ternäre biquadratische Form\par
\noindent\emph{Rusky / Russian:} тернарная форма четвёртой степени; коротко: тернарная квартика\par
\noindent\emph{Ukrajinsky / Ukrainian:} тернарна форма четвертого степеня; коротко: тернарна квартика\par
\smallskip\noindent\emph{Razlog:} Ternarna forma četvrtogo stepena jest polna opisna nazva; ternarna kvartika jest kračši moderny izraz. Arhiv dava prednost polnoj formě pri prvom pojavjenju, potom dopušča kračšu. On upominaje ukrajinsku rabotu o ternarnyh formah kako sravniteljne svjedočstvo; to ne ustanovjaje oficijalnu medžuslovjansku normu za kvartika. Vsaka nazva mora sohraniti tri prěměnne, homogennost i stepenj četyri. Ne trěba zaměnjati tu formu geometricoju krivoju bez objasnjenja prěhoda.\par
\noindent\emph{Reason:} Ternarna forma četvrtogo stepena is the full descriptive name; ternarna kvartika is a shorter modern expression. The archive prefers the full wording on first occurrence, then permits the shorter form. It cites a Ukrainian study of ternary forms as comparative evidence, not an official Interslavic norm for kvartika. Both names must preserve three variables, homogeneity, and degree four. The form must not be replaced by its geometric curve without explaining the passage between them.\par
\noindent{\footnotesize Izvorne zapisy / Original records: HIST-g001-001.}\par
\smallskip\noindent\emph{Sravnjenje s tekstom:} Kratko ime ne smi izbrisati stepenj i čislo prěměnnyh.\par
\noindent\emph{Textual comparison:} The short name must retain degree and number of variables.\par
\noindent{\footnotesize TEXT-COMP-793; pasaže / passages: 43, 53, 69; \texttt{CORPUS\_FORM\_COMPARISONS.json}.}\par
\end{minipage}\par
\par\medskip\Needspace{7\baselineskip}
\noindent\begin{minipage}{\linewidth}
\hypertarget{TERM-1737}{}
\noindent\textbf{teorija cělyh elementov}\quad{\footnotesize\texttt{TERM-1737}}\par
\noindent\emph{Iztočnik / Source:} Theorie der ganzen Größen\par
\noindent\emph{Rusky / Russian:} теория целых элементов\par
\noindent\emph{Ukrajinsky / Ukrainian:} теорія цілих елементів\par
\smallskip\noindent\emph{Razlog:} Teorija cělyh elementov modernizuje obče Größen do elementov, ale čuvaje specialny smysl cělosti. Element jest cěly po odnošenju k danomu prstenu, kogda izpolnjaje monično uravnenje s koeficientami iz nego. Zato cěly ne znači samo cělo racionalno čislo i ne govori o integraciji. Arhiv dava prednost slovjanskoj osnově prěd novym neprězračnym zajętkem. Relativno uslovje nad prstenom mora ostati dostupno čitatelju, i kogda jest radi kratkosti izpuščeno pri povtorjenju.\par
\noindent\emph{Reason:} Teorija cělyh elementov modernizes the general Größen to elements while retaining the specialized meaning of integrality. An element is integral relative to the specified ring when it satisfies a monic equation with coefficients in that ring. Cěly therefore does not mean only an ordinary integer and does not concern integration. The record favours a Slavic base over a new opaque borrowing. The reference ring must remain recoverable even when omitted in repeated wording for brevity.\par
\noindent{\footnotesize Izvorne zapisy / Original records: HIST-g108-001.}\par
\smallskip\noindent\emph{Sravnjenje s tekstom:} Monično uravnenje razlikuje cělost od običnoj cěločislenosti.\par
\noindent\emph{Textual comparison:} The monic equation distinguishes integrality from ordinary integerhood.\par
\noindent{\footnotesize TEXT-COMP-794; pasaže / passages: 16172, 16200, 16202, 16204, 16208, 16210, 16212; \texttt{CORPUS\_FORM\_COMPARISONS.json}.}\par
\end{minipage}\par
\par\medskip\Needspace{7\baselineskip}
\noindent\begin{minipage}{\linewidth}
\hypertarget{TERM-1738}{}
\noindent\textbf{totalno vzaimno prosta sistema}\quad{\footnotesize\texttt{TERM-1738}}\par
\noindent\emph{Iztočnik / Source:} total teilerfremdes System\par
\noindent\emph{Rusky / Russian:} [Bukvy ne sut zachovane / Lettering not recoverable]\par
\noindent\emph{Ukrajinsky / Ukrainian:} [Bukvy ne sut zachovane / Lettering not recoverable]\par
\smallskip\noindent\emph{Razlog:} Totalno vzaimno prosta sistema čuvaje silnějše uslovje než poparna vzaimna prostost v nekomutativnom slučaju. Za vsako razděljenje sistemy najmenše obče mnogokratno jednoj česti mora byti vzaimno prosto s odgovarjajučim mnogokratnym ostaloj česti. Pasaža daje tu definiciju i izrično ograničaje obrat od poparnej prostosti. Totalno zato ne jest retorično posiljenje i ne može byti izpuščeno. Ošteđene ruske i ukrajinske bukvy v zapisu ne dopolnjajemo kakoby ony byly zachovane.\par
\noindent\emph{Reason:} Totalno vzaimno prosta sistema preserves a stronger condition than pairwise coprimality in the noncommutative setting. For every partition of the system, the least common multiple of one part must be coprime to that of the remainder. The passage gives the definition and explicitly restricts the converse from pairwise coprimality. Totalno is therefore not rhetorical emphasis and cannot be omitted. Damaged Russian and Ukrainian lettering in the record is not silently filled in as if it had survived.\par
\noindent{\footnotesize Izvorne zapisy / Original records: HIST-g079-030.}\par
\smallskip\noindent\emph{Sravnjenje s tekstom:} Izrična definicija čuva silnějše uslovje.\par
\noindent\emph{Textual comparison:} The explicit definition preserves the stronger condition.\par
\noindent{\footnotesize TEXT-COMP-795; pasaže / passages: 11350, 11358, 11360; \texttt{CORPUS\_FORM\_COMPARISONS.json}.}\par
\end{minipage}\par
\par\medskip\Needspace{7\baselineskip}
\noindent\begin{minipage}{\linewidth}
\hypertarget{TERM-1739}{}
\noindent\textbf{determinant transformacije}\quad{\footnotesize\texttt{TERM-1739}}\par
\noindent\emph{Iztočnik / Source:} Transformationsdeterminante\par
\noindent\emph{Rusky / Russian:} определитель преобразования\par
\noindent\emph{Ukrajinsky / Ukrainian:} визначник перетворення\par
\smallskip\noindent\emph{Razlog:} Determinant transformacije imenovaje determinant matrice prěhoda medžu dvoma bazami. V dokazu o diskriminantu on nastupaje v kvadratu, zato že matrica bilinearnej formy měnjaje se v oběh argumentah. Izpustiti kvadrat ili čitati determinant kako samu transformaciju značilo by proměniti argument. Internacionalne determinant i transformacija pomogajut povezati termin s vidimoju formuloju. Prěvod ne mora izmyšljati novu nazvu za vsak takovy faktor.\par
\noindent\emph{Reason:} Determinant transformacije names the determinant of the change-of-basis matrix. In the discriminant argument it appears squared because the matrix of the bilinear form changes in both arguments. Dropping the square or treating determinant as the transformation itself would change the reasoning. International determinant and transformacija connect the term with the displayed formula. No separate new name is needed for each such factor.\par
\noindent{\footnotesize Izvorne zapisy / Original records: HIST-g137-004.}\par
\smallskip\noindent\emph{Sravnjenje s tekstom:} Faktor v diskriminantu jest kvadrat, ne sam determinant.\par
\noindent\emph{Textual comparison:} The discriminant factor is the square, not the determinant itself.\par
\noindent{\footnotesize TEXT-COMP-796; pasaže / passages: 18094, 18096, 18098, 18100, 18102, 18104; \texttt{CORPUS\_FORM\_COMPARISONS.json}.}\par
\end{minipage}\par
\par\medskip\Needspace{7\baselineskip}
\noindent\begin{minipage}{\linewidth}
\hypertarget{TERM-1740}{}
\noindent\textbf{determinant transformacije / jedinica}\quad{\footnotesize\texttt{TERM-1740}}\par
\noindent\emph{Iztočnik / Source:} Transformationsdeterminante / Einheit\par
\noindent\emph{Rusky / Russian:} определитель преобразования / единица\par
\noindent\emph{Ukrajinsky / Ukrainian:} визначник перетворення / одиниця\par
\smallskip\noindent\emph{Razlog:} Determinant transformacije mora byti jediniceju osnovnej oblasti, aby nova nezavisna sistema byla bazoju togo samogo modula. Nenulovy determinant sam po sobě ne dostačuje nad oblastju cělosti, a jedinica ne znači samo čislo $1$: jest to obratimy element. Arhiv sohranjaje kriterij v obou směrah. Kosá črta v naslovu povezuje determinant i njegov uslovje, ne predstavjaje jih sinonimami. To razlikovanje potrěbno jest takože za nezavisnost diskriminantnogo ideala od izbora bazy.\par
\noindent\emph{Reason:} The change-of-basis determinant must be a unit of the base domain for the new independent system to be a basis of the same module. Nonzero alone is insufficient over an integral domain, and unit does not mean only the number $1$: it means an invertible element. The archive preserves the criterion in both directions. The slash in the heading connects a determinant with its required property, not two synonyms. The distinction also matters to the discriminant ideal's independence of the basis.\par
\noindent{\footnotesize Izvorne zapisy / Original records: HIST-g130-008.}\par
\smallskip\noindent\emph{Sravnjenje s tekstom:} Nenulovost ne zaměnjaje obratimost v osnovnoj oblasti.\par
\noindent\emph{Textual comparison:} Being nonzero does not replace invertibility in the base domain.\par
\noindent{\footnotesize TEXT-COMP-797; pasaže / passages: 17827; \texttt{CORPUS\_FORM\_COMPARISONS.json}.}\par
\end{minipage}\par
\par\medskip\Needspace{7\baselineskip}
\noindent\begin{minipage}{\linewidth}
\hypertarget{TERM-1741}{}
\noindent\textbf{transformacijne veli?iny [Čestično nezachovane bukvy / Partly unrecoverable lettering]}\quad{\footnotesize\texttt{TERM-1741}}\par
\noindent\emph{Iztočnik / Source:} Transformationsgr??en [Čestično nezachovane bukvy / Partly unrecoverable lettering]\par
\noindent\emph{Rusky / Russian:} [Bukvy ne sut zachovane / Lettering not recoverable]\par
\noindent\emph{Ukrajinsky / Ukrainian:} [Bukvy ne sut zachovane / Lettering not recoverable]\par
\smallskip\noindent\emph{Razlog:} Transformacijne veličiny jest čitljiva nynějšna forma imenovanogo algebraičnogo objekta. V staršem zapisu čest bukv nestala; čitljivu formu zato opirajemo na pasažu, ne na ugadany popravok arhiva. Konkretno dana jest veličina $c_S^Tc_T/c_{ST}$ i uslovje, že vse te veličiny sut ravne jedinici. Internacionalny prilastok zadržaje svez s transformacijeju, a veličina ne znači samo čislovy rozměr. Jeje tip i operacije morajut byti vzete iz okolnoj algebry.\par
\noindent\emph{Reason:} Transformacijne veličiny is the readable current name for the specified algebraic object. Some letters in the earlier record were lost; the readable form is supported by the passage, not a guessed repair to the archive. The expression is $c_S^Tc_T/c_{ST}$, with the condition that all these quantities equal the identity. The international adjective preserves the connection with transformation, while quantity does not merely mean a numerical magnitude. Its type and operations come from the surrounding algebra.\par
\noindent{\footnotesize Izvorne zapisy / Original records: HIST-g211-012.}\par
\smallskip\noindent\emph{Sravnjenje s tekstom:} Čitljiva pasaža ne odstraňaje ošteđenje staršego zapisa.\par
\noindent\emph{Textual comparison:} The readable passage does not erase the earlier record's damage.\par
\noindent{\footnotesize TEXT-COMP-798; pasaže / passages: 24260, 24262, 24264, 24266; \texttt{CORPUS\_FORM\_COMPARISONS.json}.}\par
\end{minipage}\par
\par\medskip\Needspace{7\baselineskip}
\noindent\begin{minipage}{\linewidth}
\hypertarget{TERM-1742}{}
\noindent\textbf{grupa transformacij}\quad{\footnotesize\texttt{TERM-1742}}\par
\noindent\emph{Iztočnik / Source:} Transformationsgruppe\par
\noindent\emph{Rusky / Russian:} группа преобразований\par
\noindent\emph{Ukrajinsky / Ukrainian:} група перетворень\par
\smallskip\noindent\emph{Razlog:} Grupa transformacij daje kratku internacionalnu nazvu, dokolě rusko i ukrajinsko sravnjenja koristajut svoje obyčne imeniki za prěobraženje. Izbor ne dodaje novy matematičny pojem: pasaža zahtěvaje obratnu transformaciju i zatvorenost pri složenju. V tom kontekstu konečna kontinuirna grupa znači zavisimost od konečno mnogo sučstvenyh parametrov, ne konečno čislo transformacij. Pridavniki zato trěba objasniti vměstě s definicijeju, ne prevoditi izolovano po obyčnom smyslu.\par
\noindent\emph{Reason:} Grupa transformacij provides a short international name, while the Russian and Ukrainian comparisons use their ordinary transformation nouns. It adds no new mathematical concept: the passage requires inverses and closure under composition. In this context a finite continuous group depends on finitely many essential parameters; it need not contain finitely many transformations. The modifiers therefore need explanation alongside the definition, not isolated translation by their everyday senses.\par
\noindent{\footnotesize Izvorne zapisy / Original records: HIST-g075-003.}\par
\smallskip\noindent\emph{Sravnjenje s tekstom:} Konečnost parametrov ne znači konečnost množstva transformacij.\par
\noindent\emph{Textual comparison:} Finitely many parameters does not mean finitely many transformations.\par
\noindent{\footnotesize TEXT-COMP-799; pasaže / passages: 9201; \texttt{CORPUS\_FORM\_COMPARISONS.json}.}\par
\end{minipage}\par
\par\medskip\Needspace{7\baselineskip}
\noindent\begin{minipage}{\linewidth}
\hypertarget{TERM-1743}{}
\noindent\textbf{transformovany ideal}\quad{\footnotesize\texttt{TERM-1743}}\par
\noindent\emph{Iztočnik / Source:} Transformed ideal\par
\noindent\emph{Rusky / Russian:} преобразованный идеал\par
\noindent\emph{Ukrajinsky / Ukrainian:} перетворений ідеал\par
\smallskip\noindent\emph{Razlog:} Transformovany ideal jest tehnična nazva, ne samo opis ideala, s ktorym bylo něčto učinjeno. Pasaža dokazuje, že konkretny idealny kvocient ima to svojstvo, prěd priměnjenjem eliminacije. Arhivny naslov Transformed ideal jest anglijsko pojmovno nazvanje, ne njemecka citacija; opakovanje v pozdnějšem spisku ne měnjaje njegov jezyk. Sravnjenje s njemeckym transformiertes Ideal iz drugih zapisov jest opravdano tym samym definiranym objektom. Terminologična konsistencija mora čuvati definiciju, ne samo povtorjati prilastok.\par
\noindent\emph{Reason:} Transformovany ideal is a technical name, not merely an ideal to which something has been done. The passage proves that a particular ideal quotient has the property before applying elimination. The archived heading Transformed ideal is an English concept label, not a German quotation; its repetition in a later list does not change its language. Linking it with German transformiertes Ideal elsewhere is justified by the same defined object. Terminological consistency must preserve the definition, not simply repeat the adjective.\par
\noindent{\footnotesize Izvorne zapisy / Original records: HIST-g091-003, HIST-g092-072.}\par
\smallskip\noindent\emph{Sravnjenje s tekstom:} Tehnično svojstvo jest dokazano prěd eliminacijeju.\par
\noindent\emph{Textual comparison:} The technical property is proved before elimination.\par
\noindent{\footnotesize TEXT-COMP-800; pasaže / passages: 14976, 14992; \texttt{CORPUS\_FORM\_COMPARISONS.json}.}\par
\end{minipage}\par
\par\medskip\Needspace{7\baselineskip}
\noindent\begin{minipage}{\linewidth}
\hypertarget{TERM-1744}{}
\noindent\textbf{transformovane koeficienty}\quad{\footnotesize\texttt{TERM-1744}}\par
\noindent\emph{Iztočnik / Source:} transformierte Koeffizienten\par
\noindent\emph{Rusky / Russian:} преобразованные коэффициенты\par
\noindent\emph{Ukrajinsky / Ukrainian:} перетворені коефіцієнти\par
\smallskip\noindent\emph{Razlog:} Transformovane koeficienty sut koeficienty poslě točno zadanej proměny redov koordinat. Formule (47) opisuju transformaciju, a skobky okolo $(f)$ i $(\varphi)$ razlikujut transformovane koeficienty od početnyh form. Arhiv predpočita internacionalny korenj transform- dlja svezi s toj operacijeju. On ne tvrdi, že rusko preobrazovannye ili ukrajinsko peretvoreni sut odbačene kako nepravilne; to sut nacionalne sravniteljne formy. Ograničenje na danu proměnu čuva smysl participa.\par
\noindent\emph{Reason:} Transformovane koeficienty are coefficients after the specified transformation of coordinate rows. Equations (47) describe the change, and parentheses in $(f)$ and $(\varphi)$ distinguish transformed coefficients from the original forms. The record prefers international transform- to connect the word with that operation. It does not reject Russian preobrazovannye or Ukrainian peretvoreni as incorrect; they are national-language comparisons. The specified transformation gives the participle its precise scope.\par
\noindent{\footnotesize Izvorne zapisy / Original records: HIST-g036-004.}\par
\smallskip\noindent\emph{Sravnjenje s tekstom:} Skobky razlikujut koeficienty od početnyh form.\par
\noindent\emph{Textual comparison:} Parentheses distinguish coefficients from the original forms.\par
\noindent{\footnotesize TEXT-COMP-801; pasaže / passages: 4464, 4466, 4469, 4470, 4474; \texttt{CORPUS\_FORM\_COMPARISONS.json}.}\par
\end{minipage}\par
\par\medskip\Needspace{7\baselineskip}
\noindent\begin{minipage}{\linewidth}
\hypertarget{TERM-1745}{}
\noindent\textbf{transformovany ideal}\quad{\footnotesize\texttt{TERM-1745}}\par
\noindent\emph{Iztočnik / Source:} transformiertes Ideal\par
\noindent\emph{Rusky / Russian:} преобразованный идеал\par
\noindent\emph{Ukrajinsky / Ukrainian:} перетворений ідеал\par
\smallskip\noindent\emph{Razlog:} Transformovany ideal jest ponovljen čerez několiko oddělov, zato že daljše dokazy opirajut se na istu definiciju. V članku o eliminaciji ideal nastava prěz zadano linearno prěobraženje i prisojedinjenje neoprěděljenyh koeficientov k polju. Ne govori se o proizvoljnom obrazu pod kakymkoli odobraženjem. Pasivny prilastok jest čitljiv i kratok, ale mora ostati svezany s toj konstrukcijeju. Različne formule v dvou člankah ne dolžny byti nasilno imenovane jednim čislom v slovniku.\par
\noindent\emph{Reason:} Transformovany ideal recurs across several sections because later proofs depend on the same definition. In the elimination paper it is obtained by the specified linear transformation and adjunction of indeterminate coefficients to the field. It is not an arbitrary image under any map. The passive adjective is readable and concise, but must stay tied to that construction. Different equation numbers in the two papers must not be forced into a single reference in the dictionary.\par
\noindent{\footnotesize Izvorne zapisy / Original records: HIST-g087-005, HIST-g090-006, HIST-g092-025, HIST-g092-060, HIST-g093-012, HIST-g096-001.}\par
\smallskip\noindent\emph{Sravnjenje s tekstom:} Prilastok upučaje na zadanu konstrukciju.\par
\noindent\emph{Textual comparison:} The adjective points to a specified construction.\par
\noindent{\footnotesize TEXT-COMP-802; pasaže / passages: 14539, 14541, 15273, 15283, 15285, 15293; \texttt{CORPUS\_FORM\_COMPARISONS.json}.}\par
\end{minipage}\par
\par\medskip\Needspace{7\baselineskip}
\noindent\begin{minipage}{\linewidth}
\hypertarget{TERM-1746}{}
\noindent\textbf{stepenj transcendencije}\quad{\footnotesize\texttt{TERM-1746}}\par
\noindent\emph{Iztočnik / Source:} Transzendenzgrad\par
\noindent\emph{Rusky / Russian:} степень трансцендентности\par
\noindent\emph{Ukrajinsky / Ukrainian:} степінь трансцендентності\par
\smallskip\noindent\emph{Razlog:} Stepenj transcendencije vyražaje čislo algebraično nezavisnyh elementov nad imenovanym osnovnym poljem. V pasaži $k$ takyh elementov može byti najdeno, ale vsake $k+1$ sut zavisne. To ne jest obyčny algebraičny stepenj razširjenja niti stepenj pojedinogo polinoma. Imenična konstrukcija slěduje sravnjenju s ruskim i ukrajinskim izrazami i sohranjaje internacionalnu osnovu transcendenc-. Prednost jest pojmovna prěhlednost, ne dokumentovano čislovo merjenje razumělivosti.\par
\noindent\emph{Reason:} Stepenj transcendencije expresses the number of algebraically independent elements over the specified base field. The passage allows $k$ such elements but makes every $k+1$ dependent. This is not the ordinary algebraic degree of an extension or the degree of a polynomial. The noun construction follows the Russian and Ukrainian comparisons while retaining international transcendenc-. Its stated advantage is conceptual clarity, not a documented numerical measure of intelligibility.\par
\noindent{\footnotesize Izvorne zapisy / Original records: HIST-g102-009.}\par
\smallskip\noindent\emph{Sravnjenje s tekstom:} Algebraična nezavisnost jest otnosna k osnovnomu polju.\par
\noindent\emph{Textual comparison:} Algebraic independence is relative to the base field.\par
\noindent{\footnotesize TEXT-COMP-803; pasaže / passages: 15839, 15841; \texttt{CORPUS\_FORM\_COMPARISONS.json}.}\par
\end{minipage}\par
\par\medskip\Needspace{7\baselineskip}
\noindent\begin{minipage}{\linewidth}
\hypertarget{TERM-1747}{}
\noindent\textbf{transcendencijny stepen}\quad{\footnotesize\texttt{TERM-1747}}\par
\noindent\emph{Iztočnik / Source:} Transzendenzgrad\par
\noindent\emph{Rusky / Russian:} степень трансцендентности\par
\noindent\emph{Ukrajinsky / Ukrainian:} трансцендентний степінь\par
\smallskip\noindent\emph{Razlog:} Transcendencijny stepen jest prilastkova varianta stepenja transcendencije, ne nova invarianta. Pasaža definije smysl za sistem funkcij i povezuje smysl za prosty ideal s njegovym ostatkovym poljem. V dokazu pri redukciji modulo prvočislo razlučaje se poniženje togo stepena od gubitka prostosti ideala. Izpustiti prilastok ili zaměniti jeho samym čislom prěměnnyh moglo by spojiti te dva slučaje. Slovnik mora navesti variantu pod obščim pojmovnym naslovom.\par
\noindent\emph{Reason:} Transcendencijny stepen is an adjectival variant of stepenj transcendencije, not a new invariant. The passage defines it for a system of functions and relates its use for a prime ideal to the corresponding residue field. In the reduction-modulo-prime argument, a drop in transcendence degree is distinguished from loss of primality. Omitting the modifier or replacing the invariant with the number of variables could collapse those cases. The variant belongs under a shared conceptual heading.\par
\noindent{\footnotesize Izvorne zapisy / Original records: HIST-g205-009.}\par
\smallskip\noindent\emph{Sravnjenje s tekstom:} Poniženje stepena i gubitok prostosti sut različne možnosti.\par
\noindent\emph{Textual comparison:} A degree drop and loss of primality are distinct possibilities.\par
\noindent{\footnotesize TEXT-COMP-804; pasaže / passages: 22535, 22587; \texttt{CORPUS\_FORM\_COMPARISONS.json}.}\par
\end{minipage}\par
\par\medskip\Needspace{7\baselineskip}
\noindent\begin{minipage}{\linewidth}
\hypertarget{TERM-1748}{}
\noindent\textbf{transcendentny stepen / medžupolje}\quad{\footnotesize\texttt{TERM-1748}}\par
\noindent\emph{Iztočnik / Source:} Transzendenzgrad / Zwischenkörper\par
\noindent\emph{Rusky / Russian:} степень трансцендентности / промежуточное поле\par
\noindent\emph{Ukrajinsky / Ukrainian:} трансцендентний степінь / проміжне поле\par
\smallskip\noindent\emph{Razlog:} Transcendentny stepen i medžupolje imajut različne ulogi: prvo jest invarianta, drugo opisuje vloženje medžu dvoma poljami. Medžu- čini relaciju prěhlednu, no ne zahtěvaje, da obě inkluzije budut stroge. Pasaža daje $P\subseteq K\subseteq P(x_1,\ldots,x_n)$ i fiksuje stepenj transcendencije prěd zaključkom o konečnom algebraičnom razširjenju. To zaključivanje ne slěduje iz samogo slova medžupolje. Arhivna kosá črta sobiraje dva termina iz jednogo argumenta, ne dva sinonima.\par
\noindent\emph{Reason:} Transcendentny stepen and medžupolje play different roles: the first is an invariant; the second describes inclusion between fields. Medžu- makes the relation transparent without requiring both inclusions to be strict. The passage places $P\subseteq K\subseteq P(x_1,\ldots,x_n)$ and specifies the transcendence degree before concluding finiteness of an algebraic extension. That conclusion does not follow merely from the word intermediate field. The archival slash groups two terms from one argument, not two synonyms.\par
\noindent{\footnotesize Izvorne zapisy / Original records: HIST-g106-008.}\par
\smallskip\noindent\emph{Sravnjenje s tekstom:} Oba termina sodělujut v argumentu bez synonymiji.\par
\noindent\emph{Textual comparison:} The terms work together in the argument without being synonyms.\par
\noindent{\footnotesize TEXT-COMP-805; pasaže / passages: 15997, 16001; \texttt{CORPUS\_FORM\_COMPARISONS.json}.}\par
\end{minipage}\par
\par\medskip\Needspace{7\baselineskip}
\noindent\begin{minipage}{\linewidth}
\hypertarget{TERM-1752}{}
\noindent\textbf{trivialno predstavjenje}\quad{\footnotesize\texttt{TERM-1752}}\par
\noindent\emph{Iztočnik / Source:} triviale Darstellung\par
\noindent\emph{Rusky / Russian:} тривиальное представление\par
\noindent\emph{Ukrajinsky / Ukrainian:} тривіальне зображення\par
\smallskip\noindent\emph{Razlog:} Trivialno prědstavjenje tut označaje nulovo prědstavjenje: pasaža izrično priděljaje nulu vsakomu elementu. To ne jest obična grupova konvencija, pri ktorej vsaky element děluje identično. Odštěpenje nulovogo sumanda potom ostavjaje prědstavjenje, v ktorom jedinica děluje jediničnoju matricoju. Internacionalny prilastok jest jezyčno poznany, ale sam po sobě ne ustanovjaje matematičny smysl. Priložena kratka glosa nula zabrani važnomu nedorazuměnju bez zaměny avtorskogo nazvanja.\par
\noindent\emph{Reason:} Trivialno prědstavjenje here needs the gloss zero representation: the passage explicitly assigns zero to every element. This is not the usual group-theoretic convention in which every element acts as the identity. Splitting off the zero summand leaves a representation in which the identity acts by the identity matrix. The international adjective is linguistically familiar but does not determine the mathematical sense on its own. A short zero gloss prevents a serious misunderstanding without replacing the author's name.\par
\noindent{\footnotesize Izvorne zapisy / Original records: HIST-g195-005.}\par
\smallskip\noindent\emph{Sravnjenje s tekstom:} Izrična rečenica ustanovjaje specialny smysl.\par
\noindent\emph{Textual comparison:} The explicit statement determines the particular sense.\par
\noindent{\footnotesize TEXT-COMP-809; pasaže / passages: 21605, 21607; \texttt{CORPUS\_FORM\_COMPARISONS.json}.}\par
\end{minipage}\par
\par\medskip\Needspace{7\baselineskip}
\noindent\begin{minipage}{\linewidth}
\hypertarget{TERM-1753}{}
\noindent\textbf{formy inercije}\quad{\footnotesize\texttt{TERM-1753}}\par
\noindent\emph{Iztočnik / Source:} Trägheitsformen\par
\noindent\emph{Rusky / Russian:} формы инерции\par
\noindent\emph{Ukrajinsky / Ukrainian:} форми інерції\par
\smallskip\noindent\emph{Razlog:} Formy inercije jest blizky prěvod Trägheitsformen v kratkoj primětce. Članok upominaje priměnjenje formuly v disertaciji, ale ne razviva teoriju togo objekta. Arhiv zato zachovuje doslovnu nazvu i ne ustanovjaje podrobnu definiciju. Ne trěba dopolnjati zapis fizikalnym smyslom inercije ili poistovětiti njegove formy s inercijnymi poljami iz čislovoj teorije. Dodatno pojmovno utočnjenje potrěbovalo by samostojno pročitanogo primarnogo iztočnika, ne samo njemeckogo slova.\par
\noindent\emph{Reason:} Formy inercije closely renders Trägheitsformen in a brief footnote. The paper mentions an application of the formula in a dissertation but does not develop the object. The archive therefore retains the literal name without supplying a detailed definition. The entry must not be filled out with physical inertia or identified with the inertia fields of number theory. A more specific conceptual explanation would require a separately consulted primary source, not merely the German word.\par
\noindent{\footnotesize Izvorne zapisy / Original records: HIST-g078-014.}\par
\smallskip\noindent\emph{Sravnjenje s tekstom:} Upoměnka ne ustanovjaje polnu teoriju objekta.\par
\noindent\emph{Textual comparison:} The mention does not supply a full theory of the object.\par
\noindent{\footnotesize TEXT-COMP-810; pasaže / passages: 10816, 10818, 10819, 10822, 10823; \texttt{CORPUS\_FORM\_COMPARISONS.json}.}\par
\end{minipage}\par
\par\medskip\Needspace{7\baselineskip}
\noindent\begin{minipage}{\linewidth}
\hypertarget{TERM-1754}{}
\noindent\textbf{inercijno polje / razvětvenno polje}\quad{\footnotesize\texttt{TERM-1754}}\par
\noindent\emph{Iztočnik / Source:} Trägheitskörper / Verzweigungskörper\par
\noindent\emph{Rusky / Russian:} поле инерции / поле ветвления\par
\noindent\emph{Ukrajinsky / Ukrainian:} поле інерції / поле розгалуження\par
\smallskip\noindent\emph{Razlog:} Inercijno polje i razvětvenno polje pripadajut čislovo-teoretičnomu kontekstu razvětvenja. Arhiv sohranjaje dva različne pojmy, koristajuči internacionalnu inercija i slovjansku osnovu razvětven-. Pasaža ukazuje na osnovny iztočnik i kontrast s teorijeju algebraičnyh funkcij v danom izloženju. To ne daje prava identifikovati oba polja medžu soboju niti prenesti avtorske ograničenje na vse sovrěmenne teorije funkcij. Nazvy sut historijno uměščene; polna moderna definicija ne jest dobyta iz te kratkoj upoměnky.\par
\noindent\emph{Reason:} Inercijno polje and razvětvenno polje belong to the number-theoretic discussion of ramification. The record keeps two distinct concepts, using international inercija and Slavic razvětven-. The passage cites its source and contrasts the situation with algebraic function theory as treated there. This does not identify the fields with one another or extend the historical restriction to every modern function-field theory. The names are historically situated; the brief mention does not supply a complete modern definition.\par
\noindent{\footnotesize Izvorne zapisy / Original records: HIST-g076-013.}\par
\smallskip\noindent\emph{Sravnjenje s tekstom:} Nazvy ne sut dovoljne za neograničeno moderno obobčenje.\par
\noindent\emph{Textual comparison:} The names do not license an unrestricted modern generalization.\par
\noindent{\footnotesize TEXT-COMP-811; pasaže / passages: 9922, 9943; \texttt{CORPUS\_FORM\_COMPARISONS.json}.}\par
\end{minipage}\par
\par\medskip\Needspace{7\baselineskip}
\noindent\begin{minipage}{\linewidth}
\hypertarget{TERM-1755}{}
\noindent\textbf{Čirnhausova transformacija}\quad{\footnotesize\texttt{TERM-1755}}\par
\noindent\emph{Iztočnik / Source:} Tschirnhaussche Transformation\par
\noindent\emph{Rusky / Russian:} преобразование Чирнгауза\par
\noindent\emph{Ukrajinsky / Ukrainian:} перетворення Чірнгауза\par
\smallskip\noindent\emph{Razlog:} Čirnhausova transformacija prilagođaje eponim slovjanskomu čitanju, ne prevodi ime kakoby obyčno slovo. Arhiv odbacaje mehanično prěpisovanje njemeckogo Tsch- v kirilicu. Nynějšna pasaža uživa transformacija Čirnhausa, to jest isti eponim v genitivnoj konstrukciji, i označaje odstranenje drugogo najvyššego člena. Ta sintaktična varianta ne ustanovjaje novu operaciju. Originalno Tschirnhaus mora ostati nahodimo v iztočnikah i ukazatelju, aby prilagodženje ne prekinulo identifikaciju člověka.\par
\noindent\emph{Reason:} Čirnhausova transformacija adapts the eponym for Slavic reading rather than translating the name as a common word. The record rejects mechanical conversion of German Tsch- into Cyrillic. The current passage uses transformacija Čirnhausa, a genitive construction with the same name, for elimination of the next-to-highest term. The syntactic variant does not define a new operation. Original Tschirnhaus should remain findable in references and the index so adaptation does not break identification of the person.\par
\noindent{\footnotesize Izvorne zapisy / Original records: HIST-g073-018.}\par
\smallskip\noindent\emph{Sravnjenje s tekstom:} Prilastok i genitiv ukazujut na istogo člověka.\par
\noindent\emph{Textual comparison:} Adjective and genitive refer to the same person.\par
\noindent{\footnotesize TEXT-COMP-812; pasaže / passages: 8841; \texttt{CORPUS\_FORM\_COMPARISONS.json}.}\par
\end{minipage}\par
\par\medskip\Needspace{7\baselineskip}
\noindent\begin{minipage}{\linewidth}
\hypertarget{TERM-1757}{}
\noindent\textbf{okolina}\quad{\footnotesize\texttt{TERM-1757}}\par
\noindent\emph{Iztočnik / Source:} Umgebung\par
\noindent\emph{Rusky / Russian:} окрестность\par
\noindent\emph{Ukrajinsky / Ukrainian:} окіл\par
\smallskip\noindent\emph{Razlog:} Okolina jest prosta nazva dlja neighborhood, no v pasaži ima strogo lokalny smysl. Tvrdženje o ekstremnoj diskontinuirnosti govori o cěnnostjah v okolině zadanej točky i o libovoljnoj blizkosti k naprijed zadanoj cěnnosti. Ne stačaje izloženje kako funkcija někde može byti velika. Arhiv želi čitljivu prozu bez novogo topologičnogo zajętka; to ne znači upuščanje kvantifikatorov. Blizkost argumentov i blizkost cěnnostij morajut ostati različne.\par
\noindent\emph{Reason:} Okolina is a plain name for neighborhood, but has a precise local role in the passage. Extreme discontinuity concerns values in a neighborhood of a specified point and arbitrary closeness to a prescribed value. Saying the function can be large somewhere would not suffice. The record seeks readable prose without a new topological borrowing, not loss of quantifiers. Closeness of arguments and closeness of values must remain distinct.\par
\noindent{\footnotesize Izvorne zapisy / Original records: HIST-g072-002.}\par
\smallskip\noindent\emph{Sravnjenje s tekstom:} Lokalnost argumenta i cěnnosti ne možno směšati.\par
\noindent\emph{Textual comparison:} Locality of argument and value must not be conflated.\par
\noindent{\footnotesize TEXT-COMP-814; pasaže / passages: 8601, 8640; \texttt{CORPUS\_FORM\_COMPARISONS.json}.}\par
\end{minipage}\par
\par\medskip\Needspace{7\baselineskip}
\noindent\begin{minipage}{\linewidth}
\hypertarget{TERM-1758}{}
\noindent\textbf{obratno-jednoznačna funkcija}\quad{\footnotesize\texttt{TERM-1758}}\par
\noindent\emph{Iztočnik / Source:} umkehrbar eindeutig\par
\noindent\emph{Rusky / Russian:} взаимно однозначная функция\par
\noindent\emph{Ukrajinsky / Ukrainian:} взаємно однозначна функція\par
\smallskip\noindent\emph{Razlog:} Obratno-jednoznačna funkcija pojmenovuje točno dokazany krok: iz $f(x)=f(y)$ slěduje $x=y$. To jest injektivnost, odnosno jednoznačnost obratnogo pridruženja na obrazu. Pasaža potom govori o vzajemno jednoznačnoj svezi s poljem obrazov; ne tvrdi bez dodatka surjektivnost na proizvoljno večše polje. Arhivny složeny prilastok čuvaje tu logičnu poslědovatelnost. Sovrěmenna glosa injektivna jest koristna, dokolě ostaje ukazano, na ktorym množstvu definovano jest obratno pridruženje.\par
\noindent\emph{Reason:} Obratno-jednoznačna funkcija names the precise step proved: $f(x)=f(y)$ implies $x=y$. This is injectivity, or single-valuedness of the reverse correspondence on the image. The passage then describes a one-to-one correspondence with the image field, not unqualified surjectivity onto an arbitrarily larger field. The compound preserves that logical sequence. The modern gloss injective is helpful provided the domain of the inverse correspondence remains explicit.\par
\noindent{\footnotesize Izvorne zapisy / Original records: HIST-g069-004.}\par
\smallskip\noindent\emph{Sravnjenje s tekstom:} Obratno pridruženje ima jasno oprěděljenu oblast.\par
\noindent\emph{Textual comparison:} The reverse correspondence has a specified domain.\par
\noindent{\footnotesize TEXT-COMP-815; pasaže / passages: 8434, 8436, 8448, 8450, 8452; \texttt{CORPUS\_FORM\_COMPARISONS.json}.}\par
\end{minipage}\par
\par\medskip\Needspace{7\baselineskip}
\noindent\begin{minipage}{\linewidth}
\hypertarget{TERM-1759}{}
\noindent\textbf{obratno jednoznačny / vzaimno jednoznačny}\quad{\footnotesize\texttt{TERM-1759}}\par
\noindent\emph{Iztočnik / Source:} umkehrbar eindeutig / eineindeutig\par
\noindent\emph{Rusky / Russian:} обратимо однозначный / взаимно однозначный\par
\noindent\emph{Ukrajinsky / Ukrainian:} обернено однозначний / взаємно однозначний\par
\smallskip\noindent\emph{Razlog:} Obratno jednoznačny i vzaimno jednoznačny slědujut ulogam pridruženij v dokazu. Definicija homomorfije v toj pasaži už zahtěvaje izčerpanje ciljevoj oblasti; dodatna obratna jednoznačnost zato daje izomorfiju. V pridruženju klasov ostatkov imenno klasy, ne vsi posamezne prědstavitelji, odgovarjajut jednoznačno. Tymy slovami ne možno zaměnjati razliku medžu funkcijeju, bijekcijeju i vyborom prědstavitelja. Arhiv sohranjaje korespondencijny jezyk vměsto jednogo neobjasnenogo modernogo labela.\par
\noindent\emph{Reason:} Obratno jednoznačny and vzaimno jednoznačny follow the roles of the correspondences in the argument. The passage's definition of homomorphism already requires exhaustion of the target; adding inverse uniqueness therefore gives an isomorphism. In the residue-class correspondence, the classes, not all their individual representatives, correspond uniquely. The wording must not blur function, bijection, and selection of representatives. The archive retains correspondence language instead of an unexplained modern label.\par
\noindent{\footnotesize Izvorne zapisy / Original records: HIST-g111-008.}\par
\smallskip\noindent\emph{Sravnjenje s tekstom:} Prědpoloženo izčerpanje cilja mora ostati vidimo.\par
\noindent\emph{Textual comparison:} The assumed exhaustion of the target must remain visible.\par
\noindent{\footnotesize TEXT-COMP-816; pasaže / passages: 16455, 16457, 16460, 16462, 16480; \texttt{CORPUS\_FORM\_COMPARISONS.json}.}\par
\end{minipage}\par
\par\medskip\Needspace{7\baselineskip}
\noindent\begin{minipage}{\linewidth}
\hypertarget{TERM-1760}{}
\noindent\textbf{obrat}\quad{\footnotesize\texttt{TERM-1760}}\par
\noindent\emph{Iztočnik / Source:} Umkehrung\par
\noindent\emph{Rusky / Russian:} обратное утверждение / обращение\par
\noindent\emph{Ukrajinsky / Ukrainian:} обернення\par
\smallskip\noindent\emph{Razlog:} Obrat jest kračša nazva dlja obratnogo tvrdženja teorema. V oddělah o konečnyh i beskonečnyh grupah ona znači vyvod invariantnosti iz divergencijnyh relacij ili zavisimostij, ne vzjatje obratnoj grupovoj transformacije. Oba smysla obratanja pojavjajut se v blizkoj proze, zato slovnik mora razlikovati jih. Polna fraza obratny teorem ili obratno tvrdženje može služiti glosoju. Kratkost opravdana jest samo togda, kogda jasny sut hipotezy i směr implikacije.\par
\noindent\emph{Reason:} Obrat is the short name for a theorem's converse. In the sections on finite and infinite groups it means deriving invariance from divergence relations or dependencies, not taking an inverse group transformation. Both kinds of reversal occur nearby, so the dictionary must distinguish them. Converse theorem or converse statement is an appropriate expanded gloss. The short form works only when the hypotheses and direction of implication remain clear.\par
\noindent{\footnotesize Izvorne zapisy / Original records: HIST-g075-012.}\par
\smallskip\noindent\emph{Sravnjenje s tekstom:} Směr implikacije razlikuje se od obratnoj transformacije.\par
\noindent\emph{Textual comparison:} Reversing implication is distinct from inverting a transformation.\par
\noindent{\footnotesize TEXT-COMP-817; pasaže / passages: 9414, 9416, 9431, 9481, 9518; \texttt{CORPUS\_FORM\_COMPARISONS.json}.}\par
\end{minipage}\par
\par\medskip\Needspace{7\baselineskip}
\noindent\begin{minipage}{\linewidth}
\hypertarget{TERM-1761}{}
\noindent\textbf{obrat sistema ravnanj}\quad{\footnotesize\texttt{TERM-1761}}\par
\noindent\emph{Iztočnik / Source:} Umkehrung des Gleichungssystems\par
\noindent\emph{Rusky / Russian:} обращение системы уравнений\par
\noindent\emph{Ukrajinsky / Ukrainian:} обертання системи рівнянь\par
\smallskip\noindent\emph{Razlog:} Obrat sistema ravnanj označaje rěšenje obratenogo linearnogo odnosa: polarny členy najprvo izražajut se prěz druge formy, potom iz tih ravnanij dobyvajut se iskane relacije. Kratko obrat sohranjaje algebraičnu operaciju, ne obratnu implikaciju teorema. Ruska i ukrajinska fraza v arhivu podpirajut to smyslovo sravnjenje, ale same ne dokazujut obščeslovjansku ustaljenost. Polna glosa jest rěšenje sistema v obratnom směru; ona ne tvrdi, že libovoljny sistem ima jedino rěšenje.\par
\noindent\emph{Reason:} Obrat sistema ravnanj denotes solving the linear relation in the reverse direction: polar terms are first expressed through other forms, and the desired relations are then obtained from those equations. The short obrat preserves an algebraic operation, not a theorem's converse implication. The archived Russian and Ukrainian phrases support that semantic comparison, but do not themselves establish pan-Slavic usage. Solving the system in the reverse direction is the expanded gloss; it does not assert that an arbitrary system has a unique solution.\par
\noindent{\footnotesize Izvorne zapisy / Original records: HIST-g006-006.}\par
\smallskip\noindent\emph{Sravnjenje s tekstom:} Polary i linearne kombinacije oprěděljajut smysl slova obrat.\par
\noindent\emph{Textual comparison:} The polars and linear combinations determine the sense of obrat.\par
\noindent{\footnotesize TEXT-COMP-818; pasaže / passages: 713, 715, 716, 717, 719; \texttt{CORPUS\_FORM\_COMPARISONS.json}.}\par
\end{minipage}\par
\par\medskip\Needspace{7\baselineskip}
\noindent\begin{minipage}{\linewidth}
\hypertarget{TERM-1762}{}
\noindent\textbf{obrat (1) po (1')}\quad{\footnotesize\texttt{TERM-1762}}\par
\noindent\emph{Iztočnik / Source:} Umkehrung von (1) nach (1')\par
\noindent\emph{Rusky / Russian:} обращение (1) по (1')\par
\noindent\emph{Ukrajinsky / Ukrainian:} обернення (1) за (1')\par
\smallskip\noindent\emph{Razlog:} Obrat (1) po (1') upominaje konkretno obratno linearno prěobraženje. Formula (1) vyražaje y prěz x, a (1') x prěz y; koeficienty oběh prěobraženij razširjajut polje do togo samogo $P(u)=P(t)$. Zato izbor obrat trěba čitati s oběma formulami. Ne jest to obratny teorem ni kategorična definicija obratnogo morfizma. Kračša fraza slěduje zapisanomu naměru ne zatěžati dokaz novym dolgim terminom; čislena ocěna toj prednosti ne jest sohranjena.\par
\noindent\emph{Reason:} Obrat (1) po (1') refers to the particular inverse linear transformation. Equation (1) expresses y through x, while (1') expresses x through y; their coefficients extend the field to the same $P(u)=P(t)$. The choice obrat must therefore be read with both formulas. It is neither a converse theorem nor a categorical definition of inverse morphism. The short phrase follows the recorded aim of avoiding a cumbersome new term in proof prose; no numerical assessment of that advantage survives.\par
\noindent{\footnotesize Izvorne zapisy / Original records: HIST-g096-003.}\par
\smallskip\noindent\emph{Sravnjenje s tekstom:} Formula (1') daje konkretny smysl obratanju (1).\par
\noindent\emph{Textual comparison:} Equation (1') gives the concrete meaning of inverting (1).\par
\noindent{\footnotesize TEXT-COMP-819; pasaže / passages: 15275, 15276, 15278, 15280, 15281, 15283, 15374; \texttt{CORPUS\_FORM\_COMPARISONS.json}.}\par
\end{minipage}\par
\par\medskip\Needspace{7\baselineskip}
\noindent\begin{minipage}{\linewidth}
\hypertarget{TERM-1763}{}
\noindent\textbf{obrat lemy I}\quad{\footnotesize\texttt{TERM-1763}}\par
\noindent\emph{Iztočnik / Source:} Umkehrung von Hilfssatz I\par
\noindent\emph{Rusky / Russian:} обращение леммы I\par
\noindent\emph{Ukrajinsky / Ukrainian:} обернення леми I\par
\smallskip\noindent\emph{Razlog:} Obrat lemy I znači obratno logično tvrdženje. Lema I prěnosi redukovanost iz komponentnyh prědstavjenij na jih razloženje; lema IV tvrdi, že grupovanje komponentov redukovanogo prědstavjenja opět daje redukovano prědstavjenje. Ravnobrojnost komponentov služi v dokazu togo obratnogo směra. Obrat ne označaje ni obratnu matricu ni zaměnu ideala njegovym recipročnym objektom. Zachovanje nazvy lemy i rimskogo I pozvaljaje čitatelju slědovati po originalnom argumentu.\par
\noindent\emph{Reason:} Obrat lemy I means the converse logical statement. Lemma I passes reducedness from component representations to their refinement; Lemma IV says that grouping components of a reduced representation again yields a reduced representation. Equality of the number of components contributes to the converse argument. Obrat denotes neither an inverse matrix nor replacement of an ideal by a reciprocal object. Retaining the lemma name and Roman I lets the reader follow the original argument.\par
\noindent{\footnotesize Izvorne zapisy / Original records: HIST-g081-038, HIST-g082-038.}\par
\smallskip\noindent\emph{Sravnjenje s tekstom:} Dva směry izkazu o redukovanosti ne sut matrično obratanje.\par
\noindent\emph{Textual comparison:} The two directions of the reducedness statement are not matrix inversion.\par
\noindent{\footnotesize TEXT-COMP-820; pasaže / passages: 12588, 12594, 12598, 12739, 12741, 12745; \texttt{CORPUS\_FORM\_COMPARISONS.json}.}\par
\end{minipage}\par
\par\medskip\Needspace{7\baselineskip}
\noindent\begin{minipage}{\linewidth}
\hypertarget{TERM-1765}{}
\noindent\textbf{nezavisna modulna baza}\quad{\footnotesize\texttt{TERM-1765}}\par
\noindent\emph{Iztočnik / Source:} unabhängige Modulbasis\par
\noindent\emph{Rusky / Russian:} независимая модульная база\par
\noindent\emph{Ukrajinsky / Ukrainian:} незалежна модульна база\par
\smallskip\noindent\emph{Razlog:} Nezavisna modulna baza sohranjaje izričny uslov nezavisnosti, ktory v historijnoj upotrěbě slova baza ne jest izbytočny. V aktualnoj pasaži stoji nezavisima: iz summy $\sum_i h_i t_i=0$ mora slědovati $h_i=0$ za vsako i. Navadna modulna baza tam prvo znači sistem generatorov; pozdněje dopuščen jest konečny, ne potrěbno nezavisimy sistem. Zato ne možno opustiti pridavnik ili smatrjati jeho za stilističnu ukrasu. Arhivna nezavisna i tekstova nezavisima sut formy dlja sravnjenja, ne dva različne matematične pojmy.\par
\noindent\emph{Reason:} Nezavisna modulna baza preserves an explicit independence condition which is not redundant in the historical use of basis. The current passage has nezavisima: a sum $\sum_i h_i t_i=0$ must imply $h_i=0$ for every i. An ordinary module basis there initially means a generating system; a later passage permits a finite, not necessarily independent one. The adjective cannot be omitted as stylistic decoration. Archived nezavisna and textual nezavisima are forms to compare, not two different mathematical concepts.\par
\noindent{\footnotesize Izvorne zapisy / Original records: HIST-g213-008.}\par
\smallskip\noindent\emph{Sravnjenje s tekstom:} Izrična definicija razděljaje generovanje od nezavisnosti.\par
\noindent\emph{Textual comparison:} The explicit definition separates generation from independence.\par
\noindent{\footnotesize TEXT-COMP-822; pasaže / passages: 24826, 24828, 24829, 24831, 25059; \texttt{CORPUS\_FORM\_COMPARISONS.json}.}\par
\end{minipage}\par
\par\medskip\Needspace{7\baselineskip}
\noindent\begin{minipage}{\linewidth}
\hypertarget{TERM-1766}{}
\noindent\textbf{neopreděljena}\quad{\footnotesize\texttt{TERM-1766}}\par
\noindent\emph{Iztočnik / Source:} Unbestimmte\par
\noindent\emph{Rusky / Russian:} неопределенная\par
\noindent\emph{Ukrajinsky / Ukrainian:} невизначена\par
\smallskip\noindent\emph{Razlog:} Neopreděljena jest substantivizovany pridavnik za algebraičnu neoprěděljenu veličinu. V polju racionalnyh funkcij ona označaje formalny argument, ne neznanu čislenu veličinu, ktoru trěba najdti iz ravnanja. To jest zapisana pričina ne zaměnjati vsako Unbestimmte občim proměnna. Aktualny tekst piše neoprěděljena i njegove padežne formy. Razlika v pravopisu ne dokazuje noveho smysla; slovnik dolžen svjazati varianty, ostavljajuči vidljivu razliku medžu formalnym argumentom i neizvěstnoju.\par
\noindent\emph{Reason:} Neopreděljena is a substantivized adjective for an algebraic indeterminate. In a rational function field it denotes a formal argument, not an unknown numerical quantity to be found by solving an equation. This is the recorded reason not to replace every Unbestimmte with the general variable. The current text spells it neoprěděljena and inflects it accordingly. The spelling difference does not establish a new sense; the dictionary should link the variants while preserving the distinction between a formal argument and an unknown.\par
\noindent{\footnotesize Izvorne zapisy / Original records: HIST-g039-002, HIST-g090-003, HIST-g092-057.}\par
\smallskip\noindent\emph{Sravnjenje s tekstom:} Aktualna forma neoprěděljena ima algebraičnu ulogu, ne samo svakodenne značenje.\par
\noindent\emph{Textual comparison:} Current neoprěděljena has an algebraic role, not merely an everyday meaning.\par
\noindent{\footnotesize TEXT-COMP-823; pasaže / passages: 4950, 4952, 14226; \texttt{CORPUS\_FORM\_COMPARISONS.json}.}\par
\end{minipage}\par
\par\medskip\Needspace{7\baselineskip}
\noindent\begin{minipage}{\linewidth}
\hypertarget{TERM-1767}{}
\noindent\textbf{neopreděljene}\quad{\footnotesize\texttt{TERM-1767}}\par
\noindent\emph{Iztočnik / Source:} Unbestimmte\par
\noindent\emph{Rusky / Russian:} неопределенные; неопределённые\par
\noindent\emph{Ukrajinsky / Ukrainian:} невизначені\par
\smallskip\noindent\emph{Razlog:} Neopreděljene jest množinska forma togo samogo pojma. Dva priměry pojasnjajut njegovo značenje: determinanty zaměnjajut se novymi simbolami p, a algebraično nezavisne bazove čisla $\eta$ možno razgledati kako neoprěděljene. V drugom slučaju čisla ne stavajut nerěšenymi neizvěstnymi; njih nezavisnost dopušča funkcionno razgledanje. Ruski i ukrajinski komparandy izražajut istu razliku. Arhivne primětky o prověrci participialnoj formy ostajut lingvističnoju ogradoju, ne dokazom, že forma už jest potvrđena ekspertom.\par
\noindent\emph{Reason:} Neopreděljene is the plural of the same concept. Two examples clarify its use: determinants are replaced by new symbols p, and algebraically independent basis numbers $\eta$ can be treated as indeterminates. In the latter case the numbers do not become unsolved unknowns; their independence permits the functional treatment. The Russian and Ukrainian comparanda express the same distinction. Archived requests to review the participial form remain a linguistic reservation, not evidence of expert approval.\par
\noindent{\footnotesize Izvorne zapisy / Original records: HIST-g041-003, HIST-g057-011, HIST-g058-005.}\par
\smallskip\noindent\emph{Sravnjenje s tekstom:} Oba priměry pokazyvajut, začto neizvěstne ne jest dosta točna glosa.\par
\noindent\emph{Textual comparison:} Both examples show why unknowns is not an adequate gloss.\par
\noindent{\footnotesize TEXT-COMP-824; pasaže / passages: 7101, 7142; \texttt{CORPUS\_FORM\_COMPARISONS.json}.}\par
\end{minipage}\par
\par\medskip\Needspace{7\baselineskip}
\noindent\begin{minipage}{\linewidth}
\hypertarget{TERM-1768}{}
\noindent\textbf{neopreděljene / stati se neopreděljenoju}\quad{\footnotesize\texttt{TERM-1768}}\par
\noindent\emph{Iztočnik / Source:} Unbestimmte / unbestimmt werden\par
\noindent\emph{Rusky / Russian:} неопределенные / стать неопределенной\par
\noindent\emph{Ukrajinsky / Ukrainian:} невизначені / стати невизначеною\par
\smallskip\noindent\emph{Razlog:} Ta para odkriva riziko dvojakogo značenja. Neoprěděljene sut formalne veličiny; stati se neoprěděljenoju pri specializaciji znači, že kvocient ne daje oprěděljenu vrědnost. Primětka izrično pokazuje oblik $0/0$ pri jednočasnoj zaměně dvuh argumentov. Takovy kvocient ne stava novym formalnym argumentom. Zapisana potreba visokoj pozornosti zato jest opravdana. V objasnjenju treba upotrěbiti glosu ne iměti oprěděljenu vrědnost, ne prikryvati smyslovu razliku jednim automatičnym ekvivalentom.\par
\noindent\emph{Reason:} This pair exposes a risk of two meanings. Neoprěděljene are formal quantities; becoming neoprěděljena under specialization means that a quotient fails to give a defined value. The footnote explicitly exhibits $0/0$ when two arguments are specialized simultaneously. Such a quotient does not become a new formal indeterminate. The recorded high-priority concern is therefore justified. An explanation should gloss this as having no defined value rather than conceal the distinction behind a single automatic equivalent.\par
\noindent{\footnotesize Izvorne zapisy / Original records: HIST-g042-004.}\par
\smallskip\noindent\emph{Sravnjenje s tekstom:} Priměr 0/0 razkriva razliku medžu formalnym argumentom i neuspěšnoju specializacijeju.\par
\noindent\emph{Textual comparison:} The 0/0 example exposes the difference between a formal argument and an unsuccessful specialization.\par
\noindent{\footnotesize TEXT-COMP-825; pasaže / passages: 5156, 5212, 5214, 5215, 5217; \texttt{CORPUS\_FORM\_COMPARISONS.json}.}\par
\end{minipage}\par
\par\medskip\Needspace{7\baselineskip}
\noindent\begin{minipage}{\linewidth}
\hypertarget{TERM-1769}{}
\noindent\textbf{beskonečno mnogo razkladov}\quad{\footnotesize\texttt{TERM-1769}}\par
\noindent\emph{Iztočnik / Source:} unendlich viele Zerlegungen\par
\noindent\emph{Rusky / Russian:} бесконечно много разложений\par
\noindent\emph{Ukrajinsky / Ukrainian:} нескінченно багато розкладів\par
\smallskip\noindent\emph{Razlog:} Beskonečno mnogo razkladov sohranjaje i čislo možnostij, i vid operacije. Ide o različne sumné prědstavjenja grupy, ne o beskonečno mnogo faktorov v jednom produktu. Dokaz upotrěbljaje prinajmenje dvě izomorfne podgrupy i beskonečno mnogo konstant koeficientnoj oblasti, ktore komutujut s operatorami. Tih uslovij ne možno opustiti. Slovo beskonečno ne vybira medžu prebrojimoju i neprebrojimoju množinoju razkladov; takova dodatna ocěna ne jest zapisana v terminologičnom izboru.\par
\noindent\emph{Reason:} Beskonečno mnogo razkladov preserves both the number of possibilities and the operation involved. These are distinct sum representations of the group, not infinitely many factors in a single product. The proof uses at least two isomorphic subgroups and infinitely many constants of the coefficient domain commuting with the operators. Those conditions must remain. Infinitely many does not choose between a countable and an uncountable collection of decompositions; the terminology record makes no such additional cardinality claim.\par
\noindent{\footnotesize Izvorne zapisy / Original records: HIST-g079-055.}\par
\smallskip\noindent\emph{Sravnjenje s tekstom:} Izomorfne sumandy i beskonečno mnogo konstant ostajut čest teorema.\par
\noindent\emph{Textual comparison:} Isomorphic summands and infinitely many constants remain part of the theorem.\par
\noindent{\footnotesize TEXT-COMP-826; pasaže / passages: 11731, 11733, 11735, 11737, 11743; \texttt{CORPUS\_FORM\_COMPARISONS.json}.}\par
\end{minipage}\par
\par\medskip\Needspace{7\baselineskip}
\noindent\begin{minipage}{\linewidth}
\hypertarget{TERM-1770}{}
\noindent\textbf{neparne čislo defekta}\quad{\footnotesize\texttt{TERM-1770}}\par
\noindent\emph{Iztočnik / Source:} ungerade Defektzahl\par
\noindent\emph{Rusky / Russian:} нечетное число дефекта\par
\noindent\emph{Ukrajinsky / Ukrainian:} непарне число дефекту\par
\smallskip\noindent\emph{Razlog:} Neparne čislo defekta veže parnost s čislom $\lambda$. V pasaži identičnosti s neparnym $\lambda$ slědujut iz identičnostij s vyšim defektom; pozdněje isto razgledanje redukuje něktore tvorjenja na parny defekt. Neparny zato ne znači neobyčny, neregularny ili ošteđeny. Ruski nečetny i ukrajinski neparny komparand podajut aritmetičny smysl. Slovnikovo zaglavje i padežna forma v tekstu trěbujut različno soglašenje s čislo; to jest gramatika, ne nova veličina.\par
\noindent\emph{Reason:} Neparne čislo defekta attaches parity to the number $\lambda$. In the passage, identities with odd $\lambda$ follow from identities with higher defect; later the same reasoning reduces certain expressions to even defect. Odd therefore does not mean unusual, irregular, or damaged. The Russian nečetny and Ukrainian neparny comparanda supply the arithmetic sense. The dictionary headword and the inflected phrase in the text require different agreement with čislo; that is grammar, not a new quantity.\par
\noindent{\footnotesize Izvorne zapisy / Original records: HIST-g031-012.}\par
\smallskip\noindent\emph{Sravnjenje s tekstom:} Neparnost kasaje se λ, ne kakosti dokazateljstva.\par
\noindent\emph{Textual comparison:} Oddness concerns λ, not the quality of the proof.\par
\noindent{\footnotesize TEXT-COMP-827; pasaže / passages: 3865, 4336; \texttt{CORPUS\_FORM\_COMPARISONS.json}.}\par
\end{minipage}\par
\par\medskip\Needspace{7\baselineskip}
\noindent\begin{minipage}{\linewidth}
\hypertarget{TERM-1771}{}
\noindent\textbf{neparno prosto čislo}\quad{\footnotesize\texttt{TERM-1771}}\par
\noindent\emph{Iztočnik / Source:} ungerade Primzahl\par
\noindent\emph{Rusky / Russian:} нечетное простое число\par
\noindent\emph{Ukrajinsky / Ukrainian:} непарне просте число\par
\smallskip\noindent\emph{Razlog:} Neparno prosto čislo označaje prosto čislo različno od 2. To ograničenje razděljaje idealy $(2p)$, kde p jest neparno prosto čislo, od osobnogo ideala $(2)$ v razgledanoj oblasti. Primarne idealy nad 2 potom imajut drugu ulogu pri izolovanosti. Neparno nesměje otpasti pri skračenju definicije. Istovrěmenno ne trěba miješati prosto čislo p s prostym idealom: izraz $(2p)$ imenuje ideal v toj konkretnoj oblasti, ne obče pravilo o prostyh idealah cělyh čisel.\par
\noindent\emph{Reason:} Neparno prosto čislo denotes a prime number other than 2. The restriction separates ideals $(2p)$, with p odd prime, from the special ideal $(2)$ in the domain under discussion. Primary ideals over 2 then have a different role in isolation. Odd must not disappear when the definition is shortened. Nor should the prime number p be confused with a prime ideal: $(2p)$ names an ideal in this particular domain, not a general rule about prime ideals of the integers.\par
\noindent{\footnotesize Izvorne zapisy / Original records: HIST-g081-090.}\par
\smallskip\noindent\emph{Sravnjenje s tekstom:} Aritmetičny pridavnik razděljaje idealne slučaji.\par
\noindent\emph{Textual comparison:} The arithmetic adjective separates the ideal cases.\par
\noindent{\footnotesize TEXT-COMP-828; pasaže / passages: 13645, 13651; \texttt{CORPUS\_FORM\_COMPARISONS.json}.}\par
\end{minipage}\par
\par\medskip\Needspace{7\baselineskip}
\noindent\begin{minipage}{\linewidth}
\hypertarget{TERM-1772}{}
\noindent\textbf{neparny relativny stepen}\quad{\footnotesize\texttt{TERM-1772}}\par
\noindent\emph{Iztočnik / Source:} ungerader Relativgrad\par
\noindent\emph{Rusky / Russian:} нечетная относительная степень\par
\noindent\emph{Ukrajinsky / Ukrainian:} непарний відносний степінь\par
\smallskip\noindent\emph{Razlog:} Neparny relativny stepen jest neparna dimenzija razširitelnogo polja nad ukazanym osnovnym poljem. V argumentu osnovno polje jest K, a razpadno polje $P(\varepsilon)$ ima nad K stepen $(p-1)/2^n$. To izključaje preostaly indeks 2 kvaternionnogo těla. Stepen tu ne jest potega elementa ni red proizvoljnoj grupy. Ukazanje relativny sohranjaje osnovu izračuna; bez njega čitatelj mogl by pogrěšno upotrěbiti stepen nad P vměsto nad K.\par
\noindent\emph{Reason:} Neparny relativny stepen is an odd extension degree over the specified base field. In the argument the base is K, and the splitting field $P(\varepsilon)$ has degree $(p-1)/2^n$ over K. This excludes the remaining index 2 of the quaternion algebra. Degree here is neither an element's power nor an arbitrary group's order. Relative preserves the base of the calculation; without it a reader might mistakenly use the degree over P instead of K.\par
\noindent{\footnotesize Izvorne zapisy / Original records: HIST-g175-006.}\par
\smallskip\noindent\emph{Sravnjenje s tekstom:} Osnovno polje i indeks muset ostati različenymi.\par
\noindent\emph{Textual comparison:} The base field and the index must remain distinguished.\par
\noindent{\footnotesize TEXT-COMP-829; pasaže / passages: 19638; \texttt{CORPUS\_FORM\_COMPARISONS.json}.}\par
\end{minipage}\par
\par\medskip\Needspace{7\baselineskip}
\noindent\begin{minipage}{\linewidth}
\hypertarget{TERM-1773}{}
\noindent\textbf{unimodularny}\quad{\footnotesize\texttt{TERM-1773}}\par
\noindent\emph{Iztočnik / Source:} unimodular\par
\noindent\emph{Rusky / Russian:} унимодулярный\par
\noindent\emph{Ukrajinsky / Ukrainian:} унімодулярний\par
\smallskip\noindent\emph{Razlog:} Unimodularny sohranjaje medžunarodny korenj, vidljivy takože v ruskom i ukrajinskom komparandu. Arhiv govori o izběganju objašnjenja determinant jeden, ale toj glos jest prěuzky dlja zdejše cěločiselne matrice: zaměna bazy modula dopušča determinant $1$ ili $-1$. Glavno jest obratimost nad tym samym prstenom, ne samo nenulovy determinant nad poljem. Slovnik zato sohranjaje izbor termina, ale ne prěnosit nepodkrepljeno suženje sohranjenej motivacije.\par
\noindent\emph{Reason:} Unimodularny retains the international stem also visible in the Russian and Ukrainian comparanda. The archive speaks of avoiding a determinant-one paraphrase, but that gloss is too narrow for the integer matrices here: a module basis change allows determinant $1$ or $-1$. What matters is invertibility over the same ring, not merely a nonzero determinant over a field. The dictionary therefore preserves the term while declining to carry forward the unsupported narrowing in its recorded motivation.\par
\noindent{\footnotesize Izvorne zapisy / Original records: HIST-g081-097.}\par
\smallskip\noindent\emph{Sravnjenje s tekstom:} Unimodularnost zachovuje toj samy modul; glosa determinant jeden jest prěuzka.\par
\noindent\emph{Textual comparison:} Unimodularity preserves the same module; the gloss determinant one is too narrow.\par
\noindent{\footnotesize TEXT-COMP-830; pasaže / passages: 13703, 13705, 13707; \texttt{CORPUS\_FORM\_COMPARISONS.json}.}\par
\end{minipage}\par
\par\medskip\Needspace{7\baselineskip}
\noindent\begin{minipage}{\linewidth}
\hypertarget{TERM-1774}{}
\noindent\textbf{unimodularna matrica}\quad{\footnotesize\texttt{TERM-1774}}\par
\noindent\emph{Iztočnik / Source:} unimodulare Matrix\par
\noindent\emph{Rusky / Russian:} унимодулярная матрица\par
\noindent\emph{Ukrajinsky / Ukrainian:} унімодулярна матриця\par
\smallskip\noindent\emph{Razlog:} Unimodularna matrica jest matrica dopustimoj zaměny bazy. V toj prstennoj situaciji determinant P mora byti jediniceju iz Z; ne trěba, aby byl doslovno 1. Prědstavjenja potom svjazany sut formulom $R_b=P^{-1}R_aP$. Izbor medžunarodnogo pridavnika izběgaje novotvorbu, ale njegova definicija mora ostati svezana s koeficientnym prstenom. Samo nenulovost determinanta ne garantuje, že $P^{-1}$ ima koeficienty v Z.\par
\noindent\emph{Reason:} Unimodularna matrica is the matrix of an admissible basis change. In this ring setting det P must be a unit of Z; it need not literally equal 1. The representations are then related by $R_b=P^{-1}R_aP$. The international adjective avoids a new coinage, but its definition must remain tied to the coefficient ring. A merely nonzero determinant does not guarantee that $P^{-1}$ has coefficients in Z.\par
\noindent{\footnotesize Izvorne zapisy / Original records: HIST-g133-005.}\par
\smallskip\noindent\emph{Sravnjenje s tekstom:} Konkretna prědpoloženja o bazi dopolnjajut nazvu unimodularna.\par
\noindent\emph{Textual comparison:} The actual basis assumptions specify the meaning of unimodularna.\par
\noindent{\footnotesize TEXT-COMP-831; pasaže / passages: 17827, 17913, 17952, 17954; \texttt{CORPUS\_FORM\_COMPARISONS.json}.}\par
\end{minipage}\par
\par\medskip\Needspace{7\baselineskip}
\noindent\begin{minipage}{\linewidth}
\hypertarget{TERM-1775}{}
\noindent\textbf{unimodularna transformacija}\quad{\footnotesize\texttt{TERM-1775}}\par
\noindent\emph{Iztočnik / Source:} unimodulare Transformation\par
\noindent\emph{Rusky / Russian:} унимодулярное преобразование\par
\noindent\emph{Ukrajinsky / Ukrainian:} унімодулярне перетворення\par
\smallskip\noindent\emph{Razlog:} Unimodularna transformacija označaje obratimu zaměnu s koeficientami v zadanoj oblasti cělyh veličin. Tu ona jest ili oblast cělyh čisel, ili polinomy jednoj neoprěděljenej nad poljem. Formula $PAQ$ privede matricu k normalnoj formě, a $\xi=Q(\eta)$ menja formalne argumenty. Transformacija i prěobraženje v aktualnom tekstu služe tomu samomu matematičnomu odnosu. Arhivna prednost medžunarodnoj frazy ne daje osnovy zaměniti unimodularnost samo uslovjem nenulovogo determinanta.\par
\noindent\emph{Reason:} Unimodularna transformacija denotes an invertible change over the stated domain of integral quantities. Here that domain is either the integers or polynomials in one indeterminate over a field. $PAQ$ brings the matrix to normal form, while $\xi=Q(\eta)$ changes the formal arguments. Transformacija and prěobraženje in the current text serve the same mathematical relation. The archived preference for international terminology does not justify weakening unimodularity to a nonzero determinant condition.\par
\noindent{\footnotesize Izvorne zapisy / Original records: HIST-g085-010, HIST-g092-010.}\par
\smallskip\noindent\emph{Sravnjenje s tekstom:} Oblast koeficientov oprěděljaje, ktore obratne transformacije sut dopustime.\par
\noindent\emph{Textual comparison:} The coefficient domain determines which invertible transformations are admissible.\par
\noindent{\footnotesize TEXT-COMP-832; pasaže / passages: 14226, 14231, 14233, 14240, 14242, 14484; \texttt{CORPUS\_FORM\_COMPARISONS.json}.}\par
\end{minipage}\par
\par\medskip\Needspace{7\baselineskip}
\noindent\begin{minipage}{\linewidth}
\hypertarget{TERM-1776}{}
\noindent\textbf{nemožnost identičnosti}\quad{\footnotesize\texttt{TERM-1776}}\par
\noindent\emph{Iztočnik / Source:} Unmöglichkeit einer Identität\par
\noindent\emph{Rusky / Russian:} невозможность тождества\par
\noindent\emph{Ukrajinsky / Ukrainian:} неможливість тотожності\par
\smallskip\noindent\emph{Razlog:} Nemožnost identičnosti označaje nemožnost opisanogo odnosa pri dannyh uslovjah, ne trudnost jego izračunanja. Dokaz specializuje simbolične redy, aby pokazati, že predložena identičnost ne može važiti; ograničenja na indeksi i vagy redov sut čest argumenta. Zato fraza ne tvrdi, že medžu simboličnymi formami vůbec ne obstajajut identičnosti. Slovnik sohranjaje silu negacije i konkretno matematično področje; glosa těžko najdti identičnost byla by slaba i nevěrna.\par
\noindent\emph{Reason:} Nemožnost identičnosti denotes impossibility of the specified relation under its conditions, not difficulty in calculating it. The proof specializes symbolic rows to show that the proposed identity cannot hold; restrictions on indices and row weights are part of the argument. The phrase does not assert that symbolic forms admit no identities at all. The dictionary preserves the strength of negation and its mathematical scope; difficult to find an identity would be a weaker and inaccurate gloss.\par
\noindent{\footnotesize Izvorne zapisy / Original records: HIST-g032-012.}\par
\smallskip\noindent\emph{Sravnjenje s tekstom:} Specializacija izključaje konkretno opisanu identičnost.\par
\noindent\emph{Textual comparison:} Specialization rules out the particular identity described.\par
\noindent{\footnotesize TEXT-COMP-833; pasaže / passages: 4026, 4028, 4031, 4032, 4033, 4036; \texttt{CORPUS\_FORM\_COMPARISONS.json}.}\par
\end{minipage}\par
\par\medskip\Needspace{7\baselineskip}
\noindent\begin{minipage}{\linewidth}
\hypertarget{TERM-1777}{}
\noindent\textbf{diskontinuirny / ekstremno diskontinuirny / totalno diskontinuirny}\quad{\footnotesize\texttt{TERM-1777}}\par
\noindent\emph{Iztočnik / Source:} unstetig / extrem unstetig / total unstetig\par
\noindent\emph{Rusky / Russian:} разрывный / экстремально разрывный / полностью разрывный\par
\noindent\emph{Ukrajinsky / Ukrainian:} розривний / екстремально розривний / цілком розривний\par
\smallskip\noindent\emph{Razlog:} Diskontinuirny jest vybrano vměsto neprěryvny, zato že poslědnja forma sugeruje kontinuitet, tedy protivny smysl. Ta motivacija jest izrično sohranjena; ona ne jest statističny rezultat čitatelskogo testa. Ekstremno diskontinuirny ima v rabotě silnu definiciju: bliz vsakogo $z_0$ možno dostignuti vrědnosti libovoljno blizke k vsakomu zadannomu $Z_0$. Avtorska primětka upozornjaje, že Hamelovo totalno diskontinuirny indže znači něčto drugo. Te nazvy zato ne treba proglasiti bezuslovnymi sinonimami.\par
\noindent\emph{Reason:} Diskontinuirny was chosen over neprěryvny because the latter suggests continuity, the opposite meaning. That motivation is explicitly preserved; it is not a statistical result from reader testing. Ekstremno diskontinuirny has a strong definition in the work: near any $z_0$ one can attain values arbitrarily close to any prescribed $Z_0$. The author's note warns that Hamel's totalno diskontinuirny has another meaning elsewhere. These labels should therefore not be declared unconditional synonyms.\par
\noindent{\footnotesize Izvorne zapisy / Original records: HIST-g069-009.}\par
\smallskip\noindent\emph{Sravnjenje s tekstom:} Avtorska ograda o Hamelovoj nazvě ostavaje sučstvena.\par
\noindent\emph{Textual comparison:} The author's qualification concerning Hamel's label remains essential.\par
\noindent{\footnotesize TEXT-COMP-834; pasaže / passages: 8445, 8601; \texttt{CORPUS\_FORM\_COMPARISONS.json}.}\par
\end{minipage}\par
\par\medskip\Needspace{7\baselineskip}
\noindent\begin{minipage}{\linewidth}
\hypertarget{TERM-1778}{}
\noindent\textbf{diskontinuirny / kontinuirny / ekstremno diskontinuirny}\quad{\footnotesize\texttt{TERM-1778}}\par
\noindent\emph{Iztočnik / Source:} unstetig / stetig / extrem unstetig\par
\noindent\emph{Rusky / Russian:} разрывный / непрерывный / экстремально разрывный\par
\noindent\emph{Ukrajinsky / Ukrainian:} розривний / неперервний / екстремально розривний\par
\smallskip\noindent\emph{Razlog:} Para kontinuirny i diskontinuirny tvori prozračny opozitny par; ekstremno dodaje točno silnějše svojstvo, ne retorično pojačanje. V §4 sut kompleksne rěšenja aditivnogo ravnanja, ktore sut diskontinuirne, ale ne ekstremno diskontinuirne. Nakonečny teorem kasaje se rěšenij cěloj sistemy (1)–(4), s izključenjem z i jeho kompleksnogo konjugata. Razlika medžu jednim ravnjanjem i cěloju sistemoju jest neophodima, aby terminologija ne sugerovala ložnu teoremu.\par
\noindent\emph{Reason:} Kontinuirny and diskontinuirny form a transparent opposed pair; ekstremno adds a precisely stronger property, not rhetorical emphasis. Section 4 gives complex solutions of the additive equation which are discontinuous but not extremely discontinuous. The concluding theorem concerns solutions of the full system (1)–(4), except z and its complex conjugate. Distinguishing one equation from the whole system is essential if the terminology is not to suggest a false theorem.\par
\noindent{\footnotesize Izvorne zapisy / Original records: HIST-g072-001.}\par
\smallskip\noindent\emph{Sravnjenje s tekstom:} Kontrapriměry i zaključny teorem imajut različne hipotezy.\par
\noindent\emph{Textual comparison:} The counterexamples and the concluding theorem have different hypotheses.\par
\noindent{\footnotesize TEXT-COMP-835; pasaže / passages: 8605, 8609, 8690; \texttt{CORPUS\_FORM\_COMPARISONS.json}.}\par
\end{minipage}\par
\par\medskip\Needspace{7\baselineskip}
\noindent\begin{minipage}{\linewidth}
\hypertarget{TERM-1779}{}
\noindent\textbf{vrědnosti diskontinuirnosti}\quad{\footnotesize\texttt{TERM-1779}}\par
\noindent\emph{Iztočnik / Source:} Unstetigkeitswerte\par
\noindent\emph{Rusky / Russian:} значения разрывности\par
\noindent\emph{Ukrajinsky / Ukrainian:} значення розривності\par
\smallskip\noindent\emph{Razlog:} Vrědnosti diskontinuirnosti jest arhivna fraza; aktualna pasaža ima cěnnosti diskontinuirnosti. Ide o vrědnosti v ciljnoj oblasti, ktore možno libovoljno blizko aproksimovati obrazami argumentov bliz danogo $z_0$. To ne sut točky diskontinuiteta v oblasti argumentov. Opisane množiny obraznyh vrědnostij mogut byti jedno- ili dvurazměrne. Zachovanje neobyčnej avtorskoj nazvy potrěbuje glosu, ne zaměnu običnym terminom za točki prěryva. Varianty vrědnosti i cěnnosti ne měnjajut to različenje.\par
\noindent\emph{Reason:} Vrědnosti diskontinuirnosti is the archived phrase; the current passage has cěnnosti diskontinuirnosti. These are values in the target which can be approximated arbitrarily closely by images of arguments near a given $z_0$. They are not points of discontinuity in the domain. The described sets of image values may be one- or two-dimensional. Retaining the author's unusual label requires a gloss, not replacement with the ordinary term for discontinuity points. Vrědnosti and cěnnosti do not change that distinction.\par
\noindent{\footnotesize Izvorne zapisy / Original records: HIST-g072-006.}\par
\smallskip\noindent\emph{Sravnjenje s tekstom:} Jedno- i dvurazměrnost kasajut se množin obraznyh vrědnostij.\par
\noindent\emph{Textual comparison:} One- and two-dimensionality concern sets of image values.\par
\noindent{\footnotesize TEXT-COMP-836; pasaže / passages: 8640, 8642, 8644, 8646; \texttt{CORPUS\_FORM\_COMPARISONS.json}.}\par
\end{minipage}\par
\par\medskip\Needspace{7\baselineskip}
\noindent\begin{minipage}{\linewidth}
\hypertarget{TERM-1783}{}
\noindent\textbf{medžu soboju različne izomorfizmy}\quad{\footnotesize\texttt{TERM-1783}}\par
\noindent\emph{Iztočnik / Source:} unter sich verschiedene Isomorphismen\par
\noindent\emph{Rusky / Russian:} попарно различные изоморфизмы\par
\noindent\emph{Ukrajinsky / Ukrainian:} попарно різні ізоморфізми\par
\smallskip\noindent\emph{Razlog:} Medžu soboju različne izomorfizmy znači poparno različne odobraženja, ne neizomorfne polja. V pasaži vse prirědžene polja sut izomorfne k osnovnomu objektu, ale n odobraženij sut različne; nenulovy diskriminant i jego determinantno prědstavjenje ustanavjajut tu razliku. Arhivna ruska i ukrajinska forma s poparno pojasnjajut kvantifikaciju. Sohranjena primětka govori o ranějšej prověrci skana; ta slovnikova rekonstrukcija ne vydavaje ju za novu prověrku skana.\par
\noindent\emph{Reason:} Medžu soboju različne izomorfizmy means pairwise distinct maps, not nonisomorphic fields. In the passage all the associated fields are isomorphic to the object being studied, but the n maps are distinct; the nonzero discriminant and its determinant expression establish that distinction. The archived Russian and Ukrainian forms with pairwise clarify the quantification. A preserved note reports an earlier scan check; this dictionary reconstruction does not present that report as a new examination of the scan.\par
\noindent{\footnotesize Izvorne zapisy / Original records: HIST-g147-008.}\par
\smallskip\noindent\emph{Sravnjenje s tekstom:} Nenulovy diskriminant podpiraje poparnu različnost odobraženij.\par
\noindent\emph{Textual comparison:} The nonzero discriminant supports pairwise distinctness of the maps.\par
\noindent{\footnotesize TEXT-COMP-840; pasaže / passages: 18525; \texttt{CORPUS\_FORM\_COMPARISONS.json}.}\par
\end{minipage}\par
\par\medskip\Needspace{7\baselineskip}
\noindent\begin{minipage}{\linewidth}
\hypertarget{TERM-1784}{}
\noindent\textbf{minory}\quad{\footnotesize\texttt{TERM-1784}}\par
\noindent\emph{Iztočnik / Source:} Unterdeterminanten\par
\noindent\emph{Rusky / Russian:} миноры\par
\noindent\emph{Ukrajinsky / Ukrainian:} мінори\par
\smallskip\noindent\emph{Razlog:} Minory jest medžunarodna nazva determinantov kvadratnyh podmatric. V konkretnoj pasaži koordinaty prěmyh sut minory iz dvuh redov točkovyh koordinat, a obča konstrukcija upotrěbljaje $\rho$ redov. Izbor izběgaje bukvalnu novotvorbu za německo Unterdeterminanten i odgovarja sohranjenym ruskim i ukrajinskim komparandam. Minor ne jest prosto maly koeficient ili odrezok matrice; determinant ostavaje čest definicije. Polny determinant i njegov poddeterminant ne trěba smyslovo miješati.\par
\noindent\emph{Reason:} Minory is the international name for determinants of square submatrices. In the passage, line coordinates are minors formed from two rows of point coordinates, and the general construction uses $\rho$ rows. The choice avoids a literal new calque of Unterdeterminanten and agrees with the preserved Russian and Ukrainian comparanda. A minor is not simply a small coefficient or a piece of a matrix; taking a determinant remains part of the definition. Full determinants and subdeterminants must remain conceptually distinct.\par
\noindent{\footnotesize Izvorne zapisy / Original records: HIST-g028-009.}\par
\smallskip\noindent\emph{Sravnjenje s tekstom:} Poddeterminanty imajut konkretno čislo redov.\par
\noindent\emph{Textual comparison:} The subdeterminants have a specified number of rows.\par
\noindent{\footnotesize TEXT-COMP-841; pasaže / passages: 3651, 3653; \texttt{CORPUS\_FORM\_COMPARISONS.json}.}\par
\end{minipage}\par
\par\medskip\Needspace{7\baselineskip}
\noindent\begin{minipage}{\linewidth}
\hypertarget{TERM-1785}{}
\noindent\textbf{podgrupa grupy ostatkov}\quad{\footnotesize\texttt{TERM-1785}}\par
\noindent\emph{Iztočnik / Source:} Untergruppe der Restgruppe\par
\noindent\emph{Rusky / Russian:} подгруппа группы остатков\par
\noindent\emph{Ukrajinsky / Ukrainian:} підгрупа групи лишків\par
\smallskip\noindent\emph{Razlog:} Podgrupa grupy ostatkov sohranjaje avtorsku nazvu dlja specialnogo sistema klasov. Tu dozvoljeno jest dodavanje klasov i lěvo množenje polinomom; produkt dvuh libovoljnyh klasov obče ne jest definovany. Podgrupy v razloženju sut izomorfne grupam ostatkov odgovarjajučih modulov. Zato podgrupa ne znači libovoljnu podmnožinu, ale ni bez oglěda na kontekst obyčnu multiplikativnu grupu. Definicija ostavaje potrěbna pored kratkoj nazvy.\par
\noindent\emph{Reason:} Podgrupa grupy ostatkov retains the author's name for a particular system of classes. Addition of classes and left multiplication by a polynomial are permitted here; a product of two arbitrary classes is generally undefined. The subgroups in the decomposition are isomorphic to the residue groups of the corresponding modules. Subgroup therefore means neither an arbitrary subset nor, without regard to context, an ordinary multiplicative group. The definition remains necessary alongside the short name.\par
\noindent{\footnotesize Izvorne zapisy / Original records: HIST-g079-024.}\par
\smallskip\noindent\emph{Sravnjenje s tekstom:} Podgrupa ne davaje neobstoječi produkt libovoljnyh klasov.\par
\noindent\emph{Textual comparison:} Subgroup does not supply an undefined product of arbitrary classes.\par
\noindent{\footnotesize TEXT-COMP-842; pasaže / passages: 11102, 11110, 11122, 11210, 11214; \texttt{CORPUS\_FORM\_COMPARISONS.json}.}\par
\end{minipage}\par
\par\medskip\Needspace{7\baselineskip}
\noindent\begin{minipage}{\linewidth}
\hypertarget{TERM-1786}{}
\noindent\textbf{podgrupy}\quad{\footnotesize\texttt{TERM-1786}}\par
\noindent\emph{Iztočnik / Source:} Untergruppen\par
\noindent\emph{Rusky / Russian:} подгруппы\par
\noindent\emph{Ukrajinsky / Ukrainian:} підгрупи\par
\smallskip\noindent\emph{Razlog:} Podgrupy v navodnikah jest historijna nazva iz raboty Noether i Schmeidlera. Sama razgledana pasaža objasnjaje, že te sistemy klasov ostatkov sut priměry modulov v dvojnoj oblasti: razlika klasov i množenje polinomom ostajut dopustime, produkt dvuh klasov ne. Prěvod ne zaměnjaje citovanu nazvu sovrěmennym podmoduly, ale slovnik daje to smyslovo objasnjenje. To jest vernost istorii terminologije, ne predlog upotrěbljati podgrupa bez kvalifikacije vo vsakoj modulnoj situaciji.\par
\noindent\emph{Reason:} Podgrupy in quotation marks is the historical label from the work of Noether and Schmeidler. The passage itself explains that these systems of residue classes are examples of modules in a double domain: subtraction of classes and multiplication by a polynomial are allowed, but multiplication of two classes is not. The translation retains the quoted label rather than substituting modern submodules; the dictionary supplies the explanation. This preserves terminological history, not a recommendation to use subgroup without qualification in every module-theoretic setting.\par
\noindent{\footnotesize Izvorne zapisy / Original records: HIST-g081-095.}\par
\smallskip\noindent\emph{Sravnjenje s tekstom:} Historijna nazva sohranjena jest vměstě s izričnoju definicijeju.\par
\noindent\emph{Textual comparison:} The historical label is retained alongside the explicit definition.\par
\noindent{\footnotesize TEXT-COMP-843; pasaže / passages: 13668; \texttt{CORPUS\_FORM\_COMPARISONS.json}.}\par
\end{minipage}\par
\par\medskip\Needspace{7\baselineskip}
\noindent\begin{minipage}{\linewidth}
\hypertarget{TERM-1787}{}
\noindent\textbf{podgrupy projektivnoj grupy}\quad{\footnotesize\texttt{TERM-1787}}\par
\noindent\emph{Iztočnik / Source:} Untergruppen der projektiven Gruppe\par
\noindent\emph{Rusky / Russian:} подгруппы проективной группы\par
\noindent\emph{Ukrajinsky / Ukrainian:} підгрупи проєктивної групи\par
\smallskip\noindent\emph{Razlog:} Podgrupy projektivnoj grupy jest sostavny termin s produktivnym pod-. On upominaje grupy transformacij, po odnošenju k ktorym tvoret se invariantne polja. V neposrědnoj pasaži osobno vysvětlěno jest ograničenje na determinant 1, aby racionalne kombinacije invariantov znovu davali invarianty i tvorili polje. Literalna fraza ne opravdavaje potlačenje toj konvencije. Podgrupa tu ne znači podprostor projektivnogo prostora: pod- priključeno jest k grupa, ne k geometričnomu nositelju.\par
\noindent\emph{Reason:} Podgrupy projektivnoj grupy is a compound term using productive pod-. It refers to transformation groups with respect to which invariant fields are formed. The nearby passage expressly discusses restricting to determinant 1 so that rational combinations remain invariant and form a field. The literal phrase does not license suppressing that convention. Subgroup here does not mean a subspace of projective space: the prefix modifies group, not the geometric space on which it acts.\par
\noindent{\footnotesize Izvorne zapisy / Original records: HIST-g041-017.}\par
\smallskip\noindent\emph{Sravnjenje s tekstom:} Prefiks pod- ne zaměnjaje uslovja o transformacijah.\par
\noindent\emph{Textual comparison:} The prefix pod- does not replace the assumptions about transformations.\par
\noindent{\footnotesize TEXT-COMP-844; pasaže / passages: 5130, 5132; \texttt{CORPUS\_FORM\_COMPARISONS.json}.}\par
\end{minipage}\par
\par\medskip\Needspace{7\baselineskip}
\noindent\begin{minipage}{\linewidth}
\hypertarget{TERM-1788}{}
\noindent\textbf{podpolje}\quad{\footnotesize\texttt{TERM-1788}}\par
\noindent\emph{Iztočnik / Source:} Unterkörper\par
\noindent\emph{Rusky / Russian:} подполе\par
\noindent\emph{Ukrajinsky / Ukrainian:} підполе\par
\smallskip\noindent\emph{Razlog:} Podpolje prodolžaje simetričnu družinu polje, podpolje, nadpolje. V razloženju prstena komponenta sodrživaje kopiju $P_i=P\varepsilon_i$, izomorfnu k P; v kvocientnom polju podobnu ulogu ima obraz (P). Izomorfna kopija ne mora byti bukvalno isto podmnožestvo s istoju jediniceju. Kompaktna prefiksna forma sohranjaje odnos vloženja, ale konkretna definicija mora kazati, kude vloženje vodi. To ne daje osnovy zaměnjati vsako nekomutativno tělo slovom polje.\par
\noindent\emph{Reason:} Podpolje continues the symmetrical family polje, podpolje, nadpolje. In the ring decomposition a component contains a copy $P_i=P\varepsilon_i$ isomorphic to P; the image (P) has a similar role in the quotient field. An isomorphic copy need not literally be the same subset with the same identity. The compact prefix preserves the inclusion relation, while the definition must specify its target. This does not justify replacing every noncommutative division ring with field.\par
\noindent{\footnotesize Izvorne zapisy / Original records: HIST-g122-006, HIST-g124-004.}\par
\smallskip\noindent\emph{Sravnjenje s tekstom:} P\_i i (P) trěbujut svoje izomorfne identifikacije.\par
\noindent\emph{Textual comparison:} P\_i and (P) require their respective isomorphic identifications.\par
\noindent{\footnotesize TEXT-COMP-845; pasaže / passages: 17524, 17595, 17597, 17599; \texttt{CORPUS\_FORM\_COMPARISONS.json}.}\par
\end{minipage}\par
\par\medskip\Needspace{7\baselineskip}
\noindent\begin{minipage}{\linewidth}
\hypertarget{TERM-1789}{}
\noindent\textbf{podtělo}\quad{\footnotesize\texttt{TERM-1789}}\par
\noindent\emph{Iztočnik / Source:} Unterkörper\par
\noindent\emph{Rusky / Russian:} подтело\par
\noindent\emph{Ukrajinsky / Ukrainian:} підтіло\par
\smallskip\noindent\emph{Razlog:} Podtělo čuva šire značenje tělo v teoriji prědstavjenij. $\Gamma$ jest podtělo K, a K realizuje tělo avtomorfizmov; komutativnost v toj formulaciji ne jest obče zajavjena. Zato podpolje bylo by neopravdano suženje. Pozdnějša pasaža dodaje komutativnost P kako osobno prědpoloženje i svjazuje jeho s centrom K. Arhivny izbor tělo jest smyslovy, ne samo pravopisny: slovnik dolžen označiti uslovja, pri ktoryh možno govoriti o komutativnom polju.\par
\noindent\emph{Reason:} Podtělo preserves the broader meaning of tělo in representation theory. $\Gamma$ is a division subring of K, and K realizes the automorphism division ring; commutativity is not generally asserted in that formulation. Podpolje would therefore narrow the meaning without justification. A later passage adds commutativity of P as a separate assumption and relates it to the centre of K. The choice of tělo is semantic, not merely orthographic: the dictionary must state when the commutative field interpretation is available.\par
\noindent{\footnotesize Izvorne zapisy / Original records: HIST-g196-001, HIST-g197-008.}\par
\smallskip\noindent\emph{Sravnjenje s tekstom:} Osobno prědpoloženje o komutativnosti ne směje byti vneseno raněje.\par
\noindent\emph{Textual comparison:} The separate commutativity assumption must not be imported earlier.\par
\noindent{\footnotesize TEXT-COMP-846; pasaže / passages: 21691, 21707, 21729, 21755; \texttt{CORPUS\_FORM\_COMPARISONS.json}.}\par
\end{minipage}\par
\par\medskip\Needspace{7\baselineskip}
\noindent\begin{minipage}{\linewidth}
\hypertarget{TERM-1790}{}
\noindent\textbf{podpolje $P$, izomorfno k polju $\mathfrak{o}/\mathfrak{o}p$}\quad{\footnotesize\texttt{TERM-1790}}\par
\noindent\emph{Iztočnik / Source:} Unterkörper $P$ isomorph zu $\mathfrak{o}/\mathfrak{o}p$\par
\noindent\emph{Rusky / Russian:} подполе $P$, изоморфное полю $\mathfrak{o}/\mathfrak{o}p$\par
\noindent\emph{Ukrajinsky / Ukrainian:} підполе $P$, ізоморфне полю $\mathfrak{o}/\mathfrak{o}p$\par
\smallskip\noindent\emph{Razlog:} Podpolje P, izomorfno k polju $\mathfrak{o}/\mathfrak{o}p$ razděljaje dva objekta: kvocient osnovnogo prstena i njegov vloženy obraz v prstenu klasov ostatkov poredka. Aktualna formula ostvarjaje vloženje prěz druhy teorem o izomorfizmu. Slovnikovo skračenje $\mathfrak{o}/\mathfrak{o}p$ trěba čitati kako kvocient prstena o po idealu op, ne kako običny drob dvuh čisel. Prsten ostatkov sam ne mora byti polje samo zato, že sodrživaje tako podpolje.\par
\noindent\emph{Reason:} Podpolje P, izomorfno k polju $\mathfrak{o}/\mathfrak{o}p$ distinguishes two objects: the quotient of the base ring and its embedded image in the residue ring of the order. The current formula constructs that embedding through the second isomorphism theorem. The headword's shorthand $\mathfrak{o}/\mathfrak{o}p$ means the quotient of the ring o by the ideal op, not an ordinary fraction of two numbers. The residue ring itself need not be a field merely because it contains this subfield.\par
\noindent{\footnotesize Izvorne zapisy / Original records: HIST-g151-005.}\par
\smallskip\noindent\emph{Sravnjenje s tekstom:} Izomorfizm razděljaje izvorno polje ostatkov i njegov obraz.\par
\noindent\emph{Textual comparison:} The isomorphism distinguishes the residue field from its image.\par
\noindent{\footnotesize TEXT-COMP-847; pasaže / passages: 18627, 18629, 18631, 18633; \texttt{CORPUS\_FORM\_COMPARISONS.json}.}\par
\end{minipage}\par
\par\medskip\Needspace{7\baselineskip}
\noindent\begin{minipage}{\linewidth}
\hypertarget{TERM-1791}{}
\noindent\textbf{podmodul}\quad{\footnotesize\texttt{TERM-1791}}\par
\noindent\emph{Iztočnik / Source:} Untermodul\par
\noindent\emph{Rusky / Russian:} подмодуль\par
\noindent\emph{Ukrajinsky / Ukrainian:} підмодуль\par
\smallskip\noindent\emph{Razlog:} Podmodul jest kratka nazva dopustimoj podgrupy, zakrytoj pod zadannymi operatorami. V argumentu o prědstavjenjah iz podmodula U vybira se baza $z_i$, dopolnjena do bazy M. Stabilnost U daje nulovy gorny pravy blok matrice, a ostale diagonalne bloky prědstavjajut U i $M/U$. To pokazuje, začto libovoljna aditivna podgrupa ne dostatočna. Arhiv dava prednost kratkomu modulnomu terminu; konkretna operatorna definicija, ne samo podobnost slov, opravdava jego upotrěbu.\par
\noindent\emph{Reason:} Podmodul is the short name for an admissible subgroup closed under the specified operators. In the representation argument a basis $z_i$ of the submodule U is extended to a basis of M. Stability of U yields a zero upper-right matrix block, while the diagonal blocks represent U and $M/U$. This shows why an arbitrary additive subgroup is insufficient. The archive prefers the concise module term; its operator-sensitive definition, not verbal resemblance alone, justifies its use.\par
\noindent{\footnotesize Izvorne zapisy / Original records: HIST-g194-002.}\par
\smallskip\noindent\emph{Sravnjenje s tekstom:} Matričny oblik daje konkretnu prověrku operatornoj stabilnosti.\par
\noindent\emph{Textual comparison:} The matrix form gives a concrete test of operator stability.\par
\noindent{\footnotesize TEXT-COMP-848; pasaže / passages: 19835, 21488, 21494, 21496, 21503, 21517; \texttt{CORPUS\_FORM\_COMPARISONS.json}.}\par
\end{minipage}\par
\par\medskip\Needspace{7\baselineskip}
\noindent\begin{minipage}{\linewidth}
\hypertarget{TERM-1792}{}
\noindent\textbf{podkolco}\quad{\footnotesize\texttt{TERM-1792}}\par
\noindent\emph{Iztočnik / Source:} Unterring\par
\noindent\emph{Rusky / Russian:} подкольцо\par
\noindent\emph{Ukrajinsky / Ukrainian:} підкільце\par
\smallskip\noindent\emph{Razlog:} Podkolco jest ranějši izbor, ktory sohranjal prefiksnu paru podkolco i nadkolco. Aktualny tekst upotrěbljaje podprsten i prsten; slovnik ne směje predstavjati podkolco kako jedinu aktualnu normu. V razgledanom dokazu prstenovogo izomorfizma podprsten nosi operacije naslědene od nadprstena. To jest jače než prosty odnos podmnožestva. Sravnjenje sohranjaje istoriju korenja kolco, ale odděljaje ju od nyněšnjej upotrěby prsten.\par
\noindent\emph{Reason:} Podkolco is an earlier choice preserving the prefix pair podkolco and nadkolco. The current text uses podprsten and prsten; the dictionary must not present podkolco as the sole current norm. In the ring-isomorphism argument the subring carries operations inherited from the containing ring. This is stronger than mere subset inclusion. The comparison preserves the history of the stem kolco while distinguishing it from current prsten.\par
\noindent{\footnotesize Izvorne zapisy / Original records: HIST-g120-007.}\par
\smallskip\noindent\emph{Sravnjenje s tekstom:} Zaměna korenja ne měnjaje odnos naslědenyh operacij.\par
\noindent\emph{Textual comparison:} The change of stem does not alter the relation of inherited operations.\par
\noindent{\footnotesize TEXT-COMP-849; pasaže / passages: 17398; \texttt{CORPUS\_FORM\_COMPARISONS.json}.}\par
\end{minipage}\par
\par\medskip\Needspace{7\baselineskip}
\noindent\begin{minipage}{\linewidth}
\hypertarget{TERM-1793}{}
\noindent\textbf{podkolco / nadkolco / nadpolje}\quad{\footnotesize\texttt{TERM-1793}}\par
\noindent\emph{Iztočnik / Source:} Unterring / Oberring / Oberkörper\par
\noindent\emph{Rusky / Russian:} подкольцо / надкольцо / надполе\par
\noindent\emph{Ukrajinsky / Ukrainian:} підкільце / надкільце / надполе\par
\smallskip\noindent\emph{Razlog:} Podkolco, nadkolco i nadpolje byli vybrane kako simetrična prefiksna družina. V aktualnoj pasaži prve dve formy zaměnjene sut s podprsten i nadprsten, a nadpolje ostavaje. Prěhod k nadpolju služit razširjenju oblasti koeficientov; ne jest automatično tvrdženje, že vsako nadprsten jest polje. Vložene polje P i razgledano prsten imajut obšču jedinicu po izričnom prědpoloženju. Morfologična simetrija jest jezyčna prednost, ne zaměna tih uslovij.\par
\noindent\emph{Reason:} Podkolco, nadkolco and nadpolje were chosen as a symmetrical prefix family. In the current passage the first two have become podprsten and nadprsten, while nadpolje remains. Passing to an overfield extends the coefficient domain; it does not assert that every overring is a field. The embedded field P and the ring being studied have a common identity by explicit assumption. Morphological symmetry is a linguistic advantage, not a substitute for those conditions.\par
\noindent{\footnotesize Izvorne zapisy / Original records: HIST-g119-005.}\par
\smallskip\noindent\emph{Sravnjenje s tekstom:} Simetrične nazvy ostajut svezane s različenymi objektami.\par
\noindent\emph{Textual comparison:} The symmetrical names remain attached to distinct objects.\par
\noindent{\footnotesize TEXT-COMP-850; pasaže / passages: 17314, 17316, 17320, 17322, 17324; \texttt{CORPUS\_FORM\_COMPARISONS.json}.}\par
\end{minipage}\par
\par\medskip\Needspace{7\baselineskip}
\noindent\begin{minipage}{\linewidth}
\hypertarget{TERM-1794}{}
\noindent\textbf{podkolca, ktore jednovrěmenno sut idealy}\quad{\footnotesize\texttt{TERM-1794}}\par
\noindent\emph{Iztočnik / Source:} Unterringe, die zugleich Ideale sind\par
\noindent\emph{Rusky / Russian:} подкольца, которые одновременно являются идеалами\par
\noindent\emph{Ukrajinsky / Ukrainian:} підкільця, які водночас є ідеалами\par
\smallskip\noindent\emph{Razlog:} Podkolca, ktore jednovrěmenno sut idealy sohranjaje dva uslova odnosa k obhvačajučemu prstenu. Aktualno podprsteny, ktore Jednočasno sut idealy znači ne samo zatvorjenost pod vlastnymi operacijami, ale i idealnu ulogu pri množenju elementami polnogo prstena. V dokazu blokove matrice tvorjat direktnu sumu, a miješany produkt izčezaje. Prosta zaměna frazy nazvoju abstraktnogo direktnogo produkta skryla by tu konstrukciju. Ranějši korenj kolco ostavaje historijno dokumentovany, ne proglaseny aktualnym.\par
\noindent\emph{Reason:} Podkolca, ktore jednovrěmenno sut idealy preserves two conditions relative to the containing ring. Current podprsteny, ktore Jednočasno sut idealy requires not only closure under the subrings' own operations but their ideal role under multiplication by the full ring. In the proof the block matrices form a direct sum and the mixed product vanishes. Replacing the phrase with the name of an abstract direct product would hide that construction. The earlier stem kolco remains historically documented, not presented as current.\par
\noindent{\footnotesize Izvorne zapisy / Original records: HIST-g138-004.}\par
\smallskip\noindent\emph{Sravnjenje s tekstom:} Blokovy dokaz trěbuje oběh svojstv.\par
\noindent\emph{Textual comparison:} The block-matrix proof needs both properties.\par
\noindent{\footnotesize TEXT-COMP-851; pasaže / passages: 18140, 18144, 18148, 18150, 18152; \texttt{CORPUS\_FORM\_COMPARISONS.json}.}\par
\end{minipage}\par
\par\medskip\Needspace{7\baselineskip}
\noindent\begin{minipage}{\linewidth}
\hypertarget{TERM-1796}{}
\noindent\textbf{podrazděljenje}\quad{\footnotesize\texttt{TERM-1796}}\par
\noindent\emph{Iztočnik / Source:} Unterspaltung\par
\noindent\emph{Rusky / Russian:} последовательное расщепление\par
\noindent\emph{Ukrajinsky / Ukrainian:} послідовне розщеплення\par
\smallskip\noindent\emph{Razlog:} Podrazděljenje označaje postupno razdrobljenje komponentov: četyri vida razloženja pojavjajut se jedno iz drugogo. Ruski i ukrajinski komparandy podčrkujut poslědovateljnost. Kratka nazva ne znači proizvoljno poděljenje množiny ani děljenje čisla. Pasaža osobno razlikuje produktno prědstavjenje od treh redukovanyh prědstavjenij kako najmenša obšča mnogokratna. Zato podrazděljenje ne směje navesti na mysl, že vse četyri sut isti vid faktorizacije.\par
\noindent\emph{Reason:} Podrazděljenje denotes successive refinement of components: the four kinds of decomposition arise from one another. The Russian and Ukrainian comparanda emphasize succession. The short name means neither an arbitrary partition of a set nor numerical division. The passage separately distinguishes a product representation from three reduced least-common-multiple representations. Refinement must therefore not suggest that all four are the same kind of factorization.\par
\noindent{\footnotesize Izvorne zapisy / Original records: HIST-g081-014, HIST-g082-014.}\par
\smallskip\noindent\emph{Sravnjenje s tekstom:} Produkt i najmenše obšče mnogokratno ostajut različne operacije.\par
\noindent\emph{Textual comparison:} Product and least common multiple remain different operations.\par
\noindent{\footnotesize TEXT-COMP-853; pasaže / passages: 12424, 12432, 12434; \texttt{CORPUS\_FORM\_COMPARISONS.json}.}\par
\end{minipage}\par
\par\medskip\Needspace{7\baselineskip}
\noindent\begin{minipage}{\linewidth}
\hypertarget{TERM-1797}{}
\noindent\textbf{klasy nerazložimyh idealov}\quad{\footnotesize\texttt{TERM-1797}}\par
\noindent\emph{Iztočnik / Source:} unzerlegbare Idealklassen\par
\noindent\emph{Rusky / Russian:} классы неразложимых идеалов\par
\noindent\emph{Ukrajinsky / Ukrainian:} класи нерозкладних ідеалів\par
\smallskip\noindent\emph{Razlog:} Klasy nerazložimyh idealov směřuje čitatelja k idealam v klasi izomorfizma i k njihovoj nerazložimosti. V specializaciji na polno reducibilne sistemy te klasy odgovarjajut ireducibilnym klasam prědstavjenij; dvustranne nerazložime idealy osobno svjazany sut s karakterami. To ne jest nazva klas prostyh idealov v komutativnoj aritmetice. Zapisana motivacija brani imenno to smyslovo miješanje. Razlika medžu klasom i njegovym konkretnym prědstaviteljem ostavaje čest objasnjenja.\par
\noindent\emph{Reason:} Klasy nerazložimyh idealov directs the reader to ideals in an isomorphism class and their indecomposability. In the specialization to completely reducible systems these classes correspond to irreducible representation classes; two-sided indecomposable ideals are separately related to characters. This is not the name of classes of prime ideals in commutative arithmetic. Preventing that confusion is precisely the recorded motivation. The distinction between a class and a particular representative remains part of the explanation.\par
\noindent{\footnotesize Izvorne zapisy / Original records: HIST-g105-008.}\par
\smallskip\noindent\emph{Sravnjenje s tekstom:} Korespondencija ne jest definicija prostogo ideala.\par
\noindent\emph{Textual comparison:} The correspondence is not a definition of prime ideal.\par
\noindent{\footnotesize TEXT-COMP-854; pasaže / passages: 15950; \texttt{CORPUS\_FORM\_COMPARISONS.json}.}\par
\end{minipage}\par
\par\medskip\Needspace{7\baselineskip}
\noindent\begin{minipage}{\linewidth}
\hypertarget{TERM-1798}{}
\noindent\textbf{nerazložime identičnosti}\quad{\footnotesize\texttt{TERM-1798}}\par
\noindent\emph{Iztočnik / Source:} unzerlegbare Identitäten\par
\noindent\emph{Rusky / Russian:} неразложимые тождества\par
\noindent\emph{Ukrajinsky / Ukrainian:} нерозкладні тотожності\par
\smallskip\noindent\emph{Razlog:} Nerazložime identičnosti prodolžaje družinu razklad i nerazložimost, oddělenu od redukcije i reducibilnosti. V n-arnoj teoriji pytanje jest najdti polnu množinu identičnostij pri razgledanju simboličnyh redov kako nerazložimyh. To ne jest tvrdženje o nerazložimosti odnogo polinoma na faktory. Pasaže osobno govorjat o razkladnoj identičnosti i redukcijnyh teoremah; splošna zaměna vseh trěh slov s redukcijeju udalila by razliku medžu zadaniami.\par
\noindent\emph{Reason:} Nerazložime identičnosti continues the family of decomposition and indecomposability, kept separate from reduction and reducibility. In the n-ary theory the problem is to find the complete family of identities while treating the symbolic rows as indecomposable. This is not a claim that a single polynomial cannot be factored. The passages separately discuss a decomposition identity and reduction theorems; replacing all three expressions indiscriminately with reduction would remove a distinction between the problems.\par
\noindent{\footnotesize Izvorne zapisy / Original records: HIST-g028-014, HIST-g029-010.}\par
\smallskip\noindent\emph{Sravnjenje s tekstom:} Različne zadanja ne treba sliti v jedno slovo.\par
\noindent\emph{Textual comparison:} The different problems should not be collapsed into one word.\par
\noindent{\footnotesize TEXT-COMP-855; pasaže / passages: 3662, 3705, 3707, 3713; \texttt{CORPUS\_FORM\_COMPARISONS.json}.}\par
\end{minipage}\par
\par\medskip\Needspace{7\baselineskip}
\noindent\begin{minipage}{\linewidth}
\hypertarget{TERM-1799}{}
\noindent\textbf{nerazložime komponenty}\quad{\footnotesize\texttt{TERM-1799}}\par
\noindent\emph{Iztočnik / Source:} unzerlegbare Komponenten\par
\noindent\emph{Rusky / Russian:} неразложимые компоненты\par
\noindent\emph{Ukrajinsky / Ukrainian:} нерозкладні компоненти\par
\smallskip\noindent\emph{Razlog:} Nerazložime komponenty označaje komponenty direktnogo razloženja prstena. Pri ukazanyh hipotezah one nad algebraično zatvorjenym poljem stavajut prstenami ranga jeden, izomorfnymi k tomu polju. Že za opisano razloženje dosta konečno razširiti osnovno polje, jest dodatny matematičny rezultat, ne čest definicije slova komponenta. Arhiv sohranjaje medžunarodny substantiv i razlikuje nerazložimost komponentov od ireducibilnosti polinoma. Ne može se iz samoj nazvy zaključiti, že vsaka komponenta libovoljnogo prstena jest polje.\par
\noindent\emph{Reason:} Nerazložime komponenty denotes components in a direct decomposition of a ring. Under the stated hypotheses they become rank-one rings isomorphic to the algebraically closed field after extension. The fact that a finite extension already suffices for the described decomposition is an additional mathematical result, not part of the word component's definition. The archive retains the international noun and distinguishes component indecomposability from polynomial irreducibility. The label alone does not make every component of every ring a field.\par
\noindent{\footnotesize Izvorne zapisy / Original records: HIST-g129-005.}\par
\smallskip\noindent\emph{Sravnjenje s tekstom:} Konečno razširjenje i polna razložimost sut čest razgledanja.\par
\noindent\emph{Textual comparison:} Finite extension and complete decomposability belong to the argument.\par
\noindent{\footnotesize TEXT-COMP-856; pasaže / passages: 17781, 17783, 17785, 17788, 17790, 17792, 17794; \texttt{CORPUS\_FORM\_COMPARISONS.json}.}\par
\end{minipage}\par
\par\medskip\Needspace{7\baselineskip}
\noindent\begin{minipage}{\linewidth}
\hypertarget{TERM-1800}{}
\noindent\textbf{nerazložime komponenty / polja izomorfne s $\overline P$}\quad{\footnotesize\texttt{TERM-1800}}\par
\noindent\emph{Iztočnik / Source:} unzerlegbare Komponenten / zu $\overline P$ isomorphe Körper\par
\noindent\emph{Rusky / Russian:} неразложимые компоненты / поля, изоморфные $\overline P$\par
\noindent\emph{Ukrajinsky / Ukrainian:} нерозкладні компоненти / поля, ізоморфні $\overline P$\par
\smallskip\noindent\emph{Razlog:} Nerazložime komponenty i polja izomorfne s $\overline P$ opisujut iste objekte samo pod uvedenym uslovjem polnoj razložimosti prvogo roda po razširjenju koeficientov. Pasaža potom svjazuje to s diskriminantnym idealom. Izomorfno ne znači, že komponenty sut bukvalno jedno polje ili imajut obšču jedinicu; vsaka komponenta ima svoju kopiju osnovnogo polja. Arhivna fraza zato mora sohranjati i pridavnik, i uslovje izomorfizma, a matematično nadčrto treba prikazati čitljivo.\par
\noindent\emph{Reason:} Nerazložime komponenty and fields isomorphic to $\overline P$ describe the same objects only under the stated first-kind complete-decomposability condition after coefficient extension. The passage then relates this to the discriminant ideal. Isomorphic does not mean that the components are literally one field or share a single identity; each component has its copy of the base field. The archived phrase must therefore preserve both the adjective and the isomorphism condition, with the mathematical overbar displayed legibly.\par
\noindent{\footnotesize Izvorne zapisy / Original records: HIST-g145-003.}\par
\smallskip\noindent\emph{Sravnjenje s tekstom:} Komponentna jedinica i izomorfizm ne značet bukvalnu identičnost komponentov.\par
\noindent\emph{Textual comparison:} A component identity and an isomorphism do not make the components literally identical.\par
\noindent{\footnotesize TEXT-COMP-857; pasaže / passages: 18441; \texttt{CORPUS\_FORM\_COMPARISONS.json}.}\par
\end{minipage}\par
\par\medskip\Needspace{7\baselineskip}
\noindent\begin{minipage}{\linewidth}
\hypertarget{TERM-1801}{}
\noindent\textbf{prvotno razkladanje}\quad{\footnotesize\texttt{TERM-1801}}\par
\noindent\emph{Iztočnik / Source:} ursprüngliche Zerlegung\par
\noindent\emph{Rusky / Russian:} исходное разложение\par
\noindent\emph{Ukrajinsky / Ukrainian:} початкове розкладання\par
\smallskip\noindent\emph{Razlog:} Prvotno razkladanje znači početkovy izbor razloženja na redy, od ktorogo končny eksplicitny izraz ne zavisi. Pasaža sravnjuje puti prěz redy $y$ i $v$ i razděljaje vremenny rozděl posrědnyh izrazov od jedinogo zaključnogo izraza. Prvotno ne znači prvobytno ili prvotno prosto; označaje město v dokazu. Arhivny izbor čuva tu poslědovateljnost i ne tvrdi nezavisnost od vseh drugih hipotez.\par
\noindent\emph{Reason:} Prvotno razkladanje means the initial choice of decomposition into rows, on which the final explicit expression does not depend. The passage compares routes through the $y$ and $v$ rows, distinguishing the intermediate difference from their common final expression. Initial does not mean primitive or prime; it marks a position in the argument. The recorded choice preserves that sequence without asserting independence from every other hypothesis.\par
\noindent{\footnotesize Izvorne zapisy / Original records: HIST-g034-007.}\par
\smallskip\noindent\emph{Sravnjenje s tekstom:} Nezavisnost kasaje se uvedenogo razloženja, ne vseh prědpoloženij.\par
\noindent\emph{Textual comparison:} Independence concerns the introduced decomposition, not every assumption.\par
\noindent{\footnotesize TEXT-COMP-858; pasaže / passages: 4278; \texttt{CORPUS\_FORM\_COMPARISONS.json}.}\par
\end{minipage}\par
\par\medskip\Needspace{7\baselineskip}
\noindent\begin{minipage}{\linewidth}
\hypertarget{TERM-1803}{}
\noindent\textbf{ręd prěměnnyh}\quad{\footnotesize\texttt{TERM-1803}}\par
\noindent\emph{Iztočnik / Source:} Variabelnreihe\par
\noindent\emph{Rusky / Russian:} строка переменных\par
\noindent\emph{Ukrajinsky / Ukrainian:} рядок змінних\par
\smallskip\noindent\emph{Razlog:} Ręd prěměnnyh označaje red koordinat, ne čislovy red so summoju. Aktualna definicija dava za $x$ konečno mnogo elementov $x_1,\ldots,x_n$, a za $p_\rho$ elementy utvorjene iz minorov. Arhiv izrično upominaje riziko, že ręd ima takože smysl serije ili porjadka. Zato slovnik podaje konkretno koordinatno značenje i svjazuje pravopis ręd s aktualnym red. Ta glosa ne menja simbolov niti ih transformacijnogo zakona.\par
\noindent\emph{Reason:} Ręd prěměnnyh denotes a row of coordinates, not a numerical series to be summed. The current definition gives finitely many entries $x_1,\ldots,x_n$ for $x$ and entries formed from minors for $p_\rho$. The archive explicitly notes that ręd also carries series or order senses. The dictionary therefore supplies the concrete coordinate meaning and links archived ręd with current red. The gloss changes neither the symbols nor their transformation law.\par
\noindent{\footnotesize Izvorne zapisy / Original records: HIST-g030-002.}\par
\smallskip\noindent\emph{Sravnjenje s tekstom:} Definicija ne jest definicija čislovoj serije.\par
\noindent\emph{Textual comparison:} The definition is not that of a numerical series.\par
\noindent{\footnotesize TEXT-COMP-860; pasaže / passages: 3732, 3740, 3742; \texttt{CORPUS\_FORM\_COMPARISONS.json}.}\par
\end{minipage}\par
\par\medskip\Needspace{7\baselineskip}
\noindent\begin{minipage}{\linewidth}
\hypertarget{TERM-1804}{}
\noindent\textbf{rędy prěměnnyh / rędy simbolov vyšego nivoja}\quad{\footnotesize\texttt{TERM-1804}}\par
\noindent\emph{Iztočnik / Source:} Variabelnreihen / Symbolreihen höherer Stufe\par
\noindent\emph{Rusky / Russian:} строки переменных / строки символов более высокого уровня\par
\noindent\emph{Ukrajinsky / Ukrainian:} рядки змінних / символьні рядки вищого рівня\par
\smallskip\noindent\emph{Razlog:} Rędy prěměnnyh i rędy simbolov vyšego nivoja razlikujut koordinatne objekte i simbolične objekte, ktore transformujut se na odgovarjajuči način. Vyšši nivoj kasaje se konstrukcije prěz minory, ne veče težkosti teksta ili vyšego stepena vsakoj formy. Arhivna motivacija čuva matričny kontekst, ale prosta glosa red matrice ne izčerpava vse simbolične ulogy. Pasaža s $p_\rho$ i $S^\rho$ pokazuje, začto obě nazvy morajut ostati različne.\par
\noindent\emph{Reason:} Rędy prěměnnyh and rędy simbolov vyšego nivoja distinguish coordinate objects from symbolic objects with corresponding transformation behaviour. Higher level concerns their construction through minors, not greater textual difficulty or a higher degree of every form. The archived motivation preserves the matrix context, but the simple gloss matrix row does not exhaust every symbolic role. The passage involving $p_\rho$ and $S^\rho$ shows why the two names must remain distinct.\par
\noindent{\footnotesize Izvorne zapisy / Original records: HIST-g029-008.}\par
\smallskip\noindent\emph{Sravnjenje s tekstom:} Simboly i koordinaty trěbujut svjazanu, ale ne identičnu glosu.\par
\noindent\emph{Textual comparison:} Symbols and coordinates require related, but not identical, glosses.\par
\noindent{\footnotesize TEXT-COMP-861; pasaže / passages: 3701, 3703; \texttt{CORPUS\_FORM\_COMPARISONS.json}.}\par
\end{minipage}\par
\par\medskip\Needspace{7\baselineskip}
\noindent\begin{minipage}{\linewidth}
\hypertarget{TERM-1805}{}
\noindent\textbf{rędy prěměnnyh / rędy simbolov}\quad{\footnotesize\texttt{TERM-1805}}\par
\noindent\emph{Iztočnik / Source:} Variablenreihen / Symbolreihen\par
\noindent\emph{Rusky / Russian:} строки переменных / строки символов\par
\noindent\emph{Ukrajinsky / Ukrainian:} рядки змінних / символьні рядки\par
\smallskip\noindent\emph{Razlog:} Rędy prěměnnyh i rędy simbolov tvoret spojen terminologičny par. V Paper 03 elementy prvogo višego reda sut minory iz redov koordinat, a simbolične elementy ponašajut se kako minory iz odgovarjajučih redov. To opravdavaje red kako prostornu i algebraičnu glosu, ne beskonečnu seriju. Sohranjena ograda o mnogoznačnosti ręd ostavaje; rusko stroka i ukrajinsko rjadok ne dokazujut sami po sebě bezproblemnu razumělivost medžuslovjanskoj formy.\par
\noindent\emph{Reason:} Rędy prěměnnyh and rędy simbolov form a linked terminology pair. In Paper 03 the entries of the higher coordinate row are minors of coordinate rows, while symbolic entries behave like minors of corresponding rows. This supports row as a spatial and algebraic gloss, not an infinite series. The preserved reservation about ręd remains; Russian stroka and Ukrainian rjadok do not by themselves establish effortless comprehension of the Interslavic form.\par
\noindent{\footnotesize Izvorne zapisy / Original records: HIST-g028-010.}\par
\smallskip\noindent\emph{Sravnjenje s tekstom:} Matričny kontekst pojasnjaje mnogoznačno slovo red.\par
\noindent\emph{Textual comparison:} The matrix context disambiguates red.\par
\noindent{\footnotesize TEXT-COMP-862; pasaže / passages: 3651, 3653; \texttt{CORPUS\_FORM\_COMPARISONS.json}.}\par
\end{minipage}\par
\par\medskip\Needspace{7\baselineskip}
\noindent\begin{minipage}{\linewidth}
\hypertarget{TERM-1807}{}
\noindent\textbf{kompozitum polj}\quad{\footnotesize\texttt{TERM-1807}}\par
\noindent\emph{Iztočnik / Source:} Vereinigungskörper\par
\noindent\emph{Rusky / Russian:} композит полей\par
\noindent\emph{Ukrajinsky / Ukrainian:} композитум полів\par
\smallskip\noindent\emph{Razlog:} Kompozitum polj jest sovrěmenna nazva najmenšego polja, sodrživajučego razgledane polja v obščem obhvačajučem polju. V primětke ono tvorit se iz Galoisovyh razširjenij, potrěbnyh za razpad polinomov pojedinyh komponentov. Bukvalno sjedinjenje množin ne mora samo tvoriti polje, zato kompozitum pojasnjaje německu nazvu Vereinigungskörper. Izbor medžunarodnogo termina je sohranjena modernizacija; njego ne treba predstavjati kako novy matematičny objekt, izumjeny prěvodom.\par
\noindent\emph{Reason:} Kompozitum polj is the modern name for the smallest field containing the specified fields within a common overfield. In the footnote it is formed from the Galois extensions needed to split the polynomials of the separate components. A literal set union need not itself be a field, so compositum clarifies Vereinigungskörper. The international term is a recorded modernization, not a new mathematical object invented by the translation.\par
\noindent{\footnotesize Izvorne zapisy / Original records: HIST-g129-006.}\par
\smallskip\noindent\emph{Sravnjenje s tekstom:} Ne ide o prosto množinno sjedinjenje.\par
\noindent\emph{Textual comparison:} This is not a mere set union.\par
\noindent{\footnotesize TEXT-COMP-864; pasaže / passages: 17794; \texttt{CORPUS\_FORM\_COMPARISONS.json}.}\par
\end{minipage}\par
\par\medskip\Needspace{7\baselineskip}
\noindent\begin{minipage}{\linewidth}
\hypertarget{TERM-1808}{}
\noindent\textbf{ideal suženja}\quad{\footnotesize\texttt{TERM-1808}}\par
\noindent\emph{Iztočnik / Source:} Verengungsideal\par
\noindent\emph{Rusky / Russian:} идеал сужения\par
\noindent\emph{Ukrajinsky / Ukrainian:} ідеал звуження\par
\smallskip\noindent\emph{Razlog:} Ideal suženja označaje prěsěk ideala s ukazanym podprstenom. Tu jednostranny ideal nekomutativnogo prstena sužuje se do centra, i prěsěk jest modul nad centrom. Ideal razširjenja ide v obratnom směru i v toj konstrukciji jest dvustranny. Jezyčna para suženje i razširjenje jasno pokazyvaje směry, ale ne garantuje, že operacije vsude sut vzajemno obratne. Arhivna ograda o novoj tehničnoj kolokaciji ostavaje lingvističnoju nejistotoju.\par
\noindent\emph{Reason:} Ideal suženja denotes intersection of an ideal with the specified subring. Here a one-sided ideal of a noncommutative ring is contracted to the centre, and the intersection is a module over that centre. Extension goes in the other direction and yields a two-sided ideal in this construction. The linguistic pair contraction and extension makes the directions visible, but does not guarantee that the operations are everywhere inverse. The archived reservation about the technical collocation remains a linguistic uncertainty.\par
\noindent{\footnotesize Izvorne zapisy / Original records: HIST-g189-003.}\par
\smallskip\noindent\emph{Sravnjenje s tekstom:} Jednostranny izvorni ideal i dvustranno razširjenje trěbujut različenje.\par
\noindent\emph{Textual comparison:} The one-sided original ideal and two-sided extension must be distinguished.\par
\noindent{\footnotesize TEXT-COMP-865; pasaže / passages: 20724, 20726, 20728, 20733; \texttt{CORPUS\_FORM\_COMPARISONS.json}.}\par
\end{minipage}\par
\par\medskip\Needspace{7\baselineskip}
\noindent\begin{minipage}{\linewidth}
\hypertarget{TERM-1809}{}
\noindent\textbf{suženy ideal}\quad{\footnotesize\texttt{TERM-1809}}\par
\noindent\emph{Iztočnik / Source:} Verengungsideal\par
\noindent\emph{Rusky / Russian:} идеал сужения\par
\noindent\emph{Ukrajinsky / Ukrainian:} ідеал звуження\par
\smallskip\noindent\emph{Razlog:} Suženy ideal jest pridavnikova forma toj samoj operacije, ktora indže nazvana jest ideal suženja. V Paper 31 definicija jest prěsěk ideala nadprstena s osnovnym prstenom; neposrědny dokaz vračaje izvorni ideal poslě razširjenja koeficientnogo polja. To svjazano jest s konkretnoju konstrukcijeju razširjenja, ne s vsemi vloženjami prstenov. Slovnik čuva oba formulovanja i razlikuje operaciju od njenogo rezultata bez tvorjenja dvuh različnih pojmov.\par
\noindent\emph{Reason:} Suženy ideal is the adjectival form of the operation elsewhere called ideal suženja. In Paper 31 it is defined as intersection of an overring ideal with the base ring; the following proof recovers the original ideal after coefficient-field extension. That result depends on the specific extension construction, not on every ring embedding. The dictionary retains both formulations and distinguishes the operation from its result without creating two different concepts.\par
\noindent{\footnotesize Izvorne zapisy / Original records: HIST-g120-002.}\par
\smallskip\noindent\emph{Sravnjenje s tekstom:} Dokaz vračanja ideala ima konkretne hipotezy.\par
\noindent\emph{Textual comparison:} Recovery of the ideal is proved under specific hypotheses.\par
\noindent{\footnotesize TEXT-COMP-866; pasaže / passages: 17394, 17396; \texttt{CORPUS\_FORM\_COMPARISONS.json}.}\par
\end{minipage}\par
\par\medskip\Needspace{7\baselineskip}
\noindent\begin{minipage}{\linewidth}
\hypertarget{TERM-1811}{}
\noindent\textbf{operacije / vrsty spojivanja: složenje i množenje}\quad{\footnotesize\texttt{TERM-1811}}\par
\noindent\emph{Iztočnik / Source:} Verknüpfungsarten: Addition und Multiplikation\par
\noindent\emph{Rusky / Russian:} операции / виды соединения: сложение и умножение\par
\noindent\emph{Ukrajinsky / Ukrainian:} операції / способи сполучення: додавання і множення\par
\smallskip\noindent\emph{Razlog:} Operacije (vrsty spojivanja), složenje i množenje pojasnjajut starše Verknüpfungsarten, ne izpustivši jeho smysl. Dve binarne operacije davajut jednoznačno sumu i produkt elementov; potom slědujut zakony prstena. To ne označaje skladyvanje funkcij ni topologično sklejanje. Aktualna pasaža čuva obě glosy vměstě. Izbor sovrěmennogo operacije robi definiciju neposrědnějše čitljivoju, ale ne dodaje obratimost množenja ili druge zakony, ktore tekst ne požaduje.\par
\noindent\emph{Reason:} Operacije (vrsty spojivanja), složenje i množenje clarifies the older Verknüpfungsarten without discarding its meaning. Two binary operations give a unique sum and product of elements; the ring laws follow. This does not mean composition of functions or topological gluing. The current passage retains the two glosses together. Modern operations makes the definition more immediately readable, but adds neither multiplicative invertibility nor any other law not required by the text.\par
\noindent{\footnotesize Izvorne zapisy / Original records: HIST-g081-019, HIST-g082-019.}\par
\smallskip\noindent\emph{Sravnjenje s tekstom:} Definicija operacij prědhodi konkretnym zakonam prstena.\par
\noindent\emph{Textual comparison:} The definition of the operations precedes the specific ring laws.\par
\noindent{\footnotesize TEXT-COMP-868; pasaže / passages: 12476; \texttt{CORPUS\_FORM\_COMPARISONS.json}.}\par
\end{minipage}\par
\par\medskip\Needspace{7\baselineskip}
\noindent\begin{minipage}{\linewidth}
\hypertarget{TERM-1813}{}
\noindent\textbf{odpovědanje vlastnyh energijskih zakonov / obča relativnost}\quad{\footnotesize\texttt{TERM-1813}}\par
\noindent\emph{Iztočnik / Source:} Versagen eigentlicher Energiesätze / allgemeine Relativität\par
\noindent\emph{Rusky / Russian:} отказ собственных законов энергии / общая относительность\par
\noindent\emph{Ukrajinsky / Ukrainian:} відмова власних законів енергії / загальна відносність\par
\smallskip\noindent\emph{Razlog:} Odpovědanje jest ogradženy ranějši pokus prěvesti Versagen; aktualny tekst ima neuspěh vlastnyh energijskyh zakonov. Smysl ne jest otvrganje energiji ni običny otkaz, ale neobstojanje vlastnyh relacij v tehničnom smislu, gde nevlastne relacije proizhodjat iz zavisnostij i identično bezdivergencijnyh izrazov. Noether izrično razširjaje tu nazvu obča relativnost za grupovo-teoretično tvrdženje. Slovnik ne směje proměniti to v neogradženu fizičnu sloganu ili utvrditi ranějšu jezyčnu formu bez sviděteljstva.\par
\noindent\emph{Reason:} Odpovědanje is a qualified earlier attempt to translate Versagen; the current text uses neuspěh vlastnyh energijskyh zakonov. The sense is not rejection of energy or ordinary refusal, but absence of proper relations in the technical sense in which improper relations arise from dependencies and identically divergence-free expressions. Noether explicitly broadens general relativity here for the group-theoretic statement. The dictionary must not turn this into an unqualified physical slogan or endorse the earlier wording without evidence.\par
\noindent{\footnotesize Izvorne zapisy / Original records: HIST-g075-020.}\par
\smallskip\noindent\emph{Sravnjenje s tekstom:} Aktualno neuspěh i arhivno odpovědanje ne sut odinakovym jezyčnym sviděteljstvom.\par
\noindent\emph{Textual comparison:} Current neuspěh and archived odpovědanje are not the same linguistic evidence.\par
\noindent{\footnotesize TEXT-COMP-870; pasaže / passages: 9293, 9719, 9741, 9790; \texttt{CORPUS\_FORM\_COMPARISONS.json}.}\par
\end{minipage}\par
\par\medskip\Needspace{7\baselineskip}
\noindent\begin{minipage}{\linewidth}
\hypertarget{TERM-1814}{}
\noindent\textbf{grupa posunov}\quad{\footnotesize\texttt{TERM-1814}}\par
\noindent\emph{Iztočnik / Source:} Verschiebungsgruppe\par
\noindent\emph{Rusky / Russian:} группа сдвигов\par
\noindent\emph{Ukrajinsky / Ukrainian:} група зсувів\par
\smallskip\noindent\emph{Razlog:} Grupa posunov imenovana jest po transformaciji $y_i=x_i+\varepsilon_i$, s odgovarjajučim prěnosom funkcij. Posun jest prostorna translacija, ne jezyčny prěvod. Alternativa translacija izrično sohranjena jest v historijnoj primětke za lingvistično sravnjenje; nema sohranjene čislene ocěny prednosti oběh form. Konečna grupa tu znači konečno mnogo parametrov, ne konečno mnogo posunov. Definicija i energetska priměna muset ostati pri objasnjenju nazvy.\par
\noindent\emph{Reason:} Grupa posunov is named for the transformation $y_i=x_i+\varepsilon_i$, with the corresponding transport of functions. Posun is a spatial translation, not translation between languages. The alternative translacija is explicitly preserved in the historical note for linguistic comparison; no numerical ranking of the two survives. Finite group here means finitely many parameters, not finitely many shifts. The definition and energy-law application must remain attached to the term.\par
\noindent{\footnotesize Izvorne zapisy / Original records: HIST-g075-018.}\par
\smallskip\noindent\emph{Sravnjenje s tekstom:} Parametrična transformacija pojasnjaje posun i slovo konečna.\par
\noindent\emph{Textual comparison:} The parametrized transformation clarifies both shift and finite.\par
\noindent{\footnotesize TEXT-COMP-871; pasaže / passages: 9723, 9725, 9729, 9730, 9741; \texttt{CORPUS\_FORM\_COMPARISONS.json}.}\par
\end{minipage}\par
\par\medskip\Needspace{7\baselineskip}
\noindent\begin{minipage}{\linewidth}
\hypertarget{TERM-1815}{}
\noindent\textbf{skr??eno predstavjenje [Čestično nezachovane bukvy / Partly unrecoverable lettering]}\quad{\footnotesize\texttt{TERM-1815}}\par
\noindent\emph{Iztočnik / Source:} verschr?nkte Darstellung [Čestično nezachovane bukvy / Partly unrecoverable lettering]\par
\noindent\emph{Rusky / Russian:} [Bukvy ne sut zachovane / Lettering not recoverable]\par
\noindent\emph{Ukrajinsky / Ukrainian:} [Bukvy ne sut zachovane / Lettering not recoverable]\par
\smallskip\noindent\emph{Razlog:} Ošteđeny zapis skr??eno predstavjenje ne dopušča tihu rekonstrukciju vseh izvornyh bukv. Aktualna pasaža čitljivo ima skrěženo prědstavjenje i definira jego relacijeju $C_S^{\,T}C_T=C_{ST}a_{S,T}$. Grupovo dějstvo na koeficientah jest čest formulace; ne jest to obyčno matrično prědstavjenje s opuščenoju nadpisanoju operacijeju. Zapisan naměr čuva družinu so skrěženym produktom. Ruski i ukrajinski zapis v tom arhivu ošteđeny sut previše, aby služili samostojnym leksičnym sviděteljstvom.\par
\noindent\emph{Reason:} The damaged record skr??eno predstavjenje does not permit silent reconstruction of every original letter. The current passage legibly has skrěženo prědstavjenje and defines it by $C_S^{\,T}C_T=C_{ST}a_{S,T}$. The group action on coefficients belongs to the formula; this is not an ordinary matrix representation with its superscript operation omitted. The recorded intention preserves the family shared with crossed product. The Russian and Ukrainian strings in this record are too damaged to serve as independent lexical evidence.\par
\noindent{\footnotesize Izvorne zapisy / Original records: HIST-g211-006.}\par
\smallskip\noindent\emph{Sravnjenje s tekstom:} Ošteđene komparandy ne sut rekonstruovane po nagađanju.\par
\noindent\emph{Textual comparison:} The damaged comparanda have not been guessed back into existence.\par
\noindent{\footnotesize TEXT-COMP-872; pasaže / passages: 24268, 24270, 24272, 24274, 24276, 24290; \texttt{CORPUS\_FORM\_COMPARISONS.json}.}\par
\end{minipage}\par
\par\medskip\Needspace{7\baselineskip}
\noindent\begin{minipage}{\linewidth}
\hypertarget{TERM-1816}{}
\noindent\textbf{skr??eny produkt [Čestično nezachovane bukvy / Partly unrecoverable lettering]}\quad{\footnotesize\texttt{TERM-1816}}\par
\noindent\emph{Iztočnik / Source:} verschr?nktes Produkt [Čestično nezachovane bukvy / Partly unrecoverable lettering]\par
\noindent\emph{Rusky / Russian:} [Bukvy ne sut zachovane / Lettering not recoverable]\par
\noindent\emph{Ukrajinsky / Ukrainian:} [Bukvy ne sut zachovane / Lettering not recoverable]\par
\smallskip\noindent\emph{Razlog:} Skr??eny produkt jest ošteđena historijna forma; čitljiva aktualna forma jest skrěženy produkt. V konstrukciji spolno učestvujut Galoisovo polje, njegova grupa i faktorny sistem: $u_Su_T=u_{ST}a_{S,T}$ i pravilo za dějstvo na elementah polja. Zato produkt ne znači kartezijansky produkt dvuh množin ili prosto množenje dvuh čisel. Sohranjena motivacija čuva slovjansku kalku i ogradženu novu terminologiju. Čitljiva korpusna definicija dopolnjaje, ale ne popravjaje potajno ošteđene arhivne komparandy.\par
\noindent\emph{Reason:} Skr??eny produkt is a damaged historical form; the readable current form is skrěženy produkt. The construction combines a Galois field, its group and a factor system: $u_Su_T=u_{ST}a_{S,T}$ and the rule for acting on field elements. Product therefore means neither a Cartesian product of two sets nor simple multiplication of two numbers. The preserved rationale records a Slavic calque and qualified new terminology. The readable corpus definition supplements, but does not silently repair, the damaged archived comparanda.\par
\noindent{\footnotesize Izvorne zapisy / Original records: HIST-g209-004, HIST-g211-005.}\par
\smallskip\noindent\emph{Sravnjenje s tekstom:} Dějstvo i faktorne veličiny sut čest definicije.\par
\noindent\emph{Textual comparison:} The action and factor quantities belong to the definition.\par
\noindent{\footnotesize TEXT-COMP-873; pasaže / passages: 23243, 23246, 23251, 23253, 23257, 23261, 23265, 24194; \texttt{CORPUS\_FORM\_COMPARISONS.json}.}\par
\end{minipage}\par
\par\medskip\Needspace{7\baselineskip}
\noindent\begin{minipage}{\linewidth}
\hypertarget{TERM-1817}{}
\noindent\textbf{skrežene produkty}\quad{\footnotesize\texttt{TERM-1817}}\par
\noindent\emph{Iztočnik / Source:} verschrenkte Produkte\par
\noindent\emph{Rusky / Russian:} скрещенные произведения\par
\noindent\emph{Ukrajinsky / Ukrainian:} схрещені добутки\par
\smallskip\noindent\emph{Razlog:} Skrežene produkty jest množinska bibliografična forma, čitljiva tam, kde něktore ranějše terminske kartky imajut ošteđene bukvy. Nazvanje raboty svjazuje razložene skrěžene produkty s njihovymi maksimalnymi poredkami. To podpiraje identifikaciju terminologičnoj družiny, ale samo bibliografično zaglavje ne dokazuje ni definiciju ni jezyčnu normativnost. Razlika skrežene i skrěžene dolžna byti vidliva kako pravopisna variacija. Arhivna ograda o novoj matematičnoj terminologiji ostavaje.\par
\noindent\emph{Reason:} Skrežene produkty is the readable plural bibliographic form where some earlier term records have damaged lettering. The paper title links split crossed products with their maximal orders. This supports identification of the term family, but a bibliographic title alone proves neither a definition nor linguistic normativity. The difference between skrežene and skrěžene should remain visible as spelling variation. The archived reservation about the mathematical terminology remains.\par
\noindent{\footnotesize Izvorne zapisy / Original records: HIST-g215-010.}\par
\smallskip\noindent\emph{Sravnjenje s tekstom:} Zaglavje svjazano jest s terminom, ale ne zaměnjaje jego definiciju.\par
\noindent\emph{Textual comparison:} The title connects to the term but does not replace its definition.\par
\noindent{\footnotesize TEXT-COMP-874; pasaže / passages: bib-isv.tex:156, bib-isv.tex:157, 23265; \texttt{CORPUS\_FORM\_COMPARISONS.json}.}\par
\end{minipage}\par
\par\medskip\Needspace{7\baselineskip}
\noindent\begin{minipage}{\linewidth}
\hypertarget{TERM-1818}{}
\noindent\textbf{nulovy determinant}\quad{\footnotesize\texttt{TERM-1818}}\par
\noindent\emph{Iztočnik / Source:} verschwindende Determinante\par
\noindent\emph{Rusky / Russian:} нулевой детерминант\par
\noindent\emph{Ukrajinsky / Ukrainian:} нульовий детермінант\par
\smallskip\noindent\emph{Razlog:} Nulovy determinant prěnosi německo verschwindend kako algebraičnu ravnost nuli, ne kako proces pomalogo izčezanja. V pasaži priměnjenje teoremy produkta daje determinant reda $n-\lambda$, ktory jest nulovy. To ne jest determinant nulovogo reda niti tvrđenje, že vse elementy matrice sut nulove. Kratky pridavnik byl vybrany dlja sovrěmennoj čitljivosti; dokaz i razměr determinanta ostajut potrebne za točno čitanje.\par
\noindent\emph{Reason:} Nulovy determinant renders verschwindend as algebraic equality to zero, not a process of gradually disappearing. Applying the product theorem in the passage yields a determinant of order $n-\lambda$ which is zero. This is neither a determinant of order zero nor a claim that every matrix entry is zero. The short adjective was chosen for modern readability; the proof and determinant's size remain necessary to interpreting it correctly.\par
\noindent{\footnotesize Izvorne zapisy / Original records: HIST-g031-014.}\par
\smallskip\noindent\emph{Sravnjenje s tekstom:} Red determinanta i jego vrědnost sut različne veličiny.\par
\noindent\emph{Textual comparison:} The determinant's order and its value are different quantities.\par
\noindent{\footnotesize TEXT-COMP-875; pasaže / passages: 3865, 3868, 3870, 3871, 3874; \texttt{CORPUS\_FORM\_COMPARISONS.json}.}\par
\end{minipage}\par
\par\medskip\Needspace{7\baselineskip}
\noindent\begin{minipage}{\linewidth}
\hypertarget{TERM-1819}{}
\noindent\textbf{zaměna simbolov}\quad{\footnotesize\texttt{TERM-1819}}\par
\noindent\emph{Iztočnik / Source:} Vertauschung der Symbole\par
\noindent\emph{Rusky / Russian:} перестановка символов\par
\noindent\emph{Ukrajinsky / Ukrainian:} переставлення символів\par
\smallskip\noindent\emph{Razlog:} Zaměna simbolov označaje perestavjenje simboličnyh imen v danom izrazu. Ona ne znači kategoričny izomorfizm, svobodnu zaměnu libovoljnyh koeficientov ili proměnu soderžanja identičnosti. V pasaži zaměna priměnjaje se k konkretnoj forme s $a_\nu$ i $\theta$. Arhiv podaje široku razumělivost zaměna kako motivaciju, ne kak rezultat testovanja. Polny izraz i identifikacija menjanyh simbolov sut potrěbne, zato samotna obča nazva ne izčerpava operaciju.\par
\noindent\emph{Reason:} Zaměna simbolov denotes interchange of symbolic names in the specified expression. It means neither a categorical isomorphism nor free replacement of arbitrary coefficients or alteration of the identity's content. The passage applies the interchange to a particular form involving $a_\nu$ and $\theta$. The archive offers broad transparency as a motivation, not a measured reader-test result. The full expression and identification of the exchanged symbols are needed; the general label alone does not specify the operation.\par
\noindent{\footnotesize Izvorne zapisy / Original records: HIST-g014-014.}\par
\smallskip\noindent\emph{Sravnjenje s tekstom:} Obča nazva trěbuje konkretnogo argumenta.\par
\noindent\emph{Textual comparison:} The general label requires its actual operand.\par
\noindent{\footnotesize TEXT-COMP-876; pasaže / passages: 2096, 2097, 2098, 2099, 2102; \texttt{CORPUS\_FORM\_COMPARISONS.json}.}\par
\end{minipage}\par
\par\medskip\Needspace{7\baselineskip}
\noindent\begin{minipage}{\linewidth}
\hypertarget{TERM-1820}{}
\noindent\textbf{ideal razvětvenja (osnovny ideal, diferenta)}\quad{\footnotesize\texttt{TERM-1820}}\par
\noindent\emph{Iztočnik / Source:} Verzweigungsideal (Grundideal, Differente)\par
\noindent\emph{Rusky / Russian:} идеал ветвления (основной идеал, дифферента)\par
\noindent\emph{Ukrajinsky / Ukrainian:} ідеал розгалуження (основний ідеал, диферента)\par
\smallskip\noindent\emph{Razlog:} Ideal razvětvenja jest izbran po analogiji s točkoju razvětvenja. Arhivna motivacija govori o različenju od diferenty, ale sama pregledana pasaža podaje osnovny ideal i diferenta v zagradkah kako nazvy togo samogo razgledanogo ideala. To protivrěčje trěba izrično ograničiti: ne možno iz motivacije vyvesti dva različne objekta tam, gde tekst jih ne razlikuje. Norma togo ideala daje diskriminant polja; te pojmy isto ne treba sliti. Razvětvenje jest pojem, a diferenta konkretny ideal v toj formulaciji.\par
\noindent\emph{Reason:} Ideal razvětvenja was chosen by analogy with a branch point. The archived motivation says it separates ramification from the different, but the examined passage itself gives fundamental ideal and different parenthetically as names for the same ideal under discussion. That discrepancy must be qualified explicitly: the motivation cannot establish two distinct objects where the text does not distinguish them. The norm of this ideal gives the field discriminant, which likewise must not be conflated with it. Ramification is the concept; the different is the particular ideal in this formulation.\par
\noindent{\footnotesize Izvorne zapisy / Original records: HIST-g076-011.}\par
\smallskip\noindent\emph{Sravnjenje s tekstom:} Pasaža i arhivna motivacija ne podpirajut isto proste različenje nazv.\par
\noindent\emph{Textual comparison:} The passage and archived motivation do not support the same simple distinction of labels.\par
\noindent{\footnotesize TEXT-COMP-877; pasaže / passages: 9912, 9922, 9926; \texttt{CORPUS\_FORM\_COMPARISONS.json}.}\par
\end{minipage}\par
\par\medskip\Needspace{7\baselineskip}
\noindent\begin{minipage}{\linewidth}
\hypertarget{TERM-1821}{}
\noindent\textbf{idealy razvětvenja}\quad{\footnotesize\texttt{TERM-1821}}\par
\noindent\emph{Iztočnik / Source:} Verzweigungsideale\par
\noindent\emph{Rusky / Russian:} идеалы разветвления\par
\noindent\emph{Ukrajinsky / Ukrainian:} ідеали розгалуження\par
\smallskip\noindent\emph{Razlog:} Idealy razvětvenja upotrěbjajut slovjansky obraz větvi, zajedničny s shranjenymi ruskimi i ukrajinskimi formami. V toj pasaži priměry nad nedoskonalym poljem sut izrično nazvane analogami idealov razvětvenja: uspoređuje se pojavjenje vyših stepenjev prostogo ideala. Analogija ne jest tvrdženje, že vse opisane idealy sut diferenty čislovyh polj. Arhivna ograda o potrěbě jezyčnoj prověrky ostaje; podobnost slov sama ne dokazuje razumělivost.\par
\noindent\emph{Reason:} Idealy razvětvenja uses the Slavic branching metaphor shared with the preserved Russian and Ukrainian forms. The passage explicitly calls the examples over an imperfect field analogues of ramification ideals: it compares the appearance of higher powers of a prime ideal. This analogy does not identify every ideal in those examples with the different of a number field. The historical linguistic-review reservation remains; similarity of words alone does not establish reader comprehension.\par
\noindent{\footnotesize Izvorne zapisy / Original records: HIST-g100-011.}\par
\smallskip\noindent\emph{Sravnjenje s tekstom:} Nazva analogija ne identifikuje vse objekty.\par
\noindent\emph{Textual comparison:} Calling the examples analogous does not identify all the objects.\par
\noindent{\footnotesize TEXT-COMP-878; pasaže / passages: 15714; \texttt{CORPUS\_FORM\_COMPARISONS.json}.}\par
\end{minipage}\par
\par\medskip\Needspace{7\baselineskip}
\noindent\begin{minipage}{\linewidth}
\hypertarget{TERM-1822}{}
\noindent\textbf{m?sta razv?tvjenja [Čestično nezachovane bukvy / Partly unrecoverable lettering]}\quad{\footnotesize\texttt{TERM-1822}}\par
\noindent\emph{Iztočnik / Source:} Verzweigungsstellen\par
\noindent\emph{Rusky / Russian:} [Bukvy ne sut zachovane / Lettering not recoverable]\par
\noindent\emph{Ukrajinsky / Ukrainian:} [Bukvy ne sut zachovane / Lettering not recoverable]\par
\smallskip\noindent\emph{Razlog:} M?sta razv?tvjenja jest čestično ošteđeny historijny zapis. Aktualny tekst dava čitljivo města razvětvjenja, uključajuči konečne i beskonečne města. Město tu jest aritmetičny pojem, ne geometricky položaj: lokalno razloženje algebry svezano jest s normnymi ostatkami i provodnikom. Arhivna motivacija čuva toj lokalno-globalny svez; ona ne dopušča obnoviti izgubljene ruske i ukrajinske bukvy kakoby one byly pročitane.\par
\noindent\emph{Reason:} M?sta razv?tvjenja is a partly damaged historical record. The current text supplies readable města razvětvjenja, including finite and infinite places. Place is an arithmetic notion here, not a geometric location: local splitting of the algebra is linked to norm residues and the conductor. The preserved rationale retains that local–global connection; it does not license reconstruction of the lost Russian and Ukrainian lettering as though it had been read.\par
\noindent{\footnotesize Izvorne zapisy / Original records: HIST-g209-012.}\par
\smallskip\noindent\emph{Sravnjenje s tekstom:} Čitlivy aktualny tekst dopolnjaje ošteđeny zapis, ne obnovjaje jego bukvy.\par
\noindent\emph{Textual comparison:} Readable current text supplements the damaged record without recovering its letters.\par
\noindent{\footnotesize TEXT-COMP-879; pasaže / passages: 23299, 23313, 23315, 23317, 23319; \texttt{CORPUS\_FORM\_COMPARISONS.json}.}\par
\end{minipage}\par
\par\medskip\Needspace{7\baselineskip}
\noindent\begin{minipage}{\linewidth}
\hypertarget{TERM-1823}{}
\noindent\textbf{teorija razvětvenja / ideal razvětvenja}\quad{\footnotesize\texttt{TERM-1823}}\par
\noindent\emph{Iztočnik / Source:} Verzweigungstheorie / Verzweigungsideal\par
\noindent\emph{Rusky / Russian:} теория разветвления / идеал разветвления\par
\noindent\emph{Ukrajinsky / Ukrainian:} теорія розгалуження / ідеал розгалуження\par
\smallskip\noindent\emph{Razlog:} Teorija razvětvenja i ideal razvětvenja tvoret svezanu leksičnu družinu, ale nazyvajut razne stvari: oblast izslědovanja i konkretny ideal. Uvod razširjaje teoriju na poredky, a teorem o diskriminantu predstavjaje kako njejin početok. Pozdnějše točnejše oprěděljenje eksponentov ne možno iz samogo termina proglasiti dokazanym dlja vseho občego slučaja: avtorka tu piše o očekivanoj formě daljšej teorije. Nazva shranjaje kontinuitet s čislovymi i funkcijnymi poljami bez potlačenja toj razliki.\par
\noindent\emph{Reason:} Teorija razvětvenja and ideal razvětvenja form a related lexical family but name different things: a field of inquiry and a particular ideal. The introduction extends the theory to orders and presents the discriminant theorem as its starting point. The term alone cannot turn the subsequent finer determination of exponents into a proved result for the whole general case: the author describes an expected form of the further theory. The vocabulary preserves continuity with number and function fields without suppressing that distinction.\par
\noindent{\footnotesize Izvorne zapisy / Original records: HIST-g118-012.}\par
\smallskip\noindent\emph{Sravnjenje s tekstom:} Formulacija o očekivanoj teoriji ne jest gotovy obči teorem.\par
\noindent\emph{Textual comparison:} The anticipated theory is not a completed general theorem.\par
\noindent{\footnotesize TEXT-COMP-880; pasaže / passages: 17283; \texttt{CORPUS\_FORM\_COMPARISONS.json}.}\par
\end{minipage}\par
\par\medskip\Needspace{7\baselineskip}
\noindent\begin{minipage}{\linewidth}
\hypertarget{TERM-1824}{}
\noindent\textbf{teorija razvětvljenja / ideal razvětvljenja}\quad{\footnotesize\texttt{TERM-1824}}\par
\noindent\emph{Iztočnik / Source:} Verzweigungstheorie / Verzweigungsideal\par
\noindent\emph{Rusky / Russian:} теория ветвления / идеал ветвления\par
\noindent\emph{Ukrajinsky / Ukrainian:} теорія розгалуження / ідеал розгалуження\par
\smallskip\noindent\emph{Razlog:} Razvětvljenje jest shranjena varianta, izbrana zaradi čitljivogo obraza větvenja vměsto čisto internacionalnoj formy. Aktualna pasaža piše razvětvjenje i izrično identifikuje ideal razvětvjenja s diferentoju. Zato slovnik mora dati obě pisanja i jasno označiti aktualnu formu, ne izmišljati dva matematične pojma. Zapis tvrdi sravnjenje s zapadnymi i južnymi slovjanskimi formami, ale sam ne dava točne citaty; to jest shranjeno tvrđenje o metodě, ne novy prověreny dokaz. Eksponenty i izključenje neseparabilnosti ostajut pri definiciji.\par
\noindent\emph{Reason:} Razvětvljenje is a preserved variant chosen for its recognizable branching metaphor instead of a purely international form. The current passage writes razvětvjenje and explicitly identifies the ramification ideal with the different. The dictionary must therefore show both spellings and identify the current form, not invent two mathematical concepts. The record reports comparison with West and South Slavic forms but does not itself supply exact citations; that is a preserved methodological claim, not newly verified evidence. The exponent conditions and inseparability qualification remain attached to the concept.\par
\noindent{\footnotesize Izvorne zapisy / Original records: HIST-g213-013.}\par
\smallskip\noindent\emph{Sravnjenje s tekstom:} Aktualna pasaža nazyvaje ideal razvětvjenja diferentoju.\par
\noindent\emph{Textual comparison:} The current passage calls the ramification ideal the different.\par
\noindent{\footnotesize TEXT-COMP-881; pasaže / passages: 24737, 24775; \texttt{CORPUS\_FORM\_COMPARISONS.json}.}\par
\end{minipage}\par
\par\medskip\Needspace{7\baselineskip}
\noindent\begin{minipage}{\linewidth}
\hypertarget{TERM-1827}{}
\noindent\textbf{uslovje cěpov kratnyh modulo vsaky nenulovy ideal}\quad{\footnotesize\texttt{TERM-1827}}\par
\noindent\emph{Iztočnik / Source:} Vielfachen-Kettensatz mod. jedes von Null verschiedene Ideal\par
\noindent\emph{Rusky / Russian:} условие цепей кратных по модулю каждого ненулевого идеала\par
\noindent\emph{Ukrajinsky / Ukrainian:} умова ланцюгів кратних за модулем кожного ненульового ідеалу\par
\smallskip\noindent\emph{Razlog:} Uslovje cěpov kratnyh modulo vsaky nenulovy ideal čuva i směr dělimosti, i ograničenje: vsaky člen cěpa jest děliteljem fiksovanogo nenulovogo ideala. Ide o prěryvanje vlastnyh kratnyh po konečno mnogyh krokah, ne o libovoljnu cěp bez spodnjej granicy. Modulo jest shranjena internacionalna forma. Sravnjenje s odrěđenoju aksiomoju děliteljev ne dopušča sliti obě uslovja ili izpustiti nenulovy; tako by se izměnili prědpoloženja teorema.\par
\noindent\emph{Reason:} Uslovje cěpov kratnyh modulo vsaky nenulovy ideal preserves both the divisibility direction and the restriction: every member of the chain divides a fixed nonzero ideal. The condition concerns termination of successive proper multiples, not an arbitrary chain without that lower bound. Modulo is the retained international expression. Comparison with the separately stated divisor axiom does not permit merging the two conditions or dropping nonzero; doing so would alter the theorem's assumptions.\par
\noindent{\footnotesize Izvorne zapisy / Original records: HIST-g107-003.}\par
\smallskip\noindent\emph{Sravnjenje s tekstom:} Vsi členy sut děliteli fiksovanogo nenulovogo ideala.\par
\noindent\emph{Textual comparison:} All members divide a fixed nonzero ideal.\par
\noindent{\footnotesize TEXT-COMP-884; pasaže / passages: 16120, 16121, 16122, 16123; \texttt{CORPUS\_FORM\_COMPARISONS.json}.}\par
\end{minipage}\par
\par\medskip\Needspace{7\baselineskip}
\noindent\begin{minipage}{\linewidth}
\hypertarget{TERM-1828}{}
\noindent\textbf{cěp kratnyh}\quad{\footnotesize\texttt{TERM-1828}}\par
\noindent\emph{Iztočnik / Source:} Vielfachenkette\par
\noindent\emph{Rusky / Russian:} цепь кратных\par
\noindent\emph{Ukrajinsky / Ukrainian:} ланцюг кратних\par
\smallskip\noindent\emph{Razlog:} Cěp kratnyh jest druga strana dvojnego uslovja cěpov, ale arhivno objasnjenje, ktoro nazyvaje ju rastučeju, trěbuje opravy. Aktualny pokazany red ide od cělogo modula $G$ prěz podmoduly k nulovomu modulu $E$: po vključanju to jest spadajuči red. Besěda kratnyh ohranjaje historijno ponimanje dělimosti; ona ne znači rastuče vključanje. Dlina i prostost poslědovateljnyh kvocientov sut odrěđene v definiciji kompozicijnogo reda, ne v samom slově cěp.\par
\noindent\emph{Reason:} Cěp kratnyh names the other side of the double-chain condition, but the archived explanation calling it ascending requires correction. The displayed sequence runs from the whole module $G$ through submodules to the zero module $E$: it descends by inclusion. Kratnyh retains the historical divisibility viewpoint; it does not mean ascending inclusion. Length and simplicity of successive quotients are specified in the definition of a composition series, not by the word chain alone.\par
\noindent{\footnotesize Izvorne zapisy / Original records: HIST-g117-007.}\par
\smallskip\noindent\emph{Sravnjenje s tekstom:} Pokazany red spada po vključanju k nulovomu modulu.\par
\noindent\emph{Textual comparison:} The displayed sequence descends by inclusion to the zero module.\par
\noindent{\footnotesize TEXT-COMP-885; pasaže / passages: 17189, 17193, 17195, 17197, 17203, 17205, 17207; \texttt{CORPUS\_FORM\_COMPARISONS.json}.}\par
\end{minipage}\par
\par\medskip\Needspace{7\baselineskip}
\noindent\begin{minipage}{\linewidth}
\hypertarget{TERM-1829}{}
\noindent\textbf{uslovje cěpov kratnyh}\quad{\footnotesize\texttt{TERM-1829}}\par
\noindent\emph{Iztočnik / Source:} Vielfachenkettensatz\par
\noindent\emph{Rusky / Russian:} условие цепей кратных\par
\noindent\emph{Ukrajinsky / Ukrainian:} умова ланцюгів кратних\par
\smallskip\noindent\emph{Razlog:} Uslovje cěpov kratnyh jest uslovje stabilizacije, ne teorem o proizvoljnom čislu cěpov. V pomočnom argumentu cěp stepenjev glavnogo ideala stabilizuje se; potom, zaradi otsutnosti děliteljev nuly, možno skratiti stepenj $a$ i dobiti rěšenje $ax=b$. Komutativnost i otsutnost děliteljev nuly dolžne ostati v tvrdženju. Zapis uvodi kratku nazvu zaradi jej roli v dokazu, ne kakoby samo uslovje cěpov pravilo vsaky prsten poljem.\par
\noindent\emph{Reason:} Uslovje cěpov kratnyh is a stabilization condition, not a theorem about an arbitrary number of chains. In the auxiliary argument the chain of powers of a principal ideal stabilizes; absence of zero divisors then permits cancellation of the power of $a$ and yields a solution of $ax=b$. Commutativity and absence of zero divisors must remain in the statement. The short name is justified by its role in that proof, not by a claim that the chain condition alone makes every ring a field.\par
\noindent{\footnotesize Izvorne zapisy / Original records: HIST-g114-007.}\par
\smallskip\noindent\emph{Sravnjenje s tekstom:} Skratenje koristi otsutnost děliteljev nuly.\par
\noindent\emph{Textual comparison:} Cancellation uses the absence of zero divisors.\par
\noindent{\footnotesize TEXT-COMP-886; pasaže / passages: 16761, 16763, 16765, 16766, 16768, 16770, 16772; \texttt{CORPUS\_FORM\_COMPARISONS.json}.}\par
\end{minipage}\par
\par\medskip\Needspace{7\baselineskip}
\noindent\begin{minipage}{\linewidth}
\hypertarget{TERM-1830}{}
\noindent\textbf{mnogokratny teorem o cěpah / minimalno uslovje}\quad{\footnotesize\texttt{TERM-1830}}\par
\noindent\emph{Iztočnik / Source:} Vielfachenkettensatz / Minimalbedingung\par
\noindent\emph{Rusky / Russian:} теорема о множественных цепях / минимальное условие\par
\noindent\emph{Ukrajinsky / Ukrainian:} багатоланцюгова теорема / мінімальна умова\par
\smallskip\noindent\emph{Razlog:} Mnogokratny teorem o cěpah jest shranjeny prěvod, ale može mylno navoditi na mnogo cěpov ili povtorjany teorem. Točni smysl poda sama pasaža: minimalno uslovje za prave idealy, to jest bytje minimalnogo člena v razgledanoj nepustoj množině idealov. Aktualny tekst ješče sodrživaje ranějši vyraz; slovnik ne smi proglasiti jego prozračnost za ustanovjeny fakt. Sravnjenje s cěpovym uslovjem kratnyh pokazuje, začto minimalno uslovje trěba postaviti v prědny plan. Obe nazvy opisujut isto uslovje konečnosti, ne dva teorema.\par
\noindent\emph{Reason:} Mnogokratny teorem o cěpah is the preserved translation, but it can misleadingly suggest many chains or a repeatedly applied theorem. The passage itself supplies the exact meaning: the minimum condition on right ideals, stated as existence of a minimal member in a nonempty collection of ideals. The current text still contains the older phrase; the dictionary must not present its transparency as established. Comparison with the multiple-chain condition shows why the minimum-condition gloss should lead. The two names describe the same finiteness condition, not two theorems.\par
\noindent{\footnotesize Izvorne zapisy / Original records: HIST-g178-004.}\par
\smallskip\noindent\emph{Sravnjenje s tekstom:} Točny smysl jest minimalno uslovje, ne mnogo cěpov.\par
\noindent\emph{Textual comparison:} The exact meaning is the minimum condition, not a multiplicity of chains.\par
\noindent{\footnotesize TEXT-COMP-887; pasaže / passages: 19752; \texttt{CORPUS\_FORM\_COMPARISONS.json}.}\par
\end{minipage}\par
\par\medskip\Needspace{7\baselineskip}
\noindent\begin{minipage}{\linewidth}
\hypertarget{TERM-1831}{}
\noindent\textbf{četverna grupa / grupa osmogo reda}\quad{\footnotesize\texttt{TERM-1831}}\par
\noindent\emph{Iztočnik / Source:} Vierergruppe / Achtergruppe\par
\noindent\emph{Rusky / Russian:} четверная группа / группа восьмого порядка\par
\noindent\emph{Ukrajinsky / Ukrainian:} четверна група / група восьмого порядку\par
\smallskip\noindent\emph{Razlog:} Četverna grupa i grupa osmogo reda shranjajut kratke nazvy priměrov, pri čem vtora izrično nazyvaje red grupy. Kontekst uravnenij četvrtogo stepena i podstav na četyreh proměnnyh jest sučstveny: tvrdženje o grupah osmogo reda ne jest obča klasifikacija vseh abstraktnyh grup s osmju elementami. Poslědujuče dejstvje na bazovyh funkcijah objasnjaje redukciju. Poredkovy prěvod ne daje dovoljenja zaměniti konkretnu četvernu grupu libovoljnoju grupoju reda četyri.\par
\noindent\emph{Reason:} Četverna grupa and grupa osmogo reda retain compact names for the examples, with the latter making group order explicit. The setting of quartic equations and substitutions of four variables is essential: the statement about groups of order eight is not a classification of all abstract eight-element groups. The subsequent action on the basis functions explains the reduction. Order-based wording does not permit replacing the particular four-group with an arbitrary group of order four.\par
\noindent{\footnotesize Izvorne zapisy / Original records: HIST-g073-020.}\par
\smallskip\noindent\emph{Sravnjenje s tekstom:} Kontekst četvrtogo stepena ograničuje tvrdženje o grupah.\par
\noindent\emph{Textual comparison:} The quartic setting restricts the statement about the groups.\par
\noindent{\footnotesize TEXT-COMP-888; pasaže / passages: 8892, 8894, 8898, 8900; \texttt{CORPUS\_FORM\_COMPARISONS.json}.}\par
\end{minipage}\par
\par\medskip\Needspace{7\baselineskip}
\noindent\begin{minipage}{\linewidth}
\hypertarget{TERM-1832}{}
\noindent\textbf{polne sistemy invariantov}\quad{\footnotesize\texttt{TERM-1832}}\par
\noindent\emph{Iztočnik / Source:} volle Invariantensysteme\par
\noindent\emph{Rusky / Russian:} полные системы инвариантов\par
\noindent\emph{Ukrajinsky / Ukrainian:} повні системи інваріантів\par
\smallskip\noindent\emph{Razlog:} Polne sistemy invariantov vyražajut polnotu v smyslu porodženja, ne množinu, sodrživuči vsaky invariant kako odrěđeny element spiska. V izlaganju problem jest, ima li oblast konečnu integralnu bazu, iz ktorej vsaky invariant može byti vyraženy cělo i racionalno. Zato prosty pridavnik polny trěbuje to matematično objasnjenje. Hilbertovo citovano zaglavje soderži německu frazu; sam bibliografičny citat ne ustanavjaje jezyčnu normu, ale svez s definicijoju činji zmysl prověrimym.\par
\noindent\emph{Reason:} Polne sistemy invariantov expresses completeness in the sense of generation, not a list containing every invariant as a separate entry. In the exposition the problem is whether the domain has a finite integral basis from which every invariant can be expressed integrally and rationally. The plain adjective full therefore needs this mathematical explanation. Hilbert's cited title contains the German phrase; a bibliographic citation alone does not establish a linguistic norm, but its connection with the definition makes the meaning checkable.\par
\noindent{\footnotesize Izvorne zapisy / Original records: HIST-g094-022.}\par
\smallskip\noindent\emph{Sravnjenje s tekstom:} Polny sistem ne jest spis vsih pojedinyh invariantov.\par
\noindent\emph{Textual comparison:} A complete system is not a list of all individual invariants.\par
\noindent{\footnotesize TEXT-COMP-889; pasaže / passages: 15101, 15148; \texttt{CORPUS\_FORM\_COMPARISONS.json}.}\par
\end{minipage}\par
\par\medskip\Needspace{7\baselineskip}
\noindent\begin{minipage}{\linewidth}
\hypertarget{TERM-1833}{}
\noindent\textbf{polne transvekcije nad kovariantami sistema III}\quad{\footnotesize\texttt{TERM-1833}}\par
\noindent\emph{Iztočnik / Source:} volle Überschiebungen über Kovarianten von System III\par
\noindent\emph{Rusky / Russian:} полные трансвекции над ковариантами системы III\par
\noindent\emph{Ukrajinsky / Ukrainian:} повні трансвекції над коваріантами системи III\par
\smallskip\noindent\emph{Razlog:} Polne transvekcije nad kovariantami sistema III upotrěbjajut internacionalnu nazvu operacije vměsto ranějšego prěnos. Izbor opiraje se na tom, že konkretno ide o invariantno-teoretičnu transvekciju, ne proste prěměščenje objekta. Pasaža dodaje, že te transvekcije stavajut reducibilne; zato polne ne znači ireducibilne. Oznaka sistema III mora ostati ta sama v oběh pismenyh formah. Arhivna jezyčna ograda o nazvě transvekcija ostaje, a lokalny izbor ne opravdava avtomatičnu zaměnu vseho německogo Überschiebung bez konteksta.\par
\noindent\emph{Reason:} Polne transvekcije nad kovariantami sistema III uses the international operation name instead of the earlier prěnos. The reason is that the passage concerns an invariant-theoretic transvection, not simply moving an object. It adds that these transvections become reducible; full therefore does not mean irreducible. The identifier System III must remain the same in both scripts. The recorded linguistic reservation about transvekcija remains, and this contextual choice does not justify automatically replacing every occurrence of German Überschiebung without examining its role.\par
\noindent{\footnotesize Izvorne zapisy / Original records: HIST-g020-011.}\par
\smallskip\noindent\emph{Sravnjenje s tekstom:} Polny i ireducibilny ne sut sinonimy.\par
\noindent\emph{Textual comparison:} Full and irreducible are not synonyms.\par
\noindent{\footnotesize TEXT-COMP-890; pasaže / passages: 2680; \texttt{CORPUS\_FORM\_COMPARISONS.json}.}\par
\end{minipage}\par
\par\medskip\Needspace{7\baselineskip}
\noindent\begin{minipage}{\linewidth}
\hypertarget{TERM-1834}{}
\noindent\textbf{polny sistem invariantov}\quad{\footnotesize\texttt{TERM-1834}}\par
\noindent\emph{Iztočnik / Source:} volles Invariantensystem\par
\noindent\emph{Rusky / Russian:} полная система инвариантов\par
\noindent\emph{Ukrajinsky / Ukrainian:} повна система інваріантів\par
\smallskip\noindent\emph{Razlog:} Polny sistem invariantov pokračuje terminologiju sistemov form, ale tu imaje izričnu definiciju za konečnu grupu: vsaky invariant vyražaje se cělo i racionalno prěz konečno mnogo ukazanyh invariantov. Koeficienty Galoisovej rezolventy dajut takovy sistem, a druha formulacija dava granicu stepenja. Polny ne znači algebraično nezavisny ili minimalny. Zapis čuva kontinuitet nazvy, pri tom ne proglašaje vsaky ranějši Formensystem istovetnym tomu konkretnomu invariantnomu sistemu.\par
\noindent\emph{Reason:} Polny sistem invariantov continues the vocabulary of systems of forms, but here it has an explicit definition for a finite group: every invariant is expressed integrally and rationally through finitely many specified invariants. The coefficients of the Galois resolvent supply such a system, and another formulation gives a degree bound. Full does not mean algebraically independent or minimal. The recorded choice maintains continuity of terminology without identifying every earlier Formensystem with this particular invariant system.\par
\noindent{\footnotesize Izvorne zapisy / Original records: HIST-g056-007.}\par
\smallskip\noindent\emph{Sravnjenje s tekstom:} Prědstavimost ne trěbuje minimalnost ili nezavisnost.\par
\noindent\emph{Textual comparison:} Representability does not require minimality or independence.\par
\noindent{\footnotesize TEXT-COMP-891; pasaže / passages: 6658, 6701, 6703; \texttt{CORPUS\_FORM\_COMPARISONS.json}.}\par
\end{minipage}\par
\par\medskip\Needspace{7\baselineskip}
\noindent\begin{minipage}{\linewidth}
\hypertarget{TERM-1835}{}
\noindent\textbf{sovršeno / nesovršeno polje}\quad{\footnotesize\texttt{TERM-1835}}\par
\noindent\emph{Iztočnik / Source:} vollkommener / unvollkommener Körper\par
\noindent\emph{Rusky / Russian:} совершенное / несовершенное поле\par
\noindent\emph{Ukrajinsky / Ukrainian:} досконале / недосконале поле\par
\smallskip\noindent\emph{Razlog:} Sovršeno i nesovršeno polje sut historijne nazvy tehničnogo različenja, ne ocěny kačestva polja. Aktualny tekst upotrěbjaje doskonalo i nedoskonalo. Definicija govori o različnih linearnyh faktorah prostyh funkcij v vhodnom razširjenju; v karakteristice p daje i uslovje obstajanja p-togo korenja dlja vsakogo elementa. Izbor slavjanskogo pridavnika ne dopušča zaměniti doskonalost algebraičnoju zatvorjenostju. Slovnik ohranjaje ranějši korenj i aktualnu zaměnu, bez izmišljanja novogo pojma.\par
\noindent\emph{Reason:} Sovršeno and nesovršeno polje are historical names for a technical distinction, not assessments of a field's quality. The current text uses doskonalo and nedoskonalo. Its definition concerns distinct linear factors of irreducible functions in a suitable extension; in characteristic p it also gives the condition that every element has a pth root. Choosing a Slavic adjective does not permit equating perfection with algebraic closure. The dictionary preserves the earlier stem and its current replacement without inventing a new concept.\par
\noindent{\footnotesize Izvorne zapisy / Original records: HIST-g097-010, HIST-g118-008.}\par
\smallskip\noindent\emph{Sravnjenje s tekstom:} Aktualny korenj doskonal- zaměnjuje historijny sovršen-.\par
\noindent\emph{Textual comparison:} The current stem doskonal- replaces historical sovršen-.\par
\noindent{\footnotesize TEXT-COMP-892; pasaže / passages: 15363, 17275; \texttt{CORPUS\_FORM\_COMPARISONS.json}.}\par
\end{minipage}\par
\par\medskip\Needspace{7\baselineskip}
\noindent\begin{minipage}{\linewidth}
\hypertarget{TERM-1836}{}
\noindent\textbf{sovršeno osnovno polje / sovršena osnovna oblast}\quad{\footnotesize\texttt{TERM-1836}}\par
\noindent\emph{Iztočnik / Source:} vollkommener Grundkörper / vollkommener Grundbereich\par
\noindent\emph{Rusky / Russian:} совершенное основное поле / совершенная основная область\par
\noindent\emph{Ukrajinsky / Ukrainian:} досконале основне поле / досконала основна область\par
\smallskip\noindent\emph{Razlog:} Sovršeno osnovno polje i sovršena osnovna oblast čuvajut proměnu německogo Körper na Bereich v pripisci o ljudskyh prispevkah. Aktualne formy sut doskonalo osnovno polje i doskonala osnovna oblast. Pripiska razlikuje Brauerovo ograničenje na doskonalo osnovno polje od Noetherinogo odstranjenja toj ograničenosti. Zato ni zaměna oběh imen jednym, ni izpustjenje pridavnika ne jest bezsmyslovno stilistično upravjenje: ono moglo by izbrisati granice i pripisanje rezultatov. Slovnik ne vyvodi iz samej besědy oblast novu matematičnu strukturu.\par
\noindent\emph{Reason:} Sovršeno osnovno polje and sovršena osnovna oblast preserve the shift from German Körper to Bereich in the note describing the human contributions. The current forms are doskonalo osnovno polje and doskonala osnovna oblast. That note distinguishes Brauer's restriction to a perfect base field from Noether's removal of the restriction. Replacing both nouns indiscriminately, or dropping the adjective, is therefore not a harmless stylistic adjustment: it could erase the scope and attribution of the results. The dictionary does not infer a new mathematical structure from the word domain alone.\par
\noindent{\footnotesize Izvorne zapisy / Original records: HIST-g168-008.}\par
\smallskip\noindent\emph{Sravnjenje s tekstom:} Pripiska razlikuje Brauerovo prědpoloženje od Noetherinogo obobčenja.\par
\noindent\emph{Textual comparison:} The note distinguishes Brauer's assumption from Noether's generalization.\par
\noindent{\footnotesize TEXT-COMP-893; pasaže / passages: 19244; \texttt{CORPUS\_FORM\_COMPARISONS.json}.}\par
\end{minipage}\par
\par\medskip\Needspace{7\baselineskip}
\noindent\begin{minipage}{\linewidth}
\hypertarget{TERM-1837}{}
\noindent\textbf{sovršeno polje}\quad{\footnotesize\texttt{TERM-1837}}\par
\noindent\emph{Iztočnik / Source:} vollkommener Körper\par
\noindent\emph{Rusky / Russian:} совершенное поле\par
\noindent\emph{Ukrajinsky / Ukrainian:} досконале поле\par
\smallskip\noindent\emph{Razlog:} Sovršeno polje v toj pasážě svjazuje nazvu polja s razloženjem prstena konečnogo ranga. Aktualno doskonalo polje isključuje polnu razložimost drugogo roda; algebraično zatvorjeno polje jest posebno navedeny slučaj, ne sinonim. Arhivny izbor prodolžaje ranějšu terminologiju istoj sbirky i srovnjaje rusku i ukrajinsku formu. Ta kontinuitet jest jezyčna motivacija; konkretno ograničenje prvym rodom daje tekst, ne podobnost slov.\par
\noindent\emph{Reason:} Sovršeno polje in this passage links the name of the field property to decomposition of a finite-rank ring. Current doskonalo polje excludes complete decomposability of the second kind; an algebraically closed field is a specifically named case, not a synonym. The historical choice continues earlier terminology in the same collection and compares the Russian and Ukrainian forms. That continuity is a linguistic motivation; the restriction to the first kind is supplied by the text, not by similarity of words.\par
\noindent{\footnotesize Izvorne zapisy / Original records: HIST-g125-005.}\par
\smallskip\noindent\emph{Sravnjenje s tekstom:} Algebraično zatvorjeno polje jest specialny slučaj, ne sinonim.\par
\noindent\emph{Textual comparison:} An algebraically closed field is a special case, not a synonym.\par
\noindent{\footnotesize TEXT-COMP-894; pasaže / passages: 17630, 17632; \texttt{CORPUS\_FORM\_COMPARISONS.json}.}\par
\end{minipage}\par
\par\medskip\Needspace{7\baselineskip}
\noindent\begin{minipage}{\linewidth}
\hypertarget{TERM-1838}{}
\noindent\textbf{sovršeno/nesovršeno polje}\quad{\footnotesize\texttt{TERM-1838}}\par
\noindent\emph{Iztočnik / Source:} vollkommener/unvollkommener Körper\par
\noindent\emph{Rusky / Russian:} совершенное/несовершенное поле\par
\noindent\emph{Ukrajinsky / Ukrainian:} досконале/недосконале поле\par
\smallskip\noindent\emph{Razlog:} Sovršeno/nesovršeno polje jest kompaktna para s negativnym prefiksom; kosaja črta ne znači jedno proměnlivo svojstvo. V definiciji aktuálnogo teksta doskonalo polje charakteristiki nul vsegda izpolnjaje uslovje, medžu tym kako pri prostej charakteristice trěba osobno uslovje o korenjah. Daljša pasaža upozornjaje, že prisojedinjenje neoprěděljenyh može dati nedoskonalo polje, hoti osnovno polje bylo doskonalo. Zato terminy morajut byti svjazane s točno imenovanym poljem, ne avtomatično prěnesene na vsako razširjenje.\par
\noindent\emph{Reason:} Sovršeno/nesovršeno polje is a compact positive–negative pair; the slash does not denote one fluctuating property. In the current definition every characteristic-zero field satisfies the perfection condition, whereas prime characteristic requires the separate root condition. A later passage warns that adjoining indeterminates can give an imperfect field even when the base field was perfect. The terms must therefore attach to the precisely named field, not automatically extend to every extension.\par
\noindent{\footnotesize Izvorne zapisy / Original records: HIST-g095-017.}\par
\smallskip\noindent\emph{Sravnjenje s tekstom:} Doskonalost ne prěnosi se avtomatično na tuto razširjenje.\par
\noindent\emph{Textual comparison:} Perfection does not automatically pass to this extension.\par
\noindent{\footnotesize TEXT-COMP-895; pasaže / passages: 15363, 15467, 15469; \texttt{CORPUS\_FORM\_COMPARISONS.json}.}\par
\end{minipage}\par
\par\medskip\Needspace{7\baselineskip}
\noindent\begin{minipage}{\linewidth}
\hypertarget{TERM-1839}{}
\noindent\textbf{sovršeno / nesovršeno polje koeficientov}\quad{\footnotesize\texttt{TERM-1839}}\par
\noindent\emph{Iztočnik / Source:} vollkommenes / unvollkommenes Koeffizientenfeld\par
\noindent\emph{Rusky / Russian:} совершенное / несовершенное поле коэффициентов\par
\noindent\emph{Ukrajinsky / Ukrainian:} досконале / недосконале поле коефіцієнтів\par
\smallskip\noindent\emph{Razlog:} Sovršeno ili nesovršeno polje koeficientov izrično pokazuje, dlja ktorogo polja potrěbna jest hipoteza. Aktualny tekst piše doskonalo i nedoskonalo. Nad doskonalym poljem Teorema XIII dava potrěbno i dostatočno uslovje prostosti; nad nedoskonalym poljem Teorema XIV vměsto togo dava formulu s $p^g$-tym stepenjem i granicoju eksponenta. Izbor slavjanskogo korenja vměsto internacionalizmu ne jest dovoljenje opustiti tu razliku. Glosa koeficientov čuva město hipotezy v dokazě.\par
\noindent\emph{Reason:} Sovršeno or nesovršeno polje koeficientov explicitly identifies the field to which the hypothesis applies. The current text writes doskonalo and nedoskonalo. Over a perfect field Theorem XIII gives a necessary-and-sufficient primality criterion; over an imperfect field Theorem XIV instead gives a $p^g$-power formula with an exponent bound. Preference for a Slavic stem over an internationalism is not permission to omit that distinction. The coefficient-field qualifier keeps the hypothesis in its proper place in the argument.\par
\noindent{\footnotesize Izvorne zapisy / Original records: HIST-g100-006.}\par
\smallskip\noindent\emph{Sravnjenje s tekstom:} Oblast koeficientov odlučuje, ktory zaključok jest dostupy.\par
\noindent\emph{Textual comparison:} The coefficient field determines which conclusion is available.\par
\noindent{\footnotesize TEXT-COMP-896; pasaže / passages: 15661, 15663, 15667, 15669; \texttt{CORPUS\_FORM\_COMPARISONS.json}.}\par
\end{minipage}\par
\par\medskip\Needspace{7\baselineskip}
\noindent\begin{minipage}{\linewidth}
\hypertarget{TERM-1840}{}
\noindent\textbf{popolno reducibilny}\quad{\footnotesize\texttt{TERM-1840}}\par
\noindent\emph{Iztočnik / Source:} vollstaendig reduzibel\par
\noindent\emph{Rusky / Russian:} вполне приводимый\par
\noindent\emph{Ukrajinsky / Ukrainian:} повністю звідний\par
\smallskip\noindent\emph{Razlog:} Popolno reducibilny jest ranějša forma aktualnogo polno reducibilny. Internacionalny pridavnik prodolžaje terminologiju Loewyja, ale definicija v toj rabotě jest dana dlja modulov: prědstavimost kako najmenše obče kratno prostyh modulov. To ne znači prosto možnost razložiti polinom na faktory. Uvod posebno razděljaje zaslugi: obobčenje Loewyjovogo pojma i jego svjazok s izomorfijeju pripisany sut Schmeidleru. Kontinuitet termina ne smi izbrisati ni toj adaptacije, ni ljudskogo avtorstva.\par
\noindent\emph{Reason:} Popolno reducibilny is an earlier form of current polno reducibilny. The international adjective continues Loewy's terminology, but this work defines it for modules: representability as a least common multiple of simple modules. It does not mean merely that a polynomial can be factored. The introduction expressly separates the contributions, crediting Schmeidler with generalizing Loewy's notion and linking it to isomorphism. Continuity of terminology must erase neither that adaptation nor its human authorship.\par
\noindent{\footnotesize Izvorne zapisy / Original records: HIST-g079-007.}\par
\smallskip\noindent\emph{Sravnjenje s tekstom:} Modulna definicija i Schmeidlerov prispevok ostajut vidlive.\par
\noindent\emph{Textual comparison:} The module definition and Schmeidler's contribution remain visible.\par
\noindent{\footnotesize TEXT-COMP-897; pasaže / passages: 10953, 10961, 10965; \texttt{CORPUS\_FORM\_COMPARISONS.json}.}\par
\end{minipage}\par
\par\medskip\Needspace{7\baselineskip}
\noindent\begin{minipage}{\linewidth}
\hypertarget{TERM-1841}{}
\noindent\textbf{popolno reducibilna grupa}\quad{\footnotesize\texttt{TERM-1841}}\par
\noindent\emph{Iztočnik / Source:} vollstaendig reduzible Gruppe\par
\noindent\emph{Rusky / Russian:} вполне приводимая группа\par
\noindent\emph{Ukrajinsky / Ukrainian:} повністю звідна група\par
\smallskip\noindent\emph{Razlog:} Popolno reducibilna grupa prodolžaje internacionalnu družinu reducibilnost; aktualna forma jest polno reducibilna grupa. V toj rabotě grupa ostatkov tako nazyvaje se, ako dopušča prinajmenje jedno prědstavjenje kako suma konečno mnogyh prostyh grup. Prosta grupa pritom oprěděljena jest čerez prosty modul. Zato ne trěba bez objasnjenja prěnositi dana slova na proizvoljne grupy v sovrěmennom smislu. Kontinuitet nazvy čuva svjazok s ranějšimi paragrafami i s citovanoju terminologijeju Loewyja.\par
\noindent\emph{Reason:} Popolno reducibilna grupa continues the international reducibility family; the current form is polno reducibilna grupa. In this work a residue group receives that name if it admits at least one expression as a sum of finitely many simple groups. Simple group is itself defined through a simple module. These words should therefore not be transferred without explanation to arbitrary groups in the modern sense. Terminological continuity preserves the link with the earlier sections and the cited terminology of Loewy.\par
\noindent{\footnotesize Izvorne zapisy / Original records: HIST-g079-046.}\par
\smallskip\noindent\emph{Sravnjenje s tekstom:} Dostatočno jest prinajmenje jedno konečno sumno prědstavjenje.\par
\noindent\emph{Textual comparison:} At least one finite sum expression is sufficient.\par
\noindent{\footnotesize TEXT-COMP-898; pasaže / passages: 11572, 11574, 11576; \texttt{CORPUS\_FORM\_COMPARISONS.json}.}\par
\end{minipage}\par
\par\medskip\Needspace{7\baselineskip}
\noindent\begin{minipage}{\linewidth}
\hypertarget{TERM-1842}{}
\noindent\textbf{polno razložimy}\quad{\footnotesize\texttt{TERM-1842}}\par
\noindent\emph{Iztočnik / Source:} vollständig reduzibel\par
\noindent\emph{Rusky / Russian:} вполне приводимый\par
\noindent\emph{Ukrajinsky / Ukrainian:} цілком звідний\par
\smallskip\noindent\emph{Razlog:} Polno razložimy izbran jest zaradi eksplicitnoj definicije: prsten može byti prědstavjen kako direktna suma konečno mnogyh polj. Ravnocěnna formulacija govori o prostyh komponentah razloženja nulovogo ideala. Slavjanska osnova razlož- ukazuje na tu strukturu; ona ne jest prosta zaměna slova redukovany. Sovrěmenny pojem prstena bez nenulovyh nilpotentnyh elementov ne možno upotrěbiti kako neograničenu sinonimičnu definiciju togo historijnogo termina. Prědpoloženja o konkretnom prstenu ostajut potrěbne.\par
\noindent\emph{Reason:} Polno razložimy was chosen because of the explicit definition: the ring can be expressed as a direct sum of finitely many fields. The equivalent formulation concerns simple components in the decomposition of the zero ideal. The Slavic stem razlož- points to that structure; it is not a mechanical replacement for reduced. The modern notion of a ring without nonzero nilpotent elements cannot be used as an unrestricted synonymous definition of this historical term. The assumptions on the particular ring remain necessary.\par
\noindent{\footnotesize Izvorne zapisy / Original records: HIST-g125-001.}\par
\smallskip\noindent\emph{Sravnjenje s tekstom:} Definicija govori o konečnoj direktnoj sumi polj.\par
\noindent\emph{Textual comparison:} The definition concerns a finite direct sum of fields.\par
\noindent{\footnotesize TEXT-COMP-899; pasaže / passages: 17630, 17632; \texttt{CORPUS\_FORM\_COMPARISONS.json}.}\par
\end{minipage}\par
\par\medskip\Needspace{7\baselineskip}
\noindent\begin{minipage}{\linewidth}
\hypertarget{TERM-1843}{}
\noindent\textbf{popolno reducibilny}\quad{\footnotesize\texttt{TERM-1843}}\par
\noindent\emph{Iztočnik / Source:} vollständig reduzibel\par
\noindent\emph{Rusky / Russian:} вполне приводимый\par
\noindent\emph{Ukrajinsky / Ukrainian:} цілком звідний; цілком редуцибельний\par
\smallskip\noindent\emph{Razlog:} Popolno reducibilny jest obča historijna forma, nyně polno reducibilny, priměnjena k raznym objektam: sistemam, grupam, modulam i prědstavjenjam. Zapis dava prednost internacionalnoj osnově prěd novoju kalkoju; rusky i ukrajinsky terminy pokazyvajut različne tvorbe togo samogo tehničnogo pojma. Ale kriterij mora byti lokalny: v jednoj pasaži idě o vse klasy prědstavjenij nad danym poljem, v drugoj o direktny produkt prostyh grup, v tretej o ireducibilne bloky. Jedna besěda ne zaměnjaje ty definicije.\par
\noindent\emph{Reason:} Popolno reducibilny is the general historical form, now polno reducibilny, applied to different objects: systems, groups, modules and representations. The record prefers the international stem to a new calque; the Russian and Ukrainian terms illustrate different formations for the same technical family. But the criterion must be local: one passage concerns all representation classes over a given field, another a direct product of simple groups, and another irreducible blocks. One word does not replace those definitions.\par
\noindent{\footnotesize Izvorne zapisy / Original records: HIST-g118-007, HIST-g170-004, HIST-g183-001, HIST-g194-006.}\par
\smallskip\noindent\emph{Sravnjenje s tekstom:} Lokalna definicija ostaje potrěbna dlja vsakogo objekta.\par
\noindent\emph{Textual comparison:} Each object still requires its local definition.\par
\noindent{\footnotesize TEXT-COMP-900; pasaže / passages: 17273, 19280, 19314, 20149, 20151, 21549; \texttt{CORPUS\_FORM\_COMPARISONS.json}.}\par
\end{minipage}\par
\par\medskip\Needspace{7\baselineskip}
\noindent\begin{minipage}{\linewidth}
\hypertarget{TERM-1844}{}
\noindent\textbf{popolno reducibilny / popolna reducibilnost}\quad{\footnotesize\texttt{TERM-1844}}\par
\noindent\emph{Iztočnik / Source:} vollständig reduzibel\par
\noindent\emph{Rusky / Russian:} вполне приводимый / полная приводимость\par
\noindent\emph{Ukrajinsky / Ukrainian:} цілком звідний / цілковита звідність\par
\smallskip\noindent\emph{Razlog:} Popolno reducibilny i popolna reducibilnost sut pridavnikova i imennikova forma jednoj vlastnosti, nyně s polno i polna. Definicija v §13 govori o prstenu kako direktnoj sumi konečno mnogyh prostyh pravyh idealov. Pravo označuje stranu dejstvja, ne jezyčnu dopolniteljnost, ktoru možno izpustiti. Sledujuča posledica o nilpotentnyh idealah ne jest trěbovanje, že vsaky nilpotentny element jest nula. Arhivny izbor internacionalizma zato izběgaje mylno čitanje kako obyčno redukovano prsten.\par
\noindent\emph{Reason:} Popolno reducibilny and popolna reducibilnost are the adjectival and nominal forms of one property, now with polno and polna. The definition in §13 describes the ring as a direct sum of finitely many simple right ideals. Right specifies the side of the action; it is not a dispensable verbal addition. The following consequence about nilpotent ideals is not a requirement that every nilpotent element be zero. The recorded internationalism therefore avoids misleading interpretation as an ordinary reduced ring.\par
\noindent{\footnotesize Izvorne zapisy / Original records: HIST-g191-001.}\par
\smallskip\noindent\emph{Sravnjenje s tekstom:} Nilpotentny ideal i nilpotentny element ne sut isto.\par
\noindent\emph{Textual comparison:} A nilpotent ideal and a nilpotent element are not the same.\par
\noindent{\footnotesize TEXT-COMP-901; pasaže / passages: 20914, 20916, 20918; \texttt{CORPUS\_FORM\_COMPARISONS.json}.}\par
\end{minipage}\par
\par\medskip\Needspace{7\baselineskip}
\noindent\begin{minipage}{\linewidth}
\hypertarget{TERM-1845}{}
\noindent\textbf{polno razložimy prvogo roda}\quad{\footnotesize\texttt{TERM-1845}}\par
\noindent\emph{Iztočnik / Source:} vollständig reduzibel erster Art\par
\noindent\emph{Rusky / Russian:} вполне приводимый первого рода\par
\noindent\emph{Ukrajinsky / Ukrainian:} цілком звідний першого роду\par
\smallskip\noindent\emph{Razlog:} Polno razložimy prvogo roda prodolžaje slavjansku družinu razložimosti, pri tom čuva dodatno uslovje prvogo roda. Teorem v §3 svjazuje to uslovje s polnoju razložimostju poslě razširjenja do algebraično zatvorjenogo polja. Teorem v §6 dava ravnocěnny kriterij: diskriminantny ideal jest jediničny. Ni prosty razklad, ni nenulovost kakoj-nebud neukazanej veličiny ne sut zaměna toho kriterija. Prvogo roda jest matematična hipoteza, ne ocěna lepšej razložimosti.\par
\noindent\emph{Reason:} Polno razložimy prvogo roda continues the Slavic decomposition family while retaining the additional first-kind condition. The theorem in §3 connects that condition with complete decomposability after extension to an algebraically closed field. The theorem in §6 gives an equivalent criterion: the discriminant ideal is the unit ideal. Neither a mere decomposition nor nonvanishing of an unspecified quantity replaces that criterion. First kind is a mathematical hypothesis, not a judgment that the decomposition is better.\par
\noindent{\footnotesize Izvorne zapisy / Original records: HIST-g126-002, HIST-g144-003.}\par
\smallskip\noindent\emph{Sravnjenje s tekstom:} Prvogo roda ne jest neobvezny pridavok.\par
\noindent\emph{Textual comparison:} First kind is not an optional qualifier.\par
\noindent{\footnotesize TEXT-COMP-902; pasaže / passages: 17630, 17661, 18412; \texttt{CORPUS\_FORM\_COMPARISONS.json}.}\par
\end{minipage}\par
\par\medskip\Needspace{7\baselineskip}
\noindent\begin{minipage}{\linewidth}
\hypertarget{TERM-1846}{}
\noindent\textbf{polno razložimy drugogo roda}\quad{\footnotesize\texttt{TERM-1846}}\par
\noindent\emph{Iztočnik / Source:} vollständig reduzibel zweiter Art\par
\noindent\emph{Rusky / Russian:} вполне приводимый второго рода\par
\noindent\emph{Ukrajinsky / Ukrainian:} цілком звідний другого роду\par
\smallskip\noindent\emph{Razlog:} Polno razložimy drugogo roda nazyvaje drugu granu točnoj definicije, ne vtory krok procesa. Prsten ostaje polno razložimym nad početkovym poljem, ale poslě razširjenja do algebraičnogo zatvorjenja ta vlastnost može propasti; tekst dokazuje, že v razgledanom slučaju ona propadaje. Nenulovy element s nulovoju stepenju konstruuje se v razširjenom prstenu, ne prosto prědpokladaje se v početkovom polju. Takova zavisimost od oblasti koeficientov jest pričina sohraniti drugogo roda v cělom terminu.\par
\noindent\emph{Reason:} Polno razložimy drugogo roda names the other branch of the precise definition, not a second stage of a procedure. The ring remains completely decomposable over the original field, but the property can fail after extension to an algebraic closure; the text proves failure in the case under discussion. A nonzero element whose power is zero is constructed in the extended ring, not simply assumed to exist in the original field. This dependence on the coefficient field is why second kind must remain part of the term.\par
\noindent{\footnotesize Izvorne zapisy / Original records: HIST-g128-001.}\par
\smallskip\noindent\emph{Sravnjenje s tekstom:} Nenulovy nilpotentny element konstruuje se poslě razširjenja.\par
\noindent\emph{Textual comparison:} The nonzero nilpotent element is constructed after extension.\par
\noindent{\footnotesize TEXT-COMP-903; pasaže / passages: 17630, 17733, 17735, 17737, 17739, 17741, 17743, 17745, 17747, 17748, 17750; \texttt{CORPUS\_FORM\_COMPARISONS.json}.}\par
\end{minipage}\par
\par\medskip\Needspace{7\baselineskip}
\noindent\begin{minipage}{\linewidth}
\hypertarget{TERM-1847}{}
\noindent\textbf{polno redukovane idealy}\quad{\footnotesize\texttt{TERM-1847}}\par
\noindent\emph{Iztočnik / Source:} vollständig reduzible Ideale\par
\noindent\emph{Rusky / Russian:} вполне редуцируемые идеалы\par
\noindent\emph{Ukrajinsky / Ukrainian:} повністю редуковані ідеали\par
\smallskip\noindent\emph{Razlog:} Polno redukovane idealy jest shranjena historijna nazva osobnogo klasa, svjazana s rabotoju Noether i Schmeidlera. Pasaža ograničuje prěnos teoremov jednoznačnosti v nekomutativny slučaj; ne proglašaje te teoremy za vse idealy. Poslědujuče objasnjenje ukazuje na poparno vzajemno proste komponenty bez vlastnyh děliteljev. Zato redukovane ne možno čitati prosto kako ideal, ktory byl skračen računskoju proceduroju. Ranějšu ogradu o sovrěmennom prěimenovanji trěba čuvati vměstě s konkretnoju definicijeju.\par
\noindent\emph{Reason:} Polno redukovane idealy is a preserved historical label for a special class associated with the work of Noether and Schmeidler. The passage restricts transfer of uniqueness theorems to the noncommutative setting; it does not assert them for all ideals. The subsequent explanation specifies pairwise relatively prime components without proper divisors. Redukovane therefore cannot mean merely an ideal shortened by a calculation. The historical reservation about modernizing the label must remain alongside the concrete definition.\par
\noindent{\footnotesize Izvorne zapisy / Original records: HIST-g081-077.}\par
\smallskip\noindent\emph{Sravnjenje s tekstom:} Teoremy jednoznačnosti ne prěnoset se bez ograničenja.\par
\noindent\emph{Textual comparison:} The uniqueness theorems do not transfer without restriction.\par
\noindent{\footnotesize TEXT-COMP-904; pasaže / passages: 13407, 13409, 13666; \texttt{CORPUS\_FORM\_COMPARISONS.json}.}\par
\end{minipage}\par
\par\medskip\Needspace{7\baselineskip}
\noindent\begin{minipage}{\linewidth}
\hypertarget{TERM-1848}{}
\noindent\textbf{popolno reducibilno kolco}\quad{\footnotesize\texttt{TERM-1848}}\par
\noindent\emph{Iztočnik / Source:} vollständig reduzibler Ring\par
\noindent\emph{Rusky / Russian:} вполне приводимое кольцо\par
\noindent\emph{Ukrajinsky / Ukrainian:} цілком звідне кільце\par
\smallskip\noindent\emph{Razlog:} Popolno reducibilno kolco jest ranějša kombinacija, ktoroj aktualna forma jest polno reducibilno prsten. Zapis čuva korenj reducibil- zaradi svjazku s modulami i prědstavjenjami, vměsto nejasnogo redukovano. V §18 jest potrěbno takože dvustranna prostost i to, že jedinica prstena dějstvuje na modul kako jediničny operator. Širša pasaža dobivaje taki prsten čerez kvocient po radikalu. Zaměna kolco na prsten ne odstranjava ni te uslovja, ni razliku medžu idealom i elementom.\par
\noindent\emph{Reason:} Popolno reducibilno kolco is the earlier combination whose current form is polno reducibilno prsten. The record retains reducibil- for its connection with modules and representations, instead of ambiguous reduced. In §18 two-sided simplicity and action of the ring identity as the identity operator on the module are also required. The broader passage obtains such a ring by quotienting by the radical. Replacing kolco with prsten removes neither those assumptions nor the distinction between an ideal and an element.\par
\noindent{\footnotesize Izvorne zapisy / Original records: HIST-g177-006, HIST-g197-001.}\par
\smallskip\noindent\emph{Sravnjenje s tekstom:} Jedinica mora dějstvovati kako jediničny operator.\par
\noindent\emph{Textual comparison:} The identity must act as the identity operator.\par
\noindent{\footnotesize TEXT-COMP-905; pasaže / passages: 19710, 21669, 21673, 21675, 21679; \texttt{CORPUS\_FORM\_COMPARISONS.json}.}\par
\end{minipage}\par
\par\medskip\Needspace{7\baselineskip}
\noindent\begin{minipage}{\linewidth}
\hypertarget{TERM-1849}{}
\noindent\textbf{popolno reducibilno kolco / grupno kolco}\quad{\footnotesize\texttt{TERM-1849}}\par
\noindent\emph{Iztočnik / Source:} vollständig reduzibler Ring / Gruppenring\par
\noindent\emph{Rusky / Russian:} вполне приводимое кольцо / групповое кольцо\par
\noindent\emph{Ukrajinsky / Ukrainian:} цілком редуцибельне кільце / групове кільце\par
\smallskip\noindent\emph{Razlog:} Popolno reducibilno kolco i grupno kolco sut dvě svjazane nazvy, ne sinonimy. Aktualny tekst imaje polno reducibilno prsten i grupno prsten. V opisanom pristupu teorija konečnyh grupnyh prědstavjenij vhodi v teoriju idealov takogo prstena; nerazložime idealne klasy odgovarjajut ireducibilnym prědstavjenjam. Samo sostavno slovo grupno ne dokazuje polnu reducibilnost vsakego grupnogo prstena nad libovoljnymi koeficientami. Zapis koristi kontinuitet s ranějšeju modulnoju terminologijeju, a ne novo univerzalno tvrdženje.\par
\noindent\emph{Reason:} Popolno reducibilno kolco and grupno kolco are related names, not synonyms. The current text has polno reducibilno prsten and grupno prsten. In the approach described, finite-group representation theory is placed within the ideal theory of such a ring; indecomposable ideal classes correspond to irreducible representations. The compound group ring does not itself establish complete reducibility of every group ring over arbitrary coefficients. The record uses continuity with earlier module terminology, not a new universal assertion.\par
\noindent{\footnotesize Izvorne zapisy / Original records: HIST-g105-004.}\par
\smallskip\noindent\emph{Sravnjenje s tekstom:} Dvě nazvy ne sut univerzalne sinonimy.\par
\noindent\emph{Textual comparison:} The two names are not universal synonyms.\par
\noindent{\footnotesize TEXT-COMP-906; pasaže / passages: 15948, 15950; \texttt{CORPUS\_FORM\_COMPARISONS.json}.}\par
\end{minipage}\par
\par\medskip\Needspace{7\baselineskip}
\noindent\begin{minipage}{\linewidth}
\hypertarget{TERM-1850}{}
\noindent\textbf{polno redukovane idealy}\quad{\footnotesize\texttt{TERM-1850}}\par
\noindent\emph{Iztočnik / Source:} vollständig-reduzible Ideale\par
\noindent\emph{Rusky / Russian:} вполне редуцируемые идеалы\par
\noindent\emph{Ukrajinsky / Ukrainian:} повністю редуковані ідеали\par
\smallskip\noindent\emph{Razlog:} Polno redukovane idealy v toj drugoj zapisi shranjajut tože historijno ime, hoti německy termin zapisany jest s defisom. Konkretna definicija govori o razloženju na poparno vzajemno proste komponenty bez vlastnyh děliteljev; komponenty sut zato nerazložime. Iz toho vyvodima ravnost čisla komponentov pri dvuh razloženjah, zajedno s ranějšim izomorfizmom. Defis ne tvori novy pojem. Terminu potrěbno priložiti tę definiciju, ne zaměniti ju sovrěmennoju neprověrenoju frazoju iz linearnoj algebry.\par
\noindent\emph{Reason:} Polno redukovane idealy in this second record retains the same historical name although the German term is hyphenated. The concrete definition concerns decomposition into pairwise relatively prime components without proper divisors; the components are consequently indecomposable. Equality of the number of components in two decompositions follows, alongside the earlier isomorphism result. The hyphen does not create another concept. The term needs that definition, not replacement by an unchecked modern phrase from linear algebra.\par
\noindent{\footnotesize Izvorne zapisy / Original records: HIST-g081-092.}\par
\smallskip\noindent\emph{Sravnjenje s tekstom:} Definicija čuva i primarny tekst, i historijny svjazok.\par
\noindent\emph{Textual comparison:} The definition retains the primary text and the historical connection.\par
\noindent{\footnotesize TEXT-COMP-907; pasaže / passages: 13666, 13668; \texttt{CORPUS\_FORM\_COMPARISONS.json}.}\par
\end{minipage}\par
\par\medskip\Needspace{7\baselineskip}
\noindent\begin{minipage}{\linewidth}
\hypertarget{TERM-1851}{}
\noindent\textbf{polne sistemy form}\quad{\footnotesize\texttt{TERM-1851}}\par
\noindent\emph{Iztočnik / Source:} vollständige Formensysteme\par
\noindent\emph{Rusky / Russian:} полные системы форм\par
\noindent\emph{Ukrajinsky / Ukrainian:} повні системи форм\par
\smallskip\noindent\emph{Razlog:} Polne sistemy form jest kratak i doslovny prěvod zaključnoj nazvy priměnjenja. Avtorka tamo dava samo ukazanje, že priměnjenja k redukciji takyh sistemov slědujut analogično binarnomu i ternarnomu slučaju. Slovnik ne smi iz kratkoj primětky napraviti tvrdženje, že tekst detaljno konstruuje vse te sistemy. Predhodny dokaz o kontrakcijah i polarah objasnjaje, kakogo tipa jest priměnjenje; sama nazva polny ne zaměnjaje toj osnovy.\par
\noindent\emph{Reason:} Polne sistemy form is a concise literal rendering of the application named at the close of the paper. The author gives only a brief statement that applications to reduction of such systems follow analogously to the binary and ternary cases. The dictionary must not turn that note into a claim that the text constructs all those systems in detail. The preceding proof concerning contractions and polars explains the kind of application; the adjective complete does not replace that basis.\par
\noindent{\footnotesize Izvorne zapisy / Original records: HIST-g038-012.}\par
\smallskip\noindent\emph{Sravnjenje s tekstom:} Ukazano priměnjenje ne jest podrobno izvedena konstrukcija.\par
\noindent\emph{Textual comparison:} A named application is not a construction carried out in detail.\par
\noindent{\footnotesize TEXT-COMP-908; pasaže / passages: 4914, 4916, 4918; \texttt{CORPUS\_FORM\_COMPARISONS.json}.}\par
\end{minipage}\par
\par\medskip\Needspace{7\baselineskip}
\noindent\begin{minipage}{\linewidth}
\hypertarget{TERM-1852}{}
\noindent\textbf{polna matrična algebra}\quad{\footnotesize\texttt{TERM-1852}}\par
\noindent\emph{Iztočnik / Source:} vollständige Matrixalgebra\par
\noindent\emph{Rusky / Russian:} полная матричная алгебра\par
\noindent\emph{Ukrajinsky / Ukrainian:} повна матрична алгебра\par
\smallskip\noindent\emph{Razlog:} Polna matrična algebra jest internacionalna terminologija za algebra iz vseh matric danogo reda nad danym poljem, ne samo proizvoljnu podalgebru matric. Pasaža izrično objasnjaje oznaku $A\sim\Omega$ kako naleženje osobnoj klasi podobnosti, kotroj dělitevny predstavitelj jest samo osnovno polje. To ne znači bukvalnu ravnost $A$ i $\Omega$. Polna zato označuje obseg matric, a ne polnotu po někoj topologičnoj metrice.\par
\noindent\emph{Reason:} Polna matrična algebra uses international terminology for the algebra of all matrices of a specified size over the field, not merely an arbitrary matrix subalgebra. The passage explicitly explains $A\sim\Omega$ as membership in the special similarity class whose division-algebra representative is the base field itself. It does not mean literal equality of $A$ and $\Omega$. Full therefore concerns the scope of the matrices, not completeness in a topological metric.\par
\noindent{\footnotesize Izvorne zapisy / Original records: HIST-g208-006.}\par
\smallskip\noindent\emph{Sravnjenje s tekstom:} Podobnost ne znači bukvalnu jednakost algebr.\par
\noindent\emph{Textual comparison:} Similarity does not mean literal equality of the algebras.\par
\noindent{\footnotesize TEXT-COMP-909; pasaže / passages: 23013; \texttt{CORPUS\_FORM\_COMPARISONS.json}.}\par
\end{minipage}\par
\par\medskip\Needspace{7\baselineskip}
\noindent\begin{minipage}{\linewidth}
\hypertarget{TERM-1853}{}
\noindent\textbf{polna razložimost prvogo roda}\quad{\footnotesize\texttt{TERM-1853}}\par
\noindent\emph{Iztočnik / Source:} vollständige Reduzibilität erster Art\par
\noindent\emph{Rusky / Russian:} полная приводимость первого рода\par
\noindent\emph{Ukrajinsky / Ukrainian:} цілковита звідність першого роду\par
\smallskip\noindent\emph{Razlog:} Polna razložimost prvogo roda jest imennikova forma istej strukturnoj hipotezy. V §6 ona ravnosilna razloženju poslě razširjenja koeficientov na poljne komponenty izomorfne osnovnomu razširenomu polju. Dalje ta hipoteza dopušča prědstavjenje slěda, normy i diskriminanta prěz konjugovane elementy. Zato perenos nazvy iz pridavnika v imennik ne smi izpustiti prvogo roda ili uslovja o polju. Prozračnost osnovy razlož- ima smysl samo vměstě s konkretnoju strukturoju.\par
\noindent\emph{Reason:} Polna razložimost prvogo roda is the nominal form of the same structural hypothesis. In §6 it is equivalent to decomposition, after coefficient extension, into field components isomorphic to the extended base field. The hypothesis then permits representation of trace, norm and discriminant through conjugate elements. Changing from adjective to noun must therefore drop neither first kind nor the field condition. The transparency of razlož- is useful only together with the actual structure.\par
\noindent{\footnotesize Izvorne zapisy / Original records: HIST-g145-001, HIST-g146-003.}\par
\smallskip\noindent\emph{Sravnjenje s tekstom:} Tože uslovje upravjaje daljšim prědstavjenjem invariantov.\par
\noindent\emph{Textual comparison:} The same condition governs the subsequent expression of invariants.\par
\noindent{\footnotesize TEXT-COMP-910; pasaže / passages: 18441, 18470, 18472, 18474; \texttt{CORPUS\_FORM\_COMPARISONS.json}.}\par
\end{minipage}\par
\par\medskip\Needspace{7\baselineskip}
\noindent\begin{minipage}{\linewidth}
\hypertarget{TERM-1854}{}
\noindent\textbf{polnost identičnostij}\quad{\footnotesize\texttt{TERM-1854}}\par
\noindent\emph{Iztočnik / Source:} Vollständigkeit der Identitäten\par
\noindent\emph{Rusky / Russian:} полнота тождеств\par
\noindent\emph{Ukrajinsky / Ukrainian:} повнота тотожностей\par
\smallskip\noindent\emph{Razlog:} Polnost identičnostij označuje dostatočnost ukazanoj družiny dlja razgledanyh invariantno-teoretičnyh pytanj. Pasaža razlikuje dani dokaz od Grassmann–Müllerovogo i svjazuje identičnosti (20), (21) s obščimi nerazložimymi identičnostjami (36). To ne jest samo tvrdženje, že spis dolgy, ili polnost matematičnoj logiky. Slovnik sohranjaje ograničenje dlja tyh pytanj i modalnu ogradu avtorky o tom, čto raněje ne bylo zaměčeno.\par
\noindent\emph{Reason:} Polnost identičnostij means sufficiency of the specified family for the invariant-theoretic questions being considered. The passage distinguishes its proof from Grassmann–Müller's and connects identities (20) and (21) with the general indecomposable identities (36). It is not merely a statement that the list is long, or a completeness claim in mathematical logic. The dictionary preserves the restriction to those questions and the author's qualification concerning what appears not to have been noticed before.\par
\noindent{\footnotesize Izvorne zapisy / Original records: HIST-g029-017.}\par
\smallskip\noindent\emph{Sravnjenje s tekstom:} Obseg polnosti ograničen jest razgledanymi pytanjami.\par
\noindent\emph{Textual comparison:} The completeness claim is limited to the questions under consideration.\par
\noindent{\footnotesize TEXT-COMP-911; pasaže / passages: 3713; \texttt{CORPUS\_FORM\_COMPARISONS.json}.}\par
\end{minipage}\par
\par\medskip\Needspace{7\baselineskip}
\noindent\begin{minipage}{\linewidth}
\hypertarget{TERM-1860}{}
\noindent\textbf{jednakogo vida}\quad{\footnotesize\texttt{TERM-1860}}\par
\noindent\emph{Iztočnik / Source:} von gleichem Typ / von gleicher Art in Section10 divisors\par
\noindent\emph{Rusky / Russian:} того же вида\par
\noindent\emph{Ukrajinsky / Ukrainian:} того самого виду\par
\smallskip\noindent\emph{Razlog:} Jednakogo vida naměrno povtarja raněje oprěděljenu relaciju, svjazanu s izomorfizmom; ne znači proste podobje zapisa. V Teoremu IX bytje prinajmenje dvuh takyh prostyh modulov jest kriterij beskonečno mnogyh razloženj, pri uslovju beskonečno mnogyh konstant i ostalyh ukazanyh hipotezah. Soglasovanje nazvy medžu paragrafami činji zavisimost dokaza vidlivoju. Ono ne dopušča opustiti te hipotezy ili zaměniti izomorfnost bukvalnoju jednakostju modulov.\par
\noindent\emph{Reason:} Jednakogo vida deliberately repeats the previously defined relation linked to isomorphism; it does not mean mere similarity of notation. In Theorem IX the presence of at least two such simple modules is a criterion for infinitely many decompositions, assuming infinitely many constants and the other stated hypotheses. Consistency of the name across sections makes the proof's dependence visible. It permits neither removal of those hypotheses nor replacement of isomorphism by literal equality of modules.\par
\noindent{\footnotesize Izvorne zapisy / Original records: HIST-g079-061.}\par
\smallskip\noindent\emph{Sravnjenje s tekstom:} Beskonečno mnogo konstant ostaje hipotezoju.\par
\noindent\emph{Textual comparison:} Infinitely many constants remains an assumption.\par
\noindent{\footnotesize TEXT-COMP-917; pasaže / passages: 11908, 11927, 11928, 11929, 11930, 11931, 11935, 11936, 11937; \texttt{CORPUS\_FORM\_COMPARISONS.json}.}\par
\end{minipage}\par
\par\medskip\Needspace{7\baselineskip}
\noindent\begin{minipage}{\linewidth}
\hypertarget{TERM-1861}{}
\noindent\textbf{jednakogo vida}\quad{\footnotesize\texttt{TERM-1861}}\par
\noindent\emph{Iztočnik / Source:} von gleicher Art\par
\noindent\emph{Rusky / Russian:} того же вида\par
\noindent\emph{Ukrajinsky / Ukrainian:} того самого виду\par
\smallskip\noindent\emph{Razlog:} Jednakogo vida shranjaje Loewyjovu i Blumbergovu nazvu relacije, vměsto toho, aby byla od početka zaměnjena izomorfizmom. Definicija potrěbuje dva polinomy $P$ i $Q$ s uslovjami vzajemnoj prostosti i kongruencijami (25)–(28). Teorem VI teprve dokazuje ravnocěnnost s izomorfizmom grup ostatkov. Zato odrěđene nazvy morajut ostati razděljene: ih neposrědno sliti značilo by skryti matematičny sodržaj teorema. Historijna ograda o jezyčnoj formě ostaje.\par
\noindent\emph{Reason:} Jednakogo vida retains Loewy's and Blumberg's name for the relation instead of replacing it with isomorphism from the outset. The definition requires polynomials $P$ and $Q$ satisfying coprimeness conditions and congruences (25)–(28). Theorem VI then proves equivalence with isomorphism of the residue groups. The names must therefore remain distinct: merging them immediately would hide the mathematical content of the theorem. The recorded linguistic reservation remains.\par
\noindent{\footnotesize Izvorne zapisy / Original records: HIST-g079-006.}\par
\smallskip\noindent\emph{Sravnjenje s tekstom:} Zaměna nazvy izbrisala by tvrđenje o ravnocěnnosti.\par
\noindent\emph{Textual comparison:} Replacing the name would erase the equivalence being asserted.\par
\noindent{\footnotesize TEXT-COMP-918; pasaže / passages: 11648, 11650, 11652, 11653, 11655, 11657, 11658, 11660, 11664; \texttt{CORPUS\_FORM\_COMPARISONS.json}.}\par
\end{minipage}\par
\par\medskip\Needspace{7\baselineskip}
\noindent\begin{minipage}{\linewidth}
\hypertarget{TERM-1872}{}
\noindent\textbf{sistem vrědnostij}\quad{\footnotesize\texttt{TERM-1872}}\par
\noindent\emph{Iztočnik / Source:} Wertesystem\par
\noindent\emph{Rusky / Russian:} система значений\par
\noindent\emph{Ukrajinsky / Ukrainian:} набір значень\par
\smallskip\noindent\emph{Razlog:} Sistem vrědnostij tu jest spis čislovyh vrědnostij $\alpha$, ktore štečut osnovne kontrakcije. Aktualny tekst ima sistem cěnnostij i predstavjaje te vrědnosti kako koordinate rěšetkovej točky. To ne jest etičny sistem cěnnostij ili abstraktna novaja algebra. Jednoj formě odpovědaje jedna takova točka, ale razne formy mogut ležati pri istoj točki i trebujut dodatnogo uporjadkovanja. Glosa ne smi prevračati tu korespondenciju v neopravdanu bijekciju.\par
\noindent\emph{Reason:} Sistem vrědnostij here is a tuple of numerical $\alpha$ values counting elementary contractions. The current text has sistem cěnnostij and presents the values as coordinates of a lattice point. This is not an ethical value system or a newly defined abstract algebra. One form has one such point, but different forms can lie at the same point and need additional ordering. The gloss must not turn this correspondence into an unwarranted bijection.\par
\noindent{\footnotesize Izvorne zapisy / Original records: HIST-g038-005.}\par
\smallskip\noindent\emph{Sravnjenje s tekstom:} Razne formy mogut iměti tu samuju točku.\par
\noindent\emph{Textual comparison:} Different forms can have the same point.\par
\noindent{\footnotesize TEXT-COMP-929; pasaže / passages: 4864, 4906; \texttt{CORPUS\_FORM\_COMPARISONS.json}.}\par
\end{minipage}\par
\par\medskip\Needspace{7\baselineskip}
\noindent\begin{minipage}{\linewidth}
\hypertarget{TERM-1873}{}
\noindent\textbf{sistem vrědnostij}\quad{\footnotesize\texttt{TERM-1873}}\par
\noindent\emph{Iztočnik / Source:} Wertsystem\par
\noindent\emph{Rusky / Russian:} система значений\par
\noindent\emph{Ukrajinsky / Ukrainian:} система значень\par
\smallskip\noindent\emph{Razlog:} Sistem vrědnostij označuje priděljenje vrědnostij prěměnnym. Arhiv dava prednost vrědnostij prěd cennostij, aby izběgnul cěnovogo smysla; aktualna pročitana pasaža však piše sistem cěnnostij. Zato ranějša preferencija ne može byti predstavjena kako doslědno provedena norma. Pasaža govori o obščih nuljah homogennyh form i odličenju nulovogo priděljenja. Matematična glosa razjasnjaje vyraz bez tvrdženja, že samo vrědnostij jest dopustimo v celom jezyku.\par
\noindent\emph{Reason:} Sistem vrědnostij denotes an assignment of values to variables. The archive prefers vrědnostij to cennostij to avoid a price-related sense; the examined current passage nevertheless writes sistem cěnnostij. The earlier preference cannot therefore be presented as a consistently implemented norm. The passage concerns common zeros of homogeneous forms and distinguishes the zero assignment. A mathematical gloss clarifies the expression without claiming that only vrědnostij is permissible throughout the language.\par
\noindent{\footnotesize Izvorne zapisy / Original records: HIST-g051-005.}\par
\smallskip\noindent\emph{Sravnjenje s tekstom:} Aktualno cěnnostij ne soglašaje se s ranějšeju preferencijeju vrědnostij.\par
\noindent\emph{Textual comparison:} Current cěnnostij does not follow the earlier preference for vrědnostij.\par
\noindent{\footnotesize TEXT-COMP-930; pasaže / passages: 6125, 6127, 6129, 6131; \texttt{CORPUS\_FORM\_COMPARISONS.json}.}\par
\end{minipage}\par
\par\medskip\Needspace{7\baselineskip}
\noindent\begin{minipage}{\linewidth}
\hypertarget{TERM-1874}{}
\noindent\textbf{sistemy vrědnostij}\quad{\footnotesize\texttt{TERM-1874}}\par
\noindent\emph{Iztočnik / Source:} Wertsysteme\par
\noindent\emph{Rusky / Russian:} системы значений\par
\noindent\emph{Ukrajinsky / Ukrainian:} системи значень\par
\smallskip\noindent\emph{Razlog:} Sistemy vrědnostij čuva množstvenost razgledanyh priděljenj. V različnyh čestjah argumenta idě o obrazy čisel pri odobraženju i o konečne spisy vrědnostij, potrebne za uslovje ranga. Aktualne pasaže imaju sistemy cěnnostij. Shranjajemo zato ranějši izbor i aktualnu variantu, ne vstavjajemo vměsto njih novu vektorno-prostornu notaciju. Racionalne i algebraične zavisimosti i precizny rang dolžne ostati svjazane s konkretnymi priděljenjami.\par
\noindent\emph{Reason:} Sistemy vrědnostij retains the plurality of the assignments being considered. Different parts of the argument concern images of numbers under a mapping and finite tuples used in a rank condition. The current passages have sistemy cěnnostij. Both the recorded choice and the current variant are therefore preserved, without substituting new vector-space notation. Rational and algebraic dependencies and the exact rank remain attached to the particular assignments.\par
\noindent{\footnotesize Izvorne zapisy / Original records: HIST-g070-001, HIST-g072-010.}\par
\smallskip\noindent\emph{Sravnjenje s tekstom:} Razne priměnjenja frazy trebujut konkretne argumenty.\par
\noindent\emph{Textual comparison:} The different uses of the phrase require their particular arguments.\par
\noindent{\footnotesize TEXT-COMP-931; pasaže / passages: 8487, 8489, 8491, 8642, 8644, 8646; \texttt{CORPUS\_FORM\_COMPARISONS.json}.}\par
\end{minipage}\par
\par\medskip\Needspace{7\baselineskip}
\noindent\begin{minipage}{\linewidth}
\hypertarget{TERM-1875}{}
\noindent\textbf{suštno različne}\quad{\footnotesize\texttt{TERM-1875}}\par
\noindent\emph{Iztočnik / Source:} wesentlich verschieden\par
\noindent\emph{Rusky / Russian:} существенно отличные\par
\noindent\emph{Ukrajinsky / Ukrainian:} суттєво відмінні\par
\smallskip\noindent\emph{Razlog:} Suštno različne jest historijna forma aktualnogo sučstveno različne. Tu sučstvenost ne jest stilistična ocěna: tekst razděljaje oblasti po izomorfizmu i govori o oblastjah, ktore ne nastavajut tym samym konstrukcijnym principom. Daljša pasaža izrično piše neizomorfne, ili sučstveno različne. Zato slovnik može dati točnu glosu neizomorfne v tom kontekstu, ale ne zaměnjati vsako sučstveno različne v inoj rabotě bez prověrky relacije.\par
\noindent\emph{Reason:} Suštno različne is the historical form of current sučstveno različne. Essential difference is not a stylistic assessment here: the text classifies domains by isomorphism and discusses domains not produced by the same construction principle. A later passage expressly says nonisomorphic, or essentially different. The dictionary can therefore give the precise gloss nonisomorphic in this context, but cannot replace every occurrence of essentially different in another work without checking the relation.\par
\noindent{\footnotesize Izvorne zapisy / Original records: HIST-g066-002.}\par
\smallskip\noindent\emph{Sravnjenje s tekstom:} Tekst sam daje točnu glosu neizomorfne.\par
\noindent\emph{Textual comparison:} The text itself supplies the precise gloss nonisomorphic.\par
\noindent{\footnotesize TEXT-COMP-932; pasaže / passages: 8044, 8045, 8046, 8047, 8049, 8050, 8075, 8076, 8077, 8078, 8079, 8080, 8081, 8082, 8083; \texttt{CORPUS\_FORM\_COMPARISONS.json}.}\par
\end{minipage}\par
\par\medskip\Needspace{7\baselineskip}
\noindent\begin{minipage}{\linewidth}
\hypertarget{TERM-1876}{}
\noindent\textbf{suštstvene parametry / libovoljne funkcije}\quad{\footnotesize\texttt{TERM-1876}}\par
\noindent\emph{Iztočnik / Source:} wesentliche Parameter / willkürliche Funktionen\par
\noindent\emph{Rusky / Russian:} существенные параметры / произвольные функции\par
\noindent\emph{Ukrajinsky / Ukrainian:} істотні параметри / довільні функції\par
\smallskip\noindent\emph{Razlog:} Suštstvene parametry i libovoljne funkcije razlikujut dva vida zavisimosti grupy transformacij. Aktualna forma jest sučstvene parametry. Definicija izrično objasnjaje sučstvenost: zadano čislo parametrov ne možno zameniti menšim čislom. Libovoljne funkcije noset svoju funkcionalnu svobodu i potrebne uslovja regularnosti; one ne sut prosto mnogo čislovyh parametrov. Meždu obojimi slučajami tekst ostavlja osobny graničny slučaj beskonečno mnogyh parametrov. Izbor obščih slov ne smi sliti te kategorije.\par
\noindent\emph{Reason:} Suštstvene parametry and libovoljne funkcije distinguish two kinds of dependence of a transformation group. The current form is sučstvene parametry. The definition makes essential explicit: the given number of parameters cannot be replaced by a smaller number. Arbitrary functions carry functional freedom and the required regularity conditions; they are not merely many numerical parameters. Between the two cases the text retains the separate limiting case of infinitely many parameters. Choosing ordinary words must not collapse these categories.\par
\noindent{\footnotesize Izvorne zapisy / Original records: HIST-g075-005.}\par
\smallskip\noindent\emph{Sravnjenje s tekstom:} Beskonečno mnogo parametrov ne jest avtomatično libovoljna funkcija.\par
\noindent\emph{Textual comparison:} Infinitely many parameters does not automatically mean an arbitrary function.\par
\noindent{\footnotesize TEXT-COMP-933; pasaže / passages: 9201; \texttt{CORPUS\_FORM\_COMPARISONS.json}.}\par
\end{minipage}\par
\par\medskip\Needspace{7\baselineskip}
\noindent\begin{minipage}{\linewidth}
\hypertarget{TERM-1879}{}
\noindent\textbf{mimobežna prjama}\quad{\footnotesize\texttt{TERM-1879}}\par
\noindent\emph{Iztočnik / Source:} windschiefe Gerade\par
\noindent\emph{Rusky / Russian:} скрещивающаяся прямая\par
\noindent\emph{Ukrajinsky / Ukrainian:} мимобіжна пряма\par
\smallskip\noindent\emph{Razlog:} Mimobežna prjama izbrana jest za prostranstvenu geometričnu konfiguraciju; ukrajinska forma blizko srovnjaje se s tim korenjem, ruska upotrěbjaje drugu tvorbu. Tu prjamy ne peresěkajut se i ne ležat v jednoj ravnině. Relacija jest poparna: dvě prjamy, ktore sut mimobežne k tretej, v samom priměru peresěkajut se medžu soboju. Slovnik zato mora označiti, koja para jest razgledana. Arhivna ocěna prozračnosti i ograda o jezyčnoj prověrke ne sut rezultat čitatelskogo eksperimenta.\par
\noindent\emph{Reason:} Mimobežna prjama was chosen for a spatial geometric configuration; the Ukrainian form closely parallels its stem, while Russian uses a different formation. The relevant lines do not intersect and are not coplanar. The relation is pairwise: the example's two lines that are skew to a third intersect each other. The dictionary must therefore identify the pair concerned. The archived assessment of transparency and request for linguistic review are not results of a reader experiment.\par
\noindent{\footnotesize Izvorne zapisy / Original records: HIST-g081-084.}\par
\smallskip\noindent\emph{Sravnjenje s tekstom:} Dvě prjamy mogut se peresěkati, hoti obe sut mimobežne k tretej.\par
\noindent\emph{Textual comparison:} Two lines may intersect although both are skew to a third.\par
\noindent{\footnotesize TEXT-COMP-936; pasaže / passages: 13479, 13481, 13484, 13485, 13486, 13487; \texttt{CORPUS\_FORM\_COMPARISONS.json}.}\par
\end{minipage}\par
\par\medskip\Needspace{7\baselineskip}
\noindent\begin{minipage}{\linewidth}
\hypertarget{TERM-1880}{}
\noindent\textbf{dobro uporjadkovana množina / dobro uporjadkovanje}\quad{\footnotesize\texttt{TERM-1880}}\par
\noindent\emph{Iztočnik / Source:} wohlgeordnete Menge / Wohlordnung\par
\noindent\emph{Rusky / Russian:} вполне упорядоченное множество / вполне упорядочение\par
\noindent\emph{Ukrajinsky / Ukrainian:} добре впорядкована множина / добре впорядкування\par
\smallskip\noindent\emph{Razlog:} Dobro uporjadkovana množina i dobro uporjadkovanje nazyvajut tehničny poredok, ne prosto uredny spis. Primětka upotrěbjaje dobry poredok funkcij dlja poslědovateljnogo postrojenja dopustimogo sistema: funkcije prijmajut se ili izključajut se po zadanoj algebraičnoj prověrke. Procedura može byti prěkinjena na libovoljnom městu, a razgledanje vseh poredkov i měst prěkinjenja dava vse dopustime sistemy. Zato smysl ne možno skratiti do prosto izbrati funkcije; redoslednost jest čest konstrukcije.\par
\noindent\emph{Reason:} Dobro uporjadkovana množina and dobro uporjadkovanje name a technical ordering, not merely a tidy list. The footnote uses a well-order of functions to construct an admissible system successively: functions are admitted or excluded by a specified algebraic test. The procedure may stop at any point, and considering every ordering and stopping point gives all admissible systems. The meaning cannot therefore be reduced simply to choose functions; the order is part of the construction.\par
\noindent{\footnotesize Izvorne zapisy / Original records: HIST-g062-011.}\par
\smallskip\noindent\emph{Sravnjenje s tekstom:} Redoslednost i město prěkinjenja sut česti postupka.\par
\noindent\emph{Textual comparison:} The order and stopping point are parts of the procedure.\par
\noindent{\footnotesize TEXT-COMP-937; pasaže / passages: 7564, 7565, 7566, 7567, 7568, 7569, 7570, 7571, 7572, 7573, 7574; \texttt{CORPUS\_FORM\_COMPARISONS.json}.}\par
\end{minipage}\par
\par\medskip\Needspace{7\baselineskip}
\noindent\begin{minipage}{\linewidth}
\hypertarget{TERM-1881}{}
\noindent\textbf{dobro uporjadkovanje}\quad{\footnotesize\texttt{TERM-1881}}\par
\noindent\emph{Iztočnik / Source:} Wohlordnung\par
\noindent\emph{Rusky / Russian:} вполне упорядочение\par
\noindent\emph{Ukrajinsky / Ukrainian:} добре впорядкування\par
\smallskip\noindent\emph{Razlog:} Dobro uporjadkovanje prodolžaje termin iz uvoda i shranjaje tehničny smysl dobrogo poredka. V sekciji 8 ono služi izboru predstaviteljev: oblast prijmaje se, ako ne jest izomorfna nijednoj prědhodnoj oblasti. Prva vozmožna iznimka potom vodi k proturěčenju. V ranějšem priměru, naprotiv, nijedno dobro uporjadkovanje ne dava želanyh svojstv. Jedna nazva mora dopuščati oba logična konteksta; ona ne sugeruje, že libovoljny izbor daje tu samu bazu. Arhiv ostavlja jezyčnu formu dlja ocěny, bez čitateljske prověrky razumljivosti.\par
\noindent\emph{Reason:} Dobro uporjadkovanje continues the introductory term and retains the technical sense of a well-order. In section 8 it selects representatives: a domain is admitted if it is not isomorphic to any preceding domain. A first possible exception then yields a contradiction. In an earlier example, by contrast, no well-order gives the desired properties. One name must support both logical contexts; it does not suggest that arbitrary choices produce the same basis. The archive leaves the linguistic form open to assessment, without a reader test of intelligibility.\par
\noindent{\footnotesize Izvorne zapisy / Original records: HIST-g060-011, HIST-g066-006.}\par
\smallskip\noindent\emph{Sravnjenje s tekstom:} Ista nazva služi raznym kvantifikacijam.\par
\noindent\emph{Textual comparison:} The same name occurs under different quantifiers.\par
\noindent{\footnotesize TEXT-COMP-938; pasaže / passages: 7234, 7244, 8091, 8092, 8093, 8094, 8095, 8096, 8097, 8098, 8099, 8100, 8101, 8102; \texttt{CORPUS\_FORM\_COMPARISONS.json}.}\par
\end{minipage}\par
\par\medskip\Needspace{7\baselineskip}
\noindent\begin{minipage}{\linewidth}
\hypertarget{TERM-1882}{}
\noindent\textbf{dobro uporjadkovanje / porjadok elementov}\quad{\footnotesize\texttt{TERM-1882}}\par
\noindent\emph{Iztočnik / Source:} Wohlordnung / Anordnung der Elemente\par
\noindent\emph{Rusky / Russian:} [Bukvy ne sut zachovane / Lettering not recoverable]\par
\noindent\emph{Ukrajinsky / Ukrainian:} [Bukvy ne sut zachovane / Lettering not recoverable]\par
\smallskip\noindent\emph{Razlog:} Dobro uporjadkovanje i porjadok elementov ne sut dve svobodno zaměnjive frazy. Prva imenuje osobly typ poredka; druga govori o redoslědnosti elementov. Definicija racionalnoj bazy potrěbuje prinajmenje jedno dobro uporjadkovanje. Primětka pokazuje, že pri poredku $\eta,\sqrt[\nu]{\eta}$ potrěbne sut oba elementy, a pri obratnom samo korenj. Ta zavisimost od poredka jest matematičny razlog shranjenja razliky. Aktualna proza piše poredok. Poškodžene ruske i ukrajinske zapisy ne davajut čitljive sravniteljne formy; slovnik ih ne rekonstruuje iz domysla.\par
\noindent\emph{Reason:} Dobro uporjadkovanje and porjadok elementov are not freely interchangeable phrases. The former names a special kind of order; the latter refers to the sequence of the elements. The definition of a rational basis requires at least one well-order. The footnote shows that in the order $\eta,\sqrt[\nu]{\eta}$ both elements are required, whereas the reverse order needs only the root. That dependence on order is the mathematical reason for retaining the distinction. Current prose writes poredok. The damaged Russian and Ukrainian records provide no readable comparative forms; the dictionary does not reconstruct them by conjecture.\par
\noindent{\footnotesize Izvorne zapisy / Original records: HIST-g064-003.}\par
\smallskip\noindent\emph{Sravnjenje s tekstom:} Obratny poredok měnja izbor bazovyh elementov.\par
\noindent\emph{Textual comparison:} Reversing the order changes the selected basis elements.\par
\noindent{\footnotesize TEXT-COMP-939; pasaže / passages: 7718, 7719, 7720, 7727, 7728, 7729, 7730, 7731, 7732; \texttt{CORPUS\_FORM\_COMPARISONS.json}.}\par
\end{minipage}\par
\par\medskip\Needspace{7\baselineskip}
\noindent\begin{minipage}{\linewidth}
\hypertarget{TERM-1883}{}
\noindent\textbf{dobro uporjadkovanje / teorema o dobrom uporjadkovanju}\quad{\footnotesize\texttt{TERM-1883}}\par
\noindent\emph{Iztočnik / Source:} Wohlordnung / Wohlordnungssatz\par
\noindent\emph{Rusky / Russian:} вполне упорядочение / теорема о вполне упорядочении\par
\noindent\emph{Ukrajinsky / Ukrainian:} добре впорядкування / теорема про добре впорядкування\par
\smallskip\noindent\emph{Razlog:} Dobro uporjadkovanje nazyvaje objekt, a teorema o dobrom uporjadkovanju rezultat, ktory dopušča jego izbor. Arhiv naměrno prodolžaje terminologiju prědhodnoj raboty. V novom priměnjenju čisla, ktore ne možno racionalno izraziti prěz konečno mnogo prědhodnyh, tvoret bazu; poredok potom prěnosi se na obrazy i određuje polja početnyh odrězkov. Glosa proto mora shraniti i poredok, i obseg kvantifikacije. Nazva teorema ne jest sama dokazom vseh dalših tvrđenij o obrazah.\par
\noindent\emph{Reason:} Dobro uporjadkovanje names an object, while teorema o dobrom uporjadkovanju names the result permitting its choice. The record deliberately continues terminology from the preceding paper. In the new application, numbers not rationally expressible through finitely many predecessors form a basis; the order is then transported to the images and determines initial-segment fields. The gloss must therefore retain both the order and the scope of quantification. Naming the theorem does not by itself prove every subsequent claim about the images.\par
\noindent{\footnotesize Izvorne zapisy / Original records: HIST-g070-004.}\par
\smallskip\noindent\emph{Sravnjenje s tekstom:} Poredok prěnosi se na obrazy bazy.\par
\noindent\emph{Textual comparison:} The order is transported to the images of the basis.\par
\noindent{\footnotesize TEXT-COMP-940; pasaže / passages: 8489, 8491; \texttt{CORPUS\_FORM\_COMPARISONS.json}.}\par
\end{minipage}\par
\par\medskip\Needspace{7\baselineskip}
\noindent\begin{minipage}{\linewidth}
\hypertarget{TERM-1885}{}
\noindent\textbf{teorema o dobrom uporjadkovanju}\quad{\footnotesize\texttt{TERM-1885}}\par
\noindent\emph{Iztočnik / Source:} Wohlordnungssatz\par
\noindent\emph{Rusky / Russian:} теорема о вполне упорядочении\par
\noindent\emph{Ukrajinsky / Ukrainian:} теорема про добре впорядкування\par
\smallskip\noindent\emph{Razlog:} Teorema o dobrom uporjadkovanju jest razložena fraza vměsto německogo složenogo slova. Zachovany razlog sravnjava ukrajinsku konstrukciju s dobro i rusku s vpolně; ne dava čislovyh vah ni dokaza všeslovjanskoj razumljivosti. Uvod izrično pripisuje Zermelu konstrukciju, osnovanu na toj teoremě, a zatym objašnjava egzistenciju algebraičnoj bazy. Toj kredit i prelaz od poredka k bazi sut česti smisla. Kratka nazva ne smi prepisati avtorstvo konstrukcije prevoditelju.\par
\noindent\emph{Reason:} Teorema o dobrom uporjadkovanju expands the German compound into a phrase. The preserved rationale compares the Ukrainian construction with dobro and the Russian one with vpolně; it supplies neither numerical weights nor a demonstration of pan-Slavic intelligibility. The introduction explicitly credits Zermelo with a construction based on the theorem and then explains the existence of an algebraic basis. That attribution and the passage from ordering to a basis are part of the meaning. The short name must not transfer authorship of the construction to the translator.\par
\noindent{\footnotesize Izvorne zapisy / Original records: HIST-g058-003.}\par
\smallskip\noindent\emph{Sravnjenje s tekstom:} Teorem služi konstrukciji s izričnym ljudskim kreditom.\par
\noindent\emph{Textual comparison:} The theorem serves a construction with explicit human attribution.\par
\noindent{\footnotesize TEXT-COMP-942; pasaže / passages: 7134, 7142; \texttt{CORPUS\_FORM\_COMPARISONS.json}.}\par
\end{minipage}\par
\par\medskip\Needspace{7\baselineskip}
\noindent\begin{minipage}{\linewidth}
\hypertarget{TERM-1886}{}
\noindent\textbf{korenjevo polje}\quad{\footnotesize\texttt{TERM-1886}}\par
\noindent\emph{Iztočnik / Source:} Wurzelkörper\par
\noindent\emph{Rusky / Russian:} поле корней\par
\noindent\emph{Ukrajinsky / Ukrainian:} поле коренів\par
\smallskip\noindent\emph{Razlog:} Korenjevo polje izbrano jest kako prozračna tvorba, soglasna s poljem kvocientov i medžupoljem, vměsto novogo zaimstvovanja. Rusky i ukrajinsky zapisy imajo istu korenevu metaforu. V točnoj pasaži objekt jest $P^{1/p}$ pri pozitivnoj karakteristike $p$, s izričnym uslovjem konečnosti razširjenja. To ne jest prosto polje vseh korenjev libovoljnogo uravnenja; glosa mora shraniti tu konstrukciju. Tekst oddělno kaže, že pri karakteristike nulja takogo ograničenja ne potrěbuje.\par
\noindent\emph{Reason:} Korenjevo polje was chosen as a transparent formation consistent with polje kvocientov and medžupolje, rather than a new borrowing. The Russian and Ukrainian records use the same root metaphor. In the exact passage the object is $P^{1/p}$ in positive characteristic $p$, with an explicit finiteness assumption on the extension. It is not simply the field of all roots of an arbitrary equation; the gloss must preserve this construction. The text separately says that no such restriction is needed in characteristic zero.\par
\noindent{\footnotesize Izvorne zapisy / Original records: HIST-g106-011.}\par
\smallskip\noindent\emph{Sravnjenje s tekstom:} Konečnost razširjenja jest izrična hipoteza.\par
\noindent\emph{Textual comparison:} Finiteness of the extension is an explicit hypothesis.\par
\noindent{\footnotesize TEXT-COMP-943; pasaže / passages: 16019, 16021; \texttt{CORPUS\_FORM\_COMPARISONS.json}.}\par
\end{minipage}\par
\par\medskip\Needspace{7\baselineskip}
\noindent\begin{minipage}{\linewidth}
\hypertarget{TERM-1887}{}
\noindent\textbf{korenevo čislo}\quad{\footnotesize\texttt{TERM-1887}}\par
\noindent\emph{Iztočnik / Source:} Wurzelzahl\par
\noindent\emph{Rusky / Russian:} корневое число\par
\noindent\emph{Ukrajinsky / Ukrainian:} кореневе число\par
\smallskip\noindent\emph{Razlog:} Korenevo čislo zachovuje historijnu nazvu faktora, svjazanu s Lagrangevymi korenevymi čislami. Tu čislo ne znači čislenost množiny korenjev. Tekst razklada grupovu determinantu na faktory $D_\lambda$ i dokazuje, že pri dejstviju grupy oni umnožajut se korenjem jedinice; v cikličnom slučaju dobyva Lagrangevu formulu. Zato doslovna forma potrěbuje tu krátku matematičnu glosu. Pripisanje razklada Speiserovym zaključkam takože ostaje, vměsto predstavljanja nazvy kako novogo izobretenja prevoda.\par
\noindent\emph{Reason:} Korenevo čislo retains the historical name for a factor connected with Lagrange's root numbers. Here number does not mean the cardinality of a set of roots. The text factors the group determinant into factors $D_\lambda$ and proves that the group action multiplies them by roots of unity; in the cyclic case it obtains Lagrange's formula. The literal wording therefore needs this short mathematical gloss. The attribution of the decomposition to Speiser's arguments also remains, rather than presenting the name as an invention of the translation.\par
\noindent{\footnotesize Izvorne zapisy / Original records: HIST-g207-008.}\par
\smallskip\noindent\emph{Sravnjenje s tekstom:} Čislo tu ne znači čislenost korenjev.\par
\noindent\emph{Textual comparison:} Number here does not mean a count of roots.\par
\noindent{\footnotesize TEXT-COMP-944; pasaže / passages: 22884, 22886, 22888, 22890, 22892, 22894, 22895, 22897, 22899, 22901, 22948; \texttt{CORPUS\_FORM\_COMPARISONS.json}.}\par
\end{minipage}\par
\par\medskip\Needspace{7\baselineskip}
\noindent\begin{minipage}{\linewidth}
\hypertarget{TERM-1889}{}
\noindent\textbf{čislove i funkcionalne polja}\quad{\footnotesize\texttt{TERM-1889}}\par
\noindent\emph{Iztočnik / Source:} Zahl- und Funktionenkörper\par
\noindent\emph{Rusky / Russian:} числовые и функциональные поля\par
\noindent\emph{Ukrajinsky / Ukrainian:} числові та функціональні поля\par
\smallskip\noindent\emph{Razlog:} Čislove i funkcionalne polja jest historijna pridavnikova para. Zapis dava prednost funkcionalne zato, že slědujuča pasaža govori i o funkcionalnoj oblasti. To ne čini polje i oblast identičnymi objektami: tekst izhodi od prstena glavnyh idealov, tvori jego polje kvocientov i potom konečno razširjenje prvogo roda. Pojem cělyh elementov takože jest fiksovan. Alternativa polja funkcij iz drugoj raboty jezikovo jest prozračnija, ale sravnjenje dolžno pokazati oba zabilěžena izbora, ne prikazati jeden kakoby vsegda byl jedinym.\par
\noindent\emph{Reason:} Čislove i funkcionalne polja is the historical adjectival pair. The record prefers funkcionalne because the following passage also discusses a functional domain. This does not make a field and a domain identical objects: the text begins with a principal-ideal ring, forms its fraction field and then a finite extension of the first kind. The notion of integral elements is also fixed. Polja funkcij from another paper is more explicit linguistically, but the comparison must show both recorded choices rather than pretend that one was always the sole choice.\par
\noindent{\footnotesize Izvorne zapisy / Original records: HIST-g149-001.}\par
\smallskip\noindent\emph{Sravnjenje s tekstom:} Polje i funkcionalna oblast ne sut avtomatično isti objekt.\par
\noindent\emph{Textual comparison:} The field and the functional domain are not automatically the same object.\par
\noindent{\footnotesize TEXT-COMP-946; pasaže / passages: 18557, 18559; \texttt{CORPUS\_FORM\_COMPARISONS.json}.}\par
\end{minipage}\par
\par\medskip\Needspace{7\baselineskip}
\noindent\begin{minipage}{\linewidth}
\hypertarget{TERM-1890}{}
\noindent\textbf{čislove polja i polja funkcij}\quad{\footnotesize\texttt{TERM-1890}}\par
\noindent\emph{Iztočnik / Source:} Zahl- und Funktionenkörper\par
\noindent\emph{Rusky / Russian:} числовые поля и поля функций\par
\noindent\emph{Ukrajinsky / Ukrainian:} числові поля і поля функцій\par
\smallskip\noindent\emph{Razlog:} Čislove polja i polja funkcij jest druga zabilěžena forma toj samej predmetnoj pary. Izrično polja funkcij bylo izbrano vměsto staršego funkcionalne polja i soglasno s uvodom raboty. Kontekst dokazuje izpolnjenje aksiomov dlja razgledanyh konstrukcij; ne govori, že libovoljny prsten funkcij imaje te svojstva. Nazva zato ostaje povezana s koeficientnym poljem, razširjenjem i določenoju oblastju cělosti. Ta forma i pridavnikova varianta iz drugoj raboty trebujut jednoj jasno označenoj slovnikovoj svezi.\par
\noindent\emph{Reason:} Čislove polja i polja funkcij is the other recorded form of the same subject pair. Explicit polja funkcij was chosen in place of the older funkcionalne polja and in agreement with the paper's introduction. The context proves the axioms for the constructions under discussion; it does not say that any ring of functions has those properties. The name therefore remains tied to the coefficient field, extension and specified domain of integral elements. This form and the adjectival variant in another paper need a clearly marked dictionary cross-reference.\par
\noindent{\footnotesize Izvorne zapisy / Original records: HIST-g110-007.}\par
\smallskip\noindent\emph{Sravnjenje s tekstom:} Nazva dolžna ostati svjazana s ukazanymi aksiomami.\par
\noindent\emph{Textual comparison:} The name must remain connected to the stated axioms.\par
\noindent{\footnotesize TEXT-COMP-947; pasaže / passages: 16109, 16115, 16411; \texttt{CORPUS\_FORM\_COMPARISONS.json}.}\par
\end{minipage}\par
\par\medskip\Needspace{7\baselineskip}
\noindent\begin{minipage}{\linewidth}
\hypertarget{TERM-1892}{}
\noindent\textbf{čislovo-teoretičny teorem sušstvovanja}\quad{\footnotesize\texttt{TERM-1892}}\par
\noindent\emph{Iztočnik / Source:} zahlentheoretischer Existenzsatz\par
\noindent\emph{Rusky / Russian:} теоретико-числовая теорема существования\par
\noindent\emph{Ukrajinsky / Ukrainian:} числово-теоретична теорема існування\par
\smallskip\noindent\emph{Razlog:} Čislovo-teoretičny teorem sušstvovanja čuva i disciplinu, i ekzistencijalny harakter rezultata. Aktualna forma jest teorem obstajanja. Arhiv sravnjava poredok čestij v ukrajinskom i medžuslovjanskom izrazu s ruskom teoretiko-čislovaja. V tekstu dokaz silnějšego teorema pripisany jest Hassemu; potom Brauerov elementarny postupok dokazuje slabši, ale dostatočny rezultat. Slovnik ne smi sliti oba teorema ni izbrisati razděljenje doprinosa medžu Hassem, Noether i Brauerom.\par
\noindent\emph{Reason:} Čislovo-teoretičny teorem sušstvovanja preserves both the discipline and the existential character of the result. Current wording is teorem obstajanja. The record compares the order of the components in Ukrainian and Interslavic with Russian teoretiko-čislovaja. The text credits Hasse with the stronger theorem's proof; Brauer's elementary procedure then proves a weaker but sufficient result. The dictionary must not merge the two theorems or erase the distribution of contributions among Hasse, Noether and Brauer.\par
\noindent{\footnotesize Izvorne zapisy / Original records: HIST-g168-005.}\par
\smallskip\noindent\emph{Sravnjenje s tekstom:} Silnějši i slabši rezultat ne trěba sliti.\par
\noindent\emph{Textual comparison:} The stronger and weaker results must not be merged.\par
\noindent{\footnotesize TEXT-COMP-949; pasaže / passages: 19244; \texttt{CORPUS\_FORM\_COMPARISONS.json}.}\par
\end{minipage}\par
\par\medskip\Needspace{7\baselineskip}
\noindent\begin{minipage}{\linewidth}
\hypertarget{TERM-1893}{}
\noindent\textbf{čislovo polje}\quad{\footnotesize\texttt{TERM-1893}}\par
\noindent\emph{Iztočnik / Source:} Zahlkörper\par
\noindent\emph{Rusky / Russian:} числовое поле\par
\noindent\emph{Ukrajinsky / Ukrainian:} числове поле\par
\smallskip\noindent\emph{Razlog:} Čislovo polje jest ustaljena slovjanska tvorba vměsto doslovnogo čislovogo těla. Ale obseg v toj rabotě jest širši od obyčnoj sovrěmennoj definicije konečnogo razširjenja racionalnyh čisel: avtorica izrično dopušča, že polje sodrživa vse kompleksne čisla. Glosa preto ne smi bez upozorjenja suziti iztočny tekst. Nazva vyznačuje oblast koeficientov, a konečnost ili algebraičnost morajut izhoditi iz konkretnoj hipotezy. Ranějša ocěna stabilnosti slova ne ukida tu potrebu za kontekstom.\par
\noindent\emph{Reason:} Čislovo polje is the established Slavic formation rather than a literal number body. Its scope in this paper is broader, however, than the usual modern definition as a finite extension of the rationals: the author explicitly allows the field to contain all complex numbers. The gloss must therefore not silently narrow the source. The term identifies the coefficient domain; finiteness or algebraicity must come from the particular hypothesis. The earlier assessment of lexical stability does not remove that need for context.\par
\noindent{\footnotesize Izvorne zapisy / Original records: HIST-g039-003.}\par
\smallskip\noindent\emph{Sravnjenje s tekstom:} Iztočny obseg dopušča vse kompleksne čisla.\par
\noindent\emph{Textual comparison:} The source's scope permits all complex numbers.\par
\noindent{\footnotesize TEXT-COMP-950; pasaže / passages: 4950, 5061; \texttt{CORPUS\_FORM\_COMPARISONS.json}.}\par
\end{minipage}\par
\par\medskip\Needspace{7\baselineskip}
\noindent\begin{minipage}{\linewidth}
\hypertarget{TERM-1894}{}
\noindent\textbf{čislove polja najnižšego stepena / stepen nad $P$}\quad{\footnotesize\texttt{TERM-1894}}\par
\noindent\emph{Iztočnik / Source:} Zahlkörper kleinsten Grades / Grad in bezug auf $P$\par
\noindent\emph{Rusky / Russian:} числовые поля наименьшей степени / степень над $P$\par
\noindent\emph{Ukrajinsky / Ukrainian:} числові поля найменшого степеня / степінь над $P$\par
\smallskip\noindent\emph{Razlog:} Čislove polja najnižšego stepena i stepen nad $P$ rabotajut s relativnym stepenjem razširjenja, ne s redom zapisa čisel. Predlog nad čini osnovno polje jasnym i skraćaje německu odnosnu konstrukciju. V Šurovom pytanju najnižši stepenj ne garantuje jedinstveno polje do izomorfizma: tekst izrično daje protivnu mogućnost. Takože minimalno polje, definirano odsutnostju vlastnogo razpadnogo podpolja, ne jest to samo kako polje najnižšego stepena. Te dve minimalnosti dolžno razlikovati i pri kratkom slovnikovom zaglavju.\par
\noindent\emph{Reason:} Čislove polja najnižšego stepena and stepen nad $P$ concern relative extension degree, not the order in which numbers are written. Nad makes the base field explicit and shortens the German relational construction. In Schur's question, least degree does not guarantee a unique field up to isomorphism: the text expressly allows otherwise. A minimal field, defined by the absence of a proper splitting subfield, is also not the same as a field of least degree. Those two notions of minimality must remain distinct even in a short dictionary heading.\par
\noindent{\footnotesize Izvorne zapisy / Original records: HIST-g167-001.}\par
\smallskip\noindent\emph{Sravnjenje s tekstom:} Minimalno polje ne znači nužno polje najnižšego stepena.\par
\noindent\emph{Textual comparison:} A minimal field need not be a field of least degree.\par
\noindent{\footnotesize TEXT-COMP-951; pasaže / passages: 19210, 19244; \texttt{CORPUS\_FORM\_COMPARISONS.json}.}\par
\end{minipage}\par
\par\medskip\Needspace{7\baselineskip}
\noindent\begin{minipage}{\linewidth}
\hypertarget{TERM-1895}{}
\noindent\textbf{teorija čislovyh polj / algebraično čislovo polje}\quad{\footnotesize\texttt{TERM-1895}}\par
\noindent\emph{Iztočnik / Source:} Zahlkörpertheorie / algebraischer Zahlkörper\par
\noindent\emph{Rusky / Russian:} теория числовых полей / алгебраическое числовое поле\par
\noindent\emph{Ukrajinsky / Ukrainian:} теорія числових полів / алгебраїчне числове поле\par
\smallskip\noindent\emph{Razlog:} Teorija čislovyh polj i algebraično čislovo polje sohranjajut slovjansku osnovu polje vměsto bukvalnogo německogo těla. V toj rabotě opora ne jest samo terminologična analogija: avtorica izrično definira algebraično čislovo polje kako polje $n$-go stepena nad racionalnymi čislami i sravnjava jego s poljem algebraičnyh funkcij. Zato pridev algebraično i osnovno polje nose konkretny sodržaj. Citacije Hilberta, Dirichleta i Dedekinda ostajut čestju naučnogo konteksta, ne dokazom, že oni upotrebljali medžuslovjanske slovo.\par
\noindent\emph{Reason:} Teorija čislovyh polj and algebraično čislovo polje retain the Slavic field formation rather than the literal German body. In this paper the support is not merely terminological analogy: the author expressly defines an algebraic number field as a field of degree $n$ over the rationals and compares it with an algebraic function field. The adjective algebraic and the base field therefore carry concrete content. Citations of Hilbert, Dirichlet and Dedekind remain part of the scholarly context, not evidence that they used the Interslavic word.\par
\noindent{\footnotesize Izvorne zapisy / Original records: HIST-g076-002.}\par
\smallskip\noindent\emph{Sravnjenje s tekstom:} Konkretna definicija određuje obseg nazvy.\par
\noindent\emph{Textual comparison:} The explicit definition determines the term's scope.\par
\noindent{\footnotesize TEXT-COMP-952; pasaže / passages: 9830, 9849, 9879; \texttt{CORPUS\_FORM\_COMPARISONS.json}.}\par
\end{minipage}\par
\par\medskip\Needspace{7\baselineskip}
\noindent\begin{minipage}{\linewidth}
\hypertarget{TERM-1896}{}
\noindent\textbf{centar}\quad{\footnotesize\texttt{TERM-1896}}\par
\noindent\emph{Iztočnik / Source:} Zentrum\par
\noindent\emph{Rusky / Russian:} центр\par
\noindent\emph{Ukrajinsky / Ukrainian:} центр\par
\smallskip\noindent\emph{Razlog:} Centar jest internacionalna nazva, blizka ruskomu i ukrajinskomu centru. Definicija v tekstu jest točna: vsi elementy, ktore komutujut so vsimi elementami prstena. Centar jest komutativny prsten, ale ne avtomatično polje; v osoblyh teoremah on stavaje tělom ili direktnoju sumoju komutativnyh těles pri ukazanyh hipotezah. Zato stabilnost nazvy ne dovoljuje vpisati svojstva osoblogo slučaja v obču definiciju. Zabilěžene centar i centr sut jezyčne varianty, ne dva matematična objekta.\par
\noindent\emph{Reason:} Centar is an international term close to Russian and Ukrainian centr. The definition in the text is exact: all elements commuting with every element of the ring. The center is a commutative ring, but not automatically a field; in particular theorems it becomes a division ring or a direct sum of commutative fields under stated hypotheses. Stability of the name does not permit importing the special case's properties into the general definition. Recorded centar and centr are linguistic variants, not two mathematical objects.\par
\noindent{\footnotesize Izvorne zapisy / Original records: HIST-g189-001, HIST-g192-005.}\par
\smallskip\noindent\emph{Sravnjenje s tekstom:} Centar ne jest bez dodatnyh uslovij polje.\par
\noindent\emph{Textual comparison:} A center is not a field without additional conditions.\par
\noindent{\footnotesize TEXT-COMP-953; pasaže / passages: 20724, 20726, 20728, 21228, 21260; \texttt{CORPUS\_FORM\_COMPARISONS.json}.}\par
\end{minipage}\par
\par\medskip\Needspace{7\baselineskip}
\noindent\begin{minipage}{\linewidth}
\hypertarget{TERM-1897}{}
\noindent\textbf{centar / predstavjenja centra}\quad{\footnotesize\texttt{TERM-1897}}\par
\noindent\emph{Iztočnik / Source:} Zentrum / Darstellungen des Zentrums\par
\noindent\emph{Rusky / Russian:} центр / представления центра\par
\noindent\emph{Ukrajinsky / Ukrainian:} центр / зображення центра\par
\smallskip\noindent\emph{Razlog:} Centar i predstavjenja centra svjazujut jednu strukturu s jej predstavjenjami; aktualny tekst piše prědstavjenja. Arhiv upozorjaje na klasificirujuču rolu centra. Ta rola imaje izrične hipotezy: pri opisanom razširjenju bez radikala absolutno ireducibilne predstavjenja odgovarjajut predstavjenjam centra prěz skalarne matrice. Obča fraza ne smi sugerovati takvu korespondenciju dlja libovoljnogo prstena. Soglasovanje nazvy medžu sekcijami pomoga čitatelju sledovati dokaz, no ne zaměnjaje uslovja teorema.\par
\noindent\emph{Reason:} Centar and predstavjenja centra connect a structure with its representations; the current text writes prědstavjenja. The record draws attention to the center's classifying role. That role has explicit hypotheses: for the described radical-free extension, absolutely irreducible representations correspond to representations of the center through scalar matrices. The general phrase must not suggest such a correspondence for an arbitrary ring. Consistency across sections helps readers follow the proof, but does not replace its hypotheses.\par
\noindent{\footnotesize Izvorne zapisy / Original records: HIST-g199-003.}\par
\smallskip\noindent\emph{Sravnjenje s tekstom:} Hipotezy o razširjenju i radikalu ostajut suštvene.\par
\noindent\emph{Textual comparison:} The extension and radical hypotheses remain essential.\par
\noindent{\footnotesize TEXT-COMP-954; pasaže / passages: 21930, 21932, 21934, 21952, 21954, 21964, 21968; \texttt{CORPUS\_FORM\_COMPARISONS.json}.}\par
\end{minipage}\par
\par\medskip\Needspace{7\baselineskip}
\noindent\begin{minipage}{\linewidth}
\hypertarget{TERM-1898}{}
\noindent\textbf{centralny element}\quad{\footnotesize\texttt{TERM-1898}}\par
\noindent\emph{Iztočnik / Source:} Zentrumselement\par
\noindent\emph{Rusky / Russian:} центральный элемент\par
\noindent\emph{Ukrajinsky / Ukrainian:} центральний елемент\par
\smallskip\noindent\emph{Razlog:} Centralny element jest element centra, to jest element, komutujuči so vsimi elementami prstena. Tu trěba popraviti nejasnu arhivnu motivaciju: centralnost i naleženje k centru ne sut različne vlastnosti. Razlika jest medžu imenovanjem jednogo elementa i imenovanjem cěloj podstruktury. V dokazu o direktnoj sumě centralnost komponentov vyhodi iz ravnosti dvuh produktov, direktnosti sumy i dvustrannosti idealov. Pridavnik centralny čuva tu vlastnost bez vvedenja novogo pojma.\par
\noindent\emph{Reason:} A centralny element is an element of the center: one commuting with every element of the ring. An unclear archived rationale needs correction here: centrality and membership in the center are not distinct properties. The distinction is between naming one element and naming the entire substructure. In the direct-sum proof, centrality of the components follows from equality of two products, directness of the sum and the two-sided ideal condition. The adjective centralny preserves that property without introducing a new concept.\par
\noindent{\footnotesize Izvorne zapisy / Original records: HIST-g189-002.}\par
\smallskip\noindent\emph{Sravnjenje s tekstom:} To sut ravnocěnne vlastnosti elementa, ne dva pojma.\par
\noindent\emph{Textual comparison:} These are equivalent properties of an element, not two concepts.\par
\noindent{\footnotesize TEXT-COMP-955; pasaže / passages: 20724, 20747, 20749, 20751, 20753, 20755, 20757, 20759; \texttt{CORPUS\_FORM\_COMPARISONS.json}.}\par
\end{minipage}\par
\par\medskip\Needspace{7\baselineskip}
\noindent\begin{minipage}{\linewidth}
\hypertarget{TERM-1899}{}
\noindent\textbf{centralny element / centr}\quad{\footnotesize\texttt{TERM-1899}}\par
\noindent\emph{Iztočnik / Source:} Zentrumselement / Zentrum\par
\noindent\emph{Rusky / Russian:} центральный элемент / центр\par
\noindent\emph{Ukrajinsky / Ukrainian:} центральний елемент / центр\par
\smallskip\noindent\emph{Razlog:} Centralny element i centr prodolžajut tu samu algebraičnu terminologiju; centar iz ranějših sekcij ne označuje drugu strukturu. V grupovom prstenu klas $K_a$ jest suma različnih konjugovanyh elementov, ne samo množina tih elementov. Te sume komutujut s grupovymi elementami i generujut centr, zato što koeficienty centralnogo elementa sut konstantne na konjugacijnoj klasi. Nazva centralny mora sohraniti tu vezu s množenjem. Stabilny prevod ne zaměnjuje klasnu sumu centralnym elementom samej grupy.\par
\noindent\emph{Reason:} Centralny element and centr continue the same algebraic terminology; centar in earlier sections does not denote a different structure. In the group ring, the class $K_a$ is the sum of the distinct conjugate elements, not merely their set. These sums commute with the group elements and generate the center, because a central element's coefficients are constant on each conjugacy class. Centralny must retain that multiplicative relation. Consistent translation must not turn a class sum into a central element of the group itself.\par
\noindent{\footnotesize Izvorne zapisy / Original records: HIST-g204-008.}\par
\smallskip\noindent\emph{Sravnjenje s tekstom:} Klasna suma ne jest centralny grupovy element.\par
\noindent\emph{Textual comparison:} A class sum is not a central group element.\par
\noindent{\footnotesize TEXT-COMP-956; pasaže / passages: 22475, 22477, 22479, 22481, 22482, 22484, 22486; \texttt{CORPUS\_FORM\_COMPARISONS.json}.}\par
\end{minipage}\par
\par\medskip\Needspace{7\baselineskip}
\noindent\begin{minipage}{\linewidth}
\hypertarget{TERM-1900}{}
\noindent\textbf{razpadno polje}\quad{\footnotesize\texttt{TERM-1900}}\par
\noindent\emph{Iztočnik / Source:} Zerf?llungsk?rper [Čestično nezachovane bukvy / Partly unrecoverable lettering]\par
\noindent\emph{Rusky / Russian:} [Bukvy ne sut zachovane / Lettering not recoverable]\par
\noindent\emph{Ukrajinsky / Ukrainian:} [Bukvy ne sut zachovane / Lettering not recoverable]\par
\smallskip\noindent\emph{Razlog:} Razpadno polje prodolžaje zabilěženu konvenciju rabot o predstavjenjah. Rozkladno polje ostalo jest alternativnym predlogom, svjazanym s češsko-polskym sravnjenjem, ale zapis ne dava točnyh citacij; to ne jest nově potvrđena attestacija. Pasaža razlikuje odcěpjenje jednogo linearnogo faktora od polnogo razpada, a pri algebrah opira razpad na prěhod k klasu polnyh matričnyh prstenov. Zato termin treba čitati s konkretnym objektom. Poškodženy německy zapis i nečitlive ruske i ukrajinske formy ostajut označene kako nedostatočno zachovane.\par
\noindent\emph{Reason:} Razpadno polje continues the recorded convention of the representation-theory papers. Rozkladno polje remained an alternative proposal associated with Czech–Polish comparison, but the record provides no exact citations; this is not a newly verified attestation. The passage distinguishes splitting off one linear factor from complete splitting, while for algebras it defines splitting by passage to the class of full matrix rings. The term must therefore be read with its particular object. The damaged German record and unreadable Russian and Ukrainian forms remain identified as insufficiently preserved.\par
\noindent{\footnotesize Izvorne zapisy / Original records: HIST-g209-009.}\par
\smallskip\noindent\emph{Sravnjenje s tekstom:} Nazva razpada potrěbuje ukazany matematičny objekt.\par
\noindent\emph{Textual comparison:} The name of splitting needs its specified mathematical object.\par
\noindent{\footnotesize TEXT-COMP-957; pasaže / passages: 23375, 23381, 23385, 23387, 23924, 23927, 23929; \texttt{CORPUS\_FORM\_COMPARISONS.json}.}\par
\end{minipage}\par
\par\medskip\Needspace{7\baselineskip}
\noindent\begin{minipage}{\linewidth}
\hypertarget{TERM-1901}{}
\noindent\textbf{razpadno polje; odc?pno polje [Čestično nezachovane bukvy / Partly unrecoverable lettering]}\quad{\footnotesize\texttt{TERM-1901}}\par
\noindent\emph{Iztočnik / Source:} Zerf?llungsk?rper; Separationsk?rper [Čestično nezachovane bukvy / Partly unrecoverable lettering]\par
\noindent\emph{Rusky / Russian:} [Bukvy ne sut zachovane / Lettering not recoverable]\par
\noindent\emph{Ukrajinsky / Ukrainian:} [Bukvy ne sut zachovane / Lettering not recoverable]\par
\smallskip\noindent\emph{Razlog:} Razpadno polje i odcěpno polje naměrno ne sut prevedene jednoj nazvoj. V obšćej definiciji prvo daje absolutno proste vse kompozicione faktory, a drugo prinajmenje jeden. To jest razlika medžu vse i prinajmenje jedno, ne samo stilističny izbor glagola. Pri osoblyh prostyh algebrah oba pojma mogut sběgati se; dokaz to izrično objašnjava. Slovnik mora najprvo dati razne definicije i potom taj specialny slučaj. Bukvy izgubljene v arhivu ostajut označene; čitljivo odcěpno vzeto jest iz aktualnogo teksta.\par
\noindent\emph{Reason:} Razpadno polje and odcěpno polje were deliberately not given the same name. In the general definition the former makes all composition factors absolutely simple, while the latter gives at least one. This is a difference between all and at least one, not merely a stylistic choice of verb. For particular simple algebras the notions can coincide; the proof explains this explicitly. The dictionary must first give the distinct definitions and then the special case. Letters lost in the archive remain marked; readable odcěpno comes from the current text.\par
\noindent{\footnotesize Izvorne zapisy / Original records: HIST-g210-006.}\par
\smallskip\noindent\emph{Sravnjenje s tekstom:} Dve nazvy čuvajut razne kvantifikacije.\par
\noindent\emph{Textual comparison:} The two names preserve different quantifiers.\par
\noindent{\footnotesize TEXT-COMP-958; pasaže / passages: 24118, 24120, 24122, 24142, 24144; \texttt{CORPUS\_FORM\_COMPARISONS.json}.}\par
\end{minipage}\par
\par\medskip\Needspace{7\baselineskip}
\noindent\begin{minipage}{\linewidth}
\hypertarget{TERM-1902}{}
\noindent\textbf{razkladati se / razkladna forma}\quad{\footnotesize\texttt{TERM-1902}}\par
\noindent\emph{Iztočnik / Source:} zerfallen / zerfallende Form\par
\noindent\emph{Rusky / Russian:} разлагаться / разложимая форма\par
\noindent\emph{Ukrajinsky / Ukrainian:} розкладатися / розкладна форма\par
\smallskip\noindent\emph{Razlog:} Razkladati se i razkladna forma svjazujut proces s pridavnikom dlja njegovogo rezultata ili vozmožnosti. Arhiv bira tu porodinu radi kontinuiteta prědhodnyh sekcij, no ostavlja otvoreno pytanje najboljšego jezyčnogo registra. V §23 razklad vyhodi iz konkretnoj jednokratnoj kontrakcije; ostale formy mają ili reducent kako faktor, ili odgovarjajuču kontrakciju s razkladnoju formoju. To ne jest neograničeno tvrđenje, že vsaka kontrakcija polinoma razkladaje se. Značenje pridavnika određuje ranějša osobna definicija.\par
\noindent\emph{Reason:} Razkladati se and razkladna forma connect a process with an adjective for its result or possibility. The record chooses this family for continuity with preceding sections, while leaving the best linguistic register open. In section 23 decomposition follows from a specified single contraction; the remaining forms either have a reducent as a factor or the relevant contraction with a decomposable form. This is not an unrestricted claim that every contraction of a polynomial decomposes. The earlier special definition determines the adjective's meaning.\par
\noindent{\footnotesize Izvorne zapisy / Original records: HIST-g024-005.}\par
\smallskip\noindent\emph{Sravnjenje s tekstom:} Lokalna definicija ostaje potrěbna.\par
\noindent\emph{Textual comparison:} The local definition remains necessary.\par
\noindent{\footnotesize TEXT-COMP-959; pasaže / passages: 3068, 3070, 3071, 3072, 3074, 579; \texttt{CORPUS\_FORM\_COMPARISONS.json}.}\par
\end{minipage}\par
\par\medskip\Needspace{7\baselineskip}
\noindent\begin{minipage}{\linewidth}
\hypertarget{TERM-1903}{}
\noindent\textbf{razkladajut se / razkladne formy}\quad{\footnotesize\texttt{TERM-1903}}\par
\noindent\emph{Iztočnik / Source:} zerfallen / zerfallende Formen\par
\noindent\emph{Rusky / Russian:} разлагаются / разложимые формы\par
\noindent\emph{Ukrajinsky / Ukrainian:} розкладаються / розкладні форми\par
\smallskip\noindent\emph{Razlog:} Razkladajut se i razkladne formy čuvajut množinu i razliku od reducibilnosti. V §26 formuly davajut konkretne produktne izrazy s koeficientami, a pozdějša fraza odkazuje na ranějši razklad. Izbor kratkogo glagola ne dovoljuje izpustiti te koeficienty ni uslovja prědhodnyh formul. Razkladanje v toj teoriji pritom obhvača i osobno obobčenje črez vyše formy, oprěděljeno v §3. Slovnik zato ne svodi vse pojavjenja na obyčnu faktorizaciju jednogo polinoma.\par
\noindent\emph{Reason:} Razkladajut se and razkladne formy preserve the plural and the distinction from reducibility. In section 26 the formulas give specific product expressions with coefficients, followed by a reference to an earlier decomposition. Choosing a short verb does not permit omission of those coefficients or the preceding formulas' conditions. Decomposition in this theory also includes the special generalization through higher forms defined in section 3. The dictionary therefore does not reduce every occurrence to ordinary factorization of one polynomial.\par
\noindent{\footnotesize Izvorne zapisy / Original records: HIST-g027-005.}\par
\smallskip\noindent\emph{Sravnjenje s tekstom:} Glagol ne zaměnjaje koeficienty i uslovja formul.\par
\noindent\emph{Textual comparison:} The verb does not replace the formulas' coefficients and conditions.\par
\noindent{\footnotesize TEXT-COMP-960; pasaže / passages: 3383, 3385, 3387, 3389, 3391, 579; \texttt{CORPUS\_FORM\_COMPARISONS.json}.}\par
\end{minipage}\par
\par\medskip\Needspace{7\baselineskip}
\noindent\begin{minipage}{\linewidth}
\hypertarget{TERM-1904}{}
\noindent\textbf{razkladanje ili reducibilnost ostalyh form}\quad{\footnotesize\texttt{TERM-1904}}\par
\noindent\emph{Iztočnik / Source:} Zerfallen oder Reduzibilität der übrigen Formen\par
\noindent\emph{Rusky / Russian:} разложение или приводимость остальных форм\par
\noindent\emph{Ukrajinsky / Ukrainian:} розкладання або звідність решти форм\par
\smallskip\noindent\emph{Razlog:} Razkladanje ili reducibilnost sohranja dve možnosti, ktore avtorica razdělila v definicijah. Razložimost dopušča produkty nižšego stepena i vyše formy; reducibilnost ima svoju definiciju črez invariantne faktory i vyše formy. Zaključna rečenica §16 upotrěbja oba nazvanja i odkazuje na ranějši dokaz. Arhiv razgledava razložimost kako alternativu internacionalizmu, ale ih prosta zaměna zamaglila by tu lokalnu razliku. Zato rozhoduje definicija v rabotě, ne samo etimologična blizkost.\par
\noindent\emph{Reason:} Razkladanje ili reducibilnost preserves two possibilities distinguished by the author's definitions. Decomposability permits products of lower degree and higher forms; reducibility has its own definition through invariant factors and higher forms. The closing sentence of section 16 uses both names and refers to the preceding argument. The archive considers razložimost as an alternative to the internationalism, but a simple substitution would obscure this local distinction. The definitions in the paper, not merely etymological resemblance, govern the choice.\par
\noindent{\footnotesize Izvorne zapisy / Original records: HIST-g017-008.}\par
\smallskip\noindent\emph{Sravnjenje s tekstom:} Avtorica upotrěbja dve osobne definicije.\par
\noindent\emph{Textual comparison:} The author uses two specific definitions.\par
\noindent{\footnotesize TEXT-COMP-961; pasaže / passages: 547, 579, 2355; \texttt{CORPUS\_FORM\_COMPARISONS.json}.}\par
\end{minipage}\par
\par\medskip\Needspace{7\baselineskip}
\noindent\begin{minipage}{\linewidth}
\hypertarget{TERM-1905}{}
\noindent\textbf{razkladanje transvekcij s funkcionalnymi determinantami}\quad{\footnotesize\texttt{TERM-1905}}\par
\noindent\emph{Iztočnik / Source:} Zerfallen von Überschiebungen mit Funktionaldeterminanten\par
\noindent\emph{Rusky / Russian:} разложение трансвекций с функциональными детерминантами\par
\noindent\emph{Ukrajinsky / Ukrainian:} розкладання трансвекцій із функціональними детермінантами\par
\smallskip\noindent\emph{Razlog:} Razkladanje transvekcij s funkcionalnymi determinantami povezuje tri už vvedene pojmy, ne ustanovljaje novu operaciju. Zapis prodolžaje transvekcije iz prědhodnoj sekcije i ostavlja termin otevreny dlja jezyčnoj ocěny. Aktualna rečenica kaže, že određene samo-razumljive slučaje ne budut posebno imenovane. Podrobno pravilo v §3 i primeri daju smysl toj kratkosti. To ne jest dozvoljenje prevoditelju da vypusti formule ili vse složne slučaje; skraćenje pripada samomu iztočnomu izloženju.\par
\noindent\emph{Reason:} Razkladanje transvekcij s funkcionalnymi determinantami connects three already introduced concepts rather than establishing a new operation. The record continues transvekcije from the preceding section and leaves the term open to linguistic assessment. The current sentence says that certain straightforward cases will not be listed separately. The detailed rule in section 3 and the examples explain that brevity. It does not authorize the translator to omit formulas or all difficult cases; the abbreviation belongs to the source's exposition.\par
\noindent{\footnotesize Izvorne zapisy / Original records: HIST-g021-005.}\par
\smallskip\noindent\emph{Sravnjenje s tekstom:} Iztočno skraćenje ne jest dozvoljenje novogo izpuščenja.\par
\noindent\emph{Textual comparison:} The source's abbreviation does not authorize a new omission.\par
\noindent{\footnotesize TEXT-COMP-962; pasaže / passages: 2722, 595, 2837, 2856; \texttt{CORPUS\_FORM\_COMPARISONS.json}.}\par
\end{minipage}\par
\par\medskip\Needspace{7\baselineskip}
\noindent\begin{minipage}{\linewidth}
\hypertarget{TERM-1906}{}
\noindent\textbf{razlo?ena algebra [Čestično nezachovane bukvy / Partly unrecoverable lettering]}\quad{\footnotesize\texttt{TERM-1906}}\par
\noindent\emph{Iztočnik / Source:} zerfallende Algebra\par
\noindent\emph{Rusky / Russian:} [Bukvy ne sut zachovane / Lettering not recoverable]\par
\noindent\emph{Ukrajinsky / Ukrainian:} [Bukvy ne sut zachovane / Lettering not recoverable]\par
\smallskip\noindent\emph{Razlog:} Čitljiva aktualna forma razložena algebra svjazana jest s arhivnym razlo?ena algebra, čije izgubljene bukvy ne možno predstavljati kako zachovany tekst. Razlog izbora razlikuje algebru od razpadnogo polja: prvo jest objekt, ktory razpadaje se, drugo razširjenje, ktore taj razpad omogožuje. V pasaži razložena algebra jest polny matričny prsten nad centrom. To ne znači samo reducibilnu ili poluprostu algebru. Slovo potrebuje tu strukturnu glosu, ne samu analogiju s razloženjem polinoma.\par
\noindent\emph{Reason:} The readable current form razložena algebra is linked to archival razlo?ena algebra, whose lost letters must not be presented as preserved text. The rationale distinguishes the algebra from a razpadno polje: the former is the object that splits, the latter the extension enabling it. In the passage a split algebra is a full matrix ring over the center. This does not mean merely a reducible or semisimple algebra. The word needs that structural gloss, not just an analogy with polynomial factorization.\par
\noindent{\footnotesize Izvorne zapisy / Original records: HIST-g209-010.}\par
\smallskip\noindent\emph{Sravnjenje s tekstom:} Razloženost algebry ne jest prosto poluprostost.\par
\noindent\emph{Textual comparison:} Splitting of the algebra is not merely semisimplicity.\par
\noindent{\footnotesize TEXT-COMP-963; pasaže / passages: 23265, 23271; \texttt{CORPUS\_FORM\_COMPARISONS.json}.}\par
\end{minipage}\par
\par\medskip\Needspace{7\baselineskip}
\noindent\begin{minipage}{\linewidth}
\hypertarget{TERM-1907}{}
\noindent\textbf{razlo?ena algebra; povsudu razlo?ena [Čestično nezachovane bukvy / Partly unrecoverable lettering]}\quad{\footnotesize\texttt{TERM-1907}}\par
\noindent\emph{Iztočnik / Source:} zerfallende Algebra; ?berall zerfallend [Čestično nezachovane bukvy / Partly unrecoverable lettering]\par
\noindent\emph{Rusky / Russian:} [Bukvy ne sut zachovane / Lettering not recoverable]\par
\noindent\emph{Ukrajinsky / Ukrainian:} [Bukvy ne sut zachovane / Lettering not recoverable]\par
\smallskip\noindent\emph{Razlog:} Razložena algebra i povsudu razložena prodolžajut istu pridavnikovu porodinu. Povsudnost odnosi se k konečnym i beskonečnym městam čislovogo polja, ne k proizvoljnym točkam prostora. V teoremu o glavnom rodu prvo se prověrjajut razvětvjene města i uslovje na faktorny sistem; zaključok potom govori o vseh městah i globalnom razpadu. Slovnik mora razložiti tu kvantifikaciju, a ne prosto ponoviti slovo povsudu. Arhivne formy s izgubljenymi bukvami ostajut označene oddělno od čitlivogo aktualnogo teksta.\par
\noindent\emph{Reason:} Razložena algebra and povsudu razložena continue the same adjectival family. Everywhere refers to finite and infinite places of the number field, not arbitrary points in space. In the principal-genus theorem the ramified places and a condition on the factor system are checked first; the conclusion then concerns every place and global splitting. The dictionary must explain that quantification rather than merely repeat everywhere. Archived forms with lost letters remain distinguished from the readable current text.\par
\noindent{\footnotesize Izvorne zapisy / Original records: HIST-g211-004.}\par
\smallskip\noindent\emph{Sravnjenje s tekstom:} Povsudu znači vse relevantne města polja.\par
\noindent\emph{Textual comparison:} Everywhere means all relevant places of the field.\par
\noindent{\footnotesize TEXT-COMP-964; pasaže / passages: 24361, 24363, 24365, 24367, 24369, 24371, 24373, 23294; \texttt{CORPUS\_FORM\_COMPARISONS.json}.}\par
\end{minipage}\par
\par\medskip\Needspace{7\baselineskip}
\noindent\begin{minipage}{\linewidth}
\hypertarget{TERM-1908}{}
\noindent\textbf{razložima/razložime formy}\quad{\footnotesize\texttt{TERM-1908}}\par
\noindent\emph{Iztočnik / Source:} zerfallende Form(en)\par
\noindent\emph{Rusky / Russian:} разлагающаяся/разлагающиеся формы; распадающаяся/распадающиеся формы\par
\noindent\emph{Ukrajinsky / Ukrainian:} форма/форми, що розкладаються\par
\smallskip\noindent\emph{Razlog:} Razložima i razložime sut jedinstvena i množestvena forma odnogo pridavnika. Arhiv dava prednost kratkoj medžuslovjanskoj formě, togda kako sravniteljne ruske i ukrajinske izrazy često razkladyvajut značenje v pričastnu frazu. Kontekst §§12–13 kaže, kako razklad slěduje iz jednokratnogo prěnosa nad ranějše razložimoju formoju i iz funkcionalnyh determinantov. To sohranjaje svjaz s konkretnoju konstrukcijeju. Arhivna ocěna prozračnosti ostaje kvalitativna, ne izměreny čitateljsky rezultat.\par
\noindent\emph{Reason:} Razložima and razložime are singular and plural forms of one adjective. The record prefers the short Interslavic form, whereas the Russian and Ukrainian comparanda often spell out the meaning in a participial phrase. The context of sections 12–13 explains decomposition through a single transvection over an already decomposable form and through functional determinants. This retains the connection to the specific construction. The archived judgment of transparency is qualitative, not a measured reader result.\par
\noindent{\footnotesize Izvorne zapisy / Original records: HIST-g013-004, HIST-g014-007.}\par
\smallskip\noindent\emph{Sravnjenje s tekstom:} Jedinstvena i množestvena forma ne sut razne pojmy.\par
\noindent\emph{Textual comparison:} Singular and plural forms are not different concepts.\par
\noindent{\footnotesize TEXT-COMP-965; pasaže / passages: 1518, 1868, 1878, 1880, 579; \texttt{CORPUS\_FORM\_COMPARISONS.json}.}\par
\end{minipage}\par
\par\medskip\Needspace{7\baselineskip}
\noindent\begin{minipage}{\linewidth}
\hypertarget{TERM-1909}{}
\noindent\textbf{formy, koje se razkladųt na faktory}\quad{\footnotesize\texttt{TERM-1909}}\par
\noindent\emph{Iztočnik / Source:} zerfallende Formen\par
\noindent\emph{Rusky / Russian:} формы, разлагающиеся на множители\par
\noindent\emph{Ukrajinsky / Ukrainian:} форми, що розкладаються на множники\par
\smallskip\noindent\emph{Razlog:} Formy, koje se razkladųt na faktory jest rany opisny izbor; aktualna fraza piše formy, ktore razkladajut se na faktory. V arhivu jest pojem i ocěna prigodnosti dlja pilota, ale ne detaljny zapis sravnjenja alternativ. Możno zato objasniti funkcionovanje formy, ne izmisliti nezachovanu debatu: fraza izrično govori o faktorah i ne potrěbuje od čitatelja znanja novogo pridavnika. Kratko razložime iz pozdějše raboty mora byti povezano s jej osobno širšoju definicijeju, ne bezuslovno proglašeno istoznačnym.\par
\noindent\emph{Reason:} Formy, koje se razkladųt na faktory is an early descriptive choice; the current phrase reads formy, ktore razkladajut se na faktory. The archive preserves the concept and an assessment of suitability for the pilot, but no detailed comparison of alternatives. Its function can therefore be explained without inventing a lost debate: the phrase explicitly names factors and requires no knowledge of a new adjective. The shorter razložime in the later paper must be linked to that paper's specially broadened definition, not declared unconditionally synonymous.\par
\noindent{\footnotesize Izvorne zapisy / Original records: HIST-g001-011.}\par
\smallskip\noindent\emph{Sravnjenje s tekstom:} Nedostajajuči razlog ne može byti izmislen.\par
\noindent\emph{Textual comparison:} An absent rationale cannot be invented.\par
\noindent{\footnotesize TEXT-COMP-966; pasaže / passages: 132, 579; \texttt{CORPUS\_FORM\_COMPARISONS.json}.}\par
\end{minipage}\par
\par\medskip\Needspace{7\baselineskip}
\noindent\begin{minipage}{\linewidth}
\hypertarget{TERM-1910}{}
\noindent\textbf{razkladne formy}\quad{\footnotesize\texttt{TERM-1910}}\par
\noindent\emph{Iztočnik / Source:} zerfallende Formen\par
\noindent\emph{Rusky / Russian:} разложимые формы\par
\noindent\emph{Ukrajinsky / Ukrainian:} розкладні форми\par
\smallskip\noindent\emph{Razlog:} Razkladne formy jest kraća pridavnikova varianta porodiny razklad-. Arhiv ją čuva radi soglasnosti, ale ne proglašava njen idiom konačno prověrenym. V §25 avtorica razlikuje neposrědnje kontrakcije, ktore več ne nabraja, od razklada posrědstvom dvojnoj redukcije određene forme. Taka struktura dokaza daje pridavniku smisao. Njegova kratkost ne zaměnjuje odkazku na §13 i ne dopušča prevoditi redukciju i razklad jednoj neodrěđenej operacijej.\par
\noindent\emph{Reason:} Razkladne formy is the shorter adjectival variant in the razklad- family. The record retains it for consistency without declaring its idiom finally established. In section 25 the author distinguishes direct contractions, no longer enumerated, from decomposition by double reduction of a specified form. That proof structure supplies the adjective's sense. Its brevity does not replace the reference to section 13 or permit translation of reduction and decomposition as one unspecified operation.\par
\noindent{\footnotesize Izvorne zapisy / Original records: HIST-g026-004.}\par
\smallskip\noindent\emph{Sravnjenje s tekstom:} Kratka forma ostaje vezana k točnomu dokazu.\par
\noindent\emph{Textual comparison:} The short form remains tied to the exact argument.\par
\noindent{\footnotesize TEXT-COMP-967; pasaže / passages: 3258, 3260, 3262, 3264, 3265, 3266; \texttt{CORPUS\_FORM\_COMPARISONS.json}.}\par
\end{minipage}\par
\par\medskip\Needspace{7\baselineskip}
\noindent\begin{minipage}{\linewidth}
\hypertarget{TERM-1911}{}
\noindent\textbf{razložime formy}\quad{\footnotesize\texttt{TERM-1911}}\par
\noindent\emph{Iztočnik / Source:} zerfallende Formen\par
\noindent\emph{Rusky / Russian:} разлагающиеся формы; распадающиеся на множители формы; распадающиеся формы\par
\noindent\emph{Ukrajinsky / Ukrainian:} форми, що розкладаються; форми, що розкладаються на множники\par
\smallskip\noindent\emph{Razlog:} Razložime formy bylo izbrano vměsto dolge pričastnoj frazy radi čitlivosti. Sravniteljne zapisy sohranjajut i ruske varianty s razlaganjem i raspadanjem, ne jednu tobože jedinu normu. Najvažnějša ograda jest matematična: §3 izrično obobčaje običny pojem, dopuščajuči produkty nižšego stepena zajedno s vyšimi formami. Pozdějše priměry naslědujut tu definiciju. Zato zabilěžene pytanje o jezyčnom registru ne može byti rěšeno prostym zaměnjenjem slova faktorizujema bez prověrky smisla.\par
\noindent\emph{Reason:} Razložime formy was preferred to a longer participial phrase for readability. The comparative records preserve Russian variants involving both razlaganie and raspadanie, not a supposed single norm. The most important qualification is mathematical: section 3 explicitly generalizes the ordinary concept by permitting products of lower degree together with higher forms. Later examples inherit that definition. The recorded question of linguistic register therefore cannot be resolved simply by substituting factorable without checking the meaning.\par
\noindent{\footnotesize Izvorne zapisy / Original records: HIST-g005-006, HIST-g011-006, HIST-g012-009, HIST-g015-007, HIST-g016-008.}\par
\smallskip\noindent\emph{Sravnjenje s tekstom:} Obyčna faktorizacija ne pokryvaje celu definiciju.\par
\noindent\emph{Textual comparison:} Ordinary factorization does not cover the whole definition.\par
\noindent{\footnotesize TEXT-COMP-968; pasaže / passages: 579, 583, 585, 587, 589, 591, 593, 1324, 1796, 1878; \texttt{CORPUS\_FORM\_COMPARISONS.json}.}\par
\end{minipage}\par
\par\medskip\Needspace{7\baselineskip}
\noindent\begin{minipage}{\linewidth}
\hypertarget{TERM-1912}{}
\noindent\textbf{razkladny ręd form}\quad{\footnotesize\texttt{TERM-1912}}\par
\noindent\emph{Iztočnik / Source:} zerfallende Formenreihe\par
\noindent\emph{Rusky / Russian:} разложимый ряд форм\par
\noindent\emph{Ukrajinsky / Ukrainian:} розкладний ряд форм\par
\smallskip\noindent\emph{Razlog:} Razkladny ręd form pripisuje razkladnost cělomu redu form, ne samo jednogo iz njegovyh členov. Aktualna grafija jest red. V §24 redukcija takogo reda vyhodi iz formule (30.)c i slědujuče točne relacije. Arhivny razlog izrično potrěbuje zachovati tu razliku v obsegu. Slovnik zato povezuje izraz s ranějšim pojmom reda form; on ne nudi interpretaciju kakoby šlo o konvergentny beskonečny red ili samo o proizvoljny spis polinomov.\par
\noindent\emph{Reason:} Razkladny ręd form attributes decomposability to the whole row of forms, not merely to one member. The current spelling is red. In section 24 reduction of such a row follows from formula (30.)c and the ensuing exact relation. The archived rationale expressly requires preservation of that scope distinction. The dictionary therefore links the phrase to the previously defined row of forms; it does not suggest a convergent infinite series or merely an arbitrary list of polynomials.\par
\noindent{\footnotesize Izvorne zapisy / Original records: HIST-g025-009.}\par
\smallskip\noindent\emph{Sravnjenje s tekstom:} Svojstvo ne odnosi se samo k jednomu členu.\par
\noindent\emph{Textual comparison:} The property does not refer only to one member.\par
\noindent{\footnotesize TEXT-COMP-969; pasaže / passages: 3192, 3194, 3197, 3198; \texttt{CORPUS\_FORM\_COMPARISONS.json}.}\par
\end{minipage}\par
\par\medskip\Needspace{7\baselineskip}
\noindent\begin{minipage}{\linewidth}
\hypertarget{TERM-1913}{}
\noindent\textbf{razložime rědy form ili formy}\quad{\footnotesize\texttt{TERM-1913}}\par
\noindent\emph{Iztočnik / Source:} zerfallende Formenreihen oder Formen\par
\noindent\emph{Rusky / Russian:} ряды форм или формы, распадающиеся на множители\par
\noindent\emph{Ukrajinsky / Ukrainian:} ряди форм чи форми, що розкладаються на множники\par
\smallskip\noindent\emph{Razlog:} Razložime rědy form ili formy širi obseg na cěly red i na pojedinu formu. Kratka arhivna motivacija govori o kontinuitetu s §3 i faktorah, ale treba sohraniti tamo obobčenu definiciju, ne samo običnu faktorizaciju. V planě konstrukcije te slučaje označujut města, kde red konči i kde dokaz ustanovljaje izčrpanost nereducibilnyh tvorjenij. Spojka ili zato nese informaciju o predmetu redukcije. Ona ne znači, že red i jedna forma sut istovětni objekty.\par
\noindent\emph{Reason:} Razložime rědy form ili formy covers both a whole row and an individual form. The brief archived rationale refers to continuity with section 3 and to factors, but its broadened definition must be retained, not merely ordinary factorization. In the construction's plan these cases mark where the row terminates and where the argument establishes exhaustion of the irreducible constructions. Or therefore carries information about the object of reduction. It does not identify a row with one form.\par
\noindent{\footnotesize Izvorne zapisy / Original records: HIST-g008-008.}\par
\smallskip\noindent\emph{Sravnjenje s tekstom:} Spojka ili čuva obseg predmeta.\par
\noindent\emph{Textual comparison:} Or preserves the scope of the object.\par
\noindent{\footnotesize TEXT-COMP-970; pasaže / passages: 1080, 1082, 1084, 579; \texttt{CORPUS\_FORM\_COMPARISONS.json}.}\par
\end{minipage}\par
\par\medskip\Needspace{7\baselineskip}
\noindent\begin{minipage}{\linewidth}
\hypertarget{TERM-1914}{}
\noindent\textbf{razložene skrěžene produkty}\quad{\footnotesize\texttt{TERM-1914}}\par
\noindent\emph{Iztočnik / Source:} zerfallende verschränkte Produkte\par
\noindent\emph{Rusky / Russian:} расщепляющиеся скрещенные произведения\par
\noindent\emph{Ukrajinsky / Ukrainian:} розщеплені схрещені добутки\par
\smallskip\noindent\emph{Razlog:} Razložene skrěžene produkty spojujut zabilěženy termin skrěženy produkt s uslovjem razloženosti algebry. Uvod izrično povezuje to uslovje s faktornymi sistemami asocijovanymi s jedinicoju; v §1 beret se predstavitelj s faktornym sistemom jedinic. To ne jest običny produkt dvoh polinomov. Arhiv naměrno razlikuje razloženy objekt od razpadnogo polja i ostavlja jezyčnu ocěnu skrěženy otvorenoju. Pozdějši negaloisovy slučaj ima svoje hipotezy, ktore kratka nazva ne ukida.\par
\noindent\emph{Reason:} Razložene skrěžene produkty combines the recorded crossed-product term with the condition that the algebra is split. The introduction expressly connects this condition with factor systems associated with unity; section 1 chooses a representative whose factor system consists of units equal to one. This is not an ordinary product of two polynomials. The record deliberately distinguishes the split object from the splitting field and leaves the linguistic assessment of skrěženy open. The later non-Galois case has its own hypotheses, which the short name does not remove.\par
\noindent{\footnotesize Izvorne zapisy / Original records: HIST-g212-001.}\par
\smallskip\noindent\emph{Sravnjenje s tekstom:} Asocijovanje s jedinicoju daje konkretny kriterij.\par
\noindent\emph{Textual comparison:} Association with unity supplies a concrete criterion.\par
\noindent{\footnotesize TEXT-COMP-971; pasaže / passages: 24449, 24453, 24455, 24457, 24515; \texttt{CORPUS\_FORM\_COMPARISONS.json}.}\par
\end{minipage}\par
\par\medskip\Needspace{7\baselineskip}
\noindent\begin{minipage}{\linewidth}
\hypertarget{TERM-1915}{}
\noindent\textbf{razpadne polja}\quad{\footnotesize\texttt{TERM-1915}}\par
\noindent\emph{Iztočnik / Source:} zerfällende Körper / Zerfällungskörper\par
\noindent\emph{Rusky / Russian:} поля разложения\par
\noindent\emph{Ukrajinsky / Ukrainian:} поля розкладу\par
\smallskip\noindent\emph{Razlog:} Razpadne polja v bibliografiji prodolžajut istu nazvu kako v rabotah o predstavjenjah. Zapis ne hoče vvesti novu variantu samo na koncu knigy; rozkladno polje ostaje sravniteljnym predlogom, ne provedenoju zaměnoju. Bibliografična stavka naziva minimalne razpadne polja i soavtora Brauera. Ona dolžna sohraniti te identifikacijske česti i voditi k konkretnomu dělu. Sama bibliografija ne jest dodatnym nezavisnym svědectvom jezyčnoj normy ili novym dokazom minimalnosti.\par
\noindent\emph{Reason:} Razpadne polja in the bibliography continues the name used in the representation-theory papers. The record avoids introducing a different variant only at the end of the book; rozkladno polje remains a comparative proposal, not an implemented substitution. The bibliography names minimal splitting fields and Brauer as coauthor. Those identifying details must remain and lead to the particular work. The bibliography is not an additional independent attestation of linguistic norm or a new proof of minimality.\par
\noindent{\footnotesize Izvorne zapisy / Original records: HIST-g215-009.}\par
\smallskip\noindent\emph{Sravnjenje s tekstom:} Nazva v bibliografiji ne jest novy terminologičny dokaz.\par
\noindent\emph{Textual comparison:} A bibliographic name is not new terminological evidence.\par
\noindent{\footnotesize TEXT-COMP-972; pasaže / passages: bib-isv.tex:126, bib-isv.tex:127, 19760; \texttt{CORPUS\_FORM\_COMPARISONS.json}.}\par
\end{minipage}\par
\par\medskip\Needspace{7\baselineskip}
\noindent\begin{minipage}{\linewidth}
\hypertarget{TERM-1916}{}
\noindent\textbf{razkladaje se / razkladanje}\quad{\footnotesize\texttt{TERM-1916}}\par
\noindent\emph{Iztočnik / Source:} zerfällt / Zerfallen\par
\noindent\emph{Rusky / Russian:} разлагается / разложение\par
\noindent\emph{Ukrajinsky / Ukrainian:} розкладається / розкладання\par
\smallskip\noindent\emph{Razlog:} Razkladaje se i razkladanje davajut glagol i imennik za istu operaciju. Konkretny primer v §19 upotrěbja kratšu glagolnu variantu razklada se, odkazuje na formulu (7.)a i iz togo vyvodi razklad slědujuče forme. Zato različna dulžina glagola ne ukazuje novu matematičnu razliku. Arhivny razlog jest kontinuitet s determinantnymi redukcijami; on ne dava pravo zameniti specifičnu formulu obćim slovom razklad. Prěvod mora sohraniti smer togo zaključka.\par
\noindent\emph{Reason:} Razkladaje se and razkladanje give the verb and noun for the same operation. The specific example in section 19 uses the shorter verbal variant razklada se, refers to formula (7.)a and deduces decomposition of the next form. The difference in verbal length therefore does not mark a new mathematical distinction. The recorded reason is continuity with determinant reductions; it does not authorize replacing the specific formula by the generic word decomposition. Translation must preserve the direction of the inference.\par
\noindent{\footnotesize Izvorne zapisy / Original records: HIST-g020-006.}\par
\smallskip\noindent\emph{Sravnjenje s tekstom:} Razklada se i razkladaje se ne tvorjat dva matematična pojma.\par
\noindent\emph{Textual comparison:} Razklada se and razkladaje se do not create two mathematical concepts.\par
\noindent{\footnotesize TEXT-COMP-973; pasaže / passages: 2672, 2625, 2627; \texttt{CORPUS\_FORM\_COMPARISONS.json}.}\par
\end{minipage}\par
\par\medskip\Needspace{7\baselineskip}
\noindent\begin{minipage}{\linewidth}
\hypertarget{TERM-1917}{}
\noindent\textbf{razklada se / razkladna forma}\quad{\footnotesize\texttt{TERM-1917}}\par
\noindent\emph{Iztočnik / Source:} zerfällt / zerfallende Form\par
\noindent\emph{Rusky / Russian:} разлагается / разложимая форма\par
\noindent\emph{Ukrajinsky / Ukrainian:} розкладається / розкладна форма\par
\smallskip\noindent\emph{Razlog:} Razklada se i razkladna forma prodolžajut kratku glagolnu i pridavnikovu tvorbu. Arhiv upominaje poljske, češke i južnoslovjanske obrazce, ale bez točnyh citacij, zato to ostaje zabilěžena motivacija, ne nově ustanovjena attestacija. V §22 razklad dobiva se dvojnoju redukcijeju i kontrakcijeju s funkcionalnym determinantom; zaključna čest kaže, ktora točno forma čuva razkladnost pri jedno- i dvojkratnyh kontrakcijah. Taj konkretni subjekt ne možno izpustiti i potom pripisati svojstvo vsakoj formě.\par
\noindent\emph{Reason:} Razklada se and razkladna forma continue the short verbal and adjectival formation. The record mentions Polish, Czech and South Slavic patterns without exact citations, so this remains a recorded motivation rather than a newly established attestation. In section 22 decomposition is obtained by double reduction and contraction with a functional determinant; the conclusion specifies exactly which form preserves decomposability under single and double contractions. That particular subject cannot be omitted and the property then attributed to every form.\par
\noindent{\footnotesize Izvorne zapisy / Original records: HIST-g023-007.}\par
\smallskip\noindent\emph{Sravnjenje s tekstom:} Kvantifikacija ostaje vezana k konkretnomu subjektu.\par
\noindent\emph{Textual comparison:} Quantification remains attached to its particular subject.\par
\noindent{\footnotesize TEXT-COMP-974; pasaže / passages: 2952, 2954, 2966, 2968, 2970, 2972, 2982, 3000; \texttt{CORPUS\_FORM\_COMPARISONS.json}.}\par
\end{minipage}\par
\par\medskip\Needspace{7\baselineskip}
\noindent\begin{minipage}{\linewidth}
\hypertarget{TERM-1918}{}
\noindent\textbf{polje razloženja}\quad{\footnotesize\texttt{TERM-1918}}\par
\noindent\emph{Iztočnik / Source:} Zerfällungskörper\par
\noindent\emph{Rusky / Russian:} поле разложения\par
\noindent\emph{Ukrajinsky / Ukrainian:} поле розщеплення\par
\smallskip\noindent\emph{Razlog:} Polje razloženja jest zabilěžena imennička alternativa k razpadno polje. Zapis sravnjava ukrajinske rozščepennja s ruskom razloženie i biraje tu drugu tvorbu dlja medžuslovjanskogo. V članku oba jezyčna dela mają strogy smisao: $K$ razklada algebru, kogda razširjena algebra naleži klasu polnyh matric nad $K$. Lokalny kriterij dodatno potrěbuje dělimost odgovarjajučih stepenjev indeksami. Zato slovnik sohranja variantu i kriterij, ne predstavlja pojavu različnogo nazvanja kako novo matematično otkrytje.\par
\noindent\emph{Reason:} Polje razloženja is the recorded nominal alternative to razpadno polje. The record compares Ukrainian rozščepennja with Russian razloženie and chooses the latter formation for Interslavic. In the article the phrase has a strict meaning: $K$ splits the algebra when the extended algebra belongs to the class of full matrix algebras over $K$. The local criterion additionally requires divisibility of the relevant degrees by the indices. The dictionary retains both the variant and the criterion, rather than presenting a different name as a mathematical discovery.\par
\noindent{\footnotesize Izvorne zapisy / Original records: HIST-g208-007.}\par
\smallskip\noindent\emph{Sravnjenje s tekstom:} Imennikova i pridavnikova nazva trebujut jednoj svezi.\par
\noindent\emph{Textual comparison:} The nominal and adjectival names need a cross-reference.\par
\noindent{\footnotesize TEXT-COMP-975; pasaže / passages: 23013, 23097, 23099, 23104, 23107; \texttt{CORPUS\_FORM\_COMPARISONS.json}.}\par
\end{minipage}\par
\par\medskip\Needspace{7\baselineskip}
\noindent\begin{minipage}{\linewidth}
\hypertarget{TERM-1919}{}
\noindent\textbf{razpadne polja}\quad{\footnotesize\texttt{TERM-1919}}\par
\noindent\emph{Iztočnik / Source:} Zerfällungskörper\par
\noindent\emph{Rusky / Russian:} поля разложения\par
\noindent\emph{Ukrajinsky / Ukrainian:} поля розщеплення\par
\smallskip\noindent\emph{Razlog:} Razpadne polja prodolžajut ranějši izbor, a rozkladno polje ostaje imenovanoju alternativoj bez tichoj zaměny sredi raboty. Uvod određuje predmet: komutativne razširjenja oblasti koeficientov, v ktoryh predstavjenja razpadajut se na absolutno ireducibilne. Hipoteza sistema bez radikala i svez s tělom avtomorfizmov sut česti toga izloženja. Prosta pluralna nazva ne mora skrivati ni te uslovja, ni naučnu svez s Šurovymi i Brauerovo-Noetherovymi rabotami.\par
\noindent\emph{Reason:} Razpadne polja continues the earlier choice, while rozkladno polje remains a named alternative without a silent mid-paper substitution. The introduction specifies the object: commutative coefficient-field extensions over which representations split into absolutely irreducible ones. The radical-free hypothesis and connection with the automorphism division ring belong to that account. The simple plural name must not conceal those conditions or the scholarly connection with Schur's work and the Brauer–Noether paper.\par
\noindent{\footnotesize Izvorne zapisy / Original records: HIST-g178-006.}\par
\smallskip\noindent\emph{Sravnjenje s tekstom:} Hipoteza odsutnosti radikala ne jest neobvezna.\par
\noindent\emph{Textual comparison:} The radical-free hypothesis is not optional.\par
\noindent{\footnotesize TEXT-COMP-976; pasaže / passages: 19760; \texttt{CORPUS\_FORM\_COMPARISONS.json}.}\par
\end{minipage}\par
\par\medskip\Needspace{7\baselineskip}
\noindent\begin{minipage}{\linewidth}
\hypertarget{TERM-1920}{}
\noindent\textbf{razpadno polje}\quad{\footnotesize\texttt{TERM-1920}}\par
\noindent\emph{Iztočnik / Source:} Zerfällungskörper\par
\noindent\emph{Rusky / Russian:} поле расщепления\par
\noindent\emph{Ukrajinsky / Ukrainian:} поле розщеплення\par
\smallskip\noindent\emph{Razlog:} Razpadno polje zabilěženo jest v trěh kontekstah kako nazva komutativnogo razširjenja, v otličji od razložene algebry. Ta doslědnost ne sliva vse definicije: v rabotě o Šurovom pytanju govorji se o razpadu predstavjenja, a v rabotě o skrěženyh produktah o algebrě i jej maksimalnom komutativnom podpolju. Minimalnost po uključanju i najnižši stepenj takože ostajut različne. Slovo razpad- izbrano jest radi prozračnoj slovjanskoj tvorby, no arhivna ograda o jezyčnom odobrenju ne jest svědectvo provedenoj ekspertnoj ocěny.\par
\noindent\emph{Reason:} Razpadno polje is recorded in three contexts for a commutative extension, as distinct from a razložena algebra. That consistency does not merge all definitions: the paper on Schur's question concerns splitting a representation, while the crossed-product paper concerns an algebra and its maximal commutative subfield. Inclusion-minimality and least degree also remain different. The razpad- formation was chosen for its transparent Slavic construction, but the archived reservation about linguistic approval is not evidence of completed expert review.\par
\noindent{\footnotesize Izvorne zapisy / Original records: HIST-g167-005, HIST-g168-001, HIST-g212-004.}\par
\smallskip\noindent\emph{Sravnjenje s tekstom:} Doslědnost nazvy ne odstraňuje razlike konteksta.\par
\noindent\emph{Textual comparison:} Consistency of the name does not remove contextual differences.\par
\noindent{\footnotesize TEXT-COMP-977; pasaže / passages: 19210, 19244, 24449, 24515; \texttt{CORPUS\_FORM\_COMPARISONS.json}.}\par
\end{minipage}\par
\par\medskip\Needspace{7\baselineskip}
\noindent\begin{minipage}{\linewidth}
\hypertarget{TERM-1922}{}
\noindent\textbf{razložima identičnost}\quad{\footnotesize\texttt{TERM-1922}}\par
\noindent\emph{Iztočnik / Source:} zerlegbare Identität\par
\noindent\emph{Rusky / Russian:} разложимое тождество\par
\noindent\emph{Ukrajinsky / Ukrainian:} розкладна тотожність\par
\smallskip\noindent\emph{Razlog:} Razložima identičnost označaje svojstvo identičnosti; razkladna identičnost označaje imenovany vid identičnosti, ktora razkladaje simbolny red. Razlika jest matematična, ne samo stilistična. Avtorka naziva razkladnu identičnost nerazložimoju, dokolě jej redy imajut nezavisno značenje; po odgovarajučoj zaměně ona stava razložimoju. Zato te dvě nazvy ne smějut byti mehanično zamenjene jedna drugoju. Ocěna zavisi od definicije i od toga, ktore simboly jesmo ostavili nezavisnymi.\par
\noindent\emph{Reason:} Razložima identičnost describes a property of an identity; razkladna identičnost names a family of identities that decompose a symbolic row. The difference is mathematical, not merely stylistic. The author calls a decomposition identity indecomposable while its rows retain independent meaning; after the relevant substitution it becomes decomposable. The two names therefore cannot be mechanically interchanged. The classification depends on the definition and on which symbols remain independent.\par
\noindent{\footnotesize Izvorne zapisy / Original records: HIST-g034-008.}\par
\smallskip\noindent\emph{Sravnjenje s tekstom:} Razlika zavisi od nezavisnosti simbolnyh redov.\par
\noindent\emph{Textual comparison:} The distinction depends on independence of symbolic rows.\par
\noindent{\footnotesize TEXT-COMP-979; pasaže / passages: 3897, 4161, 4278; \texttt{CORPUS\_FORM\_COMPARISONS.json}.}\par
\end{minipage}\par
\par\medskip\Needspace{7\baselineskip}
\noindent\begin{minipage}{\linewidth}
\hypertarget{TERM-1923}{}
\noindent\textbf{razklad v $k$ sostavnyh ?esti / $k$ podgrup [Čestično nezachovane bukvy / Partly unrecoverable lettering]}\quad{\footnotesize\texttt{TERM-1923}}\par
\noindent\emph{Iztočnik / Source:} Zerlegung in $k$ Bestandteile / $k$ Untergruppen\par
\noindent\emph{Rusky / Russian:} ?????????? ?? $k$ ???????????? / $k$ ???????? [Čestično nezachovane bukvy / Partly unrecoverable lettering]\par
\noindent\emph{Ukrajinsky / Ukrainian:} ??????? ?? $k$ ????????? / $k$ ??????? [Čestično nezachovane bukvy / Partly unrecoverable lettering]\par
\smallskip\noindent\emph{Razlog:} Razklad v k sostavnyh česti / k podgrup imenuje razklad na k aditivnyh komponentov. Čislo k ostaje parametrom: govorimo ne neobhodimo o dvuh komponentah. Formuly potrěbujut, že suma ravna nulu povlěče nuljevost vsakej česti, a teorema svazuje razklad grupy ostatkov s najmanjšim obščim kratnym k vzajemno prostyh modulov. Podgrupa zato ne označaje proizvoljny podmnožestvo. Prva fraza zde objasnjaje smysl čestično povrědženogo zapisa; originalne nezachovane bukvy i neprovedena jezyčna ocěna ne sut prikryte obnovjenoju tablicoju.\par
\noindent\emph{Reason:} The phrase denotes decomposition into k additive components or subgroups. The symbol k remains a parameter: the decomposition need not have only two components. The formulas require a zero sum to force each component to vanish, and the theorem relates decomposition of the residue group to the least common multiple of k mutually coprime modules. Subgroup therefore does not mean an arbitrary subset. The explanatory wording clarifies a partially damaged record; it does not conceal the lost lettering or turn an outstanding linguistic assessment into completed review.\par
\noindent{\footnotesize Izvorne zapisy / Original records: HIST-g079-028.}\par
\smallskip\noindent\emph{Sravnjenje s tekstom:} Razklad uključaje uslovje nezavisnosti komponentov.\par
\noindent\emph{Textual comparison:} The decomposition includes a component-independence condition.\par
\noindent{\footnotesize TEXT-COMP-980; pasaže / passages: 11263, 11265, 11267, 11269, 11272, 11360; \texttt{CORPUS\_FORM\_COMPARISONS.json}.}\par
\end{minipage}\par
\par\medskip\Needspace{7\baselineskip}
\noindent\begin{minipage}{\linewidth}
\hypertarget{TERM-1924}{}
\noindent\textbf{razloženje na dvustranne idealy}\quad{\footnotesize\texttt{TERM-1924}}\par
\noindent\emph{Iztočnik / Source:} Zerlegung in zweiseitige Ideale\par
\noindent\emph{Rusky / Russian:} разложение на двусторонние идеалы\par
\noindent\emph{Ukrajinsky / Ukrainian:} розклад на двобічні ідеали\par
\smallskip\noindent\emph{Razlog:} Razloženje na dvustranne idealy zachovuje neobhodimy pridevnik dvustranne. V danyh uslovjah razne idealne komponenty imajut nuljevy produkt, a vsaka komponenta ravna svojemu kvadratu. Obratny argument uživa te svojstva, da ustanovi idealnu strukturu. Prsten s jediniceju jest vyrazna hipoteza izloženja, ne domysl slovnika. Krača nazva razloženje na idealy by skryla važnu informaciju v nekomutativnom kontekstě; razloženje na levostranne idealy ne jest ravnoznačno.\par
\noindent\emph{Reason:} Razloženje na dvustranne idealy preserves the essential qualification two-sided. Under the stated conditions, distinct ideal components have zero product and each component equals its square. The converse argument uses these properties to establish the ideal structure. A ring with identity is an explicit hypothesis, not an assumption supplied by the dictionary. Shortening the phrase to decomposition into ideals would hide important information in a noncommutative context; decomposition into left ideals is not equivalent.\par
\noindent{\footnotesize Izvorne zapisy / Original records: HIST-g188-001.}\par
\smallskip\noindent\emph{Sravnjenje s tekstom:} Multiplikativne uslovja ne sut zaměnljive s obščim zbirkom podmnožestv.\par
\noindent\emph{Textual comparison:} The multiplicative conditions are not interchangeable with an arbitrary collection of subsets.\par
\noindent{\footnotesize TEXT-COMP-981; pasaže / passages: 20588, 20590, 20592, 20594, 20596, 20598, 20601, 20606, 20609; \texttt{CORPUS\_FORM\_COMPARISONS.json}.}\par
\end{minipage}\par
\par\medskip\Needspace{7\baselineskip}
\noindent\begin{minipage}{\linewidth}
\hypertarget{TERM-1925}{}
\noindent\textbf{grupa razlo?enja; razv?tvjene m?sta; nerazv?tvjeny [Čestično nezachovane bukvy / Partly unrecoverable lettering]}\quad{\footnotesize\texttt{TERM-1925}}\par
\noindent\emph{Iztočnik / Source:} Zerlegungsgruppe; Verzweigungsstellen; unverzweigt\par
\noindent\emph{Rusky / Russian:} [Bukvy ne sut zachovane / Lettering not recoverable]\par
\noindent\emph{Ukrajinsky / Ukrainian:} [Bukvy ne sut zachovane / Lettering not recoverable]\par
\smallskip\noindent\emph{Razlog:} Grupa razloženja, města razvětvjenja i nerazvětvjeny tvorjat svezanu aritmetičnu grupu ponjatij. Grupa razloženja prinadleži vybranomu prostomu městu v razširjenju; ne jest proizvoljna grupa, razložena na česti. Města vključajut konečne i nekonečne města. V dokazě nerazvětvjenost jest uslovje, ktore pozvaljaje užiti normy jedinic i izvesti lokalny razpad algebry. Zapis slov v tablicě jest povrědženy; objasnjenje smysla ne udaje, že izvorna ruska i ukrajinska pisma byli obnovjene.\par
\noindent\emph{Reason:} Decomposition group, ramified places and unramified form a connected arithmetic vocabulary. The decomposition group belongs to a chosen prime place in an extension; it is not an arbitrary group decomposed into parts. Places include finite and infinite places. In the proof, unramifiedness is the condition that permits the use of norms of units and leads to local splitting of the algebra. The tabulated spellings are damaged; explaining the concepts does not pretend to recover the original Russian and Ukrainian lettering.\par
\noindent{\footnotesize Izvorne zapisy / Original records: HIST-g211-010.}\par
\smallskip\noindent\emph{Sravnjenje s tekstom:} Nerazvětvjenost jest uslovje argumenta o normah.\par
\noindent\emph{Textual comparison:} Unramifiedness is a condition of the norm argument.\par
\noindent{\footnotesize TEXT-COMP-982; pasaže / passages: 24363, 24373, 24381, 24383, 24385, 24387; \texttt{CORPUS\_FORM\_COMPARISONS.json}.}\par
\end{minipage}\par
\par\medskip\Needspace{7\baselineskip}
\noindent\begin{minipage}{\linewidth}
\hypertarget{TERM-1926}{}
\noindent\textbf{razkladna identičnost}\quad{\footnotesize\texttt{TERM-1926}}\par
\noindent\emph{Iztočnik / Source:} Zerlegungsidentität\par
\noindent\emph{Rusky / Russian:} тождество разложения\par
\noindent\emph{Ukrajinsky / Ukrainian:} тотожність розкладання\par
\smallskip\noindent\emph{Razlog:} Razkladna identičnost jest nazva izbrana dlja identičnosti, v ktorej simbolny red jest razložen na posamečne redy. Uvod svazuje to pytanje s rabotami Clebscha i Studyja; te imena označajut naučny izvor pytanja, ne jezyčno odobrenje prevoda. Kratky pridevnik razkladna imenuje funkciju formuly, dokolě razložima bi tvrdila drugu vlastnost. Arhiv ostavlja tvorbu provisionalnoju; nyněšnji tekst dava osnovu dlja semantičnoj razlike, ne svědectvo opčeho prijetja slova.\par
\noindent\emph{Reason:} Razkladna identičnost names an identity in which a symbolic row is decomposed into individual rows. The introduction connects the question with Clebsch and Study; these names identify the mathematical background, not linguistic endorsement of the translation. The concise adjective razkladna names what the formula does, whereas razložima would assert a different property. The archive leaves the formation provisional. The present text supports the semantic distinction, not a claim of general acceptance.\par
\noindent{\footnotesize Izvorne zapisy / Original records: HIST-g028-015, HIST-g029-011.}\par
\smallskip\noindent\emph{Sravnjenje s tekstom:} Nazva govori, čto formula děla, ne tvrdi jej razložimost.\par
\noindent\emph{Textual comparison:} The name states what the formula does, not that it is decomposable.\par
\noindent{\footnotesize TEXT-COMP-983; pasaže / passages: 3673, 3705, 4076; \texttt{CORPUS\_FORM\_COMPARISONS.json}.}\par
\end{minipage}\par
\par\medskip\Needspace{7\baselineskip}
\noindent\begin{minipage}{\linewidth}
\hypertarget{TERM-1927}{}
\noindent\textbf{razkladne identičnosti}\quad{\footnotesize\texttt{TERM-1927}}\par
\noindent\emph{Iztočnik / Source:} Zerlegungsidentitäten\par
\noindent\emph{Rusky / Russian:} тождества разложения\par
\noindent\emph{Ukrajinsky / Ukrainian:} тотожності розкладання\par
\smallskip\noindent\emph{Razlog:} Plural razkladne identičnosti prodolžaje tu samu nazvu v naslově četvrtoj sekcije. Definicija govori točno, koji red jest razložen i čto ostaje nezavisno; ne ide o vsako razvinenje izraza ili o zbirku slučajnyh ravnostij. Zato plural ne stvara novy matematičny pojam. Svez s nerazložimostju samyh identičnostij ostaje vidima: nazva operacije i svojstvo rezultata mogut stati jedna do drugoj bez proturěčja.\par
\noindent\emph{Reason:} The plural razkladne identičnosti continues the same term in the heading of the fourth section. The definition specifies which row is decomposed and what remains independent; it is not a name for any expansion of an expression or an arbitrary collection of equalities. The plural creates no new mathematical concept. The connection with the indecomposability of the identities remains visible: the name of the operation and a property of its result can coexist without contradiction.\par
\noindent{\footnotesize Izvorne zapisy / Original records: HIST-g033-001.}\par
\smallskip\noindent\emph{Sravnjenje s tekstom:} Plural zachovuje tu samu definiciju.\par
\noindent\emph{Textual comparison:} The plural retains the same definition.\par
\noindent{\footnotesize TEXT-COMP-984; pasaže / passages: 4074, 4076, 4157, 4158, 4161; \texttt{CORPUS\_FORM\_COMPARISONS.json}.}\par
\end{minipage}\par
\par\medskip\Needspace{7\baselineskip}
\noindent\begin{minipage}{\linewidth}
\hypertarget{TERM-1928}{}
\noindent\textbf{razkladne identičnosti v eksplicitnoj formě}\quad{\footnotesize\texttt{TERM-1928}}\par
\noindent\emph{Iztočnik / Source:} Zerlegungsidentitäten in expliziter Form\par
\noindent\emph{Rusky / Russian:} тождества разложения в явной форме\par
\noindent\emph{Ukrajinsky / Ukrainian:} тотожності розкладання в явній формі\par
\smallskip\noindent\emph{Razlog:} V eksplicitnoj formě znači, že identičnost jest izražena črez dane formy i koeficienty, bez potrěby ostaviti posamečne simbolne redy kako nezavisne pomočne objekty. Avtorka dostaje takov izraz produktovoju teoremoju i zaměnoju ranějše definicije. Eksplicitny zato ne znači samo jasno napisany, ni numerično izračunany. Naslov mora zachovati tehničny smysl i svez s ranějšimi razkladnymi identičnostjami, vključno s uslovjami, pri ktoryh postajut razložimymi.\par
\noindent\emph{Reason:} In explicit form means that the identity is expressed through the given forms and coefficients without retaining the individual symbolic rows as independent auxiliary objects. The author obtains such an expression through a product theorem and substitution of an earlier definition. Explicit therefore does not merely mean clearly written, nor does it mean numerically evaluated. The heading must retain the technical sense and the connection with the earlier decomposition identities, including the conditions under which they become decomposable.\par
\noindent{\footnotesize Izvorne zapisy / Original records: HIST-g034-001.}\par
\smallskip\noindent\emph{Sravnjenje s tekstom:} Eksplicitna forma odstranjava zavisimost od pomočnih redov.\par
\noindent\emph{Textual comparison:} The explicit form removes dependence on auxiliary rows.\par
\noindent{\footnotesize TEXT-COMP-985; pasaže / passages: 4237, 4250, 4255, 4257, 4259, 4260, 4278, 4287; \texttt{CORPUS\_FORM\_COMPARISONS.json}.}\par
\end{minipage}\par
\par\medskip\Needspace{7\baselineskip}
\noindent\begin{minipage}{\linewidth}
\hypertarget{TERM-1929}{}
\noindent\textbf{teorema razloženja}\quad{\footnotesize\texttt{TERM-1929}}\par
\noindent\emph{Iztočnik / Source:} Zerlegungssatz\par
\noindent\emph{Rusky / Russian:} теорема разложения\par
\noindent\emph{Ukrajinsky / Ukrainian:} теорема розкладу\par
\smallskip\noindent\emph{Razlog:} Teorema razloženja zde prinadleži teoriji klas polj. Dlja nerazvětvjenogo prostogo města ona izražaje stepenj ostatkov črez najmanjši eksponent, ktory daje lokalno trivijalnu klasu dlja vsakoj algebry iz nazvanoj podgrupy. Kvantifikator vsakoj i uslovje nerazvětvjenosti sut neizostavne. Nazva ne označaje ni opču teoremu o jedinstvenom razkladu na proste faktory, ni svojstvo jednoj izbranoj algebry. Arhivny izbor daje čitlivu obšču nazvu; matematično ograničenje daje konkretna teorema.\par
\noindent\emph{Reason:} Here teorema razloženja belongs to class field theory. For an unramified prime place, it expresses the residue degree through the least exponent that gives a locally trivial class for every algebra in the specified subgroup. Every and unramified are indispensable. The name denotes neither a general unique-factorization theorem nor a property of one chosen algebra. The recorded choice supplies a readable general name; the particular theorem supplies its mathematical restriction.\par
\noindent{\footnotesize Izvorne zapisy / Original records: HIST-g208-016.}\par
\smallskip\noindent\emph{Sravnjenje s tekstom:} Najmanjši eksponent odnosi se na vse algebry nazvanoj podgrupy.\par
\noindent\emph{Textual comparison:} The least exponent concerns every algebra in the specified subgroup.\par
\noindent{\footnotesize TEXT-COMP-986; pasaže / passages: 23140, 23143, 23145, 23147, 23150; \texttt{CORPUS\_FORM\_COMPARISONS.json}.}\par
\end{minipage}\par
\par\medskip\Needspace{7\baselineskip}
\noindent\begin{minipage}{\linewidth}
\hypertarget{TERM-1930}{}
\noindent\textbf{razložne teoremy}\quad{\footnotesize\texttt{TERM-1930}}\par
\noindent\emph{Iztočnik / Source:} Zerlegungssätze\par
\noindent\emph{Rusky / Russian:} теоремы о разложении\par
\noindent\emph{Ukrajinsky / Ukrainian:} теореми про розклад\par
\smallskip\noindent\emph{Razlog:} Razložne teoremy služi kako skupna nazva dlja nekoliko vidov idealnogo razloženja. Uvod oddělja produkt idealov od najmanjšego obščego kratnogo idealov i razlikuje tvrđenja o jedinstvenosti. Primarny ideal ne mora byti stepenju svojego pridruženogo prostogo ideala. Zato obšča nazva ne tvrdi, že vse teoremy sut faktorizacije s istim pravilom jedinstvenosti. Vyslovjene domněnky v podnožnoj primětke ostajut domněnkami, a ne dokazanou čestju značenja termina.\par
\noindent\emph{Reason:} Razložne teoremy serves as a collective name for several kinds of ideal decomposition. The introduction distinguishes products of ideals from least common multiples and separates the different uniqueness assertions. A primary ideal need not be a power of its associated prime ideal. The umbrella term therefore does not assert that all these theorems are factorizations with the same uniqueness rule. Conjectures in the footnote remain conjectures, not proved content smuggled into the term.\par
\noindent{\footnotesize Izvorne zapisy / Original records: HIST-g081-003, HIST-g082-003.}\par
\smallskip\noindent\emph{Sravnjenje s tekstom:} Produkt, obšči kratny i jedinstvenost imajut razne uslovja.\par
\noindent\emph{Textual comparison:} Product, common multiple and uniqueness have different conditions.\par
\noindent{\footnotesize TEXT-COMP-987; pasaže / passages: 12410, 12434, 12436, 12438; \texttt{CORPUS\_FORM\_COMPARISONS.json}.}\par
\end{minipage}\par
\par\medskip\Needspace{7\baselineskip}
\noindent\begin{minipage}{\linewidth}
\hypertarget{TERM-1931}{}
\noindent\textbf{teoremy o razloženju norm i modulov}\quad{\footnotesize\texttt{TERM-1931}}\par
\noindent\emph{Iztočnik / Source:} Zerlegungssätze für Normen und Moduln\par
\noindent\emph{Rusky / Russian:} теоремы о разложении норм и модулей\par
\noindent\emph{Ukrajinsky / Ukrainian:} теореми про розклад норм і модулів\par
\smallskip\noindent\emph{Razlog:} Teoremy o razloženju norm i modulov govori o moduloch linearnyh form, ne o geometrickoj dlžině vektora. Norma jest opisana determinantom i produktom elementarnyh děliteljov. Dlja cělobrojnyh koeficientov dava čislo klas; dlja polinomov v jednoj proměnnoj odgovarajuči razměr dava stepenj normy. Razklad normy na vzajemno proste faktory r i s vodi k određenyh modulah toga že ranga. Ne smějemo prenesti vse te značenja v pozdějšu teoriju primarnogo razloženja bez konteksta.\par
\noindent\emph{Reason:} Theorems on decomposition of norms and modules concerns modules of linear forms, not the geometric length of a vector. The norm is described by a determinant and a product of elementary divisors. For integer coefficients it counts residue classes; for univariate polynomials the relevant dimension is given by its degree. A factorization of the norm into coprime factors r and s leads to specified modules of the same rank. These meanings must not be silently replaced by the terminology of later primary decomposition.\par
\noindent{\footnotesize Izvorne zapisy / Original records: HIST-g086-001, HIST-g092-011.}\par
\smallskip\noindent\emph{Sravnjenje s tekstom:} Norma zde ne jest geometricka dlžina.\par
\noindent\emph{Textual comparison:} Here a norm is not geometric length.\par
\noindent{\footnotesize TEXT-COMP-988; pasaže / passages: 14399, 14401, 14404, 14408, 14414, 14424, 14427, 14429, 14433, 14435; \texttt{CORPUS\_FORM\_COMPARISONS.json}.}\par
\end{minipage}\par
\par\medskip\Needspace{7\baselineskip}
\noindent\begin{minipage}{\linewidth}
\hypertarget{TERM-1932}{}
\noindent\textbf{Zermelova oblast}\quad{\footnotesize\texttt{TERM-1932}}\par
\noindent\emph{Iztočnik / Source:} Zermeloscher Bereich\par
\noindent\emph{Rusky / Russian:} область Цермело\par
\noindent\emph{Ukrajinsky / Ukrainian:} область Цермело; Цермелова область\par
\smallskip\noindent\emph{Razlog:} Zermelova oblast zachovuje imeno Ernsta Zermela i označaje oblast dostanu određenoju konstrukcijeju. Ne jest drugo imeno dlja vsakoj integralnoj oblasti. Tekst ispytuje zavisimost od vybranoj bazy, integralnu zatvorenost i uključanja v daljše oblasti. Pridevnična forma dava čitlivu nazvu bez iztraty istorijnoj pripisy. Ruske i ukrajinske formy s početnym Ц sut usporedne zapisy; sami po sobě ne ustanavljajut obvezno pravilo izgovora ili transliteracije v medžuslovjanskom.\par
\noindent\emph{Reason:} Zermelova oblast preserves Ernst Zermelo's name and denotes a domain obtained by a particular construction. It is not another name for every integral domain. The text examines dependence on the chosen basis, integral closure and inclusion in further domains. The adjectival form provides a readable label without losing the historical attribution. The Russian and Ukrainian forms with initial Ц are recorded comparanda; by themselves they establish no mandatory Interslavic pronunciation or transliteration rule.\par
\noindent{\footnotesize Izvorne zapisy / Original records: HIST-g059-006, HIST-g062-004.}\par
\smallskip\noindent\emph{Sravnjenje s tekstom:} Imeno označaje konkretnu konstrukciju, ne vse oblasti.\par
\noindent\emph{Textual comparison:} The name designates a particular construction, not all domains.\par
\noindent{\footnotesize TEXT-COMP-989; pasaže / passages: 7156, 7419, 7420, 7421, 7422, 7423, 7424, 7426, 7427, 7428, 7429, 7430, 7436, 7437, 7438, 7439, 7440, 7441, 7442, 7443, 7445, 7446, 7447, 7448, 7449; \texttt{CORPUS\_FORM\_COMPARISONS.json}.}\par
\end{minipage}\par
\par\medskip\Needspace{7\baselineskip}
\noindent\begin{minipage}{\linewidth}
\hypertarget{TERM-1933}{}
\noindent\textbf{primarna komponenta naležeča k $\mathfrak{o}$}\quad{\footnotesize\texttt{TERM-1933}}\par
\noindent\emph{Iztočnik / Source:} zu $\mathfrak{o}$ gehörige Primärkomponente\par
\noindent\emph{Rusky / Russian:} первичная компонента, принадлежащая $\mathfrak{o}$\par
\noindent\emph{Ukrajinsky / Ukrainian:} первинна компонента, що належить до $\mathfrak{o}$\par
\smallskip\noindent\emph{Razlog:} Primarna komponenta naležeča k označenomu prstenu jest izključena iz prvogo tvrđenja o jedinstvenosti. Ta ograda ne jest nevažna iznimka: dokaz posebej razmatrja slučaj, v ktorom takova komponenta jest prisutna. Tekst potom dodava jedinicu i dostaje silnějši rezultat; togda komponenta pridružena cělomu prstenu ne nastupaje v najkratšem razloženju pravogo ideala. Znak prstena jest matematičny simbol, ne slovo za prevod. Ne prěnosimo hipotezu o jedinici nazad v opčenějšu teoremu.\par
\noindent\emph{Reason:} The primary component belonging to the specified ring is excluded from the first uniqueness assertion. This is a substantive restriction: the proof separately treats the case in which such a component occurs. The text then adds an identity element and obtains a stronger result; the component associated with the whole ring no longer occurs in a shortest decomposition of a proper ideal. The ring sign is mathematical notation, not a word to translate. The later unital hypothesis must not be read back into the more general theorem.\par
\noindent{\footnotesize Izvorne zapisy / Original records: HIST-g114-013.}\par
\smallskip\noindent\emph{Sravnjenje s tekstom:} Hipoteza o jedinici pojavjaje se tek pri silnějšem rezultatě.\par
\noindent\emph{Textual comparison:} The unital hypothesis enters only with the stronger result.\par
\noindent{\footnotesize TEXT-COMP-990; pasaže / passages: 16780, 16782, 16788, 16792, 16794, 16796, 16800; \texttt{CORPUS\_FORM\_COMPARISONS.json}.}\par
\end{minipage}\par
\par\medskip\Needspace{7\baselineskip}
\noindent\begin{minipage}{\linewidth}
\hypertarget{TERM-1934}{}
\noindent\textbf{element prosty k $a$ / vzajemno prosty v idealnom smyslu}\quad{\footnotesize\texttt{TERM-1934}}\par
\noindent\emph{Iztočnik / Source:} zu $a$ primes Element / prim mit teilerfremd identisch\par
\noindent\emph{Rusky / Russian:} элемент, простой к $a$ / взаимно простой в смысле идеалов\par
\noindent\emph{Ukrajinsky / Ukrainian:} елемент, простий до $a$ / взаємно простий в ідеальному сенсі\par
\smallskip\noindent\emph{Razlog:} Element prosty k a ne znači prosty element. Predlog k izražaje odnos k idealu, a podnožna primětka izrično svazuje prostotu s vzajemnoju prostotoju v idealnom smyslu. To opravdava objasnjujuču alternativu vzajemno prosty v idealnom smyslu. Denominatory lokalizacije sut izbrane po tom uslovju. Arhivny zapis govori o obnovjenoj primětke; sadašnja usporedba ustanavlja njezin smysl v tekstě, no ne predstavlja novu proverku izvornogo skena.\par
\noindent\emph{Reason:} An element prime to a is not the same thing as a prime element. The preposition k expresses a relation to an ideal, and the footnote explicitly connects this use of prime with coprimality in the ideal-theoretic sense. That supports the explanatory alternative vzajemno prosty v idealnom smyslu. The localization denominators are selected by this condition. The archived record mentions a restored footnote; the present comparison establishes its meaning in the text, not a new examination of the original scan.\par
\noindent{\footnotesize Izvorne zapisy / Original records: HIST-g159-005.}\par
\smallskip\noindent\emph{Sravnjenje s tekstom:} Predlog izražaje odnos k idealu.\par
\noindent\emph{Textual comparison:} The preposition expresses a relation to an ideal.\par
\noindent{\footnotesize TEXT-COMP-991; pasaže / passages: 18932; \texttt{CORPUS\_FORM\_COMPARISONS.json}.}\par
\end{minipage}\par
\par\medskip\Needspace{7\baselineskip}
\noindent\begin{minipage}{\linewidth}
\hypertarget{TERM-1935}{}
\noindent\textbf{primarna komponenta prinadležeča k $p$}\quad{\footnotesize\texttt{TERM-1935}}\par
\noindent\emph{Iztočnik / Source:} zu $p$ gehörige Primärkomponente\par
\noindent\emph{Rusky / Russian:} первичная компонента, принадлежащая $p$\par
\noindent\emph{Ukrajinsky / Ukrainian:} первинна компонента, що належить до $p$\par
\smallskip\noindent\emph{Razlog:} Prinadležeča k p označaje primarnu komponentu pridruženu danomu prostomu idealu, ne samo proizvoljnu komponentu vključenu v njego. V lokalnom opisu avtorka uživa dva različne znaky: frakturno p dlja prostogo ideala i obyčno p dlja elementa, ktory generuje maksimalny ideal lokalizacije. Lokalizovany ideal jest generovany odgovarajučeju primarnoju komponentoju ili postaje cělym prstenom. Razlika dvuh znakov mora ostati vidima; oběkt pridruženja ne směje byti zamenjen njegovym generatorom.\par
\noindent\emph{Reason:} Prinadležeča k p denotes the primary component associated with the specified prime ideal, not an arbitrary component merely contained in it. In the local description the author uses two distinct signs: fraktur p for the prime ideal and ordinary p for an element generating the localization's maximal ideal. The localized ideal is generated by the corresponding primary component or becomes the whole ring. The two signs must remain distinct; the associated ideal must not be confused with its generator.\par
\noindent{\footnotesize Izvorne zapisy / Original records: HIST-g160-008.}\par
\smallskip\noindent\emph{Sravnjenje s tekstom:} Frakturny i obyčny znak p označajut razne objekty.\par
\noindent\emph{Textual comparison:} Fraktur and ordinary p denote different objects.\par
\noindent{\footnotesize TEXT-COMP-992; pasaže / passages: 18970; \texttt{CORPUS\_FORM\_COMPARISONS.json}.}\par
\end{minipage}\par
\par\medskip\Needspace{7\baselineskip}
\noindent\begin{minipage}{\linewidth}
\hypertarget{TERM-1936}{}
\noindent\textbf{polja izomorfne k $Q$}\quad{\footnotesize\texttt{TERM-1936}}\par
\noindent\emph{Iztočnik / Source:} zu $Q$ isomorphe Körper\par
\noindent\emph{Rusky / Russian:} поля, изоморфные $Q$\par
\noindent\emph{Ukrajinsky / Ukrainian:} поля, ізоморфні $Q$\par
\smallskip\noindent\emph{Razlog:} Polja izomorfne k Q zachovuje izomorfizm, ne doslovnu ravnost množestv. Komponenty imajut rang jeden nad Q i svoju jedinicu; kak polja sut izomorfne k Q. Znak Q zde označaje vybrano razširjenje, ne avtomatično polje racionalnyh čisel. Arhivna primětka upomina ranějšu proverku skena; naša objasnjenja oslanjajut se na dostupny zapis i nyněšnji tekst, ne na nedokazanu novu proverku. Predlog k jest izborem doslědnogo izražanja odnosa.\par
\noindent\emph{Reason:} Polja izomorfne k Q preserves isomorphism rather than literal equality of sets. The components have rank one over Q and their own identity elements; as fields they are isomorphic to Q. Here Q denotes the chosen extension, not automatically the rational numbers. The archival note mentions an earlier scan check; this explanation uses the preserved record and present text, not an unevidenced fresh scan comparison. The preposition k expresses the relation consistently.\par
\noindent{\footnotesize Izvorne zapisy / Original records: HIST-g146-005.}\par
\smallskip\noindent\emph{Sravnjenje s tekstom:} Q ne znači avtomatično racionalne čisla.\par
\noindent\emph{Textual comparison:} Q does not automatically mean the rational numbers.\par
\noindent{\footnotesize TEXT-COMP-993; pasaže / passages: 18470, 18472, 18474; \texttt{CORPUS\_FORM\_COMPARISONS.json}.}\par
\end{minipage}\par
\par\medskip\Needspace{7\baselineskip}
\noindent\begin{minipage}{\linewidth}
\hypertarget{TERM-1937}{}
\noindent\textbf{matrične kolca homomorfne k $R$}\quad{\footnotesize\texttt{TERM-1937}}\par
\noindent\emph{Iztočnik / Source:} zu $R$ homomorphe Matrizenringe\par
\noindent\emph{Rusky / Russian:} матричные кольца, гомоморфные к $R$\par
\noindent\emph{Ukrajinsky / Ukrainian:} матричні кільця, гомоморфні до $R$\par
\smallskip\noindent\emph{Razlog:} Matrične kolca homomorfne k R zachovuje istorijno izražanje odnosa, no dokaz objasnjaje njegov točny smysl. Dějstvije elementov prstena na modulnoj bazě ideala dava matrice, a matričny obraz jest izomorfny kvocientu po anuljatoru ideala. To ne tvrdi, že sam prsten vsegda jest izomorfny matričnomu obrazu: jadro može byti nenuljevo. V citovanom prevodě stojit matričny prsten; arhivna forma kolco ostaje svědectvom izbora, ne skrytoju tvrdboju o vseobščej doslědnosti.\par
\noindent\emph{Reason:} Matrične kolca homomorfne k R preserves the historical relational wording, but the proof explains its exact meaning. The action of ring elements on a module basis of an ideal gives matrices, and the matrix image is isomorphic to the quotient by the ideal's annihilator. This does not assert that the ring itself is always isomorphic to that image: the kernel may be nonzero. The cited translation uses prsten; the archived kolco remains evidence of a choice, not a concealed claim of uniform usage.\par
\noindent{\footnotesize Izvorne zapisy / Original records: HIST-g131-001.}\par
\smallskip\noindent\emph{Sravnjenje s tekstom:} Izomorfny jest kvocient, ne neobhodimo cěly iztočny prsten.\par
\noindent\emph{Textual comparison:} The quotient is isomorphic to the image, not necessarily the entire source ring.\par
\noindent{\footnotesize TEXT-COMP-994; pasaže / passages: 17858, 17860, 17874; \texttt{CORPUS\_FORM\_COMPARISONS.json}.}\par
\end{minipage}\par
\par\medskip\Needspace{7\baselineskip}
\noindent\begin{minipage}{\linewidth}
\hypertarget{TERM-1938}{}
\noindent\textbf{sootvětni prosti idealy}\quad{\footnotesize\texttt{TERM-1938}}\par
\noindent\emph{Iztočnik / Source:} zugehörige Primideale\par
\noindent\emph{Rusky / Russian:} соответствующие простые идеалы\par
\noindent\emph{Ukrajinsky / Ukrainian:} відповідні прості ідеали\par
\smallskip\noindent\emph{Razlog:} Sootvětni prosti idealy byl zapisan kak lehča alternativa težšej tvorbě dlja odnosnoje frazy. V danym tekstě, odnako, stojit pridružene proste idealy, i ti idealy imajut određeny odnos k primarnym komponentam. Slovnik zato ne predstavi arhivnu preferenciju kak doslědno provedenu zaměnu. Sootvětni može byti čitlivy v jasnom kontekstě, no samo po sobě ne vyslovjaje definiciju pridruženogo prostogo ideala. Slědujuče dva zapisy pokazuju druge rěšenja istogo terminologičnogo pytanja.\par
\noindent\emph{Reason:} Sootvětni prosti idealy was recorded as a lighter alternative for the relational phrase. The actual passage, however, uses pridružene proste idealy, and these ideals have a defined relation to the primary components. The dictionary therefore does not present the archival preference as a consistently implemented replacement. Sootvětni may be readable in an explicit context, but alone it does not express the definition of associated prime. The following two records show other solutions to the same terminological question.\par
\noindent{\footnotesize Izvorne zapisy / Original records: HIST-g150-005.}\par
\smallskip\noindent\emph{Sravnjenje s tekstom:} Sootvětni i pridruženy ne sut prikazani kako jedna vsegda upotrěbjena forma.\par
\noindent\emph{Textual comparison:} Sootvětni and pridruženy are not represented as one uniformly implemented form.\par
\noindent{\footnotesize TEXT-COMP-995; pasaže / passages: 18593; \texttt{CORPUS\_FORM\_COMPARISONS.json}.}\par
\end{minipage}\par
\par\medskip\Needspace{7\baselineskip}
\noindent\begin{minipage}{\linewidth}
\hypertarget{TERM-1939}{}
\noindent\textbf{asociovany prost ideal}\quad{\footnotesize\texttt{TERM-1939}}\par
\noindent\emph{Iztočnik / Source:} zugehöriges Primideal\par
\noindent\emph{Rusky / Russian:} ассоциированный простой идеал\par
\noindent\emph{Ukrajinsky / Ukrainian:} асоційований простий ідеал\par
\smallskip\noindent\emph{Razlog:} Asociovany prost ideal jest svesno vybran kako tehnična nazva, v otličji od obščego prinaležny ili sootvětni. Arhiv opisuje zaměnu i v ukrajinskoj terminologiji; pet zapisov sut ponovjene upotrěby i proverky istogo izbora, ne pet nezavisnyh iznašlenij. Dlja primarnogo ideala q prosty ideal p jest jedinstveny: elementi p imajut neku stepenj v q, i tekst svazuje te idealy točnymi uslovjami dělitelnosti. Ta definicija, a ne slučajna blizost v tekstě, opravdava osobny termin.\par
\noindent\emph{Reason:} Asociovany prost ideal was deliberately chosen as a technical term, in contrast to general prinaležny or sootvětni. The archive describes a corresponding Ukrainian terminology change; the five records concern repeated uses and checks of the same choice, not five independent inventions. For a primary ideal q, the prime ideal p is unique: its elements have some power in q, and the text relates the ideals through precise divisibility conditions. That definition, rather than incidental proximity in the prose, justifies a dedicated term.\par
\noindent{\footnotesize Izvorne zapisy / Original records: HIST-g095-009, HIST-g096-009, HIST-g097-019, HIST-g098-005, HIST-g099-005.}\par
\smallskip\noindent\emph{Sravnjenje s tekstom:} Definicija oddělja termin od obščego odnosa.\par
\noindent\emph{Textual comparison:} The definition distinguishes the term from a general relationship.\par
\noindent{\footnotesize TEXT-COMP-996; pasaže / passages: 15357, 15359, 15361, 15507, 15509; \texttt{CORPUS\_FORM\_COMPARISONS.json}.}\par
\end{minipage}\par
\par\medskip\Needspace{7\baselineskip}
\noindent\begin{minipage}{\linewidth}
\hypertarget{TERM-1940}{}
\noindent\textbf{pridruženy prosty ideal}\quad{\footnotesize\texttt{TERM-1940}}\par
\noindent\emph{Iztočnik / Source:} zugehöriges Primideal\par
\noindent\emph{Rusky / Russian:} ассоциированный простой идеал\par
\noindent\emph{Ukrajinsky / Ukrainian:} належний простий ідеал\par
\smallskip\noindent\emph{Razlog:} Pridruženy prosty ideal jest drugo zachovano rěšenje za associated prime, s prozračnoju osnovoju pridruž-. Tekst dokazuje jedinstvenost i ostro razlikuje primarny ideal od pridruženogo prostogo ideala i od njegovyh stepenij. Navedeny kontraprimer takože pokaže, že samotna svez s prostym idealom ne čini ideal primarnym. U razloženyh idealov jedinstvenost pridruženih prostyh idealov ne znači jedinstvenost vsakej nerazložimej komponenty. Arhivna potrěba uskladnjenja ukrajinskogo zapisa ostaje otvorena; ne prepisujemo ju kako izvršenu ekspertnu proverku.\par
\noindent\emph{Reason:} Pridruženy prosty ideal is another preserved solution for associated prime, using the transparent pridruž- formation. The text proves uniqueness and carefully distinguishes the primary ideal from its associated prime and from powers of that prime. A counterexample also shows that the relevant relation to a prime ideal alone does not make an ideal primary. In decompositions, uniqueness of the associated primes does not mean uniqueness of every irreducible component. The recorded need to harmonize the Ukrainian wording remains open; it is not rewritten as completed expert review.\par
\noindent{\footnotesize Izvorne zapisy / Original records: HIST-g081-041, HIST-g082-041.}\par
\smallskip\noindent\emph{Sravnjenje s tekstom:} Jedinstvenost prostyh idealov ne znači jedinstvenost vseh komponentov.\par
\noindent\emph{Textual comparison:} Uniqueness of primes does not mean uniqueness of all components.\par
\noindent{\footnotesize TEXT-COMP-997; pasaže / passages: 12817, 12819, 12820, 12821, 12822, 12823, 12824, 12825, 12826, 12827, 12828, 12829, 12830, 12834, 12950; \texttt{CORPUS\_FORM\_COMPARISONS.json}.}\par
\end{minipage}\par
\par\medskip\Needspace{7\baselineskip}
\noindent\begin{minipage}{\linewidth}
\hypertarget{TERM-1941}{}
\noindent\textbf{asociovany prosty ideal / naležati k $p$}\quad{\footnotesize\texttt{TERM-1941}}\par
\noindent\emph{Iztočnik / Source:} zugehöriges Primideal / zu $p$ gehören\par
\noindent\emph{Rusky / Russian:} ассоциированный простой идеал / принадлежать $p$\par
\noindent\emph{Ukrajinsky / Ukrainian:} асоційований простий ідеал / належати до $p$\par
\smallskip\noindent\emph{Razlog:} Asociovany prosty ideal imenuje određeny oběkt, a naležati k p izražaje odnos cěloj familije primarnyh idealov k tomu oběktu. Tekst sam daje tu glosu: primarne idealy s tym samym asociovanym prostym idealom naležet k njemu. Proto nije potrěbno zaměniti vsako obyčno naležati tehničnym asociovati. Uvod razlikuje slabe i silne primarne idealy, a njihov spoj potrěbuje uslovje cěpov; jednostajna terminologija ne směje izbrisati tu hipotezu.\par
\noindent\emph{Reason:} Asociovany prosty ideal names a particular object, while naležati k p relates a whole family of primary ideals to it. The text itself supplies the gloss: primary ideals with the same associated prime belong to that prime. There is therefore no need to replace every ordinary occurrence of belong with technical associate. The passage distinguishes weak and strong primary ideals, whose identification requires a chain condition; consistent wording must not erase that hypothesis.\par
\noindent{\footnotesize Izvorne zapisy / Original records: HIST-g112-008.}\par
\smallskip\noindent\emph{Sravnjenje s tekstom:} Naležati k ne znači novy oběkt različny od asociovanogo ideala.\par
\noindent\emph{Textual comparison:} Belong to does not name a new object distinct from the associated prime.\par
\noindent{\footnotesize TEXT-COMP-998; pasaže / passages: 16607, 16609, 16615, 16619; \texttt{CORPUS\_FORM\_COMPARISONS.json}.}\par
\end{minipage}\par
\par\medskip\Needspace{7\baselineskip}
\noindent\begin{minipage}{\linewidth}
\hypertarget{TERM-1942}{}
\noindent\textbf{pridruženy prosty ideal / pridruženo prosto čislo}\quad{\footnotesize\texttt{TERM-1942}}\par
\noindent\emph{Iztočnik / Source:} zugehöriges Primideal / zugehörige Primzahl\par
\noindent\emph{Rusky / Russian:} ассоциированный простой идеал / ассоциированное простое число\par
\noindent\emph{Ukrajinsky / Ukrainian:} належний простий ідеал / належне просте число\par
\smallskip\noindent\emph{Razlog:} Pridruženy prosty ideal i pridruženo prosto čislo davajut istu odnosnu tvorbu dlja dvuh različnih oběktov. Dlja primarnogo cělogo čisla pridruženo prosto čislo jest jedinstveno do znaka; za opči primarny ideal ne slěduje, že on jest stepenju pridruženogo prostogo ideala. Pozdějši dokaz uživa dělimost medžu prostymi idealami, da doloži relativnu prostotu i grupiranje komponentov. Arhivna ukrajinska forma ostaje označena za uskladnjenje; to ne měni matematičnu razliku medžu čislom i idealom.\par
\noindent\emph{Reason:} Pridruženy prosty ideal and pridruženo prosto čislo use the same relational formation for two different objects. For a primary integer, the associated prime number is unique up to sign; a general primary ideal need not be a power of its associated prime. The later argument uses divisibility among prime ideals to characterize relative primeness and group the components. The archived Ukrainian wording remains marked for harmonization; that does not change the distinction between a number and an ideal.\par
\noindent{\footnotesize Izvorne zapisy / Original records: HIST-g081-010, HIST-g082-010.}\par
\smallskip\noindent\emph{Sravnjenje s tekstom:} Čislova analogija ne dokazuje, že vsaky primarny ideal jest stepenj.\par
\noindent\emph{Textual comparison:} The numerical analogy does not prove that every primary ideal is a power.\par
\noindent{\footnotesize TEXT-COMP-999; pasaže / passages: 12426, 13120, 13150, 13152; \texttt{CORPUS\_FORM\_COMPARISONS.json}.}\par
\end{minipage}\par
\par\medskip\Needspace{7\baselineskip}
\noindent\begin{minipage}{\linewidth}
\hypertarget{TERM-1943}{}
\noindent\textbf{pridruženy prosty ideal v nekomutativnom slučaju}\quad{\footnotesize\texttt{TERM-1943}}\par
\noindent\emph{Iztočnik / Source:} zugehöriges Primideal im nicht-kommutativen Fall\par
\noindent\emph{Rusky / Russian:} ассоциированный простой идеал в некоммутативном случае\par
\noindent\emph{Ukrajinsky / Ukrainian:} належний простий ідеал у некомутативному випадку\par
\smallskip\noindent\emph{Razlog:} V nekomutativnom slučaju jest suštna ograda, ne novo bezuslovno razširjenje termina. Avtorka pokazuje, kde obči prenos dokaza ne uspěvaje: iz nilpotentnosti dvuh elementov ne slěduje potrěbna nilpotentnost njihove razlike, a potega produkta ne može byti prosto razdělena na potegi faktorov. To ne jest tvrđenje, že nikaka teorija pridruženih prostyh idealov v nekomutativnoj algebrě ne može obstajati. Slovjansky pridruženy ostaje izborom nazvy; ne opravdava tvrditi egzistenciju, ktoru dany argument ne dokazuje.\par
\noindent\emph{Reason:} In the noncommutative case is an essential qualification, not an unconditional extension of the term. The author identifies where a formal transfer of the proof fails: nilpotence of two elements does not yield the required nilpotence of their difference, and a power of a product cannot simply be separated into powers of the factors. This does not assert that no theory of associated primes can exist in noncommutative algebra. The Slavic pridruženy remains a naming choice, not grounds for asserting existence that the argument does not prove.\par
\noindent{\footnotesize Izvorne zapisy / Original records: HIST-g081-075.}\par
\smallskip\noindent\emph{Sravnjenje s tekstom:} Neuspěh danyh dokazov ne jest teorema o nemožnosti vsej teorije.\par
\noindent\emph{Textual comparison:} Failure of these proofs is not an impossibility theorem for an entire theory.\par
\noindent{\footnotesize TEXT-COMP-1000; pasaže / passages: 13409; \texttt{CORPUS\_FORM\_COMPARISONS.json}.}\par
\end{minipage}\par
\par\medskip\Needspace{7\baselineskip}
\noindent\begin{minipage}{\linewidth}
\hypertarget{TERM-1944}{}
\noindent\textbf{pridruženy prosty ideal}\quad{\footnotesize\texttt{TERM-1944}}\par
\noindent\emph{Iztočnik / Source:} zugehöriges Primideal [upotrěba v dokazu jedinstvenosti v §7 / use in the uniqueness proof in §7]\par
\noindent\emph{Rusky / Russian:} ассоциированный простой идеал\par
\noindent\emph{Ukrajinsky / Ukrainian:} належний простий ідеал\par
\smallskip\noindent\emph{Razlog:} Tuto jest ponovjena upotrěba pridruženy prosty ideal v dokazě jedinstvenosti izolovanyh komponentov. Anglijsky dodatok o sedmoj sekciji v arhivnom polju jest opis upotrěby, ne čest německogo termina. Teorema doloži, že izolovane idealy sut određene svojimi pridruženymi prostymi idealami; mnogoznačnost ostaje pri neizolovanyh komponentah. Zato razlog za doslědnu nazvu jest možnost slědovati tom odnose, ne nova jezyčna invencija ili tvrđenje o jedinstvenosti vseh komponentov.\par
\noindent\emph{Reason:} This is a repeated use of pridruženy prosty ideal in the uniqueness argument for isolated components. The English addition about section seven in the archived field describes the use; it is not part of the German term. The theorem establishes that isolated ideals are determined by their associated primes, while ambiguity remains for nonisolated components. Consistent naming lets the reader follow that relationship. It is neither a new linguistic invention nor a claim that every component is unique.\par
\noindent{\footnotesize Izvorne zapisy / Original records: HIST-g081-059, HIST-g082-059.}\par
\smallskip\noindent\emph{Sravnjenje s tekstom:} Jedinstvenost odnosi se na određenu klasu komponentov.\par
\noindent\emph{Textual comparison:} Uniqueness concerns a specified class of components.\par
\noindent{\footnotesize TEXT-COMP-1001; pasaže / passages: 13214, 13220, 13222, 13228, 13242, 13244; \texttt{CORPUS\_FORM\_COMPARISONS.json}.}\par
\end{minipage}\par
\par\medskip\Needspace{7\baselineskip}
\noindent\begin{minipage}{\linewidth}
\hypertarget{TERM-1945}{}
\noindent\textbf{prirědžene idealy}\quad{\footnotesize\texttt{TERM-1945}}\par
\noindent\emph{Iztočnik / Source:} zugeordnete Ideale\par
\noindent\emph{Rusky / Russian:} поставленные в соответствие идеалы\par
\noindent\emph{Ukrajinsky / Ukrainian:} поставлені у відповідність ідеали\par
\smallskip\noindent\emph{Razlog:} Prirědžene idealy znači idealy, določene konkretnym prirědženjem iz modulov. V dokazě modulnomu oběktu odgovarjaje konečny sistem koeficientnyh idealov; strogy prehod medžu modulami dava prinajmenje jeden strogy prehod medžu tymi idealami. Ta svez prenosi uslovje cěpov. Zato anglijsku arhivnu glosu associated ideals ne trěba razuměti kak nazvu asociovanyh prostyh idealov. Prirědžene opisuje konstrukciju, bez pridavanja neustanovjenej prostoty.\par
\noindent\emph{Reason:} Prirědžene idealy means ideals specified by a particular assignment from modules. In the proof, a module is assigned a finite system of coefficient ideals; a strict step between modules yields at least one strict step between those ideals. This transfers the chain condition. The archive's English gloss associated ideals must therefore not be read as the technical associated-prime terminology. Prirědžene describes the construction without asserting primality that has not been established.\par
\noindent{\footnotesize Izvorne zapisy / Original records: HIST-g109-008.}\par
\smallskip\noindent\emph{Sravnjenje s tekstom:} Prirědžene idealy ne sut avtomatično proste.\par
\noindent\emph{Textual comparison:} Assigned ideals are not automatically prime.\par
\noindent{\footnotesize TEXT-COMP-1002; pasaže / passages: 16314, 16316, 16318, 16320, 16334, 16343; \texttt{CORPUS\_FORM\_COMPARISONS.json}.}\par
\end{minipage}\par
\par\medskip\Needspace{7\baselineskip}
\noindent\begin{minipage}{\linewidth}
\hypertarget{TERM-1946}{}
\noindent\textbf{prirěđene matrice}\quad{\footnotesize\texttt{TERM-1946}}\par
\noindent\emph{Iztočnik / Source:} zugeordnete Matrizen\par
\noindent\emph{Rusky / Russian:} сопоставленные матрицы\par
\noindent\emph{Ukrajinsky / Ukrainian:} матриці, поставлені у відповідність\par
\smallskip\noindent\emph{Razlog:} Prirěđene matrice sut matrice dějstvija elementov prstena po izboru modulnoj bazy. Ta baza jest čest konstrukcije, ne izpuščena podrobnost: koeficienty razloženja davajut matrice, razlici odpovědaje razlika i produktu produkt. Nuljeva matrica odpovědaje točno elementam anuljujučim ideal. Zato slovo prirěđene izražaje određeno pravilo prirědženja, ne opču podobnost matric. Aktualny tekst piše prirědžene; originalny zapis s đ ostaje istorijnym pravopisnym variantom.\par
\noindent\emph{Reason:} Prirěđene matrice are the matrices describing the action of ring elements after a module basis has been chosen. The basis is part of the construction, not an omitted detail: expansion coefficients give the matrices, and differences and products correspond. Precisely the elements annihilating the ideal give the zero matrix. The adjective therefore expresses a definite assignment, not a general resemblance between matrices. The current text writes prirědžene; the archived spelling with đ remains a historical orthographic variant.\par
\noindent{\footnotesize Izvorne zapisy / Original records: HIST-g131-004.}\par
\smallskip\noindent\emph{Sravnjenje s tekstom:} Prirědženje i njegov anuljator imajut izričnu definiciju.\par
\noindent\emph{Textual comparison:} The assignment and its annihilator have an explicit definition.\par
\noindent{\footnotesize TEXT-COMP-1003; pasaže / passages: 17858, 17860, 17874, 17880; \texttt{CORPUS\_FORM\_COMPARISONS.json}.}\par
\end{minipage}\par
\par\medskip\Needspace{7\baselineskip}
\noindent\begin{minipage}{\linewidth}
\hypertarget{TERM-1947}{}
\noindent\textbf{prirędžene rędy}\quad{\footnotesize\texttt{TERM-1947}}\par
\noindent\emph{Iztočnik / Source:} zugeordnete Reihen\par
\noindent\emph{Rusky / Russian:} сопоставленные строки\par
\noindent\emph{Ukrajinsky / Ukrainian:} зіставлені рядки\par
\smallskip\noindent\emph{Razlog:} Prirędžene rędy prodolžaje terminologiju korespondencije iz uvoda raboty. Zde red ne jest beskonečny čislovy red: elementi redov sut determinanty, a označene redy Q i P svazujut se vyraznymi formulami. Prehod od gornjih indeksov k dolnjim i obratno jest čest matematiki, ne tipografska sloboda. Aktualno priredžene redy i arhivno prirędžene rędy sut različne grafije toj nazvy; otvorena jezyčna ocěna ne jest dovršena samo tym, že relacija byla lokalizovana.\par
\noindent\emph{Reason:} Prirędžene rędy continues the introduction's correspondence terminology. Here a row is not an infinite numerical series: its entries are determinants, and the specified Q and P rows are related by explicit formulas. Moving from upper to lower indices and conversely belongs to the mathematics, not typographical discretion. Current priredžene redy and archived prirędžene rędy are different spellings of that name. Locating the relation does not by itself complete the outstanding linguistic assessment.\par
\noindent{\footnotesize Izvorne zapisy / Original records: HIST-g031-008.}\par
\smallskip\noindent\emph{Sravnjenje s tekstom:} Gornje i dolnje indeksy imajut različne ulogi.\par
\noindent\emph{Textual comparison:} Upper and lower indices have different roles.\par
\noindent{\footnotesize TEXT-COMP-1004; pasaže / passages: 3831, 3834, 3835, 3836, 3837; \texttt{CORPUS\_FORM\_COMPARISONS.json}.}\par
\end{minipage}\par
\par\medskip\Needspace{7\baselineskip}
\noindent\begin{minipage}{\linewidth}
\hypertarget{TERM-1948}{}
\noindent\textbf{dopustimy sistem}\quad{\footnotesize\texttt{TERM-1948}}\par
\noindent\emph{Iztočnik / Source:} zul?ssiges System [Čestično nezachovane bukvy / Partly unrecoverable lettering]\par
\noindent\emph{Rusky / Russian:} [Bukvy ne sut zachovane / Lettering not recoverable]\par
\noindent\emph{Ukrajinsky / Ukrainian:} [Bukvy ne sut zachovane / Lettering not recoverable]\par
\smallskip\noindent\emph{Razlog:} Dopustimy sistem v šestoj sekciji izpolnjaje točno uslovje: racionalno-cěla adjunkcija ne vvede algebraično-drobno čislo v oblast. Pridevnik znači ispolnjenje toga uslovja, ne obšču priměrnost sistema. U ranějšoj sekciji analogično město govori o algebraično-cěloj adjunkciji; tekst izrično objasnjaje zaměnu slova algebraično slovom racionalno. Ista nazva zato potrěbuje dve kontekstne glosy. Povrědžene srovniteljne zapisy ne služat svědectvom novoj ruskoj ili ukrajinskoj attestacije.\par
\noindent\emph{Reason:} In section six, dopustimy sistem satisfies an exact condition: rationally integral adjunction must introduce no algebraically fractional number into the domain. The adjective means satisfaction of that condition, not general suitability. The corresponding passage in the earlier section concerns algebraically integral adjunction; the text expressly explains the substitution of rationally for algebraically. One name therefore needs two contextual glosses. The damaged comparative records supply no fresh Russian or Ukrainian attestation.\par
\noindent{\footnotesize Izvorne zapisy / Original records: HIST-g064-011.}\par
\smallskip\noindent\emph{Sravnjenje s tekstom:} Definicija ne dopušča algebraično-drobne čisla v oblasti.\par
\noindent\emph{Textual comparison:} The definition excludes algebraically fractional numbers from the domain.\par
\noindent{\footnotesize TEXT-COMP-1005; pasaže / passages: 7785, 7786, 7787, 7788, 7789, 7790, 7791, 7792, 7793, 7794, 7795, 7802, 7803, 7804, 7805, 7806, 7807; \texttt{CORPUS\_FORM\_COMPARISONS.json}.}\par
\end{minipage}\par
\par\medskip\Needspace{7\baselineskip}
\noindent\begin{minipage}{\linewidth}
\hypertarget{TERM-1949}{}
\noindent\textbf{dopustima podgrupa}\quad{\footnotesize\texttt{TERM-1949}}\par
\noindent\emph{Iztočnik / Source:} zulässige Untergruppe\par
\noindent\emph{Rusky / Russian:} допустимая подгруппа\par
\noindent\emph{Ukrajinsky / Ukrainian:} допустима підгрупа\par
\smallskip\noindent\emph{Razlog:} Dopustima podgrupa jest podgrupa zatvorjena pod vsimi operatorami fiksovanoj oblasti operatorov. To ne jest svedeno samo k normalnoj podgrupě: pri vnutrnyh avtomorfizmah dobyvajut se normalne podgrupy, pri vsih avtomorfizmah karakteristične. Izbor dopustim- prodolžaje zabilěženu familiju i imeje blizke ruske i ukrajinske formy. Definiciju, odnako, nosi ukazany operatorny sistem. Bez njega kratka nazva ne opreděljaje, ktore podgrupy sut dopustime.\par
\noindent\emph{Reason:} A dopustima podgrupa is a subgroup closed under every operator in the fixed operator domain. It is not simply a synonym for normal subgroup: inner automorphisms give normal subgroups, while all automorphisms give characteristic subgroups. Choosing dopustim- continues the recorded family and closely parallels the Russian and Ukrainian forms. The specified operator system supplies the definition, however. Without it, the short name does not determine which subgroups are admissible.\par
\noindent{\footnotesize Izvorne zapisy / Original records: HIST-g179-002.}\par
\smallskip\noindent\emph{Sravnjenje s tekstom:} Normalnost i karakterističnost sut razne priměry.\par
\noindent\emph{Textual comparison:} Normality and characteristicity are different examples.\par
\noindent{\footnotesize TEXT-COMP-1006; pasaže / passages: 19803, 19805, 19807, 19811, 19813, 19815, 19821; \texttt{CORPUS\_FORM\_COMPARISONS.json}.}\par
\end{minipage}\par
\par\medskip\Needspace{7\baselineskip}
\noindent\begin{minipage}{\linewidth}
\hypertarget{TERM-1950}{}
\noindent\textbf{dopustimy ideal}\quad{\footnotesize\texttt{TERM-1950}}\par
\noindent\emph{Iztočnik / Source:} zulässiges Ideal\par
\noindent\emph{Rusky / Russian:} допустимый идеал\par
\noindent\emph{Ukrajinsky / Ukrainian:} допустимий ідеал\par
\smallskip\noindent\emph{Razlog:} Dopustimy ideal prodolžaje tu samu operatornu terminologiju v teoriji idealov. K obyčnoj levostrannoj, pravostrannoj ili dvustrannoj idealnoj strukture dodavaje se uslovje byti K-modulom; i homomorfizm mora soglasovati se s dějstvijem K. Ne ide o moralnu ili redakcijnu dopustimost. Iztočna primětka izrično govori, že v hiperkompleksnyh sistemah obična nazva ideal vključaje tu dopustimost po odnošenju k osnovnomu koeficientnomu polju. Ta konvencija ne smí byti bez objasnjenja prěnesena vo vse raboty.\par
\noindent\emph{Reason:} Dopustimy ideal carries the same operator vocabulary into ideal theory. Besides being a left, right or two-sided ideal, the object must be a K-module, and the homomorphism must respect the K-action. This is not moral or editorial admissibility. The source expressly says that, in hypercomplex systems, the unqualified word ideal includes admissibility relative to the underlying coefficient field. That convention must not be transferred without explanation to every paper.\par
\noindent{\footnotesize Izvorne zapisy / Original records: HIST-g186-003.}\par
\smallskip\noindent\emph{Sravnjenje s tekstom:} Dodatna struktura ograničaje i idealy i homomorfizmy.\par
\noindent\emph{Textual comparison:} The extra structure restricts both ideals and homomorphisms.\par
\noindent{\footnotesize TEXT-COMP-1007; pasaže / passages: 20408, 20410, 20412; \texttt{CORPUS\_FORM\_COMPARISONS.json}.}\par
\end{minipage}\par
\par\medskip\Needspace{7\baselineskip}
\noindent\begin{minipage}{\linewidth}
\hypertarget{TERM-1951}{}
\noindent\textbf{dopustimy sistem}\quad{\footnotesize\texttt{TERM-1951}}\par
\noindent\emph{Iztočnik / Source:} zulässiges System\par
\noindent\emph{Rusky / Russian:} допустимая система\par
\noindent\emph{Ukrajinsky / Ukrainian:} допустима система\par
\smallskip\noindent\emph{Razlog:} Dopustimy sistem v četvrtoj sekciji označaje sistem, pri ktorom algebraično-cěla adjunkcija k Zermelovoj oblasti ne vvede algebraično-drobne čisla. Izpolnjenje uslovja dava določeny vid oblasti, a obratny argument ispytuje vse oblasti nazvanoj familije. Arhivny razlog naziva dopustimy prozračnym; to jest jezyčna ocěna, ne izměrena razumljivost. Pozdějša racionalna analogija jest svezana s istim slovom, no ne dovoljuje zaměniti dve operacije adjunkcije jedna drugoju.\par
\noindent\emph{Reason:} In section four, dopustimy sistem denotes a system whose algebraically integral adjunction to the Zermelo domain introduces no algebraically fractional numbers. The condition gives the specified kind of domain, and the converse considers every domain in that family. The archived rationale calls dopustimy transparent; that is a linguistic judgment, not measured intelligibility. The later rational analogue uses the same word, but that does not make the two adjunction operations interchangeable.\par
\noindent{\footnotesize Izvorne zapisy / Original records: HIST-g062-010.}\par
\smallskip\noindent\emph{Sravnjenje s tekstom:} Ta analogija ne identifikuje dve operacije adjunkcije.\par
\noindent\emph{Textual comparison:} This analogy does not identify the two adjunction operations.\par
\noindent{\footnotesize TEXT-COMP-1008; pasaže / passages: 7523, 7524, 7525, 7526, 7527, 7528, 7529, 7531, 7532, 7533, 7534, 7535, 7536, 7537, 7539, 7540, 7541, 7542; \texttt{CORPUS\_FORM\_COMPARISONS.json}.}\par
\end{minipage}\par
\par\medskip\Needspace{7\baselineskip}
\noindent\begin{minipage}{\linewidth}
\hypertarget{TERM-1953}{}
\noindent\textbf{prirędženje}\quad{\footnotesize\texttt{TERM-1953}}\par
\noindent\emph{Iztočnik / Source:} Zuordnung\par
\noindent\emph{Rusky / Russian:} сопоставление\par
\noindent\emph{Ukrajinsky / Ukrainian:} зіставлення\par
\smallskip\noindent\emph{Razlog:} Prirędženje jest imenovnik dlja prirědženja objekta drugomu objektu, odděleny v arhivě od širšej nazvy zobraženje. V danyh pasažah prirědženje racionalnyh funkcij jest vzajemno jednoznačno i zachovuje relacije v obojih směrah; ničto z toga ne slěduje iz samogo slova. Tekst piše priredženje, bez nosovogo znaka arhivnoj formy. Razlog jezyčnogo izbora jest zachovati korespondenciju kak dějstvije, a konkretnu bijektivnost nosit teorema.\par
\noindent\emph{Reason:} Prirędženje is the noun for assigning one object to another, distinguished in the archive from the broader mapping term zobraženje. In the cited passages, the assignment of rational functions is bijective and preserves relations in both directions; none of that follows from the word alone. The current text writes priredženje without the archived nasal spelling. The lexical aim is to retain assignment as an operation; the theorem supplies the particular bijectivity.\par
\noindent{\footnotesize Izvorne zapisy / Original records: HIST-g043-004.}\par
\smallskip\noindent\emph{Sravnjenje s tekstom:} Vzajemna jednoznačnost jest vyrazna hipoteza ili rezultat.\par
\noindent\emph{Textual comparison:} Bijectivity is an explicit condition or result.\par
\noindent{\footnotesize TEXT-COMP-1010; pasaže / passages: 5325, 5374; \texttt{CORPUS\_FORM\_COMPARISONS.json}.}\par
\end{minipage}\par
\par\medskip\Needspace{7\baselineskip}
\noindent\begin{minipage}{\linewidth}
\hypertarget{TERM-1956}{}
\noindent\textbf{zvedenje procesa svijanja na osnovne svijanja}\quad{\footnotesize\texttt{TERM-1956}}\par
\noindent\emph{Iztočnik / Source:} Zurückführung des Faltungsprozesses auf Grundfaltungen\par
\noindent\emph{Rusky / Russian:} сведение процесса свертки к основным сверткам\par
\noindent\emph{Ukrajinsky / Ukrainian:} зведення процесу згортання до основних згортань\par
\smallskip\noindent\emph{Razlog:} Zvedenje procesa svijanja na osnovne svijanja bylo istorijnym naslovnym rěšenjem. Aktualny naslov piše proces kontrakcije i osnovne kontrakcije. Obe frazy potrebuju isti směr: opči proces sastavjaje se iz određenih osnovnyh operacij, ne vračaje se k ranějšej istorijnoj fazě. Teorema govori o n−1 tipah i razlikuje osnovne kontrakcije od širšej klasy kogredientnyh kontrakcij. Slovnik zato pokazuje i zabilěženy variant svijanje, i stvarnu pozdějšu formu kontrakcija, bez predstavy, že vsaka kogredientna operacija jest osnovna.\par
\noindent\emph{Reason:} Zvedenje procesa svijanja na osnovne svijanja was the historical title choice. The present heading uses proces kontrakcije and osnovne kontrakcije. Both need the same direction: the general process is composed of specified primitive operations, not returned to an earlier historical stage. The theorem gives n−1 types and distinguishes basic contractions from the larger class of cogredient contractions. The dictionary therefore shows both archived svijanje and later kontrakcija without suggesting that every cogredient operation is basic.\par
\noindent{\footnotesize Izvorne zapisy / Original records: HIST-g037-001.}\par
\smallskip\noindent\emph{Sravnjenje s tekstom:} Osnovna kontrakcija jest užša kategorija od kogredientnoj.\par
\noindent\emph{Textual comparison:} Basic contraction is narrower than cogredient contraction.\par
\noindent{\footnotesize TEXT-COMP-1013; pasaže / passages: 4647, 4649, 4651, 4653; \texttt{CORPUS\_FORM\_COMPARISONS.json}.}\par
\end{minipage}\par
\par\medskip\Needspace{7\baselineskip}
\noindent\begin{minipage}{\linewidth}
\hypertarget{TERM-1957}{}
\noindent\textbf{svedenje modula $(A,\nu)$ k modulu $(\nu)$}\quad{\footnotesize\texttt{TERM-1957}}\par
\noindent\emph{Iztočnik / Source:} Zurückführung des Moduls $(A,\nu)$ auf den Modul $(\nu)$\par
\noindent\emph{Rusky / Russian:} сведение модуля $(A,\nu)$ к модулю $(\nu)$\par
\noindent\emph{Ukrajinsky / Ukrainian:} зведення модуля $(A,\nu)$ до модуля $(\nu)$\par
\smallskip\noindent\emph{Razlog:} Svedenje modula (A,ν) k modulu (ν) zachovuje iztočny i ciljevy modul; dvě oznaky ne mogut byti spojene v neodređeno svedenje. Znak A v kratkom arhivnom zapisu v aktualnom naslově ima osobnu frakturnu formu, a ν ostaje grčskim simbolom. Nazva prodolžaje prethodnu sekciju, kde svedenje govori o kontrakcijah v konstrukciji sistema form. Ne znači zameniti vse raboty s modulom jedinoju formulačnou zaměnoju, a elipsa iz kračej susědnoj stavky ne odstraňuje konkretne argumenty.\par
\noindent\emph{Reason:} Svedenje modula (A,ν) k modulu (ν) retains both the source and target module; the two labels cannot be collapsed into an unspecified reduction. The abbreviated archive's A has a distinct fraktur form in the actual heading, while ν remains a Greek symbol. The title continues the preceding section, where reduction concerns contractions in the construction of a system of forms. It does not mean replacing every operation on a module with one symbolic substitution, and the neighbouring abbreviated entry does not remove these concrete arguments.\par
\noindent{\footnotesize Izvorne zapisy / Original records: HIST-g007-001.}\par
\smallskip\noindent\emph{Sravnjenje s tekstom:} Oba argumenty ostajut vidime.\par
\noindent\emph{Textual comparison:} Both arguments remain visible.\par
\noindent{\footnotesize TEXT-COMP-1014; pasaže / passages: 865, 700; \texttt{CORPUS\_FORM\_COMPARISONS.json}.}\par
\end{minipage}\par
\par\medskip\Needspace{7\baselineskip}
\noindent\begin{minipage}{\linewidth}
\hypertarget{TERM-1958}{}
\noindent\textbf{svedenje modula ... k ...}\quad{\footnotesize\texttt{TERM-1958}}\par
\noindent\emph{Iztočnik / Source:} Zurückführung des Moduls ... auf ...\par
\noindent\emph{Rusky / Russian:} сведение модуля ... к ...\par
\noindent\emph{Ukrajinsky / Ukrainian:} зведення модуля ... до ...\par
\smallskip\noindent\emph{Razlog:} Svedenje modula ... k ... bilo izbrano, da by internacionalny termin redukcija ne zaměnjal vsaku konstrukciju německogo Zurückführung. Iztočny naslov i tekst ukazujut točny cilj: svesti kontrakcije s jednim modulom k kontrakcijam s vyššim modulom. Slovjansky imenovnik daje čitlivy opis prehoda, no sam ne opreděljaje algebraičnu proceduru. Arhivna ocěna domačnosti i prozračnosti ne jest statističnym rezultatěm; slovnik ne dodava nedokumentovane težiny pri sravnenju s redukcijeju.\par
\noindent\emph{Reason:} Svedenje modula ... k ... was chosen so that the international term redukcija would not replace every construction expressed by German Zurückführung. The actual heading and text give the target: reduce contractions involving one module to contractions involving a higher module. The Slavic noun makes that transition readable but does not itself define the algebraic procedure. The archived assessment of familiarity and transparency is not a statistical result; no undocumented weights are added to the comparison with redukcija.\par
\noindent{\footnotesize Izvorne zapisy / Original records: HIST-g006-001.}\par
\smallskip\noindent\emph{Sravnjenje s tekstom:} Slovo svedenje ne zaměnjaje definiciju konstrukcije.\par
\noindent\emph{Textual comparison:} The word reduction does not replace the construction's definition.\par
\noindent{\footnotesize TEXT-COMP-1015; pasaže / passages: 655, 700; \texttt{CORPUS\_FORM\_COMPARISONS.json}.}\par
\end{minipage}\par
\par\medskip\Needspace{7\baselineskip}
\noindent\begin{minipage}{\linewidth}
\hypertarget{TERM-1959}{}
\noindent\textbf{prinaležne sistemi veličin / sistem nul}\quad{\footnotesize\texttt{TERM-1959}}\par
\noindent\emph{Iztočnik / Source:} zusammengehörige Wertsysteme / Nullstellensystem\par
\noindent\emph{Rusky / Russian:} связанные системы значений / система нулей\par
\noindent\emph{Ukrajinsky / Ukrainian:} пов'язані системи значень / система нулів\par
\smallskip\noindent\emph{Razlog:} Prinaležne sistemi veličin znači svazane n-tice značenij, ktore vměstě izpolnjajut polinomne uravnenja; veličiny zde ne sut fizične razměry. Pri sukcesivnoj eliminaciji jednotlivy korenj dopolnjaje se do celogo sistema, a odgovarajuči linearne faktory čuvajut tu svez. Sistem nul ne jest prosto nesvezany spis vsih korenjev po každoj proměnnoj oddělno. Arhiv označuje prinaležny kak jezyčno sporny; matematična lokalizacija objasnjaje predmet izbora, ale ne izmyslja nedostavajuču ekspertnu ocěnu.\par
\noindent\emph{Reason:} Prinaležne sistemi veličin denotes linked tuples of values satisfying the polynomial equations together; veličiny does not mean physical dimensions here. In successive elimination, an individual root is completed to a whole system, and the corresponding linear factors retain that connection. A system of zeros is not simply an unrelated list of roots for each variable separately. The archive marks prinaležny for linguistic scrutiny. Locating the mathematics explains the object of the choice without inventing the missing expert assessment.\par
\noindent{\footnotesize Izvorne zapisy / Original records: HIST-g090-012, HIST-g092-066.}\par
\smallskip\noindent\emph{Sravnjenje s tekstom:} Koordinaty dolžny vměstě izpolnjati uravnenja.\par
\noindent\emph{Textual comparison:} The coordinates must satisfy the equations together.\par
\noindent{\footnotesize TEXT-COMP-1016; pasaže / passages: 14930, 14952, 14956, 14963, 14971, 15003; \texttt{CORPUS\_FORM\_COMPARISONS.json}.}\par
\end{minipage}\par
\par\medskip\Needspace{7\baselineskip}
\noindent\begin{minipage}{\linewidth}
\hypertarget{TERM-1960}{}
\noindent\textbf{prinaležny sistem nuljev}\quad{\footnotesize\texttt{TERM-1960}}\par
\noindent\emph{Iztočnik / Source:} zusammengehöriges Nullstellensystem\par
\noindent\emph{Rusky / Russian:} соответствующая система нулей\par
\noindent\emph{Ukrajinsky / Ukrainian:} відповідна система нулів\par
\smallskip\noindent\emph{Razlog:} Prinaležny sistem nuljev naměrno ostaje obščim odnosnym izrazom, ne tehničnym asociovanym prostym idealom. V razloženju formy elementarnyh děliteljov odin indeks ν svazuje koordinaty jednogo sistema nuljev; različne indeksy odgovarjajut različnym linearnym faktoram. Slovnik mora zachovati tu společnost koordinat i ograditi ju od formalnoj terminologije asociovanyh idealov. Dodatne uslovja o doskonalom ili nedoskonalom polju upravljajut eksponentami faktorov, ne grammatikoju pridevnika.\par
\noindent\emph{Reason:} Prinaležny sistem nuljev deliberately remains a general relational expression, not the technical associated-prime term. In the factorization of the elementary-divisor form, one index ν links the coordinates of a single zero system; different indices correspond to different linear factors. The dictionary must retain that coordination and distinguish it from formal associated-ideal terminology. Further conditions concerning perfect or imperfect fields govern factor exponents, not the grammar of the adjective.\par
\noindent{\footnotesize Izvorne zapisy / Original records: HIST-g097-013.}\par
\smallskip\noindent\emph{Sravnjenje s tekstom:} Odin indeks svazuje koordinaty istogo sistema.\par
\noindent\emph{Textual comparison:} One index links coordinates of the same system.\par
\noindent{\footnotesize TEXT-COMP-1017; pasaže / passages: 15453, 15478, 15484; \texttt{CORPUS\_FORM\_COMPARISONS.json}.}\par
\end{minipage}\par
\par\medskip\Needspace{7\baselineskip}
\noindent\begin{minipage}{\linewidth}
\hypertarget{TERM-1961}{}
\noindent\textbf{sostavne identičnosti}\quad{\footnotesize\texttt{TERM-1961}}\par
\noindent\emph{Iztočnik / Source:} zusammengesetzte Identitäten\par
\noindent\emph{Rusky / Russian:} составные тождества\par
\noindent\emph{Ukrajinsky / Ukrainian:} складені тотожності\par
\smallskip\noindent\emph{Razlog:} Sostavne identičnosti označajut identičnosti dobyty određenoju specializacijeju i potom priměnjenjem už znanej identičnosti. Sostavne zato ne znači prosto složne ili težke dlja razuměnja. Avtorka dodava važnu konvenciju: identičnosti s k redami simbolov ostajut razgledane kak nerazložime, i kogda nastali podobnym postupkom iz menjšego čisla redov. Slovnik mora pokazati to ograničenje definicije, naměsto načiniti iz sostavny obšči opis vsakoj izvedene formule.\par
\noindent\emph{Reason:} Sostavne identičnosti denotes identities obtained by a specified specialization followed by application of an existing identity. Composite therefore does not merely mean complicated or difficult to understand. The author adds an important convention: identities with k symbolic rows are still treated as indecomposable even when obtained by a similar procedure from fewer rows. The dictionary must show this definitional restriction rather than turn composite into a generic description of every derived formula.\par
\noindent{\footnotesize Izvorne zapisy / Original records: HIST-g032-003.}\par
\smallskip\noindent\emph{Sravnjenje s tekstom:} Sostavny ne znači prosto težky.\par
\noindent\emph{Textual comparison:} Composite does not merely mean difficult.\par
\noindent{\footnotesize TEXT-COMP-1018; pasaže / passages: 3897; \texttt{CORPUS\_FORM\_COMPARISONS.json}.}\par
\end{minipage}\par
\par\medskip\Needspace{7\baselineskip}
\noindent\begin{minipage}{\linewidth}
\hypertarget{TERM-1962}{}
\noindent\textbf{složene polary}\quad{\footnotesize\texttt{TERM-1962}}\par
\noindent\emph{Iztočnik / Source:} zusammengesetzte Polaren\par
\noindent\emph{Rusky / Russian:} составные поляры\par
\noindent\emph{Ukrajinsky / Ukrainian:} складені поляри\par
\smallskip\noindent\emph{Razlog:} Složene polary spajaje slovjansky pridevnik s tehničnym internacionalnym imenovnikom. Zapis izrično izběgaje težšu tvorbu tipa kompozicijne polary; vybrana kratkost ne skriva, že idut o kombinovane polarne operacije. V tekstě koeficienty razvitja sut polary po dvuh različnyh redah proměnnyh; daljši dokaz uživa rekursivny sistem. Složenost oběkta i računna težkost sut zato različne stvari, iako slovo složeny može v obyčnoj prozě vyvolati obě asociacije.\par
\noindent\emph{Reason:} Složene polary combines a Slavic adjective with an international technical noun. The record expressly avoids a heavier formation such as compositional polars; the shorter wording still denotes combined polar operations. In the text, expansion coefficients are polars with respect to two different variable rows, and the subsequent argument uses a recursive system. Composition of the object and computational difficulty are therefore distinct, although an ordinary-language adjective may evoke both.\par
\noindent{\footnotesize Izvorne zapisy / Original records: HIST-g004-005.}\par
\smallskip\noindent\emph{Sravnjenje s tekstom:} Kompozicija operacij ne jest ocěna težkosti.\par
\noindent\emph{Textual comparison:} Composition of operations is not a difficulty rating.\par
\noindent{\footnotesize TEXT-COMP-1019; pasaže / passages: 422, 424, 435, 437, 439, 441; \texttt{CORPUS\_FORM\_COMPARISONS.json}.}\par
\end{minipage}\par
\par\medskip\Needspace{7\baselineskip}
\noindent\begin{minipage}{\linewidth}
\hypertarget{TERM-1963}{}
\noindent\textbf{vez medžu idealami iz polinomov i modulami iz linearnyh form}\quad{\footnotesize\texttt{TERM-1963}}\par
\noindent\emph{Iztočnik / Source:} Zusammenhang zwischen Idealen aus Polynomen und Moduln aus Linearformen\par
\noindent\emph{Rusky / Russian:} связь между идеалами из полиномов и модулями из линейных форм\par
\noindent\emph{Ukrajinsky / Ukrainian:} зв'язок між ідеалами з поліномів і модулями з лінійних форм\par
\smallskip\noindent\emph{Razlog:} Vez medžu idealami iz polinomov i modulami iz linearnyh form razkladaje dolg německy naslov v vyraznu frazu. Opakovanje iz ukazuje, iz čego sut dva vida oběktov, vměsto zgustiti vse v red pridevnikov. Most jest konkretny: polinom razglědaje se kak linearna forma v stepenovyh produktah s koeficientami v ostalyh proměnnyh. Vsaka taka forma imaje konečno mnogo členov, iako cěly sistem dopušča beskonečno mnogo monomov. Dva arhivna zapisa sut svezane kopije togo izbora, ne dva nezavisne jezyčne svědectva.\par
\noindent\emph{Reason:} Vez medžu idealami iz polinomov i modulami iz linearnyh form unfolds the long German heading into an explicit phrase. Repeating iz identifies what the two kinds of objects consist of, instead of compressing everything into an adjective chain. The bridge is concrete: a polynomial is regarded as a linear form in power products, with coefficients in the remaining variables. Each form has finitely many terms although the whole system admits infinitely many monomials. The two archived records are related copies of this choice, not two independent linguistic attestations.\par
\noindent{\footnotesize Izvorne zapisy / Original records: HIST-g088-001, HIST-g092-032.}\par
\smallskip\noindent\emph{Sravnjenje s tekstom:} Konečnost pojedinoj formy ne znači konečnost vseh monomov.\par
\noindent\emph{Textual comparison:} Finiteness of each form does not mean finiteness of all monomials.\par
\noindent{\footnotesize TEXT-COMP-1020; pasaže / passages: 14615, 14617, 14619, 14621, 14623; \texttt{CORPUS\_FORM\_COMPARISONS.json}.}\par
\end{minipage}\par
\par\medskip\Needspace{7\baselineskip}
\noindent\begin{minipage}{\linewidth}
\hypertarget{TERM-1964}{}
\noindent\textbf{sostavimost iz osnovnyh svijanj}\quad{\footnotesize\texttt{TERM-1964}}\par
\noindent\emph{Iztočnik / Source:} Zusammensetzbarkeit aus Grundfaltungen\par
\noindent\emph{Rusky / Russian:} составимость из основных сверток\par
\noindent\emph{Ukrajinsky / Ukrainian:} складність з основних згортань\par
\smallskip\noindent\emph{Razlog:} Sostavimost iz osnovnyh svijanj znači možnost sestaviti operaciju iz imenovanyh osnovnyh operacij. Ne označaje jeji težkost. Arhiv izrično upozornjaje, že ukrajinsko складність v togo zapisu bylo naměreno v prvom smyslu, i ostavlja vhodnost toj formy otvorenoju. Aktualna rabota piše osnovne kontrakcije i dokazuje dostatočnost n−1 tipov. Sravnenje mora zato čuvati i risk dvosmyslnosti, i novu kontekstnu formu, bez tvrđenja, že otvoreno jezyčno pytanje už bylo rěšeno.\par
\noindent\emph{Reason:} Sostavimost iz osnovnyh svijanj means that an operation can be composed from the specified basic operations. It does not denote difficulty. The archive explicitly warns that Ukrainian складність was intended in the former sense and leaves its suitability open. The current paper uses osnovne kontrakcije and proves the sufficiency of n−1 types. The comparison must therefore retain both the ambiguity risk and the current contextual wording without claiming that the linguistic question has been settled.\par
\noindent{\footnotesize Izvorne zapisy / Original records: HIST-g037-012.}\par
\smallskip\noindent\emph{Sravnjenje s tekstom:} Ukrajinska dvosmyslnost ostaje izrično označena.\par
\noindent\emph{Textual comparison:} The recorded Ukrainian ambiguity remains explicit.\par
\noindent{\footnotesize TEXT-COMP-1021; pasaže / passages: 4651, 4653; \texttt{CORPUS\_FORM\_COMPARISONS.json}.}\par
\end{minipage}\par
\par\medskip\Needspace{7\baselineskip}
\noindent\begin{minipage}{\linewidth}
\hypertarget{TERM-1965}{}
\noindent\textbf{skladanje}\quad{\footnotesize\texttt{TERM-1965}}\par
\noindent\emph{Iztočnik / Source:} Zusammensetzung\par
\noindent\emph{Rusky / Russian:} композиция\par
\noindent\emph{Ukrajinsky / Ukrainian:} композиція\par
\smallskip\noindent\emph{Razlog:} Skladanje bylo izbrano vměsto internacionalnogo kompozicija, ktoro jest zapisano v ruskom i ukrajinskom sravnenju. Motivacija jest slovjanska tvorba soglasna s ranějšim svijanje; aktualny tekst govori o poredku skladanja kontrakcij. V teoremě nezavisnost od poreda jest vyrazno tvrđenje o danyh osnovnyh operacijah, ne svojstvo vsakoj matematičnoj kompozicije. Arhivna prednost domačej tvorbě ne dokazuje, že internacionalna alternativa ne može byti razumljiva ili ispravna.\par
\noindent\emph{Reason:} Skladanje was chosen instead of the international kompozicija recorded in the Russian and Ukrainian comparison. The motivation is a Slavic formation consistent with earlier svijanje; the current text speaks of the order of composing contractions. Independence of order is an explicit theorem about the specified primitive operations, not a property of every mathematical composition. The archived preference for the Slavic formation does not show that the international alternative would be unintelligible or incorrect.\par
\noindent{\footnotesize Izvorne zapisy / Original records: HIST-g003-008.}\par
\smallskip\noindent\emph{Sravnjenje s tekstom:} Nezavisnost poreda jest konkrětno tvrđenje.\par
\noindent\emph{Textual comparison:} Independence of order is a particular assertion.\par
\noindent{\footnotesize TEXT-COMP-1022; pasaže / passages: 278; \texttt{CORPUS\_FORM\_COMPARISONS.json}.}\par
\end{minipage}\par
\par\medskip\Needspace{7\baselineskip}
\noindent\begin{minipage}{\linewidth}
\hypertarget{TERM-1967}{}
\noindent\textbf{dodatne členy $n_i\xi_i$ / dodatne cěločiselne kratne}\quad{\footnotesize\texttt{TERM-1967}}\par
\noindent\emph{Iztočnik / Source:} Zusatzglieder $n_i\xi_i$ / zusätzliche ganze Vielfache\par
\noindent\emph{Rusky / Russian:} дополнительные члены $n_i\xi_i$ / дополнительные целые кратные\par
\noindent\emph{Ukrajinsky / Ukrainian:} додаткові члени $n_i\xi_i$ / додаткові цілі кратні\par
\smallskip\noindent\emph{Razlog:} Dodatne členy n\_i xi\_i označajut dodatne cěločiselne kratne bazovyh elementov, ne nove neograničene koeficienty iz prstena. Oni pojavjajut se pri odstranjenju hipotezy, že prsten imaje jedinicu. Postupok sleduje dva vida idealov: idealy prstena i idealy cělyh čisel; zato čislo prirědženih idealov stava 2k. Dodatny uslovje nenuljevosti pri teoremě o cěpah takože ostaje. Kratka fraza čuva tu korekciju uslovij, a ne samo ukazuje nekake dopisane členy.\par
\noindent\emph{Reason:} Dodatne členy n\_i xi\_i means additional integer multiples of the basis elements, not unrestricted new ring coefficients. They occur when the assumption that the ring has an identity is removed. The construction tracks two kinds of ideals—ring ideals and integer ideals—so there are 2k assigned ideals. The further nonzero condition in the chain statement also remains. The short phrase preserves this adjustment of hypotheses rather than merely pointing to some added terms.\par
\noindent{\footnotesize Izvorne zapisy / Original records: HIST-g109-010.}\par
\smallskip\noindent\emph{Sravnjenje s tekstom:} Dva vida koeficientnyh idealov ne sut odvečne.\par
\noindent\emph{Textual comparison:} The two kinds of coefficient ideals are not redundant.\par
\noindent{\footnotesize TEXT-COMP-1024; pasaže / passages: 16346, 16348, 16350; \texttt{CORPUS\_FORM\_COMPARISONS.json}.}\par
\end{minipage}\par
\par\medskip\Needspace{7\baselineskip}
\noindent\begin{minipage}{\linewidth}
\hypertarget{TERM-1968}{}
\noindent\textbf{dvojnorazova polarizacija}\quad{\footnotesize\texttt{TERM-1968}}\par
\noindent\emph{Iztočnik / Source:} zweimalige Polarisation\par
\noindent\emph{Rusky / Russian:} двукратная поляризация\par
\noindent\emph{Ukrajinsky / Ukrainian:} дворазова поляризація\par
\smallskip\noindent\emph{Razlog:} Dvojnorazova polarizacija kaže, že operacija izvršava se dvakrat. Arhiv sam postavlja kraču alternativu dvojna polarizacija dlja sravnenja, no ne zabilěža dokonanu odluku v njej korist. V tekučem tekstě dlga forma ostaje pri izvedenju razklada iz određene relacije. Važno jest ne poměšati dvakratno priměnjenje s množenjem rezultata čislom dva, ni s dvoma raznymi izbornymi postupkami. Točna modalnost početnogo predloga ostaje provisionalna.\par
\noindent\emph{Reason:} Dvojnorazova polarizacija says that the operation is performed twice. The archive itself proposes the shorter dvojna polarizacija for comparison but records no completed decision in its favour. The current text retains the longer form when deriving a decomposition from a specified relation. Two applications must not be confused with multiplying the result by two, or with two alternative procedures. The original proposal remains provisional.\par
\noindent{\footnotesize Izvorne zapisy / Original records: HIST-g013-010.}\par
\smallskip\noindent\emph{Sravnjenje s tekstom:} Opakovanje operacije ne jest množenje rezultata čislom dva.\par
\noindent\emph{Textual comparison:} Repeating an operation is not multiplying its result by two.\par
\noindent{\footnotesize TEXT-COMP-1025; pasaže / passages: 1938; \texttt{CORPUS\_FORM\_COMPARISONS.json}.}\par
\end{minipage}\par
\par\medskip\Needspace{7\baselineskip}
\noindent\begin{minipage}{\linewidth}
\hypertarget{TERM-1969}{}
\noindent\textbf{dvustranno prosty}\quad{\footnotesize\texttt{TERM-1969}}\par
\noindent\emph{Iztočnik / Source:} zweiseitig einfach\par
\noindent\emph{Rusky / Russian:} двусторонне простой\par
\noindent\emph{Ukrajinsky / Ukrainian:} двобічно простий; двосторонньо простий\par
\smallskip\noindent\emph{Razlog:} Dvustranno prosty zachovuje uslovje prostoty po odnošenju k dvustrannym idealam. V teoremě iz pravostrannoj polnoj reducibilnosti i jedinice slěduje razklad na dvustranno proste komponenty; v pozdějšej sekciji takova komponenta služí osnovoju teorije modulov. Jedinica dělaje kak jediničny operator na razgledanom modulu, čto jest osobna hipoteza. Kratky pridevnik ne směje skryti, že odnostranna prostota, polna reducibilnost i dvustranna prostota sut razne tvrdženja.\par
\noindent\emph{Reason:} Dvustranno prosty retains simplicity with respect to two-sided ideals. One theorem derives decomposition into two-sided simple components from right complete reducibility and an identity; a later section uses such a component as the basis for module theory. The identity acts as the identity operator on the module under consideration, a separate hypothesis. The short adjective must not hide the distinctions between one-sided simplicity, complete reducibility and two-sided simplicity.\par
\noindent{\footnotesize Izvorne zapisy / Original records: HIST-g191-006, HIST-g197-002.}\par
\smallskip\noindent\emph{Sravnjenje s tekstom:} Jedinica i jeji dějstvije ostajut uslovjami.\par
\noindent\emph{Textual comparison:} The identity and its action remain conditions.\par
\noindent{\footnotesize TEXT-COMP-1026; pasaže / passages: 20991, 20993, 20994, 20997, 21022, 21669; \texttt{CORPUS\_FORM\_COMPARISONS.json}.}\par
\end{minipage}\par
\par\medskip\Needspace{7\baselineskip}
\noindent\begin{minipage}{\linewidth}
\hypertarget{TERM-1970}{}
\noindent\textbf{dvustranno proste kolca}\quad{\footnotesize\texttt{TERM-1970}}\par
\noindent\emph{Iztočnik / Source:} zweiseitig einfache Ringe\par
\noindent\emph{Rusky / Russian:} двусторонне простые кольца\par
\noindent\emph{Ukrajinsky / Ukrainian:} двобічно прості кільця\par
\smallskip\noindent\emph{Razlog:} Dvustranno proste kolca jest pluralna fraza za prsteny bez netrivialnyh dvustrannyh idealov, ne drugo ime nekomutativnyh těl. Uvod razděljaje njimi upravljane opče razpadne polja od polj najnižšego stepena, opisanih s pomočju pridruženih těl. Zato vsako skračenje do proste algebry potrěbuje zachovati tot odnos. Aktualny tekst ima prsteny, togda kak arhivni izbor kolca dokumentuje ranějšu terminologiju; slovnik ne pripisuje tekstu formu, ktora v citovanom městu nestojí.\par
\noindent\emph{Reason:} Dvustranno proste kolca is a plural phrase for rings with no nontrivial two-sided ideals, not another name for noncommutative division rings. The introduction distinguishes general splitting fields, described through these rings, from least-degree fields described through the associated division rings. Any shorter gloss such as simple algebra must preserve that relationship. The current text uses prsteny; archived kolca documents the earlier terminology rather than a form falsely attributed to the quoted passage.\par
\noindent{\footnotesize Izvorne zapisy / Original records: HIST-g168-003.}\par
\smallskip\noindent\emph{Sravnjenje s tekstom:} Opče razpadne polja ne sut isto čto polja najnižšego stepena.\par
\noindent\emph{Textual comparison:} General splitting fields are not the same as fields of least degree.\par
\noindent{\footnotesize TEXT-COMP-1027; pasaže / passages: 19244; \texttt{CORPUS\_FORM\_COMPARISONS.json}.}\par
\end{minipage}\par
\par\medskip\Needspace{7\baselineskip}
\noindent\begin{minipage}{\linewidth}
\hypertarget{TERM-1971}{}
\noindent\textbf{dvustranno prosto kolco}\quad{\footnotesize\texttt{TERM-1971}}\par
\noindent\emph{Iztočnik / Source:} zweiseitig einfacher Ring\par
\noindent\emph{Rusky / Russian:} двусторонне простое кольцо\par
\noindent\emph{Ukrajinsky / Ukrainian:} двобічно просте кільце\par
\smallskip\noindent\emph{Razlog:} Dvustranno prosto kolco jest jedinstvena forma istoj strukturnoj nazvy. Izbor kolco byl zabilěžen radi soglasnosti s grupno kolco i sravnenja slovjanskyh form, no nyněšnji tekst upotrěbjaje prsten. Avtorkina konstrukcija razlikuje matričny prsten od pridruženogo těla: matrice nad danym tělom mogut tvoriti dvustranno prosty prsten, ale ne sut vsegda samo tělo. Tu matematičnu razliku čuvamo nezavisno od jezyčnogo izbora kolco ili prsten.\par
\noindent\emph{Reason:} Dvustranno prosto kolco is the singular form of the same structural term. Kolco was recorded for consistency with grupno kolco and comparison of Slavic forms, but the present text uses prsten. The author's construction distinguishes the matrix ring from its associated division ring: matrices over a division ring may form a two-sided simple ring without themselves forming a division ring. That mathematical distinction must survive either lexical choice, kolco or prsten.\par
\noindent{\footnotesize Izvorne zapisy / Original records: HIST-g169-002, HIST-g171-002.}\par
\smallskip\noindent\emph{Sravnjenje s tekstom:} Matrice nad tělom ne morajut samy tvoriti tělo.\par
\noindent\emph{Textual comparison:} Matrices over a division ring need not themselves form a division ring.\par
\noindent{\footnotesize TEXT-COMP-1028; pasaže / passages: 19314, 19352; \texttt{CORPUS\_FORM\_COMPARISONS.json}.}\par
\end{minipage}\par
\par\medskip\Needspace{7\baselineskip}
\noindent\begin{minipage}{\linewidth}
\hypertarget{TERM-1972}{}
\noindent\textbf{dvustranno nerazložime kolca}\quad{\footnotesize\texttt{TERM-1972}}\par
\noindent\emph{Iztočnik / Source:} zweiseitig unzerlegbare Ringe\par
\noindent\emph{Rusky / Russian:} двусторонне неразложимые кольца\par
\noindent\emph{Ukrajinsky / Ukrainian:} двобічно нерозкладні кільця; двосторонньо нерозкладні кільця\par
\smallskip\noindent\emph{Razlog:} Dvustranno nerazložime kolca govori o nemožnosti netrivialnogo direktnogo razklada na dvustranne idealne komponenty. To ne jest sinonim prostyh prstenov i ne opisuje ireducibilnost predstavjenja. Tekst najprvo prenosi teoriju idealov na direktne komponenty, a potom podobno razkladaje modul pomoču idempotentnyh jedinic komponentov. Familija razložiti/nerazložimy zachovuje tu strukturnu operaciju. Arhivne dve ukrajinske formy sut različne jezyčne realizacije, ne dve definicije nerazložimosti.\par
\noindent\emph{Reason:} Dvustranno nerazložime kolca concerns the impossibility of a nontrivial direct decomposition into two-sided ideal components. It is not synonymous with simple rings and does not describe irreducibility of a representation. The text first transfers ideal theory to the direct components and later decomposes modules by using the components' idempotent identities. The razložiti/nerazložimy family retains that structural operation. The two recorded Ukrainian forms are linguistic variants, not two definitions of indecomposability.\par
\noindent{\footnotesize Izvorne zapisy / Original records: HIST-g188-002, HIST-g195-008.}\par
\smallskip\noindent\emph{Sravnjenje s tekstom:} Nerazložimy i prosty ne sut obšče ravnoznačne.\par
\noindent\emph{Textual comparison:} Indecomposable and simple are not generally synonymous.\par
\noindent{\footnotesize TEXT-COMP-1029; pasaže / passages: 20629, 20631, 20636, 21611, 21613, 21617, 21619, 21629, 21631; \texttt{CORPUS\_FORM\_COMPARISONS.json}.}\par
\end{minipage}\par
\par\medskip\Needspace{7\baselineskip}
\noindent\begin{minipage}{\linewidth}
\hypertarget{TERM-1973}{}
\noindent\textbf{dvustranne idealy}\quad{\footnotesize\texttt{TERM-1973}}\par
\noindent\emph{Iztočnik / Source:} zweiseitige Ideale\par
\noindent\emph{Rusky / Russian:} двусторонние идеалы\par
\noindent\emph{Ukrajinsky / Ukrainian:} двосторонні ідеали\par
\smallskip\noindent\emph{Razlog:} Dvustranne idealy jest slovjansky složeny pridevnik, paralelny zapisanym ruskim i ukrajinskim formam. Obe strany sut suštne: pri nekomutativnom prstenu ideal mora byti zatvorjen pod levym i pravym množenjem, da bi kvocient iměl potrěbnu prstenovu strukturu. V drugom městu dvustranne komponenty razkladajut modul s jediničnym operatorom. Nazva čuva toj uslovje, ne tvrdi komutativnost prstena; dvustrannost i komutativnost ne sut ravne stvari.\par
\noindent\emph{Reason:} Dvustranne idealy is a Slavic compound adjective parallel to the recorded Russian and Ukrainian forms. Both sides matter: for a noncommutative ring, closure under left and right multiplication is needed for the relevant quotient-ring structure. Elsewhere, two-sided components decompose a module on which the identity acts as the identity operator. The name preserves that condition without asserting that the ring is commutative. Two-sidedness and commutativity are not equivalent.\par
\noindent{\footnotesize Izvorne zapisy / Original records: HIST-g111-010, HIST-g195-007.}\par
\smallskip\noindent\emph{Sravnjenje s tekstom:} Dvustrannost ne znači komutativnost.\par
\noindent\emph{Textual comparison:} Two-sidedness does not mean commutativity.\par
\noindent{\footnotesize TEXT-COMP-1030; pasaže / passages: 16482, 21611, 21613, 21617, 21619, 21629; \texttt{CORPUS\_FORM\_COMPARISONS.json}.}\par
\end{minipage}\par
\par\medskip\Needspace{7\baselineskip}
\noindent\begin{minipage}{\linewidth}
\hypertarget{TERM-1974}{}
\noindent\textbf{dvustranne klasy}\quad{\footnotesize\texttt{TERM-1974}}\par
\noindent\emph{Iztočnik / Source:} zweiseitige Klassen\par
\noindent\emph{Rusky / Russian:} двусторонние классы\par
\noindent\emph{Ukrajinsky / Ukrainian:} двобічні класи\par
\smallskip\noindent\emph{Razlog:} Dvustranne klasy označajut klasy cěločislenyh matric pri levyh i pravyh unimodularnyh preobrazovanjah. Iako fraza jest analogična dvustrannym idealam, tekst naměrno prehodi od idealov k klasam, da dopušča opče elementarne dělitelje. Idealny slučaj imaje silnějše ograničenja. Zato nazva klasa ne može byti zaměnjena slovom ideal. Zakon razklada klasov potom zachovuje prosty dělitelj, eksponent i město njegova prvogo pojavjenja; to sut tri razne česti informacije.\par
\noindent\emph{Reason:} Dvustranne klasy denotes classes of integer matrices under left and right unimodular transformations. Although the phrase parallels two-sided ideals, the text deliberately passes from ideals to classes in order to allow general elementary divisors; the ideal case is more restrictive. Class therefore cannot be replaced with ideal. The class decomposition subsequently retains a prime divisor, its exponent and the position where that exponent first occurs—three separate pieces of information.\par
\noindent{\footnotesize Izvorne zapisy / Original records: HIST-g081-100.}\par
\smallskip\noindent\emph{Sravnjenje s tekstom:} Prehod od idealov k klasam pozvaljaje opče elementarne dělitelje.\par
\noindent\emph{Textual comparison:} Passing from ideals to classes permits general elementary divisors.\par
\noindent{\footnotesize TEXT-COMP-1031; pasaže / passages: 13707, 13709, 13768; \texttt{CORPUS\_FORM\_COMPARISONS.json}.}\par
\end{minipage}\par
\par\medskip\Needspace{7\baselineskip}
\noindent\begin{minipage}{\linewidth}
\hypertarget{TERM-1975}{}
\noindent\textbf{drugi teorem o izomorfizmu}\quad{\footnotesize\texttt{TERM-1975}}\par
\noindent\emph{Iztočnik / Source:} zweiter Isomorphiesatz\par
\noindent\emph{Rusky / Russian:} вторая теорема об изоморфизме\par
\noindent\emph{Ukrajinsky / Ukrainian:} друга теорема про ізоморфізм\par
\smallskip\noindent\emph{Razlog:} Drugi teorem o izomorfizmu jest izložena fraza s predlogom o vměsto německogo složenogo slova. V konkretnoj priměně ona daje podpolje izomorfno kvocientu osnovnogo prstena po prostomu idealu. Čitatelju treba pokazati sam kvocientny odnos, a ne spolehati samo na numeraciju teoremov, ktora može različati se medžu učebnikami. Arhivno drugi i drugy iz susědnyh zapisov sut jezyčne varianty istogo odkaznogo imena, ne različne teoremy.\par
\noindent\emph{Reason:} Drugi teorem o izomorfizmu expands the German compound into a phrase with o. In this application it gives a subfield isomorphic to a quotient of the base ring by a prime ideal. Readers need the actual quotient relation, not just reliance on theorem numbering, which can differ between textbooks. Archived drugi and drugy in neighbouring records are linguistic variants of the same reference name, not different theorems.\par
\noindent{\footnotesize Izvorne zapisy / Original records: HIST-g151-006.}\par
\smallskip\noindent\emph{Sravnjenje s tekstom:} Imenovany teorem upotrěblja se s određenoju formuloju.\par
\noindent\emph{Textual comparison:} The named theorem is used with a specified formula.\par
\noindent{\footnotesize TEXT-COMP-1032; pasaže / passages: 18633; \texttt{CORPUS\_FORM\_COMPARISONS.json}.}\par
\end{minipage}\par
\par\medskip\Needspace{7\baselineskip}
\noindent\begin{minipage}{\linewidth}
\hypertarget{TERM-1976}{}
\noindent\textbf{drugy teorem izomorfizma}\quad{\footnotesize\texttt{TERM-1976}}\par
\noindent\emph{Iztočnik / Source:} zweiter Isomorphiesatz\par
\noindent\emph{Rusky / Russian:} вторая теорема об изоморфизме\par
\noindent\emph{Ukrajinsky / Ukrainian:} друга теорема про ізоморфізм\par
\smallskip\noindent\emph{Razlog:} Drugy teorem izomorfizma upotrěbjaje genitiv naměsto frazy o izomorfizmu. Arhiv ho bira kako obyčno teoremno ime, ne kak novy algebraičny termin. Podnožna odkazka vodi k ranějšemu objasnjenju izomorfizma; aktualny argument identifikuje podpolje klasov predstavjenyh elementami P. Taj svez čuva smysl imena: govorenje o drugom teoremě ne smí byti zaměnjeno nepodloženym predpoloženjem, že čitatelj už zna avtorčin poredok numeracije.\par
\noindent\emph{Reason:} Drugy teorem izomorfizma uses the genitive instead of o izomorfizmu. The archive chooses it as an ordinary theorem name, not a new algebraic term. The footnote points to the earlier explanation of the isomorphism, while the present argument identifies the subfield of classes represented by elements of P. That connection supplies the name's meaning: calling it the second theorem must not assume that readers already know the author's numbering convention.\par
\noindent{\footnotesize Izvorne zapisy / Original records: HIST-g124-005.}\par
\smallskip\noindent\emph{Sravnjenje s tekstom:} Odkazka vodi k ranějšemu objasnjenju.\par
\noindent\emph{Textual comparison:} The reference leads to the earlier explanation.\par
\noindent{\footnotesize TEXT-COMP-1033; pasaže / passages: 17595, 17597, 17599; \texttt{CORPUS\_FORM\_COMPARISONS.json}.}\par
\end{minipage}\par
\par\medskip\Needspace{7\baselineskip}
\noindent\begin{minipage}{\linewidth}
\hypertarget{TERM-1977}{}
\noindent\textbf{drugy teorem o izomorfizmu}\quad{\footnotesize\texttt{TERM-1977}}\par
\noindent\emph{Iztočnik / Source:} zweiter Isomorphiesatz\par
\noindent\emph{Rusky / Russian:} вторая теорема об изоморфизме\par
\noindent\emph{Ukrajinsky / Ukrainian:} друга теорема про ізоморфізм\par
\smallskip\noindent\emph{Razlog:} Drugy teorem o izomorfizmu spajaje internacionalnu osnovu izomorf- s slovjanskim rednym čisliteljem i padežnoju frazoju. Dve zabilěžene priměnjenja ne davajut dve nazvy: jedno svazuje kvocienty v kompozicijnyh redah modulov, drugo dokazuje prstenovy izomorfizm pri prěseku podprstena s idealom. Iztočna primětka sama objasnjaje razširjenje staršej formulacije. Nazva zato mora ostati s konkretnym izomorfizmom i njegovymi uslovjami, ne samo s rednym čislom.\par
\noindent\emph{Reason:} Drugy teorem o izomorfizmu combines the international isomorf- stem with a Slavic ordinal and case construction. The two recorded uses do not create two names: one relates quotients in module composition series, while the other proves a ring isomorphism involving the intersection of a subring with an ideal. The source footnote itself explains the extension of the earlier formulation. The name must therefore remain linked to the actual isomorphism and its conditions, not merely its ordinal.\par
\noindent{\footnotesize Izvorne zapisy / Original records: HIST-g117-009, HIST-g120-004.}\par
\smallskip\noindent\emph{Sravnjenje s tekstom:} Strukturne uslovja ostajut vidime pri ponovjenoj nazvě.\par
\noindent\emph{Textual comparison:} Structural conditions remain visible under the repeated name.\par
\noindent{\footnotesize TEXT-COMP-1034; pasaže / passages: 17218, 17220, 17222, 17230, 17398; \texttt{CORPUS\_FORM\_COMPARISONS.json}.}\par
\end{minipage}\par
\par\medskip\Needspace{7\baselineskip}
\noindent\begin{minipage}{\linewidth}
\hypertarget{TERM-1978}{}
\noindent\textbf{medžupredstavjenja}\quad{\footnotesize\texttt{TERM-1978}}\par
\noindent\emph{Iztočnik / Source:} Zwischendarstellungen\par
\noindent\emph{Rusky / Russian:} промежуточные представления\par
\noindent\emph{Ukrajinsky / Ukrainian:} проміжні подання\par
\smallskip\noindent\emph{Razlog:} Medžupredstavjenja jest kratka složenica dlja predstavjenij, ktore pojavjajut se pri postupnoj zaměně komponentov v dokazě. Ne ide o matrične predstavjenja grupy, ale o izrazy ideala kako najmanjšego obščego kratnogo. Počet komponentov nosi argument: dve nejednakosti davajut ravnost k=l, a zaměnne formule ne mogu imati menjše čislo komponentov. Izbor medžu- ukazuje međukrok, bez gubljenja oběkta, ktorogo komponenty čislimo.\par
\noindent\emph{Reason:} Medžupredstavjenja is a compact compound for representations occurring during successive replacement of components in a proof. These are not matrix representations of a group but expressions of an ideal as a least common multiple. Component counts drive the argument: two inequalities give k=l, and the intermediate expressions cannot have fewer components. The medžu- formation identifies an intermediate step without losing the object whose components are counted.\par
\noindent{\footnotesize Izvorne zapisy / Original records: HIST-g081-037, HIST-g082-037.}\par
\smallskip\noindent\emph{Sravnjenje s tekstom:} Čislimo idealne komponenty, ne matrične predstavjenja.\par
\noindent\emph{Textual comparison:} The objects counted are ideal components, not matrix representations.\par
\noindent{\footnotesize TEXT-COMP-1035; pasaže / passages: 12716, 12730, 12731, 12735, 12737; \texttt{CORPUS\_FORM\_COMPARISONS.json}.}\par
\end{minipage}\par
\par\medskip\Needspace{7\baselineskip}
\noindent\begin{minipage}{\linewidth}
\hypertarget{TERM-1980}{}
\noindent\textbf{medžupolje}\quad{\footnotesize\texttt{TERM-1980}}\par
\noindent\emph{Iztočnik / Source:} Zwischenkörper\par
\noindent\emph{Rusky / Russian:} промежуточное поле\par
\noindent\emph{Ukrajinsky / Ukrainian:} проміжне поле\par
\smallskip\noindent\emph{Razlog:} Medžupolje jest osnovna složenica s predponoj medžu- i polje. Definicija, izrično svazana so Steinitzom, govori o dvuh vključanjah: polje sodržuje dolnje polje i sodržano jest v gornjem. Ona ne potrěbuje, da oba vključanja budut stroge. Pozdějši dokaz uživa tu poziciju v konečnom algebraičnom razširjenju. Arhivna ocěna prozračnosti ostaje provisionalna; ponovjena upotrěba dokazuje kontinuitet izbora, ne novu nezavisnu jezyčnu attestaciju.\par
\noindent\emph{Reason:} Medžupolje is the headword formed from medžu- and polje. The definition, explicitly connected with Steinitz, gives two inclusions: the field contains the lower field and is contained in the upper one. It does not require both inclusions to be strict. The later proof uses that position inside a finite algebraic extension. The archived transparency judgment remains provisional; repeated use demonstrates continuity of the choice, not a fresh independent linguistic attestation.\par
\noindent{\footnotesize Izvorne zapisy / Original records: HIST-g041-009, HIST-g044-009.}\par
\smallskip\noindent\emph{Sravnjenje s tekstom:} Vključanja ne morajut byti stroge.\par
\noindent\emph{Textual comparison:} The inclusions need not be strict.\par
\noindent{\footnotesize TEXT-COMP-1037; pasaže / passages: 5089, 5091, 5093, 5095, 5468, 5469, 5470; \texttt{CORPUS\_FORM\_COMPARISONS.json}.}\par
\end{minipage}\par
\par\medskip\Needspace{7\baselineskip}
\noindent\begin{minipage}{\linewidth}
\hypertarget{TERM-1981}{}
\noindent\textbf{medžupolje / polje prěsěka}\quad{\footnotesize\texttt{TERM-1981}}\par
\noindent\emph{Iztočnik / Source:} Zwischenkörper / Durchschnittskörper\par
\noindent\emph{Rusky / Russian:} промежуточное поле / поле пересечения\par
\noindent\emph{Ukrajinsky / Ukrainian:} проміжне поле / поле перетину\par
\smallskip\noindent\emph{Razlog:} Medžupolje označaje polje medžu dvoma danymi poljami; polje prěsěka označaje prěsěk dvuh polj v obščem nadpolju. Lemma VIII potrěbuje oba ponjatja: koeficienty minimalnogo polinoma ležet v prěsěku, ktory sam jest medžupoljem. Kratka tvorba medžu- + polje čini odnošenje vidimym, ale ne zaměnjuje prěsěk. Ruske i ukrajinske sravnenja takože razdělajut te dva smisly. Izbor ne jest dokazom, že vsaky čitatelj razuměje novotvor bez definicije.\par
\noindent\emph{Reason:} Medžupolje means a field between two given fields; polje prěsěka means the intersection of two fields in a common overfield. Lemma VIII needs both notions: the coefficients of the minimal polynomial lie in an intersection which is itself an intermediate field. The compact formation medžu- + polje makes the relation visible but does not replace intersection. The Russian and Ukrainian comparanda likewise distinguish the two meanings. This choice is not evidence that every reader will understand the formation without a definition.\par
\noindent{\footnotesize Izvorne zapisy / Original records: HIST-g100-009.}\par
\smallskip\noindent\emph{Sravnjenje s tekstom:} Prěsěk jest medžupoljem v danom argumentu; nazvanja ne sut sinonimy.\par
\noindent\emph{Textual comparison:} The intersection is an intermediate field in this argument; the names are not synonyms.\par
\noindent{\footnotesize TEXT-COMP-1038; pasaže / passages: 15673, 15675; \texttt{CORPUS\_FORM\_COMPARISONS.json}.}\par
\end{minipage}\par
\par\medskip\Needspace{7\baselineskip}
\noindent\begin{minipage}{\linewidth}
\hypertarget{TERM-1982}{}
\noindent\textbf{medžukolco}\quad{\footnotesize\texttt{TERM-1982}}\par
\noindent\emph{Iztočnik / Source:} Zwischenring\par
\noindent\emph{Rusky / Russian:} промежуточное кольцо\par
\noindent\emph{Ukrajinsky / Ukrainian:} проміжне кільце\par
\smallskip\noindent\emph{Razlog:} Arhivny medžukolco i nynějšny medžuprsten imajut jednaku zadaču: nazvati prsten medžu danym prstenom i jego zamknutjem. Razlika kolco/prsten prinadleži obščemu izboru nazvanja ring, ne novoj algebričnoj konstrukciji. V teoremi tvrdženje uključaje vsaky takovy medžuprsten; prevod ne směje oslabiti to do nekotorogo prstena. Analitično vyslovjenje prsten medžu nimi ostaje možno, ale jednoslovna tvorba podpiraje povtorjajuču se paralelu s medžupoljem.\par
\noindent\emph{Reason:} The archived medžukolco and present medžuprsten have the same purpose: to name a ring between a given ring and its closure. The difference kolco/prsten belongs to the general choice of the word for ring, not to a new algebraic construction. The theorem applies to every such intermediate ring; translation must not weaken this to some ring. An analytical phrase, ring between them, remains possible, but the compact formation supports the recurring parallel with intermediate field.\par
\noindent{\footnotesize Izvorne zapisy / Original records: HIST-g127-001, HIST-g129-004.}\par
\smallskip\noindent\emph{Sravnjenje s tekstom:} Nynějšny tekst ima medžuprsten; arhivny izbor ima medžukolco.\par
\noindent\emph{Textual comparison:} The present text has medžuprsten; the archived choice has medžukolco.\par
\noindent{\footnotesize TEXT-COMP-1039; pasaže / passages: 17702, 17794; \texttt{CORPUS\_FORM\_COMPARISONS.json}.}\par
\end{minipage}\par
\par\medskip\Needspace{7\baselineskip}
\noindent\begin{minipage}{\linewidth}
\hypertarget{TERM-1983}{}
\noindent\textbf{ciklična / Dicksonov tip}\quad{\footnotesize\texttt{TERM-1983}}\par
\noindent\emph{Iztočnik / Source:} zyklisch / Dicksonscher Typus\par
\noindent\emph{Rusky / Russian:} циклична / тип Диксона\par
\noindent\emph{Ukrajinsky / Ukrainian:} циклічна / тип Діксона\par
\smallskip\noindent\emph{Razlog:} Ciklična zachovuje internacionalny korenj, vidimy takože v ruskom i ukrajinskom sravnenju. V tom članku on označaje cikličnu algebru, nazvanu takože algebroju Dicksonovogo tipa. Argument prohodi črez ciklično razščěpjajuče polje i ciklično prědstavjenje; ne govori, že multiplikativna grupa vsakoj takovoj algebry jest ciklična. Dicksonovo ime jest matematična atribucija, ne slobodny ukras prevoda. Zapis o mežuslovjanskoj razumlivosti jest motivacija izbora, ne rezultat ispytanja čitateljev.\par
\noindent\emph{Reason:} Ciklična retains the international root also visible in the Russian and Ukrainian comparanda. Here it denotes a cyclic algebra, also called an algebra of Dickson type. The argument proceeds through a cyclic splitting field and a cyclic presentation; it does not say that the multiplicative group of every such algebra is cyclic. Dickson's name is mathematical attribution, not optional decoration. The recorded appeal to cross-Slavic intelligibility is a motivation for the choice, not the result of reader testing.\par
\noindent{\footnotesize Izvorne zapisy / Original records: HIST-g208-003.}\par
\smallskip\noindent\emph{Sravnjenje s tekstom:} Cikličnost algebry ne jest cikličnost jeje multiplikativnoj grupy.\par
\noindent\emph{Textual comparison:} Cyclicity of the algebra is not cyclicity of its multiplicative group.\par
\noindent{\footnotesize TEXT-COMP-1040; pasaže / passages: 22999, 23006, 23019; \texttt{CORPUS\_FORM\_COMPARISONS.json}.}\par
\end{minipage}\par
\par\medskip\Needspace{7\baselineskip}
\noindent\begin{minipage}{\linewidth}
\hypertarget{TERM-1984}{}
\noindent\textbf{ciklične grupy}\quad{\footnotesize\texttt{TERM-1984}}\par
\noindent\emph{Iztočnik / Source:} zyklische Gruppen\par
\noindent\emph{Rusky / Russian:} циклические группы\par
\noindent\emph{Ukrajinsky / Ukrainian:} циклічні групи\par
\smallskip\noindent\emph{Razlog:} Ciklične grupy nazvajut faktory v razloženju konečnoj Abelovoj grupy na direktny produkt. Internacionalny korenj odgovarjaje ruskomu cikličeskij i ukrajinskomu cykličnyj; novy domašny korenj ne byl potrěbny. Matematična granica važna jest: konečna Abelova grupa jest produkt cikličnyh grup, ne nužno sama ciklična. Podslědujuče razsmotrjenje grupovogo prstena ima dodatnu hipotezu o karakteristiki, ktora ne děli red grupy; ta hipoteza ne směje izčeznuti iz izloženja.\par
\noindent\emph{Reason:} Ciklične grupy names the factors in a direct-product decomposition of a finite abelian group. The international root corresponds to the Russian and Ukrainian comparanda; a new native-root coinage was unnecessary. The mathematical distinction matters: a finite abelian group is a product of cyclic groups, not necessarily cyclic itself. The subsequent group-ring discussion has the additional hypothesis that the characteristic does not divide the group order; this must not disappear from the exposition.\par
\noindent{\footnotesize Izvorne zapisy / Original records: HIST-g200-002.}\par
\smallskip\noindent\emph{Sravnjenje s tekstom:} Produkt cikličnyh grup ne mora byti cikličny.\par
\noindent\emph{Textual comparison:} A product of cyclic groups need not be cyclic.\par
\noindent{\footnotesize TEXT-COMP-1041; pasaže / passages: 22000, 22004, 22006, 22008, 22010, 22012, 22021; \texttt{CORPUS\_FORM\_COMPARISONS.json}.}\par
\end{minipage}\par
\par\medskip\Needspace{7\baselineskip}
\noindent\begin{minipage}{\linewidth}
\hypertarget{TERM-1985}{}
\noindent\textbf{ciklična polja $\Omega_n$ stepena $2^n$}\quad{\footnotesize\texttt{TERM-1985}}\par
\noindent\emph{Iztočnik / Source:} zyklische Koerper $\Omega_n$ des Grades $2^n$\par
\noindent\emph{Rusky / Russian:} циклические поля $\Omega_n$ степени $2^n$\par
\noindent\emph{Ukrajinsky / Ukrainian:} циклічні поля $\Omega_n$ степеня $2^n$\par
\smallskip\noindent\emph{Razlog:} Polje, ne tělo v smislu nekomutativnoj algebry děljenja, odpovida tut německomu Körper. Ciklična polja $\Omega_n$ imajut stepen $2^n$ nad racionalnym poljem. V argumentu o kvaternionah uslovje, že $-1$ jest suma dvuh kvadratov, daje razščěpjenje; minimalnost znači, že nijedno vlastno podpolje ne jest razščěpjajuče. Ne znači najmenši možny stepen medžu vsimi razščěpjajučimi poljami. Članok pripisuje dokaz suščestvovanja Hassemu; prevod i slovnik zachovujut tu atribuciju.\par
\noindent\emph{Reason:} Polje, rather than a term for a possibly noncommutative division algebra, corresponds here to Körper. The cyclic fields $\Omega_n$ have degree $2^n$ over the rationals. In the quaternion argument, the condition that $-1$ is a sum of two squares gives splitting; minimality means that no proper subfield is a splitting field. It does not mean the smallest possible degree among all splitting fields. The article credits Hasse with the existence argument, and that attribution belongs in both translation and dictionary.\par
\noindent{\footnotesize Izvorne zapisy / Original records: HIST-g172-001.}\par
\smallskip\noindent\emph{Sravnjenje s tekstom:} Minimalnost jest po vključenju podpolj.\par
\noindent\emph{Textual comparison:} Minimality is with respect to subfield inclusion.\par
\noindent{\footnotesize TEXT-COMP-1042; pasaže / passages: 19486, 19490; \texttt{CORPUS\_FORM\_COMPARISONS.json}.}\par
\end{minipage}\par
\par\medskip\Needspace{7\baselineskip}
\noindent\begin{minipage}{\linewidth}
\hypertarget{TERM-1986}{}
\noindent\textbf{ciklična polja stepena $2^n$}\quad{\footnotesize\texttt{TERM-1986}}\par
\noindent\emph{Iztočnik / Source:} zyklische Körper des Grades $2^n$\par
\noindent\emph{Rusky / Russian:} циклические поля степени $2^n$\par
\noindent\emph{Ukrajinsky / Ukrainian:} циклічні поля степеня $2^n$\par
\smallskip\noindent\emph{Razlog:} Ciklična polja stepena $2^n$ jest vtora kontekstna realizacija isteho nazvanja, ne drugo novo slovo. V §3 dokaz ustanavlja takove minimalne razščěpjajuče polja za beskonečno mnogo značenij $n$; to ne trěba směšati s ranějšim tvrdženjem za vsako $n$. Polje jest položeno v ciklotomično razširenje, ale ciklično i ciklotomično ne sut sinonimy. Nečetny stepen daljnjego razširenja i indeks kvaternionnoj algebry daju konkretny razlog razščěpjenja; samo ime polja ne nosi toj argumentacije.\par
\noindent\emph{Reason:} Ciklična polja stepena $2^n$ is a second contextual use of the same terminology, not another new word. In §3 the argument establishes such minimal splitting fields for infinitely many values of $n$; this must not be conflated with the earlier assertion for every $n$. The field is embedded in a cyclotomic extension, but cyclic and cyclotomic are not synonyms. The odd degree of the further extension and the index of the quaternion algebra provide the particular splitting argument; the name alone does not carry that reasoning.\par
\noindent{\footnotesize Izvorne zapisy / Original records: HIST-g175-001.}\par
\smallskip\noindent\emph{Sravnjenje s tekstom:} Ta nazvanja ne sut sinonimy.\par
\noindent\emph{Textual comparison:} These names are not synonyms.\par
\noindent{\footnotesize TEXT-COMP-1043; pasaže / passages: 19604, 19638; \texttt{CORPUS\_FORM\_COMPARISONS.json}.}\par
\end{minipage}\par
\par\medskip\Needspace{7\baselineskip}
\noindent\begin{minipage}{\linewidth}
\hypertarget{TERM-1987}{}
\noindent\textbf{ciklična zaměna prěměnnyh}\quad{\footnotesize\texttt{TERM-1987}}\par
\noindent\emph{Iztočnik / Source:} zyklische Vertauschung\par
\noindent\emph{Rusky / Russian:} циклическая перестановка переменных\par
\noindent\emph{Ukrajinsky / Ukrainian:} циклічна перестановка змінних\par
\smallskip\noindent\emph{Razlog:} Ciklična zaměna prěměnnyh označaje v priměčanju točno opisan postupok: iz prvogo člena nastajut ješče dva črez cikličnu permutaciju. Zato izraz ne znači proizvoljnu zaměnu ni sumu po vsim šest permutacijam trěh prěměnnyh. Dodavanje prěměnnyh razjasnjaje objekt dějstva, ktory německy izraz može ostaviti podrazuměnym. Avtorkino dodatno tvrdženje, že jednakosti važat i za oddělne članove, ostaje oddělnym matematičnym utvrždženjem.\par
\noindent\emph{Reason:} Ciklična zaměna prěměnnyh denotes the precise operation described in the footnote: two further terms are generated from the first by cyclic permutation. It therefore means neither arbitrary substitution nor summation over all six permutations of three variables. Adding prěměnnyh makes the object of the operation explicit where German can leave it understood. The author's further statement that the equalities hold for individual summands remains a separate mathematical assertion.\par
\noindent{\footnotesize Izvorne zapisy / Original records: HIST-g006-008.}\par
\smallskip\noindent\emph{Sravnjenje s tekstom:} Ciklična zaměna ne vključaje vse permutacije.\par
\noindent\emph{Textual comparison:} Cyclic permutation does not include all permutations.\par
\noindent{\footnotesize TEXT-COMP-1044; pasaže / passages: 773; \texttt{CORPUS\_FORM\_COMPARISONS.json}.}\par
\end{minipage}\par
\par\medskip\Needspace{7\baselineskip}
\noindent\begin{minipage}{\linewidth}
\hypertarget{TERM-1992}{}
\noindent\textbf{ekvivalentne predstavjenja / klas predstavjenj}\quad{\footnotesize\texttt{TERM-1992}}\par
\noindent\emph{Iztočnik / Source:} äquivalente Darstellungen / Darstellungsklasse\par
\noindent\emph{Rusky / Russian:} эквивалентные представления / класс представлений\par
\noindent\emph{Ukrajinsky / Ukrainian:} еквівалентні зображення / клас зображень\par
\smallskip\noindent\emph{Razlog:} Ekvivalentne prědstavjenja prinadležat jednomu klasu, kogda invertibilna zaměna bazy prěvede matricy odnogo v matricy drugogo. Zato ekvivalentny ne znači bukvalno jednaky i prědstavjenje ne znači ilustracija. Nynějšny tekst dava matricy $B$ i $P^{-1}BP$ kako členove istogo klasa. Srědstvo izbora jest soglasnost s klasifikacijeju modulov i predstavjenj; ruske i ukrajinske formy daju sravnenje nazvanij, ne zaměnu matematičnoj definicije.\par
\noindent\emph{Reason:} Equivalent representations belong to one class when an invertible change of basis transforms the matrices of one into those of the other. Ekvivalentny therefore does not mean literally identical, and prědstavjenje does not mean an illustration. The present text gives $B$ and $P^{-1}BP$ as members of the same class. The choice serves the classification of modules and representations; the Russian and Ukrainian forms supply terminological comparison, not a substitute for the mathematical definition.\par
\noindent{\footnotesize Izvorne zapisy / Original records: HIST-g193-005.}\par
\smallskip\noindent\emph{Sravnjenje s tekstom:} Ekvivalentne matricne predstavjenja ne sut bukvalno jednaky zapisy.\par
\noindent\emph{Textual comparison:} Equivalent matrix representations are not literally identical expressions.\par
\noindent{\footnotesize TEXT-COMP-1049; pasaže / passages: 21417, 21419, 21424; \texttt{CORPUS\_FORM\_COMPARISONS.json}.}\par
\end{minipage}\par
\par\medskip\Needspace{7\baselineskip}
\noindent\begin{minipage}{\linewidth}
\hypertarget{TERM-1993}{}
\noindent\textbf{ekvivalentne kvadratne diferencialne formy}\quad{\footnotesize\texttt{TERM-1993}}\par
\noindent\emph{Iztočnik / Source:} äquivalente quadratische Differentialformen\par
\noindent\emph{Rusky / Russian:} эквивалентные квадратичные дифференциальные формы\par
\noindent\emph{Ukrajinsky / Ukrainian:} еквівалентні квадратичні диференціальні форми\par
\smallskip\noindent\emph{Razlog:} Ekvivalentne kvadratne diferencialne formy zachovuje internacionalnu terminologiju i predmet istorijnoj teoremy. Tut diferencialna forma jest kvadratny izraz v diferencialah; ne trěba čitati to kako sovrěmenny termin za alternujuču 2-formu. Kontekst svęzuje ekvivalentnost s linearnymi transformacijami sootvětnyh invariantnyh form i s normalnymi koordinatami. On ne tvrdi, že taka že metoda rěšaje vse obobčenja: v susědnom razsmotrjenju Weylovogo obobčenja ostaje nerěšena zadača. Imena Christoffelja, Riccija, Vermeilja i drugih ostajut čestju naučnoj atribucije.\par
\noindent\emph{Reason:} Ekvivalentne kvadratne diferencialne formy preserves the international terminology and the subject of the historical theorem. Differential form here means a quadratic expression in differentials; it must not be read as the modern term for an alternating 2-form. The context relates equivalence to linear transformations of the relevant invariant forms and to normal coordinates. It does not claim that the same method settles every generalization: the adjacent discussion of Weyl's generalization leaves an equivalence question open. The names of Christoffel, Ricci, Vermeil and others remain part of the scholarly attribution.\par
\noindent{\footnotesize Izvorne zapisy / Original records: HIST-g084-011.}\par
\smallskip\noindent\emph{Sravnjenje s tekstom:} Teoremu ne trěba rasprostranjati na vse obobčenja.\par
\noindent\emph{Textual comparison:} The theorem must not be extended to every generalization.\par
\noindent{\footnotesize TEXT-COMP-1050; pasaže / passages: 14149, 14151; \texttt{CORPUS\_FORM\_COMPARISONS.json}.}\par
\end{minipage}\par
\par\medskip\Needspace{7\baselineskip}
\noindent\begin{minipage}{\linewidth}
\hypertarget{TERM-1994}{}
\noindent\textbf{ekvivalentno četverokratno simbolično predstavjenje}\quad{\footnotesize\texttt{TERM-1994}}\par
\noindent\emph{Iztočnik / Source:} äquivalente vierfache symbolische Darstellung\par
\noindent\emph{Rusky / Russian:} эквивалентное четырехкратное символическое представление\par
\noindent\emph{Ukrajinsky / Ukrainian:} еквівалентне чотириразове символічне подання\par
\smallskip\noindent\emph{Razlog:} Četverokratno prědstavjenje nazvajemo četyrma ekvivalentnymi zapisami v formuli (6). To ne znači umnoženje rezultata čislom četyri ni četyri nezavisne objekty. Srědnje dva zapisy sut skračeni determinantni izrazy; krajny prehod sodrživa znak $(-1)^{\rho(n-\rho)}$. Prěklad mora zachovati taj znak i razliku redov, a ne sojediniti vse zapisy v neodređenu frazu o podobnyh formulah. Simbolično prědstavjenje tut prinadleži simboličnomu metodu, ne teoriji predstavjenj grup.\par
\noindent\emph{Reason:} Fourfold representation names four equivalent expressions in formula (6). It does not mean multiplying the result by four or introducing four independent objects. The middle two expressions abbreviate determinants, and the final transition includes the sign $(-1)^{\rho(n-\rho)}$. Translation must preserve that sign and the distinction between rows instead of collapsing the expressions into an unspecified claim of similar formulas. Symbolic representation here belongs to the symbolic method, not to group representation theory.\par
\noindent{\footnotesize Izvorne zapisy / Original records: HIST-g030-011.}\par
\smallskip\noindent\emph{Sravnjenje s tekstom:} Znak pri zaměně redov ostaje potrěbny.\par
\noindent\emph{Textual comparison:} The sign accompanying the exchange of rows remains necessary.\par
\noindent{\footnotesize TEXT-COMP-1051; pasaže / passages: 3774, 3777, 3778, 3780, 3782; \texttt{CORPUS\_FORM\_COMPARISONS.json}.}\par
\end{minipage}\par
\par\medskip\Needspace{7\baselineskip}
\noindent\begin{minipage}{\linewidth}
\hypertarget{TERM-1995}{}
\noindent\textbf{ekvivalentno razširjeno polje / razširjeno kolco}\quad{\footnotesize\texttt{TERM-1995}}\par
\noindent\emph{Iztočnik / Source:} äquivalenter Erweiterungskörper / Erweiterungsring\par
\noindent\emph{Rusky / Russian:} эквивалентное поле расширения / кольцо расширения\par
\noindent\emph{Ukrajinsky / Ukrainian:} еквівалентне поле розширення / кільце розширення\par
\smallskip\noindent\emph{Razlog:} Ekvivalentno razširjeno polje jest kopija razširenja, svęzana izomorfizmom, ktory ostavlja bazovo polje točkovo neizměnnym. To silnějše uslovje od proizvoljnogo abstraktnogo izomorfizma jest jasno dano v priměčanju. Analogična definicija važi za prsteny. Avtorka dopuščaje vvesti nove simbole, da presěk razširenij byl točno bazovo polje; zato ekvivalentny ne znači bukvalno jednaky objekt. Arhivno kolco i nynějšno prsten trěba svęzati s obščej terminologijeju prstenov, ne s proměnoju etoj definicije.\par
\noindent\emph{Reason:} An equivalent extension field is a copy of the extension linked by an isomorphism which fixes the base field pointwise. This stronger condition than arbitrary abstract isomorphism is explicit in the footnote. The analogous definition applies to rings. The author permits new symbols to be introduced so that the extensions intersect precisely in the base field; equivalent therefore does not mean literally the same object. Archived kolco and present prsten belong to the general ring terminology, not to a change in this definition.\par
\noindent{\footnotesize Izvorne zapisy / Original records: HIST-g119-008.}\par
\smallskip\noindent\emph{Sravnjenje s tekstom:} Ekvivalentne kopije ne sut bukvalno jednaky objekty.\par
\noindent\emph{Textual comparison:} Equivalent copies are not literally identical objects.\par
\noindent{\footnotesize TEXT-COMP-1052; pasaže / passages: 17324; \texttt{CORPUS\_FORM\_COMPARISONS.json}.}\par
\end{minipage}\par
\par\medskip\Needspace{7\baselineskip}
\noindent\begin{minipage}{\linewidth}
\hypertarget{TERM-1996}{}
\noindent\textbf{ekvivalentny sistem}\quad{\footnotesize\texttt{TERM-1996}}\par
\noindent\emph{Iztočnik / Source:} äquivalentes System\par
\noindent\emph{Rusky / Russian:} эквивалентная система\par
\noindent\emph{Ukrajinsky / Ukrainian:} еквівалентна система\par
\smallskip\noindent\emph{Razlog:} Ekvivalentny sistem zachovuje vzaimnu zaměnimost sistemov v konkretnyh pravilah simboličnogo metoda. Izvod sistema (16) iz ranějših identičnostij i obratne specializacije objašnjajut, počemu jedna formulacija može zaměniti drugu. Pozdnějši sistem normalnyh form jest druga realizacija etoj ideje. Ne trěba oslablati ekvivalentny do podobny ili siliti vse upotrěby do bukvalnoj jednakosti množin: važat operacije i odnošenja, dane v každom odlomku.\par
\noindent\emph{Reason:} Ekvivalentny sistem preserves mutual replaceability of systems under the particular rules of the symbolic method. Derivation of system (16) from earlier identities and the reverse specializations explain why one formulation can replace the other. The later system of normal forms is another application of this idea. Equivalent must not be weakened to similar, nor must every occurrence be forced into literal equality of sets: the operations and relations specified in each passage determine the meaning.\par
\noindent{\footnotesize Izvorne zapisy / Original records: HIST-g032-007, HIST-g036-003.}\par
\smallskip\noindent\emph{Sravnjenje s tekstom:} Ekvivalentnost opiraje se na dane operacije.\par
\noindent\emph{Textual comparison:} Equivalence rests on the specified operations.\par
\noindent{\footnotesize TEXT-COMP-1053; pasaže / passages: 3951, 4548, 4628; \texttt{CORPUS\_FORM\_COMPARISONS.json}.}\par
\end{minipage}\par
\par\medskip\Needspace{7\baselineskip}
\noindent\begin{minipage}{\linewidth}
\hypertarget{TERM-1997}{}
\noindent\textbf{ekvivalentnost črez formy vyššego svijanja}\quad{\footnotesize\texttt{TERM-1997}}\par
\noindent\emph{Iztočnik / Source:} Äquivalenz durch Formen höherer Faltung\par
\noindent\emph{Rusky / Russian:} эквивалентность через формы более высокой свертки\par
\noindent\emph{Ukrajinsky / Ukrainian:} еквівалентність через форми вищого згортання\par
\smallskip\noindent\emph{Razlog:} Ekvivalentnost črez formy vyššego svijanja opisuje točno odnošenje: razlika form može byti izražena formami vyššego stupnja kontrakcije. Arhivno svijanje odgovarjaje nynějšnoj kontrakciji, ale uslovje o vyššosti ne jest ukras. Tekst takože govori o boljšem umenšenju težiny; to ne jest isto, čto boljša absoljutna težina. Avtorka ograničaje argument do nastoječih redov prěměnnyh i varuje pred jeho avtomatičnym prěnosom na simbolične redy. Prevod ne směje predstaviti to kako neograničene ignorovanje razlic.\par
\noindent\emph{Reason:} Equivalence through forms of higher contraction describes a specific relation: the difference of the forms can be expressed through forms of higher contraction order. Archived svijanje corresponds to present kontrakcija, but the higher-order condition is essential. The text also speaks of a greater reduction in weight, which is not the same as a greater absolute weight. The author restricts the argument to actual variable rows and warns against automatically transferring it to symbolic rows. The translation must not turn this into unrestricted disregard of differences.\par
\noindent{\footnotesize Izvorne zapisy / Original records: HIST-g037-011.}\par
\smallskip\noindent\emph{Sravnjenje s tekstom:} Bolje umenšenje težiny ne jest boljša absoljutna težina.\par
\noindent\emph{Textual comparison:} A greater weight reduction is not a greater absolute weight.\par
\noindent{\footnotesize TEXT-COMP-1054; pasaže / passages: 4872; \texttt{CORPUS\_FORM\_COMPARISONS.json}.}\par
\end{minipage}\par
\par\medskip\Needspace{7\baselineskip}
\noindent\begin{minipage}{\linewidth}
\hypertarget{TERM-1998}{}
\noindent\textbf{povsudu razložena algebra}\quad{\footnotesize\texttt{TERM-1998}}\par
\noindent\emph{Iztočnik / Source:} überall zerfallende Algebra\par
\noindent\emph{Rusky / Russian:} всюду расщепленная алгебра\par
\noindent\emph{Ukrajinsky / Ukrainian:} всюди розщеплена алгебра\par
\smallskip\noindent\emph{Razlog:} Povsudu razložena algebra prevodi tehničny smysl razščěpjenja vo vsih lokalnyh městah. To ne znači algebra, ktora prosto raspadaje se. V kontekstě centralnoj prostej algebry nad čislovym poljem lokalne algebry imajut trivialny klas; znak ekvivalentnosti s poljem ne znači bukvalnu jednakost algebry s poljem. Lokalno-globalny rezultat jest matematična teorema, ne značenje slova povsudu samo po sebi. Sravnenja rasščeplennaja i rozščeplena pomogajut otděliti tehničny smysl od običnogo razloženja.\par
\noindent\emph{Reason:} Povsudu razložena algebra renders the technical sense of splitting at every local place. It does not mean an algebra which merely falls apart. In the setting of a central simple algebra over a number field, the local algebras have trivial class; the equivalence sign with a field does not assert literal equality of the algebra and field. The local-to-global conclusion is a mathematical theorem, not something implied by everywhere alone. The Russian and Ukrainian splitting terms help distinguish this technical sense from ordinary decomposition.\par
\noindent{\footnotesize Izvorne zapisy / Original records: HIST-g208-004.}\par
\smallskip\noindent\emph{Sravnjenje s tekstom:} Ekvivalentnost s poljem ne znači bukvalnu jednakost.\par
\noindent\emph{Textual comparison:} Equivalence to the field does not mean literal equality.\par
\noindent{\footnotesize TEXT-COMP-1055; pasaže / passages: 23011, 23013, 23015, 23019; \texttt{CORPUS\_FORM\_COMPARISONS.json}.}\par
\end{minipage}\par
\par\medskip\Needspace{7\baselineskip}
\noindent\begin{minipage}{\linewidth}
\hypertarget{TERM-2001}{}
\noindent\textbf{prěhod k konečnym razširjenjam polj}\quad{\footnotesize\texttt{TERM-2001}}\par
\noindent\emph{Iztočnik / Source:} Übergang zu endlichen Erweiterungskörpern\par
\noindent\emph{Rusky / Russian:} переход к конечным расширениям полей\par
\noindent\emph{Ukrajinsky / Ukrainian:} перехід до скінченних розширень полів\par
\smallskip\noindent\emph{Razlog:} Prěhod k konečnym razširjenjam polj daje prozračnu konstrukciju razširjenje polja vměsto doslovnogo německogo složenogo slova. Konečny tut govori o konečnom stepenju razširenja, ne o konečnom čislu elementov polja. V §3 trěba prěnesti dane aksiomy na razširenje prvogo roda i na sootvětne porędky, s jasno imenovanymi izključenjami. Ruske i ukrajinske sravnenja podpirajut etu konstrukciju nazvanja; same po sobě ne dokazujut, že hipotezy možut byti opuščene.\par
\noindent\emph{Reason:} Prěhod k konečnym razširjenjam polj uses the transparent construction field extension instead of a literal German compound. Finite refers to the degree of the extension, not to the number of elements in the field. Section 3 transfers specified axioms to an extension of the first kind and to the relevant orders, with explicit exceptions. The Russian and Ukrainian comparanda support this naming construction; they do not justify omitting the hypotheses.\par
\noindent{\footnotesize Izvorne zapisy / Original records: HIST-g110-001.}\par
\smallskip\noindent\emph{Sravnjenje s tekstom:} Konečny označaje stepen razširenja.\par
\noindent\emph{Textual comparison:} Finite refers to the extension degree.\par
\noindent{\footnotesize TEXT-COMP-1058; pasaže / passages: 16387, 16389, 16391, 16395; \texttt{CORPUS\_FORM\_COMPARISONS.json}.}\par
\end{minipage}\par
\par\medskip\Needspace{7\baselineskip}
\noindent\begin{minipage}{\linewidth}
\hypertarget{TERM-2002}{}
\noindent\textbf{prěhod k razširjenym kolcam}\quad{\footnotesize\texttt{TERM-2002}}\par
\noindent\emph{Iztočnik / Source:} Übergang zu Erweiterungsringen\par
\noindent\emph{Rusky / Russian:} переход к кольцам расширения\par
\noindent\emph{Ukrajinsky / Ukrainian:} перехід до кілець розширення\par
\smallskip\noindent\emph{Razlog:} Razširjene prsteny prodolžajut terminologičnu rodinu razširjenja iz ranějših paragrafov. Nynějšny naslov ima prstenam, ne arhivno kolcam. Konstrukcija prisojedinjaje novo koeficientno prsten i sohranjaje rang v avtorkinyh uslovjah; ne jest ni količnik po idealu ni proizvoljna direktna suma. Prěhod zato označaje proměnu koeficientnoj oblasti. Formulacija mora ostaviti vidimym, po odnošenju k ktorej oblasti sut baza i rang.\par
\noindent\emph{Reason:} Extension rings continue the terminology of extension from the earlier sections. The present heading has prstenam rather than archived kolcam. The construction adjoins a new coefficient ring and preserves rank under the author's conditions; it is neither a quotient by an ideal nor an arbitrary direct sum. Passage therefore denotes change of coefficient domain. The wording must leave clear the domain relative to which basis and rank are considered.\par
\noindent{\footnotesize Izvorne zapisy / Original records: HIST-g140-001.}\par
\smallskip\noindent\emph{Sravnjenje s tekstom:} Prěhod ne jest tvorjenje količnika po idealu.\par
\noindent\emph{Textual comparison:} This passage is not formation of a quotient by an ideal.\par
\noindent{\footnotesize TEXT-COMP-1059; pasaže / passages: 18244, 18246, 18248; \texttt{CORPUS\_FORM\_COMPARISONS.json}.}\par
\end{minipage}\par
\par\medskip\Needspace{7\baselineskip}
\noindent\begin{minipage}{\linewidth}
\hypertarget{TERM-2003}{}
\noindent\textbf{prěhod k kolcu kvocientov}\quad{\footnotesize\texttt{TERM-2003}}\par
\noindent\emph{Iztočnik / Source:} Übergang zum Quotientenring\par
\noindent\emph{Rusky / Russian:} переход к кольцу частных\par
\noindent\emph{Ukrajinsky / Ukrainian:} перехід до кільця часток\par
\smallskip\noindent\emph{Razlog:} Prsten kvocientov tut jest prsten iz drobov s ograničenymi imenovateljami. Imenovatelj mora byti prosty k danomu idealu v točnom idealnom smislu, ktory objašnjaje priměčanje. Zato izraz ne trěba směšati s prstenom klasov ostatkov po idealu. Internacionalny kvocient byl zachovany radi svęzi s poljem kvocientov, medžutym ruske i ukrajinske sravnenja nazvajut droby. To različenje potrěbno jest v slovniku, zato že obšče anglijske quotient ring može zakryti razliku. Grellova atribucija dokazov ostaje zachovana.\par
\noindent\emph{Reason:} The ring of quotients here is a ring of fractions with restricted denominators. A denominator must be prime to the specified ideal in the precise ideal-theoretic sense explained in the footnote. The phrase must therefore not be confused with a residue-class ring modulo an ideal. International kvocient was retained for continuity with field of quotients, while the Russian and Ukrainian comparanda explicitly name fractions. The dictionary needs this distinction because an undifferentiated English quotient ring can obscure it. The attribution of the proofs to Grell is retained.\par
\noindent{\footnotesize Izvorne zapisy / Original records: HIST-g159-001.}\par
\smallskip\noindent\emph{Sravnjenje s tekstom:} Prsten kvocientov ne jest tut prsten klasov ostatkov.\par
\noindent\emph{Textual comparison:} Ring of quotients does not here mean a residue-class ring.\par
\noindent{\footnotesize TEXT-COMP-1060; pasaže / passages: 18932; \texttt{CORPUS\_FORM\_COMPARISONS.json}.}\par
\end{minipage}\par
\par\medskip\Needspace{7\baselineskip}
\noindent\begin{minipage}{\linewidth}
\hypertarget{TERM-2004}{}
\noindent\textbf{prěhodno prěobraženje}\quad{\footnotesize\texttt{TERM-2004}}\par
\noindent\emph{Iztočnik / Source:} Übergangssubstitution\par
\noindent\emph{Rusky / Russian:} переходное преобразование\par
\noindent\emph{Ukrajinsky / Ukrainian:} перехідне перетворення\par
\smallskip\noindent\emph{Razlog:} Prěhodno prěobraženje nazvajemo linearnym prěobraženjem, ktoro vyraža bazu jednogo modula prěz bazu jego osnovnogo modula. Norma v Definiciji III jest determinant etogo prěobraženja. Směr vyraženja baz jest važny: obratno prěobraženje by dalo obratny determinant, kogda on obstaja. Slovo prěobraženje vměsto staršego širokogo substitution bylo izbrano radi jasnogo matematičnogo registra, ne radi proměny operacije.\par
\noindent\emph{Reason:} Prěhodno prěobraženje names the linear transformation expressing a basis of one module in terms of a basis of its underlying module. In Definition III the norm is the determinant of this transformation. The direction matters: the inverse transformation would give the reciprocal determinant when it exists. Transformation was chosen instead of the older broad use of substitution for clarity of mathematical register, not to change the operation.\par
\noindent{\footnotesize Izvorne zapisy / Original records: HIST-g086-004, HIST-g092-014.}\par
\smallskip\noindent\emph{Sravnjenje s tekstom:} Směr vyraženja bazy jest suščestveny.\par
\noindent\emph{Textual comparison:} The direction in which the basis is expressed is essential.\par
\noindent{\footnotesize TEXT-COMP-1061; pasaže / passages: 14404, 14832; \texttt{CORPUS\_FORM\_COMPARISONS.json}.}\par
\end{minipage}\par
\par\medskip\Needspace{7\baselineskip}
\noindent\begin{minipage}{\linewidth}
\hypertarget{TERM-2005}{}
\noindent\textbf{vzete formy}\quad{\footnotesize\texttt{TERM-2005}}\par
\noindent\emph{Iztočnik / Source:} übernommene Formen\par
\noindent\emph{Rusky / Russian:} заимствованные формы\par
\noindent\emph{Ukrajinsky / Ukrainian:} перейняті форми\par
\smallskip\noindent\emph{Razlog:} Vzete formy sut formy prěnesene iz binarnoj oblasti po Clebschovom principu, ne proizvoljno izbrane formy. Nynějšny tekst razjasnjaje tu svęz bez dlugogo povtorjenja: govori o formah vzetyh iz binarnoj oblasti i potom pokazuje sootvětne relacije. Kratky participialny izraz oblegčaje prozu s mnogo formulami, dokolě izvor i sposob prěnosa sut blizko ukazane. Ne označaje to bibliografično zaimstvovanje bez matematičnoj operacije.\par
\noindent\emph{Reason:} Vzete formy are forms transferred from the binary domain by Clebsch's principle, not arbitrarily selected forms. The present text makes the connection explicit without lengthy repetition: it names the binary domain and then displays the corresponding relations. The short participial expression lightens formula-heavy prose provided the origin and method of transfer are stated nearby. It does not denote bibliographical borrowing without a mathematical operation.\par
\noindent{\footnotesize Izvorne zapisy / Original records: HIST-g007-002.}\par
\smallskip\noindent\emph{Sravnjenje s tekstom:} Vzete ne znači proizvoljno izbrane.\par
\noindent\emph{Textual comparison:} Taken over does not mean arbitrarily selected.\par
\noindent{\footnotesize TEXT-COMP-1062; pasaže / passages: 615, 618, 620, 622, 624, 626; \texttt{CORPUS\_FORM\_COMPARISONS.json}.}\par
\end{minipage}\par
\par\medskip\Needspace{7\baselineskip}
\noindent\begin{minipage}{\linewidth}
\hypertarget{TERM-2006}{}
\noindent\textbf{transvekcija}\quad{\footnotesize\texttt{TERM-2006}}\par
\noindent\emph{Iztočnik / Source:} Überschiebung\par
\noindent\emph{Rusky / Russian:} трансвекция / трансвектантная операция\par
\noindent\emph{Ukrajinsky / Ukrainian:} трансвекція / трансвектантна операція\par
\smallskip\noindent\emph{Razlog:} Transvekcija jest internacionalno nazvanje operacije klasičnoj teorije invariantov. Arhiv opiraje izbor na ukrajinski izraz za transvektanty, a ruski sravnjajemy izraz takože ima etoj korenj. Zachovany kratky zapis ne dava tut cěloj istorije čitanja togo izvora; ne trěba dopolnjati ju izmyslenym dokazom konsultacije. V nynějšnom tekstu operacija sojedinjaje relativno polny sistem so sistemom modula. Jeje ne trěba směšati s drugym značenjem transvection kako specialnogo linearnogo prěobraženja.\par
\noindent\emph{Reason:} Transvekcija is the international name for an operation in classical invariant theory. The archive motivates the choice by the Ukrainian word for transvectants, and the Russian comparandum uses the same root. The brief preserved note does not supply a complete history of consultation of that source; it must not be supplemented with invented evidence of reading. In the present text the operation combines a relatively complete system with the system of the modulus. It must not be confused with the other meaning of transvection as a particular linear transformation.\par
\noindent{\footnotesize Izvorne zapisy / Original records: HIST-g001-007.}\par
\smallskip\noindent\emph{Sravnjenje s tekstom:} Tut ide o invariantnu operaciju nad formami.\par
\noindent\emph{Textual comparison:} This is an invariant-theoretic operation on forms.\par
\noindent{\footnotesize TEXT-COMP-1063; pasaže / passages: 71, 182; \texttt{CORPUS\_FORM\_COMPARISONS.json}.}\par
\end{minipage}\par
\par\medskip\Needspace{7\baselineskip}
\noindent\begin{minipage}{\linewidth}
\hypertarget{TERM-2007}{}
\noindent\textbf{transvekcija črez funkcionalny determinant}\quad{\footnotesize\texttt{TERM-2007}}\par
\noindent\emph{Iztočnik / Source:} Überschiebung über eine Funktionaldeterminante\par
\noindent\emph{Rusky / Russian:} трансвекция через функциональный детерминант\par
\noindent\emph{Ukrajinsky / Ukrainian:} трансвекція через функціональний детермінант\par
\smallskip\noindent\emph{Razlog:} Transvekcija črez funkcionalny determinant zachovuje jedinstveno čislo determinant, ktoro jest važno v tom odlomku. Nynějšna forma predloga jest prěz; smysl ostaje operacija, iz ktorej slěduje razklad formy. To ne jest instrukcija poděliti čislo determinantom. Nazvanje operacije jest isto kako v množinnyh kontekstah, ale jedinstveno čislo ne směje byti avtomatično normalizovano do množiny.\par
\noindent\emph{Reason:} Transvekcija črez funkcionalny determinant preserves the singular determinant required in this passage. The present preposition is prěz; the meaning remains the operation through which the form decomposes. It is not an instruction to divide a number by a determinant. The operation has the same name as in the plural contexts, but the singular must not be automatically normalized to plural.\par
\noindent{\footnotesize Izvorne zapisy / Original records: HIST-g023-009.}\par
\smallskip\noindent\emph{Sravnjenje s tekstom:} Jedinstveno čislo ostaje zachovano.\par
\noindent\emph{Textual comparison:} The singular is preserved.\par
\noindent{\footnotesize TEXT-COMP-1064; pasaže / passages: 2982; \texttt{CORPUS\_FORM\_COMPARISONS.json}.}\par
\end{minipage}\par
\par\medskip\Needspace{7\baselineskip}
\noindent\begin{minipage}{\linewidth}
\hypertarget{TERM-2008}{}
\noindent\textbf{transvekcija črez funkcionalne determinanty}\quad{\footnotesize\texttt{TERM-2008}}\par
\noindent\emph{Iztočnik / Source:} Überschiebung über Funktionaldeterminanten\par
\noindent\emph{Rusky / Russian:} трансвекция через функциональные детерминанты\par
\noindent\emph{Ukrajinsky / Ukrainian:} трансвекція через функціональні детермінанти\par
\smallskip\noindent\emph{Razlog:} Transvekcija črez funkcionalne determinanty daje internacionalno nazvanje vměsto doslovnoj metafory posuvanja. Množina determinanty odgovarjaje dvoma zachovanyma kontekstama, dokolě inoj odlomok potrěbuje jedinstveno čislo. Arhivny črez i nynějšny prěz sut varijanty izraženja odnošenja, ne dva vida transvekcije. Slovnik svęzuje jih s edinym ponjatjem, ale ne izbrišaje razliku konkretnyh formul i razkladov.\par
\noindent\emph{Reason:} Transvekcija črez funkcionalne determinanty uses the international name instead of a literal sliding metaphor. The plural determinants corresponds to the two recorded contexts, whereas another passage requires singular. Archived črez and present prěz are variants in expressing the relation, not two kinds of transvection. The dictionary links them to one concept without erasing the distinctions between the particular formulas and decompositions.\par
\noindent{\footnotesize Izvorne zapisy / Original records: HIST-g022-003, HIST-g026-006.}\par
\smallskip\noindent\emph{Sravnjenje s tekstom:} Predložne varijanty ne nazvajut različne operacije.\par
\noindent\emph{Textual comparison:} The prepositional variants do not name different operations.\par
\noindent{\footnotesize TEXT-COMP-1065; pasaže / passages: 2837, 2856, 3260; \texttt{CORPUS\_FORM\_COMPARISONS.json}.}\par
\end{minipage}\par
\par\medskip\Needspace{7\baselineskip}
\noindent\begin{minipage}{\linewidth}
\hypertarget{TERM-2009}{}
\noindent\textbf{prěkrivanja}\quad{\footnotesize\texttt{TERM-2009}}\par
\noindent\emph{Iztočnik / Source:} Überschiebungen\par
\noindent\emph{Rusky / Russian:} перекрывания\par
\noindent\emph{Ukrajinsky / Ukrainian:} перекривання\par
\smallskip\noindent\emph{Razlog:} Prěkrivanja jest ranějši doslovny izbor dlja Überschiebungen v skobkah pored polary. Arhiv sam ostavil jego otkrytym k povtornomu razsmotrjenju, ako daljnje česti potrěbujut specialističny termin transvekcija. Takova situacija dejstvitelno jest v korpusu: nynějšny tekst sodrživa oba nazvanja. Zato etu statju ne trěba čitati kako bezuslovno priporočanje dvuh sinonimov. Obyčno prěkrivanje može skryti konkretnu invariantnu operaciju; unifikacija potrěbuje sravnenje operacij v každym kontekstě, a ne slepu globalnu zaměnu.\par
\noindent\emph{Reason:} Prěkrivanja is the earlier literal choice for Überschiebungen, parenthesized beside polars. The archive itself left it open to reconsideration if later sections required the specialist term transvekcija. That situation does occur in the corpus: the present text contains both names. This entry therefore does not unconditionally recommend two synonyms. Ordinary overlapping can obscure the particular invariant operation; standardization requires comparison of the operations in context, not blind global replacement.\par
\noindent{\footnotesize Izvorne zapisy / Original records: HIST-g005-007.}\par
\smallskip\noindent\emph{Sravnjenje s tekstom:} Arhivny izbor jest otkryty k smyslovomu sravnenju s transvekcijeju.\par
\noindent\emph{Textual comparison:} The archived choice remains subject to sense comparison with transvekcija.\par
\noindent{\footnotesize TEXT-COMP-1066; pasaže / passages: 567, 583, 585, 591, 773; \texttt{CORPUS\_FORM\_COMPARISONS.json}.}\par
\end{minipage}\par
\par\medskip\Needspace{7\baselineskip}
\noindent\begin{minipage}{\linewidth}
\hypertarget{TERM-2010}{}
\noindent\textbf{transvekcije nad funkcionalnymi determinantami}\quad{\footnotesize\texttt{TERM-2010}}\par
\noindent\emph{Iztočnik / Source:} Überschiebungen über Funktionaldeterminanten\par
\noindent\emph{Rusky / Russian:} трансвекции над функциональными детерминантами\par
\noindent\emph{Ukrajinsky / Ukrainian:} трансвекції над функціональними детермінантами\par
\smallskip\noindent\emph{Razlog:} Transvekcije nad funkcionalnymi determinantami jest množinna konstrukcija s predlogom nad. Nynějšny odlomok vyličaje konkretne formy, ktore se tak razkladajut. Nad tut ne označaje prostorno vyše: on svęzuje nazvanje operacije s objektami, nad ktorymi ona jest izvršena. Zapisana dvojnost s prěkrivanje i prěnos ne dokazuje, že vse tri slova sut zaměnime. Prěnos može nazvati cěly princip prěhoda medžu sistemami, dokolě transvekcija nazvajet konkretnu operaciju nad formami.\par
\noindent\emph{Reason:} Transvekcije nad funkcionalnymi determinantami is a plural construction using nad. The present passage lists particular forms decomposing in this way. Nad does not mean spatially above here: it relates the operation to its objects. The recorded competition with prěkrivanje and prěnos does not show that all three words are interchangeable. Prěnos can name a whole principle for passing between systems, whereas transvekcija names a particular operation on forms.\par
\noindent{\footnotesize Izvorne zapisy / Original records: HIST-g015-010.}\par
\smallskip\noindent\emph{Sravnjenje s tekstom:} Predlog nad nije prostorny tut.\par
\noindent\emph{Textual comparison:} Nad is not spatial here.\par
\noindent{\footnotesize TEXT-COMP-1067; pasaže / passages: 2184; \texttt{CORPUS\_FORM\_COMPARISONS.json}.}\par
\end{minipage}\par
\par\medskip\Needspace{7\baselineskip}
\noindent\begin{minipage}{\linewidth}
\hypertarget{TERM-2012}{}
\noindent\textbf{prěnos}\quad{\footnotesize\texttt{TERM-2012}}\par
\noindent\emph{Iztočnik / Source:} Übertragung\par
\noindent\emph{Rusky / Russian:} перенос\par
\noindent\emph{Ukrajinsky / Ukrainian:} перенесення\par
\smallskip\noindent\emph{Razlog:} Prěnos v odlomku o Weylu i Schoutenu prinadleži diferencialnoj geometriji: on označaje geometrično odnošenje, na ktorom možno osnovati prěobraženje koeficientov i teoriju invariantov. Ne jest obična transmisija ni Clebschov princip iz teorije binarnyh form. Isto slovo zato potrěbuje oddělny smysl v slovniku, a ne avtomatično sjedinenje vsih pojavjenij. Avtorka formulovala obče pitanje i navela čestične rezultaty; prevod ne směje proměniti pitanje v dokazanu obšču teoremu.\par
\noindent\emph{Reason:} Prěnos in the passage about Weyl and Schouten belongs to differential geometry: it denotes the geometrical structure used to define transformation of coefficients and an invariant theory. It is neither ordinary transmission nor Clebsch's principle from binary invariant theory. The same word consequently needs a separate dictionary sense, not automatic merging of all occurrences. The author poses a general question and cites partial results; translation must not turn that question into a proved general theorem.\par
\noindent{\footnotesize Izvorne zapisy / Original records: HIST-g094-017.}\par
\smallskip\noindent\emph{Sravnjenje s tekstom:} Obče pitanje ostaje pitanjem, ne dokazanoju teoremoju.\par
\noindent\emph{Textual comparison:} The general question remains a question, not a proved theorem.\par
\noindent{\footnotesize TEXT-COMP-1069; pasaže / passages: 15197, 15205; \texttt{CORPUS\_FORM\_COMPARISONS.json}.}\par
\end{minipage}\par
\par\medskip\Needspace{7\baselineskip}
\noindent\begin{minipage}{\linewidth}
\hypertarget{TERM-2013}{}
\noindent\textbf{prěnos iz binarnoj oblasti}\quad{\footnotesize\texttt{TERM-2013}}\par
\noindent\emph{Iztočnik / Source:} Übertragung aus dem binären Gebiet\par
\noindent\emph{Rusky / Russian:} перенос из бинарной области\par
\noindent\emph{Ukrajinsky / Ukrainian:} перенесення з бінарної області\par
\smallskip\noindent\emph{Razlog:} Prěnos iz binarnoj oblasti izrično nazvajet izvor konstrukcije. Po Clebschovom principu binarna relacija črez obrubljenje daje ternarnu relaciju i relativnu polnotu modulo determinant. Binarna oblast tut jest oblast teorije form v dvuh prěměnnyh, ne računalno dvojično kodovanje. Slovo prěnos sohranjaje svęz s principom i so vzętymi formami, ale ne znači, že formuly ostajut bukvalno neizměnnymi.\par
\noindent\emph{Reason:} Prěnos iz binarnoj oblasti explicitly names the source of the construction. By Clebsch's principle, bordering turns a binary relation into a ternary one and yields relative completeness modulo a determinant. Binary domain here means the theory of forms in two variables, not binary computer encoding. Prěnos preserves the link with the principle and the transferred forms, but does not imply that the formulas remain literally unchanged.\par
\noindent{\footnotesize Izvorne zapisy / Original records: HIST-g007-009.}\par
\smallskip\noindent\emph{Sravnjenje s tekstom:} Binarno znači s dvoma prěměnnymi.\par
\noindent\emph{Textual comparison:} Binary means involving two variables.\par
\noindent{\footnotesize TEXT-COMP-1070; pasaže / passages: 615, 618, 620, 622, 624, 626; \texttt{CORPUS\_FORM\_COMPARISONS.json}.}\par
\end{minipage}\par
\par\medskip\Needspace{7\baselineskip}
\noindent\begin{minipage}{\linewidth}
\hypertarget{TERM-2014}{}
\noindent\textbf{načelo prěnosa}\quad{\footnotesize\texttt{TERM-2014}}\par
\noindent\emph{Iztočnik / Source:} Übertragungsprinzip\par
\noindent\emph{Rusky / Russian:} принцип переноса\par
\noindent\emph{Ukrajinsky / Ukrainian:} принцип перенесення\par
\smallskip\noindent\emph{Razlog:} Načelo prěnosa jest domašny variant nazvanja princip prěnosa. V tom odlomku on svęzuje teoremy o modulah s teoremami o idealah; ne označaje geometričny prěnos Weyla ni operaciju transvekcije. Arhiv izbral načelo kako obšče slovo dlja principa. Nynějšny tekst zachovuje toj variant, zato slovnik mora uvesti oba nazvanja i jih smyslove svęzi, vměsto pretenzije, že od počętka suščestvoval samo jeden termin.\par
\noindent\emph{Reason:} Načelo prěnosa is a native-word alternative to princip prěnosa. In this passage it connects theorems about modules with theorems about ideals; it does not denote Weyl's geometric transport or the operation of transvection. The archive chose načelo as a general word for principle. The present text retains that variant, so the dictionary must include both names and their sense relations instead of pretending that only one term existed from the outset.\par
\noindent{\footnotesize Izvorne zapisy / Original records: HIST-g089-009, HIST-g092-051.}\par
\smallskip\noindent\emph{Sravnjenje s tekstom:} Načelo i princip sut dva zachovane nazvanja.\par
\noindent\emph{Textual comparison:} Načelo and princip are two preserved names.\par
\noindent{\footnotesize TEXT-COMP-1071; pasaže / passages: 14910, 14913; \texttt{CORPUS\_FORM\_COMPARISONS.json}.}\par
\end{minipage}\par
\par\medskip\Needspace{7\baselineskip}
\noindent\begin{minipage}{\linewidth}
\hypertarget{TERM-2015}{}
\noindent\textbf{princip prěnosa}\quad{\footnotesize\texttt{TERM-2015}}\par
\noindent\emph{Iztočnik / Source:} Übertragungsprinzip\par
\noindent\emph{Rusky / Russian:} принцип переноса\par
\noindent\emph{Ukrajinsky / Ukrainian:} принцип перенесення\par
\smallskip\noindent\emph{Razlog:} Princip prěnosa sojedinjaje internacionalny princip s prozračnym slovjanskim prěnos. Ranějša motivacija nazyvala konstrukciju domašnoj, ale to ne jest točno opisanje proishođenja vsih jeje čestij. Pozdnějše zapisy podčrknuli kontinuitet nazvanja. V teoriji baz princip dopuščaje prěnos dokazov iz odobražajučego sistema k početkovomu; v teoriji form ima drugu konkretnu realizaciju. Jedino nazvanje svęzuje podobnu logičnu funkciju, ne ustanavlja jedinu teoremu za vse upotrěby.\par
\noindent\emph{Reason:} Princip prěnosa combines international princip with transparent Slavic prěnos. An early rationale called the construction native, but that does not accurately describe the origin of all its components. Later records emphasize continuity of the name. In the theory of bases the principle transfers results from a representing system to the original one; in invariant theory it has a different particular realization. One name links a similar logical function without establishing a single theorem governing every use.\par
\noindent{\footnotesize Izvorne zapisy / Original records: HIST-g004-007, HIST-g040-015, HIST-g043-011, HIST-g046-014, HIST-g048-016.}\par
\smallskip\noindent\emph{Sravnjenje s tekstom:} Jedino nazvanje ne znači jedinu neograničenu teoremu.\par
\noindent\emph{Textual comparison:} One name does not imply a single unrestricted theorem.\par
\noindent{\footnotesize TEXT-COMP-1072; pasaže / passages: 5029, 5414, 5435, 5437; \texttt{CORPUS\_FORM\_COMPARISONS.json}.}\par
\end{minipage}\par
\par\medskip\Needspace{7\baselineskip}
\noindent\begin{minipage}{\linewidth}
\hypertarget{TERM-2016}{}
\noindent\textbf{princip prěnosa za racionalnu bazu}\quad{\footnotesize\texttt{TERM-2016}}\par
\noindent\emph{Iztočnik / Source:} Übertragungsprinzip für die Rationalbasis\par
\noindent\emph{Rusky / Russian:} принцип переноса для рациональной базы\par
\noindent\emph{Ukrajinsky / Ukrainian:} принцип перенесення для раціональної бази\par
\smallskip\noindent\emph{Razlog:} Princip prěnosa za racionalnu bazu jest konkretizacija predhodnogo nazvanja. Racionalna baza tut generuje funkcije racionalnymi vyraženjami, ne linearnoju kombinacijeju i ne nužno s racionalnymi čislami kako koeficientami. Avtorkin izvod prěnosi relacije za vse funkcije od odobražajučego sistema k početkovomu. Zato dopolnjenje za racionalnu bazu jest suščestveno dlja smisla, a ne izbytočny ukras. Ono ne označaje novy nezavisny princip bez svęzi s §3.\par
\noindent\emph{Reason:} Princip prěnosa za racionalnu bazu specializes the preceding name. A rational basis here generates functions by rational expressions, not by linear combinations and not necessarily with rational numbers as coefficients. The author's argument transfers the relations for all functions from the representing system to the original one. For a rational basis is therefore essential to the meaning, not ornamental wording. It does not designate a new independent principle unrelated to §3.\par
\noindent{\footnotesize Izvorne zapisy / Original records: HIST-g044-003.}\par
\smallskip\noindent\emph{Sravnjenje s tekstom:} Racionalno generovanje ne jest linearna baza.\par
\noindent\emph{Textual comparison:} Rational generation is not a linear basis.\par
\noindent{\footnotesize TEXT-COMP-1073; pasaže / passages: 5414, 5416, 5418, 5420, 5422, 5425, 5427, 5429, 5431, 5433, 5435, 5437; \texttt{CORPUS\_FORM\_COMPARISONS.json}.}\par
\end{minipage}\par
\clearpage\normalsize\section{Pisanje i redakcija / Writing and editorial conventions}
Tute odluke upravljajut objasnjenjem, oznakami i podnožnymi primětkami. One ne predstavljajut nove matematične termine.
\par\smallskip\noindent\emph{EN:} These decisions concern exposition, notation, and footnotes. They are not proposals for new mathematical terms.
\small
\par\medskip\Needspace{7\baselineskip}
\noindent\begin{minipage}{\linewidth}
\hypertarget{TERM-0219}{}
\noindent\textbf{Ponovjene podnožne primětky}\par\noindent\textbf{Repeated footnotes}\quad{\footnotesize\texttt{TERM-0219}}\par
\smallskip\noindent\emph{Razlog:} Tuto jest redakcijna odluka o podnožnoj primětkě, ne novy matematičny termin. Izvorni razlog govori, že kratka primětka nebyla ponovjena, ibo dolga pečatna forma byla prěvedena. Da bi ta odluka byla opodstatnjena, trěba sravniti obě primětky i ustanoviti, že nikakvo sodržanje ne jest izgubjeno. Sam zapis razloga ne jest dokazom takvoj polnoty.\par
\noindent\emph{Reason:} This is an editorial decision about a footnote, not a new mathematical term. The recorded reason says that the short note was not repeated because the longer printed form had been translated. Justifying that decision requires comparison of both notes and confirmation that no content was lost. The rationale alone does not establish that completeness.\par
\noindent{\footnotesize Izvorne zapisy / Original records: HIST-g154-008.}\par
\end{minipage}\par
\par\medskip\Needspace{7\baselineskip}
\noindent\begin{minipage}{\linewidth}
\hypertarget{TERM-0221}{}
\noindent\textbf{Matematične oznaky v kiriličnom tekstu}\par\noindent\textbf{Mathematical notation in Cyrillic prose}\quad{\footnotesize\texttt{TERM-0221}}\par
\smallskip\noindent\emph{Razlog:} Proměnna v formulě i rimsko čislo III ostavajut neproměnjene pri prěpisu prozy v kirilicu. Razlog jest njih identifikacijska funkcija, ne jezyk okolnoj věty. Priměr iz odděla 15 pokazuje, že ochrana oznak jest nužna i bez osobnoj formuly kongruencije; skript ne može odlučovati samo po prisutnosti jednoj TeX-komandy.\par
\noindent\emph{Reason:} A variable in a formula and the Roman numeral III remain unchanged when the prose is converted to Cyrillic. Their identifying function, rather than the surrounding sentence's language, governs the choice. The example from section 15 shows that notation needs protection even without a special congruence expression; conversion cannot be governed by the presence of one TeX command alone.\par
\noindent{\footnotesize Izvorne zapisy / Original records: HIST-g016-014.}\par
\end{minipage}\par
\par\medskip\Needspace{7\baselineskip}
\noindent\begin{minipage}{\linewidth}
\hypertarget{TERM-0242}{}
\noindent\textbf{Polna forma objasnjujučej skračenosti}\par\noindent\textbf{Expanding an explanatory abbreviation}\quad{\footnotesize\texttt{TERM-0242}}\par
\smallskip\noindent\emph{Razlog:} To jest razkriva skračenost t. j., da by čitatelj ne trěbal razpoznati osobnu skračenicu. Ta čitatelska prednost priměnjaje se i v kiriličnom tekstu. Izraz označaje objasnjenje ili preformulovanje; drugy zapisny izbor za zaključok, napriměr slědovatelno, trěbuje kontekst, a ne avtomatičnu zaměnu vsakogo d.h.\par
\noindent\emph{Reason:} To jest expands t. j. so that readers need not recognize a particular abbreviation. That reading advantage also applies in Cyrillic. The expression introduces explanation or reformulation; an inferential expression such as therefore requires contextual justification, not automatic replacement of every d.h.\par
\noindent{\footnotesize Izvorne zapisy / Original records: HIST-g099-013.}\par
\end{minipage}\par
\par\medskip\Needspace{7\baselineskip}
\noindent\begin{minipage}{\linewidth}
\hypertarget{TERM-0370}{}
\noindent\textbf{Inicial autora i TeX-komanda}\par\noindent\textbf{An author's initial and a TeX command}\quad{\footnotesize\texttt{TERM-0370}}\par
\smallskip\noindent\emph{Razlog:} Ime autora i njegov inicial morajut ostati točno identifikovane nezavisno od TeX-komand. Zapis govori, že komandny zapis mogl pokazati bukvu Ø vměsto iniciala O. pri imenu Ore, a tekst iztočnoj stranicy davaje O. Ore. To jest redakcijno pytanje identifikacije, ne novy termin ili proměna avtorskogo imena. Izvlečeny tekst jest navedeny dokaz; bez novogo vizualnogo sravnjenja ne trěba pripisyvati tutomu objasnjenju novu prověrku pečatnoj stranicy.\par
\noindent\emph{Reason:} The author's name and initial must identify the person correctly regardless of TeX commands. The record says that a command could display Ø instead of initial O. in Ore's name, while text extracted from the source page gives O. Ore. This concerns editorial identification, not a new term or a change of author name. Extracted text is the evidence reported; this explanation does not claim a fresh visual comparison with the printed page.\par
\noindent{\footnotesize Izvorne zapisy / Original records: HIST-g215-014.}\par
\end{minipage}\par
\par\medskip\Needspace{7\baselineskip}
\noindent\begin{minipage}{\linewidth}
\hypertarget{TERM-0403}{}
\noindent\textbf{Latinično skračenje časopisa v kiriličnoj proze}\par\noindent\textbf{A journal abbreviation in Cyrillic prose}\quad{\footnotesize\texttt{TERM-0403}}\par
\smallskip\noindent\emph{Razlog:} V rečenici «Podrobno predstavjenje imaje pojaviti se v Math. Ann.» skračenje Math. Ann. ostavaje v latiničnom zapisu i pri kiriličnom prěvodu okolnoj prozy. Razlog jest točna bibliografična identifikacija, ne neprevodimost cěloj rečenicy. Německo «Eine ausführliche Darstellung soll in den Math. Ann. erscheinen.» najavjaje buduču publikaciju; novy slovničny opis ne može pretvoriti tu historijnu najavu v svoje potvrđanje, že publikacija už vyšla. Zapis argumentuje ohranu skračenja, ne posebny izbor vsakego glagola.\par
\noindent\emph{Reason:} In «Podrobno predstavjenje imaje pojaviti se v Math. Ann.», Math. Ann. remains in Latin letters when the surrounding prose is converted to Cyrillic. The reason is precise bibliographical identification, not that the whole sentence is untranslatable. German «Eine ausführliche Darstellung soll in den Math. Ann. erscheinen.» announces a prospective publication; the dictionary must not turn that historical announcement into its own confirmation that publication occurred. The record argues for preserving the abbreviation, not separately for every verb choice.\par
\noindent{\footnotesize Izvorne zapisy / Original records: HIST-g105-012.}\par
\end{minipage}\par
\par\medskip\Needspace{7\baselineskip}
\noindent\begin{minipage}{\linewidth}
\hypertarget{TERM-0534}{}
\noindent\textbf{Prěhod k diskriminantnym idealam primarnyh prstenov}\par\noindent\textbf{The transition to discriminant ideals of primary rings}\quad{\footnotesize\texttt{TERM-0534}}\par
\smallskip\noindent\emph{Razlog:} Německa rečenica «Es sollen daher vorerst Diskriminantenideale primärer Ringe betrachtet werden.» ima zabilěženy prěvod «Zato najprvo treba razsmotriti diskriminantne idealy primarnyh kolc.» Razlog govori, že rečenica byla najdena v skanovanom iztočniku, ale otsutna iz čistogo TeX-teksta; uključena jest, ibo vyražaje logičny prěhod od kriterija k točkě 1. Tuto jest odluka o polnotě teksta, ne novy slovničny termin. Zapis jest svědectvom o tom sravnjenju; sam ne zaměnjaje novu prověrku pečatnoj stranicy.\par
\noindent\emph{Reason:} The German sentence «Es sollen daher vorerst Diskriminantenideale primärer Ringe betrachtet werden.» has the recorded translation «Zato najprvo treba razsmotriti diskriminantne idealy primarnyh kolc.» The rationale says that the sentence was found in the scan but absent from the clean TeX text, and included because it states the logical transition from the criterion to item 1. This is a decision about textual completeness, not a new dictionary term. The record is evidence of the reported comparison; it does not replace a fresh check of the printed page.\par
\noindent{\footnotesize Izvorne zapisy / Original records: HIST-g148-006.}\par
\smallskip\noindent\emph{Sravnjenje s tekstom:} Aktualny početok odděla 6 raboty 31 ima rečenicu Zato najprvo trěba razgleděti diskriminantne idealy primarnyh prstenov. Ona sohranjaje logičny prěhod opisany v historijnom zapisě, ale upotrěbjaje razgleděti i prstenov vměsto ranějših razsmotriti i kolc. To sravnjenje ustanavjaje prisutnost rečenicy v aktualnom prěvodu; ne ustanavjaje samo po sobě, čto jest na originalnoj pečatnoj stranicě, ni kakva byla cěla historija vključanja.\par
\noindent\emph{Textual comparison:} The current opening of section 6 of Paper 31 contains Zato najprvo trěba razgleděti diskriminantne idealy primarnyh prstenov. It retains the logical transition described in the historical record, but uses razgleděti and prstenov in place of earlier razsmotriti and kolc. This comparison establishes the sentence's presence in the current translation; it does not by itself establish the original printed reading or the full history of inclusion.\par
\noindent{\footnotesize TEXT-COMP-17; pasaže / passages: 18279; \texttt{CORPUS\_FORM\_COMPARISONS.json}.}\par
\end{minipage}\par
\par\medskip\Needspace{7\baselineskip}
\noindent\begin{minipage}{\linewidth}
\hypertarget{TERM-0568}{}
\noindent\textbf{Instrukcija k razsmotrenju produktov form}\par\noindent\textbf{An instruction to consider products of forms}\quad{\footnotesize\texttt{TERM-0568}}\par
\smallskip\noindent\emph{Razlog:} Fraza «slědujuće produkty form trěba razmatrivati» prěvodi instrukciju v dokazě. Produkt sohranjaje algebraičny objekt, a trěba čini izričnoju neobhodimost slědujučego razsmotrenja. Razlog sravnjaje konstrukciju s obyčnym ruskym i ukrajinskim naučnym registrom, ale ne navodi čitateljski test. Zapis označaje i leksiku, i konstrukciju kako provizorne; jasno vyražena instrukcija ne znači, že vse njihove morfologične detale už prověrene.\par
\noindent\emph{Reason:} The phrase «slědujuće produkty form trěba razmatrivati» translates an instruction within the proof. Produkt preserves the algebraic object, and trěba makes the necessity of the ensuing consideration explicit. The rationale compares the construction with ordinary Russian and Ukrainian mathematical register but supplies no reader test. It treats both vocabulary and construction as provisional; an intelligible instruction does not imply that every morphological detail has been verified.\par
\noindent{\footnotesize Izvorne zapisy / Original records: HIST-g019-011.}\par
\end{minipage}\par
\par\medskip\Needspace{7\baselineskip}
\noindent\begin{minipage}{\linewidth}
\hypertarget{TERM-0599}{}
\noindent\textbf{Latinične oznaky formul v dvuh pismennyh versijah}\par\noindent\textbf{Latin formula labels in two script versions}\quad{\footnotesize\texttt{TERM-0599}}\par
\smallskip\noindent\emph{Razlog:} Oznaky $\mathrm{a)}$, $\mathrm{b)}$, $\mathrm{c)}$ identifikujut formuly, nezavisno od latiničnogo ili kiriličnogo pisma okolnoj prozy. Zapis izbira prjamy matematičny šrift, da one ostanut latiničnymi oznakami i ne budut obrabotane kako slova. To jest pravilo identičnosti i pisma, ne slovničny prěvod trěh bukv. Jego priměnjenje trěbuje takože sohraniti vse upominanja tyh formul; sama tipografična komanda ne dokazuje pravilnost vsih odkazov.\par
\noindent\emph{Reason:} The labels $\mathrm{a)}$, $\mathrm{b)}$, $\mathrm{c)}$ identify formulas independently of whether the surrounding prose is Latin or Cyrillic. The record selects upright mathematical lettering so that they remain Latin labels rather than being processed as words. This is a rule about identity and script, not a dictionary translation of three letters. Applying it also requires preserving references to the formulas; the typesetting command alone does not verify every cross-reference.\par
\noindent{\footnotesize Izvorne zapisy / Original records: HIST-g021-013.}\par
\end{minipage}\par
\par\medskip\Needspace{7\baselineskip}
\noindent\begin{minipage}{\linewidth}
\hypertarget{TERM-0600}{}
\noindent\textbf{Gotične matematične oznaky i jejich šrift}\par\noindent\textbf{Fraktur mathematical identifiers and their typeface}\quad{\footnotesize\texttt{TERM-0600}}\par
\smallskip\noindent\emph{Razlog:} Gotične bukvy $\mathfrak h$ i $\mathfrak t$ ostavajut različnymi matematičnymi oznakami v pomočnom tvrđenju VI i podnožnoj primětkě. Razlog vybiraje izrično zadanje gotičnogo šrifta vměsto skračenyh komand, ktore ne sut oprěděljene v lokalnom tekstu. Tak možno sohraniti izgled oznak nezavisno od razlik v makrah. To ne dava prava zaměniti gotičnu bukvu obyčnoju latiničnoju ili tvrditi věrnost pečatnomu iztočniku bez njegovogo sravnjenja.\par
\noindent\emph{Reason:} Fraktur $\mathfrak h$ and $\mathfrak t$ remain distinct mathematical identifiers in auxiliary proposition VI and its footnote. The rationale chooses explicit fraktur styling instead of abbreviation commands not defined in the local source. This preserves the notation's appearance independently of differences in macros. It does not license replacing fraktur with ordinary Latin lettering or claiming fidelity to the printed source without comparison.\par
\noindent{\footnotesize Izvorne zapisy / Original records: HIST-g099-012.}\par
\end{minipage}\par
\par\medskip\Needspace{7\baselineskip}
\noindent\begin{minipage}{\linewidth}
\hypertarget{TERM-0691}{}
\noindent\textbf{Spolavtorstvo i sklonjenje imena}\par\noindent\textbf{Coauthorship and the inflection of a name}\quad{\footnotesize\texttt{TERM-0691}}\par
\smallskip\noindent\emph{Razlog:} Sovmestno s R. Brauerom izražaje spolavtorstvo i sohranjaje razpoznatelnost imena. Razlog izrično motivuje instrumentalnu formu Brauerom, potrebnu po s v toj frazě, i jejo priměnimost pri kiriličnom prěpisu. To jest odluka o pripisanju avtorstva i gramatyčnoj formě vlastnogo imena, ne novy matematičny termin. Aktualno zaglavje sodrživaje tu formulu; sam zapis razloga ne dokazuje polnu proverku kiriličnoj verzije.\par
\noindent\emph{Reason:} Sovmestno s R. Brauerom expresses coauthorship while keeping the name recognizable. The rationale specifically motivates the instrumental Brauerom after s in this phrase and its intended use in the Cyrillic rendering. This concerns attribution and the grammatical form of a proper name, not a new mathematical term. The current heading contains the formula; the rationale alone does not establish complete checking of the Cyrillic version.\par
\noindent{\footnotesize Izvorne zapisy / Original records: HIST-g166-003.}\par
\smallskip\noindent\emph{Sravnjenje s tekstom:} Zaglavje raboty 32 sodrživaje Sovmestno s R. Brauerom. To jest točna lokalna pojavjenost formy i pripisanja spolavtorstva; ne nova leksična attestacija matematičnogo termina. Opravdanje sklonjenja i proverka drugego pisma sut odvojene ot samego nalezenja zaglavja.\par
\noindent\emph{Textual comparison:} Paper 32's heading contains Sovmestno s R. Brauerom. This establishes the local form and coauthorship attribution, not a new mathematical lexical attestation. The justification of the inflection and checking another script remain distinct from merely locating the heading.\par
\noindent{\footnotesize TEXT-COMP-47; pasaže / passages: 19171; \texttt{CORPUS\_FORM\_COMPARISONS.json}.}\par
\end{minipage}\par
\par\medskip\Needspace{7\baselineskip}
\noindent\begin{minipage}{\linewidth}
\hypertarget{TERM-0706}{}
\noindent\textbf{Polnota klasifikacije identičnostij}\par\noindent\textbf{Completeness of the classification of identities}\quad{\footnotesize\texttt{TERM-0706}}\par
\smallskip\noindent\emph{Razlog:} Vyčerpana jest cělost nerazložimyh identičnostij izražaje polnotu klasifikacije: opisane konstrukcije daju vse takve identičnosti. Razlog svęzuje tu rečenicu s Teoremom III raboty 4. Vyčerpana ne označaje prěstanok materiala dlja počitanja, ale tvrđenje, že nema izpuščenyh slučajev. To jest izbor izloženja i kvantifikacije, ne predloženje novogo matematičnogo imenika.\par
\noindent\emph{Reason:} Vyčerpana jest cělost nerazložimyh identičnostij expresses completeness of a classification: the described constructions give all such identities. The rationale ties the sentence to Theorem III of Paper 4. Exhausted does not mean that a calculation has run out of material; it asserts that no cases have been omitted. This concerns exposition and quantification, not a newly coined mathematical noun.\par
\noindent{\footnotesize Izvorne zapisy / Original records: HIST-g034-009.}\par
\end{minipage}\par
\par\medskip\Needspace{7\baselineskip}
\noindent\begin{minipage}{\linewidth}
\hypertarget{TERM-0778}{}
\noindent\textbf{Geografično ime i pismo}\par\noindent\textbf{A place-name and its script}\quad{\footnotesize\texttt{TERM-0778}}\par
\smallskip\noindent\emph{Razlog:} Göttingen jest vlastno geografično ime, ne novy matematičny termin. Razlog zadržuje německy zapis v latiničnom prěvodu, dokolje ruska i ukrajinska sravniteljnica davajut svoje kirilične formy. Taka urednička konvencija sohranjaje razpoznatelnost města i bibliografičnyh odkazov. Bukva ö v tom imenu ne opravdava jejo upotrěbu v běžnoj medžuslovjanskoj proze. Historijny razlog govori takože o avtomatičnom kiriličnom prěpisu; sam po sobě ne dokazuje, že vsak izhod togo prěpisa byl jezyčno prověren.\par
\noindent\emph{Reason:} Göttingen is a geographical proper name, not a new mathematical term. The rationale retains the German spelling in the Latin translation, while the Russian and Ukrainian comparanda give their respective Cyrillic forms. This editorial convention preserves recognition of the city and bibliographical references. The ö in the name does not authorize its use in ordinary Interslavic prose. The historical rationale also mentions automatic Cyrillic transliteration; by itself, that statement does not establish linguistic checking of every transliterated output.\par
\noindent{\footnotesize Izvorne zapisy / Original records: HIST-g165-001.}\par
\end{minipage}\par
\par\medskip\Needspace{7\baselineskip}
\noindent\begin{minipage}{\linewidth}
\hypertarget{TERM-0838}{}
\noindent\textbf{Kogda homomorfizmy stavajut izomorfizmami}\par\noindent\textbf{When homomorphisms become isomorphisms}\quad{\footnotesize\texttt{TERM-0838}}\par
\smallskip\noindent\emph{Razlog:} Homomorfizmy prehodet v izomorfizmy jest izbor za cělo tvrdženje, ne novo slovo. On zadržuje zobraženja kak predmet rečenija vměsto nepreciznogo izjednačenja polj. Pasaž raboty 31 govorit o specialnom slučaju, kogda razgledany prsten jest polje; konstrukcija dava različne izomorfizmy s konjugovanymi poljami nad P. Ne tvrdi se, že vsaky homomorfizm bez dodatnyh uslovij jest izomorfizm, niti da izomorfne polja sut bukvalno jedna i ta sama podmnožina.\par
\noindent\emph{Reason:} Homomorfizmy prehodet v izomorfizmy is a sentence-level choice, not a new word. It preserves maps as the sentence’s subject instead of loosely equating fields. The passage in Paper 31 treats the special case in which the ring under consideration is a field; the construction gives distinct isomorphisms with conjugate fields over P. It does not assert that every homomorphism without further conditions is an isomorphism, or that isomorphic fields are literally the same subset.\par
\noindent{\footnotesize Izvorne zapisy / Original records: HIST-g147-005.}\par
\smallskip\noindent\emph{Sravnjenje s tekstom:} Kogda R jest polje, rabota 31 dobyvaje različne izomorfizmy k konjugovanym poljam Ki nad P. Konec pasaža izrično razlikuje izomorfno polje od polja už sodržanogo v zadannom Galoisovom razširjenju Q. Zato glagolna konstrukcija o homomorfizmah čuva informaciju, ktora byla by izgubjena prěz prosty zapis ravnosti polj.\par
\noindent\emph{Textual comparison:} When R is a field, Paper 31 obtains distinct isomorphisms to conjugate fields Ki over P. The end of the passage expressly distinguishes an isomorphic field from one already contained in the specified Galois extension Q. The sentence about homomorphisms therefore preserves information that a simple assertion of equality between fields would lose.\par
\noindent{\footnotesize TEXT-COMP-95; pasaže / passages: 18525; \texttt{CORPUS\_FORM\_COMPARISONS.json}.}\par
\end{minipage}\par
\par\medskip\Needspace{7\baselineskip}
\noindent\begin{minipage}{\linewidth}
\hypertarget{TERM-0870}{}
\noindent\textbf{Sohranjanje bibliografičnogo nazvanja i lokatora}\par\noindent\textbf{Preserving a bibliographic title and locator}\quad{\footnotesize\texttt{TERM-0870}}\par
\smallskip\noindent\emph{Razlog:} Idealtheorie tu jest kratko bibliografično nazvanje, zato ostaje na německom jezyku. Sohranjene sut takože §7, 3: čitatelj mora mogti najdti citovany argument. To ne jest novy medžuslovjansky termin i ne anglijsko nepreloženo objašnjenje. Zachovany opis vyboru govori o primětkě, ne prisutnoj v jednom izvlečenom tekstě, ale vidimoj v pečatnom iztočniku. To jest historično tvrdženje o osnovi vyboru; sravnjenje nynějšego korpusnogo odlomka samo po sobě ne jest novym čitanjem pečatnoj stranicy. Citacija v rabotě 31 jest svazana s teorijeju idealov poredka.\par
\noindent\emph{Reason:} Idealtheorie is a shortened bibliographic title here and therefore remains in German. The locator §7, 3 is also retained so that readers can find the cited argument. This is not a new Interslavic term or an untranslated English explanation. The preserved account says that a footnote absent from one extracted text was visible in the printed source. That is a historical claim about the basis of the choice; comparing the present corpus passage is not itself a fresh inspection of the printed page. In Paper 31 the citation accompanies the ideal theory of the order.\par
\noindent{\footnotesize Izvorne zapisy / Original records: HIST-g150-007.}\par
\smallskip\noindent\emph{Sravnjenje s tekstom:} Tri primětky raboty 31 imajut konkretne funkcije: §7, 3 pri teoriji idealov poredka; §9, 7 pri integralnoj zamknutosti; §9 i §3, 1 pri egzistenciji konečnoj modulnoj bazy. Německe kratko nazvanje Idealtheorie identifikuje iztočnik. Prěvod veznika v tretjoj primětkě ne měnja nazvanje. Točne odlomky potrvrdžajut funkciju citacij v nynějšem korpusu; ne potrvrdžajut samo soboju prijašnje tvrdženja o pečatnom svědku.\par
\noindent\emph{Textual comparison:} The three Paper 31 footnotes have specific functions: §7, 3 accompanies the ideal theory of the order; §9, 7 accompanies integral closure; §9 and §3, 1 accompany existence of a finite module basis. The shortened German title Idealtheorie identifies the source. Translating the conjunction in the third note does not change the title. These passages establish the citations' functions in the present corpus; they do not by themselves verify earlier claims about the printed witness.\par
\noindent{\footnotesize TEXT-COMP-107; pasaže / passages: 18593, 18861, 18897; \texttt{CORPUS\_FORM\_COMPARISONS.json}.}\par
\end{minipage}\par
\par\medskip\Needspace{7\baselineskip}
\noindent\begin{minipage}{\linewidth}
\hypertarget{TERM-0871}{}
\noindent\textbf{Citacija pri integralnoj zamknutosti}\par\noindent\textbf{Citation supporting integral closure}\quad{\footnotesize\texttt{TERM-0871}}\par
\smallskip\noindent\emph{Razlog:} Nazvanje Idealtheorie i lokator §9, 7 ostajut neizměnjene, zato že primětka služi bibliografičnomu odkazivanju. Ona podpiraje priměnjenje integralnoj zamknutosti pri razgledanju diskriminanta v rabotě 31. Jezykova rabota jest tu v sohranjanju funkcije citacije, ne v iznajdenju slova. Zapisana motivacija spominaje razliku medžu izvlečenym tekstom i pečatnym iztočnikom; ta istorijska osnova ne smi byti prědstavjena kak nova nezavisna prověrka skana. Ekspert može posuditi i bibliografičnu točnost, i to, dali citovany rezultat odgovarja konkretnomu kroku argumenta.\par
\noindent\emph{Reason:} The title Idealtheorie and locator §9, 7 remain unchanged because the footnote serves bibliographic reference. It supports the use of integral closure in the discussion of the discriminant in Paper 31. The linguistic task is to preserve the citation's function, not to invent a word. The recorded rationale mentions a difference between an extracted text and the printed source; that historical basis must not be presented as a new independent scan check. An expert can assess both bibliographic accuracy and whether the cited result supports this particular step of the argument.\par
\noindent{\footnotesize Izvorne zapisy / Original records: HIST-g157-010.}\par
\smallskip\noindent\emph{Sravnjenje s tekstom:} Tri primětky raboty 31 imajut konkretne funkcije: §7, 3 pri teoriji idealov poredka; §9, 7 pri integralnoj zamknutosti; §9 i §3, 1 pri egzistenciji konečnoj modulnoj bazy. Německe kratko nazvanje Idealtheorie identifikuje iztočnik. Prěvod veznika v tretjoj primětkě ne měnja nazvanje. Točne odlomky potrvrdžajut funkciju citacij v nynějšem korpusu; ne potrvrdžajut samo soboju prijašnje tvrdženja o pečatnom svědku.\par
\noindent\emph{Textual comparison:} The three Paper 31 footnotes have specific functions: §7, 3 accompanies the ideal theory of the order; §9, 7 accompanies integral closure; §9 and §3, 1 accompany existence of a finite module basis. The shortened German title Idealtheorie identifies the source. Translating the conjunction in the third note does not change the title. These passages establish the citations' functions in the present corpus; they do not by themselves verify earlier claims about the printed witness.\par
\noindent{\footnotesize TEXT-COMP-107; pasaže / passages: 18593, 18861, 18897; \texttt{CORPUS\_FORM\_COMPARISONS.json}.}\par
\end{minipage}\par
\par\medskip\Needspace{7\baselineskip}
\noindent\begin{minipage}{\linewidth}
\hypertarget{TERM-0872}{}
\noindent\textbf{Prěvod veznika v bibliografičnoj primětkě}\par\noindent\textbf{Translating a conjunction in a bibliographic note}\quad{\footnotesize\texttt{TERM-0872}}\par
\smallskip\noindent\emph{Razlog:} V Idealtheorie, §9 i §3, 1 sohranja se německe nazvanje, ale veznik und prěvodi se kak i. Granica jest smyslova: nazvanje identificiruje dělo, dokole veznik spojuje dva lokatora i prinaleži jezyku primětky. To dava čitatelju prirodnu frazu bez gubitka možlivosťi najdti iztočnik. Primětka v rabotě 31 stoji pri egzistenciji konečnoj modulnoj bazy. Zachovany zapis pripisuje ju pečatnomu svědku; ta atribucija ostaje istorijskym svedenjem, ne novym osvědčenjem stranicy. Takovy izbor prinaleži redakcionnym konvencijam, ne spisku novyh matematičnyh slov.\par
\noindent\emph{Reason:} In Idealtheorie, §9 i §3, 1 the German title is retained but the conjunction und becomes i. The boundary is functional: the title identifies the work, whereas the conjunction joins two locators and belongs to the footnote's language. This gives readers a natural phrase without sacrificing source identification. In Paper 31 the note accompanies the existence of a finite module basis. The preserved record attributes it to the printed witness; that attribution remains historical evidence, not a fresh certification of the page. This choice belongs among editorial conventions, not among new mathematical words.\par
\noindent{\footnotesize Izvorne zapisy / Original records: HIST-g158-011.}\par
\smallskip\noindent\emph{Sravnjenje s tekstom:} Tri primětky raboty 31 imajut konkretne funkcije: §7, 3 pri teoriji idealov poredka; §9, 7 pri integralnoj zamknutosti; §9 i §3, 1 pri egzistenciji konečnoj modulnoj bazy. Německe kratko nazvanje Idealtheorie identifikuje iztočnik. Prěvod veznika v tretjoj primětkě ne měnja nazvanje. Točne odlomky potrvrdžajut funkciju citacij v nynějšem korpusu; ne potrvrdžajut samo soboju prijašnje tvrdženja o pečatnom svědku.\par
\noindent\emph{Textual comparison:} The three Paper 31 footnotes have specific functions: §7, 3 accompanies the ideal theory of the order; §9, 7 accompanies integral closure; §9 and §3, 1 accompany existence of a finite module basis. The shortened German title Idealtheorie identifies the source. Translating the conjunction in the third note does not change the title. These passages establish the citations' functions in the present corpus; they do not by themselves verify earlier claims about the printed witness.\par
\noindent{\footnotesize TEXT-COMP-107; pasaže / passages: 18593, 18861, 18897; \texttt{CORPUS\_FORM\_COMPARISONS.json}.}\par
\end{minipage}\par
\par\medskip\Needspace{7\baselineskip}
\noindent\begin{minipage}{\linewidth}
\hypertarget{TERM-0882}{}
\noindent\textbf{Negacija v zaključanju dokaza}\par\noindent\textbf{Negation in a proof conclusion}\quad{\footnotesize\texttt{TERM-0882}}\par
\smallskip\noindent\emph{Razlog:} Identičnosti ne nastavajut jest věta o rezultatu konkretnej konstrukcije v dokazu, ne nove nazvanje za klasu objektov. Zapis dava prednost obyčnomu izloženju, podobnomu ruskoj i ukrajinskoj formulaciji, vměsto težšej parafrazy. Negacija mora ostati svazana s razgledanym postupkom: fraza ne znači, že nikakve identičnosti v celom predmětu ne egzistujut. Objašnjenje čitatelju ima pokazati, iz čego identičnosti ne nastavajut. Tvrdžena prozračnost jest jezykova ocěna, ne rezultat empirnogo testa.\par
\noindent\emph{Reason:} Identičnosti ne nastavajut is a sentence about the outcome of a particular construction in a proof, not a new name for a class of objects. The record favours ordinary exposition, comparable to the Russian and Ukrainian formulations, over a heavier paraphrase. The negation must remain attached to the procedure under discussion: the phrase does not mean that no identities exist anywhere in the subject. The explanation should show readers what fails to produce identities. The claimed transparency is a linguistic judgment, not an empirical test result.\par
\noindent{\footnotesize Izvorne zapisy / Original records: HIST-g012-011.}\par
\end{minipage}\par
\par\medskip\Needspace{7\baselineskip}
\noindent\begin{minipage}{\linewidth}
\hypertarget{TERM-0893}{}
\noindent\textbf{Avtorin bibliografičny odkaz}\par\noindent\textbf{The author's bibliographic cross-reference}\quad{\footnotesize\texttt{TERM-0893}}\par
\smallskip\noindent\emph{Razlog:} V rabotě citovanoj v primětce 6, str. 83 sohranja bibliografičny odkaz vměsto da razvije polnu citaciju vnutri věty. Primětka raboty 31 odkazyvaje na Grellovu rabotu kak iztočnik dokazov pri prěhodu k prstenu kvocientov. Čislo primětky i stranica sut avtorine lokatory, ne avtomatično čisla nynějšego čitateljskogo PDF. Jih možno objasniti i povezati s bibliografijeju, ale ne tiho prenumerovati. Zapis pripisuje odkaz pečatnomu tekstu; sravnjenje korpusnoj primětky potrvrdžaje jejno nynějšne sodržanje, ne novu prověrku pečatnoj stranicy.\par
\noindent\emph{Reason:} V rabotě citovanoj v primětce 6, str. 83 retains a bibliographic cross-reference rather than expanding the complete citation inside the sentence. The Paper 31 footnote directs readers to Grell's work for the proofs concerning passage to the quotient ring. Its note number and page are Noether's locators, not automatically numbers in the present reader PDF. They can be explained and linked to the bibliography, but must not be silently renumbered. The record attributes the reference to the printed text; comparison of the corpus footnote establishes its current content, not a fresh inspection of the printed page.\par
\noindent{\footnotesize Izvorne zapisy / Original records: HIST-g159-010.}\par
\smallskip\noindent\emph{Sravnjenje s tekstom:} Točka o prěhodu k prstenu kvocientov v rabotě 31 oprědělja, ktore elementy mogut byti v imenovatelju. Primětka pri toj točkě odkazyvaje na Grella i sohranja avtorin lokator primětka 6, str. 83. Definicija i odkaz imajut različne funkcije: prva oprědělja objekt, vtory ukazuje iztočnik dokazov. Ni jedno ne trěba izpustiti pri prěvodu.\par
\noindent\emph{Textual comparison:} Paper 31's paragraph on passage to the quotient ring specifies which elements may occur in the denominator. Its footnote refers to Grell and retains Noether's locator note 6, p. 83. Definition and reference have different functions: one determines the object, the other identifies the source of proofs. Neither should disappear in translation.\par
\noindent{\footnotesize TEXT-COMP-115; pasaže / passages: 18932; \texttt{CORPUS\_FORM\_COMPARISONS.json}.}\par
\end{minipage}\par
\par\medskip\Needspace{7\baselineskip}
\noindent\begin{minipage}{\linewidth}
\hypertarget{TERM-0943}{}
\noindent\textbf{Prěryv v bukvenyh označenjah tablicy}\par\noindent\textbf{Discontinuous letter labels in a table}\quad{\footnotesize\texttt{TERM-0943}}\par
\smallskip\noindent\emph{Razlog:} V tablici nereducibilnyh form raboty 2, §14, zachovany zapis ukazuje označenja A, B, C, D, E, G, dokolě spisok redukcij imaje takože F. Odluka byla sohraniti razliku, ne dodavati izmysljenu formu F ni tiho prěimenovati G. Nynějši tekst ima to samo nesovpadanje. To potvrždžaje njegovo prisutstvo v prevodu, ale ne ustanavjaje, čto bylo na originalnoj pečatnoj stranicě. Otvorjeno ostaje, jest li to pečatna nepravilnost izvora ili ošibka posrědničnogo teksta. Do razrěšenja toj razliki ni utvrdžati, ni izpravljati original bez dokaza ne jest opravdano.\par
\noindent\emph{Reason:} In the irreducible-form table of Paper 2, §14, the surviving record identifies labels A, B, C, D, E, G, while the reduction list also includes F. The decision was to preserve the difference, not add an invented form F or silently rename G. The present text contains the same mismatch. This confirms its presence in the translation, but does not establish what appeared on the original printed page. Whether it is an original typographic irregularity or an error in an intervening text remains unresolved. Neither certifying nor correcting the original without evidence is justified.\par
\noindent{\footnotesize Izvorne zapisy / Original records: HIST-g015-016.}\par
\smallskip\noindent\emph{Sravnjenje s tekstom:} Tablica §14 ima A, B, C, D, E, G; niže jest redukcija F. Sravnjenje potvrždžaje nesovpadanje v nynějšem prevodu. Ne odlučuje, jest li ono v originalnom pečatnom iztočniku; do takoj proverky bukvy ne prenumerujut se i nedostajuča matematična forma ne izmysljaje se.\par
\noindent\emph{Textual comparison:} The §14 table has A, B, C, D, E, G; below it is reduction F. Comparison confirms the mismatch in the present translation. It does not decide whether the mismatch occurs in the original printed source; without that comparison, neither relabelling nor inventing a missing mathematical form is warranted.\par
\noindent{\footnotesize TEXT-COMP-130; pasaže / passages: 2163, 2165, 2169, 2170, 2171, 2172, 2210; \texttt{CORPUS\_FORM\_COMPARISONS.json}.}\par
\end{minipage}\par
\par\medskip\Needspace{7\baselineskip}
\noindent\begin{minipage}{\linewidth}
\hypertarget{TERM-0962}{}
\noindent\textbf{Zavěr dokaza}\par\noindent\textbf{Concluding a proof}\quad{\footnotesize\texttt{TERM-0962}}\par
\smallskip\noindent\emph{Razlog:} Tym je dokazana prědava zavěr dokaza v rabotě 2, §15. Tym upominaje predhodny argument; dokazana otnosi se k dokazanoj věti ili drugomu imeniku ženskogo roda v konkretnom kontekstě. Zapis daje prědnost kratkoj, obyčnoj formulaciji pred kopirovanjem německogo poredka slov. Ne trěba avtomatično upotrěbjati isto zakončenje pri vsakom imeniku. Fraza sohranja avtorovo tvrdženje o zakončenom dokazě; ne jest osvědčenjem, že prevoditelj nezavisno proveril vsaky njegov krok.\par
\noindent\emph{Reason:} Tym je dokazana renders the conclusion of a proof in Paper 2, §15. Tym refers to the preceding argument; dokazana agrees with the proposition or other feminine noun in its particular context. The record favours a short, ordinary formulation over copying German word order. The same ending should not be applied automatically to every noun. The phrase preserves the author's claim that the proof is complete; it is not certification that the translator independently checked every step.\par
\noindent{\footnotesize Izvorne zapisy / Original records: HIST-g016-012.}\par
\end{minipage}\par
\par\medskip\Needspace{7\baselineskip}
\noindent\begin{minipage}{\linewidth}
\hypertarget{TERM-0964}{}
\noindent\textbf{Opisny prevod nazvanja časopisa}\par\noindent\textbf{Descriptive translation of a journal title}\quad{\footnotesize\texttt{TERM-0964}}\par
\smallskip\noindent\emph{Razlog:} Letopisy Njemeckogo matematičnogo družstva jest opisny prevod nazvanja časopisa v proznom zaglavju. Kogda časopis citovany jest skračenjem J. Ber. d. DMV, identifikačno skračenje ostaje sohranjeno, da čitatelj može najti iztočnik. Motivacija razlikuje čitateljsko objasnjenje od bibliografičnoj identičnosti. Družstvo naměrno izbira široku slovjansku semju za obščestvo ili združenje vměsto internacionalizovanogo imenika; jeho uměstnost ostaje otvorjena jezykovoj ocěnce. Opisny prevod ne trěba vydavati za oficialny samonazov organizacije.\par
\noindent\emph{Reason:} Letopisy Njemeckogo matematičnogo družstva is a descriptive translation of the journal title in a prose heading. Where the journal is cited as J. Ber. d. DMV, that identifying abbreviation remains so that readers can locate the source. The rationale distinguishes reader explanation from bibliographic identity. Družstvo deliberately selects a broad Slavic society-or-association family rather than an internationalized noun; its suitability remains open to linguistic appraisal. The descriptive translation should not be presented as the organization's official self-designation.\par
\noindent{\footnotesize Izvorne zapisy / Original records: HIST-g215-005.}\par
\end{minipage}\par
\par\medskip\Needspace{7\baselineskip}
\noindent\begin{minipage}{\linewidth}
\hypertarget{TERM-0965}{}
\noindent\textbf{Vsaky i samo takovy: dva směry tvrdženja}\par\noindent\textbf{Every and only such: both directions of an assertion}\quad{\footnotesize\texttt{TERM-0965}}\par
\smallskip\noindent\emph{Razlog:} Vsaky prosty element $p$, ktory vhodi v $D_{\mathfrak T}$, i samo takovy element sohranja oba směry tvrdženja raboty 31. Ne trěba skratiti jeho do vsaky bez i samo takovy: izgubila by se obratna implikacija. V sravnjenom odlomku stojit važno uslovje, že polje klasov ostatkov $\mathfrak o/\mathfrak o p$ jest doskonalo. Daljše tvrdženje o kvadratu prostogo ideala ješče vybira glavny poredok. Slovnik zato objasnjaje silu kvantifikacije vměstě s hipotezami, ne kak bezuslovno pravilo o vseh diskriminantah.\par
\noindent\emph{Reason:} Vsaky prosty element $p$, ktory vhodi v $D_{\mathfrak T}$, i samo takovy element preserves both directions of the assertion in Paper 31. Shortening it to every without and only such would lose the converse implication. The compared passage includes the important hypothesis that the residue-class field $\mathfrak o/\mathfrak o p$ is perfect. The further statement about the square of a prime ideal also selects the principal order. The dictionary therefore explains the quantification together with its hypotheses, not as an unconditional rule about all discriminants.\par
\noindent{\footnotesize Izvorne zapisy / Original records: HIST-g154-002.}\par
\smallskip\noindent\emph{Sravnjenje s tekstom:} V diskriminantnom poslědku raboty 31 vsaky i samo takovy sohranja dva směry pod uslovjem doskonalosti polja ostatkov. Drugy odlomok ograničuje uslovje bez vlastnogo dělitelja na pridružene proste idealy. V oběh slučajah kratka gramatična fraza nosi čest matematičnyh hipotez; izpuščena čest ne jest samo stilistična razlika.\par
\noindent\emph{Textual comparison:} In Paper 31's discriminant consequence, every and only such preserves both directions under the perfect-residue-field hypothesis. The other passage restricts the no-proper-divisor condition to associated prime ideals. In both cases a short grammatical phrase carries mathematical qualifications; omission would not be merely stylistic.\par
\noindent{\footnotesize TEXT-COMP-135; pasaže / passages: 18749, 18593; \texttt{CORPUS\_FORM\_COMPARISONS.json}.}\par
\end{minipage}\par
\par\medskip\Needspace{7\baselineskip}
\noindent\begin{minipage}{\linewidth}
\hypertarget{TERM-0968}{}
\noindent\textbf{Razgledati polje kak modul}\par\noindent\textbf{Regarding a field as a module}\quad{\footnotesize\texttt{TERM-0968}}\par
\smallskip\noindent\emph{Razlog:} $K$ jest razsmatrjany kako $\Omega$-modul objasnjaje, s kakoju strukturoju rabota 31 razgledaje dany objekt. Ne govori, že $K$ jest bukvalno ravno $\Omega$. V odlomku o slědu i diskriminantnom idealu pole $K$ razgledaje se kak $\Omega$-modul, s linearno nezavisnymi elementami za bazu. Sravnjeni tekst upotrěbjaje razgledany vměsto arhivnogo razsmatrjany. Motivacija vsegda ostaje izrično prědati razgledanje objekta kak modula; ona ne ustanavjaje, že imenno arhivny glagol jest jediny prirodny izbor.\par
\noindent\emph{Reason:} $K$ jest razsmatrjany kako $\Omega$-modul explains the structure under which Paper 31 considers the object. It does not say that $K$ is literally equal to $\Omega$. In the passage on trace and the discriminant ideal, the field $K$ is regarded as an $\Omega$-module, with linearly independent elements used for a basis. The compared text uses razgledany instead of archival razsmatrjany. The rationale remains explicit expression of regarded as a module; it does not establish that the archival verb is the only idiomatic choice.\par
\noindent{\footnotesize Izvorne zapisy / Original records: HIST-g157-003.}\par
\smallskip\noindent\emph{Sravnjenje s tekstom:} V rabotě 31 $K$ razgledaje se kak $\Omega$-modul pri oprěděljenju slěda. V rabotě 34 tělo ne mora byti komutativno, a pravy $K$-modul imaje izrično uslovje $ae=a$. Fraza razgledati kak i pridevnik pravy nosat različne, oba važné strukturne informacije.\par
\noindent\emph{Textual comparison:} In Paper 31, $K$ is regarded as an $\Omega$-module when defining trace. In Paper 34 the division ring need not be commutative, and the right $K$-module has the explicit condition $ae=a$. Regarded as and right carry different, equally important structural information.\par
\noindent{\footnotesize TEXT-COMP-137; pasaže / passages: 18857, 20221, 20223; \texttt{CORPUS\_FORM\_COMPARISONS.json}.}\par
\end{minipage}\par
\par\medskip\Needspace{7\baselineskip}
\noindent\begin{minipage}{\linewidth}
\hypertarget{TERM-0973}{}
\noindent\textbf{Nazvanje glavy i jeje rimsko čislo}\par\noindent\textbf{A chapter heading and its Roman numeral}\quad{\footnotesize\texttt{TERM-0973}}\par
\smallskip\noindent\emph{Razlog:} Glava III prodolžuje izbrano nazvanje glavy v rabotě 2. Ruski prěvod imaje Глава, ukrajinski Розділ, medžuslovjanski Glava. To jest odluka o strukturi i navigaciji, ne novy matematičny pojam. Rimsko čislo III sohranja identičnost glavy i mora byti soglasno s odkazami; jego ne trěba zaměnjati novym čislom samo iz-za različnogo stranovanja. Različny ukrajinski imenik ne znači, že razděljenje dela jest različnym.\par
\noindent\emph{Reason:} Glava III continues the selected chapter designation in Paper 2. Russian has Глава, Ukrainian Розділ, and Interslavic Glava. This concerns structure and navigation, not a new mathematical concept. Roman III preserves the chapter's identity and must agree with cross-references; it should not be replaced merely because pagination differs. The different Ukrainian noun does not mean that the work's division is different.\par
\noindent{\footnotesize Izvorne zapisy / Original records: HIST-g008-001.}\par
\end{minipage}\par
\par\medskip\Needspace{7\baselineskip}
\noindent\begin{minipage}{\linewidth}
\hypertarget{TERM-0974}{}
\noindent\textbf{Negacija s izključkom jediničnogo ideala}\par\noindent\textbf{Negation with the unit-ideal exception}\quad{\footnotesize\texttt{TERM-0974}}\par
\smallskip\noindent\emph{Razlog:} Ne imajut nikakogo vlastnogo dělitelja raznogo od jediničnogo ideala sohranja negaciju i izključok v rabotě 31, §7. Motivacija prodolžuje ranějše vlastny dělitelj iz §2 vměsto vvedenja novej nazvy. V sravnjenom tvrdženju uslovje otnosi se k pridruženym prostym idealam primarnyh faktorov, ne ko vsem idealam prstena bez ograničenja. Izpuščenje vlastnogo ili raznogo od jediničnogo ideala izměnilo by matematično tvrdženje. Starši jezyk idealnoj děliteljnosti mora byti razjasnjen v svojom poredku, ne avtomatično prěložen v obyčno čislenu děliteljnost.\par
\noindent\emph{Reason:} Ne imajut nikakogo vlastnogo dělitelja raznogo od jediničnogo ideala preserves the negation and exception in Paper 31, §7. The rationale carries forward vlastny dělitelj from §2 rather than introducing a new name. In the compared statement, the condition applies to the associated prime ideals of the primary factors, not to every ideal in the ring without restriction. Dropping proper or different from the unit ideal would change the assertion. The older language of ideal divisibility must be explained in its own order relation, not automatically read as ordinary numerical divisibility.\par
\noindent{\footnotesize Izvorne zapisy / Original records: HIST-g150-006.}\par
\smallskip\noindent\emph{Sravnjenje s tekstom:} V diskriminantnom poslědku raboty 31 vsaky i samo takovy sohranja dva směry pod uslovjem doskonalosti polja ostatkov. Drugy odlomok ograničuje uslovje bez vlastnogo dělitelja na pridružene proste idealy. V oběh slučajah kratka gramatična fraza nosi čest matematičnyh hipotez; izpuščena čest ne jest samo stilistična razlika.\par
\noindent\emph{Textual comparison:} In Paper 31's discriminant consequence, every and only such preserves both directions under the perfect-residue-field hypothesis. The other passage restricts the no-proper-divisor condition to associated prime ideals. In both cases a short grammatical phrase carries mathematical qualifications; omission would not be merely stylistic.\par
\noindent{\footnotesize TEXT-COMP-135; pasaže / passages: 18749, 18593; \texttt{CORPUS\_FORM\_COMPARISONS.json}.}\par
\end{minipage}\par
\par\medskip\Needspace{7\baselineskip}
\noindent\begin{minipage}{\linewidth}
\hypertarget{TERM-0993}{}
\noindent\textbf{Kogda najmenše obče mnogokratno jest produkt}\par\noindent\textbf{When the least common multiple equals the product}\quad{\footnotesize\texttt{TERM-0993}}\par
\smallskip\noindent\emph{Razlog:} Najmenše obče mnogokratno stavaje jednako produktu jest rečenica o rezultatu, ne definicija sinonimov. V rabotě 19, §8, ravnost dokazuje se za poparno vzajemno proste idealy, i Teorema XV koristi toj prěhod. Zato mnogokratno i produkt ostajut oddělne pojmy, a uslovje ne směje odpasti. Sravnjeni tekst imaje stavaje ravno produktu vměsto arhivnogo jednako. Taka proměna predikata ne měnja glavnu metodologičnu potrěbu: zachovati oba člena ravnosti i uslovje, pri ktorom ona imaje město.\par
\noindent\emph{Reason:} Najmenše obče mnogokratno stavaje jednako produktu states a result, not a definition of synonyms. Paper 19, §8 proves the equality for pairwise coprime ideals, and Theorem XV uses that transition. Least common multiple and product therefore remain distinct notions, and the hypothesis must not disappear. The compared text has stavaje ravno produktu rather than archival jednako. This predicate variation leaves the central requirement unchanged: preserve both sides of the equality and the condition under which it holds.\par
\noindent{\footnotesize Izvorne zapisy / Original records: HIST-g081-067, HIST-g082-067.}\par
\smallskip\noindent\emph{Sravnjenje s tekstom:} V rabotě 19 prehod k produktu argumentuje se pri poparnoj vzajemnoj prostosti. Pokazana jest ravnost najprvo za ideal i jego komplement, a potom postupok prodolžaje se. Izostavjenje uslovja proširilo by tvrdženje bez opravdanja.\par
\noindent\emph{Textual comparison:} Paper 19 argues for the product under pairwise coprimality. It first establishes the equality for an ideal and its complement, then continues the procedure. Removing the hypothesis would broaden the assertion without justification.\par
\noindent{\footnotesize TEXT-COMP-152; pasaže / passages: 13316, 13330; \texttt{CORPUS\_FORM\_COMPARISONS.json}.}\par
\end{minipage}\par
\par\medskip\Needspace{7\baselineskip}
\noindent\begin{minipage}{\linewidth}
\hypertarget{TERM-0994}{}
\noindent\textbf{Minus v definiciji slěda}\par\noindent\textbf{The minus sign in the trace definition}\quad{\footnotesize\texttt{TERM-0994}}\par
\smallskip\noindent\emph{Razlog:} Koeficient pri $(-t^{s-1})$ prodolžuje obyčnu konstrukciju koeficient pri v rabotě 31, §4. Minus jest čestju definicije slěda, ne tipografično ukrašenje: tekst naziva slědom koeficient pri negativnoj potegi v $|tE-C|$ i odmah dava sumu diagonalnyh elementov. Ispuščenje minusa obratilo by znak v toj konvenciji. Rusko i ukrajinsko sravnjenje izrično sohranjajut toj znak; medžuslovjansky razlog dodaje doslědnost s okolnoju rečeniceju o koeficientah.\par
\noindent\emph{Reason:} Koeficient pri $(-t^{s-1})$ continues the ordinary coefficient-at construction in Paper 31, §4. The minus sign belongs to the trace convention, not to decoration: the text names the coefficient of the negative monomial in $|tE-C|$ as the trace and immediately gives the sum of diagonal entries. Omitting the minus would reverse the sign under that convention. Both Russian and Ukrainian rationales expressly preserve it; the Interslavic rationale adds consistency with the surrounding coefficient sentence.\par
\noindent{\footnotesize Izvorne zapisy / Original records: HIST-g135-005.}\par
\smallskip\noindent\emph{Sravnjenje s tekstom:} V rabotě 31 koeficienty $|tE-C|$ sut invarianty klasa. Slěd jest koeficient pri $(-t^{s-1})$ i ravny sumě diagonalnyh elementov. Taka definicija svazuje leksičny izbor s přesnym znakom i s konkretnym smyslom klasa.\par
\noindent\emph{Textual comparison:} In Paper 31 the coefficients of $|tE-C|$ are class invariants. The trace is the coefficient of $(-t^{s-1})$ and equals the sum of diagonal entries. The definition links the lexical choice to an exact sign and a particular meaning of class.\par
\noindent{\footnotesize TEXT-COMP-145; pasaže / passages: 18024, 18026, 18028; \texttt{CORPUS\_FORM\_COMPARISONS.json}.}\par
\end{minipage}\par
\par\medskip\Needspace{7\baselineskip}
\noindent\begin{minipage}{\linewidth}
\hypertarget{TERM-1027}{}
\noindent\textbf{Komponenty kak prsteny, ne kak idealy}\par\noindent\textbf{Components as rings, not as ideals}\quad{\footnotesize\texttt{TERM-1027}}\par
\smallskip\noindent\emph{Razlog:} Komponenty kako kolca same po sebě, a ne kako idealy jest objasnjajuča primětka k diskriminantam v rabotě 31. Sravnjeni tekst imaje prsteny vměsto arhivnogo kolca i izrično razlikuje komponente kak samostojne prsteny od idealov većšego $\overline{\mathfrak R}$. To určuje, katere diskriminanty umnožajut se. Razlog kaže, že primětka byla najdena v pečatnom iztočniku; tuto ostaje tvrdženje sohranjenogo zapisa, ne nova nezavisna prověrka skana. Prosta objasnjajuča rečenica čuva razliku, ktoru sama notacija ne izrekla by.\par
\noindent\emph{Reason:} Components as rings in their own right, not as ideals is an explanatory qualification about discriminants in Paper 31. The compared text uses prsteny instead of archival kolca and distinguishes the components as independent rings from ideals of the larger $\overline{\mathfrak R}$. This determines which discriminants are multiplied. The rationale says the parenthesis was recovered from the printed source; that remains the record’s historical claim, not a fresh independent scan check. Plain explanatory prose preserves a distinction that notation alone would not express.\par
\noindent{\footnotesize Izvorne zapisy / Original records: HIST-g145-006.}\par
\smallskip\noindent\emph{Sravnjenje s tekstom:} Odlomok raboty 31 izrično razgledaje komponente kak prsteny same po sebě, ne kak idealy většego prstena. Ta razlika upravjaje interpretacijeju diskriminantnogo produkta. Arhivno pripisanje k pečatnomu iztočniku ne jest s tym samym nezavisno prověreno.\par
\noindent\emph{Textual comparison:} The Paper 31 passage explicitly considers the components as rings in their own right, not as ideals of the larger ring. That distinction governs interpretation of the discriminant product. The archival attribution to the printed witness is not independently verified by this comparison.\par
\noindent{\footnotesize TEXT-COMP-171; pasaže / passages: 18441; \texttt{CORPUS\_FORM\_COMPARISONS.json}.}\par
\end{minipage}\par
\par\medskip\Needspace{7\baselineskip}
\noindent\begin{minipage}{\linewidth}
\hypertarget{TERM-1028}{}
\noindent\textbf{Grupovanje komponentov}\par\noindent\textbf{Grouping components}\quad{\footnotesize\texttt{TERM-1028}}\par
\smallskip\noindent\emph{Razlog:} Shvatiti komponenty v grupy jest naměrny opis s sobirateljnym smyslom v rabotě 19. Sam razlog označaje glagol kak predloženy izbor v stvorjenom odbornom registru, otvorjen jezykovoj ocěně. Lema IV v nyněšnjem tekstě govori, že komponenty grupujut se, a za vsaku grupu tvori se najmenše obče mnogokratno. Zato ne ide o slučajno sobiranje ili umnoženje bez uslovij. Predmetom ocěny jest vhodnost shvatiti za taj postupok; razumlivost ne jest potvrđena čitateljskim izpytanjem.\par
\noindent\emph{Reason:} Shvatiti komponenty v grupy is an intentional collecting expression in Paper 19. The rationale itself marks the verb as a constructed-register proposal open to linguistic appraisal. Lemma IV in the present text says that the components are grouped and the least common multiple is formed for each group. This is not arbitrary collection or multiplication without conditions. The issue for appraisal is whether shvatiti suitably expresses that operation; no reader experiment establishes its intelligibility.\par
\noindent{\footnotesize Izvorne zapisy / Original records: HIST-g081-039, HIST-g082-039.}\par
\smallskip\noindent\emph{Sravnjenje s tekstom:} Lema IV govori o grupovanju komponentov redukovanogo prědstavjenja i tvorjenju najmenšego občego mnogokratnogo. Nyněšnje grupujut se čuva postupok, dokolě arhivno shvatiti ostaje predmetom jezykovej ocěny.\par
\noindent\emph{Textual comparison:} Lemma IV concerns grouping components of a reduced representation and forming the least common multiple. The present grupujut se preserves the operation; archival shvatiti remains open to linguistic appraisal.\par
\noindent{\footnotesize TEXT-COMP-172; pasaže / passages: 12741; \texttt{CORPUS\_FORM\_COMPARISONS.json}.}\par
\end{minipage}\par
\par\medskip\Needspace{7\baselineskip}
\noindent\begin{minipage}{\linewidth}
\hypertarget{TERM-1038}{}
\noindent\textbf{Sohranjenje citovanogo německogo termina}\par\noindent\textbf{Preserving a quoted German term}\quad{\footnotesize\texttt{TERM-1038}}\par
\smallskip\noindent\emph{Razlog:} Kongruenzkomposition ostaje německym slovom, zato že jest citovany historijny termin u Brandta, ne novy medžuslovjansky predlog. Podnožna primětka raboty 19 svazuje jeho s najmenšim občim mnogokratnym klasov. Sohranjenje originalnoj formy pomaga najdti imenovanu rabotu i razlikovati avtorov termin od objašnjenja prevoditelja. Okolna proza može byti polno medžuslovjanska; to ne trěbuje prevoda samogo identifikacijnogo citata.\par
\noindent\emph{Reason:} Kongruenzkomposition remains German because it is Brandt’s cited historical term, not a new Interslavic proposal. The Paper 19 footnote connects it with the least common multiple of classes. Keeping the original form lets readers locate the cited work and distinguish its terminology from the translator’s explanation. Fully localized surrounding prose does not require translating the identifying quotation itself.\par
\noindent{\footnotesize Izvorne zapisy / Original records: HIST-g081-102.}\par
\smallskip\noindent\emph{Sravnjenje s tekstom:} Primětka raboty 19 identificira Kongruenzkomposition pri Brandtu kak najmenše obče mnogokratno klasov. Německa forma služi točnomu odkazu, ne predstavjaje neprevedenu obyčnu prozu.\par
\noindent\emph{Textual comparison:} The Paper 19 footnote identifies Brandt’s Kongruenzkomposition with the least common multiple of classes. The German form serves precise reference; it is not untranslated ordinary prose.\par
\noindent{\footnotesize TEXT-COMP-176; pasaže / passages: 13711; \texttt{CORPUS\_FORM\_COMPARISONS.json}.}\par
\end{minipage}\par
\par\medskip\Needspace{7\baselineskip}
\noindent\begin{minipage}{\linewidth}
\hypertarget{TERM-1047}{}
\noindent\textbf{Konstrukcija i ranějši dokaz obstajanja}\par\noindent\textbf{Construction and a prior existence proof}\quad{\footnotesize\texttt{TERM-1047}}\par
\smallskip\noindent\emph{Razlog:} Konstrukcija po zadannom dobrom uporjadkovanju jest izbor dokaznogo izloženja, ne samo nazvanje novogo objekta. Početok §7 raboty 9 izrično razlikuje ranějši dokaz obstajanja racionalnoj bazy s dodatnym uslovjem od metody jej konstrukcije pri zadannom poredku. Zato kratka zaměna slovami baza obstaje izbrisala by pytanje, na ktore paragraf odgovarja. Dobry poredok trěba razuměti v množinovo-teoretičnom smyslu, ne kak ocěnu udobnosti. Konstrukcija tu ne znači sama po sobě konečny računalny algoritm.\par
\noindent\emph{Reason:} Konstrukcija po zadannom dobrom uporjadkovanju is a choice in proof exposition, not merely the name of a new object. Paper 9, §7 explicitly distinguishes the earlier existence proof for a rational basis with an additional condition from its construction under a specified well-ordering. Replacing the phrase with a basis exists would erase the question being answered. Well-ordering has its set-theoretic meaning, not an evaluative sense of good order; construction does not itself promise a finite computational algorithm.\par
\noindent{\footnotesize Izvorne zapisy / Original records: HIST-g065-002.}\par
\smallskip\noindent\emph{Sravnjenje s tekstom:} §7 raboty 9 razlikuje dokaz obstajanja od konstrukcije pri zadannom dobrom uporjadkovanju. §3 raboty 10 izrično prirědžaje bazove elementy s pomočju drugogo dobrogo uporjadkovanja. Reč konstrukcija ne obeščava konečny algoritm.\par
\noindent\emph{Textual comparison:} Paper 9, §7 distinguishes existence from construction under a specified well-order. Paper 10, §3 explicitly assigns basis elements using another well-order. The word construction does not promise a finite algorithm.\par
\noindent{\footnotesize TEXT-COMP-183; pasaže / passages: 7839, 7840, 7841, 7842, 8518; \texttt{CORPUS\_FORM\_COMPARISONS.json}.}\par
\end{minipage}\par
\par\medskip\Needspace{7\baselineskip}
\noindent\begin{minipage}{\linewidth}
\hypertarget{TERM-1065}{}
\noindent\textbf{Kratke komunikacije kako bibliografična kategorija}\par\noindent\textbf{Short communications as a bibliographical category}\quad{\footnotesize\texttt{TERM-1065}}\par
\smallskip\noindent\emph{Razlog:} Kratke komunikacije jest bibliografična kategorija, ne novy matematičny pojem. V spisku ona stoji pored recenzij, a pod njoju sut kratke objavy o invariantah i drugih temah. Arhivny razlog koristi obyčny naziv za kratku naučnu komunikaciju, vměsto preobrazovanja naslova kategorije v strogy termin. Točnost tu znači čuvati rod publikacije i odnošenje k recenzijam; ne znači pripisati vsakomu zapisu polny raziskovateljsky članok.\par
\noindent\emph{Reason:} Kratke komunikacije is a bibliographical category, not a new mathematical concept. It appears alongside reviews and covers short notices about invariants and other subjects. The recorded reason favors an ordinary description of a short scholarly communication rather than treating the heading as a formal mathematical term. Accuracy here means retaining the publication type and its distinction from reviews, not representing each item as a full research article.\par
\noindent{\footnotesize Izvorne zapisy / Original records: HIST-g215-002.}\par
\smallskip\noindent\emph{Sravnjenje s tekstom:} Spisok razlikuje kratke komunikacije od recenzij. Naslov trěba čitati kako žanr objavy, ne kak novy tehničny termin.\par
\noindent\emph{Textual comparison:} The list distinguishes short communications from reviews. The heading identifies a publication type rather than a new technical term.\par
\noindent{\footnotesize TEXT-COMP-195; pasaže / passages: bib-isv.tex:173, bib-isv.tex:176, bib-isv.tex:179, bib-isv.tex:180; \texttt{CORPUS\_FORM\_COMPARISONS.json}.}\par
\end{minipage}\par
\par\medskip\Needspace{7\baselineskip}
\noindent\begin{minipage}{\linewidth}
\hypertarget{TERM-1123}{}
\noindent\textbf{Lěvy gorny kut v shemě}\par\noindent\textbf{The upper-left corner in a diagram}\quad{\footnotesize\texttt{TERM-1123}}\par
\smallskip\noindent\emph{Razlog:} Lěvy gorny kut dava čitatelju směr v pravougolnoj shemě, ne definiciju novogo algebraičnogo objekta. V rabotě 2 prěhod na pravo označa kontrakciju I, a prěhod v niz kontrakciju II; forma v lěvom gornom kutu začina novy red form. Arhivny razlog jest kratky: prěmy jezyk diagramy. Sravnjenje s ruskim ugol i ukrajinskim kut pokazuje različne domače nazvy togo města. Pri novom razměščenju shemy trěba čuvati směrje i početkovu formu; sam naziv ne može napraviti obratenu diagramu pravilnoju.\par
\noindent\emph{Reason:} Lěvy gorny kut directs the reader within a rectangular arrangement; it does not define a new algebraic object. In Paper 2, moving right denotes contraction I and moving down denotes contraction II; the form in the upper-left corner starts a new series of forms. The recorded reason is simply direct diagram language. Russian ugol and Ukrainian kut supply different native names for that position. A redesigned arrangement must preserve the directions and starting form: correct wording cannot make a reversed diagram correct.\par
\noindent{\footnotesize Izvorne zapisy / Original records: HIST-g003-014.}\par
\smallskip\noindent\emph{Sravnjenje s tekstom:} Rabota 2 svezuje prěhod na pravo s kontrakcijeju I, a prěhod v niz s kontrakcijeju II. Kut označa početkovu formu podshemy; to jest navod dlja čitanja, ne nova terminologična definicija.\par
\noindent\emph{Textual comparison:} Paper 2 assigns contraction I to movement right and contraction II to movement down. The corner identifies the starting form of a subarrangement; this is a reading instruction, not a new terminological definition.\par
\noindent{\footnotesize TEXT-COMP-236; pasaže / passages: 356, 360, 362, 366; \texttt{CORPUS\_FORM\_COMPARISONS.json}.}\par
\end{minipage}\par
\par\medskip\Needspace{7\baselineskip}
\noindent\begin{minipage}{\linewidth}
\hypertarget{TERM-1155}{}
\noindent\textbf{Obseg izključenja}\par\noindent\textbf{Scope of an exception}\quad{\footnotesize\texttt{TERM-1155}}\par
\smallskip\noindent\emph{Razlog:} S izključenjem jest konstrukcija za ograničenje tvrdnje, ne novy matematičny objekt. Arhivny razlog govori o očekyvanoj prozračnosti i odělno priznaje neprověrenu idiomatičnost medžuslovjanske frazy; ne dava imenovanoj konkurenčnoj konstrukcije. V §23 raboty 2 izraz izključaje točno jednu zapisanu formu iz danogo načina redukcije i srazu opisuje, iz kakogo člena ona nastavaje kontrakcijeju. Zato nužno jest čuvati obseg izključenja i sledovajuče objasnjenje. Ne směje se iz toga učiniti izključenje iz vsakej redukcije ili iz cělogo spiska nezvodimyh form.\par
\noindent\emph{Reason:} S izključenjem restricts a statement rather than naming a new mathematical object. The archived reason invokes anticipated transparency and separately acknowledges that the idiomaticity of the Interslavic phrase is unverified; it identifies no competing construction. In §23 of Paper 2, the phrase excepts one displayed form from the stated route of reduction and immediately says which term produces it by contraction. The scope of the exception and that following explanation must both be preserved. It must not become an exception to every reduction or to the entire list of irreducible forms.\par
\noindent{\footnotesize Izvorne zapisy / Original records: HIST-g024-008.}\par
\smallskip\noindent\emph{Sravnjenje s tekstom:} Vyraženje s izključenjem ograničaje konkretny opis redukcije. Izključena forma ostaje zapisana v cělosti, vměstě s členom, iz ktorogo nastavaje kontrakcijeju.\par
\noindent\emph{Textual comparison:} The exception phrase restricts a particular account of reduction. The excepted form remains specified in full, together with the term from which contraction produces it.\par
\noindent{\footnotesize TEXT-COMP-259; pasaže / passages: 3086, 3088, 3090; \texttt{CORPUS\_FORM\_COMPARISONS.json}.}\par
\end{minipage}\par
\par\medskip\Needspace{7\baselineskip}
\noindent\begin{minipage}{\linewidth}
\hypertarget{TERM-1156}{}
\noindent\textbf{Avtorstvo dodatka}\par\noindent\textbf{Authorship of a supplement}\quad{\footnotesize\texttt{TERM-1156}}\par
\smallskip\noindent\emph{Razlog:} Zapis o dodatku k Kapfererovoj rabotě pripisyvaje sovměstnu rabotu s E. Noether dodatku, ne avtomatično cělomu članku. To jest sučstvenny razlog arhivnogo izbora: čuvati granicu avtorstva. Arhiv imaje společně s, nynějšna bibliografija spolu s; slovnik pokazuje obě formy bez prěpisyvanja bibliografičnoj zaslugy. Cěla primětka dolžna ostati vezana s dodatkom. Spisok dvoh avtorov bez toj primětki byl by kratši, ale skryl by informaciju, ktoru zapis izrično hoče čuvati.\par
\noindent\emph{Reason:} The note about the supplement to Kapferer's paper attributes joint work with E. Noether to the supplement, not automatically to the whole article. Preserving that authorship boundary is the substantive reason recorded for the choice. The archive has společně s and the current bibliography spolu s; the dictionary shows both forms without reallocating bibliographical credit. The whole note must remain attached to the supplement. A bare list of two authors would be shorter but would suppress the information the record expressly seeks to preserve.\par
\noindent{\footnotesize Izvorne zapisy / Original records: HIST-g215-012.}\par
\smallskip\noindent\emph{Sravnjenje s tekstom:} Bibliografija pripisyvaje glavnu rabotu H. Kapfereru i dodava primětku o dodatku spolu s E. Noether. Spolu jest nynějšna forma; arhivny zapis společně ne měnjaje obseg sovměstnogo avtorstva.\par
\noindent\emph{Textual comparison:} The bibliography attributes the main paper to H. Kapferer and adds a note about a supplement jointly with E. Noether. Spolu is the current wording; the archived společně does not change the scope of joint authorship.\par
\noindent{\footnotesize TEXT-COMP-260; pasaže / passages: bib-isv.tex:167, bib-isv.tex:168, bib-isv.tex:169, bib-isv.tex:170; \texttt{CORPUS\_FORM\_COMPARISONS.json}.}\par
\end{minipage}\par
\par\medskip\Needspace{7\baselineskip}
\noindent\begin{minipage}{\linewidth}
\hypertarget{TERM-1160}{}
\noindent\textbf{Konstrukcija s modulo}\par\noindent\textbf{The construction with modulo}\quad{\footnotesize\texttt{TERM-1160}}\par
\smallskip\noindent\emph{Razlog:} Modulo jest kratka medžunarodna predložna konstrukcija, izbrana vměsto odvojenyh fraz po modulju i za modulem v ruskom i ukrajinskom sravnjenju. Arhivny razlog govori o obšteslovjanskoj razumlivosti, ale ne sodrživaje čitateljskogo izslědovanja; to jest očekyvanje, ne izměreny rezultat. V uvodu raboty 2 modulo upotrěbjaje se pri odnosno polnyh sistemah form po odnošenju k $\nu$, a potom k $(s,\varrho,t)$ i $(\varrho,t)$. Zato slovnik ne směje ograničiti frazu na ostatok pri děljenju cělyh čisel. Symboli poslě modulo i lokalna definicija redukcije noset matematičny smysl.\par
\noindent\emph{Reason:} Modulo is the concise international prepositional construction selected in place of the separate Russian and Ukrainian phrases po modulyu and za modulem. The archived reason claims pan-Slavic intelligibility but contains no reader study; this is an expectation, not a measured result. In the introduction to Paper 2, modulo occurs with relatively complete systems of forms, first with respect to $\nu$ and then to $(s,\varrho,t)$ and $(\varrho,t)$. The dictionary must therefore not restrict the phrase to remainders in integer division. The symbols following modulo and the local definition of reduction carry the mathematical meaning.\par
\noindent{\footnotesize Izvorne zapisy / Original records: HIST-g002-004.}\par
\smallskip\noindent\emph{Sravnjenje s tekstom:} Uvod raboty 2 postupno měnjaje moduly pri konstrukciji odnosno polnyh sistemov form. Medžunarodna predložna konstrukcija čuva te symbole; interpretacija samo kako cěločislny ostatok byla by prěuzka.\par
\noindent\emph{Textual comparison:} The introduction to Paper 2 successively changes the moduli in constructing relatively complete systems of forms. The international prepositional construction preserves those symbols; interpreting it only as an integer remainder would be too narrow.\par
\noindent{\footnotesize TEXT-COMP-264; pasaže / passages: 182, 184, 186, 188, 190, 192, 195, 196, 198; \texttt{CORPUS\_FORM\_COMPARISONS.json}.}\par
\end{minipage}\par
\par\medskip\Needspace{7\baselineskip}
\noindent\begin{minipage}{\linewidth}
\hypertarget{TERM-1161}{}
\noindent\textbf{Kongruencija po dvom označenym modulam}\par\noindent\textbf{Congruence with respect to two specified moduli}\quad{\footnotesize\texttt{TERM-1161}}\par
\smallskip\noindent\emph{Razlog:} Kratka konstrukcija mod $(i,J)$ čuva standardny zapis kongruencije i jego dva označena modula. Arhivny razlog dava prědnost toj kompaktnosti i soglasnosti s §17 raboty 2 prěd polnymi frazami iz ruskogo i ukrajinskogo sravnjenja. V §20 ona uvodi tri relacije izvedene iz formul (24) i (26). Smysl ne jest cěločislno děljenje: $i$ i $J$ pripadajut lokalnoj teoriji form. Zato pri čitanju mora ostati vidima cela para $(i,J)$, a ne samo slovo mod. To jest konvencija matematičnogo izlaganja, ne novy naziv objekta.\par
\noindent\emph{Reason:} The compact construction mod $(i,J)$ preserves the standard congruence notation and its two specified moduli. The archived reason favours this brevity and consistency with §17 of Paper 2 over the expanded phrases in the Russian and Ukrainian comparison. In §20 it introduces three relations derived from formulas (24) and (26). The meaning is not integer division: $i$ and $J$ belong to the local theory of forms. The whole pair $(i,J)$ must therefore remain visible, not merely the word mod. This is a convention of mathematical exposition, not a newly coined object name.\par
\noindent{\footnotesize Izvorne zapisy / Original records: HIST-g021-012.}\par
\smallskip\noindent\emph{Sravnjenje s tekstom:} Uvodna fraza čuva $(i,J)$ pri trěh relacijah. Jezyk kongruencije zavisi od definicij form v rabotě, ne od obyčnogo cěločislnogo děljenja.\par
\noindent\emph{Textual comparison:} The introductory phrase preserves $(i,J)$ for three relations. The congruence language depends on the paper's definitions of forms, not on ordinary integer division.\par
\noindent{\footnotesize TEXT-COMP-265; pasaže / passages: 2766, 2767, 2771, 2775, 2776, 2777; \texttt{CORPUS\_FORM\_COMPARISONS.json}.}\par
\end{minipage}\par
\par\medskip\Needspace{7\baselineskip}
\noindent\begin{minipage}{\linewidth}
\hypertarget{TERM-1197}{}
\noindent\textbf{Odkaz i děliteljnost v dokazě}\par\noindent\textbf{Reference and divisibility in a proof}\quad{\footnotesize\texttt{TERM-1197}}\par
\smallskip\noindent\emph{Razlog:} Po punktu 1 čuva odkaz, na ktorom opiraje se tvrdženje o děliteljnosti vsih razgledanyh determinantov. V §8 raboty 31 vybrana jest lokalna modulna baza, a jej slědovy determinant děli druge determinanty v ukazanom lokalnom prstenu. Vsaky i nazvanje prstena sut suštvene: děliteljnost nad prstenom ne možno zaměniti neograničenym děljenjem v polju. Arhivny razlog tumači odkaz kako §8,1. Dostupny tekst v samoj stvari piše samo po punktu 1; §8,1 ustanovjaje lokalnu osnovu teorije, a §8,2 izrično govori o determinantah prědstavjenyh prěz bazu. Bez sravnjenja s pečatnym pasusom ne podavamo podrobnějšu identifikaciju odkaza kako dokazanu. To jest izbor dokaznoj prozy i predmet kontroly odkaza, ne samostojny novy termin.\par
\noindent\emph{Reason:} Po punktu 1 preserves a reference supporting the divisibility assertion for all the determinants under consideration. In §8 of Paper 31 a local module basis is chosen, and its trace determinant divides the other determinants in the specified local ring. Every and the ring specification are essential: divisibility in a ring must not be replaced by unrestricted division in a field. The archived reason identifies the reference as §8,1. The available text actually says only according to point 1; §8,1 establishes the local setting, whereas §8,2 expressly discusses expressing determinants through a basis. Without comparing the printed passage, the more detailed reference identification is not treated as established. This is a choice in proof exposition and a cross-reference question, not a separate new term.\par
\noindent{\footnotesize Izvorne zapisy / Original records: HIST-g162-004.}\par
\smallskip\noindent\emph{Sravnjenje s tekstom:} Dostupny pasus čuva kratky odkaz, ale ne nazyvaje paragraf. Sravnjene pasusy pokazyvajut predposylky poredkov, lokalny prsten glavnyh idealov i prědstavjenje determinantov. Podrobna identifikacija odkaza ostaje pitanjem za pečatny izvornik.\par
\noindent\emph{Textual comparison:} The available passage preserves a short reference without naming the section. The compared passages supply the assumptions on orders, a local principal ideal ring and the expression of determinants. The detailed reference identification remains a printed-witness question.\par
\noindent{\footnotesize TEXT-COMP-300; pasaže / passages: 18559, 18561, 18789, 18823, 18897, 19037; \texttt{CORPUS\_FORM\_COMPARISONS.json}.}\par
\end{minipage}\par
\par\medskip\Needspace{7\baselineskip}
\noindent\begin{minipage}{\linewidth}
\hypertarget{TERM-1198}{}
\noindent\textbf{Uvod razloženja v §20}\par\noindent\textbf{Introducing a decomposition in §20}\quad{\footnotesize\texttt{TERM-1198}}\par
\smallskip\noindent\emph{Razlog:} Po slědujućej shemě jest arhivny prěvod nach dem Ansatz v §20 raboty 2. Razlog jest ne prěnositi německo slovo, kogda formula pokazuje način razloženja. Rusko i ukrajinsko sravnjena imenno govoret o razloženju. Dostupny korpus v odgovarjajučem městě piše po slědujučem nastavku i davaje dva prědstavjenja produktov, prinaležne dvojnoj redukciji. Shema i nastavok zato sut dvě dokumentovane formulacije, ne dokazano odbačena i pobjědna alternativa s izvěstnymi vagami. Razlog proměny ne jest ustanovjen etym sravnjenjem. Čitatelj mora viděti, že izraz uvodi konkretne formule; ne mora poznavati německy termin Ansatz ni tumačiti jego kako nedokazanu matematičnu dogadku.\par
\noindent\emph{Reason:} Po slědujućej shemě is the archived rendering of nach dem Ansatz in §20 of Paper 2. Its reason is to avoid importing the German word where the formula supplies a decomposition pattern. The Russian and Ukrainian comparanda specifically refer to decomposition. At the corresponding location the available corpus says po slědujučem nastavku and gives two representations of products used in a double reduction. Shema and nastavok are therefore two documented formulations, not a demonstrated rejected/winning pair with known weights. This comparison does not establish why the wording differs. Readers need to understand that the expression introduces particular formulas; they need neither know the German word Ansatz nor interpret it as an unproved mathematical conjecture.\par
\noindent{\footnotesize Izvorne zapisy / Original records: HIST-g021-009.}\par
\smallskip\noindent\emph{Sravnjenje s tekstom:} §20 daje dvije formule za dvojnu redukciju. Nazvanje uvoda različaje se od arhivnoj glosy; samo to sravnjenje ne ustanovjaje pričinu izbora nastavok.\par
\noindent\emph{Textual comparison:} Section 20 gives two formulas for a double reduction. Its introductory wording differs from the archived gloss; this comparison alone does not establish the reason for choosing nastavok.\par
\noindent{\footnotesize TEXT-COMP-301; pasaže / passages: 2706, 2728, 2730, 2731, 2733, 2736, 2737, 2738, 2739, 2740, 2741; \texttt{CORPUS\_FORM\_COMPARISONS.json}.}\par
\end{minipage}\par
\par\medskip\Needspace{7\baselineskip}
\noindent\begin{minipage}{\linewidth}
\hypertarget{TERM-1199}{}
\noindent\textbf{Uvod vyčisljenja v §9}\par\noindent\textbf{Introducing a calculation in §9}\quad{\footnotesize\texttt{TERM-1199}}\par
\smallskip\noindent\emph{Razlog:} Po slědujučej začetnoj shemě uvodi vyčisljenje v §9 raboty 2. Pridavny izraz začetna izjasnjaje, že ide o početkovy način postavjenja vyčisljenja, a slědujuča odkazuje k bezposrědnje danym formulam. Arhivny razlog bira opis vměsto neprozračnoj německoj zaimky i ostavjaje prirodnost izraza za jezykovu ocěnku. V dostupnom tekstu izraz stojí pri dvojnoj redukciji trěh izrazov iz reducenta; ne ustanovjaje novy objekt algebraičnoj teorije. Arhivna forma s ć i korpusna forma s č čuvajut isto naměrene značenje. Izbor služi orientaciji v dokazě, a ne skraćenju samogo vyčisljenja.\par
\noindent\emph{Reason:} Po slědujučej začetnoj shemě introduces a calculation in §9 of Paper 2. Začetna explains that this is the initial setup of the calculation, while slědujuča points to the immediately following formulas. The archived reason chooses a description instead of an opaque German loan and leaves the expression's naturalness open to linguistic assessment. In the available text it accompanies a double reduction of three expressions from a reducent; it does not introduce a new object of algebraic theory. The archived ć and the corpus č retain the same intended meaning. The choice helps readers navigate the proof; it does not abbreviate the calculation itself.\par
\noindent{\footnotesize Izvorne zapisy / Original records: HIST-g010-005.}\par
\smallskip\noindent\emph{Sravnjenje s tekstom:} Uvodna fraza v §9 i §11 ne zaměnjaje vyčisljenja. Primětka v §11 izrično zabranjaje čitati sim kako ravnost pri izpuščenyh produktah.\par
\noindent\emph{Textual comparison:} The introductory phrase in §§9 and 11 does not replace the calculations. The §11 footnote expressly rules out reading sim as equality when products have been omitted.\par
\noindent{\footnotesize TEXT-COMP-302; pasaže / passages: 1229, 1231, 1232, 1233, 1234, 1235, 1236, 1237, 1238, 1239, 1560, 1562, 1563, 1564, 1565, 1566, 1567, 1568, 1569, 1570, 1464; \texttt{CORPUS\_FORM\_COMPARISONS.json}.}\par
\end{minipage}\par
\par\medskip\Needspace{7\baselineskip}
\noindent\begin{minipage}{\linewidth}
\hypertarget{TERM-1200}{}
\noindent\textbf{Uvod vyčisljenja i skraćena notacija v §11}\par\noindent\textbf{Introducing a calculation and abbreviated notation in §11}\quad{\footnotesize\texttt{TERM-1200}}\par
\smallskip\noindent\emph{Razlog:} Po začetnoj shemě v §11 raboty 2 prodolžaje opisnu strategiju iz §9 za německo nach dem Ansatz. Arhivna forma začętnoj imaje nosovu bukvu; dostupny korpus piše začetnoj. Razlog izbora jest objasniti početkovu formu vyčisljenja bez potrěby znati německy izraz; prikladnost same formulacije ostala jest otvorjena. Osoblivo važno jest čitati primětku pri etom městě: znak $\sim$ ne označa ravnost, zato že opreděljene produkty s $u_x$ i vyšimi formami sut izpuščene. Prěvod uvodnoj frazy ne može čitatelju sugerovati točnu identičnost tam, kde avtor uživa skraćeno prědstavjenje. Notacija i njeno objasnjenje morajut ostati vměstě. Arhivna oznaka jedinicy odsyla k §12; citovany naslov dostupnogo teksta odvojeno ustanovjaje, že tu razgledano město naleži §11.\par
\noindent\emph{Reason:} Po začetnoj shemě in §11 of Paper 2 continues the descriptive treatment of German nach dem Ansatz used in §9. The archived začętnoj has a nasal vowel letter; the available corpus writes začetnoj. The choice aims to explain the initial form of the calculation without requiring knowledge of the German expression; the suitability of the wording itself remains open. The accompanying footnote is particularly important: $\sim$ is not equality, because specified products involving $u_x$ and higher forms are omitted. Translating the introductory phrase must not suggest an exact identity where the author uses an abbreviated representation. The notation and its explanation must remain together. The archival unit label points to §12; the quoted heading of the available text separately establishes that the passage discussed here belongs to §11.\par
\noindent{\footnotesize Izvorne zapisy / Original records: HIST-g013-009.}\par
\smallskip\noindent\emph{Sravnjenje s tekstom:} Uvodna fraza v §9 i §11 ne zaměnjaje vyčisljenja. Primětka v §11 izrično zabranjaje čitati sim kako ravnost pri izpuščenyh produktah.\par
\noindent\emph{Textual comparison:} The introductory phrase in §§9 and 11 does not replace the calculations. The §11 footnote expressly rules out reading sim as equality when products have been omitted.\par
\noindent{\footnotesize TEXT-COMP-302; pasaže / passages: 1229, 1231, 1232, 1233, 1234, 1235, 1236, 1237, 1238, 1239, 1560, 1562, 1563, 1564, 1565, 1566, 1567, 1568, 1569, 1570, 1464; \texttt{CORPUS\_FORM\_COMPARISONS.json}.}\par
\end{minipage}\par
\par\medskip\Needspace{7\baselineskip}
\noindent\begin{minipage}{\linewidth}
\hypertarget{TERM-1201}{}
\noindent\textbf{Uvod rezultata vyčisljenja}\par\noindent\textbf{Introducing a calculated result}\quad{\footnotesize\texttt{TERM-1201}}\par
\smallskip\noindent\emph{Razlog:} Fraza po sootvětnom računu uvodi rezultat vyčisljenja, a ne novy matematičny pojem. Arhivny razlog davaje prědnost obyčnoj dokazovoj prozě; rusky i ukrajinsky sravnjajemy izraz takože govori o odgovarjajučem vyčisljenju. Izbor računu ostal jest privremenny: zapis izrično potrěbuje sravnjenja s izračunanje i računati, bez čislovyh vah. V §11 raboty 2 dostupny tekst piše po odgovarjajučem računu kako vyše i potom navodi formuly. Kako vyše odsyla k podobnomu postupu, ne k doslovnomu povtorjenju vsego dokaza. Fraza ne zaměnjaje formuly i ne dokazuje, že izračun byl nezavisno prověren. Sravnjenje pokazyvaje kontinuitet izloženja i različje medžu arhivnym sootvětnom i tekstovym odgovarjajučem.\par
\noindent\emph{Reason:} The phrase po sootvětnom računu introduces the result of a calculation, not a new mathematical concept. The recorded reason favors ordinary proof prose; the Russian and Ukrainian comparanda likewise refer to a corresponding calculation. Računu remained provisional: the record expressly calls for comparison with izračunanje and računati, without numerical weights. In §11 of Paper 2 the available text reads po odgovarjajučem računu kako vyše and then gives formulas. Kako vyše points to an analogous procedure, not a verbatim repetition of the entire proof. The phrase neither replaces the formulas nor establishes that the calculation has been independently checked. The comparison documents continuity of exposition and the difference between archived sootvětnom and textual odgovarjajučem.\par
\noindent{\footnotesize Izvorne zapisy / Original records: HIST-g012-015.}\par
\smallskip\noindent\emph{Sravnjenje s tekstom:} Fraza po odgovarjajučem računu kako vyše odsyla k analogičnomu postupu i neposrědnje prědhodi konkretnym formulam. Naslov paragrafa potvrđuje §11.\par
\noindent\emph{Textual comparison:} Po odgovarjajučem računu kako vyše points to an analogous procedure and immediately precedes concrete formulas. The section heading confirms §11.\par
\noindent{\footnotesize TEXT-COMP-303; pasaže / passages: 1464, 1668, 1671, 1672, 1673; \texttt{CORPUS\_FORM\_COMPARISONS.json}.}\par
\end{minipage}\par
\par\medskip\Needspace{7\baselineskip}
\noindent\begin{minipage}{\linewidth}
\hypertarget{TERM-1202}{}
\noindent\textbf{Avtorova dodatna primětka}\par\noindent\textbf{The author's supplementary note}\quad{\footnotesize\texttt{TERM-1202}}\par
\smallskip\noindent\emph{Razlog:} Dodatna primětka jest naslov avtorovoj primětky poslě teoremy, ne novo objasnjenje prěvoditelja. Zapis daje prědnost prirodnoj opisnoj frazě prěd doslovnym prěvodom německogo Nachbemerkung; rusky i ukrajinsky ekvivalent takože podpirajut tu funkciju naslova. V §1 raboty 2 primětka sodrživaje izključenje k teoremě I, kogda dva ukazane parametry stanut nula. Ona utočnjaje, ktore kontrakcije ostavajut, kako razuměti red form i začto teoremy o reducentah ostavajut v silě. Dodatna zato ne znači nepotrěbna: primětka ograničava obseg prědhodečego tvrđenja. Kratky naslov primětka v drugom městě ima podobnu izložiteljnu funkciju, ale ne davaje pravo izgubiti tuty uslovja ili směšati avtorov tekst s redakcijnym komentarjem.\par
\noindent\emph{Reason:} Dodatna primětka labels the author's subsequent note, not an explanation added by the translator. The record favors a natural descriptive phrase over a literal rendering of German Nachbemerkung; the Russian and Ukrainian comparanda support that function. In §1 of Paper 2 the note gives an exception to Theorem I when two specified parameters become zero. It explains which contractions remain, how to understand the series of forms, and why the theorems concerning reducents still hold. Additional therefore does not mean dispensable: the note qualifies the preceding assertion. The shorter label primětka elsewhere serves a similar expository function, but does not justify dropping these conditions or confusing the author's text with editorial commentary.\par
\noindent{\footnotesize Izvorne zapisy / Original records: HIST-g003-015.}\par
\smallskip\noindent\emph{Sravnjenje s tekstom:} Dodatna primětka ograničava teoremu I; kratka primětka uvodi daljše relacije. Izložiteljny naslov ne davaje povoljenje izpustiti matematične uslovja.\par
\noindent\emph{Textual comparison:} The supplementary note qualifies Theorem I; the short note introduces further relations. An expository heading does not license omission of mathematical conditions.\par
\noindent{\footnotesize TEXT-COMP-304; pasaže / passages: 368, 378, 987, 991, 992, 993; \texttt{CORPUS\_FORM\_COMPARISONS.json}.}\par
\end{minipage}\par
\par\medskip\Needspace{7\baselineskip}
\noindent\begin{minipage}{\linewidth}
\hypertarget{TERM-1203}{}
\noindent\textbf{Kratky naslov primětky}\par\noindent\textbf{The short note heading}\quad{\footnotesize\texttt{TERM-1203}}\par
\smallskip\noindent\emph{Razlog:} Primětka jest kratky izložiteljny naslov za německo Nachbemerkung. Arhivny izbor odklanja doslovnu konstrukciju so značenjem poslje-zaměčanje, da by čitatelj neposrědnje poznal funkciju odstavca. Rusko primečanije i ukrajinska prymitka sut zapisane sravnjajeme formy; to ne jest čislovo glasovanje medžu jezykami. V rabotě 2 tekst pod tym naslovom govori o analogičnom izračunjenju prěkrivanj $f$ nad $j$, reducibilnyh po §5, i davaje relacije (4). Naslov ne měnjaje jih dokazovogo statusa i ne čini jih komentarjem prěvoditelja. Različje medžu kratkym primětka i razširjenym dodatna primětka trěba ocěnjati po funkciji konkretnoj primětky; zapis ne ustanovjaje dvě različne matematične kategorije.\par
\noindent\emph{Reason:} Primětka is a short expository label for German Nachbemerkung. The recorded choice rejects a literal construction meaning after-remark so that readers can recognize the paragraph's function directly. Russian primečanie and Ukrainian prymitka are the recorded comparanda; this is not a numerical vote among languages. In Paper 2 the paragraph under this label describes an analogous calculation of overlaps of $f$ over $j$, reducible by §5, and supplies relations (4). The label changes neither their argumentative role nor their authorship. The difference between primětka and the longer dodatna primětka should be assessed by the function of the particular note; the record establishes no two distinct mathematical categories.\par
\noindent{\footnotesize Izvorne zapisy / Original records: HIST-g007-010.}\par
\smallskip\noindent\emph{Sravnjenje s tekstom:} Dodatna primětka ograničava teoremu I; kratka primětka uvodi daljše relacije. Izložiteljny naslov ne davaje povoljenje izpustiti matematične uslovja.\par
\noindent\emph{Textual comparison:} The supplementary note qualifies Theorem I; the short note introduces further relations. An expository heading does not license omission of mathematical conditions.\par
\noindent{\footnotesize TEXT-COMP-304; pasaže / passages: 368, 378, 987, 991, 992, 993; \texttt{CORPUS\_FORM\_COMPARISONS.json}.}\par
\end{minipage}\par
\par\medskip\Needspace{7\baselineskip}
\noindent\begin{minipage}{\linewidth}
\hypertarget{TERM-1207}{}
\noindent\textbf{Matematične tablice i znak modulo}\par\noindent\textbf{Mathematical arrays and modulo notation}\quad{\footnotesize\texttt{TERM-1207}}\par
\smallskip\noindent\emph{Razlog:} To jest pravilo pisma i notacije, ne slovnikova para dvou jezykov. Pri kiriličnom prěpisu prozy matematične bukvy, grečske znaky i tehnične edinice razstupa ostavajut neizměnjene. V §14 raboty 2 jedna tablica stoi v drugoj; končanje vnutrěnnoj tablicy ne konča matematičny kontekst vněšnjej. Formulacija uključaje takože izraz modulo: znak $\nu$ v kongruenciji jest matematičny simbol, ne slovo za transliteraciju. Arhivny razlog opiraje se na ohranu identičnosti notacije, ne na sravnjenje ruskyh i ukrajinskyh terminov. Prozni tekst vnutri formuly trěba razlikovati od simbolov; jedna obalna matematična struktura sama ne znači, že vse vidime slova sut netransliterovateljne. To pravilo opisuje, čto se mora zachovati, ale ne jest dokazom pravilnosti cělogo kiriličnogo prěpisa.\par
\noindent\emph{Reason:} This is a script and notation rule, not a bilingual lexical pair. When prose is projected into Cyrillic, mathematical letters, Greek symbols, and technical spacing units remain unchanged. In §14 of Paper 2 one array sits inside another; closing the inner array does not end the surrounding mathematical context. The rule also covers modulo expressions: $\nu$ in the congruence is a mathematical symbol, not a word to transliterate. The recorded reason is preservation of notation, not a comparison of Russian and Ukrainian terms. Prose inside a formula still has to be distinguished from symbols; an enclosing mathematical structure does not by itself make every visible word exempt from transliteration. The rule specifies what must be preserved, but is not evidence that the entire Cyrillic projection has been checked.\par
\noindent{\footnotesize Izvorne zapisy / Original records: HIST-g015-017.}\par
\smallskip\noindent\emph{Sravnjenje s tekstom:} Končanje vnutrěnnoj tablicy ne izključaje slědujuče simboly iz vněšnjej matematičnoj struktury. Simbol pri modulo takože ostavaje matematičnym znakom.\par
\noindent\emph{Textual comparison:} Closing an inner array does not remove the following symbols from the outer mathematical structure. The symbol in the modulo expression likewise remains mathematical notation.\par
\noindent{\footnotesize TEXT-COMP-308; pasaže / passages: 2037, 2038, 2039, 2040, 2041, 2042, 2043, 2044, 2162, 2163, 2164, 2165, 2166, 2167, 2168, 2169, 2170, 2171, 2172, 2173, 2212, 2214; \texttt{CORPUS\_FORM\_COMPARISONS.json}.}\par
\end{minipage}\par
\par\medskip\Needspace{7\baselineskip}
\noindent\begin{minipage}{\linewidth}
\hypertarget{TERM-1208}{}
\noindent\textbf{Vložene tablice i identičnost simbolov}\par\noindent\textbf{Nested arrays and symbol identity}\quad{\footnotesize\texttt{TERM-1208}}\par
\smallskip\noindent\emph{Razlog:} Vložena tablica ne měnjaje identičnost bukv izvan sebe. V §13 raboty 2 vnutrěnna tablica konča se prěd izrazami s $H$, $a$ i $j$, ale vse te izrazy ostavajut v obemlječej matematičnoj tablici. One ne smějut stati kiriličnymi dvojnicami podobnogo vida; takože oznaka edinicy em v zapisu razstupa jest tehnična edinica, ne proza. Arhiv dokumentuje razlog različenja vnutrěnnoj i vněšnjej struktury pri prěpisu. Iz togo vyhodi obče pravilo: matematičny kontekst konča se tek, kogda konči se odgovarjajuča obemlječa struktura, ne pri prvom znaku konca vnutri nje. Tuta odluka pripada k vernosti notacije. Ona ne tvori novy matematičny termin i ne davaje kiriličnoj versiji status nezavisnogo prěvoda.\par
\noindent\emph{Reason:} An embedded array does not change the identity of letters outside it. In §13 of Paper 2 the inner array ends before expressions involving $H$, $a$, and $j$, yet all remain inside the enclosing mathematical array. They must not become visually similar Cyrillic letters; likewise em in a spacing instruction is a technical unit, not prose. The archive records the reason for distinguishing inner and outer structure during script conversion. The general rule is that a mathematical context ends only when its corresponding enclosing structure ends, not at the first internal closing marker. This decision concerns fidelity of notation. It coins no mathematical term and does not make the Cyrillic version an independent translation.\par
\noindent{\footnotesize Izvorne zapisy / Original records: HIST-g014-018.}\par
\smallskip\noindent\emph{Sravnjenje s tekstom:} Končanje vnutrěnnoj tablicy ne izključaje slědujuče simboly iz vněšnjej matematičnoj struktury. Simbol pri modulo takože ostavaje matematičnym znakom.\par
\noindent\emph{Textual comparison:} Closing an inner array does not remove the following symbols from the outer mathematical structure. The symbol in the modulo expression likewise remains mathematical notation.\par
\noindent{\footnotesize TEXT-COMP-308; pasaže / passages: 2037, 2038, 2039, 2040, 2041, 2042, 2043, 2044, 2162, 2163, 2164, 2165, 2166, 2167, 2168, 2169, 2170, 2171, 2172, 2173, 2212, 2214; \texttt{CORPUS\_FORM\_COMPARISONS.json}.}\par
\end{minipage}\par
\par\medskip\Needspace{7\baselineskip}
\noindent\begin{minipage}{\linewidth}
\hypertarget{TERM-1212}{}
\noindent\textbf{Kritika dokaza bez zaměny tvrđenja}\par\noindent\textbf{Criticizing a proof without changing the claim}\quad{\footnotesize\texttt{TERM-1212}}\par
\smallskip\noindent\emph{Razlog:} Ne jest dostatočno obosnovany odnosi se k dokazu, ne k istinnosti teorema. Arhivny razlog požadaje izběgati nepotrěbno polemičnogo tona, ne oslabevati konkretnu logičnu kritiku. V primětce raboty 7 Noether kaže, že iz invariantnosti koeficientov jednej funkcije ješče ne jest pokazana jih racionalna izrazimost prěz koeficienty drugoj. Potom daje postup s diferencialovanjem, ktory tu prazninu obhodi. Rusky izraz o nesostojateljnosti i ukrajinsky o nekorektnosti pokazujut različne sposoby izrečti ocěnku dokaza, ale ne opravdajut prěvod teorem jest nepravdivy. Takože nedostatok podrobnosti byl by preslaby, ako skryvaje imenno nevyvedeny krok. Ovdě se objasnjaje kritika, ktora jest v avtorovoj primětce; to ne jest nezavisna prověrka citovanogo Weberovogo dokaza.\par
\noindent\emph{Reason:} Ne jest dostatočno obosnovany concerns the proof, not the truth of the theorem. The recorded reason calls for avoiding unnecessary polemic, not weakening a specific logical criticism. In a footnote to Paper 7, Noether says that invariance of the coefficients of one function has not yet established their rational expressibility through the coefficients of another. She then gives a differentiation procedure that avoids the gap. The Russian expression of invalidity and the Ukrainian expression of incorrectness illustrate different ways of evaluating the proof, but do not warrant translating the statement as the theorem is false. Mere lack of detail would also be too weak if it concealed the missing inference. This explains the criticism in the author's footnote; it is not an independent examination of the cited proof by Weber.\par
\noindent{\footnotesize Izvorne zapisy / Original records: HIST-g056-012.}\par
\smallskip\noindent\emph{Sravnjenje s tekstom:} Primětka razlikuje invariantnost koeficientov od jih racionalnoj izrazimosti prěz druge koeficienty. Kritika odsyla k nevyvedenomu kroku, ne k nepravdivosti teorema.\par
\noindent\emph{Textual comparison:} The footnote distinguishes invariance of coefficients from their rational expressibility through other coefficients. The criticism concerns a missing inference, not falsity of the theorem.\par
\noindent{\footnotesize TEXT-COMP-312; pasaže / passages: 6706, 6708, 6709, 6711, 6713, 6715, 6716, 6717, 6718, 6720; \texttt{CORPUS\_FORM\_COMPARISONS.json}.}\par
\end{minipage}\par
\par\medskip\Needspace{7\baselineskip}
\noindent\begin{minipage}{\linewidth}
\hypertarget{TERM-1223}{}
\noindent\textbf{Komutativnost koeficientov v nekomutativnom prstenu}\par\noindent\textbf{Commuting coefficients in a noncommutative ring}\quad{\footnotesize\texttt{TERM-1223}}\par
\smallskip\noindent\emph{Razlog:} To jest odluka o formulovanju uslovja, ne samostojny novy termin. Německo vertauschbar prěvedeno jest glagolom komutovati, da fraza jasno kaže: elementy $\Sigma$ komutujut so vsimi elementami prstena. Zapis objašnjaje, že ta formulacija čuva smysl avtorovoj podčrtnoj primětky: dani zapisi i dokaz ostavajut uporabimy takože v ukazanom nekomutativnom slučaju. Uslovje ne tvrdži, že prsten v cělosti jest komutativny. Ono dopušča potrěbnu prěměnu poredka koeficientov, bez dopuščanja libovoljnoj prěměny vseh faktorov. V dostupnom tekstu ta primětka stoji v dokazě korespondencije klasov idealov i matričnyh prědstavjenij. Arhivno kolca i nynějšnje prsteny sut različne leksične formy; matematičny symbol $\Sigma$ ne jest slovo za prěvod.\par
\noindent\emph{Reason:} This is a decision about stating a hypothesis, not a separate new term. German vertauschbar is rendered by komutovati so that the sentence says explicitly that the elements of $\Sigma$ commute with every ring element. The record explains the purpose of the author's footnote: the notation and argument also remain applicable in the specified noncommutative setting. The condition does not assert that the whole ring is commutative. It permits the required interchange with coefficients, not arbitrary rearrangement of all factors. In the available text the footnote occurs in the correspondence between ideal classes and matrix-representation classes. Archival kolca and current prsteny are different lexical forms; the mathematical symbol $\Sigma$ is not a word to translate.\par
\noindent{\footnotesize Izvorne zapisy / Original records: HIST-g134-006.}\par
\smallskip\noindent\emph{Sravnjenje s tekstom:} Podčrtna primětka ograničava nekomutativno razširjenje dokaza prědpoloženjem o $\Sigma$; ona ne razširja komutativnost na cěly prsten.\par
\noindent\emph{Textual comparison:} The footnote limits the noncommutative extension of the proof by a hypothesis on $\Sigma$; it does not extend commutativity to the whole ring.\par
\noindent{\footnotesize TEXT-COMP-321; pasaže / passages: 17985, 17989, 17993, 17995; \texttt{CORPUS\_FORM\_COMPARISONS.json}.}\par
\end{minipage}\par
\par\medskip\Needspace{7\baselineskip}
\noindent\begin{minipage}{\linewidth}
\hypertarget{TERM-1267}{}
\noindent\textbf{Latinica i kirilica pri nulj-}\par\noindent\textbf{Latin and Cyrillic treatment of nulj-}\quad{\footnotesize\texttt{TERM-1267}}\par
\smallskip\noindent\emph{Razlog:} Toj zapis prodolžaje terminologiju raboty 24, ale njegov osobny razlog tyče se pismenosti. Latinske nulje i polje nuljev služe osnovoju, a kiriličny zapis ima slědovati jednoj povtorimoj prěobrazbi; arhiv izrično ukazuje lj v ль, bez samovoljnogo pravopisnogo popravljanja v kirilici. Tako dvě pismenosti ne stavajut dvoma nezavisnymi prěvodami. To jest zasada zapisana pri prvonačalnom izborě, ne tvrdženje, že v toj rekonstrukciji byl iznova prověren cěly konverter. Izbor samogo korenja, njegovogo sklanjanja i njegove matematičnoj priměny ostavaje oddělnym lingvističnym voprosom. Rusky i ukrajinsky zapisy mogut podpirati sravnenje, ale avtomatična kirilična kopija medžuslovjanskogo teksta ne može byti tretjim nezavisnym svědkom.\par
\noindent\emph{Reason:} This record continues Paper 24's terminology, but its distinctive reason concerns script. Latin nulje and polje nuljev are the source forms, while Cyrillic is to follow one reproducible conversion; the archive expressly specifies lj to ль, without discretionary spelling changes in Cyrillic. The two scripts therefore do not become independent translations. This is the policy recorded with the original choice, not a claim that the whole converter was freshly tested in this reconstruction. Selection of the stem, its inflection and its mathematical application remain separate linguistic questions. Russian and Ukrainian records can inform comparison, but an automatic Cyrillic copy of the Interslavic text cannot count as a third independent witness.\par
\noindent{\footnotesize Izvorne zapisy / Original records: HIST-g097-001.}\par
\smallskip\noindent\emph{Sravnjenje s tekstom:} Rabota 24, §1, razgleda algebraične vrijednosti nad poljem s ostalymi proměnnymi kako parametrami. Produkt linearnyh faktorov rezultanta čuva takože kratnosti. Zapis o pismenosti v §3 prodolžaje termin nulj-; sam latinicky kontekst ne jest novym testom kiriličnogo konvertera ni nezavisnym svědčenjem jezykovoj upotrěby.\par
\noindent\emph{Textual comparison:} Paper 24, §1, considers algebraic values over a field retaining the other variables as parameters. The resultant's product of linear factors also carries multiplicities. The §3 script record continues nulj- terminology; the Latin context is not itself a new test of the Cyrillic converter or independent language attestation.\par
\noindent{\footnotesize TEXT-COMP-359; pasaže / passages: 15263, 15341, 15343, 15344, 15346; \texttt{CORPUS\_FORM\_COMPARISONS.json}.}\par
\end{minipage}\par
\par\medskip\Needspace{7\baselineskip}
\noindent\begin{minipage}{\linewidth}
\hypertarget{TERM-1272}{}
\noindent\textbf{Samo jedin: jedinstvenost s točnymi uslovjami}\par\noindent\textbf{Only one: uniqueness with its exact hypotheses}\quad{\footnotesize\texttt{TERM-1272}}\par
\smallskip\noindent\emph{Razlog:} Samo jedin izrično čuva jedinstvenost, ne jest stilističny dodatok. Arhiv podčrtaje, že na tom uslovju zavisi zaključok. V rabotě 31, §8, govori se o multiplikacionnom prstenu s jedinym prostym idealom različnym od nulovogo i jediničnogo. V toj klasě prstenov to vodi k prstenu glavnyh idealov. Ne možno izostaviti ni klasu prstenov, ni izključene idealy, i poslě tvrdžiti to o libovoljnom prstenu s jedinym prostym idealom. Fraza takože ne zaměnjaje postojanje samo gornjoju graniceju od jednogo ideala. Ruska i ukrajinska paralela imajut tu istu omeđujuču česticu. To jest odluka o logičnoj točnosti rečenice, ne novy nazov matematičnogo objekta.\par
\noindent\emph{Reason:} Samo jedin explicitly preserves uniqueness; it is not stylistic padding. The record stresses that the consequence depends on this condition. Paper 31, §8, concerns a multiplication ring with just one prime ideal other than the zero and unit ideals. Within that class of rings the conclusion is a principal ideal ring. Neither the ring class nor the excluded ideals may be dropped to assert this of an arbitrary ring with one prime ideal. The phrase also does not replace existence by merely an upper bound of one ideal. The Russian and Ukrainian parallels have the same restrictive particle. This is a decision about a sentence's logical precision, not a new mathematical object-name.\par
\noindent{\footnotesize Izvorne zapisy / Original records: HIST-g156-001.}\par
\smallskip\noindent\emph{Sravnjenje s tekstom:} Rabota 31, §8, najprvo ograničavaje klasu prstenov, potom davaje zaključok pri jedinom prostom idealu različnom od nulovogo i jediničnogo. Pri lokalizaciji povtorjajut se i ta jedinstvenost i svojstva multiplikacionnogo prstena. To objašnjaje začto povtorjenje uslovja jest matematično korisno, a ne izbytočno.\par
\noindent\emph{Textual comparison:} Paper 31, §8, first specifies the ring class, then draws the consequence when there is only one prime ideal other than zero and the unit ideal. The localization passage repeats both uniqueness and the multiplication-ring properties. This explains why repeating the hypothesis is mathematically useful rather than redundant.\par
\noindent{\footnotesize TEXT-COMP-363; pasaže / passages: 18787, 18789, 18823, 18932, 18970; \texttt{CORPUS\_FORM\_COMPARISONS.json}.}\par
\end{minipage}\par
\par\medskip\Needspace{7\baselineskip}
\noindent\begin{minipage}{\linewidth}
\hypertarget{TERM-1275}{}
\noindent\textbf{Povtorjenje uslovja poslě lokalizacije}\par\noindent\textbf{Repeating the hypothesis after localization}\quad{\footnotesize\texttt{TERM-1275}}\par
\smallskip\noindent\emph{Razlog:} To jest povtorjenje nužnogo uslovja v novom městě dokaza. Zapis izrično objašnjaje, že rečenica o jedinom prostom idealu v lokalizovanom prstenu ima pokazati, začto možno priměniti rannějši zaključok o glavnyh idealah. V rabotě 31, §8, lokalizacija čuva potrěbne svojstva multiplikacionnogo prstena; izključene ostavajut nulovy i jediničny idealy. Z toho ne slěduje, že vsaky lokalny prsten jest prstenom glavnyh idealov. Arhivno $o_p$ jest oznaka prstena, ne čislo 0 s indeksom; dostupny tekst koristi frakturne bukvy za prsten i prosty ideal. Pismenost oznaki i matematičny smysl uslovja trěba čuvati oddělno. Ponavljanje tu pomaga čitatelju proslediti priměnu teoremy.\par
\noindent\emph{Reason:} This is a repetition of a necessary hypothesis at a new point in the proof. The record expressly explains that the localized ring's one-prime-ideal sentence makes the earlier principal-ideal conclusion applicable. In Paper 31, §8, localization preserves the required multiplication-ring properties, and the zero and unit ideals remain excluded. It does not follow that every local ring is a principal ideal ring. Archival $o_p$ denotes the ring, not the numeral 0 with a subscript; the available text uses fraktur letters for the ring and prime ideal. The typography of the identifier and the hypothesis's mathematical meaning must be preserved separately. Repetition here helps readers follow the application of a theorem.\par
\noindent{\footnotesize Izvorne zapisy / Original records: HIST-g160-002.}\par
\smallskip\noindent\emph{Sravnjenje s tekstom:} Rabota 31, §8, najprvo ograničavaje klasu prstenov, potom davaje zaključok pri jedinom prostom idealu različnom od nulovogo i jediničnogo. Pri lokalizaciji povtorjajut se i ta jedinstvenost i svojstva multiplikacionnogo prstena. To objašnjaje začto povtorjenje uslovja jest matematično korisno, a ne izbytočno.\par
\noindent\emph{Textual comparison:} Paper 31, §8, first specifies the ring class, then draws the consequence when there is only one prime ideal other than zero and the unit ideal. The localization passage repeats both uniqueness and the multiplication-ring properties. This explains why repeating the hypothesis is mathematically useful rather than redundant.\par
\noindent{\footnotesize TEXT-COMP-363; pasaže / passages: 18787, 18789, 18823, 18932, 18970; \texttt{CORPUS\_FORM\_COMPARISONS.json}.}\par
\end{minipage}\par
\par\medskip\Needspace{7\baselineskip}
\noindent\begin{minipage}{\linewidth}
\hypertarget{TERM-1277}{}
\noindent\textbf{Bez povtorjenja i točno jedin raz}\par\noindent\textbf{Without repetition and exactly once}\quad{\footnotesize\texttt{TERM-1277}}\par
\smallskip\noindent\emph{Razlog:} Bez povtorjenja i jedin raz i samo jedin raz čuvajut različne, povězane česti tvrdženja: nepovtorjanje i cělo objimanje. Arhiv izrično odvrgaje njih razgledanje kako pustogo stila. V rabotě 9, §8, osnovny seznam ima po jednom polju iz vsakoj izomorfnej klasy. Potom konstrukcija pro jednu klasu davaje vse različne polja v toj klasě, každo samo jedin raz. Jedin prědstavitelj klasy zato ne znači isto kako jedin pojav vsakogo polja v klasě. Izraz bez povtorjenja sam ne garantuje, že žadno polje ne jest izpuščeno; jedin raz i samo jedin raz vyražaje i postojanje i jedinstvenost pojavu. Ta razlika potrěbuje čitanje cěloj konstrukcije, ne samo zaměnu sinonimov.\par
\noindent\emph{Reason:} Bez povtorjenja and jedin raz i samo jedin raz preserve related but distinct parts of the assertion: non-repetition and exhaustive coverage. The record expressly rejects treating them as mere style. In Paper 9, §8, the initial list has one field from each isomorphism class. The subsequent construction for a class gives all distinct fields in that class, each once. One representative of a class therefore does not mean the same thing as one occurrence of every field in it. Without repetition alone does not guarantee that no field is omitted; once and only once expresses both existence and uniqueness of occurrence. The distinction requires reading the whole construction, not merely substituting synonyms.\par
\noindent{\footnotesize Izvorne zapisy / Original records: HIST-g066-011.}\par
\smallskip\noindent\emph{Sravnjenje s tekstom:} Rabota 9, §8, počinaje iz prědstaviteljev izomorfnyh klasov i dalje izčisljaje različne polja v každoj klasě. Nepovtorjanje i cělost izčislenja trěba prověriti na pravom urovnju. Zato objašnjenje ne sklada vse pojavy slova prědstavitelj v jeden nerazličen smysl.\par
\noindent\emph{Textual comparison:} Paper 9, §8, starts from representatives of isomorphism classes and proceeds to enumerate distinct fields within each class. Non-repetition and completeness must be checked at the correct level. The explanation therefore does not collapse every representative-related formulation into one undifferentiated sense.\par
\noindent{\footnotesize TEXT-COMP-367; pasaže / passages: 8121, 8122, 8123, 8124, 8125, 8126, 8127, 8128, 8129, 8130, 8131, 8132, 8133, 8135, 8136, 8137, 8138, 8139, 8140, 8141, 8142, 8143, 8144, 8145, 8146, 8147, 8148, 8150, 8151; \texttt{CORPUS\_FORM\_COMPARISONS.json}.}\par
\end{minipage}\par
\par\medskip\Needspace{7\baselineskip}
\noindent\begin{minipage}{\linewidth}
\hypertarget{TERM-1279}{}
\noindent\textbf{Konec povtorjanoj operacije}\par\noindent\textbf{Termination of a repeated operation}\quad{\footnotesize\texttt{TERM-1279}}\par
\smallskip\noindent\emph{Razlog:} Operacija prěryvaje se opisuje konec povtorjanogo koraka dokaza poslě oprěděljenogo čisla priměn. Arhivny razlog daje prědnost procedurnoj formulaciji, v paraleli s ruskoju i ukrajinskoju frazoju. V rabotě 4, §4, granično čislo povtorjenij zavisi od razgledanogo slučaja i od indeksov: v jednom městě $\rho$, v drugom $\sigma$. To ne jest slučajny tehničny neuspěh ni prěryvanje od vněšnjego dějatelja. Takože ne govori se o približnom sbliženju nekonečnogo procesa. Pri prěnosu takoj frazy v drugu matematičnu oblast trěba oddělno ustanoviti, ktory korak ponavlja se, ktora granica čisla korakov jest dokazana i pod ktorymi uslovjami. Sam glagol ne dokazuje končanje.\par
\noindent\emph{Reason:} Operacija prěryvaje se describes the end of a repeated proof operation after a specified number of applications. The recorded reason prefers procedural wording, parallel to the Russian and Ukrainian phrases. In Paper 4, §4, the repetition limit depends on the case and indices: $\rho$ at one point, $\sigma$ at another. This is not an accidental technical failure or an interruption by an external actor. Nor is it approximate convergence of an infinite process. Transferring such wording to another mathematical subject requires identifying the repeated step, the proved bound on repetitions and its hypotheses. The verb alone does not establish termination.\par
\noindent{\footnotesize Izvorne zapisy / Original records: HIST-g033-004.}\par
\smallskip\noindent\emph{Sravnjenje s tekstom:} Rabota 4, §4, povězuje konec operacije s indeksami i granicami sum. Odvisno od slučaja, čislo povtorjenij jest ograničeno s $\rho$ ili $\sigma$. Slovny prěvod mora čuvati čislo i slučaj; samo prěryvaje se ne govori, ktora granica jest pravdiva.\par
\noindent\emph{Textual comparison:} Paper 4, §4, links termination to indices and summation limits. Depending on the case, repetitions are bounded by $\rho$ or $\sigma$. The verbal translation must preserve both number and case; terminates alone does not say which bound is valid.\par
\noindent{\footnotesize TEXT-COMP-369; pasaže / passages: 4109, 4122, 4163; \texttt{CORPUS\_FORM\_COMPARISONS.json}.}\par
\end{minipage}\par
\par\medskip\Needspace{7\baselineskip}
\noindent\begin{minipage}{\linewidth}
\hypertarget{TERM-1298}{}
\noindent\textbf{Oznaka razširjenogo ideala}\par\noindent\textbf{Notation for an extended ideal}\quad{\footnotesize\texttt{TERM-1298}}\par
\smallskip\noindent\emph{Razlog:} Formula $P=Tp$ v zapisu čuva identitet razširjenogo ideala i njegov okolny prsten. Arhivny razlog jest doslědnost s rannějšimi odlomkami o razširenju idealov. V dostupnom tekstu raboty 31, §8, oznaki za idealy i prsten pišut se frakturno: dolno p imenuje prosty ideal osnovnogo prstena, ne prosty element. Konečne sumy produktov s elementami nadprstena tvoret razširjeny ideal; ne potrěbno, že vsaky njegov element jest jeden takovy produkt. Ne trěba prědpolagati, že razširjeny ideal ostavaje prosty. Isti odlomok ima nerazrěšenu oznaku osnovy pri slovu poredok, kde piše idealno p, hoti bližši tekst govori o lokalizovanom osnovnom prstenu. To potrěbuje izvorno sravnenje, ne molčalivu zaměnu.\par
\noindent\emph{Reason:} The recorded formula $P=Tp$ preserves the extended ideal's identity and ambient ring. Its stated reason is continuity with the earlier ideal-extension passages. In Paper 31, §8, the current text uses fraktur for the ideals and ring: lower-case p denotes a prime ideal of the base ring, not a prime element. The extended ideal consists of finite sums of products with elements of the overring; not every element need be a single such product. The extended ideal must not be assumed prime. The same passage has an unresolved base label beside order, writing ideal p although the surrounding text concerns the localized base ring. This needs source comparison, not silent substitution.\par
\noindent{\footnotesize Izvorne zapisy / Original records: HIST-g161-002.}\par
\smallskip\noindent\emph{Sravnjenje s tekstom:} Rabota 31, §8, izrično naznačaje razširjeny ideal v nadprstenu, a potom koristi jeho primarne komponenty. Formula ne govori, že on jest prosty. V istom odlomku oznaka pri konečnom poredku jest zapisana idealnym p, ne navedenym lokalizovanym osnovnym prstenom; to ostavaje otvorenym voprosom iztočnogo teksta.\par
\noindent\emph{Textual comparison:} Paper 31, §8, explicitly identifies the extended ideal in the overring and then uses its primary components. The formula does not assert that it is prime. In the same passage the finite-order label uses ideal p rather than the stated localized base ring; this remains an open source-text question.\par
\noindent{\footnotesize TEXT-COMP-385; pasaže / passages: 18932, 18970, 19003, 19071; \texttt{CORPUS\_FORM\_COMPARISONS.json}.}\par
\end{minipage}\par
\par\medskip\Needspace{7\baselineskip}
\noindent\begin{minipage}{\linewidth}
\hypertarget{TERM-1299}{}
\noindent\textbf{Vhodi v kako dělitelnost}\par\noindent\textbf{Vhodi v as divisibility}\quad{\footnotesize\texttt{TERM-1299}}\par
\smallskip\noindent\emph{Razlog:} Vhodi v v toj rečenici vyražaje dělitelnost, ne pripadnost elementa množstvu. Arhivny razlog čuva tu glagolnu formulaciju, čoby dokaz ostal paralelny s teoremnym tvrdženjem. V rabotě 31, §8, ono povězuje prosty ideal v osnovnom prstenu s prostym elementom $p$ v lokalizovanom prstenu; $p$ děli tamošnji diskriminantny element $D_{T_P}$. Ne možno zaměniti ideal i njegov lokalny generator ili čitati vhodi kako prosto figuriranje symbola v formulě. Arhivne prostye bukvy v oznakah ne kažut vsih tipografičnyh razlik nynějšego teksta. Smysl čuva se točnym určenjem objekta, osnovnogo prstena i napravjenja dělitelnosti, ne samym povtorjenjem glagola.\par
\noindent\emph{Reason:} Vhodi v in this sentence expresses divisibility, not set membership. The recorded reason preserves that verbal formulation to keep the proof parallel to the theorem statement. In Paper 31, §8, it connects a prime ideal of the base ring with a prime element $p$ in the localized ring: $p$ divides the local discriminant element $D_{T_P}$. The ideal and its local generator must not be interchanged, nor vhodi read as a symbol merely appearing in a formula. The archival plain letters do not preserve every typographic distinction of the current text. Meaning is retained by specifying the object, base ring and direction of divisibility, not by repeating the verb alone.\par
\noindent{\footnotesize Izvorne zapisy / Original records: HIST-g164-003.}\par
\smallskip\noindent\emph{Sravnjenje s tekstom:} Dokaz raboty 31, §8, prěhodi od dělitelnosti diskriminantnogo ideala prostym idealom do dělitelnosti lokalnogo diskriminantnogo elementa generatorom $p$. Glagol vhodi čuva tu svezu, ako čitatelj različi ideal, element i prsten. Smer dělitelnosti ne smije byti obratny.\par
\noindent\emph{Textual comparison:} The proof in Paper 31, §8, passes from divisibility of the discriminant ideal by a prime ideal to divisibility of a local discriminant element by the generator $p$. Vhodi preserves that connection only if the ideal, element and ring are distinguished. The direction of divisibility must not be reversed.\par
\noindent{\footnotesize TEXT-COMP-386; pasaže / passages: 18715, 19037, 19071, 19105, 19107; \texttt{CORPUS\_FORM\_COMPARISONS.json}.}\par
\end{minipage}\par
\par\medskip\Needspace{7\baselineskip}
\noindent\begin{minipage}{\linewidth}
\hypertarget{TERM-1300}{}
\noindent\textbf{Jedin element poraždaje ideal}\par\noindent\textbf{One element generates the ideal}\quad{\footnotesize\texttt{TERM-1300}}\par
\smallskip\noindent\emph{Razlog:} Poraždaje čuva izbrany glagol za algebraično generovanje. Zapis ne prědlagaje novu konkurentnu formu, ale prodolžaje uže upotrěbjany jezyk dokaza. V rabotě 31, §8, sam $p$ poraždaje prosty ideal, zato že jego glavny ideal ima točno prvu stepen jedinogo dopuščenogo prostogo idealnogo faktora. To znači ravnost idealov, ne samo pripadnost $p$ idealu. Rečenica zavisi od navedenogo svojstva multiplikacionnogo prstena i jedinstvenosti prostogo ideala različnogo od nulovogo i jediničnogo. Slovo sam podčrtaje, že ne potrěbny druge generatory; ne tvrdži, že generator jest jedinstveny. Ruska i ukrajinska paralela takože čuvajut glagol za generovanje.\par
\noindent\emph{Reason:} Poraždaje retains the selected verb for algebraic generation. The record proposes no new competing form; it continues the proof's existing terminology. In Paper 31, §8, $p$ itself generates the prime ideal because its principal ideal has exactly the first power of the only allowed prime-ideal factor. This asserts equality of ideals, not merely membership of $p$ in the ideal. The sentence depends on the stated multiplication-ring properties and uniqueness of the prime ideal other than zero and the unit ideal. Itself emphasizes that no additional generators are needed; it does not assert uniqueness of the generator. The Russian and Ukrainian parallels likewise retain the generation verb.\par
\noindent{\footnotesize Izvorne zapisy / Original records: HIST-g156-005.}\par
\smallskip\noindent\emph{Sravnjenje s tekstom:} Uslovja raboty 31, §8, dajut, že glavny ideal od $p$ jest ravny razgledanomu prostomu idealu. Poraždaje vyražaje tu ravnost. Pripadnost jednogo elementa idealu ne byla by dovoljna; ni slovo sam ne ustanavljaje jedinstvenost generatora.\par
\noindent\emph{Textual comparison:} The hypotheses of Paper 31, §8, make the principal ideal of $p$ equal to the prime ideal in question. Generates expresses that equality. Membership of one element would not suffice; itself does not establish uniqueness of the generator.\par
\noindent{\footnotesize TEXT-COMP-387; pasaže / passages: 18789, 18823; \texttt{CORPUS\_FORM\_COMPARISONS.json}.}\par
\end{minipage}\par
\par\medskip\Needspace{7\baselineskip}
\noindent\begin{minipage}{\linewidth}
\hypertarget{TERM-1310}{}
\noindent\textbf{Avtorovo popravjenje: polja s konjugatami}\par\noindent\textbf{Authorial correction: fields with conjugates}\quad{\footnotesize\texttt{TERM-1310}}\par
\smallskip\noindent\emph{Razlog:} To jest redakcijna odluka o priměnjenju avtorovogo popravjenja, ne novo nazvanje za polje. Arhiv govori o razširjenju izbranyh $(H,y)$ i $(H,Y)$ vsimi konjugatami do indeksa $\lambda_0-1$. Dostupna rabota 16 v prvom popravjenju navodi točny naslov ranějšej raboty, pečatne lokatory i zaměny; rabota 9, §7, už sodrživaje razširjene polja. Črty i gornje indekse tu označajut konjugate, ne derivacije ni obične stepeni. Zaměna měni osnovu razsudženja o medžupoljah i prěsěkah. Diplomatično izdanje ranějšej raboty i izdanje s priměnjenym popravjenjem morajut byti različene. Nynešnja sravniteljna prověrka čita prěvod popravjenja; ne predstavlja novu prověrku pečatnogo německogo iztočnika.\par
\noindent\emph{Reason:} This is an editorial decision to apply an authorial correction, not a new name for a field. The record says selected $(H,y)$ and $(H,Y)$ are enlarged by all conjugates through index $\lambda_0-1$. The available Paper 16’s first correction names the earlier work, printed locators and replacements; Paper 9, §7, already contains the enlarged fields. Primes and parenthesized superscripts here label conjugates, not derivatives or ordinary powers. The change affects the argument about intermediate fields and intersections. A diplomatic edition of the earlier work and an edition incorporating the correction must be distinguished. The present comparison reads the translated correction; it does not claim a new examination of the printed German witness.\par
\noindent{\footnotesize Izvorne zapisy / Original records: HIST-g065-018.}\par
\smallskip\noindent\emph{Sravnjenje s tekstom:} Prvo popravjenje raboty 16 naznačaje ranějši naslov, pečatne lokatory i proširjene polja. Te liste konjugatov sut už prisutne v rabotě 9, §7. To podpiraje redakcijnu istoriju zapisanej zaměny, ale ne novu diplomatičnu prověrku pečatnogo originala.\par
\noindent\emph{Textual comparison:} Paper 16’s first correction specifies the earlier title, printed locators and enlarged fields. Those conjugate lists are already present in Paper 9, §7. This supports the editorial history of the recorded replacement, but is not a new diplomatic comparison with the printed original.\par
\noindent{\footnotesize TEXT-COMP-396; pasaže / passages: 7913, 7914, 7915, 7916, 7917, 7918, 7975, 7976, 7977, 7978, 7979, 7980, 7981, 7982, 7983, 7984, 7985, 7986, 8010, 8011, 8012, 8013, 8014, 8015, 10893, 10896; \texttt{CORPUS\_FORM\_COMPARISONS.json}.}\par
\end{minipage}\par
\par\medskip\Needspace{7\baselineskip}
\noindent\begin{minipage}{\linewidth}
\hypertarget{TERM-1328}{}
\noindent\textbf{V polinomnom smyslu: ograda odnosa}\par\noindent\textbf{In the polynomial sense: qualifying a relation}\quad{\footnotesize\texttt{TERM-1328}}\par
\smallskip\noindent\emph{Razlog:} V polinomnom smyslu jest ograda smysla, ne novo nazvanje za matematičny objekt. Zapis ju dodržuje, da by čitatelj ne směšal děliteljnost polinomov s děliteljnostju idealov ili modulov. Ruska i ukrajinska paralela takože izrično označajut smysl. V rabotě 22 ta ograda nastupaje pri vzajemno prostyh faktorah i pri največšem občem dělitelju polinomov, dočim okolny tekst upotrěbljaje idealne kvocienty i modulne normy. Samo prisutstvo idealov v odlomku zato ne měni tip dělitelja. Prěvod mora čuvati objekt i odnos: polinomni obči dělitelj ne jest avtomatično ideal generovany vsimi danymi polinomami.\par
\noindent\emph{Reason:} V polinomnom smyslu is a qualification of sense, not a new mathematical object name. The record retains it to prevent polynomial divisibility from being confused with ideal or module divisibility. Russian and Ukrainian also mark the sense explicitly. In Paper 22 it accompanies relatively prime polynomial factors and a greatest common divisor of polynomials, while the surrounding argument uses ideal quotients and module norms. The presence of ideals in the passage does not change the divisor's type. The translation must preserve both object and relation: a polynomial gcd is not automatically the ideal generated by all the polynomials.\par
\noindent{\footnotesize Izvorne zapisy / Original records: HIST-g091-005, HIST-g092-074.}\par
\smallskip\noindent\emph{Sravnjenje s tekstom:} Rabota 22 odlišava največši obči dělitelj polinomov od idealnogo kvocienta i modulnoj normy. Ograda v polinomnom smyslu čuva tip odnosa, ne jest odvečny ukras prozy.\par
\noindent\emph{Textual comparison:} Paper 22 distinguishes a polynomial greatest common divisor from an ideal quotient and a module norm. In the polynomial sense preserves the relation's type; it is not dispensable prose decoration.\par
\noindent{\footnotesize TEXT-COMP-413; pasaže / passages: 14913, 14915, 14918, 14920, 14922, 14925, 14992, 14994, 14996, 14998; \texttt{CORPUS\_FORM\_COMPARISONS.json}.}\par
\end{minipage}\par
\par\medskip\Needspace{7\baselineskip}
\noindent\begin{minipage}{\linewidth}
\hypertarget{TERM-1383}{}
\noindent\textbf{Kogda vse primarne komponenty sut proste}\par\noindent\textbf{When all primary components are prime}\quad{\footnotesize\texttt{TERM-1383}}\par
\smallskip\noindent\emph{Razlog:} Primarne komponenty stavajut prostymi idealami jest čest definicije polnoj razložimosti, ne novy termin ni tvrdženje, že vsaka primarna komponenta vsegda jest prosta. Arhiv žela čuvati razliku primarny/prosty, potrěbnu za pozdnějši diskriminantny kriterij. Rabota 31, §3, veže uslovje k produktnomu prědstavjenju nulovogo ideala i k vsim njegovym komponentam. Ravnoznačno prědstavjenje jest direktna suma konečno mnogyh polj. Vse ne može byti zaměnjeno s někoje; glagol stavajut tu izražaje uslovje definicije, ne časovy proces. Tak samo se oddělja prvy rod, ktory trěbuje daljšego uslovja na proste idealy.\par
\noindent\emph{Reason:} Primarne komponenty stavajut prostymi idealami is part of the definition of complete reducibility, not a new term or an assertion that every primary component is always prime. The recorded reason is to preserve the primary/prime distinction needed for the later discriminant criterion. Paper 31, §3, attaches the condition to the product representation of the zero ideal and to all its components. The equivalent representation is a direct sum of finitely many fields. All cannot become some; become expresses a defining condition here, not a temporal process. First-kind complete reducibility adds a further condition on the primes.\par
\noindent{\footnotesize Izvorne zapisy / Original records: HIST-g125-003.}\par
\smallskip\noindent\emph{Sravnjenje s tekstom:} Produkt idealov i direktna suma prstena imajut razne objekty. Polna razložimost trěbuje prostotu vsih primarnyh komponentov.\par
\noindent\emph{Textual comparison:} The product of ideals and the direct sum of the ring concern different objects. Complete reducibility requires every primary component to be prime.\par
\noindent{\footnotesize TEXT-COMP-461; pasaže / passages: 17434, 17436, 17438, 17440, 17469, 17470, 17471, 17472, 17473, 17474, 17475, 17476, 17477, 17478, 17479, 17480, 17630, 17632; \texttt{CORPUS\_FORM\_COMPARISONS.json}.}\par
\end{minipage}\par
\par\medskip\Needspace{7\baselineskip}
\noindent\begin{minipage}{\linewidth}
\hypertarget{TERM-1388}{}
\noindent\textbf{Granica prěnosa: nedostajuča operacija}\par\noindent\textbf{A limit on transfer: a missing operation}\quad{\footnotesize\texttt{TERM-1388}}\par
\smallskip\noindent\emph{Razlog:} Pojem produkta ne prěnosi se na moduly čuva ograničenje metoda, ne izkazuje, že ne može obstajati nikaky produkt nikakih modulov. Arhiv ostavlja produkt izrično, zato že v razgleđanoj oblasti $T$ ne jest definirano množenje dvoh jej elementov. Zato danu definiciju primarnogo i prostogo ideala ne možno neposrědnje prěnesti. Rabota 19, §9, od toga razlikuje teoremy osnovane samo na děliteljnosti i najmenšem občem mnogokratnom, ktore se prěnoset. Posebno razgleđaje nekomutativne prsteny i neuspěh potrěbnyh dokazov. Izraz ne směje byti skračen do bezuslovnogo zakaza, ktory by obhvačal i drugače definirane konstrukcije.\par
\noindent\emph{Reason:} Pojem produkta ne prěnosi se na moduly preserves a limitation of the method, not a claim that no kind of product of modules can exist. The archive keeps product explicit because multiplication of two elements of the domain $T$ under discussion is undefined. The given definitions of primary and prime ideals therefore cannot transfer directly. Paper 19, §9, contrasts this with theorems based only on divisibility and ideal least common multiples, which do transfer. It separately discusses noncommutative rings and failure of the required proofs. The phrase must not become an unconditional prohibition covering differently defined constructions.\par
\noindent{\footnotesize Izvorne zapisy / Original records: HIST-g081-074.}\par
\smallskip\noindent\emph{Sravnjenje s tekstom:} Teoremy o děliteljnosti mogut ostati, hoti rezultat potrěbujuči nedefinirany produkt ne može prějti. Negacija jest lokalna k danoj strukturi.\par
\noindent\emph{Textual comparison:} Divisibility theorems can survive when a result requiring an undefined product cannot transfer. The negative statement is local to the given structure.\par
\noindent{\footnotesize TEXT-COMP-466; pasaže / passages: 13360, 13362, 13364, 13366, 13403, 13405, 13407, 13409; \texttt{CORPUS\_FORM\_COMPARISONS.json}.}\par
\end{minipage}\par
\par\medskip\Needspace{7\baselineskip}
\noindent\begin{minipage}{\linewidth}
\hypertarget{TERM-1403}{}
\noindent\textbf{Imena i bibliografične nazvanja v rabotě 12}\par\noindent\textbf{Names and bibliographic titles in Paper 12}\quad{\footnotesize\texttt{TERM-1403}}\par
\smallskip\noindent\emph{Razlog:} Pravilo za rabotu 12 razlikuje imena v běžnom izloženju od točnyh bibliografičnyh nazvanij. Arhiv izbral slovjanske formy Kristofel, Riči, Riman, Lipšic i Lagranž za izloženje, ktore može byti prěvedeno v kirilicu; nazvanja Crelle, Math. Annalen, Werke i Mécanique analytique ostavajut latinskymi bibliografičnymi ostrovami. Razlog jest čitlivost teksta zajedno s možnostju najdti citovanu publikaciju. Nynějšny latinsky tekst pokazuje tuto razliku pri imenah i citacijah v uvodu. To jest pravilo pisanja i identifikacije, ne novo matematično slovo; anglijsky opis v starom spisku ne jest medžuslovjansky termin. Tuto sravnjenje samo po sobě ne proverjaje vsakogo kiriličnogo výstupa.\par
\noindent\emph{Reason:} The Paper 12 convention distinguishes names in running exposition from exact bibliographic titles. The record selects Slavicized forms such as Kristofel, Riči, Riman, Lipšic and Lagranž for prose that can be converted to Cyrillic, while Crelle, Math. Annalen, Werke and Mécanique analytique remain Latin-script bibliographic islands. Its rationale combines readable prose with identifiable references. The current Latin text shows that distinction in the introduction's names and citations. This is a writing and identification convention, not a mathematical coinage; the older English description is not an Interslavic term. This comparison alone does not verify every Cyrillic output.\par
\noindent{\footnotesize Izvorne zapisy / Original records: HIST-g074-020.}\par
\smallskip\noindent\emph{Sravnjenje s tekstom:} Prisposobljene imena ostavajut različne od izvorno napisanyh bibliografičnyh nazvanij.\par
\noindent\emph{Textual comparison:} Adapted names remain distinct from bibliographic titles in their original spelling.\par
\noindent{\footnotesize TEXT-COMP-479; pasaže / passages: 8976, 8978, 9146; \texttt{CORPUS\_FORM\_COMPARISONS.json}.}\par
\end{minipage}\par
\par\medskip\Needspace{7\baselineskip}
\noindent\begin{minipage}{\linewidth}
\hypertarget{TERM-1404}{}
\noindent\textbf{Latinske citacije v kiriličnom izloženju}\par\noindent\textbf{Latin citations in Cyrillic exposition}\quad{\footnotesize\texttt{TERM-1404}}\par
\smallskip\noindent\emph{Razlog:} V rabotě 11 zaščita citacij čuva izvorno pisanje inostrannyh nazvanij, dokolě medžuslovjansko izloženje i prisposobljene imena podležat prěvodu v kirilicu. Arhiv povezuje tuto razliku s kurzivoju i zaštitoju bibliografičnyh fragmentov; cilj byl izběgnuti slučajnu směs latinskyh i kiriličnyh bukv v jednom slově. Nynějšny tekst, napr., razlikuje běžno imeno Lurot od německogo nazvanja njegove raboty, a Weber Algebra ohranjaje kako citaciju. Kurziva sama ne dokazuje, že tekst jest inostranny: v matematičnom izloženju ona može iměti druge funkcije. Zato pravilo trěbuje smyslovo označene bibliografične česti, ne slěpu zaščitu vsego kurzivnogo teksta. Kirilična versija ostavaje projekcijeju latinskoj, ne nezavisnym prěvodnym svědkom.\par
\noindent\emph{Reason:} In Paper 11, citation protection preserves the original spelling of foreign titles while Interslavic exposition and adapted names are eligible for Cyrillic conversion. The record associates this distinction with italics and protected bibliographic spans, aiming to avoid accidental Latin–Cyrillic mixtures within a word. The current text, for example, distinguishes the running name Lurot from the German title of his work and retains Weber Algebra as a citation. Italics alone do not establish foreign-language status: they have other mathematical uses. The rule therefore requires semantically identified bibliographic spans, not blanket protection of italic text. The Cyrillic edition remains a projection of the Latin edition, not an independent translation witness.\par
\noindent{\footnotesize Izvorne zapisy / Original records: HIST-g073-022.}\par
\smallskip\noindent\emph{Sravnjenje s tekstom:} Latinsky tekst razlikuje prisposobljene imena od citacij. Sravnjenje ne jest proverka vsego kiriličnogo izdanja.\par
\noindent\emph{Textual comparison:} The Latin text distinguishes adapted names from citations. This comparison is not verification of the entire Cyrillic reader.\par
\noindent{\footnotesize TEXT-COMP-480; pasaže / passages: 8892, 8900, 8904; \texttt{CORPUS\_FORM\_COMPARISONS.json}.}\par
\end{minipage}\par
\par\medskip\Needspace{7\baselineskip}
\noindent\begin{minipage}{\linewidth}
\hypertarget{TERM-1414}{}
\noindent\textbf{Upozorenje o neopreděljenom kvocientu}\par\noindent\textbf{The warning about an undefined quotient}\quad{\footnotesize\texttt{TERM-1414}}\par
\smallskip\noindent\emph{Razlog:} Tuto izložiteljsko rěšenje čuva upozorenje, že racionalny kvocient po jednočasnoj substituciji dvoh prěměnnyh može stati neopreděljenym. Rabota 6, §2, pokazyvaje priměr, kde čislitelj i imenitelj oba izčezajut pri $x_3=x_4=0$. Zapis $0/0$ označaje nedopustimu podělu, ne novu čislovu cěnnost funkcije. Razlog arhiva jest čuvanje ograničenja dokaza: zaključanje za jednu sukcesivnu substituciju ne može bez dodatnogo opravdanja važiti za dvě zajedno. Takova rečenica ne jest ukrasna primětka, ktorej možno dati kračši nejasny opis. Zapis neopreděljenym jest padežna forma iz konkretnoj rečenice, ne osnovna slovnična forma pridavnogo imena.\par
\noindent\emph{Reason:} This expository decision preserves the warning that a rational quotient may become undefined after simultaneous substitution of two variables. Paper 6, §2, gives an example whose numerator and denominator both vanish at $x_3=x_4=0$. The notation $0/0$ marks an inadmissible division, not a newly assigned numerical value of the function. The recorded reason is to retain a limitation of the proof: an argument for a single successive substitution does not automatically apply to two substitutions together. The sentence is substantive, not dispensable commentary. Neopreděljenym is an inflected form from the sentence, not the adjective's dictionary citation form.\par
\noindent{\footnotesize Izvorne zapisy / Original records: HIST-g042-015.}\par
\smallskip\noindent\emph{Sravnjenje s tekstom:} Čislitelj i imenitelj mogut izčeznuti zajedno. Zapis $0/0$ ne davaje funkciji cěnnost.\par
\noindent\emph{Textual comparison:} Numerator and denominator can vanish together. The notation $0/0$ assigns no value to the function.\par
\noindent{\footnotesize TEXT-COMP-488; pasaže / passages: 5212, 5213, 5214, 5215, 5216, 5217; \texttt{CORPUS\_FORM\_COMPARISONS.json}.}\par
\end{minipage}\par
\par\medskip\Needspace{7\baselineskip}
\noindent\begin{minipage}{\linewidth}
\hypertarget{TERM-1427}{}
\noindent\textbf{Komponentne odobraženja i vyraženje invariantov}\par\noindent\textbf{Component maps and the expression of invariants}\quad{\footnotesize\texttt{TERM-1427}}\par
\smallskip\noindent\emph{Razlog:} Konstrukcija $R$ jest homomorfno k $K_i$ čuva konkretne odobraženja v komponenty, potrěbne za vyraženje slěda, normy i diskriminanta. Arhiv bira internacionalno homomorfno zajedno s izričnym izraziti črez elementy; nynějšny tekst piše prěz. V rabotě 31, §6, komponenty imajut vid $K_i e_i$, kde $K_i$ leži v obhvatnom polju $Q$. Homomorfizm zato ne jest samo neoprěděljena svez dvoh objektov. Ne možno bez daljšego uslovja zaměniti homomorfno slovom izomorfno: tekst čini tuty prěhod samo v osobnom slučaju, kogda izvorno prsten samo jest polje. Tuto jest izložiteljsko rěšenje o smyslu odobraženja, ne novy slovničny korenj.\par
\noindent\emph{Reason:} The construction $R$ jest homomorfno k $K_i$ retains the concrete component maps needed to express trace, norm and discriminant. The archive selects international homomorfno with the explicit phrase izraziti črez elementy; the current text uses prěz. In Paper 31, §6, the components have the form $K_i e_i$, with $K_i$ contained in the overfield $Q$. A homomorphism is therefore not merely an unspecified association of objects. Homomorphic cannot be changed to isomorphic without further conditions: the text makes that transition in the special case where the original ring is itself a field. This is an expository decision about the map's meaning, not a newly coined stem.\par
\noindent{\footnotesize Izvorne zapisy / Original records: HIST-g146-007.}\par
\smallskip\noindent\emph{Sravnjenje s tekstom:} Komponentne odobraženja dajut izraz slěda, normy i diskriminanta. Izomorfizm potrěbuje osobny slučaj polja.\par
\noindent\emph{Textual comparison:} Component maps give expressions for trace, norm and discriminant. Isomorphism uses the special field case.\par
\noindent{\footnotesize TEXT-COMP-499; pasaže / passages: 18470, 18471, 18472, 18473, 18474, 18525; \texttt{CORPUS\_FORM\_COMPARISONS.json}.}\par
\end{minipage}\par
\par\medskip\Needspace{7\baselineskip}
\noindent\begin{minipage}{\linewidth}
\hypertarget{TERM-1428}{}
\noindent\textbf{Izomorfija ne znači vključanje}\par\noindent\textbf{Isomorphism does not assert inclusion}\quad{\footnotesize\texttt{TERM-1428}}\par
\smallskip\noindent\emph{Razlog:} Ograničenje samo kako ekvivalentno polje, a ne kako polje už ležeče v $Q$, čuva razliku medžu izomorfijeju i vključanjem. To jest cěly matematičny razlog zachovanogo izbora. Rabota 31, §6, davaje izomorfne obrazy $K_i$ v jednom Galoisovom polju $Q$, ale izvorno polje ne prědpokladaje už tamo směstjenym. Slova samo i a ne ne sut stilističny dodatok: ona ograđajut tvrđenje. V tutom kontekstu ekvivalentno objasnjaje se pomoču navedenyh izomorfizmov; ne trěba go čitati kako libovoljnu neoprěděljenu ekvivalenciju. Pri prěnosu v logiku razlikovanje abstraktnogo objekta i njegove realizacije mora ostati jednako eksplicitno.\par
\noindent\emph{Reason:} The qualification 'only as an equivalent field, not as a field already lying in $Q$' preserves the difference between isomorphism and inclusion. That is the full mathematical reason for the recorded choice. Paper 31, §6, supplies isomorphic images $K_i$ in one Galois overfield $Q$, but does not assume the original field is already situated there. Only and not are not stylistic additions: they delimit the claim. Equivalent is explained here by the stated isomorphisms, not an unspecified equivalence relation. Transfer to logic must preserve the distinction between an abstract object and its realization just as explicitly.\par
\noindent{\footnotesize Izvorne zapisy / Original records: HIST-g147-007.}\par
\smallskip\noindent\emph{Sravnjenje s tekstom:} Tekst ne prědpokladaje, že izvorno polje samo už leži v $Q$.\par
\noindent\emph{Textual comparison:} The text does not assume the original field is already contained in $Q$.\par
\noindent{\footnotesize TEXT-COMP-500; pasaže / passages: 18525; \texttt{CORPUS\_FORM\_COMPARISONS.json}.}\par
\end{minipage}\par
\par\medskip\Needspace{7\baselineskip}
\noindent\begin{minipage}{\linewidth}
\hypertarget{TERM-1449}{}
\noindent\textbf{Jednako nazvanje pri upučenju v slědujuči paragraf}\par\noindent\textbf{Consistent naming across a forward reference}\quad{\footnotesize\texttt{TERM-1449}}\par
\smallskip\noindent\emph{Razlog:} Zapis čuva svez konca §5 s naslovom i dokazom §6, ne novu različnu matematičnu rěč. §5 obeščaje prostějšu konstrukciju pomoču racionalnoj bazy; §6 pod tym samym nazvanjem davaje definiciju i konstrukciju. Pri rěšenju užito jest kontinuitet upučenja: zaměna sinonimom samo v jednom iz tih měst mogla by stvoriti vtis novego objekta. Arhivny znak prěd čislom 6 ne jest zachovany dostatočno za ego samostojnu rekonstrukciju, ale směr upučenja jest čitliv iz sosednyh odstavkov i naslova. To jest izložiteljska konvencija, ktora potrěbuje prověrjenje oboju koncev upučenja, ne novu slovničnu jedinicu.\par
\noindent\emph{Reason:} The record preserves the connection between the end of §5 and the heading and argument of §6, rather than introducing a distinct mathematical word. Section 5 promises a simpler construction using a rational basis; §6 supplies its definition and construction under that same name. Referential continuity governs the choice: replacing the term with a synonym at only one endpoint could suggest a new object. The archival character before the numeral 6 is not sufficiently preserved to reconstruct on its own, but the destination is readable from the adjoining paragraphs and heading. This is an expository convention requiring both ends of the reference to be checked, not a new lexical item.\par
\noindent{\footnotesize Izvorne zapisy / Original records: HIST-g063-012.}\par
\smallskip\noindent\emph{Sravnjenje s tekstom:} Upučenje i naslov upotrěbljajut isto nazvanje za objavljenu konstrukciju.\par
\noindent\emph{Textual comparison:} The forward reference and heading use the same name for the promised construction.\par
\noindent{\footnotesize TEXT-COMP-518; pasaže / passages: 7694, 7695, 7696, 7697, 7716, 7718, 7719; \texttt{CORPUS\_FORM\_COMPARISONS.json}.}\par
\end{minipage}\par
\par\medskip\Needspace{7\baselineskip}
\noindent\begin{minipage}{\linewidth}
\hypertarget{TERM-1471}{}
\noindent\textbf{Dva objekty v jednom naslovu}\par\noindent\textbf{Two objects in one heading}\quad{\footnotesize\texttt{TERM-1471}}\par
\smallskip\noindent\emph{Razlog:} Naslov čuva dvě različne česti zadanja: redukciju modula $(ss'u)^4$ i opis sistema od $s$. §17 raboty 2 izrično razděljaje dokaz na A) reducibilnost modula i B) ostavanje jedinoj formy $s$ v sistemu pri vyšem modulu. V arhivu od $s$ označeno jest kako vrěmenna prozirna konstrukcija, čekajuča jezyčno ocěnjenje; ne trěba iz povtorjenja vyvoditi, že to ocěnjenje nastupilo. Četvrta stepenj, apostrof i razlikovanje $s$ od njegovogo sistema morajut ostati točne. To jest rěšenje o strukturi naslova i upravljanju, ne samostojno novo značenje slova redukcija.\par
\noindent\emph{Reason:} The heading preserves two parts of the task: reduction of the modulus $(ss'u)^4$ and description of the system of $s$. Paper 2, §17, explicitly divides the argument into A) reducibility of the modulus and B) the system's reduction to the single form $s$ under the higher modulus. The archive labels od $s$ a provisional transparent construction pending linguistic review; recurrence does not show that such review occurred. The fourth power, prime mark and distinction between $s$ and its system must remain exact. This is a decision about heading structure and grammatical government, not a separate new sense of redukcija.\par
\noindent{\footnotesize Izvorne zapisy / Original records: HIST-g018-002.}\par
\smallskip\noindent\emph{Sravnjenje s tekstom:} Redukcija modula i opis sistema ostavajut različne česti zadanja.\par
\noindent\emph{Textual comparison:} Reduction of the modulus and description of the system remain distinct parts of the task.\par
\noindent{\footnotesize TEXT-COMP-539; pasaže / passages: 2384, 2387, 2389, 2391, 2394, 2395, 2399; \texttt{CORPUS\_FORM\_COMPARISONS.json}.}\par
\end{minipage}\par
\par\medskip\Needspace{7\baselineskip}
\noindent\begin{minipage}{\linewidth}
\hypertarget{TERM-1475}{}
\noindent\textbf{Směr upučenja i čislo operacij}\par\noindent\textbf{Reference direction and number of operations}\quad{\footnotesize\texttt{TERM-1475}}\par
\smallskip\noindent\emph{Razlog:} Věta o neposrědnom slědovanju redukcij čuva zavisimost od §21 i točno jednu kontrakciju. Arhiv pozvoljaje novy poredok slov, ale ne odstranjenje tih dvoh podatkov. Staru konstrukciju črez jednokratno svijanje nynějša pasaža §22 prědaje kak prěz jednokratnu kontrakciju; ide o upučenje na konkretnu operaciju, ne o volnu zaměnu vsih slov s podobnym běžnym značenjem. Neposrědno ili prosto opisuje kratkost vyvoda iz nazvanyh rezultatov; ne znači, že rezultat ne potrěbuje osnovy. Tuto jest izložiteljsko rěšenje o zavisimosti dokaza, s čislom operacij i iztočnikom vyvoda kako nepoměnnymi čestjami smysla.\par
\noindent\emph{Reason:} The sentence about reductions following directly preserves their dependence on §21 and the use of exactly one contraction. The archive allows a modernized word order, but not removal of those two facts. The older črez jednokratno svijanje appears in the current §22 passage as prěz jednokratnu kontrakciju: this identifies a particular operation, not a licence to interchange all words with similar everyday meanings. Directly or simply describes the short derivation from named results; it does not mean that the result needs no basis. This is an expository decision about proof dependence, with the operation count and cited source of the inference preserved.\par
\noindent{\footnotesize Izvorne zapisy / Original records: HIST-g023-003.}\par
\smallskip\noindent\emph{Sravnjenje s tekstom:} Upučenje na §21 i jedna kontrakcija nosut smysl vyvoda.\par
\noindent\emph{Textual comparison:} The reference to §21 and a single contraction carry the inference's meaning.\par
\noindent{\footnotesize TEXT-COMP-543; pasaže / passages: 2950, 2952, 2954, 2957; \texttt{CORPUS\_FORM\_COMPARISONS.json}.}\par
\end{minipage}\par
\par\medskip\Needspace{7\baselineskip}
\noindent\begin{minipage}{\linewidth}
\hypertarget{TERM-1515}{}
\noindent\textbf{Konstrukcija i historičny metod}\par\noindent\textbf{Construction and historical method}\quad{\footnotesize\texttt{TERM-1515}}\par
\smallskip\noindent\emph{Razlog:} Rekursivna konstrukcija specialnyh razvitij v rędy jest metodna fraza, ne nazvanje novogo matematičnogo objekta. Uvod raboty 2 razlikuje točnu formulaciju redukcijnyh teoremov od računnogo sredstva za jih priměnjenje. Arhiv trěbuje čuvati postepenny, rekursivny postupok bez anahronističnogo prěimenovanja v sovrěmenny algoritm. To ne zabranjaje izjasniti način računa: §2 pokazuje, kako koeficienty slědujut iz sistema uravnenij. No pojem rekursivno sam ne dodava čislenu složenost, formalny računny model ili ocenjenje potrebnyh resursov. Fraza jest zato predstavjena medžu izložiteljskymi odlukami, s terminologičnymi svezami k razloženju i polaram, a ne kako odrěčena ili pozabylena slovnička stavka.\par
\noindent\emph{Reason:} Rekursivna konstrukcija specialnyh razvitij v rędy is a methodological phrase, not the name of a new mathematical object. Paper 2's introduction distinguishes precise formulation of reduction theorems from the computational means of applying them. The archive requires retaining a stepwise recurrent procedure without anachronistically renaming it as a modern algorithm. This does not prohibit explaining the calculation: §2 shows how coefficients follow from a system of equations. But recursive alone adds no complexity bound, formal model of computation, or resource estimate. The phrase is therefore presented among expository decisions, with links to expansion and polar terminology, rather than omitted from the dictionary.\par
\noindent{\footnotesize Izvorne zapisy / Original records: HIST-g002-008.}\par
\smallskip\noindent\emph{Sravnjenje s tekstom:} Uvodna metodna fraza imaje konkretno prodolženje v sistemu za koeficienty. Jeje dokumentovanje ne trěbuje pripisati historičnomu tekstu sovrěmennu teoriju algoritmov.\par
\noindent\emph{Textual comparison:} The introductory methodological phrase has a concrete continuation in the coefficient system. Documenting it does not require attributing a modern theory of algorithms to the historical text.\par
\noindent{\footnotesize TEXT-COMP-577; pasaže / passages: 180, 437, 439, 441, 515, 518, 519, 520, 521, 523; \texttt{CORPUS\_FORM\_COMPARISONS.json}.}\par
\end{minipage}\par
\par\medskip\Needspace{7\baselineskip}
\noindent\begin{minipage}{\linewidth}
\hypertarget{TERM-1610}{}
\noindent\textbf{Ravnost idealov iz vseh lokalizacij}\par\noindent\textbf{Equality of ideals from all localizations}\quad{\footnotesize\texttt{TERM-1610}}\par
\smallskip\noindent\emph{Razlog:} To jest cělo tvrdženje o ravnosti idealov, ne novo leksično zaglavje. Arhiv govori o dopolnjenju po pečatnom svědku; toj opis ostavaje historijna informacija, ne nova vizualna prověrka skana. Nynějšny konec §8, točka 3 raboty 31 izrično trěbuje sběganje razširjenyh idealov v vsih tam opreděljenyh lokalizacijah. Dokaz se opira na iste proste idealy i iste eksponenty. Vsi ne možno oslabiti na jednu izbranu lokalizaciju; identični ne znači samo izomorfni. Izraz kolca kvocientov tu označaje prsteny s dozvoljenymi imenovateljami, ne prsteny ostatkovyh klasov. Tute kvantifikacija i tip konstrukcije sut glavno obrazloženje prěvoda.\par
\noindent\emph{Reason:} This is a complete assertion about equality of ideals, not a newly coined headword. The archive reports expansion from a printed witness; that remains historical information, not a fresh visual check of the scan. The current end of Paper 31, §8, item 3 expressly requires equality of the extended ideals in all the localizations defined there. Its argument uses the same prime ideals with the same exponents. All cannot be weakened to one selected localization, and identical does not merely mean isomorphic. Kolca kvocientov here denotes rings with permitted denominators, not residue-class rings. Quantifier scope and construction type are the central translation choices.\par
\noindent{\footnotesize Izvorne zapisy / Original records: HIST-g160-010.}\par
\smallskip\noindent\emph{Sravnjenje s tekstom:} Vsi lokalizacije i ravnost idealov ne možno zaměniti jednoj lokalizacijeju i izomorfijeju.\par
\noindent\emph{Textual comparison:} All localizations and equality of ideals cannot be replaced by one localization and isomorphism.\par
\noindent{\footnotesize TEXT-COMP-667; pasaže / passages: 18932, 18970; \texttt{CORPUS\_FORM\_COMPARISONS.json}.}\par
\end{minipage}\par
\par\medskip\Needspace{7\baselineskip}
\noindent\begin{minipage}{\linewidth}
\hypertarget{TERM-1613}{}
\noindent\textbf{Rimske oznaky VI i VII i svez teoremov}\par\noindent\textbf{Roman VI and VII and the connection between theorems}\quad{\footnotesize\texttt{TERM-1613}}\par
\smallskip\noindent\emph{Razlog:} Izbor v tom zapisu ohranjaje rimsky VI i VII kako iste oznaky v oběh pismach. To jest pravilo identifikacije, ne tvorjenje dvoh novyh terminov. V rabotě 17 teorem VI veže jednakost vida modulov s izomorfijeju grup ostatkov; prěd teoremom VII tekst izrično kaže, že teorem V bude prěformulovany prěz VI. Zato oznaky nesut dokaznu zavisnost i ne smějut byti zaměnjene podobnymi kiriličnymi bukvami. Arhiv popisuje kiriličnu ochranu; nynějšne sravnjenje potvrđuje latinične lokatory i svez, ne novu polnu prověrku kiriličnogo čitatelja.\par
\noindent\emph{Reason:} The choice preserves Roman VI and VII as the same identifiers in both scripts. It is an identification rule, not the creation of two new terms. In Paper 17, Theorem VI connects sameness of type of modules with isomorphism of residue groups; immediately before Theorem VII, the text expressly says that Theorem V is being reformulated through VI. The labels therefore carry a proof dependency and must not be replaced by visually similar Cyrillic letters. The archive describes Cyrillic protection; the present comparison confirms the Latin locations and relation, not a fresh full check of the Cyrillic reader.\par
\noindent{\footnotesize Izvorne zapisy / Original records: HIST-g079-054.}\par
\smallskip\noindent\emph{Sravnjenje s tekstom:} Teorem VI služi prěformulaciji V v VII; rimsky znak jest identifikator.\par
\noindent\emph{Textual comparison:} Theorem VI supports reformulating V as VII; the Roman numeral is an identifier.\par
\noindent{\footnotesize TEXT-COMP-670; pasaže / passages: 11664, 11697, 11699; \texttt{CORPUS\_FORM\_COMPARISONS.json}.}\par
\end{minipage}\par
\par\medskip\Needspace{7\baselineskip}
\noindent\begin{minipage}{\linewidth}
\hypertarget{TERM-1616}{}
\noindent\textbf{Adaptovane imena i točna avtorska identičnost}\par\noindent\textbf{Adapted names and exact author identification}\quad{\footnotesize\texttt{TERM-1616}}\par
\smallskip\noindent\emph{Razlog:} Lurot, Kastelnuovo i Enrikes sut adaptovane osnovy vlastnyh imen, s genitivom pri teorem. Arhiv objašnjaje Lurot kako način izběgnuti mehanično prěpisano th v kirilici. To jest redakcijna pravopisna odluka, ne dokaz, že izvorno ime Lüroth trěba izbrisati iz bibliografije. V rabotě 11 Lurot i Kastelnuovo podpirajut pozitivne tvrdženja o minimalnoj bazi, dokole Enrikes jest citovany za protivny priměr pri vyšem čislu neizvěstnyh. Tri imena v arhivnom spisku zato ne označajut tri sinonimne teoremy. Bibliografija mora dati točnu identičnost citovanyh avtorov; jezyčna adaptacija v prozě ne směje odstraniti avtorsky prinos.\par
\noindent\emph{Reason:} Lurot, Kastelnuovo, and Enrikes are adapted proper-name stems, used in the genitive after teorem. The archive explains Lurot as avoiding a mechanical th rendering in Cyrillic. This is an editorial spelling choice, not a reason to erase Lüroth's original name from bibliographic identification. In Paper 11, Lüroth and Castelnuovo support positive minimal-basis assertions, while Enriques is cited for a counterexample with more indeterminates. The three names in the archival list therefore do not designate three synonymous theorems. Bibliography must identify the actual cited authors; linguistic adaptation in prose must not remove their contributions.\par
\noindent{\footnotesize Izvorne zapisy / Original records: HIST-g073-019.}\par
\smallskip\noindent\emph{Sravnjenje s tekstom:} Imena v jednom spisku ne činęt citovane rezultaty zaměnjivymi.\par
\noindent\emph{Textual comparison:} Names in one list do not make their cited results interchangeable.\par
\noindent{\footnotesize TEXT-COMP-673; pasaže / passages: 8892; \texttt{CORPUS\_FORM\_COMPARISONS.json}.}\par
\end{minipage}\par
\par\medskip\Needspace{7\baselineskip}
\noindent\begin{minipage}{\linewidth}
\hypertarget{TERM-1617}{}
\noindent\textbf{Ostrowski: ime v prozě i bibliografiji}\par\noindent\textbf{Ostrowski: the name in prose and bibliography}\quad{\footnotesize\texttt{TERM-1617}}\par
\smallskip\noindent\emph{Razlog:} Zapis o teoremu Ostrowski ohranjaje historijno latinično napisanje i posebno opisuje kirilično Островски v prozě. Bibliografsky citat i obyčna rečenica ne morajut iměti jednaku grafičnu obrabotku; važno jest, aby oboje vodilo k istomu avtoru. Rusky i ukrajinsky zapisy tu sami ohranjajut latinično ime, zato ne davajut nezavisny dokaz za kiriličnu normu. V rabotě 20 jde o zachovanje absolutnoj nerazložimosti polinoma modulo vsi kromě konečno mnogyh prostyh idealov. Ne trěba to zaměniti drugym poznatym teoremom togo avtora. Nynějšny latiničny korpus potvrđuje ime i konkretny citat; polna kirilična prověrka tu ne jest tvrdžena.\par
\noindent\emph{Reason:} The Ostrowski record preserves the historical Latin spelling and separately describes Cyrillic Островски in running prose. A bibliographic quotation and an ordinary sentence need not receive identical script treatment, but both must identify the same author. The Russian and Ukrainian records here themselves retain the Latin name, so they supply no independent authority for Cyrillic spelling. In Paper 20 the result concerns preservation of absolute polynomial irreducibility modulo all but finitely many prime ideals. It must not be replaced by another well-known theorem of the same author. The current Latin text confirms the name and specific citation; no exhaustive Cyrillic check is claimed here.\par
\noindent{\footnotesize Izvorne zapisy / Original records: HIST-g083-012.}\par
\smallskip\noindent\emph{Sravnjenje s tekstom:} Izključene sut najvyše konečno mnoge proste idealy; citat identifikuje rezultat.\par
\noindent\emph{Textual comparison:} At most finitely many prime ideals are excluded; the citation identifies the result.\par
\noindent{\footnotesize TEXT-COMP-674; pasaže / passages: 13809, 13811, 14011, 14016, 14022; \texttt{CORPUS\_FORM\_COMPARISONS.json}.}\par
\end{minipage}\par
\par\medskip\Needspace{7\baselineskip}
\noindent\begin{minipage}{\linewidth}
\hypertarget{TERM-1618}{}
\noindent\textbf{Odkaz na prvi slučaj teorema X}\par\noindent\textbf{Reference to the first case of Theorem X}\quad{\footnotesize\texttt{TERM-1618}}\par
\smallskip\noindent\emph{Razlog:} Teorema X v tom zapisu jest odkaz, ne novy pojem. Arhiv govori, že pri prěpisu v kirilicu byla opravjena pokvarjena oznaka; samo zabilěženje ne ustanavjaje, že byla provedena ljudska redakcija, ni ne vosstanavjaje nezachovane bukvy zapisane kako ?? v razlogu. Točno rimsky X mora ostati X. V nynějšnom §12 raboty 17 primětka ukazuje na prvi slučaj teorema X; sam teorem razlikuje dvě alternativy a trěbuje beskonečno mnogo konstant v oblasti racionalnosti. Prověrka odkaza zato obhvača cělj, slučaj i prědpoloženja, ne samo podobnost znaka.\par
\noindent\emph{Reason:} Teorema X is a reference in this record, not a new concept. The archive reports correction of a damaged identifier during Cyrillic conversion; that report neither establishes a human editor nor recovers the lost lettering represented as ?? in the rationale. The exact Roman X must remain X. In the current Paper 17, §12, the footnote points to the first case of Theorem X; the theorem distinguishes two alternatives and assumes infinitely many constants in the rationality domain. Checking the reference therefore includes its destination, case, and hypotheses, not merely the shape of the numeral.\par
\noindent{\footnotesize Izvorne zapisy / Original records: HIST-g079-079.}\par
\smallskip\noindent\emph{Sravnjenje s tekstom:} Odkaz trěbuje pravilny rimsky znak i pravilny slučaj teorema.\par
\noindent\emph{Textual comparison:} The reference requires the right Roman numeral and the right case of the theorem.\par
\noindent{\footnotesize TEXT-COMP-675; pasaže / passages: 11951, 12302; \texttt{CORPUS\_FORM\_COMPARISONS.json}.}\par
\end{minipage}\par
\par\medskip\Needspace{7\baselineskip}
\noindent\begin{minipage}{\linewidth}
\hypertarget{TERM-1619}{}
\noindent\textbf{Različne tvrdženja pod oznakami XI i XII}\par\noindent\textbf{Distinct assertions labelled XI and XII}\quad{\footnotesize\texttt{TERM-1619}}\par
\smallskip\noindent\emph{Razlog:} XI i XII sut stabilne rimske identifikatory, ohranjene takože v kiriličnom tekstu po arhivnom pravilu. Prěvod okolnogo naziva ne měnjaje kodovne znaky oznaky. V rabotě 17 XI veže red konečnoj grupy ostatkov s čislom linearno nezavisnyh integralov diferencialnogo sistema; XII govori o diferencialnyh izrazah, ktore prěvodęt integraly medžu modulami jednakogo vida. Čisla zato razlikujut dve tvrdženja i jih odkazy, ne dvě formy istogo termina. Nynějšne pasaže potvrđajut tu razliku. Iz pravilno prikazanoj latiničnoj oznaky same ne možno zaključiti, že vsegda byla pravilna i kirilična projekcija.\par
\noindent\emph{Reason:} XI and XII are stable Roman identifiers, retained in Cyrillic text under the recorded rule. Translating the surrounding label does not change the identifier's characters. In Paper 17, XI relates the order of a finite residue group to the number of linearly independent integrals of a differential system; XII concerns differential expressions carrying integrals between modules of the same type. The numbers therefore distinguish two assertions and their references, not two forms of one term. Current passages confirm that distinction. A correctly displayed Latin identifier alone does not prove that the Cyrillic projection has always been correct.\par
\noindent{\footnotesize Izvorne zapisy / Original records: HIST-g079-070.}\par
\smallskip\noindent\emph{Sravnjenje s tekstom:} Dve oznaky razlikujut čislovy rezultat i odnos medžu sistemami integralov.\par
\noindent\emph{Textual comparison:} The identifiers distinguish a counting result from a relation between systems of integrals.\par
\noindent{\footnotesize TEXT-COMP-676; pasaže / passages: 11986, 11988, 11989, 12036, 12038; \texttt{CORPUS\_FORM\_COMPARISONS.json}.}\par
\end{minipage}\par
\par\medskip\Needspace{7\baselineskip}
\noindent\begin{minipage}{\linewidth}
\hypertarget{TERM-1620}{}
\noindent\textbf{Vosstanovjena iztočna věta o determinantu}\par\noindent\textbf{Restored source sentence about a determinant}\quad{\footnotesize\texttt{TERM-1620}}\par
\smallskip\noindent\emph{Razlog:} Razširjenje po pečatnom skanu jest redakcijna primětka o vosstanovjenju iztočnogo teksta, ne matematično slovo avtora. Arhiv svęzuje ju s větoju, že determinant matrice proměny bazy jest jedinica; ta věta nedostavala v tadašnjom prěpisanom odlomku. Nynějšny §4 raboty 31 ju sodržuje i upotrěbjaje obratnu matricu v relaciji ekvivalentnosti. Jedinica tu jest silnějše uslovje než nenulovy determinant nad oblastju cělosti: ona dava obratnu matricu nad toj samoj oblastju. Novo sravnjenje potvrđuje prisutnost věty, ne povtorjaje arhivnu prověrku pečatnogo svědka. Takovo vosstanovjenje trěba dokazati sravnjenjem iztočnikov; ne směje se vydavati za avtorskogo novogo termina ili za tihu matematičnu popravu.\par
\noindent\emph{Reason:} Razširjenje po pečatnom skanu is an editorial note about restoring source text, not the author's mathematical vocabulary. The archive links it to the sentence saying that the determinant of a basis-change matrix is a unit, missing from the then-current transcription excerpt. Paper 31, §4 now contains that sentence and uses the inverse matrix in the equivalence relation. Unit is stronger here than nonzero determinant over an integral domain: it permits inversion over that same domain. The new comparison confirms presence of the sentence, not a repeat of the archival printed-witness check. Such restoration needs source comparison; it must not be presented as a newly coined authorial term or a silent mathematical correction.\par
\noindent{\footnotesize Izvorne zapisy / Original records: HIST-g132-005.}\par
\smallskip\noindent\emph{Sravnjenje s tekstom:} Determinant-jedinica opravdava upotrěbu obratnoj matrice nad istoju oblastju.\par
\noindent\emph{Textual comparison:} A unit determinant licenses use of the inverse matrix over the same base.\par
\noindent{\footnotesize TEXT-COMP-677; pasaže / passages: 17827, 17909, 17911, 17913, 17918, 17921; \texttt{CORPUS\_FORM\_COMPARISONS.json}.}\par
\end{minipage}\par
\par\medskip\Needspace{7\baselineskip}
\noindent\begin{minipage}{\linewidth}
\hypertarget{TERM-1638}{}
\noindent\textbf{Bibliografičny titul i lokalna oznaka stranic}\par\noindent\textbf{Bibliographic title and localized page label}\quad{\footnotesize\texttt{TERM-1638}}\par
\smallskip\noindent\emph{Razlog:} Sitz. Ber. d. Preuß. Akad. d. Wiss. jest skračeny titul periodičnogo izdanja, ne matematičny termin za prěvod. Zapisany izbor čuva titul i bibliografične čisla, ale lokalizuje S. kako str.; ruska i ukrajinska verzija imajut c. za stranicu. Nynějšny početok raboty 32 potvrđuje 1927, str. 221--228 i sovavtorstvo R. Brauera. Čitatelj tako može najdti original, a obyčna oznaka stranic ne trěbuje znanja němčiny. Redakcijne razsmotrenje jedinstvenogo bibliografičnogo stila ne dava prava odstraniti ime sovavtora ili zameniti iztočnik novym izdanjem bez objašnjenja.\par
\noindent\emph{Reason:} Sitz. Ber. d. Preuß. Akad. d. Wiss. is an abbreviated periodical title, not mathematical vocabulary to translate. The recorded choice preserves the title and bibliographic numbers but localizes S. as str.; the Russian and Ukrainian versions use their page abbreviation c. The current opening of paper 32 confirms 1927, pp. 221--228 and coauthorship with R. Brauer. A reader can therefore locate the original without needing German to understand the ordinary page label. Review for a consistent bibliographic style does not authorize removing a coauthor or substituting a later edition without explanation.\par
\noindent{\footnotesize Izvorne zapisy / Original records: HIST-g166-005.}\par
\smallskip\noindent\emph{Sravnjenje s tekstom:} Lokalizovana oznaka stranic čuva identitet publikacije i sovavtorstvo.\par
\noindent\emph{Textual comparison:} The localized page label preserves the publication's identity and coauthorship.\par
\noindent{\footnotesize TEXT-COMP-695; pasaže / passages: 19171, 19173, 19176; \texttt{CORPUS\_FORM\_COMPARISONS.json}.}\par
\end{minipage}\par
\par\medskip\Needspace{7\baselineskip}
\noindent\begin{minipage}{\linewidth}
\hypertarget{TERM-1640}{}
\noindent\textbf{Avtorske ograničenje povtorjanih redukcij}\par\noindent\textbf{Authorial restriction on repeating routine reductions}\quad{\footnotesize\texttt{TERM-1640}}\par
\smallskip\noindent\emph{Razlog:} Dalej ne budut spominjane čuva avtorskym glasom oznajmenje o sposobu izloženja. V §25 raboty 2 ide točno o redukcijah, ktore neposrědnje nastavajut kontrakcijeju s reducentami ili razkladnymi formami; potrěbne daljše konstrukcije slědujut dalje. Zapisany razlog jest modernizacija prozy bez proměny toj ograničenoj odluki avtorky. Fraza ne jest instrukcija prěvoditelju, da može izpustiti težke dokazy, formulne kroky ili dokumentaciju svoji izborov. Njen obseg mora ostati svezany s navedenym razredom rutinnyh redukcij. Zato ona pripada pravilam izloženja v slovniku izborov, ne spisku novyh matematičnyh imenikov.\par
\noindent\emph{Reason:} Dalej ne budut spominjane preserves an announcement, in the author's voice, about the exposition. Paper 2, §25 refers specifically to reductions obtained directly by contraction with reducers or decomposable forms; further constructions still follow. The recorded reason is to modernize the prose without changing this limited authorial decision. The sentence is not an instruction permitting a translator to omit difficult proofs, formula steps or documentation of translation choices. Its scope must remain attached to the specified class of routine reductions. It therefore belongs among expository choices, not in a list of newly coined mathematical nouns.\par
\noindent{\footnotesize Izvorne zapisy / Original records: HIST-g026-005.}\par
\smallskip\noindent\emph{Sravnjenje s tekstom:} Odluka ne ponavljati rutinne redukcije ne jest dozvoljenje izpustiti drugy tekst.\par
\noindent\emph{Textual comparison:} The decision not to repeat routine reductions is not permission to omit other text.\par
\noindent{\footnotesize TEXT-COMP-697; pasaže / passages: 3258, 3260; \texttt{CORPUS\_FORM\_COMPARISONS.json}.}\par
\end{minipage}\par
\par\medskip\Needspace{7\baselineskip}
\noindent\begin{minipage}{\linewidth}
\hypertarget{TERM-1641}{}
\noindent\textbf{Povtorjena formula i iztočna věrnost}\par\noindent\textbf{Repeated formula and source fidelity}\quad{\footnotesize\texttt{TERM-1641}}\par
\smallskip\noindent\emph{Razlog:} Povtor $H_j^2$, $H_j^2$ jest odluka o věrnosti prěpisanja, ne nova matematična nazva. Arhiv govori, že redukcijna věta povtorjaje formulu i že povtor byl ohranjen do sravnjenja s pečatnym iztočnikom. Nynějšny §16 raboty 2 ješče ima oba jednaka izrazy. To potvrđuje sostojanje teksta, ne to, že avtorka naměrno hotěla povtor. Bez iztočnogo dokazanja nije opravdano proizvoljno změniti indeks ili stepen. Ako pečatny original sam ima pogrěšku, diplomatski tekst ju čuva, a predložena poprava pripada oddělnomu kritičnomu izdanju.\par
\noindent\emph{Reason:} The repetition $H_j^2$, $H_j^2$ records a transcription-fidelity decision, not a new mathematical name. The archive says that the reduction sentence repeats the formula and that the repetition was retained pending comparison with print. Current paper 2, §16 still has both identical expressions. This establishes the text's state, not that the author intended the repetition. An index or exponent must not be changed arbitrarily without source evidence. If the printed original itself contains an error, a diplomatic text preserves it and a proposed emendation belongs in a separate critical edition.\par
\noindent{\footnotesize Izvorne zapisy / Original records: HIST-g017-014.}\par
\smallskip\noindent\emph{Sravnjenje s tekstom:} Dva jednaka izrazy ostavajut v tekstu; pravilnost povtora ne jest tym dokazana.\par
\noindent\emph{Textual comparison:} Both identical expressions remain in the text; this does not establish that the repetition is correct.\par
\noindent{\footnotesize TEXT-COMP-698; pasaže / passages: 2351, 2353, 2355; \texttt{CORPUS\_FORM\_COMPARISONS.json}.}\par
\end{minipage}\par
\par\medskip\Needspace{7\baselineskip}
\noindent\begin{minipage}{\linewidth}
\hypertarget{TERM-1642}{}
\noindent\textbf{Pomoćne polinomy: prěvod i kritična poprava}\par\noindent\textbf{Auxiliary polynomials: translation and critical correction}\quad{\footnotesize\texttt{TERM-1642}}\par
\smallskip\noindent\emph{Razlog:} Prěvod $A$ i $B$ jest zabeležena kritična poprava němčinskogo zapisa $A$ und psi. Razlog jest matematičny: v identičnosti $K=A\varphi+B\psi$ iskane pomoćne polinomy sut $A$ i $B$, a $\varphi$ i $\psi$ sut dane. Nynějšna rabota 45 ima $A$ i $B$. Ta soglasnost s formuloju objasnjaje motivaciju, ale ne dokazuje, čto točno stoji v pečatnom originalu; arhiv odkazuje na tadašnji prěpis i stranovy tekst. Izbor zato ne možno skryti pod nazvoju obyčnogo prěvoda. Diplomatska i popravjena verzija morajut jasno razlikovati iztočno čitanje i redakcijnu intervenciju.\par
\noindent\emph{Reason:} The rendering $A$ and $B$ is a recorded critical correction of German wording reported as $A$ und psi. Its reason is mathematical: in $K=A\varphi+B\psi$, the auxiliary polynomials sought are $A$ and $B$, whereas $\varphi$ and $\psi$ are given. Current paper 45 has $A$ and $B$. Agreement with the formula explains the motivation but does not establish what the printed original says; the archive refers to its then-staged transcription and page text. This choice must not be hidden as ordinary translation. Diplomatic and corrected versions must distinguish the source reading from the editorial intervention.\par
\noindent{\footnotesize Izvorne zapisy / Original records: HIST-g214-014.}\par
\smallskip\noindent\emph{Sravnjenje s tekstom:} Prěvod $A$ i $B$ soglašaje se s identičnostju, ale predstavlja zabeleženu intervenciju.\par
\noindent\emph{Textual comparison:} The wording $A$ and $B$ agrees with the identity but represents a recorded intervention.\par
\noindent{\footnotesize TEXT-COMP-699; pasaže / passages: 45-isv.tex:78, 45-isv.tex:79, 45-isv.tex:80, 45-isv.tex:81, 45-isv.tex:83, 45-isv.tex:85; \texttt{CORPUS\_FORM\_COMPARISONS.json}.}\par
\end{minipage}\par
\par\medskip\Needspace{7\baselineskip}
\noindent\begin{minipage}{\linewidth}
\hypertarget{TERM-1676}{}
\noindent\textbf{Simbolične formy s faktorom H}\par\noindent\textbf{Symbolic forms with factor H}\quad{\footnotesize\texttt{TERM-1676}}\par
\smallskip\noindent\emph{Razlog:} Arhivno theta-H / f-H jest kratky zapis simboličnyh form, ne para novyh slov. Nynějšne pasaže pišut $\theta\cdot H$ i $f\cdot H$ v kontekstu jednokratnoj kontrakcije. V čitateljskom prědstavjenju zato pišemo grčsku bukvu i točku množenja, a ne puščajemo, da arhivny defis vypada kako znak odimanja. To jest objasnjeno čitanje matematičnoj notacije; izvorna arhivna grafija ostaje dostupna v zapisu. Razlog prěvoda čuva odkaz na simboličny sistem, vměsto zaměniti go neoprědělennym slovnym opisom.\par
\noindent\emph{Reason:} Archived theta-H / f-H is shorthand for symbolic forms, not a pair of new words. The current passages write $\theta\cdot H$ and $f\cdot H$ in the context of onefold contraction. The reading form therefore uses the Greek letter and multiplication dot rather than allowing the archival hyphen to look like subtraction. This is an explained reading of mathematical notation; the original archival spelling remains available in the record. The translation rationale preserves the reference to the symbolic system instead of replacing it with an indefinite verbal description.\par
\noindent{\footnotesize Izvorne zapisy / Original records: HIST-g025-011.}\par
\smallskip\noindent\emph{Sravnjenje s tekstom:} Arhivny defis ne čita se kako minus.\par
\noindent\emph{Textual comparison:} The archival hyphen is not read as a minus sign.\par
\noindent{\footnotesize TEXT-COMP-733; pasaže / passages: 3176, 3192; \texttt{CORPUS\_FORM\_COMPARISONS.json}.}\par
\end{minipage}\par
\par\medskip\Needspace{7\baselineskip}
\noindent\begin{minipage}{\linewidth}
\hypertarget{TERM-1689}{}
\noindent\textbf{Čuvanje imena Sistem III}\par\noindent\textbf{Preserving the name System III}\quad{\footnotesize\texttt{TERM-1689}}\par
\smallskip\noindent\emph{Razlog:} Sistem III jest strukturno ime v rabotě, povtorjeno v několiko paragrafah. Prevod imenika ne opravdava zaměnu rimskogo III arabskym 3, ni kiriličnymi bukvami, ktore slično izgledajut. Zaglavje §14 ukazuje, že vnutri Sistema III obgovarjajut se formy devetogo poreda; III zato ne jest poredok formy. Zaključny pregled pripisuje tomu sistemu 117 form. Razlog jezyčnogo izbora jest kontinuitet odkazov, a ne tvorjenje nove matematične klasifikacije.\par
\noindent\emph{Reason:} Sistem III is a structural name recurring across several sections. Translating the noun does not justify replacing Roman III by Arabic 3 or by similar-looking Cyrillic letters. The heading of section 14 discusses forms of order nine within System III, so III is not the order of a form. The closing synopsis assigns the system 117 forms. The reason for the language choice is continuity of references, not the creation of a new mathematical classification.\par
\noindent{\footnotesize Izvorne zapisy / Original records: HIST-g014-002, HIST-g015-002, HIST-g016-002.}\par
\smallskip\noindent\emph{Sravnjenje s tekstom:} Rimski numer identifikuje sistem, a ne stepen ili poredok formy.\par
\noindent\emph{Textual comparison:} The Roman numeral identifies the system, not the degree or order of a form.\par
\noindent{\footnotesize TEXT-COMP-746; pasaže / passages: 2149, 3512; \texttt{CORPUS\_FORM\_COMPARISONS.json}.}\par
\end{minipage}\par
\par\medskip\Needspace{7\baselineskip}
\noindent\begin{minipage}{\linewidth}
\hypertarget{TERM-1699}{}
\noindent\textbf{Navod k čitanju tabely}\par\noindent\textbf{An instruction for reading the table}\quad{\footnotesize\texttt{TERM-1699}}\par
\smallskip\noindent\emph{Razlog:} Provodimo tabelarno govori čitatelju, kako bude izložena konstrukcija form jednakogo poreda. To jest izbor obyčnogo dokaznogo registra, ne novy algebraičny termin. Imenno tabelarno svazuje rečenicu s nadslědujučimi strukturnymi podatkami; izpustiti jeho značilo by oslabiti navod dlja čitanja. Prva osoba množstvena prěnosi golos avtorskogo izloženja, ne tvrdženje, že prevoditelj ili čitatelj izvršil naučno izslědovanje. Polna tabela mora ostati čitljiva zajedno s rečenicoju.\par
\noindent\emph{Reason:} Provodimo tabelarno tells the reader how the construction of forms of equal order will be presented. This is ordinary expository language, not a new algebraic term. Tabelarno specifically connects the sentence with the structured material that follows; omitting it would weaken the reading instruction. The first-person plural conveys the voice of the exposition, not a claim that the translator or reader conducted the research. The complete table must remain readable together with the sentence.\par
\noindent{\footnotesize Izvorne zapisy / Original records: HIST-g008-007.}\par
\smallskip\noindent\emph{Sravnjenje s tekstom:} Navod k tabeli jest čest avtorskogo izloženja.\par
\noindent\emph{Textual comparison:} The instruction referring to the table is part of the author's exposition.\par
\noindent{\footnotesize TEXT-COMP-756; pasaže / passages: 1076; \texttt{CORPUS\_FORM\_COMPARISONS.json}.}\par
\end{minipage}\par
\par\medskip\Needspace{7\baselineskip}
\noindent\begin{minipage}{\linewidth}
\hypertarget{TERM-1701}{}
\noindent\textbf{Česti časopisa v bibliografiji}\par\noindent\textbf{Journal parts in bibliographic references}\quad{\footnotesize\texttt{TERM-1701}}\par
\smallskip\noindent\emph{Razlog:} Čest I i čest II označajut česti izdanja časopisa, ne odděly dokazov. Vybrano kratko čest čuvaje bibliografičnu informaciju vměstě s tomom, rokom i stranicami. Zaměna s glava ili s tom mogla by poslati čitatelja k drugomu městu. Ruske i ukrajinske sravnjenja podpirajut smysl česti; iz njih samyh ne slěduje, že vybrana medžuslovjanska forma jest jedinym možnym izrazom. V nynějšnom spisku komunikacij čest I i čest II ostajut čitljivi, a originalne čisla ne sut lokalizovane.\par
\noindent\emph{Reason:} Čest I and čest II identify parts of the journal publication, not divisions of an argument. The short word čest preserves the bibliographic distinction alongside the volume, year, and pages. Replacing it with chapter or volume could direct readers to the wrong place. The Russian and Ukrainian comparisons support the part sense, not the claim that the chosen Interslavic wording is the only possible form. The current list of communications retains both part labels without altering the original numbers.\par
\noindent{\footnotesize Izvorne zapisy / Original records: HIST-g215-011.}\par
\smallskip\noindent\emph{Sravnjenje s tekstom:} Nazvy čest I i čest II ostajut bibliografičnymi lokatorami.\par
\noindent\emph{Textual comparison:} The part labels remain bibliographic locators.\par
\noindent{\footnotesize TEXT-COMP-758; pasaže / passages: bib-isv.tex:179, bib-isv.tex:180, bib-isv.tex:181; \texttt{CORPUS\_FORM\_COMPARISONS.json}.}\par
\end{minipage}\par
\par\medskip\Needspace{7\baselineskip}
\noindent\begin{minipage}{\linewidth}
\hypertarget{TERM-1734}{}
\noindent\textbf{Podsistemy i rimske indekse}\par\noindent\textbf{Subsystems and Roman subscripts}\quad{\footnotesize\texttt{TERM-1734}}\par
\smallskip\noindent\emph{Razlog:} Podsistemy jest prěvodimy imenik, ale $s_I$, $s_{II}$, $s_{III}$ i $s_{IV}$ sut stabilne matematične nazvy. Rimske indekse razlučajut konstrukcije prěz različne kontrakcije; ne sut čisla stepena, ktore možno proizvoljno prepismovati. Zato v kiriličnom čitanju mora měnjati se okolna proza, ne ty identifikatory. Arhiv izrično čuvaje rimske labely. To jest pravilo notacije povezano s terminom, ne dodatny slovnikovy sinonim.\par
\noindent\emph{Reason:} Podsistemy is a translatable noun, but $s_I$, $s_{II}$, $s_{III}$, and $s_{IV}$ are stable mathematical names. Their Roman subscripts distinguish different contraction constructions; they are not degree numerals to be freely respelled. A Cyrillic reading must therefore change the surrounding prose, not those identifiers. The record explicitly protects the Roman labels. This is a notation rule associated with the term, not another lexical synonym.\par
\noindent{\footnotesize Izvorne zapisy / Original records: HIST-g018-010.}\par
\smallskip\noindent\emph{Sravnjenje s tekstom:} Indekse razlikujut konkretne konstrukcije.\par
\noindent\emph{Textual comparison:} The subscripts distinguish particular constructions.\par
\noindent{\footnotesize TEXT-COMP-791; pasaže / passages: 2474, 2477, 2478, 2479, 2480; \texttt{CORPUS\_FORM\_COMPARISONS.json}.}\par
\end{minipage}\par
\par\medskip\Needspace{7\baselineskip}
\noindent\begin{minipage}{\linewidth}
\hypertarget{TERM-1749}{}
\noindent\textbf{Učestvovanje v kontrakciji}\par\noindent\textbf{Participation in a contraction}\quad{\footnotesize\texttt{TERM-1749}}\par
\smallskip\noindent\emph{Razlog:} Vhodi v svijanje jest glagolna konstrukcija za učestvovanje formy v konkretnoj invariantnoj operaciji. Arhiv ne sozdava novy tehničny glagol; obyčno vhodi povezuje subjekt s nazvoju operacije. Nynějšne pasaže pišut vhodi v kontrakciju. Taka zaměna imenika može uskladniti terminologiju, ale ne smi izbrisati uslovje, že forma nastupaje samo v produktah ili stepenjah. Arhivno tvrđenje o širokoj razumělivosti jest motivacija izbora, ne rezultat čitatelskogo testa.\par
\noindent\emph{Reason:} Vhodi v svijanje is a verbal construction for a form's participation in a particular invariant-theoretic operation. The record avoids coining a new technical verb: ordinary vhodi connects the subject with the operation's name. The current passages use vhodi v kontrakciju. Changing the noun can align terminology, but must not erase restrictions to products or powers. The archival claim of broad intelligibility is a motivation, not a reader-test result.\par
\noindent{\footnotesize Izvorne zapisy / Original records: HIST-g015-014.}\par
\smallskip\noindent\emph{Sravnjenje s tekstom:} Glagolna konstrukcija ne zaměnjaje uslovja operacije.\par
\noindent\emph{Textual comparison:} The verbal construction does not replace the operation's conditions.\par
\noindent{\footnotesize TEXT-COMP-806; pasaže / passages: 1802, 2217; \texttt{CORPUS\_FORM\_COMPARISONS.json}.}\par
\end{minipage}\par
\par\medskip\Needspace{7\baselineskip}
\noindent\begin{minipage}{\linewidth}
\hypertarget{TERM-1750}{}
\noindent\textbf{Negacija v opisu kontrakcije}\par\noindent\textbf{Negation in a contraction statement}\quad{\footnotesize\texttt{TERM-1750}}\par
\smallskip\noindent\emph{Razlog:} Ne vhodi v svijanje jest negativna forma toj samej konstrukcije. Negacija odnosi se k učestvovanju v imenovanoj kontrakciji, ne k obstavanju formy ili k tomu, že ona jest nula. Pasaže razlikujut nevhodenje v obče od nevhodenja v produktah i od razgledanja kako kovarianta ili kontravarianta. Zato kopirovanje jednogo kratkogo izraza bez njegovyh ograničenij moglo by izměniti dokaz. Vhodi i vhodet dodatno zahtěvajut soglasovanje s jedinym ili množestvenym subjektom, vključajuči simbolične izrazy.\par
\noindent\emph{Reason:} Ne vhodi v svijanje negates the same construction. It negates participation in the named contraction, not the form's existence or its being nonzero. The passages distinguish exclusion altogether from exclusion in products, and distinguish treatment as a covariant or contravariant. Copying the short expression without those qualifications could change the proof. Vhodi and vhodet also require agreement with a singular or plural subject, including symbolic expressions.\par
\noindent{\footnotesize Izvorne zapisy / Original records: HIST-g014-017, HIST-g015-015, HIST-g016-010.}\par
\smallskip\noindent\emph{Sravnjenje s tekstom:} Subjekt, partner i ograničenje morajut ostati jasne.\par
\noindent\emph{Textual comparison:} Subject, partner, and qualification must remain clear.\par
\noindent{\footnotesize TEXT-COMP-807; pasaže / passages: 2004, 2126, 2128, 2219, 2271; \texttt{CORPUS\_FORM\_COMPARISONS.json}.}\par
\end{minipage}\par
\par\medskip\Needspace{7\baselineskip}
\noindent\begin{minipage}{\linewidth}
\hypertarget{TERM-1751}{}
\noindent\textbf{Negacija ograničena na stepenje}\par\noindent\textbf{Negation restricted to powers}\quad{\footnotesize\texttt{TERM-1751}}\par
\smallskip\noindent\emph{Razlog:} V stepenjah ne vhodi v svijanje prěměšča ograničenje k početku rečenice, aby jego oblast byla vidima. Ne znači to, že forma nikogda ne može učestvovati v kontrakciji. Nynějšna pasaža o $s$ razlučaje potencije i produkty s kontragredientnymi formami, i iz toga izvodi, ktore produkty ješče trěba razgleděti. Slovnik zato mora čuvati uslovje o stepenjah i imenovane partnery kontrakcije. Slědovanje njemeckomu poredku slov ne jest važnějše než ta jasna oblast negacije.\par
\noindent\emph{Reason:} V stepenjah ne vhodi v svijanje moves the qualification to the beginning so its scope is visible. It does not mean the form can never enter a contraction. The current passage about $s$ distinguishes powers from products with contragredient forms, then determines which products remain to be considered. The dictionary must preserve the restriction to powers and the named contraction partners. Clear negation scope matters more than reproducing German word order.\par
\noindent{\footnotesize Izvorne zapisy / Original records: HIST-g019-009.}\par
\smallskip\noindent\emph{Sravnjenje s tekstom:} Ograničenje na potencije jest čest tvrdženja.\par
\noindent\emph{Textual comparison:} The restriction to powers is part of the assertion.\par
\noindent{\footnotesize TEXT-COMP-808; pasaže / passages: 2617, 2625, 2676; \texttt{CORPUS\_FORM\_COMPARISONS.json}.}\par
\end{minipage}\par
\par\medskip\Needspace{7\baselineskip}
\noindent\begin{minipage}{\linewidth}
\hypertarget{TERM-1756}{}
\noindent\textbf{Tvrdženje i dva pismena oblika}\par\noindent\textbf{Tvrdženje across the two scripts}\quad{\footnotesize\texttt{TERM-1756}}\par
\smallskip\noindent\emph{Razlog:} Tvrdženje jest izbor jezyčnoj i pismennej formy za statement, ne novy termin Noether. Arhiv razlučaje ranje tvrđenje s đ od vybranej formy s dž i svazuje izbor s pravilnym prěhodom v kirilicu. Jeho sravniteljne ruske i ukrajinske slova ne sut citati iz originala. Nynějšny latinski tekst ima tvrdženje v običnoj dokaznoj proze. Tehnično uspešno prěpisovanje samo po sobě ne dokazuje pravilnost jezyčnogo izbora; ono objasnjaje jedinu zabilježenu pričinu normalizacije.\par
\noindent\emph{Reason:} Tvrdženje is a linguistic and script choice for statement, not a new Noether term. The archive distinguishes earlier tvrđenje with đ from the selected dž form and connects the choice with consistent Cyrillic conversion. Its Russian and Ukrainian comparisons are not quotations from the original. Current Latin text uses tvrdženje in ordinary proof prose. Successful script conversion does not by itself establish linguistic correctness; it explains one recorded reason for normalization.\par
\noindent{\footnotesize Izvorne zapisy / Original records: HIST-g062-013.}\par
\smallskip\noindent\emph{Sravnjenje s tekstom:} Nynějšna latinska forma ne jest njemecky izvorni termin.\par
\noindent\emph{Textual comparison:} The current Latin form is not a German source term.\par
\noindent{\footnotesize TEXT-COMP-813; pasaže / passages: 7539, 7540, 7541, 7542; \texttt{CORPUS\_FORM\_COMPARISONS.json}.}\par
\end{minipage}\par
\par\medskip\Needspace{7\baselineskip}
\noindent\begin{minipage}{\linewidth}
\hypertarget{TERM-1764}{}
\noindent\textbf{Nezavisnost od vybranoj klasi}\par\noindent\textbf{Independence from the chosen class}\quad{\footnotesize\texttt{TERM-1764}}\par
\smallskip\noindent\emph{Razlog:} Ne zavisi od klasa položenogo v osnovu jest tvrdženje o nezavisnosti od vybranoj klasi, ne o statističnoj nezavisnosti. V razgledanom prstenu bez nulovyh děliteljev ono kasaje se diskriminanta ideala: vsak element determinanta ima tu nezavisnost. To ne odstraňaje drugu, osobno opisanu zavisnost od bazy do kvadrata jedinice. Fraza položenogo v osnovu sohranjaje predmet izbora; kračše ne zavisi od ničego bylo by matematično ložno.\par
\noindent\emph{Reason:} Ne zavisi od klasa položenogo v osnovu states independence from the chosen class, not statistical independence. For the ring without zero divisors under consideration it concerns the ideal's discriminant: each determinant entry has this independence. It does not remove the separately described change by the square of a unit under change of basis. Položenogo v osnovu preserves what is being chosen; shortening the phrase to independent of everything would be mathematically false.\par
\noindent{\footnotesize Izvorne zapisy / Original records: HIST-g137-006.}\par
\smallskip\noindent\emph{Sravnjenje s tekstom:} Nezavisnost od klasa ne odstraňaje faktor kvadrata jedinice pri zaměně bazy.\par
\noindent\emph{Textual comparison:} Independence from the class does not remove the square-unit factor under a basis change.\par
\noindent{\footnotesize TEXT-COMP-821; pasaže / passages: 18094, 18102, 18104, 18106; \texttt{CORPUS\_FORM\_COMPARISONS.json}.}\par
\end{minipage}\par
\par\medskip\Needspace{7\baselineskip}
\noindent\begin{minipage}{\linewidth}
\hypertarget{TERM-1780}{}
\noindent\textbf{Odkaz na ranějše redukcije}\par\noindent\textbf{Referring to earlier reductions}\quad{\footnotesize\texttt{TERM-1780}}\par
\smallskip\noindent\emph{Razlog:} S učtenjem redukcij togo paragrafa jest spojna fraza, ktora upominaje raněje opravdane redukcije v tym samym oddělu. Aktualny tekst ima s uvažanjem redukcij tutogo paragrafa. V oběh formah zavisnost vyvoda mora ostati vidljiva: formuly za red form $(\theta gu)^2$ dobyvajut se iz (27) s tymi redukcijami, ne iz samoj formuly bez dodatnyh krokov. Izbor modernogo porjadka slov jest stilističny; ne daje dozvoljenje izpustiti referenciju ili matematiku.\par
\noindent\emph{Reason:} S učtenjem redukcij togo paragrafa is a connective phrase referring to reductions already justified in the same section. The current text has s uvažanjem redukcij tutogo paragrafa. In both forms the dependency must remain visible: the formulas for the family $(\theta gu)^2$ follow from (27) together with those reductions, not from the formula alone without additional steps. Choosing modern word order is stylistic; it does not license omission of the reference or the mathematics.\par
\noindent{\footnotesize Izvorne zapisy / Original records: HIST-g022-007.}\par
\smallskip\noindent\emph{Sravnjenje s tekstom:} Aktualno s uvažanjem i arhivno s učtenjem imajut istu dokaznu ulogu.\par
\noindent\emph{Textual comparison:} Current s uvažanjem and archived s učtenjem have the same role in the deduction.\par
\noindent{\footnotesize TEXT-COMP-837; pasaže / passages: 2883; \texttt{CORPUS\_FORM\_COMPARISONS.json}.}\par
\end{minipage}\par
\par\medskip\Needspace{7\baselineskip}
\noindent\begin{minipage}{\linewidth}
\hypertarget{TERM-1781}{}
\noindent\textbf{Jednočasno razgledanje form}\par\noindent\textbf{Considering the forms together}\quad{\footnotesize\texttt{TERM-1781}}\par
\smallskip\noindent\emph{Razlog:} Pri jednovrěmennom razsmotrenju jest uslov postupka, ne samo časovy ukazatelj. V redukcijnoj primětke trěba razgledati skupaj navedene formy i njihove razklady; ignorovanje drugoj skupiny moglo by proměniti ono, čto se reducira. Aktualna fraza pri jednočasnom razgledanju jest kračša i znači to samo. Arhiv označuje izvorno formulovanje kako potrěbujuče jezyčnu prověrku. Slovnik sohranjaje tu ogradu, ale objasnjaje dokaznu ulogu bez služebnogo jezyka.\par
\noindent\emph{Reason:} Pri jednovrěmennom razsmotrenju is a condition on the procedure, not merely a time marker. In the reduction note the listed forms and their decompositions must be considered together; ignoring the second group could change what is being reduced. The current phrase pri jednočasnom razgledanju is shorter and expresses the same relation. The archive marks the earlier phrasing as needing linguistic review. The dictionary preserves that reservation while explaining the phrase's role in the argument.\par
\noindent{\footnotesize Izvorne zapisy / Original records: HIST-g026-009.}\par
\smallskip\noindent\emph{Sravnjenje s tekstom:} Jednočasno razgledanje spojuje navedene formy s njihovymi razkladami.\par
\noindent\emph{Textual comparison:} Considering the forms together connects them with their decompositions.\par
\noindent{\footnotesize TEXT-COMP-838; pasaže / passages: 3309; \texttt{CORPUS\_FORM\_COMPARISONS.json}.}\par
\end{minipage}\par
\par\medskip\Needspace{7\baselineskip}
\noindent\begin{minipage}{\linewidth}
\hypertarget{TERM-1782}{}
\noindent\textbf{Bibliografičny kredit za učestje}\par\noindent\textbf{Bibliographic credit for participation}\quad{\footnotesize\texttt{TERM-1782}}\par
\smallskip\noindent\emph{Razlog:} S učestjem znači s prinosom ili učestjem imenovanoj osoby. Bibliografična rubrika obhvačaje knjigy drugih avtorov, v izdanju ktoryh Emmy Noether učestvovala: zapisy za Dedekinda i Cantora osobno imenujut vydavatelje. Fraza ne čini ju jedinoju avtoriceju tih knjig ni jedinym redaktorom. Izbor čuva točnu razliku medžu avtorstvom i učestjem. Takovy kredit dolžen ostati pri prěvodu zaglavij i pri preuređenju bibliografije.\par
\noindent\emph{Reason:} S učestjem means with the contribution or participation of the named person. The bibliographic category includes books by other authors in whose publication Emmy Noether participated: the Dedekind and Cantor entries separately name the editors. The phrase does not make her the sole author of those books or their sole editor. The choice preserves the distinction between authorship and participation. That credit must survive translation of headings and rearrangement of the bibliography.\par
\noindent{\footnotesize Izvorne zapisy / Original records: HIST-g215-013.}\par
\smallskip\noindent\emph{Sravnjenje s tekstom:} Rubrika i konkretne kredity imenujut različne ulogy.\par
\noindent\emph{Textual comparison:} The heading and individual credits name different roles.\par
\noindent{\footnotesize TEXT-COMP-839; pasaže / passages: bib-isv.tex:213, bib-isv.tex:217, bib-isv.tex:220; \texttt{CORPUS\_FORM\_COMPARISONS.json}.}\par
\end{minipage}\par
\par\medskip\Needspace{7\baselineskip}
\noindent\begin{minipage}{\linewidth}
\hypertarget{TERM-1795}{}
\noindent\textbf{Metafora cělogo i drobnogo}\par\noindent\textbf{The integral–fractional metaphor}\quad{\footnotesize\texttt{TERM-1795}}\par
\smallskip\noindent\emph{Razlog:} Razlika med cělym i drobnym razmyvaje se sohranjaje avtorsku završnu formulaciju. V karakteristice p prvotno polje ne sodrživaje drobnyh veličin v smyslu uvedenej definicije; zato něktore dodatne uslovja otpadajut i klasa dopustimyh oblastij cělosti razširjaje se. To ne znači, že vsaka racionalna funkcija stava polinomom ili že imenovatelje možno ignorovati. Fraza ima strukturny smysl o definicijah, ne bezuslovnu algebraičnu ravnost cělogo i drobnogo.\par
\noindent\emph{Reason:} Razlika med cělym i drobnym razmyvaje se preserves the author's closing formulation. In characteristic p the prime field contains no fractional quantities in the sense of the definition introduced, so certain supplementary conditions disappear and the admissible class of domains of integral quantities broadens. This does not mean that every rational function becomes a polynomial or that denominators can be ignored. The phrase makes a structural observation about definitions, not an unconditional algebraic identification of integral and fractional.\par
\noindent{\footnotesize Izvorne zapisy / Original records: HIST-g068-014.}\par
\smallskip\noindent\emph{Sravnjenje s tekstom:} Prosta karakteristika měnjaje definicijne uslovja, ne izključaje vse imenovatelje.\par
\noindent\emph{Textual comparison:} Prime characteristic changes the defining conditions; it does not eliminate every denominator.\par
\noindent{\footnotesize TEXT-COMP-852; pasaže / passages: 8382, 8383, 8384, 8385, 8386, 8387, 8388, 8389, 8390, 8391, 8392, 8393, 8394, 8395, 8396; \texttt{CORPUS\_FORM\_COMPARISONS.json}.}\par
\end{minipage}\par
\par\medskip\Needspace{7\baselineskip}
\noindent\begin{minipage}{\linewidth}
\hypertarget{TERM-1802}{}
\noindent\textbf{Polna fraza vměsto skračenja}\par\noindent\textbf{Writing the phrase out instead of abbreviating}\quad{\footnotesize\texttt{TERM-1802}}\par
\smallskip\noindent\emph{Razlog:} I podobne bylo vybrano kako polna fraza vměsto skračenja i t. d., da by avtomatično prěpisyvanje v kirilicu ne ostavilo směšany zapis. To jest tehnično motivovana pisemna konvencija, ne dokaz, že skračenje jest jezyčno nepravilno. V trěh pregledanyh aktualnyh pasažah jednake serije form však končat s i t. d. Zato historijno potrěbovanje ne možno opisati kako vsude ostvarjeno. Navedena slova upominajut prodolženje obrazca, ne dozvoljajut izpustiti novo matematično sodržanje.\par
\noindent\emph{Reason:} I podobne was chosen as a written-out phrase instead of i t. d. to avoid mixed-script output during automatic Cyrillic conversion. This is a technically motivated writing convention, not evidence that the abbreviation is linguistically wrong. Yet the three examined current passages end the corresponding families of forms with i t. d. The historical requirement therefore cannot be described as uniformly implemented. The phrase signals continuation of a pattern; it does not authorize omission of new mathematical content.\par
\noindent{\footnotesize Izvorne zapisy / Original records: HIST-g022-008, HIST-g023-011, HIST-g024-013.}\par
\smallskip\noindent\emph{Sravnjenje s tekstom:} Vse tri prověrene pasaže imajo i t. d.; polna fraza nebyla doslědno provedena.\par
\noindent\emph{Textual comparison:} All three examined passages have i t. d.; the written-out phrase was not consistently implemented.\par
\noindent{\footnotesize TEXT-COMP-859; pasaže / passages: 2881, 3002, 3118; \texttt{CORPUS\_FORM\_COMPARISONS.json}.}\par
\end{minipage}\par
\par\medskip\Needspace{7\baselineskip}
\noindent\begin{minipage}{\linewidth}
\hypertarget{TERM-1806}{}
\noindent\textbf{Poděkovanje s konkretnym kreditom}\par\noindent\textbf{Acknowledgment with specific credit}\quad{\footnotesize\texttt{TERM-1806}}\par
\smallskip\noindent\emph{Razlog:} Zahvaljujem Hentzeltu za zaměčenja vyražaje priznanje naučnogo prinosa. Sohranjena motivacija kaže, že konstrukcija zahvaljujem zaměnila grub prěvod glagola dlga. V aktualnom odstavku Noether pripisuje Hentzeltu i konkretne pobudky k obobčenju bazovyh teorem, i daljše menše zaměčenja. To jest glas avtorice v prěvodu, ne izjava prevoditelja. Stilistična obrabotka ne směje izbrisati imenovanogo čelověka ni zameniti konkretne zaslugy obščim poděkovanjem.\par
\noindent\emph{Reason:} Zahvaljujem Hentzeltu za zaměčenja expresses acknowledgment of a scholarly contribution. The preserved motivation says that zahvaljujem replaced a rough calque of owing. In the current paragraph Noether credits Hentzelt both with specific prompts toward generalizing the basis theorems and with further smaller observations. This is the author's voice in translation, not a statement by the translator. Stylistic adjustment must not remove the named person or replace specific contributions with generic thanks.\par
\noindent{\footnotesize Izvorne zapisy / Original records: HIST-g040-020.}\par
\smallskip\noindent\emph{Sravnjenje s tekstom:} Prěvod priznanja ne směje stratiti konkretne zaslugy.\par
\noindent\emph{Textual comparison:} Translation of the acknowledgment must not lose the specific contributions.\par
\noindent{\footnotesize TEXT-COMP-863; pasaže / passages: 5033; \texttt{CORPUS\_FORM\_COMPARISONS.json}.}\par
\end{minipage}\par
\par\medskip\Needspace{7\baselineskip}
\noindent\begin{minipage}{\linewidth}
\hypertarget{TERM-1810}{}
\noindent\textbf{Zaglavje s matematičnoju hipotezoju}\par\noindent\textbf{A heading with a mathematical hypothesis}\quad{\footnotesize\texttt{TERM-1810}}\par
\smallskip\noindent\emph{Razlog:} Ponašanje slěda i diskriminanta pri direktnoj sumě klasov jest zaglavje s matematičnoju hipotezoju. Slěd elementa po odnošenju k klasu stava summoju slědov po odnošenju k komponentnym klasam; diskriminantny ideal stava produktom komponentnyh diskriminantnyh idealov. Nije pravdivo zameniti obě operacije jednim občim skladanjem. Klasy i baza, prirěđena direktnoj sumě, sut čest argumenta. Ponašanje tu znači zakon proměny invariantov, ne neformalnu analogiju.\par
\noindent\emph{Reason:} Ponašanje slěda i diskriminanta pri direktnoj sumě klasov is a heading carrying a mathematical hypothesis. The trace relative to the class becomes the sum of traces relative to component classes; the discriminant ideal becomes the product of the component discriminant ideals. The two operations cannot be replaced by a single vague combination. The classes and the basis adapted to the direct sum belong to the argument. Behaviour here means a law governing the invariants, not an informal analogy.\par
\noindent{\footnotesize Izvorne zapisy / Original records: HIST-g139-001.}\par
\smallskip\noindent\emph{Sravnjenje s tekstom:} Ne jedna i ta sama operacija za oba invarianta.\par
\noindent\emph{Textual comparison:} The two invariants do not obey the same operation.\par
\noindent{\footnotesize TEXT-COMP-867; pasaže / passages: 18181, 18183, 18185, 18193, 18195; \texttt{CORPUS\_FORM\_COMPARISONS.json}.}\par
\end{minipage}\par
\par\medskip\Needspace{7\baselineskip}
\noindent\begin{minipage}{\linewidth}
\hypertarget{TERM-1812}{}
\noindent\textbf{Prěnos bazy prěz izomorfizm}\par\noindent\textbf{Transporting a basis through an isomorphism}\quad{\footnotesize\texttt{TERM-1812}}\par
\smallskip\noindent\emph{Razlog:} Baza prirěđena posrědstvom izomorfizma označaje prěnesenu bazu, ne novu libovoljno vybranu bazu. Elementy $\beta_i$ odgovarjajut $\alpha_i$ pod modulnym izomorfizmom; zato množenje tym samym $\gamma$ ima iste koeficienty v oběh bazah. To daje identične matrične prsteny za te konkretno svezane bazy. Ako bazy samo nezavisno vybrati, iz takoj frazy nelzja vyvesti ravnost pojedinyh matric. Relacionna konstrukcija zato sohranjaje ključny krok dokaza.\par
\noindent\emph{Reason:} Baza prirěđena posrědstvom izomorfizma denotes a transported basis, not a newly chosen arbitrary basis. The $\beta_i$ correspond to the $\alpha_i$ under a module isomorphism, so multiplication by the same $\gamma$ has identical coefficients in both bases. This yields identical matrix rings for those specifically related bases. Independent choices of bases would not justify equality of the individual matrices. The relational phrase therefore preserves an essential step of the proof.\par
\noindent{\footnotesize Izvorne zapisy / Original records: HIST-g134-003.}\par
\smallskip\noindent\emph{Sravnjenje s tekstom:} Istovetne koeficienty proizhodjat iz modulnogo izomorfizma.\par
\noindent\emph{Textual comparison:} Identical coefficients result from the module isomorphism.\par
\noindent{\footnotesize TEXT-COMP-869; pasaže / passages: 17985; \texttt{CORPUS\_FORM\_COMPARISONS.json}.}\par
\end{minipage}\par
\par\medskip\Needspace{7\baselineskip}
\noindent\begin{minipage}{\linewidth}
\hypertarget{TERM-1825}{}
\noindent\textbf{Odkazka na konec sekcije}\par\noindent\textbf{A reference to the end of a section}\quad{\footnotesize\texttt{TERM-1825}}\par
\smallskip\noindent\emph{Razlog:} Srav. konec § 5 jest kratka odkazna fraza: ona posylaje čitatelja na konec odrěđenej sekcije, ne na vse jeje sodržanje. V aktualnom zapisu fraza stoji v tekstovom vstavjenju pri formuli. Bukvalny znak § čuva identitet odkazki pri prěpisu prozy v kirilicu; broj 5 ostaje neproměnjeny. To jest konvencija odkazovanja i pisanja, ne novy matematičny termin. Samo sohranjenje odkazki ne jest prověrka vsih povjazanyh formul.\par
\noindent\emph{Reason:} Srav. konec § 5 is a compact cross-reference directing the reader to the end of a specified section, not indiscriminately to everything in it. In the current source it appears in a text insertion beside the formula. The literal section sign preserves the reference when prose is converted to Cyrillic; the number 5 is unchanged. This is a referencing and writing convention, not a new mathematical term. Preserving the reference is not in itself verification of all the formulas it connects.\par
\noindent{\footnotesize Izvorne zapisy / Original records: HIST-g011-007.}\par
\smallskip\noindent\emph{Sravnjenje s tekstom:} Odkazka ukazuje na konec sekcije 5.\par
\noindent\emph{Textual comparison:} The reference specifies the end of section 5.\par
\noindent{\footnotesize TEXT-COMP-882; pasaže / passages: 1340, 1341, 1342, 1343, 1344; \texttt{CORPUS\_FORM\_COMPARISONS.json}.}\par
\end{minipage}\par
\par\medskip\Needspace{7\baselineskip}
\noindent\begin{minipage}{\linewidth}
\hypertarget{TERM-1826}{}
\noindent\textbf{Soglasovanje pri simboličnyh subjektah}\par\noindent\textbf{Agreement with symbolic subjects}\quad{\footnotesize\texttt{TERM-1826}}\par
\smallskip\noindent\emph{Razlog:} Vhodi i vhodet razděljajut jedinstveny i množstveny subjekt v toj konkretnoj rečenici. $H^2L$ samo ne vhodi v kontrakciju s $K$; dva imenovane izraza, $H^3$ i $H^2L$, ne vhodet s $\theta$. Zato prisutnost dvuh glagolnyh form ne jest sama po sobě neporednost. Historijny zapis ostavil skladnju pri simboličnyh subjektah otvorenoju dlja jezyčnoj ocěny. Slovnik čuva to ograničenje; razgledanje jednoj rečenice ne ustanavjaje obče pravilo dlja vseh formulnih subjektov.\par
\noindent\emph{Reason:} Vhodi and vhodet distinguish a singular from a plural subject in this particular sentence. $H^2L$ alone does not enter a contraction with $K$; the two named expressions $H^3$ and $H^2L$ do not enter one with $\theta$. The presence of the two verb forms is therefore not by itself inconsistency. The historical record left agreement with symbolic subjects open to linguistic assessment. The dictionary preserves that qualification; one examined sentence does not establish a general rule for every formula used as a subject.\par
\noindent{\footnotesize Izvorne zapisy / Original records: HIST-g016-013.}\par
\smallskip\noindent\emph{Sravnjenje s tekstom:} Rečenica razděljaje jedinstveny i množstveny subjekt.\par
\noindent\emph{Textual comparison:} The sentence distinguishes singular and plural subjects.\par
\noindent{\footnotesize TEXT-COMP-883; pasaže / passages: 2271; \texttt{CORPUS\_FORM\_COMPARISONS.json}.}\par
\end{minipage}\par
\par\medskip\Needspace{7\baselineskip}
\noindent\begin{minipage}{\linewidth}
\hypertarget{TERM-1855}{}
\noindent\textbf{Dvojno izključenje v hipoteze}\par\noindent\textbf{A double exclusion in a hypothesis}\quad{\footnotesize\texttt{TERM-1855}}\par
\smallskip\noindent\emph{Razlog:} Različny od nulovogo i jediničnogo ideala ostavlja oba krajne ideala izvan razgledanja. V konkretnoj posledici prsten ima jedin prosty ideal medžu ostalymi idealami; odtud dobiva se prsten glavnyh idealov. Izpustjenje jediničnogo bi izměnilo obseg jedinstvenosti. To jest jasno napisana hipoteza, ne nova nazva klasa ili neobvezna slovna dodatka. Arhivny izbor naměrno povtarja formulaciju predhodnyh aksiomov, aby zavisimost byla vidliva.\par
\noindent\emph{Reason:} Različny od nulovogo i jediničnogo ideala excludes both extreme ideals. In the particular consequence the ring has one prime ideal among the remaining ideals, from which the principal-ideal conclusion is obtained. Dropping unit would alter the scope of uniqueness. This is an explicitly written hypothesis, not a new class name or optional verbal addition. The recorded choice deliberately repeats the preceding axiomatic wording to make the dependence visible.\par
\noindent{\footnotesize Izvorne zapisy / Original records: HIST-g156-002.}\par
\smallskip\noindent\emph{Sravnjenje s tekstom:} Jediničny i nulovy ideal ne vhodet v obseg tvrdženja.\par
\noindent\emph{Textual comparison:} The unit and zero ideals lie outside the statement's scope.\par
\noindent{\footnotesize TEXT-COMP-912; pasaže / passages: 18823; \texttt{CORPUS\_FORM\_COMPARISONS.json}.}\par
\end{minipage}\par
\par\medskip\Needspace{7\baselineskip}
\noindent\begin{minipage}{\linewidth}
\hypertarget{TERM-1856}{}
\noindent\textbf{Nulovy i jediničny ideal pri množenju}\par\noindent\textbf{Zero and unit ideals in multiplication}\quad{\footnotesize\texttt{TERM-1856}}\par
\smallskip\noindent\emph{Razlog:} Idealy različne od nulovogo i jediničnogo idealov čuva množstveno čitanje i obě izključenja v tvrđenju o idealnom produktu. Predhodna lema iz jednakosti $\mathfrak a=\mathfrak a\mathfrak b$ i nenulovosti $\mathfrak a$ vyvodi, že $\mathfrak b$ jest jediničny ideal; poslědujuča neravnost zato potrěbuje te granice. Frazu ne možno skratiti prosto do nenulove idealy, ibo by dopuštjala jediničny faktor. Ruskaja i ukrajinskaja sravniteljnaja forma čuvajut isto dvojno izključenje.\par
\noindent\emph{Reason:} Idealy različne od nulovogo i jediničnogo idealov preserves the plural reading and both exclusions in the statement about ideal products. The preceding lemma deduces that $\mathfrak b$ is the unit ideal from $\mathfrak a=\mathfrak a\mathfrak b$ and nonzero $\mathfrak a$; the following inequality therefore needs those boundaries. The phrase cannot be shortened simply to nonzero ideals, which would allow the unit factor. The Russian and Ukrainian comparanda preserve the same double exclusion.\par
\noindent{\footnotesize Izvorne zapisy / Original records: HIST-g115-007.}\par
\smallskip\noindent\emph{Sravnjenje s tekstom:} Izključenje jediničnogo faktora jest sučstveno.\par
\noindent\emph{Textual comparison:} Exclusion of the unit factor is essential.\par
\noindent{\footnotesize TEXT-COMP-913; pasaže / passages: 16849, 16851, 16853; \texttt{CORPUS\_FORM\_COMPARISONS.json}.}\par
\end{minipage}\par
\par\medskip\Needspace{7\baselineskip}
\noindent\begin{minipage}{\linewidth}
\hypertarget{TERM-1857}{}
\noindent\textbf{Nenulovost diskriminantnogo ideala}\par\noindent\textbf{Nonvanishing of the discriminant ideal}\quad{\footnotesize\texttt{TERM-1857}}\par
\smallskip\noindent\emph{Razlog:} Različny od nulovogo ideala na koncu definicije govori o samom diskriminantnom idealu. On jest porodženy determinantami slědov, a nenulovost pri linearno nezavisnom izboru elementov daje nenulovost togo ideala. Slabše tvrdženje, že někoja matrica može byti nenulova, ne prědaje togo zaključka. Arhivna motivacija govori o prověrke skana; zde nově prověrena jest aktualna tekstova pasaža, ne povtorjena nezavisna prověrka skana.\par
\noindent\emph{Reason:} Različny od nulovogo ideala at the end of the definition concerns the discriminant ideal itself. It is generated by trace determinants, and nonvanishing for a linearly independent choice of elements establishes that the ideal is nonzero. A weaker statement that some matrix might be nonzero would not convey that conclusion. The historical rationale reports a scan check; the new examination here is of the current textual passage, not a repeated independent inspection of the scan.\par
\noindent{\footnotesize Izvorne zapisy / Original records: HIST-g158-005.}\par
\smallskip\noindent\emph{Sravnjenje s tekstom:} Zaključok kasaje se ideala, ne samo matrice.\par
\noindent\emph{Textual comparison:} The conclusion concerns the ideal, not merely a matrix.\par
\noindent{\footnotesize TEXT-COMP-914; pasaže / passages: 18895, 18897; \texttt{CORPUS\_FORM\_COMPARISONS.json}.}\par
\end{minipage}\par
\par\medskip\Needspace{7\baselineskip}
\noindent\begin{minipage}{\linewidth}
\hypertarget{TERM-1858}{}
\noindent\textbf{Izrična hipoteza o nenulovom idealu}\par\noindent\textbf{An explicit nonzero-ideal hypothesis}\quad{\footnotesize\texttt{TERM-1858}}\par
\smallskip\noindent\emph{Razlog:} Ideal razny od nulovogo ideala jest izrična forma korotšego nenulovy ideal. Byla izbrana, aby kontrast nulovogo i jediničnogo ideala ostal čitlivy pri prěhode do prstena ostatkov. Pasaža prědstavjaje takovy ideal kako produkt konečno mnogyh vzajemno prostyh primarnyh idealov. Ona odrěđuje ješče svojstvo jih pridruženih prostyh idealov; ta čest ne dolžna propasti pri skračenju frazy. Razny i različny sut varianty jezyčnogo vyraženja, ne dva različne matematične uslovja.\par
\noindent\emph{Reason:} Ideal razny od nulovogo ideala is an explicit form of the shorter nonzero ideal. It was chosen to keep the zero–unit contrast legible when passing to residue rings. The passage expresses such an ideal as a product of finitely many relatively prime primary ideals. It also specifies a property of their associated prime ideals; that part must survive any shortening of the phrase. Razny and različny are variants of expression, not two different mathematical conditions.\par
\noindent{\footnotesize Izvorne zapisy / Original records: HIST-g150-002.}\par
\smallskip\noindent\emph{Sravnjenje s tekstom:} Obě česti formulacije noset matematične uslovja.\par
\noindent\emph{Textual comparison:} Both parts of the formulation carry mathematical conditions.\par
\noindent{\footnotesize TEXT-COMP-915; pasaže / passages: 18593, 18627; \texttt{CORPUS\_FORM\_COMPARISONS.json}.}\par
\end{minipage}\par
\par\medskip\Needspace{7\baselineskip}
\noindent\begin{minipage}{\linewidth}
\hypertarget{TERM-1859}{}
\noindent\textbf{Rang jedan kako predikat}\par\noindent\textbf{Rank one as a predicate}\quad{\footnotesize\texttt{TERM-1859}}\par
\smallskip\noindent\emph{Razlog:} Imaje rang jedan prěvodi staršu imennu konstrukciju prostym predikatom. Rang jest relativny k ukazanomu podpolju $P$; pri danoj hipoteze algebraičnogo zatvorjenja i prostosti nulovogo ideala tekst zaključuje, že prsten jest identičny s $P$. Jediničny element togda služi kako modulna baza. To ne jest tvrdženje, že libovoljny prsten ima jedin element. Zapis dava prednost obyčnoj sovrěmennoj formulaciji, ale široka razumělivost ne byla změrjena čitatelskym testom.\par
\noindent\emph{Reason:} Imaje rang jedan renders an older nominal construction as a direct predicate. Rank is relative to the specified subfield $P$; under the algebraic-closure and zero-ideal primality hypotheses, the text concludes that the ring is identical with $P$. The identity element then serves as a module basis. This does not say that an arbitrary ring has only one element. The record prefers ordinary contemporary phrasing, but broad intelligibility was not measured by a reader test.\par
\noindent{\footnotesize Izvorne zapisy / Original records: HIST-g142-004.}\par
\smallskip\noindent\emph{Sravnjenje s tekstom:} Jedinica tvori modulnu bazu pri ukazanyh hipotezah.\par
\noindent\emph{Textual comparison:} The identity forms a module basis under the stated hypotheses.\par
\noindent{\footnotesize TEXT-COMP-916; pasaže / passages: 18350, 18352, 18354; \texttt{CORPUS\_FORM\_COMPARISONS.json}.}\par
\end{minipage}\par
\par\medskip\Needspace{7\baselineskip}
\noindent\begin{minipage}{\linewidth}
\hypertarget{TERM-1862}{}
\noindent\textbf{Obseg nepisanogo modula}\par\noindent\textbf{The scope of an implicit modulus}\quad{\footnotesize\texttt{TERM-1862}}\par
\smallskip\noindent\emph{Razlog:} Od nyně vse formuly trěba razuměti mod označuje obseg dalšego čitanja, ne novu matematičnu operaciju. Aktualna rečenica najprvo dava modul $(\varrho,t)$, a od formuly (27) širši modul, vključajuči i, J i vyše formy. Kratko zaglavje slovnika ne dolžno izbrisati tu drugu čest. Izrična konstrukcija trěba razuměti bila izbrana zaradi jasnosti; zapis čuva otvoreno pytanje o modalnosti i porjadku slov. Simboly samogo modula ne prevodet se pri změně pisma.\par
\noindent\emph{Reason:} Od nyně vse formuly trěba razuměti mod specifies the scope of subsequent reading, not a new mathematical operation. The current sentence first gives modulus $(\varrho,t)$, then from formula (27) a larger modulus including i, J and higher forms. The short dictionary heading must not erase that second part. The explicit trěba razuměti construction was chosen for clarity; the record preserves an open question about modality and word order. The modulus symbols themselves are not translated when the script changes.\par
\noindent{\footnotesize Izvorne zapisy / Original records: HIST-g018-011.}\par
\smallskip\noindent\emph{Sravnjenje s tekstom:} Od formuly (27) modul jest širši.\par
\noindent\emph{Textual comparison:} The modulus becomes larger from formula (27).\par
\noindent{\footnotesize TEXT-COMP-919; pasaže / passages: 2485; \texttt{CORPUS\_FORM\_COMPARISONS.json}.}\par
\end{minipage}\par
\par\medskip\Needspace{7\baselineskip}
\noindent\begin{minipage}{\linewidth}
\hypertarget{TERM-1863}{}
\noindent\textbf{Konečny i nenulovy skalar}\par\noindent\textbf{A finite nonzero scalar}\quad{\footnotesize\texttt{TERM-1863}}\par
\smallskip\noindent\emph{Razlog:} Od nuly različny i konečny jest dvojno ograničenje realnogo skalara $c$. Ono stojit pri prěpisu relacije medžu obrazami i početkovymi veličinami; obe vlastnosti trěba sohraniti. Konečny ne kaže, že množina možlivyh $c$ jest konečna, a nenulovy ne znači pozitivny. Prosta spojka i činji hipotezy čitlivymi bez preměny jih matematičnogo obsega. Tuto jest formulacija uslovja, ne slovnikov novy tehničny imennik.\par
\noindent\emph{Reason:} Od nuly različny i konečny places two restrictions on the real scalar $c$. They occur when rewriting a relation between images and original quantities; both must remain. Finite does not say that the set of possible $c$ is finite, and nonzero does not mean positive. The simple conjunction makes the hypotheses legible without changing their mathematical scope. This is the wording of a condition, not a newly coined technical noun.\par
\noindent{\footnotesize Izvorne zapisy / Original records: HIST-g072-012.}\par
\smallskip\noindent\emph{Sravnjenje s tekstom:} Konečnost skalara ne jest konečnost množiny izborov.\par
\noindent\emph{Textual comparison:} Finiteness of the scalar is not finiteness of the set of choices.\par
\noindent{\footnotesize TEXT-COMP-920; pasaže / passages: 8678, 8680, 8682, 8684, 8686; \texttt{CORPUS\_FORM\_COMPARISONS.json}.}\par
\end{minipage}\par
\par\medskip\Needspace{7\baselineskip}
\noindent\begin{minipage}{\linewidth}
\hypertarget{TERM-1864}{}
\noindent\textbf{Hipoteza pri nenulovosti determinanta}\par\noindent\textbf{The hypothesis for determinant nonvanishing}\quad{\footnotesize\texttt{TERM-1864}}\par
\smallskip\noindent\emph{Razlog:} Nenulovy, skoro tau sut linearno nezavisne pripojuje nenulovost determinanta k izričnoj hipoteze. Skoro tu znači kogda uslovje izpolnjeno, ne tvrdženje o skorom časovom nastupanju. Aktualny tekst nazyvaje elementy $\tau_1,\ldots,\tau_n$ i potom vyvodi nenulovost diskriminantnogo ideala. Samo n-spis elementov bez nezavisnosti ne dava togo zaključka. Zapis upominaje ranějšu skanovu prověrku; novym svědectvom tu jest pročitana aktualna pasaža, ne povtoreny pregled skana.\par
\noindent\emph{Reason:} Nenulovy, skoro tau sut linearno nezavisne attaches nonvanishing of the determinant to an explicit hypothesis. Skoro here means when the condition is satisfied, not that something will happen soon in time. The current text names the elements $\tau_1,\ldots,\tau_n$ and then deduces nonvanishing of the discriminant ideal. An n-tuple without independence does not supply that conclusion. The record mentions an earlier scan check; the new evidence here is the examined current passage, not another inspection of the scan.\par
\noindent{\footnotesize Izvorne zapisy / Original records: HIST-g158-004.}\par
\smallskip\noindent\emph{Sravnjenje s tekstom:} Hipoteza odnosi se k izbranym elementam.\par
\noindent\emph{Textual comparison:} The hypothesis concerns the selected elements.\par
\noindent{\footnotesize TEXT-COMP-921; pasaže / passages: 18895; \texttt{CORPUS\_FORM\_COMPARISONS.json}.}\par
\end{minipage}\par
\par\medskip\Needspace{7\baselineskip}
\noindent\begin{minipage}{\linewidth}
\hypertarget{TERM-1865}{}
\noindent\textbf{Nenulovy element s nulovoju stepenju}\par\noindent\textbf{A nonzero element with a zero power}\quad{\footnotesize\texttt{TERM-1865}}\par
\smallskip\noindent\emph{Razlog:} Nenulovy element i někoja stepenj izčezaje opisujut nilpotentny svědok, ne protivrěčje. Element sam ne jest nula, ale někoja jego pozitivna stepenj jest nula. V dokazě element $G(\gamma)$ nastavaje po razširjenju koeficientov; jego nenulovost vyhodi iz linearnoj nezavisnosti nižših stepenjev, a p-ta stepenj izčezaje po definicijnom uravnenju. Glagol izčezaje zato znači algebraičnu ravnost nuli, ne postupno umenšenje veličiny.\par
\noindent\emph{Reason:} Nenulovy element and někoja stepenj izčezaje describe a nilpotent witness, not a contradiction. The element itself is nonzero but some positive power is zero. In the proof $G(\gamma)$ is constructed after coefficient extension; its nonvanishing follows from linear independence of lower powers, while its pth power vanishes by the defining equation. Izčezaje therefore means algebraic equality to zero, not gradual decrease in magnitude.\par
\noindent{\footnotesize Izvorne zapisy / Original records: HIST-g128-003.}\par
\smallskip\noindent\emph{Sravnjenje s tekstom:} Nenulovost elementa i nulovost jeho stepenja imajut razne dokazy.\par
\noindent\emph{Textual comparison:} Nonvanishing of the element and vanishing of its power have different justifications.\par
\noindent{\footnotesize TEXT-COMP-922; pasaže / passages: 17735, 17737, 17739, 17741, 17743, 17745, 17747, 17748, 17750; \texttt{CORPUS\_FORM\_COMPARISONS.json}.}\par
\end{minipage}\par
\par\medskip\Needspace{7\baselineskip}
\noindent\begin{minipage}{\linewidth}
\hypertarget{TERM-1866}{}
\noindent\textbf{Prědpoloženje v obratnom dokazu}\par\noindent\textbf{An assumption in a converse proof}\quad{\footnotesize\texttt{TERM-1866}}\par
\smallskip\noindent\emph{Razlog:} Predpostavjeno razloženje čuva logičny směr obrata: razloženje nyně jest hipoteza, iz ktorej dolžne slědovati aksiomy. Aktualny tekst piše prědpoloženo razloženje i izrično ogrančuje jego formu. Egzistencija jediničnogo elementa tam ne jest prědpoložena; ona ne smi byti neprimětno dodana slovnikom. Změna pridavnika ne měnja dokaznu rolu, ale skračenje do prosto razloženje by skrylo, čto jest dano i čto trěba vyvesti.\par
\noindent\emph{Reason:} Predpostavjeno razloženje preserves the direction of the converse: decomposition is now an assumption from which the axioms are to follow. The current text writes prědpoloženo razloženje and explicitly restricts its form. Existence of an identity element is not assumed there; the dictionary must not silently add it. Changing the adjective does not change its logical role, but shortening it simply to decomposition would hide what is given and what must be derived.\par
\noindent{\footnotesize Izvorne zapisy / Original records: HIST-g116-007.}\par
\smallskip\noindent\emph{Sravnjenje s tekstom:} Jediničny element ne jest početkova hipoteza.\par
\noindent\emph{Textual comparison:} An identity element is not an initial hypothesis.\par
\noindent{\footnotesize TEXT-COMP-923; pasaže / passages: 17013, 17015, 17017, 17019, 17021, 17023; \texttt{CORPUS\_FORM\_COMPARISONS.json}.}\par
\end{minipage}\par
\par\medskip\Needspace{7\baselineskip}
\noindent\begin{minipage}{\linewidth}
\hypertarget{TERM-1867}{}
\noindent\textbf{Kredit osobě, ktora predložila rabotu}\par\noindent\textbf{Credit to the person presenting the paper}\quad{\footnotesize\texttt{TERM-1867}}\par
\smallskip\noindent\emph{Razlog:} Predloženo g. Šurom označuje osobu, ktora predložila rabotu akademiji. To bibliografičny kredit, ne pripisanje jej avtorstva Šuru: zaglavje odrěđeno nazyvaje Noether i Brauera. Inštrumentalna forma Šurom i latinsko Š izbrany sut, aby kiriličny prěpis dal Шуром. To motivacija pisma; ona ne zaměnjaje originalno ime Schur v bibliografičnoj identifikaciji. Kratka formula dolžna sohraniti i rolu, i konkretnu osobu.\par
\noindent\emph{Reason:} Predloženo g. Šurom identifies the person who presented the paper to the academy. It is bibliographic credit, not attribution of authorship to Schur: the heading separately identifies Noether and Brauer. Instrumental Šurom and Latin Š were chosen so that Cyrillic conversion yields Шуром. That is a script-related motivation; it does not replace the original name Schur in bibliographic identification. The compact formula must retain both the role and the particular person.\par
\noindent{\footnotesize Izvorne zapisy / Original records: HIST-g166-004.}\par
\smallskip\noindent\emph{Sravnjenje s tekstom:} Predloženo ne znači napisano.\par
\noindent\emph{Textual comparison:} Presented does not mean written.\par
\noindent{\footnotesize TEXT-COMP-924; pasaže / passages: 19171, 19173, 19176; \texttt{CORPUS\_FORM\_COMPARISONS.json}.}\par
\end{minipage}\par
\par\medskip\Needspace{7\baselineskip}
\noindent\begin{minipage}{\linewidth}
\hypertarget{TERM-1868}{}
\noindent\textbf{Dopustimost lokalnogo izbora bazy}\par\noindent\textbf{Possibility of choosing a local basis}\quad{\footnotesize\texttt{TERM-1868}}\par
\smallskip\noindent\emph{Razlog:} Što jest vsegda možno čuva tvrdženje o dopustimosti izbora, ne obeščanje proizvoljnoj konstrukcije. V pasaži možno izbrati elementy iz početkovogo poredka $\mathfrak T$ kako modulnu bazu lokalizovanogo poredka. Aktualna fraza kako jest vsegda možno ima tu istu rolu. Ona ne govori, že početkovy modul ima globalnu svobodnu bazu nad početkovym prstenom. Arhivna motivacija upominaje obnovu po skanu; naš aktualny pregled podtvrđaje tekstovu formulaciju, ne samostojnu skanovu prověrku.\par
\noindent\emph{Reason:} Što jest vsegda možno preserves an assertion that a particular choice is possible, not a promise of arbitrary construction. The passage permits elements of the original order $\mathfrak T$ to be chosen as a module basis of the localized order. Current kako jest vsegda možno plays the same role. It does not assert that the original module has a global free basis over the original ring. The historical rationale mentions recovery from a scan; the present examination confirms the textual formulation, not an independent scan check.\par
\noindent{\footnotesize Izvorne zapisy / Original records: HIST-g162-003.}\par
\smallskip\noindent\emph{Sravnjenje s tekstom:} Lokalna baza ne ustanavjaje globalnu svobodnu bazu.\par
\noindent\emph{Textual comparison:} A local basis does not establish a global free basis.\par
\noindent{\footnotesize TEXT-COMP-925; pasaže / passages: 19037; \texttt{CORPUS\_FORM\_COMPARISONS.json}.}\par
\end{minipage}\par
\par\medskip\Needspace{7\baselineskip}
\noindent\begin{minipage}{\linewidth}
\hypertarget{TERM-1869}{}
\noindent\textbf{Adaptacija imen Weber i Steinitz}\par\noindent\textbf{Adaptation of the names Weber and Steinitz}\quad{\footnotesize\texttt{TERM-1869}}\par
\smallskip\noindent\emph{Razlog:} Veber i E. Štejnic sut adaptovane pisanja vlastnyh imen, ne novy matematičny leksikon. Historijny izbor měl tehničnu pričinu: neprěvedene Weber moglo dati směšane pismo pri avtomatičnoj kirilizaciji. Aktualna pasaža doista ima Veberu i E. Štejnicu pri připisaniu abstraktnoj definicije polja. Pisanje dolžno sohraniti prepoznavanje Webera i Steinitza i njihove kredity. Kontrola pisma sama ne ustanavjaje, že adaptacija jest jedina jezyčno dopustima forma.\par
\noindent\emph{Reason:} Veber and E. Štejnic are adapted spellings of proper names, not new mathematical vocabulary. The recorded choice had a technical reason: unadapted Weber could produce mixed script during automatic Cyrillic conversion. The current passage does have Veberu and E. Štejnicu when crediting the abstract definition of a field. The spelling must preserve identification of Weber and Steinitz and their scholarly credit. Script checking alone does not establish the adaptation as the only linguistically permissible form.\par
\noindent{\footnotesize Izvorne zapisy / Original records: HIST-g068-015.}\par
\smallskip\noindent\emph{Sravnjenje s tekstom:} Adaptovano pisanje ne odstranjaje ljudsky kredit.\par
\noindent\emph{Textual comparison:} Adapted spelling does not remove human credit.\par
\noindent{\footnotesize TEXT-COMP-926; pasaže / passages: 8346, 8347, 8348, 8349, 8350, 8351; \texttt{CORPUS\_FORM\_COMPARISONS.json}.}\par
\end{minipage}\par
\par\medskip\Needspace{7\baselineskip}
\noindent\begin{minipage}{\linewidth}
\hypertarget{TERM-1870}{}
\noindent\textbf{Izomorfizm kako pričina zaključka}\par\noindent\textbf{Isomorphism as the reason for an inference}\quad{\footnotesize\texttt{TERM-1870}}\par
\smallskip\noindent\emph{Razlog:} Zaradi izomorfizma kolc klasov ostatkov jest pričinna svez, ne samo pomenanje srodnyh objektov. Aktualna forma zaměnjuje kolc slovom prstenov. Pasaža upotrěbjaje izomorfizm dlja prenosa svojstva multiplikacijnogo prstena na lokalizovany prsten, potom priměnjaje ranějši rezultat o glavnyh idealah. Izpustiti zaradi značilo by oslabiti vidlivu osnovu zaključka. Samo izomorfizm kvocientov ne nazyva početkovy i lokalizovany prsten bukvalno jednakymi.\par
\noindent\emph{Reason:} Zaradi izomorfizma kolc klasov ostatkov is a causal link, not merely a mention of related objects. The current form replaces kolc with prstenov. The passage uses the isomorphism to transfer the multiplication-ring property to the localized ring, then applies the earlier result concerning principal ideals. Dropping zaradi would obscure the stated basis of the inference. Isomorphism of the residue rings does not call the original and localized rings literally identical.\par
\noindent{\footnotesize Izvorne zapisy / Original records: HIST-g160-003.}\par
\smallskip\noindent\emph{Sravnjenje s tekstom:} Pričinna svez ostaje čest dokaza.\par
\noindent\emph{Textual comparison:} The causal connection remains part of the proof.\par
\noindent{\footnotesize TEXT-COMP-927; pasaže / passages: 18970; \texttt{CORPUS\_FORM\_COMPARISONS.json}.}\par
\end{minipage}\par
\par\medskip\Needspace{7\baselineskip}
\noindent\begin{minipage}{\linewidth}
\hypertarget{TERM-1871}{}
\noindent\textbf{Prěhod k reducibilnosti i soglasovanje}\par\noindent\textbf{Becoming reducible and agreement}\quad{\footnotesize\texttt{TERM-1871}}\par
\smallskip\noindent\emph{Razlog:} Stavajut reducibilne kaže, že razgledane polne transvekcije postavajut reducibilne v danoj konstrukciji. Glagol čuva prěhod, a množstvena forma odnosi se k transvekcijam, ne k jedinom funkcionalnomu determinantu v toj samoj rečenici. Internacionalna osnova reducibil- nastavlja obšču terminologiju. Arhivny zapis ostavlja morfologiju i soglasovanje otvorenymi dlja jezyčnoj ocěny; matematicky kontekst sam ne jest dokazom polnoj gramatičnoj normativnosti.\par
\noindent\emph{Reason:} Stavajut reducibilne says that the full transvections under discussion become reducible in the specified construction. The verb retains the transition, and the plural refers to the transvections, not to the singular functional determinant in the same sentence. The international reducibil- stem continues the terminology. The historical note leaves morphology and agreement open to linguistic assessment; mathematical context alone is not proof of complete grammatical normativity.\par
\noindent{\footnotesize Izvorne zapisy / Original records: HIST-g020-012.}\par
\smallskip\noindent\emph{Sravnjenje s tekstom:} Soglasovanje prati transvekcije, ne susědny jedinstveny imenik.\par
\noindent\emph{Textual comparison:} Agreement follows the transvections, not the neighbouring singular noun.\par
\noindent{\footnotesize TEXT-COMP-928; pasaže / passages: 2680; \texttt{CORPUS\_FORM\_COMPARISONS.json}.}\par
\end{minipage}\par
\par\medskip\Needspace{7\baselineskip}
\noindent\begin{minipage}{\linewidth}
\hypertarget{TERM-1877}{}
\noindent\textbf{Odkazka na razširjenje istogo ranga}\par\noindent\textbf{Reference to a same-rank extension}\quad{\footnotesize\texttt{TERM-1877}}\par
\smallskip\noindent\emph{Razlog:} Definirano, kako v §1, 2, kako kolco togo samogo ranga čuva i odkazku, i točny tip razširjenja. Aktualny tekst piše prsten i osobno vyžaduje otsutnost děliteljev nuly i zadany prěsěk. Sledujuči argument o diskriminante koristi sohranjenje bazy. Zato fraza ne označuje libovoljno razširjenje koeficientov bez tyh prědpoloženj. Ranějša motivacija o oprave tiskanoj odkazki ostaje historijnym zapisom; zde prověřeny jest aktualny tekst i jego dokazna rola.\par
\noindent\emph{Reason:} Definirano, kako v §1, 2, kako kolco togo samogo ranga preserves both the reference and the exact type of extension. The current text writes prsten and separately requires absence of zero divisors and a specified intersection. The subsequent discriminant argument uses preservation of the basis. The phrase therefore does not denote an arbitrary coefficient extension without those assumptions. The earlier rationale concerning restoration of the printed reference remains a historical record; the present check concerns the current text and its role in the proof.\par
\noindent{\footnotesize Izvorne zapisy / Original records: HIST-g140-006.}\par
\smallskip\noindent\emph{Sravnjenje s tekstom:} Odkazka svjazuje formulaciju s opredělenoju konstrukcijeju.\par
\noindent\emph{Textual comparison:} The reference ties the formulation to a specified construction.\par
\noindent{\footnotesize TEXT-COMP-934; pasaže / passages: 18244, 18246; \texttt{CORPUS\_FORM\_COMPARISONS.json}.}\par
\end{minipage}\par
\par\medskip\Needspace{7\baselineskip}
\noindent\begin{minipage}{\linewidth}
\hypertarget{TERM-1878}{}
\noindent\textbf{Matematična libovoljnost}\par\noindent\textbf{Mathematical arbitrariness}\quad{\footnotesize\texttt{TERM-1878}}\par
\smallskip\noindent\emph{Razlog:} Libovolne veličiny, nyně libovoljne veličiny, označuje matematičnu libovoljnost $p$ i $q$, ne nahodne ili kapriciozne izbory. Rečenica dava čitatelju pokyn, kako razuměti veličiny v identičnosti; ona ne kaže, že možno bez uslovja izměniti vse simbole formula. Shranjene ruske i ukrajinske formy imajut tože kvantifikacijno značenje. Korotka fraza ima smysl samo s ukazanymi argumentami, ne kako neograničeno pravilo svobody.\par
\noindent\emph{Reason:} Libovolne veličiny, now libovoljne veličiny, expresses mathematical arbitrariness of $p$ and $q$, not random or capricious choices. The sentence tells the reader how to treat the quantities in the identity; it does not say that all symbols in the formula can be altered without restriction. The preserved Russian and Ukrainian forms have the same quantificational meaning. The short phrase is meaningful with its specified arguments, not as an unrestricted rule of freedom.\par
\noindent{\footnotesize Izvorne zapisy / Original records: HIST-g031-011.}\par
\smallskip\noindent\emph{Sravnjenje s tekstom:} Libovoljnost ima ukazanyh nositeljev.\par
\noindent\emph{Textual comparison:} Arbitrariness has specified arguments.\par
\noindent{\footnotesize TEXT-COMP-935; pasaže / passages: 3865; \texttt{CORPUS\_FORM\_COMPARISONS.json}.}\par
\end{minipage}\par
\par\medskip\Needspace{7\baselineskip}
\noindent\begin{minipage}{\linewidth}
\hypertarget{TERM-1884}{}
\noindent\textbf{Črta v oznakě dobrogo poredka}\par\noindent\textbf{The prime in the notation for a well-order}\quad{\footnotesize\texttt{TERM-1884}}\par
\smallskip\noindent\emph{Razlog:} Znak $\Omega'$ označuje drugo dobro uporjadkovanje, ne proizvodnu od prvogo. Rěči Omega prime v historijnom zapisu sut opis simbola, ne čest medžuslovjanskogo termina. Aktualny dokaz izbira prvy ješče neprirěđeny element po tomu drugomu poredku; tím nastaje vzajemno jednoznačno prirěđenje medžu bazami. Zachovanje črty jest zato potrěbno dlja razlikovanja dvoh izborov. Jezyk okolnoj frazy može se změniti, ale matematična oznaka ostaje ista.\par
\noindent\emph{Reason:} The symbol $\Omega'$ denotes a second well-order, not a derivative of the first. Omega prime in the historical record is a description of the symbol, not part of the Interslavic term. The current proof selects the first still-unassigned element in that second order, thereby constructing a bijective assignment between the bases. Retaining the prime is necessary to distinguish the two choices. The language surrounding it may change, but the mathematical identifier remains the same.\par
\noindent{\footnotesize Izvorne zapisy / Original records: HIST-g071-003.}\par
\smallskip\noindent\emph{Sravnjenje s tekstom:} Črta razlikuje dva poredky; ne označuje derivaciju.\par
\noindent\emph{Textual comparison:} The prime distinguishes two orders; it does not indicate differentiation.\par
\noindent{\footnotesize TEXT-COMP-941; pasaže / passages: 8518, 8522, 8523, 8524, 8525; \texttt{CORPUS\_FORM\_COMPARISONS.json}.}\par
\end{minipage}\par
\par\medskip\Needspace{7\baselineskip}
\noindent\begin{minipage}{\linewidth}
\hypertarget{TERM-1888}{}
\noindent\textbf{Odkazka na numerovano svojstvo}\par\noindent\textbf{Reference to a numbered property}\quad{\footnotesize\texttt{TERM-1888}}\par
\smallskip\noindent\emph{Razlog:} Gubi svojstvo 1 jest odkazka na točno numerovano uslovje, ne obča fraza o gubljenju informaciji. V definiciji veličiny prvogo tipa pri adjunkciji vvodet algebraično-drobno čislo v oblast cělosti. Dokaz potom pomnožuje $y(\eta)$ polinomom, aby to svojstvo prestalo iměti město. Simbol Y iz kratkogo zapisa ne smi zaměniti aktualne $y$ i $\eta$ ni izbrisati broj 1. Zachovanje te odkazki jest razlog izbora frazy i uslovje sledovanja dokazovoj argumentaciji.\par
\noindent\emph{Reason:} Gubi svojstvo 1 refers to a precisely numbered condition, not to a general loss of information. In the definition, adjoining quantities of the first type introduces an algebraically fractional number into the integral domain. The proof then multiplies $y(\eta)$ by a polynomial so that this property ceases to hold. The Y of the abbreviated record must not replace the actual $y$ and $\eta$ or remove the number 1. Preserving that reference is the reason for the phrase and is necessary for following the argument.\par
\noindent{\footnotesize Izvorne zapisy / Original records: HIST-g065-017.}\par
\smallskip\noindent\emph{Sravnjenje s tekstom:} Broj odkazki i aktualne oznaky trebujut sohranjenja.\par
\noindent\emph{Textual comparison:} The reference number and actual symbols must be preserved.\par
\noindent{\footnotesize TEXT-COMP-945; pasaže / passages: 7857, 7858, 7859, 7860, 7861, 7862, 7863, 7864, 7865, 7871, 7872, 7903, 7904, 7905, 7906, 7907, 7908, 7909, 7910, 7911; \texttt{CORPUS\_FORM\_COMPARISONS.json}.}\par
\end{minipage}\par
\par\medskip\Needspace{7\baselineskip}
\noindent\begin{minipage}{\linewidth}
\hypertarget{TERM-1891}{}
\noindent\textbf{Čislovo fiksovanje koeficientov}\par\noindent\textbf{Fixing coefficients numerically}\quad{\footnotesize\texttt{TERM-1891}}\par
\smallskip\noindent\emph{Razlog:} Čislovo fiksovati označuje zadanje čislovyh cěnnostij někoj česti koeficientov posrědstvom linearnoj transformacije. Arhiv izrično odlišuje to od algebraičnogo izključenja koeficientov. Pasaža dodava ključno uslovje: relacije, platne dlja specialnoj formy, ostajut platne i v občem slučaju. Zato čislovo ne znači, že ide samo o približnom numeričnom izračunu. Glagol i njegova dopolnjenja čuvajut točnu rolu specializacije v argumentu.\par
\noindent\emph{Reason:} Čislovo fiksovati denotes assigning numerical values to some coefficients by a linear transformation. The record explicitly distinguishes this from algebraic elimination of the coefficients. The passage adds the crucial condition that relations valid for the special form remain valid in the general case. Numerically therefore does not mean merely an approximate numerical computation. The verb and its complements preserve the exact role of specialization in the argument.\par
\noindent{\footnotesize Izvorne zapisy / Original records: HIST-g043-014.}\par
\smallskip\noindent\emph{Sravnjenje s tekstom:} Fiksovanje cěnnostij ne jest algebraično izključenje.\par
\noindent\emph{Textual comparison:} Fixing values is not algebraic elimination.\par
\noindent{\footnotesize TEXT-COMP-948; pasaže / passages: 5384; \texttt{CORPUS\_FORM\_COMPARISONS.json}.}\par
\end{minipage}\par
\par\medskip\Needspace{7\baselineskip}
\noindent\begin{minipage}{\linewidth}
\hypertarget{TERM-1921}{}
\noindent\textbf{Padežne formy nazvy v naslově}\par\noindent\textbf{Case forms of a term in a title}\quad{\footnotesize\texttt{TERM-1921}}\par
\smallskip\noindent\emph{Razlog:} Razpadno polje jest osnovna nazva; razpadne polja i razpadnyh poljah sut jej pluralne padežne formy, ne tri razne termina. V naslově sovmestnoj raboty s Richardom Brauerom predlog o potrěbuje lokativ: o minimalnyh razpadnyh poljah. Arhivny komentar izbral slovjansku osnovu razpad- naměsto germanskogo zaimovanja i upominal anglijsku kontrolu, no bez dostatočnogo citata dlja novogo nezavisnogo podtvrđenja. Gramatika naslova i priznanje sovavtorstva ostajut oddělne od izbora leksikona.\par
\noindent\emph{Reason:} Razpadno polje is the headword; razpadne polja and razpadnyh poljah are its plural case forms, not three different terms. In the title of the joint paper with Richard Brauer, o takes the locative: o minimalnyh razpadnyh poljah. The archived note preferred the Slavic razpad- formation to a Germanic borrowing and mentioned an English control, but supplied no sufficiently specific citation for a new independent corroboration. Title grammar and acknowledgement of coauthorship are separate from the lexical choice.\par
\noindent{\footnotesize Izvorne zapisy / Original records: HIST-g166-001.}\par
\smallskip\noindent\emph{Sravnjenje s tekstom:} Lokativ v naslově ne jest novy termin.\par
\noindent\emph{Textual comparison:} The title's locative is not a new term.\par
\noindent{\footnotesize TEXT-COMP-978; pasaže / passages: 19171, 19173; \texttt{CORPUS\_FORM\_COMPARISONS.json}.}\par
\end{minipage}\par
\par\medskip\Needspace{7\baselineskip}
\noindent\begin{minipage}{\linewidth}
\hypertarget{TERM-1952}{}
\noindent\textbf{Kratka odkazka na izvor dokaza}\par\noindent\textbf{A short reference to a proof source}\quad{\footnotesize\texttt{TERM-1952}}\par
\smallskip\noindent\emph{Razlog:} Za dokaz sravn. H. Grell jest kratka bibliografična odkazka. Ona pripisuje izvor dokaza imenovanomu matematikovi, ne tvrdi, že prevoditelj sam izvršil novy nezavisny dokaz. Aktualna podnožna primětka piše Za dokaz sravni. H. Grell.; razlika meždu skračenjem i polnějšim glagolom ne měni adresata odkazky. Arhiv opisuje obnovjenje pečatnoj primětky, no nyněšnje objasnjenje ne tvrdí novu vizualnu proverku skena. Ime, směr odkazky i njezina uloga morajut ostati vidime.\par
\noindent\emph{Reason:} Za dokaz sravn. H. Grell is a short bibliographic reference. It identifies the named mathematician as the source of a proof; it does not claim that the translator supplied a fresh independent proof. The present footnote writes Za dokaz sravni. H. Grell. The abbreviation and fuller verb have the same addressee. The archive describes restoration of the printed note, but this explanation claims no fresh visual scan check. The name, direction of reference and function must remain visible.\par
\noindent{\footnotesize Izvorne zapisy / Original records: HIST-g161-008.}\par
\smallskip\noindent\emph{Sravnjenje s tekstom:} Bibliografična odkazka zachovuje čužy naučny doprinos.\par
\noindent\emph{Textual comparison:} The bibliographic reference preserves another scholar's contribution.\par
\noindent{\footnotesize TEXT-COMP-1009; pasaže / passages: 19003, 18932; \texttt{CORPUS\_FORM\_COMPARISONS.json}.}\par
\end{minipage}\par
\par\medskip\Needspace{7\baselineskip}
\noindent\begin{minipage}{\linewidth}
\hypertarget{TERM-1954}{}
\noindent\textbf{Prevedena proza i zachovany bibliografičny naslov}\par\noindent\textbf{Translated prose and a retained bibliographic title}\quad{\footnotesize\texttt{TERM-1954}}\par
\smallskip\noindent\emph{Razlog:} Arhivna fraza K aksiomatice ... Idealtheorie ne jest polny bibliografičny naslov. Aktualna podnožna primětka prevodi uvodno objasnjenje o aksiomatice multiplikacijskyh prstenov, a potom odkazuje na dělo skračeno imenovano Idealtheorie. Razlog v arhivě, že cěly red jest naslov, zato trěbuje popravu. Zachovany německy naslov pomaga najdti citovanu rabotu; prevedena okolna proza objasnjaje, dlja čego citat služi. To ne jest směšany medžuslovjansko-německy termin dlja aksiomatiki.\par
\noindent\emph{Reason:} The archived phrase K aksiomatice ... Idealtheorie is not a complete bibliographic title. The actual footnote translates the introductory explanation about the axiomatics of multiplication rings and then cites the work under the short title Idealtheorie. The archived rationale treating the whole string as a title therefore needs correction. The retained German title helps identify the cited work; the translated surrounding prose explains why it is cited. This is not a hybrid Interslavic–German technical term for axiomatics.\par
\noindent{\footnotesize Izvorne zapisy / Original records: HIST-g155-009.}\par
\smallskip\noindent\emph{Sravnjenje s tekstom:} Okolna proza ne jest čestju bibliografičnogo naslova.\par
\noindent\emph{Textual comparison:} The surrounding prose is not part of the bibliographic title.\par
\noindent{\footnotesize TEXT-COMP-1011; pasaže / passages: 18789, 18861, 18897; \texttt{CORPUS\_FORM\_COMPARISONS.json}.}\par
\end{minipage}\par
\par\medskip\Needspace{7\baselineskip}
\noindent\begin{minipage}{\linewidth}
\hypertarget{TERM-1955}{}
\noindent\textbf{Avtorkina primětka o proverki izračuna}\par\noindent\textbf{The author's note on checking a calculation}\quad{\footnotesize\texttt{TERM-1955}}\par
\smallskip\noindent\emph{Razlog:} Za kontrolu računa uvodi avtorkin sposob proverky izračuna. V sledovajučej tablicě stojat dve specialne formy i njihove vrednosti; one dopolnjajut dvojnu redukciju. Račun znači matematičny izračun, ne račun za platbu. Tuto jest jezyk samej raboty, ne komentar o proverki AI-prevoda. Citovanje dvuh testovyh slučajev ne daje osnovu tvrdženju, že tym byla nezavisno dokazana vsaka formula v cěloj rabotě.\par
\noindent\emph{Reason:} Za kontrolu računa introduces the author's way of checking a computation. The following table gives two special forms and their values, supplementing double reduction. Račun means mathematical calculation, not a bill. This is the language of the original work, not commentary on checking the AI translation. The two test cases do not justify a claim that every formula in the paper has thereby received an independent proof.\par
\noindent{\footnotesize Izvorne zapisy / Original records: HIST-g008-010.}\par
\smallskip\noindent\emph{Sravnjenje s tekstom:} Dva slučaja ne zaměnjajut vse dokazy.\par
\noindent\emph{Textual comparison:} Two cases do not replace all proofs.\par
\noindent{\footnotesize TEXT-COMP-1012; pasaže / passages: 1105, 1109, 1110, 1111, 1114, 1115, 1116; \texttt{CORPUS\_FORM\_COMPARISONS.json}.}\par
\end{minipage}\par
\par\medskip\Needspace{7\baselineskip}
\noindent\begin{minipage}{\linewidth}
\hypertarget{TERM-1966}{}
\noindent\textbf{Oznaka lokalnogo dopolnjenja teoremy}\par\noindent\textbf{A label for a local addition to a theorem}\quad{\footnotesize\texttt{TERM-1966}}\par
\smallskip\noindent\emph{Razlog:} Dodatok jest kratka oznaka lokalnogo dopolnjenja teoremy. Ne označaje bibliografičny dodatok na koncu knigy i ne dodavaje novy razděl samym slovom. Zde primětka oslabjaje potrěbu redukovanogo predstavjenja do najkratšego i potom objasnjaje, čto ostaje jedinstveno. Oznaka mora ostati pri tom argumentě, da čitatelj vidi jeji zavisimost od teoremy IX. Dodatok i korolar mogut iměti blizku izložiteljnu ulogu, no arhiv nedava razloga zaměniti avtorkinu oznaku bez potrěby.\par
\noindent\emph{Reason:} Dodatok is a short label for a local addition to a theorem. It is not a bibliographic appendix at the end of the book and does not by itself create a new section. Here the note weakens reduced representation to shortest representation and explains what remains unique. The label belongs beside that argument so readers can see its dependence on theorem IX. An addendum and a corollary may have similar expository functions, but the archive supplies no reason to replace the author's label unnecessarily.\par
\noindent{\footnotesize Izvorne zapisy / Original records: HIST-g081-049, HIST-g082-049.}\par
\smallskip\noindent\emph{Sravnjenje s tekstom:} Lokalna oznaka ostaje pri svojem dokazu.\par
\noindent\emph{Textual comparison:} The local label stays with its proof.\par
\noindent{\footnotesize TEXT-COMP-1023; pasaže / passages: 13028, 13030, 13032, 13034, 13036, 13038; \texttt{CORPUS\_FORM\_COMPARISONS.json}.}\par
\end{minipage}\par
\par\medskip\Needspace{7\baselineskip}
\noindent\begin{minipage}{\linewidth}
\hypertarget{TERM-1979}{}
\noindent\textbf{Padežna forma nazvy medžupolje}\par\noindent\textbf{An inflected form of medžupolje}\quad{\footnotesize\texttt{TERM-1979}}\par
\smallskip\noindent\emph{Razlog:} Medžupolja i medžupolj ne treba čitati kak dve nove jedinstvene nazvy. V lokalizovanoj větě konečno mnogo medžupolj poslednja forma služi kak genitiv množiny. Arhivny komentar, ktory nazyvaje ju kratkoju nazvoju i daje kiriličny prepis, zato ne jest dostatočny osnov dlja napravjenja iz nje samostalnogo osnovnogo slova. Značenje ostaje polje medžu dvoma ukazanymi poljami. Determinističny prepis v kirilicu ne daje nezavisno podtvrđenje ni gramatiki, ni razumljivosti složenicy.\par
\noindent\emph{Reason:} Medžupolja and medžupolj should not be read as two newly coined singular headwords. In the located sentence konečno mnogo medžupolj, the latter functions as a genitive plural. The archived comment calling it a compact term and giving a Cyrillic rendering is therefore insufficient grounds for treating it as a separate lemma. The meaning remains a field between two specified fields. Deterministic Cyrillic conversion independently verifies neither grammar nor intelligibility of the compound.\par
\noindent{\footnotesize Izvorne zapisy / Original records: HIST-g065-012.}\par
\smallskip\noindent\emph{Sravnjenje s tekstom:} Kirilica ne tvori vtoru nezavisnu jezyčnu proverku.\par
\noindent\emph{Textual comparison:} Cyrillic does not supply a second independent linguistic check.\par
\noindent{\footnotesize TEXT-COMP-1036; pasaže / passages: 7913, 7914, 7915, 7916, 7917, 7918, 7919, 7920; \texttt{CORPUS\_FORM\_COMPARISONS.json}.}\par
\end{minipage}\par
\par\medskip\Needspace{7\baselineskip}
\noindent\begin{minipage}{\linewidth}
\hypertarget{TERM-1988}{}
\noindent\textbf{Diskriminant kako kvadrat determinanta}\par\noindent\textbf{Discriminant as a determinant square}\quad{\footnotesize\texttt{TERM-1988}}\par
\smallskip\noindent\emph{Razlog:} Opis diskriminant kako kvadrat determinanta razjasnjaje formulu, a ne vvodit nezavisny termin. $S$ tut jest slěd, a alfa-elementy tvorjat bazu. Kvadrat odnaša se k cělomu determinantu matricy konjugovanyh baznyh elementov, ne k oddělnym elementam matricy. Vertikalne črty označajut determinant, ne apsolutnu veličinu ni kvadrat kompleksnogo modula. Arhiv spominaje ranějšu prověrku izvornogo obraza; nynějšno objašnjenje opiraje se na zachovanu formulu i kontekst, ne tvrdi novu prověrku obraza.\par
\noindent\emph{Reason:} The description discriminant as the square of a determinant explains the formula rather than introducing an independent term. Here $S$ is the trace and the alpha elements form a basis. The square applies to the entire determinant of the matrix of conjugate basis elements, not to its individual entries. The vertical bars denote a determinant, not absolute value or squared complex modulus. The archive mentions an earlier image comparison; this explanation relies on the preserved formula and context and does not claim a new image check.\par
\noindent{\footnotesize Izvorne zapisy / Original records: HIST-g147-003.}\par
\smallskip\noindent\emph{Sravnjenje s tekstom:} Črty označajut determinant, ne kompleksny modul.\par
\noindent\emph{Textual comparison:} The bars denote a determinant, not complex modulus.\par
\noindent{\footnotesize TEXT-COMP-1045; pasaže / passages: 18513, 18515, 18519, 18520, 18521, 18522, 18523, 18525; \texttt{CORPUS\_FORM\_COMPARISONS.json}.}\par
\end{minipage}\par
\par\medskip\Needspace{7\baselineskip}
\noindent\begin{minipage}{\linewidth}
\hypertarget{TERM-1989}{}
\noindent\textbf{Neizměnny zapis slědovogo determinanta}\par\noindent\textbf{Unchanged trace-determinant notation}\quad{\footnotesize\texttt{TERM-1989}}\par
\smallskip\noindent\emph{Razlog:} To jest izbor zapisa, ne prevod slova: izraz $|S(\alpha_i\alpha_j)|$ ostaje matematičnym identifikatorom v različnyh jazykah i pismenah. Latinsko $S$, grčke bukvy i indeksy ne trěba transliterovati kako običnu prozu. Kontekst govori o nezavisnoj bazi polja prvogo roda i nenulovom diskriminantu; ne dava takovo tvrdženje za proizvoljnu zavisnu množinu. Arhivna naměra zachovati formulu bajt za bajtom ne zaměnjuje prověrku vsakoj gotovoj versije.\par
\noindent\emph{Reason:} This is a notation choice, not a word translation: $|S(\alpha_i\alpha_j)|$ remains a mathematical identifier across languages and scripts. Latin $S$, Greek letters and indices must not be transliterated as ordinary prose. The context concerns an independent basis of a field of the first kind and a nonzero discriminant; it makes no such assertion for an arbitrary dependent set. The archived intention to preserve the formula byte for byte does not substitute for checking each resulting edition.\par
\noindent{\footnotesize Izvorne zapisy / Original records: HIST-g157-007.}\par
\smallskip\noindent\emph{Sravnjenje s tekstom:} Nezavisnost bazy jest čestju konteksta.\par
\noindent\emph{Textual comparison:} Independence of the basis is part of the context.\par
\noindent{\footnotesize TEXT-COMP-1046; pasaže / passages: 18857, 18859, 18861; \texttt{CORPUS\_FORM\_COMPARISONS.json}.}\par
\end{minipage}\par
\par\medskip\Needspace{7\baselineskip}
\noindent\begin{minipage}{\linewidth}
\hypertarget{TERM-1990}{}
\noindent\textbf{Determinanty n-torkov i jih ideal}\par\noindent\textbf{Determinants of n-tuples and their ideal}\quad{\footnotesize\texttt{TERM-1990}}\par
\smallskip\noindent\emph{Razlog:} Slovo determinant lokalizuje objašnjenje, medžutym $|S(\tau_i\tau_j)|$ ostaje formulom. Tau-elementy proběgajut vse razsmatrjane n-torky; ne sut od počętka vsi nazvane bazami. Determinant jest nenulovy pri nezavisnosti, a ideal jest generovany determinantami tyh n-torkov. Zaměna tau na alfa radi vněšne jednolikosťi by zakryla razliku dvuh kontekstov. Takže trěba zachovati razliku medžu množinoju generatorov ideala i jednym generatorom v specialnom slučaju.\par
\noindent\emph{Reason:} The word determinant localizes the explanation while $|S(\tau_i\tau_j)|$ remains a formula. The tau elements range over the relevant n-tuples; they are not all assumed at the outset to be bases. The determinant is nonzero when the elements are independent, and the ideal is generated by the determinants of these tuples. Replacing tau by alpha for visual uniformity would hide the distinction between the contexts. The set of ideal generators must also be distinguished from the single generator in the special case.\par
\noindent{\footnotesize Izvorne zapisy / Original records: HIST-g158-003.}\par
\smallskip\noindent\emph{Sravnjenje s tekstom:} Ne vsaka razsmatrjana n-torka jest baza.\par
\noindent\emph{Textual comparison:} Not every tuple considered is a basis.\par
\noindent{\footnotesize TEXT-COMP-1047; pasaže / passages: 18895, 18897; \texttt{CORPUS\_FORM\_COMPARISONS.json}.}\par
\end{minipage}\par
\par\medskip\Needspace{7\baselineskip}
\noindent\begin{minipage}{\linewidth}
\hypertarget{TERM-1991}{}
\noindent\textbf{Latinska bukva v odsylkě na §12 C)}\par\noindent\textbf{Latin letter in the reference to §12 C)}\quad{\footnotesize\texttt{TERM-1991}}\par
\smallskip\noindent\emph{Razlog:} Odsylka §12 C) označaje konkretno podrazděljenje i ne jest tekst dlja jazykovoj zaměny. Latinska bukva C mora ostati C, takože v kiriličnom tekstu; podobny vid kiriličnoj bukvy ne garantuje jednaky identifikator. Arhiv predpisuje zaščitny TeX-zapis s mbox, ale nynějšna latinska odsylka ima običny zapis C. Zato tut dokumentujemo naměru i identitet odsylky, ne tvrdimo, že arhivny obvoj jest v nynějšnom tekstu realizovany ili že vsaka kirilična odsylka byla nyně prověrena.\par
\noindent\emph{Reason:} The reference §12 C) identifies a particular subsection and is not prose to be translated. Latin C must remain C, including in Cyrillic text; a similar-looking Cyrillic letter does not guarantee the same identifier. The archive prescribes a protective TeX mbox wrapper, but the present Latin reference uses plain C. This entry therefore documents the intended protection and reference identity, not a claim that the archived wrapper is implemented or that every Cyrillic reference has now been checked.\par
\noindent{\footnotesize Izvorne zapisy / Original records: HIST-g022-009.}\par
\smallskip\noindent\emph{Sravnjenje s tekstom:} Nynějšny zapis ne ima arhivno predpisanogo mbox.\par
\noindent\emph{Textual comparison:} The present notation lacks the mbox prescribed in the archive.\par
\noindent{\footnotesize TEXT-COMP-1048; pasaže / passages: 2858; \texttt{CORPUS\_FORM\_COMPARISONS.json}.}\par
\end{minipage}\par
\par\medskip\Needspace{7\baselineskip}
\noindent\begin{minipage}{\linewidth}
\hypertarget{TERM-1999}{}
\noindent\textbf{Naslov obzora konstrukcije}\par\noindent\textbf{Heading for an overview of the construction}\quad{\footnotesize\texttt{TERM-1999}}\par
\smallskip\noindent\emph{Razlog:} Obzor tvorjenja sistema jest orientujuči naslov, ne novy matematičny termin. On vvodit tabličny opis konstrukcije: kontrakcije, nerazložime formy, redukcije i nove reduktanty. Tvorjenje bylo izbrano kako obšče nazvanje postupka, soglasno drugim městam prevoda. Obzor ne dava dozvoljenje skratiti avtorkinu konstrukciju do rezjume: tablica i jeje vysvětljenja sut predmetom polnogo prevoda. Takovy izbor trěba dokumentovati s terminologijeju, ale ne čisliti kako novo slovo matematičnogo slovnika.\par
\noindent\emph{Reason:} Obzor tvorjenja sistema is an orienting heading, not a new mathematical term. It introduces a tabular account of the construction: contractions, indecomposable forms, reductions and new reductants. Tvorjenje was chosen as a general name for the process, consistently with other passages. Overview does not license replacing the author's construction with a summary: the table and its explanations require complete translation. Such a choice belongs in the documentation but should not be counted as a new mathematical word.\par
\noindent{\footnotesize Izvorne zapisy / Original records: HIST-g008-002.}\par
\smallskip\noindent\emph{Sravnjenje s tekstom:} Naslov orientuje čitatelja v polnoj konstrukciji.\par
\noindent\emph{Textual comparison:} The heading guides the reader through the full construction.\par
\noindent{\footnotesize TEXT-COMP-1056; pasaže / passages: 1043, 1048, 1076, 1078, 1080, 1082, 1084; \texttt{CORPUS\_FORM\_COMPARISONS.json}.}\par
\end{minipage}\par
\par\medskip\Needspace{7\baselineskip}
\noindent\begin{minipage}{\linewidth}
\hypertarget{TERM-2000}{}
\noindent\textbf{Od porędka k prstenam nad poljem}\par\noindent\textbf{From an order to rings over a field}\quad{\footnotesize\texttt{TERM-2000}}\par
\smallskip\noindent\emph{Razlog:} Naslov sohranjaje prěhod od algebričnogo porędka k prstenam konečnogo ranga nad poljem, ne prěhod od običnogo raspoređenja. Kontekst tvori količnik $R=T/Tp$ nad poljem, izomorfnym sootvětnomu količniku bazovogo prstena. Arhivny zapis porjadka, kolcam, vzhodno i nynějšny poredka, prstenam, po odnošenju ne sut jedna i ta že površnost; slovnik ih pokazuje kako svęzane realizacije, ne pripisuje staru frazu nynějšnomu tekstu. Predlog mora izraziti bazovo polje otnosno ktoroho jest rang konečny.\par
\noindent\emph{Reason:} The heading preserves the passage from an algebraic order to rings of finite rank over a field, not from an ordinary arrangement. The context forms the quotient $R=T/Tp$ over a field isomorphic to the corresponding quotient of the base ring. The archived wording porjadka, kolcam, vzhodno and the present poredka, prstenam, po odnošenju are not the same surface form; the dictionary relates them without attributing the old wording to the current text. The prepositional phrase must identify the base field relative to which the rank is finite.\par
\noindent{\footnotesize Izvorne zapisy / Original records: HIST-g151-001.}\par
\smallskip\noindent\emph{Sravnjenje s tekstom:} Naslov govori o algebričnom porędku.\par
\noindent\emph{Textual comparison:} The heading concerns an algebraic order.\par
\noindent{\footnotesize TEXT-COMP-1057; pasaže / passages: 18626, 18627, 18629, 18631, 18633; \texttt{CORPUS\_FORM\_COMPARISONS.json}.}\par
\end{minipage}\par
\par\medskip\Needspace{7\baselineskip}
\noindent\begin{minipage}{\linewidth}
\hypertarget{TERM-2011}{}
\noindent\textbf{Naslov obzora redukcij}\par\noindent\textbf{Heading for an overview of reductions}\quad{\footnotesize\texttt{TERM-2011}}\par
\smallskip\noindent\emph{Razlog:} Obzor redukcij jest kratky orientujuči naslov pred podrobnym spiskom slučajev. Izbor obzor soglasuje se s drugymi naslovami obzora; on ne zaměnjuje sam spisok ni jego formuly. Arhiv protivopostavlja obzor slovu pregled i dava tvrdenje o proishođenju poslědnjego bez priloženogo etimologičnogo dokaza. To tvrdenje ne prinimamo kako obosnovany lingvističny rezultat. Dostatočny lokalny razlog izbora jest jedinoobrazje naslovov i jasna funkcija obzora.\par
\noindent\emph{Reason:} Obzor redukcij is a short orienting heading before a detailed list of cases. Obzor is consistent with other overview headings; it does not replace the list or its formulas. The archive contrasts obzor with pregled and makes an etymological claim about the latter without supporting evidence. That claim is not adopted as an established linguistic result. The adequate local reason for the choice is consistency of headings and the clear function of an overview.\par
\noindent{\footnotesize Izvorne zapisy / Original records: HIST-g012-004.}\par
\smallskip\noindent\emph{Sravnjenje s tekstom:} Kratky naslov ne skračaje matematičny spisok.\par
\noindent\emph{Textual comparison:} A brief heading does not abbreviate the mathematical list.\par
\noindent{\footnotesize TEXT-COMP-1068; pasaže / passages: 1520, 1522, 1524, 1526, 1528, 1530; \texttt{CORPUS\_FORM\_COMPARISONS.json}.}\par
\end{minipage}\par
\normalsize

\clearpage\section{Leksični, gramatični i redakcijni izbory / Lexical, grammatical, and editorial choices}
Slědujuče sravnjenja dokumentujut konkretne izbory. Redy v tabelah odkazujut na izdanja, opreděljene v podatkah, ne avtomatično na najnovějši prěvod. Izvorna formula, historijna prěvodna forma i redakcijna poprava ne sut jedno i to samo.
\par\emph{EN:} These comparisons document particular choices. Table line numbers refer to the editions identified in the accompanying data, not automatically the latest translation. A source formula, historical translation, and editorial emendation are distinguished explicitly.
\small
\par\medskip\Needspace{8\baselineskip}
\hypertarget{DEC-0001}{}
\noindent\textbf{Neter- v běgu prozy; Noether v citacijah / Neter- in prose; Noether in citations}\quad{\footnotesize\texttt{DEC-0001}}\par
\smallskip\noindent\emph{Razlog:} Neter- jest izbrana osnova dlja sklanjanja i pridavnyh form v prevodě; Noether ostaje v avtorskyh imenah, bibliografičnyh naslovah i citacijah. Kratka osnova dopuščaje Neterov i Netera bez směšanja latinskogo imena s kiriličnymi koncovkami. Odbačene byli nesoglasne nacionalne adaptacije i avtomatično prěpisyvanje bibliografije. To jest redakcijna konvencija, ne izměreny optimum razumljivosti. V rabotě 45 Neterov fundamentalny teorem odnaša se k Maksu Noetheru: ime ne trěba avtomatično poistovětiti s Emmy Noether, avtorom korpusa. Sohranjanje identiteta osoby jest važnějše od jednolikosťi pisanja.\par
\smallskip\noindent\emph{Reason:} Neter- is the chosen stem for inflection and adjectival forms in the translation; Noether remains in author names, bibliographical titles and citations. The compact stem permits Neterov and Netera without mixing a Latin surname with Cyrillic endings. Inconsistent national adaptations and wholesale rewriting of bibliographical names were rejected. This is an editorial convention, not a measured optimum of intelligibility. In paper 45, the fundamental theorem refers to Max Noether: the name must not automatically be identified with Emmy Noether, the corpus author. Preserving the person's identity takes precedence over spelling uniformity.\par
\smallskip\noindent\emph{Izvory:} Rabota 45, naslov i uvod, svęz s Maksom Noetherom.\par
\noindent\emph{Sources:} Paper 45, heading and introduction, referring to Max Noether.\par
\par\medskip\Needspace{8\baselineskip}
\hypertarget{DEC-0002}{}
\noindent\textbf{Prědstavjenje tipa Weierstrass / Representation of Weierstrass type}\quad{\footnotesize\texttt{DEC-0002}}\par
\smallskip\noindent\emph{Razlog:} Konstrukcija prědstavjenje tipa Weierstrass ostavlja prezime točno čitlivym i daje odnošenje črez imenovnik tipa, vměsto tvorjenja nacionalno izgovorjenogo pridavnogo slova. TeX-označenje foreign čuva to prezime pri prěpisu v kirilicu; ono ne ustanovjaje lingvističnu normu samo po sobě. Alternativy byli adaptovati prezime ili dodati koncovku k neadaptovanoj osnovi. Izbrana fraza izběgaje oba problema bez proměny parametričnogo prědstavjenja, o ktorom govori teorema. Nynějšny odlomok ima upravo tu konstrukciju; ne jest to tvrdženje o vsih matematičnyh eponimah.\par
\smallskip\noindent\emph{Reason:} Prědstavjenje tipa Weierstrass leaves the surname exactly recognizable and expresses the relation through of the type instead of constructing a nationally adapted adjective. The TeX foreign marker protects the surname during Cyrillic conversion; the marker does not establish a linguistic norm by itself. The alternatives were to adapt the surname or attach an ending to an unadapted stem. The chosen phrase avoids both problems without changing the parametric representation discussed by the theorem. The present passage uses this construction; it is not a claim about every mathematical eponym.\par
\smallskip\noindent\emph{Izvory:} Rabota 13, priměr parametričnogo prědstavjenja tipa Weierstrass.\par
\noindent\emph{Sources:} Paper 13, example of a parametric representation of Weierstrass type.\par
\par\medskip\Needspace{8\baselineskip}
\hypertarget{DEC-0003}{}
\noindent\textbf{Prěz i črěz / Prěz and črěz}\quad{\footnotesize\texttt{DEC-0003}}\par
\smallskip\noindent\emph{Razlog:} Prěz jest edinym predlogom etoj redakcije za through, across, via. Oba prěz i črěz stojat v konsultovanoj kopiji Medžuslovjanskogo slovnika i oba upravljajut akuzativom. Zato izbor prěz ne oprědělaje črěz kako nepravilen. Zapisana motivacija jest rasprostranjenje poznavajemosti mimo ranějšego vostočnoslovjanskogo težišča. Slovničny profil prěz ima pozitivne i negativne ocenky; črěz ima drugu kombinaciju. Eto služi sravnenju, ne čitatelskomu ispytanju. Kriterije tvorjenja Medžuslovjanskogo davajut šest ravnověsnyh skupin, ale dlja etoj odluki ne jest zachovana čislena tablica glasov. Izloženje ne rekonstruuje iz nje izmysleny rezultat.\par
\smallskip\noindent\emph{Reason:} Prěz is this edition's consistent preposition for through, across and via. Both prěz and črěz occur in the consulted Interslavic Dictionary copy, and both govern the accusative. Choosing prěz therefore does not make črěz incorrect. The recorded motivation is to distribute recognizability beyond the translation's earlier East-Slavic emphasis. The dictionary profile for prěz contains both positive and negative assessments, and črěz has a different combination. These support comparison, not a reader experiment. The Interslavic design criteria specify six equally weighted groups, but no numerical vote table survives for this decision; no score is invented from that general policy.\par
\smallskip\noindent\emph{Izvory:} Medžuslovjansky slovnik, zaglavja 2284 i 4473; Jan van Steenbergen, Kriterije tvorjenja, izbor korenjev.\par
\noindent\emph{Sources:} Interslavic Dictionary, entries 2284 and 4473; Jan van Steenbergen, Design criteria, Which roots to use?\par
\par\medskip\Needspace{8\baselineskip}
\hypertarget{DEC-0004}{}
\noindent\textbf{Sut pravdive: istinnost tvrdženij / Sut pravdive: truth of assertions}\quad{\footnotesize\texttt{DEC-0004}}\par
\smallskip\noindent\emph{Razlog:} Za tvrdženja, relacije i zakony izbrana jest fraza sut pravdive vměsto neujasnjenogo matematičnogo značenja važiti. Konsultovany slovnik dava važiti v smislu váženja i uravnověšenja, a pravdivy vključaje značenje true. To podpiraje izrično vyslovjenje istinnosti, ale sam nedostatok značenja v jednom slovniku ne dokazuje, že vsaka srodna nacionalna upotrěba jest nelegitimna. Zapisana odluka važila za množinne predikaty; jedinstveno čislo i participy potrěbujut sootvětnu gramatiku. Cena izričnosti jest težša fraza. Pravdivy takože ne trěba avtomatično proširiti na vse smysly německogo gültig, napr. dopustimost konstrukcije ili obseg priměnjenja.\par
\smallskip\noindent\emph{Reason:} For assertions, relations and laws, sut pravdive was chosen instead of leaving the mathematical sense of važiti unexplained. The consulted dictionary gives važiti in the sense of weighing and balancing, while pravdivy includes true. This supports an explicit assertion of truth, but absence of a sense from one dictionary does not prove that every cognate national usage is illegitimate. The recorded decision concerned plural predicates; singulars and participles require their own grammatical forms. Explicitness costs a heavier phrase. Pravdivy must not automatically be extended to every sense of German gültig, such as admissibility of a construction or scope of application.\par
\smallskip\noindent\emph{Izvory:} Medžuslovjansky slovnik, zaglavja 3307 i 6005.\par
\noindent\emph{Sources:} Interslavic Dictionary, entries 3307 and 6005.\par
\par\medskip\Needspace{8\baselineskip}
\hypertarget{DEC-0005}{}
\noindent\textbf{Nyně / Now}\quad{\footnotesize\texttt{DEC-0005}}\par
\smallskip\noindent\emph{Razlog:} Nyně bylo izbrano kako jedino slovo za now vměsto směsi s teper. Konsultovany slovnik dava nyně i sravnenja s ruskim nyne, ukrajinskim nyni i češskim nyní; njegova značenjova glosa vključaje nynějšny čas. Izbor ima razlog v terminologičnoj konsistenciji i oficialnom zaglavju, ne v tvrdženju, že vsaka slovjanska skupina uvaža jego za običnu formu. Slovničny profil izrično sodrživa negativne ocenky za nekoliko južnyh jazykov. Zaměna vsěh form slovom sada by prěsunula prednost k drugoj skupině. Ni etoj izbor ni zapisano čislo pojavjenij ne zaměnjujut izměrjenje razumljivosti.\par
\smallskip\noindent\emph{Reason:} Nyně was selected as the consistent word for now instead of a mixture with teper. The consulted dictionary records nyně with Russian nyne, Ukrainian nyni and Czech nyní among its comparanda; its gloss includes present time. The justification is consistency and an established dictionary headword, not a claim that every Slavic group regards it as the ordinary form. The dictionary profile explicitly has negative assessments for several South Slavic languages. Replacing everything with sada would shift the preference to another group. Neither the choice nor its recorded occurrence count substitutes for a measurement of comprehension.\par
\smallskip\noindent\emph{Izvory:} Medžuslovjansky slovnik, zaglavje 5698.\par
\noindent\emph{Sources:} Interslavic Dictionary, entry 5698.\par
\par\medskip\Needspace{8\baselineskip}
\hypertarget{DEC-0006}{}
\noindent\textbf{Kde / Where}\quad{\footnotesize\texttt{DEC-0006}}\par
\smallskip\noindent\emph{Razlog:} Kde zjedinjuje ranějše gde i gdě v jednu slovničnu formu. Sravljajemy zapis pokazuje češske i slovačske kde, poljske gdzie, ruske gde, srbske gde i druge srodne formy. Pozitivne ocenky v slovničnom profilu davajut široku osnovu dlja etoj konvencije, ale ne značat, že vsi jazyky pišut to slovo jednako. Izbor ne proměnjaje smysl města ili uvod odnošenej klauzy; zadanje jest odstraniti neobosnovano kolebanje pisanja v odnom prevodě. Različne nacionalne formy ostajut legitimne v svojih jazykah.\par
\smallskip\noindent\emph{Reason:} Kde brings the earlier gde and gdě into one dictionary form. The comparative record includes Czech and Slovak kde, Polish gdzie, Russian gde, Serbian gde and other cognates. Positive assessments in the dictionary profile provide broad support for the convention but do not mean that all languages spell the word identically. The choice changes neither the spatial meaning nor its function in introducing a relative clause; it removes unmotivated spelling variation within one translation. The national forms remain legitimate in their own languages.\par
\smallskip\noindent\emph{Izvory:} Medžuslovjansky slovnik, zaglavje 942.\par
\noindent\emph{Sources:} Interslavic Dictionary, entry 942.\par
\par\medskip\Needspace{8\baselineskip}
\hypertarget{DEC-0007}{}
\noindent\textbf{Ješče / Still, yet, even}\quad{\footnotesize\texttt{DEC-0007}}\par
\smallskip\noindent\emph{Razlog:} Ješče bylo izbrano vměsto paralelnogo još, kogda kontekst ne dava razliku smysla. Slovnik povezuje ješče s vostočnymi i zapadnymi formami, ale takože s južnymi še, još, jošte, ušte i ošte. Etymologična srodnost ne znači bukvalnu grafičnu jednakost. Arhiv govori o pozitivnoj ocěnkě vsih jazykov, no sam profil ima približne ocenky dlja slovenskogo i srbskogo; to ograničenje trěba sohraniti. Upotrěba jednogo zaglavja jest redakcijna konvencija, ne dokaz, že vsak čitatelj bez usilja razlikuje vse smysly still, yet i even.\par
\smallskip\noindent\emph{Reason:} Ješče was preferred to parallel još where the context supplied no difference in meaning. The dictionary relates ješče to East and West Slavic forms but also to southern še, još, jošte, ušte and ošte. Etymological kinship does not mean literal identity of spelling. The archive speaks of positive assessments for every language, but the profile itself marks Slovene and Serbian only approximately; that qualification must be retained. Using one headword is an editorial convention, not proof that every reader effortlessly distinguishes all the senses still, yet and even.\par
\smallskip\noindent\emph{Izvory:} Medžuslovjansky slovnik, zaglavje 1076.\par
\noindent\emph{Sources:} Interslavic Dictionary, entry 1076.\par
\par\medskip\Needspace{8\baselineskip}
\hypertarget{DEC-0008}{}
\noindent\textbf{Jednolike formy i kontekstne izključki / Consistent forms and contextual exceptions}\quad{\footnotesize\texttt{DEC-0008}}\par
\smallskip\noindent\emph{Razlog:} V redakciji byli izbrany takože (vměsto takodže), togda (tada), trěba (treba), nehaj (neka v postavjenjah), ješče (jošče), už (uže), tu (tut), ponovno (znovu), poręd- (porjad-), libovoljny (matematično proizvoljny) i možno (možno jest/jest možno). To ne sut jedinnadset jednakovyh pravopisnyh popravok. Už i uže stojat v jednom slovničnom zaglavju, a tu i tut imajut odvojene zaglavja: kratše formy sut prednosti redakcije, ne zabrana drugih form. Pri nehaj trěba razlikovati postavjenje od něktory, ktory izražaje neoprěděljenost. Libovoljny ima slovnične značenja optional i arbitrary, no tip 3 i směšany jazykovy profil; proizvoljny takože jest zaglavje. Izbor v matematičnom smislu ne daje pravo nazvati libovoljny opće priznatym ili najbolje razumljivym. Možno jest kratka bezosobna konstrukcija, ne novo slovo. Čislo 1433 v zapisu opisuje obseg intervencije, ne rezultat vagovanja. Nynějšny osnovny pravopis može pisati pored- vměsto etymologičnogo poręd-.\par
\smallskip\noindent\emph{Reason:} The edition selected takože (for takodže), togda (tada), trěba (treba), nehaj (neka in stipulations), ješče (jošče), už (uže), tu (tut), ponovno (znovu), poręd- (porjad-), libovoljny (mathematical proizvoljny), and možno (možno jest/jest možno). These are not eleven equivalent spelling corrections. Už and uže share a dictionary entry, while tu and tut have separate entries: the shorter forms are editorial preferences, not prohibitions on the alternatives. Nehaj requires a distinction between stipulation and the indefinite něktory. Libovoljny is glossed optional and arbitrary but has type 3 and a mixed language profile; proizvoljny is also a headword. The mathematical preference therefore cannot be described as universally established or demonstrably best understood. Možno is a shorter impersonal construction, not a newly coined word. The recorded total of 1,433 describes the intervention's extent, not a weighting result. Present basic orthography may write pored- instead of etymological poręd-.\par
\smallskip\noindent\emph{Izvory:} Slovnik, zaglavja 1716, 2120, 13042, 9558, 1076, 114, 1758, 15606, 5961, 3145, 22300, 22318, 926; razlikovanje nehaj/něktory: odgovarjajuče zaglavje niže.\par
\noindent\emph{Sources:} Dictionary entries 1716, 2120, 13042, 9558, 1076, 114, 1758, 15606, 5961, 3145, 22300, 22318, 926; see the separate entry on nehaj/něktory.\par
\par\medskip\Needspace{8\baselineskip}
\hypertarget{DEC-0009}{}
\noindent\textbf{Po odnošenju k: odnos i padež / With respect to: relation and case}\quad{\footnotesize\texttt{DEC-0009}}\par
\smallskip\noindent\emph{Razlog:} Po odnošenju k vyražaje odnos bazy, polja, operatora, grupy, ranga, cělosti ili invariantnosti k danomu objektu. Slovnikovo zaglavje po odnošeńju k trěbuje dativ; osnovny pravopis piše ń kako n. Historijni razlog dava prednost toj izričnoj predložnoj konstrukciji prěd vzhodno k i prěd prěnosom nacionalnyh względem, vzhledem, glede ili u odnosu. On upominaje 401 pasažu i sravniteljne formy iz vostočnoj, zapadnoj i južnoj slovjanskoj skupiny, ale ne dava čislene težiny za izbor. Dolžina frazy jest priznana cena obščej prozračnosti; razumělivost nebyla izměrena čitatelskym testom. Točna starša kopija slovnika takože ima v odnošeńju k s tym samym smyslom i padežom. Zato jeho historijno odbacanje ne možno opirati na tvrdženje, že samo po odnošeńju k jest oficialno. Prednost po može byti konsistentna domova konvencija, ne izključna jezyčna norma. Dve točno lokalizovane pasaže iz diskriminantnoj teorije potvrđajut aktualno po odnošenju k; one ne zaměnjajut novu prověrku vseh 401 pojavjenj.\par
\smallskip\noindent\emph{Reason:} Po odnošenju k expresses the relation of a basis, field, operator, group, rank, integrality, or invariance to a specified object. Dictionary headword po odnošeńju k governs the dative; basic spelling writes ń as n. The historical rationale favours this explicit prepositional construction over vzhodno k and transfers of national względem, vzhledem, glede, or u odnosu. It reports 401 passages and East, West, and South Slavic comparanda but supplies no numerical selection weights. The phrase's length is acknowledged as a cost of shared transparency; intelligibility was not measured by a reader test. The exact earlier dictionary copy also contains v odnošeńju k with the same meaning and government. Its historical rejection therefore cannot rest on a claim that only po odnošeńju k is official. Preferring po can be a consistent house convention, not an exclusive linguistic norm. Two precisely located discriminant-theory passages attest current po odnošenju k; they do not substitute for a fresh check of all 401 occurrences.\par
\smallskip\noindent\emph{Izvory:} Medžuslovjansky slovnik: po odnošeńju k (18227), v odnošeńju k (19113), odnosno (4107), vȯzhodno (4185); oficialny pravopis, opcionalne bukvy; sravnjenje TEXT-COMP-06.\par
\noindent\emph{Sources:} Interslavic Dictionary: po odnošeńju k (18227), v odnošeńju k (19113), odnosno (4107), vȯzhodno (4185); official orthography, optional letters; comparison TEXT-COMP-06.\par
\smallskip\noindent\emph{Izvory / Sources:} \href{https://interslavic-dictionary.com/}{Medžuslovjansky slovnik / Interslavic Dictionary}, 18227, 19113, 4107, 4185; \href{https://interslavic.fun/learn/orthography/}{pravopis / orthography}, opcionalne bukvy / optional letters.
\par\noindent\emph{Sravnjenje / Comparison:} \hyperlink{TERM-0318}{TERM-0318}, \hyperlink{TERM-0327}{TERM-0327}; \texttt{TEXT-COMP-06}, \texttt{CORPUS\_FORM\_COMPARISONS.json}.
\par\medskip\Needspace{8\baselineskip}
\hypertarget{DEC-0010}{}
\noindent\textbf{Něktory i nehaj: něšto jest, ili nehaj bude / Něktory and nehaj: existence versus stipulation}\quad{\footnotesize\texttt{DEC-0010}}\par
\smallskip\noindent\emph{Razlog:} Nehaj uvodi želanje ili matematično postavjenje, napriměr nehaj x bude element. Něktory označaje neoprěděljeny objekt, napriměr něktora potenca jest děljiva. To različenje jest dano i v slovniku: nehaj jest častica let/may, něktory jest pridavno slovo some. Dve zachovane věty potrěbujut vtory smysl, ne prvu konstrukciju. Prva govori o suščestvovanji potencie ideala; vtora o stepenu asociovanogo prostogo ideala, dělimom primarnym idealom. Opisateljna fraza ktory imaje něktory stepenj takože čuva soglasje roda. Takovy razbor pokazuje, počemu mnogokratna forma neka ne dopuščaje edinu bezuslovnu zaměnu. Zapisane sedm pisanj przez pripadaju oddělnomu pravopisnomu izboru prěz.\par
\smallskip\noindent\emph{Reason:} Nehaj introduces a wish or mathematical stipulation, for example let x be an element. Něktory denotes an unspecified object, as in some power is divisible. The dictionary makes the distinction too: nehaj is the particle let/may, whereas něktory is the adjective some. The two preserved sentences require the second meaning, not the first construction. One asserts existence of a power of an ideal; the other concerns a power of its associated prime divisible by the primary ideal. The descriptive phrase which has some power also maintains gender agreement. This analysis explains why repeated neka did not admit one unconditional replacement. The seven recorded przez spellings belong to the separate orthographic choice prěz.\par
\smallskip\noindent\emph{Izvory:} Medžuslovjansky slovnik, zaglavja 9558 i 5692; rabota 19, priměčanje o nulah idealov; rabota 25, uvod.\par
\noindent\emph{Sources:} Interslavic Dictionary, entries 9558 and 5692; paper 19, footnote on zeros of ideals; paper 25, introduction.\par
\par\medskip\Needspace{8\baselineskip}
\hypertarget{DEC-0011}{}
\noindent\textbf{Pravdivost, izpolnjenje i priměnimost / Truth, satisfaction and applicability}\quad{\footnotesize\texttt{DEC-0011}}\par
\smallskip\noindent\emph{Razlog:} Německo gelten ne ima jedno medžuslovjansko sootvětny slovo za vse matematične konteksty. Tvrdženje jest pravdivo; uslovje jest izpolnjeno; teorema može byti priměnjena; pravilo ostavaje v silě; formula može služiti kako analog. Obična uvodna fraza jest imaje město, s množinnoj formoj imajut město. Izbor izběgaje směšanje s važiti (vážiti težinu) i platiti (platiti děngami), no ne otricaje legitimnost češskogo platí ili hrvatskogo vrijedi v ih vlastnyh jazykah. Ranějše sut pravdive jest dobro samo pri sootvětnom množinnom subjekte i istinnostnom smislu. Zachovany plan razděljuje 42 kontekstne pererabotky od 311 prostyh predikatov; iz 87 pojavjenij sut pravdive 73 ostajut bez proměny, 14 potrěbujut soglasje ili drugy smysl. Tako iz slovnika slěduje ne avtomatična zaměna, ale potrěba razbrati, čto jest pravdivo, izpolnjeno ili priměnimo.\par
\smallskip\noindent\emph{Reason:} German gelten has no single Interslavic equivalent across all mathematical contexts. An assertion is true; a condition is satisfied; a theorem applies; a rule remains in force; a formula may serve as an analogue. The ordinary introductory phrase is imaje město, with plural imajut město. The choice avoids confusion with važiti (weigh) and platiti (pay), without denying the legitimacy of Czech platí or Croatian vrijedi in their own languages. The earlier sut pravdive is appropriate only with the corresponding plural subject and truth sense. The preserved plan distinguishes 42 contextual recastings from 311 simple predicates; of 87 occurrences of sut pravdive, 73 remain and 14 require agreement or a different sense. The dictionary thus motivates an analysis of what is true, satisfied or applicable, not an automatic substitution.\par
\smallskip\noindent\emph{Izvory:} Slovnik, zaglavja 6005, 2114, 3307, 597, 875, 2721; polne kontekstne sravnenja predikatov pravdivosti i priměnimosti v prilogu s podatkami.\par
\noindent\emph{Sources:} Dictionary entries 6005, 2114, 3307, 597, 875, 2721; complete contextual comparisons of truth and applicability predicates in the accompanying data.\par
\par\medskip\Needspace{8\baselineskip}
\hypertarget{DEC-0012}{}
\noindent\textbf{Zato že; pričina ne jest dopuščenje / Zato že: cause is not concession}\quad{\footnotesize\texttt{DEC-0012}}\par
\smallskip\noindent\emph{Razlog:} Zato že jest izbrana obična pričinna spojka, kogda německy tekst dava razlog. Kratko bo, opisateljna konstrukcija i porędok klauz ostajut možnosti po kontekstu. Hoti označaje dopuščenje; jednako protivostavjenje; prinajmenje ograničenu minimalnu tvrdnju. Zato je ne možno směšati samo zato, že vsi stojat pri počętku věty. Izvory ne sut polno soglasne o poneže: gramatičny spisok dava whereas, no slovnično zaglavje 10909 dava because, since, inasmuch as. To nesoglasje oslablja ranějše tvrdženje, že pričinno poneže jest jednostavno neoficialno. Gramatika takože dopuščaje ibo i jer/jerbo. Zato že ostaje obosnovana jednolika prednost etoj redakcije, ne jedina možna norma. Zapis odlišaje dve istinne dopustiteljne konstrukcije s hoti, jednu s potrebnym prinajmenje, 21 nepravilno protivostavjeny kontekst s jednako i dve istinne protivostavjenja; čitatelsko ispytanje ne bylo dokazano.\par
\smallskip\noindent\emph{Reason:} Zato že is the edition's ordinary causal conjunction where the German text supplies a reason. Short bo, an analytic construction and different clause order remain contextual options. Hoti marks concession, jednako contrast, and prinajmenje a restricted minimum claim. They are not interchangeable merely because each can introduce a sentence. The authorities disagree about poneže: the grammar's list gives whereas, but dictionary entry 10909 gives because, since and inasmuch as. That disagreement weakens the earlier assertion that causal poneže was simply unofficial. The grammar also allows ibo and jer/jerbo. Zato že remains a defensible consistent preference, not the only possible norm. The decision distinguishes two genuine hoti concessions, one context requiring prinajmenje, 21 non-adversative jednako contexts and two genuine contrasts; it supplies no reader experiment.\par
\smallskip\noindent\emph{Izvory:} Jan van Steenbergen, Gramatika, Spojky; Slovnik, zaglavja 10909 i 8296.\par
\noindent\emph{Sources:} Jan van Steenbergen, Grammar, Conjunctions; Dictionary entries 10909 and 8296.\par
\par\medskip\Needspace{8\baselineskip}
\hypertarget{DEC-0013}{}
\noindent\textbf{Prsten: algebraična struktura / Prsten: an algebraic ring}\quad{\footnotesize\texttt{DEC-0013}}\par
\smallskip\noindent\emph{Razlog:} Prsten jest nazvanje algebraičnoj struktury s dodavanjem i množenjem, ne nov znak dlja fizičnogo kružka. Sravnenje imalo četiri kandidaty: prsten/pierścień, kolco, okruh i kolobar. Zapisane šest ravnověsnyh skupin davajut: ruska 1 dlja kolco; ukrajinsko-běloruska 1 dlja kolco; poljska 1 dlja prsten/pierścień; češko-slovačska 1 dlja okruh; slovensko-BCS skupina 0,5 dlja kolobar i 0,5 dlja prsten; bulgarsko-makedonska 1 dlja prsten. Suma jest 2,5, 2, 1 i 0,5 sootvětno. Prsten ima najvyši rezultat, no ne absolutnu bolšinu iz šesti. Čislo slov v dokumentah ne jest vagoju jazykov; sravnenja imajut različnu dlžinu i pokrytje. Običny slovnik podpiraje fizičny korenj, a matematične izvory potrěbny sut dlja tehničnogo smysla. Genitiv jedinstvenogo čisla i množinne formy byli razlišane kontekstno; bibliografičny original ne byl preimenovany. Vagi sut redakcijna heuristika, ne statističny dokaz razumljivosti.\par
\smallskip\noindent\emph{Reason:} Prsten names the algebraic structure with addition and multiplication, not a new symbol for a physical ring. Four candidates were compared: prsten/pierścień, kolco, okruh and kolobar. The six equally weighted groups were assigned as follows: Russian, 1 to kolco; Ukrainian–Belarusian, 1 to kolco; Polish, 1 to prsten/pierścień; Czech–Slovak, 1 to okruh; Slovene–BCS, 0.5 to kolobar and 0.5 to prsten; Bulgarian–Macedonian, 1 to prsten. The totals are respectively 2.5, 2, 1 and 0.5. Prsten has the largest share, not an absolute majority of six. Document token counts are not language weights; the comparators differ in length and coverage. The general dictionary supports the physical root, while mathematical sources are needed for the technical meaning. Singular genitives and plural forms were distinguished contextually, and an original bibliographical title was not renamed. The weights are an editorial heuristic, not a statistical demonstration of comprehension.\par
\smallskip\noindent\emph{Izvory:} Jan van Steenbergen, Kriterije tvorjenja, šest skupin; polna tablica vagov i sravnenje pri zaglavju prsten.\par
\noindent\emph{Sources:} Jan van Steenbergen, Design criteria, six groups; complete weight table and comparison under prsten.\par
\par\medskip\Needspace{8\baselineskip}
\hypertarget{DEC-0014}{}
\noindent\textbf{Sostojati se iz: sostav / Sostojati se iz: composition}\quad{\footnotesize\texttt{DEC-0014}}\par
\smallskip\noindent\emph{Razlog:} Sostojati se iz bylo izbrano za sostav objektov iz častej. Slovnik dava sȯstojati sę iz s osnovoj sȯstoji, no takože dopuščaje skladati sę iz. Izbor prvogo jest konsistencija, ne dokaz nepravosťi vtorogo. Predlog iz upravljaje genitivom; povratna častica se jest čast izbranej frazy. No německo bestehen ne znači vsegda consist of: tvrdženje o tom, v čem leži problema ili v čem jest značenje metoda, potrěbuje drugu konstrukciju. Zapis razlišaje 143 proměny i sedm sohranjenyh kontekstov, uključajuči pet smyslovyh pererabotok. Zato pravilo čuva sostav, ali ne nalaga jego na vsako podobno slovo. Nepravilnu prezentnu osnovu ne popravjaje prosto dodavanje se.\par
\smallskip\noindent\emph{Reason:} Sostojati se iz was chosen for composition of an object from parts. The dictionary gives sȯstojati sę iz with present stem sȯstoji, but also permits skladati sę iz. Preferring the former is a consistency choice, not proof that the latter is wrong. Iz governs the genitive, and reflexive se belongs to the selected expression. German bestehen, however, does not always mean consist of: a statement about where a problem lies or what a method consists in needs a different construction. The record distinguishes 143 changes and seven retained contexts, including five semantic recastings among the changes. The rule preserves composition rather than imposing it on every similar-looking word. Merely adding se does not repair an incorrect present stem.\par
\smallskip\noindent\emph{Izvory:} Slovnik, zaglavja 1874 i 4463; Gramatika, Predlogy, iz s genitivom.\par
\noindent\emph{Sources:} Dictionary entries 1874 and 4463; Grammar, Prepositions, iz with the genitive.\par
\par\medskip\Needspace{8\baselineskip}
\hypertarget{DEC-0015}{}
\noindent\textbf{Obči, obobčiti i soobčenje: razny smysl / Obči, obobčiti and soobčenje: different senses}\quad{\footnotesize\texttt{DEC-0015}}\par
\smallskip\noindent\emph{Razlog:} Obči označaje obče svojstvo ili obseg; obobčiti/obobčati označajut generalizaciju; voobče jest prislovje. Soobčenje može govoriti o informaciji, komunikaciji, publikaciji ili soobčenom priměru, čto potrěbuje razbor německogo konteksta. Zato bylo izbrano 11 opisateljnyh kontekstnyh rěšenij s publikacija, informacija, komunikacija, prědstaviti ili izvěstiti, vměsto edinogo mehaničnogo korenja. Slovnik potvrđaje i obć-, i sȯobć- familije; on ne zapreščaje poslědnju. Razlika obć/obč jest drugy vopros: oficialny pravopis dopuščaje zapis opcionalnogo ć črez osnovno č. Ranějše izdanje upotrěbjalo ć, nynějšny osnovny zapis može dati obč. Eto ne jest novo semantično slovo i ne opravdyvaje nazvati vsako č ošibkoju. Gramatično soglasje i značenje morajut byti prověrjene oddělno od izbora bukvy.\par
\smallskip\noindent\emph{Reason:} Obči expresses a common property or general scope; obobčiti/obobčati express generalization; voobče is adverbial. Soobčenje may concern information, communication, publication or a communicated example, requiring examination of the German context. Eleven contextual solutions therefore use publikacija, informacija, komunikacija, prědstaviti or izvěstiti instead of one mechanical root substitution. The dictionary attests both obć- and sȯobć- families; it does not prohibit the latter. The difference between obć and obč is a separate matter: official orthography permits optional ć to be written as basic č. An earlier edition used ć; the present basic spelling may use obč. This is not a new semantic word and does not make every č erroneous. Agreement and meaning must be checked independently of the letter choice.\par
\smallskip\noindent\emph{Izvory:} Slovnik, zaglavja 1000, 35425, 35426, 35428, 20817, 21169, 21171, 21172, 23055, 4979, 769, 736, 2422, 3688; Pravopis, opcionalne bukvy; Gramatika, Pridavne slova.\par
\noindent\emph{Sources:} Dictionary entries 1000, 35425, 35426, 35428, 20817, 21169, 21171, 21172, 23055, 4979, 769, 736, 2422, 3688; Orthography, optional letters; Grammar, Adjectives.\par
\par\medskip\Needspace{8\baselineskip}
\hypertarget{DEC-0016}{}
\noindent\textbf{Prědstaviti, prědstavjati, prědstavjenje / Present, represent, representation}\quad{\footnotesize\texttt{DEC-0016}}\par
\smallskip\noindent\emph{Razlog:} Prědstav- jest jednolika osnova izbrana iz slovničnyh prědstaviti, prědstavjati i prědstavjeńje. Glagoly razlikujut aspekt; koncovky morajut sootvětstvovati licu, čislu i ulozi participija. Zato predstav- → prědstav- ne izčerpava odluku: prědstavja, prědstavjajut, prědstavjaje i prědstavjene imaju razne funkcije. Slovnično zaglavje prědstavjeńje samo po sobě dava prezentaciju i teatralno prědstavjenje, ne specifikaciju teorije prědstavjenij. Matematično značenje ustanovjaje kontekst korpusa i sravnenje s disciplinarnymi izvorami. Izgovorjena forma ne zaměnjuje definiciju, napr. homomorfizma grupy v linearnu grupu. Odluka zato objedinjaje leksičnu osnovu i gramatiku, ale ne tvrdže, že vsaka nova tehnična upotrěba už jest kodifikovana slovnikom.\par
\smallskip\noindent\emph{Reason:} Prědstav- is the consistent stem selected from dictionary prědstaviti, prědstavjati and prědstavjeńje. The verbs distinguish aspect; endings must match person, number and participial function. Thus predstav- → prědstav- does not exhaust the decision: prědstavja, prědstavjajut, prědstavjaje and prědstavjene have different functions. The noun entry itself gives presentation and theatrical performance, not a specification of representation theory. The technical sense comes from the mathematical context and disciplinary comparison. A pronounceable form does not replace a definition, such as a homomorphism from a group into a linear group. The decision combines a lexical stem with grammar without pretending that every technical extension is already codified in the dictionary.\par
\smallskip\noindent\emph{Izvory:} Slovnik, zaglavja 2422, 1524, 21766, 11264; Gramatika, Glagoly, prezentne osnovy i kratke formy -aje.\par
\noindent\emph{Sources:} Dictionary entries 2422, 1524, 21766, 11264; Grammar, Verbs, present stems and contracted -aje forms.\par
\par\medskip\Needspace{8\baselineskip}
\hypertarget{DEC-0017}{}
\noindent\textbf{Prědstav- v složenyh slovah i koncovkah / Prědstav- in compounds and endings}\quad{\footnotesize\texttt{DEC-0017}}\par
\smallskip\noindent\emph{Razlog:} Jednolika osnova potrěbna jest i v medžuprědstavjenje, ne samo na počętku slova. Formy prědstavljam i prědstavlja byli izravnene s izbranym prědstavjam i prědstavja. Gramatika dopuščaje kratke prezentne koncovky -am, -aš, -a, -amo, -ate pri glagolah s -aje-; množinno tretje lice ostaje -ajut. Zato prědstavjajų → prědstavjajut jest popravka ličnoj koncovky, ne prosto odstranjenje diakritiki. Pet zapisanyh slučajev pokazujut dva razna uslovja: složeno slovo čuva svoj korenj, a lična forma čuva broj i ulogu subjekta. Vse, čto samo slučajno sodrživa podobnu poslědovatelnost bukv, ne pripada avtomatično istoj familiji.\par
\smallskip\noindent\emph{Reason:} The consistent stem is needed inside medžuprědstavjenje, not only at the beginning of a word. Prědstavljam and prědstavlja were aligned with selected prědstavjam and prědstavja. The grammar permits contracted present endings -am, -aš, -a, -amo and -ate in verbs with -aje-, while the third-person plural remains -ajut. Prědstavjajų → prědstavjajut therefore repairs a personal ending, not merely a diacritic. The five recorded cases illustrate two requirements: a compound retains its lexical stem, and a finite form retains its subject's number and grammatical role. A coincidentally similar letter sequence does not automatically belong to the same family.\par
\smallskip\noindent\emph{Izvory:} Slovnik, prědstavjati, 1524; Gramatika, Glagoly, kontrakcija prezentnyh koncovok.\par
\noindent\emph{Sources:} Dictionary entry 1524, prědstavjati; Grammar, Verbs, contraction of present endings.\par
\par\medskip\Needspace{8\baselineskip}
\hypertarget{DEC-0018}{}
\noindent\textbf{Kde, nehaj, kako: funkcija ostaje važna / Kde, nehaj, kako: function still matters}\quad{\footnotesize\texttt{DEC-0018}}\par
\smallskip\noindent\emph{Razlog:} Kde bylo upotrěbjeno vměsto gdje, nehaj vměsto neha pri želanju ili postavjenju, i kako vměsto kak v sootvětnyh poredbnyh i opisateljnyh konstrukcijah. Prva zaměna čuva městno odnošenje; vtora čuva želateljnu ili postavljajuču funkciju; tretja potrebuje kontekst načina ili sravnenja. Eto tri različne gramatične odluke, ne jedna nacionalna bukvenna zaměna. Pravilo nehaj ne vključaje neoprěděljeno něktory; pravilo kako ne menja citaty, imena ili matematične identifikatory. Zapisanyh 136 pojavjenij daje obseg priměnjenja. On ne dokazuje, že kratke ili nacionalne formy ne mogu byti razumljive.\par
\smallskip\noindent\emph{Reason:} Kde replaced gdje, nehaj replaced neha in wishes or stipulations, and kako replaced kak in the relevant comparative and descriptive constructions. The first preserves a spatial relation, the second a jussive or stipulative function, and the third requires a manner or comparison context. These are three grammatical decisions, not one national letter substitution. The nehaj rule excludes indefinite něktory, while kako does not alter quotations, names or mathematical identifiers. The recorded 136 occurrences describe application scope, not proof that shorter or national forms cannot be understood.\par
\smallskip\noindent\emph{Izvory:} Slovnik, kde 942 i nehaj 9558; Gramatika, Spojky, kako … tako i.\par
\noindent\emph{Sources:} Dictionary entries kde 942 and nehaj 9558; Grammar, Conjunctions, kako … tako i.\par
\par\medskip\Needspace{8\baselineskip}
\hypertarget{DEC-0019}{}
\noindent\textbf{Prědstavjenij; tym jest dokazano / Prědstavjenij; thus it is proved}\quad{\footnotesize\texttt{DEC-0019}}\par
\smallskip\noindent\emph{Razlog:} Prědstavjenij jest prijety genitiv množiny za prědstavjenje; prědstavjenja jest genitiv jedinstvenogo čisla. Arhiv pripisyvaje prvu formu oficialnomu generatoru i zapisuje 196 kontekstov, ale sam reproducirujemy zapusk togo generatora ne jest najdeny pri etoj rekonstrukciji. Obča gramatična stranica, ktoru bylo možno pročitati, govori o nulevoj koncovkě neutra i epentetičnom e/o; ona sama ne dokazuje točnu formu -nij. Zato prědstavjenij jest jasno zabeležena konvencija s nepolnym morfologičnym dokazom, ne tvrdo ustanovjena ošibka vsakoj alternativy. Oddělno, prědstava v slovniku znači mentalny obraz; matematično prědstavjenje v naslově potrěbuje drugy tehničny smysl. Zaključna fraza tym jest dokazano izražaje dovršenost dokaza vměsto nejasnogo čto bylo dokazati. Te tri odluke imajut različne osnovy i različnu silu dokazov.\par
\smallskip\noindent\emph{Reason:} Prědstavjenij is the adopted genitive plural of prědstavjenje; prědstavjenja is the genitive singular. The archive attributes the plural to an official generator and records 196 contexts, but a reproducible execution of that generator has not been recovered in this reconstruction. The general grammar page that could be consulted describes a zero neuter ending with epenthetic e/o; it does not itself establish the exact -nij form. Prědstavjenij is therefore documented as an adopted convention with incomplete morphological evidence, not as conclusive proof that every alternative is erroneous. Separately, dictionary prědstava denotes a mental image; mathematical representation in a title requires the different technical sense. The proof ending tym jest dokazano asserts a completed proof instead of the unclear čto bylo dokazati. These three decisions rest on different grounds and carry different evidential strength.\par
\smallskip\noindent\emph{Izvory:} Slovnik, zaglavja 11264 i 21769; Jan van Steenbergen, Gramatika, Imenovniki, sklanjanje srednjego roda; arhivny rezultat generatora ostaje neopakovany.\par
\noindent\emph{Sources:} Dictionary entries 11264 and 21769; Jan van Steenbergen, Grammar, Nouns, neuter declension; the archived generator result has not been independently reproduced.\par
\par\medskip\Needspace{8\baselineskip}
\hypertarget{DEC-0020}{}
\noindent\textbf{Čto v pytanju; sut v prozě / Čto in questions; sut in prose}\quad{\footnotesize\texttt{DEC-0020}}\par
\smallskip\noindent\emph{Razlog:} Čto jest slovnično pytateljno slovo, napr. pytanje, čto možno rěkti. Ono ne potrěbuje odstranjenja iz vsakej věty. Ale posledstvo, odnošena klauza, ravnoznačnost, parentetično prědpoloženje i svęz diskursa ne vyražajut jednu i tu že relaciju: po kontekstu byly izbrany to, ktoro, tym, pri čem i ravnoznačno. Avtomatično čto → že by skrylo ty različenja. Odluka vključaje 64 takovyh kontekstov i jedno sohranjeno pytanje; 23 zaključne dokazne frazy obgovorjene sut oddělno. Podobno, sut jest množinna kopula v prozě, ale su ili Su može byti matematičny znak. Zapis razlišaje pet kopul i 245 sohranjenyh znakovyh pojavjenij. Gramatika teksta ne dava pravo proměnjati avtorsku notaciju.\par
\smallskip\noindent\emph{Reason:} Čto is a dictionary interrogative, as in the question what can be said. It need not be removed from every sentence. But consequence, a relative clause, equivalence, a parenthetical assumption and discourse linkage do not express the same relation: contextual choices include to, ktoro, tym, pri čem and ravnoznačno. Automatic čto → že would conceal those distinctions. The decision covers 64 such contexts and one retained question; 23 proof-ending formulas are discussed separately. Similarly, sut is a plural copula in prose, while su or Su may be a mathematical symbol. The record distinguishes five copulas from 245 preserved symbolic occurrences. The grammar of the exposition does not authorize changing the author's notation.\par
\smallskip\noindent\emph{Izvory:} Slovnik, čto 19670; Gramatika, Glagoly, formy byti; oddělno zaglavje Tym jest dokazano.\par
\noindent\emph{Sources:} Dictionary entry čto 19670; Grammar, Verbs, forms of byti; see the separate entry Tym jest dokazano.\par
\par\medskip\Needspace{8\baselineskip}
\hypertarget{DEC-0021}{}
\noindent\textbf{Odpovědati, odgovarjajuči, korespondencija, neuspěh / Correspond, corresponding, correspondence, failure}\quad{\footnotesize\texttt{DEC-0021}}\par
\smallskip\noindent\emph{Razlog:} Odpovědati može značiti odpovědati na pytanje ili sootvětstvovati něčemu; slovnik dava odvojene zaglavja, pri vtorom s dativom. Zato 182 glagola v smislu korespondencije ne byli zaměnjene samo zaradi podobnosti s odpovědju. Modifikatory dobyli formu odgovarjajuči, a imenovniki otnošenja korespondencija. V tehničnom kontekstě to znači zadano otnošenje medžu objektami, ne poštovu korespondenciju; sam obči slovnik ne specifikuje vsaku takovu strukturu. Tri pasaže s německym Versagen govoriat o neuspěhu vlastnyh energijskyh zakonov, ne o ih sootvětnosti. Eto semantična popravka, ne stilistična sinonimija. Dva diskursne slučaje vyražajut Analogično. Polny izbor tak razděljuje odpověď, korespondenciju, analogiju i neuspěh, vměsto mehaničnogo izravnjenja cěloj familije.\par
\smallskip\noindent\emph{Reason:} Odpovědati can mean answer a question or correspond to something; the dictionary gives separate entries, the latter with the dative. The 182 correspondence verbs were therefore not replaced merely because they resemble answer vocabulary. Modifiers use odgovarjajuči, while relation nouns use korespondencija. In a technical context the noun denotes a specified relation between objects, not postal correspondence; the general dictionary does not define every such structure. Three passages with German Versagen concern the failure of proper energy laws, not their correspondence. That is a semantic correction, not stylistic synonymy. Two discourse contexts instead use Analogično. The complete choice distinguishes answer, correspondence, analogy and failure rather than mechanically leveling the entire family.\par
\smallskip\noindent\emph{Izvory:} Medžuslovjansky slovnik, zaglavja 2189, 3161, 35938, 6063, 6550.\par
\noindent\emph{Sources:} Interslavic Dictionary, entries 2189, 3161, 35938, 6063, 6550.\par
\par\medskip\Needspace{8\baselineskip}
\hypertarget{DEC-0022}{}
\noindent\textbf{Dobyti i dobyvati; dobyty rezultat / Obtain: aspect and the resulting object}\quad{\footnotesize\texttt{DEC-0022}}\par
\smallskip\noindent\emph{Razlog:} Dobyti jest dokonany, dobyvati nedokonany glagol dlja dobyvanja rezultata. Slovnik dava takove značenje i srodne nacionalne formy, ale ih poznavajemost ne byla izměrena čitateljskim testom. Získati bylo odbačeno kako druga nacionalno ograničena prednost; dostati ostaje slovničny comparator, ne zakazano slovo. Vosem kontekstov ne dopuščalo jedinu koncovku: budučna forma dobudemo, prezentny povratny predikat i prošlo dobil imajut razne funkcije. Dva participne konteksty byli jasnějše izražene prěmoj větoju ili ukazanjem na razloženje iz punkta 2. Prva zabeležena forma dobiva- i nynějšna dobyva- takože potrěbujut različenje pravopisnogo sloja; ne trěba potajno vydavati istoričnu formu za nynějšnu. Objekty, dokaz i vremenny smysl ostajut te že.\par
\smallskip\noindent\emph{Reason:} Dobyti is perfective and dobyvati imperfective for obtaining a result. The dictionary supplies that meaning and related national forms, but reader recognition was not measured. Získati was rejected as shifting the preference to another national group; dostati remains a dictionary comparator, not a forbidden word. The eight contexts did not permit one ending: future dobudemo, a present reflexive predicate and past dobil serve different functions. Two participial contexts were clearer as direct clauses or a reference to the decomposition in item 2. Recorded dobiva- and current dobyva- also require an orthographic-layer distinction; a historical form must not silently be presented as current. The objects, argument and temporal meaning remain unchanged.\par
\smallskip\noindent\emph{Izvory:} Medžuslovjansky slovnik, zaglavja 3192, 3191, 1900.\par
\noindent\emph{Sources:} Interslavic Dictionary, entries 3192, 3191, 1900.\par
\par\medskip\Needspace{8\baselineskip}
\hypertarget{DEC-0023}{}
\noindent\textbf{Pytanje; vopros i zapytanje / Question: pytanje, vopros and zapytanje}\quad{\footnotesize\texttt{DEC-0023}}\par
\smallskip\noindent\emph{Razlog:} Pytanje jest izbrano dlja pytanja v matematičnom izlaganju. V konsultovanoj kopiji pytańje ima tip 1, vȯpros tip 3; zapytańje takože jest tip 1. Izbor korotkogo pytanje usoglasnjaje tekst s už produktivnoj familijej i ima srodne formy v poljskom, ukrajinskom i BCS, s glagolnymi sravnenjami v zapadnyh jazykah. To ne znači, že vsaka nacija upotrěbjaje isto slovo: češsko otázka i slovensko vprašanje pokazujut razliku. Tip jest slovnična klasifikacija, ne procent razuměnja. V devetih slučajah pytanje, pytanja, pytanjah i pytanjam čuvajut padež i funkciju. Vopros ostaje dokumentovany comparator; ne jest predstavjen kako neobstajuče slovo.\par
\smallskip\noindent\emph{Reason:} Pytanje is the selected noun for a question in mathematical exposition. In the consulted copy pytańje is type 1 and vȯpros type 3; zapytańje is also type 1. The shorter pytanje aligns the text with an already productive family and has cognate forms in Polish, Ukrainian and BCS, with related verbs in western languages. This does not mean every language uses the same noun: Czech otázka and Slovene vprašanje illustrate the difference. A type is a dictionary classification, not a percentage of comprehension. Across nine cases, pytanje, pytanja, pytanjah and pytanjam preserve case and function. Vopros remains a documented comparator, not a supposedly nonexistent word.\par
\smallskip\noindent\emph{Izvory:} Medžuslovjansky slovnik, zaglavja 5768, 554, 4766.\par
\noindent\emph{Sources:} Interslavic Dictionary, entries 5768, 554, 4766.\par
\par\medskip\Needspace{8\baselineskip}
\hypertarget{DEC-0024}{}
\noindent\textbf{Sodrživati; sodržanje; držati / Contain, content and hold}\quad{\footnotesize\texttt{DEC-0024}}\par
\smallskip\noindent\emph{Razlog:} Sodrživati izražaje matematično sodržanje, a sodržanje sodržaj teksta ili objekta. Obči slovnik razlišaje sȯdŕživati contain od sȯdŕžati hold back/inhibit. Zato korotši glagol ne jest avtomatično aspektny par s tym že značenjem. Držati i zadržati mogut vyražati držanje ili sohranjenje i ne pripadajut toj že zaměně. Slovnična osnova s prezentnymi pravilami daje sodrživaje, sodrživajut i odgovarjajuče participy. V matematičnoj upotrěbě ostaje potrěbno čitati definiciju, da se ne směšajut členstvo i vključanje. Jedna fraza o soglasju po smislu potrěbovala smyslovo soglasje, ne sodržanje. Zapis razděljuje 387 oblikovyh proměn, jednu takovu frazu, devet sohranjenyh imenovnikov i šestnadset odlišnyh glagolov.\par
\smallskip\noindent\emph{Reason:} Sodrživati expresses mathematical containment, and sodržanje the content of a text or object. The general dictionary distinguishes sȯdŕživati, contain, from sȯdŕžati, hold back or inhibit. The shorter verb is therefore not automatically an aspectual counterpart with the same meaning. Držati and zadržati may mean holding or retaining and do not belong to the same substitution. The dictionary stem and present-tense rules yield sodrživaje, sodrživajut and corresponding participles. Mathematical use still requires the definition to distinguish membership from inclusion. One phrase about agreement in meaning required smyslovo soglasje rather than containment. The record separates 387 inflectional changes, that one phrase, nine retained nouns and sixteen distinct verbs.\par
\smallskip\noindent\emph{Izvory:} Medžuslovjansky slovnik, zaglavja 2163, 3942, 72, 2617; Gramatika, Glagoly.\par
\noindent\emph{Sources:} Interslavic Dictionary, entries 2163, 3942, 72, 2617; Grammar, Verbs.\par
\par\medskip\Needspace{8\baselineskip}
\hypertarget{DEC-0025}{}
\noindent\textbf{Obstajati, suči, sučstveny, sučnost / Exist, present, essential and essence}\quad{\footnotesize\texttt{DEC-0025}}\par
\smallskip\noindent\emph{Razlog:} Suščestv-/sušt- ne pokryvaje jedno matematično značenje. Izbrana razdělenja sut obstajati dlja obstajanja, suči dlja prěsutnogo ili nynějšnogo, sučstveny dlja sučstvenosti i v sučnosti dlja sučnosti kako pojma. Napriměr O sučej literaturi govori o dostupnoj literaturi, ne o dokazě eksistencije literatury. Slovnik ima obstajati tipa 1, ale jego profil sodrživa negativne i približne ocenky; ocenka tipa ne dokazuje všeobču prednost nad sųtstvovati. Sųćstveny i sųćnosť takože imajut različne profile i značenja. Osnovny pravopis daje č vměsto opcionalnogo ć. Soglasje, prislovje na -o i stepenovanje razbirajut se po gramatikě. Izbor jest različenje smyslov s opisanym kompromisom, ne odstranjenje vsego, čto pripominaje rusky jazyk.\par
\smallskip\noindent\emph{Reason:} The suščestv-/sušt- family does not cover just one mathematical meaning. The selected distinctions are obstajati for existence, suči for present or current, sučstveny for essential, and v sučnosti for essence as a concept. For example, O sučej literaturi concerns available literature, not a proof of the literature's existence. Obstajati is dictionary type 1, but its profile includes negative and approximate assessments; the type does not establish a universal advantage over sųtstvovati. Sųćstveny and sųćnosť also have different profiles and senses. Basic orthography writes č for optional ć. Agreement, adverbs in -o and comparison follow the grammar. This is semantic differentiation with a stated compromise, not the removal of everything reminiscent of Russian.\par
\smallskip\noindent\emph{Izvory:} Medžuslovjansky slovnik, zaglavja 9895, 17048, 20193, 20192, 36616, 36821, 36822; Gramatika, Pridavne slova; Pravopis, opcionalne bukvy.\par
\noindent\emph{Sources:} Interslavic Dictionary, entries 9895, 17048, 20193, 20192, 36616, 36821, 36822; Grammar, Adjectives; Orthography, optional letters.\par
\par\medskip\Needspace{8\baselineskip}
\hypertarget{DEC-0026}{}
\noindent\textbf{Sravniti; identifikovati; prirěditi / Compare, identify and assign}\quad{\footnotesize\texttt{DEC-0026}}\par
\smallskip\noindent\emph{Razlog:} Sravniti/sravnjati i sravnjenje sut izbrana familija dlja sravnenja. Zapadne porówn-/porovn-/srovn- i južne primerj-/uspored-/spored- ne davajut očividnogo jednogo kandidata; arhiv zato upotrěbjaje slovnično zaglavje kako redakcijnu oporu. To ne jest opublikovana čislena poběda: polna tablica vagov dlja etoj odluki ne jest zachovana. Četyri konteksty potrěbujut identifikovati ili prirěditi, zato že identifikacija elementov i zadanje sootvětnosti ne sut isto kako ih sravnenje. Sravni izražaje čitlivu citatnu instrukciju. Ravn-/rovn- v ravnosti, uravnenijah, ploskosti ili ravnověsju ne mogu byti popravjene kak by pripadali vsegda istoj familiji; 589 takyh tokenov byli odlišene v zapisě.\par
\smallskip\noindent\emph{Reason:} Sravniti/sravnjati and sravnjenje are the selected comparison family. Western porówn-/porovn-/srovn- and southern primerj-/uspored-/spored- do not supply one obvious candidate; the record therefore uses the dictionary headword as an editorial anchor. This is not a published numerical victory: a complete candidate-weight table has not been preserved for this decision. Four contexts require identifikovati or prirěditi because identifying elements and assigning a correspondence are not the same as comparing them. Sravni provides an explicit cross-reference instruction. Ravn-/rovn- in equality, equations, flatness or balance cannot be treated as the same family; the record distinguishes 589 such tokens.\par
\smallskip\noindent\emph{Izvory:} Medžuslovjansky slovnik, zaglavja 20168, 35220, 1142, 679, 3510, 15906.\par
\noindent\emph{Sources:} Interslavic Dictionary, entries 20168, 35220, 1142, 679, 3510, 15906.\par
\par\medskip\Needspace{8\baselineskip}
\hypertarget{DEC-0027}{}
\noindent\textbf{Sběgati se, identičnost i soglasnost / Coincide, identity and agreement}\quad{\footnotesize\texttt{DEC-0027}}\par
\smallskip\noindent\emph{Razlog:} Sběgati se jest izbrany predikat za objektov, ktore sběgajut se v matematičnom smislu. Pri imenovnikah trěba vyraziti, jest li rěč o identičnosti objektov ili o soglasnosti definicij. Zapis dava 88 predikatov, trinadset imenovnikov identičnosti i dva naslova s soglasnostju. Slovnično sȯvpadeńje ima glosu coincidence i sravnenja s náhoda/shoda okolností; to motivuje izběganje slučajnostnogo čitanja, no ne dokazuje, že glosa coincidence nikogda ne može iměti matematičny smysl. Izbrano sběgati se takože ima glosu converge: bez konteksta ono može vyzvati drugu asociaciju. V toh pasažah ne jest govora o graničnoj shodimosti ili geometričnom prěkrivanju. Prednost prědstavjaje kontrolovany kontekstny izbor, ne zasvědčenu bezdvojznačnost vsih form.\par
\smallskip\noindent\emph{Reason:} Sběgati se is the selected predicate for objects coinciding mathematically. The nouns must distinguish identity of objects from agreement between definitions. The record gives 88 predicates, thirteen identity nouns and two agreement headings. Dictionary sȯvpadeńje is glossed coincidence and compared with náhoda/shoda okolností; this motivates avoiding a chance reading, but does not prove that coincidence can never have a mathematical sense. Selected sběgati se is also glossed converge and may evoke another association without context. These passages concern neither limiting convergence nor geometric overlap. The preference is a controlled contextual choice, not demonstrated ambiguity-free wording.\par
\smallskip\noindent\emph{Izvory:} Medžuslovjansky slovnik, zaglavja 18677, 35748, 6301, 35382.\par
\noindent\emph{Sources:} Interslavic Dictionary, entries 18677, 35748, 6301, 35382.\par
\par\medskip\Needspace{8\baselineskip}
\hypertarget{DEC-0028}{}
\noindent\textbf{Sodrživati vměsto obsah- i vsebu- / Sodrživati alongside national containment terms}\quad{\footnotesize\texttt{DEC-0028}}\par
\smallskip\noindent\emph{Razlog:} Češko-slovačsko obsah- i slovensko vsebu- označajut sodržanje v svojih jazykah. Slovnično zaglavje sȯdŕživati izrično dava obsahovat i vsebovati kako nacionalne sravnenja. V medžuslovjanskom korpusě ih střědanje po rabotah ne predstavjalo posebno matematično različenje, zato 97 predikatov byli usoglasnjene s izbranym sodrživati. To rozširjaje isto pravilo, ktoro odlišaje sodržanje od držanja i sadržaja; ne izmysljaje novo značenje. Izkvarjene množinne formy vsebujeju i vsebijejut potrěbovaly pravilne lične koncovky. Imati moglo byti širše, ale ne izrično by vyrazilo odnošenje sodržanja; vključati trěbuje prověrku konkretnogo smisla, ne kategoričnu zabranu. Nacionalne slova ne sut ošibky v nacionalnyh izvorah.\par
\smallskip\noindent\emph{Reason:} Czech–Slovak obsah- and Slovene vsebu- express containment in their own languages. The dictionary entry sȯdŕživati explicitly gives obsahovat and vsebovati as national comparators. Their alternation by paper in the Interslavic corpus did not mark a distinct mathematical concept, so 97 predicates were aligned with selected sodrživati. This extends the same distinction between containment, holding and content rather than inventing a new meaning. Malformed plurals vsebujeju and vsebijejut required appropriate personal endings. Imati could be broader but would not explicitly name containment; vključati requires examination of the particular sense, not categorical prohibition. The national words are not errors in their national sources.\par
\smallskip\noindent\emph{Izvory:} Medžuslovjansky slovnik, zaglavja 2163.\par
\noindent\emph{Sources:} Interslavic Dictionary, entries 2163.\par
\par\medskip\Needspace{8\baselineskip}
\hypertarget{DEC-0029}{}
\noindent\textbf{Vzajemny, vzajemno, vzajemnost / Mutual, mutually, mutuality}\quad{\footnotesize\texttt{DEC-0029}}\par
\smallskip\noindent\emph{Razlog:} Vzajem- jest osnovny zapis slovničnogo vzajėm-. Vzajemny modifikuje imenovnik, vzajemno jest prislovje, a vzajemnost nazvanje otnošenja. Zapisany vzaim- ne čini drugu teoriju reciprocity; zaměna čuva padež, broj, kapitalizaciju i sintaksu v deveti formah. Nacionalne sravnenja uključajut poljsko wzajemny, češko vzájemný, rusko vzaimnyj i hrvatsko uzajaman. Jih srodnost obosnovaje izbor familije, no ne dava identičny pravopis vo vsih jazykah. Internacionalno reciprocal ne bylo potrěbno vměsto uže obosnovanogo korenja. Kontekst mora ješče razlikovati vzajemno otnošenje od drugyh tehničnyh značenij reciprociteta; samo pravopisno pravilo ne definuje teorem reciprocity.\par
\smallskip\noindent\emph{Reason:} Vzajem- is the basic spelling of dictionary vzajėm-. Vzajemny modifies a noun, vzajemno is adverbial, and vzajemnost names a relation. Recorded vzaim- did not encode a different theory of reciprocity; replacement preserves case, number, capitalization and syntax across nine forms. Comparators include Polish wzajemny, Czech vzájemný, Russian vzaimnyj and Croatian uzajaman. Their relationship motivates the family without implying identical spelling across languages. An international reciprocal term was unnecessary in place of an already supported root. Context must still distinguish a mutual relation from specialized senses of reciprocity; the spelling rule alone does not define a reciprocity theorem.\par
\smallskip\noindent\emph{Izvory:} Medžuslovjansky slovnik, zaglavja 13986, 23052, 1804.\par
\noindent\emph{Sources:} Interslavic Dictionary, entries 13986, 23052, 1804.\par
\par\medskip\Needspace{8\baselineskip}
\hypertarget{DEC-0030}{}
\noindent\textbf{Razgledati i razvažati / Examine and reflect}\quad{\footnotesize\texttt{DEC-0030}}\par
\smallskip\noindent\emph{Razlog:} Razgledati/razgleděti označajut razgled objekta ili proučenje; razvažati označaje razmyšljenje i obdumavanje. Odluka ne razdělja vsako německo slovo avtomatično: v šestnadsetih zapisanyh kontekstah Überlegung bylo razuměno kako razvaženje, a mnogo drugih Betrachtung ili in Betracht kommen kako razgled. Te smysly častično prěkrivajut se, zato izbrany imenovnik potrebuje cělu větu. Slovnik razlišaje råzględati tipa 2, råzvažati tipa 1 i råzsmatrjati tipa 3; takova klasifikacija ne dava pravilo, že najnižše čislo može zaměniti vse druge smysly. Fraza bez vsakogo razgleda napravi soglasje genitiva po bez, vměsto neobosnovanogo ženskogo razgledky. Rěč jest o razboru i padežu, ne prosto o odstranjanju s iz korenja.\par
\smallskip\noindent\emph{Reason:} Razgledati/razgleděti express examining or considering an object; razvažati expresses reflection or deliberation. The choice does not mechanically assign each German word: sixteen recorded Überlegung contexts were understood as deliberation, while many Betrachtung or in Betracht kommen contexts involved examination or consideration. These senses overlap, so the selected noun needs its whole sentence. The dictionary distinguishes type-2 råzględati, type-1 råzvažati and type-3 råzsmatrjati; that classification does not allow the lowest number to replace all other senses. Bez vsakogo razgleda gives genitive agreement after bez instead of unsupported feminine razgledky. This concerns semantic analysis and case, not merely deleting s from a root.\par
\smallskip\noindent\emph{Izvory:} Medžuslovjansky slovnik, zaglavja 11776, 11777, 6010, 34013.\par
\noindent\emph{Sources:} Interslavic Dictionary, entries 11776, 11777, 6010, 34013.\par
\par\medskip\Needspace{8\baselineskip}
\hypertarget{DEC-0031}{}
\noindent\textbf{Izpolniti, izpolnjati, izpolnjenje / Fulfil, carry out, fulfilment}\quad{\footnotesize\texttt{DEC-0031}}\par
\smallskip\noindent\emph{Razlog:} Izpoln- jest osnovny zapis slovničnogo izpȯln-, ne ispoln-. Familija ima dokonany izpolniti, nedokonany izpolnjati i imenovnik izpolnjenje. V kontekstah uslovij ili operacij govori o izpolnjenju ili izvršenju; srodne nacionalne slova mogut iměti druge zvukove realizacije. Zapis sodrživa točno 56 popravenyh pojavjenij i 129 uže pravilnyh form. Čisla dokazujut tolko obseg zapisanogo usoglasnjenja. Alternativy izvršiti, realizovati i završiti nebyly potrěbne, kogda smysl i sintaksa uže sootvětstvovali izpoln-. Drugy kontekst, napr. završiti dokaz, ne mora upotrěbiti isto slovo samo zaradi edinogo stilja.\par
\smallskip\noindent\emph{Reason:} Izpoln- is the basic spelling of dictionary izpȯln-, not ispoln-. The family has perfective izpolniti, imperfective izpolnjati and noun izpolnjenje. In contexts of conditions or operations it expresses fulfilment or execution; cognate national forms may realize the sounds differently. The record contains exactly 56 changed occurrences and 129 already matching forms; those numbers describe only the recorded normalization. Alternatives izvršiti, realizovati and završiti were unnecessary where meaning and syntax already fitted izpoln-. A different context, such as finishing a proof, need not use the same word merely for stylistic uniformity.\par
\smallskip\noindent\emph{Izvory:} Medžuslovjansky slovnik, zaglavja 597, 2501, 36609; Gramatika, Glagoly.\par
\noindent\emph{Sources:} Interslavic Dictionary, entries 597, 2501, 36609; Grammar, Verbs.\par
\par\medskip\Needspace{8\baselineskip}
\hypertarget{DEC-0032}{}
\noindent\textbf{Proizhoditi iz, slědovati, nastavati, početny / Origin, consequence, arising and initial state}\quad{\footnotesize\texttt{DEC-0032}}\par
\smallskip\noindent\emph{Razlog:} Proizhoditi iz izražaje istinny izvor ili poreklo, napr. pripisanje rezultata avtoru. Iz togo ne slěduje, že vsako německo ausgehen, hervorgehen ili Ausgangs- znači poreklo. Logično posledstvo potrěbuje slěduje; nastavanje objekta nastavaje/nastavajut; početna veličina početny; odpravna točka odpravny punkt/točka. Dokaz, izvedenje, členstvo, bazis, vyhodeči luč i buduča publikacija potrěbujut svoje kontekstne konstrukcije. Slovnik potvrđaje proizhoditi iz, izhoditi v smislu vyhoda, početkovy i izvesti/izvoditi. Ego osnovne glosy ne sut iscrpny opis vsej matematičnoj upotrěby: važna jest logična relacija v konkretnoj věti. Iz 114 zapisanyh slučajev petnadset čuvajut istinny smysl porekla. Opis ne pretvarjaje vsaku posledicu v isto obče iz-odnošenje.\par
\smallskip\noindent\emph{Reason:} Proizhoditi iz expresses genuine source or origin, including attribution of a result to an author. It does not follow that every German ausgehen, hervorgehen or Ausgangs- denotes origin. A logical consequence requires slěduje; an object arising, nastavaje/nastavajut; an initial quantity, početny; a starting point, odpravny punkt/točka. Proof, derivation, membership, a basis, an emanating ray and forthcoming publication require their own constructions. The dictionary attests proizhoditi iz, izhoditi for exiting, početkovy and izvesti/izvoditi. These basic glosses do not exhaust mathematical usage: the logical relation in the particular sentence matters. Fifteen of the 114 recorded cases retain genuine origin. The account does not flatten every consequence into one generic from-relation.\par
\smallskip\noindent\emph{Izvory:} Medžuslovjansky slovnik, zaglavja 11566, 5355, 29170, 10670, 10669, 36764, 36763.\par
\noindent\emph{Sources:} Interslavic Dictionary, entries 11566, 5355, 29170, 10670, 10669, 36764, 36763.\par
\par\medskip\Needspace{8\baselineskip}
\hypertarget{DEC-0033}{}
\noindent\textbf{Vsegda / Always}\quad{\footnotesize\texttt{DEC-0033}}\par
\smallskip\noindent\emph{Razlog:} Vsegda bylo izbrano kako jednolika forma dlja always. Slovnik dava tak vsegda, kako vsegdy; nijedna ne jest jedina dozvoljena koncovka. Sravnenja s ruskim vsegda, češskim vždy, slovenskym vedno i hrvatskim svagda/uvijek pokazuju srodnost i takože različenje. Profil sodrživa približne ocenky za několik jazykov, ne vseslovjansku identičnost. Trideset pet zapisanyh vsagda ne izražajut poseben časovy ili logičny smysl, zato usoglasnjenje ne měnjaje kvantifikaciju. V matematičnom tvrdženju vsegda vse ravno zavisi od zadanyh uslovij i obsega; slovnično zaglavje ne daje pravo proměniti ograničeno tvrdženje v bezuslovno.\par
\smallskip\noindent\emph{Reason:} Vsegda was selected as the consistent form for always. The dictionary gives both vsegda and vsegdy; neither is the only permitted ending. Russian vsegda, Czech vždy, Slovene vedno and Croatian svagda/uvijek show both kinship and difference. The profile includes approximate assessments for several languages, not pan-Slavic identity. The 35 recorded vsagda forms did not express a separate temporal or logical sense, so alignment leaves quantification unchanged. In a mathematical statement always remains governed by the stated hypotheses and domain; a dictionary headword does not license turning a conditional claim into an unconditional one.\par
\smallskip\noindent\emph{Izvory:} Medžuslovjansky slovnik, zaglavja 780.\par
\noindent\emph{Sources:} Interslavic Dictionary, entries 780.\par
\par\medskip\Needspace{8\baselineskip}
\hypertarget{DEC-0034}{}
\noindent\textbf{Slědujuči i slědny / Following and next}\quad{\footnotesize\texttt{DEC-0034}}\par
\smallskip\noindent\emph{Razlog:} Slědujuči jest izbrany modifikator za slědujuči objekt ili daljše tvrdženje; slědny jest drugy slovničny pridavnik. Ih razlika ne dopuščaje zaměnu bez prověrky sintaksisa. Ranějše pravilo izbralo slědujući s ć i izravnilo slědeč- i slědujuč-; no nynějšny osnovny pravopis dopuščaje slědujuči s č. Zato argument, že č samo po sobě jest nepravilen, ne važi za drugy pravopisny sloj. Stvarna gramatična popravka slědujučnoj → slědujućej, odnosno slědujučej, tiče se mekkoj pridavnoj osnovy i soglasja, ne samo bukvy. Jasno različenje pravopisnoj konvencije i padežnoj koncovky brani slovnik od predstavljanja supersedovanoj grafije kako jedine normy.\par
\smallskip\noindent\emph{Reason:} Slědujuči is the selected modifier for a following object or subsequent assertion; slědny is another dictionary adjective. The distinction does not permit substitution without checking syntax. An earlier policy chose slědujući with ć and aligned slědeč- and slědujuč-, but current basic orthography permits slědujuči with č. The argument that č is inherently wrong therefore does not apply across orthographic layers. The substantive repair slědujučnoj → slědujućej, or basic slědujučej, concerns a soft adjectival stem and agreement, not merely a letter. Separating orthographic policy from the case ending prevents a superseded spelling from being presented as the only norm.\par
\smallskip\noindent\emph{Izvory:} Medžuslovjansky slovnik, zaglavja 1378, 12301; Gramatika, Pridavne slova; Pravopis, opcionalne bukvy.\par
\noindent\emph{Sources:} Interslavic Dictionary, entries 1378, 12301; Grammar, Adjectives; Orthography, optional letters.\par
\par\medskip\Needspace{8\baselineskip}
\hypertarget{DEC-0035}{}
\noindent\textbf{Toj, togo; sej i tutoj / Demonstratives: ordinary reference and explicit contrast}\quad{\footnotesize\texttt{DEC-0035}}\par
\smallskip\noindent\emph{Razlog:} Toj jest osnovny pokazateljny zaimek v konsultovanoj gramatikě, kogda ne potrěbujemo izrično razlikovati blizky i daleky objekt. Tutoj i tamtoj služat toj razlikě. Sej jest istorično obosnovany, no gramatika upozornjaje, že njegova forma može byti nejasna govoritěljam mnogih nynějšnyh jazykov. To jest kvalitativno upozornjenje, ne rezultat ispytanja etyh deseti vět. Devet kontekstov dobylo toj, ale genitiv mužskogo roda pri modulu potrěboval togo. Město sklonjenogo slova ne može postaviti nominativny slovničny zaglavok. Ne potrěbna proximalna emfaza tutoj ne byla dodana. Avtorsko ukazanje na točno oprěděljenu stvar, ne etymologičny ukus, vladaje izborom.\par
\smallskip\noindent\emph{Reason:} Toj is the basic demonstrative in the consulted grammar where explicit near–far contrast is unnecessary. Tutoj and tamtoj express that contrast. Sej has historical grounds, but the grammar warns that it may be unclear to speakers of many modern languages. This is a qualitative warning, not a test of these ten sentences. Nine contexts received toj, while the masculine genitive modifying a module required togo. An inflected word cannot simply be replaced by the nominative headword. Unneeded proximal emphasis through tutoj was not added. The author's reference to a particular object, rather than etymological preference, governs the choice.\par
\smallskip\noindent\emph{Izvory:} Jan van Steenbergen, Gramatika, Zaimena, pokazateljne zaimena.\par
\noindent\emph{Sources:} Jan van Steenbergen, Grammar, Pronouns, demonstratives.\par
\par\medskip\Needspace{8\baselineskip}
\hypertarget{DEC-0036}{}
\noindent\textbf{Poslědnji, poslědny / Last}\quad{\footnotesize\texttt{DEC-0036}}\par
\smallskip\noindent\emph{Razlog:} Poslědnji i poslědny sut dva pridavne oblika v slovniku. Redakcija čuva ě v osnově poslědn-, vměsto střědanja s posledn-. Petnadset razlišnyh kontekstov ne davajut poseben smysl formě bez ě; padež, rod, broj i velka bukva ostajut odvisne od věty. Zadnji ili ostatni byli nacionalne alternativy, ne potrěbne v uže ustalenej familiji. Slovničny profil ima negativnu bělorusku i približne poljsku i slovensku ocenku: široka srodnost ne jest soglasje vsih jazykov. Eto pravopisna konsistencija v izbranej sistemě, ne novo nazvanje matematičnogo objekta.\par
\smallskip\noindent\emph{Reason:} Poslědnji and poslědny are both dictionary adjective forms. The edition keeps ě in poslědn- instead of alternating with posledn-. The fifteen recorded contexts supply no separate sense for the spelling without ě; case, gender, number and capitalization remain sentence-dependent. National alternatives zadnji or ostatni were unnecessary within the established family. The dictionary profile has a negative Belarusian assessment and approximate Polish and Slovene assessments: broad kinship is not unanimity. This is spelling consistency within the chosen system, not a new mathematical term.\par
\smallskip\noindent\emph{Izvory:} Medžuslovjansky slovnik, zaglavja 2323.\par
\noindent\emph{Sources:} Interslavic Dictionary, entries 2323.\par
\par\medskip\Needspace{8\baselineskip}
\hypertarget{DEC-0037}{}
\noindent\textbf{Cělo čislo; cěločislny / Integer; integer-valued or integral}\quad{\footnotesize\texttt{DEC-0037}}\par
\smallskip\noindent\emph{Razlog:} Cělo čislo jest slovnično nazvanje integer. Složena osnova cěločisl- dava pridavne formy v korpusě, s istim ě i koncovkami po kontekstu. Dvoslovno cělo čislo ne mora v vsakom slučaju zaměniti pridavnik: to by změnilo sintaksu. No obča etymologija whole ne rěšaje vse značenja německogo ganz. Cěločislna funkcija, algebraična cělost i običny cěly objekt potrěbujut svoje matematične definicije; pravopisno izravnjenje ne ustanovjaje ih ravnoznačnost. Zapis čuva 19 sklonjenyh tipov v 119 zaměnah i 138 uže soglasnyh form. Slovnik podpiraje celočislny osnovny smysl, no ne osvobodžaje od prověrky tehničnoj polisemije.\par
\smallskip\noindent\emph{Reason:} Cělo čislo is the dictionary term for integer. The compound stem cěločisl- supplies adjectival forms in the corpus, retaining ě and contextual endings. The two-word noun need not replace an adjective everywhere, since that would alter syntax. But the shared etymology whole does not settle every sense of German ganz. Integer-valued functions, algebraic integrality and an ordinary whole object require their respective mathematical definitions; spelling normalization does not establish their equivalence. The record preserves nineteen inflectional types in 119 substitutions alongside 138 already matching forms. The dictionary supports the basic integer sense without removing the need to examine technical polysemy.\par
\smallskip\noindent\emph{Izvory:} Medžuslovjansky slovnik, zaglavja 36585.\par
\noindent\emph{Sources:} Interslavic Dictionary, entries 36585.\par
\par\medskip\Needspace{8\baselineskip}
\hypertarget{DEC-0038}{}
\noindent\textbf{Izpolnjati uslovje / Satisfy a condition}\quad{\footnotesize\texttt{DEC-0038}}\par
\smallskip\noindent\emph{Razlog:} Formalny objekt izpolnjaje uslovje, relaciju ili aksiom. Izpolnjati/izpolniti byli izbrany dlja etogo tehničnogo smysla vměsto obširnogo satisfy/please u zadovaljati/zadovoliti. Poslědne sut stvarne slovnične slova: ih emocijno značenje samo po sobě ne dokazuje nemogućnost tehničnogo proširenja. Razlog jest jasnost v dannom matematičnom registru i soglasje s uže izbranoj familijej. Dostavati značit dostatočnost, ne izpolnjenje točno zadannogo uslovja. S novym transitivnym glagolom objekt potrěbuje akuzativ; ostavjenje starogo dativa ili relativnoj konstrukcije moglo by skvariti větu. Zato 88 predikatov uključajut 32 redove s kontekstno proměnjenym upravjanjem. Ni strokovy broj ni široka glosa ne zaměnjujut to proverjanje.\par
\smallskip\noindent\emph{Reason:} A formal object satisfies a condition, relation or axiom. Izpolnjati/izpolniti were chosen for that technical sense instead of the broader satisfy/please of zadovaljati/zadovoliti. The latter are genuine dictionary words: their emotional sense does not itself prove a technical extension impossible. The rationale is clarity in this mathematical register and consistency with the selected family. Dostavati denotes sufficiency, not satisfaction of a specified condition. With the new transitive verb the object requires the accusative; retaining the former dative or relative construction could damage the sentence. The 88 predicates therefore include 32 lines with contextual government changes. Neither line counts nor a broad gloss replaces that check.\par
\smallskip\noindent\emph{Izvory:} Medžuslovjansky slovnik, zaglavja 597, 2501, 36609, 444, 3356, 21715.\par
\noindent\emph{Sources:} Interslavic Dictionary, entries 597, 2501, 36609, 444, 3356, 21715.\par
\par\medskip\Needspace{8\baselineskip}
\hypertarget{DEC-0039}{}
\noindent\textbf{Potrěbny, potrěbno; ne potrěbno / Necessary, necessarily; scope of negation}\quad{\footnotesize\texttt{DEC-0039}}\par
\smallskip\noindent\emph{Razlog:} Potrěbny modifikuje uslovje, potrěbno može vyražati nužnost tvrdženja, a trěba jest bezosobna konstrukcija. Pri negaciji trěba najprvo rěšiti logičny obseg: ne jest potrěbno, že P ne znači isto kako P ne potrěbuje se. Zato ne nužno ne dopuščaje avtomatično ne trěba; neizběžno by dodalo smysl neizběžnosti, a morati/musěti by změnilo konstrukciju. Potrěbno i dostatočno uslovje ima jedinstveno čislo, kogda original govori o jednoj Bedingung. Slovnik potvrđaje potrěbny/potrěbno, ale takože nužny; izbor pervoj familije jest izloženaja prednost, ne dokaz neobstajanja drugoj. Trěbujut slěduje iz -ovati → -uj- i množinnoj koncovky -ut. Ty gramatične i logične proměny ne možna zaměniti jedinym poiskom korenja.\par
\smallskip\noindent\emph{Reason:} Potrěbny modifies a condition, potrěbno can express the necessity of an assertion, and trěba is impersonal. With negation, logical scope must be settled first: it is not necessary that P is not the same as P is not needed. Thus ne nužno cannot automatically become ne trěba; neizběžno would introduce inevitability, while morati/musěti would change the construction. A necessary and sufficient condition remains singular where the original has one Bedingung. The dictionary attests potrěbny/potrěbno but also nužny; preferring the former does not make the latter nonexistent. Trěbujut follows -ovati → -uj- with plural -ut. These grammatical and logical changes cannot be reduced to a root search.\par
\smallskip\noindent\emph{Izvory:} Medžuslovjansky slovnik, zaglavja 29955, 1296, 13042, 5786, 1942, 28707.\par
\noindent\emph{Sources:} Interslavic Dictionary, entries 29955, 1296, 13042, 5786, 1942, 28707.\par
\par\medskip\Needspace{8\baselineskip}
\hypertarget{DEC-0040}{}
\noindent\textbf{Daljši, dalje, daljny / Further, further on, remote}\quad{\footnotesize\texttt{DEC-0040}}\par
\smallskip\noindent\emph{Razlog:} Daljši označaje daljši argument, dodatny slučaj ili pokračovanje; dalje jest prislovje. Slovnik odlišaje daljnji/daljny v smislu remote. V zapisanyh kontekstah německy weitere i im folgenden ne govorjat o geografičnoj vzdalenosti. Zato 79 pridavnyh form i vosem srodnyh slučajev dobyli daljš-, s rodovym i padežnym soglasjem, a 212 prislovij dalje ostali. Slědujuči bylo by preuzko, kogda avtor znači additional, ne samo next in order. Nacionalno hrvatsko daljn- ne jest nepravilno v hrvatskom; redakcija izbiraje medžuslovjansko različenje iz slovnika. To pravilo ne izključaje prostorny smysl daljši, kogda jego stvarno potrěbuje kontekst.\par
\smallskip\noindent\emph{Reason:} Daljši marks a further argument, additional case or continuation; dalje is the adverb. The dictionary distinguishes daljnji/daljny in the sense remote. In the recorded contexts German weitere and im folgenden do not denote geographical distance. The 79 adjective forms and eight adjacent cases therefore received daljš- with appropriate agreement, while 212 adverbial dalje forms remained. Slědujuči would be too narrow where the author means additional rather than merely next in order. Croatian daljn- is not wrong in Croatian; the edition selects the Interslavic dictionary distinction. The rule does not exclude a spatial sense of daljši where context actually requires it.\par
\smallskip\noindent\emph{Izvory:} Medžuslovjansky slovnik, zaglavja 3032, 1177, 25102.\par
\noindent\emph{Sources:} Interslavic Dictionary, entries 3032, 1177, 25102.\par
\par\medskip\Needspace{8\baselineskip}
\hypertarget{DEC-0041}{}
\noindent\textbf{Razgledati objekt kako něšto / Consider an object as something}\quad{\footnotesize\texttt{DEC-0041}}\par
\smallskip\noindent\emph{Razlog:} Razgledati može uvoditi matematično razgledanje ili konstrukciju, v ktoroj objekt razumějemo kako objekt drugogo vida. Slovnično smatrjati znači watch/look; ględati takože ne dava samo po sobě tehničnu klasifikaciju. Zato v 93 zapisanyh neperceptivnyh kontekstah izbrano jest razgledati. Razvažati by zdůraznilo obdumavanje, prěgledati pregled ili inspekciju; ne sut avtomatične zaměny za regard as. Nedokonany aspekt čuva obči ili postavljajuči smysl, a neprěmenna zaměna na razgleděti by dodala drugu aspektnu perspektivu. Sto procentov dokazannoj razumljivosti iz togo ne slěduje; negativne i približne slovnične ocenky ostajut časťju dokazov. Smysl rečenija, ne podobnost slova k nacionalnomu glagolu, opreděljaje izbor.\par
\smallskip\noindent\emph{Reason:} Razgledati can introduce mathematical examination or treating an object as an object of another kind. Dictionary smatrjati means watch/look; ględati likewise does not by itself express technical classification. Razgledati was therefore selected in the 93 recorded non-perceptual contexts. Razvažati would foreground reflection, while prěgledati suggests review or inspection; neither automatically replaces regard as. Imperfective aspect preserves the general or stipulative sense, whereas uniform razgleděti would introduce a different aspectual perspective. This does not establish perfect comprehension; negative and approximate dictionary assessments remain part of the evidence. The sentence's meaning, not resemblance to a national verb, governs the choice.\par
\smallskip\noindent\emph{Izvory:} Medžuslovjansky slovnik, zaglavja 11776, 11777, 20742.\par
\noindent\emph{Sources:} Interslavic Dictionary, entries 11776, 11777, 20742.\par
\par\medskip\Needspace{8\baselineskip}
\hypertarget{DEC-0042}{}
\noindent\textbf{Kontrakcija; ne konvolucija / Contraction, not convolution}\quad{\footnotesize\texttt{DEC-0042}}\par
\smallskip\noindent\emph{Razlog:} Kontrakcija označaje operaciju invariantnoj teorije, v ktoroj spajanje simboličnyh redov daje sootvětny matričny produkt i izměnjenje težiny. V toh kontekstah německo Faltung ne označaje fizično svivanje i ne potrěbuje konvoluciju. Poslědnja jest drugy matematičny termin; transvekcija i prěnos takože označajut odlišne operacije v korpusě. Slovnično kontrakcija i nacionalne matematične sravnenja podpirajut internacionalny naziv, ale sama glosa contraction ne jest dokaz tožestva vseh operacij s tym imenem. Definicija i formula imajut prvenstvo. Prěhod od neutra svijanje k ženskomu kontrakcija potrěbuje soglasje pridavnikov, zaimen i predikatov. Glagolne slučaje izrěčeny sut črez izvršiti kontrakciju; samokontrakcija črez kontrakciju samej s soboju. Takova pererabotka čuva operaciju, ne samo prěvodi bukvalno metaforu.\par
\smallskip\noindent\emph{Reason:} Kontrakcija denotes the invariant-theoretic operation in which pairing symbolic rows gives the corresponding matrix product and change of weight. In these contexts German Faltung is neither physical folding nor convolution. Konvolucija is a different mathematical term, and transvection and transfer also name distinct operations in the corpus. Dictionary kontrakcija and national mathematical comparators support the international name, but the gloss contraction alone does not establish identity among every operation bearing it. The definition and formula take precedence. Moving from neuter svijanje to feminine kontrakcija requires agreement in adjectives, pronouns and predicates. Verbal cases use izvršiti kontrakciju; self-contraction is expressed by contraction with itself. The recasting preserves the operation rather than literally translating its metaphor.\par
\smallskip\noindent\emph{Izvory:} Medžuslovjansky slovnik, zaglavja 19795.\par
\noindent\emph{Sources:} Interslavic Dictionary, entries 19795.\par
\par\medskip\Needspace{8\baselineskip}
\hypertarget{DEC-0043}{}
\noindent\textbf{Prědpokladati i prědpoložiti / Assume: imperfective and perfective}\quad{\footnotesize\texttt{DEC-0043}}\par
\smallskip\noindent\emph{Razlog:} Prědpokladati jest izbrany nedokonany glagol; prědpoložiti dokonany, s participnymi i imenovnymi formami prědpolož-. Prědpolagati takože jest gramatično i slovnično zaglavje tipa 1: njegovo odbačenje jest prednost jednolike redakcije, ne normativna zabrana. Zaměna cěloj familije prědpolož- na prědpoklad- by narušila aspekt i tvorjenje imenovnikov. Prěd- s ě jest pravopisna konvencija izdanje. Zapisano razdělenje 107 proměn vključaje 36 nedokonanyh glagolov, 17 pasivnyh ili pridavnyh form, 11 dokonanyh i 43 imenovnika. Suma čuva razny gramatičny raboty pri odnom matematičnom dějstvu postavjenja uslovja; ne poravnjava ih v jednu mehaničnu slovničnu zaměnu.\par
\smallskip\noindent\emph{Reason:} Prědpokladati is the chosen imperfective; prědpoložiti is perfective, with participial and nominal prědpolož- forms. Prědpolagati is also grammatical and a type-1 headword: rejecting it is an editorial consistency preference, not a normative ban. Replacing the whole prědpolož- family with prědpoklad- would damage aspect and noun formation. Prěd- with ě follows the edition's spelling policy. The recorded 107 changes comprise 36 imperfective verbs, seventeen passive or adjectival forms, eleven perfectives and 43 nouns. These counts preserve different grammatical jobs within mathematical assumption rather than flattening them into one mechanical substitution.\par
\smallskip\noindent\emph{Izvory:} Medžuslovjansky slovnik, zaglavja 3130, 5324, 11247, 18400, 20062.\par
\noindent\emph{Sources:} Interslavic Dictionary, entries 3130, 5324, 11247, 18400, 20062.\par
\par\medskip\Needspace{8\baselineskip}
\hypertarget{DEC-0044}{}
\noindent\textbf{Zato; poslědovateljno; zato že / Therefore, successively and because}\quad{\footnotesize\texttt{DEC-0044}}\par
\smallskip\noindent\emph{Razlog:} Zato vyražaje rezultat ili vyvod; poslědovateljno časovy ili operativny porędok; zato že pričinu. Eto tri različne logične veze. Iz 15 pojavjenij poslědovateljno tri stvarno govorjat o sukcesivnoj kontrakciji ili porědku v kompozicijnom redu i ostajut; dvanadset označajut vyvod. Slědovateljno takože jest slovničny rezultatny marker, a njegovo zaměnjenje ne čini slovo negramatičnym. Redakcija preferuje kratko zato. Jedna konstrukcija, v ktoroj zato vměstě s že moglo by smestiti pričinu i posledok, dobyla tym samym. Ty dva slova ne sut obči ključ za vsaky vyvod, zato že v drugih kontekstah mogu značiť the same. Čitatelj mora viděti napravljenje argumenta, ne tolko pravilno izbranu bukvu.\par
\smallskip\noindent\emph{Reason:} Zato marks a result or inference; poslědovateljno temporal or operational succession; zato že a cause. These are three different logical connections. Of fifteen poslědovateljno occurrences, three genuinely concern successive contraction or order in a composition series and remain; twelve express inference. Slědovateljno is also a dictionary result marker, so replacing it does not make it ungrammatical. The edition prefers shorter zato. One construction in which zato beside že could confuse cause with consequence instead uses tym samym. That phrase is not a universal inference marker, since elsewhere it can mean the same. Readers need to see the direction of the argument, not merely an approved spelling.\par
\smallskip\noindent\emph{Izvory:} Medžuslovjansky slovnik, zaglavja 35139, 3056, 10980.\par
\noindent\emph{Sources:} Interslavic Dictionary, entries 35139, 3056, 10980.\par
\par\medskip\Needspace{8\baselineskip}
\hypertarget{DEC-0045}{}
\noindent\textbf{Prinaležati; prinaležet / Belong; plural agreement}\quad{\footnotesize\texttt{DEC-0045}}\par
\smallskip\noindent\emph{Razlog:} Prinaležati jest izbrany glagol členstva ili prinaležnosti, s prezentom prinaleži. Osnova nema d v prinadlež-, zato izravnjenje usoglasnjaje familiju s slovnikom. Ale deset množinnyh prinadležat potrěbuju prinaležet, ne samo odstranjenje d: slovnična prezentna osnova na -i- bere koncovku -et. Prinaležnost v občnoj glosě znači affiliation; matematično členstvo potrebuje kontekst i ne izravnjava se avtomatično s vključanjem podmnožiny. Zapis razlišaje 64 lične predikaty, 23 modifikatory/participy, devet pridavnikov, dva infinitiva i pet imenovnikov. Čislo 103 zato ne daje edinogo paradigmatskogo rěšenja bez gramatiky. Znak členstva možno sohraniti, ale nema nuždy zaměnjati im cělu avtorsku prozu.\par
\smallskip\noindent\emph{Reason:} Prinaležati is the chosen verb of membership or belonging, with present prinaleži. Its stem lacks the d of prinadlež-, so alignment follows the dictionary. But ten plural prinadležat forms require prinaležet, not merely deletion of d: the dictionary present stem in -i- takes plural -et. Prinaležnost is generally glossed affiliation; mathematical membership needs context and must not automatically be equated with subset inclusion. The record distinguishes 64 finite predicates, 23 modifiers or participles, nine adjectives, two infinitives and five nouns. The total of 103 therefore does not supply one grammatical replacement. Membership notation may be preserved without replacing all the author's prose by a symbol.\par
\smallskip\noindent\emph{Izvory:} Medžuslovjansky slovnik, zaglavja 20656, 20657.\par
\noindent\emph{Sources:} Interslavic Dictionary, entries 20656, 20657.\par
\par\medskip\Needspace{8\baselineskip}
\hypertarget{DEC-0046}{}
\noindent\textbf{Srodne formy ne imajut jednu gramatiku / Related forms do not share one grammar}\quad{\footnotesize\texttt{DEC-0046}}\par
\smallskip\noindent\emph{Razlog:} Toj, slědujuči, slědujut, potrěbno, trěbujut, javjajut, poslědica i poslědovateljny pripadajut odlišnym pravilam. Pokazateljno toj čuva običnu referenciju; slědujuči i slědujut razlikujut particip od ličnogo glagola; potrěbno čuva modalen smysl; trěbujut potrěbuje -uj-; javjajut čuva slovničnu osnovu javjati, vměsto avtomatičnogo -lj-. Poslědica jest posledok, a poslědovateljny govori o porědku. Zato objedinjena kontrolja tih 21 forma ne jest nova edinaja lingvistična pravila. I tukaj ć/č zavisi od pravopisnogo sloja, no odsutstvo potrebnogo j ili pogrešno lice jest gramatičny problem. Srodnost familij pomaga najti pojavjenja, ne rěšava ih funkciju.\par
\smallskip\noindent\emph{Reason:} Toj, slědujuči, slědujut, potrěbno, trěbujut, javjajut, poslědica and poslědovateljny belong to different rules. Demonstrative toj preserves ordinary reference; slědujuči and slědujut distinguish participle from finite verb; potrěbno preserves modal meaning; trěbujut requires -uj-; javjajut follows dictionary javjati rather than automatically adding -lj-. Poslědica is a consequence, while poslědovateljny concerns sequence. The combined examination of these 21 forms therefore does not constitute a new single linguistic rule. Here too ć/č depends on orthographic policy, but a missing required j or wrong person is grammatical. Family resemblance helps locate occurrences without deciding their function.\par
\smallskip\noindent\emph{Izvory:} Medžuslovjansky slovnik, zaglavja 1939, 1378, 1942, 29955, 3612, 18328, 2333.\par
\noindent\emph{Sources:} Interslavic Dictionary, entries 1939, 1378, 1942, 29955, 3612, 18328, 2333.\par
\par\medskip\Needspace{8\baselineskip}
\hypertarget{DEC-0047}{}
\noindent\textbf{Obstajanje: dokaz, že objekt jest / Existence: a proof that an object exists}\quad{\footnotesize\texttt{DEC-0047}}\par
\smallskip\noindent\emph{Razlog:} Obstajanje jest izbrany imenovnik k obstajati za eksistenciju funkcije, konstante, ideala ili polja. V zapisanyh 25 imenovnyh i šesti glagolnyh kontekstah sušstv- ne označaje sučnost, sučstvenost ili nynějšnu dostupnost. Zato sučnost ili sučstveny ne byli priměrne zaměny, hoti izgledaju srodno. Imenovnik i glagol morajut sohraniti padež, infinitiv, lice i broj. Odluka razširjaje istu konvenciju kako pri obstajati; ne dava nov nezavisen dokaz razumljivosti. Slabosti slovničnogo profila obstajati ostajut vidlivy i ovdě. V teoremě o obstajanju smysl jest istinnost eksistencijnogo tvrdženja, ne prosto fizično prisutstvo i ne ocenka važnosti objekta.\par
\smallskip\noindent\emph{Reason:} Obstajanje is the selected noun corresponding to obstajati for the existence of a function, constant, ideal or field. In the recorded 25 nominal and six verbal contexts, sušstv- denotes neither essence, essentiality nor present availability. Sučnost or sučstveny would therefore not be appropriate replacements despite their resemblance. Noun and verb must preserve case, infinitive, person and number. The choice extends the same convention as obstajati; it supplies no new independent evidence of comprehension. The weaknesses of the dictionary profile for obstajati remain visible here. An existence theorem asserts existence, not merely physical presence or the object's importance.\par
\smallskip\noindent\emph{Izvory:} Medžuslovjansky slovnik, zaglavja 9895, 36821, 36822.\par
\noindent\emph{Sources:} Interslavic Dictionary, entries 9895, 36821, 36822.\par
\par\medskip\Needspace{8\baselineskip}
\hypertarget{DEC-0048}{}
\noindent\textbf{Doskonalo polje; nedoskonalo polje / Perfect and imperfect fields}\quad{\footnotesize\texttt{DEC-0048}}\par
\smallskip\noindent\emph{Razlog:} Doskonalo polje jest tehničny termin, ne ocenka krasoty ili polnoty. V korpusnoj definiciji vsaka nerazložima funkcija nad takym poljem v vhodnom razširjenju razpada se na različne linearne faktory. To različenje vedet izbor, a ne obično značenje perfect. Sravnenja sut razděljene: češske i slovenske matematične izvory citovane v zapisě davajut perfekt-, bosansky savršen-, a obči slovnik podpiraje doskonal-/dokonal- s poljskimi, češko-slovačskimi, ukrajinskimi i běloruskimi sravnenjami. Profil doskonaly ima i negativne ocenky, osoblivo dlja několikyh južnyh jazykov i ruskogo; ne jest neutralna vseobča poběda. Izbor jest provisionalna tehnična konvencija s definicijeju. Dokončeny by značilo completed, popolny full/complete, a perfekt može takože nazvati gramatičny čas. Nedoskonaly čuva negaciju istoj točnoj definicije.\par
\smallskip\noindent\emph{Reason:} Doskonalo polje is a technical term, not a judgment of beauty or completeness. In the corpus definition every irreducible function over such a field splits into distinct linear factors in a suitable extension. That distinction governs the choice, not the ordinary meaning of perfect. The evidence is divided: the Czech and Slovene mathematical sources cited in the record use perfekt-, the Bosnian source savršen-, while the general dictionary supports doskonal-/dokonal- through Polish, Czech–Slovak, Ukrainian and Belarusian comparators. Doskonaly also has negative profile assessments, notably for several southern languages and Russian; it is not a neutral universal winner. The choice is a provisional technical convention anchored in the definition. Dokončeny would mean completed, popolny full/complete, and perfekt may also name a grammatical tense. Nedoskonaly preserves negation of the same precise definition.\par
\smallskip\noindent\emph{Izvory:} Medžuslovjansky slovnik, zaglavja 4207, 15052, 35256; korpus, rabota 24, § 1, punkt 14, definicija doskonalogo polja.\par
\noindent\emph{Sources:} Interslavic Dictionary, entries 4207, 15052, 35256; corpus, paper 24, § 1, item 14, definition of a perfect field.\par
\par\medskip\Needspace{8\baselineskip}
\hypertarget{DEC-0049}{}
\noindent\textbf{Razgledajmo; razgleděti; razgledanje / Let us consider; consider; examination}\quad{\footnotesize\texttt{DEC-0049}}\par
\smallskip\noindent\emph{Razlog:} Razgledajmo uvodi avtorsko razgledanje; razgleděti jest dokonany infinitiv; razgledanje imenovnik dějstva. Ty formy nemajut jednu aspektnu ili sintaktičnu ulogu, ale pripadajut odnoj izbranej familiji vměsto razsmotr-. Slovnik sodrživa i råzsmotriti tipa 3, zato izbor drugoj familije ne dokazuje, že prvy glagol neobstaje. Važnější problem jest različenje examination od reasoning a idiomu take into account. V zachovanom sravnenju 39 slučajev pet imenovnikov s německym Überlegung potrěbuje razvažanje, a dva s in Betracht konstrukciju v razgled. Ostale 32 ostajut pri razgled-. Tako celostna familijna tablica sama po sobě ne jest dostatočny překlad: kontekst može izključiti činovne najbližšu formu.\par
\smallskip\noindent\emph{Reason:} Razgledajmo introduces the author's examination, razgleděti is a perfective infinitive, and razgledanje an action noun. These forms have different aspectual and syntactic roles while belonging to the selected family instead of razsmotr-. The dictionary also contains type-3 råzsmotriti, so choosing another family does not prove the first verb nonexistent. The more important issue is distinguishing examination from reasoning and the idiom take into account. In the preserved comparison of 39 cases, five nouns aligned with German Überlegung require razvažanje and two in-Betracht constructions require v razgled. The other 32 retain razgled-. A complete family table is thus not sufficient translation by itself: context can exclude the morphologically nearest form.\par
\smallskip\noindent\emph{Izvory:} Medžuslovjansky slovnik, zaglavja 30464, 11776, 11777, 6010.\par
\noindent\emph{Sources:} Interslavic Dictionary, entries 30464, 11776, 11777, 6010.\par
\par\medskip\Needspace{8\baselineskip}
\hypertarget{DEC-0050}{}
\noindent\textbf{Razvažanje i v razgled / Reasoning and taking into account}\quad{\footnotesize\texttt{DEC-0050}}\par
\smallskip\noindent\emph{Razlog:} Pri frazah entsprechende Überlegung, Alle benutzten Überlegungen, Nach unseren allgemeinen Überlegungen, keine Überlegungen über kommutative Körper i rein algebraisch-arithmetische Überlegungen izbrany sut padežne formy razvažanje, razvažanja, razvažanjam i razvažanjami. Rěč jest o toku razsuždenja, ne samo o prědmetě izslědovanja. Dve frazy in Betracht kommende i kommt nicht in Betracht potrěbuju idiom vyražajuči, čto prihodi v razgled ili ne prihodi v razgled; bukvalno v razgledanje ne prenosi to automatično. Eto kontekstna motivacija, ne tvrdženje, že německe Betrachtung i Überlegung nikogda ne prěkrivajut se. Slovo može byti pravilno samo po sobě i nepodhodno v konkretnoj frazě. Dokazateljny prilog sohranjaje vse sedm svęzanyh izvorno-cělovyh fraz.\par
\smallskip\noindent\emph{Reason:} For entsprechende Überlegung, Alle benutzten Überlegungen, Nach unseren allgemeinen Überlegungen, keine Überlegungen über kommutative Körper and rein algebraisch-arithmetische Überlegungen, the selected case forms are razvažanje, razvažanja, razvažanjam and razvažanjami. The issue is a course of reasoning rather than merely an object of investigation. In Betracht kommende and kommt nicht in Betracht require an idiom conveying coming, or not coming, into consideration; literal v razgledanje does not automatically convey that. This is contextual reasoning, not a claim that German Betrachtung and Überlegung never overlap. A word can be well formed yet unsuitable in a particular phrase. The evidence supplement preserves all seven linked source–target expressions.\par
\smallskip\noindent\emph{Izvory:} Medžuslovjansky slovnik, zaglavja 6010, 11776, 11777.\par
\noindent\emph{Sources:} Interslavic Dictionary, entries 6010, 11776, 11777.\par
\par\medskip\Needspace{8\baselineskip}
\hypertarget{DEC-0051}{}
\noindent\textbf{Odgovarjati, analogično i po redu / Correspondence, analogy and ordered pairing}\quad{\footnotesize\texttt{DEC-0051}}\par
\smallskip\noindent\emph{Razlog:} Odgovarjati i odgovarjajuči označajut sootvětnost; korespondencija imenovaje otnošenje; analogično uvodi sravnimy argument. Německe beziehungsweise i jeweils mogut vměsto togo opreděljati točno sparjenje: prvy s prvym, vtory s vtorym, po redu ili vsak od ih. Zato 23 glagola i 23 modifikatory, dva imenovnika, šest analogičnyh diskursnyh upotrěb i tri porědkove konstrukcije ne dopuščajut jednu zaměnu sootv-. Osoblivo pri dvuh spiskah izgubljeny porędok menja matematično sootvětnost. Odnosno ili respektivno ne sut avtomatično nevozmožne, ale zapis izbiraje izrično sparjenje tam, kde obči marker ne dostava. Izbrana familija sootvětstvuje drugim člankam o korespondenciji; slovnične glosy ne zaměnjujut konkretno zadano otnošenje.\par
\smallskip\noindent\emph{Reason:} Odgovarjati and odgovarjajuči express correspondence; korespondencija names the relation; analogično introduces a comparable argument. German beziehungsweise and jeweils may instead specify an exact pairing: first with first, second with second, in order, or each of them. The 23 verbs, 23 modifiers, two nouns, six analogous discourse uses and three ordered constructions therefore do not admit one sootv- replacement. With two lists, losing order changes the mathematical correspondence. Odnosno or respektivno are not inherently impossible, but the record chooses explicit pairing where a general marker is insufficient. The family is consistent with the other correspondence entries; dictionary glosses do not replace the specified relation.\par
\smallskip\noindent\emph{Izvory:} Slovnik, zaglavja 35938, 6063, 23209.\par
\noindent\emph{Sources:} Dictionary entries 35938, 6063, 23209.\par
\par\medskip\Needspace{8\baselineskip}
\hypertarget{DEC-0052}{}
\noindent\textbf{Potrěbny i neobhodimy / Necessary and indispensable}\quad{\footnotesize\texttt{DEC-0052}}\par
\smallskip\noindent\emph{Razlog:} Neobhodimy jest stvarno slovnično zaglavje tipa 1 za necessary/indispensable. Njegovo odbačenje v redakciji jest izbor jednoj familije potrěbn-, ne dokaz lingvistično nepravilnogo slova. Dvadeset pet imenovnyh i predikativnyh slučajev dopuščajut izravnjenje, no vosemnadset potrěbuju polnu modalnu ili vyvodnu klauzu. Frazy must necessarily, not necessarily i follows necessarily nemajut istu logičnu konstrukciju. Dvojny modal možno skratiti samo ako sila tvrdženja ostaje ta že. Ne trěba ne znači avtomatično not necessarily; neizběžno by dodalo neizběžnost. Prednost terminologične jednolikosti zato podčinena jest obsegu negacije i razlicě medžu uslovjem i tvrđenjem o vyvodu. Bibliografično preloženy naslov sleduje tu že semantiku, no originalny naslov ne potrěbuje prepisyvanja.\par
\smallskip\noindent\emph{Reason:} Neobhodimy is a genuine type-1 headword for necessary/indispensable. Rejecting it in this edition selects a single potrěbn- family; it does not prove the word linguistically wrong. Twenty-five nominal and predicative cases permit alignment, but eighteen require a complete modal or inferential clause. Must necessarily, not necessarily and follows necessarily have different logical constructions. A doubled modal can be shortened only if the assertion retains its force. Ne trěba does not automatically mean not necessarily, and neizběžno would add inevitability. Terminological consistency is therefore subordinate to negation scope and the difference between a condition and a claim of inference. A translated bibliographical title follows the same semantics without requiring alteration of the original title.\par
\smallskip\noindent\emph{Izvory:} Slovnik, zaglavja 33105, 29955, 1296.\par
\noindent\emph{Sources:} Dictionary entries 33105, 29955, 1296.\par
\par\medskip\Needspace{8\baselineskip}
\hypertarget{DEC-0053}{}
\noindent\textbf{Shranjati, shraniti, shranjenje / Preserve, preservation and conservation}\quad{\footnotesize\texttt{DEC-0053}}\par
\smallskip\noindent\emph{Razlog:} Shranjati jest nedokonany glagol, shraniti dokonany, a shranjenje imenovnik dlja sohranjenja svojstva ili zakona. Slovnik daje preserve/preservation/conservation, v odličenji od guard/protect pri ohranjati. Eto podpira izbrany register invariantnosti i zakonov sohranjenja; ne dokazuje, že nacionalno slovensko ohranjati jest v tom jazykě nepravilno. Trideset šest kontekstov vključajut rang, regularnost, relacije, formy i fizikalne zakony. Lična glagolna forma, pasivny particip i imenovnik ne mogu dobyti jednu koncovku. Ostavaje vyrazilo by tolko ostavanje, ne vsaku aktyvnu operaciju sohranjenja. Slovo čuva uslovja avtorskoj invariantnosti; samo njegov naziv ne dokazuje, že veličina jest sohranjena pri vsakom prěobraženju.\par
\smallskip\noindent\emph{Reason:} Shranjati is imperfective, shraniti perfective, and shranjenje the noun for preservation of a property or conservation law. The dictionary gives preserve/preservation/conservation, distinct from guard/protect for ohranjati. This supports the selected register of invariance and conservation without making national Slovene ohranjati wrong in Slovene. The 36 contexts include rank, regularity, relations, forms and physical laws. Finite verb, passive participle and noun cannot take one ending. Ostavaje would express remaining, not every active preservation operation. The word retains the author's conditions of invariance; naming conservation does not prove a quantity invariant under every transformation.\par
\smallskip\noindent\emph{Izvory:} Slovnik, zaglavja 18694, 18695, 18696, 1563.\par
\noindent\emph{Sources:} Dictionary entries 18694, 18695, 18696, 1563.\par
\par\medskip\Needspace{8\baselineskip}
\hypertarget{DEC-0054}{}
\noindent\textbf{Izpolnjati: objekt i padež / Satisfy: object and case}\quad{\footnotesize\texttt{DEC-0054}}\par
\smallskip\noindent\emph{Razlog:} Uslovje, uravnenje ili relacija jest objekt izpolnjenja; vybran element ili bazis može byti subjektom. Zato naslědžene udovletvor-/udovolja konstrukcije potrěbujut nejen novu osnovu izpoln-, ale i sootvětno upravjanje objektom. Dativ ili relativna klauza, pravilna pri drugom glagolu, ne može ostati bez razbora. V 34 zapisanyh slučajah ne ide o ludsku zadovoljenost: izbor izpolnjati/izpolniti jest obosnovany formalnym smyslom. Zadovaljati ostaje slovnično slovo s širšimi značenjami, a njegovo izključanje iz etogo registra jest kontekstna prednost. To prodolžaje razlikovanje iz članka Izpolnjati uslovje; ne dava nov dokaz pan-slovjanskoj prednosti po čislu popravok. Vse kvantifikatory i ograničenja na objekt morajut ostati v věti.\par
\smallskip\noindent\emph{Reason:} A condition, equation or relation is the object of satisfaction; a selected element or basis may be its subject. The inherited udovletvor-/udovolja constructions therefore need not only the izpoln- stem but corresponding object government. A dative or relative clause appropriate with a different verb cannot remain without examination. None of the 34 recorded cases concerns human contentment; the formal sense motivates izpolnjati/izpolniti. Zadovaljati remains a dictionary word with broader meanings, and its exclusion from this register is contextual. This continues the distinction under Satisfy a condition; the number of repairs supplies no additional evidence of pan-Slavic superiority. All quantifiers and restrictions on the object must remain.\par
\smallskip\noindent\emph{Izvory:} Slovnik, zaglavja 597, 2501, 36609, 444, 3356.\par
\noindent\emph{Sources:} Dictionary entries 597, 2501, 36609, 444, 3356.\par
\par\medskip\Needspace{8\baselineskip}
\hypertarget{DEC-0055}{}
\noindent\textbf{Zaštita citatnogo prezimena pri proměně pisma / Protecting a cited surname during script conversion}\quad{\footnotesize\texttt{DEC-0055}}\par
\smallskip\noindent\emph{Razlog:} Prezime Weierstrass ostaje točno napisano i kogda okružajuča medžuslovjanska proza prěhodi iz latinice v kirilicu. V izvornoj LaTeX-razmetkě oboločka foreign označaje tu granicu. Njejna identična tekstova definicija dava upravo svoj argument: ne dodaje preklad, font ni novo matematično značenje. Vključeny tekst mora iměti tu definiciju dostupnu do prvogo upotrěbjenja; definicija v pozdnějšoj lokalnoj grupě ne vladaje ranějšim tekstom. Odstranjenje oznaky moglo bi ostaviti vidlivy PDF neproměnjenym, ale izgubilo bi granicu dlja prěpisa. Zato sohranjenje čitlivogo rezultata i sohranjenje mašinno razpoznavajemogo imena sut odvojene potrěby. Ta oznaka ne jest dokaz lingvistično pravilnogo eponima.\par
\smallskip\noindent\emph{Reason:} The surname Weierstrass remains exactly spelled when the surrounding Interslavic prose changes from Latin to Cyrillic. In the editable LaTeX, the foreign wrapper marks that boundary. Its identity definition returns precisely its argument: it adds neither a translation, a font nor mathematical meaning. An included text must have that definition available before first use; a definition in a later local group does not govern earlier text. Removing the marker could leave the visible PDF unchanged while losing the boundary needed by script conversion. Preserving the readable result and preserving the machine-recognizable name are therefore separate requirements. The marker does not certify the linguistic correctness of an eponym.\par
\smallskip\noindent\emph{Izvory:} Korpus, rabota 13, eponim Weierstrass i lokalna LaTeX-definicija oboločky foreign.\par
\noindent\emph{Sources:} Corpus, paper 13, Weierstrass eponym and the local LaTeX identity wrapper foreign.\par
\par\medskip\Needspace{8\baselineskip}
\hypertarget{DEC-0056}{}
\noindent\textbf{Pored i poredok; poręd i porędok / Beside and order: basic and etymological spelling}\quad{\footnotesize\texttt{DEC-0056}}\par
\smallskip\noindent\emph{Razlog:} Pored, predlog, i poredok, imenovnik, nemajut isto značenje, hoti imajut sroden korenj. Slovnik piše etymologično poręd (+genitiv) za beside/next to i porędȯk za order/sequence. Osnovny pravopis dopuščaje odstranjenje opcionalnyh diakritik, zato pored ne jest samo po sobě nepravilno. Ranějše pravilo izravnjavalo 34 predloga i pet imenovnikov po tehdy izbranoj etymologičnoj sistemě; nynějšny osnovny zapis mora byti razlišen od togo istorijnogo izbora. Kromě može vnovesti izključanje vměsto susědstva ili dodatka i ne jest avtomatična zaměna. Pri predlogu važen jest padež; pri imenovniku njegov broj i sklonjenje. Oba zadanja ne može rěšiti jedna bukva.\par
\smallskip\noindent\emph{Reason:} Pored the preposition and poredok the noun have different meanings despite their related root. The dictionary writes etymological poręd, with the genitive, for beside/next to and porędȯk for order/sequence. Basic orthography permits removal of optional diacritics, so pored is not inherently wrong. The earlier rule aligned 34 prepositions and five nouns with the etymological system then selected; current basic spelling must be distinguished from that historical choice. Kromě may introduce exclusion instead of adjacency or addition and is not an automatic substitute. The preposition requires the appropriate governed case, while the noun requires number and inflection. One letter cannot settle both tasks.\par
\smallskip\noindent\emph{Izvory:} Slovnik, zaglavja 5198, 3145; Pravopis, opcionalne bukvy.\par
\noindent\emph{Sources:} Dictionary entries 5198, 3145; Orthography, optional letters.\par
\par\medskip\Needspace{8\baselineskip}
\hypertarget{DEC-0057}{}
\noindent\textbf{Vsih, vsim, vsimi: vseobče otnošenje / Vsih, vsim, vsimi: universal reference}\quad{\footnotesize\texttt{DEC-0057}}\par
\smallskip\noindent\emph{Razlog:} Vsih, vsim i vsimi sut pluralne kosye formy zaimena ves/vsi. Polna opublikovana tablica dava vsih v genitivě i lokativě, vsim v dativě i vsimi v instrumentalě. Ona prěmo podpiraje izbrane formy, kde kratak lokalny tekst imal jedino odkaz na tablicu. Eto jest dokaz konsultovany pri rekonstrukciji, ne tvrdženje o ranějšem čitanju. Istoričny odkaz na vśaky ne byl sam po sobě dostatočny: each/every i all ne sut jedno zaglavje. Zapis razlikuje 193 vseh, jedno vsem i jedno vsemi s množinnym smyslom. Vsem v lokativě jedinstvenogo čisla jest legitimna forma v toj že paradigmě; pravilo ne směje odstraniti ju bez konteksta. Každy takože ne može zaměniti vse takove pojavjenja bez proměny distributivnogo smysla.\par
\smallskip\noindent\emph{Reason:} Vsih, vsim and vsimi are plural oblique forms of ves/vsi. The full published table gives vsih in the genitive and locative, vsim in the dative and vsimi in the instrumental. It directly supports the selected forms where the short local text contained only a component reference. This is evidence consulted during reconstruction, not a claim of earlier consultation. The historical citation of vśaky alone was insufficient: each/every and all are not one headword. The record distinguishes 193 vseh, one vsem and one vsemi with plural meaning. Vsem is a legitimate locative singular in the same paradigm; a rule must not remove it without context. Nor can každy replace every such occurrence without changing distributive meaning.\par
\smallskip\noindent\emph{Izvory:} Jan van Steenbergen, Gramatika, Zaimena, polna opublikovana tablica ves; slovnik, zaglavje 4326.\par
\noindent\emph{Sources:} Jan van Steenbergen, Grammar, Pronouns, complete published ves table; dictionary entry 4326.\par
\par\medskip\Needspace{8\baselineskip}
\hypertarget{DEC-0058}{}
\noindent\textbf{Genitiv množiny glagolnyh imenovnikov / Genitive plurals of verbal nouns}\quad{\footnotesize\texttt{DEC-0058}}\par
\smallskip\noindent\emph{Razlog:} Formy na -anij/-enij byli izbrany dlja genitiva množiny glagolnyh imenovnikov vměsto naslědženyh -anj/-enj. Obseg vključaje 182 pojavjenja i 39 leksičnyh tipov, napr. pytanj, rěšenj, jednačenj i uravnenj. Razlog v zapisě odvolajet se na tot že generator, ktory byl citovany pri prědstavjenij. Njegova reproducirujema evidencija ne jest v etoj rekonstrukciji nezavisno potvrdžena; obča gramatična stranica o nulevom genitivě neutra sama ne dokazuje točnoj koncovky -nij. To ograničenje mora važiti i dlja proširenogo pravila, ne tolko prvogo priměra. Stepenj, stupenj i korenj sut odlišne imenovniki, ne pokvarjene genitivy, zato automat na vsako -enj byl by pogrešny. Čisla pojavjenij ne proměnjajut slabost istogo morfologičnogo dokazateljstva.\par
\smallskip\noindent\emph{Reason:} Forms in -anij/-enij were adopted for genitive plurals of verbal nouns instead of inherited -anj/-enj. The scope includes 182 occurrences across 39 lexical types, such as pytanj, rěšenj, jednačenj and uravnenj. The record appeals to the same generator cited for prědstavjenij. Its reproducible evidence has not been independently recovered in this reconstruction; the general grammar's account of a zero neuter genitive does not itself establish the exact -nij form. That qualification must accompany the extended rule, not only the first example. Stepenj, stupenj and korenj are distinct lexical nouns, not damaged genitives, so rewriting every -enj would be wrong. Occurrence counts cannot strengthen the same incomplete morphological evidence.\par
\smallskip\noindent\emph{Izvory:} Gramatika, Imenovniki, sklanjanje srednjego roda; odkaz na generator pri prědstavjenij ostaje nepotvrdženy.\par
\noindent\emph{Sources:} Grammar, Nouns, neuter declension; the generator reference under prědstavjenij remains unverified.\par
\par\medskip\Needspace{8\baselineskip}
\hypertarget{DEC-0059}{}
\noindent\textbf{-ajut i -aju: kto jest subjekt? / -ajut and -aju: who is the subject?}\quad{\footnotesize\texttt{DEC-0059}}\par
\smallskip\noindent\emph{Razlog:} Tretje lice množiny ima koncovky -ut/-et, pri kontrakciji -ajut. Prvo lice jedinstvenogo čisla može završati se na -aju. Zato dodavanje t potrěbuje subjekt ili odnosnu predhodnu frazu: avtorovo vybiraju, naměrjaju, podaju, pripominaju ili dodaju ne jest množinna ošibka. Imennye formy slučaju i algebraju takože ne sut glagoly. Zapis odlišaje 107 množinnyh predikatov od peti avtorskyh jedinstvenyh i od imennogo materijala. Kogda subjekt jest koordinovany, padežna forma slova okolo glagola sama ne dostava za opreděljenje broja. Pravilo popravjaje soglasje, čuvajuči korenj, aspekt, originalnu notaciju i avtorsku perspektivu. To jest kontekstna gramatika, ne zaměna vsakej koncovky -aju.\par
\smallskip\noindent\emph{Reason:} Third-person plurals end in -ut/-et, or contracted -ajut; first-person singulars can end in -aju. Adding t therefore requires the subject or relative antecedent: the author's vybiraju, naměrjaju, podaju, pripominaju or dodaju are not plural errors. Nominal slučaju and algebraju are not verbs either. The record distinguishes 107 plural predicates from five authorial singulars and from nominal material. With a coordinated subject, the case of a nearby word alone is insufficient to determine number. The rule repairs agreement while retaining stem, aspect, original notation and authorial perspective. This is contextual grammar, not substitution of every -aju ending.\par
\smallskip\noindent\emph{Izvory:} Gramatika, Glagoly, prezentne koncovky i kontrakcija.\par
\noindent\emph{Sources:} Grammar, Verbs, present endings and contraction.\par
\par\medskip\Needspace{8\baselineskip}
\hypertarget{DEC-0060}{}
\noindent\textbf{Vsaky i každy; each ne jest prosto all / Vsaky and každy: each is not simply all}\quad{\footnotesize\texttt{DEC-0060}}\par
\smallskip\noindent\emph{Razlog:} Vsaky i každy sut oba slovnične zaglavja tipa 1 za each/every. Profil každy ima pozitivne zapadne i vostočne ocenky, ale negativne za několik južnyh jazykov; profil vsaky ima širšu južnu podpiru, s približnymi ocenkami v poljskom i slovačskom. Eto motivuje preferenciju vsaky, ne izměreny dobytok razumljivosti. Nacionalne každý/każdy i vsak/svaki ostajut sravnenja, ne alternativy za glasovanje po čislu tokenov. Sto pojavjenij s každy ima distributivnu ulogu; zaměna na vse/vsi mogla by změniti obseg, kolektivno čitanje ili soglasje. Tablica 19 forma vključaje popravljanje každome → vsakom i každų → vsaku. Uže mnogočislna upotrěba vsaky podpira konsistenciju izdanje, no sama ne daje jezyčnu vlast nad izvorami.\par
\smallskip\noindent\emph{Reason:} Vsaky and každy are both type-1 dictionary headwords for each/every. Každy has positive western and eastern profile assessments but negatives in several southern languages; vsaky has broader southern support, with approximate Polish and Slovak assessments. That motivates preferring vsaky, not a measured improvement in comprehension. National každý/każdy and vsak/svaki remain comparators rather than candidates voted on by token volume. The hundred každy occurrences have a distributive role; replacing them with vse/vsi could alter scope, a collective reading or agreement. The nineteen-form table includes každome → vsakom and každų → vsaku. Existing extensive use of vsaky supports consistency within the edition without giving the edition linguistic authority over its sources.\par
\smallskip\noindent\emph{Izvory:} Slovnik, zaglavja 4326, 13864, 937.\par
\noindent\emph{Sources:} Dictionary entries 4326, 13864, 937.\par
\par\medskip\Needspace{8\baselineskip}
\hypertarget{DEC-0061}{}
\noindent\textbf{Mogu i mogut: čislo predikata / Mogu and mogut: number of the predicate}\quad{\footnotesize\texttt{DEC-0061}}\par
\smallskip\noindent\emph{Razlog:} Koncovka -u može označati prvo lice jedinstvenogo čisla; -ut označaje tretje lice množiny. Trideset sedm predikatov v zapisě ima množinny subjekt ili odnosno prědhodno slovo, zato dostavajut t. Avtorovo mogu i budu ostajut jedinstveny. Osoblivo poučny jest priměr, v ktorom jedna stroka imaje prvo mogu s avtorom kako subjektom i drugo mogu s množinnym subjektom: jednu formu trěba ostaviti, drugu změniti. Izbor zavisi od sintakse, ne od samogo slova. Korenj, aspekt i pravopis ostajut otdělne pytanja; takym ograničenjem se ne maskuje leksična změna kako popravka soglasja.\par
\smallskip\noindent\emph{Reason:} The ending -u can mark first-person singular, while -ut marks third-person plural. The recorded thirty-seven predicates have plural subjects or relative antecedents and therefore take final t. The author's mogu and budu remain singular. One particularly instructive line contains a first mogu referring to the author and a second with a plural subject: one is retained and the other changed. Syntax, not the word alone, determines the choice. Stem, aspect and spelling remain separate questions, preventing a lexical revision from being disguised as agreement repair.\par
\smallskip\noindent\emph{Izvory:} Gramatika, Glagoly, koncovky prezenta; zapis razlikuje 29 pojavjenij v člankah i osme v knigě.\par
\noindent\emph{Sources:} Grammar, Verbs, present endings; the record distinguishes 29 occurrences in the papers and eight in the book.\par
\par\medskip\Needspace{8\baselineskip}
\hypertarget{DEC-0062}{}
\noindent\textbf{Polny, cělkom, doskonaly / Polny, cělkom, doskonaly}\quad{\footnotesize\texttt{DEC-0062}}\par
\smallskip\noindent\emph{Razlog:} Poln- v polnoj reducibilnosti označaje tehnično polnotu; cělkom znači celostno ili v celosti; doskonal- označaje osoblivu vlastnost polja. Zapis razděljaje 94 pojavjenja popȯl-familije na 80 terminov polnoj reducibilnosti, 13 običnyh prislovnikov i jedno doskonalo polje. Německe vollständig a vollkommen ne možno sjediniti jedinym ekvivalentom bez čitanja definicije. Slovnik imaje popȯlno takože: njegova prisutnost ne znači ošibku, a tip 2 ne jest zabrana. Preferencija poln- pokračuje uže upotrěbjany termin; cělkom imaje razny profil po jazykah, ne všeobče pozitivne ocenky. Doskonaly ne znači dopolnjeny ili topologično polny. Smyslove razděljenje jest važnějše od povrhnostnoj jednobraznosti.\par
\smallskip\noindent\emph{Reason:} Poln- in complete reducibility expresses technical completeness; cělkom means entirely; doskonal- names the particular property of a perfect field. The record partitions 94 popȯl-family occurrences into 80 complete-reducibility terms, thirteen ordinary adverbs and one perfect-field predicate. German vollständig and vollkommen cannot share an automatic equivalent without reading the definition. Popȯlno is also in the dictionary: its presence is not an error and type 2 is not a prohibition. Poln- continues established terminology; cělkom has a mixed language profile rather than universally positive assessments. Doskonaly does not mean completed or topologically complete. The semantic distinction matters more than superficial uniformity.\par
\smallskip\noindent\emph{Izvory:} Slovnik, zaglavja 10921, 3141, 4391, 4207; članok o doskonalom polju.\par
\noindent\emph{Sources:} Dictionary entries 10921, 3141, 4391, 4207; the entry on perfect fields.\par
\par\medskip\Needspace{8\baselineskip}
\hypertarget{DEC-0063}{}
\noindent\textbf{Tvoret, zaviset, vhodet, vodet / Tvoret, zaviset, vhodet, vodet}\quad{\footnotesize\texttt{DEC-0063}}\par
\smallskip\noindent\emph{Razlog:} Tretje lice množiny vtorego sprezanja završaje se na -et. Formy tvoret, zaviset, vhodet i vodet sootvětstvujut množinnym subjektam v 51 zabilěženom kontekstě. Bez t formy tvore, zavise, vhode i vode ne sut čislovno neutralne alternativy etogo prezenta. No analogična zaměna vsakego konečnogo -e byla by pogrešna: prvo potrěbno ustanoviti subjekt, sprezanje i glagolnu ulogu. Otdělno ostajut slučaji leže, kde jedinstveno i množinno čislo potrěbujut razne popravky. Ni leksičny smysl, ni notacija ne daju povod za prěpisyvanje celej frazy, kogda dostava popraviti soglasje.\par
\smallskip\noindent\emph{Reason:} Third-person plurals of the second conjugation end in -et. Tvoret, zaviset, vhodet and vodet agree with plural subjects in the fifty-one recorded contexts. Bare tvore, zavise, vhode and vode are not number-neutral alternatives in this present-tense paradigm. But changing every final -e would be wrong: subject, conjugation and verbal function must first be established. Leže requires separate treatment because singular and plural contexts call for different corrections. Neither lexical meaning nor notation calls for rewriting the whole sentence when agreement alone needs repair.\par
\smallskip\noindent\emph{Izvory:} Gramatika, Glagoly, vtoro sprezanje, tretje lice množiny.\par
\noindent\emph{Sources:} Grammar, Verbs, second conjugation, third-person plural.\par
\par\medskip\Needspace{8\baselineskip}
\hypertarget{DEC-0064}{}
\noindent\textbf{Zato kako vyvod / Zato as an inference marker}\quad{\footnotesize\texttt{DEC-0064}}\par
\smallskip\noindent\emph{Razlog:} Zato označaje vyvod ili rezultat. Jedinnadset tedy/teda v članku o Maxu Noetheru imajut tu ulogu po zabilěženom sravnenju konteksta. Togda by prineslo vrěmenny ili uslovny smysl, a imenno by predstavilo pojasnjenje vměsto vyvoda. Eto ne jest razlog zabraniti tedy vo vsakoj situaciji: druge članki razlikujut posledicu, ekvivalenciju, apoziciju i uslovnu konstrukciju. Slovnik daje zato kako therefore, no ne daje izměrenu razumljivost vsej frazy. Konsistencija v izdani jest dodatny razlog za izbor, ne zaměna za ustanovjenje logičnoj ulogi.\par
\smallskip\noindent\emph{Reason:} Zato marks a conclusion or result. Eleven tedy/teda occurrences in the paper about Max Noether have this function according to the recorded contextual comparison. Togda would introduce temporal or conditional meaning, while imenno would turn a conclusion into an explanation. This does not justify banning tedy in every setting: other entries distinguish consequence, equivalence, apposition and conditional constructions. The dictionary gives zato as therefore, not as measured comprehension of the whole sentence. Consistency within the edition supports the choice only after its logical function has been established.\par
\smallskip\noindent\emph{Izvory:} Slovnik, zaglavje 3056; članok o logičnoj posledici.\par
\noindent\emph{Sources:} Dictionary entry 3056; the entry on logical consequence.\par
\par\medskip\Needspace{8\baselineskip}
\hypertarget{DEC-0065}{}
\noindent\textbf{Leži ili ležet / Leži or ležet}\quad{\footnotesize\texttt{DEC-0065}}\par
\smallskip\noindent\emph{Razlog:} Jedinstveny subjekt potrěbuje leži, množinny ležet. Zato isti naslědženy oblik leže ne može iměti jednu avtomatičnu zaměnu. Zapis razlikuje dva jedinstvene simbolične subjekty od jednogo množinnogo; dvanadset drugih predikatov s množinnym subjektom takože dostavajut -et. Čislo simbola ne slěduje iz togo, koliko bukv ili indeksov ima njegov zapis: potrěbno čitati, čto symbol označaje i kak on vstopaje v frazu. Proizhode, drže, različe, vyhode, proizvode i shvate sut vključeny po toj že sintaktičnoj proverci. Pravilo jest soglasje, ne povrhnostno izravnanje vseh podobnyh koncovok.\par
\smallskip\noindent\emph{Reason:} A singular subject requires leži and a plural subject ležet. The inherited form leže therefore cannot have one automatic replacement. The record distinguishes two singular symbolic subjects from one plural subject; twelve other plural predicates also take -et. A symbol's grammatical number is not determined by how many letters or subscripts it contains: what it denotes and how it functions in the sentence must be read. Proizhode, drže, različe, vyhode, proizvode and shvate were included through the same syntactic test. The rule is agreement, not surface uniformity among similar endings.\par
\smallskip\noindent\emph{Izvory:} Gramatika, Glagoly, singularno -i i pluralno -et.\par
\noindent\emph{Sources:} Grammar, Verbs, singular -i and plural -et.\par
\par\medskip\Needspace{8\baselineskip}
\hypertarget{DEC-0066}{}
\noindent\textbf{Oprěděliti, oprěděljati: definicija i opreděljenje / Oprěděliti, oprěděljati: definition and determination}\quad{\footnotesize\texttt{DEC-0066}}\par
\smallskip\noindent\emph{Razlog:} Oprěděliti i oprěděljati pokryvajut define, determine i specify; oprěděljeny i neoprěděljeny služat odvisno od konteksta. Izbor oprěděl- vměsto opreděl- slěduje pravopisu slovničnyh zaglavij. To ne dava pravo zaměniti cělu familiju na definiovati: opreděleny koeficient, neoprěděljena veličina i izrična definicija ne imajut jednu izključnu ulogu. Zapis odlišaje 544 forma, ktore dopuščajut pravopisnu popravku, od 15 s otdělnym problemom morfologije. Zaglavja sut tipa 2 i imajut raznu podpiru po jazykah; opis ne směje tvrditi, že vse izbrane formy sut tipa 1 ili jednako blizke vsem Slovanam.\par
\smallskip\noindent\emph{Reason:} Oprěděliti and oprěděljati cover define, determine and specify; oprěděljeny and neoprěděljeny serve context-dependent functions. Oprěděl- rather than opreděl- follows the dictionary headwords. This does not justify replacing the whole family with definiovati: a determined coefficient, an indeterminate quantity and an explicit definition do not have one exclusive function. The record separates 544 forms admitting spelling correction from fifteen with an additional morphological problem. These headwords are type 2 with varying language support; the account must not imply that every selected form is type 1 or equally accessible to all Slavic speakers.\par
\smallskip\noindent\emph{Izvory:} Slovnik, zaglavja 29544, 32606, 9608, 9609, 10237, 10238.\par
\noindent\emph{Sources:} Dictionary entries 29544, 32606, 9608, 9609, 10237, 10238.\par
\par\medskip\Needspace{8\baselineskip}
\hypertarget{DEC-0067}{}
\noindent\textbf{Oprěděljeny: oblik zavisi od konstrukcije / Oprěděljeny: form depends on the construction}\quad{\footnotesize\texttt{DEC-0067}}\par
\smallskip\noindent\emph{Razlog:} Padež i čislo nominalizovanyh neoprěděljenyh, pasivny particip i finitny glagol ne sut jedna kategorija. Zato 15 problematičnyh form oprěděl-familije potrěbuje čitanje frazy, ne tolko dodavanje ě. Glagolny imenovnik ostaje imenovnikom, a particip mora soglasiti se s opisyvanym objektom. V jednom slučaju nezasviděteljeno soopreděljeno bylo razloženo kako takože oprěděljeno. Analitična forma izbjegaje nejasny prefiks, ale njena věrnost zavisi od togo, jest li smysl ‘takože opreděljeno’, a ne jaka společna determinacija: istoričny zapis pripisyvaje izbrany smysl německomu kontekstu. Ovo jest razloženo kontekstno rěšenje, ne odobrenje produktivnogo so- vo vsakej tehničnoj složenici.\par
\smallskip\noindent\emph{Reason:} Nominalized indeterminates, passive participles and finite verbs belong to different grammatical categories. The fifteen problematic oprěděl-family forms therefore require sentence reading, not merely an added ě. A verbal noun remains a noun, and a participle must agree with the object described. One unattested soopreděljeno was rendered analytically as takože oprěděljeno. This avoids an opaque prefix, but its fidelity depends on the meaning being ‘also determined’, not a stronger claim of joint determination; the historical record attributes that reading to the German context. The decision is a contextual analysis, not approval of productive so- in every technical compound.\par
\smallskip\noindent\emph{Izvory:} Slovnik, oprěděl-familija; Gramatika, Glagoly i Pridavniky; soopreděljeno → takože oprěděljeno jest zabilěžena kontekstna para.\par
\noindent\emph{Sources:} Dictionary, oprěděl-family; Grammar, Verbs and Adjectives; soopreděljeno → takože oprěděljeno is the recorded contextual pair.\par
\par\medskip\Needspace{8\baselineskip}
\hypertarget{DEC-0068}{}
\noindent\textbf{Imenno: pojasnjenje termina / Imenno: explaining a term}\quad{\footnotesize\texttt{DEC-0068}}\par
\smallskip\noindent\emph{Razlog:} Imenno v nerazložimo vloženje, imenno prědstavjenje uvodi pojasnjenje predhodnogo termina. Zabilěžena německa fraza eine irreduzible Einbettung (Darstellung) podpiraje apoziciju: skobky nazyvajut istu konstrukciju drugim imenem. Zato by tvrdilo vyvod, a togda vrěmenny sled; ni jedno ne izražaje tu ulogu. Imenno jest slovnično namely, no je v tom članku izbrano po sintakse i smyslu skobok, ne po samomu nalěpku tipa 1. Ne slěduje iz togo, že vsako also ili tedy v matematičnom dokazě znači imenno.\par
\smallskip\noindent\emph{Reason:} In nerazložimo vloženje, imenno prědstavjenje, imenno explains the preceding term. The recorded German phrase eine irreduzible Einbettung (Darstellung) supports apposition: the parentheses name the same construction in another way. Zato would assert an inference and togda a temporal sequence; neither expresses this function. Imenno is the dictionary's namely, but the choice rests on syntax and the parentheses' meaning, not the type-1 label alone. It does not follow that every also or tedy in a proof means imenno.\par
\smallskip\noindent\emph{Izvory:} Slovnik, zaglavje 17486; německa fraza zachovana v zapisě odluki.\par
\noindent\emph{Sources:} Dictionary entry 17486; the German phrase preserved in the decision record.\par
\par\medskip\Needspace{8\baselineskip}
\hypertarget{DEC-0069}{}
\noindent\textbf{To jest: ekvivalencija i preformulovanje / To jest: equivalence and reformulation}\quad{\footnotesize\texttt{DEC-0069}}\par
\smallskip\noindent\emph{Razlog:} To jest uvodi ekvivalentno tvrdženje, definiciju ili pojasnjenje. Dvadeset sedm pojavjenij tedy/teda ima tu ulogu v zabilěženoj analizi; osme blizke pojavjenja v istih strokah ostajut markery vyvoda. Zato jedinica odluki ne jest stroka, ale konkretna část frazy. P → Q jest odlišna uloga od P, to jest Q: vtora konstrukcija predstavlja Q kako drugo izraženje P. Imenno može byti uže po smyslu, a togda dodaje vrěme. Sam německy also ne razrěšaje razliku; potrěbno čitati, jest li nasledna fraza vyvodom ili razjasnjenjem.\par
\smallskip\noindent\emph{Reason:} To jest introduces an equivalent statement, definition or reformulation. The recorded analysis assigns this function to twenty-seven tedy/teda occurrences while preserving eight neighbouring consequence markers on the same lines. The unit of decision is therefore the particular clause, not the line. P → Q differs from P, that is Q: the latter presents Q as another expression of P. Imenno can be narrower, while togda adds time. German also alone does not resolve the distinction; the following clause must be read as either conclusion or clarification.\par
\smallskip\noindent\emph{Izvory:} Kontekstna analiza sojuzov; sravnenje s imenno i zato v susědnyh člankah.\par
\noindent\emph{Sources:} Contextual analysis of connectives; compare imenno and zato in the related entries.\par
\par\medskip\Needspace{8\baselineskip}
\hypertarget{DEC-0070}{}
\noindent\textbf{To jest v algebričnyh definicijah / To jest in algebraic definitions}\quad{\footnotesize\texttt{DEC-0070}}\par
\smallskip\noindent\emph{Razlog:} Dvadset četyri konstrukcije v knigě upotrěbljajut to jest pri ekvivalenciji, apoziciji ili objašnjenju. Priměry v zapisě zahrnjajut rang jeden nad K i rang t² nad P, prěmo proizvedenje i prěmy dodavajemy člen, Galoisovo polje razloženja i matričny prsten. Rangi nad raznymi bazovymi poljami ne možno sravniti bez proměny bazy; to jest samo označaje načto razjasnjenje, ne dokazuje matematiku. Zato mora ostati tam, kde nasledno tvrdženje slěduje iz predhodnogo, ale ne povtarjaje jego. Osme takyh susědnyh markerov bylo zabilěženo otdělno. Nišči od etyh kontekstov ne byl klassifikovany kako vrěmenny sled.\par
\smallskip\noindent\emph{Reason:} Twenty-four constructions in the book use to jest for equivalence, apposition or explanation. Recorded examples concern rank one over K versus rank t² over P, direct products and direct summands, Galois splitting fields and matrix rings. Ranks over different base fields cannot be compared without accounting for the change of scalars; to jest signals clarification but does not itself prove the mathematics. Zato remains appropriate where a statement follows from its predecessor rather than restating it. Eight neighbouring markers were distinguished in this way. None of these contexts was classified as a sequence of events.\par
\smallskip\noindent\emph{Izvory:} Zabilěžene algebrične konteksty knigi; obče različenje vyvoda i ekvivalencije.\par
\noindent\emph{Sources:} Recorded algebraic contexts in the book; the general distinction between inference and equivalence.\par
\par\medskip\Needspace{8\baselineskip}
\hypertarget{DEC-0071}{}
\noindent\textbf{Libovoljny: libovoljno izbrany objekt / Libovoljny: an arbitrarily chosen object}\quad{\footnotesize\texttt{DEC-0071}}\par
\smallskip\noindent\emph{Razlog:} Libovoljny označaje objekt, ktory možno izbrati bez dodatnogo ograničenja v danom obsegu. To ne znači svobodny: svobodny modul imaje osoblivu algebričnu definiciju. Šestdeset form libovol-/proizvol- bylo sjedinjeno s uže upotrěbjanym libovolj-. Slovnik priznava libovoljny kako tip 3, proizvoljny kako tip 2; prvy ne jest zato oficijalny poběditelj vseobčego glasovanja. Češsky libovoln- i slověnsky poljubn- sut zabilěženy sravnenja, ale bez čislovoj tablicy vah ne možno pripisati im izračunany rezultat. Razlog izdanje jest smysl arbitrary spolu s konsistenciej padeža, čisla i roda; tvrditi izměrenu pan-slovjansku prěvagu bylo by silnějše od dokazov.\par
\smallskip\noindent\emph{Reason:} Libovoljny denotes an object chosen without further restriction within the stated domain. It does not mean free: a free module has a particular algebraic definition. Sixty libovol-/proizvol- forms were brought into line with existing libovolj-. The dictionary lists libovoljny as type 3 and proizvoljny as type 2; the former is not thereby an official winner of a universal vote. Czech libovoln- and Slovene poljubn- are recorded comparators, but without a numerical weighting table they cannot be assigned a computed result. The edition's reasons are the arbitrary sense and consistency of case, number and gender. A measured pan-Slavic advantage would be a stronger claim than the evidence supports.\par
\smallskip\noindent\emph{Izvory:} Slovnik, zaglavja 22300, 22318; česke i slověnske sravnenja citovane v zapisě.\par
\noindent\emph{Sources:} Dictionary entries 22300, 22318; Czech and Slovene comparisons cited in the record.\par
\par\medskip\Needspace{8\baselineskip}
\hypertarget{DEC-0072}{}
\noindent\textbf{Ako … to: uslovje i posledica / Ako … to: condition and consequence}\quad{\footnotesize\texttt{DEC-0072}}\par
\smallskip\noindent\emph{Razlog:} Ako … to označaje uslovnu konstrukciju: kogda jest izpolnjeno predpoloženje, slěduje posledica. Devet zabilěženih tedy/teda imajut ulogu to v apodozi, ne samostojnogo therefore. Togda može sugerovati vrěme, ktore tu ne jest trěbno. Zato može pridati silu samostojnogo vyvoda vměsto vtoroj časti uslovja. Konsultovana tablica sojuzov izrično dava ako kako if; ona sama ne dava polno sravnenje svih možnyh apodoz. Izbor to jest kontekstna interpretacija zabilěženogo německogo wenn/ist … so. Čitatelj mora razlikovati logično uslovje od vrěmennogo when, osoblivo kogda matematičny tekst upotrěbljaje dynamičny glagol.\par
\smallskip\noindent\emph{Reason:} Ako … to marks a conditional: the consequence holds under the hypothesis. Nine recorded tedy/teda occurrences function as the apodotic to, rather than an independent therefore. Togda may suggest time where none is needed; zato can give the force of a separate inference instead of the second half of the conditional. The consulted conjunction list explicitly gives ako as if, but does not itself compare every possible apodosis. Selecting to is a contextual interpretation of the recorded German wenn/ist … so constructions. Logical condition and temporal when must remain distinct, especially when mathematical prose uses a dynamic verb.\par
\smallskip\noindent\emph{Izvory:} Gramatika, Sojuzy, ako; slovnik, zaglavje 17076; zabilěžene uslovne frazy.\par
\noindent\emph{Sources:} Grammar, Conjunctions, ako; dictionary entry 17076; the recorded conditional sentences.\par
\par\medskip\Needspace{8\baselineskip}
\hypertarget{DEC-0073}{}
\noindent\textbf{Dobyvati i dobyty: dobivanje rezultata / Dobyvati and dobyty: obtaining a result}\quad{\footnotesize\texttt{DEC-0073}}\par
\smallskip\noindent\emph{Razlog:} Dobyvati jest nesovršeny glagol obtain; dobyti jest sovršeny parnjak. Dobyty, dobyta i dobyte sut participialne formy, ne finitne formy dobyvati. Zato 188 finitnyh i devet pasivnyh dobiv- form ne dopuščajut prosto i → y vo vsakom městě. Slovnik podpiraje oba aspekta, s širšimi sravnenjami v slovjanskyh jazykah; ni profil ni čislo pojavjenij ne jest test čitateljskoj razumljivosti. Zabilěženy konteksty sut polučenje matematičnogo rezultata, ne poběda, dostavka ili dostatočnost. Dobiv- jest normalno v nekotoryh nacionalnyh jazykah, ale rěšenje eto izdanje privodit k izbranomu medžuslovjanskomu zaglavju.\par
\smallskip\noindent\emph{Reason:} Dobyvati is imperfective obtain and dobyti its perfective counterpart. Dobyty, dobyta and dobyte are participial forms, not finite forms of dobyvati. The 188 finite and nine passive dobiv- forms therefore cannot all be handled by changing i to y. The dictionary supports both aspects with comparators across Slavic languages; neither those profiles nor occurrence counts constitute a reader-comprehension test. The recorded contexts concern obtaining mathematical results, not winning, delivery or sufficiency. Dobiv- is normal in some national languages, but this edition adopts the selected Interslavic headword.\par
\smallskip\noindent\emph{Izvory:} Slovnik, zaglavja 3191, 3192; Gramatika, Glagoly, aspekt i participy.\par
\noindent\emph{Sources:} Dictionary entries 3191, 3192; Grammar, Verbs, aspect and participles.\par
\par\medskip\Needspace{8\baselineskip}
\hypertarget{DEC-0074}{}
\noindent\textbf{Zato, to jest i tym samym / Zato, to jest and tym samym}\quad{\footnotesize\texttt{DEC-0074}}\par
\smallskip\noindent\emph{Razlog:} Pri razborě 155 tedy v knigě zapis razlikuje 151 vyvod, jednu ekvivalenciju i tri frazy s blizkym vyraženjem pričiny. Prvy slučaj dava zato, vtory to jest. V tretjem izbrano tym samym, aby ne skryti otnošenje medžu ‘slědovateljno’ i ‘poněže’ v nejasnoj posledovnosti zato … zato že. To ne jest zabranjene povtorjenje po estetike, ale potrěba razděljati pričinu i rezultat. Jedna ekvivalencija o porědku grupy i eksponentě ne mora postati dopolnitelno silnym tvrdženjem prosto zato, že drugy 151 token označaje vyvod. Čislovno bolša klassa ne upravjaje iskľučenjem.\par
\smallskip\noindent\emph{Reason:} Among 155 tedy occurrences in the book, the record distinguishes 151 conclusions, one equivalence and three sentences with an adjacent causal expression. These yield zato, to jest and, in the third setting, tym samym. The last choice prevents the relation between therefore and because from becoming obscure in a zato … zato že sequence. This is not an aesthetic ban on repetition but a distinction between reason and result. The group-order/exponent equivalence must not become a stronger claim merely because the other 151 tokens mark conclusions. The largest category does not govern its exceptions.\par
\smallskip\noindent\emph{Izvory:} Zabilěžena razdělba kontekstov; slovnik, zato; članok o pričinnym i rezultatnym sojuzah.\par
\noindent\emph{Sources:} The recorded contextual partition; dictionary, zato; the entry on causal and result connectives.\par
\par\medskip\Needspace{8\baselineskip}
\hypertarget{DEC-0075}{}
\noindent\textbf{Dobyvati ne jest dostavjati / Dobyvati is not dostavjati}\quad{\footnotesize\texttt{DEC-0075}}\par
\smallskip\noindent\emph{Razlog:} Dobyvati označaje dobivanje rezultata; dostaviti i dostavjati označajut dostavku ili davanje něčego. Slovnik takože priznava dostavati kako obtain, ale izdanje izbiraje jednu familiju dlja etogo smysla. Od 190 zabilěženih dostav-/dostaj- pojavjenij, 185 znamenajut dobivanje. Pet drugih ostajut: tri daty predloženja raboty, jedno dostavjeno dokazateljstvo i jedna fraza o davanyh modulah. Avtomatična zaměna celej familije by porušila eti smysly. Tip 2 pri dostavati ne čini glagol nelegalnym; semantično razděljenje i terminologična konsistencija objasnjajut preferenciju dobyvati.\par
\smallskip\noindent\emph{Reason:} Dobyvati expresses obtaining a result; dostaviti and dostavjati express supplying or providing something. The dictionary also accepts dostavati for obtain, but the edition chooses one family for that sense. Of 190 recorded dostav-/dostaj- occurrences, 185 mean obtaining. Five others are retained: three submission dates, a supplied proof and a clause about supplied modules. Replacing the whole family would damage those meanings. Type 2 does not make dostavati illegitimate; semantic separation and terminological consistency explain the preference for dobyvati.\par
\smallskip\noindent\emph{Izvory:} Slovnik, zaglavja 3191, 2511, 3105, 20.\par
\noindent\emph{Sources:} Dictionary entries 3191, 2511, 3105, 20.\par
\par\medskip\Needspace{8\baselineskip}
\hypertarget{DEC-0076}{}
\noindent\textbf{Največši: nepravilno stepenovanje / Največši: irregular comparison}\quad{\footnotesize\texttt{DEC-0076}}\par
\smallskip\noindent\emph{Razlog:} Največši slěduje paradigmi veliky → večši → največši. Gramatika dopuščaje takože bolši: izbrany oblik ne jest jedina legitimna možnost. Slovnično najvęćši daje v osnovnom pravopisu največši; ę i ć ne trebujut ostati v čitateljskom osnovnom zapisě. Četyrdeset pet pridavnikovyh form potrěbuje soglasje s objektom, ne prosto zaměnu korenja. Tri neutra množiny podpolja ne smějut dostati oblik ženskogo jedinstvenogo čisla. Prislovnik najčestěje prinaleži drugoj familiji i ne znači greatest. Slovničny profil največši ne jest univerzalno pozitiven; doslědnost paradigmy ne dokazuje, že vsaky govoriteli razumějut ju jednako.\par
\smallskip\noindent\emph{Reason:} Največši follows veliky → večši → največši. The grammar also permits bolši, so the selected form is not the only legitimate choice. Dictionary najvęćši becomes največši in basic spelling; ę and ć need not be retained there. Forty-five adjectival forms require agreement with their nouns, not merely a stem substitution. Three neuter-plural podpolja contexts must not receive a feminine singular form. Frequency najčestěje belongs to another family and does not mean greatest. Največši's dictionary profile is not universally positive; consistency of paradigm does not establish equal comprehension by all speakers.\par
\smallskip\noindent\emph{Izvory:} Slovnik, zaglavje 9364; Gramatika, Pridavniky, veliky/večši/bolši i naj-; Pravopis, osnovne i opcionalne bukvy.\par
\noindent\emph{Sources:} Dictionary entry 9364; Grammar, Adjectives, veliky/večši/bolši and naj-; Orthography, basic and optional letters.\par
\par\medskip\Needspace{8\baselineskip}
\hypertarget{DEC-0077}{}
\noindent\textbf{Najčestěje: čestota, ne veličina / Najčestěje: frequency, not magnitude}\quad{\footnotesize\texttt{DEC-0077}}\par
\smallskip\noindent\emph{Razlog:} Najčestěje znači most often. Naslědženo največestěje sjedinjaje časti familije največši i familije čestoty, ale to ne jest zabilěženo slovnično zaglavje. Jedna takova forma v korpusu ima prislovničnu ulogu o čestotě. Slovnik dava najčęstěje i čęstěje; osnovny zapis odstranjava opcionalno ę. Popravka mora ostaviti razliku medžu generally/most often i always: najčestěje ne čini tvrdženje bez iskľučenij. Izbor ne jest stylistično skraćenje, ale razděljenje dvuh smyslov — kako često i koliko veliko.\par
\smallskip\noindent\emph{Reason:} Najčestěje means most often. Inherited največestěje blends parts of the greatest and frequency families but is not an attested dictionary headword. The one corpus occurrence has an adverbial frequency function. The dictionary gives najčęstěje and čęstěje; basic spelling removes optional ę. The correction must preserve the difference between generally/most often and always: najčestěje does not make an exceptionless claim. This is not stylistic shortening but separation of two meanings—how often and how large.\par
\smallskip\noindent\emph{Izvory:} Slovnik, zaglavja 20519, 19669; Pravopis, opcionalno ę.\par
\noindent\emph{Sources:} Dictionary entries 20519, 19669; Orthography, optional ę.\par
\par\medskip\Needspace{8\baselineskip}
\hypertarget{DEC-0078}{}
\noindent\textbf{Prinajmenje i najmenše obče mnogokratno / Prinajmenje and the least common multiple}\quad{\footnotesize\texttt{DEC-0078}}\par
\smallskip\noindent\emph{Razlog:} Prinajmenje označaje dolnu granicu: at least. Najmenše obče mnogokratno nazyvaje nejmenši objekt, ktory izpolnjaje uslovje byti obče mnogokratno. Zapis razděljaje 172 granične prislovniky i 20 pridavnikovyh terminov; jedno uže prisutno prinajmenje služi kako sravnenje. Německe mindestens/wenigstens ne jest to samo kako kleinstes gemeinsames Vielfaches. Forma najmenše sootvětstvuje srednjemu rodu mnogokratno; slovnično najmėnje jest prislovnik, ne vměsto pridavnika. Pravilo maly → menši i naj- daje oblik, ale matematična definicija opreděljaje, kogda treba upotrěbiti jego. Zato nikakova globalna zaměna vseh minimum-form ne može zaměniti analiz.\par
\smallskip\noindent\emph{Reason:} Prinajmenje expresses a lower bound: at least. Najmenše obče mnogokratno names the least object satisfying the common-multiple condition. The record separates 172 bounding adverbs from twenty adjectival technical terms, with one existing prinajmenje as a comparator. German mindestens/wenigstens is not kleinstes gemeinsames Vielfaches. Najmenše agrees with neuter mnogokratno; dictionary najmėnje is an adverb, not a substitute adjective. Maly → menši with naj- supplies the form, but the mathematical definition determines when to use it. No global replacement of minimum-related forms can replace that analysis.\par
\smallskip\noindent\emph{Izvory:} Slovnik, zaglavja 390, 4718; Gramatika, Pridavniky, maly → menši i naj-.\par
\noindent\emph{Sources:} Dictionary entries 390, 4718; Grammar, Adjectives, maly → menši and naj-.\par
\par\medskip\Needspace{8\baselineskip}
\hypertarget{DEC-0079}{}
\noindent\textbf{To jest pri formulah i definicijah / To jest with formulas and definitions}\quad{\footnotesize\texttt{DEC-0079}}\par
\smallskip\noindent\emph{Razlog:} Petnadset dodatnyh konstrukcij upotrěbljajut to jest dlja ekvivalentnyh formula, definicij, indeksov, razloženih algebr, prstenov, idealov ili porědkov. Zapis razlikuje ih od vyvodov: sledujuča fraza ne pridaje novo slědstvo, ale imenovaje ili preformuluje isto. Eto različenje jest lingvistično i matematično; formalna ekvivalencija ne směje byti izmišljena, kogda německy tekst dava jedino implikaciju. Zato i togda ostajut nevhodne dlja etoj klassy, ale ne sut zabranjene slova. Vse ostatne pojavjenja istogo povrhnostnogo markera potrěbujut svoju klassifikaciju.\par
\smallskip\noindent\emph{Reason:} Fifteen further constructions use to jest for equivalent formulas, definitions, indices, split algebras, rings, ideals or orders. The record distinguishes them from inferences: the following phrase names or reformulates the same thing rather than adding a consequence. This is both a linguistic and a mathematical distinction; formal equivalence must not be invented where the German gives only an implication. Zato and togda are unsuitable for this class, not prohibited words. Other occurrences of the same surface marker require their own classification.\par
\smallskip\noindent\emph{Izvory:} Zabilěžena analiza definicij i ekvivalencij; članok To jest.\par
\noindent\emph{Sources:} The recorded analysis of definitions and equivalences; the To jest entry.\par
\par\medskip\Needspace{8\baselineskip}
\hypertarget{DEC-0080}{}
\noindent\textbf{Vyše i najvyše: sravnenje i gorna granica / Vyše and najvyše: comparison and an upper bound}\quad{\footnotesize\texttt{DEC-0080}}\par
\smallskip\noindent\emph{Razlog:} Vyše imaje v slovniku smysly higher i more; fraza opreděljaje, ktory jest aktivny. Najvyše može v količestvenom kontekstě označati gornu granicu, kako najvyše ρ+1 funkcij, ne najveću dosěgnutu veličinu. Četyrdeset devet samostojnyh više/najviše form bylo privedeno k vyše/najvyše; nacionalne sravnenja ostajut validne v svojih jazykah. Gramatika dopuščaje takože bolje/najbolje pri analitičnom stepenovanju. Zato vyše ne jest jedina možnost, a slovnično most samo po sobě ne dokazuje smysl at most. Ograničenje na prislovniky čuvaje nepričastne pridavnike i glagolne koreni. Pri negativnoj frazě treba otdělno razlikovati no more od never.\par
\smallskip\noindent\emph{Reason:} The dictionary gives vyše both higher and more; the sentence selects the active sense. In a quantitative setting, najvyše can express an upper bound, as in at most ρ+1 functions, rather than an attained maximum. Forty-nine standalone više/najviše forms were standardized as vyše/najvyše; national comparators remain valid in their languages. The grammar also allows bolje/najbolje for analytic comparison. Vyše is therefore not the sole option, and dictionary most alone does not establish at most. Restricting the rule to adverbs protects unrelated adjectives and verbal roots. Negative sentences must distinguish no more from never.\par
\smallskip\noindent\emph{Izvory:} Slovnik, zaglavja 2508, 2680, 1488; Gramatika, Pridavniky, analitično stepenovanje; Teorema II o racionalnoj bazě polja.\par
\noindent\emph{Sources:} Dictionary entries 2508, 2680, 1488; Grammar, Adjectives, analytic comparison; Theorem II on a field's rational basis.\par
\par\medskip\Needspace{8\baselineskip}
\hypertarget{DEC-0081}{}
\noindent\textbf{Zato v dokazě / Zato in a proof}\quad{\footnotesize\texttt{DEC-0081}}\par
\smallskip\noindent\emph{Razlog:} Zato uvodi tvrdženje, ktore slěduje iz prědhodnyh predpoloženij ili argumenta. Dvadeset pet zabilěženih kontekstov zahrnjaje takože konstrukcije s Nehaj: sam počatok s proizvoljnym objektom ne čini nasledny vyvod definiciej. To jest by nepravilno predstavilo rezultat kako preformulovanje; togda by moglo suggerovati vrěme. Velika bukva Zato na počatku frazy ne jest novy termin. Pravidlo čuva avtorsku strukturu dokazanja i ne proglašaje formalnu ekvivalenciju tam, kde izvor dava jedino slědstvo.\par
\smallskip\noindent\emph{Reason:} Zato introduces a statement following from preceding assumptions or argument. The twenty-five recorded contexts include Nehaj constructions: starting with an arbitrary object does not turn a subsequent inference into a definition. To jest would wrongly present the result as a reformulation; togda could suggest time. Sentence-initial Zato is not a new term. The rule preserves the author's proof structure without claiming equivalence where the source gives only a consequence.\par
\smallskip\noindent\emph{Izvory:} Slovnik, zaglavje 3056; zabilěžene německe also/so/folglich konteksty.\par
\noindent\emph{Sources:} Dictionary entry 3056; the recorded German also/so/folglich contexts.\par
\par\medskip\Needspace{8\baselineskip}
\hypertarget{DEC-0082}{}
\noindent\textbf{Bolje ne jest bolši / Bolje is not bolši}\quad{\footnotesize\texttt{DEC-0082}}\par
\smallskip\noindent\emph{Razlog:} Bolje jest prislovnik stepena, bolši jest pridavnik. Pet bolj i jedno boljše v zapisě modifikujut stepenj ili sravnenje, zato izbrano bolje. Formy bolšim, bolšego i bolšej ostajut pridavnikove padežne formy; globalna proměna korenja by porušila soglasje. Gramatika izrično dopuščaje bolje i vyše pri analitičnom sravnenju. Slovničny profil bolje imaje pozitivne vostočne i slověnske, ale mnogo negativnyh ocenok drugde. Zato morfologična opravdanost ne mora byti opisana kako universalna leksična prěvaga. Izbor jest doslědny oblik v konkretnom prislovničnom upotrěbjenju.\par
\smallskip\noindent\emph{Reason:} Bolje is a degree adverb; bolši is an adjective. Five bolj occurrences and one boljše modify degree or comparison and were rendered as bolje. Bolšim, bolšego and bolšej remain adjectival case forms; a root-wide replacement would damage agreement. The grammar explicitly permits both bolje and vyše in analytic comparison. Bolje's dictionary profile has eastern and Slovene positives but numerous negatives elsewhere. Morphological justification must therefore not be presented as universal lexical superiority. The decision establishes the appropriate form in these adverbial contexts.\par
\smallskip\noindent\emph{Izvory:} Slovnik, zaglavje 16915; Gramatika, Pridavniky, analitično i nepravilno stepenovanje.\par
\noindent\emph{Sources:} Dictionary entry 16915; Grammar, Adjectives, analytic and irregular comparison.\par
\par\medskip\Needspace{8\baselineskip}
\hypertarget{DEC-0083}{}
\noindent\textbf{Vyvod i fraza na že / Inference and a že clause}\quad{\footnotesize\texttt{DEC-0083}}\par
\smallskip\noindent\emph{Razlog:} Trideset tri markery vyvoda v etoj gruppě izbirajut zato. Dvě blizke konstrukcije tedy, že ne podležat tomu prostomu pravilu: tam samostojna zaměna ne objasnjaje, kak že-fraza zavisi od glavnogo predikata. Potrěbno razlikovati ‘iz etogo slěduje tvrdženje’ od ‘požadanje znači, že…’ ili ‘prědpokladajemo, že…’. Eto objasnjaje, začto podobna okolica ili isti německy sojuz ne dostava dlja mehaničnoj zaměny. Razlika jest v sintaktičnoj konstrukciji, ne v raznom slovniku dlja raznyh stranic.\par
\smallskip\noindent\emph{Reason:} Thirty-three inference markers in this group take zato. Two neighbouring tedy, že constructions do not admit the same simple rule: replacing the token alone would not explain how the že-clause depends on the main predicate. A conclusion must be distinguished from ‘the requirement means that…’ or ‘we assume that…’. Similar surrounding wording or the same German connective is insufficient for mechanical substitution. The difference lies in syntax, not in having different vocabularies for different pages.\par
\smallskip\noindent\emph{Izvory:} Zabilěžene německe vyvody; članok o Zato požadanje znači, že.\par
\noindent\emph{Sources:} The recorded German inferences; the entry on Zato požadanje znači, že.\par
\par\medskip\Needspace{8\baselineskip}
\hypertarget{DEC-0084}{}
\noindent\textbf{Zato požadanje znači, že / Zato požadanje znači, že}\quad{\footnotesize\texttt{DEC-0084}}\par
\smallskip\noindent\emph{Razlog:} Německe Die Forderung besagt also, daß dava Zato požadanje znači, že; Wir nehmen also an dava Zato prědpokladajemo, že. Peršej frazě vlada znači, vtoroj prědpokladajemo. Přesun markera vyvoda k počatku frazy izjasnjaje konstrukciju, ne dodava novo predpoloženje. Dvě takove konstrukcije treba razlikovati od 45 običnyh zato/Zato, kde změna porědka ne jest potrěbna. To jest ne može zaměniti zato, kogda avtor vyvodit, čto treba prědpoložiti; samo slovo assume ne čini celu frazu definiciej.\par
\smallskip\noindent\emph{Reason:} German Die Forderung besagt also, daß is rendered as Zato požadanje znači, že; Wir nehmen also an as Zato prědpokladajemo, že. The first subordinate clause depends on means, the second on assume. Moving the inference marker to the beginning clarifies the construction without adding a new assumption. These two sentences differ from forty-five ordinary zato/Zato cases needing no rearrangement. To jest cannot replace zato where the author concludes what must be assumed; assume alone does not make the whole sentence a definition.\par
\smallskip\noindent\emph{Izvory:} Dvě německe izvorne frazy zachovane v zapisě; slovnik, zato.\par
\noindent\emph{Sources:} The two German source phrases preserved in the record; dictionary, zato.\par
\par\medskip\Needspace{8\baselineskip}
\hypertarget{DEC-0085}{}
\noindent\textbf{Ponovno: povtorjenje / Ponovno: repetition}\quad{\footnotesize\texttt{DEC-0085}}\par
\smallskip\noindent\emph{Razlog:} Ponovno označaje again/once more. Dvadeset četyri znovo konteksty zahrnjajut povtorno označenje, podstavjenje, priměnjenje, predpoloženje ili povtorno pojavjenje vlastnosti. Ni jeden ne byl opisan kako obnova od samogo počatka. Zato prěklad ne mora uvesti silnějši smysl ‘nanovo’. Slovnik priznava ponovno s razprostranenymi sravnenjami; znovu i iznova ne treba objevjati nevozmožnymi sinonimami. Jedna predpočtana forma v izdani jest terminologična politika, ne dokaz, že tolko ona možet izražati povtorjenje.\par
\smallskip\noindent\emph{Reason:} Ponovno expresses again/once more. The twenty-four znovo contexts include repeated notation, substitution, application, assumption or recurrence of a property. None was described as starting afresh from the beginning, so the translation must not introduce that stronger sense. The dictionary supports ponovno with widespread comparators; znovu and iznova need not be declared impossible synonyms. A preferred form within the edition is a consistency policy, not proof that only this word can express repetition.\par
\smallskip\noindent\emph{Izvory:} Slovnik, zaglavje 5961; člen o ponovno v osnovnyh normalizacijah.\par
\noindent\emph{Sources:} Dictionary entry 5961; ponovno in the basic normalization conventions.\par
\par\medskip\Needspace{8\baselineskip}
\hypertarget{DEC-0086}{}
\noindent\textbf{V tom slučaju: uslovna uloga Dann / V tom slučaju: conditional Dann}\quad{\footnotesize\texttt{DEC-0086}}\par
\smallskip\noindent\emph{Razlog:} V tom slučaju v frazě o I. Šuru znači ‘pod tym predpoloženjem’. Německy tekst najprvo zadava hypotetično Es enthalte P … i potom govori Dann zeigte I. Schur, daß …. Danne dokazanje pripisano jest Šuru, ale Dann v etoj konstrukciji ne označaje datu ili sled istorijnyh sobytij. Analitična fraza čuva uslovnu ulogu bez dvojsmyslenosti togda. Zato by moglo prestaviti prosto kauzalno ili dokazateljno slědstvo, a ne obseg predpoloženja. Ime, avtorsko pripisanje i prošly čas ostajut častju izvornoj prozy.\par
\smallskip\noindent\emph{Reason:} V tom slučaju in the Schur sentence means under that hypothesis. The German first introduces hypothetical Es enthalte P … and then Dann zeigte I. Schur, daß …. The result is attributed to Schur, but Dann here does not date a historical event. The analytic phrase retains the conditional scope without the ambiguity of togda. Zato could suggest an ordinary causal or argumentative transition instead of the hypothesis's scope. The name, attribution and past tense remain part of the source prose.\par
\smallskip\noindent\emph{Izvory:} Zabilěžena německa fraza Es enthalte P … Dann zeigte I. Schur, daß …; analiza uslovnoj konstrukcije.\par
\noindent\emph{Sources:} The recorded German Es enthalte P … Dann zeigte I. Schur, daß …; analysis of the conditional construction.\par
\par\medskip\Needspace{8\baselineskip}
\hypertarget{DEC-0087}{}
\noindent\textbf{Najprvo: porędok izlaganja i prvenstvo / Najprvo: order of exposition and priority}\quad{\footnotesize\texttt{DEC-0087}}\par
\smallskip\noindent\emph{Razlog:} Najprvo znači first of all i služi porědku izlaganja. Od 26 zabilěženih pojavjenij, 25 označajut postupok ili diskursny porędok; odkaz na Ostrowskogo označaje, kto prvy postavil pytanje. To poslednje ne jest jedino stilističny marker: avtorsko prvenstvo mora ostati čitateljno. Raněje by značilo earlier, naprěd in advance, a ne nužno first. Prvo može v nekotoryh frazah byti dovoljno, zato tvrdženje, že ono vsegda gubi smysl, bylo by prěsilno. Izbor najprvo slěduje slovniku i doslědnosti, s očuvanjem razlike medžu porědkom argumenta i istorijnym prvenstvom.\par
\smallskip\noindent\emph{Reason:} Najprvo means first of all and orders exposition. Of twenty-six recorded occurrences, twenty-five concern procedure or discourse order; the Ostrowski attribution identifies who first posed the question. That last use is not merely a stylistic marker: scholarly priority must remain legible. Raněje would mean earlier and naprěd in advance, neither necessarily first. Prvo can suffice in some sentences, so claiming it always loses the meaning would be too strong. Najprvo follows the dictionary and consistency policy while preserving the distinction between argumentative order and historical priority.\par
\smallskip\noindent\emph{Izvory:} Slovnik, zaglavje 19868; zabilěžen odkaz na Ostrowskogo.\par
\noindent\emph{Sources:} Dictionary entry 19868; the recorded Ostrowski attribution.\par
\par\medskip\Needspace{8\baselineskip}
\hypertarget{DEC-0088}{}
\noindent\textbf{Različnyh, različne, različnoj / Različnyh, različne, različnoj}\quad{\footnotesize\texttt{DEC-0088}}\par
\smallskip\noindent\emph{Razlog:} Različny jest tvrdokorenny pridavnik. Devetnadset genitivov množiny dava različnyh; tri nominativy množiny bez mužskogo oduševljenogo subjekta davajut različne; tri ženske kosye formy davajut različnoj. Značenje different/various ostaje, ale koncovka zavisi od imenovnika. Jednako pravilo dlja svih 25 forma by izgubilo rod, čislo ili padež. Slovnik dava råzličny, a polna gramatična tablica podpiraje etu razliku koncovok. Različnost jest otdělny imenovnik i ne podleži pridavnikovomu pravilu.\par
\smallskip\noindent\emph{Reason:} Različny is a hard-stem adjective. Nineteen genitive plurals yield različnyh, three non-masculine-animate nominative plurals različne, and three feminine oblique forms različnoj. The meaning different/various is retained, but the ending depends on the noun. One replacement for all twenty-five forms would lose gender, number or case. The dictionary gives råzličny, and the complete grammatical table supports the ending distinctions. Različnost is a separate noun, not subject to the adjective rule.\par
\smallskip\noindent\emph{Izvory:} Slovnik, zaglavje 2211; Gramatika, Pridavniky, polna tablica koncovok.\par
\noindent\emph{Sources:} Dictionary entry 2211; Grammar, Adjectives, complete table of endings.\par
\par\medskip\Needspace{8\baselineskip}
\hypertarget{DEC-0089}{}
\noindent\textbf{Tym samym pri pričinnom objašnjenju / Tym samym alongside a causal explanation}\quad{\footnotesize\texttt{DEC-0089}}\par
\smallskip\noindent\emph{Razlog:} Tym samym v dvuh frazah čuva ulogu vyvoda, kogda blizka že-konstrukcija dava pričinu. Druge 44 markery v toj gruppě davajut zato/Zato. Izbor ne dava novo logično pravilo: on razděljuje ‘iz etogo slěduje’ od ‘poněže’, ktore mogut byti v jednoj frazě. Povtorjenje zato samo po sobě ne jest gramatična ošibka; potrěbno pokazati nejasnost konkretnogo otnošenja. Zapis motivuje analitično razjasnjenje, ne univerzalnu zabranu povtorjenja. Togda ne izražaje tuto dokazateljnu ulogu.\par
\smallskip\noindent\emph{Reason:} Tym samym preserves inference in two sentences where an adjacent že-construction supplies the reason. The other forty-four markers in that group take zato/Zato. This introduces no new logical rule: it separates follows from this from because, which can occur in one sentence. Repetition of zato is not intrinsically ungrammatical; the particular relation must be shown to be unclear. The recorded rationale supports analytic clarification, not a universal ban on repetition. Togda does not express the argumentative function here.\par
\smallskip\noindent\emph{Izvory:} Zabilěžene pričinne frazy; slovnik, zato; sravnenje s tym samym.\par
\noindent\emph{Sources:} The recorded causal sentences; dictionary, zato; comparison with tym samym.\par
\par\medskip\Needspace{8\baselineskip}
\hypertarget{DEC-0090}{}
\noindent\textbf{Raznyh razkladov / Raznyh razkladov}\quad{\footnotesize\texttt{DEC-0090}}\par
\smallskip\noindent\emph{Razlog:} V izomorfija raznyh razkladov pridavnik soglasuje se s genitivom množiny razkladov. Tvrdokorenny razny zato dava raznyh, ne raznih. Izbor ne menja, čto sut razne razklady, i ne tvrdi, že oni sut po vseh matematičnyh vlastnostjah različno ustrojene. Slovnik dava råzny kako different/various; koncovku dava sklonjenje. Popravka jednoj frazy ne podpiraje zaměnu vseh -ih: mekke pridavniky i zaimena imajut svoje legitimne i-formy.\par
\smallskip\noindent\emph{Reason:} In izomorfija raznyh razkladov, the adjective agrees with genitive-plural razkladov. Hard-stem razny therefore takes raznyh, not raznih. This does not alter which decompositions are distinct or claim that they differ in every mathematical property. The dictionary supplies råzny as different/various; declension supplies the ending. Correcting this phrase does not justify replacing all -ih endings: soft adjectives and pronouns have legitimate i-forms.\par
\smallskip\noindent\emph{Izvory:} Slovnik, zaglavje 2034; Gramatika, Pridavniky, genitiv množiny.\par
\noindent\emph{Sources:} Dictionary entry 2034; Grammar, Adjectives, genitive plural.\par
\par\medskip\Needspace{8\baselineskip}
\hypertarget{DEC-0091}{}
\noindent\textbf{Danyh i oddělnyh: -nyh pri tvrdem n / Danyh and oddělnyh: -nyh after hard n}\quad{\footnotesize\texttt{DEC-0091}}\par
\smallskip\noindent\emph{Razlog:} Tvrdokorenny pridavnik ili pasivny particip na -ny/-eny/-jeny sklanjaje se kako pridavnik: v genitivě i lokativě množiny dava -nyh/-enyh/-jenyh. Zapis zahrnjaje 63 pojavjenja v 19 formah. Zato završno n ne opravdava nacionalno -nih v etoj paradigmě. No -ih posle mekkogo korenja ostaje pravilo v drugih slovah. Gramatika glagolov objasnjaje tvorjenje pasivnogo participa, a pridavnikova tablica njegovo soglasje; ni jedna od tih dvuh informacij sama ne dostava dlja čitanja cěloj frazy.\par
\smallskip\noindent\emph{Reason:} A hard-stem adjective or passive participle in -ny/-eny/-jeny declines adjectivally, giving genitive/locative-plural -nyh/-enyh/-jenyh. The record covers 63 occurrences across nineteen forms. Final n does not license inherited -nih in this paradigm. But -ih after a soft stem remains regular elsewhere. The verb grammar explains passive-participle formation and the adjective table its agreement; neither piece of information alone replaces reading the sentence.\par
\smallskip\noindent\emph{Izvory:} Gramatika, Glagoly, pasivny particip; Pridavniky, tvrdo i mekko sklonjenje.\par
\noindent\emph{Sources:} Grammar, Verbs, passive participles; Adjectives, hard and soft declension.\par
\par\medskip\Needspace{8\baselineskip}
\hypertarget{DEC-0092}{}
\noindent\textbf{Mnogyh: množina, ne novy kvantifikator / Mnogyh: inflection, not a new quantifier}\quad{\footnotesize\texttt{DEC-0092}}\par
\smallskip\noindent\emph{Razlog:} Mnogyh jest genitiv ili lokativ množiny od mnogy. Završno g ne jest medžu mekkymi korenjami pridavnikov š/ž/č/j, zato i ne zaměnjuje y. Zapis dava 142 konteksty s mnogih, vse s množinnym otnošenjem. Popravka ne prevodi many na all, ni na ‘bezkoněčno mnogo’: količestveny smysl ostaje izvorno ograničeny. Mnogo kako prislovnik ili količestveno slovo imaje drugu konstrukciju i ne směje byti složeno s paradigmom mnogy tolko zaradi podobnogo korenja.\par
\smallskip\noindent\emph{Reason:} Mnogyh is the genitive or locative plural of mnogy. Final g is not among the adjective soft stems š/ž/č/j, so it does not change y to i. The record gives 142 mnogih contexts, all with plural reference. The correction does not turn many into all or infinitely many: the original quantitative scope remains. Adverbial or quantifying mnogo has a different construction and must not be merged with the mnogy paradigm merely because the root resembles it.\par
\smallskip\noindent\emph{Izvory:} Slovnik, zaglavje 17077; Gramatika, Pridavniky, tvrde koreni.\par
\noindent\emph{Sources:} Dictionary entry 17077; Grammar, Adjectives, hard stems.\par
\par\medskip\Needspace{8\baselineskip}
\hypertarget{DEC-0093}{}
\noindent\textbf{Vyvod, ekvivalencija, uslovje: tri ulogi / Inference, equivalence, condition: three functions}\quad{\footnotesize\texttt{DEC-0093}}\par
\smallskip\noindent\emph{Razlog:} Zabilěžena polna razdělba 478 tedy/teda razlikuje 396 vyvodov, 72 ekvivalencije ili pojasnjenja i deset uslovnyh apodoz. Nijedna ne byla klassifikovana kako sled sobytij ili prazdny diskursny marker. Izbrany obče sootvětnosti sut zato, to jest i to/v tom slučaju, s analitičnymi frazami tam, kde pričina i rezultat susědujut. Čisla opisujut etot korpus, ne vse upotrěbjenja slov v medžuslovjanskom. Poslednja gruppa 84 slučajev sodrži 80 vyvodov i četyri pojasnjenja. Metoda jest ustanoviti ulogu frazy do izbora markera; največša klassa ne može zaměniti razbor ostatnyh.\par
\smallskip\noindent\emph{Reason:} The recorded full partition of 478 tedy/teda occurrences distinguishes 396 inferences, 72 equivalences or explanations, and ten conditional apodoses. None was classified as event sequence or empty discourse filler. The principal equivalents are zato, to jest and to/v tom slučaju, with analytic wording where reason and result sit together. These counts describe this corpus, not every possible Interslavic use. The final group of 84 contains eighty inferences and four explanations. The method determines the clause's function before selecting a marker; the largest category cannot replace analysis of the others.\par
\smallskip\noindent\emph{Izvory:} Polna zabilěžena razdělba 478 kontekstov; članki Zato, To jest i Ako … to.\par
\noindent\emph{Sources:} The complete recorded partition of 478 contexts; the Zato, To jest and Ako … to entries.\par
\par\medskip\Needspace{8\baselineskip}
\hypertarget{DEC-0094}{}
\noindent\textbf{Tyh: pokazateljno, ne tihy / Tyh: demonstrative, not tihy}\quad{\footnotesize\texttt{DEC-0094}}\par
\smallskip\noindent\emph{Razlog:} Tyh jest genitiv/lokativ množiny zaimena toj. Polna gramatična tablica dava takože tym i tymi; propratna nota dopuščaje alternaty těh, těm i těmi. Zapis razlikuje 42 demonstrativne tih/Tih otnošenja od pridavnika tihy ‘quiet’. Zato ovo ne jest pravilo, že bukvy tih vsegda sut ošibka. Izbrana forma čuva ukazanje na iste matematične objekty. Otdělna možnost těh ne jest odvržena: ona ne byla cilj etoj normalizacije i ne protivoreči gramatičnomu izvoru.\par
\smallskip\noindent\emph{Reason:} Tyh is the genitive/locative plural of toj. The complete table also gives tym and tymi; its note allows těh, těm and těmi as alternatives. The record distinguishes forty-two demonstrative tih/Tih occurrences from the adjective tihy, quiet. This is therefore not a rule that the letters tih are always erroneous. The selected form retains reference to the same mathematical objects. The alternative těh is not rejected: it was outside this normalization and is compatible with the grammar.\par
\smallskip\noindent\emph{Izvory:} Jan van Steenbergen, Gramatika, Zaimena, toj i alternativy; slovnik, zaglavje 1844.\par
\noindent\emph{Sources:} Jan van Steenbergen, Grammar, Pronouns, toj and its alternatives; dictionary entry 1844.\par
\par\medskip\Needspace{8\baselineskip}
\hypertarget{DEC-0095}{}
\noindent\textbf{Kromě: iskľučenje i dopolnjenje / Kromě: exception and addition}\quad{\footnotesize\texttt{DEC-0095}}\par
\smallskip\noindent\emph{Razlog:} Kromě s genitivom označaje except/apart from. Petnadset osim i dva krome byli privedeny k tomu zaglavju. No apart from može v kontekstě izključati ili pridavati: ‘kromě X, takože Y’ ne znači, že X jest isključen iz vseh razgledanyh objektov. Tuto razliku mora čuvati cela fraza. Poręd ‘beside’ ne jest prosta sinonimna zaměna pri matematičnom iskľučenju. Slovnik podpira kromě, ale jego profil ne jest jednako pozitiven dlja vseh jazykov. Konsistencija i upravjenje genitivom objasnjajut izbor bez izmišljena merenja razumljivosti.\par
\smallskip\noindent\emph{Reason:} Kromě governs the genitive and means except/apart from. Fifteen osim occurrences and two krome were standardized to that headword. But apart from may exclude or add: ‘apart from X, also Y’ does not exclude X from every object under consideration. The full sentence must preserve that distinction. Poręd, beside, is not an automatic synonym in mathematical exception clauses. The dictionary supports kromě, but its language profile is not uniformly positive. Consistency and genitive government explain the choice without inventing a comprehension measurement.\par
\smallskip\noindent\emph{Izvory:} Slovnik, zaglavje 318; kontekstna razlika medžu iskľučenjem i dodatkom.\par
\noindent\emph{Sources:} Dictionary entry 318; the contextual distinction between exception and addition.\par
\par\medskip\Needspace{8\baselineskip}
\hypertarget{DEC-0096}{}
\noindent\textbf{-yh v pridavnikah, zaimenah i participah / -yh in adjectives, determiners and participles}\quad{\footnotesize\texttt{DEC-0096}}\par
\smallskip\noindent\emph{Razlog:} Osemdeset pojavjenij v dvadeseti formah slěduje tvrdemu genitivu/lokativu množiny. Sravnenja vključajut drugih → drugyh, ktorih → ktoryh, prvih → prvyh, takih → takyh i vzętih → vzętyh. Drugy/prvy sut porědkove čislovniky s pridavnikovym sklonjenjem; ktory jest pridavnikovo zaimeno; vzęty jest pasivny particip. Jednaka koncovka zato ne znači jednaku leksičnu kategoriju. Mekke pridavniky, osobne zaimena, někoj- i oddělne čislovnične paradigmy ostajut izvan pravila. Polna tablica vseh 20 par ostaje v sravnitelnyh dannyh, aby priměry ne zaměnili obseg.\par
\smallskip\noindent\emph{Reason:} Eighty occurrences across twenty forms follow hard-stem genitive/locative plural inflection. Examples include drugih → drugyh, ktorih → ktoryh, prvih → prvyh, takih → takyh and vzętih → vzętyh. Drugy/prvy are adjectivally declined ordinals; ktory is an adjectival determiner; vzęty is a passive participle. A shared ending therefore does not imply one lexical category. Soft adjectives, personal pronouns, někoj- and distinct numeral paradigms remain outside the rule. The full twenty-pair mapping is retained in the comparison data so the examples do not substitute for coverage.\par
\smallskip\noindent\emph{Izvory:} Gramatika, Pridavniky; Čislovniky, porědkove; Zaimena, ktory/kaky; Glagoly, particip na -ty.\par
\noindent\emph{Sources:} Grammar, Adjectives; Numerals, ordinals; Pronouns, ktory/kaky; Verbs, participles in -ty.\par
\par\medskip\Needspace{8\baselineskip}
\hypertarget{DEC-0097}{}
\noindent\textbf{Priměr i napriměr / Priměr and napriměr}\quad{\footnotesize\texttt{DEC-0097}}\par
\smallskip\noindent\emph{Razlog:} Priměr jest imenovnik ‘example’; napriměr jest prislovnik ‘for example’. Zato nominativ, padežne formy i adverbialna fraza ne imajut jednu završnu formu. Zapis zahrnjaje 31 frazu na primjer/na priměr/na primer, jedno naprimer i 53 imenovne formy s primjer-. Slovnik daje priměr i napriměr; nacionalno ije ne dava osnovny oblik etogo zaglavja. Kontrapriměr ostaje poseben složeny termin, a priměnjenje znači application, ne example. Zaměna mora takože čuvati logičnu ulogu priměra: ilustracija ne stava dokazom vseobčego tvrdženja.\par
\smallskip\noindent\emph{Reason:} Priměr is the noun example; napriměr is the adverb for example. Nominatives, inflected nouns and adverbial phrases therefore cannot all share one resulting form. The record covers 31 na primjer/na priměr/na primer phrases, one naprimer and 53 primjer- noun forms. The dictionary supplies priměr and napriměr; national ije is not this headword's basic form. Kontrapriměr remains a distinct compound and priměnjenje means application, not example. The change must also retain the example's logical role: an illustration does not become a proof of a universal statement.\par
\smallskip\noindent\emph{Izvory:} Slovnik, zaglavja 3590, 28131; polna razdělba imennogo i prislovničnogo upotrěbjenja.\par
\noindent\emph{Sources:} Dictionary entries 3590, 28131; the complete noun/adverb partition.\par
\par\medskip\Needspace{8\baselineskip}
\hypertarget{DEC-0098}{}
\noindent\textbf{Priměry i priměry: kako čitati čisla pojavjenij / Priměry and priměry: interpreting occurrence counts}\quad{\footnotesize\texttt{DEC-0098}}\par
\smallskip\noindent\emph{Razlog:} Priměry na počatku zaglavja i priměry vnutri frazy sut isti leksičny oblik s raznym registrom bukvy. V zabilěženom rezultatě jest 133 akcentovane imenovne formy priměr, uključajuči devet Priměry; 49 adverbialnyh napriměr i dva posebe ohranjena kontrapriměr. Eti čisla objasnjajut obseg, ne stvarjajut novu leksičnu odluku ili vahy pri izborě. Poisk, ktory ne uključaje veliku bukvu, može dati drugo čislo bez kakoj-koli lingvističnoj razlike. Zato čisla pojavjenij treba prikladyvati k jasno opredělenoj familiji, ne izdavati ih za čislo raznyh terminov.\par
\smallskip\noindent\emph{Reason:} Heading-initial Priměry and sentence-internal priměry are the same lexical form with different capitalization. The recorded result has 133 accented priměr noun forms, including nine Priměry, 49 adverbial napriměr forms and two separately retained kontrapriměr compounds. These numbers explain scope; they create neither a new lexical decision nor weights for selection. A search excluding capitals can produce a different count without any linguistic difference. Occurrence totals must therefore be attached to a defined family, not presented as numbers of distinct terms.\par
\smallskip\noindent\emph{Izvory:} Zabilěžena polna inventarizacija priměr-familije; leksično rěšenje jest v članku Priměr i napriměr.\par
\noindent\emph{Sources:} The recorded full priměr-family inventory; the lexical decision is in Priměr and napriměr.\par
\par\medskip\Needspace{8\baselineskip}
\hypertarget{DEC-0099}{}
\noindent\textbf{Znakah, krokah, izrazah: lokativ množiny / Znakah, krokah, izrazah: locative plural}\quad{\footnotesize\texttt{DEC-0099}}\par
\smallskip\noindent\emph{Razlog:} Lokativ množiny imenovnikov znak, krok, zaključok, rezultat i izraz završaje se na -ah: znakah, krokah, zaključkah, rezultatah, izrazah. Dvanadset zabilěženih pojavjenij upravjany sut v, pri ili lokativnym po. Polna tablica prvogo sklonjenja dava -ah; ona ne opravdava automatično -yh, ktore prinaleži pridavnikam. Konečne -ih zato treba razpoznati po klasi slova i upravjenju, ne tolko po bukvah. Zaměna ne menja padeže drugih pojavjenij, a po samo po sobě ne vsegda zahtěvaje isti padež vo vsakoj konstrukciji.\par
\smallskip\noindent\emph{Reason:} The locative plural of znak, krok, zaključok, rezultat and izraz ends in -ah: znakah, krokah, zaključkah, rezultatah, izrazah. Twelve recorded occurrences are governed by v, pri or locative po. The complete first-declension table gives -ah; it does not justify the adjectival ending -yh. A final -ih must therefore be classified by word class and government, not letters alone. This rule does not change other cases, and po by itself does not require the same case in every construction.\par
\smallskip\noindent\emph{Izvory:} Jan van Steenbergen, Gramatika, Imenovniki, polna tablica prvogo sklonjenja; pet izvorno-cělevyh par.\par
\noindent\emph{Sources:} Jan van Steenbergen, Grammar, Nouns, complete first-declension table; the five source–target pairs.\par
\par\medskip\Needspace{8\baselineskip}
\hypertarget{DEC-0100}{}
\noindent\textbf{Vměsto: zaměna, ne město / Vměsto: substitution, not place}\quad{\footnotesize\texttt{DEC-0100}}\par
\smallskip\noindent\emph{Razlog:} Vměsto s genitivom označaje ‘instead of’. Trideset tri konteksty zahrnjajut zaměnu formuly, simbola, predpoloženja, konstrukcije ili anaforično vměsto toga. Zato vmesto, naměsto i razděljeno na město davajut v tom smyslu vměsto. Fizično město i imaje město ostajut drugi smysly; zaměna korenja město porušila by stotine nepričastnyh pojavjenij. Slovnik podpira predlog, s nacionalnymi statt-ekvivalentami v sravnitelnyh kolonah. Jedna pravopisna forma ne znači, že vsako ‘place’ slěduje tomu že pravidlu.\par
\smallskip\noindent\emph{Reason:} Vměsto governs the genitive and means instead of. Thirty-three contexts concern replacement of a formula, symbol, assumption or construction, or anaphoric vměsto toga. Vmesto, naměsto and spaced na město therefore yield vměsto in this sense. Physical place and imaje město retain different meanings; replacing the root město would damage hundreds of unrelated occurrences. The dictionary supports the preposition with national instead-of equivalents in its comparison columns. One spelling does not mean that every occurrence of place follows the same rule.\par
\smallskip\noindent\emph{Izvory:} Slovnik, zaglavje 2531; německe statt/anstelle/an Stelle von konteksty.\par
\noindent\emph{Sources:} Dictionary entry 2531; the German statt/anstelle/an Stelle von contexts.\par
\par\medskip\Needspace{8\baselineskip}
\hypertarget{DEC-0101}{}
\noindent\textbf{Vyšših: korenj komparativa / Vyšših: the comparative stem}\quad{\footnotesize\texttt{DEC-0101}}\par
\smallskip\noindent\emph{Razlog:} Vyšših v etom izdani označaje genitiv/lokativ množiny komparativa ‘higher’. Slovnik i gramatika davajut vysši; zapis izbiraje asimilovanu povrhnost vyšš-, uže upotrěbjanu v korpusu. Zato treba razlikovati dva suda: vyših gubi čast komparativnogo korenja; vyšših vměsto vysših jest redakcijny izbor realizacije, ne doslovna forma slovničnogo zaglavja. Deset kontekstov govorit o vyšših formah, derivacijah ili podobnyh objektah, ne o prislovniku vyše. Mekko -ih ostaje opravdano; tvrdokorenny -yh byl by druga ošibka.\par
\smallskip\noindent\emph{Reason:} Vyšših is this edition's genitive/locative plural comparative higher. The dictionary and grammar give vysši; the record chooses the assimilated surface vyšš-, already used in the corpus. Two judgments must be separated: vyših loses part of the comparative stem, whereas vyšših rather than vysših is an editorial realization, not the literal dictionary headword. Ten contexts concern higher forms, derivatives or related objects, not adverbial vyše. Soft-stem -ih remains appropriate; hard-stem -yh would introduce another error.\par
\smallskip\noindent\emph{Izvory:} Slovnik, zaglavje 20302; Gramatika, Pridavniky, vysoky → vysši.\par
\noindent\emph{Sources:} Dictionary entry 20302; Grammar, Adjectives, vysoky → vysši.\par
\par\medskip\Needspace{8\baselineskip}
\hypertarget{DEC-0102}{}
\noindent\textbf{Jih: čije objekty? / Jih: whose objects?}\quad{\footnotesize\texttt{DEC-0102}}\par
\smallskip\noindent\emph{Razlog:} Jih jest nesklanjajemo posvojno zaimeno ‘their’. Gramatika predstavlja jego, jej i jih kako najčestějše tretjelicne formy, a jihny dopuščaje kako pridavnikovu alternativu. Zapis izbiraje jih v 22 njihov- posvojenijah. Svoj ne jest zaměna: ono bi povezalo posvojitelja s subjektom frazy, čego eti konteksty po zapisě ne izražajut. Rod i čislo posědanogo objekta ne menja jih. Nacionalno njihov nije ‘loš jazyk’; rěšenje bira medžuslovjansku konstrukciju. Slovničny tip 1 podpiraje preferenciju, ale ne daje universalno pozitivnu razumljivost.\par
\smallskip\noindent\emph{Reason:} Jih is the indeclinable possessive their. The grammar presents jego, jej and jih as the commonest third-person forms and permits adjectival jihny as an alternative. The record selects jih for twenty-two njihov- possessives. Svoj is not interchangeable: it would identify the possessor with the clause's subject, which the recorded contexts do not express. The possessed object's gender and number do not alter jih. National njihov is not ‘bad language’; this chooses an Interslavic construction. Type 1 supports the preference but does not imply universally positive comprehension.\par
\smallskip\noindent\emph{Izvory:} Slovnik, zaglavja 2601, 6258; Gramatika, Zaimena, posvojne i refleksivno svoj.\par
\noindent\emph{Sources:} Dictionary entries 2601, 6258; Grammar, Pronouns, possessives and reflexive svoj.\par
\par\medskip\Needspace{8\baselineskip}
\hypertarget{DEC-0103}{}
\noindent\textbf{Libovoljny i ljubezno: razne familije / Libovoljny and ljubezno: different families}\quad{\footnotesize\texttt{DEC-0103}}\par
\smallskip\noindent\emph{Razlog:} Dvadset jedno ljub- forma v matematičnyh kvantifikacijah znači arbitrary/any, zato izbrano sklonjeno libovoljny. To ne jest ‘dear’, ‘loving’ ili ‘free’. Otdělny ljubazno znači kindly i ne podleži etomu pravidlu. Razlika jest semantična, ne prosto korenjova: vsako izbrano polje, element ili modul može byti libovoljny bez togo, aby byl svobodny modul. Libovoljny ostaje redakcijny termin s omeženym slovničnym profilom tipa 3; ni veličina korpusa, ni mnogo ranějših pojavjenij ne povyšaje jego do izměrenogo vseobčego standarda.\par
\smallskip\noindent\emph{Reason:} Twenty-one ljub- forms in mathematical quantification mean arbitrary/any and take inflected libovoljny. They do not mean dear, loving or free. Separate ljubazno means kindly and is outside this rule. The distinction is semantic rather than merely root-based: an arbitrarily chosen module need not be a free module. Libovoljny remains an editorial term with a limited type-3 dictionary profile; corpus size and numerous earlier occurrences do not promote it to a measured universal standard.\par
\smallskip\noindent\emph{Izvory:} Slovnik, zaglavje 22300; sravnenje s ljubezno, zaglavje 8908.\par
\noindent\emph{Sources:} Dictionary entry 22300; compare ljubezno, entry 8908.\par
\par\medskip\Needspace{8\baselineskip}
\hypertarget{DEC-0104}{}
\noindent\textbf{Ljubezno: priznanije čužoj pomoci / Ljubezno: acknowledging another scholar's help}\quad{\footnotesize\texttt{DEC-0104}}\par
\smallskip\noindent\emph{Razlog:} Ljubezno prevodi německo freundlichst v priznanju E. Fischeru, ktory dozvolil Noether upotrěbiti neopublikovane teoremy. To jest ljubaznost v naučnom odnosu, ne proizvoljnost matematičnogo objekta. Slovnik izrično dava ljubezno; ljubazno pojavjaje se v hrvatskoj i srbskoj sravnitelnoj koloni. Izbor čuva izraz zahvalnosti i čužu zaslugu. Odstraniti prislovnik kako ‘nepotrěbnu dekoraciju’ by izgubilo avtorsku informaciju. Tuto rěšenje ne tvrdi novo dopuščenje za današnje izdanje: ono věrno prenosi istorijno pripisanje.\par
\smallskip\noindent\emph{Reason:} Ljubezno renders German freundlichst in the acknowledgment of E. Fischer, who permitted Noether to use unpublished theorems. It expresses scholarly courtesy, not arbitrary mathematical selection. The dictionary gives ljubezno; ljubazno appears in its Croatian and Serbian comparison columns. The choice preserves gratitude and another scholar's contribution. Deleting the adverb as unnecessary decoration would lose authorial information. This records historical attribution, not new permission for the present edition.\par
\smallskip\noindent\emph{Izvory:} Slovnik, zaglavje 8908; německa fraza deren Verwendung mir dieser freundlichst gestattete.\par
\noindent\emph{Sources:} Dictionary entry 8908; German deren Verwendung mir dieser freundlichst gestattete.\par
\par\medskip\Needspace{8\baselineskip}
\hypertarget{DEC-0105}{}
\noindent\textbf{Byti i biti / Byti and biti}\quad{\footnotesize\texttt{DEC-0105}}\par
\smallskip\noindent\emph{Razlog:} Byti jest kopula ‘be’; biti jest drugy glagol ‘hit/beat’. Četyrdeset tri infinitivy biti v zapisě imajut kopulativnu, modalnu ili pasivnu ulogu, zato davajut byti. To ne isključaje legitimno biti v drugom kontekstě. Forma byti ne menja mora, može ili drugy modalny izraz, ni pasivny particip s kotorym se povezaje. Nacionalna kopula biti ne jest dokaz, že oba medžuslovjanske zaglavja sut sinonimy. Finite jest/sut/bude i prošle formy podležat vlastnoj paradigmi; pravilo o infinitivě ne raspisuje ih nanovo.\par
\smallskip\noindent\emph{Reason:} Byti is the copula be; biti is the separate verb hit/beat. The forty-three recorded biti infinitives have copular, modal or passive functions and therefore become byti. This does not exclude legitimate biti elsewhere. Byti preserves the surrounding modal expression and passive participle. National-language copular biti is not evidence that the two Interslavic headwords are synonyms. Finite jest/sut/bude and past forms follow their own paradigm; an infinitive rule does not rewrite them.\par
\smallskip\noindent\emph{Izvory:} Slovnik, zaglavja 593, 4049; Gramatika, Glagoly, byti i biti → bij-/bity.\par
\noindent\emph{Sources:} Dictionary entries 593, 4049; Grammar, Verbs, byti and biti → bij-/bity.\par
\par\medskip\Needspace{8\baselineskip}
\hypertarget{DEC-0106}{}
\noindent\textbf{Tu, ta i takože / Tu, ta and takože}\quad{\footnotesize\texttt{DEC-0106}}\par
\smallskip\noindent\emph{Razlog:} Tu jest ‘here’; ta jest ženska forma demonstrativa toj. Osem tuta/Tuta pojavjenij razděljaje se na četyri prislovnika i četyri pridavnikove ukazanja: ta razlika, ta sistema, ta invariantna definicija i ta razširjena sistema. Německy rod ne upravjaje cělevym soglasjem: System jest německe neutrum, ale sistema jest ženska. Takozhe odvolavaje se na takože ‘also’, ne na jednu iz dvuh predhodnyh kategorij. Slovnično tut takože obstaja; to ne opravdava nerazličenu tuta. Jedna povrhnostna forma može skryvati dvě sintaktične ulogi.\par
\smallskip\noindent\emph{Reason:} Tu means here; ta is the feminine demonstrative form of toj. Eight tuta/Tuta occurrences split into four adverbs and four attributive demonstratives: ta razlika, ta sistema, ta invariantna definicija and ta razširjena sistema. German gender does not control target agreement: System is German neuter, but sistema is feminine. Takozhe represents takože, also, not either preceding category. Dictionary tut also exists, but does not justify undifferentiated tuta. One surface form can conceal two syntactic functions.\par
\smallskip\noindent\emph{Izvory:} Slovnik, zaglavja 1758, 2893, 1716; Gramatika, Zaimena, toj/ta/to.\par
\noindent\emph{Sources:} Dictionary entries 1758, 2893, 1716; Grammar, Pronouns, toj/ta/to.\par
\par\medskip\Needspace{8\baselineskip}
\hypertarget{DEC-0107}{}
\noindent\textbf{Harakterne svojstva i stepenne redy / Harakterne svojstva and stepenne redy}\quad{\footnotesize\texttt{DEC-0107}}\par
\smallskip\noindent\emph{Razlog:} Harakterne svojstva i v stepenne redy potrěbujut -e pri neživotnoj množině. Svojstva sut neutrum množiny, redy mužsky neživotny akuzativ množiny. Zato harakternye i stepennye ne sut formy etoj paradigmy. Izbor popravjaje koncovku, ne odlučuje mežu harakterny i harakterističny ili medžu raznymi slovami dlja power series. To mora byti otdělna leksična odluka. Genitiv/akuzativ oduševljenogo -yh byl by nepravilny, poneže avtor ne govori o osobah.\par
\smallskip\noindent\emph{Reason:} Harakterne svojstva and v stepenne redy take -e with non-animate plurals. Svojstva is neuter plural; redy is masculine-inanimate accusative plural. Harakternye and stepennye are not forms of this paradigm. The decision repairs the ending without deciding between harakterny and harakterističny or among terms for power series; those are separate lexical choices. Animate genitive/accusative -yh would be inappropriate because the author is not referring to persons.\par
\smallskip\noindent\emph{Izvory:} Gramatika, Pridavniky, nominativ/akuzativ neživotnoj množiny; německe charakteristischen Eigenschaften i Potenzreihenentwicklung.\par
\noindent\emph{Sources:} Grammar, Adjectives, non-animate nominative/accusative plural; German charakteristischen Eigenschaften and Potenzreihenentwicklung.\par
\par\medskip\Needspace{8\baselineskip}
\hypertarget{DEC-0108}{}
\noindent\textbf{Najkratši: najmenša dolžina / Najkratši: shortest length}\quad{\footnotesize\texttt{DEC-0108}}\par
\smallskip\noindent\emph{Razlog:} Najkratši slěduje kratky → kratši i prefiksu naj-. Trinadset najkrajš- form v zapisě znači shortest pri prědstavjenju ili dolžině. Krajny ‘extreme/utmost’ ne jest tot že korenj, i minimalny može v matematike značiti minimalen po častičnom porědku bez najmenej dolžiny. Zato treba čuvati točny kriterij německogo kürzest-. Nacionalno najkrajši može legitimno značiti shortest, ale izdanje izbiraje medžuslovjansku paradigmu krat-. Padež, rod i čislo vse ravno zavisjut od frazy.\par
\smallskip\noindent\emph{Reason:} Najkratši follows kratky → kratši with naj-. Thirteen recorded najkrajš- forms mean shortest in relation to representation or length. Krajny, extreme/utmost, has a different root, and mathematical minimal can mean minimal under a partial order without least length. The exact criterion of German kürzest- must therefore be retained. National najkrajši can legitimately mean shortest, but the edition adopts the Interslavic krat- paradigm. Case, gender and number still depend on the sentence.\par
\smallskip\noindent\emph{Izvory:} Slovnik, zaglavja 2641, 757; Gramatika, Pridavniky, kratky → kratši i naj-.\par
\noindent\emph{Sources:} Dictionary entries 2641, 757; Grammar, Adjectives, kratky → kratši and naj-.\par
\par\medskip\Needspace{8\baselineskip}
\hypertarget{DEC-0109}{}
\noindent\textbf{Daljši, ranši, pozdnějši, najobčejši / Daljši, ranši, pozdnějši, najobčejši}\quad{\footnotesize\texttt{DEC-0109}}\par
\smallskip\noindent\emph{Razlog:} Daljši označaje further/subsequent v izslědovanju; ranši i pozdnějši sut pridavniky earlier/later. Raněje ostaje prislovnik. Šestnadset slučajev zahrnjaje jedno dalših, dvě pozdnejš- formy, šest ranějš-/ranjejš- pridavnikov i sedm skraćenyh najobć- superlativov. Poslednje davajut najobćejš- v proširenom ili najobčejš- v osnovnom pravopisu. Německe allgemeinste zahtěvaje most general, ne prosto common. Gramatika dopuščaje daleky → dalši, takže dalš- ne jest univerzalno nevěrny; lokalny izbor daljši jest smyslovy i slovničny, ne osuda vsej paradigmy.\par
\smallskip\noindent\emph{Reason:} Daljši expresses further/subsequent in investigation; ranši and pozdnějši are adjectival earlier/later, while raněje remains an adverb. Sixteen cases comprise one dalših, two pozdnejš- forms, six ranějš-/ranjejš- adjectives and seven truncated najobć- superlatives. The last yield extended-spelling najobćejš- or basic najobčejš-. German allgemeinste requires most general, not merely common. The grammar permits daleky → dalši, so dalš- is not universally invalid; the contextual preference for daljši is semantic and lexical, not a condemnation of the entire paradigm.\par
\smallskip\noindent\emph{Izvory:} Slovnik, zaglavja 3032, 4060, 11106, 1000; Gramatika, Pridavniky, komparativ i superlativ.\par
\noindent\emph{Sources:} Dictionary entries 3032, 4060, 11106, 1000; Grammar, Adjectives, comparative and superlative.\par
\par\medskip\Needspace{8\baselineskip}
\hypertarget{DEC-0110}{}
\noindent\textbf{Tretji, treťji, trětji / Tretji, treťji, trětji}\quad{\footnotesize\texttt{DEC-0110}}\par
\smallskip\noindent\emph{Razlog:} Tretji v osnovnom pravopisu i slovnično treťji v proširenom označajut third. Gramatična stranica piše trětji: eto jest realno razhodenje izvorov, ne razlika matematičnogo smysla. Istorična odluka predpočitala slovnično treťji i mekko sklonjenje; tretjega dalo treťjego, ne prosto dodato dijakritično znamenje. Petnadset pojavjenij zahrnjaje porędok, stepenj, integral i koreni jedinice. Kardinal tri i frakcija tretjina sut druge kategorije. Kasnija osnovna grafija ne potrebuje ť, zato nemarkovano tretj- ne može v novom izdani byti proglašeno samo po sobě ošibkoj.\par
\smallskip\noindent\emph{Reason:} Basic-spelling tretji and dictionary extended treťji mean third. The grammar page spells trětji: a real source disagreement, not a mathematical distinction. The historical choice preferred dictionary treťji with soft declension; tretjega became treťjego, not merely an added diacritic. Fifteen occurrences cover order, degree, an integral and roots of unity. Cardinal tri and fractional tretjina are separate categories. Later basic orthography does not require ť, so unmarked tretj- cannot be called intrinsically erroneous in that edition.\par
\smallskip\noindent\emph{Izvory:} Slovnik, zaglavje 439; Gramatika, Čislovniky, porědkove, s razhodenjem trětji; Pravopis, opcionalne bukvy.\par
\noindent\emph{Sources:} Dictionary entry 439; Grammar, Numerals, ordinals, including conflicting trětji; Orthography, optional letters.\par
\par\medskip\Needspace{8\baselineskip}
\hypertarget{DEC-0111}{}
\noindent\textbf{Iduči i poprědny / Iduči and poprědny}\quad{\footnotesize\texttt{DEC-0111}}\par
\smallskip\noindent\emph{Razlog:} Iduči jest reguljarny aktivny particip idti; ne treba odstranjati jego prosto zato, že opisuje abstraktny objekt. Sedm zabilěženih upotrěbjenij znači ‘iduči prěz punkt’, ‘iduči od p do q’ ili podobno razširenje argumenta. Jedno prědidučih znači predhodne čisla v fiksovanom dobrom porědku, zato izbrano poprědnyh. Prědhodny jest takože slovnično zaglavje; razlika tipa motivuje redakcijnu preferenciju, ne zabranu. Ranših by moglo suggerovati vrěme vměsto porědka bazovyh členov. Participska metafora i pridavnik ‘previous’ potrěbujut raznu analizu.\par
\smallskip\noindent\emph{Reason:} Iduči is a regular active participle of idti and need not be removed merely because it describes an abstract object. Seven recorded uses mean going through a point, running from p to q or comparable extension of an argument. One prědidučih instead means preceding numbers in a fixed well-order, giving poprědnyh. Prědhodny is also a dictionary headword; its different type motivates editorial preference, not prohibition. Ranših could suggest time rather than the order of basis members. Participial metaphor and adjectival previous require different analyses.\par
\smallskip\noindent\emph{Izvory:} Slovnik, zaglavja 2104, 20046, 14924; Gramatika, Glagoly, aktivny particip.\par
\noindent\emph{Sources:} Dictionary entries 2104, 20046, 14924; Grammar, Verbs, active participles.\par
\par\medskip\Needspace{8\baselineskip}
\hypertarget{DEC-0112}{}
\noindent\textbf{Četyrěh i šesti / Četyrěh and šesti}\quad{\footnotesize\texttt{DEC-0112}}\par
\smallskip\noindent\emph{Razlog:} Četyrěh jest genitiv/lokativ četyri, kako pokazuje polna tablica 2–4. Šesti slěduje sklonjenju čislovnikov od pet nagore po vzoru kost. Zato tri četyrih i dvě šestih v genitivnyh frazah ne trebujut jednakej koncovky. Oni označajut točno četyri ili šest objektov, ne porědkove četvrty/šesty. Gramatika dopuščaje nesklanjane čislovniky kako prostějšu možnost, ale izdanje izbiraje izrično sklonjenje. To jest obrazovany registr, ne tvrdženje, že alternativna olahčena konstrukcija jest nerazumljiva.\par
\smallskip\noindent\emph{Reason:} Četyrěh is genitive/locative četyri in the complete 2–4 table. Šesti follows the kost pattern used for numbers from five upwards. Three četyrih and two šestih in genitive phrases therefore require different endings. They denote exactly four or six objects, not the ordinals fourth/sixth. The grammar permits undeclined numerals as an easier option, but the edition chooses explicit inflection. That is a register choice, not a claim that the simplified construction is unintelligible.\par
\smallskip\noindent\emph{Izvory:} Slovnik, zaglavja 2856, 1739; Jan van Steenbergen, Gramatika, Čislovniky, tablice 2–4 i 5–99.\par
\noindent\emph{Sources:} Dictionary entries 2856, 1739; Jan van Steenbergen, Grammar, Numerals, tables for 2–4 and 5–99.\par
\par\medskip\Needspace{8\baselineskip}
\hypertarget{DEC-0113}{}
\noindent\textbf{Idučih: otnošenje k Gaussu / Idučih: a relation to Gauss}\quad{\footnotesize\texttt{DEC-0113}}\par
\smallskip\noindent\emph{Razlog:} Idučih v ‘pytanjah idučih nazad k Gaussu’ izražaje německo auf Gauss zurückgehend. To jest participialno opisanje izhoda pytanja, ne tolko poprědny po vrěmeni. Slovnično idti i gramatično iduči podpirajut č; mekky genitiv množiny dava -ih. Istorična zaměna idućih → idučih zato ima konkretnu paradigmu. Ona ne znači, že vsako ć v drugih slovah bylo ošibkoj: opcionalna etimologična grafija i drugi leksične familije morajut byti razlikovane. Osnovny pravopis čini č čitateljno doslědnym.\par
\smallskip\noindent\emph{Reason:} Idučih in questions going back to Gauss expresses German auf Gauss zurückgehend. This participial description concerns the questions' origin, not merely temporal priority. Dictionary idti and grammatical iduči support č; the soft genitive plural supplies -ih. The historical idućih → idučih change thus has a specific paradigm. It does not prove every ć elsewhere erroneous: optional etymological spelling and different lexical families must be distinguished. Basic orthography makes č consistent for readers.\par
\smallskip\noindent\emph{Izvory:} Gramatika, Glagoly, iduči; německy izraz auf Gauss zurückgehenden Fragestellungen.\par
\noindent\emph{Sources:} Grammar, Verbs, iduči; German auf Gauss zurückgehenden Fragestellungen.\par
\par\medskip\Needspace{8\baselineskip}
\hypertarget{DEC-0114}{}
\noindent\textbf{Poprědny i raněje / Poprědny and raněje}\quad{\footnotesize\texttt{DEC-0114}}\par
\smallskip\noindent\emph{Razlog:} Poprědny jest pridavnik ‘previous’; raněje jest prislovnik ‘previously/earlier’. Zapis razděljaje 49 predhodn-/prědhodn- pojavjenij na 47 pridavnikov i dvě prislovnične konstrukcije: raněje prijetym i raněje spomenuty. Poslednje modifikujut přijemanje ili spominanje, ne sut neutra pridavnika. Poprědnji i poprědny sut oba slovnične formy, a prědhodny ostaje dopustimy sinonim. Izbor tvrdogo poprědny v izdani podpira jednu paradigmu; prědnje ‘in advance’ by v dvuh prislovnikah moglo pridati planovanu anticipaciju, kotroj izvor ne zahtěvaje.\par
\smallskip\noindent\emph{Reason:} Poprědny is adjectival previous; raněje is adverbial previously/earlier. The record divides 49 predhodn-/prědhodn- occurrences into 47 adjectives and two adverbial constructions: raněje prijetym and raněje spomenuty. The latter modify accepting or mentioning, not neuter adjective slots. Both poprědnji and poprědny are dictionary forms, while prědhodny remains an available synonym. Selecting hard-stem poprědny supports one edition paradigm; prědnje, in advance, could add planned anticipation that the two adverbial source contexts do not require.\par
\smallskip\noindent\emph{Izvory:} Slovnik, zaglavja 20046, 14924, 4061; Gramatika, Pridavniky.\par
\noindent\emph{Sources:} Dictionary entries 20046, 14924, 4061; Grammar, Adjectives.\par
\par\medskip\Needspace{8\baselineskip}
\hypertarget{DEC-0115}{}
\noindent\textbf{Prědhoditi: objekt prědhodi drugomu / Prědhoditi: one object precedes another}\quad{\footnotesize\texttt{DEC-0115}}\par
\smallskip\noindent\emph{Razlog:} Prědhoditi jest glagol ‘precede’, ne pridavnik ‘previous’. Šest glagolnyh prethod-/predhod- form davajut prědhoditi-familiju, šest pridavnikov poprědny-familiju; sedm uže reguljarnyh glagolnyh form ostajut. Finite prědhodi, prošlo prědhodila i aktivno prědhodeči imajut razne gramatične ulogi, ale čuvajut otnošenje predhodnosti. Proměniti ih vse na pridavnik by izgubilo predikat ili participialnu konstrukciju. Obči korenj v dvuh jazykah zato ne jest dovoljno opreděljenje jedinice prěklada.\par
\smallskip\noindent\emph{Reason:} Prědhoditi is the verb precede, not the adjective previous. Six prethod-/predhod- verb forms yield the prědhoditi family, six adjectives the poprědny family, while seven already regular verb forms remain. Finite prědhodi, past prědhodila and active prědhodeči serve different grammatical functions but preserve precedence. Turning them all into adjectives would lose a predicate or participial construction. A shared root across languages is therefore insufficient to define the translation unit.\par
\smallskip\noindent\emph{Izvory:} Slovnik, zaglavja 501, 20046; Gramatika, Glagoly, prezent, prošly čas i aktivny particip.\par
\noindent\emph{Sources:} Dictionary entries 501, 20046; Grammar, Verbs, present, past and active participles.\par
\par\medskip\Needspace{8\baselineskip}
\hypertarget{DEC-0116}{}
\noindent\textbf{Prědpokladati i prědpoloženje / Prědpokladati and prědpoloženje}\quad{\footnotesize\texttt{DEC-0116}}\par
\smallskip\noindent\emph{Razlog:} Formalno predpoloženje izražajut nesovršeno prědpokladati, sovršeno prědpoložiti i imenovnik prědpoloženje. Zapis razlikuje 27 nesovršenyh i 23 sovršene glagoly, 56 imenovnikov na -ka, 67 na -jenje i 31 pasivnu ili pridavnikovu formu. Prědpostavjati jest v slovniku ‘prefer’, zato podobna forma ne garantuje potrebny smysl. Pri zaměně ženskogo predpostavka na neutrum prědpoloženje treba takože iz toj → iz togo i byla upotrěbljena dodatna → bylo upotrěbljeno dodatno. Uslovje i domněvati ne sut avtomatične ekvivalenty formalnoj premisy. Prědpoklad i prědpolagati sut legitimne alternativy; preferencija drugoj familije jest redakcijna, ne izmišljena zabrana. Negacija ne- mora ostati.\par
\smallskip\noindent\emph{Reason:} Formal assumption is expressed by imperfective prědpokladati, perfective prědpoložiti and noun prědpoloženje. The record distinguishes 27 imperfective verbs, 23 perfective verbs, 56 -ka nouns, 67 -jenje nouns and 31 passive/adjectival forms. Dictionary prědpostavjati means prefer, so resemblance does not establish the required sense. Changing feminine predpostavka to neuter prědpoloženje also requires iz toj → iz togo and byla upotrěbljena dodatna → bylo upotrěbljeno dodatno. Uslovje and domněvati are not automatic equivalents of a formal premise. Prědpoklad and prědpolagati are legitimate alternatives; preferring another family is editorial, not an invented prohibition. Negation ne- must remain.\par
\smallskip\noindent\emph{Izvory:} Slovnik, zaglavja 5324, 3130, 11247, 18400, 20062, 16759; sedm polnyh frazovyh sravnenij v dannyh.\par
\noindent\emph{Sources:} Dictionary entries 5324, 3130, 11247, 18400, 20062, 16759; seven complete phrase comparisons in the data.\par
\par\medskip\Needspace{8\baselineskip}
\hypertarget{DEC-0117}{}
\noindent\textbf{Ješče i ide / Ješče and ide}\quad{\footnotesize\texttt{DEC-0117}}\par
\smallskip\noindent\emph{Razlog:} Ješče znači still/even/yet, a ide jest tretje lice jedinstvenogo čisla idti. Dvadeset sedm ještě i 29 jde/idé form v zapisě davajut eti slovnične obliky. Češske jde jest pravilno v češskom; ono ne jest izbrana medžuslovjanska forma. Ide o v matematičnoj prozě znači ‘reč jest o’, ne obyčno fizično idenje. Ješče može označati trajanje, dodatok ili usiljenje; konkretny obseg ne směje byti izravnen do jednogo ‘still’. Izbor slovničnogo oblika čuva tu kontekstnu polisemiju.\par
\smallskip\noindent\emph{Reason:} Ješče means still/even/yet, while ide is third-person singular idti. Twenty-seven ještě occurrences and 29 jde/idé forms yield those dictionary forms. Czech jde is correct Czech, but not the selected Interslavic form. Mathematical ide o means concerns/is about rather than ordinarily denoting physical motion. Ješče may express continuation, addition or emphasis; its contextual scope must not be flattened into a single still. Selecting the dictionary form preserves that contextual polysemy.\par
\smallskip\noindent\emph{Izvory:} Slovnik, zaglavja 1076, 2104; Gramatika, Glagoly, idti.\par
\noindent\emph{Sources:} Dictionary entries 1076, 2104; Grammar, Verbs, idti.\par
\par\medskip\Needspace{8\baselineskip}
\hypertarget{DEC-0118}{}
\noindent\textbf{Prěhodi: pravopis i finitna forma / Prěhodi: spelling and finite form}\quad{\footnotesize\texttt{DEC-0118}}\par
\smallskip\noindent\emph{Razlog:} Prěhoditi dava tretje lice jedinstvenogo čisla prěhodi, ne prěhodit. Petnadset takyh finitnyh form potrěbuje morfologičnu popravku. Trideset jedno plain-e prehod- pojavjenje znači drugu rabotu: doslědny izbor prěhod- v togdašnjem pravopisu, ne dokaz, že pojednostavljeno e jest samo po sobě zabranjeno. Prěhodim jest dopustima alternativa prvogo lica i ostaje; množinno prěhodet takože ne mora byti sg prěhodi. Particip prěhodeči slěduje vtoromu sprezanju. Medžu 286 članami pregledanoj familije 240 služilo kak neproměnjene sravnenja, ne kak dodatne glasove za izbrany korenj.\par
\smallskip\noindent\emph{Reason:} Prěhoditi has third-person singular prěhodi, not prěhodit; fifteen finite forms require that morphological correction. Thirty-one plain-e prehod- occurrences concern something different: consistency with the then-selected prěhod- spelling, not proof that simplified e is intrinsically forbidden. Prěhodim is a permitted first-person alternative and remains; plural prěhodet must not become singular prěhodi. Participial prěhodeči follows the second conjugation. Of 286 examined family members, 240 served as unchanged comparators, not additional votes for a root.\par
\smallskip\noindent\emph{Izvory:} Slovnik, zaglavja 19506, 2425; Gramatika, Glagoly, vtoro sprezanje i alternativno prvo lice.\par
\noindent\emph{Sources:} Dictionary entries 19506, 2425; Grammar, Verbs, second conjugation and alternative first person.\par
\par\medskip\Needspace{8\baselineskip}
\hypertarget{DEC-0119}{}
\noindent\textbf{Red, porędok i rang / Red, porędok and rang}\quad{\footnotesize\texttt{DEC-0119}}\par
\smallskip\noindent\emph{Razlog:} Red, etimologično ręd, pokryvaje row/sequence/order; to ne daje avtomatično vse izvedene matematične pridavnike. Slovnično rędovy znači vojenny private, a rędny regular. Zato 35 zabilěženih pridavnikovyh ili rednost form potrěbuje frazu po konkretnom smyslu: serija, rang ili porędok. Drugy 356 pojavjenij sjedinjaje pravopis row/order-familije, a dvě rędku davajut rędu. Pět matematičnyh identifikatorov red ostajut symbolami, ne prěkladanym tekstom. Obyčny slovničny smysl ne dokazuje, že tehnično razširenje jest nemožno; izbrana izrična fraza jest opatrna politika, ktora ne polagaje nejasny novy smysl bez vysvětlenja.\par
\smallskip\noindent\emph{Reason:} Red, etymological ręd, covers row/sequence/order but does not automatically establish every mathematical adjective derived from it. Dictionary rędovy means military private and rędny regular. Thirty-five recorded adjectival or rednost forms therefore require phrases selected for the particular sense: series, rank or order. Another 356 occurrences standardize the row/order family's spelling, and two rędku forms become rędu. Five mathematical red identifiers remain symbols, not translated prose. An ordinary dictionary sense does not prove technical extension impossible; the explicit phrase is a cautious policy avoiding an unexplained new meaning.\par
\par\smallskip\noindent\emph{Sravnjenje formov / Comparison of forms}
\begin{longtable}{@{}P{0.43\textwidth}P{0.43\textwidth}P{0.08\textwidth}@{}}
\toprule Sravniteljnaja forma / Compared form & Izbrana forma ili funkcija / Selected form or function & Pojav. / Uses\\\midrule\endfirsthead
\toprule Sravniteljnaja forma / Compared form & Izbrana forma ili funkcija / Selected form or function & Pojav. / Uses\\\midrule\endhead
\bottomrule\endfoot
rędovo razloženje & razloženje v rędy & 3\\\addlinespace
rędovogo razloženja & razloženja v rędy & 3\\\addlinespace
rędovyh razloženij & razloženij v rędy & 1\\\addlinespace
redovymi razloženjami & razloženjami v rędy & 1\\\addlinespace
determinanty $\rho$-rędnoj matrice & determinanty matrice ręda $\rho$ & 1\\\addlinespace
Vsakomu $\rho$-rędnomu matričnomu produktu & Vsakomu matričnomu produktu ręda $\rho$ & 1\\\addlinespace
nulovy $(n-\lambda)$-rędny determinant & nulovy determinant ręda $(n-\lambda)$ & 1\\\addlinespace
Teorema produkta, utvorena jeden raz za $\rho$-rędne, a drugi raz za $(\rho-1)$-rędne matrice & Teorema produkta, utvorena jeden raz za matrice ręda $\rho$, a drugi raz za matrice ręda $(\rho-1)$ & 1\\\addlinespace
$n$-rędnu determinantu & determinantu ręda $n$ & 1\\\addlinespace
dane sut jedno-, dvo-, do $(n-1)$-rędnyh determinantov & dane sut prěz determinanty ręda 1, 2 až do $(n-1)$ & 1\\\addlinespace
s $\bar k_i$-rędnymi iz $\overline{\mathfrak A}_i$ & s determinantami ręda $\bar k_i$ iz $\overline{\mathfrak A}_i$ & 1\\\addlinespace
v někoj rednosti & v někom porędku & 1\\\addlinespace
s odgovarjajučimi $t$-rednymi matricami $A$ & s odgovarjajučimi matricami $A$ ręda $t$ & 1\\\addlinespace
$t$-redne determinanty & determinanty ręda $t$ & 1\\\addlinespace
$\rho$-rędnyh determinantov & determinantov ręda $\rho$ & 3\\\addlinespace
$k_1$-rędnyh determinantov & determinantov ręda $k_1$ & 1\\\addlinespace
$k_2$-rędnyh determinantov & determinantov ręda $k_2$ & 1\\\addlinespace
$\bar k_i$-rędnyh determinantov & determinantov ręda $\bar k_i$ & 1\\\addlinespace
$k_i$-rędnyh determinantov & determinantov ręda $k_i$ & 2\\\addlinespace
$\varrho$-rednyh determinantov & determinantov ręda $\varrho$ & 1\\\addlinespace
$\rho$-rędna determinanta & determinanta ręda $\rho$ & 1\\\addlinespace
$n$-rędnyh matric & matric ręda $n$ & 1\\\addlinespace
$n$-rędnyh matric & matric ręda $n$ & 3\\\addlinespace
$r$-rednyh matric & matric ręda $r$ & 2\\\addlinespace
\end{longtable}
24 različne frazne sravnjenja imajut 34 priměnjenja. Jedna fraza sodržuje dva prilagateljnika; zato 35 proizvodnyh formov ne znači 35 različnyh fraz. Dvě formy rędku sut odvojene od tųh fraz.\par\emph{EN:} The 24 distinct phrase comparisons have 34 applications. One phrase contains two adjectives, so 35 derived forms do not mean 35 distinct phrases. The two rędku forms are separate from those phrases.\par
\smallskip\noindent\emph{Izvory:} Slovnik, zaglavja 2283, 3798, 3145, 17027, 36025; německe Reihe/Reihenfolge/Reihenentwicklung/reihig konteksty.\par
\noindent\emph{Sources:} Dictionary entries 2283, 3798, 3145, 17027, 36025; German Reihe/Reihenfolge/Reihenentwicklung/reihig contexts.\par
\par\medskip\Needspace{8\baselineskip}
\hypertarget{DEC-0120}{}
\noindent\textbf{Razširiti, razširjeny, razširjenje / Razširiti, razširjeny, razširjenje}\quad{\footnotesize\texttt{DEC-0120}}\par
\smallskip\noindent\emph{Razlog:} Razširiti dava pasivno razširjeny i imenovnik razširjenje. Sto devetnadset razširen- form v zapisě slěduje toj derivaciji; 222 uže s j ostajut sravnenjem, ne nezavisnym jezyčnym glasom. Posebna razširenej dava razširjenoj, poneže particip završaje se na n i sklanjaje se tvrdo; prědhodno j ne čini njegov korenj mekym pridavnikom. Arhivny opis ‘soft j-stem’ zato ne objasnjaje tu koncovku pravilno. Korpusna podpora deveti kandidatnyh slověnskih fajlov s 18 -jen- tokenami ne jest devet nezavisnyh ekspertov: četiri repozitorijske familije i nepotvrdženo avtorstvo ostajut ograničenja. AI-generovany material ne směje dokazovati vlastny izbor.\par
\smallskip\noindent\emph{Reason:} Razširiti yields passive razširjeny and noun razširjenje. The 119 recorded razširen- forms follow that derivation; 222 existing j-forms are comparators, not independent linguistic votes. Exceptional razširenej becomes razširjenoj because the participle ends in n and declines as a hard adjective; the earlier j does not make its adjectival stem soft. The archived description ‘soft j-stem’ therefore does not correctly explain this ending. Nine candidate Slovene files with eighteen -jen- tokens are not nine independent experts: four repository families and unconfirmed authorship limit the evidence. AI-generated material cannot validate its own choice.\par
\smallskip\noindent\emph{Izvory:} Slovnik, zaglavja 1162, 3601, 11873; Gramatika, Glagoly, -iti → -jeny → -jenje; Pridavniky, tvrdo n.\par
\noindent\emph{Sources:} Dictionary entries 1162, 3601, 11873; Grammar, Verbs, -iti → -jeny → -jenje; Adjectives, hard n.\par
\par\medskip\Needspace{8\baselineskip}
\hypertarget{DEC-0121}{}
\noindent\textbf{Razširjenje: razširenje objekta / Razširjenje: extension of an object}\quad{\footnotesize\texttt{DEC-0121}}\par
\smallskip\noindent\emph{Razlog:} Tri forme prošir- označajut matematično razširenje, ne prostorno širenje bez dodatnogo smisla. Slovnik dava razširiti, razširjeny i razširjenje; hrvatsko proširenje stojí v sravniteljnoj koloni, ne kako izbrano medžuslovjansko zaglavje. Zato dva nominativa/akuzativa prijmajut razširjenje, a instrumental razširjenjem. Izbor korenja ne směje prominiti objekt razširenja, jego bazu ili otnošenje vključenja.\par
\smallskip\noindent\emph{Reason:} The three prošir- forms denote mathematical extension. The dictionary supplies razširiti, razširjeny and razširjenje; Croatian proširenje appears in a comparison column, not as the selected Interslavic headword. Two nominative/accusative forms therefore take razširjenje and the instrumental takes razširjenjem. The lexical choice must not change the object extended, its base or the inclusion relation.\par
\smallskip\noindent\emph{Izvory:} Slovnik, zaglavja 1162, 3601, 11873\par
\noindent\emph{Sources:} Dictionary entries 1162, 3601, 11873\par
\par\medskip\Needspace{8\baselineskip}
\hypertarget{DEC-0122}{}
\noindent\textbf{Tako že i takože: slědstvo i dodavanje / Tako že and takože: consequence and addition}\quad{\footnotesize\texttt{DEC-0122}}\par
\smallskip\noindent\emph{Razlog:} Razděljenje slov jest smyslovno: tako že vvodí slědstvo ili cilj, takože dodava ‘takže, i to’. Četrnadset tak že i devetnadset takže pojavjenij potrěbujut prvu konstrukciju; jedno takže i pět takođe potrěbujut drugu. Německe nicht nur … sondern sogar podpiraje dodatny smysl togo iznimnogo takže. Jednaka tipografična zaměna vsej familije by izbrisala razliku medžu logičnym slědstvom i dodatnym tvrđenjem.\par
\smallskip\noindent\emph{Reason:} Word division carries meaning: tako že introduces a consequence or purpose, whereas takože adds ‘also’. Fourteen tak že and nineteen takže occurrences require the first construction; one takže and five takođe occurrences require the second. German nicht nur … sondern sogar supports the additive interpretation of the exceptional takže. A uniform spelling replacement would erase the distinction between logical consequence and an additional assertion.\par
\smallskip\noindent\emph{Izvory:} Slovnik, zaglavja 12906, 1716\par
\noindent\emph{Sources:} Dictionary entries 12906, 1716\par
\par\medskip\Needspace{8\baselineskip}
\hypertarget{DEC-0123}{}
\noindent\textbf{Obstajati: bytije matematičnogo objekta / Obstajati: mathematical existence}\quad{\footnotesize\texttt{DEC-0123}}\par
\smallskip\noindent\emph{Razlog:} Obstajati v značenju ‘egzistovati’ jest posebno slovnično zaglavje, različno od homonima ‘ostavati, provoditi vrěme’. Četyrdeset četiri postoj- forme prěvedene sut v obstaj-familiju: obstaja, obstajajut, obstajati, obstajanje, obstajanjem i obstajajuči. To jest doslědnost redakcijnogo termina, ne dokaz jeho najvyššej razumljivosti vo vsih slovjanskyh jezykah. Slovničny status i omeženy sravniteljny profil treba čitati oddělno. Kvantifikator i zavisimost objekta od danih parametrov ostajut čestju tvrđenja.\par
\smallskip\noindent\emph{Reason:} Obstajati meaning ‘exist’ is a separate dictionary entry from its homonym meaning ‘stay, spend time’. Forty-four postoj- forms were rendered with the obstaj- family: obstaja, obstajajut, obstajati, obstajanje, obstajanjem and obstajajuči. This is consistency within an editorial terminology, not evidence that the family is maximally intelligible across all Slavic languages. Dictionary status and its limited comparison profile are different questions. The quantifier and the object's dependence on given parameters remain part of the assertion.\par
\par\smallskip\noindent\emph{Sravnjenje formov / Comparison of forms}
\begin{longtable}{@{}P{0.43\textwidth}P{0.43\textwidth}P{0.08\textwidth}@{}}
\toprule Sravniteljnaja forma / Compared form & Izbrana forma ili funkcija / Selected form or function & Red / Line\\\midrule\endfirsthead
\toprule Sravniteljnaja forma / Compared form & Izbrana forma ili funkcija / Selected form or function & Red / Line\\\midrule\endhead
\bottomrule\endfoot
Postojí & Obstaja & \\\addlinespace
postojí & obstaja & \\\addlinespace
postoji & obstaja & \\\addlinespace
postojut & obstajajut & \\\addlinespace
postojat & obstajajut & \\\addlinespace
postojati & obstajati & \\\addlinespace
postojanje & obstajanje & \\\addlinespace
postojanjem & obstajanjem & \\\addlinespace
postojači & obstajajuči & \\\addlinespace
postoječe & obstajajuče & \\\addlinespace
\end{longtable}
\smallskip\noindent\emph{Izvory:} Slovnik, zaglavja 9895, 9897\par
\noindent\emph{Sources:} Dictionary entries 9895, 9897\par
\par\medskip\Needspace{8\baselineskip}
\hypertarget{DEC-0124}{}
\noindent\textbf{Što, ktory, to i tym: referencija v prědloženju / Što, ktory, to and tym: reference within the sentence}\quad{\footnotesize\texttt{DEC-0124}}\par
\smallskip\noindent\emph{Razlog:} Trideset nepytatelnyh što konstrukcij potrěbovalo analizu cělogo prědloženja. Ktory jest zaimeno s imenovnym antecedentom; to može ukazati na cělo tvrđenje; tym može izraziti slědstvo; ravnoznačno izrično označaje ekvivalenciju. Ni jedno ne jest univerzalna zaměna za što. Gramatika dopušćaje čto i što v pytanju, zato zapis ne ustanavja zabranu što. Trideset konkretnyh sravnenij jest priloženo; fraza o eksponentah ima vlastnu slovničnu analizu v naslědujučem zaglavju.\par
\smallskip\noindent\emph{Reason:} Thirty non-interrogative što constructions required sentence-level analysis. Ktory has a nominal antecedent; to may refer to a whole assertion; tym may express a consequence; ravnoznačno explicitly expresses equivalence. None is a universal replacement for što. The grammar permits both čto and što in questions, so this record does not establish a ban on što. All thirty individual comparisons accompany the entry. The expression concerning exponents receives its own lexical analysis below.\par
\par\smallskip\noindent\emph{Sravnjenje formov / Comparison of forms}
\begin{longtable}{@{}P{0.43\textwidth}P{0.43\textwidth}P{0.08\textwidth}@{}}
\toprule Sravniteljnaja forma / Compared form & Izbrana forma ili funkcija / Selected form or function & Red / Line\\\midrule\endfirsthead
\toprule Sravniteljnaja forma / Compared form & Izbrana forma ili funkcija / Selected form or function & Red / Line\\\midrule\endhead
\bottomrule\endfoot
(što v istině jest slučaj & (to v istině jest slučaj & 259\\\addlinespace
identičnostij, što za pytanja & identičnostij; to za pytanja & 3690\\\addlinespace
, što imaje označiti vyše dano pisanje & ; to imaje označiti vyše dano pisanje & 3759\\\addlinespace
i togda što polary & i zato že polary & 4562\\\addlinespace
substituciju --- što možno & substituciju --- to možno & 5273\\\addlinespace
prěz što želana redukcija & tym želana redukcija & 8867\\\addlinespace
funkcij, prěz što nastavaje & funkcij; tym nastavaje & 8877\\\addlinespace
, što daje protivrěčje. & ; to daje protivrěčje. & 11386\\\addlinespace
derivatov, što protivorěčilo & derivatov; to protivorěčilo & 12009\\\addlinespace
ili, što jest to samo, jednakogo & ili, ravnoznačno, jednakogo & 12013\\\addlinespace
što davaje obratno pridruženje & to davaje obratno pridruženje & 12175\\\addlinespace
što možno takože pisati & to možno takože pisati & 14115\\\addlinespace
-- što (146) osoblivo & -- to (146) osoblivo & 14124\\\addlinespace
abo, što jest s tym identično, nuly & abo, ravnoznačno, nuly & 14186\\\addlinespace
što po (18) jest ravnoznačno s: & to po (18) jest ravnoznačno s: & 14578\\\addlinespace
togo, što bylo dokazano vyšše & togo, ktoro bylo dokazano vyšše & 14784\\\addlinespace
faktorov -- što vodi & faktorov -- to vodi & 14902\\\addlinespace
, što možno opustiti s obzirom na slědujući paragraf & ; to možno opustiti s obzirom na slědujući paragraf & 14948\\\addlinespace
nulovogo ideala, što imaje město & nulovogo ideala; to imaje město & 14971\\\addlinespace
pojedinačno --- što tu ponovno & pojedinačno --- to tu ponovno & 15076\\\addlinespace
--- na što Fischer & --- na to Fischer & 15139\\\addlinespace
abo --- što zaradi vzajemnoj racionalnoj izrazimosti jest ravnoznačno --- & abo --- ravnoznačno zaradi vzajemnoj racionalnoj izrazimosti --- & 15260\\\addlinespace
abo, što vyhodi na to samo, & abo, ravnoznačno, & 15270\\\addlinespace
ograničiti --- što togda ne označaje & ograničiti --- to togda ne označaje & 15282\\\addlinespace
--- abo što po 6. jest ravnoznačno, ako & --- abo, ravnoznačno po 6., ako & 15316\\\addlinespace
prostyh idealov, što ponovno bude nezavisno od vvedenja transformovanyh idealov, trěba & prostyh idealov (to ponovno bude nezavisno od vvedenja transformovanyh idealov), trěba & 15488\\\addlinespace
komponent po I, što tu vyhodi & komponent po I, ktoro tu vyhodi & 15826\\\addlinespace
Imenno, što jest vsegda možno, vozmimo & Imenno, kako jest vsegda možno, vozmimo & 19014\\\addlinespace
abo, što zbog & abo, zbog & 23104\\\addlinespace
ne vhodi), što ponovno daje & ne vhodi); to ponovno daje & 23114\\\addlinespace
\end{longtable}
\smallskip\noindent\emph{Izvory:} Jan van Steenbergen, Gramatika: zaimena, relativne i pytateljne formy.\par
\noindent\emph{Sources:} Grammar: Pronouns, relative and interrogative pronouns\par
\par\medskip\Needspace{8\baselineskip}
\hypertarget{DEC-0125}{}
\noindent\textbf{Znak prěd jednoj tretinoju / The sign before one third}\quad{\footnotesize\texttt{DEC-0125}}\par
\smallskip\noindent\emph{Razlog:} V Paper 2, §11, arhivna kolacija ukazuje odlišen minus prěd prvym členom: ∼−⅓ … −⅙ … vměsto ∼⅓ … −⅙ …. Razlog ne jest jezyčna preferencija: prethodna negativna veličina sama ima koeficient −⅓, i otčitanje dava plus. Zapis upoređaje německy tekst i tiskany izvor; ovde se objasnjaje ta sačuvana kolacija, ne tvrdi nova vizualna prověrka izvora. Takovu korekciju treba razlišiti od avtorskogo znaka i od izbora termina.\par
\smallskip\noindent\emph{Reason:} In Paper 2, §11, the archived collation identifies an extra minus before the first term: ∼−⅓ … −⅙ … instead of ∼⅓ … −⅙ …. The rationale is not linguistic preference: the preceding quantity being subtracted itself has coefficient −⅓, so subtraction gives a positive term. The record compares German text and print. This entry explains that preserved collation; it does not claim a new visual examination of the primary source. Such a correction must be distinguished from the author's notation and from terminology.\par
\par\smallskip\noindent\emph{Sravnjenje formov / Comparison of forms}
\begin{longtable}{@{}P{0.43\textwidth}P{0.43\textwidth}P{0.08\textwidth}@{}}
\toprule Sravniteljnaja forma / Compared form & Izbrana forma ili funkcija / Selected form or function & Red / Line\\\midrule\endfirsthead
\toprule Sravniteljnaja forma / Compared form & Izbrana forma ili funkcija / Selected form or function & Red / Line\\\midrule\endhead
\bottomrule\endfoot
\textbackslash{}\allowbreak{}sim-\textbackslash{}\allowbreak{}frac13(\textbackslash{}\allowbreak{}vartheta\textbackslash{}\allowbreak{}R x)-\textbackslash{}\allowbreak{}frac16s\_\textbackslash{}\allowbreak{}vartheta(\textbackslash{}\allowbreak{}theta su), & \textbackslash{}\allowbreak{}sim\textbackslash{}\allowbreak{}frac13(\textbackslash{}\allowbreak{}vartheta\textbackslash{}\allowbreak{}R x)-\textbackslash{}\allowbreak{}frac16s\_\textbackslash{}\allowbreak{}vartheta(\textbackslash{}\allowbreak{}theta su), & \\\addlinespace
\end{longtable}
\smallskip\noindent\emph{Izvory:} Dělo 2, §11; německy tekst ED0014, red 1796; zachovano sravnjenje matematičnyh znakov.\par
\noindent\emph{Sources:} Paper 2, §11; German ED0014, line 1796; archived sign collation\par
\par\medskip\Needspace{8\baselineskip}
\hypertarget{DEC-0126}{}
\noindent\textbf{Jednaky tekst v samostojnoj i zbirnoj formě / The same text in separate and collected form}\quad{\footnotesize\texttt{DEC-0126}}\par
\smallskip\noindent\emph{Razlog:} Samostojna statja i zbirna kniga morajut davati jednaku prěkladanu frazu i formulu. Zapis odkril, že odvojeny tekst Paper 2, §11, ne byl točno toj tekst, ktory stajal v zbirnoj knigě. To ne jest novy termin ili nova matematična odluka: potrěbno bylo sjediniti dvě predstavy istego teksta, bez vračanja starějšeho prěklada. Povtorimost zato zahtěvaje sravnenje samogo teksta, ne toliko jednakogo nazvanja dokumenta.\par
\smallskip\noindent\emph{Reason:} The separate article and collected volume must give the same translated sentence and formula. The record found that the separate text of Paper 2, §11 was not exactly the text used in the collection. This is neither a new term nor another mathematical decision: the two representations of the same text needed reconciliation without reinstating an older translation. Reproducibility therefore requires comparison of the text itself, not merely matching document titles.\par
\smallskip\noindent\emph{Izvory:} Dělo 2, §11; sravnjenje samostojnogo teksta i sbornogo izdanja.\par
\noindent\emph{Sources:} Paper 2, §11; comparison of separate and collected texts\par
\par\medskip\Needspace{8\baselineskip}
\hypertarget{DEC-0127}{}
\noindent\textbf{Čto se tyče: o čem govorit fraza / Čto se tyče: what the phrase concerns}\quad{\footnotesize\texttt{DEC-0127}}\par
\smallskip\noindent\emph{Razlog:} Čto se tyče eksponentov znači ‘o eksponentah, v otnošenju k eksponentam’. Slovnik ima cělu konstrukciju čto sę tyče, tip 2; hrvatsko što se tiče jest sravniteljna forma. Izbor cěloj konstrukcije jest silnějši od zaměny jednoj bukvy bez analizy. Osnovny pravopis čto se tyče i razširjeny čto sę tyče ne sut dva raznyh termina. Dozvoljeno pytateljno što ne ustanavjaje imenno tu predložnu frazu. Oznaka tip 2 ostaje vidima: fraza ne predstavljaje obče potvrđeno edino rěšenje.\par
\smallskip\noindent\emph{Reason:} Čto se tyče eksponentov means ‘regarding the exponents’. The dictionary records the whole construction čto sę tyče as type 2; Croatian što se tiče is a comparison form. Selecting a complete construction is stronger evidence than changing one letter without analysis. Basic čto se tyče and extended čto sę tyče are not different terms. Permitted interrogative što does not by itself establish this particular prepositional phrase. The type-2 classification remains explicit: the expression is not presented as a universally established sole solution.\par
\par\smallskip\noindent\emph{Sravnjenje formov / Comparison of forms}
\begin{longtable}{@{}P{0.43\textwidth}P{0.43\textwidth}P{0.08\textwidth}@{}}
\toprule Sravniteljnaja forma / Compared form & Izbrana forma ili funkcija / Selected form or function & Red / Line\\\midrule\endfirsthead
\toprule Sravniteljnaja forma / Compared form & Izbrana forma ili funkcija / Selected form or function & Red / Line\\\midrule\endhead
\bottomrule\endfoot
što se tiče & čto sę tyče & 11403\\\addlinespace
\end{longtable}
\smallskip\noindent\emph{Izvory:} Slovnik, zaglavja 21268\par
\noindent\emph{Sources:} Dictionary entries 21268\par
\par\medskip\Needspace{8\baselineskip}
\hypertarget{DEC-0128}{}
\noindent\textbf{Racionalny koeficient i zavisimy indeks / Rational coefficients and a bound index}\quad{\footnotesize\texttt{DEC-0128}}\par
\smallskip\noindent\emph{Razlog:} Paper 8 potrěbuje racionalnyh $c_1,c_2$, ne racionalno-cělyh $\theta_1,\theta_2$. Pridavnik racionalno-cěly by dodal uslovje cělosti, a zaměna symbolov by dala drugy objekty. V Paper 22 sredny člen jednoj λ-indeksovanoj poslědovnosti jest $B_\lambda$, ne $B_i$. Sosednje definicije i citovany německy tekst podpirajut oba čitanja. To ne opravdava obču zaměnu i s λ niti odstranjenje racionalno-cěly iz mest, kde avtor zaista govori o cělyh veličinah.\par
\smallskip\noindent\emph{Reason:} Paper 8 requires rational $c_1,c_2$, not rational-integer $\theta_1,\theta_2$. Racionalno-cěly would add an integrality condition, while the substituted symbols identify different objects. In Paper 22 the middle member of a λ-indexed sequence is $B_\lambda$, not $B_i$. The surrounding definitions and cited German text support both readings. Neither licenses replacing every i by λ or removing racionalno-cěly where the author really does require integers.\par
\par\smallskip\noindent\emph{Sravnjenje formov / Comparison of forms}
\begin{longtable}{@{}P{0.43\textwidth}P{0.43\textwidth}P{0.08\textwidth}@{}}
\toprule Sravniteljnaja forma / Compared form & Izbrana forma ili funkcija / Selected form or function & Red / Line\\\midrule\endfirsthead
\toprule Sravniteljnaja forma / Compared form & Izbrana forma ili funkcija / Selected form or function & Red / Line\\\midrule\endhead
\bottomrule\endfoot
racionalno-cělyh $\theta_1,\theta_2$ & racionalnyh $c_1,c_2$ & 6797\\\addlinespace
\textbackslash{}\allowbreak{}Bmod\_i & \textbackslash{}\allowbreak{}Bmod\_\textbackslash{}\allowbreak{}lambda & 14315\\\addlinespace
\end{longtable}
\smallskip\noindent\emph{Izvory:} Dělo 8, německy ED0014 vrsta 6043; Dělo 22, vrsty 12839–12843\par
\noindent\emph{Sources:} Paper 8, German ED0014 line 6043; Paper 22, lines 12839–12843\par
\par\medskip\Needspace{8\baselineskip}
\hypertarget{DEC-0129}{}
\noindent\textbf{K čemu odnosi se citacija teorema / What a theorem citation qualifies}\quad{\footnotesize\texttt{DEC-0129}}\par
\smallskip\noindent\emph{Razlog:} Citacija H.-N., teorem VIII, odnosi se k drugoj vrsti formule (4) v Paper 24, ne k slědujučoj prozě. Jej poměščenje jest čest smysla: čitatelj mora znati, ktoro tvrđenje podpiraje citovany rezultat. Zato citacija i jej koma stojat pri formuli, a prodolženje počinaje s malym i dalje. Samo nazvanje teorema ne bylo potrěbno prěvoditi inako. Položenje nomera formule jest oddělno tipografično pytanje.\par
\smallskip\noindent\emph{Reason:} The citation H.-N., theorem VIII qualifies the second row of formula (4) in Paper 24, not the following prose. Placement affects meaning: readers must know which assertion the cited result supports. The citation and its comma therefore belong with the formula, and the continuation begins with lowercase i dalje. The theorem's wording did not require a different translation. The equation number's position is a separate typographic question.\par
\par\smallskip\noindent\emph{Sravnjenje formov / Comparison of forms}
\begin{longtable}{@{}P{0.43\textwidth}P{0.43\textwidth}P{0.08\textwidth}@{}}
\toprule Sravniteljnaja forma / Compared form & Izbrana forma ili funkcija / Selected form or function & Red / Line\\\midrule\endfirsthead
\toprule Sravniteljnaja forma / Compared form & Izbrana forma ili funkcija / Selected form or function & Red / Line\\\midrule\endhead
\bottomrule\endfoot
\& R\textasciicircum{}\{(n)\}\allowbreak{}\textbackslash{}\allowbreak{}cdots R\textasciicircum{}\{(i)\}\allowbreak{}\textbackslash{}\allowbreak{}frakg\_\{i-1\}\allowbreak{}\&\textbackslash{}\allowbreak{}equiv0\textbackslash{}\allowbreak{}pmod\{\textbackslash{}\allowbreak{}frakm\}\allowbreak{}, \textbackslash{}\allowbreak{}tag\{4\}\allowbreak{} & \& R\textasciicircum{}\{(n)\}\allowbreak{}\textbackslash{}\allowbreak{}cdots R\textasciicircum{}\{(i)\}\allowbreak{}\textbackslash{}\allowbreak{}frakg\_\{i-1\}\allowbreak{}\&\textbackslash{}\allowbreak{}equiv0\textbackslash{}\allowbreak{}pmod\{\textbackslash{}\allowbreak{}frakm\}\allowbreak{},\textbackslash{}\allowbreak{}quad\textbackslash{}\allowbreak{}text\{(H.-N., teorem VIII),\}\allowbreak{} \textbackslash{}\allowbreak{}tag\{4\}\allowbreak{} & 15306\\\addlinespace
(H.-N., teorem VIII), i dalje: & i dalje: & 15308\\\addlinespace
\end{longtable}
\smallskip\noindent\emph{Izvory:} Dělo 24, německy ED0014 vrsty 13729–13731\par
\noindent\emph{Sources:} Paper 24, German ED0014 lines 13729–13731\par
\par\medskip\Needspace{8\baselineskip}
\hypertarget{DEC-0130}{}
\noindent\textbf{Individualna norma i i zato / The individual norm and i zato}\quad{\footnotesize\texttt{DEC-0130}}\par
\smallskip\noindent\emph{Razlog:} V Paper 24 dva izpisy $E_0^{(i)}$ měli byti $E^{(i)}$; individualna norma i pozdnějši produkt ne imajut vlastnyh avtorskyh nomerov. Podlě sačuvanoj kolacije jediny istinny nomer (5) ostaje pri pozdnějšej formuli. Fraza i zato posrěd premise i zaključka izrično čuva deduktivno otnošenje. Slovnično zato jest izbrany izraz za slědstvo, ne avtomatičny prěklad vsego německogo also vo vsih kontekstah. Korekcija indeksa, nomerovanje i logična spojka potrěbujut razne vrste dokazov.\par
\smallskip\noindent\emph{Reason:} In Paper 24 two occurrences of $E_0^{(i)}$ should read $E^{(i)}$; the individual norm and a subsequent product have no equation numbers of their own. The preserved collation leaves the genuine (5) with the later formula. I zato between premises and conclusion explicitly preserves the inference. Dictionary zato is the selected result connective, not an automatic translation of every German also in every context. Index correction, numbering and a logical connective require different kinds of evidence.\par
\par\smallskip\noindent\emph{Sravnjenje formov / Comparison of forms}
\begin{longtable}{@{}P{0.43\textwidth}P{0.43\textwidth}P{0.08\textwidth}@{}}
\toprule Sravniteljnaja forma / Compared form & Izbrana forma ili funkcija / Selected form or function & Red / Line\\\midrule\endfirsthead
\toprule Sravniteljnaja forma / Compared form & Izbrana forma ili funkcija / Selected form or function & Red / Line\\\midrule\endhead
\bottomrule\endfoot
E\_0\textasciicircum{}\{(i)\}\allowbreak{}\textbackslash{}\allowbreak{}zeta\_0 & E\textasciicircum{}\{(i)\}\allowbreak{}\textbackslash{}\allowbreak{}zeta\_0 & 15287\\\addlinespace
\textbackslash{}\allowbreak{}begin\{equation\}\allowbreak{} & \textbackslash{}\allowbreak{}begin\{equation*\}\allowbreak{} & 15295\\\addlinespace
E\_0\textasciicircum{}\{(i)\}\allowbreak{}E\_1\textasciicircum{}\{(i)\}\allowbreak{} & E\textasciicircum{}\{(i)\}\allowbreak{}E\_1\textasciicircum{}\{(i)\}\allowbreak{} & 15296\\\addlinespace
\textbackslash{}\allowbreak{}end\{equation\}\allowbreak{} & \textbackslash{}\allowbreak{}end\{equation*\}\allowbreak{} & 15298\\\addlinespace
\textbackslash{}\allowbreak{}tag\{5\}\allowbreak{} & \textbackslash{}\allowbreak{}notag & 15321\\\addlinespace
\end{longtable}
\smallskip\noindent\emph{Izvory:} Dělo 24; pečatny iztočnik; zato, DEC-0044.\par
\noindent\emph{Sources:} Paper 24; printed source; zato, DEC-0044\par
\par\medskip\Needspace{8\baselineskip}
\hypertarget{DEC-0131}{}
\noindent\textbf{Prosty lěvy ideal i jego indeks / A simple left ideal and its index}\quad{\footnotesize\texttt{DEC-0131}}\par
\smallskip\noindent\emph{Razlog:} V Paper 34 multiplicitet pripada $|W_\nu|$, a odpovědajuči ideal jest imenno $\ell_\nu$. Dvě frazy morajut imenovati toj ideal; fraza o kompozicijnom faktoru mora takože čuvati lěvy. Prosty ideal bez opreděljenja strany može změniti matematično tvrđenje. Německa definicija $\ell_\nu$ kako einfaches Linksideal podpiraje prosty lěvy ideal. Dva prethodna generična $W_i$ ostajut pravilna: zavisime indekse treba razsuditi po definiciji, ne po podobnosti bukv.\par
\smallskip\noindent\emph{Reason:} In Paper 34 the multiplicity belongs to $|W_\nu|$ and the corresponding ideal is specifically $\ell_\nu$. Both relevant phrases must name that ideal; the composition-factor sentence must also retain ‘left’. A simple ideal without its handedness can change the mathematical assertion. The German definition of $\ell_\nu$ as an einfaches Linksideal supports prosty lěvy ideal. Two preceding generic $W_i$ occurrences remain correct: bound indices must be interpreted through their definitions, not through visual similarity.\par
\par\smallskip\noindent\emph{Sravnjenje formov / Comparison of forms}
\begin{longtable}{@{}P{0.43\textwidth}P{0.43\textwidth}P{0.08\textwidth}@{}}
\toprule Sravniteljnaja forma / Compared form & Izbrana forma ili funkcija / Selected form or function & Red / Line\\\midrule\endfirsthead
\toprule Sravniteljnaja forma / Compared form & Izbrana forma ili funkcija / Selected form or function & Red / Line\\\midrule\endhead
\bottomrule\endfoot
\textbackslash{}\allowbreak{}providecommand\{\textbackslash{}\allowbreak{}mc\}\allowbreak{}\{\textbackslash{}\allowbreak{}mathfrak\{c\}\allowbreak{}\}\allowbreak{} & \textbackslash{}\allowbreak{}providecommand\{\textbackslash{}\allowbreak{}mc\}\allowbreak{}\{\textbackslash{}\allowbreak{}mathfrak\{c\}\allowbreak{}\}\allowbreak{}\textbackslash{}\allowbreak{}providecommand\{\textbackslash{}\allowbreak{}ml\}\allowbreak{}\{\textbackslash{}\allowbreak{}mathfrak\{l\}\allowbreak{}\}\allowbreak{} & 22041\\\addlinespace
odgovarjajuči prosty ideal & odgovarjajuči prosty lěvy ideal \$\textbackslash{}\allowbreak{}ml\_\textbackslash{}\allowbreak{}nu\$ & 22116\\\addlinespace
\end{longtable}
\smallskip\noindent\emph{Izvory:} Dělo 34; německy tekst ED0014, redy 17850 i 18065–18069.\par
\noindent\emph{Sources:} Paper 34; German ED0014 lines 17850 and 18065–18069\par
\par\medskip\Needspace{8\baselineskip}
\hypertarget{DEC-0132}{}
\noindent\textbf{Jedno izdanje, ne směs dvuh redakcij / One edition, not a mixture of two redactions}\quad{\footnotesize\texttt{DEC-0132}}\par
\smallskip\noindent\emph{Razlog:} Paper 45 měl četyri svjazane odlomky pozdnějšej redakcije pod bibliografičnym nazvanjem statje iz Math. Annalen 97 (1927), 559–567. Izbor izvora dotyka teoremy II, lemmy II, dodatka o Bertini, eksponenta specializacije i teoremy IV s dokazami. Zaměniti toliko formulu (3) stvorilo by tretju, nesviděčenu směs. Arhivna odluka takože predlagala $xK_2/xK_3$ i $q^{(i)}$ vměsto tiskanych čitanij: to sut redakcijne emendacije, ne věrny doslovny prěpis. Poslědnja ima protivno tiskano sviděteljstvo $q_i$; bez nove kolacije ne směje byti predstavljena kako nesporna. Diplomatično izdanje i korigovano izdanje morajut byti razdělene.\par
\smallskip\noindent\emph{Reason:} Paper 45 contained four connected passages from a later redaction under the bibliographic identity of the article in Math. Annalen 97 (1927), 559–567. Choosing a witness affects Theorem II, Lemma II, the Bertini addition, the specialization exponent and Theorem IV with their proofs. Replacing equation (3) alone would create an unattested third mixture. The archived decision also proposed $xK_2/xK_3$ and $q^{(i)}$ against printed readings: these are editorial emendations, not diplomatic transcription. The latter has adverse printed evidence $q_i$ and must not be presented as settled without renewed collation. Diplomatic and corrected editions must remain distinct.\par
\smallskip\noindent\emph{Izvory:} Dělo 45, Math. Annalen 97 (1927), 559–567; sravnjenje četyreh čestij različnyh redakcij i zachovano protivno čitanje.\par
\noindent\emph{Sources:} Paper 45, Math. Annalen 97 (1927), 559–567; four-span redaction comparison and retained adverse reading\par
\par\medskip\Needspace{8\baselineskip}
\hypertarget{DEC-0133}{}
\noindent\textbf{Pet kolon i dva znaka primětky / Five columns and two note markers}\quad{\footnotesize\texttt{DEC-0133}}\par
\smallskip\noindent\emph{Razlog:} Shema v Paper 2 potrěbuje pet kolon, završnu vrstu gornjej česti i dva svjazane rozety, ktore ukazujut na primětku. Četyri kolony ne predstavjajut drukše legitimno čitanje: gubi se čest shemy. Primětka i jej priključenje morajut ostati čitateljne. Istorična izgotovka upotrěbila jest malu rasternu sliku rozety; to čuva izgled, ale ne dava polno uredimy znak. Za iskano uredimo LaTeX izdanje potrěbno jest věrno vektorno ili tipografično kodovanje. Sam slovnik ne tvrdi, že to ograničenje v korpusě uže popravjeno.\par
\smallskip\noindent\emph{Reason:} The schema in Paper 2 requires five columns, the terminal row of its upper part and two related rosettes marking a note. Four columns are not another defensible reading: part of the schema is lost. The note and its attachment must remain intelligible. The historical implementation used a small raster image of the rosette. That preserves appearance but does not provide a fully editable symbol. The requested editable LaTeX edition requires faithful vector or typographic encoding. This dictionary does not claim that the corpus has already been repaired in that respect.\par
\smallskip\noindent\emph{Izvory:} Dělo 2; německy tekst ED0014, redy 606–658; zachovana pečatna shema.\par
\noindent\emph{Sources:} Paper 2; German ED0014 lines 606–658; retained printed schema\par
\par\medskip\Needspace{8\baselineskip}
\hypertarget{DEC-0134}{}
\noindent\textbf{Nomery formul: čitatelj mora viděti cěly blok / Equation numbers identify the whole display}\quad{\footnotesize\texttt{DEC-0134}}\par
\smallskip\noindent\emph{Razlog:} Nomery (2)–(5) v Paper 24 stojat slěva; (3) i (4) sut vertikalno centrane k svojim mnogovrstnym blokam. Nomer položeny pri jednoj vrsti može nepravilno suziti obseg citacije. Izbor centrane metky podpiraje tiskany raspored, ne predpočinanje novej redakcije. Tekst citacije i matematične členy ostajut jednaki. Razměr i položenje treba proverjati na stranici, poneže pravilny LaTeX kod sam po sobě ne dokazuje čitateljny rezultat.\par
\smallskip\noindent\emph{Reason:} Numbers (2)–(5) in Paper 24 stand on the left; (3) and (4) are vertically centered on their multiline displays. A number attached to just one row can misleadingly narrow a reference's scope. Centered labels follow the printed arrangement rather than a preference for modern redesign. Citation text and mathematical terms remain unchanged. Size and placement require inspection on the page: syntactically correct LaTeX alone does not establish a readable result.\par
\smallskip\noindent\emph{Izvory:} Dělo 24; pečatny iztočnik.\par
\noindent\emph{Sources:} Paper 24; printed source\par
\par\medskip\Needspace{8\baselineskip}
\hypertarget{DEC-0135}{}
\noindent\textbf{Egzistovati i obstajati / Egzistovati and obstajati}\quad{\footnotesize\texttt{DEC-0135}}\par
\smallskip\noindent\emph{Razlog:} Slovnik dava egzistovati, egzistovanje, egzistencija i egzistencijalny. Zaměna 133 eksist- form na egzist- byla pravopisna, ne odstranjenje internacionalnoj familije v korist obstajati. Dvě familije mogu izražati bytije; njih razpreděljenje jest redakcijna konvencija, ne dva matematične smysly po nužnosti. Sačuvany sufiks egzistencijny ne jest doslovno slovnično egzistencijalny. Ta tvorba potrěbuje oddělno morfologično opravdanje i ne može dobiti status slovničnogo zaglavja samym pravilnym početkom egzist-.\par
\smallskip\noindent\emph{Reason:} The dictionary supplies egzistovati, egzistovanje, egzistencija and egzistencijalny. Replacing 133 eksist- forms with egzist- was an orthographic decision, not elimination of the international family in favor of obstajati. Both families can express existence; their distribution is an editorial convention, not necessarily a distinction between mathematical concepts. Retained egzistencijny is not the dictionary's exact egzistencijalny. That derivation needs separate morphological justification and cannot acquire headword status merely from its correctly spelled prefix.\par
\par\smallskip\noindent\emph{Sravnjenje formov / Comparison of forms}
\begin{longtable}{@{}P{0.43\textwidth}P{0.43\textwidth}P{0.08\textwidth}@{}}
\toprule Sravniteljnaja forma / Compared form & Izbrana forma ili funkcija / Selected form or function & Red / Line\\\midrule\endfirsthead
\toprule Sravniteljnaja forma / Compared form & Izbrana forma ili funkcija / Selected form or function & Red / Line\\\midrule\endhead
\bottomrule\endfoot
rule & eksist- -> egzist-; Eksist- -> Egzist- & \\\addlinespace
suffix\_policy & byte-preserve every suffix after the six-character prefix & \\\addlinespace
retained\_house\_surface & egzistencijny retains the edition's productive -cijny suffix & \\\addlinespace
\end{longtable}
\smallskip\noindent\emph{Izvory:} Slovnik, zaglavja 3600, 37680, 15832, 4087\par
\noindent\emph{Sources:} Dictionary entries 3600, 37680, 15832, 4087\par
\par\medskip\Needspace{8\baselineskip}
\hypertarget{DEC-0136}{}
\noindent\textbf{Pojedinny: jedinica razsmatrana oddělno / Pojedinny: an item considered separately}\quad{\footnotesize\texttt{DEC-0136}}\par
\smallskip\noindent\emph{Razlog:} Slovnik ima pojedinny ‘single, individual’, s češskym jednotlivý i slovačkym jednotlivý/pojediný v sravnenju. Petnadset jednotliv- fraz odnosi se k německomu einzelnen i zahtěvaje pridavnik za oddělne členy. Istorična redakcija izbrala pojedinyh, pojedinymi, pojedinoj i pojedine s jednim n. To ne jest točno citovano pojedinny s dvuma n: dokaz za značenje ne dokazuje tu pravopisnu razliku. Slovnik zato otkryto razlišaje preporučeno zaglavje i naslědovanu formu, ktora potrěbuje doslědnu prověrku.\par
\smallskip\noindent\emph{Reason:} The dictionary has pojedinny, ‘single, individual’, with Czech jednotlivý and Slovak jednotlivý/pojediný as comparators. Fifteen jednotliv- phrases correspond to German einzelnen and require an adjective for separately considered members. The historical edition selected pojedinyh, pojedinymi, pojedinoj and pojedine with one n. These are not exact forms of cited pojedinny with two n: evidence for the sense does not establish that spelling difference. The dictionary therefore distinguishes the supported headword from the inherited form requiring consistent examination.\par
\par\smallskip\noindent\emph{Sravnjenje formov / Comparison of forms}
\begin{longtable}{@{}P{0.43\textwidth}P{0.43\textwidth}P{0.08\textwidth}@{}}
\toprule Sravniteljnaja forma / Compared form & Izbrana forma ili funkcija / Selected form or function & Red / Line\\\midrule\endfirsthead
\toprule Sravniteljnaja forma / Compared form & Izbrana forma ili funkcija / Selected form or function & Red / Line\\\midrule\endhead
\bottomrule\endfoot
jednotlivyh & pojedinyh & \\\addlinespace
jednotlivymi & pojedinymi & \\\addlinespace
jednotlivoj & pojedinoj & \\\addlinespace
jednotlive & pojedine & \\\addlinespace
jednotlivé & pojedine & \\\addlinespace
\end{longtable}
\smallskip\noindent\emph{Izvory:} Slovnik, zaglavja 37972\par
\noindent\emph{Sources:} Dictionary entries 37972\par
\par\medskip\Needspace{8\baselineskip}
\hypertarget{DEC-0137}{}
\noindent\textbf{Jednako: protivopostavjenje / Jednako: contrast}\quad{\footnotesize\texttt{DEC-0137}}\par
\smallskip\noindent\emph{Razlog:} Jednako ima slovničny smysl ‘however, yet’ i stojí v gramatike medžu protivnymi spojkami. Dvadeset dva však pojavjenja odpovadajut dvadeseti aber, jednomu allerdings i jednomu dagegen. Izbor jednako čuva protivopostavjenje, ale ne dokazuje, že vsi ti německe izrazy vsigdá sut sinonimy. Však jest legitimna nacionalna forma, ne zato izbrana medžuslovjanska forma. Jednaky, jednakost i jednakovo imajut druge funkcii; podobnost korenja ne opravdava njih zaměnu so spojkoju.\par
\smallskip\noindent\emph{Reason:} Jednako has dictionary sense ‘however, yet’ and is listed as a contrastive conjunction in the grammar. The twenty-two však occurrences correspond to twenty aber, one allerdings and one dagegen. Selecting jednako preserves contrast but does not establish that these German expressions are always interchangeable. Však is a legitimate national-language form without thereby becoming the selected Interslavic form. Jednaky, jednakost and jednakovo have other functions; a shared root does not justify replacing them with the conjunction.\par
\par\smallskip\noindent\emph{Sravnjenje formov / Comparison of forms}
\begin{longtable}{@{}P{0.43\textwidth}P{0.43\textwidth}P{0.08\textwidth}@{}}
\toprule Sravniteljnaja forma / Compared form & Izbrana forma ili funkcija / Selected form or function & Red / Line\\\midrule\endfirsthead
\toprule Sravniteljnaja forma / Compared form & Izbrana forma ili funkcija / Selected form or function & Red / Line\\\midrule\endhead
\bottomrule\endfoot
však & jednako & \\\addlinespace
\end{longtable}
\smallskip\noindent\emph{Izvory:} Slovnik, zaglavje 8296; Jan van Steenbergen, Gramatika: spojky.\par
\noindent\emph{Sources:} Dictionary entries 8296; Grammar: Conjunctions\par
\par\medskip\Needspace{8\baselineskip}
\hypertarget{DEC-0138}{}
\noindent\textbf{Indeks pri koeficientu / The coefficient's index}\quad{\footnotesize\texttt{DEC-0138}}\par
\smallskip\noindent\emph{Razlog:} Dva d(y) i jedno f(0,y) v Paper 45 byly čitane kako $d_i(y)$, $d_i(y)$ i $f_i(0,y)$. Arhivna argumentacija svjazuje ih s istom i-indeksovanoju konstrukcijeju; to jest redakcijno razsudženje konkretnoj formule, ne jezyčna normalizacija. Veliko P i $h_i y^{(K)}$ ostali sut bez zaměny, poneže njih nejednoznačnost ne opravdavala samovoljnu korekciju. Istorično prijetje ne zaměnjaje izvornu kolaciju. Diplomatičny tekst mora čuvati tiskany znak, a argumentovana emendacija mora byti označena v oddělnom korigovanom izdanju.\par
\smallskip\noindent\emph{Reason:} Two d(y) occurrences and one f(0,y) in Paper 45 were read as $d_i(y)$, $d_i(y)$ and $f_i(0,y)$. The archived argument connects them to the same i-indexed construction. This is an editorial judgment about particular formulas, not linguistic normalization. Capital P and $h_i y^{(K)}$ were retained because their ambiguity did not justify an arbitrary correction. Historical acceptance does not replace primary collation. Diplomatic text must retain the printed sign; an argued emendation belongs in a separately identified corrected edition.\par
\par\smallskip\noindent\emph{Sravnjenje formov / Comparison of forms}
\begin{longtable}{@{}P{0.43\textwidth}P{0.43\textwidth}P{0.08\textwidth}@{}}
\toprule Sravniteljnaja forma / Compared form & Izbrana forma ili funkcija / Selected form or function & Red / Line\\\midrule\endfirsthead
\toprule Sravniteljnaja forma / Compared form & Izbrana forma ili funkcija / Selected form or function & Red / Line\\\midrule\endhead
\bottomrule\endfoot
d(y) & d\_\{i\}\allowbreak{}(y) at lines 479 and 501 & \\\addlinespace
f(0,y) & f\_\{i\}\allowbreak{}(0,y) at line 512 & \\\addlinespace
\end{longtable}
\smallskip\noindent\emph{Izvory:} Dělo 45; razbor četyreh nejasnyh čitanij.\par
\noindent\emph{Sources:} Paper 45; four-uncertainty adjudication\par
\par\medskip\Needspace{8\baselineskip}
\hypertarget{DEC-0139}{}
\noindent\textbf{Grupovati i pojmati: dva smysly shvat- / Grupovati and pojmati: two senses of shvat-}\quad{\footnotesize\texttt{DEC-0139}}\par
\smallskip\noindent\emph{Razlog:} Sedm shvat- pojavjenij opisuje grupovanje objektov, a šest abstraktno pojmanje. Slovnik dava grupovati i pojmati; shvatiti ima smysl fizičnogo hvatanya i jest svršeny glagol. Německy kontekst razděljuje zusammenfassen i pojmovno auffassen: jedina zaměna ne bi čuvala oba. Nesvršeny grupovati i pojmati takože trebujut pravilne finitne i participialne formy, ne prosto zaměnu korenja. Metaforično ‘grasp’ može v drugih jezykah byti legitimno; ovde izrično razděljenje uklanja nejasnost.\par
\smallskip\noindent\emph{Reason:} Seven shvat- occurrences describe grouping objects, while six concern abstract conception. The dictionary supplies grupovati and pojmati; shvatiti denotes physical grabbing and is perfective. German context distinguishes grouping or combining from conceptual apprehension: one substitution would not preserve both. Imperfective grupovati and pojmati also require the appropriate finite and participial forms rather than a root swap. Metaphorical ‘grasp’ can be legitimate in other languages; the explicit distinction here avoids ambiguity.\par
\par\smallskip\noindent\emph{Sravnjenje formov / Comparison of forms}
\begin{longtable}{@{}P{0.43\textwidth}P{0.43\textwidth}P{0.08\textwidth}@{}}
\toprule Sravniteljnaja forma / Compared form & Izbrana forma ili funkcija / Selected form or function & Red / Line\\\midrule\endfirsthead
\toprule Sravniteljnaja forma / Compared form & Izbrana forma ili funkcija / Selected form or function & Red / Line\\\midrule\endhead
\bottomrule\endfoot
shvatet & grupujut & 12738\\\addlinespace
shvaćanje & grupovanje & 13019\\\addlinespace
shvaćajut & grupujut & 13019\\\addlinespace
Shvatimo & Grupujemo & 13147\\\addlinespace
shvačene & grupovane & 13155\\\addlinespace
Shvatimo & Grupujemo & 13301\\\addlinespace
shvatiti & pojmati & 13800\\\addlinespace
shvatiti & grupovati & 14401\\\addlinespace
shvati & pojmaje & 15304\\\addlinespace
shvatiti & pojmati & 15387\\\addlinespace
shvačeno & pojmano & 15450\\\addlinespace
shvatiti & pojmati & 15914\\\addlinespace
shvata & pojmaje & 15945\\\addlinespace
\end{longtable}
\smallskip\noindent\emph{Izvory:} Slovnik, zaglavja 3950, 18292, 1051\par
\noindent\emph{Sources:} Dictionary entries 3950, 18292, 1051\par
\par\medskip\Needspace{8\baselineskip}
\hypertarget{DEC-0140}{}
\noindent\textbf{Kongruencija, ne nekongruencija / Congruence, not noncongruence}\quad{\footnotesize\texttt{DEC-0140}}\par
\smallskip\noindent\emph{Razlog:} Dva otnošenja v Paper 30, §8, sut $n\equiv0$ modulo p, ne $n\not\equiv0$ modulo p. Arhivna kolacija tiskanoj stranice 52 i německogo teksta podpiraje pozitivny znak; dokaz dalje potrěbuje imenno tu kongruenciju. Četyri sosednje negativne sravnenja ostajut negativne. Zato nije dozvoljeno odstraniti vse znaky negacije ili opraviti tekst samo po očekivanoj teoremě. Izbor mora byti lokalny, s oběma čitanjami i izvorom dostupnymi recenzentu.\par
\smallskip\noindent\emph{Reason:} Two relations in Paper 30, §8 read $n\equiv0$ modulo p, not $n\not\equiv0$ modulo p. The archived comparison of printed page 52 and German text supports congruence; the subsequent argument requires it. Four neighboring negative controls remain negative. Removing all negations, or changing the text merely to fit an expected theorem, would therefore be unjustified. The decision must be local, with both readings and the source available for review.\par
\smallskip\noindent\emph{Izvory:} Emmy Noether, Abstrakter Aufbau der Idealtheorie in algebraischen Zahl- und Funktionenkörpern, Math. Annalen 96 (1927), str. 52, DOI 10.1007/BF01209152.\par
\noindent\emph{Sources:} Emmy Noether, Abstrakter Aufbau der Idealtheorie in algebraischen Zahl- und Funktionenkörpern, Math. Annalen 96 (1927), p.52, DOI 10.1007/BF01209152\par
\par\medskip\Needspace{8\baselineskip}
\hypertarget{DEC-0141}{}
\noindent\textbf{Zaradi: příčina, ne cilj / Zaradi: a reason, not a purpose}\quad{\footnotesize\texttt{DEC-0141}}\par
\smallskip\noindent\emph{Razlog:} Devetnadset fraz s zbog vyražaje příčinu ili logičnu osnovu: německo wegen i Infolge. Slovnik dava zaradi s genitivom; jeho dva smysly ‘radi čegoto’ i ‘po příčině čegoto’ ne smějut byti poměšane. Ovde kontekst izbiraje vtory. Hrvatsko zbog i srbsko због ostajut legitimno sravniteljno sviděteljstvo. Profil razumljivosti zaradi ima i negativne i čestične ocěny; tip 1 ne znači edinoglasje. Predlog pred formuloju ne měnjaje samogo matematičnogo izraza.\par
\smallskip\noindent\emph{Reason:} Nineteen zbog phrases express a reason or logical ground, corresponding to German wegen and Infolge. The dictionary supplies zaradi with the genitive; its senses ‘for the sake of’ and ‘on account of’ must not be conflated. Context selects the latter here. Croatian zbog and Serbian због remain legitimate comparison evidence. Zaradi's intelligibility profile includes negative and partial assessments: type 1 does not mean unanimity. Before a formula the preposition does not alter the mathematical expression.\par
\par\smallskip\noindent\emph{Sravnjenje formov / Comparison of forms}
\begin{longtable}{@{}P{0.43\textwidth}P{0.43\textwidth}P{0.08\textwidth}@{}}
\toprule Sravniteljnaja forma / Compared form & Izbrana forma ili funkcija / Selected form or function & Red / Line\\\midrule\endfirsthead
\toprule Sravniteljnaja forma / Compared form & Izbrana forma ili funkcija / Selected form or function & Red / Line\\\midrule\endhead
\bottomrule\endfoot
Zbog & Zaradi & 4582\\\addlinespace
zbog & zaradi & 9919\\\addlinespace
zbog & zaradi & 9966\\\addlinespace
Zbog & Zaradi & 10025\\\addlinespace
zbog & zaradi & 11115\\\addlinespace
Zbog & Zaradi & 13467\\\addlinespace
zbog & zaradi & 13540\\\addlinespace
Zbog & Zaradi & 15998\\\addlinespace
zbog & zaradi & 16909\\\addlinespace
zbog & zaradi & 17040\\\addlinespace
zbog & zaradi & 17066\\\addlinespace
zbog & zaradi & 17066\\\addlinespace
Zbog & Zaradi & 17126\\\addlinespace
Zbog & Zaradi & 17360\\\addlinespace
zbog & zaradi & 19625\\\addlinespace
zbog & zaradi & 19707\\\addlinespace
zbog & zaradi & 22651\\\addlinespace
zbog & zaradi & 23124\\\addlinespace
zbog & zaradi & 23124\\\addlinespace
\end{longtable}
\smallskip\noindent\emph{Izvory:} Slovnik, zaglavja 21663\par
\noindent\emph{Sources:} Dictionary entries 21663\par
\par\medskip\Needspace{8\baselineskip}
\hypertarget{DEC-0142}{}
\noindent\textbf{Sostojati se iz: cělo i jego česti / Sostojati se iz: a whole and its parts}\quad{\footnotesize\texttt{DEC-0142}}\par
\smallskip\noindent\emph{Razlog:} Slovnična konstrukcija sostojati se iz imaje prezent sostoji i zahtěvaje genitiv po iz. Dvanadset finitnyh i dva participialne sostaj- pojavjenja odnose se k sostavu objekta. Sostoječu se i sostoječa se čuvajut rod i padež frazy. Jedno mesto uže imělo se, zato vměsto dodavanja vtorogo se potrěbovalo samo sostoji. Německo bestehen može takože značiti obstajanje; sam glagol ne ustanavjaje sostav bez analizy dopolnjenja. Slovničny tip 2 ostaje vidimy.\par
\smallskip\noindent\emph{Reason:} The dictionary construction sostojati se iz has present sostoji and requires the genitive after iz. Twelve finite and two participial sostaj- occurrences concern an object's composition. Sostoječu se and sostoječa se preserve agreement. One site already had se and therefore needed sostoji alone, not another reflexive particle. German bestehen can also express existence: the verb alone does not establish composition without examination of its complement. The dictionary's type-2 classification remains explicit.\par
\par\smallskip\noindent\emph{Sravnjenje formov / Comparison of forms}
\begin{longtable}{@{}P{0.43\textwidth}P{0.43\textwidth}P{0.08\textwidth}@{}}
\toprule Sravniteljnaja forma / Compared form & Izbrana forma ili funkcija / Selected form or function & Red / Line\\\midrule\endfirsthead
\toprule Sravniteljnaja forma / Compared form & Izbrana forma ili funkcija / Selected form or function & Red / Line\\\midrule\endhead
\bottomrule\endfoot
sostaji & sostoji se & 623\\\addlinespace
sostaji & sostoji se & 677\\\addlinespace
sostaji & sostoji se & 914\\\addlinespace
sostaji & sostoji se & 2392\\\addlinespace
sostaji & sostoji se & 3744\\\addlinespace
sostaji & sostoji se & 3781\\\addlinespace
sostaji & sostoji se & 3781\\\addlinespace
sostaji & sostoji & 4383\\\addlinespace
sostaji & sostoji se & 5984\\\addlinespace
sostaječu & sostoječu se & 6021\\\addlinespace
sostaji & sostoji se & 6030\\\addlinespace
sostajuča & sostoječa se & 7168\\\addlinespace
sostaji & sostoji se & 11162\\\addlinespace
sostaji & sostoji se & 21961\\\addlinespace
\end{longtable}
\smallskip\noindent\emph{Izvory:} Slovnik, zaglavje 1874; Jan van Steenbergen, Gramatika: glagoly i predlogy.\par
\noindent\emph{Sources:} Dictionary entries 1874; Grammar: Verbs, Prepositions\par
\par\medskip\Needspace{8\baselineskip}
\hypertarget{DEC-0143}{}
\noindent\textbf{Javno, jasno i jednoznačno / Javno, jasno and jednoznačno}\quad{\footnotesize\texttt{DEC-0143}}\par
\smallskip\noindent\emph{Razlog:} Javno v slovniku znači ‘publicly’; jasno znači ‘clearly’. Osm německyh offenbar/ersichtlich kontekstov potrěbuje jasno, ne javno. Sedm explizit kontekstov istorija prěvodi familijeju jednoznačn-, ktorej slovničny smysl vključaje explicit i unambiguous. Ale eksplicitnost i jednoznačnost matematično ne sut to že: formula može byti jednoznačno oprěděljena bez eksplicitnogo izraza. Zato ta polisemija potrěbuje kontekstualnu prověrku, ne opiru samo na anglijskoj glosi. Dvanadset glagolnyh jav- form jest druga kategorija i ne podleži toj zaměně.\par
\smallskip\noindent\emph{Reason:} Dictionary javno means ‘publicly’, while jasno means ‘clearly’. Eight German offenbar/ersichtlich contexts require jasno rather than javno. Seven explizit contexts were historically rendered with jednoznačn-, whose dictionary gloss includes both explicit and unambiguous. But mathematical explicitness and uniqueness are not the same: a formula can be uniquely determined without an explicit expression. This polysemy needs contextual checking, not reliance on the English gloss alone. Twelve jav- verbs belong to another category and are outside this substitution.\par
\par\smallskip\noindent\emph{Sravnjenje formov / Comparison of forms}
\begin{longtable}{@{}P{0.43\textwidth}P{0.43\textwidth}P{0.08\textwidth}@{}}
\toprule Sravniteljnaja forma / Compared form & Izbrana forma ili funkcija / Selected form or function & Red / Line\\\midrule\endfirsthead
\toprule Sravniteljnaja forma / Compared form & Izbrana forma ili funkcija / Selected form or function & Red / Line\\\midrule\endhead
\bottomrule\endfoot
javne & jednoznačne & 58\\\addlinespace
javně & jednoznačno & 173\\\addlinespace
javno & jednoznačno & 6172\\\addlinespace
javno & jednoznačno & 6172\\\addlinespace
javnyh & jednoznačnyh & 9531\\\addlinespace
javno & jednoznačno & 9698\\\addlinespace
javno & jednoznačno & 9733\\\addlinespace
javno & jasno & 11045\\\addlinespace
javno & jasno & 16014\\\addlinespace
javno & jasno & 17826\\\addlinespace
javno & jasno & 19847\\\addlinespace
javno & jasno & 21540\\\addlinespace
javno & jasno & 22132\\\addlinespace
javno & jasno & 22532\\\addlinespace
javno & jasno & 23402\\\addlinespace
\end{longtable}
\smallskip\noindent\emph{Izvory:} Slovnik, zaglavja 17607, 17620, 17625, 26327, 33082, 734, 8241\par
\noindent\emph{Sources:} Dictionary entries 17607, 17620, 17625, 26327, 33082, 734, 8241\par
\par\medskip\Needspace{8\baselineskip}
\hypertarget{DEC-0144}{}
\noindent\textbf{Očevidny i očevidno / Očevidny and očevidno}\quad{\footnotesize\texttt{DEC-0144}}\par
\smallskip\noindent\emph{Razlog:} Oba upotrěbljene slovnične izdanja davajut očevidny i očevidno; očivid- jest v něktoryh nacionalnyh sravniteljnyh formah. Trideset sedm pojavjenij prěvzima očevid-, s čuvanjem pridavnikovoj ili prislovnikovoj funkcii. To ne jest izbor medžu ‘očevidno’, ‘samorazumno’ i ‘prividno’: posljednje može značiti samo navidnu istinnost, čto by oslabilo tvrđenje. Prilaganje pravopisa ne opravdava novu ocěnu težkosti avtorskogo dokaza.\par
\smallskip\noindent\emph{Reason:} Both consulted dictionary editions give očevidny and očevidno; očivid- occurs in some national-language comparison forms. Thirty-seven occurrences take očevid- while preserving adjectival or adverbial function. This is not a choice among ‘obvious’, ‘self-evident’ and ‘apparent’: the last may weaken the truth claim. Orthographic alignment does not license a new judgment about the difficulty of the author's argument.\par
\smallskip\noindent\emph{Izvory:} Slovnik, zaglavja 3960, 3961\par
\noindent\emph{Sources:} Dictionary entries 3960, 3961\par
\par\medskip\Needspace{8\baselineskip}
\hypertarget{DEC-0145}{}
\noindent\textbf{Izvor mora byti dohvatny / A source must be retrievable}\quad{\footnotesize\texttt{DEC-0145}}\par
\smallskip\noindent\emph{Razlog:} $B_i$bliografična ili tekstova referencija mora privesti drugogo izslědovatelja k istomu dokumentu. Lokalna adresa na jednom kompjutoru ne jest dostatočna. Zapis o 219 tekstovyh edinicah razlišaje istoričny izvorny dokument od njegovogo odlomka v zbirnoj knigě. Njih razne kontrolne sumy ne sut nužno proturěčje: měrjat razne teksty. Prěnosimy ukazatelji i zachranjene avtorskyh priznanja omogožňajut povtorno sravnenje. Rezervny font čuva čitateljnost, ale ne ustanavjaje jezyčnu ili matematičnu správnost.\par
\smallskip\noindent\emph{Reason:} A bibliographic or textual reference must let another researcher retrieve the same document. An address on one computer is insufficient. The record for 219 textual units distinguishes a historical source document from its excerpt in the collection. Their different checksums are not necessarily contradictory: they measure different texts. Portable references and preserved scholarly credits make comparison repeatable. A fallback font preserves readability but establishes neither linguistic nor mathematical correctness.\par
\smallskip\noindent\emph{Izvory:} Bibliografične identičnosti 219 tekstovyh jedinic.\par
\noindent\emph{Sources:} Source identities of the 219 assembled textual units\par
\par\medskip\Needspace{8\baselineskip}
\hypertarget{DEC-0146}{}
\noindent\textbf{Imenno i ime / Imenno and ime}\quad{\footnotesize\texttt{DEC-0146}}\par
\smallskip\noindent\emph{Razlog:} Imenno jest prislovnik ‘imenno, to jest’, s německym nämlich v sravniteljnym sloupci. Dvadeset devet prislovnikovyh imeno pojavjenij potrěbuje dvojno n. Gramatika spominaje imeno v teoretično regularizovanoj paradigmě imenovnika ime; to ne jest sviděteljstvo za prislovnik. Jednaka bukvova forma v drugoj slovnoj kategoriji ne opravdava prěklad. Izbor imenno čuva objašnjenje ili nabrajanje, ne dodava novo logično slědstvo.\par
\smallskip\noindent\emph{Reason:} Imenno is the adverb ‘namely’, with German nämlich in its comparison column. Twenty-nine adverbial imeno occurrences require double n. The grammar mentions imeno in a theoretically regularized paradigm of the noun ime, ‘name’; that is not evidence for the adverb. An identical spelling in another word class does not justify a translation. Imenno preserves explanation or enumeration rather than adding a logical consequence.\par
\smallskip\noindent\emph{Izvory:} Slovnik, zaglavje 17486; Jan van Steenbergen, Gramatika: imenovniki, ime.\par
\noindent\emph{Sources:} Dictionary entries 17486; Grammar: Nouns, ime\par
\par\medskip\Needspace{8\baselineskip}
\hypertarget{DEC-0147}{}
\noindent\textbf{Jednoj: ženska forma čislovnika jedin / Jednoj: the feminine oblique of jedin}\quad{\footnotesize\texttt{DEC-0147}}\par
\smallskip\noindent\emph{Razlog:} Čislovnik jedin imaje ženske oblike jedna, jednoj i jednu. Sedmdeset pet jednej pojavjenij s ženskymi objektami potrěbuje jednoj v odpovědajučem kosom padežu. To jest paradígma, ne novy korenj i ne dokaz, že německe eines samo po sobě izbiraje žensky rod: rod kontroluje medžuslovjanski imenovnik. Šest povtorjenij v samostojnyh statjah sut kopije istih izborov, ne šest nezavisnyh jezyčnyh sviděteljstv.\par
\smallskip\noindent\emph{Reason:} The numeral jedin has feminine forms jedna, jednoj and jednu. Seventy-five jednej occurrences with feminine referents require jednoj in the relevant oblique case. This is a paradigm, not a new root; German eines alone does not select feminine gender, which is controlled by the Interslavic noun. Six repetitions in separate articles are copies of the same decisions, not six independent linguistic attestations.\par
\smallskip\noindent\emph{Izvory:} Slovnik, zaglavje 531; Jan van Steenbergen, Gramatika: čislovniki, jedin.\par
\noindent\emph{Sources:} Dictionary entries 531; Grammar: Numerals, jedin\par
\par\medskip\Needspace{8\baselineskip}
\hypertarget{DEC-0148}{}
\noindent\textbf{Citovany jezyk ostaje citovanym jezykom / A quoted language remains a quoted language}\quad{\footnotesize\texttt{DEC-0148}}\par
\smallskip\noindent\emph{Razlog:} Originalne nazvanja, imena i bibliografične izrazy ne sut neprěvedene medžuslovjanske slova. Pri prěvodu v kirilicu treba ih označiti kako cěle citovane izrazy i čuvati njihov originalny pravopis. To objašnjaje 147 takyh odlomkov v zbirnom tekstě. Označiti samo bukvu, ktora ne jest v medžuslovjanskoj azbukě, razbilo by ime na směs pism. Označiti vse latinske slova ukrylo by pravě neprěvedenu prozu. Odluka jest semantična, ne prosta proverka bukv.\par
\smallskip\noindent\emph{Reason:} Original titles, names and bibliographic expressions are not untranslated Interslavic words. Conversion to Cyrillic must identify complete quoted expressions and preserve their original spelling. This accounts for 147 such passages in the collected text. Protecting only a letter absent from the Interslavic alphabet would split a name between scripts. Protecting every Latin-looking word would conceal genuinely untranslated prose. The decision is semantic, not merely an alphabet check.\par
\smallskip\noindent\emph{Izvory:} Polny spisok izvornyh naslovov i imen; pravopisna razlika medžu citatom i prěvodnoju prozoju.\par
\noindent\emph{Sources:} Original-language title and name inventory; orthographic distinction between quotation and translated prose\par
\par\medskip\Needspace{8\baselineskip}
\hypertarget{DEC-0149}{}
\noindent\textbf{Uslovje: neutrum i soglasje / Uslovje: neuter gender and agreement}\quad{\footnotesize\texttt{DEC-0149}}\par
\smallskip\noindent\emph{Razlog:} Uslovje jest neutrum: uslovje, uslovja, uslovjem, uslovjami. Pri zaměně mužskogo ili ženskogo imenovnika treba proměniti takože njegovu okolinu: stranskoju → stranskym, ta → to, dostatočna → dostatočno, pobočnoj → pobočnogo, ktora → ktoro, toj → to, jednaky → jednake. Relativno ktoro v frazě ‘pomoču pobočnogo uslovja …, ktoro … fiksuje integral’ jest podmet svojej klauzy: genitiv antecedenta ne čini samo zaimeno genitivom. Istorična glosa ‘genitive-controlled’ byla v tom smyslu netočna. Pridavnik uslovny ostaje oddělno slovnično zaglavje.\par
\smallskip\noindent\emph{Reason:} Uslovje is neuter: uslovje, uslovja, uslovjem, uslovjami. Replacing a masculine or feminine noun also affects its surroundings: stranskoju → stranskym, ta → to, dostatočna → dostatočno, pobočnoj → pobočnogo, ktora → ktoro, toj → to, jednaky → jednake. In ‘by means of the subsidiary condition …, which … determines the integral’, ktoro is the subject of its own clause: the antecedent's genitive does not make the relative pronoun genitive. The historical gloss ‘genitive-controlled’ was inaccurate in that respect. Uslovny remains a separate adjectival headword.\par
\smallskip\noindent\emph{Izvory:} Slovnik, zaglavja 3094 i 1618; Dělo 12, aktualny tekst, red 9555; Jan van Steenbergen, Gramatika: imenovniki, pridavniky i zaimena.\par
\noindent\emph{Sources:} Dictionary entries 3094, 1618; current Paper 12 context, base line 9555; Grammar: Nouns, Adjectives, Pronouns\par
\par\medskip\Needspace{8\baselineskip}
\hypertarget{DEC-0150}{}
\noindent\textbf{Dozvaljati, jednočasno i doslědny pravopis / Dozvaljati, jednočasno and consistent spelling}\quad{\footnotesize\texttt{DEC-0150}}\par
\smallskip\noindent\emph{Razlog:} Trinadset pojavjenij obuhvata nekoliko raznyh problemov. Dozvalja znači ‘dopuščaje’ i zaměnjaje pozwalja; jednočasno znači ‘v istom vrěmeni’ v dvuh zároveň frazah; največe/veče i najvyšše/vysše morajut čuvati razliku veličiny i vysoty ili stepenja. V zapisanom tekstě stoji največše i najvyššego/najvyššej. Ostale zaměny, odpovědajut, bikvadratičnu, Pokaže, homogenna, prěvodi i sobrane, sjedinjajut pravopis ili konkretnu glagolnu formu. Polska ili češka bukva ne jest sama po sobě dokaz drugačšego matematičnogo smysla; kazdy par jest dan v tablici.\par
\smallskip\noindent\emph{Reason:} The thirteen occurrences concern several different problems. Dozvalja, ‘permits’, replaces pozwalja; jednočasno, ‘simultaneously’, replaces two zároveň occurrences. Comparatives of size must remain distinct from those of height or degree: the recorded forms are največše and najvyššego/najvyššej. The other replacements—odpovědajut, bikvadratičnu, Pokaže, homogenna, prěvodi and sobrane—standardize spelling or a particular verb form. A Polish or Czech letter alone does not prove a different mathematical meaning. Every pair is given in the comparison table.\par
\par\smallskip\noindent\emph{Sravnjenje formov / Comparison of forms}
\begin{longtable}{@{}P{0.43\textwidth}P{0.43\textwidth}P{0.08\textwidth}@{}}
\toprule Sravniteljnaja forma / Compared form & Izbrana forma ili funkcija / Selected form or function & Red / Line\\\midrule\endfirsthead
\toprule Sravniteljnaja forma / Compared form & Izbrana forma ili funkcija / Selected form or function & Red / Line\\\midrule\endhead
\bottomrule\endfoot
odpowědajut & odpovědajut & 513\\\addlinespace
biquadratičnu & bikvadratičnu & 626\\\addlinespace
Pokáže & Pokaže & 682\\\addlinespace
największe & največše & 3876\\\addlinespace
pozwalja & dozvalja & 6230\\\addlinespace
zároveň & jednočasno & 7199\\\addlinespace
zároveň & jednočasno & 8458\\\addlinespace
homogenná & homogenna & 9687\\\addlinespace
prevodí & prěvodi & 10651\\\addlinespace
najwyšego & najvyššego & 15160\\\addlinespace
najwyššej & najvyššej & 16407\\\addlinespace
sobrané & sobrane & 16493\\\addlinespace
Pokáže & Pokaže & 19395\\\addlinespace
\end{longtable}
\smallskip\noindent\emph{Izvory:} Slovnik, zaglavja 435 i 3852; Jan van Steenbergen, Gramatika: pridavniky i glagoly; Pravopis.\par
\noindent\emph{Sources:} Dictionary entries 435, 3852; Grammar: Adjectives, Verbs, Orthography\par
\par\medskip\Needspace{8\baselineskip}
\hypertarget{DEC-0151}{}
\noindent\textbf{Prědtym ili točno; izpolnjeny ili dějajuči / Prědtym or točno; fulfilled or acting}\quad{\footnotesize\texttt{DEC-0151}}\par
\smallskip\noindent\emph{Razlog:} Češke právě ne ima jednogo prěklada. Tri odvolanja na prědtym razsmotreny ili dokazany rezultat potrěbujut prědtym; pět fokusnyh izrazov, ktore ukazujut imenno na danu redukciju ili teoremu, potrěbuje točno. Działa v frazě o izpolnjenom uslovju ne znači, že uslovje ‘dějaje’: potrěbno jest jest izpolnjeny s odpovědajučim soglasjem. Użyvanjem prijma upotrěboju, instrumental od upotrěba. Korpusne paralely podpirajut izbranu formu, ale ne sut nezavisna potvrda smysla.\par
\smallskip\noindent\emph{Reason:} Czech právě has no single translation. Three references to previously considered or proved results require prědtym; five focus expressions identifying the particular reduction or theorem require točno. Działa in a statement about a fulfilled condition does not say that the condition ‘acts’: the required construction is jest izpolnjeny with the relevant agreement. Użyvanjem becomes upotrěboju, instrumental of upotrěba. Corpus parallels support the chosen form but are not independent confirmation of its meaning.\par
\par\smallskip\noindent\emph{Sravnjenje formov / Comparison of forms}
\begin{longtable}{@{}P{0.43\textwidth}P{0.43\textwidth}P{0.08\textwidth}@{}}
\toprule Sravniteljnaja forma / Compared form & Izbrana forma ili funkcija / Selected form or function & Red / Line\\\midrule\endfirsthead
\toprule Sravniteljnaja forma / Compared form & Izbrana forma ili funkcija / Selected form or function & Red / Line\\\midrule\endhead
\bottomrule\endfoot
działa & jest izpolnjeny & 495\\\addlinespace
právě & prědtym & 10509\\\addlinespace
právě & točno & 10578\\\addlinespace
právě & prědtym & 12909\\\addlinespace
právě & prědtym & 13102\\\addlinespace
działa & jest izpolnjeny & 15704\\\addlinespace
użyvanjem & upotrěboju & 16511\\\addlinespace
právě & točno & 23044\\\addlinespace
právě & točno & 23088\\\addlinespace
právě & točno & 23296\\\addlinespace
právě & točno & 143\\\addlinespace
\end{longtable}
\smallskip\noindent\emph{Izvory:} Slovnik, zaglavja 36623 i 35654, tȯčno i izpȯlniti; sravnjenja německyh vět.\par
\noindent\emph{Sources:} Dictionary entries 36623, 35654; dictionary tȯčno and izpȯlniti; German sentence comparisons\par
\par\medskip\Needspace{8\baselineskip}
\hypertarget{DEC-0152}{}
\noindent\textbf{Polno ime i město izdavanja / A complete name and place of publication}\quad{\footnotesize\texttt{DEC-0152}}\par
\smallskip\noindent\emph{Razlog:} A. Loewy, B. L. van der Waerden, Max Noether i A. Ostrowski v bibliografičnoj funkciji ostajut cělymi originalnymi imenami. Isto pravilo obejma W. Schmeidler, R. Weitzenböck, Øystein Ore, J. Cavaillès, Braunschweig i Vieweg. Dvanadset pojavjenij zahranja takože povtorjene města i izdatelja. Inicijaly i partikuly van der pripadajut imenu. To ne ustanavjaje pravopis za sklonjene medžuslovjanske izvedenice: citacija i gramatično vključenje sut razne uloge.\par
\smallskip\noindent\emph{Reason:} A. Loewy, B. L. van der Waerden, Max Noether and A. Ostrowski retain their complete original names when serving as bibliographic identities. The same rule covers W. Schmeidler, R. Weitzenböck, Øystein Ore, J. Cavaillès, Braunschweig and Vieweg. Twelve occurrences include repeated places and publisher names. Initials and particles such as van der belong to the name. This does not prescribe spelling for inflected Interslavic derivatives: citation and grammatical integration are different functions.\par
\par\smallskip\noindent\emph{Sravnjenje formov / Comparison of forms}
\begin{longtable}{@{}P{0.43\textwidth}P{0.43\textwidth}P{0.08\textwidth}@{}}
\toprule Sravniteljnaja forma / Compared form & Izbrana forma ili funkcija / Selected form or function & Red / Line\\\midrule\endfirsthead
\toprule Sravniteljnaja forma / Compared form & Izbrana forma ili funkcija / Selected form or function & Red / Line\\\midrule\endhead
\bottomrule\endfoot
A. Loewy & zachovany citat / retained citation & 121\\\addlinespace
B. L. van der Waerden & zachovany citat / retained citation & 124\\\addlinespace
Max Noether & zachovany citat / retained citation & 134\\\addlinespace
A. Ostrowski & zachovany citat / retained citation & 381\\\addlinespace
W. Schmeidler & zachovany citat / retained citation & 80\\\addlinespace
R. Weitzenböck & zachovany citat / retained citation & 94\\\addlinespace
Braunschweig & zachovany citat / retained citation & 196\\\addlinespace
Vieweg & zachovany citat / retained citation & 196\\\addlinespace
Øystein Ore & zachovany citat / retained citation & 217\\\addlinespace
Vieweg & zachovany citat / retained citation & 217\\\addlinespace
Braunschweig & zachovany citat / retained citation & 217\\\addlinespace
J. Cavaillès & zachovany citat / retained citation & 220\\\addlinespace
\end{longtable}
\smallskip\noindent\emph{Izvory:} Dělo 45 i bibliografija: polny spisok imen ljudij i měst.\par
\noindent\emph{Sources:} Paper 45 and bibliography: complete name/place inventory\par
\par\medskip\Needspace{8\baselineskip}
\hypertarget{DEC-0153}{}
\noindent\textbf{Curišsko: pridavnik od Zürich / Curišsko: an adjective referring to Zurich}\quad{\footnotesize\texttt{DEC-0153}}\par
\smallskip\noindent\emph{Razlog:} Moje curišsko predavanje prěvodi meinen Züricher Vortrag: ide o grad Zürich, ne o Švejcariju. Slovnik nema potvrđeno zaglavje za toj pridavnik. Curišsko jest zato predlog izdanja, osnovany na dopuščenoj azbukě, početnom c v ruskoj i ukrajinskoj adaptaciji i analogiji Čeh → češsky. Ta analogija ne jest obče pravilo za vse tuđozemne imena. Zürihsko čuva nedopuščeno ü v domačej prozě; švejcarsko bi změnilo referent. Sačuvana ocěna 0,80 jest redakcijna uvěrenost, ne izměrjena věrojatnost razuměnja. Nacionalne AI prěklady ne sut nezavisni ljudski svidětelji.\par
\smallskip\noindent\emph{Reason:} Moje curišsko predavanje renders meinen Züricher Vortrag: the reference is to Zurich, not Switzerland. No established dictionary headword was found for the adjective. Curišsko is therefore an editorial proposal based on the permitted alphabet, initial c in the Russian/Ukrainian adaptations and the analogy Čeh → češsky. That analogy is not a universal rule for foreign names. Zürihsko retains non-Interslavic ü in integrated prose; švejcarsko changes the referent. The recorded 0.80 is editorial confidence, not a measured probability of comprehension. AI-translated national-language witnesses are not independent human attestations.\par
\smallskip\noindent\emph{Izvory:} Slovnik, zaglavja 4730 i 2417; německy tekst ED0015, red 19798; Pravopis: inostranna bukva ü.\par
\noindent\emph{Sources:} Dictionary entries 4730, 2417; German ED0015 line 19798; orthography on foreign ü\par
\par\medskip\Needspace{8\baselineskip}
\hypertarget{DEC-0154}{}
\noindent\textbf{Rimsky nomer ne jest slovo / A Roman numeral is not a word}\quad{\footnotesize\texttt{DEC-0154}}\par
\smallskip\noindent\emph{Razlog:} I, V, X i njihove kombinacije mogu byti rimsky nomer, matematičny symbol, inicijal ili čest slova. Nijedna obča zaměna bukv ne čuva vse te uloge. Istorična analiza razlišila 722 točno uměstjene strukturne nomery od imen i prozy; 692 jednobukvene I/V/M pojavjenja potrěbovala sut dodatnu kontekstualnu prověrku. Čislo nepodržanyh znakova zato ne jest měra jezyčnoj kačestva. Kiriličny tekst jest prěpis pismena iz latinskogo prěklada, ne novy nezavisny prěklad.\par
\smallskip\noindent\emph{Reason:} I, V, X and their combinations can be Roman numerals, mathematical symbols, initials or parts of words. No indiscriminate letter substitution preserves all those roles. The historical analysis distinguished 722 specifically located structural numerals from names and prose; 692 single-letter I/V/M occurrences required additional contextual examination. A count of unsupported characters is therefore not a measure of linguistic quality. The Cyrillic text is a script conversion of the Latin translation, not an independent translation.\par
\smallskip\noindent\emph{Izvory:} Spisok rimskyh čisel i jednobukvenyh izrazov s funkcijeju každogo pojavjenja; pravilo pisemnogo prěpisa.\par
\noindent\emph{Sources:} Roman-numeral and one-letter role inventories; script-conversion policy\par
\par\medskip\Needspace{8\baselineskip}
\hypertarget{DEC-0158}{}
\noindent\textbf{Imena v prěkladanoj prozě / Names integrated into translated prose}\quad{\footnotesize\texttt{DEC-0158}}\par
\smallskip\noindent\emph{Razlog:} Dvadset imenovyh familij imaje 66 sklonjenyh ili izvedenyh pojavjenij. Predložene forme zahranjajut Enrikes, Furtvengler, Grasman, Getingen, Hurvic, Helder, Kenig, Kenigsberg, Levi, Lurot, Ostrovsk-, Švarc, Suškevič, Silov, Varden, Veber, Vederbern, Vajerštras, Velštajn i Vejl. Polna tablica čuva istorične argumenty i ocěny za vsaku familiju. Osoblivo Levi za Loewy može byti poměšano s drugim prezimenem; odsutnost w/ö v azbukě sama ne izbiraje samoglasnik. Ruske/ukrajinske prěklady i notranja doslědnost podpirajut predlog, ale ne dokazujut obče medžuslovjansko pravilo. Čislene ocěny uvěrjenosti sut historijne redakcijne sude, ne kalibrovane pravdopodobnosti. Posebne objasnjenja v tabeli sut pozdnějša kritična analiza zachovanyh razlogov.\par
\smallskip\noindent\emph{Reason:} Twenty name families account for 66 inflected or derived occurrences. Proposed forms include Enrikes, Furtvengler, Grasman, Getingen, Hurvic, Helder, Kenig, Kenigsberg, Levi, Lurot, Ostrovsk-, Švarc, Suškevič, Silov, Varden, Veber, Vederbern, Vajerštras, Velštajn and Vejl. The complete table preserves the historical grounds and confidence for every family. In particular, Levi for Loewy can be confused with another surname; an alphabet's exclusion of w/ö does not itself select a vowel. Russian/Ukrainian translations and internal consistency support a proposal, not a general Interslavic standard. The numerical confidence estimates are historical editorial judgments, not calibrated probabilities. The individual table explanations are a later critical examination of the preserved reasons.\par
\par\smallskip\noindent\emph{Sravnjenje formov / Comparison of forms}
\begin{longtable}{@{}P{0.43\textwidth}P{0.43\textwidth}P{0.08\textwidth}@{}}
\toprule Sravniteljnaja forma / Compared form & Izbrana forma ili funkcija / Selected form or function & Pojav. / Uses\\\midrule\endfirsthead
\toprule Sravniteljnaja forma / Compared form & Izbrana forma ili funkcija / Selected form or function & Pojav. / Uses\\\midrule\endhead
\bottomrule\endfoot
Enriques → Enrikes; 2 pojav. / uses; 0.99 & Prěvod imaje kontrolnu formu Enrikes, podobnu formam v ruskoj i ukrajinskoj verziji. To podpiraje vnutrišne soglasje, ne nezavisno zasvědočenje v ljudskoj literaturě. EN: The existing Interslavic Enrikes and comparable forms in the Russian/Ukrainian translations support internal consistency, not independent attestation in human scholarship. & \\\addlinespace
Furtwängler → Furtvengler; 1 pojav. / uses; 0.96 & Ruska i ukrajinska verzije dajut Furtvengler. Vyljučenje w i ä iz osnovnogo alfabeta ne oprěděljaje samo po sobě izbor e za ä; to jest redakcijno prědloženje. EN: The Russian/Ukrainian versions use Furtvengler. Excluding w and ä from the basic alphabet does not itself determine e for ä; that vowel is an editorial proposal. & \\\addlinespace
Graßmann → Grasman; 1 pojav. / uses; 0.94 & Izbor zaměnjaje ß jednim s i uproščaje dvojno nn. Zapis odkazuje na slovjansko uproščenje soglasnyh, ale ne daje samostojny oficialny vzorec za to ime. EN: The choice writes ß as a single s and simplifies nn. The record invokes Slavic consonant simplification but supplies no independent official paradigm for this name. & \\\addlinespace
Göttingen → Getingen; 2 pojav. / uses; 0.99 & Kontrolna forma Getingen jest v medžuslovjanskom tekstu i v dvoh sravnitelnyh prěvodah. Podobnost podpiraje doslědnost izbora, ne obču normu za německo ö. EN: Getingen occurs in the Interslavic text and both comparison translations. Their agreement supports consistency, not a general norm for German ö. & \\\addlinespace
Hurwitz → Hurvic; 3 pojav. / uses; 0.97 & Sravnitelne formy Gurvic ne sut prosto kopirovane: Hurvic zachovuje h i prědaje w i tz kako v i c. Razlika početnogo soglasnogo mora ostati vidima. EN: The comparanda Gurvic are not copied wholesale: Hurvic retains h and represents w and tz as v and c. The difference in the initial consonant must remain visible. & \\\addlinespace
Hölder → Helder; 3 pojav. / uses; 0.91 & Sravnjenja Gelder i Gyolder ne dajut jedinu samoglasnu normu. Helder zachovuje h; izbor e jest prědložena prilagodba, ne slědstvo samogo vyljučenja ö. EN: The Gelder/Gyolder comparanda do not establish one vowel norm. Helder retains h; e is a proposed adaptation, not a consequence of merely excluding ö. & \\\addlinespace
König → Kenig; 1 pojav. / uses; 0.97 & Kenig odgovarjaje formam v ruskoj i ukrajinskoj verziji. Otsustvje ö v osnovnom alfabetě objasnjaje potrěbu prilagodby, ale ne dokazuje, že e jest jediny izbor. EN: Kenig agrees with the Russian/Ukrainian versions. The absence of ö in the basic alphabet motivates adaptation but does not prove e is the only choice. & \\\addlinespace
Königsberg → Kenigsberg; 1 pojav. / uses; 0.98 & Kenigsberg razširjaje to samo pravilo, kako Kenig. Srodnost imen podpiraje doslědnost, ale ne tvori dodatny nezavisny dokaz za samoglasnu. EN: Kenigsberg extends the same convention as Kenig. The names' relationship supports consistency but provides no additional independent evidence for the vowel. & \\\addlinespace
Loewy → Levi; 2 pojav. / uses; 0.84 & Rusko Lyovi i ukrajinsko Levi razhode se. Izbor Levi jest redakcijna preferencija; tvrdženje največšej razumlivosti nebylo izměreno i isto pisanje može označati drugo prezvime. EN: Russian Lyovi and Ukrainian Levi diverge. Levi is an editorial preference; maximal intelligibility was not measured, and the spelling can coincide with another surname. & \\\addlinespace
Lüroth → Lurot; 15 pojav. / uses; 0.99 & Postoječe kontrolne formy Lurot podpirajut doslědnost. Izbor u za ü i t za th jest prilagodba imena; sam alfabet ne dokazuje vsaku čest toj prilagodby. EN: Existing Lurot controls support consistency. The u for ü and t for th are name adaptations; the alphabet alone does not establish every part of them. & \\\addlinespace
Ostrowsk → Ostrovsk; 1 pojav. / uses; 0.99 & Ostrovski jest kontrolna forma v korpusu i v dvoh sravnitelnyh verzijah. Osnova Ostrovsk- zachovuje sklonjenje imena pri zaměně w na v. EN: Ostrovski is an existing corpus form shared by both comparison versions. The stem Ostrovsk- retains inflection while replacing w with v. & \\\addlinespace
Schwarz → Švarc; 1 pojav. / uses; 0.99 & Kontrolno Švarcovomu i sravnitelno Švarc podpirajut osnovu Švarc-. Sklonjeni eponim jest različen od izvorno zapisanogo imena Schwarz v citatu. EN: Existing Švarcovomu and comparative Švarc support Švarc-. The inflected eponym is distinct from the source spelling Schwarz in a citation. & \\\addlinespace
Suschkewitsch → Suškevič; 1 pojav. / uses; 0.99 & Suškevič sleduje formam v ruskoj i ukrajinskoj verziji i prědaje slovjansko ime bez německoj ortografije. Tute prěvody ne sut nezavisni ljudski slovniki. EN: Suškevič follows the Russian/Ukrainian versions and represents the Slavic name without German orthography. Those translations are not independent human dictionaries. & \\\addlinespace
Sylow → Silov; 4 pojav. / uses; 0.99 & Sravnitelne prěvody imajut eponim Silov. Historijno tvrdženje o ustanovjenoj matematičnoj upotrěbě trěbuje konkretny odborny izvor; same sravnitelne prěvody go ne zaměnjajut. EN: The comparison translations use Silov. The historical claim of established mathematical usage needs a specific disciplinary source; the comparison translations do not supply one by themselves. & \\\addlinespace
Waerden → Varden; 3 pojav. / uses; 0.98 & Varden odgovarjaje sravnitelnyma prěvodama i zapisu v vměsto w. Česti van der i sklonjeni konec imena trěbujut byti razgledane v cěloj frazě. EN: Varden agrees with the comparison translations and the use of v for w. The particles van der and the inflection must be assessed within the whole phrase. & \\\addlinespace
Weber → Veber; 11 pojav. / uses; 0.99 & Tri korpusne kontrolne formy Veber i ruska/ukrajinska sravnjenja podpirajut jednotny zapis. Broj kontrol ne jest jezyčno glasovanje. EN: Three Veber corpus controls and Russian/Ukrainian comparanda support a consistent spelling. The number of controls is not a linguistic vote. & \\\addlinespace
Wedderburn → Vederbern; 7 pojav. / uses; 0.96 & Sravnjenja Vedderbern podpirajut samoglasnu e; Vederbern dopolnitelno uproščaje dd i zaměnjaje w. To sut oddělne odluke, ne jedna mehanična zaměna bukvy. EN: The comparanda Vedderbern support e; Vederbern additionally simplifies dd and replaces w. These are separate decisions, not one mechanical letter replacement. & \\\addlinespace
Weierstrass → Vajerštras; 5 pojav. / uses; 0.93 & Historijni zapis navodi Klajn i Štejnic kako dokaz za ei → aj, ale Štejnic sodržuje ej, ne aj. Tute kontrole zato ne dokazujut jedino pravilo. Vajerštras ostavaje redakcijno prědloženje s v i jednim s. EN: The historical record cites Klajn and Štejnic for ei → aj, but Štejnic contains ej, not aj. Those controls therefore do not establish one rule. Vajerštras remains an editorial proposal with v and a single s. & \\\addlinespace
Wellstein → Velštajn; 1 pojav. / uses; 0.87 & Velštajn zaměnjaje w i uproščaje ll. I zde kontrolne Klajn i Štejnic ne dajut jednaku prilagodbu ei; izbor aj ne mora byti predstavjeny kako jednoznačno izvedeny iz obojih formov. EN: Velštajn replaces w and simplifies ll. Here too Klajn and Štejnic exemplify different treatments of ei; aj must not be presented as uniquely derived from both controls. & \\\addlinespace
Wejl → Vejl; 1 pojav. / uses; 0.99 & Naslědžena sklonjena osnova Wejl već sodržuje ejl. Odluka Vejl měnjaje početno w na v, a ne rekonstruuje novo vsi samoglasny. EN: The inherited inflected stem Wejl already contains ejl. Vejl changes initial w to v rather than reconstructing the vowels anew. & \\\addlinespace
\end{longtable}
\smallskip\noindent\emph{Izvory:} Razbor dvadeseti familij vlastnyh imen; izvorne formy i historijne ocěny uvěrjenosti sut zachovane v prilogu podatkov.\par
\noindent\emph{Sources:} Twenty-family proper-name adjudication; original forms and historical confidence preserved in the accompanying table\par
\par\medskip\Needspace{8\baselineskip}
\hypertarget{DEC-0159}{}
\noindent\textbf{Ime v kurzivě ostaje cělo ime / An italicized name is still a complete name}\quad{\footnotesize\texttt{DEC-0159}}\par
\smallskip\noindent\emph{Razlog:} Tipografija ne određaje jezyk. Petnadset uže odvojenyh imen, nazvanij i podpisa, vključajuči E. Müller, Weber, Du Pasquier, Loewy, Poincaré i Emmy Noether, Göttingen, čuvajut cělu originalnu formu v kurzivě ili drugom slogovom obliku. Město i inicijal ne smějut ostati napol prěpisane. Odluka se odnosi k cělomu izrazu, ne samo k bukvě ü, w ili é.\par
\smallskip\noindent\emph{Reason:} Typography does not determine language. Fifteen already delimited names, titles and signatures—including E. Müller, Weber, Du Pasquier, Loewy, Poincaré and Emmy Noether, Göttingen—retain their complete original form inside their existing typography. Places and initials must not be half-converted. The decision concerns the whole expression, not just ü, w or é.\par
\par\smallskip\noindent\emph{Sravnjenje formov / Comparison of forms}
\begin{longtable}{@{}P{0.43\textwidth}P{0.43\textwidth}P{0.08\textwidth}@{}}
\toprule Sravniteljnaja forma / Compared form & Izbrana forma ili funkcija / Selected form or function & Red / Line\\\midrule\endfirsthead
\toprule Sravniteljnaja forma / Compared form & Izbrana forma ili funkcija / Selected form or function & Red / Line\\\midrule\endhead
\bottomrule\endfoot
E. Müller & zachovany citat / retained citation & 3709\\\addlinespace
Müller & zachovany citat / retained citation & 3711\\\addlinespace
Weber & zachovany citat / retained citation & 7232\\\addlinespace
Göttingen & zachovany citat / retained citation & 8906\\\addlinespace
Nachr. v. d. Ges. d. Wiss. zu Göttingen 1918, s. 37--44 & zachovany citat / retained citation & 8933\\\addlinespace
Nachr. v. d. Ges. d. Wiss. zu Göttingen 1918, s. 235--257 & zachovany citat / retained citation & 9187\\\addlinespace
Du Pasquier & zachovany citat / retained citation & 11090\\\addlinespace
Loewy & zachovany citat / retained citation & 11350\\\addlinespace
Poincaré & zachovany citat / retained citation & 12036\\\addlinespace
A.\textasciitilde{}Ostrowski & zachovany citat / retained citation & 13809\\\addlinespace
A.\textasciitilde{}Ostrowski & zachovany citat / retained citation & 13861\\\addlinespace
Emmy Noether, Göttingen & zachovany citat / retained citation & 23227\\\addlinespace
\end{longtable}
\smallskip\noindent\emph{Izvory:} Polny spisok sravnjenyh inostrannyh izrazov, odděl 159.\par
\noindent\emph{Sources:} Complete foreign-expression comparison inventory, group 159\par
\par\medskip\Needspace{8\baselineskip}
\hypertarget{DEC-0160}{}
\noindent\textbf{Hurwitz v citaciji, Hurvica v prozě / Hurwitz in a citation, Hurvica in prose}\quad{\footnotesize\texttt{DEC-0160}}\par
\smallskip\noindent\emph{Razlog:} Dvadset nesklonjenyh imen čuva originalny zapis: medžu nimi Enriques, $W_i$rtinger, Weierstrass, Loewy, Weyl i Du Pasquier. Poslědnje jest cělo prezimeno s partikuloju, ne samo Pasquier. V izrazu Hurwitz tamže ime ima citacijsku funkciju; sklonjeno Hurvica jest gramatično uključeno v medžuslovjansku frazu. Zato dvě forme mogu legitimno obstajati v jednoj knigě, ako njih razlika ostaje jasna.\par
\smallskip\noindent\emph{Reason:} Twenty uninflected names retain their original spelling, among them Enriques, $W_i$rtinger, Weierstrass, Loewy, Weyl and Du Pasquier. The last is the complete surname with its particle, not Pasquier alone. Hurwitz in Hurwitz tamže functions as a citation, whereas inflected Hurvica is integrated into an Interslavic phrase. Both forms can therefore legitimately occur in one book when their functions remain clear.\par
\smallskip\noindent\emph{Izvory:} Polny spisok sravnjenyh vlastnyh imen, odděl 160.\par
\noindent\emph{Sources:} Complete personal-name comparison inventory, group 160\par
\par\medskip\Needspace{8\baselineskip}
\hypertarget{DEC-0161}{}
\noindent\textbf{Inicijaly i partikuly v imenah / Initials and particles in names}\quad{\footnotesize\texttt{DEC-0161}}\par
\smallskip\noindent\emph{Razlog:} J. König, A. Hurwitz, G. Köthe, A. Châtelet i B. L. van der Waerden čuvajut inicijaly i partikuly skupa s prezimenem. Dvadset zabilěženih pojavjenij ne ustanavjaje dvadset novih terminov: priměnjaje jedno pravilo k točno identifikovanym imenam. Skryti samo Waerden, ostavljajuči van der za avtomatičny prěpis, dalo by směšano ime.\par
\smallskip\noindent\emph{Reason:} J. König, A. Hurwitz, G. Köthe, A. Châtelet and B. L. van der Waerden retain initials and particles with their surnames. The twenty recorded occurrences do not establish twenty new terms: they apply one rule to identified names. Protecting Waerden alone while automatically converting van der would produce a mixed-script name.\par
\smallskip\noindent\emph{Izvory:} Polny spisok sravnjenyh vlastnyh imen, odděl 161.\par
\noindent\emph{Sources:} Complete personal-name comparison inventory, group 161\par
\par\medskip\Needspace{8\baselineskip}
\hypertarget{DEC-0162}{}
\noindent\textbf{Svjazane imena dvuh avtorov / Paired authors' names}\quad{\footnotesize\texttt{DEC-0162}}\par
\smallskip\noindent\emph{Razlog:} G. Wallenberg–A. Guldberg i Artin–van der Waerden sut svjazane avtorske identifikacije. Ih česti, inicijaly i spojny znak treba čuvati skupa. Isto pravilo kontroluje ostale zabilěžene nesklonjene imena, kako W. Schmeidler, R. Weitzenböck i H. Prüfer. To čuva naučno priznanje i možnost najdti izvor; ne govori, kak treba tvoriti sklonjenu izvedenicu prezimena.\par
\smallskip\noindent\emph{Reason:} G. Wallenberg–A. Guldberg and Artin–van der Waerden are paired author identities. Components, initials and connecting punctuation belong together. The same rule governs the other recorded uninflected names, including W. Schmeidler, R. Weitzenböck and H. Prüfer. It preserves scholarly credit and source retrieval; it does not prescribe inflected derivatives of those surnames.\par
\smallskip\noindent\emph{Izvory:} Polny spisok sravnjenyh vlastnyh imen, odděl 162.\par
\noindent\emph{Sources:} Complete personal-name comparison inventory, group 162\par
\par\medskip\Needspace{8\baselineskip}
\hypertarget{DEC-0163}{}
\noindent\textbf{Originalno nazvanje i město / An original title and a place name}\quad{\footnotesize\texttt{DEC-0163}}\par
\smallskip\noindent\emph{Razlog:} Algebraische Theorie der Körper i Rationale Funktionenkörper ostajut originalnymi nazvanjami, kako i polno nazvanje s Lorentzgruppe. Göttinger Nachrichten i $W_i$ener Sitzungsberichte sut časopisne ili institucionalne nazvanja, ne prislovnikove frazy prěklada. Göttingen i Dixmuiden označajut města. Dvadset izrazov v tablici jest klasifikovano po toj funkcii; čuvanje nazvanja ne směje obejmati sosednju medžuslovjansku frazu.\par
\smallskip\noindent\emph{Reason:} Algebraische Theorie der Körper and Rationale Funktionenkörper remain original titles, as does the complete title containing Lorentzgruppe. Göttinger Nachrichten and $W_i$ener Sitzungsberichte are publication or institutional names, not translated adverbial phrases. Göttingen and Dixmuiden identify places. The twenty expressions in the table are classified by function; preserving a title must not absorb neighboring Interslavic prose.\par
\smallskip\noindent\emph{Izvory:} Polny spisok naslovov, publikacij i měst, odděl 163.\par
\noindent\emph{Sources:} Complete title, venue and place inventory, group 163\par
\par\medskip\Needspace{8\baselineskip}
\hypertarget{DEC-0164}{}
\noindent\textbf{Jezyk bibliografije ne zavisi od jezyka knigy / Bibliographic language need not match the book's language}\quad{\footnotesize\texttt{DEC-0164}}\par
\smallskip\noindent\emph{Razlog:} The Madison Colloquium 1913 čuva anglijsko nazvanje, Neue Begründung der algebraischen Zahlentheorie německo. R. Hölzer i W. Krull sut imena, Zürich i Göttingen města. Dvadset takyh izrazov čuva se v originalnoj formě zato, že ih funkcija jest identifikacija, ne objašnjenje matematiky. Slovjansko objašnjenje možno dodati oddělno, ale ne vměsto originalnogo nazvanja.\par
\smallskip\noindent\emph{Reason:} The Madison Colloquium 1913 retains its English title and Neue Begründung der algebraischen Zahlentheorie its German title. R. Hölzer and W. Krull identify people; Zürich and Göttingen identify places. Twenty expressions retain their original forms because their function is identification, not mathematical explanation. A Slavic gloss can be supplied separately, but should not replace the original title.\par
\smallskip\noindent\emph{Izvory:} Polny spisok naslovov, publikacij, imen i měst, odděl 164.\par
\noindent\emph{Sources:} Complete title, venue, name and place inventory, group 164\par
\par\medskip\Needspace{8\baselineskip}
\hypertarget{DEC-0165}{}
\noindent\textbf{Jordan–Hölder: citovano ime teoremy / Jordan–Hölder: a cited theorem name}\quad{\footnotesize\texttt{DEC-0165}}\par
\smallskip\noindent\emph{Razlog:} Jordan–Hölder i Hilbert–Furtwängler v citovanoj formě čuvajut oba prezimena. Isto pravilo obejma nazvanja i časopisne izrazy v ostalyh petnadseti zapisanyh cělinah. Dva nestandardne znaky v jednom nazvanju ne činjut dva termina. Cěla bibliografična fraza može zahranjati godinu i stranicu; njen obseg treba določiti po čitanju, ne po čislu bukv, ktore program ne poznaje.\par
\smallskip\noindent\emph{Reason:} Jordan–Hölder and Hilbert–Furtwängler retain both surnames in their cited forms. The rule also covers titles and publication expressions among the fifteen recorded units. Two nonstandard letters within one title do not make two terms. A complete bibliographic expression may include its year and pages; its extent must be determined by reading, not by counting characters a program does not recognize.\par
\smallskip\noindent\emph{Izvory:} Polny spisok imen, naslovov i citatov, odděl 165.\par
\noindent\emph{Sources:} Complete name/title/citation inventory, group 165\par
\par\medskip\Needspace{8\baselineskip}
\hypertarget{DEC-0166}{}
\noindent\textbf{Nazvanje knigy i nomer glavy / A book title and a chapter number}\quad{\footnotesize\texttt{DEC-0166}}\par
\smallskip\noindent\emph{Razlog:} Polna citacija Lürotha, Castelnuova ili Enriquesa jest jedna čitateljska cělina, daže ako sodržuje nekoliko tuđych slov. Pri Theorie d. algebr. Größen originalno nazvanje završaje se prěd Kap. VI. Rimsky nomer glavy ostaje strukturna oznaka, ne prěkladano slovo. Jedenadset takyh cělin pokazuje, čemu granica citacije potrěbuje analizu sadržanja.\par
\smallskip\noindent\emph{Reason:} A complete citation to Lüroth, Castelnuovo or Enriques is one reading unit even when it contains several foreign words. In Theorie d. algebr. Größen the original title ends before Kap. VI. The chapter numeral remains a structural reference rather than a translated word. These eleven units show why quotation boundaries require interpretation of content.\par
\smallskip\noindent\emph{Izvory:} Polny spisok citatov, odděl 166.\par
\noindent\emph{Sources:} Complete citation inventory, group 166\par
\par\medskip\Needspace{8\baselineskip}
\hypertarget{DEC-0167}{}
\noindent\textbf{Formula v nazvanju statje / A formula within an article title}\quad{\footnotesize\texttt{DEC-0167}}\par
\smallskip\noindent\emph{Razlog:} Hassevo nazvanje Über p-adische Schiefkörper sodržuje matematičny symbol p, izvorno v gotičnom obliku. Citat mora čuvati i tekst i symbol. Drugy priměr, Schurovo nazvanje, završaje se prěd Satz VII: to jest odvolanje na teoremu, ne čest nazvanja. Osm cělyh citacij vključaje takože odvojenje treh Fraenkelovyh rabot; ih ne treba sjediniti samo zato, že imajut jednogo avtora.\par
\smallskip\noindent\emph{Reason:} Hasse's title Über p-adische Schiefkörper contains the mathematical symbol p, originally in fraktur. Both title text and symbol must be preserved. A Schur title ends before Satz VII: that is a theorem reference, not part of the title. The eight complete citations also keep three Fraenkel works distinct; common authorship does not make them one citation.\par
\smallskip\noindent\emph{Izvory:} Polny spisok citatov, odděl 167; matematičny zapis v citatah ostavaje netknjeny.\par
\noindent\emph{Sources:} Complete citation inventory, group 167; literal mathematical payload retained\par
\par\medskip\Needspace{8\baselineskip}
\hypertarget{DEC-0168}{}
\noindent\textbf{Povtorjeno ime ne znači isto upotrěbjenje / A repeated name need not have the same function}\quad{\footnotesize\texttt{DEC-0168}}\par
\smallskip\noindent\emph{Razlog:} Devetnadset izrazov zahranja skračenje D.–W., imena Ostrowski, Weitzenböck i W. Krull i nazvanje Encyklopädie d. math. $W_i$ss. Ako jedno ime nastupaje dvakrat v istoj frazě, kazdo pojavjenje mora byti svjazano so svojim kontekstom. Inače možno pomylno zaščititi sklonjenu prěkladanu formu vměsto citacije. Sosednje medžuslovjanske spojky ostajut čestju prěkladanoj prozy.\par
\smallskip\noindent\emph{Reason:} Nineteen expressions include D.–W., the names Ostrowski, Weitzenböck and W. Krull, and Encyklopädie d. math. $W_i$ss. When a name appears twice in one sentence, each occurrence must be tied to its own context. Otherwise an inflected translated form could be protected instead of the citation. Neighboring Interslavic connectives remain translated prose.\par
\smallskip\noindent\emph{Izvory:} Polny spisok povtorjenyh imen i citatov, odděl 168.\par
\noindent\emph{Sources:} Complete repeated-name and citation inventory, group 168\par
\par\medskip\Needspace{8\baselineskip}
\hypertarget{DEC-0169}{}
\noindent\textbf{D.–W. i neměnjane bibliografične frazy / D.–W. and retained bibliographic expressions}\quad{\footnotesize\texttt{DEC-0169}}\par
\smallskip\noindent\emph{Razlog:} Jedinnadset upotrěbjenij D.–W. sut odvolanja na avtorske inicijaly, ne rimsky nomery ili nepoznate slova. Drugyh deset cělin zahranja nazvanja, Braunschweig, Vieweg i německe Vgl. etwa. Posljednje jest bibliografičny ukazatel ‘sravni, napr.ʼ: jeho čuvanje jest odluka o citacijskom sloju, ne tvrdženje, že jest medžuslovjansky. V lokalizovanoj navigaciji takove frazy potrěbujut razumljivo objašnjenje.\par
\smallskip\noindent\emph{Reason:} Eleven D.–W. occurrences are author abbreviations, not Roman numerals or unknown words. Ten other units include titles, Braunschweig, Vieweg and German Vgl. etwa. The last is a bibliographic cue meaning ‘compare, for example’: retaining it is a choice about the citation layer, not a claim that it is Interslavic. Localized navigation still needs an intelligible explanation.\par
\smallskip\noindent\emph{Izvory:} Polny spisok skračenij, naslovov i citatov, odděl 169.\par
\noindent\emph{Sources:} Complete abbreviation/title/citation inventory, group 169\par
\par\medskip\Needspace{8\baselineskip}
\hypertarget{DEC-0170}{}
\noindent\textbf{Citovane W. Schmeidler i sklonjeno V. Šmajdlera / Cited W. Schmeidler and inflected V. Šmajdlera}\quad{\footnotesize\texttt{DEC-0170}}\par
\smallskip\noindent\emph{Razlog:} Originalno W. Schmeidler v bibliografiji ostaje neproměnjeno. Fraza V. Šmajdlera jest, naprotiv, medžuslovjansko sklonjenje; istorija izdanja prijela V. po uživanoj formě V. Šmajdlerom. Ta odluka o inicijalu mora byti vidima, ne prikryta kako čisto tipografična rabota. Četyrdeset tri citacijske cěliny čuvajut druge originalne forme. Dva nazvanja s matematičnym n razděljajut se okolo symbola; Bd. II der prěd Gesammelten Werke ostaje posebnym strukturnym odvolanjem.\par
\smallskip\noindent\emph{Reason:} Original W. Schmeidler remains unchanged in bibliographic citation. V. Šmajdlera, by contrast, is an inflected Interslavic form; the edition adopted V. following its existing V. Šmajdlerom. That decision about an initial must be visible, not disguised as merely typographic work. Forty-three citation units retain other original forms. Two titles containing mathematical n are divided around the symbol; Bd. II der before Gesammelten Werke remains a separate structural reference.\par
\par\smallskip\noindent\emph{Sravnjenje formov / Comparison of forms}
\begin{longtable}{@{}P{0.43\textwidth}P{0.43\textwidth}P{0.08\textwidth}@{}}
\toprule Sravniteljnaja forma / Compared form & Izbrana forma ili funkcija / Selected form or function & Red / Line\\\midrule\endfirsthead
\toprule Sravniteljnaja forma / Compared form & Izbrana forma ili funkcija / Selected form or function & Red / Line\\\midrule\endhead
\bottomrule\endfoot
Über einige algebraische Körper 8. Grades & zachovany citat / retained citation & 10031\\\addlinespace
Neuer Beweis des Hilbertschen Satzes über homogene Funktionen, Gött. Nachr. 1899 & zachovany citat / retained citation & 10122\\\addlinespace
Über die Existenz einer endlichen Basis bei gewissen Funktionensystemen & zachovany citat / retained citation & 10122\\\addlinespace
A. Loewy & zachovany citat / retained citation & 10953\\\addlinespace
A. Loewy & zachovany citat / retained citation & 10953\\\addlinespace
Loewy & zachovany citat / retained citation & 10953\\\addlinespace
Göttingen & zachovany citat / retained citation & 10953\\\addlinespace
W. Šmajdlera & zachovany citat / retained citation & 10965\\\addlinespace
Hurwitz & zachovany citat / retained citation & 13390\\\addlinespace
Du Pasquier & zachovany citat / retained citation & 13390\\\addlinespace
Weyl & zachovany citat / retained citation & 14151\\\addlinespace
Weyl & zachovany citat / retained citation & 14151\\\addlinespace
Wien. Ber. & zachovany citat / retained citation & 14151\\\addlinespace
Weyl & zachovany citat / retained citation & 15205\\\addlinespace
Weitzenböck & zachovany citat / retained citation & 15205\\\addlinespace
Weyl & zachovany citat / retained citation & 15205\\\addlinespace
Weyl & zachovany citat / retained citation & 15205\\\addlinespace
Weitzenböck & zachovany citat / retained citation & 15205\\\addlinespace
Wiener Akademie & zachovany citat / retained citation & 15205\\\addlinespace
Zur Theorie der aus & zachovany citat / retained citation & 17273\\\addlinespace
Haupteinheiten gebildeten komplexen Größen & zachovany citat / retained citation & 17273\\\addlinespace
Gött. Nachr. & zachovany citat / retained citation & 17273\\\addlinespace
Naturf. Gesellschaft Zürich & zachovany citat / retained citation & 17273\\\addlinespace
Ostrowski & zachovany citat / retained citation & 17277\\\addlinespace
A. Ostrowski: Zur arithmetischen Theorie der algebraischen Größen. Nachr. Gött. Ges. d. Wissensch. 1919. & zachovany citat / retained citation & 17277\\\addlinespace
W. Krull & zachovany citat / retained citation & 17277\\\addlinespace
Zur Theorie der aus & zachovany citat / retained citation & 17279\\\addlinespace
Haupteinheiten gebildeten komplexen Größen. Gött. Nachr. 1885. & zachovany citat / retained citation & 17279\\\addlinespace
Mathematische Probleme (Gött. Nachr. 1900, s. 253--297) & zachovany citat / retained citation & 22523\\\addlinespace
Über die Konstruktion relativ-Abelscher Zahlkörper durch Modulfunktionen von zwei Variabeln & zachovany citat / retained citation & 22523\\\addlinespace
Die Endlichkeit des Systems der ganzzahligen Invarianten binärer Formen (Gött. Nachr. 1919, s. 138--156) & zachovany citat / retained citation & 22533\\\addlinespace
B. L. van der Waerden & zachovany citat / retained citation & 22535\\\addlinespace
A. Ostrowski, Zur arithmetischen Theorie der algebraischen Größen, Gött. Nachr. 1919 & zachovany citat / retained citation & 22535\\\addlinespace
Ein algebraisches Kriterium für absolute Irreduzibilität & zachovany citat / retained citation & 22535\\\addlinespace
Zürich & zachovany citat / retained citation & 23236\\\addlinespace
van der Waerden & zachovany citat / retained citation & 23389\\\addlinespace
van der Waerden & zachovany citat / retained citation & 23389\\\addlinespace
G. Köthe & zachovany citat / retained citation & 23389\\\addlinespace
van der Waerden & zachovany citat / retained citation & 23389\\\addlinespace
G. Köthe & zachovany citat / retained citation & 23391\\\addlinespace
van der Waerden & zachovany citat / retained citation & 23391\\\addlinespace
Gesammelten Werke & zachovany citat / retained citation & 24103\\\addlinespace
Verhandlungen des Internationalen Mathematiker-Kongresses Zürich 1932, Bd. I. & zachovany citat / retained citation & 24190\\\addlinespace
Göttingen & zachovany citat / retained citation & 24694\\\addlinespace
\end{longtable}
\smallskip\noindent\emph{Izvory:} Polny spisok citatov, odděl 170; izbor inicijaly v prilagođenom imenu.\par
\noindent\emph{Sources:} Complete citation inventory, group 170; integrated-name initial decision\par
\par\medskip\Needspace{8\baselineskip}
\hypertarget{DEC-0171}{}
\noindent\textbf{Citacija v statji i v zbirnoj knigě / A citation in an article and in the collection}\quad{\footnotesize\texttt{DEC-0171}}\par
\smallskip\noindent\emph{Razlog:} Četrnadset citacij v Paper 2, 8, 22 i 24 mora imati jednaku originalnu formu v samostojnoj statji i v zbirnoj knigě. Zahranjajut Fischerovo nazvanje o modulsystemah, König, Dixmuiden, Ostrowski i nazvanje o absolutnoj irreducibilnosti. Nula nepoznatyh znakova v toj istorijnoj proverce dokazuje toliko zaključenu klasifikaciju tih izrazov. Ne dokazuje polnu terminologičnu pravilnost, razumljivost ili věrnost každoj formule. Te poslednje ocěny potrěbujut svoji dokazy.\par
\smallskip\noindent\emph{Reason:} Fourteen citations in Papers 2, 8, 22 and 24 must have the same original forms in the separate articles and the collection. They include Fischer's title on module systems, König, Dixmuiden, Ostrowski and a title concerning absolute irreducibility. Zero unknown characters in that historical check establishes only completion of that expression classification. It does not establish complete terminological correctness, intelligibility or fidelity of every formula. Those judgments require their own evidence.\par
\par\smallskip\noindent\emph{Sravnjenje formov / Comparison of forms}
\begin{longtable}{@{}P{0.43\textwidth}P{0.43\textwidth}P{0.08\textwidth}@{}}
\toprule Sravniteljnaja forma / Compared form & Izbrana forma ili funkcija / Selected form or function & Red / Line\\\midrule\endfirsthead
\toprule Sravniteljnaja forma / Compared form & Izbrana forma ili funkcija / Selected form or function & Red / Line\\\midrule\endhead
\bottomrule\endfoot
Abh. der Gött. Ges. d. Wiss. & zachovany citat / retained citation & \\\addlinespace
Sitzb. d. Ak. d. Wiss. Wien & zachovany citat / retained citation & \\\addlinespace
Mém. d. 1. Soc. Roy. des Sciences de Liège & zachovany citat / retained citation & \\\addlinespace
Über algebraische Modulsysteme und lineare homogene partielle Differentialgleichungen mit konstanten Koeffizienten & zachovany citat / retained citation & \\\addlinespace
Dixmuiden & zachovany citat / retained citation & \\\addlinespace
König & zachovany citat / retained citation & \\\addlinespace
Rechteckige Systeme und Moduln in algebraischen Zahlkörpern & zachovany citat / retained citation & \\\addlinespace
Algebraische Theorie der Körper & zachovany citat / retained citation & \\\addlinespace
A. Ostrowski & zachovany citat / retained citation & \\\addlinespace
A. Ostrowski, Zur arithmetischen Theorie der algebraischen Größen, Gött. Nachr. 1919, s. 279--298 & zachovany citat / retained citation & \\\addlinespace
A. Ostrowski & zachovany citat / retained citation & \\\addlinespace
Ein algebraisches Kriterium für absolute Irreduzibilität & zachovany citat / retained citation & \\\addlinespace
Göttingen & zachovany citat / retained citation & \\\addlinespace
\end{longtable}
\smallskip\noindent\emph{Izvory:} Polny spisok citatov iz statij i sbornika, odděl 171.\par
\noindent\emph{Sources:} Complete article/collection citation inventory, group 171\par
\normalsize

\clearpage\section{Dvě pisma, jeden prěvod / Two scripts, one translation}
Latinica i kirilica ne sut dva nezavisne prěvody. Kirilična verzija prěpisuje medžuslovjansky tekst, ale zachovuje formule, bibliografične identičnosti i izvorne citaty. Razlika zavisi od funkcije v konkretnom městě: \emph{On} može byti medžuslovjansky zaimennik, a \emph{II} rimsko čislo. Ime v sklonjenoj frazě može iměti drugu funkciju od togo samogo imena v bibliografiji.
\par\emph{EN:} Latin and Cyrillic are not independent translations. The Cyrillic version transliterates Interslavic text while retaining formulas, bibliographic identities, and source-language quotations. Function is contextual: \emph{On} may be an Interslavic pronoun, whereas \emph{II} is a Roman numeral. A name integrated into an inflected phrase may have a different function from the same name in a bibliography.
\par Odluka ne ustanavjaje opče pravilo za vsako pojavjenje istogo napisa. Za vse niže navedene slučaje tablice podatkov zachovujut datoteku, red i stupac. / These choices do not establish a global rule for every occurrence of the same spelling. The data tables retain file, line, and column for every case listed below.
\small\begin{longtable}{@{}P{0.28\textwidth}P{0.54\textwidth}P{0.10\textwidth}@{}}
\toprule Forma / Form & Funkcija / Function & Pojav. / Uses\\\midrule\endhead
\bottomrule\endfoot
Algebra & zachovany izvorni zapis / retained source spelling & 3\\
Grassmann & jezyčny tekst za prěpis / transliterated language text & 1\\
II & zachovany izvorni zapis / retained source spelling & 215\\
IIa & zachovany izvorni zapis / retained source spelling & 3\\
Ideale & zachovany izvorni zapis / retained source spelling & 2\\
Jordan & jezyčny tekst za prěpis / transliterated language text & 1\\
No & zachovany izvorni zapis / retained source spelling & 2\\
On & jezyčny tekst za prěpis / transliterated language text & 16\\
XIII & zachovany izvorni zapis / retained source spelling & 10\\
XXII & zachovany izvorni zapis / retained source spelling & 1\\
der & jezyčny tekst za prěpis / transliterated language text & 3\\
in & zachovany izvorni zapis / retained source spelling & 7\\
jan & jezyčny tekst za prěpis / transliterated language text & 1\\
teoria & zachovany izvorni zapis / retained source spelling & 2\\
van & jezyčny tekst za prěpis / transliterated language text & 3\\
I & rimsko čislo / Roman numeral & 73\\
I & jezyčny tekst za prěpis / transliterated language text & 1\\
V & rimsko čislo / Roman numeral & 19\\
in inf. & polny izvorni citat / complete source quotation & 1\\
Hensel & bibliografična identičnost / bibliographic identity & 1\\
Clebsch & bibliografična identičnost / bibliographic identity & 1\\
Macaulay & bibliografična identičnost / bibliographic identity & 1\\
Masazo Sono & bibliografična identičnost / bibliographic identity & 1\\
Sylvester & bibliografična identičnost / bibliographic identity & 1\\
Leopold Kronecker & bibliografična identičnost / bibliographic identity & 1\\
Leopold Kronecker & bibliografična identičnost / bibliographic identity & 1\\
Leopold Kronecker & bibliografična identičnost / bibliographic identity & 1\\
\end{longtable}\normalsize
\noindent Imena \emph{Grassmann} i \emph{Jordan} v tabeli sut česti prilagođenyh fraz Grassmann--Mullerovogo i Jordan-Helderov. Podobno, van der v frazah van der Vardena/Vardenu jest vključeno v jezyčny prěpis. To ne ustanavjaje novy način zapisa imena v bibliografiji. Dvě citatne skračenja Jan i pet nesklonjenyh pojavjenj Noether zachovujut izvorni zapis.
\par\emph{EN:} The Grassmann and Jordan cases occur inside the integrated phrases Grassmann--Mullerovogo and Jordan-Helderov. Likewise, van der is included in transliteration in van der Vardena/Vardenu. This does not prescribe a new bibliographic spelling. Two citation-month abbreviations Jan and five undeclined Noether identities retain their source spelling.

\begin{thebibliography}{9}

\bibitem{designcriteria}
Jan van Steenbergen. ``Design criteria.'' Interslavic language project,
section ``Which roots to use?''
\href{https://interslavic.fun/learn/misc/design-criteria/}{Project documentation}.

\bibitem{merunka2019}
Vojtěch Merunka, Jan van Steenbergen, Lina Yordanova, and Maria Kocór.
``The Interslavic language as a tool for supporting e-democracy in Central and Eastern Europe.''
\emph{International Journal of Electronic Governance} 11, nos. 3--4 (2019): 260--288.
\href{https://doi.org/10.1504/IJEG.2019.103710}{doi:10.1504/IJEG.2019.103710}.

\bibitem{iso704}
International Organization for Standardization.
\emph{ISO 704:2022, Terminology work---Principles and methods}. 4th ed. Geneva: ISO, 2022.
\href{https://www.iso.org/standard/79077.html}{Official record}.

\bibitem{nazar2023}
Rogelio Nazar and Nicolás Acosta.
``Termout: a tool for the semi-automatic creation of term databases.''
In \emph{Proceedings of the Workshop on Computational Terminology in NLP and Translation Studies}, 9--18. 2023.
\href{https://aclanthology.org/2023.contents-1.2/}{ACL Anthology}.

\bibitem{court2024}
Sara Court and Micha Elsner.
``Shortcomings of LLMs for Low-Resource Translation: Retrieval and Understanding Are Both the Problem.''
In \emph{Proceedings of the Ninth Conference on Machine Translation}, 1332--1354. 2024.
\href{https://doi.org/10.18653/v1/2024.wmt-1.125}{doi:10.18653/v1/2024.wmt-1.125}.

\bibitem{bane2023}
Fred Bane, Anna Zaretskaya, Tània Blanch Miró, Celia Soler Uguet, and João Torres.
``Coming to Terms with Glossary Enforcement: A Study of Three Approaches to Enforcing Terminology in NMT.''
In \emph{Proceedings of the 24th Annual Conference of the European Association for Machine Translation}, 345--353. 2023.
\href{https://aclanthology.org/2023.eamt-1.34/}{ACL Anthology}.

\bibitem{freitag2021}
Markus Freitag, George Foster, David Grangier, Viresh Ratnakar, Qijun Tan, and Wolfgang Macherey.
``Experts, Errors, and Context: A Large-Scale Study of Human Evaluation for Machine Translation.''
\emph{Transactions of the Association for Computational Linguistics} 9 (2021): 1460--1474.
\href{https://doi.org/10.1162/tacl_a_00437}{doi:10.1162/tacl\_a\_00437}.

\bibitem{agrawal2024}
Sweta Agrawal, António Farinhas, Ricardo Rei, and Andre Martins.
``Can Automatic Metrics Assess High-Quality Translations?''
In \emph{Proceedings of the 2024 Conference on Empirical Methods in Natural Language Processing}, 14491--14502. 2024.
\href{https://doi.org/10.18653/v1/2024.emnlp-main.802}{doi:10.18653/v1/2024.emnlp-main.802}.

\end{thebibliography}

\end{document}



